% \def\fileversion{3.27} \def\filedate{2019/11/05}
% \MakeShortVerb{\|}
% \title{A \LaTeX{} document class for the Revista Mexicana de
% Astronom\'\i a y Astrof\'\i sica\thanks{This file has version
% \fileversion{} last revised \filedate{}.}}
% \author{William J. Henney\thanks{w.henney@irya.unam.mx}}
% \date{\filedate}
% \maketitle
% \paragraph{Abstract}
% This document describes the \texttt{rmaa} \LaTeX{}
% document class for the journal ``Revista Mexicana de
% Astronom\'\i a y Astrof\'\i sica''. The macros can be used for 
% \begin{enumerate}
% \item Preparing author manuscripts for the journal (including
% conference proceedings to be published in the ``Serie de
% Conferencias'')
% \item Typesetting an entire volume of the journal by inputting the
% individual article files from a container file. Facilities are
% included for automatic generation of TOC and Author Index. 
% \end{enumerate}
% 
% \tableofcontents
% \section{The documentation driver file}
% This is the part that generates the documentation that you may or
% may not be reading, which is done with the help of the Makefile
% discussed in section~\ref{sec:driver_make}. 
%    \begin{macrocode}
%<*driver>
\documentclass{ltxdoc}
\usepackage{hyperref}
\usepackage{palatino}
\EnableCrossrefs \CodelineIndex \RecordChanges
\setlength{\parindent}{0pt}
\def\MacroFont{\footnotesize\tt}
\begin{document} \DocInput{rmaa.dtx} \PrintIndex \PrintChanges \end{document}
%</driver>
%    \end{macrocode}

% \section{The implementation}
% 
% \changes{3.24}{2007/09/14}{Added user access to the macro version
% and date via \cs{rmVersion} and \cs{rmDate}}
% 
%    \begin{macrocode}
%<*package>
\newcommand{\rmaa@version}{3.26}
\newcommand{\rmaa@date}{2011/08/15}
\newcommand\rmVersion{\rmaa@version}
\newcommand\rmDate{\rmaa@date}
\NeedsTeXFormat{LaTeX2e}[1995/06/01]
%    \end{macrocode}
% \changes{3.27}{2019/11/02}{Change abstract to appear before
% resumen. Change RevMexAA to RMxAA for consistency with ADS.}
%
% \changes{3.22}{2007/03/22}{Allow proceedings articles to have
%   received/accepted dates} \changes{3.21}{2006/05/15}{Changes for
%   Silvia - abstracts are now shipped out to ADS.}
% \changes{3.19}{2006/01/19}{Changes for Silvia - remove adresses from
%   end, new abstract style} \changes{2.2beta}{2001/10/30}{ Fixed bug
%   whereby twocolumn option was not passed through to book
%   class. Finally sorted out the ``label on last page'' issue. }
% \changes{2.1}{2001/01/23}{ Fixed some inconsistencies in
%   \cs{nocite} and \cs{bibitem} uniqification Fixed the definition of
%   \cs{ion} so it will work in math mode and in contexts where the
%   typesize is different from normal (e.g., title, fig captions).
%   Modified the definition of \cs{@caption}, so we don't write any
%   list-of-figures data (we never use it anyway, and it just causes
%   problems with fragile commands in fig captions). }
% \changes{2.0}{2000/11/30}{ Added `newstyle' option, which should be
%   used in all new proceedings volumes (wider abstract, smaller
%   typeface in fig captions, etc.) Added commands for formatting
%   table captions. New demonstation document rmpostsamp.tex for
%   1-page posters. This will be the last release under a PRI
%   government (for a while...)}  \changes{1.9}{2000/06/06}{
%   Corrected position of page number in footer of articles that open
%   on even pages. Ensured that first article in each session opens on
%   an odd page, as do list of abstracts and author index. Added extra
%   space in TOC before the backmatter. Added command for editors,
%   \cs{BoldSessionTitles}, that makes produces bold titles for
%   session in the TOC.}  \changes{1.8}{2000/05/16}{ Fixed bug in
%   last page label with \cs{suppressfulladdresses} }
% \changes{1.7}{2000/05/11}{ Fixed more little things. Page labels
%   now work, even if you have no bibliography. You can output the
%   full addresses by hand with the \cs{outputfulladdresses}
%   command. Editors' Guide updated but still not finished. Revamped
%   splitbook.sh now produces PS file of each article for web page and
%   also an EPS file of each page for sending to the printers.}
% \changes{1.6}{2000/04/15}{ More cleaning up: full stops after
%   titles in index, running headers and footers in frontmatter,
%   fontsize in title of index, etc} \changes{1.5}{2000/04/15}{
%   Eliminated superscripts from titles in TOC and index}
% \changes{1.4}{2000/04/14}{ Added copyright and license info (GPL)}
% \changes{1.3}{2000/03/01}{ Corrected \cs{figcaption} to do the
%   caption formatting right Corrected font size of sec headings in
%   front matter} \changes{1.2}{2000/02/16}{ Added \cs{plate}
%   command} \changes{1.1}{2000/02/14}{ Added manuscript option,
%   corrected bugs} \changes{1.0}{2000/02/08}{ Now does a whole book
%   - lots of stuff tidied up} \changes{0.7}{1999/10/04}{ Increased
%   spacing between sections in single column styles. Suppressed
%   floats at bottom of pages (didn't work though) }
% \changes{0.6}{1999/09/24}{ Tweaked \cs{adjustfinalcols} (still not
%   perfect) Added proceedings option Added onecolumn/twocolumn
%   options (orthogonal to journal/proceedings) }
% \changes{0.5}{1999/09/23}{ Added \cs{fulladdresses} and
%   \cs{adjustfinalcols} (just about working)}
% \changes{0.4}{1999/09/22}{ Minor stylistic changes}
% \changes{0.3}{1999/06/28}{ Better mechanism for alternative
%   affiliations Added `guide' option for user guide}
% \changes{0.2}{1999/05/27}{ All basic functionality present, I
%   think} \changes{0.1}{1999/05/26}{ Just about working - still
%   missing Fig/Tab/Refs} \changes{0.0}{1999/05/21}{ So bored today,
%   decided to rewrite RMAA LaTeX macros from scratch}
%    \begin{macrocode}
\ProvidesClass{rmaa}[\rmaa@date\space v\rmaa@version\space RMxAA article
class]
%    \end{macrocode}
% \subsection{Option Handling}
% \label{sec:options}
% Declarations that are needed in options
% \changes{3.05}{2003/10/22}{Added natbib option}
% \changes{3.06}{2003/01/03}{Natbib option turned on permanently}
% \changes{3.07}{2003/01/11}{Added proof option}
% \changes{3.14}{2003/01/11}{Option to turn off natbib}
% \changes{3.18}{2003/10/22}{trying out doctex mode} 
% \changes{3.23}{2007/09/09}{Manuscript is now on by default.}
% \changes{3.26}{2011/08/15}{New option for review articles}
%    \begin{macrocode}
\newif\if@reviewarticle
\@reviewarticlefalse
\newif\if@proof
\@prooffalse
\newif\if@journal
\@journaltrue
\newif\if@printdates
\@printdatesfalse
\newif\if@proceedings
\@proceedingsfalse
\newif\if@manuscript
\@manuscripttrue
\newif\if@bigfont
\@bigfonttrue
\newif\if@guide
\@guidefalse
\newif\if@preprint
\@preprintfalse
\newif\if@rmtwocol
\@rmtwocoltrue
\newif\if@biblist
\@biblistfalse
\newif\if@book
\@bookfalse
\newif\if@natbib
\@natbibtrue
\newif\if@rmcite
\@rmcitefalse
\newif\if@titlevisible
\@titlevisibletrue
\newif\if@abstractpresent
\@abstractpresentfalse
\newif\if@printbanner
\@printbannerfalse
\newif\if@suppressfulladdresses
\@suppressfulladdressesfalse
\newif\if@newstyle
\@newstylefalse
\newif\if@debug
\@debugfalse
\newif\if@boldkeywords
\@boldkeywordsfalse
%    \end{macrocode}
% The option handling---mainly setting various booleans.
% \changes{2.2beta11}{2001/11/26}{Made twocol and
% newstyle the default for proceedings.}
% \changes{3.23}{2007/09/09}{Manuscript is now 10 pt instead of 12pt}
% \changes{3.23}{2007/09/10}{Full addresses are now automatically
% suppressed for proceedings style.}
% \changes{3.24}{2007/09/14}{New option \texttt{boldkeywords} that
% does the keywords in the old ugly way. Now off by default.}
%    \begin{macrocode}
\DeclareOption{manuscript}{\typeout{Manuscript for referee} 
  \@manuscripttrue\@journalfalse\@rmtwocolfalse
  \PassOptionsToClass{10pt}{article}}
\DeclareOption{bigfont}{\typeout{Big font manuscript requested by Gary
  Ferland} 
  \@manuscripttrue\@journalfalse\@rmtwocolfalse\@bigfonttrue
  \PassOptionsToClass{11pt}{article}}
\DeclareOption{journal}{\typeout{Main journal option}%
  \@journaltrue\@rmtwocoltrue\@printdatestrue\@manuscriptfalse
  \PassOptionsToClass{10pt}{article}}
\DeclareOption{guide}{\typeout{User guide option}%
  \@guidetrue\@journalfalse\@rmtwocoltrue\@manuscriptfalse
  \PassOptionsToClass{10pt}{article}}
\DeclareOption{preprint}{\typeout{Preprint option}%
  \@preprinttrue\@manuscriptfalse\PassOptionsToClass{10pt}{article}}
\DeclareOption{proof}{\typeout{Author proof option}%
  \@prooftrue\@manuscriptfalse}
\DeclareOption{proceedings}{\typeout{Proceedings option}%
  \@proceedingstrue\@journalfalse\@rmtwocoltrue\@newstyletrue
  \@suppressfulladdressestrue
  \@printdatesfalse\@manuscriptfalse\PassOptionsToClass{10pt}{article}} 
\DeclareOption{printdates}{\typeout{Force printing of received/accepted dates}%
  \@printdatestrue} 
\DeclareOption{book}{\typeout{Whole issue option}%
  \@booktrue} 
\DeclareOption{nonatbib}{\typeout{Turn off natbib support (required
    for compatibility with old rmaacite package).}%
  \@natbibfalse} 
\DeclareOption{newstyle}{\typeout{Use new layout style}%
  \@newstyletrue}
\DeclareOption{oldstyle}{\typeout{Use old layout style}%
  \@newstylefalse}
\DeclareOption{debug}{\typeout{Loads of annoying messages}%
  \@debugtrue}
%    \end{macrocode}
% These two need to be declared explicitly so that |\ExecuteOptions|
% actually does something.
%    \begin{macrocode}
\DeclareOption{twocolumn}{\@rmtwocoltrue}
\DeclareOption{onecolumn}{\@rmtwocolfalse}
\DeclareOption{twoside}{\PassOptionsToClass{\CurrentOption}{article}}
\DeclareOption{boldkeywords}{\@boldkeywordstrue}
%    \end{macrocode}
% Bold, upper-case keywords used to be the default. Now they must be
% turned on explicitly with the \texttt{boldkeywords} option. 
%    \begin{macrocode}
\DeclareOption{boldkeywords}{\@boldkeywordstrue}
%    \end{macrocode}
% Pass any other option along to the \texttt{article} class.
%    \begin{macrocode}
\DeclareOption*{\PassOptionsToClass{\CurrentOption}{article}}
%    \end{macrocode}
% Default options are to be two-sided and based on the standard
% \LaTeX{} \texttt{article} class,
% unless we are processing the entire book, in which case we base
% ourselves on the standard \texttt{book} class. Default is to use
% manuscript option unless otherwise specified. 
%    \begin{macrocode}
\ExecuteOptions{twoside,manuscript} \ProcessOptions\relax
\if@book
\if@rmtwocol\LoadClass[twocolumn]{book}\else\LoadClass{book}\fi
\newcommand{\abstract}{\relax}
\newcommand{\@maketitle}{\relax}
\else
\if@rmtwocol\LoadClass[twocolumn]{article}\else\LoadClass{article}\fi
\fi
%    \end{macrocode}
% Attempt to use natbib for citations 02 Jan 2003. 
% \changes{3.05}{2003/01/02}{If the natbib option is used, then set
% things up so that natbib's version of the bibliography is identical
% to the standard one.}
% \changes{3.09}{2003/01/23}{Explicitly require a recent version of
% natbib, so that we get a decent error message if the package is too
% old.} 
% \changes{3.23}{2007/09/09}{Now allow Bibtex to work properly} 
% \changes{3.25}{2009/10/01}{Change \cs{newcommand} tp \cs{def} when setting \cs{bibfont}. This allows it to work with all versions of natbib.} 
% \changes{3.25}{2010/02/06}{Updated copyright year.} 
% \changes{3.25}{2010/02/06}{Updated email and web addresses to \texttt{crya.unam.mx}.} 
%    \begin{macrocode}
\if@natbib
\RequirePackage{natbib}[1999/02/25]
\bibpunct{(}{)}{;}{a}{}{,}
\def\bibfont{\small}
\renewcommand\bibsection{\section*{REFERENCES}}
\setlength\bibsep{0pt}
\setlength\bibhang{1.5em}
\bibliographystyle{rmaa}
\fi
%    \end{macrocode}
% We want the \texttt{graphicx} package for the extended syntax of
% |\includegraphics|. We want the \texttt{ifthen} package for using in
% these macros (although we don't much yet). We explicitly load
% \texttt{babel} since some \LaTeX{} distributions messed up the
% hyphenation if we didn't.  
%    \begin{macrocode}
\RequirePackage{graphicx}
\RequirePackage{ifthen}
\IfFileExists{babel.sty}{%
%   \let\@restoreindent=\@afterindentfalse
  \RequirePackage[english]{babel}
%   \@restoreindent
  %% reinstate the definition of \@roman so as not to break .idx files
  \@ifundefined{org@@roman}{\relax}{\renewcommand{\@roman}[1]{\org@@roman{##1}}}
  }{%
  \typeout{Package babel.sty not found - don't worry about this}}
%    \end{macrocode}
% Try to load the \texttt{dcolumn} and \texttt{booktabs} packages if
% they can be found. These can produce nicer-looking tables. 
%    \begin{macrocode}
\IfFileExists{dcolumn.sty}{%
  \RequirePackage{dcolumn}
  }{%
  \typeout{Package dcolumn.sty not found - you will not be able to
    align columns on decimal points in tables}}
\IfFileExists{booktabs.sty}{%
  \RequirePackage{booktabs}
  }{%
  \typeout{Package booktabs.sty not found - your tables will not look
    as pretty as they ought too.}
  }
\IfFileExists{makeidx.sty}{%
  \RequirePackage{makeidx}
}{%
  \typeout{Package makeidx.sty not found - Author Index may not work.}
  }
%% commented out for moment since requires more modern version of LaTeX
%%\RequirePackage[latin1]{inputenc}
\IfFileExists{hyperref.sty}{%
  %\RequirePackage[draft,implicit=false]{hyperref}
  %\hypersetup{hidelinks=true}
  \RequirePackage{url}
  \urlstyle{same}
  \newcommand\email{\begingroup \urlstyle{rm}\Url}
  }{%
  \typeout{Package hyperref.sty not found - line breaks in email addresses
    will not work properly.}
  \newcommand\email[1]{#1}
  }
%    \end{macrocode}
% \subsection{Other setup stuff}\label{sec:setup}
% Some declarations: 
%    \begin{macrocode}
\newif\if@needaddresses
\@needaddressestrue
\newif\if@frontmatter
\@frontmatterfalse
\newif\if@inpreamble
\@inpreambletrue
\newif\if@firstauthor
\@firstauthortrue
\newif\if@firstaffil
\@firstaffiltrue
\newif\if@firstaltaffil
\@firstaltaffiltrue
\newif\if@firstkwd
\@firstkwdtrue
\newif\if@noexpandADS
\@noexpandADSfalse
\newif\if@sessiontitlepages
\@sessiontitlepagesfalse
%    \end{macrocode}
% \begin{macro}{\@rmaaEmpty}
% \begin{macro}{\@rmaaArg}
% These are used in various commands to check if arg was left blank
%    \begin{macrocode}
\newcommand{\@rmaaEmpty}{}
\newcommand{\@rmaaArg}{}
%    \end{macrocode}  
% \end{macro}
% \end{macro}
%
% \subsection{Compatibility with old \LaTeX{} versions}
% \label{sec:compat}
% This package is based on the 1998/12/01 release of \LaTeX{}.
% Older releases may not define all the commands that we need.
% Luckily, we can use |\providecommand| to define these if necessary. 
% The following commands are not defined in the 1994/06/01 release. 
%
% \begin{macro}{\MakeUpperCase}
% \begin{macro}{\textsuperscript}
% \begin{macro}{\textbackslash}
% User commands:
%    \begin{macrocode}
\providecommand{\MakeUpperCase}[1]{\uppercase{#1}}
\providecommand{\textsuperscript}[1]{\ensuremath{^\textrm{#1}}}
\providecommand{\textbackslash}{\ensuremath{\backslash}}
%    \end{macrocode}  
% \end{macro}
% \end{macro}
% \end{macro}
% \begin{macro}{\protected@xdef}
% \begin{macro}{\restore@protect}
% \begin{macro}{\@unexpandable@protect}
% \begin{macro}{\hb@xt@}
% \LaTeX{} internal commands:
%    \begin{macrocode}
\providecommand{\protected@xdef}{%
        \let\@@protect\protect
        \let\protect\@unexpandable@protect
        \afterassignment\restore@protect
        \xdef}
\providecommand{\restore@protect}{\let\protect\@@protect}
\providecommand{\@unexpandable@protect}{\noexpand\protect\noexpand}
\providecommand{\hb@xt@}{\hbox to}
%    \end{macrocode}
% \end{macro}
% \end{macro}
% \end{macro}
% \end{macro}
%
% \subsection{Page layout}\label{sec:page-layout}
%
% The dimensions of the text were defined by Elizabeth. Text width is
% 169\,mm with equal margins on the two sides: 
%    \begin{macrocode}
\setlength{\textwidth}{169mm}
\setlength{\@tempdima}{\paperwidth}
\addtolength{\@tempdima}{-\textwidth}
\setlength{\oddsidemargin}{0.5\@tempdima}
\addtolength{\oddsidemargin}{-1in}
\setlength{\evensidemargin}{\oddsidemargin}
%    \end{macrocode}
% Text height is 225\,mm, top margin is 14\,mm. 
%    \begin{macrocode}
\setlength{\textheight}{225mm} 
\setlength{\topmargin}{-1in} 
\addtolength{\topmargin}{14mm}
%    \end{macrocode}
% Space between columns in two-column styles, header, footer,
% paragraph indentation:
%    \begin{macrocode}
\setlength{\columnsep}{9.5mm}
\setlength{\headsep}{0.7cm}
\setlength{\footskip}{0.75cm}
\setlength{\parindent}{1.5em}
%    \end{macrocode}
% 
% \begin{macro}{\@set@scalable@length}
% Try out a new scheme for dealing with various spacing dimensions and
% allowing the user to easily modify them. The command
% \begin{quote}
% |\set@scalable@length|\marg{name}\marg{defaultlength}\marg{stretch}\marg{shrink}
% \end{quote}
% does a 
% |\setlength|\marg{\cs{name}}\marg{defaultlength plus stretch minus shrink}
% and also defines the user command
% \textbackslash\meta{name}|Scale|\marg{factor} that can 
% be used to rescale length by \meta{factor}. Additionally, now
% defines the package command
% \textbackslash\meta{name}|@set@scale|\marg{FACTOR} 
% that also rescales the length by \meta{FACTOR}, so that the length
% becomes \meta{factor} $\times$ \meta{FACTOR} $\times$
% \meta{defaultlength}.   
%    \begin{macrocode}
\newcommand\@set@scalable@length[4]{%
  %% arg #1: NAME of length without leading `\'
  %%
  %% arg #2: default VALUE  
  \expandafter\newlength\csname @#1@default\endcsname
  \expandafter\setlength\csname @#1@default\endcsname{#2}%
  \expandafter\newlength\csname @@#1@@default\endcsname
  \expandafter\setlength\csname @@#1@@default\endcsname{#2}%
  %% arg #3: default STRETCH  
  \expandafter\newlength\csname @#1@stretch\endcsname
  \expandafter\setlength\csname @#1@stretch\endcsname{#3}%
  \expandafter\newlength\csname @@#1@@stretch\endcsname
  \expandafter\setlength\csname @@#1@@stretch\endcsname{#3}%
  %% arg #4: default SHRINK  
  \expandafter\newlength\csname @#1@shrink\endcsname
  \expandafter\setlength\csname @#1@shrink\endcsname{#4}%
  \expandafter\newlength\csname @@#1@@shrink\endcsname
  \expandafter\setlength\csname @@#1@@shrink\endcsname{#4}%
  %%
  %% User-defined scale \@NAME@scale is unity by default
  \expandafter\gdef\csname @#1@scale\endcsname{1}%
  %% Package-defined scale \@NAME@SCALE is unity by default
  \expandafter\gdef\csname @#1@SCALE\endcsname{1}%
  %%
  %% internal command \@reset@NAME should expand to 
  %% \setlength\@@NAME@@default{\@NAME@SCALE\@NAME@default}
  %% \setlength\NAME{\@NAME@scale\@@NAME@@default}
  %% in other words: scale x SCALE x default 
  \expandafter\gdef\csname @reset@#1\endcsname{%
    \expandafter\setlength\csname @@#1@@default\endcsname{%
      \csname @#1@SCALE\endcsname \csname @#1@default\endcsname
      }
    \expandafter\setlength\csname @@#1@@stretch\endcsname{%
      \csname @#1@SCALE\endcsname \csname @#1@stretch\endcsname
      }
    \expandafter\setlength\csname @@#1@@shrink\endcsname{%
      \csname @#1@SCALE\endcsname \csname @#1@shrink\endcsname
      }
    \expandafter\setlength\csname #1\endcsname{%
      \csname @#1@scale\endcsname \csname @@#1@@default\endcsname
      \@plus
      \csname @#1@scale\endcsname \csname @@#1@@stretch\endcsname
      \@minus
      \csname @#1@scale\endcsname \csname @@#1@@shrink\endcsname
      }}
  %%
  %% User command \NAMEScale should expand to 
  %% {\renewcommand\@NAME@scale{#1}\@reset@NAME}
  \expandafter\gdef\csname #1Scale\endcsname##1{%
    \expandafter\gdef\csname @#1@scale\endcsname{##1}%
    \csname @reset@#1\endcsname}
  %%
  %% Package command \NAME@set@SCALE should expand to 
  %% {\renewcommand\@NAME@SCALE{#1}\@reset@NAME}
  \expandafter\gdef\csname #1@set@SCALE\endcsname##1{%
    \expandafter\gdef\csname @#1@SCALE\endcsname{##1}%
    \csname @reset@#1\endcsname}
  %%
  %% initialize at default values with \@reset@NAME
  \csname @reset@#1\endcsname
  }
%    \end{macrocode}
% \end{macro}
% \begin{macro}{\parskip}
% First use of this new infrastructure: set interparagraph space to be
% zero but with a small amount of stretchiness.
%    \begin{macrocode}
%   \ifthenelse{\boolean{@manuscript}}{%
%     \@set@scalable@length{parskip}{0.5\baselineskip}{0.3\p@}{0.1\p@}
%   }{
    \@set@scalable@length{parskip}{0\p@}{0.3\p@}{0\p@}
%   }
%    \end{macrocode}
% \end{macro}
% 
% \begin{macro}{\abovesectionskip}
% \begin{macro}{\belowsectionskip}
% \begin{macro}{\aboveparaskip}
% \begin{macro}{\belowparaskip}
% Define some utility rubber lengths for use in
% section~\ref{sec:section-commands}. 
%    \begin{macrocode}
\newlength{\abovesectionskip}
\newlength{\belowsectionskip}
\newlength{\aboveparaskip}
\newlength{\belowparaskip}
%    \end{macrocode}
% The default values for these lengths are those traditionally used
% for the onecolumn style. 
%    \begin{macrocode}
\@set@scalable@length{abovesectionskip}{3.5ex}{4ex}{2ex}
\@set@scalable@length{belowsectionskip}{2.3ex}{0ex}{2ex}
\@set@scalable@length{aboveparaskip}{3.25ex}{2ex}{0.5ex}
\@set@scalable@length{belowparaskip}{-1em}{\z@}{\z@}
%    \end{macrocode}
% \end{macro}\end{macro}\end{macro}\end{macro}
% We also make scalable some of the standard float parameters. These
% are the same defaults as in the standard article classes, but
% slightly more stretchy. For the twocolumn style, they are currently
% halved, see below. Note that I have also lumped |\abovecaptionskip|
% in with these, since it plays a similar role. There is also a
% |\belowcaptionskip| but it is set to zero in the standard classes. 
%    \begin{macrocode}
\@set@scalable@length{floatsep}{12\p@}{4\p@}{4\p@}
\@set@scalable@length{textfloatsep}{20\p@}{4\p@}{8\p@}
\@set@scalable@length{intextsep}{12\p@}{4\p@}{4\p@}
\@set@scalable@length{dblfloatsep}{12\p@}{4\p@}{4\p@}
\@set@scalable@length{dbltextfloatsep}{20\p@}{4\p@}{8\p@}
\@set@scalable@length{abovecaptionskip}{10\p@}{\z@}{\z@}
%    \end{macrocode}
% \begin{macro}{\@before@title@skip}         
% \begin{macro}{\@before@author@skip}        
% \begin{macro}{\@between@author@skip}       
% \begin{macro}{\@before@abstract@skip}      
% \begin{macro}{\@after@abstracthead@skip}   
% \begin{macro}{\@between@abstract@skip}     
% \begin{macro}{\@before@kwd@skip}           
% \begin{macro}{\@before@dates@skip}           
% Spacing parameters used in |\@maketitle|. These are not changed by
% default except by the user, so they don't need to go into
% |\@package@default@lengths|.  
%    \begin{macrocode}
\newlength{\@before@title@skip}         
\newlength{\@before@author@skip}        
\newlength{\@between@author@skip}       
\newlength{\@before@abstract@skip}      
\newlength{\@after@abstracthead@skip}   
\newlength{\@between@abstract@skip}     
\newlength{\@before@kwd@skip}           
\newlength{\@before@dates@skip}           
\@set@scalable@length{@before@title@skip}{2em}{\z@}{\z@}
\@set@scalable@length{@before@author@skip}{1.5em}{\z@}{\z@}
\@set@scalable@length{@between@author@skip}{0.5em}{\z@}{\z@}
\@set@scalable@length{@before@abstract@skip}{3ex}{1ex}{0.5ex}
\@set@scalable@length{@after@abstracthead@skip}{1ex}{\z@}{\z@}
\@set@scalable@length{@between@abstract@skip}{2ex}{1ex}{0.5ex}
\@set@scalable@length{@before@kwd@skip}{1ex}{0.5ex}{0.2ex}
\@set@scalable@length{@before@dates@skip}{2ex}{0.5ex}{0.2ex}
%    \end{macrocode}
% \end{macro}\end{macro}\end{macro}\end{macro}
% \end{macro}\end{macro}\end{macro}\end{macro}
% 
% Spacing above and below displayed equations also wants to be done
% like this, but it is more complicated since the standard article
% classes make these change with the fontsize. Something to do later,
% maybe. 
% 
%    \begin{macrocode}
%    \end{macrocode}
% \begin{macro}{\@package@default@lengths}
% The value actually used for these lengths is controlled by the
% \meta{SCALE} parameter, which depends on the style option that is in
% effect. This is set using the |\@package@default@lengths| command,
% which must be called every time the style changes (e.g., changing
% from onecolumn to twocolumn). 
%    \begin{macrocode}
\newcommand{\@package@default@lengths}{%
  \ifthenelse{\boolean{@rmtwocol}}{%
    %% twocolumn mode
    \abovesectionskip@set@SCALE{0.5}
    \belowsectionskip@set@SCALE{0.5}
    \aboveparaskip@set@SCALE{0.5}
    \belowparaskip@set@SCALE{1}
    \floatsep@set@SCALE{0.5}
    \textfloatsep@set@SCALE{0.5}
    \intextsep@set@SCALE{0.5}
    \dblfloatsep@set@SCALE{0.5}
    \dbltextfloatsep@set@SCALE{0.5}
    }{%
    %% onecolumn mode
    \abovesectionskip@set@SCALE{1}
    \belowsectionskip@set@SCALE{1}
    \aboveparaskip@set@SCALE{1}
    \belowparaskip@set@SCALE{1}
    \floatsep@set@SCALE{1}
    \textfloatsep@set@SCALE{1}
    \intextsep@set@SCALE{1}
    \dblfloatsep@set@SCALE{1}
    \dbltextfloatsep@set@SCALE{1}
    }%
  }
\@package@default@lengths
%    \end{macrocode}
% \end{macro}
% \begin{macro}{\RescaleLengths}
% \begin{macro}{\RescaleSecLengths}
% \begin{macro}{\RescaleFloatLengths}
% \begin{macro}{\RescaleTitleLengths}
% User commands that will fiddle with all the scalable lengths. 
%    \begin{macrocode}
\newcommand{\RescaleLengths}[1]{%
  \RescaleSecLengths{#1}\RescaleFloatLengths{#1}\RescaleTitleLengths{#1}}
\newcommand{\RescaleSecLengths}[1]{%
%% sections
  \abovesectionskipScale{#1}
  \belowsectionskipScale{#1}
  \aboveparaskipScale{#1}
  }
\newcommand{\RescaleFloatLengths}[1]{%
%% floats
  \floatsepScale{#1}
  \textfloatsepScale{#1}
  \intextsepScale{#1}
  \dblfloatsepScale{#1}
  \dbltextfloatsepScale{#1}
  \abovecaptionskipScale{#1}
  }
\newcommand{\RescaleTitleLengths}[1]{%
%% maketitle
  \@before@title@skipScale{#1}           
  \@before@author@skipScale{#1}        
  \@between@author@skipScale{#1}       
  \@before@abstract@skipScale{#1}      
  \@after@abstracthead@skipScale{#1}   
  \@between@abstract@skipScale{#1}     
  \@before@kwd@skipScale{#1}           
  \@before@dates@skipScale{#1}           
  }
%    \end{macrocode}
% \end{macro}\end{macro}\end{macro}\end{macro}
% \changes{2.2beta14}{2002/03/17}{Finished implementation of new
% infrastructure for scalable lengths. Converted the above/below
% section skips to use it. Also, all the floatsep parameters. Also,
% all the vertical lengths that are used by |\@maketitle|. }

% \begin{macro}{\MSextramargin}
% For the \texttt{manuscript} option we have an extra 0.3\,inch
% vertical margin and 1\,inch horizontal margin for marginal notes (by
% hand or with |\marginpar|. 
%
% \changes{3.23}{2007/09/09}{Reduce baselinestretch in
%   \texttt{manuscript} style to 1.2. No longer use raggedright for
%   manuscript.}
%
% \changes{3.24}{2007/09/14}{Further reduction of baselinestretch in
%   \texttt{manuscript} style to 1.1 since we now have wider margins.}
%
% \changes{3.24}{2007/09/14}{Split \texttt{manuscript}'s extra margin
%   into separate horizontal and vertical parts.}
% 
%    \begin{macrocode}
\newlength{\MSextraHmargin}\newlength{\MSextraVmargin}
\setlength{\MSextraHmargin}{1in} 
\setlength{\MSextraVmargin}{0.3in} 
\if@manuscript
   \renewcommand{\baselinestretch}{1.1}
   \addtolength{\textwidth}{-2\MSextraHmargin}
   \addtolength{\oddsidemargin}{\MSextraHmargin}
   \setlength{\evensidemargin}{\oddsidemargin}
   \addtolength{\textheight}{-2\MSextraVmargin}
   \addtolength{\topmargin}{\MSextraVmargin}
%    \raggedright
\fi
%    \end{macrocode}
% \end{macro}
% Float parameters: don't allow bottom floats by default. 
%    \begin{macrocode}
\renewcommand{\fps@figure}{tp}
\renewcommand{\fps@table}{tp}
%    \end{macrocode}
%
% \begin{macro}{\rmaa@AuthorWidthFrac}
% \begin{macro}{\rmaa@AbstractWidthFrac}
% These parameters control what fraction of the total text width can
% be used by the list of authors and the abstract/resumen,
% respectively.  
%    \begin{macrocode}
\newcommand{\rmaa@AuthorWidthFrac}{1.0}
\newcommand{\rmaa@AbstractWidthFrac}{0.75}
%    \end{macrocode}
% Use a wide abstract to save space in the new style.
%    \begin{macrocode}
\if@newstyle\renewcommand{\rmaa@AbstractWidthFrac}{1.0}\fi
%    \end{macrocode}
% \end{macro}
% \end{macro}
%
% Some things to consider in future versions. Trying to avoid too many
% widows, etc. Apparently, the following will work best with
% \cs{flushbottom}. 
%    \begin{macrocode}
%%\widowpenalty=10000
%%\clubpenalty=10000
%    \end{macrocode}
% Also, should perhaps try getting rid of underfull/overfull box warnings. 
% \changes{3.10}{2003/01/24}{Set large values for \cs{hbadness},
% \cs{vbadness}, \cs{hfuzz}, \cs{vfuzz}. To reduce amount of stuff
% that gets written to the terminal every time you compile something.}
%    \begin{macrocode}
\hbadness=10000
\vbadness=10000
\setlength\hfuzz{\maxdimen}
\setlength\vfuzz{\maxdimen}
%    \end{macrocode}
% \subsection{Various text styles for different contexts}
% \label{sec:text-styles} 
% \begin{macro}{\rmaa@AffilStyle}
% \begin{macro}{\rmaa@PagesStyle}
% \begin{macro}{\rmaa@YearStyle}
% \begin{macro}{\rmaa@VolumeStyle}
% \begin{macro}{\rmaa@VolumenStyle}
% \begin{macro}{\rmaa@JournalNameStyle}
% \begin{macro}{\rmaa@ReceivedStyle}
% \begin{macro}{\rmaa@AbstractHeaderStyle}
% \begin{macro}{\rmaa@JournalNameStyle}
% Style commands for the different fields used by |\maketitle| in
% constructing the title page of the article. These have the form
% |\rmaa@XYZStyle{TEXT}|. I'm not sure these were really a good idea,
% but never mind \dots 
%    \begin{macrocode}
\newcommand{\rmaa@AffilStyle}[1]{{\normalsize\rmfamily#1}}
\newcommand{\rmaa@PagesStyle}[1]{{\normalsize\rmfamily#1}}
\newcommand{\rmaa@YearStyle}[1]{{\normalsize\rmfamily#1}}
\newcommand{\rmaa@VolumeStyle}[1]{{\normalsize\rmfamily\bfseries#1}}
\newcommand{\rmaa@VolumenStyle}[1]{{\Large\sffamily\bfseries#1}}
\newcommand{\rmaa@JournalNameStyle}[1]{{\normalsize\itshape#1}}
\newcommand{\rmaa@ReceivedStyle}[1]{{\normalsize\rmfamily\itshape#1}}
\newcommand{\rmaa@AbstractHeaderStyle}[1]{%
  {\normalsize\rmfamily\MakeUpperCase{#1}}}
\newcommand{\rmaa@TitleStyle}[1]{%
  {\large\bfseries\rmfamily\boldmath\MakeUpperCase{#1}\par}}
%    \end{macrocode}
% \end{macro}
% \end{macro}
% \end{macro}
% \end{macro}
% \end{macro}
% \end{macro}
% \end{macro}
% \end{macro}
% \end{macro}
%
% \begin{macro}{\rmaa@FrontSecStyle}
% Headings of ``front sections'' (see section~\ref{sec:front-sections})
%    \begin{macrocode}
\newcommand{\rmaa@FrontSecStyle}[1]{%
  {\normalsize\bfseries\rmfamily\MakeUpperCase{#1}}}
%    \end{macrocode}
% \end{macro}
%
% \begin{macro}{\rmaa@TOCTitleStyle}
% \begin{macro}{\rmaa@TOCPartStyle}
% \begin{macro}{\rmaa@IndexTitleStyle}
% \begin{macro}{\rmaaTOCAuthorStyle}
% \begin{macro}{\rmaaIndexAuthorStyle}
% Commands to set style of TOC and index entries.
% Rewritten to avoid unnecessary font commands in the .toc and .ind files
%    \begin{macrocode}
\newcommand{\rmaa@TOCTitleStyle}[1]{\MakeUpperCase{#1}}
\newcommand{\rmaa@TOCPartStyle}[1]{\rmaa@TOCTitleStyle{#1}}
\newcommand{\rmaa@IndexTitleStyle}[1]{#1}
\newcommand{\rmaaTOCAuthorStyle}[1]{%
  {\itshape#1}}
\newcommand{\rmaaIndexAuthorStyle}[1]{%
  \protect\mbox{\bfseries#1}}
%    \end{macrocode}
% \end{macro}
% \end{macro}
% \end{macro}
% \end{macro}
% \end{macro}
%
% \begin{macro}{\rmaa@KeywordStyle}
% \begin{macro}{\rmaa@KeywordHeaderStyle}
% The style of the keywords depends on whether we are using the
% \texttt{newstyle} option or not. If we are, then they are small. 
%    \begin{macrocode}
\ifthenelse{\boolean{@newstyle}}{
  \newcommand{\rmaa@KeywordStyle}[1]{%
    {\small%
      \if@boldkeywords 
      \bfseries\MakeUpperCase{\raggedright #1}
      \else
      \raggedright #1
      \fi
    }}
  \newcommand{\rmaa@KeywordHeaderStyle}[1]{%
    {\small\itshape#1}}}
{
  \newcommand{\rmaa@KeywordStyle}[1]{%
    {\normalsize%
      \if@boldkeywords\bfseries\fi\raggedright #1}}
  \newcommand{\rmaa@KeywordHeaderStyle}[1]{%
    {\normalsize\itshape#1}}}
%    \end{macrocode}
% \end{macro}
% \end{macro}
%
% \subsection{Page styles}\label{sec:page-styles}
% \begin{macro}{\ps@rmaa}
% The running headers for normal pages:
%    \begin{macrocode}
\newcommand{\ps@rmaa}{%
%    \end{macrocode}
% nothing in either foot; 
%    \begin{macrocode}
  \renewcommand\@oddfoot{\relax}%
  \renewcommand\@evenfoot{\relax}%
%    \end{macrocode}
% page number and short author list in the even heads; 
%    \begin{macrocode}
  \renewcommand\@evenhead{\rlap{\thepage}\hfil\@shortauthor\hfil}%
%    \end{macrocode}
% short title and page number in the odd heads.
%    \begin{macrocode}
  \renewcommand\@oddhead{\hfil\@shorttitle\hfil\llap{\thepage}}
  }
%    \end{macrocode}
% \begin{macro}{\ps@rmaafirstpage}
% The header of the first page of each
% article:
%    \begin{macrocode}
\newcommand{\ps@rmaafirstpage}{%
  \renewcommand\@oddfoot{\hfil\thepage}%
%    \end{macrocode}
% footer corrected (06 Jun 2000) to have page on outside edge;
%    \begin{macrocode}
  \renewcommand\@evenfoot{\thepage\hfil}%
%    \end{macrocode}
% Header has journal name, volume, and page numbers for a regular
% article, or similar for special cases. 
% \changes{3.27}{2019/11/03}{Change to three-line header with copyright info and DOI.}
%    \begin{macrocode}
\if@preprint
\renewcommand\@oddhead{%
  \if@proceedings 
  \textsc{To appear in ``\@conftitle{} (\rmaa@YearStyle{\@year})''}\hfil
  \textit{RMxAC}
  \else
  \textsc{To appear in} \@journalname\hfil
  \fi}\fi%
\if@guide
\renewcommand\@oddhead{\textsc{Documentation} (\rmaa@YearStyle{\@year})\hfil}
\fi%
\if@manuscript
\renewcommand\@oddhead{Manuscript for \@journalname{} (\rmaa@YearStyle{\@year})\hfil}
\fi%
\if@journal
\renewcommand\@oddhead{%
  \parbox{\textwidth}{
    \@journalname, \rmaa@VolumeStyle{\@volume},
    \rmaa@PagesStyle{\@artpages}\
    (\rmaa@YearStyle{\@year})\hfil\\[1pt]
    {\small\textcopyright{} \@year:
      Instituto de Astronom\'\i{}a,
      Universidad Nacional Aut\'onoma de M\'exico}\hfil\\[1pt]
    {\small\@DOI@url}\hfil
  }%
}\fi%
  \renewcommand\@evenhead{\@oddhead}}
%    \end{macrocode}
% By default, use the |\ps@rmaa| page style.
%    \begin{macrocode}
\pagestyle{rmaa}
%    \end{macrocode}
% \end{macro}
% \end{macro}
%
% \subsection{Preamble commands}\label{sec:preamble} 
% Everything you need for setting the title page, TOC and index
% entries, etc.
% \subsubsection{Commands for the title page}
% \label{sec:title-commands}
% \begin{macro}{\resumen}
% \begin{macro}{\@resumen}
% \begin{macro}{\abstract}
% \begin{macro}{\@abstract}
% User commands for setting the resumen and abstract. 
% These should be used in the document header. The internal commands
% |\@resumen| and |\@abstract| are used to hold the actual text
% pending output via the |\maketitle| command. 
%    \begin{macrocode}
\newif\if@setresumen
\newcommand{\resumen}[1]{%
%%  \renewcommand{\@resumen}{{\selectlanguage{spanish}#1}}%
  \renewcommand{\@resumen}{{\ignorespaces#1}}%
  \@setresumentrue}
\renewcommand{\abstract}[1]{%
  \renewcommand{\@abstract}{{\ignorespaces#1}}
  \renewcommand\ADS@abstract{#1}
%% Not a good idea, since then author macros are not expanded and
%% cause errors later. However, by not doing it we are vulnerable to
%% fragile commands.....
  \if@noexpandADS\@onelevel@sanitize{\ADS@abstract}\fi
  \@abstractpresenttrue
  }
\newcommand{\ADSnoexpand}{\@noexpandADStrue}
%    \end{macrocode}
% \begin{macro}{\boldabstract}
% This is the version for doing 1-page conference proceedings. Still
% save the text for ADS but then write it out immediately as a bold
% paragraph. 
%    \begin{macrocode}
\newcommand{\boldabstract}[1]{%
  \gdef\ADS@abstract{#1}%
  \renewcommand{\@abstract}{#1}%
  \if@noexpandADS\@onelevel@sanitize{\ADS@abstract}\fi%
  \noindent\textbf{\boldmath #1}
  }
%    \end{macrocode}
% \end{macro}
% Default value in each case is to print out a warning if not set
%    \begin{macrocode}
  \newcommand{\@resumen}{%
    \@setresumenfalse \ClassWarningNoLine{rmaa}{No Spanish abstract
      (resumen) specified - if you do not speak Spanish, this will be
      translated by the editors}%
    \textbf{Favor de proporcionar un resumen en espa\~nol. If you are
      unable to translate your abstract into Spanish, the editors will
      do it for you.} }
  \newcommand{\@abstract}{\ClassWarningNoLine{rmaa}{No abstract
      specified}}
%    \end{macrocode}
% \end{macro}
% \end{macro}
% \end{macro}
% \end{macro}
%
% \begin{macro}{\title}
% \begin{macro}{\@rawtitle}
% \begin{macro}{\@TOCtitle}
% \begin{macro}{\rmaaIndexTitle}
% Similarly, |\title| is the preamble command to set the title of the
% article. This is used in various places, typeset in different
% styles, hence the use of the internal |\@rawtitle|, |\@title|,
% |\rmaaIndexTitle|, and |\@TOCtitle| commands, each of which uses the
% relevant style from section~\ref{sec:text-styles}
%    \begin{macrocode}
\newcommand{\@rawtitle}{Raw Title}
\newcommand{\@TOCtitle}{\rmaa@TOCTitleStyle{Title for TOC}}
\newcommand{\rmaaIndexTitle}{\rmaa@IndexTitleStyle{Title for index}.}
\newcommand{\@TOCaftertitle}{}
\renewcommand{\title}[1]{%
%%  {\ADS@sanitize\xdef\@rawtitle{#1}}
  \renewcommand\@rawtitle{\noexpand#1}
  \let\ADS@title=\@rawtitle
  \renewcommand{\@title}{\rmaa@TitleStyle{#1}}%
  \renewcommand{\@TOCtitle}{\rmaa@TOCTitleStyle{#1}\@TOCaftertitle}%
  \renewcommand{\rmaaIndexTitle}{\rmaa@IndexTitleStyle{#1\@TOCaftertitle}.}%
  }
%    \end{macrocode}
% \end{macro}
% \end{macro}
% \end{macro}
% \end{macro}
% 
% All the weird and whacky ways of indicating the author
% affiliations. 
% \begin{macro}{\affil}
% First, the way that sets it directly underneath the authors name.
% Note that |\affil| trumps all other affiliations for the purposes of
% getting into the ADS stuff (we use |\gdef|).   
% \dots 
%    \begin{macrocode}
\newcommand{\affil}[1]{%
  \\[1.5ex]\rmaa@AffilStyle{#1}\\[1.5ex]%
  \if@firstaffil\gdef\ADS@affil{#1}\@firstaffilfalse\fi%
  }
\renewcommand{\and}{and\\[1.5ex]}
%    \end{macrocode}
% \begin{macro}{\altaffil}
% Alternative affiliations are footnoted. 
%    \begin{macrocode}
\newcommand{\altaffil}[1]{%
  \if@firstaltaffil\gdef\ADS@affil{#1}\@firstaltaffilfalse\fi%
  \thanks{#1}}
%    \end{macrocode}
% \begin{macro}{\altaffilmark}
% \begin{macro}{\altaffiltext}
% Alternative-alternative mechanism for affiliations---should not be
% mixed with |\altaffil| (unless you know what you are doing) or there
% will be problems with numbering.
%    \begin{macrocode}
\newcommand{\altaffilmark}[1]{\textsuperscript{#1}}
\newcommand{\altaffiltext}[2]{%
  \if@firstaltaffil\gdef\ADS@affil{#2}\@firstaltaffilfalse\fi
  \stepcounter{footnote}%
  \protected@xdef\@thanks{\@thanks
    \protect\footnotetext[#1]{#2}}} 
%    \end{macrocode}
% \end{macro}
% \end{macro}
% \end{macro}
% \end{macro}
%
% \begin{macro}{\keywords}
% \begin{macro}{\@keywords}
% \begin{macro}{\@keywordheader}
% \label{keywords}Sundry minor preamble commands. All
% self-explanatory, I hope. |\keywords| is the old user interface
% (deprecated). New articles should use |\addkeyword| once for each
% keyword. That way the ADS record comes out right too (comma rather
% than em-dash separator).
% \changes{2.2beta12}{2001/12/07}{\cs{addkeyword} New user interface
%   for keywords.} 
% \changes{2.2beta13}{2002/03/11}{\cs{addkeyword} Improved behaviour
%   with .ads file. Use
%   \cs{ADS@sanitize} to change definitions of `~' and \cs{ion}.}
% \changes{3.23}{2007/09/09}{\cs{addkeyword} now signals a compilation
% error if the author tries to add more than 6 keywords.}
% \changes{3.24}{2007/09/14}{kwdcounter is now reset in
% \cs{@begofarticlehook} to prevent error in book style} 
%    \begin{macrocode}
\newcommand{\@kwdsep}{ --- }
\newcommand{\ADS@kwdsep}{, }
\newcounter{kwdcounter}\setcounter{kwdcounter}{0}
\newcommand{\addkeyword}[1]{%
  \addtocounter{kwdcounter}{1}
  \ifthenelse{\thekwdcounter > 6}{
    \ClassError{rmaa}{Too many keywords. A maximum of 6 keywords is allowed.}{}
  }{}
  \if@firstkwd%
  \renewcommand\@keywords{\if@boldkeywords\uppercase{#1}\else #1\fi}
  {\ADS@sanitize\def\ion{\plainion}\gdef\ADS@keywords{#1}}
  \@firstkwdfalse
  \else%
  \edef\@keywords{\@keywords\@kwdsep \if@boldkeywords\uppercase{#1}\else #1\fi}%
  {\ADS@sanitize\def\ion{\plainion}\xdef\ADS@keywords{\ADS@keywords\ADS@kwdsep #1}}%
  \fi}
\newcommand{\ADS@sanitize}{%
  \def\mathrm{}
  \def\@{\ }
  \def~{\ }
  \def\langle{<}
  \def\rangle{>}
  \def\log{log }
  \def\sin{sin }
  \def\cos{cos }
  \def\tan{tan }
  \def\fd{^d.}
  \def\fh{^h.}
  \def\fm{^m.}
  \def\fs{^s.}
  \def\fdg{�.}
  \def\farcm{'.}
  \def\farcs{''.}
  \def\arcdeg{�}
  \def\,{ }
  }
\newcommand{\keywords}[1]{%
  \renewcommand{\ADS@keywords}{#1}
  \renewcommand{\@keywords}{\uppercase{#1}}
  }
\newcommand{\@keywordheader}{\rmaa@KeywordHeaderStyle{Key Words:}}
\newcommand{\@keywords}{%
  \ClassWarningNoLine{rmaa}{No keywords specified}%
  Keywords go here
  }
%    \end{macrocode}
% \begin{macro}{\shorttitle}
% \begin{macro}{\shortauthor}
% \begin{macro}{\@shorttitle}
% \begin{macro}{\@shortauthor}
% Commands for setting the running headers.
%    \begin{macrocode}
\newcommand{\shorttitle}[1]{%
  \renewcommand\@shorttitle{\MakeUpperCase{#1}}}
\newcommand{\shortauthor}[1]{%
  \renewcommand\@shortauthor{\MakeUpperCase{#1}}}
\newcommand{\@shorttitle}{SHORT TITLE GOES HERE}
\newcommand{\@shortauthor}{SHORT AUTHOR GOES HERE}
%    \end{macrocode}
% \end{macro}
% \end{macro}
% \end{macro}
% \end{macro}
% \end{macro}
% \end{macro}
% \end{macro}


% \subsubsection{The full addresses}
% \label{sec:full-addresses}
% \begin{macro}{\fulladdresses}
% \begin{macro}{\@fulladdresses}
% \begin{macro}{\@fulladdheight}
% \begin{macro}{\@fulladddepth}
% \begin{macro}{\@fulladdbox}
% Aaghh, the full addresses! A perennial source of problems. 
%    \begin{macrocode}
\newlength{\@fulladdheight}
\newlength{\@fulladddepth}
\newsavebox{\@fulladdbox}
%% \fulladdresses command modified (11 May 2000) so that we can put
%% \fulladdresses{} in the article to suppress the output of the
%% addresses
\newcommand{\@FinalFigure}{\relax}
\newcommand{\FinalFigure}[1]{\renewcommand{\@FinalFigure}{#1}}
\newcommand{\fulladdresses}[1]{%
  \renewcommand{\@rmaaArg}{#1}
  \ifx\@rmaaArg\@rmaaEmpty
     \renewcommand{\@fulladdresses}{}
  \else
     \renewcommand{\@fulladdresses}{
       \normalsize%
       \ifx\@FinalFigure\relax\else\vfill\@FinalFigure\vfill\fi
%%       \vspace*{1ex}
       %% do addresses as a list (copied from thebibliography environment) 
       \list{\null}{\leftmargin 1.5em\topsep\z@\labelwidth\z@\labelsep\z@\itemindent -1.5em
         \itemsep\z@\parsep\z@\usecounter{enumi}}
       \sloppy\clubpenalty4000\widowpenalty4000
       \sfcode`\.=1000\relax
       #1
       \endlist
       }
  \fi
  }
\newcommand{\@fulladdresses}{Full addresses go here}
%    \end{macrocode}
% \begin{macro}{\suppressfulladdresses}
% New user-level command for suppressing full addresses should you so
% desire. Currently used in the single-page-poster style. 
%    \begin{macrocode}
\newcommand{\suppressfulladdresses}{\@suppressfulladdressestrue\@needaddressesfalse}
%    \end{macrocode}
% \end{macro}
% \end{macro}
% \end{macro}
% \end{macro}
% \end{macro}
% \end{macro}

% \begin{macro}{\printdates}
% \begin{macro}{\noprintdates}
% Command to activate the printing out of the received/accepted
% dates. This can also be done using the \texttt{printdates} option,
% which is the default for \texttt{journal}, but not for \texttt{proceedings}.
%    \begin{macrocode}
\newcommand{\printdates}{\@printdatestrue}
\newcommand{\noprintdates}{\@printdatesfalse}
%    \end{macrocode}
% \end{macro}
% \end{macro}
% \begin{macro}{\@outputaddressesifnecessary}
% Command to output the addresses if we haven't already done so
% This is called by: 
% \begin{enumerate}
% \item The end of |thebibliography| environment
% \item The |\outputfulladdresses| command
% \end{enumerate}
%    \begin{macrocode}
\newcommand{\@outputaddressesifnecessary}{%
  \ClassInfo{rmaa}{Calling @outputaddressesifnecessary}
  \ifthenelse{\boolean{@needaddresses}}{%
    %% If in one column style, or in the first column of a two column
    %% style then we can try to put it on this page 
    \ifthenelse{\boolean{@firstcolumn} \or \not\boolean{@rmtwocol}}{%
      \@outputaddresses}{%
      %% otherwise give warning it will be on a separate page
      \ClassWarning{rmaa}{Full addresses will go on a separate
        page (please see User Guide for how to remedy this)}}%
    %% If we haven't managed yet, put it on its own page
    \ifthenelse{\boolean{@needaddresses}}{%
      \if@debug{\ClassWarning{rmaa}{In this silly bit!}}\fi
      \clearpage\vfill\@outputaddresses}{}%
    }{
    %% If we don't need them, we still might need to do the eoa hook 
    \ifthenelse{\boolean{@suppressfulladdresses}}{%
      \if@debug{\ClassWarning{rmaa}{XXX}}\fi
%%      \if@book\else\@endofarticlehook\null\fi
      }{}
    }%
  }  
%    \end{macrocode}
% \end{macro}

% \begin{macro}{\outputfulladdresses}
% A user interface to the above command, in case there is no biblio
% (11 May 2000) 
%    \begin{macrocode}
\newcommand{\outputfulladdresses}{\@outputaddressesifnecessary}
%    \end{macrocode}
% \end{macro}

% \begin{macro}{\adjustfinalcols}
% This command starts a new column and should be 
% used to ensure that the columns balance on the last page. It also
% outputs the full 
% author addresses to ensure that this occurs in the left hand
% column. 
% The |\par| command is so that the current line will be taken account
% of when working out |\pagetotal|
%    \begin{macrocode}
\newcommand{\adjustfinalcols}{\if@twocolumn\goodbreak\@outputaddresses\newpage\noindent\fi}
%    \end{macrocode}
% \end{macro}

% \begin{macro}{\@outputaddresses}
% Command to do the work of outputting the addresses. 
% This can be called by 
% \begin{enumerate}
% \item an explicit |\adjustfinalcols|
% \item via |\@outputaddressesifnecessary|, q.v.
% \end{enumerate}
%
%    \begin{macrocode}
\newcommand{\@outputaddresses}{%
  % Allow option of suppressing the full addresses, in which case
  % we just run the hook if we are not in a book style
  \if@suppressfulladdresses
%%     \if@book\else\@endofarticlehook\null\fi
  \else
  %% save the full addresses in a box for output later
  \sbox{\@fulladdbox}{\parbox{\textwidth}{\@fulladdresses}}
  % calculate height of address box  
  \settoheight{\@fulladdheight}{\usebox{\@fulladdbox}}
  \settodepth{\@fulladddepth}{\usebox{\@fulladdbox}}
  \addtolength{\@fulladdheight}{\@fulladddepth}
  \if@debug
  { \typeout{Address total height: \the\@fulladdheight}
    \typeout{Text height: \the\pagetotal}
    \typeout{Goal height: \the\pagegoal}
    }
  \fi
  % reuse \@fulladdheight to mean the total height of text+addresses
  \addtolength{\@fulladdheight}{\pagetotal}
%%  \addtolength{\@fulladdheight}{-\baselineskip}
  \if@debug\typeout{Adjusted text + addresses height: \the\@fulladdheight}\fi
  % check it will fit on current page
  \ifdim\@fulladdheight<\pagegoal
     %% must shift it to the left
     %% to compensate for the hanging indentation in the reference list. 
     %% NOTE: this is not foolproof!
     %% NB. This will give overfull hbox warnings, which is a bit ugly.
     \vfill\noindent\if@biblist\hspace*{-\leftmargin}\fi
     \usebox{\@fulladdbox}
     % Set this to show we've done it
     \@needaddressesfalse
     % and assume we are at the end of the article
     % so that we do the EOA hook if not in the book style
%%     \if@book\else\@endofarticlehook\fi
  \else
     % if not, just do nothing. With luck, this function will get 
     % called again on the next page
     \if@debug{\typeout{Holding over full addresses to next page}}\fi
  \fi\fi
  }
%    \end{macrocode}
% \end{macro}

% \subsubsection{Commands for editors}
% \label{sec:editor-commands}
% Some user/editor commands for doing the TOC and index in the book
% option. 
% They are defined here rather in the book-specific part since they
% need to be always available in order to avoid breaking documents. 
% \begin{macro}{\listofauthors}
% \begin{macro}{\@rawauthors}
% \begin{macro}{\rmaaTOCauthors}
% The command |\listofauthors| formats an author list for use in the
% TOC. It should contain the same names as in the |\author| command,
% but without the affiliation garbage and using |\&| instead of
% |and|. 
%    \begin{macrocode}
\newcommand{\@rawauthors}{Raw Authors}
\newcommand{\rmaaTOCauthors}{\noexpand\rmaaTOCAuthorStyle{Author list for TOC}}
\newcommand{\listofauthors}[1]{%
  \renewcommand{\@rawauthors}{#1}
  \renewcommand\rmaaTOCauthors{\noexpand\rmaaTOCAuthorStyle{#1}}}
%    \end{macrocode}
% \begin{macro}{\indexauthor}\label{indexauthor}
% The command |\indexauthor| produces an index entry for each
% author. 
% This command must come \emph{after} the
% |\title| and the |\listofauthors| commands. Also, there should be one 
% |\indexauthor| command for \emph{each} author. We also add the
% author onto the end of the |\ADS@authors| macro and construct the
% |\ADS@authorinitial| by a most cunning method (see below), but only
% for the first author. 
% 
% No, No, NO! This was all wrong. Backed out completely: 
% \begin{quote}
% Completely rewritten (19 Dec 2002) since we now use the index class, so 
% as to be able to alphabetize accented characters. A side effect of this is
% that |\index| no longer expands out the argument after the |@|, so we 
% instead add all the |\index| commands to a list (|\indexauthor@deferlist|), 
% which is executed in the |\@begofarticlehook|. 
% 
% Note the stuff with |\toks0| is copied from Exercise 20.15 of the \TeX{}book---it is
% to make sure the latin1-accented name gets passed into the idx
% file. 
% \end{quote}
% We now go back to the old way of doing things, using an optional
% argument to the \cs{indexauthor} command that has the name without
% accents. 
% \changes{3.08}{2003/01/16}{Remerged changes from 3.06 that got
% dropped from 3.07 after a synchronization disaster. Basically, the
% accentless optional argument to \cs{indexauthor}. Must add the macro
% directories to my unison set-up to stop this sort of thing from happening.}
%    \begin{macrocode}
\newcommand\@isortkey{}
\newcommand\@emptykey{empty}
\newcommand\indexauthor[2][empty]{
  \renewcommand\@isortkey{#1}
  \ifx\@isortkey\@emptykey \renewcommand\@isortkey{#2}\fi
  \index{\@isortkey @{{\rmaaIndexAuthorStyle{#2}}\hskip1em plus1em minus0.4em
      \rmaaIndexTitle\hskip1em plus2em minus0.4em\penalty-100
      \rmaaTOCauthors}}
  \indexauthorADS{#2}
}
%%\def\indexauthor@deferlist{\relax}
%%\long\def\indexauthor#1{
%%  \def\@@name{#1} \toks0=\expandafter{\@@name}
%%  \toks1=\expandafter{\indexauthor@deferlist}
%%  \protected@xdef\indexauthor@deferlist{%
%%    \the\toks1
%%    \noexpand\index{\the\toks0@{\noexpand\rmaaIndexAuthorStyle{\the\toks0}\noexpand\hskip1em plus1em minus0.4em 
%%        \rmaaIndexTitle\noexpand\hskip1em plus2em minus0.4em\noexpand\penalty-100
%%        \rmaaTOCauthors}}
%%  }
%%  \indexauthorADS{#1}
%%  }
%%\xdef\indexauthor@deferlist{}
\newcommand\indexauthorADS[1]{%
  \if@firstauthor\protected@xdef\ADS@authorinitial{\firstletter{#1}}\fi
  \protected@xdef\ADS@authors{\ADS@authors\if@firstauthor\else ; \fi
    #1}
  \@firstauthorfalse
}
%    \end{macrocode}
% Note that the first |#1| sets the key for the index, making sure it
% actually starts with  
% first letter of the surname. Meanwhile, the |@{...}| expression,
% which starts with some  
% control sequences is what is written to the |.ind| file. 
% See the makeindex man page for details. 
% \end{macro}
% \end{macro}
% \end{macro}
% \end{macro}

% \begin{macro}{\KillAllAccents}
% \begin{macro}{\firstletter}
% These are utility functions used in the above. |lett| is a counter
% indicating the current letter in the string. |\firstletter| just
% prints the first letter of its argument.
%    \begin{macrocode}
\newcommand{\KillAllAccents}{\def\'{}\def\`{}\def\~{}\def\.{}\def\i{i}}
\newcommand\firstletter[1]{\@car #1\@nil}
%    \end{macrocode}
% \end{macro}
% \end{macro}
% \begin{macro}{\SetVolume}
% \begin{macro}{\SetMSnumber}
% \begin{macro}{\SetFirstPage}
% \begin{macro}{\SetYear}
% \begin{macro}{\ReceivedDate}
% \begin{macro}{\AcceptedDate}
% These are for the editors to set the volume number, etc. 
%    \begin{macrocode}
\newcommand{\SetVolume}[1]{%
  \renewcommand{\@volume}{#1}
  \edef\ADS@volume{\PadToFourWithDots{#1}}
  \renewcommand{\@volumen}{\rmaa@VolumenStyle{volumen #1}}%
  }
\newcommand{\SetMSnumber}[1]{\renewcommand{\@MSnum}{#1}}
%    \end{macrocode}
% Note that there is no real reason to use |\SetFirstPage| since all
% the page numbering is done automatically in the book style. 
%    \begin{macrocode}
\newcommand{\SetFirstPage}[1]{\setcounter{page}{#1}}
\newcommand{\SetYear}[1]{%
  \renewcommand{\@year}{#1}
  \renewcommand{\ADS@year}{#1}
  \renewcommand{\@anyo}{\rmaa@VolumenStyle{#1}}
  }
\newcommand{\ReceivedDate}[1]{\renewcommand{\@received}{#1}}
\newcommand{\AcceptedDate}[1]{\renewcommand{\@accepted}{#1}}
%    \end{macrocode}
%
% \begin{macro}{\DOI}
% User command to enter the DOI, which will be typeset as a link.
% \changes{3.27}{2019/11/02}{New user command |\DOI{2019.55.02.21}| to typeset
% DOI in the header of first page.}
%    \begin{macrocode}
\newcommand\DOI[1]{\def\@DOI@url{\url{https://doi.org/10.22201/ia.01851101p.#1}}}
%    \end{macrocode}
% \end{macro}
% \begin{macro}{\@volume}
% \begin{macro}{\@volumen}
% \begin{macro}{\@artpages}
% \begin{macro}{\@year}
% \begin{macro}{\@anyo}
% \begin{macro}{\@received}
% \begin{macro}{\@accepted}
% \begin{macro}{\@MSnum}
% Default values for the containers used by above commands, plus a
% few more useful ones. 
%    \begin{macrocode}
\newcommand{\@volume}{00}
\newcommand{\@volumen}{\rmaa@VolumenStyle{00}}
\newcommand{\@artpages}{%
  \oldpageref{\articlename:RMAAFirstPage}--\oldpageref{\articlename:RMAALastPage}}
\newcommand{\@year}{\the\year}
\newcommand{\@anyo}{\rmaa@VolumenStyle{2007}}
\newcommand{\@journalname}{\rmaa@JournalNameStyle{%
    \if@proceedings\RMAASClong\else\RMAAlong\fi}}
\newcommand{\@received}{}
\newcommand{\@accepted}{}
\newcommand{\@MSnum}{0000}
\newcommand{\@DOI@prefix}{https://doi.org/10.22201/ia.01851101p}
\def\@DOI@url{\url{https://doi.org/10.22201/ia.01851101p.2099.00.00.00}}
%    \end{macrocode}
% \end{macro}
% \end{macro}
% \end{macro}
% \end{macro}
% \end{macro}
% \end{macro}
% \end{macro}
% \end{macro}
% \end{macro}
% \end{macro}
% \end{macro}
% \end{macro}
% \end{macro}
% \end{macro}
% 
% The following commands are most useful in the \texttt{book}
% style. However, the |\SetConfTitle| command can now be used in
% individual articles to set the header for the \texttt{preprint}
% style. 
%    \begin{macrocode}
\newcommand{\@reunionplace}{Lugar de la reuni\'on}
\newcommand{\@reuniondate}{Fecha de la reuni\'on}
\newcommand{\@conftitle}{Conference Title}
\newcommand{\SetReunion}[2]{\renewcommand{\@reunionplace}{#1}\renewcommand{\@reuniondate}{#2}}
\newcommand{\SetConfTitle}[1]{\renewcommand{\@conftitle}{#1}}
\newcommand{\@editors}{Juan Ram\'on de la Fuente and ``El Mosh''}
\newcommand{\SetEditors}[1]{\renewcommand{\@editors}{#1}}
\newcommand{\@editores}{Juan Ram\'on de la Fuente y ``El Mosh''}
\newcommand{\SetEditores}[1]{\renewcommand{\@editores}{#1}}
\newcommand{\@mes}{mes de la publicaci\'on}
\newcommand{\SetMes}[1]{\renewcommand{\@mes}{\rmaa@VolumenStyle{#1}}}
%    \end{macrocode}

% \subsection{Notes for Author Proofs and Communication between Editors}
% \label{sec:notes}
% We introduce a command \cs{Note} that should be used to describe
% any change made to the manuscript. These are a bit tricky to get
% right. In general we use \cs{marginpar} but this only works in outer
% paragraph mode, so we just bung the box where we are if we are in a
% figure caption or something. If we are in the preamble, then we just
% use a \marginpar as part of the \maketitle command.
%    \begin{macrocode}
\newif\if@arenotes
\@arenotesfalse
\newcounter{notecounter}
\setcounter{notecounter}{0}
\newcommand\NotesList{}
\newcommand\Note[1]{%
  \if@proof{%
    \@arenotestrue%
    \stepcounter{notecounter}%
    \if@inpreamble\else{%
      \NoteMark{\arabic{notecounter}}%
    }\fi
    \protected@xdef\NotesList{\NotesList 
      \protect\item #1}%
  }\fi%
}
\newcommand\Changed[2]{Changed ``#1'' to ``#2''.}
\newcommand\NoteMark[1]{%
  \ifinner
  \raisebox{1ex}[0pt][0pt]{\makebox[0pt][l]{\framebox{\footnotesize\bfseries
        N#1}}}%
  \else
  \marginpar{\framebox{\footnotesize\bfseries N#1}}%
  \fi
}
\newcommand\RefNote[1]{\if@proof\NoteMark{\ref{#1}}\relax\fi}
\newcommand\NoteLabel{\protect\label}
\newcommand\OutputNotes{%
  \bgroup
  \onecolumn
  \section*{EDITORIAL NOTES}
  \begin{enumerate}
    \NotesList
  \end{enumerate}
  \egroup
}
\DeclareRobustCommand\cs[1]{\texttt{\char`\\#1}}

%    \end{macrocode}
% \subsection{Bibliographic Records for ADS}
% \label{sec:ADS}
% \changes{2.2beta11}{2001/11/27}{Just about finished all the ADS
% stuff now. Except for the references, which I have given up on for
% the time being.}
% \begin{macro}{\ADS@BibCode}
% \begin{macro}{\ADS@year}
% \begin{macro}{\ADS@journal}
% \begin{macro}{\ADS@volume}
% \begin{macro}{\ADS@firstpage}
% \begin{macro}{\ADS@authorinitial}
% We use the |\ADS@BibCode| command to hold the bibliographic code for
% ADS. It is 
% used in the |%R| 
% field of the bibliographic record for each article. 
% The format is 
% \begin{quote}
% \texttt{YYYYJJJJJVVVVMPPPPA}  
% \end{quote}
% where \texttt{YYYY} is the year, \texttt{JJJJJ} is the journal code,
% which is \texttt{RMxAA} or \texttt{RMxAC}, \texttt{VVVV} is the
% volume number (right-justified), \texttt{M} is not used by us,
% \texttt{PPPP} is the first page number (right-justified), and
% \texttt{A} is the initial of the first author's surname. The volume
% and page fields are padded with periods (`\texttt{.}'), which causes
% us some fun (this is done with the macro |\PadToFourWithDots| described
% below) . 
% 
% All these individual bits are held in the commands |\ADS@year|,
% |\ADS@journal|, |\ADS@volume|, |\ADS@firstpage|,
% |\ADS@authorinitial|. 
%    \begin{macrocode}
\newcommand{\SetADSBibCode}[1]{\renewcommand{\ADS@BibCode}{#1}}
\newcommand{\ADS@BibCode}
{\ADS@year\ADS@journal\ADS@volume.\PadToFourWithDots{\ADS@firstpage}\ADS@authorinitial}
\newcommand{\ADS@year}{2001} 
\newcommand{\ADS@journal}{\if@proceedings RMxAC\else RMxAA\fi}
\newcommand{\ADS@volume}{....} 
%    \end{macrocode}
% These two are just set to default values here, |\ADS@firstpage| is set
% in the |document| environment for included articles, in the
% |frontsection| environment, |\tableofcontents| command,  |index| and
% |abstracts| environments. Also by |\AtBeginDocument| for stand-alone
% articles. 
% 
% God knows how we are going to do |\ADS@authorinitial|!
%    \begin{macrocode}
\newcommand{\ADS@firstpage}{....}
\newcommand{\ADS@authorinitial}{A}
%    \end{macrocode}
% \end{macro}
% \end{macro}
% \end{macro}
% \end{macro}
% \end{macro}
% \end{macro}
% We also have a whole slew of other fields that can be set. These are
% all listed in Table~\ref{tab:ADS}. Most of them still have to be
% defined. 
% \begin{table}[!ht]\raggedleft
% \caption{ADS Bibliographic Records}\label{tab:ADS}
% \begin{tabular}{lp{8em}lp{8em}}\hline
% Code & Description & Macro Name & Comment \\ \hline
% |%R| & Bibliographic Code & |\ADS@BibCode| & required \\
% |%A| & Author list & |\ADS@authors| & required, generate from |\indexauthor| \\
% |%F| & Author affiliation & |\ADS@affil| & format: AA(\dots)
% AB(\dots) \dots \\
% |%J| & Journal Name & |\ADS@journalname| &  \\
% |%D| & Publication Date & |\ADS@fulldate| &  \\
% |%L| & Last Page & |\ADS@lastpage| & generate from label \\
% |%T| & Title & |\ADS@title| & required, generate from |\title| \\
% |%C| & Abstract Copyright & |\ADS@copyright| & ask UNAM lawyers! \\
% |%U| & URL for electronic document & |\ADS@url| & hmmmm \\
% |%K| & Keywords & |\ADS@keywords| & set equal to |\@keywords| \\
% |%A| & Abstract Text & |\ADS@abstract| & set equal to |\@abstract| \\
% |%R| & References & |\ADS@references| & somehow generate from the
% bibliography  \\ \hline
% \end{tabular}
% \end{table}
% A lot of these macros are set by other commands, as mentioned in the
% Table, but some need to be set  explicitly here. 
% 
% \begin{macro}{\ADS@journalname}
% This
% does not contain any accents at the moment so as not to cause
% problems---this maybe should be changed though. 
%    \begin{macrocode}
\newcommand{\ADS@journalname}{%
  \if@proceedings\@conftitle\space (Eds. \@editors) \fi%
  Revista Mexicana de Astronom\'\i{}a y Astrof\'\i{}sica%
  \if@proceedings\space (Serie de Conferencias)\fi%
  \space Vol. \@volume, pp. \ADS@firstpage--\ADS@lastpage\space (\@year)%
  \space (http://www.astroscu.unam.mx/\~{}rmaa/)
  }
%    \end{macrocode}
% \end{macro}
% \begin{macro}{\ADS@keywords}
% \begin{macro}{\ADS@abstract}
% \begin{macro}{\ADS@title}
% |\ADS@abstract| is straightforward, since we already have the text
% saved as |\@abstract|. We could do the same with |\ADS@keywords| but
% to avoid unwanted formatting we explicitly set it in |\keywords|
% (p.~\pageref{keywords}). |\ADS@title| is just set to |\@rawtitle|,
% which we were already using for another purpose. 
%    \begin{macrocode}
\let\ADS@title=\@rawtitle
\newcommand{\ADS@keywords}{No keywords specified}
\newcommand{\ADS@abstract}{No abstract specified}
%    \end{macrocode}
% \end{macro}\end{macro}\end{macro}
% \begin{macro}{\ADS@authors}
% |\ADS@authors| wants to be in ``Other, A. N.'' format, but with
% semi-colons between authors. The most natural way to do this seems to
% be with |\indexauthor| (p.~\pageref{indexauthor}). We append each
% author to this macro, which means we must explicitly clear it for
% each new article (done in |\article|). .   
%    \begin{macrocode}
\newcommand{\ADS@authors}{}
%    \end{macrocode}
% \end{macro}
% \begin{macro}{\ADS@fulldate}
% \begin{macro}{\ADS@lastpage}
% \begin{macro}{\ADS@copyright}
% \begin{macro}{\ADS@website}
% \begin{macro}{\ADS@url}
% \begin{macro}{\SetWebsite}
% \begin{macro}{\SetMonthNumeric}
% Other sundry data, together with user commands for bits that can be
% set by the Editors. The |\ADS@url| is constructed from the website
% (currently |http://www.astrosmo.unam.mx/rmaa/|,
% the journal descriptor (e.g., |RMxAC..10/|), followed by the
% path to the PDF file (e.g., |PDF/RMxAC..10_HEW1.pdf|). 
% 
%    \begin{macrocode}
\newcommand{\ADS@website}{http://www.astrosmo.unam.mx/rmaa/}
\newcommand{\ADS@url}
{\ADS@website\ADS@journal\ADS@volume\ADS@number/PDF/\ADS@journal\ADS@volume\ADS@number\_\articlename.pdf}
\newcommand{\ADS@number}{}
\newcommand{\ADS@fulldate}{00/\ADS@year} 
\newcommand{\ADS@lastpage}{00} 
\newcommand{\ADS@copyright}{(C) \ADS@year; Instituto de Astronom\'\i{}a,
  Universidad Nacional Aut\'onoma de M\'exico} 
\newcommand{\SetNumber}[1]{\renewcommand{\ADS@number}{-#1}}
\newcommand{\SetWebsite}[1]{\renewcommand{\ADS@website}{#1}}
\newcommand{\SetMonthNumeric}[1]{\renewcommand{\ADS@fulldate}{#1/\ADS@year}}
%    \end{macrocode}
% \end{macro}\end{macro}\end{macro}\end{macro}\end{macro}\end{macro}\end{macro}
% \begin{macro}{\ADS@affil}
% \begin{macro}{\ADS@references}
% These two are somewhat more problematical. Affiliation currently
% only works for the first author (see \dots) pending redesign of the
% user interface. References are now done by wrapping the
% |thebibliography| environment to open another input stream on the
% same file and search through for the |\bibitems|. 
%    \begin{macrocode}
\newcommand{\ADS@references}{}
\newcommand{\ADS@affil}{}
%    \end{macrocode}
% \end{macro}\end{macro}
% 
% Define two macros that expand to the strings
% |\begin{thebibliography}| and |\end{thebibliography}|. Note that the
% idea for this and much else is inspired by the code for the docstrip
% package. Also define |\StartBibX| in case authors use the more
% standard |\begin{thebibliography}{}| syntax. 
%    \begin{macrocode}
\edef\EndBib{\expandafter\@gobble\string\\end\string{thebibliography\string}}
\edef\StartBib{\expandafter\@gobble\string\\begin\string{thebibliography\string}}
\edef\StartBibX{\expandafter\@gobble\string\\begin\string{thebibliography\string}\string{\string}}
%    \end{macrocode}
% Define two new logical variables. |\ifAnotherline| is true if we
% want to read another line from the input stream. It is set to false
% when we encounter the |\end{thebibliography}|. |\ifInBib| is set to
% true when we frind the |\begin{thebibliography}| and tells us to
% start saving the following lines. 
%    \begin{macrocode}
\newif\ifAnotherLine
\newif\ifInBib
%    \end{macrocode}
% \begin{macro}{makeOther}
% We need a macro to change the catcodes of lots of things to other. 
%    \begin{macrocode}
\def\makeOther#1{\catcode`#1=12\relax}
%    \end{macrocode}
% \end{macro}
% \begin{macro}{\RereadBiblio}
% Define a macro to do all the fiddly stuff. This is used at the
% beginning of |thebibliography| environment (Sec~\ref{sec:biblio}).
%    \begin{macrocode}
\newread\tf@biblio
\newcommand{\RereadBiblio}{%
%    \end{macrocode}
% Open the current file again for reading
%    \begin{macrocode}
  \begingroup
  \immediate\openin\tf@biblio=\articlename
%    \end{macrocode}
% Everything in a group to be safe. Initialize logicals. Black magic
% to change
% catcodes and avoid blank lines being interpreted as |\par|. 
%    \begin{macrocode}
  \AnotherLinetrue\InBibfalse
  \makeOther\ \makeOther\\\makeOther\$%
  \makeOther\#\makeOther\^\makeOther\^^K%
  \makeOther\_\makeOther\^^A\makeOther\%%
  \makeOther\~\makeOther\{\makeOther\}\makeOther\&%
  \endlinechar-1\relax
%    \end{macrocode}
% Start looping through the file. Read each line to the macro
% |\Line|. If we have got to EOF or end of biblio then set
% |\AnotherLinefalse|.
%    \begin{macrocode}
  \loop
  \read\tf@biblio to \Line
  \ifx\Line\EndBib\AnotherLinefalse\InBibfalse\fi
  \ifeof\tf@biblio \AnotherLinefalse\fi
%    \end{macrocode}
% If we are in the biblio and the line is not blank then append it to
% the |\ADS@references| macro together with a newline. 
%    \begin{macrocode}
  \ifInBib
  \ifx\Line\empty\else\xdef\ADS@references{\ADS@references\Line^^J}\fi
  \fi
%    \end{macrocode}
% Check if the line we just saw was the start of biblio, then repeat
% the loop if necessary. 
%    \begin{macrocode}
  \ifx\Line\StartBib\InBibtrue\fi
  \ifx\Line\StartBibX\InBibtrue\fi
  \ifAnotherLine\repeat
  \immediate\closein\tf@biblio
  \endgroup
  }
%    \end{macrocode}
% \end{macro}
% 
% \begin{macro}{\PadToFourWithDots}
% Utility function for left-padding a number to 4 characters with
% periods. Used in 
% constructing |\ADS@BibCode|. This is a very simple algorithm and
% \emph{requires} that the argument actually be a number---if it is
% not then dire consequences will result.  
%    \begin{macrocode}
\newcommand{\PadToFourWithDots}[1]{%
  \ifnum #1 > 999 #1%
  \else\ifnum #1 > 99 .#1%
  \else\ifnum #1 > 9 ..#1%
  \else ...#1%
  \fi\fi\fi}
%%  \ifnum \count0 > 999 #1%
%%  \else\ifnum \count0 > 99 .#1%
%%  \else\ifnum \count0 > 9 ..#1%
%%  \else ...#1%
%%  \fi\fi\fi}
%    \end{macrocode}
% \end{macro}
% 
% \begin{macro}{\ADS@dumpdata}
% \begin{macro}{\ADS@write}
% \begin{macro}{\ADS@openfile}
% |\ADS@dumpdata| is the macro that actually does the job of writing out all the
% data to an auxiliary file, using |\ADS@write| to write each
% field. The file is set up with |\ADS@openfile|, which currently is
% done in the TOC of a book. 
%    \begin{macrocode}
\newcommand{\ADSfield}[2]{\relax}
\newcommand{\ADSvfield}[2]{\relax}
%%\renewcommand{\ADSfield}[12]{\%#1 #2}
\newcommand{\ADS@dumpdata}
{ 
  \ADS@write{\string\begin{ADSarticle}}
  \ADS@WRITE{\string\ADSfield{R}{\ADS@BibCode}}
  \ADS@write{\string\ADSfield{A}{\ADS@authors}}
  {\ADS@sanitize\ADS@write{\string\ADSfield{F}{AA(\ADS@affil)}}}
  \ADS@write{\string\ADSfield{J}{\ADS@journalname}}
  \ADS@write{\string\ADSfield{D}{\ADS@fulldate}}
  \ADS@write{\string\ADSfield{L}{\ADS@lastpage}}
  {\ADS@sanitize\ADS@write{\string\ADSfield{T}{\ADS@title}}}
  \ADS@write{\string\ADSfield{C}{\ADS@copyright}}
  \ADS@write{\string\ADSfield{U}{\ADS@url}}
  \ADS@write{\string\ADSfield{K}{\ADS@keywords}}
  {\ADS@sanitize\ADS@write{\string\ADSfield{B}{\ADS@abstract}}}
  \ADS@write{\string\ADSfield{Z}{\ADS@references}}
  \ADS@write{\string\end{ADSarticle}}
  }
%%\newcommand{\ADS@write}[1]{\if@book\addtocontents{ads}{#1}\fi}
%%\newcommand{\ADS@openfile}{\@starttoc{ads}}
\newcommand{\ADS@write}[1]{\set@display@protect\immediate\write\tf@ads{#1}}
\newcommand{\ADS@WRITE}[1]{\immediate\write\tf@ads{#1}}
\newcommand{\ADS@openfile}{%
  \newwrite\tf@ads%
  \immediate\openout\tf@ads\jobname.ads\relax
  \ADS@WRITE{\ADS@preamble}}
\newcommand{\ADS@closefile}{\ADS@WRITE{\ADS@postamble}}
\def\ADS@preamble{%
  %%  \string\documentclass{article}
  %%  \string\begin{document}
  }
\def\ADS@postamble{%
  %%  \string\end{document}
  \string%\string% Local Variables:^J 
  \string%\string% mode: LaTeX^J
  \string%\string% mode: Font-lock^J
  \string%\string% End:^J
  }
%    \end{macrocode}
% \end{macro}
% \end{macro}
% \end{macro}
% 
% \subsection{Automatic calculation of the page range}
% \label{sec:page-range}
% Another issue that has cost me a lot of grief. Still, it finally
% seems to be sorted now. All hail Jeff Goldberg. 
%
% Cunning trick to get the page numbers from |\pageref|. 
% Put a label on the first and last page of the document using
% |\AtBeginDocument| and |\AtEndDocument|. . 
% This way won't work for the book option. 
% In that case, we do it by defining |\@endofarticlehook| and
% |\@begofarticlehook| as parallels
% to |\AtEndDocument| and |\AtBeginDocument|. 
%    \begin{macrocode}
\if@book\AtEndDocument{\ADS@closefile}\else
\AtBeginDocument{\@begofarticlehook}
\AtEndDocument{\@endofarticlehook}
\fi
%    \end{macrocode}
% \begin{macro}{\@begofarticlehook}
% Since we are getting more and more things that need to be done at
% the beginning of every article, frontsection, etc we define a sister
% hook to 
% |\@endofarticlehook|. This hook does the following: 
% \begin{enumerate}
% \item Labels the first page
% \item Sets |\ADS@firstpage|
% \item runs and then clears |\indexauthor@deferlist|
% \item resets the keyword counter
% \end{enumerate}
% It is executed in the following circumstances: 
% \begin{enumerate}
% \item At |\begin{document}| for standalone files
% \item At |\begin{document}| for each included file
% \item At beginning of front sections 
% \item At beginning of TOC
% \item At beginning of Author Index
% \item At beginning of abstracts
% \end{enumerate}
% For some reason, this is not always within the right group, so we
% use |\xdef| just to be sure. 
%    \begin{macrocode}
\newcommand{\@begofarticlehook}{%
  \if@frontmatter\frontlabel{RMAAFirstPage}\else\label{RMAAFirstPage}\fi%
  \xdef\ADS@firstpage{\@arabic{\c@page}}%
  \if@book\else\makeindex\fi % needed so default index is defined in stand-alone articles
%%  \indexauthor@deferlist
%%  \xdef\indexauthor@deferlist{}
  \@inpreamblefalse
  \setcounter{kwdcounter}{0}
  }
%    \end{macrocode}
% \end{macro}
% \begin{macro}{\@endofarticlehook}
% This is a hook to be run at the end of each article
% currently, it just labels the last page and writes out all the ADS
% data to an auxiliary file.  
% It should not produce any printed output. 
% It is modelled on the |\AtEndDocument| command, to which it reduces
% in the case of a single article. 
%    \begin{macrocode}
\newcommand{\@endofarticlehook}{%
  \if@debug{\ClassWarning{rmaa}{***Doing Hook***}}\fi
  \if@debug{\ClassWarning{rmaa}{ADSabstract is: \ADS@abstract}}\fi
  \clearpage\rmaa@lastpagelabel\if@book\ADS@dumpdata\fi
  \ifx\NotesList\@rmaaEmpty\else\OutputNotes\fi
  \if@debug{\ClassWarning{rmaa}{Notes list is: \NotesList}}\fi
}
%    \end{macrocode}
% \begin{macro}{\rmaa@lastpagelabel}
% New version of last page label command copied from Jeff Goldberg's
% lastpage.sty package - we can't just use |\label| here in case the
% output routine for the last page has already been called. Instead,
% we first explicitly clear the page, step back the page counter,
% explicitly write to the .aux file, then restep the page counter
% again. 
% 
% Also set the |\ADS@LastPage| macro to the appropriate number. 
%    \begin{macrocode}
\newcommand{\rmaa@lastpagelabel}{%
  \addtocounter{page}{-1}%
  \immediate\write\@auxout{\string
    \newlabel{\articlename:RMAALastPage}{{}{\thepage}}}%
  \edef\ADS@lastpage{\thepage}%
  \addtocounter{page}{1}}
%    \end{macrocode}
% \end{macro}
% \end{macro}

% \subsection{Deal with explicit change in number of columns}
% \label{sec:columns}
% Sometimes, when formatting a book, we need to switch between
% one-column articles and two-column articles. For this to work, we
% need to wrap the standard \LaTeX{} commands to set the |\@rmtwocol|
% flag appropriately and also reset the package default values of all
% scalable lengths. .  
%    \begin{macrocode}
\let\latexonecolumn=\onecolumn
\let\latextwocolumn=\twocolumn
\renewcommand{\onecolumn}{\@rmtwocolfalse\@package@default@lengths\latexonecolumn}
\renewcommand{\twocolumn}{\@rmtwocoltrue\@package@default@lengths\latextwocolumn}
%    \end{macrocode}
% 
% \subsection{The actual formatting of the title}
% \label{sec:title-format}

% Note, we have to redefine the |\maketitle| command to use the
% numbered footnotes. This means we needn't use the atbeginend stuff. 
%    \begin{macrocode}
\renewcommand{\maketitle}{\par
  \begingroup
%    \end{macrocode}
% Don't switch to marked footnotes. 
% This is consistent with existing style but I'm not sure I like it
%    \begin{macrocode}
  %%  \renewcommand\thefootnote{\fnsymbol{footnote}}%
%%  \def\@makefnmark{\hbox to\z@{$\m@th^{\@thefnmark}$\hss}}%
  \def\@makefnmark{\textsuperscript{\@thefnmark}\hspace{-0.2em}}%
  \long\def\@makefntext##1{\parindent 1em\noindent
    \hbox to1.8em{\hss$\m@th^{\@thefnmark}$}##1}%
  \if@rmtwocol
  { \if@debug{\ClassWarning{rmaa}{TWO COLUMN}}\fi
    \twocolumn[\@maketitle]%
    }
  \else
  { \if@debug{\ClassWarning{rmaa}{ONE COLUMN}}\fi
    \newpage
    \global\@topnum\z@   % Prevents figures from going at top of page.
    \if@manuscript \@MSmaketitle \else \@maketitle \fi
    }
  \fi
  \if@manuscript\else\thispagestyle{rmaafirstpage}\@thanks\fi
  \ifx\NotesList\@rmaaEmpty\else{%
    \NoteMark{1--\arabic{notecounter}}\vspace*{-2\baselineskip}%
  }\fi%
  \endgroup
%    \end{macrocode}
% Don't reset footnote counter (not sure whether this is good thing
% or not).
%    \begin{macrocode}
  %%  \setcounter{footnote}{0}%
  \let\thanks\relax\let\affil\relax\let\altaffil\relax
  \let\maketitle\relax\let\@maketitle\relax
  \gdef\@thanks{}\gdef\@author{}\gdef\@title{}}
%    \end{macrocode}

% \begin{macro}{\@kwhwidth}
% \begin{macro}{\@kwtwidth}
% Redefine \cs{@maketitle} to include the abstract, resumen and
% keywords. Note that in the manuscripty style we use
% \cs{@MSmaketitle} instead.  
% \changes{3.16}{2003/04/07}{
%  Added extra vertical space between authors and accepted date.}
% \changes{3.26}{2011/08/15}{Add extra label to review articles}
%    \begin{macrocode}
\newlength{\@kwhwidth}\newlength{\@kwtwidth}
\newcommand\reviewname{Review}
\renewcommand{\@maketitle}{%
  \newpage
  \null
  \if@reviewarticle
  \smallskip
  {\large\bfseries\itshape\reviewname}%
  \bigskip
  \else
  \vskip\@before@title@skip%
  \fi
  \begin{center}%
    {\@title \par}%
    \vskip\@before@author@skip%
    \parbox{\rmaa@AuthorWidthFrac\textwidth}{\centering
      \lineskip\@between@author@skip%
      \@author
      \par}
    %% this should be replaced with `Submitted to...' for preprints
    \if@printdates
    \vskip\@before@dates@skip%
    \if@preprint
    \rmaa@ReceivedStyle{Submitted to RMxAA \@received}
    \else
    \rmaa@ReceivedStyle{Received \@received; accepted \@accepted}
    \fi
    \fi
    \if@manuscript
    \rmaa@ReceivedStyle{Submitted to RMxAA \@received}
    \fi
    \if@guide
    \rmaa@ReceivedStyle{Version \rmaa@version, \rmaa@date}
    \fi
    \vskip\@before@abstract@skip%
    \if@abstractpresent
    \parbox{\rmaa@AbstractWidthFrac\textwidth}{%
      \setlength{\parindent}{0.9cm}%
      \centerline{\rmaa@AbstractHeaderStyle{Resumen}}\par%
      \vskip\@after@abstracthead@skip%
      \if@newstyle\noindent\fi\@resumen
      \if@setresumen\else\@abstract\fi % append english abstract if
                                       % spanish one not set
      \vskip\@between@abstract@skip%
      \centerline{\rmaa@AbstractHeaderStyle{Abstract}}\par%
      \vskip\@after@abstracthead@skip%
      \if@newstyle\noindent\fi\@abstract
      \if@guide\else            %do not use keywords for guide
      \vskip\@before@kwd@skip%
      % Format the Key Words with hanging indentation:
      % They are put in a \parbox with width equal to the line width
      % minus the width of the header. 
      { % Widths of keyword
        \settowidth{\@kwhwidth}{\@keywordheader~}   % header and
        \setlength{\@kwtwidth}{\linewidth}          % keyword text
        \addtolength{\@kwtwidth}{-\@kwhwidth}       %
        \noindent\@keywordheader~\parbox[t]{\@kwtwidth}{%
          \rmaa@KeywordStyle{\boldmath\uppercase{\@keywords}}}}%
      \fi
      }
      \fi
  \end{center}%
  \vspace*{\afterkwdskip}
  \par
  }
\newlength{\afterkwdskip}\setlength{\afterkwdskip}{\z@}
%    \end{macrocode}
% \end{macro}
% \end{macro}
%
%
% \begin{macro}{\reviewarticle}
% \begin{macro}{\reviewtoc}
% \begin{macro}{\reviewlof}
%   User command to specify a review article. Must be used in
%   preamble. Main effect is to place the text ``Review'' before the
%   article title.  Also add command for setting a table of contents
%   and a list of figures at the beginning of the article.
%    \begin{macrocode}
\newcommand{\reviewarticle}{\@reviewarticletrue}
\newcommand\reviewtoc{%
  \section*{CONTENTS}%
  \@starttoc{rtc}%
}
\newcommand\reviewlof{%
  \section*{LIST OF FIGURES}%
  \@starttoc{lof}%
}
%    \end{macrocode}
% \end{macro}
% \end{macro}
% \end{macro}
% \begin{macro}{\l@section}
%    \begin{macrocode}
\renewcommand*\l@section[2]{%
  \ifnum \c@tocdepth >\z@
    \addpenalty\@secpenalty
    \addvspace{0.8em \@plus\p@}%
    \setlength\@tempdima{1.5em}%
    \begingroup
      \parindent \z@ \rightskip \@pnumwidth
      \parfillskip -\@pnumwidth
      \leavevmode %\bfseries
      \advance\leftskip\@tempdima
      \hskip -\leftskip
      #1\nobreak\hfil \nobreak\hb@xt@\@pnumwidth{\hss #2}\par
    \endgroup
  \fi}
%    \end{macrocode}
% \end{macro}
%
% \begin{macro}{\@MSmaketitle}
% A special version of |\@maketitle| for the \texttt{manuscript}
% option. This can deal with a multi-page abstract/resumen. Not very
% well, mind. 
%    \begin{macrocode}
\newcommand{\@MSmaketitle}{%
  \if@debug\ClassInfo{rmaa}{Using special MS maketitle command}\fi
  \newpage
  {
    \thispagestyle{rmaafirstpage}
%     \setlength{\parindent}{0pt}
    \null\vskip 2ex
    \begin{center}
      \@title\par
      \@author\par
      \rmaa@ReceivedStyle{%
        \ifx\@received\@rmaaEmpty
          Draft version: \@date
        \else
          Received: \@received
          \ifx\@accepted\@rmaaEmpty
          \else
            ; Accepted: \@accepted
          \fi
        \fi}\par
    \end{center}
    \@thanks
    \centerline{\rmaa@AbstractHeaderStyle{Resumen}}
    \noindent\@resumen\par\vskip 2ex
    \centerline{\rmaa@AbstractHeaderStyle{Abstract}}
    \noindent\@abstract\par\vskip 2ex
    \noindent\@keywordheader~\rmaa@KeywordStyle{\@keywords}\par
%     \@keywordheader~\rmaa@KeywordStyle{\uppercase{\@keywords}}\par
  }%
  }
%    \end{macrocode}
% \end{macro}

% \subsection{The sectioning commands}
% \label{sec:section-commands}

% \begin{macro}{\section}
% \begin{macro}{\subsection}
% \begin{macro}{\subsubsection}
% \begin{macro}{\paragraph}
% \begin{macro}{\subparagraph}
% Redefine all the standard commands to use these lengths. 
%    \begin{macrocode}
\renewcommand\section{\@startsection {section}{1}{\z@}%
                                   {\abovesectionskip}%
                                   {\belowsectionskip}%
                                   %% how do we get upper case here
                                   {\centering\reset@font\normalsize}}
\renewcommand\subsection{\@startsection{subsection}{2}{\z@}%
                                     {\abovesectionskip}%
                                     {\belowsectionskip}%
                                     {\centering\reset@font\normalsize\itshape}}
\renewcommand\subsubsection{\@startsection{subsubsection}{3}{\z@}%
                                     {\abovesectionskip}%
                                     {\belowsectionskip}%
                                     {\centering\reset@font\normalsize\itshape}}
\renewcommand\paragraph{\@startsection{paragraph}{4}{\z@}%
                                    {\aboveparaskip}%
                                    {\belowparaskip}%
                                    {\reset@font\normalsize\bfseries}}
\renewcommand\subparagraph{\@startsection{subparagraph}{5}{\parindent}%
                                      {\aboveparaskip}%
                                      {\belowparaskip}%
                                      {\reset@font\normalsize\bfseries}}
%    \end{macrocode}
% \end{macro}
% \end{macro}
% \end{macro}
% \end{macro}
% \end{macro}

% \begin{macro}{secnumdepth}
% Make sure that we number the subsubsections---this is done by
% default in the |article| class but not the |book| class, which is
% why we must set it explicitly here.
%    \begin{macrocode}
\setcounter{secnumdepth}{3}
%    \end{macrocode}
% \end{macro}
% \changes{2.2beta10}{2001/11/07}{Fixed \cs{secnumdepth} for \cs{book}
% style.}

% \begin{macro}{\thesubsection}
% \begin{macro}{\thesubsubsection}
% For subsection and subsubsection, we want the numbers to be roman
% but the title to be italic. This seems the easiest way to go about it.
%    \begin{macrocode}
\renewcommand{\thesubsection}{{\upshape\thesection.\@arabic\c@subsection}}
\renewcommand{\thesubsubsection}{\thesubsection{\upshape.\@arabic\c@subsubsection}}
%    \end{macrocode}
% \end{macro}
% \end{macro}

% \begin{macro}{\@seccntformat}
% Add ``.'' to the number here so we can remove it from |\thesection|,
% etc. Reduce gap between number and title WJH 27 Jan 1995
%    \begin{macrocode}
\def\@seccntformat#1{\csname the#1\endcsname .\hskip 0.4em}
%    \end{macrocode}
% \end{macro}

% \begin{macro}{\@sect}
% I didn't really want to mess with the |\@sect| command, but it was
% the only way I could manage to force upper case. Other than the
% commented parts, it is identical to the standard version. 
%    \begin{macrocode}
\def\@sect#1#2#3#4#5#6[#7]#8{
  \ifnum #2>\c@secnumdepth
     \let\@svsec\@empty
  \else
     \refstepcounter{#1}%
     \let\@@protect\protect
     \def\protect{\noexpand\protect\noexpand}%
     \edef\@svsec{\@seccntformat{#1}}%
     \let\protect\@@protect
  \fi
%    \end{macrocode}
% This is the bit I changed (27 Jan 1995). Use upper case for
% |\section|, use mixed case for the rest. 
%    \begin{macrocode}
  \ifnum #2=1                   
     \def\sec@upcase##1{\uppercase{##1}}
     \def\sec@shape{\upshape}
  \else                         
     \def\sec@upcase##1{\relax{##1}}
     \def\sec@shape{\itshape}
  \fi
  \@tempskipa #5\relax
  \ifdim \@tempskipa>\z@
     \begingroup #6\relax
          \@hangfrom{\hskip #3\relax\@svsec}%
                    {\interlinepenalty \@M \sec@upcase{#8}\par}%
     \endgroup
     \csname #1mark\endcsname{#7}
%    \end{macrocode}
% Only write TOC entries if doing the Guide style. 
%    \begin{macrocode}
     \if@guide 
        \addcontentsline{toc}{#1}{%
          \ifnum #2>\c@secnumdepth \else
             \protect\numberline{\csname the#1\endcsname}
          \fi
          #7}%
     \fi
%    \end{macrocode}
% Also allow per-article TOC for review papers
%    \begin{macrocode}
     \if@reviewarticle
        \ifnum #2>2 \else
           \addcontentsline{rtc}{#1}{%
             \protect\numberline{\csname the#1\endcsname}
             \sec@upcase{\sec@shape #7}}%
        \fi
     \fi
  \else
     \def\@svsechd{#6\hskip #3\relax  %% \relax added 2 May 90
       \@svsec #8\csname #1mark\endcsname{#7}
       \if@guide
          \addcontentsline{toc}{#1}{%
            \ifnum #2>\c@secnumdepth \else
               \protect\numberline{\csname the#1\endcsname}%
            \fi
            #7}%
       \fi
       }%
  \fi
  \@xsect{#5}}
%    \end{macrocode}
% \end{macro}


% \subsection{Appendices}
% \label{sec:appendices}
% Make sure that appendices in individual articles work properly. We
% copy the definition from the article class (the book one doesn't do
% what we want), then we make the equation numbering be subordinate to
% the appendix (equations A12, B6, etc). Finally, we have a medskip 
% followed by APPENDICES to introduce them (but what if there is only
% one!). 
% \changes{3.17}{2002/04/08}{ 
%  Added appendices environment that does the right thing.}
% 
%    \begin{macrocode}
\newenvironment{appendices}{\par
  \begingroup
  \setcounter{section}{0}%
  \setcounter{subsection}{0}%
  \renewcommand\thesection{\@Alph\c@section}
  \renewcommand\theequation{\thesection\arabic{equation}}
  \medskip
  \section*{APPENDICES}}
{\endgroup}
%    \end{macrocode}

% \subsection{Floating thingies}
% \label{sec:floats}   

%% 
%    \begin{macrocode}
\renewcommand{\thetable}{\@arabic\c@table}
\renewcommand{\thefigure}{\@arabic\c@figure}
\renewcommand{\figurename}{Fig.}
\renewcommand{\tablename}{TABLE}
%\renewcommand{\fnum@table}{\MakeUpperCase{\tablename}~\thetable}
%    \end{macrocode}
% For some reason, I have to write `TABLE' explicitly here
% If I use |\tablename|, it gets converted to mixed case---God knows
% why. 
% |\figurename| is also being ignored. 
%    \begin{macrocode}
\renewcommand{\fnum@table}{TABLE~\thetable}
\renewcommand{\fnum@figure}{Fig.~\thefigure}
%    \end{macrocode}

% Do captions in 9 pt face in the main journal style and in new-style
% anything. 
%    \begin{macrocode}
\long\def\@makefigcap#1#2{%
  \if@proceedings\else\small\fi
  \if@newstyle\small\fi 
  \vskip\abovecaptionskip
  \sbox\@tempboxa{#1. #2}%
  \ifdim \wd\@tempboxa >\hsize
    #1. #2 \par
  \else
    \global \@minipagefalse
    \hb@xt@\hsize{\hfil\box\@tempboxa\hfil}%
  \fi
  \vskip\belowcaptionskip}
%    \end{macrocode}

% Make sure table caption is upper case. 
%    \begin{macrocode}
\long\def\@maketabcap#1#2{%
  \hbox to\hsize{\hfil #1 \hfil}
  \vskip 6pt 
  \centering\MakeUpperCase{#2}\par
  \vskip 6pt}
%    \end{macrocode}

%    \begin{macrocode}
\newlength{\@rottabcapwidth}
\long\def\@makerottabcap#1#2{%
  \setbox\@tempboxa\hbox{\MakeUpperCase{#2}}%
  \ifdim\wd\@tempboxa >.8\vsize
     \setlength\@rottabcapwidth{.8\vsize}
  \else
     \setlength\@rottabcapwidth{\wd\@tempboxa}
  \fi
  \rotatebox{90}{%
    \begin{minipage}{\vsize}\centering
      \begin{minipage}{\@rottabcapwidth}\centering
        #1\\ \vskip 6pt \MakeUpperCase{#2}
      \end{minipage}%
    \end{minipage}
    }%
  \hskip 6pt
  }

\newenvironment{Sideways}
    {\begin{sideways}\begin{minipage}{\vsize}\centering}
    {\end{minipage}\end{sideways}}

\newenvironment{Sidewaystable}
{%
  \let\oldcaption=\caption%
  \def\caption{
    \def\@captype{table}\renewcommand{\@makecaption}{\@maketabcap}
    \oldcaption}
  \begin{sidewaystable}\centering}
  {\end{sidewaystable}\let\caption=\oldcaption}
\newenvironment{Sidewaystable*}
{%
  \let\oldcaption=\caption%
  \def\caption{
    \def\@captype{table}\renewcommand{\@makecaption}{\@maketabcap}
    \oldcaption}
  \begin{sidewaystable*}\centering}
  {\end{sidewaystable*}\let\caption=\oldcaption}
%    \end{macrocode}

% Not sure why we need to stick these lines in. 
%    \begin{macrocode}
\renewenvironment{figure}
                 {\renewcommand{\@makecaption}{\@makefigcap}\@float{figure}}
                 {\end@float}
\renewenvironment{figure*}
                 {\renewcommand{\@makecaption}{\@makefigcap}\@dblfloat{figure}}
                 {\end@dblfloat}
\renewenvironment{table}
                 {\renewcommand{\@makecaption}{\@maketabcap}\@float{table}}
                 {\end@float}
\renewenvironment{table*}
                 {\renewcommand{\@makecaption}{\@maketabcap}\@dblfloat{table}}
                 {\end@dblfloat}
%    \end{macrocode}

% Redefine |\@caption| so that we don't bother with the list of
% figures. This saves us having to worry about fragile commands in the
% captions. 
%    \begin{macrocode}
\long\def\@caption#1[#2]#3{%
  \par
%%  \addcontentsline{\csname ext@#1\endcsname}{#1}%
%%    {\protect\numberline{\csname the#1\endcsname}{\ignorespaces #2}}%
  \begingroup
    \@parboxrestore
    \if@minipage
      \@setminipage
    \fi
    \normalsize
    \@makecaption{\csname fnum@#1\endcsname}{\ignorespaces #3}\par
  \endgroup}
%    \end{macrocode}

% Make tables not quite so cramped. Allow |\hline\hline| to give a
% poor man's thick rule. 
%    \begin{macrocode}
\renewcommand{\arraystretch}{1.15}
\setlength{\doublerulesep}{0.4pt}
%    \end{macrocode}

% \begin{macro}{\tablecols}
% \begin{macro}{\@tablecols}
% \begin{macro}{\tabcaptioncontinued}
% |\tablecols| is a user command for setting the number of columns. It
% is used by commands that try and span the entire table (I think
% there must be a better way of acheiving
% this). |\tabcaptioncontinued| is 
% a user command for setting the caption of a continued table on the
% next page. 
%    \begin{macrocode}
\newcommand{\tablecols}[1]{\renewcommand{\@tablecols}{#1}}
\newcommand{\@tablecols}{1}
\newcommand{\tabcaptioncontinued}{%
  \multicolumn{\@tablecols}{c}{\addtocounter{table}{-1}TABLE~\thetable\ (CONTINUED)} \\ \\
  }
%    \end{macrocode}
% \end{macro}
% \end{macro}
% \end{macro}

% \subsubsection{Footnotes in tables}
% \label{sec:table-footnotes}
% \begin{macro}{\tablenotemark}
% \begin{macro}{\tablenote}
% \begin{macro}{\tablenoteskip}
% Old user commands |\tablenotemark{MARK}|,
% |\tablenote{MARK}{TEXT}|---now deprecated. 
%    \begin{macrocode}
\newcommand\tablenotemark[1]{$^\mathrm{#1}$} 
\newcommand\tablenote[2]{\if@proceedings\else\small\fi ${}^\mathrm{#1}$ #2} 
\newcommand\tablenoteskip{\\[0.5\baselineskip]} 
%    \end{macrocode}
% \end{macro}
% \end{macro}
% \end{macro}

% \begin{macro}{\tabnotemark}
% \begin{macro}{\tabnotetext}
% \begin{macro}{\tabnotetextTWO}
% Alternative way of doing footnotes in tables. 
%    \begin{macrocode}
\newcommand{\tabnotemark}[1]{\textsuperscript{\lowercase{#1}}}
%\newcommand{\tabnotetextTWO}[4]{\tabnotetext{#1}{#2}\tabnotetext{#3}{#4}}
\newcommand{\tabnotetextTWO}[4]{%
  \multicolumn{\@tablecols}{l}{%
    \parbox{0.5\tabnotewidth}{\small\textsuperscript{#1}#2}%
    \parbox{0.5\tabnotewidth}{\small\textsuperscript{#3}#4}%
    }%
  \\}
\newcommand{\tabnotetext}[2]{%
  \multicolumn{\@tablecols}{@{}p{\tabnotewidth}@{}}{\small\textsuperscript{#1}#2}\\}
\newlength{\tabnotewidth}
%    \end{macrocode}
% \end{macro}
% \end{macro}
% \end{macro}


% \begin{macro}{changemargin}
%    \begin{macrocode}
% Command to allow margins to be changed on the fly. Useful to allow
% tables and Figures to protrude into the margins. 
%
% \changes{3.24}{2007/09/14}{Added \texttt{changemargin} environment
% (copied from TeX FAQ).}
\newenvironment{changemargin}[2]{% 
  % From TeX FAQ http://www.tex.ac.uk/cgi-bin/texfaq2html?label=chngmargonfly
  \begin{list}{}{%
    \setlength{\topsep}{0pt}%
    \setlength{\leftmargin}{#1}%
    \setlength{\rightmargin}{#2}%
    \setlength{\listparindent}{\parindent}%
    \setlength{\itemindent}{\parindent}%
    \setlength{\parsep}{\parskip}%
  }%
  \item[]}{\end{list}}
%    \end{macrocode}
% \end{macro}

% \subsection{The bibliography and references}
% \label{sec:biblio}
% \begin{macro}{\beforebiblioskip}
% \begin{macro}{\afterbiblioskip}
% I'm not entirely sure why I introduced these two commands, since
% directly putting |\vspace| before and after
% |\begin{thebibliography}| seems to work well enough. 
%    \begin{macrocode}
\newlength{\beforebiblioskip}
\setlength{\beforebiblioskip}{0pt}
\newlength{\afterbiblioskip}
\setlength{\afterbiblioskip}{0pt}
%    \end{macrocode}
% \end{macro}
% \end{macro}
% \begin{macro}{thebibliography}
% The environment itself is done as starred section with extra
% vertical space after heading. This space is currently zero. 
% In the book style we do the cunning trick with rereading the file in
% order to get the references into the |\ADS@references| macro.
% 
% \changes{3.05}{2003/01/02}{Removed the \cs{beforebiblioskip} and
% \cs{afterbiblioskip} commands, since they cause unwanted vertical
% space in some instances.}  
%    \begin{macrocode}
\renewenvironment{thebibliography}{%
  \if@book\RereadBiblio\fi
%%  \vspace*{\beforebiblioskip}
  \section*{REFERENCES} \frenchspacing
%%  \vspace*{\afterbiblioskip}
  \small
%    \end{macrocode} 
% The references themselves are set as a list in order 
% to get the hanging indentation. They are done in a smaller font and
% with french spacing so we needn't worry about the distinction
% between sentence-ending periods and abbreviations. 
%    \begin{macrocode} 
  \list{\null}{\leftmargin 1.5em\labelwidth\z@\labelsep\z@\itemindent -1.5em
    \itemsep\z@\parsep\z@\usecounter{enumi}}
  \sloppy\clubpenalty4000\widowpenalty4000
  \sfcode`\.=1000\relax\@biblisttrue}
{\endlist\@biblistfalse\normalsize\@outputaddressesifnecessary}
%    \end{macrocode}
% \end{macro}
% \subsection{Uniquification of cross-references}
% \label{sec:uniquification}
% 
% \begin{macro}{\articlename}
% In order that the |\label{}| command in individual articles be
% guaranteed to give a unique anchor, we append the article name (held
% in the |\articlename| command) to each key. This requires
% redefinition of the following commands: |\cite|, |\nocite|, |\ref|,
% |\pageref|, |\label|. 
% \begin{macro}{\oldcite}
% \begin{macro}{\oldnocite}
% \begin{macro}{\oldref}
% \begin{macro}{\oldpageref}
% \begin{macro}{\oldlabel}
% The original versions of these commands are
% saved as |\oldcite|, |\oldnocite|, etc. 
%
% The uniquifying of cross-references is now always done, so 
% we don't get multiply-defined labels in the single article case
%    \begin{macrocode}
\newcommand{\articlename}{anon}
\let\oldcite=\cite 
\let\oldnocite=\nocite 
\let\oldref=\ref 
\let\oldpageref=\pageref
\let\oldlabel=\label
%    \end{macrocode}
% \end{macro}
% \end{macro}
% \end{macro}
% \end{macro}
% \end{macro}
% \end{macro}
% Uniquify all cross-references by prepending the article name
%    \begin{macrocode}
\renewcommand{\ref}[1]{\oldref{\articlename:#1}%
  \if@debug{\ClassInfo{rmaa}{Reference to \articlename:#1}}\fi}
\renewcommand{\pageref}[1]{\oldpageref{\articlename:#1}%
  \if@debug{\ClassInfo{rmaa}{Page reference to \articlename:#1}}\fi}
\renewcommand{\label}[1]{\oldlabel{\articlename:#1}%
  \if@debug{\ClassInfo{rmaa}{Label for \articlename:#1}}\fi}
%    \end{macrocode}
% Citations are more complicated since they can contain a list of
% keys. These are really only here for use with the rmaacite package,
% which is utterly deprecated (02 Jan 2003) now that I have the natbib
% code working. The command \cs{cite} should \emph{not} be used with
% natbib: use \cs{citep}, \cs{citet}, etc.\@ instead. 
%    \begin{macrocode}
\renewcommand{\nocite}[1]{\@bsphack
  \@for\@citeb:=#1\do{%
    \edef\@citeb{\expandafter\@firstofone\@citeb:\articlename}%
    \if@filesw\immediate\write\@auxout{\string\citation{\@citeb}}\fi
    \@ifundefined{b@\@citeb:\articlename}{\G@refundefinedtrue
      \@latex@warning{Citation `\@citeb' undefined}}{}}%
  \@esphack}
\renewcommand{\cite}[1]{\oldcite{#1:\articlename}}
%   \end{macrocode}


% Horrible fiddling with the bibliography citation mechanism. 
% This has to be globally available in case rmaacite puts it in the
% \texttt{.aux} file. 
%    \begin{macrocode}
\def\lbibcite#1#2#3{\global\@namedef{b@#1}{#2}
                    \global\@namedef{by@#1}{#3}}
%    \end{macrocode}

% \begin{macro}{\oldbibitem}
% \begin{macro}{\@rmaaBibArg}
% \begin{macro}{\@rmaaSpace}
% \begin{macro}{rmaabib}
% Of course, |\bibitem| needs to get the uniquification treatment as
% well. While we are at it, we redefine it to use consecutive integers
% for the key if the user leaves it blank (using the counter
% |rmaabib|). This allows one to put 
% |\bibitem{}| for all your references without getting annoying
% messages about multiply-defined labels.
% 
% We also take the opportunity to insert the optional argument to
% \cs{bibitem} if it is missing in the form Author(DATE) so that natbib
% doesn't try to do numerical citations. Natbib already provides a
% hook \cs{@extra@b@citeb} to help uniquify the citation key, so we
% just plug into that. However, it puts it \emph{after} the original
% key, not before as I had had it. Easiest solution is to change the
% way that I do it. 02 Jan 2003. Note that the \cs{@extra@b@citeb}
% does not get written to the .aux  file in the \cs{citation} command
% but it does work. 
% \changes{3.14}{06 Feb 2003}{Reinstalled old definition of \cs{bibitem}
%  for the case where the nonatbib option is requested, so that the
%  rmaacite package will work. This means that you can't use rmaacite
%  and natbib in the same volume.}
%    \begin{macrocode}
\if@natbib\def\@extra@b@citeb{:\articlename}\fi
%%\else
\let\oldbibitem\bibitem
\newcommand{\@rmaaBibArg}{}
\newcommand{\@rmaaSpace}{ }
\newcounter{rmaabib}
\if@natbib
\renewcommand{\bibitem}[2][ABC(2003)]{%
  \renewcommand{\@rmaaBibArg}{#2}
  \ifx\@rmaaBibArg\@rmaaEmpty
     \addtocounter{rmaabib}{1}
     \oldbibitem[#1]{\thermaabib:\articlename}
  \else\ifx\@rmaaBibArg\@rmaaSpace
     \addtocounter{rmaabib}{1}
     \oldbibitem[#1]{\thermaabib:\articlename}
  \else 
     \oldbibitem[#1]{#2:\articlename}
  \fi\fi
  \renewcommand{\@rmaaBibArg}{}
  }
\else
\renewcommand{\bibitem}[1]{%
  \renewcommand{\@rmaaBibArg}{#1}
  \ifx\@rmaaBibArg\@rmaaEmpty
     \addtocounter{rmaabib}{1}
     \oldbibitem{\articlename:\thermaabib}
  \else\ifx\@rmaaBibArg\@rmaaSpace
     \addtocounter{rmaabib}{1}
     \oldbibitem{\articlename:\thermaabib}
  \else
     \oldbibitem{\articlename:#1}
  \fi\fi
  \renewcommand{\@rmaaBibArg}{}
  }
\fi
%%\fi
%    \end{macrocode}
% \end{macro}
% \end{macro}
% \end{macro}
% \end{macro}

% \begin{macro}{oldpcite}
% \begin{macro}{oldscite}
% Then, we need to repeat the above scheme, but for the |rmaacite|
% definition. 
%    \begin{macrocode}
\newcommand{\@redefinebibitem}{%
  \let\oldbibitem\bibitem
  \let\oldcite\cite
  \let\oldpcite\pcite
  \let\oldscite\scite
  \def\bibitem[##1<##2>]##3{%
    \renewcommand{\@rmaaBibArg}{##3}
    \ifx\@rmaaBibArg\@rmaaEmpty
       \addtocounter{rmaabib}{1}
       \oldbibitem[##1<##2>]{\thermaabib:\articlename}
    \else\ifx\@rmaaBibArg\@rmaaSpace
       \addtocounter{rmaabib}{1}
       \oldbibitem[##1<##2>]{\thermaabib:\articlename}
    \else 
       \oldbibitem[##1<##2>]{##3:\articlename}
    % trust the author here
%%       \oldbibitem[##1<##2>]{##3}
    \fi\fi
    }
  }
%    \end{macrocode}
% \end{macro}
% \end{macro}


% \subsection{HTML Table of Contents}
% \label{sec:hyper-TOC}
% \begin{macro}{\HyperTOCEntry}
% \begin{macro}{\HyperTOCpart}
% \begin{macro}{\HyperTOCmeeting}
% \begin{macro}{\HyperTOCeditors}
% \begin{macro}{\HyperTOCpreamble}
% These commands are used to write information to the \texttt{.htoc}
% file, which is used in generating a Web Page of the Table of
% Contents. Since they only ever appear preceded by |\protect| in the
% argument to an |\addtocontents| command,  they should never need to
% be expanded by the macros. However, we define them here just in case
% \dots
%    \begin{macrocode}
\newcommand{\HyperTOCEntry}[3]{\relax}
\newcommand{\HyperTOCpart}[1]{\relax}
\newcommand{\HyperTOCmeeting}[3]{\relax}
\newcommand{\HyperTOCeditors}[1]{\relax}
\newcommand{\HyperTOCpreamble}{\relax}
%    \end{macrocode}
% \end{macro}
% \end{macro}
% \end{macro}
% \end{macro}
% \end{macro}
%

% \subsection{Formatting the poster abstracts} 
% These can now be used without using the book option.
%    \begin{macrocode}

\newcommand{\@makeminititle}{%
  \newpage
  \null 
  \vskip 3em%
  \begin{center}%
    {\@title \par}%
  \end{center}
  \vskip 3em%
  }

\newif\if@abstractsads
\@abstractsadsfalse
\newcommand\abstractsADS{\@abstractsadstrue}

\if@book\RequirePackage{multicol}\fi
\newenvironment{abstracts}[1][Abstracts of Contributed Papers]{%
  %% First, make sure we start on an odd page (06 Jun 2000)
  \if@book\clearpage\thispagestyle{empty}\cleardoublepage\fi
  %%
  \renewcommand{\articlename}{ABSTRACTS}
  \title{#1}\shortauthor{Abstracts}\shorttitle{Abstracts}
  \listofauthors{}
  \if@book
  \addcontentsline{toc}{chapter}{\protect\@fiddleauthors{\@TOCtitle}{\rmaaTOCauthors}}%
  \addtocontents{htc}{\protect\HyperTOCEntry{\@rawtitle}{\@rawauthors}{\articlename}}
  \fi
  \pagestyle{rmaa} % this now seems necessary (06 Jun 2000)
  \thispagestyle{rmaafirstpage}
%%  \twocolumn[\@makeminititle]
  \renewcommand{\abstract}[1]{%
    \renewcommand\ADS@abstract{##1}%
    \renewcommand{\@abstract}{##1}\@writeabstract%
  }
  \renewcommand{\title}[2][]{%
    \gdef\ADS@authors{}%
    \gdef\ADS@affil{}%
    \gdef\ADS@references{}%
    \gdef\ADS@keywords{}%
    \renewcommand{\@rmaaArg}{##1}
    \ifx\@rmaaArg\@rmaaEmpty
       \renewcommand{\@TOCaftertitle}{}
    \else
       \renewcommand{\@TOCaftertitle}{ (\@rmaaArg)}
    \fi
    \renewcommand{\@title}{\rmaa@TOCTitleStyle{##2}}%
    \renewcommand{\@TOCtitle}{\rmaa@TOCTitleStyle{##2}\@TOCaftertitle}%
    \renewcommand{\rmaaIndexTitle}{\rmaa@IndexTitleStyle{##2\@TOCaftertitle.}}%
    \renewcommand\ADS@title{##2}%
    \xdef\ADS@firstpage{\@arabic{\c@page}}%
    }
  \renewcommand{\author}[1]{\let\@thanks\@empty\renewcommand\@author{##1} }
  \renewcommand{\footnotetext}[2][XX]{\par\noindent\textsuperscript{##1}~##2}
  \renewcommand{\footnoterule}{\kern-3\p@ \hrule \@width 3em \kern 2.6\p@}
  \if@book\@begofarticlehook\fi
  \if@book\begin{multicols}{2}[\@makeminititle]\else\@makeminititle\fi
  \begingroup
  \setcounter{footnote}{0}
  }% and at the end of the environment....
  {\endgroup\if@book\end{multicols}\listofabstracts\label{RMAALastPage}\clearpage\fi}
%% auxiliary commands used:
%% \@writeabstract does all the actual formatting 
\newcommand{\@writeabstract}{%
  \if@book%
  \addcontentsline{loa}{abstract}{\protect\@fiddleauthors{\@TOCtitle}{\rmaaTOCauthors}}%
  \fi%
  \begin{center}\@title\par\@author\end{center}
  \@abstract
  \par\noindent\footnoterule\nobreak\footnotesize\@thanks
  \par\noindent\hrulefill\goodbreak
  \vskip 1cm \@plus0.5cm
  %% write ADS entry
  \if@book\if@abstractsads\ADS@dumpdata\fi\fi%
  \endgroup % this pair is to make each abstract a self-contained group
  \begingroup
  \setcounter{footnote}{0}
  }

\if@book
%% List of abstracts (stuff written to .loa file)
\newcommand{\listofabstracts}{%
  \begin{multicols}{2}[%
    \vskip 1.5cm \vfill % leave at least 1.5 cm and go to bottom of page if possible
    \centerline{\MakeUpperCase{List of Abstracts}}
    \vskip 1.5cm \@plus0.5cm
    ][6cm] % this makes it start a new page if there is not at least 6cm of space
    \renewcommand{\@fiddleauthors}[2]{##1\hskip1em plus2em minus0.4em\penalty-100  ##2}
    \parindent\z@
%%    \parskip.5\p@ \@plus .7\p@\relax 
    \parskip3\p@ \@plus 1\p@ \@minus 1\p@\relax % allow quite a lot of
                                % leeway here to avoid breaks inside
                                % entries 
    \renewcommand{\\}{\relax} % temporarily turn off explicit line breaks just in case
    \renewcommand{\altaffilmark}[1]{\relax} % ditto with this
    \@starttoc{loa}
  \end{multicols}
  }
\newcommand*{\l@abstract}[2]{%
  \par\penalty-1000
  #1\hskip 1em\penalty0\null\nobreak\hfill#2}
\fi

%    \end{macrocode}
 
% \subsection{Code for typesetting an entire volume of conference
% proceedings} \label{sec:whole-volume}
%    \begin{macrocode}

\if@book
%% environment for including the individual articles
\newcommand{\@fiddleauthors}[2]{\relax}
\renewcommand{\theequation}{\@arabic\c@equation}
\newenvironment{articles}%
{ \begingroup % make it a group so that author definitions do not
              % spill over
  \gdef\ADS@authors{}
  \pagestyle{rmaa}
  \renewenvironment{document}{%
    \@begofarticlehook
    \addtocounter{chapter}{1}%
    \addcontentsline{toc}{chapter}{%
      \protect\@fiddleauthors{\@TOCtitle}{\rmaaTOCauthors}}%
    \addtocontents{htc}{\protect\HyperTOCEntry{\@rawtitle}{\@rawauthors}{\articlename}}
    }{\relax}
  \renewcommand{\documentclass}[2][foo]{\relax} 
  \renewcommand{\usepackage}[2][foo]{%
    \def\@packagename{##2}\def\@rmaacitename{rmaacite}%
    \ifx\@packagename\@rmaacitename
    \@rmcitetrue\ClassInfo{rmaa}{This paper uses rmaacite}
    \makeatletter%%
%% This is file `rmaacite.sty',
%% generated with the docstrip utility.
%%
%% The original source files were:
%%
%% rmaa.dtx  (with options: `rmaacite')
%% 
%% PLEASE DO NOT EDIT THIS FILE! It is generated automatically from
%% the file rmaa.dtx using the accompanying Makefile and rmaa.ins
%% Any hand-edited changes will be lost
%% 
%% A LaTeX document class for the RMxAA
%% 
%% Author: William Henney (w.henney@crya.unam.mx)
%% 
%% Latest version available from: http://www.crya.unam.mx/rmaa
%% 
%% Copyright (C) 1999-2010 Instituto de Astronomia, UNAM, Mexico
%% 
%% This program is free software; you can redistribute it and/or
%% modify it under the terms of the GNU General Public License
%% as published by the Free Software Foundation; either version 2
%% of the License, or (at your option) any later version.
%% 
%% This program is distributed in the hope that it will be useful,
%% but WITHOUT ANY WARRANTY; without even the implied warranty of
%% MERCHANTABILITY or FITNESS FOR A PARTICULAR PURPOSE.  See the
%% GNU General Public License for more details.
%% 
%% You should have received a copy of the GNU General Public License
%% along with this program; if not, write to the Free Software
%% Foundation, Inc., 59 Temple Place - Suite 330, Boston, MA  02111-1307, USA.






\def\bibitem[#1<#2>]#3{\par\hangindent=1.5em\hangafter=1
  \if@filesw
  {\def\protect##1{\string ##1\space}
    \immediate\write\@auxout{\string\lbibcite{#3}{#1}{#2}}}
  \fi}
\@redefinebibitem % what it says
\def\lbibcite#1#2#3{\global\@namedef{b@#1}{#2}
                    \global\@namedef{by@#1}{#3}}

\def\@bcite#1#2{({#1\if@tempswa , #2\fi})}
\def\@pcite#1#2{#1\if@tempswa , #2\fi}
\def\@brcite{\ (}
\def\@ebrcite{)}
\def\@fmttop#1{\@fmtcite[\csname #1\endcsname]}
\def\@citex[#1]#2{%
%%  \if@filesw\immediate\write\@auxout{\string\citation{#2}}\fi
%% WJH added 16 Jan 2002
  \@for\@citeb:=#2\do{%
    \edef\@citeb{\expandafter\@firstofone\@citeb:\articlename}%
    \if@filesw\immediate\write\@auxout{\string\citation{\@citeb}}\fi}%
%% to here
  \def\@citea{}%
  \@cite{\@for\@citeb:=#2\do
    {%
      \edef\@citeb{\expandafter\@firstofone\@citeb:\articlename}%
      \@citea\def\@citea{;\ }\@ifundefined
      {b@\@citeb}{{\bf ?}\@warning {Citation `\@citeb' on page \thepage \space undefined}}%
      {\csname b@\@citeb\endcsname\@citebro\csname by@\@citeb\endcsname\@citebrc}}}%
  {#1}}

\def\cite{\@ifnextchar [{\let\@citebro=\space
                          \let\@citebrc=\relax
                          \let\@cite=\@bcite\@tempswatrue\@citex}
                        {\let\@citebro=\space
                          \let\@citebrc=\relax
                          \let\@cite=\@bcite\@tempswafalse\@citex[]}}
\def\pcite{\@ifnextchar [{\let\@citebro=\space
                          \let\@citebrc=\relax
                          \let\@cite=\@pcite\@tempswatrue\@citex}
                        {\let\@citebro=\space
                          \let\@citebrc=\relax
                          \let\@cite=\@pcite\@tempswafalse\@citex[]}}
\def\scite{\@ifnextchar [{\let\@citebro=\@brcite
                          \let\@citebrc=\@ebrcite
                          \let\@cite=\@pcite\@tempswatrue\@citex}
                        {\let\@citebro=\@brcite
                          \let\@citebrc=\@ebrcite
                          \let\@cite=\@pcite\@tempswafalse\@citex[]}}

\renewenvironment{thebibliography}
  {\section*{REFERENCES}\vspace*{\afterbiblioskip}\small
    \if@book\RereadBiblio\fi
    \sloppy\clubpenalty4000\widowpenalty4000
    \sfcode`\.=1000\relax
    \setlength\parindent{\z@}\setlength\parskip{\z@}}
  {\relax\normalsize\@outputaddressesifnecessary}








\endinput
%%
%% End of file `rmaacite.sty'.
\makeatother % read it in here so that all changes
    \else                                         % are local to paper
    \if@debug\ClassWarning{rmaa}{Package \@packagename\space was requested by \articlename\space}\fi
    \fi} 
  %% Don't allow individual articles to mess with the pages, volume or year
  \renewcommand{\SetFirstPage}[1]{\relax}
  \renewcommand{\SetYear}[1]{\relax}
  \renewcommand{\SetVolume}[1]{\relax}
  \renewcommand{\SetConfTitle}[1]{\relax}
  %% Don't allow individual articles to change these
  \renewcommand{\nonstopmode}{\relax}
  \renewcommand{\scrollmode}{\relax}
  \renewcommand{\outputplates}{\relax}
  %% XXXX need to check if using rmaacite and act accordingly
  %% avoid multiple declarations of the same length
  \let\oldnewlength=\newlength
  \renewcommand{\newlength}[1]{%
    \edef\reserved@a{\expandafter\@gobble\string ##1}%
    \@ifundefined{\reserved@a}{%
      \oldnewlength{##1}}{%
      \if@debug{\ClassInfo{rmaa}{\reserved@a\space already defined: doing nothing}}\fi%
      \relax}}%
  \renewcommand\newcommand{\providecommand}
  }%
{\endgroup\onecolumn}
     

%    \end{macrocode}
% Command to be used in `articles' environment for including an
% article
%    \begin{macrocode}
\newcommand{\article}[2][]{%
  % optional first argument is extra text to put in brackets after title in TOC
  % e.g. ``Invited Talk'' or ``Poster''
  \renewcommand{\articlename}{#2}
  % reset loads of counters
  \setcounter{footnote}{0}
  \setcounter{figure}{0}
  \setcounter{table}{0}
  \setcounter{section}{0}
  \setcounter{equation}{0}
  \setcounter{rmaabib}{0}
  \gdef\ADS@authors{}%
  \gdef\ADS@affil{}%
  \gdef\ADS@references{}%
  \@firstauthortrue
  \@firstaffiltrue
  \@firstaltaffiltrue
  \@firstkwdtrue
  \@rmcitefalse
  \if@suppressfulladdresses\else\@needaddressestrue\fi
  \renewcommand{\@rmaaArg}{#1}
  \ifx\@rmaaArg\@rmaaEmpty
     \renewcommand{\@TOCaftertitle}{}
  \else
     \renewcommand{\@TOCaftertitle}{ (#1)}
  \fi
%    \end{macrocode}
% In the main journal articles must start on an odd page. 
%    \begin{macrocode}
  \if@journal\clearpage\thispagestyle{empty}\cleardoublepage\fi
%% Now input the file, then run the end-of-article hook (reinstated 28
  %% Sep 2001)
  {\input{\@articledir/\articlename.tex}\@endofarticlehook}}
%    \end{macrocode}
% Command to separate articles into sessions. Write an entry in the
% TOC, makes sure we start on an odd page (added 06 Jun 2000), then do
% a title page if \cs{SessionTitlePages} was called. 
% \changes{3.12}{04 Feb 2003}{Added option to have title page for each
% session. Turn it on with \cs{SessionTitlePages}.}
%    \begin{macrocode}
\newcommand{\session}[1]{
  %% this does nothing except write an entry in the TOC
  \addcontentsline{toc}{part}{\rmaa@TOCPartStyle{#1}}
  \addtocontents{htc}{\protect\HyperTOCpart{#1}}
  %% and also makes sure we start on an odd page (added 06 Jun 2000)
  \clearpage\thispagestyle{empty}\cleardoublepage
  \stepcounter{part}
  \if@sessiontitlepages{%
    \onecolumn
    \thispagestyle{empty}
    \renewcommand{\articlename}{Section\thepart TitlePage}
    \label{RMAAFirstPage}
    \make@session@title{#1}
    \label{RMAALastPage}
    \clearpage\thispagestyle{empty}\cleardoublepage
  }\fi
  }
\newcommand{\make@session@title}[1]{%
  {%
    \setlength\parindent{0pt}
    \vspace*{\stretch{1}}
    \rule{\linewidth}{1mm}
    \begin{flushright}\Huge\bfseries\boldmath
      #1
    \end{flushright}
    \rule{\linewidth}{1mm}
    \vspace*{\stretch{2}}
  }
}
%    \end{macrocode}
% Command to turn on separate title pages for each session. 
%    \begin{macrocode}
\newcommand{\SessionTitlePages}{\@sessiontitlepagestrue}
\newcommand{\BoldSessionTitles}{%
  \renewcommand{\rmaa@TOCPartStyle}[1]{\textbf{\rmaa@TOCTitleStyle{##1}}}
  }
\newcommand{\ArticleDir}[1]{\renewcommand{\@articledir}{#1/}\graphicspath{{1/}}}
\newcommand{\@articledir}{.}
%\newcommand{\ps@rmaafront}{\ps@rmaa}% page style for front matter (same as ever, actually)
\newcommand{\ps@rmaafront}{% page style for front matter:  page # in each foot
  \renewcommand\@oddfoot{\hfil{\thepage}}%
  \renewcommand\@evenfoot{{\thepage}\hfil}%
  % shorttitle in both heads
  \renewcommand\@evenhead{\hfil\@shorttitle\hfil}%
  \renewcommand\@oddhead{\@evenhead}
  }
\newcommand{\ps@rmaafrontfirstpage}{% page style for first page of front matter
  % page # in each foot, on outside edge
  \renewcommand\@oddfoot{\hfil\thepage}%
  \renewcommand\@evenfoot{\thepage\hfil}%
  % nothing in the even heads
  \renewcommand\@evenhead{\relax}%
  \renewcommand\@oddhead{\relax}%
  }
%% slight modification to the standard \frontmatter and \mainmatter
%% commands, just to set the \pagestyle
\renewcommand\frontmatter{%
  \latexonecolumn
  \cleardoublepage
  \@frontmattertrue
  \@mainmatterfalse
  \pagenumbering{roman}
  \pagestyle{rmaafront}
  \addtocounter{page}{2} % leave space for the title page
  }
\renewcommand\mainmatter{%
  \pagestyle{empty}
  \if@rmtwocol\twocolumn\fi
  \cleardoublepage
  \@frontmatterfalse
  \@mainmattertrue
  \pagenumbering{arabic}
  \pagestyle{rmaa}
  }
\renewcommand\backmatter{%
  \clearpage\pagestyle{empty}\cleardoublepage
  \@mainmatterfalse
  %% put extra space in the TOC
  \addtocontents{toc}{\bigskip}
  }
%    \end{macrocode}
%
% \subsection{Front sections}\label{sec:front-sections}
% \begin{macro}{\notitle}
% \begin{macro}{\shortname}
% \begin{macro}{frontsection}
% The environment \texttt{frontsection} is for doing a section in the
% front matter (preface, list of participants, or whatever). It has 
% one optional argument:
% \begin{enumerate}
% \item[\#1] Short title (no embedded spaces) for
% use in making PS files. E.g., ``list\_of\_participants''.   
% \end{enumerate}
% and two mandatory arguments: 
% \begin{enumerate}
% \item[\#2] Title of section (e.g., `Preface')
% \item[\#3] Commands to be evaluated before processing the contents
% of the environment. This is mainly useful for putting in the things
% that would go in the preamble of a normal article  (e.g.,
% |\listofauthors|, |\indexauthor|, |\notitle|) 
% \end{enumerate}
%    \begin{macrocode}
\newcommand{\notitle}{\@titlevisiblefalse} \newcommand{\shortname}{}
\newenvironment{frontsection}[3][]{%
  \begingroup
  \renewcommand{\shortname}{#1}
  \ifx\shortname\@rmaaEmpty
     \renewcommand{\articlename}{#2}
  \else
     \renewcommand{\articlename}{\shortname}
  \fi
  \@begofarticlehook
  \@titlevisibletrue % default
  \thispagestyle{rmaafrontfirstpage}
  \listofauthors{} % blank out this in case it is not specified
  \title{{#2}}\shorttitle{#2}\shortauthor{#2}
  #3 % evaluate the command given as second mandatory arg
  \addcontentsline{toc}{chapter}{\protect\@fiddleauthors{\@TOCtitle}{\rmaaTOCauthors}}%
  \addtocontents{htc}{\protect\HyperTOCEntry{\@rawtitle}{\@rawauthors}{\articlename}}
  \if@titlevisible
     \vspace*{1cm}\centerline{\rmaa@FrontSecStyle{\@title}}\vspace*{1cm}
  \fi
  \setlength{\parskip}{3ex}
  }{\frontlabel{RMAALastPage}\clearpage\endgroup}
%    \end{macrocode}
% \end{macro}
% \end{macro}
% \end{macro}
%    \begin{macrocode}
%% Definition of the command \frontlabel, which rights out something like `F4' to
%% the .aux file, instead of `\textlatin  {i}'. This is so it is easy to post-process
%% the file for making PS files of individual sections
\newcommand{\frontlabel}[1]{%
  \@bsphack
  \protected@write\@auxout{}%
     {\string\newlabel{\articlename:#1}{{}{F\@arabic{\c@page}}}}
  \@esphack
  }
%% User command to say X, Y, Z,\\ Editors of the Proceedings
\newcommand{\EditorSignature}{\vspace*{\parskip}\rightline{\@editors}%
  \rightline{Editors of the proceedings}}
%% Fiddle with \thesection, etc. so that they don't include the chapter number
%% Basically go back to the `article' definitions of these
\renewcommand{\thesection}{\@arabic\c@section}
\renewcommand{\thefigure}{\@arabic\c@figure}
\renewcommand{\thetable}{\@arabic\c@table}

\newcommand{\@oldargs}{}
\newcommand{\@newargs}{}
%% add TOC entry for each paper
%\renewcommand{\title}[1]%
%{\renewcommand{\@title}{\rmaa@TitleStyle{#1}}%
%\addtocounter{chapter}{1}\addcontentsline{toc}{chapter}{#1}%
%}

%% Formatting of the table of contents
%% dimensions of various parts
\newlength{\TOCtextwidth}\setlength{\TOCtextwidth}{12.65cm}
\newlength{\TOCpagenumwidth}\setlength{\TOCpagenumwidth}{1.3cm}
\newlength{\TOCentrywidth}\setlength{\TOCentrywidth}{\TOCtextwidth}
\addtolength{\TOCentrywidth}{-\TOCpagenumwidth} 
\newlength{\TOCextramargin}\setlength{\TOCextramargin}{1.25cm}
\newcommand{\TOCbox}[1]{\hspace{\TOCextramargin}\parbox{\TOCtextwidth}{#1}\par}
\newlength{\TOCextramarginx}\setlength{\TOCextramarginx}{\TOCextramargin}\addtolength{\TOCextramarginx}{4.3cm}
\newlength{\TOCtextwidthx}\setlength{\TOCtextwidthx}{\TOCtextwidth}\addtolength{\TOCtextwidthx}{-4.3cm}
\newcommand{\TOCboxx}[1]{\hspace{\TOCextramarginx}\parbox{\TOCtextwidthx}{#1}\par}

\newcommand{\RMAAbannerimage}{rmsc}
\renewcommand\tableofcontents{%
  \renewcommand{\articlename}{CONTENTS}
  \begingroup
  \@begofarticlehook
  \if@twocolumn
     \@restonecoltrue\onecolumn
  \else
     \@restonecolfalse
  \fi
  %% 9pt for journal style
  \if@proceedings\else\small\fi
  \if@newstyle\small\fi %% same for new style
  \if@proceedings\thispagestyle{rmaafrontfirstpage}\else\thispagestyle{empty}\fi
  \shortauthor{Contents}\shorttitle{Contents}
  \vspace*{18mm}
  \if@printbanner
  %% put in the banner header (need to get better image of this)
     \TOCbox{\hspace*{-0.15cm}\includegraphics[width=0.9\TOCtextwidth]{\RMAAbannerimage}}
     \vspace*{8mm}
  %% volume and publication date (should really be in funny font)
     \TOCboxx{\@volumen\\[2mm] \@mes\ \@anyo}
  \else
  %% put in the banner header (need to get better image of this)
     \TOCbox{\hspace*{-0.15cm}\includegraphics[width=0.9\TOCtextwidth]{\RMAAbannerimage}}
     \vspace*{12.5mm}
  %% volume and publication date (should really be in funny font)
     \TOCboxx{~\\[2mm] ~}
  \fi
  \vspace*{16mm}
  %% This is the command that fiddles with the TOC entry so that the
  %% title is flushleft but the authors are flushright
  \renewcommand{\@fiddleauthors}[2]{%
    % first put the paper title
    \renewcommand{\\}{\relax} % temporarily turn off explicit line breaks just in case
    \renewcommand{\altaffilmark}[1]{\relax} % ditto with this
    \renewcommand{\textsuperscript}[1]{\relax} % ditto with this 
    ##1
    % now save the author list in a box
    \sbox\@tempboxa{##2}%
    % check that box is shorter than the total width available for the entry
    \ifdim \wd\@tempboxa <\TOCentrywidth
       % if so, then put it right here, flush against right margin
       \penalty10000\hskip1em plus10em\penalty0\hbox{}\nobreak\hfill\box\@tempboxa
       % Note that we force a space of at least 1em between title and authorlist, 
       % stretchable to 11em since we may have very short titles but still with an 
       % authorlist too long to go on the same line.  The only linebreak allowed is 
       % directly before the authorlist.
    \else
       % if it is too long, we set the original authorlist text in its own parbox
       \par\vskip1ex\parbox[b]{\TOCentrywidth}{\raggedleft ##2}
    \fi
    }
  \@starttoc{toc}%
  \if@restonecol\twocolumn\fi
  \frontlabel{RMAALastPage}
  \clearpage
  \thispagestyle{empty}
  \cleardoublepage % make sure we start on an odd page
  \endgroup

  %%% ALSO, do the HTML TOC
  %% The file that will be used to make the HTML TOC
  \@starttoc{htc}
  %% All the stuff in this file does nothing in the book style
  \addtocontents{htc}{\protect\HyperTOCmeeting{\@conftitle}{\@reunionplace}{\@reuniondate}}
  \addtocontents{htc}{\protect\HyperTOCeditors{\@editors}}
  \addtocontents{htc}{\protect\HyperTOCpreamble}
  \addtocontents{htc}{\protect\HyperTOCEntry{Table of Contents}{}{\articlename}}

  %% and the ADS file
  \ADS@openfile
  }
%% This is now much simpler
\renewcommand\l@chapter[2]{%TOC entry for articles
  \ifnum \c@tocdepth >\m@ne
    \vskip 1.0em \@plus\p@
    \TOCbox{%
      \parbox[b]{\TOCentrywidth}{#1}\parbox[b]{\TOCpagenumwidth}{\hfill#2}
    }
  \fi} 
\renewcommand\l@part[2]{%TOC entry for sessions: centered, no page number
  \ifnum \c@tocdepth >\m@ne
    \vskip 1.5em \@plus1.5em \goodbreak
    \TOCbox{%
      \parbox[b]{\TOCentrywidth}{\hfil#1\hfil}\hfill
    }
  \fi\nobreak} 

%% Formatting theindex
%% strongly encourage page breaks between index entries
\renewcommand{\@idxitem}{\par\penalty-1000\hangindent.45cm} 
\renewcommand{\indexspace}{\par\penalty-1000\vskip 10\p@ \@plus15\p@ \@minus3\p@\relax}
\renewenvironment{theindex}{%
  %% first, make sure we start on an odd page (06 Jun 2000)
  \clearpage\thispagestyle{empty}\cleardoublepage
  %% this looks a bit like an article but is two columns and without abstract, etc
  \renewcommand{\articlename}{INDEX}
  \title{Author Index}\shortauthor{Author Index}\shorttitle{Author Index}
  \listofauthors{}
  \addcontentsline{toc}{chapter}{\protect\@fiddleauthors{\@TOCtitle}{\rmaaTOCauthors}}%
  \addtocontents{htc}{\protect\HyperTOCEntry{\@rawtitle}{\@rawauthors}{\articlename}}
  \pagestyle{rmaa}
  \thispagestyle{rmaafirstpage}
  \@begofarticlehook
  %% 9pt for journal style
  \if@proceedings\else\small\fi
  \if@newstyle\small\fi
  \twocolumn[\@makeminititle]
  %% vertical spacing (see also definition of \@idxitem and \indexspace
  \parskip\z@ \@plus .6\p@\relax % allow quite a lot of leaway here to avoid breaks inside entries
  %% horizontal spacing, breaking and glue parameters (see also the stuff in the definition 
  %% of \indexauthor
  \parindent\z@ 
  \exhyphenpenalty10000 % don't break on explicit hyphens, it always looks bad
  \rightskip=0pt plus.2em % allow a tiny bit of leeway at the ends of lines
  \spaceskip.3333em plus.5em minus.15em % a lot of leeway between words
  \xspaceskip.5em plus.5em minus.4em % this seems to have no effect
  \let\item\@idxitem
  \renewcommand{\\}{\relax} % temporarily turn off explicit line breaks just in case
  \renewcommand{\altaffilmark}[1]{\relax} % ditto with this
  \renewcommand{\textsuperscript}[1]{\relax} % ditto with this
  }{\label{RMAALastPage}\clearpage} % clearpage necessary so that we know about \shortauthor



%% end definitions for book styles
\fi

%    \end{macrocode}
% In order to do away with plates, we sometimes want a figure to
% appear in color only in the online version. The following command
% can bve used in a figure caption to typeset a note to this effect. 
%    \begin{macrocode}
\newcommand\ColorOnlineOnly{%
  \textsc{Note: this figure is available in color in the electronic
    version of this article, obtainable from}
  \texttt{http://\discretionary{}{}{}www.\discretionary{}{}{}astroscu.\discretionary{}{}{}unam.\discretionary{}{}{}mx/\discretionary{}{}{}\~{}rmaa/}.}
%% Plates
%% Define 2 counters: plate is the number in the whole document, plateloop is used 
%% internally by \outputplates
\newcounter{plate}
\newcounter{plateloop}
%% Currently limited to 15 plates in document
%% when we do a label, we need to write to \newlabel entries to the .aux file, one 
%% that says RMAAPlateX so that it can be referred to in \ouputplates, and one that
%% says articlename:RMAAPlateX so that the dvips program can find it and asociate it
%% with the right article
\newcommand{\@platelabeli}{\oldlabel{RMAAPlate1}\label{RMAAPlate1}}
\newcommand{\@platelabelii}{\oldlabel{RMAAPlate2}\label{RMAAPlate2}}
\newcommand{\@platelabeliii}{\oldlabel{RMAAPlate3}\label{RMAAPlate3}}
\newcommand{\@platelabeliv}{\oldlabel{RMAAPlate4}\label{RMAAPlate4}}
\newcommand{\@platelabelv}{\oldlabel{RMAAPlate5}\label{RMAAPlate5}}
\newcommand{\@platelabelvi}{\oldlabel{RMAAPlate6}\label{RMAAPlate6}}
\newcommand{\@platelabelvii}{\oldlabel{RMAAPlate7}\label{RMAAPlate7}}
\newcommand{\@platelabelviii}{\oldlabel{RMAAPlate8}\label{RMAAPlate8}}
\newcommand{\@platelabelix}{\oldlabel{RMAAPlate9}\label{RMAAPlate9}}
\newcommand{\@platelabelx}{\oldlabel{RMAAPlate10}\label{RMAAPlate10}}
\newcommand{\@platelabelxi}{\oldlabel{RMAAPlate11}\label{RMAAPlate11}}
\newcommand{\@platelabelxii}{\oldlabel{RMAAPlate12}\label{RMAAPlate12}}
\newcommand{\@platelabelxiii}{\oldlabel{RMAAPlate13}\label{RMAAPlate13}}
\newcommand{\@platelabelxiv}{\oldlabel{RMAAPlate14}\label{RMAAPlate14}}
\newcommand{\@platelabelxv}{\oldlabel{RMAAPlate15}\label{RMAAPlate15}}
\newcommand{\@outputplatelabeli}{\oldlabel{RMAAOutputPlate1}}
\newcommand{\@outputplatelabelii}{\oldlabel{RMAAOutputPlate2}}
\newcommand{\@outputplatelabeliii}{\oldlabel{RMAAOutputPlate3}}
\newcommand{\@outputplatelabeliv}{\oldlabel{RMAAOutputPlate4}}
\newcommand{\@outputplatelabelv}{\oldlabel{RMAAOutputPlate5}}
\newcommand{\@outputplatelabelvi}{\oldlabel{RMAAOutputPlate6}}
\newcommand{\@outputplatelabelvii}{\oldlabel{RMAAOutputPlate7}}
\newcommand{\@outputplatelabelviii}{\oldlabel{RMAAOutputPlate8}}
\newcommand{\@outputplatelabelix}{\oldlabel{RMAAOutputPlate9}}
\newcommand{\@outputplatelabelx}{\oldlabel{RMAAOutputPlate10}}
\newcommand{\@outputplatelabelxi}{\oldlabel{RMAAOutputPlate11}}
\newcommand{\@outputplatelabelxii}{\oldlabel{RMAAOutputPlate12}}
\newcommand{\@outputplatelabelxiii}{\oldlabel{RMAAOutputPlate13}}
\newcommand{\@outputplatelabelxiv}{\oldlabel{RMAAOutputPlate14}}
\newcommand{\@outputplatelabelxv}{\oldlabel{RMAAOutputPlate15}}
\newcommand{\@platerefi}{\oldpageref{RMAAPlate1}}
\newcommand{\@platerefii}{\oldpageref{RMAAPlate2}}
\newcommand{\@platerefiii}{\oldpageref{RMAAPlate3}}
\newcommand{\@platerefiv}{\oldpageref{RMAAPlate4}}
\newcommand{\@platerefv}{\oldpageref{RMAAPlate5}}
\newcommand{\@platerefvi}{\oldpageref{RMAAPlate6}}
\newcommand{\@platerefvii}{\oldpageref{RMAAPlate7}}
\newcommand{\@platerefviii}{\oldpageref{RMAAPlate8}}
\newcommand{\@platerefix}{\oldpageref{RMAAPlate9}}
\newcommand{\@platerefx}{\oldpageref{RMAAPlate10}}
\newcommand{\@platerefxi}{\oldpageref{RMAAPlate11}}
\newcommand{\@platerefxii}{\oldpageref{RMAAPlate12}}
\newcommand{\@platerefxiii}{\oldpageref{RMAAPlate13}}
\newcommand{\@platerefxiv}{\oldpageref{RMAAPlate14}}
\newcommand{\@platerefxv}{\oldpageref{RMAAPlate15}}
\newcounter{figplatei}
\newcounter{figplateii}
\newcounter{figplateiii}
\newcounter{figplateiv}
\newcounter{figplatev}
\newcounter{figplatevi}
\newcounter{figplatevii}
\newcounter{figplateviii}
\newcounter{figplateix}
\newcounter{figplatex}
\newcounter{figplatexi}
\newcounter{figplatexii}
\newcounter{figplatexiii}
\newcounter{figplatexiv}
\newcounter{figplatexv}
%% Printed output should use the counter `plate'
\renewcommand{\theplate}{\@arabic\c@plate}
\renewcommand{\theplateloop}{\@arabic\c@plateloop}
%% Command to output all the saved plates in this paper
\newcommand{\outputplates}{%
  \setcounter{plateloop}{0}
  \if@book\else\@endofarticlehook\fi % not entirely sure why, but this
                                     % seems to be needed
  \clearpage % first output the last page of text
  % now set pagestyle for blank even pages and make sure we start on an odd page
  \pagestyle{empty}
  \cleardoublepage
  \onecolumn
  \parindent\z@
  % first put in an extra sheet
  % the main purpose of this is so there is always a blank even page before 
  % each plate for filling in the PS file if the article has an odd number
  % of pages (see splitbook.sh)
  \null\vfil\centerline{PLATES}\vfil\clearpage\pagestyle{rmaaplate}\cleardoublepage
  % output each plate in turn
  \@whilenum\value{plateloop}<\value{plate}\do{%
    \addtocounter{plateloop}{1}
    % restore the figure counter from \@figcounti, \@figcountii, etc
    \setcounter{figure}{\value{figplate\roman{plateloop}}}
    % spew out the plate (in its own group for safety)
    {\csname @plate\roman{plateloop}\endcsname\par}
    }
  \pagestyle{rmaa}
  % reset the page number as if nothing had happened
  \addtocounter{page}{-\value{plate}}
  \addtocounter{page}{-\value{plate}}
  }
\newcommand\figcaption{%
  \def\@captype{figure}
  \renewcommand{\@makecaption}{\@makefigcap}
  \caption}
\newcommand\tabcaption{%
  \def\@captype{table}
  \renewcommand{\@makecaption}{\@maketabcap}
  \caption}

\newcommand{\plate}[2][1]{%
  %% User command to make a plate page: 
  %% Mandatory argument is what goes on the plate - this should not 
  %% include any floating environments (e.g. figure). Instead, \figcaption 
  %% should be used to do the captions. 
  %% Optional argument is the number of figures (i.e. captions) on the plate
  %% page (default is 1). This is necessary for incrementing the figure counter. 
  %%
  %% Note that we have to use a \gdef rather than an \xdef when we 
  %% save the plate output in a command. This means that in the book styles when it 
  %% comes to spewing out the plate, we will have forgotten about any user macros 
  %% that may have been defined. This can be got round by repeating the macros inside the plate.
  %% (Presumably, one would use a \newcommand rather than \renewcommand if they were new.)
  % increment global and local counters
  \addtocounter{plate}{1}
  % write appropriate label to .aux file
  \csname @platelabel\roman{plate}\endcsname
  % define a global macro called \@platei, \@plateii, \@plateiii, \@plateiv, ...
  \expandafter\gdef\csname @plate\roman{plate}\endcsname{%
    \csname @outputplatelabel\roman{plateloop}\endcsname
    \null\vfil #2 \vfil
    \cleardoublepage
    }
  % save the current shortauthor and shorttitle
  % need to be \xdef to take current values of \@shortauthor, etc 
  \expandafter\xdef\csname @shorta\roman{plate}\endcsname{\@shortauthor}
  \expandafter\xdef\csname @shortt\roman{plate}\endcsname{\@shorttitle}
  % save the current value of the figure counter for later use in \outputplates
  \setcounter{figplate\roman{plate}}{\value{figure}}
  % increment the figure counter according to the optional argument
  \addtocounter{figure}{#1}
  }
%% page style for plate pages 
\newcommand{\ps@rmaaplate}{%
  % Plate X (see page Y) in the foot
  \renewcommand\@oddfoot{\hfil\csname @shorta\roman{plateloop}\endcsname\ 
    (see page~\csname @plateref\roman{plateloop}\endcsname)\hfil}%
  \renewcommand\@evenfoot{\relax}%
  % nothing in even heads (blank pages)
  \renewcommand\@evenhead{\relax}
  % short title and page num in the odd heads
  \renewcommand\@oddhead{\hfil\csname @shortt\roman{plateloop}\endcsname%
    \hfil\llap{PLATE~\theplateloop}}
  }





%% Things for the User Guide
\if@guide
%% stolen from doc.sty
\def\MakeShortVerb#1{%
  \expandafter\ifx\csname cc\string#1\endcsname\relax
  \@shortvrbinfo{Made }{#1}%
  \add@special{#1}%
  \expandafter
  \xdef\csname cc\string#1\endcsname{\the\catcode`#1}%
  \begingroup
  \catcode`\~\active  \lccode`\~`#1%
  \lowercase{%
    \global\expandafter\let
    \csname ac\string#1\endcsname~%
    \gdef~{\verb~}}%
  \endgroup
  \global\catcode`#1\active
  \else
  \@shortvrbinfo\@empty{#1 already}%
  \fi}
\def\DeleteShortVerb#1{%
  \expandafter\ifx\csname cc\string#1\endcsname\relax
  \else
  \@shortvrbinfo{Deleted }{#1 as}%
  \rem@special{#1}%
  \global\catcode`#1\csname cc\string#1\endcsname
  \global \expandafter\let \csname cc\string#1\endcsname \relax
  \ifnum\catcode`#1=\active
  \begingroup
  \catcode`\~\active   \lccode`\~`#1%
  \lowercase{%
    \global\expandafter\let\expandafter~%
    \csname ac\string#1\endcsname}%
  \endgroup \fi \fi}
\def\@shortvrbinfo#1#2{%
  \PackageInfo{doc}{%
    #1\expandafter\@gobble\string#2 a short reference
    for \string\verb}}
\def\add@special#1{%
  \rem@special{#1}%
  \expandafter\gdef\expandafter\dospecials\expandafter
  {\dospecials \do #1}%
  \expandafter\gdef\expandafter\@sanitize\expandafter
  {\@sanitize \@makeother #1}}
\def\rem@special#1{%
  \def\do##1{%
    \ifnum`#1=`##1 \else \noexpand\do\noexpand##1\fi}%
  \xdef\dospecials{\dospecials}%
  \begingroup
  \def\@makeother##1{%
    \ifnum`#1=`##1 \else \noexpand\@makeother\noexpand##1\fi}%
  \xdef\@sanitize{\@sanitize}%
  \endgroup}
%% define | as shortcut for \verbatim
\MakeShortVerb{\|}
%% comand to specify arguments to LaTeX commands
\newcommand{\rmArg}[1]{%
  \ensuremath{\langle}{\scriptsize \textsf{\MakeUpperCase{#1}}}\ensuremath{\rangle}}
%\newcommand{\rmArg}[1]{\ensuremath{\langle}\textsc{#1}\ensuremath{\rangle}}
\providecommand{\BibTeX}{{\rmfamily B\kern-.05em%
    \textsc{i\kern-.025em b}\kern-.08em%
    T\kern-.1667em\lower.7ex\hbox{E}\kern-.125emX}}
%% Environment for code examples
\newenvironment{Code}{%
  \begin{list}{}{\setlength{\leftmargin}{0.05\columnwidth}} 
  \item[]}{\end{list}}
\fi
%</package>
%    \end{macrocode}

% \section{User commands for astronomical articles}
% \label{sec:user-commands}
% Some random definitions. These go both into the main class file and
% also into |adshtml.sty|.
% 
% Definitions of the journal name in different incarnations.  
%    \begin{macrocode}
%<*package|adshtml-style>
\newcommand{\RMAA}{RMxAA}
\newcommand{\RMAAlong}{Revista Mexicana de Astronom\'{\i}a y Astrof\'{\i}sica}
\newcommand{\RMAASC}{RMxAC}
\newcommand{\RMAASClong}{\RMAA\ (Serie de Conferencias)}
%    \end{macrocode}
% \begin{macro}{\ibidrule}
% This puts a horizontal rule with a period after it. It is used when
% the author list of a reference is identical to the preceding one. 
%    \begin{macrocode}
\newcommand{\ibidrule}{\rule[-0.03cm]{4.7em}{.01cm}.\ }
%    \end{macrocode}
% \end{macro}
% \begin{macro}{\ThisVolume}
% References to other articles in the same volume can be coded as,
% e.g., |\bibitem{} Bloggs, J. \ThisVolume{bloggs}|.
%    \begin{macrocode}
\newcommand{\ThisVolume}[1]{%
  \@year, \if@proceedings\RMAASC\else\RMAA\fi, \@volume,
  \if@book\oldpageref{#1:RMAAFirstPage}\else 000\fi\ (this volume)}
%    \end{macrocode}
% \end{macro}
% Other stuff that may be useful in ref lists.
%    \begin{macrocode}
\newcommand{\etal}{et al.\@}
%    \end{macrocode}
% Various new commands from AAS macros
%    \begin{macrocode}
\def\acknowledgments{\vskip 3ex plus .8ex minus .4ex}
\let\acknowledgements=\acknowledgments                  % second spelling
\let\jnl@style=\rm
\def\ref@jnl#1{{\jnl@style#1}}
\DeclareRobustCommand\aj{\ref@jnl{AJ}}                   
\DeclareRobustCommand\araa{\ref@jnl{ARA\&A}}             
\DeclareRobustCommand\apj{\ref@jnl{ApJ}}                 
\DeclareRobustCommand\apjl{\ref@jnl{ApJ}}                
\DeclareRobustCommand\apjs{\ref@jnl{ApJS}}               
\DeclareRobustCommand\ao{\ref@jnl{Appl.~Opt.}}           
\DeclareRobustCommand\apss{\ref@jnl{Ap\&SS}}             
\DeclareRobustCommand\aap{\ref@jnl{A\&A}}                
\DeclareRobustCommand\aapr{\ref@jnl{A\&A~Rev.}}          
\DeclareRobustCommand\aaps{\ref@jnl{A\&AS}}              
\DeclareRobustCommand\azh{\ref@jnl{AZh}}                 
\DeclareRobustCommand\baas{\ref@jnl{BAAS}}               
\DeclareRobustCommand\jrasc{\ref@jnl{JRASC}}             
\DeclareRobustCommand\memras{\ref@jnl{MmRAS}}            
\DeclareRobustCommand\mnras{\ref@jnl{MNRAS}}             
\DeclareRobustCommand\pra{\ref@jnl{Phys.~Rev.~A}}        
\DeclareRobustCommand\prb{\ref@jnl{Phys.~Rev.~B}}        
\DeclareRobustCommand\prc{\ref@jnl{Phys.~Rev.~C}}        
\DeclareRobustCommand\prd{\ref@jnl{Phys.~Rev.~D}}        
\DeclareRobustCommand\pre{\ref@jnl{Phys.~Rev.~E}}        
\DeclareRobustCommand\prl{\ref@jnl{Phys.~Rev.~Lett.}}    
\DeclareRobustCommand\pasp{\ref@jnl{PASP}}               
\DeclareRobustCommand\pasj{\ref@jnl{PASJ}}               
\DeclareRobustCommand\qjras{\ref@jnl{QJRAS}}             
\DeclareRobustCommand\skytel{\ref@jnl{S\&T}}             
\DeclareRobustCommand\solphys{\ref@jnl{Sol.~Phys.}}      
\DeclareRobustCommand\sovast{\ref@jnl{Soviet~Ast.}}      
\DeclareRobustCommand\ssr{\ref@jnl{Space~Sci.~Rev.}}     
\DeclareRobustCommand\zap{\ref@jnl{ZAp}}                 
\DeclareRobustCommand\nat{\ref@jnl{Nature}}              
\DeclareRobustCommand\iaucirc{\ref@jnl{IAU~Circ.}}
\DeclareRobustCommand\aplett{\ref@jnl{Astrophys.~Lett.}}
\DeclareRobustCommand\apspr{\ref@jnl{Astrophys.~Space~Phys.~Res.}}
\DeclareRobustCommand\bain{\ref@jnl{Bull.~Astron.~Inst.~Netherlands}}
\DeclareRobustCommand\fcp{\ref@jnl{Fund.~Cosmic~Phys.}}
\DeclareRobustCommand\gca{\ref@jnl{Geochim.~Cosmochim.~Acta}}
\DeclareRobustCommand\grl{\ref@jnl{Geophys.~Res.~Lett.}}
\DeclareRobustCommand\jcp{\ref@jnl{J.~Chem.~Phys.}}      
\DeclareRobustCommand\jgr{\ref@jnl{J.~Geophys.~Res.}}    
\DeclareRobustCommand\jqsrt{\ref@jnl{J.~Quant.~Spec.~Radiat.~Transf.}}
\DeclareRobustCommand\memsai{\ref@jnl{Mem.~Soc.~Astron.~Italiana}}
\DeclareRobustCommand\nphysa{\ref@jnl{Nucl.~Phys.~A}}
\DeclareRobustCommand\physrep{\ref@jnl{Phys.~Rep.}}
\DeclareRobustCommand\physscr{\ref@jnl{Phys.~Scr}}
\DeclareRobustCommand\planss{\ref@jnl{Planet.~Space~Sci.}}
\DeclareRobustCommand\procspie{\ref@jnl{Proc.~SPIE}}
\let\astap=\aap
\let\apjlett=\apjl
\let\apjsupp=\apjs
\let\applopt=\ao
\def\phn{\phantom{0}}
\def\phd{\phantom{.}}
\def\phs{\phantom{\ensuremath{-}}}
\def\phm#1{\phantom{#1}}
%    \end{macrocode}
% All these are redone as robust and using more modern commands where
% possible. The robustness is so they will not get expanded if used in
% abstracts and titles when these are passed to the ADS file. 
% \changes{3.16}{2003/04/07}{
%  Added \cs{Color} command for doing photometric colors.}
%    \begin{macrocode}
%%\def\sun{\hbox{$\odot$}}
\DeclareRobustCommand\sun{\ensuremath{\odot}}
%%\def\earth{\hbox{$\oplus$}}
\DeclareRobustCommand\earth{\ensuremath{\oplus}}
%%\def\lesssim{\mathrel{\hbox{\rlap{\hbox{\lower4pt\hbox{$\sim$}}}\hbox{$<$}}}}
%%\def\gtrsim{\mathrel{\hbox{\rlap{\hbox{\lower4pt\hbox{$\sim$}}}\hbox{$>$}}}}
\DeclareRobustCommand\lesssim{\mathrel{\hbox{\smaller\smaller\rlap{\hbox{\lower3pt\hbox{$\sim$}}}\raise1.5pt\hbox{$<$}}}}
\DeclareRobustCommand\gtrsim{\mathrel{\hbox{\smaller\smaller\rlap{\hbox{\lower3pt\hbox{$\sim$}}}\raise1.5pt\hbox{$>$}}}}
%%\def\lta{\mathrel{\spose{\lower 3pt\hbox{$\mathchar"218$}}
%%     \raise 2.0pt\hbox{$\mathchar"13C$}}}
%%\def\gta{\mathrel{\spose{\lower 3pt\hbox{$\mathchar"218$}}
%%     \raise 2.0pt\hbox{$\mathchar"13E$}}}
\DeclareRobustCommand\sq{\hbox{\rlap{$\sqcap$}$\sqcup$}}
%%\def\arcdeg{\hbox{$^\circ$}}
%%\def\arcmin{\hbox{$^\prime$}}
%%\def\arcsec{\hbox{$^{\prime\prime}$}}
\DeclareRobustCommand\arcdeg{\ensuremath{^\circ}}
\DeclareRobustCommand\arcmin{\ensuremath{^\prime}}
\DeclareRobustCommand\arcsec{\ensuremath{^{\prime\prime}}}
%%\def\fd{\hbox{$.\!\!^{\rm d}$}}
%%\def\fh{\hbox{$.\!\!^{\rm h}$}}
%%\def\fm{\hbox{$.\!\!^{\rm m}$}}
%%\def\fs{\hbox{$.\!\!^{\rm s}$}}
%%\def\fdg{\hbox{$.\!\!^\circ$}}
%%\def\farcm{\hbox{$.\mkern-4mu^\prime$}}
%%\def\farcs{\hbox{$.\!\!^{\prime\prime}$}}
%%\def\fp{\hbox{$.\!\!^{\scriptscriptstyle\rm p}$}}
\DeclareRobustCommand\fd{\ensuremath{.\!\!^\mathrm{d}}}
\DeclareRobustCommand\fh{\ensuremath{.\!\!^\mathrm{h}}}
\DeclareRobustCommand\fm{\ensuremath{.\!\!^\mathrm{m}}}
\DeclareRobustCommand\fs{\ensuremath{\mkern0.5mu.\mkern-5mu^\mathrm{s}\mkern-1mu}}
\DeclareRobustCommand\fdg{\ensuremath{.\!\!^\circ}}
\DeclareRobustCommand\farcm{\ensuremath{.\mkern-4mu^\prime}}
\DeclareRobustCommand\farcs{\ensuremath{\mkern1mu.\mkern-6mu^{\prime\prime}\mkern-1.5mu}}
%%\DeclareRobustCommand\farcs{\ensuremath{.\!\!^{\prime\prime}}}
\DeclareRobustCommand\fp{\ensuremath{.\!\!^{\scriptscriptstyle\rm p}}}
%%\def\micron{\hbox{$\mu$m}}
\DeclareRobustCommand\micron{\ensuremath{\mu\mathrm{m}}}
\let\la=\lesssim                        
\let\ga=\gtrsim
%%\def\case#1#2{\hbox{$\frac{#1}{#2}$}}
%%\def\slantfrac#1#2{\hbox{$\,^#1\!/_#2$}}
%%\def\onehalf{\slantfrac{1}{2}}
%%\def\onethird{\slantfrac{1}{3}}
%%\def\twothirds{\slantfrac{2}{3}}
%%\def\onequarter{\slantfrac{1}{4}}
%%\def\threequarters{\slantfrac{3}{4}}
\DeclareRobustCommand\case[2]{\ensuremath{\frac{#1}{#2}}}
\DeclareRobustCommand\slantfrac[2]{\ensuremath{\,^#1\!/_#2}}
\DeclareRobustCommand\onehalf{\slantfrac{1}{2}}
\DeclareRobustCommand\onethird{\slantfrac{1}{3}}
\DeclareRobustCommand\twothirds{\slantfrac{2}{3}}
\DeclareRobustCommand\onequarter{\slantfrac{1}{4}}
\DeclareRobustCommand\threequarters{\slantfrac{3}{4}}
%%\def\ubvr{\hbox{$U\!BV\!R$}}            
%%\def\ub{\hbox{$U\!-\!B$}}               
%%\def\bv{\hbox{$B\!-\!V$}}               
%%\def\vr{\hbox{$V\!-\!R$}}               
%%\def\ur{\hbox{$U\!-\!R$}}
%    \end{macrocode}
% New commands for ``photometric'' colors 
%    \begin{macrocode}
\DeclareRobustCommand\ubvr{\ensuremath{U\!BV\!R}}            
\DeclareRobustCommand\Color[2]{%
  \ensuremath{%
    #1\!\ifx#1V\!\fi-\!#2
  }}
\DeclareRobustCommand\ub{\Color{U}{B}}
\DeclareRobustCommand\bv{\Color{B}{V}}
\DeclareRobustCommand\vr{\Color{V}{R}}
\DeclareRobustCommand\ur{\Color{U}{R}}
%% New version of the \ion command, that actually works! 
%    \end{macrocode}
% \begin{macro}{\ion}
% This version of the |\ion| command allows you to write either
% |\ion{N}{4}| or |\ion{N}{iv}| in order to get N~IV. 
%    \begin{macrocode}
\DeclareRobustCommand\ion[2]{%
  \mbox{#1\kern0.2em%
    \smaller\rmfamily%
    \edef\@tempa{\@car#2\@nil}%
    \ifcat1\@tempa%
    \@Roman{#2}%
    \else%
    \uppercase{#2}%
    \fi}}
\newcounter{IonStage}
\DeclareRobustCommand\plainion[2]{\setcounter{IonStage}{#2}#1 \Roman{IonStage}}
%%\def\ion#1#2{\ensuremath{\mathrm{#1}\;}{\protect\small\rm\@Roman{#2}}}
%    \end{macrocode}
% \end{macro}
%    \begin{macrocode}
\DeclareRobustCommand\nodata{$\cdots$}
%    \end{macrocode}
% These are needed by the |\ion| command. Isn't there an easier way of
% doing this?
%    \begin{macrocode}
\let\@small=\small
\let\@normalsize=\normalsize
\let\@footnotesize=\footnotesize
\let\@scriptsize=\scriptsize
\let\@tiny=\tiny
\let\@large=\large
\let\@Large=\Large
\let\@LARGE=\LARGE
\let\@huge=\huge
\let\@Huge=\Huge
\newcommand{\smaller}{\small}
\renewcommand{\normalsize}{\renewcommand{\smaller}{\small}\@normalsize}
\renewcommand{\small}{\renewcommand{\smaller}{\footnotesize}\@small}
\renewcommand{\footnotesize}{\renewcommand{\smaller}{\scriptsize}\@footnotesize}
\renewcommand{\scriptsize}{\renewcommand{\smaller}{\tiny}\@scriptsize}
\renewcommand{\tiny}{\renewcommand{\smaller}{\tiny}\@tiny}
\renewcommand{\large}{\renewcommand{\smaller}{\normalsize}\@large}  
\renewcommand{\Large}{\renewcommand{\smaller}{\large}\@Large}  
\renewcommand{\LARGE}{\renewcommand{\smaller}{\Large}\@LARGE}  
\renewcommand{\huge}{\renewcommand{\smaller}{\LARGE}\@huge}  
\renewcommand{\Huge}{\renewcommand{\smaller}{\huge}\@Huge}  
%</package|adshtml-style>
%    \end{macrocode}


% \section{Other files}

% \subsection{The driver and make files}
% \label{sec:driver_make}
% These are put here for completeness, although they still need to be
% distributed separately since otherwise the user won't be able to
% unpack this file into its constituent parts.
% 
% The |Makefile| has two principal rules: |docs|, and |macros|. The
% |docs| rule produces a PDF and gzipped PS file of the documentation.
% I still haven't perfected the combination of rules required to get
% it to run \LaTeX{} enough times in all circumstances. The |macros|
% rule runs \LaTeX{} on |rmaa.ins| in order to unpack |rmaa.cls| and
% all the other little files (including the |Makefile| and the
% |rmaa.ins| files themselves\dots 
% 
% First set up lists of various types of file. 
%    \begin{macrocode}
%<*makefile>
DOCFILES=rmaa.pdf rmaa.ps.gz
MACRO_LTX_FILES=rmaa.cls rmaa.ins rmaacite.sty adshtml.sty ADS.tex
SHELLSCRIPTS=makeads.sh
MACRO_OTHER_FILES=adshtml.perl .latex2html-init
MAKEFILE=Makefile
MACROFILES=$(MACRO_LTX_FILES) $(SHELLSCRIPTS) $(MACRO_OTHER_FILES)         
SKELETONFILES=ADS.tex 
%    \end{macrocode}
% Install the files read only in the parent directory (idea is that
% the source files are in a subdir of the installation dir).
%    \begin{macrocode}
INSTALL=install 
INSTALL_DIR=..
%    \end{macrocode}
% The default rule is to make everything. Installation must be
% explicitly asked for.
%    \begin{macrocode}
all: docs macros

install: install-docs install-macros

%    \end{macrocode}
% Most files are installed read-only to prevent editing. Skeleton
% files must be writeable. Must be careful not to overwrite the
% Makefile in the parent directory. 
%    \begin{macrocode}
install-macros: macros
        $(INSTALL) -m 444 $(MACRO_LTX_FILES) $(MACRO_OTHER_FILES) $(INSTALL_DIR)
        $(INSTALL) -m 555 $(SHELLSCRIPTS) $(INSTALL_DIR)
        $(INSTALL) -m 644 $(SKELETONFILES) $(INSTALL_DIR)

install-docs: docs
        $(INSTALL) rmaa.pdf $(INSTALL_DIR)/rm-fulldocs.pdf

%    \end{macrocode}
% Make the PDF and gzipped PS files. 
%    \begin{macrocode}
docs: $(DOCFILES)

%    \end{macrocode}
% Make everything else if either the |.dtx| or |.ins| files have
% changed. 
%    \begin{macrocode}
macros: rmaa.dtx rmaa.ins
%    \end{macrocode}
% Ensure that all the macro files can be written to then recreate
% them. The |touch| is there in case any of the files don't exist
% yet. 
%    \begin{macrocode}
        touch $(MACROFILES) $(MAKEFILE)
        chmod u+w $(MACROFILES) $(MAKEFILE)
        latex "%%
%% This is file `rmaa.ins',
%% generated with the docstrip utility.
%%
%% The original source files were:
%%
%% rmaa.dtx  (with options: `ins')
%% 
%% PLEASE DO NOT EDIT THIS FILE! It is generated automatically from
%% the file rmaa.dtx using the accompanying Makefile and rmaa.ins
%% Any hand-edited changes will be lost
%% 
%% A LaTeX document class for the RMxAA
%% 
%% Author: William Henney (w.henney@crya.unam.mx)
%% 
%% Latest version available from: http://www.crya.unam.mx/rmaa
%% 
%% Copyright (C) 1999-2010 Instituto de Astronomia, UNAM, Mexico
%% 
%% This program is free software; you can redistribute it and/or
%% modify it under the terms of the GNU General Public License
%% as published by the Free Software Foundation; either version 2
%% of the License, or (at your option) any later version.
%% 
%% This program is distributed in the hope that it will be useful,
%% but WITHOUT ANY WARRANTY; without even the implied warranty of
%% MERCHANTABILITY or FITNESS FOR A PARTICULAR PURPOSE.  See the
%% GNU General Public License for more details.
%% 
%% You should have received a copy of the GNU General Public License
%% along with this program; if not, write to the Free Software
%% Foundation, Inc., 59 Temple Place - Suite 330, Boston, MA  02111-1307, USA.





%%
%% Installation batch file for the RMAA LaTeX package
%%
\def\batchfile{rmaa.ins}
\input docstrip.tex
\preamble

PLEASE DO NOT EDIT THIS FILE! It is generated automatically from
the file rmaa.dtx using the accompanying Makefile and rmaa.ins
Any hand-edited changes will be lost

A LaTeX document class for the RMxAA

Author: William Henney (w.henney@crya.unam.mx)

Latest version available from: http://www.crya.unam.mx/rmaa

Copyright (C) 1999-2010 Instituto de Astronomia, UNAM, Mexico

This program is free software; you can redistribute it and/or
modify it under the terms of the GNU General Public License
as published by the Free Software Foundation; either version 2
of the License, or (at your option) any later version.

This program is distributed in the hope that it will be useful,
but WITHOUT ANY WARRANTY; without even the implied warranty of
MERCHANTABILITY or FITNESS FOR A PARTICULAR PURPOSE.  See the
GNU General Public License for more details.

You should have received a copy of the GNU General Public License
along with this program; if not, write to the Free Software
Foundation, Inc., 59 Temple Place - Suite 330, Boston, MA  02111-1307, USA.
\endpreamble
\askforoverwritefalse
\generate{%
  \file{rmaa.cls}{\from{rmaa.dtx}{package}}
  \file{rmaa.ins}{\from{rmaa.dtx}{ins}}
  \file{rmaacite.sty}{\from{rmaa.dtx}{rmaacite}}
  \file{adshtml.sty}{\from{rmaa.dtx}{adshtml-style}}
  \file{ADS.tex}{\from{rmaa.dtx}{ads-container}}
  }
\generate{\nopreamble\nopostamble%
  \file{Makefile}{\from{rmaa.dtx}{makefile}}
  \file{adshtml.perl}{\from{rmaa.dtx}{adshtml-perl}}
  \file{makeads.sh}{\from{rmaa.dtx}{ads-shell-script}}
  \file{sample.latex2html-init}{\from{rmaa.dtx}{latex2html-init}}
  }

%% Local Variables:
%% mode: LaTeX
%% mode: Font-lock
%% End:









\endinput
%%
%% End of file `rmaa.ins'.
"
%    \end{macrocode}
% \TeX{} automatically adds the |.tex| extension to extensionless
% files, bless it. We fix this and take the opportunity to convert
% back to TABs all that it has lovingly converted to spaces. 
%    \begin{macrocode}
        cp -f $(MAKEFILE) $(MAKEFILE).BCK && unexpand $(MAKEFILE).tex >| $(MAKEFILE)
%    \end{macrocode}
% Make the shell scripts executable and everything unwriteable (so
% they don't get hand-edited by mistake). However, undo this for the
% files that the user is \emph{supposed} to modify. 
%    \begin{macrocode}
        chmod a+x $(SHELLSCRIPTS)
        chmod a-w $(MACROFILES) $(MAKEFILE)
        chmod u+w $(SKELETONFILES)

%    \end{macrocode}
% We now use PDF\LaTeX{} for this.
%    \begin{macrocode}
rmaa.pdf: rmaa.aux rmaa.ind rmaa.gls
        @echo "REGENERATING .pdf FILE"
        pdflatex "\nonstopmode% \def\fileversion{4.0} \def\filedate{2019/11/12}
% \MakeShortVerb{\|}
% \title{A \LaTeX{} document class for the Revista Mexicana de
% Astronom\'\i a y Astrof\'\i sica\thanks{This file has version
% \fileversion{} last revised \filedate{}.}}
% \author{William J. Henney\thanks{w.henney@irya.unam.mx}}
% \date{\filedate}
% \maketitle
% \paragraph{Abstract}
% This document describes the \texttt{rmaa} \LaTeX{}
% document class for the journal ``Revista Mexicana de
% Astronom\'\i a y Astrof\'\i sica''. The macros can be used for 
% \begin{enumerate}
% \item Preparing author manuscripts for the journal (including
% conference proceedings to be published in the ``Serie de
% Conferencias'')
% \item Typesetting an entire volume of the journal by inputting the
% individual article files from a container file. Facilities are
% included for automatic generation of TOC and Author Index. 
% \end{enumerate}
% 
% \tableofcontents
% \section{The documentation driver file}
% This is the part that generates the documentation that you may or
% may not be reading, which is done with the help of the Makefile
% discussed in section~\ref{sec:driver_make}. 
%    \begin{macrocode}
%<*driver>
\documentclass{ltxdoc}
\usepackage{hyperref}
\usepackage{palatino}
\EnableCrossrefs \CodelineIndex \RecordChanges
\setlength{\parindent}{0pt}
\def\MacroFont{\footnotesize\tt}
\begin{document} \DocInput{rmaa.dtx} \PrintIndex \PrintChanges \end{document}
%</driver>
%    \end{macrocode}

% \section{The implementation}
% 
% \changes{3.24}{2007/09/14}{Added user access to the macro version
% and date via \cs{rmVersion} and \cs{rmDate}}
% 
%    \begin{macrocode}
%<*package>
\newcommand{\rmaa@version}{3.26}
\newcommand{\rmaa@date}{2011/08/15}
\newcommand\rmVersion{\rmaa@version}
\newcommand\rmDate{\rmaa@date}
\NeedsTeXFormat{LaTeX2e}[1995/06/01]
%    \end{macrocode}
% \changes{4.0}{2019/11/02}{Change abstract to appear before
% resumen. Change RevMexAA to RMxAA for consistency with ADS.}
%
% \changes{3.22}{2007/03/22}{Allow proceedings articles to have
%   received/accepted dates} \changes{3.21}{2006/05/15}{Changes for
%   Silvia - abstracts are now shipped out to ADS.}
% \changes{3.19}{2006/01/19}{Changes for Silvia - remove adresses from
%   end, new abstract style} \changes{2.2beta}{2001/10/30}{ Fixed bug
%   whereby twocolumn option was not passed through to book
%   class. Finally sorted out the ``label on last page'' issue. }
% \changes{2.1}{2001/01/23}{ Fixed some inconsistencies in
%   \cs{nocite} and \cs{bibitem} uniqification Fixed the definition of
%   \cs{ion} so it will work in math mode and in contexts where the
%   typesize is different from normal (e.g., title, fig captions).
%   Modified the definition of \cs{@caption}, so we don't write any
%   list-of-figures data (we never use it anyway, and it just causes
%   problems with fragile commands in fig captions). }
% \changes{2.0}{2000/11/30}{ Added `newstyle' option, which should be
%   used in all new proceedings volumes (wider abstract, smaller
%   typeface in fig captions, etc.) Added commands for formatting
%   table captions. New demonstation document rmpostsamp.tex for
%   1-page posters. This will be the last release under a PRI
%   government (for a while...)}  \changes{1.9}{2000/06/06}{
%   Corrected position of page number in footer of articles that open
%   on even pages. Ensured that first article in each session opens on
%   an odd page, as do list of abstracts and author index. Added extra
%   space in TOC before the backmatter. Added command for editors,
%   \cs{BoldSessionTitles}, that makes produces bold titles for
%   session in the TOC.}  \changes{1.8}{2000/05/16}{ Fixed bug in
%   last page label with \cs{suppressfulladdresses} }
% \changes{1.7}{2000/05/11}{ Fixed more little things. Page labels
%   now work, even if you have no bibliography. You can output the
%   full addresses by hand with the \cs{outputfulladdresses}
%   command. Editors' Guide updated but still not finished. Revamped
%   splitbook.sh now produces PS file of each article for web page and
%   also an EPS file of each page for sending to the printers.}
% \changes{1.6}{2000/04/15}{ More cleaning up: full stops after
%   titles in index, running headers and footers in frontmatter,
%   fontsize in title of index, etc} \changes{1.5}{2000/04/15}{
%   Eliminated superscripts from titles in TOC and index}
% \changes{1.4}{2000/04/14}{ Added copyright and license info (GPL)}
% \changes{1.3}{2000/03/01}{ Corrected \cs{figcaption} to do the
%   caption formatting right Corrected font size of sec headings in
%   front matter} \changes{1.2}{2000/02/16}{ Added \cs{plate}
%   command} \changes{1.1}{2000/02/14}{ Added manuscript option,
%   corrected bugs} \changes{1.0}{2000/02/08}{ Now does a whole book
%   - lots of stuff tidied up} \changes{0.7}{1999/10/04}{ Increased
%   spacing between sections in single column styles. Suppressed
%   floats at bottom of pages (didn't work though) }
% \changes{0.6}{1999/09/24}{ Tweaked \cs{adjustfinalcols} (still not
%   perfect) Added proceedings option Added onecolumn/twocolumn
%   options (orthogonal to journal/proceedings) }
% \changes{0.5}{1999/09/23}{ Added \cs{fulladdresses} and
%   \cs{adjustfinalcols} (just about working)}
% \changes{0.4}{1999/09/22}{ Minor stylistic changes}
% \changes{0.3}{1999/06/28}{ Better mechanism for alternative
%   affiliations Added `guide' option for user guide}
% \changes{0.2}{1999/05/27}{ All basic functionality present, I
%   think} \changes{0.1}{1999/05/26}{ Just about working - still
%   missing Fig/Tab/Refs} \changes{0.0}{1999/05/21}{ So bored today,
%   decided to rewrite RMAA LaTeX macros from scratch}
%    \begin{macrocode}
\ProvidesClass{rmaa}[\rmaa@date\space v\rmaa@version\space RMxAA article
class]
%    \end{macrocode}
% \subsection{Option Handling}
% \label{sec:options}
% Declarations that are needed in options
% \changes{3.05}{2003/10/22}{Added natbib option}
% \changes{3.06}{2003/01/03}{Natbib option turned on permanently}
% \changes{3.07}{2003/01/11}{Added proof option}
% \changes{3.14}{2003/01/11}{Option to turn off natbib}
% \changes{3.18}{2003/10/22}{trying out doctex mode} 
% \changes{3.23}{2007/09/09}{Manuscript is now on by default.}
% \changes{3.26}{2011/08/15}{New option for review articles}
%    \begin{macrocode}
\newif\if@reviewarticle
\@reviewarticlefalse
\newif\if@proof
\@prooffalse
\newif\if@journal
\@journaltrue
\newif\if@printdates
\@printdatesfalse
\newif\if@proceedings
\@proceedingsfalse
\newif\if@manuscript
\@manuscripttrue
\newif\if@bigfont
\@bigfonttrue
\newif\if@guide
\@guidefalse
\newif\if@preprint
\@preprintfalse
\newif\if@rmtwocol
\@rmtwocoltrue
\newif\if@biblist
\@biblistfalse
\newif\if@book
\@bookfalse
\newif\if@natbib
\@natbibtrue
\newif\if@rmcite
\@rmcitefalse
\newif\if@titlevisible
\@titlevisibletrue
\newif\if@abstractpresent
\@abstractpresentfalse
\newif\if@printbanner
\@printbannerfalse
\newif\if@suppressfulladdresses
\@suppressfulladdressesfalse
\newif\if@newstyle
\@newstylefalse
\newif\if@debug
\@debugfalse
\newif\if@boldkeywords
\@boldkeywordsfalse
%    \end{macrocode}
% The option handling---mainly setting various booleans.
% \changes{2.2beta11}{2001/11/26}{Made twocol and
% newstyle the default for proceedings.}
% \changes{3.23}{2007/09/09}{Manuscript is now 10 pt instead of 12pt}
% \changes{3.23}{2007/09/10}{Full addresses are now automatically
% suppressed for proceedings style.}
% \changes{3.24}{2007/09/14}{New option \texttt{boldkeywords} that
% does the keywords in the old ugly way. Now off by default.}
%    \begin{macrocode}
\DeclareOption{manuscript}{\typeout{Manuscript for referee} 
  \@manuscripttrue\@journalfalse\@rmtwocolfalse
  \PassOptionsToClass{10pt}{article}}
\DeclareOption{bigfont}{\typeout{Big font manuscript requested by Gary
  Ferland} 
  \@manuscripttrue\@journalfalse\@rmtwocolfalse\@bigfonttrue
  \PassOptionsToClass{11pt}{article}}
\DeclareOption{journal}{\typeout{Main journal option}%
  \@journaltrue\@rmtwocoltrue\@printdatestrue\@manuscriptfalse
  \PassOptionsToClass{10pt}{article}}
\DeclareOption{guide}{\typeout{User guide option}%
  \@guidetrue\@journalfalse\@rmtwocoltrue\@manuscriptfalse
  \PassOptionsToClass{10pt}{article}}
\DeclareOption{preprint}{\typeout{Preprint option}%
  \@preprinttrue\@manuscriptfalse\PassOptionsToClass{10pt}{article}}
\DeclareOption{proof}{\typeout{Author proof option}%
  \@prooftrue\@manuscriptfalse}
\DeclareOption{proceedings}{\typeout{Proceedings option}%
  \@proceedingstrue\@journalfalse\@rmtwocoltrue\@newstyletrue
  \@suppressfulladdressestrue
  \@printdatesfalse\@manuscriptfalse\PassOptionsToClass{10pt}{article}} 
\DeclareOption{printdates}{\typeout{Force printing of received/accepted dates}%
  \@printdatestrue} 
\DeclareOption{book}{\typeout{Whole issue option}%
  \@booktrue} 
\DeclareOption{nonatbib}{\typeout{Turn off natbib support (required
    for compatibility with old rmaacite package).}%
  \@natbibfalse} 
\DeclareOption{newstyle}{\typeout{Use new layout style}%
  \@newstyletrue}
\DeclareOption{oldstyle}{\typeout{Use old layout style}%
  \@newstylefalse}
\DeclareOption{debug}{\typeout{Loads of annoying messages}%
  \@debugtrue}
%    \end{macrocode}
% These two need to be declared explicitly so that |\ExecuteOptions|
% actually does something.
%    \begin{macrocode}
\DeclareOption{twocolumn}{\@rmtwocoltrue}
\DeclareOption{onecolumn}{\@rmtwocolfalse}
\DeclareOption{twoside}{\PassOptionsToClass{\CurrentOption}{article}}
\DeclareOption{boldkeywords}{\@boldkeywordstrue}
%    \end{macrocode}
% Bold, upper-case keywords used to be the default. Now they must be
% turned on explicitly with the \texttt{boldkeywords} option. 
%    \begin{macrocode}
\DeclareOption{boldkeywords}{\@boldkeywordstrue}
%    \end{macrocode}
% Pass any other option along to the \texttt{article} class.
%    \begin{macrocode}
\DeclareOption*{\PassOptionsToClass{\CurrentOption}{article}}
%    \end{macrocode}
% Default options are to be two-sided and based on the standard
% \LaTeX{} \texttt{article} class,
% unless we are processing the entire book, in which case we base
% ourselves on the standard \texttt{book} class. Default is to use
% manuscript option unless otherwise specified. 
%    \begin{macrocode}
\ExecuteOptions{twoside,manuscript} \ProcessOptions\relax
\if@book
\if@rmtwocol\LoadClass[twocolumn]{book}\else\LoadClass{book}\fi
\newcommand{\abstract}{\relax}
\newcommand{\@maketitle}{\relax}
\else
\if@rmtwocol\LoadClass[twocolumn]{article}\else\LoadClass{article}\fi
\fi
%    \end{macrocode}
% Attempt to use natbib for citations 02 Jan 2003. 
% \changes{3.05}{2003/01/02}{If the natbib option is used, then set
% things up so that natbib's version of the bibliography is identical
% to the standard one.}
% \changes{3.09}{2003/01/23}{Explicitly require a recent version of
% natbib, so that we get a decent error message if the package is too
% old.} 
% \changes{3.23}{2007/09/09}{Now allow Bibtex to work properly} 
% \changes{3.25}{2009/10/01}{Change \cs{newcommand} tp \cs{def} when setting \cs{bibfont}. This allows it to work with all versions of natbib.} 
% \changes{3.25}{2010/02/06}{Updated copyright year.} 
% \changes{3.25}{2010/02/06}{Updated email and web addresses to \texttt{crya.unam.mx}.} 
%    \begin{macrocode}
\if@natbib
\RequirePackage{natbib}[1999/02/25]
\bibpunct{(}{)}{;}{a}{}{,}
\def\bibfont{\small}
\renewcommand\bibsection{\section*{REFERENCES}}
\setlength\bibsep{0pt}
\setlength\bibhang{1.5em}
\bibliographystyle{rmaa}
\fi
%    \end{macrocode}
% We want the \texttt{graphicx} package for the extended syntax of
% |\includegraphics|. We want the \texttt{ifthen} package for using in
% these macros (although we don't much yet). We explicitly load
% \texttt{babel} since some \LaTeX{} distributions messed up the
% hyphenation if we didn't.  
%    \begin{macrocode}
\RequirePackage{graphicx}
\RequirePackage{ifthen}
\IfFileExists{babel.sty}{%
%   \let\@restoreindent=\@afterindentfalse
  \RequirePackage[english]{babel}
%   \@restoreindent
  %% reinstate the definition of \@roman so as not to break .idx files
  \@ifundefined{org@@roman}{\relax}{\renewcommand{\@roman}[1]{\org@@roman{##1}}}
  }{%
  \typeout{Package babel.sty not found - don't worry about this}}
%    \end{macrocode}
% Try to load the \texttt{dcolumn} and \texttt{booktabs} packages if
% they can be found. These can produce nicer-looking tables. 
%    \begin{macrocode}
\IfFileExists{dcolumn.sty}{%
  \RequirePackage{dcolumn}
  }{%
  \typeout{Package dcolumn.sty not found - you will not be able to
    align columns on decimal points in tables}}
\IfFileExists{booktabs.sty}{%
  \RequirePackage{booktabs}
  }{%
  \typeout{Package booktabs.sty not found - your tables will not look
    as pretty as they ought too.}
  }
\IfFileExists{makeidx.sty}{%
  \RequirePackage{makeidx}
}{%
  \typeout{Package makeidx.sty not found - Author Index may not work.}
  }
%% commented out for moment since requires more modern version of LaTeX
%% \RequirePackage[latin1]{inputenc}
%    \end{macrocode}
% \changes{4.0}{2019/11/13}{Make \texttt{hyperref} and
% \texttt{etoolbox} into hard requirements, since we now use them
% extensively.}  We use hyperref to get internal and external links.
% But I don't like the colored boxes so I turn those off.  We use
% etoolbox for patching and robustifying at the moment.  In the
% future, I want to port a lot of other stuff over to it.  For
% instance, replace the ifthen package and the TeX ifs too.
%
% We target the latest version (1.7) of PDF, rather than the default
% of 1.5.  This reduces the chance of getting warnings from included
% graphics.  This is not strictly related to hyperref, but we put it
% here anyway.
%    \begin{macrocode}
\RequirePackage{etoolbox}
\RequirePackage{hyperref}
\hypersetup{hidelinks}
\pdfminorversion=7
%    \end{macrocode}
% We patch the \cs{@combinedblfloats} following a suggestion by Heiko
% Oberdiek at \url{https://tex.stackexchange.com/a/249743/3843}, which
% avoids problems when hyperlinks are split across
% pages. \changes{4.0}{2019/11/12}{Fix for issue where breaking
% hyperlink across pages gave an error.}
%    \begin{macrocode}
\urlstyle{same}
\newcommand\email{\begingroup \urlstyle{rm}\Url}
\patchcmd\@combinedblfloats{\box\@outputbox}{\unvbox\@outputbox}{}{%
  \errmessage{\noexpand\@combinedblfloats could not be patched}%
}%
%    \end{macrocode}
% \subsection{Other setup stuff}\label{sec:setup}
% Some declarations: 
%    \begin{macrocode}
\newif\if@needaddresses
\@needaddressestrue
\newif\if@frontmatter
\@frontmatterfalse
\newif\if@inpreamble
\@inpreambletrue
\newif\if@firstauthor
\@firstauthortrue
\newif\if@firstaffil
\@firstaffiltrue
\newif\if@firstaltaffil
\@firstaltaffiltrue
\newif\if@firstkwd
\@firstkwdtrue
\newif\if@noexpandADS
\@noexpandADSfalse
\newif\if@sessiontitlepages
\@sessiontitlepagesfalse
%    \end{macrocode}
% \begin{macro}{\@rmaaEmpty}
% \begin{macro}{\@rmaaArg}
% These are used in various commands to check if arg was left blank
%    \begin{macrocode}
\newcommand{\@rmaaEmpty}{}
\newcommand{\@rmaaArg}{}
%    \end{macrocode}  
% \end{macro}
% \end{macro}
%
% \subsection{Compatibility with old \LaTeX{} versions}
% \label{sec:compat}
% This package is based on the 1998/12/01 release of \LaTeX{}.
% Older releases may not define all the commands that we need.
% Luckily, we can use |\providecommand| to define these if necessary. 
% The following commands are not defined in the 1994/06/01 release. 
%
% \begin{macro}{\MakeUpperCase}
% \begin{macro}{\textsuperscript}
% \begin{macro}{\textbackslash}
% User commands:
%    \begin{macrocode}
\providecommand{\MakeUpperCase}[1]{\uppercase{#1}}
\providecommand{\textsuperscript}[1]{\ensuremath{^\textrm{#1}}}
\providecommand{\textbackslash}{\ensuremath{\backslash}}
%    \end{macrocode}  
% \end{macro}
% \end{macro}
% \end{macro}
% \begin{macro}{\protected@xdef}
% \begin{macro}{\restore@protect}
% \begin{macro}{\@unexpandable@protect}
% \begin{macro}{\hb@xt@}
% \LaTeX{} internal commands:
%    \begin{macrocode}
\providecommand{\protected@xdef}{%
        \let\@@protect\protect
        \let\protect\@unexpandable@protect
        \afterassignment\restore@protect
        \xdef}
\providecommand{\restore@protect}{\let\protect\@@protect}
\providecommand{\@unexpandable@protect}{\noexpand\protect\noexpand}
\providecommand{\hb@xt@}{\hbox to}
%    \end{macrocode}
% \end{macro}
% \end{macro}
% \end{macro}
% \end{macro}
%
% \subsection{Page layout}\label{sec:page-layout}
%
% The dimensions of the text were defined by Elizabeth. Text width is
% 169\,mm with equal margins on the two sides: 
%    \begin{macrocode}
\setlength{\textwidth}{169mm}
\setlength{\@tempdima}{\paperwidth}
\addtolength{\@tempdima}{-\textwidth}
\setlength{\oddsidemargin}{0.5\@tempdima}
\addtolength{\oddsidemargin}{-1in}
\setlength{\evensidemargin}{\oddsidemargin}
%    \end{macrocode}
% Text height is 225\,mm, top margin is 14\,mm. 
%    \begin{macrocode}
\setlength{\textheight}{225mm} 
\setlength{\topmargin}{-1in} 
\addtolength{\topmargin}{14mm}
%    \end{macrocode}
% Space between columns in two-column styles, header, footer,
% paragraph indentation:
%    \begin{macrocode}
\setlength{\columnsep}{9.5mm}
\setlength{\headsep}{0.7cm}
\setlength{\footskip}{0.75cm}
\setlength{\parindent}{1.5em}
%    \end{macrocode}
% 
% \begin{macro}{\@set@scalable@length}
% Try out a new scheme for dealing with various spacing dimensions and
% allowing the user to easily modify them. The command
% \begin{quote}
% |\set@scalable@length|\marg{name}\marg{defaultlength}\marg{stretch}\marg{shrink}
% \end{quote}
% does a 
% |\setlength|\marg{\cs{name}}\marg{defaultlength plus stretch minus shrink}
% and also defines the user command
% \textbackslash\meta{name}|Scale|\marg{factor} that can 
% be used to rescale length by \meta{factor}. Additionally, now
% defines the package command
% \textbackslash\meta{name}|@set@scale|\marg{FACTOR} 
% that also rescales the length by \meta{FACTOR}, so that the length
% becomes \meta{factor} $\times$ \meta{FACTOR} $\times$
% \meta{defaultlength}.   
%    \begin{macrocode}
\newcommand\@set@scalable@length[4]{%
  %% arg #1: NAME of length without leading `\'
  %%
  %% arg #2: default VALUE  
  \expandafter\newlength\csname @#1@default\endcsname
  \expandafter\setlength\csname @#1@default\endcsname{#2}%
  \expandafter\newlength\csname @@#1@@default\endcsname
  \expandafter\setlength\csname @@#1@@default\endcsname{#2}%
  %% arg #3: default STRETCH  
  \expandafter\newlength\csname @#1@stretch\endcsname
  \expandafter\setlength\csname @#1@stretch\endcsname{#3}%
  \expandafter\newlength\csname @@#1@@stretch\endcsname
  \expandafter\setlength\csname @@#1@@stretch\endcsname{#3}%
  %% arg #4: default SHRINK  
  \expandafter\newlength\csname @#1@shrink\endcsname
  \expandafter\setlength\csname @#1@shrink\endcsname{#4}%
  \expandafter\newlength\csname @@#1@@shrink\endcsname
  \expandafter\setlength\csname @@#1@@shrink\endcsname{#4}%
  %%
  %% User-defined scale \@NAME@scale is unity by default
  \expandafter\gdef\csname @#1@scale\endcsname{1}%
  %% Package-defined scale \@NAME@SCALE is unity by default
  \expandafter\gdef\csname @#1@SCALE\endcsname{1}%
  %%
  %% internal command \@reset@NAME should expand to 
  %% \setlength\@@NAME@@default{\@NAME@SCALE\@NAME@default}
  %% \setlength\NAME{\@NAME@scale\@@NAME@@default}
  %% in other words: scale x SCALE x default 
  \expandafter\gdef\csname @reset@#1\endcsname{%
    \expandafter\setlength\csname @@#1@@default\endcsname{%
      \csname @#1@SCALE\endcsname \csname @#1@default\endcsname
      }
    \expandafter\setlength\csname @@#1@@stretch\endcsname{%
      \csname @#1@SCALE\endcsname \csname @#1@stretch\endcsname
      }
    \expandafter\setlength\csname @@#1@@shrink\endcsname{%
      \csname @#1@SCALE\endcsname \csname @#1@shrink\endcsname
      }
    \expandafter\setlength\csname #1\endcsname{%
      \csname @#1@scale\endcsname \csname @@#1@@default\endcsname
      \@plus
      \csname @#1@scale\endcsname \csname @@#1@@stretch\endcsname
      \@minus
      \csname @#1@scale\endcsname \csname @@#1@@shrink\endcsname
      }}
  %%
  %% User command \NAMEScale should expand to 
  %% {\renewcommand\@NAME@scale{#1}\@reset@NAME}
  \expandafter\gdef\csname #1Scale\endcsname##1{%
    \expandafter\gdef\csname @#1@scale\endcsname{##1}%
    \csname @reset@#1\endcsname}
  %%
  %% Package command \NAME@set@SCALE should expand to 
  %% {\renewcommand\@NAME@SCALE{#1}\@reset@NAME}
  \expandafter\gdef\csname #1@set@SCALE\endcsname##1{%
    \expandafter\gdef\csname @#1@SCALE\endcsname{##1}%
    \csname @reset@#1\endcsname}
  %%
  %% initialize at default values with \@reset@NAME
  \csname @reset@#1\endcsname
  }
%    \end{macrocode}
% \end{macro}
% \begin{macro}{\parskip}
% First use of this new infrastructure: set interparagraph space to be
% zero but with a small amount of stretchiness.
%    \begin{macrocode}
%   \ifthenelse{\boolean{@manuscript}}{%
%     \@set@scalable@length{parskip}{0.5\baselineskip}{0.3\p@}{0.1\p@}
%   }{
    \@set@scalable@length{parskip}{0\p@}{0.3\p@}{0\p@}
%   }
%    \end{macrocode}
% \end{macro}
% 
% \begin{macro}{\abovesectionskip}
% \changes{4.0}{2019/11/12}{Use negative values for skip so that first paragraph after section headings is not indented.  This is something that I have been trying to fix for years!  It turns out that the anser was in the documentation for the standard command \cs{@startsection.}}
% \begin{macro}{\belowsectionskip}
% \begin{macro}{\aboveparaskip}
% \begin{macro}{\belowparaskip}
% Define some utility rubber lengths for use in
% section~\ref{sec:section-commands}. 
%    \begin{macrocode}
\newlength{\abovesectionskip}
\newlength{\belowsectionskip}
\newlength{\aboveparaskip}
\newlength{\belowparaskip}
%    \end{macrocode}
% The default values for these lengths are those traditionally used
% for the onecolumn style. 
%    \begin{macrocode}
\@set@scalable@length{abovesectionskip}{-3.5ex}{-4ex}{-2ex}
\@set@scalable@length{belowsectionskip}{2.3ex}{0ex}{2ex}
\@set@scalable@length{aboveparaskip}{3.25ex}{2ex}{0.5ex}
\@set@scalable@length{belowparaskip}{-1em}{\z@}{\z@}
%    \end{macrocode}
% \end{macro}\end{macro}\end{macro}\end{macro}
% We also make scalable some of the standard float parameters. These
% are the same defaults as in the standard article classes, but
% slightly more stretchy. For the twocolumn style, they are currently
% halved, see below. Note that I have also lumped |\abovecaptionskip|
% in with these, since it plays a similar role. There is also a
% |\belowcaptionskip| but it is set to zero in the standard classes. 
%    \begin{macrocode}
\@set@scalable@length{floatsep}{12\p@}{4\p@}{4\p@}
\@set@scalable@length{textfloatsep}{20\p@}{4\p@}{8\p@}
\@set@scalable@length{intextsep}{12\p@}{4\p@}{4\p@}
\@set@scalable@length{dblfloatsep}{12\p@}{4\p@}{4\p@}
\@set@scalable@length{dbltextfloatsep}{20\p@}{4\p@}{8\p@}
\@set@scalable@length{abovecaptionskip}{10\p@}{\z@}{\z@}
%    \end{macrocode}
% \begin{macro}{\@before@title@skip}         
% \begin{macro}{\@before@author@skip}        
% \begin{macro}{\@between@author@skip}       
% \begin{macro}{\@before@abstract@skip}      
% \begin{macro}{\@after@abstracthead@skip}   
% \begin{macro}{\@between@abstract@skip}     
% \begin{macro}{\@before@kwd@skip}           
% \begin{macro}{\@before@dates@skip}           
% Spacing parameters used in |\@maketitle|. These are not changed by
% default except by the user, so they don't need to go into
% |\@package@default@lengths|.  
%    \begin{macrocode}
\newlength{\@before@title@skip}         
\newlength{\@before@author@skip}        
\newlength{\@between@author@skip}       
\newlength{\@before@abstract@skip}      
\newlength{\@after@abstracthead@skip}   
\newlength{\@between@abstract@skip}     
\newlength{\@before@kwd@skip}           
\newlength{\@before@dates@skip}           
\@set@scalable@length{@before@title@skip}{2em}{\z@}{\z@}
\@set@scalable@length{@before@author@skip}{1.5em}{\z@}{\z@}
\@set@scalable@length{@between@author@skip}{0.5em}{\z@}{\z@}
\@set@scalable@length{@before@abstract@skip}{3ex}{1ex}{0.5ex}
\@set@scalable@length{@after@abstracthead@skip}{1ex}{\z@}{\z@}
\@set@scalable@length{@between@abstract@skip}{2ex}{1ex}{0.5ex}
\@set@scalable@length{@before@kwd@skip}{1ex}{0.5ex}{0.2ex}
\@set@scalable@length{@before@dates@skip}{2ex}{0.5ex}{0.2ex}
%    \end{macrocode}
% \end{macro}\end{macro}\end{macro}\end{macro}
% \end{macro}\end{macro}\end{macro}\end{macro}
% 
% Spacing above and below displayed equations also wants to be done
% like this, but it is more complicated since the standard article
% classes make these change with the fontsize. Something to do later,
% maybe. 
% 
%    \begin{macrocode}
%    \end{macrocode}
% \begin{macro}{\@package@default@lengths}
% The value actually used for these lengths is controlled by the
% \meta{SCALE} parameter, which depends on the style option that is in
% effect. This is set using the |\@package@default@lengths| command,
% which must be called every time the style changes (e.g., changing
% from onecolumn to twocolumn). 
%    \begin{macrocode}
\newcommand{\@package@default@lengths}{%
  \ifthenelse{\boolean{@rmtwocol}}{%
    %% twocolumn mode
    \abovesectionskip@set@SCALE{0.5}
    \belowsectionskip@set@SCALE{0.5}
    \aboveparaskip@set@SCALE{0.5}
    \belowparaskip@set@SCALE{1}
    \floatsep@set@SCALE{0.5}
    \textfloatsep@set@SCALE{0.5}
    \intextsep@set@SCALE{0.5}
    \dblfloatsep@set@SCALE{0.5}
    \dbltextfloatsep@set@SCALE{0.5}
    }{%
    %% onecolumn mode
    \abovesectionskip@set@SCALE{1}
    \belowsectionskip@set@SCALE{1}
    \aboveparaskip@set@SCALE{1}
    \belowparaskip@set@SCALE{1}
    \floatsep@set@SCALE{1}
    \textfloatsep@set@SCALE{1}
    \intextsep@set@SCALE{1}
    \dblfloatsep@set@SCALE{1}
    \dbltextfloatsep@set@SCALE{1}
    }%
  }
\@package@default@lengths
%    \end{macrocode}
% \end{macro}
% \begin{macro}{\RescaleLengths}
% \begin{macro}{\RescaleSecLengths}
% \begin{macro}{\RescaleFloatLengths}
% \begin{macro}{\RescaleTitleLengths}
% User commands that will fiddle with all the scalable lengths. 
%    \begin{macrocode}
\newcommand{\RescaleLengths}[1]{%
  \RescaleSecLengths{#1}\RescaleFloatLengths{#1}\RescaleTitleLengths{#1}}
\newcommand{\RescaleSecLengths}[1]{%
%% sections
  \abovesectionskipScale{#1}
  \belowsectionskipScale{#1}
  \aboveparaskipScale{#1}
  }
\newcommand{\RescaleFloatLengths}[1]{%
%% floats
  \floatsepScale{#1}
  \textfloatsepScale{#1}
  \intextsepScale{#1}
  \dblfloatsepScale{#1}
  \dbltextfloatsepScale{#1}
  \abovecaptionskipScale{#1}
  }
\newcommand{\RescaleTitleLengths}[1]{%
%% maketitle
  \@before@title@skipScale{#1}           
  \@before@author@skipScale{#1}        
  \@between@author@skipScale{#1}       
  \@before@abstract@skipScale{#1}      
  \@after@abstracthead@skipScale{#1}   
  \@between@abstract@skipScale{#1}     
  \@before@kwd@skipScale{#1}           
  \@before@dates@skipScale{#1}           
  }
%    \end{macrocode}
% \end{macro}\end{macro}\end{macro}\end{macro}
% \changes{2.2beta14}{2002/03/17}{Finished implementation of new
% infrastructure for scalable lengths. Converted the above/below
% section skips to use it. Also, all the floatsep parameters. Also,
% all the vertical lengths that are used by |\@maketitle|. }

% \begin{macro}{\MSextramargin}
% For the \texttt{manuscript} option we have an extra 0.3\,inch
% vertical margin and 1\,inch horizontal margin for marginal notes (by
% hand or with |\marginpar|. 
%
% \changes{3.23}{2007/09/09}{Reduce baselinestretch in
%   \texttt{manuscript} style to 1.2. No longer use raggedright for
%   manuscript.}
%
% \changes{3.24}{2007/09/14}{Further reduction of baselinestretch in
%   \texttt{manuscript} style to 1.1 since we now have wider margins.}
%
% \changes{3.24}{2007/09/14}{Split \texttt{manuscript}'s extra margin
%   into separate horizontal and vertical parts.}
% 
%    \begin{macrocode}
\newlength{\MSextraHmargin}\newlength{\MSextraVmargin}
\setlength{\MSextraHmargin}{1in} 
\setlength{\MSextraVmargin}{0.3in} 
\if@manuscript
   \renewcommand{\baselinestretch}{1.1}
   \addtolength{\textwidth}{-2\MSextraHmargin}
   \addtolength{\oddsidemargin}{\MSextraHmargin}
   \setlength{\evensidemargin}{\oddsidemargin}
   \addtolength{\textheight}{-2\MSextraVmargin}
   \addtolength{\topmargin}{\MSextraVmargin}
%    \raggedright
\fi
%    \end{macrocode}
% \end{macro}
% Float parameters: don't allow bottom floats by default. 
%    \begin{macrocode}
\renewcommand{\fps@figure}{tp}
\renewcommand{\fps@table}{tp}
%    \end{macrocode}
%
% \begin{macro}{\rmaa@AuthorWidthFrac}
% \begin{macro}{\rmaa@AbstractWidthFrac}
% These parameters control what fraction of the total text width can
% be used by the list of authors and the abstract/resumen,
% respectively.  
%    \begin{macrocode}
\newcommand{\rmaa@AuthorWidthFrac}{1.0}
\newcommand{\rmaa@AbstractWidthFrac}{0.75}
%    \end{macrocode}
% Use a wide abstract to save space in the new style.
%    \begin{macrocode}
\if@newstyle\renewcommand{\rmaa@AbstractWidthFrac}{1.0}\fi
%    \end{macrocode}
% \end{macro}
% \end{macro}
%
% Some things to consider in future versions. Trying to avoid too many
% widows, etc. Apparently, the following will work best with
% \cs{flushbottom}. 
%    \begin{macrocode}
%%\widowpenalty=10000
%%\clubpenalty=10000
%    \end{macrocode}
% Also, should perhaps try getting rid of underfull/overfull box warnings. 
% \changes{3.10}{2003/01/24}{Set large values for \cs{hbadness},
% \cs{vbadness}, \cs{hfuzz}, \cs{vfuzz}. To reduce amount of stuff
% that gets written to the terminal every time you compile something.}
%    \begin{macrocode}
\hbadness=10000
\vbadness=10000
\setlength\hfuzz{\maxdimen}
\setlength\vfuzz{\maxdimen}
%    \end{macrocode}
% \subsection{Various text styles for different contexts}
% \label{sec:text-styles} 
% \begin{macro}{\rmaa@AffilStyle}
% \begin{macro}{\rmaa@PagesStyle}
% \begin{macro}{\rmaa@YearStyle}
% \begin{macro}{\rmaa@VolumeStyle}
% \begin{macro}{\rmaa@VolumenStyle}
% \begin{macro}{\rmaa@JournalNameStyle}
% \begin{macro}{\rmaa@ReceivedStyle}
% \begin{macro}{\rmaa@AbstractHeaderStyle}
% \begin{macro}{\rmaa@JournalNameStyle}
% Style commands for the different fields used by |\maketitle| in
% constructing the title page of the article. These have the form
% |\rmaa@XYZStyle{TEXT}|. I'm not sure these were really a good idea,
% but never mind \dots 
%    \begin{macrocode}
\newcommand{\rmaa@AffilStyle}[1]{{\normalsize\rmfamily#1}}
\newcommand{\rmaa@PagesStyle}[1]{{\normalsize\rmfamily#1}}
\newcommand{\rmaa@YearStyle}[1]{{\normalsize\rmfamily#1}}
\newcommand{\rmaa@VolumeStyle}[1]{{\normalsize\rmfamily\bfseries#1}}
\newcommand{\rmaa@VolumenStyle}[1]{{\Large\sffamily\bfseries#1}}
\newcommand{\rmaa@JournalNameStyle}[1]{{\normalsize\itshape#1}}
\newcommand{\rmaa@ReceivedStyle}[1]{{\normalsize\rmfamily\itshape#1}}
\newcommand{\rmaa@AbstractHeaderStyle}[1]{%
  {\normalsize\rmfamily\MakeUpperCase{#1}}}
\newcommand{\rmaa@TitleStyle}[1]{%
  {\large\bfseries\rmfamily\boldmath\MakeUpperCase{#1}\par}}
%    \end{macrocode}
% \end{macro}
% \end{macro}
% \end{macro}
% \end{macro}
% \end{macro}
% \end{macro}
% \end{macro}
% \end{macro}
% \end{macro}
%
% \begin{macro}{\rmaa@FrontSecStyle}
% Headings of ``front sections'' (see section~\ref{sec:front-sections})
%    \begin{macrocode}
\newcommand{\rmaa@FrontSecStyle}[1]{%
  {\normalsize\bfseries\rmfamily\MakeUpperCase{#1}}}
%    \end{macrocode}
% \end{macro}
%
% \begin{macro}{\rmaa@TOCTitleStyle}
% \begin{macro}{\rmaa@TOCPartStyle}
% \begin{macro}{\rmaa@IndexTitleStyle}
% \begin{macro}{\rmaaTOCAuthorStyle}
% \begin{macro}{\rmaaIndexAuthorStyle}
% Commands to set style of TOC and index entries.
% Rewritten to avoid unnecessary font commands in the .toc and .ind files
%    \begin{macrocode}
\newcommand{\rmaa@TOCTitleStyle}[1]{\MakeUpperCase{#1}}
\newcommand{\rmaa@TOCPartStyle}[1]{\rmaa@TOCTitleStyle{#1}}
\newcommand{\rmaa@IndexTitleStyle}[1]{#1}
\newcommand{\rmaaTOCAuthorStyle}[1]{%
  {\itshape#1}}
\newcommand{\rmaaIndexAuthorStyle}[1]{%
  \protect\mbox{\bfseries#1}}
%    \end{macrocode}
% \end{macro}
% \end{macro}
% \end{macro}
% \end{macro}
% \end{macro}
%
% \begin{macro}{\rmaa@KeywordStyle}
% \begin{macro}{\rmaa@KeywordHeaderStyle}
% The style of the keywords depends on whether we are using the
% \texttt{newstyle} option or not. If we are, then they are small. 
%    \begin{macrocode}
\ifthenelse{\boolean{@newstyle}}{
  \newcommand{\rmaa@KeywordStyle}[1]{%
    {\small%
      \if@boldkeywords 
      \bfseries\MakeUpperCase{\raggedright #1}
      \else
      \raggedright #1
      \fi
    }}
  \newcommand{\rmaa@KeywordHeaderStyle}[1]{%
    {\small\itshape#1}}}
{
  \newcommand{\rmaa@KeywordStyle}[1]{%
    {\normalsize%
      \if@boldkeywords\bfseries\fi\raggedright #1}}
  \newcommand{\rmaa@KeywordHeaderStyle}[1]{%
    {\normalsize\itshape#1}}}
%    \end{macrocode}
% \end{macro}
% \end{macro}
%
% \subsection{Page styles}\label{sec:page-styles}
% \begin{macro}{\ps@rmaa}
% The running headers for normal pages:
%    \begin{macrocode}
\newcommand{\ps@rmaa}{%
%    \end{macrocode}
% nothing in either foot; 
%    \begin{macrocode}
  \renewcommand\@oddfoot{\relax}%
  \renewcommand\@evenfoot{\relax}%
%    \end{macrocode}
% page number and short author list in the even heads; 
%    \begin{macrocode}
  \renewcommand\@evenhead{\rlap{\thepage}\hfil\@shortauthor\hfil}%
%    \end{macrocode}
% short title and page number in the odd heads.
%    \begin{macrocode}
  \renewcommand\@oddhead{\hfil\@shorttitle\hfil\llap{\thepage}}
  }
%    \end{macrocode}
% \begin{macro}{\ps@rmaafirstpage}
% The header of the first page of each
% article:
%    \begin{macrocode}
\newcommand{\ps@rmaafirstpage}{%
  \renewcommand\@oddfoot{\hfil\thepage}%
%    \end{macrocode}
% footer corrected (06 Jun 2000) to have page on outside edge;
%    \begin{macrocode}
  \renewcommand\@evenfoot{\thepage\hfil}%
%    \end{macrocode}
% Header has journal name, volume, and page numbers for a regular
% article, or similar for special cases. 
% \changes{4.0}{2019/11/03}{Change to three-line header with copyright info and DOI.}
%    \begin{macrocode}
\if@preprint
\renewcommand\@oddhead{%
  \if@proceedings 
  \textsc{To appear in ``\@conftitle{} (\rmaa@YearStyle{\@year})''}\hfil
  \textit{RMxAC}
  \else
  \textsc{To appear in} \@journalname\hfil
  \fi}\fi%
\if@guide
\renewcommand\@oddhead{\textsc{Documentation} (\rmaa@YearStyle{\@year})\hfil}
\fi%
\if@manuscript
\renewcommand\@oddhead{Manuscript for \@journalname{} (\rmaa@YearStyle{\@year})\hfil}
\fi%
\if@journal
\renewcommand\@oddhead{%
  \parbox{\textwidth}{
    \@journalname, \rmaa@VolumeStyle{\@volume},
    \rmaa@PagesStyle{\@artpages}\
    (\rmaa@YearStyle{\@year})\hfil\\[1pt]
    {\small\textcopyright{} \@year:
      Instituto de Astronom\'\i{}a,
      Universidad Nacional Aut\'onoma de M\'exico}\hfil\\[1pt]
    {\small\@DOI@url}\hfil
  }%
}\fi%
  \renewcommand\@evenhead{\@oddhead}}
%    \end{macrocode}
% By default, use the |\ps@rmaa| page style.
%    \begin{macrocode}
\pagestyle{rmaa}
%    \end{macrocode}
% \end{macro}
% \end{macro}
%
% \subsection{Preamble commands}\label{sec:preamble} 
% Everything you need for setting the title page, TOC and index
% entries, etc.
% \subsubsection{Commands for the title page}
% \label{sec:title-commands}
% \begin{macro}{\resumen}
% \begin{macro}{\@resumen}
% \begin{macro}{\abstract}
% \begin{macro}{\@abstract}
% User commands for setting the resumen and abstract. 
% These should be used in the document header. The internal commands
% |\@resumen| and |\@abstract| are used to hold the actual text
% pending output via the |\maketitle| command. 
%    \begin{macrocode}
\newif\if@setresumen
\newcommand{\resumen}[1]{%
%%  \renewcommand{\@resumen}{{\selectlanguage{spanish}#1}}%
  \renewcommand{\@resumen}{{\ignorespaces#1}}%
  \@setresumentrue}
\renewcommand{\abstract}[1]{%
  \renewcommand{\@abstract}{{\ignorespaces#1}}
  \renewcommand\ADS@abstract{#1}
%% Not a good idea, since then author macros are not expanded and
%% cause errors later. However, by not doing it we are vulnerable to
%% fragile commands.....
  \if@noexpandADS\@onelevel@sanitize{\ADS@abstract}\fi
  \@abstractpresenttrue
  }
\newcommand{\ADSnoexpand}{\@noexpandADStrue}
%    \end{macrocode}
% \begin{macro}{\boldabstract}
% This is the version for doing 1-page conference proceedings. Still
% save the text for ADS but then write it out immediately as a bold
% paragraph. 
%    \begin{macrocode}
\newcommand{\boldabstract}[1]{%
  \gdef\ADS@abstract{#1}%
  \renewcommand{\@abstract}{#1}%
  \if@noexpandADS\@onelevel@sanitize{\ADS@abstract}\fi%
  \noindent\textbf{\boldmath #1}
  }
%    \end{macrocode}
% \end{macro}
% Default value in each case is to print out a warning if not set
%    \begin{macrocode}
  \newcommand{\@resumen}{%
    \@setresumenfalse \ClassWarningNoLine{rmaa}{No Spanish abstract
      (resumen) specified - if you do not speak Spanish, this will be
      translated by the editors}%
    \textbf{Favor de proporcionar un resumen en espa\~nol. If you are
      unable to translate your abstract into Spanish, the editors will
      do it for you.} }
  \newcommand{\@abstract}{\ClassWarningNoLine{rmaa}{No abstract
      specified}}
%    \end{macrocode}
% \end{macro}
% \end{macro}
% \end{macro}
% \end{macro}
%
% \begin{macro}{\title}
% \begin{macro}{\@rawtitle}
% \begin{macro}{\@TOCtitle}
% \begin{macro}{\rmaaIndexTitle}
% Similarly, |\title| is the preamble command to set the title of the
% article. This is used in various places, typeset in different
% styles, hence the use of the internal |\@rawtitle|, |\@title|,
% |\rmaaIndexTitle|, and |\@TOCtitle| commands, each of which uses the
% relevant style from section~\ref{sec:text-styles}
%    \begin{macrocode}
\newcommand{\@rawtitle}{Raw Title}
\newcommand{\@TOCtitle}{\rmaa@TOCTitleStyle{Title for TOC}}
\newcommand{\rmaaIndexTitle}{\rmaa@IndexTitleStyle{Title for index}.}
\newcommand{\@TOCaftertitle}{}
\renewcommand{\title}[1]{%
%%  {\ADS@sanitize\xdef\@rawtitle{#1}}
  \renewcommand\@rawtitle{\noexpand#1}
  \let\ADS@title=\@rawtitle
  \renewcommand{\@title}{\rmaa@TitleStyle{#1}}%
  \renewcommand{\@TOCtitle}{\rmaa@TOCTitleStyle{#1}\@TOCaftertitle}%
  \renewcommand{\rmaaIndexTitle}{\rmaa@IndexTitleStyle{#1\@TOCaftertitle}.}%
  }
%    \end{macrocode}
% \end{macro}
% \end{macro}
% \end{macro}
% \end{macro}
% 
% All the weird and whacky ways of indicating the author
% affiliations. 
% \begin{macro}{\affil}
% First, the way that sets it directly underneath the authors name.
% Note that |\affil| trumps all other affiliations for the purposes of
% getting into the ADS stuff (we use |\gdef|).   
% \dots 
%    \begin{macrocode}
\newcommand{\affil}[1]{%
  \\[1.5ex]\rmaa@AffilStyle{#1}\\[1.5ex]%
  \if@firstaffil\gdef\ADS@affil{#1}\@firstaffilfalse\fi%
  }
\renewcommand{\and}{and\\[1.5ex]}
%    \end{macrocode}
% \begin{macro}{\altaffil}
% Alternative affiliations are footnoted. 
%    \begin{macrocode}
\newcommand{\altaffil}[1]{%
  \if@firstaltaffil\gdef\ADS@affil{#1}\@firstaltaffilfalse\fi%
  \thanks{#1}}
%    \end{macrocode}
% \begin{macro}{\altaffilmark}
% \begin{macro}{\altaffiltext}
% Alternative-alternative mechanism for affiliations---should not be
% mixed with |\altaffil| (unless you know what you are doing) or there
% will be problems with numbering.
%    \begin{macrocode}
\newcommand{\altaffilmark}[1]{\textsuperscript{#1}}
\newcommand{\altaffiltext}[2]{%
  \if@firstaltaffil\gdef\ADS@affil{#2}\@firstaltaffilfalse\fi
  \stepcounter{footnote}%
  \protected@xdef\@thanks{\@thanks
    \protect\footnotetext[#1]{#2}}} 
%    \end{macrocode}
% \end{macro}
% \end{macro}
% \end{macro}
% \end{macro}
%
% \begin{macro}{\keywords}
% \begin{macro}{\@keywords}
% \begin{macro}{\@keywordheader}
% \label{keywords}Sundry minor preamble commands. All
% self-explanatory, I hope. |\keywords| is the old user interface
% (deprecated). New articles should use |\addkeyword| once for each
% keyword. That way the ADS record comes out right too (comma rather
% than em-dash separator).
% \changes{2.2beta12}{2001/12/07}{\cs{addkeyword} New user interface
%   for keywords.} 
% \changes{2.2beta13}{2002/03/11}{\cs{addkeyword} Improved behaviour
%   with .ads file. Use
%   \cs{ADS@sanitize} to change definitions of `~' and \cs{ion}.}
% \changes{3.23}{2007/09/09}{\cs{addkeyword} now signals a compilation
% error if the author tries to add more than 6 keywords.}
% \changes{3.24}{2007/09/14}{kwdcounter is now reset in
% \cs{@begofarticlehook} to prevent error in book style} 
%    \begin{macrocode}
\newcommand{\@kwdsep}{ --- }
\newcommand{\ADS@kwdsep}{, }
\newcounter{kwdcounter}\setcounter{kwdcounter}{0}
\newcommand{\addkeyword}[1]{%
  \addtocounter{kwdcounter}{1}
  \ifthenelse{\thekwdcounter > 6}{
    \ClassError{rmaa}{Too many keywords. A maximum of 6 keywords is allowed.}{}
  }{}
  \if@firstkwd%
  \renewcommand\@keywords{\if@boldkeywords\uppercase{#1}\else #1\fi}
  {\ADS@sanitize\def\ion{\plainion}\gdef\ADS@keywords{#1}}
  \@firstkwdfalse
  \else%
  \edef\@keywords{\@keywords\@kwdsep \if@boldkeywords\uppercase{#1}\else #1\fi}%
  {\ADS@sanitize\def\ion{\plainion}\xdef\ADS@keywords{\ADS@keywords\ADS@kwdsep #1}}%
  \fi}
\newcommand{\ADS@sanitize}{%
  \def\mathrm{}
  \def\@{\ }
  \def~{\ }
  \def\langle{<}
  \def\rangle{>}
  \def\log{log }
  \def\sin{sin }
  \def\cos{cos }
  \def\tan{tan }
  \def\fd{^d.}
  \def\fh{^h.}
  \def\fm{^m.}
  \def\fs{^s.}
  \def\fdg{�.}
  \def\farcm{'.}
  \def\farcs{''.}
  \def\arcdeg{�}
  \def\,{ }
  }
\newcommand{\keywords}[1]{%
  \renewcommand{\ADS@keywords}{#1}
  \renewcommand{\@keywords}{\uppercase{#1}}
  }
\newcommand{\@keywordheader}{\rmaa@KeywordHeaderStyle{Key Words:}}
\newcommand{\@keywords}{%
  \ClassWarningNoLine{rmaa}{No keywords specified}%
  Keywords go here
  }
%    \end{macrocode}
% \begin{macro}{\shorttitle}
% \begin{macro}{\shortauthor}
% \begin{macro}{\@shorttitle}
% \begin{macro}{\@shortauthor}
% Commands for setting the running headers.
%    \begin{macrocode}
\newcommand{\shorttitle}[1]{%
  \renewcommand\@shorttitle{\MakeUpperCase{#1}}}
\newcommand{\shortauthor}[1]{%
  \renewcommand\@shortauthor{\MakeUpperCase{#1}}}
\newcommand{\@shorttitle}{SHORT TITLE GOES HERE}
\newcommand{\@shortauthor}{SHORT AUTHOR GOES HERE}
%    \end{macrocode}
% \end{macro}
% \end{macro}
% \end{macro}
% \end{macro}
% \end{macro}
% \end{macro}
% \end{macro}


% \subsubsection{The full addresses}
% \label{sec:full-addresses}
% \begin{macro}{\fulladdresses}
% \begin{macro}{\@fulladdresses}
% \begin{macro}{\@fulladdheight}
% \begin{macro}{\@fulladddepth}
% \begin{macro}{\@fulladdbox}
% Aaghh, the full addresses! A perennial source of problems. 
%    \begin{macrocode}
\newlength{\@fulladdheight}
\newlength{\@fulladddepth}
\newsavebox{\@fulladdbox}
%% \fulladdresses command modified (11 May 2000) so that we can put
%% \fulladdresses{} in the article to suppress the output of the
%% addresses
\newcommand{\@FinalFigure}{\relax}
\newcommand{\FinalFigure}[1]{\renewcommand{\@FinalFigure}{#1}}
\newcommand{\fulladdresses}[1]{%
  \renewcommand{\@rmaaArg}{#1}
  \ifx\@rmaaArg\@rmaaEmpty
     \renewcommand{\@fulladdresses}{}
  \else
     \renewcommand{\@fulladdresses}{
       \normalsize%
       \ifx\@FinalFigure\relax\else\vfill\@FinalFigure\vfill\fi
%%       \vspace*{1ex}
       %% do addresses as a list (copied from thebibliography environment) 
       \list{\null}{\leftmargin 1.5em\topsep\z@\labelwidth\z@\labelsep\z@\itemindent -1.5em
         \itemsep\z@\parsep\z@\usecounter{enumi}}
       \sloppy\clubpenalty4000\widowpenalty4000
       \sfcode`\.=1000\relax
       #1
       \endlist
       }
  \fi
  }
\newcommand{\@fulladdresses}{Full addresses go here}
%    \end{macrocode}
% \begin{macro}{\suppressfulladdresses}
% New user-level command for suppressing full addresses should you so
% desire. Currently used in the single-page-poster style. 
%    \begin{macrocode}
\newcommand{\suppressfulladdresses}{\@suppressfulladdressestrue\@needaddressesfalse}
%    \end{macrocode}
% \end{macro}
% \end{macro}
% \end{macro}
% \end{macro}
% \end{macro}
% \end{macro}

% \begin{macro}{\printdates}
% \begin{macro}{\noprintdates}
% Command to activate the printing out of the received/accepted
% dates. This can also be done using the \texttt{printdates} option,
% which is the default for \texttt{journal}, but not for \texttt{proceedings}.
%    \begin{macrocode}
\newcommand{\printdates}{\@printdatestrue}
\newcommand{\noprintdates}{\@printdatesfalse}
%    \end{macrocode}
% \end{macro}
% \end{macro}
% \begin{macro}{\@outputaddressesifnecessary}
% Command to output the addresses if we haven't already done so
% This is called by: 
% \begin{enumerate}
% \item The end of |thebibliography| environment
% \item The |\outputfulladdresses| command
% \end{enumerate}
%    \begin{macrocode}
\newcommand{\@outputaddressesifnecessary}{%
  \ClassInfo{rmaa}{Calling @outputaddressesifnecessary}
  \ifthenelse{\boolean{@needaddresses}}{%
    %% If in one column style, or in the first column of a two column
    %% style then we can try to put it on this page 
    \ifthenelse{\boolean{@firstcolumn} \or \not\boolean{@rmtwocol}}{%
      \@outputaddresses}{%
      %% otherwise give warning it will be on a separate page
      \ClassWarning{rmaa}{Full addresses will go on a separate
        page (please see User Guide for how to remedy this)}}%
    %% If we haven't managed yet, put it on its own page
    \ifthenelse{\boolean{@needaddresses}}{%
      \if@debug{\ClassWarning{rmaa}{In this silly bit!}}\fi
      \clearpage\vfill\@outputaddresses}{}%
    }{
    %% If we don't need them, we still might need to do the eoa hook 
    \ifthenelse{\boolean{@suppressfulladdresses}}{%
      \if@debug{\ClassWarning{rmaa}{XXX}}\fi
%%      \if@book\else\@endofarticlehook\null\fi
      }{}
    }%
  }  
%    \end{macrocode}
% \end{macro}

% \begin{macro}{\outputfulladdresses}
% A user interface to the above command, in case there is no biblio
% (11 May 2000) 
%    \begin{macrocode}
\newcommand{\outputfulladdresses}{\@outputaddressesifnecessary}
%    \end{macrocode}
% \end{macro}

% \begin{macro}{\adjustfinalcols}
% This command starts a new column and should be 
% used to ensure that the columns balance on the last page. It also
% outputs the full 
% author addresses to ensure that this occurs in the left hand
% column. 
% The |\par| command is so that the current line will be taken account
% of when working out |\pagetotal|
%    \begin{macrocode}
\newcommand{\adjustfinalcols}{\if@twocolumn\goodbreak\@outputaddresses\newpage\noindent\fi}
%    \end{macrocode}
% \end{macro}

% \begin{macro}{\@outputaddresses}
% Command to do the work of outputting the addresses. 
% This can be called by 
% \begin{enumerate}
% \item an explicit |\adjustfinalcols|
% \item via |\@outputaddressesifnecessary|, q.v.
% \end{enumerate}
%
%    \begin{macrocode}
\newcommand{\@outputaddresses}{%
  % Allow option of suppressing the full addresses, in which case
  % we just run the hook if we are not in a book style
  \if@suppressfulladdresses
%%     \if@book\else\@endofarticlehook\null\fi
  \else
  %% save the full addresses in a box for output later
  \sbox{\@fulladdbox}{\parbox{\textwidth}{\@fulladdresses}}
  % calculate height of address box  
  \settoheight{\@fulladdheight}{\usebox{\@fulladdbox}}
  \settodepth{\@fulladddepth}{\usebox{\@fulladdbox}}
  \addtolength{\@fulladdheight}{\@fulladddepth}
  \if@debug
  { \typeout{Address total height: \the\@fulladdheight}
    \typeout{Text height: \the\pagetotal}
    \typeout{Goal height: \the\pagegoal}
    }
  \fi
  % reuse \@fulladdheight to mean the total height of text+addresses
  \addtolength{\@fulladdheight}{\pagetotal}
%%  \addtolength{\@fulladdheight}{-\baselineskip}
  \if@debug\typeout{Adjusted text + addresses height: \the\@fulladdheight}\fi
  % check it will fit on current page
  \ifdim\@fulladdheight<\pagegoal
     %% must shift it to the left
     %% to compensate for the hanging indentation in the reference list. 
     %% NOTE: this is not foolproof!
     %% NB. This will give overfull hbox warnings, which is a bit ugly.
     \vfill\noindent\if@biblist\hspace*{-\leftmargin}\fi
     \usebox{\@fulladdbox}
     % Set this to show we've done it
     \@needaddressesfalse
     % and assume we are at the end of the article
     % so that we do the EOA hook if not in the book style
%%     \if@book\else\@endofarticlehook\fi
  \else
     % if not, just do nothing. With luck, this function will get 
     % called again on the next page
     \if@debug{\typeout{Holding over full addresses to next page}}\fi
  \fi\fi
  }
%    \end{macrocode}
% \end{macro}

% \subsubsection{Commands for editors}
% \label{sec:editor-commands}
% Some user/editor commands for doing the TOC and index in the book
% option. 
% They are defined here rather in the book-specific part since they
% need to be always available in order to avoid breaking documents. 
% \begin{macro}{\listofauthors}
% \begin{macro}{\@rawauthors}
% \begin{macro}{\rmaaTOCauthors}
% The command |\listofauthors| formats an author list for use in the
% TOC. It should contain the same names as in the |\author| command,
% but without the affiliation garbage and using |\&| instead of
% |and|. 
%    \begin{macrocode}
\newcommand{\@rawauthors}{Raw Authors}
\newcommand{\rmaaTOCauthors}{\noexpand\rmaaTOCAuthorStyle{Author list for TOC}}
\newcommand{\listofauthors}[1]{%
  \renewcommand{\@rawauthors}{#1}
  \renewcommand\rmaaTOCauthors{\noexpand\rmaaTOCAuthorStyle{#1}}}
%    \end{macrocode}
% \begin{macro}{\indexauthor}\label{indexauthor}
% The command |\indexauthor| produces an index entry for each
% author. 
% This command must come \emph{after} the
% |\title| and the |\listofauthors| commands. Also, there should be one 
% |\indexauthor| command for \emph{each} author. We also add the
% author onto the end of the |\ADS@authors| macro and construct the
% |\ADS@authorinitial| by a most cunning method (see below), but only
% for the first author. 
% 
% No, No, NO! This was all wrong. Backed out completely: 
% \begin{quote}
% Completely rewritten (19 Dec 2002) since we now use the index class, so 
% as to be able to alphabetize accented characters. A side effect of this is
% that |\index| no longer expands out the argument after the |@|, so we 
% instead add all the |\index| commands to a list (|\indexauthor@deferlist|), 
% which is executed in the |\@begofarticlehook|. 
% 
% Note the stuff with |\toks0| is copied from Exercise 20.15 of the \TeX{}book---it is
% to make sure the latin1-accented name gets passed into the idx
% file. 
% \end{quote}
% We now go back to the old way of doing things, using an optional
% argument to the \cs{indexauthor} command that has the name without
% accents. 
% \changes{3.08}{2003/01/16}{Remerged changes from 3.06 that got
% dropped from 3.07 after a synchronization disaster. Basically, the
% accentless optional argument to \cs{indexauthor}. Must add the macro
% directories to my unison set-up to stop this sort of thing from happening.}
%    \begin{macrocode}
\newcommand\@isortkey{}
\newcommand\@emptykey{empty}
\newcommand\indexauthor[2][empty]{
  \renewcommand\@isortkey{#1}
  \ifx\@isortkey\@emptykey \renewcommand\@isortkey{#2}\fi
  \index{\@isortkey @{{\rmaaIndexAuthorStyle{#2}}\hskip1em plus1em minus0.4em
      \rmaaIndexTitle\hskip1em plus2em minus0.4em\penalty-100
      \rmaaTOCauthors}}
  \indexauthorADS{#2}
}
%%\def\indexauthor@deferlist{\relax}
%%\long\def\indexauthor#1{
%%  \def\@@name{#1} \toks0=\expandafter{\@@name}
%%  \toks1=\expandafter{\indexauthor@deferlist}
%%  \protected@xdef\indexauthor@deferlist{%
%%    \the\toks1
%%    \noexpand\index{\the\toks0@{\noexpand\rmaaIndexAuthorStyle{\the\toks0}\noexpand\hskip1em plus1em minus0.4em 
%%        \rmaaIndexTitle\noexpand\hskip1em plus2em minus0.4em\noexpand\penalty-100
%%        \rmaaTOCauthors}}
%%  }
%%  \indexauthorADS{#1}
%%  }
%%\xdef\indexauthor@deferlist{}
\newcommand\indexauthorADS[1]{%
  \if@firstauthor\protected@xdef\ADS@authorinitial{\firstletter{#1}}\fi
  \protected@xdef\ADS@authors{\ADS@authors\if@firstauthor\else ; \fi
    #1}
  \@firstauthorfalse
}
%    \end{macrocode}
% Note that the first |#1| sets the key for the index, making sure it
% actually starts with  
% first letter of the surname. Meanwhile, the |@{...}| expression,
% which starts with some  
% control sequences is what is written to the |.ind| file. 
% See the makeindex man page for details. 
% \end{macro}
% \end{macro}
% \end{macro}
% \end{macro}

% \begin{macro}{\KillAllAccents}
% \begin{macro}{\firstletter}
% These are utility functions used in the above. |lett| is a counter
% indicating the current letter in the string. |\firstletter| just
% prints the first letter of its argument.
%    \begin{macrocode}
\newcommand{\KillAllAccents}{\def\'{}\def\`{}\def\~{}\def\.{}\def\i{i}}
\newcommand\firstletter[1]{\@car #1\@nil}
%    \end{macrocode}
% \end{macro}
% \end{macro}
% \begin{macro}{\SetVolume}
% \begin{macro}{\SetMSnumber}
% \begin{macro}{\SetFirstPage}
% \begin{macro}{\SetYear}
% \begin{macro}{\ReceivedDate}
% \begin{macro}{\AcceptedDate}
% These are for the editors to set the volume number, etc. 
%    \begin{macrocode}
\newcommand{\SetVolume}[1]{%
  \renewcommand{\@volume}{#1}
  \edef\ADS@volume{\PadToFourWithDots{#1}}
  \renewcommand{\@volumen}{\rmaa@VolumenStyle{volumen #1}}%
  }
\newcommand{\SetMSnumber}[1]{\renewcommand{\@MSnum}{#1}}
%    \end{macrocode}
% Note that there is no real reason to use |\SetFirstPage| since all
% the page numbering is done automatically in the book style. 
%    \begin{macrocode}
\newcommand{\SetFirstPage}[1]{\setcounter{page}{#1}}
\newcommand{\SetYear}[1]{%
  \renewcommand{\@year}{#1}
  \renewcommand{\ADS@year}{#1}
  \renewcommand{\@anyo}{\rmaa@VolumenStyle{#1}}
  }
\newcommand{\ReceivedDate}[1]{\renewcommand{\@received}{#1}}
\newcommand{\AcceptedDate}[1]{\renewcommand{\@accepted}{#1}}
%    \end{macrocode}
%
% \begin{macro}{\DOI}
% User/editor command to enter the DOI, which will be typeset as a link.
% \changes{4.0}{2019/11/02}{New user command \cs{DOI} to typeset
% DOI in the header of first page and add it to ADS info.}
%    \begin{macrocode}
\newcommand\DOI[1]{%
  \def\@DOI@url{\url{https://doi.org/10.22201/ia.01851101p.#1}}%
  \def\ADS@doi{https://doi.org/10.22201/ia.01851101p.#1}%
}%
%    \end{macrocode}
% \end{macro}
%
% \begin{macro}{\ORCIDiD}
% \changes{4.0}{2019/11/12}{Add mechanism for giving ORCID iD links}
% \begin{macro}{\ORCIDiDfull}
% \begin{macro}{\@orcid@icon}
% \begin{macro}{\@orcid@iconbig}
% Commands to add icon links to ORCID iDs of authors. For example
% \begin{quote}
% |\ORCIDiD|\marg{0000-0000-0000-0000}
% \end{quote}
% to be used after each author name in the |\author|
% command and in the |\fulladdresses|.  Alternatively,
% |\OORCIDiDfull| will typeset the icon plus the full text of the
% link, but we don't use that at present.  The icon image is stored in |ORCIDiD_iconvector.pdf|, which needs to be on the input path.  In principle, we could make EPS versions if necessary.
%    \begin{macrocode}
\newcommand\@orcid@icon{\includegraphics[height=1.8ex]{ORCIDiD_iconvector}}
\newcommand\@orcid@iconbig{\includegraphics[height=2ex]{ORCIDiD_iconvector}}
\newcommand\ORCIDiD[1]{\href{https://orcid.org/#1}{\,\@orcid@icon}}
\newcommand\ORCIDiDfull[1]{%
  \href{https://orcid.org/#1}{\@orcid@iconbig}\,\url{https://orcid.org/#1}}
%    \end{macrocode}
% \end{macro}
% \end{macro}
% \end{macro}
% \end{macro}
%
% \begin{macro}{\@volume}
% \begin{macro}{\@volumen}
% \begin{macro}{\@artpages}
% \begin{macro}{\@year}
% \begin{macro}{\@anyo}
% \begin{macro}{\@received}
% \begin{macro}{\@accepted}
% \begin{macro}{\@MSnum}
% Default values for the containers used by above commands, plus a
% few more useful ones. 
%    \begin{macrocode}
\newcommand{\@volume}{00}
\newcommand{\@volumen}{\rmaa@VolumenStyle{00}}
\newcommand{\@artpages}{%
  \oldpageref{\articlename:RMAAFirstPage}--\oldpageref{\articlename:RMAALastPage}}
\newcommand{\@year}{\the\year}
\newcommand{\@anyo}{\rmaa@VolumenStyle{2007}}
\newcommand{\@journalname}{\rmaa@JournalNameStyle{%
    \if@proceedings\RMAASClong\else\RMAAlong\fi}}
\newcommand{\@received}{}
\newcommand{\@accepted}{}
\newcommand{\@MSnum}{0000}
\newcommand{\@DOI@prefix}{https://doi.org/10.22201/ia.01851101p}
\def\@DOI@url{\url{https://doi.org/10.22201/ia.01851101p.2099.00.00.00}}
%    \end{macrocode}
% \end{macro}
% \end{macro}
% \end{macro}
% \end{macro}
% \end{macro}
% \end{macro}
% \end{macro}
% \end{macro}
% \end{macro}
% \end{macro}
% \end{macro}
% \end{macro}
% \end{macro}
% \end{macro}
% 
% The following commands are most useful in the \texttt{book}
% style. However, the |\SetConfTitle| command can now be used in
% individual articles to set the header for the \texttt{preprint}
% style. 
%    \begin{macrocode}
\newcommand{\@reunionplace}{Lugar de la reuni\'on}
\newcommand{\@reuniondate}{Fecha de la reuni\'on}
\newcommand{\@conftitle}{Conference Title}
\newcommand{\SetReunion}[2]{\renewcommand{\@reunionplace}{#1}\renewcommand{\@reuniondate}{#2}}
\newcommand{\SetConfTitle}[1]{\renewcommand{\@conftitle}{#1}}
\newcommand{\@editors}{Juan Ram\'on de la Fuente and ``El Mosh''}
\newcommand{\SetEditors}[1]{\renewcommand{\@editors}{#1}}
\newcommand{\@editores}{Juan Ram\'on de la Fuente y ``El Mosh''}
\newcommand{\SetEditores}[1]{\renewcommand{\@editores}{#1}}
\newcommand{\@mes}{mes de la publicaci\'on}
\newcommand{\SetMes}[1]{\renewcommand{\@mes}{\rmaa@VolumenStyle{#1}}}
%    \end{macrocode}

% \subsection{Notes for Author Proofs and Communication between Editors}
% \label{sec:notes}
% We introduce a command \cs{Note} that should be used to describe
% any change made to the manuscript. These are a bit tricky to get
% right. In general we use \cs{marginpar} but this only works in outer
% paragraph mode, so we just bung the box where we are if we are in a
% figure caption or something. If we are in the preamble, then we just
% use a \marginpar as part of the \maketitle command.
%    \begin{macrocode}
\newif\if@arenotes
\@arenotesfalse
\newcounter{notecounter}
\setcounter{notecounter}{0}
\newcommand\NotesList{}
\newcommand\Note[1]{%
  \if@proof{%
    \@arenotestrue%
    \stepcounter{notecounter}%
    \if@inpreamble\else{%
      \NoteMark{\arabic{notecounter}}%
    }\fi
    \protected@xdef\NotesList{\NotesList 
      \protect\item #1}%
  }\fi%
}
\newcommand\Changed[2]{Changed ``#1'' to ``#2''.}
\newcommand\NoteMark[1]{%
  \ifinner
  \raisebox{1ex}[0pt][0pt]{\makebox[0pt][l]{\framebox{\footnotesize\bfseries
        N#1}}}%
  \else
  \marginpar{\framebox{\footnotesize\bfseries N#1}}%
  \fi
}
\newcommand\RefNote[1]{\if@proof\NoteMark{\ref{#1}}\relax\fi}
\newcommand\NoteLabel{\protect\label}
\newcommand\OutputNotes{%
  \bgroup
  \onecolumn
  \section*{EDITORIAL NOTES}
  \begin{enumerate}
    \NotesList
  \end{enumerate}
  \egroup
}
\DeclareRobustCommand\cs[1]{\texttt{\char`\\#1}}

%    \end{macrocode}
% \subsection{Bibliographic Records for ADS}
% \label{sec:ADS}
% \changes{2.2beta11}{2001/11/27}{Just about finished all the ADS
% stuff now. Except for the references, which I have given up on for
% the time being.}
% \begin{macro}{\ADS@BibCode}
% \begin{macro}{\ADS@year}
% \begin{macro}{\ADS@journal}
% \begin{macro}{\ADS@volume}
% \begin{macro}{\ADS@firstpage}
% \begin{macro}{\ADS@authorinitial}
% We use the |\ADS@BibCode| command to hold the bibliographic code for
% ADS. It is 
% used in the |%R| 
% field of the bibliographic record for each article. 
% The format is 
% \begin{quote}
% \texttt{YYYYJJJJJVVVVMPPPPA}  
% \end{quote}
% where \texttt{YYYY} is the year, \texttt{JJJJJ} is the journal code,
% which is \texttt{RMxAA} or \texttt{RMxAC}, \texttt{VVVV} is the
% volume number (right-justified), \texttt{M} is not used by us,
% \texttt{PPPP} is the first page number (right-justified), and
% \texttt{A} is the initial of the first author's surname. The volume
% and page fields are padded with periods (`\texttt{.}'), which causes
% us some fun (this is done with the macro |\PadToFourWithDots| described
% below) . 
% 
% All these individual bits are held in the commands |\ADS@year|,
% |\ADS@journal|, |\ADS@volume|, |\ADS@firstpage|,
% |\ADS@authorinitial|. 
%    \begin{macrocode}
\newcommand{\SetADSBibCode}[1]{\renewcommand{\ADS@BibCode}{#1}}
\newcommand{\ADS@BibCode}
{\ADS@year\ADS@journal\ADS@volume.\PadToFourWithDots{\ADS@firstpage}\ADS@authorinitial}
\newcommand{\ADS@year}{2001} 
\newcommand{\ADS@journal}{\if@proceedings RMxAC\else RMxAA\fi}
\newcommand{\ADS@volume}{....} 
%    \end{macrocode}
% These two are just set to default values here, |\ADS@firstpage| is set
% in the |document| environment for included articles, in the
% |frontsection| environment, |\tableofcontents| command,  |index| and
% |abstracts| environments. Also by |\AtBeginDocument| for stand-alone
% articles. 
% 
% God knows how we are going to do |\ADS@authorinitial|!
%    \begin{macrocode}
\newcommand{\ADS@firstpage}{....}
\newcommand{\ADS@authorinitial}{A}
%    \end{macrocode}
% \end{macro}
% \end{macro}
% \end{macro}
% \end{macro}
% \end{macro}
% \end{macro}
% We also have a whole slew of other fields that can be set. These are
% all listed in Table~\ref{tab:ADS}. Most of them still have to be
% defined. 
% \begin{table}[!ht]\raggedleft
% \caption{ADS Bibliographic Records}\label{tab:ADS}
% \begin{tabular}{lp{8em}lp{8em}}\hline
% Code & Description & Macro Name & Comment \\ \hline
% |%R| & Bibliographic Code & |\ADS@BibCode| & required \\
% |%A| & Author list & |\ADS@authors| & required, generate from |\indexauthor| \\
% |%F| & Author affiliation & |\ADS@affil| & format: AA(\dots)
% AB(\dots) \dots \\
% |%J| & Journal Name & |\ADS@journalname| &  \\
% |%D| & Publication Date & |\ADS@fulldate| &  \\
% |%L| & Last Page & |\ADS@lastpage| & generate from label \\
% |%T| & Title & |\ADS@title| & required, generate from |\title| \\
% |%C| & Abstract Copyright & |\ADS@copyright| & set to IA-UNAM \\
% |%U| & URL for electronic document & |\ADS@url| & set to PDF file \\
% |%d| & DOI for article & |\ADS@doi| & set by |\DOI| \\
% |%K| & Keywords & |\ADS@keywords| & set equal to |\@keywords| \\
% |%A| & Abstract Text & |\ADS@abstract| & set equal to |\@abstract| \\
% |%R| & References & |\ADS@references| & generate from the
% bibliography  \\ \hline
% \end{tabular}
% \end{table}
% A lot of these macros are set by other commands, as mentioned in the
% Table, but some need to be set  explicitly here. 
% 
% \begin{macro}{\ADS@doi}
% \changes{4.0}{2019/11/08}{Save DOI for ADS}
% The user/editor command \cs{DOI} is used to set this. 
%    \begin{macrocode}
\newcommand{\ADS@doi}{https://doi.org/10.22201/ia.01851101p.0000.00.00.00}
%    \end{macrocode}
% \end{macro}
% 
% \begin{macro}{\ADS@journalname}
% This
% does not contain any accents at the moment so as not to cause
% problems---this maybe should be changed though. 
%    \begin{macrocode}
\newcommand{\ADS@journalname}{%
  \if@proceedings\@conftitle\space (Eds. \@editors) \fi%
  Revista Mexicana de Astronom\'\i{}a y Astrof\'\i{}sica%
  \if@proceedings\space (Serie de Conferencias)\fi%
  \space Vol. \@volume, pp. \ADS@firstpage--\ADS@lastpage\space (\@year)%
  \space (http://www.astroscu.unam.mx/\~{}rmaa/)
  }
%    \end{macrocode}
% \end{macro}
% \begin{macro}{\ADS@keywords}
% \begin{macro}{\ADS@abstract}
% \begin{macro}{\ADS@title}
% |\ADS@abstract| is straightforward, since we already have the text
% saved as |\@abstract|. We could do the same with |\ADS@keywords| but
% to avoid unwanted formatting we explicitly set it in |\keywords|
% (p.~\pageref{keywords}). |\ADS@title| is just set to |\@rawtitle|,
% which we were already using for another purpose. 
%    \begin{macrocode}
\let\ADS@title=\@rawtitle
\newcommand{\ADS@keywords}{No keywords specified}
\newcommand{\ADS@abstract}{No abstract specified}
%    \end{macrocode}
% \end{macro}\end{macro}\end{macro}
% \begin{macro}{\ADS@authors}
% |\ADS@authors| wants to be in ``Other, A. N.'' format, but with
% semi-colons between authors. The most natural way to do this seems to
% be with |\indexauthor| (p.~\pageref{indexauthor}). We append each
% author to this macro, which means we must explicitly clear it for
% each new article (done in |\article|). .   
%    \begin{macrocode}
\newcommand{\ADS@authors}{}
%    \end{macrocode}
% \end{macro}
% \begin{macro}{\ADS@fulldate}
% \begin{macro}{\ADS@lastpage}
% \begin{macro}{\ADS@copyright}
% \begin{macro}{\ADS@website}
% \begin{macro}{\ADS@url}
% \begin{macro}{\SetWebsite}
% \begin{macro}{\SetMonthNumeric}
% Other sundry data, together with user commands for bits that can be
% set by the Editors. The |\ADS@url| is constructed from the website
% (currently |http://www.astrosmo.unam.mx/rmaa/|,
% the journal descriptor (e.g., |RMxAC..10/|), followed by the
% path to the PDF file (e.g., |PDF/RMxAC..10_HEW1.pdf|). 
% 
%    \begin{macrocode}
\newcommand{\ADS@website}{http://www.astrosmo.unam.mx/rmaa/}
\newcommand{\ADS@url}
{\ADS@website\ADS@journal\ADS@volume\ADS@number/PDF/\ADS@journal\ADS@volume\ADS@number\_\articlename.pdf}
\newcommand{\ADS@number}{}
\newcommand{\ADS@fulldate}{00/\ADS@year} 
\newcommand{\ADS@lastpage}{00} 
\newcommand{\ADS@copyright}{(C) \ADS@year; Instituto de Astronom\'\i{}a,
  Universidad Nacional Aut\'onoma de M\'exico} 
\newcommand{\SetNumber}[1]{\renewcommand{\ADS@number}{-#1}}
\newcommand{\SetWebsite}[1]{\renewcommand{\ADS@website}{#1}}
\newcommand{\SetMonthNumeric}[1]{\renewcommand{\ADS@fulldate}{#1/\ADS@year}}
%    \end{macrocode}
% \end{macro}\end{macro}\end{macro}\end{macro}\end{macro}\end{macro}\end{macro}
% \begin{macro}{\ADS@affil}
% \begin{macro}{\ADS@references}
% These two are somewhat more problematical. Affiliation currently
% only works for the first author (see \dots) pending redesign of the
% user interface. References are now done by wrapping the
% |thebibliography| environment to open another input stream on the
% same file and search through for the |\bibitems|. 
%    \begin{macrocode}
\newcommand{\ADS@references}{}
\newcommand{\ADS@affil}{}
%    \end{macrocode}
% \end{macro}\end{macro}
% 
% Define two macros that expand to the strings
% |\begin{thebibliography}| and |\end{thebibliography}|. Note that the
% idea for this and much else is inspired by the code for the docstrip
% package. Also define |\StartBibX| in case authors use the more
% standard |\begin{thebibliography}{}| syntax. 
%    \begin{macrocode}
\edef\EndBib{\expandafter\@gobble\string\\end\string{thebibliography\string}}
\edef\StartBib{\expandafter\@gobble\string\\begin\string{thebibliography\string}}
\edef\StartBibX{\expandafter\@gobble\string\\begin\string{thebibliography\string}\string{\string}}
%    \end{macrocode}
% Define two new logical variables. |\ifAnotherline| is true if we
% want to read another line from the input stream. It is set to false
% when we encounter the |\end{thebibliography}|. |\ifInBib| is set to
% true when we frind the |\begin{thebibliography}| and tells us to
% start saving the following lines. 
%    \begin{macrocode}
\newif\ifAnotherLine
\newif\ifInBib
%    \end{macrocode}
% \begin{macro}{makeOther}
% We need a macro to change the catcodes of lots of things to other. 
%    \begin{macrocode}
\def\makeOther#1{\catcode`#1=12\relax}
%    \end{macrocode}
% \end{macro}
% \begin{macro}{\RereadBiblio}
% Define a macro to do all the fiddly stuff. This is used at the
% beginning of |thebibliography| environment (Sec~\ref{sec:biblio}).
%    \begin{macrocode}
\newread\tf@biblio
\newcommand{\RereadBiblio}{%
%    \end{macrocode}
% Open the current file again for reading
%    \begin{macrocode}
  \begingroup
  \immediate\openin\tf@biblio=\articlename
%    \end{macrocode}
% Everything in a group to be safe. Initialize logicals. Black magic
% to change
% catcodes and avoid blank lines being interpreted as |\par|. 
%    \begin{macrocode}
  \AnotherLinetrue\InBibfalse
  \makeOther\ \makeOther\\\makeOther\$%
  \makeOther\#\makeOther\^\makeOther\^^K%
  \makeOther\_\makeOther\^^A\makeOther\%%
  \makeOther\~\makeOther\{\makeOther\}\makeOther\&%
  \endlinechar-1\relax
%    \end{macrocode}
% Start looping through the file. Read each line to the macro
% |\Line|. If we have got to EOF or end of biblio then set
% |\AnotherLinefalse|.
%    \begin{macrocode}
  \loop
  \read\tf@biblio to \Line
  \ifx\Line\EndBib\AnotherLinefalse\InBibfalse\fi
  \ifeof\tf@biblio \AnotherLinefalse\fi
%    \end{macrocode}
% If we are in the biblio and the line is not blank then append it to
% the |\ADS@references| macro together with a newline. 
%    \begin{macrocode}
  \ifInBib
  \ifx\Line\empty\else\xdef\ADS@references{\ADS@references\Line^^J}\fi
  \fi
%    \end{macrocode}
% Check if the line we just saw was the start of biblio, then repeat
% the loop if necessary. 
%    \begin{macrocode}
  \ifx\Line\StartBib\InBibtrue\fi
  \ifx\Line\StartBibX\InBibtrue\fi
  \ifAnotherLine\repeat
  \immediate\closein\tf@biblio
  \endgroup
  }
%    \end{macrocode}
% \end{macro}
% 
% \begin{macro}{\PadToFourWithDots}
% Utility function for left-padding a number to 4 characters with
% periods. Used in 
% constructing |\ADS@BibCode|. This is a very simple algorithm and
% \emph{requires} that the argument actually be a number---if it is
% not then dire consequences will result.  
%    \begin{macrocode}
\newcommand{\PadToFourWithDots}[1]{%
  \ifnum #1 > 999 #1%
  \else\ifnum #1 > 99 .#1%
  \else\ifnum #1 > 9 ..#1%
  \else ...#1%
  \fi\fi\fi}
%%  \ifnum \count0 > 999 #1%
%%  \else\ifnum \count0 > 99 .#1%
%%  \else\ifnum \count0 > 9 ..#1%
%%  \else ...#1%
%%  \fi\fi\fi}
%    \end{macrocode}
% \end{macro}
%
% \begin{macro}{\AuthorADS}
% \begin{macro}{\ADS@affiliations}
% New mechanism for setting the affiliations for ADS. The author labels (AA, AB, \dots) are generated using two nested counters.  We use the |\robustify| command from \texttt{etoolbox} so that accented characters in the affiliation addresses do not blow up.  The list of accents may need to be expanded. 
%    \begin{macrocode}
\newcounter{AUTHOR}
\newcounter{author}[AUTHOR]
\stepcounter{AUTHOR}
\newcommand\ADS@affiliations{}
\robustify{\'}\robustify{\i}\robustify{\~}\robustify{\`}
\def\AuthorADS#1#2{%
  \begingroup
  \def\ORCIDiD##1{<ID system="ORCID">##1</ID>}
  \def\Email##1{<EMAIL>##1</EMAIL>}
  \def\Affil##1{##1}
  \stepcounter{author}
  \ifthenelse{\value{author} > 26}{\stepcounter{AUTHOR}\stepcounter{author}}{}
  \ifthenelse{\equal{\ADS@affiliations}{\@rmaaEmpty}}{%
    \xdef\ADS@affiliations{\Alph{AUTHOR}\Alph{author}(#2)}
  }{%
    \xdef\ADS@affiliations{\ADS@affiliations, \Alph{AUTHOR}\Alph{author}(#2)}
  }%
  \endgroup
}
%    \end{macrocode}
% \end{macro}
% \end{macro}
% %
% \begin{macro}{\ADS@dumpdata}
% \changes{4.0}{2019/11/13}{Write all author affiliations for ADS}
% \begin{macro}{\ADS@write}
% \begin{macro}{\ADS@openfile}
% |\ADS@dumpdata| is the macro that actually does the job of writing out all the
% data to an auxiliary file, using |\ADS@write| to write each
% field. The file is set up with |\ADS@openfile|, which currently is
% done in the TOC of a book.  The new multi-author affiliations are used if |\ADS@affiliations| has been set, otherwise it falls back to the single-author version |\ADS@affil|. 
%    \begin{macrocode}
\newcommand{\ADSfield}[2]{\relax}
\newcommand{\ADSvfield}[2]{\relax}
%%\renewcommand{\ADSfield}[12]{\%#1 #2}
\newcommand{\ADS@dumpdata}
{ 
  \ADS@write{\string\begin{ADSarticle}}
  \ADS@WRITE{\string\ADSfield{R}{\ADS@BibCode}}
  \ADS@write{\string\ADSfield{A}{\ADS@authors}}
  \ifthenelse{\equal{\ADS@affiliations}{\@rmaaEmpty}}{%
    {\ADS@sanitize\ADS@write{\string\ADSfield{F}{AA(\ADS@affil)}}}
  }{%
    {\ADS@sanitize\ADS@write{\string\ADSfield{F}{\ADS@affiliations}}}
  }
  \ADS@write{\string\ADSfield{J}{\ADS@journalname}}
  \ADS@write{\string\ADSfield{D}{\ADS@fulldate}}
  \ADS@write{\string\ADSfield{L}{\ADS@lastpage}}
  {\ADS@sanitize\ADS@write{\string\ADSfield{T}{\ADS@title}}}
  \ADS@write{\string\ADSfield{C}{\ADS@copyright}}
  \ADS@write{\string\ADSfield{d}{\ADS@doi}}
  \ADS@write{\string\ADSfield{U}{\ADS@url}}
  \ADS@write{\string\ADSfield{K}{\ADS@keywords}}
  {\ADS@sanitize\ADS@write{\string\ADSfield{B}{\ADS@abstract}}}
  \ADS@write{\string\ADSfield{Z}{\ADS@references}}
  \ADS@write{\string\end{ADSarticle}}
  }
%%\newcommand{\ADS@write}[1]{\if@book\addtocontents{ads}{#1}\fi}
%%\newcommand{\ADS@openfile}{\@starttoc{ads}}
\newcommand{\ADS@write}[1]{\set@display@protect\immediate\write\tf@ads{#1}}
\newcommand{\ADS@WRITE}[1]{\immediate\write\tf@ads{#1}}
\newcommand{\ADS@openfile}{%
  \newwrite\tf@ads%
  \immediate\openout\tf@ads\jobname.ads\relax
  \ADS@WRITE{\ADS@preamble}}
\newcommand{\ADS@closefile}{\ADS@WRITE{\ADS@postamble}}
\def\ADS@preamble{%
  %%  \string\documentclass{article}
  %%  \string\begin{document}
  }
\def\ADS@postamble{%
  %%  \string\end{document}
  \string%\string% Local Variables:^J 
  \string%\string% mode: LaTeX^J
  \string%\string% mode: Font-lock^J
  \string%\string% End:^J
  }
%    \end{macrocode}
% \end{macro}
% \end{macro}
% \end{macro}
% 
% \subsection{Automatic calculation of the page range}
% \label{sec:page-range}
% Another issue that has cost me a lot of grief. Still, it finally
% seems to be sorted now. All hail Jeff Goldberg. 
%
% Cunning trick to get the page numbers from |\pageref|.  Put a label
% on the first and last page of the document using |\AtBeginDocument|
% and |\AtEndDocument|. .  This way won't work for the book option.
% In that case, we do it by defining |\@endofarticlehook| and
% |\@begofarticlehook| as parallels to |\AtEndDocument| and
% |\AtBeginDocument|.  Note that we have to defer adding
% |\@begofarticlehook| to |\AtBeginDocument| to make sure it comes
% after the redefinitions of |\ref| and |\label|.
%    \begin{macrocode}
\if@book\AtEndDocument{\ADS@closefile}\else
 %\AtBeginDocument{\@begofarticlehook}
\AtEndDocument{\@endofarticlehook}
\fi
%    \end{macrocode}
% \begin{macro}{\@begofarticlehook}
% Since we are getting more and more things that need to be done at
% the beginning of every article, frontsection, etc we define a sister
% hook to 
% |\@endofarticlehook|. This hook does the following: 
% \begin{enumerate}
% \item Labels the first page
% \item Sets |\ADS@firstpage|
% \item runs and then clears |\indexauthor@deferlist|
% \item resets the keyword counter
% \end{enumerate}
% It is executed in the following circumstances: 
% \begin{enumerate}
% \item At |\begin{document}| for standalone files
% \item At |\begin{document}| for each included file
% \item At beginning of front sections 
% \item At beginning of TOC
% \item At beginning of Author Index
% \item At beginning of abstracts
% \end{enumerate}
% For some reason, this is not always within the right group, so we
% use |\xdef| just to be sure. 
%    \begin{macrocode}
\newcommand{\@begofarticlehook}{%
  \if@frontmatter\frontlabel{RMAAFirstPage}\else\label{RMAAFirstPage}\fi%
  \xdef\ADS@firstpage{\@arabic{\c@page}}%
  \if@book\else\makeindex\fi % needed so default index is defined in stand-alone articles
%%  \indexauthor@deferlist
%%  \xdef\indexauthor@deferlist{}
  \@inpreamblefalse
  \setcounter{kwdcounter}{0}
  }
%    \end{macrocode}
% \end{macro}
% \begin{macro}{\@endofarticlehook}
% This is a hook to be run at the end of each article
% currently, it just labels the last page and writes out all the ADS
% data to an auxiliary file.  
% It should not produce any printed output. 
% It is modelled on the |\AtEndDocument| command, to which it reduces
% in the case of a single article. 
%    \begin{macrocode}
\newcommand{\@endofarticlehook}{%
  \if@debug{\ClassWarning{rmaa}{***Doing Hook***}}\fi
  \if@debug{\ClassWarning{rmaa}{ADSabstract is: \ADS@abstract}}\fi
  \clearpage\rmaa@lastpagelabel\if@book\ADS@dumpdata\fi
  \ifx\NotesList\@rmaaEmpty\else\OutputNotes\fi
  \if@debug{\ClassWarning{rmaa}{Notes list is: \NotesList}}\fi
}
%    \end{macrocode}
% \begin{macro}{\rmaa@lastpagelabel}
% New version of last page label command copied from Jeff Goldberg's
% lastpage.sty package - we can't just use |\label| here in case the
% output routine for the last page has already been called. Instead,
% we first explicitly clear the page, step back the page counter,
% explicitly write to the .aux file, then restep the page counter
% again.
%
% But this will conflict with the hyperref package, which does labels differently.  Compensate for that by using the five-argument form.
% 
% Also set the |\ADS@LastPage| macro to the appropriate number. 
%    \begin{macrocode}
\newcommand{\rmaa@lastpagelabel}{%
  \addtocounter{page}{-1}%
  \immediate\write\@auxout{\string
    \newlabel{\articlename:RMAALastPage}{{}{\thepage}{}{}{}}}%
  \edef\ADS@lastpage{\thepage}%
  \addtocounter{page}{1}}
%    \end{macrocode}
% \end{macro}
% \end{macro}

% \subsection{Deal with explicit change in number of columns}
% \label{sec:columns}
% Sometimes, when formatting a book, we need to switch between
% one-column articles and two-column articles. For this to work, we
% need to wrap the standard \LaTeX{} commands to set the |\@rmtwocol|
% flag appropriately and also reset the package default values of all
% scalable lengths. .  
%    \begin{macrocode}
\let\latexonecolumn=\onecolumn
\let\latextwocolumn=\twocolumn
\renewcommand{\onecolumn}{\@rmtwocolfalse\@package@default@lengths\latexonecolumn}
\renewcommand{\twocolumn}{\@rmtwocoltrue\@package@default@lengths\latextwocolumn}
%    \end{macrocode}
% 
% \subsection{The actual formatting of the title}
% \label{sec:title-format}

% Note, we have to redefine the |\maketitle| command to use the
% numbered footnotes. This means we needn't use the atbeginend stuff. 
%    \begin{macrocode}
\renewcommand{\maketitle}{\par
  \begingroup
%    \end{macrocode}
% Don't switch to marked footnotes. 
% This is consistent with existing style but I'm not sure I like it
%    \begin{macrocode}
  %%  \renewcommand\thefootnote{\fnsymbol{footnote}}%
%%  \def\@makefnmark{\hbox to\z@{$\m@th^{\@thefnmark}$\hss}}%
  \def\@makefnmark{\textsuperscript{\@thefnmark}\hspace{-0.2em}}%
  \long\def\@makefntext##1{\parindent 1em\noindent
    \hbox to1.8em{\hss$\m@th^{\@thefnmark}$}##1}%
  \if@rmtwocol
  { \if@debug{\ClassWarning{rmaa}{TWO COLUMN}}\fi
    \twocolumn[\@maketitle]%
    }
  \else
  { \if@debug{\ClassWarning{rmaa}{ONE COLUMN}}\fi
    \newpage
    \global\@topnum\z@   % Prevents figures from going at top of page.
    \if@manuscript \@MSmaketitle \else \@maketitle \fi
    }
  \fi
  \if@manuscript\else\thispagestyle{rmaafirstpage}\@thanks\fi
  \ifx\NotesList\@rmaaEmpty\else{%
    \NoteMark{1--\arabic{notecounter}}\vspace*{-2\baselineskip}%
  }\fi%
  \endgroup
%    \end{macrocode}
% Don't reset footnote counter (not sure whether this is good thing
% or not).
%    \begin{macrocode}
  %%  \setcounter{footnote}{0}%
  \let\thanks\relax\let\affil\relax\let\altaffil\relax
  \let\maketitle\relax\let\@maketitle\relax
  \gdef\@thanks{}\gdef\@author{}\gdef\@title{}}
%    \end{macrocode}

% \begin{macro}{\@kwhwidth}
% \begin{macro}{\@kwtwidth}
% Redefine \cs{@maketitle} to include the abstract, resumen and
% keywords. Note that in the manuscripty style we use
% \cs{@MSmaketitle} instead.  
% \changes{3.16}{2003/04/07}{
%  Added extra vertical space between authors and accepted date.}
% \changes{3.26}{2011/08/15}{Add extra label to review articles}
%    \begin{macrocode}
\newlength{\@kwhwidth}\newlength{\@kwtwidth}
\newcommand\reviewname{Review}
\renewcommand{\@maketitle}{%
  \newpage
  \null
  \if@reviewarticle
  \smallskip
  {\large\bfseries\itshape\reviewname}%
  \bigskip
  \else
  \vskip\@before@title@skip%
  \fi
  \begin{center}%
    {\@title \par}%
    \vskip\@before@author@skip%
    \parbox{\rmaa@AuthorWidthFrac\textwidth}{\centering
      \lineskip\@between@author@skip%
      \@author
      \par}
    %% this should be replaced with `Submitted to...' for preprints
    \if@printdates
    \vskip\@before@dates@skip%
    \if@preprint
    \rmaa@ReceivedStyle{Submitted to RMxAA \@received}
    \else
    \rmaa@ReceivedStyle{Received \@received; accepted \@accepted}
    \fi
    \fi
    \if@manuscript
    \rmaa@ReceivedStyle{Submitted to RMxAA \@received}
    \fi
    \if@guide
    \rmaa@ReceivedStyle{Version \rmaa@version, \rmaa@date}
    \fi
    \vskip\@before@abstract@skip%
    \if@abstractpresent
    \parbox{\rmaa@AbstractWidthFrac\textwidth}{%
      \setlength{\parindent}{0.9cm}%
      \centerline{\rmaa@AbstractHeaderStyle{Abstract}}\par%
      \vskip\@after@abstracthead@skip%
      \if@newstyle\noindent\fi\@abstract
      \vskip\@between@abstract@skip%
      \centerline{\rmaa@AbstractHeaderStyle{Resumen}}\par%
      \vskip\@after@abstracthead@skip%
      \if@newstyle\noindent\fi\@resumen
      \if@setresumen\else\@abstract\fi % append english abstract if
                                       % spanish one not set
      \if@guide\else            %do not use keywords for guide
      \vskip\@before@kwd@skip%
      % Format the Key Words with hanging indentation:
      % They are put in a \parbox with width equal to the line width
      % minus the width of the header. 
      { % Widths of keyword
        \settowidth{\@kwhwidth}{\@keywordheader~}   % header and
        \setlength{\@kwtwidth}{\linewidth}          % keyword text
        \addtolength{\@kwtwidth}{-\@kwhwidth}       %
        \noindent\@keywordheader~\parbox[t]{\@kwtwidth}{%
          \rmaa@KeywordStyle{\boldmath\uppercase{\@keywords}}}}%
      \fi
      }
      \fi
  \end{center}%
  \vspace*{\afterkwdskip}
  \par
  }
\newlength{\afterkwdskip}\setlength{\afterkwdskip}{\z@}
%    \end{macrocode}
% \end{macro}
% \end{macro}
%
%
% \begin{macro}{\reviewarticle}
% \begin{macro}{\reviewtoc}
% \begin{macro}{\reviewlof}
%   User command to specify a review article. Must be used in
%   preamble. Main effect is to place the text ``Review'' before the
%   article title.  Also add command for setting a table of contents
%   and a list of figures at the beginning of the article.
%    \begin{macrocode}
\newcommand{\reviewarticle}{\@reviewarticletrue}
\newcommand\reviewtoc{%
  \section*{CONTENTS}%
  \@starttoc{rtc}%
}
\newcommand\reviewlof{%
  \section*{LIST OF FIGURES}%
  \@starttoc{lof}%
}
%    \end{macrocode}
% \end{macro}
% \end{macro}
% \end{macro}
% \begin{macro}{\l@section}
%    \begin{macrocode}
\renewcommand*\l@section[2]{%
  \ifnum \c@tocdepth >\z@
    \addpenalty\@secpenalty
    \addvspace{0.8em \@plus\p@}%
    \setlength\@tempdima{1.5em}%
    \begingroup
      \parindent \z@ \rightskip \@pnumwidth
      \parfillskip -\@pnumwidth
      \leavevmode %\bfseries
      \advance\leftskip\@tempdima
      \hskip -\leftskip
      #1\nobreak\hfil \nobreak\hb@xt@\@pnumwidth{\hss #2}\par
    \endgroup
  \fi}
%    \end{macrocode}
% \end{macro}
%
% \begin{macro}{\@MSmaketitle}
% A special version of |\@maketitle| for the \texttt{manuscript}
% option. This can deal with a multi-page abstract/resumen. Not very
% well, mind. 
%    \begin{macrocode}
\newcommand{\@MSmaketitle}{%
  \if@debug\ClassInfo{rmaa}{Using special MS maketitle command}\fi
  \newpage
  {
    \thispagestyle{rmaafirstpage}
%     \setlength{\parindent}{0pt}
    \null\vskip 2ex
    \begin{center}
      \@title\par
      \@author\par
      \rmaa@ReceivedStyle{%
        \ifx\@received\@rmaaEmpty
          Draft version: \@date
        \else
          Received: \@received
          \ifx\@accepted\@rmaaEmpty
          \else
            ; Accepted: \@accepted
          \fi
        \fi}\par
    \end{center}
    \@thanks
    \centerline{\rmaa@AbstractHeaderStyle{Resumen}}
    \noindent\@resumen\par\vskip 2ex
    \centerline{\rmaa@AbstractHeaderStyle{Abstract}}
    \noindent\@abstract\par\vskip 2ex
    \noindent\@keywordheader~\rmaa@KeywordStyle{\@keywords}\par
%     \@keywordheader~\rmaa@KeywordStyle{\uppercase{\@keywords}}\par
  }%
  }
%    \end{macrocode}
% \end{macro}

% \subsection{The sectioning commands}
% \label{sec:section-commands}

% \begin{macro}{\section}
% \begin{macro}{\subsection}
% \begin{macro}{\subsubsection}
% \begin{macro}{\paragraph}
% \begin{macro}{\subparagraph}
% Redefine all the standard commands to use these lengths. 
%    \begin{macrocode}
\renewcommand\section{\@startsection {section}{1}{\z@}%
                                   {\abovesectionskip}%
                                   {\belowsectionskip}%
                                   %% how do we get upper case here
                                   {\centering\reset@font\normalsize}}
\renewcommand\subsection{\@startsection{subsection}{2}{\z@}%
                                     {\abovesectionskip}%
                                     {\belowsectionskip}%
                                     {\centering\reset@font\normalsize\itshape}}
\renewcommand\subsubsection{\@startsection{subsubsection}{3}{\z@}%
                                     {\abovesectionskip}%
                                     {\belowsectionskip}%
                                     {\centering\reset@font\normalsize\itshape}}
\renewcommand\paragraph{\@startsection{paragraph}{4}{\z@}%
                                    {\aboveparaskip}%
                                    {\belowparaskip}%
                                    {\reset@font\normalsize\bfseries}}
\renewcommand\subparagraph{\@startsection{subparagraph}{5}{\parindent}%
                                      {\aboveparaskip}%
                                      {\belowparaskip}%
                                      {\reset@font\normalsize\bfseries}}
%    \end{macrocode}
% \end{macro}
% \end{macro}
% \end{macro}
% \end{macro}
% \end{macro}

% \begin{macro}{secnumdepth}
% Make sure that we number the subsubsections---this is done by
% default in the |article| class but not the |book| class, which is
% why we must set it explicitly here.
%    \begin{macrocode}
\setcounter{secnumdepth}{3}
%    \end{macrocode}
% \end{macro}
% \changes{2.2beta10}{2001/11/07}{Fixed \cs{secnumdepth} for \cs{book}
% style.}

% \begin{macro}{\thesubsection}
% \begin{macro}{\thesubsubsection}
% For subsection and subsubsection, we want the numbers to be roman
% but the title to be italic. This seems the easiest way to go about it.
%    \begin{macrocode}
\renewcommand{\thesubsection}{{\upshape\thesection.\@arabic\c@subsection}}
\renewcommand{\thesubsubsection}{\thesubsection{\upshape.\@arabic\c@subsubsection}}
%    \end{macrocode}
% \end{macro}
% \end{macro}

% \begin{macro}{\@seccntformat}
% Add ``.'' to the number here so we can remove it from |\thesection|,
% etc. Reduce gap between number and title WJH 27 Jan 1995
%    \begin{macrocode}
\def\@seccntformat#1{\csname the#1\endcsname .\hskip 0.4em}
%    \end{macrocode}
% \end{macro}

% \begin{macro}{\@sect}
% I didn't really want to mess with the |\@sect| command, but it was
% the only way I could manage to force upper case. Other than the
% commented parts, it is identical to the standard version. 
%    \begin{macrocode}
\def\@sect#1#2#3#4#5#6[#7]#8{
  \ifnum #2>\c@secnumdepth
     \let\@svsec\@empty
  \else
     \refstepcounter{#1}%
     \let\@@protect\protect
     \def\protect{\noexpand\protect\noexpand}%
     \edef\@svsec{\@seccntformat{#1}}%
     \let\protect\@@protect
  \fi
%    \end{macrocode}
% This is the bit I changed (27 Jan 1995). Use upper case for
% |\section|, use mixed case for the rest. 
%    \begin{macrocode}
  \ifnum #2=1                   
     \def\sec@upcase##1{\uppercase{##1}}%
     \def\sec@shape{\upshape}%
  \else                         
     \def\sec@upcase##1{\relax{##1}}%
     \def\sec@shape{\itshape}%
  \fi
  \@tempskipa #5\relax
  \ifdim \@tempskipa>\z@
     \begingroup #6\relax
          \@hangfrom{\hskip #3\relax\@svsec}%
                    {\interlinepenalty \@M \sec@upcase{#8}\par}%
     \endgroup
     \csname #1mark\endcsname{#7}%
%    \end{macrocode}
% Only write TOC entries if doing the Guide style. 
%    \begin{macrocode}
     \if@guide 
        \addcontentsline{toc}{#1}{%
          \ifnum #2>\c@secnumdepth \else
             \protect\numberline{\csname the#1\endcsname}%
          \fi
          #7}%
     \fi
%    \end{macrocode}
% Also allow per-article TOC for review papers
%    \begin{macrocode}
     \if@reviewarticle
        \ifnum #2>2 \else
           \addcontentsline{rtc}{#1}{%
             \protect\numberline{\csname the#1\endcsname}%
             \sec@upcase{\sec@shape #7}}%
        \fi
     \fi
  \else
     \def\@svsechd{#6\hskip #3\relax  %% \relax added 2 May 90
       \@svsec #8\csname #1mark\endcsname{#7}%
       \if@guide
          \addcontentsline{toc}{#1}{%
            \ifnum #2>\c@secnumdepth \else
               \protect\numberline{\csname the#1\endcsname}%
            \fi
            #7}%
       \fi
       }%
  \fi
  \@xsect{#5}}
%    \end{macrocode}
% \end{macro}


% \subsection{Appendices}
% \label{sec:appendices}
% Make sure that appendices in individual articles work properly. We
% copy the definition from the article class (the book one doesn't do
% what we want).  
% \begin{environment}{appendices}
%   This uses an environment |appendices| inside of which each
%   |\section| will become a separate appendix.  Note that the
%   appendices must come before the references (otherwise the full
%   addresses mechanism will not work properly). We make the equation
%   numbering be subordinate to the appendix (equations A12, B6,
%   etc). Finally, we have a medskip followed by APPENDICES to
%   introduce them. \changes{3.17}{2002/04/08}{ Added appendices environment
%   that does the right thing.}
%    \begin{macrocode}
\newenvironment{appendices}{\par
  \begingroup
  \setcounter{section}{0}%
  \setcounter{subsection}{0}%
  \renewcommand\thesection{\@Alph\c@section}
  \renewcommand\theequation{\thesection\arabic{equation}}
  \medskip
  \section*{APPENDICES}}
{\endgroup}
%    \end{macrocode}
% \end{environment}
% \begin{environment}{appendix}
%   If there is only one appendix, then use |appendix| environment
%   instead.  \changes{4.0}{2019/11/12}{ Added
%   appendix environment for case where there is only one section.}
%    \begin{macrocode}
\renewenvironment{appendix}{\par
  \begingroup
  \setcounter{section}{0}%
  \setcounter{subsection}{0}%
  \renewcommand\thesection{\@Alph\c@section}
  \renewcommand\theequation{\thesection\arabic{equation}}
  \medskip
  \section*{APPENDIX}}
{\endgroup}
%    \end{macrocode}
% \end{environment}

% \subsection{Floating thingies}
% \label{sec:floats}   

%% 
%    \begin{macrocode}
\renewcommand{\thetable}{\@arabic\c@table}
\renewcommand{\thefigure}{\@arabic\c@figure}
\renewcommand{\figurename}{Fig.}
\renewcommand{\tablename}{TABLE}
%\renewcommand{\fnum@table}{\MakeUpperCase{\tablename}~\thetable}
%    \end{macrocode}
% For some reason, I have to write `TABLE' explicitly here
% If I use |\tablename|, it gets converted to mixed case---God knows
% why. 
% |\figurename| is also being ignored. 
%    \begin{macrocode}
\renewcommand{\fnum@table}{TABLE~\thetable}
\renewcommand{\fnum@figure}{Fig.~\thefigure}
%    \end{macrocode}

% Do captions in 9 pt face in the main journal style and in new-style
% anything. 
%    \begin{macrocode}
\long\def\@makefigcap#1#2{%
  \if@proceedings\else\small\fi
  \if@newstyle\small\fi 
  \vskip\abovecaptionskip
  \sbox\@tempboxa{#1. #2}%
  \ifdim \wd\@tempboxa >\hsize
    #1. #2 \par
  \else
    \global \@minipagefalse
    \hb@xt@\hsize{\hfil\box\@tempboxa\hfil}%
  \fi
  \vskip\belowcaptionskip}
%    \end{macrocode}

% Make sure table caption is upper case. 
%    \begin{macrocode}
\long\def\@maketabcap#1#2{%
  \hbox to\hsize{\hfil #1 \hfil}
  \vskip 6pt 
  \centering\MakeUpperCase{#2}\par
  \vskip 6pt}
%    \end{macrocode}

%    \begin{macrocode}
\newlength{\@rottabcapwidth}
\long\def\@makerottabcap#1#2{%
  \setbox\@tempboxa\hbox{\MakeUpperCase{#2}}%
  \ifdim\wd\@tempboxa >.8\vsize
     \setlength\@rottabcapwidth{.8\vsize}
  \else
     \setlength\@rottabcapwidth{\wd\@tempboxa}
  \fi
  \rotatebox{90}{%
    \begin{minipage}{\vsize}\centering
      \begin{minipage}{\@rottabcapwidth}\centering
        #1\\ \vskip 6pt \MakeUpperCase{#2}
      \end{minipage}%
    \end{minipage}
    }%
  \hskip 6pt
  }

\newenvironment{Sideways}
    {\begin{sideways}\begin{minipage}{\vsize}\centering}
    {\end{minipage}\end{sideways}}

\newenvironment{Sidewaystable}
{%
  \let\oldcaption=\caption%
  \def\caption{
    \def\@captype{table}\renewcommand{\@makecaption}{\@maketabcap}
    \oldcaption}
  \begin{sidewaystable}\centering}
  {\end{sidewaystable}\let\caption=\oldcaption}
\newenvironment{Sidewaystable*}
{%
  \let\oldcaption=\caption%
  \def\caption{
    \def\@captype{table}\renewcommand{\@makecaption}{\@maketabcap}
    \oldcaption}
  \begin{sidewaystable*}\centering}
  {\end{sidewaystable*}\let\caption=\oldcaption}
%    \end{macrocode}

% Not sure why we need to stick these lines in. 
%    \begin{macrocode}
\renewenvironment{figure}
                 {\renewcommand{\@makecaption}{\@makefigcap}\@float{figure}}
                 {\end@float}
\renewenvironment{figure*}
                 {\renewcommand{\@makecaption}{\@makefigcap}\@dblfloat{figure}}
                 {\end@dblfloat}
\renewenvironment{table}
                 {\renewcommand{\@makecaption}{\@maketabcap}\@float{table}}
                 {\end@float}
\renewenvironment{table*}
                 {\renewcommand{\@makecaption}{\@maketabcap}\@dblfloat{table}}
                 {\end@dblfloat}
%    \end{macrocode}

% Redefine |\@caption| so that we don't bother with the list of
% figures. This saves us having to worry about fragile commands in the
% captions. 
%    \begin{macrocode}
\long\def\@caption#1[#2]#3{%
  \par
%%  \addcontentsline{\csname ext@#1\endcsname}{#1}%
%%    {\protect\numberline{\csname the#1\endcsname}{\ignorespaces #2}}%
  \begingroup
    \@parboxrestore
    \if@minipage
      \@setminipage
    \fi
    \normalsize
    \@makecaption{\csname fnum@#1\endcsname}{\ignorespaces #3}\par
  \endgroup}
%    \end{macrocode}

% Make tables not quite so cramped. Allow |\hline\hline| to give a
% poor man's thick rule. 
%    \begin{macrocode}
\renewcommand{\arraystretch}{1.15}
\setlength{\doublerulesep}{0.4pt}
%    \end{macrocode}

% \begin{macro}{\tablecols}
% \begin{macro}{\@tablecols}
% \begin{macro}{\tabcaptioncontinued}
% |\tablecols| is a user command for setting the number of columns. It
% is used by commands that try and span the entire table (I think
% there must be a better way of acheiving
% this). |\tabcaptioncontinued| is 
% a user command for setting the caption of a continued table on the
% next page. 
%    \begin{macrocode}
\newcommand{\tablecols}[1]{\renewcommand{\@tablecols}{#1}}
\newcommand{\@tablecols}{1}
\newcommand{\tabcaptioncontinued}{%
  \multicolumn{\@tablecols}{c}{\addtocounter{table}{-1}TABLE~\thetable\ (CONTINUED)} \\ \\
  }
%    \end{macrocode}
% \end{macro}
% \end{macro}
% \end{macro}

% \subsubsection{Footnotes in tables}
% \label{sec:table-footnotes}
% \begin{macro}{\tablenotemark}
% \begin{macro}{\tablenote}
% \begin{macro}{\tablenoteskip}
% Old user commands |\tablenotemark{MARK}|,
% |\tablenote{MARK}{TEXT}|---now deprecated. 
%    \begin{macrocode}
\newcommand\tablenotemark[1]{$^\mathrm{#1}$} 
\newcommand\tablenote[2]{\if@proceedings\else\small\fi ${}^\mathrm{#1}$ #2} 
\newcommand\tablenoteskip{\\[0.5\baselineskip]} 
%    \end{macrocode}
% \end{macro}
% \end{macro}
% \end{macro}

% \begin{macro}{\tabnotemark}
% \begin{macro}{\tabnotetext}
% \begin{macro}{\tabnotetextTWO}
% Alternative way of doing footnotes in tables. 
%    \begin{macrocode}
\newcommand{\tabnotemark}[1]{\textsuperscript{\lowercase{#1}}}
%\newcommand{\tabnotetextTWO}[4]{\tabnotetext{#1}{#2}\tabnotetext{#3}{#4}}
\newcommand{\tabnotetextTWO}[4]{%
  \multicolumn{\@tablecols}{l}{%
    \parbox{0.5\tabnotewidth}{\small\textsuperscript{#1}#2}%
    \parbox{0.5\tabnotewidth}{\small\textsuperscript{#3}#4}%
    }%
  \\}
\newcommand{\tabnotetext}[2]{%
  \multicolumn{\@tablecols}{@{}p{\tabnotewidth}@{}}{\small\textsuperscript{#1}#2}\\}
\newlength{\tabnotewidth}
%    \end{macrocode}
% \end{macro}
% \end{macro}
% \end{macro}


% \begin{macro}{changemargin}
%    \begin{macrocode}
% Command to allow margins to be changed on the fly. Useful to allow
% tables and Figures to protrude into the margins. 
%
% \changes{3.24}{2007/09/14}{Added \texttt{changemargin} environment
% (copied from TeX FAQ).}
\newenvironment{changemargin}[2]{% 
  % From TeX FAQ http://www.tex.ac.uk/cgi-bin/texfaq2html?label=chngmargonfly
  \begin{list}{}{%
    \setlength{\topsep}{0pt}%
    \setlength{\leftmargin}{#1}%
    \setlength{\rightmargin}{#2}%
    \setlength{\listparindent}{\parindent}%
    \setlength{\itemindent}{\parindent}%
    \setlength{\parsep}{\parskip}%
  }%
  \item[]}{\end{list}}
%    \end{macrocode}
% \end{macro}

% \subsection{The bibliography and references}
% \label{sec:biblio}
% \begin{macro}{\beforebiblioskip}
% \begin{macro}{\afterbiblioskip}
% I'm not entirely sure why I introduced these two commands, since
% directly putting |\vspace| before and after
% |\begin{thebibliography}| seems to work well enough. 
%    \begin{macrocode}
\newlength{\beforebiblioskip}
\setlength{\beforebiblioskip}{0pt}
\newlength{\afterbiblioskip}
\setlength{\afterbiblioskip}{0pt}
%    \end{macrocode}
% \end{macro}
% \end{macro}
% \begin{macro}{thebibliography}
% The environment itself is done as starred section with extra
% vertical space after heading. This space is currently zero. 
% In the book style we do the cunning trick with rereading the file in
% order to get the references into the |\ADS@references| macro.
% 
% \changes{3.05}{2003/01/02}{Removed the \cs{beforebiblioskip} and
% \cs{afterbiblioskip} commands, since they cause unwanted vertical
% space in some instances.}  
%    \begin{macrocode}
\renewenvironment{thebibliography}{%
  \if@book\RereadBiblio\fi
%%  \vspace*{\beforebiblioskip}
  \section*{REFERENCES} \frenchspacing
%%  \vspace*{\afterbiblioskip}
  \small
%    \end{macrocode} 
% The references themselves are set as a list in order 
% to get the hanging indentation. They are done in a smaller font and
% with french spacing so we needn't worry about the distinction
% between sentence-ending periods and abbreviations. 
%    \begin{macrocode} 
  \list{\null}{\leftmargin 1.5em\labelwidth\z@\labelsep\z@\itemindent -1.5em
    \itemsep\z@\parsep\z@\usecounter{enumi}}
  \sloppy\clubpenalty4000\widowpenalty4000
  \sfcode`\.=1000\relax\@biblisttrue}
{\endlist\@biblistfalse\normalsize\@outputaddressesifnecessary}
%    \end{macrocode}
% \end{macro}
% \subsection{Uniquification of cross-references}
% \label{sec:uniquification}
% 
% \begin{macro}{\articlename}
% In order that the |\label{}| command in individual articles be
% guaranteed to give a unique anchor, we append the article name (held
% in the |\articlename| command) to each key. This requires
% redefinition of the following commands: |\cite|, |\nocite|, |\ref|,
% |\pageref|, |\label|. 
% \begin{macro}{\oldcite}
% \begin{macro}{\oldnocite}
% \begin{macro}{\oldref}
% \begin{macro}{\oldpageref}
% \begin{macro}{\oldlabel}
% The original versions of these commands are
% saved as |\oldcite|, |\oldnocite|, etc. 
%
% The uniquifying of cross-references is now always done, so 
% we don't get multiply-defined labels in the single article case
%    \begin{macrocode}
\newcommand{\articlename}{anon}
\AtBeginDocument{
  \let\oldcite=\cite 
  \let\oldnocite=\nocite 
  \let\oldref=\ref 
  \let\oldpageref=\pageref
  \let\oldlabel=\label
}
%    \end{macrocode}
% \end{macro}
% \end{macro}
% \end{macro}
% \end{macro}
% \end{macro}
% \end{macro}
%
% Uniquify all cross-references by prepending the article name.  This
% now has to go into |\AtBeginDocument| so it is working on the new
% versions that hyperref.sty has set up.  Then we need to put the
% |\@begofarticlehook| afterwards, so that the first page label is using
% the new definitions.  If amsmath is loaded, then we also need to patch
% the special label command that they use. \changes{4.0}{2019/11/18}{Fix labels for amsmath display equations, in case that package has been loaded.}
%    \begin{macrocode}
\AtBeginDocument{
  \renewcommand{\ref}[1]{\oldref{\articlename:#1}%
    \if@debug{\ClassInfo{rmaa}{Reference to \articlename:#1}}\fi}
  \renewcommand{\pageref}[1]{\oldpageref{\articlename:#1}%
    \if@debug{\ClassInfo{rmaa}{Page reference to \articlename:#1}}\fi}
  \renewcommand{\label}[1]{\oldlabel{\articlename:#1}%
    \if@debug{\ClassInfo{rmaa}{Label for \articlename:#1}}\fi}%
  \ifdef{\label@in@display}{%
    \let\oldlabel@in@display=\label@in@display
    \renewcommand{\label@in@display}[1]{\oldlabel@in@display{\articlename:#1}%
      \if@debug{\ClassInfo{rmaa}{AMSTeX Label for \articlename:#1}}\fi}%
  }{}
  \@begofarticlehook
}%
%    \end{macrocode}
% Citations are more complicated since they can contain a list of
% keys. These are really only here for use with the rmaacite package,
% which is utterly deprecated (02 Jan 2003) now that I have the natbib
% code working. The command \cs{cite} should \emph{not} be used with
% natbib: use \cs{citep}, \cs{citet}, etc.\@ instead. 
%    \begin{macrocode}
\renewcommand{\nocite}[1]{\@bsphack
  \@for\@citeb:=#1\do{%
    \edef\@citeb{\expandafter\@firstofone\@citeb:\articlename}%
    \if@filesw\immediate\write\@auxout{\string\citation{\@citeb}}\fi
    \@ifundefined{b@\@citeb:\articlename}{\G@refundefinedtrue
      \@latex@warning{Citation `\@citeb' undefined}}{}}%
  \@esphack}
\renewcommand{\cite}[1]{\oldcite{#1:\articlename}}
%   \end{macrocode}


% Horrible fiddling with the bibliography citation mechanism. 
% This has to be globally available in case rmaacite puts it in the
% \texttt{.aux} file. 
%    \begin{macrocode}
\def\lbibcite#1#2#3{\global\@namedef{b@#1}{#2}
                    \global\@namedef{by@#1}{#3}}
%    \end{macrocode}

% \begin{macro}{\oldbibitem}
% \begin{macro}{\@rmaaBibArg}
% \begin{macro}{\@rmaaSpace}
% \begin{macro}{rmaabib}
% Of course, |\bibitem| needs to get the uniquification treatment as
% well. While we are at it, we redefine it to use consecutive integers
% for the key if the user leaves it blank (using the counter
% |rmaabib|). This allows one to put 
% |\bibitem{}| for all your references without getting annoying
% messages about multiply-defined labels.
% 
% We also take the opportunity to insert the optional argument to
% \cs{bibitem} if it is missing in the form Author(DATE) so that natbib
% doesn't try to do numerical citations. Natbib already provides a
% hook \cs{@extra@b@citeb} to help uniquify the citation key, so we
% just plug into that. However, it puts it \emph{after} the original
% key, not before as I had had it. Easiest solution is to change the
% way that I do it. 02 Jan 2003. Note that the \cs{@extra@b@citeb}
% does not get written to the .aux  file in the \cs{citation} command
% but it does work. 
% \changes{3.14}{06 Feb 2003}{Reinstalled old definition of \cs{bibitem}
%  for the case where the nonatbib option is requested, so that the
%  rmaacite package will work. This means that you can't use rmaacite
%  and natbib in the same volume.}
%    \begin{macrocode}
\if@natbib\def\@extra@b@citeb{:\articlename}\fi
%%\else
\let\oldbibitem\bibitem
\newcommand{\@rmaaBibArg}{}
\newcommand{\@rmaaSpace}{ }
\newcounter{rmaabib}
\if@natbib
\renewcommand{\bibitem}[2][ABC(2003)]{%
  \renewcommand{\@rmaaBibArg}{#2}
  \ifx\@rmaaBibArg\@rmaaEmpty
     \addtocounter{rmaabib}{1}
     \oldbibitem[#1]{\thermaabib:\articlename}
  \else\ifx\@rmaaBibArg\@rmaaSpace
     \addtocounter{rmaabib}{1}
     \oldbibitem[#1]{\thermaabib:\articlename}
  \else 
     \oldbibitem[#1]{#2:\articlename}
  \fi\fi
  \renewcommand{\@rmaaBibArg}{}
  }
\else
\renewcommand{\bibitem}[1]{%
  \renewcommand{\@rmaaBibArg}{#1}
  \ifx\@rmaaBibArg\@rmaaEmpty
     \addtocounter{rmaabib}{1}
     \oldbibitem{\articlename:\thermaabib}
  \else\ifx\@rmaaBibArg\@rmaaSpace
     \addtocounter{rmaabib}{1}
     \oldbibitem{\articlename:\thermaabib}
  \else
     \oldbibitem{\articlename:#1}
  \fi\fi
  \renewcommand{\@rmaaBibArg}{}
  }
\fi
%%\fi
%    \end{macrocode}
% \end{macro}
% \end{macro}
% \end{macro}
% \end{macro}

% \begin{macro}{oldpcite}
% \begin{macro}{oldscite}
% Then, we need to repeat the above scheme, but for the |rmaacite|
% definition. 
%    \begin{macrocode}
\newcommand{\@redefinebibitem}{%
  \let\oldbibitem\bibitem
  \let\oldcite\cite
  \let\oldpcite\pcite
  \let\oldscite\scite
  \def\bibitem[##1<##2>]##3{%
    \renewcommand{\@rmaaBibArg}{##3}
    \ifx\@rmaaBibArg\@rmaaEmpty
       \addtocounter{rmaabib}{1}
       \oldbibitem[##1<##2>]{\thermaabib:\articlename}
    \else\ifx\@rmaaBibArg\@rmaaSpace
       \addtocounter{rmaabib}{1}
       \oldbibitem[##1<##2>]{\thermaabib:\articlename}
    \else 
       \oldbibitem[##1<##2>]{##3:\articlename}
    % trust the author here
%%       \oldbibitem[##1<##2>]{##3}
    \fi\fi
    }
  }
%    \end{macrocode}
% \end{macro}
% \end{macro}


% \subsection{HTML Table of Contents}
% \label{sec:hyper-TOC}
% \begin{macro}{\HyperTOCEntry}
% \begin{macro}{\HyperTOCpart}
% \begin{macro}{\HyperTOCmeeting}
% \begin{macro}{\HyperTOCeditors}
% \begin{macro}{\HyperTOCpreamble}
% These commands are used to write information to the \texttt{.htoc}
% file, which is used in generating a Web Page of the Table of
% Contents. Since they only ever appear preceded by |\protect| in the
% argument to an |\addtocontents| command,  they should never need to
% be expanded by the macros. However, we define them here just in case
% \dots
%    \begin{macrocode}
\newcommand{\HyperTOCEntry}[3]{\relax}
\newcommand{\HyperTOCpart}[1]{\relax}
\newcommand{\HyperTOCmeeting}[3]{\relax}
\newcommand{\HyperTOCeditors}[1]{\relax}
\newcommand{\HyperTOCpreamble}{\relax}
%    \end{macrocode}
% \end{macro}
% \end{macro}
% \end{macro}
% \end{macro}
% \end{macro}
%

% \subsection{Formatting the poster abstracts} 
% These can now be used without using the book option.
%    \begin{macrocode}

\newcommand{\@makeminititle}{%
  \newpage
  \null 
  \vskip 3em%
  \begin{center}%
    {\@title \par}%
  \end{center}
  \vskip 3em%
  }

\newif\if@abstractsads
\@abstractsadsfalse
\newcommand\abstractsADS{\@abstractsadstrue}

\if@book\RequirePackage{multicol}\fi
\newenvironment{abstracts}[1][Abstracts of Contributed Papers]{%
  %% First, make sure we start on an odd page (06 Jun 2000)
  \if@book\clearpage\thispagestyle{empty}\cleardoublepage\fi
  %%
  \renewcommand{\articlename}{ABSTRACTS}
  \title{#1}\shortauthor{Abstracts}\shorttitle{Abstracts}
  \listofauthors{}
  \if@book
  \addcontentsline{toc}{chapter}{\protect\@fiddleauthors{\@TOCtitle}{\rmaaTOCauthors}}%
  \addtocontents{htc}{\protect\HyperTOCEntry{\@rawtitle}{\@rawauthors}{\articlename}}
  \fi
  \pagestyle{rmaa} % this now seems necessary (06 Jun 2000)
  \thispagestyle{rmaafirstpage}
%%  \twocolumn[\@makeminititle]
  \renewcommand{\abstract}[1]{%
    \renewcommand\ADS@abstract{##1}%
    \renewcommand{\@abstract}{##1}\@writeabstract%
  }
  \renewcommand{\title}[2][]{%
    \gdef\ADS@authors{}%
    \gdef\ADS@affil{}%
    \gdef\ADS@references{}%
    \gdef\ADS@keywords{}%
    \renewcommand{\@rmaaArg}{##1}
    \ifx\@rmaaArg\@rmaaEmpty
       \renewcommand{\@TOCaftertitle}{}
    \else
       \renewcommand{\@TOCaftertitle}{ (\@rmaaArg)}
    \fi
    \renewcommand{\@title}{\rmaa@TOCTitleStyle{##2}}%
    \renewcommand{\@TOCtitle}{\rmaa@TOCTitleStyle{##2}\@TOCaftertitle}%
    \renewcommand{\rmaaIndexTitle}{\rmaa@IndexTitleStyle{##2\@TOCaftertitle.}}%
    \renewcommand\ADS@title{##2}%
    \xdef\ADS@firstpage{\@arabic{\c@page}}%
    }
  \renewcommand{\author}[1]{\let\@thanks\@empty\renewcommand\@author{##1} }
  \renewcommand{\footnotetext}[2][XX]{\par\noindent\textsuperscript{##1}~##2}
  \renewcommand{\footnoterule}{\kern-3\p@ \hrule \@width 3em \kern 2.6\p@}
  \if@book\@begofarticlehook\fi
  \if@book\begin{multicols}{2}[\@makeminititle]\else\@makeminititle\fi
  \begingroup
  \setcounter{footnote}{0}
  }% and at the end of the environment....
  {\endgroup\if@book\end{multicols}\listofabstracts\label{RMAALastPage}\clearpage\fi}
%% auxiliary commands used:
%% \@writeabstract does all the actual formatting 
\newcommand{\@writeabstract}{%
  \if@book%
  \addcontentsline{loa}{abstract}{\protect\@fiddleauthors{\@TOCtitle}{\rmaaTOCauthors}}%
  \fi%
  \begin{center}\@title\par\@author\end{center}
  \@abstract
  \par\noindent\footnoterule\nobreak\footnotesize\@thanks
  \par\noindent\hrulefill\goodbreak
  \vskip 1cm \@plus0.5cm
  %% write ADS entry
  \if@book\if@abstractsads\ADS@dumpdata\fi\fi%
  \endgroup % this pair is to make each abstract a self-contained group
  \begingroup
  \setcounter{footnote}{0}
  }

\if@book
%% List of abstracts (stuff written to .loa file)
\newcommand{\listofabstracts}{%
  \begin{multicols}{2}[%
    \vskip 1.5cm \vfill % leave at least 1.5 cm and go to bottom of page if possible
    \centerline{\MakeUpperCase{List of Abstracts}}
    \vskip 1.5cm \@plus0.5cm
    ][6cm] % this makes it start a new page if there is not at least 6cm of space
    \renewcommand{\@fiddleauthors}[2]{##1\hskip1em plus2em minus0.4em\penalty-100  ##2}
    \parindent\z@
%%    \parskip.5\p@ \@plus .7\p@\relax 
    \parskip3\p@ \@plus 1\p@ \@minus 1\p@\relax % allow quite a lot of
                                % leeway here to avoid breaks inside
                                % entries 
    \renewcommand{\\}{\relax} % temporarily turn off explicit line breaks just in case
    \renewcommand{\altaffilmark}[1]{\relax} % ditto with this
    \@starttoc{loa}
  \end{multicols}
  }
\newcommand*{\l@abstract}[2]{%
  \par\penalty-1000
  #1\hskip 1em\penalty0\null\nobreak\hfill#2}
\fi

%    \end{macrocode}
 
% \subsection{Code for typesetting an entire volume of conference
% proceedings} \label{sec:whole-volume}
%    \begin{macrocode}

\if@book
%% environment for including the individual articles
\newcommand{\@fiddleauthors}[2]{\relax}
\renewcommand{\theequation}{\@arabic\c@equation}
\newenvironment{articles}%
{ \begingroup % make it a group so that author definitions do not
              % spill over
  \gdef\ADS@authors{}
  \pagestyle{rmaa}
  \renewenvironment{document}{%
    \@begofarticlehook
    \addtocounter{chapter}{1}%
    \addcontentsline{toc}{chapter}{%
      \protect\@fiddleauthors{\@TOCtitle}{\rmaaTOCauthors}}%
    \addtocontents{htc}{\protect\HyperTOCEntry{\@rawtitle}{\@rawauthors}{\articlename}}
    }{\relax}
  \renewcommand{\documentclass}[2][foo]{\relax} 
  \renewcommand{\usepackage}[2][foo]{%
    \def\@packagename{##2}\def\@rmaacitename{rmaacite}%
    \ifx\@packagename\@rmaacitename
    \@rmcitetrue\ClassInfo{rmaa}{This paper uses rmaacite}
    \makeatletter%%
%% This is file `rmaacite.sty',
%% generated with the docstrip utility.
%%
%% The original source files were:
%%
%% rmaa.dtx  (with options: `rmaacite')
%% 
%% PLEASE DO NOT EDIT THIS FILE! It is generated automatically from
%% the file rmaa.dtx using the accompanying Makefile and rmaa.ins
%% Any hand-edited changes will be lost
%% 
%% A LaTeX document class for the RMxAA
%% 
%% Author: William Henney (w.henney@crya.unam.mx)
%% 
%% Latest version available from: http://www.crya.unam.mx/rmaa
%% 
%% Copyright (C) 1999-2010 Instituto de Astronomia, UNAM, Mexico
%% 
%% This program is free software; you can redistribute it and/or
%% modify it under the terms of the GNU General Public License
%% as published by the Free Software Foundation; either version 2
%% of the License, or (at your option) any later version.
%% 
%% This program is distributed in the hope that it will be useful,
%% but WITHOUT ANY WARRANTY; without even the implied warranty of
%% MERCHANTABILITY or FITNESS FOR A PARTICULAR PURPOSE.  See the
%% GNU General Public License for more details.
%% 
%% You should have received a copy of the GNU General Public License
%% along with this program; if not, write to the Free Software
%% Foundation, Inc., 59 Temple Place - Suite 330, Boston, MA  02111-1307, USA.






\def\bibitem[#1<#2>]#3{\par\hangindent=1.5em\hangafter=1
  \if@filesw
  {\def\protect##1{\string ##1\space}
    \immediate\write\@auxout{\string\lbibcite{#3}{#1}{#2}}}
  \fi}
\@redefinebibitem % what it says
\def\lbibcite#1#2#3{\global\@namedef{b@#1}{#2}
                    \global\@namedef{by@#1}{#3}}

\def\@bcite#1#2{({#1\if@tempswa , #2\fi})}
\def\@pcite#1#2{#1\if@tempswa , #2\fi}
\def\@brcite{\ (}
\def\@ebrcite{)}
\def\@fmttop#1{\@fmtcite[\csname #1\endcsname]}
\def\@citex[#1]#2{%
%%  \if@filesw\immediate\write\@auxout{\string\citation{#2}}\fi
%% WJH added 16 Jan 2002
  \@for\@citeb:=#2\do{%
    \edef\@citeb{\expandafter\@firstofone\@citeb:\articlename}%
    \if@filesw\immediate\write\@auxout{\string\citation{\@citeb}}\fi}%
%% to here
  \def\@citea{}%
  \@cite{\@for\@citeb:=#2\do
    {%
      \edef\@citeb{\expandafter\@firstofone\@citeb:\articlename}%
      \@citea\def\@citea{;\ }\@ifundefined
      {b@\@citeb}{{\bf ?}\@warning {Citation `\@citeb' on page \thepage \space undefined}}%
      {\csname b@\@citeb\endcsname\@citebro\csname by@\@citeb\endcsname\@citebrc}}}%
  {#1}}

\def\cite{\@ifnextchar [{\let\@citebro=\space
                          \let\@citebrc=\relax
                          \let\@cite=\@bcite\@tempswatrue\@citex}
                        {\let\@citebro=\space
                          \let\@citebrc=\relax
                          \let\@cite=\@bcite\@tempswafalse\@citex[]}}
\def\pcite{\@ifnextchar [{\let\@citebro=\space
                          \let\@citebrc=\relax
                          \let\@cite=\@pcite\@tempswatrue\@citex}
                        {\let\@citebro=\space
                          \let\@citebrc=\relax
                          \let\@cite=\@pcite\@tempswafalse\@citex[]}}
\def\scite{\@ifnextchar [{\let\@citebro=\@brcite
                          \let\@citebrc=\@ebrcite
                          \let\@cite=\@pcite\@tempswatrue\@citex}
                        {\let\@citebro=\@brcite
                          \let\@citebrc=\@ebrcite
                          \let\@cite=\@pcite\@tempswafalse\@citex[]}}

\renewenvironment{thebibliography}
  {\section*{REFERENCES}\vspace*{\afterbiblioskip}\small
    \if@book\RereadBiblio\fi
    \sloppy\clubpenalty4000\widowpenalty4000
    \sfcode`\.=1000\relax
    \setlength\parindent{\z@}\setlength\parskip{\z@}}
  {\relax\normalsize\@outputaddressesifnecessary}








\endinput
%%
%% End of file `rmaacite.sty'.
\makeatother % read it in here so that all changes
    \else                                         % are local to paper
    \if@debug\ClassWarning{rmaa}{Package \@packagename\space was requested by \articlename\space}\fi
    \fi} 
  %% Don't allow individual articles to mess with the pages, volume or year
  \renewcommand{\SetFirstPage}[1]{\relax}
  \renewcommand{\SetYear}[1]{\relax}
  \renewcommand{\SetVolume}[1]{\relax}
  \renewcommand{\SetConfTitle}[1]{\relax}
  %% Don't allow individual articles to change these
  \renewcommand{\nonstopmode}{\relax}
  \renewcommand{\scrollmode}{\relax}
  \renewcommand{\outputplates}{\relax}
  %% XXXX need to check if using rmaacite and act accordingly
  %% avoid multiple declarations of the same length
  \let\oldnewlength=\newlength
  \renewcommand{\newlength}[1]{%
    \edef\reserved@a{\expandafter\@gobble\string ##1}%
    \@ifundefined{\reserved@a}{%
      \oldnewlength{##1}}{%
      \if@debug{\ClassInfo{rmaa}{\reserved@a\space already defined: doing nothing}}\fi%
      \relax}}%
  \renewcommand\newcommand{\providecommand}
  }%
{\endgroup\onecolumn}
     

%    \end{macrocode}
% Command to be used in `articles' environment for including an
% article
%    \begin{macrocode}
\newcommand{\article}[2][]{%
  % optional first argument is extra text to put in brackets after title in TOC
  % e.g. ``Invited Talk'' or ``Poster''
  \renewcommand{\articlename}{#2}
  % reset loads of counters
  \setcounter{footnote}{0}
  \setcounter{figure}{0}
  \setcounter{table}{0}
  \setcounter{section}{0}
  \setcounter{equation}{0}
  \setcounter{rmaabib}{0}
  \setcounter{author}{0}
  \setcounter{AUTHOR}{1}
  \gdef\ADS@authors{}%
  \gdef\ADS@affil{}%
  \gdef\ADS@affiliations{}%
  \gdef\ADS@references{}%
  \@firstauthortrue
  \@firstaffiltrue
  \@firstaltaffiltrue
  \@firstkwdtrue
  \@rmcitefalse
  \if@suppressfulladdresses\else\@needaddressestrue\fi
  \renewcommand{\@rmaaArg}{#1}
  \ifx\@rmaaArg\@rmaaEmpty
     \renewcommand{\@TOCaftertitle}{}
  \else
     \renewcommand{\@TOCaftertitle}{ (#1)}
  \fi
%    \end{macrocode}
% In the main journal articles must start on an odd page. 
%    \begin{macrocode}
  \if@journal\clearpage\thispagestyle{empty}\cleardoublepage\fi
%% Now input the file, then run the end-of-article hook (reinstated 28
  %% Sep 2001)
  {\input{\@articledir/\articlename.tex}\@endofarticlehook}}
%    \end{macrocode}
% Command to separate articles into sessions. Write an entry in the
% TOC, makes sure we start on an odd page (added 06 Jun 2000), then do
% a title page if \cs{SessionTitlePages} was called. 
% \changes{3.12}{04 Feb 2003}{Added option to have title page for each
% session. Turn it on with \cs{SessionTitlePages}.}
%    \begin{macrocode}
\newcommand{\session}[1]{
  %% this does nothing except write an entry in the TOC
  \addcontentsline{toc}{part}{\rmaa@TOCPartStyle{#1}}
  \addtocontents{htc}{\protect\HyperTOCpart{#1}}
  %% and also makes sure we start on an odd page (added 06 Jun 2000)
  \clearpage\thispagestyle{empty}\cleardoublepage
  \stepcounter{part}
  \if@sessiontitlepages{%
    \onecolumn
    \thispagestyle{empty}
    \renewcommand{\articlename}{Section\thepart TitlePage}
    \label{RMAAFirstPage}
    \make@session@title{#1}
    \label{RMAALastPage}
    \clearpage\thispagestyle{empty}\cleardoublepage
  }\fi
  }
\newcommand{\make@session@title}[1]{%
  {%
    \setlength\parindent{0pt}
    \vspace*{\stretch{1}}
    \rule{\linewidth}{1mm}
    \begin{flushright}\Huge\bfseries\boldmath
      #1
    \end{flushright}
    \rule{\linewidth}{1mm}
    \vspace*{\stretch{2}}
  }
}
%    \end{macrocode}
% Command to turn on separate title pages for each session. 
%    \begin{macrocode}
\newcommand{\SessionTitlePages}{\@sessiontitlepagestrue}
\newcommand{\BoldSessionTitles}{%
  \renewcommand{\rmaa@TOCPartStyle}[1]{\textbf{\rmaa@TOCTitleStyle{##1}}}
  }
\newcommand{\ArticleDir}[1]{\renewcommand{\@articledir}{#1/}\graphicspath{{1/}}}
\newcommand{\@articledir}{.}
%\newcommand{\ps@rmaafront}{\ps@rmaa}% page style for front matter (same as ever, actually)
\newcommand{\ps@rmaafront}{% page style for front matter:  page # in each foot
  \renewcommand\@oddfoot{\hfil{\thepage}}%
  \renewcommand\@evenfoot{{\thepage}\hfil}%
  % shorttitle in both heads
  \renewcommand\@evenhead{\hfil\@shorttitle\hfil}%
  \renewcommand\@oddhead{\@evenhead}
  }
\newcommand{\ps@rmaafrontfirstpage}{% page style for first page of front matter
  % page # in each foot, on outside edge
  \renewcommand\@oddfoot{\hfil\thepage}%
  \renewcommand\@evenfoot{\thepage\hfil}%
  % nothing in the even heads
  \renewcommand\@evenhead{\relax}%
  \renewcommand\@oddhead{\relax}%
  }
%% slight modification to the standard \frontmatter and \mainmatter
%% commands, just to set the \pagestyle
\renewcommand\frontmatter{%
  \latexonecolumn
  \cleardoublepage
  \@frontmattertrue
  \@mainmatterfalse
  \pagenumbering{roman}
  \pagestyle{rmaafront}
  \addtocounter{page}{2} % leave space for the title page
  }
\renewcommand\mainmatter{%
  \pagestyle{empty}
  \if@rmtwocol\twocolumn\fi
  \cleardoublepage
  \@frontmatterfalse
  \@mainmattertrue
  \pagenumbering{arabic}
  \pagestyle{rmaa}
  }
\renewcommand\backmatter{%
  \clearpage\pagestyle{empty}\cleardoublepage
  \@mainmatterfalse
  %% put extra space in the TOC
  \addtocontents{toc}{\bigskip}
  }
%    \end{macrocode}
%
% \subsection{Front sections}\label{sec:front-sections}
% \begin{macro}{\notitle}
% \begin{macro}{\shortname}
% \begin{macro}{frontsection}
% The environment \texttt{frontsection} is for doing a section in the
% front matter (preface, list of participants, or whatever). It has 
% one optional argument:
% \begin{enumerate}
% \item[\#1] Short title (no embedded spaces) for
% use in making PS files. E.g., ``list\_of\_participants''.   
% \end{enumerate}
% and two mandatory arguments: 
% \begin{enumerate}
% \item[\#2] Title of section (e.g., `Preface')
% \item[\#3] Commands to be evaluated before processing the contents
% of the environment. This is mainly useful for putting in the things
% that would go in the preamble of a normal article  (e.g.,
% |\listofauthors|, |\indexauthor|, |\notitle|) 
% \end{enumerate}
%    \begin{macrocode}
\newcommand{\notitle}{\@titlevisiblefalse} \newcommand{\shortname}{}
\newenvironment{frontsection}[3][]{%
  \begingroup
  \renewcommand{\shortname}{#1}
  \ifx\shortname\@rmaaEmpty
     \renewcommand{\articlename}{#2}
  \else
     \renewcommand{\articlename}{\shortname}
  \fi
  \@begofarticlehook
  \@titlevisibletrue % default
  \thispagestyle{rmaafrontfirstpage}
  \listofauthors{} % blank out this in case it is not specified
  \title{{#2}}\shorttitle{#2}\shortauthor{#2}
  #3 % evaluate the command given as second mandatory arg
  \addcontentsline{toc}{chapter}{\protect\@fiddleauthors{\@TOCtitle}{\rmaaTOCauthors}}%
  \addtocontents{htc}{\protect\HyperTOCEntry{\@rawtitle}{\@rawauthors}{\articlename}}
  \if@titlevisible
     \vspace*{1cm}\centerline{\rmaa@FrontSecStyle{\@title}}\vspace*{1cm}
  \fi
  \setlength{\parskip}{3ex}
  }{\frontlabel{RMAALastPage}\clearpage\endgroup}
%    \end{macrocode}
% \end{macro}
% \end{macro}
% \end{macro}
%    \begin{macrocode}
%% Definition of the command \frontlabel, which rights out something like `F4' to
%% the .aux file, instead of `\textlatin  {i}'. This is so it is easy to post-process
%% the file for making PS files of individual sections
\newcommand{\frontlabel}[1]{%
  \@bsphack
  \protected@write\@auxout{}%
     {\string\newlabel{\articlename:#1}{{}{F\@arabic{\c@page}}}}
  \@esphack
  }
%% User command to say X, Y, Z,\\ Editors of the Proceedings
\newcommand{\EditorSignature}{\vspace*{\parskip}\rightline{\@editors}%
  \rightline{Editors of the proceedings}}
%% Fiddle with \thesection, etc. so that they don't include the chapter number
%% Basically go back to the `article' definitions of these
\renewcommand{\thesection}{\@arabic\c@section}
\renewcommand{\thefigure}{\@arabic\c@figure}
\renewcommand{\thetable}{\@arabic\c@table}

\newcommand{\@oldargs}{}
\newcommand{\@newargs}{}
%% add TOC entry for each paper
%\renewcommand{\title}[1]%
%{\renewcommand{\@title}{\rmaa@TitleStyle{#1}}%
%\addtocounter{chapter}{1}\addcontentsline{toc}{chapter}{#1}%
%}

%% Formatting of the table of contents
%% dimensions of various parts
\newlength{\TOCtextwidth}\setlength{\TOCtextwidth}{12.65cm}
\newlength{\TOCpagenumwidth}\setlength{\TOCpagenumwidth}{1.3cm}
\newlength{\TOCentrywidth}\setlength{\TOCentrywidth}{\TOCtextwidth}
\addtolength{\TOCentrywidth}{-\TOCpagenumwidth} 
\newlength{\TOCextramargin}\setlength{\TOCextramargin}{1.25cm}
\newcommand{\TOCbox}[1]{\hspace{\TOCextramargin}\parbox{\TOCtextwidth}{#1}\par}
\newlength{\TOCextramarginx}\setlength{\TOCextramarginx}{\TOCextramargin}\addtolength{\TOCextramarginx}{4.3cm}
\newlength{\TOCtextwidthx}\setlength{\TOCtextwidthx}{\TOCtextwidth}\addtolength{\TOCtextwidthx}{-4.3cm}
\newcommand{\TOCboxx}[1]{\hspace{\TOCextramarginx}\parbox{\TOCtextwidthx}{#1}\par}

\newcommand{\RMAAbannerimage}{rmsc}
\renewcommand\tableofcontents{%
  \renewcommand{\articlename}{CONTENTS}
  \begingroup
  \@begofarticlehook
  \if@twocolumn
     \@restonecoltrue\onecolumn
  \else
     \@restonecolfalse
  \fi
  %% 9pt for journal style
  \if@proceedings\else\small\fi
  \if@newstyle\small\fi %% same for new style
  \if@proceedings\thispagestyle{rmaafrontfirstpage}\else\thispagestyle{empty}\fi
  \shortauthor{Contents}\shorttitle{Contents}
  \vspace*{18mm}
  \if@printbanner
  %% put in the banner header (need to get better image of this)
     \TOCbox{\hspace*{-0.15cm}\includegraphics[width=0.9\TOCtextwidth]{\RMAAbannerimage}}
     \vspace*{8mm}
  %% volume and publication date (should really be in funny font)
     \TOCboxx{\@volumen\\[2mm] \@mes\ \@anyo}
  \else
  %% put in the banner header (need to get better image of this)
     \TOCbox{\hspace*{-0.15cm}\includegraphics[width=0.9\TOCtextwidth]{\RMAAbannerimage}}
     \vspace*{12.5mm}
  %% volume and publication date (should really be in funny font)
     \TOCboxx{~\\[2mm] ~}
  \fi
  \vspace*{16mm}
  %% This is the command that fiddles with the TOC entry so that the
  %% title is flushleft but the authors are flushright
  \renewcommand{\@fiddleauthors}[2]{%
    % first put the paper title
    \renewcommand{\\}{\relax} % temporarily turn off explicit line breaks just in case
    \renewcommand{\altaffilmark}[1]{\relax} % ditto with this
    \renewcommand{\textsuperscript}[1]{\relax} % ditto with this 
    ##1
    % now save the author list in a box
    \sbox\@tempboxa{##2}%
    % check that box is shorter than the total width available for the entry
    \ifdim \wd\@tempboxa <\TOCentrywidth
       % if so, then put it right here, flush against right margin
       \penalty10000\hskip1em plus10em\penalty0\hbox{}\nobreak\hfill\box\@tempboxa
       % Note that we force a space of at least 1em between title and authorlist, 
       % stretchable to 11em since we may have very short titles but still with an 
       % authorlist too long to go on the same line.  The only linebreak allowed is 
       % directly before the authorlist.
    \else
       % if it is too long, we set the original authorlist text in its own parbox
       \par\vskip1ex\parbox[b]{\TOCentrywidth}{\raggedleft ##2}
    \fi
    }
  \@starttoc{toc}%
  \if@restonecol\twocolumn\fi
  \frontlabel{RMAALastPage}
  \clearpage
  \thispagestyle{empty}
  \cleardoublepage % make sure we start on an odd page
  \endgroup

  %%% ALSO, do the HTML TOC
  %% The file that will be used to make the HTML TOC
  \@starttoc{htc}
  %% All the stuff in this file does nothing in the book style
  \addtocontents{htc}{\protect\HyperTOCmeeting{\@conftitle}{\@reunionplace}{\@reuniondate}}
  \addtocontents{htc}{\protect\HyperTOCeditors{\@editors}}
  \addtocontents{htc}{\protect\HyperTOCpreamble}
  \addtocontents{htc}{\protect\HyperTOCEntry{Table of Contents}{}{\articlename}}

  %% and the ADS file
  \ADS@openfile
  }
%% This is now much simpler
\renewcommand\l@chapter[2]{%TOC entry for articles
  \ifnum \c@tocdepth >\m@ne
    \vskip 1.0em \@plus\p@
    \TOCbox{%
      \parbox[b]{\TOCentrywidth}{#1}\parbox[b]{\TOCpagenumwidth}{\hfill#2}
    }
  \fi} 
\renewcommand\l@part[2]{%TOC entry for sessions: centered, no page number
  \ifnum \c@tocdepth >\m@ne
    \vskip 1.5em \@plus1.5em \goodbreak
    \TOCbox{%
      \parbox[b]{\TOCentrywidth}{\hfil#1\hfil}\hfill
    }
  \fi\nobreak} 

%% Formatting theindex
%% strongly encourage page breaks between index entries
\renewcommand{\@idxitem}{\par\penalty-1000\hangindent.45cm} 
\renewcommand{\indexspace}{\par\penalty-1000\vskip 10\p@ \@plus15\p@ \@minus3\p@\relax}
\renewenvironment{theindex}{%
  %% first, make sure we start on an odd page (06 Jun 2000)
  \clearpage\thispagestyle{empty}\cleardoublepage
  %% this looks a bit like an article but is two columns and without abstract, etc
  \renewcommand{\articlename}{INDEX}
  \title{Author Index}\shortauthor{Author Index}\shorttitle{Author Index}
  \listofauthors{}
  \addcontentsline{toc}{chapter}{\protect\@fiddleauthors{\@TOCtitle}{\rmaaTOCauthors}}%
  \addtocontents{htc}{\protect\HyperTOCEntry{\@rawtitle}{\@rawauthors}{\articlename}}
  \pagestyle{rmaa}
  \thispagestyle{rmaafirstpage}
  \@begofarticlehook
  %% 9pt for journal style
  \if@proceedings\else\small\fi
  \if@newstyle\small\fi
  \twocolumn[\@makeminititle]
  %% vertical spacing (see also definition of \@idxitem and \indexspace
  \parskip\z@ \@plus .6\p@\relax % allow quite a lot of leaway here to avoid breaks inside entries
  %% horizontal spacing, breaking and glue parameters (see also the stuff in the definition 
  %% of \indexauthor
  \parindent\z@ 
  \exhyphenpenalty10000 % don't break on explicit hyphens, it always looks bad
  \rightskip=0pt plus.2em % allow a tiny bit of leeway at the ends of lines
  \spaceskip.3333em plus.5em minus.15em % a lot of leeway between words
  \xspaceskip.5em plus.5em minus.4em % this seems to have no effect
  \let\item\@idxitem
  \renewcommand{\\}{\relax} % temporarily turn off explicit line breaks just in case
  \renewcommand{\altaffilmark}[1]{\relax} % ditto with this
  \renewcommand{\textsuperscript}[1]{\relax} % ditto with this
  }{\label{RMAALastPage}\clearpage} % clearpage necessary so that we know about \shortauthor



%% end definitions for book styles
\fi

%    \end{macrocode}
% In order to do away with plates, we sometimes want a figure to
% appear in color only in the online version. The following command
% can bve used in a figure caption to typeset a note to this effect. 
%    \begin{macrocode}
\newcommand\ColorOnlineOnly{%
  \textsc{Note: this figure is available in color in the electronic
    version of this article, obtainable from}
  \texttt{http://\discretionary{}{}{}www.\discretionary{}{}{}astroscu.\discretionary{}{}{}unam.\discretionary{}{}{}mx/\discretionary{}{}{}\~{}rmaa/}.}
%% Plates
%% Define 2 counters: plate is the number in the whole document, plateloop is used 
%% internally by \outputplates
\newcounter{plate}
\newcounter{plateloop}
%% Currently limited to 15 plates in document
%% when we do a label, we need to write to \newlabel entries to the .aux file, one 
%% that says RMAAPlateX so that it can be referred to in \ouputplates, and one that
%% says articlename:RMAAPlateX so that the dvips program can find it and asociate it
%% with the right article
\newcommand{\@platelabeli}{\oldlabel{RMAAPlate1}\label{RMAAPlate1}}
\newcommand{\@platelabelii}{\oldlabel{RMAAPlate2}\label{RMAAPlate2}}
\newcommand{\@platelabeliii}{\oldlabel{RMAAPlate3}\label{RMAAPlate3}}
\newcommand{\@platelabeliv}{\oldlabel{RMAAPlate4}\label{RMAAPlate4}}
\newcommand{\@platelabelv}{\oldlabel{RMAAPlate5}\label{RMAAPlate5}}
\newcommand{\@platelabelvi}{\oldlabel{RMAAPlate6}\label{RMAAPlate6}}
\newcommand{\@platelabelvii}{\oldlabel{RMAAPlate7}\label{RMAAPlate7}}
\newcommand{\@platelabelviii}{\oldlabel{RMAAPlate8}\label{RMAAPlate8}}
\newcommand{\@platelabelix}{\oldlabel{RMAAPlate9}\label{RMAAPlate9}}
\newcommand{\@platelabelx}{\oldlabel{RMAAPlate10}\label{RMAAPlate10}}
\newcommand{\@platelabelxi}{\oldlabel{RMAAPlate11}\label{RMAAPlate11}}
\newcommand{\@platelabelxii}{\oldlabel{RMAAPlate12}\label{RMAAPlate12}}
\newcommand{\@platelabelxiii}{\oldlabel{RMAAPlate13}\label{RMAAPlate13}}
\newcommand{\@platelabelxiv}{\oldlabel{RMAAPlate14}\label{RMAAPlate14}}
\newcommand{\@platelabelxv}{\oldlabel{RMAAPlate15}\label{RMAAPlate15}}
\newcommand{\@outputplatelabeli}{\oldlabel{RMAAOutputPlate1}}
\newcommand{\@outputplatelabelii}{\oldlabel{RMAAOutputPlate2}}
\newcommand{\@outputplatelabeliii}{\oldlabel{RMAAOutputPlate3}}
\newcommand{\@outputplatelabeliv}{\oldlabel{RMAAOutputPlate4}}
\newcommand{\@outputplatelabelv}{\oldlabel{RMAAOutputPlate5}}
\newcommand{\@outputplatelabelvi}{\oldlabel{RMAAOutputPlate6}}
\newcommand{\@outputplatelabelvii}{\oldlabel{RMAAOutputPlate7}}
\newcommand{\@outputplatelabelviii}{\oldlabel{RMAAOutputPlate8}}
\newcommand{\@outputplatelabelix}{\oldlabel{RMAAOutputPlate9}}
\newcommand{\@outputplatelabelx}{\oldlabel{RMAAOutputPlate10}}
\newcommand{\@outputplatelabelxi}{\oldlabel{RMAAOutputPlate11}}
\newcommand{\@outputplatelabelxii}{\oldlabel{RMAAOutputPlate12}}
\newcommand{\@outputplatelabelxiii}{\oldlabel{RMAAOutputPlate13}}
\newcommand{\@outputplatelabelxiv}{\oldlabel{RMAAOutputPlate14}}
\newcommand{\@outputplatelabelxv}{\oldlabel{RMAAOutputPlate15}}
\newcommand{\@platerefi}{\oldpageref{RMAAPlate1}}
\newcommand{\@platerefii}{\oldpageref{RMAAPlate2}}
\newcommand{\@platerefiii}{\oldpageref{RMAAPlate3}}
\newcommand{\@platerefiv}{\oldpageref{RMAAPlate4}}
\newcommand{\@platerefv}{\oldpageref{RMAAPlate5}}
\newcommand{\@platerefvi}{\oldpageref{RMAAPlate6}}
\newcommand{\@platerefvii}{\oldpageref{RMAAPlate7}}
\newcommand{\@platerefviii}{\oldpageref{RMAAPlate8}}
\newcommand{\@platerefix}{\oldpageref{RMAAPlate9}}
\newcommand{\@platerefx}{\oldpageref{RMAAPlate10}}
\newcommand{\@platerefxi}{\oldpageref{RMAAPlate11}}
\newcommand{\@platerefxii}{\oldpageref{RMAAPlate12}}
\newcommand{\@platerefxiii}{\oldpageref{RMAAPlate13}}
\newcommand{\@platerefxiv}{\oldpageref{RMAAPlate14}}
\newcommand{\@platerefxv}{\oldpageref{RMAAPlate15}}
\newcounter{figplatei}
\newcounter{figplateii}
\newcounter{figplateiii}
\newcounter{figplateiv}
\newcounter{figplatev}
\newcounter{figplatevi}
\newcounter{figplatevii}
\newcounter{figplateviii}
\newcounter{figplateix}
\newcounter{figplatex}
\newcounter{figplatexi}
\newcounter{figplatexii}
\newcounter{figplatexiii}
\newcounter{figplatexiv}
\newcounter{figplatexv}
%% Printed output should use the counter `plate'
\renewcommand{\theplate}{\@arabic\c@plate}
\renewcommand{\theplateloop}{\@arabic\c@plateloop}
%% Command to output all the saved plates in this paper
\newcommand{\outputplates}{%
  \setcounter{plateloop}{0}
  \if@book\else\@endofarticlehook\fi % not entirely sure why, but this
                                     % seems to be needed
  \clearpage % first output the last page of text
  % now set pagestyle for blank even pages and make sure we start on an odd page
  \pagestyle{empty}
  \cleardoublepage
  \onecolumn
  \parindent\z@
  % first put in an extra sheet
  % the main purpose of this is so there is always a blank even page before 
  % each plate for filling in the PS file if the article has an odd number
  % of pages (see splitbook.sh)
  \null\vfil\centerline{PLATES}\vfil\clearpage\pagestyle{rmaaplate}\cleardoublepage
  % output each plate in turn
  \@whilenum\value{plateloop}<\value{plate}\do{%
    \addtocounter{plateloop}{1}
    % restore the figure counter from \@figcounti, \@figcountii, etc
    \setcounter{figure}{\value{figplate\roman{plateloop}}}
    % spew out the plate (in its own group for safety)
    {\csname @plate\roman{plateloop}\endcsname\par}
    }
  \pagestyle{rmaa}
  % reset the page number as if nothing had happened
  \addtocounter{page}{-\value{plate}}
  \addtocounter{page}{-\value{plate}}
  }
\newcommand\figcaption{%
  \def\@captype{figure}
  \renewcommand{\@makecaption}{\@makefigcap}
  \caption}
\newcommand\tabcaption{%
  \def\@captype{table}
  \renewcommand{\@makecaption}{\@maketabcap}
  \caption}

\newcommand{\plate}[2][1]{%
  %% User command to make a plate page: 
  %% Mandatory argument is what goes on the plate - this should not 
  %% include any floating environments (e.g. figure). Instead, \figcaption 
  %% should be used to do the captions. 
  %% Optional argument is the number of figures (i.e. captions) on the plate
  %% page (default is 1). This is necessary for incrementing the figure counter. 
  %%
  %% Note that we have to use a \gdef rather than an \xdef when we 
  %% save the plate output in a command. This means that in the book styles when it 
  %% comes to spewing out the plate, we will have forgotten about any user macros 
  %% that may have been defined. This can be got round by repeating the macros inside the plate.
  %% (Presumably, one would use a \newcommand rather than \renewcommand if they were new.)
  % increment global and local counters
  \addtocounter{plate}{1}
  % write appropriate label to .aux file
  \csname @platelabel\roman{plate}\endcsname
  % define a global macro called \@platei, \@plateii, \@plateiii, \@plateiv, ...
  \expandafter\gdef\csname @plate\roman{plate}\endcsname{%
    \csname @outputplatelabel\roman{plateloop}\endcsname
    \null\vfil #2 \vfil
    \cleardoublepage
    }
  % save the current shortauthor and shorttitle
  % need to be \xdef to take current values of \@shortauthor, etc 
  \expandafter\xdef\csname @shorta\roman{plate}\endcsname{\@shortauthor}
  \expandafter\xdef\csname @shortt\roman{plate}\endcsname{\@shorttitle}
  % save the current value of the figure counter for later use in \outputplates
  \setcounter{figplate\roman{plate}}{\value{figure}}
  % increment the figure counter according to the optional argument
  \addtocounter{figure}{#1}
  }
%% page style for plate pages 
\newcommand{\ps@rmaaplate}{%
  % Plate X (see page Y) in the foot
  \renewcommand\@oddfoot{\hfil\csname @shorta\roman{plateloop}\endcsname\ 
    (see page~\csname @plateref\roman{plateloop}\endcsname)\hfil}%
  \renewcommand\@evenfoot{\relax}%
  % nothing in even heads (blank pages)
  \renewcommand\@evenhead{\relax}
  % short title and page num in the odd heads
  \renewcommand\@oddhead{\hfil\csname @shortt\roman{plateloop}\endcsname%
    \hfil\llap{PLATE~\theplateloop}}
  }





%% Things for the User Guide
\if@guide
%% stolen from doc.sty
\def\MakeShortVerb#1{%
  \expandafter\ifx\csname cc\string#1\endcsname\relax
  \@shortvrbinfo{Made }{#1}%
  \add@special{#1}%
  \expandafter
  \xdef\csname cc\string#1\endcsname{\the\catcode`#1}%
  \begingroup
  \catcode`\~\active  \lccode`\~`#1%
  \lowercase{%
    \global\expandafter\let
    \csname ac\string#1\endcsname~%
    \gdef~{\verb~}}%
  \endgroup
  \global\catcode`#1\active
  \else
  \@shortvrbinfo\@empty{#1 already}%
  \fi}
\def\DeleteShortVerb#1{%
  \expandafter\ifx\csname cc\string#1\endcsname\relax
  \else
  \@shortvrbinfo{Deleted }{#1 as}%
  \rem@special{#1}%
  \global\catcode`#1\csname cc\string#1\endcsname
  \global \expandafter\let \csname cc\string#1\endcsname \relax
  \ifnum\catcode`#1=\active
  \begingroup
  \catcode`\~\active   \lccode`\~`#1%
  \lowercase{%
    \global\expandafter\let\expandafter~%
    \csname ac\string#1\endcsname}%
  \endgroup \fi \fi}
\def\@shortvrbinfo#1#2{%
  \PackageInfo{doc}{%
    #1\expandafter\@gobble\string#2 a short reference
    for \string\verb}}
\def\add@special#1{%
  \rem@special{#1}%
  \expandafter\gdef\expandafter\dospecials\expandafter
  {\dospecials \do #1}%
  \expandafter\gdef\expandafter\@sanitize\expandafter
  {\@sanitize \@makeother #1}}
\def\rem@special#1{%
  \def\do##1{%
    \ifnum`#1=`##1 \else \noexpand\do\noexpand##1\fi}%
  \xdef\dospecials{\dospecials}%
  \begingroup
  \def\@makeother##1{%
    \ifnum`#1=`##1 \else \noexpand\@makeother\noexpand##1\fi}%
  \xdef\@sanitize{\@sanitize}%
  \endgroup}
%% define | as shortcut for \verbatim
\MakeShortVerb{\|}
%% comand to specify arguments to LaTeX commands
\newcommand{\rmArg}[1]{%
  \ensuremath{\langle}{\scriptsize \textsf{\MakeUpperCase{#1}}}\ensuremath{\rangle}}
%\newcommand{\rmArg}[1]{\ensuremath{\langle}\textsc{#1}\ensuremath{\rangle}}
\providecommand{\BibTeX}{{\rmfamily B\kern-.05em%
    \textsc{i\kern-.025em b}\kern-.08em%
    T\kern-.1667em\lower.7ex\hbox{E}\kern-.125emX}}
%% Environment for code examples
\newenvironment{Code}{%
  \begin{list}{}{\setlength{\leftmargin}{0.05\columnwidth}} 
  \item[]}{\end{list}}
\fi
%</package>
%    \end{macrocode}

% \section{User commands for astronomical articles}
% \label{sec:user-commands}
% Some random definitions. These go both into the main class file and
% also into |adshtml.sty|.
% 
% Definitions of the journal name in different incarnations.  
%    \begin{macrocode}
%<*package|adshtml-style>
\newcommand{\RMAA}{RMxAA}
\newcommand{\RMAAlong}{Revista Mexicana de Astronom\'{\i}a y Astrof\'{\i}sica}
\newcommand{\RMAASC}{RMxAC}
\newcommand{\RMAASClong}{\RMAA\ (Serie de Conferencias)}
%    \end{macrocode}
% \begin{macro}{\ibidrule}
% This puts a horizontal rule with a period after it. It is used when
% the author list of a reference is identical to the preceding one. 
%    \begin{macrocode}
\newcommand{\ibidrule}{\rule[-0.03cm]{4.7em}{.01cm}.\ }
%    \end{macrocode}
% \end{macro}
% \begin{macro}{\ThisVolume}
% References to other articles in the same volume can be coded as,
% e.g., |\bibitem{} Bloggs, J. \ThisVolume{bloggs}|.
%    \begin{macrocode}
\newcommand{\ThisVolume}[1]{%
  \@year, \if@proceedings\RMAASC\else\RMAA\fi, \@volume,
  \if@book\oldpageref{#1:RMAAFirstPage}\else 000\fi\ (this volume)}
%    \end{macrocode}
% \end{macro}
% Other stuff that may be useful in ref lists.
%    \begin{macrocode}
\newcommand{\etal}{et al.\@}
%    \end{macrocode}
% Various new commands from AAS macros
%    \begin{macrocode}
\def\acknowledgments{\vskip 3ex plus .8ex minus .4ex}
\let\acknowledgements=\acknowledgments                  % second spelling
\let\jnl@style=\rm
\def\ref@jnl#1{{\jnl@style#1}}
\DeclareRobustCommand\aj{\ref@jnl{AJ}}                   
\DeclareRobustCommand\araa{\ref@jnl{ARA\&A}}             
\DeclareRobustCommand\apj{\ref@jnl{ApJ}}                 
\DeclareRobustCommand\apjl{\ref@jnl{ApJ}}                
\DeclareRobustCommand\apjs{\ref@jnl{ApJS}}               
\DeclareRobustCommand\ao{\ref@jnl{Appl.~Opt.}}           
\DeclareRobustCommand\apss{\ref@jnl{Ap\&SS}}             
\DeclareRobustCommand\aap{\ref@jnl{A\&A}}                
\DeclareRobustCommand\aapr{\ref@jnl{A\&A~Rev.}}          
\DeclareRobustCommand\aaps{\ref@jnl{A\&AS}}              
\DeclareRobustCommand\azh{\ref@jnl{AZh}}                 
\DeclareRobustCommand\baas{\ref@jnl{BAAS}}               
\DeclareRobustCommand\jrasc{\ref@jnl{JRASC}}             
\DeclareRobustCommand\memras{\ref@jnl{MmRAS}}            
\DeclareRobustCommand\mnras{\ref@jnl{MNRAS}}             
\DeclareRobustCommand\pra{\ref@jnl{Phys.~Rev.~A}}        
\DeclareRobustCommand\prb{\ref@jnl{Phys.~Rev.~B}}        
\DeclareRobustCommand\prc{\ref@jnl{Phys.~Rev.~C}}        
\DeclareRobustCommand\prd{\ref@jnl{Phys.~Rev.~D}}        
\DeclareRobustCommand\pre{\ref@jnl{Phys.~Rev.~E}}        
\DeclareRobustCommand\prl{\ref@jnl{Phys.~Rev.~Lett.}}    
\DeclareRobustCommand\pasp{\ref@jnl{PASP}}               
\DeclareRobustCommand\pasj{\ref@jnl{PASJ}}               
\DeclareRobustCommand\qjras{\ref@jnl{QJRAS}}             
\DeclareRobustCommand\skytel{\ref@jnl{S\&T}}             
\DeclareRobustCommand\solphys{\ref@jnl{Sol.~Phys.}}      
\DeclareRobustCommand\sovast{\ref@jnl{Soviet~Ast.}}      
\DeclareRobustCommand\ssr{\ref@jnl{Space~Sci.~Rev.}}     
\DeclareRobustCommand\zap{\ref@jnl{ZAp}}                 
\DeclareRobustCommand\nat{\ref@jnl{Nature}}              
\DeclareRobustCommand\iaucirc{\ref@jnl{IAU~Circ.}}
\DeclareRobustCommand\aplett{\ref@jnl{Astrophys.~Lett.}}
\DeclareRobustCommand\apspr{\ref@jnl{Astrophys.~Space~Phys.~Res.}}
\DeclareRobustCommand\bain{\ref@jnl{Bull.~Astron.~Inst.~Netherlands}}
\DeclareRobustCommand\fcp{\ref@jnl{Fund.~Cosmic~Phys.}}
\DeclareRobustCommand\gca{\ref@jnl{Geochim.~Cosmochim.~Acta}}
\DeclareRobustCommand\grl{\ref@jnl{Geophys.~Res.~Lett.}}
\DeclareRobustCommand\jcp{\ref@jnl{J.~Chem.~Phys.}}      
\DeclareRobustCommand\jgr{\ref@jnl{J.~Geophys.~Res.}}    
\DeclareRobustCommand\jqsrt{\ref@jnl{J.~Quant.~Spec.~Radiat.~Transf.}}
\DeclareRobustCommand\memsai{\ref@jnl{Mem.~Soc.~Astron.~Italiana}}
\DeclareRobustCommand\nphysa{\ref@jnl{Nucl.~Phys.~A}}
\DeclareRobustCommand\physrep{\ref@jnl{Phys.~Rep.}}
\DeclareRobustCommand\physscr{\ref@jnl{Phys.~Scr}}
\DeclareRobustCommand\planss{\ref@jnl{Planet.~Space~Sci.}}
\DeclareRobustCommand\procspie{\ref@jnl{Proc.~SPIE}}
\let\astap=\aap
\let\apjlett=\apjl
\let\apjsupp=\apjs
\let\applopt=\ao
\def\phn{\phantom{0}}
\def\phd{\phantom{.}}
\def\phs{\phantom{\ensuremath{-}}}
\def\phm#1{\phantom{#1}}
%    \end{macrocode}
% All these are redone as robust and using more modern commands where
% possible. The robustness is so they will not get expanded if used in
% abstracts and titles when these are passed to the ADS file. 
% \changes{3.16}{2003/04/07}{
%  Added \cs{Color} command for doing photometric colors.}
%    \begin{macrocode}
%%\def\sun{\hbox{$\odot$}}
\DeclareRobustCommand\sun{\ensuremath{\odot}}
%%\def\earth{\hbox{$\oplus$}}
\DeclareRobustCommand\earth{\ensuremath{\oplus}}
%%\def\lesssim{\mathrel{\hbox{\rlap{\hbox{\lower4pt\hbox{$\sim$}}}\hbox{$<$}}}}
%%\def\gtrsim{\mathrel{\hbox{\rlap{\hbox{\lower4pt\hbox{$\sim$}}}\hbox{$>$}}}}
\DeclareRobustCommand\lesssim{\mathrel{\hbox{\smaller\smaller\rlap{\hbox{\lower3pt\hbox{$\sim$}}}\raise1.5pt\hbox{$<$}}}}
\DeclareRobustCommand\gtrsim{\mathrel{\hbox{\smaller\smaller\rlap{\hbox{\lower3pt\hbox{$\sim$}}}\raise1.5pt\hbox{$>$}}}}
%%\def\lta{\mathrel{\spose{\lower 3pt\hbox{$\mathchar"218$}}
%%     \raise 2.0pt\hbox{$\mathchar"13C$}}}
%%\def\gta{\mathrel{\spose{\lower 3pt\hbox{$\mathchar"218$}}
%%     \raise 2.0pt\hbox{$\mathchar"13E$}}}
\DeclareRobustCommand\sq{\hbox{\rlap{$\sqcap$}$\sqcup$}}
%%\def\arcdeg{\hbox{$^\circ$}}
%%\def\arcmin{\hbox{$^\prime$}}
%%\def\arcsec{\hbox{$^{\prime\prime}$}}
\DeclareRobustCommand\arcdeg{\ensuremath{^\circ}}
\DeclareRobustCommand\arcmin{\ensuremath{^\prime}}
\DeclareRobustCommand\arcsec{\ensuremath{^{\prime\prime}}}
%%\def\fd{\hbox{$.\!\!^{\rm d}$}}
%%\def\fh{\hbox{$.\!\!^{\rm h}$}}
%%\def\fm{\hbox{$.\!\!^{\rm m}$}}
%%\def\fs{\hbox{$.\!\!^{\rm s}$}}
%%\def\fdg{\hbox{$.\!\!^\circ$}}
%%\def\farcm{\hbox{$.\mkern-4mu^\prime$}}
%%\def\farcs{\hbox{$.\!\!^{\prime\prime}$}}
%%\def\fp{\hbox{$.\!\!^{\scriptscriptstyle\rm p}$}}
\DeclareRobustCommand\fd{\ensuremath{.\!\!^\mathrm{d}}}
\DeclareRobustCommand\fh{\ensuremath{.\!\!^\mathrm{h}}}
\DeclareRobustCommand\fm{\ensuremath{.\!\!^\mathrm{m}}}
\DeclareRobustCommand\fs{\ensuremath{\mkern0.5mu.\mkern-5mu^\mathrm{s}\mkern-1mu}}
\DeclareRobustCommand\fdg{\ensuremath{.\!\!^\circ}}
\DeclareRobustCommand\farcm{\ensuremath{.\mkern-4mu^\prime}}
\DeclareRobustCommand\farcs{\ensuremath{\mkern1mu.\mkern-6mu^{\prime\prime}\mkern-1.5mu}}
%%\DeclareRobustCommand\farcs{\ensuremath{.\!\!^{\prime\prime}}}
\DeclareRobustCommand\fp{\ensuremath{.\!\!^{\scriptscriptstyle\rm p}}}
%%\def\micron{\hbox{$\mu$m}}
\DeclareRobustCommand\micron{\ensuremath{\mu\mathrm{m}}}
\let\la=\lesssim                        
\let\ga=\gtrsim
%%\def\case#1#2{\hbox{$\frac{#1}{#2}$}}
%%\def\slantfrac#1#2{\hbox{$\,^#1\!/_#2$}}
%%\def\onehalf{\slantfrac{1}{2}}
%%\def\onethird{\slantfrac{1}{3}}
%%\def\twothirds{\slantfrac{2}{3}}
%%\def\onequarter{\slantfrac{1}{4}}
%%\def\threequarters{\slantfrac{3}{4}}
\DeclareRobustCommand\case[2]{\ensuremath{\frac{#1}{#2}}}
\DeclareRobustCommand\slantfrac[2]{\ensuremath{\,^#1\!/_#2}}
\DeclareRobustCommand\onehalf{\slantfrac{1}{2}}
\DeclareRobustCommand\onethird{\slantfrac{1}{3}}
\DeclareRobustCommand\twothirds{\slantfrac{2}{3}}
\DeclareRobustCommand\onequarter{\slantfrac{1}{4}}
\DeclareRobustCommand\threequarters{\slantfrac{3}{4}}
%%\def\ubvr{\hbox{$U\!BV\!R$}}            
%%\def\ub{\hbox{$U\!-\!B$}}               
%%\def\bv{\hbox{$B\!-\!V$}}               
%%\def\vr{\hbox{$V\!-\!R$}}               
%%\def\ur{\hbox{$U\!-\!R$}}
%    \end{macrocode}
% New commands for ``photometric'' colors 
%    \begin{macrocode}
\DeclareRobustCommand\ubvr{\ensuremath{U\!BV\!R}}            
\DeclareRobustCommand\Color[2]{%
  \ensuremath{%
    #1\!\ifx#1V\!\fi-\!#2
  }}
\DeclareRobustCommand\ub{\Color{U}{B}}
\DeclareRobustCommand\bv{\Color{B}{V}}
\DeclareRobustCommand\vr{\Color{V}{R}}
\DeclareRobustCommand\ur{\Color{U}{R}}
%% New version of the \ion command, that actually works! 
%    \end{macrocode}
% \begin{macro}{\ion}
% This version of the |\ion| command allows you to write either
% |\ion{N}{4}| or |\ion{N}{iv}| in order to get N~IV. 
%    \begin{macrocode}
\DeclareRobustCommand\ion[2]{%
  \mbox{#1\kern0.2em%
    \smaller\rmfamily%
    \edef\@tempa{\@car#2\@nil}%
    \ifcat1\@tempa%
    \@Roman{#2}%
    \else%
    \uppercase{#2}%
    \fi}}
\newcounter{IonStage}
\DeclareRobustCommand\plainion[2]{\setcounter{IonStage}{#2}#1 \Roman{IonStage}}
%%\def\ion#1#2{\ensuremath{\mathrm{#1}\;}{\protect\small\rm\@Roman{#2}}}
%    \end{macrocode}
% \end{macro}
%    \begin{macrocode}
\DeclareRobustCommand\nodata{$\cdots$}
%    \end{macrocode}
% These are needed by the |\ion| command. Isn't there an easier way of
% doing this?
%    \begin{macrocode}
\let\@small=\small
\let\@normalsize=\normalsize
\let\@footnotesize=\footnotesize
\let\@scriptsize=\scriptsize
\let\@tiny=\tiny
\let\@large=\large
\let\@Large=\Large
\let\@LARGE=\LARGE
\let\@huge=\huge
\let\@Huge=\Huge
\newcommand{\smaller}{\small}
\renewcommand{\normalsize}{\renewcommand{\smaller}{\small}\@normalsize}
\renewcommand{\small}{\renewcommand{\smaller}{\footnotesize}\@small}
\renewcommand{\footnotesize}{\renewcommand{\smaller}{\scriptsize}\@footnotesize}
\renewcommand{\scriptsize}{\renewcommand{\smaller}{\tiny}\@scriptsize}
\renewcommand{\tiny}{\renewcommand{\smaller}{\tiny}\@tiny}
\renewcommand{\large}{\renewcommand{\smaller}{\normalsize}\@large}  
\renewcommand{\Large}{\renewcommand{\smaller}{\large}\@Large}  
\renewcommand{\LARGE}{\renewcommand{\smaller}{\Large}\@LARGE}  
\renewcommand{\huge}{\renewcommand{\smaller}{\LARGE}\@huge}  
\renewcommand{\Huge}{\renewcommand{\smaller}{\huge}\@Huge}  
%</package|adshtml-style>
%    \end{macrocode}


% \section{Other files}

% \subsection{The driver and make files}
% \label{sec:driver_make}
% These are put here for completeness, although they still need to be
% distributed separately since otherwise the user won't be able to
% unpack this file into its constituent parts.
% 
% The |Makefile| has two principal rules: |docs|, and |macros|. The
% |docs| rule produces a PDF and gzipped PS file of the documentation.
% I still haven't perfected the combination of rules required to get
% it to run \LaTeX{} enough times in all circumstances. The |macros|
% rule runs \LaTeX{} on |rmaa.ins| in order to unpack |rmaa.cls| and
% all the other little files (including the |Makefile| and the
% |rmaa.ins| files themselves\dots 
% 
% First set up lists of various types of file. 
%    \begin{macrocode}
%<*makefile>
DOCFILES=rmaa.pdf
MACRO_LTX_FILES=rmaa.cls rmaacite.sty adshtml.sty
SHELLSCRIPTS=makeads.sh
MACRO_OTHER_FILES=adshtml.perl sample.latex2html-init
MAKEFILE=Makefile
MACROFILES=$(MACRO_LTX_FILES) $(SHELLSCRIPTS) $(MACRO_OTHER_FILES)         
SKELETONFILES=ADS.tex 
%    \end{macrocode}
% Install the files read only into the package directory, specified via a path relative to the source directory. \changes{4.0}{2019/11/14}{Reorganize folder layout so that hardly anything is in the top level.  That way, things are clearer on github.}
%    \begin{macrocode}
INSTALL=install 
INSTALL_DIR=../packages
%    \end{macrocode}
% The default rule is to make everything. Installation must be
% explicitly asked for.
%    \begin{macrocode}
all: docs macros

install: install-docs install-macros

%    \end{macrocode}
% Most files are installed read-only to prevent editing. Skeleton
% files must be writeable. Must be careful not to overwrite the
% Makefile in the parent directory. 
%    \begin{macrocode}
install-macros: macros
        $(INSTALL) -m 444 $(MACRO_LTX_FILES) $(MACRO_OTHER_FILES) $(INSTALL_DIR)
        $(INSTALL) -m 555 $(SHELLSCRIPTS) $(INSTALL_DIR)
        $(INSTALL) -m 644 $(SKELETONFILES) $(INSTALL_DIR)

install-docs: docs
        $(INSTALL) rmaa.pdf $(INSTALL_DIR)/rm-fulldocs.pdf

%    \end{macrocode}
% Make the PDF and gzipped PS files. 
%    \begin{macrocode}
docs: $(DOCFILES)

%    \end{macrocode}
% Make everything else if either the |.dtx| or |.ins| files have
% changed. 
%    \begin{macrocode}
macros: rmaa.dtx rmaa.ins
%    \end{macrocode}
% Ensure that all the macro files can be written to then recreate
% them. The |touch| is there in case any of the files don't exist
% yet. 
%    \begin{macrocode}
        touch $(MACROFILES) $(MAKEFILE)
        chmod u+w $(MACROFILES) $(MAKEFILE) rmaa.ins
        latex "%%
%% This is file `rmaa.ins',
%% generated with the docstrip utility.
%%
%% The original source files were:
%%
%% rmaa.dtx  (with options: `ins')
%% 
%% PLEASE DO NOT EDIT THIS FILE! It is generated automatically from
%% the file rmaa.dtx using the accompanying Makefile and rmaa.ins
%% Any hand-edited changes will be lost
%% 
%% A LaTeX document class for the RMxAA
%% 
%% Author: William Henney (w.henney@crya.unam.mx)
%% 
%% Latest version available from: http://www.crya.unam.mx/rmaa
%% 
%% Copyright (C) 1999-2010 Instituto de Astronomia, UNAM, Mexico
%% 
%% This program is free software; you can redistribute it and/or
%% modify it under the terms of the GNU General Public License
%% as published by the Free Software Foundation; either version 2
%% of the License, or (at your option) any later version.
%% 
%% This program is distributed in the hope that it will be useful,
%% but WITHOUT ANY WARRANTY; without even the implied warranty of
%% MERCHANTABILITY or FITNESS FOR A PARTICULAR PURPOSE.  See the
%% GNU General Public License for more details.
%% 
%% You should have received a copy of the GNU General Public License
%% along with this program; if not, write to the Free Software
%% Foundation, Inc., 59 Temple Place - Suite 330, Boston, MA  02111-1307, USA.





%%
%% Installation batch file for the RMAA LaTeX package
%%
\def\batchfile{rmaa.ins}
\input docstrip.tex
\preamble

PLEASE DO NOT EDIT THIS FILE! It is generated automatically from
the file rmaa.dtx using the accompanying Makefile and rmaa.ins
Any hand-edited changes will be lost

A LaTeX document class for the RMxAA

Author: William Henney (w.henney@crya.unam.mx)

Latest version available from: http://www.crya.unam.mx/rmaa

Copyright (C) 1999-2010 Instituto de Astronomia, UNAM, Mexico

This program is free software; you can redistribute it and/or
modify it under the terms of the GNU General Public License
as published by the Free Software Foundation; either version 2
of the License, or (at your option) any later version.

This program is distributed in the hope that it will be useful,
but WITHOUT ANY WARRANTY; without even the implied warranty of
MERCHANTABILITY or FITNESS FOR A PARTICULAR PURPOSE.  See the
GNU General Public License for more details.

You should have received a copy of the GNU General Public License
along with this program; if not, write to the Free Software
Foundation, Inc., 59 Temple Place - Suite 330, Boston, MA  02111-1307, USA.
\endpreamble
\askforoverwritefalse
\generate{%
  \file{rmaa.cls}{\from{rmaa.dtx}{package}}
  \file{rmaa.ins}{\from{rmaa.dtx}{ins}}
  \file{rmaacite.sty}{\from{rmaa.dtx}{rmaacite}}
  \file{adshtml.sty}{\from{rmaa.dtx}{adshtml-style}}
  \file{ADS.tex}{\from{rmaa.dtx}{ads-container}}
  }
\generate{\nopreamble\nopostamble%
  \file{Makefile}{\from{rmaa.dtx}{makefile}}
  \file{adshtml.perl}{\from{rmaa.dtx}{adshtml-perl}}
  \file{makeads.sh}{\from{rmaa.dtx}{ads-shell-script}}
  \file{sample.latex2html-init}{\from{rmaa.dtx}{latex2html-init}}
  }

%% Local Variables:
%% mode: LaTeX
%% mode: Font-lock
%% End:









\endinput
%%
%% End of file `rmaa.ins'.
"
%    \end{macrocode}
% \TeX{} automatically adds the |.tex| extension to extensionless
% files. We fix this and take the opportunity to convert
% back to TABs all that it has lovingly converted to spaces. 
%    \begin{macrocode}
        cp -f $(MAKEFILE) $(MAKEFILE).BCK && unexpand $(MAKEFILE).tex >| $(MAKEFILE) && rm $(MAKEFILE).tex
%    \end{macrocode}
% Make the shell scripts executable and everything unwriteable (so
% they don't get hand-edited by mistake). However, undo this for the
% files that the user is \emph{supposed} to modify. 
%    \begin{macrocode}
        chmod a+x $(SHELLSCRIPTS)
        chmod a-w $(MACROFILES) $(MAKEFILE)
        chmod u+w $(SKELETONFILES)

%    \end{macrocode}
% We now use PDF\LaTeX{} for this.
%    \begin{macrocode}
rmaa.idx: rmaa.dtx
        @echo "REGENERATING .idx FILE"
        pdflatex "\nonstopmode% \def\fileversion{4.0} \def\filedate{2019/11/12}
% \MakeShortVerb{\|}
% \title{A \LaTeX{} document class for the Revista Mexicana de
% Astronom\'\i a y Astrof\'\i sica\thanks{This file has version
% \fileversion{} last revised \filedate{}.}}
% \author{William J. Henney\thanks{w.henney@irya.unam.mx}}
% \date{\filedate}
% \maketitle
% \paragraph{Abstract}
% This document describes the \texttt{rmaa} \LaTeX{}
% document class for the journal ``Revista Mexicana de
% Astronom\'\i a y Astrof\'\i sica''. The macros can be used for 
% \begin{enumerate}
% \item Preparing author manuscripts for the journal (including
% conference proceedings to be published in the ``Serie de
% Conferencias'')
% \item Typesetting an entire volume of the journal by inputting the
% individual article files from a container file. Facilities are
% included for automatic generation of TOC and Author Index. 
% \end{enumerate}
% 
% \tableofcontents
% \section{The documentation driver file}
% This is the part that generates the documentation that you may or
% may not be reading, which is done with the help of the Makefile
% discussed in section~\ref{sec:driver_make}. 
%    \begin{macrocode}
%<*driver>
\documentclass{ltxdoc}
\usepackage{hyperref}
\usepackage{palatino}
\EnableCrossrefs \CodelineIndex \RecordChanges
\setlength{\parindent}{0pt}
\def\MacroFont{\footnotesize\tt}
\begin{document} \DocInput{rmaa.dtx} \PrintIndex \PrintChanges \end{document}
%</driver>
%    \end{macrocode}

% \section{The implementation}
% 
% \changes{3.24}{2007/09/14}{Added user access to the macro version
% and date via \cs{rmVersion} and \cs{rmDate}}
% 
%    \begin{macrocode}
%<*package>
\newcommand{\rmaa@version}{3.26}
\newcommand{\rmaa@date}{2011/08/15}
\newcommand\rmVersion{\rmaa@version}
\newcommand\rmDate{\rmaa@date}
\NeedsTeXFormat{LaTeX2e}[1995/06/01]
%    \end{macrocode}
% \changes{4.0}{2019/11/02}{Change abstract to appear before
% resumen. Change RevMexAA to RMxAA for consistency with ADS.}
%
% \changes{3.22}{2007/03/22}{Allow proceedings articles to have
%   received/accepted dates} \changes{3.21}{2006/05/15}{Changes for
%   Silvia - abstracts are now shipped out to ADS.}
% \changes{3.19}{2006/01/19}{Changes for Silvia - remove adresses from
%   end, new abstract style} \changes{2.2beta}{2001/10/30}{ Fixed bug
%   whereby twocolumn option was not passed through to book
%   class. Finally sorted out the ``label on last page'' issue. }
% \changes{2.1}{2001/01/23}{ Fixed some inconsistencies in
%   \cs{nocite} and \cs{bibitem} uniqification Fixed the definition of
%   \cs{ion} so it will work in math mode and in contexts where the
%   typesize is different from normal (e.g., title, fig captions).
%   Modified the definition of \cs{@caption}, so we don't write any
%   list-of-figures data (we never use it anyway, and it just causes
%   problems with fragile commands in fig captions). }
% \changes{2.0}{2000/11/30}{ Added `newstyle' option, which should be
%   used in all new proceedings volumes (wider abstract, smaller
%   typeface in fig captions, etc.) Added commands for formatting
%   table captions. New demonstation document rmpostsamp.tex for
%   1-page posters. This will be the last release under a PRI
%   government (for a while...)}  \changes{1.9}{2000/06/06}{
%   Corrected position of page number in footer of articles that open
%   on even pages. Ensured that first article in each session opens on
%   an odd page, as do list of abstracts and author index. Added extra
%   space in TOC before the backmatter. Added command for editors,
%   \cs{BoldSessionTitles}, that makes produces bold titles for
%   session in the TOC.}  \changes{1.8}{2000/05/16}{ Fixed bug in
%   last page label with \cs{suppressfulladdresses} }
% \changes{1.7}{2000/05/11}{ Fixed more little things. Page labels
%   now work, even if you have no bibliography. You can output the
%   full addresses by hand with the \cs{outputfulladdresses}
%   command. Editors' Guide updated but still not finished. Revamped
%   splitbook.sh now produces PS file of each article for web page and
%   also an EPS file of each page for sending to the printers.}
% \changes{1.6}{2000/04/15}{ More cleaning up: full stops after
%   titles in index, running headers and footers in frontmatter,
%   fontsize in title of index, etc} \changes{1.5}{2000/04/15}{
%   Eliminated superscripts from titles in TOC and index}
% \changes{1.4}{2000/04/14}{ Added copyright and license info (GPL)}
% \changes{1.3}{2000/03/01}{ Corrected \cs{figcaption} to do the
%   caption formatting right Corrected font size of sec headings in
%   front matter} \changes{1.2}{2000/02/16}{ Added \cs{plate}
%   command} \changes{1.1}{2000/02/14}{ Added manuscript option,
%   corrected bugs} \changes{1.0}{2000/02/08}{ Now does a whole book
%   - lots of stuff tidied up} \changes{0.7}{1999/10/04}{ Increased
%   spacing between sections in single column styles. Suppressed
%   floats at bottom of pages (didn't work though) }
% \changes{0.6}{1999/09/24}{ Tweaked \cs{adjustfinalcols} (still not
%   perfect) Added proceedings option Added onecolumn/twocolumn
%   options (orthogonal to journal/proceedings) }
% \changes{0.5}{1999/09/23}{ Added \cs{fulladdresses} and
%   \cs{adjustfinalcols} (just about working)}
% \changes{0.4}{1999/09/22}{ Minor stylistic changes}
% \changes{0.3}{1999/06/28}{ Better mechanism for alternative
%   affiliations Added `guide' option for user guide}
% \changes{0.2}{1999/05/27}{ All basic functionality present, I
%   think} \changes{0.1}{1999/05/26}{ Just about working - still
%   missing Fig/Tab/Refs} \changes{0.0}{1999/05/21}{ So bored today,
%   decided to rewrite RMAA LaTeX macros from scratch}
%    \begin{macrocode}
\ProvidesClass{rmaa}[\rmaa@date\space v\rmaa@version\space RMxAA article
class]
%    \end{macrocode}
% \subsection{Option Handling}
% \label{sec:options}
% Declarations that are needed in options
% \changes{3.05}{2003/10/22}{Added natbib option}
% \changes{3.06}{2003/01/03}{Natbib option turned on permanently}
% \changes{3.07}{2003/01/11}{Added proof option}
% \changes{3.14}{2003/01/11}{Option to turn off natbib}
% \changes{3.18}{2003/10/22}{trying out doctex mode} 
% \changes{3.23}{2007/09/09}{Manuscript is now on by default.}
% \changes{3.26}{2011/08/15}{New option for review articles}
%    \begin{macrocode}
\newif\if@reviewarticle
\@reviewarticlefalse
\newif\if@proof
\@prooffalse
\newif\if@journal
\@journaltrue
\newif\if@printdates
\@printdatesfalse
\newif\if@proceedings
\@proceedingsfalse
\newif\if@manuscript
\@manuscripttrue
\newif\if@bigfont
\@bigfonttrue
\newif\if@guide
\@guidefalse
\newif\if@preprint
\@preprintfalse
\newif\if@rmtwocol
\@rmtwocoltrue
\newif\if@biblist
\@biblistfalse
\newif\if@book
\@bookfalse
\newif\if@natbib
\@natbibtrue
\newif\if@rmcite
\@rmcitefalse
\newif\if@titlevisible
\@titlevisibletrue
\newif\if@abstractpresent
\@abstractpresentfalse
\newif\if@printbanner
\@printbannerfalse
\newif\if@suppressfulladdresses
\@suppressfulladdressesfalse
\newif\if@newstyle
\@newstylefalse
\newif\if@debug
\@debugfalse
\newif\if@boldkeywords
\@boldkeywordsfalse
%    \end{macrocode}
% The option handling---mainly setting various booleans.
% \changes{2.2beta11}{2001/11/26}{Made twocol and
% newstyle the default for proceedings.}
% \changes{3.23}{2007/09/09}{Manuscript is now 10 pt instead of 12pt}
% \changes{3.23}{2007/09/10}{Full addresses are now automatically
% suppressed for proceedings style.}
% \changes{3.24}{2007/09/14}{New option \texttt{boldkeywords} that
% does the keywords in the old ugly way. Now off by default.}
%    \begin{macrocode}
\DeclareOption{manuscript}{\typeout{Manuscript for referee} 
  \@manuscripttrue\@journalfalse\@rmtwocolfalse
  \PassOptionsToClass{10pt}{article}}
\DeclareOption{bigfont}{\typeout{Big font manuscript requested by Gary
  Ferland} 
  \@manuscripttrue\@journalfalse\@rmtwocolfalse\@bigfonttrue
  \PassOptionsToClass{11pt}{article}}
\DeclareOption{journal}{\typeout{Main journal option}%
  \@journaltrue\@rmtwocoltrue\@printdatestrue\@manuscriptfalse
  \PassOptionsToClass{10pt}{article}}
\DeclareOption{guide}{\typeout{User guide option}%
  \@guidetrue\@journalfalse\@rmtwocoltrue\@manuscriptfalse
  \PassOptionsToClass{10pt}{article}}
\DeclareOption{preprint}{\typeout{Preprint option}%
  \@preprinttrue\@manuscriptfalse\PassOptionsToClass{10pt}{article}}
\DeclareOption{proof}{\typeout{Author proof option}%
  \@prooftrue\@manuscriptfalse}
\DeclareOption{proceedings}{\typeout{Proceedings option}%
  \@proceedingstrue\@journalfalse\@rmtwocoltrue\@newstyletrue
  \@suppressfulladdressestrue
  \@printdatesfalse\@manuscriptfalse\PassOptionsToClass{10pt}{article}} 
\DeclareOption{printdates}{\typeout{Force printing of received/accepted dates}%
  \@printdatestrue} 
\DeclareOption{book}{\typeout{Whole issue option}%
  \@booktrue} 
\DeclareOption{nonatbib}{\typeout{Turn off natbib support (required
    for compatibility with old rmaacite package).}%
  \@natbibfalse} 
\DeclareOption{newstyle}{\typeout{Use new layout style}%
  \@newstyletrue}
\DeclareOption{oldstyle}{\typeout{Use old layout style}%
  \@newstylefalse}
\DeclareOption{debug}{\typeout{Loads of annoying messages}%
  \@debugtrue}
%    \end{macrocode}
% These two need to be declared explicitly so that |\ExecuteOptions|
% actually does something.
%    \begin{macrocode}
\DeclareOption{twocolumn}{\@rmtwocoltrue}
\DeclareOption{onecolumn}{\@rmtwocolfalse}
\DeclareOption{twoside}{\PassOptionsToClass{\CurrentOption}{article}}
\DeclareOption{boldkeywords}{\@boldkeywordstrue}
%    \end{macrocode}
% Bold, upper-case keywords used to be the default. Now they must be
% turned on explicitly with the \texttt{boldkeywords} option. 
%    \begin{macrocode}
\DeclareOption{boldkeywords}{\@boldkeywordstrue}
%    \end{macrocode}
% Pass any other option along to the \texttt{article} class.
%    \begin{macrocode}
\DeclareOption*{\PassOptionsToClass{\CurrentOption}{article}}
%    \end{macrocode}
% Default options are to be two-sided and based on the standard
% \LaTeX{} \texttt{article} class,
% unless we are processing the entire book, in which case we base
% ourselves on the standard \texttt{book} class. Default is to use
% manuscript option unless otherwise specified. 
%    \begin{macrocode}
\ExecuteOptions{twoside,manuscript} \ProcessOptions\relax
\if@book
\if@rmtwocol\LoadClass[twocolumn]{book}\else\LoadClass{book}\fi
\newcommand{\abstract}{\relax}
\newcommand{\@maketitle}{\relax}
\else
\if@rmtwocol\LoadClass[twocolumn]{article}\else\LoadClass{article}\fi
\fi
%    \end{macrocode}
% Attempt to use natbib for citations 02 Jan 2003. 
% \changes{3.05}{2003/01/02}{If the natbib option is used, then set
% things up so that natbib's version of the bibliography is identical
% to the standard one.}
% \changes{3.09}{2003/01/23}{Explicitly require a recent version of
% natbib, so that we get a decent error message if the package is too
% old.} 
% \changes{3.23}{2007/09/09}{Now allow Bibtex to work properly} 
% \changes{3.25}{2009/10/01}{Change \cs{newcommand} tp \cs{def} when setting \cs{bibfont}. This allows it to work with all versions of natbib.} 
% \changes{3.25}{2010/02/06}{Updated copyright year.} 
% \changes{3.25}{2010/02/06}{Updated email and web addresses to \texttt{crya.unam.mx}.} 
%    \begin{macrocode}
\if@natbib
\RequirePackage{natbib}[1999/02/25]
\bibpunct{(}{)}{;}{a}{}{,}
\def\bibfont{\small}
\renewcommand\bibsection{\section*{REFERENCES}}
\setlength\bibsep{0pt}
\setlength\bibhang{1.5em}
\bibliographystyle{rmaa}
\fi
%    \end{macrocode}
% We want the \texttt{graphicx} package for the extended syntax of
% |\includegraphics|. We want the \texttt{ifthen} package for using in
% these macros (although we don't much yet). We explicitly load
% \texttt{babel} since some \LaTeX{} distributions messed up the
% hyphenation if we didn't.  
%    \begin{macrocode}
\RequirePackage{graphicx}
\RequirePackage{ifthen}
\IfFileExists{babel.sty}{%
%   \let\@restoreindent=\@afterindentfalse
  \RequirePackage[english]{babel}
%   \@restoreindent
  %% reinstate the definition of \@roman so as not to break .idx files
  \@ifundefined{org@@roman}{\relax}{\renewcommand{\@roman}[1]{\org@@roman{##1}}}
  }{%
  \typeout{Package babel.sty not found - don't worry about this}}
%    \end{macrocode}
% Try to load the \texttt{dcolumn} and \texttt{booktabs} packages if
% they can be found. These can produce nicer-looking tables. 
%    \begin{macrocode}
\IfFileExists{dcolumn.sty}{%
  \RequirePackage{dcolumn}
  }{%
  \typeout{Package dcolumn.sty not found - you will not be able to
    align columns on decimal points in tables}}
\IfFileExists{booktabs.sty}{%
  \RequirePackage{booktabs}
  }{%
  \typeout{Package booktabs.sty not found - your tables will not look
    as pretty as they ought too.}
  }
\IfFileExists{makeidx.sty}{%
  \RequirePackage{makeidx}
}{%
  \typeout{Package makeidx.sty not found - Author Index may not work.}
  }
%% commented out for moment since requires more modern version of LaTeX
%% \RequirePackage[latin1]{inputenc}
%    \end{macrocode}
% \changes{4.0}{2019/11/13}{Make \texttt{hyperref} and
% \texttt{etoolbox} into hard requirements, since we now use them
% extensively.}  We use hyperref to get internal and external links.
% But I don't like the colored boxes so I turn those off.  We use
% etoolbox for patching and robustifying at the moment.  In the
% future, I want to port a lot of other stuff over to it.  For
% instance, replace the ifthen package and the TeX ifs too.
%
% We target the latest version (1.7) of PDF, rather than the default
% of 1.5.  This reduces the chance of getting warnings from included
% graphics.  This is not strictly related to hyperref, but we put it
% here anyway.
%    \begin{macrocode}
\RequirePackage{etoolbox}
\RequirePackage{hyperref}
\hypersetup{hidelinks}
\pdfminorversion=7
%    \end{macrocode}
% We patch the \cs{@combinedblfloats} following a suggestion by Heiko
% Oberdiek at \url{https://tex.stackexchange.com/a/249743/3843}, which
% avoids problems when hyperlinks are split across
% pages. \changes{4.0}{2019/11/12}{Fix for issue where breaking
% hyperlink across pages gave an error.}
%    \begin{macrocode}
\urlstyle{same}
\newcommand\email{\begingroup \urlstyle{rm}\Url}
\patchcmd\@combinedblfloats{\box\@outputbox}{\unvbox\@outputbox}{}{%
  \errmessage{\noexpand\@combinedblfloats could not be patched}%
}%
%    \end{macrocode}
% \subsection{Other setup stuff}\label{sec:setup}
% Some declarations: 
%    \begin{macrocode}
\newif\if@needaddresses
\@needaddressestrue
\newif\if@frontmatter
\@frontmatterfalse
\newif\if@inpreamble
\@inpreambletrue
\newif\if@firstauthor
\@firstauthortrue
\newif\if@firstaffil
\@firstaffiltrue
\newif\if@firstaltaffil
\@firstaltaffiltrue
\newif\if@firstkwd
\@firstkwdtrue
\newif\if@noexpandADS
\@noexpandADSfalse
\newif\if@sessiontitlepages
\@sessiontitlepagesfalse
%    \end{macrocode}
% \begin{macro}{\@rmaaEmpty}
% \begin{macro}{\@rmaaArg}
% These are used in various commands to check if arg was left blank
%    \begin{macrocode}
\newcommand{\@rmaaEmpty}{}
\newcommand{\@rmaaArg}{}
%    \end{macrocode}  
% \end{macro}
% \end{macro}
%
% \subsection{Compatibility with old \LaTeX{} versions}
% \label{sec:compat}
% This package is based on the 1998/12/01 release of \LaTeX{}.
% Older releases may not define all the commands that we need.
% Luckily, we can use |\providecommand| to define these if necessary. 
% The following commands are not defined in the 1994/06/01 release. 
%
% \begin{macro}{\MakeUpperCase}
% \begin{macro}{\textsuperscript}
% \begin{macro}{\textbackslash}
% User commands:
%    \begin{macrocode}
\providecommand{\MakeUpperCase}[1]{\uppercase{#1}}
\providecommand{\textsuperscript}[1]{\ensuremath{^\textrm{#1}}}
\providecommand{\textbackslash}{\ensuremath{\backslash}}
%    \end{macrocode}  
% \end{macro}
% \end{macro}
% \end{macro}
% \begin{macro}{\protected@xdef}
% \begin{macro}{\restore@protect}
% \begin{macro}{\@unexpandable@protect}
% \begin{macro}{\hb@xt@}
% \LaTeX{} internal commands:
%    \begin{macrocode}
\providecommand{\protected@xdef}{%
        \let\@@protect\protect
        \let\protect\@unexpandable@protect
        \afterassignment\restore@protect
        \xdef}
\providecommand{\restore@protect}{\let\protect\@@protect}
\providecommand{\@unexpandable@protect}{\noexpand\protect\noexpand}
\providecommand{\hb@xt@}{\hbox to}
%    \end{macrocode}
% \end{macro}
% \end{macro}
% \end{macro}
% \end{macro}
%
% \subsection{Page layout}\label{sec:page-layout}
%
% The dimensions of the text were defined by Elizabeth. Text width is
% 169\,mm with equal margins on the two sides: 
%    \begin{macrocode}
\setlength{\textwidth}{169mm}
\setlength{\@tempdima}{\paperwidth}
\addtolength{\@tempdima}{-\textwidth}
\setlength{\oddsidemargin}{0.5\@tempdima}
\addtolength{\oddsidemargin}{-1in}
\setlength{\evensidemargin}{\oddsidemargin}
%    \end{macrocode}
% Text height is 225\,mm, top margin is 14\,mm. 
%    \begin{macrocode}
\setlength{\textheight}{225mm} 
\setlength{\topmargin}{-1in} 
\addtolength{\topmargin}{14mm}
%    \end{macrocode}
% Space between columns in two-column styles, header, footer,
% paragraph indentation:
%    \begin{macrocode}
\setlength{\columnsep}{9.5mm}
\setlength{\headsep}{0.7cm}
\setlength{\footskip}{0.75cm}
\setlength{\parindent}{1.5em}
%    \end{macrocode}
% 
% \begin{macro}{\@set@scalable@length}
% Try out a new scheme for dealing with various spacing dimensions and
% allowing the user to easily modify them. The command
% \begin{quote}
% |\set@scalable@length|\marg{name}\marg{defaultlength}\marg{stretch}\marg{shrink}
% \end{quote}
% does a 
% |\setlength|\marg{\cs{name}}\marg{defaultlength plus stretch minus shrink}
% and also defines the user command
% \textbackslash\meta{name}|Scale|\marg{factor} that can 
% be used to rescale length by \meta{factor}. Additionally, now
% defines the package command
% \textbackslash\meta{name}|@set@scale|\marg{FACTOR} 
% that also rescales the length by \meta{FACTOR}, so that the length
% becomes \meta{factor} $\times$ \meta{FACTOR} $\times$
% \meta{defaultlength}.   
%    \begin{macrocode}
\newcommand\@set@scalable@length[4]{%
  %% arg #1: NAME of length without leading `\'
  %%
  %% arg #2: default VALUE  
  \expandafter\newlength\csname @#1@default\endcsname
  \expandafter\setlength\csname @#1@default\endcsname{#2}%
  \expandafter\newlength\csname @@#1@@default\endcsname
  \expandafter\setlength\csname @@#1@@default\endcsname{#2}%
  %% arg #3: default STRETCH  
  \expandafter\newlength\csname @#1@stretch\endcsname
  \expandafter\setlength\csname @#1@stretch\endcsname{#3}%
  \expandafter\newlength\csname @@#1@@stretch\endcsname
  \expandafter\setlength\csname @@#1@@stretch\endcsname{#3}%
  %% arg #4: default SHRINK  
  \expandafter\newlength\csname @#1@shrink\endcsname
  \expandafter\setlength\csname @#1@shrink\endcsname{#4}%
  \expandafter\newlength\csname @@#1@@shrink\endcsname
  \expandafter\setlength\csname @@#1@@shrink\endcsname{#4}%
  %%
  %% User-defined scale \@NAME@scale is unity by default
  \expandafter\gdef\csname @#1@scale\endcsname{1}%
  %% Package-defined scale \@NAME@SCALE is unity by default
  \expandafter\gdef\csname @#1@SCALE\endcsname{1}%
  %%
  %% internal command \@reset@NAME should expand to 
  %% \setlength\@@NAME@@default{\@NAME@SCALE\@NAME@default}
  %% \setlength\NAME{\@NAME@scale\@@NAME@@default}
  %% in other words: scale x SCALE x default 
  \expandafter\gdef\csname @reset@#1\endcsname{%
    \expandafter\setlength\csname @@#1@@default\endcsname{%
      \csname @#1@SCALE\endcsname \csname @#1@default\endcsname
      }
    \expandafter\setlength\csname @@#1@@stretch\endcsname{%
      \csname @#1@SCALE\endcsname \csname @#1@stretch\endcsname
      }
    \expandafter\setlength\csname @@#1@@shrink\endcsname{%
      \csname @#1@SCALE\endcsname \csname @#1@shrink\endcsname
      }
    \expandafter\setlength\csname #1\endcsname{%
      \csname @#1@scale\endcsname \csname @@#1@@default\endcsname
      \@plus
      \csname @#1@scale\endcsname \csname @@#1@@stretch\endcsname
      \@minus
      \csname @#1@scale\endcsname \csname @@#1@@shrink\endcsname
      }}
  %%
  %% User command \NAMEScale should expand to 
  %% {\renewcommand\@NAME@scale{#1}\@reset@NAME}
  \expandafter\gdef\csname #1Scale\endcsname##1{%
    \expandafter\gdef\csname @#1@scale\endcsname{##1}%
    \csname @reset@#1\endcsname}
  %%
  %% Package command \NAME@set@SCALE should expand to 
  %% {\renewcommand\@NAME@SCALE{#1}\@reset@NAME}
  \expandafter\gdef\csname #1@set@SCALE\endcsname##1{%
    \expandafter\gdef\csname @#1@SCALE\endcsname{##1}%
    \csname @reset@#1\endcsname}
  %%
  %% initialize at default values with \@reset@NAME
  \csname @reset@#1\endcsname
  }
%    \end{macrocode}
% \end{macro}
% \begin{macro}{\parskip}
% First use of this new infrastructure: set interparagraph space to be
% zero but with a small amount of stretchiness.
%    \begin{macrocode}
%   \ifthenelse{\boolean{@manuscript}}{%
%     \@set@scalable@length{parskip}{0.5\baselineskip}{0.3\p@}{0.1\p@}
%   }{
    \@set@scalable@length{parskip}{0\p@}{0.3\p@}{0\p@}
%   }
%    \end{macrocode}
% \end{macro}
% 
% \begin{macro}{\abovesectionskip}
% \changes{4.0}{2019/11/12}{Use negative values for skip so that first paragraph after section headings is not indented.  This is something that I have been trying to fix for years!  It turns out that the anser was in the documentation for the standard command \cs{@startsection.}}
% \begin{macro}{\belowsectionskip}
% \begin{macro}{\aboveparaskip}
% \begin{macro}{\belowparaskip}
% Define some utility rubber lengths for use in
% section~\ref{sec:section-commands}. 
%    \begin{macrocode}
\newlength{\abovesectionskip}
\newlength{\belowsectionskip}
\newlength{\aboveparaskip}
\newlength{\belowparaskip}
%    \end{macrocode}
% The default values for these lengths are those traditionally used
% for the onecolumn style. 
%    \begin{macrocode}
\@set@scalable@length{abovesectionskip}{-3.5ex}{-4ex}{-2ex}
\@set@scalable@length{belowsectionskip}{2.3ex}{0ex}{2ex}
\@set@scalable@length{aboveparaskip}{3.25ex}{2ex}{0.5ex}
\@set@scalable@length{belowparaskip}{-1em}{\z@}{\z@}
%    \end{macrocode}
% \end{macro}\end{macro}\end{macro}\end{macro}
% We also make scalable some of the standard float parameters. These
% are the same defaults as in the standard article classes, but
% slightly more stretchy. For the twocolumn style, they are currently
% halved, see below. Note that I have also lumped |\abovecaptionskip|
% in with these, since it plays a similar role. There is also a
% |\belowcaptionskip| but it is set to zero in the standard classes. 
%    \begin{macrocode}
\@set@scalable@length{floatsep}{12\p@}{4\p@}{4\p@}
\@set@scalable@length{textfloatsep}{20\p@}{4\p@}{8\p@}
\@set@scalable@length{intextsep}{12\p@}{4\p@}{4\p@}
\@set@scalable@length{dblfloatsep}{12\p@}{4\p@}{4\p@}
\@set@scalable@length{dbltextfloatsep}{20\p@}{4\p@}{8\p@}
\@set@scalable@length{abovecaptionskip}{10\p@}{\z@}{\z@}
%    \end{macrocode}
% \begin{macro}{\@before@title@skip}         
% \begin{macro}{\@before@author@skip}        
% \begin{macro}{\@between@author@skip}       
% \begin{macro}{\@before@abstract@skip}      
% \begin{macro}{\@after@abstracthead@skip}   
% \begin{macro}{\@between@abstract@skip}     
% \begin{macro}{\@before@kwd@skip}           
% \begin{macro}{\@before@dates@skip}           
% Spacing parameters used in |\@maketitle|. These are not changed by
% default except by the user, so they don't need to go into
% |\@package@default@lengths|.  
%    \begin{macrocode}
\newlength{\@before@title@skip}         
\newlength{\@before@author@skip}        
\newlength{\@between@author@skip}       
\newlength{\@before@abstract@skip}      
\newlength{\@after@abstracthead@skip}   
\newlength{\@between@abstract@skip}     
\newlength{\@before@kwd@skip}           
\newlength{\@before@dates@skip}           
\@set@scalable@length{@before@title@skip}{2em}{\z@}{\z@}
\@set@scalable@length{@before@author@skip}{1.5em}{\z@}{\z@}
\@set@scalable@length{@between@author@skip}{0.5em}{\z@}{\z@}
\@set@scalable@length{@before@abstract@skip}{3ex}{1ex}{0.5ex}
\@set@scalable@length{@after@abstracthead@skip}{1ex}{\z@}{\z@}
\@set@scalable@length{@between@abstract@skip}{2ex}{1ex}{0.5ex}
\@set@scalable@length{@before@kwd@skip}{1ex}{0.5ex}{0.2ex}
\@set@scalable@length{@before@dates@skip}{2ex}{0.5ex}{0.2ex}
%    \end{macrocode}
% \end{macro}\end{macro}\end{macro}\end{macro}
% \end{macro}\end{macro}\end{macro}\end{macro}
% 
% Spacing above and below displayed equations also wants to be done
% like this, but it is more complicated since the standard article
% classes make these change with the fontsize. Something to do later,
% maybe. 
% 
%    \begin{macrocode}
%    \end{macrocode}
% \begin{macro}{\@package@default@lengths}
% The value actually used for these lengths is controlled by the
% \meta{SCALE} parameter, which depends on the style option that is in
% effect. This is set using the |\@package@default@lengths| command,
% which must be called every time the style changes (e.g., changing
% from onecolumn to twocolumn). 
%    \begin{macrocode}
\newcommand{\@package@default@lengths}{%
  \ifthenelse{\boolean{@rmtwocol}}{%
    %% twocolumn mode
    \abovesectionskip@set@SCALE{0.5}
    \belowsectionskip@set@SCALE{0.5}
    \aboveparaskip@set@SCALE{0.5}
    \belowparaskip@set@SCALE{1}
    \floatsep@set@SCALE{0.5}
    \textfloatsep@set@SCALE{0.5}
    \intextsep@set@SCALE{0.5}
    \dblfloatsep@set@SCALE{0.5}
    \dbltextfloatsep@set@SCALE{0.5}
    }{%
    %% onecolumn mode
    \abovesectionskip@set@SCALE{1}
    \belowsectionskip@set@SCALE{1}
    \aboveparaskip@set@SCALE{1}
    \belowparaskip@set@SCALE{1}
    \floatsep@set@SCALE{1}
    \textfloatsep@set@SCALE{1}
    \intextsep@set@SCALE{1}
    \dblfloatsep@set@SCALE{1}
    \dbltextfloatsep@set@SCALE{1}
    }%
  }
\@package@default@lengths
%    \end{macrocode}
% \end{macro}
% \begin{macro}{\RescaleLengths}
% \begin{macro}{\RescaleSecLengths}
% \begin{macro}{\RescaleFloatLengths}
% \begin{macro}{\RescaleTitleLengths}
% User commands that will fiddle with all the scalable lengths. 
%    \begin{macrocode}
\newcommand{\RescaleLengths}[1]{%
  \RescaleSecLengths{#1}\RescaleFloatLengths{#1}\RescaleTitleLengths{#1}}
\newcommand{\RescaleSecLengths}[1]{%
%% sections
  \abovesectionskipScale{#1}
  \belowsectionskipScale{#1}
  \aboveparaskipScale{#1}
  }
\newcommand{\RescaleFloatLengths}[1]{%
%% floats
  \floatsepScale{#1}
  \textfloatsepScale{#1}
  \intextsepScale{#1}
  \dblfloatsepScale{#1}
  \dbltextfloatsepScale{#1}
  \abovecaptionskipScale{#1}
  }
\newcommand{\RescaleTitleLengths}[1]{%
%% maketitle
  \@before@title@skipScale{#1}           
  \@before@author@skipScale{#1}        
  \@between@author@skipScale{#1}       
  \@before@abstract@skipScale{#1}      
  \@after@abstracthead@skipScale{#1}   
  \@between@abstract@skipScale{#1}     
  \@before@kwd@skipScale{#1}           
  \@before@dates@skipScale{#1}           
  }
%    \end{macrocode}
% \end{macro}\end{macro}\end{macro}\end{macro}
% \changes{2.2beta14}{2002/03/17}{Finished implementation of new
% infrastructure for scalable lengths. Converted the above/below
% section skips to use it. Also, all the floatsep parameters. Also,
% all the vertical lengths that are used by |\@maketitle|. }

% \begin{macro}{\MSextramargin}
% For the \texttt{manuscript} option we have an extra 0.3\,inch
% vertical margin and 1\,inch horizontal margin for marginal notes (by
% hand or with |\marginpar|. 
%
% \changes{3.23}{2007/09/09}{Reduce baselinestretch in
%   \texttt{manuscript} style to 1.2. No longer use raggedright for
%   manuscript.}
%
% \changes{3.24}{2007/09/14}{Further reduction of baselinestretch in
%   \texttt{manuscript} style to 1.1 since we now have wider margins.}
%
% \changes{3.24}{2007/09/14}{Split \texttt{manuscript}'s extra margin
%   into separate horizontal and vertical parts.}
% 
%    \begin{macrocode}
\newlength{\MSextraHmargin}\newlength{\MSextraVmargin}
\setlength{\MSextraHmargin}{1in} 
\setlength{\MSextraVmargin}{0.3in} 
\if@manuscript
   \renewcommand{\baselinestretch}{1.1}
   \addtolength{\textwidth}{-2\MSextraHmargin}
   \addtolength{\oddsidemargin}{\MSextraHmargin}
   \setlength{\evensidemargin}{\oddsidemargin}
   \addtolength{\textheight}{-2\MSextraVmargin}
   \addtolength{\topmargin}{\MSextraVmargin}
%    \raggedright
\fi
%    \end{macrocode}
% \end{macro}
% Float parameters: don't allow bottom floats by default. 
%    \begin{macrocode}
\renewcommand{\fps@figure}{tp}
\renewcommand{\fps@table}{tp}
%    \end{macrocode}
%
% \begin{macro}{\rmaa@AuthorWidthFrac}
% \begin{macro}{\rmaa@AbstractWidthFrac}
% These parameters control what fraction of the total text width can
% be used by the list of authors and the abstract/resumen,
% respectively.  
%    \begin{macrocode}
\newcommand{\rmaa@AuthorWidthFrac}{1.0}
\newcommand{\rmaa@AbstractWidthFrac}{0.75}
%    \end{macrocode}
% Use a wide abstract to save space in the new style.
%    \begin{macrocode}
\if@newstyle\renewcommand{\rmaa@AbstractWidthFrac}{1.0}\fi
%    \end{macrocode}
% \end{macro}
% \end{macro}
%
% Some things to consider in future versions. Trying to avoid too many
% widows, etc. Apparently, the following will work best with
% \cs{flushbottom}. 
%    \begin{macrocode}
%%\widowpenalty=10000
%%\clubpenalty=10000
%    \end{macrocode}
% Also, should perhaps try getting rid of underfull/overfull box warnings. 
% \changes{3.10}{2003/01/24}{Set large values for \cs{hbadness},
% \cs{vbadness}, \cs{hfuzz}, \cs{vfuzz}. To reduce amount of stuff
% that gets written to the terminal every time you compile something.}
%    \begin{macrocode}
\hbadness=10000
\vbadness=10000
\setlength\hfuzz{\maxdimen}
\setlength\vfuzz{\maxdimen}
%    \end{macrocode}
% \subsection{Various text styles for different contexts}
% \label{sec:text-styles} 
% \begin{macro}{\rmaa@AffilStyle}
% \begin{macro}{\rmaa@PagesStyle}
% \begin{macro}{\rmaa@YearStyle}
% \begin{macro}{\rmaa@VolumeStyle}
% \begin{macro}{\rmaa@VolumenStyle}
% \begin{macro}{\rmaa@JournalNameStyle}
% \begin{macro}{\rmaa@ReceivedStyle}
% \begin{macro}{\rmaa@AbstractHeaderStyle}
% \begin{macro}{\rmaa@JournalNameStyle}
% Style commands for the different fields used by |\maketitle| in
% constructing the title page of the article. These have the form
% |\rmaa@XYZStyle{TEXT}|. I'm not sure these were really a good idea,
% but never mind \dots 
%    \begin{macrocode}
\newcommand{\rmaa@AffilStyle}[1]{{\normalsize\rmfamily#1}}
\newcommand{\rmaa@PagesStyle}[1]{{\normalsize\rmfamily#1}}
\newcommand{\rmaa@YearStyle}[1]{{\normalsize\rmfamily#1}}
\newcommand{\rmaa@VolumeStyle}[1]{{\normalsize\rmfamily\bfseries#1}}
\newcommand{\rmaa@VolumenStyle}[1]{{\Large\sffamily\bfseries#1}}
\newcommand{\rmaa@JournalNameStyle}[1]{{\normalsize\itshape#1}}
\newcommand{\rmaa@ReceivedStyle}[1]{{\normalsize\rmfamily\itshape#1}}
\newcommand{\rmaa@AbstractHeaderStyle}[1]{%
  {\normalsize\rmfamily\MakeUpperCase{#1}}}
\newcommand{\rmaa@TitleStyle}[1]{%
  {\large\bfseries\rmfamily\boldmath\MakeUpperCase{#1}\par}}
%    \end{macrocode}
% \end{macro}
% \end{macro}
% \end{macro}
% \end{macro}
% \end{macro}
% \end{macro}
% \end{macro}
% \end{macro}
% \end{macro}
%
% \begin{macro}{\rmaa@FrontSecStyle}
% Headings of ``front sections'' (see section~\ref{sec:front-sections})
%    \begin{macrocode}
\newcommand{\rmaa@FrontSecStyle}[1]{%
  {\normalsize\bfseries\rmfamily\MakeUpperCase{#1}}}
%    \end{macrocode}
% \end{macro}
%
% \begin{macro}{\rmaa@TOCTitleStyle}
% \begin{macro}{\rmaa@TOCPartStyle}
% \begin{macro}{\rmaa@IndexTitleStyle}
% \begin{macro}{\rmaaTOCAuthorStyle}
% \begin{macro}{\rmaaIndexAuthorStyle}
% Commands to set style of TOC and index entries.
% Rewritten to avoid unnecessary font commands in the .toc and .ind files
%    \begin{macrocode}
\newcommand{\rmaa@TOCTitleStyle}[1]{\MakeUpperCase{#1}}
\newcommand{\rmaa@TOCPartStyle}[1]{\rmaa@TOCTitleStyle{#1}}
\newcommand{\rmaa@IndexTitleStyle}[1]{#1}
\newcommand{\rmaaTOCAuthorStyle}[1]{%
  {\itshape#1}}
\newcommand{\rmaaIndexAuthorStyle}[1]{%
  \protect\mbox{\bfseries#1}}
%    \end{macrocode}
% \end{macro}
% \end{macro}
% \end{macro}
% \end{macro}
% \end{macro}
%
% \begin{macro}{\rmaa@KeywordStyle}
% \begin{macro}{\rmaa@KeywordHeaderStyle}
% The style of the keywords depends on whether we are using the
% \texttt{newstyle} option or not. If we are, then they are small. 
%    \begin{macrocode}
\ifthenelse{\boolean{@newstyle}}{
  \newcommand{\rmaa@KeywordStyle}[1]{%
    {\small%
      \if@boldkeywords 
      \bfseries\MakeUpperCase{\raggedright #1}
      \else
      \raggedright #1
      \fi
    }}
  \newcommand{\rmaa@KeywordHeaderStyle}[1]{%
    {\small\itshape#1}}}
{
  \newcommand{\rmaa@KeywordStyle}[1]{%
    {\normalsize%
      \if@boldkeywords\bfseries\fi\raggedright #1}}
  \newcommand{\rmaa@KeywordHeaderStyle}[1]{%
    {\normalsize\itshape#1}}}
%    \end{macrocode}
% \end{macro}
% \end{macro}
%
% \subsection{Page styles}\label{sec:page-styles}
% \begin{macro}{\ps@rmaa}
% The running headers for normal pages:
%    \begin{macrocode}
\newcommand{\ps@rmaa}{%
%    \end{macrocode}
% nothing in either foot; 
%    \begin{macrocode}
  \renewcommand\@oddfoot{\relax}%
  \renewcommand\@evenfoot{\relax}%
%    \end{macrocode}
% page number and short author list in the even heads; 
%    \begin{macrocode}
  \renewcommand\@evenhead{\rlap{\thepage}\hfil\@shortauthor\hfil}%
%    \end{macrocode}
% short title and page number in the odd heads.
%    \begin{macrocode}
  \renewcommand\@oddhead{\hfil\@shorttitle\hfil\llap{\thepage}}
  }
%    \end{macrocode}
% \begin{macro}{\ps@rmaafirstpage}
% The header of the first page of each
% article:
%    \begin{macrocode}
\newcommand{\ps@rmaafirstpage}{%
  \renewcommand\@oddfoot{\hfil\thepage}%
%    \end{macrocode}
% footer corrected (06 Jun 2000) to have page on outside edge;
%    \begin{macrocode}
  \renewcommand\@evenfoot{\thepage\hfil}%
%    \end{macrocode}
% Header has journal name, volume, and page numbers for a regular
% article, or similar for special cases. 
% \changes{4.0}{2019/11/03}{Change to three-line header with copyright info and DOI.}
%    \begin{macrocode}
\if@preprint
\renewcommand\@oddhead{%
  \if@proceedings 
  \textsc{To appear in ``\@conftitle{} (\rmaa@YearStyle{\@year})''}\hfil
  \textit{RMxAC}
  \else
  \textsc{To appear in} \@journalname\hfil
  \fi}\fi%
\if@guide
\renewcommand\@oddhead{\textsc{Documentation} (\rmaa@YearStyle{\@year})\hfil}
\fi%
\if@manuscript
\renewcommand\@oddhead{Manuscript for \@journalname{} (\rmaa@YearStyle{\@year})\hfil}
\fi%
\if@journal
\renewcommand\@oddhead{%
  \parbox{\textwidth}{
    \@journalname, \rmaa@VolumeStyle{\@volume},
    \rmaa@PagesStyle{\@artpages}\
    (\rmaa@YearStyle{\@year})\hfil\\[1pt]
    {\small\textcopyright{} \@year:
      Instituto de Astronom\'\i{}a,
      Universidad Nacional Aut\'onoma de M\'exico}\hfil\\[1pt]
    {\small\@DOI@url}\hfil
  }%
}\fi%
  \renewcommand\@evenhead{\@oddhead}}
%    \end{macrocode}
% By default, use the |\ps@rmaa| page style.
%    \begin{macrocode}
\pagestyle{rmaa}
%    \end{macrocode}
% \end{macro}
% \end{macro}
%
% \subsection{Preamble commands}\label{sec:preamble} 
% Everything you need for setting the title page, TOC and index
% entries, etc.
% \subsubsection{Commands for the title page}
% \label{sec:title-commands}
% \begin{macro}{\resumen}
% \begin{macro}{\@resumen}
% \begin{macro}{\abstract}
% \begin{macro}{\@abstract}
% User commands for setting the resumen and abstract. 
% These should be used in the document header. The internal commands
% |\@resumen| and |\@abstract| are used to hold the actual text
% pending output via the |\maketitle| command. 
%    \begin{macrocode}
\newif\if@setresumen
\newcommand{\resumen}[1]{%
%%  \renewcommand{\@resumen}{{\selectlanguage{spanish}#1}}%
  \renewcommand{\@resumen}{{\ignorespaces#1}}%
  \@setresumentrue}
\renewcommand{\abstract}[1]{%
  \renewcommand{\@abstract}{{\ignorespaces#1}}
  \renewcommand\ADS@abstract{#1}
%% Not a good idea, since then author macros are not expanded and
%% cause errors later. However, by not doing it we are vulnerable to
%% fragile commands.....
  \if@noexpandADS\@onelevel@sanitize{\ADS@abstract}\fi
  \@abstractpresenttrue
  }
\newcommand{\ADSnoexpand}{\@noexpandADStrue}
%    \end{macrocode}
% \begin{macro}{\boldabstract}
% This is the version for doing 1-page conference proceedings. Still
% save the text for ADS but then write it out immediately as a bold
% paragraph. 
%    \begin{macrocode}
\newcommand{\boldabstract}[1]{%
  \gdef\ADS@abstract{#1}%
  \renewcommand{\@abstract}{#1}%
  \if@noexpandADS\@onelevel@sanitize{\ADS@abstract}\fi%
  \noindent\textbf{\boldmath #1}
  }
%    \end{macrocode}
% \end{macro}
% Default value in each case is to print out a warning if not set
%    \begin{macrocode}
  \newcommand{\@resumen}{%
    \@setresumenfalse \ClassWarningNoLine{rmaa}{No Spanish abstract
      (resumen) specified - if you do not speak Spanish, this will be
      translated by the editors}%
    \textbf{Favor de proporcionar un resumen en espa\~nol. If you are
      unable to translate your abstract into Spanish, the editors will
      do it for you.} }
  \newcommand{\@abstract}{\ClassWarningNoLine{rmaa}{No abstract
      specified}}
%    \end{macrocode}
% \end{macro}
% \end{macro}
% \end{macro}
% \end{macro}
%
% \begin{macro}{\title}
% \begin{macro}{\@rawtitle}
% \begin{macro}{\@TOCtitle}
% \begin{macro}{\rmaaIndexTitle}
% Similarly, |\title| is the preamble command to set the title of the
% article. This is used in various places, typeset in different
% styles, hence the use of the internal |\@rawtitle|, |\@title|,
% |\rmaaIndexTitle|, and |\@TOCtitle| commands, each of which uses the
% relevant style from section~\ref{sec:text-styles}
%    \begin{macrocode}
\newcommand{\@rawtitle}{Raw Title}
\newcommand{\@TOCtitle}{\rmaa@TOCTitleStyle{Title for TOC}}
\newcommand{\rmaaIndexTitle}{\rmaa@IndexTitleStyle{Title for index}.}
\newcommand{\@TOCaftertitle}{}
\renewcommand{\title}[1]{%
%%  {\ADS@sanitize\xdef\@rawtitle{#1}}
  \renewcommand\@rawtitle{\noexpand#1}
  \let\ADS@title=\@rawtitle
  \renewcommand{\@title}{\rmaa@TitleStyle{#1}}%
  \renewcommand{\@TOCtitle}{\rmaa@TOCTitleStyle{#1}\@TOCaftertitle}%
  \renewcommand{\rmaaIndexTitle}{\rmaa@IndexTitleStyle{#1\@TOCaftertitle}.}%
  }
%    \end{macrocode}
% \end{macro}
% \end{macro}
% \end{macro}
% \end{macro}
% 
% All the weird and whacky ways of indicating the author
% affiliations. 
% \begin{macro}{\affil}
% First, the way that sets it directly underneath the authors name.
% Note that |\affil| trumps all other affiliations for the purposes of
% getting into the ADS stuff (we use |\gdef|).   
% \dots 
%    \begin{macrocode}
\newcommand{\affil}[1]{%
  \\[1.5ex]\rmaa@AffilStyle{#1}\\[1.5ex]%
  \if@firstaffil\gdef\ADS@affil{#1}\@firstaffilfalse\fi%
  }
\renewcommand{\and}{and\\[1.5ex]}
%    \end{macrocode}
% \begin{macro}{\altaffil}
% Alternative affiliations are footnoted. 
%    \begin{macrocode}
\newcommand{\altaffil}[1]{%
  \if@firstaltaffil\gdef\ADS@affil{#1}\@firstaltaffilfalse\fi%
  \thanks{#1}}
%    \end{macrocode}
% \begin{macro}{\altaffilmark}
% \begin{macro}{\altaffiltext}
% Alternative-alternative mechanism for affiliations---should not be
% mixed with |\altaffil| (unless you know what you are doing) or there
% will be problems with numbering.
%    \begin{macrocode}
\newcommand{\altaffilmark}[1]{\textsuperscript{#1}}
\newcommand{\altaffiltext}[2]{%
  \if@firstaltaffil\gdef\ADS@affil{#2}\@firstaltaffilfalse\fi
  \stepcounter{footnote}%
  \protected@xdef\@thanks{\@thanks
    \protect\footnotetext[#1]{#2}}} 
%    \end{macrocode}
% \end{macro}
% \end{macro}
% \end{macro}
% \end{macro}
%
% \begin{macro}{\keywords}
% \begin{macro}{\@keywords}
% \begin{macro}{\@keywordheader}
% \label{keywords}Sundry minor preamble commands. All
% self-explanatory, I hope. |\keywords| is the old user interface
% (deprecated). New articles should use |\addkeyword| once for each
% keyword. That way the ADS record comes out right too (comma rather
% than em-dash separator).
% \changes{2.2beta12}{2001/12/07}{\cs{addkeyword} New user interface
%   for keywords.} 
% \changes{2.2beta13}{2002/03/11}{\cs{addkeyword} Improved behaviour
%   with .ads file. Use
%   \cs{ADS@sanitize} to change definitions of `~' and \cs{ion}.}
% \changes{3.23}{2007/09/09}{\cs{addkeyword} now signals a compilation
% error if the author tries to add more than 6 keywords.}
% \changes{3.24}{2007/09/14}{kwdcounter is now reset in
% \cs{@begofarticlehook} to prevent error in book style} 
%    \begin{macrocode}
\newcommand{\@kwdsep}{ --- }
\newcommand{\ADS@kwdsep}{, }
\newcounter{kwdcounter}\setcounter{kwdcounter}{0}
\newcommand{\addkeyword}[1]{%
  \addtocounter{kwdcounter}{1}
  \ifthenelse{\thekwdcounter > 6}{
    \ClassError{rmaa}{Too many keywords. A maximum of 6 keywords is allowed.}{}
  }{}
  \if@firstkwd%
  \renewcommand\@keywords{\if@boldkeywords\uppercase{#1}\else #1\fi}
  {\ADS@sanitize\def\ion{\plainion}\gdef\ADS@keywords{#1}}
  \@firstkwdfalse
  \else%
  \edef\@keywords{\@keywords\@kwdsep \if@boldkeywords\uppercase{#1}\else #1\fi}%
  {\ADS@sanitize\def\ion{\plainion}\xdef\ADS@keywords{\ADS@keywords\ADS@kwdsep #1}}%
  \fi}
\newcommand{\ADS@sanitize}{%
  \def\mathrm{}
  \def\@{\ }
  \def~{\ }
  \def\langle{<}
  \def\rangle{>}
  \def\log{log }
  \def\sin{sin }
  \def\cos{cos }
  \def\tan{tan }
  \def\fd{^d.}
  \def\fh{^h.}
  \def\fm{^m.}
  \def\fs{^s.}
  \def\fdg{�.}
  \def\farcm{'.}
  \def\farcs{''.}
  \def\arcdeg{�}
  \def\,{ }
  }
\newcommand{\keywords}[1]{%
  \renewcommand{\ADS@keywords}{#1}
  \renewcommand{\@keywords}{\uppercase{#1}}
  }
\newcommand{\@keywordheader}{\rmaa@KeywordHeaderStyle{Key Words:}}
\newcommand{\@keywords}{%
  \ClassWarningNoLine{rmaa}{No keywords specified}%
  Keywords go here
  }
%    \end{macrocode}
% \begin{macro}{\shorttitle}
% \begin{macro}{\shortauthor}
% \begin{macro}{\@shorttitle}
% \begin{macro}{\@shortauthor}
% Commands for setting the running headers.
%    \begin{macrocode}
\newcommand{\shorttitle}[1]{%
  \renewcommand\@shorttitle{\MakeUpperCase{#1}}}
\newcommand{\shortauthor}[1]{%
  \renewcommand\@shortauthor{\MakeUpperCase{#1}}}
\newcommand{\@shorttitle}{SHORT TITLE GOES HERE}
\newcommand{\@shortauthor}{SHORT AUTHOR GOES HERE}
%    \end{macrocode}
% \end{macro}
% \end{macro}
% \end{macro}
% \end{macro}
% \end{macro}
% \end{macro}
% \end{macro}


% \subsubsection{The full addresses}
% \label{sec:full-addresses}
% \begin{macro}{\fulladdresses}
% \begin{macro}{\@fulladdresses}
% \begin{macro}{\@fulladdheight}
% \begin{macro}{\@fulladddepth}
% \begin{macro}{\@fulladdbox}
% Aaghh, the full addresses! A perennial source of problems. 
%    \begin{macrocode}
\newlength{\@fulladdheight}
\newlength{\@fulladddepth}
\newsavebox{\@fulladdbox}
%% \fulladdresses command modified (11 May 2000) so that we can put
%% \fulladdresses{} in the article to suppress the output of the
%% addresses
\newcommand{\@FinalFigure}{\relax}
\newcommand{\FinalFigure}[1]{\renewcommand{\@FinalFigure}{#1}}
\newcommand{\fulladdresses}[1]{%
  \renewcommand{\@rmaaArg}{#1}
  \ifx\@rmaaArg\@rmaaEmpty
     \renewcommand{\@fulladdresses}{}
  \else
     \renewcommand{\@fulladdresses}{
       \normalsize%
       \ifx\@FinalFigure\relax\else\vfill\@FinalFigure\vfill\fi
%%       \vspace*{1ex}
       %% do addresses as a list (copied from thebibliography environment) 
       \list{\null}{\leftmargin 1.5em\topsep\z@\labelwidth\z@\labelsep\z@\itemindent -1.5em
         \itemsep\z@\parsep\z@\usecounter{enumi}}
       \sloppy\clubpenalty4000\widowpenalty4000
       \sfcode`\.=1000\relax
       #1
       \endlist
       }
  \fi
  }
\newcommand{\@fulladdresses}{Full addresses go here}
%    \end{macrocode}
% \begin{macro}{\suppressfulladdresses}
% New user-level command for suppressing full addresses should you so
% desire. Currently used in the single-page-poster style. 
%    \begin{macrocode}
\newcommand{\suppressfulladdresses}{\@suppressfulladdressestrue\@needaddressesfalse}
%    \end{macrocode}
% \end{macro}
% \end{macro}
% \end{macro}
% \end{macro}
% \end{macro}
% \end{macro}

% \begin{macro}{\printdates}
% \begin{macro}{\noprintdates}
% Command to activate the printing out of the received/accepted
% dates. This can also be done using the \texttt{printdates} option,
% which is the default for \texttt{journal}, but not for \texttt{proceedings}.
%    \begin{macrocode}
\newcommand{\printdates}{\@printdatestrue}
\newcommand{\noprintdates}{\@printdatesfalse}
%    \end{macrocode}
% \end{macro}
% \end{macro}
% \begin{macro}{\@outputaddressesifnecessary}
% Command to output the addresses if we haven't already done so
% This is called by: 
% \begin{enumerate}
% \item The end of |thebibliography| environment
% \item The |\outputfulladdresses| command
% \end{enumerate}
%    \begin{macrocode}
\newcommand{\@outputaddressesifnecessary}{%
  \ClassInfo{rmaa}{Calling @outputaddressesifnecessary}
  \ifthenelse{\boolean{@needaddresses}}{%
    %% If in one column style, or in the first column of a two column
    %% style then we can try to put it on this page 
    \ifthenelse{\boolean{@firstcolumn} \or \not\boolean{@rmtwocol}}{%
      \@outputaddresses}{%
      %% otherwise give warning it will be on a separate page
      \ClassWarning{rmaa}{Full addresses will go on a separate
        page (please see User Guide for how to remedy this)}}%
    %% If we haven't managed yet, put it on its own page
    \ifthenelse{\boolean{@needaddresses}}{%
      \if@debug{\ClassWarning{rmaa}{In this silly bit!}}\fi
      \clearpage\vfill\@outputaddresses}{}%
    }{
    %% If we don't need them, we still might need to do the eoa hook 
    \ifthenelse{\boolean{@suppressfulladdresses}}{%
      \if@debug{\ClassWarning{rmaa}{XXX}}\fi
%%      \if@book\else\@endofarticlehook\null\fi
      }{}
    }%
  }  
%    \end{macrocode}
% \end{macro}

% \begin{macro}{\outputfulladdresses}
% A user interface to the above command, in case there is no biblio
% (11 May 2000) 
%    \begin{macrocode}
\newcommand{\outputfulladdresses}{\@outputaddressesifnecessary}
%    \end{macrocode}
% \end{macro}

% \begin{macro}{\adjustfinalcols}
% This command starts a new column and should be 
% used to ensure that the columns balance on the last page. It also
% outputs the full 
% author addresses to ensure that this occurs in the left hand
% column. 
% The |\par| command is so that the current line will be taken account
% of when working out |\pagetotal|
%    \begin{macrocode}
\newcommand{\adjustfinalcols}{\if@twocolumn\goodbreak\@outputaddresses\newpage\noindent\fi}
%    \end{macrocode}
% \end{macro}

% \begin{macro}{\@outputaddresses}
% Command to do the work of outputting the addresses. 
% This can be called by 
% \begin{enumerate}
% \item an explicit |\adjustfinalcols|
% \item via |\@outputaddressesifnecessary|, q.v.
% \end{enumerate}
%
%    \begin{macrocode}
\newcommand{\@outputaddresses}{%
  % Allow option of suppressing the full addresses, in which case
  % we just run the hook if we are not in a book style
  \if@suppressfulladdresses
%%     \if@book\else\@endofarticlehook\null\fi
  \else
  %% save the full addresses in a box for output later
  \sbox{\@fulladdbox}{\parbox{\textwidth}{\@fulladdresses}}
  % calculate height of address box  
  \settoheight{\@fulladdheight}{\usebox{\@fulladdbox}}
  \settodepth{\@fulladddepth}{\usebox{\@fulladdbox}}
  \addtolength{\@fulladdheight}{\@fulladddepth}
  \if@debug
  { \typeout{Address total height: \the\@fulladdheight}
    \typeout{Text height: \the\pagetotal}
    \typeout{Goal height: \the\pagegoal}
    }
  \fi
  % reuse \@fulladdheight to mean the total height of text+addresses
  \addtolength{\@fulladdheight}{\pagetotal}
%%  \addtolength{\@fulladdheight}{-\baselineskip}
  \if@debug\typeout{Adjusted text + addresses height: \the\@fulladdheight}\fi
  % check it will fit on current page
  \ifdim\@fulladdheight<\pagegoal
     %% must shift it to the left
     %% to compensate for the hanging indentation in the reference list. 
     %% NOTE: this is not foolproof!
     %% NB. This will give overfull hbox warnings, which is a bit ugly.
     \vfill\noindent\if@biblist\hspace*{-\leftmargin}\fi
     \usebox{\@fulladdbox}
     % Set this to show we've done it
     \@needaddressesfalse
     % and assume we are at the end of the article
     % so that we do the EOA hook if not in the book style
%%     \if@book\else\@endofarticlehook\fi
  \else
     % if not, just do nothing. With luck, this function will get 
     % called again on the next page
     \if@debug{\typeout{Holding over full addresses to next page}}\fi
  \fi\fi
  }
%    \end{macrocode}
% \end{macro}

% \subsubsection{Commands for editors}
% \label{sec:editor-commands}
% Some user/editor commands for doing the TOC and index in the book
% option. 
% They are defined here rather in the book-specific part since they
% need to be always available in order to avoid breaking documents. 
% \begin{macro}{\listofauthors}
% \begin{macro}{\@rawauthors}
% \begin{macro}{\rmaaTOCauthors}
% The command |\listofauthors| formats an author list for use in the
% TOC. It should contain the same names as in the |\author| command,
% but without the affiliation garbage and using |\&| instead of
% |and|. 
%    \begin{macrocode}
\newcommand{\@rawauthors}{Raw Authors}
\newcommand{\rmaaTOCauthors}{\noexpand\rmaaTOCAuthorStyle{Author list for TOC}}
\newcommand{\listofauthors}[1]{%
  \renewcommand{\@rawauthors}{#1}
  \renewcommand\rmaaTOCauthors{\noexpand\rmaaTOCAuthorStyle{#1}}}
%    \end{macrocode}
% \begin{macro}{\indexauthor}\label{indexauthor}
% The command |\indexauthor| produces an index entry for each
% author. 
% This command must come \emph{after} the
% |\title| and the |\listofauthors| commands. Also, there should be one 
% |\indexauthor| command for \emph{each} author. We also add the
% author onto the end of the |\ADS@authors| macro and construct the
% |\ADS@authorinitial| by a most cunning method (see below), but only
% for the first author. 
% 
% No, No, NO! This was all wrong. Backed out completely: 
% \begin{quote}
% Completely rewritten (19 Dec 2002) since we now use the index class, so 
% as to be able to alphabetize accented characters. A side effect of this is
% that |\index| no longer expands out the argument after the |@|, so we 
% instead add all the |\index| commands to a list (|\indexauthor@deferlist|), 
% which is executed in the |\@begofarticlehook|. 
% 
% Note the stuff with |\toks0| is copied from Exercise 20.15 of the \TeX{}book---it is
% to make sure the latin1-accented name gets passed into the idx
% file. 
% \end{quote}
% We now go back to the old way of doing things, using an optional
% argument to the \cs{indexauthor} command that has the name without
% accents. 
% \changes{3.08}{2003/01/16}{Remerged changes from 3.06 that got
% dropped from 3.07 after a synchronization disaster. Basically, the
% accentless optional argument to \cs{indexauthor}. Must add the macro
% directories to my unison set-up to stop this sort of thing from happening.}
%    \begin{macrocode}
\newcommand\@isortkey{}
\newcommand\@emptykey{empty}
\newcommand\indexauthor[2][empty]{
  \renewcommand\@isortkey{#1}
  \ifx\@isortkey\@emptykey \renewcommand\@isortkey{#2}\fi
  \index{\@isortkey @{{\rmaaIndexAuthorStyle{#2}}\hskip1em plus1em minus0.4em
      \rmaaIndexTitle\hskip1em plus2em minus0.4em\penalty-100
      \rmaaTOCauthors}}
  \indexauthorADS{#2}
}
%%\def\indexauthor@deferlist{\relax}
%%\long\def\indexauthor#1{
%%  \def\@@name{#1} \toks0=\expandafter{\@@name}
%%  \toks1=\expandafter{\indexauthor@deferlist}
%%  \protected@xdef\indexauthor@deferlist{%
%%    \the\toks1
%%    \noexpand\index{\the\toks0@{\noexpand\rmaaIndexAuthorStyle{\the\toks0}\noexpand\hskip1em plus1em minus0.4em 
%%        \rmaaIndexTitle\noexpand\hskip1em plus2em minus0.4em\noexpand\penalty-100
%%        \rmaaTOCauthors}}
%%  }
%%  \indexauthorADS{#1}
%%  }
%%\xdef\indexauthor@deferlist{}
\newcommand\indexauthorADS[1]{%
  \if@firstauthor\protected@xdef\ADS@authorinitial{\firstletter{#1}}\fi
  \protected@xdef\ADS@authors{\ADS@authors\if@firstauthor\else ; \fi
    #1}
  \@firstauthorfalse
}
%    \end{macrocode}
% Note that the first |#1| sets the key for the index, making sure it
% actually starts with  
% first letter of the surname. Meanwhile, the |@{...}| expression,
% which starts with some  
% control sequences is what is written to the |.ind| file. 
% See the makeindex man page for details. 
% \end{macro}
% \end{macro}
% \end{macro}
% \end{macro}

% \begin{macro}{\KillAllAccents}
% \begin{macro}{\firstletter}
% These are utility functions used in the above. |lett| is a counter
% indicating the current letter in the string. |\firstletter| just
% prints the first letter of its argument.
%    \begin{macrocode}
\newcommand{\KillAllAccents}{\def\'{}\def\`{}\def\~{}\def\.{}\def\i{i}}
\newcommand\firstletter[1]{\@car #1\@nil}
%    \end{macrocode}
% \end{macro}
% \end{macro}
% \begin{macro}{\SetVolume}
% \begin{macro}{\SetMSnumber}
% \begin{macro}{\SetFirstPage}
% \begin{macro}{\SetYear}
% \begin{macro}{\ReceivedDate}
% \begin{macro}{\AcceptedDate}
% These are for the editors to set the volume number, etc. 
%    \begin{macrocode}
\newcommand{\SetVolume}[1]{%
  \renewcommand{\@volume}{#1}
  \edef\ADS@volume{\PadToFourWithDots{#1}}
  \renewcommand{\@volumen}{\rmaa@VolumenStyle{volumen #1}}%
  }
\newcommand{\SetMSnumber}[1]{\renewcommand{\@MSnum}{#1}}
%    \end{macrocode}
% Note that there is no real reason to use |\SetFirstPage| since all
% the page numbering is done automatically in the book style. 
%    \begin{macrocode}
\newcommand{\SetFirstPage}[1]{\setcounter{page}{#1}}
\newcommand{\SetYear}[1]{%
  \renewcommand{\@year}{#1}
  \renewcommand{\ADS@year}{#1}
  \renewcommand{\@anyo}{\rmaa@VolumenStyle{#1}}
  }
\newcommand{\ReceivedDate}[1]{\renewcommand{\@received}{#1}}
\newcommand{\AcceptedDate}[1]{\renewcommand{\@accepted}{#1}}
%    \end{macrocode}
%
% \begin{macro}{\DOI}
% User/editor command to enter the DOI, which will be typeset as a link.
% \changes{4.0}{2019/11/02}{New user command \cs{DOI} to typeset
% DOI in the header of first page and add it to ADS info.}
%    \begin{macrocode}
\newcommand\DOI[1]{%
  \def\@DOI@url{\url{https://doi.org/10.22201/ia.01851101p.#1}}%
  \def\ADS@doi{https://doi.org/10.22201/ia.01851101p.#1}%
}%
%    \end{macrocode}
% \end{macro}
%
% \begin{macro}{\ORCIDiD}
% \changes{4.0}{2019/11/12}{Add mechanism for giving ORCID iD links}
% \begin{macro}{\ORCIDiDfull}
% \begin{macro}{\@orcid@icon}
% \begin{macro}{\@orcid@iconbig}
% Commands to add icon links to ORCID iDs of authors. For example
% \begin{quote}
% |\ORCIDiD|\marg{0000-0000-0000-0000}
% \end{quote}
% to be used after each author name in the |\author|
% command and in the |\fulladdresses|.  Alternatively,
% |\OORCIDiDfull| will typeset the icon plus the full text of the
% link, but we don't use that at present.  The icon image is stored in |ORCIDiD_iconvector.pdf|, which needs to be on the input path.  In principle, we could make EPS versions if necessary.
%    \begin{macrocode}
\newcommand\@orcid@icon{\includegraphics[height=1.8ex]{ORCIDiD_iconvector}}
\newcommand\@orcid@iconbig{\includegraphics[height=2ex]{ORCIDiD_iconvector}}
\newcommand\ORCIDiD[1]{\href{https://orcid.org/#1}{\,\@orcid@icon}}
\newcommand\ORCIDiDfull[1]{%
  \href{https://orcid.org/#1}{\@orcid@iconbig}\,\url{https://orcid.org/#1}}
%    \end{macrocode}
% \end{macro}
% \end{macro}
% \end{macro}
% \end{macro}
%
% \begin{macro}{\@volume}
% \begin{macro}{\@volumen}
% \begin{macro}{\@artpages}
% \begin{macro}{\@year}
% \begin{macro}{\@anyo}
% \begin{macro}{\@received}
% \begin{macro}{\@accepted}
% \begin{macro}{\@MSnum}
% Default values for the containers used by above commands, plus a
% few more useful ones. 
%    \begin{macrocode}
\newcommand{\@volume}{00}
\newcommand{\@volumen}{\rmaa@VolumenStyle{00}}
\newcommand{\@artpages}{%
  \oldpageref{\articlename:RMAAFirstPage}--\oldpageref{\articlename:RMAALastPage}}
\newcommand{\@year}{\the\year}
\newcommand{\@anyo}{\rmaa@VolumenStyle{2007}}
\newcommand{\@journalname}{\rmaa@JournalNameStyle{%
    \if@proceedings\RMAASClong\else\RMAAlong\fi}}
\newcommand{\@received}{}
\newcommand{\@accepted}{}
\newcommand{\@MSnum}{0000}
\newcommand{\@DOI@prefix}{https://doi.org/10.22201/ia.01851101p}
\def\@DOI@url{\url{https://doi.org/10.22201/ia.01851101p.2099.00.00.00}}
%    \end{macrocode}
% \end{macro}
% \end{macro}
% \end{macro}
% \end{macro}
% \end{macro}
% \end{macro}
% \end{macro}
% \end{macro}
% \end{macro}
% \end{macro}
% \end{macro}
% \end{macro}
% \end{macro}
% \end{macro}
% 
% The following commands are most useful in the \texttt{book}
% style. However, the |\SetConfTitle| command can now be used in
% individual articles to set the header for the \texttt{preprint}
% style. 
%    \begin{macrocode}
\newcommand{\@reunionplace}{Lugar de la reuni\'on}
\newcommand{\@reuniondate}{Fecha de la reuni\'on}
\newcommand{\@conftitle}{Conference Title}
\newcommand{\SetReunion}[2]{\renewcommand{\@reunionplace}{#1}\renewcommand{\@reuniondate}{#2}}
\newcommand{\SetConfTitle}[1]{\renewcommand{\@conftitle}{#1}}
\newcommand{\@editors}{Juan Ram\'on de la Fuente and ``El Mosh''}
\newcommand{\SetEditors}[1]{\renewcommand{\@editors}{#1}}
\newcommand{\@editores}{Juan Ram\'on de la Fuente y ``El Mosh''}
\newcommand{\SetEditores}[1]{\renewcommand{\@editores}{#1}}
\newcommand{\@mes}{mes de la publicaci\'on}
\newcommand{\SetMes}[1]{\renewcommand{\@mes}{\rmaa@VolumenStyle{#1}}}
%    \end{macrocode}

% \subsection{Notes for Author Proofs and Communication between Editors}
% \label{sec:notes}
% We introduce a command \cs{Note} that should be used to describe
% any change made to the manuscript. These are a bit tricky to get
% right. In general we use \cs{marginpar} but this only works in outer
% paragraph mode, so we just bung the box where we are if we are in a
% figure caption or something. If we are in the preamble, then we just
% use a \marginpar as part of the \maketitle command.
%    \begin{macrocode}
\newif\if@arenotes
\@arenotesfalse
\newcounter{notecounter}
\setcounter{notecounter}{0}
\newcommand\NotesList{}
\newcommand\Note[1]{%
  \if@proof{%
    \@arenotestrue%
    \stepcounter{notecounter}%
    \if@inpreamble\else{%
      \NoteMark{\arabic{notecounter}}%
    }\fi
    \protected@xdef\NotesList{\NotesList 
      \protect\item #1}%
  }\fi%
}
\newcommand\Changed[2]{Changed ``#1'' to ``#2''.}
\newcommand\NoteMark[1]{%
  \ifinner
  \raisebox{1ex}[0pt][0pt]{\makebox[0pt][l]{\framebox{\footnotesize\bfseries
        N#1}}}%
  \else
  \marginpar{\framebox{\footnotesize\bfseries N#1}}%
  \fi
}
\newcommand\RefNote[1]{\if@proof\NoteMark{\ref{#1}}\relax\fi}
\newcommand\NoteLabel{\protect\label}
\newcommand\OutputNotes{%
  \bgroup
  \onecolumn
  \section*{EDITORIAL NOTES}
  \begin{enumerate}
    \NotesList
  \end{enumerate}
  \egroup
}
\DeclareRobustCommand\cs[1]{\texttt{\char`\\#1}}

%    \end{macrocode}
% \subsection{Bibliographic Records for ADS}
% \label{sec:ADS}
% \changes{2.2beta11}{2001/11/27}{Just about finished all the ADS
% stuff now. Except for the references, which I have given up on for
% the time being.}
% \begin{macro}{\ADS@BibCode}
% \begin{macro}{\ADS@year}
% \begin{macro}{\ADS@journal}
% \begin{macro}{\ADS@volume}
% \begin{macro}{\ADS@firstpage}
% \begin{macro}{\ADS@authorinitial}
% We use the |\ADS@BibCode| command to hold the bibliographic code for
% ADS. It is 
% used in the |%R| 
% field of the bibliographic record for each article. 
% The format is 
% \begin{quote}
% \texttt{YYYYJJJJJVVVVMPPPPA}  
% \end{quote}
% where \texttt{YYYY} is the year, \texttt{JJJJJ} is the journal code,
% which is \texttt{RMxAA} or \texttt{RMxAC}, \texttt{VVVV} is the
% volume number (right-justified), \texttt{M} is not used by us,
% \texttt{PPPP} is the first page number (right-justified), and
% \texttt{A} is the initial of the first author's surname. The volume
% and page fields are padded with periods (`\texttt{.}'), which causes
% us some fun (this is done with the macro |\PadToFourWithDots| described
% below) . 
% 
% All these individual bits are held in the commands |\ADS@year|,
% |\ADS@journal|, |\ADS@volume|, |\ADS@firstpage|,
% |\ADS@authorinitial|. 
%    \begin{macrocode}
\newcommand{\SetADSBibCode}[1]{\renewcommand{\ADS@BibCode}{#1}}
\newcommand{\ADS@BibCode}
{\ADS@year\ADS@journal\ADS@volume.\PadToFourWithDots{\ADS@firstpage}\ADS@authorinitial}
\newcommand{\ADS@year}{2001} 
\newcommand{\ADS@journal}{\if@proceedings RMxAC\else RMxAA\fi}
\newcommand{\ADS@volume}{....} 
%    \end{macrocode}
% These two are just set to default values here, |\ADS@firstpage| is set
% in the |document| environment for included articles, in the
% |frontsection| environment, |\tableofcontents| command,  |index| and
% |abstracts| environments. Also by |\AtBeginDocument| for stand-alone
% articles. 
% 
% God knows how we are going to do |\ADS@authorinitial|!
%    \begin{macrocode}
\newcommand{\ADS@firstpage}{....}
\newcommand{\ADS@authorinitial}{A}
%    \end{macrocode}
% \end{macro}
% \end{macro}
% \end{macro}
% \end{macro}
% \end{macro}
% \end{macro}
% We also have a whole slew of other fields that can be set. These are
% all listed in Table~\ref{tab:ADS}. Most of them still have to be
% defined. 
% \begin{table}[!ht]\raggedleft
% \caption{ADS Bibliographic Records}\label{tab:ADS}
% \begin{tabular}{lp{8em}lp{8em}}\hline
% Code & Description & Macro Name & Comment \\ \hline
% |%R| & Bibliographic Code & |\ADS@BibCode| & required \\
% |%A| & Author list & |\ADS@authors| & required, generate from |\indexauthor| \\
% |%F| & Author affiliation & |\ADS@affil| & format: AA(\dots)
% AB(\dots) \dots \\
% |%J| & Journal Name & |\ADS@journalname| &  \\
% |%D| & Publication Date & |\ADS@fulldate| &  \\
% |%L| & Last Page & |\ADS@lastpage| & generate from label \\
% |%T| & Title & |\ADS@title| & required, generate from |\title| \\
% |%C| & Abstract Copyright & |\ADS@copyright| & set to IA-UNAM \\
% |%U| & URL for electronic document & |\ADS@url| & set to PDF file \\
% |%d| & DOI for article & |\ADS@doi| & set by |\DOI| \\
% |%K| & Keywords & |\ADS@keywords| & set equal to |\@keywords| \\
% |%A| & Abstract Text & |\ADS@abstract| & set equal to |\@abstract| \\
% |%R| & References & |\ADS@references| & generate from the
% bibliography  \\ \hline
% \end{tabular}
% \end{table}
% A lot of these macros are set by other commands, as mentioned in the
% Table, but some need to be set  explicitly here. 
% 
% \begin{macro}{\ADS@doi}
% \changes{4.0}{2019/11/08}{Save DOI for ADS}
% The user/editor command \cs{DOI} is used to set this. 
%    \begin{macrocode}
\newcommand{\ADS@doi}{https://doi.org/10.22201/ia.01851101p.0000.00.00.00}
%    \end{macrocode}
% \end{macro}
% 
% \begin{macro}{\ADS@journalname}
% This
% does not contain any accents at the moment so as not to cause
% problems---this maybe should be changed though. 
%    \begin{macrocode}
\newcommand{\ADS@journalname}{%
  \if@proceedings\@conftitle\space (Eds. \@editors) \fi%
  Revista Mexicana de Astronom\'\i{}a y Astrof\'\i{}sica%
  \if@proceedings\space (Serie de Conferencias)\fi%
  \space Vol. \@volume, pp. \ADS@firstpage--\ADS@lastpage\space (\@year)%
  \space (http://www.astroscu.unam.mx/\~{}rmaa/)
  }
%    \end{macrocode}
% \end{macro}
% \begin{macro}{\ADS@keywords}
% \begin{macro}{\ADS@abstract}
% \begin{macro}{\ADS@title}
% |\ADS@abstract| is straightforward, since we already have the text
% saved as |\@abstract|. We could do the same with |\ADS@keywords| but
% to avoid unwanted formatting we explicitly set it in |\keywords|
% (p.~\pageref{keywords}). |\ADS@title| is just set to |\@rawtitle|,
% which we were already using for another purpose. 
%    \begin{macrocode}
\let\ADS@title=\@rawtitle
\newcommand{\ADS@keywords}{No keywords specified}
\newcommand{\ADS@abstract}{No abstract specified}
%    \end{macrocode}
% \end{macro}\end{macro}\end{macro}
% \begin{macro}{\ADS@authors}
% |\ADS@authors| wants to be in ``Other, A. N.'' format, but with
% semi-colons between authors. The most natural way to do this seems to
% be with |\indexauthor| (p.~\pageref{indexauthor}). We append each
% author to this macro, which means we must explicitly clear it for
% each new article (done in |\article|). .   
%    \begin{macrocode}
\newcommand{\ADS@authors}{}
%    \end{macrocode}
% \end{macro}
% \begin{macro}{\ADS@fulldate}
% \begin{macro}{\ADS@lastpage}
% \begin{macro}{\ADS@copyright}
% \begin{macro}{\ADS@website}
% \begin{macro}{\ADS@url}
% \begin{macro}{\SetWebsite}
% \begin{macro}{\SetMonthNumeric}
% Other sundry data, together with user commands for bits that can be
% set by the Editors. The |\ADS@url| is constructed from the website
% (currently |http://www.astrosmo.unam.mx/rmaa/|,
% the journal descriptor (e.g., |RMxAC..10/|), followed by the
% path to the PDF file (e.g., |PDF/RMxAC..10_HEW1.pdf|). 
% 
%    \begin{macrocode}
\newcommand{\ADS@website}{http://www.astrosmo.unam.mx/rmaa/}
\newcommand{\ADS@url}
{\ADS@website\ADS@journal\ADS@volume\ADS@number/PDF/\ADS@journal\ADS@volume\ADS@number\_\articlename.pdf}
\newcommand{\ADS@number}{}
\newcommand{\ADS@fulldate}{00/\ADS@year} 
\newcommand{\ADS@lastpage}{00} 
\newcommand{\ADS@copyright}{(C) \ADS@year; Instituto de Astronom\'\i{}a,
  Universidad Nacional Aut\'onoma de M\'exico} 
\newcommand{\SetNumber}[1]{\renewcommand{\ADS@number}{-#1}}
\newcommand{\SetWebsite}[1]{\renewcommand{\ADS@website}{#1}}
\newcommand{\SetMonthNumeric}[1]{\renewcommand{\ADS@fulldate}{#1/\ADS@year}}
%    \end{macrocode}
% \end{macro}\end{macro}\end{macro}\end{macro}\end{macro}\end{macro}\end{macro}
% \begin{macro}{\ADS@affil}
% \begin{macro}{\ADS@references}
% These two are somewhat more problematical. Affiliation currently
% only works for the first author (see \dots) pending redesign of the
% user interface. References are now done by wrapping the
% |thebibliography| environment to open another input stream on the
% same file and search through for the |\bibitems|. 
%    \begin{macrocode}
\newcommand{\ADS@references}{}
\newcommand{\ADS@affil}{}
%    \end{macrocode}
% \end{macro}\end{macro}
% 
% Define two macros that expand to the strings
% |\begin{thebibliography}| and |\end{thebibliography}|. Note that the
% idea for this and much else is inspired by the code for the docstrip
% package. Also define |\StartBibX| in case authors use the more
% standard |\begin{thebibliography}{}| syntax. 
%    \begin{macrocode}
\edef\EndBib{\expandafter\@gobble\string\\end\string{thebibliography\string}}
\edef\StartBib{\expandafter\@gobble\string\\begin\string{thebibliography\string}}
\edef\StartBibX{\expandafter\@gobble\string\\begin\string{thebibliography\string}\string{\string}}
%    \end{macrocode}
% Define two new logical variables. |\ifAnotherline| is true if we
% want to read another line from the input stream. It is set to false
% when we encounter the |\end{thebibliography}|. |\ifInBib| is set to
% true when we frind the |\begin{thebibliography}| and tells us to
% start saving the following lines. 
%    \begin{macrocode}
\newif\ifAnotherLine
\newif\ifInBib
%    \end{macrocode}
% \begin{macro}{makeOther}
% We need a macro to change the catcodes of lots of things to other. 
%    \begin{macrocode}
\def\makeOther#1{\catcode`#1=12\relax}
%    \end{macrocode}
% \end{macro}
% \begin{macro}{\RereadBiblio}
% Define a macro to do all the fiddly stuff. This is used at the
% beginning of |thebibliography| environment (Sec~\ref{sec:biblio}).
%    \begin{macrocode}
\newread\tf@biblio
\newcommand{\RereadBiblio}{%
%    \end{macrocode}
% Open the current file again for reading
%    \begin{macrocode}
  \begingroup
  \immediate\openin\tf@biblio=\articlename
%    \end{macrocode}
% Everything in a group to be safe. Initialize logicals. Black magic
% to change
% catcodes and avoid blank lines being interpreted as |\par|. 
%    \begin{macrocode}
  \AnotherLinetrue\InBibfalse
  \makeOther\ \makeOther\\\makeOther\$%
  \makeOther\#\makeOther\^\makeOther\^^K%
  \makeOther\_\makeOther\^^A\makeOther\%%
  \makeOther\~\makeOther\{\makeOther\}\makeOther\&%
  \endlinechar-1\relax
%    \end{macrocode}
% Start looping through the file. Read each line to the macro
% |\Line|. If we have got to EOF or end of biblio then set
% |\AnotherLinefalse|.
%    \begin{macrocode}
  \loop
  \read\tf@biblio to \Line
  \ifx\Line\EndBib\AnotherLinefalse\InBibfalse\fi
  \ifeof\tf@biblio \AnotherLinefalse\fi
%    \end{macrocode}
% If we are in the biblio and the line is not blank then append it to
% the |\ADS@references| macro together with a newline. 
%    \begin{macrocode}
  \ifInBib
  \ifx\Line\empty\else\xdef\ADS@references{\ADS@references\Line^^J}\fi
  \fi
%    \end{macrocode}
% Check if the line we just saw was the start of biblio, then repeat
% the loop if necessary. 
%    \begin{macrocode}
  \ifx\Line\StartBib\InBibtrue\fi
  \ifx\Line\StartBibX\InBibtrue\fi
  \ifAnotherLine\repeat
  \immediate\closein\tf@biblio
  \endgroup
  }
%    \end{macrocode}
% \end{macro}
% 
% \begin{macro}{\PadToFourWithDots}
% Utility function for left-padding a number to 4 characters with
% periods. Used in 
% constructing |\ADS@BibCode|. This is a very simple algorithm and
% \emph{requires} that the argument actually be a number---if it is
% not then dire consequences will result.  
%    \begin{macrocode}
\newcommand{\PadToFourWithDots}[1]{%
  \ifnum #1 > 999 #1%
  \else\ifnum #1 > 99 .#1%
  \else\ifnum #1 > 9 ..#1%
  \else ...#1%
  \fi\fi\fi}
%%  \ifnum \count0 > 999 #1%
%%  \else\ifnum \count0 > 99 .#1%
%%  \else\ifnum \count0 > 9 ..#1%
%%  \else ...#1%
%%  \fi\fi\fi}
%    \end{macrocode}
% \end{macro}
%
% \begin{macro}{\AuthorADS}
% \begin{macro}{\ADS@affiliations}
% New mechanism for setting the affiliations for ADS. The author labels (AA, AB, \dots) are generated using two nested counters.  We use the |\robustify| command from \texttt{etoolbox} so that accented characters in the affiliation addresses do not blow up.  The list of accents may need to be expanded. 
%    \begin{macrocode}
\newcounter{AUTHOR}
\newcounter{author}[AUTHOR]
\stepcounter{AUTHOR}
\newcommand\ADS@affiliations{}
\robustify{\'}\robustify{\i}\robustify{\~}\robustify{\`}
\def\AuthorADS#1#2{%
  \begingroup
  \def\ORCIDiD##1{<ID system="ORCID">##1</ID>}
  \def\Email##1{<EMAIL>##1</EMAIL>}
  \def\Affil##1{##1}
  \stepcounter{author}
  \ifthenelse{\value{author} > 26}{\stepcounter{AUTHOR}\stepcounter{author}}{}
  \ifthenelse{\equal{\ADS@affiliations}{\@rmaaEmpty}}{%
    \xdef\ADS@affiliations{\Alph{AUTHOR}\Alph{author}(#2)}
  }{%
    \xdef\ADS@affiliations{\ADS@affiliations, \Alph{AUTHOR}\Alph{author}(#2)}
  }%
  \endgroup
}
%    \end{macrocode}
% \end{macro}
% \end{macro}
% %
% \begin{macro}{\ADS@dumpdata}
% \changes{4.0}{2019/11/13}{Write all author affiliations for ADS}
% \begin{macro}{\ADS@write}
% \begin{macro}{\ADS@openfile}
% |\ADS@dumpdata| is the macro that actually does the job of writing out all the
% data to an auxiliary file, using |\ADS@write| to write each
% field. The file is set up with |\ADS@openfile|, which currently is
% done in the TOC of a book.  The new multi-author affiliations are used if |\ADS@affiliations| has been set, otherwise it falls back to the single-author version |\ADS@affil|. 
%    \begin{macrocode}
\newcommand{\ADSfield}[2]{\relax}
\newcommand{\ADSvfield}[2]{\relax}
%%\renewcommand{\ADSfield}[12]{\%#1 #2}
\newcommand{\ADS@dumpdata}
{ 
  \ADS@write{\string\begin{ADSarticle}}
  \ADS@WRITE{\string\ADSfield{R}{\ADS@BibCode}}
  \ADS@write{\string\ADSfield{A}{\ADS@authors}}
  \ifthenelse{\equal{\ADS@affiliations}{\@rmaaEmpty}}{%
    {\ADS@sanitize\ADS@write{\string\ADSfield{F}{AA(\ADS@affil)}}}
  }{%
    {\ADS@sanitize\ADS@write{\string\ADSfield{F}{\ADS@affiliations}}}
  }
  \ADS@write{\string\ADSfield{J}{\ADS@journalname}}
  \ADS@write{\string\ADSfield{D}{\ADS@fulldate}}
  \ADS@write{\string\ADSfield{L}{\ADS@lastpage}}
  {\ADS@sanitize\ADS@write{\string\ADSfield{T}{\ADS@title}}}
  \ADS@write{\string\ADSfield{C}{\ADS@copyright}}
  \ADS@write{\string\ADSfield{d}{\ADS@doi}}
  \ADS@write{\string\ADSfield{U}{\ADS@url}}
  \ADS@write{\string\ADSfield{K}{\ADS@keywords}}
  {\ADS@sanitize\ADS@write{\string\ADSfield{B}{\ADS@abstract}}}
  \ADS@write{\string\ADSfield{Z}{\ADS@references}}
  \ADS@write{\string\end{ADSarticle}}
  }
%%\newcommand{\ADS@write}[1]{\if@book\addtocontents{ads}{#1}\fi}
%%\newcommand{\ADS@openfile}{\@starttoc{ads}}
\newcommand{\ADS@write}[1]{\set@display@protect\immediate\write\tf@ads{#1}}
\newcommand{\ADS@WRITE}[1]{\immediate\write\tf@ads{#1}}
\newcommand{\ADS@openfile}{%
  \newwrite\tf@ads%
  \immediate\openout\tf@ads\jobname.ads\relax
  \ADS@WRITE{\ADS@preamble}}
\newcommand{\ADS@closefile}{\ADS@WRITE{\ADS@postamble}}
\def\ADS@preamble{%
  %%  \string\documentclass{article}
  %%  \string\begin{document}
  }
\def\ADS@postamble{%
  %%  \string\end{document}
  \string%\string% Local Variables:^J 
  \string%\string% mode: LaTeX^J
  \string%\string% mode: Font-lock^J
  \string%\string% End:^J
  }
%    \end{macrocode}
% \end{macro}
% \end{macro}
% \end{macro}
% 
% \subsection{Automatic calculation of the page range}
% \label{sec:page-range}
% Another issue that has cost me a lot of grief. Still, it finally
% seems to be sorted now. All hail Jeff Goldberg. 
%
% Cunning trick to get the page numbers from |\pageref|.  Put a label
% on the first and last page of the document using |\AtBeginDocument|
% and |\AtEndDocument|. .  This way won't work for the book option.
% In that case, we do it by defining |\@endofarticlehook| and
% |\@begofarticlehook| as parallels to |\AtEndDocument| and
% |\AtBeginDocument|.  Note that we have to defer adding
% |\@begofarticlehook| to |\AtBeginDocument| to make sure it comes
% after the redefinitions of |\ref| and |\label|.
%    \begin{macrocode}
\if@book\AtEndDocument{\ADS@closefile}\else
 %\AtBeginDocument{\@begofarticlehook}
\AtEndDocument{\@endofarticlehook}
\fi
%    \end{macrocode}
% \begin{macro}{\@begofarticlehook}
% Since we are getting more and more things that need to be done at
% the beginning of every article, frontsection, etc we define a sister
% hook to 
% |\@endofarticlehook|. This hook does the following: 
% \begin{enumerate}
% \item Labels the first page
% \item Sets |\ADS@firstpage|
% \item runs and then clears |\indexauthor@deferlist|
% \item resets the keyword counter
% \end{enumerate}
% It is executed in the following circumstances: 
% \begin{enumerate}
% \item At |\begin{document}| for standalone files
% \item At |\begin{document}| for each included file
% \item At beginning of front sections 
% \item At beginning of TOC
% \item At beginning of Author Index
% \item At beginning of abstracts
% \end{enumerate}
% For some reason, this is not always within the right group, so we
% use |\xdef| just to be sure. 
%    \begin{macrocode}
\newcommand{\@begofarticlehook}{%
  \if@frontmatter\frontlabel{RMAAFirstPage}\else\label{RMAAFirstPage}\fi%
  \xdef\ADS@firstpage{\@arabic{\c@page}}%
  \if@book\else\makeindex\fi % needed so default index is defined in stand-alone articles
%%  \indexauthor@deferlist
%%  \xdef\indexauthor@deferlist{}
  \@inpreamblefalse
  \setcounter{kwdcounter}{0}
  }
%    \end{macrocode}
% \end{macro}
% \begin{macro}{\@endofarticlehook}
% This is a hook to be run at the end of each article
% currently, it just labels the last page and writes out all the ADS
% data to an auxiliary file.  
% It should not produce any printed output. 
% It is modelled on the |\AtEndDocument| command, to which it reduces
% in the case of a single article. 
%    \begin{macrocode}
\newcommand{\@endofarticlehook}{%
  \if@debug{\ClassWarning{rmaa}{***Doing Hook***}}\fi
  \if@debug{\ClassWarning{rmaa}{ADSabstract is: \ADS@abstract}}\fi
  \clearpage\rmaa@lastpagelabel\if@book\ADS@dumpdata\fi
  \ifx\NotesList\@rmaaEmpty\else\OutputNotes\fi
  \if@debug{\ClassWarning{rmaa}{Notes list is: \NotesList}}\fi
}
%    \end{macrocode}
% \begin{macro}{\rmaa@lastpagelabel}
% New version of last page label command copied from Jeff Goldberg's
% lastpage.sty package - we can't just use |\label| here in case the
% output routine for the last page has already been called. Instead,
% we first explicitly clear the page, step back the page counter,
% explicitly write to the .aux file, then restep the page counter
% again.
%
% But this will conflict with the hyperref package, which does labels differently.  Compensate for that by using the five-argument form.
% 
% Also set the |\ADS@LastPage| macro to the appropriate number. 
%    \begin{macrocode}
\newcommand{\rmaa@lastpagelabel}{%
  \addtocounter{page}{-1}%
  \immediate\write\@auxout{\string
    \newlabel{\articlename:RMAALastPage}{{}{\thepage}{}{}{}}}%
  \edef\ADS@lastpage{\thepage}%
  \addtocounter{page}{1}}
%    \end{macrocode}
% \end{macro}
% \end{macro}

% \subsection{Deal with explicit change in number of columns}
% \label{sec:columns}
% Sometimes, when formatting a book, we need to switch between
% one-column articles and two-column articles. For this to work, we
% need to wrap the standard \LaTeX{} commands to set the |\@rmtwocol|
% flag appropriately and also reset the package default values of all
% scalable lengths. .  
%    \begin{macrocode}
\let\latexonecolumn=\onecolumn
\let\latextwocolumn=\twocolumn
\renewcommand{\onecolumn}{\@rmtwocolfalse\@package@default@lengths\latexonecolumn}
\renewcommand{\twocolumn}{\@rmtwocoltrue\@package@default@lengths\latextwocolumn}
%    \end{macrocode}
% 
% \subsection{The actual formatting of the title}
% \label{sec:title-format}

% Note, we have to redefine the |\maketitle| command to use the
% numbered footnotes. This means we needn't use the atbeginend stuff. 
%    \begin{macrocode}
\renewcommand{\maketitle}{\par
  \begingroup
%    \end{macrocode}
% Don't switch to marked footnotes. 
% This is consistent with existing style but I'm not sure I like it
%    \begin{macrocode}
  %%  \renewcommand\thefootnote{\fnsymbol{footnote}}%
%%  \def\@makefnmark{\hbox to\z@{$\m@th^{\@thefnmark}$\hss}}%
  \def\@makefnmark{\textsuperscript{\@thefnmark}\hspace{-0.2em}}%
  \long\def\@makefntext##1{\parindent 1em\noindent
    \hbox to1.8em{\hss$\m@th^{\@thefnmark}$}##1}%
  \if@rmtwocol
  { \if@debug{\ClassWarning{rmaa}{TWO COLUMN}}\fi
    \twocolumn[\@maketitle]%
    }
  \else
  { \if@debug{\ClassWarning{rmaa}{ONE COLUMN}}\fi
    \newpage
    \global\@topnum\z@   % Prevents figures from going at top of page.
    \if@manuscript \@MSmaketitle \else \@maketitle \fi
    }
  \fi
  \if@manuscript\else\thispagestyle{rmaafirstpage}\@thanks\fi
  \ifx\NotesList\@rmaaEmpty\else{%
    \NoteMark{1--\arabic{notecounter}}\vspace*{-2\baselineskip}%
  }\fi%
  \endgroup
%    \end{macrocode}
% Don't reset footnote counter (not sure whether this is good thing
% or not).
%    \begin{macrocode}
  %%  \setcounter{footnote}{0}%
  \let\thanks\relax\let\affil\relax\let\altaffil\relax
  \let\maketitle\relax\let\@maketitle\relax
  \gdef\@thanks{}\gdef\@author{}\gdef\@title{}}
%    \end{macrocode}

% \begin{macro}{\@kwhwidth}
% \begin{macro}{\@kwtwidth}
% Redefine \cs{@maketitle} to include the abstract, resumen and
% keywords. Note that in the manuscripty style we use
% \cs{@MSmaketitle} instead.  
% \changes{3.16}{2003/04/07}{
%  Added extra vertical space between authors and accepted date.}
% \changes{3.26}{2011/08/15}{Add extra label to review articles}
%    \begin{macrocode}
\newlength{\@kwhwidth}\newlength{\@kwtwidth}
\newcommand\reviewname{Review}
\renewcommand{\@maketitle}{%
  \newpage
  \null
  \if@reviewarticle
  \smallskip
  {\large\bfseries\itshape\reviewname}%
  \bigskip
  \else
  \vskip\@before@title@skip%
  \fi
  \begin{center}%
    {\@title \par}%
    \vskip\@before@author@skip%
    \parbox{\rmaa@AuthorWidthFrac\textwidth}{\centering
      \lineskip\@between@author@skip%
      \@author
      \par}
    %% this should be replaced with `Submitted to...' for preprints
    \if@printdates
    \vskip\@before@dates@skip%
    \if@preprint
    \rmaa@ReceivedStyle{Submitted to RMxAA \@received}
    \else
    \rmaa@ReceivedStyle{Received \@received; accepted \@accepted}
    \fi
    \fi
    \if@manuscript
    \rmaa@ReceivedStyle{Submitted to RMxAA \@received}
    \fi
    \if@guide
    \rmaa@ReceivedStyle{Version \rmaa@version, \rmaa@date}
    \fi
    \vskip\@before@abstract@skip%
    \if@abstractpresent
    \parbox{\rmaa@AbstractWidthFrac\textwidth}{%
      \setlength{\parindent}{0.9cm}%
      \centerline{\rmaa@AbstractHeaderStyle{Abstract}}\par%
      \vskip\@after@abstracthead@skip%
      \if@newstyle\noindent\fi\@abstract
      \vskip\@between@abstract@skip%
      \centerline{\rmaa@AbstractHeaderStyle{Resumen}}\par%
      \vskip\@after@abstracthead@skip%
      \if@newstyle\noindent\fi\@resumen
      \if@setresumen\else\@abstract\fi % append english abstract if
                                       % spanish one not set
      \if@guide\else            %do not use keywords for guide
      \vskip\@before@kwd@skip%
      % Format the Key Words with hanging indentation:
      % They are put in a \parbox with width equal to the line width
      % minus the width of the header. 
      { % Widths of keyword
        \settowidth{\@kwhwidth}{\@keywordheader~}   % header and
        \setlength{\@kwtwidth}{\linewidth}          % keyword text
        \addtolength{\@kwtwidth}{-\@kwhwidth}       %
        \noindent\@keywordheader~\parbox[t]{\@kwtwidth}{%
          \rmaa@KeywordStyle{\boldmath\uppercase{\@keywords}}}}%
      \fi
      }
      \fi
  \end{center}%
  \vspace*{\afterkwdskip}
  \par
  }
\newlength{\afterkwdskip}\setlength{\afterkwdskip}{\z@}
%    \end{macrocode}
% \end{macro}
% \end{macro}
%
%
% \begin{macro}{\reviewarticle}
% \begin{macro}{\reviewtoc}
% \begin{macro}{\reviewlof}
%   User command to specify a review article. Must be used in
%   preamble. Main effect is to place the text ``Review'' before the
%   article title.  Also add command for setting a table of contents
%   and a list of figures at the beginning of the article.
%    \begin{macrocode}
\newcommand{\reviewarticle}{\@reviewarticletrue}
\newcommand\reviewtoc{%
  \section*{CONTENTS}%
  \@starttoc{rtc}%
}
\newcommand\reviewlof{%
  \section*{LIST OF FIGURES}%
  \@starttoc{lof}%
}
%    \end{macrocode}
% \end{macro}
% \end{macro}
% \end{macro}
% \begin{macro}{\l@section}
%    \begin{macrocode}
\renewcommand*\l@section[2]{%
  \ifnum \c@tocdepth >\z@
    \addpenalty\@secpenalty
    \addvspace{0.8em \@plus\p@}%
    \setlength\@tempdima{1.5em}%
    \begingroup
      \parindent \z@ \rightskip \@pnumwidth
      \parfillskip -\@pnumwidth
      \leavevmode %\bfseries
      \advance\leftskip\@tempdima
      \hskip -\leftskip
      #1\nobreak\hfil \nobreak\hb@xt@\@pnumwidth{\hss #2}\par
    \endgroup
  \fi}
%    \end{macrocode}
% \end{macro}
%
% \begin{macro}{\@MSmaketitle}
% A special version of |\@maketitle| for the \texttt{manuscript}
% option. This can deal with a multi-page abstract/resumen. Not very
% well, mind. 
%    \begin{macrocode}
\newcommand{\@MSmaketitle}{%
  \if@debug\ClassInfo{rmaa}{Using special MS maketitle command}\fi
  \newpage
  {
    \thispagestyle{rmaafirstpage}
%     \setlength{\parindent}{0pt}
    \null\vskip 2ex
    \begin{center}
      \@title\par
      \@author\par
      \rmaa@ReceivedStyle{%
        \ifx\@received\@rmaaEmpty
          Draft version: \@date
        \else
          Received: \@received
          \ifx\@accepted\@rmaaEmpty
          \else
            ; Accepted: \@accepted
          \fi
        \fi}\par
    \end{center}
    \@thanks
    \centerline{\rmaa@AbstractHeaderStyle{Resumen}}
    \noindent\@resumen\par\vskip 2ex
    \centerline{\rmaa@AbstractHeaderStyle{Abstract}}
    \noindent\@abstract\par\vskip 2ex
    \noindent\@keywordheader~\rmaa@KeywordStyle{\@keywords}\par
%     \@keywordheader~\rmaa@KeywordStyle{\uppercase{\@keywords}}\par
  }%
  }
%    \end{macrocode}
% \end{macro}

% \subsection{The sectioning commands}
% \label{sec:section-commands}

% \begin{macro}{\section}
% \begin{macro}{\subsection}
% \begin{macro}{\subsubsection}
% \begin{macro}{\paragraph}
% \begin{macro}{\subparagraph}
% Redefine all the standard commands to use these lengths. 
%    \begin{macrocode}
\renewcommand\section{\@startsection {section}{1}{\z@}%
                                   {\abovesectionskip}%
                                   {\belowsectionskip}%
                                   %% how do we get upper case here
                                   {\centering\reset@font\normalsize}}
\renewcommand\subsection{\@startsection{subsection}{2}{\z@}%
                                     {\abovesectionskip}%
                                     {\belowsectionskip}%
                                     {\centering\reset@font\normalsize\itshape}}
\renewcommand\subsubsection{\@startsection{subsubsection}{3}{\z@}%
                                     {\abovesectionskip}%
                                     {\belowsectionskip}%
                                     {\centering\reset@font\normalsize\itshape}}
\renewcommand\paragraph{\@startsection{paragraph}{4}{\z@}%
                                    {\aboveparaskip}%
                                    {\belowparaskip}%
                                    {\reset@font\normalsize\bfseries}}
\renewcommand\subparagraph{\@startsection{subparagraph}{5}{\parindent}%
                                      {\aboveparaskip}%
                                      {\belowparaskip}%
                                      {\reset@font\normalsize\bfseries}}
%    \end{macrocode}
% \end{macro}
% \end{macro}
% \end{macro}
% \end{macro}
% \end{macro}

% \begin{macro}{secnumdepth}
% Make sure that we number the subsubsections---this is done by
% default in the |article| class but not the |book| class, which is
% why we must set it explicitly here.
%    \begin{macrocode}
\setcounter{secnumdepth}{3}
%    \end{macrocode}
% \end{macro}
% \changes{2.2beta10}{2001/11/07}{Fixed \cs{secnumdepth} for \cs{book}
% style.}

% \begin{macro}{\thesubsection}
% \begin{macro}{\thesubsubsection}
% For subsection and subsubsection, we want the numbers to be roman
% but the title to be italic. This seems the easiest way to go about it.
%    \begin{macrocode}
\renewcommand{\thesubsection}{{\upshape\thesection.\@arabic\c@subsection}}
\renewcommand{\thesubsubsection}{\thesubsection{\upshape.\@arabic\c@subsubsection}}
%    \end{macrocode}
% \end{macro}
% \end{macro}

% \begin{macro}{\@seccntformat}
% Add ``.'' to the number here so we can remove it from |\thesection|,
% etc. Reduce gap between number and title WJH 27 Jan 1995
%    \begin{macrocode}
\def\@seccntformat#1{\csname the#1\endcsname .\hskip 0.4em}
%    \end{macrocode}
% \end{macro}

% \begin{macro}{\@sect}
% I didn't really want to mess with the |\@sect| command, but it was
% the only way I could manage to force upper case. Other than the
% commented parts, it is identical to the standard version. 
%    \begin{macrocode}
\def\@sect#1#2#3#4#5#6[#7]#8{
  \ifnum #2>\c@secnumdepth
     \let\@svsec\@empty
  \else
     \refstepcounter{#1}%
     \let\@@protect\protect
     \def\protect{\noexpand\protect\noexpand}%
     \edef\@svsec{\@seccntformat{#1}}%
     \let\protect\@@protect
  \fi
%    \end{macrocode}
% This is the bit I changed (27 Jan 1995). Use upper case for
% |\section|, use mixed case for the rest. 
%    \begin{macrocode}
  \ifnum #2=1                   
     \def\sec@upcase##1{\uppercase{##1}}%
     \def\sec@shape{\upshape}%
  \else                         
     \def\sec@upcase##1{\relax{##1}}%
     \def\sec@shape{\itshape}%
  \fi
  \@tempskipa #5\relax
  \ifdim \@tempskipa>\z@
     \begingroup #6\relax
          \@hangfrom{\hskip #3\relax\@svsec}%
                    {\interlinepenalty \@M \sec@upcase{#8}\par}%
     \endgroup
     \csname #1mark\endcsname{#7}%
%    \end{macrocode}
% Only write TOC entries if doing the Guide style. 
%    \begin{macrocode}
     \if@guide 
        \addcontentsline{toc}{#1}{%
          \ifnum #2>\c@secnumdepth \else
             \protect\numberline{\csname the#1\endcsname}%
          \fi
          #7}%
     \fi
%    \end{macrocode}
% Also allow per-article TOC for review papers
%    \begin{macrocode}
     \if@reviewarticle
        \ifnum #2>2 \else
           \addcontentsline{rtc}{#1}{%
             \protect\numberline{\csname the#1\endcsname}%
             \sec@upcase{\sec@shape #7}}%
        \fi
     \fi
  \else
     \def\@svsechd{#6\hskip #3\relax  %% \relax added 2 May 90
       \@svsec #8\csname #1mark\endcsname{#7}%
       \if@guide
          \addcontentsline{toc}{#1}{%
            \ifnum #2>\c@secnumdepth \else
               \protect\numberline{\csname the#1\endcsname}%
            \fi
            #7}%
       \fi
       }%
  \fi
  \@xsect{#5}}
%    \end{macrocode}
% \end{macro}


% \subsection{Appendices}
% \label{sec:appendices}
% Make sure that appendices in individual articles work properly. We
% copy the definition from the article class (the book one doesn't do
% what we want).  
% \begin{environment}{appendices}
%   This uses an environment |appendices| inside of which each
%   |\section| will become a separate appendix.  Note that the
%   appendices must come before the references (otherwise the full
%   addresses mechanism will not work properly). We make the equation
%   numbering be subordinate to the appendix (equations A12, B6,
%   etc). Finally, we have a medskip followed by APPENDICES to
%   introduce them. \changes{3.17}{2002/04/08}{ Added appendices environment
%   that does the right thing.}
%    \begin{macrocode}
\newenvironment{appendices}{\par
  \begingroup
  \setcounter{section}{0}%
  \setcounter{subsection}{0}%
  \renewcommand\thesection{\@Alph\c@section}
  \renewcommand\theequation{\thesection\arabic{equation}}
  \medskip
  \section*{APPENDICES}}
{\endgroup}
%    \end{macrocode}
% \end{environment}
% \begin{environment}{appendix}
%   If there is only one appendix, then use |appendix| environment
%   instead.  \changes{4.0}{2019/11/12}{ Added
%   appendix environment for case where there is only one section.}
%    \begin{macrocode}
\renewenvironment{appendix}{\par
  \begingroup
  \setcounter{section}{0}%
  \setcounter{subsection}{0}%
  \renewcommand\thesection{\@Alph\c@section}
  \renewcommand\theequation{\thesection\arabic{equation}}
  \medskip
  \section*{APPENDIX}}
{\endgroup}
%    \end{macrocode}
% \end{environment}

% \subsection{Floating thingies}
% \label{sec:floats}   

%% 
%    \begin{macrocode}
\renewcommand{\thetable}{\@arabic\c@table}
\renewcommand{\thefigure}{\@arabic\c@figure}
\renewcommand{\figurename}{Fig.}
\renewcommand{\tablename}{TABLE}
%\renewcommand{\fnum@table}{\MakeUpperCase{\tablename}~\thetable}
%    \end{macrocode}
% For some reason, I have to write `TABLE' explicitly here
% If I use |\tablename|, it gets converted to mixed case---God knows
% why. 
% |\figurename| is also being ignored. 
%    \begin{macrocode}
\renewcommand{\fnum@table}{TABLE~\thetable}
\renewcommand{\fnum@figure}{Fig.~\thefigure}
%    \end{macrocode}

% Do captions in 9 pt face in the main journal style and in new-style
% anything. 
%    \begin{macrocode}
\long\def\@makefigcap#1#2{%
  \if@proceedings\else\small\fi
  \if@newstyle\small\fi 
  \vskip\abovecaptionskip
  \sbox\@tempboxa{#1. #2}%
  \ifdim \wd\@tempboxa >\hsize
    #1. #2 \par
  \else
    \global \@minipagefalse
    \hb@xt@\hsize{\hfil\box\@tempboxa\hfil}%
  \fi
  \vskip\belowcaptionskip}
%    \end{macrocode}

% Make sure table caption is upper case. 
%    \begin{macrocode}
\long\def\@maketabcap#1#2{%
  \hbox to\hsize{\hfil #1 \hfil}
  \vskip 6pt 
  \centering\MakeUpperCase{#2}\par
  \vskip 6pt}
%    \end{macrocode}

%    \begin{macrocode}
\newlength{\@rottabcapwidth}
\long\def\@makerottabcap#1#2{%
  \setbox\@tempboxa\hbox{\MakeUpperCase{#2}}%
  \ifdim\wd\@tempboxa >.8\vsize
     \setlength\@rottabcapwidth{.8\vsize}
  \else
     \setlength\@rottabcapwidth{\wd\@tempboxa}
  \fi
  \rotatebox{90}{%
    \begin{minipage}{\vsize}\centering
      \begin{minipage}{\@rottabcapwidth}\centering
        #1\\ \vskip 6pt \MakeUpperCase{#2}
      \end{minipage}%
    \end{minipage}
    }%
  \hskip 6pt
  }

\newenvironment{Sideways}
    {\begin{sideways}\begin{minipage}{\vsize}\centering}
    {\end{minipage}\end{sideways}}

\newenvironment{Sidewaystable}
{%
  \let\oldcaption=\caption%
  \def\caption{
    \def\@captype{table}\renewcommand{\@makecaption}{\@maketabcap}
    \oldcaption}
  \begin{sidewaystable}\centering}
  {\end{sidewaystable}\let\caption=\oldcaption}
\newenvironment{Sidewaystable*}
{%
  \let\oldcaption=\caption%
  \def\caption{
    \def\@captype{table}\renewcommand{\@makecaption}{\@maketabcap}
    \oldcaption}
  \begin{sidewaystable*}\centering}
  {\end{sidewaystable*}\let\caption=\oldcaption}
%    \end{macrocode}

% Not sure why we need to stick these lines in. 
%    \begin{macrocode}
\renewenvironment{figure}
                 {\renewcommand{\@makecaption}{\@makefigcap}\@float{figure}}
                 {\end@float}
\renewenvironment{figure*}
                 {\renewcommand{\@makecaption}{\@makefigcap}\@dblfloat{figure}}
                 {\end@dblfloat}
\renewenvironment{table}
                 {\renewcommand{\@makecaption}{\@maketabcap}\@float{table}}
                 {\end@float}
\renewenvironment{table*}
                 {\renewcommand{\@makecaption}{\@maketabcap}\@dblfloat{table}}
                 {\end@dblfloat}
%    \end{macrocode}

% Redefine |\@caption| so that we don't bother with the list of
% figures. This saves us having to worry about fragile commands in the
% captions. 
%    \begin{macrocode}
\long\def\@caption#1[#2]#3{%
  \par
%%  \addcontentsline{\csname ext@#1\endcsname}{#1}%
%%    {\protect\numberline{\csname the#1\endcsname}{\ignorespaces #2}}%
  \begingroup
    \@parboxrestore
    \if@minipage
      \@setminipage
    \fi
    \normalsize
    \@makecaption{\csname fnum@#1\endcsname}{\ignorespaces #3}\par
  \endgroup}
%    \end{macrocode}

% Make tables not quite so cramped. Allow |\hline\hline| to give a
% poor man's thick rule. 
%    \begin{macrocode}
\renewcommand{\arraystretch}{1.15}
\setlength{\doublerulesep}{0.4pt}
%    \end{macrocode}

% \begin{macro}{\tablecols}
% \begin{macro}{\@tablecols}
% \begin{macro}{\tabcaptioncontinued}
% |\tablecols| is a user command for setting the number of columns. It
% is used by commands that try and span the entire table (I think
% there must be a better way of acheiving
% this). |\tabcaptioncontinued| is 
% a user command for setting the caption of a continued table on the
% next page. 
%    \begin{macrocode}
\newcommand{\tablecols}[1]{\renewcommand{\@tablecols}{#1}}
\newcommand{\@tablecols}{1}
\newcommand{\tabcaptioncontinued}{%
  \multicolumn{\@tablecols}{c}{\addtocounter{table}{-1}TABLE~\thetable\ (CONTINUED)} \\ \\
  }
%    \end{macrocode}
% \end{macro}
% \end{macro}
% \end{macro}

% \subsubsection{Footnotes in tables}
% \label{sec:table-footnotes}
% \begin{macro}{\tablenotemark}
% \begin{macro}{\tablenote}
% \begin{macro}{\tablenoteskip}
% Old user commands |\tablenotemark{MARK}|,
% |\tablenote{MARK}{TEXT}|---now deprecated. 
%    \begin{macrocode}
\newcommand\tablenotemark[1]{$^\mathrm{#1}$} 
\newcommand\tablenote[2]{\if@proceedings\else\small\fi ${}^\mathrm{#1}$ #2} 
\newcommand\tablenoteskip{\\[0.5\baselineskip]} 
%    \end{macrocode}
% \end{macro}
% \end{macro}
% \end{macro}

% \begin{macro}{\tabnotemark}
% \begin{macro}{\tabnotetext}
% \begin{macro}{\tabnotetextTWO}
% Alternative way of doing footnotes in tables. 
%    \begin{macrocode}
\newcommand{\tabnotemark}[1]{\textsuperscript{\lowercase{#1}}}
%\newcommand{\tabnotetextTWO}[4]{\tabnotetext{#1}{#2}\tabnotetext{#3}{#4}}
\newcommand{\tabnotetextTWO}[4]{%
  \multicolumn{\@tablecols}{l}{%
    \parbox{0.5\tabnotewidth}{\small\textsuperscript{#1}#2}%
    \parbox{0.5\tabnotewidth}{\small\textsuperscript{#3}#4}%
    }%
  \\}
\newcommand{\tabnotetext}[2]{%
  \multicolumn{\@tablecols}{@{}p{\tabnotewidth}@{}}{\small\textsuperscript{#1}#2}\\}
\newlength{\tabnotewidth}
%    \end{macrocode}
% \end{macro}
% \end{macro}
% \end{macro}


% \begin{macro}{changemargin}
%    \begin{macrocode}
% Command to allow margins to be changed on the fly. Useful to allow
% tables and Figures to protrude into the margins. 
%
% \changes{3.24}{2007/09/14}{Added \texttt{changemargin} environment
% (copied from TeX FAQ).}
\newenvironment{changemargin}[2]{% 
  % From TeX FAQ http://www.tex.ac.uk/cgi-bin/texfaq2html?label=chngmargonfly
  \begin{list}{}{%
    \setlength{\topsep}{0pt}%
    \setlength{\leftmargin}{#1}%
    \setlength{\rightmargin}{#2}%
    \setlength{\listparindent}{\parindent}%
    \setlength{\itemindent}{\parindent}%
    \setlength{\parsep}{\parskip}%
  }%
  \item[]}{\end{list}}
%    \end{macrocode}
% \end{macro}

% \subsection{The bibliography and references}
% \label{sec:biblio}
% \begin{macro}{\beforebiblioskip}
% \begin{macro}{\afterbiblioskip}
% I'm not entirely sure why I introduced these two commands, since
% directly putting |\vspace| before and after
% |\begin{thebibliography}| seems to work well enough. 
%    \begin{macrocode}
\newlength{\beforebiblioskip}
\setlength{\beforebiblioskip}{0pt}
\newlength{\afterbiblioskip}
\setlength{\afterbiblioskip}{0pt}
%    \end{macrocode}
% \end{macro}
% \end{macro}
% \begin{macro}{thebibliography}
% The environment itself is done as starred section with extra
% vertical space after heading. This space is currently zero. 
% In the book style we do the cunning trick with rereading the file in
% order to get the references into the |\ADS@references| macro.
% 
% \changes{3.05}{2003/01/02}{Removed the \cs{beforebiblioskip} and
% \cs{afterbiblioskip} commands, since they cause unwanted vertical
% space in some instances.}  
%    \begin{macrocode}
\renewenvironment{thebibliography}{%
  \if@book\RereadBiblio\fi
%%  \vspace*{\beforebiblioskip}
  \section*{REFERENCES} \frenchspacing
%%  \vspace*{\afterbiblioskip}
  \small
%    \end{macrocode} 
% The references themselves are set as a list in order 
% to get the hanging indentation. They are done in a smaller font and
% with french spacing so we needn't worry about the distinction
% between sentence-ending periods and abbreviations. 
%    \begin{macrocode} 
  \list{\null}{\leftmargin 1.5em\labelwidth\z@\labelsep\z@\itemindent -1.5em
    \itemsep\z@\parsep\z@\usecounter{enumi}}
  \sloppy\clubpenalty4000\widowpenalty4000
  \sfcode`\.=1000\relax\@biblisttrue}
{\endlist\@biblistfalse\normalsize\@outputaddressesifnecessary}
%    \end{macrocode}
% \end{macro}
% \subsection{Uniquification of cross-references}
% \label{sec:uniquification}
% 
% \begin{macro}{\articlename}
% In order that the |\label{}| command in individual articles be
% guaranteed to give a unique anchor, we append the article name (held
% in the |\articlename| command) to each key. This requires
% redefinition of the following commands: |\cite|, |\nocite|, |\ref|,
% |\pageref|, |\label|. 
% \begin{macro}{\oldcite}
% \begin{macro}{\oldnocite}
% \begin{macro}{\oldref}
% \begin{macro}{\oldpageref}
% \begin{macro}{\oldlabel}
% The original versions of these commands are
% saved as |\oldcite|, |\oldnocite|, etc. 
%
% The uniquifying of cross-references is now always done, so 
% we don't get multiply-defined labels in the single article case
%    \begin{macrocode}
\newcommand{\articlename}{anon}
\AtBeginDocument{
  \let\oldcite=\cite 
  \let\oldnocite=\nocite 
  \let\oldref=\ref 
  \let\oldpageref=\pageref
  \let\oldlabel=\label
}
%    \end{macrocode}
% \end{macro}
% \end{macro}
% \end{macro}
% \end{macro}
% \end{macro}
% \end{macro}
%
% Uniquify all cross-references by prepending the article name.  This
% now has to go into |\AtBeginDocument| so it is working on the new
% versions that hyperref.sty has set up.  Then we need to put the
% |\@begofarticlehook| afterwards, so that the first page label is using
% the new definitions.  If amsmath is loaded, then we also need to patch
% the special label command that they use. \changes{4.0}{2019/11/18}{Fix labels for amsmath display equations, in case that package has been loaded.}
%    \begin{macrocode}
\AtBeginDocument{
  \renewcommand{\ref}[1]{\oldref{\articlename:#1}%
    \if@debug{\ClassInfo{rmaa}{Reference to \articlename:#1}}\fi}
  \renewcommand{\pageref}[1]{\oldpageref{\articlename:#1}%
    \if@debug{\ClassInfo{rmaa}{Page reference to \articlename:#1}}\fi}
  \renewcommand{\label}[1]{\oldlabel{\articlename:#1}%
    \if@debug{\ClassInfo{rmaa}{Label for \articlename:#1}}\fi}%
  \ifdef{\label@in@display}{%
    \let\oldlabel@in@display=\label@in@display
    \renewcommand{\label@in@display}[1]{\oldlabel@in@display{\articlename:#1}%
      \if@debug{\ClassInfo{rmaa}{AMSTeX Label for \articlename:#1}}\fi}%
  }{}
  \@begofarticlehook
}%
%    \end{macrocode}
% Citations are more complicated since they can contain a list of
% keys. These are really only here for use with the rmaacite package,
% which is utterly deprecated (02 Jan 2003) now that I have the natbib
% code working. The command \cs{cite} should \emph{not} be used with
% natbib: use \cs{citep}, \cs{citet}, etc.\@ instead. 
%    \begin{macrocode}
\renewcommand{\nocite}[1]{\@bsphack
  \@for\@citeb:=#1\do{%
    \edef\@citeb{\expandafter\@firstofone\@citeb:\articlename}%
    \if@filesw\immediate\write\@auxout{\string\citation{\@citeb}}\fi
    \@ifundefined{b@\@citeb:\articlename}{\G@refundefinedtrue
      \@latex@warning{Citation `\@citeb' undefined}}{}}%
  \@esphack}
\renewcommand{\cite}[1]{\oldcite{#1:\articlename}}
%   \end{macrocode}


% Horrible fiddling with the bibliography citation mechanism. 
% This has to be globally available in case rmaacite puts it in the
% \texttt{.aux} file. 
%    \begin{macrocode}
\def\lbibcite#1#2#3{\global\@namedef{b@#1}{#2}
                    \global\@namedef{by@#1}{#3}}
%    \end{macrocode}

% \begin{macro}{\oldbibitem}
% \begin{macro}{\@rmaaBibArg}
% \begin{macro}{\@rmaaSpace}
% \begin{macro}{rmaabib}
% Of course, |\bibitem| needs to get the uniquification treatment as
% well. While we are at it, we redefine it to use consecutive integers
% for the key if the user leaves it blank (using the counter
% |rmaabib|). This allows one to put 
% |\bibitem{}| for all your references without getting annoying
% messages about multiply-defined labels.
% 
% We also take the opportunity to insert the optional argument to
% \cs{bibitem} if it is missing in the form Author(DATE) so that natbib
% doesn't try to do numerical citations. Natbib already provides a
% hook \cs{@extra@b@citeb} to help uniquify the citation key, so we
% just plug into that. However, it puts it \emph{after} the original
% key, not before as I had had it. Easiest solution is to change the
% way that I do it. 02 Jan 2003. Note that the \cs{@extra@b@citeb}
% does not get written to the .aux  file in the \cs{citation} command
% but it does work. 
% \changes{3.14}{06 Feb 2003}{Reinstalled old definition of \cs{bibitem}
%  for the case where the nonatbib option is requested, so that the
%  rmaacite package will work. This means that you can't use rmaacite
%  and natbib in the same volume.}
%    \begin{macrocode}
\if@natbib\def\@extra@b@citeb{:\articlename}\fi
%%\else
\let\oldbibitem\bibitem
\newcommand{\@rmaaBibArg}{}
\newcommand{\@rmaaSpace}{ }
\newcounter{rmaabib}
\if@natbib
\renewcommand{\bibitem}[2][ABC(2003)]{%
  \renewcommand{\@rmaaBibArg}{#2}
  \ifx\@rmaaBibArg\@rmaaEmpty
     \addtocounter{rmaabib}{1}
     \oldbibitem[#1]{\thermaabib:\articlename}
  \else\ifx\@rmaaBibArg\@rmaaSpace
     \addtocounter{rmaabib}{1}
     \oldbibitem[#1]{\thermaabib:\articlename}
  \else 
     \oldbibitem[#1]{#2:\articlename}
  \fi\fi
  \renewcommand{\@rmaaBibArg}{}
  }
\else
\renewcommand{\bibitem}[1]{%
  \renewcommand{\@rmaaBibArg}{#1}
  \ifx\@rmaaBibArg\@rmaaEmpty
     \addtocounter{rmaabib}{1}
     \oldbibitem{\articlename:\thermaabib}
  \else\ifx\@rmaaBibArg\@rmaaSpace
     \addtocounter{rmaabib}{1}
     \oldbibitem{\articlename:\thermaabib}
  \else
     \oldbibitem{\articlename:#1}
  \fi\fi
  \renewcommand{\@rmaaBibArg}{}
  }
\fi
%%\fi
%    \end{macrocode}
% \end{macro}
% \end{macro}
% \end{macro}
% \end{macro}

% \begin{macro}{oldpcite}
% \begin{macro}{oldscite}
% Then, we need to repeat the above scheme, but for the |rmaacite|
% definition. 
%    \begin{macrocode}
\newcommand{\@redefinebibitem}{%
  \let\oldbibitem\bibitem
  \let\oldcite\cite
  \let\oldpcite\pcite
  \let\oldscite\scite
  \def\bibitem[##1<##2>]##3{%
    \renewcommand{\@rmaaBibArg}{##3}
    \ifx\@rmaaBibArg\@rmaaEmpty
       \addtocounter{rmaabib}{1}
       \oldbibitem[##1<##2>]{\thermaabib:\articlename}
    \else\ifx\@rmaaBibArg\@rmaaSpace
       \addtocounter{rmaabib}{1}
       \oldbibitem[##1<##2>]{\thermaabib:\articlename}
    \else 
       \oldbibitem[##1<##2>]{##3:\articlename}
    % trust the author here
%%       \oldbibitem[##1<##2>]{##3}
    \fi\fi
    }
  }
%    \end{macrocode}
% \end{macro}
% \end{macro}


% \subsection{HTML Table of Contents}
% \label{sec:hyper-TOC}
% \begin{macro}{\HyperTOCEntry}
% \begin{macro}{\HyperTOCpart}
% \begin{macro}{\HyperTOCmeeting}
% \begin{macro}{\HyperTOCeditors}
% \begin{macro}{\HyperTOCpreamble}
% These commands are used to write information to the \texttt{.htoc}
% file, which is used in generating a Web Page of the Table of
% Contents. Since they only ever appear preceded by |\protect| in the
% argument to an |\addtocontents| command,  they should never need to
% be expanded by the macros. However, we define them here just in case
% \dots
%    \begin{macrocode}
\newcommand{\HyperTOCEntry}[3]{\relax}
\newcommand{\HyperTOCpart}[1]{\relax}
\newcommand{\HyperTOCmeeting}[3]{\relax}
\newcommand{\HyperTOCeditors}[1]{\relax}
\newcommand{\HyperTOCpreamble}{\relax}
%    \end{macrocode}
% \end{macro}
% \end{macro}
% \end{macro}
% \end{macro}
% \end{macro}
%

% \subsection{Formatting the poster abstracts} 
% These can now be used without using the book option.
%    \begin{macrocode}

\newcommand{\@makeminititle}{%
  \newpage
  \null 
  \vskip 3em%
  \begin{center}%
    {\@title \par}%
  \end{center}
  \vskip 3em%
  }

\newif\if@abstractsads
\@abstractsadsfalse
\newcommand\abstractsADS{\@abstractsadstrue}

\if@book\RequirePackage{multicol}\fi
\newenvironment{abstracts}[1][Abstracts of Contributed Papers]{%
  %% First, make sure we start on an odd page (06 Jun 2000)
  \if@book\clearpage\thispagestyle{empty}\cleardoublepage\fi
  %%
  \renewcommand{\articlename}{ABSTRACTS}
  \title{#1}\shortauthor{Abstracts}\shorttitle{Abstracts}
  \listofauthors{}
  \if@book
  \addcontentsline{toc}{chapter}{\protect\@fiddleauthors{\@TOCtitle}{\rmaaTOCauthors}}%
  \addtocontents{htc}{\protect\HyperTOCEntry{\@rawtitle}{\@rawauthors}{\articlename}}
  \fi
  \pagestyle{rmaa} % this now seems necessary (06 Jun 2000)
  \thispagestyle{rmaafirstpage}
%%  \twocolumn[\@makeminititle]
  \renewcommand{\abstract}[1]{%
    \renewcommand\ADS@abstract{##1}%
    \renewcommand{\@abstract}{##1}\@writeabstract%
  }
  \renewcommand{\title}[2][]{%
    \gdef\ADS@authors{}%
    \gdef\ADS@affil{}%
    \gdef\ADS@references{}%
    \gdef\ADS@keywords{}%
    \renewcommand{\@rmaaArg}{##1}
    \ifx\@rmaaArg\@rmaaEmpty
       \renewcommand{\@TOCaftertitle}{}
    \else
       \renewcommand{\@TOCaftertitle}{ (\@rmaaArg)}
    \fi
    \renewcommand{\@title}{\rmaa@TOCTitleStyle{##2}}%
    \renewcommand{\@TOCtitle}{\rmaa@TOCTitleStyle{##2}\@TOCaftertitle}%
    \renewcommand{\rmaaIndexTitle}{\rmaa@IndexTitleStyle{##2\@TOCaftertitle.}}%
    \renewcommand\ADS@title{##2}%
    \xdef\ADS@firstpage{\@arabic{\c@page}}%
    }
  \renewcommand{\author}[1]{\let\@thanks\@empty\renewcommand\@author{##1} }
  \renewcommand{\footnotetext}[2][XX]{\par\noindent\textsuperscript{##1}~##2}
  \renewcommand{\footnoterule}{\kern-3\p@ \hrule \@width 3em \kern 2.6\p@}
  \if@book\@begofarticlehook\fi
  \if@book\begin{multicols}{2}[\@makeminititle]\else\@makeminititle\fi
  \begingroup
  \setcounter{footnote}{0}
  }% and at the end of the environment....
  {\endgroup\if@book\end{multicols}\listofabstracts\label{RMAALastPage}\clearpage\fi}
%% auxiliary commands used:
%% \@writeabstract does all the actual formatting 
\newcommand{\@writeabstract}{%
  \if@book%
  \addcontentsline{loa}{abstract}{\protect\@fiddleauthors{\@TOCtitle}{\rmaaTOCauthors}}%
  \fi%
  \begin{center}\@title\par\@author\end{center}
  \@abstract
  \par\noindent\footnoterule\nobreak\footnotesize\@thanks
  \par\noindent\hrulefill\goodbreak
  \vskip 1cm \@plus0.5cm
  %% write ADS entry
  \if@book\if@abstractsads\ADS@dumpdata\fi\fi%
  \endgroup % this pair is to make each abstract a self-contained group
  \begingroup
  \setcounter{footnote}{0}
  }

\if@book
%% List of abstracts (stuff written to .loa file)
\newcommand{\listofabstracts}{%
  \begin{multicols}{2}[%
    \vskip 1.5cm \vfill % leave at least 1.5 cm and go to bottom of page if possible
    \centerline{\MakeUpperCase{List of Abstracts}}
    \vskip 1.5cm \@plus0.5cm
    ][6cm] % this makes it start a new page if there is not at least 6cm of space
    \renewcommand{\@fiddleauthors}[2]{##1\hskip1em plus2em minus0.4em\penalty-100  ##2}
    \parindent\z@
%%    \parskip.5\p@ \@plus .7\p@\relax 
    \parskip3\p@ \@plus 1\p@ \@minus 1\p@\relax % allow quite a lot of
                                % leeway here to avoid breaks inside
                                % entries 
    \renewcommand{\\}{\relax} % temporarily turn off explicit line breaks just in case
    \renewcommand{\altaffilmark}[1]{\relax} % ditto with this
    \@starttoc{loa}
  \end{multicols}
  }
\newcommand*{\l@abstract}[2]{%
  \par\penalty-1000
  #1\hskip 1em\penalty0\null\nobreak\hfill#2}
\fi

%    \end{macrocode}
 
% \subsection{Code for typesetting an entire volume of conference
% proceedings} \label{sec:whole-volume}
%    \begin{macrocode}

\if@book
%% environment for including the individual articles
\newcommand{\@fiddleauthors}[2]{\relax}
\renewcommand{\theequation}{\@arabic\c@equation}
\newenvironment{articles}%
{ \begingroup % make it a group so that author definitions do not
              % spill over
  \gdef\ADS@authors{}
  \pagestyle{rmaa}
  \renewenvironment{document}{%
    \@begofarticlehook
    \addtocounter{chapter}{1}%
    \addcontentsline{toc}{chapter}{%
      \protect\@fiddleauthors{\@TOCtitle}{\rmaaTOCauthors}}%
    \addtocontents{htc}{\protect\HyperTOCEntry{\@rawtitle}{\@rawauthors}{\articlename}}
    }{\relax}
  \renewcommand{\documentclass}[2][foo]{\relax} 
  \renewcommand{\usepackage}[2][foo]{%
    \def\@packagename{##2}\def\@rmaacitename{rmaacite}%
    \ifx\@packagename\@rmaacitename
    \@rmcitetrue\ClassInfo{rmaa}{This paper uses rmaacite}
    \makeatletter%%
%% This is file `rmaacite.sty',
%% generated with the docstrip utility.
%%
%% The original source files were:
%%
%% rmaa.dtx  (with options: `rmaacite')
%% 
%% PLEASE DO NOT EDIT THIS FILE! It is generated automatically from
%% the file rmaa.dtx using the accompanying Makefile and rmaa.ins
%% Any hand-edited changes will be lost
%% 
%% A LaTeX document class for the RMxAA
%% 
%% Author: William Henney (w.henney@crya.unam.mx)
%% 
%% Latest version available from: http://www.crya.unam.mx/rmaa
%% 
%% Copyright (C) 1999-2010 Instituto de Astronomia, UNAM, Mexico
%% 
%% This program is free software; you can redistribute it and/or
%% modify it under the terms of the GNU General Public License
%% as published by the Free Software Foundation; either version 2
%% of the License, or (at your option) any later version.
%% 
%% This program is distributed in the hope that it will be useful,
%% but WITHOUT ANY WARRANTY; without even the implied warranty of
%% MERCHANTABILITY or FITNESS FOR A PARTICULAR PURPOSE.  See the
%% GNU General Public License for more details.
%% 
%% You should have received a copy of the GNU General Public License
%% along with this program; if not, write to the Free Software
%% Foundation, Inc., 59 Temple Place - Suite 330, Boston, MA  02111-1307, USA.






\def\bibitem[#1<#2>]#3{\par\hangindent=1.5em\hangafter=1
  \if@filesw
  {\def\protect##1{\string ##1\space}
    \immediate\write\@auxout{\string\lbibcite{#3}{#1}{#2}}}
  \fi}
\@redefinebibitem % what it says
\def\lbibcite#1#2#3{\global\@namedef{b@#1}{#2}
                    \global\@namedef{by@#1}{#3}}

\def\@bcite#1#2{({#1\if@tempswa , #2\fi})}
\def\@pcite#1#2{#1\if@tempswa , #2\fi}
\def\@brcite{\ (}
\def\@ebrcite{)}
\def\@fmttop#1{\@fmtcite[\csname #1\endcsname]}
\def\@citex[#1]#2{%
%%  \if@filesw\immediate\write\@auxout{\string\citation{#2}}\fi
%% WJH added 16 Jan 2002
  \@for\@citeb:=#2\do{%
    \edef\@citeb{\expandafter\@firstofone\@citeb:\articlename}%
    \if@filesw\immediate\write\@auxout{\string\citation{\@citeb}}\fi}%
%% to here
  \def\@citea{}%
  \@cite{\@for\@citeb:=#2\do
    {%
      \edef\@citeb{\expandafter\@firstofone\@citeb:\articlename}%
      \@citea\def\@citea{;\ }\@ifundefined
      {b@\@citeb}{{\bf ?}\@warning {Citation `\@citeb' on page \thepage \space undefined}}%
      {\csname b@\@citeb\endcsname\@citebro\csname by@\@citeb\endcsname\@citebrc}}}%
  {#1}}

\def\cite{\@ifnextchar [{\let\@citebro=\space
                          \let\@citebrc=\relax
                          \let\@cite=\@bcite\@tempswatrue\@citex}
                        {\let\@citebro=\space
                          \let\@citebrc=\relax
                          \let\@cite=\@bcite\@tempswafalse\@citex[]}}
\def\pcite{\@ifnextchar [{\let\@citebro=\space
                          \let\@citebrc=\relax
                          \let\@cite=\@pcite\@tempswatrue\@citex}
                        {\let\@citebro=\space
                          \let\@citebrc=\relax
                          \let\@cite=\@pcite\@tempswafalse\@citex[]}}
\def\scite{\@ifnextchar [{\let\@citebro=\@brcite
                          \let\@citebrc=\@ebrcite
                          \let\@cite=\@pcite\@tempswatrue\@citex}
                        {\let\@citebro=\@brcite
                          \let\@citebrc=\@ebrcite
                          \let\@cite=\@pcite\@tempswafalse\@citex[]}}

\renewenvironment{thebibliography}
  {\section*{REFERENCES}\vspace*{\afterbiblioskip}\small
    \if@book\RereadBiblio\fi
    \sloppy\clubpenalty4000\widowpenalty4000
    \sfcode`\.=1000\relax
    \setlength\parindent{\z@}\setlength\parskip{\z@}}
  {\relax\normalsize\@outputaddressesifnecessary}








\endinput
%%
%% End of file `rmaacite.sty'.
\makeatother % read it in here so that all changes
    \else                                         % are local to paper
    \if@debug\ClassWarning{rmaa}{Package \@packagename\space was requested by \articlename\space}\fi
    \fi} 
  %% Don't allow individual articles to mess with the pages, volume or year
  \renewcommand{\SetFirstPage}[1]{\relax}
  \renewcommand{\SetYear}[1]{\relax}
  \renewcommand{\SetVolume}[1]{\relax}
  \renewcommand{\SetConfTitle}[1]{\relax}
  %% Don't allow individual articles to change these
  \renewcommand{\nonstopmode}{\relax}
  \renewcommand{\scrollmode}{\relax}
  \renewcommand{\outputplates}{\relax}
  %% XXXX need to check if using rmaacite and act accordingly
  %% avoid multiple declarations of the same length
  \let\oldnewlength=\newlength
  \renewcommand{\newlength}[1]{%
    \edef\reserved@a{\expandafter\@gobble\string ##1}%
    \@ifundefined{\reserved@a}{%
      \oldnewlength{##1}}{%
      \if@debug{\ClassInfo{rmaa}{\reserved@a\space already defined: doing nothing}}\fi%
      \relax}}%
  \renewcommand\newcommand{\providecommand}
  }%
{\endgroup\onecolumn}
     

%    \end{macrocode}
% Command to be used in `articles' environment for including an
% article
%    \begin{macrocode}
\newcommand{\article}[2][]{%
  % optional first argument is extra text to put in brackets after title in TOC
  % e.g. ``Invited Talk'' or ``Poster''
  \renewcommand{\articlename}{#2}
  % reset loads of counters
  \setcounter{footnote}{0}
  \setcounter{figure}{0}
  \setcounter{table}{0}
  \setcounter{section}{0}
  \setcounter{equation}{0}
  \setcounter{rmaabib}{0}
  \setcounter{author}{0}
  \setcounter{AUTHOR}{1}
  \gdef\ADS@authors{}%
  \gdef\ADS@affil{}%
  \gdef\ADS@affiliations{}%
  \gdef\ADS@references{}%
  \@firstauthortrue
  \@firstaffiltrue
  \@firstaltaffiltrue
  \@firstkwdtrue
  \@rmcitefalse
  \if@suppressfulladdresses\else\@needaddressestrue\fi
  \renewcommand{\@rmaaArg}{#1}
  \ifx\@rmaaArg\@rmaaEmpty
     \renewcommand{\@TOCaftertitle}{}
  \else
     \renewcommand{\@TOCaftertitle}{ (#1)}
  \fi
%    \end{macrocode}
% In the main journal articles must start on an odd page. 
%    \begin{macrocode}
  \if@journal\clearpage\thispagestyle{empty}\cleardoublepage\fi
%% Now input the file, then run the end-of-article hook (reinstated 28
  %% Sep 2001)
  {\input{\@articledir/\articlename.tex}\@endofarticlehook}}
%    \end{macrocode}
% Command to separate articles into sessions. Write an entry in the
% TOC, makes sure we start on an odd page (added 06 Jun 2000), then do
% a title page if \cs{SessionTitlePages} was called. 
% \changes{3.12}{04 Feb 2003}{Added option to have title page for each
% session. Turn it on with \cs{SessionTitlePages}.}
%    \begin{macrocode}
\newcommand{\session}[1]{
  %% this does nothing except write an entry in the TOC
  \addcontentsline{toc}{part}{\rmaa@TOCPartStyle{#1}}
  \addtocontents{htc}{\protect\HyperTOCpart{#1}}
  %% and also makes sure we start on an odd page (added 06 Jun 2000)
  \clearpage\thispagestyle{empty}\cleardoublepage
  \stepcounter{part}
  \if@sessiontitlepages{%
    \onecolumn
    \thispagestyle{empty}
    \renewcommand{\articlename}{Section\thepart TitlePage}
    \label{RMAAFirstPage}
    \make@session@title{#1}
    \label{RMAALastPage}
    \clearpage\thispagestyle{empty}\cleardoublepage
  }\fi
  }
\newcommand{\make@session@title}[1]{%
  {%
    \setlength\parindent{0pt}
    \vspace*{\stretch{1}}
    \rule{\linewidth}{1mm}
    \begin{flushright}\Huge\bfseries\boldmath
      #1
    \end{flushright}
    \rule{\linewidth}{1mm}
    \vspace*{\stretch{2}}
  }
}
%    \end{macrocode}
% Command to turn on separate title pages for each session. 
%    \begin{macrocode}
\newcommand{\SessionTitlePages}{\@sessiontitlepagestrue}
\newcommand{\BoldSessionTitles}{%
  \renewcommand{\rmaa@TOCPartStyle}[1]{\textbf{\rmaa@TOCTitleStyle{##1}}}
  }
\newcommand{\ArticleDir}[1]{\renewcommand{\@articledir}{#1/}\graphicspath{{1/}}}
\newcommand{\@articledir}{.}
%\newcommand{\ps@rmaafront}{\ps@rmaa}% page style for front matter (same as ever, actually)
\newcommand{\ps@rmaafront}{% page style for front matter:  page # in each foot
  \renewcommand\@oddfoot{\hfil{\thepage}}%
  \renewcommand\@evenfoot{{\thepage}\hfil}%
  % shorttitle in both heads
  \renewcommand\@evenhead{\hfil\@shorttitle\hfil}%
  \renewcommand\@oddhead{\@evenhead}
  }
\newcommand{\ps@rmaafrontfirstpage}{% page style for first page of front matter
  % page # in each foot, on outside edge
  \renewcommand\@oddfoot{\hfil\thepage}%
  \renewcommand\@evenfoot{\thepage\hfil}%
  % nothing in the even heads
  \renewcommand\@evenhead{\relax}%
  \renewcommand\@oddhead{\relax}%
  }
%% slight modification to the standard \frontmatter and \mainmatter
%% commands, just to set the \pagestyle
\renewcommand\frontmatter{%
  \latexonecolumn
  \cleardoublepage
  \@frontmattertrue
  \@mainmatterfalse
  \pagenumbering{roman}
  \pagestyle{rmaafront}
  \addtocounter{page}{2} % leave space for the title page
  }
\renewcommand\mainmatter{%
  \pagestyle{empty}
  \if@rmtwocol\twocolumn\fi
  \cleardoublepage
  \@frontmatterfalse
  \@mainmattertrue
  \pagenumbering{arabic}
  \pagestyle{rmaa}
  }
\renewcommand\backmatter{%
  \clearpage\pagestyle{empty}\cleardoublepage
  \@mainmatterfalse
  %% put extra space in the TOC
  \addtocontents{toc}{\bigskip}
  }
%    \end{macrocode}
%
% \subsection{Front sections}\label{sec:front-sections}
% \begin{macro}{\notitle}
% \begin{macro}{\shortname}
% \begin{macro}{frontsection}
% The environment \texttt{frontsection} is for doing a section in the
% front matter (preface, list of participants, or whatever). It has 
% one optional argument:
% \begin{enumerate}
% \item[\#1] Short title (no embedded spaces) for
% use in making PS files. E.g., ``list\_of\_participants''.   
% \end{enumerate}
% and two mandatory arguments: 
% \begin{enumerate}
% \item[\#2] Title of section (e.g., `Preface')
% \item[\#3] Commands to be evaluated before processing the contents
% of the environment. This is mainly useful for putting in the things
% that would go in the preamble of a normal article  (e.g.,
% |\listofauthors|, |\indexauthor|, |\notitle|) 
% \end{enumerate}
%    \begin{macrocode}
\newcommand{\notitle}{\@titlevisiblefalse} \newcommand{\shortname}{}
\newenvironment{frontsection}[3][]{%
  \begingroup
  \renewcommand{\shortname}{#1}
  \ifx\shortname\@rmaaEmpty
     \renewcommand{\articlename}{#2}
  \else
     \renewcommand{\articlename}{\shortname}
  \fi
  \@begofarticlehook
  \@titlevisibletrue % default
  \thispagestyle{rmaafrontfirstpage}
  \listofauthors{} % blank out this in case it is not specified
  \title{{#2}}\shorttitle{#2}\shortauthor{#2}
  #3 % evaluate the command given as second mandatory arg
  \addcontentsline{toc}{chapter}{\protect\@fiddleauthors{\@TOCtitle}{\rmaaTOCauthors}}%
  \addtocontents{htc}{\protect\HyperTOCEntry{\@rawtitle}{\@rawauthors}{\articlename}}
  \if@titlevisible
     \vspace*{1cm}\centerline{\rmaa@FrontSecStyle{\@title}}\vspace*{1cm}
  \fi
  \setlength{\parskip}{3ex}
  }{\frontlabel{RMAALastPage}\clearpage\endgroup}
%    \end{macrocode}
% \end{macro}
% \end{macro}
% \end{macro}
%    \begin{macrocode}
%% Definition of the command \frontlabel, which rights out something like `F4' to
%% the .aux file, instead of `\textlatin  {i}'. This is so it is easy to post-process
%% the file for making PS files of individual sections
\newcommand{\frontlabel}[1]{%
  \@bsphack
  \protected@write\@auxout{}%
     {\string\newlabel{\articlename:#1}{{}{F\@arabic{\c@page}}}}
  \@esphack
  }
%% User command to say X, Y, Z,\\ Editors of the Proceedings
\newcommand{\EditorSignature}{\vspace*{\parskip}\rightline{\@editors}%
  \rightline{Editors of the proceedings}}
%% Fiddle with \thesection, etc. so that they don't include the chapter number
%% Basically go back to the `article' definitions of these
\renewcommand{\thesection}{\@arabic\c@section}
\renewcommand{\thefigure}{\@arabic\c@figure}
\renewcommand{\thetable}{\@arabic\c@table}

\newcommand{\@oldargs}{}
\newcommand{\@newargs}{}
%% add TOC entry for each paper
%\renewcommand{\title}[1]%
%{\renewcommand{\@title}{\rmaa@TitleStyle{#1}}%
%\addtocounter{chapter}{1}\addcontentsline{toc}{chapter}{#1}%
%}

%% Formatting of the table of contents
%% dimensions of various parts
\newlength{\TOCtextwidth}\setlength{\TOCtextwidth}{12.65cm}
\newlength{\TOCpagenumwidth}\setlength{\TOCpagenumwidth}{1.3cm}
\newlength{\TOCentrywidth}\setlength{\TOCentrywidth}{\TOCtextwidth}
\addtolength{\TOCentrywidth}{-\TOCpagenumwidth} 
\newlength{\TOCextramargin}\setlength{\TOCextramargin}{1.25cm}
\newcommand{\TOCbox}[1]{\hspace{\TOCextramargin}\parbox{\TOCtextwidth}{#1}\par}
\newlength{\TOCextramarginx}\setlength{\TOCextramarginx}{\TOCextramargin}\addtolength{\TOCextramarginx}{4.3cm}
\newlength{\TOCtextwidthx}\setlength{\TOCtextwidthx}{\TOCtextwidth}\addtolength{\TOCtextwidthx}{-4.3cm}
\newcommand{\TOCboxx}[1]{\hspace{\TOCextramarginx}\parbox{\TOCtextwidthx}{#1}\par}

\newcommand{\RMAAbannerimage}{rmsc}
\renewcommand\tableofcontents{%
  \renewcommand{\articlename}{CONTENTS}
  \begingroup
  \@begofarticlehook
  \if@twocolumn
     \@restonecoltrue\onecolumn
  \else
     \@restonecolfalse
  \fi
  %% 9pt for journal style
  \if@proceedings\else\small\fi
  \if@newstyle\small\fi %% same for new style
  \if@proceedings\thispagestyle{rmaafrontfirstpage}\else\thispagestyle{empty}\fi
  \shortauthor{Contents}\shorttitle{Contents}
  \vspace*{18mm}
  \if@printbanner
  %% put in the banner header (need to get better image of this)
     \TOCbox{\hspace*{-0.15cm}\includegraphics[width=0.9\TOCtextwidth]{\RMAAbannerimage}}
     \vspace*{8mm}
  %% volume and publication date (should really be in funny font)
     \TOCboxx{\@volumen\\[2mm] \@mes\ \@anyo}
  \else
  %% put in the banner header (need to get better image of this)
     \TOCbox{\hspace*{-0.15cm}\includegraphics[width=0.9\TOCtextwidth]{\RMAAbannerimage}}
     \vspace*{12.5mm}
  %% volume and publication date (should really be in funny font)
     \TOCboxx{~\\[2mm] ~}
  \fi
  \vspace*{16mm}
  %% This is the command that fiddles with the TOC entry so that the
  %% title is flushleft but the authors are flushright
  \renewcommand{\@fiddleauthors}[2]{%
    % first put the paper title
    \renewcommand{\\}{\relax} % temporarily turn off explicit line breaks just in case
    \renewcommand{\altaffilmark}[1]{\relax} % ditto with this
    \renewcommand{\textsuperscript}[1]{\relax} % ditto with this 
    ##1
    % now save the author list in a box
    \sbox\@tempboxa{##2}%
    % check that box is shorter than the total width available for the entry
    \ifdim \wd\@tempboxa <\TOCentrywidth
       % if so, then put it right here, flush against right margin
       \penalty10000\hskip1em plus10em\penalty0\hbox{}\nobreak\hfill\box\@tempboxa
       % Note that we force a space of at least 1em between title and authorlist, 
       % stretchable to 11em since we may have very short titles but still with an 
       % authorlist too long to go on the same line.  The only linebreak allowed is 
       % directly before the authorlist.
    \else
       % if it is too long, we set the original authorlist text in its own parbox
       \par\vskip1ex\parbox[b]{\TOCentrywidth}{\raggedleft ##2}
    \fi
    }
  \@starttoc{toc}%
  \if@restonecol\twocolumn\fi
  \frontlabel{RMAALastPage}
  \clearpage
  \thispagestyle{empty}
  \cleardoublepage % make sure we start on an odd page
  \endgroup

  %%% ALSO, do the HTML TOC
  %% The file that will be used to make the HTML TOC
  \@starttoc{htc}
  %% All the stuff in this file does nothing in the book style
  \addtocontents{htc}{\protect\HyperTOCmeeting{\@conftitle}{\@reunionplace}{\@reuniondate}}
  \addtocontents{htc}{\protect\HyperTOCeditors{\@editors}}
  \addtocontents{htc}{\protect\HyperTOCpreamble}
  \addtocontents{htc}{\protect\HyperTOCEntry{Table of Contents}{}{\articlename}}

  %% and the ADS file
  \ADS@openfile
  }
%% This is now much simpler
\renewcommand\l@chapter[2]{%TOC entry for articles
  \ifnum \c@tocdepth >\m@ne
    \vskip 1.0em \@plus\p@
    \TOCbox{%
      \parbox[b]{\TOCentrywidth}{#1}\parbox[b]{\TOCpagenumwidth}{\hfill#2}
    }
  \fi} 
\renewcommand\l@part[2]{%TOC entry for sessions: centered, no page number
  \ifnum \c@tocdepth >\m@ne
    \vskip 1.5em \@plus1.5em \goodbreak
    \TOCbox{%
      \parbox[b]{\TOCentrywidth}{\hfil#1\hfil}\hfill
    }
  \fi\nobreak} 

%% Formatting theindex
%% strongly encourage page breaks between index entries
\renewcommand{\@idxitem}{\par\penalty-1000\hangindent.45cm} 
\renewcommand{\indexspace}{\par\penalty-1000\vskip 10\p@ \@plus15\p@ \@minus3\p@\relax}
\renewenvironment{theindex}{%
  %% first, make sure we start on an odd page (06 Jun 2000)
  \clearpage\thispagestyle{empty}\cleardoublepage
  %% this looks a bit like an article but is two columns and without abstract, etc
  \renewcommand{\articlename}{INDEX}
  \title{Author Index}\shortauthor{Author Index}\shorttitle{Author Index}
  \listofauthors{}
  \addcontentsline{toc}{chapter}{\protect\@fiddleauthors{\@TOCtitle}{\rmaaTOCauthors}}%
  \addtocontents{htc}{\protect\HyperTOCEntry{\@rawtitle}{\@rawauthors}{\articlename}}
  \pagestyle{rmaa}
  \thispagestyle{rmaafirstpage}
  \@begofarticlehook
  %% 9pt for journal style
  \if@proceedings\else\small\fi
  \if@newstyle\small\fi
  \twocolumn[\@makeminititle]
  %% vertical spacing (see also definition of \@idxitem and \indexspace
  \parskip\z@ \@plus .6\p@\relax % allow quite a lot of leaway here to avoid breaks inside entries
  %% horizontal spacing, breaking and glue parameters (see also the stuff in the definition 
  %% of \indexauthor
  \parindent\z@ 
  \exhyphenpenalty10000 % don't break on explicit hyphens, it always looks bad
  \rightskip=0pt plus.2em % allow a tiny bit of leeway at the ends of lines
  \spaceskip.3333em plus.5em minus.15em % a lot of leeway between words
  \xspaceskip.5em plus.5em minus.4em % this seems to have no effect
  \let\item\@idxitem
  \renewcommand{\\}{\relax} % temporarily turn off explicit line breaks just in case
  \renewcommand{\altaffilmark}[1]{\relax} % ditto with this
  \renewcommand{\textsuperscript}[1]{\relax} % ditto with this
  }{\label{RMAALastPage}\clearpage} % clearpage necessary so that we know about \shortauthor



%% end definitions for book styles
\fi

%    \end{macrocode}
% In order to do away with plates, we sometimes want a figure to
% appear in color only in the online version. The following command
% can bve used in a figure caption to typeset a note to this effect. 
%    \begin{macrocode}
\newcommand\ColorOnlineOnly{%
  \textsc{Note: this figure is available in color in the electronic
    version of this article, obtainable from}
  \texttt{http://\discretionary{}{}{}www.\discretionary{}{}{}astroscu.\discretionary{}{}{}unam.\discretionary{}{}{}mx/\discretionary{}{}{}\~{}rmaa/}.}
%% Plates
%% Define 2 counters: plate is the number in the whole document, plateloop is used 
%% internally by \outputplates
\newcounter{plate}
\newcounter{plateloop}
%% Currently limited to 15 plates in document
%% when we do a label, we need to write to \newlabel entries to the .aux file, one 
%% that says RMAAPlateX so that it can be referred to in \ouputplates, and one that
%% says articlename:RMAAPlateX so that the dvips program can find it and asociate it
%% with the right article
\newcommand{\@platelabeli}{\oldlabel{RMAAPlate1}\label{RMAAPlate1}}
\newcommand{\@platelabelii}{\oldlabel{RMAAPlate2}\label{RMAAPlate2}}
\newcommand{\@platelabeliii}{\oldlabel{RMAAPlate3}\label{RMAAPlate3}}
\newcommand{\@platelabeliv}{\oldlabel{RMAAPlate4}\label{RMAAPlate4}}
\newcommand{\@platelabelv}{\oldlabel{RMAAPlate5}\label{RMAAPlate5}}
\newcommand{\@platelabelvi}{\oldlabel{RMAAPlate6}\label{RMAAPlate6}}
\newcommand{\@platelabelvii}{\oldlabel{RMAAPlate7}\label{RMAAPlate7}}
\newcommand{\@platelabelviii}{\oldlabel{RMAAPlate8}\label{RMAAPlate8}}
\newcommand{\@platelabelix}{\oldlabel{RMAAPlate9}\label{RMAAPlate9}}
\newcommand{\@platelabelx}{\oldlabel{RMAAPlate10}\label{RMAAPlate10}}
\newcommand{\@platelabelxi}{\oldlabel{RMAAPlate11}\label{RMAAPlate11}}
\newcommand{\@platelabelxii}{\oldlabel{RMAAPlate12}\label{RMAAPlate12}}
\newcommand{\@platelabelxiii}{\oldlabel{RMAAPlate13}\label{RMAAPlate13}}
\newcommand{\@platelabelxiv}{\oldlabel{RMAAPlate14}\label{RMAAPlate14}}
\newcommand{\@platelabelxv}{\oldlabel{RMAAPlate15}\label{RMAAPlate15}}
\newcommand{\@outputplatelabeli}{\oldlabel{RMAAOutputPlate1}}
\newcommand{\@outputplatelabelii}{\oldlabel{RMAAOutputPlate2}}
\newcommand{\@outputplatelabeliii}{\oldlabel{RMAAOutputPlate3}}
\newcommand{\@outputplatelabeliv}{\oldlabel{RMAAOutputPlate4}}
\newcommand{\@outputplatelabelv}{\oldlabel{RMAAOutputPlate5}}
\newcommand{\@outputplatelabelvi}{\oldlabel{RMAAOutputPlate6}}
\newcommand{\@outputplatelabelvii}{\oldlabel{RMAAOutputPlate7}}
\newcommand{\@outputplatelabelviii}{\oldlabel{RMAAOutputPlate8}}
\newcommand{\@outputplatelabelix}{\oldlabel{RMAAOutputPlate9}}
\newcommand{\@outputplatelabelx}{\oldlabel{RMAAOutputPlate10}}
\newcommand{\@outputplatelabelxi}{\oldlabel{RMAAOutputPlate11}}
\newcommand{\@outputplatelabelxii}{\oldlabel{RMAAOutputPlate12}}
\newcommand{\@outputplatelabelxiii}{\oldlabel{RMAAOutputPlate13}}
\newcommand{\@outputplatelabelxiv}{\oldlabel{RMAAOutputPlate14}}
\newcommand{\@outputplatelabelxv}{\oldlabel{RMAAOutputPlate15}}
\newcommand{\@platerefi}{\oldpageref{RMAAPlate1}}
\newcommand{\@platerefii}{\oldpageref{RMAAPlate2}}
\newcommand{\@platerefiii}{\oldpageref{RMAAPlate3}}
\newcommand{\@platerefiv}{\oldpageref{RMAAPlate4}}
\newcommand{\@platerefv}{\oldpageref{RMAAPlate5}}
\newcommand{\@platerefvi}{\oldpageref{RMAAPlate6}}
\newcommand{\@platerefvii}{\oldpageref{RMAAPlate7}}
\newcommand{\@platerefviii}{\oldpageref{RMAAPlate8}}
\newcommand{\@platerefix}{\oldpageref{RMAAPlate9}}
\newcommand{\@platerefx}{\oldpageref{RMAAPlate10}}
\newcommand{\@platerefxi}{\oldpageref{RMAAPlate11}}
\newcommand{\@platerefxii}{\oldpageref{RMAAPlate12}}
\newcommand{\@platerefxiii}{\oldpageref{RMAAPlate13}}
\newcommand{\@platerefxiv}{\oldpageref{RMAAPlate14}}
\newcommand{\@platerefxv}{\oldpageref{RMAAPlate15}}
\newcounter{figplatei}
\newcounter{figplateii}
\newcounter{figplateiii}
\newcounter{figplateiv}
\newcounter{figplatev}
\newcounter{figplatevi}
\newcounter{figplatevii}
\newcounter{figplateviii}
\newcounter{figplateix}
\newcounter{figplatex}
\newcounter{figplatexi}
\newcounter{figplatexii}
\newcounter{figplatexiii}
\newcounter{figplatexiv}
\newcounter{figplatexv}
%% Printed output should use the counter `plate'
\renewcommand{\theplate}{\@arabic\c@plate}
\renewcommand{\theplateloop}{\@arabic\c@plateloop}
%% Command to output all the saved plates in this paper
\newcommand{\outputplates}{%
  \setcounter{plateloop}{0}
  \if@book\else\@endofarticlehook\fi % not entirely sure why, but this
                                     % seems to be needed
  \clearpage % first output the last page of text
  % now set pagestyle for blank even pages and make sure we start on an odd page
  \pagestyle{empty}
  \cleardoublepage
  \onecolumn
  \parindent\z@
  % first put in an extra sheet
  % the main purpose of this is so there is always a blank even page before 
  % each plate for filling in the PS file if the article has an odd number
  % of pages (see splitbook.sh)
  \null\vfil\centerline{PLATES}\vfil\clearpage\pagestyle{rmaaplate}\cleardoublepage
  % output each plate in turn
  \@whilenum\value{plateloop}<\value{plate}\do{%
    \addtocounter{plateloop}{1}
    % restore the figure counter from \@figcounti, \@figcountii, etc
    \setcounter{figure}{\value{figplate\roman{plateloop}}}
    % spew out the plate (in its own group for safety)
    {\csname @plate\roman{plateloop}\endcsname\par}
    }
  \pagestyle{rmaa}
  % reset the page number as if nothing had happened
  \addtocounter{page}{-\value{plate}}
  \addtocounter{page}{-\value{plate}}
  }
\newcommand\figcaption{%
  \def\@captype{figure}
  \renewcommand{\@makecaption}{\@makefigcap}
  \caption}
\newcommand\tabcaption{%
  \def\@captype{table}
  \renewcommand{\@makecaption}{\@maketabcap}
  \caption}

\newcommand{\plate}[2][1]{%
  %% User command to make a plate page: 
  %% Mandatory argument is what goes on the plate - this should not 
  %% include any floating environments (e.g. figure). Instead, \figcaption 
  %% should be used to do the captions. 
  %% Optional argument is the number of figures (i.e. captions) on the plate
  %% page (default is 1). This is necessary for incrementing the figure counter. 
  %%
  %% Note that we have to use a \gdef rather than an \xdef when we 
  %% save the plate output in a command. This means that in the book styles when it 
  %% comes to spewing out the plate, we will have forgotten about any user macros 
  %% that may have been defined. This can be got round by repeating the macros inside the plate.
  %% (Presumably, one would use a \newcommand rather than \renewcommand if they were new.)
  % increment global and local counters
  \addtocounter{plate}{1}
  % write appropriate label to .aux file
  \csname @platelabel\roman{plate}\endcsname
  % define a global macro called \@platei, \@plateii, \@plateiii, \@plateiv, ...
  \expandafter\gdef\csname @plate\roman{plate}\endcsname{%
    \csname @outputplatelabel\roman{plateloop}\endcsname
    \null\vfil #2 \vfil
    \cleardoublepage
    }
  % save the current shortauthor and shorttitle
  % need to be \xdef to take current values of \@shortauthor, etc 
  \expandafter\xdef\csname @shorta\roman{plate}\endcsname{\@shortauthor}
  \expandafter\xdef\csname @shortt\roman{plate}\endcsname{\@shorttitle}
  % save the current value of the figure counter for later use in \outputplates
  \setcounter{figplate\roman{plate}}{\value{figure}}
  % increment the figure counter according to the optional argument
  \addtocounter{figure}{#1}
  }
%% page style for plate pages 
\newcommand{\ps@rmaaplate}{%
  % Plate X (see page Y) in the foot
  \renewcommand\@oddfoot{\hfil\csname @shorta\roman{plateloop}\endcsname\ 
    (see page~\csname @plateref\roman{plateloop}\endcsname)\hfil}%
  \renewcommand\@evenfoot{\relax}%
  % nothing in even heads (blank pages)
  \renewcommand\@evenhead{\relax}
  % short title and page num in the odd heads
  \renewcommand\@oddhead{\hfil\csname @shortt\roman{plateloop}\endcsname%
    \hfil\llap{PLATE~\theplateloop}}
  }





%% Things for the User Guide
\if@guide
%% stolen from doc.sty
\def\MakeShortVerb#1{%
  \expandafter\ifx\csname cc\string#1\endcsname\relax
  \@shortvrbinfo{Made }{#1}%
  \add@special{#1}%
  \expandafter
  \xdef\csname cc\string#1\endcsname{\the\catcode`#1}%
  \begingroup
  \catcode`\~\active  \lccode`\~`#1%
  \lowercase{%
    \global\expandafter\let
    \csname ac\string#1\endcsname~%
    \gdef~{\verb~}}%
  \endgroup
  \global\catcode`#1\active
  \else
  \@shortvrbinfo\@empty{#1 already}%
  \fi}
\def\DeleteShortVerb#1{%
  \expandafter\ifx\csname cc\string#1\endcsname\relax
  \else
  \@shortvrbinfo{Deleted }{#1 as}%
  \rem@special{#1}%
  \global\catcode`#1\csname cc\string#1\endcsname
  \global \expandafter\let \csname cc\string#1\endcsname \relax
  \ifnum\catcode`#1=\active
  \begingroup
  \catcode`\~\active   \lccode`\~`#1%
  \lowercase{%
    \global\expandafter\let\expandafter~%
    \csname ac\string#1\endcsname}%
  \endgroup \fi \fi}
\def\@shortvrbinfo#1#2{%
  \PackageInfo{doc}{%
    #1\expandafter\@gobble\string#2 a short reference
    for \string\verb}}
\def\add@special#1{%
  \rem@special{#1}%
  \expandafter\gdef\expandafter\dospecials\expandafter
  {\dospecials \do #1}%
  \expandafter\gdef\expandafter\@sanitize\expandafter
  {\@sanitize \@makeother #1}}
\def\rem@special#1{%
  \def\do##1{%
    \ifnum`#1=`##1 \else \noexpand\do\noexpand##1\fi}%
  \xdef\dospecials{\dospecials}%
  \begingroup
  \def\@makeother##1{%
    \ifnum`#1=`##1 \else \noexpand\@makeother\noexpand##1\fi}%
  \xdef\@sanitize{\@sanitize}%
  \endgroup}
%% define | as shortcut for \verbatim
\MakeShortVerb{\|}
%% comand to specify arguments to LaTeX commands
\newcommand{\rmArg}[1]{%
  \ensuremath{\langle}{\scriptsize \textsf{\MakeUpperCase{#1}}}\ensuremath{\rangle}}
%\newcommand{\rmArg}[1]{\ensuremath{\langle}\textsc{#1}\ensuremath{\rangle}}
\providecommand{\BibTeX}{{\rmfamily B\kern-.05em%
    \textsc{i\kern-.025em b}\kern-.08em%
    T\kern-.1667em\lower.7ex\hbox{E}\kern-.125emX}}
%% Environment for code examples
\newenvironment{Code}{%
  \begin{list}{}{\setlength{\leftmargin}{0.05\columnwidth}} 
  \item[]}{\end{list}}
\fi
%</package>
%    \end{macrocode}

% \section{User commands for astronomical articles}
% \label{sec:user-commands}
% Some random definitions. These go both into the main class file and
% also into |adshtml.sty|.
% 
% Definitions of the journal name in different incarnations.  
%    \begin{macrocode}
%<*package|adshtml-style>
\newcommand{\RMAA}{RMxAA}
\newcommand{\RMAAlong}{Revista Mexicana de Astronom\'{\i}a y Astrof\'{\i}sica}
\newcommand{\RMAASC}{RMxAC}
\newcommand{\RMAASClong}{\RMAA\ (Serie de Conferencias)}
%    \end{macrocode}
% \begin{macro}{\ibidrule}
% This puts a horizontal rule with a period after it. It is used when
% the author list of a reference is identical to the preceding one. 
%    \begin{macrocode}
\newcommand{\ibidrule}{\rule[-0.03cm]{4.7em}{.01cm}.\ }
%    \end{macrocode}
% \end{macro}
% \begin{macro}{\ThisVolume}
% References to other articles in the same volume can be coded as,
% e.g., |\bibitem{} Bloggs, J. \ThisVolume{bloggs}|.
%    \begin{macrocode}
\newcommand{\ThisVolume}[1]{%
  \@year, \if@proceedings\RMAASC\else\RMAA\fi, \@volume,
  \if@book\oldpageref{#1:RMAAFirstPage}\else 000\fi\ (this volume)}
%    \end{macrocode}
% \end{macro}
% Other stuff that may be useful in ref lists.
%    \begin{macrocode}
\newcommand{\etal}{et al.\@}
%    \end{macrocode}
% Various new commands from AAS macros
%    \begin{macrocode}
\def\acknowledgments{\vskip 3ex plus .8ex minus .4ex}
\let\acknowledgements=\acknowledgments                  % second spelling
\let\jnl@style=\rm
\def\ref@jnl#1{{\jnl@style#1}}
\DeclareRobustCommand\aj{\ref@jnl{AJ}}                   
\DeclareRobustCommand\araa{\ref@jnl{ARA\&A}}             
\DeclareRobustCommand\apj{\ref@jnl{ApJ}}                 
\DeclareRobustCommand\apjl{\ref@jnl{ApJ}}                
\DeclareRobustCommand\apjs{\ref@jnl{ApJS}}               
\DeclareRobustCommand\ao{\ref@jnl{Appl.~Opt.}}           
\DeclareRobustCommand\apss{\ref@jnl{Ap\&SS}}             
\DeclareRobustCommand\aap{\ref@jnl{A\&A}}                
\DeclareRobustCommand\aapr{\ref@jnl{A\&A~Rev.}}          
\DeclareRobustCommand\aaps{\ref@jnl{A\&AS}}              
\DeclareRobustCommand\azh{\ref@jnl{AZh}}                 
\DeclareRobustCommand\baas{\ref@jnl{BAAS}}               
\DeclareRobustCommand\jrasc{\ref@jnl{JRASC}}             
\DeclareRobustCommand\memras{\ref@jnl{MmRAS}}            
\DeclareRobustCommand\mnras{\ref@jnl{MNRAS}}             
\DeclareRobustCommand\pra{\ref@jnl{Phys.~Rev.~A}}        
\DeclareRobustCommand\prb{\ref@jnl{Phys.~Rev.~B}}        
\DeclareRobustCommand\prc{\ref@jnl{Phys.~Rev.~C}}        
\DeclareRobustCommand\prd{\ref@jnl{Phys.~Rev.~D}}        
\DeclareRobustCommand\pre{\ref@jnl{Phys.~Rev.~E}}        
\DeclareRobustCommand\prl{\ref@jnl{Phys.~Rev.~Lett.}}    
\DeclareRobustCommand\pasp{\ref@jnl{PASP}}               
\DeclareRobustCommand\pasj{\ref@jnl{PASJ}}               
\DeclareRobustCommand\qjras{\ref@jnl{QJRAS}}             
\DeclareRobustCommand\skytel{\ref@jnl{S\&T}}             
\DeclareRobustCommand\solphys{\ref@jnl{Sol.~Phys.}}      
\DeclareRobustCommand\sovast{\ref@jnl{Soviet~Ast.}}      
\DeclareRobustCommand\ssr{\ref@jnl{Space~Sci.~Rev.}}     
\DeclareRobustCommand\zap{\ref@jnl{ZAp}}                 
\DeclareRobustCommand\nat{\ref@jnl{Nature}}              
\DeclareRobustCommand\iaucirc{\ref@jnl{IAU~Circ.}}
\DeclareRobustCommand\aplett{\ref@jnl{Astrophys.~Lett.}}
\DeclareRobustCommand\apspr{\ref@jnl{Astrophys.~Space~Phys.~Res.}}
\DeclareRobustCommand\bain{\ref@jnl{Bull.~Astron.~Inst.~Netherlands}}
\DeclareRobustCommand\fcp{\ref@jnl{Fund.~Cosmic~Phys.}}
\DeclareRobustCommand\gca{\ref@jnl{Geochim.~Cosmochim.~Acta}}
\DeclareRobustCommand\grl{\ref@jnl{Geophys.~Res.~Lett.}}
\DeclareRobustCommand\jcp{\ref@jnl{J.~Chem.~Phys.}}      
\DeclareRobustCommand\jgr{\ref@jnl{J.~Geophys.~Res.}}    
\DeclareRobustCommand\jqsrt{\ref@jnl{J.~Quant.~Spec.~Radiat.~Transf.}}
\DeclareRobustCommand\memsai{\ref@jnl{Mem.~Soc.~Astron.~Italiana}}
\DeclareRobustCommand\nphysa{\ref@jnl{Nucl.~Phys.~A}}
\DeclareRobustCommand\physrep{\ref@jnl{Phys.~Rep.}}
\DeclareRobustCommand\physscr{\ref@jnl{Phys.~Scr}}
\DeclareRobustCommand\planss{\ref@jnl{Planet.~Space~Sci.}}
\DeclareRobustCommand\procspie{\ref@jnl{Proc.~SPIE}}
\let\astap=\aap
\let\apjlett=\apjl
\let\apjsupp=\apjs
\let\applopt=\ao
\def\phn{\phantom{0}}
\def\phd{\phantom{.}}
\def\phs{\phantom{\ensuremath{-}}}
\def\phm#1{\phantom{#1}}
%    \end{macrocode}
% All these are redone as robust and using more modern commands where
% possible. The robustness is so they will not get expanded if used in
% abstracts and titles when these are passed to the ADS file. 
% \changes{3.16}{2003/04/07}{
%  Added \cs{Color} command for doing photometric colors.}
%    \begin{macrocode}
%%\def\sun{\hbox{$\odot$}}
\DeclareRobustCommand\sun{\ensuremath{\odot}}
%%\def\earth{\hbox{$\oplus$}}
\DeclareRobustCommand\earth{\ensuremath{\oplus}}
%%\def\lesssim{\mathrel{\hbox{\rlap{\hbox{\lower4pt\hbox{$\sim$}}}\hbox{$<$}}}}
%%\def\gtrsim{\mathrel{\hbox{\rlap{\hbox{\lower4pt\hbox{$\sim$}}}\hbox{$>$}}}}
\DeclareRobustCommand\lesssim{\mathrel{\hbox{\smaller\smaller\rlap{\hbox{\lower3pt\hbox{$\sim$}}}\raise1.5pt\hbox{$<$}}}}
\DeclareRobustCommand\gtrsim{\mathrel{\hbox{\smaller\smaller\rlap{\hbox{\lower3pt\hbox{$\sim$}}}\raise1.5pt\hbox{$>$}}}}
%%\def\lta{\mathrel{\spose{\lower 3pt\hbox{$\mathchar"218$}}
%%     \raise 2.0pt\hbox{$\mathchar"13C$}}}
%%\def\gta{\mathrel{\spose{\lower 3pt\hbox{$\mathchar"218$}}
%%     \raise 2.0pt\hbox{$\mathchar"13E$}}}
\DeclareRobustCommand\sq{\hbox{\rlap{$\sqcap$}$\sqcup$}}
%%\def\arcdeg{\hbox{$^\circ$}}
%%\def\arcmin{\hbox{$^\prime$}}
%%\def\arcsec{\hbox{$^{\prime\prime}$}}
\DeclareRobustCommand\arcdeg{\ensuremath{^\circ}}
\DeclareRobustCommand\arcmin{\ensuremath{^\prime}}
\DeclareRobustCommand\arcsec{\ensuremath{^{\prime\prime}}}
%%\def\fd{\hbox{$.\!\!^{\rm d}$}}
%%\def\fh{\hbox{$.\!\!^{\rm h}$}}
%%\def\fm{\hbox{$.\!\!^{\rm m}$}}
%%\def\fs{\hbox{$.\!\!^{\rm s}$}}
%%\def\fdg{\hbox{$.\!\!^\circ$}}
%%\def\farcm{\hbox{$.\mkern-4mu^\prime$}}
%%\def\farcs{\hbox{$.\!\!^{\prime\prime}$}}
%%\def\fp{\hbox{$.\!\!^{\scriptscriptstyle\rm p}$}}
\DeclareRobustCommand\fd{\ensuremath{.\!\!^\mathrm{d}}}
\DeclareRobustCommand\fh{\ensuremath{.\!\!^\mathrm{h}}}
\DeclareRobustCommand\fm{\ensuremath{.\!\!^\mathrm{m}}}
\DeclareRobustCommand\fs{\ensuremath{\mkern0.5mu.\mkern-5mu^\mathrm{s}\mkern-1mu}}
\DeclareRobustCommand\fdg{\ensuremath{.\!\!^\circ}}
\DeclareRobustCommand\farcm{\ensuremath{.\mkern-4mu^\prime}}
\DeclareRobustCommand\farcs{\ensuremath{\mkern1mu.\mkern-6mu^{\prime\prime}\mkern-1.5mu}}
%%\DeclareRobustCommand\farcs{\ensuremath{.\!\!^{\prime\prime}}}
\DeclareRobustCommand\fp{\ensuremath{.\!\!^{\scriptscriptstyle\rm p}}}
%%\def\micron{\hbox{$\mu$m}}
\DeclareRobustCommand\micron{\ensuremath{\mu\mathrm{m}}}
\let\la=\lesssim                        
\let\ga=\gtrsim
%%\def\case#1#2{\hbox{$\frac{#1}{#2}$}}
%%\def\slantfrac#1#2{\hbox{$\,^#1\!/_#2$}}
%%\def\onehalf{\slantfrac{1}{2}}
%%\def\onethird{\slantfrac{1}{3}}
%%\def\twothirds{\slantfrac{2}{3}}
%%\def\onequarter{\slantfrac{1}{4}}
%%\def\threequarters{\slantfrac{3}{4}}
\DeclareRobustCommand\case[2]{\ensuremath{\frac{#1}{#2}}}
\DeclareRobustCommand\slantfrac[2]{\ensuremath{\,^#1\!/_#2}}
\DeclareRobustCommand\onehalf{\slantfrac{1}{2}}
\DeclareRobustCommand\onethird{\slantfrac{1}{3}}
\DeclareRobustCommand\twothirds{\slantfrac{2}{3}}
\DeclareRobustCommand\onequarter{\slantfrac{1}{4}}
\DeclareRobustCommand\threequarters{\slantfrac{3}{4}}
%%\def\ubvr{\hbox{$U\!BV\!R$}}            
%%\def\ub{\hbox{$U\!-\!B$}}               
%%\def\bv{\hbox{$B\!-\!V$}}               
%%\def\vr{\hbox{$V\!-\!R$}}               
%%\def\ur{\hbox{$U\!-\!R$}}
%    \end{macrocode}
% New commands for ``photometric'' colors 
%    \begin{macrocode}
\DeclareRobustCommand\ubvr{\ensuremath{U\!BV\!R}}            
\DeclareRobustCommand\Color[2]{%
  \ensuremath{%
    #1\!\ifx#1V\!\fi-\!#2
  }}
\DeclareRobustCommand\ub{\Color{U}{B}}
\DeclareRobustCommand\bv{\Color{B}{V}}
\DeclareRobustCommand\vr{\Color{V}{R}}
\DeclareRobustCommand\ur{\Color{U}{R}}
%% New version of the \ion command, that actually works! 
%    \end{macrocode}
% \begin{macro}{\ion}
% This version of the |\ion| command allows you to write either
% |\ion{N}{4}| or |\ion{N}{iv}| in order to get N~IV. 
%    \begin{macrocode}
\DeclareRobustCommand\ion[2]{%
  \mbox{#1\kern0.2em%
    \smaller\rmfamily%
    \edef\@tempa{\@car#2\@nil}%
    \ifcat1\@tempa%
    \@Roman{#2}%
    \else%
    \uppercase{#2}%
    \fi}}
\newcounter{IonStage}
\DeclareRobustCommand\plainion[2]{\setcounter{IonStage}{#2}#1 \Roman{IonStage}}
%%\def\ion#1#2{\ensuremath{\mathrm{#1}\;}{\protect\small\rm\@Roman{#2}}}
%    \end{macrocode}
% \end{macro}
%    \begin{macrocode}
\DeclareRobustCommand\nodata{$\cdots$}
%    \end{macrocode}
% These are needed by the |\ion| command. Isn't there an easier way of
% doing this?
%    \begin{macrocode}
\let\@small=\small
\let\@normalsize=\normalsize
\let\@footnotesize=\footnotesize
\let\@scriptsize=\scriptsize
\let\@tiny=\tiny
\let\@large=\large
\let\@Large=\Large
\let\@LARGE=\LARGE
\let\@huge=\huge
\let\@Huge=\Huge
\newcommand{\smaller}{\small}
\renewcommand{\normalsize}{\renewcommand{\smaller}{\small}\@normalsize}
\renewcommand{\small}{\renewcommand{\smaller}{\footnotesize}\@small}
\renewcommand{\footnotesize}{\renewcommand{\smaller}{\scriptsize}\@footnotesize}
\renewcommand{\scriptsize}{\renewcommand{\smaller}{\tiny}\@scriptsize}
\renewcommand{\tiny}{\renewcommand{\smaller}{\tiny}\@tiny}
\renewcommand{\large}{\renewcommand{\smaller}{\normalsize}\@large}  
\renewcommand{\Large}{\renewcommand{\smaller}{\large}\@Large}  
\renewcommand{\LARGE}{\renewcommand{\smaller}{\Large}\@LARGE}  
\renewcommand{\huge}{\renewcommand{\smaller}{\LARGE}\@huge}  
\renewcommand{\Huge}{\renewcommand{\smaller}{\huge}\@Huge}  
%</package|adshtml-style>
%    \end{macrocode}


% \section{Other files}

% \subsection{The driver and make files}
% \label{sec:driver_make}
% These are put here for completeness, although they still need to be
% distributed separately since otherwise the user won't be able to
% unpack this file into its constituent parts.
% 
% The |Makefile| has two principal rules: |docs|, and |macros|. The
% |docs| rule produces a PDF and gzipped PS file of the documentation.
% I still haven't perfected the combination of rules required to get
% it to run \LaTeX{} enough times in all circumstances. The |macros|
% rule runs \LaTeX{} on |rmaa.ins| in order to unpack |rmaa.cls| and
% all the other little files (including the |Makefile| and the
% |rmaa.ins| files themselves\dots 
% 
% First set up lists of various types of file. 
%    \begin{macrocode}
%<*makefile>
DOCFILES=rmaa.pdf
MACRO_LTX_FILES=rmaa.cls rmaacite.sty adshtml.sty
SHELLSCRIPTS=makeads.sh
MACRO_OTHER_FILES=adshtml.perl sample.latex2html-init
MAKEFILE=Makefile
MACROFILES=$(MACRO_LTX_FILES) $(SHELLSCRIPTS) $(MACRO_OTHER_FILES)         
SKELETONFILES=ADS.tex 
%    \end{macrocode}
% Install the files read only into the package directory, specified via a path relative to the source directory. \changes{4.0}{2019/11/14}{Reorganize folder layout so that hardly anything is in the top level.  That way, things are clearer on github.}
%    \begin{macrocode}
INSTALL=install 
INSTALL_DIR=../packages
%    \end{macrocode}
% The default rule is to make everything. Installation must be
% explicitly asked for.
%    \begin{macrocode}
all: docs macros

install: install-docs install-macros

%    \end{macrocode}
% Most files are installed read-only to prevent editing. Skeleton
% files must be writeable. Must be careful not to overwrite the
% Makefile in the parent directory. 
%    \begin{macrocode}
install-macros: macros
        $(INSTALL) -m 444 $(MACRO_LTX_FILES) $(MACRO_OTHER_FILES) $(INSTALL_DIR)
        $(INSTALL) -m 555 $(SHELLSCRIPTS) $(INSTALL_DIR)
        $(INSTALL) -m 644 $(SKELETONFILES) $(INSTALL_DIR)

install-docs: docs
        $(INSTALL) rmaa.pdf $(INSTALL_DIR)/rm-fulldocs.pdf

%    \end{macrocode}
% Make the PDF and gzipped PS files. 
%    \begin{macrocode}
docs: $(DOCFILES)

%    \end{macrocode}
% Make everything else if either the |.dtx| or |.ins| files have
% changed. 
%    \begin{macrocode}
macros: rmaa.dtx rmaa.ins
%    \end{macrocode}
% Ensure that all the macro files can be written to then recreate
% them. The |touch| is there in case any of the files don't exist
% yet. 
%    \begin{macrocode}
        touch $(MACROFILES) $(MAKEFILE)
        chmod u+w $(MACROFILES) $(MAKEFILE) rmaa.ins
        latex "%%
%% This is file `rmaa.ins',
%% generated with the docstrip utility.
%%
%% The original source files were:
%%
%% rmaa.dtx  (with options: `ins')
%% 
%% PLEASE DO NOT EDIT THIS FILE! It is generated automatically from
%% the file rmaa.dtx using the accompanying Makefile and rmaa.ins
%% Any hand-edited changes will be lost
%% 
%% A LaTeX document class for the RMxAA
%% 
%% Author: William Henney (w.henney@crya.unam.mx)
%% 
%% Latest version available from: http://www.crya.unam.mx/rmaa
%% 
%% Copyright (C) 1999-2010 Instituto de Astronomia, UNAM, Mexico
%% 
%% This program is free software; you can redistribute it and/or
%% modify it under the terms of the GNU General Public License
%% as published by the Free Software Foundation; either version 2
%% of the License, or (at your option) any later version.
%% 
%% This program is distributed in the hope that it will be useful,
%% but WITHOUT ANY WARRANTY; without even the implied warranty of
%% MERCHANTABILITY or FITNESS FOR A PARTICULAR PURPOSE.  See the
%% GNU General Public License for more details.
%% 
%% You should have received a copy of the GNU General Public License
%% along with this program; if not, write to the Free Software
%% Foundation, Inc., 59 Temple Place - Suite 330, Boston, MA  02111-1307, USA.





%%
%% Installation batch file for the RMAA LaTeX package
%%
\def\batchfile{rmaa.ins}
\input docstrip.tex
\preamble

PLEASE DO NOT EDIT THIS FILE! It is generated automatically from
the file rmaa.dtx using the accompanying Makefile and rmaa.ins
Any hand-edited changes will be lost

A LaTeX document class for the RMxAA

Author: William Henney (w.henney@crya.unam.mx)

Latest version available from: http://www.crya.unam.mx/rmaa

Copyright (C) 1999-2010 Instituto de Astronomia, UNAM, Mexico

This program is free software; you can redistribute it and/or
modify it under the terms of the GNU General Public License
as published by the Free Software Foundation; either version 2
of the License, or (at your option) any later version.

This program is distributed in the hope that it will be useful,
but WITHOUT ANY WARRANTY; without even the implied warranty of
MERCHANTABILITY or FITNESS FOR A PARTICULAR PURPOSE.  See the
GNU General Public License for more details.

You should have received a copy of the GNU General Public License
along with this program; if not, write to the Free Software
Foundation, Inc., 59 Temple Place - Suite 330, Boston, MA  02111-1307, USA.
\endpreamble
\askforoverwritefalse
\generate{%
  \file{rmaa.cls}{\from{rmaa.dtx}{package}}
  \file{rmaa.ins}{\from{rmaa.dtx}{ins}}
  \file{rmaacite.sty}{\from{rmaa.dtx}{rmaacite}}
  \file{adshtml.sty}{\from{rmaa.dtx}{adshtml-style}}
  \file{ADS.tex}{\from{rmaa.dtx}{ads-container}}
  }
\generate{\nopreamble\nopostamble%
  \file{Makefile}{\from{rmaa.dtx}{makefile}}
  \file{adshtml.perl}{\from{rmaa.dtx}{adshtml-perl}}
  \file{makeads.sh}{\from{rmaa.dtx}{ads-shell-script}}
  \file{sample.latex2html-init}{\from{rmaa.dtx}{latex2html-init}}
  }

%% Local Variables:
%% mode: LaTeX
%% mode: Font-lock
%% End:









\endinput
%%
%% End of file `rmaa.ins'.
"
%    \end{macrocode}
% \TeX{} automatically adds the |.tex| extension to extensionless
% files. We fix this and take the opportunity to convert
% back to TABs all that it has lovingly converted to spaces. 
%    \begin{macrocode}
        cp -f $(MAKEFILE) $(MAKEFILE).BCK && unexpand $(MAKEFILE).tex >| $(MAKEFILE) && rm $(MAKEFILE).tex
%    \end{macrocode}
% Make the shell scripts executable and everything unwriteable (so
% they don't get hand-edited by mistake). However, undo this for the
% files that the user is \emph{supposed} to modify. 
%    \begin{macrocode}
        chmod a+x $(SHELLSCRIPTS)
        chmod a-w $(MACROFILES) $(MAKEFILE)
        chmod u+w $(SKELETONFILES)

%    \end{macrocode}
% We now use PDF\LaTeX{} for this.
%    \begin{macrocode}
rmaa.idx: rmaa.dtx
        @echo "REGENERATING .idx FILE"
        pdflatex "\nonstopmode% \def\fileversion{4.0} \def\filedate{2019/11/12}
% \MakeShortVerb{\|}
% \title{A \LaTeX{} document class for the Revista Mexicana de
% Astronom\'\i a y Astrof\'\i sica\thanks{This file has version
% \fileversion{} last revised \filedate{}.}}
% \author{William J. Henney\thanks{w.henney@irya.unam.mx}}
% \date{\filedate}
% \maketitle
% \paragraph{Abstract}
% This document describes the \texttt{rmaa} \LaTeX{}
% document class for the journal ``Revista Mexicana de
% Astronom\'\i a y Astrof\'\i sica''. The macros can be used for 
% \begin{enumerate}
% \item Preparing author manuscripts for the journal (including
% conference proceedings to be published in the ``Serie de
% Conferencias'')
% \item Typesetting an entire volume of the journal by inputting the
% individual article files from a container file. Facilities are
% included for automatic generation of TOC and Author Index. 
% \end{enumerate}
% 
% \tableofcontents
% \section{The documentation driver file}
% This is the part that generates the documentation that you may or
% may not be reading, which is done with the help of the Makefile
% discussed in section~\ref{sec:driver_make}. 
%    \begin{macrocode}
%<*driver>
\documentclass{ltxdoc}
\usepackage{hyperref}
\usepackage{palatino}
\EnableCrossrefs \CodelineIndex \RecordChanges
\setlength{\parindent}{0pt}
\def\MacroFont{\footnotesize\tt}
\begin{document} \DocInput{rmaa.dtx} \PrintIndex \PrintChanges \end{document}
%</driver>
%    \end{macrocode}

% \section{The implementation}
% 
% \changes{3.24}{2007/09/14}{Added user access to the macro version
% and date via \cs{rmVersion} and \cs{rmDate}}
% 
%    \begin{macrocode}
%<*package>
\newcommand{\rmaa@version}{3.26}
\newcommand{\rmaa@date}{2011/08/15}
\newcommand\rmVersion{\rmaa@version}
\newcommand\rmDate{\rmaa@date}
\NeedsTeXFormat{LaTeX2e}[1995/06/01]
%    \end{macrocode}
% \changes{4.0}{2019/11/02}{Change abstract to appear before
% resumen. Change RevMexAA to RMxAA for consistency with ADS.}
%
% \changes{3.22}{2007/03/22}{Allow proceedings articles to have
%   received/accepted dates} \changes{3.21}{2006/05/15}{Changes for
%   Silvia - abstracts are now shipped out to ADS.}
% \changes{3.19}{2006/01/19}{Changes for Silvia - remove adresses from
%   end, new abstract style} \changes{2.2beta}{2001/10/30}{ Fixed bug
%   whereby twocolumn option was not passed through to book
%   class. Finally sorted out the ``label on last page'' issue. }
% \changes{2.1}{2001/01/23}{ Fixed some inconsistencies in
%   \cs{nocite} and \cs{bibitem} uniqification Fixed the definition of
%   \cs{ion} so it will work in math mode and in contexts where the
%   typesize is different from normal (e.g., title, fig captions).
%   Modified the definition of \cs{@caption}, so we don't write any
%   list-of-figures data (we never use it anyway, and it just causes
%   problems with fragile commands in fig captions). }
% \changes{2.0}{2000/11/30}{ Added `newstyle' option, which should be
%   used in all new proceedings volumes (wider abstract, smaller
%   typeface in fig captions, etc.) Added commands for formatting
%   table captions. New demonstation document rmpostsamp.tex for
%   1-page posters. This will be the last release under a PRI
%   government (for a while...)}  \changes{1.9}{2000/06/06}{
%   Corrected position of page number in footer of articles that open
%   on even pages. Ensured that first article in each session opens on
%   an odd page, as do list of abstracts and author index. Added extra
%   space in TOC before the backmatter. Added command for editors,
%   \cs{BoldSessionTitles}, that makes produces bold titles for
%   session in the TOC.}  \changes{1.8}{2000/05/16}{ Fixed bug in
%   last page label with \cs{suppressfulladdresses} }
% \changes{1.7}{2000/05/11}{ Fixed more little things. Page labels
%   now work, even if you have no bibliography. You can output the
%   full addresses by hand with the \cs{outputfulladdresses}
%   command. Editors' Guide updated but still not finished. Revamped
%   splitbook.sh now produces PS file of each article for web page and
%   also an EPS file of each page for sending to the printers.}
% \changes{1.6}{2000/04/15}{ More cleaning up: full stops after
%   titles in index, running headers and footers in frontmatter,
%   fontsize in title of index, etc} \changes{1.5}{2000/04/15}{
%   Eliminated superscripts from titles in TOC and index}
% \changes{1.4}{2000/04/14}{ Added copyright and license info (GPL)}
% \changes{1.3}{2000/03/01}{ Corrected \cs{figcaption} to do the
%   caption formatting right Corrected font size of sec headings in
%   front matter} \changes{1.2}{2000/02/16}{ Added \cs{plate}
%   command} \changes{1.1}{2000/02/14}{ Added manuscript option,
%   corrected bugs} \changes{1.0}{2000/02/08}{ Now does a whole book
%   - lots of stuff tidied up} \changes{0.7}{1999/10/04}{ Increased
%   spacing between sections in single column styles. Suppressed
%   floats at bottom of pages (didn't work though) }
% \changes{0.6}{1999/09/24}{ Tweaked \cs{adjustfinalcols} (still not
%   perfect) Added proceedings option Added onecolumn/twocolumn
%   options (orthogonal to journal/proceedings) }
% \changes{0.5}{1999/09/23}{ Added \cs{fulladdresses} and
%   \cs{adjustfinalcols} (just about working)}
% \changes{0.4}{1999/09/22}{ Minor stylistic changes}
% \changes{0.3}{1999/06/28}{ Better mechanism for alternative
%   affiliations Added `guide' option for user guide}
% \changes{0.2}{1999/05/27}{ All basic functionality present, I
%   think} \changes{0.1}{1999/05/26}{ Just about working - still
%   missing Fig/Tab/Refs} \changes{0.0}{1999/05/21}{ So bored today,
%   decided to rewrite RMAA LaTeX macros from scratch}
%    \begin{macrocode}
\ProvidesClass{rmaa}[\rmaa@date\space v\rmaa@version\space RMxAA article
class]
%    \end{macrocode}
% \subsection{Option Handling}
% \label{sec:options}
% Declarations that are needed in options
% \changes{3.05}{2003/10/22}{Added natbib option}
% \changes{3.06}{2003/01/03}{Natbib option turned on permanently}
% \changes{3.07}{2003/01/11}{Added proof option}
% \changes{3.14}{2003/01/11}{Option to turn off natbib}
% \changes{3.18}{2003/10/22}{trying out doctex mode} 
% \changes{3.23}{2007/09/09}{Manuscript is now on by default.}
% \changes{3.26}{2011/08/15}{New option for review articles}
%    \begin{macrocode}
\newif\if@reviewarticle
\@reviewarticlefalse
\newif\if@proof
\@prooffalse
\newif\if@journal
\@journaltrue
\newif\if@printdates
\@printdatesfalse
\newif\if@proceedings
\@proceedingsfalse
\newif\if@manuscript
\@manuscripttrue
\newif\if@bigfont
\@bigfonttrue
\newif\if@guide
\@guidefalse
\newif\if@preprint
\@preprintfalse
\newif\if@rmtwocol
\@rmtwocoltrue
\newif\if@biblist
\@biblistfalse
\newif\if@book
\@bookfalse
\newif\if@natbib
\@natbibtrue
\newif\if@rmcite
\@rmcitefalse
\newif\if@titlevisible
\@titlevisibletrue
\newif\if@abstractpresent
\@abstractpresentfalse
\newif\if@printbanner
\@printbannerfalse
\newif\if@suppressfulladdresses
\@suppressfulladdressesfalse
\newif\if@newstyle
\@newstylefalse
\newif\if@debug
\@debugfalse
\newif\if@boldkeywords
\@boldkeywordsfalse
%    \end{macrocode}
% The option handling---mainly setting various booleans.
% \changes{2.2beta11}{2001/11/26}{Made twocol and
% newstyle the default for proceedings.}
% \changes{3.23}{2007/09/09}{Manuscript is now 10 pt instead of 12pt}
% \changes{3.23}{2007/09/10}{Full addresses are now automatically
% suppressed for proceedings style.}
% \changes{3.24}{2007/09/14}{New option \texttt{boldkeywords} that
% does the keywords in the old ugly way. Now off by default.}
%    \begin{macrocode}
\DeclareOption{manuscript}{\typeout{Manuscript for referee} 
  \@manuscripttrue\@journalfalse\@rmtwocolfalse
  \PassOptionsToClass{10pt}{article}}
\DeclareOption{bigfont}{\typeout{Big font manuscript requested by Gary
  Ferland} 
  \@manuscripttrue\@journalfalse\@rmtwocolfalse\@bigfonttrue
  \PassOptionsToClass{11pt}{article}}
\DeclareOption{journal}{\typeout{Main journal option}%
  \@journaltrue\@rmtwocoltrue\@printdatestrue\@manuscriptfalse
  \PassOptionsToClass{10pt}{article}}
\DeclareOption{guide}{\typeout{User guide option}%
  \@guidetrue\@journalfalse\@rmtwocoltrue\@manuscriptfalse
  \PassOptionsToClass{10pt}{article}}
\DeclareOption{preprint}{\typeout{Preprint option}%
  \@preprinttrue\@manuscriptfalse\PassOptionsToClass{10pt}{article}}
\DeclareOption{proof}{\typeout{Author proof option}%
  \@prooftrue\@manuscriptfalse}
\DeclareOption{proceedings}{\typeout{Proceedings option}%
  \@proceedingstrue\@journalfalse\@rmtwocoltrue\@newstyletrue
  \@suppressfulladdressestrue
  \@printdatesfalse\@manuscriptfalse\PassOptionsToClass{10pt}{article}} 
\DeclareOption{printdates}{\typeout{Force printing of received/accepted dates}%
  \@printdatestrue} 
\DeclareOption{book}{\typeout{Whole issue option}%
  \@booktrue} 
\DeclareOption{nonatbib}{\typeout{Turn off natbib support (required
    for compatibility with old rmaacite package).}%
  \@natbibfalse} 
\DeclareOption{newstyle}{\typeout{Use new layout style}%
  \@newstyletrue}
\DeclareOption{oldstyle}{\typeout{Use old layout style}%
  \@newstylefalse}
\DeclareOption{debug}{\typeout{Loads of annoying messages}%
  \@debugtrue}
%    \end{macrocode}
% These two need to be declared explicitly so that |\ExecuteOptions|
% actually does something.
%    \begin{macrocode}
\DeclareOption{twocolumn}{\@rmtwocoltrue}
\DeclareOption{onecolumn}{\@rmtwocolfalse}
\DeclareOption{twoside}{\PassOptionsToClass{\CurrentOption}{article}}
\DeclareOption{boldkeywords}{\@boldkeywordstrue}
%    \end{macrocode}
% Bold, upper-case keywords used to be the default. Now they must be
% turned on explicitly with the \texttt{boldkeywords} option. 
%    \begin{macrocode}
\DeclareOption{boldkeywords}{\@boldkeywordstrue}
%    \end{macrocode}
% Pass any other option along to the \texttt{article} class.
%    \begin{macrocode}
\DeclareOption*{\PassOptionsToClass{\CurrentOption}{article}}
%    \end{macrocode}
% Default options are to be two-sided and based on the standard
% \LaTeX{} \texttt{article} class,
% unless we are processing the entire book, in which case we base
% ourselves on the standard \texttt{book} class. Default is to use
% manuscript option unless otherwise specified. 
%    \begin{macrocode}
\ExecuteOptions{twoside,manuscript} \ProcessOptions\relax
\if@book
\if@rmtwocol\LoadClass[twocolumn]{book}\else\LoadClass{book}\fi
\newcommand{\abstract}{\relax}
\newcommand{\@maketitle}{\relax}
\else
\if@rmtwocol\LoadClass[twocolumn]{article}\else\LoadClass{article}\fi
\fi
%    \end{macrocode}
% Attempt to use natbib for citations 02 Jan 2003. 
% \changes{3.05}{2003/01/02}{If the natbib option is used, then set
% things up so that natbib's version of the bibliography is identical
% to the standard one.}
% \changes{3.09}{2003/01/23}{Explicitly require a recent version of
% natbib, so that we get a decent error message if the package is too
% old.} 
% \changes{3.23}{2007/09/09}{Now allow Bibtex to work properly} 
% \changes{3.25}{2009/10/01}{Change \cs{newcommand} tp \cs{def} when setting \cs{bibfont}. This allows it to work with all versions of natbib.} 
% \changes{3.25}{2010/02/06}{Updated copyright year.} 
% \changes{3.25}{2010/02/06}{Updated email and web addresses to \texttt{crya.unam.mx}.} 
%    \begin{macrocode}
\if@natbib
\RequirePackage{natbib}[1999/02/25]
\bibpunct{(}{)}{;}{a}{}{,}
\def\bibfont{\small}
\renewcommand\bibsection{\section*{REFERENCES}}
\setlength\bibsep{0pt}
\setlength\bibhang{1.5em}
\bibliographystyle{rmaa}
\fi
%    \end{macrocode}
% We want the \texttt{graphicx} package for the extended syntax of
% |\includegraphics|. We want the \texttt{ifthen} package for using in
% these macros (although we don't much yet). We explicitly load
% \texttt{babel} since some \LaTeX{} distributions messed up the
% hyphenation if we didn't.  
%    \begin{macrocode}
\RequirePackage{graphicx}
\RequirePackage{ifthen}
\IfFileExists{babel.sty}{%
%   \let\@restoreindent=\@afterindentfalse
  \RequirePackage[english]{babel}
%   \@restoreindent
  %% reinstate the definition of \@roman so as not to break .idx files
  \@ifundefined{org@@roman}{\relax}{\renewcommand{\@roman}[1]{\org@@roman{##1}}}
  }{%
  \typeout{Package babel.sty not found - don't worry about this}}
%    \end{macrocode}
% Try to load the \texttt{dcolumn} and \texttt{booktabs} packages if
% they can be found. These can produce nicer-looking tables. 
%    \begin{macrocode}
\IfFileExists{dcolumn.sty}{%
  \RequirePackage{dcolumn}
  }{%
  \typeout{Package dcolumn.sty not found - you will not be able to
    align columns on decimal points in tables}}
\IfFileExists{booktabs.sty}{%
  \RequirePackage{booktabs}
  }{%
  \typeout{Package booktabs.sty not found - your tables will not look
    as pretty as they ought too.}
  }
\IfFileExists{makeidx.sty}{%
  \RequirePackage{makeidx}
}{%
  \typeout{Package makeidx.sty not found - Author Index may not work.}
  }
%% commented out for moment since requires more modern version of LaTeX
%% \RequirePackage[latin1]{inputenc}
%    \end{macrocode}
% \changes{4.0}{2019/11/13}{Make \texttt{hyperref} and
% \texttt{etoolbox} into hard requirements, since we now use them
% extensively.}  We use hyperref to get internal and external links.
% But I don't like the colored boxes so I turn those off.  We use
% etoolbox for patching and robustifying at the moment.  In the
% future, I want to port a lot of other stuff over to it.  For
% instance, replace the ifthen package and the TeX ifs too.
%
% We target the latest version (1.7) of PDF, rather than the default
% of 1.5.  This reduces the chance of getting warnings from included
% graphics.  This is not strictly related to hyperref, but we put it
% here anyway.
%    \begin{macrocode}
\RequirePackage{etoolbox}
\RequirePackage{hyperref}
\hypersetup{hidelinks}
\pdfminorversion=7
%    \end{macrocode}
% We patch the \cs{@combinedblfloats} following a suggestion by Heiko
% Oberdiek at \url{https://tex.stackexchange.com/a/249743/3843}, which
% avoids problems when hyperlinks are split across
% pages. \changes{4.0}{2019/11/12}{Fix for issue where breaking
% hyperlink across pages gave an error.}
%    \begin{macrocode}
\urlstyle{same}
\newcommand\email{\begingroup \urlstyle{rm}\Url}
\patchcmd\@combinedblfloats{\box\@outputbox}{\unvbox\@outputbox}{}{%
  \errmessage{\noexpand\@combinedblfloats could not be patched}%
}%
%    \end{macrocode}
% \subsection{Other setup stuff}\label{sec:setup}
% Some declarations: 
%    \begin{macrocode}
\newif\if@needaddresses
\@needaddressestrue
\newif\if@frontmatter
\@frontmatterfalse
\newif\if@inpreamble
\@inpreambletrue
\newif\if@firstauthor
\@firstauthortrue
\newif\if@firstaffil
\@firstaffiltrue
\newif\if@firstaltaffil
\@firstaltaffiltrue
\newif\if@firstkwd
\@firstkwdtrue
\newif\if@noexpandADS
\@noexpandADSfalse
\newif\if@sessiontitlepages
\@sessiontitlepagesfalse
%    \end{macrocode}
% \begin{macro}{\@rmaaEmpty}
% \begin{macro}{\@rmaaArg}
% These are used in various commands to check if arg was left blank
%    \begin{macrocode}
\newcommand{\@rmaaEmpty}{}
\newcommand{\@rmaaArg}{}
%    \end{macrocode}  
% \end{macro}
% \end{macro}
%
% \subsection{Compatibility with old \LaTeX{} versions}
% \label{sec:compat}
% This package is based on the 1998/12/01 release of \LaTeX{}.
% Older releases may not define all the commands that we need.
% Luckily, we can use |\providecommand| to define these if necessary. 
% The following commands are not defined in the 1994/06/01 release. 
%
% \begin{macro}{\MakeUpperCase}
% \begin{macro}{\textsuperscript}
% \begin{macro}{\textbackslash}
% User commands:
%    \begin{macrocode}
\providecommand{\MakeUpperCase}[1]{\uppercase{#1}}
\providecommand{\textsuperscript}[1]{\ensuremath{^\textrm{#1}}}
\providecommand{\textbackslash}{\ensuremath{\backslash}}
%    \end{macrocode}  
% \end{macro}
% \end{macro}
% \end{macro}
% \begin{macro}{\protected@xdef}
% \begin{macro}{\restore@protect}
% \begin{macro}{\@unexpandable@protect}
% \begin{macro}{\hb@xt@}
% \LaTeX{} internal commands:
%    \begin{macrocode}
\providecommand{\protected@xdef}{%
        \let\@@protect\protect
        \let\protect\@unexpandable@protect
        \afterassignment\restore@protect
        \xdef}
\providecommand{\restore@protect}{\let\protect\@@protect}
\providecommand{\@unexpandable@protect}{\noexpand\protect\noexpand}
\providecommand{\hb@xt@}{\hbox to}
%    \end{macrocode}
% \end{macro}
% \end{macro}
% \end{macro}
% \end{macro}
%
% \subsection{Page layout}\label{sec:page-layout}
%
% The dimensions of the text were defined by Elizabeth. Text width is
% 169\,mm with equal margins on the two sides: 
%    \begin{macrocode}
\setlength{\textwidth}{169mm}
\setlength{\@tempdima}{\paperwidth}
\addtolength{\@tempdima}{-\textwidth}
\setlength{\oddsidemargin}{0.5\@tempdima}
\addtolength{\oddsidemargin}{-1in}
\setlength{\evensidemargin}{\oddsidemargin}
%    \end{macrocode}
% Text height is 225\,mm, top margin is 14\,mm. 
%    \begin{macrocode}
\setlength{\textheight}{225mm} 
\setlength{\topmargin}{-1in} 
\addtolength{\topmargin}{14mm}
%    \end{macrocode}
% Space between columns in two-column styles, header, footer,
% paragraph indentation:
%    \begin{macrocode}
\setlength{\columnsep}{9.5mm}
\setlength{\headsep}{0.7cm}
\setlength{\footskip}{0.75cm}
\setlength{\parindent}{1.5em}
%    \end{macrocode}
% 
% \begin{macro}{\@set@scalable@length}
% Try out a new scheme for dealing with various spacing dimensions and
% allowing the user to easily modify them. The command
% \begin{quote}
% |\set@scalable@length|\marg{name}\marg{defaultlength}\marg{stretch}\marg{shrink}
% \end{quote}
% does a 
% |\setlength|\marg{\cs{name}}\marg{defaultlength plus stretch minus shrink}
% and also defines the user command
% \textbackslash\meta{name}|Scale|\marg{factor} that can 
% be used to rescale length by \meta{factor}. Additionally, now
% defines the package command
% \textbackslash\meta{name}|@set@scale|\marg{FACTOR} 
% that also rescales the length by \meta{FACTOR}, so that the length
% becomes \meta{factor} $\times$ \meta{FACTOR} $\times$
% \meta{defaultlength}.   
%    \begin{macrocode}
\newcommand\@set@scalable@length[4]{%
  %% arg #1: NAME of length without leading `\'
  %%
  %% arg #2: default VALUE  
  \expandafter\newlength\csname @#1@default\endcsname
  \expandafter\setlength\csname @#1@default\endcsname{#2}%
  \expandafter\newlength\csname @@#1@@default\endcsname
  \expandafter\setlength\csname @@#1@@default\endcsname{#2}%
  %% arg #3: default STRETCH  
  \expandafter\newlength\csname @#1@stretch\endcsname
  \expandafter\setlength\csname @#1@stretch\endcsname{#3}%
  \expandafter\newlength\csname @@#1@@stretch\endcsname
  \expandafter\setlength\csname @@#1@@stretch\endcsname{#3}%
  %% arg #4: default SHRINK  
  \expandafter\newlength\csname @#1@shrink\endcsname
  \expandafter\setlength\csname @#1@shrink\endcsname{#4}%
  \expandafter\newlength\csname @@#1@@shrink\endcsname
  \expandafter\setlength\csname @@#1@@shrink\endcsname{#4}%
  %%
  %% User-defined scale \@NAME@scale is unity by default
  \expandafter\gdef\csname @#1@scale\endcsname{1}%
  %% Package-defined scale \@NAME@SCALE is unity by default
  \expandafter\gdef\csname @#1@SCALE\endcsname{1}%
  %%
  %% internal command \@reset@NAME should expand to 
  %% \setlength\@@NAME@@default{\@NAME@SCALE\@NAME@default}
  %% \setlength\NAME{\@NAME@scale\@@NAME@@default}
  %% in other words: scale x SCALE x default 
  \expandafter\gdef\csname @reset@#1\endcsname{%
    \expandafter\setlength\csname @@#1@@default\endcsname{%
      \csname @#1@SCALE\endcsname \csname @#1@default\endcsname
      }
    \expandafter\setlength\csname @@#1@@stretch\endcsname{%
      \csname @#1@SCALE\endcsname \csname @#1@stretch\endcsname
      }
    \expandafter\setlength\csname @@#1@@shrink\endcsname{%
      \csname @#1@SCALE\endcsname \csname @#1@shrink\endcsname
      }
    \expandafter\setlength\csname #1\endcsname{%
      \csname @#1@scale\endcsname \csname @@#1@@default\endcsname
      \@plus
      \csname @#1@scale\endcsname \csname @@#1@@stretch\endcsname
      \@minus
      \csname @#1@scale\endcsname \csname @@#1@@shrink\endcsname
      }}
  %%
  %% User command \NAMEScale should expand to 
  %% {\renewcommand\@NAME@scale{#1}\@reset@NAME}
  \expandafter\gdef\csname #1Scale\endcsname##1{%
    \expandafter\gdef\csname @#1@scale\endcsname{##1}%
    \csname @reset@#1\endcsname}
  %%
  %% Package command \NAME@set@SCALE should expand to 
  %% {\renewcommand\@NAME@SCALE{#1}\@reset@NAME}
  \expandafter\gdef\csname #1@set@SCALE\endcsname##1{%
    \expandafter\gdef\csname @#1@SCALE\endcsname{##1}%
    \csname @reset@#1\endcsname}
  %%
  %% initialize at default values with \@reset@NAME
  \csname @reset@#1\endcsname
  }
%    \end{macrocode}
% \end{macro}
% \begin{macro}{\parskip}
% First use of this new infrastructure: set interparagraph space to be
% zero but with a small amount of stretchiness.
%    \begin{macrocode}
%   \ifthenelse{\boolean{@manuscript}}{%
%     \@set@scalable@length{parskip}{0.5\baselineskip}{0.3\p@}{0.1\p@}
%   }{
    \@set@scalable@length{parskip}{0\p@}{0.3\p@}{0\p@}
%   }
%    \end{macrocode}
% \end{macro}
% 
% \begin{macro}{\abovesectionskip}
% \changes{4.0}{2019/11/12}{Use negative values for skip so that first paragraph after section headings is not indented.  This is something that I have been trying to fix for years!  It turns out that the anser was in the documentation for the standard command \cs{@startsection.}}
% \begin{macro}{\belowsectionskip}
% \begin{macro}{\aboveparaskip}
% \begin{macro}{\belowparaskip}
% Define some utility rubber lengths for use in
% section~\ref{sec:section-commands}. 
%    \begin{macrocode}
\newlength{\abovesectionskip}
\newlength{\belowsectionskip}
\newlength{\aboveparaskip}
\newlength{\belowparaskip}
%    \end{macrocode}
% The default values for these lengths are those traditionally used
% for the onecolumn style. 
%    \begin{macrocode}
\@set@scalable@length{abovesectionskip}{-3.5ex}{-4ex}{-2ex}
\@set@scalable@length{belowsectionskip}{2.3ex}{0ex}{2ex}
\@set@scalable@length{aboveparaskip}{3.25ex}{2ex}{0.5ex}
\@set@scalable@length{belowparaskip}{-1em}{\z@}{\z@}
%    \end{macrocode}
% \end{macro}\end{macro}\end{macro}\end{macro}
% We also make scalable some of the standard float parameters. These
% are the same defaults as in the standard article classes, but
% slightly more stretchy. For the twocolumn style, they are currently
% halved, see below. Note that I have also lumped |\abovecaptionskip|
% in with these, since it plays a similar role. There is also a
% |\belowcaptionskip| but it is set to zero in the standard classes. 
%    \begin{macrocode}
\@set@scalable@length{floatsep}{12\p@}{4\p@}{4\p@}
\@set@scalable@length{textfloatsep}{20\p@}{4\p@}{8\p@}
\@set@scalable@length{intextsep}{12\p@}{4\p@}{4\p@}
\@set@scalable@length{dblfloatsep}{12\p@}{4\p@}{4\p@}
\@set@scalable@length{dbltextfloatsep}{20\p@}{4\p@}{8\p@}
\@set@scalable@length{abovecaptionskip}{10\p@}{\z@}{\z@}
%    \end{macrocode}
% \begin{macro}{\@before@title@skip}         
% \begin{macro}{\@before@author@skip}        
% \begin{macro}{\@between@author@skip}       
% \begin{macro}{\@before@abstract@skip}      
% \begin{macro}{\@after@abstracthead@skip}   
% \begin{macro}{\@between@abstract@skip}     
% \begin{macro}{\@before@kwd@skip}           
% \begin{macro}{\@before@dates@skip}           
% Spacing parameters used in |\@maketitle|. These are not changed by
% default except by the user, so they don't need to go into
% |\@package@default@lengths|.  
%    \begin{macrocode}
\newlength{\@before@title@skip}         
\newlength{\@before@author@skip}        
\newlength{\@between@author@skip}       
\newlength{\@before@abstract@skip}      
\newlength{\@after@abstracthead@skip}   
\newlength{\@between@abstract@skip}     
\newlength{\@before@kwd@skip}           
\newlength{\@before@dates@skip}           
\@set@scalable@length{@before@title@skip}{2em}{\z@}{\z@}
\@set@scalable@length{@before@author@skip}{1.5em}{\z@}{\z@}
\@set@scalable@length{@between@author@skip}{0.5em}{\z@}{\z@}
\@set@scalable@length{@before@abstract@skip}{3ex}{1ex}{0.5ex}
\@set@scalable@length{@after@abstracthead@skip}{1ex}{\z@}{\z@}
\@set@scalable@length{@between@abstract@skip}{2ex}{1ex}{0.5ex}
\@set@scalable@length{@before@kwd@skip}{1ex}{0.5ex}{0.2ex}
\@set@scalable@length{@before@dates@skip}{2ex}{0.5ex}{0.2ex}
%    \end{macrocode}
% \end{macro}\end{macro}\end{macro}\end{macro}
% \end{macro}\end{macro}\end{macro}\end{macro}
% 
% Spacing above and below displayed equations also wants to be done
% like this, but it is more complicated since the standard article
% classes make these change with the fontsize. Something to do later,
% maybe. 
% 
%    \begin{macrocode}
%    \end{macrocode}
% \begin{macro}{\@package@default@lengths}
% The value actually used for these lengths is controlled by the
% \meta{SCALE} parameter, which depends on the style option that is in
% effect. This is set using the |\@package@default@lengths| command,
% which must be called every time the style changes (e.g., changing
% from onecolumn to twocolumn). 
%    \begin{macrocode}
\newcommand{\@package@default@lengths}{%
  \ifthenelse{\boolean{@rmtwocol}}{%
    %% twocolumn mode
    \abovesectionskip@set@SCALE{0.5}
    \belowsectionskip@set@SCALE{0.5}
    \aboveparaskip@set@SCALE{0.5}
    \belowparaskip@set@SCALE{1}
    \floatsep@set@SCALE{0.5}
    \textfloatsep@set@SCALE{0.5}
    \intextsep@set@SCALE{0.5}
    \dblfloatsep@set@SCALE{0.5}
    \dbltextfloatsep@set@SCALE{0.5}
    }{%
    %% onecolumn mode
    \abovesectionskip@set@SCALE{1}
    \belowsectionskip@set@SCALE{1}
    \aboveparaskip@set@SCALE{1}
    \belowparaskip@set@SCALE{1}
    \floatsep@set@SCALE{1}
    \textfloatsep@set@SCALE{1}
    \intextsep@set@SCALE{1}
    \dblfloatsep@set@SCALE{1}
    \dbltextfloatsep@set@SCALE{1}
    }%
  }
\@package@default@lengths
%    \end{macrocode}
% \end{macro}
% \begin{macro}{\RescaleLengths}
% \begin{macro}{\RescaleSecLengths}
% \begin{macro}{\RescaleFloatLengths}
% \begin{macro}{\RescaleTitleLengths}
% User commands that will fiddle with all the scalable lengths. 
%    \begin{macrocode}
\newcommand{\RescaleLengths}[1]{%
  \RescaleSecLengths{#1}\RescaleFloatLengths{#1}\RescaleTitleLengths{#1}}
\newcommand{\RescaleSecLengths}[1]{%
%% sections
  \abovesectionskipScale{#1}
  \belowsectionskipScale{#1}
  \aboveparaskipScale{#1}
  }
\newcommand{\RescaleFloatLengths}[1]{%
%% floats
  \floatsepScale{#1}
  \textfloatsepScale{#1}
  \intextsepScale{#1}
  \dblfloatsepScale{#1}
  \dbltextfloatsepScale{#1}
  \abovecaptionskipScale{#1}
  }
\newcommand{\RescaleTitleLengths}[1]{%
%% maketitle
  \@before@title@skipScale{#1}           
  \@before@author@skipScale{#1}        
  \@between@author@skipScale{#1}       
  \@before@abstract@skipScale{#1}      
  \@after@abstracthead@skipScale{#1}   
  \@between@abstract@skipScale{#1}     
  \@before@kwd@skipScale{#1}           
  \@before@dates@skipScale{#1}           
  }
%    \end{macrocode}
% \end{macro}\end{macro}\end{macro}\end{macro}
% \changes{2.2beta14}{2002/03/17}{Finished implementation of new
% infrastructure for scalable lengths. Converted the above/below
% section skips to use it. Also, all the floatsep parameters. Also,
% all the vertical lengths that are used by |\@maketitle|. }

% \begin{macro}{\MSextramargin}
% For the \texttt{manuscript} option we have an extra 0.3\,inch
% vertical margin and 1\,inch horizontal margin for marginal notes (by
% hand or with |\marginpar|. 
%
% \changes{3.23}{2007/09/09}{Reduce baselinestretch in
%   \texttt{manuscript} style to 1.2. No longer use raggedright for
%   manuscript.}
%
% \changes{3.24}{2007/09/14}{Further reduction of baselinestretch in
%   \texttt{manuscript} style to 1.1 since we now have wider margins.}
%
% \changes{3.24}{2007/09/14}{Split \texttt{manuscript}'s extra margin
%   into separate horizontal and vertical parts.}
% 
%    \begin{macrocode}
\newlength{\MSextraHmargin}\newlength{\MSextraVmargin}
\setlength{\MSextraHmargin}{1in} 
\setlength{\MSextraVmargin}{0.3in} 
\if@manuscript
   \renewcommand{\baselinestretch}{1.1}
   \addtolength{\textwidth}{-2\MSextraHmargin}
   \addtolength{\oddsidemargin}{\MSextraHmargin}
   \setlength{\evensidemargin}{\oddsidemargin}
   \addtolength{\textheight}{-2\MSextraVmargin}
   \addtolength{\topmargin}{\MSextraVmargin}
%    \raggedright
\fi
%    \end{macrocode}
% \end{macro}
% Float parameters: don't allow bottom floats by default. 
%    \begin{macrocode}
\renewcommand{\fps@figure}{tp}
\renewcommand{\fps@table}{tp}
%    \end{macrocode}
%
% \begin{macro}{\rmaa@AuthorWidthFrac}
% \begin{macro}{\rmaa@AbstractWidthFrac}
% These parameters control what fraction of the total text width can
% be used by the list of authors and the abstract/resumen,
% respectively.  
%    \begin{macrocode}
\newcommand{\rmaa@AuthorWidthFrac}{1.0}
\newcommand{\rmaa@AbstractWidthFrac}{0.75}
%    \end{macrocode}
% Use a wide abstract to save space in the new style.
%    \begin{macrocode}
\if@newstyle\renewcommand{\rmaa@AbstractWidthFrac}{1.0}\fi
%    \end{macrocode}
% \end{macro}
% \end{macro}
%
% Some things to consider in future versions. Trying to avoid too many
% widows, etc. Apparently, the following will work best with
% \cs{flushbottom}. 
%    \begin{macrocode}
%%\widowpenalty=10000
%%\clubpenalty=10000
%    \end{macrocode}
% Also, should perhaps try getting rid of underfull/overfull box warnings. 
% \changes{3.10}{2003/01/24}{Set large values for \cs{hbadness},
% \cs{vbadness}, \cs{hfuzz}, \cs{vfuzz}. To reduce amount of stuff
% that gets written to the terminal every time you compile something.}
%    \begin{macrocode}
\hbadness=10000
\vbadness=10000
\setlength\hfuzz{\maxdimen}
\setlength\vfuzz{\maxdimen}
%    \end{macrocode}
% \subsection{Various text styles for different contexts}
% \label{sec:text-styles} 
% \begin{macro}{\rmaa@AffilStyle}
% \begin{macro}{\rmaa@PagesStyle}
% \begin{macro}{\rmaa@YearStyle}
% \begin{macro}{\rmaa@VolumeStyle}
% \begin{macro}{\rmaa@VolumenStyle}
% \begin{macro}{\rmaa@JournalNameStyle}
% \begin{macro}{\rmaa@ReceivedStyle}
% \begin{macro}{\rmaa@AbstractHeaderStyle}
% \begin{macro}{\rmaa@JournalNameStyle}
% Style commands for the different fields used by |\maketitle| in
% constructing the title page of the article. These have the form
% |\rmaa@XYZStyle{TEXT}|. I'm not sure these were really a good idea,
% but never mind \dots 
%    \begin{macrocode}
\newcommand{\rmaa@AffilStyle}[1]{{\normalsize\rmfamily#1}}
\newcommand{\rmaa@PagesStyle}[1]{{\normalsize\rmfamily#1}}
\newcommand{\rmaa@YearStyle}[1]{{\normalsize\rmfamily#1}}
\newcommand{\rmaa@VolumeStyle}[1]{{\normalsize\rmfamily\bfseries#1}}
\newcommand{\rmaa@VolumenStyle}[1]{{\Large\sffamily\bfseries#1}}
\newcommand{\rmaa@JournalNameStyle}[1]{{\normalsize\itshape#1}}
\newcommand{\rmaa@ReceivedStyle}[1]{{\normalsize\rmfamily\itshape#1}}
\newcommand{\rmaa@AbstractHeaderStyle}[1]{%
  {\normalsize\rmfamily\MakeUpperCase{#1}}}
\newcommand{\rmaa@TitleStyle}[1]{%
  {\large\bfseries\rmfamily\boldmath\MakeUpperCase{#1}\par}}
%    \end{macrocode}
% \end{macro}
% \end{macro}
% \end{macro}
% \end{macro}
% \end{macro}
% \end{macro}
% \end{macro}
% \end{macro}
% \end{macro}
%
% \begin{macro}{\rmaa@FrontSecStyle}
% Headings of ``front sections'' (see section~\ref{sec:front-sections})
%    \begin{macrocode}
\newcommand{\rmaa@FrontSecStyle}[1]{%
  {\normalsize\bfseries\rmfamily\MakeUpperCase{#1}}}
%    \end{macrocode}
% \end{macro}
%
% \begin{macro}{\rmaa@TOCTitleStyle}
% \begin{macro}{\rmaa@TOCPartStyle}
% \begin{macro}{\rmaa@IndexTitleStyle}
% \begin{macro}{\rmaaTOCAuthorStyle}
% \begin{macro}{\rmaaIndexAuthorStyle}
% Commands to set style of TOC and index entries.
% Rewritten to avoid unnecessary font commands in the .toc and .ind files
%    \begin{macrocode}
\newcommand{\rmaa@TOCTitleStyle}[1]{\MakeUpperCase{#1}}
\newcommand{\rmaa@TOCPartStyle}[1]{\rmaa@TOCTitleStyle{#1}}
\newcommand{\rmaa@IndexTitleStyle}[1]{#1}
\newcommand{\rmaaTOCAuthorStyle}[1]{%
  {\itshape#1}}
\newcommand{\rmaaIndexAuthorStyle}[1]{%
  \protect\mbox{\bfseries#1}}
%    \end{macrocode}
% \end{macro}
% \end{macro}
% \end{macro}
% \end{macro}
% \end{macro}
%
% \begin{macro}{\rmaa@KeywordStyle}
% \begin{macro}{\rmaa@KeywordHeaderStyle}
% The style of the keywords depends on whether we are using the
% \texttt{newstyle} option or not. If we are, then they are small. 
%    \begin{macrocode}
\ifthenelse{\boolean{@newstyle}}{
  \newcommand{\rmaa@KeywordStyle}[1]{%
    {\small%
      \if@boldkeywords 
      \bfseries\MakeUpperCase{\raggedright #1}
      \else
      \raggedright #1
      \fi
    }}
  \newcommand{\rmaa@KeywordHeaderStyle}[1]{%
    {\small\itshape#1}}}
{
  \newcommand{\rmaa@KeywordStyle}[1]{%
    {\normalsize%
      \if@boldkeywords\bfseries\fi\raggedright #1}}
  \newcommand{\rmaa@KeywordHeaderStyle}[1]{%
    {\normalsize\itshape#1}}}
%    \end{macrocode}
% \end{macro}
% \end{macro}
%
% \subsection{Page styles}\label{sec:page-styles}
% \begin{macro}{\ps@rmaa}
% The running headers for normal pages:
%    \begin{macrocode}
\newcommand{\ps@rmaa}{%
%    \end{macrocode}
% nothing in either foot; 
%    \begin{macrocode}
  \renewcommand\@oddfoot{\relax}%
  \renewcommand\@evenfoot{\relax}%
%    \end{macrocode}
% page number and short author list in the even heads; 
%    \begin{macrocode}
  \renewcommand\@evenhead{\rlap{\thepage}\hfil\@shortauthor\hfil}%
%    \end{macrocode}
% short title and page number in the odd heads.
%    \begin{macrocode}
  \renewcommand\@oddhead{\hfil\@shorttitle\hfil\llap{\thepage}}
  }
%    \end{macrocode}
% \begin{macro}{\ps@rmaafirstpage}
% The header of the first page of each
% article:
%    \begin{macrocode}
\newcommand{\ps@rmaafirstpage}{%
  \renewcommand\@oddfoot{\hfil\thepage}%
%    \end{macrocode}
% footer corrected (06 Jun 2000) to have page on outside edge;
%    \begin{macrocode}
  \renewcommand\@evenfoot{\thepage\hfil}%
%    \end{macrocode}
% Header has journal name, volume, and page numbers for a regular
% article, or similar for special cases. 
% \changes{4.0}{2019/11/03}{Change to three-line header with copyright info and DOI.}
%    \begin{macrocode}
\if@preprint
\renewcommand\@oddhead{%
  \if@proceedings 
  \textsc{To appear in ``\@conftitle{} (\rmaa@YearStyle{\@year})''}\hfil
  \textit{RMxAC}
  \else
  \textsc{To appear in} \@journalname\hfil
  \fi}\fi%
\if@guide
\renewcommand\@oddhead{\textsc{Documentation} (\rmaa@YearStyle{\@year})\hfil}
\fi%
\if@manuscript
\renewcommand\@oddhead{Manuscript for \@journalname{} (\rmaa@YearStyle{\@year})\hfil}
\fi%
\if@journal
\renewcommand\@oddhead{%
  \parbox{\textwidth}{
    \@journalname, \rmaa@VolumeStyle{\@volume},
    \rmaa@PagesStyle{\@artpages}\
    (\rmaa@YearStyle{\@year})\hfil\\[1pt]
    {\small\textcopyright{} \@year:
      Instituto de Astronom\'\i{}a,
      Universidad Nacional Aut\'onoma de M\'exico}\hfil\\[1pt]
    {\small\@DOI@url}\hfil
  }%
}\fi%
  \renewcommand\@evenhead{\@oddhead}}
%    \end{macrocode}
% By default, use the |\ps@rmaa| page style.
%    \begin{macrocode}
\pagestyle{rmaa}
%    \end{macrocode}
% \end{macro}
% \end{macro}
%
% \subsection{Preamble commands}\label{sec:preamble} 
% Everything you need for setting the title page, TOC and index
% entries, etc.
% \subsubsection{Commands for the title page}
% \label{sec:title-commands}
% \begin{macro}{\resumen}
% \begin{macro}{\@resumen}
% \begin{macro}{\abstract}
% \begin{macro}{\@abstract}
% User commands for setting the resumen and abstract. 
% These should be used in the document header. The internal commands
% |\@resumen| and |\@abstract| are used to hold the actual text
% pending output via the |\maketitle| command. 
%    \begin{macrocode}
\newif\if@setresumen
\newcommand{\resumen}[1]{%
%%  \renewcommand{\@resumen}{{\selectlanguage{spanish}#1}}%
  \renewcommand{\@resumen}{{\ignorespaces#1}}%
  \@setresumentrue}
\renewcommand{\abstract}[1]{%
  \renewcommand{\@abstract}{{\ignorespaces#1}}
  \renewcommand\ADS@abstract{#1}
%% Not a good idea, since then author macros are not expanded and
%% cause errors later. However, by not doing it we are vulnerable to
%% fragile commands.....
  \if@noexpandADS\@onelevel@sanitize{\ADS@abstract}\fi
  \@abstractpresenttrue
  }
\newcommand{\ADSnoexpand}{\@noexpandADStrue}
%    \end{macrocode}
% \begin{macro}{\boldabstract}
% This is the version for doing 1-page conference proceedings. Still
% save the text for ADS but then write it out immediately as a bold
% paragraph. 
%    \begin{macrocode}
\newcommand{\boldabstract}[1]{%
  \gdef\ADS@abstract{#1}%
  \renewcommand{\@abstract}{#1}%
  \if@noexpandADS\@onelevel@sanitize{\ADS@abstract}\fi%
  \noindent\textbf{\boldmath #1}
  }
%    \end{macrocode}
% \end{macro}
% Default value in each case is to print out a warning if not set
%    \begin{macrocode}
  \newcommand{\@resumen}{%
    \@setresumenfalse \ClassWarningNoLine{rmaa}{No Spanish abstract
      (resumen) specified - if you do not speak Spanish, this will be
      translated by the editors}%
    \textbf{Favor de proporcionar un resumen en espa\~nol. If you are
      unable to translate your abstract into Spanish, the editors will
      do it for you.} }
  \newcommand{\@abstract}{\ClassWarningNoLine{rmaa}{No abstract
      specified}}
%    \end{macrocode}
% \end{macro}
% \end{macro}
% \end{macro}
% \end{macro}
%
% \begin{macro}{\title}
% \begin{macro}{\@rawtitle}
% \begin{macro}{\@TOCtitle}
% \begin{macro}{\rmaaIndexTitle}
% Similarly, |\title| is the preamble command to set the title of the
% article. This is used in various places, typeset in different
% styles, hence the use of the internal |\@rawtitle|, |\@title|,
% |\rmaaIndexTitle|, and |\@TOCtitle| commands, each of which uses the
% relevant style from section~\ref{sec:text-styles}
%    \begin{macrocode}
\newcommand{\@rawtitle}{Raw Title}
\newcommand{\@TOCtitle}{\rmaa@TOCTitleStyle{Title for TOC}}
\newcommand{\rmaaIndexTitle}{\rmaa@IndexTitleStyle{Title for index}.}
\newcommand{\@TOCaftertitle}{}
\renewcommand{\title}[1]{%
%%  {\ADS@sanitize\xdef\@rawtitle{#1}}
  \renewcommand\@rawtitle{\noexpand#1}
  \let\ADS@title=\@rawtitle
  \renewcommand{\@title}{\rmaa@TitleStyle{#1}}%
  \renewcommand{\@TOCtitle}{\rmaa@TOCTitleStyle{#1}\@TOCaftertitle}%
  \renewcommand{\rmaaIndexTitle}{\rmaa@IndexTitleStyle{#1\@TOCaftertitle}.}%
  }
%    \end{macrocode}
% \end{macro}
% \end{macro}
% \end{macro}
% \end{macro}
% 
% All the weird and whacky ways of indicating the author
% affiliations. 
% \begin{macro}{\affil}
% First, the way that sets it directly underneath the authors name.
% Note that |\affil| trumps all other affiliations for the purposes of
% getting into the ADS stuff (we use |\gdef|).   
% \dots 
%    \begin{macrocode}
\newcommand{\affil}[1]{%
  \\[1.5ex]\rmaa@AffilStyle{#1}\\[1.5ex]%
  \if@firstaffil\gdef\ADS@affil{#1}\@firstaffilfalse\fi%
  }
\renewcommand{\and}{and\\[1.5ex]}
%    \end{macrocode}
% \begin{macro}{\altaffil}
% Alternative affiliations are footnoted. 
%    \begin{macrocode}
\newcommand{\altaffil}[1]{%
  \if@firstaltaffil\gdef\ADS@affil{#1}\@firstaltaffilfalse\fi%
  \thanks{#1}}
%    \end{macrocode}
% \begin{macro}{\altaffilmark}
% \begin{macro}{\altaffiltext}
% Alternative-alternative mechanism for affiliations---should not be
% mixed with |\altaffil| (unless you know what you are doing) or there
% will be problems with numbering.
%    \begin{macrocode}
\newcommand{\altaffilmark}[1]{\textsuperscript{#1}}
\newcommand{\altaffiltext}[2]{%
  \if@firstaltaffil\gdef\ADS@affil{#2}\@firstaltaffilfalse\fi
  \stepcounter{footnote}%
  \protected@xdef\@thanks{\@thanks
    \protect\footnotetext[#1]{#2}}} 
%    \end{macrocode}
% \end{macro}
% \end{macro}
% \end{macro}
% \end{macro}
%
% \begin{macro}{\keywords}
% \begin{macro}{\@keywords}
% \begin{macro}{\@keywordheader}
% \label{keywords}Sundry minor preamble commands. All
% self-explanatory, I hope. |\keywords| is the old user interface
% (deprecated). New articles should use |\addkeyword| once for each
% keyword. That way the ADS record comes out right too (comma rather
% than em-dash separator).
% \changes{2.2beta12}{2001/12/07}{\cs{addkeyword} New user interface
%   for keywords.} 
% \changes{2.2beta13}{2002/03/11}{\cs{addkeyword} Improved behaviour
%   with .ads file. Use
%   \cs{ADS@sanitize} to change definitions of `~' and \cs{ion}.}
% \changes{3.23}{2007/09/09}{\cs{addkeyword} now signals a compilation
% error if the author tries to add more than 6 keywords.}
% \changes{3.24}{2007/09/14}{kwdcounter is now reset in
% \cs{@begofarticlehook} to prevent error in book style} 
%    \begin{macrocode}
\newcommand{\@kwdsep}{ --- }
\newcommand{\ADS@kwdsep}{, }
\newcounter{kwdcounter}\setcounter{kwdcounter}{0}
\newcommand{\addkeyword}[1]{%
  \addtocounter{kwdcounter}{1}
  \ifthenelse{\thekwdcounter > 6}{
    \ClassError{rmaa}{Too many keywords. A maximum of 6 keywords is allowed.}{}
  }{}
  \if@firstkwd%
  \renewcommand\@keywords{\if@boldkeywords\uppercase{#1}\else #1\fi}
  {\ADS@sanitize\def\ion{\plainion}\gdef\ADS@keywords{#1}}
  \@firstkwdfalse
  \else%
  \edef\@keywords{\@keywords\@kwdsep \if@boldkeywords\uppercase{#1}\else #1\fi}%
  {\ADS@sanitize\def\ion{\plainion}\xdef\ADS@keywords{\ADS@keywords\ADS@kwdsep #1}}%
  \fi}
\newcommand{\ADS@sanitize}{%
  \def\mathrm{}
  \def\@{\ }
  \def~{\ }
  \def\langle{<}
  \def\rangle{>}
  \def\log{log }
  \def\sin{sin }
  \def\cos{cos }
  \def\tan{tan }
  \def\fd{^d.}
  \def\fh{^h.}
  \def\fm{^m.}
  \def\fs{^s.}
  \def\fdg{�.}
  \def\farcm{'.}
  \def\farcs{''.}
  \def\arcdeg{�}
  \def\,{ }
  }
\newcommand{\keywords}[1]{%
  \renewcommand{\ADS@keywords}{#1}
  \renewcommand{\@keywords}{\uppercase{#1}}
  }
\newcommand{\@keywordheader}{\rmaa@KeywordHeaderStyle{Key Words:}}
\newcommand{\@keywords}{%
  \ClassWarningNoLine{rmaa}{No keywords specified}%
  Keywords go here
  }
%    \end{macrocode}
% \begin{macro}{\shorttitle}
% \begin{macro}{\shortauthor}
% \begin{macro}{\@shorttitle}
% \begin{macro}{\@shortauthor}
% Commands for setting the running headers.
%    \begin{macrocode}
\newcommand{\shorttitle}[1]{%
  \renewcommand\@shorttitle{\MakeUpperCase{#1}}}
\newcommand{\shortauthor}[1]{%
  \renewcommand\@shortauthor{\MakeUpperCase{#1}}}
\newcommand{\@shorttitle}{SHORT TITLE GOES HERE}
\newcommand{\@shortauthor}{SHORT AUTHOR GOES HERE}
%    \end{macrocode}
% \end{macro}
% \end{macro}
% \end{macro}
% \end{macro}
% \end{macro}
% \end{macro}
% \end{macro}


% \subsubsection{The full addresses}
% \label{sec:full-addresses}
% \begin{macro}{\fulladdresses}
% \begin{macro}{\@fulladdresses}
% \begin{macro}{\@fulladdheight}
% \begin{macro}{\@fulladddepth}
% \begin{macro}{\@fulladdbox}
% Aaghh, the full addresses! A perennial source of problems. 
%    \begin{macrocode}
\newlength{\@fulladdheight}
\newlength{\@fulladddepth}
\newsavebox{\@fulladdbox}
%% \fulladdresses command modified (11 May 2000) so that we can put
%% \fulladdresses{} in the article to suppress the output of the
%% addresses
\newcommand{\@FinalFigure}{\relax}
\newcommand{\FinalFigure}[1]{\renewcommand{\@FinalFigure}{#1}}
\newcommand{\fulladdresses}[1]{%
  \renewcommand{\@rmaaArg}{#1}
  \ifx\@rmaaArg\@rmaaEmpty
     \renewcommand{\@fulladdresses}{}
  \else
     \renewcommand{\@fulladdresses}{
       \normalsize%
       \ifx\@FinalFigure\relax\else\vfill\@FinalFigure\vfill\fi
%%       \vspace*{1ex}
       %% do addresses as a list (copied from thebibliography environment) 
       \list{\null}{\leftmargin 1.5em\topsep\z@\labelwidth\z@\labelsep\z@\itemindent -1.5em
         \itemsep\z@\parsep\z@\usecounter{enumi}}
       \sloppy\clubpenalty4000\widowpenalty4000
       \sfcode`\.=1000\relax
       #1
       \endlist
       }
  \fi
  }
\newcommand{\@fulladdresses}{Full addresses go here}
%    \end{macrocode}
% \begin{macro}{\suppressfulladdresses}
% New user-level command for suppressing full addresses should you so
% desire. Currently used in the single-page-poster style. 
%    \begin{macrocode}
\newcommand{\suppressfulladdresses}{\@suppressfulladdressestrue\@needaddressesfalse}
%    \end{macrocode}
% \end{macro}
% \end{macro}
% \end{macro}
% \end{macro}
% \end{macro}
% \end{macro}

% \begin{macro}{\printdates}
% \begin{macro}{\noprintdates}
% Command to activate the printing out of the received/accepted
% dates. This can also be done using the \texttt{printdates} option,
% which is the default for \texttt{journal}, but not for \texttt{proceedings}.
%    \begin{macrocode}
\newcommand{\printdates}{\@printdatestrue}
\newcommand{\noprintdates}{\@printdatesfalse}
%    \end{macrocode}
% \end{macro}
% \end{macro}
% \begin{macro}{\@outputaddressesifnecessary}
% Command to output the addresses if we haven't already done so
% This is called by: 
% \begin{enumerate}
% \item The end of |thebibliography| environment
% \item The |\outputfulladdresses| command
% \end{enumerate}
%    \begin{macrocode}
\newcommand{\@outputaddressesifnecessary}{%
  \ClassInfo{rmaa}{Calling @outputaddressesifnecessary}
  \ifthenelse{\boolean{@needaddresses}}{%
    %% If in one column style, or in the first column of a two column
    %% style then we can try to put it on this page 
    \ifthenelse{\boolean{@firstcolumn} \or \not\boolean{@rmtwocol}}{%
      \@outputaddresses}{%
      %% otherwise give warning it will be on a separate page
      \ClassWarning{rmaa}{Full addresses will go on a separate
        page (please see User Guide for how to remedy this)}}%
    %% If we haven't managed yet, put it on its own page
    \ifthenelse{\boolean{@needaddresses}}{%
      \if@debug{\ClassWarning{rmaa}{In this silly bit!}}\fi
      \clearpage\vfill\@outputaddresses}{}%
    }{
    %% If we don't need them, we still might need to do the eoa hook 
    \ifthenelse{\boolean{@suppressfulladdresses}}{%
      \if@debug{\ClassWarning{rmaa}{XXX}}\fi
%%      \if@book\else\@endofarticlehook\null\fi
      }{}
    }%
  }  
%    \end{macrocode}
% \end{macro}

% \begin{macro}{\outputfulladdresses}
% A user interface to the above command, in case there is no biblio
% (11 May 2000) 
%    \begin{macrocode}
\newcommand{\outputfulladdresses}{\@outputaddressesifnecessary}
%    \end{macrocode}
% \end{macro}

% \begin{macro}{\adjustfinalcols}
% This command starts a new column and should be 
% used to ensure that the columns balance on the last page. It also
% outputs the full 
% author addresses to ensure that this occurs in the left hand
% column. 
% The |\par| command is so that the current line will be taken account
% of when working out |\pagetotal|
%    \begin{macrocode}
\newcommand{\adjustfinalcols}{\if@twocolumn\goodbreak\@outputaddresses\newpage\noindent\fi}
%    \end{macrocode}
% \end{macro}

% \begin{macro}{\@outputaddresses}
% Command to do the work of outputting the addresses. 
% This can be called by 
% \begin{enumerate}
% \item an explicit |\adjustfinalcols|
% \item via |\@outputaddressesifnecessary|, q.v.
% \end{enumerate}
%
%    \begin{macrocode}
\newcommand{\@outputaddresses}{%
  % Allow option of suppressing the full addresses, in which case
  % we just run the hook if we are not in a book style
  \if@suppressfulladdresses
%%     \if@book\else\@endofarticlehook\null\fi
  \else
  %% save the full addresses in a box for output later
  \sbox{\@fulladdbox}{\parbox{\textwidth}{\@fulladdresses}}
  % calculate height of address box  
  \settoheight{\@fulladdheight}{\usebox{\@fulladdbox}}
  \settodepth{\@fulladddepth}{\usebox{\@fulladdbox}}
  \addtolength{\@fulladdheight}{\@fulladddepth}
  \if@debug
  { \typeout{Address total height: \the\@fulladdheight}
    \typeout{Text height: \the\pagetotal}
    \typeout{Goal height: \the\pagegoal}
    }
  \fi
  % reuse \@fulladdheight to mean the total height of text+addresses
  \addtolength{\@fulladdheight}{\pagetotal}
%%  \addtolength{\@fulladdheight}{-\baselineskip}
  \if@debug\typeout{Adjusted text + addresses height: \the\@fulladdheight}\fi
  % check it will fit on current page
  \ifdim\@fulladdheight<\pagegoal
     %% must shift it to the left
     %% to compensate for the hanging indentation in the reference list. 
     %% NOTE: this is not foolproof!
     %% NB. This will give overfull hbox warnings, which is a bit ugly.
     \vfill\noindent\if@biblist\hspace*{-\leftmargin}\fi
     \usebox{\@fulladdbox}
     % Set this to show we've done it
     \@needaddressesfalse
     % and assume we are at the end of the article
     % so that we do the EOA hook if not in the book style
%%     \if@book\else\@endofarticlehook\fi
  \else
     % if not, just do nothing. With luck, this function will get 
     % called again on the next page
     \if@debug{\typeout{Holding over full addresses to next page}}\fi
  \fi\fi
  }
%    \end{macrocode}
% \end{macro}

% \subsubsection{Commands for editors}
% \label{sec:editor-commands}
% Some user/editor commands for doing the TOC and index in the book
% option. 
% They are defined here rather in the book-specific part since they
% need to be always available in order to avoid breaking documents. 
% \begin{macro}{\listofauthors}
% \begin{macro}{\@rawauthors}
% \begin{macro}{\rmaaTOCauthors}
% The command |\listofauthors| formats an author list for use in the
% TOC. It should contain the same names as in the |\author| command,
% but without the affiliation garbage and using |\&| instead of
% |and|. 
%    \begin{macrocode}
\newcommand{\@rawauthors}{Raw Authors}
\newcommand{\rmaaTOCauthors}{\noexpand\rmaaTOCAuthorStyle{Author list for TOC}}
\newcommand{\listofauthors}[1]{%
  \renewcommand{\@rawauthors}{#1}
  \renewcommand\rmaaTOCauthors{\noexpand\rmaaTOCAuthorStyle{#1}}}
%    \end{macrocode}
% \begin{macro}{\indexauthor}\label{indexauthor}
% The command |\indexauthor| produces an index entry for each
% author. 
% This command must come \emph{after} the
% |\title| and the |\listofauthors| commands. Also, there should be one 
% |\indexauthor| command for \emph{each} author. We also add the
% author onto the end of the |\ADS@authors| macro and construct the
% |\ADS@authorinitial| by a most cunning method (see below), but only
% for the first author. 
% 
% No, No, NO! This was all wrong. Backed out completely: 
% \begin{quote}
% Completely rewritten (19 Dec 2002) since we now use the index class, so 
% as to be able to alphabetize accented characters. A side effect of this is
% that |\index| no longer expands out the argument after the |@|, so we 
% instead add all the |\index| commands to a list (|\indexauthor@deferlist|), 
% which is executed in the |\@begofarticlehook|. 
% 
% Note the stuff with |\toks0| is copied from Exercise 20.15 of the \TeX{}book---it is
% to make sure the latin1-accented name gets passed into the idx
% file. 
% \end{quote}
% We now go back to the old way of doing things, using an optional
% argument to the \cs{indexauthor} command that has the name without
% accents. 
% \changes{3.08}{2003/01/16}{Remerged changes from 3.06 that got
% dropped from 3.07 after a synchronization disaster. Basically, the
% accentless optional argument to \cs{indexauthor}. Must add the macro
% directories to my unison set-up to stop this sort of thing from happening.}
%    \begin{macrocode}
\newcommand\@isortkey{}
\newcommand\@emptykey{empty}
\newcommand\indexauthor[2][empty]{
  \renewcommand\@isortkey{#1}
  \ifx\@isortkey\@emptykey \renewcommand\@isortkey{#2}\fi
  \index{\@isortkey @{{\rmaaIndexAuthorStyle{#2}}\hskip1em plus1em minus0.4em
      \rmaaIndexTitle\hskip1em plus2em minus0.4em\penalty-100
      \rmaaTOCauthors}}
  \indexauthorADS{#2}
}
%%\def\indexauthor@deferlist{\relax}
%%\long\def\indexauthor#1{
%%  \def\@@name{#1} \toks0=\expandafter{\@@name}
%%  \toks1=\expandafter{\indexauthor@deferlist}
%%  \protected@xdef\indexauthor@deferlist{%
%%    \the\toks1
%%    \noexpand\index{\the\toks0@{\noexpand\rmaaIndexAuthorStyle{\the\toks0}\noexpand\hskip1em plus1em minus0.4em 
%%        \rmaaIndexTitle\noexpand\hskip1em plus2em minus0.4em\noexpand\penalty-100
%%        \rmaaTOCauthors}}
%%  }
%%  \indexauthorADS{#1}
%%  }
%%\xdef\indexauthor@deferlist{}
\newcommand\indexauthorADS[1]{%
  \if@firstauthor\protected@xdef\ADS@authorinitial{\firstletter{#1}}\fi
  \protected@xdef\ADS@authors{\ADS@authors\if@firstauthor\else ; \fi
    #1}
  \@firstauthorfalse
}
%    \end{macrocode}
% Note that the first |#1| sets the key for the index, making sure it
% actually starts with  
% first letter of the surname. Meanwhile, the |@{...}| expression,
% which starts with some  
% control sequences is what is written to the |.ind| file. 
% See the makeindex man page for details. 
% \end{macro}
% \end{macro}
% \end{macro}
% \end{macro}

% \begin{macro}{\KillAllAccents}
% \begin{macro}{\firstletter}
% These are utility functions used in the above. |lett| is a counter
% indicating the current letter in the string. |\firstletter| just
% prints the first letter of its argument.
%    \begin{macrocode}
\newcommand{\KillAllAccents}{\def\'{}\def\`{}\def\~{}\def\.{}\def\i{i}}
\newcommand\firstletter[1]{\@car #1\@nil}
%    \end{macrocode}
% \end{macro}
% \end{macro}
% \begin{macro}{\SetVolume}
% \begin{macro}{\SetMSnumber}
% \begin{macro}{\SetFirstPage}
% \begin{macro}{\SetYear}
% \begin{macro}{\ReceivedDate}
% \begin{macro}{\AcceptedDate}
% These are for the editors to set the volume number, etc. 
%    \begin{macrocode}
\newcommand{\SetVolume}[1]{%
  \renewcommand{\@volume}{#1}
  \edef\ADS@volume{\PadToFourWithDots{#1}}
  \renewcommand{\@volumen}{\rmaa@VolumenStyle{volumen #1}}%
  }
\newcommand{\SetMSnumber}[1]{\renewcommand{\@MSnum}{#1}}
%    \end{macrocode}
% Note that there is no real reason to use |\SetFirstPage| since all
% the page numbering is done automatically in the book style. 
%    \begin{macrocode}
\newcommand{\SetFirstPage}[1]{\setcounter{page}{#1}}
\newcommand{\SetYear}[1]{%
  \renewcommand{\@year}{#1}
  \renewcommand{\ADS@year}{#1}
  \renewcommand{\@anyo}{\rmaa@VolumenStyle{#1}}
  }
\newcommand{\ReceivedDate}[1]{\renewcommand{\@received}{#1}}
\newcommand{\AcceptedDate}[1]{\renewcommand{\@accepted}{#1}}
%    \end{macrocode}
%
% \begin{macro}{\DOI}
% User/editor command to enter the DOI, which will be typeset as a link.
% \changes{4.0}{2019/11/02}{New user command \cs{DOI} to typeset
% DOI in the header of first page and add it to ADS info.}
%    \begin{macrocode}
\newcommand\DOI[1]{%
  \def\@DOI@url{\url{https://doi.org/10.22201/ia.01851101p.#1}}%
  \def\ADS@doi{https://doi.org/10.22201/ia.01851101p.#1}%
}%
%    \end{macrocode}
% \end{macro}
%
% \begin{macro}{\ORCIDiD}
% \changes{4.0}{2019/11/12}{Add mechanism for giving ORCID iD links}
% \begin{macro}{\ORCIDiDfull}
% \begin{macro}{\@orcid@icon}
% \begin{macro}{\@orcid@iconbig}
% Commands to add icon links to ORCID iDs of authors. For example
% \begin{quote}
% |\ORCIDiD|\marg{0000-0000-0000-0000}
% \end{quote}
% to be used after each author name in the |\author|
% command and in the |\fulladdresses|.  Alternatively,
% |\OORCIDiDfull| will typeset the icon plus the full text of the
% link, but we don't use that at present.  The icon image is stored in |ORCIDiD_iconvector.pdf|, which needs to be on the input path.  In principle, we could make EPS versions if necessary.
%    \begin{macrocode}
\newcommand\@orcid@icon{\includegraphics[height=1.8ex]{ORCIDiD_iconvector}}
\newcommand\@orcid@iconbig{\includegraphics[height=2ex]{ORCIDiD_iconvector}}
\newcommand\ORCIDiD[1]{\href{https://orcid.org/#1}{\,\@orcid@icon}}
\newcommand\ORCIDiDfull[1]{%
  \href{https://orcid.org/#1}{\@orcid@iconbig}\,\url{https://orcid.org/#1}}
%    \end{macrocode}
% \end{macro}
% \end{macro}
% \end{macro}
% \end{macro}
%
% \begin{macro}{\@volume}
% \begin{macro}{\@volumen}
% \begin{macro}{\@artpages}
% \begin{macro}{\@year}
% \begin{macro}{\@anyo}
% \begin{macro}{\@received}
% \begin{macro}{\@accepted}
% \begin{macro}{\@MSnum}
% Default values for the containers used by above commands, plus a
% few more useful ones. 
%    \begin{macrocode}
\newcommand{\@volume}{00}
\newcommand{\@volumen}{\rmaa@VolumenStyle{00}}
\newcommand{\@artpages}{%
  \oldpageref{\articlename:RMAAFirstPage}--\oldpageref{\articlename:RMAALastPage}}
\newcommand{\@year}{\the\year}
\newcommand{\@anyo}{\rmaa@VolumenStyle{2007}}
\newcommand{\@journalname}{\rmaa@JournalNameStyle{%
    \if@proceedings\RMAASClong\else\RMAAlong\fi}}
\newcommand{\@received}{}
\newcommand{\@accepted}{}
\newcommand{\@MSnum}{0000}
\newcommand{\@DOI@prefix}{https://doi.org/10.22201/ia.01851101p}
\def\@DOI@url{\url{https://doi.org/10.22201/ia.01851101p.2099.00.00.00}}
%    \end{macrocode}
% \end{macro}
% \end{macro}
% \end{macro}
% \end{macro}
% \end{macro}
% \end{macro}
% \end{macro}
% \end{macro}
% \end{macro}
% \end{macro}
% \end{macro}
% \end{macro}
% \end{macro}
% \end{macro}
% 
% The following commands are most useful in the \texttt{book}
% style. However, the |\SetConfTitle| command can now be used in
% individual articles to set the header for the \texttt{preprint}
% style. 
%    \begin{macrocode}
\newcommand{\@reunionplace}{Lugar de la reuni\'on}
\newcommand{\@reuniondate}{Fecha de la reuni\'on}
\newcommand{\@conftitle}{Conference Title}
\newcommand{\SetReunion}[2]{\renewcommand{\@reunionplace}{#1}\renewcommand{\@reuniondate}{#2}}
\newcommand{\SetConfTitle}[1]{\renewcommand{\@conftitle}{#1}}
\newcommand{\@editors}{Juan Ram\'on de la Fuente and ``El Mosh''}
\newcommand{\SetEditors}[1]{\renewcommand{\@editors}{#1}}
\newcommand{\@editores}{Juan Ram\'on de la Fuente y ``El Mosh''}
\newcommand{\SetEditores}[1]{\renewcommand{\@editores}{#1}}
\newcommand{\@mes}{mes de la publicaci\'on}
\newcommand{\SetMes}[1]{\renewcommand{\@mes}{\rmaa@VolumenStyle{#1}}}
%    \end{macrocode}

% \subsection{Notes for Author Proofs and Communication between Editors}
% \label{sec:notes}
% We introduce a command \cs{Note} that should be used to describe
% any change made to the manuscript. These are a bit tricky to get
% right. In general we use \cs{marginpar} but this only works in outer
% paragraph mode, so we just bung the box where we are if we are in a
% figure caption or something. If we are in the preamble, then we just
% use a \marginpar as part of the \maketitle command.
%    \begin{macrocode}
\newif\if@arenotes
\@arenotesfalse
\newcounter{notecounter}
\setcounter{notecounter}{0}
\newcommand\NotesList{}
\newcommand\Note[1]{%
  \if@proof{%
    \@arenotestrue%
    \stepcounter{notecounter}%
    \if@inpreamble\else{%
      \NoteMark{\arabic{notecounter}}%
    }\fi
    \protected@xdef\NotesList{\NotesList 
      \protect\item #1}%
  }\fi%
}
\newcommand\Changed[2]{Changed ``#1'' to ``#2''.}
\newcommand\NoteMark[1]{%
  \ifinner
  \raisebox{1ex}[0pt][0pt]{\makebox[0pt][l]{\framebox{\footnotesize\bfseries
        N#1}}}%
  \else
  \marginpar{\framebox{\footnotesize\bfseries N#1}}%
  \fi
}
\newcommand\RefNote[1]{\if@proof\NoteMark{\ref{#1}}\relax\fi}
\newcommand\NoteLabel{\protect\label}
\newcommand\OutputNotes{%
  \bgroup
  \onecolumn
  \section*{EDITORIAL NOTES}
  \begin{enumerate}
    \NotesList
  \end{enumerate}
  \egroup
}
\DeclareRobustCommand\cs[1]{\texttt{\char`\\#1}}

%    \end{macrocode}
% \subsection{Bibliographic Records for ADS}
% \label{sec:ADS}
% \changes{2.2beta11}{2001/11/27}{Just about finished all the ADS
% stuff now. Except for the references, which I have given up on for
% the time being.}
% \begin{macro}{\ADS@BibCode}
% \begin{macro}{\ADS@year}
% \begin{macro}{\ADS@journal}
% \begin{macro}{\ADS@volume}
% \begin{macro}{\ADS@firstpage}
% \begin{macro}{\ADS@authorinitial}
% We use the |\ADS@BibCode| command to hold the bibliographic code for
% ADS. It is 
% used in the |%R| 
% field of the bibliographic record for each article. 
% The format is 
% \begin{quote}
% \texttt{YYYYJJJJJVVVVMPPPPA}  
% \end{quote}
% where \texttt{YYYY} is the year, \texttt{JJJJJ} is the journal code,
% which is \texttt{RMxAA} or \texttt{RMxAC}, \texttt{VVVV} is the
% volume number (right-justified), \texttt{M} is not used by us,
% \texttt{PPPP} is the first page number (right-justified), and
% \texttt{A} is the initial of the first author's surname. The volume
% and page fields are padded with periods (`\texttt{.}'), which causes
% us some fun (this is done with the macro |\PadToFourWithDots| described
% below) . 
% 
% All these individual bits are held in the commands |\ADS@year|,
% |\ADS@journal|, |\ADS@volume|, |\ADS@firstpage|,
% |\ADS@authorinitial|. 
%    \begin{macrocode}
\newcommand{\SetADSBibCode}[1]{\renewcommand{\ADS@BibCode}{#1}}
\newcommand{\ADS@BibCode}
{\ADS@year\ADS@journal\ADS@volume.\PadToFourWithDots{\ADS@firstpage}\ADS@authorinitial}
\newcommand{\ADS@year}{2001} 
\newcommand{\ADS@journal}{\if@proceedings RMxAC\else RMxAA\fi}
\newcommand{\ADS@volume}{....} 
%    \end{macrocode}
% These two are just set to default values here, |\ADS@firstpage| is set
% in the |document| environment for included articles, in the
% |frontsection| environment, |\tableofcontents| command,  |index| and
% |abstracts| environments. Also by |\AtBeginDocument| for stand-alone
% articles. 
% 
% God knows how we are going to do |\ADS@authorinitial|!
%    \begin{macrocode}
\newcommand{\ADS@firstpage}{....}
\newcommand{\ADS@authorinitial}{A}
%    \end{macrocode}
% \end{macro}
% \end{macro}
% \end{macro}
% \end{macro}
% \end{macro}
% \end{macro}
% We also have a whole slew of other fields that can be set. These are
% all listed in Table~\ref{tab:ADS}. Most of them still have to be
% defined. 
% \begin{table}[!ht]\raggedleft
% \caption{ADS Bibliographic Records}\label{tab:ADS}
% \begin{tabular}{lp{8em}lp{8em}}\hline
% Code & Description & Macro Name & Comment \\ \hline
% |%R| & Bibliographic Code & |\ADS@BibCode| & required \\
% |%A| & Author list & |\ADS@authors| & required, generate from |\indexauthor| \\
% |%F| & Author affiliation & |\ADS@affil| & format: AA(\dots)
% AB(\dots) \dots \\
% |%J| & Journal Name & |\ADS@journalname| &  \\
% |%D| & Publication Date & |\ADS@fulldate| &  \\
% |%L| & Last Page & |\ADS@lastpage| & generate from label \\
% |%T| & Title & |\ADS@title| & required, generate from |\title| \\
% |%C| & Abstract Copyright & |\ADS@copyright| & set to IA-UNAM \\
% |%U| & URL for electronic document & |\ADS@url| & set to PDF file \\
% |%d| & DOI for article & |\ADS@doi| & set by |\DOI| \\
% |%K| & Keywords & |\ADS@keywords| & set equal to |\@keywords| \\
% |%A| & Abstract Text & |\ADS@abstract| & set equal to |\@abstract| \\
% |%R| & References & |\ADS@references| & generate from the
% bibliography  \\ \hline
% \end{tabular}
% \end{table}
% A lot of these macros are set by other commands, as mentioned in the
% Table, but some need to be set  explicitly here. 
% 
% \begin{macro}{\ADS@doi}
% \changes{4.0}{2019/11/08}{Save DOI for ADS}
% The user/editor command \cs{DOI} is used to set this. 
%    \begin{macrocode}
\newcommand{\ADS@doi}{https://doi.org/10.22201/ia.01851101p.0000.00.00.00}
%    \end{macrocode}
% \end{macro}
% 
% \begin{macro}{\ADS@journalname}
% This
% does not contain any accents at the moment so as not to cause
% problems---this maybe should be changed though. 
%    \begin{macrocode}
\newcommand{\ADS@journalname}{%
  \if@proceedings\@conftitle\space (Eds. \@editors) \fi%
  Revista Mexicana de Astronom\'\i{}a y Astrof\'\i{}sica%
  \if@proceedings\space (Serie de Conferencias)\fi%
  \space Vol. \@volume, pp. \ADS@firstpage--\ADS@lastpage\space (\@year)%
  \space (http://www.astroscu.unam.mx/\~{}rmaa/)
  }
%    \end{macrocode}
% \end{macro}
% \begin{macro}{\ADS@keywords}
% \begin{macro}{\ADS@abstract}
% \begin{macro}{\ADS@title}
% |\ADS@abstract| is straightforward, since we already have the text
% saved as |\@abstract|. We could do the same with |\ADS@keywords| but
% to avoid unwanted formatting we explicitly set it in |\keywords|
% (p.~\pageref{keywords}). |\ADS@title| is just set to |\@rawtitle|,
% which we were already using for another purpose. 
%    \begin{macrocode}
\let\ADS@title=\@rawtitle
\newcommand{\ADS@keywords}{No keywords specified}
\newcommand{\ADS@abstract}{No abstract specified}
%    \end{macrocode}
% \end{macro}\end{macro}\end{macro}
% \begin{macro}{\ADS@authors}
% |\ADS@authors| wants to be in ``Other, A. N.'' format, but with
% semi-colons between authors. The most natural way to do this seems to
% be with |\indexauthor| (p.~\pageref{indexauthor}). We append each
% author to this macro, which means we must explicitly clear it for
% each new article (done in |\article|). .   
%    \begin{macrocode}
\newcommand{\ADS@authors}{}
%    \end{macrocode}
% \end{macro}
% \begin{macro}{\ADS@fulldate}
% \begin{macro}{\ADS@lastpage}
% \begin{macro}{\ADS@copyright}
% \begin{macro}{\ADS@website}
% \begin{macro}{\ADS@url}
% \begin{macro}{\SetWebsite}
% \begin{macro}{\SetMonthNumeric}
% Other sundry data, together with user commands for bits that can be
% set by the Editors. The |\ADS@url| is constructed from the website
% (currently |http://www.astrosmo.unam.mx/rmaa/|,
% the journal descriptor (e.g., |RMxAC..10/|), followed by the
% path to the PDF file (e.g., |PDF/RMxAC..10_HEW1.pdf|). 
% 
%    \begin{macrocode}
\newcommand{\ADS@website}{http://www.astrosmo.unam.mx/rmaa/}
\newcommand{\ADS@url}
{\ADS@website\ADS@journal\ADS@volume\ADS@number/PDF/\ADS@journal\ADS@volume\ADS@number\_\articlename.pdf}
\newcommand{\ADS@number}{}
\newcommand{\ADS@fulldate}{00/\ADS@year} 
\newcommand{\ADS@lastpage}{00} 
\newcommand{\ADS@copyright}{(C) \ADS@year; Instituto de Astronom\'\i{}a,
  Universidad Nacional Aut\'onoma de M\'exico} 
\newcommand{\SetNumber}[1]{\renewcommand{\ADS@number}{-#1}}
\newcommand{\SetWebsite}[1]{\renewcommand{\ADS@website}{#1}}
\newcommand{\SetMonthNumeric}[1]{\renewcommand{\ADS@fulldate}{#1/\ADS@year}}
%    \end{macrocode}
% \end{macro}\end{macro}\end{macro}\end{macro}\end{macro}\end{macro}\end{macro}
% \begin{macro}{\ADS@affil}
% \begin{macro}{\ADS@references}
% These two are somewhat more problematical. Affiliation currently
% only works for the first author (see \dots) pending redesign of the
% user interface. References are now done by wrapping the
% |thebibliography| environment to open another input stream on the
% same file and search through for the |\bibitems|. 
%    \begin{macrocode}
\newcommand{\ADS@references}{}
\newcommand{\ADS@affil}{}
%    \end{macrocode}
% \end{macro}\end{macro}
% 
% Define two macros that expand to the strings
% |\begin{thebibliography}| and |\end{thebibliography}|. Note that the
% idea for this and much else is inspired by the code for the docstrip
% package. Also define |\StartBibX| in case authors use the more
% standard |\begin{thebibliography}{}| syntax. 
%    \begin{macrocode}
\edef\EndBib{\expandafter\@gobble\string\\end\string{thebibliography\string}}
\edef\StartBib{\expandafter\@gobble\string\\begin\string{thebibliography\string}}
\edef\StartBibX{\expandafter\@gobble\string\\begin\string{thebibliography\string}\string{\string}}
%    \end{macrocode}
% Define two new logical variables. |\ifAnotherline| is true if we
% want to read another line from the input stream. It is set to false
% when we encounter the |\end{thebibliography}|. |\ifInBib| is set to
% true when we frind the |\begin{thebibliography}| and tells us to
% start saving the following lines. 
%    \begin{macrocode}
\newif\ifAnotherLine
\newif\ifInBib
%    \end{macrocode}
% \begin{macro}{makeOther}
% We need a macro to change the catcodes of lots of things to other. 
%    \begin{macrocode}
\def\makeOther#1{\catcode`#1=12\relax}
%    \end{macrocode}
% \end{macro}
% \begin{macro}{\RereadBiblio}
% Define a macro to do all the fiddly stuff. This is used at the
% beginning of |thebibliography| environment (Sec~\ref{sec:biblio}).
%    \begin{macrocode}
\newread\tf@biblio
\newcommand{\RereadBiblio}{%
%    \end{macrocode}
% Open the current file again for reading
%    \begin{macrocode}
  \begingroup
  \immediate\openin\tf@biblio=\articlename
%    \end{macrocode}
% Everything in a group to be safe. Initialize logicals. Black magic
% to change
% catcodes and avoid blank lines being interpreted as |\par|. 
%    \begin{macrocode}
  \AnotherLinetrue\InBibfalse
  \makeOther\ \makeOther\\\makeOther\$%
  \makeOther\#\makeOther\^\makeOther\^^K%
  \makeOther\_\makeOther\^^A\makeOther\%%
  \makeOther\~\makeOther\{\makeOther\}\makeOther\&%
  \endlinechar-1\relax
%    \end{macrocode}
% Start looping through the file. Read each line to the macro
% |\Line|. If we have got to EOF or end of biblio then set
% |\AnotherLinefalse|.
%    \begin{macrocode}
  \loop
  \read\tf@biblio to \Line
  \ifx\Line\EndBib\AnotherLinefalse\InBibfalse\fi
  \ifeof\tf@biblio \AnotherLinefalse\fi
%    \end{macrocode}
% If we are in the biblio and the line is not blank then append it to
% the |\ADS@references| macro together with a newline. 
%    \begin{macrocode}
  \ifInBib
  \ifx\Line\empty\else\xdef\ADS@references{\ADS@references\Line^^J}\fi
  \fi
%    \end{macrocode}
% Check if the line we just saw was the start of biblio, then repeat
% the loop if necessary. 
%    \begin{macrocode}
  \ifx\Line\StartBib\InBibtrue\fi
  \ifx\Line\StartBibX\InBibtrue\fi
  \ifAnotherLine\repeat
  \immediate\closein\tf@biblio
  \endgroup
  }
%    \end{macrocode}
% \end{macro}
% 
% \begin{macro}{\PadToFourWithDots}
% Utility function for left-padding a number to 4 characters with
% periods. Used in 
% constructing |\ADS@BibCode|. This is a very simple algorithm and
% \emph{requires} that the argument actually be a number---if it is
% not then dire consequences will result.  
%    \begin{macrocode}
\newcommand{\PadToFourWithDots}[1]{%
  \ifnum #1 > 999 #1%
  \else\ifnum #1 > 99 .#1%
  \else\ifnum #1 > 9 ..#1%
  \else ...#1%
  \fi\fi\fi}
%%  \ifnum \count0 > 999 #1%
%%  \else\ifnum \count0 > 99 .#1%
%%  \else\ifnum \count0 > 9 ..#1%
%%  \else ...#1%
%%  \fi\fi\fi}
%    \end{macrocode}
% \end{macro}
%
% \begin{macro}{\AuthorADS}
% \begin{macro}{\ADS@affiliations}
% New mechanism for setting the affiliations for ADS. The author labels (AA, AB, \dots) are generated using two nested counters.  We use the |\robustify| command from \texttt{etoolbox} so that accented characters in the affiliation addresses do not blow up.  The list of accents may need to be expanded. 
%    \begin{macrocode}
\newcounter{AUTHOR}
\newcounter{author}[AUTHOR]
\stepcounter{AUTHOR}
\newcommand\ADS@affiliations{}
\robustify{\'}\robustify{\i}\robustify{\~}\robustify{\`}
\def\AuthorADS#1#2{%
  \begingroup
  \def\ORCIDiD##1{<ID system="ORCID">##1</ID>}
  \def\Email##1{<EMAIL>##1</EMAIL>}
  \def\Affil##1{##1}
  \stepcounter{author}
  \ifthenelse{\value{author} > 26}{\stepcounter{AUTHOR}\stepcounter{author}}{}
  \ifthenelse{\equal{\ADS@affiliations}{\@rmaaEmpty}}{%
    \xdef\ADS@affiliations{\Alph{AUTHOR}\Alph{author}(#2)}
  }{%
    \xdef\ADS@affiliations{\ADS@affiliations, \Alph{AUTHOR}\Alph{author}(#2)}
  }%
  \endgroup
}
%    \end{macrocode}
% \end{macro}
% \end{macro}
% %
% \begin{macro}{\ADS@dumpdata}
% \changes{4.0}{2019/11/13}{Write all author affiliations for ADS}
% \begin{macro}{\ADS@write}
% \begin{macro}{\ADS@openfile}
% |\ADS@dumpdata| is the macro that actually does the job of writing out all the
% data to an auxiliary file, using |\ADS@write| to write each
% field. The file is set up with |\ADS@openfile|, which currently is
% done in the TOC of a book.  The new multi-author affiliations are used if |\ADS@affiliations| has been set, otherwise it falls back to the single-author version |\ADS@affil|. 
%    \begin{macrocode}
\newcommand{\ADSfield}[2]{\relax}
\newcommand{\ADSvfield}[2]{\relax}
%%\renewcommand{\ADSfield}[12]{\%#1 #2}
\newcommand{\ADS@dumpdata}
{ 
  \ADS@write{\string\begin{ADSarticle}}
  \ADS@WRITE{\string\ADSfield{R}{\ADS@BibCode}}
  \ADS@write{\string\ADSfield{A}{\ADS@authors}}
  \ifthenelse{\equal{\ADS@affiliations}{\@rmaaEmpty}}{%
    {\ADS@sanitize\ADS@write{\string\ADSfield{F}{AA(\ADS@affil)}}}
  }{%
    {\ADS@sanitize\ADS@write{\string\ADSfield{F}{\ADS@affiliations}}}
  }
  \ADS@write{\string\ADSfield{J}{\ADS@journalname}}
  \ADS@write{\string\ADSfield{D}{\ADS@fulldate}}
  \ADS@write{\string\ADSfield{L}{\ADS@lastpage}}
  {\ADS@sanitize\ADS@write{\string\ADSfield{T}{\ADS@title}}}
  \ADS@write{\string\ADSfield{C}{\ADS@copyright}}
  \ADS@write{\string\ADSfield{d}{\ADS@doi}}
  \ADS@write{\string\ADSfield{U}{\ADS@url}}
  \ADS@write{\string\ADSfield{K}{\ADS@keywords}}
  {\ADS@sanitize\ADS@write{\string\ADSfield{B}{\ADS@abstract}}}
  \ADS@write{\string\ADSfield{Z}{\ADS@references}}
  \ADS@write{\string\end{ADSarticle}}
  }
%%\newcommand{\ADS@write}[1]{\if@book\addtocontents{ads}{#1}\fi}
%%\newcommand{\ADS@openfile}{\@starttoc{ads}}
\newcommand{\ADS@write}[1]{\set@display@protect\immediate\write\tf@ads{#1}}
\newcommand{\ADS@WRITE}[1]{\immediate\write\tf@ads{#1}}
\newcommand{\ADS@openfile}{%
  \newwrite\tf@ads%
  \immediate\openout\tf@ads\jobname.ads\relax
  \ADS@WRITE{\ADS@preamble}}
\newcommand{\ADS@closefile}{\ADS@WRITE{\ADS@postamble}}
\def\ADS@preamble{%
  %%  \string\documentclass{article}
  %%  \string\begin{document}
  }
\def\ADS@postamble{%
  %%  \string\end{document}
  \string%\string% Local Variables:^J 
  \string%\string% mode: LaTeX^J
  \string%\string% mode: Font-lock^J
  \string%\string% End:^J
  }
%    \end{macrocode}
% \end{macro}
% \end{macro}
% \end{macro}
% 
% \subsection{Automatic calculation of the page range}
% \label{sec:page-range}
% Another issue that has cost me a lot of grief. Still, it finally
% seems to be sorted now. All hail Jeff Goldberg. 
%
% Cunning trick to get the page numbers from |\pageref|.  Put a label
% on the first and last page of the document using |\AtBeginDocument|
% and |\AtEndDocument|. .  This way won't work for the book option.
% In that case, we do it by defining |\@endofarticlehook| and
% |\@begofarticlehook| as parallels to |\AtEndDocument| and
% |\AtBeginDocument|.  Note that we have to defer adding
% |\@begofarticlehook| to |\AtBeginDocument| to make sure it comes
% after the redefinitions of |\ref| and |\label|.
%    \begin{macrocode}
\if@book\AtEndDocument{\ADS@closefile}\else
 %\AtBeginDocument{\@begofarticlehook}
\AtEndDocument{\@endofarticlehook}
\fi
%    \end{macrocode}
% \begin{macro}{\@begofarticlehook}
% Since we are getting more and more things that need to be done at
% the beginning of every article, frontsection, etc we define a sister
% hook to 
% |\@endofarticlehook|. This hook does the following: 
% \begin{enumerate}
% \item Labels the first page
% \item Sets |\ADS@firstpage|
% \item runs and then clears |\indexauthor@deferlist|
% \item resets the keyword counter
% \end{enumerate}
% It is executed in the following circumstances: 
% \begin{enumerate}
% \item At |\begin{document}| for standalone files
% \item At |\begin{document}| for each included file
% \item At beginning of front sections 
% \item At beginning of TOC
% \item At beginning of Author Index
% \item At beginning of abstracts
% \end{enumerate}
% For some reason, this is not always within the right group, so we
% use |\xdef| just to be sure. 
%    \begin{macrocode}
\newcommand{\@begofarticlehook}{%
  \if@frontmatter\frontlabel{RMAAFirstPage}\else\label{RMAAFirstPage}\fi%
  \xdef\ADS@firstpage{\@arabic{\c@page}}%
  \if@book\else\makeindex\fi % needed so default index is defined in stand-alone articles
%%  \indexauthor@deferlist
%%  \xdef\indexauthor@deferlist{}
  \@inpreamblefalse
  \setcounter{kwdcounter}{0}
  }
%    \end{macrocode}
% \end{macro}
% \begin{macro}{\@endofarticlehook}
% This is a hook to be run at the end of each article
% currently, it just labels the last page and writes out all the ADS
% data to an auxiliary file.  
% It should not produce any printed output. 
% It is modelled on the |\AtEndDocument| command, to which it reduces
% in the case of a single article. 
%    \begin{macrocode}
\newcommand{\@endofarticlehook}{%
  \if@debug{\ClassWarning{rmaa}{***Doing Hook***}}\fi
  \if@debug{\ClassWarning{rmaa}{ADSabstract is: \ADS@abstract}}\fi
  \clearpage\rmaa@lastpagelabel\if@book\ADS@dumpdata\fi
  \ifx\NotesList\@rmaaEmpty\else\OutputNotes\fi
  \if@debug{\ClassWarning{rmaa}{Notes list is: \NotesList}}\fi
}
%    \end{macrocode}
% \begin{macro}{\rmaa@lastpagelabel}
% New version of last page label command copied from Jeff Goldberg's
% lastpage.sty package - we can't just use |\label| here in case the
% output routine for the last page has already been called. Instead,
% we first explicitly clear the page, step back the page counter,
% explicitly write to the .aux file, then restep the page counter
% again.
%
% But this will conflict with the hyperref package, which does labels differently.  Compensate for that by using the five-argument form.
% 
% Also set the |\ADS@LastPage| macro to the appropriate number. 
%    \begin{macrocode}
\newcommand{\rmaa@lastpagelabel}{%
  \addtocounter{page}{-1}%
  \immediate\write\@auxout{\string
    \newlabel{\articlename:RMAALastPage}{{}{\thepage}{}{}{}}}%
  \edef\ADS@lastpage{\thepage}%
  \addtocounter{page}{1}}
%    \end{macrocode}
% \end{macro}
% \end{macro}

% \subsection{Deal with explicit change in number of columns}
% \label{sec:columns}
% Sometimes, when formatting a book, we need to switch between
% one-column articles and two-column articles. For this to work, we
% need to wrap the standard \LaTeX{} commands to set the |\@rmtwocol|
% flag appropriately and also reset the package default values of all
% scalable lengths. .  
%    \begin{macrocode}
\let\latexonecolumn=\onecolumn
\let\latextwocolumn=\twocolumn
\renewcommand{\onecolumn}{\@rmtwocolfalse\@package@default@lengths\latexonecolumn}
\renewcommand{\twocolumn}{\@rmtwocoltrue\@package@default@lengths\latextwocolumn}
%    \end{macrocode}
% 
% \subsection{The actual formatting of the title}
% \label{sec:title-format}

% Note, we have to redefine the |\maketitle| command to use the
% numbered footnotes. This means we needn't use the atbeginend stuff. 
%    \begin{macrocode}
\renewcommand{\maketitle}{\par
  \begingroup
%    \end{macrocode}
% Don't switch to marked footnotes. 
% This is consistent with existing style but I'm not sure I like it
%    \begin{macrocode}
  %%  \renewcommand\thefootnote{\fnsymbol{footnote}}%
%%  \def\@makefnmark{\hbox to\z@{$\m@th^{\@thefnmark}$\hss}}%
  \def\@makefnmark{\textsuperscript{\@thefnmark}\hspace{-0.2em}}%
  \long\def\@makefntext##1{\parindent 1em\noindent
    \hbox to1.8em{\hss$\m@th^{\@thefnmark}$}##1}%
  \if@rmtwocol
  { \if@debug{\ClassWarning{rmaa}{TWO COLUMN}}\fi
    \twocolumn[\@maketitle]%
    }
  \else
  { \if@debug{\ClassWarning{rmaa}{ONE COLUMN}}\fi
    \newpage
    \global\@topnum\z@   % Prevents figures from going at top of page.
    \if@manuscript \@MSmaketitle \else \@maketitle \fi
    }
  \fi
  \if@manuscript\else\thispagestyle{rmaafirstpage}\@thanks\fi
  \ifx\NotesList\@rmaaEmpty\else{%
    \NoteMark{1--\arabic{notecounter}}\vspace*{-2\baselineskip}%
  }\fi%
  \endgroup
%    \end{macrocode}
% Don't reset footnote counter (not sure whether this is good thing
% or not).
%    \begin{macrocode}
  %%  \setcounter{footnote}{0}%
  \let\thanks\relax\let\affil\relax\let\altaffil\relax
  \let\maketitle\relax\let\@maketitle\relax
  \gdef\@thanks{}\gdef\@author{}\gdef\@title{}}
%    \end{macrocode}

% \begin{macro}{\@kwhwidth}
% \begin{macro}{\@kwtwidth}
% Redefine \cs{@maketitle} to include the abstract, resumen and
% keywords. Note that in the manuscripty style we use
% \cs{@MSmaketitle} instead.  
% \changes{3.16}{2003/04/07}{
%  Added extra vertical space between authors and accepted date.}
% \changes{3.26}{2011/08/15}{Add extra label to review articles}
%    \begin{macrocode}
\newlength{\@kwhwidth}\newlength{\@kwtwidth}
\newcommand\reviewname{Review}
\renewcommand{\@maketitle}{%
  \newpage
  \null
  \if@reviewarticle
  \smallskip
  {\large\bfseries\itshape\reviewname}%
  \bigskip
  \else
  \vskip\@before@title@skip%
  \fi
  \begin{center}%
    {\@title \par}%
    \vskip\@before@author@skip%
    \parbox{\rmaa@AuthorWidthFrac\textwidth}{\centering
      \lineskip\@between@author@skip%
      \@author
      \par}
    %% this should be replaced with `Submitted to...' for preprints
    \if@printdates
    \vskip\@before@dates@skip%
    \if@preprint
    \rmaa@ReceivedStyle{Submitted to RMxAA \@received}
    \else
    \rmaa@ReceivedStyle{Received \@received; accepted \@accepted}
    \fi
    \fi
    \if@manuscript
    \rmaa@ReceivedStyle{Submitted to RMxAA \@received}
    \fi
    \if@guide
    \rmaa@ReceivedStyle{Version \rmaa@version, \rmaa@date}
    \fi
    \vskip\@before@abstract@skip%
    \if@abstractpresent
    \parbox{\rmaa@AbstractWidthFrac\textwidth}{%
      \setlength{\parindent}{0.9cm}%
      \centerline{\rmaa@AbstractHeaderStyle{Abstract}}\par%
      \vskip\@after@abstracthead@skip%
      \if@newstyle\noindent\fi\@abstract
      \vskip\@between@abstract@skip%
      \centerline{\rmaa@AbstractHeaderStyle{Resumen}}\par%
      \vskip\@after@abstracthead@skip%
      \if@newstyle\noindent\fi\@resumen
      \if@setresumen\else\@abstract\fi % append english abstract if
                                       % spanish one not set
      \if@guide\else            %do not use keywords for guide
      \vskip\@before@kwd@skip%
      % Format the Key Words with hanging indentation:
      % They are put in a \parbox with width equal to the line width
      % minus the width of the header. 
      { % Widths of keyword
        \settowidth{\@kwhwidth}{\@keywordheader~}   % header and
        \setlength{\@kwtwidth}{\linewidth}          % keyword text
        \addtolength{\@kwtwidth}{-\@kwhwidth}       %
        \noindent\@keywordheader~\parbox[t]{\@kwtwidth}{%
          \rmaa@KeywordStyle{\boldmath\uppercase{\@keywords}}}}%
      \fi
      }
      \fi
  \end{center}%
  \vspace*{\afterkwdskip}
  \par
  }
\newlength{\afterkwdskip}\setlength{\afterkwdskip}{\z@}
%    \end{macrocode}
% \end{macro}
% \end{macro}
%
%
% \begin{macro}{\reviewarticle}
% \begin{macro}{\reviewtoc}
% \begin{macro}{\reviewlof}
%   User command to specify a review article. Must be used in
%   preamble. Main effect is to place the text ``Review'' before the
%   article title.  Also add command for setting a table of contents
%   and a list of figures at the beginning of the article.
%    \begin{macrocode}
\newcommand{\reviewarticle}{\@reviewarticletrue}
\newcommand\reviewtoc{%
  \section*{CONTENTS}%
  \@starttoc{rtc}%
}
\newcommand\reviewlof{%
  \section*{LIST OF FIGURES}%
  \@starttoc{lof}%
}
%    \end{macrocode}
% \end{macro}
% \end{macro}
% \end{macro}
% \begin{macro}{\l@section}
%    \begin{macrocode}
\renewcommand*\l@section[2]{%
  \ifnum \c@tocdepth >\z@
    \addpenalty\@secpenalty
    \addvspace{0.8em \@plus\p@}%
    \setlength\@tempdima{1.5em}%
    \begingroup
      \parindent \z@ \rightskip \@pnumwidth
      \parfillskip -\@pnumwidth
      \leavevmode %\bfseries
      \advance\leftskip\@tempdima
      \hskip -\leftskip
      #1\nobreak\hfil \nobreak\hb@xt@\@pnumwidth{\hss #2}\par
    \endgroup
  \fi}
%    \end{macrocode}
% \end{macro}
%
% \begin{macro}{\@MSmaketitle}
% A special version of |\@maketitle| for the \texttt{manuscript}
% option. This can deal with a multi-page abstract/resumen. Not very
% well, mind. 
%    \begin{macrocode}
\newcommand{\@MSmaketitle}{%
  \if@debug\ClassInfo{rmaa}{Using special MS maketitle command}\fi
  \newpage
  {
    \thispagestyle{rmaafirstpage}
%     \setlength{\parindent}{0pt}
    \null\vskip 2ex
    \begin{center}
      \@title\par
      \@author\par
      \rmaa@ReceivedStyle{%
        \ifx\@received\@rmaaEmpty
          Draft version: \@date
        \else
          Received: \@received
          \ifx\@accepted\@rmaaEmpty
          \else
            ; Accepted: \@accepted
          \fi
        \fi}\par
    \end{center}
    \@thanks
    \centerline{\rmaa@AbstractHeaderStyle{Resumen}}
    \noindent\@resumen\par\vskip 2ex
    \centerline{\rmaa@AbstractHeaderStyle{Abstract}}
    \noindent\@abstract\par\vskip 2ex
    \noindent\@keywordheader~\rmaa@KeywordStyle{\@keywords}\par
%     \@keywordheader~\rmaa@KeywordStyle{\uppercase{\@keywords}}\par
  }%
  }
%    \end{macrocode}
% \end{macro}

% \subsection{The sectioning commands}
% \label{sec:section-commands}

% \begin{macro}{\section}
% \begin{macro}{\subsection}
% \begin{macro}{\subsubsection}
% \begin{macro}{\paragraph}
% \begin{macro}{\subparagraph}
% Redefine all the standard commands to use these lengths. 
%    \begin{macrocode}
\renewcommand\section{\@startsection {section}{1}{\z@}%
                                   {\abovesectionskip}%
                                   {\belowsectionskip}%
                                   %% how do we get upper case here
                                   {\centering\reset@font\normalsize}}
\renewcommand\subsection{\@startsection{subsection}{2}{\z@}%
                                     {\abovesectionskip}%
                                     {\belowsectionskip}%
                                     {\centering\reset@font\normalsize\itshape}}
\renewcommand\subsubsection{\@startsection{subsubsection}{3}{\z@}%
                                     {\abovesectionskip}%
                                     {\belowsectionskip}%
                                     {\centering\reset@font\normalsize\itshape}}
\renewcommand\paragraph{\@startsection{paragraph}{4}{\z@}%
                                    {\aboveparaskip}%
                                    {\belowparaskip}%
                                    {\reset@font\normalsize\bfseries}}
\renewcommand\subparagraph{\@startsection{subparagraph}{5}{\parindent}%
                                      {\aboveparaskip}%
                                      {\belowparaskip}%
                                      {\reset@font\normalsize\bfseries}}
%    \end{macrocode}
% \end{macro}
% \end{macro}
% \end{macro}
% \end{macro}
% \end{macro}

% \begin{macro}{secnumdepth}
% Make sure that we number the subsubsections---this is done by
% default in the |article| class but not the |book| class, which is
% why we must set it explicitly here.
%    \begin{macrocode}
\setcounter{secnumdepth}{3}
%    \end{macrocode}
% \end{macro}
% \changes{2.2beta10}{2001/11/07}{Fixed \cs{secnumdepth} for \cs{book}
% style.}

% \begin{macro}{\thesubsection}
% \begin{macro}{\thesubsubsection}
% For subsection and subsubsection, we want the numbers to be roman
% but the title to be italic. This seems the easiest way to go about it.
%    \begin{macrocode}
\renewcommand{\thesubsection}{{\upshape\thesection.\@arabic\c@subsection}}
\renewcommand{\thesubsubsection}{\thesubsection{\upshape.\@arabic\c@subsubsection}}
%    \end{macrocode}
% \end{macro}
% \end{macro}

% \begin{macro}{\@seccntformat}
% Add ``.'' to the number here so we can remove it from |\thesection|,
% etc. Reduce gap between number and title WJH 27 Jan 1995
%    \begin{macrocode}
\def\@seccntformat#1{\csname the#1\endcsname .\hskip 0.4em}
%    \end{macrocode}
% \end{macro}

% \begin{macro}{\@sect}
% I didn't really want to mess with the |\@sect| command, but it was
% the only way I could manage to force upper case. Other than the
% commented parts, it is identical to the standard version. 
%    \begin{macrocode}
\def\@sect#1#2#3#4#5#6[#7]#8{
  \ifnum #2>\c@secnumdepth
     \let\@svsec\@empty
  \else
     \refstepcounter{#1}%
     \let\@@protect\protect
     \def\protect{\noexpand\protect\noexpand}%
     \edef\@svsec{\@seccntformat{#1}}%
     \let\protect\@@protect
  \fi
%    \end{macrocode}
% This is the bit I changed (27 Jan 1995). Use upper case for
% |\section|, use mixed case for the rest. 
%    \begin{macrocode}
  \ifnum #2=1                   
     \def\sec@upcase##1{\uppercase{##1}}%
     \def\sec@shape{\upshape}%
  \else                         
     \def\sec@upcase##1{\relax{##1}}%
     \def\sec@shape{\itshape}%
  \fi
  \@tempskipa #5\relax
  \ifdim \@tempskipa>\z@
     \begingroup #6\relax
          \@hangfrom{\hskip #3\relax\@svsec}%
                    {\interlinepenalty \@M \sec@upcase{#8}\par}%
     \endgroup
     \csname #1mark\endcsname{#7}%
%    \end{macrocode}
% Only write TOC entries if doing the Guide style. 
%    \begin{macrocode}
     \if@guide 
        \addcontentsline{toc}{#1}{%
          \ifnum #2>\c@secnumdepth \else
             \protect\numberline{\csname the#1\endcsname}%
          \fi
          #7}%
     \fi
%    \end{macrocode}
% Also allow per-article TOC for review papers
%    \begin{macrocode}
     \if@reviewarticle
        \ifnum #2>2 \else
           \addcontentsline{rtc}{#1}{%
             \protect\numberline{\csname the#1\endcsname}%
             \sec@upcase{\sec@shape #7}}%
        \fi
     \fi
  \else
     \def\@svsechd{#6\hskip #3\relax  %% \relax added 2 May 90
       \@svsec #8\csname #1mark\endcsname{#7}%
       \if@guide
          \addcontentsline{toc}{#1}{%
            \ifnum #2>\c@secnumdepth \else
               \protect\numberline{\csname the#1\endcsname}%
            \fi
            #7}%
       \fi
       }%
  \fi
  \@xsect{#5}}
%    \end{macrocode}
% \end{macro}


% \subsection{Appendices}
% \label{sec:appendices}
% Make sure that appendices in individual articles work properly. We
% copy the definition from the article class (the book one doesn't do
% what we want).  
% \begin{environment}{appendices}
%   This uses an environment |appendices| inside of which each
%   |\section| will become a separate appendix.  Note that the
%   appendices must come before the references (otherwise the full
%   addresses mechanism will not work properly). We make the equation
%   numbering be subordinate to the appendix (equations A12, B6,
%   etc). Finally, we have a medskip followed by APPENDICES to
%   introduce them. \changes{3.17}{2002/04/08}{ Added appendices environment
%   that does the right thing.}
%    \begin{macrocode}
\newenvironment{appendices}{\par
  \begingroup
  \setcounter{section}{0}%
  \setcounter{subsection}{0}%
  \renewcommand\thesection{\@Alph\c@section}
  \renewcommand\theequation{\thesection\arabic{equation}}
  \medskip
  \section*{APPENDICES}}
{\endgroup}
%    \end{macrocode}
% \end{environment}
% \begin{environment}{appendix}
%   If there is only one appendix, then use |appendix| environment
%   instead.  \changes{4.0}{2019/11/12}{ Added
%   appendix environment for case where there is only one section.}
%    \begin{macrocode}
\renewenvironment{appendix}{\par
  \begingroup
  \setcounter{section}{0}%
  \setcounter{subsection}{0}%
  \renewcommand\thesection{\@Alph\c@section}
  \renewcommand\theequation{\thesection\arabic{equation}}
  \medskip
  \section*{APPENDIX}}
{\endgroup}
%    \end{macrocode}
% \end{environment}

% \subsection{Floating thingies}
% \label{sec:floats}   

%% 
%    \begin{macrocode}
\renewcommand{\thetable}{\@arabic\c@table}
\renewcommand{\thefigure}{\@arabic\c@figure}
\renewcommand{\figurename}{Fig.}
\renewcommand{\tablename}{TABLE}
%\renewcommand{\fnum@table}{\MakeUpperCase{\tablename}~\thetable}
%    \end{macrocode}
% For some reason, I have to write `TABLE' explicitly here
% If I use |\tablename|, it gets converted to mixed case---God knows
% why. 
% |\figurename| is also being ignored. 
%    \begin{macrocode}
\renewcommand{\fnum@table}{TABLE~\thetable}
\renewcommand{\fnum@figure}{Fig.~\thefigure}
%    \end{macrocode}

% Do captions in 9 pt face in the main journal style and in new-style
% anything. 
%    \begin{macrocode}
\long\def\@makefigcap#1#2{%
  \if@proceedings\else\small\fi
  \if@newstyle\small\fi 
  \vskip\abovecaptionskip
  \sbox\@tempboxa{#1. #2}%
  \ifdim \wd\@tempboxa >\hsize
    #1. #2 \par
  \else
    \global \@minipagefalse
    \hb@xt@\hsize{\hfil\box\@tempboxa\hfil}%
  \fi
  \vskip\belowcaptionskip}
%    \end{macrocode}

% Make sure table caption is upper case. 
%    \begin{macrocode}
\long\def\@maketabcap#1#2{%
  \hbox to\hsize{\hfil #1 \hfil}
  \vskip 6pt 
  \centering\MakeUpperCase{#2}\par
  \vskip 6pt}
%    \end{macrocode}

%    \begin{macrocode}
\newlength{\@rottabcapwidth}
\long\def\@makerottabcap#1#2{%
  \setbox\@tempboxa\hbox{\MakeUpperCase{#2}}%
  \ifdim\wd\@tempboxa >.8\vsize
     \setlength\@rottabcapwidth{.8\vsize}
  \else
     \setlength\@rottabcapwidth{\wd\@tempboxa}
  \fi
  \rotatebox{90}{%
    \begin{minipage}{\vsize}\centering
      \begin{minipage}{\@rottabcapwidth}\centering
        #1\\ \vskip 6pt \MakeUpperCase{#2}
      \end{minipage}%
    \end{minipage}
    }%
  \hskip 6pt
  }

\newenvironment{Sideways}
    {\begin{sideways}\begin{minipage}{\vsize}\centering}
    {\end{minipage}\end{sideways}}

\newenvironment{Sidewaystable}
{%
  \let\oldcaption=\caption%
  \def\caption{
    \def\@captype{table}\renewcommand{\@makecaption}{\@maketabcap}
    \oldcaption}
  \begin{sidewaystable}\centering}
  {\end{sidewaystable}\let\caption=\oldcaption}
\newenvironment{Sidewaystable*}
{%
  \let\oldcaption=\caption%
  \def\caption{
    \def\@captype{table}\renewcommand{\@makecaption}{\@maketabcap}
    \oldcaption}
  \begin{sidewaystable*}\centering}
  {\end{sidewaystable*}\let\caption=\oldcaption}
%    \end{macrocode}

% Not sure why we need to stick these lines in. 
%    \begin{macrocode}
\renewenvironment{figure}
                 {\renewcommand{\@makecaption}{\@makefigcap}\@float{figure}}
                 {\end@float}
\renewenvironment{figure*}
                 {\renewcommand{\@makecaption}{\@makefigcap}\@dblfloat{figure}}
                 {\end@dblfloat}
\renewenvironment{table}
                 {\renewcommand{\@makecaption}{\@maketabcap}\@float{table}}
                 {\end@float}
\renewenvironment{table*}
                 {\renewcommand{\@makecaption}{\@maketabcap}\@dblfloat{table}}
                 {\end@dblfloat}
%    \end{macrocode}

% Redefine |\@caption| so that we don't bother with the list of
% figures. This saves us having to worry about fragile commands in the
% captions. 
%    \begin{macrocode}
\long\def\@caption#1[#2]#3{%
  \par
%%  \addcontentsline{\csname ext@#1\endcsname}{#1}%
%%    {\protect\numberline{\csname the#1\endcsname}{\ignorespaces #2}}%
  \begingroup
    \@parboxrestore
    \if@minipage
      \@setminipage
    \fi
    \normalsize
    \@makecaption{\csname fnum@#1\endcsname}{\ignorespaces #3}\par
  \endgroup}
%    \end{macrocode}

% Make tables not quite so cramped. Allow |\hline\hline| to give a
% poor man's thick rule. 
%    \begin{macrocode}
\renewcommand{\arraystretch}{1.15}
\setlength{\doublerulesep}{0.4pt}
%    \end{macrocode}

% \begin{macro}{\tablecols}
% \begin{macro}{\@tablecols}
% \begin{macro}{\tabcaptioncontinued}
% |\tablecols| is a user command for setting the number of columns. It
% is used by commands that try and span the entire table (I think
% there must be a better way of acheiving
% this). |\tabcaptioncontinued| is 
% a user command for setting the caption of a continued table on the
% next page. 
%    \begin{macrocode}
\newcommand{\tablecols}[1]{\renewcommand{\@tablecols}{#1}}
\newcommand{\@tablecols}{1}
\newcommand{\tabcaptioncontinued}{%
  \multicolumn{\@tablecols}{c}{\addtocounter{table}{-1}TABLE~\thetable\ (CONTINUED)} \\ \\
  }
%    \end{macrocode}
% \end{macro}
% \end{macro}
% \end{macro}

% \subsubsection{Footnotes in tables}
% \label{sec:table-footnotes}
% \begin{macro}{\tablenotemark}
% \begin{macro}{\tablenote}
% \begin{macro}{\tablenoteskip}
% Old user commands |\tablenotemark{MARK}|,
% |\tablenote{MARK}{TEXT}|---now deprecated. 
%    \begin{macrocode}
\newcommand\tablenotemark[1]{$^\mathrm{#1}$} 
\newcommand\tablenote[2]{\if@proceedings\else\small\fi ${}^\mathrm{#1}$ #2} 
\newcommand\tablenoteskip{\\[0.5\baselineskip]} 
%    \end{macrocode}
% \end{macro}
% \end{macro}
% \end{macro}

% \begin{macro}{\tabnotemark}
% \begin{macro}{\tabnotetext}
% \begin{macro}{\tabnotetextTWO}
% Alternative way of doing footnotes in tables. 
%    \begin{macrocode}
\newcommand{\tabnotemark}[1]{\textsuperscript{\lowercase{#1}}}
%\newcommand{\tabnotetextTWO}[4]{\tabnotetext{#1}{#2}\tabnotetext{#3}{#4}}
\newcommand{\tabnotetextTWO}[4]{%
  \multicolumn{\@tablecols}{l}{%
    \parbox{0.5\tabnotewidth}{\small\textsuperscript{#1}#2}%
    \parbox{0.5\tabnotewidth}{\small\textsuperscript{#3}#4}%
    }%
  \\}
\newcommand{\tabnotetext}[2]{%
  \multicolumn{\@tablecols}{@{}p{\tabnotewidth}@{}}{\small\textsuperscript{#1}#2}\\}
\newlength{\tabnotewidth}
%    \end{macrocode}
% \end{macro}
% \end{macro}
% \end{macro}


% \begin{macro}{changemargin}
%    \begin{macrocode}
% Command to allow margins to be changed on the fly. Useful to allow
% tables and Figures to protrude into the margins. 
%
% \changes{3.24}{2007/09/14}{Added \texttt{changemargin} environment
% (copied from TeX FAQ).}
\newenvironment{changemargin}[2]{% 
  % From TeX FAQ http://www.tex.ac.uk/cgi-bin/texfaq2html?label=chngmargonfly
  \begin{list}{}{%
    \setlength{\topsep}{0pt}%
    \setlength{\leftmargin}{#1}%
    \setlength{\rightmargin}{#2}%
    \setlength{\listparindent}{\parindent}%
    \setlength{\itemindent}{\parindent}%
    \setlength{\parsep}{\parskip}%
  }%
  \item[]}{\end{list}}
%    \end{macrocode}
% \end{macro}

% \subsection{The bibliography and references}
% \label{sec:biblio}
% \begin{macro}{\beforebiblioskip}
% \begin{macro}{\afterbiblioskip}
% I'm not entirely sure why I introduced these two commands, since
% directly putting |\vspace| before and after
% |\begin{thebibliography}| seems to work well enough. 
%    \begin{macrocode}
\newlength{\beforebiblioskip}
\setlength{\beforebiblioskip}{0pt}
\newlength{\afterbiblioskip}
\setlength{\afterbiblioskip}{0pt}
%    \end{macrocode}
% \end{macro}
% \end{macro}
% \begin{macro}{thebibliography}
% The environment itself is done as starred section with extra
% vertical space after heading. This space is currently zero. 
% In the book style we do the cunning trick with rereading the file in
% order to get the references into the |\ADS@references| macro.
% 
% \changes{3.05}{2003/01/02}{Removed the \cs{beforebiblioskip} and
% \cs{afterbiblioskip} commands, since they cause unwanted vertical
% space in some instances.}  
%    \begin{macrocode}
\renewenvironment{thebibliography}{%
  \if@book\RereadBiblio\fi
%%  \vspace*{\beforebiblioskip}
  \section*{REFERENCES} \frenchspacing
%%  \vspace*{\afterbiblioskip}
  \small
%    \end{macrocode} 
% The references themselves are set as a list in order 
% to get the hanging indentation. They are done in a smaller font and
% with french spacing so we needn't worry about the distinction
% between sentence-ending periods and abbreviations. 
%    \begin{macrocode} 
  \list{\null}{\leftmargin 1.5em\labelwidth\z@\labelsep\z@\itemindent -1.5em
    \itemsep\z@\parsep\z@\usecounter{enumi}}
  \sloppy\clubpenalty4000\widowpenalty4000
  \sfcode`\.=1000\relax\@biblisttrue}
{\endlist\@biblistfalse\normalsize\@outputaddressesifnecessary}
%    \end{macrocode}
% \end{macro}
% \subsection{Uniquification of cross-references}
% \label{sec:uniquification}
% 
% \begin{macro}{\articlename}
% In order that the |\label{}| command in individual articles be
% guaranteed to give a unique anchor, we append the article name (held
% in the |\articlename| command) to each key. This requires
% redefinition of the following commands: |\cite|, |\nocite|, |\ref|,
% |\pageref|, |\label|. 
% \begin{macro}{\oldcite}
% \begin{macro}{\oldnocite}
% \begin{macro}{\oldref}
% \begin{macro}{\oldpageref}
% \begin{macro}{\oldlabel}
% The original versions of these commands are
% saved as |\oldcite|, |\oldnocite|, etc. 
%
% The uniquifying of cross-references is now always done, so 
% we don't get multiply-defined labels in the single article case
%    \begin{macrocode}
\newcommand{\articlename}{anon}
\AtBeginDocument{
  \let\oldcite=\cite 
  \let\oldnocite=\nocite 
  \let\oldref=\ref 
  \let\oldpageref=\pageref
  \let\oldlabel=\label
}
%    \end{macrocode}
% \end{macro}
% \end{macro}
% \end{macro}
% \end{macro}
% \end{macro}
% \end{macro}
%
% Uniquify all cross-references by prepending the article name.  This
% now has to go into |\AtBeginDocument| so it is working on the new
% versions that hyperref.sty has set up.  Then we need to put the
% |\@begofarticlehook| afterwards, so that the first page label is using
% the new definitions.  If amsmath is loaded, then we also need to patch
% the special label command that they use. \changes{4.0}{2019/11/18}{Fix labels for amsmath display equations, in case that package has been loaded.}
%    \begin{macrocode}
\AtBeginDocument{
  \renewcommand{\ref}[1]{\oldref{\articlename:#1}%
    \if@debug{\ClassInfo{rmaa}{Reference to \articlename:#1}}\fi}
  \renewcommand{\pageref}[1]{\oldpageref{\articlename:#1}%
    \if@debug{\ClassInfo{rmaa}{Page reference to \articlename:#1}}\fi}
  \renewcommand{\label}[1]{\oldlabel{\articlename:#1}%
    \if@debug{\ClassInfo{rmaa}{Label for \articlename:#1}}\fi}%
  \ifdef{\label@in@display}{%
    \let\oldlabel@in@display=\label@in@display
    \renewcommand{\label@in@display}[1]{\oldlabel@in@display{\articlename:#1}%
      \if@debug{\ClassInfo{rmaa}{AMSTeX Label for \articlename:#1}}\fi}%
  }{}
  \@begofarticlehook
}%
%    \end{macrocode}
% Citations are more complicated since they can contain a list of
% keys. These are really only here for use with the rmaacite package,
% which is utterly deprecated (02 Jan 2003) now that I have the natbib
% code working. The command \cs{cite} should \emph{not} be used with
% natbib: use \cs{citep}, \cs{citet}, etc.\@ instead. 
%    \begin{macrocode}
\renewcommand{\nocite}[1]{\@bsphack
  \@for\@citeb:=#1\do{%
    \edef\@citeb{\expandafter\@firstofone\@citeb:\articlename}%
    \if@filesw\immediate\write\@auxout{\string\citation{\@citeb}}\fi
    \@ifundefined{b@\@citeb:\articlename}{\G@refundefinedtrue
      \@latex@warning{Citation `\@citeb' undefined}}{}}%
  \@esphack}
\renewcommand{\cite}[1]{\oldcite{#1:\articlename}}
%   \end{macrocode}


% Horrible fiddling with the bibliography citation mechanism. 
% This has to be globally available in case rmaacite puts it in the
% \texttt{.aux} file. 
%    \begin{macrocode}
\def\lbibcite#1#2#3{\global\@namedef{b@#1}{#2}
                    \global\@namedef{by@#1}{#3}}
%    \end{macrocode}

% \begin{macro}{\oldbibitem}
% \begin{macro}{\@rmaaBibArg}
% \begin{macro}{\@rmaaSpace}
% \begin{macro}{rmaabib}
% Of course, |\bibitem| needs to get the uniquification treatment as
% well. While we are at it, we redefine it to use consecutive integers
% for the key if the user leaves it blank (using the counter
% |rmaabib|). This allows one to put 
% |\bibitem{}| for all your references without getting annoying
% messages about multiply-defined labels.
% 
% We also take the opportunity to insert the optional argument to
% \cs{bibitem} if it is missing in the form Author(DATE) so that natbib
% doesn't try to do numerical citations. Natbib already provides a
% hook \cs{@extra@b@citeb} to help uniquify the citation key, so we
% just plug into that. However, it puts it \emph{after} the original
% key, not before as I had had it. Easiest solution is to change the
% way that I do it. 02 Jan 2003. Note that the \cs{@extra@b@citeb}
% does not get written to the .aux  file in the \cs{citation} command
% but it does work. 
% \changes{3.14}{06 Feb 2003}{Reinstalled old definition of \cs{bibitem}
%  for the case where the nonatbib option is requested, so that the
%  rmaacite package will work. This means that you can't use rmaacite
%  and natbib in the same volume.}
%    \begin{macrocode}
\if@natbib\def\@extra@b@citeb{:\articlename}\fi
%%\else
\let\oldbibitem\bibitem
\newcommand{\@rmaaBibArg}{}
\newcommand{\@rmaaSpace}{ }
\newcounter{rmaabib}
\if@natbib
\renewcommand{\bibitem}[2][ABC(2003)]{%
  \renewcommand{\@rmaaBibArg}{#2}
  \ifx\@rmaaBibArg\@rmaaEmpty
     \addtocounter{rmaabib}{1}
     \oldbibitem[#1]{\thermaabib:\articlename}
  \else\ifx\@rmaaBibArg\@rmaaSpace
     \addtocounter{rmaabib}{1}
     \oldbibitem[#1]{\thermaabib:\articlename}
  \else 
     \oldbibitem[#1]{#2:\articlename}
  \fi\fi
  \renewcommand{\@rmaaBibArg}{}
  }
\else
\renewcommand{\bibitem}[1]{%
  \renewcommand{\@rmaaBibArg}{#1}
  \ifx\@rmaaBibArg\@rmaaEmpty
     \addtocounter{rmaabib}{1}
     \oldbibitem{\articlename:\thermaabib}
  \else\ifx\@rmaaBibArg\@rmaaSpace
     \addtocounter{rmaabib}{1}
     \oldbibitem{\articlename:\thermaabib}
  \else
     \oldbibitem{\articlename:#1}
  \fi\fi
  \renewcommand{\@rmaaBibArg}{}
  }
\fi
%%\fi
%    \end{macrocode}
% \end{macro}
% \end{macro}
% \end{macro}
% \end{macro}

% \begin{macro}{oldpcite}
% \begin{macro}{oldscite}
% Then, we need to repeat the above scheme, but for the |rmaacite|
% definition. 
%    \begin{macrocode}
\newcommand{\@redefinebibitem}{%
  \let\oldbibitem\bibitem
  \let\oldcite\cite
  \let\oldpcite\pcite
  \let\oldscite\scite
  \def\bibitem[##1<##2>]##3{%
    \renewcommand{\@rmaaBibArg}{##3}
    \ifx\@rmaaBibArg\@rmaaEmpty
       \addtocounter{rmaabib}{1}
       \oldbibitem[##1<##2>]{\thermaabib:\articlename}
    \else\ifx\@rmaaBibArg\@rmaaSpace
       \addtocounter{rmaabib}{1}
       \oldbibitem[##1<##2>]{\thermaabib:\articlename}
    \else 
       \oldbibitem[##1<##2>]{##3:\articlename}
    % trust the author here
%%       \oldbibitem[##1<##2>]{##3}
    \fi\fi
    }
  }
%    \end{macrocode}
% \end{macro}
% \end{macro}


% \subsection{HTML Table of Contents}
% \label{sec:hyper-TOC}
% \begin{macro}{\HyperTOCEntry}
% \begin{macro}{\HyperTOCpart}
% \begin{macro}{\HyperTOCmeeting}
% \begin{macro}{\HyperTOCeditors}
% \begin{macro}{\HyperTOCpreamble}
% These commands are used to write information to the \texttt{.htoc}
% file, which is used in generating a Web Page of the Table of
% Contents. Since they only ever appear preceded by |\protect| in the
% argument to an |\addtocontents| command,  they should never need to
% be expanded by the macros. However, we define them here just in case
% \dots
%    \begin{macrocode}
\newcommand{\HyperTOCEntry}[3]{\relax}
\newcommand{\HyperTOCpart}[1]{\relax}
\newcommand{\HyperTOCmeeting}[3]{\relax}
\newcommand{\HyperTOCeditors}[1]{\relax}
\newcommand{\HyperTOCpreamble}{\relax}
%    \end{macrocode}
% \end{macro}
% \end{macro}
% \end{macro}
% \end{macro}
% \end{macro}
%

% \subsection{Formatting the poster abstracts} 
% These can now be used without using the book option.
%    \begin{macrocode}

\newcommand{\@makeminititle}{%
  \newpage
  \null 
  \vskip 3em%
  \begin{center}%
    {\@title \par}%
  \end{center}
  \vskip 3em%
  }

\newif\if@abstractsads
\@abstractsadsfalse
\newcommand\abstractsADS{\@abstractsadstrue}

\if@book\RequirePackage{multicol}\fi
\newenvironment{abstracts}[1][Abstracts of Contributed Papers]{%
  %% First, make sure we start on an odd page (06 Jun 2000)
  \if@book\clearpage\thispagestyle{empty}\cleardoublepage\fi
  %%
  \renewcommand{\articlename}{ABSTRACTS}
  \title{#1}\shortauthor{Abstracts}\shorttitle{Abstracts}
  \listofauthors{}
  \if@book
  \addcontentsline{toc}{chapter}{\protect\@fiddleauthors{\@TOCtitle}{\rmaaTOCauthors}}%
  \addtocontents{htc}{\protect\HyperTOCEntry{\@rawtitle}{\@rawauthors}{\articlename}}
  \fi
  \pagestyle{rmaa} % this now seems necessary (06 Jun 2000)
  \thispagestyle{rmaafirstpage}
%%  \twocolumn[\@makeminititle]
  \renewcommand{\abstract}[1]{%
    \renewcommand\ADS@abstract{##1}%
    \renewcommand{\@abstract}{##1}\@writeabstract%
  }
  \renewcommand{\title}[2][]{%
    \gdef\ADS@authors{}%
    \gdef\ADS@affil{}%
    \gdef\ADS@references{}%
    \gdef\ADS@keywords{}%
    \renewcommand{\@rmaaArg}{##1}
    \ifx\@rmaaArg\@rmaaEmpty
       \renewcommand{\@TOCaftertitle}{}
    \else
       \renewcommand{\@TOCaftertitle}{ (\@rmaaArg)}
    \fi
    \renewcommand{\@title}{\rmaa@TOCTitleStyle{##2}}%
    \renewcommand{\@TOCtitle}{\rmaa@TOCTitleStyle{##2}\@TOCaftertitle}%
    \renewcommand{\rmaaIndexTitle}{\rmaa@IndexTitleStyle{##2\@TOCaftertitle.}}%
    \renewcommand\ADS@title{##2}%
    \xdef\ADS@firstpage{\@arabic{\c@page}}%
    }
  \renewcommand{\author}[1]{\let\@thanks\@empty\renewcommand\@author{##1} }
  \renewcommand{\footnotetext}[2][XX]{\par\noindent\textsuperscript{##1}~##2}
  \renewcommand{\footnoterule}{\kern-3\p@ \hrule \@width 3em \kern 2.6\p@}
  \if@book\@begofarticlehook\fi
  \if@book\begin{multicols}{2}[\@makeminititle]\else\@makeminititle\fi
  \begingroup
  \setcounter{footnote}{0}
  }% and at the end of the environment....
  {\endgroup\if@book\end{multicols}\listofabstracts\label{RMAALastPage}\clearpage\fi}
%% auxiliary commands used:
%% \@writeabstract does all the actual formatting 
\newcommand{\@writeabstract}{%
  \if@book%
  \addcontentsline{loa}{abstract}{\protect\@fiddleauthors{\@TOCtitle}{\rmaaTOCauthors}}%
  \fi%
  \begin{center}\@title\par\@author\end{center}
  \@abstract
  \par\noindent\footnoterule\nobreak\footnotesize\@thanks
  \par\noindent\hrulefill\goodbreak
  \vskip 1cm \@plus0.5cm
  %% write ADS entry
  \if@book\if@abstractsads\ADS@dumpdata\fi\fi%
  \endgroup % this pair is to make each abstract a self-contained group
  \begingroup
  \setcounter{footnote}{0}
  }

\if@book
%% List of abstracts (stuff written to .loa file)
\newcommand{\listofabstracts}{%
  \begin{multicols}{2}[%
    \vskip 1.5cm \vfill % leave at least 1.5 cm and go to bottom of page if possible
    \centerline{\MakeUpperCase{List of Abstracts}}
    \vskip 1.5cm \@plus0.5cm
    ][6cm] % this makes it start a new page if there is not at least 6cm of space
    \renewcommand{\@fiddleauthors}[2]{##1\hskip1em plus2em minus0.4em\penalty-100  ##2}
    \parindent\z@
%%    \parskip.5\p@ \@plus .7\p@\relax 
    \parskip3\p@ \@plus 1\p@ \@minus 1\p@\relax % allow quite a lot of
                                % leeway here to avoid breaks inside
                                % entries 
    \renewcommand{\\}{\relax} % temporarily turn off explicit line breaks just in case
    \renewcommand{\altaffilmark}[1]{\relax} % ditto with this
    \@starttoc{loa}
  \end{multicols}
  }
\newcommand*{\l@abstract}[2]{%
  \par\penalty-1000
  #1\hskip 1em\penalty0\null\nobreak\hfill#2}
\fi

%    \end{macrocode}
 
% \subsection{Code for typesetting an entire volume of conference
% proceedings} \label{sec:whole-volume}
%    \begin{macrocode}

\if@book
%% environment for including the individual articles
\newcommand{\@fiddleauthors}[2]{\relax}
\renewcommand{\theequation}{\@arabic\c@equation}
\newenvironment{articles}%
{ \begingroup % make it a group so that author definitions do not
              % spill over
  \gdef\ADS@authors{}
  \pagestyle{rmaa}
  \renewenvironment{document}{%
    \@begofarticlehook
    \addtocounter{chapter}{1}%
    \addcontentsline{toc}{chapter}{%
      \protect\@fiddleauthors{\@TOCtitle}{\rmaaTOCauthors}}%
    \addtocontents{htc}{\protect\HyperTOCEntry{\@rawtitle}{\@rawauthors}{\articlename}}
    }{\relax}
  \renewcommand{\documentclass}[2][foo]{\relax} 
  \renewcommand{\usepackage}[2][foo]{%
    \def\@packagename{##2}\def\@rmaacitename{rmaacite}%
    \ifx\@packagename\@rmaacitename
    \@rmcitetrue\ClassInfo{rmaa}{This paper uses rmaacite}
    \makeatletter\input{rmaacite.sty}\makeatother % read it in here so that all changes
    \else                                         % are local to paper
    \if@debug\ClassWarning{rmaa}{Package \@packagename\space was requested by \articlename\space}\fi
    \fi} 
  %% Don't allow individual articles to mess with the pages, volume or year
  \renewcommand{\SetFirstPage}[1]{\relax}
  \renewcommand{\SetYear}[1]{\relax}
  \renewcommand{\SetVolume}[1]{\relax}
  \renewcommand{\SetConfTitle}[1]{\relax}
  %% Don't allow individual articles to change these
  \renewcommand{\nonstopmode}{\relax}
  \renewcommand{\scrollmode}{\relax}
  \renewcommand{\outputplates}{\relax}
  %% XXXX need to check if using rmaacite and act accordingly
  %% avoid multiple declarations of the same length
  \let\oldnewlength=\newlength
  \renewcommand{\newlength}[1]{%
    \edef\reserved@a{\expandafter\@gobble\string ##1}%
    \@ifundefined{\reserved@a}{%
      \oldnewlength{##1}}{%
      \if@debug{\ClassInfo{rmaa}{\reserved@a\space already defined: doing nothing}}\fi%
      \relax}}%
  \renewcommand\newcommand{\providecommand}
  }%
{\endgroup\onecolumn}
     

%    \end{macrocode}
% Command to be used in `articles' environment for including an
% article
%    \begin{macrocode}
\newcommand{\article}[2][]{%
  % optional first argument is extra text to put in brackets after title in TOC
  % e.g. ``Invited Talk'' or ``Poster''
  \renewcommand{\articlename}{#2}
  % reset loads of counters
  \setcounter{footnote}{0}
  \setcounter{figure}{0}
  \setcounter{table}{0}
  \setcounter{section}{0}
  \setcounter{equation}{0}
  \setcounter{rmaabib}{0}
  \setcounter{author}{0}
  \setcounter{AUTHOR}{1}
  \gdef\ADS@authors{}%
  \gdef\ADS@affil{}%
  \gdef\ADS@affiliations{}%
  \gdef\ADS@references{}%
  \@firstauthortrue
  \@firstaffiltrue
  \@firstaltaffiltrue
  \@firstkwdtrue
  \@rmcitefalse
  \if@suppressfulladdresses\else\@needaddressestrue\fi
  \renewcommand{\@rmaaArg}{#1}
  \ifx\@rmaaArg\@rmaaEmpty
     \renewcommand{\@TOCaftertitle}{}
  \else
     \renewcommand{\@TOCaftertitle}{ (#1)}
  \fi
%    \end{macrocode}
% In the main journal articles must start on an odd page. 
%    \begin{macrocode}
  \if@journal\clearpage\thispagestyle{empty}\cleardoublepage\fi
%% Now input the file, then run the end-of-article hook (reinstated 28
  %% Sep 2001)
  {\input{\@articledir/\articlename.tex}\@endofarticlehook}}
%    \end{macrocode}
% Command to separate articles into sessions. Write an entry in the
% TOC, makes sure we start on an odd page (added 06 Jun 2000), then do
% a title page if \cs{SessionTitlePages} was called. 
% \changes{3.12}{04 Feb 2003}{Added option to have title page for each
% session. Turn it on with \cs{SessionTitlePages}.}
%    \begin{macrocode}
\newcommand{\session}[1]{
  %% this does nothing except write an entry in the TOC
  \addcontentsline{toc}{part}{\rmaa@TOCPartStyle{#1}}
  \addtocontents{htc}{\protect\HyperTOCpart{#1}}
  %% and also makes sure we start on an odd page (added 06 Jun 2000)
  \clearpage\thispagestyle{empty}\cleardoublepage
  \stepcounter{part}
  \if@sessiontitlepages{%
    \onecolumn
    \thispagestyle{empty}
    \renewcommand{\articlename}{Section\thepart TitlePage}
    \label{RMAAFirstPage}
    \make@session@title{#1}
    \label{RMAALastPage}
    \clearpage\thispagestyle{empty}\cleardoublepage
  }\fi
  }
\newcommand{\make@session@title}[1]{%
  {%
    \setlength\parindent{0pt}
    \vspace*{\stretch{1}}
    \rule{\linewidth}{1mm}
    \begin{flushright}\Huge\bfseries\boldmath
      #1
    \end{flushright}
    \rule{\linewidth}{1mm}
    \vspace*{\stretch{2}}
  }
}
%    \end{macrocode}
% Command to turn on separate title pages for each session. 
%    \begin{macrocode}
\newcommand{\SessionTitlePages}{\@sessiontitlepagestrue}
\newcommand{\BoldSessionTitles}{%
  \renewcommand{\rmaa@TOCPartStyle}[1]{\textbf{\rmaa@TOCTitleStyle{##1}}}
  }
\newcommand{\ArticleDir}[1]{\renewcommand{\@articledir}{#1/}\graphicspath{{1/}}}
\newcommand{\@articledir}{.}
%\newcommand{\ps@rmaafront}{\ps@rmaa}% page style for front matter (same as ever, actually)
\newcommand{\ps@rmaafront}{% page style for front matter:  page # in each foot
  \renewcommand\@oddfoot{\hfil{\thepage}}%
  \renewcommand\@evenfoot{{\thepage}\hfil}%
  % shorttitle in both heads
  \renewcommand\@evenhead{\hfil\@shorttitle\hfil}%
  \renewcommand\@oddhead{\@evenhead}
  }
\newcommand{\ps@rmaafrontfirstpage}{% page style for first page of front matter
  % page # in each foot, on outside edge
  \renewcommand\@oddfoot{\hfil\thepage}%
  \renewcommand\@evenfoot{\thepage\hfil}%
  % nothing in the even heads
  \renewcommand\@evenhead{\relax}%
  \renewcommand\@oddhead{\relax}%
  }
%% slight modification to the standard \frontmatter and \mainmatter
%% commands, just to set the \pagestyle
\renewcommand\frontmatter{%
  \latexonecolumn
  \cleardoublepage
  \@frontmattertrue
  \@mainmatterfalse
  \pagenumbering{roman}
  \pagestyle{rmaafront}
  \addtocounter{page}{2} % leave space for the title page
  }
\renewcommand\mainmatter{%
  \pagestyle{empty}
  \if@rmtwocol\twocolumn\fi
  \cleardoublepage
  \@frontmatterfalse
  \@mainmattertrue
  \pagenumbering{arabic}
  \pagestyle{rmaa}
  }
\renewcommand\backmatter{%
  \clearpage\pagestyle{empty}\cleardoublepage
  \@mainmatterfalse
  %% put extra space in the TOC
  \addtocontents{toc}{\bigskip}
  }
%    \end{macrocode}
%
% \subsection{Front sections}\label{sec:front-sections}
% \begin{macro}{\notitle}
% \begin{macro}{\shortname}
% \begin{macro}{frontsection}
% The environment \texttt{frontsection} is for doing a section in the
% front matter (preface, list of participants, or whatever). It has 
% one optional argument:
% \begin{enumerate}
% \item[\#1] Short title (no embedded spaces) for
% use in making PS files. E.g., ``list\_of\_participants''.   
% \end{enumerate}
% and two mandatory arguments: 
% \begin{enumerate}
% \item[\#2] Title of section (e.g., `Preface')
% \item[\#3] Commands to be evaluated before processing the contents
% of the environment. This is mainly useful for putting in the things
% that would go in the preamble of a normal article  (e.g.,
% |\listofauthors|, |\indexauthor|, |\notitle|) 
% \end{enumerate}
%    \begin{macrocode}
\newcommand{\notitle}{\@titlevisiblefalse} \newcommand{\shortname}{}
\newenvironment{frontsection}[3][]{%
  \begingroup
  \renewcommand{\shortname}{#1}
  \ifx\shortname\@rmaaEmpty
     \renewcommand{\articlename}{#2}
  \else
     \renewcommand{\articlename}{\shortname}
  \fi
  \@begofarticlehook
  \@titlevisibletrue % default
  \thispagestyle{rmaafrontfirstpage}
  \listofauthors{} % blank out this in case it is not specified
  \title{{#2}}\shorttitle{#2}\shortauthor{#2}
  #3 % evaluate the command given as second mandatory arg
  \addcontentsline{toc}{chapter}{\protect\@fiddleauthors{\@TOCtitle}{\rmaaTOCauthors}}%
  \addtocontents{htc}{\protect\HyperTOCEntry{\@rawtitle}{\@rawauthors}{\articlename}}
  \if@titlevisible
     \vspace*{1cm}\centerline{\rmaa@FrontSecStyle{\@title}}\vspace*{1cm}
  \fi
  \setlength{\parskip}{3ex}
  }{\frontlabel{RMAALastPage}\clearpage\endgroup}
%    \end{macrocode}
% \end{macro}
% \end{macro}
% \end{macro}
%    \begin{macrocode}
%% Definition of the command \frontlabel, which rights out something like `F4' to
%% the .aux file, instead of `\textlatin  {i}'. This is so it is easy to post-process
%% the file for making PS files of individual sections
\newcommand{\frontlabel}[1]{%
  \@bsphack
  \protected@write\@auxout{}%
     {\string\newlabel{\articlename:#1}{{}{F\@arabic{\c@page}}}}
  \@esphack
  }
%% User command to say X, Y, Z,\\ Editors of the Proceedings
\newcommand{\EditorSignature}{\vspace*{\parskip}\rightline{\@editors}%
  \rightline{Editors of the proceedings}}
%% Fiddle with \thesection, etc. so that they don't include the chapter number
%% Basically go back to the `article' definitions of these
\renewcommand{\thesection}{\@arabic\c@section}
\renewcommand{\thefigure}{\@arabic\c@figure}
\renewcommand{\thetable}{\@arabic\c@table}

\newcommand{\@oldargs}{}
\newcommand{\@newargs}{}
%% add TOC entry for each paper
%\renewcommand{\title}[1]%
%{\renewcommand{\@title}{\rmaa@TitleStyle{#1}}%
%\addtocounter{chapter}{1}\addcontentsline{toc}{chapter}{#1}%
%}

%% Formatting of the table of contents
%% dimensions of various parts
\newlength{\TOCtextwidth}\setlength{\TOCtextwidth}{12.65cm}
\newlength{\TOCpagenumwidth}\setlength{\TOCpagenumwidth}{1.3cm}
\newlength{\TOCentrywidth}\setlength{\TOCentrywidth}{\TOCtextwidth}
\addtolength{\TOCentrywidth}{-\TOCpagenumwidth} 
\newlength{\TOCextramargin}\setlength{\TOCextramargin}{1.25cm}
\newcommand{\TOCbox}[1]{\hspace{\TOCextramargin}\parbox{\TOCtextwidth}{#1}\par}
\newlength{\TOCextramarginx}\setlength{\TOCextramarginx}{\TOCextramargin}\addtolength{\TOCextramarginx}{4.3cm}
\newlength{\TOCtextwidthx}\setlength{\TOCtextwidthx}{\TOCtextwidth}\addtolength{\TOCtextwidthx}{-4.3cm}
\newcommand{\TOCboxx}[1]{\hspace{\TOCextramarginx}\parbox{\TOCtextwidthx}{#1}\par}

\newcommand{\RMAAbannerimage}{rmsc}
\renewcommand\tableofcontents{%
  \renewcommand{\articlename}{CONTENTS}
  \begingroup
  \@begofarticlehook
  \if@twocolumn
     \@restonecoltrue\onecolumn
  \else
     \@restonecolfalse
  \fi
  %% 9pt for journal style
  \if@proceedings\else\small\fi
  \if@newstyle\small\fi %% same for new style
  \if@proceedings\thispagestyle{rmaafrontfirstpage}\else\thispagestyle{empty}\fi
  \shortauthor{Contents}\shorttitle{Contents}
  \vspace*{18mm}
  \if@printbanner
  %% put in the banner header (need to get better image of this)
     \TOCbox{\hspace*{-0.15cm}\includegraphics[width=0.9\TOCtextwidth]{\RMAAbannerimage}}
     \vspace*{8mm}
  %% volume and publication date (should really be in funny font)
     \TOCboxx{\@volumen\\[2mm] \@mes\ \@anyo}
  \else
  %% put in the banner header (need to get better image of this)
     \TOCbox{\hspace*{-0.15cm}\includegraphics[width=0.9\TOCtextwidth]{\RMAAbannerimage}}
     \vspace*{12.5mm}
  %% volume and publication date (should really be in funny font)
     \TOCboxx{~\\[2mm] ~}
  \fi
  \vspace*{16mm}
  %% This is the command that fiddles with the TOC entry so that the
  %% title is flushleft but the authors are flushright
  \renewcommand{\@fiddleauthors}[2]{%
    % first put the paper title
    \renewcommand{\\}{\relax} % temporarily turn off explicit line breaks just in case
    \renewcommand{\altaffilmark}[1]{\relax} % ditto with this
    \renewcommand{\textsuperscript}[1]{\relax} % ditto with this 
    ##1
    % now save the author list in a box
    \sbox\@tempboxa{##2}%
    % check that box is shorter than the total width available for the entry
    \ifdim \wd\@tempboxa <\TOCentrywidth
       % if so, then put it right here, flush against right margin
       \penalty10000\hskip1em plus10em\penalty0\hbox{}\nobreak\hfill\box\@tempboxa
       % Note that we force a space of at least 1em between title and authorlist, 
       % stretchable to 11em since we may have very short titles but still with an 
       % authorlist too long to go on the same line.  The only linebreak allowed is 
       % directly before the authorlist.
    \else
       % if it is too long, we set the original authorlist text in its own parbox
       \par\vskip1ex\parbox[b]{\TOCentrywidth}{\raggedleft ##2}
    \fi
    }
  \@starttoc{toc}%
  \if@restonecol\twocolumn\fi
  \frontlabel{RMAALastPage}
  \clearpage
  \thispagestyle{empty}
  \cleardoublepage % make sure we start on an odd page
  \endgroup

  %%% ALSO, do the HTML TOC
  %% The file that will be used to make the HTML TOC
  \@starttoc{htc}
  %% All the stuff in this file does nothing in the book style
  \addtocontents{htc}{\protect\HyperTOCmeeting{\@conftitle}{\@reunionplace}{\@reuniondate}}
  \addtocontents{htc}{\protect\HyperTOCeditors{\@editors}}
  \addtocontents{htc}{\protect\HyperTOCpreamble}
  \addtocontents{htc}{\protect\HyperTOCEntry{Table of Contents}{}{\articlename}}

  %% and the ADS file
  \ADS@openfile
  }
%% This is now much simpler
\renewcommand\l@chapter[2]{%TOC entry for articles
  \ifnum \c@tocdepth >\m@ne
    \vskip 1.0em \@plus\p@
    \TOCbox{%
      \parbox[b]{\TOCentrywidth}{#1}\parbox[b]{\TOCpagenumwidth}{\hfill#2}
    }
  \fi} 
\renewcommand\l@part[2]{%TOC entry for sessions: centered, no page number
  \ifnum \c@tocdepth >\m@ne
    \vskip 1.5em \@plus1.5em \goodbreak
    \TOCbox{%
      \parbox[b]{\TOCentrywidth}{\hfil#1\hfil}\hfill
    }
  \fi\nobreak} 

%% Formatting theindex
%% strongly encourage page breaks between index entries
\renewcommand{\@idxitem}{\par\penalty-1000\hangindent.45cm} 
\renewcommand{\indexspace}{\par\penalty-1000\vskip 10\p@ \@plus15\p@ \@minus3\p@\relax}
\renewenvironment{theindex}{%
  %% first, make sure we start on an odd page (06 Jun 2000)
  \clearpage\thispagestyle{empty}\cleardoublepage
  %% this looks a bit like an article but is two columns and without abstract, etc
  \renewcommand{\articlename}{INDEX}
  \title{Author Index}\shortauthor{Author Index}\shorttitle{Author Index}
  \listofauthors{}
  \addcontentsline{toc}{chapter}{\protect\@fiddleauthors{\@TOCtitle}{\rmaaTOCauthors}}%
  \addtocontents{htc}{\protect\HyperTOCEntry{\@rawtitle}{\@rawauthors}{\articlename}}
  \pagestyle{rmaa}
  \thispagestyle{rmaafirstpage}
  \@begofarticlehook
  %% 9pt for journal style
  \if@proceedings\else\small\fi
  \if@newstyle\small\fi
  \twocolumn[\@makeminititle]
  %% vertical spacing (see also definition of \@idxitem and \indexspace
  \parskip\z@ \@plus .6\p@\relax % allow quite a lot of leaway here to avoid breaks inside entries
  %% horizontal spacing, breaking and glue parameters (see also the stuff in the definition 
  %% of \indexauthor
  \parindent\z@ 
  \exhyphenpenalty10000 % don't break on explicit hyphens, it always looks bad
  \rightskip=0pt plus.2em % allow a tiny bit of leeway at the ends of lines
  \spaceskip.3333em plus.5em minus.15em % a lot of leeway between words
  \xspaceskip.5em plus.5em minus.4em % this seems to have no effect
  \let\item\@idxitem
  \renewcommand{\\}{\relax} % temporarily turn off explicit line breaks just in case
  \renewcommand{\altaffilmark}[1]{\relax} % ditto with this
  \renewcommand{\textsuperscript}[1]{\relax} % ditto with this
  }{\label{RMAALastPage}\clearpage} % clearpage necessary so that we know about \shortauthor



%% end definitions for book styles
\fi

%    \end{macrocode}
% In order to do away with plates, we sometimes want a figure to
% appear in color only in the online version. The following command
% can bve used in a figure caption to typeset a note to this effect. 
%    \begin{macrocode}
\newcommand\ColorOnlineOnly{%
  \textsc{Note: this figure is available in color in the electronic
    version of this article, obtainable from}
  \texttt{http://\discretionary{}{}{}www.\discretionary{}{}{}astroscu.\discretionary{}{}{}unam.\discretionary{}{}{}mx/\discretionary{}{}{}\~{}rmaa/}.}
%% Plates
%% Define 2 counters: plate is the number in the whole document, plateloop is used 
%% internally by \outputplates
\newcounter{plate}
\newcounter{plateloop}
%% Currently limited to 15 plates in document
%% when we do a label, we need to write to \newlabel entries to the .aux file, one 
%% that says RMAAPlateX so that it can be referred to in \ouputplates, and one that
%% says articlename:RMAAPlateX so that the dvips program can find it and asociate it
%% with the right article
\newcommand{\@platelabeli}{\oldlabel{RMAAPlate1}\label{RMAAPlate1}}
\newcommand{\@platelabelii}{\oldlabel{RMAAPlate2}\label{RMAAPlate2}}
\newcommand{\@platelabeliii}{\oldlabel{RMAAPlate3}\label{RMAAPlate3}}
\newcommand{\@platelabeliv}{\oldlabel{RMAAPlate4}\label{RMAAPlate4}}
\newcommand{\@platelabelv}{\oldlabel{RMAAPlate5}\label{RMAAPlate5}}
\newcommand{\@platelabelvi}{\oldlabel{RMAAPlate6}\label{RMAAPlate6}}
\newcommand{\@platelabelvii}{\oldlabel{RMAAPlate7}\label{RMAAPlate7}}
\newcommand{\@platelabelviii}{\oldlabel{RMAAPlate8}\label{RMAAPlate8}}
\newcommand{\@platelabelix}{\oldlabel{RMAAPlate9}\label{RMAAPlate9}}
\newcommand{\@platelabelx}{\oldlabel{RMAAPlate10}\label{RMAAPlate10}}
\newcommand{\@platelabelxi}{\oldlabel{RMAAPlate11}\label{RMAAPlate11}}
\newcommand{\@platelabelxii}{\oldlabel{RMAAPlate12}\label{RMAAPlate12}}
\newcommand{\@platelabelxiii}{\oldlabel{RMAAPlate13}\label{RMAAPlate13}}
\newcommand{\@platelabelxiv}{\oldlabel{RMAAPlate14}\label{RMAAPlate14}}
\newcommand{\@platelabelxv}{\oldlabel{RMAAPlate15}\label{RMAAPlate15}}
\newcommand{\@outputplatelabeli}{\oldlabel{RMAAOutputPlate1}}
\newcommand{\@outputplatelabelii}{\oldlabel{RMAAOutputPlate2}}
\newcommand{\@outputplatelabeliii}{\oldlabel{RMAAOutputPlate3}}
\newcommand{\@outputplatelabeliv}{\oldlabel{RMAAOutputPlate4}}
\newcommand{\@outputplatelabelv}{\oldlabel{RMAAOutputPlate5}}
\newcommand{\@outputplatelabelvi}{\oldlabel{RMAAOutputPlate6}}
\newcommand{\@outputplatelabelvii}{\oldlabel{RMAAOutputPlate7}}
\newcommand{\@outputplatelabelviii}{\oldlabel{RMAAOutputPlate8}}
\newcommand{\@outputplatelabelix}{\oldlabel{RMAAOutputPlate9}}
\newcommand{\@outputplatelabelx}{\oldlabel{RMAAOutputPlate10}}
\newcommand{\@outputplatelabelxi}{\oldlabel{RMAAOutputPlate11}}
\newcommand{\@outputplatelabelxii}{\oldlabel{RMAAOutputPlate12}}
\newcommand{\@outputplatelabelxiii}{\oldlabel{RMAAOutputPlate13}}
\newcommand{\@outputplatelabelxiv}{\oldlabel{RMAAOutputPlate14}}
\newcommand{\@outputplatelabelxv}{\oldlabel{RMAAOutputPlate15}}
\newcommand{\@platerefi}{\oldpageref{RMAAPlate1}}
\newcommand{\@platerefii}{\oldpageref{RMAAPlate2}}
\newcommand{\@platerefiii}{\oldpageref{RMAAPlate3}}
\newcommand{\@platerefiv}{\oldpageref{RMAAPlate4}}
\newcommand{\@platerefv}{\oldpageref{RMAAPlate5}}
\newcommand{\@platerefvi}{\oldpageref{RMAAPlate6}}
\newcommand{\@platerefvii}{\oldpageref{RMAAPlate7}}
\newcommand{\@platerefviii}{\oldpageref{RMAAPlate8}}
\newcommand{\@platerefix}{\oldpageref{RMAAPlate9}}
\newcommand{\@platerefx}{\oldpageref{RMAAPlate10}}
\newcommand{\@platerefxi}{\oldpageref{RMAAPlate11}}
\newcommand{\@platerefxii}{\oldpageref{RMAAPlate12}}
\newcommand{\@platerefxiii}{\oldpageref{RMAAPlate13}}
\newcommand{\@platerefxiv}{\oldpageref{RMAAPlate14}}
\newcommand{\@platerefxv}{\oldpageref{RMAAPlate15}}
\newcounter{figplatei}
\newcounter{figplateii}
\newcounter{figplateiii}
\newcounter{figplateiv}
\newcounter{figplatev}
\newcounter{figplatevi}
\newcounter{figplatevii}
\newcounter{figplateviii}
\newcounter{figplateix}
\newcounter{figplatex}
\newcounter{figplatexi}
\newcounter{figplatexii}
\newcounter{figplatexiii}
\newcounter{figplatexiv}
\newcounter{figplatexv}
%% Printed output should use the counter `plate'
\renewcommand{\theplate}{\@arabic\c@plate}
\renewcommand{\theplateloop}{\@arabic\c@plateloop}
%% Command to output all the saved plates in this paper
\newcommand{\outputplates}{%
  \setcounter{plateloop}{0}
  \if@book\else\@endofarticlehook\fi % not entirely sure why, but this
                                     % seems to be needed
  \clearpage % first output the last page of text
  % now set pagestyle for blank even pages and make sure we start on an odd page
  \pagestyle{empty}
  \cleardoublepage
  \onecolumn
  \parindent\z@
  % first put in an extra sheet
  % the main purpose of this is so there is always a blank even page before 
  % each plate for filling in the PS file if the article has an odd number
  % of pages (see splitbook.sh)
  \null\vfil\centerline{PLATES}\vfil\clearpage\pagestyle{rmaaplate}\cleardoublepage
  % output each plate in turn
  \@whilenum\value{plateloop}<\value{plate}\do{%
    \addtocounter{plateloop}{1}
    % restore the figure counter from \@figcounti, \@figcountii, etc
    \setcounter{figure}{\value{figplate\roman{plateloop}}}
    % spew out the plate (in its own group for safety)
    {\csname @plate\roman{plateloop}\endcsname\par}
    }
  \pagestyle{rmaa}
  % reset the page number as if nothing had happened
  \addtocounter{page}{-\value{plate}}
  \addtocounter{page}{-\value{plate}}
  }
\newcommand\figcaption{%
  \def\@captype{figure}
  \renewcommand{\@makecaption}{\@makefigcap}
  \caption}
\newcommand\tabcaption{%
  \def\@captype{table}
  \renewcommand{\@makecaption}{\@maketabcap}
  \caption}

\newcommand{\plate}[2][1]{%
  %% User command to make a plate page: 
  %% Mandatory argument is what goes on the plate - this should not 
  %% include any floating environments (e.g. figure). Instead, \figcaption 
  %% should be used to do the captions. 
  %% Optional argument is the number of figures (i.e. captions) on the plate
  %% page (default is 1). This is necessary for incrementing the figure counter. 
  %%
  %% Note that we have to use a \gdef rather than an \xdef when we 
  %% save the plate output in a command. This means that in the book styles when it 
  %% comes to spewing out the plate, we will have forgotten about any user macros 
  %% that may have been defined. This can be got round by repeating the macros inside the plate.
  %% (Presumably, one would use a \newcommand rather than \renewcommand if they were new.)
  % increment global and local counters
  \addtocounter{plate}{1}
  % write appropriate label to .aux file
  \csname @platelabel\roman{plate}\endcsname
  % define a global macro called \@platei, \@plateii, \@plateiii, \@plateiv, ...
  \expandafter\gdef\csname @plate\roman{plate}\endcsname{%
    \csname @outputplatelabel\roman{plateloop}\endcsname
    \null\vfil #2 \vfil
    \cleardoublepage
    }
  % save the current shortauthor and shorttitle
  % need to be \xdef to take current values of \@shortauthor, etc 
  \expandafter\xdef\csname @shorta\roman{plate}\endcsname{\@shortauthor}
  \expandafter\xdef\csname @shortt\roman{plate}\endcsname{\@shorttitle}
  % save the current value of the figure counter for later use in \outputplates
  \setcounter{figplate\roman{plate}}{\value{figure}}
  % increment the figure counter according to the optional argument
  \addtocounter{figure}{#1}
  }
%% page style for plate pages 
\newcommand{\ps@rmaaplate}{%
  % Plate X (see page Y) in the foot
  \renewcommand\@oddfoot{\hfil\csname @shorta\roman{plateloop}\endcsname\ 
    (see page~\csname @plateref\roman{plateloop}\endcsname)\hfil}%
  \renewcommand\@evenfoot{\relax}%
  % nothing in even heads (blank pages)
  \renewcommand\@evenhead{\relax}
  % short title and page num in the odd heads
  \renewcommand\@oddhead{\hfil\csname @shortt\roman{plateloop}\endcsname%
    \hfil\llap{PLATE~\theplateloop}}
  }





%% Things for the User Guide
\if@guide
%% stolen from doc.sty
\def\MakeShortVerb#1{%
  \expandafter\ifx\csname cc\string#1\endcsname\relax
  \@shortvrbinfo{Made }{#1}%
  \add@special{#1}%
  \expandafter
  \xdef\csname cc\string#1\endcsname{\the\catcode`#1}%
  \begingroup
  \catcode`\~\active  \lccode`\~`#1%
  \lowercase{%
    \global\expandafter\let
    \csname ac\string#1\endcsname~%
    \gdef~{\verb~}}%
  \endgroup
  \global\catcode`#1\active
  \else
  \@shortvrbinfo\@empty{#1 already}%
  \fi}
\def\DeleteShortVerb#1{%
  \expandafter\ifx\csname cc\string#1\endcsname\relax
  \else
  \@shortvrbinfo{Deleted }{#1 as}%
  \rem@special{#1}%
  \global\catcode`#1\csname cc\string#1\endcsname
  \global \expandafter\let \csname cc\string#1\endcsname \relax
  \ifnum\catcode`#1=\active
  \begingroup
  \catcode`\~\active   \lccode`\~`#1%
  \lowercase{%
    \global\expandafter\let\expandafter~%
    \csname ac\string#1\endcsname}%
  \endgroup \fi \fi}
\def\@shortvrbinfo#1#2{%
  \PackageInfo{doc}{%
    #1\expandafter\@gobble\string#2 a short reference
    for \string\verb}}
\def\add@special#1{%
  \rem@special{#1}%
  \expandafter\gdef\expandafter\dospecials\expandafter
  {\dospecials \do #1}%
  \expandafter\gdef\expandafter\@sanitize\expandafter
  {\@sanitize \@makeother #1}}
\def\rem@special#1{%
  \def\do##1{%
    \ifnum`#1=`##1 \else \noexpand\do\noexpand##1\fi}%
  \xdef\dospecials{\dospecials}%
  \begingroup
  \def\@makeother##1{%
    \ifnum`#1=`##1 \else \noexpand\@makeother\noexpand##1\fi}%
  \xdef\@sanitize{\@sanitize}%
  \endgroup}
%% define | as shortcut for \verbatim
\MakeShortVerb{\|}
%% comand to specify arguments to LaTeX commands
\newcommand{\rmArg}[1]{%
  \ensuremath{\langle}{\scriptsize \textsf{\MakeUpperCase{#1}}}\ensuremath{\rangle}}
%\newcommand{\rmArg}[1]{\ensuremath{\langle}\textsc{#1}\ensuremath{\rangle}}
\providecommand{\BibTeX}{{\rmfamily B\kern-.05em%
    \textsc{i\kern-.025em b}\kern-.08em%
    T\kern-.1667em\lower.7ex\hbox{E}\kern-.125emX}}
%% Environment for code examples
\newenvironment{Code}{%
  \begin{list}{}{\setlength{\leftmargin}{0.05\columnwidth}} 
  \item[]}{\end{list}}
\fi
%</package>
%    \end{macrocode}

% \section{User commands for astronomical articles}
% \label{sec:user-commands}
% Some random definitions. These go both into the main class file and
% also into |adshtml.sty|.
% 
% Definitions of the journal name in different incarnations.  
%    \begin{macrocode}
%<*package|adshtml-style>
\newcommand{\RMAA}{RMxAA}
\newcommand{\RMAAlong}{Revista Mexicana de Astronom\'{\i}a y Astrof\'{\i}sica}
\newcommand{\RMAASC}{RMxAC}
\newcommand{\RMAASClong}{\RMAA\ (Serie de Conferencias)}
%    \end{macrocode}
% \begin{macro}{\ibidrule}
% This puts a horizontal rule with a period after it. It is used when
% the author list of a reference is identical to the preceding one. 
%    \begin{macrocode}
\newcommand{\ibidrule}{\rule[-0.03cm]{4.7em}{.01cm}.\ }
%    \end{macrocode}
% \end{macro}
% \begin{macro}{\ThisVolume}
% References to other articles in the same volume can be coded as,
% e.g., |\bibitem{} Bloggs, J. \ThisVolume{bloggs}|.
%    \begin{macrocode}
\newcommand{\ThisVolume}[1]{%
  \@year, \if@proceedings\RMAASC\else\RMAA\fi, \@volume,
  \if@book\oldpageref{#1:RMAAFirstPage}\else 000\fi\ (this volume)}
%    \end{macrocode}
% \end{macro}
% Other stuff that may be useful in ref lists.
%    \begin{macrocode}
\newcommand{\etal}{et al.\@}
%    \end{macrocode}
% Various new commands from AAS macros
%    \begin{macrocode}
\def\acknowledgments{\vskip 3ex plus .8ex minus .4ex}
\let\acknowledgements=\acknowledgments                  % second spelling
\let\jnl@style=\rm
\def\ref@jnl#1{{\jnl@style#1}}
\DeclareRobustCommand\aj{\ref@jnl{AJ}}                   
\DeclareRobustCommand\araa{\ref@jnl{ARA\&A}}             
\DeclareRobustCommand\apj{\ref@jnl{ApJ}}                 
\DeclareRobustCommand\apjl{\ref@jnl{ApJ}}                
\DeclareRobustCommand\apjs{\ref@jnl{ApJS}}               
\DeclareRobustCommand\ao{\ref@jnl{Appl.~Opt.}}           
\DeclareRobustCommand\apss{\ref@jnl{Ap\&SS}}             
\DeclareRobustCommand\aap{\ref@jnl{A\&A}}                
\DeclareRobustCommand\aapr{\ref@jnl{A\&A~Rev.}}          
\DeclareRobustCommand\aaps{\ref@jnl{A\&AS}}              
\DeclareRobustCommand\azh{\ref@jnl{AZh}}                 
\DeclareRobustCommand\baas{\ref@jnl{BAAS}}               
\DeclareRobustCommand\jrasc{\ref@jnl{JRASC}}             
\DeclareRobustCommand\memras{\ref@jnl{MmRAS}}            
\DeclareRobustCommand\mnras{\ref@jnl{MNRAS}}             
\DeclareRobustCommand\pra{\ref@jnl{Phys.~Rev.~A}}        
\DeclareRobustCommand\prb{\ref@jnl{Phys.~Rev.~B}}        
\DeclareRobustCommand\prc{\ref@jnl{Phys.~Rev.~C}}        
\DeclareRobustCommand\prd{\ref@jnl{Phys.~Rev.~D}}        
\DeclareRobustCommand\pre{\ref@jnl{Phys.~Rev.~E}}        
\DeclareRobustCommand\prl{\ref@jnl{Phys.~Rev.~Lett.}}    
\DeclareRobustCommand\pasp{\ref@jnl{PASP}}               
\DeclareRobustCommand\pasj{\ref@jnl{PASJ}}               
\DeclareRobustCommand\qjras{\ref@jnl{QJRAS}}             
\DeclareRobustCommand\skytel{\ref@jnl{S\&T}}             
\DeclareRobustCommand\solphys{\ref@jnl{Sol.~Phys.}}      
\DeclareRobustCommand\sovast{\ref@jnl{Soviet~Ast.}}      
\DeclareRobustCommand\ssr{\ref@jnl{Space~Sci.~Rev.}}     
\DeclareRobustCommand\zap{\ref@jnl{ZAp}}                 
\DeclareRobustCommand\nat{\ref@jnl{Nature}}              
\DeclareRobustCommand\iaucirc{\ref@jnl{IAU~Circ.}}
\DeclareRobustCommand\aplett{\ref@jnl{Astrophys.~Lett.}}
\DeclareRobustCommand\apspr{\ref@jnl{Astrophys.~Space~Phys.~Res.}}
\DeclareRobustCommand\bain{\ref@jnl{Bull.~Astron.~Inst.~Netherlands}}
\DeclareRobustCommand\fcp{\ref@jnl{Fund.~Cosmic~Phys.}}
\DeclareRobustCommand\gca{\ref@jnl{Geochim.~Cosmochim.~Acta}}
\DeclareRobustCommand\grl{\ref@jnl{Geophys.~Res.~Lett.}}
\DeclareRobustCommand\jcp{\ref@jnl{J.~Chem.~Phys.}}      
\DeclareRobustCommand\jgr{\ref@jnl{J.~Geophys.~Res.}}    
\DeclareRobustCommand\jqsrt{\ref@jnl{J.~Quant.~Spec.~Radiat.~Transf.}}
\DeclareRobustCommand\memsai{\ref@jnl{Mem.~Soc.~Astron.~Italiana}}
\DeclareRobustCommand\nphysa{\ref@jnl{Nucl.~Phys.~A}}
\DeclareRobustCommand\physrep{\ref@jnl{Phys.~Rep.}}
\DeclareRobustCommand\physscr{\ref@jnl{Phys.~Scr}}
\DeclareRobustCommand\planss{\ref@jnl{Planet.~Space~Sci.}}
\DeclareRobustCommand\procspie{\ref@jnl{Proc.~SPIE}}
\let\astap=\aap
\let\apjlett=\apjl
\let\apjsupp=\apjs
\let\applopt=\ao
\def\phn{\phantom{0}}
\def\phd{\phantom{.}}
\def\phs{\phantom{\ensuremath{-}}}
\def\phm#1{\phantom{#1}}
%    \end{macrocode}
% All these are redone as robust and using more modern commands where
% possible. The robustness is so they will not get expanded if used in
% abstracts and titles when these are passed to the ADS file. 
% \changes{3.16}{2003/04/07}{
%  Added \cs{Color} command for doing photometric colors.}
%    \begin{macrocode}
%%\def\sun{\hbox{$\odot$}}
\DeclareRobustCommand\sun{\ensuremath{\odot}}
%%\def\earth{\hbox{$\oplus$}}
\DeclareRobustCommand\earth{\ensuremath{\oplus}}
%%\def\lesssim{\mathrel{\hbox{\rlap{\hbox{\lower4pt\hbox{$\sim$}}}\hbox{$<$}}}}
%%\def\gtrsim{\mathrel{\hbox{\rlap{\hbox{\lower4pt\hbox{$\sim$}}}\hbox{$>$}}}}
\DeclareRobustCommand\lesssim{\mathrel{\hbox{\smaller\smaller\rlap{\hbox{\lower3pt\hbox{$\sim$}}}\raise1.5pt\hbox{$<$}}}}
\DeclareRobustCommand\gtrsim{\mathrel{\hbox{\smaller\smaller\rlap{\hbox{\lower3pt\hbox{$\sim$}}}\raise1.5pt\hbox{$>$}}}}
%%\def\lta{\mathrel{\spose{\lower 3pt\hbox{$\mathchar"218$}}
%%     \raise 2.0pt\hbox{$\mathchar"13C$}}}
%%\def\gta{\mathrel{\spose{\lower 3pt\hbox{$\mathchar"218$}}
%%     \raise 2.0pt\hbox{$\mathchar"13E$}}}
\DeclareRobustCommand\sq{\hbox{\rlap{$\sqcap$}$\sqcup$}}
%%\def\arcdeg{\hbox{$^\circ$}}
%%\def\arcmin{\hbox{$^\prime$}}
%%\def\arcsec{\hbox{$^{\prime\prime}$}}
\DeclareRobustCommand\arcdeg{\ensuremath{^\circ}}
\DeclareRobustCommand\arcmin{\ensuremath{^\prime}}
\DeclareRobustCommand\arcsec{\ensuremath{^{\prime\prime}}}
%%\def\fd{\hbox{$.\!\!^{\rm d}$}}
%%\def\fh{\hbox{$.\!\!^{\rm h}$}}
%%\def\fm{\hbox{$.\!\!^{\rm m}$}}
%%\def\fs{\hbox{$.\!\!^{\rm s}$}}
%%\def\fdg{\hbox{$.\!\!^\circ$}}
%%\def\farcm{\hbox{$.\mkern-4mu^\prime$}}
%%\def\farcs{\hbox{$.\!\!^{\prime\prime}$}}
%%\def\fp{\hbox{$.\!\!^{\scriptscriptstyle\rm p}$}}
\DeclareRobustCommand\fd{\ensuremath{.\!\!^\mathrm{d}}}
\DeclareRobustCommand\fh{\ensuremath{.\!\!^\mathrm{h}}}
\DeclareRobustCommand\fm{\ensuremath{.\!\!^\mathrm{m}}}
\DeclareRobustCommand\fs{\ensuremath{\mkern0.5mu.\mkern-5mu^\mathrm{s}\mkern-1mu}}
\DeclareRobustCommand\fdg{\ensuremath{.\!\!^\circ}}
\DeclareRobustCommand\farcm{\ensuremath{.\mkern-4mu^\prime}}
\DeclareRobustCommand\farcs{\ensuremath{\mkern1mu.\mkern-6mu^{\prime\prime}\mkern-1.5mu}}
%%\DeclareRobustCommand\farcs{\ensuremath{.\!\!^{\prime\prime}}}
\DeclareRobustCommand\fp{\ensuremath{.\!\!^{\scriptscriptstyle\rm p}}}
%%\def\micron{\hbox{$\mu$m}}
\DeclareRobustCommand\micron{\ensuremath{\mu\mathrm{m}}}
\let\la=\lesssim                        
\let\ga=\gtrsim
%%\def\case#1#2{\hbox{$\frac{#1}{#2}$}}
%%\def\slantfrac#1#2{\hbox{$\,^#1\!/_#2$}}
%%\def\onehalf{\slantfrac{1}{2}}
%%\def\onethird{\slantfrac{1}{3}}
%%\def\twothirds{\slantfrac{2}{3}}
%%\def\onequarter{\slantfrac{1}{4}}
%%\def\threequarters{\slantfrac{3}{4}}
\DeclareRobustCommand\case[2]{\ensuremath{\frac{#1}{#2}}}
\DeclareRobustCommand\slantfrac[2]{\ensuremath{\,^#1\!/_#2}}
\DeclareRobustCommand\onehalf{\slantfrac{1}{2}}
\DeclareRobustCommand\onethird{\slantfrac{1}{3}}
\DeclareRobustCommand\twothirds{\slantfrac{2}{3}}
\DeclareRobustCommand\onequarter{\slantfrac{1}{4}}
\DeclareRobustCommand\threequarters{\slantfrac{3}{4}}
%%\def\ubvr{\hbox{$U\!BV\!R$}}            
%%\def\ub{\hbox{$U\!-\!B$}}               
%%\def\bv{\hbox{$B\!-\!V$}}               
%%\def\vr{\hbox{$V\!-\!R$}}               
%%\def\ur{\hbox{$U\!-\!R$}}
%    \end{macrocode}
% New commands for ``photometric'' colors 
%    \begin{macrocode}
\DeclareRobustCommand\ubvr{\ensuremath{U\!BV\!R}}            
\DeclareRobustCommand\Color[2]{%
  \ensuremath{%
    #1\!\ifx#1V\!\fi-\!#2
  }}
\DeclareRobustCommand\ub{\Color{U}{B}}
\DeclareRobustCommand\bv{\Color{B}{V}}
\DeclareRobustCommand\vr{\Color{V}{R}}
\DeclareRobustCommand\ur{\Color{U}{R}}
%% New version of the \ion command, that actually works! 
%    \end{macrocode}
% \begin{macro}{\ion}
% This version of the |\ion| command allows you to write either
% |\ion{N}{4}| or |\ion{N}{iv}| in order to get N~IV. 
%    \begin{macrocode}
\DeclareRobustCommand\ion[2]{%
  \mbox{#1\kern0.2em%
    \smaller\rmfamily%
    \edef\@tempa{\@car#2\@nil}%
    \ifcat1\@tempa%
    \@Roman{#2}%
    \else%
    \uppercase{#2}%
    \fi}}
\newcounter{IonStage}
\DeclareRobustCommand\plainion[2]{\setcounter{IonStage}{#2}#1 \Roman{IonStage}}
%%\def\ion#1#2{\ensuremath{\mathrm{#1}\;}{\protect\small\rm\@Roman{#2}}}
%    \end{macrocode}
% \end{macro}
%    \begin{macrocode}
\DeclareRobustCommand\nodata{$\cdots$}
%    \end{macrocode}
% These are needed by the |\ion| command. Isn't there an easier way of
% doing this?
%    \begin{macrocode}
\let\@small=\small
\let\@normalsize=\normalsize
\let\@footnotesize=\footnotesize
\let\@scriptsize=\scriptsize
\let\@tiny=\tiny
\let\@large=\large
\let\@Large=\Large
\let\@LARGE=\LARGE
\let\@huge=\huge
\let\@Huge=\Huge
\newcommand{\smaller}{\small}
\renewcommand{\normalsize}{\renewcommand{\smaller}{\small}\@normalsize}
\renewcommand{\small}{\renewcommand{\smaller}{\footnotesize}\@small}
\renewcommand{\footnotesize}{\renewcommand{\smaller}{\scriptsize}\@footnotesize}
\renewcommand{\scriptsize}{\renewcommand{\smaller}{\tiny}\@scriptsize}
\renewcommand{\tiny}{\renewcommand{\smaller}{\tiny}\@tiny}
\renewcommand{\large}{\renewcommand{\smaller}{\normalsize}\@large}  
\renewcommand{\Large}{\renewcommand{\smaller}{\large}\@Large}  
\renewcommand{\LARGE}{\renewcommand{\smaller}{\Large}\@LARGE}  
\renewcommand{\huge}{\renewcommand{\smaller}{\LARGE}\@huge}  
\renewcommand{\Huge}{\renewcommand{\smaller}{\huge}\@Huge}  
%</package|adshtml-style>
%    \end{macrocode}


% \section{Other files}

% \subsection{The driver and make files}
% \label{sec:driver_make}
% These are put here for completeness, although they still need to be
% distributed separately since otherwise the user won't be able to
% unpack this file into its constituent parts.
% 
% The |Makefile| has two principal rules: |docs|, and |macros|. The
% |docs| rule produces a PDF and gzipped PS file of the documentation.
% I still haven't perfected the combination of rules required to get
% it to run \LaTeX{} enough times in all circumstances. The |macros|
% rule runs \LaTeX{} on |rmaa.ins| in order to unpack |rmaa.cls| and
% all the other little files (including the |Makefile| and the
% |rmaa.ins| files themselves\dots 
% 
% First set up lists of various types of file. 
%    \begin{macrocode}
%<*makefile>
DOCFILES=rmaa.pdf
MACRO_LTX_FILES=rmaa.cls rmaacite.sty adshtml.sty
SHELLSCRIPTS=makeads.sh
MACRO_OTHER_FILES=adshtml.perl sample.latex2html-init
MAKEFILE=Makefile
MACROFILES=$(MACRO_LTX_FILES) $(SHELLSCRIPTS) $(MACRO_OTHER_FILES)         
SKELETONFILES=ADS.tex 
%    \end{macrocode}
% Install the files read only into the package directory, specified via a path relative to the source directory. \changes{4.0}{2019/11/14}{Reorganize folder layout so that hardly anything is in the top level.  That way, things are clearer on github.}
%    \begin{macrocode}
INSTALL=install 
INSTALL_DIR=../packages
%    \end{macrocode}
% The default rule is to make everything. Installation must be
% explicitly asked for.
%    \begin{macrocode}
all: docs macros

install: install-docs install-macros

%    \end{macrocode}
% Most files are installed read-only to prevent editing. Skeleton
% files must be writeable. Must be careful not to overwrite the
% Makefile in the parent directory. 
%    \begin{macrocode}
install-macros: macros
        $(INSTALL) -m 444 $(MACRO_LTX_FILES) $(MACRO_OTHER_FILES) $(INSTALL_DIR)
        $(INSTALL) -m 555 $(SHELLSCRIPTS) $(INSTALL_DIR)
        $(INSTALL) -m 644 $(SKELETONFILES) $(INSTALL_DIR)

install-docs: docs
        $(INSTALL) rmaa.pdf $(INSTALL_DIR)/rm-fulldocs.pdf

%    \end{macrocode}
% Make the PDF and gzipped PS files. 
%    \begin{macrocode}
docs: $(DOCFILES)

%    \end{macrocode}
% Make everything else if either the |.dtx| or |.ins| files have
% changed. 
%    \begin{macrocode}
macros: rmaa.dtx rmaa.ins
%    \end{macrocode}
% Ensure that all the macro files can be written to then recreate
% them. The |touch| is there in case any of the files don't exist
% yet. 
%    \begin{macrocode}
        touch $(MACROFILES) $(MAKEFILE)
        chmod u+w $(MACROFILES) $(MAKEFILE) rmaa.ins
        latex "\input{rmaa.ins}"
%    \end{macrocode}
% \TeX{} automatically adds the |.tex| extension to extensionless
% files. We fix this and take the opportunity to convert
% back to TABs all that it has lovingly converted to spaces. 
%    \begin{macrocode}
        cp -f $(MAKEFILE) $(MAKEFILE).BCK && unexpand $(MAKEFILE).tex >| $(MAKEFILE) && rm $(MAKEFILE).tex
%    \end{macrocode}
% Make the shell scripts executable and everything unwriteable (so
% they don't get hand-edited by mistake). However, undo this for the
% files that the user is \emph{supposed} to modify. 
%    \begin{macrocode}
        chmod a+x $(SHELLSCRIPTS)
        chmod a-w $(MACROFILES) $(MAKEFILE)
        chmod u+w $(SKELETONFILES)

%    \end{macrocode}
% We now use PDF\LaTeX{} for this.
%    \begin{macrocode}
rmaa.idx: rmaa.dtx
        @echo "REGENERATING .idx FILE"
        pdflatex "\nonstopmode\input{rmaa.dtx}"

rmaa.aux: rmaa.dtx rmaa.ind rmaa.gls
        @echo "REGENERATING .aux FILE"
        pdflatex "\nonstopmode\input{rmaa.dtx}"

rmaa.ind: rmaa.idx
        @echo "REGENERATING INDEX"
        makeindex -s gind.ist rmaa

rmaa.gls: rmaa.glo
        @echo "REGENERATING CHANGELOG"
        makeindex -s gglo.ist -o rmaa.gls rmaa.glo

rmaa.pdf: rmaa.aux rmaa.ind rmaa.gls
        @echo "REGENERATING .pdf FILE"
        pdflatex "\nonstopmode\input{rmaa.dtx}"

clean: 
        rm -f rmaa.aux rmaa.dvi rmaa.ind rmaa.gls
%</makefile>
%    \end{macrocode}

% The |rmaa.ins| file uses the \textsf{docstrip} package to split out
% all the files and strip off this documentation.
%    \begin{macrocode}
%<*ins>
%%
%% Installation batch file for the RMAA LaTeX package
%%
\def\batchfile{rmaa.ins}
\input docstrip.tex
\preamble

PLEASE DO NOT EDIT THIS FILE! It is generated automatically from
the file rmaa.dtx using the accompanying Makefile and rmaa.ins 
Any hand-edited changes will be lost

A LaTeX document class for the RMxAA

Author: William Henney (w.henney@irya.unam.mx)

Latest version available from: http://www.irya.unam.mx/rmaa

Copyright (C) 1999-2019 Instituto de Astronomia, UNAM, Mexico

Permission is hereby granted, free of charge, to any person obtaining
a copy of this software and associated documentation files (the
"Software"), to deal in the Software without restriction, including
without limitation the rights to use, copy, modify, merge, publish,
distribute, sublicense, and/or sell copies of the Software, and to
permit persons to whom the Software is furnished to do so, subject to
the following conditions:

The above copyright notice and this permission notice shall be
included in all copies or substantial portions of the Software.

THE SOFTWARE IS PROVIDED "AS IS", WITHOUT WARRANTY OF ANY KIND,
EXPRESS OR IMPLIED, INCLUDING BUT NOT LIMITED TO THE WARRANTIES OF
MERCHANTABILITY, FITNESS FOR A PARTICULAR PURPOSE AND
NONINFRINGEMENT. IN NO EVENT SHALL THE AUTHORS OR COPYRIGHT HOLDERS BE
LIABLE FOR ANY CLAIM, DAMAGES OR OTHER LIABILITY, WHETHER IN AN ACTION
OF CONTRACT, TORT OR OTHERWISE, ARISING FROM, OUT OF OR IN CONNECTION
WITH THE SOFTWARE OR THE USE OR OTHER DEALINGS IN THE SOFTWARE.
\endpreamble
%    \end{macrocode}
% Clobber files without asking
%    \begin{macrocode}
\askforoverwritefalse
%    \end{macrocode}
% Make the \LaTeX{} macro files (class and styles)
%    \begin{macrocode}
\generate{% 
  \file{rmaa.cls}{\from{rmaa.dtx}{package}}
  \file{rmaa.ins}{\from{rmaa.dtx}{ins}}
  \file{rmaacite.sty}{\from{rmaa.dtx}{rmaacite}}
  \file{adshtml.sty}{\from{rmaa.dtx}{adshtml-style}}
  \file{ADS.tex}{\from{rmaa.dtx}{ads-container}}
  }
%    \end{macrocode}
% Make the other non-\LaTeX{} files. These are done separately since
% they can't use `\%\%' as a comment character for the
% preamble/postamble (and Perl files can't have the postamble I don't
% think). 
%    \begin{macrocode}
\generate{\nopreamble\nopostamble% 
  \file{Makefile}{\from{rmaa.dtx}{makefile}}
  \file{adshtml.perl}{\from{rmaa.dtx}{adshtml-perl}}
  \file{makeads.sh}{\from{rmaa.dtx}{ads-shell-script}}
  \file{sample.latex2html-init}{\from{rmaa.dtx}{latex2html-init}}
  }

%% Local Variables: 
%% mode: LaTeX
%% mode: Font-lock
%% End:
%</ins>
%    \end{macrocode}

% \subsection{Extended citation syntax}
% Modified from a file called mnrascite.sty
% I no longer remember where I stole the original version of this from
% 
% Alterations to the cite commands to make them do more sensible things.
% N.B. the bibliography style should produce entries which have labels where
% the date is delimited by angle brackets '<>' for these macros to work
% properly. They will not give very sensible error messages if you do not do 
% this properly.
% 
% A new command |\scite{key}| is defined which produces a citation which is 
% suitable for insertion in a sentence. i.e., it does not put brackets around 
% the whole citation, only the date. Also |\pcite{key}| produce output with no
% parentheses whatsoever. 
% 
% Should redefine |\@bibitem| really to have any hope of compatibility
% the command |\lbibcite| is defined to deal with the citations.
%    \begin{macrocode}
%<*rmaacite>
\def\bibitem[#1<#2>]#3{\par\hangindent=1.5em\hangafter=1
  \if@filesw
  {\def\protect##1{\string ##1\space}
    \immediate\write\@auxout{\string\lbibcite{#3}{#1}{#2}}}
  \fi}
\@redefinebibitem % what it says
%    \end{macrocode}
% Note that the |\lbibcite| defines two types of citation one contains the names,
% the other the year
%    \begin{macrocode}
\def\lbibcite#1#2#3{\global\@namedef{b@#1}{#2}
                    \global\@namedef{by@#1}{#3}}

\def\@bcite#1#2{({#1\if@tempswa , #2\fi})}
\def\@pcite#1#2{#1\if@tempswa , #2\fi}
%    \end{macrocode}
% This declaration contains what is put between the name and the date
% \changes{3.11}{2003/01/31}{Moved the \cs{articlename} to after the key for consistency
% with the new natbib stuff.}
%    \begin{macrocode}
\def\@brcite{\ (}
\def\@ebrcite{)}
\def\@fmttop#1{\@fmtcite[\csname #1\endcsname]}
\def\@citex[#1]#2{%
%%  \if@filesw\immediate\write\@auxout{\string\citation{#2}}\fi
%% WJH added 16 Jan 2002
  \@for\@citeb:=#2\do{%
    \edef\@citeb{\expandafter\@firstofone\@citeb:\articlename}%
    \if@filesw\immediate\write\@auxout{\string\citation{\@citeb}}\fi}%
%% to here
  \def\@citea{}%
  \@cite{\@for\@citeb:=#2\do
    {%
      \edef\@citeb{\expandafter\@firstofone\@citeb:\articlename}%
      \@citea\def\@citea{;\ }\@ifundefined
      {b@\@citeb}{{\bf ?}\@warning {Citation `\@citeb' on page \thepage \space undefined}}%
      {\csname b@\@citeb\endcsname\@citebro\csname by@\@citeb\endcsname\@citebrc}}}%
  {#1}}

\def\cite{\@ifnextchar [{\let\@citebro=\space
                          \let\@citebrc=\relax
                          \let\@cite=\@bcite\@tempswatrue\@citex}
                        {\let\@citebro=\space
                          \let\@citebrc=\relax
                          \let\@cite=\@bcite\@tempswafalse\@citex[]}}
\def\pcite{\@ifnextchar [{\let\@citebro=\space
                          \let\@citebrc=\relax
                          \let\@cite=\@pcite\@tempswatrue\@citex}
                        {\let\@citebro=\space
                          \let\@citebrc=\relax
                          \let\@cite=\@pcite\@tempswafalse\@citex[]}}
\def\scite{\@ifnextchar [{\let\@citebro=\@brcite
                          \let\@citebrc=\@ebrcite
                          \let\@cite=\@pcite\@tempswatrue\@citex}
                        {\let\@citebro=\@brcite
                          \let\@citebrc=\@ebrcite
                          \let\@cite=\@pcite\@tempswafalse\@citex[]}}

\renewenvironment{thebibliography}
  {\section*{REFERENCES}\vspace*{\afterbiblioskip}\small
    \if@book\RereadBiblio\fi
    \sloppy\clubpenalty4000\widowpenalty4000
    \sfcode`\.=1000\relax
    \setlength\parindent{\z@}\setlength\parskip{\z@}}
  {\relax\normalsize\@outputaddressesifnecessary} 
%</rmaacite>
%    \end{macrocode}

% \subsection{Auxilliary files for generating ADS data}

% The ADS data are written by \LaTeX{} to a |.ads| file. This file
% then has to be included in the following file (|ADS.tex|),
% which is run through \LaTeX{}2HTML (see below). 

%    \begin{macrocode}
%<*ads-container>
%%
%% Instructions:
%% Edit the \input{} line below to reflect the name of your file. 
\documentclass{article}
\usepackage[latin1]{inputenc}
\usepackage{adshtml}
\newcommand{\ADSfield}[2]
{\%#1  #2\\}
%%\newcommand{\ADSvfield}[2]
%%{\%#1 \begin{verbatim}#2\end{verbatim}}
\newenvironment{ADSarticle}{}{\par\bigskip}
\renewcommand{\bibitem}[2][]{\par}
\setlength{\parindent}{0pt}
\renewcommand{\ThisVolume}[1]{%  
  2001, \RMAASC, 12, \oldpageref{#1:RMAAFirstPage}}  
 

\begin{document}
\makeatletter
\input{rmbooksamp.ads}
\end{document}
%</ads-container>
%    \end{macrocode}

% This is the |adshtml.sty| style file used by |ADS.tex|. It just
% defines a few 
% macros that might be used in abstracts, etc (really it should define
% them all). This now gets all its code from
% Sec~\ref{sec:user-commands}. 
%    \begin{macrocode}
%<*adshtml-style>
%% Nothing else needs to go here now. 
%</adshtml-style>
%    \end{macrocode}
% The init file for latex2html.
%    \begin{macrocode}
%<*latex2html-init>
$PDFTOCAIRO = "pdftocairo";
$PDFLATEX = "pdflatex";
$LATEX = "latex";
$MAX_SPLIT_DEPTH = 0;
$HTML_VERSION = "4.0,math";
$NO_SIMPLE_MATH = 1;
$ADDRESS = "<p><address><a href=\"mailto:rmaa\@astro.unam.mx\">rmaa\@astro.unam.mx</address></p>"; 
$INFO = "<p>This document was created using latex2html and the rmaa latex macros</p>"; 
$NO_NAVIGATION = 1;
1; 
%</latex2html-init>
%    \end{macrocode}
% This is the |adshtml.perl| file that is used by \LaTeX{}2HTML. Some
% bizarre \TeX{}y command sometimes show up so we tell it to ignore
% these. All the astronomy macros that come from AAS are processed in
% \TeX{}. That way they get made to images but with an |ALT| tag that
% has the original code, which is what will be picked up by |lynx|,
% with the result that they get passed intact into the final |.txt|
% file. 
%    \begin{macrocode}
%<*adshtml-perl>
package main;

$dum = &do_cmd_makeatletter();

sub do_cmd_RMAA {
    local($_) = @_;
    "RMxAA" . $_;
}

sub do_cmd_RMAASC {
    local($_) = @_;
    "RMxAASC" . $_;
}

sub do_cmd_etal {
    local($_) = @_;
    "et al\." . $_;
}

sub do_cmd_ibidrule {
    local($_) = @_;
    "<img src=\"foo\" alt=\"\ibidrule\">" . $_;
}


sub do_cmd_bibitem {
    local($_) = @_;
    local($key);
    local($opt, $dummy)=&get_next_optional_argument; 
    $key = &missing_braces unless
        ((s/$next_pair_pr_rx/$key = $2; ''/eo)
        ||(s/$next_pair_rx/$key = $2; ''/eo));
    "<p>" . $_; 
}

&ignore_commands (<<_RAW_ARG_CMDS_);
lowercase
adjustfinalcols
textsuperscript # {}
newline # []
thinmuskip
mskip
smallskip
medskip
bigskip
hspace # {}
hspacestar # {}
vspace # {}
vspacestar # {}
mathrm
_at_mathrm
spacefactor
discretionary
unhbox 
voidb_at_x 
penalty 
_at_M
_at_m
null
IeC
_RAW_ARG_CMDS_

&process_commands_in_tex (<<_RAW_ARG_CMDS_);
sun
earth
lesssim
gtrsim
sq
argdeg
arcmin
arcsec
fd
fh
fm
fs
fdg
farcm
farcs
fp
micron
la
ga
case # {} # {}
slantfrac # {} # {}
onehalf
onethird
twothirds
onequarter
threequarters
ubvr
ub
bv
vr
ur
aj
araa
apj
apjl
apjs
ao
apss
aap
aapr
aaps
azh
baas
jrasc
memras
mnras
pra
prb
prc
prd
pre
prl
pasp
pasj
qjras
skytel
solphys
sovast
ssr
zap
nat
iaucirc
aplett
apspr
bain
fcp
gca
grl
jcp
jgr
jqsrt
memsai
nphysa
physrep
physscr
planss
procspie
_RAW_ARG_CMDS_

1;                              # This must be the last line
%</adshtml-perl> 
%    \end{macrocode}

% This is a shell script |makeads.sh| that generates the |_ads.txt| file from
% |_ads.tex| by using \LaTeX{}2HTML and |lynx|.  \changes{4.0}{2019/11/13}{Make sure that lynx writes utf-8 output instead of latin-1.}
% \changes{2.2beta13}{2002/03/12}{Complete rewrite of |makeads.sh|. Now much 
% improved I think}
% \changes{4.0}{2019/11/13}{Fix the at-sign substitution so that it only eliminates \cs{@} and does not mess with email addresses.}

%    \begin{macrocode}
%<*ads-shell-script>
#! /bin/sh
if [ -z $1 ]; then
    echo "Usage: $0 basename"
    exit
fi
BASENAME=$1
ADSBASENAME=${BASENAME}_ads
AUXFILE=$BASENAME.aux
TMPFILE=tmp.$$


echo "Generating $ADSBASENAME.tex"

# For some reason, must double-up \\ to \\\\ but not \ to \\ - odd!
cat > $ADSBASENAME.tex <<EOF
%% PLEASE DO NOT EDIT THIS FILE!  It is generated automatically from
%% the script makeads.sh.  Any hand-edited changes will be lost
%% 
\documentclass{article}
\usepackage[latin1]{inputenc}
\usepackage{html}
\usepackage{adshtml}
\newcommand{\ADSfield}[2]{\%#1  #2\\\\}
\newenvironment{ADSarticle}{}{\par\bigskip}
\setlength{\parindent}{0pt}
\newcounter{istage}
\renewcommand{\ion}[2]{\setcounter{istage}{#2}#1 \Roman{istage}}
\begin{document}
\makeatletter
EOF
echo >> $ADSBASENAME.tex "\input{$BASENAME.ads2}" 
echo >> $ADSBASENAME.tex "\end{document}"



echo "Generating $ADSBASENAME.sed" 

# Find type of journal
grep '\\documentclass\[.*proceedings.*\]' $BASENAME.tex >/dev/null \
     && JOURNAL='RMxAC' \
     || JOURNAL='RMxAA'

# Find the year and volume number
YEAR=`grep '\\SetYear' $BASENAME.tex | sed 's/.*\\SetYear{\(.*\)}.*/\1/'`
VOLUME=`grep '\\SetVolume' $BASENAME.tex | sed 's/.*\\SetVolume{\(.*\)}.*/\1/'`

# Generate an auxilliary sed script 
cat > $TMPFILE <<EOF
s|\\\\newlabel{\(.*\):RMAAFirstPage}{{.*}{\(.*\)}}|s/\\\\\\\\ThisVolume{\1}/$YEAR, RMxAC, $VOLUME, \2/g|g
EOF

# Now go through the .aux file, looking for \newlabels and use the sed 
# script just generated in order to make another sed script, one that 
# will replace \ThisVolume{LABEL} with the correct reference
grep '^\\newlabel{.*:RMAAFirstPage}' $AUXFILE \
     | sed -f $TMPFILE > $ADSBASENAME.sed && rm $TMPFILE

# Apend a further sed command to get rid of \@, since redefining 
# it in _ads.tex didn't work
echo >> $ADSBASENAME.sed 's/\\\\@//g'

echo "Generating $BASENAME.ads2 from $ADSBASENAME.sed and $BASENAME.ads" 

# Apply this to the .ads file, generating the .ads2 file, which is included
# by the _ads.tex file.
sed -f $ADSBASENAME.sed $BASENAME.ads > $BASENAME.ads2

echo "Generating $ADSBASENAME.txt from $ADSBASENAME.tex and $BASENAME.ads2"

# Now run latex2html on the _ads.tex file
latex2html -no_navigation -info 0 -address 0 -html_version 4.0,math $ADSBASENAME

# Then, convert to utf8 with lynx and tidy it up a bit
lynx --justify=off --width=100 --dump --nolist --display_charset=utf-8 \
     $ADSBASENAME/$ADSBASENAME.html | cut -c4- \
     | sed 's/ibidrule/\\ibidrule/g' > $ADSBASENAME.txt
%</ads-shell-script>
%    \end{macrocode}

% \section{Epilogue}
% Err, that's it \dots

\endinput


% \endinput
% Local Variables:
% mode: doctex
% TeX-master: t
% End:
"

rmaa.aux: rmaa.dtx rmaa.ind rmaa.gls
        @echo "REGENERATING .aux FILE"
        pdflatex "\nonstopmode% \def\fileversion{4.0} \def\filedate{2019/11/12}
% \MakeShortVerb{\|}
% \title{A \LaTeX{} document class for the Revista Mexicana de
% Astronom\'\i a y Astrof\'\i sica\thanks{This file has version
% \fileversion{} last revised \filedate{}.}}
% \author{William J. Henney\thanks{w.henney@irya.unam.mx}}
% \date{\filedate}
% \maketitle
% \paragraph{Abstract}
% This document describes the \texttt{rmaa} \LaTeX{}
% document class for the journal ``Revista Mexicana de
% Astronom\'\i a y Astrof\'\i sica''. The macros can be used for 
% \begin{enumerate}
% \item Preparing author manuscripts for the journal (including
% conference proceedings to be published in the ``Serie de
% Conferencias'')
% \item Typesetting an entire volume of the journal by inputting the
% individual article files from a container file. Facilities are
% included for automatic generation of TOC and Author Index. 
% \end{enumerate}
% 
% \tableofcontents
% \section{The documentation driver file}
% This is the part that generates the documentation that you may or
% may not be reading, which is done with the help of the Makefile
% discussed in section~\ref{sec:driver_make}. 
%    \begin{macrocode}
%<*driver>
\documentclass{ltxdoc}
\usepackage{hyperref}
\usepackage{palatino}
\EnableCrossrefs \CodelineIndex \RecordChanges
\setlength{\parindent}{0pt}
\def\MacroFont{\footnotesize\tt}
\begin{document} \DocInput{rmaa.dtx} \PrintIndex \PrintChanges \end{document}
%</driver>
%    \end{macrocode}

% \section{The implementation}
% 
% \changes{3.24}{2007/09/14}{Added user access to the macro version
% and date via \cs{rmVersion} and \cs{rmDate}}
% 
%    \begin{macrocode}
%<*package>
\newcommand{\rmaa@version}{3.26}
\newcommand{\rmaa@date}{2011/08/15}
\newcommand\rmVersion{\rmaa@version}
\newcommand\rmDate{\rmaa@date}
\NeedsTeXFormat{LaTeX2e}[1995/06/01]
%    \end{macrocode}
% \changes{4.0}{2019/11/02}{Change abstract to appear before
% resumen. Change RevMexAA to RMxAA for consistency with ADS.}
%
% \changes{3.22}{2007/03/22}{Allow proceedings articles to have
%   received/accepted dates} \changes{3.21}{2006/05/15}{Changes for
%   Silvia - abstracts are now shipped out to ADS.}
% \changes{3.19}{2006/01/19}{Changes for Silvia - remove adresses from
%   end, new abstract style} \changes{2.2beta}{2001/10/30}{ Fixed bug
%   whereby twocolumn option was not passed through to book
%   class. Finally sorted out the ``label on last page'' issue. }
% \changes{2.1}{2001/01/23}{ Fixed some inconsistencies in
%   \cs{nocite} and \cs{bibitem} uniqification Fixed the definition of
%   \cs{ion} so it will work in math mode and in contexts where the
%   typesize is different from normal (e.g., title, fig captions).
%   Modified the definition of \cs{@caption}, so we don't write any
%   list-of-figures data (we never use it anyway, and it just causes
%   problems with fragile commands in fig captions). }
% \changes{2.0}{2000/11/30}{ Added `newstyle' option, which should be
%   used in all new proceedings volumes (wider abstract, smaller
%   typeface in fig captions, etc.) Added commands for formatting
%   table captions. New demonstation document rmpostsamp.tex for
%   1-page posters. This will be the last release under a PRI
%   government (for a while...)}  \changes{1.9}{2000/06/06}{
%   Corrected position of page number in footer of articles that open
%   on even pages. Ensured that first article in each session opens on
%   an odd page, as do list of abstracts and author index. Added extra
%   space in TOC before the backmatter. Added command for editors,
%   \cs{BoldSessionTitles}, that makes produces bold titles for
%   session in the TOC.}  \changes{1.8}{2000/05/16}{ Fixed bug in
%   last page label with \cs{suppressfulladdresses} }
% \changes{1.7}{2000/05/11}{ Fixed more little things. Page labels
%   now work, even if you have no bibliography. You can output the
%   full addresses by hand with the \cs{outputfulladdresses}
%   command. Editors' Guide updated but still not finished. Revamped
%   splitbook.sh now produces PS file of each article for web page and
%   also an EPS file of each page for sending to the printers.}
% \changes{1.6}{2000/04/15}{ More cleaning up: full stops after
%   titles in index, running headers and footers in frontmatter,
%   fontsize in title of index, etc} \changes{1.5}{2000/04/15}{
%   Eliminated superscripts from titles in TOC and index}
% \changes{1.4}{2000/04/14}{ Added copyright and license info (GPL)}
% \changes{1.3}{2000/03/01}{ Corrected \cs{figcaption} to do the
%   caption formatting right Corrected font size of sec headings in
%   front matter} \changes{1.2}{2000/02/16}{ Added \cs{plate}
%   command} \changes{1.1}{2000/02/14}{ Added manuscript option,
%   corrected bugs} \changes{1.0}{2000/02/08}{ Now does a whole book
%   - lots of stuff tidied up} \changes{0.7}{1999/10/04}{ Increased
%   spacing between sections in single column styles. Suppressed
%   floats at bottom of pages (didn't work though) }
% \changes{0.6}{1999/09/24}{ Tweaked \cs{adjustfinalcols} (still not
%   perfect) Added proceedings option Added onecolumn/twocolumn
%   options (orthogonal to journal/proceedings) }
% \changes{0.5}{1999/09/23}{ Added \cs{fulladdresses} and
%   \cs{adjustfinalcols} (just about working)}
% \changes{0.4}{1999/09/22}{ Minor stylistic changes}
% \changes{0.3}{1999/06/28}{ Better mechanism for alternative
%   affiliations Added `guide' option for user guide}
% \changes{0.2}{1999/05/27}{ All basic functionality present, I
%   think} \changes{0.1}{1999/05/26}{ Just about working - still
%   missing Fig/Tab/Refs} \changes{0.0}{1999/05/21}{ So bored today,
%   decided to rewrite RMAA LaTeX macros from scratch}
%    \begin{macrocode}
\ProvidesClass{rmaa}[\rmaa@date\space v\rmaa@version\space RMxAA article
class]
%    \end{macrocode}
% \subsection{Option Handling}
% \label{sec:options}
% Declarations that are needed in options
% \changes{3.05}{2003/10/22}{Added natbib option}
% \changes{3.06}{2003/01/03}{Natbib option turned on permanently}
% \changes{3.07}{2003/01/11}{Added proof option}
% \changes{3.14}{2003/01/11}{Option to turn off natbib}
% \changes{3.18}{2003/10/22}{trying out doctex mode} 
% \changes{3.23}{2007/09/09}{Manuscript is now on by default.}
% \changes{3.26}{2011/08/15}{New option for review articles}
%    \begin{macrocode}
\newif\if@reviewarticle
\@reviewarticlefalse
\newif\if@proof
\@prooffalse
\newif\if@journal
\@journaltrue
\newif\if@printdates
\@printdatesfalse
\newif\if@proceedings
\@proceedingsfalse
\newif\if@manuscript
\@manuscripttrue
\newif\if@bigfont
\@bigfonttrue
\newif\if@guide
\@guidefalse
\newif\if@preprint
\@preprintfalse
\newif\if@rmtwocol
\@rmtwocoltrue
\newif\if@biblist
\@biblistfalse
\newif\if@book
\@bookfalse
\newif\if@natbib
\@natbibtrue
\newif\if@rmcite
\@rmcitefalse
\newif\if@titlevisible
\@titlevisibletrue
\newif\if@abstractpresent
\@abstractpresentfalse
\newif\if@printbanner
\@printbannerfalse
\newif\if@suppressfulladdresses
\@suppressfulladdressesfalse
\newif\if@newstyle
\@newstylefalse
\newif\if@debug
\@debugfalse
\newif\if@boldkeywords
\@boldkeywordsfalse
%    \end{macrocode}
% The option handling---mainly setting various booleans.
% \changes{2.2beta11}{2001/11/26}{Made twocol and
% newstyle the default for proceedings.}
% \changes{3.23}{2007/09/09}{Manuscript is now 10 pt instead of 12pt}
% \changes{3.23}{2007/09/10}{Full addresses are now automatically
% suppressed for proceedings style.}
% \changes{3.24}{2007/09/14}{New option \texttt{boldkeywords} that
% does the keywords in the old ugly way. Now off by default.}
%    \begin{macrocode}
\DeclareOption{manuscript}{\typeout{Manuscript for referee} 
  \@manuscripttrue\@journalfalse\@rmtwocolfalse
  \PassOptionsToClass{10pt}{article}}
\DeclareOption{bigfont}{\typeout{Big font manuscript requested by Gary
  Ferland} 
  \@manuscripttrue\@journalfalse\@rmtwocolfalse\@bigfonttrue
  \PassOptionsToClass{11pt}{article}}
\DeclareOption{journal}{\typeout{Main journal option}%
  \@journaltrue\@rmtwocoltrue\@printdatestrue\@manuscriptfalse
  \PassOptionsToClass{10pt}{article}}
\DeclareOption{guide}{\typeout{User guide option}%
  \@guidetrue\@journalfalse\@rmtwocoltrue\@manuscriptfalse
  \PassOptionsToClass{10pt}{article}}
\DeclareOption{preprint}{\typeout{Preprint option}%
  \@preprinttrue\@manuscriptfalse\PassOptionsToClass{10pt}{article}}
\DeclareOption{proof}{\typeout{Author proof option}%
  \@prooftrue\@manuscriptfalse}
\DeclareOption{proceedings}{\typeout{Proceedings option}%
  \@proceedingstrue\@journalfalse\@rmtwocoltrue\@newstyletrue
  \@suppressfulladdressestrue
  \@printdatesfalse\@manuscriptfalse\PassOptionsToClass{10pt}{article}} 
\DeclareOption{printdates}{\typeout{Force printing of received/accepted dates}%
  \@printdatestrue} 
\DeclareOption{book}{\typeout{Whole issue option}%
  \@booktrue} 
\DeclareOption{nonatbib}{\typeout{Turn off natbib support (required
    for compatibility with old rmaacite package).}%
  \@natbibfalse} 
\DeclareOption{newstyle}{\typeout{Use new layout style}%
  \@newstyletrue}
\DeclareOption{oldstyle}{\typeout{Use old layout style}%
  \@newstylefalse}
\DeclareOption{debug}{\typeout{Loads of annoying messages}%
  \@debugtrue}
%    \end{macrocode}
% These two need to be declared explicitly so that |\ExecuteOptions|
% actually does something.
%    \begin{macrocode}
\DeclareOption{twocolumn}{\@rmtwocoltrue}
\DeclareOption{onecolumn}{\@rmtwocolfalse}
\DeclareOption{twoside}{\PassOptionsToClass{\CurrentOption}{article}}
\DeclareOption{boldkeywords}{\@boldkeywordstrue}
%    \end{macrocode}
% Bold, upper-case keywords used to be the default. Now they must be
% turned on explicitly with the \texttt{boldkeywords} option. 
%    \begin{macrocode}
\DeclareOption{boldkeywords}{\@boldkeywordstrue}
%    \end{macrocode}
% Pass any other option along to the \texttt{article} class.
%    \begin{macrocode}
\DeclareOption*{\PassOptionsToClass{\CurrentOption}{article}}
%    \end{macrocode}
% Default options are to be two-sided and based on the standard
% \LaTeX{} \texttt{article} class,
% unless we are processing the entire book, in which case we base
% ourselves on the standard \texttt{book} class. Default is to use
% manuscript option unless otherwise specified. 
%    \begin{macrocode}
\ExecuteOptions{twoside,manuscript} \ProcessOptions\relax
\if@book
\if@rmtwocol\LoadClass[twocolumn]{book}\else\LoadClass{book}\fi
\newcommand{\abstract}{\relax}
\newcommand{\@maketitle}{\relax}
\else
\if@rmtwocol\LoadClass[twocolumn]{article}\else\LoadClass{article}\fi
\fi
%    \end{macrocode}
% Attempt to use natbib for citations 02 Jan 2003. 
% \changes{3.05}{2003/01/02}{If the natbib option is used, then set
% things up so that natbib's version of the bibliography is identical
% to the standard one.}
% \changes{3.09}{2003/01/23}{Explicitly require a recent version of
% natbib, so that we get a decent error message if the package is too
% old.} 
% \changes{3.23}{2007/09/09}{Now allow Bibtex to work properly} 
% \changes{3.25}{2009/10/01}{Change \cs{newcommand} tp \cs{def} when setting \cs{bibfont}. This allows it to work with all versions of natbib.} 
% \changes{3.25}{2010/02/06}{Updated copyright year.} 
% \changes{3.25}{2010/02/06}{Updated email and web addresses to \texttt{crya.unam.mx}.} 
%    \begin{macrocode}
\if@natbib
\RequirePackage{natbib}[1999/02/25]
\bibpunct{(}{)}{;}{a}{}{,}
\def\bibfont{\small}
\renewcommand\bibsection{\section*{REFERENCES}}
\setlength\bibsep{0pt}
\setlength\bibhang{1.5em}
\bibliographystyle{rmaa}
\fi
%    \end{macrocode}
% We want the \texttt{graphicx} package for the extended syntax of
% |\includegraphics|. We want the \texttt{ifthen} package for using in
% these macros (although we don't much yet). We explicitly load
% \texttt{babel} since some \LaTeX{} distributions messed up the
% hyphenation if we didn't.  
%    \begin{macrocode}
\RequirePackage{graphicx}
\RequirePackage{ifthen}
\IfFileExists{babel.sty}{%
%   \let\@restoreindent=\@afterindentfalse
  \RequirePackage[english]{babel}
%   \@restoreindent
  %% reinstate the definition of \@roman so as not to break .idx files
  \@ifundefined{org@@roman}{\relax}{\renewcommand{\@roman}[1]{\org@@roman{##1}}}
  }{%
  \typeout{Package babel.sty not found - don't worry about this}}
%    \end{macrocode}
% Try to load the \texttt{dcolumn} and \texttt{booktabs} packages if
% they can be found. These can produce nicer-looking tables. 
%    \begin{macrocode}
\IfFileExists{dcolumn.sty}{%
  \RequirePackage{dcolumn}
  }{%
  \typeout{Package dcolumn.sty not found - you will not be able to
    align columns on decimal points in tables}}
\IfFileExists{booktabs.sty}{%
  \RequirePackage{booktabs}
  }{%
  \typeout{Package booktabs.sty not found - your tables will not look
    as pretty as they ought too.}
  }
\IfFileExists{makeidx.sty}{%
  \RequirePackage{makeidx}
}{%
  \typeout{Package makeidx.sty not found - Author Index may not work.}
  }
%% commented out for moment since requires more modern version of LaTeX
%% \RequirePackage[latin1]{inputenc}
%    \end{macrocode}
% \changes{4.0}{2019/11/13}{Make \texttt{hyperref} and
% \texttt{etoolbox} into hard requirements, since we now use them
% extensively.}  We use hyperref to get internal and external links.
% But I don't like the colored boxes so I turn those off.  We use
% etoolbox for patching and robustifying at the moment.  In the
% future, I want to port a lot of other stuff over to it.  For
% instance, replace the ifthen package and the TeX ifs too.
%
% We target the latest version (1.7) of PDF, rather than the default
% of 1.5.  This reduces the chance of getting warnings from included
% graphics.  This is not strictly related to hyperref, but we put it
% here anyway.
%    \begin{macrocode}
\RequirePackage{etoolbox}
\RequirePackage{hyperref}
\hypersetup{hidelinks}
\pdfminorversion=7
%    \end{macrocode}
% We patch the \cs{@combinedblfloats} following a suggestion by Heiko
% Oberdiek at \url{https://tex.stackexchange.com/a/249743/3843}, which
% avoids problems when hyperlinks are split across
% pages. \changes{4.0}{2019/11/12}{Fix for issue where breaking
% hyperlink across pages gave an error.}
%    \begin{macrocode}
\urlstyle{same}
\newcommand\email{\begingroup \urlstyle{rm}\Url}
\patchcmd\@combinedblfloats{\box\@outputbox}{\unvbox\@outputbox}{}{%
  \errmessage{\noexpand\@combinedblfloats could not be patched}%
}%
%    \end{macrocode}
% \subsection{Other setup stuff}\label{sec:setup}
% Some declarations: 
%    \begin{macrocode}
\newif\if@needaddresses
\@needaddressestrue
\newif\if@frontmatter
\@frontmatterfalse
\newif\if@inpreamble
\@inpreambletrue
\newif\if@firstauthor
\@firstauthortrue
\newif\if@firstaffil
\@firstaffiltrue
\newif\if@firstaltaffil
\@firstaltaffiltrue
\newif\if@firstkwd
\@firstkwdtrue
\newif\if@noexpandADS
\@noexpandADSfalse
\newif\if@sessiontitlepages
\@sessiontitlepagesfalse
%    \end{macrocode}
% \begin{macro}{\@rmaaEmpty}
% \begin{macro}{\@rmaaArg}
% These are used in various commands to check if arg was left blank
%    \begin{macrocode}
\newcommand{\@rmaaEmpty}{}
\newcommand{\@rmaaArg}{}
%    \end{macrocode}  
% \end{macro}
% \end{macro}
%
% \subsection{Compatibility with old \LaTeX{} versions}
% \label{sec:compat}
% This package is based on the 1998/12/01 release of \LaTeX{}.
% Older releases may not define all the commands that we need.
% Luckily, we can use |\providecommand| to define these if necessary. 
% The following commands are not defined in the 1994/06/01 release. 
%
% \begin{macro}{\MakeUpperCase}
% \begin{macro}{\textsuperscript}
% \begin{macro}{\textbackslash}
% User commands:
%    \begin{macrocode}
\providecommand{\MakeUpperCase}[1]{\uppercase{#1}}
\providecommand{\textsuperscript}[1]{\ensuremath{^\textrm{#1}}}
\providecommand{\textbackslash}{\ensuremath{\backslash}}
%    \end{macrocode}  
% \end{macro}
% \end{macro}
% \end{macro}
% \begin{macro}{\protected@xdef}
% \begin{macro}{\restore@protect}
% \begin{macro}{\@unexpandable@protect}
% \begin{macro}{\hb@xt@}
% \LaTeX{} internal commands:
%    \begin{macrocode}
\providecommand{\protected@xdef}{%
        \let\@@protect\protect
        \let\protect\@unexpandable@protect
        \afterassignment\restore@protect
        \xdef}
\providecommand{\restore@protect}{\let\protect\@@protect}
\providecommand{\@unexpandable@protect}{\noexpand\protect\noexpand}
\providecommand{\hb@xt@}{\hbox to}
%    \end{macrocode}
% \end{macro}
% \end{macro}
% \end{macro}
% \end{macro}
%
% \subsection{Page layout}\label{sec:page-layout}
%
% The dimensions of the text were defined by Elizabeth. Text width is
% 169\,mm with equal margins on the two sides: 
%    \begin{macrocode}
\setlength{\textwidth}{169mm}
\setlength{\@tempdima}{\paperwidth}
\addtolength{\@tempdima}{-\textwidth}
\setlength{\oddsidemargin}{0.5\@tempdima}
\addtolength{\oddsidemargin}{-1in}
\setlength{\evensidemargin}{\oddsidemargin}
%    \end{macrocode}
% Text height is 225\,mm, top margin is 14\,mm. 
%    \begin{macrocode}
\setlength{\textheight}{225mm} 
\setlength{\topmargin}{-1in} 
\addtolength{\topmargin}{14mm}
%    \end{macrocode}
% Space between columns in two-column styles, header, footer,
% paragraph indentation:
%    \begin{macrocode}
\setlength{\columnsep}{9.5mm}
\setlength{\headsep}{0.7cm}
\setlength{\footskip}{0.75cm}
\setlength{\parindent}{1.5em}
%    \end{macrocode}
% 
% \begin{macro}{\@set@scalable@length}
% Try out a new scheme for dealing with various spacing dimensions and
% allowing the user to easily modify them. The command
% \begin{quote}
% |\set@scalable@length|\marg{name}\marg{defaultlength}\marg{stretch}\marg{shrink}
% \end{quote}
% does a 
% |\setlength|\marg{\cs{name}}\marg{defaultlength plus stretch minus shrink}
% and also defines the user command
% \textbackslash\meta{name}|Scale|\marg{factor} that can 
% be used to rescale length by \meta{factor}. Additionally, now
% defines the package command
% \textbackslash\meta{name}|@set@scale|\marg{FACTOR} 
% that also rescales the length by \meta{FACTOR}, so that the length
% becomes \meta{factor} $\times$ \meta{FACTOR} $\times$
% \meta{defaultlength}.   
%    \begin{macrocode}
\newcommand\@set@scalable@length[4]{%
  %% arg #1: NAME of length without leading `\'
  %%
  %% arg #2: default VALUE  
  \expandafter\newlength\csname @#1@default\endcsname
  \expandafter\setlength\csname @#1@default\endcsname{#2}%
  \expandafter\newlength\csname @@#1@@default\endcsname
  \expandafter\setlength\csname @@#1@@default\endcsname{#2}%
  %% arg #3: default STRETCH  
  \expandafter\newlength\csname @#1@stretch\endcsname
  \expandafter\setlength\csname @#1@stretch\endcsname{#3}%
  \expandafter\newlength\csname @@#1@@stretch\endcsname
  \expandafter\setlength\csname @@#1@@stretch\endcsname{#3}%
  %% arg #4: default SHRINK  
  \expandafter\newlength\csname @#1@shrink\endcsname
  \expandafter\setlength\csname @#1@shrink\endcsname{#4}%
  \expandafter\newlength\csname @@#1@@shrink\endcsname
  \expandafter\setlength\csname @@#1@@shrink\endcsname{#4}%
  %%
  %% User-defined scale \@NAME@scale is unity by default
  \expandafter\gdef\csname @#1@scale\endcsname{1}%
  %% Package-defined scale \@NAME@SCALE is unity by default
  \expandafter\gdef\csname @#1@SCALE\endcsname{1}%
  %%
  %% internal command \@reset@NAME should expand to 
  %% \setlength\@@NAME@@default{\@NAME@SCALE\@NAME@default}
  %% \setlength\NAME{\@NAME@scale\@@NAME@@default}
  %% in other words: scale x SCALE x default 
  \expandafter\gdef\csname @reset@#1\endcsname{%
    \expandafter\setlength\csname @@#1@@default\endcsname{%
      \csname @#1@SCALE\endcsname \csname @#1@default\endcsname
      }
    \expandafter\setlength\csname @@#1@@stretch\endcsname{%
      \csname @#1@SCALE\endcsname \csname @#1@stretch\endcsname
      }
    \expandafter\setlength\csname @@#1@@shrink\endcsname{%
      \csname @#1@SCALE\endcsname \csname @#1@shrink\endcsname
      }
    \expandafter\setlength\csname #1\endcsname{%
      \csname @#1@scale\endcsname \csname @@#1@@default\endcsname
      \@plus
      \csname @#1@scale\endcsname \csname @@#1@@stretch\endcsname
      \@minus
      \csname @#1@scale\endcsname \csname @@#1@@shrink\endcsname
      }}
  %%
  %% User command \NAMEScale should expand to 
  %% {\renewcommand\@NAME@scale{#1}\@reset@NAME}
  \expandafter\gdef\csname #1Scale\endcsname##1{%
    \expandafter\gdef\csname @#1@scale\endcsname{##1}%
    \csname @reset@#1\endcsname}
  %%
  %% Package command \NAME@set@SCALE should expand to 
  %% {\renewcommand\@NAME@SCALE{#1}\@reset@NAME}
  \expandafter\gdef\csname #1@set@SCALE\endcsname##1{%
    \expandafter\gdef\csname @#1@SCALE\endcsname{##1}%
    \csname @reset@#1\endcsname}
  %%
  %% initialize at default values with \@reset@NAME
  \csname @reset@#1\endcsname
  }
%    \end{macrocode}
% \end{macro}
% \begin{macro}{\parskip}
% First use of this new infrastructure: set interparagraph space to be
% zero but with a small amount of stretchiness.
%    \begin{macrocode}
%   \ifthenelse{\boolean{@manuscript}}{%
%     \@set@scalable@length{parskip}{0.5\baselineskip}{0.3\p@}{0.1\p@}
%   }{
    \@set@scalable@length{parskip}{0\p@}{0.3\p@}{0\p@}
%   }
%    \end{macrocode}
% \end{macro}
% 
% \begin{macro}{\abovesectionskip}
% \changes{4.0}{2019/11/12}{Use negative values for skip so that first paragraph after section headings is not indented.  This is something that I have been trying to fix for years!  It turns out that the anser was in the documentation for the standard command \cs{@startsection.}}
% \begin{macro}{\belowsectionskip}
% \begin{macro}{\aboveparaskip}
% \begin{macro}{\belowparaskip}
% Define some utility rubber lengths for use in
% section~\ref{sec:section-commands}. 
%    \begin{macrocode}
\newlength{\abovesectionskip}
\newlength{\belowsectionskip}
\newlength{\aboveparaskip}
\newlength{\belowparaskip}
%    \end{macrocode}
% The default values for these lengths are those traditionally used
% for the onecolumn style. 
%    \begin{macrocode}
\@set@scalable@length{abovesectionskip}{-3.5ex}{-4ex}{-2ex}
\@set@scalable@length{belowsectionskip}{2.3ex}{0ex}{2ex}
\@set@scalable@length{aboveparaskip}{3.25ex}{2ex}{0.5ex}
\@set@scalable@length{belowparaskip}{-1em}{\z@}{\z@}
%    \end{macrocode}
% \end{macro}\end{macro}\end{macro}\end{macro}
% We also make scalable some of the standard float parameters. These
% are the same defaults as in the standard article classes, but
% slightly more stretchy. For the twocolumn style, they are currently
% halved, see below. Note that I have also lumped |\abovecaptionskip|
% in with these, since it plays a similar role. There is also a
% |\belowcaptionskip| but it is set to zero in the standard classes. 
%    \begin{macrocode}
\@set@scalable@length{floatsep}{12\p@}{4\p@}{4\p@}
\@set@scalable@length{textfloatsep}{20\p@}{4\p@}{8\p@}
\@set@scalable@length{intextsep}{12\p@}{4\p@}{4\p@}
\@set@scalable@length{dblfloatsep}{12\p@}{4\p@}{4\p@}
\@set@scalable@length{dbltextfloatsep}{20\p@}{4\p@}{8\p@}
\@set@scalable@length{abovecaptionskip}{10\p@}{\z@}{\z@}
%    \end{macrocode}
% \begin{macro}{\@before@title@skip}         
% \begin{macro}{\@before@author@skip}        
% \begin{macro}{\@between@author@skip}       
% \begin{macro}{\@before@abstract@skip}      
% \begin{macro}{\@after@abstracthead@skip}   
% \begin{macro}{\@between@abstract@skip}     
% \begin{macro}{\@before@kwd@skip}           
% \begin{macro}{\@before@dates@skip}           
% Spacing parameters used in |\@maketitle|. These are not changed by
% default except by the user, so they don't need to go into
% |\@package@default@lengths|.  
%    \begin{macrocode}
\newlength{\@before@title@skip}         
\newlength{\@before@author@skip}        
\newlength{\@between@author@skip}       
\newlength{\@before@abstract@skip}      
\newlength{\@after@abstracthead@skip}   
\newlength{\@between@abstract@skip}     
\newlength{\@before@kwd@skip}           
\newlength{\@before@dates@skip}           
\@set@scalable@length{@before@title@skip}{2em}{\z@}{\z@}
\@set@scalable@length{@before@author@skip}{1.5em}{\z@}{\z@}
\@set@scalable@length{@between@author@skip}{0.5em}{\z@}{\z@}
\@set@scalable@length{@before@abstract@skip}{3ex}{1ex}{0.5ex}
\@set@scalable@length{@after@abstracthead@skip}{1ex}{\z@}{\z@}
\@set@scalable@length{@between@abstract@skip}{2ex}{1ex}{0.5ex}
\@set@scalable@length{@before@kwd@skip}{1ex}{0.5ex}{0.2ex}
\@set@scalable@length{@before@dates@skip}{2ex}{0.5ex}{0.2ex}
%    \end{macrocode}
% \end{macro}\end{macro}\end{macro}\end{macro}
% \end{macro}\end{macro}\end{macro}\end{macro}
% 
% Spacing above and below displayed equations also wants to be done
% like this, but it is more complicated since the standard article
% classes make these change with the fontsize. Something to do later,
% maybe. 
% 
%    \begin{macrocode}
%    \end{macrocode}
% \begin{macro}{\@package@default@lengths}
% The value actually used for these lengths is controlled by the
% \meta{SCALE} parameter, which depends on the style option that is in
% effect. This is set using the |\@package@default@lengths| command,
% which must be called every time the style changes (e.g., changing
% from onecolumn to twocolumn). 
%    \begin{macrocode}
\newcommand{\@package@default@lengths}{%
  \ifthenelse{\boolean{@rmtwocol}}{%
    %% twocolumn mode
    \abovesectionskip@set@SCALE{0.5}
    \belowsectionskip@set@SCALE{0.5}
    \aboveparaskip@set@SCALE{0.5}
    \belowparaskip@set@SCALE{1}
    \floatsep@set@SCALE{0.5}
    \textfloatsep@set@SCALE{0.5}
    \intextsep@set@SCALE{0.5}
    \dblfloatsep@set@SCALE{0.5}
    \dbltextfloatsep@set@SCALE{0.5}
    }{%
    %% onecolumn mode
    \abovesectionskip@set@SCALE{1}
    \belowsectionskip@set@SCALE{1}
    \aboveparaskip@set@SCALE{1}
    \belowparaskip@set@SCALE{1}
    \floatsep@set@SCALE{1}
    \textfloatsep@set@SCALE{1}
    \intextsep@set@SCALE{1}
    \dblfloatsep@set@SCALE{1}
    \dbltextfloatsep@set@SCALE{1}
    }%
  }
\@package@default@lengths
%    \end{macrocode}
% \end{macro}
% \begin{macro}{\RescaleLengths}
% \begin{macro}{\RescaleSecLengths}
% \begin{macro}{\RescaleFloatLengths}
% \begin{macro}{\RescaleTitleLengths}
% User commands that will fiddle with all the scalable lengths. 
%    \begin{macrocode}
\newcommand{\RescaleLengths}[1]{%
  \RescaleSecLengths{#1}\RescaleFloatLengths{#1}\RescaleTitleLengths{#1}}
\newcommand{\RescaleSecLengths}[1]{%
%% sections
  \abovesectionskipScale{#1}
  \belowsectionskipScale{#1}
  \aboveparaskipScale{#1}
  }
\newcommand{\RescaleFloatLengths}[1]{%
%% floats
  \floatsepScale{#1}
  \textfloatsepScale{#1}
  \intextsepScale{#1}
  \dblfloatsepScale{#1}
  \dbltextfloatsepScale{#1}
  \abovecaptionskipScale{#1}
  }
\newcommand{\RescaleTitleLengths}[1]{%
%% maketitle
  \@before@title@skipScale{#1}           
  \@before@author@skipScale{#1}        
  \@between@author@skipScale{#1}       
  \@before@abstract@skipScale{#1}      
  \@after@abstracthead@skipScale{#1}   
  \@between@abstract@skipScale{#1}     
  \@before@kwd@skipScale{#1}           
  \@before@dates@skipScale{#1}           
  }
%    \end{macrocode}
% \end{macro}\end{macro}\end{macro}\end{macro}
% \changes{2.2beta14}{2002/03/17}{Finished implementation of new
% infrastructure for scalable lengths. Converted the above/below
% section skips to use it. Also, all the floatsep parameters. Also,
% all the vertical lengths that are used by |\@maketitle|. }

% \begin{macro}{\MSextramargin}
% For the \texttt{manuscript} option we have an extra 0.3\,inch
% vertical margin and 1\,inch horizontal margin for marginal notes (by
% hand or with |\marginpar|. 
%
% \changes{3.23}{2007/09/09}{Reduce baselinestretch in
%   \texttt{manuscript} style to 1.2. No longer use raggedright for
%   manuscript.}
%
% \changes{3.24}{2007/09/14}{Further reduction of baselinestretch in
%   \texttt{manuscript} style to 1.1 since we now have wider margins.}
%
% \changes{3.24}{2007/09/14}{Split \texttt{manuscript}'s extra margin
%   into separate horizontal and vertical parts.}
% 
%    \begin{macrocode}
\newlength{\MSextraHmargin}\newlength{\MSextraVmargin}
\setlength{\MSextraHmargin}{1in} 
\setlength{\MSextraVmargin}{0.3in} 
\if@manuscript
   \renewcommand{\baselinestretch}{1.1}
   \addtolength{\textwidth}{-2\MSextraHmargin}
   \addtolength{\oddsidemargin}{\MSextraHmargin}
   \setlength{\evensidemargin}{\oddsidemargin}
   \addtolength{\textheight}{-2\MSextraVmargin}
   \addtolength{\topmargin}{\MSextraVmargin}
%    \raggedright
\fi
%    \end{macrocode}
% \end{macro}
% Float parameters: don't allow bottom floats by default. 
%    \begin{macrocode}
\renewcommand{\fps@figure}{tp}
\renewcommand{\fps@table}{tp}
%    \end{macrocode}
%
% \begin{macro}{\rmaa@AuthorWidthFrac}
% \begin{macro}{\rmaa@AbstractWidthFrac}
% These parameters control what fraction of the total text width can
% be used by the list of authors and the abstract/resumen,
% respectively.  
%    \begin{macrocode}
\newcommand{\rmaa@AuthorWidthFrac}{1.0}
\newcommand{\rmaa@AbstractWidthFrac}{0.75}
%    \end{macrocode}
% Use a wide abstract to save space in the new style.
%    \begin{macrocode}
\if@newstyle\renewcommand{\rmaa@AbstractWidthFrac}{1.0}\fi
%    \end{macrocode}
% \end{macro}
% \end{macro}
%
% Some things to consider in future versions. Trying to avoid too many
% widows, etc. Apparently, the following will work best with
% \cs{flushbottom}. 
%    \begin{macrocode}
%%\widowpenalty=10000
%%\clubpenalty=10000
%    \end{macrocode}
% Also, should perhaps try getting rid of underfull/overfull box warnings. 
% \changes{3.10}{2003/01/24}{Set large values for \cs{hbadness},
% \cs{vbadness}, \cs{hfuzz}, \cs{vfuzz}. To reduce amount of stuff
% that gets written to the terminal every time you compile something.}
%    \begin{macrocode}
\hbadness=10000
\vbadness=10000
\setlength\hfuzz{\maxdimen}
\setlength\vfuzz{\maxdimen}
%    \end{macrocode}
% \subsection{Various text styles for different contexts}
% \label{sec:text-styles} 
% \begin{macro}{\rmaa@AffilStyle}
% \begin{macro}{\rmaa@PagesStyle}
% \begin{macro}{\rmaa@YearStyle}
% \begin{macro}{\rmaa@VolumeStyle}
% \begin{macro}{\rmaa@VolumenStyle}
% \begin{macro}{\rmaa@JournalNameStyle}
% \begin{macro}{\rmaa@ReceivedStyle}
% \begin{macro}{\rmaa@AbstractHeaderStyle}
% \begin{macro}{\rmaa@JournalNameStyle}
% Style commands for the different fields used by |\maketitle| in
% constructing the title page of the article. These have the form
% |\rmaa@XYZStyle{TEXT}|. I'm not sure these were really a good idea,
% but never mind \dots 
%    \begin{macrocode}
\newcommand{\rmaa@AffilStyle}[1]{{\normalsize\rmfamily#1}}
\newcommand{\rmaa@PagesStyle}[1]{{\normalsize\rmfamily#1}}
\newcommand{\rmaa@YearStyle}[1]{{\normalsize\rmfamily#1}}
\newcommand{\rmaa@VolumeStyle}[1]{{\normalsize\rmfamily\bfseries#1}}
\newcommand{\rmaa@VolumenStyle}[1]{{\Large\sffamily\bfseries#1}}
\newcommand{\rmaa@JournalNameStyle}[1]{{\normalsize\itshape#1}}
\newcommand{\rmaa@ReceivedStyle}[1]{{\normalsize\rmfamily\itshape#1}}
\newcommand{\rmaa@AbstractHeaderStyle}[1]{%
  {\normalsize\rmfamily\MakeUpperCase{#1}}}
\newcommand{\rmaa@TitleStyle}[1]{%
  {\large\bfseries\rmfamily\boldmath\MakeUpperCase{#1}\par}}
%    \end{macrocode}
% \end{macro}
% \end{macro}
% \end{macro}
% \end{macro}
% \end{macro}
% \end{macro}
% \end{macro}
% \end{macro}
% \end{macro}
%
% \begin{macro}{\rmaa@FrontSecStyle}
% Headings of ``front sections'' (see section~\ref{sec:front-sections})
%    \begin{macrocode}
\newcommand{\rmaa@FrontSecStyle}[1]{%
  {\normalsize\bfseries\rmfamily\MakeUpperCase{#1}}}
%    \end{macrocode}
% \end{macro}
%
% \begin{macro}{\rmaa@TOCTitleStyle}
% \begin{macro}{\rmaa@TOCPartStyle}
% \begin{macro}{\rmaa@IndexTitleStyle}
% \begin{macro}{\rmaaTOCAuthorStyle}
% \begin{macro}{\rmaaIndexAuthorStyle}
% Commands to set style of TOC and index entries.
% Rewritten to avoid unnecessary font commands in the .toc and .ind files
%    \begin{macrocode}
\newcommand{\rmaa@TOCTitleStyle}[1]{\MakeUpperCase{#1}}
\newcommand{\rmaa@TOCPartStyle}[1]{\rmaa@TOCTitleStyle{#1}}
\newcommand{\rmaa@IndexTitleStyle}[1]{#1}
\newcommand{\rmaaTOCAuthorStyle}[1]{%
  {\itshape#1}}
\newcommand{\rmaaIndexAuthorStyle}[1]{%
  \protect\mbox{\bfseries#1}}
%    \end{macrocode}
% \end{macro}
% \end{macro}
% \end{macro}
% \end{macro}
% \end{macro}
%
% \begin{macro}{\rmaa@KeywordStyle}
% \begin{macro}{\rmaa@KeywordHeaderStyle}
% The style of the keywords depends on whether we are using the
% \texttt{newstyle} option or not. If we are, then they are small. 
%    \begin{macrocode}
\ifthenelse{\boolean{@newstyle}}{
  \newcommand{\rmaa@KeywordStyle}[1]{%
    {\small%
      \if@boldkeywords 
      \bfseries\MakeUpperCase{\raggedright #1}
      \else
      \raggedright #1
      \fi
    }}
  \newcommand{\rmaa@KeywordHeaderStyle}[1]{%
    {\small\itshape#1}}}
{
  \newcommand{\rmaa@KeywordStyle}[1]{%
    {\normalsize%
      \if@boldkeywords\bfseries\fi\raggedright #1}}
  \newcommand{\rmaa@KeywordHeaderStyle}[1]{%
    {\normalsize\itshape#1}}}
%    \end{macrocode}
% \end{macro}
% \end{macro}
%
% \subsection{Page styles}\label{sec:page-styles}
% \begin{macro}{\ps@rmaa}
% The running headers for normal pages:
%    \begin{macrocode}
\newcommand{\ps@rmaa}{%
%    \end{macrocode}
% nothing in either foot; 
%    \begin{macrocode}
  \renewcommand\@oddfoot{\relax}%
  \renewcommand\@evenfoot{\relax}%
%    \end{macrocode}
% page number and short author list in the even heads; 
%    \begin{macrocode}
  \renewcommand\@evenhead{\rlap{\thepage}\hfil\@shortauthor\hfil}%
%    \end{macrocode}
% short title and page number in the odd heads.
%    \begin{macrocode}
  \renewcommand\@oddhead{\hfil\@shorttitle\hfil\llap{\thepage}}
  }
%    \end{macrocode}
% \begin{macro}{\ps@rmaafirstpage}
% The header of the first page of each
% article:
%    \begin{macrocode}
\newcommand{\ps@rmaafirstpage}{%
  \renewcommand\@oddfoot{\hfil\thepage}%
%    \end{macrocode}
% footer corrected (06 Jun 2000) to have page on outside edge;
%    \begin{macrocode}
  \renewcommand\@evenfoot{\thepage\hfil}%
%    \end{macrocode}
% Header has journal name, volume, and page numbers for a regular
% article, or similar for special cases. 
% \changes{4.0}{2019/11/03}{Change to three-line header with copyright info and DOI.}
%    \begin{macrocode}
\if@preprint
\renewcommand\@oddhead{%
  \if@proceedings 
  \textsc{To appear in ``\@conftitle{} (\rmaa@YearStyle{\@year})''}\hfil
  \textit{RMxAC}
  \else
  \textsc{To appear in} \@journalname\hfil
  \fi}\fi%
\if@guide
\renewcommand\@oddhead{\textsc{Documentation} (\rmaa@YearStyle{\@year})\hfil}
\fi%
\if@manuscript
\renewcommand\@oddhead{Manuscript for \@journalname{} (\rmaa@YearStyle{\@year})\hfil}
\fi%
\if@journal
\renewcommand\@oddhead{%
  \parbox{\textwidth}{
    \@journalname, \rmaa@VolumeStyle{\@volume},
    \rmaa@PagesStyle{\@artpages}\
    (\rmaa@YearStyle{\@year})\hfil\\[1pt]
    {\small\textcopyright{} \@year:
      Instituto de Astronom\'\i{}a,
      Universidad Nacional Aut\'onoma de M\'exico}\hfil\\[1pt]
    {\small\@DOI@url}\hfil
  }%
}\fi%
  \renewcommand\@evenhead{\@oddhead}}
%    \end{macrocode}
% By default, use the |\ps@rmaa| page style.
%    \begin{macrocode}
\pagestyle{rmaa}
%    \end{macrocode}
% \end{macro}
% \end{macro}
%
% \subsection{Preamble commands}\label{sec:preamble} 
% Everything you need for setting the title page, TOC and index
% entries, etc.
% \subsubsection{Commands for the title page}
% \label{sec:title-commands}
% \begin{macro}{\resumen}
% \begin{macro}{\@resumen}
% \begin{macro}{\abstract}
% \begin{macro}{\@abstract}
% User commands for setting the resumen and abstract. 
% These should be used in the document header. The internal commands
% |\@resumen| and |\@abstract| are used to hold the actual text
% pending output via the |\maketitle| command. 
%    \begin{macrocode}
\newif\if@setresumen
\newcommand{\resumen}[1]{%
%%  \renewcommand{\@resumen}{{\selectlanguage{spanish}#1}}%
  \renewcommand{\@resumen}{{\ignorespaces#1}}%
  \@setresumentrue}
\renewcommand{\abstract}[1]{%
  \renewcommand{\@abstract}{{\ignorespaces#1}}
  \renewcommand\ADS@abstract{#1}
%% Not a good idea, since then author macros are not expanded and
%% cause errors later. However, by not doing it we are vulnerable to
%% fragile commands.....
  \if@noexpandADS\@onelevel@sanitize{\ADS@abstract}\fi
  \@abstractpresenttrue
  }
\newcommand{\ADSnoexpand}{\@noexpandADStrue}
%    \end{macrocode}
% \begin{macro}{\boldabstract}
% This is the version for doing 1-page conference proceedings. Still
% save the text for ADS but then write it out immediately as a bold
% paragraph. 
%    \begin{macrocode}
\newcommand{\boldabstract}[1]{%
  \gdef\ADS@abstract{#1}%
  \renewcommand{\@abstract}{#1}%
  \if@noexpandADS\@onelevel@sanitize{\ADS@abstract}\fi%
  \noindent\textbf{\boldmath #1}
  }
%    \end{macrocode}
% \end{macro}
% Default value in each case is to print out a warning if not set
%    \begin{macrocode}
  \newcommand{\@resumen}{%
    \@setresumenfalse \ClassWarningNoLine{rmaa}{No Spanish abstract
      (resumen) specified - if you do not speak Spanish, this will be
      translated by the editors}%
    \textbf{Favor de proporcionar un resumen en espa\~nol. If you are
      unable to translate your abstract into Spanish, the editors will
      do it for you.} }
  \newcommand{\@abstract}{\ClassWarningNoLine{rmaa}{No abstract
      specified}}
%    \end{macrocode}
% \end{macro}
% \end{macro}
% \end{macro}
% \end{macro}
%
% \begin{macro}{\title}
% \begin{macro}{\@rawtitle}
% \begin{macro}{\@TOCtitle}
% \begin{macro}{\rmaaIndexTitle}
% Similarly, |\title| is the preamble command to set the title of the
% article. This is used in various places, typeset in different
% styles, hence the use of the internal |\@rawtitle|, |\@title|,
% |\rmaaIndexTitle|, and |\@TOCtitle| commands, each of which uses the
% relevant style from section~\ref{sec:text-styles}
%    \begin{macrocode}
\newcommand{\@rawtitle}{Raw Title}
\newcommand{\@TOCtitle}{\rmaa@TOCTitleStyle{Title for TOC}}
\newcommand{\rmaaIndexTitle}{\rmaa@IndexTitleStyle{Title for index}.}
\newcommand{\@TOCaftertitle}{}
\renewcommand{\title}[1]{%
%%  {\ADS@sanitize\xdef\@rawtitle{#1}}
  \renewcommand\@rawtitle{\noexpand#1}
  \let\ADS@title=\@rawtitle
  \renewcommand{\@title}{\rmaa@TitleStyle{#1}}%
  \renewcommand{\@TOCtitle}{\rmaa@TOCTitleStyle{#1}\@TOCaftertitle}%
  \renewcommand{\rmaaIndexTitle}{\rmaa@IndexTitleStyle{#1\@TOCaftertitle}.}%
  }
%    \end{macrocode}
% \end{macro}
% \end{macro}
% \end{macro}
% \end{macro}
% 
% All the weird and whacky ways of indicating the author
% affiliations. 
% \begin{macro}{\affil}
% First, the way that sets it directly underneath the authors name.
% Note that |\affil| trumps all other affiliations for the purposes of
% getting into the ADS stuff (we use |\gdef|).   
% \dots 
%    \begin{macrocode}
\newcommand{\affil}[1]{%
  \\[1.5ex]\rmaa@AffilStyle{#1}\\[1.5ex]%
  \if@firstaffil\gdef\ADS@affil{#1}\@firstaffilfalse\fi%
  }
\renewcommand{\and}{and\\[1.5ex]}
%    \end{macrocode}
% \begin{macro}{\altaffil}
% Alternative affiliations are footnoted. 
%    \begin{macrocode}
\newcommand{\altaffil}[1]{%
  \if@firstaltaffil\gdef\ADS@affil{#1}\@firstaltaffilfalse\fi%
  \thanks{#1}}
%    \end{macrocode}
% \begin{macro}{\altaffilmark}
% \begin{macro}{\altaffiltext}
% Alternative-alternative mechanism for affiliations---should not be
% mixed with |\altaffil| (unless you know what you are doing) or there
% will be problems with numbering.
%    \begin{macrocode}
\newcommand{\altaffilmark}[1]{\textsuperscript{#1}}
\newcommand{\altaffiltext}[2]{%
  \if@firstaltaffil\gdef\ADS@affil{#2}\@firstaltaffilfalse\fi
  \stepcounter{footnote}%
  \protected@xdef\@thanks{\@thanks
    \protect\footnotetext[#1]{#2}}} 
%    \end{macrocode}
% \end{macro}
% \end{macro}
% \end{macro}
% \end{macro}
%
% \begin{macro}{\keywords}
% \begin{macro}{\@keywords}
% \begin{macro}{\@keywordheader}
% \label{keywords}Sundry minor preamble commands. All
% self-explanatory, I hope. |\keywords| is the old user interface
% (deprecated). New articles should use |\addkeyword| once for each
% keyword. That way the ADS record comes out right too (comma rather
% than em-dash separator).
% \changes{2.2beta12}{2001/12/07}{\cs{addkeyword} New user interface
%   for keywords.} 
% \changes{2.2beta13}{2002/03/11}{\cs{addkeyword} Improved behaviour
%   with .ads file. Use
%   \cs{ADS@sanitize} to change definitions of `~' and \cs{ion}.}
% \changes{3.23}{2007/09/09}{\cs{addkeyword} now signals a compilation
% error if the author tries to add more than 6 keywords.}
% \changes{3.24}{2007/09/14}{kwdcounter is now reset in
% \cs{@begofarticlehook} to prevent error in book style} 
%    \begin{macrocode}
\newcommand{\@kwdsep}{ --- }
\newcommand{\ADS@kwdsep}{, }
\newcounter{kwdcounter}\setcounter{kwdcounter}{0}
\newcommand{\addkeyword}[1]{%
  \addtocounter{kwdcounter}{1}
  \ifthenelse{\thekwdcounter > 6}{
    \ClassError{rmaa}{Too many keywords. A maximum of 6 keywords is allowed.}{}
  }{}
  \if@firstkwd%
  \renewcommand\@keywords{\if@boldkeywords\uppercase{#1}\else #1\fi}
  {\ADS@sanitize\def\ion{\plainion}\gdef\ADS@keywords{#1}}
  \@firstkwdfalse
  \else%
  \edef\@keywords{\@keywords\@kwdsep \if@boldkeywords\uppercase{#1}\else #1\fi}%
  {\ADS@sanitize\def\ion{\plainion}\xdef\ADS@keywords{\ADS@keywords\ADS@kwdsep #1}}%
  \fi}
\newcommand{\ADS@sanitize}{%
  \def\mathrm{}
  \def\@{\ }
  \def~{\ }
  \def\langle{<}
  \def\rangle{>}
  \def\log{log }
  \def\sin{sin }
  \def\cos{cos }
  \def\tan{tan }
  \def\fd{^d.}
  \def\fh{^h.}
  \def\fm{^m.}
  \def\fs{^s.}
  \def\fdg{�.}
  \def\farcm{'.}
  \def\farcs{''.}
  \def\arcdeg{�}
  \def\,{ }
  }
\newcommand{\keywords}[1]{%
  \renewcommand{\ADS@keywords}{#1}
  \renewcommand{\@keywords}{\uppercase{#1}}
  }
\newcommand{\@keywordheader}{\rmaa@KeywordHeaderStyle{Key Words:}}
\newcommand{\@keywords}{%
  \ClassWarningNoLine{rmaa}{No keywords specified}%
  Keywords go here
  }
%    \end{macrocode}
% \begin{macro}{\shorttitle}
% \begin{macro}{\shortauthor}
% \begin{macro}{\@shorttitle}
% \begin{macro}{\@shortauthor}
% Commands for setting the running headers.
%    \begin{macrocode}
\newcommand{\shorttitle}[1]{%
  \renewcommand\@shorttitle{\MakeUpperCase{#1}}}
\newcommand{\shortauthor}[1]{%
  \renewcommand\@shortauthor{\MakeUpperCase{#1}}}
\newcommand{\@shorttitle}{SHORT TITLE GOES HERE}
\newcommand{\@shortauthor}{SHORT AUTHOR GOES HERE}
%    \end{macrocode}
% \end{macro}
% \end{macro}
% \end{macro}
% \end{macro}
% \end{macro}
% \end{macro}
% \end{macro}


% \subsubsection{The full addresses}
% \label{sec:full-addresses}
% \begin{macro}{\fulladdresses}
% \begin{macro}{\@fulladdresses}
% \begin{macro}{\@fulladdheight}
% \begin{macro}{\@fulladddepth}
% \begin{macro}{\@fulladdbox}
% Aaghh, the full addresses! A perennial source of problems. 
%    \begin{macrocode}
\newlength{\@fulladdheight}
\newlength{\@fulladddepth}
\newsavebox{\@fulladdbox}
%% \fulladdresses command modified (11 May 2000) so that we can put
%% \fulladdresses{} in the article to suppress the output of the
%% addresses
\newcommand{\@FinalFigure}{\relax}
\newcommand{\FinalFigure}[1]{\renewcommand{\@FinalFigure}{#1}}
\newcommand{\fulladdresses}[1]{%
  \renewcommand{\@rmaaArg}{#1}
  \ifx\@rmaaArg\@rmaaEmpty
     \renewcommand{\@fulladdresses}{}
  \else
     \renewcommand{\@fulladdresses}{
       \normalsize%
       \ifx\@FinalFigure\relax\else\vfill\@FinalFigure\vfill\fi
%%       \vspace*{1ex}
       %% do addresses as a list (copied from thebibliography environment) 
       \list{\null}{\leftmargin 1.5em\topsep\z@\labelwidth\z@\labelsep\z@\itemindent -1.5em
         \itemsep\z@\parsep\z@\usecounter{enumi}}
       \sloppy\clubpenalty4000\widowpenalty4000
       \sfcode`\.=1000\relax
       #1
       \endlist
       }
  \fi
  }
\newcommand{\@fulladdresses}{Full addresses go here}
%    \end{macrocode}
% \begin{macro}{\suppressfulladdresses}
% New user-level command for suppressing full addresses should you so
% desire. Currently used in the single-page-poster style. 
%    \begin{macrocode}
\newcommand{\suppressfulladdresses}{\@suppressfulladdressestrue\@needaddressesfalse}
%    \end{macrocode}
% \end{macro}
% \end{macro}
% \end{macro}
% \end{macro}
% \end{macro}
% \end{macro}

% \begin{macro}{\printdates}
% \begin{macro}{\noprintdates}
% Command to activate the printing out of the received/accepted
% dates. This can also be done using the \texttt{printdates} option,
% which is the default for \texttt{journal}, but not for \texttt{proceedings}.
%    \begin{macrocode}
\newcommand{\printdates}{\@printdatestrue}
\newcommand{\noprintdates}{\@printdatesfalse}
%    \end{macrocode}
% \end{macro}
% \end{macro}
% \begin{macro}{\@outputaddressesifnecessary}
% Command to output the addresses if we haven't already done so
% This is called by: 
% \begin{enumerate}
% \item The end of |thebibliography| environment
% \item The |\outputfulladdresses| command
% \end{enumerate}
%    \begin{macrocode}
\newcommand{\@outputaddressesifnecessary}{%
  \ClassInfo{rmaa}{Calling @outputaddressesifnecessary}
  \ifthenelse{\boolean{@needaddresses}}{%
    %% If in one column style, or in the first column of a two column
    %% style then we can try to put it on this page 
    \ifthenelse{\boolean{@firstcolumn} \or \not\boolean{@rmtwocol}}{%
      \@outputaddresses}{%
      %% otherwise give warning it will be on a separate page
      \ClassWarning{rmaa}{Full addresses will go on a separate
        page (please see User Guide for how to remedy this)}}%
    %% If we haven't managed yet, put it on its own page
    \ifthenelse{\boolean{@needaddresses}}{%
      \if@debug{\ClassWarning{rmaa}{In this silly bit!}}\fi
      \clearpage\vfill\@outputaddresses}{}%
    }{
    %% If we don't need them, we still might need to do the eoa hook 
    \ifthenelse{\boolean{@suppressfulladdresses}}{%
      \if@debug{\ClassWarning{rmaa}{XXX}}\fi
%%      \if@book\else\@endofarticlehook\null\fi
      }{}
    }%
  }  
%    \end{macrocode}
% \end{macro}

% \begin{macro}{\outputfulladdresses}
% A user interface to the above command, in case there is no biblio
% (11 May 2000) 
%    \begin{macrocode}
\newcommand{\outputfulladdresses}{\@outputaddressesifnecessary}
%    \end{macrocode}
% \end{macro}

% \begin{macro}{\adjustfinalcols}
% This command starts a new column and should be 
% used to ensure that the columns balance on the last page. It also
% outputs the full 
% author addresses to ensure that this occurs in the left hand
% column. 
% The |\par| command is so that the current line will be taken account
% of when working out |\pagetotal|
%    \begin{macrocode}
\newcommand{\adjustfinalcols}{\if@twocolumn\goodbreak\@outputaddresses\newpage\noindent\fi}
%    \end{macrocode}
% \end{macro}

% \begin{macro}{\@outputaddresses}
% Command to do the work of outputting the addresses. 
% This can be called by 
% \begin{enumerate}
% \item an explicit |\adjustfinalcols|
% \item via |\@outputaddressesifnecessary|, q.v.
% \end{enumerate}
%
%    \begin{macrocode}
\newcommand{\@outputaddresses}{%
  % Allow option of suppressing the full addresses, in which case
  % we just run the hook if we are not in a book style
  \if@suppressfulladdresses
%%     \if@book\else\@endofarticlehook\null\fi
  \else
  %% save the full addresses in a box for output later
  \sbox{\@fulladdbox}{\parbox{\textwidth}{\@fulladdresses}}
  % calculate height of address box  
  \settoheight{\@fulladdheight}{\usebox{\@fulladdbox}}
  \settodepth{\@fulladddepth}{\usebox{\@fulladdbox}}
  \addtolength{\@fulladdheight}{\@fulladddepth}
  \if@debug
  { \typeout{Address total height: \the\@fulladdheight}
    \typeout{Text height: \the\pagetotal}
    \typeout{Goal height: \the\pagegoal}
    }
  \fi
  % reuse \@fulladdheight to mean the total height of text+addresses
  \addtolength{\@fulladdheight}{\pagetotal}
%%  \addtolength{\@fulladdheight}{-\baselineskip}
  \if@debug\typeout{Adjusted text + addresses height: \the\@fulladdheight}\fi
  % check it will fit on current page
  \ifdim\@fulladdheight<\pagegoal
     %% must shift it to the left
     %% to compensate for the hanging indentation in the reference list. 
     %% NOTE: this is not foolproof!
     %% NB. This will give overfull hbox warnings, which is a bit ugly.
     \vfill\noindent\if@biblist\hspace*{-\leftmargin}\fi
     \usebox{\@fulladdbox}
     % Set this to show we've done it
     \@needaddressesfalse
     % and assume we are at the end of the article
     % so that we do the EOA hook if not in the book style
%%     \if@book\else\@endofarticlehook\fi
  \else
     % if not, just do nothing. With luck, this function will get 
     % called again on the next page
     \if@debug{\typeout{Holding over full addresses to next page}}\fi
  \fi\fi
  }
%    \end{macrocode}
% \end{macro}

% \subsubsection{Commands for editors}
% \label{sec:editor-commands}
% Some user/editor commands for doing the TOC and index in the book
% option. 
% They are defined here rather in the book-specific part since they
% need to be always available in order to avoid breaking documents. 
% \begin{macro}{\listofauthors}
% \begin{macro}{\@rawauthors}
% \begin{macro}{\rmaaTOCauthors}
% The command |\listofauthors| formats an author list for use in the
% TOC. It should contain the same names as in the |\author| command,
% but without the affiliation garbage and using |\&| instead of
% |and|. 
%    \begin{macrocode}
\newcommand{\@rawauthors}{Raw Authors}
\newcommand{\rmaaTOCauthors}{\noexpand\rmaaTOCAuthorStyle{Author list for TOC}}
\newcommand{\listofauthors}[1]{%
  \renewcommand{\@rawauthors}{#1}
  \renewcommand\rmaaTOCauthors{\noexpand\rmaaTOCAuthorStyle{#1}}}
%    \end{macrocode}
% \begin{macro}{\indexauthor}\label{indexauthor}
% The command |\indexauthor| produces an index entry for each
% author. 
% This command must come \emph{after} the
% |\title| and the |\listofauthors| commands. Also, there should be one 
% |\indexauthor| command for \emph{each} author. We also add the
% author onto the end of the |\ADS@authors| macro and construct the
% |\ADS@authorinitial| by a most cunning method (see below), but only
% for the first author. 
% 
% No, No, NO! This was all wrong. Backed out completely: 
% \begin{quote}
% Completely rewritten (19 Dec 2002) since we now use the index class, so 
% as to be able to alphabetize accented characters. A side effect of this is
% that |\index| no longer expands out the argument after the |@|, so we 
% instead add all the |\index| commands to a list (|\indexauthor@deferlist|), 
% which is executed in the |\@begofarticlehook|. 
% 
% Note the stuff with |\toks0| is copied from Exercise 20.15 of the \TeX{}book---it is
% to make sure the latin1-accented name gets passed into the idx
% file. 
% \end{quote}
% We now go back to the old way of doing things, using an optional
% argument to the \cs{indexauthor} command that has the name without
% accents. 
% \changes{3.08}{2003/01/16}{Remerged changes from 3.06 that got
% dropped from 3.07 after a synchronization disaster. Basically, the
% accentless optional argument to \cs{indexauthor}. Must add the macro
% directories to my unison set-up to stop this sort of thing from happening.}
%    \begin{macrocode}
\newcommand\@isortkey{}
\newcommand\@emptykey{empty}
\newcommand\indexauthor[2][empty]{
  \renewcommand\@isortkey{#1}
  \ifx\@isortkey\@emptykey \renewcommand\@isortkey{#2}\fi
  \index{\@isortkey @{{\rmaaIndexAuthorStyle{#2}}\hskip1em plus1em minus0.4em
      \rmaaIndexTitle\hskip1em plus2em minus0.4em\penalty-100
      \rmaaTOCauthors}}
  \indexauthorADS{#2}
}
%%\def\indexauthor@deferlist{\relax}
%%\long\def\indexauthor#1{
%%  \def\@@name{#1} \toks0=\expandafter{\@@name}
%%  \toks1=\expandafter{\indexauthor@deferlist}
%%  \protected@xdef\indexauthor@deferlist{%
%%    \the\toks1
%%    \noexpand\index{\the\toks0@{\noexpand\rmaaIndexAuthorStyle{\the\toks0}\noexpand\hskip1em plus1em minus0.4em 
%%        \rmaaIndexTitle\noexpand\hskip1em plus2em minus0.4em\noexpand\penalty-100
%%        \rmaaTOCauthors}}
%%  }
%%  \indexauthorADS{#1}
%%  }
%%\xdef\indexauthor@deferlist{}
\newcommand\indexauthorADS[1]{%
  \if@firstauthor\protected@xdef\ADS@authorinitial{\firstletter{#1}}\fi
  \protected@xdef\ADS@authors{\ADS@authors\if@firstauthor\else ; \fi
    #1}
  \@firstauthorfalse
}
%    \end{macrocode}
% Note that the first |#1| sets the key for the index, making sure it
% actually starts with  
% first letter of the surname. Meanwhile, the |@{...}| expression,
% which starts with some  
% control sequences is what is written to the |.ind| file. 
% See the makeindex man page for details. 
% \end{macro}
% \end{macro}
% \end{macro}
% \end{macro}

% \begin{macro}{\KillAllAccents}
% \begin{macro}{\firstletter}
% These are utility functions used in the above. |lett| is a counter
% indicating the current letter in the string. |\firstletter| just
% prints the first letter of its argument.
%    \begin{macrocode}
\newcommand{\KillAllAccents}{\def\'{}\def\`{}\def\~{}\def\.{}\def\i{i}}
\newcommand\firstletter[1]{\@car #1\@nil}
%    \end{macrocode}
% \end{macro}
% \end{macro}
% \begin{macro}{\SetVolume}
% \begin{macro}{\SetMSnumber}
% \begin{macro}{\SetFirstPage}
% \begin{macro}{\SetYear}
% \begin{macro}{\ReceivedDate}
% \begin{macro}{\AcceptedDate}
% These are for the editors to set the volume number, etc. 
%    \begin{macrocode}
\newcommand{\SetVolume}[1]{%
  \renewcommand{\@volume}{#1}
  \edef\ADS@volume{\PadToFourWithDots{#1}}
  \renewcommand{\@volumen}{\rmaa@VolumenStyle{volumen #1}}%
  }
\newcommand{\SetMSnumber}[1]{\renewcommand{\@MSnum}{#1}}
%    \end{macrocode}
% Note that there is no real reason to use |\SetFirstPage| since all
% the page numbering is done automatically in the book style. 
%    \begin{macrocode}
\newcommand{\SetFirstPage}[1]{\setcounter{page}{#1}}
\newcommand{\SetYear}[1]{%
  \renewcommand{\@year}{#1}
  \renewcommand{\ADS@year}{#1}
  \renewcommand{\@anyo}{\rmaa@VolumenStyle{#1}}
  }
\newcommand{\ReceivedDate}[1]{\renewcommand{\@received}{#1}}
\newcommand{\AcceptedDate}[1]{\renewcommand{\@accepted}{#1}}
%    \end{macrocode}
%
% \begin{macro}{\DOI}
% User/editor command to enter the DOI, which will be typeset as a link.
% \changes{4.0}{2019/11/02}{New user command \cs{DOI} to typeset
% DOI in the header of first page and add it to ADS info.}
%    \begin{macrocode}
\newcommand\DOI[1]{%
  \def\@DOI@url{\url{https://doi.org/10.22201/ia.01851101p.#1}}%
  \def\ADS@doi{https://doi.org/10.22201/ia.01851101p.#1}%
}%
%    \end{macrocode}
% \end{macro}
%
% \begin{macro}{\ORCIDiD}
% \changes{4.0}{2019/11/12}{Add mechanism for giving ORCID iD links}
% \begin{macro}{\ORCIDiDfull}
% \begin{macro}{\@orcid@icon}
% \begin{macro}{\@orcid@iconbig}
% Commands to add icon links to ORCID iDs of authors. For example
% \begin{quote}
% |\ORCIDiD|\marg{0000-0000-0000-0000}
% \end{quote}
% to be used after each author name in the |\author|
% command and in the |\fulladdresses|.  Alternatively,
% |\OORCIDiDfull| will typeset the icon plus the full text of the
% link, but we don't use that at present.  The icon image is stored in |ORCIDiD_iconvector.pdf|, which needs to be on the input path.  In principle, we could make EPS versions if necessary.
%    \begin{macrocode}
\newcommand\@orcid@icon{\includegraphics[height=1.8ex]{ORCIDiD_iconvector}}
\newcommand\@orcid@iconbig{\includegraphics[height=2ex]{ORCIDiD_iconvector}}
\newcommand\ORCIDiD[1]{\href{https://orcid.org/#1}{\,\@orcid@icon}}
\newcommand\ORCIDiDfull[1]{%
  \href{https://orcid.org/#1}{\@orcid@iconbig}\,\url{https://orcid.org/#1}}
%    \end{macrocode}
% \end{macro}
% \end{macro}
% \end{macro}
% \end{macro}
%
% \begin{macro}{\@volume}
% \begin{macro}{\@volumen}
% \begin{macro}{\@artpages}
% \begin{macro}{\@year}
% \begin{macro}{\@anyo}
% \begin{macro}{\@received}
% \begin{macro}{\@accepted}
% \begin{macro}{\@MSnum}
% Default values for the containers used by above commands, plus a
% few more useful ones. 
%    \begin{macrocode}
\newcommand{\@volume}{00}
\newcommand{\@volumen}{\rmaa@VolumenStyle{00}}
\newcommand{\@artpages}{%
  \oldpageref{\articlename:RMAAFirstPage}--\oldpageref{\articlename:RMAALastPage}}
\newcommand{\@year}{\the\year}
\newcommand{\@anyo}{\rmaa@VolumenStyle{2007}}
\newcommand{\@journalname}{\rmaa@JournalNameStyle{%
    \if@proceedings\RMAASClong\else\RMAAlong\fi}}
\newcommand{\@received}{}
\newcommand{\@accepted}{}
\newcommand{\@MSnum}{0000}
\newcommand{\@DOI@prefix}{https://doi.org/10.22201/ia.01851101p}
\def\@DOI@url{\url{https://doi.org/10.22201/ia.01851101p.2099.00.00.00}}
%    \end{macrocode}
% \end{macro}
% \end{macro}
% \end{macro}
% \end{macro}
% \end{macro}
% \end{macro}
% \end{macro}
% \end{macro}
% \end{macro}
% \end{macro}
% \end{macro}
% \end{macro}
% \end{macro}
% \end{macro}
% 
% The following commands are most useful in the \texttt{book}
% style. However, the |\SetConfTitle| command can now be used in
% individual articles to set the header for the \texttt{preprint}
% style. 
%    \begin{macrocode}
\newcommand{\@reunionplace}{Lugar de la reuni\'on}
\newcommand{\@reuniondate}{Fecha de la reuni\'on}
\newcommand{\@conftitle}{Conference Title}
\newcommand{\SetReunion}[2]{\renewcommand{\@reunionplace}{#1}\renewcommand{\@reuniondate}{#2}}
\newcommand{\SetConfTitle}[1]{\renewcommand{\@conftitle}{#1}}
\newcommand{\@editors}{Juan Ram\'on de la Fuente and ``El Mosh''}
\newcommand{\SetEditors}[1]{\renewcommand{\@editors}{#1}}
\newcommand{\@editores}{Juan Ram\'on de la Fuente y ``El Mosh''}
\newcommand{\SetEditores}[1]{\renewcommand{\@editores}{#1}}
\newcommand{\@mes}{mes de la publicaci\'on}
\newcommand{\SetMes}[1]{\renewcommand{\@mes}{\rmaa@VolumenStyle{#1}}}
%    \end{macrocode}

% \subsection{Notes for Author Proofs and Communication between Editors}
% \label{sec:notes}
% We introduce a command \cs{Note} that should be used to describe
% any change made to the manuscript. These are a bit tricky to get
% right. In general we use \cs{marginpar} but this only works in outer
% paragraph mode, so we just bung the box where we are if we are in a
% figure caption or something. If we are in the preamble, then we just
% use a \marginpar as part of the \maketitle command.
%    \begin{macrocode}
\newif\if@arenotes
\@arenotesfalse
\newcounter{notecounter}
\setcounter{notecounter}{0}
\newcommand\NotesList{}
\newcommand\Note[1]{%
  \if@proof{%
    \@arenotestrue%
    \stepcounter{notecounter}%
    \if@inpreamble\else{%
      \NoteMark{\arabic{notecounter}}%
    }\fi
    \protected@xdef\NotesList{\NotesList 
      \protect\item #1}%
  }\fi%
}
\newcommand\Changed[2]{Changed ``#1'' to ``#2''.}
\newcommand\NoteMark[1]{%
  \ifinner
  \raisebox{1ex}[0pt][0pt]{\makebox[0pt][l]{\framebox{\footnotesize\bfseries
        N#1}}}%
  \else
  \marginpar{\framebox{\footnotesize\bfseries N#1}}%
  \fi
}
\newcommand\RefNote[1]{\if@proof\NoteMark{\ref{#1}}\relax\fi}
\newcommand\NoteLabel{\protect\label}
\newcommand\OutputNotes{%
  \bgroup
  \onecolumn
  \section*{EDITORIAL NOTES}
  \begin{enumerate}
    \NotesList
  \end{enumerate}
  \egroup
}
\DeclareRobustCommand\cs[1]{\texttt{\char`\\#1}}

%    \end{macrocode}
% \subsection{Bibliographic Records for ADS}
% \label{sec:ADS}
% \changes{2.2beta11}{2001/11/27}{Just about finished all the ADS
% stuff now. Except for the references, which I have given up on for
% the time being.}
% \begin{macro}{\ADS@BibCode}
% \begin{macro}{\ADS@year}
% \begin{macro}{\ADS@journal}
% \begin{macro}{\ADS@volume}
% \begin{macro}{\ADS@firstpage}
% \begin{macro}{\ADS@authorinitial}
% We use the |\ADS@BibCode| command to hold the bibliographic code for
% ADS. It is 
% used in the |%R| 
% field of the bibliographic record for each article. 
% The format is 
% \begin{quote}
% \texttt{YYYYJJJJJVVVVMPPPPA}  
% \end{quote}
% where \texttt{YYYY} is the year, \texttt{JJJJJ} is the journal code,
% which is \texttt{RMxAA} or \texttt{RMxAC}, \texttt{VVVV} is the
% volume number (right-justified), \texttt{M} is not used by us,
% \texttt{PPPP} is the first page number (right-justified), and
% \texttt{A} is the initial of the first author's surname. The volume
% and page fields are padded with periods (`\texttt{.}'), which causes
% us some fun (this is done with the macro |\PadToFourWithDots| described
% below) . 
% 
% All these individual bits are held in the commands |\ADS@year|,
% |\ADS@journal|, |\ADS@volume|, |\ADS@firstpage|,
% |\ADS@authorinitial|. 
%    \begin{macrocode}
\newcommand{\SetADSBibCode}[1]{\renewcommand{\ADS@BibCode}{#1}}
\newcommand{\ADS@BibCode}
{\ADS@year\ADS@journal\ADS@volume.\PadToFourWithDots{\ADS@firstpage}\ADS@authorinitial}
\newcommand{\ADS@year}{2001} 
\newcommand{\ADS@journal}{\if@proceedings RMxAC\else RMxAA\fi}
\newcommand{\ADS@volume}{....} 
%    \end{macrocode}
% These two are just set to default values here, |\ADS@firstpage| is set
% in the |document| environment for included articles, in the
% |frontsection| environment, |\tableofcontents| command,  |index| and
% |abstracts| environments. Also by |\AtBeginDocument| for stand-alone
% articles. 
% 
% God knows how we are going to do |\ADS@authorinitial|!
%    \begin{macrocode}
\newcommand{\ADS@firstpage}{....}
\newcommand{\ADS@authorinitial}{A}
%    \end{macrocode}
% \end{macro}
% \end{macro}
% \end{macro}
% \end{macro}
% \end{macro}
% \end{macro}
% We also have a whole slew of other fields that can be set. These are
% all listed in Table~\ref{tab:ADS}. Most of them still have to be
% defined. 
% \begin{table}[!ht]\raggedleft
% \caption{ADS Bibliographic Records}\label{tab:ADS}
% \begin{tabular}{lp{8em}lp{8em}}\hline
% Code & Description & Macro Name & Comment \\ \hline
% |%R| & Bibliographic Code & |\ADS@BibCode| & required \\
% |%A| & Author list & |\ADS@authors| & required, generate from |\indexauthor| \\
% |%F| & Author affiliation & |\ADS@affil| & format: AA(\dots)
% AB(\dots) \dots \\
% |%J| & Journal Name & |\ADS@journalname| &  \\
% |%D| & Publication Date & |\ADS@fulldate| &  \\
% |%L| & Last Page & |\ADS@lastpage| & generate from label \\
% |%T| & Title & |\ADS@title| & required, generate from |\title| \\
% |%C| & Abstract Copyright & |\ADS@copyright| & set to IA-UNAM \\
% |%U| & URL for electronic document & |\ADS@url| & set to PDF file \\
% |%d| & DOI for article & |\ADS@doi| & set by |\DOI| \\
% |%K| & Keywords & |\ADS@keywords| & set equal to |\@keywords| \\
% |%A| & Abstract Text & |\ADS@abstract| & set equal to |\@abstract| \\
% |%R| & References & |\ADS@references| & generate from the
% bibliography  \\ \hline
% \end{tabular}
% \end{table}
% A lot of these macros are set by other commands, as mentioned in the
% Table, but some need to be set  explicitly here. 
% 
% \begin{macro}{\ADS@doi}
% \changes{4.0}{2019/11/08}{Save DOI for ADS}
% The user/editor command \cs{DOI} is used to set this. 
%    \begin{macrocode}
\newcommand{\ADS@doi}{https://doi.org/10.22201/ia.01851101p.0000.00.00.00}
%    \end{macrocode}
% \end{macro}
% 
% \begin{macro}{\ADS@journalname}
% This
% does not contain any accents at the moment so as not to cause
% problems---this maybe should be changed though. 
%    \begin{macrocode}
\newcommand{\ADS@journalname}{%
  \if@proceedings\@conftitle\space (Eds. \@editors) \fi%
  Revista Mexicana de Astronom\'\i{}a y Astrof\'\i{}sica%
  \if@proceedings\space (Serie de Conferencias)\fi%
  \space Vol. \@volume, pp. \ADS@firstpage--\ADS@lastpage\space (\@year)%
  \space (http://www.astroscu.unam.mx/\~{}rmaa/)
  }
%    \end{macrocode}
% \end{macro}
% \begin{macro}{\ADS@keywords}
% \begin{macro}{\ADS@abstract}
% \begin{macro}{\ADS@title}
% |\ADS@abstract| is straightforward, since we already have the text
% saved as |\@abstract|. We could do the same with |\ADS@keywords| but
% to avoid unwanted formatting we explicitly set it in |\keywords|
% (p.~\pageref{keywords}). |\ADS@title| is just set to |\@rawtitle|,
% which we were already using for another purpose. 
%    \begin{macrocode}
\let\ADS@title=\@rawtitle
\newcommand{\ADS@keywords}{No keywords specified}
\newcommand{\ADS@abstract}{No abstract specified}
%    \end{macrocode}
% \end{macro}\end{macro}\end{macro}
% \begin{macro}{\ADS@authors}
% |\ADS@authors| wants to be in ``Other, A. N.'' format, but with
% semi-colons between authors. The most natural way to do this seems to
% be with |\indexauthor| (p.~\pageref{indexauthor}). We append each
% author to this macro, which means we must explicitly clear it for
% each new article (done in |\article|). .   
%    \begin{macrocode}
\newcommand{\ADS@authors}{}
%    \end{macrocode}
% \end{macro}
% \begin{macro}{\ADS@fulldate}
% \begin{macro}{\ADS@lastpage}
% \begin{macro}{\ADS@copyright}
% \begin{macro}{\ADS@website}
% \begin{macro}{\ADS@url}
% \begin{macro}{\SetWebsite}
% \begin{macro}{\SetMonthNumeric}
% Other sundry data, together with user commands for bits that can be
% set by the Editors. The |\ADS@url| is constructed from the website
% (currently |http://www.astrosmo.unam.mx/rmaa/|,
% the journal descriptor (e.g., |RMxAC..10/|), followed by the
% path to the PDF file (e.g., |PDF/RMxAC..10_HEW1.pdf|). 
% 
%    \begin{macrocode}
\newcommand{\ADS@website}{http://www.astrosmo.unam.mx/rmaa/}
\newcommand{\ADS@url}
{\ADS@website\ADS@journal\ADS@volume\ADS@number/PDF/\ADS@journal\ADS@volume\ADS@number\_\articlename.pdf}
\newcommand{\ADS@number}{}
\newcommand{\ADS@fulldate}{00/\ADS@year} 
\newcommand{\ADS@lastpage}{00} 
\newcommand{\ADS@copyright}{(C) \ADS@year; Instituto de Astronom\'\i{}a,
  Universidad Nacional Aut\'onoma de M\'exico} 
\newcommand{\SetNumber}[1]{\renewcommand{\ADS@number}{-#1}}
\newcommand{\SetWebsite}[1]{\renewcommand{\ADS@website}{#1}}
\newcommand{\SetMonthNumeric}[1]{\renewcommand{\ADS@fulldate}{#1/\ADS@year}}
%    \end{macrocode}
% \end{macro}\end{macro}\end{macro}\end{macro}\end{macro}\end{macro}\end{macro}
% \begin{macro}{\ADS@affil}
% \begin{macro}{\ADS@references}
% These two are somewhat more problematical. Affiliation currently
% only works for the first author (see \dots) pending redesign of the
% user interface. References are now done by wrapping the
% |thebibliography| environment to open another input stream on the
% same file and search through for the |\bibitems|. 
%    \begin{macrocode}
\newcommand{\ADS@references}{}
\newcommand{\ADS@affil}{}
%    \end{macrocode}
% \end{macro}\end{macro}
% 
% Define two macros that expand to the strings
% |\begin{thebibliography}| and |\end{thebibliography}|. Note that the
% idea for this and much else is inspired by the code for the docstrip
% package. Also define |\StartBibX| in case authors use the more
% standard |\begin{thebibliography}{}| syntax. 
%    \begin{macrocode}
\edef\EndBib{\expandafter\@gobble\string\\end\string{thebibliography\string}}
\edef\StartBib{\expandafter\@gobble\string\\begin\string{thebibliography\string}}
\edef\StartBibX{\expandafter\@gobble\string\\begin\string{thebibliography\string}\string{\string}}
%    \end{macrocode}
% Define two new logical variables. |\ifAnotherline| is true if we
% want to read another line from the input stream. It is set to false
% when we encounter the |\end{thebibliography}|. |\ifInBib| is set to
% true when we frind the |\begin{thebibliography}| and tells us to
% start saving the following lines. 
%    \begin{macrocode}
\newif\ifAnotherLine
\newif\ifInBib
%    \end{macrocode}
% \begin{macro}{makeOther}
% We need a macro to change the catcodes of lots of things to other. 
%    \begin{macrocode}
\def\makeOther#1{\catcode`#1=12\relax}
%    \end{macrocode}
% \end{macro}
% \begin{macro}{\RereadBiblio}
% Define a macro to do all the fiddly stuff. This is used at the
% beginning of |thebibliography| environment (Sec~\ref{sec:biblio}).
%    \begin{macrocode}
\newread\tf@biblio
\newcommand{\RereadBiblio}{%
%    \end{macrocode}
% Open the current file again for reading
%    \begin{macrocode}
  \begingroup
  \immediate\openin\tf@biblio=\articlename
%    \end{macrocode}
% Everything in a group to be safe. Initialize logicals. Black magic
% to change
% catcodes and avoid blank lines being interpreted as |\par|. 
%    \begin{macrocode}
  \AnotherLinetrue\InBibfalse
  \makeOther\ \makeOther\\\makeOther\$%
  \makeOther\#\makeOther\^\makeOther\^^K%
  \makeOther\_\makeOther\^^A\makeOther\%%
  \makeOther\~\makeOther\{\makeOther\}\makeOther\&%
  \endlinechar-1\relax
%    \end{macrocode}
% Start looping through the file. Read each line to the macro
% |\Line|. If we have got to EOF or end of biblio then set
% |\AnotherLinefalse|.
%    \begin{macrocode}
  \loop
  \read\tf@biblio to \Line
  \ifx\Line\EndBib\AnotherLinefalse\InBibfalse\fi
  \ifeof\tf@biblio \AnotherLinefalse\fi
%    \end{macrocode}
% If we are in the biblio and the line is not blank then append it to
% the |\ADS@references| macro together with a newline. 
%    \begin{macrocode}
  \ifInBib
  \ifx\Line\empty\else\xdef\ADS@references{\ADS@references\Line^^J}\fi
  \fi
%    \end{macrocode}
% Check if the line we just saw was the start of biblio, then repeat
% the loop if necessary. 
%    \begin{macrocode}
  \ifx\Line\StartBib\InBibtrue\fi
  \ifx\Line\StartBibX\InBibtrue\fi
  \ifAnotherLine\repeat
  \immediate\closein\tf@biblio
  \endgroup
  }
%    \end{macrocode}
% \end{macro}
% 
% \begin{macro}{\PadToFourWithDots}
% Utility function for left-padding a number to 4 characters with
% periods. Used in 
% constructing |\ADS@BibCode|. This is a very simple algorithm and
% \emph{requires} that the argument actually be a number---if it is
% not then dire consequences will result.  
%    \begin{macrocode}
\newcommand{\PadToFourWithDots}[1]{%
  \ifnum #1 > 999 #1%
  \else\ifnum #1 > 99 .#1%
  \else\ifnum #1 > 9 ..#1%
  \else ...#1%
  \fi\fi\fi}
%%  \ifnum \count0 > 999 #1%
%%  \else\ifnum \count0 > 99 .#1%
%%  \else\ifnum \count0 > 9 ..#1%
%%  \else ...#1%
%%  \fi\fi\fi}
%    \end{macrocode}
% \end{macro}
%
% \begin{macro}{\AuthorADS}
% \begin{macro}{\ADS@affiliations}
% New mechanism for setting the affiliations for ADS. The author labels (AA, AB, \dots) are generated using two nested counters.  We use the |\robustify| command from \texttt{etoolbox} so that accented characters in the affiliation addresses do not blow up.  The list of accents may need to be expanded. 
%    \begin{macrocode}
\newcounter{AUTHOR}
\newcounter{author}[AUTHOR]
\stepcounter{AUTHOR}
\newcommand\ADS@affiliations{}
\robustify{\'}\robustify{\i}\robustify{\~}\robustify{\`}
\def\AuthorADS#1#2{%
  \begingroup
  \def\ORCIDiD##1{<ID system="ORCID">##1</ID>}
  \def\Email##1{<EMAIL>##1</EMAIL>}
  \def\Affil##1{##1}
  \stepcounter{author}
  \ifthenelse{\value{author} > 26}{\stepcounter{AUTHOR}\stepcounter{author}}{}
  \ifthenelse{\equal{\ADS@affiliations}{\@rmaaEmpty}}{%
    \xdef\ADS@affiliations{\Alph{AUTHOR}\Alph{author}(#2)}
  }{%
    \xdef\ADS@affiliations{\ADS@affiliations, \Alph{AUTHOR}\Alph{author}(#2)}
  }%
  \endgroup
}
%    \end{macrocode}
% \end{macro}
% \end{macro}
% %
% \begin{macro}{\ADS@dumpdata}
% \changes{4.0}{2019/11/13}{Write all author affiliations for ADS}
% \begin{macro}{\ADS@write}
% \begin{macro}{\ADS@openfile}
% |\ADS@dumpdata| is the macro that actually does the job of writing out all the
% data to an auxiliary file, using |\ADS@write| to write each
% field. The file is set up with |\ADS@openfile|, which currently is
% done in the TOC of a book.  The new multi-author affiliations are used if |\ADS@affiliations| has been set, otherwise it falls back to the single-author version |\ADS@affil|. 
%    \begin{macrocode}
\newcommand{\ADSfield}[2]{\relax}
\newcommand{\ADSvfield}[2]{\relax}
%%\renewcommand{\ADSfield}[12]{\%#1 #2}
\newcommand{\ADS@dumpdata}
{ 
  \ADS@write{\string\begin{ADSarticle}}
  \ADS@WRITE{\string\ADSfield{R}{\ADS@BibCode}}
  \ADS@write{\string\ADSfield{A}{\ADS@authors}}
  \ifthenelse{\equal{\ADS@affiliations}{\@rmaaEmpty}}{%
    {\ADS@sanitize\ADS@write{\string\ADSfield{F}{AA(\ADS@affil)}}}
  }{%
    {\ADS@sanitize\ADS@write{\string\ADSfield{F}{\ADS@affiliations}}}
  }
  \ADS@write{\string\ADSfield{J}{\ADS@journalname}}
  \ADS@write{\string\ADSfield{D}{\ADS@fulldate}}
  \ADS@write{\string\ADSfield{L}{\ADS@lastpage}}
  {\ADS@sanitize\ADS@write{\string\ADSfield{T}{\ADS@title}}}
  \ADS@write{\string\ADSfield{C}{\ADS@copyright}}
  \ADS@write{\string\ADSfield{d}{\ADS@doi}}
  \ADS@write{\string\ADSfield{U}{\ADS@url}}
  \ADS@write{\string\ADSfield{K}{\ADS@keywords}}
  {\ADS@sanitize\ADS@write{\string\ADSfield{B}{\ADS@abstract}}}
  \ADS@write{\string\ADSfield{Z}{\ADS@references}}
  \ADS@write{\string\end{ADSarticle}}
  }
%%\newcommand{\ADS@write}[1]{\if@book\addtocontents{ads}{#1}\fi}
%%\newcommand{\ADS@openfile}{\@starttoc{ads}}
\newcommand{\ADS@write}[1]{\set@display@protect\immediate\write\tf@ads{#1}}
\newcommand{\ADS@WRITE}[1]{\immediate\write\tf@ads{#1}}
\newcommand{\ADS@openfile}{%
  \newwrite\tf@ads%
  \immediate\openout\tf@ads\jobname.ads\relax
  \ADS@WRITE{\ADS@preamble}}
\newcommand{\ADS@closefile}{\ADS@WRITE{\ADS@postamble}}
\def\ADS@preamble{%
  %%  \string\documentclass{article}
  %%  \string\begin{document}
  }
\def\ADS@postamble{%
  %%  \string\end{document}
  \string%\string% Local Variables:^J 
  \string%\string% mode: LaTeX^J
  \string%\string% mode: Font-lock^J
  \string%\string% End:^J
  }
%    \end{macrocode}
% \end{macro}
% \end{macro}
% \end{macro}
% 
% \subsection{Automatic calculation of the page range}
% \label{sec:page-range}
% Another issue that has cost me a lot of grief. Still, it finally
% seems to be sorted now. All hail Jeff Goldberg. 
%
% Cunning trick to get the page numbers from |\pageref|.  Put a label
% on the first and last page of the document using |\AtBeginDocument|
% and |\AtEndDocument|. .  This way won't work for the book option.
% In that case, we do it by defining |\@endofarticlehook| and
% |\@begofarticlehook| as parallels to |\AtEndDocument| and
% |\AtBeginDocument|.  Note that we have to defer adding
% |\@begofarticlehook| to |\AtBeginDocument| to make sure it comes
% after the redefinitions of |\ref| and |\label|.
%    \begin{macrocode}
\if@book\AtEndDocument{\ADS@closefile}\else
 %\AtBeginDocument{\@begofarticlehook}
\AtEndDocument{\@endofarticlehook}
\fi
%    \end{macrocode}
% \begin{macro}{\@begofarticlehook}
% Since we are getting more and more things that need to be done at
% the beginning of every article, frontsection, etc we define a sister
% hook to 
% |\@endofarticlehook|. This hook does the following: 
% \begin{enumerate}
% \item Labels the first page
% \item Sets |\ADS@firstpage|
% \item runs and then clears |\indexauthor@deferlist|
% \item resets the keyword counter
% \end{enumerate}
% It is executed in the following circumstances: 
% \begin{enumerate}
% \item At |\begin{document}| for standalone files
% \item At |\begin{document}| for each included file
% \item At beginning of front sections 
% \item At beginning of TOC
% \item At beginning of Author Index
% \item At beginning of abstracts
% \end{enumerate}
% For some reason, this is not always within the right group, so we
% use |\xdef| just to be sure. 
%    \begin{macrocode}
\newcommand{\@begofarticlehook}{%
  \if@frontmatter\frontlabel{RMAAFirstPage}\else\label{RMAAFirstPage}\fi%
  \xdef\ADS@firstpage{\@arabic{\c@page}}%
  \if@book\else\makeindex\fi % needed so default index is defined in stand-alone articles
%%  \indexauthor@deferlist
%%  \xdef\indexauthor@deferlist{}
  \@inpreamblefalse
  \setcounter{kwdcounter}{0}
  }
%    \end{macrocode}
% \end{macro}
% \begin{macro}{\@endofarticlehook}
% This is a hook to be run at the end of each article
% currently, it just labels the last page and writes out all the ADS
% data to an auxiliary file.  
% It should not produce any printed output. 
% It is modelled on the |\AtEndDocument| command, to which it reduces
% in the case of a single article. 
%    \begin{macrocode}
\newcommand{\@endofarticlehook}{%
  \if@debug{\ClassWarning{rmaa}{***Doing Hook***}}\fi
  \if@debug{\ClassWarning{rmaa}{ADSabstract is: \ADS@abstract}}\fi
  \clearpage\rmaa@lastpagelabel\if@book\ADS@dumpdata\fi
  \ifx\NotesList\@rmaaEmpty\else\OutputNotes\fi
  \if@debug{\ClassWarning{rmaa}{Notes list is: \NotesList}}\fi
}
%    \end{macrocode}
% \begin{macro}{\rmaa@lastpagelabel}
% New version of last page label command copied from Jeff Goldberg's
% lastpage.sty package - we can't just use |\label| here in case the
% output routine for the last page has already been called. Instead,
% we first explicitly clear the page, step back the page counter,
% explicitly write to the .aux file, then restep the page counter
% again.
%
% But this will conflict with the hyperref package, which does labels differently.  Compensate for that by using the five-argument form.
% 
% Also set the |\ADS@LastPage| macro to the appropriate number. 
%    \begin{macrocode}
\newcommand{\rmaa@lastpagelabel}{%
  \addtocounter{page}{-1}%
  \immediate\write\@auxout{\string
    \newlabel{\articlename:RMAALastPage}{{}{\thepage}{}{}{}}}%
  \edef\ADS@lastpage{\thepage}%
  \addtocounter{page}{1}}
%    \end{macrocode}
% \end{macro}
% \end{macro}

% \subsection{Deal with explicit change in number of columns}
% \label{sec:columns}
% Sometimes, when formatting a book, we need to switch between
% one-column articles and two-column articles. For this to work, we
% need to wrap the standard \LaTeX{} commands to set the |\@rmtwocol|
% flag appropriately and also reset the package default values of all
% scalable lengths. .  
%    \begin{macrocode}
\let\latexonecolumn=\onecolumn
\let\latextwocolumn=\twocolumn
\renewcommand{\onecolumn}{\@rmtwocolfalse\@package@default@lengths\latexonecolumn}
\renewcommand{\twocolumn}{\@rmtwocoltrue\@package@default@lengths\latextwocolumn}
%    \end{macrocode}
% 
% \subsection{The actual formatting of the title}
% \label{sec:title-format}

% Note, we have to redefine the |\maketitle| command to use the
% numbered footnotes. This means we needn't use the atbeginend stuff. 
%    \begin{macrocode}
\renewcommand{\maketitle}{\par
  \begingroup
%    \end{macrocode}
% Don't switch to marked footnotes. 
% This is consistent with existing style but I'm not sure I like it
%    \begin{macrocode}
  %%  \renewcommand\thefootnote{\fnsymbol{footnote}}%
%%  \def\@makefnmark{\hbox to\z@{$\m@th^{\@thefnmark}$\hss}}%
  \def\@makefnmark{\textsuperscript{\@thefnmark}\hspace{-0.2em}}%
  \long\def\@makefntext##1{\parindent 1em\noindent
    \hbox to1.8em{\hss$\m@th^{\@thefnmark}$}##1}%
  \if@rmtwocol
  { \if@debug{\ClassWarning{rmaa}{TWO COLUMN}}\fi
    \twocolumn[\@maketitle]%
    }
  \else
  { \if@debug{\ClassWarning{rmaa}{ONE COLUMN}}\fi
    \newpage
    \global\@topnum\z@   % Prevents figures from going at top of page.
    \if@manuscript \@MSmaketitle \else \@maketitle \fi
    }
  \fi
  \if@manuscript\else\thispagestyle{rmaafirstpage}\@thanks\fi
  \ifx\NotesList\@rmaaEmpty\else{%
    \NoteMark{1--\arabic{notecounter}}\vspace*{-2\baselineskip}%
  }\fi%
  \endgroup
%    \end{macrocode}
% Don't reset footnote counter (not sure whether this is good thing
% or not).
%    \begin{macrocode}
  %%  \setcounter{footnote}{0}%
  \let\thanks\relax\let\affil\relax\let\altaffil\relax
  \let\maketitle\relax\let\@maketitle\relax
  \gdef\@thanks{}\gdef\@author{}\gdef\@title{}}
%    \end{macrocode}

% \begin{macro}{\@kwhwidth}
% \begin{macro}{\@kwtwidth}
% Redefine \cs{@maketitle} to include the abstract, resumen and
% keywords. Note that in the manuscripty style we use
% \cs{@MSmaketitle} instead.  
% \changes{3.16}{2003/04/07}{
%  Added extra vertical space between authors and accepted date.}
% \changes{3.26}{2011/08/15}{Add extra label to review articles}
%    \begin{macrocode}
\newlength{\@kwhwidth}\newlength{\@kwtwidth}
\newcommand\reviewname{Review}
\renewcommand{\@maketitle}{%
  \newpage
  \null
  \if@reviewarticle
  \smallskip
  {\large\bfseries\itshape\reviewname}%
  \bigskip
  \else
  \vskip\@before@title@skip%
  \fi
  \begin{center}%
    {\@title \par}%
    \vskip\@before@author@skip%
    \parbox{\rmaa@AuthorWidthFrac\textwidth}{\centering
      \lineskip\@between@author@skip%
      \@author
      \par}
    %% this should be replaced with `Submitted to...' for preprints
    \if@printdates
    \vskip\@before@dates@skip%
    \if@preprint
    \rmaa@ReceivedStyle{Submitted to RMxAA \@received}
    \else
    \rmaa@ReceivedStyle{Received \@received; accepted \@accepted}
    \fi
    \fi
    \if@manuscript
    \rmaa@ReceivedStyle{Submitted to RMxAA \@received}
    \fi
    \if@guide
    \rmaa@ReceivedStyle{Version \rmaa@version, \rmaa@date}
    \fi
    \vskip\@before@abstract@skip%
    \if@abstractpresent
    \parbox{\rmaa@AbstractWidthFrac\textwidth}{%
      \setlength{\parindent}{0.9cm}%
      \centerline{\rmaa@AbstractHeaderStyle{Abstract}}\par%
      \vskip\@after@abstracthead@skip%
      \if@newstyle\noindent\fi\@abstract
      \vskip\@between@abstract@skip%
      \centerline{\rmaa@AbstractHeaderStyle{Resumen}}\par%
      \vskip\@after@abstracthead@skip%
      \if@newstyle\noindent\fi\@resumen
      \if@setresumen\else\@abstract\fi % append english abstract if
                                       % spanish one not set
      \if@guide\else            %do not use keywords for guide
      \vskip\@before@kwd@skip%
      % Format the Key Words with hanging indentation:
      % They are put in a \parbox with width equal to the line width
      % minus the width of the header. 
      { % Widths of keyword
        \settowidth{\@kwhwidth}{\@keywordheader~}   % header and
        \setlength{\@kwtwidth}{\linewidth}          % keyword text
        \addtolength{\@kwtwidth}{-\@kwhwidth}       %
        \noindent\@keywordheader~\parbox[t]{\@kwtwidth}{%
          \rmaa@KeywordStyle{\boldmath\uppercase{\@keywords}}}}%
      \fi
      }
      \fi
  \end{center}%
  \vspace*{\afterkwdskip}
  \par
  }
\newlength{\afterkwdskip}\setlength{\afterkwdskip}{\z@}
%    \end{macrocode}
% \end{macro}
% \end{macro}
%
%
% \begin{macro}{\reviewarticle}
% \begin{macro}{\reviewtoc}
% \begin{macro}{\reviewlof}
%   User command to specify a review article. Must be used in
%   preamble. Main effect is to place the text ``Review'' before the
%   article title.  Also add command for setting a table of contents
%   and a list of figures at the beginning of the article.
%    \begin{macrocode}
\newcommand{\reviewarticle}{\@reviewarticletrue}
\newcommand\reviewtoc{%
  \section*{CONTENTS}%
  \@starttoc{rtc}%
}
\newcommand\reviewlof{%
  \section*{LIST OF FIGURES}%
  \@starttoc{lof}%
}
%    \end{macrocode}
% \end{macro}
% \end{macro}
% \end{macro}
% \begin{macro}{\l@section}
%    \begin{macrocode}
\renewcommand*\l@section[2]{%
  \ifnum \c@tocdepth >\z@
    \addpenalty\@secpenalty
    \addvspace{0.8em \@plus\p@}%
    \setlength\@tempdima{1.5em}%
    \begingroup
      \parindent \z@ \rightskip \@pnumwidth
      \parfillskip -\@pnumwidth
      \leavevmode %\bfseries
      \advance\leftskip\@tempdima
      \hskip -\leftskip
      #1\nobreak\hfil \nobreak\hb@xt@\@pnumwidth{\hss #2}\par
    \endgroup
  \fi}
%    \end{macrocode}
% \end{macro}
%
% \begin{macro}{\@MSmaketitle}
% A special version of |\@maketitle| for the \texttt{manuscript}
% option. This can deal with a multi-page abstract/resumen. Not very
% well, mind. 
%    \begin{macrocode}
\newcommand{\@MSmaketitle}{%
  \if@debug\ClassInfo{rmaa}{Using special MS maketitle command}\fi
  \newpage
  {
    \thispagestyle{rmaafirstpage}
%     \setlength{\parindent}{0pt}
    \null\vskip 2ex
    \begin{center}
      \@title\par
      \@author\par
      \rmaa@ReceivedStyle{%
        \ifx\@received\@rmaaEmpty
          Draft version: \@date
        \else
          Received: \@received
          \ifx\@accepted\@rmaaEmpty
          \else
            ; Accepted: \@accepted
          \fi
        \fi}\par
    \end{center}
    \@thanks
    \centerline{\rmaa@AbstractHeaderStyle{Resumen}}
    \noindent\@resumen\par\vskip 2ex
    \centerline{\rmaa@AbstractHeaderStyle{Abstract}}
    \noindent\@abstract\par\vskip 2ex
    \noindent\@keywordheader~\rmaa@KeywordStyle{\@keywords}\par
%     \@keywordheader~\rmaa@KeywordStyle{\uppercase{\@keywords}}\par
  }%
  }
%    \end{macrocode}
% \end{macro}

% \subsection{The sectioning commands}
% \label{sec:section-commands}

% \begin{macro}{\section}
% \begin{macro}{\subsection}
% \begin{macro}{\subsubsection}
% \begin{macro}{\paragraph}
% \begin{macro}{\subparagraph}
% Redefine all the standard commands to use these lengths. 
%    \begin{macrocode}
\renewcommand\section{\@startsection {section}{1}{\z@}%
                                   {\abovesectionskip}%
                                   {\belowsectionskip}%
                                   %% how do we get upper case here
                                   {\centering\reset@font\normalsize}}
\renewcommand\subsection{\@startsection{subsection}{2}{\z@}%
                                     {\abovesectionskip}%
                                     {\belowsectionskip}%
                                     {\centering\reset@font\normalsize\itshape}}
\renewcommand\subsubsection{\@startsection{subsubsection}{3}{\z@}%
                                     {\abovesectionskip}%
                                     {\belowsectionskip}%
                                     {\centering\reset@font\normalsize\itshape}}
\renewcommand\paragraph{\@startsection{paragraph}{4}{\z@}%
                                    {\aboveparaskip}%
                                    {\belowparaskip}%
                                    {\reset@font\normalsize\bfseries}}
\renewcommand\subparagraph{\@startsection{subparagraph}{5}{\parindent}%
                                      {\aboveparaskip}%
                                      {\belowparaskip}%
                                      {\reset@font\normalsize\bfseries}}
%    \end{macrocode}
% \end{macro}
% \end{macro}
% \end{macro}
% \end{macro}
% \end{macro}

% \begin{macro}{secnumdepth}
% Make sure that we number the subsubsections---this is done by
% default in the |article| class but not the |book| class, which is
% why we must set it explicitly here.
%    \begin{macrocode}
\setcounter{secnumdepth}{3}
%    \end{macrocode}
% \end{macro}
% \changes{2.2beta10}{2001/11/07}{Fixed \cs{secnumdepth} for \cs{book}
% style.}

% \begin{macro}{\thesubsection}
% \begin{macro}{\thesubsubsection}
% For subsection and subsubsection, we want the numbers to be roman
% but the title to be italic. This seems the easiest way to go about it.
%    \begin{macrocode}
\renewcommand{\thesubsection}{{\upshape\thesection.\@arabic\c@subsection}}
\renewcommand{\thesubsubsection}{\thesubsection{\upshape.\@arabic\c@subsubsection}}
%    \end{macrocode}
% \end{macro}
% \end{macro}

% \begin{macro}{\@seccntformat}
% Add ``.'' to the number here so we can remove it from |\thesection|,
% etc. Reduce gap between number and title WJH 27 Jan 1995
%    \begin{macrocode}
\def\@seccntformat#1{\csname the#1\endcsname .\hskip 0.4em}
%    \end{macrocode}
% \end{macro}

% \begin{macro}{\@sect}
% I didn't really want to mess with the |\@sect| command, but it was
% the only way I could manage to force upper case. Other than the
% commented parts, it is identical to the standard version. 
%    \begin{macrocode}
\def\@sect#1#2#3#4#5#6[#7]#8{
  \ifnum #2>\c@secnumdepth
     \let\@svsec\@empty
  \else
     \refstepcounter{#1}%
     \let\@@protect\protect
     \def\protect{\noexpand\protect\noexpand}%
     \edef\@svsec{\@seccntformat{#1}}%
     \let\protect\@@protect
  \fi
%    \end{macrocode}
% This is the bit I changed (27 Jan 1995). Use upper case for
% |\section|, use mixed case for the rest. 
%    \begin{macrocode}
  \ifnum #2=1                   
     \def\sec@upcase##1{\uppercase{##1}}%
     \def\sec@shape{\upshape}%
  \else                         
     \def\sec@upcase##1{\relax{##1}}%
     \def\sec@shape{\itshape}%
  \fi
  \@tempskipa #5\relax
  \ifdim \@tempskipa>\z@
     \begingroup #6\relax
          \@hangfrom{\hskip #3\relax\@svsec}%
                    {\interlinepenalty \@M \sec@upcase{#8}\par}%
     \endgroup
     \csname #1mark\endcsname{#7}%
%    \end{macrocode}
% Only write TOC entries if doing the Guide style. 
%    \begin{macrocode}
     \if@guide 
        \addcontentsline{toc}{#1}{%
          \ifnum #2>\c@secnumdepth \else
             \protect\numberline{\csname the#1\endcsname}%
          \fi
          #7}%
     \fi
%    \end{macrocode}
% Also allow per-article TOC for review papers
%    \begin{macrocode}
     \if@reviewarticle
        \ifnum #2>2 \else
           \addcontentsline{rtc}{#1}{%
             \protect\numberline{\csname the#1\endcsname}%
             \sec@upcase{\sec@shape #7}}%
        \fi
     \fi
  \else
     \def\@svsechd{#6\hskip #3\relax  %% \relax added 2 May 90
       \@svsec #8\csname #1mark\endcsname{#7}%
       \if@guide
          \addcontentsline{toc}{#1}{%
            \ifnum #2>\c@secnumdepth \else
               \protect\numberline{\csname the#1\endcsname}%
            \fi
            #7}%
       \fi
       }%
  \fi
  \@xsect{#5}}
%    \end{macrocode}
% \end{macro}


% \subsection{Appendices}
% \label{sec:appendices}
% Make sure that appendices in individual articles work properly. We
% copy the definition from the article class (the book one doesn't do
% what we want).  
% \begin{environment}{appendices}
%   This uses an environment |appendices| inside of which each
%   |\section| will become a separate appendix.  Note that the
%   appendices must come before the references (otherwise the full
%   addresses mechanism will not work properly). We make the equation
%   numbering be subordinate to the appendix (equations A12, B6,
%   etc). Finally, we have a medskip followed by APPENDICES to
%   introduce them. \changes{3.17}{2002/04/08}{ Added appendices environment
%   that does the right thing.}
%    \begin{macrocode}
\newenvironment{appendices}{\par
  \begingroup
  \setcounter{section}{0}%
  \setcounter{subsection}{0}%
  \renewcommand\thesection{\@Alph\c@section}
  \renewcommand\theequation{\thesection\arabic{equation}}
  \medskip
  \section*{APPENDICES}}
{\endgroup}
%    \end{macrocode}
% \end{environment}
% \begin{environment}{appendix}
%   If there is only one appendix, then use |appendix| environment
%   instead.  \changes{4.0}{2019/11/12}{ Added
%   appendix environment for case where there is only one section.}
%    \begin{macrocode}
\renewenvironment{appendix}{\par
  \begingroup
  \setcounter{section}{0}%
  \setcounter{subsection}{0}%
  \renewcommand\thesection{\@Alph\c@section}
  \renewcommand\theequation{\thesection\arabic{equation}}
  \medskip
  \section*{APPENDIX}}
{\endgroup}
%    \end{macrocode}
% \end{environment}

% \subsection{Floating thingies}
% \label{sec:floats}   

%% 
%    \begin{macrocode}
\renewcommand{\thetable}{\@arabic\c@table}
\renewcommand{\thefigure}{\@arabic\c@figure}
\renewcommand{\figurename}{Fig.}
\renewcommand{\tablename}{TABLE}
%\renewcommand{\fnum@table}{\MakeUpperCase{\tablename}~\thetable}
%    \end{macrocode}
% For some reason, I have to write `TABLE' explicitly here
% If I use |\tablename|, it gets converted to mixed case---God knows
% why. 
% |\figurename| is also being ignored. 
%    \begin{macrocode}
\renewcommand{\fnum@table}{TABLE~\thetable}
\renewcommand{\fnum@figure}{Fig.~\thefigure}
%    \end{macrocode}

% Do captions in 9 pt face in the main journal style and in new-style
% anything. 
%    \begin{macrocode}
\long\def\@makefigcap#1#2{%
  \if@proceedings\else\small\fi
  \if@newstyle\small\fi 
  \vskip\abovecaptionskip
  \sbox\@tempboxa{#1. #2}%
  \ifdim \wd\@tempboxa >\hsize
    #1. #2 \par
  \else
    \global \@minipagefalse
    \hb@xt@\hsize{\hfil\box\@tempboxa\hfil}%
  \fi
  \vskip\belowcaptionskip}
%    \end{macrocode}

% Make sure table caption is upper case. 
%    \begin{macrocode}
\long\def\@maketabcap#1#2{%
  \hbox to\hsize{\hfil #1 \hfil}
  \vskip 6pt 
  \centering\MakeUpperCase{#2}\par
  \vskip 6pt}
%    \end{macrocode}

%    \begin{macrocode}
\newlength{\@rottabcapwidth}
\long\def\@makerottabcap#1#2{%
  \setbox\@tempboxa\hbox{\MakeUpperCase{#2}}%
  \ifdim\wd\@tempboxa >.8\vsize
     \setlength\@rottabcapwidth{.8\vsize}
  \else
     \setlength\@rottabcapwidth{\wd\@tempboxa}
  \fi
  \rotatebox{90}{%
    \begin{minipage}{\vsize}\centering
      \begin{minipage}{\@rottabcapwidth}\centering
        #1\\ \vskip 6pt \MakeUpperCase{#2}
      \end{minipage}%
    \end{minipage}
    }%
  \hskip 6pt
  }

\newenvironment{Sideways}
    {\begin{sideways}\begin{minipage}{\vsize}\centering}
    {\end{minipage}\end{sideways}}

\newenvironment{Sidewaystable}
{%
  \let\oldcaption=\caption%
  \def\caption{
    \def\@captype{table}\renewcommand{\@makecaption}{\@maketabcap}
    \oldcaption}
  \begin{sidewaystable}\centering}
  {\end{sidewaystable}\let\caption=\oldcaption}
\newenvironment{Sidewaystable*}
{%
  \let\oldcaption=\caption%
  \def\caption{
    \def\@captype{table}\renewcommand{\@makecaption}{\@maketabcap}
    \oldcaption}
  \begin{sidewaystable*}\centering}
  {\end{sidewaystable*}\let\caption=\oldcaption}
%    \end{macrocode}

% Not sure why we need to stick these lines in. 
%    \begin{macrocode}
\renewenvironment{figure}
                 {\renewcommand{\@makecaption}{\@makefigcap}\@float{figure}}
                 {\end@float}
\renewenvironment{figure*}
                 {\renewcommand{\@makecaption}{\@makefigcap}\@dblfloat{figure}}
                 {\end@dblfloat}
\renewenvironment{table}
                 {\renewcommand{\@makecaption}{\@maketabcap}\@float{table}}
                 {\end@float}
\renewenvironment{table*}
                 {\renewcommand{\@makecaption}{\@maketabcap}\@dblfloat{table}}
                 {\end@dblfloat}
%    \end{macrocode}

% Redefine |\@caption| so that we don't bother with the list of
% figures. This saves us having to worry about fragile commands in the
% captions. 
%    \begin{macrocode}
\long\def\@caption#1[#2]#3{%
  \par
%%  \addcontentsline{\csname ext@#1\endcsname}{#1}%
%%    {\protect\numberline{\csname the#1\endcsname}{\ignorespaces #2}}%
  \begingroup
    \@parboxrestore
    \if@minipage
      \@setminipage
    \fi
    \normalsize
    \@makecaption{\csname fnum@#1\endcsname}{\ignorespaces #3}\par
  \endgroup}
%    \end{macrocode}

% Make tables not quite so cramped. Allow |\hline\hline| to give a
% poor man's thick rule. 
%    \begin{macrocode}
\renewcommand{\arraystretch}{1.15}
\setlength{\doublerulesep}{0.4pt}
%    \end{macrocode}

% \begin{macro}{\tablecols}
% \begin{macro}{\@tablecols}
% \begin{macro}{\tabcaptioncontinued}
% |\tablecols| is a user command for setting the number of columns. It
% is used by commands that try and span the entire table (I think
% there must be a better way of acheiving
% this). |\tabcaptioncontinued| is 
% a user command for setting the caption of a continued table on the
% next page. 
%    \begin{macrocode}
\newcommand{\tablecols}[1]{\renewcommand{\@tablecols}{#1}}
\newcommand{\@tablecols}{1}
\newcommand{\tabcaptioncontinued}{%
  \multicolumn{\@tablecols}{c}{\addtocounter{table}{-1}TABLE~\thetable\ (CONTINUED)} \\ \\
  }
%    \end{macrocode}
% \end{macro}
% \end{macro}
% \end{macro}

% \subsubsection{Footnotes in tables}
% \label{sec:table-footnotes}
% \begin{macro}{\tablenotemark}
% \begin{macro}{\tablenote}
% \begin{macro}{\tablenoteskip}
% Old user commands |\tablenotemark{MARK}|,
% |\tablenote{MARK}{TEXT}|---now deprecated. 
%    \begin{macrocode}
\newcommand\tablenotemark[1]{$^\mathrm{#1}$} 
\newcommand\tablenote[2]{\if@proceedings\else\small\fi ${}^\mathrm{#1}$ #2} 
\newcommand\tablenoteskip{\\[0.5\baselineskip]} 
%    \end{macrocode}
% \end{macro}
% \end{macro}
% \end{macro}

% \begin{macro}{\tabnotemark}
% \begin{macro}{\tabnotetext}
% \begin{macro}{\tabnotetextTWO}
% Alternative way of doing footnotes in tables. 
%    \begin{macrocode}
\newcommand{\tabnotemark}[1]{\textsuperscript{\lowercase{#1}}}
%\newcommand{\tabnotetextTWO}[4]{\tabnotetext{#1}{#2}\tabnotetext{#3}{#4}}
\newcommand{\tabnotetextTWO}[4]{%
  \multicolumn{\@tablecols}{l}{%
    \parbox{0.5\tabnotewidth}{\small\textsuperscript{#1}#2}%
    \parbox{0.5\tabnotewidth}{\small\textsuperscript{#3}#4}%
    }%
  \\}
\newcommand{\tabnotetext}[2]{%
  \multicolumn{\@tablecols}{@{}p{\tabnotewidth}@{}}{\small\textsuperscript{#1}#2}\\}
\newlength{\tabnotewidth}
%    \end{macrocode}
% \end{macro}
% \end{macro}
% \end{macro}


% \begin{macro}{changemargin}
%    \begin{macrocode}
% Command to allow margins to be changed on the fly. Useful to allow
% tables and Figures to protrude into the margins. 
%
% \changes{3.24}{2007/09/14}{Added \texttt{changemargin} environment
% (copied from TeX FAQ).}
\newenvironment{changemargin}[2]{% 
  % From TeX FAQ http://www.tex.ac.uk/cgi-bin/texfaq2html?label=chngmargonfly
  \begin{list}{}{%
    \setlength{\topsep}{0pt}%
    \setlength{\leftmargin}{#1}%
    \setlength{\rightmargin}{#2}%
    \setlength{\listparindent}{\parindent}%
    \setlength{\itemindent}{\parindent}%
    \setlength{\parsep}{\parskip}%
  }%
  \item[]}{\end{list}}
%    \end{macrocode}
% \end{macro}

% \subsection{The bibliography and references}
% \label{sec:biblio}
% \begin{macro}{\beforebiblioskip}
% \begin{macro}{\afterbiblioskip}
% I'm not entirely sure why I introduced these two commands, since
% directly putting |\vspace| before and after
% |\begin{thebibliography}| seems to work well enough. 
%    \begin{macrocode}
\newlength{\beforebiblioskip}
\setlength{\beforebiblioskip}{0pt}
\newlength{\afterbiblioskip}
\setlength{\afterbiblioskip}{0pt}
%    \end{macrocode}
% \end{macro}
% \end{macro}
% \begin{macro}{thebibliography}
% The environment itself is done as starred section with extra
% vertical space after heading. This space is currently zero. 
% In the book style we do the cunning trick with rereading the file in
% order to get the references into the |\ADS@references| macro.
% 
% \changes{3.05}{2003/01/02}{Removed the \cs{beforebiblioskip} and
% \cs{afterbiblioskip} commands, since they cause unwanted vertical
% space in some instances.}  
%    \begin{macrocode}
\renewenvironment{thebibliography}{%
  \if@book\RereadBiblio\fi
%%  \vspace*{\beforebiblioskip}
  \section*{REFERENCES} \frenchspacing
%%  \vspace*{\afterbiblioskip}
  \small
%    \end{macrocode} 
% The references themselves are set as a list in order 
% to get the hanging indentation. They are done in a smaller font and
% with french spacing so we needn't worry about the distinction
% between sentence-ending periods and abbreviations. 
%    \begin{macrocode} 
  \list{\null}{\leftmargin 1.5em\labelwidth\z@\labelsep\z@\itemindent -1.5em
    \itemsep\z@\parsep\z@\usecounter{enumi}}
  \sloppy\clubpenalty4000\widowpenalty4000
  \sfcode`\.=1000\relax\@biblisttrue}
{\endlist\@biblistfalse\normalsize\@outputaddressesifnecessary}
%    \end{macrocode}
% \end{macro}
% \subsection{Uniquification of cross-references}
% \label{sec:uniquification}
% 
% \begin{macro}{\articlename}
% In order that the |\label{}| command in individual articles be
% guaranteed to give a unique anchor, we append the article name (held
% in the |\articlename| command) to each key. This requires
% redefinition of the following commands: |\cite|, |\nocite|, |\ref|,
% |\pageref|, |\label|. 
% \begin{macro}{\oldcite}
% \begin{macro}{\oldnocite}
% \begin{macro}{\oldref}
% \begin{macro}{\oldpageref}
% \begin{macro}{\oldlabel}
% The original versions of these commands are
% saved as |\oldcite|, |\oldnocite|, etc. 
%
% The uniquifying of cross-references is now always done, so 
% we don't get multiply-defined labels in the single article case
%    \begin{macrocode}
\newcommand{\articlename}{anon}
\AtBeginDocument{
  \let\oldcite=\cite 
  \let\oldnocite=\nocite 
  \let\oldref=\ref 
  \let\oldpageref=\pageref
  \let\oldlabel=\label
}
%    \end{macrocode}
% \end{macro}
% \end{macro}
% \end{macro}
% \end{macro}
% \end{macro}
% \end{macro}
%
% Uniquify all cross-references by prepending the article name.  This
% now has to go into |\AtBeginDocument| so it is working on the new
% versions that hyperref.sty has set up.  Then we need to put the
% |\@begofarticlehook| afterwards, so that the first page label is using
% the new definitions.  If amsmath is loaded, then we also need to patch
% the special label command that they use. \changes{4.0}{2019/11/18}{Fix labels for amsmath display equations, in case that package has been loaded.}
%    \begin{macrocode}
\AtBeginDocument{
  \renewcommand{\ref}[1]{\oldref{\articlename:#1}%
    \if@debug{\ClassInfo{rmaa}{Reference to \articlename:#1}}\fi}
  \renewcommand{\pageref}[1]{\oldpageref{\articlename:#1}%
    \if@debug{\ClassInfo{rmaa}{Page reference to \articlename:#1}}\fi}
  \renewcommand{\label}[1]{\oldlabel{\articlename:#1}%
    \if@debug{\ClassInfo{rmaa}{Label for \articlename:#1}}\fi}%
  \ifdef{\label@in@display}{%
    \let\oldlabel@in@display=\label@in@display
    \renewcommand{\label@in@display}[1]{\oldlabel@in@display{\articlename:#1}%
      \if@debug{\ClassInfo{rmaa}{AMSTeX Label for \articlename:#1}}\fi}%
  }{}
  \@begofarticlehook
}%
%    \end{macrocode}
% Citations are more complicated since they can contain a list of
% keys. These are really only here for use with the rmaacite package,
% which is utterly deprecated (02 Jan 2003) now that I have the natbib
% code working. The command \cs{cite} should \emph{not} be used with
% natbib: use \cs{citep}, \cs{citet}, etc.\@ instead. 
%    \begin{macrocode}
\renewcommand{\nocite}[1]{\@bsphack
  \@for\@citeb:=#1\do{%
    \edef\@citeb{\expandafter\@firstofone\@citeb:\articlename}%
    \if@filesw\immediate\write\@auxout{\string\citation{\@citeb}}\fi
    \@ifundefined{b@\@citeb:\articlename}{\G@refundefinedtrue
      \@latex@warning{Citation `\@citeb' undefined}}{}}%
  \@esphack}
\renewcommand{\cite}[1]{\oldcite{#1:\articlename}}
%   \end{macrocode}


% Horrible fiddling with the bibliography citation mechanism. 
% This has to be globally available in case rmaacite puts it in the
% \texttt{.aux} file. 
%    \begin{macrocode}
\def\lbibcite#1#2#3{\global\@namedef{b@#1}{#2}
                    \global\@namedef{by@#1}{#3}}
%    \end{macrocode}

% \begin{macro}{\oldbibitem}
% \begin{macro}{\@rmaaBibArg}
% \begin{macro}{\@rmaaSpace}
% \begin{macro}{rmaabib}
% Of course, |\bibitem| needs to get the uniquification treatment as
% well. While we are at it, we redefine it to use consecutive integers
% for the key if the user leaves it blank (using the counter
% |rmaabib|). This allows one to put 
% |\bibitem{}| for all your references without getting annoying
% messages about multiply-defined labels.
% 
% We also take the opportunity to insert the optional argument to
% \cs{bibitem} if it is missing in the form Author(DATE) so that natbib
% doesn't try to do numerical citations. Natbib already provides a
% hook \cs{@extra@b@citeb} to help uniquify the citation key, so we
% just plug into that. However, it puts it \emph{after} the original
% key, not before as I had had it. Easiest solution is to change the
% way that I do it. 02 Jan 2003. Note that the \cs{@extra@b@citeb}
% does not get written to the .aux  file in the \cs{citation} command
% but it does work. 
% \changes{3.14}{06 Feb 2003}{Reinstalled old definition of \cs{bibitem}
%  for the case where the nonatbib option is requested, so that the
%  rmaacite package will work. This means that you can't use rmaacite
%  and natbib in the same volume.}
%    \begin{macrocode}
\if@natbib\def\@extra@b@citeb{:\articlename}\fi
%%\else
\let\oldbibitem\bibitem
\newcommand{\@rmaaBibArg}{}
\newcommand{\@rmaaSpace}{ }
\newcounter{rmaabib}
\if@natbib
\renewcommand{\bibitem}[2][ABC(2003)]{%
  \renewcommand{\@rmaaBibArg}{#2}
  \ifx\@rmaaBibArg\@rmaaEmpty
     \addtocounter{rmaabib}{1}
     \oldbibitem[#1]{\thermaabib:\articlename}
  \else\ifx\@rmaaBibArg\@rmaaSpace
     \addtocounter{rmaabib}{1}
     \oldbibitem[#1]{\thermaabib:\articlename}
  \else 
     \oldbibitem[#1]{#2:\articlename}
  \fi\fi
  \renewcommand{\@rmaaBibArg}{}
  }
\else
\renewcommand{\bibitem}[1]{%
  \renewcommand{\@rmaaBibArg}{#1}
  \ifx\@rmaaBibArg\@rmaaEmpty
     \addtocounter{rmaabib}{1}
     \oldbibitem{\articlename:\thermaabib}
  \else\ifx\@rmaaBibArg\@rmaaSpace
     \addtocounter{rmaabib}{1}
     \oldbibitem{\articlename:\thermaabib}
  \else
     \oldbibitem{\articlename:#1}
  \fi\fi
  \renewcommand{\@rmaaBibArg}{}
  }
\fi
%%\fi
%    \end{macrocode}
% \end{macro}
% \end{macro}
% \end{macro}
% \end{macro}

% \begin{macro}{oldpcite}
% \begin{macro}{oldscite}
% Then, we need to repeat the above scheme, but for the |rmaacite|
% definition. 
%    \begin{macrocode}
\newcommand{\@redefinebibitem}{%
  \let\oldbibitem\bibitem
  \let\oldcite\cite
  \let\oldpcite\pcite
  \let\oldscite\scite
  \def\bibitem[##1<##2>]##3{%
    \renewcommand{\@rmaaBibArg}{##3}
    \ifx\@rmaaBibArg\@rmaaEmpty
       \addtocounter{rmaabib}{1}
       \oldbibitem[##1<##2>]{\thermaabib:\articlename}
    \else\ifx\@rmaaBibArg\@rmaaSpace
       \addtocounter{rmaabib}{1}
       \oldbibitem[##1<##2>]{\thermaabib:\articlename}
    \else 
       \oldbibitem[##1<##2>]{##3:\articlename}
    % trust the author here
%%       \oldbibitem[##1<##2>]{##3}
    \fi\fi
    }
  }
%    \end{macrocode}
% \end{macro}
% \end{macro}


% \subsection{HTML Table of Contents}
% \label{sec:hyper-TOC}
% \begin{macro}{\HyperTOCEntry}
% \begin{macro}{\HyperTOCpart}
% \begin{macro}{\HyperTOCmeeting}
% \begin{macro}{\HyperTOCeditors}
% \begin{macro}{\HyperTOCpreamble}
% These commands are used to write information to the \texttt{.htoc}
% file, which is used in generating a Web Page of the Table of
% Contents. Since they only ever appear preceded by |\protect| in the
% argument to an |\addtocontents| command,  they should never need to
% be expanded by the macros. However, we define them here just in case
% \dots
%    \begin{macrocode}
\newcommand{\HyperTOCEntry}[3]{\relax}
\newcommand{\HyperTOCpart}[1]{\relax}
\newcommand{\HyperTOCmeeting}[3]{\relax}
\newcommand{\HyperTOCeditors}[1]{\relax}
\newcommand{\HyperTOCpreamble}{\relax}
%    \end{macrocode}
% \end{macro}
% \end{macro}
% \end{macro}
% \end{macro}
% \end{macro}
%

% \subsection{Formatting the poster abstracts} 
% These can now be used without using the book option.
%    \begin{macrocode}

\newcommand{\@makeminititle}{%
  \newpage
  \null 
  \vskip 3em%
  \begin{center}%
    {\@title \par}%
  \end{center}
  \vskip 3em%
  }

\newif\if@abstractsads
\@abstractsadsfalse
\newcommand\abstractsADS{\@abstractsadstrue}

\if@book\RequirePackage{multicol}\fi
\newenvironment{abstracts}[1][Abstracts of Contributed Papers]{%
  %% First, make sure we start on an odd page (06 Jun 2000)
  \if@book\clearpage\thispagestyle{empty}\cleardoublepage\fi
  %%
  \renewcommand{\articlename}{ABSTRACTS}
  \title{#1}\shortauthor{Abstracts}\shorttitle{Abstracts}
  \listofauthors{}
  \if@book
  \addcontentsline{toc}{chapter}{\protect\@fiddleauthors{\@TOCtitle}{\rmaaTOCauthors}}%
  \addtocontents{htc}{\protect\HyperTOCEntry{\@rawtitle}{\@rawauthors}{\articlename}}
  \fi
  \pagestyle{rmaa} % this now seems necessary (06 Jun 2000)
  \thispagestyle{rmaafirstpage}
%%  \twocolumn[\@makeminititle]
  \renewcommand{\abstract}[1]{%
    \renewcommand\ADS@abstract{##1}%
    \renewcommand{\@abstract}{##1}\@writeabstract%
  }
  \renewcommand{\title}[2][]{%
    \gdef\ADS@authors{}%
    \gdef\ADS@affil{}%
    \gdef\ADS@references{}%
    \gdef\ADS@keywords{}%
    \renewcommand{\@rmaaArg}{##1}
    \ifx\@rmaaArg\@rmaaEmpty
       \renewcommand{\@TOCaftertitle}{}
    \else
       \renewcommand{\@TOCaftertitle}{ (\@rmaaArg)}
    \fi
    \renewcommand{\@title}{\rmaa@TOCTitleStyle{##2}}%
    \renewcommand{\@TOCtitle}{\rmaa@TOCTitleStyle{##2}\@TOCaftertitle}%
    \renewcommand{\rmaaIndexTitle}{\rmaa@IndexTitleStyle{##2\@TOCaftertitle.}}%
    \renewcommand\ADS@title{##2}%
    \xdef\ADS@firstpage{\@arabic{\c@page}}%
    }
  \renewcommand{\author}[1]{\let\@thanks\@empty\renewcommand\@author{##1} }
  \renewcommand{\footnotetext}[2][XX]{\par\noindent\textsuperscript{##1}~##2}
  \renewcommand{\footnoterule}{\kern-3\p@ \hrule \@width 3em \kern 2.6\p@}
  \if@book\@begofarticlehook\fi
  \if@book\begin{multicols}{2}[\@makeminititle]\else\@makeminititle\fi
  \begingroup
  \setcounter{footnote}{0}
  }% and at the end of the environment....
  {\endgroup\if@book\end{multicols}\listofabstracts\label{RMAALastPage}\clearpage\fi}
%% auxiliary commands used:
%% \@writeabstract does all the actual formatting 
\newcommand{\@writeabstract}{%
  \if@book%
  \addcontentsline{loa}{abstract}{\protect\@fiddleauthors{\@TOCtitle}{\rmaaTOCauthors}}%
  \fi%
  \begin{center}\@title\par\@author\end{center}
  \@abstract
  \par\noindent\footnoterule\nobreak\footnotesize\@thanks
  \par\noindent\hrulefill\goodbreak
  \vskip 1cm \@plus0.5cm
  %% write ADS entry
  \if@book\if@abstractsads\ADS@dumpdata\fi\fi%
  \endgroup % this pair is to make each abstract a self-contained group
  \begingroup
  \setcounter{footnote}{0}
  }

\if@book
%% List of abstracts (stuff written to .loa file)
\newcommand{\listofabstracts}{%
  \begin{multicols}{2}[%
    \vskip 1.5cm \vfill % leave at least 1.5 cm and go to bottom of page if possible
    \centerline{\MakeUpperCase{List of Abstracts}}
    \vskip 1.5cm \@plus0.5cm
    ][6cm] % this makes it start a new page if there is not at least 6cm of space
    \renewcommand{\@fiddleauthors}[2]{##1\hskip1em plus2em minus0.4em\penalty-100  ##2}
    \parindent\z@
%%    \parskip.5\p@ \@plus .7\p@\relax 
    \parskip3\p@ \@plus 1\p@ \@minus 1\p@\relax % allow quite a lot of
                                % leeway here to avoid breaks inside
                                % entries 
    \renewcommand{\\}{\relax} % temporarily turn off explicit line breaks just in case
    \renewcommand{\altaffilmark}[1]{\relax} % ditto with this
    \@starttoc{loa}
  \end{multicols}
  }
\newcommand*{\l@abstract}[2]{%
  \par\penalty-1000
  #1\hskip 1em\penalty0\null\nobreak\hfill#2}
\fi

%    \end{macrocode}
 
% \subsection{Code for typesetting an entire volume of conference
% proceedings} \label{sec:whole-volume}
%    \begin{macrocode}

\if@book
%% environment for including the individual articles
\newcommand{\@fiddleauthors}[2]{\relax}
\renewcommand{\theequation}{\@arabic\c@equation}
\newenvironment{articles}%
{ \begingroup % make it a group so that author definitions do not
              % spill over
  \gdef\ADS@authors{}
  \pagestyle{rmaa}
  \renewenvironment{document}{%
    \@begofarticlehook
    \addtocounter{chapter}{1}%
    \addcontentsline{toc}{chapter}{%
      \protect\@fiddleauthors{\@TOCtitle}{\rmaaTOCauthors}}%
    \addtocontents{htc}{\protect\HyperTOCEntry{\@rawtitle}{\@rawauthors}{\articlename}}
    }{\relax}
  \renewcommand{\documentclass}[2][foo]{\relax} 
  \renewcommand{\usepackage}[2][foo]{%
    \def\@packagename{##2}\def\@rmaacitename{rmaacite}%
    \ifx\@packagename\@rmaacitename
    \@rmcitetrue\ClassInfo{rmaa}{This paper uses rmaacite}
    \makeatletter\input{rmaacite.sty}\makeatother % read it in here so that all changes
    \else                                         % are local to paper
    \if@debug\ClassWarning{rmaa}{Package \@packagename\space was requested by \articlename\space}\fi
    \fi} 
  %% Don't allow individual articles to mess with the pages, volume or year
  \renewcommand{\SetFirstPage}[1]{\relax}
  \renewcommand{\SetYear}[1]{\relax}
  \renewcommand{\SetVolume}[1]{\relax}
  \renewcommand{\SetConfTitle}[1]{\relax}
  %% Don't allow individual articles to change these
  \renewcommand{\nonstopmode}{\relax}
  \renewcommand{\scrollmode}{\relax}
  \renewcommand{\outputplates}{\relax}
  %% XXXX need to check if using rmaacite and act accordingly
  %% avoid multiple declarations of the same length
  \let\oldnewlength=\newlength
  \renewcommand{\newlength}[1]{%
    \edef\reserved@a{\expandafter\@gobble\string ##1}%
    \@ifundefined{\reserved@a}{%
      \oldnewlength{##1}}{%
      \if@debug{\ClassInfo{rmaa}{\reserved@a\space already defined: doing nothing}}\fi%
      \relax}}%
  \renewcommand\newcommand{\providecommand}
  }%
{\endgroup\onecolumn}
     

%    \end{macrocode}
% Command to be used in `articles' environment for including an
% article
%    \begin{macrocode}
\newcommand{\article}[2][]{%
  % optional first argument is extra text to put in brackets after title in TOC
  % e.g. ``Invited Talk'' or ``Poster''
  \renewcommand{\articlename}{#2}
  % reset loads of counters
  \setcounter{footnote}{0}
  \setcounter{figure}{0}
  \setcounter{table}{0}
  \setcounter{section}{0}
  \setcounter{equation}{0}
  \setcounter{rmaabib}{0}
  \setcounter{author}{0}
  \setcounter{AUTHOR}{1}
  \gdef\ADS@authors{}%
  \gdef\ADS@affil{}%
  \gdef\ADS@affiliations{}%
  \gdef\ADS@references{}%
  \@firstauthortrue
  \@firstaffiltrue
  \@firstaltaffiltrue
  \@firstkwdtrue
  \@rmcitefalse
  \if@suppressfulladdresses\else\@needaddressestrue\fi
  \renewcommand{\@rmaaArg}{#1}
  \ifx\@rmaaArg\@rmaaEmpty
     \renewcommand{\@TOCaftertitle}{}
  \else
     \renewcommand{\@TOCaftertitle}{ (#1)}
  \fi
%    \end{macrocode}
% In the main journal articles must start on an odd page. 
%    \begin{macrocode}
  \if@journal\clearpage\thispagestyle{empty}\cleardoublepage\fi
%% Now input the file, then run the end-of-article hook (reinstated 28
  %% Sep 2001)
  {\input{\@articledir/\articlename.tex}\@endofarticlehook}}
%    \end{macrocode}
% Command to separate articles into sessions. Write an entry in the
% TOC, makes sure we start on an odd page (added 06 Jun 2000), then do
% a title page if \cs{SessionTitlePages} was called. 
% \changes{3.12}{04 Feb 2003}{Added option to have title page for each
% session. Turn it on with \cs{SessionTitlePages}.}
%    \begin{macrocode}
\newcommand{\session}[1]{
  %% this does nothing except write an entry in the TOC
  \addcontentsline{toc}{part}{\rmaa@TOCPartStyle{#1}}
  \addtocontents{htc}{\protect\HyperTOCpart{#1}}
  %% and also makes sure we start on an odd page (added 06 Jun 2000)
  \clearpage\thispagestyle{empty}\cleardoublepage
  \stepcounter{part}
  \if@sessiontitlepages{%
    \onecolumn
    \thispagestyle{empty}
    \renewcommand{\articlename}{Section\thepart TitlePage}
    \label{RMAAFirstPage}
    \make@session@title{#1}
    \label{RMAALastPage}
    \clearpage\thispagestyle{empty}\cleardoublepage
  }\fi
  }
\newcommand{\make@session@title}[1]{%
  {%
    \setlength\parindent{0pt}
    \vspace*{\stretch{1}}
    \rule{\linewidth}{1mm}
    \begin{flushright}\Huge\bfseries\boldmath
      #1
    \end{flushright}
    \rule{\linewidth}{1mm}
    \vspace*{\stretch{2}}
  }
}
%    \end{macrocode}
% Command to turn on separate title pages for each session. 
%    \begin{macrocode}
\newcommand{\SessionTitlePages}{\@sessiontitlepagestrue}
\newcommand{\BoldSessionTitles}{%
  \renewcommand{\rmaa@TOCPartStyle}[1]{\textbf{\rmaa@TOCTitleStyle{##1}}}
  }
\newcommand{\ArticleDir}[1]{\renewcommand{\@articledir}{#1/}\graphicspath{{1/}}}
\newcommand{\@articledir}{.}
%\newcommand{\ps@rmaafront}{\ps@rmaa}% page style for front matter (same as ever, actually)
\newcommand{\ps@rmaafront}{% page style for front matter:  page # in each foot
  \renewcommand\@oddfoot{\hfil{\thepage}}%
  \renewcommand\@evenfoot{{\thepage}\hfil}%
  % shorttitle in both heads
  \renewcommand\@evenhead{\hfil\@shorttitle\hfil}%
  \renewcommand\@oddhead{\@evenhead}
  }
\newcommand{\ps@rmaafrontfirstpage}{% page style for first page of front matter
  % page # in each foot, on outside edge
  \renewcommand\@oddfoot{\hfil\thepage}%
  \renewcommand\@evenfoot{\thepage\hfil}%
  % nothing in the even heads
  \renewcommand\@evenhead{\relax}%
  \renewcommand\@oddhead{\relax}%
  }
%% slight modification to the standard \frontmatter and \mainmatter
%% commands, just to set the \pagestyle
\renewcommand\frontmatter{%
  \latexonecolumn
  \cleardoublepage
  \@frontmattertrue
  \@mainmatterfalse
  \pagenumbering{roman}
  \pagestyle{rmaafront}
  \addtocounter{page}{2} % leave space for the title page
  }
\renewcommand\mainmatter{%
  \pagestyle{empty}
  \if@rmtwocol\twocolumn\fi
  \cleardoublepage
  \@frontmatterfalse
  \@mainmattertrue
  \pagenumbering{arabic}
  \pagestyle{rmaa}
  }
\renewcommand\backmatter{%
  \clearpage\pagestyle{empty}\cleardoublepage
  \@mainmatterfalse
  %% put extra space in the TOC
  \addtocontents{toc}{\bigskip}
  }
%    \end{macrocode}
%
% \subsection{Front sections}\label{sec:front-sections}
% \begin{macro}{\notitle}
% \begin{macro}{\shortname}
% \begin{macro}{frontsection}
% The environment \texttt{frontsection} is for doing a section in the
% front matter (preface, list of participants, or whatever). It has 
% one optional argument:
% \begin{enumerate}
% \item[\#1] Short title (no embedded spaces) for
% use in making PS files. E.g., ``list\_of\_participants''.   
% \end{enumerate}
% and two mandatory arguments: 
% \begin{enumerate}
% \item[\#2] Title of section (e.g., `Preface')
% \item[\#3] Commands to be evaluated before processing the contents
% of the environment. This is mainly useful for putting in the things
% that would go in the preamble of a normal article  (e.g.,
% |\listofauthors|, |\indexauthor|, |\notitle|) 
% \end{enumerate}
%    \begin{macrocode}
\newcommand{\notitle}{\@titlevisiblefalse} \newcommand{\shortname}{}
\newenvironment{frontsection}[3][]{%
  \begingroup
  \renewcommand{\shortname}{#1}
  \ifx\shortname\@rmaaEmpty
     \renewcommand{\articlename}{#2}
  \else
     \renewcommand{\articlename}{\shortname}
  \fi
  \@begofarticlehook
  \@titlevisibletrue % default
  \thispagestyle{rmaafrontfirstpage}
  \listofauthors{} % blank out this in case it is not specified
  \title{{#2}}\shorttitle{#2}\shortauthor{#2}
  #3 % evaluate the command given as second mandatory arg
  \addcontentsline{toc}{chapter}{\protect\@fiddleauthors{\@TOCtitle}{\rmaaTOCauthors}}%
  \addtocontents{htc}{\protect\HyperTOCEntry{\@rawtitle}{\@rawauthors}{\articlename}}
  \if@titlevisible
     \vspace*{1cm}\centerline{\rmaa@FrontSecStyle{\@title}}\vspace*{1cm}
  \fi
  \setlength{\parskip}{3ex}
  }{\frontlabel{RMAALastPage}\clearpage\endgroup}
%    \end{macrocode}
% \end{macro}
% \end{macro}
% \end{macro}
%    \begin{macrocode}
%% Definition of the command \frontlabel, which rights out something like `F4' to
%% the .aux file, instead of `\textlatin  {i}'. This is so it is easy to post-process
%% the file for making PS files of individual sections
\newcommand{\frontlabel}[1]{%
  \@bsphack
  \protected@write\@auxout{}%
     {\string\newlabel{\articlename:#1}{{}{F\@arabic{\c@page}}}}
  \@esphack
  }
%% User command to say X, Y, Z,\\ Editors of the Proceedings
\newcommand{\EditorSignature}{\vspace*{\parskip}\rightline{\@editors}%
  \rightline{Editors of the proceedings}}
%% Fiddle with \thesection, etc. so that they don't include the chapter number
%% Basically go back to the `article' definitions of these
\renewcommand{\thesection}{\@arabic\c@section}
\renewcommand{\thefigure}{\@arabic\c@figure}
\renewcommand{\thetable}{\@arabic\c@table}

\newcommand{\@oldargs}{}
\newcommand{\@newargs}{}
%% add TOC entry for each paper
%\renewcommand{\title}[1]%
%{\renewcommand{\@title}{\rmaa@TitleStyle{#1}}%
%\addtocounter{chapter}{1}\addcontentsline{toc}{chapter}{#1}%
%}

%% Formatting of the table of contents
%% dimensions of various parts
\newlength{\TOCtextwidth}\setlength{\TOCtextwidth}{12.65cm}
\newlength{\TOCpagenumwidth}\setlength{\TOCpagenumwidth}{1.3cm}
\newlength{\TOCentrywidth}\setlength{\TOCentrywidth}{\TOCtextwidth}
\addtolength{\TOCentrywidth}{-\TOCpagenumwidth} 
\newlength{\TOCextramargin}\setlength{\TOCextramargin}{1.25cm}
\newcommand{\TOCbox}[1]{\hspace{\TOCextramargin}\parbox{\TOCtextwidth}{#1}\par}
\newlength{\TOCextramarginx}\setlength{\TOCextramarginx}{\TOCextramargin}\addtolength{\TOCextramarginx}{4.3cm}
\newlength{\TOCtextwidthx}\setlength{\TOCtextwidthx}{\TOCtextwidth}\addtolength{\TOCtextwidthx}{-4.3cm}
\newcommand{\TOCboxx}[1]{\hspace{\TOCextramarginx}\parbox{\TOCtextwidthx}{#1}\par}

\newcommand{\RMAAbannerimage}{rmsc}
\renewcommand\tableofcontents{%
  \renewcommand{\articlename}{CONTENTS}
  \begingroup
  \@begofarticlehook
  \if@twocolumn
     \@restonecoltrue\onecolumn
  \else
     \@restonecolfalse
  \fi
  %% 9pt for journal style
  \if@proceedings\else\small\fi
  \if@newstyle\small\fi %% same for new style
  \if@proceedings\thispagestyle{rmaafrontfirstpage}\else\thispagestyle{empty}\fi
  \shortauthor{Contents}\shorttitle{Contents}
  \vspace*{18mm}
  \if@printbanner
  %% put in the banner header (need to get better image of this)
     \TOCbox{\hspace*{-0.15cm}\includegraphics[width=0.9\TOCtextwidth]{\RMAAbannerimage}}
     \vspace*{8mm}
  %% volume and publication date (should really be in funny font)
     \TOCboxx{\@volumen\\[2mm] \@mes\ \@anyo}
  \else
  %% put in the banner header (need to get better image of this)
     \TOCbox{\hspace*{-0.15cm}\includegraphics[width=0.9\TOCtextwidth]{\RMAAbannerimage}}
     \vspace*{12.5mm}
  %% volume and publication date (should really be in funny font)
     \TOCboxx{~\\[2mm] ~}
  \fi
  \vspace*{16mm}
  %% This is the command that fiddles with the TOC entry so that the
  %% title is flushleft but the authors are flushright
  \renewcommand{\@fiddleauthors}[2]{%
    % first put the paper title
    \renewcommand{\\}{\relax} % temporarily turn off explicit line breaks just in case
    \renewcommand{\altaffilmark}[1]{\relax} % ditto with this
    \renewcommand{\textsuperscript}[1]{\relax} % ditto with this 
    ##1
    % now save the author list in a box
    \sbox\@tempboxa{##2}%
    % check that box is shorter than the total width available for the entry
    \ifdim \wd\@tempboxa <\TOCentrywidth
       % if so, then put it right here, flush against right margin
       \penalty10000\hskip1em plus10em\penalty0\hbox{}\nobreak\hfill\box\@tempboxa
       % Note that we force a space of at least 1em between title and authorlist, 
       % stretchable to 11em since we may have very short titles but still with an 
       % authorlist too long to go on the same line.  The only linebreak allowed is 
       % directly before the authorlist.
    \else
       % if it is too long, we set the original authorlist text in its own parbox
       \par\vskip1ex\parbox[b]{\TOCentrywidth}{\raggedleft ##2}
    \fi
    }
  \@starttoc{toc}%
  \if@restonecol\twocolumn\fi
  \frontlabel{RMAALastPage}
  \clearpage
  \thispagestyle{empty}
  \cleardoublepage % make sure we start on an odd page
  \endgroup

  %%% ALSO, do the HTML TOC
  %% The file that will be used to make the HTML TOC
  \@starttoc{htc}
  %% All the stuff in this file does nothing in the book style
  \addtocontents{htc}{\protect\HyperTOCmeeting{\@conftitle}{\@reunionplace}{\@reuniondate}}
  \addtocontents{htc}{\protect\HyperTOCeditors{\@editors}}
  \addtocontents{htc}{\protect\HyperTOCpreamble}
  \addtocontents{htc}{\protect\HyperTOCEntry{Table of Contents}{}{\articlename}}

  %% and the ADS file
  \ADS@openfile
  }
%% This is now much simpler
\renewcommand\l@chapter[2]{%TOC entry for articles
  \ifnum \c@tocdepth >\m@ne
    \vskip 1.0em \@plus\p@
    \TOCbox{%
      \parbox[b]{\TOCentrywidth}{#1}\parbox[b]{\TOCpagenumwidth}{\hfill#2}
    }
  \fi} 
\renewcommand\l@part[2]{%TOC entry for sessions: centered, no page number
  \ifnum \c@tocdepth >\m@ne
    \vskip 1.5em \@plus1.5em \goodbreak
    \TOCbox{%
      \parbox[b]{\TOCentrywidth}{\hfil#1\hfil}\hfill
    }
  \fi\nobreak} 

%% Formatting theindex
%% strongly encourage page breaks between index entries
\renewcommand{\@idxitem}{\par\penalty-1000\hangindent.45cm} 
\renewcommand{\indexspace}{\par\penalty-1000\vskip 10\p@ \@plus15\p@ \@minus3\p@\relax}
\renewenvironment{theindex}{%
  %% first, make sure we start on an odd page (06 Jun 2000)
  \clearpage\thispagestyle{empty}\cleardoublepage
  %% this looks a bit like an article but is two columns and without abstract, etc
  \renewcommand{\articlename}{INDEX}
  \title{Author Index}\shortauthor{Author Index}\shorttitle{Author Index}
  \listofauthors{}
  \addcontentsline{toc}{chapter}{\protect\@fiddleauthors{\@TOCtitle}{\rmaaTOCauthors}}%
  \addtocontents{htc}{\protect\HyperTOCEntry{\@rawtitle}{\@rawauthors}{\articlename}}
  \pagestyle{rmaa}
  \thispagestyle{rmaafirstpage}
  \@begofarticlehook
  %% 9pt for journal style
  \if@proceedings\else\small\fi
  \if@newstyle\small\fi
  \twocolumn[\@makeminititle]
  %% vertical spacing (see also definition of \@idxitem and \indexspace
  \parskip\z@ \@plus .6\p@\relax % allow quite a lot of leaway here to avoid breaks inside entries
  %% horizontal spacing, breaking and glue parameters (see also the stuff in the definition 
  %% of \indexauthor
  \parindent\z@ 
  \exhyphenpenalty10000 % don't break on explicit hyphens, it always looks bad
  \rightskip=0pt plus.2em % allow a tiny bit of leeway at the ends of lines
  \spaceskip.3333em plus.5em minus.15em % a lot of leeway between words
  \xspaceskip.5em plus.5em minus.4em % this seems to have no effect
  \let\item\@idxitem
  \renewcommand{\\}{\relax} % temporarily turn off explicit line breaks just in case
  \renewcommand{\altaffilmark}[1]{\relax} % ditto with this
  \renewcommand{\textsuperscript}[1]{\relax} % ditto with this
  }{\label{RMAALastPage}\clearpage} % clearpage necessary so that we know about \shortauthor



%% end definitions for book styles
\fi

%    \end{macrocode}
% In order to do away with plates, we sometimes want a figure to
% appear in color only in the online version. The following command
% can bve used in a figure caption to typeset a note to this effect. 
%    \begin{macrocode}
\newcommand\ColorOnlineOnly{%
  \textsc{Note: this figure is available in color in the electronic
    version of this article, obtainable from}
  \texttt{http://\discretionary{}{}{}www.\discretionary{}{}{}astroscu.\discretionary{}{}{}unam.\discretionary{}{}{}mx/\discretionary{}{}{}\~{}rmaa/}.}
%% Plates
%% Define 2 counters: plate is the number in the whole document, plateloop is used 
%% internally by \outputplates
\newcounter{plate}
\newcounter{plateloop}
%% Currently limited to 15 plates in document
%% when we do a label, we need to write to \newlabel entries to the .aux file, one 
%% that says RMAAPlateX so that it can be referred to in \ouputplates, and one that
%% says articlename:RMAAPlateX so that the dvips program can find it and asociate it
%% with the right article
\newcommand{\@platelabeli}{\oldlabel{RMAAPlate1}\label{RMAAPlate1}}
\newcommand{\@platelabelii}{\oldlabel{RMAAPlate2}\label{RMAAPlate2}}
\newcommand{\@platelabeliii}{\oldlabel{RMAAPlate3}\label{RMAAPlate3}}
\newcommand{\@platelabeliv}{\oldlabel{RMAAPlate4}\label{RMAAPlate4}}
\newcommand{\@platelabelv}{\oldlabel{RMAAPlate5}\label{RMAAPlate5}}
\newcommand{\@platelabelvi}{\oldlabel{RMAAPlate6}\label{RMAAPlate6}}
\newcommand{\@platelabelvii}{\oldlabel{RMAAPlate7}\label{RMAAPlate7}}
\newcommand{\@platelabelviii}{\oldlabel{RMAAPlate8}\label{RMAAPlate8}}
\newcommand{\@platelabelix}{\oldlabel{RMAAPlate9}\label{RMAAPlate9}}
\newcommand{\@platelabelx}{\oldlabel{RMAAPlate10}\label{RMAAPlate10}}
\newcommand{\@platelabelxi}{\oldlabel{RMAAPlate11}\label{RMAAPlate11}}
\newcommand{\@platelabelxii}{\oldlabel{RMAAPlate12}\label{RMAAPlate12}}
\newcommand{\@platelabelxiii}{\oldlabel{RMAAPlate13}\label{RMAAPlate13}}
\newcommand{\@platelabelxiv}{\oldlabel{RMAAPlate14}\label{RMAAPlate14}}
\newcommand{\@platelabelxv}{\oldlabel{RMAAPlate15}\label{RMAAPlate15}}
\newcommand{\@outputplatelabeli}{\oldlabel{RMAAOutputPlate1}}
\newcommand{\@outputplatelabelii}{\oldlabel{RMAAOutputPlate2}}
\newcommand{\@outputplatelabeliii}{\oldlabel{RMAAOutputPlate3}}
\newcommand{\@outputplatelabeliv}{\oldlabel{RMAAOutputPlate4}}
\newcommand{\@outputplatelabelv}{\oldlabel{RMAAOutputPlate5}}
\newcommand{\@outputplatelabelvi}{\oldlabel{RMAAOutputPlate6}}
\newcommand{\@outputplatelabelvii}{\oldlabel{RMAAOutputPlate7}}
\newcommand{\@outputplatelabelviii}{\oldlabel{RMAAOutputPlate8}}
\newcommand{\@outputplatelabelix}{\oldlabel{RMAAOutputPlate9}}
\newcommand{\@outputplatelabelx}{\oldlabel{RMAAOutputPlate10}}
\newcommand{\@outputplatelabelxi}{\oldlabel{RMAAOutputPlate11}}
\newcommand{\@outputplatelabelxii}{\oldlabel{RMAAOutputPlate12}}
\newcommand{\@outputplatelabelxiii}{\oldlabel{RMAAOutputPlate13}}
\newcommand{\@outputplatelabelxiv}{\oldlabel{RMAAOutputPlate14}}
\newcommand{\@outputplatelabelxv}{\oldlabel{RMAAOutputPlate15}}
\newcommand{\@platerefi}{\oldpageref{RMAAPlate1}}
\newcommand{\@platerefii}{\oldpageref{RMAAPlate2}}
\newcommand{\@platerefiii}{\oldpageref{RMAAPlate3}}
\newcommand{\@platerefiv}{\oldpageref{RMAAPlate4}}
\newcommand{\@platerefv}{\oldpageref{RMAAPlate5}}
\newcommand{\@platerefvi}{\oldpageref{RMAAPlate6}}
\newcommand{\@platerefvii}{\oldpageref{RMAAPlate7}}
\newcommand{\@platerefviii}{\oldpageref{RMAAPlate8}}
\newcommand{\@platerefix}{\oldpageref{RMAAPlate9}}
\newcommand{\@platerefx}{\oldpageref{RMAAPlate10}}
\newcommand{\@platerefxi}{\oldpageref{RMAAPlate11}}
\newcommand{\@platerefxii}{\oldpageref{RMAAPlate12}}
\newcommand{\@platerefxiii}{\oldpageref{RMAAPlate13}}
\newcommand{\@platerefxiv}{\oldpageref{RMAAPlate14}}
\newcommand{\@platerefxv}{\oldpageref{RMAAPlate15}}
\newcounter{figplatei}
\newcounter{figplateii}
\newcounter{figplateiii}
\newcounter{figplateiv}
\newcounter{figplatev}
\newcounter{figplatevi}
\newcounter{figplatevii}
\newcounter{figplateviii}
\newcounter{figplateix}
\newcounter{figplatex}
\newcounter{figplatexi}
\newcounter{figplatexii}
\newcounter{figplatexiii}
\newcounter{figplatexiv}
\newcounter{figplatexv}
%% Printed output should use the counter `plate'
\renewcommand{\theplate}{\@arabic\c@plate}
\renewcommand{\theplateloop}{\@arabic\c@plateloop}
%% Command to output all the saved plates in this paper
\newcommand{\outputplates}{%
  \setcounter{plateloop}{0}
  \if@book\else\@endofarticlehook\fi % not entirely sure why, but this
                                     % seems to be needed
  \clearpage % first output the last page of text
  % now set pagestyle for blank even pages and make sure we start on an odd page
  \pagestyle{empty}
  \cleardoublepage
  \onecolumn
  \parindent\z@
  % first put in an extra sheet
  % the main purpose of this is so there is always a blank even page before 
  % each plate for filling in the PS file if the article has an odd number
  % of pages (see splitbook.sh)
  \null\vfil\centerline{PLATES}\vfil\clearpage\pagestyle{rmaaplate}\cleardoublepage
  % output each plate in turn
  \@whilenum\value{plateloop}<\value{plate}\do{%
    \addtocounter{plateloop}{1}
    % restore the figure counter from \@figcounti, \@figcountii, etc
    \setcounter{figure}{\value{figplate\roman{plateloop}}}
    % spew out the plate (in its own group for safety)
    {\csname @plate\roman{plateloop}\endcsname\par}
    }
  \pagestyle{rmaa}
  % reset the page number as if nothing had happened
  \addtocounter{page}{-\value{plate}}
  \addtocounter{page}{-\value{plate}}
  }
\newcommand\figcaption{%
  \def\@captype{figure}
  \renewcommand{\@makecaption}{\@makefigcap}
  \caption}
\newcommand\tabcaption{%
  \def\@captype{table}
  \renewcommand{\@makecaption}{\@maketabcap}
  \caption}

\newcommand{\plate}[2][1]{%
  %% User command to make a plate page: 
  %% Mandatory argument is what goes on the plate - this should not 
  %% include any floating environments (e.g. figure). Instead, \figcaption 
  %% should be used to do the captions. 
  %% Optional argument is the number of figures (i.e. captions) on the plate
  %% page (default is 1). This is necessary for incrementing the figure counter. 
  %%
  %% Note that we have to use a \gdef rather than an \xdef when we 
  %% save the plate output in a command. This means that in the book styles when it 
  %% comes to spewing out the plate, we will have forgotten about any user macros 
  %% that may have been defined. This can be got round by repeating the macros inside the plate.
  %% (Presumably, one would use a \newcommand rather than \renewcommand if they were new.)
  % increment global and local counters
  \addtocounter{plate}{1}
  % write appropriate label to .aux file
  \csname @platelabel\roman{plate}\endcsname
  % define a global macro called \@platei, \@plateii, \@plateiii, \@plateiv, ...
  \expandafter\gdef\csname @plate\roman{plate}\endcsname{%
    \csname @outputplatelabel\roman{plateloop}\endcsname
    \null\vfil #2 \vfil
    \cleardoublepage
    }
  % save the current shortauthor and shorttitle
  % need to be \xdef to take current values of \@shortauthor, etc 
  \expandafter\xdef\csname @shorta\roman{plate}\endcsname{\@shortauthor}
  \expandafter\xdef\csname @shortt\roman{plate}\endcsname{\@shorttitle}
  % save the current value of the figure counter for later use in \outputplates
  \setcounter{figplate\roman{plate}}{\value{figure}}
  % increment the figure counter according to the optional argument
  \addtocounter{figure}{#1}
  }
%% page style for plate pages 
\newcommand{\ps@rmaaplate}{%
  % Plate X (see page Y) in the foot
  \renewcommand\@oddfoot{\hfil\csname @shorta\roman{plateloop}\endcsname\ 
    (see page~\csname @plateref\roman{plateloop}\endcsname)\hfil}%
  \renewcommand\@evenfoot{\relax}%
  % nothing in even heads (blank pages)
  \renewcommand\@evenhead{\relax}
  % short title and page num in the odd heads
  \renewcommand\@oddhead{\hfil\csname @shortt\roman{plateloop}\endcsname%
    \hfil\llap{PLATE~\theplateloop}}
  }





%% Things for the User Guide
\if@guide
%% stolen from doc.sty
\def\MakeShortVerb#1{%
  \expandafter\ifx\csname cc\string#1\endcsname\relax
  \@shortvrbinfo{Made }{#1}%
  \add@special{#1}%
  \expandafter
  \xdef\csname cc\string#1\endcsname{\the\catcode`#1}%
  \begingroup
  \catcode`\~\active  \lccode`\~`#1%
  \lowercase{%
    \global\expandafter\let
    \csname ac\string#1\endcsname~%
    \gdef~{\verb~}}%
  \endgroup
  \global\catcode`#1\active
  \else
  \@shortvrbinfo\@empty{#1 already}%
  \fi}
\def\DeleteShortVerb#1{%
  \expandafter\ifx\csname cc\string#1\endcsname\relax
  \else
  \@shortvrbinfo{Deleted }{#1 as}%
  \rem@special{#1}%
  \global\catcode`#1\csname cc\string#1\endcsname
  \global \expandafter\let \csname cc\string#1\endcsname \relax
  \ifnum\catcode`#1=\active
  \begingroup
  \catcode`\~\active   \lccode`\~`#1%
  \lowercase{%
    \global\expandafter\let\expandafter~%
    \csname ac\string#1\endcsname}%
  \endgroup \fi \fi}
\def\@shortvrbinfo#1#2{%
  \PackageInfo{doc}{%
    #1\expandafter\@gobble\string#2 a short reference
    for \string\verb}}
\def\add@special#1{%
  \rem@special{#1}%
  \expandafter\gdef\expandafter\dospecials\expandafter
  {\dospecials \do #1}%
  \expandafter\gdef\expandafter\@sanitize\expandafter
  {\@sanitize \@makeother #1}}
\def\rem@special#1{%
  \def\do##1{%
    \ifnum`#1=`##1 \else \noexpand\do\noexpand##1\fi}%
  \xdef\dospecials{\dospecials}%
  \begingroup
  \def\@makeother##1{%
    \ifnum`#1=`##1 \else \noexpand\@makeother\noexpand##1\fi}%
  \xdef\@sanitize{\@sanitize}%
  \endgroup}
%% define | as shortcut for \verbatim
\MakeShortVerb{\|}
%% comand to specify arguments to LaTeX commands
\newcommand{\rmArg}[1]{%
  \ensuremath{\langle}{\scriptsize \textsf{\MakeUpperCase{#1}}}\ensuremath{\rangle}}
%\newcommand{\rmArg}[1]{\ensuremath{\langle}\textsc{#1}\ensuremath{\rangle}}
\providecommand{\BibTeX}{{\rmfamily B\kern-.05em%
    \textsc{i\kern-.025em b}\kern-.08em%
    T\kern-.1667em\lower.7ex\hbox{E}\kern-.125emX}}
%% Environment for code examples
\newenvironment{Code}{%
  \begin{list}{}{\setlength{\leftmargin}{0.05\columnwidth}} 
  \item[]}{\end{list}}
\fi
%</package>
%    \end{macrocode}

% \section{User commands for astronomical articles}
% \label{sec:user-commands}
% Some random definitions. These go both into the main class file and
% also into |adshtml.sty|.
% 
% Definitions of the journal name in different incarnations.  
%    \begin{macrocode}
%<*package|adshtml-style>
\newcommand{\RMAA}{RMxAA}
\newcommand{\RMAAlong}{Revista Mexicana de Astronom\'{\i}a y Astrof\'{\i}sica}
\newcommand{\RMAASC}{RMxAC}
\newcommand{\RMAASClong}{\RMAA\ (Serie de Conferencias)}
%    \end{macrocode}
% \begin{macro}{\ibidrule}
% This puts a horizontal rule with a period after it. It is used when
% the author list of a reference is identical to the preceding one. 
%    \begin{macrocode}
\newcommand{\ibidrule}{\rule[-0.03cm]{4.7em}{.01cm}.\ }
%    \end{macrocode}
% \end{macro}
% \begin{macro}{\ThisVolume}
% References to other articles in the same volume can be coded as,
% e.g., |\bibitem{} Bloggs, J. \ThisVolume{bloggs}|.
%    \begin{macrocode}
\newcommand{\ThisVolume}[1]{%
  \@year, \if@proceedings\RMAASC\else\RMAA\fi, \@volume,
  \if@book\oldpageref{#1:RMAAFirstPage}\else 000\fi\ (this volume)}
%    \end{macrocode}
% \end{macro}
% Other stuff that may be useful in ref lists.
%    \begin{macrocode}
\newcommand{\etal}{et al.\@}
%    \end{macrocode}
% Various new commands from AAS macros
%    \begin{macrocode}
\def\acknowledgments{\vskip 3ex plus .8ex minus .4ex}
\let\acknowledgements=\acknowledgments                  % second spelling
\let\jnl@style=\rm
\def\ref@jnl#1{{\jnl@style#1}}
\DeclareRobustCommand\aj{\ref@jnl{AJ}}                   
\DeclareRobustCommand\araa{\ref@jnl{ARA\&A}}             
\DeclareRobustCommand\apj{\ref@jnl{ApJ}}                 
\DeclareRobustCommand\apjl{\ref@jnl{ApJ}}                
\DeclareRobustCommand\apjs{\ref@jnl{ApJS}}               
\DeclareRobustCommand\ao{\ref@jnl{Appl.~Opt.}}           
\DeclareRobustCommand\apss{\ref@jnl{Ap\&SS}}             
\DeclareRobustCommand\aap{\ref@jnl{A\&A}}                
\DeclareRobustCommand\aapr{\ref@jnl{A\&A~Rev.}}          
\DeclareRobustCommand\aaps{\ref@jnl{A\&AS}}              
\DeclareRobustCommand\azh{\ref@jnl{AZh}}                 
\DeclareRobustCommand\baas{\ref@jnl{BAAS}}               
\DeclareRobustCommand\jrasc{\ref@jnl{JRASC}}             
\DeclareRobustCommand\memras{\ref@jnl{MmRAS}}            
\DeclareRobustCommand\mnras{\ref@jnl{MNRAS}}             
\DeclareRobustCommand\pra{\ref@jnl{Phys.~Rev.~A}}        
\DeclareRobustCommand\prb{\ref@jnl{Phys.~Rev.~B}}        
\DeclareRobustCommand\prc{\ref@jnl{Phys.~Rev.~C}}        
\DeclareRobustCommand\prd{\ref@jnl{Phys.~Rev.~D}}        
\DeclareRobustCommand\pre{\ref@jnl{Phys.~Rev.~E}}        
\DeclareRobustCommand\prl{\ref@jnl{Phys.~Rev.~Lett.}}    
\DeclareRobustCommand\pasp{\ref@jnl{PASP}}               
\DeclareRobustCommand\pasj{\ref@jnl{PASJ}}               
\DeclareRobustCommand\qjras{\ref@jnl{QJRAS}}             
\DeclareRobustCommand\skytel{\ref@jnl{S\&T}}             
\DeclareRobustCommand\solphys{\ref@jnl{Sol.~Phys.}}      
\DeclareRobustCommand\sovast{\ref@jnl{Soviet~Ast.}}      
\DeclareRobustCommand\ssr{\ref@jnl{Space~Sci.~Rev.}}     
\DeclareRobustCommand\zap{\ref@jnl{ZAp}}                 
\DeclareRobustCommand\nat{\ref@jnl{Nature}}              
\DeclareRobustCommand\iaucirc{\ref@jnl{IAU~Circ.}}
\DeclareRobustCommand\aplett{\ref@jnl{Astrophys.~Lett.}}
\DeclareRobustCommand\apspr{\ref@jnl{Astrophys.~Space~Phys.~Res.}}
\DeclareRobustCommand\bain{\ref@jnl{Bull.~Astron.~Inst.~Netherlands}}
\DeclareRobustCommand\fcp{\ref@jnl{Fund.~Cosmic~Phys.}}
\DeclareRobustCommand\gca{\ref@jnl{Geochim.~Cosmochim.~Acta}}
\DeclareRobustCommand\grl{\ref@jnl{Geophys.~Res.~Lett.}}
\DeclareRobustCommand\jcp{\ref@jnl{J.~Chem.~Phys.}}      
\DeclareRobustCommand\jgr{\ref@jnl{J.~Geophys.~Res.}}    
\DeclareRobustCommand\jqsrt{\ref@jnl{J.~Quant.~Spec.~Radiat.~Transf.}}
\DeclareRobustCommand\memsai{\ref@jnl{Mem.~Soc.~Astron.~Italiana}}
\DeclareRobustCommand\nphysa{\ref@jnl{Nucl.~Phys.~A}}
\DeclareRobustCommand\physrep{\ref@jnl{Phys.~Rep.}}
\DeclareRobustCommand\physscr{\ref@jnl{Phys.~Scr}}
\DeclareRobustCommand\planss{\ref@jnl{Planet.~Space~Sci.}}
\DeclareRobustCommand\procspie{\ref@jnl{Proc.~SPIE}}
\let\astap=\aap
\let\apjlett=\apjl
\let\apjsupp=\apjs
\let\applopt=\ao
\def\phn{\phantom{0}}
\def\phd{\phantom{.}}
\def\phs{\phantom{\ensuremath{-}}}
\def\phm#1{\phantom{#1}}
%    \end{macrocode}
% All these are redone as robust and using more modern commands where
% possible. The robustness is so they will not get expanded if used in
% abstracts and titles when these are passed to the ADS file. 
% \changes{3.16}{2003/04/07}{
%  Added \cs{Color} command for doing photometric colors.}
%    \begin{macrocode}
%%\def\sun{\hbox{$\odot$}}
\DeclareRobustCommand\sun{\ensuremath{\odot}}
%%\def\earth{\hbox{$\oplus$}}
\DeclareRobustCommand\earth{\ensuremath{\oplus}}
%%\def\lesssim{\mathrel{\hbox{\rlap{\hbox{\lower4pt\hbox{$\sim$}}}\hbox{$<$}}}}
%%\def\gtrsim{\mathrel{\hbox{\rlap{\hbox{\lower4pt\hbox{$\sim$}}}\hbox{$>$}}}}
\DeclareRobustCommand\lesssim{\mathrel{\hbox{\smaller\smaller\rlap{\hbox{\lower3pt\hbox{$\sim$}}}\raise1.5pt\hbox{$<$}}}}
\DeclareRobustCommand\gtrsim{\mathrel{\hbox{\smaller\smaller\rlap{\hbox{\lower3pt\hbox{$\sim$}}}\raise1.5pt\hbox{$>$}}}}
%%\def\lta{\mathrel{\spose{\lower 3pt\hbox{$\mathchar"218$}}
%%     \raise 2.0pt\hbox{$\mathchar"13C$}}}
%%\def\gta{\mathrel{\spose{\lower 3pt\hbox{$\mathchar"218$}}
%%     \raise 2.0pt\hbox{$\mathchar"13E$}}}
\DeclareRobustCommand\sq{\hbox{\rlap{$\sqcap$}$\sqcup$}}
%%\def\arcdeg{\hbox{$^\circ$}}
%%\def\arcmin{\hbox{$^\prime$}}
%%\def\arcsec{\hbox{$^{\prime\prime}$}}
\DeclareRobustCommand\arcdeg{\ensuremath{^\circ}}
\DeclareRobustCommand\arcmin{\ensuremath{^\prime}}
\DeclareRobustCommand\arcsec{\ensuremath{^{\prime\prime}}}
%%\def\fd{\hbox{$.\!\!^{\rm d}$}}
%%\def\fh{\hbox{$.\!\!^{\rm h}$}}
%%\def\fm{\hbox{$.\!\!^{\rm m}$}}
%%\def\fs{\hbox{$.\!\!^{\rm s}$}}
%%\def\fdg{\hbox{$.\!\!^\circ$}}
%%\def\farcm{\hbox{$.\mkern-4mu^\prime$}}
%%\def\farcs{\hbox{$.\!\!^{\prime\prime}$}}
%%\def\fp{\hbox{$.\!\!^{\scriptscriptstyle\rm p}$}}
\DeclareRobustCommand\fd{\ensuremath{.\!\!^\mathrm{d}}}
\DeclareRobustCommand\fh{\ensuremath{.\!\!^\mathrm{h}}}
\DeclareRobustCommand\fm{\ensuremath{.\!\!^\mathrm{m}}}
\DeclareRobustCommand\fs{\ensuremath{\mkern0.5mu.\mkern-5mu^\mathrm{s}\mkern-1mu}}
\DeclareRobustCommand\fdg{\ensuremath{.\!\!^\circ}}
\DeclareRobustCommand\farcm{\ensuremath{.\mkern-4mu^\prime}}
\DeclareRobustCommand\farcs{\ensuremath{\mkern1mu.\mkern-6mu^{\prime\prime}\mkern-1.5mu}}
%%\DeclareRobustCommand\farcs{\ensuremath{.\!\!^{\prime\prime}}}
\DeclareRobustCommand\fp{\ensuremath{.\!\!^{\scriptscriptstyle\rm p}}}
%%\def\micron{\hbox{$\mu$m}}
\DeclareRobustCommand\micron{\ensuremath{\mu\mathrm{m}}}
\let\la=\lesssim                        
\let\ga=\gtrsim
%%\def\case#1#2{\hbox{$\frac{#1}{#2}$}}
%%\def\slantfrac#1#2{\hbox{$\,^#1\!/_#2$}}
%%\def\onehalf{\slantfrac{1}{2}}
%%\def\onethird{\slantfrac{1}{3}}
%%\def\twothirds{\slantfrac{2}{3}}
%%\def\onequarter{\slantfrac{1}{4}}
%%\def\threequarters{\slantfrac{3}{4}}
\DeclareRobustCommand\case[2]{\ensuremath{\frac{#1}{#2}}}
\DeclareRobustCommand\slantfrac[2]{\ensuremath{\,^#1\!/_#2}}
\DeclareRobustCommand\onehalf{\slantfrac{1}{2}}
\DeclareRobustCommand\onethird{\slantfrac{1}{3}}
\DeclareRobustCommand\twothirds{\slantfrac{2}{3}}
\DeclareRobustCommand\onequarter{\slantfrac{1}{4}}
\DeclareRobustCommand\threequarters{\slantfrac{3}{4}}
%%\def\ubvr{\hbox{$U\!BV\!R$}}            
%%\def\ub{\hbox{$U\!-\!B$}}               
%%\def\bv{\hbox{$B\!-\!V$}}               
%%\def\vr{\hbox{$V\!-\!R$}}               
%%\def\ur{\hbox{$U\!-\!R$}}
%    \end{macrocode}
% New commands for ``photometric'' colors 
%    \begin{macrocode}
\DeclareRobustCommand\ubvr{\ensuremath{U\!BV\!R}}            
\DeclareRobustCommand\Color[2]{%
  \ensuremath{%
    #1\!\ifx#1V\!\fi-\!#2
  }}
\DeclareRobustCommand\ub{\Color{U}{B}}
\DeclareRobustCommand\bv{\Color{B}{V}}
\DeclareRobustCommand\vr{\Color{V}{R}}
\DeclareRobustCommand\ur{\Color{U}{R}}
%% New version of the \ion command, that actually works! 
%    \end{macrocode}
% \begin{macro}{\ion}
% This version of the |\ion| command allows you to write either
% |\ion{N}{4}| or |\ion{N}{iv}| in order to get N~IV. 
%    \begin{macrocode}
\DeclareRobustCommand\ion[2]{%
  \mbox{#1\kern0.2em%
    \smaller\rmfamily%
    \edef\@tempa{\@car#2\@nil}%
    \ifcat1\@tempa%
    \@Roman{#2}%
    \else%
    \uppercase{#2}%
    \fi}}
\newcounter{IonStage}
\DeclareRobustCommand\plainion[2]{\setcounter{IonStage}{#2}#1 \Roman{IonStage}}
%%\def\ion#1#2{\ensuremath{\mathrm{#1}\;}{\protect\small\rm\@Roman{#2}}}
%    \end{macrocode}
% \end{macro}
%    \begin{macrocode}
\DeclareRobustCommand\nodata{$\cdots$}
%    \end{macrocode}
% These are needed by the |\ion| command. Isn't there an easier way of
% doing this?
%    \begin{macrocode}
\let\@small=\small
\let\@normalsize=\normalsize
\let\@footnotesize=\footnotesize
\let\@scriptsize=\scriptsize
\let\@tiny=\tiny
\let\@large=\large
\let\@Large=\Large
\let\@LARGE=\LARGE
\let\@huge=\huge
\let\@Huge=\Huge
\newcommand{\smaller}{\small}
\renewcommand{\normalsize}{\renewcommand{\smaller}{\small}\@normalsize}
\renewcommand{\small}{\renewcommand{\smaller}{\footnotesize}\@small}
\renewcommand{\footnotesize}{\renewcommand{\smaller}{\scriptsize}\@footnotesize}
\renewcommand{\scriptsize}{\renewcommand{\smaller}{\tiny}\@scriptsize}
\renewcommand{\tiny}{\renewcommand{\smaller}{\tiny}\@tiny}
\renewcommand{\large}{\renewcommand{\smaller}{\normalsize}\@large}  
\renewcommand{\Large}{\renewcommand{\smaller}{\large}\@Large}  
\renewcommand{\LARGE}{\renewcommand{\smaller}{\Large}\@LARGE}  
\renewcommand{\huge}{\renewcommand{\smaller}{\LARGE}\@huge}  
\renewcommand{\Huge}{\renewcommand{\smaller}{\huge}\@Huge}  
%</package|adshtml-style>
%    \end{macrocode}


% \section{Other files}

% \subsection{The driver and make files}
% \label{sec:driver_make}
% These are put here for completeness, although they still need to be
% distributed separately since otherwise the user won't be able to
% unpack this file into its constituent parts.
% 
% The |Makefile| has two principal rules: |docs|, and |macros|. The
% |docs| rule produces a PDF and gzipped PS file of the documentation.
% I still haven't perfected the combination of rules required to get
% it to run \LaTeX{} enough times in all circumstances. The |macros|
% rule runs \LaTeX{} on |rmaa.ins| in order to unpack |rmaa.cls| and
% all the other little files (including the |Makefile| and the
% |rmaa.ins| files themselves\dots 
% 
% First set up lists of various types of file. 
%    \begin{macrocode}
%<*makefile>
DOCFILES=rmaa.pdf
MACRO_LTX_FILES=rmaa.cls rmaacite.sty adshtml.sty
SHELLSCRIPTS=makeads.sh
MACRO_OTHER_FILES=adshtml.perl sample.latex2html-init
MAKEFILE=Makefile
MACROFILES=$(MACRO_LTX_FILES) $(SHELLSCRIPTS) $(MACRO_OTHER_FILES)         
SKELETONFILES=ADS.tex 
%    \end{macrocode}
% Install the files read only into the package directory, specified via a path relative to the source directory. \changes{4.0}{2019/11/14}{Reorganize folder layout so that hardly anything is in the top level.  That way, things are clearer on github.}
%    \begin{macrocode}
INSTALL=install 
INSTALL_DIR=../packages
%    \end{macrocode}
% The default rule is to make everything. Installation must be
% explicitly asked for.
%    \begin{macrocode}
all: docs macros

install: install-docs install-macros

%    \end{macrocode}
% Most files are installed read-only to prevent editing. Skeleton
% files must be writeable. Must be careful not to overwrite the
% Makefile in the parent directory. 
%    \begin{macrocode}
install-macros: macros
        $(INSTALL) -m 444 $(MACRO_LTX_FILES) $(MACRO_OTHER_FILES) $(INSTALL_DIR)
        $(INSTALL) -m 555 $(SHELLSCRIPTS) $(INSTALL_DIR)
        $(INSTALL) -m 644 $(SKELETONFILES) $(INSTALL_DIR)

install-docs: docs
        $(INSTALL) rmaa.pdf $(INSTALL_DIR)/rm-fulldocs.pdf

%    \end{macrocode}
% Make the PDF and gzipped PS files. 
%    \begin{macrocode}
docs: $(DOCFILES)

%    \end{macrocode}
% Make everything else if either the |.dtx| or |.ins| files have
% changed. 
%    \begin{macrocode}
macros: rmaa.dtx rmaa.ins
%    \end{macrocode}
% Ensure that all the macro files can be written to then recreate
% them. The |touch| is there in case any of the files don't exist
% yet. 
%    \begin{macrocode}
        touch $(MACROFILES) $(MAKEFILE)
        chmod u+w $(MACROFILES) $(MAKEFILE) rmaa.ins
        latex "\input{rmaa.ins}"
%    \end{macrocode}
% \TeX{} automatically adds the |.tex| extension to extensionless
% files. We fix this and take the opportunity to convert
% back to TABs all that it has lovingly converted to spaces. 
%    \begin{macrocode}
        cp -f $(MAKEFILE) $(MAKEFILE).BCK && unexpand $(MAKEFILE).tex >| $(MAKEFILE) && rm $(MAKEFILE).tex
%    \end{macrocode}
% Make the shell scripts executable and everything unwriteable (so
% they don't get hand-edited by mistake). However, undo this for the
% files that the user is \emph{supposed} to modify. 
%    \begin{macrocode}
        chmod a+x $(SHELLSCRIPTS)
        chmod a-w $(MACROFILES) $(MAKEFILE)
        chmod u+w $(SKELETONFILES)

%    \end{macrocode}
% We now use PDF\LaTeX{} for this.
%    \begin{macrocode}
rmaa.idx: rmaa.dtx
        @echo "REGENERATING .idx FILE"
        pdflatex "\nonstopmode\input{rmaa.dtx}"

rmaa.aux: rmaa.dtx rmaa.ind rmaa.gls
        @echo "REGENERATING .aux FILE"
        pdflatex "\nonstopmode\input{rmaa.dtx}"

rmaa.ind: rmaa.idx
        @echo "REGENERATING INDEX"
        makeindex -s gind.ist rmaa

rmaa.gls: rmaa.glo
        @echo "REGENERATING CHANGELOG"
        makeindex -s gglo.ist -o rmaa.gls rmaa.glo

rmaa.pdf: rmaa.aux rmaa.ind rmaa.gls
        @echo "REGENERATING .pdf FILE"
        pdflatex "\nonstopmode\input{rmaa.dtx}"

clean: 
        rm -f rmaa.aux rmaa.dvi rmaa.ind rmaa.gls
%</makefile>
%    \end{macrocode}

% The |rmaa.ins| file uses the \textsf{docstrip} package to split out
% all the files and strip off this documentation.
%    \begin{macrocode}
%<*ins>
%%
%% Installation batch file for the RMAA LaTeX package
%%
\def\batchfile{rmaa.ins}
\input docstrip.tex
\preamble

PLEASE DO NOT EDIT THIS FILE! It is generated automatically from
the file rmaa.dtx using the accompanying Makefile and rmaa.ins 
Any hand-edited changes will be lost

A LaTeX document class for the RMxAA

Author: William Henney (w.henney@irya.unam.mx)

Latest version available from: http://www.irya.unam.mx/rmaa

Copyright (C) 1999-2019 Instituto de Astronomia, UNAM, Mexico

Permission is hereby granted, free of charge, to any person obtaining
a copy of this software and associated documentation files (the
"Software"), to deal in the Software without restriction, including
without limitation the rights to use, copy, modify, merge, publish,
distribute, sublicense, and/or sell copies of the Software, and to
permit persons to whom the Software is furnished to do so, subject to
the following conditions:

The above copyright notice and this permission notice shall be
included in all copies or substantial portions of the Software.

THE SOFTWARE IS PROVIDED "AS IS", WITHOUT WARRANTY OF ANY KIND,
EXPRESS OR IMPLIED, INCLUDING BUT NOT LIMITED TO THE WARRANTIES OF
MERCHANTABILITY, FITNESS FOR A PARTICULAR PURPOSE AND
NONINFRINGEMENT. IN NO EVENT SHALL THE AUTHORS OR COPYRIGHT HOLDERS BE
LIABLE FOR ANY CLAIM, DAMAGES OR OTHER LIABILITY, WHETHER IN AN ACTION
OF CONTRACT, TORT OR OTHERWISE, ARISING FROM, OUT OF OR IN CONNECTION
WITH THE SOFTWARE OR THE USE OR OTHER DEALINGS IN THE SOFTWARE.
\endpreamble
%    \end{macrocode}
% Clobber files without asking
%    \begin{macrocode}
\askforoverwritefalse
%    \end{macrocode}
% Make the \LaTeX{} macro files (class and styles)
%    \begin{macrocode}
\generate{% 
  \file{rmaa.cls}{\from{rmaa.dtx}{package}}
  \file{rmaa.ins}{\from{rmaa.dtx}{ins}}
  \file{rmaacite.sty}{\from{rmaa.dtx}{rmaacite}}
  \file{adshtml.sty}{\from{rmaa.dtx}{adshtml-style}}
  \file{ADS.tex}{\from{rmaa.dtx}{ads-container}}
  }
%    \end{macrocode}
% Make the other non-\LaTeX{} files. These are done separately since
% they can't use `\%\%' as a comment character for the
% preamble/postamble (and Perl files can't have the postamble I don't
% think). 
%    \begin{macrocode}
\generate{\nopreamble\nopostamble% 
  \file{Makefile}{\from{rmaa.dtx}{makefile}}
  \file{adshtml.perl}{\from{rmaa.dtx}{adshtml-perl}}
  \file{makeads.sh}{\from{rmaa.dtx}{ads-shell-script}}
  \file{sample.latex2html-init}{\from{rmaa.dtx}{latex2html-init}}
  }

%% Local Variables: 
%% mode: LaTeX
%% mode: Font-lock
%% End:
%</ins>
%    \end{macrocode}

% \subsection{Extended citation syntax}
% Modified from a file called mnrascite.sty
% I no longer remember where I stole the original version of this from
% 
% Alterations to the cite commands to make them do more sensible things.
% N.B. the bibliography style should produce entries which have labels where
% the date is delimited by angle brackets '<>' for these macros to work
% properly. They will not give very sensible error messages if you do not do 
% this properly.
% 
% A new command |\scite{key}| is defined which produces a citation which is 
% suitable for insertion in a sentence. i.e., it does not put brackets around 
% the whole citation, only the date. Also |\pcite{key}| produce output with no
% parentheses whatsoever. 
% 
% Should redefine |\@bibitem| really to have any hope of compatibility
% the command |\lbibcite| is defined to deal with the citations.
%    \begin{macrocode}
%<*rmaacite>
\def\bibitem[#1<#2>]#3{\par\hangindent=1.5em\hangafter=1
  \if@filesw
  {\def\protect##1{\string ##1\space}
    \immediate\write\@auxout{\string\lbibcite{#3}{#1}{#2}}}
  \fi}
\@redefinebibitem % what it says
%    \end{macrocode}
% Note that the |\lbibcite| defines two types of citation one contains the names,
% the other the year
%    \begin{macrocode}
\def\lbibcite#1#2#3{\global\@namedef{b@#1}{#2}
                    \global\@namedef{by@#1}{#3}}

\def\@bcite#1#2{({#1\if@tempswa , #2\fi})}
\def\@pcite#1#2{#1\if@tempswa , #2\fi}
%    \end{macrocode}
% This declaration contains what is put between the name and the date
% \changes{3.11}{2003/01/31}{Moved the \cs{articlename} to after the key for consistency
% with the new natbib stuff.}
%    \begin{macrocode}
\def\@brcite{\ (}
\def\@ebrcite{)}
\def\@fmttop#1{\@fmtcite[\csname #1\endcsname]}
\def\@citex[#1]#2{%
%%  \if@filesw\immediate\write\@auxout{\string\citation{#2}}\fi
%% WJH added 16 Jan 2002
  \@for\@citeb:=#2\do{%
    \edef\@citeb{\expandafter\@firstofone\@citeb:\articlename}%
    \if@filesw\immediate\write\@auxout{\string\citation{\@citeb}}\fi}%
%% to here
  \def\@citea{}%
  \@cite{\@for\@citeb:=#2\do
    {%
      \edef\@citeb{\expandafter\@firstofone\@citeb:\articlename}%
      \@citea\def\@citea{;\ }\@ifundefined
      {b@\@citeb}{{\bf ?}\@warning {Citation `\@citeb' on page \thepage \space undefined}}%
      {\csname b@\@citeb\endcsname\@citebro\csname by@\@citeb\endcsname\@citebrc}}}%
  {#1}}

\def\cite{\@ifnextchar [{\let\@citebro=\space
                          \let\@citebrc=\relax
                          \let\@cite=\@bcite\@tempswatrue\@citex}
                        {\let\@citebro=\space
                          \let\@citebrc=\relax
                          \let\@cite=\@bcite\@tempswafalse\@citex[]}}
\def\pcite{\@ifnextchar [{\let\@citebro=\space
                          \let\@citebrc=\relax
                          \let\@cite=\@pcite\@tempswatrue\@citex}
                        {\let\@citebro=\space
                          \let\@citebrc=\relax
                          \let\@cite=\@pcite\@tempswafalse\@citex[]}}
\def\scite{\@ifnextchar [{\let\@citebro=\@brcite
                          \let\@citebrc=\@ebrcite
                          \let\@cite=\@pcite\@tempswatrue\@citex}
                        {\let\@citebro=\@brcite
                          \let\@citebrc=\@ebrcite
                          \let\@cite=\@pcite\@tempswafalse\@citex[]}}

\renewenvironment{thebibliography}
  {\section*{REFERENCES}\vspace*{\afterbiblioskip}\small
    \if@book\RereadBiblio\fi
    \sloppy\clubpenalty4000\widowpenalty4000
    \sfcode`\.=1000\relax
    \setlength\parindent{\z@}\setlength\parskip{\z@}}
  {\relax\normalsize\@outputaddressesifnecessary} 
%</rmaacite>
%    \end{macrocode}

% \subsection{Auxilliary files for generating ADS data}

% The ADS data are written by \LaTeX{} to a |.ads| file. This file
% then has to be included in the following file (|ADS.tex|),
% which is run through \LaTeX{}2HTML (see below). 

%    \begin{macrocode}
%<*ads-container>
%%
%% Instructions:
%% Edit the \input{} line below to reflect the name of your file. 
\documentclass{article}
\usepackage[latin1]{inputenc}
\usepackage{adshtml}
\newcommand{\ADSfield}[2]
{\%#1  #2\\}
%%\newcommand{\ADSvfield}[2]
%%{\%#1 \begin{verbatim}#2\end{verbatim}}
\newenvironment{ADSarticle}{}{\par\bigskip}
\renewcommand{\bibitem}[2][]{\par}
\setlength{\parindent}{0pt}
\renewcommand{\ThisVolume}[1]{%  
  2001, \RMAASC, 12, \oldpageref{#1:RMAAFirstPage}}  
 

\begin{document}
\makeatletter
\input{rmbooksamp.ads}
\end{document}
%</ads-container>
%    \end{macrocode}

% This is the |adshtml.sty| style file used by |ADS.tex|. It just
% defines a few 
% macros that might be used in abstracts, etc (really it should define
% them all). This now gets all its code from
% Sec~\ref{sec:user-commands}. 
%    \begin{macrocode}
%<*adshtml-style>
%% Nothing else needs to go here now. 
%</adshtml-style>
%    \end{macrocode}
% The init file for latex2html.
%    \begin{macrocode}
%<*latex2html-init>
$PDFTOCAIRO = "pdftocairo";
$PDFLATEX = "pdflatex";
$LATEX = "latex";
$MAX_SPLIT_DEPTH = 0;
$HTML_VERSION = "4.0,math";
$NO_SIMPLE_MATH = 1;
$ADDRESS = "<p><address><a href=\"mailto:rmaa\@astro.unam.mx\">rmaa\@astro.unam.mx</address></p>"; 
$INFO = "<p>This document was created using latex2html and the rmaa latex macros</p>"; 
$NO_NAVIGATION = 1;
1; 
%</latex2html-init>
%    \end{macrocode}
% This is the |adshtml.perl| file that is used by \LaTeX{}2HTML. Some
% bizarre \TeX{}y command sometimes show up so we tell it to ignore
% these. All the astronomy macros that come from AAS are processed in
% \TeX{}. That way they get made to images but with an |ALT| tag that
% has the original code, which is what will be picked up by |lynx|,
% with the result that they get passed intact into the final |.txt|
% file. 
%    \begin{macrocode}
%<*adshtml-perl>
package main;

$dum = &do_cmd_makeatletter();

sub do_cmd_RMAA {
    local($_) = @_;
    "RMxAA" . $_;
}

sub do_cmd_RMAASC {
    local($_) = @_;
    "RMxAASC" . $_;
}

sub do_cmd_etal {
    local($_) = @_;
    "et al\." . $_;
}

sub do_cmd_ibidrule {
    local($_) = @_;
    "<img src=\"foo\" alt=\"\ibidrule\">" . $_;
}


sub do_cmd_bibitem {
    local($_) = @_;
    local($key);
    local($opt, $dummy)=&get_next_optional_argument; 
    $key = &missing_braces unless
        ((s/$next_pair_pr_rx/$key = $2; ''/eo)
        ||(s/$next_pair_rx/$key = $2; ''/eo));
    "<p>" . $_; 
}

&ignore_commands (<<_RAW_ARG_CMDS_);
lowercase
adjustfinalcols
textsuperscript # {}
newline # []
thinmuskip
mskip
smallskip
medskip
bigskip
hspace # {}
hspacestar # {}
vspace # {}
vspacestar # {}
mathrm
_at_mathrm
spacefactor
discretionary
unhbox 
voidb_at_x 
penalty 
_at_M
_at_m
null
IeC
_RAW_ARG_CMDS_

&process_commands_in_tex (<<_RAW_ARG_CMDS_);
sun
earth
lesssim
gtrsim
sq
argdeg
arcmin
arcsec
fd
fh
fm
fs
fdg
farcm
farcs
fp
micron
la
ga
case # {} # {}
slantfrac # {} # {}
onehalf
onethird
twothirds
onequarter
threequarters
ubvr
ub
bv
vr
ur
aj
araa
apj
apjl
apjs
ao
apss
aap
aapr
aaps
azh
baas
jrasc
memras
mnras
pra
prb
prc
prd
pre
prl
pasp
pasj
qjras
skytel
solphys
sovast
ssr
zap
nat
iaucirc
aplett
apspr
bain
fcp
gca
grl
jcp
jgr
jqsrt
memsai
nphysa
physrep
physscr
planss
procspie
_RAW_ARG_CMDS_

1;                              # This must be the last line
%</adshtml-perl> 
%    \end{macrocode}

% This is a shell script |makeads.sh| that generates the |_ads.txt| file from
% |_ads.tex| by using \LaTeX{}2HTML and |lynx|.  \changes{4.0}{2019/11/13}{Make sure that lynx writes utf-8 output instead of latin-1.}
% \changes{2.2beta13}{2002/03/12}{Complete rewrite of |makeads.sh|. Now much 
% improved I think}
% \changes{4.0}{2019/11/13}{Fix the at-sign substitution so that it only eliminates \cs{@} and does not mess with email addresses.}

%    \begin{macrocode}
%<*ads-shell-script>
#! /bin/sh
if [ -z $1 ]; then
    echo "Usage: $0 basename"
    exit
fi
BASENAME=$1
ADSBASENAME=${BASENAME}_ads
AUXFILE=$BASENAME.aux
TMPFILE=tmp.$$


echo "Generating $ADSBASENAME.tex"

# For some reason, must double-up \\ to \\\\ but not \ to \\ - odd!
cat > $ADSBASENAME.tex <<EOF
%% PLEASE DO NOT EDIT THIS FILE!  It is generated automatically from
%% the script makeads.sh.  Any hand-edited changes will be lost
%% 
\documentclass{article}
\usepackage[latin1]{inputenc}
\usepackage{html}
\usepackage{adshtml}
\newcommand{\ADSfield}[2]{\%#1  #2\\\\}
\newenvironment{ADSarticle}{}{\par\bigskip}
\setlength{\parindent}{0pt}
\newcounter{istage}
\renewcommand{\ion}[2]{\setcounter{istage}{#2}#1 \Roman{istage}}
\begin{document}
\makeatletter
EOF
echo >> $ADSBASENAME.tex "\input{$BASENAME.ads2}" 
echo >> $ADSBASENAME.tex "\end{document}"



echo "Generating $ADSBASENAME.sed" 

# Find type of journal
grep '\\documentclass\[.*proceedings.*\]' $BASENAME.tex >/dev/null \
     && JOURNAL='RMxAC' \
     || JOURNAL='RMxAA'

# Find the year and volume number
YEAR=`grep '\\SetYear' $BASENAME.tex | sed 's/.*\\SetYear{\(.*\)}.*/\1/'`
VOLUME=`grep '\\SetVolume' $BASENAME.tex | sed 's/.*\\SetVolume{\(.*\)}.*/\1/'`

# Generate an auxilliary sed script 
cat > $TMPFILE <<EOF
s|\\\\newlabel{\(.*\):RMAAFirstPage}{{.*}{\(.*\)}}|s/\\\\\\\\ThisVolume{\1}/$YEAR, RMxAC, $VOLUME, \2/g|g
EOF

# Now go through the .aux file, looking for \newlabels and use the sed 
# script just generated in order to make another sed script, one that 
# will replace \ThisVolume{LABEL} with the correct reference
grep '^\\newlabel{.*:RMAAFirstPage}' $AUXFILE \
     | sed -f $TMPFILE > $ADSBASENAME.sed && rm $TMPFILE

# Apend a further sed command to get rid of \@, since redefining 
# it in _ads.tex didn't work
echo >> $ADSBASENAME.sed 's/\\\\@//g'

echo "Generating $BASENAME.ads2 from $ADSBASENAME.sed and $BASENAME.ads" 

# Apply this to the .ads file, generating the .ads2 file, which is included
# by the _ads.tex file.
sed -f $ADSBASENAME.sed $BASENAME.ads > $BASENAME.ads2

echo "Generating $ADSBASENAME.txt from $ADSBASENAME.tex and $BASENAME.ads2"

# Now run latex2html on the _ads.tex file
latex2html -no_navigation -info 0 -address 0 -html_version 4.0,math $ADSBASENAME

# Then, convert to utf8 with lynx and tidy it up a bit
lynx --justify=off --width=100 --dump --nolist --display_charset=utf-8 \
     $ADSBASENAME/$ADSBASENAME.html | cut -c4- \
     | sed 's/ibidrule/\\ibidrule/g' > $ADSBASENAME.txt
%</ads-shell-script>
%    \end{macrocode}

% \section{Epilogue}
% Err, that's it \dots

\endinput


% \endinput
% Local Variables:
% mode: doctex
% TeX-master: t
% End:
"

rmaa.ind: rmaa.idx
        @echo "REGENERATING INDEX"
        makeindex -s gind.ist rmaa

rmaa.gls: rmaa.glo
        @echo "REGENERATING CHANGELOG"
        makeindex -s gglo.ist -o rmaa.gls rmaa.glo

rmaa.pdf: rmaa.aux rmaa.ind rmaa.gls
        @echo "REGENERATING .pdf FILE"
        pdflatex "\nonstopmode% \def\fileversion{4.0} \def\filedate{2019/11/12}
% \MakeShortVerb{\|}
% \title{A \LaTeX{} document class for the Revista Mexicana de
% Astronom\'\i a y Astrof\'\i sica\thanks{This file has version
% \fileversion{} last revised \filedate{}.}}
% \author{William J. Henney\thanks{w.henney@irya.unam.mx}}
% \date{\filedate}
% \maketitle
% \paragraph{Abstract}
% This document describes the \texttt{rmaa} \LaTeX{}
% document class for the journal ``Revista Mexicana de
% Astronom\'\i a y Astrof\'\i sica''. The macros can be used for 
% \begin{enumerate}
% \item Preparing author manuscripts for the journal (including
% conference proceedings to be published in the ``Serie de
% Conferencias'')
% \item Typesetting an entire volume of the journal by inputting the
% individual article files from a container file. Facilities are
% included for automatic generation of TOC and Author Index. 
% \end{enumerate}
% 
% \tableofcontents
% \section{The documentation driver file}
% This is the part that generates the documentation that you may or
% may not be reading, which is done with the help of the Makefile
% discussed in section~\ref{sec:driver_make}. 
%    \begin{macrocode}
%<*driver>
\documentclass{ltxdoc}
\usepackage{hyperref}
\usepackage{palatino}
\EnableCrossrefs \CodelineIndex \RecordChanges
\setlength{\parindent}{0pt}
\def\MacroFont{\footnotesize\tt}
\begin{document} \DocInput{rmaa.dtx} \PrintIndex \PrintChanges \end{document}
%</driver>
%    \end{macrocode}

% \section{The implementation}
% 
% \changes{3.24}{2007/09/14}{Added user access to the macro version
% and date via \cs{rmVersion} and \cs{rmDate}}
% 
%    \begin{macrocode}
%<*package>
\newcommand{\rmaa@version}{3.26}
\newcommand{\rmaa@date}{2011/08/15}
\newcommand\rmVersion{\rmaa@version}
\newcommand\rmDate{\rmaa@date}
\NeedsTeXFormat{LaTeX2e}[1995/06/01]
%    \end{macrocode}
% \changes{4.0}{2019/11/02}{Change abstract to appear before
% resumen. Change RevMexAA to RMxAA for consistency with ADS.}
%
% \changes{3.22}{2007/03/22}{Allow proceedings articles to have
%   received/accepted dates} \changes{3.21}{2006/05/15}{Changes for
%   Silvia - abstracts are now shipped out to ADS.}
% \changes{3.19}{2006/01/19}{Changes for Silvia - remove adresses from
%   end, new abstract style} \changes{2.2beta}{2001/10/30}{ Fixed bug
%   whereby twocolumn option was not passed through to book
%   class. Finally sorted out the ``label on last page'' issue. }
% \changes{2.1}{2001/01/23}{ Fixed some inconsistencies in
%   \cs{nocite} and \cs{bibitem} uniqification Fixed the definition of
%   \cs{ion} so it will work in math mode and in contexts where the
%   typesize is different from normal (e.g., title, fig captions).
%   Modified the definition of \cs{@caption}, so we don't write any
%   list-of-figures data (we never use it anyway, and it just causes
%   problems with fragile commands in fig captions). }
% \changes{2.0}{2000/11/30}{ Added `newstyle' option, which should be
%   used in all new proceedings volumes (wider abstract, smaller
%   typeface in fig captions, etc.) Added commands for formatting
%   table captions. New demonstation document rmpostsamp.tex for
%   1-page posters. This will be the last release under a PRI
%   government (for a while...)}  \changes{1.9}{2000/06/06}{
%   Corrected position of page number in footer of articles that open
%   on even pages. Ensured that first article in each session opens on
%   an odd page, as do list of abstracts and author index. Added extra
%   space in TOC before the backmatter. Added command for editors,
%   \cs{BoldSessionTitles}, that makes produces bold titles for
%   session in the TOC.}  \changes{1.8}{2000/05/16}{ Fixed bug in
%   last page label with \cs{suppressfulladdresses} }
% \changes{1.7}{2000/05/11}{ Fixed more little things. Page labels
%   now work, even if you have no bibliography. You can output the
%   full addresses by hand with the \cs{outputfulladdresses}
%   command. Editors' Guide updated but still not finished. Revamped
%   splitbook.sh now produces PS file of each article for web page and
%   also an EPS file of each page for sending to the printers.}
% \changes{1.6}{2000/04/15}{ More cleaning up: full stops after
%   titles in index, running headers and footers in frontmatter,
%   fontsize in title of index, etc} \changes{1.5}{2000/04/15}{
%   Eliminated superscripts from titles in TOC and index}
% \changes{1.4}{2000/04/14}{ Added copyright and license info (GPL)}
% \changes{1.3}{2000/03/01}{ Corrected \cs{figcaption} to do the
%   caption formatting right Corrected font size of sec headings in
%   front matter} \changes{1.2}{2000/02/16}{ Added \cs{plate}
%   command} \changes{1.1}{2000/02/14}{ Added manuscript option,
%   corrected bugs} \changes{1.0}{2000/02/08}{ Now does a whole book
%   - lots of stuff tidied up} \changes{0.7}{1999/10/04}{ Increased
%   spacing between sections in single column styles. Suppressed
%   floats at bottom of pages (didn't work though) }
% \changes{0.6}{1999/09/24}{ Tweaked \cs{adjustfinalcols} (still not
%   perfect) Added proceedings option Added onecolumn/twocolumn
%   options (orthogonal to journal/proceedings) }
% \changes{0.5}{1999/09/23}{ Added \cs{fulladdresses} and
%   \cs{adjustfinalcols} (just about working)}
% \changes{0.4}{1999/09/22}{ Minor stylistic changes}
% \changes{0.3}{1999/06/28}{ Better mechanism for alternative
%   affiliations Added `guide' option for user guide}
% \changes{0.2}{1999/05/27}{ All basic functionality present, I
%   think} \changes{0.1}{1999/05/26}{ Just about working - still
%   missing Fig/Tab/Refs} \changes{0.0}{1999/05/21}{ So bored today,
%   decided to rewrite RMAA LaTeX macros from scratch}
%    \begin{macrocode}
\ProvidesClass{rmaa}[\rmaa@date\space v\rmaa@version\space RMxAA article
class]
%    \end{macrocode}
% \subsection{Option Handling}
% \label{sec:options}
% Declarations that are needed in options
% \changes{3.05}{2003/10/22}{Added natbib option}
% \changes{3.06}{2003/01/03}{Natbib option turned on permanently}
% \changes{3.07}{2003/01/11}{Added proof option}
% \changes{3.14}{2003/01/11}{Option to turn off natbib}
% \changes{3.18}{2003/10/22}{trying out doctex mode} 
% \changes{3.23}{2007/09/09}{Manuscript is now on by default.}
% \changes{3.26}{2011/08/15}{New option for review articles}
%    \begin{macrocode}
\newif\if@reviewarticle
\@reviewarticlefalse
\newif\if@proof
\@prooffalse
\newif\if@journal
\@journaltrue
\newif\if@printdates
\@printdatesfalse
\newif\if@proceedings
\@proceedingsfalse
\newif\if@manuscript
\@manuscripttrue
\newif\if@bigfont
\@bigfonttrue
\newif\if@guide
\@guidefalse
\newif\if@preprint
\@preprintfalse
\newif\if@rmtwocol
\@rmtwocoltrue
\newif\if@biblist
\@biblistfalse
\newif\if@book
\@bookfalse
\newif\if@natbib
\@natbibtrue
\newif\if@rmcite
\@rmcitefalse
\newif\if@titlevisible
\@titlevisibletrue
\newif\if@abstractpresent
\@abstractpresentfalse
\newif\if@printbanner
\@printbannerfalse
\newif\if@suppressfulladdresses
\@suppressfulladdressesfalse
\newif\if@newstyle
\@newstylefalse
\newif\if@debug
\@debugfalse
\newif\if@boldkeywords
\@boldkeywordsfalse
%    \end{macrocode}
% The option handling---mainly setting various booleans.
% \changes{2.2beta11}{2001/11/26}{Made twocol and
% newstyle the default for proceedings.}
% \changes{3.23}{2007/09/09}{Manuscript is now 10 pt instead of 12pt}
% \changes{3.23}{2007/09/10}{Full addresses are now automatically
% suppressed for proceedings style.}
% \changes{3.24}{2007/09/14}{New option \texttt{boldkeywords} that
% does the keywords in the old ugly way. Now off by default.}
%    \begin{macrocode}
\DeclareOption{manuscript}{\typeout{Manuscript for referee} 
  \@manuscripttrue\@journalfalse\@rmtwocolfalse
  \PassOptionsToClass{10pt}{article}}
\DeclareOption{bigfont}{\typeout{Big font manuscript requested by Gary
  Ferland} 
  \@manuscripttrue\@journalfalse\@rmtwocolfalse\@bigfonttrue
  \PassOptionsToClass{11pt}{article}}
\DeclareOption{journal}{\typeout{Main journal option}%
  \@journaltrue\@rmtwocoltrue\@printdatestrue\@manuscriptfalse
  \PassOptionsToClass{10pt}{article}}
\DeclareOption{guide}{\typeout{User guide option}%
  \@guidetrue\@journalfalse\@rmtwocoltrue\@manuscriptfalse
  \PassOptionsToClass{10pt}{article}}
\DeclareOption{preprint}{\typeout{Preprint option}%
  \@preprinttrue\@manuscriptfalse\PassOptionsToClass{10pt}{article}}
\DeclareOption{proof}{\typeout{Author proof option}%
  \@prooftrue\@manuscriptfalse}
\DeclareOption{proceedings}{\typeout{Proceedings option}%
  \@proceedingstrue\@journalfalse\@rmtwocoltrue\@newstyletrue
  \@suppressfulladdressestrue
  \@printdatesfalse\@manuscriptfalse\PassOptionsToClass{10pt}{article}} 
\DeclareOption{printdates}{\typeout{Force printing of received/accepted dates}%
  \@printdatestrue} 
\DeclareOption{book}{\typeout{Whole issue option}%
  \@booktrue} 
\DeclareOption{nonatbib}{\typeout{Turn off natbib support (required
    for compatibility with old rmaacite package).}%
  \@natbibfalse} 
\DeclareOption{newstyle}{\typeout{Use new layout style}%
  \@newstyletrue}
\DeclareOption{oldstyle}{\typeout{Use old layout style}%
  \@newstylefalse}
\DeclareOption{debug}{\typeout{Loads of annoying messages}%
  \@debugtrue}
%    \end{macrocode}
% These two need to be declared explicitly so that |\ExecuteOptions|
% actually does something.
%    \begin{macrocode}
\DeclareOption{twocolumn}{\@rmtwocoltrue}
\DeclareOption{onecolumn}{\@rmtwocolfalse}
\DeclareOption{twoside}{\PassOptionsToClass{\CurrentOption}{article}}
\DeclareOption{boldkeywords}{\@boldkeywordstrue}
%    \end{macrocode}
% Bold, upper-case keywords used to be the default. Now they must be
% turned on explicitly with the \texttt{boldkeywords} option. 
%    \begin{macrocode}
\DeclareOption{boldkeywords}{\@boldkeywordstrue}
%    \end{macrocode}
% Pass any other option along to the \texttt{article} class.
%    \begin{macrocode}
\DeclareOption*{\PassOptionsToClass{\CurrentOption}{article}}
%    \end{macrocode}
% Default options are to be two-sided and based on the standard
% \LaTeX{} \texttt{article} class,
% unless we are processing the entire book, in which case we base
% ourselves on the standard \texttt{book} class. Default is to use
% manuscript option unless otherwise specified. 
%    \begin{macrocode}
\ExecuteOptions{twoside,manuscript} \ProcessOptions\relax
\if@book
\if@rmtwocol\LoadClass[twocolumn]{book}\else\LoadClass{book}\fi
\newcommand{\abstract}{\relax}
\newcommand{\@maketitle}{\relax}
\else
\if@rmtwocol\LoadClass[twocolumn]{article}\else\LoadClass{article}\fi
\fi
%    \end{macrocode}
% Attempt to use natbib for citations 02 Jan 2003. 
% \changes{3.05}{2003/01/02}{If the natbib option is used, then set
% things up so that natbib's version of the bibliography is identical
% to the standard one.}
% \changes{3.09}{2003/01/23}{Explicitly require a recent version of
% natbib, so that we get a decent error message if the package is too
% old.} 
% \changes{3.23}{2007/09/09}{Now allow Bibtex to work properly} 
% \changes{3.25}{2009/10/01}{Change \cs{newcommand} tp \cs{def} when setting \cs{bibfont}. This allows it to work with all versions of natbib.} 
% \changes{3.25}{2010/02/06}{Updated copyright year.} 
% \changes{3.25}{2010/02/06}{Updated email and web addresses to \texttt{crya.unam.mx}.} 
%    \begin{macrocode}
\if@natbib
\RequirePackage{natbib}[1999/02/25]
\bibpunct{(}{)}{;}{a}{}{,}
\def\bibfont{\small}
\renewcommand\bibsection{\section*{REFERENCES}}
\setlength\bibsep{0pt}
\setlength\bibhang{1.5em}
\bibliographystyle{rmaa}
\fi
%    \end{macrocode}
% We want the \texttt{graphicx} package for the extended syntax of
% |\includegraphics|. We want the \texttt{ifthen} package for using in
% these macros (although we don't much yet). We explicitly load
% \texttt{babel} since some \LaTeX{} distributions messed up the
% hyphenation if we didn't.  
%    \begin{macrocode}
\RequirePackage{graphicx}
\RequirePackage{ifthen}
\IfFileExists{babel.sty}{%
%   \let\@restoreindent=\@afterindentfalse
  \RequirePackage[english]{babel}
%   \@restoreindent
  %% reinstate the definition of \@roman so as not to break .idx files
  \@ifundefined{org@@roman}{\relax}{\renewcommand{\@roman}[1]{\org@@roman{##1}}}
  }{%
  \typeout{Package babel.sty not found - don't worry about this}}
%    \end{macrocode}
% Try to load the \texttt{dcolumn} and \texttt{booktabs} packages if
% they can be found. These can produce nicer-looking tables. 
%    \begin{macrocode}
\IfFileExists{dcolumn.sty}{%
  \RequirePackage{dcolumn}
  }{%
  \typeout{Package dcolumn.sty not found - you will not be able to
    align columns on decimal points in tables}}
\IfFileExists{booktabs.sty}{%
  \RequirePackage{booktabs}
  }{%
  \typeout{Package booktabs.sty not found - your tables will not look
    as pretty as they ought too.}
  }
\IfFileExists{makeidx.sty}{%
  \RequirePackage{makeidx}
}{%
  \typeout{Package makeidx.sty not found - Author Index may not work.}
  }
%% commented out for moment since requires more modern version of LaTeX
%% \RequirePackage[latin1]{inputenc}
%    \end{macrocode}
% \changes{4.0}{2019/11/13}{Make \texttt{hyperref} and
% \texttt{etoolbox} into hard requirements, since we now use them
% extensively.}  We use hyperref to get internal and external links.
% But I don't like the colored boxes so I turn those off.  We use
% etoolbox for patching and robustifying at the moment.  In the
% future, I want to port a lot of other stuff over to it.  For
% instance, replace the ifthen package and the TeX ifs too.
%
% We target the latest version (1.7) of PDF, rather than the default
% of 1.5.  This reduces the chance of getting warnings from included
% graphics.  This is not strictly related to hyperref, but we put it
% here anyway.
%    \begin{macrocode}
\RequirePackage{etoolbox}
\RequirePackage{hyperref}
\hypersetup{hidelinks}
\pdfminorversion=7
%    \end{macrocode}
% We patch the \cs{@combinedblfloats} following a suggestion by Heiko
% Oberdiek at \url{https://tex.stackexchange.com/a/249743/3843}, which
% avoids problems when hyperlinks are split across
% pages. \changes{4.0}{2019/11/12}{Fix for issue where breaking
% hyperlink across pages gave an error.}
%    \begin{macrocode}
\urlstyle{same}
\newcommand\email{\begingroup \urlstyle{rm}\Url}
\patchcmd\@combinedblfloats{\box\@outputbox}{\unvbox\@outputbox}{}{%
  \errmessage{\noexpand\@combinedblfloats could not be patched}%
}%
%    \end{macrocode}
% \subsection{Other setup stuff}\label{sec:setup}
% Some declarations: 
%    \begin{macrocode}
\newif\if@needaddresses
\@needaddressestrue
\newif\if@frontmatter
\@frontmatterfalse
\newif\if@inpreamble
\@inpreambletrue
\newif\if@firstauthor
\@firstauthortrue
\newif\if@firstaffil
\@firstaffiltrue
\newif\if@firstaltaffil
\@firstaltaffiltrue
\newif\if@firstkwd
\@firstkwdtrue
\newif\if@noexpandADS
\@noexpandADSfalse
\newif\if@sessiontitlepages
\@sessiontitlepagesfalse
%    \end{macrocode}
% \begin{macro}{\@rmaaEmpty}
% \begin{macro}{\@rmaaArg}
% These are used in various commands to check if arg was left blank
%    \begin{macrocode}
\newcommand{\@rmaaEmpty}{}
\newcommand{\@rmaaArg}{}
%    \end{macrocode}  
% \end{macro}
% \end{macro}
%
% \subsection{Compatibility with old \LaTeX{} versions}
% \label{sec:compat}
% This package is based on the 1998/12/01 release of \LaTeX{}.
% Older releases may not define all the commands that we need.
% Luckily, we can use |\providecommand| to define these if necessary. 
% The following commands are not defined in the 1994/06/01 release. 
%
% \begin{macro}{\MakeUpperCase}
% \begin{macro}{\textsuperscript}
% \begin{macro}{\textbackslash}
% User commands:
%    \begin{macrocode}
\providecommand{\MakeUpperCase}[1]{\uppercase{#1}}
\providecommand{\textsuperscript}[1]{\ensuremath{^\textrm{#1}}}
\providecommand{\textbackslash}{\ensuremath{\backslash}}
%    \end{macrocode}  
% \end{macro}
% \end{macro}
% \end{macro}
% \begin{macro}{\protected@xdef}
% \begin{macro}{\restore@protect}
% \begin{macro}{\@unexpandable@protect}
% \begin{macro}{\hb@xt@}
% \LaTeX{} internal commands:
%    \begin{macrocode}
\providecommand{\protected@xdef}{%
        \let\@@protect\protect
        \let\protect\@unexpandable@protect
        \afterassignment\restore@protect
        \xdef}
\providecommand{\restore@protect}{\let\protect\@@protect}
\providecommand{\@unexpandable@protect}{\noexpand\protect\noexpand}
\providecommand{\hb@xt@}{\hbox to}
%    \end{macrocode}
% \end{macro}
% \end{macro}
% \end{macro}
% \end{macro}
%
% \subsection{Page layout}\label{sec:page-layout}
%
% The dimensions of the text were defined by Elizabeth. Text width is
% 169\,mm with equal margins on the two sides: 
%    \begin{macrocode}
\setlength{\textwidth}{169mm}
\setlength{\@tempdima}{\paperwidth}
\addtolength{\@tempdima}{-\textwidth}
\setlength{\oddsidemargin}{0.5\@tempdima}
\addtolength{\oddsidemargin}{-1in}
\setlength{\evensidemargin}{\oddsidemargin}
%    \end{macrocode}
% Text height is 225\,mm, top margin is 14\,mm. 
%    \begin{macrocode}
\setlength{\textheight}{225mm} 
\setlength{\topmargin}{-1in} 
\addtolength{\topmargin}{14mm}
%    \end{macrocode}
% Space between columns in two-column styles, header, footer,
% paragraph indentation:
%    \begin{macrocode}
\setlength{\columnsep}{9.5mm}
\setlength{\headsep}{0.7cm}
\setlength{\footskip}{0.75cm}
\setlength{\parindent}{1.5em}
%    \end{macrocode}
% 
% \begin{macro}{\@set@scalable@length}
% Try out a new scheme for dealing with various spacing dimensions and
% allowing the user to easily modify them. The command
% \begin{quote}
% |\set@scalable@length|\marg{name}\marg{defaultlength}\marg{stretch}\marg{shrink}
% \end{quote}
% does a 
% |\setlength|\marg{\cs{name}}\marg{defaultlength plus stretch minus shrink}
% and also defines the user command
% \textbackslash\meta{name}|Scale|\marg{factor} that can 
% be used to rescale length by \meta{factor}. Additionally, now
% defines the package command
% \textbackslash\meta{name}|@set@scale|\marg{FACTOR} 
% that also rescales the length by \meta{FACTOR}, so that the length
% becomes \meta{factor} $\times$ \meta{FACTOR} $\times$
% \meta{defaultlength}.   
%    \begin{macrocode}
\newcommand\@set@scalable@length[4]{%
  %% arg #1: NAME of length without leading `\'
  %%
  %% arg #2: default VALUE  
  \expandafter\newlength\csname @#1@default\endcsname
  \expandafter\setlength\csname @#1@default\endcsname{#2}%
  \expandafter\newlength\csname @@#1@@default\endcsname
  \expandafter\setlength\csname @@#1@@default\endcsname{#2}%
  %% arg #3: default STRETCH  
  \expandafter\newlength\csname @#1@stretch\endcsname
  \expandafter\setlength\csname @#1@stretch\endcsname{#3}%
  \expandafter\newlength\csname @@#1@@stretch\endcsname
  \expandafter\setlength\csname @@#1@@stretch\endcsname{#3}%
  %% arg #4: default SHRINK  
  \expandafter\newlength\csname @#1@shrink\endcsname
  \expandafter\setlength\csname @#1@shrink\endcsname{#4}%
  \expandafter\newlength\csname @@#1@@shrink\endcsname
  \expandafter\setlength\csname @@#1@@shrink\endcsname{#4}%
  %%
  %% User-defined scale \@NAME@scale is unity by default
  \expandafter\gdef\csname @#1@scale\endcsname{1}%
  %% Package-defined scale \@NAME@SCALE is unity by default
  \expandafter\gdef\csname @#1@SCALE\endcsname{1}%
  %%
  %% internal command \@reset@NAME should expand to 
  %% \setlength\@@NAME@@default{\@NAME@SCALE\@NAME@default}
  %% \setlength\NAME{\@NAME@scale\@@NAME@@default}
  %% in other words: scale x SCALE x default 
  \expandafter\gdef\csname @reset@#1\endcsname{%
    \expandafter\setlength\csname @@#1@@default\endcsname{%
      \csname @#1@SCALE\endcsname \csname @#1@default\endcsname
      }
    \expandafter\setlength\csname @@#1@@stretch\endcsname{%
      \csname @#1@SCALE\endcsname \csname @#1@stretch\endcsname
      }
    \expandafter\setlength\csname @@#1@@shrink\endcsname{%
      \csname @#1@SCALE\endcsname \csname @#1@shrink\endcsname
      }
    \expandafter\setlength\csname #1\endcsname{%
      \csname @#1@scale\endcsname \csname @@#1@@default\endcsname
      \@plus
      \csname @#1@scale\endcsname \csname @@#1@@stretch\endcsname
      \@minus
      \csname @#1@scale\endcsname \csname @@#1@@shrink\endcsname
      }}
  %%
  %% User command \NAMEScale should expand to 
  %% {\renewcommand\@NAME@scale{#1}\@reset@NAME}
  \expandafter\gdef\csname #1Scale\endcsname##1{%
    \expandafter\gdef\csname @#1@scale\endcsname{##1}%
    \csname @reset@#1\endcsname}
  %%
  %% Package command \NAME@set@SCALE should expand to 
  %% {\renewcommand\@NAME@SCALE{#1}\@reset@NAME}
  \expandafter\gdef\csname #1@set@SCALE\endcsname##1{%
    \expandafter\gdef\csname @#1@SCALE\endcsname{##1}%
    \csname @reset@#1\endcsname}
  %%
  %% initialize at default values with \@reset@NAME
  \csname @reset@#1\endcsname
  }
%    \end{macrocode}
% \end{macro}
% \begin{macro}{\parskip}
% First use of this new infrastructure: set interparagraph space to be
% zero but with a small amount of stretchiness.
%    \begin{macrocode}
%   \ifthenelse{\boolean{@manuscript}}{%
%     \@set@scalable@length{parskip}{0.5\baselineskip}{0.3\p@}{0.1\p@}
%   }{
    \@set@scalable@length{parskip}{0\p@}{0.3\p@}{0\p@}
%   }
%    \end{macrocode}
% \end{macro}
% 
% \begin{macro}{\abovesectionskip}
% \changes{4.0}{2019/11/12}{Use negative values for skip so that first paragraph after section headings is not indented.  This is something that I have been trying to fix for years!  It turns out that the anser was in the documentation for the standard command \cs{@startsection.}}
% \begin{macro}{\belowsectionskip}
% \begin{macro}{\aboveparaskip}
% \begin{macro}{\belowparaskip}
% Define some utility rubber lengths for use in
% section~\ref{sec:section-commands}. 
%    \begin{macrocode}
\newlength{\abovesectionskip}
\newlength{\belowsectionskip}
\newlength{\aboveparaskip}
\newlength{\belowparaskip}
%    \end{macrocode}
% The default values for these lengths are those traditionally used
% for the onecolumn style. 
%    \begin{macrocode}
\@set@scalable@length{abovesectionskip}{-3.5ex}{-4ex}{-2ex}
\@set@scalable@length{belowsectionskip}{2.3ex}{0ex}{2ex}
\@set@scalable@length{aboveparaskip}{3.25ex}{2ex}{0.5ex}
\@set@scalable@length{belowparaskip}{-1em}{\z@}{\z@}
%    \end{macrocode}
% \end{macro}\end{macro}\end{macro}\end{macro}
% We also make scalable some of the standard float parameters. These
% are the same defaults as in the standard article classes, but
% slightly more stretchy. For the twocolumn style, they are currently
% halved, see below. Note that I have also lumped |\abovecaptionskip|
% in with these, since it plays a similar role. There is also a
% |\belowcaptionskip| but it is set to zero in the standard classes. 
%    \begin{macrocode}
\@set@scalable@length{floatsep}{12\p@}{4\p@}{4\p@}
\@set@scalable@length{textfloatsep}{20\p@}{4\p@}{8\p@}
\@set@scalable@length{intextsep}{12\p@}{4\p@}{4\p@}
\@set@scalable@length{dblfloatsep}{12\p@}{4\p@}{4\p@}
\@set@scalable@length{dbltextfloatsep}{20\p@}{4\p@}{8\p@}
\@set@scalable@length{abovecaptionskip}{10\p@}{\z@}{\z@}
%    \end{macrocode}
% \begin{macro}{\@before@title@skip}         
% \begin{macro}{\@before@author@skip}        
% \begin{macro}{\@between@author@skip}       
% \begin{macro}{\@before@abstract@skip}      
% \begin{macro}{\@after@abstracthead@skip}   
% \begin{macro}{\@between@abstract@skip}     
% \begin{macro}{\@before@kwd@skip}           
% \begin{macro}{\@before@dates@skip}           
% Spacing parameters used in |\@maketitle|. These are not changed by
% default except by the user, so they don't need to go into
% |\@package@default@lengths|.  
%    \begin{macrocode}
\newlength{\@before@title@skip}         
\newlength{\@before@author@skip}        
\newlength{\@between@author@skip}       
\newlength{\@before@abstract@skip}      
\newlength{\@after@abstracthead@skip}   
\newlength{\@between@abstract@skip}     
\newlength{\@before@kwd@skip}           
\newlength{\@before@dates@skip}           
\@set@scalable@length{@before@title@skip}{2em}{\z@}{\z@}
\@set@scalable@length{@before@author@skip}{1.5em}{\z@}{\z@}
\@set@scalable@length{@between@author@skip}{0.5em}{\z@}{\z@}
\@set@scalable@length{@before@abstract@skip}{3ex}{1ex}{0.5ex}
\@set@scalable@length{@after@abstracthead@skip}{1ex}{\z@}{\z@}
\@set@scalable@length{@between@abstract@skip}{2ex}{1ex}{0.5ex}
\@set@scalable@length{@before@kwd@skip}{1ex}{0.5ex}{0.2ex}
\@set@scalable@length{@before@dates@skip}{2ex}{0.5ex}{0.2ex}
%    \end{macrocode}
% \end{macro}\end{macro}\end{macro}\end{macro}
% \end{macro}\end{macro}\end{macro}\end{macro}
% 
% Spacing above and below displayed equations also wants to be done
% like this, but it is more complicated since the standard article
% classes make these change with the fontsize. Something to do later,
% maybe. 
% 
%    \begin{macrocode}
%    \end{macrocode}
% \begin{macro}{\@package@default@lengths}
% The value actually used for these lengths is controlled by the
% \meta{SCALE} parameter, which depends on the style option that is in
% effect. This is set using the |\@package@default@lengths| command,
% which must be called every time the style changes (e.g., changing
% from onecolumn to twocolumn). 
%    \begin{macrocode}
\newcommand{\@package@default@lengths}{%
  \ifthenelse{\boolean{@rmtwocol}}{%
    %% twocolumn mode
    \abovesectionskip@set@SCALE{0.5}
    \belowsectionskip@set@SCALE{0.5}
    \aboveparaskip@set@SCALE{0.5}
    \belowparaskip@set@SCALE{1}
    \floatsep@set@SCALE{0.5}
    \textfloatsep@set@SCALE{0.5}
    \intextsep@set@SCALE{0.5}
    \dblfloatsep@set@SCALE{0.5}
    \dbltextfloatsep@set@SCALE{0.5}
    }{%
    %% onecolumn mode
    \abovesectionskip@set@SCALE{1}
    \belowsectionskip@set@SCALE{1}
    \aboveparaskip@set@SCALE{1}
    \belowparaskip@set@SCALE{1}
    \floatsep@set@SCALE{1}
    \textfloatsep@set@SCALE{1}
    \intextsep@set@SCALE{1}
    \dblfloatsep@set@SCALE{1}
    \dbltextfloatsep@set@SCALE{1}
    }%
  }
\@package@default@lengths
%    \end{macrocode}
% \end{macro}
% \begin{macro}{\RescaleLengths}
% \begin{macro}{\RescaleSecLengths}
% \begin{macro}{\RescaleFloatLengths}
% \begin{macro}{\RescaleTitleLengths}
% User commands that will fiddle with all the scalable lengths. 
%    \begin{macrocode}
\newcommand{\RescaleLengths}[1]{%
  \RescaleSecLengths{#1}\RescaleFloatLengths{#1}\RescaleTitleLengths{#1}}
\newcommand{\RescaleSecLengths}[1]{%
%% sections
  \abovesectionskipScale{#1}
  \belowsectionskipScale{#1}
  \aboveparaskipScale{#1}
  }
\newcommand{\RescaleFloatLengths}[1]{%
%% floats
  \floatsepScale{#1}
  \textfloatsepScale{#1}
  \intextsepScale{#1}
  \dblfloatsepScale{#1}
  \dbltextfloatsepScale{#1}
  \abovecaptionskipScale{#1}
  }
\newcommand{\RescaleTitleLengths}[1]{%
%% maketitle
  \@before@title@skipScale{#1}           
  \@before@author@skipScale{#1}        
  \@between@author@skipScale{#1}       
  \@before@abstract@skipScale{#1}      
  \@after@abstracthead@skipScale{#1}   
  \@between@abstract@skipScale{#1}     
  \@before@kwd@skipScale{#1}           
  \@before@dates@skipScale{#1}           
  }
%    \end{macrocode}
% \end{macro}\end{macro}\end{macro}\end{macro}
% \changes{2.2beta14}{2002/03/17}{Finished implementation of new
% infrastructure for scalable lengths. Converted the above/below
% section skips to use it. Also, all the floatsep parameters. Also,
% all the vertical lengths that are used by |\@maketitle|. }

% \begin{macro}{\MSextramargin}
% For the \texttt{manuscript} option we have an extra 0.3\,inch
% vertical margin and 1\,inch horizontal margin for marginal notes (by
% hand or with |\marginpar|. 
%
% \changes{3.23}{2007/09/09}{Reduce baselinestretch in
%   \texttt{manuscript} style to 1.2. No longer use raggedright for
%   manuscript.}
%
% \changes{3.24}{2007/09/14}{Further reduction of baselinestretch in
%   \texttt{manuscript} style to 1.1 since we now have wider margins.}
%
% \changes{3.24}{2007/09/14}{Split \texttt{manuscript}'s extra margin
%   into separate horizontal and vertical parts.}
% 
%    \begin{macrocode}
\newlength{\MSextraHmargin}\newlength{\MSextraVmargin}
\setlength{\MSextraHmargin}{1in} 
\setlength{\MSextraVmargin}{0.3in} 
\if@manuscript
   \renewcommand{\baselinestretch}{1.1}
   \addtolength{\textwidth}{-2\MSextraHmargin}
   \addtolength{\oddsidemargin}{\MSextraHmargin}
   \setlength{\evensidemargin}{\oddsidemargin}
   \addtolength{\textheight}{-2\MSextraVmargin}
   \addtolength{\topmargin}{\MSextraVmargin}
%    \raggedright
\fi
%    \end{macrocode}
% \end{macro}
% Float parameters: don't allow bottom floats by default. 
%    \begin{macrocode}
\renewcommand{\fps@figure}{tp}
\renewcommand{\fps@table}{tp}
%    \end{macrocode}
%
% \begin{macro}{\rmaa@AuthorWidthFrac}
% \begin{macro}{\rmaa@AbstractWidthFrac}
% These parameters control what fraction of the total text width can
% be used by the list of authors and the abstract/resumen,
% respectively.  
%    \begin{macrocode}
\newcommand{\rmaa@AuthorWidthFrac}{1.0}
\newcommand{\rmaa@AbstractWidthFrac}{0.75}
%    \end{macrocode}
% Use a wide abstract to save space in the new style.
%    \begin{macrocode}
\if@newstyle\renewcommand{\rmaa@AbstractWidthFrac}{1.0}\fi
%    \end{macrocode}
% \end{macro}
% \end{macro}
%
% Some things to consider in future versions. Trying to avoid too many
% widows, etc. Apparently, the following will work best with
% \cs{flushbottom}. 
%    \begin{macrocode}
%%\widowpenalty=10000
%%\clubpenalty=10000
%    \end{macrocode}
% Also, should perhaps try getting rid of underfull/overfull box warnings. 
% \changes{3.10}{2003/01/24}{Set large values for \cs{hbadness},
% \cs{vbadness}, \cs{hfuzz}, \cs{vfuzz}. To reduce amount of stuff
% that gets written to the terminal every time you compile something.}
%    \begin{macrocode}
\hbadness=10000
\vbadness=10000
\setlength\hfuzz{\maxdimen}
\setlength\vfuzz{\maxdimen}
%    \end{macrocode}
% \subsection{Various text styles for different contexts}
% \label{sec:text-styles} 
% \begin{macro}{\rmaa@AffilStyle}
% \begin{macro}{\rmaa@PagesStyle}
% \begin{macro}{\rmaa@YearStyle}
% \begin{macro}{\rmaa@VolumeStyle}
% \begin{macro}{\rmaa@VolumenStyle}
% \begin{macro}{\rmaa@JournalNameStyle}
% \begin{macro}{\rmaa@ReceivedStyle}
% \begin{macro}{\rmaa@AbstractHeaderStyle}
% \begin{macro}{\rmaa@JournalNameStyle}
% Style commands for the different fields used by |\maketitle| in
% constructing the title page of the article. These have the form
% |\rmaa@XYZStyle{TEXT}|. I'm not sure these were really a good idea,
% but never mind \dots 
%    \begin{macrocode}
\newcommand{\rmaa@AffilStyle}[1]{{\normalsize\rmfamily#1}}
\newcommand{\rmaa@PagesStyle}[1]{{\normalsize\rmfamily#1}}
\newcommand{\rmaa@YearStyle}[1]{{\normalsize\rmfamily#1}}
\newcommand{\rmaa@VolumeStyle}[1]{{\normalsize\rmfamily\bfseries#1}}
\newcommand{\rmaa@VolumenStyle}[1]{{\Large\sffamily\bfseries#1}}
\newcommand{\rmaa@JournalNameStyle}[1]{{\normalsize\itshape#1}}
\newcommand{\rmaa@ReceivedStyle}[1]{{\normalsize\rmfamily\itshape#1}}
\newcommand{\rmaa@AbstractHeaderStyle}[1]{%
  {\normalsize\rmfamily\MakeUpperCase{#1}}}
\newcommand{\rmaa@TitleStyle}[1]{%
  {\large\bfseries\rmfamily\boldmath\MakeUpperCase{#1}\par}}
%    \end{macrocode}
% \end{macro}
% \end{macro}
% \end{macro}
% \end{macro}
% \end{macro}
% \end{macro}
% \end{macro}
% \end{macro}
% \end{macro}
%
% \begin{macro}{\rmaa@FrontSecStyle}
% Headings of ``front sections'' (see section~\ref{sec:front-sections})
%    \begin{macrocode}
\newcommand{\rmaa@FrontSecStyle}[1]{%
  {\normalsize\bfseries\rmfamily\MakeUpperCase{#1}}}
%    \end{macrocode}
% \end{macro}
%
% \begin{macro}{\rmaa@TOCTitleStyle}
% \begin{macro}{\rmaa@TOCPartStyle}
% \begin{macro}{\rmaa@IndexTitleStyle}
% \begin{macro}{\rmaaTOCAuthorStyle}
% \begin{macro}{\rmaaIndexAuthorStyle}
% Commands to set style of TOC and index entries.
% Rewritten to avoid unnecessary font commands in the .toc and .ind files
%    \begin{macrocode}
\newcommand{\rmaa@TOCTitleStyle}[1]{\MakeUpperCase{#1}}
\newcommand{\rmaa@TOCPartStyle}[1]{\rmaa@TOCTitleStyle{#1}}
\newcommand{\rmaa@IndexTitleStyle}[1]{#1}
\newcommand{\rmaaTOCAuthorStyle}[1]{%
  {\itshape#1}}
\newcommand{\rmaaIndexAuthorStyle}[1]{%
  \protect\mbox{\bfseries#1}}
%    \end{macrocode}
% \end{macro}
% \end{macro}
% \end{macro}
% \end{macro}
% \end{macro}
%
% \begin{macro}{\rmaa@KeywordStyle}
% \begin{macro}{\rmaa@KeywordHeaderStyle}
% The style of the keywords depends on whether we are using the
% \texttt{newstyle} option or not. If we are, then they are small. 
%    \begin{macrocode}
\ifthenelse{\boolean{@newstyle}}{
  \newcommand{\rmaa@KeywordStyle}[1]{%
    {\small%
      \if@boldkeywords 
      \bfseries\MakeUpperCase{\raggedright #1}
      \else
      \raggedright #1
      \fi
    }}
  \newcommand{\rmaa@KeywordHeaderStyle}[1]{%
    {\small\itshape#1}}}
{
  \newcommand{\rmaa@KeywordStyle}[1]{%
    {\normalsize%
      \if@boldkeywords\bfseries\fi\raggedright #1}}
  \newcommand{\rmaa@KeywordHeaderStyle}[1]{%
    {\normalsize\itshape#1}}}
%    \end{macrocode}
% \end{macro}
% \end{macro}
%
% \subsection{Page styles}\label{sec:page-styles}
% \begin{macro}{\ps@rmaa}
% The running headers for normal pages:
%    \begin{macrocode}
\newcommand{\ps@rmaa}{%
%    \end{macrocode}
% nothing in either foot; 
%    \begin{macrocode}
  \renewcommand\@oddfoot{\relax}%
  \renewcommand\@evenfoot{\relax}%
%    \end{macrocode}
% page number and short author list in the even heads; 
%    \begin{macrocode}
  \renewcommand\@evenhead{\rlap{\thepage}\hfil\@shortauthor\hfil}%
%    \end{macrocode}
% short title and page number in the odd heads.
%    \begin{macrocode}
  \renewcommand\@oddhead{\hfil\@shorttitle\hfil\llap{\thepage}}
  }
%    \end{macrocode}
% \begin{macro}{\ps@rmaafirstpage}
% The header of the first page of each
% article:
%    \begin{macrocode}
\newcommand{\ps@rmaafirstpage}{%
  \renewcommand\@oddfoot{\hfil\thepage}%
%    \end{macrocode}
% footer corrected (06 Jun 2000) to have page on outside edge;
%    \begin{macrocode}
  \renewcommand\@evenfoot{\thepage\hfil}%
%    \end{macrocode}
% Header has journal name, volume, and page numbers for a regular
% article, or similar for special cases. 
% \changes{4.0}{2019/11/03}{Change to three-line header with copyright info and DOI.}
%    \begin{macrocode}
\if@preprint
\renewcommand\@oddhead{%
  \if@proceedings 
  \textsc{To appear in ``\@conftitle{} (\rmaa@YearStyle{\@year})''}\hfil
  \textit{RMxAC}
  \else
  \textsc{To appear in} \@journalname\hfil
  \fi}\fi%
\if@guide
\renewcommand\@oddhead{\textsc{Documentation} (\rmaa@YearStyle{\@year})\hfil}
\fi%
\if@manuscript
\renewcommand\@oddhead{Manuscript for \@journalname{} (\rmaa@YearStyle{\@year})\hfil}
\fi%
\if@journal
\renewcommand\@oddhead{%
  \parbox{\textwidth}{
    \@journalname, \rmaa@VolumeStyle{\@volume},
    \rmaa@PagesStyle{\@artpages}\
    (\rmaa@YearStyle{\@year})\hfil\\[1pt]
    {\small\textcopyright{} \@year:
      Instituto de Astronom\'\i{}a,
      Universidad Nacional Aut\'onoma de M\'exico}\hfil\\[1pt]
    {\small\@DOI@url}\hfil
  }%
}\fi%
  \renewcommand\@evenhead{\@oddhead}}
%    \end{macrocode}
% By default, use the |\ps@rmaa| page style.
%    \begin{macrocode}
\pagestyle{rmaa}
%    \end{macrocode}
% \end{macro}
% \end{macro}
%
% \subsection{Preamble commands}\label{sec:preamble} 
% Everything you need for setting the title page, TOC and index
% entries, etc.
% \subsubsection{Commands for the title page}
% \label{sec:title-commands}
% \begin{macro}{\resumen}
% \begin{macro}{\@resumen}
% \begin{macro}{\abstract}
% \begin{macro}{\@abstract}
% User commands for setting the resumen and abstract. 
% These should be used in the document header. The internal commands
% |\@resumen| and |\@abstract| are used to hold the actual text
% pending output via the |\maketitle| command. 
%    \begin{macrocode}
\newif\if@setresumen
\newcommand{\resumen}[1]{%
%%  \renewcommand{\@resumen}{{\selectlanguage{spanish}#1}}%
  \renewcommand{\@resumen}{{\ignorespaces#1}}%
  \@setresumentrue}
\renewcommand{\abstract}[1]{%
  \renewcommand{\@abstract}{{\ignorespaces#1}}
  \renewcommand\ADS@abstract{#1}
%% Not a good idea, since then author macros are not expanded and
%% cause errors later. However, by not doing it we are vulnerable to
%% fragile commands.....
  \if@noexpandADS\@onelevel@sanitize{\ADS@abstract}\fi
  \@abstractpresenttrue
  }
\newcommand{\ADSnoexpand}{\@noexpandADStrue}
%    \end{macrocode}
% \begin{macro}{\boldabstract}
% This is the version for doing 1-page conference proceedings. Still
% save the text for ADS but then write it out immediately as a bold
% paragraph. 
%    \begin{macrocode}
\newcommand{\boldabstract}[1]{%
  \gdef\ADS@abstract{#1}%
  \renewcommand{\@abstract}{#1}%
  \if@noexpandADS\@onelevel@sanitize{\ADS@abstract}\fi%
  \noindent\textbf{\boldmath #1}
  }
%    \end{macrocode}
% \end{macro}
% Default value in each case is to print out a warning if not set
%    \begin{macrocode}
  \newcommand{\@resumen}{%
    \@setresumenfalse \ClassWarningNoLine{rmaa}{No Spanish abstract
      (resumen) specified - if you do not speak Spanish, this will be
      translated by the editors}%
    \textbf{Favor de proporcionar un resumen en espa\~nol. If you are
      unable to translate your abstract into Spanish, the editors will
      do it for you.} }
  \newcommand{\@abstract}{\ClassWarningNoLine{rmaa}{No abstract
      specified}}
%    \end{macrocode}
% \end{macro}
% \end{macro}
% \end{macro}
% \end{macro}
%
% \begin{macro}{\title}
% \begin{macro}{\@rawtitle}
% \begin{macro}{\@TOCtitle}
% \begin{macro}{\rmaaIndexTitle}
% Similarly, |\title| is the preamble command to set the title of the
% article. This is used in various places, typeset in different
% styles, hence the use of the internal |\@rawtitle|, |\@title|,
% |\rmaaIndexTitle|, and |\@TOCtitle| commands, each of which uses the
% relevant style from section~\ref{sec:text-styles}
%    \begin{macrocode}
\newcommand{\@rawtitle}{Raw Title}
\newcommand{\@TOCtitle}{\rmaa@TOCTitleStyle{Title for TOC}}
\newcommand{\rmaaIndexTitle}{\rmaa@IndexTitleStyle{Title for index}.}
\newcommand{\@TOCaftertitle}{}
\renewcommand{\title}[1]{%
%%  {\ADS@sanitize\xdef\@rawtitle{#1}}
  \renewcommand\@rawtitle{\noexpand#1}
  \let\ADS@title=\@rawtitle
  \renewcommand{\@title}{\rmaa@TitleStyle{#1}}%
  \renewcommand{\@TOCtitle}{\rmaa@TOCTitleStyle{#1}\@TOCaftertitle}%
  \renewcommand{\rmaaIndexTitle}{\rmaa@IndexTitleStyle{#1\@TOCaftertitle}.}%
  }
%    \end{macrocode}
% \end{macro}
% \end{macro}
% \end{macro}
% \end{macro}
% 
% All the weird and whacky ways of indicating the author
% affiliations. 
% \begin{macro}{\affil}
% First, the way that sets it directly underneath the authors name.
% Note that |\affil| trumps all other affiliations for the purposes of
% getting into the ADS stuff (we use |\gdef|).   
% \dots 
%    \begin{macrocode}
\newcommand{\affil}[1]{%
  \\[1.5ex]\rmaa@AffilStyle{#1}\\[1.5ex]%
  \if@firstaffil\gdef\ADS@affil{#1}\@firstaffilfalse\fi%
  }
\renewcommand{\and}{and\\[1.5ex]}
%    \end{macrocode}
% \begin{macro}{\altaffil}
% Alternative affiliations are footnoted. 
%    \begin{macrocode}
\newcommand{\altaffil}[1]{%
  \if@firstaltaffil\gdef\ADS@affil{#1}\@firstaltaffilfalse\fi%
  \thanks{#1}}
%    \end{macrocode}
% \begin{macro}{\altaffilmark}
% \begin{macro}{\altaffiltext}
% Alternative-alternative mechanism for affiliations---should not be
% mixed with |\altaffil| (unless you know what you are doing) or there
% will be problems with numbering.
%    \begin{macrocode}
\newcommand{\altaffilmark}[1]{\textsuperscript{#1}}
\newcommand{\altaffiltext}[2]{%
  \if@firstaltaffil\gdef\ADS@affil{#2}\@firstaltaffilfalse\fi
  \stepcounter{footnote}%
  \protected@xdef\@thanks{\@thanks
    \protect\footnotetext[#1]{#2}}} 
%    \end{macrocode}
% \end{macro}
% \end{macro}
% \end{macro}
% \end{macro}
%
% \begin{macro}{\keywords}
% \begin{macro}{\@keywords}
% \begin{macro}{\@keywordheader}
% \label{keywords}Sundry minor preamble commands. All
% self-explanatory, I hope. |\keywords| is the old user interface
% (deprecated). New articles should use |\addkeyword| once for each
% keyword. That way the ADS record comes out right too (comma rather
% than em-dash separator).
% \changes{2.2beta12}{2001/12/07}{\cs{addkeyword} New user interface
%   for keywords.} 
% \changes{2.2beta13}{2002/03/11}{\cs{addkeyword} Improved behaviour
%   with .ads file. Use
%   \cs{ADS@sanitize} to change definitions of `~' and \cs{ion}.}
% \changes{3.23}{2007/09/09}{\cs{addkeyword} now signals a compilation
% error if the author tries to add more than 6 keywords.}
% \changes{3.24}{2007/09/14}{kwdcounter is now reset in
% \cs{@begofarticlehook} to prevent error in book style} 
%    \begin{macrocode}
\newcommand{\@kwdsep}{ --- }
\newcommand{\ADS@kwdsep}{, }
\newcounter{kwdcounter}\setcounter{kwdcounter}{0}
\newcommand{\addkeyword}[1]{%
  \addtocounter{kwdcounter}{1}
  \ifthenelse{\thekwdcounter > 6}{
    \ClassError{rmaa}{Too many keywords. A maximum of 6 keywords is allowed.}{}
  }{}
  \if@firstkwd%
  \renewcommand\@keywords{\if@boldkeywords\uppercase{#1}\else #1\fi}
  {\ADS@sanitize\def\ion{\plainion}\gdef\ADS@keywords{#1}}
  \@firstkwdfalse
  \else%
  \edef\@keywords{\@keywords\@kwdsep \if@boldkeywords\uppercase{#1}\else #1\fi}%
  {\ADS@sanitize\def\ion{\plainion}\xdef\ADS@keywords{\ADS@keywords\ADS@kwdsep #1}}%
  \fi}
\newcommand{\ADS@sanitize}{%
  \def\mathrm{}
  \def\@{\ }
  \def~{\ }
  \def\langle{<}
  \def\rangle{>}
  \def\log{log }
  \def\sin{sin }
  \def\cos{cos }
  \def\tan{tan }
  \def\fd{^d.}
  \def\fh{^h.}
  \def\fm{^m.}
  \def\fs{^s.}
  \def\fdg{�.}
  \def\farcm{'.}
  \def\farcs{''.}
  \def\arcdeg{�}
  \def\,{ }
  }
\newcommand{\keywords}[1]{%
  \renewcommand{\ADS@keywords}{#1}
  \renewcommand{\@keywords}{\uppercase{#1}}
  }
\newcommand{\@keywordheader}{\rmaa@KeywordHeaderStyle{Key Words:}}
\newcommand{\@keywords}{%
  \ClassWarningNoLine{rmaa}{No keywords specified}%
  Keywords go here
  }
%    \end{macrocode}
% \begin{macro}{\shorttitle}
% \begin{macro}{\shortauthor}
% \begin{macro}{\@shorttitle}
% \begin{macro}{\@shortauthor}
% Commands for setting the running headers.
%    \begin{macrocode}
\newcommand{\shorttitle}[1]{%
  \renewcommand\@shorttitle{\MakeUpperCase{#1}}}
\newcommand{\shortauthor}[1]{%
  \renewcommand\@shortauthor{\MakeUpperCase{#1}}}
\newcommand{\@shorttitle}{SHORT TITLE GOES HERE}
\newcommand{\@shortauthor}{SHORT AUTHOR GOES HERE}
%    \end{macrocode}
% \end{macro}
% \end{macro}
% \end{macro}
% \end{macro}
% \end{macro}
% \end{macro}
% \end{macro}


% \subsubsection{The full addresses}
% \label{sec:full-addresses}
% \begin{macro}{\fulladdresses}
% \begin{macro}{\@fulladdresses}
% \begin{macro}{\@fulladdheight}
% \begin{macro}{\@fulladddepth}
% \begin{macro}{\@fulladdbox}
% Aaghh, the full addresses! A perennial source of problems. 
%    \begin{macrocode}
\newlength{\@fulladdheight}
\newlength{\@fulladddepth}
\newsavebox{\@fulladdbox}
%% \fulladdresses command modified (11 May 2000) so that we can put
%% \fulladdresses{} in the article to suppress the output of the
%% addresses
\newcommand{\@FinalFigure}{\relax}
\newcommand{\FinalFigure}[1]{\renewcommand{\@FinalFigure}{#1}}
\newcommand{\fulladdresses}[1]{%
  \renewcommand{\@rmaaArg}{#1}
  \ifx\@rmaaArg\@rmaaEmpty
     \renewcommand{\@fulladdresses}{}
  \else
     \renewcommand{\@fulladdresses}{
       \normalsize%
       \ifx\@FinalFigure\relax\else\vfill\@FinalFigure\vfill\fi
%%       \vspace*{1ex}
       %% do addresses as a list (copied from thebibliography environment) 
       \list{\null}{\leftmargin 1.5em\topsep\z@\labelwidth\z@\labelsep\z@\itemindent -1.5em
         \itemsep\z@\parsep\z@\usecounter{enumi}}
       \sloppy\clubpenalty4000\widowpenalty4000
       \sfcode`\.=1000\relax
       #1
       \endlist
       }
  \fi
  }
\newcommand{\@fulladdresses}{Full addresses go here}
%    \end{macrocode}
% \begin{macro}{\suppressfulladdresses}
% New user-level command for suppressing full addresses should you so
% desire. Currently used in the single-page-poster style. 
%    \begin{macrocode}
\newcommand{\suppressfulladdresses}{\@suppressfulladdressestrue\@needaddressesfalse}
%    \end{macrocode}
% \end{macro}
% \end{macro}
% \end{macro}
% \end{macro}
% \end{macro}
% \end{macro}

% \begin{macro}{\printdates}
% \begin{macro}{\noprintdates}
% Command to activate the printing out of the received/accepted
% dates. This can also be done using the \texttt{printdates} option,
% which is the default for \texttt{journal}, but not for \texttt{proceedings}.
%    \begin{macrocode}
\newcommand{\printdates}{\@printdatestrue}
\newcommand{\noprintdates}{\@printdatesfalse}
%    \end{macrocode}
% \end{macro}
% \end{macro}
% \begin{macro}{\@outputaddressesifnecessary}
% Command to output the addresses if we haven't already done so
% This is called by: 
% \begin{enumerate}
% \item The end of |thebibliography| environment
% \item The |\outputfulladdresses| command
% \end{enumerate}
%    \begin{macrocode}
\newcommand{\@outputaddressesifnecessary}{%
  \ClassInfo{rmaa}{Calling @outputaddressesifnecessary}
  \ifthenelse{\boolean{@needaddresses}}{%
    %% If in one column style, or in the first column of a two column
    %% style then we can try to put it on this page 
    \ifthenelse{\boolean{@firstcolumn} \or \not\boolean{@rmtwocol}}{%
      \@outputaddresses}{%
      %% otherwise give warning it will be on a separate page
      \ClassWarning{rmaa}{Full addresses will go on a separate
        page (please see User Guide for how to remedy this)}}%
    %% If we haven't managed yet, put it on its own page
    \ifthenelse{\boolean{@needaddresses}}{%
      \if@debug{\ClassWarning{rmaa}{In this silly bit!}}\fi
      \clearpage\vfill\@outputaddresses}{}%
    }{
    %% If we don't need them, we still might need to do the eoa hook 
    \ifthenelse{\boolean{@suppressfulladdresses}}{%
      \if@debug{\ClassWarning{rmaa}{XXX}}\fi
%%      \if@book\else\@endofarticlehook\null\fi
      }{}
    }%
  }  
%    \end{macrocode}
% \end{macro}

% \begin{macro}{\outputfulladdresses}
% A user interface to the above command, in case there is no biblio
% (11 May 2000) 
%    \begin{macrocode}
\newcommand{\outputfulladdresses}{\@outputaddressesifnecessary}
%    \end{macrocode}
% \end{macro}

% \begin{macro}{\adjustfinalcols}
% This command starts a new column and should be 
% used to ensure that the columns balance on the last page. It also
% outputs the full 
% author addresses to ensure that this occurs in the left hand
% column. 
% The |\par| command is so that the current line will be taken account
% of when working out |\pagetotal|
%    \begin{macrocode}
\newcommand{\adjustfinalcols}{\if@twocolumn\goodbreak\@outputaddresses\newpage\noindent\fi}
%    \end{macrocode}
% \end{macro}

% \begin{macro}{\@outputaddresses}
% Command to do the work of outputting the addresses. 
% This can be called by 
% \begin{enumerate}
% \item an explicit |\adjustfinalcols|
% \item via |\@outputaddressesifnecessary|, q.v.
% \end{enumerate}
%
%    \begin{macrocode}
\newcommand{\@outputaddresses}{%
  % Allow option of suppressing the full addresses, in which case
  % we just run the hook if we are not in a book style
  \if@suppressfulladdresses
%%     \if@book\else\@endofarticlehook\null\fi
  \else
  %% save the full addresses in a box for output later
  \sbox{\@fulladdbox}{\parbox{\textwidth}{\@fulladdresses}}
  % calculate height of address box  
  \settoheight{\@fulladdheight}{\usebox{\@fulladdbox}}
  \settodepth{\@fulladddepth}{\usebox{\@fulladdbox}}
  \addtolength{\@fulladdheight}{\@fulladddepth}
  \if@debug
  { \typeout{Address total height: \the\@fulladdheight}
    \typeout{Text height: \the\pagetotal}
    \typeout{Goal height: \the\pagegoal}
    }
  \fi
  % reuse \@fulladdheight to mean the total height of text+addresses
  \addtolength{\@fulladdheight}{\pagetotal}
%%  \addtolength{\@fulladdheight}{-\baselineskip}
  \if@debug\typeout{Adjusted text + addresses height: \the\@fulladdheight}\fi
  % check it will fit on current page
  \ifdim\@fulladdheight<\pagegoal
     %% must shift it to the left
     %% to compensate for the hanging indentation in the reference list. 
     %% NOTE: this is not foolproof!
     %% NB. This will give overfull hbox warnings, which is a bit ugly.
     \vfill\noindent\if@biblist\hspace*{-\leftmargin}\fi
     \usebox{\@fulladdbox}
     % Set this to show we've done it
     \@needaddressesfalse
     % and assume we are at the end of the article
     % so that we do the EOA hook if not in the book style
%%     \if@book\else\@endofarticlehook\fi
  \else
     % if not, just do nothing. With luck, this function will get 
     % called again on the next page
     \if@debug{\typeout{Holding over full addresses to next page}}\fi
  \fi\fi
  }
%    \end{macrocode}
% \end{macro}

% \subsubsection{Commands for editors}
% \label{sec:editor-commands}
% Some user/editor commands for doing the TOC and index in the book
% option. 
% They are defined here rather in the book-specific part since they
% need to be always available in order to avoid breaking documents. 
% \begin{macro}{\listofauthors}
% \begin{macro}{\@rawauthors}
% \begin{macro}{\rmaaTOCauthors}
% The command |\listofauthors| formats an author list for use in the
% TOC. It should contain the same names as in the |\author| command,
% but without the affiliation garbage and using |\&| instead of
% |and|. 
%    \begin{macrocode}
\newcommand{\@rawauthors}{Raw Authors}
\newcommand{\rmaaTOCauthors}{\noexpand\rmaaTOCAuthorStyle{Author list for TOC}}
\newcommand{\listofauthors}[1]{%
  \renewcommand{\@rawauthors}{#1}
  \renewcommand\rmaaTOCauthors{\noexpand\rmaaTOCAuthorStyle{#1}}}
%    \end{macrocode}
% \begin{macro}{\indexauthor}\label{indexauthor}
% The command |\indexauthor| produces an index entry for each
% author. 
% This command must come \emph{after} the
% |\title| and the |\listofauthors| commands. Also, there should be one 
% |\indexauthor| command for \emph{each} author. We also add the
% author onto the end of the |\ADS@authors| macro and construct the
% |\ADS@authorinitial| by a most cunning method (see below), but only
% for the first author. 
% 
% No, No, NO! This was all wrong. Backed out completely: 
% \begin{quote}
% Completely rewritten (19 Dec 2002) since we now use the index class, so 
% as to be able to alphabetize accented characters. A side effect of this is
% that |\index| no longer expands out the argument after the |@|, so we 
% instead add all the |\index| commands to a list (|\indexauthor@deferlist|), 
% which is executed in the |\@begofarticlehook|. 
% 
% Note the stuff with |\toks0| is copied from Exercise 20.15 of the \TeX{}book---it is
% to make sure the latin1-accented name gets passed into the idx
% file. 
% \end{quote}
% We now go back to the old way of doing things, using an optional
% argument to the \cs{indexauthor} command that has the name without
% accents. 
% \changes{3.08}{2003/01/16}{Remerged changes from 3.06 that got
% dropped from 3.07 after a synchronization disaster. Basically, the
% accentless optional argument to \cs{indexauthor}. Must add the macro
% directories to my unison set-up to stop this sort of thing from happening.}
%    \begin{macrocode}
\newcommand\@isortkey{}
\newcommand\@emptykey{empty}
\newcommand\indexauthor[2][empty]{
  \renewcommand\@isortkey{#1}
  \ifx\@isortkey\@emptykey \renewcommand\@isortkey{#2}\fi
  \index{\@isortkey @{{\rmaaIndexAuthorStyle{#2}}\hskip1em plus1em minus0.4em
      \rmaaIndexTitle\hskip1em plus2em minus0.4em\penalty-100
      \rmaaTOCauthors}}
  \indexauthorADS{#2}
}
%%\def\indexauthor@deferlist{\relax}
%%\long\def\indexauthor#1{
%%  \def\@@name{#1} \toks0=\expandafter{\@@name}
%%  \toks1=\expandafter{\indexauthor@deferlist}
%%  \protected@xdef\indexauthor@deferlist{%
%%    \the\toks1
%%    \noexpand\index{\the\toks0@{\noexpand\rmaaIndexAuthorStyle{\the\toks0}\noexpand\hskip1em plus1em minus0.4em 
%%        \rmaaIndexTitle\noexpand\hskip1em plus2em minus0.4em\noexpand\penalty-100
%%        \rmaaTOCauthors}}
%%  }
%%  \indexauthorADS{#1}
%%  }
%%\xdef\indexauthor@deferlist{}
\newcommand\indexauthorADS[1]{%
  \if@firstauthor\protected@xdef\ADS@authorinitial{\firstletter{#1}}\fi
  \protected@xdef\ADS@authors{\ADS@authors\if@firstauthor\else ; \fi
    #1}
  \@firstauthorfalse
}
%    \end{macrocode}
% Note that the first |#1| sets the key for the index, making sure it
% actually starts with  
% first letter of the surname. Meanwhile, the |@{...}| expression,
% which starts with some  
% control sequences is what is written to the |.ind| file. 
% See the makeindex man page for details. 
% \end{macro}
% \end{macro}
% \end{macro}
% \end{macro}

% \begin{macro}{\KillAllAccents}
% \begin{macro}{\firstletter}
% These are utility functions used in the above. |lett| is a counter
% indicating the current letter in the string. |\firstletter| just
% prints the first letter of its argument.
%    \begin{macrocode}
\newcommand{\KillAllAccents}{\def\'{}\def\`{}\def\~{}\def\.{}\def\i{i}}
\newcommand\firstletter[1]{\@car #1\@nil}
%    \end{macrocode}
% \end{macro}
% \end{macro}
% \begin{macro}{\SetVolume}
% \begin{macro}{\SetMSnumber}
% \begin{macro}{\SetFirstPage}
% \begin{macro}{\SetYear}
% \begin{macro}{\ReceivedDate}
% \begin{macro}{\AcceptedDate}
% These are for the editors to set the volume number, etc. 
%    \begin{macrocode}
\newcommand{\SetVolume}[1]{%
  \renewcommand{\@volume}{#1}
  \edef\ADS@volume{\PadToFourWithDots{#1}}
  \renewcommand{\@volumen}{\rmaa@VolumenStyle{volumen #1}}%
  }
\newcommand{\SetMSnumber}[1]{\renewcommand{\@MSnum}{#1}}
%    \end{macrocode}
% Note that there is no real reason to use |\SetFirstPage| since all
% the page numbering is done automatically in the book style. 
%    \begin{macrocode}
\newcommand{\SetFirstPage}[1]{\setcounter{page}{#1}}
\newcommand{\SetYear}[1]{%
  \renewcommand{\@year}{#1}
  \renewcommand{\ADS@year}{#1}
  \renewcommand{\@anyo}{\rmaa@VolumenStyle{#1}}
  }
\newcommand{\ReceivedDate}[1]{\renewcommand{\@received}{#1}}
\newcommand{\AcceptedDate}[1]{\renewcommand{\@accepted}{#1}}
%    \end{macrocode}
%
% \begin{macro}{\DOI}
% User/editor command to enter the DOI, which will be typeset as a link.
% \changes{4.0}{2019/11/02}{New user command \cs{DOI} to typeset
% DOI in the header of first page and add it to ADS info.}
%    \begin{macrocode}
\newcommand\DOI[1]{%
  \def\@DOI@url{\url{https://doi.org/10.22201/ia.01851101p.#1}}%
  \def\ADS@doi{https://doi.org/10.22201/ia.01851101p.#1}%
}%
%    \end{macrocode}
% \end{macro}
%
% \begin{macro}{\ORCIDiD}
% \changes{4.0}{2019/11/12}{Add mechanism for giving ORCID iD links}
% \begin{macro}{\ORCIDiDfull}
% \begin{macro}{\@orcid@icon}
% \begin{macro}{\@orcid@iconbig}
% Commands to add icon links to ORCID iDs of authors. For example
% \begin{quote}
% |\ORCIDiD|\marg{0000-0000-0000-0000}
% \end{quote}
% to be used after each author name in the |\author|
% command and in the |\fulladdresses|.  Alternatively,
% |\OORCIDiDfull| will typeset the icon plus the full text of the
% link, but we don't use that at present.  The icon image is stored in |ORCIDiD_iconvector.pdf|, which needs to be on the input path.  In principle, we could make EPS versions if necessary.
%    \begin{macrocode}
\newcommand\@orcid@icon{\includegraphics[height=1.8ex]{ORCIDiD_iconvector}}
\newcommand\@orcid@iconbig{\includegraphics[height=2ex]{ORCIDiD_iconvector}}
\newcommand\ORCIDiD[1]{\href{https://orcid.org/#1}{\,\@orcid@icon}}
\newcommand\ORCIDiDfull[1]{%
  \href{https://orcid.org/#1}{\@orcid@iconbig}\,\url{https://orcid.org/#1}}
%    \end{macrocode}
% \end{macro}
% \end{macro}
% \end{macro}
% \end{macro}
%
% \begin{macro}{\@volume}
% \begin{macro}{\@volumen}
% \begin{macro}{\@artpages}
% \begin{macro}{\@year}
% \begin{macro}{\@anyo}
% \begin{macro}{\@received}
% \begin{macro}{\@accepted}
% \begin{macro}{\@MSnum}
% Default values for the containers used by above commands, plus a
% few more useful ones. 
%    \begin{macrocode}
\newcommand{\@volume}{00}
\newcommand{\@volumen}{\rmaa@VolumenStyle{00}}
\newcommand{\@artpages}{%
  \oldpageref{\articlename:RMAAFirstPage}--\oldpageref{\articlename:RMAALastPage}}
\newcommand{\@year}{\the\year}
\newcommand{\@anyo}{\rmaa@VolumenStyle{2007}}
\newcommand{\@journalname}{\rmaa@JournalNameStyle{%
    \if@proceedings\RMAASClong\else\RMAAlong\fi}}
\newcommand{\@received}{}
\newcommand{\@accepted}{}
\newcommand{\@MSnum}{0000}
\newcommand{\@DOI@prefix}{https://doi.org/10.22201/ia.01851101p}
\def\@DOI@url{\url{https://doi.org/10.22201/ia.01851101p.2099.00.00.00}}
%    \end{macrocode}
% \end{macro}
% \end{macro}
% \end{macro}
% \end{macro}
% \end{macro}
% \end{macro}
% \end{macro}
% \end{macro}
% \end{macro}
% \end{macro}
% \end{macro}
% \end{macro}
% \end{macro}
% \end{macro}
% 
% The following commands are most useful in the \texttt{book}
% style. However, the |\SetConfTitle| command can now be used in
% individual articles to set the header for the \texttt{preprint}
% style. 
%    \begin{macrocode}
\newcommand{\@reunionplace}{Lugar de la reuni\'on}
\newcommand{\@reuniondate}{Fecha de la reuni\'on}
\newcommand{\@conftitle}{Conference Title}
\newcommand{\SetReunion}[2]{\renewcommand{\@reunionplace}{#1}\renewcommand{\@reuniondate}{#2}}
\newcommand{\SetConfTitle}[1]{\renewcommand{\@conftitle}{#1}}
\newcommand{\@editors}{Juan Ram\'on de la Fuente and ``El Mosh''}
\newcommand{\SetEditors}[1]{\renewcommand{\@editors}{#1}}
\newcommand{\@editores}{Juan Ram\'on de la Fuente y ``El Mosh''}
\newcommand{\SetEditores}[1]{\renewcommand{\@editores}{#1}}
\newcommand{\@mes}{mes de la publicaci\'on}
\newcommand{\SetMes}[1]{\renewcommand{\@mes}{\rmaa@VolumenStyle{#1}}}
%    \end{macrocode}

% \subsection{Notes for Author Proofs and Communication between Editors}
% \label{sec:notes}
% We introduce a command \cs{Note} that should be used to describe
% any change made to the manuscript. These are a bit tricky to get
% right. In general we use \cs{marginpar} but this only works in outer
% paragraph mode, so we just bung the box where we are if we are in a
% figure caption or something. If we are in the preamble, then we just
% use a \marginpar as part of the \maketitle command.
%    \begin{macrocode}
\newif\if@arenotes
\@arenotesfalse
\newcounter{notecounter}
\setcounter{notecounter}{0}
\newcommand\NotesList{}
\newcommand\Note[1]{%
  \if@proof{%
    \@arenotestrue%
    \stepcounter{notecounter}%
    \if@inpreamble\else{%
      \NoteMark{\arabic{notecounter}}%
    }\fi
    \protected@xdef\NotesList{\NotesList 
      \protect\item #1}%
  }\fi%
}
\newcommand\Changed[2]{Changed ``#1'' to ``#2''.}
\newcommand\NoteMark[1]{%
  \ifinner
  \raisebox{1ex}[0pt][0pt]{\makebox[0pt][l]{\framebox{\footnotesize\bfseries
        N#1}}}%
  \else
  \marginpar{\framebox{\footnotesize\bfseries N#1}}%
  \fi
}
\newcommand\RefNote[1]{\if@proof\NoteMark{\ref{#1}}\relax\fi}
\newcommand\NoteLabel{\protect\label}
\newcommand\OutputNotes{%
  \bgroup
  \onecolumn
  \section*{EDITORIAL NOTES}
  \begin{enumerate}
    \NotesList
  \end{enumerate}
  \egroup
}
\DeclareRobustCommand\cs[1]{\texttt{\char`\\#1}}

%    \end{macrocode}
% \subsection{Bibliographic Records for ADS}
% \label{sec:ADS}
% \changes{2.2beta11}{2001/11/27}{Just about finished all the ADS
% stuff now. Except for the references, which I have given up on for
% the time being.}
% \begin{macro}{\ADS@BibCode}
% \begin{macro}{\ADS@year}
% \begin{macro}{\ADS@journal}
% \begin{macro}{\ADS@volume}
% \begin{macro}{\ADS@firstpage}
% \begin{macro}{\ADS@authorinitial}
% We use the |\ADS@BibCode| command to hold the bibliographic code for
% ADS. It is 
% used in the |%R| 
% field of the bibliographic record for each article. 
% The format is 
% \begin{quote}
% \texttt{YYYYJJJJJVVVVMPPPPA}  
% \end{quote}
% where \texttt{YYYY} is the year, \texttt{JJJJJ} is the journal code,
% which is \texttt{RMxAA} or \texttt{RMxAC}, \texttt{VVVV} is the
% volume number (right-justified), \texttt{M} is not used by us,
% \texttt{PPPP} is the first page number (right-justified), and
% \texttt{A} is the initial of the first author's surname. The volume
% and page fields are padded with periods (`\texttt{.}'), which causes
% us some fun (this is done with the macro |\PadToFourWithDots| described
% below) . 
% 
% All these individual bits are held in the commands |\ADS@year|,
% |\ADS@journal|, |\ADS@volume|, |\ADS@firstpage|,
% |\ADS@authorinitial|. 
%    \begin{macrocode}
\newcommand{\SetADSBibCode}[1]{\renewcommand{\ADS@BibCode}{#1}}
\newcommand{\ADS@BibCode}
{\ADS@year\ADS@journal\ADS@volume.\PadToFourWithDots{\ADS@firstpage}\ADS@authorinitial}
\newcommand{\ADS@year}{2001} 
\newcommand{\ADS@journal}{\if@proceedings RMxAC\else RMxAA\fi}
\newcommand{\ADS@volume}{....} 
%    \end{macrocode}
% These two are just set to default values here, |\ADS@firstpage| is set
% in the |document| environment for included articles, in the
% |frontsection| environment, |\tableofcontents| command,  |index| and
% |abstracts| environments. Also by |\AtBeginDocument| for stand-alone
% articles. 
% 
% God knows how we are going to do |\ADS@authorinitial|!
%    \begin{macrocode}
\newcommand{\ADS@firstpage}{....}
\newcommand{\ADS@authorinitial}{A}
%    \end{macrocode}
% \end{macro}
% \end{macro}
% \end{macro}
% \end{macro}
% \end{macro}
% \end{macro}
% We also have a whole slew of other fields that can be set. These are
% all listed in Table~\ref{tab:ADS}. Most of them still have to be
% defined. 
% \begin{table}[!ht]\raggedleft
% \caption{ADS Bibliographic Records}\label{tab:ADS}
% \begin{tabular}{lp{8em}lp{8em}}\hline
% Code & Description & Macro Name & Comment \\ \hline
% |%R| & Bibliographic Code & |\ADS@BibCode| & required \\
% |%A| & Author list & |\ADS@authors| & required, generate from |\indexauthor| \\
% |%F| & Author affiliation & |\ADS@affil| & format: AA(\dots)
% AB(\dots) \dots \\
% |%J| & Journal Name & |\ADS@journalname| &  \\
% |%D| & Publication Date & |\ADS@fulldate| &  \\
% |%L| & Last Page & |\ADS@lastpage| & generate from label \\
% |%T| & Title & |\ADS@title| & required, generate from |\title| \\
% |%C| & Abstract Copyright & |\ADS@copyright| & set to IA-UNAM \\
% |%U| & URL for electronic document & |\ADS@url| & set to PDF file \\
% |%d| & DOI for article & |\ADS@doi| & set by |\DOI| \\
% |%K| & Keywords & |\ADS@keywords| & set equal to |\@keywords| \\
% |%A| & Abstract Text & |\ADS@abstract| & set equal to |\@abstract| \\
% |%R| & References & |\ADS@references| & generate from the
% bibliography  \\ \hline
% \end{tabular}
% \end{table}
% A lot of these macros are set by other commands, as mentioned in the
% Table, but some need to be set  explicitly here. 
% 
% \begin{macro}{\ADS@doi}
% \changes{4.0}{2019/11/08}{Save DOI for ADS}
% The user/editor command \cs{DOI} is used to set this. 
%    \begin{macrocode}
\newcommand{\ADS@doi}{https://doi.org/10.22201/ia.01851101p.0000.00.00.00}
%    \end{macrocode}
% \end{macro}
% 
% \begin{macro}{\ADS@journalname}
% This
% does not contain any accents at the moment so as not to cause
% problems---this maybe should be changed though. 
%    \begin{macrocode}
\newcommand{\ADS@journalname}{%
  \if@proceedings\@conftitle\space (Eds. \@editors) \fi%
  Revista Mexicana de Astronom\'\i{}a y Astrof\'\i{}sica%
  \if@proceedings\space (Serie de Conferencias)\fi%
  \space Vol. \@volume, pp. \ADS@firstpage--\ADS@lastpage\space (\@year)%
  \space (http://www.astroscu.unam.mx/\~{}rmaa/)
  }
%    \end{macrocode}
% \end{macro}
% \begin{macro}{\ADS@keywords}
% \begin{macro}{\ADS@abstract}
% \begin{macro}{\ADS@title}
% |\ADS@abstract| is straightforward, since we already have the text
% saved as |\@abstract|. We could do the same with |\ADS@keywords| but
% to avoid unwanted formatting we explicitly set it in |\keywords|
% (p.~\pageref{keywords}). |\ADS@title| is just set to |\@rawtitle|,
% which we were already using for another purpose. 
%    \begin{macrocode}
\let\ADS@title=\@rawtitle
\newcommand{\ADS@keywords}{No keywords specified}
\newcommand{\ADS@abstract}{No abstract specified}
%    \end{macrocode}
% \end{macro}\end{macro}\end{macro}
% \begin{macro}{\ADS@authors}
% |\ADS@authors| wants to be in ``Other, A. N.'' format, but with
% semi-colons between authors. The most natural way to do this seems to
% be with |\indexauthor| (p.~\pageref{indexauthor}). We append each
% author to this macro, which means we must explicitly clear it for
% each new article (done in |\article|). .   
%    \begin{macrocode}
\newcommand{\ADS@authors}{}
%    \end{macrocode}
% \end{macro}
% \begin{macro}{\ADS@fulldate}
% \begin{macro}{\ADS@lastpage}
% \begin{macro}{\ADS@copyright}
% \begin{macro}{\ADS@website}
% \begin{macro}{\ADS@url}
% \begin{macro}{\SetWebsite}
% \begin{macro}{\SetMonthNumeric}
% Other sundry data, together with user commands for bits that can be
% set by the Editors. The |\ADS@url| is constructed from the website
% (currently |http://www.astrosmo.unam.mx/rmaa/|,
% the journal descriptor (e.g., |RMxAC..10/|), followed by the
% path to the PDF file (e.g., |PDF/RMxAC..10_HEW1.pdf|). 
% 
%    \begin{macrocode}
\newcommand{\ADS@website}{http://www.astrosmo.unam.mx/rmaa/}
\newcommand{\ADS@url}
{\ADS@website\ADS@journal\ADS@volume\ADS@number/PDF/\ADS@journal\ADS@volume\ADS@number\_\articlename.pdf}
\newcommand{\ADS@number}{}
\newcommand{\ADS@fulldate}{00/\ADS@year} 
\newcommand{\ADS@lastpage}{00} 
\newcommand{\ADS@copyright}{(C) \ADS@year; Instituto de Astronom\'\i{}a,
  Universidad Nacional Aut\'onoma de M\'exico} 
\newcommand{\SetNumber}[1]{\renewcommand{\ADS@number}{-#1}}
\newcommand{\SetWebsite}[1]{\renewcommand{\ADS@website}{#1}}
\newcommand{\SetMonthNumeric}[1]{\renewcommand{\ADS@fulldate}{#1/\ADS@year}}
%    \end{macrocode}
% \end{macro}\end{macro}\end{macro}\end{macro}\end{macro}\end{macro}\end{macro}
% \begin{macro}{\ADS@affil}
% \begin{macro}{\ADS@references}
% These two are somewhat more problematical. Affiliation currently
% only works for the first author (see \dots) pending redesign of the
% user interface. References are now done by wrapping the
% |thebibliography| environment to open another input stream on the
% same file and search through for the |\bibitems|. 
%    \begin{macrocode}
\newcommand{\ADS@references}{}
\newcommand{\ADS@affil}{}
%    \end{macrocode}
% \end{macro}\end{macro}
% 
% Define two macros that expand to the strings
% |\begin{thebibliography}| and |\end{thebibliography}|. Note that the
% idea for this and much else is inspired by the code for the docstrip
% package. Also define |\StartBibX| in case authors use the more
% standard |\begin{thebibliography}{}| syntax. 
%    \begin{macrocode}
\edef\EndBib{\expandafter\@gobble\string\\end\string{thebibliography\string}}
\edef\StartBib{\expandafter\@gobble\string\\begin\string{thebibliography\string}}
\edef\StartBibX{\expandafter\@gobble\string\\begin\string{thebibliography\string}\string{\string}}
%    \end{macrocode}
% Define two new logical variables. |\ifAnotherline| is true if we
% want to read another line from the input stream. It is set to false
% when we encounter the |\end{thebibliography}|. |\ifInBib| is set to
% true when we frind the |\begin{thebibliography}| and tells us to
% start saving the following lines. 
%    \begin{macrocode}
\newif\ifAnotherLine
\newif\ifInBib
%    \end{macrocode}
% \begin{macro}{makeOther}
% We need a macro to change the catcodes of lots of things to other. 
%    \begin{macrocode}
\def\makeOther#1{\catcode`#1=12\relax}
%    \end{macrocode}
% \end{macro}
% \begin{macro}{\RereadBiblio}
% Define a macro to do all the fiddly stuff. This is used at the
% beginning of |thebibliography| environment (Sec~\ref{sec:biblio}).
%    \begin{macrocode}
\newread\tf@biblio
\newcommand{\RereadBiblio}{%
%    \end{macrocode}
% Open the current file again for reading
%    \begin{macrocode}
  \begingroup
  \immediate\openin\tf@biblio=\articlename
%    \end{macrocode}
% Everything in a group to be safe. Initialize logicals. Black magic
% to change
% catcodes and avoid blank lines being interpreted as |\par|. 
%    \begin{macrocode}
  \AnotherLinetrue\InBibfalse
  \makeOther\ \makeOther\\\makeOther\$%
  \makeOther\#\makeOther\^\makeOther\^^K%
  \makeOther\_\makeOther\^^A\makeOther\%%
  \makeOther\~\makeOther\{\makeOther\}\makeOther\&%
  \endlinechar-1\relax
%    \end{macrocode}
% Start looping through the file. Read each line to the macro
% |\Line|. If we have got to EOF or end of biblio then set
% |\AnotherLinefalse|.
%    \begin{macrocode}
  \loop
  \read\tf@biblio to \Line
  \ifx\Line\EndBib\AnotherLinefalse\InBibfalse\fi
  \ifeof\tf@biblio \AnotherLinefalse\fi
%    \end{macrocode}
% If we are in the biblio and the line is not blank then append it to
% the |\ADS@references| macro together with a newline. 
%    \begin{macrocode}
  \ifInBib
  \ifx\Line\empty\else\xdef\ADS@references{\ADS@references\Line^^J}\fi
  \fi
%    \end{macrocode}
% Check if the line we just saw was the start of biblio, then repeat
% the loop if necessary. 
%    \begin{macrocode}
  \ifx\Line\StartBib\InBibtrue\fi
  \ifx\Line\StartBibX\InBibtrue\fi
  \ifAnotherLine\repeat
  \immediate\closein\tf@biblio
  \endgroup
  }
%    \end{macrocode}
% \end{macro}
% 
% \begin{macro}{\PadToFourWithDots}
% Utility function for left-padding a number to 4 characters with
% periods. Used in 
% constructing |\ADS@BibCode|. This is a very simple algorithm and
% \emph{requires} that the argument actually be a number---if it is
% not then dire consequences will result.  
%    \begin{macrocode}
\newcommand{\PadToFourWithDots}[1]{%
  \ifnum #1 > 999 #1%
  \else\ifnum #1 > 99 .#1%
  \else\ifnum #1 > 9 ..#1%
  \else ...#1%
  \fi\fi\fi}
%%  \ifnum \count0 > 999 #1%
%%  \else\ifnum \count0 > 99 .#1%
%%  \else\ifnum \count0 > 9 ..#1%
%%  \else ...#1%
%%  \fi\fi\fi}
%    \end{macrocode}
% \end{macro}
%
% \begin{macro}{\AuthorADS}
% \begin{macro}{\ADS@affiliations}
% New mechanism for setting the affiliations for ADS. The author labels (AA, AB, \dots) are generated using two nested counters.  We use the |\robustify| command from \texttt{etoolbox} so that accented characters in the affiliation addresses do not blow up.  The list of accents may need to be expanded. 
%    \begin{macrocode}
\newcounter{AUTHOR}
\newcounter{author}[AUTHOR]
\stepcounter{AUTHOR}
\newcommand\ADS@affiliations{}
\robustify{\'}\robustify{\i}\robustify{\~}\robustify{\`}
\def\AuthorADS#1#2{%
  \begingroup
  \def\ORCIDiD##1{<ID system="ORCID">##1</ID>}
  \def\Email##1{<EMAIL>##1</EMAIL>}
  \def\Affil##1{##1}
  \stepcounter{author}
  \ifthenelse{\value{author} > 26}{\stepcounter{AUTHOR}\stepcounter{author}}{}
  \ifthenelse{\equal{\ADS@affiliations}{\@rmaaEmpty}}{%
    \xdef\ADS@affiliations{\Alph{AUTHOR}\Alph{author}(#2)}
  }{%
    \xdef\ADS@affiliations{\ADS@affiliations, \Alph{AUTHOR}\Alph{author}(#2)}
  }%
  \endgroup
}
%    \end{macrocode}
% \end{macro}
% \end{macro}
% %
% \begin{macro}{\ADS@dumpdata}
% \changes{4.0}{2019/11/13}{Write all author affiliations for ADS}
% \begin{macro}{\ADS@write}
% \begin{macro}{\ADS@openfile}
% |\ADS@dumpdata| is the macro that actually does the job of writing out all the
% data to an auxiliary file, using |\ADS@write| to write each
% field. The file is set up with |\ADS@openfile|, which currently is
% done in the TOC of a book.  The new multi-author affiliations are used if |\ADS@affiliations| has been set, otherwise it falls back to the single-author version |\ADS@affil|. 
%    \begin{macrocode}
\newcommand{\ADSfield}[2]{\relax}
\newcommand{\ADSvfield}[2]{\relax}
%%\renewcommand{\ADSfield}[12]{\%#1 #2}
\newcommand{\ADS@dumpdata}
{ 
  \ADS@write{\string\begin{ADSarticle}}
  \ADS@WRITE{\string\ADSfield{R}{\ADS@BibCode}}
  \ADS@write{\string\ADSfield{A}{\ADS@authors}}
  \ifthenelse{\equal{\ADS@affiliations}{\@rmaaEmpty}}{%
    {\ADS@sanitize\ADS@write{\string\ADSfield{F}{AA(\ADS@affil)}}}
  }{%
    {\ADS@sanitize\ADS@write{\string\ADSfield{F}{\ADS@affiliations}}}
  }
  \ADS@write{\string\ADSfield{J}{\ADS@journalname}}
  \ADS@write{\string\ADSfield{D}{\ADS@fulldate}}
  \ADS@write{\string\ADSfield{L}{\ADS@lastpage}}
  {\ADS@sanitize\ADS@write{\string\ADSfield{T}{\ADS@title}}}
  \ADS@write{\string\ADSfield{C}{\ADS@copyright}}
  \ADS@write{\string\ADSfield{d}{\ADS@doi}}
  \ADS@write{\string\ADSfield{U}{\ADS@url}}
  \ADS@write{\string\ADSfield{K}{\ADS@keywords}}
  {\ADS@sanitize\ADS@write{\string\ADSfield{B}{\ADS@abstract}}}
  \ADS@write{\string\ADSfield{Z}{\ADS@references}}
  \ADS@write{\string\end{ADSarticle}}
  }
%%\newcommand{\ADS@write}[1]{\if@book\addtocontents{ads}{#1}\fi}
%%\newcommand{\ADS@openfile}{\@starttoc{ads}}
\newcommand{\ADS@write}[1]{\set@display@protect\immediate\write\tf@ads{#1}}
\newcommand{\ADS@WRITE}[1]{\immediate\write\tf@ads{#1}}
\newcommand{\ADS@openfile}{%
  \newwrite\tf@ads%
  \immediate\openout\tf@ads\jobname.ads\relax
  \ADS@WRITE{\ADS@preamble}}
\newcommand{\ADS@closefile}{\ADS@WRITE{\ADS@postamble}}
\def\ADS@preamble{%
  %%  \string\documentclass{article}
  %%  \string\begin{document}
  }
\def\ADS@postamble{%
  %%  \string\end{document}
  \string%\string% Local Variables:^J 
  \string%\string% mode: LaTeX^J
  \string%\string% mode: Font-lock^J
  \string%\string% End:^J
  }
%    \end{macrocode}
% \end{macro}
% \end{macro}
% \end{macro}
% 
% \subsection{Automatic calculation of the page range}
% \label{sec:page-range}
% Another issue that has cost me a lot of grief. Still, it finally
% seems to be sorted now. All hail Jeff Goldberg. 
%
% Cunning trick to get the page numbers from |\pageref|.  Put a label
% on the first and last page of the document using |\AtBeginDocument|
% and |\AtEndDocument|. .  This way won't work for the book option.
% In that case, we do it by defining |\@endofarticlehook| and
% |\@begofarticlehook| as parallels to |\AtEndDocument| and
% |\AtBeginDocument|.  Note that we have to defer adding
% |\@begofarticlehook| to |\AtBeginDocument| to make sure it comes
% after the redefinitions of |\ref| and |\label|.
%    \begin{macrocode}
\if@book\AtEndDocument{\ADS@closefile}\else
 %\AtBeginDocument{\@begofarticlehook}
\AtEndDocument{\@endofarticlehook}
\fi
%    \end{macrocode}
% \begin{macro}{\@begofarticlehook}
% Since we are getting more and more things that need to be done at
% the beginning of every article, frontsection, etc we define a sister
% hook to 
% |\@endofarticlehook|. This hook does the following: 
% \begin{enumerate}
% \item Labels the first page
% \item Sets |\ADS@firstpage|
% \item runs and then clears |\indexauthor@deferlist|
% \item resets the keyword counter
% \end{enumerate}
% It is executed in the following circumstances: 
% \begin{enumerate}
% \item At |\begin{document}| for standalone files
% \item At |\begin{document}| for each included file
% \item At beginning of front sections 
% \item At beginning of TOC
% \item At beginning of Author Index
% \item At beginning of abstracts
% \end{enumerate}
% For some reason, this is not always within the right group, so we
% use |\xdef| just to be sure. 
%    \begin{macrocode}
\newcommand{\@begofarticlehook}{%
  \if@frontmatter\frontlabel{RMAAFirstPage}\else\label{RMAAFirstPage}\fi%
  \xdef\ADS@firstpage{\@arabic{\c@page}}%
  \if@book\else\makeindex\fi % needed so default index is defined in stand-alone articles
%%  \indexauthor@deferlist
%%  \xdef\indexauthor@deferlist{}
  \@inpreamblefalse
  \setcounter{kwdcounter}{0}
  }
%    \end{macrocode}
% \end{macro}
% \begin{macro}{\@endofarticlehook}
% This is a hook to be run at the end of each article
% currently, it just labels the last page and writes out all the ADS
% data to an auxiliary file.  
% It should not produce any printed output. 
% It is modelled on the |\AtEndDocument| command, to which it reduces
% in the case of a single article. 
%    \begin{macrocode}
\newcommand{\@endofarticlehook}{%
  \if@debug{\ClassWarning{rmaa}{***Doing Hook***}}\fi
  \if@debug{\ClassWarning{rmaa}{ADSabstract is: \ADS@abstract}}\fi
  \clearpage\rmaa@lastpagelabel\if@book\ADS@dumpdata\fi
  \ifx\NotesList\@rmaaEmpty\else\OutputNotes\fi
  \if@debug{\ClassWarning{rmaa}{Notes list is: \NotesList}}\fi
}
%    \end{macrocode}
% \begin{macro}{\rmaa@lastpagelabel}
% New version of last page label command copied from Jeff Goldberg's
% lastpage.sty package - we can't just use |\label| here in case the
% output routine for the last page has already been called. Instead,
% we first explicitly clear the page, step back the page counter,
% explicitly write to the .aux file, then restep the page counter
% again.
%
% But this will conflict with the hyperref package, which does labels differently.  Compensate for that by using the five-argument form.
% 
% Also set the |\ADS@LastPage| macro to the appropriate number. 
%    \begin{macrocode}
\newcommand{\rmaa@lastpagelabel}{%
  \addtocounter{page}{-1}%
  \immediate\write\@auxout{\string
    \newlabel{\articlename:RMAALastPage}{{}{\thepage}{}{}{}}}%
  \edef\ADS@lastpage{\thepage}%
  \addtocounter{page}{1}}
%    \end{macrocode}
% \end{macro}
% \end{macro}

% \subsection{Deal with explicit change in number of columns}
% \label{sec:columns}
% Sometimes, when formatting a book, we need to switch between
% one-column articles and two-column articles. For this to work, we
% need to wrap the standard \LaTeX{} commands to set the |\@rmtwocol|
% flag appropriately and also reset the package default values of all
% scalable lengths. .  
%    \begin{macrocode}
\let\latexonecolumn=\onecolumn
\let\latextwocolumn=\twocolumn
\renewcommand{\onecolumn}{\@rmtwocolfalse\@package@default@lengths\latexonecolumn}
\renewcommand{\twocolumn}{\@rmtwocoltrue\@package@default@lengths\latextwocolumn}
%    \end{macrocode}
% 
% \subsection{The actual formatting of the title}
% \label{sec:title-format}

% Note, we have to redefine the |\maketitle| command to use the
% numbered footnotes. This means we needn't use the atbeginend stuff. 
%    \begin{macrocode}
\renewcommand{\maketitle}{\par
  \begingroup
%    \end{macrocode}
% Don't switch to marked footnotes. 
% This is consistent with existing style but I'm not sure I like it
%    \begin{macrocode}
  %%  \renewcommand\thefootnote{\fnsymbol{footnote}}%
%%  \def\@makefnmark{\hbox to\z@{$\m@th^{\@thefnmark}$\hss}}%
  \def\@makefnmark{\textsuperscript{\@thefnmark}\hspace{-0.2em}}%
  \long\def\@makefntext##1{\parindent 1em\noindent
    \hbox to1.8em{\hss$\m@th^{\@thefnmark}$}##1}%
  \if@rmtwocol
  { \if@debug{\ClassWarning{rmaa}{TWO COLUMN}}\fi
    \twocolumn[\@maketitle]%
    }
  \else
  { \if@debug{\ClassWarning{rmaa}{ONE COLUMN}}\fi
    \newpage
    \global\@topnum\z@   % Prevents figures from going at top of page.
    \if@manuscript \@MSmaketitle \else \@maketitle \fi
    }
  \fi
  \if@manuscript\else\thispagestyle{rmaafirstpage}\@thanks\fi
  \ifx\NotesList\@rmaaEmpty\else{%
    \NoteMark{1--\arabic{notecounter}}\vspace*{-2\baselineskip}%
  }\fi%
  \endgroup
%    \end{macrocode}
% Don't reset footnote counter (not sure whether this is good thing
% or not).
%    \begin{macrocode}
  %%  \setcounter{footnote}{0}%
  \let\thanks\relax\let\affil\relax\let\altaffil\relax
  \let\maketitle\relax\let\@maketitle\relax
  \gdef\@thanks{}\gdef\@author{}\gdef\@title{}}
%    \end{macrocode}

% \begin{macro}{\@kwhwidth}
% \begin{macro}{\@kwtwidth}
% Redefine \cs{@maketitle} to include the abstract, resumen and
% keywords. Note that in the manuscripty style we use
% \cs{@MSmaketitle} instead.  
% \changes{3.16}{2003/04/07}{
%  Added extra vertical space between authors and accepted date.}
% \changes{3.26}{2011/08/15}{Add extra label to review articles}
%    \begin{macrocode}
\newlength{\@kwhwidth}\newlength{\@kwtwidth}
\newcommand\reviewname{Review}
\renewcommand{\@maketitle}{%
  \newpage
  \null
  \if@reviewarticle
  \smallskip
  {\large\bfseries\itshape\reviewname}%
  \bigskip
  \else
  \vskip\@before@title@skip%
  \fi
  \begin{center}%
    {\@title \par}%
    \vskip\@before@author@skip%
    \parbox{\rmaa@AuthorWidthFrac\textwidth}{\centering
      \lineskip\@between@author@skip%
      \@author
      \par}
    %% this should be replaced with `Submitted to...' for preprints
    \if@printdates
    \vskip\@before@dates@skip%
    \if@preprint
    \rmaa@ReceivedStyle{Submitted to RMxAA \@received}
    \else
    \rmaa@ReceivedStyle{Received \@received; accepted \@accepted}
    \fi
    \fi
    \if@manuscript
    \rmaa@ReceivedStyle{Submitted to RMxAA \@received}
    \fi
    \if@guide
    \rmaa@ReceivedStyle{Version \rmaa@version, \rmaa@date}
    \fi
    \vskip\@before@abstract@skip%
    \if@abstractpresent
    \parbox{\rmaa@AbstractWidthFrac\textwidth}{%
      \setlength{\parindent}{0.9cm}%
      \centerline{\rmaa@AbstractHeaderStyle{Abstract}}\par%
      \vskip\@after@abstracthead@skip%
      \if@newstyle\noindent\fi\@abstract
      \vskip\@between@abstract@skip%
      \centerline{\rmaa@AbstractHeaderStyle{Resumen}}\par%
      \vskip\@after@abstracthead@skip%
      \if@newstyle\noindent\fi\@resumen
      \if@setresumen\else\@abstract\fi % append english abstract if
                                       % spanish one not set
      \if@guide\else            %do not use keywords for guide
      \vskip\@before@kwd@skip%
      % Format the Key Words with hanging indentation:
      % They are put in a \parbox with width equal to the line width
      % minus the width of the header. 
      { % Widths of keyword
        \settowidth{\@kwhwidth}{\@keywordheader~}   % header and
        \setlength{\@kwtwidth}{\linewidth}          % keyword text
        \addtolength{\@kwtwidth}{-\@kwhwidth}       %
        \noindent\@keywordheader~\parbox[t]{\@kwtwidth}{%
          \rmaa@KeywordStyle{\boldmath\uppercase{\@keywords}}}}%
      \fi
      }
      \fi
  \end{center}%
  \vspace*{\afterkwdskip}
  \par
  }
\newlength{\afterkwdskip}\setlength{\afterkwdskip}{\z@}
%    \end{macrocode}
% \end{macro}
% \end{macro}
%
%
% \begin{macro}{\reviewarticle}
% \begin{macro}{\reviewtoc}
% \begin{macro}{\reviewlof}
%   User command to specify a review article. Must be used in
%   preamble. Main effect is to place the text ``Review'' before the
%   article title.  Also add command for setting a table of contents
%   and a list of figures at the beginning of the article.
%    \begin{macrocode}
\newcommand{\reviewarticle}{\@reviewarticletrue}
\newcommand\reviewtoc{%
  \section*{CONTENTS}%
  \@starttoc{rtc}%
}
\newcommand\reviewlof{%
  \section*{LIST OF FIGURES}%
  \@starttoc{lof}%
}
%    \end{macrocode}
% \end{macro}
% \end{macro}
% \end{macro}
% \begin{macro}{\l@section}
%    \begin{macrocode}
\renewcommand*\l@section[2]{%
  \ifnum \c@tocdepth >\z@
    \addpenalty\@secpenalty
    \addvspace{0.8em \@plus\p@}%
    \setlength\@tempdima{1.5em}%
    \begingroup
      \parindent \z@ \rightskip \@pnumwidth
      \parfillskip -\@pnumwidth
      \leavevmode %\bfseries
      \advance\leftskip\@tempdima
      \hskip -\leftskip
      #1\nobreak\hfil \nobreak\hb@xt@\@pnumwidth{\hss #2}\par
    \endgroup
  \fi}
%    \end{macrocode}
% \end{macro}
%
% \begin{macro}{\@MSmaketitle}
% A special version of |\@maketitle| for the \texttt{manuscript}
% option. This can deal with a multi-page abstract/resumen. Not very
% well, mind. 
%    \begin{macrocode}
\newcommand{\@MSmaketitle}{%
  \if@debug\ClassInfo{rmaa}{Using special MS maketitle command}\fi
  \newpage
  {
    \thispagestyle{rmaafirstpage}
%     \setlength{\parindent}{0pt}
    \null\vskip 2ex
    \begin{center}
      \@title\par
      \@author\par
      \rmaa@ReceivedStyle{%
        \ifx\@received\@rmaaEmpty
          Draft version: \@date
        \else
          Received: \@received
          \ifx\@accepted\@rmaaEmpty
          \else
            ; Accepted: \@accepted
          \fi
        \fi}\par
    \end{center}
    \@thanks
    \centerline{\rmaa@AbstractHeaderStyle{Resumen}}
    \noindent\@resumen\par\vskip 2ex
    \centerline{\rmaa@AbstractHeaderStyle{Abstract}}
    \noindent\@abstract\par\vskip 2ex
    \noindent\@keywordheader~\rmaa@KeywordStyle{\@keywords}\par
%     \@keywordheader~\rmaa@KeywordStyle{\uppercase{\@keywords}}\par
  }%
  }
%    \end{macrocode}
% \end{macro}

% \subsection{The sectioning commands}
% \label{sec:section-commands}

% \begin{macro}{\section}
% \begin{macro}{\subsection}
% \begin{macro}{\subsubsection}
% \begin{macro}{\paragraph}
% \begin{macro}{\subparagraph}
% Redefine all the standard commands to use these lengths. 
%    \begin{macrocode}
\renewcommand\section{\@startsection {section}{1}{\z@}%
                                   {\abovesectionskip}%
                                   {\belowsectionskip}%
                                   %% how do we get upper case here
                                   {\centering\reset@font\normalsize}}
\renewcommand\subsection{\@startsection{subsection}{2}{\z@}%
                                     {\abovesectionskip}%
                                     {\belowsectionskip}%
                                     {\centering\reset@font\normalsize\itshape}}
\renewcommand\subsubsection{\@startsection{subsubsection}{3}{\z@}%
                                     {\abovesectionskip}%
                                     {\belowsectionskip}%
                                     {\centering\reset@font\normalsize\itshape}}
\renewcommand\paragraph{\@startsection{paragraph}{4}{\z@}%
                                    {\aboveparaskip}%
                                    {\belowparaskip}%
                                    {\reset@font\normalsize\bfseries}}
\renewcommand\subparagraph{\@startsection{subparagraph}{5}{\parindent}%
                                      {\aboveparaskip}%
                                      {\belowparaskip}%
                                      {\reset@font\normalsize\bfseries}}
%    \end{macrocode}
% \end{macro}
% \end{macro}
% \end{macro}
% \end{macro}
% \end{macro}

% \begin{macro}{secnumdepth}
% Make sure that we number the subsubsections---this is done by
% default in the |article| class but not the |book| class, which is
% why we must set it explicitly here.
%    \begin{macrocode}
\setcounter{secnumdepth}{3}
%    \end{macrocode}
% \end{macro}
% \changes{2.2beta10}{2001/11/07}{Fixed \cs{secnumdepth} for \cs{book}
% style.}

% \begin{macro}{\thesubsection}
% \begin{macro}{\thesubsubsection}
% For subsection and subsubsection, we want the numbers to be roman
% but the title to be italic. This seems the easiest way to go about it.
%    \begin{macrocode}
\renewcommand{\thesubsection}{{\upshape\thesection.\@arabic\c@subsection}}
\renewcommand{\thesubsubsection}{\thesubsection{\upshape.\@arabic\c@subsubsection}}
%    \end{macrocode}
% \end{macro}
% \end{macro}

% \begin{macro}{\@seccntformat}
% Add ``.'' to the number here so we can remove it from |\thesection|,
% etc. Reduce gap between number and title WJH 27 Jan 1995
%    \begin{macrocode}
\def\@seccntformat#1{\csname the#1\endcsname .\hskip 0.4em}
%    \end{macrocode}
% \end{macro}

% \begin{macro}{\@sect}
% I didn't really want to mess with the |\@sect| command, but it was
% the only way I could manage to force upper case. Other than the
% commented parts, it is identical to the standard version. 
%    \begin{macrocode}
\def\@sect#1#2#3#4#5#6[#7]#8{
  \ifnum #2>\c@secnumdepth
     \let\@svsec\@empty
  \else
     \refstepcounter{#1}%
     \let\@@protect\protect
     \def\protect{\noexpand\protect\noexpand}%
     \edef\@svsec{\@seccntformat{#1}}%
     \let\protect\@@protect
  \fi
%    \end{macrocode}
% This is the bit I changed (27 Jan 1995). Use upper case for
% |\section|, use mixed case for the rest. 
%    \begin{macrocode}
  \ifnum #2=1                   
     \def\sec@upcase##1{\uppercase{##1}}%
     \def\sec@shape{\upshape}%
  \else                         
     \def\sec@upcase##1{\relax{##1}}%
     \def\sec@shape{\itshape}%
  \fi
  \@tempskipa #5\relax
  \ifdim \@tempskipa>\z@
     \begingroup #6\relax
          \@hangfrom{\hskip #3\relax\@svsec}%
                    {\interlinepenalty \@M \sec@upcase{#8}\par}%
     \endgroup
     \csname #1mark\endcsname{#7}%
%    \end{macrocode}
% Only write TOC entries if doing the Guide style. 
%    \begin{macrocode}
     \if@guide 
        \addcontentsline{toc}{#1}{%
          \ifnum #2>\c@secnumdepth \else
             \protect\numberline{\csname the#1\endcsname}%
          \fi
          #7}%
     \fi
%    \end{macrocode}
% Also allow per-article TOC for review papers
%    \begin{macrocode}
     \if@reviewarticle
        \ifnum #2>2 \else
           \addcontentsline{rtc}{#1}{%
             \protect\numberline{\csname the#1\endcsname}%
             \sec@upcase{\sec@shape #7}}%
        \fi
     \fi
  \else
     \def\@svsechd{#6\hskip #3\relax  %% \relax added 2 May 90
       \@svsec #8\csname #1mark\endcsname{#7}%
       \if@guide
          \addcontentsline{toc}{#1}{%
            \ifnum #2>\c@secnumdepth \else
               \protect\numberline{\csname the#1\endcsname}%
            \fi
            #7}%
       \fi
       }%
  \fi
  \@xsect{#5}}
%    \end{macrocode}
% \end{macro}


% \subsection{Appendices}
% \label{sec:appendices}
% Make sure that appendices in individual articles work properly. We
% copy the definition from the article class (the book one doesn't do
% what we want).  
% \begin{environment}{appendices}
%   This uses an environment |appendices| inside of which each
%   |\section| will become a separate appendix.  Note that the
%   appendices must come before the references (otherwise the full
%   addresses mechanism will not work properly). We make the equation
%   numbering be subordinate to the appendix (equations A12, B6,
%   etc). Finally, we have a medskip followed by APPENDICES to
%   introduce them. \changes{3.17}{2002/04/08}{ Added appendices environment
%   that does the right thing.}
%    \begin{macrocode}
\newenvironment{appendices}{\par
  \begingroup
  \setcounter{section}{0}%
  \setcounter{subsection}{0}%
  \renewcommand\thesection{\@Alph\c@section}
  \renewcommand\theequation{\thesection\arabic{equation}}
  \medskip
  \section*{APPENDICES}}
{\endgroup}
%    \end{macrocode}
% \end{environment}
% \begin{environment}{appendix}
%   If there is only one appendix, then use |appendix| environment
%   instead.  \changes{4.0}{2019/11/12}{ Added
%   appendix environment for case where there is only one section.}
%    \begin{macrocode}
\renewenvironment{appendix}{\par
  \begingroup
  \setcounter{section}{0}%
  \setcounter{subsection}{0}%
  \renewcommand\thesection{\@Alph\c@section}
  \renewcommand\theequation{\thesection\arabic{equation}}
  \medskip
  \section*{APPENDIX}}
{\endgroup}
%    \end{macrocode}
% \end{environment}

% \subsection{Floating thingies}
% \label{sec:floats}   

%% 
%    \begin{macrocode}
\renewcommand{\thetable}{\@arabic\c@table}
\renewcommand{\thefigure}{\@arabic\c@figure}
\renewcommand{\figurename}{Fig.}
\renewcommand{\tablename}{TABLE}
%\renewcommand{\fnum@table}{\MakeUpperCase{\tablename}~\thetable}
%    \end{macrocode}
% For some reason, I have to write `TABLE' explicitly here
% If I use |\tablename|, it gets converted to mixed case---God knows
% why. 
% |\figurename| is also being ignored. 
%    \begin{macrocode}
\renewcommand{\fnum@table}{TABLE~\thetable}
\renewcommand{\fnum@figure}{Fig.~\thefigure}
%    \end{macrocode}

% Do captions in 9 pt face in the main journal style and in new-style
% anything. 
%    \begin{macrocode}
\long\def\@makefigcap#1#2{%
  \if@proceedings\else\small\fi
  \if@newstyle\small\fi 
  \vskip\abovecaptionskip
  \sbox\@tempboxa{#1. #2}%
  \ifdim \wd\@tempboxa >\hsize
    #1. #2 \par
  \else
    \global \@minipagefalse
    \hb@xt@\hsize{\hfil\box\@tempboxa\hfil}%
  \fi
  \vskip\belowcaptionskip}
%    \end{macrocode}

% Make sure table caption is upper case. 
%    \begin{macrocode}
\long\def\@maketabcap#1#2{%
  \hbox to\hsize{\hfil #1 \hfil}
  \vskip 6pt 
  \centering\MakeUpperCase{#2}\par
  \vskip 6pt}
%    \end{macrocode}

%    \begin{macrocode}
\newlength{\@rottabcapwidth}
\long\def\@makerottabcap#1#2{%
  \setbox\@tempboxa\hbox{\MakeUpperCase{#2}}%
  \ifdim\wd\@tempboxa >.8\vsize
     \setlength\@rottabcapwidth{.8\vsize}
  \else
     \setlength\@rottabcapwidth{\wd\@tempboxa}
  \fi
  \rotatebox{90}{%
    \begin{minipage}{\vsize}\centering
      \begin{minipage}{\@rottabcapwidth}\centering
        #1\\ \vskip 6pt \MakeUpperCase{#2}
      \end{minipage}%
    \end{minipage}
    }%
  \hskip 6pt
  }

\newenvironment{Sideways}
    {\begin{sideways}\begin{minipage}{\vsize}\centering}
    {\end{minipage}\end{sideways}}

\newenvironment{Sidewaystable}
{%
  \let\oldcaption=\caption%
  \def\caption{
    \def\@captype{table}\renewcommand{\@makecaption}{\@maketabcap}
    \oldcaption}
  \begin{sidewaystable}\centering}
  {\end{sidewaystable}\let\caption=\oldcaption}
\newenvironment{Sidewaystable*}
{%
  \let\oldcaption=\caption%
  \def\caption{
    \def\@captype{table}\renewcommand{\@makecaption}{\@maketabcap}
    \oldcaption}
  \begin{sidewaystable*}\centering}
  {\end{sidewaystable*}\let\caption=\oldcaption}
%    \end{macrocode}

% Not sure why we need to stick these lines in. 
%    \begin{macrocode}
\renewenvironment{figure}
                 {\renewcommand{\@makecaption}{\@makefigcap}\@float{figure}}
                 {\end@float}
\renewenvironment{figure*}
                 {\renewcommand{\@makecaption}{\@makefigcap}\@dblfloat{figure}}
                 {\end@dblfloat}
\renewenvironment{table}
                 {\renewcommand{\@makecaption}{\@maketabcap}\@float{table}}
                 {\end@float}
\renewenvironment{table*}
                 {\renewcommand{\@makecaption}{\@maketabcap}\@dblfloat{table}}
                 {\end@dblfloat}
%    \end{macrocode}

% Redefine |\@caption| so that we don't bother with the list of
% figures. This saves us having to worry about fragile commands in the
% captions. 
%    \begin{macrocode}
\long\def\@caption#1[#2]#3{%
  \par
%%  \addcontentsline{\csname ext@#1\endcsname}{#1}%
%%    {\protect\numberline{\csname the#1\endcsname}{\ignorespaces #2}}%
  \begingroup
    \@parboxrestore
    \if@minipage
      \@setminipage
    \fi
    \normalsize
    \@makecaption{\csname fnum@#1\endcsname}{\ignorespaces #3}\par
  \endgroup}
%    \end{macrocode}

% Make tables not quite so cramped. Allow |\hline\hline| to give a
% poor man's thick rule. 
%    \begin{macrocode}
\renewcommand{\arraystretch}{1.15}
\setlength{\doublerulesep}{0.4pt}
%    \end{macrocode}

% \begin{macro}{\tablecols}
% \begin{macro}{\@tablecols}
% \begin{macro}{\tabcaptioncontinued}
% |\tablecols| is a user command for setting the number of columns. It
% is used by commands that try and span the entire table (I think
% there must be a better way of acheiving
% this). |\tabcaptioncontinued| is 
% a user command for setting the caption of a continued table on the
% next page. 
%    \begin{macrocode}
\newcommand{\tablecols}[1]{\renewcommand{\@tablecols}{#1}}
\newcommand{\@tablecols}{1}
\newcommand{\tabcaptioncontinued}{%
  \multicolumn{\@tablecols}{c}{\addtocounter{table}{-1}TABLE~\thetable\ (CONTINUED)} \\ \\
  }
%    \end{macrocode}
% \end{macro}
% \end{macro}
% \end{macro}

% \subsubsection{Footnotes in tables}
% \label{sec:table-footnotes}
% \begin{macro}{\tablenotemark}
% \begin{macro}{\tablenote}
% \begin{macro}{\tablenoteskip}
% Old user commands |\tablenotemark{MARK}|,
% |\tablenote{MARK}{TEXT}|---now deprecated. 
%    \begin{macrocode}
\newcommand\tablenotemark[1]{$^\mathrm{#1}$} 
\newcommand\tablenote[2]{\if@proceedings\else\small\fi ${}^\mathrm{#1}$ #2} 
\newcommand\tablenoteskip{\\[0.5\baselineskip]} 
%    \end{macrocode}
% \end{macro}
% \end{macro}
% \end{macro}

% \begin{macro}{\tabnotemark}
% \begin{macro}{\tabnotetext}
% \begin{macro}{\tabnotetextTWO}
% Alternative way of doing footnotes in tables. 
%    \begin{macrocode}
\newcommand{\tabnotemark}[1]{\textsuperscript{\lowercase{#1}}}
%\newcommand{\tabnotetextTWO}[4]{\tabnotetext{#1}{#2}\tabnotetext{#3}{#4}}
\newcommand{\tabnotetextTWO}[4]{%
  \multicolumn{\@tablecols}{l}{%
    \parbox{0.5\tabnotewidth}{\small\textsuperscript{#1}#2}%
    \parbox{0.5\tabnotewidth}{\small\textsuperscript{#3}#4}%
    }%
  \\}
\newcommand{\tabnotetext}[2]{%
  \multicolumn{\@tablecols}{@{}p{\tabnotewidth}@{}}{\small\textsuperscript{#1}#2}\\}
\newlength{\tabnotewidth}
%    \end{macrocode}
% \end{macro}
% \end{macro}
% \end{macro}


% \begin{macro}{changemargin}
%    \begin{macrocode}
% Command to allow margins to be changed on the fly. Useful to allow
% tables and Figures to protrude into the margins. 
%
% \changes{3.24}{2007/09/14}{Added \texttt{changemargin} environment
% (copied from TeX FAQ).}
\newenvironment{changemargin}[2]{% 
  % From TeX FAQ http://www.tex.ac.uk/cgi-bin/texfaq2html?label=chngmargonfly
  \begin{list}{}{%
    \setlength{\topsep}{0pt}%
    \setlength{\leftmargin}{#1}%
    \setlength{\rightmargin}{#2}%
    \setlength{\listparindent}{\parindent}%
    \setlength{\itemindent}{\parindent}%
    \setlength{\parsep}{\parskip}%
  }%
  \item[]}{\end{list}}
%    \end{macrocode}
% \end{macro}

% \subsection{The bibliography and references}
% \label{sec:biblio}
% \begin{macro}{\beforebiblioskip}
% \begin{macro}{\afterbiblioskip}
% I'm not entirely sure why I introduced these two commands, since
% directly putting |\vspace| before and after
% |\begin{thebibliography}| seems to work well enough. 
%    \begin{macrocode}
\newlength{\beforebiblioskip}
\setlength{\beforebiblioskip}{0pt}
\newlength{\afterbiblioskip}
\setlength{\afterbiblioskip}{0pt}
%    \end{macrocode}
% \end{macro}
% \end{macro}
% \begin{macro}{thebibliography}
% The environment itself is done as starred section with extra
% vertical space after heading. This space is currently zero. 
% In the book style we do the cunning trick with rereading the file in
% order to get the references into the |\ADS@references| macro.
% 
% \changes{3.05}{2003/01/02}{Removed the \cs{beforebiblioskip} and
% \cs{afterbiblioskip} commands, since they cause unwanted vertical
% space in some instances.}  
%    \begin{macrocode}
\renewenvironment{thebibliography}{%
  \if@book\RereadBiblio\fi
%%  \vspace*{\beforebiblioskip}
  \section*{REFERENCES} \frenchspacing
%%  \vspace*{\afterbiblioskip}
  \small
%    \end{macrocode} 
% The references themselves are set as a list in order 
% to get the hanging indentation. They are done in a smaller font and
% with french spacing so we needn't worry about the distinction
% between sentence-ending periods and abbreviations. 
%    \begin{macrocode} 
  \list{\null}{\leftmargin 1.5em\labelwidth\z@\labelsep\z@\itemindent -1.5em
    \itemsep\z@\parsep\z@\usecounter{enumi}}
  \sloppy\clubpenalty4000\widowpenalty4000
  \sfcode`\.=1000\relax\@biblisttrue}
{\endlist\@biblistfalse\normalsize\@outputaddressesifnecessary}
%    \end{macrocode}
% \end{macro}
% \subsection{Uniquification of cross-references}
% \label{sec:uniquification}
% 
% \begin{macro}{\articlename}
% In order that the |\label{}| command in individual articles be
% guaranteed to give a unique anchor, we append the article name (held
% in the |\articlename| command) to each key. This requires
% redefinition of the following commands: |\cite|, |\nocite|, |\ref|,
% |\pageref|, |\label|. 
% \begin{macro}{\oldcite}
% \begin{macro}{\oldnocite}
% \begin{macro}{\oldref}
% \begin{macro}{\oldpageref}
% \begin{macro}{\oldlabel}
% The original versions of these commands are
% saved as |\oldcite|, |\oldnocite|, etc. 
%
% The uniquifying of cross-references is now always done, so 
% we don't get multiply-defined labels in the single article case
%    \begin{macrocode}
\newcommand{\articlename}{anon}
\AtBeginDocument{
  \let\oldcite=\cite 
  \let\oldnocite=\nocite 
  \let\oldref=\ref 
  \let\oldpageref=\pageref
  \let\oldlabel=\label
}
%    \end{macrocode}
% \end{macro}
% \end{macro}
% \end{macro}
% \end{macro}
% \end{macro}
% \end{macro}
%
% Uniquify all cross-references by prepending the article name.  This
% now has to go into |\AtBeginDocument| so it is working on the new
% versions that hyperref.sty has set up.  Then we need to put the
% |\@begofarticlehook| afterwards, so that the first page label is using
% the new definitions.  If amsmath is loaded, then we also need to patch
% the special label command that they use. \changes{4.0}{2019/11/18}{Fix labels for amsmath display equations, in case that package has been loaded.}
%    \begin{macrocode}
\AtBeginDocument{
  \renewcommand{\ref}[1]{\oldref{\articlename:#1}%
    \if@debug{\ClassInfo{rmaa}{Reference to \articlename:#1}}\fi}
  \renewcommand{\pageref}[1]{\oldpageref{\articlename:#1}%
    \if@debug{\ClassInfo{rmaa}{Page reference to \articlename:#1}}\fi}
  \renewcommand{\label}[1]{\oldlabel{\articlename:#1}%
    \if@debug{\ClassInfo{rmaa}{Label for \articlename:#1}}\fi}%
  \ifdef{\label@in@display}{%
    \let\oldlabel@in@display=\label@in@display
    \renewcommand{\label@in@display}[1]{\oldlabel@in@display{\articlename:#1}%
      \if@debug{\ClassInfo{rmaa}{AMSTeX Label for \articlename:#1}}\fi}%
  }{}
  \@begofarticlehook
}%
%    \end{macrocode}
% Citations are more complicated since they can contain a list of
% keys. These are really only here for use with the rmaacite package,
% which is utterly deprecated (02 Jan 2003) now that I have the natbib
% code working. The command \cs{cite} should \emph{not} be used with
% natbib: use \cs{citep}, \cs{citet}, etc.\@ instead. 
%    \begin{macrocode}
\renewcommand{\nocite}[1]{\@bsphack
  \@for\@citeb:=#1\do{%
    \edef\@citeb{\expandafter\@firstofone\@citeb:\articlename}%
    \if@filesw\immediate\write\@auxout{\string\citation{\@citeb}}\fi
    \@ifundefined{b@\@citeb:\articlename}{\G@refundefinedtrue
      \@latex@warning{Citation `\@citeb' undefined}}{}}%
  \@esphack}
\renewcommand{\cite}[1]{\oldcite{#1:\articlename}}
%   \end{macrocode}


% Horrible fiddling with the bibliography citation mechanism. 
% This has to be globally available in case rmaacite puts it in the
% \texttt{.aux} file. 
%    \begin{macrocode}
\def\lbibcite#1#2#3{\global\@namedef{b@#1}{#2}
                    \global\@namedef{by@#1}{#3}}
%    \end{macrocode}

% \begin{macro}{\oldbibitem}
% \begin{macro}{\@rmaaBibArg}
% \begin{macro}{\@rmaaSpace}
% \begin{macro}{rmaabib}
% Of course, |\bibitem| needs to get the uniquification treatment as
% well. While we are at it, we redefine it to use consecutive integers
% for the key if the user leaves it blank (using the counter
% |rmaabib|). This allows one to put 
% |\bibitem{}| for all your references without getting annoying
% messages about multiply-defined labels.
% 
% We also take the opportunity to insert the optional argument to
% \cs{bibitem} if it is missing in the form Author(DATE) so that natbib
% doesn't try to do numerical citations. Natbib already provides a
% hook \cs{@extra@b@citeb} to help uniquify the citation key, so we
% just plug into that. However, it puts it \emph{after} the original
% key, not before as I had had it. Easiest solution is to change the
% way that I do it. 02 Jan 2003. Note that the \cs{@extra@b@citeb}
% does not get written to the .aux  file in the \cs{citation} command
% but it does work. 
% \changes{3.14}{06 Feb 2003}{Reinstalled old definition of \cs{bibitem}
%  for the case where the nonatbib option is requested, so that the
%  rmaacite package will work. This means that you can't use rmaacite
%  and natbib in the same volume.}
%    \begin{macrocode}
\if@natbib\def\@extra@b@citeb{:\articlename}\fi
%%\else
\let\oldbibitem\bibitem
\newcommand{\@rmaaBibArg}{}
\newcommand{\@rmaaSpace}{ }
\newcounter{rmaabib}
\if@natbib
\renewcommand{\bibitem}[2][ABC(2003)]{%
  \renewcommand{\@rmaaBibArg}{#2}
  \ifx\@rmaaBibArg\@rmaaEmpty
     \addtocounter{rmaabib}{1}
     \oldbibitem[#1]{\thermaabib:\articlename}
  \else\ifx\@rmaaBibArg\@rmaaSpace
     \addtocounter{rmaabib}{1}
     \oldbibitem[#1]{\thermaabib:\articlename}
  \else 
     \oldbibitem[#1]{#2:\articlename}
  \fi\fi
  \renewcommand{\@rmaaBibArg}{}
  }
\else
\renewcommand{\bibitem}[1]{%
  \renewcommand{\@rmaaBibArg}{#1}
  \ifx\@rmaaBibArg\@rmaaEmpty
     \addtocounter{rmaabib}{1}
     \oldbibitem{\articlename:\thermaabib}
  \else\ifx\@rmaaBibArg\@rmaaSpace
     \addtocounter{rmaabib}{1}
     \oldbibitem{\articlename:\thermaabib}
  \else
     \oldbibitem{\articlename:#1}
  \fi\fi
  \renewcommand{\@rmaaBibArg}{}
  }
\fi
%%\fi
%    \end{macrocode}
% \end{macro}
% \end{macro}
% \end{macro}
% \end{macro}

% \begin{macro}{oldpcite}
% \begin{macro}{oldscite}
% Then, we need to repeat the above scheme, but for the |rmaacite|
% definition. 
%    \begin{macrocode}
\newcommand{\@redefinebibitem}{%
  \let\oldbibitem\bibitem
  \let\oldcite\cite
  \let\oldpcite\pcite
  \let\oldscite\scite
  \def\bibitem[##1<##2>]##3{%
    \renewcommand{\@rmaaBibArg}{##3}
    \ifx\@rmaaBibArg\@rmaaEmpty
       \addtocounter{rmaabib}{1}
       \oldbibitem[##1<##2>]{\thermaabib:\articlename}
    \else\ifx\@rmaaBibArg\@rmaaSpace
       \addtocounter{rmaabib}{1}
       \oldbibitem[##1<##2>]{\thermaabib:\articlename}
    \else 
       \oldbibitem[##1<##2>]{##3:\articlename}
    % trust the author here
%%       \oldbibitem[##1<##2>]{##3}
    \fi\fi
    }
  }
%    \end{macrocode}
% \end{macro}
% \end{macro}


% \subsection{HTML Table of Contents}
% \label{sec:hyper-TOC}
% \begin{macro}{\HyperTOCEntry}
% \begin{macro}{\HyperTOCpart}
% \begin{macro}{\HyperTOCmeeting}
% \begin{macro}{\HyperTOCeditors}
% \begin{macro}{\HyperTOCpreamble}
% These commands are used to write information to the \texttt{.htoc}
% file, which is used in generating a Web Page of the Table of
% Contents. Since they only ever appear preceded by |\protect| in the
% argument to an |\addtocontents| command,  they should never need to
% be expanded by the macros. However, we define them here just in case
% \dots
%    \begin{macrocode}
\newcommand{\HyperTOCEntry}[3]{\relax}
\newcommand{\HyperTOCpart}[1]{\relax}
\newcommand{\HyperTOCmeeting}[3]{\relax}
\newcommand{\HyperTOCeditors}[1]{\relax}
\newcommand{\HyperTOCpreamble}{\relax}
%    \end{macrocode}
% \end{macro}
% \end{macro}
% \end{macro}
% \end{macro}
% \end{macro}
%

% \subsection{Formatting the poster abstracts} 
% These can now be used without using the book option.
%    \begin{macrocode}

\newcommand{\@makeminititle}{%
  \newpage
  \null 
  \vskip 3em%
  \begin{center}%
    {\@title \par}%
  \end{center}
  \vskip 3em%
  }

\newif\if@abstractsads
\@abstractsadsfalse
\newcommand\abstractsADS{\@abstractsadstrue}

\if@book\RequirePackage{multicol}\fi
\newenvironment{abstracts}[1][Abstracts of Contributed Papers]{%
  %% First, make sure we start on an odd page (06 Jun 2000)
  \if@book\clearpage\thispagestyle{empty}\cleardoublepage\fi
  %%
  \renewcommand{\articlename}{ABSTRACTS}
  \title{#1}\shortauthor{Abstracts}\shorttitle{Abstracts}
  \listofauthors{}
  \if@book
  \addcontentsline{toc}{chapter}{\protect\@fiddleauthors{\@TOCtitle}{\rmaaTOCauthors}}%
  \addtocontents{htc}{\protect\HyperTOCEntry{\@rawtitle}{\@rawauthors}{\articlename}}
  \fi
  \pagestyle{rmaa} % this now seems necessary (06 Jun 2000)
  \thispagestyle{rmaafirstpage}
%%  \twocolumn[\@makeminititle]
  \renewcommand{\abstract}[1]{%
    \renewcommand\ADS@abstract{##1}%
    \renewcommand{\@abstract}{##1}\@writeabstract%
  }
  \renewcommand{\title}[2][]{%
    \gdef\ADS@authors{}%
    \gdef\ADS@affil{}%
    \gdef\ADS@references{}%
    \gdef\ADS@keywords{}%
    \renewcommand{\@rmaaArg}{##1}
    \ifx\@rmaaArg\@rmaaEmpty
       \renewcommand{\@TOCaftertitle}{}
    \else
       \renewcommand{\@TOCaftertitle}{ (\@rmaaArg)}
    \fi
    \renewcommand{\@title}{\rmaa@TOCTitleStyle{##2}}%
    \renewcommand{\@TOCtitle}{\rmaa@TOCTitleStyle{##2}\@TOCaftertitle}%
    \renewcommand{\rmaaIndexTitle}{\rmaa@IndexTitleStyle{##2\@TOCaftertitle.}}%
    \renewcommand\ADS@title{##2}%
    \xdef\ADS@firstpage{\@arabic{\c@page}}%
    }
  \renewcommand{\author}[1]{\let\@thanks\@empty\renewcommand\@author{##1} }
  \renewcommand{\footnotetext}[2][XX]{\par\noindent\textsuperscript{##1}~##2}
  \renewcommand{\footnoterule}{\kern-3\p@ \hrule \@width 3em \kern 2.6\p@}
  \if@book\@begofarticlehook\fi
  \if@book\begin{multicols}{2}[\@makeminititle]\else\@makeminititle\fi
  \begingroup
  \setcounter{footnote}{0}
  }% and at the end of the environment....
  {\endgroup\if@book\end{multicols}\listofabstracts\label{RMAALastPage}\clearpage\fi}
%% auxiliary commands used:
%% \@writeabstract does all the actual formatting 
\newcommand{\@writeabstract}{%
  \if@book%
  \addcontentsline{loa}{abstract}{\protect\@fiddleauthors{\@TOCtitle}{\rmaaTOCauthors}}%
  \fi%
  \begin{center}\@title\par\@author\end{center}
  \@abstract
  \par\noindent\footnoterule\nobreak\footnotesize\@thanks
  \par\noindent\hrulefill\goodbreak
  \vskip 1cm \@plus0.5cm
  %% write ADS entry
  \if@book\if@abstractsads\ADS@dumpdata\fi\fi%
  \endgroup % this pair is to make each abstract a self-contained group
  \begingroup
  \setcounter{footnote}{0}
  }

\if@book
%% List of abstracts (stuff written to .loa file)
\newcommand{\listofabstracts}{%
  \begin{multicols}{2}[%
    \vskip 1.5cm \vfill % leave at least 1.5 cm and go to bottom of page if possible
    \centerline{\MakeUpperCase{List of Abstracts}}
    \vskip 1.5cm \@plus0.5cm
    ][6cm] % this makes it start a new page if there is not at least 6cm of space
    \renewcommand{\@fiddleauthors}[2]{##1\hskip1em plus2em minus0.4em\penalty-100  ##2}
    \parindent\z@
%%    \parskip.5\p@ \@plus .7\p@\relax 
    \parskip3\p@ \@plus 1\p@ \@minus 1\p@\relax % allow quite a lot of
                                % leeway here to avoid breaks inside
                                % entries 
    \renewcommand{\\}{\relax} % temporarily turn off explicit line breaks just in case
    \renewcommand{\altaffilmark}[1]{\relax} % ditto with this
    \@starttoc{loa}
  \end{multicols}
  }
\newcommand*{\l@abstract}[2]{%
  \par\penalty-1000
  #1\hskip 1em\penalty0\null\nobreak\hfill#2}
\fi

%    \end{macrocode}
 
% \subsection{Code for typesetting an entire volume of conference
% proceedings} \label{sec:whole-volume}
%    \begin{macrocode}

\if@book
%% environment for including the individual articles
\newcommand{\@fiddleauthors}[2]{\relax}
\renewcommand{\theequation}{\@arabic\c@equation}
\newenvironment{articles}%
{ \begingroup % make it a group so that author definitions do not
              % spill over
  \gdef\ADS@authors{}
  \pagestyle{rmaa}
  \renewenvironment{document}{%
    \@begofarticlehook
    \addtocounter{chapter}{1}%
    \addcontentsline{toc}{chapter}{%
      \protect\@fiddleauthors{\@TOCtitle}{\rmaaTOCauthors}}%
    \addtocontents{htc}{\protect\HyperTOCEntry{\@rawtitle}{\@rawauthors}{\articlename}}
    }{\relax}
  \renewcommand{\documentclass}[2][foo]{\relax} 
  \renewcommand{\usepackage}[2][foo]{%
    \def\@packagename{##2}\def\@rmaacitename{rmaacite}%
    \ifx\@packagename\@rmaacitename
    \@rmcitetrue\ClassInfo{rmaa}{This paper uses rmaacite}
    \makeatletter\input{rmaacite.sty}\makeatother % read it in here so that all changes
    \else                                         % are local to paper
    \if@debug\ClassWarning{rmaa}{Package \@packagename\space was requested by \articlename\space}\fi
    \fi} 
  %% Don't allow individual articles to mess with the pages, volume or year
  \renewcommand{\SetFirstPage}[1]{\relax}
  \renewcommand{\SetYear}[1]{\relax}
  \renewcommand{\SetVolume}[1]{\relax}
  \renewcommand{\SetConfTitle}[1]{\relax}
  %% Don't allow individual articles to change these
  \renewcommand{\nonstopmode}{\relax}
  \renewcommand{\scrollmode}{\relax}
  \renewcommand{\outputplates}{\relax}
  %% XXXX need to check if using rmaacite and act accordingly
  %% avoid multiple declarations of the same length
  \let\oldnewlength=\newlength
  \renewcommand{\newlength}[1]{%
    \edef\reserved@a{\expandafter\@gobble\string ##1}%
    \@ifundefined{\reserved@a}{%
      \oldnewlength{##1}}{%
      \if@debug{\ClassInfo{rmaa}{\reserved@a\space already defined: doing nothing}}\fi%
      \relax}}%
  \renewcommand\newcommand{\providecommand}
  }%
{\endgroup\onecolumn}
     

%    \end{macrocode}
% Command to be used in `articles' environment for including an
% article
%    \begin{macrocode}
\newcommand{\article}[2][]{%
  % optional first argument is extra text to put in brackets after title in TOC
  % e.g. ``Invited Talk'' or ``Poster''
  \renewcommand{\articlename}{#2}
  % reset loads of counters
  \setcounter{footnote}{0}
  \setcounter{figure}{0}
  \setcounter{table}{0}
  \setcounter{section}{0}
  \setcounter{equation}{0}
  \setcounter{rmaabib}{0}
  \setcounter{author}{0}
  \setcounter{AUTHOR}{1}
  \gdef\ADS@authors{}%
  \gdef\ADS@affil{}%
  \gdef\ADS@affiliations{}%
  \gdef\ADS@references{}%
  \@firstauthortrue
  \@firstaffiltrue
  \@firstaltaffiltrue
  \@firstkwdtrue
  \@rmcitefalse
  \if@suppressfulladdresses\else\@needaddressestrue\fi
  \renewcommand{\@rmaaArg}{#1}
  \ifx\@rmaaArg\@rmaaEmpty
     \renewcommand{\@TOCaftertitle}{}
  \else
     \renewcommand{\@TOCaftertitle}{ (#1)}
  \fi
%    \end{macrocode}
% In the main journal articles must start on an odd page. 
%    \begin{macrocode}
  \if@journal\clearpage\thispagestyle{empty}\cleardoublepage\fi
%% Now input the file, then run the end-of-article hook (reinstated 28
  %% Sep 2001)
  {\input{\@articledir/\articlename.tex}\@endofarticlehook}}
%    \end{macrocode}
% Command to separate articles into sessions. Write an entry in the
% TOC, makes sure we start on an odd page (added 06 Jun 2000), then do
% a title page if \cs{SessionTitlePages} was called. 
% \changes{3.12}{04 Feb 2003}{Added option to have title page for each
% session. Turn it on with \cs{SessionTitlePages}.}
%    \begin{macrocode}
\newcommand{\session}[1]{
  %% this does nothing except write an entry in the TOC
  \addcontentsline{toc}{part}{\rmaa@TOCPartStyle{#1}}
  \addtocontents{htc}{\protect\HyperTOCpart{#1}}
  %% and also makes sure we start on an odd page (added 06 Jun 2000)
  \clearpage\thispagestyle{empty}\cleardoublepage
  \stepcounter{part}
  \if@sessiontitlepages{%
    \onecolumn
    \thispagestyle{empty}
    \renewcommand{\articlename}{Section\thepart TitlePage}
    \label{RMAAFirstPage}
    \make@session@title{#1}
    \label{RMAALastPage}
    \clearpage\thispagestyle{empty}\cleardoublepage
  }\fi
  }
\newcommand{\make@session@title}[1]{%
  {%
    \setlength\parindent{0pt}
    \vspace*{\stretch{1}}
    \rule{\linewidth}{1mm}
    \begin{flushright}\Huge\bfseries\boldmath
      #1
    \end{flushright}
    \rule{\linewidth}{1mm}
    \vspace*{\stretch{2}}
  }
}
%    \end{macrocode}
% Command to turn on separate title pages for each session. 
%    \begin{macrocode}
\newcommand{\SessionTitlePages}{\@sessiontitlepagestrue}
\newcommand{\BoldSessionTitles}{%
  \renewcommand{\rmaa@TOCPartStyle}[1]{\textbf{\rmaa@TOCTitleStyle{##1}}}
  }
\newcommand{\ArticleDir}[1]{\renewcommand{\@articledir}{#1/}\graphicspath{{1/}}}
\newcommand{\@articledir}{.}
%\newcommand{\ps@rmaafront}{\ps@rmaa}% page style for front matter (same as ever, actually)
\newcommand{\ps@rmaafront}{% page style for front matter:  page # in each foot
  \renewcommand\@oddfoot{\hfil{\thepage}}%
  \renewcommand\@evenfoot{{\thepage}\hfil}%
  % shorttitle in both heads
  \renewcommand\@evenhead{\hfil\@shorttitle\hfil}%
  \renewcommand\@oddhead{\@evenhead}
  }
\newcommand{\ps@rmaafrontfirstpage}{% page style for first page of front matter
  % page # in each foot, on outside edge
  \renewcommand\@oddfoot{\hfil\thepage}%
  \renewcommand\@evenfoot{\thepage\hfil}%
  % nothing in the even heads
  \renewcommand\@evenhead{\relax}%
  \renewcommand\@oddhead{\relax}%
  }
%% slight modification to the standard \frontmatter and \mainmatter
%% commands, just to set the \pagestyle
\renewcommand\frontmatter{%
  \latexonecolumn
  \cleardoublepage
  \@frontmattertrue
  \@mainmatterfalse
  \pagenumbering{roman}
  \pagestyle{rmaafront}
  \addtocounter{page}{2} % leave space for the title page
  }
\renewcommand\mainmatter{%
  \pagestyle{empty}
  \if@rmtwocol\twocolumn\fi
  \cleardoublepage
  \@frontmatterfalse
  \@mainmattertrue
  \pagenumbering{arabic}
  \pagestyle{rmaa}
  }
\renewcommand\backmatter{%
  \clearpage\pagestyle{empty}\cleardoublepage
  \@mainmatterfalse
  %% put extra space in the TOC
  \addtocontents{toc}{\bigskip}
  }
%    \end{macrocode}
%
% \subsection{Front sections}\label{sec:front-sections}
% \begin{macro}{\notitle}
% \begin{macro}{\shortname}
% \begin{macro}{frontsection}
% The environment \texttt{frontsection} is for doing a section in the
% front matter (preface, list of participants, or whatever). It has 
% one optional argument:
% \begin{enumerate}
% \item[\#1] Short title (no embedded spaces) for
% use in making PS files. E.g., ``list\_of\_participants''.   
% \end{enumerate}
% and two mandatory arguments: 
% \begin{enumerate}
% \item[\#2] Title of section (e.g., `Preface')
% \item[\#3] Commands to be evaluated before processing the contents
% of the environment. This is mainly useful for putting in the things
% that would go in the preamble of a normal article  (e.g.,
% |\listofauthors|, |\indexauthor|, |\notitle|) 
% \end{enumerate}
%    \begin{macrocode}
\newcommand{\notitle}{\@titlevisiblefalse} \newcommand{\shortname}{}
\newenvironment{frontsection}[3][]{%
  \begingroup
  \renewcommand{\shortname}{#1}
  \ifx\shortname\@rmaaEmpty
     \renewcommand{\articlename}{#2}
  \else
     \renewcommand{\articlename}{\shortname}
  \fi
  \@begofarticlehook
  \@titlevisibletrue % default
  \thispagestyle{rmaafrontfirstpage}
  \listofauthors{} % blank out this in case it is not specified
  \title{{#2}}\shorttitle{#2}\shortauthor{#2}
  #3 % evaluate the command given as second mandatory arg
  \addcontentsline{toc}{chapter}{\protect\@fiddleauthors{\@TOCtitle}{\rmaaTOCauthors}}%
  \addtocontents{htc}{\protect\HyperTOCEntry{\@rawtitle}{\@rawauthors}{\articlename}}
  \if@titlevisible
     \vspace*{1cm}\centerline{\rmaa@FrontSecStyle{\@title}}\vspace*{1cm}
  \fi
  \setlength{\parskip}{3ex}
  }{\frontlabel{RMAALastPage}\clearpage\endgroup}
%    \end{macrocode}
% \end{macro}
% \end{macro}
% \end{macro}
%    \begin{macrocode}
%% Definition of the command \frontlabel, which rights out something like `F4' to
%% the .aux file, instead of `\textlatin  {i}'. This is so it is easy to post-process
%% the file for making PS files of individual sections
\newcommand{\frontlabel}[1]{%
  \@bsphack
  \protected@write\@auxout{}%
     {\string\newlabel{\articlename:#1}{{}{F\@arabic{\c@page}}}}
  \@esphack
  }
%% User command to say X, Y, Z,\\ Editors of the Proceedings
\newcommand{\EditorSignature}{\vspace*{\parskip}\rightline{\@editors}%
  \rightline{Editors of the proceedings}}
%% Fiddle with \thesection, etc. so that they don't include the chapter number
%% Basically go back to the `article' definitions of these
\renewcommand{\thesection}{\@arabic\c@section}
\renewcommand{\thefigure}{\@arabic\c@figure}
\renewcommand{\thetable}{\@arabic\c@table}

\newcommand{\@oldargs}{}
\newcommand{\@newargs}{}
%% add TOC entry for each paper
%\renewcommand{\title}[1]%
%{\renewcommand{\@title}{\rmaa@TitleStyle{#1}}%
%\addtocounter{chapter}{1}\addcontentsline{toc}{chapter}{#1}%
%}

%% Formatting of the table of contents
%% dimensions of various parts
\newlength{\TOCtextwidth}\setlength{\TOCtextwidth}{12.65cm}
\newlength{\TOCpagenumwidth}\setlength{\TOCpagenumwidth}{1.3cm}
\newlength{\TOCentrywidth}\setlength{\TOCentrywidth}{\TOCtextwidth}
\addtolength{\TOCentrywidth}{-\TOCpagenumwidth} 
\newlength{\TOCextramargin}\setlength{\TOCextramargin}{1.25cm}
\newcommand{\TOCbox}[1]{\hspace{\TOCextramargin}\parbox{\TOCtextwidth}{#1}\par}
\newlength{\TOCextramarginx}\setlength{\TOCextramarginx}{\TOCextramargin}\addtolength{\TOCextramarginx}{4.3cm}
\newlength{\TOCtextwidthx}\setlength{\TOCtextwidthx}{\TOCtextwidth}\addtolength{\TOCtextwidthx}{-4.3cm}
\newcommand{\TOCboxx}[1]{\hspace{\TOCextramarginx}\parbox{\TOCtextwidthx}{#1}\par}

\newcommand{\RMAAbannerimage}{rmsc}
\renewcommand\tableofcontents{%
  \renewcommand{\articlename}{CONTENTS}
  \begingroup
  \@begofarticlehook
  \if@twocolumn
     \@restonecoltrue\onecolumn
  \else
     \@restonecolfalse
  \fi
  %% 9pt for journal style
  \if@proceedings\else\small\fi
  \if@newstyle\small\fi %% same for new style
  \if@proceedings\thispagestyle{rmaafrontfirstpage}\else\thispagestyle{empty}\fi
  \shortauthor{Contents}\shorttitle{Contents}
  \vspace*{18mm}
  \if@printbanner
  %% put in the banner header (need to get better image of this)
     \TOCbox{\hspace*{-0.15cm}\includegraphics[width=0.9\TOCtextwidth]{\RMAAbannerimage}}
     \vspace*{8mm}
  %% volume and publication date (should really be in funny font)
     \TOCboxx{\@volumen\\[2mm] \@mes\ \@anyo}
  \else
  %% put in the banner header (need to get better image of this)
     \TOCbox{\hspace*{-0.15cm}\includegraphics[width=0.9\TOCtextwidth]{\RMAAbannerimage}}
     \vspace*{12.5mm}
  %% volume and publication date (should really be in funny font)
     \TOCboxx{~\\[2mm] ~}
  \fi
  \vspace*{16mm}
  %% This is the command that fiddles with the TOC entry so that the
  %% title is flushleft but the authors are flushright
  \renewcommand{\@fiddleauthors}[2]{%
    % first put the paper title
    \renewcommand{\\}{\relax} % temporarily turn off explicit line breaks just in case
    \renewcommand{\altaffilmark}[1]{\relax} % ditto with this
    \renewcommand{\textsuperscript}[1]{\relax} % ditto with this 
    ##1
    % now save the author list in a box
    \sbox\@tempboxa{##2}%
    % check that box is shorter than the total width available for the entry
    \ifdim \wd\@tempboxa <\TOCentrywidth
       % if so, then put it right here, flush against right margin
       \penalty10000\hskip1em plus10em\penalty0\hbox{}\nobreak\hfill\box\@tempboxa
       % Note that we force a space of at least 1em between title and authorlist, 
       % stretchable to 11em since we may have very short titles but still with an 
       % authorlist too long to go on the same line.  The only linebreak allowed is 
       % directly before the authorlist.
    \else
       % if it is too long, we set the original authorlist text in its own parbox
       \par\vskip1ex\parbox[b]{\TOCentrywidth}{\raggedleft ##2}
    \fi
    }
  \@starttoc{toc}%
  \if@restonecol\twocolumn\fi
  \frontlabel{RMAALastPage}
  \clearpage
  \thispagestyle{empty}
  \cleardoublepage % make sure we start on an odd page
  \endgroup

  %%% ALSO, do the HTML TOC
  %% The file that will be used to make the HTML TOC
  \@starttoc{htc}
  %% All the stuff in this file does nothing in the book style
  \addtocontents{htc}{\protect\HyperTOCmeeting{\@conftitle}{\@reunionplace}{\@reuniondate}}
  \addtocontents{htc}{\protect\HyperTOCeditors{\@editors}}
  \addtocontents{htc}{\protect\HyperTOCpreamble}
  \addtocontents{htc}{\protect\HyperTOCEntry{Table of Contents}{}{\articlename}}

  %% and the ADS file
  \ADS@openfile
  }
%% This is now much simpler
\renewcommand\l@chapter[2]{%TOC entry for articles
  \ifnum \c@tocdepth >\m@ne
    \vskip 1.0em \@plus\p@
    \TOCbox{%
      \parbox[b]{\TOCentrywidth}{#1}\parbox[b]{\TOCpagenumwidth}{\hfill#2}
    }
  \fi} 
\renewcommand\l@part[2]{%TOC entry for sessions: centered, no page number
  \ifnum \c@tocdepth >\m@ne
    \vskip 1.5em \@plus1.5em \goodbreak
    \TOCbox{%
      \parbox[b]{\TOCentrywidth}{\hfil#1\hfil}\hfill
    }
  \fi\nobreak} 

%% Formatting theindex
%% strongly encourage page breaks between index entries
\renewcommand{\@idxitem}{\par\penalty-1000\hangindent.45cm} 
\renewcommand{\indexspace}{\par\penalty-1000\vskip 10\p@ \@plus15\p@ \@minus3\p@\relax}
\renewenvironment{theindex}{%
  %% first, make sure we start on an odd page (06 Jun 2000)
  \clearpage\thispagestyle{empty}\cleardoublepage
  %% this looks a bit like an article but is two columns and without abstract, etc
  \renewcommand{\articlename}{INDEX}
  \title{Author Index}\shortauthor{Author Index}\shorttitle{Author Index}
  \listofauthors{}
  \addcontentsline{toc}{chapter}{\protect\@fiddleauthors{\@TOCtitle}{\rmaaTOCauthors}}%
  \addtocontents{htc}{\protect\HyperTOCEntry{\@rawtitle}{\@rawauthors}{\articlename}}
  \pagestyle{rmaa}
  \thispagestyle{rmaafirstpage}
  \@begofarticlehook
  %% 9pt for journal style
  \if@proceedings\else\small\fi
  \if@newstyle\small\fi
  \twocolumn[\@makeminititle]
  %% vertical spacing (see also definition of \@idxitem and \indexspace
  \parskip\z@ \@plus .6\p@\relax % allow quite a lot of leaway here to avoid breaks inside entries
  %% horizontal spacing, breaking and glue parameters (see also the stuff in the definition 
  %% of \indexauthor
  \parindent\z@ 
  \exhyphenpenalty10000 % don't break on explicit hyphens, it always looks bad
  \rightskip=0pt plus.2em % allow a tiny bit of leeway at the ends of lines
  \spaceskip.3333em plus.5em minus.15em % a lot of leeway between words
  \xspaceskip.5em plus.5em minus.4em % this seems to have no effect
  \let\item\@idxitem
  \renewcommand{\\}{\relax} % temporarily turn off explicit line breaks just in case
  \renewcommand{\altaffilmark}[1]{\relax} % ditto with this
  \renewcommand{\textsuperscript}[1]{\relax} % ditto with this
  }{\label{RMAALastPage}\clearpage} % clearpage necessary so that we know about \shortauthor



%% end definitions for book styles
\fi

%    \end{macrocode}
% In order to do away with plates, we sometimes want a figure to
% appear in color only in the online version. The following command
% can bve used in a figure caption to typeset a note to this effect. 
%    \begin{macrocode}
\newcommand\ColorOnlineOnly{%
  \textsc{Note: this figure is available in color in the electronic
    version of this article, obtainable from}
  \texttt{http://\discretionary{}{}{}www.\discretionary{}{}{}astroscu.\discretionary{}{}{}unam.\discretionary{}{}{}mx/\discretionary{}{}{}\~{}rmaa/}.}
%% Plates
%% Define 2 counters: plate is the number in the whole document, plateloop is used 
%% internally by \outputplates
\newcounter{plate}
\newcounter{plateloop}
%% Currently limited to 15 plates in document
%% when we do a label, we need to write to \newlabel entries to the .aux file, one 
%% that says RMAAPlateX so that it can be referred to in \ouputplates, and one that
%% says articlename:RMAAPlateX so that the dvips program can find it and asociate it
%% with the right article
\newcommand{\@platelabeli}{\oldlabel{RMAAPlate1}\label{RMAAPlate1}}
\newcommand{\@platelabelii}{\oldlabel{RMAAPlate2}\label{RMAAPlate2}}
\newcommand{\@platelabeliii}{\oldlabel{RMAAPlate3}\label{RMAAPlate3}}
\newcommand{\@platelabeliv}{\oldlabel{RMAAPlate4}\label{RMAAPlate4}}
\newcommand{\@platelabelv}{\oldlabel{RMAAPlate5}\label{RMAAPlate5}}
\newcommand{\@platelabelvi}{\oldlabel{RMAAPlate6}\label{RMAAPlate6}}
\newcommand{\@platelabelvii}{\oldlabel{RMAAPlate7}\label{RMAAPlate7}}
\newcommand{\@platelabelviii}{\oldlabel{RMAAPlate8}\label{RMAAPlate8}}
\newcommand{\@platelabelix}{\oldlabel{RMAAPlate9}\label{RMAAPlate9}}
\newcommand{\@platelabelx}{\oldlabel{RMAAPlate10}\label{RMAAPlate10}}
\newcommand{\@platelabelxi}{\oldlabel{RMAAPlate11}\label{RMAAPlate11}}
\newcommand{\@platelabelxii}{\oldlabel{RMAAPlate12}\label{RMAAPlate12}}
\newcommand{\@platelabelxiii}{\oldlabel{RMAAPlate13}\label{RMAAPlate13}}
\newcommand{\@platelabelxiv}{\oldlabel{RMAAPlate14}\label{RMAAPlate14}}
\newcommand{\@platelabelxv}{\oldlabel{RMAAPlate15}\label{RMAAPlate15}}
\newcommand{\@outputplatelabeli}{\oldlabel{RMAAOutputPlate1}}
\newcommand{\@outputplatelabelii}{\oldlabel{RMAAOutputPlate2}}
\newcommand{\@outputplatelabeliii}{\oldlabel{RMAAOutputPlate3}}
\newcommand{\@outputplatelabeliv}{\oldlabel{RMAAOutputPlate4}}
\newcommand{\@outputplatelabelv}{\oldlabel{RMAAOutputPlate5}}
\newcommand{\@outputplatelabelvi}{\oldlabel{RMAAOutputPlate6}}
\newcommand{\@outputplatelabelvii}{\oldlabel{RMAAOutputPlate7}}
\newcommand{\@outputplatelabelviii}{\oldlabel{RMAAOutputPlate8}}
\newcommand{\@outputplatelabelix}{\oldlabel{RMAAOutputPlate9}}
\newcommand{\@outputplatelabelx}{\oldlabel{RMAAOutputPlate10}}
\newcommand{\@outputplatelabelxi}{\oldlabel{RMAAOutputPlate11}}
\newcommand{\@outputplatelabelxii}{\oldlabel{RMAAOutputPlate12}}
\newcommand{\@outputplatelabelxiii}{\oldlabel{RMAAOutputPlate13}}
\newcommand{\@outputplatelabelxiv}{\oldlabel{RMAAOutputPlate14}}
\newcommand{\@outputplatelabelxv}{\oldlabel{RMAAOutputPlate15}}
\newcommand{\@platerefi}{\oldpageref{RMAAPlate1}}
\newcommand{\@platerefii}{\oldpageref{RMAAPlate2}}
\newcommand{\@platerefiii}{\oldpageref{RMAAPlate3}}
\newcommand{\@platerefiv}{\oldpageref{RMAAPlate4}}
\newcommand{\@platerefv}{\oldpageref{RMAAPlate5}}
\newcommand{\@platerefvi}{\oldpageref{RMAAPlate6}}
\newcommand{\@platerefvii}{\oldpageref{RMAAPlate7}}
\newcommand{\@platerefviii}{\oldpageref{RMAAPlate8}}
\newcommand{\@platerefix}{\oldpageref{RMAAPlate9}}
\newcommand{\@platerefx}{\oldpageref{RMAAPlate10}}
\newcommand{\@platerefxi}{\oldpageref{RMAAPlate11}}
\newcommand{\@platerefxii}{\oldpageref{RMAAPlate12}}
\newcommand{\@platerefxiii}{\oldpageref{RMAAPlate13}}
\newcommand{\@platerefxiv}{\oldpageref{RMAAPlate14}}
\newcommand{\@platerefxv}{\oldpageref{RMAAPlate15}}
\newcounter{figplatei}
\newcounter{figplateii}
\newcounter{figplateiii}
\newcounter{figplateiv}
\newcounter{figplatev}
\newcounter{figplatevi}
\newcounter{figplatevii}
\newcounter{figplateviii}
\newcounter{figplateix}
\newcounter{figplatex}
\newcounter{figplatexi}
\newcounter{figplatexii}
\newcounter{figplatexiii}
\newcounter{figplatexiv}
\newcounter{figplatexv}
%% Printed output should use the counter `plate'
\renewcommand{\theplate}{\@arabic\c@plate}
\renewcommand{\theplateloop}{\@arabic\c@plateloop}
%% Command to output all the saved plates in this paper
\newcommand{\outputplates}{%
  \setcounter{plateloop}{0}
  \if@book\else\@endofarticlehook\fi % not entirely sure why, but this
                                     % seems to be needed
  \clearpage % first output the last page of text
  % now set pagestyle for blank even pages and make sure we start on an odd page
  \pagestyle{empty}
  \cleardoublepage
  \onecolumn
  \parindent\z@
  % first put in an extra sheet
  % the main purpose of this is so there is always a blank even page before 
  % each plate for filling in the PS file if the article has an odd number
  % of pages (see splitbook.sh)
  \null\vfil\centerline{PLATES}\vfil\clearpage\pagestyle{rmaaplate}\cleardoublepage
  % output each plate in turn
  \@whilenum\value{plateloop}<\value{plate}\do{%
    \addtocounter{plateloop}{1}
    % restore the figure counter from \@figcounti, \@figcountii, etc
    \setcounter{figure}{\value{figplate\roman{plateloop}}}
    % spew out the plate (in its own group for safety)
    {\csname @plate\roman{plateloop}\endcsname\par}
    }
  \pagestyle{rmaa}
  % reset the page number as if nothing had happened
  \addtocounter{page}{-\value{plate}}
  \addtocounter{page}{-\value{plate}}
  }
\newcommand\figcaption{%
  \def\@captype{figure}
  \renewcommand{\@makecaption}{\@makefigcap}
  \caption}
\newcommand\tabcaption{%
  \def\@captype{table}
  \renewcommand{\@makecaption}{\@maketabcap}
  \caption}

\newcommand{\plate}[2][1]{%
  %% User command to make a plate page: 
  %% Mandatory argument is what goes on the plate - this should not 
  %% include any floating environments (e.g. figure). Instead, \figcaption 
  %% should be used to do the captions. 
  %% Optional argument is the number of figures (i.e. captions) on the plate
  %% page (default is 1). This is necessary for incrementing the figure counter. 
  %%
  %% Note that we have to use a \gdef rather than an \xdef when we 
  %% save the plate output in a command. This means that in the book styles when it 
  %% comes to spewing out the plate, we will have forgotten about any user macros 
  %% that may have been defined. This can be got round by repeating the macros inside the plate.
  %% (Presumably, one would use a \newcommand rather than \renewcommand if they were new.)
  % increment global and local counters
  \addtocounter{plate}{1}
  % write appropriate label to .aux file
  \csname @platelabel\roman{plate}\endcsname
  % define a global macro called \@platei, \@plateii, \@plateiii, \@plateiv, ...
  \expandafter\gdef\csname @plate\roman{plate}\endcsname{%
    \csname @outputplatelabel\roman{plateloop}\endcsname
    \null\vfil #2 \vfil
    \cleardoublepage
    }
  % save the current shortauthor and shorttitle
  % need to be \xdef to take current values of \@shortauthor, etc 
  \expandafter\xdef\csname @shorta\roman{plate}\endcsname{\@shortauthor}
  \expandafter\xdef\csname @shortt\roman{plate}\endcsname{\@shorttitle}
  % save the current value of the figure counter for later use in \outputplates
  \setcounter{figplate\roman{plate}}{\value{figure}}
  % increment the figure counter according to the optional argument
  \addtocounter{figure}{#1}
  }
%% page style for plate pages 
\newcommand{\ps@rmaaplate}{%
  % Plate X (see page Y) in the foot
  \renewcommand\@oddfoot{\hfil\csname @shorta\roman{plateloop}\endcsname\ 
    (see page~\csname @plateref\roman{plateloop}\endcsname)\hfil}%
  \renewcommand\@evenfoot{\relax}%
  % nothing in even heads (blank pages)
  \renewcommand\@evenhead{\relax}
  % short title and page num in the odd heads
  \renewcommand\@oddhead{\hfil\csname @shortt\roman{plateloop}\endcsname%
    \hfil\llap{PLATE~\theplateloop}}
  }





%% Things for the User Guide
\if@guide
%% stolen from doc.sty
\def\MakeShortVerb#1{%
  \expandafter\ifx\csname cc\string#1\endcsname\relax
  \@shortvrbinfo{Made }{#1}%
  \add@special{#1}%
  \expandafter
  \xdef\csname cc\string#1\endcsname{\the\catcode`#1}%
  \begingroup
  \catcode`\~\active  \lccode`\~`#1%
  \lowercase{%
    \global\expandafter\let
    \csname ac\string#1\endcsname~%
    \gdef~{\verb~}}%
  \endgroup
  \global\catcode`#1\active
  \else
  \@shortvrbinfo\@empty{#1 already}%
  \fi}
\def\DeleteShortVerb#1{%
  \expandafter\ifx\csname cc\string#1\endcsname\relax
  \else
  \@shortvrbinfo{Deleted }{#1 as}%
  \rem@special{#1}%
  \global\catcode`#1\csname cc\string#1\endcsname
  \global \expandafter\let \csname cc\string#1\endcsname \relax
  \ifnum\catcode`#1=\active
  \begingroup
  \catcode`\~\active   \lccode`\~`#1%
  \lowercase{%
    \global\expandafter\let\expandafter~%
    \csname ac\string#1\endcsname}%
  \endgroup \fi \fi}
\def\@shortvrbinfo#1#2{%
  \PackageInfo{doc}{%
    #1\expandafter\@gobble\string#2 a short reference
    for \string\verb}}
\def\add@special#1{%
  \rem@special{#1}%
  \expandafter\gdef\expandafter\dospecials\expandafter
  {\dospecials \do #1}%
  \expandafter\gdef\expandafter\@sanitize\expandafter
  {\@sanitize \@makeother #1}}
\def\rem@special#1{%
  \def\do##1{%
    \ifnum`#1=`##1 \else \noexpand\do\noexpand##1\fi}%
  \xdef\dospecials{\dospecials}%
  \begingroup
  \def\@makeother##1{%
    \ifnum`#1=`##1 \else \noexpand\@makeother\noexpand##1\fi}%
  \xdef\@sanitize{\@sanitize}%
  \endgroup}
%% define | as shortcut for \verbatim
\MakeShortVerb{\|}
%% comand to specify arguments to LaTeX commands
\newcommand{\rmArg}[1]{%
  \ensuremath{\langle}{\scriptsize \textsf{\MakeUpperCase{#1}}}\ensuremath{\rangle}}
%\newcommand{\rmArg}[1]{\ensuremath{\langle}\textsc{#1}\ensuremath{\rangle}}
\providecommand{\BibTeX}{{\rmfamily B\kern-.05em%
    \textsc{i\kern-.025em b}\kern-.08em%
    T\kern-.1667em\lower.7ex\hbox{E}\kern-.125emX}}
%% Environment for code examples
\newenvironment{Code}{%
  \begin{list}{}{\setlength{\leftmargin}{0.05\columnwidth}} 
  \item[]}{\end{list}}
\fi
%</package>
%    \end{macrocode}

% \section{User commands for astronomical articles}
% \label{sec:user-commands}
% Some random definitions. These go both into the main class file and
% also into |adshtml.sty|.
% 
% Definitions of the journal name in different incarnations.  
%    \begin{macrocode}
%<*package|adshtml-style>
\newcommand{\RMAA}{RMxAA}
\newcommand{\RMAAlong}{Revista Mexicana de Astronom\'{\i}a y Astrof\'{\i}sica}
\newcommand{\RMAASC}{RMxAC}
\newcommand{\RMAASClong}{\RMAA\ (Serie de Conferencias)}
%    \end{macrocode}
% \begin{macro}{\ibidrule}
% This puts a horizontal rule with a period after it. It is used when
% the author list of a reference is identical to the preceding one. 
%    \begin{macrocode}
\newcommand{\ibidrule}{\rule[-0.03cm]{4.7em}{.01cm}.\ }
%    \end{macrocode}
% \end{macro}
% \begin{macro}{\ThisVolume}
% References to other articles in the same volume can be coded as,
% e.g., |\bibitem{} Bloggs, J. \ThisVolume{bloggs}|.
%    \begin{macrocode}
\newcommand{\ThisVolume}[1]{%
  \@year, \if@proceedings\RMAASC\else\RMAA\fi, \@volume,
  \if@book\oldpageref{#1:RMAAFirstPage}\else 000\fi\ (this volume)}
%    \end{macrocode}
% \end{macro}
% Other stuff that may be useful in ref lists.
%    \begin{macrocode}
\newcommand{\etal}{et al.\@}
%    \end{macrocode}
% Various new commands from AAS macros
%    \begin{macrocode}
\def\acknowledgments{\vskip 3ex plus .8ex minus .4ex}
\let\acknowledgements=\acknowledgments                  % second spelling
\let\jnl@style=\rm
\def\ref@jnl#1{{\jnl@style#1}}
\DeclareRobustCommand\aj{\ref@jnl{AJ}}                   
\DeclareRobustCommand\araa{\ref@jnl{ARA\&A}}             
\DeclareRobustCommand\apj{\ref@jnl{ApJ}}                 
\DeclareRobustCommand\apjl{\ref@jnl{ApJ}}                
\DeclareRobustCommand\apjs{\ref@jnl{ApJS}}               
\DeclareRobustCommand\ao{\ref@jnl{Appl.~Opt.}}           
\DeclareRobustCommand\apss{\ref@jnl{Ap\&SS}}             
\DeclareRobustCommand\aap{\ref@jnl{A\&A}}                
\DeclareRobustCommand\aapr{\ref@jnl{A\&A~Rev.}}          
\DeclareRobustCommand\aaps{\ref@jnl{A\&AS}}              
\DeclareRobustCommand\azh{\ref@jnl{AZh}}                 
\DeclareRobustCommand\baas{\ref@jnl{BAAS}}               
\DeclareRobustCommand\jrasc{\ref@jnl{JRASC}}             
\DeclareRobustCommand\memras{\ref@jnl{MmRAS}}            
\DeclareRobustCommand\mnras{\ref@jnl{MNRAS}}             
\DeclareRobustCommand\pra{\ref@jnl{Phys.~Rev.~A}}        
\DeclareRobustCommand\prb{\ref@jnl{Phys.~Rev.~B}}        
\DeclareRobustCommand\prc{\ref@jnl{Phys.~Rev.~C}}        
\DeclareRobustCommand\prd{\ref@jnl{Phys.~Rev.~D}}        
\DeclareRobustCommand\pre{\ref@jnl{Phys.~Rev.~E}}        
\DeclareRobustCommand\prl{\ref@jnl{Phys.~Rev.~Lett.}}    
\DeclareRobustCommand\pasp{\ref@jnl{PASP}}               
\DeclareRobustCommand\pasj{\ref@jnl{PASJ}}               
\DeclareRobustCommand\qjras{\ref@jnl{QJRAS}}             
\DeclareRobustCommand\skytel{\ref@jnl{S\&T}}             
\DeclareRobustCommand\solphys{\ref@jnl{Sol.~Phys.}}      
\DeclareRobustCommand\sovast{\ref@jnl{Soviet~Ast.}}      
\DeclareRobustCommand\ssr{\ref@jnl{Space~Sci.~Rev.}}     
\DeclareRobustCommand\zap{\ref@jnl{ZAp}}                 
\DeclareRobustCommand\nat{\ref@jnl{Nature}}              
\DeclareRobustCommand\iaucirc{\ref@jnl{IAU~Circ.}}
\DeclareRobustCommand\aplett{\ref@jnl{Astrophys.~Lett.}}
\DeclareRobustCommand\apspr{\ref@jnl{Astrophys.~Space~Phys.~Res.}}
\DeclareRobustCommand\bain{\ref@jnl{Bull.~Astron.~Inst.~Netherlands}}
\DeclareRobustCommand\fcp{\ref@jnl{Fund.~Cosmic~Phys.}}
\DeclareRobustCommand\gca{\ref@jnl{Geochim.~Cosmochim.~Acta}}
\DeclareRobustCommand\grl{\ref@jnl{Geophys.~Res.~Lett.}}
\DeclareRobustCommand\jcp{\ref@jnl{J.~Chem.~Phys.}}      
\DeclareRobustCommand\jgr{\ref@jnl{J.~Geophys.~Res.}}    
\DeclareRobustCommand\jqsrt{\ref@jnl{J.~Quant.~Spec.~Radiat.~Transf.}}
\DeclareRobustCommand\memsai{\ref@jnl{Mem.~Soc.~Astron.~Italiana}}
\DeclareRobustCommand\nphysa{\ref@jnl{Nucl.~Phys.~A}}
\DeclareRobustCommand\physrep{\ref@jnl{Phys.~Rep.}}
\DeclareRobustCommand\physscr{\ref@jnl{Phys.~Scr}}
\DeclareRobustCommand\planss{\ref@jnl{Planet.~Space~Sci.}}
\DeclareRobustCommand\procspie{\ref@jnl{Proc.~SPIE}}
\let\astap=\aap
\let\apjlett=\apjl
\let\apjsupp=\apjs
\let\applopt=\ao
\def\phn{\phantom{0}}
\def\phd{\phantom{.}}
\def\phs{\phantom{\ensuremath{-}}}
\def\phm#1{\phantom{#1}}
%    \end{macrocode}
% All these are redone as robust and using more modern commands where
% possible. The robustness is so they will not get expanded if used in
% abstracts and titles when these are passed to the ADS file. 
% \changes{3.16}{2003/04/07}{
%  Added \cs{Color} command for doing photometric colors.}
%    \begin{macrocode}
%%\def\sun{\hbox{$\odot$}}
\DeclareRobustCommand\sun{\ensuremath{\odot}}
%%\def\earth{\hbox{$\oplus$}}
\DeclareRobustCommand\earth{\ensuremath{\oplus}}
%%\def\lesssim{\mathrel{\hbox{\rlap{\hbox{\lower4pt\hbox{$\sim$}}}\hbox{$<$}}}}
%%\def\gtrsim{\mathrel{\hbox{\rlap{\hbox{\lower4pt\hbox{$\sim$}}}\hbox{$>$}}}}
\DeclareRobustCommand\lesssim{\mathrel{\hbox{\smaller\smaller\rlap{\hbox{\lower3pt\hbox{$\sim$}}}\raise1.5pt\hbox{$<$}}}}
\DeclareRobustCommand\gtrsim{\mathrel{\hbox{\smaller\smaller\rlap{\hbox{\lower3pt\hbox{$\sim$}}}\raise1.5pt\hbox{$>$}}}}
%%\def\lta{\mathrel{\spose{\lower 3pt\hbox{$\mathchar"218$}}
%%     \raise 2.0pt\hbox{$\mathchar"13C$}}}
%%\def\gta{\mathrel{\spose{\lower 3pt\hbox{$\mathchar"218$}}
%%     \raise 2.0pt\hbox{$\mathchar"13E$}}}
\DeclareRobustCommand\sq{\hbox{\rlap{$\sqcap$}$\sqcup$}}
%%\def\arcdeg{\hbox{$^\circ$}}
%%\def\arcmin{\hbox{$^\prime$}}
%%\def\arcsec{\hbox{$^{\prime\prime}$}}
\DeclareRobustCommand\arcdeg{\ensuremath{^\circ}}
\DeclareRobustCommand\arcmin{\ensuremath{^\prime}}
\DeclareRobustCommand\arcsec{\ensuremath{^{\prime\prime}}}
%%\def\fd{\hbox{$.\!\!^{\rm d}$}}
%%\def\fh{\hbox{$.\!\!^{\rm h}$}}
%%\def\fm{\hbox{$.\!\!^{\rm m}$}}
%%\def\fs{\hbox{$.\!\!^{\rm s}$}}
%%\def\fdg{\hbox{$.\!\!^\circ$}}
%%\def\farcm{\hbox{$.\mkern-4mu^\prime$}}
%%\def\farcs{\hbox{$.\!\!^{\prime\prime}$}}
%%\def\fp{\hbox{$.\!\!^{\scriptscriptstyle\rm p}$}}
\DeclareRobustCommand\fd{\ensuremath{.\!\!^\mathrm{d}}}
\DeclareRobustCommand\fh{\ensuremath{.\!\!^\mathrm{h}}}
\DeclareRobustCommand\fm{\ensuremath{.\!\!^\mathrm{m}}}
\DeclareRobustCommand\fs{\ensuremath{\mkern0.5mu.\mkern-5mu^\mathrm{s}\mkern-1mu}}
\DeclareRobustCommand\fdg{\ensuremath{.\!\!^\circ}}
\DeclareRobustCommand\farcm{\ensuremath{.\mkern-4mu^\prime}}
\DeclareRobustCommand\farcs{\ensuremath{\mkern1mu.\mkern-6mu^{\prime\prime}\mkern-1.5mu}}
%%\DeclareRobustCommand\farcs{\ensuremath{.\!\!^{\prime\prime}}}
\DeclareRobustCommand\fp{\ensuremath{.\!\!^{\scriptscriptstyle\rm p}}}
%%\def\micron{\hbox{$\mu$m}}
\DeclareRobustCommand\micron{\ensuremath{\mu\mathrm{m}}}
\let\la=\lesssim                        
\let\ga=\gtrsim
%%\def\case#1#2{\hbox{$\frac{#1}{#2}$}}
%%\def\slantfrac#1#2{\hbox{$\,^#1\!/_#2$}}
%%\def\onehalf{\slantfrac{1}{2}}
%%\def\onethird{\slantfrac{1}{3}}
%%\def\twothirds{\slantfrac{2}{3}}
%%\def\onequarter{\slantfrac{1}{4}}
%%\def\threequarters{\slantfrac{3}{4}}
\DeclareRobustCommand\case[2]{\ensuremath{\frac{#1}{#2}}}
\DeclareRobustCommand\slantfrac[2]{\ensuremath{\,^#1\!/_#2}}
\DeclareRobustCommand\onehalf{\slantfrac{1}{2}}
\DeclareRobustCommand\onethird{\slantfrac{1}{3}}
\DeclareRobustCommand\twothirds{\slantfrac{2}{3}}
\DeclareRobustCommand\onequarter{\slantfrac{1}{4}}
\DeclareRobustCommand\threequarters{\slantfrac{3}{4}}
%%\def\ubvr{\hbox{$U\!BV\!R$}}            
%%\def\ub{\hbox{$U\!-\!B$}}               
%%\def\bv{\hbox{$B\!-\!V$}}               
%%\def\vr{\hbox{$V\!-\!R$}}               
%%\def\ur{\hbox{$U\!-\!R$}}
%    \end{macrocode}
% New commands for ``photometric'' colors 
%    \begin{macrocode}
\DeclareRobustCommand\ubvr{\ensuremath{U\!BV\!R}}            
\DeclareRobustCommand\Color[2]{%
  \ensuremath{%
    #1\!\ifx#1V\!\fi-\!#2
  }}
\DeclareRobustCommand\ub{\Color{U}{B}}
\DeclareRobustCommand\bv{\Color{B}{V}}
\DeclareRobustCommand\vr{\Color{V}{R}}
\DeclareRobustCommand\ur{\Color{U}{R}}
%% New version of the \ion command, that actually works! 
%    \end{macrocode}
% \begin{macro}{\ion}
% This version of the |\ion| command allows you to write either
% |\ion{N}{4}| or |\ion{N}{iv}| in order to get N~IV. 
%    \begin{macrocode}
\DeclareRobustCommand\ion[2]{%
  \mbox{#1\kern0.2em%
    \smaller\rmfamily%
    \edef\@tempa{\@car#2\@nil}%
    \ifcat1\@tempa%
    \@Roman{#2}%
    \else%
    \uppercase{#2}%
    \fi}}
\newcounter{IonStage}
\DeclareRobustCommand\plainion[2]{\setcounter{IonStage}{#2}#1 \Roman{IonStage}}
%%\def\ion#1#2{\ensuremath{\mathrm{#1}\;}{\protect\small\rm\@Roman{#2}}}
%    \end{macrocode}
% \end{macro}
%    \begin{macrocode}
\DeclareRobustCommand\nodata{$\cdots$}
%    \end{macrocode}
% These are needed by the |\ion| command. Isn't there an easier way of
% doing this?
%    \begin{macrocode}
\let\@small=\small
\let\@normalsize=\normalsize
\let\@footnotesize=\footnotesize
\let\@scriptsize=\scriptsize
\let\@tiny=\tiny
\let\@large=\large
\let\@Large=\Large
\let\@LARGE=\LARGE
\let\@huge=\huge
\let\@Huge=\Huge
\newcommand{\smaller}{\small}
\renewcommand{\normalsize}{\renewcommand{\smaller}{\small}\@normalsize}
\renewcommand{\small}{\renewcommand{\smaller}{\footnotesize}\@small}
\renewcommand{\footnotesize}{\renewcommand{\smaller}{\scriptsize}\@footnotesize}
\renewcommand{\scriptsize}{\renewcommand{\smaller}{\tiny}\@scriptsize}
\renewcommand{\tiny}{\renewcommand{\smaller}{\tiny}\@tiny}
\renewcommand{\large}{\renewcommand{\smaller}{\normalsize}\@large}  
\renewcommand{\Large}{\renewcommand{\smaller}{\large}\@Large}  
\renewcommand{\LARGE}{\renewcommand{\smaller}{\Large}\@LARGE}  
\renewcommand{\huge}{\renewcommand{\smaller}{\LARGE}\@huge}  
\renewcommand{\Huge}{\renewcommand{\smaller}{\huge}\@Huge}  
%</package|adshtml-style>
%    \end{macrocode}


% \section{Other files}

% \subsection{The driver and make files}
% \label{sec:driver_make}
% These are put here for completeness, although they still need to be
% distributed separately since otherwise the user won't be able to
% unpack this file into its constituent parts.
% 
% The |Makefile| has two principal rules: |docs|, and |macros|. The
% |docs| rule produces a PDF and gzipped PS file of the documentation.
% I still haven't perfected the combination of rules required to get
% it to run \LaTeX{} enough times in all circumstances. The |macros|
% rule runs \LaTeX{} on |rmaa.ins| in order to unpack |rmaa.cls| and
% all the other little files (including the |Makefile| and the
% |rmaa.ins| files themselves\dots 
% 
% First set up lists of various types of file. 
%    \begin{macrocode}
%<*makefile>
DOCFILES=rmaa.pdf
MACRO_LTX_FILES=rmaa.cls rmaacite.sty adshtml.sty
SHELLSCRIPTS=makeads.sh
MACRO_OTHER_FILES=adshtml.perl sample.latex2html-init
MAKEFILE=Makefile
MACROFILES=$(MACRO_LTX_FILES) $(SHELLSCRIPTS) $(MACRO_OTHER_FILES)         
SKELETONFILES=ADS.tex 
%    \end{macrocode}
% Install the files read only into the package directory, specified via a path relative to the source directory. \changes{4.0}{2019/11/14}{Reorganize folder layout so that hardly anything is in the top level.  That way, things are clearer on github.}
%    \begin{macrocode}
INSTALL=install 
INSTALL_DIR=../packages
%    \end{macrocode}
% The default rule is to make everything. Installation must be
% explicitly asked for.
%    \begin{macrocode}
all: docs macros

install: install-docs install-macros

%    \end{macrocode}
% Most files are installed read-only to prevent editing. Skeleton
% files must be writeable. Must be careful not to overwrite the
% Makefile in the parent directory. 
%    \begin{macrocode}
install-macros: macros
        $(INSTALL) -m 444 $(MACRO_LTX_FILES) $(MACRO_OTHER_FILES) $(INSTALL_DIR)
        $(INSTALL) -m 555 $(SHELLSCRIPTS) $(INSTALL_DIR)
        $(INSTALL) -m 644 $(SKELETONFILES) $(INSTALL_DIR)

install-docs: docs
        $(INSTALL) rmaa.pdf $(INSTALL_DIR)/rm-fulldocs.pdf

%    \end{macrocode}
% Make the PDF and gzipped PS files. 
%    \begin{macrocode}
docs: $(DOCFILES)

%    \end{macrocode}
% Make everything else if either the |.dtx| or |.ins| files have
% changed. 
%    \begin{macrocode}
macros: rmaa.dtx rmaa.ins
%    \end{macrocode}
% Ensure that all the macro files can be written to then recreate
% them. The |touch| is there in case any of the files don't exist
% yet. 
%    \begin{macrocode}
        touch $(MACROFILES) $(MAKEFILE)
        chmod u+w $(MACROFILES) $(MAKEFILE) rmaa.ins
        latex "\input{rmaa.ins}"
%    \end{macrocode}
% \TeX{} automatically adds the |.tex| extension to extensionless
% files. We fix this and take the opportunity to convert
% back to TABs all that it has lovingly converted to spaces. 
%    \begin{macrocode}
        cp -f $(MAKEFILE) $(MAKEFILE).BCK && unexpand $(MAKEFILE).tex >| $(MAKEFILE) && rm $(MAKEFILE).tex
%    \end{macrocode}
% Make the shell scripts executable and everything unwriteable (so
% they don't get hand-edited by mistake). However, undo this for the
% files that the user is \emph{supposed} to modify. 
%    \begin{macrocode}
        chmod a+x $(SHELLSCRIPTS)
        chmod a-w $(MACROFILES) $(MAKEFILE)
        chmod u+w $(SKELETONFILES)

%    \end{macrocode}
% We now use PDF\LaTeX{} for this.
%    \begin{macrocode}
rmaa.idx: rmaa.dtx
        @echo "REGENERATING .idx FILE"
        pdflatex "\nonstopmode\input{rmaa.dtx}"

rmaa.aux: rmaa.dtx rmaa.ind rmaa.gls
        @echo "REGENERATING .aux FILE"
        pdflatex "\nonstopmode\input{rmaa.dtx}"

rmaa.ind: rmaa.idx
        @echo "REGENERATING INDEX"
        makeindex -s gind.ist rmaa

rmaa.gls: rmaa.glo
        @echo "REGENERATING CHANGELOG"
        makeindex -s gglo.ist -o rmaa.gls rmaa.glo

rmaa.pdf: rmaa.aux rmaa.ind rmaa.gls
        @echo "REGENERATING .pdf FILE"
        pdflatex "\nonstopmode\input{rmaa.dtx}"

clean: 
        rm -f rmaa.aux rmaa.dvi rmaa.ind rmaa.gls
%</makefile>
%    \end{macrocode}

% The |rmaa.ins| file uses the \textsf{docstrip} package to split out
% all the files and strip off this documentation.
%    \begin{macrocode}
%<*ins>
%%
%% Installation batch file for the RMAA LaTeX package
%%
\def\batchfile{rmaa.ins}
\input docstrip.tex
\preamble

PLEASE DO NOT EDIT THIS FILE! It is generated automatically from
the file rmaa.dtx using the accompanying Makefile and rmaa.ins 
Any hand-edited changes will be lost

A LaTeX document class for the RMxAA

Author: William Henney (w.henney@irya.unam.mx)

Latest version available from: http://www.irya.unam.mx/rmaa

Copyright (C) 1999-2019 Instituto de Astronomia, UNAM, Mexico

Permission is hereby granted, free of charge, to any person obtaining
a copy of this software and associated documentation files (the
"Software"), to deal in the Software without restriction, including
without limitation the rights to use, copy, modify, merge, publish,
distribute, sublicense, and/or sell copies of the Software, and to
permit persons to whom the Software is furnished to do so, subject to
the following conditions:

The above copyright notice and this permission notice shall be
included in all copies or substantial portions of the Software.

THE SOFTWARE IS PROVIDED "AS IS", WITHOUT WARRANTY OF ANY KIND,
EXPRESS OR IMPLIED, INCLUDING BUT NOT LIMITED TO THE WARRANTIES OF
MERCHANTABILITY, FITNESS FOR A PARTICULAR PURPOSE AND
NONINFRINGEMENT. IN NO EVENT SHALL THE AUTHORS OR COPYRIGHT HOLDERS BE
LIABLE FOR ANY CLAIM, DAMAGES OR OTHER LIABILITY, WHETHER IN AN ACTION
OF CONTRACT, TORT OR OTHERWISE, ARISING FROM, OUT OF OR IN CONNECTION
WITH THE SOFTWARE OR THE USE OR OTHER DEALINGS IN THE SOFTWARE.
\endpreamble
%    \end{macrocode}
% Clobber files without asking
%    \begin{macrocode}
\askforoverwritefalse
%    \end{macrocode}
% Make the \LaTeX{} macro files (class and styles)
%    \begin{macrocode}
\generate{% 
  \file{rmaa.cls}{\from{rmaa.dtx}{package}}
  \file{rmaa.ins}{\from{rmaa.dtx}{ins}}
  \file{rmaacite.sty}{\from{rmaa.dtx}{rmaacite}}
  \file{adshtml.sty}{\from{rmaa.dtx}{adshtml-style}}
  \file{ADS.tex}{\from{rmaa.dtx}{ads-container}}
  }
%    \end{macrocode}
% Make the other non-\LaTeX{} files. These are done separately since
% they can't use `\%\%' as a comment character for the
% preamble/postamble (and Perl files can't have the postamble I don't
% think). 
%    \begin{macrocode}
\generate{\nopreamble\nopostamble% 
  \file{Makefile}{\from{rmaa.dtx}{makefile}}
  \file{adshtml.perl}{\from{rmaa.dtx}{adshtml-perl}}
  \file{makeads.sh}{\from{rmaa.dtx}{ads-shell-script}}
  \file{sample.latex2html-init}{\from{rmaa.dtx}{latex2html-init}}
  }

%% Local Variables: 
%% mode: LaTeX
%% mode: Font-lock
%% End:
%</ins>
%    \end{macrocode}

% \subsection{Extended citation syntax}
% Modified from a file called mnrascite.sty
% I no longer remember where I stole the original version of this from
% 
% Alterations to the cite commands to make them do more sensible things.
% N.B. the bibliography style should produce entries which have labels where
% the date is delimited by angle brackets '<>' for these macros to work
% properly. They will not give very sensible error messages if you do not do 
% this properly.
% 
% A new command |\scite{key}| is defined which produces a citation which is 
% suitable for insertion in a sentence. i.e., it does not put brackets around 
% the whole citation, only the date. Also |\pcite{key}| produce output with no
% parentheses whatsoever. 
% 
% Should redefine |\@bibitem| really to have any hope of compatibility
% the command |\lbibcite| is defined to deal with the citations.
%    \begin{macrocode}
%<*rmaacite>
\def\bibitem[#1<#2>]#3{\par\hangindent=1.5em\hangafter=1
  \if@filesw
  {\def\protect##1{\string ##1\space}
    \immediate\write\@auxout{\string\lbibcite{#3}{#1}{#2}}}
  \fi}
\@redefinebibitem % what it says
%    \end{macrocode}
% Note that the |\lbibcite| defines two types of citation one contains the names,
% the other the year
%    \begin{macrocode}
\def\lbibcite#1#2#3{\global\@namedef{b@#1}{#2}
                    \global\@namedef{by@#1}{#3}}

\def\@bcite#1#2{({#1\if@tempswa , #2\fi})}
\def\@pcite#1#2{#1\if@tempswa , #2\fi}
%    \end{macrocode}
% This declaration contains what is put between the name and the date
% \changes{3.11}{2003/01/31}{Moved the \cs{articlename} to after the key for consistency
% with the new natbib stuff.}
%    \begin{macrocode}
\def\@brcite{\ (}
\def\@ebrcite{)}
\def\@fmttop#1{\@fmtcite[\csname #1\endcsname]}
\def\@citex[#1]#2{%
%%  \if@filesw\immediate\write\@auxout{\string\citation{#2}}\fi
%% WJH added 16 Jan 2002
  \@for\@citeb:=#2\do{%
    \edef\@citeb{\expandafter\@firstofone\@citeb:\articlename}%
    \if@filesw\immediate\write\@auxout{\string\citation{\@citeb}}\fi}%
%% to here
  \def\@citea{}%
  \@cite{\@for\@citeb:=#2\do
    {%
      \edef\@citeb{\expandafter\@firstofone\@citeb:\articlename}%
      \@citea\def\@citea{;\ }\@ifundefined
      {b@\@citeb}{{\bf ?}\@warning {Citation `\@citeb' on page \thepage \space undefined}}%
      {\csname b@\@citeb\endcsname\@citebro\csname by@\@citeb\endcsname\@citebrc}}}%
  {#1}}

\def\cite{\@ifnextchar [{\let\@citebro=\space
                          \let\@citebrc=\relax
                          \let\@cite=\@bcite\@tempswatrue\@citex}
                        {\let\@citebro=\space
                          \let\@citebrc=\relax
                          \let\@cite=\@bcite\@tempswafalse\@citex[]}}
\def\pcite{\@ifnextchar [{\let\@citebro=\space
                          \let\@citebrc=\relax
                          \let\@cite=\@pcite\@tempswatrue\@citex}
                        {\let\@citebro=\space
                          \let\@citebrc=\relax
                          \let\@cite=\@pcite\@tempswafalse\@citex[]}}
\def\scite{\@ifnextchar [{\let\@citebro=\@brcite
                          \let\@citebrc=\@ebrcite
                          \let\@cite=\@pcite\@tempswatrue\@citex}
                        {\let\@citebro=\@brcite
                          \let\@citebrc=\@ebrcite
                          \let\@cite=\@pcite\@tempswafalse\@citex[]}}

\renewenvironment{thebibliography}
  {\section*{REFERENCES}\vspace*{\afterbiblioskip}\small
    \if@book\RereadBiblio\fi
    \sloppy\clubpenalty4000\widowpenalty4000
    \sfcode`\.=1000\relax
    \setlength\parindent{\z@}\setlength\parskip{\z@}}
  {\relax\normalsize\@outputaddressesifnecessary} 
%</rmaacite>
%    \end{macrocode}

% \subsection{Auxilliary files for generating ADS data}

% The ADS data are written by \LaTeX{} to a |.ads| file. This file
% then has to be included in the following file (|ADS.tex|),
% which is run through \LaTeX{}2HTML (see below). 

%    \begin{macrocode}
%<*ads-container>
%%
%% Instructions:
%% Edit the \input{} line below to reflect the name of your file. 
\documentclass{article}
\usepackage[latin1]{inputenc}
\usepackage{adshtml}
\newcommand{\ADSfield}[2]
{\%#1  #2\\}
%%\newcommand{\ADSvfield}[2]
%%{\%#1 \begin{verbatim}#2\end{verbatim}}
\newenvironment{ADSarticle}{}{\par\bigskip}
\renewcommand{\bibitem}[2][]{\par}
\setlength{\parindent}{0pt}
\renewcommand{\ThisVolume}[1]{%  
  2001, \RMAASC, 12, \oldpageref{#1:RMAAFirstPage}}  
 

\begin{document}
\makeatletter
\input{rmbooksamp.ads}
\end{document}
%</ads-container>
%    \end{macrocode}

% This is the |adshtml.sty| style file used by |ADS.tex|. It just
% defines a few 
% macros that might be used in abstracts, etc (really it should define
% them all). This now gets all its code from
% Sec~\ref{sec:user-commands}. 
%    \begin{macrocode}
%<*adshtml-style>
%% Nothing else needs to go here now. 
%</adshtml-style>
%    \end{macrocode}
% The init file for latex2html.
%    \begin{macrocode}
%<*latex2html-init>
$PDFTOCAIRO = "pdftocairo";
$PDFLATEX = "pdflatex";
$LATEX = "latex";
$MAX_SPLIT_DEPTH = 0;
$HTML_VERSION = "4.0,math";
$NO_SIMPLE_MATH = 1;
$ADDRESS = "<p><address><a href=\"mailto:rmaa\@astro.unam.mx\">rmaa\@astro.unam.mx</address></p>"; 
$INFO = "<p>This document was created using latex2html and the rmaa latex macros</p>"; 
$NO_NAVIGATION = 1;
1; 
%</latex2html-init>
%    \end{macrocode}
% This is the |adshtml.perl| file that is used by \LaTeX{}2HTML. Some
% bizarre \TeX{}y command sometimes show up so we tell it to ignore
% these. All the astronomy macros that come from AAS are processed in
% \TeX{}. That way they get made to images but with an |ALT| tag that
% has the original code, which is what will be picked up by |lynx|,
% with the result that they get passed intact into the final |.txt|
% file. 
%    \begin{macrocode}
%<*adshtml-perl>
package main;

$dum = &do_cmd_makeatletter();

sub do_cmd_RMAA {
    local($_) = @_;
    "RMxAA" . $_;
}

sub do_cmd_RMAASC {
    local($_) = @_;
    "RMxAASC" . $_;
}

sub do_cmd_etal {
    local($_) = @_;
    "et al\." . $_;
}

sub do_cmd_ibidrule {
    local($_) = @_;
    "<img src=\"foo\" alt=\"\ibidrule\">" . $_;
}


sub do_cmd_bibitem {
    local($_) = @_;
    local($key);
    local($opt, $dummy)=&get_next_optional_argument; 
    $key = &missing_braces unless
        ((s/$next_pair_pr_rx/$key = $2; ''/eo)
        ||(s/$next_pair_rx/$key = $2; ''/eo));
    "<p>" . $_; 
}

&ignore_commands (<<_RAW_ARG_CMDS_);
lowercase
adjustfinalcols
textsuperscript # {}
newline # []
thinmuskip
mskip
smallskip
medskip
bigskip
hspace # {}
hspacestar # {}
vspace # {}
vspacestar # {}
mathrm
_at_mathrm
spacefactor
discretionary
unhbox 
voidb_at_x 
penalty 
_at_M
_at_m
null
IeC
_RAW_ARG_CMDS_

&process_commands_in_tex (<<_RAW_ARG_CMDS_);
sun
earth
lesssim
gtrsim
sq
argdeg
arcmin
arcsec
fd
fh
fm
fs
fdg
farcm
farcs
fp
micron
la
ga
case # {} # {}
slantfrac # {} # {}
onehalf
onethird
twothirds
onequarter
threequarters
ubvr
ub
bv
vr
ur
aj
araa
apj
apjl
apjs
ao
apss
aap
aapr
aaps
azh
baas
jrasc
memras
mnras
pra
prb
prc
prd
pre
prl
pasp
pasj
qjras
skytel
solphys
sovast
ssr
zap
nat
iaucirc
aplett
apspr
bain
fcp
gca
grl
jcp
jgr
jqsrt
memsai
nphysa
physrep
physscr
planss
procspie
_RAW_ARG_CMDS_

1;                              # This must be the last line
%</adshtml-perl> 
%    \end{macrocode}

% This is a shell script |makeads.sh| that generates the |_ads.txt| file from
% |_ads.tex| by using \LaTeX{}2HTML and |lynx|.  \changes{4.0}{2019/11/13}{Make sure that lynx writes utf-8 output instead of latin-1.}
% \changes{2.2beta13}{2002/03/12}{Complete rewrite of |makeads.sh|. Now much 
% improved I think}
% \changes{4.0}{2019/11/13}{Fix the at-sign substitution so that it only eliminates \cs{@} and does not mess with email addresses.}

%    \begin{macrocode}
%<*ads-shell-script>
#! /bin/sh
if [ -z $1 ]; then
    echo "Usage: $0 basename"
    exit
fi
BASENAME=$1
ADSBASENAME=${BASENAME}_ads
AUXFILE=$BASENAME.aux
TMPFILE=tmp.$$


echo "Generating $ADSBASENAME.tex"

# For some reason, must double-up \\ to \\\\ but not \ to \\ - odd!
cat > $ADSBASENAME.tex <<EOF
%% PLEASE DO NOT EDIT THIS FILE!  It is generated automatically from
%% the script makeads.sh.  Any hand-edited changes will be lost
%% 
\documentclass{article}
\usepackage[latin1]{inputenc}
\usepackage{html}
\usepackage{adshtml}
\newcommand{\ADSfield}[2]{\%#1  #2\\\\}
\newenvironment{ADSarticle}{}{\par\bigskip}
\setlength{\parindent}{0pt}
\newcounter{istage}
\renewcommand{\ion}[2]{\setcounter{istage}{#2}#1 \Roman{istage}}
\begin{document}
\makeatletter
EOF
echo >> $ADSBASENAME.tex "\input{$BASENAME.ads2}" 
echo >> $ADSBASENAME.tex "\end{document}"



echo "Generating $ADSBASENAME.sed" 

# Find type of journal
grep '\\documentclass\[.*proceedings.*\]' $BASENAME.tex >/dev/null \
     && JOURNAL='RMxAC' \
     || JOURNAL='RMxAA'

# Find the year and volume number
YEAR=`grep '\\SetYear' $BASENAME.tex | sed 's/.*\\SetYear{\(.*\)}.*/\1/'`
VOLUME=`grep '\\SetVolume' $BASENAME.tex | sed 's/.*\\SetVolume{\(.*\)}.*/\1/'`

# Generate an auxilliary sed script 
cat > $TMPFILE <<EOF
s|\\\\newlabel{\(.*\):RMAAFirstPage}{{.*}{\(.*\)}}|s/\\\\\\\\ThisVolume{\1}/$YEAR, RMxAC, $VOLUME, \2/g|g
EOF

# Now go through the .aux file, looking for \newlabels and use the sed 
# script just generated in order to make another sed script, one that 
# will replace \ThisVolume{LABEL} with the correct reference
grep '^\\newlabel{.*:RMAAFirstPage}' $AUXFILE \
     | sed -f $TMPFILE > $ADSBASENAME.sed && rm $TMPFILE

# Apend a further sed command to get rid of \@, since redefining 
# it in _ads.tex didn't work
echo >> $ADSBASENAME.sed 's/\\\\@//g'

echo "Generating $BASENAME.ads2 from $ADSBASENAME.sed and $BASENAME.ads" 

# Apply this to the .ads file, generating the .ads2 file, which is included
# by the _ads.tex file.
sed -f $ADSBASENAME.sed $BASENAME.ads > $BASENAME.ads2

echo "Generating $ADSBASENAME.txt from $ADSBASENAME.tex and $BASENAME.ads2"

# Now run latex2html on the _ads.tex file
latex2html -no_navigation -info 0 -address 0 -html_version 4.0,math $ADSBASENAME

# Then, convert to utf8 with lynx and tidy it up a bit
lynx --justify=off --width=100 --dump --nolist --display_charset=utf-8 \
     $ADSBASENAME/$ADSBASENAME.html | cut -c4- \
     | sed 's/ibidrule/\\ibidrule/g' > $ADSBASENAME.txt
%</ads-shell-script>
%    \end{macrocode}

% \section{Epilogue}
% Err, that's it \dots

\endinput


% \endinput
% Local Variables:
% mode: doctex
% TeX-master: t
% End:
"

clean: 
        rm -f rmaa.aux rmaa.dvi rmaa.ind rmaa.gls
%</makefile>
%    \end{macrocode}

% The |rmaa.ins| file uses the \textsf{docstrip} package to split out
% all the files and strip off this documentation.
%    \begin{macrocode}
%<*ins>
%%
%% Installation batch file for the RMAA LaTeX package
%%
\def\batchfile{rmaa.ins}
\input docstrip.tex
\preamble

PLEASE DO NOT EDIT THIS FILE! It is generated automatically from
the file rmaa.dtx using the accompanying Makefile and rmaa.ins 
Any hand-edited changes will be lost

A LaTeX document class for the RMxAA

Author: William Henney (w.henney@irya.unam.mx)

Latest version available from: http://www.irya.unam.mx/rmaa

Copyright (C) 1999-2019 Instituto de Astronomia, UNAM, Mexico

Permission is hereby granted, free of charge, to any person obtaining
a copy of this software and associated documentation files (the
"Software"), to deal in the Software without restriction, including
without limitation the rights to use, copy, modify, merge, publish,
distribute, sublicense, and/or sell copies of the Software, and to
permit persons to whom the Software is furnished to do so, subject to
the following conditions:

The above copyright notice and this permission notice shall be
included in all copies or substantial portions of the Software.

THE SOFTWARE IS PROVIDED "AS IS", WITHOUT WARRANTY OF ANY KIND,
EXPRESS OR IMPLIED, INCLUDING BUT NOT LIMITED TO THE WARRANTIES OF
MERCHANTABILITY, FITNESS FOR A PARTICULAR PURPOSE AND
NONINFRINGEMENT. IN NO EVENT SHALL THE AUTHORS OR COPYRIGHT HOLDERS BE
LIABLE FOR ANY CLAIM, DAMAGES OR OTHER LIABILITY, WHETHER IN AN ACTION
OF CONTRACT, TORT OR OTHERWISE, ARISING FROM, OUT OF OR IN CONNECTION
WITH THE SOFTWARE OR THE USE OR OTHER DEALINGS IN THE SOFTWARE.
\endpreamble
%    \end{macrocode}
% Clobber files without asking
%    \begin{macrocode}
\askforoverwritefalse
%    \end{macrocode}
% Make the \LaTeX{} macro files (class and styles)
%    \begin{macrocode}
\generate{% 
  \file{rmaa.cls}{\from{rmaa.dtx}{package}}
  \file{rmaa.ins}{\from{rmaa.dtx}{ins}}
  \file{rmaacite.sty}{\from{rmaa.dtx}{rmaacite}}
  \file{adshtml.sty}{\from{rmaa.dtx}{adshtml-style}}
  \file{ADS.tex}{\from{rmaa.dtx}{ads-container}}
  }
%    \end{macrocode}
% Make the other non-\LaTeX{} files. These are done separately since
% they can't use `\%\%' as a comment character for the
% preamble/postamble (and Perl files can't have the postamble I don't
% think). 
%    \begin{macrocode}
\generate{\nopreamble\nopostamble% 
  \file{Makefile}{\from{rmaa.dtx}{makefile}}
  \file{adshtml.perl}{\from{rmaa.dtx}{adshtml-perl}}
  \file{makeads.sh}{\from{rmaa.dtx}{ads-shell-script}}
  \file{sample.latex2html-init}{\from{rmaa.dtx}{latex2html-init}}
  }

%% Local Variables: 
%% mode: LaTeX
%% mode: Font-lock
%% End:
%</ins>
%    \end{macrocode}

% \subsection{Extended citation syntax}
% Modified from a file called mnrascite.sty
% I no longer remember where I stole the original version of this from
% 
% Alterations to the cite commands to make them do more sensible things.
% N.B. the bibliography style should produce entries which have labels where
% the date is delimited by angle brackets '<>' for these macros to work
% properly. They will not give very sensible error messages if you do not do 
% this properly.
% 
% A new command |\scite{key}| is defined which produces a citation which is 
% suitable for insertion in a sentence. i.e., it does not put brackets around 
% the whole citation, only the date. Also |\pcite{key}| produce output with no
% parentheses whatsoever. 
% 
% Should redefine |\@bibitem| really to have any hope of compatibility
% the command |\lbibcite| is defined to deal with the citations.
%    \begin{macrocode}
%<*rmaacite>
\def\bibitem[#1<#2>]#3{\par\hangindent=1.5em\hangafter=1
  \if@filesw
  {\def\protect##1{\string ##1\space}
    \immediate\write\@auxout{\string\lbibcite{#3}{#1}{#2}}}
  \fi}
\@redefinebibitem % what it says
%    \end{macrocode}
% Note that the |\lbibcite| defines two types of citation one contains the names,
% the other the year
%    \begin{macrocode}
\def\lbibcite#1#2#3{\global\@namedef{b@#1}{#2}
                    \global\@namedef{by@#1}{#3}}

\def\@bcite#1#2{({#1\if@tempswa , #2\fi})}
\def\@pcite#1#2{#1\if@tempswa , #2\fi}
%    \end{macrocode}
% This declaration contains what is put between the name and the date
% \changes{3.11}{2003/01/31}{Moved the \cs{articlename} to after the key for consistency
% with the new natbib stuff.}
%    \begin{macrocode}
\def\@brcite{\ (}
\def\@ebrcite{)}
\def\@fmttop#1{\@fmtcite[\csname #1\endcsname]}
\def\@citex[#1]#2{%
%%  \if@filesw\immediate\write\@auxout{\string\citation{#2}}\fi
%% WJH added 16 Jan 2002
  \@for\@citeb:=#2\do{%
    \edef\@citeb{\expandafter\@firstofone\@citeb:\articlename}%
    \if@filesw\immediate\write\@auxout{\string\citation{\@citeb}}\fi}%
%% to here
  \def\@citea{}%
  \@cite{\@for\@citeb:=#2\do
    {%
      \edef\@citeb{\expandafter\@firstofone\@citeb:\articlename}%
      \@citea\def\@citea{;\ }\@ifundefined
      {b@\@citeb}{{\bf ?}\@warning {Citation `\@citeb' on page \thepage \space undefined}}%
      {\csname b@\@citeb\endcsname\@citebro\csname by@\@citeb\endcsname\@citebrc}}}%
  {#1}}

\def\cite{\@ifnextchar [{\let\@citebro=\space
                          \let\@citebrc=\relax
                          \let\@cite=\@bcite\@tempswatrue\@citex}
                        {\let\@citebro=\space
                          \let\@citebrc=\relax
                          \let\@cite=\@bcite\@tempswafalse\@citex[]}}
\def\pcite{\@ifnextchar [{\let\@citebro=\space
                          \let\@citebrc=\relax
                          \let\@cite=\@pcite\@tempswatrue\@citex}
                        {\let\@citebro=\space
                          \let\@citebrc=\relax
                          \let\@cite=\@pcite\@tempswafalse\@citex[]}}
\def\scite{\@ifnextchar [{\let\@citebro=\@brcite
                          \let\@citebrc=\@ebrcite
                          \let\@cite=\@pcite\@tempswatrue\@citex}
                        {\let\@citebro=\@brcite
                          \let\@citebrc=\@ebrcite
                          \let\@cite=\@pcite\@tempswafalse\@citex[]}}

\renewenvironment{thebibliography}
  {\section*{REFERENCES}\vspace*{\afterbiblioskip}\small
    \if@book\RereadBiblio\fi
    \sloppy\clubpenalty4000\widowpenalty4000
    \sfcode`\.=1000\relax
    \setlength\parindent{\z@}\setlength\parskip{\z@}}
  {\relax\normalsize\@outputaddressesifnecessary} 
%</rmaacite>
%    \end{macrocode}

% \subsection{Auxilliary files for generating ADS data}

% The ADS data are written by \LaTeX{} to a |.ads| file. This file
% then has to be included in the following file (|ADS.tex|),
% which is run through \LaTeX{}2HTML (see below). 

%    \begin{macrocode}
%<*ads-container>
%%
%% Instructions:
%% Edit the \input{} line below to reflect the name of your file. 
\documentclass{article}
\usepackage[latin1]{inputenc}
\usepackage{adshtml}
\newcommand{\ADSfield}[2]
{\%#1  #2\\}
%%\newcommand{\ADSvfield}[2]
%%{\%#1 \begin{verbatim}#2\end{verbatim}}
\newenvironment{ADSarticle}{}{\par\bigskip}
\renewcommand{\bibitem}[2][]{\par}
\setlength{\parindent}{0pt}
\renewcommand{\ThisVolume}[1]{%  
  2001, \RMAASC, 12, \oldpageref{#1:RMAAFirstPage}}  
 

\begin{document}
\makeatletter
%% Demonstration file for formatting an entire volume
\documentclass[book,journal,twocolumn]{rmaa}
%%%%%%%%%%%%%%%%%%%%%%%%%%%%%%%%%%%%%%%%%%%%%%%%%%%%%%%%%%%%%%%%
%%%%%%%%%%%%%%%%            PREAMBLE            %%%%%%%%%%%%%%%%
%%%%%%%%%%%%%%%%%%%%%%%%%%%%%%%%%%%%%%%%%%%%%%%%%%%%%%%%%%%%%%%%

\usepackage{paralist}
\usepackage{psfrag,color}
% Needed for landscape tables
\usepackage{pdflscape}

% Used in rm-minimal-test
\usepackage{kantlipsum}

% publication details of volume
\SetVolume{9}
\SetYear{2000}
\SetMonthNumeric{06}
\SetMes{junio}
% details of meeting
\SetConfTitle{Astrophysical Plasmas: Codes, Models, \& Observations}
\SetReunion{M\'exico, D. F.}{octubre 25--29, 1999}
\SetEditors{S. Jane Arthur, Nancy S. Brickhouse, \& Jos\'e Franco}

% this writes the .idx file that is used to generate the index 
\makeindex

%% This command makes the session headings come out bold in the
%% Table-of-Contents. Comment it out if you don't like it.
\BoldSessionTitles

%% This command makes a title page for each session
%% Comment it out if you don't like it.
\SessionTitlePages

%% Use this command if you want the abstract-only papers to be
%% indexed by ADS
\abstractsADS

%% At the special request of Alberto Carramin~ana, the following is a
%% way to suppress the author adresses at the end of each article
%\suppressfulladdresses

%% Use this command to add Received/accepted dates for each paper
% \printdates


\begin{document}

%%%%%%%%%%%%%%%%%%%%%%%%%%%%%%%%%%%%%%%%%%%%%%%%%%%%%%%%%%%%%%%%b
%%%%%%%%%%%%%%%%          FRONT MATTER          %%%%%%%%%%%%%%%%
%%%%%%%%%%%%%%%%%%%%%%%%%%%%%%%%%%%%%%%%%%%%%%%%%%%%%%%%%%%%%%%%
\frontmatter % start roman numbering

% format the TOC
\renewcommand{\RMAAbannerimage}{rmsc_blank}
\tableofcontents

%% preface, group photo, list of participants go here
%% use the `frontsection' environment (see editors' guide)
\begin{frontsection}{Preface}{%
    \listofauthors{S.~J.~Arthur, N.~S.~Brickhouse, J.~Franco}
    \indexauthor{Arthur, S. J.}
    \indexauthor{Brickhouse, N.}
    \indexauthor{Franco, J.}
    }
  These are the proceedings of a very nice conference that Jane and
  Pepe organised. This is the preface, which explains what a wonderful
  time was had by all and how grateful they are for all the people who
  gave them money or helped out.
  
  This is another paragraph of the preface, in case they have anything
  else to say.  For example, the Scientific Organizing Committee
  consisted of
  \begin{quote}
    J. Franco (Mexico, Chair), D. Cox (USA) and various other notables
  \end{quote}
  while the Local Organizing Committee comprised
  \begin{quote}
    S. J. Arthur (Chair), S. Kurtz, etc.
  \end{quote}

  \EditorSignature

\end{frontsection}

\begin{frontsection}[Photo]{Group Photograph}{\notitle}
\null\vfill
\centerline{Here goes the photograph.}
\vfill
\end{frontsection}

\begin{frontsection}[Participants]{List of Participants}{}
%% Here you put a formatted list of participants
\vfil
\dotfill First page of list of participants \dotfill
\vfil\clearpage

\null\vfil
\dotfill Second page of list of participants \dotfill
\vfil

\end{frontsection}

\begin{frontsection}[XXX]{Another front section}{}

\end{frontsection}

%%%%%%%%%%%%%%%%%%%%%%%%%%%%%%%%%%%%%%%%%%%%%%%%%%%%%%%%%%%%%%%%
%%%%%%%%%%%%%%%%          MAIN MATTER           %%%%%%%%%%%%%%%%
%%%%%%%%%%%%%%%%%%%%%%%%%%%%%%%%%%%%%%%%%%%%%%%%%%%%%%%%%%%%%%%%
\mainmatter % start arabic numbering

\begin{articles}
  %% Main articles go here
  %% Use \article{ARTICLENAME} for each individual article
  %% If you want any text to appear in brackets after the article
  %% title in the TOC and index, then put it as an optional first
  %% argument, e.g.,  \article[Invited Talk]{ARTICLENAME} 

  \session{The Contributed Papers}
  \article{rm-journal-example} % An illustrative article
  \article{rm-onepage} % another article, which has a plate
  \article{rm-minimal-test}
  \session{Conference Summary}
  \article[Invited Review]{rmsummsamp} % An illustrative conference summary

\end{articles}

%%%%%%%%%%%%%%%%%%%%%%%%%%%%%%%%%%%%%%%%%%%%%%%%%%%%%%%%%%%%%%%%
%%%%%%%%%%%%%%%%          BACK MATTER           %%%%%%%%%%%%%%%%
%%%%%%%%%%%%%%%%%%%%%%%%%%%%%%%%%%%%%%%%%%%%%%%%%%%%%%%%%%%%%%%%
\backmatter

%% Abstracts of presentations for which the
%% authors have not submitted a paper
\begin{abstracts}

  \title{An abstract of a poster}
  \author{W. J. Henney\altaffilmark{1}}
  \listofauthors{W.~J.~Henney}
  \indexauthor{Henney, W. J}
  \altaffiltext{1}{Where I am from done with altaffilmark/text}
  \abstract{This is a test abstract.   This is a test abstract.   This
    is a test abstract.   This is a test abstract.}
  

  \title{An abstract of another  poster}
  \author{W.~J.~Henney\altaffil{This is done with altaffil} and
    A.~N.~Other\altaffil{So is this} }
  \listofauthors{W.~J.~Henney \& A.~N.~Other}
  \indexauthor{Henney, W. J}
  \indexauthor{Other, A. N.}
  \abstract{This is a test abstract.   This is a test abstract.   This
    is a test abstract.   This is a test abstract.}

  \title{An abstract of yet another  poster}
  \author{W.~J.~Henney\altaffil{This is done with altaffil} and
    A.~N.~Other\altaffil{So is this} }
  \listofauthors{W.~J.~Henney \& A.~N.~Other}
  \indexauthor{Henney, W. J}
  \indexauthor{Other, A. N.}
  \abstract{This is a test abstract.   This is a test abstract.   This
    is a test abstract.   This is a test abstract.}

  \title{The abstract of a poster with a very long title to see how
    well it gets split up into lines}
  \author{A.~N.~Other\altaffil{The affiliation of A.~N.~Other is also
      very long so that it will go over several lines. His email
      address is anon@nowhere.org}}
  \listofauthors{A.~N.~Other}
  \indexauthor{Other, A. N.}
  \abstract{We will try and make this abstract longer so there is more
    stuff to put on the page and it doesn't look too empty. However, I
    have nothing whatsoever to say so will have to stop here. Let's
    repeat that a few times. We will try and make this abstract longer
    so there is more stuff to put on the page and it doesn't look too
    empty. However, I have nothing whatsoever to say so will have to
    stop here. Let's repeat that a few times. We will try and make
    this abstract longer so there is more stuff to put on the page and
    it doesn't look too empty. However, I have nothing whatsoever to
    say so will have to stop here. Let's repeat that a few times. We
    will try and make this abstract longer so there is more stuff to
    put on the page and it doesn't look too empty. However, I have
    nothing whatsoever to say so will have to stop here. Let's repeat
    that a few times.}
    
  \title[Comment to go in TOC]{The abstract of a second poster with a
    very long title to see how well it gets split up into lines}
  \author{A.~N.~Other\altaffil{The affiliation of A.~N.~Other is also
      very long so that it will go over several lines. His email
      address is anon@nowhere.org}}
  \listofauthors{A.~N.~Other}
  \indexauthor{Other, A. N.}
  \abstract{We will try and make this abstract longer so there is more
    stuff to put on the page and it doesn't look too empty. However, I
    have nothing whatsoever to say so will have to stop here. Let's
    repeat that a few times. We will try and make this abstract longer
    so there is more stuff to put on the page and it doesn't look too
    empty. However, I have nothing whatsoever to say so will have to
    stop here. Let's repeat that a few times. We will try and make
    this abstract longer so there is more stuff to put on the page and
    it doesn't look too empty. However, I have nothing whatsoever to
    say so will have to stop here. Let's repeat that a few times. We
    will try and make this abstract longer so there is more stuff to
    put on the page and it doesn't look too empty. However, I have
    nothing whatsoever to say so will have to stop here. Let's repeat
    that a few times.}
    
  \title{The abstract of a third poster with a very long title to see how
    well it gets split up into lines}
  \author{A.~N.~Other\altaffil{The affiliation of A.~N.~Other is also
      very long so that it will go over several lines. His email
      address is anon@nowhere.org}}
  \listofauthors{A.~N.~Other}
  \indexauthor{Other, A. N.}
  \abstract{We will try and make this abstract longer so there is more
    stuff to put on the page and it doesn't look too empty. However, I
    have nothing whatsoever to say so will have to stop here. Let's
    repeat that a few times. We will try and make this abstract longer
    so there is more stuff to put on the page and it doesn't look too
    empty. However, I have nothing whatsoever to say so will have to
    stop here. Let's repeat that a few times. We will try and make
    this abstract longer so there is more stuff to put on the page and
    it doesn't look too empty. However, I have nothing whatsoever to
    say so will have to stop here. Let's repeat that a few times. We
    will try and make this abstract longer so there is more stuff to
    put on the page and it doesn't look too empty. However, I have
    nothing whatsoever to say so will have to stop here. Let's repeat
    that a few times.}
    
  \title{The abstract of a fourth poster with a very long title to see how
    well it gets split up into lines}
  \author{A.~N.~Other\altaffil{The affiliation of A.~N.~Other is also
      very long so that it will go over several lines. His email
      address is anon@nowhere.org}}
  \listofauthors{A.~N.~Other}
  \indexauthor{Other, A. N.}
  \abstract{We will try and make this abstract longer so there is more
    stuff to put on the page and it doesn't look too empty. However, I
    have nothing whatsoever to say so will have to stop here. Let's
    repeat that a few times. We will try and make this abstract longer
    so there is more stuff to put on the page and it doesn't look too
    empty. However, I have nothing whatsoever to say so will have to
    stop here. Let's repeat that a few times. We will try and make
    this abstract longer so there is more stuff to put on the page and
    it doesn't look too empty. However, I have nothing whatsoever to
    say so will have to stop here. Let's repeat that a few times. We
    will try and make this abstract longer so there is more stuff to
    put on the page and it doesn't look too empty. However, I have
    nothing whatsoever to say so will have to stop here. Let's repeat
    that a few times.}
    
  \title{The abstract of a fifth poster with a very long title to see how
    well it gets split up into lines}
  \author{A.~N.~Other\altaffil{The affiliation of A.~N.~Other is also
      very long so that it will go over several lines. His email
      address is anon@nowhere.org}}
  \listofauthors{A.~N.~Other}
  \indexauthor{Other, A. N.}
  \abstract{We will try and make this abstract longer so there is more
    stuff to put on the page and it doesn't look too empty. However, I
    have nothing whatsoever to say so will have to stop here. Let's
    repeat that a few times. We will try and make this abstract longer
    so there is more stuff to put on the page and it doesn't look too
    empty. However, I have nothing whatsoever to say so will have to
    stop here. Let's repeat that a few times. We will try and make
    this abstract longer so there is more stuff to put on the page and
    it doesn't look too empty. However, I have nothing whatsoever to
    say so will have to stop here. Let's repeat that a few times. We
    will try and make this abstract longer so there is more stuff to
    put on the page and it doesn't look too empty. However, I have
    nothing whatsoever to say so will have to stop here. Let's repeat
    that a few times.}
    
  \title{The abstract of a sixth poster with a very long title to see how
    well it gets split up into lines}
  \author{A.~N.~Other\altaffil{The affiliation of A.~N.~Other is also
      very long so that it will go over several lines. His email
      address is anon@nowhere.org}}
  \listofauthors{A.~N.~Other}
  \indexauthor{Other, A. N.}
  \abstract{We will try and make this abstract longer so there is more
    stuff to put on the page and it doesn't look too empty. However, I
    have nothing whatsoever to say so will have to stop here. Let's
    repeat that a few times. We will try and make this abstract longer
    so there is more stuff to put on the page and it doesn't look too
    empty. However, I have nothing whatsoever to say so will have to
    stop here. Let's repeat that a few times. We will try and make
    this abstract longer so there is more stuff to put on the page and
    it doesn't look too empty. However, I have nothing whatsoever to
    say so will have to stop here. Let's repeat that a few times. We
    will try and make this abstract longer so there is more stuff to
    put on the page and it doesn't look too empty. However, I have
    nothing whatsoever to say so will have to stop here. Let's repeat
    that a few times.}
    
  \title{The abstract of a seventh poster with a very long title to see how
    well it gets split up into lines}
  \author{A.~N.~Other\altaffil{The affiliation of A.~N.~Other is also
      very long so that it will go over several lines. His email
      address is anon@nowhere.org}}
  \listofauthors{A.~N.~Other}
  \indexauthor{Other, A. N.}
  \abstract{We will try and make this abstract longer so there is more
    stuff to put on the page and it doesn't look too empty. However, I
    have nothing whatsoever to say so will have to stop here. Let's
    repeat that a few times. We will try and make this abstract longer
    so there is more stuff to put on the page and it doesn't look too
    empty. However, I have nothing whatsoever to say so will have to
    stop here. Let's repeat that a few times. We will try and make
    this abstract longer so there is more stuff to put on the page and
    it doesn't look too empty. However, I have nothing whatsoever to
    say so will have to stop here. Let's repeat that a few times. We
    will try and make this abstract longer so there is more stuff to
    put on the page and it doesn't look too empty. However, I have
    nothing whatsoever to say so will have to stop here. Let's repeat
    that a few times.}
  
  \title{An abstract to go in the list of abstracts, but not in the
    index of authors}
  \author{I.~A.~M. Hidden\altaffil{Puerto Escondido}}
  \listofauthors{I.~A.~M. Hidden}
  \abstract{The idea is that by omitting the
    \textbackslash{}indexauthor command, we can avoid an entry in the
    author index.}

\end{abstracts}


%% The Index: after latexing twice, run `makeindex -s rmaa' on this
%% file, then latex twice more
\printindex

%% if you rename this file to something other than `rmbooksamp', then
%% you must change the above line accordingly

%%
%% Plates: If you have any, then the following line will print
%% them out:

\outputplates

\end{document}

%%% Local Variables:
%%% mode: latex
%%% TeX-master: t
%%% End:

\end{document}
%</ads-container>
%    \end{macrocode}

% This is the |adshtml.sty| style file used by |ADS.tex|. It just
% defines a few 
% macros that might be used in abstracts, etc (really it should define
% them all). This now gets all its code from
% Sec~\ref{sec:user-commands}. 
%    \begin{macrocode}
%<*adshtml-style>
%% Nothing else needs to go here now. 
%</adshtml-style>
%    \end{macrocode}
% The init file for latex2html.
%    \begin{macrocode}
%<*latex2html-init>
$PDFTOCAIRO = "pdftocairo";
$PDFLATEX = "pdflatex";
$LATEX = "latex";
$MAX_SPLIT_DEPTH = 0;
$HTML_VERSION = "4.0,math";
$NO_SIMPLE_MATH = 1;
$ADDRESS = "<p><address><a href=\"mailto:rmaa\@astro.unam.mx\">rmaa\@astro.unam.mx</address></p>"; 
$INFO = "<p>This document was created using latex2html and the rmaa latex macros</p>"; 
$NO_NAVIGATION = 1;
1; 
%</latex2html-init>
%    \end{macrocode}
% This is the |adshtml.perl| file that is used by \LaTeX{}2HTML. Some
% bizarre \TeX{}y command sometimes show up so we tell it to ignore
% these. All the astronomy macros that come from AAS are processed in
% \TeX{}. That way they get made to images but with an |ALT| tag that
% has the original code, which is what will be picked up by |lynx|,
% with the result that they get passed intact into the final |.txt|
% file. 
%    \begin{macrocode}
%<*adshtml-perl>
package main;

$dum = &do_cmd_makeatletter();

sub do_cmd_RMAA {
    local($_) = @_;
    "RMxAA" . $_;
}

sub do_cmd_RMAASC {
    local($_) = @_;
    "RMxAASC" . $_;
}

sub do_cmd_etal {
    local($_) = @_;
    "et al\." . $_;
}

sub do_cmd_ibidrule {
    local($_) = @_;
    "<img src=\"foo\" alt=\"\ibidrule\">" . $_;
}


sub do_cmd_bibitem {
    local($_) = @_;
    local($key);
    local($opt, $dummy)=&get_next_optional_argument; 
    $key = &missing_braces unless
        ((s/$next_pair_pr_rx/$key = $2; ''/eo)
        ||(s/$next_pair_rx/$key = $2; ''/eo));
    "<p>" . $_; 
}

&ignore_commands (<<_RAW_ARG_CMDS_);
lowercase
adjustfinalcols
textsuperscript # {}
newline # []
thinmuskip
mskip
smallskip
medskip
bigskip
hspace # {}
hspacestar # {}
vspace # {}
vspacestar # {}
mathrm
_at_mathrm
spacefactor
discretionary
unhbox 
voidb_at_x 
penalty 
_at_M
_at_m
null
IeC
_RAW_ARG_CMDS_

&process_commands_in_tex (<<_RAW_ARG_CMDS_);
sun
earth
lesssim
gtrsim
sq
argdeg
arcmin
arcsec
fd
fh
fm
fs
fdg
farcm
farcs
fp
micron
la
ga
case # {} # {}
slantfrac # {} # {}
onehalf
onethird
twothirds
onequarter
threequarters
ubvr
ub
bv
vr
ur
aj
araa
apj
apjl
apjs
ao
apss
aap
aapr
aaps
azh
baas
jrasc
memras
mnras
pra
prb
prc
prd
pre
prl
pasp
pasj
qjras
skytel
solphys
sovast
ssr
zap
nat
iaucirc
aplett
apspr
bain
fcp
gca
grl
jcp
jgr
jqsrt
memsai
nphysa
physrep
physscr
planss
procspie
_RAW_ARG_CMDS_

1;                              # This must be the last line
%</adshtml-perl> 
%    \end{macrocode}

% This is a shell script |makeads.sh| that generates the |_ads.txt| file from
% |_ads.tex| by using \LaTeX{}2HTML and |lynx|.  \changes{4.0}{2019/11/13}{Make sure that lynx writes utf-8 output instead of latin-1.}
% \changes{2.2beta13}{2002/03/12}{Complete rewrite of |makeads.sh|. Now much 
% improved I think}
% \changes{4.0}{2019/11/13}{Fix the at-sign substitution so that it only eliminates \cs{@} and does not mess with email addresses.}

%    \begin{macrocode}
%<*ads-shell-script>
#! /bin/sh
if [ -z $1 ]; then
    echo "Usage: $0 basename"
    exit
fi
BASENAME=$1
ADSBASENAME=${BASENAME}_ads
AUXFILE=$BASENAME.aux
TMPFILE=tmp.$$


echo "Generating $ADSBASENAME.tex"

# For some reason, must double-up \\ to \\\\ but not \ to \\ - odd!
cat > $ADSBASENAME.tex <<EOF
%% PLEASE DO NOT EDIT THIS FILE!  It is generated automatically from
%% the script makeads.sh.  Any hand-edited changes will be lost
%% 
\documentclass{article}
\usepackage[latin1]{inputenc}
\usepackage{html}
\usepackage{adshtml}
\newcommand{\ADSfield}[2]{\%#1  #2\\\\}
\newenvironment{ADSarticle}{}{\par\bigskip}
\setlength{\parindent}{0pt}
\newcounter{istage}
\renewcommand{\ion}[2]{\setcounter{istage}{#2}#1 \Roman{istage}}
\begin{document}
\makeatletter
EOF
echo >> $ADSBASENAME.tex "\input{$BASENAME.ads2}" 
echo >> $ADSBASENAME.tex "\end{document}"



echo "Generating $ADSBASENAME.sed" 

# Find type of journal
grep '\\documentclass\[.*proceedings.*\]' $BASENAME.tex >/dev/null \
     && JOURNAL='RMxAC' \
     || JOURNAL='RMxAA'

# Find the year and volume number
YEAR=`grep '\\SetYear' $BASENAME.tex | sed 's/.*\\SetYear{\(.*\)}.*/\1/'`
VOLUME=`grep '\\SetVolume' $BASENAME.tex | sed 's/.*\\SetVolume{\(.*\)}.*/\1/'`

# Generate an auxilliary sed script 
cat > $TMPFILE <<EOF
s|\\\\newlabel{\(.*\):RMAAFirstPage}{{.*}{\(.*\)}}|s/\\\\\\\\ThisVolume{\1}/$YEAR, RMxAC, $VOLUME, \2/g|g
EOF

# Now go through the .aux file, looking for \newlabels and use the sed 
# script just generated in order to make another sed script, one that 
# will replace \ThisVolume{LABEL} with the correct reference
grep '^\\newlabel{.*:RMAAFirstPage}' $AUXFILE \
     | sed -f $TMPFILE > $ADSBASENAME.sed && rm $TMPFILE

# Apend a further sed command to get rid of \@, since redefining 
# it in _ads.tex didn't work
echo >> $ADSBASENAME.sed 's/\\\\@//g'

echo "Generating $BASENAME.ads2 from $ADSBASENAME.sed and $BASENAME.ads" 

# Apply this to the .ads file, generating the .ads2 file, which is included
# by the _ads.tex file.
sed -f $ADSBASENAME.sed $BASENAME.ads > $BASENAME.ads2

echo "Generating $ADSBASENAME.txt from $ADSBASENAME.tex and $BASENAME.ads2"

# Now run latex2html on the _ads.tex file
latex2html -no_navigation -info 0 -address 0 -html_version 4.0,math $ADSBASENAME

# Then, convert to utf8 with lynx and tidy it up a bit
lynx --justify=off --width=100 --dump --nolist --display_charset=utf-8 \
     $ADSBASENAME/$ADSBASENAME.html | cut -c4- \
     | sed 's/ibidrule/\\ibidrule/g' > $ADSBASENAME.txt
%</ads-shell-script>
%    \end{macrocode}

% \section{Epilogue}
% Err, that's it \dots

\endinput


% \endinput
% Local Variables:
% mode: doctex
% TeX-master: t
% End:
"

rmaa.aux: rmaa.dtx rmaa.ind rmaa.gls
        @echo "REGENERATING .aux FILE"
        pdflatex "\nonstopmode% \def\fileversion{4.0} \def\filedate{2019/11/12}
% \MakeShortVerb{\|}
% \title{A \LaTeX{} document class for the Revista Mexicana de
% Astronom\'\i a y Astrof\'\i sica\thanks{This file has version
% \fileversion{} last revised \filedate{}.}}
% \author{William J. Henney\thanks{w.henney@irya.unam.mx}}
% \date{\filedate}
% \maketitle
% \paragraph{Abstract}
% This document describes the \texttt{rmaa} \LaTeX{}
% document class for the journal ``Revista Mexicana de
% Astronom\'\i a y Astrof\'\i sica''. The macros can be used for 
% \begin{enumerate}
% \item Preparing author manuscripts for the journal (including
% conference proceedings to be published in the ``Serie de
% Conferencias'')
% \item Typesetting an entire volume of the journal by inputting the
% individual article files from a container file. Facilities are
% included for automatic generation of TOC and Author Index. 
% \end{enumerate}
% 
% \tableofcontents
% \section{The documentation driver file}
% This is the part that generates the documentation that you may or
% may not be reading, which is done with the help of the Makefile
% discussed in section~\ref{sec:driver_make}. 
%    \begin{macrocode}
%<*driver>
\documentclass{ltxdoc}
\usepackage{hyperref}
\usepackage{palatino}
\EnableCrossrefs \CodelineIndex \RecordChanges
\setlength{\parindent}{0pt}
\def\MacroFont{\footnotesize\tt}
\begin{document} \DocInput{rmaa.dtx} \PrintIndex \PrintChanges \end{document}
%</driver>
%    \end{macrocode}

% \section{The implementation}
% 
% \changes{3.24}{2007/09/14}{Added user access to the macro version
% and date via \cs{rmVersion} and \cs{rmDate}}
% 
%    \begin{macrocode}
%<*package>
\newcommand{\rmaa@version}{3.26}
\newcommand{\rmaa@date}{2011/08/15}
\newcommand\rmVersion{\rmaa@version}
\newcommand\rmDate{\rmaa@date}
\NeedsTeXFormat{LaTeX2e}[1995/06/01]
%    \end{macrocode}
% \changes{4.0}{2019/11/02}{Change abstract to appear before
% resumen. Change RevMexAA to RMxAA for consistency with ADS.}
%
% \changes{3.22}{2007/03/22}{Allow proceedings articles to have
%   received/accepted dates} \changes{3.21}{2006/05/15}{Changes for
%   Silvia - abstracts are now shipped out to ADS.}
% \changes{3.19}{2006/01/19}{Changes for Silvia - remove adresses from
%   end, new abstract style} \changes{2.2beta}{2001/10/30}{ Fixed bug
%   whereby twocolumn option was not passed through to book
%   class. Finally sorted out the ``label on last page'' issue. }
% \changes{2.1}{2001/01/23}{ Fixed some inconsistencies in
%   \cs{nocite} and \cs{bibitem} uniqification Fixed the definition of
%   \cs{ion} so it will work in math mode and in contexts where the
%   typesize is different from normal (e.g., title, fig captions).
%   Modified the definition of \cs{@caption}, so we don't write any
%   list-of-figures data (we never use it anyway, and it just causes
%   problems with fragile commands in fig captions). }
% \changes{2.0}{2000/11/30}{ Added `newstyle' option, which should be
%   used in all new proceedings volumes (wider abstract, smaller
%   typeface in fig captions, etc.) Added commands for formatting
%   table captions. New demonstation document rmpostsamp.tex for
%   1-page posters. This will be the last release under a PRI
%   government (for a while...)}  \changes{1.9}{2000/06/06}{
%   Corrected position of page number in footer of articles that open
%   on even pages. Ensured that first article in each session opens on
%   an odd page, as do list of abstracts and author index. Added extra
%   space in TOC before the backmatter. Added command for editors,
%   \cs{BoldSessionTitles}, that makes produces bold titles for
%   session in the TOC.}  \changes{1.8}{2000/05/16}{ Fixed bug in
%   last page label with \cs{suppressfulladdresses} }
% \changes{1.7}{2000/05/11}{ Fixed more little things. Page labels
%   now work, even if you have no bibliography. You can output the
%   full addresses by hand with the \cs{outputfulladdresses}
%   command. Editors' Guide updated but still not finished. Revamped
%   splitbook.sh now produces PS file of each article for web page and
%   also an EPS file of each page for sending to the printers.}
% \changes{1.6}{2000/04/15}{ More cleaning up: full stops after
%   titles in index, running headers and footers in frontmatter,
%   fontsize in title of index, etc} \changes{1.5}{2000/04/15}{
%   Eliminated superscripts from titles in TOC and index}
% \changes{1.4}{2000/04/14}{ Added copyright and license info (GPL)}
% \changes{1.3}{2000/03/01}{ Corrected \cs{figcaption} to do the
%   caption formatting right Corrected font size of sec headings in
%   front matter} \changes{1.2}{2000/02/16}{ Added \cs{plate}
%   command} \changes{1.1}{2000/02/14}{ Added manuscript option,
%   corrected bugs} \changes{1.0}{2000/02/08}{ Now does a whole book
%   - lots of stuff tidied up} \changes{0.7}{1999/10/04}{ Increased
%   spacing between sections in single column styles. Suppressed
%   floats at bottom of pages (didn't work though) }
% \changes{0.6}{1999/09/24}{ Tweaked \cs{adjustfinalcols} (still not
%   perfect) Added proceedings option Added onecolumn/twocolumn
%   options (orthogonal to journal/proceedings) }
% \changes{0.5}{1999/09/23}{ Added \cs{fulladdresses} and
%   \cs{adjustfinalcols} (just about working)}
% \changes{0.4}{1999/09/22}{ Minor stylistic changes}
% \changes{0.3}{1999/06/28}{ Better mechanism for alternative
%   affiliations Added `guide' option for user guide}
% \changes{0.2}{1999/05/27}{ All basic functionality present, I
%   think} \changes{0.1}{1999/05/26}{ Just about working - still
%   missing Fig/Tab/Refs} \changes{0.0}{1999/05/21}{ So bored today,
%   decided to rewrite RMAA LaTeX macros from scratch}
%    \begin{macrocode}
\ProvidesClass{rmaa}[\rmaa@date\space v\rmaa@version\space RMxAA article
class]
%    \end{macrocode}
% \subsection{Option Handling}
% \label{sec:options}
% Declarations that are needed in options
% \changes{3.05}{2003/10/22}{Added natbib option}
% \changes{3.06}{2003/01/03}{Natbib option turned on permanently}
% \changes{3.07}{2003/01/11}{Added proof option}
% \changes{3.14}{2003/01/11}{Option to turn off natbib}
% \changes{3.18}{2003/10/22}{trying out doctex mode} 
% \changes{3.23}{2007/09/09}{Manuscript is now on by default.}
% \changes{3.26}{2011/08/15}{New option for review articles}
%    \begin{macrocode}
\newif\if@reviewarticle
\@reviewarticlefalse
\newif\if@proof
\@prooffalse
\newif\if@journal
\@journaltrue
\newif\if@printdates
\@printdatesfalse
\newif\if@proceedings
\@proceedingsfalse
\newif\if@manuscript
\@manuscripttrue
\newif\if@bigfont
\@bigfonttrue
\newif\if@guide
\@guidefalse
\newif\if@preprint
\@preprintfalse
\newif\if@rmtwocol
\@rmtwocoltrue
\newif\if@biblist
\@biblistfalse
\newif\if@book
\@bookfalse
\newif\if@natbib
\@natbibtrue
\newif\if@rmcite
\@rmcitefalse
\newif\if@titlevisible
\@titlevisibletrue
\newif\if@abstractpresent
\@abstractpresentfalse
\newif\if@printbanner
\@printbannerfalse
\newif\if@suppressfulladdresses
\@suppressfulladdressesfalse
\newif\if@newstyle
\@newstylefalse
\newif\if@debug
\@debugfalse
\newif\if@boldkeywords
\@boldkeywordsfalse
%    \end{macrocode}
% The option handling---mainly setting various booleans.
% \changes{2.2beta11}{2001/11/26}{Made twocol and
% newstyle the default for proceedings.}
% \changes{3.23}{2007/09/09}{Manuscript is now 10 pt instead of 12pt}
% \changes{3.23}{2007/09/10}{Full addresses are now automatically
% suppressed for proceedings style.}
% \changes{3.24}{2007/09/14}{New option \texttt{boldkeywords} that
% does the keywords in the old ugly way. Now off by default.}
%    \begin{macrocode}
\DeclareOption{manuscript}{\typeout{Manuscript for referee} 
  \@manuscripttrue\@journalfalse\@rmtwocolfalse
  \PassOptionsToClass{10pt}{article}}
\DeclareOption{bigfont}{\typeout{Big font manuscript requested by Gary
  Ferland} 
  \@manuscripttrue\@journalfalse\@rmtwocolfalse\@bigfonttrue
  \PassOptionsToClass{11pt}{article}}
\DeclareOption{journal}{\typeout{Main journal option}%
  \@journaltrue\@rmtwocoltrue\@printdatestrue\@manuscriptfalse
  \PassOptionsToClass{10pt}{article}}
\DeclareOption{guide}{\typeout{User guide option}%
  \@guidetrue\@journalfalse\@rmtwocoltrue\@manuscriptfalse
  \PassOptionsToClass{10pt}{article}}
\DeclareOption{preprint}{\typeout{Preprint option}%
  \@preprinttrue\@manuscriptfalse\PassOptionsToClass{10pt}{article}}
\DeclareOption{proof}{\typeout{Author proof option}%
  \@prooftrue\@manuscriptfalse}
\DeclareOption{proceedings}{\typeout{Proceedings option}%
  \@proceedingstrue\@journalfalse\@rmtwocoltrue\@newstyletrue
  \@suppressfulladdressestrue
  \@printdatesfalse\@manuscriptfalse\PassOptionsToClass{10pt}{article}} 
\DeclareOption{printdates}{\typeout{Force printing of received/accepted dates}%
  \@printdatestrue} 
\DeclareOption{book}{\typeout{Whole issue option}%
  \@booktrue} 
\DeclareOption{nonatbib}{\typeout{Turn off natbib support (required
    for compatibility with old rmaacite package).}%
  \@natbibfalse} 
\DeclareOption{newstyle}{\typeout{Use new layout style}%
  \@newstyletrue}
\DeclareOption{oldstyle}{\typeout{Use old layout style}%
  \@newstylefalse}
\DeclareOption{debug}{\typeout{Loads of annoying messages}%
  \@debugtrue}
%    \end{macrocode}
% These two need to be declared explicitly so that |\ExecuteOptions|
% actually does something.
%    \begin{macrocode}
\DeclareOption{twocolumn}{\@rmtwocoltrue}
\DeclareOption{onecolumn}{\@rmtwocolfalse}
\DeclareOption{twoside}{\PassOptionsToClass{\CurrentOption}{article}}
\DeclareOption{boldkeywords}{\@boldkeywordstrue}
%    \end{macrocode}
% Bold, upper-case keywords used to be the default. Now they must be
% turned on explicitly with the \texttt{boldkeywords} option. 
%    \begin{macrocode}
\DeclareOption{boldkeywords}{\@boldkeywordstrue}
%    \end{macrocode}
% Pass any other option along to the \texttt{article} class.
%    \begin{macrocode}
\DeclareOption*{\PassOptionsToClass{\CurrentOption}{article}}
%    \end{macrocode}
% Default options are to be two-sided and based on the standard
% \LaTeX{} \texttt{article} class,
% unless we are processing the entire book, in which case we base
% ourselves on the standard \texttt{book} class. Default is to use
% manuscript option unless otherwise specified. 
%    \begin{macrocode}
\ExecuteOptions{twoside,manuscript} \ProcessOptions\relax
\if@book
\if@rmtwocol\LoadClass[twocolumn]{book}\else\LoadClass{book}\fi
\newcommand{\abstract}{\relax}
\newcommand{\@maketitle}{\relax}
\else
\if@rmtwocol\LoadClass[twocolumn]{article}\else\LoadClass{article}\fi
\fi
%    \end{macrocode}
% Attempt to use natbib for citations 02 Jan 2003. 
% \changes{3.05}{2003/01/02}{If the natbib option is used, then set
% things up so that natbib's version of the bibliography is identical
% to the standard one.}
% \changes{3.09}{2003/01/23}{Explicitly require a recent version of
% natbib, so that we get a decent error message if the package is too
% old.} 
% \changes{3.23}{2007/09/09}{Now allow Bibtex to work properly} 
% \changes{3.25}{2009/10/01}{Change \cs{newcommand} tp \cs{def} when setting \cs{bibfont}. This allows it to work with all versions of natbib.} 
% \changes{3.25}{2010/02/06}{Updated copyright year.} 
% \changes{3.25}{2010/02/06}{Updated email and web addresses to \texttt{crya.unam.mx}.} 
%    \begin{macrocode}
\if@natbib
\RequirePackage{natbib}[1999/02/25]
\bibpunct{(}{)}{;}{a}{}{,}
\def\bibfont{\small}
\renewcommand\bibsection{\section*{REFERENCES}}
\setlength\bibsep{0pt}
\setlength\bibhang{1.5em}
\bibliographystyle{rmaa}
\fi
%    \end{macrocode}
% We want the \texttt{graphicx} package for the extended syntax of
% |\includegraphics|. We want the \texttt{ifthen} package for using in
% these macros (although we don't much yet). We explicitly load
% \texttt{babel} since some \LaTeX{} distributions messed up the
% hyphenation if we didn't.  
%    \begin{macrocode}
\RequirePackage{graphicx}
\RequirePackage{ifthen}
\IfFileExists{babel.sty}{%
%   \let\@restoreindent=\@afterindentfalse
  \RequirePackage[english]{babel}
%   \@restoreindent
  %% reinstate the definition of \@roman so as not to break .idx files
  \@ifundefined{org@@roman}{\relax}{\renewcommand{\@roman}[1]{\org@@roman{##1}}}
  }{%
  \typeout{Package babel.sty not found - don't worry about this}}
%    \end{macrocode}
% Try to load the \texttt{dcolumn} and \texttt{booktabs} packages if
% they can be found. These can produce nicer-looking tables. 
%    \begin{macrocode}
\IfFileExists{dcolumn.sty}{%
  \RequirePackage{dcolumn}
  }{%
  \typeout{Package dcolumn.sty not found - you will not be able to
    align columns on decimal points in tables}}
\IfFileExists{booktabs.sty}{%
  \RequirePackage{booktabs}
  }{%
  \typeout{Package booktabs.sty not found - your tables will not look
    as pretty as they ought too.}
  }
\IfFileExists{makeidx.sty}{%
  \RequirePackage{makeidx}
}{%
  \typeout{Package makeidx.sty not found - Author Index may not work.}
  }
%% commented out for moment since requires more modern version of LaTeX
%% \RequirePackage[latin1]{inputenc}
%    \end{macrocode}
% \changes{4.0}{2019/11/13}{Make \texttt{hyperref} and
% \texttt{etoolbox} into hard requirements, since we now use them
% extensively.}  We use hyperref to get internal and external links.
% But I don't like the colored boxes so I turn those off.  We use
% etoolbox for patching and robustifying at the moment.  In the
% future, I want to port a lot of other stuff over to it.  For
% instance, replace the ifthen package and the TeX ifs too.
%
% We target the latest version (1.7) of PDF, rather than the default
% of 1.5.  This reduces the chance of getting warnings from included
% graphics.  This is not strictly related to hyperref, but we put it
% here anyway.
%    \begin{macrocode}
\RequirePackage{etoolbox}
\RequirePackage{hyperref}
\hypersetup{hidelinks}
\pdfminorversion=7
%    \end{macrocode}
% We patch the \cs{@combinedblfloats} following a suggestion by Heiko
% Oberdiek at \url{https://tex.stackexchange.com/a/249743/3843}, which
% avoids problems when hyperlinks are split across
% pages. \changes{4.0}{2019/11/12}{Fix for issue where breaking
% hyperlink across pages gave an error.}
%    \begin{macrocode}
\urlstyle{same}
\newcommand\email{\begingroup \urlstyle{rm}\Url}
\patchcmd\@combinedblfloats{\box\@outputbox}{\unvbox\@outputbox}{}{%
  \errmessage{\noexpand\@combinedblfloats could not be patched}%
}%
%    \end{macrocode}
% \subsection{Other setup stuff}\label{sec:setup}
% Some declarations: 
%    \begin{macrocode}
\newif\if@needaddresses
\@needaddressestrue
\newif\if@frontmatter
\@frontmatterfalse
\newif\if@inpreamble
\@inpreambletrue
\newif\if@firstauthor
\@firstauthortrue
\newif\if@firstaffil
\@firstaffiltrue
\newif\if@firstaltaffil
\@firstaltaffiltrue
\newif\if@firstkwd
\@firstkwdtrue
\newif\if@noexpandADS
\@noexpandADSfalse
\newif\if@sessiontitlepages
\@sessiontitlepagesfalse
%    \end{macrocode}
% \begin{macro}{\@rmaaEmpty}
% \begin{macro}{\@rmaaArg}
% These are used in various commands to check if arg was left blank
%    \begin{macrocode}
\newcommand{\@rmaaEmpty}{}
\newcommand{\@rmaaArg}{}
%    \end{macrocode}  
% \end{macro}
% \end{macro}
%
% \subsection{Compatibility with old \LaTeX{} versions}
% \label{sec:compat}
% This package is based on the 1998/12/01 release of \LaTeX{}.
% Older releases may not define all the commands that we need.
% Luckily, we can use |\providecommand| to define these if necessary. 
% The following commands are not defined in the 1994/06/01 release. 
%
% \begin{macro}{\MakeUpperCase}
% \begin{macro}{\textsuperscript}
% \begin{macro}{\textbackslash}
% User commands:
%    \begin{macrocode}
\providecommand{\MakeUpperCase}[1]{\uppercase{#1}}
\providecommand{\textsuperscript}[1]{\ensuremath{^\textrm{#1}}}
\providecommand{\textbackslash}{\ensuremath{\backslash}}
%    \end{macrocode}  
% \end{macro}
% \end{macro}
% \end{macro}
% \begin{macro}{\protected@xdef}
% \begin{macro}{\restore@protect}
% \begin{macro}{\@unexpandable@protect}
% \begin{macro}{\hb@xt@}
% \LaTeX{} internal commands:
%    \begin{macrocode}
\providecommand{\protected@xdef}{%
        \let\@@protect\protect
        \let\protect\@unexpandable@protect
        \afterassignment\restore@protect
        \xdef}
\providecommand{\restore@protect}{\let\protect\@@protect}
\providecommand{\@unexpandable@protect}{\noexpand\protect\noexpand}
\providecommand{\hb@xt@}{\hbox to}
%    \end{macrocode}
% \end{macro}
% \end{macro}
% \end{macro}
% \end{macro}
%
% \subsection{Page layout}\label{sec:page-layout}
%
% The dimensions of the text were defined by Elizabeth. Text width is
% 169\,mm with equal margins on the two sides: 
%    \begin{macrocode}
\setlength{\textwidth}{169mm}
\setlength{\@tempdima}{\paperwidth}
\addtolength{\@tempdima}{-\textwidth}
\setlength{\oddsidemargin}{0.5\@tempdima}
\addtolength{\oddsidemargin}{-1in}
\setlength{\evensidemargin}{\oddsidemargin}
%    \end{macrocode}
% Text height is 225\,mm, top margin is 14\,mm. 
%    \begin{macrocode}
\setlength{\textheight}{225mm} 
\setlength{\topmargin}{-1in} 
\addtolength{\topmargin}{14mm}
%    \end{macrocode}
% Space between columns in two-column styles, header, footer,
% paragraph indentation:
%    \begin{macrocode}
\setlength{\columnsep}{9.5mm}
\setlength{\headsep}{0.7cm}
\setlength{\footskip}{0.75cm}
\setlength{\parindent}{1.5em}
%    \end{macrocode}
% 
% \begin{macro}{\@set@scalable@length}
% Try out a new scheme for dealing with various spacing dimensions and
% allowing the user to easily modify them. The command
% \begin{quote}
% |\set@scalable@length|\marg{name}\marg{defaultlength}\marg{stretch}\marg{shrink}
% \end{quote}
% does a 
% |\setlength|\marg{\cs{name}}\marg{defaultlength plus stretch minus shrink}
% and also defines the user command
% \textbackslash\meta{name}|Scale|\marg{factor} that can 
% be used to rescale length by \meta{factor}. Additionally, now
% defines the package command
% \textbackslash\meta{name}|@set@scale|\marg{FACTOR} 
% that also rescales the length by \meta{FACTOR}, so that the length
% becomes \meta{factor} $\times$ \meta{FACTOR} $\times$
% \meta{defaultlength}.   
%    \begin{macrocode}
\newcommand\@set@scalable@length[4]{%
  %% arg #1: NAME of length without leading `\'
  %%
  %% arg #2: default VALUE  
  \expandafter\newlength\csname @#1@default\endcsname
  \expandafter\setlength\csname @#1@default\endcsname{#2}%
  \expandafter\newlength\csname @@#1@@default\endcsname
  \expandafter\setlength\csname @@#1@@default\endcsname{#2}%
  %% arg #3: default STRETCH  
  \expandafter\newlength\csname @#1@stretch\endcsname
  \expandafter\setlength\csname @#1@stretch\endcsname{#3}%
  \expandafter\newlength\csname @@#1@@stretch\endcsname
  \expandafter\setlength\csname @@#1@@stretch\endcsname{#3}%
  %% arg #4: default SHRINK  
  \expandafter\newlength\csname @#1@shrink\endcsname
  \expandafter\setlength\csname @#1@shrink\endcsname{#4}%
  \expandafter\newlength\csname @@#1@@shrink\endcsname
  \expandafter\setlength\csname @@#1@@shrink\endcsname{#4}%
  %%
  %% User-defined scale \@NAME@scale is unity by default
  \expandafter\gdef\csname @#1@scale\endcsname{1}%
  %% Package-defined scale \@NAME@SCALE is unity by default
  \expandafter\gdef\csname @#1@SCALE\endcsname{1}%
  %%
  %% internal command \@reset@NAME should expand to 
  %% \setlength\@@NAME@@default{\@NAME@SCALE\@NAME@default}
  %% \setlength\NAME{\@NAME@scale\@@NAME@@default}
  %% in other words: scale x SCALE x default 
  \expandafter\gdef\csname @reset@#1\endcsname{%
    \expandafter\setlength\csname @@#1@@default\endcsname{%
      \csname @#1@SCALE\endcsname \csname @#1@default\endcsname
      }
    \expandafter\setlength\csname @@#1@@stretch\endcsname{%
      \csname @#1@SCALE\endcsname \csname @#1@stretch\endcsname
      }
    \expandafter\setlength\csname @@#1@@shrink\endcsname{%
      \csname @#1@SCALE\endcsname \csname @#1@shrink\endcsname
      }
    \expandafter\setlength\csname #1\endcsname{%
      \csname @#1@scale\endcsname \csname @@#1@@default\endcsname
      \@plus
      \csname @#1@scale\endcsname \csname @@#1@@stretch\endcsname
      \@minus
      \csname @#1@scale\endcsname \csname @@#1@@shrink\endcsname
      }}
  %%
  %% User command \NAMEScale should expand to 
  %% {\renewcommand\@NAME@scale{#1}\@reset@NAME}
  \expandafter\gdef\csname #1Scale\endcsname##1{%
    \expandafter\gdef\csname @#1@scale\endcsname{##1}%
    \csname @reset@#1\endcsname}
  %%
  %% Package command \NAME@set@SCALE should expand to 
  %% {\renewcommand\@NAME@SCALE{#1}\@reset@NAME}
  \expandafter\gdef\csname #1@set@SCALE\endcsname##1{%
    \expandafter\gdef\csname @#1@SCALE\endcsname{##1}%
    \csname @reset@#1\endcsname}
  %%
  %% initialize at default values with \@reset@NAME
  \csname @reset@#1\endcsname
  }
%    \end{macrocode}
% \end{macro}
% \begin{macro}{\parskip}
% First use of this new infrastructure: set interparagraph space to be
% zero but with a small amount of stretchiness.
%    \begin{macrocode}
%   \ifthenelse{\boolean{@manuscript}}{%
%     \@set@scalable@length{parskip}{0.5\baselineskip}{0.3\p@}{0.1\p@}
%   }{
    \@set@scalable@length{parskip}{0\p@}{0.3\p@}{0\p@}
%   }
%    \end{macrocode}
% \end{macro}
% 
% \begin{macro}{\abovesectionskip}
% \changes{4.0}{2019/11/12}{Use negative values for skip so that first paragraph after section headings is not indented.  This is something that I have been trying to fix for years!  It turns out that the anser was in the documentation for the standard command \cs{@startsection.}}
% \begin{macro}{\belowsectionskip}
% \begin{macro}{\aboveparaskip}
% \begin{macro}{\belowparaskip}
% Define some utility rubber lengths for use in
% section~\ref{sec:section-commands}. 
%    \begin{macrocode}
\newlength{\abovesectionskip}
\newlength{\belowsectionskip}
\newlength{\aboveparaskip}
\newlength{\belowparaskip}
%    \end{macrocode}
% The default values for these lengths are those traditionally used
% for the onecolumn style. 
%    \begin{macrocode}
\@set@scalable@length{abovesectionskip}{-3.5ex}{-4ex}{-2ex}
\@set@scalable@length{belowsectionskip}{2.3ex}{0ex}{2ex}
\@set@scalable@length{aboveparaskip}{3.25ex}{2ex}{0.5ex}
\@set@scalable@length{belowparaskip}{-1em}{\z@}{\z@}
%    \end{macrocode}
% \end{macro}\end{macro}\end{macro}\end{macro}
% We also make scalable some of the standard float parameters. These
% are the same defaults as in the standard article classes, but
% slightly more stretchy. For the twocolumn style, they are currently
% halved, see below. Note that I have also lumped |\abovecaptionskip|
% in with these, since it plays a similar role. There is also a
% |\belowcaptionskip| but it is set to zero in the standard classes. 
%    \begin{macrocode}
\@set@scalable@length{floatsep}{12\p@}{4\p@}{4\p@}
\@set@scalable@length{textfloatsep}{20\p@}{4\p@}{8\p@}
\@set@scalable@length{intextsep}{12\p@}{4\p@}{4\p@}
\@set@scalable@length{dblfloatsep}{12\p@}{4\p@}{4\p@}
\@set@scalable@length{dbltextfloatsep}{20\p@}{4\p@}{8\p@}
\@set@scalable@length{abovecaptionskip}{10\p@}{\z@}{\z@}
%    \end{macrocode}
% \begin{macro}{\@before@title@skip}         
% \begin{macro}{\@before@author@skip}        
% \begin{macro}{\@between@author@skip}       
% \begin{macro}{\@before@abstract@skip}      
% \begin{macro}{\@after@abstracthead@skip}   
% \begin{macro}{\@between@abstract@skip}     
% \begin{macro}{\@before@kwd@skip}           
% \begin{macro}{\@before@dates@skip}           
% Spacing parameters used in |\@maketitle|. These are not changed by
% default except by the user, so they don't need to go into
% |\@package@default@lengths|.  
%    \begin{macrocode}
\newlength{\@before@title@skip}         
\newlength{\@before@author@skip}        
\newlength{\@between@author@skip}       
\newlength{\@before@abstract@skip}      
\newlength{\@after@abstracthead@skip}   
\newlength{\@between@abstract@skip}     
\newlength{\@before@kwd@skip}           
\newlength{\@before@dates@skip}           
\@set@scalable@length{@before@title@skip}{2em}{\z@}{\z@}
\@set@scalable@length{@before@author@skip}{1.5em}{\z@}{\z@}
\@set@scalable@length{@between@author@skip}{0.5em}{\z@}{\z@}
\@set@scalable@length{@before@abstract@skip}{3ex}{1ex}{0.5ex}
\@set@scalable@length{@after@abstracthead@skip}{1ex}{\z@}{\z@}
\@set@scalable@length{@between@abstract@skip}{2ex}{1ex}{0.5ex}
\@set@scalable@length{@before@kwd@skip}{1ex}{0.5ex}{0.2ex}
\@set@scalable@length{@before@dates@skip}{2ex}{0.5ex}{0.2ex}
%    \end{macrocode}
% \end{macro}\end{macro}\end{macro}\end{macro}
% \end{macro}\end{macro}\end{macro}\end{macro}
% 
% Spacing above and below displayed equations also wants to be done
% like this, but it is more complicated since the standard article
% classes make these change with the fontsize. Something to do later,
% maybe. 
% 
%    \begin{macrocode}
%    \end{macrocode}
% \begin{macro}{\@package@default@lengths}
% The value actually used for these lengths is controlled by the
% \meta{SCALE} parameter, which depends on the style option that is in
% effect. This is set using the |\@package@default@lengths| command,
% which must be called every time the style changes (e.g., changing
% from onecolumn to twocolumn). 
%    \begin{macrocode}
\newcommand{\@package@default@lengths}{%
  \ifthenelse{\boolean{@rmtwocol}}{%
    %% twocolumn mode
    \abovesectionskip@set@SCALE{0.5}
    \belowsectionskip@set@SCALE{0.5}
    \aboveparaskip@set@SCALE{0.5}
    \belowparaskip@set@SCALE{1}
    \floatsep@set@SCALE{0.5}
    \textfloatsep@set@SCALE{0.5}
    \intextsep@set@SCALE{0.5}
    \dblfloatsep@set@SCALE{0.5}
    \dbltextfloatsep@set@SCALE{0.5}
    }{%
    %% onecolumn mode
    \abovesectionskip@set@SCALE{1}
    \belowsectionskip@set@SCALE{1}
    \aboveparaskip@set@SCALE{1}
    \belowparaskip@set@SCALE{1}
    \floatsep@set@SCALE{1}
    \textfloatsep@set@SCALE{1}
    \intextsep@set@SCALE{1}
    \dblfloatsep@set@SCALE{1}
    \dbltextfloatsep@set@SCALE{1}
    }%
  }
\@package@default@lengths
%    \end{macrocode}
% \end{macro}
% \begin{macro}{\RescaleLengths}
% \begin{macro}{\RescaleSecLengths}
% \begin{macro}{\RescaleFloatLengths}
% \begin{macro}{\RescaleTitleLengths}
% User commands that will fiddle with all the scalable lengths. 
%    \begin{macrocode}
\newcommand{\RescaleLengths}[1]{%
  \RescaleSecLengths{#1}\RescaleFloatLengths{#1}\RescaleTitleLengths{#1}}
\newcommand{\RescaleSecLengths}[1]{%
%% sections
  \abovesectionskipScale{#1}
  \belowsectionskipScale{#1}
  \aboveparaskipScale{#1}
  }
\newcommand{\RescaleFloatLengths}[1]{%
%% floats
  \floatsepScale{#1}
  \textfloatsepScale{#1}
  \intextsepScale{#1}
  \dblfloatsepScale{#1}
  \dbltextfloatsepScale{#1}
  \abovecaptionskipScale{#1}
  }
\newcommand{\RescaleTitleLengths}[1]{%
%% maketitle
  \@before@title@skipScale{#1}           
  \@before@author@skipScale{#1}        
  \@between@author@skipScale{#1}       
  \@before@abstract@skipScale{#1}      
  \@after@abstracthead@skipScale{#1}   
  \@between@abstract@skipScale{#1}     
  \@before@kwd@skipScale{#1}           
  \@before@dates@skipScale{#1}           
  }
%    \end{macrocode}
% \end{macro}\end{macro}\end{macro}\end{macro}
% \changes{2.2beta14}{2002/03/17}{Finished implementation of new
% infrastructure for scalable lengths. Converted the above/below
% section skips to use it. Also, all the floatsep parameters. Also,
% all the vertical lengths that are used by |\@maketitle|. }

% \begin{macro}{\MSextramargin}
% For the \texttt{manuscript} option we have an extra 0.3\,inch
% vertical margin and 1\,inch horizontal margin for marginal notes (by
% hand or with |\marginpar|. 
%
% \changes{3.23}{2007/09/09}{Reduce baselinestretch in
%   \texttt{manuscript} style to 1.2. No longer use raggedright for
%   manuscript.}
%
% \changes{3.24}{2007/09/14}{Further reduction of baselinestretch in
%   \texttt{manuscript} style to 1.1 since we now have wider margins.}
%
% \changes{3.24}{2007/09/14}{Split \texttt{manuscript}'s extra margin
%   into separate horizontal and vertical parts.}
% 
%    \begin{macrocode}
\newlength{\MSextraHmargin}\newlength{\MSextraVmargin}
\setlength{\MSextraHmargin}{1in} 
\setlength{\MSextraVmargin}{0.3in} 
\if@manuscript
   \renewcommand{\baselinestretch}{1.1}
   \addtolength{\textwidth}{-2\MSextraHmargin}
   \addtolength{\oddsidemargin}{\MSextraHmargin}
   \setlength{\evensidemargin}{\oddsidemargin}
   \addtolength{\textheight}{-2\MSextraVmargin}
   \addtolength{\topmargin}{\MSextraVmargin}
%    \raggedright
\fi
%    \end{macrocode}
% \end{macro}
% Float parameters: don't allow bottom floats by default. 
%    \begin{macrocode}
\renewcommand{\fps@figure}{tp}
\renewcommand{\fps@table}{tp}
%    \end{macrocode}
%
% \begin{macro}{\rmaa@AuthorWidthFrac}
% \begin{macro}{\rmaa@AbstractWidthFrac}
% These parameters control what fraction of the total text width can
% be used by the list of authors and the abstract/resumen,
% respectively.  
%    \begin{macrocode}
\newcommand{\rmaa@AuthorWidthFrac}{1.0}
\newcommand{\rmaa@AbstractWidthFrac}{0.75}
%    \end{macrocode}
% Use a wide abstract to save space in the new style.
%    \begin{macrocode}
\if@newstyle\renewcommand{\rmaa@AbstractWidthFrac}{1.0}\fi
%    \end{macrocode}
% \end{macro}
% \end{macro}
%
% Some things to consider in future versions. Trying to avoid too many
% widows, etc. Apparently, the following will work best with
% \cs{flushbottom}. 
%    \begin{macrocode}
%%\widowpenalty=10000
%%\clubpenalty=10000
%    \end{macrocode}
% Also, should perhaps try getting rid of underfull/overfull box warnings. 
% \changes{3.10}{2003/01/24}{Set large values for \cs{hbadness},
% \cs{vbadness}, \cs{hfuzz}, \cs{vfuzz}. To reduce amount of stuff
% that gets written to the terminal every time you compile something.}
%    \begin{macrocode}
\hbadness=10000
\vbadness=10000
\setlength\hfuzz{\maxdimen}
\setlength\vfuzz{\maxdimen}
%    \end{macrocode}
% \subsection{Various text styles for different contexts}
% \label{sec:text-styles} 
% \begin{macro}{\rmaa@AffilStyle}
% \begin{macro}{\rmaa@PagesStyle}
% \begin{macro}{\rmaa@YearStyle}
% \begin{macro}{\rmaa@VolumeStyle}
% \begin{macro}{\rmaa@VolumenStyle}
% \begin{macro}{\rmaa@JournalNameStyle}
% \begin{macro}{\rmaa@ReceivedStyle}
% \begin{macro}{\rmaa@AbstractHeaderStyle}
% \begin{macro}{\rmaa@JournalNameStyle}
% Style commands for the different fields used by |\maketitle| in
% constructing the title page of the article. These have the form
% |\rmaa@XYZStyle{TEXT}|. I'm not sure these were really a good idea,
% but never mind \dots 
%    \begin{macrocode}
\newcommand{\rmaa@AffilStyle}[1]{{\normalsize\rmfamily#1}}
\newcommand{\rmaa@PagesStyle}[1]{{\normalsize\rmfamily#1}}
\newcommand{\rmaa@YearStyle}[1]{{\normalsize\rmfamily#1}}
\newcommand{\rmaa@VolumeStyle}[1]{{\normalsize\rmfamily\bfseries#1}}
\newcommand{\rmaa@VolumenStyle}[1]{{\Large\sffamily\bfseries#1}}
\newcommand{\rmaa@JournalNameStyle}[1]{{\normalsize\itshape#1}}
\newcommand{\rmaa@ReceivedStyle}[1]{{\normalsize\rmfamily\itshape#1}}
\newcommand{\rmaa@AbstractHeaderStyle}[1]{%
  {\normalsize\rmfamily\MakeUpperCase{#1}}}
\newcommand{\rmaa@TitleStyle}[1]{%
  {\large\bfseries\rmfamily\boldmath\MakeUpperCase{#1}\par}}
%    \end{macrocode}
% \end{macro}
% \end{macro}
% \end{macro}
% \end{macro}
% \end{macro}
% \end{macro}
% \end{macro}
% \end{macro}
% \end{macro}
%
% \begin{macro}{\rmaa@FrontSecStyle}
% Headings of ``front sections'' (see section~\ref{sec:front-sections})
%    \begin{macrocode}
\newcommand{\rmaa@FrontSecStyle}[1]{%
  {\normalsize\bfseries\rmfamily\MakeUpperCase{#1}}}
%    \end{macrocode}
% \end{macro}
%
% \begin{macro}{\rmaa@TOCTitleStyle}
% \begin{macro}{\rmaa@TOCPartStyle}
% \begin{macro}{\rmaa@IndexTitleStyle}
% \begin{macro}{\rmaaTOCAuthorStyle}
% \begin{macro}{\rmaaIndexAuthorStyle}
% Commands to set style of TOC and index entries.
% Rewritten to avoid unnecessary font commands in the .toc and .ind files
%    \begin{macrocode}
\newcommand{\rmaa@TOCTitleStyle}[1]{\MakeUpperCase{#1}}
\newcommand{\rmaa@TOCPartStyle}[1]{\rmaa@TOCTitleStyle{#1}}
\newcommand{\rmaa@IndexTitleStyle}[1]{#1}
\newcommand{\rmaaTOCAuthorStyle}[1]{%
  {\itshape#1}}
\newcommand{\rmaaIndexAuthorStyle}[1]{%
  \protect\mbox{\bfseries#1}}
%    \end{macrocode}
% \end{macro}
% \end{macro}
% \end{macro}
% \end{macro}
% \end{macro}
%
% \begin{macro}{\rmaa@KeywordStyle}
% \begin{macro}{\rmaa@KeywordHeaderStyle}
% The style of the keywords depends on whether we are using the
% \texttt{newstyle} option or not. If we are, then they are small. 
%    \begin{macrocode}
\ifthenelse{\boolean{@newstyle}}{
  \newcommand{\rmaa@KeywordStyle}[1]{%
    {\small%
      \if@boldkeywords 
      \bfseries\MakeUpperCase{\raggedright #1}
      \else
      \raggedright #1
      \fi
    }}
  \newcommand{\rmaa@KeywordHeaderStyle}[1]{%
    {\small\itshape#1}}}
{
  \newcommand{\rmaa@KeywordStyle}[1]{%
    {\normalsize%
      \if@boldkeywords\bfseries\fi\raggedright #1}}
  \newcommand{\rmaa@KeywordHeaderStyle}[1]{%
    {\normalsize\itshape#1}}}
%    \end{macrocode}
% \end{macro}
% \end{macro}
%
% \subsection{Page styles}\label{sec:page-styles}
% \begin{macro}{\ps@rmaa}
% The running headers for normal pages:
%    \begin{macrocode}
\newcommand{\ps@rmaa}{%
%    \end{macrocode}
% nothing in either foot; 
%    \begin{macrocode}
  \renewcommand\@oddfoot{\relax}%
  \renewcommand\@evenfoot{\relax}%
%    \end{macrocode}
% page number and short author list in the even heads; 
%    \begin{macrocode}
  \renewcommand\@evenhead{\rlap{\thepage}\hfil\@shortauthor\hfil}%
%    \end{macrocode}
% short title and page number in the odd heads.
%    \begin{macrocode}
  \renewcommand\@oddhead{\hfil\@shorttitle\hfil\llap{\thepage}}
  }
%    \end{macrocode}
% \begin{macro}{\ps@rmaafirstpage}
% The header of the first page of each
% article:
%    \begin{macrocode}
\newcommand{\ps@rmaafirstpage}{%
  \renewcommand\@oddfoot{\hfil\thepage}%
%    \end{macrocode}
% footer corrected (06 Jun 2000) to have page on outside edge;
%    \begin{macrocode}
  \renewcommand\@evenfoot{\thepage\hfil}%
%    \end{macrocode}
% Header has journal name, volume, and page numbers for a regular
% article, or similar for special cases. 
% \changes{4.0}{2019/11/03}{Change to three-line header with copyright info and DOI.}
%    \begin{macrocode}
\if@preprint
\renewcommand\@oddhead{%
  \if@proceedings 
  \textsc{To appear in ``\@conftitle{} (\rmaa@YearStyle{\@year})''}\hfil
  \textit{RMxAC}
  \else
  \textsc{To appear in} \@journalname\hfil
  \fi}\fi%
\if@guide
\renewcommand\@oddhead{\textsc{Documentation} (\rmaa@YearStyle{\@year})\hfil}
\fi%
\if@manuscript
\renewcommand\@oddhead{Manuscript for \@journalname{} (\rmaa@YearStyle{\@year})\hfil}
\fi%
\if@journal
\renewcommand\@oddhead{%
  \parbox{\textwidth}{
    \@journalname, \rmaa@VolumeStyle{\@volume},
    \rmaa@PagesStyle{\@artpages}\
    (\rmaa@YearStyle{\@year})\hfil\\[1pt]
    {\small\textcopyright{} \@year:
      Instituto de Astronom\'\i{}a,
      Universidad Nacional Aut\'onoma de M\'exico}\hfil\\[1pt]
    {\small\@DOI@url}\hfil
  }%
}\fi%
  \renewcommand\@evenhead{\@oddhead}}
%    \end{macrocode}
% By default, use the |\ps@rmaa| page style.
%    \begin{macrocode}
\pagestyle{rmaa}
%    \end{macrocode}
% \end{macro}
% \end{macro}
%
% \subsection{Preamble commands}\label{sec:preamble} 
% Everything you need for setting the title page, TOC and index
% entries, etc.
% \subsubsection{Commands for the title page}
% \label{sec:title-commands}
% \begin{macro}{\resumen}
% \begin{macro}{\@resumen}
% \begin{macro}{\abstract}
% \begin{macro}{\@abstract}
% User commands for setting the resumen and abstract. 
% These should be used in the document header. The internal commands
% |\@resumen| and |\@abstract| are used to hold the actual text
% pending output via the |\maketitle| command. 
%    \begin{macrocode}
\newif\if@setresumen
\newcommand{\resumen}[1]{%
%%  \renewcommand{\@resumen}{{\selectlanguage{spanish}#1}}%
  \renewcommand{\@resumen}{{\ignorespaces#1}}%
  \@setresumentrue}
\renewcommand{\abstract}[1]{%
  \renewcommand{\@abstract}{{\ignorespaces#1}}
  \renewcommand\ADS@abstract{#1}
%% Not a good idea, since then author macros are not expanded and
%% cause errors later. However, by not doing it we are vulnerable to
%% fragile commands.....
  \if@noexpandADS\@onelevel@sanitize{\ADS@abstract}\fi
  \@abstractpresenttrue
  }
\newcommand{\ADSnoexpand}{\@noexpandADStrue}
%    \end{macrocode}
% \begin{macro}{\boldabstract}
% This is the version for doing 1-page conference proceedings. Still
% save the text for ADS but then write it out immediately as a bold
% paragraph. 
%    \begin{macrocode}
\newcommand{\boldabstract}[1]{%
  \gdef\ADS@abstract{#1}%
  \renewcommand{\@abstract}{#1}%
  \if@noexpandADS\@onelevel@sanitize{\ADS@abstract}\fi%
  \noindent\textbf{\boldmath #1}
  }
%    \end{macrocode}
% \end{macro}
% Default value in each case is to print out a warning if not set
%    \begin{macrocode}
  \newcommand{\@resumen}{%
    \@setresumenfalse \ClassWarningNoLine{rmaa}{No Spanish abstract
      (resumen) specified - if you do not speak Spanish, this will be
      translated by the editors}%
    \textbf{Favor de proporcionar un resumen en espa\~nol. If you are
      unable to translate your abstract into Spanish, the editors will
      do it for you.} }
  \newcommand{\@abstract}{\ClassWarningNoLine{rmaa}{No abstract
      specified}}
%    \end{macrocode}
% \end{macro}
% \end{macro}
% \end{macro}
% \end{macro}
%
% \begin{macro}{\title}
% \begin{macro}{\@rawtitle}
% \begin{macro}{\@TOCtitle}
% \begin{macro}{\rmaaIndexTitle}
% Similarly, |\title| is the preamble command to set the title of the
% article. This is used in various places, typeset in different
% styles, hence the use of the internal |\@rawtitle|, |\@title|,
% |\rmaaIndexTitle|, and |\@TOCtitle| commands, each of which uses the
% relevant style from section~\ref{sec:text-styles}
%    \begin{macrocode}
\newcommand{\@rawtitle}{Raw Title}
\newcommand{\@TOCtitle}{\rmaa@TOCTitleStyle{Title for TOC}}
\newcommand{\rmaaIndexTitle}{\rmaa@IndexTitleStyle{Title for index}.}
\newcommand{\@TOCaftertitle}{}
\renewcommand{\title}[1]{%
%%  {\ADS@sanitize\xdef\@rawtitle{#1}}
  \renewcommand\@rawtitle{\noexpand#1}
  \let\ADS@title=\@rawtitle
  \renewcommand{\@title}{\rmaa@TitleStyle{#1}}%
  \renewcommand{\@TOCtitle}{\rmaa@TOCTitleStyle{#1}\@TOCaftertitle}%
  \renewcommand{\rmaaIndexTitle}{\rmaa@IndexTitleStyle{#1\@TOCaftertitle}.}%
  }
%    \end{macrocode}
% \end{macro}
% \end{macro}
% \end{macro}
% \end{macro}
% 
% All the weird and whacky ways of indicating the author
% affiliations. 
% \begin{macro}{\affil}
% First, the way that sets it directly underneath the authors name.
% Note that |\affil| trumps all other affiliations for the purposes of
% getting into the ADS stuff (we use |\gdef|).   
% \dots 
%    \begin{macrocode}
\newcommand{\affil}[1]{%
  \\[1.5ex]\rmaa@AffilStyle{#1}\\[1.5ex]%
  \if@firstaffil\gdef\ADS@affil{#1}\@firstaffilfalse\fi%
  }
\renewcommand{\and}{and\\[1.5ex]}
%    \end{macrocode}
% \begin{macro}{\altaffil}
% Alternative affiliations are footnoted. 
%    \begin{macrocode}
\newcommand{\altaffil}[1]{%
  \if@firstaltaffil\gdef\ADS@affil{#1}\@firstaltaffilfalse\fi%
  \thanks{#1}}
%    \end{macrocode}
% \begin{macro}{\altaffilmark}
% \begin{macro}{\altaffiltext}
% Alternative-alternative mechanism for affiliations---should not be
% mixed with |\altaffil| (unless you know what you are doing) or there
% will be problems with numbering.
%    \begin{macrocode}
\newcommand{\altaffilmark}[1]{\textsuperscript{#1}}
\newcommand{\altaffiltext}[2]{%
  \if@firstaltaffil\gdef\ADS@affil{#2}\@firstaltaffilfalse\fi
  \stepcounter{footnote}%
  \protected@xdef\@thanks{\@thanks
    \protect\footnotetext[#1]{#2}}} 
%    \end{macrocode}
% \end{macro}
% \end{macro}
% \end{macro}
% \end{macro}
%
% \begin{macro}{\keywords}
% \begin{macro}{\@keywords}
% \begin{macro}{\@keywordheader}
% \label{keywords}Sundry minor preamble commands. All
% self-explanatory, I hope. |\keywords| is the old user interface
% (deprecated). New articles should use |\addkeyword| once for each
% keyword. That way the ADS record comes out right too (comma rather
% than em-dash separator).
% \changes{2.2beta12}{2001/12/07}{\cs{addkeyword} New user interface
%   for keywords.} 
% \changes{2.2beta13}{2002/03/11}{\cs{addkeyword} Improved behaviour
%   with .ads file. Use
%   \cs{ADS@sanitize} to change definitions of `~' and \cs{ion}.}
% \changes{3.23}{2007/09/09}{\cs{addkeyword} now signals a compilation
% error if the author tries to add more than 6 keywords.}
% \changes{3.24}{2007/09/14}{kwdcounter is now reset in
% \cs{@begofarticlehook} to prevent error in book style} 
%    \begin{macrocode}
\newcommand{\@kwdsep}{ --- }
\newcommand{\ADS@kwdsep}{, }
\newcounter{kwdcounter}\setcounter{kwdcounter}{0}
\newcommand{\addkeyword}[1]{%
  \addtocounter{kwdcounter}{1}
  \ifthenelse{\thekwdcounter > 6}{
    \ClassError{rmaa}{Too many keywords. A maximum of 6 keywords is allowed.}{}
  }{}
  \if@firstkwd%
  \renewcommand\@keywords{\if@boldkeywords\uppercase{#1}\else #1\fi}
  {\ADS@sanitize\def\ion{\plainion}\gdef\ADS@keywords{#1}}
  \@firstkwdfalse
  \else%
  \edef\@keywords{\@keywords\@kwdsep \if@boldkeywords\uppercase{#1}\else #1\fi}%
  {\ADS@sanitize\def\ion{\plainion}\xdef\ADS@keywords{\ADS@keywords\ADS@kwdsep #1}}%
  \fi}
\newcommand{\ADS@sanitize}{%
  \def\mathrm{}
  \def\@{\ }
  \def~{\ }
  \def\langle{<}
  \def\rangle{>}
  \def\log{log }
  \def\sin{sin }
  \def\cos{cos }
  \def\tan{tan }
  \def\fd{^d.}
  \def\fh{^h.}
  \def\fm{^m.}
  \def\fs{^s.}
  \def\fdg{�.}
  \def\farcm{'.}
  \def\farcs{''.}
  \def\arcdeg{�}
  \def\,{ }
  }
\newcommand{\keywords}[1]{%
  \renewcommand{\ADS@keywords}{#1}
  \renewcommand{\@keywords}{\uppercase{#1}}
  }
\newcommand{\@keywordheader}{\rmaa@KeywordHeaderStyle{Key Words:}}
\newcommand{\@keywords}{%
  \ClassWarningNoLine{rmaa}{No keywords specified}%
  Keywords go here
  }
%    \end{macrocode}
% \begin{macro}{\shorttitle}
% \begin{macro}{\shortauthor}
% \begin{macro}{\@shorttitle}
% \begin{macro}{\@shortauthor}
% Commands for setting the running headers.
%    \begin{macrocode}
\newcommand{\shorttitle}[1]{%
  \renewcommand\@shorttitle{\MakeUpperCase{#1}}}
\newcommand{\shortauthor}[1]{%
  \renewcommand\@shortauthor{\MakeUpperCase{#1}}}
\newcommand{\@shorttitle}{SHORT TITLE GOES HERE}
\newcommand{\@shortauthor}{SHORT AUTHOR GOES HERE}
%    \end{macrocode}
% \end{macro}
% \end{macro}
% \end{macro}
% \end{macro}
% \end{macro}
% \end{macro}
% \end{macro}


% \subsubsection{The full addresses}
% \label{sec:full-addresses}
% \begin{macro}{\fulladdresses}
% \begin{macro}{\@fulladdresses}
% \begin{macro}{\@fulladdheight}
% \begin{macro}{\@fulladddepth}
% \begin{macro}{\@fulladdbox}
% Aaghh, the full addresses! A perennial source of problems. 
%    \begin{macrocode}
\newlength{\@fulladdheight}
\newlength{\@fulladddepth}
\newsavebox{\@fulladdbox}
%% \fulladdresses command modified (11 May 2000) so that we can put
%% \fulladdresses{} in the article to suppress the output of the
%% addresses
\newcommand{\@FinalFigure}{\relax}
\newcommand{\FinalFigure}[1]{\renewcommand{\@FinalFigure}{#1}}
\newcommand{\fulladdresses}[1]{%
  \renewcommand{\@rmaaArg}{#1}
  \ifx\@rmaaArg\@rmaaEmpty
     \renewcommand{\@fulladdresses}{}
  \else
     \renewcommand{\@fulladdresses}{
       \normalsize%
       \ifx\@FinalFigure\relax\else\vfill\@FinalFigure\vfill\fi
%%       \vspace*{1ex}
       %% do addresses as a list (copied from thebibliography environment) 
       \list{\null}{\leftmargin 1.5em\topsep\z@\labelwidth\z@\labelsep\z@\itemindent -1.5em
         \itemsep\z@\parsep\z@\usecounter{enumi}}
       \sloppy\clubpenalty4000\widowpenalty4000
       \sfcode`\.=1000\relax
       #1
       \endlist
       }
  \fi
  }
\newcommand{\@fulladdresses}{Full addresses go here}
%    \end{macrocode}
% \begin{macro}{\suppressfulladdresses}
% New user-level command for suppressing full addresses should you so
% desire. Currently used in the single-page-poster style. 
%    \begin{macrocode}
\newcommand{\suppressfulladdresses}{\@suppressfulladdressestrue\@needaddressesfalse}
%    \end{macrocode}
% \end{macro}
% \end{macro}
% \end{macro}
% \end{macro}
% \end{macro}
% \end{macro}

% \begin{macro}{\printdates}
% \begin{macro}{\noprintdates}
% Command to activate the printing out of the received/accepted
% dates. This can also be done using the \texttt{printdates} option,
% which is the default for \texttt{journal}, but not for \texttt{proceedings}.
%    \begin{macrocode}
\newcommand{\printdates}{\@printdatestrue}
\newcommand{\noprintdates}{\@printdatesfalse}
%    \end{macrocode}
% \end{macro}
% \end{macro}
% \begin{macro}{\@outputaddressesifnecessary}
% Command to output the addresses if we haven't already done so
% This is called by: 
% \begin{enumerate}
% \item The end of |thebibliography| environment
% \item The |\outputfulladdresses| command
% \end{enumerate}
%    \begin{macrocode}
\newcommand{\@outputaddressesifnecessary}{%
  \ClassInfo{rmaa}{Calling @outputaddressesifnecessary}
  \ifthenelse{\boolean{@needaddresses}}{%
    %% If in one column style, or in the first column of a two column
    %% style then we can try to put it on this page 
    \ifthenelse{\boolean{@firstcolumn} \or \not\boolean{@rmtwocol}}{%
      \@outputaddresses}{%
      %% otherwise give warning it will be on a separate page
      \ClassWarning{rmaa}{Full addresses will go on a separate
        page (please see User Guide for how to remedy this)}}%
    %% If we haven't managed yet, put it on its own page
    \ifthenelse{\boolean{@needaddresses}}{%
      \if@debug{\ClassWarning{rmaa}{In this silly bit!}}\fi
      \clearpage\vfill\@outputaddresses}{}%
    }{
    %% If we don't need them, we still might need to do the eoa hook 
    \ifthenelse{\boolean{@suppressfulladdresses}}{%
      \if@debug{\ClassWarning{rmaa}{XXX}}\fi
%%      \if@book\else\@endofarticlehook\null\fi
      }{}
    }%
  }  
%    \end{macrocode}
% \end{macro}

% \begin{macro}{\outputfulladdresses}
% A user interface to the above command, in case there is no biblio
% (11 May 2000) 
%    \begin{macrocode}
\newcommand{\outputfulladdresses}{\@outputaddressesifnecessary}
%    \end{macrocode}
% \end{macro}

% \begin{macro}{\adjustfinalcols}
% This command starts a new column and should be 
% used to ensure that the columns balance on the last page. It also
% outputs the full 
% author addresses to ensure that this occurs in the left hand
% column. 
% The |\par| command is so that the current line will be taken account
% of when working out |\pagetotal|
%    \begin{macrocode}
\newcommand{\adjustfinalcols}{\if@twocolumn\goodbreak\@outputaddresses\newpage\noindent\fi}
%    \end{macrocode}
% \end{macro}

% \begin{macro}{\@outputaddresses}
% Command to do the work of outputting the addresses. 
% This can be called by 
% \begin{enumerate}
% \item an explicit |\adjustfinalcols|
% \item via |\@outputaddressesifnecessary|, q.v.
% \end{enumerate}
%
%    \begin{macrocode}
\newcommand{\@outputaddresses}{%
  % Allow option of suppressing the full addresses, in which case
  % we just run the hook if we are not in a book style
  \if@suppressfulladdresses
%%     \if@book\else\@endofarticlehook\null\fi
  \else
  %% save the full addresses in a box for output later
  \sbox{\@fulladdbox}{\parbox{\textwidth}{\@fulladdresses}}
  % calculate height of address box  
  \settoheight{\@fulladdheight}{\usebox{\@fulladdbox}}
  \settodepth{\@fulladddepth}{\usebox{\@fulladdbox}}
  \addtolength{\@fulladdheight}{\@fulladddepth}
  \if@debug
  { \typeout{Address total height: \the\@fulladdheight}
    \typeout{Text height: \the\pagetotal}
    \typeout{Goal height: \the\pagegoal}
    }
  \fi
  % reuse \@fulladdheight to mean the total height of text+addresses
  \addtolength{\@fulladdheight}{\pagetotal}
%%  \addtolength{\@fulladdheight}{-\baselineskip}
  \if@debug\typeout{Adjusted text + addresses height: \the\@fulladdheight}\fi
  % check it will fit on current page
  \ifdim\@fulladdheight<\pagegoal
     %% must shift it to the left
     %% to compensate for the hanging indentation in the reference list. 
     %% NOTE: this is not foolproof!
     %% NB. This will give overfull hbox warnings, which is a bit ugly.
     \vfill\noindent\if@biblist\hspace*{-\leftmargin}\fi
     \usebox{\@fulladdbox}
     % Set this to show we've done it
     \@needaddressesfalse
     % and assume we are at the end of the article
     % so that we do the EOA hook if not in the book style
%%     \if@book\else\@endofarticlehook\fi
  \else
     % if not, just do nothing. With luck, this function will get 
     % called again on the next page
     \if@debug{\typeout{Holding over full addresses to next page}}\fi
  \fi\fi
  }
%    \end{macrocode}
% \end{macro}

% \subsubsection{Commands for editors}
% \label{sec:editor-commands}
% Some user/editor commands for doing the TOC and index in the book
% option. 
% They are defined here rather in the book-specific part since they
% need to be always available in order to avoid breaking documents. 
% \begin{macro}{\listofauthors}
% \begin{macro}{\@rawauthors}
% \begin{macro}{\rmaaTOCauthors}
% The command |\listofauthors| formats an author list for use in the
% TOC. It should contain the same names as in the |\author| command,
% but without the affiliation garbage and using |\&| instead of
% |and|. 
%    \begin{macrocode}
\newcommand{\@rawauthors}{Raw Authors}
\newcommand{\rmaaTOCauthors}{\noexpand\rmaaTOCAuthorStyle{Author list for TOC}}
\newcommand{\listofauthors}[1]{%
  \renewcommand{\@rawauthors}{#1}
  \renewcommand\rmaaTOCauthors{\noexpand\rmaaTOCAuthorStyle{#1}}}
%    \end{macrocode}
% \begin{macro}{\indexauthor}\label{indexauthor}
% The command |\indexauthor| produces an index entry for each
% author. 
% This command must come \emph{after} the
% |\title| and the |\listofauthors| commands. Also, there should be one 
% |\indexauthor| command for \emph{each} author. We also add the
% author onto the end of the |\ADS@authors| macro and construct the
% |\ADS@authorinitial| by a most cunning method (see below), but only
% for the first author. 
% 
% No, No, NO! This was all wrong. Backed out completely: 
% \begin{quote}
% Completely rewritten (19 Dec 2002) since we now use the index class, so 
% as to be able to alphabetize accented characters. A side effect of this is
% that |\index| no longer expands out the argument after the |@|, so we 
% instead add all the |\index| commands to a list (|\indexauthor@deferlist|), 
% which is executed in the |\@begofarticlehook|. 
% 
% Note the stuff with |\toks0| is copied from Exercise 20.15 of the \TeX{}book---it is
% to make sure the latin1-accented name gets passed into the idx
% file. 
% \end{quote}
% We now go back to the old way of doing things, using an optional
% argument to the \cs{indexauthor} command that has the name without
% accents. 
% \changes{3.08}{2003/01/16}{Remerged changes from 3.06 that got
% dropped from 3.07 after a synchronization disaster. Basically, the
% accentless optional argument to \cs{indexauthor}. Must add the macro
% directories to my unison set-up to stop this sort of thing from happening.}
%    \begin{macrocode}
\newcommand\@isortkey{}
\newcommand\@emptykey{empty}
\newcommand\indexauthor[2][empty]{
  \renewcommand\@isortkey{#1}
  \ifx\@isortkey\@emptykey \renewcommand\@isortkey{#2}\fi
  \index{\@isortkey @{{\rmaaIndexAuthorStyle{#2}}\hskip1em plus1em minus0.4em
      \rmaaIndexTitle\hskip1em plus2em minus0.4em\penalty-100
      \rmaaTOCauthors}}
  \indexauthorADS{#2}
}
%%\def\indexauthor@deferlist{\relax}
%%\long\def\indexauthor#1{
%%  \def\@@name{#1} \toks0=\expandafter{\@@name}
%%  \toks1=\expandafter{\indexauthor@deferlist}
%%  \protected@xdef\indexauthor@deferlist{%
%%    \the\toks1
%%    \noexpand\index{\the\toks0@{\noexpand\rmaaIndexAuthorStyle{\the\toks0}\noexpand\hskip1em plus1em minus0.4em 
%%        \rmaaIndexTitle\noexpand\hskip1em plus2em minus0.4em\noexpand\penalty-100
%%        \rmaaTOCauthors}}
%%  }
%%  \indexauthorADS{#1}
%%  }
%%\xdef\indexauthor@deferlist{}
\newcommand\indexauthorADS[1]{%
  \if@firstauthor\protected@xdef\ADS@authorinitial{\firstletter{#1}}\fi
  \protected@xdef\ADS@authors{\ADS@authors\if@firstauthor\else ; \fi
    #1}
  \@firstauthorfalse
}
%    \end{macrocode}
% Note that the first |#1| sets the key for the index, making sure it
% actually starts with  
% first letter of the surname. Meanwhile, the |@{...}| expression,
% which starts with some  
% control sequences is what is written to the |.ind| file. 
% See the makeindex man page for details. 
% \end{macro}
% \end{macro}
% \end{macro}
% \end{macro}

% \begin{macro}{\KillAllAccents}
% \begin{macro}{\firstletter}
% These are utility functions used in the above. |lett| is a counter
% indicating the current letter in the string. |\firstletter| just
% prints the first letter of its argument.
%    \begin{macrocode}
\newcommand{\KillAllAccents}{\def\'{}\def\`{}\def\~{}\def\.{}\def\i{i}}
\newcommand\firstletter[1]{\@car #1\@nil}
%    \end{macrocode}
% \end{macro}
% \end{macro}
% \begin{macro}{\SetVolume}
% \begin{macro}{\SetMSnumber}
% \begin{macro}{\SetFirstPage}
% \begin{macro}{\SetYear}
% \begin{macro}{\ReceivedDate}
% \begin{macro}{\AcceptedDate}
% These are for the editors to set the volume number, etc. 
%    \begin{macrocode}
\newcommand{\SetVolume}[1]{%
  \renewcommand{\@volume}{#1}
  \edef\ADS@volume{\PadToFourWithDots{#1}}
  \renewcommand{\@volumen}{\rmaa@VolumenStyle{volumen #1}}%
  }
\newcommand{\SetMSnumber}[1]{\renewcommand{\@MSnum}{#1}}
%    \end{macrocode}
% Note that there is no real reason to use |\SetFirstPage| since all
% the page numbering is done automatically in the book style. 
%    \begin{macrocode}
\newcommand{\SetFirstPage}[1]{\setcounter{page}{#1}}
\newcommand{\SetYear}[1]{%
  \renewcommand{\@year}{#1}
  \renewcommand{\ADS@year}{#1}
  \renewcommand{\@anyo}{\rmaa@VolumenStyle{#1}}
  }
\newcommand{\ReceivedDate}[1]{\renewcommand{\@received}{#1}}
\newcommand{\AcceptedDate}[1]{\renewcommand{\@accepted}{#1}}
%    \end{macrocode}
%
% \begin{macro}{\DOI}
% User/editor command to enter the DOI, which will be typeset as a link.
% \changes{4.0}{2019/11/02}{New user command \cs{DOI} to typeset
% DOI in the header of first page and add it to ADS info.}
%    \begin{macrocode}
\newcommand\DOI[1]{%
  \def\@DOI@url{\url{https://doi.org/10.22201/ia.01851101p.#1}}%
  \def\ADS@doi{https://doi.org/10.22201/ia.01851101p.#1}%
}%
%    \end{macrocode}
% \end{macro}
%
% \begin{macro}{\ORCIDiD}
% \changes{4.0}{2019/11/12}{Add mechanism for giving ORCID iD links}
% \begin{macro}{\ORCIDiDfull}
% \begin{macro}{\@orcid@icon}
% \begin{macro}{\@orcid@iconbig}
% Commands to add icon links to ORCID iDs of authors. For example
% \begin{quote}
% |\ORCIDiD|\marg{0000-0000-0000-0000}
% \end{quote}
% to be used after each author name in the |\author|
% command and in the |\fulladdresses|.  Alternatively,
% |\OORCIDiDfull| will typeset the icon plus the full text of the
% link, but we don't use that at present.  The icon image is stored in |ORCIDiD_iconvector.pdf|, which needs to be on the input path.  In principle, we could make EPS versions if necessary.
%    \begin{macrocode}
\newcommand\@orcid@icon{\includegraphics[height=1.8ex]{ORCIDiD_iconvector}}
\newcommand\@orcid@iconbig{\includegraphics[height=2ex]{ORCIDiD_iconvector}}
\newcommand\ORCIDiD[1]{\href{https://orcid.org/#1}{\,\@orcid@icon}}
\newcommand\ORCIDiDfull[1]{%
  \href{https://orcid.org/#1}{\@orcid@iconbig}\,\url{https://orcid.org/#1}}
%    \end{macrocode}
% \end{macro}
% \end{macro}
% \end{macro}
% \end{macro}
%
% \begin{macro}{\@volume}
% \begin{macro}{\@volumen}
% \begin{macro}{\@artpages}
% \begin{macro}{\@year}
% \begin{macro}{\@anyo}
% \begin{macro}{\@received}
% \begin{macro}{\@accepted}
% \begin{macro}{\@MSnum}
% Default values for the containers used by above commands, plus a
% few more useful ones. 
%    \begin{macrocode}
\newcommand{\@volume}{00}
\newcommand{\@volumen}{\rmaa@VolumenStyle{00}}
\newcommand{\@artpages}{%
  \oldpageref{\articlename:RMAAFirstPage}--\oldpageref{\articlename:RMAALastPage}}
\newcommand{\@year}{\the\year}
\newcommand{\@anyo}{\rmaa@VolumenStyle{2007}}
\newcommand{\@journalname}{\rmaa@JournalNameStyle{%
    \if@proceedings\RMAASClong\else\RMAAlong\fi}}
\newcommand{\@received}{}
\newcommand{\@accepted}{}
\newcommand{\@MSnum}{0000}
\newcommand{\@DOI@prefix}{https://doi.org/10.22201/ia.01851101p}
\def\@DOI@url{\url{https://doi.org/10.22201/ia.01851101p.2099.00.00.00}}
%    \end{macrocode}
% \end{macro}
% \end{macro}
% \end{macro}
% \end{macro}
% \end{macro}
% \end{macro}
% \end{macro}
% \end{macro}
% \end{macro}
% \end{macro}
% \end{macro}
% \end{macro}
% \end{macro}
% \end{macro}
% 
% The following commands are most useful in the \texttt{book}
% style. However, the |\SetConfTitle| command can now be used in
% individual articles to set the header for the \texttt{preprint}
% style. 
%    \begin{macrocode}
\newcommand{\@reunionplace}{Lugar de la reuni\'on}
\newcommand{\@reuniondate}{Fecha de la reuni\'on}
\newcommand{\@conftitle}{Conference Title}
\newcommand{\SetReunion}[2]{\renewcommand{\@reunionplace}{#1}\renewcommand{\@reuniondate}{#2}}
\newcommand{\SetConfTitle}[1]{\renewcommand{\@conftitle}{#1}}
\newcommand{\@editors}{Juan Ram\'on de la Fuente and ``El Mosh''}
\newcommand{\SetEditors}[1]{\renewcommand{\@editors}{#1}}
\newcommand{\@editores}{Juan Ram\'on de la Fuente y ``El Mosh''}
\newcommand{\SetEditores}[1]{\renewcommand{\@editores}{#1}}
\newcommand{\@mes}{mes de la publicaci\'on}
\newcommand{\SetMes}[1]{\renewcommand{\@mes}{\rmaa@VolumenStyle{#1}}}
%    \end{macrocode}

% \subsection{Notes for Author Proofs and Communication between Editors}
% \label{sec:notes}
% We introduce a command \cs{Note} that should be used to describe
% any change made to the manuscript. These are a bit tricky to get
% right. In general we use \cs{marginpar} but this only works in outer
% paragraph mode, so we just bung the box where we are if we are in a
% figure caption or something. If we are in the preamble, then we just
% use a \marginpar as part of the \maketitle command.
%    \begin{macrocode}
\newif\if@arenotes
\@arenotesfalse
\newcounter{notecounter}
\setcounter{notecounter}{0}
\newcommand\NotesList{}
\newcommand\Note[1]{%
  \if@proof{%
    \@arenotestrue%
    \stepcounter{notecounter}%
    \if@inpreamble\else{%
      \NoteMark{\arabic{notecounter}}%
    }\fi
    \protected@xdef\NotesList{\NotesList 
      \protect\item #1}%
  }\fi%
}
\newcommand\Changed[2]{Changed ``#1'' to ``#2''.}
\newcommand\NoteMark[1]{%
  \ifinner
  \raisebox{1ex}[0pt][0pt]{\makebox[0pt][l]{\framebox{\footnotesize\bfseries
        N#1}}}%
  \else
  \marginpar{\framebox{\footnotesize\bfseries N#1}}%
  \fi
}
\newcommand\RefNote[1]{\if@proof\NoteMark{\ref{#1}}\relax\fi}
\newcommand\NoteLabel{\protect\label}
\newcommand\OutputNotes{%
  \bgroup
  \onecolumn
  \section*{EDITORIAL NOTES}
  \begin{enumerate}
    \NotesList
  \end{enumerate}
  \egroup
}
\DeclareRobustCommand\cs[1]{\texttt{\char`\\#1}}

%    \end{macrocode}
% \subsection{Bibliographic Records for ADS}
% \label{sec:ADS}
% \changes{2.2beta11}{2001/11/27}{Just about finished all the ADS
% stuff now. Except for the references, which I have given up on for
% the time being.}
% \begin{macro}{\ADS@BibCode}
% \begin{macro}{\ADS@year}
% \begin{macro}{\ADS@journal}
% \begin{macro}{\ADS@volume}
% \begin{macro}{\ADS@firstpage}
% \begin{macro}{\ADS@authorinitial}
% We use the |\ADS@BibCode| command to hold the bibliographic code for
% ADS. It is 
% used in the |%R| 
% field of the bibliographic record for each article. 
% The format is 
% \begin{quote}
% \texttt{YYYYJJJJJVVVVMPPPPA}  
% \end{quote}
% where \texttt{YYYY} is the year, \texttt{JJJJJ} is the journal code,
% which is \texttt{RMxAA} or \texttt{RMxAC}, \texttt{VVVV} is the
% volume number (right-justified), \texttt{M} is not used by us,
% \texttt{PPPP} is the first page number (right-justified), and
% \texttt{A} is the initial of the first author's surname. The volume
% and page fields are padded with periods (`\texttt{.}'), which causes
% us some fun (this is done with the macro |\PadToFourWithDots| described
% below) . 
% 
% All these individual bits are held in the commands |\ADS@year|,
% |\ADS@journal|, |\ADS@volume|, |\ADS@firstpage|,
% |\ADS@authorinitial|. 
%    \begin{macrocode}
\newcommand{\SetADSBibCode}[1]{\renewcommand{\ADS@BibCode}{#1}}
\newcommand{\ADS@BibCode}
{\ADS@year\ADS@journal\ADS@volume.\PadToFourWithDots{\ADS@firstpage}\ADS@authorinitial}
\newcommand{\ADS@year}{2001} 
\newcommand{\ADS@journal}{\if@proceedings RMxAC\else RMxAA\fi}
\newcommand{\ADS@volume}{....} 
%    \end{macrocode}
% These two are just set to default values here, |\ADS@firstpage| is set
% in the |document| environment for included articles, in the
% |frontsection| environment, |\tableofcontents| command,  |index| and
% |abstracts| environments. Also by |\AtBeginDocument| for stand-alone
% articles. 
% 
% God knows how we are going to do |\ADS@authorinitial|!
%    \begin{macrocode}
\newcommand{\ADS@firstpage}{....}
\newcommand{\ADS@authorinitial}{A}
%    \end{macrocode}
% \end{macro}
% \end{macro}
% \end{macro}
% \end{macro}
% \end{macro}
% \end{macro}
% We also have a whole slew of other fields that can be set. These are
% all listed in Table~\ref{tab:ADS}. Most of them still have to be
% defined. 
% \begin{table}[!ht]\raggedleft
% \caption{ADS Bibliographic Records}\label{tab:ADS}
% \begin{tabular}{lp{8em}lp{8em}}\hline
% Code & Description & Macro Name & Comment \\ \hline
% |%R| & Bibliographic Code & |\ADS@BibCode| & required \\
% |%A| & Author list & |\ADS@authors| & required, generate from |\indexauthor| \\
% |%F| & Author affiliation & |\ADS@affil| & format: AA(\dots)
% AB(\dots) \dots \\
% |%J| & Journal Name & |\ADS@journalname| &  \\
% |%D| & Publication Date & |\ADS@fulldate| &  \\
% |%L| & Last Page & |\ADS@lastpage| & generate from label \\
% |%T| & Title & |\ADS@title| & required, generate from |\title| \\
% |%C| & Abstract Copyright & |\ADS@copyright| & set to IA-UNAM \\
% |%U| & URL for electronic document & |\ADS@url| & set to PDF file \\
% |%d| & DOI for article & |\ADS@doi| & set by |\DOI| \\
% |%K| & Keywords & |\ADS@keywords| & set equal to |\@keywords| \\
% |%A| & Abstract Text & |\ADS@abstract| & set equal to |\@abstract| \\
% |%R| & References & |\ADS@references| & generate from the
% bibliography  \\ \hline
% \end{tabular}
% \end{table}
% A lot of these macros are set by other commands, as mentioned in the
% Table, but some need to be set  explicitly here. 
% 
% \begin{macro}{\ADS@doi}
% \changes{4.0}{2019/11/08}{Save DOI for ADS}
% The user/editor command \cs{DOI} is used to set this. 
%    \begin{macrocode}
\newcommand{\ADS@doi}{https://doi.org/10.22201/ia.01851101p.0000.00.00.00}
%    \end{macrocode}
% \end{macro}
% 
% \begin{macro}{\ADS@journalname}
% This
% does not contain any accents at the moment so as not to cause
% problems---this maybe should be changed though. 
%    \begin{macrocode}
\newcommand{\ADS@journalname}{%
  \if@proceedings\@conftitle\space (Eds. \@editors) \fi%
  Revista Mexicana de Astronom\'\i{}a y Astrof\'\i{}sica%
  \if@proceedings\space (Serie de Conferencias)\fi%
  \space Vol. \@volume, pp. \ADS@firstpage--\ADS@lastpage\space (\@year)%
  \space (http://www.astroscu.unam.mx/\~{}rmaa/)
  }
%    \end{macrocode}
% \end{macro}
% \begin{macro}{\ADS@keywords}
% \begin{macro}{\ADS@abstract}
% \begin{macro}{\ADS@title}
% |\ADS@abstract| is straightforward, since we already have the text
% saved as |\@abstract|. We could do the same with |\ADS@keywords| but
% to avoid unwanted formatting we explicitly set it in |\keywords|
% (p.~\pageref{keywords}). |\ADS@title| is just set to |\@rawtitle|,
% which we were already using for another purpose. 
%    \begin{macrocode}
\let\ADS@title=\@rawtitle
\newcommand{\ADS@keywords}{No keywords specified}
\newcommand{\ADS@abstract}{No abstract specified}
%    \end{macrocode}
% \end{macro}\end{macro}\end{macro}
% \begin{macro}{\ADS@authors}
% |\ADS@authors| wants to be in ``Other, A. N.'' format, but with
% semi-colons between authors. The most natural way to do this seems to
% be with |\indexauthor| (p.~\pageref{indexauthor}). We append each
% author to this macro, which means we must explicitly clear it for
% each new article (done in |\article|). .   
%    \begin{macrocode}
\newcommand{\ADS@authors}{}
%    \end{macrocode}
% \end{macro}
% \begin{macro}{\ADS@fulldate}
% \begin{macro}{\ADS@lastpage}
% \begin{macro}{\ADS@copyright}
% \begin{macro}{\ADS@website}
% \begin{macro}{\ADS@url}
% \begin{macro}{\SetWebsite}
% \begin{macro}{\SetMonthNumeric}
% Other sundry data, together with user commands for bits that can be
% set by the Editors. The |\ADS@url| is constructed from the website
% (currently |http://www.astrosmo.unam.mx/rmaa/|,
% the journal descriptor (e.g., |RMxAC..10/|), followed by the
% path to the PDF file (e.g., |PDF/RMxAC..10_HEW1.pdf|). 
% 
%    \begin{macrocode}
\newcommand{\ADS@website}{http://www.astrosmo.unam.mx/rmaa/}
\newcommand{\ADS@url}
{\ADS@website\ADS@journal\ADS@volume\ADS@number/PDF/\ADS@journal\ADS@volume\ADS@number\_\articlename.pdf}
\newcommand{\ADS@number}{}
\newcommand{\ADS@fulldate}{00/\ADS@year} 
\newcommand{\ADS@lastpage}{00} 
\newcommand{\ADS@copyright}{(C) \ADS@year; Instituto de Astronom\'\i{}a,
  Universidad Nacional Aut\'onoma de M\'exico} 
\newcommand{\SetNumber}[1]{\renewcommand{\ADS@number}{-#1}}
\newcommand{\SetWebsite}[1]{\renewcommand{\ADS@website}{#1}}
\newcommand{\SetMonthNumeric}[1]{\renewcommand{\ADS@fulldate}{#1/\ADS@year}}
%    \end{macrocode}
% \end{macro}\end{macro}\end{macro}\end{macro}\end{macro}\end{macro}\end{macro}
% \begin{macro}{\ADS@affil}
% \begin{macro}{\ADS@references}
% These two are somewhat more problematical. Affiliation currently
% only works for the first author (see \dots) pending redesign of the
% user interface. References are now done by wrapping the
% |thebibliography| environment to open another input stream on the
% same file and search through for the |\bibitems|. 
%    \begin{macrocode}
\newcommand{\ADS@references}{}
\newcommand{\ADS@affil}{}
%    \end{macrocode}
% \end{macro}\end{macro}
% 
% Define two macros that expand to the strings
% |\begin{thebibliography}| and |\end{thebibliography}|. Note that the
% idea for this and much else is inspired by the code for the docstrip
% package. Also define |\StartBibX| in case authors use the more
% standard |\begin{thebibliography}{}| syntax. 
%    \begin{macrocode}
\edef\EndBib{\expandafter\@gobble\string\\end\string{thebibliography\string}}
\edef\StartBib{\expandafter\@gobble\string\\begin\string{thebibliography\string}}
\edef\StartBibX{\expandafter\@gobble\string\\begin\string{thebibliography\string}\string{\string}}
%    \end{macrocode}
% Define two new logical variables. |\ifAnotherline| is true if we
% want to read another line from the input stream. It is set to false
% when we encounter the |\end{thebibliography}|. |\ifInBib| is set to
% true when we frind the |\begin{thebibliography}| and tells us to
% start saving the following lines. 
%    \begin{macrocode}
\newif\ifAnotherLine
\newif\ifInBib
%    \end{macrocode}
% \begin{macro}{makeOther}
% We need a macro to change the catcodes of lots of things to other. 
%    \begin{macrocode}
\def\makeOther#1{\catcode`#1=12\relax}
%    \end{macrocode}
% \end{macro}
% \begin{macro}{\RereadBiblio}
% Define a macro to do all the fiddly stuff. This is used at the
% beginning of |thebibliography| environment (Sec~\ref{sec:biblio}).
%    \begin{macrocode}
\newread\tf@biblio
\newcommand{\RereadBiblio}{%
%    \end{macrocode}
% Open the current file again for reading
%    \begin{macrocode}
  \begingroup
  \immediate\openin\tf@biblio=\articlename
%    \end{macrocode}
% Everything in a group to be safe. Initialize logicals. Black magic
% to change
% catcodes and avoid blank lines being interpreted as |\par|. 
%    \begin{macrocode}
  \AnotherLinetrue\InBibfalse
  \makeOther\ \makeOther\\\makeOther\$%
  \makeOther\#\makeOther\^\makeOther\^^K%
  \makeOther\_\makeOther\^^A\makeOther\%%
  \makeOther\~\makeOther\{\makeOther\}\makeOther\&%
  \endlinechar-1\relax
%    \end{macrocode}
% Start looping through the file. Read each line to the macro
% |\Line|. If we have got to EOF or end of biblio then set
% |\AnotherLinefalse|.
%    \begin{macrocode}
  \loop
  \read\tf@biblio to \Line
  \ifx\Line\EndBib\AnotherLinefalse\InBibfalse\fi
  \ifeof\tf@biblio \AnotherLinefalse\fi
%    \end{macrocode}
% If we are in the biblio and the line is not blank then append it to
% the |\ADS@references| macro together with a newline. 
%    \begin{macrocode}
  \ifInBib
  \ifx\Line\empty\else\xdef\ADS@references{\ADS@references\Line^^J}\fi
  \fi
%    \end{macrocode}
% Check if the line we just saw was the start of biblio, then repeat
% the loop if necessary. 
%    \begin{macrocode}
  \ifx\Line\StartBib\InBibtrue\fi
  \ifx\Line\StartBibX\InBibtrue\fi
  \ifAnotherLine\repeat
  \immediate\closein\tf@biblio
  \endgroup
  }
%    \end{macrocode}
% \end{macro}
% 
% \begin{macro}{\PadToFourWithDots}
% Utility function for left-padding a number to 4 characters with
% periods. Used in 
% constructing |\ADS@BibCode|. This is a very simple algorithm and
% \emph{requires} that the argument actually be a number---if it is
% not then dire consequences will result.  
%    \begin{macrocode}
\newcommand{\PadToFourWithDots}[1]{%
  \ifnum #1 > 999 #1%
  \else\ifnum #1 > 99 .#1%
  \else\ifnum #1 > 9 ..#1%
  \else ...#1%
  \fi\fi\fi}
%%  \ifnum \count0 > 999 #1%
%%  \else\ifnum \count0 > 99 .#1%
%%  \else\ifnum \count0 > 9 ..#1%
%%  \else ...#1%
%%  \fi\fi\fi}
%    \end{macrocode}
% \end{macro}
%
% \begin{macro}{\AuthorADS}
% \begin{macro}{\ADS@affiliations}
% New mechanism for setting the affiliations for ADS. The author labels (AA, AB, \dots) are generated using two nested counters.  We use the |\robustify| command from \texttt{etoolbox} so that accented characters in the affiliation addresses do not blow up.  The list of accents may need to be expanded. 
%    \begin{macrocode}
\newcounter{AUTHOR}
\newcounter{author}[AUTHOR]
\stepcounter{AUTHOR}
\newcommand\ADS@affiliations{}
\robustify{\'}\robustify{\i}\robustify{\~}\robustify{\`}
\def\AuthorADS#1#2{%
  \begingroup
  \def\ORCIDiD##1{<ID system="ORCID">##1</ID>}
  \def\Email##1{<EMAIL>##1</EMAIL>}
  \def\Affil##1{##1}
  \stepcounter{author}
  \ifthenelse{\value{author} > 26}{\stepcounter{AUTHOR}\stepcounter{author}}{}
  \ifthenelse{\equal{\ADS@affiliations}{\@rmaaEmpty}}{%
    \xdef\ADS@affiliations{\Alph{AUTHOR}\Alph{author}(#2)}
  }{%
    \xdef\ADS@affiliations{\ADS@affiliations, \Alph{AUTHOR}\Alph{author}(#2)}
  }%
  \endgroup
}
%    \end{macrocode}
% \end{macro}
% \end{macro}
% %
% \begin{macro}{\ADS@dumpdata}
% \changes{4.0}{2019/11/13}{Write all author affiliations for ADS}
% \begin{macro}{\ADS@write}
% \begin{macro}{\ADS@openfile}
% |\ADS@dumpdata| is the macro that actually does the job of writing out all the
% data to an auxiliary file, using |\ADS@write| to write each
% field. The file is set up with |\ADS@openfile|, which currently is
% done in the TOC of a book.  The new multi-author affiliations are used if |\ADS@affiliations| has been set, otherwise it falls back to the single-author version |\ADS@affil|. 
%    \begin{macrocode}
\newcommand{\ADSfield}[2]{\relax}
\newcommand{\ADSvfield}[2]{\relax}
%%\renewcommand{\ADSfield}[12]{\%#1 #2}
\newcommand{\ADS@dumpdata}
{ 
  \ADS@write{\string\begin{ADSarticle}}
  \ADS@WRITE{\string\ADSfield{R}{\ADS@BibCode}}
  \ADS@write{\string\ADSfield{A}{\ADS@authors}}
  \ifthenelse{\equal{\ADS@affiliations}{\@rmaaEmpty}}{%
    {\ADS@sanitize\ADS@write{\string\ADSfield{F}{AA(\ADS@affil)}}}
  }{%
    {\ADS@sanitize\ADS@write{\string\ADSfield{F}{\ADS@affiliations}}}
  }
  \ADS@write{\string\ADSfield{J}{\ADS@journalname}}
  \ADS@write{\string\ADSfield{D}{\ADS@fulldate}}
  \ADS@write{\string\ADSfield{L}{\ADS@lastpage}}
  {\ADS@sanitize\ADS@write{\string\ADSfield{T}{\ADS@title}}}
  \ADS@write{\string\ADSfield{C}{\ADS@copyright}}
  \ADS@write{\string\ADSfield{d}{\ADS@doi}}
  \ADS@write{\string\ADSfield{U}{\ADS@url}}
  \ADS@write{\string\ADSfield{K}{\ADS@keywords}}
  {\ADS@sanitize\ADS@write{\string\ADSfield{B}{\ADS@abstract}}}
  \ADS@write{\string\ADSfield{Z}{\ADS@references}}
  \ADS@write{\string\end{ADSarticle}}
  }
%%\newcommand{\ADS@write}[1]{\if@book\addtocontents{ads}{#1}\fi}
%%\newcommand{\ADS@openfile}{\@starttoc{ads}}
\newcommand{\ADS@write}[1]{\set@display@protect\immediate\write\tf@ads{#1}}
\newcommand{\ADS@WRITE}[1]{\immediate\write\tf@ads{#1}}
\newcommand{\ADS@openfile}{%
  \newwrite\tf@ads%
  \immediate\openout\tf@ads\jobname.ads\relax
  \ADS@WRITE{\ADS@preamble}}
\newcommand{\ADS@closefile}{\ADS@WRITE{\ADS@postamble}}
\def\ADS@preamble{%
  %%  \string\documentclass{article}
  %%  \string\begin{document}
  }
\def\ADS@postamble{%
  %%  \string\end{document}
  \string%\string% Local Variables:^J 
  \string%\string% mode: LaTeX^J
  \string%\string% mode: Font-lock^J
  \string%\string% End:^J
  }
%    \end{macrocode}
% \end{macro}
% \end{macro}
% \end{macro}
% 
% \subsection{Automatic calculation of the page range}
% \label{sec:page-range}
% Another issue that has cost me a lot of grief. Still, it finally
% seems to be sorted now. All hail Jeff Goldberg. 
%
% Cunning trick to get the page numbers from |\pageref|.  Put a label
% on the first and last page of the document using |\AtBeginDocument|
% and |\AtEndDocument|. .  This way won't work for the book option.
% In that case, we do it by defining |\@endofarticlehook| and
% |\@begofarticlehook| as parallels to |\AtEndDocument| and
% |\AtBeginDocument|.  Note that we have to defer adding
% |\@begofarticlehook| to |\AtBeginDocument| to make sure it comes
% after the redefinitions of |\ref| and |\label|.
%    \begin{macrocode}
\if@book\AtEndDocument{\ADS@closefile}\else
 %\AtBeginDocument{\@begofarticlehook}
\AtEndDocument{\@endofarticlehook}
\fi
%    \end{macrocode}
% \begin{macro}{\@begofarticlehook}
% Since we are getting more and more things that need to be done at
% the beginning of every article, frontsection, etc we define a sister
% hook to 
% |\@endofarticlehook|. This hook does the following: 
% \begin{enumerate}
% \item Labels the first page
% \item Sets |\ADS@firstpage|
% \item runs and then clears |\indexauthor@deferlist|
% \item resets the keyword counter
% \end{enumerate}
% It is executed in the following circumstances: 
% \begin{enumerate}
% \item At |\begin{document}| for standalone files
% \item At |\begin{document}| for each included file
% \item At beginning of front sections 
% \item At beginning of TOC
% \item At beginning of Author Index
% \item At beginning of abstracts
% \end{enumerate}
% For some reason, this is not always within the right group, so we
% use |\xdef| just to be sure. 
%    \begin{macrocode}
\newcommand{\@begofarticlehook}{%
  \if@frontmatter\frontlabel{RMAAFirstPage}\else\label{RMAAFirstPage}\fi%
  \xdef\ADS@firstpage{\@arabic{\c@page}}%
  \if@book\else\makeindex\fi % needed so default index is defined in stand-alone articles
%%  \indexauthor@deferlist
%%  \xdef\indexauthor@deferlist{}
  \@inpreamblefalse
  \setcounter{kwdcounter}{0}
  }
%    \end{macrocode}
% \end{macro}
% \begin{macro}{\@endofarticlehook}
% This is a hook to be run at the end of each article
% currently, it just labels the last page and writes out all the ADS
% data to an auxiliary file.  
% It should not produce any printed output. 
% It is modelled on the |\AtEndDocument| command, to which it reduces
% in the case of a single article. 
%    \begin{macrocode}
\newcommand{\@endofarticlehook}{%
  \if@debug{\ClassWarning{rmaa}{***Doing Hook***}}\fi
  \if@debug{\ClassWarning{rmaa}{ADSabstract is: \ADS@abstract}}\fi
  \clearpage\rmaa@lastpagelabel\if@book\ADS@dumpdata\fi
  \ifx\NotesList\@rmaaEmpty\else\OutputNotes\fi
  \if@debug{\ClassWarning{rmaa}{Notes list is: \NotesList}}\fi
}
%    \end{macrocode}
% \begin{macro}{\rmaa@lastpagelabel}
% New version of last page label command copied from Jeff Goldberg's
% lastpage.sty package - we can't just use |\label| here in case the
% output routine for the last page has already been called. Instead,
% we first explicitly clear the page, step back the page counter,
% explicitly write to the .aux file, then restep the page counter
% again.
%
% But this will conflict with the hyperref package, which does labels differently.  Compensate for that by using the five-argument form.
% 
% Also set the |\ADS@LastPage| macro to the appropriate number. 
%    \begin{macrocode}
\newcommand{\rmaa@lastpagelabel}{%
  \addtocounter{page}{-1}%
  \immediate\write\@auxout{\string
    \newlabel{\articlename:RMAALastPage}{{}{\thepage}{}{}{}}}%
  \edef\ADS@lastpage{\thepage}%
  \addtocounter{page}{1}}
%    \end{macrocode}
% \end{macro}
% \end{macro}

% \subsection{Deal with explicit change in number of columns}
% \label{sec:columns}
% Sometimes, when formatting a book, we need to switch between
% one-column articles and two-column articles. For this to work, we
% need to wrap the standard \LaTeX{} commands to set the |\@rmtwocol|
% flag appropriately and also reset the package default values of all
% scalable lengths. .  
%    \begin{macrocode}
\let\latexonecolumn=\onecolumn
\let\latextwocolumn=\twocolumn
\renewcommand{\onecolumn}{\@rmtwocolfalse\@package@default@lengths\latexonecolumn}
\renewcommand{\twocolumn}{\@rmtwocoltrue\@package@default@lengths\latextwocolumn}
%    \end{macrocode}
% 
% \subsection{The actual formatting of the title}
% \label{sec:title-format}

% Note, we have to redefine the |\maketitle| command to use the
% numbered footnotes. This means we needn't use the atbeginend stuff. 
%    \begin{macrocode}
\renewcommand{\maketitle}{\par
  \begingroup
%    \end{macrocode}
% Don't switch to marked footnotes. 
% This is consistent with existing style but I'm not sure I like it
%    \begin{macrocode}
  %%  \renewcommand\thefootnote{\fnsymbol{footnote}}%
%%  \def\@makefnmark{\hbox to\z@{$\m@th^{\@thefnmark}$\hss}}%
  \def\@makefnmark{\textsuperscript{\@thefnmark}\hspace{-0.2em}}%
  \long\def\@makefntext##1{\parindent 1em\noindent
    \hbox to1.8em{\hss$\m@th^{\@thefnmark}$}##1}%
  \if@rmtwocol
  { \if@debug{\ClassWarning{rmaa}{TWO COLUMN}}\fi
    \twocolumn[\@maketitle]%
    }
  \else
  { \if@debug{\ClassWarning{rmaa}{ONE COLUMN}}\fi
    \newpage
    \global\@topnum\z@   % Prevents figures from going at top of page.
    \if@manuscript \@MSmaketitle \else \@maketitle \fi
    }
  \fi
  \if@manuscript\else\thispagestyle{rmaafirstpage}\@thanks\fi
  \ifx\NotesList\@rmaaEmpty\else{%
    \NoteMark{1--\arabic{notecounter}}\vspace*{-2\baselineskip}%
  }\fi%
  \endgroup
%    \end{macrocode}
% Don't reset footnote counter (not sure whether this is good thing
% or not).
%    \begin{macrocode}
  %%  \setcounter{footnote}{0}%
  \let\thanks\relax\let\affil\relax\let\altaffil\relax
  \let\maketitle\relax\let\@maketitle\relax
  \gdef\@thanks{}\gdef\@author{}\gdef\@title{}}
%    \end{macrocode}

% \begin{macro}{\@kwhwidth}
% \begin{macro}{\@kwtwidth}
% Redefine \cs{@maketitle} to include the abstract, resumen and
% keywords. Note that in the manuscripty style we use
% \cs{@MSmaketitle} instead.  
% \changes{3.16}{2003/04/07}{
%  Added extra vertical space between authors and accepted date.}
% \changes{3.26}{2011/08/15}{Add extra label to review articles}
%    \begin{macrocode}
\newlength{\@kwhwidth}\newlength{\@kwtwidth}
\newcommand\reviewname{Review}
\renewcommand{\@maketitle}{%
  \newpage
  \null
  \if@reviewarticle
  \smallskip
  {\large\bfseries\itshape\reviewname}%
  \bigskip
  \else
  \vskip\@before@title@skip%
  \fi
  \begin{center}%
    {\@title \par}%
    \vskip\@before@author@skip%
    \parbox{\rmaa@AuthorWidthFrac\textwidth}{\centering
      \lineskip\@between@author@skip%
      \@author
      \par}
    %% this should be replaced with `Submitted to...' for preprints
    \if@printdates
    \vskip\@before@dates@skip%
    \if@preprint
    \rmaa@ReceivedStyle{Submitted to RMxAA \@received}
    \else
    \rmaa@ReceivedStyle{Received \@received; accepted \@accepted}
    \fi
    \fi
    \if@manuscript
    \rmaa@ReceivedStyle{Submitted to RMxAA \@received}
    \fi
    \if@guide
    \rmaa@ReceivedStyle{Version \rmaa@version, \rmaa@date}
    \fi
    \vskip\@before@abstract@skip%
    \if@abstractpresent
    \parbox{\rmaa@AbstractWidthFrac\textwidth}{%
      \setlength{\parindent}{0.9cm}%
      \centerline{\rmaa@AbstractHeaderStyle{Abstract}}\par%
      \vskip\@after@abstracthead@skip%
      \if@newstyle\noindent\fi\@abstract
      \vskip\@between@abstract@skip%
      \centerline{\rmaa@AbstractHeaderStyle{Resumen}}\par%
      \vskip\@after@abstracthead@skip%
      \if@newstyle\noindent\fi\@resumen
      \if@setresumen\else\@abstract\fi % append english abstract if
                                       % spanish one not set
      \if@guide\else            %do not use keywords for guide
      \vskip\@before@kwd@skip%
      % Format the Key Words with hanging indentation:
      % They are put in a \parbox with width equal to the line width
      % minus the width of the header. 
      { % Widths of keyword
        \settowidth{\@kwhwidth}{\@keywordheader~}   % header and
        \setlength{\@kwtwidth}{\linewidth}          % keyword text
        \addtolength{\@kwtwidth}{-\@kwhwidth}       %
        \noindent\@keywordheader~\parbox[t]{\@kwtwidth}{%
          \rmaa@KeywordStyle{\boldmath\uppercase{\@keywords}}}}%
      \fi
      }
      \fi
  \end{center}%
  \vspace*{\afterkwdskip}
  \par
  }
\newlength{\afterkwdskip}\setlength{\afterkwdskip}{\z@}
%    \end{macrocode}
% \end{macro}
% \end{macro}
%
%
% \begin{macro}{\reviewarticle}
% \begin{macro}{\reviewtoc}
% \begin{macro}{\reviewlof}
%   User command to specify a review article. Must be used in
%   preamble. Main effect is to place the text ``Review'' before the
%   article title.  Also add command for setting a table of contents
%   and a list of figures at the beginning of the article.
%    \begin{macrocode}
\newcommand{\reviewarticle}{\@reviewarticletrue}
\newcommand\reviewtoc{%
  \section*{CONTENTS}%
  \@starttoc{rtc}%
}
\newcommand\reviewlof{%
  \section*{LIST OF FIGURES}%
  \@starttoc{lof}%
}
%    \end{macrocode}
% \end{macro}
% \end{macro}
% \end{macro}
% \begin{macro}{\l@section}
%    \begin{macrocode}
\renewcommand*\l@section[2]{%
  \ifnum \c@tocdepth >\z@
    \addpenalty\@secpenalty
    \addvspace{0.8em \@plus\p@}%
    \setlength\@tempdima{1.5em}%
    \begingroup
      \parindent \z@ \rightskip \@pnumwidth
      \parfillskip -\@pnumwidth
      \leavevmode %\bfseries
      \advance\leftskip\@tempdima
      \hskip -\leftskip
      #1\nobreak\hfil \nobreak\hb@xt@\@pnumwidth{\hss #2}\par
    \endgroup
  \fi}
%    \end{macrocode}
% \end{macro}
%
% \begin{macro}{\@MSmaketitle}
% A special version of |\@maketitle| for the \texttt{manuscript}
% option. This can deal with a multi-page abstract/resumen. Not very
% well, mind. 
%    \begin{macrocode}
\newcommand{\@MSmaketitle}{%
  \if@debug\ClassInfo{rmaa}{Using special MS maketitle command}\fi
  \newpage
  {
    \thispagestyle{rmaafirstpage}
%     \setlength{\parindent}{0pt}
    \null\vskip 2ex
    \begin{center}
      \@title\par
      \@author\par
      \rmaa@ReceivedStyle{%
        \ifx\@received\@rmaaEmpty
          Draft version: \@date
        \else
          Received: \@received
          \ifx\@accepted\@rmaaEmpty
          \else
            ; Accepted: \@accepted
          \fi
        \fi}\par
    \end{center}
    \@thanks
    \centerline{\rmaa@AbstractHeaderStyle{Resumen}}
    \noindent\@resumen\par\vskip 2ex
    \centerline{\rmaa@AbstractHeaderStyle{Abstract}}
    \noindent\@abstract\par\vskip 2ex
    \noindent\@keywordheader~\rmaa@KeywordStyle{\@keywords}\par
%     \@keywordheader~\rmaa@KeywordStyle{\uppercase{\@keywords}}\par
  }%
  }
%    \end{macrocode}
% \end{macro}

% \subsection{The sectioning commands}
% \label{sec:section-commands}

% \begin{macro}{\section}
% \begin{macro}{\subsection}
% \begin{macro}{\subsubsection}
% \begin{macro}{\paragraph}
% \begin{macro}{\subparagraph}
% Redefine all the standard commands to use these lengths. 
%    \begin{macrocode}
\renewcommand\section{\@startsection {section}{1}{\z@}%
                                   {\abovesectionskip}%
                                   {\belowsectionskip}%
                                   %% how do we get upper case here
                                   {\centering\reset@font\normalsize}}
\renewcommand\subsection{\@startsection{subsection}{2}{\z@}%
                                     {\abovesectionskip}%
                                     {\belowsectionskip}%
                                     {\centering\reset@font\normalsize\itshape}}
\renewcommand\subsubsection{\@startsection{subsubsection}{3}{\z@}%
                                     {\abovesectionskip}%
                                     {\belowsectionskip}%
                                     {\centering\reset@font\normalsize\itshape}}
\renewcommand\paragraph{\@startsection{paragraph}{4}{\z@}%
                                    {\aboveparaskip}%
                                    {\belowparaskip}%
                                    {\reset@font\normalsize\bfseries}}
\renewcommand\subparagraph{\@startsection{subparagraph}{5}{\parindent}%
                                      {\aboveparaskip}%
                                      {\belowparaskip}%
                                      {\reset@font\normalsize\bfseries}}
%    \end{macrocode}
% \end{macro}
% \end{macro}
% \end{macro}
% \end{macro}
% \end{macro}

% \begin{macro}{secnumdepth}
% Make sure that we number the subsubsections---this is done by
% default in the |article| class but not the |book| class, which is
% why we must set it explicitly here.
%    \begin{macrocode}
\setcounter{secnumdepth}{3}
%    \end{macrocode}
% \end{macro}
% \changes{2.2beta10}{2001/11/07}{Fixed \cs{secnumdepth} for \cs{book}
% style.}

% \begin{macro}{\thesubsection}
% \begin{macro}{\thesubsubsection}
% For subsection and subsubsection, we want the numbers to be roman
% but the title to be italic. This seems the easiest way to go about it.
%    \begin{macrocode}
\renewcommand{\thesubsection}{{\upshape\thesection.\@arabic\c@subsection}}
\renewcommand{\thesubsubsection}{\thesubsection{\upshape.\@arabic\c@subsubsection}}
%    \end{macrocode}
% \end{macro}
% \end{macro}

% \begin{macro}{\@seccntformat}
% Add ``.'' to the number here so we can remove it from |\thesection|,
% etc. Reduce gap between number and title WJH 27 Jan 1995
%    \begin{macrocode}
\def\@seccntformat#1{\csname the#1\endcsname .\hskip 0.4em}
%    \end{macrocode}
% \end{macro}

% \begin{macro}{\@sect}
% I didn't really want to mess with the |\@sect| command, but it was
% the only way I could manage to force upper case. Other than the
% commented parts, it is identical to the standard version. 
%    \begin{macrocode}
\def\@sect#1#2#3#4#5#6[#7]#8{
  \ifnum #2>\c@secnumdepth
     \let\@svsec\@empty
  \else
     \refstepcounter{#1}%
     \let\@@protect\protect
     \def\protect{\noexpand\protect\noexpand}%
     \edef\@svsec{\@seccntformat{#1}}%
     \let\protect\@@protect
  \fi
%    \end{macrocode}
% This is the bit I changed (27 Jan 1995). Use upper case for
% |\section|, use mixed case for the rest. 
%    \begin{macrocode}
  \ifnum #2=1                   
     \def\sec@upcase##1{\uppercase{##1}}%
     \def\sec@shape{\upshape}%
  \else                         
     \def\sec@upcase##1{\relax{##1}}%
     \def\sec@shape{\itshape}%
  \fi
  \@tempskipa #5\relax
  \ifdim \@tempskipa>\z@
     \begingroup #6\relax
          \@hangfrom{\hskip #3\relax\@svsec}%
                    {\interlinepenalty \@M \sec@upcase{#8}\par}%
     \endgroup
     \csname #1mark\endcsname{#7}%
%    \end{macrocode}
% Only write TOC entries if doing the Guide style. 
%    \begin{macrocode}
     \if@guide 
        \addcontentsline{toc}{#1}{%
          \ifnum #2>\c@secnumdepth \else
             \protect\numberline{\csname the#1\endcsname}%
          \fi
          #7}%
     \fi
%    \end{macrocode}
% Also allow per-article TOC for review papers
%    \begin{macrocode}
     \if@reviewarticle
        \ifnum #2>2 \else
           \addcontentsline{rtc}{#1}{%
             \protect\numberline{\csname the#1\endcsname}%
             \sec@upcase{\sec@shape #7}}%
        \fi
     \fi
  \else
     \def\@svsechd{#6\hskip #3\relax  %% \relax added 2 May 90
       \@svsec #8\csname #1mark\endcsname{#7}%
       \if@guide
          \addcontentsline{toc}{#1}{%
            \ifnum #2>\c@secnumdepth \else
               \protect\numberline{\csname the#1\endcsname}%
            \fi
            #7}%
       \fi
       }%
  \fi
  \@xsect{#5}}
%    \end{macrocode}
% \end{macro}


% \subsection{Appendices}
% \label{sec:appendices}
% Make sure that appendices in individual articles work properly. We
% copy the definition from the article class (the book one doesn't do
% what we want).  
% \begin{environment}{appendices}
%   This uses an environment |appendices| inside of which each
%   |\section| will become a separate appendix.  Note that the
%   appendices must come before the references (otherwise the full
%   addresses mechanism will not work properly). We make the equation
%   numbering be subordinate to the appendix (equations A12, B6,
%   etc). Finally, we have a medskip followed by APPENDICES to
%   introduce them. \changes{3.17}{2002/04/08}{ Added appendices environment
%   that does the right thing.}
%    \begin{macrocode}
\newenvironment{appendices}{\par
  \begingroup
  \setcounter{section}{0}%
  \setcounter{subsection}{0}%
  \renewcommand\thesection{\@Alph\c@section}
  \renewcommand\theequation{\thesection\arabic{equation}}
  \medskip
  \section*{APPENDICES}}
{\endgroup}
%    \end{macrocode}
% \end{environment}
% \begin{environment}{appendix}
%   If there is only one appendix, then use |appendix| environment
%   instead.  \changes{4.0}{2019/11/12}{ Added
%   appendix environment for case where there is only one section.}
%    \begin{macrocode}
\renewenvironment{appendix}{\par
  \begingroup
  \setcounter{section}{0}%
  \setcounter{subsection}{0}%
  \renewcommand\thesection{\@Alph\c@section}
  \renewcommand\theequation{\thesection\arabic{equation}}
  \medskip
  \section*{APPENDIX}}
{\endgroup}
%    \end{macrocode}
% \end{environment}

% \subsection{Floating thingies}
% \label{sec:floats}   

%% 
%    \begin{macrocode}
\renewcommand{\thetable}{\@arabic\c@table}
\renewcommand{\thefigure}{\@arabic\c@figure}
\renewcommand{\figurename}{Fig.}
\renewcommand{\tablename}{TABLE}
%\renewcommand{\fnum@table}{\MakeUpperCase{\tablename}~\thetable}
%    \end{macrocode}
% For some reason, I have to write `TABLE' explicitly here
% If I use |\tablename|, it gets converted to mixed case---God knows
% why. 
% |\figurename| is also being ignored. 
%    \begin{macrocode}
\renewcommand{\fnum@table}{TABLE~\thetable}
\renewcommand{\fnum@figure}{Fig.~\thefigure}
%    \end{macrocode}

% Do captions in 9 pt face in the main journal style and in new-style
% anything. 
%    \begin{macrocode}
\long\def\@makefigcap#1#2{%
  \if@proceedings\else\small\fi
  \if@newstyle\small\fi 
  \vskip\abovecaptionskip
  \sbox\@tempboxa{#1. #2}%
  \ifdim \wd\@tempboxa >\hsize
    #1. #2 \par
  \else
    \global \@minipagefalse
    \hb@xt@\hsize{\hfil\box\@tempboxa\hfil}%
  \fi
  \vskip\belowcaptionskip}
%    \end{macrocode}

% Make sure table caption is upper case. 
%    \begin{macrocode}
\long\def\@maketabcap#1#2{%
  \hbox to\hsize{\hfil #1 \hfil}
  \vskip 6pt 
  \centering\MakeUpperCase{#2}\par
  \vskip 6pt}
%    \end{macrocode}

%    \begin{macrocode}
\newlength{\@rottabcapwidth}
\long\def\@makerottabcap#1#2{%
  \setbox\@tempboxa\hbox{\MakeUpperCase{#2}}%
  \ifdim\wd\@tempboxa >.8\vsize
     \setlength\@rottabcapwidth{.8\vsize}
  \else
     \setlength\@rottabcapwidth{\wd\@tempboxa}
  \fi
  \rotatebox{90}{%
    \begin{minipage}{\vsize}\centering
      \begin{minipage}{\@rottabcapwidth}\centering
        #1\\ \vskip 6pt \MakeUpperCase{#2}
      \end{minipage}%
    \end{minipage}
    }%
  \hskip 6pt
  }

\newenvironment{Sideways}
    {\begin{sideways}\begin{minipage}{\vsize}\centering}
    {\end{minipage}\end{sideways}}

\newenvironment{Sidewaystable}
{%
  \let\oldcaption=\caption%
  \def\caption{
    \def\@captype{table}\renewcommand{\@makecaption}{\@maketabcap}
    \oldcaption}
  \begin{sidewaystable}\centering}
  {\end{sidewaystable}\let\caption=\oldcaption}
\newenvironment{Sidewaystable*}
{%
  \let\oldcaption=\caption%
  \def\caption{
    \def\@captype{table}\renewcommand{\@makecaption}{\@maketabcap}
    \oldcaption}
  \begin{sidewaystable*}\centering}
  {\end{sidewaystable*}\let\caption=\oldcaption}
%    \end{macrocode}

% Not sure why we need to stick these lines in. 
%    \begin{macrocode}
\renewenvironment{figure}
                 {\renewcommand{\@makecaption}{\@makefigcap}\@float{figure}}
                 {\end@float}
\renewenvironment{figure*}
                 {\renewcommand{\@makecaption}{\@makefigcap}\@dblfloat{figure}}
                 {\end@dblfloat}
\renewenvironment{table}
                 {\renewcommand{\@makecaption}{\@maketabcap}\@float{table}}
                 {\end@float}
\renewenvironment{table*}
                 {\renewcommand{\@makecaption}{\@maketabcap}\@dblfloat{table}}
                 {\end@dblfloat}
%    \end{macrocode}

% Redefine |\@caption| so that we don't bother with the list of
% figures. This saves us having to worry about fragile commands in the
% captions. 
%    \begin{macrocode}
\long\def\@caption#1[#2]#3{%
  \par
%%  \addcontentsline{\csname ext@#1\endcsname}{#1}%
%%    {\protect\numberline{\csname the#1\endcsname}{\ignorespaces #2}}%
  \begingroup
    \@parboxrestore
    \if@minipage
      \@setminipage
    \fi
    \normalsize
    \@makecaption{\csname fnum@#1\endcsname}{\ignorespaces #3}\par
  \endgroup}
%    \end{macrocode}

% Make tables not quite so cramped. Allow |\hline\hline| to give a
% poor man's thick rule. 
%    \begin{macrocode}
\renewcommand{\arraystretch}{1.15}
\setlength{\doublerulesep}{0.4pt}
%    \end{macrocode}

% \begin{macro}{\tablecols}
% \begin{macro}{\@tablecols}
% \begin{macro}{\tabcaptioncontinued}
% |\tablecols| is a user command for setting the number of columns. It
% is used by commands that try and span the entire table (I think
% there must be a better way of acheiving
% this). |\tabcaptioncontinued| is 
% a user command for setting the caption of a continued table on the
% next page. 
%    \begin{macrocode}
\newcommand{\tablecols}[1]{\renewcommand{\@tablecols}{#1}}
\newcommand{\@tablecols}{1}
\newcommand{\tabcaptioncontinued}{%
  \multicolumn{\@tablecols}{c}{\addtocounter{table}{-1}TABLE~\thetable\ (CONTINUED)} \\ \\
  }
%    \end{macrocode}
% \end{macro}
% \end{macro}
% \end{macro}

% \subsubsection{Footnotes in tables}
% \label{sec:table-footnotes}
% \begin{macro}{\tablenotemark}
% \begin{macro}{\tablenote}
% \begin{macro}{\tablenoteskip}
% Old user commands |\tablenotemark{MARK}|,
% |\tablenote{MARK}{TEXT}|---now deprecated. 
%    \begin{macrocode}
\newcommand\tablenotemark[1]{$^\mathrm{#1}$} 
\newcommand\tablenote[2]{\if@proceedings\else\small\fi ${}^\mathrm{#1}$ #2} 
\newcommand\tablenoteskip{\\[0.5\baselineskip]} 
%    \end{macrocode}
% \end{macro}
% \end{macro}
% \end{macro}

% \begin{macro}{\tabnotemark}
% \begin{macro}{\tabnotetext}
% \begin{macro}{\tabnotetextTWO}
% Alternative way of doing footnotes in tables. 
%    \begin{macrocode}
\newcommand{\tabnotemark}[1]{\textsuperscript{\lowercase{#1}}}
%\newcommand{\tabnotetextTWO}[4]{\tabnotetext{#1}{#2}\tabnotetext{#3}{#4}}
\newcommand{\tabnotetextTWO}[4]{%
  \multicolumn{\@tablecols}{l}{%
    \parbox{0.5\tabnotewidth}{\small\textsuperscript{#1}#2}%
    \parbox{0.5\tabnotewidth}{\small\textsuperscript{#3}#4}%
    }%
  \\}
\newcommand{\tabnotetext}[2]{%
  \multicolumn{\@tablecols}{@{}p{\tabnotewidth}@{}}{\small\textsuperscript{#1}#2}\\}
\newlength{\tabnotewidth}
%    \end{macrocode}
% \end{macro}
% \end{macro}
% \end{macro}


% \begin{macro}{changemargin}
%    \begin{macrocode}
% Command to allow margins to be changed on the fly. Useful to allow
% tables and Figures to protrude into the margins. 
%
% \changes{3.24}{2007/09/14}{Added \texttt{changemargin} environment
% (copied from TeX FAQ).}
\newenvironment{changemargin}[2]{% 
  % From TeX FAQ http://www.tex.ac.uk/cgi-bin/texfaq2html?label=chngmargonfly
  \begin{list}{}{%
    \setlength{\topsep}{0pt}%
    \setlength{\leftmargin}{#1}%
    \setlength{\rightmargin}{#2}%
    \setlength{\listparindent}{\parindent}%
    \setlength{\itemindent}{\parindent}%
    \setlength{\parsep}{\parskip}%
  }%
  \item[]}{\end{list}}
%    \end{macrocode}
% \end{macro}

% \subsection{The bibliography and references}
% \label{sec:biblio}
% \begin{macro}{\beforebiblioskip}
% \begin{macro}{\afterbiblioskip}
% I'm not entirely sure why I introduced these two commands, since
% directly putting |\vspace| before and after
% |\begin{thebibliography}| seems to work well enough. 
%    \begin{macrocode}
\newlength{\beforebiblioskip}
\setlength{\beforebiblioskip}{0pt}
\newlength{\afterbiblioskip}
\setlength{\afterbiblioskip}{0pt}
%    \end{macrocode}
% \end{macro}
% \end{macro}
% \begin{macro}{thebibliography}
% The environment itself is done as starred section with extra
% vertical space after heading. This space is currently zero. 
% In the book style we do the cunning trick with rereading the file in
% order to get the references into the |\ADS@references| macro.
% 
% \changes{3.05}{2003/01/02}{Removed the \cs{beforebiblioskip} and
% \cs{afterbiblioskip} commands, since they cause unwanted vertical
% space in some instances.}  
%    \begin{macrocode}
\renewenvironment{thebibliography}{%
  \if@book\RereadBiblio\fi
%%  \vspace*{\beforebiblioskip}
  \section*{REFERENCES} \frenchspacing
%%  \vspace*{\afterbiblioskip}
  \small
%    \end{macrocode} 
% The references themselves are set as a list in order 
% to get the hanging indentation. They are done in a smaller font and
% with french spacing so we needn't worry about the distinction
% between sentence-ending periods and abbreviations. 
%    \begin{macrocode} 
  \list{\null}{\leftmargin 1.5em\labelwidth\z@\labelsep\z@\itemindent -1.5em
    \itemsep\z@\parsep\z@\usecounter{enumi}}
  \sloppy\clubpenalty4000\widowpenalty4000
  \sfcode`\.=1000\relax\@biblisttrue}
{\endlist\@biblistfalse\normalsize\@outputaddressesifnecessary}
%    \end{macrocode}
% \end{macro}
% \subsection{Uniquification of cross-references}
% \label{sec:uniquification}
% 
% \begin{macro}{\articlename}
% In order that the |\label{}| command in individual articles be
% guaranteed to give a unique anchor, we append the article name (held
% in the |\articlename| command) to each key. This requires
% redefinition of the following commands: |\cite|, |\nocite|, |\ref|,
% |\pageref|, |\label|. 
% \begin{macro}{\oldcite}
% \begin{macro}{\oldnocite}
% \begin{macro}{\oldref}
% \begin{macro}{\oldpageref}
% \begin{macro}{\oldlabel}
% The original versions of these commands are
% saved as |\oldcite|, |\oldnocite|, etc. 
%
% The uniquifying of cross-references is now always done, so 
% we don't get multiply-defined labels in the single article case
%    \begin{macrocode}
\newcommand{\articlename}{anon}
\AtBeginDocument{
  \let\oldcite=\cite 
  \let\oldnocite=\nocite 
  \let\oldref=\ref 
  \let\oldpageref=\pageref
  \let\oldlabel=\label
}
%    \end{macrocode}
% \end{macro}
% \end{macro}
% \end{macro}
% \end{macro}
% \end{macro}
% \end{macro}
%
% Uniquify all cross-references by prepending the article name.  This
% now has to go into |\AtBeginDocument| so it is working on the new
% versions that hyperref.sty has set up.  Then we need to put the
% |\@begofarticlehook| afterwards, so that the first page label is using
% the new definitions.  If amsmath is loaded, then we also need to patch
% the special label command that they use. \changes{4.0}{2019/11/18}{Fix labels for amsmath display equations, in case that package has been loaded.}
%    \begin{macrocode}
\AtBeginDocument{
  \renewcommand{\ref}[1]{\oldref{\articlename:#1}%
    \if@debug{\ClassInfo{rmaa}{Reference to \articlename:#1}}\fi}
  \renewcommand{\pageref}[1]{\oldpageref{\articlename:#1}%
    \if@debug{\ClassInfo{rmaa}{Page reference to \articlename:#1}}\fi}
  \renewcommand{\label}[1]{\oldlabel{\articlename:#1}%
    \if@debug{\ClassInfo{rmaa}{Label for \articlename:#1}}\fi}%
  \ifdef{\label@in@display}{%
    \let\oldlabel@in@display=\label@in@display
    \renewcommand{\label@in@display}[1]{\oldlabel@in@display{\articlename:#1}%
      \if@debug{\ClassInfo{rmaa}{AMSTeX Label for \articlename:#1}}\fi}%
  }{}
  \@begofarticlehook
}%
%    \end{macrocode}
% Citations are more complicated since they can contain a list of
% keys. These are really only here for use with the rmaacite package,
% which is utterly deprecated (02 Jan 2003) now that I have the natbib
% code working. The command \cs{cite} should \emph{not} be used with
% natbib: use \cs{citep}, \cs{citet}, etc.\@ instead. 
%    \begin{macrocode}
\renewcommand{\nocite}[1]{\@bsphack
  \@for\@citeb:=#1\do{%
    \edef\@citeb{\expandafter\@firstofone\@citeb:\articlename}%
    \if@filesw\immediate\write\@auxout{\string\citation{\@citeb}}\fi
    \@ifundefined{b@\@citeb:\articlename}{\G@refundefinedtrue
      \@latex@warning{Citation `\@citeb' undefined}}{}}%
  \@esphack}
\renewcommand{\cite}[1]{\oldcite{#1:\articlename}}
%   \end{macrocode}


% Horrible fiddling with the bibliography citation mechanism. 
% This has to be globally available in case rmaacite puts it in the
% \texttt{.aux} file. 
%    \begin{macrocode}
\def\lbibcite#1#2#3{\global\@namedef{b@#1}{#2}
                    \global\@namedef{by@#1}{#3}}
%    \end{macrocode}

% \begin{macro}{\oldbibitem}
% \begin{macro}{\@rmaaBibArg}
% \begin{macro}{\@rmaaSpace}
% \begin{macro}{rmaabib}
% Of course, |\bibitem| needs to get the uniquification treatment as
% well. While we are at it, we redefine it to use consecutive integers
% for the key if the user leaves it blank (using the counter
% |rmaabib|). This allows one to put 
% |\bibitem{}| for all your references without getting annoying
% messages about multiply-defined labels.
% 
% We also take the opportunity to insert the optional argument to
% \cs{bibitem} if it is missing in the form Author(DATE) so that natbib
% doesn't try to do numerical citations. Natbib already provides a
% hook \cs{@extra@b@citeb} to help uniquify the citation key, so we
% just plug into that. However, it puts it \emph{after} the original
% key, not before as I had had it. Easiest solution is to change the
% way that I do it. 02 Jan 2003. Note that the \cs{@extra@b@citeb}
% does not get written to the .aux  file in the \cs{citation} command
% but it does work. 
% \changes{3.14}{06 Feb 2003}{Reinstalled old definition of \cs{bibitem}
%  for the case where the nonatbib option is requested, so that the
%  rmaacite package will work. This means that you can't use rmaacite
%  and natbib in the same volume.}
%    \begin{macrocode}
\if@natbib\def\@extra@b@citeb{:\articlename}\fi
%%\else
\let\oldbibitem\bibitem
\newcommand{\@rmaaBibArg}{}
\newcommand{\@rmaaSpace}{ }
\newcounter{rmaabib}
\if@natbib
\renewcommand{\bibitem}[2][ABC(2003)]{%
  \renewcommand{\@rmaaBibArg}{#2}
  \ifx\@rmaaBibArg\@rmaaEmpty
     \addtocounter{rmaabib}{1}
     \oldbibitem[#1]{\thermaabib:\articlename}
  \else\ifx\@rmaaBibArg\@rmaaSpace
     \addtocounter{rmaabib}{1}
     \oldbibitem[#1]{\thermaabib:\articlename}
  \else 
     \oldbibitem[#1]{#2:\articlename}
  \fi\fi
  \renewcommand{\@rmaaBibArg}{}
  }
\else
\renewcommand{\bibitem}[1]{%
  \renewcommand{\@rmaaBibArg}{#1}
  \ifx\@rmaaBibArg\@rmaaEmpty
     \addtocounter{rmaabib}{1}
     \oldbibitem{\articlename:\thermaabib}
  \else\ifx\@rmaaBibArg\@rmaaSpace
     \addtocounter{rmaabib}{1}
     \oldbibitem{\articlename:\thermaabib}
  \else
     \oldbibitem{\articlename:#1}
  \fi\fi
  \renewcommand{\@rmaaBibArg}{}
  }
\fi
%%\fi
%    \end{macrocode}
% \end{macro}
% \end{macro}
% \end{macro}
% \end{macro}

% \begin{macro}{oldpcite}
% \begin{macro}{oldscite}
% Then, we need to repeat the above scheme, but for the |rmaacite|
% definition. 
%    \begin{macrocode}
\newcommand{\@redefinebibitem}{%
  \let\oldbibitem\bibitem
  \let\oldcite\cite
  \let\oldpcite\pcite
  \let\oldscite\scite
  \def\bibitem[##1<##2>]##3{%
    \renewcommand{\@rmaaBibArg}{##3}
    \ifx\@rmaaBibArg\@rmaaEmpty
       \addtocounter{rmaabib}{1}
       \oldbibitem[##1<##2>]{\thermaabib:\articlename}
    \else\ifx\@rmaaBibArg\@rmaaSpace
       \addtocounter{rmaabib}{1}
       \oldbibitem[##1<##2>]{\thermaabib:\articlename}
    \else 
       \oldbibitem[##1<##2>]{##3:\articlename}
    % trust the author here
%%       \oldbibitem[##1<##2>]{##3}
    \fi\fi
    }
  }
%    \end{macrocode}
% \end{macro}
% \end{macro}


% \subsection{HTML Table of Contents}
% \label{sec:hyper-TOC}
% \begin{macro}{\HyperTOCEntry}
% \begin{macro}{\HyperTOCpart}
% \begin{macro}{\HyperTOCmeeting}
% \begin{macro}{\HyperTOCeditors}
% \begin{macro}{\HyperTOCpreamble}
% These commands are used to write information to the \texttt{.htoc}
% file, which is used in generating a Web Page of the Table of
% Contents. Since they only ever appear preceded by |\protect| in the
% argument to an |\addtocontents| command,  they should never need to
% be expanded by the macros. However, we define them here just in case
% \dots
%    \begin{macrocode}
\newcommand{\HyperTOCEntry}[3]{\relax}
\newcommand{\HyperTOCpart}[1]{\relax}
\newcommand{\HyperTOCmeeting}[3]{\relax}
\newcommand{\HyperTOCeditors}[1]{\relax}
\newcommand{\HyperTOCpreamble}{\relax}
%    \end{macrocode}
% \end{macro}
% \end{macro}
% \end{macro}
% \end{macro}
% \end{macro}
%

% \subsection{Formatting the poster abstracts} 
% These can now be used without using the book option.
%    \begin{macrocode}

\newcommand{\@makeminititle}{%
  \newpage
  \null 
  \vskip 3em%
  \begin{center}%
    {\@title \par}%
  \end{center}
  \vskip 3em%
  }

\newif\if@abstractsads
\@abstractsadsfalse
\newcommand\abstractsADS{\@abstractsadstrue}

\if@book\RequirePackage{multicol}\fi
\newenvironment{abstracts}[1][Abstracts of Contributed Papers]{%
  %% First, make sure we start on an odd page (06 Jun 2000)
  \if@book\clearpage\thispagestyle{empty}\cleardoublepage\fi
  %%
  \renewcommand{\articlename}{ABSTRACTS}
  \title{#1}\shortauthor{Abstracts}\shorttitle{Abstracts}
  \listofauthors{}
  \if@book
  \addcontentsline{toc}{chapter}{\protect\@fiddleauthors{\@TOCtitle}{\rmaaTOCauthors}}%
  \addtocontents{htc}{\protect\HyperTOCEntry{\@rawtitle}{\@rawauthors}{\articlename}}
  \fi
  \pagestyle{rmaa} % this now seems necessary (06 Jun 2000)
  \thispagestyle{rmaafirstpage}
%%  \twocolumn[\@makeminititle]
  \renewcommand{\abstract}[1]{%
    \renewcommand\ADS@abstract{##1}%
    \renewcommand{\@abstract}{##1}\@writeabstract%
  }
  \renewcommand{\title}[2][]{%
    \gdef\ADS@authors{}%
    \gdef\ADS@affil{}%
    \gdef\ADS@references{}%
    \gdef\ADS@keywords{}%
    \renewcommand{\@rmaaArg}{##1}
    \ifx\@rmaaArg\@rmaaEmpty
       \renewcommand{\@TOCaftertitle}{}
    \else
       \renewcommand{\@TOCaftertitle}{ (\@rmaaArg)}
    \fi
    \renewcommand{\@title}{\rmaa@TOCTitleStyle{##2}}%
    \renewcommand{\@TOCtitle}{\rmaa@TOCTitleStyle{##2}\@TOCaftertitle}%
    \renewcommand{\rmaaIndexTitle}{\rmaa@IndexTitleStyle{##2\@TOCaftertitle.}}%
    \renewcommand\ADS@title{##2}%
    \xdef\ADS@firstpage{\@arabic{\c@page}}%
    }
  \renewcommand{\author}[1]{\let\@thanks\@empty\renewcommand\@author{##1} }
  \renewcommand{\footnotetext}[2][XX]{\par\noindent\textsuperscript{##1}~##2}
  \renewcommand{\footnoterule}{\kern-3\p@ \hrule \@width 3em \kern 2.6\p@}
  \if@book\@begofarticlehook\fi
  \if@book\begin{multicols}{2}[\@makeminititle]\else\@makeminititle\fi
  \begingroup
  \setcounter{footnote}{0}
  }% and at the end of the environment....
  {\endgroup\if@book\end{multicols}\listofabstracts\label{RMAALastPage}\clearpage\fi}
%% auxiliary commands used:
%% \@writeabstract does all the actual formatting 
\newcommand{\@writeabstract}{%
  \if@book%
  \addcontentsline{loa}{abstract}{\protect\@fiddleauthors{\@TOCtitle}{\rmaaTOCauthors}}%
  \fi%
  \begin{center}\@title\par\@author\end{center}
  \@abstract
  \par\noindent\footnoterule\nobreak\footnotesize\@thanks
  \par\noindent\hrulefill\goodbreak
  \vskip 1cm \@plus0.5cm
  %% write ADS entry
  \if@book\if@abstractsads\ADS@dumpdata\fi\fi%
  \endgroup % this pair is to make each abstract a self-contained group
  \begingroup
  \setcounter{footnote}{0}
  }

\if@book
%% List of abstracts (stuff written to .loa file)
\newcommand{\listofabstracts}{%
  \begin{multicols}{2}[%
    \vskip 1.5cm \vfill % leave at least 1.5 cm and go to bottom of page if possible
    \centerline{\MakeUpperCase{List of Abstracts}}
    \vskip 1.5cm \@plus0.5cm
    ][6cm] % this makes it start a new page if there is not at least 6cm of space
    \renewcommand{\@fiddleauthors}[2]{##1\hskip1em plus2em minus0.4em\penalty-100  ##2}
    \parindent\z@
%%    \parskip.5\p@ \@plus .7\p@\relax 
    \parskip3\p@ \@plus 1\p@ \@minus 1\p@\relax % allow quite a lot of
                                % leeway here to avoid breaks inside
                                % entries 
    \renewcommand{\\}{\relax} % temporarily turn off explicit line breaks just in case
    \renewcommand{\altaffilmark}[1]{\relax} % ditto with this
    \@starttoc{loa}
  \end{multicols}
  }
\newcommand*{\l@abstract}[2]{%
  \par\penalty-1000
  #1\hskip 1em\penalty0\null\nobreak\hfill#2}
\fi

%    \end{macrocode}
 
% \subsection{Code for typesetting an entire volume of conference
% proceedings} \label{sec:whole-volume}
%    \begin{macrocode}

\if@book
%% environment for including the individual articles
\newcommand{\@fiddleauthors}[2]{\relax}
\renewcommand{\theequation}{\@arabic\c@equation}
\newenvironment{articles}%
{ \begingroup % make it a group so that author definitions do not
              % spill over
  \gdef\ADS@authors{}
  \pagestyle{rmaa}
  \renewenvironment{document}{%
    \@begofarticlehook
    \addtocounter{chapter}{1}%
    \addcontentsline{toc}{chapter}{%
      \protect\@fiddleauthors{\@TOCtitle}{\rmaaTOCauthors}}%
    \addtocontents{htc}{\protect\HyperTOCEntry{\@rawtitle}{\@rawauthors}{\articlename}}
    }{\relax}
  \renewcommand{\documentclass}[2][foo]{\relax} 
  \renewcommand{\usepackage}[2][foo]{%
    \def\@packagename{##2}\def\@rmaacitename{rmaacite}%
    \ifx\@packagename\@rmaacitename
    \@rmcitetrue\ClassInfo{rmaa}{This paper uses rmaacite}
    \makeatletter%%
%% This is file `rmaacite.sty',
%% generated with the docstrip utility.
%%
%% The original source files were:
%%
%% rmaa.dtx  (with options: `rmaacite')
%% 
%% PLEASE DO NOT EDIT THIS FILE! It is generated automatically from
%% the file rmaa.dtx using the accompanying Makefile and rmaa.ins
%% Any hand-edited changes will be lost
%% 
%% A LaTeX document class for the RMxAA
%% 
%% Author: William Henney (w.henney@crya.unam.mx)
%% 
%% Latest version available from: http://www.crya.unam.mx/rmaa
%% 
%% Copyright (C) 1999-2010 Instituto de Astronomia, UNAM, Mexico
%% 
%% This program is free software; you can redistribute it and/or
%% modify it under the terms of the GNU General Public License
%% as published by the Free Software Foundation; either version 2
%% of the License, or (at your option) any later version.
%% 
%% This program is distributed in the hope that it will be useful,
%% but WITHOUT ANY WARRANTY; without even the implied warranty of
%% MERCHANTABILITY or FITNESS FOR A PARTICULAR PURPOSE.  See the
%% GNU General Public License for more details.
%% 
%% You should have received a copy of the GNU General Public License
%% along with this program; if not, write to the Free Software
%% Foundation, Inc., 59 Temple Place - Suite 330, Boston, MA  02111-1307, USA.






\def\bibitem[#1<#2>]#3{\par\hangindent=1.5em\hangafter=1
  \if@filesw
  {\def\protect##1{\string ##1\space}
    \immediate\write\@auxout{\string\lbibcite{#3}{#1}{#2}}}
  \fi}
\@redefinebibitem % what it says
\def\lbibcite#1#2#3{\global\@namedef{b@#1}{#2}
                    \global\@namedef{by@#1}{#3}}

\def\@bcite#1#2{({#1\if@tempswa , #2\fi})}
\def\@pcite#1#2{#1\if@tempswa , #2\fi}
\def\@brcite{\ (}
\def\@ebrcite{)}
\def\@fmttop#1{\@fmtcite[\csname #1\endcsname]}
\def\@citex[#1]#2{%
%%  \if@filesw\immediate\write\@auxout{\string\citation{#2}}\fi
%% WJH added 16 Jan 2002
  \@for\@citeb:=#2\do{%
    \edef\@citeb{\expandafter\@firstofone\@citeb:\articlename}%
    \if@filesw\immediate\write\@auxout{\string\citation{\@citeb}}\fi}%
%% to here
  \def\@citea{}%
  \@cite{\@for\@citeb:=#2\do
    {%
      \edef\@citeb{\expandafter\@firstofone\@citeb:\articlename}%
      \@citea\def\@citea{;\ }\@ifundefined
      {b@\@citeb}{{\bf ?}\@warning {Citation `\@citeb' on page \thepage \space undefined}}%
      {\csname b@\@citeb\endcsname\@citebro\csname by@\@citeb\endcsname\@citebrc}}}%
  {#1}}

\def\cite{\@ifnextchar [{\let\@citebro=\space
                          \let\@citebrc=\relax
                          \let\@cite=\@bcite\@tempswatrue\@citex}
                        {\let\@citebro=\space
                          \let\@citebrc=\relax
                          \let\@cite=\@bcite\@tempswafalse\@citex[]}}
\def\pcite{\@ifnextchar [{\let\@citebro=\space
                          \let\@citebrc=\relax
                          \let\@cite=\@pcite\@tempswatrue\@citex}
                        {\let\@citebro=\space
                          \let\@citebrc=\relax
                          \let\@cite=\@pcite\@tempswafalse\@citex[]}}
\def\scite{\@ifnextchar [{\let\@citebro=\@brcite
                          \let\@citebrc=\@ebrcite
                          \let\@cite=\@pcite\@tempswatrue\@citex}
                        {\let\@citebro=\@brcite
                          \let\@citebrc=\@ebrcite
                          \let\@cite=\@pcite\@tempswafalse\@citex[]}}

\renewenvironment{thebibliography}
  {\section*{REFERENCES}\vspace*{\afterbiblioskip}\small
    \if@book\RereadBiblio\fi
    \sloppy\clubpenalty4000\widowpenalty4000
    \sfcode`\.=1000\relax
    \setlength\parindent{\z@}\setlength\parskip{\z@}}
  {\relax\normalsize\@outputaddressesifnecessary}








\endinput
%%
%% End of file `rmaacite.sty'.
\makeatother % read it in here so that all changes
    \else                                         % are local to paper
    \if@debug\ClassWarning{rmaa}{Package \@packagename\space was requested by \articlename\space}\fi
    \fi} 
  %% Don't allow individual articles to mess with the pages, volume or year
  \renewcommand{\SetFirstPage}[1]{\relax}
  \renewcommand{\SetYear}[1]{\relax}
  \renewcommand{\SetVolume}[1]{\relax}
  \renewcommand{\SetConfTitle}[1]{\relax}
  %% Don't allow individual articles to change these
  \renewcommand{\nonstopmode}{\relax}
  \renewcommand{\scrollmode}{\relax}
  \renewcommand{\outputplates}{\relax}
  %% XXXX need to check if using rmaacite and act accordingly
  %% avoid multiple declarations of the same length
  \let\oldnewlength=\newlength
  \renewcommand{\newlength}[1]{%
    \edef\reserved@a{\expandafter\@gobble\string ##1}%
    \@ifundefined{\reserved@a}{%
      \oldnewlength{##1}}{%
      \if@debug{\ClassInfo{rmaa}{\reserved@a\space already defined: doing nothing}}\fi%
      \relax}}%
  \renewcommand\newcommand{\providecommand}
  }%
{\endgroup\onecolumn}
     

%    \end{macrocode}
% Command to be used in `articles' environment for including an
% article
%    \begin{macrocode}
\newcommand{\article}[2][]{%
  % optional first argument is extra text to put in brackets after title in TOC
  % e.g. ``Invited Talk'' or ``Poster''
  \renewcommand{\articlename}{#2}
  % reset loads of counters
  \setcounter{footnote}{0}
  \setcounter{figure}{0}
  \setcounter{table}{0}
  \setcounter{section}{0}
  \setcounter{equation}{0}
  \setcounter{rmaabib}{0}
  \setcounter{author}{0}
  \setcounter{AUTHOR}{1}
  \gdef\ADS@authors{}%
  \gdef\ADS@affil{}%
  \gdef\ADS@affiliations{}%
  \gdef\ADS@references{}%
  \@firstauthortrue
  \@firstaffiltrue
  \@firstaltaffiltrue
  \@firstkwdtrue
  \@rmcitefalse
  \if@suppressfulladdresses\else\@needaddressestrue\fi
  \renewcommand{\@rmaaArg}{#1}
  \ifx\@rmaaArg\@rmaaEmpty
     \renewcommand{\@TOCaftertitle}{}
  \else
     \renewcommand{\@TOCaftertitle}{ (#1)}
  \fi
%    \end{macrocode}
% In the main journal articles must start on an odd page. 
%    \begin{macrocode}
  \if@journal\clearpage\thispagestyle{empty}\cleardoublepage\fi
%% Now input the file, then run the end-of-article hook (reinstated 28
  %% Sep 2001)
  {\input{\@articledir/\articlename.tex}\@endofarticlehook}}
%    \end{macrocode}
% Command to separate articles into sessions. Write an entry in the
% TOC, makes sure we start on an odd page (added 06 Jun 2000), then do
% a title page if \cs{SessionTitlePages} was called. 
% \changes{3.12}{04 Feb 2003}{Added option to have title page for each
% session. Turn it on with \cs{SessionTitlePages}.}
%    \begin{macrocode}
\newcommand{\session}[1]{
  %% this does nothing except write an entry in the TOC
  \addcontentsline{toc}{part}{\rmaa@TOCPartStyle{#1}}
  \addtocontents{htc}{\protect\HyperTOCpart{#1}}
  %% and also makes sure we start on an odd page (added 06 Jun 2000)
  \clearpage\thispagestyle{empty}\cleardoublepage
  \stepcounter{part}
  \if@sessiontitlepages{%
    \onecolumn
    \thispagestyle{empty}
    \renewcommand{\articlename}{Section\thepart TitlePage}
    \label{RMAAFirstPage}
    \make@session@title{#1}
    \label{RMAALastPage}
    \clearpage\thispagestyle{empty}\cleardoublepage
  }\fi
  }
\newcommand{\make@session@title}[1]{%
  {%
    \setlength\parindent{0pt}
    \vspace*{\stretch{1}}
    \rule{\linewidth}{1mm}
    \begin{flushright}\Huge\bfseries\boldmath
      #1
    \end{flushright}
    \rule{\linewidth}{1mm}
    \vspace*{\stretch{2}}
  }
}
%    \end{macrocode}
% Command to turn on separate title pages for each session. 
%    \begin{macrocode}
\newcommand{\SessionTitlePages}{\@sessiontitlepagestrue}
\newcommand{\BoldSessionTitles}{%
  \renewcommand{\rmaa@TOCPartStyle}[1]{\textbf{\rmaa@TOCTitleStyle{##1}}}
  }
\newcommand{\ArticleDir}[1]{\renewcommand{\@articledir}{#1/}\graphicspath{{1/}}}
\newcommand{\@articledir}{.}
%\newcommand{\ps@rmaafront}{\ps@rmaa}% page style for front matter (same as ever, actually)
\newcommand{\ps@rmaafront}{% page style for front matter:  page # in each foot
  \renewcommand\@oddfoot{\hfil{\thepage}}%
  \renewcommand\@evenfoot{{\thepage}\hfil}%
  % shorttitle in both heads
  \renewcommand\@evenhead{\hfil\@shorttitle\hfil}%
  \renewcommand\@oddhead{\@evenhead}
  }
\newcommand{\ps@rmaafrontfirstpage}{% page style for first page of front matter
  % page # in each foot, on outside edge
  \renewcommand\@oddfoot{\hfil\thepage}%
  \renewcommand\@evenfoot{\thepage\hfil}%
  % nothing in the even heads
  \renewcommand\@evenhead{\relax}%
  \renewcommand\@oddhead{\relax}%
  }
%% slight modification to the standard \frontmatter and \mainmatter
%% commands, just to set the \pagestyle
\renewcommand\frontmatter{%
  \latexonecolumn
  \cleardoublepage
  \@frontmattertrue
  \@mainmatterfalse
  \pagenumbering{roman}
  \pagestyle{rmaafront}
  \addtocounter{page}{2} % leave space for the title page
  }
\renewcommand\mainmatter{%
  \pagestyle{empty}
  \if@rmtwocol\twocolumn\fi
  \cleardoublepage
  \@frontmatterfalse
  \@mainmattertrue
  \pagenumbering{arabic}
  \pagestyle{rmaa}
  }
\renewcommand\backmatter{%
  \clearpage\pagestyle{empty}\cleardoublepage
  \@mainmatterfalse
  %% put extra space in the TOC
  \addtocontents{toc}{\bigskip}
  }
%    \end{macrocode}
%
% \subsection{Front sections}\label{sec:front-sections}
% \begin{macro}{\notitle}
% \begin{macro}{\shortname}
% \begin{macro}{frontsection}
% The environment \texttt{frontsection} is for doing a section in the
% front matter (preface, list of participants, or whatever). It has 
% one optional argument:
% \begin{enumerate}
% \item[\#1] Short title (no embedded spaces) for
% use in making PS files. E.g., ``list\_of\_participants''.   
% \end{enumerate}
% and two mandatory arguments: 
% \begin{enumerate}
% \item[\#2] Title of section (e.g., `Preface')
% \item[\#3] Commands to be evaluated before processing the contents
% of the environment. This is mainly useful for putting in the things
% that would go in the preamble of a normal article  (e.g.,
% |\listofauthors|, |\indexauthor|, |\notitle|) 
% \end{enumerate}
%    \begin{macrocode}
\newcommand{\notitle}{\@titlevisiblefalse} \newcommand{\shortname}{}
\newenvironment{frontsection}[3][]{%
  \begingroup
  \renewcommand{\shortname}{#1}
  \ifx\shortname\@rmaaEmpty
     \renewcommand{\articlename}{#2}
  \else
     \renewcommand{\articlename}{\shortname}
  \fi
  \@begofarticlehook
  \@titlevisibletrue % default
  \thispagestyle{rmaafrontfirstpage}
  \listofauthors{} % blank out this in case it is not specified
  \title{{#2}}\shorttitle{#2}\shortauthor{#2}
  #3 % evaluate the command given as second mandatory arg
  \addcontentsline{toc}{chapter}{\protect\@fiddleauthors{\@TOCtitle}{\rmaaTOCauthors}}%
  \addtocontents{htc}{\protect\HyperTOCEntry{\@rawtitle}{\@rawauthors}{\articlename}}
  \if@titlevisible
     \vspace*{1cm}\centerline{\rmaa@FrontSecStyle{\@title}}\vspace*{1cm}
  \fi
  \setlength{\parskip}{3ex}
  }{\frontlabel{RMAALastPage}\clearpage\endgroup}
%    \end{macrocode}
% \end{macro}
% \end{macro}
% \end{macro}
%    \begin{macrocode}
%% Definition of the command \frontlabel, which rights out something like `F4' to
%% the .aux file, instead of `\textlatin  {i}'. This is so it is easy to post-process
%% the file for making PS files of individual sections
\newcommand{\frontlabel}[1]{%
  \@bsphack
  \protected@write\@auxout{}%
     {\string\newlabel{\articlename:#1}{{}{F\@arabic{\c@page}}}}
  \@esphack
  }
%% User command to say X, Y, Z,\\ Editors of the Proceedings
\newcommand{\EditorSignature}{\vspace*{\parskip}\rightline{\@editors}%
  \rightline{Editors of the proceedings}}
%% Fiddle with \thesection, etc. so that they don't include the chapter number
%% Basically go back to the `article' definitions of these
\renewcommand{\thesection}{\@arabic\c@section}
\renewcommand{\thefigure}{\@arabic\c@figure}
\renewcommand{\thetable}{\@arabic\c@table}

\newcommand{\@oldargs}{}
\newcommand{\@newargs}{}
%% add TOC entry for each paper
%\renewcommand{\title}[1]%
%{\renewcommand{\@title}{\rmaa@TitleStyle{#1}}%
%\addtocounter{chapter}{1}\addcontentsline{toc}{chapter}{#1}%
%}

%% Formatting of the table of contents
%% dimensions of various parts
\newlength{\TOCtextwidth}\setlength{\TOCtextwidth}{12.65cm}
\newlength{\TOCpagenumwidth}\setlength{\TOCpagenumwidth}{1.3cm}
\newlength{\TOCentrywidth}\setlength{\TOCentrywidth}{\TOCtextwidth}
\addtolength{\TOCentrywidth}{-\TOCpagenumwidth} 
\newlength{\TOCextramargin}\setlength{\TOCextramargin}{1.25cm}
\newcommand{\TOCbox}[1]{\hspace{\TOCextramargin}\parbox{\TOCtextwidth}{#1}\par}
\newlength{\TOCextramarginx}\setlength{\TOCextramarginx}{\TOCextramargin}\addtolength{\TOCextramarginx}{4.3cm}
\newlength{\TOCtextwidthx}\setlength{\TOCtextwidthx}{\TOCtextwidth}\addtolength{\TOCtextwidthx}{-4.3cm}
\newcommand{\TOCboxx}[1]{\hspace{\TOCextramarginx}\parbox{\TOCtextwidthx}{#1}\par}

\newcommand{\RMAAbannerimage}{rmsc}
\renewcommand\tableofcontents{%
  \renewcommand{\articlename}{CONTENTS}
  \begingroup
  \@begofarticlehook
  \if@twocolumn
     \@restonecoltrue\onecolumn
  \else
     \@restonecolfalse
  \fi
  %% 9pt for journal style
  \if@proceedings\else\small\fi
  \if@newstyle\small\fi %% same for new style
  \if@proceedings\thispagestyle{rmaafrontfirstpage}\else\thispagestyle{empty}\fi
  \shortauthor{Contents}\shorttitle{Contents}
  \vspace*{18mm}
  \if@printbanner
  %% put in the banner header (need to get better image of this)
     \TOCbox{\hspace*{-0.15cm}\includegraphics[width=0.9\TOCtextwidth]{\RMAAbannerimage}}
     \vspace*{8mm}
  %% volume and publication date (should really be in funny font)
     \TOCboxx{\@volumen\\[2mm] \@mes\ \@anyo}
  \else
  %% put in the banner header (need to get better image of this)
     \TOCbox{\hspace*{-0.15cm}\includegraphics[width=0.9\TOCtextwidth]{\RMAAbannerimage}}
     \vspace*{12.5mm}
  %% volume and publication date (should really be in funny font)
     \TOCboxx{~\\[2mm] ~}
  \fi
  \vspace*{16mm}
  %% This is the command that fiddles with the TOC entry so that the
  %% title is flushleft but the authors are flushright
  \renewcommand{\@fiddleauthors}[2]{%
    % first put the paper title
    \renewcommand{\\}{\relax} % temporarily turn off explicit line breaks just in case
    \renewcommand{\altaffilmark}[1]{\relax} % ditto with this
    \renewcommand{\textsuperscript}[1]{\relax} % ditto with this 
    ##1
    % now save the author list in a box
    \sbox\@tempboxa{##2}%
    % check that box is shorter than the total width available for the entry
    \ifdim \wd\@tempboxa <\TOCentrywidth
       % if so, then put it right here, flush against right margin
       \penalty10000\hskip1em plus10em\penalty0\hbox{}\nobreak\hfill\box\@tempboxa
       % Note that we force a space of at least 1em between title and authorlist, 
       % stretchable to 11em since we may have very short titles but still with an 
       % authorlist too long to go on the same line.  The only linebreak allowed is 
       % directly before the authorlist.
    \else
       % if it is too long, we set the original authorlist text in its own parbox
       \par\vskip1ex\parbox[b]{\TOCentrywidth}{\raggedleft ##2}
    \fi
    }
  \@starttoc{toc}%
  \if@restonecol\twocolumn\fi
  \frontlabel{RMAALastPage}
  \clearpage
  \thispagestyle{empty}
  \cleardoublepage % make sure we start on an odd page
  \endgroup

  %%% ALSO, do the HTML TOC
  %% The file that will be used to make the HTML TOC
  \@starttoc{htc}
  %% All the stuff in this file does nothing in the book style
  \addtocontents{htc}{\protect\HyperTOCmeeting{\@conftitle}{\@reunionplace}{\@reuniondate}}
  \addtocontents{htc}{\protect\HyperTOCeditors{\@editors}}
  \addtocontents{htc}{\protect\HyperTOCpreamble}
  \addtocontents{htc}{\protect\HyperTOCEntry{Table of Contents}{}{\articlename}}

  %% and the ADS file
  \ADS@openfile
  }
%% This is now much simpler
\renewcommand\l@chapter[2]{%TOC entry for articles
  \ifnum \c@tocdepth >\m@ne
    \vskip 1.0em \@plus\p@
    \TOCbox{%
      \parbox[b]{\TOCentrywidth}{#1}\parbox[b]{\TOCpagenumwidth}{\hfill#2}
    }
  \fi} 
\renewcommand\l@part[2]{%TOC entry for sessions: centered, no page number
  \ifnum \c@tocdepth >\m@ne
    \vskip 1.5em \@plus1.5em \goodbreak
    \TOCbox{%
      \parbox[b]{\TOCentrywidth}{\hfil#1\hfil}\hfill
    }
  \fi\nobreak} 

%% Formatting theindex
%% strongly encourage page breaks between index entries
\renewcommand{\@idxitem}{\par\penalty-1000\hangindent.45cm} 
\renewcommand{\indexspace}{\par\penalty-1000\vskip 10\p@ \@plus15\p@ \@minus3\p@\relax}
\renewenvironment{theindex}{%
  %% first, make sure we start on an odd page (06 Jun 2000)
  \clearpage\thispagestyle{empty}\cleardoublepage
  %% this looks a bit like an article but is two columns and without abstract, etc
  \renewcommand{\articlename}{INDEX}
  \title{Author Index}\shortauthor{Author Index}\shorttitle{Author Index}
  \listofauthors{}
  \addcontentsline{toc}{chapter}{\protect\@fiddleauthors{\@TOCtitle}{\rmaaTOCauthors}}%
  \addtocontents{htc}{\protect\HyperTOCEntry{\@rawtitle}{\@rawauthors}{\articlename}}
  \pagestyle{rmaa}
  \thispagestyle{rmaafirstpage}
  \@begofarticlehook
  %% 9pt for journal style
  \if@proceedings\else\small\fi
  \if@newstyle\small\fi
  \twocolumn[\@makeminititle]
  %% vertical spacing (see also definition of \@idxitem and \indexspace
  \parskip\z@ \@plus .6\p@\relax % allow quite a lot of leaway here to avoid breaks inside entries
  %% horizontal spacing, breaking and glue parameters (see also the stuff in the definition 
  %% of \indexauthor
  \parindent\z@ 
  \exhyphenpenalty10000 % don't break on explicit hyphens, it always looks bad
  \rightskip=0pt plus.2em % allow a tiny bit of leeway at the ends of lines
  \spaceskip.3333em plus.5em minus.15em % a lot of leeway between words
  \xspaceskip.5em plus.5em minus.4em % this seems to have no effect
  \let\item\@idxitem
  \renewcommand{\\}{\relax} % temporarily turn off explicit line breaks just in case
  \renewcommand{\altaffilmark}[1]{\relax} % ditto with this
  \renewcommand{\textsuperscript}[1]{\relax} % ditto with this
  }{\label{RMAALastPage}\clearpage} % clearpage necessary so that we know about \shortauthor



%% end definitions for book styles
\fi

%    \end{macrocode}
% In order to do away with plates, we sometimes want a figure to
% appear in color only in the online version. The following command
% can bve used in a figure caption to typeset a note to this effect. 
%    \begin{macrocode}
\newcommand\ColorOnlineOnly{%
  \textsc{Note: this figure is available in color in the electronic
    version of this article, obtainable from}
  \texttt{http://\discretionary{}{}{}www.\discretionary{}{}{}astroscu.\discretionary{}{}{}unam.\discretionary{}{}{}mx/\discretionary{}{}{}\~{}rmaa/}.}
%% Plates
%% Define 2 counters: plate is the number in the whole document, plateloop is used 
%% internally by \outputplates
\newcounter{plate}
\newcounter{plateloop}
%% Currently limited to 15 plates in document
%% when we do a label, we need to write to \newlabel entries to the .aux file, one 
%% that says RMAAPlateX so that it can be referred to in \ouputplates, and one that
%% says articlename:RMAAPlateX so that the dvips program can find it and asociate it
%% with the right article
\newcommand{\@platelabeli}{\oldlabel{RMAAPlate1}\label{RMAAPlate1}}
\newcommand{\@platelabelii}{\oldlabel{RMAAPlate2}\label{RMAAPlate2}}
\newcommand{\@platelabeliii}{\oldlabel{RMAAPlate3}\label{RMAAPlate3}}
\newcommand{\@platelabeliv}{\oldlabel{RMAAPlate4}\label{RMAAPlate4}}
\newcommand{\@platelabelv}{\oldlabel{RMAAPlate5}\label{RMAAPlate5}}
\newcommand{\@platelabelvi}{\oldlabel{RMAAPlate6}\label{RMAAPlate6}}
\newcommand{\@platelabelvii}{\oldlabel{RMAAPlate7}\label{RMAAPlate7}}
\newcommand{\@platelabelviii}{\oldlabel{RMAAPlate8}\label{RMAAPlate8}}
\newcommand{\@platelabelix}{\oldlabel{RMAAPlate9}\label{RMAAPlate9}}
\newcommand{\@platelabelx}{\oldlabel{RMAAPlate10}\label{RMAAPlate10}}
\newcommand{\@platelabelxi}{\oldlabel{RMAAPlate11}\label{RMAAPlate11}}
\newcommand{\@platelabelxii}{\oldlabel{RMAAPlate12}\label{RMAAPlate12}}
\newcommand{\@platelabelxiii}{\oldlabel{RMAAPlate13}\label{RMAAPlate13}}
\newcommand{\@platelabelxiv}{\oldlabel{RMAAPlate14}\label{RMAAPlate14}}
\newcommand{\@platelabelxv}{\oldlabel{RMAAPlate15}\label{RMAAPlate15}}
\newcommand{\@outputplatelabeli}{\oldlabel{RMAAOutputPlate1}}
\newcommand{\@outputplatelabelii}{\oldlabel{RMAAOutputPlate2}}
\newcommand{\@outputplatelabeliii}{\oldlabel{RMAAOutputPlate3}}
\newcommand{\@outputplatelabeliv}{\oldlabel{RMAAOutputPlate4}}
\newcommand{\@outputplatelabelv}{\oldlabel{RMAAOutputPlate5}}
\newcommand{\@outputplatelabelvi}{\oldlabel{RMAAOutputPlate6}}
\newcommand{\@outputplatelabelvii}{\oldlabel{RMAAOutputPlate7}}
\newcommand{\@outputplatelabelviii}{\oldlabel{RMAAOutputPlate8}}
\newcommand{\@outputplatelabelix}{\oldlabel{RMAAOutputPlate9}}
\newcommand{\@outputplatelabelx}{\oldlabel{RMAAOutputPlate10}}
\newcommand{\@outputplatelabelxi}{\oldlabel{RMAAOutputPlate11}}
\newcommand{\@outputplatelabelxii}{\oldlabel{RMAAOutputPlate12}}
\newcommand{\@outputplatelabelxiii}{\oldlabel{RMAAOutputPlate13}}
\newcommand{\@outputplatelabelxiv}{\oldlabel{RMAAOutputPlate14}}
\newcommand{\@outputplatelabelxv}{\oldlabel{RMAAOutputPlate15}}
\newcommand{\@platerefi}{\oldpageref{RMAAPlate1}}
\newcommand{\@platerefii}{\oldpageref{RMAAPlate2}}
\newcommand{\@platerefiii}{\oldpageref{RMAAPlate3}}
\newcommand{\@platerefiv}{\oldpageref{RMAAPlate4}}
\newcommand{\@platerefv}{\oldpageref{RMAAPlate5}}
\newcommand{\@platerefvi}{\oldpageref{RMAAPlate6}}
\newcommand{\@platerefvii}{\oldpageref{RMAAPlate7}}
\newcommand{\@platerefviii}{\oldpageref{RMAAPlate8}}
\newcommand{\@platerefix}{\oldpageref{RMAAPlate9}}
\newcommand{\@platerefx}{\oldpageref{RMAAPlate10}}
\newcommand{\@platerefxi}{\oldpageref{RMAAPlate11}}
\newcommand{\@platerefxii}{\oldpageref{RMAAPlate12}}
\newcommand{\@platerefxiii}{\oldpageref{RMAAPlate13}}
\newcommand{\@platerefxiv}{\oldpageref{RMAAPlate14}}
\newcommand{\@platerefxv}{\oldpageref{RMAAPlate15}}
\newcounter{figplatei}
\newcounter{figplateii}
\newcounter{figplateiii}
\newcounter{figplateiv}
\newcounter{figplatev}
\newcounter{figplatevi}
\newcounter{figplatevii}
\newcounter{figplateviii}
\newcounter{figplateix}
\newcounter{figplatex}
\newcounter{figplatexi}
\newcounter{figplatexii}
\newcounter{figplatexiii}
\newcounter{figplatexiv}
\newcounter{figplatexv}
%% Printed output should use the counter `plate'
\renewcommand{\theplate}{\@arabic\c@plate}
\renewcommand{\theplateloop}{\@arabic\c@plateloop}
%% Command to output all the saved plates in this paper
\newcommand{\outputplates}{%
  \setcounter{plateloop}{0}
  \if@book\else\@endofarticlehook\fi % not entirely sure why, but this
                                     % seems to be needed
  \clearpage % first output the last page of text
  % now set pagestyle for blank even pages and make sure we start on an odd page
  \pagestyle{empty}
  \cleardoublepage
  \onecolumn
  \parindent\z@
  % first put in an extra sheet
  % the main purpose of this is so there is always a blank even page before 
  % each plate for filling in the PS file if the article has an odd number
  % of pages (see splitbook.sh)
  \null\vfil\centerline{PLATES}\vfil\clearpage\pagestyle{rmaaplate}\cleardoublepage
  % output each plate in turn
  \@whilenum\value{plateloop}<\value{plate}\do{%
    \addtocounter{plateloop}{1}
    % restore the figure counter from \@figcounti, \@figcountii, etc
    \setcounter{figure}{\value{figplate\roman{plateloop}}}
    % spew out the plate (in its own group for safety)
    {\csname @plate\roman{plateloop}\endcsname\par}
    }
  \pagestyle{rmaa}
  % reset the page number as if nothing had happened
  \addtocounter{page}{-\value{plate}}
  \addtocounter{page}{-\value{plate}}
  }
\newcommand\figcaption{%
  \def\@captype{figure}
  \renewcommand{\@makecaption}{\@makefigcap}
  \caption}
\newcommand\tabcaption{%
  \def\@captype{table}
  \renewcommand{\@makecaption}{\@maketabcap}
  \caption}

\newcommand{\plate}[2][1]{%
  %% User command to make a plate page: 
  %% Mandatory argument is what goes on the plate - this should not 
  %% include any floating environments (e.g. figure). Instead, \figcaption 
  %% should be used to do the captions. 
  %% Optional argument is the number of figures (i.e. captions) on the plate
  %% page (default is 1). This is necessary for incrementing the figure counter. 
  %%
  %% Note that we have to use a \gdef rather than an \xdef when we 
  %% save the plate output in a command. This means that in the book styles when it 
  %% comes to spewing out the plate, we will have forgotten about any user macros 
  %% that may have been defined. This can be got round by repeating the macros inside the plate.
  %% (Presumably, one would use a \newcommand rather than \renewcommand if they were new.)
  % increment global and local counters
  \addtocounter{plate}{1}
  % write appropriate label to .aux file
  \csname @platelabel\roman{plate}\endcsname
  % define a global macro called \@platei, \@plateii, \@plateiii, \@plateiv, ...
  \expandafter\gdef\csname @plate\roman{plate}\endcsname{%
    \csname @outputplatelabel\roman{plateloop}\endcsname
    \null\vfil #2 \vfil
    \cleardoublepage
    }
  % save the current shortauthor and shorttitle
  % need to be \xdef to take current values of \@shortauthor, etc 
  \expandafter\xdef\csname @shorta\roman{plate}\endcsname{\@shortauthor}
  \expandafter\xdef\csname @shortt\roman{plate}\endcsname{\@shorttitle}
  % save the current value of the figure counter for later use in \outputplates
  \setcounter{figplate\roman{plate}}{\value{figure}}
  % increment the figure counter according to the optional argument
  \addtocounter{figure}{#1}
  }
%% page style for plate pages 
\newcommand{\ps@rmaaplate}{%
  % Plate X (see page Y) in the foot
  \renewcommand\@oddfoot{\hfil\csname @shorta\roman{plateloop}\endcsname\ 
    (see page~\csname @plateref\roman{plateloop}\endcsname)\hfil}%
  \renewcommand\@evenfoot{\relax}%
  % nothing in even heads (blank pages)
  \renewcommand\@evenhead{\relax}
  % short title and page num in the odd heads
  \renewcommand\@oddhead{\hfil\csname @shortt\roman{plateloop}\endcsname%
    \hfil\llap{PLATE~\theplateloop}}
  }





%% Things for the User Guide
\if@guide
%% stolen from doc.sty
\def\MakeShortVerb#1{%
  \expandafter\ifx\csname cc\string#1\endcsname\relax
  \@shortvrbinfo{Made }{#1}%
  \add@special{#1}%
  \expandafter
  \xdef\csname cc\string#1\endcsname{\the\catcode`#1}%
  \begingroup
  \catcode`\~\active  \lccode`\~`#1%
  \lowercase{%
    \global\expandafter\let
    \csname ac\string#1\endcsname~%
    \gdef~{\verb~}}%
  \endgroup
  \global\catcode`#1\active
  \else
  \@shortvrbinfo\@empty{#1 already}%
  \fi}
\def\DeleteShortVerb#1{%
  \expandafter\ifx\csname cc\string#1\endcsname\relax
  \else
  \@shortvrbinfo{Deleted }{#1 as}%
  \rem@special{#1}%
  \global\catcode`#1\csname cc\string#1\endcsname
  \global \expandafter\let \csname cc\string#1\endcsname \relax
  \ifnum\catcode`#1=\active
  \begingroup
  \catcode`\~\active   \lccode`\~`#1%
  \lowercase{%
    \global\expandafter\let\expandafter~%
    \csname ac\string#1\endcsname}%
  \endgroup \fi \fi}
\def\@shortvrbinfo#1#2{%
  \PackageInfo{doc}{%
    #1\expandafter\@gobble\string#2 a short reference
    for \string\verb}}
\def\add@special#1{%
  \rem@special{#1}%
  \expandafter\gdef\expandafter\dospecials\expandafter
  {\dospecials \do #1}%
  \expandafter\gdef\expandafter\@sanitize\expandafter
  {\@sanitize \@makeother #1}}
\def\rem@special#1{%
  \def\do##1{%
    \ifnum`#1=`##1 \else \noexpand\do\noexpand##1\fi}%
  \xdef\dospecials{\dospecials}%
  \begingroup
  \def\@makeother##1{%
    \ifnum`#1=`##1 \else \noexpand\@makeother\noexpand##1\fi}%
  \xdef\@sanitize{\@sanitize}%
  \endgroup}
%% define | as shortcut for \verbatim
\MakeShortVerb{\|}
%% comand to specify arguments to LaTeX commands
\newcommand{\rmArg}[1]{%
  \ensuremath{\langle}{\scriptsize \textsf{\MakeUpperCase{#1}}}\ensuremath{\rangle}}
%\newcommand{\rmArg}[1]{\ensuremath{\langle}\textsc{#1}\ensuremath{\rangle}}
\providecommand{\BibTeX}{{\rmfamily B\kern-.05em%
    \textsc{i\kern-.025em b}\kern-.08em%
    T\kern-.1667em\lower.7ex\hbox{E}\kern-.125emX}}
%% Environment for code examples
\newenvironment{Code}{%
  \begin{list}{}{\setlength{\leftmargin}{0.05\columnwidth}} 
  \item[]}{\end{list}}
\fi
%</package>
%    \end{macrocode}

% \section{User commands for astronomical articles}
% \label{sec:user-commands}
% Some random definitions. These go both into the main class file and
% also into |adshtml.sty|.
% 
% Definitions of the journal name in different incarnations.  
%    \begin{macrocode}
%<*package|adshtml-style>
\newcommand{\RMAA}{RMxAA}
\newcommand{\RMAAlong}{Revista Mexicana de Astronom\'{\i}a y Astrof\'{\i}sica}
\newcommand{\RMAASC}{RMxAC}
\newcommand{\RMAASClong}{\RMAA\ (Serie de Conferencias)}
%    \end{macrocode}
% \begin{macro}{\ibidrule}
% This puts a horizontal rule with a period after it. It is used when
% the author list of a reference is identical to the preceding one. 
%    \begin{macrocode}
\newcommand{\ibidrule}{\rule[-0.03cm]{4.7em}{.01cm}.\ }
%    \end{macrocode}
% \end{macro}
% \begin{macro}{\ThisVolume}
% References to other articles in the same volume can be coded as,
% e.g., |\bibitem{} Bloggs, J. \ThisVolume{bloggs}|.
%    \begin{macrocode}
\newcommand{\ThisVolume}[1]{%
  \@year, \if@proceedings\RMAASC\else\RMAA\fi, \@volume,
  \if@book\oldpageref{#1:RMAAFirstPage}\else 000\fi\ (this volume)}
%    \end{macrocode}
% \end{macro}
% Other stuff that may be useful in ref lists.
%    \begin{macrocode}
\newcommand{\etal}{et al.\@}
%    \end{macrocode}
% Various new commands from AAS macros
%    \begin{macrocode}
\def\acknowledgments{\vskip 3ex plus .8ex minus .4ex}
\let\acknowledgements=\acknowledgments                  % second spelling
\let\jnl@style=\rm
\def\ref@jnl#1{{\jnl@style#1}}
\DeclareRobustCommand\aj{\ref@jnl{AJ}}                   
\DeclareRobustCommand\araa{\ref@jnl{ARA\&A}}             
\DeclareRobustCommand\apj{\ref@jnl{ApJ}}                 
\DeclareRobustCommand\apjl{\ref@jnl{ApJ}}                
\DeclareRobustCommand\apjs{\ref@jnl{ApJS}}               
\DeclareRobustCommand\ao{\ref@jnl{Appl.~Opt.}}           
\DeclareRobustCommand\apss{\ref@jnl{Ap\&SS}}             
\DeclareRobustCommand\aap{\ref@jnl{A\&A}}                
\DeclareRobustCommand\aapr{\ref@jnl{A\&A~Rev.}}          
\DeclareRobustCommand\aaps{\ref@jnl{A\&AS}}              
\DeclareRobustCommand\azh{\ref@jnl{AZh}}                 
\DeclareRobustCommand\baas{\ref@jnl{BAAS}}               
\DeclareRobustCommand\jrasc{\ref@jnl{JRASC}}             
\DeclareRobustCommand\memras{\ref@jnl{MmRAS}}            
\DeclareRobustCommand\mnras{\ref@jnl{MNRAS}}             
\DeclareRobustCommand\pra{\ref@jnl{Phys.~Rev.~A}}        
\DeclareRobustCommand\prb{\ref@jnl{Phys.~Rev.~B}}        
\DeclareRobustCommand\prc{\ref@jnl{Phys.~Rev.~C}}        
\DeclareRobustCommand\prd{\ref@jnl{Phys.~Rev.~D}}        
\DeclareRobustCommand\pre{\ref@jnl{Phys.~Rev.~E}}        
\DeclareRobustCommand\prl{\ref@jnl{Phys.~Rev.~Lett.}}    
\DeclareRobustCommand\pasp{\ref@jnl{PASP}}               
\DeclareRobustCommand\pasj{\ref@jnl{PASJ}}               
\DeclareRobustCommand\qjras{\ref@jnl{QJRAS}}             
\DeclareRobustCommand\skytel{\ref@jnl{S\&T}}             
\DeclareRobustCommand\solphys{\ref@jnl{Sol.~Phys.}}      
\DeclareRobustCommand\sovast{\ref@jnl{Soviet~Ast.}}      
\DeclareRobustCommand\ssr{\ref@jnl{Space~Sci.~Rev.}}     
\DeclareRobustCommand\zap{\ref@jnl{ZAp}}                 
\DeclareRobustCommand\nat{\ref@jnl{Nature}}              
\DeclareRobustCommand\iaucirc{\ref@jnl{IAU~Circ.}}
\DeclareRobustCommand\aplett{\ref@jnl{Astrophys.~Lett.}}
\DeclareRobustCommand\apspr{\ref@jnl{Astrophys.~Space~Phys.~Res.}}
\DeclareRobustCommand\bain{\ref@jnl{Bull.~Astron.~Inst.~Netherlands}}
\DeclareRobustCommand\fcp{\ref@jnl{Fund.~Cosmic~Phys.}}
\DeclareRobustCommand\gca{\ref@jnl{Geochim.~Cosmochim.~Acta}}
\DeclareRobustCommand\grl{\ref@jnl{Geophys.~Res.~Lett.}}
\DeclareRobustCommand\jcp{\ref@jnl{J.~Chem.~Phys.}}      
\DeclareRobustCommand\jgr{\ref@jnl{J.~Geophys.~Res.}}    
\DeclareRobustCommand\jqsrt{\ref@jnl{J.~Quant.~Spec.~Radiat.~Transf.}}
\DeclareRobustCommand\memsai{\ref@jnl{Mem.~Soc.~Astron.~Italiana}}
\DeclareRobustCommand\nphysa{\ref@jnl{Nucl.~Phys.~A}}
\DeclareRobustCommand\physrep{\ref@jnl{Phys.~Rep.}}
\DeclareRobustCommand\physscr{\ref@jnl{Phys.~Scr}}
\DeclareRobustCommand\planss{\ref@jnl{Planet.~Space~Sci.}}
\DeclareRobustCommand\procspie{\ref@jnl{Proc.~SPIE}}
\let\astap=\aap
\let\apjlett=\apjl
\let\apjsupp=\apjs
\let\applopt=\ao
\def\phn{\phantom{0}}
\def\phd{\phantom{.}}
\def\phs{\phantom{\ensuremath{-}}}
\def\phm#1{\phantom{#1}}
%    \end{macrocode}
% All these are redone as robust and using more modern commands where
% possible. The robustness is so they will not get expanded if used in
% abstracts and titles when these are passed to the ADS file. 
% \changes{3.16}{2003/04/07}{
%  Added \cs{Color} command for doing photometric colors.}
%    \begin{macrocode}
%%\def\sun{\hbox{$\odot$}}
\DeclareRobustCommand\sun{\ensuremath{\odot}}
%%\def\earth{\hbox{$\oplus$}}
\DeclareRobustCommand\earth{\ensuremath{\oplus}}
%%\def\lesssim{\mathrel{\hbox{\rlap{\hbox{\lower4pt\hbox{$\sim$}}}\hbox{$<$}}}}
%%\def\gtrsim{\mathrel{\hbox{\rlap{\hbox{\lower4pt\hbox{$\sim$}}}\hbox{$>$}}}}
\DeclareRobustCommand\lesssim{\mathrel{\hbox{\smaller\smaller\rlap{\hbox{\lower3pt\hbox{$\sim$}}}\raise1.5pt\hbox{$<$}}}}
\DeclareRobustCommand\gtrsim{\mathrel{\hbox{\smaller\smaller\rlap{\hbox{\lower3pt\hbox{$\sim$}}}\raise1.5pt\hbox{$>$}}}}
%%\def\lta{\mathrel{\spose{\lower 3pt\hbox{$\mathchar"218$}}
%%     \raise 2.0pt\hbox{$\mathchar"13C$}}}
%%\def\gta{\mathrel{\spose{\lower 3pt\hbox{$\mathchar"218$}}
%%     \raise 2.0pt\hbox{$\mathchar"13E$}}}
\DeclareRobustCommand\sq{\hbox{\rlap{$\sqcap$}$\sqcup$}}
%%\def\arcdeg{\hbox{$^\circ$}}
%%\def\arcmin{\hbox{$^\prime$}}
%%\def\arcsec{\hbox{$^{\prime\prime}$}}
\DeclareRobustCommand\arcdeg{\ensuremath{^\circ}}
\DeclareRobustCommand\arcmin{\ensuremath{^\prime}}
\DeclareRobustCommand\arcsec{\ensuremath{^{\prime\prime}}}
%%\def\fd{\hbox{$.\!\!^{\rm d}$}}
%%\def\fh{\hbox{$.\!\!^{\rm h}$}}
%%\def\fm{\hbox{$.\!\!^{\rm m}$}}
%%\def\fs{\hbox{$.\!\!^{\rm s}$}}
%%\def\fdg{\hbox{$.\!\!^\circ$}}
%%\def\farcm{\hbox{$.\mkern-4mu^\prime$}}
%%\def\farcs{\hbox{$.\!\!^{\prime\prime}$}}
%%\def\fp{\hbox{$.\!\!^{\scriptscriptstyle\rm p}$}}
\DeclareRobustCommand\fd{\ensuremath{.\!\!^\mathrm{d}}}
\DeclareRobustCommand\fh{\ensuremath{.\!\!^\mathrm{h}}}
\DeclareRobustCommand\fm{\ensuremath{.\!\!^\mathrm{m}}}
\DeclareRobustCommand\fs{\ensuremath{\mkern0.5mu.\mkern-5mu^\mathrm{s}\mkern-1mu}}
\DeclareRobustCommand\fdg{\ensuremath{.\!\!^\circ}}
\DeclareRobustCommand\farcm{\ensuremath{.\mkern-4mu^\prime}}
\DeclareRobustCommand\farcs{\ensuremath{\mkern1mu.\mkern-6mu^{\prime\prime}\mkern-1.5mu}}
%%\DeclareRobustCommand\farcs{\ensuremath{.\!\!^{\prime\prime}}}
\DeclareRobustCommand\fp{\ensuremath{.\!\!^{\scriptscriptstyle\rm p}}}
%%\def\micron{\hbox{$\mu$m}}
\DeclareRobustCommand\micron{\ensuremath{\mu\mathrm{m}}}
\let\la=\lesssim                        
\let\ga=\gtrsim
%%\def\case#1#2{\hbox{$\frac{#1}{#2}$}}
%%\def\slantfrac#1#2{\hbox{$\,^#1\!/_#2$}}
%%\def\onehalf{\slantfrac{1}{2}}
%%\def\onethird{\slantfrac{1}{3}}
%%\def\twothirds{\slantfrac{2}{3}}
%%\def\onequarter{\slantfrac{1}{4}}
%%\def\threequarters{\slantfrac{3}{4}}
\DeclareRobustCommand\case[2]{\ensuremath{\frac{#1}{#2}}}
\DeclareRobustCommand\slantfrac[2]{\ensuremath{\,^#1\!/_#2}}
\DeclareRobustCommand\onehalf{\slantfrac{1}{2}}
\DeclareRobustCommand\onethird{\slantfrac{1}{3}}
\DeclareRobustCommand\twothirds{\slantfrac{2}{3}}
\DeclareRobustCommand\onequarter{\slantfrac{1}{4}}
\DeclareRobustCommand\threequarters{\slantfrac{3}{4}}
%%\def\ubvr{\hbox{$U\!BV\!R$}}            
%%\def\ub{\hbox{$U\!-\!B$}}               
%%\def\bv{\hbox{$B\!-\!V$}}               
%%\def\vr{\hbox{$V\!-\!R$}}               
%%\def\ur{\hbox{$U\!-\!R$}}
%    \end{macrocode}
% New commands for ``photometric'' colors 
%    \begin{macrocode}
\DeclareRobustCommand\ubvr{\ensuremath{U\!BV\!R}}            
\DeclareRobustCommand\Color[2]{%
  \ensuremath{%
    #1\!\ifx#1V\!\fi-\!#2
  }}
\DeclareRobustCommand\ub{\Color{U}{B}}
\DeclareRobustCommand\bv{\Color{B}{V}}
\DeclareRobustCommand\vr{\Color{V}{R}}
\DeclareRobustCommand\ur{\Color{U}{R}}
%% New version of the \ion command, that actually works! 
%    \end{macrocode}
% \begin{macro}{\ion}
% This version of the |\ion| command allows you to write either
% |\ion{N}{4}| or |\ion{N}{iv}| in order to get N~IV. 
%    \begin{macrocode}
\DeclareRobustCommand\ion[2]{%
  \mbox{#1\kern0.2em%
    \smaller\rmfamily%
    \edef\@tempa{\@car#2\@nil}%
    \ifcat1\@tempa%
    \@Roman{#2}%
    \else%
    \uppercase{#2}%
    \fi}}
\newcounter{IonStage}
\DeclareRobustCommand\plainion[2]{\setcounter{IonStage}{#2}#1 \Roman{IonStage}}
%%\def\ion#1#2{\ensuremath{\mathrm{#1}\;}{\protect\small\rm\@Roman{#2}}}
%    \end{macrocode}
% \end{macro}
%    \begin{macrocode}
\DeclareRobustCommand\nodata{$\cdots$}
%    \end{macrocode}
% These are needed by the |\ion| command. Isn't there an easier way of
% doing this?
%    \begin{macrocode}
\let\@small=\small
\let\@normalsize=\normalsize
\let\@footnotesize=\footnotesize
\let\@scriptsize=\scriptsize
\let\@tiny=\tiny
\let\@large=\large
\let\@Large=\Large
\let\@LARGE=\LARGE
\let\@huge=\huge
\let\@Huge=\Huge
\newcommand{\smaller}{\small}
\renewcommand{\normalsize}{\renewcommand{\smaller}{\small}\@normalsize}
\renewcommand{\small}{\renewcommand{\smaller}{\footnotesize}\@small}
\renewcommand{\footnotesize}{\renewcommand{\smaller}{\scriptsize}\@footnotesize}
\renewcommand{\scriptsize}{\renewcommand{\smaller}{\tiny}\@scriptsize}
\renewcommand{\tiny}{\renewcommand{\smaller}{\tiny}\@tiny}
\renewcommand{\large}{\renewcommand{\smaller}{\normalsize}\@large}  
\renewcommand{\Large}{\renewcommand{\smaller}{\large}\@Large}  
\renewcommand{\LARGE}{\renewcommand{\smaller}{\Large}\@LARGE}  
\renewcommand{\huge}{\renewcommand{\smaller}{\LARGE}\@huge}  
\renewcommand{\Huge}{\renewcommand{\smaller}{\huge}\@Huge}  
%</package|adshtml-style>
%    \end{macrocode}


% \section{Other files}

% \subsection{The driver and make files}
% \label{sec:driver_make}
% These are put here for completeness, although they still need to be
% distributed separately since otherwise the user won't be able to
% unpack this file into its constituent parts.
% 
% The |Makefile| has two principal rules: |docs|, and |macros|. The
% |docs| rule produces a PDF and gzipped PS file of the documentation.
% I still haven't perfected the combination of rules required to get
% it to run \LaTeX{} enough times in all circumstances. The |macros|
% rule runs \LaTeX{} on |rmaa.ins| in order to unpack |rmaa.cls| and
% all the other little files (including the |Makefile| and the
% |rmaa.ins| files themselves\dots 
% 
% First set up lists of various types of file. 
%    \begin{macrocode}
%<*makefile>
DOCFILES=rmaa.pdf
MACRO_LTX_FILES=rmaa.cls rmaacite.sty adshtml.sty
SHELLSCRIPTS=makeads.sh
MACRO_OTHER_FILES=adshtml.perl sample.latex2html-init
MAKEFILE=Makefile
MACROFILES=$(MACRO_LTX_FILES) $(SHELLSCRIPTS) $(MACRO_OTHER_FILES)         
SKELETONFILES=ADS.tex 
%    \end{macrocode}
% Install the files read only into the package directory, specified via a path relative to the source directory. \changes{4.0}{2019/11/14}{Reorganize folder layout so that hardly anything is in the top level.  That way, things are clearer on github.}
%    \begin{macrocode}
INSTALL=install 
INSTALL_DIR=../packages
%    \end{macrocode}
% The default rule is to make everything. Installation must be
% explicitly asked for.
%    \begin{macrocode}
all: docs macros

install: install-docs install-macros

%    \end{macrocode}
% Most files are installed read-only to prevent editing. Skeleton
% files must be writeable. Must be careful not to overwrite the
% Makefile in the parent directory. 
%    \begin{macrocode}
install-macros: macros
        $(INSTALL) -m 444 $(MACRO_LTX_FILES) $(MACRO_OTHER_FILES) $(INSTALL_DIR)
        $(INSTALL) -m 555 $(SHELLSCRIPTS) $(INSTALL_DIR)
        $(INSTALL) -m 644 $(SKELETONFILES) $(INSTALL_DIR)

install-docs: docs
        $(INSTALL) rmaa.pdf $(INSTALL_DIR)/rm-fulldocs.pdf

%    \end{macrocode}
% Make the PDF and gzipped PS files. 
%    \begin{macrocode}
docs: $(DOCFILES)

%    \end{macrocode}
% Make everything else if either the |.dtx| or |.ins| files have
% changed. 
%    \begin{macrocode}
macros: rmaa.dtx rmaa.ins
%    \end{macrocode}
% Ensure that all the macro files can be written to then recreate
% them. The |touch| is there in case any of the files don't exist
% yet. 
%    \begin{macrocode}
        touch $(MACROFILES) $(MAKEFILE)
        chmod u+w $(MACROFILES) $(MAKEFILE) rmaa.ins
        latex "%%
%% This is file `rmaa.ins',
%% generated with the docstrip utility.
%%
%% The original source files were:
%%
%% rmaa.dtx  (with options: `ins')
%% 
%% PLEASE DO NOT EDIT THIS FILE! It is generated automatically from
%% the file rmaa.dtx using the accompanying Makefile and rmaa.ins
%% Any hand-edited changes will be lost
%% 
%% A LaTeX document class for the RMxAA
%% 
%% Author: William Henney (w.henney@crya.unam.mx)
%% 
%% Latest version available from: http://www.crya.unam.mx/rmaa
%% 
%% Copyright (C) 1999-2010 Instituto de Astronomia, UNAM, Mexico
%% 
%% This program is free software; you can redistribute it and/or
%% modify it under the terms of the GNU General Public License
%% as published by the Free Software Foundation; either version 2
%% of the License, or (at your option) any later version.
%% 
%% This program is distributed in the hope that it will be useful,
%% but WITHOUT ANY WARRANTY; without even the implied warranty of
%% MERCHANTABILITY or FITNESS FOR A PARTICULAR PURPOSE.  See the
%% GNU General Public License for more details.
%% 
%% You should have received a copy of the GNU General Public License
%% along with this program; if not, write to the Free Software
%% Foundation, Inc., 59 Temple Place - Suite 330, Boston, MA  02111-1307, USA.





%%
%% Installation batch file for the RMAA LaTeX package
%%
\def\batchfile{rmaa.ins}
\input docstrip.tex
\preamble

PLEASE DO NOT EDIT THIS FILE! It is generated automatically from
the file rmaa.dtx using the accompanying Makefile and rmaa.ins
Any hand-edited changes will be lost

A LaTeX document class for the RMxAA

Author: William Henney (w.henney@crya.unam.mx)

Latest version available from: http://www.crya.unam.mx/rmaa

Copyright (C) 1999-2010 Instituto de Astronomia, UNAM, Mexico

This program is free software; you can redistribute it and/or
modify it under the terms of the GNU General Public License
as published by the Free Software Foundation; either version 2
of the License, or (at your option) any later version.

This program is distributed in the hope that it will be useful,
but WITHOUT ANY WARRANTY; without even the implied warranty of
MERCHANTABILITY or FITNESS FOR A PARTICULAR PURPOSE.  See the
GNU General Public License for more details.

You should have received a copy of the GNU General Public License
along with this program; if not, write to the Free Software
Foundation, Inc., 59 Temple Place - Suite 330, Boston, MA  02111-1307, USA.
\endpreamble
\askforoverwritefalse
\generate{%
  \file{rmaa.cls}{\from{rmaa.dtx}{package}}
  \file{rmaa.ins}{\from{rmaa.dtx}{ins}}
  \file{rmaacite.sty}{\from{rmaa.dtx}{rmaacite}}
  \file{adshtml.sty}{\from{rmaa.dtx}{adshtml-style}}
  \file{ADS.tex}{\from{rmaa.dtx}{ads-container}}
  }
\generate{\nopreamble\nopostamble%
  \file{Makefile}{\from{rmaa.dtx}{makefile}}
  \file{adshtml.perl}{\from{rmaa.dtx}{adshtml-perl}}
  \file{makeads.sh}{\from{rmaa.dtx}{ads-shell-script}}
  \file{sample.latex2html-init}{\from{rmaa.dtx}{latex2html-init}}
  }

%% Local Variables:
%% mode: LaTeX
%% mode: Font-lock
%% End:









\endinput
%%
%% End of file `rmaa.ins'.
"
%    \end{macrocode}
% \TeX{} automatically adds the |.tex| extension to extensionless
% files. We fix this and take the opportunity to convert
% back to TABs all that it has lovingly converted to spaces. 
%    \begin{macrocode}
        cp -f $(MAKEFILE) $(MAKEFILE).BCK && unexpand $(MAKEFILE).tex >| $(MAKEFILE) && rm $(MAKEFILE).tex
%    \end{macrocode}
% Make the shell scripts executable and everything unwriteable (so
% they don't get hand-edited by mistake). However, undo this for the
% files that the user is \emph{supposed} to modify. 
%    \begin{macrocode}
        chmod a+x $(SHELLSCRIPTS)
        chmod a-w $(MACROFILES) $(MAKEFILE)
        chmod u+w $(SKELETONFILES)

%    \end{macrocode}
% We now use PDF\LaTeX{} for this.
%    \begin{macrocode}
rmaa.idx: rmaa.dtx
        @echo "REGENERATING .idx FILE"
        pdflatex "\nonstopmode% \def\fileversion{4.0} \def\filedate{2019/11/12}
% \MakeShortVerb{\|}
% \title{A \LaTeX{} document class for the Revista Mexicana de
% Astronom\'\i a y Astrof\'\i sica\thanks{This file has version
% \fileversion{} last revised \filedate{}.}}
% \author{William J. Henney\thanks{w.henney@irya.unam.mx}}
% \date{\filedate}
% \maketitle
% \paragraph{Abstract}
% This document describes the \texttt{rmaa} \LaTeX{}
% document class for the journal ``Revista Mexicana de
% Astronom\'\i a y Astrof\'\i sica''. The macros can be used for 
% \begin{enumerate}
% \item Preparing author manuscripts for the journal (including
% conference proceedings to be published in the ``Serie de
% Conferencias'')
% \item Typesetting an entire volume of the journal by inputting the
% individual article files from a container file. Facilities are
% included for automatic generation of TOC and Author Index. 
% \end{enumerate}
% 
% \tableofcontents
% \section{The documentation driver file}
% This is the part that generates the documentation that you may or
% may not be reading, which is done with the help of the Makefile
% discussed in section~\ref{sec:driver_make}. 
%    \begin{macrocode}
%<*driver>
\documentclass{ltxdoc}
\usepackage{hyperref}
\usepackage{palatino}
\EnableCrossrefs \CodelineIndex \RecordChanges
\setlength{\parindent}{0pt}
\def\MacroFont{\footnotesize\tt}
\begin{document} \DocInput{rmaa.dtx} \PrintIndex \PrintChanges \end{document}
%</driver>
%    \end{macrocode}

% \section{The implementation}
% 
% \changes{3.24}{2007/09/14}{Added user access to the macro version
% and date via \cs{rmVersion} and \cs{rmDate}}
% 
%    \begin{macrocode}
%<*package>
\newcommand{\rmaa@version}{3.26}
\newcommand{\rmaa@date}{2011/08/15}
\newcommand\rmVersion{\rmaa@version}
\newcommand\rmDate{\rmaa@date}
\NeedsTeXFormat{LaTeX2e}[1995/06/01]
%    \end{macrocode}
% \changes{4.0}{2019/11/02}{Change abstract to appear before
% resumen. Change RevMexAA to RMxAA for consistency with ADS.}
%
% \changes{3.22}{2007/03/22}{Allow proceedings articles to have
%   received/accepted dates} \changes{3.21}{2006/05/15}{Changes for
%   Silvia - abstracts are now shipped out to ADS.}
% \changes{3.19}{2006/01/19}{Changes for Silvia - remove adresses from
%   end, new abstract style} \changes{2.2beta}{2001/10/30}{ Fixed bug
%   whereby twocolumn option was not passed through to book
%   class. Finally sorted out the ``label on last page'' issue. }
% \changes{2.1}{2001/01/23}{ Fixed some inconsistencies in
%   \cs{nocite} and \cs{bibitem} uniqification Fixed the definition of
%   \cs{ion} so it will work in math mode and in contexts where the
%   typesize is different from normal (e.g., title, fig captions).
%   Modified the definition of \cs{@caption}, so we don't write any
%   list-of-figures data (we never use it anyway, and it just causes
%   problems with fragile commands in fig captions). }
% \changes{2.0}{2000/11/30}{ Added `newstyle' option, which should be
%   used in all new proceedings volumes (wider abstract, smaller
%   typeface in fig captions, etc.) Added commands for formatting
%   table captions. New demonstation document rmpostsamp.tex for
%   1-page posters. This will be the last release under a PRI
%   government (for a while...)}  \changes{1.9}{2000/06/06}{
%   Corrected position of page number in footer of articles that open
%   on even pages. Ensured that first article in each session opens on
%   an odd page, as do list of abstracts and author index. Added extra
%   space in TOC before the backmatter. Added command for editors,
%   \cs{BoldSessionTitles}, that makes produces bold titles for
%   session in the TOC.}  \changes{1.8}{2000/05/16}{ Fixed bug in
%   last page label with \cs{suppressfulladdresses} }
% \changes{1.7}{2000/05/11}{ Fixed more little things. Page labels
%   now work, even if you have no bibliography. You can output the
%   full addresses by hand with the \cs{outputfulladdresses}
%   command. Editors' Guide updated but still not finished. Revamped
%   splitbook.sh now produces PS file of each article for web page and
%   also an EPS file of each page for sending to the printers.}
% \changes{1.6}{2000/04/15}{ More cleaning up: full stops after
%   titles in index, running headers and footers in frontmatter,
%   fontsize in title of index, etc} \changes{1.5}{2000/04/15}{
%   Eliminated superscripts from titles in TOC and index}
% \changes{1.4}{2000/04/14}{ Added copyright and license info (GPL)}
% \changes{1.3}{2000/03/01}{ Corrected \cs{figcaption} to do the
%   caption formatting right Corrected font size of sec headings in
%   front matter} \changes{1.2}{2000/02/16}{ Added \cs{plate}
%   command} \changes{1.1}{2000/02/14}{ Added manuscript option,
%   corrected bugs} \changes{1.0}{2000/02/08}{ Now does a whole book
%   - lots of stuff tidied up} \changes{0.7}{1999/10/04}{ Increased
%   spacing between sections in single column styles. Suppressed
%   floats at bottom of pages (didn't work though) }
% \changes{0.6}{1999/09/24}{ Tweaked \cs{adjustfinalcols} (still not
%   perfect) Added proceedings option Added onecolumn/twocolumn
%   options (orthogonal to journal/proceedings) }
% \changes{0.5}{1999/09/23}{ Added \cs{fulladdresses} and
%   \cs{adjustfinalcols} (just about working)}
% \changes{0.4}{1999/09/22}{ Minor stylistic changes}
% \changes{0.3}{1999/06/28}{ Better mechanism for alternative
%   affiliations Added `guide' option for user guide}
% \changes{0.2}{1999/05/27}{ All basic functionality present, I
%   think} \changes{0.1}{1999/05/26}{ Just about working - still
%   missing Fig/Tab/Refs} \changes{0.0}{1999/05/21}{ So bored today,
%   decided to rewrite RMAA LaTeX macros from scratch}
%    \begin{macrocode}
\ProvidesClass{rmaa}[\rmaa@date\space v\rmaa@version\space RMxAA article
class]
%    \end{macrocode}
% \subsection{Option Handling}
% \label{sec:options}
% Declarations that are needed in options
% \changes{3.05}{2003/10/22}{Added natbib option}
% \changes{3.06}{2003/01/03}{Natbib option turned on permanently}
% \changes{3.07}{2003/01/11}{Added proof option}
% \changes{3.14}{2003/01/11}{Option to turn off natbib}
% \changes{3.18}{2003/10/22}{trying out doctex mode} 
% \changes{3.23}{2007/09/09}{Manuscript is now on by default.}
% \changes{3.26}{2011/08/15}{New option for review articles}
%    \begin{macrocode}
\newif\if@reviewarticle
\@reviewarticlefalse
\newif\if@proof
\@prooffalse
\newif\if@journal
\@journaltrue
\newif\if@printdates
\@printdatesfalse
\newif\if@proceedings
\@proceedingsfalse
\newif\if@manuscript
\@manuscripttrue
\newif\if@bigfont
\@bigfonttrue
\newif\if@guide
\@guidefalse
\newif\if@preprint
\@preprintfalse
\newif\if@rmtwocol
\@rmtwocoltrue
\newif\if@biblist
\@biblistfalse
\newif\if@book
\@bookfalse
\newif\if@natbib
\@natbibtrue
\newif\if@rmcite
\@rmcitefalse
\newif\if@titlevisible
\@titlevisibletrue
\newif\if@abstractpresent
\@abstractpresentfalse
\newif\if@printbanner
\@printbannerfalse
\newif\if@suppressfulladdresses
\@suppressfulladdressesfalse
\newif\if@newstyle
\@newstylefalse
\newif\if@debug
\@debugfalse
\newif\if@boldkeywords
\@boldkeywordsfalse
%    \end{macrocode}
% The option handling---mainly setting various booleans.
% \changes{2.2beta11}{2001/11/26}{Made twocol and
% newstyle the default for proceedings.}
% \changes{3.23}{2007/09/09}{Manuscript is now 10 pt instead of 12pt}
% \changes{3.23}{2007/09/10}{Full addresses are now automatically
% suppressed for proceedings style.}
% \changes{3.24}{2007/09/14}{New option \texttt{boldkeywords} that
% does the keywords in the old ugly way. Now off by default.}
%    \begin{macrocode}
\DeclareOption{manuscript}{\typeout{Manuscript for referee} 
  \@manuscripttrue\@journalfalse\@rmtwocolfalse
  \PassOptionsToClass{10pt}{article}}
\DeclareOption{bigfont}{\typeout{Big font manuscript requested by Gary
  Ferland} 
  \@manuscripttrue\@journalfalse\@rmtwocolfalse\@bigfonttrue
  \PassOptionsToClass{11pt}{article}}
\DeclareOption{journal}{\typeout{Main journal option}%
  \@journaltrue\@rmtwocoltrue\@printdatestrue\@manuscriptfalse
  \PassOptionsToClass{10pt}{article}}
\DeclareOption{guide}{\typeout{User guide option}%
  \@guidetrue\@journalfalse\@rmtwocoltrue\@manuscriptfalse
  \PassOptionsToClass{10pt}{article}}
\DeclareOption{preprint}{\typeout{Preprint option}%
  \@preprinttrue\@manuscriptfalse\PassOptionsToClass{10pt}{article}}
\DeclareOption{proof}{\typeout{Author proof option}%
  \@prooftrue\@manuscriptfalse}
\DeclareOption{proceedings}{\typeout{Proceedings option}%
  \@proceedingstrue\@journalfalse\@rmtwocoltrue\@newstyletrue
  \@suppressfulladdressestrue
  \@printdatesfalse\@manuscriptfalse\PassOptionsToClass{10pt}{article}} 
\DeclareOption{printdates}{\typeout{Force printing of received/accepted dates}%
  \@printdatestrue} 
\DeclareOption{book}{\typeout{Whole issue option}%
  \@booktrue} 
\DeclareOption{nonatbib}{\typeout{Turn off natbib support (required
    for compatibility with old rmaacite package).}%
  \@natbibfalse} 
\DeclareOption{newstyle}{\typeout{Use new layout style}%
  \@newstyletrue}
\DeclareOption{oldstyle}{\typeout{Use old layout style}%
  \@newstylefalse}
\DeclareOption{debug}{\typeout{Loads of annoying messages}%
  \@debugtrue}
%    \end{macrocode}
% These two need to be declared explicitly so that |\ExecuteOptions|
% actually does something.
%    \begin{macrocode}
\DeclareOption{twocolumn}{\@rmtwocoltrue}
\DeclareOption{onecolumn}{\@rmtwocolfalse}
\DeclareOption{twoside}{\PassOptionsToClass{\CurrentOption}{article}}
\DeclareOption{boldkeywords}{\@boldkeywordstrue}
%    \end{macrocode}
% Bold, upper-case keywords used to be the default. Now they must be
% turned on explicitly with the \texttt{boldkeywords} option. 
%    \begin{macrocode}
\DeclareOption{boldkeywords}{\@boldkeywordstrue}
%    \end{macrocode}
% Pass any other option along to the \texttt{article} class.
%    \begin{macrocode}
\DeclareOption*{\PassOptionsToClass{\CurrentOption}{article}}
%    \end{macrocode}
% Default options are to be two-sided and based on the standard
% \LaTeX{} \texttt{article} class,
% unless we are processing the entire book, in which case we base
% ourselves on the standard \texttt{book} class. Default is to use
% manuscript option unless otherwise specified. 
%    \begin{macrocode}
\ExecuteOptions{twoside,manuscript} \ProcessOptions\relax
\if@book
\if@rmtwocol\LoadClass[twocolumn]{book}\else\LoadClass{book}\fi
\newcommand{\abstract}{\relax}
\newcommand{\@maketitle}{\relax}
\else
\if@rmtwocol\LoadClass[twocolumn]{article}\else\LoadClass{article}\fi
\fi
%    \end{macrocode}
% Attempt to use natbib for citations 02 Jan 2003. 
% \changes{3.05}{2003/01/02}{If the natbib option is used, then set
% things up so that natbib's version of the bibliography is identical
% to the standard one.}
% \changes{3.09}{2003/01/23}{Explicitly require a recent version of
% natbib, so that we get a decent error message if the package is too
% old.} 
% \changes{3.23}{2007/09/09}{Now allow Bibtex to work properly} 
% \changes{3.25}{2009/10/01}{Change \cs{newcommand} tp \cs{def} when setting \cs{bibfont}. This allows it to work with all versions of natbib.} 
% \changes{3.25}{2010/02/06}{Updated copyright year.} 
% \changes{3.25}{2010/02/06}{Updated email and web addresses to \texttt{crya.unam.mx}.} 
%    \begin{macrocode}
\if@natbib
\RequirePackage{natbib}[1999/02/25]
\bibpunct{(}{)}{;}{a}{}{,}
\def\bibfont{\small}
\renewcommand\bibsection{\section*{REFERENCES}}
\setlength\bibsep{0pt}
\setlength\bibhang{1.5em}
\bibliographystyle{rmaa}
\fi
%    \end{macrocode}
% We want the \texttt{graphicx} package for the extended syntax of
% |\includegraphics|. We want the \texttt{ifthen} package for using in
% these macros (although we don't much yet). We explicitly load
% \texttt{babel} since some \LaTeX{} distributions messed up the
% hyphenation if we didn't.  
%    \begin{macrocode}
\RequirePackage{graphicx}
\RequirePackage{ifthen}
\IfFileExists{babel.sty}{%
%   \let\@restoreindent=\@afterindentfalse
  \RequirePackage[english]{babel}
%   \@restoreindent
  %% reinstate the definition of \@roman so as not to break .idx files
  \@ifundefined{org@@roman}{\relax}{\renewcommand{\@roman}[1]{\org@@roman{##1}}}
  }{%
  \typeout{Package babel.sty not found - don't worry about this}}
%    \end{macrocode}
% Try to load the \texttt{dcolumn} and \texttt{booktabs} packages if
% they can be found. These can produce nicer-looking tables. 
%    \begin{macrocode}
\IfFileExists{dcolumn.sty}{%
  \RequirePackage{dcolumn}
  }{%
  \typeout{Package dcolumn.sty not found - you will not be able to
    align columns on decimal points in tables}}
\IfFileExists{booktabs.sty}{%
  \RequirePackage{booktabs}
  }{%
  \typeout{Package booktabs.sty not found - your tables will not look
    as pretty as they ought too.}
  }
\IfFileExists{makeidx.sty}{%
  \RequirePackage{makeidx}
}{%
  \typeout{Package makeidx.sty not found - Author Index may not work.}
  }
%% commented out for moment since requires more modern version of LaTeX
%% \RequirePackage[latin1]{inputenc}
%    \end{macrocode}
% \changes{4.0}{2019/11/13}{Make \texttt{hyperref} and
% \texttt{etoolbox} into hard requirements, since we now use them
% extensively.}  We use hyperref to get internal and external links.
% But I don't like the colored boxes so I turn those off.  We use
% etoolbox for patching and robustifying at the moment.  In the
% future, I want to port a lot of other stuff over to it.  For
% instance, replace the ifthen package and the TeX ifs too.
%
% We target the latest version (1.7) of PDF, rather than the default
% of 1.5.  This reduces the chance of getting warnings from included
% graphics.  This is not strictly related to hyperref, but we put it
% here anyway.
%    \begin{macrocode}
\RequirePackage{etoolbox}
\RequirePackage{hyperref}
\hypersetup{hidelinks}
\pdfminorversion=7
%    \end{macrocode}
% We patch the \cs{@combinedblfloats} following a suggestion by Heiko
% Oberdiek at \url{https://tex.stackexchange.com/a/249743/3843}, which
% avoids problems when hyperlinks are split across
% pages. \changes{4.0}{2019/11/12}{Fix for issue where breaking
% hyperlink across pages gave an error.}
%    \begin{macrocode}
\urlstyle{same}
\newcommand\email{\begingroup \urlstyle{rm}\Url}
\patchcmd\@combinedblfloats{\box\@outputbox}{\unvbox\@outputbox}{}{%
  \errmessage{\noexpand\@combinedblfloats could not be patched}%
}%
%    \end{macrocode}
% \subsection{Other setup stuff}\label{sec:setup}
% Some declarations: 
%    \begin{macrocode}
\newif\if@needaddresses
\@needaddressestrue
\newif\if@frontmatter
\@frontmatterfalse
\newif\if@inpreamble
\@inpreambletrue
\newif\if@firstauthor
\@firstauthortrue
\newif\if@firstaffil
\@firstaffiltrue
\newif\if@firstaltaffil
\@firstaltaffiltrue
\newif\if@firstkwd
\@firstkwdtrue
\newif\if@noexpandADS
\@noexpandADSfalse
\newif\if@sessiontitlepages
\@sessiontitlepagesfalse
%    \end{macrocode}
% \begin{macro}{\@rmaaEmpty}
% \begin{macro}{\@rmaaArg}
% These are used in various commands to check if arg was left blank
%    \begin{macrocode}
\newcommand{\@rmaaEmpty}{}
\newcommand{\@rmaaArg}{}
%    \end{macrocode}  
% \end{macro}
% \end{macro}
%
% \subsection{Compatibility with old \LaTeX{} versions}
% \label{sec:compat}
% This package is based on the 1998/12/01 release of \LaTeX{}.
% Older releases may not define all the commands that we need.
% Luckily, we can use |\providecommand| to define these if necessary. 
% The following commands are not defined in the 1994/06/01 release. 
%
% \begin{macro}{\MakeUpperCase}
% \begin{macro}{\textsuperscript}
% \begin{macro}{\textbackslash}
% User commands:
%    \begin{macrocode}
\providecommand{\MakeUpperCase}[1]{\uppercase{#1}}
\providecommand{\textsuperscript}[1]{\ensuremath{^\textrm{#1}}}
\providecommand{\textbackslash}{\ensuremath{\backslash}}
%    \end{macrocode}  
% \end{macro}
% \end{macro}
% \end{macro}
% \begin{macro}{\protected@xdef}
% \begin{macro}{\restore@protect}
% \begin{macro}{\@unexpandable@protect}
% \begin{macro}{\hb@xt@}
% \LaTeX{} internal commands:
%    \begin{macrocode}
\providecommand{\protected@xdef}{%
        \let\@@protect\protect
        \let\protect\@unexpandable@protect
        \afterassignment\restore@protect
        \xdef}
\providecommand{\restore@protect}{\let\protect\@@protect}
\providecommand{\@unexpandable@protect}{\noexpand\protect\noexpand}
\providecommand{\hb@xt@}{\hbox to}
%    \end{macrocode}
% \end{macro}
% \end{macro}
% \end{macro}
% \end{macro}
%
% \subsection{Page layout}\label{sec:page-layout}
%
% The dimensions of the text were defined by Elizabeth. Text width is
% 169\,mm with equal margins on the two sides: 
%    \begin{macrocode}
\setlength{\textwidth}{169mm}
\setlength{\@tempdima}{\paperwidth}
\addtolength{\@tempdima}{-\textwidth}
\setlength{\oddsidemargin}{0.5\@tempdima}
\addtolength{\oddsidemargin}{-1in}
\setlength{\evensidemargin}{\oddsidemargin}
%    \end{macrocode}
% Text height is 225\,mm, top margin is 14\,mm. 
%    \begin{macrocode}
\setlength{\textheight}{225mm} 
\setlength{\topmargin}{-1in} 
\addtolength{\topmargin}{14mm}
%    \end{macrocode}
% Space between columns in two-column styles, header, footer,
% paragraph indentation:
%    \begin{macrocode}
\setlength{\columnsep}{9.5mm}
\setlength{\headsep}{0.7cm}
\setlength{\footskip}{0.75cm}
\setlength{\parindent}{1.5em}
%    \end{macrocode}
% 
% \begin{macro}{\@set@scalable@length}
% Try out a new scheme for dealing with various spacing dimensions and
% allowing the user to easily modify them. The command
% \begin{quote}
% |\set@scalable@length|\marg{name}\marg{defaultlength}\marg{stretch}\marg{shrink}
% \end{quote}
% does a 
% |\setlength|\marg{\cs{name}}\marg{defaultlength plus stretch minus shrink}
% and also defines the user command
% \textbackslash\meta{name}|Scale|\marg{factor} that can 
% be used to rescale length by \meta{factor}. Additionally, now
% defines the package command
% \textbackslash\meta{name}|@set@scale|\marg{FACTOR} 
% that also rescales the length by \meta{FACTOR}, so that the length
% becomes \meta{factor} $\times$ \meta{FACTOR} $\times$
% \meta{defaultlength}.   
%    \begin{macrocode}
\newcommand\@set@scalable@length[4]{%
  %% arg #1: NAME of length without leading `\'
  %%
  %% arg #2: default VALUE  
  \expandafter\newlength\csname @#1@default\endcsname
  \expandafter\setlength\csname @#1@default\endcsname{#2}%
  \expandafter\newlength\csname @@#1@@default\endcsname
  \expandafter\setlength\csname @@#1@@default\endcsname{#2}%
  %% arg #3: default STRETCH  
  \expandafter\newlength\csname @#1@stretch\endcsname
  \expandafter\setlength\csname @#1@stretch\endcsname{#3}%
  \expandafter\newlength\csname @@#1@@stretch\endcsname
  \expandafter\setlength\csname @@#1@@stretch\endcsname{#3}%
  %% arg #4: default SHRINK  
  \expandafter\newlength\csname @#1@shrink\endcsname
  \expandafter\setlength\csname @#1@shrink\endcsname{#4}%
  \expandafter\newlength\csname @@#1@@shrink\endcsname
  \expandafter\setlength\csname @@#1@@shrink\endcsname{#4}%
  %%
  %% User-defined scale \@NAME@scale is unity by default
  \expandafter\gdef\csname @#1@scale\endcsname{1}%
  %% Package-defined scale \@NAME@SCALE is unity by default
  \expandafter\gdef\csname @#1@SCALE\endcsname{1}%
  %%
  %% internal command \@reset@NAME should expand to 
  %% \setlength\@@NAME@@default{\@NAME@SCALE\@NAME@default}
  %% \setlength\NAME{\@NAME@scale\@@NAME@@default}
  %% in other words: scale x SCALE x default 
  \expandafter\gdef\csname @reset@#1\endcsname{%
    \expandafter\setlength\csname @@#1@@default\endcsname{%
      \csname @#1@SCALE\endcsname \csname @#1@default\endcsname
      }
    \expandafter\setlength\csname @@#1@@stretch\endcsname{%
      \csname @#1@SCALE\endcsname \csname @#1@stretch\endcsname
      }
    \expandafter\setlength\csname @@#1@@shrink\endcsname{%
      \csname @#1@SCALE\endcsname \csname @#1@shrink\endcsname
      }
    \expandafter\setlength\csname #1\endcsname{%
      \csname @#1@scale\endcsname \csname @@#1@@default\endcsname
      \@plus
      \csname @#1@scale\endcsname \csname @@#1@@stretch\endcsname
      \@minus
      \csname @#1@scale\endcsname \csname @@#1@@shrink\endcsname
      }}
  %%
  %% User command \NAMEScale should expand to 
  %% {\renewcommand\@NAME@scale{#1}\@reset@NAME}
  \expandafter\gdef\csname #1Scale\endcsname##1{%
    \expandafter\gdef\csname @#1@scale\endcsname{##1}%
    \csname @reset@#1\endcsname}
  %%
  %% Package command \NAME@set@SCALE should expand to 
  %% {\renewcommand\@NAME@SCALE{#1}\@reset@NAME}
  \expandafter\gdef\csname #1@set@SCALE\endcsname##1{%
    \expandafter\gdef\csname @#1@SCALE\endcsname{##1}%
    \csname @reset@#1\endcsname}
  %%
  %% initialize at default values with \@reset@NAME
  \csname @reset@#1\endcsname
  }
%    \end{macrocode}
% \end{macro}
% \begin{macro}{\parskip}
% First use of this new infrastructure: set interparagraph space to be
% zero but with a small amount of stretchiness.
%    \begin{macrocode}
%   \ifthenelse{\boolean{@manuscript}}{%
%     \@set@scalable@length{parskip}{0.5\baselineskip}{0.3\p@}{0.1\p@}
%   }{
    \@set@scalable@length{parskip}{0\p@}{0.3\p@}{0\p@}
%   }
%    \end{macrocode}
% \end{macro}
% 
% \begin{macro}{\abovesectionskip}
% \changes{4.0}{2019/11/12}{Use negative values for skip so that first paragraph after section headings is not indented.  This is something that I have been trying to fix for years!  It turns out that the anser was in the documentation for the standard command \cs{@startsection.}}
% \begin{macro}{\belowsectionskip}
% \begin{macro}{\aboveparaskip}
% \begin{macro}{\belowparaskip}
% Define some utility rubber lengths for use in
% section~\ref{sec:section-commands}. 
%    \begin{macrocode}
\newlength{\abovesectionskip}
\newlength{\belowsectionskip}
\newlength{\aboveparaskip}
\newlength{\belowparaskip}
%    \end{macrocode}
% The default values for these lengths are those traditionally used
% for the onecolumn style. 
%    \begin{macrocode}
\@set@scalable@length{abovesectionskip}{-3.5ex}{-4ex}{-2ex}
\@set@scalable@length{belowsectionskip}{2.3ex}{0ex}{2ex}
\@set@scalable@length{aboveparaskip}{3.25ex}{2ex}{0.5ex}
\@set@scalable@length{belowparaskip}{-1em}{\z@}{\z@}
%    \end{macrocode}
% \end{macro}\end{macro}\end{macro}\end{macro}
% We also make scalable some of the standard float parameters. These
% are the same defaults as in the standard article classes, but
% slightly more stretchy. For the twocolumn style, they are currently
% halved, see below. Note that I have also lumped |\abovecaptionskip|
% in with these, since it plays a similar role. There is also a
% |\belowcaptionskip| but it is set to zero in the standard classes. 
%    \begin{macrocode}
\@set@scalable@length{floatsep}{12\p@}{4\p@}{4\p@}
\@set@scalable@length{textfloatsep}{20\p@}{4\p@}{8\p@}
\@set@scalable@length{intextsep}{12\p@}{4\p@}{4\p@}
\@set@scalable@length{dblfloatsep}{12\p@}{4\p@}{4\p@}
\@set@scalable@length{dbltextfloatsep}{20\p@}{4\p@}{8\p@}
\@set@scalable@length{abovecaptionskip}{10\p@}{\z@}{\z@}
%    \end{macrocode}
% \begin{macro}{\@before@title@skip}         
% \begin{macro}{\@before@author@skip}        
% \begin{macro}{\@between@author@skip}       
% \begin{macro}{\@before@abstract@skip}      
% \begin{macro}{\@after@abstracthead@skip}   
% \begin{macro}{\@between@abstract@skip}     
% \begin{macro}{\@before@kwd@skip}           
% \begin{macro}{\@before@dates@skip}           
% Spacing parameters used in |\@maketitle|. These are not changed by
% default except by the user, so they don't need to go into
% |\@package@default@lengths|.  
%    \begin{macrocode}
\newlength{\@before@title@skip}         
\newlength{\@before@author@skip}        
\newlength{\@between@author@skip}       
\newlength{\@before@abstract@skip}      
\newlength{\@after@abstracthead@skip}   
\newlength{\@between@abstract@skip}     
\newlength{\@before@kwd@skip}           
\newlength{\@before@dates@skip}           
\@set@scalable@length{@before@title@skip}{2em}{\z@}{\z@}
\@set@scalable@length{@before@author@skip}{1.5em}{\z@}{\z@}
\@set@scalable@length{@between@author@skip}{0.5em}{\z@}{\z@}
\@set@scalable@length{@before@abstract@skip}{3ex}{1ex}{0.5ex}
\@set@scalable@length{@after@abstracthead@skip}{1ex}{\z@}{\z@}
\@set@scalable@length{@between@abstract@skip}{2ex}{1ex}{0.5ex}
\@set@scalable@length{@before@kwd@skip}{1ex}{0.5ex}{0.2ex}
\@set@scalable@length{@before@dates@skip}{2ex}{0.5ex}{0.2ex}
%    \end{macrocode}
% \end{macro}\end{macro}\end{macro}\end{macro}
% \end{macro}\end{macro}\end{macro}\end{macro}
% 
% Spacing above and below displayed equations also wants to be done
% like this, but it is more complicated since the standard article
% classes make these change with the fontsize. Something to do later,
% maybe. 
% 
%    \begin{macrocode}
%    \end{macrocode}
% \begin{macro}{\@package@default@lengths}
% The value actually used for these lengths is controlled by the
% \meta{SCALE} parameter, which depends on the style option that is in
% effect. This is set using the |\@package@default@lengths| command,
% which must be called every time the style changes (e.g., changing
% from onecolumn to twocolumn). 
%    \begin{macrocode}
\newcommand{\@package@default@lengths}{%
  \ifthenelse{\boolean{@rmtwocol}}{%
    %% twocolumn mode
    \abovesectionskip@set@SCALE{0.5}
    \belowsectionskip@set@SCALE{0.5}
    \aboveparaskip@set@SCALE{0.5}
    \belowparaskip@set@SCALE{1}
    \floatsep@set@SCALE{0.5}
    \textfloatsep@set@SCALE{0.5}
    \intextsep@set@SCALE{0.5}
    \dblfloatsep@set@SCALE{0.5}
    \dbltextfloatsep@set@SCALE{0.5}
    }{%
    %% onecolumn mode
    \abovesectionskip@set@SCALE{1}
    \belowsectionskip@set@SCALE{1}
    \aboveparaskip@set@SCALE{1}
    \belowparaskip@set@SCALE{1}
    \floatsep@set@SCALE{1}
    \textfloatsep@set@SCALE{1}
    \intextsep@set@SCALE{1}
    \dblfloatsep@set@SCALE{1}
    \dbltextfloatsep@set@SCALE{1}
    }%
  }
\@package@default@lengths
%    \end{macrocode}
% \end{macro}
% \begin{macro}{\RescaleLengths}
% \begin{macro}{\RescaleSecLengths}
% \begin{macro}{\RescaleFloatLengths}
% \begin{macro}{\RescaleTitleLengths}
% User commands that will fiddle with all the scalable lengths. 
%    \begin{macrocode}
\newcommand{\RescaleLengths}[1]{%
  \RescaleSecLengths{#1}\RescaleFloatLengths{#1}\RescaleTitleLengths{#1}}
\newcommand{\RescaleSecLengths}[1]{%
%% sections
  \abovesectionskipScale{#1}
  \belowsectionskipScale{#1}
  \aboveparaskipScale{#1}
  }
\newcommand{\RescaleFloatLengths}[1]{%
%% floats
  \floatsepScale{#1}
  \textfloatsepScale{#1}
  \intextsepScale{#1}
  \dblfloatsepScale{#1}
  \dbltextfloatsepScale{#1}
  \abovecaptionskipScale{#1}
  }
\newcommand{\RescaleTitleLengths}[1]{%
%% maketitle
  \@before@title@skipScale{#1}           
  \@before@author@skipScale{#1}        
  \@between@author@skipScale{#1}       
  \@before@abstract@skipScale{#1}      
  \@after@abstracthead@skipScale{#1}   
  \@between@abstract@skipScale{#1}     
  \@before@kwd@skipScale{#1}           
  \@before@dates@skipScale{#1}           
  }
%    \end{macrocode}
% \end{macro}\end{macro}\end{macro}\end{macro}
% \changes{2.2beta14}{2002/03/17}{Finished implementation of new
% infrastructure for scalable lengths. Converted the above/below
% section skips to use it. Also, all the floatsep parameters. Also,
% all the vertical lengths that are used by |\@maketitle|. }

% \begin{macro}{\MSextramargin}
% For the \texttt{manuscript} option we have an extra 0.3\,inch
% vertical margin and 1\,inch horizontal margin for marginal notes (by
% hand or with |\marginpar|. 
%
% \changes{3.23}{2007/09/09}{Reduce baselinestretch in
%   \texttt{manuscript} style to 1.2. No longer use raggedright for
%   manuscript.}
%
% \changes{3.24}{2007/09/14}{Further reduction of baselinestretch in
%   \texttt{manuscript} style to 1.1 since we now have wider margins.}
%
% \changes{3.24}{2007/09/14}{Split \texttt{manuscript}'s extra margin
%   into separate horizontal and vertical parts.}
% 
%    \begin{macrocode}
\newlength{\MSextraHmargin}\newlength{\MSextraVmargin}
\setlength{\MSextraHmargin}{1in} 
\setlength{\MSextraVmargin}{0.3in} 
\if@manuscript
   \renewcommand{\baselinestretch}{1.1}
   \addtolength{\textwidth}{-2\MSextraHmargin}
   \addtolength{\oddsidemargin}{\MSextraHmargin}
   \setlength{\evensidemargin}{\oddsidemargin}
   \addtolength{\textheight}{-2\MSextraVmargin}
   \addtolength{\topmargin}{\MSextraVmargin}
%    \raggedright
\fi
%    \end{macrocode}
% \end{macro}
% Float parameters: don't allow bottom floats by default. 
%    \begin{macrocode}
\renewcommand{\fps@figure}{tp}
\renewcommand{\fps@table}{tp}
%    \end{macrocode}
%
% \begin{macro}{\rmaa@AuthorWidthFrac}
% \begin{macro}{\rmaa@AbstractWidthFrac}
% These parameters control what fraction of the total text width can
% be used by the list of authors and the abstract/resumen,
% respectively.  
%    \begin{macrocode}
\newcommand{\rmaa@AuthorWidthFrac}{1.0}
\newcommand{\rmaa@AbstractWidthFrac}{0.75}
%    \end{macrocode}
% Use a wide abstract to save space in the new style.
%    \begin{macrocode}
\if@newstyle\renewcommand{\rmaa@AbstractWidthFrac}{1.0}\fi
%    \end{macrocode}
% \end{macro}
% \end{macro}
%
% Some things to consider in future versions. Trying to avoid too many
% widows, etc. Apparently, the following will work best with
% \cs{flushbottom}. 
%    \begin{macrocode}
%%\widowpenalty=10000
%%\clubpenalty=10000
%    \end{macrocode}
% Also, should perhaps try getting rid of underfull/overfull box warnings. 
% \changes{3.10}{2003/01/24}{Set large values for \cs{hbadness},
% \cs{vbadness}, \cs{hfuzz}, \cs{vfuzz}. To reduce amount of stuff
% that gets written to the terminal every time you compile something.}
%    \begin{macrocode}
\hbadness=10000
\vbadness=10000
\setlength\hfuzz{\maxdimen}
\setlength\vfuzz{\maxdimen}
%    \end{macrocode}
% \subsection{Various text styles for different contexts}
% \label{sec:text-styles} 
% \begin{macro}{\rmaa@AffilStyle}
% \begin{macro}{\rmaa@PagesStyle}
% \begin{macro}{\rmaa@YearStyle}
% \begin{macro}{\rmaa@VolumeStyle}
% \begin{macro}{\rmaa@VolumenStyle}
% \begin{macro}{\rmaa@JournalNameStyle}
% \begin{macro}{\rmaa@ReceivedStyle}
% \begin{macro}{\rmaa@AbstractHeaderStyle}
% \begin{macro}{\rmaa@JournalNameStyle}
% Style commands for the different fields used by |\maketitle| in
% constructing the title page of the article. These have the form
% |\rmaa@XYZStyle{TEXT}|. I'm not sure these were really a good idea,
% but never mind \dots 
%    \begin{macrocode}
\newcommand{\rmaa@AffilStyle}[1]{{\normalsize\rmfamily#1}}
\newcommand{\rmaa@PagesStyle}[1]{{\normalsize\rmfamily#1}}
\newcommand{\rmaa@YearStyle}[1]{{\normalsize\rmfamily#1}}
\newcommand{\rmaa@VolumeStyle}[1]{{\normalsize\rmfamily\bfseries#1}}
\newcommand{\rmaa@VolumenStyle}[1]{{\Large\sffamily\bfseries#1}}
\newcommand{\rmaa@JournalNameStyle}[1]{{\normalsize\itshape#1}}
\newcommand{\rmaa@ReceivedStyle}[1]{{\normalsize\rmfamily\itshape#1}}
\newcommand{\rmaa@AbstractHeaderStyle}[1]{%
  {\normalsize\rmfamily\MakeUpperCase{#1}}}
\newcommand{\rmaa@TitleStyle}[1]{%
  {\large\bfseries\rmfamily\boldmath\MakeUpperCase{#1}\par}}
%    \end{macrocode}
% \end{macro}
% \end{macro}
% \end{macro}
% \end{macro}
% \end{macro}
% \end{macro}
% \end{macro}
% \end{macro}
% \end{macro}
%
% \begin{macro}{\rmaa@FrontSecStyle}
% Headings of ``front sections'' (see section~\ref{sec:front-sections})
%    \begin{macrocode}
\newcommand{\rmaa@FrontSecStyle}[1]{%
  {\normalsize\bfseries\rmfamily\MakeUpperCase{#1}}}
%    \end{macrocode}
% \end{macro}
%
% \begin{macro}{\rmaa@TOCTitleStyle}
% \begin{macro}{\rmaa@TOCPartStyle}
% \begin{macro}{\rmaa@IndexTitleStyle}
% \begin{macro}{\rmaaTOCAuthorStyle}
% \begin{macro}{\rmaaIndexAuthorStyle}
% Commands to set style of TOC and index entries.
% Rewritten to avoid unnecessary font commands in the .toc and .ind files
%    \begin{macrocode}
\newcommand{\rmaa@TOCTitleStyle}[1]{\MakeUpperCase{#1}}
\newcommand{\rmaa@TOCPartStyle}[1]{\rmaa@TOCTitleStyle{#1}}
\newcommand{\rmaa@IndexTitleStyle}[1]{#1}
\newcommand{\rmaaTOCAuthorStyle}[1]{%
  {\itshape#1}}
\newcommand{\rmaaIndexAuthorStyle}[1]{%
  \protect\mbox{\bfseries#1}}
%    \end{macrocode}
% \end{macro}
% \end{macro}
% \end{macro}
% \end{macro}
% \end{macro}
%
% \begin{macro}{\rmaa@KeywordStyle}
% \begin{macro}{\rmaa@KeywordHeaderStyle}
% The style of the keywords depends on whether we are using the
% \texttt{newstyle} option or not. If we are, then they are small. 
%    \begin{macrocode}
\ifthenelse{\boolean{@newstyle}}{
  \newcommand{\rmaa@KeywordStyle}[1]{%
    {\small%
      \if@boldkeywords 
      \bfseries\MakeUpperCase{\raggedright #1}
      \else
      \raggedright #1
      \fi
    }}
  \newcommand{\rmaa@KeywordHeaderStyle}[1]{%
    {\small\itshape#1}}}
{
  \newcommand{\rmaa@KeywordStyle}[1]{%
    {\normalsize%
      \if@boldkeywords\bfseries\fi\raggedright #1}}
  \newcommand{\rmaa@KeywordHeaderStyle}[1]{%
    {\normalsize\itshape#1}}}
%    \end{macrocode}
% \end{macro}
% \end{macro}
%
% \subsection{Page styles}\label{sec:page-styles}
% \begin{macro}{\ps@rmaa}
% The running headers for normal pages:
%    \begin{macrocode}
\newcommand{\ps@rmaa}{%
%    \end{macrocode}
% nothing in either foot; 
%    \begin{macrocode}
  \renewcommand\@oddfoot{\relax}%
  \renewcommand\@evenfoot{\relax}%
%    \end{macrocode}
% page number and short author list in the even heads; 
%    \begin{macrocode}
  \renewcommand\@evenhead{\rlap{\thepage}\hfil\@shortauthor\hfil}%
%    \end{macrocode}
% short title and page number in the odd heads.
%    \begin{macrocode}
  \renewcommand\@oddhead{\hfil\@shorttitle\hfil\llap{\thepage}}
  }
%    \end{macrocode}
% \begin{macro}{\ps@rmaafirstpage}
% The header of the first page of each
% article:
%    \begin{macrocode}
\newcommand{\ps@rmaafirstpage}{%
  \renewcommand\@oddfoot{\hfil\thepage}%
%    \end{macrocode}
% footer corrected (06 Jun 2000) to have page on outside edge;
%    \begin{macrocode}
  \renewcommand\@evenfoot{\thepage\hfil}%
%    \end{macrocode}
% Header has journal name, volume, and page numbers for a regular
% article, or similar for special cases. 
% \changes{4.0}{2019/11/03}{Change to three-line header with copyright info and DOI.}
%    \begin{macrocode}
\if@preprint
\renewcommand\@oddhead{%
  \if@proceedings 
  \textsc{To appear in ``\@conftitle{} (\rmaa@YearStyle{\@year})''}\hfil
  \textit{RMxAC}
  \else
  \textsc{To appear in} \@journalname\hfil
  \fi}\fi%
\if@guide
\renewcommand\@oddhead{\textsc{Documentation} (\rmaa@YearStyle{\@year})\hfil}
\fi%
\if@manuscript
\renewcommand\@oddhead{Manuscript for \@journalname{} (\rmaa@YearStyle{\@year})\hfil}
\fi%
\if@journal
\renewcommand\@oddhead{%
  \parbox{\textwidth}{
    \@journalname, \rmaa@VolumeStyle{\@volume},
    \rmaa@PagesStyle{\@artpages}\
    (\rmaa@YearStyle{\@year})\hfil\\[1pt]
    {\small\textcopyright{} \@year:
      Instituto de Astronom\'\i{}a,
      Universidad Nacional Aut\'onoma de M\'exico}\hfil\\[1pt]
    {\small\@DOI@url}\hfil
  }%
}\fi%
  \renewcommand\@evenhead{\@oddhead}}
%    \end{macrocode}
% By default, use the |\ps@rmaa| page style.
%    \begin{macrocode}
\pagestyle{rmaa}
%    \end{macrocode}
% \end{macro}
% \end{macro}
%
% \subsection{Preamble commands}\label{sec:preamble} 
% Everything you need for setting the title page, TOC and index
% entries, etc.
% \subsubsection{Commands for the title page}
% \label{sec:title-commands}
% \begin{macro}{\resumen}
% \begin{macro}{\@resumen}
% \begin{macro}{\abstract}
% \begin{macro}{\@abstract}
% User commands for setting the resumen and abstract. 
% These should be used in the document header. The internal commands
% |\@resumen| and |\@abstract| are used to hold the actual text
% pending output via the |\maketitle| command. 
%    \begin{macrocode}
\newif\if@setresumen
\newcommand{\resumen}[1]{%
%%  \renewcommand{\@resumen}{{\selectlanguage{spanish}#1}}%
  \renewcommand{\@resumen}{{\ignorespaces#1}}%
  \@setresumentrue}
\renewcommand{\abstract}[1]{%
  \renewcommand{\@abstract}{{\ignorespaces#1}}
  \renewcommand\ADS@abstract{#1}
%% Not a good idea, since then author macros are not expanded and
%% cause errors later. However, by not doing it we are vulnerable to
%% fragile commands.....
  \if@noexpandADS\@onelevel@sanitize{\ADS@abstract}\fi
  \@abstractpresenttrue
  }
\newcommand{\ADSnoexpand}{\@noexpandADStrue}
%    \end{macrocode}
% \begin{macro}{\boldabstract}
% This is the version for doing 1-page conference proceedings. Still
% save the text for ADS but then write it out immediately as a bold
% paragraph. 
%    \begin{macrocode}
\newcommand{\boldabstract}[1]{%
  \gdef\ADS@abstract{#1}%
  \renewcommand{\@abstract}{#1}%
  \if@noexpandADS\@onelevel@sanitize{\ADS@abstract}\fi%
  \noindent\textbf{\boldmath #1}
  }
%    \end{macrocode}
% \end{macro}
% Default value in each case is to print out a warning if not set
%    \begin{macrocode}
  \newcommand{\@resumen}{%
    \@setresumenfalse \ClassWarningNoLine{rmaa}{No Spanish abstract
      (resumen) specified - if you do not speak Spanish, this will be
      translated by the editors}%
    \textbf{Favor de proporcionar un resumen en espa\~nol. If you are
      unable to translate your abstract into Spanish, the editors will
      do it for you.} }
  \newcommand{\@abstract}{\ClassWarningNoLine{rmaa}{No abstract
      specified}}
%    \end{macrocode}
% \end{macro}
% \end{macro}
% \end{macro}
% \end{macro}
%
% \begin{macro}{\title}
% \begin{macro}{\@rawtitle}
% \begin{macro}{\@TOCtitle}
% \begin{macro}{\rmaaIndexTitle}
% Similarly, |\title| is the preamble command to set the title of the
% article. This is used in various places, typeset in different
% styles, hence the use of the internal |\@rawtitle|, |\@title|,
% |\rmaaIndexTitle|, and |\@TOCtitle| commands, each of which uses the
% relevant style from section~\ref{sec:text-styles}
%    \begin{macrocode}
\newcommand{\@rawtitle}{Raw Title}
\newcommand{\@TOCtitle}{\rmaa@TOCTitleStyle{Title for TOC}}
\newcommand{\rmaaIndexTitle}{\rmaa@IndexTitleStyle{Title for index}.}
\newcommand{\@TOCaftertitle}{}
\renewcommand{\title}[1]{%
%%  {\ADS@sanitize\xdef\@rawtitle{#1}}
  \renewcommand\@rawtitle{\noexpand#1}
  \let\ADS@title=\@rawtitle
  \renewcommand{\@title}{\rmaa@TitleStyle{#1}}%
  \renewcommand{\@TOCtitle}{\rmaa@TOCTitleStyle{#1}\@TOCaftertitle}%
  \renewcommand{\rmaaIndexTitle}{\rmaa@IndexTitleStyle{#1\@TOCaftertitle}.}%
  }
%    \end{macrocode}
% \end{macro}
% \end{macro}
% \end{macro}
% \end{macro}
% 
% All the weird and whacky ways of indicating the author
% affiliations. 
% \begin{macro}{\affil}
% First, the way that sets it directly underneath the authors name.
% Note that |\affil| trumps all other affiliations for the purposes of
% getting into the ADS stuff (we use |\gdef|).   
% \dots 
%    \begin{macrocode}
\newcommand{\affil}[1]{%
  \\[1.5ex]\rmaa@AffilStyle{#1}\\[1.5ex]%
  \if@firstaffil\gdef\ADS@affil{#1}\@firstaffilfalse\fi%
  }
\renewcommand{\and}{and\\[1.5ex]}
%    \end{macrocode}
% \begin{macro}{\altaffil}
% Alternative affiliations are footnoted. 
%    \begin{macrocode}
\newcommand{\altaffil}[1]{%
  \if@firstaltaffil\gdef\ADS@affil{#1}\@firstaltaffilfalse\fi%
  \thanks{#1}}
%    \end{macrocode}
% \begin{macro}{\altaffilmark}
% \begin{macro}{\altaffiltext}
% Alternative-alternative mechanism for affiliations---should not be
% mixed with |\altaffil| (unless you know what you are doing) or there
% will be problems with numbering.
%    \begin{macrocode}
\newcommand{\altaffilmark}[1]{\textsuperscript{#1}}
\newcommand{\altaffiltext}[2]{%
  \if@firstaltaffil\gdef\ADS@affil{#2}\@firstaltaffilfalse\fi
  \stepcounter{footnote}%
  \protected@xdef\@thanks{\@thanks
    \protect\footnotetext[#1]{#2}}} 
%    \end{macrocode}
% \end{macro}
% \end{macro}
% \end{macro}
% \end{macro}
%
% \begin{macro}{\keywords}
% \begin{macro}{\@keywords}
% \begin{macro}{\@keywordheader}
% \label{keywords}Sundry minor preamble commands. All
% self-explanatory, I hope. |\keywords| is the old user interface
% (deprecated). New articles should use |\addkeyword| once for each
% keyword. That way the ADS record comes out right too (comma rather
% than em-dash separator).
% \changes{2.2beta12}{2001/12/07}{\cs{addkeyword} New user interface
%   for keywords.} 
% \changes{2.2beta13}{2002/03/11}{\cs{addkeyword} Improved behaviour
%   with .ads file. Use
%   \cs{ADS@sanitize} to change definitions of `~' and \cs{ion}.}
% \changes{3.23}{2007/09/09}{\cs{addkeyword} now signals a compilation
% error if the author tries to add more than 6 keywords.}
% \changes{3.24}{2007/09/14}{kwdcounter is now reset in
% \cs{@begofarticlehook} to prevent error in book style} 
%    \begin{macrocode}
\newcommand{\@kwdsep}{ --- }
\newcommand{\ADS@kwdsep}{, }
\newcounter{kwdcounter}\setcounter{kwdcounter}{0}
\newcommand{\addkeyword}[1]{%
  \addtocounter{kwdcounter}{1}
  \ifthenelse{\thekwdcounter > 6}{
    \ClassError{rmaa}{Too many keywords. A maximum of 6 keywords is allowed.}{}
  }{}
  \if@firstkwd%
  \renewcommand\@keywords{\if@boldkeywords\uppercase{#1}\else #1\fi}
  {\ADS@sanitize\def\ion{\plainion}\gdef\ADS@keywords{#1}}
  \@firstkwdfalse
  \else%
  \edef\@keywords{\@keywords\@kwdsep \if@boldkeywords\uppercase{#1}\else #1\fi}%
  {\ADS@sanitize\def\ion{\plainion}\xdef\ADS@keywords{\ADS@keywords\ADS@kwdsep #1}}%
  \fi}
\newcommand{\ADS@sanitize}{%
  \def\mathrm{}
  \def\@{\ }
  \def~{\ }
  \def\langle{<}
  \def\rangle{>}
  \def\log{log }
  \def\sin{sin }
  \def\cos{cos }
  \def\tan{tan }
  \def\fd{^d.}
  \def\fh{^h.}
  \def\fm{^m.}
  \def\fs{^s.}
  \def\fdg{�.}
  \def\farcm{'.}
  \def\farcs{''.}
  \def\arcdeg{�}
  \def\,{ }
  }
\newcommand{\keywords}[1]{%
  \renewcommand{\ADS@keywords}{#1}
  \renewcommand{\@keywords}{\uppercase{#1}}
  }
\newcommand{\@keywordheader}{\rmaa@KeywordHeaderStyle{Key Words:}}
\newcommand{\@keywords}{%
  \ClassWarningNoLine{rmaa}{No keywords specified}%
  Keywords go here
  }
%    \end{macrocode}
% \begin{macro}{\shorttitle}
% \begin{macro}{\shortauthor}
% \begin{macro}{\@shorttitle}
% \begin{macro}{\@shortauthor}
% Commands for setting the running headers.
%    \begin{macrocode}
\newcommand{\shorttitle}[1]{%
  \renewcommand\@shorttitle{\MakeUpperCase{#1}}}
\newcommand{\shortauthor}[1]{%
  \renewcommand\@shortauthor{\MakeUpperCase{#1}}}
\newcommand{\@shorttitle}{SHORT TITLE GOES HERE}
\newcommand{\@shortauthor}{SHORT AUTHOR GOES HERE}
%    \end{macrocode}
% \end{macro}
% \end{macro}
% \end{macro}
% \end{macro}
% \end{macro}
% \end{macro}
% \end{macro}


% \subsubsection{The full addresses}
% \label{sec:full-addresses}
% \begin{macro}{\fulladdresses}
% \begin{macro}{\@fulladdresses}
% \begin{macro}{\@fulladdheight}
% \begin{macro}{\@fulladddepth}
% \begin{macro}{\@fulladdbox}
% Aaghh, the full addresses! A perennial source of problems. 
%    \begin{macrocode}
\newlength{\@fulladdheight}
\newlength{\@fulladddepth}
\newsavebox{\@fulladdbox}
%% \fulladdresses command modified (11 May 2000) so that we can put
%% \fulladdresses{} in the article to suppress the output of the
%% addresses
\newcommand{\@FinalFigure}{\relax}
\newcommand{\FinalFigure}[1]{\renewcommand{\@FinalFigure}{#1}}
\newcommand{\fulladdresses}[1]{%
  \renewcommand{\@rmaaArg}{#1}
  \ifx\@rmaaArg\@rmaaEmpty
     \renewcommand{\@fulladdresses}{}
  \else
     \renewcommand{\@fulladdresses}{
       \normalsize%
       \ifx\@FinalFigure\relax\else\vfill\@FinalFigure\vfill\fi
%%       \vspace*{1ex}
       %% do addresses as a list (copied from thebibliography environment) 
       \list{\null}{\leftmargin 1.5em\topsep\z@\labelwidth\z@\labelsep\z@\itemindent -1.5em
         \itemsep\z@\parsep\z@\usecounter{enumi}}
       \sloppy\clubpenalty4000\widowpenalty4000
       \sfcode`\.=1000\relax
       #1
       \endlist
       }
  \fi
  }
\newcommand{\@fulladdresses}{Full addresses go here}
%    \end{macrocode}
% \begin{macro}{\suppressfulladdresses}
% New user-level command for suppressing full addresses should you so
% desire. Currently used in the single-page-poster style. 
%    \begin{macrocode}
\newcommand{\suppressfulladdresses}{\@suppressfulladdressestrue\@needaddressesfalse}
%    \end{macrocode}
% \end{macro}
% \end{macro}
% \end{macro}
% \end{macro}
% \end{macro}
% \end{macro}

% \begin{macro}{\printdates}
% \begin{macro}{\noprintdates}
% Command to activate the printing out of the received/accepted
% dates. This can also be done using the \texttt{printdates} option,
% which is the default for \texttt{journal}, but not for \texttt{proceedings}.
%    \begin{macrocode}
\newcommand{\printdates}{\@printdatestrue}
\newcommand{\noprintdates}{\@printdatesfalse}
%    \end{macrocode}
% \end{macro}
% \end{macro}
% \begin{macro}{\@outputaddressesifnecessary}
% Command to output the addresses if we haven't already done so
% This is called by: 
% \begin{enumerate}
% \item The end of |thebibliography| environment
% \item The |\outputfulladdresses| command
% \end{enumerate}
%    \begin{macrocode}
\newcommand{\@outputaddressesifnecessary}{%
  \ClassInfo{rmaa}{Calling @outputaddressesifnecessary}
  \ifthenelse{\boolean{@needaddresses}}{%
    %% If in one column style, or in the first column of a two column
    %% style then we can try to put it on this page 
    \ifthenelse{\boolean{@firstcolumn} \or \not\boolean{@rmtwocol}}{%
      \@outputaddresses}{%
      %% otherwise give warning it will be on a separate page
      \ClassWarning{rmaa}{Full addresses will go on a separate
        page (please see User Guide for how to remedy this)}}%
    %% If we haven't managed yet, put it on its own page
    \ifthenelse{\boolean{@needaddresses}}{%
      \if@debug{\ClassWarning{rmaa}{In this silly bit!}}\fi
      \clearpage\vfill\@outputaddresses}{}%
    }{
    %% If we don't need them, we still might need to do the eoa hook 
    \ifthenelse{\boolean{@suppressfulladdresses}}{%
      \if@debug{\ClassWarning{rmaa}{XXX}}\fi
%%      \if@book\else\@endofarticlehook\null\fi
      }{}
    }%
  }  
%    \end{macrocode}
% \end{macro}

% \begin{macro}{\outputfulladdresses}
% A user interface to the above command, in case there is no biblio
% (11 May 2000) 
%    \begin{macrocode}
\newcommand{\outputfulladdresses}{\@outputaddressesifnecessary}
%    \end{macrocode}
% \end{macro}

% \begin{macro}{\adjustfinalcols}
% This command starts a new column and should be 
% used to ensure that the columns balance on the last page. It also
% outputs the full 
% author addresses to ensure that this occurs in the left hand
% column. 
% The |\par| command is so that the current line will be taken account
% of when working out |\pagetotal|
%    \begin{macrocode}
\newcommand{\adjustfinalcols}{\if@twocolumn\goodbreak\@outputaddresses\newpage\noindent\fi}
%    \end{macrocode}
% \end{macro}

% \begin{macro}{\@outputaddresses}
% Command to do the work of outputting the addresses. 
% This can be called by 
% \begin{enumerate}
% \item an explicit |\adjustfinalcols|
% \item via |\@outputaddressesifnecessary|, q.v.
% \end{enumerate}
%
%    \begin{macrocode}
\newcommand{\@outputaddresses}{%
  % Allow option of suppressing the full addresses, in which case
  % we just run the hook if we are not in a book style
  \if@suppressfulladdresses
%%     \if@book\else\@endofarticlehook\null\fi
  \else
  %% save the full addresses in a box for output later
  \sbox{\@fulladdbox}{\parbox{\textwidth}{\@fulladdresses}}
  % calculate height of address box  
  \settoheight{\@fulladdheight}{\usebox{\@fulladdbox}}
  \settodepth{\@fulladddepth}{\usebox{\@fulladdbox}}
  \addtolength{\@fulladdheight}{\@fulladddepth}
  \if@debug
  { \typeout{Address total height: \the\@fulladdheight}
    \typeout{Text height: \the\pagetotal}
    \typeout{Goal height: \the\pagegoal}
    }
  \fi
  % reuse \@fulladdheight to mean the total height of text+addresses
  \addtolength{\@fulladdheight}{\pagetotal}
%%  \addtolength{\@fulladdheight}{-\baselineskip}
  \if@debug\typeout{Adjusted text + addresses height: \the\@fulladdheight}\fi
  % check it will fit on current page
  \ifdim\@fulladdheight<\pagegoal
     %% must shift it to the left
     %% to compensate for the hanging indentation in the reference list. 
     %% NOTE: this is not foolproof!
     %% NB. This will give overfull hbox warnings, which is a bit ugly.
     \vfill\noindent\if@biblist\hspace*{-\leftmargin}\fi
     \usebox{\@fulladdbox}
     % Set this to show we've done it
     \@needaddressesfalse
     % and assume we are at the end of the article
     % so that we do the EOA hook if not in the book style
%%     \if@book\else\@endofarticlehook\fi
  \else
     % if not, just do nothing. With luck, this function will get 
     % called again on the next page
     \if@debug{\typeout{Holding over full addresses to next page}}\fi
  \fi\fi
  }
%    \end{macrocode}
% \end{macro}

% \subsubsection{Commands for editors}
% \label{sec:editor-commands}
% Some user/editor commands for doing the TOC and index in the book
% option. 
% They are defined here rather in the book-specific part since they
% need to be always available in order to avoid breaking documents. 
% \begin{macro}{\listofauthors}
% \begin{macro}{\@rawauthors}
% \begin{macro}{\rmaaTOCauthors}
% The command |\listofauthors| formats an author list for use in the
% TOC. It should contain the same names as in the |\author| command,
% but without the affiliation garbage and using |\&| instead of
% |and|. 
%    \begin{macrocode}
\newcommand{\@rawauthors}{Raw Authors}
\newcommand{\rmaaTOCauthors}{\noexpand\rmaaTOCAuthorStyle{Author list for TOC}}
\newcommand{\listofauthors}[1]{%
  \renewcommand{\@rawauthors}{#1}
  \renewcommand\rmaaTOCauthors{\noexpand\rmaaTOCAuthorStyle{#1}}}
%    \end{macrocode}
% \begin{macro}{\indexauthor}\label{indexauthor}
% The command |\indexauthor| produces an index entry for each
% author. 
% This command must come \emph{after} the
% |\title| and the |\listofauthors| commands. Also, there should be one 
% |\indexauthor| command for \emph{each} author. We also add the
% author onto the end of the |\ADS@authors| macro and construct the
% |\ADS@authorinitial| by a most cunning method (see below), but only
% for the first author. 
% 
% No, No, NO! This was all wrong. Backed out completely: 
% \begin{quote}
% Completely rewritten (19 Dec 2002) since we now use the index class, so 
% as to be able to alphabetize accented characters. A side effect of this is
% that |\index| no longer expands out the argument after the |@|, so we 
% instead add all the |\index| commands to a list (|\indexauthor@deferlist|), 
% which is executed in the |\@begofarticlehook|. 
% 
% Note the stuff with |\toks0| is copied from Exercise 20.15 of the \TeX{}book---it is
% to make sure the latin1-accented name gets passed into the idx
% file. 
% \end{quote}
% We now go back to the old way of doing things, using an optional
% argument to the \cs{indexauthor} command that has the name without
% accents. 
% \changes{3.08}{2003/01/16}{Remerged changes from 3.06 that got
% dropped from 3.07 after a synchronization disaster. Basically, the
% accentless optional argument to \cs{indexauthor}. Must add the macro
% directories to my unison set-up to stop this sort of thing from happening.}
%    \begin{macrocode}
\newcommand\@isortkey{}
\newcommand\@emptykey{empty}
\newcommand\indexauthor[2][empty]{
  \renewcommand\@isortkey{#1}
  \ifx\@isortkey\@emptykey \renewcommand\@isortkey{#2}\fi
  \index{\@isortkey @{{\rmaaIndexAuthorStyle{#2}}\hskip1em plus1em minus0.4em
      \rmaaIndexTitle\hskip1em plus2em minus0.4em\penalty-100
      \rmaaTOCauthors}}
  \indexauthorADS{#2}
}
%%\def\indexauthor@deferlist{\relax}
%%\long\def\indexauthor#1{
%%  \def\@@name{#1} \toks0=\expandafter{\@@name}
%%  \toks1=\expandafter{\indexauthor@deferlist}
%%  \protected@xdef\indexauthor@deferlist{%
%%    \the\toks1
%%    \noexpand\index{\the\toks0@{\noexpand\rmaaIndexAuthorStyle{\the\toks0}\noexpand\hskip1em plus1em minus0.4em 
%%        \rmaaIndexTitle\noexpand\hskip1em plus2em minus0.4em\noexpand\penalty-100
%%        \rmaaTOCauthors}}
%%  }
%%  \indexauthorADS{#1}
%%  }
%%\xdef\indexauthor@deferlist{}
\newcommand\indexauthorADS[1]{%
  \if@firstauthor\protected@xdef\ADS@authorinitial{\firstletter{#1}}\fi
  \protected@xdef\ADS@authors{\ADS@authors\if@firstauthor\else ; \fi
    #1}
  \@firstauthorfalse
}
%    \end{macrocode}
% Note that the first |#1| sets the key for the index, making sure it
% actually starts with  
% first letter of the surname. Meanwhile, the |@{...}| expression,
% which starts with some  
% control sequences is what is written to the |.ind| file. 
% See the makeindex man page for details. 
% \end{macro}
% \end{macro}
% \end{macro}
% \end{macro}

% \begin{macro}{\KillAllAccents}
% \begin{macro}{\firstletter}
% These are utility functions used in the above. |lett| is a counter
% indicating the current letter in the string. |\firstletter| just
% prints the first letter of its argument.
%    \begin{macrocode}
\newcommand{\KillAllAccents}{\def\'{}\def\`{}\def\~{}\def\.{}\def\i{i}}
\newcommand\firstletter[1]{\@car #1\@nil}
%    \end{macrocode}
% \end{macro}
% \end{macro}
% \begin{macro}{\SetVolume}
% \begin{macro}{\SetMSnumber}
% \begin{macro}{\SetFirstPage}
% \begin{macro}{\SetYear}
% \begin{macro}{\ReceivedDate}
% \begin{macro}{\AcceptedDate}
% These are for the editors to set the volume number, etc. 
%    \begin{macrocode}
\newcommand{\SetVolume}[1]{%
  \renewcommand{\@volume}{#1}
  \edef\ADS@volume{\PadToFourWithDots{#1}}
  \renewcommand{\@volumen}{\rmaa@VolumenStyle{volumen #1}}%
  }
\newcommand{\SetMSnumber}[1]{\renewcommand{\@MSnum}{#1}}
%    \end{macrocode}
% Note that there is no real reason to use |\SetFirstPage| since all
% the page numbering is done automatically in the book style. 
%    \begin{macrocode}
\newcommand{\SetFirstPage}[1]{\setcounter{page}{#1}}
\newcommand{\SetYear}[1]{%
  \renewcommand{\@year}{#1}
  \renewcommand{\ADS@year}{#1}
  \renewcommand{\@anyo}{\rmaa@VolumenStyle{#1}}
  }
\newcommand{\ReceivedDate}[1]{\renewcommand{\@received}{#1}}
\newcommand{\AcceptedDate}[1]{\renewcommand{\@accepted}{#1}}
%    \end{macrocode}
%
% \begin{macro}{\DOI}
% User/editor command to enter the DOI, which will be typeset as a link.
% \changes{4.0}{2019/11/02}{New user command \cs{DOI} to typeset
% DOI in the header of first page and add it to ADS info.}
%    \begin{macrocode}
\newcommand\DOI[1]{%
  \def\@DOI@url{\url{https://doi.org/10.22201/ia.01851101p.#1}}%
  \def\ADS@doi{https://doi.org/10.22201/ia.01851101p.#1}%
}%
%    \end{macrocode}
% \end{macro}
%
% \begin{macro}{\ORCIDiD}
% \changes{4.0}{2019/11/12}{Add mechanism for giving ORCID iD links}
% \begin{macro}{\ORCIDiDfull}
% \begin{macro}{\@orcid@icon}
% \begin{macro}{\@orcid@iconbig}
% Commands to add icon links to ORCID iDs of authors. For example
% \begin{quote}
% |\ORCIDiD|\marg{0000-0000-0000-0000}
% \end{quote}
% to be used after each author name in the |\author|
% command and in the |\fulladdresses|.  Alternatively,
% |\OORCIDiDfull| will typeset the icon plus the full text of the
% link, but we don't use that at present.  The icon image is stored in |ORCIDiD_iconvector.pdf|, which needs to be on the input path.  In principle, we could make EPS versions if necessary.
%    \begin{macrocode}
\newcommand\@orcid@icon{\includegraphics[height=1.8ex]{ORCIDiD_iconvector}}
\newcommand\@orcid@iconbig{\includegraphics[height=2ex]{ORCIDiD_iconvector}}
\newcommand\ORCIDiD[1]{\href{https://orcid.org/#1}{\,\@orcid@icon}}
\newcommand\ORCIDiDfull[1]{%
  \href{https://orcid.org/#1}{\@orcid@iconbig}\,\url{https://orcid.org/#1}}
%    \end{macrocode}
% \end{macro}
% \end{macro}
% \end{macro}
% \end{macro}
%
% \begin{macro}{\@volume}
% \begin{macro}{\@volumen}
% \begin{macro}{\@artpages}
% \begin{macro}{\@year}
% \begin{macro}{\@anyo}
% \begin{macro}{\@received}
% \begin{macro}{\@accepted}
% \begin{macro}{\@MSnum}
% Default values for the containers used by above commands, plus a
% few more useful ones. 
%    \begin{macrocode}
\newcommand{\@volume}{00}
\newcommand{\@volumen}{\rmaa@VolumenStyle{00}}
\newcommand{\@artpages}{%
  \oldpageref{\articlename:RMAAFirstPage}--\oldpageref{\articlename:RMAALastPage}}
\newcommand{\@year}{\the\year}
\newcommand{\@anyo}{\rmaa@VolumenStyle{2007}}
\newcommand{\@journalname}{\rmaa@JournalNameStyle{%
    \if@proceedings\RMAASClong\else\RMAAlong\fi}}
\newcommand{\@received}{}
\newcommand{\@accepted}{}
\newcommand{\@MSnum}{0000}
\newcommand{\@DOI@prefix}{https://doi.org/10.22201/ia.01851101p}
\def\@DOI@url{\url{https://doi.org/10.22201/ia.01851101p.2099.00.00.00}}
%    \end{macrocode}
% \end{macro}
% \end{macro}
% \end{macro}
% \end{macro}
% \end{macro}
% \end{macro}
% \end{macro}
% \end{macro}
% \end{macro}
% \end{macro}
% \end{macro}
% \end{macro}
% \end{macro}
% \end{macro}
% 
% The following commands are most useful in the \texttt{book}
% style. However, the |\SetConfTitle| command can now be used in
% individual articles to set the header for the \texttt{preprint}
% style. 
%    \begin{macrocode}
\newcommand{\@reunionplace}{Lugar de la reuni\'on}
\newcommand{\@reuniondate}{Fecha de la reuni\'on}
\newcommand{\@conftitle}{Conference Title}
\newcommand{\SetReunion}[2]{\renewcommand{\@reunionplace}{#1}\renewcommand{\@reuniondate}{#2}}
\newcommand{\SetConfTitle}[1]{\renewcommand{\@conftitle}{#1}}
\newcommand{\@editors}{Juan Ram\'on de la Fuente and ``El Mosh''}
\newcommand{\SetEditors}[1]{\renewcommand{\@editors}{#1}}
\newcommand{\@editores}{Juan Ram\'on de la Fuente y ``El Mosh''}
\newcommand{\SetEditores}[1]{\renewcommand{\@editores}{#1}}
\newcommand{\@mes}{mes de la publicaci\'on}
\newcommand{\SetMes}[1]{\renewcommand{\@mes}{\rmaa@VolumenStyle{#1}}}
%    \end{macrocode}

% \subsection{Notes for Author Proofs and Communication between Editors}
% \label{sec:notes}
% We introduce a command \cs{Note} that should be used to describe
% any change made to the manuscript. These are a bit tricky to get
% right. In general we use \cs{marginpar} but this only works in outer
% paragraph mode, so we just bung the box where we are if we are in a
% figure caption or something. If we are in the preamble, then we just
% use a \marginpar as part of the \maketitle command.
%    \begin{macrocode}
\newif\if@arenotes
\@arenotesfalse
\newcounter{notecounter}
\setcounter{notecounter}{0}
\newcommand\NotesList{}
\newcommand\Note[1]{%
  \if@proof{%
    \@arenotestrue%
    \stepcounter{notecounter}%
    \if@inpreamble\else{%
      \NoteMark{\arabic{notecounter}}%
    }\fi
    \protected@xdef\NotesList{\NotesList 
      \protect\item #1}%
  }\fi%
}
\newcommand\Changed[2]{Changed ``#1'' to ``#2''.}
\newcommand\NoteMark[1]{%
  \ifinner
  \raisebox{1ex}[0pt][0pt]{\makebox[0pt][l]{\framebox{\footnotesize\bfseries
        N#1}}}%
  \else
  \marginpar{\framebox{\footnotesize\bfseries N#1}}%
  \fi
}
\newcommand\RefNote[1]{\if@proof\NoteMark{\ref{#1}}\relax\fi}
\newcommand\NoteLabel{\protect\label}
\newcommand\OutputNotes{%
  \bgroup
  \onecolumn
  \section*{EDITORIAL NOTES}
  \begin{enumerate}
    \NotesList
  \end{enumerate}
  \egroup
}
\DeclareRobustCommand\cs[1]{\texttt{\char`\\#1}}

%    \end{macrocode}
% \subsection{Bibliographic Records for ADS}
% \label{sec:ADS}
% \changes{2.2beta11}{2001/11/27}{Just about finished all the ADS
% stuff now. Except for the references, which I have given up on for
% the time being.}
% \begin{macro}{\ADS@BibCode}
% \begin{macro}{\ADS@year}
% \begin{macro}{\ADS@journal}
% \begin{macro}{\ADS@volume}
% \begin{macro}{\ADS@firstpage}
% \begin{macro}{\ADS@authorinitial}
% We use the |\ADS@BibCode| command to hold the bibliographic code for
% ADS. It is 
% used in the |%R| 
% field of the bibliographic record for each article. 
% The format is 
% \begin{quote}
% \texttt{YYYYJJJJJVVVVMPPPPA}  
% \end{quote}
% where \texttt{YYYY} is the year, \texttt{JJJJJ} is the journal code,
% which is \texttt{RMxAA} or \texttt{RMxAC}, \texttt{VVVV} is the
% volume number (right-justified), \texttt{M} is not used by us,
% \texttt{PPPP} is the first page number (right-justified), and
% \texttt{A} is the initial of the first author's surname. The volume
% and page fields are padded with periods (`\texttt{.}'), which causes
% us some fun (this is done with the macro |\PadToFourWithDots| described
% below) . 
% 
% All these individual bits are held in the commands |\ADS@year|,
% |\ADS@journal|, |\ADS@volume|, |\ADS@firstpage|,
% |\ADS@authorinitial|. 
%    \begin{macrocode}
\newcommand{\SetADSBibCode}[1]{\renewcommand{\ADS@BibCode}{#1}}
\newcommand{\ADS@BibCode}
{\ADS@year\ADS@journal\ADS@volume.\PadToFourWithDots{\ADS@firstpage}\ADS@authorinitial}
\newcommand{\ADS@year}{2001} 
\newcommand{\ADS@journal}{\if@proceedings RMxAC\else RMxAA\fi}
\newcommand{\ADS@volume}{....} 
%    \end{macrocode}
% These two are just set to default values here, |\ADS@firstpage| is set
% in the |document| environment for included articles, in the
% |frontsection| environment, |\tableofcontents| command,  |index| and
% |abstracts| environments. Also by |\AtBeginDocument| for stand-alone
% articles. 
% 
% God knows how we are going to do |\ADS@authorinitial|!
%    \begin{macrocode}
\newcommand{\ADS@firstpage}{....}
\newcommand{\ADS@authorinitial}{A}
%    \end{macrocode}
% \end{macro}
% \end{macro}
% \end{macro}
% \end{macro}
% \end{macro}
% \end{macro}
% We also have a whole slew of other fields that can be set. These are
% all listed in Table~\ref{tab:ADS}. Most of them still have to be
% defined. 
% \begin{table}[!ht]\raggedleft
% \caption{ADS Bibliographic Records}\label{tab:ADS}
% \begin{tabular}{lp{8em}lp{8em}}\hline
% Code & Description & Macro Name & Comment \\ \hline
% |%R| & Bibliographic Code & |\ADS@BibCode| & required \\
% |%A| & Author list & |\ADS@authors| & required, generate from |\indexauthor| \\
% |%F| & Author affiliation & |\ADS@affil| & format: AA(\dots)
% AB(\dots) \dots \\
% |%J| & Journal Name & |\ADS@journalname| &  \\
% |%D| & Publication Date & |\ADS@fulldate| &  \\
% |%L| & Last Page & |\ADS@lastpage| & generate from label \\
% |%T| & Title & |\ADS@title| & required, generate from |\title| \\
% |%C| & Abstract Copyright & |\ADS@copyright| & set to IA-UNAM \\
% |%U| & URL for electronic document & |\ADS@url| & set to PDF file \\
% |%d| & DOI for article & |\ADS@doi| & set by |\DOI| \\
% |%K| & Keywords & |\ADS@keywords| & set equal to |\@keywords| \\
% |%A| & Abstract Text & |\ADS@abstract| & set equal to |\@abstract| \\
% |%R| & References & |\ADS@references| & generate from the
% bibliography  \\ \hline
% \end{tabular}
% \end{table}
% A lot of these macros are set by other commands, as mentioned in the
% Table, but some need to be set  explicitly here. 
% 
% \begin{macro}{\ADS@doi}
% \changes{4.0}{2019/11/08}{Save DOI for ADS}
% The user/editor command \cs{DOI} is used to set this. 
%    \begin{macrocode}
\newcommand{\ADS@doi}{https://doi.org/10.22201/ia.01851101p.0000.00.00.00}
%    \end{macrocode}
% \end{macro}
% 
% \begin{macro}{\ADS@journalname}
% This
% does not contain any accents at the moment so as not to cause
% problems---this maybe should be changed though. 
%    \begin{macrocode}
\newcommand{\ADS@journalname}{%
  \if@proceedings\@conftitle\space (Eds. \@editors) \fi%
  Revista Mexicana de Astronom\'\i{}a y Astrof\'\i{}sica%
  \if@proceedings\space (Serie de Conferencias)\fi%
  \space Vol. \@volume, pp. \ADS@firstpage--\ADS@lastpage\space (\@year)%
  \space (http://www.astroscu.unam.mx/\~{}rmaa/)
  }
%    \end{macrocode}
% \end{macro}
% \begin{macro}{\ADS@keywords}
% \begin{macro}{\ADS@abstract}
% \begin{macro}{\ADS@title}
% |\ADS@abstract| is straightforward, since we already have the text
% saved as |\@abstract|. We could do the same with |\ADS@keywords| but
% to avoid unwanted formatting we explicitly set it in |\keywords|
% (p.~\pageref{keywords}). |\ADS@title| is just set to |\@rawtitle|,
% which we were already using for another purpose. 
%    \begin{macrocode}
\let\ADS@title=\@rawtitle
\newcommand{\ADS@keywords}{No keywords specified}
\newcommand{\ADS@abstract}{No abstract specified}
%    \end{macrocode}
% \end{macro}\end{macro}\end{macro}
% \begin{macro}{\ADS@authors}
% |\ADS@authors| wants to be in ``Other, A. N.'' format, but with
% semi-colons between authors. The most natural way to do this seems to
% be with |\indexauthor| (p.~\pageref{indexauthor}). We append each
% author to this macro, which means we must explicitly clear it for
% each new article (done in |\article|). .   
%    \begin{macrocode}
\newcommand{\ADS@authors}{}
%    \end{macrocode}
% \end{macro}
% \begin{macro}{\ADS@fulldate}
% \begin{macro}{\ADS@lastpage}
% \begin{macro}{\ADS@copyright}
% \begin{macro}{\ADS@website}
% \begin{macro}{\ADS@url}
% \begin{macro}{\SetWebsite}
% \begin{macro}{\SetMonthNumeric}
% Other sundry data, together with user commands for bits that can be
% set by the Editors. The |\ADS@url| is constructed from the website
% (currently |http://www.astrosmo.unam.mx/rmaa/|,
% the journal descriptor (e.g., |RMxAC..10/|), followed by the
% path to the PDF file (e.g., |PDF/RMxAC..10_HEW1.pdf|). 
% 
%    \begin{macrocode}
\newcommand{\ADS@website}{http://www.astrosmo.unam.mx/rmaa/}
\newcommand{\ADS@url}
{\ADS@website\ADS@journal\ADS@volume\ADS@number/PDF/\ADS@journal\ADS@volume\ADS@number\_\articlename.pdf}
\newcommand{\ADS@number}{}
\newcommand{\ADS@fulldate}{00/\ADS@year} 
\newcommand{\ADS@lastpage}{00} 
\newcommand{\ADS@copyright}{(C) \ADS@year; Instituto de Astronom\'\i{}a,
  Universidad Nacional Aut\'onoma de M\'exico} 
\newcommand{\SetNumber}[1]{\renewcommand{\ADS@number}{-#1}}
\newcommand{\SetWebsite}[1]{\renewcommand{\ADS@website}{#1}}
\newcommand{\SetMonthNumeric}[1]{\renewcommand{\ADS@fulldate}{#1/\ADS@year}}
%    \end{macrocode}
% \end{macro}\end{macro}\end{macro}\end{macro}\end{macro}\end{macro}\end{macro}
% \begin{macro}{\ADS@affil}
% \begin{macro}{\ADS@references}
% These two are somewhat more problematical. Affiliation currently
% only works for the first author (see \dots) pending redesign of the
% user interface. References are now done by wrapping the
% |thebibliography| environment to open another input stream on the
% same file and search through for the |\bibitems|. 
%    \begin{macrocode}
\newcommand{\ADS@references}{}
\newcommand{\ADS@affil}{}
%    \end{macrocode}
% \end{macro}\end{macro}
% 
% Define two macros that expand to the strings
% |\begin{thebibliography}| and |\end{thebibliography}|. Note that the
% idea for this and much else is inspired by the code for the docstrip
% package. Also define |\StartBibX| in case authors use the more
% standard |\begin{thebibliography}{}| syntax. 
%    \begin{macrocode}
\edef\EndBib{\expandafter\@gobble\string\\end\string{thebibliography\string}}
\edef\StartBib{\expandafter\@gobble\string\\begin\string{thebibliography\string}}
\edef\StartBibX{\expandafter\@gobble\string\\begin\string{thebibliography\string}\string{\string}}
%    \end{macrocode}
% Define two new logical variables. |\ifAnotherline| is true if we
% want to read another line from the input stream. It is set to false
% when we encounter the |\end{thebibliography}|. |\ifInBib| is set to
% true when we frind the |\begin{thebibliography}| and tells us to
% start saving the following lines. 
%    \begin{macrocode}
\newif\ifAnotherLine
\newif\ifInBib
%    \end{macrocode}
% \begin{macro}{makeOther}
% We need a macro to change the catcodes of lots of things to other. 
%    \begin{macrocode}
\def\makeOther#1{\catcode`#1=12\relax}
%    \end{macrocode}
% \end{macro}
% \begin{macro}{\RereadBiblio}
% Define a macro to do all the fiddly stuff. This is used at the
% beginning of |thebibliography| environment (Sec~\ref{sec:biblio}).
%    \begin{macrocode}
\newread\tf@biblio
\newcommand{\RereadBiblio}{%
%    \end{macrocode}
% Open the current file again for reading
%    \begin{macrocode}
  \begingroup
  \immediate\openin\tf@biblio=\articlename
%    \end{macrocode}
% Everything in a group to be safe. Initialize logicals. Black magic
% to change
% catcodes and avoid blank lines being interpreted as |\par|. 
%    \begin{macrocode}
  \AnotherLinetrue\InBibfalse
  \makeOther\ \makeOther\\\makeOther\$%
  \makeOther\#\makeOther\^\makeOther\^^K%
  \makeOther\_\makeOther\^^A\makeOther\%%
  \makeOther\~\makeOther\{\makeOther\}\makeOther\&%
  \endlinechar-1\relax
%    \end{macrocode}
% Start looping through the file. Read each line to the macro
% |\Line|. If we have got to EOF or end of biblio then set
% |\AnotherLinefalse|.
%    \begin{macrocode}
  \loop
  \read\tf@biblio to \Line
  \ifx\Line\EndBib\AnotherLinefalse\InBibfalse\fi
  \ifeof\tf@biblio \AnotherLinefalse\fi
%    \end{macrocode}
% If we are in the biblio and the line is not blank then append it to
% the |\ADS@references| macro together with a newline. 
%    \begin{macrocode}
  \ifInBib
  \ifx\Line\empty\else\xdef\ADS@references{\ADS@references\Line^^J}\fi
  \fi
%    \end{macrocode}
% Check if the line we just saw was the start of biblio, then repeat
% the loop if necessary. 
%    \begin{macrocode}
  \ifx\Line\StartBib\InBibtrue\fi
  \ifx\Line\StartBibX\InBibtrue\fi
  \ifAnotherLine\repeat
  \immediate\closein\tf@biblio
  \endgroup
  }
%    \end{macrocode}
% \end{macro}
% 
% \begin{macro}{\PadToFourWithDots}
% Utility function for left-padding a number to 4 characters with
% periods. Used in 
% constructing |\ADS@BibCode|. This is a very simple algorithm and
% \emph{requires} that the argument actually be a number---if it is
% not then dire consequences will result.  
%    \begin{macrocode}
\newcommand{\PadToFourWithDots}[1]{%
  \ifnum #1 > 999 #1%
  \else\ifnum #1 > 99 .#1%
  \else\ifnum #1 > 9 ..#1%
  \else ...#1%
  \fi\fi\fi}
%%  \ifnum \count0 > 999 #1%
%%  \else\ifnum \count0 > 99 .#1%
%%  \else\ifnum \count0 > 9 ..#1%
%%  \else ...#1%
%%  \fi\fi\fi}
%    \end{macrocode}
% \end{macro}
%
% \begin{macro}{\AuthorADS}
% \begin{macro}{\ADS@affiliations}
% New mechanism for setting the affiliations for ADS. The author labels (AA, AB, \dots) are generated using two nested counters.  We use the |\robustify| command from \texttt{etoolbox} so that accented characters in the affiliation addresses do not blow up.  The list of accents may need to be expanded. 
%    \begin{macrocode}
\newcounter{AUTHOR}
\newcounter{author}[AUTHOR]
\stepcounter{AUTHOR}
\newcommand\ADS@affiliations{}
\robustify{\'}\robustify{\i}\robustify{\~}\robustify{\`}
\def\AuthorADS#1#2{%
  \begingroup
  \def\ORCIDiD##1{<ID system="ORCID">##1</ID>}
  \def\Email##1{<EMAIL>##1</EMAIL>}
  \def\Affil##1{##1}
  \stepcounter{author}
  \ifthenelse{\value{author} > 26}{\stepcounter{AUTHOR}\stepcounter{author}}{}
  \ifthenelse{\equal{\ADS@affiliations}{\@rmaaEmpty}}{%
    \xdef\ADS@affiliations{\Alph{AUTHOR}\Alph{author}(#2)}
  }{%
    \xdef\ADS@affiliations{\ADS@affiliations, \Alph{AUTHOR}\Alph{author}(#2)}
  }%
  \endgroup
}
%    \end{macrocode}
% \end{macro}
% \end{macro}
% %
% \begin{macro}{\ADS@dumpdata}
% \changes{4.0}{2019/11/13}{Write all author affiliations for ADS}
% \begin{macro}{\ADS@write}
% \begin{macro}{\ADS@openfile}
% |\ADS@dumpdata| is the macro that actually does the job of writing out all the
% data to an auxiliary file, using |\ADS@write| to write each
% field. The file is set up with |\ADS@openfile|, which currently is
% done in the TOC of a book.  The new multi-author affiliations are used if |\ADS@affiliations| has been set, otherwise it falls back to the single-author version |\ADS@affil|. 
%    \begin{macrocode}
\newcommand{\ADSfield}[2]{\relax}
\newcommand{\ADSvfield}[2]{\relax}
%%\renewcommand{\ADSfield}[12]{\%#1 #2}
\newcommand{\ADS@dumpdata}
{ 
  \ADS@write{\string\begin{ADSarticle}}
  \ADS@WRITE{\string\ADSfield{R}{\ADS@BibCode}}
  \ADS@write{\string\ADSfield{A}{\ADS@authors}}
  \ifthenelse{\equal{\ADS@affiliations}{\@rmaaEmpty}}{%
    {\ADS@sanitize\ADS@write{\string\ADSfield{F}{AA(\ADS@affil)}}}
  }{%
    {\ADS@sanitize\ADS@write{\string\ADSfield{F}{\ADS@affiliations}}}
  }
  \ADS@write{\string\ADSfield{J}{\ADS@journalname}}
  \ADS@write{\string\ADSfield{D}{\ADS@fulldate}}
  \ADS@write{\string\ADSfield{L}{\ADS@lastpage}}
  {\ADS@sanitize\ADS@write{\string\ADSfield{T}{\ADS@title}}}
  \ADS@write{\string\ADSfield{C}{\ADS@copyright}}
  \ADS@write{\string\ADSfield{d}{\ADS@doi}}
  \ADS@write{\string\ADSfield{U}{\ADS@url}}
  \ADS@write{\string\ADSfield{K}{\ADS@keywords}}
  {\ADS@sanitize\ADS@write{\string\ADSfield{B}{\ADS@abstract}}}
  \ADS@write{\string\ADSfield{Z}{\ADS@references}}
  \ADS@write{\string\end{ADSarticle}}
  }
%%\newcommand{\ADS@write}[1]{\if@book\addtocontents{ads}{#1}\fi}
%%\newcommand{\ADS@openfile}{\@starttoc{ads}}
\newcommand{\ADS@write}[1]{\set@display@protect\immediate\write\tf@ads{#1}}
\newcommand{\ADS@WRITE}[1]{\immediate\write\tf@ads{#1}}
\newcommand{\ADS@openfile}{%
  \newwrite\tf@ads%
  \immediate\openout\tf@ads\jobname.ads\relax
  \ADS@WRITE{\ADS@preamble}}
\newcommand{\ADS@closefile}{\ADS@WRITE{\ADS@postamble}}
\def\ADS@preamble{%
  %%  \string\documentclass{article}
  %%  \string\begin{document}
  }
\def\ADS@postamble{%
  %%  \string\end{document}
  \string%\string% Local Variables:^J 
  \string%\string% mode: LaTeX^J
  \string%\string% mode: Font-lock^J
  \string%\string% End:^J
  }
%    \end{macrocode}
% \end{macro}
% \end{macro}
% \end{macro}
% 
% \subsection{Automatic calculation of the page range}
% \label{sec:page-range}
% Another issue that has cost me a lot of grief. Still, it finally
% seems to be sorted now. All hail Jeff Goldberg. 
%
% Cunning trick to get the page numbers from |\pageref|.  Put a label
% on the first and last page of the document using |\AtBeginDocument|
% and |\AtEndDocument|. .  This way won't work for the book option.
% In that case, we do it by defining |\@endofarticlehook| and
% |\@begofarticlehook| as parallels to |\AtEndDocument| and
% |\AtBeginDocument|.  Note that we have to defer adding
% |\@begofarticlehook| to |\AtBeginDocument| to make sure it comes
% after the redefinitions of |\ref| and |\label|.
%    \begin{macrocode}
\if@book\AtEndDocument{\ADS@closefile}\else
 %\AtBeginDocument{\@begofarticlehook}
\AtEndDocument{\@endofarticlehook}
\fi
%    \end{macrocode}
% \begin{macro}{\@begofarticlehook}
% Since we are getting more and more things that need to be done at
% the beginning of every article, frontsection, etc we define a sister
% hook to 
% |\@endofarticlehook|. This hook does the following: 
% \begin{enumerate}
% \item Labels the first page
% \item Sets |\ADS@firstpage|
% \item runs and then clears |\indexauthor@deferlist|
% \item resets the keyword counter
% \end{enumerate}
% It is executed in the following circumstances: 
% \begin{enumerate}
% \item At |\begin{document}| for standalone files
% \item At |\begin{document}| for each included file
% \item At beginning of front sections 
% \item At beginning of TOC
% \item At beginning of Author Index
% \item At beginning of abstracts
% \end{enumerate}
% For some reason, this is not always within the right group, so we
% use |\xdef| just to be sure. 
%    \begin{macrocode}
\newcommand{\@begofarticlehook}{%
  \if@frontmatter\frontlabel{RMAAFirstPage}\else\label{RMAAFirstPage}\fi%
  \xdef\ADS@firstpage{\@arabic{\c@page}}%
  \if@book\else\makeindex\fi % needed so default index is defined in stand-alone articles
%%  \indexauthor@deferlist
%%  \xdef\indexauthor@deferlist{}
  \@inpreamblefalse
  \setcounter{kwdcounter}{0}
  }
%    \end{macrocode}
% \end{macro}
% \begin{macro}{\@endofarticlehook}
% This is a hook to be run at the end of each article
% currently, it just labels the last page and writes out all the ADS
% data to an auxiliary file.  
% It should not produce any printed output. 
% It is modelled on the |\AtEndDocument| command, to which it reduces
% in the case of a single article. 
%    \begin{macrocode}
\newcommand{\@endofarticlehook}{%
  \if@debug{\ClassWarning{rmaa}{***Doing Hook***}}\fi
  \if@debug{\ClassWarning{rmaa}{ADSabstract is: \ADS@abstract}}\fi
  \clearpage\rmaa@lastpagelabel\if@book\ADS@dumpdata\fi
  \ifx\NotesList\@rmaaEmpty\else\OutputNotes\fi
  \if@debug{\ClassWarning{rmaa}{Notes list is: \NotesList}}\fi
}
%    \end{macrocode}
% \begin{macro}{\rmaa@lastpagelabel}
% New version of last page label command copied from Jeff Goldberg's
% lastpage.sty package - we can't just use |\label| here in case the
% output routine for the last page has already been called. Instead,
% we first explicitly clear the page, step back the page counter,
% explicitly write to the .aux file, then restep the page counter
% again.
%
% But this will conflict with the hyperref package, which does labels differently.  Compensate for that by using the five-argument form.
% 
% Also set the |\ADS@LastPage| macro to the appropriate number. 
%    \begin{macrocode}
\newcommand{\rmaa@lastpagelabel}{%
  \addtocounter{page}{-1}%
  \immediate\write\@auxout{\string
    \newlabel{\articlename:RMAALastPage}{{}{\thepage}{}{}{}}}%
  \edef\ADS@lastpage{\thepage}%
  \addtocounter{page}{1}}
%    \end{macrocode}
% \end{macro}
% \end{macro}

% \subsection{Deal with explicit change in number of columns}
% \label{sec:columns}
% Sometimes, when formatting a book, we need to switch between
% one-column articles and two-column articles. For this to work, we
% need to wrap the standard \LaTeX{} commands to set the |\@rmtwocol|
% flag appropriately and also reset the package default values of all
% scalable lengths. .  
%    \begin{macrocode}
\let\latexonecolumn=\onecolumn
\let\latextwocolumn=\twocolumn
\renewcommand{\onecolumn}{\@rmtwocolfalse\@package@default@lengths\latexonecolumn}
\renewcommand{\twocolumn}{\@rmtwocoltrue\@package@default@lengths\latextwocolumn}
%    \end{macrocode}
% 
% \subsection{The actual formatting of the title}
% \label{sec:title-format}

% Note, we have to redefine the |\maketitle| command to use the
% numbered footnotes. This means we needn't use the atbeginend stuff. 
%    \begin{macrocode}
\renewcommand{\maketitle}{\par
  \begingroup
%    \end{macrocode}
% Don't switch to marked footnotes. 
% This is consistent with existing style but I'm not sure I like it
%    \begin{macrocode}
  %%  \renewcommand\thefootnote{\fnsymbol{footnote}}%
%%  \def\@makefnmark{\hbox to\z@{$\m@th^{\@thefnmark}$\hss}}%
  \def\@makefnmark{\textsuperscript{\@thefnmark}\hspace{-0.2em}}%
  \long\def\@makefntext##1{\parindent 1em\noindent
    \hbox to1.8em{\hss$\m@th^{\@thefnmark}$}##1}%
  \if@rmtwocol
  { \if@debug{\ClassWarning{rmaa}{TWO COLUMN}}\fi
    \twocolumn[\@maketitle]%
    }
  \else
  { \if@debug{\ClassWarning{rmaa}{ONE COLUMN}}\fi
    \newpage
    \global\@topnum\z@   % Prevents figures from going at top of page.
    \if@manuscript \@MSmaketitle \else \@maketitle \fi
    }
  \fi
  \if@manuscript\else\thispagestyle{rmaafirstpage}\@thanks\fi
  \ifx\NotesList\@rmaaEmpty\else{%
    \NoteMark{1--\arabic{notecounter}}\vspace*{-2\baselineskip}%
  }\fi%
  \endgroup
%    \end{macrocode}
% Don't reset footnote counter (not sure whether this is good thing
% or not).
%    \begin{macrocode}
  %%  \setcounter{footnote}{0}%
  \let\thanks\relax\let\affil\relax\let\altaffil\relax
  \let\maketitle\relax\let\@maketitle\relax
  \gdef\@thanks{}\gdef\@author{}\gdef\@title{}}
%    \end{macrocode}

% \begin{macro}{\@kwhwidth}
% \begin{macro}{\@kwtwidth}
% Redefine \cs{@maketitle} to include the abstract, resumen and
% keywords. Note that in the manuscripty style we use
% \cs{@MSmaketitle} instead.  
% \changes{3.16}{2003/04/07}{
%  Added extra vertical space between authors and accepted date.}
% \changes{3.26}{2011/08/15}{Add extra label to review articles}
%    \begin{macrocode}
\newlength{\@kwhwidth}\newlength{\@kwtwidth}
\newcommand\reviewname{Review}
\renewcommand{\@maketitle}{%
  \newpage
  \null
  \if@reviewarticle
  \smallskip
  {\large\bfseries\itshape\reviewname}%
  \bigskip
  \else
  \vskip\@before@title@skip%
  \fi
  \begin{center}%
    {\@title \par}%
    \vskip\@before@author@skip%
    \parbox{\rmaa@AuthorWidthFrac\textwidth}{\centering
      \lineskip\@between@author@skip%
      \@author
      \par}
    %% this should be replaced with `Submitted to...' for preprints
    \if@printdates
    \vskip\@before@dates@skip%
    \if@preprint
    \rmaa@ReceivedStyle{Submitted to RMxAA \@received}
    \else
    \rmaa@ReceivedStyle{Received \@received; accepted \@accepted}
    \fi
    \fi
    \if@manuscript
    \rmaa@ReceivedStyle{Submitted to RMxAA \@received}
    \fi
    \if@guide
    \rmaa@ReceivedStyle{Version \rmaa@version, \rmaa@date}
    \fi
    \vskip\@before@abstract@skip%
    \if@abstractpresent
    \parbox{\rmaa@AbstractWidthFrac\textwidth}{%
      \setlength{\parindent}{0.9cm}%
      \centerline{\rmaa@AbstractHeaderStyle{Abstract}}\par%
      \vskip\@after@abstracthead@skip%
      \if@newstyle\noindent\fi\@abstract
      \vskip\@between@abstract@skip%
      \centerline{\rmaa@AbstractHeaderStyle{Resumen}}\par%
      \vskip\@after@abstracthead@skip%
      \if@newstyle\noindent\fi\@resumen
      \if@setresumen\else\@abstract\fi % append english abstract if
                                       % spanish one not set
      \if@guide\else            %do not use keywords for guide
      \vskip\@before@kwd@skip%
      % Format the Key Words with hanging indentation:
      % They are put in a \parbox with width equal to the line width
      % minus the width of the header. 
      { % Widths of keyword
        \settowidth{\@kwhwidth}{\@keywordheader~}   % header and
        \setlength{\@kwtwidth}{\linewidth}          % keyword text
        \addtolength{\@kwtwidth}{-\@kwhwidth}       %
        \noindent\@keywordheader~\parbox[t]{\@kwtwidth}{%
          \rmaa@KeywordStyle{\boldmath\uppercase{\@keywords}}}}%
      \fi
      }
      \fi
  \end{center}%
  \vspace*{\afterkwdskip}
  \par
  }
\newlength{\afterkwdskip}\setlength{\afterkwdskip}{\z@}
%    \end{macrocode}
% \end{macro}
% \end{macro}
%
%
% \begin{macro}{\reviewarticle}
% \begin{macro}{\reviewtoc}
% \begin{macro}{\reviewlof}
%   User command to specify a review article. Must be used in
%   preamble. Main effect is to place the text ``Review'' before the
%   article title.  Also add command for setting a table of contents
%   and a list of figures at the beginning of the article.
%    \begin{macrocode}
\newcommand{\reviewarticle}{\@reviewarticletrue}
\newcommand\reviewtoc{%
  \section*{CONTENTS}%
  \@starttoc{rtc}%
}
\newcommand\reviewlof{%
  \section*{LIST OF FIGURES}%
  \@starttoc{lof}%
}
%    \end{macrocode}
% \end{macro}
% \end{macro}
% \end{macro}
% \begin{macro}{\l@section}
%    \begin{macrocode}
\renewcommand*\l@section[2]{%
  \ifnum \c@tocdepth >\z@
    \addpenalty\@secpenalty
    \addvspace{0.8em \@plus\p@}%
    \setlength\@tempdima{1.5em}%
    \begingroup
      \parindent \z@ \rightskip \@pnumwidth
      \parfillskip -\@pnumwidth
      \leavevmode %\bfseries
      \advance\leftskip\@tempdima
      \hskip -\leftskip
      #1\nobreak\hfil \nobreak\hb@xt@\@pnumwidth{\hss #2}\par
    \endgroup
  \fi}
%    \end{macrocode}
% \end{macro}
%
% \begin{macro}{\@MSmaketitle}
% A special version of |\@maketitle| for the \texttt{manuscript}
% option. This can deal with a multi-page abstract/resumen. Not very
% well, mind. 
%    \begin{macrocode}
\newcommand{\@MSmaketitle}{%
  \if@debug\ClassInfo{rmaa}{Using special MS maketitle command}\fi
  \newpage
  {
    \thispagestyle{rmaafirstpage}
%     \setlength{\parindent}{0pt}
    \null\vskip 2ex
    \begin{center}
      \@title\par
      \@author\par
      \rmaa@ReceivedStyle{%
        \ifx\@received\@rmaaEmpty
          Draft version: \@date
        \else
          Received: \@received
          \ifx\@accepted\@rmaaEmpty
          \else
            ; Accepted: \@accepted
          \fi
        \fi}\par
    \end{center}
    \@thanks
    \centerline{\rmaa@AbstractHeaderStyle{Resumen}}
    \noindent\@resumen\par\vskip 2ex
    \centerline{\rmaa@AbstractHeaderStyle{Abstract}}
    \noindent\@abstract\par\vskip 2ex
    \noindent\@keywordheader~\rmaa@KeywordStyle{\@keywords}\par
%     \@keywordheader~\rmaa@KeywordStyle{\uppercase{\@keywords}}\par
  }%
  }
%    \end{macrocode}
% \end{macro}

% \subsection{The sectioning commands}
% \label{sec:section-commands}

% \begin{macro}{\section}
% \begin{macro}{\subsection}
% \begin{macro}{\subsubsection}
% \begin{macro}{\paragraph}
% \begin{macro}{\subparagraph}
% Redefine all the standard commands to use these lengths. 
%    \begin{macrocode}
\renewcommand\section{\@startsection {section}{1}{\z@}%
                                   {\abovesectionskip}%
                                   {\belowsectionskip}%
                                   %% how do we get upper case here
                                   {\centering\reset@font\normalsize}}
\renewcommand\subsection{\@startsection{subsection}{2}{\z@}%
                                     {\abovesectionskip}%
                                     {\belowsectionskip}%
                                     {\centering\reset@font\normalsize\itshape}}
\renewcommand\subsubsection{\@startsection{subsubsection}{3}{\z@}%
                                     {\abovesectionskip}%
                                     {\belowsectionskip}%
                                     {\centering\reset@font\normalsize\itshape}}
\renewcommand\paragraph{\@startsection{paragraph}{4}{\z@}%
                                    {\aboveparaskip}%
                                    {\belowparaskip}%
                                    {\reset@font\normalsize\bfseries}}
\renewcommand\subparagraph{\@startsection{subparagraph}{5}{\parindent}%
                                      {\aboveparaskip}%
                                      {\belowparaskip}%
                                      {\reset@font\normalsize\bfseries}}
%    \end{macrocode}
% \end{macro}
% \end{macro}
% \end{macro}
% \end{macro}
% \end{macro}

% \begin{macro}{secnumdepth}
% Make sure that we number the subsubsections---this is done by
% default in the |article| class but not the |book| class, which is
% why we must set it explicitly here.
%    \begin{macrocode}
\setcounter{secnumdepth}{3}
%    \end{macrocode}
% \end{macro}
% \changes{2.2beta10}{2001/11/07}{Fixed \cs{secnumdepth} for \cs{book}
% style.}

% \begin{macro}{\thesubsection}
% \begin{macro}{\thesubsubsection}
% For subsection and subsubsection, we want the numbers to be roman
% but the title to be italic. This seems the easiest way to go about it.
%    \begin{macrocode}
\renewcommand{\thesubsection}{{\upshape\thesection.\@arabic\c@subsection}}
\renewcommand{\thesubsubsection}{\thesubsection{\upshape.\@arabic\c@subsubsection}}
%    \end{macrocode}
% \end{macro}
% \end{macro}

% \begin{macro}{\@seccntformat}
% Add ``.'' to the number here so we can remove it from |\thesection|,
% etc. Reduce gap between number and title WJH 27 Jan 1995
%    \begin{macrocode}
\def\@seccntformat#1{\csname the#1\endcsname .\hskip 0.4em}
%    \end{macrocode}
% \end{macro}

% \begin{macro}{\@sect}
% I didn't really want to mess with the |\@sect| command, but it was
% the only way I could manage to force upper case. Other than the
% commented parts, it is identical to the standard version. 
%    \begin{macrocode}
\def\@sect#1#2#3#4#5#6[#7]#8{
  \ifnum #2>\c@secnumdepth
     \let\@svsec\@empty
  \else
     \refstepcounter{#1}%
     \let\@@protect\protect
     \def\protect{\noexpand\protect\noexpand}%
     \edef\@svsec{\@seccntformat{#1}}%
     \let\protect\@@protect
  \fi
%    \end{macrocode}
% This is the bit I changed (27 Jan 1995). Use upper case for
% |\section|, use mixed case for the rest. 
%    \begin{macrocode}
  \ifnum #2=1                   
     \def\sec@upcase##1{\uppercase{##1}}%
     \def\sec@shape{\upshape}%
  \else                         
     \def\sec@upcase##1{\relax{##1}}%
     \def\sec@shape{\itshape}%
  \fi
  \@tempskipa #5\relax
  \ifdim \@tempskipa>\z@
     \begingroup #6\relax
          \@hangfrom{\hskip #3\relax\@svsec}%
                    {\interlinepenalty \@M \sec@upcase{#8}\par}%
     \endgroup
     \csname #1mark\endcsname{#7}%
%    \end{macrocode}
% Only write TOC entries if doing the Guide style. 
%    \begin{macrocode}
     \if@guide 
        \addcontentsline{toc}{#1}{%
          \ifnum #2>\c@secnumdepth \else
             \protect\numberline{\csname the#1\endcsname}%
          \fi
          #7}%
     \fi
%    \end{macrocode}
% Also allow per-article TOC for review papers
%    \begin{macrocode}
     \if@reviewarticle
        \ifnum #2>2 \else
           \addcontentsline{rtc}{#1}{%
             \protect\numberline{\csname the#1\endcsname}%
             \sec@upcase{\sec@shape #7}}%
        \fi
     \fi
  \else
     \def\@svsechd{#6\hskip #3\relax  %% \relax added 2 May 90
       \@svsec #8\csname #1mark\endcsname{#7}%
       \if@guide
          \addcontentsline{toc}{#1}{%
            \ifnum #2>\c@secnumdepth \else
               \protect\numberline{\csname the#1\endcsname}%
            \fi
            #7}%
       \fi
       }%
  \fi
  \@xsect{#5}}
%    \end{macrocode}
% \end{macro}


% \subsection{Appendices}
% \label{sec:appendices}
% Make sure that appendices in individual articles work properly. We
% copy the definition from the article class (the book one doesn't do
% what we want).  
% \begin{environment}{appendices}
%   This uses an environment |appendices| inside of which each
%   |\section| will become a separate appendix.  Note that the
%   appendices must come before the references (otherwise the full
%   addresses mechanism will not work properly). We make the equation
%   numbering be subordinate to the appendix (equations A12, B6,
%   etc). Finally, we have a medskip followed by APPENDICES to
%   introduce them. \changes{3.17}{2002/04/08}{ Added appendices environment
%   that does the right thing.}
%    \begin{macrocode}
\newenvironment{appendices}{\par
  \begingroup
  \setcounter{section}{0}%
  \setcounter{subsection}{0}%
  \renewcommand\thesection{\@Alph\c@section}
  \renewcommand\theequation{\thesection\arabic{equation}}
  \medskip
  \section*{APPENDICES}}
{\endgroup}
%    \end{macrocode}
% \end{environment}
% \begin{environment}{appendix}
%   If there is only one appendix, then use |appendix| environment
%   instead.  \changes{4.0}{2019/11/12}{ Added
%   appendix environment for case where there is only one section.}
%    \begin{macrocode}
\renewenvironment{appendix}{\par
  \begingroup
  \setcounter{section}{0}%
  \setcounter{subsection}{0}%
  \renewcommand\thesection{\@Alph\c@section}
  \renewcommand\theequation{\thesection\arabic{equation}}
  \medskip
  \section*{APPENDIX}}
{\endgroup}
%    \end{macrocode}
% \end{environment}

% \subsection{Floating thingies}
% \label{sec:floats}   

%% 
%    \begin{macrocode}
\renewcommand{\thetable}{\@arabic\c@table}
\renewcommand{\thefigure}{\@arabic\c@figure}
\renewcommand{\figurename}{Fig.}
\renewcommand{\tablename}{TABLE}
%\renewcommand{\fnum@table}{\MakeUpperCase{\tablename}~\thetable}
%    \end{macrocode}
% For some reason, I have to write `TABLE' explicitly here
% If I use |\tablename|, it gets converted to mixed case---God knows
% why. 
% |\figurename| is also being ignored. 
%    \begin{macrocode}
\renewcommand{\fnum@table}{TABLE~\thetable}
\renewcommand{\fnum@figure}{Fig.~\thefigure}
%    \end{macrocode}

% Do captions in 9 pt face in the main journal style and in new-style
% anything. 
%    \begin{macrocode}
\long\def\@makefigcap#1#2{%
  \if@proceedings\else\small\fi
  \if@newstyle\small\fi 
  \vskip\abovecaptionskip
  \sbox\@tempboxa{#1. #2}%
  \ifdim \wd\@tempboxa >\hsize
    #1. #2 \par
  \else
    \global \@minipagefalse
    \hb@xt@\hsize{\hfil\box\@tempboxa\hfil}%
  \fi
  \vskip\belowcaptionskip}
%    \end{macrocode}

% Make sure table caption is upper case. 
%    \begin{macrocode}
\long\def\@maketabcap#1#2{%
  \hbox to\hsize{\hfil #1 \hfil}
  \vskip 6pt 
  \centering\MakeUpperCase{#2}\par
  \vskip 6pt}
%    \end{macrocode}

%    \begin{macrocode}
\newlength{\@rottabcapwidth}
\long\def\@makerottabcap#1#2{%
  \setbox\@tempboxa\hbox{\MakeUpperCase{#2}}%
  \ifdim\wd\@tempboxa >.8\vsize
     \setlength\@rottabcapwidth{.8\vsize}
  \else
     \setlength\@rottabcapwidth{\wd\@tempboxa}
  \fi
  \rotatebox{90}{%
    \begin{minipage}{\vsize}\centering
      \begin{minipage}{\@rottabcapwidth}\centering
        #1\\ \vskip 6pt \MakeUpperCase{#2}
      \end{minipage}%
    \end{minipage}
    }%
  \hskip 6pt
  }

\newenvironment{Sideways}
    {\begin{sideways}\begin{minipage}{\vsize}\centering}
    {\end{minipage}\end{sideways}}

\newenvironment{Sidewaystable}
{%
  \let\oldcaption=\caption%
  \def\caption{
    \def\@captype{table}\renewcommand{\@makecaption}{\@maketabcap}
    \oldcaption}
  \begin{sidewaystable}\centering}
  {\end{sidewaystable}\let\caption=\oldcaption}
\newenvironment{Sidewaystable*}
{%
  \let\oldcaption=\caption%
  \def\caption{
    \def\@captype{table}\renewcommand{\@makecaption}{\@maketabcap}
    \oldcaption}
  \begin{sidewaystable*}\centering}
  {\end{sidewaystable*}\let\caption=\oldcaption}
%    \end{macrocode}

% Not sure why we need to stick these lines in. 
%    \begin{macrocode}
\renewenvironment{figure}
                 {\renewcommand{\@makecaption}{\@makefigcap}\@float{figure}}
                 {\end@float}
\renewenvironment{figure*}
                 {\renewcommand{\@makecaption}{\@makefigcap}\@dblfloat{figure}}
                 {\end@dblfloat}
\renewenvironment{table}
                 {\renewcommand{\@makecaption}{\@maketabcap}\@float{table}}
                 {\end@float}
\renewenvironment{table*}
                 {\renewcommand{\@makecaption}{\@maketabcap}\@dblfloat{table}}
                 {\end@dblfloat}
%    \end{macrocode}

% Redefine |\@caption| so that we don't bother with the list of
% figures. This saves us having to worry about fragile commands in the
% captions. 
%    \begin{macrocode}
\long\def\@caption#1[#2]#3{%
  \par
%%  \addcontentsline{\csname ext@#1\endcsname}{#1}%
%%    {\protect\numberline{\csname the#1\endcsname}{\ignorespaces #2}}%
  \begingroup
    \@parboxrestore
    \if@minipage
      \@setminipage
    \fi
    \normalsize
    \@makecaption{\csname fnum@#1\endcsname}{\ignorespaces #3}\par
  \endgroup}
%    \end{macrocode}

% Make tables not quite so cramped. Allow |\hline\hline| to give a
% poor man's thick rule. 
%    \begin{macrocode}
\renewcommand{\arraystretch}{1.15}
\setlength{\doublerulesep}{0.4pt}
%    \end{macrocode}

% \begin{macro}{\tablecols}
% \begin{macro}{\@tablecols}
% \begin{macro}{\tabcaptioncontinued}
% |\tablecols| is a user command for setting the number of columns. It
% is used by commands that try and span the entire table (I think
% there must be a better way of acheiving
% this). |\tabcaptioncontinued| is 
% a user command for setting the caption of a continued table on the
% next page. 
%    \begin{macrocode}
\newcommand{\tablecols}[1]{\renewcommand{\@tablecols}{#1}}
\newcommand{\@tablecols}{1}
\newcommand{\tabcaptioncontinued}{%
  \multicolumn{\@tablecols}{c}{\addtocounter{table}{-1}TABLE~\thetable\ (CONTINUED)} \\ \\
  }
%    \end{macrocode}
% \end{macro}
% \end{macro}
% \end{macro}

% \subsubsection{Footnotes in tables}
% \label{sec:table-footnotes}
% \begin{macro}{\tablenotemark}
% \begin{macro}{\tablenote}
% \begin{macro}{\tablenoteskip}
% Old user commands |\tablenotemark{MARK}|,
% |\tablenote{MARK}{TEXT}|---now deprecated. 
%    \begin{macrocode}
\newcommand\tablenotemark[1]{$^\mathrm{#1}$} 
\newcommand\tablenote[2]{\if@proceedings\else\small\fi ${}^\mathrm{#1}$ #2} 
\newcommand\tablenoteskip{\\[0.5\baselineskip]} 
%    \end{macrocode}
% \end{macro}
% \end{macro}
% \end{macro}

% \begin{macro}{\tabnotemark}
% \begin{macro}{\tabnotetext}
% \begin{macro}{\tabnotetextTWO}
% Alternative way of doing footnotes in tables. 
%    \begin{macrocode}
\newcommand{\tabnotemark}[1]{\textsuperscript{\lowercase{#1}}}
%\newcommand{\tabnotetextTWO}[4]{\tabnotetext{#1}{#2}\tabnotetext{#3}{#4}}
\newcommand{\tabnotetextTWO}[4]{%
  \multicolumn{\@tablecols}{l}{%
    \parbox{0.5\tabnotewidth}{\small\textsuperscript{#1}#2}%
    \parbox{0.5\tabnotewidth}{\small\textsuperscript{#3}#4}%
    }%
  \\}
\newcommand{\tabnotetext}[2]{%
  \multicolumn{\@tablecols}{@{}p{\tabnotewidth}@{}}{\small\textsuperscript{#1}#2}\\}
\newlength{\tabnotewidth}
%    \end{macrocode}
% \end{macro}
% \end{macro}
% \end{macro}


% \begin{macro}{changemargin}
%    \begin{macrocode}
% Command to allow margins to be changed on the fly. Useful to allow
% tables and Figures to protrude into the margins. 
%
% \changes{3.24}{2007/09/14}{Added \texttt{changemargin} environment
% (copied from TeX FAQ).}
\newenvironment{changemargin}[2]{% 
  % From TeX FAQ http://www.tex.ac.uk/cgi-bin/texfaq2html?label=chngmargonfly
  \begin{list}{}{%
    \setlength{\topsep}{0pt}%
    \setlength{\leftmargin}{#1}%
    \setlength{\rightmargin}{#2}%
    \setlength{\listparindent}{\parindent}%
    \setlength{\itemindent}{\parindent}%
    \setlength{\parsep}{\parskip}%
  }%
  \item[]}{\end{list}}
%    \end{macrocode}
% \end{macro}

% \subsection{The bibliography and references}
% \label{sec:biblio}
% \begin{macro}{\beforebiblioskip}
% \begin{macro}{\afterbiblioskip}
% I'm not entirely sure why I introduced these two commands, since
% directly putting |\vspace| before and after
% |\begin{thebibliography}| seems to work well enough. 
%    \begin{macrocode}
\newlength{\beforebiblioskip}
\setlength{\beforebiblioskip}{0pt}
\newlength{\afterbiblioskip}
\setlength{\afterbiblioskip}{0pt}
%    \end{macrocode}
% \end{macro}
% \end{macro}
% \begin{macro}{thebibliography}
% The environment itself is done as starred section with extra
% vertical space after heading. This space is currently zero. 
% In the book style we do the cunning trick with rereading the file in
% order to get the references into the |\ADS@references| macro.
% 
% \changes{3.05}{2003/01/02}{Removed the \cs{beforebiblioskip} and
% \cs{afterbiblioskip} commands, since they cause unwanted vertical
% space in some instances.}  
%    \begin{macrocode}
\renewenvironment{thebibliography}{%
  \if@book\RereadBiblio\fi
%%  \vspace*{\beforebiblioskip}
  \section*{REFERENCES} \frenchspacing
%%  \vspace*{\afterbiblioskip}
  \small
%    \end{macrocode} 
% The references themselves are set as a list in order 
% to get the hanging indentation. They are done in a smaller font and
% with french spacing so we needn't worry about the distinction
% between sentence-ending periods and abbreviations. 
%    \begin{macrocode} 
  \list{\null}{\leftmargin 1.5em\labelwidth\z@\labelsep\z@\itemindent -1.5em
    \itemsep\z@\parsep\z@\usecounter{enumi}}
  \sloppy\clubpenalty4000\widowpenalty4000
  \sfcode`\.=1000\relax\@biblisttrue}
{\endlist\@biblistfalse\normalsize\@outputaddressesifnecessary}
%    \end{macrocode}
% \end{macro}
% \subsection{Uniquification of cross-references}
% \label{sec:uniquification}
% 
% \begin{macro}{\articlename}
% In order that the |\label{}| command in individual articles be
% guaranteed to give a unique anchor, we append the article name (held
% in the |\articlename| command) to each key. This requires
% redefinition of the following commands: |\cite|, |\nocite|, |\ref|,
% |\pageref|, |\label|. 
% \begin{macro}{\oldcite}
% \begin{macro}{\oldnocite}
% \begin{macro}{\oldref}
% \begin{macro}{\oldpageref}
% \begin{macro}{\oldlabel}
% The original versions of these commands are
% saved as |\oldcite|, |\oldnocite|, etc. 
%
% The uniquifying of cross-references is now always done, so 
% we don't get multiply-defined labels in the single article case
%    \begin{macrocode}
\newcommand{\articlename}{anon}
\AtBeginDocument{
  \let\oldcite=\cite 
  \let\oldnocite=\nocite 
  \let\oldref=\ref 
  \let\oldpageref=\pageref
  \let\oldlabel=\label
}
%    \end{macrocode}
% \end{macro}
% \end{macro}
% \end{macro}
% \end{macro}
% \end{macro}
% \end{macro}
%
% Uniquify all cross-references by prepending the article name.  This
% now has to go into |\AtBeginDocument| so it is working on the new
% versions that hyperref.sty has set up.  Then we need to put the
% |\@begofarticlehook| afterwards, so that the first page label is using
% the new definitions.  If amsmath is loaded, then we also need to patch
% the special label command that they use. \changes{4.0}{2019/11/18}{Fix labels for amsmath display equations, in case that package has been loaded.}
%    \begin{macrocode}
\AtBeginDocument{
  \renewcommand{\ref}[1]{\oldref{\articlename:#1}%
    \if@debug{\ClassInfo{rmaa}{Reference to \articlename:#1}}\fi}
  \renewcommand{\pageref}[1]{\oldpageref{\articlename:#1}%
    \if@debug{\ClassInfo{rmaa}{Page reference to \articlename:#1}}\fi}
  \renewcommand{\label}[1]{\oldlabel{\articlename:#1}%
    \if@debug{\ClassInfo{rmaa}{Label for \articlename:#1}}\fi}%
  \ifdef{\label@in@display}{%
    \let\oldlabel@in@display=\label@in@display
    \renewcommand{\label@in@display}[1]{\oldlabel@in@display{\articlename:#1}%
      \if@debug{\ClassInfo{rmaa}{AMSTeX Label for \articlename:#1}}\fi}%
  }{}
  \@begofarticlehook
}%
%    \end{macrocode}
% Citations are more complicated since they can contain a list of
% keys. These are really only here for use with the rmaacite package,
% which is utterly deprecated (02 Jan 2003) now that I have the natbib
% code working. The command \cs{cite} should \emph{not} be used with
% natbib: use \cs{citep}, \cs{citet}, etc.\@ instead. 
%    \begin{macrocode}
\renewcommand{\nocite}[1]{\@bsphack
  \@for\@citeb:=#1\do{%
    \edef\@citeb{\expandafter\@firstofone\@citeb:\articlename}%
    \if@filesw\immediate\write\@auxout{\string\citation{\@citeb}}\fi
    \@ifundefined{b@\@citeb:\articlename}{\G@refundefinedtrue
      \@latex@warning{Citation `\@citeb' undefined}}{}}%
  \@esphack}
\renewcommand{\cite}[1]{\oldcite{#1:\articlename}}
%   \end{macrocode}


% Horrible fiddling with the bibliography citation mechanism. 
% This has to be globally available in case rmaacite puts it in the
% \texttt{.aux} file. 
%    \begin{macrocode}
\def\lbibcite#1#2#3{\global\@namedef{b@#1}{#2}
                    \global\@namedef{by@#1}{#3}}
%    \end{macrocode}

% \begin{macro}{\oldbibitem}
% \begin{macro}{\@rmaaBibArg}
% \begin{macro}{\@rmaaSpace}
% \begin{macro}{rmaabib}
% Of course, |\bibitem| needs to get the uniquification treatment as
% well. While we are at it, we redefine it to use consecutive integers
% for the key if the user leaves it blank (using the counter
% |rmaabib|). This allows one to put 
% |\bibitem{}| for all your references without getting annoying
% messages about multiply-defined labels.
% 
% We also take the opportunity to insert the optional argument to
% \cs{bibitem} if it is missing in the form Author(DATE) so that natbib
% doesn't try to do numerical citations. Natbib already provides a
% hook \cs{@extra@b@citeb} to help uniquify the citation key, so we
% just plug into that. However, it puts it \emph{after} the original
% key, not before as I had had it. Easiest solution is to change the
% way that I do it. 02 Jan 2003. Note that the \cs{@extra@b@citeb}
% does not get written to the .aux  file in the \cs{citation} command
% but it does work. 
% \changes{3.14}{06 Feb 2003}{Reinstalled old definition of \cs{bibitem}
%  for the case where the nonatbib option is requested, so that the
%  rmaacite package will work. This means that you can't use rmaacite
%  and natbib in the same volume.}
%    \begin{macrocode}
\if@natbib\def\@extra@b@citeb{:\articlename}\fi
%%\else
\let\oldbibitem\bibitem
\newcommand{\@rmaaBibArg}{}
\newcommand{\@rmaaSpace}{ }
\newcounter{rmaabib}
\if@natbib
\renewcommand{\bibitem}[2][ABC(2003)]{%
  \renewcommand{\@rmaaBibArg}{#2}
  \ifx\@rmaaBibArg\@rmaaEmpty
     \addtocounter{rmaabib}{1}
     \oldbibitem[#1]{\thermaabib:\articlename}
  \else\ifx\@rmaaBibArg\@rmaaSpace
     \addtocounter{rmaabib}{1}
     \oldbibitem[#1]{\thermaabib:\articlename}
  \else 
     \oldbibitem[#1]{#2:\articlename}
  \fi\fi
  \renewcommand{\@rmaaBibArg}{}
  }
\else
\renewcommand{\bibitem}[1]{%
  \renewcommand{\@rmaaBibArg}{#1}
  \ifx\@rmaaBibArg\@rmaaEmpty
     \addtocounter{rmaabib}{1}
     \oldbibitem{\articlename:\thermaabib}
  \else\ifx\@rmaaBibArg\@rmaaSpace
     \addtocounter{rmaabib}{1}
     \oldbibitem{\articlename:\thermaabib}
  \else
     \oldbibitem{\articlename:#1}
  \fi\fi
  \renewcommand{\@rmaaBibArg}{}
  }
\fi
%%\fi
%    \end{macrocode}
% \end{macro}
% \end{macro}
% \end{macro}
% \end{macro}

% \begin{macro}{oldpcite}
% \begin{macro}{oldscite}
% Then, we need to repeat the above scheme, but for the |rmaacite|
% definition. 
%    \begin{macrocode}
\newcommand{\@redefinebibitem}{%
  \let\oldbibitem\bibitem
  \let\oldcite\cite
  \let\oldpcite\pcite
  \let\oldscite\scite
  \def\bibitem[##1<##2>]##3{%
    \renewcommand{\@rmaaBibArg}{##3}
    \ifx\@rmaaBibArg\@rmaaEmpty
       \addtocounter{rmaabib}{1}
       \oldbibitem[##1<##2>]{\thermaabib:\articlename}
    \else\ifx\@rmaaBibArg\@rmaaSpace
       \addtocounter{rmaabib}{1}
       \oldbibitem[##1<##2>]{\thermaabib:\articlename}
    \else 
       \oldbibitem[##1<##2>]{##3:\articlename}
    % trust the author here
%%       \oldbibitem[##1<##2>]{##3}
    \fi\fi
    }
  }
%    \end{macrocode}
% \end{macro}
% \end{macro}


% \subsection{HTML Table of Contents}
% \label{sec:hyper-TOC}
% \begin{macro}{\HyperTOCEntry}
% \begin{macro}{\HyperTOCpart}
% \begin{macro}{\HyperTOCmeeting}
% \begin{macro}{\HyperTOCeditors}
% \begin{macro}{\HyperTOCpreamble}
% These commands are used to write information to the \texttt{.htoc}
% file, which is used in generating a Web Page of the Table of
% Contents. Since they only ever appear preceded by |\protect| in the
% argument to an |\addtocontents| command,  they should never need to
% be expanded by the macros. However, we define them here just in case
% \dots
%    \begin{macrocode}
\newcommand{\HyperTOCEntry}[3]{\relax}
\newcommand{\HyperTOCpart}[1]{\relax}
\newcommand{\HyperTOCmeeting}[3]{\relax}
\newcommand{\HyperTOCeditors}[1]{\relax}
\newcommand{\HyperTOCpreamble}{\relax}
%    \end{macrocode}
% \end{macro}
% \end{macro}
% \end{macro}
% \end{macro}
% \end{macro}
%

% \subsection{Formatting the poster abstracts} 
% These can now be used without using the book option.
%    \begin{macrocode}

\newcommand{\@makeminititle}{%
  \newpage
  \null 
  \vskip 3em%
  \begin{center}%
    {\@title \par}%
  \end{center}
  \vskip 3em%
  }

\newif\if@abstractsads
\@abstractsadsfalse
\newcommand\abstractsADS{\@abstractsadstrue}

\if@book\RequirePackage{multicol}\fi
\newenvironment{abstracts}[1][Abstracts of Contributed Papers]{%
  %% First, make sure we start on an odd page (06 Jun 2000)
  \if@book\clearpage\thispagestyle{empty}\cleardoublepage\fi
  %%
  \renewcommand{\articlename}{ABSTRACTS}
  \title{#1}\shortauthor{Abstracts}\shorttitle{Abstracts}
  \listofauthors{}
  \if@book
  \addcontentsline{toc}{chapter}{\protect\@fiddleauthors{\@TOCtitle}{\rmaaTOCauthors}}%
  \addtocontents{htc}{\protect\HyperTOCEntry{\@rawtitle}{\@rawauthors}{\articlename}}
  \fi
  \pagestyle{rmaa} % this now seems necessary (06 Jun 2000)
  \thispagestyle{rmaafirstpage}
%%  \twocolumn[\@makeminititle]
  \renewcommand{\abstract}[1]{%
    \renewcommand\ADS@abstract{##1}%
    \renewcommand{\@abstract}{##1}\@writeabstract%
  }
  \renewcommand{\title}[2][]{%
    \gdef\ADS@authors{}%
    \gdef\ADS@affil{}%
    \gdef\ADS@references{}%
    \gdef\ADS@keywords{}%
    \renewcommand{\@rmaaArg}{##1}
    \ifx\@rmaaArg\@rmaaEmpty
       \renewcommand{\@TOCaftertitle}{}
    \else
       \renewcommand{\@TOCaftertitle}{ (\@rmaaArg)}
    \fi
    \renewcommand{\@title}{\rmaa@TOCTitleStyle{##2}}%
    \renewcommand{\@TOCtitle}{\rmaa@TOCTitleStyle{##2}\@TOCaftertitle}%
    \renewcommand{\rmaaIndexTitle}{\rmaa@IndexTitleStyle{##2\@TOCaftertitle.}}%
    \renewcommand\ADS@title{##2}%
    \xdef\ADS@firstpage{\@arabic{\c@page}}%
    }
  \renewcommand{\author}[1]{\let\@thanks\@empty\renewcommand\@author{##1} }
  \renewcommand{\footnotetext}[2][XX]{\par\noindent\textsuperscript{##1}~##2}
  \renewcommand{\footnoterule}{\kern-3\p@ \hrule \@width 3em \kern 2.6\p@}
  \if@book\@begofarticlehook\fi
  \if@book\begin{multicols}{2}[\@makeminititle]\else\@makeminititle\fi
  \begingroup
  \setcounter{footnote}{0}
  }% and at the end of the environment....
  {\endgroup\if@book\end{multicols}\listofabstracts\label{RMAALastPage}\clearpage\fi}
%% auxiliary commands used:
%% \@writeabstract does all the actual formatting 
\newcommand{\@writeabstract}{%
  \if@book%
  \addcontentsline{loa}{abstract}{\protect\@fiddleauthors{\@TOCtitle}{\rmaaTOCauthors}}%
  \fi%
  \begin{center}\@title\par\@author\end{center}
  \@abstract
  \par\noindent\footnoterule\nobreak\footnotesize\@thanks
  \par\noindent\hrulefill\goodbreak
  \vskip 1cm \@plus0.5cm
  %% write ADS entry
  \if@book\if@abstractsads\ADS@dumpdata\fi\fi%
  \endgroup % this pair is to make each abstract a self-contained group
  \begingroup
  \setcounter{footnote}{0}
  }

\if@book
%% List of abstracts (stuff written to .loa file)
\newcommand{\listofabstracts}{%
  \begin{multicols}{2}[%
    \vskip 1.5cm \vfill % leave at least 1.5 cm and go to bottom of page if possible
    \centerline{\MakeUpperCase{List of Abstracts}}
    \vskip 1.5cm \@plus0.5cm
    ][6cm] % this makes it start a new page if there is not at least 6cm of space
    \renewcommand{\@fiddleauthors}[2]{##1\hskip1em plus2em minus0.4em\penalty-100  ##2}
    \parindent\z@
%%    \parskip.5\p@ \@plus .7\p@\relax 
    \parskip3\p@ \@plus 1\p@ \@minus 1\p@\relax % allow quite a lot of
                                % leeway here to avoid breaks inside
                                % entries 
    \renewcommand{\\}{\relax} % temporarily turn off explicit line breaks just in case
    \renewcommand{\altaffilmark}[1]{\relax} % ditto with this
    \@starttoc{loa}
  \end{multicols}
  }
\newcommand*{\l@abstract}[2]{%
  \par\penalty-1000
  #1\hskip 1em\penalty0\null\nobreak\hfill#2}
\fi

%    \end{macrocode}
 
% \subsection{Code for typesetting an entire volume of conference
% proceedings} \label{sec:whole-volume}
%    \begin{macrocode}

\if@book
%% environment for including the individual articles
\newcommand{\@fiddleauthors}[2]{\relax}
\renewcommand{\theequation}{\@arabic\c@equation}
\newenvironment{articles}%
{ \begingroup % make it a group so that author definitions do not
              % spill over
  \gdef\ADS@authors{}
  \pagestyle{rmaa}
  \renewenvironment{document}{%
    \@begofarticlehook
    \addtocounter{chapter}{1}%
    \addcontentsline{toc}{chapter}{%
      \protect\@fiddleauthors{\@TOCtitle}{\rmaaTOCauthors}}%
    \addtocontents{htc}{\protect\HyperTOCEntry{\@rawtitle}{\@rawauthors}{\articlename}}
    }{\relax}
  \renewcommand{\documentclass}[2][foo]{\relax} 
  \renewcommand{\usepackage}[2][foo]{%
    \def\@packagename{##2}\def\@rmaacitename{rmaacite}%
    \ifx\@packagename\@rmaacitename
    \@rmcitetrue\ClassInfo{rmaa}{This paper uses rmaacite}
    \makeatletter\input{rmaacite.sty}\makeatother % read it in here so that all changes
    \else                                         % are local to paper
    \if@debug\ClassWarning{rmaa}{Package \@packagename\space was requested by \articlename\space}\fi
    \fi} 
  %% Don't allow individual articles to mess with the pages, volume or year
  \renewcommand{\SetFirstPage}[1]{\relax}
  \renewcommand{\SetYear}[1]{\relax}
  \renewcommand{\SetVolume}[1]{\relax}
  \renewcommand{\SetConfTitle}[1]{\relax}
  %% Don't allow individual articles to change these
  \renewcommand{\nonstopmode}{\relax}
  \renewcommand{\scrollmode}{\relax}
  \renewcommand{\outputplates}{\relax}
  %% XXXX need to check if using rmaacite and act accordingly
  %% avoid multiple declarations of the same length
  \let\oldnewlength=\newlength
  \renewcommand{\newlength}[1]{%
    \edef\reserved@a{\expandafter\@gobble\string ##1}%
    \@ifundefined{\reserved@a}{%
      \oldnewlength{##1}}{%
      \if@debug{\ClassInfo{rmaa}{\reserved@a\space already defined: doing nothing}}\fi%
      \relax}}%
  \renewcommand\newcommand{\providecommand}
  }%
{\endgroup\onecolumn}
     

%    \end{macrocode}
% Command to be used in `articles' environment for including an
% article
%    \begin{macrocode}
\newcommand{\article}[2][]{%
  % optional first argument is extra text to put in brackets after title in TOC
  % e.g. ``Invited Talk'' or ``Poster''
  \renewcommand{\articlename}{#2}
  % reset loads of counters
  \setcounter{footnote}{0}
  \setcounter{figure}{0}
  \setcounter{table}{0}
  \setcounter{section}{0}
  \setcounter{equation}{0}
  \setcounter{rmaabib}{0}
  \setcounter{author}{0}
  \setcounter{AUTHOR}{1}
  \gdef\ADS@authors{}%
  \gdef\ADS@affil{}%
  \gdef\ADS@affiliations{}%
  \gdef\ADS@references{}%
  \@firstauthortrue
  \@firstaffiltrue
  \@firstaltaffiltrue
  \@firstkwdtrue
  \@rmcitefalse
  \if@suppressfulladdresses\else\@needaddressestrue\fi
  \renewcommand{\@rmaaArg}{#1}
  \ifx\@rmaaArg\@rmaaEmpty
     \renewcommand{\@TOCaftertitle}{}
  \else
     \renewcommand{\@TOCaftertitle}{ (#1)}
  \fi
%    \end{macrocode}
% In the main journal articles must start on an odd page. 
%    \begin{macrocode}
  \if@journal\clearpage\thispagestyle{empty}\cleardoublepage\fi
%% Now input the file, then run the end-of-article hook (reinstated 28
  %% Sep 2001)
  {\input{\@articledir/\articlename.tex}\@endofarticlehook}}
%    \end{macrocode}
% Command to separate articles into sessions. Write an entry in the
% TOC, makes sure we start on an odd page (added 06 Jun 2000), then do
% a title page if \cs{SessionTitlePages} was called. 
% \changes{3.12}{04 Feb 2003}{Added option to have title page for each
% session. Turn it on with \cs{SessionTitlePages}.}
%    \begin{macrocode}
\newcommand{\session}[1]{
  %% this does nothing except write an entry in the TOC
  \addcontentsline{toc}{part}{\rmaa@TOCPartStyle{#1}}
  \addtocontents{htc}{\protect\HyperTOCpart{#1}}
  %% and also makes sure we start on an odd page (added 06 Jun 2000)
  \clearpage\thispagestyle{empty}\cleardoublepage
  \stepcounter{part}
  \if@sessiontitlepages{%
    \onecolumn
    \thispagestyle{empty}
    \renewcommand{\articlename}{Section\thepart TitlePage}
    \label{RMAAFirstPage}
    \make@session@title{#1}
    \label{RMAALastPage}
    \clearpage\thispagestyle{empty}\cleardoublepage
  }\fi
  }
\newcommand{\make@session@title}[1]{%
  {%
    \setlength\parindent{0pt}
    \vspace*{\stretch{1}}
    \rule{\linewidth}{1mm}
    \begin{flushright}\Huge\bfseries\boldmath
      #1
    \end{flushright}
    \rule{\linewidth}{1mm}
    \vspace*{\stretch{2}}
  }
}
%    \end{macrocode}
% Command to turn on separate title pages for each session. 
%    \begin{macrocode}
\newcommand{\SessionTitlePages}{\@sessiontitlepagestrue}
\newcommand{\BoldSessionTitles}{%
  \renewcommand{\rmaa@TOCPartStyle}[1]{\textbf{\rmaa@TOCTitleStyle{##1}}}
  }
\newcommand{\ArticleDir}[1]{\renewcommand{\@articledir}{#1/}\graphicspath{{1/}}}
\newcommand{\@articledir}{.}
%\newcommand{\ps@rmaafront}{\ps@rmaa}% page style for front matter (same as ever, actually)
\newcommand{\ps@rmaafront}{% page style for front matter:  page # in each foot
  \renewcommand\@oddfoot{\hfil{\thepage}}%
  \renewcommand\@evenfoot{{\thepage}\hfil}%
  % shorttitle in both heads
  \renewcommand\@evenhead{\hfil\@shorttitle\hfil}%
  \renewcommand\@oddhead{\@evenhead}
  }
\newcommand{\ps@rmaafrontfirstpage}{% page style for first page of front matter
  % page # in each foot, on outside edge
  \renewcommand\@oddfoot{\hfil\thepage}%
  \renewcommand\@evenfoot{\thepage\hfil}%
  % nothing in the even heads
  \renewcommand\@evenhead{\relax}%
  \renewcommand\@oddhead{\relax}%
  }
%% slight modification to the standard \frontmatter and \mainmatter
%% commands, just to set the \pagestyle
\renewcommand\frontmatter{%
  \latexonecolumn
  \cleardoublepage
  \@frontmattertrue
  \@mainmatterfalse
  \pagenumbering{roman}
  \pagestyle{rmaafront}
  \addtocounter{page}{2} % leave space for the title page
  }
\renewcommand\mainmatter{%
  \pagestyle{empty}
  \if@rmtwocol\twocolumn\fi
  \cleardoublepage
  \@frontmatterfalse
  \@mainmattertrue
  \pagenumbering{arabic}
  \pagestyle{rmaa}
  }
\renewcommand\backmatter{%
  \clearpage\pagestyle{empty}\cleardoublepage
  \@mainmatterfalse
  %% put extra space in the TOC
  \addtocontents{toc}{\bigskip}
  }
%    \end{macrocode}
%
% \subsection{Front sections}\label{sec:front-sections}
% \begin{macro}{\notitle}
% \begin{macro}{\shortname}
% \begin{macro}{frontsection}
% The environment \texttt{frontsection} is for doing a section in the
% front matter (preface, list of participants, or whatever). It has 
% one optional argument:
% \begin{enumerate}
% \item[\#1] Short title (no embedded spaces) for
% use in making PS files. E.g., ``list\_of\_participants''.   
% \end{enumerate}
% and two mandatory arguments: 
% \begin{enumerate}
% \item[\#2] Title of section (e.g., `Preface')
% \item[\#3] Commands to be evaluated before processing the contents
% of the environment. This is mainly useful for putting in the things
% that would go in the preamble of a normal article  (e.g.,
% |\listofauthors|, |\indexauthor|, |\notitle|) 
% \end{enumerate}
%    \begin{macrocode}
\newcommand{\notitle}{\@titlevisiblefalse} \newcommand{\shortname}{}
\newenvironment{frontsection}[3][]{%
  \begingroup
  \renewcommand{\shortname}{#1}
  \ifx\shortname\@rmaaEmpty
     \renewcommand{\articlename}{#2}
  \else
     \renewcommand{\articlename}{\shortname}
  \fi
  \@begofarticlehook
  \@titlevisibletrue % default
  \thispagestyle{rmaafrontfirstpage}
  \listofauthors{} % blank out this in case it is not specified
  \title{{#2}}\shorttitle{#2}\shortauthor{#2}
  #3 % evaluate the command given as second mandatory arg
  \addcontentsline{toc}{chapter}{\protect\@fiddleauthors{\@TOCtitle}{\rmaaTOCauthors}}%
  \addtocontents{htc}{\protect\HyperTOCEntry{\@rawtitle}{\@rawauthors}{\articlename}}
  \if@titlevisible
     \vspace*{1cm}\centerline{\rmaa@FrontSecStyle{\@title}}\vspace*{1cm}
  \fi
  \setlength{\parskip}{3ex}
  }{\frontlabel{RMAALastPage}\clearpage\endgroup}
%    \end{macrocode}
% \end{macro}
% \end{macro}
% \end{macro}
%    \begin{macrocode}
%% Definition of the command \frontlabel, which rights out something like `F4' to
%% the .aux file, instead of `\textlatin  {i}'. This is so it is easy to post-process
%% the file for making PS files of individual sections
\newcommand{\frontlabel}[1]{%
  \@bsphack
  \protected@write\@auxout{}%
     {\string\newlabel{\articlename:#1}{{}{F\@arabic{\c@page}}}}
  \@esphack
  }
%% User command to say X, Y, Z,\\ Editors of the Proceedings
\newcommand{\EditorSignature}{\vspace*{\parskip}\rightline{\@editors}%
  \rightline{Editors of the proceedings}}
%% Fiddle with \thesection, etc. so that they don't include the chapter number
%% Basically go back to the `article' definitions of these
\renewcommand{\thesection}{\@arabic\c@section}
\renewcommand{\thefigure}{\@arabic\c@figure}
\renewcommand{\thetable}{\@arabic\c@table}

\newcommand{\@oldargs}{}
\newcommand{\@newargs}{}
%% add TOC entry for each paper
%\renewcommand{\title}[1]%
%{\renewcommand{\@title}{\rmaa@TitleStyle{#1}}%
%\addtocounter{chapter}{1}\addcontentsline{toc}{chapter}{#1}%
%}

%% Formatting of the table of contents
%% dimensions of various parts
\newlength{\TOCtextwidth}\setlength{\TOCtextwidth}{12.65cm}
\newlength{\TOCpagenumwidth}\setlength{\TOCpagenumwidth}{1.3cm}
\newlength{\TOCentrywidth}\setlength{\TOCentrywidth}{\TOCtextwidth}
\addtolength{\TOCentrywidth}{-\TOCpagenumwidth} 
\newlength{\TOCextramargin}\setlength{\TOCextramargin}{1.25cm}
\newcommand{\TOCbox}[1]{\hspace{\TOCextramargin}\parbox{\TOCtextwidth}{#1}\par}
\newlength{\TOCextramarginx}\setlength{\TOCextramarginx}{\TOCextramargin}\addtolength{\TOCextramarginx}{4.3cm}
\newlength{\TOCtextwidthx}\setlength{\TOCtextwidthx}{\TOCtextwidth}\addtolength{\TOCtextwidthx}{-4.3cm}
\newcommand{\TOCboxx}[1]{\hspace{\TOCextramarginx}\parbox{\TOCtextwidthx}{#1}\par}

\newcommand{\RMAAbannerimage}{rmsc}
\renewcommand\tableofcontents{%
  \renewcommand{\articlename}{CONTENTS}
  \begingroup
  \@begofarticlehook
  \if@twocolumn
     \@restonecoltrue\onecolumn
  \else
     \@restonecolfalse
  \fi
  %% 9pt for journal style
  \if@proceedings\else\small\fi
  \if@newstyle\small\fi %% same for new style
  \if@proceedings\thispagestyle{rmaafrontfirstpage}\else\thispagestyle{empty}\fi
  \shortauthor{Contents}\shorttitle{Contents}
  \vspace*{18mm}
  \if@printbanner
  %% put in the banner header (need to get better image of this)
     \TOCbox{\hspace*{-0.15cm}\includegraphics[width=0.9\TOCtextwidth]{\RMAAbannerimage}}
     \vspace*{8mm}
  %% volume and publication date (should really be in funny font)
     \TOCboxx{\@volumen\\[2mm] \@mes\ \@anyo}
  \else
  %% put in the banner header (need to get better image of this)
     \TOCbox{\hspace*{-0.15cm}\includegraphics[width=0.9\TOCtextwidth]{\RMAAbannerimage}}
     \vspace*{12.5mm}
  %% volume and publication date (should really be in funny font)
     \TOCboxx{~\\[2mm] ~}
  \fi
  \vspace*{16mm}
  %% This is the command that fiddles with the TOC entry so that the
  %% title is flushleft but the authors are flushright
  \renewcommand{\@fiddleauthors}[2]{%
    % first put the paper title
    \renewcommand{\\}{\relax} % temporarily turn off explicit line breaks just in case
    \renewcommand{\altaffilmark}[1]{\relax} % ditto with this
    \renewcommand{\textsuperscript}[1]{\relax} % ditto with this 
    ##1
    % now save the author list in a box
    \sbox\@tempboxa{##2}%
    % check that box is shorter than the total width available for the entry
    \ifdim \wd\@tempboxa <\TOCentrywidth
       % if so, then put it right here, flush against right margin
       \penalty10000\hskip1em plus10em\penalty0\hbox{}\nobreak\hfill\box\@tempboxa
       % Note that we force a space of at least 1em between title and authorlist, 
       % stretchable to 11em since we may have very short titles but still with an 
       % authorlist too long to go on the same line.  The only linebreak allowed is 
       % directly before the authorlist.
    \else
       % if it is too long, we set the original authorlist text in its own parbox
       \par\vskip1ex\parbox[b]{\TOCentrywidth}{\raggedleft ##2}
    \fi
    }
  \@starttoc{toc}%
  \if@restonecol\twocolumn\fi
  \frontlabel{RMAALastPage}
  \clearpage
  \thispagestyle{empty}
  \cleardoublepage % make sure we start on an odd page
  \endgroup

  %%% ALSO, do the HTML TOC
  %% The file that will be used to make the HTML TOC
  \@starttoc{htc}
  %% All the stuff in this file does nothing in the book style
  \addtocontents{htc}{\protect\HyperTOCmeeting{\@conftitle}{\@reunionplace}{\@reuniondate}}
  \addtocontents{htc}{\protect\HyperTOCeditors{\@editors}}
  \addtocontents{htc}{\protect\HyperTOCpreamble}
  \addtocontents{htc}{\protect\HyperTOCEntry{Table of Contents}{}{\articlename}}

  %% and the ADS file
  \ADS@openfile
  }
%% This is now much simpler
\renewcommand\l@chapter[2]{%TOC entry for articles
  \ifnum \c@tocdepth >\m@ne
    \vskip 1.0em \@plus\p@
    \TOCbox{%
      \parbox[b]{\TOCentrywidth}{#1}\parbox[b]{\TOCpagenumwidth}{\hfill#2}
    }
  \fi} 
\renewcommand\l@part[2]{%TOC entry for sessions: centered, no page number
  \ifnum \c@tocdepth >\m@ne
    \vskip 1.5em \@plus1.5em \goodbreak
    \TOCbox{%
      \parbox[b]{\TOCentrywidth}{\hfil#1\hfil}\hfill
    }
  \fi\nobreak} 

%% Formatting theindex
%% strongly encourage page breaks between index entries
\renewcommand{\@idxitem}{\par\penalty-1000\hangindent.45cm} 
\renewcommand{\indexspace}{\par\penalty-1000\vskip 10\p@ \@plus15\p@ \@minus3\p@\relax}
\renewenvironment{theindex}{%
  %% first, make sure we start on an odd page (06 Jun 2000)
  \clearpage\thispagestyle{empty}\cleardoublepage
  %% this looks a bit like an article but is two columns and without abstract, etc
  \renewcommand{\articlename}{INDEX}
  \title{Author Index}\shortauthor{Author Index}\shorttitle{Author Index}
  \listofauthors{}
  \addcontentsline{toc}{chapter}{\protect\@fiddleauthors{\@TOCtitle}{\rmaaTOCauthors}}%
  \addtocontents{htc}{\protect\HyperTOCEntry{\@rawtitle}{\@rawauthors}{\articlename}}
  \pagestyle{rmaa}
  \thispagestyle{rmaafirstpage}
  \@begofarticlehook
  %% 9pt for journal style
  \if@proceedings\else\small\fi
  \if@newstyle\small\fi
  \twocolumn[\@makeminititle]
  %% vertical spacing (see also definition of \@idxitem and \indexspace
  \parskip\z@ \@plus .6\p@\relax % allow quite a lot of leaway here to avoid breaks inside entries
  %% horizontal spacing, breaking and glue parameters (see also the stuff in the definition 
  %% of \indexauthor
  \parindent\z@ 
  \exhyphenpenalty10000 % don't break on explicit hyphens, it always looks bad
  \rightskip=0pt plus.2em % allow a tiny bit of leeway at the ends of lines
  \spaceskip.3333em plus.5em minus.15em % a lot of leeway between words
  \xspaceskip.5em plus.5em minus.4em % this seems to have no effect
  \let\item\@idxitem
  \renewcommand{\\}{\relax} % temporarily turn off explicit line breaks just in case
  \renewcommand{\altaffilmark}[1]{\relax} % ditto with this
  \renewcommand{\textsuperscript}[1]{\relax} % ditto with this
  }{\label{RMAALastPage}\clearpage} % clearpage necessary so that we know about \shortauthor



%% end definitions for book styles
\fi

%    \end{macrocode}
% In order to do away with plates, we sometimes want a figure to
% appear in color only in the online version. The following command
% can bve used in a figure caption to typeset a note to this effect. 
%    \begin{macrocode}
\newcommand\ColorOnlineOnly{%
  \textsc{Note: this figure is available in color in the electronic
    version of this article, obtainable from}
  \texttt{http://\discretionary{}{}{}www.\discretionary{}{}{}astroscu.\discretionary{}{}{}unam.\discretionary{}{}{}mx/\discretionary{}{}{}\~{}rmaa/}.}
%% Plates
%% Define 2 counters: plate is the number in the whole document, plateloop is used 
%% internally by \outputplates
\newcounter{plate}
\newcounter{plateloop}
%% Currently limited to 15 plates in document
%% when we do a label, we need to write to \newlabel entries to the .aux file, one 
%% that says RMAAPlateX so that it can be referred to in \ouputplates, and one that
%% says articlename:RMAAPlateX so that the dvips program can find it and asociate it
%% with the right article
\newcommand{\@platelabeli}{\oldlabel{RMAAPlate1}\label{RMAAPlate1}}
\newcommand{\@platelabelii}{\oldlabel{RMAAPlate2}\label{RMAAPlate2}}
\newcommand{\@platelabeliii}{\oldlabel{RMAAPlate3}\label{RMAAPlate3}}
\newcommand{\@platelabeliv}{\oldlabel{RMAAPlate4}\label{RMAAPlate4}}
\newcommand{\@platelabelv}{\oldlabel{RMAAPlate5}\label{RMAAPlate5}}
\newcommand{\@platelabelvi}{\oldlabel{RMAAPlate6}\label{RMAAPlate6}}
\newcommand{\@platelabelvii}{\oldlabel{RMAAPlate7}\label{RMAAPlate7}}
\newcommand{\@platelabelviii}{\oldlabel{RMAAPlate8}\label{RMAAPlate8}}
\newcommand{\@platelabelix}{\oldlabel{RMAAPlate9}\label{RMAAPlate9}}
\newcommand{\@platelabelx}{\oldlabel{RMAAPlate10}\label{RMAAPlate10}}
\newcommand{\@platelabelxi}{\oldlabel{RMAAPlate11}\label{RMAAPlate11}}
\newcommand{\@platelabelxii}{\oldlabel{RMAAPlate12}\label{RMAAPlate12}}
\newcommand{\@platelabelxiii}{\oldlabel{RMAAPlate13}\label{RMAAPlate13}}
\newcommand{\@platelabelxiv}{\oldlabel{RMAAPlate14}\label{RMAAPlate14}}
\newcommand{\@platelabelxv}{\oldlabel{RMAAPlate15}\label{RMAAPlate15}}
\newcommand{\@outputplatelabeli}{\oldlabel{RMAAOutputPlate1}}
\newcommand{\@outputplatelabelii}{\oldlabel{RMAAOutputPlate2}}
\newcommand{\@outputplatelabeliii}{\oldlabel{RMAAOutputPlate3}}
\newcommand{\@outputplatelabeliv}{\oldlabel{RMAAOutputPlate4}}
\newcommand{\@outputplatelabelv}{\oldlabel{RMAAOutputPlate5}}
\newcommand{\@outputplatelabelvi}{\oldlabel{RMAAOutputPlate6}}
\newcommand{\@outputplatelabelvii}{\oldlabel{RMAAOutputPlate7}}
\newcommand{\@outputplatelabelviii}{\oldlabel{RMAAOutputPlate8}}
\newcommand{\@outputplatelabelix}{\oldlabel{RMAAOutputPlate9}}
\newcommand{\@outputplatelabelx}{\oldlabel{RMAAOutputPlate10}}
\newcommand{\@outputplatelabelxi}{\oldlabel{RMAAOutputPlate11}}
\newcommand{\@outputplatelabelxii}{\oldlabel{RMAAOutputPlate12}}
\newcommand{\@outputplatelabelxiii}{\oldlabel{RMAAOutputPlate13}}
\newcommand{\@outputplatelabelxiv}{\oldlabel{RMAAOutputPlate14}}
\newcommand{\@outputplatelabelxv}{\oldlabel{RMAAOutputPlate15}}
\newcommand{\@platerefi}{\oldpageref{RMAAPlate1}}
\newcommand{\@platerefii}{\oldpageref{RMAAPlate2}}
\newcommand{\@platerefiii}{\oldpageref{RMAAPlate3}}
\newcommand{\@platerefiv}{\oldpageref{RMAAPlate4}}
\newcommand{\@platerefv}{\oldpageref{RMAAPlate5}}
\newcommand{\@platerefvi}{\oldpageref{RMAAPlate6}}
\newcommand{\@platerefvii}{\oldpageref{RMAAPlate7}}
\newcommand{\@platerefviii}{\oldpageref{RMAAPlate8}}
\newcommand{\@platerefix}{\oldpageref{RMAAPlate9}}
\newcommand{\@platerefx}{\oldpageref{RMAAPlate10}}
\newcommand{\@platerefxi}{\oldpageref{RMAAPlate11}}
\newcommand{\@platerefxii}{\oldpageref{RMAAPlate12}}
\newcommand{\@platerefxiii}{\oldpageref{RMAAPlate13}}
\newcommand{\@platerefxiv}{\oldpageref{RMAAPlate14}}
\newcommand{\@platerefxv}{\oldpageref{RMAAPlate15}}
\newcounter{figplatei}
\newcounter{figplateii}
\newcounter{figplateiii}
\newcounter{figplateiv}
\newcounter{figplatev}
\newcounter{figplatevi}
\newcounter{figplatevii}
\newcounter{figplateviii}
\newcounter{figplateix}
\newcounter{figplatex}
\newcounter{figplatexi}
\newcounter{figplatexii}
\newcounter{figplatexiii}
\newcounter{figplatexiv}
\newcounter{figplatexv}
%% Printed output should use the counter `plate'
\renewcommand{\theplate}{\@arabic\c@plate}
\renewcommand{\theplateloop}{\@arabic\c@plateloop}
%% Command to output all the saved plates in this paper
\newcommand{\outputplates}{%
  \setcounter{plateloop}{0}
  \if@book\else\@endofarticlehook\fi % not entirely sure why, but this
                                     % seems to be needed
  \clearpage % first output the last page of text
  % now set pagestyle for blank even pages and make sure we start on an odd page
  \pagestyle{empty}
  \cleardoublepage
  \onecolumn
  \parindent\z@
  % first put in an extra sheet
  % the main purpose of this is so there is always a blank even page before 
  % each plate for filling in the PS file if the article has an odd number
  % of pages (see splitbook.sh)
  \null\vfil\centerline{PLATES}\vfil\clearpage\pagestyle{rmaaplate}\cleardoublepage
  % output each plate in turn
  \@whilenum\value{plateloop}<\value{plate}\do{%
    \addtocounter{plateloop}{1}
    % restore the figure counter from \@figcounti, \@figcountii, etc
    \setcounter{figure}{\value{figplate\roman{plateloop}}}
    % spew out the plate (in its own group for safety)
    {\csname @plate\roman{plateloop}\endcsname\par}
    }
  \pagestyle{rmaa}
  % reset the page number as if nothing had happened
  \addtocounter{page}{-\value{plate}}
  \addtocounter{page}{-\value{plate}}
  }
\newcommand\figcaption{%
  \def\@captype{figure}
  \renewcommand{\@makecaption}{\@makefigcap}
  \caption}
\newcommand\tabcaption{%
  \def\@captype{table}
  \renewcommand{\@makecaption}{\@maketabcap}
  \caption}

\newcommand{\plate}[2][1]{%
  %% User command to make a plate page: 
  %% Mandatory argument is what goes on the plate - this should not 
  %% include any floating environments (e.g. figure). Instead, \figcaption 
  %% should be used to do the captions. 
  %% Optional argument is the number of figures (i.e. captions) on the plate
  %% page (default is 1). This is necessary for incrementing the figure counter. 
  %%
  %% Note that we have to use a \gdef rather than an \xdef when we 
  %% save the plate output in a command. This means that in the book styles when it 
  %% comes to spewing out the plate, we will have forgotten about any user macros 
  %% that may have been defined. This can be got round by repeating the macros inside the plate.
  %% (Presumably, one would use a \newcommand rather than \renewcommand if they were new.)
  % increment global and local counters
  \addtocounter{plate}{1}
  % write appropriate label to .aux file
  \csname @platelabel\roman{plate}\endcsname
  % define a global macro called \@platei, \@plateii, \@plateiii, \@plateiv, ...
  \expandafter\gdef\csname @plate\roman{plate}\endcsname{%
    \csname @outputplatelabel\roman{plateloop}\endcsname
    \null\vfil #2 \vfil
    \cleardoublepage
    }
  % save the current shortauthor and shorttitle
  % need to be \xdef to take current values of \@shortauthor, etc 
  \expandafter\xdef\csname @shorta\roman{plate}\endcsname{\@shortauthor}
  \expandafter\xdef\csname @shortt\roman{plate}\endcsname{\@shorttitle}
  % save the current value of the figure counter for later use in \outputplates
  \setcounter{figplate\roman{plate}}{\value{figure}}
  % increment the figure counter according to the optional argument
  \addtocounter{figure}{#1}
  }
%% page style for plate pages 
\newcommand{\ps@rmaaplate}{%
  % Plate X (see page Y) in the foot
  \renewcommand\@oddfoot{\hfil\csname @shorta\roman{plateloop}\endcsname\ 
    (see page~\csname @plateref\roman{plateloop}\endcsname)\hfil}%
  \renewcommand\@evenfoot{\relax}%
  % nothing in even heads (blank pages)
  \renewcommand\@evenhead{\relax}
  % short title and page num in the odd heads
  \renewcommand\@oddhead{\hfil\csname @shortt\roman{plateloop}\endcsname%
    \hfil\llap{PLATE~\theplateloop}}
  }





%% Things for the User Guide
\if@guide
%% stolen from doc.sty
\def\MakeShortVerb#1{%
  \expandafter\ifx\csname cc\string#1\endcsname\relax
  \@shortvrbinfo{Made }{#1}%
  \add@special{#1}%
  \expandafter
  \xdef\csname cc\string#1\endcsname{\the\catcode`#1}%
  \begingroup
  \catcode`\~\active  \lccode`\~`#1%
  \lowercase{%
    \global\expandafter\let
    \csname ac\string#1\endcsname~%
    \gdef~{\verb~}}%
  \endgroup
  \global\catcode`#1\active
  \else
  \@shortvrbinfo\@empty{#1 already}%
  \fi}
\def\DeleteShortVerb#1{%
  \expandafter\ifx\csname cc\string#1\endcsname\relax
  \else
  \@shortvrbinfo{Deleted }{#1 as}%
  \rem@special{#1}%
  \global\catcode`#1\csname cc\string#1\endcsname
  \global \expandafter\let \csname cc\string#1\endcsname \relax
  \ifnum\catcode`#1=\active
  \begingroup
  \catcode`\~\active   \lccode`\~`#1%
  \lowercase{%
    \global\expandafter\let\expandafter~%
    \csname ac\string#1\endcsname}%
  \endgroup \fi \fi}
\def\@shortvrbinfo#1#2{%
  \PackageInfo{doc}{%
    #1\expandafter\@gobble\string#2 a short reference
    for \string\verb}}
\def\add@special#1{%
  \rem@special{#1}%
  \expandafter\gdef\expandafter\dospecials\expandafter
  {\dospecials \do #1}%
  \expandafter\gdef\expandafter\@sanitize\expandafter
  {\@sanitize \@makeother #1}}
\def\rem@special#1{%
  \def\do##1{%
    \ifnum`#1=`##1 \else \noexpand\do\noexpand##1\fi}%
  \xdef\dospecials{\dospecials}%
  \begingroup
  \def\@makeother##1{%
    \ifnum`#1=`##1 \else \noexpand\@makeother\noexpand##1\fi}%
  \xdef\@sanitize{\@sanitize}%
  \endgroup}
%% define | as shortcut for \verbatim
\MakeShortVerb{\|}
%% comand to specify arguments to LaTeX commands
\newcommand{\rmArg}[1]{%
  \ensuremath{\langle}{\scriptsize \textsf{\MakeUpperCase{#1}}}\ensuremath{\rangle}}
%\newcommand{\rmArg}[1]{\ensuremath{\langle}\textsc{#1}\ensuremath{\rangle}}
\providecommand{\BibTeX}{{\rmfamily B\kern-.05em%
    \textsc{i\kern-.025em b}\kern-.08em%
    T\kern-.1667em\lower.7ex\hbox{E}\kern-.125emX}}
%% Environment for code examples
\newenvironment{Code}{%
  \begin{list}{}{\setlength{\leftmargin}{0.05\columnwidth}} 
  \item[]}{\end{list}}
\fi
%</package>
%    \end{macrocode}

% \section{User commands for astronomical articles}
% \label{sec:user-commands}
% Some random definitions. These go both into the main class file and
% also into |adshtml.sty|.
% 
% Definitions of the journal name in different incarnations.  
%    \begin{macrocode}
%<*package|adshtml-style>
\newcommand{\RMAA}{RMxAA}
\newcommand{\RMAAlong}{Revista Mexicana de Astronom\'{\i}a y Astrof\'{\i}sica}
\newcommand{\RMAASC}{RMxAC}
\newcommand{\RMAASClong}{\RMAA\ (Serie de Conferencias)}
%    \end{macrocode}
% \begin{macro}{\ibidrule}
% This puts a horizontal rule with a period after it. It is used when
% the author list of a reference is identical to the preceding one. 
%    \begin{macrocode}
\newcommand{\ibidrule}{\rule[-0.03cm]{4.7em}{.01cm}.\ }
%    \end{macrocode}
% \end{macro}
% \begin{macro}{\ThisVolume}
% References to other articles in the same volume can be coded as,
% e.g., |\bibitem{} Bloggs, J. \ThisVolume{bloggs}|.
%    \begin{macrocode}
\newcommand{\ThisVolume}[1]{%
  \@year, \if@proceedings\RMAASC\else\RMAA\fi, \@volume,
  \if@book\oldpageref{#1:RMAAFirstPage}\else 000\fi\ (this volume)}
%    \end{macrocode}
% \end{macro}
% Other stuff that may be useful in ref lists.
%    \begin{macrocode}
\newcommand{\etal}{et al.\@}
%    \end{macrocode}
% Various new commands from AAS macros
%    \begin{macrocode}
\def\acknowledgments{\vskip 3ex plus .8ex minus .4ex}
\let\acknowledgements=\acknowledgments                  % second spelling
\let\jnl@style=\rm
\def\ref@jnl#1{{\jnl@style#1}}
\DeclareRobustCommand\aj{\ref@jnl{AJ}}                   
\DeclareRobustCommand\araa{\ref@jnl{ARA\&A}}             
\DeclareRobustCommand\apj{\ref@jnl{ApJ}}                 
\DeclareRobustCommand\apjl{\ref@jnl{ApJ}}                
\DeclareRobustCommand\apjs{\ref@jnl{ApJS}}               
\DeclareRobustCommand\ao{\ref@jnl{Appl.~Opt.}}           
\DeclareRobustCommand\apss{\ref@jnl{Ap\&SS}}             
\DeclareRobustCommand\aap{\ref@jnl{A\&A}}                
\DeclareRobustCommand\aapr{\ref@jnl{A\&A~Rev.}}          
\DeclareRobustCommand\aaps{\ref@jnl{A\&AS}}              
\DeclareRobustCommand\azh{\ref@jnl{AZh}}                 
\DeclareRobustCommand\baas{\ref@jnl{BAAS}}               
\DeclareRobustCommand\jrasc{\ref@jnl{JRASC}}             
\DeclareRobustCommand\memras{\ref@jnl{MmRAS}}            
\DeclareRobustCommand\mnras{\ref@jnl{MNRAS}}             
\DeclareRobustCommand\pra{\ref@jnl{Phys.~Rev.~A}}        
\DeclareRobustCommand\prb{\ref@jnl{Phys.~Rev.~B}}        
\DeclareRobustCommand\prc{\ref@jnl{Phys.~Rev.~C}}        
\DeclareRobustCommand\prd{\ref@jnl{Phys.~Rev.~D}}        
\DeclareRobustCommand\pre{\ref@jnl{Phys.~Rev.~E}}        
\DeclareRobustCommand\prl{\ref@jnl{Phys.~Rev.~Lett.}}    
\DeclareRobustCommand\pasp{\ref@jnl{PASP}}               
\DeclareRobustCommand\pasj{\ref@jnl{PASJ}}               
\DeclareRobustCommand\qjras{\ref@jnl{QJRAS}}             
\DeclareRobustCommand\skytel{\ref@jnl{S\&T}}             
\DeclareRobustCommand\solphys{\ref@jnl{Sol.~Phys.}}      
\DeclareRobustCommand\sovast{\ref@jnl{Soviet~Ast.}}      
\DeclareRobustCommand\ssr{\ref@jnl{Space~Sci.~Rev.}}     
\DeclareRobustCommand\zap{\ref@jnl{ZAp}}                 
\DeclareRobustCommand\nat{\ref@jnl{Nature}}              
\DeclareRobustCommand\iaucirc{\ref@jnl{IAU~Circ.}}
\DeclareRobustCommand\aplett{\ref@jnl{Astrophys.~Lett.}}
\DeclareRobustCommand\apspr{\ref@jnl{Astrophys.~Space~Phys.~Res.}}
\DeclareRobustCommand\bain{\ref@jnl{Bull.~Astron.~Inst.~Netherlands}}
\DeclareRobustCommand\fcp{\ref@jnl{Fund.~Cosmic~Phys.}}
\DeclareRobustCommand\gca{\ref@jnl{Geochim.~Cosmochim.~Acta}}
\DeclareRobustCommand\grl{\ref@jnl{Geophys.~Res.~Lett.}}
\DeclareRobustCommand\jcp{\ref@jnl{J.~Chem.~Phys.}}      
\DeclareRobustCommand\jgr{\ref@jnl{J.~Geophys.~Res.}}    
\DeclareRobustCommand\jqsrt{\ref@jnl{J.~Quant.~Spec.~Radiat.~Transf.}}
\DeclareRobustCommand\memsai{\ref@jnl{Mem.~Soc.~Astron.~Italiana}}
\DeclareRobustCommand\nphysa{\ref@jnl{Nucl.~Phys.~A}}
\DeclareRobustCommand\physrep{\ref@jnl{Phys.~Rep.}}
\DeclareRobustCommand\physscr{\ref@jnl{Phys.~Scr}}
\DeclareRobustCommand\planss{\ref@jnl{Planet.~Space~Sci.}}
\DeclareRobustCommand\procspie{\ref@jnl{Proc.~SPIE}}
\let\astap=\aap
\let\apjlett=\apjl
\let\apjsupp=\apjs
\let\applopt=\ao
\def\phn{\phantom{0}}
\def\phd{\phantom{.}}
\def\phs{\phantom{\ensuremath{-}}}
\def\phm#1{\phantom{#1}}
%    \end{macrocode}
% All these are redone as robust and using more modern commands where
% possible. The robustness is so they will not get expanded if used in
% abstracts and titles when these are passed to the ADS file. 
% \changes{3.16}{2003/04/07}{
%  Added \cs{Color} command for doing photometric colors.}
%    \begin{macrocode}
%%\def\sun{\hbox{$\odot$}}
\DeclareRobustCommand\sun{\ensuremath{\odot}}
%%\def\earth{\hbox{$\oplus$}}
\DeclareRobustCommand\earth{\ensuremath{\oplus}}
%%\def\lesssim{\mathrel{\hbox{\rlap{\hbox{\lower4pt\hbox{$\sim$}}}\hbox{$<$}}}}
%%\def\gtrsim{\mathrel{\hbox{\rlap{\hbox{\lower4pt\hbox{$\sim$}}}\hbox{$>$}}}}
\DeclareRobustCommand\lesssim{\mathrel{\hbox{\smaller\smaller\rlap{\hbox{\lower3pt\hbox{$\sim$}}}\raise1.5pt\hbox{$<$}}}}
\DeclareRobustCommand\gtrsim{\mathrel{\hbox{\smaller\smaller\rlap{\hbox{\lower3pt\hbox{$\sim$}}}\raise1.5pt\hbox{$>$}}}}
%%\def\lta{\mathrel{\spose{\lower 3pt\hbox{$\mathchar"218$}}
%%     \raise 2.0pt\hbox{$\mathchar"13C$}}}
%%\def\gta{\mathrel{\spose{\lower 3pt\hbox{$\mathchar"218$}}
%%     \raise 2.0pt\hbox{$\mathchar"13E$}}}
\DeclareRobustCommand\sq{\hbox{\rlap{$\sqcap$}$\sqcup$}}
%%\def\arcdeg{\hbox{$^\circ$}}
%%\def\arcmin{\hbox{$^\prime$}}
%%\def\arcsec{\hbox{$^{\prime\prime}$}}
\DeclareRobustCommand\arcdeg{\ensuremath{^\circ}}
\DeclareRobustCommand\arcmin{\ensuremath{^\prime}}
\DeclareRobustCommand\arcsec{\ensuremath{^{\prime\prime}}}
%%\def\fd{\hbox{$.\!\!^{\rm d}$}}
%%\def\fh{\hbox{$.\!\!^{\rm h}$}}
%%\def\fm{\hbox{$.\!\!^{\rm m}$}}
%%\def\fs{\hbox{$.\!\!^{\rm s}$}}
%%\def\fdg{\hbox{$.\!\!^\circ$}}
%%\def\farcm{\hbox{$.\mkern-4mu^\prime$}}
%%\def\farcs{\hbox{$.\!\!^{\prime\prime}$}}
%%\def\fp{\hbox{$.\!\!^{\scriptscriptstyle\rm p}$}}
\DeclareRobustCommand\fd{\ensuremath{.\!\!^\mathrm{d}}}
\DeclareRobustCommand\fh{\ensuremath{.\!\!^\mathrm{h}}}
\DeclareRobustCommand\fm{\ensuremath{.\!\!^\mathrm{m}}}
\DeclareRobustCommand\fs{\ensuremath{\mkern0.5mu.\mkern-5mu^\mathrm{s}\mkern-1mu}}
\DeclareRobustCommand\fdg{\ensuremath{.\!\!^\circ}}
\DeclareRobustCommand\farcm{\ensuremath{.\mkern-4mu^\prime}}
\DeclareRobustCommand\farcs{\ensuremath{\mkern1mu.\mkern-6mu^{\prime\prime}\mkern-1.5mu}}
%%\DeclareRobustCommand\farcs{\ensuremath{.\!\!^{\prime\prime}}}
\DeclareRobustCommand\fp{\ensuremath{.\!\!^{\scriptscriptstyle\rm p}}}
%%\def\micron{\hbox{$\mu$m}}
\DeclareRobustCommand\micron{\ensuremath{\mu\mathrm{m}}}
\let\la=\lesssim                        
\let\ga=\gtrsim
%%\def\case#1#2{\hbox{$\frac{#1}{#2}$}}
%%\def\slantfrac#1#2{\hbox{$\,^#1\!/_#2$}}
%%\def\onehalf{\slantfrac{1}{2}}
%%\def\onethird{\slantfrac{1}{3}}
%%\def\twothirds{\slantfrac{2}{3}}
%%\def\onequarter{\slantfrac{1}{4}}
%%\def\threequarters{\slantfrac{3}{4}}
\DeclareRobustCommand\case[2]{\ensuremath{\frac{#1}{#2}}}
\DeclareRobustCommand\slantfrac[2]{\ensuremath{\,^#1\!/_#2}}
\DeclareRobustCommand\onehalf{\slantfrac{1}{2}}
\DeclareRobustCommand\onethird{\slantfrac{1}{3}}
\DeclareRobustCommand\twothirds{\slantfrac{2}{3}}
\DeclareRobustCommand\onequarter{\slantfrac{1}{4}}
\DeclareRobustCommand\threequarters{\slantfrac{3}{4}}
%%\def\ubvr{\hbox{$U\!BV\!R$}}            
%%\def\ub{\hbox{$U\!-\!B$}}               
%%\def\bv{\hbox{$B\!-\!V$}}               
%%\def\vr{\hbox{$V\!-\!R$}}               
%%\def\ur{\hbox{$U\!-\!R$}}
%    \end{macrocode}
% New commands for ``photometric'' colors 
%    \begin{macrocode}
\DeclareRobustCommand\ubvr{\ensuremath{U\!BV\!R}}            
\DeclareRobustCommand\Color[2]{%
  \ensuremath{%
    #1\!\ifx#1V\!\fi-\!#2
  }}
\DeclareRobustCommand\ub{\Color{U}{B}}
\DeclareRobustCommand\bv{\Color{B}{V}}
\DeclareRobustCommand\vr{\Color{V}{R}}
\DeclareRobustCommand\ur{\Color{U}{R}}
%% New version of the \ion command, that actually works! 
%    \end{macrocode}
% \begin{macro}{\ion}
% This version of the |\ion| command allows you to write either
% |\ion{N}{4}| or |\ion{N}{iv}| in order to get N~IV. 
%    \begin{macrocode}
\DeclareRobustCommand\ion[2]{%
  \mbox{#1\kern0.2em%
    \smaller\rmfamily%
    \edef\@tempa{\@car#2\@nil}%
    \ifcat1\@tempa%
    \@Roman{#2}%
    \else%
    \uppercase{#2}%
    \fi}}
\newcounter{IonStage}
\DeclareRobustCommand\plainion[2]{\setcounter{IonStage}{#2}#1 \Roman{IonStage}}
%%\def\ion#1#2{\ensuremath{\mathrm{#1}\;}{\protect\small\rm\@Roman{#2}}}
%    \end{macrocode}
% \end{macro}
%    \begin{macrocode}
\DeclareRobustCommand\nodata{$\cdots$}
%    \end{macrocode}
% These are needed by the |\ion| command. Isn't there an easier way of
% doing this?
%    \begin{macrocode}
\let\@small=\small
\let\@normalsize=\normalsize
\let\@footnotesize=\footnotesize
\let\@scriptsize=\scriptsize
\let\@tiny=\tiny
\let\@large=\large
\let\@Large=\Large
\let\@LARGE=\LARGE
\let\@huge=\huge
\let\@Huge=\Huge
\newcommand{\smaller}{\small}
\renewcommand{\normalsize}{\renewcommand{\smaller}{\small}\@normalsize}
\renewcommand{\small}{\renewcommand{\smaller}{\footnotesize}\@small}
\renewcommand{\footnotesize}{\renewcommand{\smaller}{\scriptsize}\@footnotesize}
\renewcommand{\scriptsize}{\renewcommand{\smaller}{\tiny}\@scriptsize}
\renewcommand{\tiny}{\renewcommand{\smaller}{\tiny}\@tiny}
\renewcommand{\large}{\renewcommand{\smaller}{\normalsize}\@large}  
\renewcommand{\Large}{\renewcommand{\smaller}{\large}\@Large}  
\renewcommand{\LARGE}{\renewcommand{\smaller}{\Large}\@LARGE}  
\renewcommand{\huge}{\renewcommand{\smaller}{\LARGE}\@huge}  
\renewcommand{\Huge}{\renewcommand{\smaller}{\huge}\@Huge}  
%</package|adshtml-style>
%    \end{macrocode}


% \section{Other files}

% \subsection{The driver and make files}
% \label{sec:driver_make}
% These are put here for completeness, although they still need to be
% distributed separately since otherwise the user won't be able to
% unpack this file into its constituent parts.
% 
% The |Makefile| has two principal rules: |docs|, and |macros|. The
% |docs| rule produces a PDF and gzipped PS file of the documentation.
% I still haven't perfected the combination of rules required to get
% it to run \LaTeX{} enough times in all circumstances. The |macros|
% rule runs \LaTeX{} on |rmaa.ins| in order to unpack |rmaa.cls| and
% all the other little files (including the |Makefile| and the
% |rmaa.ins| files themselves\dots 
% 
% First set up lists of various types of file. 
%    \begin{macrocode}
%<*makefile>
DOCFILES=rmaa.pdf
MACRO_LTX_FILES=rmaa.cls rmaacite.sty adshtml.sty
SHELLSCRIPTS=makeads.sh
MACRO_OTHER_FILES=adshtml.perl sample.latex2html-init
MAKEFILE=Makefile
MACROFILES=$(MACRO_LTX_FILES) $(SHELLSCRIPTS) $(MACRO_OTHER_FILES)         
SKELETONFILES=ADS.tex 
%    \end{macrocode}
% Install the files read only into the package directory, specified via a path relative to the source directory. \changes{4.0}{2019/11/14}{Reorganize folder layout so that hardly anything is in the top level.  That way, things are clearer on github.}
%    \begin{macrocode}
INSTALL=install 
INSTALL_DIR=../packages
%    \end{macrocode}
% The default rule is to make everything. Installation must be
% explicitly asked for.
%    \begin{macrocode}
all: docs macros

install: install-docs install-macros

%    \end{macrocode}
% Most files are installed read-only to prevent editing. Skeleton
% files must be writeable. Must be careful not to overwrite the
% Makefile in the parent directory. 
%    \begin{macrocode}
install-macros: macros
        $(INSTALL) -m 444 $(MACRO_LTX_FILES) $(MACRO_OTHER_FILES) $(INSTALL_DIR)
        $(INSTALL) -m 555 $(SHELLSCRIPTS) $(INSTALL_DIR)
        $(INSTALL) -m 644 $(SKELETONFILES) $(INSTALL_DIR)

install-docs: docs
        $(INSTALL) rmaa.pdf $(INSTALL_DIR)/rm-fulldocs.pdf

%    \end{macrocode}
% Make the PDF and gzipped PS files. 
%    \begin{macrocode}
docs: $(DOCFILES)

%    \end{macrocode}
% Make everything else if either the |.dtx| or |.ins| files have
% changed. 
%    \begin{macrocode}
macros: rmaa.dtx rmaa.ins
%    \end{macrocode}
% Ensure that all the macro files can be written to then recreate
% them. The |touch| is there in case any of the files don't exist
% yet. 
%    \begin{macrocode}
        touch $(MACROFILES) $(MAKEFILE)
        chmod u+w $(MACROFILES) $(MAKEFILE) rmaa.ins
        latex "\input{rmaa.ins}"
%    \end{macrocode}
% \TeX{} automatically adds the |.tex| extension to extensionless
% files. We fix this and take the opportunity to convert
% back to TABs all that it has lovingly converted to spaces. 
%    \begin{macrocode}
        cp -f $(MAKEFILE) $(MAKEFILE).BCK && unexpand $(MAKEFILE).tex >| $(MAKEFILE) && rm $(MAKEFILE).tex
%    \end{macrocode}
% Make the shell scripts executable and everything unwriteable (so
% they don't get hand-edited by mistake). However, undo this for the
% files that the user is \emph{supposed} to modify. 
%    \begin{macrocode}
        chmod a+x $(SHELLSCRIPTS)
        chmod a-w $(MACROFILES) $(MAKEFILE)
        chmod u+w $(SKELETONFILES)

%    \end{macrocode}
% We now use PDF\LaTeX{} for this.
%    \begin{macrocode}
rmaa.idx: rmaa.dtx
        @echo "REGENERATING .idx FILE"
        pdflatex "\nonstopmode\input{rmaa.dtx}"

rmaa.aux: rmaa.dtx rmaa.ind rmaa.gls
        @echo "REGENERATING .aux FILE"
        pdflatex "\nonstopmode\input{rmaa.dtx}"

rmaa.ind: rmaa.idx
        @echo "REGENERATING INDEX"
        makeindex -s gind.ist rmaa

rmaa.gls: rmaa.glo
        @echo "REGENERATING CHANGELOG"
        makeindex -s gglo.ist -o rmaa.gls rmaa.glo

rmaa.pdf: rmaa.aux rmaa.ind rmaa.gls
        @echo "REGENERATING .pdf FILE"
        pdflatex "\nonstopmode\input{rmaa.dtx}"

clean: 
        rm -f rmaa.aux rmaa.dvi rmaa.ind rmaa.gls
%</makefile>
%    \end{macrocode}

% The |rmaa.ins| file uses the \textsf{docstrip} package to split out
% all the files and strip off this documentation.
%    \begin{macrocode}
%<*ins>
%%
%% Installation batch file for the RMAA LaTeX package
%%
\def\batchfile{rmaa.ins}
\input docstrip.tex
\preamble

PLEASE DO NOT EDIT THIS FILE! It is generated automatically from
the file rmaa.dtx using the accompanying Makefile and rmaa.ins 
Any hand-edited changes will be lost

A LaTeX document class for the RMxAA

Author: William Henney (w.henney@irya.unam.mx)

Latest version available from: http://www.irya.unam.mx/rmaa

Copyright (C) 1999-2019 Instituto de Astronomia, UNAM, Mexico

Permission is hereby granted, free of charge, to any person obtaining
a copy of this software and associated documentation files (the
"Software"), to deal in the Software without restriction, including
without limitation the rights to use, copy, modify, merge, publish,
distribute, sublicense, and/or sell copies of the Software, and to
permit persons to whom the Software is furnished to do so, subject to
the following conditions:

The above copyright notice and this permission notice shall be
included in all copies or substantial portions of the Software.

THE SOFTWARE IS PROVIDED "AS IS", WITHOUT WARRANTY OF ANY KIND,
EXPRESS OR IMPLIED, INCLUDING BUT NOT LIMITED TO THE WARRANTIES OF
MERCHANTABILITY, FITNESS FOR A PARTICULAR PURPOSE AND
NONINFRINGEMENT. IN NO EVENT SHALL THE AUTHORS OR COPYRIGHT HOLDERS BE
LIABLE FOR ANY CLAIM, DAMAGES OR OTHER LIABILITY, WHETHER IN AN ACTION
OF CONTRACT, TORT OR OTHERWISE, ARISING FROM, OUT OF OR IN CONNECTION
WITH THE SOFTWARE OR THE USE OR OTHER DEALINGS IN THE SOFTWARE.
\endpreamble
%    \end{macrocode}
% Clobber files without asking
%    \begin{macrocode}
\askforoverwritefalse
%    \end{macrocode}
% Make the \LaTeX{} macro files (class and styles)
%    \begin{macrocode}
\generate{% 
  \file{rmaa.cls}{\from{rmaa.dtx}{package}}
  \file{rmaa.ins}{\from{rmaa.dtx}{ins}}
  \file{rmaacite.sty}{\from{rmaa.dtx}{rmaacite}}
  \file{adshtml.sty}{\from{rmaa.dtx}{adshtml-style}}
  \file{ADS.tex}{\from{rmaa.dtx}{ads-container}}
  }
%    \end{macrocode}
% Make the other non-\LaTeX{} files. These are done separately since
% they can't use `\%\%' as a comment character for the
% preamble/postamble (and Perl files can't have the postamble I don't
% think). 
%    \begin{macrocode}
\generate{\nopreamble\nopostamble% 
  \file{Makefile}{\from{rmaa.dtx}{makefile}}
  \file{adshtml.perl}{\from{rmaa.dtx}{adshtml-perl}}
  \file{makeads.sh}{\from{rmaa.dtx}{ads-shell-script}}
  \file{sample.latex2html-init}{\from{rmaa.dtx}{latex2html-init}}
  }

%% Local Variables: 
%% mode: LaTeX
%% mode: Font-lock
%% End:
%</ins>
%    \end{macrocode}

% \subsection{Extended citation syntax}
% Modified from a file called mnrascite.sty
% I no longer remember where I stole the original version of this from
% 
% Alterations to the cite commands to make them do more sensible things.
% N.B. the bibliography style should produce entries which have labels where
% the date is delimited by angle brackets '<>' for these macros to work
% properly. They will not give very sensible error messages if you do not do 
% this properly.
% 
% A new command |\scite{key}| is defined which produces a citation which is 
% suitable for insertion in a sentence. i.e., it does not put brackets around 
% the whole citation, only the date. Also |\pcite{key}| produce output with no
% parentheses whatsoever. 
% 
% Should redefine |\@bibitem| really to have any hope of compatibility
% the command |\lbibcite| is defined to deal with the citations.
%    \begin{macrocode}
%<*rmaacite>
\def\bibitem[#1<#2>]#3{\par\hangindent=1.5em\hangafter=1
  \if@filesw
  {\def\protect##1{\string ##1\space}
    \immediate\write\@auxout{\string\lbibcite{#3}{#1}{#2}}}
  \fi}
\@redefinebibitem % what it says
%    \end{macrocode}
% Note that the |\lbibcite| defines two types of citation one contains the names,
% the other the year
%    \begin{macrocode}
\def\lbibcite#1#2#3{\global\@namedef{b@#1}{#2}
                    \global\@namedef{by@#1}{#3}}

\def\@bcite#1#2{({#1\if@tempswa , #2\fi})}
\def\@pcite#1#2{#1\if@tempswa , #2\fi}
%    \end{macrocode}
% This declaration contains what is put between the name and the date
% \changes{3.11}{2003/01/31}{Moved the \cs{articlename} to after the key for consistency
% with the new natbib stuff.}
%    \begin{macrocode}
\def\@brcite{\ (}
\def\@ebrcite{)}
\def\@fmttop#1{\@fmtcite[\csname #1\endcsname]}
\def\@citex[#1]#2{%
%%  \if@filesw\immediate\write\@auxout{\string\citation{#2}}\fi
%% WJH added 16 Jan 2002
  \@for\@citeb:=#2\do{%
    \edef\@citeb{\expandafter\@firstofone\@citeb:\articlename}%
    \if@filesw\immediate\write\@auxout{\string\citation{\@citeb}}\fi}%
%% to here
  \def\@citea{}%
  \@cite{\@for\@citeb:=#2\do
    {%
      \edef\@citeb{\expandafter\@firstofone\@citeb:\articlename}%
      \@citea\def\@citea{;\ }\@ifundefined
      {b@\@citeb}{{\bf ?}\@warning {Citation `\@citeb' on page \thepage \space undefined}}%
      {\csname b@\@citeb\endcsname\@citebro\csname by@\@citeb\endcsname\@citebrc}}}%
  {#1}}

\def\cite{\@ifnextchar [{\let\@citebro=\space
                          \let\@citebrc=\relax
                          \let\@cite=\@bcite\@tempswatrue\@citex}
                        {\let\@citebro=\space
                          \let\@citebrc=\relax
                          \let\@cite=\@bcite\@tempswafalse\@citex[]}}
\def\pcite{\@ifnextchar [{\let\@citebro=\space
                          \let\@citebrc=\relax
                          \let\@cite=\@pcite\@tempswatrue\@citex}
                        {\let\@citebro=\space
                          \let\@citebrc=\relax
                          \let\@cite=\@pcite\@tempswafalse\@citex[]}}
\def\scite{\@ifnextchar [{\let\@citebro=\@brcite
                          \let\@citebrc=\@ebrcite
                          \let\@cite=\@pcite\@tempswatrue\@citex}
                        {\let\@citebro=\@brcite
                          \let\@citebrc=\@ebrcite
                          \let\@cite=\@pcite\@tempswafalse\@citex[]}}

\renewenvironment{thebibliography}
  {\section*{REFERENCES}\vspace*{\afterbiblioskip}\small
    \if@book\RereadBiblio\fi
    \sloppy\clubpenalty4000\widowpenalty4000
    \sfcode`\.=1000\relax
    \setlength\parindent{\z@}\setlength\parskip{\z@}}
  {\relax\normalsize\@outputaddressesifnecessary} 
%</rmaacite>
%    \end{macrocode}

% \subsection{Auxilliary files for generating ADS data}

% The ADS data are written by \LaTeX{} to a |.ads| file. This file
% then has to be included in the following file (|ADS.tex|),
% which is run through \LaTeX{}2HTML (see below). 

%    \begin{macrocode}
%<*ads-container>
%%
%% Instructions:
%% Edit the \input{} line below to reflect the name of your file. 
\documentclass{article}
\usepackage[latin1]{inputenc}
\usepackage{adshtml}
\newcommand{\ADSfield}[2]
{\%#1  #2\\}
%%\newcommand{\ADSvfield}[2]
%%{\%#1 \begin{verbatim}#2\end{verbatim}}
\newenvironment{ADSarticle}{}{\par\bigskip}
\renewcommand{\bibitem}[2][]{\par}
\setlength{\parindent}{0pt}
\renewcommand{\ThisVolume}[1]{%  
  2001, \RMAASC, 12, \oldpageref{#1:RMAAFirstPage}}  
 

\begin{document}
\makeatletter
\input{rmbooksamp.ads}
\end{document}
%</ads-container>
%    \end{macrocode}

% This is the |adshtml.sty| style file used by |ADS.tex|. It just
% defines a few 
% macros that might be used in abstracts, etc (really it should define
% them all). This now gets all its code from
% Sec~\ref{sec:user-commands}. 
%    \begin{macrocode}
%<*adshtml-style>
%% Nothing else needs to go here now. 
%</adshtml-style>
%    \end{macrocode}
% The init file for latex2html.
%    \begin{macrocode}
%<*latex2html-init>
$PDFTOCAIRO = "pdftocairo";
$PDFLATEX = "pdflatex";
$LATEX = "latex";
$MAX_SPLIT_DEPTH = 0;
$HTML_VERSION = "4.0,math";
$NO_SIMPLE_MATH = 1;
$ADDRESS = "<p><address><a href=\"mailto:rmaa\@astro.unam.mx\">rmaa\@astro.unam.mx</address></p>"; 
$INFO = "<p>This document was created using latex2html and the rmaa latex macros</p>"; 
$NO_NAVIGATION = 1;
1; 
%</latex2html-init>
%    \end{macrocode}
% This is the |adshtml.perl| file that is used by \LaTeX{}2HTML. Some
% bizarre \TeX{}y command sometimes show up so we tell it to ignore
% these. All the astronomy macros that come from AAS are processed in
% \TeX{}. That way they get made to images but with an |ALT| tag that
% has the original code, which is what will be picked up by |lynx|,
% with the result that they get passed intact into the final |.txt|
% file. 
%    \begin{macrocode}
%<*adshtml-perl>
package main;

$dum = &do_cmd_makeatletter();

sub do_cmd_RMAA {
    local($_) = @_;
    "RMxAA" . $_;
}

sub do_cmd_RMAASC {
    local($_) = @_;
    "RMxAASC" . $_;
}

sub do_cmd_etal {
    local($_) = @_;
    "et al\." . $_;
}

sub do_cmd_ibidrule {
    local($_) = @_;
    "<img src=\"foo\" alt=\"\ibidrule\">" . $_;
}


sub do_cmd_bibitem {
    local($_) = @_;
    local($key);
    local($opt, $dummy)=&get_next_optional_argument; 
    $key = &missing_braces unless
        ((s/$next_pair_pr_rx/$key = $2; ''/eo)
        ||(s/$next_pair_rx/$key = $2; ''/eo));
    "<p>" . $_; 
}

&ignore_commands (<<_RAW_ARG_CMDS_);
lowercase
adjustfinalcols
textsuperscript # {}
newline # []
thinmuskip
mskip
smallskip
medskip
bigskip
hspace # {}
hspacestar # {}
vspace # {}
vspacestar # {}
mathrm
_at_mathrm
spacefactor
discretionary
unhbox 
voidb_at_x 
penalty 
_at_M
_at_m
null
IeC
_RAW_ARG_CMDS_

&process_commands_in_tex (<<_RAW_ARG_CMDS_);
sun
earth
lesssim
gtrsim
sq
argdeg
arcmin
arcsec
fd
fh
fm
fs
fdg
farcm
farcs
fp
micron
la
ga
case # {} # {}
slantfrac # {} # {}
onehalf
onethird
twothirds
onequarter
threequarters
ubvr
ub
bv
vr
ur
aj
araa
apj
apjl
apjs
ao
apss
aap
aapr
aaps
azh
baas
jrasc
memras
mnras
pra
prb
prc
prd
pre
prl
pasp
pasj
qjras
skytel
solphys
sovast
ssr
zap
nat
iaucirc
aplett
apspr
bain
fcp
gca
grl
jcp
jgr
jqsrt
memsai
nphysa
physrep
physscr
planss
procspie
_RAW_ARG_CMDS_

1;                              # This must be the last line
%</adshtml-perl> 
%    \end{macrocode}

% This is a shell script |makeads.sh| that generates the |_ads.txt| file from
% |_ads.tex| by using \LaTeX{}2HTML and |lynx|.  \changes{4.0}{2019/11/13}{Make sure that lynx writes utf-8 output instead of latin-1.}
% \changes{2.2beta13}{2002/03/12}{Complete rewrite of |makeads.sh|. Now much 
% improved I think}
% \changes{4.0}{2019/11/13}{Fix the at-sign substitution so that it only eliminates \cs{@} and does not mess with email addresses.}

%    \begin{macrocode}
%<*ads-shell-script>
#! /bin/sh
if [ -z $1 ]; then
    echo "Usage: $0 basename"
    exit
fi
BASENAME=$1
ADSBASENAME=${BASENAME}_ads
AUXFILE=$BASENAME.aux
TMPFILE=tmp.$$


echo "Generating $ADSBASENAME.tex"

# For some reason, must double-up \\ to \\\\ but not \ to \\ - odd!
cat > $ADSBASENAME.tex <<EOF
%% PLEASE DO NOT EDIT THIS FILE!  It is generated automatically from
%% the script makeads.sh.  Any hand-edited changes will be lost
%% 
\documentclass{article}
\usepackage[latin1]{inputenc}
\usepackage{html}
\usepackage{adshtml}
\newcommand{\ADSfield}[2]{\%#1  #2\\\\}
\newenvironment{ADSarticle}{}{\par\bigskip}
\setlength{\parindent}{0pt}
\newcounter{istage}
\renewcommand{\ion}[2]{\setcounter{istage}{#2}#1 \Roman{istage}}
\begin{document}
\makeatletter
EOF
echo >> $ADSBASENAME.tex "\input{$BASENAME.ads2}" 
echo >> $ADSBASENAME.tex "\end{document}"



echo "Generating $ADSBASENAME.sed" 

# Find type of journal
grep '\\documentclass\[.*proceedings.*\]' $BASENAME.tex >/dev/null \
     && JOURNAL='RMxAC' \
     || JOURNAL='RMxAA'

# Find the year and volume number
YEAR=`grep '\\SetYear' $BASENAME.tex | sed 's/.*\\SetYear{\(.*\)}.*/\1/'`
VOLUME=`grep '\\SetVolume' $BASENAME.tex | sed 's/.*\\SetVolume{\(.*\)}.*/\1/'`

# Generate an auxilliary sed script 
cat > $TMPFILE <<EOF
s|\\\\newlabel{\(.*\):RMAAFirstPage}{{.*}{\(.*\)}}|s/\\\\\\\\ThisVolume{\1}/$YEAR, RMxAC, $VOLUME, \2/g|g
EOF

# Now go through the .aux file, looking for \newlabels and use the sed 
# script just generated in order to make another sed script, one that 
# will replace \ThisVolume{LABEL} with the correct reference
grep '^\\newlabel{.*:RMAAFirstPage}' $AUXFILE \
     | sed -f $TMPFILE > $ADSBASENAME.sed && rm $TMPFILE

# Apend a further sed command to get rid of \@, since redefining 
# it in _ads.tex didn't work
echo >> $ADSBASENAME.sed 's/\\\\@//g'

echo "Generating $BASENAME.ads2 from $ADSBASENAME.sed and $BASENAME.ads" 

# Apply this to the .ads file, generating the .ads2 file, which is included
# by the _ads.tex file.
sed -f $ADSBASENAME.sed $BASENAME.ads > $BASENAME.ads2

echo "Generating $ADSBASENAME.txt from $ADSBASENAME.tex and $BASENAME.ads2"

# Now run latex2html on the _ads.tex file
latex2html -no_navigation -info 0 -address 0 -html_version 4.0,math $ADSBASENAME

# Then, convert to utf8 with lynx and tidy it up a bit
lynx --justify=off --width=100 --dump --nolist --display_charset=utf-8 \
     $ADSBASENAME/$ADSBASENAME.html | cut -c4- \
     | sed 's/ibidrule/\\ibidrule/g' > $ADSBASENAME.txt
%</ads-shell-script>
%    \end{macrocode}

% \section{Epilogue}
% Err, that's it \dots

\endinput


% \endinput
% Local Variables:
% mode: doctex
% TeX-master: t
% End:
"

rmaa.aux: rmaa.dtx rmaa.ind rmaa.gls
        @echo "REGENERATING .aux FILE"
        pdflatex "\nonstopmode% \def\fileversion{4.0} \def\filedate{2019/11/12}
% \MakeShortVerb{\|}
% \title{A \LaTeX{} document class for the Revista Mexicana de
% Astronom\'\i a y Astrof\'\i sica\thanks{This file has version
% \fileversion{} last revised \filedate{}.}}
% \author{William J. Henney\thanks{w.henney@irya.unam.mx}}
% \date{\filedate}
% \maketitle
% \paragraph{Abstract}
% This document describes the \texttt{rmaa} \LaTeX{}
% document class for the journal ``Revista Mexicana de
% Astronom\'\i a y Astrof\'\i sica''. The macros can be used for 
% \begin{enumerate}
% \item Preparing author manuscripts for the journal (including
% conference proceedings to be published in the ``Serie de
% Conferencias'')
% \item Typesetting an entire volume of the journal by inputting the
% individual article files from a container file. Facilities are
% included for automatic generation of TOC and Author Index. 
% \end{enumerate}
% 
% \tableofcontents
% \section{The documentation driver file}
% This is the part that generates the documentation that you may or
% may not be reading, which is done with the help of the Makefile
% discussed in section~\ref{sec:driver_make}. 
%    \begin{macrocode}
%<*driver>
\documentclass{ltxdoc}
\usepackage{hyperref}
\usepackage{palatino}
\EnableCrossrefs \CodelineIndex \RecordChanges
\setlength{\parindent}{0pt}
\def\MacroFont{\footnotesize\tt}
\begin{document} \DocInput{rmaa.dtx} \PrintIndex \PrintChanges \end{document}
%</driver>
%    \end{macrocode}

% \section{The implementation}
% 
% \changes{3.24}{2007/09/14}{Added user access to the macro version
% and date via \cs{rmVersion} and \cs{rmDate}}
% 
%    \begin{macrocode}
%<*package>
\newcommand{\rmaa@version}{3.26}
\newcommand{\rmaa@date}{2011/08/15}
\newcommand\rmVersion{\rmaa@version}
\newcommand\rmDate{\rmaa@date}
\NeedsTeXFormat{LaTeX2e}[1995/06/01]
%    \end{macrocode}
% \changes{4.0}{2019/11/02}{Change abstract to appear before
% resumen. Change RevMexAA to RMxAA for consistency with ADS.}
%
% \changes{3.22}{2007/03/22}{Allow proceedings articles to have
%   received/accepted dates} \changes{3.21}{2006/05/15}{Changes for
%   Silvia - abstracts are now shipped out to ADS.}
% \changes{3.19}{2006/01/19}{Changes for Silvia - remove adresses from
%   end, new abstract style} \changes{2.2beta}{2001/10/30}{ Fixed bug
%   whereby twocolumn option was not passed through to book
%   class. Finally sorted out the ``label on last page'' issue. }
% \changes{2.1}{2001/01/23}{ Fixed some inconsistencies in
%   \cs{nocite} and \cs{bibitem} uniqification Fixed the definition of
%   \cs{ion} so it will work in math mode and in contexts where the
%   typesize is different from normal (e.g., title, fig captions).
%   Modified the definition of \cs{@caption}, so we don't write any
%   list-of-figures data (we never use it anyway, and it just causes
%   problems with fragile commands in fig captions). }
% \changes{2.0}{2000/11/30}{ Added `newstyle' option, which should be
%   used in all new proceedings volumes (wider abstract, smaller
%   typeface in fig captions, etc.) Added commands for formatting
%   table captions. New demonstation document rmpostsamp.tex for
%   1-page posters. This will be the last release under a PRI
%   government (for a while...)}  \changes{1.9}{2000/06/06}{
%   Corrected position of page number in footer of articles that open
%   on even pages. Ensured that first article in each session opens on
%   an odd page, as do list of abstracts and author index. Added extra
%   space in TOC before the backmatter. Added command for editors,
%   \cs{BoldSessionTitles}, that makes produces bold titles for
%   session in the TOC.}  \changes{1.8}{2000/05/16}{ Fixed bug in
%   last page label with \cs{suppressfulladdresses} }
% \changes{1.7}{2000/05/11}{ Fixed more little things. Page labels
%   now work, even if you have no bibliography. You can output the
%   full addresses by hand with the \cs{outputfulladdresses}
%   command. Editors' Guide updated but still not finished. Revamped
%   splitbook.sh now produces PS file of each article for web page and
%   also an EPS file of each page for sending to the printers.}
% \changes{1.6}{2000/04/15}{ More cleaning up: full stops after
%   titles in index, running headers and footers in frontmatter,
%   fontsize in title of index, etc} \changes{1.5}{2000/04/15}{
%   Eliminated superscripts from titles in TOC and index}
% \changes{1.4}{2000/04/14}{ Added copyright and license info (GPL)}
% \changes{1.3}{2000/03/01}{ Corrected \cs{figcaption} to do the
%   caption formatting right Corrected font size of sec headings in
%   front matter} \changes{1.2}{2000/02/16}{ Added \cs{plate}
%   command} \changes{1.1}{2000/02/14}{ Added manuscript option,
%   corrected bugs} \changes{1.0}{2000/02/08}{ Now does a whole book
%   - lots of stuff tidied up} \changes{0.7}{1999/10/04}{ Increased
%   spacing between sections in single column styles. Suppressed
%   floats at bottom of pages (didn't work though) }
% \changes{0.6}{1999/09/24}{ Tweaked \cs{adjustfinalcols} (still not
%   perfect) Added proceedings option Added onecolumn/twocolumn
%   options (orthogonal to journal/proceedings) }
% \changes{0.5}{1999/09/23}{ Added \cs{fulladdresses} and
%   \cs{adjustfinalcols} (just about working)}
% \changes{0.4}{1999/09/22}{ Minor stylistic changes}
% \changes{0.3}{1999/06/28}{ Better mechanism for alternative
%   affiliations Added `guide' option for user guide}
% \changes{0.2}{1999/05/27}{ All basic functionality present, I
%   think} \changes{0.1}{1999/05/26}{ Just about working - still
%   missing Fig/Tab/Refs} \changes{0.0}{1999/05/21}{ So bored today,
%   decided to rewrite RMAA LaTeX macros from scratch}
%    \begin{macrocode}
\ProvidesClass{rmaa}[\rmaa@date\space v\rmaa@version\space RMxAA article
class]
%    \end{macrocode}
% \subsection{Option Handling}
% \label{sec:options}
% Declarations that are needed in options
% \changes{3.05}{2003/10/22}{Added natbib option}
% \changes{3.06}{2003/01/03}{Natbib option turned on permanently}
% \changes{3.07}{2003/01/11}{Added proof option}
% \changes{3.14}{2003/01/11}{Option to turn off natbib}
% \changes{3.18}{2003/10/22}{trying out doctex mode} 
% \changes{3.23}{2007/09/09}{Manuscript is now on by default.}
% \changes{3.26}{2011/08/15}{New option for review articles}
%    \begin{macrocode}
\newif\if@reviewarticle
\@reviewarticlefalse
\newif\if@proof
\@prooffalse
\newif\if@journal
\@journaltrue
\newif\if@printdates
\@printdatesfalse
\newif\if@proceedings
\@proceedingsfalse
\newif\if@manuscript
\@manuscripttrue
\newif\if@bigfont
\@bigfonttrue
\newif\if@guide
\@guidefalse
\newif\if@preprint
\@preprintfalse
\newif\if@rmtwocol
\@rmtwocoltrue
\newif\if@biblist
\@biblistfalse
\newif\if@book
\@bookfalse
\newif\if@natbib
\@natbibtrue
\newif\if@rmcite
\@rmcitefalse
\newif\if@titlevisible
\@titlevisibletrue
\newif\if@abstractpresent
\@abstractpresentfalse
\newif\if@printbanner
\@printbannerfalse
\newif\if@suppressfulladdresses
\@suppressfulladdressesfalse
\newif\if@newstyle
\@newstylefalse
\newif\if@debug
\@debugfalse
\newif\if@boldkeywords
\@boldkeywordsfalse
%    \end{macrocode}
% The option handling---mainly setting various booleans.
% \changes{2.2beta11}{2001/11/26}{Made twocol and
% newstyle the default for proceedings.}
% \changes{3.23}{2007/09/09}{Manuscript is now 10 pt instead of 12pt}
% \changes{3.23}{2007/09/10}{Full addresses are now automatically
% suppressed for proceedings style.}
% \changes{3.24}{2007/09/14}{New option \texttt{boldkeywords} that
% does the keywords in the old ugly way. Now off by default.}
%    \begin{macrocode}
\DeclareOption{manuscript}{\typeout{Manuscript for referee} 
  \@manuscripttrue\@journalfalse\@rmtwocolfalse
  \PassOptionsToClass{10pt}{article}}
\DeclareOption{bigfont}{\typeout{Big font manuscript requested by Gary
  Ferland} 
  \@manuscripttrue\@journalfalse\@rmtwocolfalse\@bigfonttrue
  \PassOptionsToClass{11pt}{article}}
\DeclareOption{journal}{\typeout{Main journal option}%
  \@journaltrue\@rmtwocoltrue\@printdatestrue\@manuscriptfalse
  \PassOptionsToClass{10pt}{article}}
\DeclareOption{guide}{\typeout{User guide option}%
  \@guidetrue\@journalfalse\@rmtwocoltrue\@manuscriptfalse
  \PassOptionsToClass{10pt}{article}}
\DeclareOption{preprint}{\typeout{Preprint option}%
  \@preprinttrue\@manuscriptfalse\PassOptionsToClass{10pt}{article}}
\DeclareOption{proof}{\typeout{Author proof option}%
  \@prooftrue\@manuscriptfalse}
\DeclareOption{proceedings}{\typeout{Proceedings option}%
  \@proceedingstrue\@journalfalse\@rmtwocoltrue\@newstyletrue
  \@suppressfulladdressestrue
  \@printdatesfalse\@manuscriptfalse\PassOptionsToClass{10pt}{article}} 
\DeclareOption{printdates}{\typeout{Force printing of received/accepted dates}%
  \@printdatestrue} 
\DeclareOption{book}{\typeout{Whole issue option}%
  \@booktrue} 
\DeclareOption{nonatbib}{\typeout{Turn off natbib support (required
    for compatibility with old rmaacite package).}%
  \@natbibfalse} 
\DeclareOption{newstyle}{\typeout{Use new layout style}%
  \@newstyletrue}
\DeclareOption{oldstyle}{\typeout{Use old layout style}%
  \@newstylefalse}
\DeclareOption{debug}{\typeout{Loads of annoying messages}%
  \@debugtrue}
%    \end{macrocode}
% These two need to be declared explicitly so that |\ExecuteOptions|
% actually does something.
%    \begin{macrocode}
\DeclareOption{twocolumn}{\@rmtwocoltrue}
\DeclareOption{onecolumn}{\@rmtwocolfalse}
\DeclareOption{twoside}{\PassOptionsToClass{\CurrentOption}{article}}
\DeclareOption{boldkeywords}{\@boldkeywordstrue}
%    \end{macrocode}
% Bold, upper-case keywords used to be the default. Now they must be
% turned on explicitly with the \texttt{boldkeywords} option. 
%    \begin{macrocode}
\DeclareOption{boldkeywords}{\@boldkeywordstrue}
%    \end{macrocode}
% Pass any other option along to the \texttt{article} class.
%    \begin{macrocode}
\DeclareOption*{\PassOptionsToClass{\CurrentOption}{article}}
%    \end{macrocode}
% Default options are to be two-sided and based on the standard
% \LaTeX{} \texttt{article} class,
% unless we are processing the entire book, in which case we base
% ourselves on the standard \texttt{book} class. Default is to use
% manuscript option unless otherwise specified. 
%    \begin{macrocode}
\ExecuteOptions{twoside,manuscript} \ProcessOptions\relax
\if@book
\if@rmtwocol\LoadClass[twocolumn]{book}\else\LoadClass{book}\fi
\newcommand{\abstract}{\relax}
\newcommand{\@maketitle}{\relax}
\else
\if@rmtwocol\LoadClass[twocolumn]{article}\else\LoadClass{article}\fi
\fi
%    \end{macrocode}
% Attempt to use natbib for citations 02 Jan 2003. 
% \changes{3.05}{2003/01/02}{If the natbib option is used, then set
% things up so that natbib's version of the bibliography is identical
% to the standard one.}
% \changes{3.09}{2003/01/23}{Explicitly require a recent version of
% natbib, so that we get a decent error message if the package is too
% old.} 
% \changes{3.23}{2007/09/09}{Now allow Bibtex to work properly} 
% \changes{3.25}{2009/10/01}{Change \cs{newcommand} tp \cs{def} when setting \cs{bibfont}. This allows it to work with all versions of natbib.} 
% \changes{3.25}{2010/02/06}{Updated copyright year.} 
% \changes{3.25}{2010/02/06}{Updated email and web addresses to \texttt{crya.unam.mx}.} 
%    \begin{macrocode}
\if@natbib
\RequirePackage{natbib}[1999/02/25]
\bibpunct{(}{)}{;}{a}{}{,}
\def\bibfont{\small}
\renewcommand\bibsection{\section*{REFERENCES}}
\setlength\bibsep{0pt}
\setlength\bibhang{1.5em}
\bibliographystyle{rmaa}
\fi
%    \end{macrocode}
% We want the \texttt{graphicx} package for the extended syntax of
% |\includegraphics|. We want the \texttt{ifthen} package for using in
% these macros (although we don't much yet). We explicitly load
% \texttt{babel} since some \LaTeX{} distributions messed up the
% hyphenation if we didn't.  
%    \begin{macrocode}
\RequirePackage{graphicx}
\RequirePackage{ifthen}
\IfFileExists{babel.sty}{%
%   \let\@restoreindent=\@afterindentfalse
  \RequirePackage[english]{babel}
%   \@restoreindent
  %% reinstate the definition of \@roman so as not to break .idx files
  \@ifundefined{org@@roman}{\relax}{\renewcommand{\@roman}[1]{\org@@roman{##1}}}
  }{%
  \typeout{Package babel.sty not found - don't worry about this}}
%    \end{macrocode}
% Try to load the \texttt{dcolumn} and \texttt{booktabs} packages if
% they can be found. These can produce nicer-looking tables. 
%    \begin{macrocode}
\IfFileExists{dcolumn.sty}{%
  \RequirePackage{dcolumn}
  }{%
  \typeout{Package dcolumn.sty not found - you will not be able to
    align columns on decimal points in tables}}
\IfFileExists{booktabs.sty}{%
  \RequirePackage{booktabs}
  }{%
  \typeout{Package booktabs.sty not found - your tables will not look
    as pretty as they ought too.}
  }
\IfFileExists{makeidx.sty}{%
  \RequirePackage{makeidx}
}{%
  \typeout{Package makeidx.sty not found - Author Index may not work.}
  }
%% commented out for moment since requires more modern version of LaTeX
%% \RequirePackage[latin1]{inputenc}
%    \end{macrocode}
% \changes{4.0}{2019/11/13}{Make \texttt{hyperref} and
% \texttt{etoolbox} into hard requirements, since we now use them
% extensively.}  We use hyperref to get internal and external links.
% But I don't like the colored boxes so I turn those off.  We use
% etoolbox for patching and robustifying at the moment.  In the
% future, I want to port a lot of other stuff over to it.  For
% instance, replace the ifthen package and the TeX ifs too.
%
% We target the latest version (1.7) of PDF, rather than the default
% of 1.5.  This reduces the chance of getting warnings from included
% graphics.  This is not strictly related to hyperref, but we put it
% here anyway.
%    \begin{macrocode}
\RequirePackage{etoolbox}
\RequirePackage{hyperref}
\hypersetup{hidelinks}
\pdfminorversion=7
%    \end{macrocode}
% We patch the \cs{@combinedblfloats} following a suggestion by Heiko
% Oberdiek at \url{https://tex.stackexchange.com/a/249743/3843}, which
% avoids problems when hyperlinks are split across
% pages. \changes{4.0}{2019/11/12}{Fix for issue where breaking
% hyperlink across pages gave an error.}
%    \begin{macrocode}
\urlstyle{same}
\newcommand\email{\begingroup \urlstyle{rm}\Url}
\patchcmd\@combinedblfloats{\box\@outputbox}{\unvbox\@outputbox}{}{%
  \errmessage{\noexpand\@combinedblfloats could not be patched}%
}%
%    \end{macrocode}
% \subsection{Other setup stuff}\label{sec:setup}
% Some declarations: 
%    \begin{macrocode}
\newif\if@needaddresses
\@needaddressestrue
\newif\if@frontmatter
\@frontmatterfalse
\newif\if@inpreamble
\@inpreambletrue
\newif\if@firstauthor
\@firstauthortrue
\newif\if@firstaffil
\@firstaffiltrue
\newif\if@firstaltaffil
\@firstaltaffiltrue
\newif\if@firstkwd
\@firstkwdtrue
\newif\if@noexpandADS
\@noexpandADSfalse
\newif\if@sessiontitlepages
\@sessiontitlepagesfalse
%    \end{macrocode}
% \begin{macro}{\@rmaaEmpty}
% \begin{macro}{\@rmaaArg}
% These are used in various commands to check if arg was left blank
%    \begin{macrocode}
\newcommand{\@rmaaEmpty}{}
\newcommand{\@rmaaArg}{}
%    \end{macrocode}  
% \end{macro}
% \end{macro}
%
% \subsection{Compatibility with old \LaTeX{} versions}
% \label{sec:compat}
% This package is based on the 1998/12/01 release of \LaTeX{}.
% Older releases may not define all the commands that we need.
% Luckily, we can use |\providecommand| to define these if necessary. 
% The following commands are not defined in the 1994/06/01 release. 
%
% \begin{macro}{\MakeUpperCase}
% \begin{macro}{\textsuperscript}
% \begin{macro}{\textbackslash}
% User commands:
%    \begin{macrocode}
\providecommand{\MakeUpperCase}[1]{\uppercase{#1}}
\providecommand{\textsuperscript}[1]{\ensuremath{^\textrm{#1}}}
\providecommand{\textbackslash}{\ensuremath{\backslash}}
%    \end{macrocode}  
% \end{macro}
% \end{macro}
% \end{macro}
% \begin{macro}{\protected@xdef}
% \begin{macro}{\restore@protect}
% \begin{macro}{\@unexpandable@protect}
% \begin{macro}{\hb@xt@}
% \LaTeX{} internal commands:
%    \begin{macrocode}
\providecommand{\protected@xdef}{%
        \let\@@protect\protect
        \let\protect\@unexpandable@protect
        \afterassignment\restore@protect
        \xdef}
\providecommand{\restore@protect}{\let\protect\@@protect}
\providecommand{\@unexpandable@protect}{\noexpand\protect\noexpand}
\providecommand{\hb@xt@}{\hbox to}
%    \end{macrocode}
% \end{macro}
% \end{macro}
% \end{macro}
% \end{macro}
%
% \subsection{Page layout}\label{sec:page-layout}
%
% The dimensions of the text were defined by Elizabeth. Text width is
% 169\,mm with equal margins on the two sides: 
%    \begin{macrocode}
\setlength{\textwidth}{169mm}
\setlength{\@tempdima}{\paperwidth}
\addtolength{\@tempdima}{-\textwidth}
\setlength{\oddsidemargin}{0.5\@tempdima}
\addtolength{\oddsidemargin}{-1in}
\setlength{\evensidemargin}{\oddsidemargin}
%    \end{macrocode}
% Text height is 225\,mm, top margin is 14\,mm. 
%    \begin{macrocode}
\setlength{\textheight}{225mm} 
\setlength{\topmargin}{-1in} 
\addtolength{\topmargin}{14mm}
%    \end{macrocode}
% Space between columns in two-column styles, header, footer,
% paragraph indentation:
%    \begin{macrocode}
\setlength{\columnsep}{9.5mm}
\setlength{\headsep}{0.7cm}
\setlength{\footskip}{0.75cm}
\setlength{\parindent}{1.5em}
%    \end{macrocode}
% 
% \begin{macro}{\@set@scalable@length}
% Try out a new scheme for dealing with various spacing dimensions and
% allowing the user to easily modify them. The command
% \begin{quote}
% |\set@scalable@length|\marg{name}\marg{defaultlength}\marg{stretch}\marg{shrink}
% \end{quote}
% does a 
% |\setlength|\marg{\cs{name}}\marg{defaultlength plus stretch minus shrink}
% and also defines the user command
% \textbackslash\meta{name}|Scale|\marg{factor} that can 
% be used to rescale length by \meta{factor}. Additionally, now
% defines the package command
% \textbackslash\meta{name}|@set@scale|\marg{FACTOR} 
% that also rescales the length by \meta{FACTOR}, so that the length
% becomes \meta{factor} $\times$ \meta{FACTOR} $\times$
% \meta{defaultlength}.   
%    \begin{macrocode}
\newcommand\@set@scalable@length[4]{%
  %% arg #1: NAME of length without leading `\'
  %%
  %% arg #2: default VALUE  
  \expandafter\newlength\csname @#1@default\endcsname
  \expandafter\setlength\csname @#1@default\endcsname{#2}%
  \expandafter\newlength\csname @@#1@@default\endcsname
  \expandafter\setlength\csname @@#1@@default\endcsname{#2}%
  %% arg #3: default STRETCH  
  \expandafter\newlength\csname @#1@stretch\endcsname
  \expandafter\setlength\csname @#1@stretch\endcsname{#3}%
  \expandafter\newlength\csname @@#1@@stretch\endcsname
  \expandafter\setlength\csname @@#1@@stretch\endcsname{#3}%
  %% arg #4: default SHRINK  
  \expandafter\newlength\csname @#1@shrink\endcsname
  \expandafter\setlength\csname @#1@shrink\endcsname{#4}%
  \expandafter\newlength\csname @@#1@@shrink\endcsname
  \expandafter\setlength\csname @@#1@@shrink\endcsname{#4}%
  %%
  %% User-defined scale \@NAME@scale is unity by default
  \expandafter\gdef\csname @#1@scale\endcsname{1}%
  %% Package-defined scale \@NAME@SCALE is unity by default
  \expandafter\gdef\csname @#1@SCALE\endcsname{1}%
  %%
  %% internal command \@reset@NAME should expand to 
  %% \setlength\@@NAME@@default{\@NAME@SCALE\@NAME@default}
  %% \setlength\NAME{\@NAME@scale\@@NAME@@default}
  %% in other words: scale x SCALE x default 
  \expandafter\gdef\csname @reset@#1\endcsname{%
    \expandafter\setlength\csname @@#1@@default\endcsname{%
      \csname @#1@SCALE\endcsname \csname @#1@default\endcsname
      }
    \expandafter\setlength\csname @@#1@@stretch\endcsname{%
      \csname @#1@SCALE\endcsname \csname @#1@stretch\endcsname
      }
    \expandafter\setlength\csname @@#1@@shrink\endcsname{%
      \csname @#1@SCALE\endcsname \csname @#1@shrink\endcsname
      }
    \expandafter\setlength\csname #1\endcsname{%
      \csname @#1@scale\endcsname \csname @@#1@@default\endcsname
      \@plus
      \csname @#1@scale\endcsname \csname @@#1@@stretch\endcsname
      \@minus
      \csname @#1@scale\endcsname \csname @@#1@@shrink\endcsname
      }}
  %%
  %% User command \NAMEScale should expand to 
  %% {\renewcommand\@NAME@scale{#1}\@reset@NAME}
  \expandafter\gdef\csname #1Scale\endcsname##1{%
    \expandafter\gdef\csname @#1@scale\endcsname{##1}%
    \csname @reset@#1\endcsname}
  %%
  %% Package command \NAME@set@SCALE should expand to 
  %% {\renewcommand\@NAME@SCALE{#1}\@reset@NAME}
  \expandafter\gdef\csname #1@set@SCALE\endcsname##1{%
    \expandafter\gdef\csname @#1@SCALE\endcsname{##1}%
    \csname @reset@#1\endcsname}
  %%
  %% initialize at default values with \@reset@NAME
  \csname @reset@#1\endcsname
  }
%    \end{macrocode}
% \end{macro}
% \begin{macro}{\parskip}
% First use of this new infrastructure: set interparagraph space to be
% zero but with a small amount of stretchiness.
%    \begin{macrocode}
%   \ifthenelse{\boolean{@manuscript}}{%
%     \@set@scalable@length{parskip}{0.5\baselineskip}{0.3\p@}{0.1\p@}
%   }{
    \@set@scalable@length{parskip}{0\p@}{0.3\p@}{0\p@}
%   }
%    \end{macrocode}
% \end{macro}
% 
% \begin{macro}{\abovesectionskip}
% \changes{4.0}{2019/11/12}{Use negative values for skip so that first paragraph after section headings is not indented.  This is something that I have been trying to fix for years!  It turns out that the anser was in the documentation for the standard command \cs{@startsection.}}
% \begin{macro}{\belowsectionskip}
% \begin{macro}{\aboveparaskip}
% \begin{macro}{\belowparaskip}
% Define some utility rubber lengths for use in
% section~\ref{sec:section-commands}. 
%    \begin{macrocode}
\newlength{\abovesectionskip}
\newlength{\belowsectionskip}
\newlength{\aboveparaskip}
\newlength{\belowparaskip}
%    \end{macrocode}
% The default values for these lengths are those traditionally used
% for the onecolumn style. 
%    \begin{macrocode}
\@set@scalable@length{abovesectionskip}{-3.5ex}{-4ex}{-2ex}
\@set@scalable@length{belowsectionskip}{2.3ex}{0ex}{2ex}
\@set@scalable@length{aboveparaskip}{3.25ex}{2ex}{0.5ex}
\@set@scalable@length{belowparaskip}{-1em}{\z@}{\z@}
%    \end{macrocode}
% \end{macro}\end{macro}\end{macro}\end{macro}
% We also make scalable some of the standard float parameters. These
% are the same defaults as in the standard article classes, but
% slightly more stretchy. For the twocolumn style, they are currently
% halved, see below. Note that I have also lumped |\abovecaptionskip|
% in with these, since it plays a similar role. There is also a
% |\belowcaptionskip| but it is set to zero in the standard classes. 
%    \begin{macrocode}
\@set@scalable@length{floatsep}{12\p@}{4\p@}{4\p@}
\@set@scalable@length{textfloatsep}{20\p@}{4\p@}{8\p@}
\@set@scalable@length{intextsep}{12\p@}{4\p@}{4\p@}
\@set@scalable@length{dblfloatsep}{12\p@}{4\p@}{4\p@}
\@set@scalable@length{dbltextfloatsep}{20\p@}{4\p@}{8\p@}
\@set@scalable@length{abovecaptionskip}{10\p@}{\z@}{\z@}
%    \end{macrocode}
% \begin{macro}{\@before@title@skip}         
% \begin{macro}{\@before@author@skip}        
% \begin{macro}{\@between@author@skip}       
% \begin{macro}{\@before@abstract@skip}      
% \begin{macro}{\@after@abstracthead@skip}   
% \begin{macro}{\@between@abstract@skip}     
% \begin{macro}{\@before@kwd@skip}           
% \begin{macro}{\@before@dates@skip}           
% Spacing parameters used in |\@maketitle|. These are not changed by
% default except by the user, so they don't need to go into
% |\@package@default@lengths|.  
%    \begin{macrocode}
\newlength{\@before@title@skip}         
\newlength{\@before@author@skip}        
\newlength{\@between@author@skip}       
\newlength{\@before@abstract@skip}      
\newlength{\@after@abstracthead@skip}   
\newlength{\@between@abstract@skip}     
\newlength{\@before@kwd@skip}           
\newlength{\@before@dates@skip}           
\@set@scalable@length{@before@title@skip}{2em}{\z@}{\z@}
\@set@scalable@length{@before@author@skip}{1.5em}{\z@}{\z@}
\@set@scalable@length{@between@author@skip}{0.5em}{\z@}{\z@}
\@set@scalable@length{@before@abstract@skip}{3ex}{1ex}{0.5ex}
\@set@scalable@length{@after@abstracthead@skip}{1ex}{\z@}{\z@}
\@set@scalable@length{@between@abstract@skip}{2ex}{1ex}{0.5ex}
\@set@scalable@length{@before@kwd@skip}{1ex}{0.5ex}{0.2ex}
\@set@scalable@length{@before@dates@skip}{2ex}{0.5ex}{0.2ex}
%    \end{macrocode}
% \end{macro}\end{macro}\end{macro}\end{macro}
% \end{macro}\end{macro}\end{macro}\end{macro}
% 
% Spacing above and below displayed equations also wants to be done
% like this, but it is more complicated since the standard article
% classes make these change with the fontsize. Something to do later,
% maybe. 
% 
%    \begin{macrocode}
%    \end{macrocode}
% \begin{macro}{\@package@default@lengths}
% The value actually used for these lengths is controlled by the
% \meta{SCALE} parameter, which depends on the style option that is in
% effect. This is set using the |\@package@default@lengths| command,
% which must be called every time the style changes (e.g., changing
% from onecolumn to twocolumn). 
%    \begin{macrocode}
\newcommand{\@package@default@lengths}{%
  \ifthenelse{\boolean{@rmtwocol}}{%
    %% twocolumn mode
    \abovesectionskip@set@SCALE{0.5}
    \belowsectionskip@set@SCALE{0.5}
    \aboveparaskip@set@SCALE{0.5}
    \belowparaskip@set@SCALE{1}
    \floatsep@set@SCALE{0.5}
    \textfloatsep@set@SCALE{0.5}
    \intextsep@set@SCALE{0.5}
    \dblfloatsep@set@SCALE{0.5}
    \dbltextfloatsep@set@SCALE{0.5}
    }{%
    %% onecolumn mode
    \abovesectionskip@set@SCALE{1}
    \belowsectionskip@set@SCALE{1}
    \aboveparaskip@set@SCALE{1}
    \belowparaskip@set@SCALE{1}
    \floatsep@set@SCALE{1}
    \textfloatsep@set@SCALE{1}
    \intextsep@set@SCALE{1}
    \dblfloatsep@set@SCALE{1}
    \dbltextfloatsep@set@SCALE{1}
    }%
  }
\@package@default@lengths
%    \end{macrocode}
% \end{macro}
% \begin{macro}{\RescaleLengths}
% \begin{macro}{\RescaleSecLengths}
% \begin{macro}{\RescaleFloatLengths}
% \begin{macro}{\RescaleTitleLengths}
% User commands that will fiddle with all the scalable lengths. 
%    \begin{macrocode}
\newcommand{\RescaleLengths}[1]{%
  \RescaleSecLengths{#1}\RescaleFloatLengths{#1}\RescaleTitleLengths{#1}}
\newcommand{\RescaleSecLengths}[1]{%
%% sections
  \abovesectionskipScale{#1}
  \belowsectionskipScale{#1}
  \aboveparaskipScale{#1}
  }
\newcommand{\RescaleFloatLengths}[1]{%
%% floats
  \floatsepScale{#1}
  \textfloatsepScale{#1}
  \intextsepScale{#1}
  \dblfloatsepScale{#1}
  \dbltextfloatsepScale{#1}
  \abovecaptionskipScale{#1}
  }
\newcommand{\RescaleTitleLengths}[1]{%
%% maketitle
  \@before@title@skipScale{#1}           
  \@before@author@skipScale{#1}        
  \@between@author@skipScale{#1}       
  \@before@abstract@skipScale{#1}      
  \@after@abstracthead@skipScale{#1}   
  \@between@abstract@skipScale{#1}     
  \@before@kwd@skipScale{#1}           
  \@before@dates@skipScale{#1}           
  }
%    \end{macrocode}
% \end{macro}\end{macro}\end{macro}\end{macro}
% \changes{2.2beta14}{2002/03/17}{Finished implementation of new
% infrastructure for scalable lengths. Converted the above/below
% section skips to use it. Also, all the floatsep parameters. Also,
% all the vertical lengths that are used by |\@maketitle|. }

% \begin{macro}{\MSextramargin}
% For the \texttt{manuscript} option we have an extra 0.3\,inch
% vertical margin and 1\,inch horizontal margin for marginal notes (by
% hand or with |\marginpar|. 
%
% \changes{3.23}{2007/09/09}{Reduce baselinestretch in
%   \texttt{manuscript} style to 1.2. No longer use raggedright for
%   manuscript.}
%
% \changes{3.24}{2007/09/14}{Further reduction of baselinestretch in
%   \texttt{manuscript} style to 1.1 since we now have wider margins.}
%
% \changes{3.24}{2007/09/14}{Split \texttt{manuscript}'s extra margin
%   into separate horizontal and vertical parts.}
% 
%    \begin{macrocode}
\newlength{\MSextraHmargin}\newlength{\MSextraVmargin}
\setlength{\MSextraHmargin}{1in} 
\setlength{\MSextraVmargin}{0.3in} 
\if@manuscript
   \renewcommand{\baselinestretch}{1.1}
   \addtolength{\textwidth}{-2\MSextraHmargin}
   \addtolength{\oddsidemargin}{\MSextraHmargin}
   \setlength{\evensidemargin}{\oddsidemargin}
   \addtolength{\textheight}{-2\MSextraVmargin}
   \addtolength{\topmargin}{\MSextraVmargin}
%    \raggedright
\fi
%    \end{macrocode}
% \end{macro}
% Float parameters: don't allow bottom floats by default. 
%    \begin{macrocode}
\renewcommand{\fps@figure}{tp}
\renewcommand{\fps@table}{tp}
%    \end{macrocode}
%
% \begin{macro}{\rmaa@AuthorWidthFrac}
% \begin{macro}{\rmaa@AbstractWidthFrac}
% These parameters control what fraction of the total text width can
% be used by the list of authors and the abstract/resumen,
% respectively.  
%    \begin{macrocode}
\newcommand{\rmaa@AuthorWidthFrac}{1.0}
\newcommand{\rmaa@AbstractWidthFrac}{0.75}
%    \end{macrocode}
% Use a wide abstract to save space in the new style.
%    \begin{macrocode}
\if@newstyle\renewcommand{\rmaa@AbstractWidthFrac}{1.0}\fi
%    \end{macrocode}
% \end{macro}
% \end{macro}
%
% Some things to consider in future versions. Trying to avoid too many
% widows, etc. Apparently, the following will work best with
% \cs{flushbottom}. 
%    \begin{macrocode}
%%\widowpenalty=10000
%%\clubpenalty=10000
%    \end{macrocode}
% Also, should perhaps try getting rid of underfull/overfull box warnings. 
% \changes{3.10}{2003/01/24}{Set large values for \cs{hbadness},
% \cs{vbadness}, \cs{hfuzz}, \cs{vfuzz}. To reduce amount of stuff
% that gets written to the terminal every time you compile something.}
%    \begin{macrocode}
\hbadness=10000
\vbadness=10000
\setlength\hfuzz{\maxdimen}
\setlength\vfuzz{\maxdimen}
%    \end{macrocode}
% \subsection{Various text styles for different contexts}
% \label{sec:text-styles} 
% \begin{macro}{\rmaa@AffilStyle}
% \begin{macro}{\rmaa@PagesStyle}
% \begin{macro}{\rmaa@YearStyle}
% \begin{macro}{\rmaa@VolumeStyle}
% \begin{macro}{\rmaa@VolumenStyle}
% \begin{macro}{\rmaa@JournalNameStyle}
% \begin{macro}{\rmaa@ReceivedStyle}
% \begin{macro}{\rmaa@AbstractHeaderStyle}
% \begin{macro}{\rmaa@JournalNameStyle}
% Style commands for the different fields used by |\maketitle| in
% constructing the title page of the article. These have the form
% |\rmaa@XYZStyle{TEXT}|. I'm not sure these were really a good idea,
% but never mind \dots 
%    \begin{macrocode}
\newcommand{\rmaa@AffilStyle}[1]{{\normalsize\rmfamily#1}}
\newcommand{\rmaa@PagesStyle}[1]{{\normalsize\rmfamily#1}}
\newcommand{\rmaa@YearStyle}[1]{{\normalsize\rmfamily#1}}
\newcommand{\rmaa@VolumeStyle}[1]{{\normalsize\rmfamily\bfseries#1}}
\newcommand{\rmaa@VolumenStyle}[1]{{\Large\sffamily\bfseries#1}}
\newcommand{\rmaa@JournalNameStyle}[1]{{\normalsize\itshape#1}}
\newcommand{\rmaa@ReceivedStyle}[1]{{\normalsize\rmfamily\itshape#1}}
\newcommand{\rmaa@AbstractHeaderStyle}[1]{%
  {\normalsize\rmfamily\MakeUpperCase{#1}}}
\newcommand{\rmaa@TitleStyle}[1]{%
  {\large\bfseries\rmfamily\boldmath\MakeUpperCase{#1}\par}}
%    \end{macrocode}
% \end{macro}
% \end{macro}
% \end{macro}
% \end{macro}
% \end{macro}
% \end{macro}
% \end{macro}
% \end{macro}
% \end{macro}
%
% \begin{macro}{\rmaa@FrontSecStyle}
% Headings of ``front sections'' (see section~\ref{sec:front-sections})
%    \begin{macrocode}
\newcommand{\rmaa@FrontSecStyle}[1]{%
  {\normalsize\bfseries\rmfamily\MakeUpperCase{#1}}}
%    \end{macrocode}
% \end{macro}
%
% \begin{macro}{\rmaa@TOCTitleStyle}
% \begin{macro}{\rmaa@TOCPartStyle}
% \begin{macro}{\rmaa@IndexTitleStyle}
% \begin{macro}{\rmaaTOCAuthorStyle}
% \begin{macro}{\rmaaIndexAuthorStyle}
% Commands to set style of TOC and index entries.
% Rewritten to avoid unnecessary font commands in the .toc and .ind files
%    \begin{macrocode}
\newcommand{\rmaa@TOCTitleStyle}[1]{\MakeUpperCase{#1}}
\newcommand{\rmaa@TOCPartStyle}[1]{\rmaa@TOCTitleStyle{#1}}
\newcommand{\rmaa@IndexTitleStyle}[1]{#1}
\newcommand{\rmaaTOCAuthorStyle}[1]{%
  {\itshape#1}}
\newcommand{\rmaaIndexAuthorStyle}[1]{%
  \protect\mbox{\bfseries#1}}
%    \end{macrocode}
% \end{macro}
% \end{macro}
% \end{macro}
% \end{macro}
% \end{macro}
%
% \begin{macro}{\rmaa@KeywordStyle}
% \begin{macro}{\rmaa@KeywordHeaderStyle}
% The style of the keywords depends on whether we are using the
% \texttt{newstyle} option or not. If we are, then they are small. 
%    \begin{macrocode}
\ifthenelse{\boolean{@newstyle}}{
  \newcommand{\rmaa@KeywordStyle}[1]{%
    {\small%
      \if@boldkeywords 
      \bfseries\MakeUpperCase{\raggedright #1}
      \else
      \raggedright #1
      \fi
    }}
  \newcommand{\rmaa@KeywordHeaderStyle}[1]{%
    {\small\itshape#1}}}
{
  \newcommand{\rmaa@KeywordStyle}[1]{%
    {\normalsize%
      \if@boldkeywords\bfseries\fi\raggedright #1}}
  \newcommand{\rmaa@KeywordHeaderStyle}[1]{%
    {\normalsize\itshape#1}}}
%    \end{macrocode}
% \end{macro}
% \end{macro}
%
% \subsection{Page styles}\label{sec:page-styles}
% \begin{macro}{\ps@rmaa}
% The running headers for normal pages:
%    \begin{macrocode}
\newcommand{\ps@rmaa}{%
%    \end{macrocode}
% nothing in either foot; 
%    \begin{macrocode}
  \renewcommand\@oddfoot{\relax}%
  \renewcommand\@evenfoot{\relax}%
%    \end{macrocode}
% page number and short author list in the even heads; 
%    \begin{macrocode}
  \renewcommand\@evenhead{\rlap{\thepage}\hfil\@shortauthor\hfil}%
%    \end{macrocode}
% short title and page number in the odd heads.
%    \begin{macrocode}
  \renewcommand\@oddhead{\hfil\@shorttitle\hfil\llap{\thepage}}
  }
%    \end{macrocode}
% \begin{macro}{\ps@rmaafirstpage}
% The header of the first page of each
% article:
%    \begin{macrocode}
\newcommand{\ps@rmaafirstpage}{%
  \renewcommand\@oddfoot{\hfil\thepage}%
%    \end{macrocode}
% footer corrected (06 Jun 2000) to have page on outside edge;
%    \begin{macrocode}
  \renewcommand\@evenfoot{\thepage\hfil}%
%    \end{macrocode}
% Header has journal name, volume, and page numbers for a regular
% article, or similar for special cases. 
% \changes{4.0}{2019/11/03}{Change to three-line header with copyright info and DOI.}
%    \begin{macrocode}
\if@preprint
\renewcommand\@oddhead{%
  \if@proceedings 
  \textsc{To appear in ``\@conftitle{} (\rmaa@YearStyle{\@year})''}\hfil
  \textit{RMxAC}
  \else
  \textsc{To appear in} \@journalname\hfil
  \fi}\fi%
\if@guide
\renewcommand\@oddhead{\textsc{Documentation} (\rmaa@YearStyle{\@year})\hfil}
\fi%
\if@manuscript
\renewcommand\@oddhead{Manuscript for \@journalname{} (\rmaa@YearStyle{\@year})\hfil}
\fi%
\if@journal
\renewcommand\@oddhead{%
  \parbox{\textwidth}{
    \@journalname, \rmaa@VolumeStyle{\@volume},
    \rmaa@PagesStyle{\@artpages}\
    (\rmaa@YearStyle{\@year})\hfil\\[1pt]
    {\small\textcopyright{} \@year:
      Instituto de Astronom\'\i{}a,
      Universidad Nacional Aut\'onoma de M\'exico}\hfil\\[1pt]
    {\small\@DOI@url}\hfil
  }%
}\fi%
  \renewcommand\@evenhead{\@oddhead}}
%    \end{macrocode}
% By default, use the |\ps@rmaa| page style.
%    \begin{macrocode}
\pagestyle{rmaa}
%    \end{macrocode}
% \end{macro}
% \end{macro}
%
% \subsection{Preamble commands}\label{sec:preamble} 
% Everything you need for setting the title page, TOC and index
% entries, etc.
% \subsubsection{Commands for the title page}
% \label{sec:title-commands}
% \begin{macro}{\resumen}
% \begin{macro}{\@resumen}
% \begin{macro}{\abstract}
% \begin{macro}{\@abstract}
% User commands for setting the resumen and abstract. 
% These should be used in the document header. The internal commands
% |\@resumen| and |\@abstract| are used to hold the actual text
% pending output via the |\maketitle| command. 
%    \begin{macrocode}
\newif\if@setresumen
\newcommand{\resumen}[1]{%
%%  \renewcommand{\@resumen}{{\selectlanguage{spanish}#1}}%
  \renewcommand{\@resumen}{{\ignorespaces#1}}%
  \@setresumentrue}
\renewcommand{\abstract}[1]{%
  \renewcommand{\@abstract}{{\ignorespaces#1}}
  \renewcommand\ADS@abstract{#1}
%% Not a good idea, since then author macros are not expanded and
%% cause errors later. However, by not doing it we are vulnerable to
%% fragile commands.....
  \if@noexpandADS\@onelevel@sanitize{\ADS@abstract}\fi
  \@abstractpresenttrue
  }
\newcommand{\ADSnoexpand}{\@noexpandADStrue}
%    \end{macrocode}
% \begin{macro}{\boldabstract}
% This is the version for doing 1-page conference proceedings. Still
% save the text for ADS but then write it out immediately as a bold
% paragraph. 
%    \begin{macrocode}
\newcommand{\boldabstract}[1]{%
  \gdef\ADS@abstract{#1}%
  \renewcommand{\@abstract}{#1}%
  \if@noexpandADS\@onelevel@sanitize{\ADS@abstract}\fi%
  \noindent\textbf{\boldmath #1}
  }
%    \end{macrocode}
% \end{macro}
% Default value in each case is to print out a warning if not set
%    \begin{macrocode}
  \newcommand{\@resumen}{%
    \@setresumenfalse \ClassWarningNoLine{rmaa}{No Spanish abstract
      (resumen) specified - if you do not speak Spanish, this will be
      translated by the editors}%
    \textbf{Favor de proporcionar un resumen en espa\~nol. If you are
      unable to translate your abstract into Spanish, the editors will
      do it for you.} }
  \newcommand{\@abstract}{\ClassWarningNoLine{rmaa}{No abstract
      specified}}
%    \end{macrocode}
% \end{macro}
% \end{macro}
% \end{macro}
% \end{macro}
%
% \begin{macro}{\title}
% \begin{macro}{\@rawtitle}
% \begin{macro}{\@TOCtitle}
% \begin{macro}{\rmaaIndexTitle}
% Similarly, |\title| is the preamble command to set the title of the
% article. This is used in various places, typeset in different
% styles, hence the use of the internal |\@rawtitle|, |\@title|,
% |\rmaaIndexTitle|, and |\@TOCtitle| commands, each of which uses the
% relevant style from section~\ref{sec:text-styles}
%    \begin{macrocode}
\newcommand{\@rawtitle}{Raw Title}
\newcommand{\@TOCtitle}{\rmaa@TOCTitleStyle{Title for TOC}}
\newcommand{\rmaaIndexTitle}{\rmaa@IndexTitleStyle{Title for index}.}
\newcommand{\@TOCaftertitle}{}
\renewcommand{\title}[1]{%
%%  {\ADS@sanitize\xdef\@rawtitle{#1}}
  \renewcommand\@rawtitle{\noexpand#1}
  \let\ADS@title=\@rawtitle
  \renewcommand{\@title}{\rmaa@TitleStyle{#1}}%
  \renewcommand{\@TOCtitle}{\rmaa@TOCTitleStyle{#1}\@TOCaftertitle}%
  \renewcommand{\rmaaIndexTitle}{\rmaa@IndexTitleStyle{#1\@TOCaftertitle}.}%
  }
%    \end{macrocode}
% \end{macro}
% \end{macro}
% \end{macro}
% \end{macro}
% 
% All the weird and whacky ways of indicating the author
% affiliations. 
% \begin{macro}{\affil}
% First, the way that sets it directly underneath the authors name.
% Note that |\affil| trumps all other affiliations for the purposes of
% getting into the ADS stuff (we use |\gdef|).   
% \dots 
%    \begin{macrocode}
\newcommand{\affil}[1]{%
  \\[1.5ex]\rmaa@AffilStyle{#1}\\[1.5ex]%
  \if@firstaffil\gdef\ADS@affil{#1}\@firstaffilfalse\fi%
  }
\renewcommand{\and}{and\\[1.5ex]}
%    \end{macrocode}
% \begin{macro}{\altaffil}
% Alternative affiliations are footnoted. 
%    \begin{macrocode}
\newcommand{\altaffil}[1]{%
  \if@firstaltaffil\gdef\ADS@affil{#1}\@firstaltaffilfalse\fi%
  \thanks{#1}}
%    \end{macrocode}
% \begin{macro}{\altaffilmark}
% \begin{macro}{\altaffiltext}
% Alternative-alternative mechanism for affiliations---should not be
% mixed with |\altaffil| (unless you know what you are doing) or there
% will be problems with numbering.
%    \begin{macrocode}
\newcommand{\altaffilmark}[1]{\textsuperscript{#1}}
\newcommand{\altaffiltext}[2]{%
  \if@firstaltaffil\gdef\ADS@affil{#2}\@firstaltaffilfalse\fi
  \stepcounter{footnote}%
  \protected@xdef\@thanks{\@thanks
    \protect\footnotetext[#1]{#2}}} 
%    \end{macrocode}
% \end{macro}
% \end{macro}
% \end{macro}
% \end{macro}
%
% \begin{macro}{\keywords}
% \begin{macro}{\@keywords}
% \begin{macro}{\@keywordheader}
% \label{keywords}Sundry minor preamble commands. All
% self-explanatory, I hope. |\keywords| is the old user interface
% (deprecated). New articles should use |\addkeyword| once for each
% keyword. That way the ADS record comes out right too (comma rather
% than em-dash separator).
% \changes{2.2beta12}{2001/12/07}{\cs{addkeyword} New user interface
%   for keywords.} 
% \changes{2.2beta13}{2002/03/11}{\cs{addkeyword} Improved behaviour
%   with .ads file. Use
%   \cs{ADS@sanitize} to change definitions of `~' and \cs{ion}.}
% \changes{3.23}{2007/09/09}{\cs{addkeyword} now signals a compilation
% error if the author tries to add more than 6 keywords.}
% \changes{3.24}{2007/09/14}{kwdcounter is now reset in
% \cs{@begofarticlehook} to prevent error in book style} 
%    \begin{macrocode}
\newcommand{\@kwdsep}{ --- }
\newcommand{\ADS@kwdsep}{, }
\newcounter{kwdcounter}\setcounter{kwdcounter}{0}
\newcommand{\addkeyword}[1]{%
  \addtocounter{kwdcounter}{1}
  \ifthenelse{\thekwdcounter > 6}{
    \ClassError{rmaa}{Too many keywords. A maximum of 6 keywords is allowed.}{}
  }{}
  \if@firstkwd%
  \renewcommand\@keywords{\if@boldkeywords\uppercase{#1}\else #1\fi}
  {\ADS@sanitize\def\ion{\plainion}\gdef\ADS@keywords{#1}}
  \@firstkwdfalse
  \else%
  \edef\@keywords{\@keywords\@kwdsep \if@boldkeywords\uppercase{#1}\else #1\fi}%
  {\ADS@sanitize\def\ion{\plainion}\xdef\ADS@keywords{\ADS@keywords\ADS@kwdsep #1}}%
  \fi}
\newcommand{\ADS@sanitize}{%
  \def\mathrm{}
  \def\@{\ }
  \def~{\ }
  \def\langle{<}
  \def\rangle{>}
  \def\log{log }
  \def\sin{sin }
  \def\cos{cos }
  \def\tan{tan }
  \def\fd{^d.}
  \def\fh{^h.}
  \def\fm{^m.}
  \def\fs{^s.}
  \def\fdg{�.}
  \def\farcm{'.}
  \def\farcs{''.}
  \def\arcdeg{�}
  \def\,{ }
  }
\newcommand{\keywords}[1]{%
  \renewcommand{\ADS@keywords}{#1}
  \renewcommand{\@keywords}{\uppercase{#1}}
  }
\newcommand{\@keywordheader}{\rmaa@KeywordHeaderStyle{Key Words:}}
\newcommand{\@keywords}{%
  \ClassWarningNoLine{rmaa}{No keywords specified}%
  Keywords go here
  }
%    \end{macrocode}
% \begin{macro}{\shorttitle}
% \begin{macro}{\shortauthor}
% \begin{macro}{\@shorttitle}
% \begin{macro}{\@shortauthor}
% Commands for setting the running headers.
%    \begin{macrocode}
\newcommand{\shorttitle}[1]{%
  \renewcommand\@shorttitle{\MakeUpperCase{#1}}}
\newcommand{\shortauthor}[1]{%
  \renewcommand\@shortauthor{\MakeUpperCase{#1}}}
\newcommand{\@shorttitle}{SHORT TITLE GOES HERE}
\newcommand{\@shortauthor}{SHORT AUTHOR GOES HERE}
%    \end{macrocode}
% \end{macro}
% \end{macro}
% \end{macro}
% \end{macro}
% \end{macro}
% \end{macro}
% \end{macro}


% \subsubsection{The full addresses}
% \label{sec:full-addresses}
% \begin{macro}{\fulladdresses}
% \begin{macro}{\@fulladdresses}
% \begin{macro}{\@fulladdheight}
% \begin{macro}{\@fulladddepth}
% \begin{macro}{\@fulladdbox}
% Aaghh, the full addresses! A perennial source of problems. 
%    \begin{macrocode}
\newlength{\@fulladdheight}
\newlength{\@fulladddepth}
\newsavebox{\@fulladdbox}
%% \fulladdresses command modified (11 May 2000) so that we can put
%% \fulladdresses{} in the article to suppress the output of the
%% addresses
\newcommand{\@FinalFigure}{\relax}
\newcommand{\FinalFigure}[1]{\renewcommand{\@FinalFigure}{#1}}
\newcommand{\fulladdresses}[1]{%
  \renewcommand{\@rmaaArg}{#1}
  \ifx\@rmaaArg\@rmaaEmpty
     \renewcommand{\@fulladdresses}{}
  \else
     \renewcommand{\@fulladdresses}{
       \normalsize%
       \ifx\@FinalFigure\relax\else\vfill\@FinalFigure\vfill\fi
%%       \vspace*{1ex}
       %% do addresses as a list (copied from thebibliography environment) 
       \list{\null}{\leftmargin 1.5em\topsep\z@\labelwidth\z@\labelsep\z@\itemindent -1.5em
         \itemsep\z@\parsep\z@\usecounter{enumi}}
       \sloppy\clubpenalty4000\widowpenalty4000
       \sfcode`\.=1000\relax
       #1
       \endlist
       }
  \fi
  }
\newcommand{\@fulladdresses}{Full addresses go here}
%    \end{macrocode}
% \begin{macro}{\suppressfulladdresses}
% New user-level command for suppressing full addresses should you so
% desire. Currently used in the single-page-poster style. 
%    \begin{macrocode}
\newcommand{\suppressfulladdresses}{\@suppressfulladdressestrue\@needaddressesfalse}
%    \end{macrocode}
% \end{macro}
% \end{macro}
% \end{macro}
% \end{macro}
% \end{macro}
% \end{macro}

% \begin{macro}{\printdates}
% \begin{macro}{\noprintdates}
% Command to activate the printing out of the received/accepted
% dates. This can also be done using the \texttt{printdates} option,
% which is the default for \texttt{journal}, but not for \texttt{proceedings}.
%    \begin{macrocode}
\newcommand{\printdates}{\@printdatestrue}
\newcommand{\noprintdates}{\@printdatesfalse}
%    \end{macrocode}
% \end{macro}
% \end{macro}
% \begin{macro}{\@outputaddressesifnecessary}
% Command to output the addresses if we haven't already done so
% This is called by: 
% \begin{enumerate}
% \item The end of |thebibliography| environment
% \item The |\outputfulladdresses| command
% \end{enumerate}
%    \begin{macrocode}
\newcommand{\@outputaddressesifnecessary}{%
  \ClassInfo{rmaa}{Calling @outputaddressesifnecessary}
  \ifthenelse{\boolean{@needaddresses}}{%
    %% If in one column style, or in the first column of a two column
    %% style then we can try to put it on this page 
    \ifthenelse{\boolean{@firstcolumn} \or \not\boolean{@rmtwocol}}{%
      \@outputaddresses}{%
      %% otherwise give warning it will be on a separate page
      \ClassWarning{rmaa}{Full addresses will go on a separate
        page (please see User Guide for how to remedy this)}}%
    %% If we haven't managed yet, put it on its own page
    \ifthenelse{\boolean{@needaddresses}}{%
      \if@debug{\ClassWarning{rmaa}{In this silly bit!}}\fi
      \clearpage\vfill\@outputaddresses}{}%
    }{
    %% If we don't need them, we still might need to do the eoa hook 
    \ifthenelse{\boolean{@suppressfulladdresses}}{%
      \if@debug{\ClassWarning{rmaa}{XXX}}\fi
%%      \if@book\else\@endofarticlehook\null\fi
      }{}
    }%
  }  
%    \end{macrocode}
% \end{macro}

% \begin{macro}{\outputfulladdresses}
% A user interface to the above command, in case there is no biblio
% (11 May 2000) 
%    \begin{macrocode}
\newcommand{\outputfulladdresses}{\@outputaddressesifnecessary}
%    \end{macrocode}
% \end{macro}

% \begin{macro}{\adjustfinalcols}
% This command starts a new column and should be 
% used to ensure that the columns balance on the last page. It also
% outputs the full 
% author addresses to ensure that this occurs in the left hand
% column. 
% The |\par| command is so that the current line will be taken account
% of when working out |\pagetotal|
%    \begin{macrocode}
\newcommand{\adjustfinalcols}{\if@twocolumn\goodbreak\@outputaddresses\newpage\noindent\fi}
%    \end{macrocode}
% \end{macro}

% \begin{macro}{\@outputaddresses}
% Command to do the work of outputting the addresses. 
% This can be called by 
% \begin{enumerate}
% \item an explicit |\adjustfinalcols|
% \item via |\@outputaddressesifnecessary|, q.v.
% \end{enumerate}
%
%    \begin{macrocode}
\newcommand{\@outputaddresses}{%
  % Allow option of suppressing the full addresses, in which case
  % we just run the hook if we are not in a book style
  \if@suppressfulladdresses
%%     \if@book\else\@endofarticlehook\null\fi
  \else
  %% save the full addresses in a box for output later
  \sbox{\@fulladdbox}{\parbox{\textwidth}{\@fulladdresses}}
  % calculate height of address box  
  \settoheight{\@fulladdheight}{\usebox{\@fulladdbox}}
  \settodepth{\@fulladddepth}{\usebox{\@fulladdbox}}
  \addtolength{\@fulladdheight}{\@fulladddepth}
  \if@debug
  { \typeout{Address total height: \the\@fulladdheight}
    \typeout{Text height: \the\pagetotal}
    \typeout{Goal height: \the\pagegoal}
    }
  \fi
  % reuse \@fulladdheight to mean the total height of text+addresses
  \addtolength{\@fulladdheight}{\pagetotal}
%%  \addtolength{\@fulladdheight}{-\baselineskip}
  \if@debug\typeout{Adjusted text + addresses height: \the\@fulladdheight}\fi
  % check it will fit on current page
  \ifdim\@fulladdheight<\pagegoal
     %% must shift it to the left
     %% to compensate for the hanging indentation in the reference list. 
     %% NOTE: this is not foolproof!
     %% NB. This will give overfull hbox warnings, which is a bit ugly.
     \vfill\noindent\if@biblist\hspace*{-\leftmargin}\fi
     \usebox{\@fulladdbox}
     % Set this to show we've done it
     \@needaddressesfalse
     % and assume we are at the end of the article
     % so that we do the EOA hook if not in the book style
%%     \if@book\else\@endofarticlehook\fi
  \else
     % if not, just do nothing. With luck, this function will get 
     % called again on the next page
     \if@debug{\typeout{Holding over full addresses to next page}}\fi
  \fi\fi
  }
%    \end{macrocode}
% \end{macro}

% \subsubsection{Commands for editors}
% \label{sec:editor-commands}
% Some user/editor commands for doing the TOC and index in the book
% option. 
% They are defined here rather in the book-specific part since they
% need to be always available in order to avoid breaking documents. 
% \begin{macro}{\listofauthors}
% \begin{macro}{\@rawauthors}
% \begin{macro}{\rmaaTOCauthors}
% The command |\listofauthors| formats an author list for use in the
% TOC. It should contain the same names as in the |\author| command,
% but without the affiliation garbage and using |\&| instead of
% |and|. 
%    \begin{macrocode}
\newcommand{\@rawauthors}{Raw Authors}
\newcommand{\rmaaTOCauthors}{\noexpand\rmaaTOCAuthorStyle{Author list for TOC}}
\newcommand{\listofauthors}[1]{%
  \renewcommand{\@rawauthors}{#1}
  \renewcommand\rmaaTOCauthors{\noexpand\rmaaTOCAuthorStyle{#1}}}
%    \end{macrocode}
% \begin{macro}{\indexauthor}\label{indexauthor}
% The command |\indexauthor| produces an index entry for each
% author. 
% This command must come \emph{after} the
% |\title| and the |\listofauthors| commands. Also, there should be one 
% |\indexauthor| command for \emph{each} author. We also add the
% author onto the end of the |\ADS@authors| macro and construct the
% |\ADS@authorinitial| by a most cunning method (see below), but only
% for the first author. 
% 
% No, No, NO! This was all wrong. Backed out completely: 
% \begin{quote}
% Completely rewritten (19 Dec 2002) since we now use the index class, so 
% as to be able to alphabetize accented characters. A side effect of this is
% that |\index| no longer expands out the argument after the |@|, so we 
% instead add all the |\index| commands to a list (|\indexauthor@deferlist|), 
% which is executed in the |\@begofarticlehook|. 
% 
% Note the stuff with |\toks0| is copied from Exercise 20.15 of the \TeX{}book---it is
% to make sure the latin1-accented name gets passed into the idx
% file. 
% \end{quote}
% We now go back to the old way of doing things, using an optional
% argument to the \cs{indexauthor} command that has the name without
% accents. 
% \changes{3.08}{2003/01/16}{Remerged changes from 3.06 that got
% dropped from 3.07 after a synchronization disaster. Basically, the
% accentless optional argument to \cs{indexauthor}. Must add the macro
% directories to my unison set-up to stop this sort of thing from happening.}
%    \begin{macrocode}
\newcommand\@isortkey{}
\newcommand\@emptykey{empty}
\newcommand\indexauthor[2][empty]{
  \renewcommand\@isortkey{#1}
  \ifx\@isortkey\@emptykey \renewcommand\@isortkey{#2}\fi
  \index{\@isortkey @{{\rmaaIndexAuthorStyle{#2}}\hskip1em plus1em minus0.4em
      \rmaaIndexTitle\hskip1em plus2em minus0.4em\penalty-100
      \rmaaTOCauthors}}
  \indexauthorADS{#2}
}
%%\def\indexauthor@deferlist{\relax}
%%\long\def\indexauthor#1{
%%  \def\@@name{#1} \toks0=\expandafter{\@@name}
%%  \toks1=\expandafter{\indexauthor@deferlist}
%%  \protected@xdef\indexauthor@deferlist{%
%%    \the\toks1
%%    \noexpand\index{\the\toks0@{\noexpand\rmaaIndexAuthorStyle{\the\toks0}\noexpand\hskip1em plus1em minus0.4em 
%%        \rmaaIndexTitle\noexpand\hskip1em plus2em minus0.4em\noexpand\penalty-100
%%        \rmaaTOCauthors}}
%%  }
%%  \indexauthorADS{#1}
%%  }
%%\xdef\indexauthor@deferlist{}
\newcommand\indexauthorADS[1]{%
  \if@firstauthor\protected@xdef\ADS@authorinitial{\firstletter{#1}}\fi
  \protected@xdef\ADS@authors{\ADS@authors\if@firstauthor\else ; \fi
    #1}
  \@firstauthorfalse
}
%    \end{macrocode}
% Note that the first |#1| sets the key for the index, making sure it
% actually starts with  
% first letter of the surname. Meanwhile, the |@{...}| expression,
% which starts with some  
% control sequences is what is written to the |.ind| file. 
% See the makeindex man page for details. 
% \end{macro}
% \end{macro}
% \end{macro}
% \end{macro}

% \begin{macro}{\KillAllAccents}
% \begin{macro}{\firstletter}
% These are utility functions used in the above. |lett| is a counter
% indicating the current letter in the string. |\firstletter| just
% prints the first letter of its argument.
%    \begin{macrocode}
\newcommand{\KillAllAccents}{\def\'{}\def\`{}\def\~{}\def\.{}\def\i{i}}
\newcommand\firstletter[1]{\@car #1\@nil}
%    \end{macrocode}
% \end{macro}
% \end{macro}
% \begin{macro}{\SetVolume}
% \begin{macro}{\SetMSnumber}
% \begin{macro}{\SetFirstPage}
% \begin{macro}{\SetYear}
% \begin{macro}{\ReceivedDate}
% \begin{macro}{\AcceptedDate}
% These are for the editors to set the volume number, etc. 
%    \begin{macrocode}
\newcommand{\SetVolume}[1]{%
  \renewcommand{\@volume}{#1}
  \edef\ADS@volume{\PadToFourWithDots{#1}}
  \renewcommand{\@volumen}{\rmaa@VolumenStyle{volumen #1}}%
  }
\newcommand{\SetMSnumber}[1]{\renewcommand{\@MSnum}{#1}}
%    \end{macrocode}
% Note that there is no real reason to use |\SetFirstPage| since all
% the page numbering is done automatically in the book style. 
%    \begin{macrocode}
\newcommand{\SetFirstPage}[1]{\setcounter{page}{#1}}
\newcommand{\SetYear}[1]{%
  \renewcommand{\@year}{#1}
  \renewcommand{\ADS@year}{#1}
  \renewcommand{\@anyo}{\rmaa@VolumenStyle{#1}}
  }
\newcommand{\ReceivedDate}[1]{\renewcommand{\@received}{#1}}
\newcommand{\AcceptedDate}[1]{\renewcommand{\@accepted}{#1}}
%    \end{macrocode}
%
% \begin{macro}{\DOI}
% User/editor command to enter the DOI, which will be typeset as a link.
% \changes{4.0}{2019/11/02}{New user command \cs{DOI} to typeset
% DOI in the header of first page and add it to ADS info.}
%    \begin{macrocode}
\newcommand\DOI[1]{%
  \def\@DOI@url{\url{https://doi.org/10.22201/ia.01851101p.#1}}%
  \def\ADS@doi{https://doi.org/10.22201/ia.01851101p.#1}%
}%
%    \end{macrocode}
% \end{macro}
%
% \begin{macro}{\ORCIDiD}
% \changes{4.0}{2019/11/12}{Add mechanism for giving ORCID iD links}
% \begin{macro}{\ORCIDiDfull}
% \begin{macro}{\@orcid@icon}
% \begin{macro}{\@orcid@iconbig}
% Commands to add icon links to ORCID iDs of authors. For example
% \begin{quote}
% |\ORCIDiD|\marg{0000-0000-0000-0000}
% \end{quote}
% to be used after each author name in the |\author|
% command and in the |\fulladdresses|.  Alternatively,
% |\OORCIDiDfull| will typeset the icon plus the full text of the
% link, but we don't use that at present.  The icon image is stored in |ORCIDiD_iconvector.pdf|, which needs to be on the input path.  In principle, we could make EPS versions if necessary.
%    \begin{macrocode}
\newcommand\@orcid@icon{\includegraphics[height=1.8ex]{ORCIDiD_iconvector}}
\newcommand\@orcid@iconbig{\includegraphics[height=2ex]{ORCIDiD_iconvector}}
\newcommand\ORCIDiD[1]{\href{https://orcid.org/#1}{\,\@orcid@icon}}
\newcommand\ORCIDiDfull[1]{%
  \href{https://orcid.org/#1}{\@orcid@iconbig}\,\url{https://orcid.org/#1}}
%    \end{macrocode}
% \end{macro}
% \end{macro}
% \end{macro}
% \end{macro}
%
% \begin{macro}{\@volume}
% \begin{macro}{\@volumen}
% \begin{macro}{\@artpages}
% \begin{macro}{\@year}
% \begin{macro}{\@anyo}
% \begin{macro}{\@received}
% \begin{macro}{\@accepted}
% \begin{macro}{\@MSnum}
% Default values for the containers used by above commands, plus a
% few more useful ones. 
%    \begin{macrocode}
\newcommand{\@volume}{00}
\newcommand{\@volumen}{\rmaa@VolumenStyle{00}}
\newcommand{\@artpages}{%
  \oldpageref{\articlename:RMAAFirstPage}--\oldpageref{\articlename:RMAALastPage}}
\newcommand{\@year}{\the\year}
\newcommand{\@anyo}{\rmaa@VolumenStyle{2007}}
\newcommand{\@journalname}{\rmaa@JournalNameStyle{%
    \if@proceedings\RMAASClong\else\RMAAlong\fi}}
\newcommand{\@received}{}
\newcommand{\@accepted}{}
\newcommand{\@MSnum}{0000}
\newcommand{\@DOI@prefix}{https://doi.org/10.22201/ia.01851101p}
\def\@DOI@url{\url{https://doi.org/10.22201/ia.01851101p.2099.00.00.00}}
%    \end{macrocode}
% \end{macro}
% \end{macro}
% \end{macro}
% \end{macro}
% \end{macro}
% \end{macro}
% \end{macro}
% \end{macro}
% \end{macro}
% \end{macro}
% \end{macro}
% \end{macro}
% \end{macro}
% \end{macro}
% 
% The following commands are most useful in the \texttt{book}
% style. However, the |\SetConfTitle| command can now be used in
% individual articles to set the header for the \texttt{preprint}
% style. 
%    \begin{macrocode}
\newcommand{\@reunionplace}{Lugar de la reuni\'on}
\newcommand{\@reuniondate}{Fecha de la reuni\'on}
\newcommand{\@conftitle}{Conference Title}
\newcommand{\SetReunion}[2]{\renewcommand{\@reunionplace}{#1}\renewcommand{\@reuniondate}{#2}}
\newcommand{\SetConfTitle}[1]{\renewcommand{\@conftitle}{#1}}
\newcommand{\@editors}{Juan Ram\'on de la Fuente and ``El Mosh''}
\newcommand{\SetEditors}[1]{\renewcommand{\@editors}{#1}}
\newcommand{\@editores}{Juan Ram\'on de la Fuente y ``El Mosh''}
\newcommand{\SetEditores}[1]{\renewcommand{\@editores}{#1}}
\newcommand{\@mes}{mes de la publicaci\'on}
\newcommand{\SetMes}[1]{\renewcommand{\@mes}{\rmaa@VolumenStyle{#1}}}
%    \end{macrocode}

% \subsection{Notes for Author Proofs and Communication between Editors}
% \label{sec:notes}
% We introduce a command \cs{Note} that should be used to describe
% any change made to the manuscript. These are a bit tricky to get
% right. In general we use \cs{marginpar} but this only works in outer
% paragraph mode, so we just bung the box where we are if we are in a
% figure caption or something. If we are in the preamble, then we just
% use a \marginpar as part of the \maketitle command.
%    \begin{macrocode}
\newif\if@arenotes
\@arenotesfalse
\newcounter{notecounter}
\setcounter{notecounter}{0}
\newcommand\NotesList{}
\newcommand\Note[1]{%
  \if@proof{%
    \@arenotestrue%
    \stepcounter{notecounter}%
    \if@inpreamble\else{%
      \NoteMark{\arabic{notecounter}}%
    }\fi
    \protected@xdef\NotesList{\NotesList 
      \protect\item #1}%
  }\fi%
}
\newcommand\Changed[2]{Changed ``#1'' to ``#2''.}
\newcommand\NoteMark[1]{%
  \ifinner
  \raisebox{1ex}[0pt][0pt]{\makebox[0pt][l]{\framebox{\footnotesize\bfseries
        N#1}}}%
  \else
  \marginpar{\framebox{\footnotesize\bfseries N#1}}%
  \fi
}
\newcommand\RefNote[1]{\if@proof\NoteMark{\ref{#1}}\relax\fi}
\newcommand\NoteLabel{\protect\label}
\newcommand\OutputNotes{%
  \bgroup
  \onecolumn
  \section*{EDITORIAL NOTES}
  \begin{enumerate}
    \NotesList
  \end{enumerate}
  \egroup
}
\DeclareRobustCommand\cs[1]{\texttt{\char`\\#1}}

%    \end{macrocode}
% \subsection{Bibliographic Records for ADS}
% \label{sec:ADS}
% \changes{2.2beta11}{2001/11/27}{Just about finished all the ADS
% stuff now. Except for the references, which I have given up on for
% the time being.}
% \begin{macro}{\ADS@BibCode}
% \begin{macro}{\ADS@year}
% \begin{macro}{\ADS@journal}
% \begin{macro}{\ADS@volume}
% \begin{macro}{\ADS@firstpage}
% \begin{macro}{\ADS@authorinitial}
% We use the |\ADS@BibCode| command to hold the bibliographic code for
% ADS. It is 
% used in the |%R| 
% field of the bibliographic record for each article. 
% The format is 
% \begin{quote}
% \texttt{YYYYJJJJJVVVVMPPPPA}  
% \end{quote}
% where \texttt{YYYY} is the year, \texttt{JJJJJ} is the journal code,
% which is \texttt{RMxAA} or \texttt{RMxAC}, \texttt{VVVV} is the
% volume number (right-justified), \texttt{M} is not used by us,
% \texttt{PPPP} is the first page number (right-justified), and
% \texttt{A} is the initial of the first author's surname. The volume
% and page fields are padded with periods (`\texttt{.}'), which causes
% us some fun (this is done with the macro |\PadToFourWithDots| described
% below) . 
% 
% All these individual bits are held in the commands |\ADS@year|,
% |\ADS@journal|, |\ADS@volume|, |\ADS@firstpage|,
% |\ADS@authorinitial|. 
%    \begin{macrocode}
\newcommand{\SetADSBibCode}[1]{\renewcommand{\ADS@BibCode}{#1}}
\newcommand{\ADS@BibCode}
{\ADS@year\ADS@journal\ADS@volume.\PadToFourWithDots{\ADS@firstpage}\ADS@authorinitial}
\newcommand{\ADS@year}{2001} 
\newcommand{\ADS@journal}{\if@proceedings RMxAC\else RMxAA\fi}
\newcommand{\ADS@volume}{....} 
%    \end{macrocode}
% These two are just set to default values here, |\ADS@firstpage| is set
% in the |document| environment for included articles, in the
% |frontsection| environment, |\tableofcontents| command,  |index| and
% |abstracts| environments. Also by |\AtBeginDocument| for stand-alone
% articles. 
% 
% God knows how we are going to do |\ADS@authorinitial|!
%    \begin{macrocode}
\newcommand{\ADS@firstpage}{....}
\newcommand{\ADS@authorinitial}{A}
%    \end{macrocode}
% \end{macro}
% \end{macro}
% \end{macro}
% \end{macro}
% \end{macro}
% \end{macro}
% We also have a whole slew of other fields that can be set. These are
% all listed in Table~\ref{tab:ADS}. Most of them still have to be
% defined. 
% \begin{table}[!ht]\raggedleft
% \caption{ADS Bibliographic Records}\label{tab:ADS}
% \begin{tabular}{lp{8em}lp{8em}}\hline
% Code & Description & Macro Name & Comment \\ \hline
% |%R| & Bibliographic Code & |\ADS@BibCode| & required \\
% |%A| & Author list & |\ADS@authors| & required, generate from |\indexauthor| \\
% |%F| & Author affiliation & |\ADS@affil| & format: AA(\dots)
% AB(\dots) \dots \\
% |%J| & Journal Name & |\ADS@journalname| &  \\
% |%D| & Publication Date & |\ADS@fulldate| &  \\
% |%L| & Last Page & |\ADS@lastpage| & generate from label \\
% |%T| & Title & |\ADS@title| & required, generate from |\title| \\
% |%C| & Abstract Copyright & |\ADS@copyright| & set to IA-UNAM \\
% |%U| & URL for electronic document & |\ADS@url| & set to PDF file \\
% |%d| & DOI for article & |\ADS@doi| & set by |\DOI| \\
% |%K| & Keywords & |\ADS@keywords| & set equal to |\@keywords| \\
% |%A| & Abstract Text & |\ADS@abstract| & set equal to |\@abstract| \\
% |%R| & References & |\ADS@references| & generate from the
% bibliography  \\ \hline
% \end{tabular}
% \end{table}
% A lot of these macros are set by other commands, as mentioned in the
% Table, but some need to be set  explicitly here. 
% 
% \begin{macro}{\ADS@doi}
% \changes{4.0}{2019/11/08}{Save DOI for ADS}
% The user/editor command \cs{DOI} is used to set this. 
%    \begin{macrocode}
\newcommand{\ADS@doi}{https://doi.org/10.22201/ia.01851101p.0000.00.00.00}
%    \end{macrocode}
% \end{macro}
% 
% \begin{macro}{\ADS@journalname}
% This
% does not contain any accents at the moment so as not to cause
% problems---this maybe should be changed though. 
%    \begin{macrocode}
\newcommand{\ADS@journalname}{%
  \if@proceedings\@conftitle\space (Eds. \@editors) \fi%
  Revista Mexicana de Astronom\'\i{}a y Astrof\'\i{}sica%
  \if@proceedings\space (Serie de Conferencias)\fi%
  \space Vol. \@volume, pp. \ADS@firstpage--\ADS@lastpage\space (\@year)%
  \space (http://www.astroscu.unam.mx/\~{}rmaa/)
  }
%    \end{macrocode}
% \end{macro}
% \begin{macro}{\ADS@keywords}
% \begin{macro}{\ADS@abstract}
% \begin{macro}{\ADS@title}
% |\ADS@abstract| is straightforward, since we already have the text
% saved as |\@abstract|. We could do the same with |\ADS@keywords| but
% to avoid unwanted formatting we explicitly set it in |\keywords|
% (p.~\pageref{keywords}). |\ADS@title| is just set to |\@rawtitle|,
% which we were already using for another purpose. 
%    \begin{macrocode}
\let\ADS@title=\@rawtitle
\newcommand{\ADS@keywords}{No keywords specified}
\newcommand{\ADS@abstract}{No abstract specified}
%    \end{macrocode}
% \end{macro}\end{macro}\end{macro}
% \begin{macro}{\ADS@authors}
% |\ADS@authors| wants to be in ``Other, A. N.'' format, but with
% semi-colons between authors. The most natural way to do this seems to
% be with |\indexauthor| (p.~\pageref{indexauthor}). We append each
% author to this macro, which means we must explicitly clear it for
% each new article (done in |\article|). .   
%    \begin{macrocode}
\newcommand{\ADS@authors}{}
%    \end{macrocode}
% \end{macro}
% \begin{macro}{\ADS@fulldate}
% \begin{macro}{\ADS@lastpage}
% \begin{macro}{\ADS@copyright}
% \begin{macro}{\ADS@website}
% \begin{macro}{\ADS@url}
% \begin{macro}{\SetWebsite}
% \begin{macro}{\SetMonthNumeric}
% Other sundry data, together with user commands for bits that can be
% set by the Editors. The |\ADS@url| is constructed from the website
% (currently |http://www.astrosmo.unam.mx/rmaa/|,
% the journal descriptor (e.g., |RMxAC..10/|), followed by the
% path to the PDF file (e.g., |PDF/RMxAC..10_HEW1.pdf|). 
% 
%    \begin{macrocode}
\newcommand{\ADS@website}{http://www.astrosmo.unam.mx/rmaa/}
\newcommand{\ADS@url}
{\ADS@website\ADS@journal\ADS@volume\ADS@number/PDF/\ADS@journal\ADS@volume\ADS@number\_\articlename.pdf}
\newcommand{\ADS@number}{}
\newcommand{\ADS@fulldate}{00/\ADS@year} 
\newcommand{\ADS@lastpage}{00} 
\newcommand{\ADS@copyright}{(C) \ADS@year; Instituto de Astronom\'\i{}a,
  Universidad Nacional Aut\'onoma de M\'exico} 
\newcommand{\SetNumber}[1]{\renewcommand{\ADS@number}{-#1}}
\newcommand{\SetWebsite}[1]{\renewcommand{\ADS@website}{#1}}
\newcommand{\SetMonthNumeric}[1]{\renewcommand{\ADS@fulldate}{#1/\ADS@year}}
%    \end{macrocode}
% \end{macro}\end{macro}\end{macro}\end{macro}\end{macro}\end{macro}\end{macro}
% \begin{macro}{\ADS@affil}
% \begin{macro}{\ADS@references}
% These two are somewhat more problematical. Affiliation currently
% only works for the first author (see \dots) pending redesign of the
% user interface. References are now done by wrapping the
% |thebibliography| environment to open another input stream on the
% same file and search through for the |\bibitems|. 
%    \begin{macrocode}
\newcommand{\ADS@references}{}
\newcommand{\ADS@affil}{}
%    \end{macrocode}
% \end{macro}\end{macro}
% 
% Define two macros that expand to the strings
% |\begin{thebibliography}| and |\end{thebibliography}|. Note that the
% idea for this and much else is inspired by the code for the docstrip
% package. Also define |\StartBibX| in case authors use the more
% standard |\begin{thebibliography}{}| syntax. 
%    \begin{macrocode}
\edef\EndBib{\expandafter\@gobble\string\\end\string{thebibliography\string}}
\edef\StartBib{\expandafter\@gobble\string\\begin\string{thebibliography\string}}
\edef\StartBibX{\expandafter\@gobble\string\\begin\string{thebibliography\string}\string{\string}}
%    \end{macrocode}
% Define two new logical variables. |\ifAnotherline| is true if we
% want to read another line from the input stream. It is set to false
% when we encounter the |\end{thebibliography}|. |\ifInBib| is set to
% true when we frind the |\begin{thebibliography}| and tells us to
% start saving the following lines. 
%    \begin{macrocode}
\newif\ifAnotherLine
\newif\ifInBib
%    \end{macrocode}
% \begin{macro}{makeOther}
% We need a macro to change the catcodes of lots of things to other. 
%    \begin{macrocode}
\def\makeOther#1{\catcode`#1=12\relax}
%    \end{macrocode}
% \end{macro}
% \begin{macro}{\RereadBiblio}
% Define a macro to do all the fiddly stuff. This is used at the
% beginning of |thebibliography| environment (Sec~\ref{sec:biblio}).
%    \begin{macrocode}
\newread\tf@biblio
\newcommand{\RereadBiblio}{%
%    \end{macrocode}
% Open the current file again for reading
%    \begin{macrocode}
  \begingroup
  \immediate\openin\tf@biblio=\articlename
%    \end{macrocode}
% Everything in a group to be safe. Initialize logicals. Black magic
% to change
% catcodes and avoid blank lines being interpreted as |\par|. 
%    \begin{macrocode}
  \AnotherLinetrue\InBibfalse
  \makeOther\ \makeOther\\\makeOther\$%
  \makeOther\#\makeOther\^\makeOther\^^K%
  \makeOther\_\makeOther\^^A\makeOther\%%
  \makeOther\~\makeOther\{\makeOther\}\makeOther\&%
  \endlinechar-1\relax
%    \end{macrocode}
% Start looping through the file. Read each line to the macro
% |\Line|. If we have got to EOF or end of biblio then set
% |\AnotherLinefalse|.
%    \begin{macrocode}
  \loop
  \read\tf@biblio to \Line
  \ifx\Line\EndBib\AnotherLinefalse\InBibfalse\fi
  \ifeof\tf@biblio \AnotherLinefalse\fi
%    \end{macrocode}
% If we are in the biblio and the line is not blank then append it to
% the |\ADS@references| macro together with a newline. 
%    \begin{macrocode}
  \ifInBib
  \ifx\Line\empty\else\xdef\ADS@references{\ADS@references\Line^^J}\fi
  \fi
%    \end{macrocode}
% Check if the line we just saw was the start of biblio, then repeat
% the loop if necessary. 
%    \begin{macrocode}
  \ifx\Line\StartBib\InBibtrue\fi
  \ifx\Line\StartBibX\InBibtrue\fi
  \ifAnotherLine\repeat
  \immediate\closein\tf@biblio
  \endgroup
  }
%    \end{macrocode}
% \end{macro}
% 
% \begin{macro}{\PadToFourWithDots}
% Utility function for left-padding a number to 4 characters with
% periods. Used in 
% constructing |\ADS@BibCode|. This is a very simple algorithm and
% \emph{requires} that the argument actually be a number---if it is
% not then dire consequences will result.  
%    \begin{macrocode}
\newcommand{\PadToFourWithDots}[1]{%
  \ifnum #1 > 999 #1%
  \else\ifnum #1 > 99 .#1%
  \else\ifnum #1 > 9 ..#1%
  \else ...#1%
  \fi\fi\fi}
%%  \ifnum \count0 > 999 #1%
%%  \else\ifnum \count0 > 99 .#1%
%%  \else\ifnum \count0 > 9 ..#1%
%%  \else ...#1%
%%  \fi\fi\fi}
%    \end{macrocode}
% \end{macro}
%
% \begin{macro}{\AuthorADS}
% \begin{macro}{\ADS@affiliations}
% New mechanism for setting the affiliations for ADS. The author labels (AA, AB, \dots) are generated using two nested counters.  We use the |\robustify| command from \texttt{etoolbox} so that accented characters in the affiliation addresses do not blow up.  The list of accents may need to be expanded. 
%    \begin{macrocode}
\newcounter{AUTHOR}
\newcounter{author}[AUTHOR]
\stepcounter{AUTHOR}
\newcommand\ADS@affiliations{}
\robustify{\'}\robustify{\i}\robustify{\~}\robustify{\`}
\def\AuthorADS#1#2{%
  \begingroup
  \def\ORCIDiD##1{<ID system="ORCID">##1</ID>}
  \def\Email##1{<EMAIL>##1</EMAIL>}
  \def\Affil##1{##1}
  \stepcounter{author}
  \ifthenelse{\value{author} > 26}{\stepcounter{AUTHOR}\stepcounter{author}}{}
  \ifthenelse{\equal{\ADS@affiliations}{\@rmaaEmpty}}{%
    \xdef\ADS@affiliations{\Alph{AUTHOR}\Alph{author}(#2)}
  }{%
    \xdef\ADS@affiliations{\ADS@affiliations, \Alph{AUTHOR}\Alph{author}(#2)}
  }%
  \endgroup
}
%    \end{macrocode}
% \end{macro}
% \end{macro}
% %
% \begin{macro}{\ADS@dumpdata}
% \changes{4.0}{2019/11/13}{Write all author affiliations for ADS}
% \begin{macro}{\ADS@write}
% \begin{macro}{\ADS@openfile}
% |\ADS@dumpdata| is the macro that actually does the job of writing out all the
% data to an auxiliary file, using |\ADS@write| to write each
% field. The file is set up with |\ADS@openfile|, which currently is
% done in the TOC of a book.  The new multi-author affiliations are used if |\ADS@affiliations| has been set, otherwise it falls back to the single-author version |\ADS@affil|. 
%    \begin{macrocode}
\newcommand{\ADSfield}[2]{\relax}
\newcommand{\ADSvfield}[2]{\relax}
%%\renewcommand{\ADSfield}[12]{\%#1 #2}
\newcommand{\ADS@dumpdata}
{ 
  \ADS@write{\string\begin{ADSarticle}}
  \ADS@WRITE{\string\ADSfield{R}{\ADS@BibCode}}
  \ADS@write{\string\ADSfield{A}{\ADS@authors}}
  \ifthenelse{\equal{\ADS@affiliations}{\@rmaaEmpty}}{%
    {\ADS@sanitize\ADS@write{\string\ADSfield{F}{AA(\ADS@affil)}}}
  }{%
    {\ADS@sanitize\ADS@write{\string\ADSfield{F}{\ADS@affiliations}}}
  }
  \ADS@write{\string\ADSfield{J}{\ADS@journalname}}
  \ADS@write{\string\ADSfield{D}{\ADS@fulldate}}
  \ADS@write{\string\ADSfield{L}{\ADS@lastpage}}
  {\ADS@sanitize\ADS@write{\string\ADSfield{T}{\ADS@title}}}
  \ADS@write{\string\ADSfield{C}{\ADS@copyright}}
  \ADS@write{\string\ADSfield{d}{\ADS@doi}}
  \ADS@write{\string\ADSfield{U}{\ADS@url}}
  \ADS@write{\string\ADSfield{K}{\ADS@keywords}}
  {\ADS@sanitize\ADS@write{\string\ADSfield{B}{\ADS@abstract}}}
  \ADS@write{\string\ADSfield{Z}{\ADS@references}}
  \ADS@write{\string\end{ADSarticle}}
  }
%%\newcommand{\ADS@write}[1]{\if@book\addtocontents{ads}{#1}\fi}
%%\newcommand{\ADS@openfile}{\@starttoc{ads}}
\newcommand{\ADS@write}[1]{\set@display@protect\immediate\write\tf@ads{#1}}
\newcommand{\ADS@WRITE}[1]{\immediate\write\tf@ads{#1}}
\newcommand{\ADS@openfile}{%
  \newwrite\tf@ads%
  \immediate\openout\tf@ads\jobname.ads\relax
  \ADS@WRITE{\ADS@preamble}}
\newcommand{\ADS@closefile}{\ADS@WRITE{\ADS@postamble}}
\def\ADS@preamble{%
  %%  \string\documentclass{article}
  %%  \string\begin{document}
  }
\def\ADS@postamble{%
  %%  \string\end{document}
  \string%\string% Local Variables:^J 
  \string%\string% mode: LaTeX^J
  \string%\string% mode: Font-lock^J
  \string%\string% End:^J
  }
%    \end{macrocode}
% \end{macro}
% \end{macro}
% \end{macro}
% 
% \subsection{Automatic calculation of the page range}
% \label{sec:page-range}
% Another issue that has cost me a lot of grief. Still, it finally
% seems to be sorted now. All hail Jeff Goldberg. 
%
% Cunning trick to get the page numbers from |\pageref|.  Put a label
% on the first and last page of the document using |\AtBeginDocument|
% and |\AtEndDocument|. .  This way won't work for the book option.
% In that case, we do it by defining |\@endofarticlehook| and
% |\@begofarticlehook| as parallels to |\AtEndDocument| and
% |\AtBeginDocument|.  Note that we have to defer adding
% |\@begofarticlehook| to |\AtBeginDocument| to make sure it comes
% after the redefinitions of |\ref| and |\label|.
%    \begin{macrocode}
\if@book\AtEndDocument{\ADS@closefile}\else
 %\AtBeginDocument{\@begofarticlehook}
\AtEndDocument{\@endofarticlehook}
\fi
%    \end{macrocode}
% \begin{macro}{\@begofarticlehook}
% Since we are getting more and more things that need to be done at
% the beginning of every article, frontsection, etc we define a sister
% hook to 
% |\@endofarticlehook|. This hook does the following: 
% \begin{enumerate}
% \item Labels the first page
% \item Sets |\ADS@firstpage|
% \item runs and then clears |\indexauthor@deferlist|
% \item resets the keyword counter
% \end{enumerate}
% It is executed in the following circumstances: 
% \begin{enumerate}
% \item At |\begin{document}| for standalone files
% \item At |\begin{document}| for each included file
% \item At beginning of front sections 
% \item At beginning of TOC
% \item At beginning of Author Index
% \item At beginning of abstracts
% \end{enumerate}
% For some reason, this is not always within the right group, so we
% use |\xdef| just to be sure. 
%    \begin{macrocode}
\newcommand{\@begofarticlehook}{%
  \if@frontmatter\frontlabel{RMAAFirstPage}\else\label{RMAAFirstPage}\fi%
  \xdef\ADS@firstpage{\@arabic{\c@page}}%
  \if@book\else\makeindex\fi % needed so default index is defined in stand-alone articles
%%  \indexauthor@deferlist
%%  \xdef\indexauthor@deferlist{}
  \@inpreamblefalse
  \setcounter{kwdcounter}{0}
  }
%    \end{macrocode}
% \end{macro}
% \begin{macro}{\@endofarticlehook}
% This is a hook to be run at the end of each article
% currently, it just labels the last page and writes out all the ADS
% data to an auxiliary file.  
% It should not produce any printed output. 
% It is modelled on the |\AtEndDocument| command, to which it reduces
% in the case of a single article. 
%    \begin{macrocode}
\newcommand{\@endofarticlehook}{%
  \if@debug{\ClassWarning{rmaa}{***Doing Hook***}}\fi
  \if@debug{\ClassWarning{rmaa}{ADSabstract is: \ADS@abstract}}\fi
  \clearpage\rmaa@lastpagelabel\if@book\ADS@dumpdata\fi
  \ifx\NotesList\@rmaaEmpty\else\OutputNotes\fi
  \if@debug{\ClassWarning{rmaa}{Notes list is: \NotesList}}\fi
}
%    \end{macrocode}
% \begin{macro}{\rmaa@lastpagelabel}
% New version of last page label command copied from Jeff Goldberg's
% lastpage.sty package - we can't just use |\label| here in case the
% output routine for the last page has already been called. Instead,
% we first explicitly clear the page, step back the page counter,
% explicitly write to the .aux file, then restep the page counter
% again.
%
% But this will conflict with the hyperref package, which does labels differently.  Compensate for that by using the five-argument form.
% 
% Also set the |\ADS@LastPage| macro to the appropriate number. 
%    \begin{macrocode}
\newcommand{\rmaa@lastpagelabel}{%
  \addtocounter{page}{-1}%
  \immediate\write\@auxout{\string
    \newlabel{\articlename:RMAALastPage}{{}{\thepage}{}{}{}}}%
  \edef\ADS@lastpage{\thepage}%
  \addtocounter{page}{1}}
%    \end{macrocode}
% \end{macro}
% \end{macro}

% \subsection{Deal with explicit change in number of columns}
% \label{sec:columns}
% Sometimes, when formatting a book, we need to switch between
% one-column articles and two-column articles. For this to work, we
% need to wrap the standard \LaTeX{} commands to set the |\@rmtwocol|
% flag appropriately and also reset the package default values of all
% scalable lengths. .  
%    \begin{macrocode}
\let\latexonecolumn=\onecolumn
\let\latextwocolumn=\twocolumn
\renewcommand{\onecolumn}{\@rmtwocolfalse\@package@default@lengths\latexonecolumn}
\renewcommand{\twocolumn}{\@rmtwocoltrue\@package@default@lengths\latextwocolumn}
%    \end{macrocode}
% 
% \subsection{The actual formatting of the title}
% \label{sec:title-format}

% Note, we have to redefine the |\maketitle| command to use the
% numbered footnotes. This means we needn't use the atbeginend stuff. 
%    \begin{macrocode}
\renewcommand{\maketitle}{\par
  \begingroup
%    \end{macrocode}
% Don't switch to marked footnotes. 
% This is consistent with existing style but I'm not sure I like it
%    \begin{macrocode}
  %%  \renewcommand\thefootnote{\fnsymbol{footnote}}%
%%  \def\@makefnmark{\hbox to\z@{$\m@th^{\@thefnmark}$\hss}}%
  \def\@makefnmark{\textsuperscript{\@thefnmark}\hspace{-0.2em}}%
  \long\def\@makefntext##1{\parindent 1em\noindent
    \hbox to1.8em{\hss$\m@th^{\@thefnmark}$}##1}%
  \if@rmtwocol
  { \if@debug{\ClassWarning{rmaa}{TWO COLUMN}}\fi
    \twocolumn[\@maketitle]%
    }
  \else
  { \if@debug{\ClassWarning{rmaa}{ONE COLUMN}}\fi
    \newpage
    \global\@topnum\z@   % Prevents figures from going at top of page.
    \if@manuscript \@MSmaketitle \else \@maketitle \fi
    }
  \fi
  \if@manuscript\else\thispagestyle{rmaafirstpage}\@thanks\fi
  \ifx\NotesList\@rmaaEmpty\else{%
    \NoteMark{1--\arabic{notecounter}}\vspace*{-2\baselineskip}%
  }\fi%
  \endgroup
%    \end{macrocode}
% Don't reset footnote counter (not sure whether this is good thing
% or not).
%    \begin{macrocode}
  %%  \setcounter{footnote}{0}%
  \let\thanks\relax\let\affil\relax\let\altaffil\relax
  \let\maketitle\relax\let\@maketitle\relax
  \gdef\@thanks{}\gdef\@author{}\gdef\@title{}}
%    \end{macrocode}

% \begin{macro}{\@kwhwidth}
% \begin{macro}{\@kwtwidth}
% Redefine \cs{@maketitle} to include the abstract, resumen and
% keywords. Note that in the manuscripty style we use
% \cs{@MSmaketitle} instead.  
% \changes{3.16}{2003/04/07}{
%  Added extra vertical space between authors and accepted date.}
% \changes{3.26}{2011/08/15}{Add extra label to review articles}
%    \begin{macrocode}
\newlength{\@kwhwidth}\newlength{\@kwtwidth}
\newcommand\reviewname{Review}
\renewcommand{\@maketitle}{%
  \newpage
  \null
  \if@reviewarticle
  \smallskip
  {\large\bfseries\itshape\reviewname}%
  \bigskip
  \else
  \vskip\@before@title@skip%
  \fi
  \begin{center}%
    {\@title \par}%
    \vskip\@before@author@skip%
    \parbox{\rmaa@AuthorWidthFrac\textwidth}{\centering
      \lineskip\@between@author@skip%
      \@author
      \par}
    %% this should be replaced with `Submitted to...' for preprints
    \if@printdates
    \vskip\@before@dates@skip%
    \if@preprint
    \rmaa@ReceivedStyle{Submitted to RMxAA \@received}
    \else
    \rmaa@ReceivedStyle{Received \@received; accepted \@accepted}
    \fi
    \fi
    \if@manuscript
    \rmaa@ReceivedStyle{Submitted to RMxAA \@received}
    \fi
    \if@guide
    \rmaa@ReceivedStyle{Version \rmaa@version, \rmaa@date}
    \fi
    \vskip\@before@abstract@skip%
    \if@abstractpresent
    \parbox{\rmaa@AbstractWidthFrac\textwidth}{%
      \setlength{\parindent}{0.9cm}%
      \centerline{\rmaa@AbstractHeaderStyle{Abstract}}\par%
      \vskip\@after@abstracthead@skip%
      \if@newstyle\noindent\fi\@abstract
      \vskip\@between@abstract@skip%
      \centerline{\rmaa@AbstractHeaderStyle{Resumen}}\par%
      \vskip\@after@abstracthead@skip%
      \if@newstyle\noindent\fi\@resumen
      \if@setresumen\else\@abstract\fi % append english abstract if
                                       % spanish one not set
      \if@guide\else            %do not use keywords for guide
      \vskip\@before@kwd@skip%
      % Format the Key Words with hanging indentation:
      % They are put in a \parbox with width equal to the line width
      % minus the width of the header. 
      { % Widths of keyword
        \settowidth{\@kwhwidth}{\@keywordheader~}   % header and
        \setlength{\@kwtwidth}{\linewidth}          % keyword text
        \addtolength{\@kwtwidth}{-\@kwhwidth}       %
        \noindent\@keywordheader~\parbox[t]{\@kwtwidth}{%
          \rmaa@KeywordStyle{\boldmath\uppercase{\@keywords}}}}%
      \fi
      }
      \fi
  \end{center}%
  \vspace*{\afterkwdskip}
  \par
  }
\newlength{\afterkwdskip}\setlength{\afterkwdskip}{\z@}
%    \end{macrocode}
% \end{macro}
% \end{macro}
%
%
% \begin{macro}{\reviewarticle}
% \begin{macro}{\reviewtoc}
% \begin{macro}{\reviewlof}
%   User command to specify a review article. Must be used in
%   preamble. Main effect is to place the text ``Review'' before the
%   article title.  Also add command for setting a table of contents
%   and a list of figures at the beginning of the article.
%    \begin{macrocode}
\newcommand{\reviewarticle}{\@reviewarticletrue}
\newcommand\reviewtoc{%
  \section*{CONTENTS}%
  \@starttoc{rtc}%
}
\newcommand\reviewlof{%
  \section*{LIST OF FIGURES}%
  \@starttoc{lof}%
}
%    \end{macrocode}
% \end{macro}
% \end{macro}
% \end{macro}
% \begin{macro}{\l@section}
%    \begin{macrocode}
\renewcommand*\l@section[2]{%
  \ifnum \c@tocdepth >\z@
    \addpenalty\@secpenalty
    \addvspace{0.8em \@plus\p@}%
    \setlength\@tempdima{1.5em}%
    \begingroup
      \parindent \z@ \rightskip \@pnumwidth
      \parfillskip -\@pnumwidth
      \leavevmode %\bfseries
      \advance\leftskip\@tempdima
      \hskip -\leftskip
      #1\nobreak\hfil \nobreak\hb@xt@\@pnumwidth{\hss #2}\par
    \endgroup
  \fi}
%    \end{macrocode}
% \end{macro}
%
% \begin{macro}{\@MSmaketitle}
% A special version of |\@maketitle| for the \texttt{manuscript}
% option. This can deal with a multi-page abstract/resumen. Not very
% well, mind. 
%    \begin{macrocode}
\newcommand{\@MSmaketitle}{%
  \if@debug\ClassInfo{rmaa}{Using special MS maketitle command}\fi
  \newpage
  {
    \thispagestyle{rmaafirstpage}
%     \setlength{\parindent}{0pt}
    \null\vskip 2ex
    \begin{center}
      \@title\par
      \@author\par
      \rmaa@ReceivedStyle{%
        \ifx\@received\@rmaaEmpty
          Draft version: \@date
        \else
          Received: \@received
          \ifx\@accepted\@rmaaEmpty
          \else
            ; Accepted: \@accepted
          \fi
        \fi}\par
    \end{center}
    \@thanks
    \centerline{\rmaa@AbstractHeaderStyle{Resumen}}
    \noindent\@resumen\par\vskip 2ex
    \centerline{\rmaa@AbstractHeaderStyle{Abstract}}
    \noindent\@abstract\par\vskip 2ex
    \noindent\@keywordheader~\rmaa@KeywordStyle{\@keywords}\par
%     \@keywordheader~\rmaa@KeywordStyle{\uppercase{\@keywords}}\par
  }%
  }
%    \end{macrocode}
% \end{macro}

% \subsection{The sectioning commands}
% \label{sec:section-commands}

% \begin{macro}{\section}
% \begin{macro}{\subsection}
% \begin{macro}{\subsubsection}
% \begin{macro}{\paragraph}
% \begin{macro}{\subparagraph}
% Redefine all the standard commands to use these lengths. 
%    \begin{macrocode}
\renewcommand\section{\@startsection {section}{1}{\z@}%
                                   {\abovesectionskip}%
                                   {\belowsectionskip}%
                                   %% how do we get upper case here
                                   {\centering\reset@font\normalsize}}
\renewcommand\subsection{\@startsection{subsection}{2}{\z@}%
                                     {\abovesectionskip}%
                                     {\belowsectionskip}%
                                     {\centering\reset@font\normalsize\itshape}}
\renewcommand\subsubsection{\@startsection{subsubsection}{3}{\z@}%
                                     {\abovesectionskip}%
                                     {\belowsectionskip}%
                                     {\centering\reset@font\normalsize\itshape}}
\renewcommand\paragraph{\@startsection{paragraph}{4}{\z@}%
                                    {\aboveparaskip}%
                                    {\belowparaskip}%
                                    {\reset@font\normalsize\bfseries}}
\renewcommand\subparagraph{\@startsection{subparagraph}{5}{\parindent}%
                                      {\aboveparaskip}%
                                      {\belowparaskip}%
                                      {\reset@font\normalsize\bfseries}}
%    \end{macrocode}
% \end{macro}
% \end{macro}
% \end{macro}
% \end{macro}
% \end{macro}

% \begin{macro}{secnumdepth}
% Make sure that we number the subsubsections---this is done by
% default in the |article| class but not the |book| class, which is
% why we must set it explicitly here.
%    \begin{macrocode}
\setcounter{secnumdepth}{3}
%    \end{macrocode}
% \end{macro}
% \changes{2.2beta10}{2001/11/07}{Fixed \cs{secnumdepth} for \cs{book}
% style.}

% \begin{macro}{\thesubsection}
% \begin{macro}{\thesubsubsection}
% For subsection and subsubsection, we want the numbers to be roman
% but the title to be italic. This seems the easiest way to go about it.
%    \begin{macrocode}
\renewcommand{\thesubsection}{{\upshape\thesection.\@arabic\c@subsection}}
\renewcommand{\thesubsubsection}{\thesubsection{\upshape.\@arabic\c@subsubsection}}
%    \end{macrocode}
% \end{macro}
% \end{macro}

% \begin{macro}{\@seccntformat}
% Add ``.'' to the number here so we can remove it from |\thesection|,
% etc. Reduce gap between number and title WJH 27 Jan 1995
%    \begin{macrocode}
\def\@seccntformat#1{\csname the#1\endcsname .\hskip 0.4em}
%    \end{macrocode}
% \end{macro}

% \begin{macro}{\@sect}
% I didn't really want to mess with the |\@sect| command, but it was
% the only way I could manage to force upper case. Other than the
% commented parts, it is identical to the standard version. 
%    \begin{macrocode}
\def\@sect#1#2#3#4#5#6[#7]#8{
  \ifnum #2>\c@secnumdepth
     \let\@svsec\@empty
  \else
     \refstepcounter{#1}%
     \let\@@protect\protect
     \def\protect{\noexpand\protect\noexpand}%
     \edef\@svsec{\@seccntformat{#1}}%
     \let\protect\@@protect
  \fi
%    \end{macrocode}
% This is the bit I changed (27 Jan 1995). Use upper case for
% |\section|, use mixed case for the rest. 
%    \begin{macrocode}
  \ifnum #2=1                   
     \def\sec@upcase##1{\uppercase{##1}}%
     \def\sec@shape{\upshape}%
  \else                         
     \def\sec@upcase##1{\relax{##1}}%
     \def\sec@shape{\itshape}%
  \fi
  \@tempskipa #5\relax
  \ifdim \@tempskipa>\z@
     \begingroup #6\relax
          \@hangfrom{\hskip #3\relax\@svsec}%
                    {\interlinepenalty \@M \sec@upcase{#8}\par}%
     \endgroup
     \csname #1mark\endcsname{#7}%
%    \end{macrocode}
% Only write TOC entries if doing the Guide style. 
%    \begin{macrocode}
     \if@guide 
        \addcontentsline{toc}{#1}{%
          \ifnum #2>\c@secnumdepth \else
             \protect\numberline{\csname the#1\endcsname}%
          \fi
          #7}%
     \fi
%    \end{macrocode}
% Also allow per-article TOC for review papers
%    \begin{macrocode}
     \if@reviewarticle
        \ifnum #2>2 \else
           \addcontentsline{rtc}{#1}{%
             \protect\numberline{\csname the#1\endcsname}%
             \sec@upcase{\sec@shape #7}}%
        \fi
     \fi
  \else
     \def\@svsechd{#6\hskip #3\relax  %% \relax added 2 May 90
       \@svsec #8\csname #1mark\endcsname{#7}%
       \if@guide
          \addcontentsline{toc}{#1}{%
            \ifnum #2>\c@secnumdepth \else
               \protect\numberline{\csname the#1\endcsname}%
            \fi
            #7}%
       \fi
       }%
  \fi
  \@xsect{#5}}
%    \end{macrocode}
% \end{macro}


% \subsection{Appendices}
% \label{sec:appendices}
% Make sure that appendices in individual articles work properly. We
% copy the definition from the article class (the book one doesn't do
% what we want).  
% \begin{environment}{appendices}
%   This uses an environment |appendices| inside of which each
%   |\section| will become a separate appendix.  Note that the
%   appendices must come before the references (otherwise the full
%   addresses mechanism will not work properly). We make the equation
%   numbering be subordinate to the appendix (equations A12, B6,
%   etc). Finally, we have a medskip followed by APPENDICES to
%   introduce them. \changes{3.17}{2002/04/08}{ Added appendices environment
%   that does the right thing.}
%    \begin{macrocode}
\newenvironment{appendices}{\par
  \begingroup
  \setcounter{section}{0}%
  \setcounter{subsection}{0}%
  \renewcommand\thesection{\@Alph\c@section}
  \renewcommand\theequation{\thesection\arabic{equation}}
  \medskip
  \section*{APPENDICES}}
{\endgroup}
%    \end{macrocode}
% \end{environment}
% \begin{environment}{appendix}
%   If there is only one appendix, then use |appendix| environment
%   instead.  \changes{4.0}{2019/11/12}{ Added
%   appendix environment for case where there is only one section.}
%    \begin{macrocode}
\renewenvironment{appendix}{\par
  \begingroup
  \setcounter{section}{0}%
  \setcounter{subsection}{0}%
  \renewcommand\thesection{\@Alph\c@section}
  \renewcommand\theequation{\thesection\arabic{equation}}
  \medskip
  \section*{APPENDIX}}
{\endgroup}
%    \end{macrocode}
% \end{environment}

% \subsection{Floating thingies}
% \label{sec:floats}   

%% 
%    \begin{macrocode}
\renewcommand{\thetable}{\@arabic\c@table}
\renewcommand{\thefigure}{\@arabic\c@figure}
\renewcommand{\figurename}{Fig.}
\renewcommand{\tablename}{TABLE}
%\renewcommand{\fnum@table}{\MakeUpperCase{\tablename}~\thetable}
%    \end{macrocode}
% For some reason, I have to write `TABLE' explicitly here
% If I use |\tablename|, it gets converted to mixed case---God knows
% why. 
% |\figurename| is also being ignored. 
%    \begin{macrocode}
\renewcommand{\fnum@table}{TABLE~\thetable}
\renewcommand{\fnum@figure}{Fig.~\thefigure}
%    \end{macrocode}

% Do captions in 9 pt face in the main journal style and in new-style
% anything. 
%    \begin{macrocode}
\long\def\@makefigcap#1#2{%
  \if@proceedings\else\small\fi
  \if@newstyle\small\fi 
  \vskip\abovecaptionskip
  \sbox\@tempboxa{#1. #2}%
  \ifdim \wd\@tempboxa >\hsize
    #1. #2 \par
  \else
    \global \@minipagefalse
    \hb@xt@\hsize{\hfil\box\@tempboxa\hfil}%
  \fi
  \vskip\belowcaptionskip}
%    \end{macrocode}

% Make sure table caption is upper case. 
%    \begin{macrocode}
\long\def\@maketabcap#1#2{%
  \hbox to\hsize{\hfil #1 \hfil}
  \vskip 6pt 
  \centering\MakeUpperCase{#2}\par
  \vskip 6pt}
%    \end{macrocode}

%    \begin{macrocode}
\newlength{\@rottabcapwidth}
\long\def\@makerottabcap#1#2{%
  \setbox\@tempboxa\hbox{\MakeUpperCase{#2}}%
  \ifdim\wd\@tempboxa >.8\vsize
     \setlength\@rottabcapwidth{.8\vsize}
  \else
     \setlength\@rottabcapwidth{\wd\@tempboxa}
  \fi
  \rotatebox{90}{%
    \begin{minipage}{\vsize}\centering
      \begin{minipage}{\@rottabcapwidth}\centering
        #1\\ \vskip 6pt \MakeUpperCase{#2}
      \end{minipage}%
    \end{minipage}
    }%
  \hskip 6pt
  }

\newenvironment{Sideways}
    {\begin{sideways}\begin{minipage}{\vsize}\centering}
    {\end{minipage}\end{sideways}}

\newenvironment{Sidewaystable}
{%
  \let\oldcaption=\caption%
  \def\caption{
    \def\@captype{table}\renewcommand{\@makecaption}{\@maketabcap}
    \oldcaption}
  \begin{sidewaystable}\centering}
  {\end{sidewaystable}\let\caption=\oldcaption}
\newenvironment{Sidewaystable*}
{%
  \let\oldcaption=\caption%
  \def\caption{
    \def\@captype{table}\renewcommand{\@makecaption}{\@maketabcap}
    \oldcaption}
  \begin{sidewaystable*}\centering}
  {\end{sidewaystable*}\let\caption=\oldcaption}
%    \end{macrocode}

% Not sure why we need to stick these lines in. 
%    \begin{macrocode}
\renewenvironment{figure}
                 {\renewcommand{\@makecaption}{\@makefigcap}\@float{figure}}
                 {\end@float}
\renewenvironment{figure*}
                 {\renewcommand{\@makecaption}{\@makefigcap}\@dblfloat{figure}}
                 {\end@dblfloat}
\renewenvironment{table}
                 {\renewcommand{\@makecaption}{\@maketabcap}\@float{table}}
                 {\end@float}
\renewenvironment{table*}
                 {\renewcommand{\@makecaption}{\@maketabcap}\@dblfloat{table}}
                 {\end@dblfloat}
%    \end{macrocode}

% Redefine |\@caption| so that we don't bother with the list of
% figures. This saves us having to worry about fragile commands in the
% captions. 
%    \begin{macrocode}
\long\def\@caption#1[#2]#3{%
  \par
%%  \addcontentsline{\csname ext@#1\endcsname}{#1}%
%%    {\protect\numberline{\csname the#1\endcsname}{\ignorespaces #2}}%
  \begingroup
    \@parboxrestore
    \if@minipage
      \@setminipage
    \fi
    \normalsize
    \@makecaption{\csname fnum@#1\endcsname}{\ignorespaces #3}\par
  \endgroup}
%    \end{macrocode}

% Make tables not quite so cramped. Allow |\hline\hline| to give a
% poor man's thick rule. 
%    \begin{macrocode}
\renewcommand{\arraystretch}{1.15}
\setlength{\doublerulesep}{0.4pt}
%    \end{macrocode}

% \begin{macro}{\tablecols}
% \begin{macro}{\@tablecols}
% \begin{macro}{\tabcaptioncontinued}
% |\tablecols| is a user command for setting the number of columns. It
% is used by commands that try and span the entire table (I think
% there must be a better way of acheiving
% this). |\tabcaptioncontinued| is 
% a user command for setting the caption of a continued table on the
% next page. 
%    \begin{macrocode}
\newcommand{\tablecols}[1]{\renewcommand{\@tablecols}{#1}}
\newcommand{\@tablecols}{1}
\newcommand{\tabcaptioncontinued}{%
  \multicolumn{\@tablecols}{c}{\addtocounter{table}{-1}TABLE~\thetable\ (CONTINUED)} \\ \\
  }
%    \end{macrocode}
% \end{macro}
% \end{macro}
% \end{macro}

% \subsubsection{Footnotes in tables}
% \label{sec:table-footnotes}
% \begin{macro}{\tablenotemark}
% \begin{macro}{\tablenote}
% \begin{macro}{\tablenoteskip}
% Old user commands |\tablenotemark{MARK}|,
% |\tablenote{MARK}{TEXT}|---now deprecated. 
%    \begin{macrocode}
\newcommand\tablenotemark[1]{$^\mathrm{#1}$} 
\newcommand\tablenote[2]{\if@proceedings\else\small\fi ${}^\mathrm{#1}$ #2} 
\newcommand\tablenoteskip{\\[0.5\baselineskip]} 
%    \end{macrocode}
% \end{macro}
% \end{macro}
% \end{macro}

% \begin{macro}{\tabnotemark}
% \begin{macro}{\tabnotetext}
% \begin{macro}{\tabnotetextTWO}
% Alternative way of doing footnotes in tables. 
%    \begin{macrocode}
\newcommand{\tabnotemark}[1]{\textsuperscript{\lowercase{#1}}}
%\newcommand{\tabnotetextTWO}[4]{\tabnotetext{#1}{#2}\tabnotetext{#3}{#4}}
\newcommand{\tabnotetextTWO}[4]{%
  \multicolumn{\@tablecols}{l}{%
    \parbox{0.5\tabnotewidth}{\small\textsuperscript{#1}#2}%
    \parbox{0.5\tabnotewidth}{\small\textsuperscript{#3}#4}%
    }%
  \\}
\newcommand{\tabnotetext}[2]{%
  \multicolumn{\@tablecols}{@{}p{\tabnotewidth}@{}}{\small\textsuperscript{#1}#2}\\}
\newlength{\tabnotewidth}
%    \end{macrocode}
% \end{macro}
% \end{macro}
% \end{macro}


% \begin{macro}{changemargin}
%    \begin{macrocode}
% Command to allow margins to be changed on the fly. Useful to allow
% tables and Figures to protrude into the margins. 
%
% \changes{3.24}{2007/09/14}{Added \texttt{changemargin} environment
% (copied from TeX FAQ).}
\newenvironment{changemargin}[2]{% 
  % From TeX FAQ http://www.tex.ac.uk/cgi-bin/texfaq2html?label=chngmargonfly
  \begin{list}{}{%
    \setlength{\topsep}{0pt}%
    \setlength{\leftmargin}{#1}%
    \setlength{\rightmargin}{#2}%
    \setlength{\listparindent}{\parindent}%
    \setlength{\itemindent}{\parindent}%
    \setlength{\parsep}{\parskip}%
  }%
  \item[]}{\end{list}}
%    \end{macrocode}
% \end{macro}

% \subsection{The bibliography and references}
% \label{sec:biblio}
% \begin{macro}{\beforebiblioskip}
% \begin{macro}{\afterbiblioskip}
% I'm not entirely sure why I introduced these two commands, since
% directly putting |\vspace| before and after
% |\begin{thebibliography}| seems to work well enough. 
%    \begin{macrocode}
\newlength{\beforebiblioskip}
\setlength{\beforebiblioskip}{0pt}
\newlength{\afterbiblioskip}
\setlength{\afterbiblioskip}{0pt}
%    \end{macrocode}
% \end{macro}
% \end{macro}
% \begin{macro}{thebibliography}
% The environment itself is done as starred section with extra
% vertical space after heading. This space is currently zero. 
% In the book style we do the cunning trick with rereading the file in
% order to get the references into the |\ADS@references| macro.
% 
% \changes{3.05}{2003/01/02}{Removed the \cs{beforebiblioskip} and
% \cs{afterbiblioskip} commands, since they cause unwanted vertical
% space in some instances.}  
%    \begin{macrocode}
\renewenvironment{thebibliography}{%
  \if@book\RereadBiblio\fi
%%  \vspace*{\beforebiblioskip}
  \section*{REFERENCES} \frenchspacing
%%  \vspace*{\afterbiblioskip}
  \small
%    \end{macrocode} 
% The references themselves are set as a list in order 
% to get the hanging indentation. They are done in a smaller font and
% with french spacing so we needn't worry about the distinction
% between sentence-ending periods and abbreviations. 
%    \begin{macrocode} 
  \list{\null}{\leftmargin 1.5em\labelwidth\z@\labelsep\z@\itemindent -1.5em
    \itemsep\z@\parsep\z@\usecounter{enumi}}
  \sloppy\clubpenalty4000\widowpenalty4000
  \sfcode`\.=1000\relax\@biblisttrue}
{\endlist\@biblistfalse\normalsize\@outputaddressesifnecessary}
%    \end{macrocode}
% \end{macro}
% \subsection{Uniquification of cross-references}
% \label{sec:uniquification}
% 
% \begin{macro}{\articlename}
% In order that the |\label{}| command in individual articles be
% guaranteed to give a unique anchor, we append the article name (held
% in the |\articlename| command) to each key. This requires
% redefinition of the following commands: |\cite|, |\nocite|, |\ref|,
% |\pageref|, |\label|. 
% \begin{macro}{\oldcite}
% \begin{macro}{\oldnocite}
% \begin{macro}{\oldref}
% \begin{macro}{\oldpageref}
% \begin{macro}{\oldlabel}
% The original versions of these commands are
% saved as |\oldcite|, |\oldnocite|, etc. 
%
% The uniquifying of cross-references is now always done, so 
% we don't get multiply-defined labels in the single article case
%    \begin{macrocode}
\newcommand{\articlename}{anon}
\AtBeginDocument{
  \let\oldcite=\cite 
  \let\oldnocite=\nocite 
  \let\oldref=\ref 
  \let\oldpageref=\pageref
  \let\oldlabel=\label
}
%    \end{macrocode}
% \end{macro}
% \end{macro}
% \end{macro}
% \end{macro}
% \end{macro}
% \end{macro}
%
% Uniquify all cross-references by prepending the article name.  This
% now has to go into |\AtBeginDocument| so it is working on the new
% versions that hyperref.sty has set up.  Then we need to put the
% |\@begofarticlehook| afterwards, so that the first page label is using
% the new definitions.  If amsmath is loaded, then we also need to patch
% the special label command that they use. \changes{4.0}{2019/11/18}{Fix labels for amsmath display equations, in case that package has been loaded.}
%    \begin{macrocode}
\AtBeginDocument{
  \renewcommand{\ref}[1]{\oldref{\articlename:#1}%
    \if@debug{\ClassInfo{rmaa}{Reference to \articlename:#1}}\fi}
  \renewcommand{\pageref}[1]{\oldpageref{\articlename:#1}%
    \if@debug{\ClassInfo{rmaa}{Page reference to \articlename:#1}}\fi}
  \renewcommand{\label}[1]{\oldlabel{\articlename:#1}%
    \if@debug{\ClassInfo{rmaa}{Label for \articlename:#1}}\fi}%
  \ifdef{\label@in@display}{%
    \let\oldlabel@in@display=\label@in@display
    \renewcommand{\label@in@display}[1]{\oldlabel@in@display{\articlename:#1}%
      \if@debug{\ClassInfo{rmaa}{AMSTeX Label for \articlename:#1}}\fi}%
  }{}
  \@begofarticlehook
}%
%    \end{macrocode}
% Citations are more complicated since they can contain a list of
% keys. These are really only here for use with the rmaacite package,
% which is utterly deprecated (02 Jan 2003) now that I have the natbib
% code working. The command \cs{cite} should \emph{not} be used with
% natbib: use \cs{citep}, \cs{citet}, etc.\@ instead. 
%    \begin{macrocode}
\renewcommand{\nocite}[1]{\@bsphack
  \@for\@citeb:=#1\do{%
    \edef\@citeb{\expandafter\@firstofone\@citeb:\articlename}%
    \if@filesw\immediate\write\@auxout{\string\citation{\@citeb}}\fi
    \@ifundefined{b@\@citeb:\articlename}{\G@refundefinedtrue
      \@latex@warning{Citation `\@citeb' undefined}}{}}%
  \@esphack}
\renewcommand{\cite}[1]{\oldcite{#1:\articlename}}
%   \end{macrocode}


% Horrible fiddling with the bibliography citation mechanism. 
% This has to be globally available in case rmaacite puts it in the
% \texttt{.aux} file. 
%    \begin{macrocode}
\def\lbibcite#1#2#3{\global\@namedef{b@#1}{#2}
                    \global\@namedef{by@#1}{#3}}
%    \end{macrocode}

% \begin{macro}{\oldbibitem}
% \begin{macro}{\@rmaaBibArg}
% \begin{macro}{\@rmaaSpace}
% \begin{macro}{rmaabib}
% Of course, |\bibitem| needs to get the uniquification treatment as
% well. While we are at it, we redefine it to use consecutive integers
% for the key if the user leaves it blank (using the counter
% |rmaabib|). This allows one to put 
% |\bibitem{}| for all your references without getting annoying
% messages about multiply-defined labels.
% 
% We also take the opportunity to insert the optional argument to
% \cs{bibitem} if it is missing in the form Author(DATE) so that natbib
% doesn't try to do numerical citations. Natbib already provides a
% hook \cs{@extra@b@citeb} to help uniquify the citation key, so we
% just plug into that. However, it puts it \emph{after} the original
% key, not before as I had had it. Easiest solution is to change the
% way that I do it. 02 Jan 2003. Note that the \cs{@extra@b@citeb}
% does not get written to the .aux  file in the \cs{citation} command
% but it does work. 
% \changes{3.14}{06 Feb 2003}{Reinstalled old definition of \cs{bibitem}
%  for the case where the nonatbib option is requested, so that the
%  rmaacite package will work. This means that you can't use rmaacite
%  and natbib in the same volume.}
%    \begin{macrocode}
\if@natbib\def\@extra@b@citeb{:\articlename}\fi
%%\else
\let\oldbibitem\bibitem
\newcommand{\@rmaaBibArg}{}
\newcommand{\@rmaaSpace}{ }
\newcounter{rmaabib}
\if@natbib
\renewcommand{\bibitem}[2][ABC(2003)]{%
  \renewcommand{\@rmaaBibArg}{#2}
  \ifx\@rmaaBibArg\@rmaaEmpty
     \addtocounter{rmaabib}{1}
     \oldbibitem[#1]{\thermaabib:\articlename}
  \else\ifx\@rmaaBibArg\@rmaaSpace
     \addtocounter{rmaabib}{1}
     \oldbibitem[#1]{\thermaabib:\articlename}
  \else 
     \oldbibitem[#1]{#2:\articlename}
  \fi\fi
  \renewcommand{\@rmaaBibArg}{}
  }
\else
\renewcommand{\bibitem}[1]{%
  \renewcommand{\@rmaaBibArg}{#1}
  \ifx\@rmaaBibArg\@rmaaEmpty
     \addtocounter{rmaabib}{1}
     \oldbibitem{\articlename:\thermaabib}
  \else\ifx\@rmaaBibArg\@rmaaSpace
     \addtocounter{rmaabib}{1}
     \oldbibitem{\articlename:\thermaabib}
  \else
     \oldbibitem{\articlename:#1}
  \fi\fi
  \renewcommand{\@rmaaBibArg}{}
  }
\fi
%%\fi
%    \end{macrocode}
% \end{macro}
% \end{macro}
% \end{macro}
% \end{macro}

% \begin{macro}{oldpcite}
% \begin{macro}{oldscite}
% Then, we need to repeat the above scheme, but for the |rmaacite|
% definition. 
%    \begin{macrocode}
\newcommand{\@redefinebibitem}{%
  \let\oldbibitem\bibitem
  \let\oldcite\cite
  \let\oldpcite\pcite
  \let\oldscite\scite
  \def\bibitem[##1<##2>]##3{%
    \renewcommand{\@rmaaBibArg}{##3}
    \ifx\@rmaaBibArg\@rmaaEmpty
       \addtocounter{rmaabib}{1}
       \oldbibitem[##1<##2>]{\thermaabib:\articlename}
    \else\ifx\@rmaaBibArg\@rmaaSpace
       \addtocounter{rmaabib}{1}
       \oldbibitem[##1<##2>]{\thermaabib:\articlename}
    \else 
       \oldbibitem[##1<##2>]{##3:\articlename}
    % trust the author here
%%       \oldbibitem[##1<##2>]{##3}
    \fi\fi
    }
  }
%    \end{macrocode}
% \end{macro}
% \end{macro}


% \subsection{HTML Table of Contents}
% \label{sec:hyper-TOC}
% \begin{macro}{\HyperTOCEntry}
% \begin{macro}{\HyperTOCpart}
% \begin{macro}{\HyperTOCmeeting}
% \begin{macro}{\HyperTOCeditors}
% \begin{macro}{\HyperTOCpreamble}
% These commands are used to write information to the \texttt{.htoc}
% file, which is used in generating a Web Page of the Table of
% Contents. Since they only ever appear preceded by |\protect| in the
% argument to an |\addtocontents| command,  they should never need to
% be expanded by the macros. However, we define them here just in case
% \dots
%    \begin{macrocode}
\newcommand{\HyperTOCEntry}[3]{\relax}
\newcommand{\HyperTOCpart}[1]{\relax}
\newcommand{\HyperTOCmeeting}[3]{\relax}
\newcommand{\HyperTOCeditors}[1]{\relax}
\newcommand{\HyperTOCpreamble}{\relax}
%    \end{macrocode}
% \end{macro}
% \end{macro}
% \end{macro}
% \end{macro}
% \end{macro}
%

% \subsection{Formatting the poster abstracts} 
% These can now be used without using the book option.
%    \begin{macrocode}

\newcommand{\@makeminititle}{%
  \newpage
  \null 
  \vskip 3em%
  \begin{center}%
    {\@title \par}%
  \end{center}
  \vskip 3em%
  }

\newif\if@abstractsads
\@abstractsadsfalse
\newcommand\abstractsADS{\@abstractsadstrue}

\if@book\RequirePackage{multicol}\fi
\newenvironment{abstracts}[1][Abstracts of Contributed Papers]{%
  %% First, make sure we start on an odd page (06 Jun 2000)
  \if@book\clearpage\thispagestyle{empty}\cleardoublepage\fi
  %%
  \renewcommand{\articlename}{ABSTRACTS}
  \title{#1}\shortauthor{Abstracts}\shorttitle{Abstracts}
  \listofauthors{}
  \if@book
  \addcontentsline{toc}{chapter}{\protect\@fiddleauthors{\@TOCtitle}{\rmaaTOCauthors}}%
  \addtocontents{htc}{\protect\HyperTOCEntry{\@rawtitle}{\@rawauthors}{\articlename}}
  \fi
  \pagestyle{rmaa} % this now seems necessary (06 Jun 2000)
  \thispagestyle{rmaafirstpage}
%%  \twocolumn[\@makeminititle]
  \renewcommand{\abstract}[1]{%
    \renewcommand\ADS@abstract{##1}%
    \renewcommand{\@abstract}{##1}\@writeabstract%
  }
  \renewcommand{\title}[2][]{%
    \gdef\ADS@authors{}%
    \gdef\ADS@affil{}%
    \gdef\ADS@references{}%
    \gdef\ADS@keywords{}%
    \renewcommand{\@rmaaArg}{##1}
    \ifx\@rmaaArg\@rmaaEmpty
       \renewcommand{\@TOCaftertitle}{}
    \else
       \renewcommand{\@TOCaftertitle}{ (\@rmaaArg)}
    \fi
    \renewcommand{\@title}{\rmaa@TOCTitleStyle{##2}}%
    \renewcommand{\@TOCtitle}{\rmaa@TOCTitleStyle{##2}\@TOCaftertitle}%
    \renewcommand{\rmaaIndexTitle}{\rmaa@IndexTitleStyle{##2\@TOCaftertitle.}}%
    \renewcommand\ADS@title{##2}%
    \xdef\ADS@firstpage{\@arabic{\c@page}}%
    }
  \renewcommand{\author}[1]{\let\@thanks\@empty\renewcommand\@author{##1} }
  \renewcommand{\footnotetext}[2][XX]{\par\noindent\textsuperscript{##1}~##2}
  \renewcommand{\footnoterule}{\kern-3\p@ \hrule \@width 3em \kern 2.6\p@}
  \if@book\@begofarticlehook\fi
  \if@book\begin{multicols}{2}[\@makeminititle]\else\@makeminititle\fi
  \begingroup
  \setcounter{footnote}{0}
  }% and at the end of the environment....
  {\endgroup\if@book\end{multicols}\listofabstracts\label{RMAALastPage}\clearpage\fi}
%% auxiliary commands used:
%% \@writeabstract does all the actual formatting 
\newcommand{\@writeabstract}{%
  \if@book%
  \addcontentsline{loa}{abstract}{\protect\@fiddleauthors{\@TOCtitle}{\rmaaTOCauthors}}%
  \fi%
  \begin{center}\@title\par\@author\end{center}
  \@abstract
  \par\noindent\footnoterule\nobreak\footnotesize\@thanks
  \par\noindent\hrulefill\goodbreak
  \vskip 1cm \@plus0.5cm
  %% write ADS entry
  \if@book\if@abstractsads\ADS@dumpdata\fi\fi%
  \endgroup % this pair is to make each abstract a self-contained group
  \begingroup
  \setcounter{footnote}{0}
  }

\if@book
%% List of abstracts (stuff written to .loa file)
\newcommand{\listofabstracts}{%
  \begin{multicols}{2}[%
    \vskip 1.5cm \vfill % leave at least 1.5 cm and go to bottom of page if possible
    \centerline{\MakeUpperCase{List of Abstracts}}
    \vskip 1.5cm \@plus0.5cm
    ][6cm] % this makes it start a new page if there is not at least 6cm of space
    \renewcommand{\@fiddleauthors}[2]{##1\hskip1em plus2em minus0.4em\penalty-100  ##2}
    \parindent\z@
%%    \parskip.5\p@ \@plus .7\p@\relax 
    \parskip3\p@ \@plus 1\p@ \@minus 1\p@\relax % allow quite a lot of
                                % leeway here to avoid breaks inside
                                % entries 
    \renewcommand{\\}{\relax} % temporarily turn off explicit line breaks just in case
    \renewcommand{\altaffilmark}[1]{\relax} % ditto with this
    \@starttoc{loa}
  \end{multicols}
  }
\newcommand*{\l@abstract}[2]{%
  \par\penalty-1000
  #1\hskip 1em\penalty0\null\nobreak\hfill#2}
\fi

%    \end{macrocode}
 
% \subsection{Code for typesetting an entire volume of conference
% proceedings} \label{sec:whole-volume}
%    \begin{macrocode}

\if@book
%% environment for including the individual articles
\newcommand{\@fiddleauthors}[2]{\relax}
\renewcommand{\theequation}{\@arabic\c@equation}
\newenvironment{articles}%
{ \begingroup % make it a group so that author definitions do not
              % spill over
  \gdef\ADS@authors{}
  \pagestyle{rmaa}
  \renewenvironment{document}{%
    \@begofarticlehook
    \addtocounter{chapter}{1}%
    \addcontentsline{toc}{chapter}{%
      \protect\@fiddleauthors{\@TOCtitle}{\rmaaTOCauthors}}%
    \addtocontents{htc}{\protect\HyperTOCEntry{\@rawtitle}{\@rawauthors}{\articlename}}
    }{\relax}
  \renewcommand{\documentclass}[2][foo]{\relax} 
  \renewcommand{\usepackage}[2][foo]{%
    \def\@packagename{##2}\def\@rmaacitename{rmaacite}%
    \ifx\@packagename\@rmaacitename
    \@rmcitetrue\ClassInfo{rmaa}{This paper uses rmaacite}
    \makeatletter\input{rmaacite.sty}\makeatother % read it in here so that all changes
    \else                                         % are local to paper
    \if@debug\ClassWarning{rmaa}{Package \@packagename\space was requested by \articlename\space}\fi
    \fi} 
  %% Don't allow individual articles to mess with the pages, volume or year
  \renewcommand{\SetFirstPage}[1]{\relax}
  \renewcommand{\SetYear}[1]{\relax}
  \renewcommand{\SetVolume}[1]{\relax}
  \renewcommand{\SetConfTitle}[1]{\relax}
  %% Don't allow individual articles to change these
  \renewcommand{\nonstopmode}{\relax}
  \renewcommand{\scrollmode}{\relax}
  \renewcommand{\outputplates}{\relax}
  %% XXXX need to check if using rmaacite and act accordingly
  %% avoid multiple declarations of the same length
  \let\oldnewlength=\newlength
  \renewcommand{\newlength}[1]{%
    \edef\reserved@a{\expandafter\@gobble\string ##1}%
    \@ifundefined{\reserved@a}{%
      \oldnewlength{##1}}{%
      \if@debug{\ClassInfo{rmaa}{\reserved@a\space already defined: doing nothing}}\fi%
      \relax}}%
  \renewcommand\newcommand{\providecommand}
  }%
{\endgroup\onecolumn}
     

%    \end{macrocode}
% Command to be used in `articles' environment for including an
% article
%    \begin{macrocode}
\newcommand{\article}[2][]{%
  % optional first argument is extra text to put in brackets after title in TOC
  % e.g. ``Invited Talk'' or ``Poster''
  \renewcommand{\articlename}{#2}
  % reset loads of counters
  \setcounter{footnote}{0}
  \setcounter{figure}{0}
  \setcounter{table}{0}
  \setcounter{section}{0}
  \setcounter{equation}{0}
  \setcounter{rmaabib}{0}
  \setcounter{author}{0}
  \setcounter{AUTHOR}{1}
  \gdef\ADS@authors{}%
  \gdef\ADS@affil{}%
  \gdef\ADS@affiliations{}%
  \gdef\ADS@references{}%
  \@firstauthortrue
  \@firstaffiltrue
  \@firstaltaffiltrue
  \@firstkwdtrue
  \@rmcitefalse
  \if@suppressfulladdresses\else\@needaddressestrue\fi
  \renewcommand{\@rmaaArg}{#1}
  \ifx\@rmaaArg\@rmaaEmpty
     \renewcommand{\@TOCaftertitle}{}
  \else
     \renewcommand{\@TOCaftertitle}{ (#1)}
  \fi
%    \end{macrocode}
% In the main journal articles must start on an odd page. 
%    \begin{macrocode}
  \if@journal\clearpage\thispagestyle{empty}\cleardoublepage\fi
%% Now input the file, then run the end-of-article hook (reinstated 28
  %% Sep 2001)
  {\input{\@articledir/\articlename.tex}\@endofarticlehook}}
%    \end{macrocode}
% Command to separate articles into sessions. Write an entry in the
% TOC, makes sure we start on an odd page (added 06 Jun 2000), then do
% a title page if \cs{SessionTitlePages} was called. 
% \changes{3.12}{04 Feb 2003}{Added option to have title page for each
% session. Turn it on with \cs{SessionTitlePages}.}
%    \begin{macrocode}
\newcommand{\session}[1]{
  %% this does nothing except write an entry in the TOC
  \addcontentsline{toc}{part}{\rmaa@TOCPartStyle{#1}}
  \addtocontents{htc}{\protect\HyperTOCpart{#1}}
  %% and also makes sure we start on an odd page (added 06 Jun 2000)
  \clearpage\thispagestyle{empty}\cleardoublepage
  \stepcounter{part}
  \if@sessiontitlepages{%
    \onecolumn
    \thispagestyle{empty}
    \renewcommand{\articlename}{Section\thepart TitlePage}
    \label{RMAAFirstPage}
    \make@session@title{#1}
    \label{RMAALastPage}
    \clearpage\thispagestyle{empty}\cleardoublepage
  }\fi
  }
\newcommand{\make@session@title}[1]{%
  {%
    \setlength\parindent{0pt}
    \vspace*{\stretch{1}}
    \rule{\linewidth}{1mm}
    \begin{flushright}\Huge\bfseries\boldmath
      #1
    \end{flushright}
    \rule{\linewidth}{1mm}
    \vspace*{\stretch{2}}
  }
}
%    \end{macrocode}
% Command to turn on separate title pages for each session. 
%    \begin{macrocode}
\newcommand{\SessionTitlePages}{\@sessiontitlepagestrue}
\newcommand{\BoldSessionTitles}{%
  \renewcommand{\rmaa@TOCPartStyle}[1]{\textbf{\rmaa@TOCTitleStyle{##1}}}
  }
\newcommand{\ArticleDir}[1]{\renewcommand{\@articledir}{#1/}\graphicspath{{1/}}}
\newcommand{\@articledir}{.}
%\newcommand{\ps@rmaafront}{\ps@rmaa}% page style for front matter (same as ever, actually)
\newcommand{\ps@rmaafront}{% page style for front matter:  page # in each foot
  \renewcommand\@oddfoot{\hfil{\thepage}}%
  \renewcommand\@evenfoot{{\thepage}\hfil}%
  % shorttitle in both heads
  \renewcommand\@evenhead{\hfil\@shorttitle\hfil}%
  \renewcommand\@oddhead{\@evenhead}
  }
\newcommand{\ps@rmaafrontfirstpage}{% page style for first page of front matter
  % page # in each foot, on outside edge
  \renewcommand\@oddfoot{\hfil\thepage}%
  \renewcommand\@evenfoot{\thepage\hfil}%
  % nothing in the even heads
  \renewcommand\@evenhead{\relax}%
  \renewcommand\@oddhead{\relax}%
  }
%% slight modification to the standard \frontmatter and \mainmatter
%% commands, just to set the \pagestyle
\renewcommand\frontmatter{%
  \latexonecolumn
  \cleardoublepage
  \@frontmattertrue
  \@mainmatterfalse
  \pagenumbering{roman}
  \pagestyle{rmaafront}
  \addtocounter{page}{2} % leave space for the title page
  }
\renewcommand\mainmatter{%
  \pagestyle{empty}
  \if@rmtwocol\twocolumn\fi
  \cleardoublepage
  \@frontmatterfalse
  \@mainmattertrue
  \pagenumbering{arabic}
  \pagestyle{rmaa}
  }
\renewcommand\backmatter{%
  \clearpage\pagestyle{empty}\cleardoublepage
  \@mainmatterfalse
  %% put extra space in the TOC
  \addtocontents{toc}{\bigskip}
  }
%    \end{macrocode}
%
% \subsection{Front sections}\label{sec:front-sections}
% \begin{macro}{\notitle}
% \begin{macro}{\shortname}
% \begin{macro}{frontsection}
% The environment \texttt{frontsection} is for doing a section in the
% front matter (preface, list of participants, or whatever). It has 
% one optional argument:
% \begin{enumerate}
% \item[\#1] Short title (no embedded spaces) for
% use in making PS files. E.g., ``list\_of\_participants''.   
% \end{enumerate}
% and two mandatory arguments: 
% \begin{enumerate}
% \item[\#2] Title of section (e.g., `Preface')
% \item[\#3] Commands to be evaluated before processing the contents
% of the environment. This is mainly useful for putting in the things
% that would go in the preamble of a normal article  (e.g.,
% |\listofauthors|, |\indexauthor|, |\notitle|) 
% \end{enumerate}
%    \begin{macrocode}
\newcommand{\notitle}{\@titlevisiblefalse} \newcommand{\shortname}{}
\newenvironment{frontsection}[3][]{%
  \begingroup
  \renewcommand{\shortname}{#1}
  \ifx\shortname\@rmaaEmpty
     \renewcommand{\articlename}{#2}
  \else
     \renewcommand{\articlename}{\shortname}
  \fi
  \@begofarticlehook
  \@titlevisibletrue % default
  \thispagestyle{rmaafrontfirstpage}
  \listofauthors{} % blank out this in case it is not specified
  \title{{#2}}\shorttitle{#2}\shortauthor{#2}
  #3 % evaluate the command given as second mandatory arg
  \addcontentsline{toc}{chapter}{\protect\@fiddleauthors{\@TOCtitle}{\rmaaTOCauthors}}%
  \addtocontents{htc}{\protect\HyperTOCEntry{\@rawtitle}{\@rawauthors}{\articlename}}
  \if@titlevisible
     \vspace*{1cm}\centerline{\rmaa@FrontSecStyle{\@title}}\vspace*{1cm}
  \fi
  \setlength{\parskip}{3ex}
  }{\frontlabel{RMAALastPage}\clearpage\endgroup}
%    \end{macrocode}
% \end{macro}
% \end{macro}
% \end{macro}
%    \begin{macrocode}
%% Definition of the command \frontlabel, which rights out something like `F4' to
%% the .aux file, instead of `\textlatin  {i}'. This is so it is easy to post-process
%% the file for making PS files of individual sections
\newcommand{\frontlabel}[1]{%
  \@bsphack
  \protected@write\@auxout{}%
     {\string\newlabel{\articlename:#1}{{}{F\@arabic{\c@page}}}}
  \@esphack
  }
%% User command to say X, Y, Z,\\ Editors of the Proceedings
\newcommand{\EditorSignature}{\vspace*{\parskip}\rightline{\@editors}%
  \rightline{Editors of the proceedings}}
%% Fiddle with \thesection, etc. so that they don't include the chapter number
%% Basically go back to the `article' definitions of these
\renewcommand{\thesection}{\@arabic\c@section}
\renewcommand{\thefigure}{\@arabic\c@figure}
\renewcommand{\thetable}{\@arabic\c@table}

\newcommand{\@oldargs}{}
\newcommand{\@newargs}{}
%% add TOC entry for each paper
%\renewcommand{\title}[1]%
%{\renewcommand{\@title}{\rmaa@TitleStyle{#1}}%
%\addtocounter{chapter}{1}\addcontentsline{toc}{chapter}{#1}%
%}

%% Formatting of the table of contents
%% dimensions of various parts
\newlength{\TOCtextwidth}\setlength{\TOCtextwidth}{12.65cm}
\newlength{\TOCpagenumwidth}\setlength{\TOCpagenumwidth}{1.3cm}
\newlength{\TOCentrywidth}\setlength{\TOCentrywidth}{\TOCtextwidth}
\addtolength{\TOCentrywidth}{-\TOCpagenumwidth} 
\newlength{\TOCextramargin}\setlength{\TOCextramargin}{1.25cm}
\newcommand{\TOCbox}[1]{\hspace{\TOCextramargin}\parbox{\TOCtextwidth}{#1}\par}
\newlength{\TOCextramarginx}\setlength{\TOCextramarginx}{\TOCextramargin}\addtolength{\TOCextramarginx}{4.3cm}
\newlength{\TOCtextwidthx}\setlength{\TOCtextwidthx}{\TOCtextwidth}\addtolength{\TOCtextwidthx}{-4.3cm}
\newcommand{\TOCboxx}[1]{\hspace{\TOCextramarginx}\parbox{\TOCtextwidthx}{#1}\par}

\newcommand{\RMAAbannerimage}{rmsc}
\renewcommand\tableofcontents{%
  \renewcommand{\articlename}{CONTENTS}
  \begingroup
  \@begofarticlehook
  \if@twocolumn
     \@restonecoltrue\onecolumn
  \else
     \@restonecolfalse
  \fi
  %% 9pt for journal style
  \if@proceedings\else\small\fi
  \if@newstyle\small\fi %% same for new style
  \if@proceedings\thispagestyle{rmaafrontfirstpage}\else\thispagestyle{empty}\fi
  \shortauthor{Contents}\shorttitle{Contents}
  \vspace*{18mm}
  \if@printbanner
  %% put in the banner header (need to get better image of this)
     \TOCbox{\hspace*{-0.15cm}\includegraphics[width=0.9\TOCtextwidth]{\RMAAbannerimage}}
     \vspace*{8mm}
  %% volume and publication date (should really be in funny font)
     \TOCboxx{\@volumen\\[2mm] \@mes\ \@anyo}
  \else
  %% put in the banner header (need to get better image of this)
     \TOCbox{\hspace*{-0.15cm}\includegraphics[width=0.9\TOCtextwidth]{\RMAAbannerimage}}
     \vspace*{12.5mm}
  %% volume and publication date (should really be in funny font)
     \TOCboxx{~\\[2mm] ~}
  \fi
  \vspace*{16mm}
  %% This is the command that fiddles with the TOC entry so that the
  %% title is flushleft but the authors are flushright
  \renewcommand{\@fiddleauthors}[2]{%
    % first put the paper title
    \renewcommand{\\}{\relax} % temporarily turn off explicit line breaks just in case
    \renewcommand{\altaffilmark}[1]{\relax} % ditto with this
    \renewcommand{\textsuperscript}[1]{\relax} % ditto with this 
    ##1
    % now save the author list in a box
    \sbox\@tempboxa{##2}%
    % check that box is shorter than the total width available for the entry
    \ifdim \wd\@tempboxa <\TOCentrywidth
       % if so, then put it right here, flush against right margin
       \penalty10000\hskip1em plus10em\penalty0\hbox{}\nobreak\hfill\box\@tempboxa
       % Note that we force a space of at least 1em between title and authorlist, 
       % stretchable to 11em since we may have very short titles but still with an 
       % authorlist too long to go on the same line.  The only linebreak allowed is 
       % directly before the authorlist.
    \else
       % if it is too long, we set the original authorlist text in its own parbox
       \par\vskip1ex\parbox[b]{\TOCentrywidth}{\raggedleft ##2}
    \fi
    }
  \@starttoc{toc}%
  \if@restonecol\twocolumn\fi
  \frontlabel{RMAALastPage}
  \clearpage
  \thispagestyle{empty}
  \cleardoublepage % make sure we start on an odd page
  \endgroup

  %%% ALSO, do the HTML TOC
  %% The file that will be used to make the HTML TOC
  \@starttoc{htc}
  %% All the stuff in this file does nothing in the book style
  \addtocontents{htc}{\protect\HyperTOCmeeting{\@conftitle}{\@reunionplace}{\@reuniondate}}
  \addtocontents{htc}{\protect\HyperTOCeditors{\@editors}}
  \addtocontents{htc}{\protect\HyperTOCpreamble}
  \addtocontents{htc}{\protect\HyperTOCEntry{Table of Contents}{}{\articlename}}

  %% and the ADS file
  \ADS@openfile
  }
%% This is now much simpler
\renewcommand\l@chapter[2]{%TOC entry for articles
  \ifnum \c@tocdepth >\m@ne
    \vskip 1.0em \@plus\p@
    \TOCbox{%
      \parbox[b]{\TOCentrywidth}{#1}\parbox[b]{\TOCpagenumwidth}{\hfill#2}
    }
  \fi} 
\renewcommand\l@part[2]{%TOC entry for sessions: centered, no page number
  \ifnum \c@tocdepth >\m@ne
    \vskip 1.5em \@plus1.5em \goodbreak
    \TOCbox{%
      \parbox[b]{\TOCentrywidth}{\hfil#1\hfil}\hfill
    }
  \fi\nobreak} 

%% Formatting theindex
%% strongly encourage page breaks between index entries
\renewcommand{\@idxitem}{\par\penalty-1000\hangindent.45cm} 
\renewcommand{\indexspace}{\par\penalty-1000\vskip 10\p@ \@plus15\p@ \@minus3\p@\relax}
\renewenvironment{theindex}{%
  %% first, make sure we start on an odd page (06 Jun 2000)
  \clearpage\thispagestyle{empty}\cleardoublepage
  %% this looks a bit like an article but is two columns and without abstract, etc
  \renewcommand{\articlename}{INDEX}
  \title{Author Index}\shortauthor{Author Index}\shorttitle{Author Index}
  \listofauthors{}
  \addcontentsline{toc}{chapter}{\protect\@fiddleauthors{\@TOCtitle}{\rmaaTOCauthors}}%
  \addtocontents{htc}{\protect\HyperTOCEntry{\@rawtitle}{\@rawauthors}{\articlename}}
  \pagestyle{rmaa}
  \thispagestyle{rmaafirstpage}
  \@begofarticlehook
  %% 9pt for journal style
  \if@proceedings\else\small\fi
  \if@newstyle\small\fi
  \twocolumn[\@makeminititle]
  %% vertical spacing (see also definition of \@idxitem and \indexspace
  \parskip\z@ \@plus .6\p@\relax % allow quite a lot of leaway here to avoid breaks inside entries
  %% horizontal spacing, breaking and glue parameters (see also the stuff in the definition 
  %% of \indexauthor
  \parindent\z@ 
  \exhyphenpenalty10000 % don't break on explicit hyphens, it always looks bad
  \rightskip=0pt plus.2em % allow a tiny bit of leeway at the ends of lines
  \spaceskip.3333em plus.5em minus.15em % a lot of leeway between words
  \xspaceskip.5em plus.5em minus.4em % this seems to have no effect
  \let\item\@idxitem
  \renewcommand{\\}{\relax} % temporarily turn off explicit line breaks just in case
  \renewcommand{\altaffilmark}[1]{\relax} % ditto with this
  \renewcommand{\textsuperscript}[1]{\relax} % ditto with this
  }{\label{RMAALastPage}\clearpage} % clearpage necessary so that we know about \shortauthor



%% end definitions for book styles
\fi

%    \end{macrocode}
% In order to do away with plates, we sometimes want a figure to
% appear in color only in the online version. The following command
% can bve used in a figure caption to typeset a note to this effect. 
%    \begin{macrocode}
\newcommand\ColorOnlineOnly{%
  \textsc{Note: this figure is available in color in the electronic
    version of this article, obtainable from}
  \texttt{http://\discretionary{}{}{}www.\discretionary{}{}{}astroscu.\discretionary{}{}{}unam.\discretionary{}{}{}mx/\discretionary{}{}{}\~{}rmaa/}.}
%% Plates
%% Define 2 counters: plate is the number in the whole document, plateloop is used 
%% internally by \outputplates
\newcounter{plate}
\newcounter{plateloop}
%% Currently limited to 15 plates in document
%% when we do a label, we need to write to \newlabel entries to the .aux file, one 
%% that says RMAAPlateX so that it can be referred to in \ouputplates, and one that
%% says articlename:RMAAPlateX so that the dvips program can find it and asociate it
%% with the right article
\newcommand{\@platelabeli}{\oldlabel{RMAAPlate1}\label{RMAAPlate1}}
\newcommand{\@platelabelii}{\oldlabel{RMAAPlate2}\label{RMAAPlate2}}
\newcommand{\@platelabeliii}{\oldlabel{RMAAPlate3}\label{RMAAPlate3}}
\newcommand{\@platelabeliv}{\oldlabel{RMAAPlate4}\label{RMAAPlate4}}
\newcommand{\@platelabelv}{\oldlabel{RMAAPlate5}\label{RMAAPlate5}}
\newcommand{\@platelabelvi}{\oldlabel{RMAAPlate6}\label{RMAAPlate6}}
\newcommand{\@platelabelvii}{\oldlabel{RMAAPlate7}\label{RMAAPlate7}}
\newcommand{\@platelabelviii}{\oldlabel{RMAAPlate8}\label{RMAAPlate8}}
\newcommand{\@platelabelix}{\oldlabel{RMAAPlate9}\label{RMAAPlate9}}
\newcommand{\@platelabelx}{\oldlabel{RMAAPlate10}\label{RMAAPlate10}}
\newcommand{\@platelabelxi}{\oldlabel{RMAAPlate11}\label{RMAAPlate11}}
\newcommand{\@platelabelxii}{\oldlabel{RMAAPlate12}\label{RMAAPlate12}}
\newcommand{\@platelabelxiii}{\oldlabel{RMAAPlate13}\label{RMAAPlate13}}
\newcommand{\@platelabelxiv}{\oldlabel{RMAAPlate14}\label{RMAAPlate14}}
\newcommand{\@platelabelxv}{\oldlabel{RMAAPlate15}\label{RMAAPlate15}}
\newcommand{\@outputplatelabeli}{\oldlabel{RMAAOutputPlate1}}
\newcommand{\@outputplatelabelii}{\oldlabel{RMAAOutputPlate2}}
\newcommand{\@outputplatelabeliii}{\oldlabel{RMAAOutputPlate3}}
\newcommand{\@outputplatelabeliv}{\oldlabel{RMAAOutputPlate4}}
\newcommand{\@outputplatelabelv}{\oldlabel{RMAAOutputPlate5}}
\newcommand{\@outputplatelabelvi}{\oldlabel{RMAAOutputPlate6}}
\newcommand{\@outputplatelabelvii}{\oldlabel{RMAAOutputPlate7}}
\newcommand{\@outputplatelabelviii}{\oldlabel{RMAAOutputPlate8}}
\newcommand{\@outputplatelabelix}{\oldlabel{RMAAOutputPlate9}}
\newcommand{\@outputplatelabelx}{\oldlabel{RMAAOutputPlate10}}
\newcommand{\@outputplatelabelxi}{\oldlabel{RMAAOutputPlate11}}
\newcommand{\@outputplatelabelxii}{\oldlabel{RMAAOutputPlate12}}
\newcommand{\@outputplatelabelxiii}{\oldlabel{RMAAOutputPlate13}}
\newcommand{\@outputplatelabelxiv}{\oldlabel{RMAAOutputPlate14}}
\newcommand{\@outputplatelabelxv}{\oldlabel{RMAAOutputPlate15}}
\newcommand{\@platerefi}{\oldpageref{RMAAPlate1}}
\newcommand{\@platerefii}{\oldpageref{RMAAPlate2}}
\newcommand{\@platerefiii}{\oldpageref{RMAAPlate3}}
\newcommand{\@platerefiv}{\oldpageref{RMAAPlate4}}
\newcommand{\@platerefv}{\oldpageref{RMAAPlate5}}
\newcommand{\@platerefvi}{\oldpageref{RMAAPlate6}}
\newcommand{\@platerefvii}{\oldpageref{RMAAPlate7}}
\newcommand{\@platerefviii}{\oldpageref{RMAAPlate8}}
\newcommand{\@platerefix}{\oldpageref{RMAAPlate9}}
\newcommand{\@platerefx}{\oldpageref{RMAAPlate10}}
\newcommand{\@platerefxi}{\oldpageref{RMAAPlate11}}
\newcommand{\@platerefxii}{\oldpageref{RMAAPlate12}}
\newcommand{\@platerefxiii}{\oldpageref{RMAAPlate13}}
\newcommand{\@platerefxiv}{\oldpageref{RMAAPlate14}}
\newcommand{\@platerefxv}{\oldpageref{RMAAPlate15}}
\newcounter{figplatei}
\newcounter{figplateii}
\newcounter{figplateiii}
\newcounter{figplateiv}
\newcounter{figplatev}
\newcounter{figplatevi}
\newcounter{figplatevii}
\newcounter{figplateviii}
\newcounter{figplateix}
\newcounter{figplatex}
\newcounter{figplatexi}
\newcounter{figplatexii}
\newcounter{figplatexiii}
\newcounter{figplatexiv}
\newcounter{figplatexv}
%% Printed output should use the counter `plate'
\renewcommand{\theplate}{\@arabic\c@plate}
\renewcommand{\theplateloop}{\@arabic\c@plateloop}
%% Command to output all the saved plates in this paper
\newcommand{\outputplates}{%
  \setcounter{plateloop}{0}
  \if@book\else\@endofarticlehook\fi % not entirely sure why, but this
                                     % seems to be needed
  \clearpage % first output the last page of text
  % now set pagestyle for blank even pages and make sure we start on an odd page
  \pagestyle{empty}
  \cleardoublepage
  \onecolumn
  \parindent\z@
  % first put in an extra sheet
  % the main purpose of this is so there is always a blank even page before 
  % each plate for filling in the PS file if the article has an odd number
  % of pages (see splitbook.sh)
  \null\vfil\centerline{PLATES}\vfil\clearpage\pagestyle{rmaaplate}\cleardoublepage
  % output each plate in turn
  \@whilenum\value{plateloop}<\value{plate}\do{%
    \addtocounter{plateloop}{1}
    % restore the figure counter from \@figcounti, \@figcountii, etc
    \setcounter{figure}{\value{figplate\roman{plateloop}}}
    % spew out the plate (in its own group for safety)
    {\csname @plate\roman{plateloop}\endcsname\par}
    }
  \pagestyle{rmaa}
  % reset the page number as if nothing had happened
  \addtocounter{page}{-\value{plate}}
  \addtocounter{page}{-\value{plate}}
  }
\newcommand\figcaption{%
  \def\@captype{figure}
  \renewcommand{\@makecaption}{\@makefigcap}
  \caption}
\newcommand\tabcaption{%
  \def\@captype{table}
  \renewcommand{\@makecaption}{\@maketabcap}
  \caption}

\newcommand{\plate}[2][1]{%
  %% User command to make a plate page: 
  %% Mandatory argument is what goes on the plate - this should not 
  %% include any floating environments (e.g. figure). Instead, \figcaption 
  %% should be used to do the captions. 
  %% Optional argument is the number of figures (i.e. captions) on the plate
  %% page (default is 1). This is necessary for incrementing the figure counter. 
  %%
  %% Note that we have to use a \gdef rather than an \xdef when we 
  %% save the plate output in a command. This means that in the book styles when it 
  %% comes to spewing out the plate, we will have forgotten about any user macros 
  %% that may have been defined. This can be got round by repeating the macros inside the plate.
  %% (Presumably, one would use a \newcommand rather than \renewcommand if they were new.)
  % increment global and local counters
  \addtocounter{plate}{1}
  % write appropriate label to .aux file
  \csname @platelabel\roman{plate}\endcsname
  % define a global macro called \@platei, \@plateii, \@plateiii, \@plateiv, ...
  \expandafter\gdef\csname @plate\roman{plate}\endcsname{%
    \csname @outputplatelabel\roman{plateloop}\endcsname
    \null\vfil #2 \vfil
    \cleardoublepage
    }
  % save the current shortauthor and shorttitle
  % need to be \xdef to take current values of \@shortauthor, etc 
  \expandafter\xdef\csname @shorta\roman{plate}\endcsname{\@shortauthor}
  \expandafter\xdef\csname @shortt\roman{plate}\endcsname{\@shorttitle}
  % save the current value of the figure counter for later use in \outputplates
  \setcounter{figplate\roman{plate}}{\value{figure}}
  % increment the figure counter according to the optional argument
  \addtocounter{figure}{#1}
  }
%% page style for plate pages 
\newcommand{\ps@rmaaplate}{%
  % Plate X (see page Y) in the foot
  \renewcommand\@oddfoot{\hfil\csname @shorta\roman{plateloop}\endcsname\ 
    (see page~\csname @plateref\roman{plateloop}\endcsname)\hfil}%
  \renewcommand\@evenfoot{\relax}%
  % nothing in even heads (blank pages)
  \renewcommand\@evenhead{\relax}
  % short title and page num in the odd heads
  \renewcommand\@oddhead{\hfil\csname @shortt\roman{plateloop}\endcsname%
    \hfil\llap{PLATE~\theplateloop}}
  }





%% Things for the User Guide
\if@guide
%% stolen from doc.sty
\def\MakeShortVerb#1{%
  \expandafter\ifx\csname cc\string#1\endcsname\relax
  \@shortvrbinfo{Made }{#1}%
  \add@special{#1}%
  \expandafter
  \xdef\csname cc\string#1\endcsname{\the\catcode`#1}%
  \begingroup
  \catcode`\~\active  \lccode`\~`#1%
  \lowercase{%
    \global\expandafter\let
    \csname ac\string#1\endcsname~%
    \gdef~{\verb~}}%
  \endgroup
  \global\catcode`#1\active
  \else
  \@shortvrbinfo\@empty{#1 already}%
  \fi}
\def\DeleteShortVerb#1{%
  \expandafter\ifx\csname cc\string#1\endcsname\relax
  \else
  \@shortvrbinfo{Deleted }{#1 as}%
  \rem@special{#1}%
  \global\catcode`#1\csname cc\string#1\endcsname
  \global \expandafter\let \csname cc\string#1\endcsname \relax
  \ifnum\catcode`#1=\active
  \begingroup
  \catcode`\~\active   \lccode`\~`#1%
  \lowercase{%
    \global\expandafter\let\expandafter~%
    \csname ac\string#1\endcsname}%
  \endgroup \fi \fi}
\def\@shortvrbinfo#1#2{%
  \PackageInfo{doc}{%
    #1\expandafter\@gobble\string#2 a short reference
    for \string\verb}}
\def\add@special#1{%
  \rem@special{#1}%
  \expandafter\gdef\expandafter\dospecials\expandafter
  {\dospecials \do #1}%
  \expandafter\gdef\expandafter\@sanitize\expandafter
  {\@sanitize \@makeother #1}}
\def\rem@special#1{%
  \def\do##1{%
    \ifnum`#1=`##1 \else \noexpand\do\noexpand##1\fi}%
  \xdef\dospecials{\dospecials}%
  \begingroup
  \def\@makeother##1{%
    \ifnum`#1=`##1 \else \noexpand\@makeother\noexpand##1\fi}%
  \xdef\@sanitize{\@sanitize}%
  \endgroup}
%% define | as shortcut for \verbatim
\MakeShortVerb{\|}
%% comand to specify arguments to LaTeX commands
\newcommand{\rmArg}[1]{%
  \ensuremath{\langle}{\scriptsize \textsf{\MakeUpperCase{#1}}}\ensuremath{\rangle}}
%\newcommand{\rmArg}[1]{\ensuremath{\langle}\textsc{#1}\ensuremath{\rangle}}
\providecommand{\BibTeX}{{\rmfamily B\kern-.05em%
    \textsc{i\kern-.025em b}\kern-.08em%
    T\kern-.1667em\lower.7ex\hbox{E}\kern-.125emX}}
%% Environment for code examples
\newenvironment{Code}{%
  \begin{list}{}{\setlength{\leftmargin}{0.05\columnwidth}} 
  \item[]}{\end{list}}
\fi
%</package>
%    \end{macrocode}

% \section{User commands for astronomical articles}
% \label{sec:user-commands}
% Some random definitions. These go both into the main class file and
% also into |adshtml.sty|.
% 
% Definitions of the journal name in different incarnations.  
%    \begin{macrocode}
%<*package|adshtml-style>
\newcommand{\RMAA}{RMxAA}
\newcommand{\RMAAlong}{Revista Mexicana de Astronom\'{\i}a y Astrof\'{\i}sica}
\newcommand{\RMAASC}{RMxAC}
\newcommand{\RMAASClong}{\RMAA\ (Serie de Conferencias)}
%    \end{macrocode}
% \begin{macro}{\ibidrule}
% This puts a horizontal rule with a period after it. It is used when
% the author list of a reference is identical to the preceding one. 
%    \begin{macrocode}
\newcommand{\ibidrule}{\rule[-0.03cm]{4.7em}{.01cm}.\ }
%    \end{macrocode}
% \end{macro}
% \begin{macro}{\ThisVolume}
% References to other articles in the same volume can be coded as,
% e.g., |\bibitem{} Bloggs, J. \ThisVolume{bloggs}|.
%    \begin{macrocode}
\newcommand{\ThisVolume}[1]{%
  \@year, \if@proceedings\RMAASC\else\RMAA\fi, \@volume,
  \if@book\oldpageref{#1:RMAAFirstPage}\else 000\fi\ (this volume)}
%    \end{macrocode}
% \end{macro}
% Other stuff that may be useful in ref lists.
%    \begin{macrocode}
\newcommand{\etal}{et al.\@}
%    \end{macrocode}
% Various new commands from AAS macros
%    \begin{macrocode}
\def\acknowledgments{\vskip 3ex plus .8ex minus .4ex}
\let\acknowledgements=\acknowledgments                  % second spelling
\let\jnl@style=\rm
\def\ref@jnl#1{{\jnl@style#1}}
\DeclareRobustCommand\aj{\ref@jnl{AJ}}                   
\DeclareRobustCommand\araa{\ref@jnl{ARA\&A}}             
\DeclareRobustCommand\apj{\ref@jnl{ApJ}}                 
\DeclareRobustCommand\apjl{\ref@jnl{ApJ}}                
\DeclareRobustCommand\apjs{\ref@jnl{ApJS}}               
\DeclareRobustCommand\ao{\ref@jnl{Appl.~Opt.}}           
\DeclareRobustCommand\apss{\ref@jnl{Ap\&SS}}             
\DeclareRobustCommand\aap{\ref@jnl{A\&A}}                
\DeclareRobustCommand\aapr{\ref@jnl{A\&A~Rev.}}          
\DeclareRobustCommand\aaps{\ref@jnl{A\&AS}}              
\DeclareRobustCommand\azh{\ref@jnl{AZh}}                 
\DeclareRobustCommand\baas{\ref@jnl{BAAS}}               
\DeclareRobustCommand\jrasc{\ref@jnl{JRASC}}             
\DeclareRobustCommand\memras{\ref@jnl{MmRAS}}            
\DeclareRobustCommand\mnras{\ref@jnl{MNRAS}}             
\DeclareRobustCommand\pra{\ref@jnl{Phys.~Rev.~A}}        
\DeclareRobustCommand\prb{\ref@jnl{Phys.~Rev.~B}}        
\DeclareRobustCommand\prc{\ref@jnl{Phys.~Rev.~C}}        
\DeclareRobustCommand\prd{\ref@jnl{Phys.~Rev.~D}}        
\DeclareRobustCommand\pre{\ref@jnl{Phys.~Rev.~E}}        
\DeclareRobustCommand\prl{\ref@jnl{Phys.~Rev.~Lett.}}    
\DeclareRobustCommand\pasp{\ref@jnl{PASP}}               
\DeclareRobustCommand\pasj{\ref@jnl{PASJ}}               
\DeclareRobustCommand\qjras{\ref@jnl{QJRAS}}             
\DeclareRobustCommand\skytel{\ref@jnl{S\&T}}             
\DeclareRobustCommand\solphys{\ref@jnl{Sol.~Phys.}}      
\DeclareRobustCommand\sovast{\ref@jnl{Soviet~Ast.}}      
\DeclareRobustCommand\ssr{\ref@jnl{Space~Sci.~Rev.}}     
\DeclareRobustCommand\zap{\ref@jnl{ZAp}}                 
\DeclareRobustCommand\nat{\ref@jnl{Nature}}              
\DeclareRobustCommand\iaucirc{\ref@jnl{IAU~Circ.}}
\DeclareRobustCommand\aplett{\ref@jnl{Astrophys.~Lett.}}
\DeclareRobustCommand\apspr{\ref@jnl{Astrophys.~Space~Phys.~Res.}}
\DeclareRobustCommand\bain{\ref@jnl{Bull.~Astron.~Inst.~Netherlands}}
\DeclareRobustCommand\fcp{\ref@jnl{Fund.~Cosmic~Phys.}}
\DeclareRobustCommand\gca{\ref@jnl{Geochim.~Cosmochim.~Acta}}
\DeclareRobustCommand\grl{\ref@jnl{Geophys.~Res.~Lett.}}
\DeclareRobustCommand\jcp{\ref@jnl{J.~Chem.~Phys.}}      
\DeclareRobustCommand\jgr{\ref@jnl{J.~Geophys.~Res.}}    
\DeclareRobustCommand\jqsrt{\ref@jnl{J.~Quant.~Spec.~Radiat.~Transf.}}
\DeclareRobustCommand\memsai{\ref@jnl{Mem.~Soc.~Astron.~Italiana}}
\DeclareRobustCommand\nphysa{\ref@jnl{Nucl.~Phys.~A}}
\DeclareRobustCommand\physrep{\ref@jnl{Phys.~Rep.}}
\DeclareRobustCommand\physscr{\ref@jnl{Phys.~Scr}}
\DeclareRobustCommand\planss{\ref@jnl{Planet.~Space~Sci.}}
\DeclareRobustCommand\procspie{\ref@jnl{Proc.~SPIE}}
\let\astap=\aap
\let\apjlett=\apjl
\let\apjsupp=\apjs
\let\applopt=\ao
\def\phn{\phantom{0}}
\def\phd{\phantom{.}}
\def\phs{\phantom{\ensuremath{-}}}
\def\phm#1{\phantom{#1}}
%    \end{macrocode}
% All these are redone as robust and using more modern commands where
% possible. The robustness is so they will not get expanded if used in
% abstracts and titles when these are passed to the ADS file. 
% \changes{3.16}{2003/04/07}{
%  Added \cs{Color} command for doing photometric colors.}
%    \begin{macrocode}
%%\def\sun{\hbox{$\odot$}}
\DeclareRobustCommand\sun{\ensuremath{\odot}}
%%\def\earth{\hbox{$\oplus$}}
\DeclareRobustCommand\earth{\ensuremath{\oplus}}
%%\def\lesssim{\mathrel{\hbox{\rlap{\hbox{\lower4pt\hbox{$\sim$}}}\hbox{$<$}}}}
%%\def\gtrsim{\mathrel{\hbox{\rlap{\hbox{\lower4pt\hbox{$\sim$}}}\hbox{$>$}}}}
\DeclareRobustCommand\lesssim{\mathrel{\hbox{\smaller\smaller\rlap{\hbox{\lower3pt\hbox{$\sim$}}}\raise1.5pt\hbox{$<$}}}}
\DeclareRobustCommand\gtrsim{\mathrel{\hbox{\smaller\smaller\rlap{\hbox{\lower3pt\hbox{$\sim$}}}\raise1.5pt\hbox{$>$}}}}
%%\def\lta{\mathrel{\spose{\lower 3pt\hbox{$\mathchar"218$}}
%%     \raise 2.0pt\hbox{$\mathchar"13C$}}}
%%\def\gta{\mathrel{\spose{\lower 3pt\hbox{$\mathchar"218$}}
%%     \raise 2.0pt\hbox{$\mathchar"13E$}}}
\DeclareRobustCommand\sq{\hbox{\rlap{$\sqcap$}$\sqcup$}}
%%\def\arcdeg{\hbox{$^\circ$}}
%%\def\arcmin{\hbox{$^\prime$}}
%%\def\arcsec{\hbox{$^{\prime\prime}$}}
\DeclareRobustCommand\arcdeg{\ensuremath{^\circ}}
\DeclareRobustCommand\arcmin{\ensuremath{^\prime}}
\DeclareRobustCommand\arcsec{\ensuremath{^{\prime\prime}}}
%%\def\fd{\hbox{$.\!\!^{\rm d}$}}
%%\def\fh{\hbox{$.\!\!^{\rm h}$}}
%%\def\fm{\hbox{$.\!\!^{\rm m}$}}
%%\def\fs{\hbox{$.\!\!^{\rm s}$}}
%%\def\fdg{\hbox{$.\!\!^\circ$}}
%%\def\farcm{\hbox{$.\mkern-4mu^\prime$}}
%%\def\farcs{\hbox{$.\!\!^{\prime\prime}$}}
%%\def\fp{\hbox{$.\!\!^{\scriptscriptstyle\rm p}$}}
\DeclareRobustCommand\fd{\ensuremath{.\!\!^\mathrm{d}}}
\DeclareRobustCommand\fh{\ensuremath{.\!\!^\mathrm{h}}}
\DeclareRobustCommand\fm{\ensuremath{.\!\!^\mathrm{m}}}
\DeclareRobustCommand\fs{\ensuremath{\mkern0.5mu.\mkern-5mu^\mathrm{s}\mkern-1mu}}
\DeclareRobustCommand\fdg{\ensuremath{.\!\!^\circ}}
\DeclareRobustCommand\farcm{\ensuremath{.\mkern-4mu^\prime}}
\DeclareRobustCommand\farcs{\ensuremath{\mkern1mu.\mkern-6mu^{\prime\prime}\mkern-1.5mu}}
%%\DeclareRobustCommand\farcs{\ensuremath{.\!\!^{\prime\prime}}}
\DeclareRobustCommand\fp{\ensuremath{.\!\!^{\scriptscriptstyle\rm p}}}
%%\def\micron{\hbox{$\mu$m}}
\DeclareRobustCommand\micron{\ensuremath{\mu\mathrm{m}}}
\let\la=\lesssim                        
\let\ga=\gtrsim
%%\def\case#1#2{\hbox{$\frac{#1}{#2}$}}
%%\def\slantfrac#1#2{\hbox{$\,^#1\!/_#2$}}
%%\def\onehalf{\slantfrac{1}{2}}
%%\def\onethird{\slantfrac{1}{3}}
%%\def\twothirds{\slantfrac{2}{3}}
%%\def\onequarter{\slantfrac{1}{4}}
%%\def\threequarters{\slantfrac{3}{4}}
\DeclareRobustCommand\case[2]{\ensuremath{\frac{#1}{#2}}}
\DeclareRobustCommand\slantfrac[2]{\ensuremath{\,^#1\!/_#2}}
\DeclareRobustCommand\onehalf{\slantfrac{1}{2}}
\DeclareRobustCommand\onethird{\slantfrac{1}{3}}
\DeclareRobustCommand\twothirds{\slantfrac{2}{3}}
\DeclareRobustCommand\onequarter{\slantfrac{1}{4}}
\DeclareRobustCommand\threequarters{\slantfrac{3}{4}}
%%\def\ubvr{\hbox{$U\!BV\!R$}}            
%%\def\ub{\hbox{$U\!-\!B$}}               
%%\def\bv{\hbox{$B\!-\!V$}}               
%%\def\vr{\hbox{$V\!-\!R$}}               
%%\def\ur{\hbox{$U\!-\!R$}}
%    \end{macrocode}
% New commands for ``photometric'' colors 
%    \begin{macrocode}
\DeclareRobustCommand\ubvr{\ensuremath{U\!BV\!R}}            
\DeclareRobustCommand\Color[2]{%
  \ensuremath{%
    #1\!\ifx#1V\!\fi-\!#2
  }}
\DeclareRobustCommand\ub{\Color{U}{B}}
\DeclareRobustCommand\bv{\Color{B}{V}}
\DeclareRobustCommand\vr{\Color{V}{R}}
\DeclareRobustCommand\ur{\Color{U}{R}}
%% New version of the \ion command, that actually works! 
%    \end{macrocode}
% \begin{macro}{\ion}
% This version of the |\ion| command allows you to write either
% |\ion{N}{4}| or |\ion{N}{iv}| in order to get N~IV. 
%    \begin{macrocode}
\DeclareRobustCommand\ion[2]{%
  \mbox{#1\kern0.2em%
    \smaller\rmfamily%
    \edef\@tempa{\@car#2\@nil}%
    \ifcat1\@tempa%
    \@Roman{#2}%
    \else%
    \uppercase{#2}%
    \fi}}
\newcounter{IonStage}
\DeclareRobustCommand\plainion[2]{\setcounter{IonStage}{#2}#1 \Roman{IonStage}}
%%\def\ion#1#2{\ensuremath{\mathrm{#1}\;}{\protect\small\rm\@Roman{#2}}}
%    \end{macrocode}
% \end{macro}
%    \begin{macrocode}
\DeclareRobustCommand\nodata{$\cdots$}
%    \end{macrocode}
% These are needed by the |\ion| command. Isn't there an easier way of
% doing this?
%    \begin{macrocode}
\let\@small=\small
\let\@normalsize=\normalsize
\let\@footnotesize=\footnotesize
\let\@scriptsize=\scriptsize
\let\@tiny=\tiny
\let\@large=\large
\let\@Large=\Large
\let\@LARGE=\LARGE
\let\@huge=\huge
\let\@Huge=\Huge
\newcommand{\smaller}{\small}
\renewcommand{\normalsize}{\renewcommand{\smaller}{\small}\@normalsize}
\renewcommand{\small}{\renewcommand{\smaller}{\footnotesize}\@small}
\renewcommand{\footnotesize}{\renewcommand{\smaller}{\scriptsize}\@footnotesize}
\renewcommand{\scriptsize}{\renewcommand{\smaller}{\tiny}\@scriptsize}
\renewcommand{\tiny}{\renewcommand{\smaller}{\tiny}\@tiny}
\renewcommand{\large}{\renewcommand{\smaller}{\normalsize}\@large}  
\renewcommand{\Large}{\renewcommand{\smaller}{\large}\@Large}  
\renewcommand{\LARGE}{\renewcommand{\smaller}{\Large}\@LARGE}  
\renewcommand{\huge}{\renewcommand{\smaller}{\LARGE}\@huge}  
\renewcommand{\Huge}{\renewcommand{\smaller}{\huge}\@Huge}  
%</package|adshtml-style>
%    \end{macrocode}


% \section{Other files}

% \subsection{The driver and make files}
% \label{sec:driver_make}
% These are put here for completeness, although they still need to be
% distributed separately since otherwise the user won't be able to
% unpack this file into its constituent parts.
% 
% The |Makefile| has two principal rules: |docs|, and |macros|. The
% |docs| rule produces a PDF and gzipped PS file of the documentation.
% I still haven't perfected the combination of rules required to get
% it to run \LaTeX{} enough times in all circumstances. The |macros|
% rule runs \LaTeX{} on |rmaa.ins| in order to unpack |rmaa.cls| and
% all the other little files (including the |Makefile| and the
% |rmaa.ins| files themselves\dots 
% 
% First set up lists of various types of file. 
%    \begin{macrocode}
%<*makefile>
DOCFILES=rmaa.pdf
MACRO_LTX_FILES=rmaa.cls rmaacite.sty adshtml.sty
SHELLSCRIPTS=makeads.sh
MACRO_OTHER_FILES=adshtml.perl sample.latex2html-init
MAKEFILE=Makefile
MACROFILES=$(MACRO_LTX_FILES) $(SHELLSCRIPTS) $(MACRO_OTHER_FILES)         
SKELETONFILES=ADS.tex 
%    \end{macrocode}
% Install the files read only into the package directory, specified via a path relative to the source directory. \changes{4.0}{2019/11/14}{Reorganize folder layout so that hardly anything is in the top level.  That way, things are clearer on github.}
%    \begin{macrocode}
INSTALL=install 
INSTALL_DIR=../packages
%    \end{macrocode}
% The default rule is to make everything. Installation must be
% explicitly asked for.
%    \begin{macrocode}
all: docs macros

install: install-docs install-macros

%    \end{macrocode}
% Most files are installed read-only to prevent editing. Skeleton
% files must be writeable. Must be careful not to overwrite the
% Makefile in the parent directory. 
%    \begin{macrocode}
install-macros: macros
        $(INSTALL) -m 444 $(MACRO_LTX_FILES) $(MACRO_OTHER_FILES) $(INSTALL_DIR)
        $(INSTALL) -m 555 $(SHELLSCRIPTS) $(INSTALL_DIR)
        $(INSTALL) -m 644 $(SKELETONFILES) $(INSTALL_DIR)

install-docs: docs
        $(INSTALL) rmaa.pdf $(INSTALL_DIR)/rm-fulldocs.pdf

%    \end{macrocode}
% Make the PDF and gzipped PS files. 
%    \begin{macrocode}
docs: $(DOCFILES)

%    \end{macrocode}
% Make everything else if either the |.dtx| or |.ins| files have
% changed. 
%    \begin{macrocode}
macros: rmaa.dtx rmaa.ins
%    \end{macrocode}
% Ensure that all the macro files can be written to then recreate
% them. The |touch| is there in case any of the files don't exist
% yet. 
%    \begin{macrocode}
        touch $(MACROFILES) $(MAKEFILE)
        chmod u+w $(MACROFILES) $(MAKEFILE) rmaa.ins
        latex "\input{rmaa.ins}"
%    \end{macrocode}
% \TeX{} automatically adds the |.tex| extension to extensionless
% files. We fix this and take the opportunity to convert
% back to TABs all that it has lovingly converted to spaces. 
%    \begin{macrocode}
        cp -f $(MAKEFILE) $(MAKEFILE).BCK && unexpand $(MAKEFILE).tex >| $(MAKEFILE) && rm $(MAKEFILE).tex
%    \end{macrocode}
% Make the shell scripts executable and everything unwriteable (so
% they don't get hand-edited by mistake). However, undo this for the
% files that the user is \emph{supposed} to modify. 
%    \begin{macrocode}
        chmod a+x $(SHELLSCRIPTS)
        chmod a-w $(MACROFILES) $(MAKEFILE)
        chmod u+w $(SKELETONFILES)

%    \end{macrocode}
% We now use PDF\LaTeX{} for this.
%    \begin{macrocode}
rmaa.idx: rmaa.dtx
        @echo "REGENERATING .idx FILE"
        pdflatex "\nonstopmode\input{rmaa.dtx}"

rmaa.aux: rmaa.dtx rmaa.ind rmaa.gls
        @echo "REGENERATING .aux FILE"
        pdflatex "\nonstopmode\input{rmaa.dtx}"

rmaa.ind: rmaa.idx
        @echo "REGENERATING INDEX"
        makeindex -s gind.ist rmaa

rmaa.gls: rmaa.glo
        @echo "REGENERATING CHANGELOG"
        makeindex -s gglo.ist -o rmaa.gls rmaa.glo

rmaa.pdf: rmaa.aux rmaa.ind rmaa.gls
        @echo "REGENERATING .pdf FILE"
        pdflatex "\nonstopmode\input{rmaa.dtx}"

clean: 
        rm -f rmaa.aux rmaa.dvi rmaa.ind rmaa.gls
%</makefile>
%    \end{macrocode}

% The |rmaa.ins| file uses the \textsf{docstrip} package to split out
% all the files and strip off this documentation.
%    \begin{macrocode}
%<*ins>
%%
%% Installation batch file for the RMAA LaTeX package
%%
\def\batchfile{rmaa.ins}
\input docstrip.tex
\preamble

PLEASE DO NOT EDIT THIS FILE! It is generated automatically from
the file rmaa.dtx using the accompanying Makefile and rmaa.ins 
Any hand-edited changes will be lost

A LaTeX document class for the RMxAA

Author: William Henney (w.henney@irya.unam.mx)

Latest version available from: http://www.irya.unam.mx/rmaa

Copyright (C) 1999-2019 Instituto de Astronomia, UNAM, Mexico

Permission is hereby granted, free of charge, to any person obtaining
a copy of this software and associated documentation files (the
"Software"), to deal in the Software without restriction, including
without limitation the rights to use, copy, modify, merge, publish,
distribute, sublicense, and/or sell copies of the Software, and to
permit persons to whom the Software is furnished to do so, subject to
the following conditions:

The above copyright notice and this permission notice shall be
included in all copies or substantial portions of the Software.

THE SOFTWARE IS PROVIDED "AS IS", WITHOUT WARRANTY OF ANY KIND,
EXPRESS OR IMPLIED, INCLUDING BUT NOT LIMITED TO THE WARRANTIES OF
MERCHANTABILITY, FITNESS FOR A PARTICULAR PURPOSE AND
NONINFRINGEMENT. IN NO EVENT SHALL THE AUTHORS OR COPYRIGHT HOLDERS BE
LIABLE FOR ANY CLAIM, DAMAGES OR OTHER LIABILITY, WHETHER IN AN ACTION
OF CONTRACT, TORT OR OTHERWISE, ARISING FROM, OUT OF OR IN CONNECTION
WITH THE SOFTWARE OR THE USE OR OTHER DEALINGS IN THE SOFTWARE.
\endpreamble
%    \end{macrocode}
% Clobber files without asking
%    \begin{macrocode}
\askforoverwritefalse
%    \end{macrocode}
% Make the \LaTeX{} macro files (class and styles)
%    \begin{macrocode}
\generate{% 
  \file{rmaa.cls}{\from{rmaa.dtx}{package}}
  \file{rmaa.ins}{\from{rmaa.dtx}{ins}}
  \file{rmaacite.sty}{\from{rmaa.dtx}{rmaacite}}
  \file{adshtml.sty}{\from{rmaa.dtx}{adshtml-style}}
  \file{ADS.tex}{\from{rmaa.dtx}{ads-container}}
  }
%    \end{macrocode}
% Make the other non-\LaTeX{} files. These are done separately since
% they can't use `\%\%' as a comment character for the
% preamble/postamble (and Perl files can't have the postamble I don't
% think). 
%    \begin{macrocode}
\generate{\nopreamble\nopostamble% 
  \file{Makefile}{\from{rmaa.dtx}{makefile}}
  \file{adshtml.perl}{\from{rmaa.dtx}{adshtml-perl}}
  \file{makeads.sh}{\from{rmaa.dtx}{ads-shell-script}}
  \file{sample.latex2html-init}{\from{rmaa.dtx}{latex2html-init}}
  }

%% Local Variables: 
%% mode: LaTeX
%% mode: Font-lock
%% End:
%</ins>
%    \end{macrocode}

% \subsection{Extended citation syntax}
% Modified from a file called mnrascite.sty
% I no longer remember where I stole the original version of this from
% 
% Alterations to the cite commands to make them do more sensible things.
% N.B. the bibliography style should produce entries which have labels where
% the date is delimited by angle brackets '<>' for these macros to work
% properly. They will not give very sensible error messages if you do not do 
% this properly.
% 
% A new command |\scite{key}| is defined which produces a citation which is 
% suitable for insertion in a sentence. i.e., it does not put brackets around 
% the whole citation, only the date. Also |\pcite{key}| produce output with no
% parentheses whatsoever. 
% 
% Should redefine |\@bibitem| really to have any hope of compatibility
% the command |\lbibcite| is defined to deal with the citations.
%    \begin{macrocode}
%<*rmaacite>
\def\bibitem[#1<#2>]#3{\par\hangindent=1.5em\hangafter=1
  \if@filesw
  {\def\protect##1{\string ##1\space}
    \immediate\write\@auxout{\string\lbibcite{#3}{#1}{#2}}}
  \fi}
\@redefinebibitem % what it says
%    \end{macrocode}
% Note that the |\lbibcite| defines two types of citation one contains the names,
% the other the year
%    \begin{macrocode}
\def\lbibcite#1#2#3{\global\@namedef{b@#1}{#2}
                    \global\@namedef{by@#1}{#3}}

\def\@bcite#1#2{({#1\if@tempswa , #2\fi})}
\def\@pcite#1#2{#1\if@tempswa , #2\fi}
%    \end{macrocode}
% This declaration contains what is put between the name and the date
% \changes{3.11}{2003/01/31}{Moved the \cs{articlename} to after the key for consistency
% with the new natbib stuff.}
%    \begin{macrocode}
\def\@brcite{\ (}
\def\@ebrcite{)}
\def\@fmttop#1{\@fmtcite[\csname #1\endcsname]}
\def\@citex[#1]#2{%
%%  \if@filesw\immediate\write\@auxout{\string\citation{#2}}\fi
%% WJH added 16 Jan 2002
  \@for\@citeb:=#2\do{%
    \edef\@citeb{\expandafter\@firstofone\@citeb:\articlename}%
    \if@filesw\immediate\write\@auxout{\string\citation{\@citeb}}\fi}%
%% to here
  \def\@citea{}%
  \@cite{\@for\@citeb:=#2\do
    {%
      \edef\@citeb{\expandafter\@firstofone\@citeb:\articlename}%
      \@citea\def\@citea{;\ }\@ifundefined
      {b@\@citeb}{{\bf ?}\@warning {Citation `\@citeb' on page \thepage \space undefined}}%
      {\csname b@\@citeb\endcsname\@citebro\csname by@\@citeb\endcsname\@citebrc}}}%
  {#1}}

\def\cite{\@ifnextchar [{\let\@citebro=\space
                          \let\@citebrc=\relax
                          \let\@cite=\@bcite\@tempswatrue\@citex}
                        {\let\@citebro=\space
                          \let\@citebrc=\relax
                          \let\@cite=\@bcite\@tempswafalse\@citex[]}}
\def\pcite{\@ifnextchar [{\let\@citebro=\space
                          \let\@citebrc=\relax
                          \let\@cite=\@pcite\@tempswatrue\@citex}
                        {\let\@citebro=\space
                          \let\@citebrc=\relax
                          \let\@cite=\@pcite\@tempswafalse\@citex[]}}
\def\scite{\@ifnextchar [{\let\@citebro=\@brcite
                          \let\@citebrc=\@ebrcite
                          \let\@cite=\@pcite\@tempswatrue\@citex}
                        {\let\@citebro=\@brcite
                          \let\@citebrc=\@ebrcite
                          \let\@cite=\@pcite\@tempswafalse\@citex[]}}

\renewenvironment{thebibliography}
  {\section*{REFERENCES}\vspace*{\afterbiblioskip}\small
    \if@book\RereadBiblio\fi
    \sloppy\clubpenalty4000\widowpenalty4000
    \sfcode`\.=1000\relax
    \setlength\parindent{\z@}\setlength\parskip{\z@}}
  {\relax\normalsize\@outputaddressesifnecessary} 
%</rmaacite>
%    \end{macrocode}

% \subsection{Auxilliary files for generating ADS data}

% The ADS data are written by \LaTeX{} to a |.ads| file. This file
% then has to be included in the following file (|ADS.tex|),
% which is run through \LaTeX{}2HTML (see below). 

%    \begin{macrocode}
%<*ads-container>
%%
%% Instructions:
%% Edit the \input{} line below to reflect the name of your file. 
\documentclass{article}
\usepackage[latin1]{inputenc}
\usepackage{adshtml}
\newcommand{\ADSfield}[2]
{\%#1  #2\\}
%%\newcommand{\ADSvfield}[2]
%%{\%#1 \begin{verbatim}#2\end{verbatim}}
\newenvironment{ADSarticle}{}{\par\bigskip}
\renewcommand{\bibitem}[2][]{\par}
\setlength{\parindent}{0pt}
\renewcommand{\ThisVolume}[1]{%  
  2001, \RMAASC, 12, \oldpageref{#1:RMAAFirstPage}}  
 

\begin{document}
\makeatletter
\input{rmbooksamp.ads}
\end{document}
%</ads-container>
%    \end{macrocode}

% This is the |adshtml.sty| style file used by |ADS.tex|. It just
% defines a few 
% macros that might be used in abstracts, etc (really it should define
% them all). This now gets all its code from
% Sec~\ref{sec:user-commands}. 
%    \begin{macrocode}
%<*adshtml-style>
%% Nothing else needs to go here now. 
%</adshtml-style>
%    \end{macrocode}
% The init file for latex2html.
%    \begin{macrocode}
%<*latex2html-init>
$PDFTOCAIRO = "pdftocairo";
$PDFLATEX = "pdflatex";
$LATEX = "latex";
$MAX_SPLIT_DEPTH = 0;
$HTML_VERSION = "4.0,math";
$NO_SIMPLE_MATH = 1;
$ADDRESS = "<p><address><a href=\"mailto:rmaa\@astro.unam.mx\">rmaa\@astro.unam.mx</address></p>"; 
$INFO = "<p>This document was created using latex2html and the rmaa latex macros</p>"; 
$NO_NAVIGATION = 1;
1; 
%</latex2html-init>
%    \end{macrocode}
% This is the |adshtml.perl| file that is used by \LaTeX{}2HTML. Some
% bizarre \TeX{}y command sometimes show up so we tell it to ignore
% these. All the astronomy macros that come from AAS are processed in
% \TeX{}. That way they get made to images but with an |ALT| tag that
% has the original code, which is what will be picked up by |lynx|,
% with the result that they get passed intact into the final |.txt|
% file. 
%    \begin{macrocode}
%<*adshtml-perl>
package main;

$dum = &do_cmd_makeatletter();

sub do_cmd_RMAA {
    local($_) = @_;
    "RMxAA" . $_;
}

sub do_cmd_RMAASC {
    local($_) = @_;
    "RMxAASC" . $_;
}

sub do_cmd_etal {
    local($_) = @_;
    "et al\." . $_;
}

sub do_cmd_ibidrule {
    local($_) = @_;
    "<img src=\"foo\" alt=\"\ibidrule\">" . $_;
}


sub do_cmd_bibitem {
    local($_) = @_;
    local($key);
    local($opt, $dummy)=&get_next_optional_argument; 
    $key = &missing_braces unless
        ((s/$next_pair_pr_rx/$key = $2; ''/eo)
        ||(s/$next_pair_rx/$key = $2; ''/eo));
    "<p>" . $_; 
}

&ignore_commands (<<_RAW_ARG_CMDS_);
lowercase
adjustfinalcols
textsuperscript # {}
newline # []
thinmuskip
mskip
smallskip
medskip
bigskip
hspace # {}
hspacestar # {}
vspace # {}
vspacestar # {}
mathrm
_at_mathrm
spacefactor
discretionary
unhbox 
voidb_at_x 
penalty 
_at_M
_at_m
null
IeC
_RAW_ARG_CMDS_

&process_commands_in_tex (<<_RAW_ARG_CMDS_);
sun
earth
lesssim
gtrsim
sq
argdeg
arcmin
arcsec
fd
fh
fm
fs
fdg
farcm
farcs
fp
micron
la
ga
case # {} # {}
slantfrac # {} # {}
onehalf
onethird
twothirds
onequarter
threequarters
ubvr
ub
bv
vr
ur
aj
araa
apj
apjl
apjs
ao
apss
aap
aapr
aaps
azh
baas
jrasc
memras
mnras
pra
prb
prc
prd
pre
prl
pasp
pasj
qjras
skytel
solphys
sovast
ssr
zap
nat
iaucirc
aplett
apspr
bain
fcp
gca
grl
jcp
jgr
jqsrt
memsai
nphysa
physrep
physscr
planss
procspie
_RAW_ARG_CMDS_

1;                              # This must be the last line
%</adshtml-perl> 
%    \end{macrocode}

% This is a shell script |makeads.sh| that generates the |_ads.txt| file from
% |_ads.tex| by using \LaTeX{}2HTML and |lynx|.  \changes{4.0}{2019/11/13}{Make sure that lynx writes utf-8 output instead of latin-1.}
% \changes{2.2beta13}{2002/03/12}{Complete rewrite of |makeads.sh|. Now much 
% improved I think}
% \changes{4.0}{2019/11/13}{Fix the at-sign substitution so that it only eliminates \cs{@} and does not mess with email addresses.}

%    \begin{macrocode}
%<*ads-shell-script>
#! /bin/sh
if [ -z $1 ]; then
    echo "Usage: $0 basename"
    exit
fi
BASENAME=$1
ADSBASENAME=${BASENAME}_ads
AUXFILE=$BASENAME.aux
TMPFILE=tmp.$$


echo "Generating $ADSBASENAME.tex"

# For some reason, must double-up \\ to \\\\ but not \ to \\ - odd!
cat > $ADSBASENAME.tex <<EOF
%% PLEASE DO NOT EDIT THIS FILE!  It is generated automatically from
%% the script makeads.sh.  Any hand-edited changes will be lost
%% 
\documentclass{article}
\usepackage[latin1]{inputenc}
\usepackage{html}
\usepackage{adshtml}
\newcommand{\ADSfield}[2]{\%#1  #2\\\\}
\newenvironment{ADSarticle}{}{\par\bigskip}
\setlength{\parindent}{0pt}
\newcounter{istage}
\renewcommand{\ion}[2]{\setcounter{istage}{#2}#1 \Roman{istage}}
\begin{document}
\makeatletter
EOF
echo >> $ADSBASENAME.tex "\input{$BASENAME.ads2}" 
echo >> $ADSBASENAME.tex "\end{document}"



echo "Generating $ADSBASENAME.sed" 

# Find type of journal
grep '\\documentclass\[.*proceedings.*\]' $BASENAME.tex >/dev/null \
     && JOURNAL='RMxAC' \
     || JOURNAL='RMxAA'

# Find the year and volume number
YEAR=`grep '\\SetYear' $BASENAME.tex | sed 's/.*\\SetYear{\(.*\)}.*/\1/'`
VOLUME=`grep '\\SetVolume' $BASENAME.tex | sed 's/.*\\SetVolume{\(.*\)}.*/\1/'`

# Generate an auxilliary sed script 
cat > $TMPFILE <<EOF
s|\\\\newlabel{\(.*\):RMAAFirstPage}{{.*}{\(.*\)}}|s/\\\\\\\\ThisVolume{\1}/$YEAR, RMxAC, $VOLUME, \2/g|g
EOF

# Now go through the .aux file, looking for \newlabels and use the sed 
# script just generated in order to make another sed script, one that 
# will replace \ThisVolume{LABEL} with the correct reference
grep '^\\newlabel{.*:RMAAFirstPage}' $AUXFILE \
     | sed -f $TMPFILE > $ADSBASENAME.sed && rm $TMPFILE

# Apend a further sed command to get rid of \@, since redefining 
# it in _ads.tex didn't work
echo >> $ADSBASENAME.sed 's/\\\\@//g'

echo "Generating $BASENAME.ads2 from $ADSBASENAME.sed and $BASENAME.ads" 

# Apply this to the .ads file, generating the .ads2 file, which is included
# by the _ads.tex file.
sed -f $ADSBASENAME.sed $BASENAME.ads > $BASENAME.ads2

echo "Generating $ADSBASENAME.txt from $ADSBASENAME.tex and $BASENAME.ads2"

# Now run latex2html on the _ads.tex file
latex2html -no_navigation -info 0 -address 0 -html_version 4.0,math $ADSBASENAME

# Then, convert to utf8 with lynx and tidy it up a bit
lynx --justify=off --width=100 --dump --nolist --display_charset=utf-8 \
     $ADSBASENAME/$ADSBASENAME.html | cut -c4- \
     | sed 's/ibidrule/\\ibidrule/g' > $ADSBASENAME.txt
%</ads-shell-script>
%    \end{macrocode}

% \section{Epilogue}
% Err, that's it \dots

\endinput


% \endinput
% Local Variables:
% mode: doctex
% TeX-master: t
% End:
"

rmaa.ind: rmaa.idx
        @echo "REGENERATING INDEX"
        makeindex -s gind.ist rmaa

rmaa.gls: rmaa.glo
        @echo "REGENERATING CHANGELOG"
        makeindex -s gglo.ist -o rmaa.gls rmaa.glo

rmaa.pdf: rmaa.aux rmaa.ind rmaa.gls
        @echo "REGENERATING .pdf FILE"
        pdflatex "\nonstopmode% \def\fileversion{4.0} \def\filedate{2019/11/12}
% \MakeShortVerb{\|}
% \title{A \LaTeX{} document class for the Revista Mexicana de
% Astronom\'\i a y Astrof\'\i sica\thanks{This file has version
% \fileversion{} last revised \filedate{}.}}
% \author{William J. Henney\thanks{w.henney@irya.unam.mx}}
% \date{\filedate}
% \maketitle
% \paragraph{Abstract}
% This document describes the \texttt{rmaa} \LaTeX{}
% document class for the journal ``Revista Mexicana de
% Astronom\'\i a y Astrof\'\i sica''. The macros can be used for 
% \begin{enumerate}
% \item Preparing author manuscripts for the journal (including
% conference proceedings to be published in the ``Serie de
% Conferencias'')
% \item Typesetting an entire volume of the journal by inputting the
% individual article files from a container file. Facilities are
% included for automatic generation of TOC and Author Index. 
% \end{enumerate}
% 
% \tableofcontents
% \section{The documentation driver file}
% This is the part that generates the documentation that you may or
% may not be reading, which is done with the help of the Makefile
% discussed in section~\ref{sec:driver_make}. 
%    \begin{macrocode}
%<*driver>
\documentclass{ltxdoc}
\usepackage{hyperref}
\usepackage{palatino}
\EnableCrossrefs \CodelineIndex \RecordChanges
\setlength{\parindent}{0pt}
\def\MacroFont{\footnotesize\tt}
\begin{document} \DocInput{rmaa.dtx} \PrintIndex \PrintChanges \end{document}
%</driver>
%    \end{macrocode}

% \section{The implementation}
% 
% \changes{3.24}{2007/09/14}{Added user access to the macro version
% and date via \cs{rmVersion} and \cs{rmDate}}
% 
%    \begin{macrocode}
%<*package>
\newcommand{\rmaa@version}{3.26}
\newcommand{\rmaa@date}{2011/08/15}
\newcommand\rmVersion{\rmaa@version}
\newcommand\rmDate{\rmaa@date}
\NeedsTeXFormat{LaTeX2e}[1995/06/01]
%    \end{macrocode}
% \changes{4.0}{2019/11/02}{Change abstract to appear before
% resumen. Change RevMexAA to RMxAA for consistency with ADS.}
%
% \changes{3.22}{2007/03/22}{Allow proceedings articles to have
%   received/accepted dates} \changes{3.21}{2006/05/15}{Changes for
%   Silvia - abstracts are now shipped out to ADS.}
% \changes{3.19}{2006/01/19}{Changes for Silvia - remove adresses from
%   end, new abstract style} \changes{2.2beta}{2001/10/30}{ Fixed bug
%   whereby twocolumn option was not passed through to book
%   class. Finally sorted out the ``label on last page'' issue. }
% \changes{2.1}{2001/01/23}{ Fixed some inconsistencies in
%   \cs{nocite} and \cs{bibitem} uniqification Fixed the definition of
%   \cs{ion} so it will work in math mode and in contexts where the
%   typesize is different from normal (e.g., title, fig captions).
%   Modified the definition of \cs{@caption}, so we don't write any
%   list-of-figures data (we never use it anyway, and it just causes
%   problems with fragile commands in fig captions). }
% \changes{2.0}{2000/11/30}{ Added `newstyle' option, which should be
%   used in all new proceedings volumes (wider abstract, smaller
%   typeface in fig captions, etc.) Added commands for formatting
%   table captions. New demonstation document rmpostsamp.tex for
%   1-page posters. This will be the last release under a PRI
%   government (for a while...)}  \changes{1.9}{2000/06/06}{
%   Corrected position of page number in footer of articles that open
%   on even pages. Ensured that first article in each session opens on
%   an odd page, as do list of abstracts and author index. Added extra
%   space in TOC before the backmatter. Added command for editors,
%   \cs{BoldSessionTitles}, that makes produces bold titles for
%   session in the TOC.}  \changes{1.8}{2000/05/16}{ Fixed bug in
%   last page label with \cs{suppressfulladdresses} }
% \changes{1.7}{2000/05/11}{ Fixed more little things. Page labels
%   now work, even if you have no bibliography. You can output the
%   full addresses by hand with the \cs{outputfulladdresses}
%   command. Editors' Guide updated but still not finished. Revamped
%   splitbook.sh now produces PS file of each article for web page and
%   also an EPS file of each page for sending to the printers.}
% \changes{1.6}{2000/04/15}{ More cleaning up: full stops after
%   titles in index, running headers and footers in frontmatter,
%   fontsize in title of index, etc} \changes{1.5}{2000/04/15}{
%   Eliminated superscripts from titles in TOC and index}
% \changes{1.4}{2000/04/14}{ Added copyright and license info (GPL)}
% \changes{1.3}{2000/03/01}{ Corrected \cs{figcaption} to do the
%   caption formatting right Corrected font size of sec headings in
%   front matter} \changes{1.2}{2000/02/16}{ Added \cs{plate}
%   command} \changes{1.1}{2000/02/14}{ Added manuscript option,
%   corrected bugs} \changes{1.0}{2000/02/08}{ Now does a whole book
%   - lots of stuff tidied up} \changes{0.7}{1999/10/04}{ Increased
%   spacing between sections in single column styles. Suppressed
%   floats at bottom of pages (didn't work though) }
% \changes{0.6}{1999/09/24}{ Tweaked \cs{adjustfinalcols} (still not
%   perfect) Added proceedings option Added onecolumn/twocolumn
%   options (orthogonal to journal/proceedings) }
% \changes{0.5}{1999/09/23}{ Added \cs{fulladdresses} and
%   \cs{adjustfinalcols} (just about working)}
% \changes{0.4}{1999/09/22}{ Minor stylistic changes}
% \changes{0.3}{1999/06/28}{ Better mechanism for alternative
%   affiliations Added `guide' option for user guide}
% \changes{0.2}{1999/05/27}{ All basic functionality present, I
%   think} \changes{0.1}{1999/05/26}{ Just about working - still
%   missing Fig/Tab/Refs} \changes{0.0}{1999/05/21}{ So bored today,
%   decided to rewrite RMAA LaTeX macros from scratch}
%    \begin{macrocode}
\ProvidesClass{rmaa}[\rmaa@date\space v\rmaa@version\space RMxAA article
class]
%    \end{macrocode}
% \subsection{Option Handling}
% \label{sec:options}
% Declarations that are needed in options
% \changes{3.05}{2003/10/22}{Added natbib option}
% \changes{3.06}{2003/01/03}{Natbib option turned on permanently}
% \changes{3.07}{2003/01/11}{Added proof option}
% \changes{3.14}{2003/01/11}{Option to turn off natbib}
% \changes{3.18}{2003/10/22}{trying out doctex mode} 
% \changes{3.23}{2007/09/09}{Manuscript is now on by default.}
% \changes{3.26}{2011/08/15}{New option for review articles}
%    \begin{macrocode}
\newif\if@reviewarticle
\@reviewarticlefalse
\newif\if@proof
\@prooffalse
\newif\if@journal
\@journaltrue
\newif\if@printdates
\@printdatesfalse
\newif\if@proceedings
\@proceedingsfalse
\newif\if@manuscript
\@manuscripttrue
\newif\if@bigfont
\@bigfonttrue
\newif\if@guide
\@guidefalse
\newif\if@preprint
\@preprintfalse
\newif\if@rmtwocol
\@rmtwocoltrue
\newif\if@biblist
\@biblistfalse
\newif\if@book
\@bookfalse
\newif\if@natbib
\@natbibtrue
\newif\if@rmcite
\@rmcitefalse
\newif\if@titlevisible
\@titlevisibletrue
\newif\if@abstractpresent
\@abstractpresentfalse
\newif\if@printbanner
\@printbannerfalse
\newif\if@suppressfulladdresses
\@suppressfulladdressesfalse
\newif\if@newstyle
\@newstylefalse
\newif\if@debug
\@debugfalse
\newif\if@boldkeywords
\@boldkeywordsfalse
%    \end{macrocode}
% The option handling---mainly setting various booleans.
% \changes{2.2beta11}{2001/11/26}{Made twocol and
% newstyle the default for proceedings.}
% \changes{3.23}{2007/09/09}{Manuscript is now 10 pt instead of 12pt}
% \changes{3.23}{2007/09/10}{Full addresses are now automatically
% suppressed for proceedings style.}
% \changes{3.24}{2007/09/14}{New option \texttt{boldkeywords} that
% does the keywords in the old ugly way. Now off by default.}
%    \begin{macrocode}
\DeclareOption{manuscript}{\typeout{Manuscript for referee} 
  \@manuscripttrue\@journalfalse\@rmtwocolfalse
  \PassOptionsToClass{10pt}{article}}
\DeclareOption{bigfont}{\typeout{Big font manuscript requested by Gary
  Ferland} 
  \@manuscripttrue\@journalfalse\@rmtwocolfalse\@bigfonttrue
  \PassOptionsToClass{11pt}{article}}
\DeclareOption{journal}{\typeout{Main journal option}%
  \@journaltrue\@rmtwocoltrue\@printdatestrue\@manuscriptfalse
  \PassOptionsToClass{10pt}{article}}
\DeclareOption{guide}{\typeout{User guide option}%
  \@guidetrue\@journalfalse\@rmtwocoltrue\@manuscriptfalse
  \PassOptionsToClass{10pt}{article}}
\DeclareOption{preprint}{\typeout{Preprint option}%
  \@preprinttrue\@manuscriptfalse\PassOptionsToClass{10pt}{article}}
\DeclareOption{proof}{\typeout{Author proof option}%
  \@prooftrue\@manuscriptfalse}
\DeclareOption{proceedings}{\typeout{Proceedings option}%
  \@proceedingstrue\@journalfalse\@rmtwocoltrue\@newstyletrue
  \@suppressfulladdressestrue
  \@printdatesfalse\@manuscriptfalse\PassOptionsToClass{10pt}{article}} 
\DeclareOption{printdates}{\typeout{Force printing of received/accepted dates}%
  \@printdatestrue} 
\DeclareOption{book}{\typeout{Whole issue option}%
  \@booktrue} 
\DeclareOption{nonatbib}{\typeout{Turn off natbib support (required
    for compatibility with old rmaacite package).}%
  \@natbibfalse} 
\DeclareOption{newstyle}{\typeout{Use new layout style}%
  \@newstyletrue}
\DeclareOption{oldstyle}{\typeout{Use old layout style}%
  \@newstylefalse}
\DeclareOption{debug}{\typeout{Loads of annoying messages}%
  \@debugtrue}
%    \end{macrocode}
% These two need to be declared explicitly so that |\ExecuteOptions|
% actually does something.
%    \begin{macrocode}
\DeclareOption{twocolumn}{\@rmtwocoltrue}
\DeclareOption{onecolumn}{\@rmtwocolfalse}
\DeclareOption{twoside}{\PassOptionsToClass{\CurrentOption}{article}}
\DeclareOption{boldkeywords}{\@boldkeywordstrue}
%    \end{macrocode}
% Bold, upper-case keywords used to be the default. Now they must be
% turned on explicitly with the \texttt{boldkeywords} option. 
%    \begin{macrocode}
\DeclareOption{boldkeywords}{\@boldkeywordstrue}
%    \end{macrocode}
% Pass any other option along to the \texttt{article} class.
%    \begin{macrocode}
\DeclareOption*{\PassOptionsToClass{\CurrentOption}{article}}
%    \end{macrocode}
% Default options are to be two-sided and based on the standard
% \LaTeX{} \texttt{article} class,
% unless we are processing the entire book, in which case we base
% ourselves on the standard \texttt{book} class. Default is to use
% manuscript option unless otherwise specified. 
%    \begin{macrocode}
\ExecuteOptions{twoside,manuscript} \ProcessOptions\relax
\if@book
\if@rmtwocol\LoadClass[twocolumn]{book}\else\LoadClass{book}\fi
\newcommand{\abstract}{\relax}
\newcommand{\@maketitle}{\relax}
\else
\if@rmtwocol\LoadClass[twocolumn]{article}\else\LoadClass{article}\fi
\fi
%    \end{macrocode}
% Attempt to use natbib for citations 02 Jan 2003. 
% \changes{3.05}{2003/01/02}{If the natbib option is used, then set
% things up so that natbib's version of the bibliography is identical
% to the standard one.}
% \changes{3.09}{2003/01/23}{Explicitly require a recent version of
% natbib, so that we get a decent error message if the package is too
% old.} 
% \changes{3.23}{2007/09/09}{Now allow Bibtex to work properly} 
% \changes{3.25}{2009/10/01}{Change \cs{newcommand} tp \cs{def} when setting \cs{bibfont}. This allows it to work with all versions of natbib.} 
% \changes{3.25}{2010/02/06}{Updated copyright year.} 
% \changes{3.25}{2010/02/06}{Updated email and web addresses to \texttt{crya.unam.mx}.} 
%    \begin{macrocode}
\if@natbib
\RequirePackage{natbib}[1999/02/25]
\bibpunct{(}{)}{;}{a}{}{,}
\def\bibfont{\small}
\renewcommand\bibsection{\section*{REFERENCES}}
\setlength\bibsep{0pt}
\setlength\bibhang{1.5em}
\bibliographystyle{rmaa}
\fi
%    \end{macrocode}
% We want the \texttt{graphicx} package for the extended syntax of
% |\includegraphics|. We want the \texttt{ifthen} package for using in
% these macros (although we don't much yet). We explicitly load
% \texttt{babel} since some \LaTeX{} distributions messed up the
% hyphenation if we didn't.  
%    \begin{macrocode}
\RequirePackage{graphicx}
\RequirePackage{ifthen}
\IfFileExists{babel.sty}{%
%   \let\@restoreindent=\@afterindentfalse
  \RequirePackage[english]{babel}
%   \@restoreindent
  %% reinstate the definition of \@roman so as not to break .idx files
  \@ifundefined{org@@roman}{\relax}{\renewcommand{\@roman}[1]{\org@@roman{##1}}}
  }{%
  \typeout{Package babel.sty not found - don't worry about this}}
%    \end{macrocode}
% Try to load the \texttt{dcolumn} and \texttt{booktabs} packages if
% they can be found. These can produce nicer-looking tables. 
%    \begin{macrocode}
\IfFileExists{dcolumn.sty}{%
  \RequirePackage{dcolumn}
  }{%
  \typeout{Package dcolumn.sty not found - you will not be able to
    align columns on decimal points in tables}}
\IfFileExists{booktabs.sty}{%
  \RequirePackage{booktabs}
  }{%
  \typeout{Package booktabs.sty not found - your tables will not look
    as pretty as they ought too.}
  }
\IfFileExists{makeidx.sty}{%
  \RequirePackage{makeidx}
}{%
  \typeout{Package makeidx.sty not found - Author Index may not work.}
  }
%% commented out for moment since requires more modern version of LaTeX
%% \RequirePackage[latin1]{inputenc}
%    \end{macrocode}
% \changes{4.0}{2019/11/13}{Make \texttt{hyperref} and
% \texttt{etoolbox} into hard requirements, since we now use them
% extensively.}  We use hyperref to get internal and external links.
% But I don't like the colored boxes so I turn those off.  We use
% etoolbox for patching and robustifying at the moment.  In the
% future, I want to port a lot of other stuff over to it.  For
% instance, replace the ifthen package and the TeX ifs too.
%
% We target the latest version (1.7) of PDF, rather than the default
% of 1.5.  This reduces the chance of getting warnings from included
% graphics.  This is not strictly related to hyperref, but we put it
% here anyway.
%    \begin{macrocode}
\RequirePackage{etoolbox}
\RequirePackage{hyperref}
\hypersetup{hidelinks}
\pdfminorversion=7
%    \end{macrocode}
% We patch the \cs{@combinedblfloats} following a suggestion by Heiko
% Oberdiek at \url{https://tex.stackexchange.com/a/249743/3843}, which
% avoids problems when hyperlinks are split across
% pages. \changes{4.0}{2019/11/12}{Fix for issue where breaking
% hyperlink across pages gave an error.}
%    \begin{macrocode}
\urlstyle{same}
\newcommand\email{\begingroup \urlstyle{rm}\Url}
\patchcmd\@combinedblfloats{\box\@outputbox}{\unvbox\@outputbox}{}{%
  \errmessage{\noexpand\@combinedblfloats could not be patched}%
}%
%    \end{macrocode}
% \subsection{Other setup stuff}\label{sec:setup}
% Some declarations: 
%    \begin{macrocode}
\newif\if@needaddresses
\@needaddressestrue
\newif\if@frontmatter
\@frontmatterfalse
\newif\if@inpreamble
\@inpreambletrue
\newif\if@firstauthor
\@firstauthortrue
\newif\if@firstaffil
\@firstaffiltrue
\newif\if@firstaltaffil
\@firstaltaffiltrue
\newif\if@firstkwd
\@firstkwdtrue
\newif\if@noexpandADS
\@noexpandADSfalse
\newif\if@sessiontitlepages
\@sessiontitlepagesfalse
%    \end{macrocode}
% \begin{macro}{\@rmaaEmpty}
% \begin{macro}{\@rmaaArg}
% These are used in various commands to check if arg was left blank
%    \begin{macrocode}
\newcommand{\@rmaaEmpty}{}
\newcommand{\@rmaaArg}{}
%    \end{macrocode}  
% \end{macro}
% \end{macro}
%
% \subsection{Compatibility with old \LaTeX{} versions}
% \label{sec:compat}
% This package is based on the 1998/12/01 release of \LaTeX{}.
% Older releases may not define all the commands that we need.
% Luckily, we can use |\providecommand| to define these if necessary. 
% The following commands are not defined in the 1994/06/01 release. 
%
% \begin{macro}{\MakeUpperCase}
% \begin{macro}{\textsuperscript}
% \begin{macro}{\textbackslash}
% User commands:
%    \begin{macrocode}
\providecommand{\MakeUpperCase}[1]{\uppercase{#1}}
\providecommand{\textsuperscript}[1]{\ensuremath{^\textrm{#1}}}
\providecommand{\textbackslash}{\ensuremath{\backslash}}
%    \end{macrocode}  
% \end{macro}
% \end{macro}
% \end{macro}
% \begin{macro}{\protected@xdef}
% \begin{macro}{\restore@protect}
% \begin{macro}{\@unexpandable@protect}
% \begin{macro}{\hb@xt@}
% \LaTeX{} internal commands:
%    \begin{macrocode}
\providecommand{\protected@xdef}{%
        \let\@@protect\protect
        \let\protect\@unexpandable@protect
        \afterassignment\restore@protect
        \xdef}
\providecommand{\restore@protect}{\let\protect\@@protect}
\providecommand{\@unexpandable@protect}{\noexpand\protect\noexpand}
\providecommand{\hb@xt@}{\hbox to}
%    \end{macrocode}
% \end{macro}
% \end{macro}
% \end{macro}
% \end{macro}
%
% \subsection{Page layout}\label{sec:page-layout}
%
% The dimensions of the text were defined by Elizabeth. Text width is
% 169\,mm with equal margins on the two sides: 
%    \begin{macrocode}
\setlength{\textwidth}{169mm}
\setlength{\@tempdima}{\paperwidth}
\addtolength{\@tempdima}{-\textwidth}
\setlength{\oddsidemargin}{0.5\@tempdima}
\addtolength{\oddsidemargin}{-1in}
\setlength{\evensidemargin}{\oddsidemargin}
%    \end{macrocode}
% Text height is 225\,mm, top margin is 14\,mm. 
%    \begin{macrocode}
\setlength{\textheight}{225mm} 
\setlength{\topmargin}{-1in} 
\addtolength{\topmargin}{14mm}
%    \end{macrocode}
% Space between columns in two-column styles, header, footer,
% paragraph indentation:
%    \begin{macrocode}
\setlength{\columnsep}{9.5mm}
\setlength{\headsep}{0.7cm}
\setlength{\footskip}{0.75cm}
\setlength{\parindent}{1.5em}
%    \end{macrocode}
% 
% \begin{macro}{\@set@scalable@length}
% Try out a new scheme for dealing with various spacing dimensions and
% allowing the user to easily modify them. The command
% \begin{quote}
% |\set@scalable@length|\marg{name}\marg{defaultlength}\marg{stretch}\marg{shrink}
% \end{quote}
% does a 
% |\setlength|\marg{\cs{name}}\marg{defaultlength plus stretch minus shrink}
% and also defines the user command
% \textbackslash\meta{name}|Scale|\marg{factor} that can 
% be used to rescale length by \meta{factor}. Additionally, now
% defines the package command
% \textbackslash\meta{name}|@set@scale|\marg{FACTOR} 
% that also rescales the length by \meta{FACTOR}, so that the length
% becomes \meta{factor} $\times$ \meta{FACTOR} $\times$
% \meta{defaultlength}.   
%    \begin{macrocode}
\newcommand\@set@scalable@length[4]{%
  %% arg #1: NAME of length without leading `\'
  %%
  %% arg #2: default VALUE  
  \expandafter\newlength\csname @#1@default\endcsname
  \expandafter\setlength\csname @#1@default\endcsname{#2}%
  \expandafter\newlength\csname @@#1@@default\endcsname
  \expandafter\setlength\csname @@#1@@default\endcsname{#2}%
  %% arg #3: default STRETCH  
  \expandafter\newlength\csname @#1@stretch\endcsname
  \expandafter\setlength\csname @#1@stretch\endcsname{#3}%
  \expandafter\newlength\csname @@#1@@stretch\endcsname
  \expandafter\setlength\csname @@#1@@stretch\endcsname{#3}%
  %% arg #4: default SHRINK  
  \expandafter\newlength\csname @#1@shrink\endcsname
  \expandafter\setlength\csname @#1@shrink\endcsname{#4}%
  \expandafter\newlength\csname @@#1@@shrink\endcsname
  \expandafter\setlength\csname @@#1@@shrink\endcsname{#4}%
  %%
  %% User-defined scale \@NAME@scale is unity by default
  \expandafter\gdef\csname @#1@scale\endcsname{1}%
  %% Package-defined scale \@NAME@SCALE is unity by default
  \expandafter\gdef\csname @#1@SCALE\endcsname{1}%
  %%
  %% internal command \@reset@NAME should expand to 
  %% \setlength\@@NAME@@default{\@NAME@SCALE\@NAME@default}
  %% \setlength\NAME{\@NAME@scale\@@NAME@@default}
  %% in other words: scale x SCALE x default 
  \expandafter\gdef\csname @reset@#1\endcsname{%
    \expandafter\setlength\csname @@#1@@default\endcsname{%
      \csname @#1@SCALE\endcsname \csname @#1@default\endcsname
      }
    \expandafter\setlength\csname @@#1@@stretch\endcsname{%
      \csname @#1@SCALE\endcsname \csname @#1@stretch\endcsname
      }
    \expandafter\setlength\csname @@#1@@shrink\endcsname{%
      \csname @#1@SCALE\endcsname \csname @#1@shrink\endcsname
      }
    \expandafter\setlength\csname #1\endcsname{%
      \csname @#1@scale\endcsname \csname @@#1@@default\endcsname
      \@plus
      \csname @#1@scale\endcsname \csname @@#1@@stretch\endcsname
      \@minus
      \csname @#1@scale\endcsname \csname @@#1@@shrink\endcsname
      }}
  %%
  %% User command \NAMEScale should expand to 
  %% {\renewcommand\@NAME@scale{#1}\@reset@NAME}
  \expandafter\gdef\csname #1Scale\endcsname##1{%
    \expandafter\gdef\csname @#1@scale\endcsname{##1}%
    \csname @reset@#1\endcsname}
  %%
  %% Package command \NAME@set@SCALE should expand to 
  %% {\renewcommand\@NAME@SCALE{#1}\@reset@NAME}
  \expandafter\gdef\csname #1@set@SCALE\endcsname##1{%
    \expandafter\gdef\csname @#1@SCALE\endcsname{##1}%
    \csname @reset@#1\endcsname}
  %%
  %% initialize at default values with \@reset@NAME
  \csname @reset@#1\endcsname
  }
%    \end{macrocode}
% \end{macro}
% \begin{macro}{\parskip}
% First use of this new infrastructure: set interparagraph space to be
% zero but with a small amount of stretchiness.
%    \begin{macrocode}
%   \ifthenelse{\boolean{@manuscript}}{%
%     \@set@scalable@length{parskip}{0.5\baselineskip}{0.3\p@}{0.1\p@}
%   }{
    \@set@scalable@length{parskip}{0\p@}{0.3\p@}{0\p@}
%   }
%    \end{macrocode}
% \end{macro}
% 
% \begin{macro}{\abovesectionskip}
% \changes{4.0}{2019/11/12}{Use negative values for skip so that first paragraph after section headings is not indented.  This is something that I have been trying to fix for years!  It turns out that the anser was in the documentation for the standard command \cs{@startsection.}}
% \begin{macro}{\belowsectionskip}
% \begin{macro}{\aboveparaskip}
% \begin{macro}{\belowparaskip}
% Define some utility rubber lengths for use in
% section~\ref{sec:section-commands}. 
%    \begin{macrocode}
\newlength{\abovesectionskip}
\newlength{\belowsectionskip}
\newlength{\aboveparaskip}
\newlength{\belowparaskip}
%    \end{macrocode}
% The default values for these lengths are those traditionally used
% for the onecolumn style. 
%    \begin{macrocode}
\@set@scalable@length{abovesectionskip}{-3.5ex}{-4ex}{-2ex}
\@set@scalable@length{belowsectionskip}{2.3ex}{0ex}{2ex}
\@set@scalable@length{aboveparaskip}{3.25ex}{2ex}{0.5ex}
\@set@scalable@length{belowparaskip}{-1em}{\z@}{\z@}
%    \end{macrocode}
% \end{macro}\end{macro}\end{macro}\end{macro}
% We also make scalable some of the standard float parameters. These
% are the same defaults as in the standard article classes, but
% slightly more stretchy. For the twocolumn style, they are currently
% halved, see below. Note that I have also lumped |\abovecaptionskip|
% in with these, since it plays a similar role. There is also a
% |\belowcaptionskip| but it is set to zero in the standard classes. 
%    \begin{macrocode}
\@set@scalable@length{floatsep}{12\p@}{4\p@}{4\p@}
\@set@scalable@length{textfloatsep}{20\p@}{4\p@}{8\p@}
\@set@scalable@length{intextsep}{12\p@}{4\p@}{4\p@}
\@set@scalable@length{dblfloatsep}{12\p@}{4\p@}{4\p@}
\@set@scalable@length{dbltextfloatsep}{20\p@}{4\p@}{8\p@}
\@set@scalable@length{abovecaptionskip}{10\p@}{\z@}{\z@}
%    \end{macrocode}
% \begin{macro}{\@before@title@skip}         
% \begin{macro}{\@before@author@skip}        
% \begin{macro}{\@between@author@skip}       
% \begin{macro}{\@before@abstract@skip}      
% \begin{macro}{\@after@abstracthead@skip}   
% \begin{macro}{\@between@abstract@skip}     
% \begin{macro}{\@before@kwd@skip}           
% \begin{macro}{\@before@dates@skip}           
% Spacing parameters used in |\@maketitle|. These are not changed by
% default except by the user, so they don't need to go into
% |\@package@default@lengths|.  
%    \begin{macrocode}
\newlength{\@before@title@skip}         
\newlength{\@before@author@skip}        
\newlength{\@between@author@skip}       
\newlength{\@before@abstract@skip}      
\newlength{\@after@abstracthead@skip}   
\newlength{\@between@abstract@skip}     
\newlength{\@before@kwd@skip}           
\newlength{\@before@dates@skip}           
\@set@scalable@length{@before@title@skip}{2em}{\z@}{\z@}
\@set@scalable@length{@before@author@skip}{1.5em}{\z@}{\z@}
\@set@scalable@length{@between@author@skip}{0.5em}{\z@}{\z@}
\@set@scalable@length{@before@abstract@skip}{3ex}{1ex}{0.5ex}
\@set@scalable@length{@after@abstracthead@skip}{1ex}{\z@}{\z@}
\@set@scalable@length{@between@abstract@skip}{2ex}{1ex}{0.5ex}
\@set@scalable@length{@before@kwd@skip}{1ex}{0.5ex}{0.2ex}
\@set@scalable@length{@before@dates@skip}{2ex}{0.5ex}{0.2ex}
%    \end{macrocode}
% \end{macro}\end{macro}\end{macro}\end{macro}
% \end{macro}\end{macro}\end{macro}\end{macro}
% 
% Spacing above and below displayed equations also wants to be done
% like this, but it is more complicated since the standard article
% classes make these change with the fontsize. Something to do later,
% maybe. 
% 
%    \begin{macrocode}
%    \end{macrocode}
% \begin{macro}{\@package@default@lengths}
% The value actually used for these lengths is controlled by the
% \meta{SCALE} parameter, which depends on the style option that is in
% effect. This is set using the |\@package@default@lengths| command,
% which must be called every time the style changes (e.g., changing
% from onecolumn to twocolumn). 
%    \begin{macrocode}
\newcommand{\@package@default@lengths}{%
  \ifthenelse{\boolean{@rmtwocol}}{%
    %% twocolumn mode
    \abovesectionskip@set@SCALE{0.5}
    \belowsectionskip@set@SCALE{0.5}
    \aboveparaskip@set@SCALE{0.5}
    \belowparaskip@set@SCALE{1}
    \floatsep@set@SCALE{0.5}
    \textfloatsep@set@SCALE{0.5}
    \intextsep@set@SCALE{0.5}
    \dblfloatsep@set@SCALE{0.5}
    \dbltextfloatsep@set@SCALE{0.5}
    }{%
    %% onecolumn mode
    \abovesectionskip@set@SCALE{1}
    \belowsectionskip@set@SCALE{1}
    \aboveparaskip@set@SCALE{1}
    \belowparaskip@set@SCALE{1}
    \floatsep@set@SCALE{1}
    \textfloatsep@set@SCALE{1}
    \intextsep@set@SCALE{1}
    \dblfloatsep@set@SCALE{1}
    \dbltextfloatsep@set@SCALE{1}
    }%
  }
\@package@default@lengths
%    \end{macrocode}
% \end{macro}
% \begin{macro}{\RescaleLengths}
% \begin{macro}{\RescaleSecLengths}
% \begin{macro}{\RescaleFloatLengths}
% \begin{macro}{\RescaleTitleLengths}
% User commands that will fiddle with all the scalable lengths. 
%    \begin{macrocode}
\newcommand{\RescaleLengths}[1]{%
  \RescaleSecLengths{#1}\RescaleFloatLengths{#1}\RescaleTitleLengths{#1}}
\newcommand{\RescaleSecLengths}[1]{%
%% sections
  \abovesectionskipScale{#1}
  \belowsectionskipScale{#1}
  \aboveparaskipScale{#1}
  }
\newcommand{\RescaleFloatLengths}[1]{%
%% floats
  \floatsepScale{#1}
  \textfloatsepScale{#1}
  \intextsepScale{#1}
  \dblfloatsepScale{#1}
  \dbltextfloatsepScale{#1}
  \abovecaptionskipScale{#1}
  }
\newcommand{\RescaleTitleLengths}[1]{%
%% maketitle
  \@before@title@skipScale{#1}           
  \@before@author@skipScale{#1}        
  \@between@author@skipScale{#1}       
  \@before@abstract@skipScale{#1}      
  \@after@abstracthead@skipScale{#1}   
  \@between@abstract@skipScale{#1}     
  \@before@kwd@skipScale{#1}           
  \@before@dates@skipScale{#1}           
  }
%    \end{macrocode}
% \end{macro}\end{macro}\end{macro}\end{macro}
% \changes{2.2beta14}{2002/03/17}{Finished implementation of new
% infrastructure for scalable lengths. Converted the above/below
% section skips to use it. Also, all the floatsep parameters. Also,
% all the vertical lengths that are used by |\@maketitle|. }

% \begin{macro}{\MSextramargin}
% For the \texttt{manuscript} option we have an extra 0.3\,inch
% vertical margin and 1\,inch horizontal margin for marginal notes (by
% hand or with |\marginpar|. 
%
% \changes{3.23}{2007/09/09}{Reduce baselinestretch in
%   \texttt{manuscript} style to 1.2. No longer use raggedright for
%   manuscript.}
%
% \changes{3.24}{2007/09/14}{Further reduction of baselinestretch in
%   \texttt{manuscript} style to 1.1 since we now have wider margins.}
%
% \changes{3.24}{2007/09/14}{Split \texttt{manuscript}'s extra margin
%   into separate horizontal and vertical parts.}
% 
%    \begin{macrocode}
\newlength{\MSextraHmargin}\newlength{\MSextraVmargin}
\setlength{\MSextraHmargin}{1in} 
\setlength{\MSextraVmargin}{0.3in} 
\if@manuscript
   \renewcommand{\baselinestretch}{1.1}
   \addtolength{\textwidth}{-2\MSextraHmargin}
   \addtolength{\oddsidemargin}{\MSextraHmargin}
   \setlength{\evensidemargin}{\oddsidemargin}
   \addtolength{\textheight}{-2\MSextraVmargin}
   \addtolength{\topmargin}{\MSextraVmargin}
%    \raggedright
\fi
%    \end{macrocode}
% \end{macro}
% Float parameters: don't allow bottom floats by default. 
%    \begin{macrocode}
\renewcommand{\fps@figure}{tp}
\renewcommand{\fps@table}{tp}
%    \end{macrocode}
%
% \begin{macro}{\rmaa@AuthorWidthFrac}
% \begin{macro}{\rmaa@AbstractWidthFrac}
% These parameters control what fraction of the total text width can
% be used by the list of authors and the abstract/resumen,
% respectively.  
%    \begin{macrocode}
\newcommand{\rmaa@AuthorWidthFrac}{1.0}
\newcommand{\rmaa@AbstractWidthFrac}{0.75}
%    \end{macrocode}
% Use a wide abstract to save space in the new style.
%    \begin{macrocode}
\if@newstyle\renewcommand{\rmaa@AbstractWidthFrac}{1.0}\fi
%    \end{macrocode}
% \end{macro}
% \end{macro}
%
% Some things to consider in future versions. Trying to avoid too many
% widows, etc. Apparently, the following will work best with
% \cs{flushbottom}. 
%    \begin{macrocode}
%%\widowpenalty=10000
%%\clubpenalty=10000
%    \end{macrocode}
% Also, should perhaps try getting rid of underfull/overfull box warnings. 
% \changes{3.10}{2003/01/24}{Set large values for \cs{hbadness},
% \cs{vbadness}, \cs{hfuzz}, \cs{vfuzz}. To reduce amount of stuff
% that gets written to the terminal every time you compile something.}
%    \begin{macrocode}
\hbadness=10000
\vbadness=10000
\setlength\hfuzz{\maxdimen}
\setlength\vfuzz{\maxdimen}
%    \end{macrocode}
% \subsection{Various text styles for different contexts}
% \label{sec:text-styles} 
% \begin{macro}{\rmaa@AffilStyle}
% \begin{macro}{\rmaa@PagesStyle}
% \begin{macro}{\rmaa@YearStyle}
% \begin{macro}{\rmaa@VolumeStyle}
% \begin{macro}{\rmaa@VolumenStyle}
% \begin{macro}{\rmaa@JournalNameStyle}
% \begin{macro}{\rmaa@ReceivedStyle}
% \begin{macro}{\rmaa@AbstractHeaderStyle}
% \begin{macro}{\rmaa@JournalNameStyle}
% Style commands for the different fields used by |\maketitle| in
% constructing the title page of the article. These have the form
% |\rmaa@XYZStyle{TEXT}|. I'm not sure these were really a good idea,
% but never mind \dots 
%    \begin{macrocode}
\newcommand{\rmaa@AffilStyle}[1]{{\normalsize\rmfamily#1}}
\newcommand{\rmaa@PagesStyle}[1]{{\normalsize\rmfamily#1}}
\newcommand{\rmaa@YearStyle}[1]{{\normalsize\rmfamily#1}}
\newcommand{\rmaa@VolumeStyle}[1]{{\normalsize\rmfamily\bfseries#1}}
\newcommand{\rmaa@VolumenStyle}[1]{{\Large\sffamily\bfseries#1}}
\newcommand{\rmaa@JournalNameStyle}[1]{{\normalsize\itshape#1}}
\newcommand{\rmaa@ReceivedStyle}[1]{{\normalsize\rmfamily\itshape#1}}
\newcommand{\rmaa@AbstractHeaderStyle}[1]{%
  {\normalsize\rmfamily\MakeUpperCase{#1}}}
\newcommand{\rmaa@TitleStyle}[1]{%
  {\large\bfseries\rmfamily\boldmath\MakeUpperCase{#1}\par}}
%    \end{macrocode}
% \end{macro}
% \end{macro}
% \end{macro}
% \end{macro}
% \end{macro}
% \end{macro}
% \end{macro}
% \end{macro}
% \end{macro}
%
% \begin{macro}{\rmaa@FrontSecStyle}
% Headings of ``front sections'' (see section~\ref{sec:front-sections})
%    \begin{macrocode}
\newcommand{\rmaa@FrontSecStyle}[1]{%
  {\normalsize\bfseries\rmfamily\MakeUpperCase{#1}}}
%    \end{macrocode}
% \end{macro}
%
% \begin{macro}{\rmaa@TOCTitleStyle}
% \begin{macro}{\rmaa@TOCPartStyle}
% \begin{macro}{\rmaa@IndexTitleStyle}
% \begin{macro}{\rmaaTOCAuthorStyle}
% \begin{macro}{\rmaaIndexAuthorStyle}
% Commands to set style of TOC and index entries.
% Rewritten to avoid unnecessary font commands in the .toc and .ind files
%    \begin{macrocode}
\newcommand{\rmaa@TOCTitleStyle}[1]{\MakeUpperCase{#1}}
\newcommand{\rmaa@TOCPartStyle}[1]{\rmaa@TOCTitleStyle{#1}}
\newcommand{\rmaa@IndexTitleStyle}[1]{#1}
\newcommand{\rmaaTOCAuthorStyle}[1]{%
  {\itshape#1}}
\newcommand{\rmaaIndexAuthorStyle}[1]{%
  \protect\mbox{\bfseries#1}}
%    \end{macrocode}
% \end{macro}
% \end{macro}
% \end{macro}
% \end{macro}
% \end{macro}
%
% \begin{macro}{\rmaa@KeywordStyle}
% \begin{macro}{\rmaa@KeywordHeaderStyle}
% The style of the keywords depends on whether we are using the
% \texttt{newstyle} option or not. If we are, then they are small. 
%    \begin{macrocode}
\ifthenelse{\boolean{@newstyle}}{
  \newcommand{\rmaa@KeywordStyle}[1]{%
    {\small%
      \if@boldkeywords 
      \bfseries\MakeUpperCase{\raggedright #1}
      \else
      \raggedright #1
      \fi
    }}
  \newcommand{\rmaa@KeywordHeaderStyle}[1]{%
    {\small\itshape#1}}}
{
  \newcommand{\rmaa@KeywordStyle}[1]{%
    {\normalsize%
      \if@boldkeywords\bfseries\fi\raggedright #1}}
  \newcommand{\rmaa@KeywordHeaderStyle}[1]{%
    {\normalsize\itshape#1}}}
%    \end{macrocode}
% \end{macro}
% \end{macro}
%
% \subsection{Page styles}\label{sec:page-styles}
% \begin{macro}{\ps@rmaa}
% The running headers for normal pages:
%    \begin{macrocode}
\newcommand{\ps@rmaa}{%
%    \end{macrocode}
% nothing in either foot; 
%    \begin{macrocode}
  \renewcommand\@oddfoot{\relax}%
  \renewcommand\@evenfoot{\relax}%
%    \end{macrocode}
% page number and short author list in the even heads; 
%    \begin{macrocode}
  \renewcommand\@evenhead{\rlap{\thepage}\hfil\@shortauthor\hfil}%
%    \end{macrocode}
% short title and page number in the odd heads.
%    \begin{macrocode}
  \renewcommand\@oddhead{\hfil\@shorttitle\hfil\llap{\thepage}}
  }
%    \end{macrocode}
% \begin{macro}{\ps@rmaafirstpage}
% The header of the first page of each
% article:
%    \begin{macrocode}
\newcommand{\ps@rmaafirstpage}{%
  \renewcommand\@oddfoot{\hfil\thepage}%
%    \end{macrocode}
% footer corrected (06 Jun 2000) to have page on outside edge;
%    \begin{macrocode}
  \renewcommand\@evenfoot{\thepage\hfil}%
%    \end{macrocode}
% Header has journal name, volume, and page numbers for a regular
% article, or similar for special cases. 
% \changes{4.0}{2019/11/03}{Change to three-line header with copyright info and DOI.}
%    \begin{macrocode}
\if@preprint
\renewcommand\@oddhead{%
  \if@proceedings 
  \textsc{To appear in ``\@conftitle{} (\rmaa@YearStyle{\@year})''}\hfil
  \textit{RMxAC}
  \else
  \textsc{To appear in} \@journalname\hfil
  \fi}\fi%
\if@guide
\renewcommand\@oddhead{\textsc{Documentation} (\rmaa@YearStyle{\@year})\hfil}
\fi%
\if@manuscript
\renewcommand\@oddhead{Manuscript for \@journalname{} (\rmaa@YearStyle{\@year})\hfil}
\fi%
\if@journal
\renewcommand\@oddhead{%
  \parbox{\textwidth}{
    \@journalname, \rmaa@VolumeStyle{\@volume},
    \rmaa@PagesStyle{\@artpages}\
    (\rmaa@YearStyle{\@year})\hfil\\[1pt]
    {\small\textcopyright{} \@year:
      Instituto de Astronom\'\i{}a,
      Universidad Nacional Aut\'onoma de M\'exico}\hfil\\[1pt]
    {\small\@DOI@url}\hfil
  }%
}\fi%
  \renewcommand\@evenhead{\@oddhead}}
%    \end{macrocode}
% By default, use the |\ps@rmaa| page style.
%    \begin{macrocode}
\pagestyle{rmaa}
%    \end{macrocode}
% \end{macro}
% \end{macro}
%
% \subsection{Preamble commands}\label{sec:preamble} 
% Everything you need for setting the title page, TOC and index
% entries, etc.
% \subsubsection{Commands for the title page}
% \label{sec:title-commands}
% \begin{macro}{\resumen}
% \begin{macro}{\@resumen}
% \begin{macro}{\abstract}
% \begin{macro}{\@abstract}
% User commands for setting the resumen and abstract. 
% These should be used in the document header. The internal commands
% |\@resumen| and |\@abstract| are used to hold the actual text
% pending output via the |\maketitle| command. 
%    \begin{macrocode}
\newif\if@setresumen
\newcommand{\resumen}[1]{%
%%  \renewcommand{\@resumen}{{\selectlanguage{spanish}#1}}%
  \renewcommand{\@resumen}{{\ignorespaces#1}}%
  \@setresumentrue}
\renewcommand{\abstract}[1]{%
  \renewcommand{\@abstract}{{\ignorespaces#1}}
  \renewcommand\ADS@abstract{#1}
%% Not a good idea, since then author macros are not expanded and
%% cause errors later. However, by not doing it we are vulnerable to
%% fragile commands.....
  \if@noexpandADS\@onelevel@sanitize{\ADS@abstract}\fi
  \@abstractpresenttrue
  }
\newcommand{\ADSnoexpand}{\@noexpandADStrue}
%    \end{macrocode}
% \begin{macro}{\boldabstract}
% This is the version for doing 1-page conference proceedings. Still
% save the text for ADS but then write it out immediately as a bold
% paragraph. 
%    \begin{macrocode}
\newcommand{\boldabstract}[1]{%
  \gdef\ADS@abstract{#1}%
  \renewcommand{\@abstract}{#1}%
  \if@noexpandADS\@onelevel@sanitize{\ADS@abstract}\fi%
  \noindent\textbf{\boldmath #1}
  }
%    \end{macrocode}
% \end{macro}
% Default value in each case is to print out a warning if not set
%    \begin{macrocode}
  \newcommand{\@resumen}{%
    \@setresumenfalse \ClassWarningNoLine{rmaa}{No Spanish abstract
      (resumen) specified - if you do not speak Spanish, this will be
      translated by the editors}%
    \textbf{Favor de proporcionar un resumen en espa\~nol. If you are
      unable to translate your abstract into Spanish, the editors will
      do it for you.} }
  \newcommand{\@abstract}{\ClassWarningNoLine{rmaa}{No abstract
      specified}}
%    \end{macrocode}
% \end{macro}
% \end{macro}
% \end{macro}
% \end{macro}
%
% \begin{macro}{\title}
% \begin{macro}{\@rawtitle}
% \begin{macro}{\@TOCtitle}
% \begin{macro}{\rmaaIndexTitle}
% Similarly, |\title| is the preamble command to set the title of the
% article. This is used in various places, typeset in different
% styles, hence the use of the internal |\@rawtitle|, |\@title|,
% |\rmaaIndexTitle|, and |\@TOCtitle| commands, each of which uses the
% relevant style from section~\ref{sec:text-styles}
%    \begin{macrocode}
\newcommand{\@rawtitle}{Raw Title}
\newcommand{\@TOCtitle}{\rmaa@TOCTitleStyle{Title for TOC}}
\newcommand{\rmaaIndexTitle}{\rmaa@IndexTitleStyle{Title for index}.}
\newcommand{\@TOCaftertitle}{}
\renewcommand{\title}[1]{%
%%  {\ADS@sanitize\xdef\@rawtitle{#1}}
  \renewcommand\@rawtitle{\noexpand#1}
  \let\ADS@title=\@rawtitle
  \renewcommand{\@title}{\rmaa@TitleStyle{#1}}%
  \renewcommand{\@TOCtitle}{\rmaa@TOCTitleStyle{#1}\@TOCaftertitle}%
  \renewcommand{\rmaaIndexTitle}{\rmaa@IndexTitleStyle{#1\@TOCaftertitle}.}%
  }
%    \end{macrocode}
% \end{macro}
% \end{macro}
% \end{macro}
% \end{macro}
% 
% All the weird and whacky ways of indicating the author
% affiliations. 
% \begin{macro}{\affil}
% First, the way that sets it directly underneath the authors name.
% Note that |\affil| trumps all other affiliations for the purposes of
% getting into the ADS stuff (we use |\gdef|).   
% \dots 
%    \begin{macrocode}
\newcommand{\affil}[1]{%
  \\[1.5ex]\rmaa@AffilStyle{#1}\\[1.5ex]%
  \if@firstaffil\gdef\ADS@affil{#1}\@firstaffilfalse\fi%
  }
\renewcommand{\and}{and\\[1.5ex]}
%    \end{macrocode}
% \begin{macro}{\altaffil}
% Alternative affiliations are footnoted. 
%    \begin{macrocode}
\newcommand{\altaffil}[1]{%
  \if@firstaltaffil\gdef\ADS@affil{#1}\@firstaltaffilfalse\fi%
  \thanks{#1}}
%    \end{macrocode}
% \begin{macro}{\altaffilmark}
% \begin{macro}{\altaffiltext}
% Alternative-alternative mechanism for affiliations---should not be
% mixed with |\altaffil| (unless you know what you are doing) or there
% will be problems with numbering.
%    \begin{macrocode}
\newcommand{\altaffilmark}[1]{\textsuperscript{#1}}
\newcommand{\altaffiltext}[2]{%
  \if@firstaltaffil\gdef\ADS@affil{#2}\@firstaltaffilfalse\fi
  \stepcounter{footnote}%
  \protected@xdef\@thanks{\@thanks
    \protect\footnotetext[#1]{#2}}} 
%    \end{macrocode}
% \end{macro}
% \end{macro}
% \end{macro}
% \end{macro}
%
% \begin{macro}{\keywords}
% \begin{macro}{\@keywords}
% \begin{macro}{\@keywordheader}
% \label{keywords}Sundry minor preamble commands. All
% self-explanatory, I hope. |\keywords| is the old user interface
% (deprecated). New articles should use |\addkeyword| once for each
% keyword. That way the ADS record comes out right too (comma rather
% than em-dash separator).
% \changes{2.2beta12}{2001/12/07}{\cs{addkeyword} New user interface
%   for keywords.} 
% \changes{2.2beta13}{2002/03/11}{\cs{addkeyword} Improved behaviour
%   with .ads file. Use
%   \cs{ADS@sanitize} to change definitions of `~' and \cs{ion}.}
% \changes{3.23}{2007/09/09}{\cs{addkeyword} now signals a compilation
% error if the author tries to add more than 6 keywords.}
% \changes{3.24}{2007/09/14}{kwdcounter is now reset in
% \cs{@begofarticlehook} to prevent error in book style} 
%    \begin{macrocode}
\newcommand{\@kwdsep}{ --- }
\newcommand{\ADS@kwdsep}{, }
\newcounter{kwdcounter}\setcounter{kwdcounter}{0}
\newcommand{\addkeyword}[1]{%
  \addtocounter{kwdcounter}{1}
  \ifthenelse{\thekwdcounter > 6}{
    \ClassError{rmaa}{Too many keywords. A maximum of 6 keywords is allowed.}{}
  }{}
  \if@firstkwd%
  \renewcommand\@keywords{\if@boldkeywords\uppercase{#1}\else #1\fi}
  {\ADS@sanitize\def\ion{\plainion}\gdef\ADS@keywords{#1}}
  \@firstkwdfalse
  \else%
  \edef\@keywords{\@keywords\@kwdsep \if@boldkeywords\uppercase{#1}\else #1\fi}%
  {\ADS@sanitize\def\ion{\plainion}\xdef\ADS@keywords{\ADS@keywords\ADS@kwdsep #1}}%
  \fi}
\newcommand{\ADS@sanitize}{%
  \def\mathrm{}
  \def\@{\ }
  \def~{\ }
  \def\langle{<}
  \def\rangle{>}
  \def\log{log }
  \def\sin{sin }
  \def\cos{cos }
  \def\tan{tan }
  \def\fd{^d.}
  \def\fh{^h.}
  \def\fm{^m.}
  \def\fs{^s.}
  \def\fdg{�.}
  \def\farcm{'.}
  \def\farcs{''.}
  \def\arcdeg{�}
  \def\,{ }
  }
\newcommand{\keywords}[1]{%
  \renewcommand{\ADS@keywords}{#1}
  \renewcommand{\@keywords}{\uppercase{#1}}
  }
\newcommand{\@keywordheader}{\rmaa@KeywordHeaderStyle{Key Words:}}
\newcommand{\@keywords}{%
  \ClassWarningNoLine{rmaa}{No keywords specified}%
  Keywords go here
  }
%    \end{macrocode}
% \begin{macro}{\shorttitle}
% \begin{macro}{\shortauthor}
% \begin{macro}{\@shorttitle}
% \begin{macro}{\@shortauthor}
% Commands for setting the running headers.
%    \begin{macrocode}
\newcommand{\shorttitle}[1]{%
  \renewcommand\@shorttitle{\MakeUpperCase{#1}}}
\newcommand{\shortauthor}[1]{%
  \renewcommand\@shortauthor{\MakeUpperCase{#1}}}
\newcommand{\@shorttitle}{SHORT TITLE GOES HERE}
\newcommand{\@shortauthor}{SHORT AUTHOR GOES HERE}
%    \end{macrocode}
% \end{macro}
% \end{macro}
% \end{macro}
% \end{macro}
% \end{macro}
% \end{macro}
% \end{macro}


% \subsubsection{The full addresses}
% \label{sec:full-addresses}
% \begin{macro}{\fulladdresses}
% \begin{macro}{\@fulladdresses}
% \begin{macro}{\@fulladdheight}
% \begin{macro}{\@fulladddepth}
% \begin{macro}{\@fulladdbox}
% Aaghh, the full addresses! A perennial source of problems. 
%    \begin{macrocode}
\newlength{\@fulladdheight}
\newlength{\@fulladddepth}
\newsavebox{\@fulladdbox}
%% \fulladdresses command modified (11 May 2000) so that we can put
%% \fulladdresses{} in the article to suppress the output of the
%% addresses
\newcommand{\@FinalFigure}{\relax}
\newcommand{\FinalFigure}[1]{\renewcommand{\@FinalFigure}{#1}}
\newcommand{\fulladdresses}[1]{%
  \renewcommand{\@rmaaArg}{#1}
  \ifx\@rmaaArg\@rmaaEmpty
     \renewcommand{\@fulladdresses}{}
  \else
     \renewcommand{\@fulladdresses}{
       \normalsize%
       \ifx\@FinalFigure\relax\else\vfill\@FinalFigure\vfill\fi
%%       \vspace*{1ex}
       %% do addresses as a list (copied from thebibliography environment) 
       \list{\null}{\leftmargin 1.5em\topsep\z@\labelwidth\z@\labelsep\z@\itemindent -1.5em
         \itemsep\z@\parsep\z@\usecounter{enumi}}
       \sloppy\clubpenalty4000\widowpenalty4000
       \sfcode`\.=1000\relax
       #1
       \endlist
       }
  \fi
  }
\newcommand{\@fulladdresses}{Full addresses go here}
%    \end{macrocode}
% \begin{macro}{\suppressfulladdresses}
% New user-level command for suppressing full addresses should you so
% desire. Currently used in the single-page-poster style. 
%    \begin{macrocode}
\newcommand{\suppressfulladdresses}{\@suppressfulladdressestrue\@needaddressesfalse}
%    \end{macrocode}
% \end{macro}
% \end{macro}
% \end{macro}
% \end{macro}
% \end{macro}
% \end{macro}

% \begin{macro}{\printdates}
% \begin{macro}{\noprintdates}
% Command to activate the printing out of the received/accepted
% dates. This can also be done using the \texttt{printdates} option,
% which is the default for \texttt{journal}, but not for \texttt{proceedings}.
%    \begin{macrocode}
\newcommand{\printdates}{\@printdatestrue}
\newcommand{\noprintdates}{\@printdatesfalse}
%    \end{macrocode}
% \end{macro}
% \end{macro}
% \begin{macro}{\@outputaddressesifnecessary}
% Command to output the addresses if we haven't already done so
% This is called by: 
% \begin{enumerate}
% \item The end of |thebibliography| environment
% \item The |\outputfulladdresses| command
% \end{enumerate}
%    \begin{macrocode}
\newcommand{\@outputaddressesifnecessary}{%
  \ClassInfo{rmaa}{Calling @outputaddressesifnecessary}
  \ifthenelse{\boolean{@needaddresses}}{%
    %% If in one column style, or in the first column of a two column
    %% style then we can try to put it on this page 
    \ifthenelse{\boolean{@firstcolumn} \or \not\boolean{@rmtwocol}}{%
      \@outputaddresses}{%
      %% otherwise give warning it will be on a separate page
      \ClassWarning{rmaa}{Full addresses will go on a separate
        page (please see User Guide for how to remedy this)}}%
    %% If we haven't managed yet, put it on its own page
    \ifthenelse{\boolean{@needaddresses}}{%
      \if@debug{\ClassWarning{rmaa}{In this silly bit!}}\fi
      \clearpage\vfill\@outputaddresses}{}%
    }{
    %% If we don't need them, we still might need to do the eoa hook 
    \ifthenelse{\boolean{@suppressfulladdresses}}{%
      \if@debug{\ClassWarning{rmaa}{XXX}}\fi
%%      \if@book\else\@endofarticlehook\null\fi
      }{}
    }%
  }  
%    \end{macrocode}
% \end{macro}

% \begin{macro}{\outputfulladdresses}
% A user interface to the above command, in case there is no biblio
% (11 May 2000) 
%    \begin{macrocode}
\newcommand{\outputfulladdresses}{\@outputaddressesifnecessary}
%    \end{macrocode}
% \end{macro}

% \begin{macro}{\adjustfinalcols}
% This command starts a new column and should be 
% used to ensure that the columns balance on the last page. It also
% outputs the full 
% author addresses to ensure that this occurs in the left hand
% column. 
% The |\par| command is so that the current line will be taken account
% of when working out |\pagetotal|
%    \begin{macrocode}
\newcommand{\adjustfinalcols}{\if@twocolumn\goodbreak\@outputaddresses\newpage\noindent\fi}
%    \end{macrocode}
% \end{macro}

% \begin{macro}{\@outputaddresses}
% Command to do the work of outputting the addresses. 
% This can be called by 
% \begin{enumerate}
% \item an explicit |\adjustfinalcols|
% \item via |\@outputaddressesifnecessary|, q.v.
% \end{enumerate}
%
%    \begin{macrocode}
\newcommand{\@outputaddresses}{%
  % Allow option of suppressing the full addresses, in which case
  % we just run the hook if we are not in a book style
  \if@suppressfulladdresses
%%     \if@book\else\@endofarticlehook\null\fi
  \else
  %% save the full addresses in a box for output later
  \sbox{\@fulladdbox}{\parbox{\textwidth}{\@fulladdresses}}
  % calculate height of address box  
  \settoheight{\@fulladdheight}{\usebox{\@fulladdbox}}
  \settodepth{\@fulladddepth}{\usebox{\@fulladdbox}}
  \addtolength{\@fulladdheight}{\@fulladddepth}
  \if@debug
  { \typeout{Address total height: \the\@fulladdheight}
    \typeout{Text height: \the\pagetotal}
    \typeout{Goal height: \the\pagegoal}
    }
  \fi
  % reuse \@fulladdheight to mean the total height of text+addresses
  \addtolength{\@fulladdheight}{\pagetotal}
%%  \addtolength{\@fulladdheight}{-\baselineskip}
  \if@debug\typeout{Adjusted text + addresses height: \the\@fulladdheight}\fi
  % check it will fit on current page
  \ifdim\@fulladdheight<\pagegoal
     %% must shift it to the left
     %% to compensate for the hanging indentation in the reference list. 
     %% NOTE: this is not foolproof!
     %% NB. This will give overfull hbox warnings, which is a bit ugly.
     \vfill\noindent\if@biblist\hspace*{-\leftmargin}\fi
     \usebox{\@fulladdbox}
     % Set this to show we've done it
     \@needaddressesfalse
     % and assume we are at the end of the article
     % so that we do the EOA hook if not in the book style
%%     \if@book\else\@endofarticlehook\fi
  \else
     % if not, just do nothing. With luck, this function will get 
     % called again on the next page
     \if@debug{\typeout{Holding over full addresses to next page}}\fi
  \fi\fi
  }
%    \end{macrocode}
% \end{macro}

% \subsubsection{Commands for editors}
% \label{sec:editor-commands}
% Some user/editor commands for doing the TOC and index in the book
% option. 
% They are defined here rather in the book-specific part since they
% need to be always available in order to avoid breaking documents. 
% \begin{macro}{\listofauthors}
% \begin{macro}{\@rawauthors}
% \begin{macro}{\rmaaTOCauthors}
% The command |\listofauthors| formats an author list for use in the
% TOC. It should contain the same names as in the |\author| command,
% but without the affiliation garbage and using |\&| instead of
% |and|. 
%    \begin{macrocode}
\newcommand{\@rawauthors}{Raw Authors}
\newcommand{\rmaaTOCauthors}{\noexpand\rmaaTOCAuthorStyle{Author list for TOC}}
\newcommand{\listofauthors}[1]{%
  \renewcommand{\@rawauthors}{#1}
  \renewcommand\rmaaTOCauthors{\noexpand\rmaaTOCAuthorStyle{#1}}}
%    \end{macrocode}
% \begin{macro}{\indexauthor}\label{indexauthor}
% The command |\indexauthor| produces an index entry for each
% author. 
% This command must come \emph{after} the
% |\title| and the |\listofauthors| commands. Also, there should be one 
% |\indexauthor| command for \emph{each} author. We also add the
% author onto the end of the |\ADS@authors| macro and construct the
% |\ADS@authorinitial| by a most cunning method (see below), but only
% for the first author. 
% 
% No, No, NO! This was all wrong. Backed out completely: 
% \begin{quote}
% Completely rewritten (19 Dec 2002) since we now use the index class, so 
% as to be able to alphabetize accented characters. A side effect of this is
% that |\index| no longer expands out the argument after the |@|, so we 
% instead add all the |\index| commands to a list (|\indexauthor@deferlist|), 
% which is executed in the |\@begofarticlehook|. 
% 
% Note the stuff with |\toks0| is copied from Exercise 20.15 of the \TeX{}book---it is
% to make sure the latin1-accented name gets passed into the idx
% file. 
% \end{quote}
% We now go back to the old way of doing things, using an optional
% argument to the \cs{indexauthor} command that has the name without
% accents. 
% \changes{3.08}{2003/01/16}{Remerged changes from 3.06 that got
% dropped from 3.07 after a synchronization disaster. Basically, the
% accentless optional argument to \cs{indexauthor}. Must add the macro
% directories to my unison set-up to stop this sort of thing from happening.}
%    \begin{macrocode}
\newcommand\@isortkey{}
\newcommand\@emptykey{empty}
\newcommand\indexauthor[2][empty]{
  \renewcommand\@isortkey{#1}
  \ifx\@isortkey\@emptykey \renewcommand\@isortkey{#2}\fi
  \index{\@isortkey @{{\rmaaIndexAuthorStyle{#2}}\hskip1em plus1em minus0.4em
      \rmaaIndexTitle\hskip1em plus2em minus0.4em\penalty-100
      \rmaaTOCauthors}}
  \indexauthorADS{#2}
}
%%\def\indexauthor@deferlist{\relax}
%%\long\def\indexauthor#1{
%%  \def\@@name{#1} \toks0=\expandafter{\@@name}
%%  \toks1=\expandafter{\indexauthor@deferlist}
%%  \protected@xdef\indexauthor@deferlist{%
%%    \the\toks1
%%    \noexpand\index{\the\toks0@{\noexpand\rmaaIndexAuthorStyle{\the\toks0}\noexpand\hskip1em plus1em minus0.4em 
%%        \rmaaIndexTitle\noexpand\hskip1em plus2em minus0.4em\noexpand\penalty-100
%%        \rmaaTOCauthors}}
%%  }
%%  \indexauthorADS{#1}
%%  }
%%\xdef\indexauthor@deferlist{}
\newcommand\indexauthorADS[1]{%
  \if@firstauthor\protected@xdef\ADS@authorinitial{\firstletter{#1}}\fi
  \protected@xdef\ADS@authors{\ADS@authors\if@firstauthor\else ; \fi
    #1}
  \@firstauthorfalse
}
%    \end{macrocode}
% Note that the first |#1| sets the key for the index, making sure it
% actually starts with  
% first letter of the surname. Meanwhile, the |@{...}| expression,
% which starts with some  
% control sequences is what is written to the |.ind| file. 
% See the makeindex man page for details. 
% \end{macro}
% \end{macro}
% \end{macro}
% \end{macro}

% \begin{macro}{\KillAllAccents}
% \begin{macro}{\firstletter}
% These are utility functions used in the above. |lett| is a counter
% indicating the current letter in the string. |\firstletter| just
% prints the first letter of its argument.
%    \begin{macrocode}
\newcommand{\KillAllAccents}{\def\'{}\def\`{}\def\~{}\def\.{}\def\i{i}}
\newcommand\firstletter[1]{\@car #1\@nil}
%    \end{macrocode}
% \end{macro}
% \end{macro}
% \begin{macro}{\SetVolume}
% \begin{macro}{\SetMSnumber}
% \begin{macro}{\SetFirstPage}
% \begin{macro}{\SetYear}
% \begin{macro}{\ReceivedDate}
% \begin{macro}{\AcceptedDate}
% These are for the editors to set the volume number, etc. 
%    \begin{macrocode}
\newcommand{\SetVolume}[1]{%
  \renewcommand{\@volume}{#1}
  \edef\ADS@volume{\PadToFourWithDots{#1}}
  \renewcommand{\@volumen}{\rmaa@VolumenStyle{volumen #1}}%
  }
\newcommand{\SetMSnumber}[1]{\renewcommand{\@MSnum}{#1}}
%    \end{macrocode}
% Note that there is no real reason to use |\SetFirstPage| since all
% the page numbering is done automatically in the book style. 
%    \begin{macrocode}
\newcommand{\SetFirstPage}[1]{\setcounter{page}{#1}}
\newcommand{\SetYear}[1]{%
  \renewcommand{\@year}{#1}
  \renewcommand{\ADS@year}{#1}
  \renewcommand{\@anyo}{\rmaa@VolumenStyle{#1}}
  }
\newcommand{\ReceivedDate}[1]{\renewcommand{\@received}{#1}}
\newcommand{\AcceptedDate}[1]{\renewcommand{\@accepted}{#1}}
%    \end{macrocode}
%
% \begin{macro}{\DOI}
% User/editor command to enter the DOI, which will be typeset as a link.
% \changes{4.0}{2019/11/02}{New user command \cs{DOI} to typeset
% DOI in the header of first page and add it to ADS info.}
%    \begin{macrocode}
\newcommand\DOI[1]{%
  \def\@DOI@url{\url{https://doi.org/10.22201/ia.01851101p.#1}}%
  \def\ADS@doi{https://doi.org/10.22201/ia.01851101p.#1}%
}%
%    \end{macrocode}
% \end{macro}
%
% \begin{macro}{\ORCIDiD}
% \changes{4.0}{2019/11/12}{Add mechanism for giving ORCID iD links}
% \begin{macro}{\ORCIDiDfull}
% \begin{macro}{\@orcid@icon}
% \begin{macro}{\@orcid@iconbig}
% Commands to add icon links to ORCID iDs of authors. For example
% \begin{quote}
% |\ORCIDiD|\marg{0000-0000-0000-0000}
% \end{quote}
% to be used after each author name in the |\author|
% command and in the |\fulladdresses|.  Alternatively,
% |\OORCIDiDfull| will typeset the icon plus the full text of the
% link, but we don't use that at present.  The icon image is stored in |ORCIDiD_iconvector.pdf|, which needs to be on the input path.  In principle, we could make EPS versions if necessary.
%    \begin{macrocode}
\newcommand\@orcid@icon{\includegraphics[height=1.8ex]{ORCIDiD_iconvector}}
\newcommand\@orcid@iconbig{\includegraphics[height=2ex]{ORCIDiD_iconvector}}
\newcommand\ORCIDiD[1]{\href{https://orcid.org/#1}{\,\@orcid@icon}}
\newcommand\ORCIDiDfull[1]{%
  \href{https://orcid.org/#1}{\@orcid@iconbig}\,\url{https://orcid.org/#1}}
%    \end{macrocode}
% \end{macro}
% \end{macro}
% \end{macro}
% \end{macro}
%
% \begin{macro}{\@volume}
% \begin{macro}{\@volumen}
% \begin{macro}{\@artpages}
% \begin{macro}{\@year}
% \begin{macro}{\@anyo}
% \begin{macro}{\@received}
% \begin{macro}{\@accepted}
% \begin{macro}{\@MSnum}
% Default values for the containers used by above commands, plus a
% few more useful ones. 
%    \begin{macrocode}
\newcommand{\@volume}{00}
\newcommand{\@volumen}{\rmaa@VolumenStyle{00}}
\newcommand{\@artpages}{%
  \oldpageref{\articlename:RMAAFirstPage}--\oldpageref{\articlename:RMAALastPage}}
\newcommand{\@year}{\the\year}
\newcommand{\@anyo}{\rmaa@VolumenStyle{2007}}
\newcommand{\@journalname}{\rmaa@JournalNameStyle{%
    \if@proceedings\RMAASClong\else\RMAAlong\fi}}
\newcommand{\@received}{}
\newcommand{\@accepted}{}
\newcommand{\@MSnum}{0000}
\newcommand{\@DOI@prefix}{https://doi.org/10.22201/ia.01851101p}
\def\@DOI@url{\url{https://doi.org/10.22201/ia.01851101p.2099.00.00.00}}
%    \end{macrocode}
% \end{macro}
% \end{macro}
% \end{macro}
% \end{macro}
% \end{macro}
% \end{macro}
% \end{macro}
% \end{macro}
% \end{macro}
% \end{macro}
% \end{macro}
% \end{macro}
% \end{macro}
% \end{macro}
% 
% The following commands are most useful in the \texttt{book}
% style. However, the |\SetConfTitle| command can now be used in
% individual articles to set the header for the \texttt{preprint}
% style. 
%    \begin{macrocode}
\newcommand{\@reunionplace}{Lugar de la reuni\'on}
\newcommand{\@reuniondate}{Fecha de la reuni\'on}
\newcommand{\@conftitle}{Conference Title}
\newcommand{\SetReunion}[2]{\renewcommand{\@reunionplace}{#1}\renewcommand{\@reuniondate}{#2}}
\newcommand{\SetConfTitle}[1]{\renewcommand{\@conftitle}{#1}}
\newcommand{\@editors}{Juan Ram\'on de la Fuente and ``El Mosh''}
\newcommand{\SetEditors}[1]{\renewcommand{\@editors}{#1}}
\newcommand{\@editores}{Juan Ram\'on de la Fuente y ``El Mosh''}
\newcommand{\SetEditores}[1]{\renewcommand{\@editores}{#1}}
\newcommand{\@mes}{mes de la publicaci\'on}
\newcommand{\SetMes}[1]{\renewcommand{\@mes}{\rmaa@VolumenStyle{#1}}}
%    \end{macrocode}

% \subsection{Notes for Author Proofs and Communication between Editors}
% \label{sec:notes}
% We introduce a command \cs{Note} that should be used to describe
% any change made to the manuscript. These are a bit tricky to get
% right. In general we use \cs{marginpar} but this only works in outer
% paragraph mode, so we just bung the box where we are if we are in a
% figure caption or something. If we are in the preamble, then we just
% use a \marginpar as part of the \maketitle command.
%    \begin{macrocode}
\newif\if@arenotes
\@arenotesfalse
\newcounter{notecounter}
\setcounter{notecounter}{0}
\newcommand\NotesList{}
\newcommand\Note[1]{%
  \if@proof{%
    \@arenotestrue%
    \stepcounter{notecounter}%
    \if@inpreamble\else{%
      \NoteMark{\arabic{notecounter}}%
    }\fi
    \protected@xdef\NotesList{\NotesList 
      \protect\item #1}%
  }\fi%
}
\newcommand\Changed[2]{Changed ``#1'' to ``#2''.}
\newcommand\NoteMark[1]{%
  \ifinner
  \raisebox{1ex}[0pt][0pt]{\makebox[0pt][l]{\framebox{\footnotesize\bfseries
        N#1}}}%
  \else
  \marginpar{\framebox{\footnotesize\bfseries N#1}}%
  \fi
}
\newcommand\RefNote[1]{\if@proof\NoteMark{\ref{#1}}\relax\fi}
\newcommand\NoteLabel{\protect\label}
\newcommand\OutputNotes{%
  \bgroup
  \onecolumn
  \section*{EDITORIAL NOTES}
  \begin{enumerate}
    \NotesList
  \end{enumerate}
  \egroup
}
\DeclareRobustCommand\cs[1]{\texttt{\char`\\#1}}

%    \end{macrocode}
% \subsection{Bibliographic Records for ADS}
% \label{sec:ADS}
% \changes{2.2beta11}{2001/11/27}{Just about finished all the ADS
% stuff now. Except for the references, which I have given up on for
% the time being.}
% \begin{macro}{\ADS@BibCode}
% \begin{macro}{\ADS@year}
% \begin{macro}{\ADS@journal}
% \begin{macro}{\ADS@volume}
% \begin{macro}{\ADS@firstpage}
% \begin{macro}{\ADS@authorinitial}
% We use the |\ADS@BibCode| command to hold the bibliographic code for
% ADS. It is 
% used in the |%R| 
% field of the bibliographic record for each article. 
% The format is 
% \begin{quote}
% \texttt{YYYYJJJJJVVVVMPPPPA}  
% \end{quote}
% where \texttt{YYYY} is the year, \texttt{JJJJJ} is the journal code,
% which is \texttt{RMxAA} or \texttt{RMxAC}, \texttt{VVVV} is the
% volume number (right-justified), \texttt{M} is not used by us,
% \texttt{PPPP} is the first page number (right-justified), and
% \texttt{A} is the initial of the first author's surname. The volume
% and page fields are padded with periods (`\texttt{.}'), which causes
% us some fun (this is done with the macro |\PadToFourWithDots| described
% below) . 
% 
% All these individual bits are held in the commands |\ADS@year|,
% |\ADS@journal|, |\ADS@volume|, |\ADS@firstpage|,
% |\ADS@authorinitial|. 
%    \begin{macrocode}
\newcommand{\SetADSBibCode}[1]{\renewcommand{\ADS@BibCode}{#1}}
\newcommand{\ADS@BibCode}
{\ADS@year\ADS@journal\ADS@volume.\PadToFourWithDots{\ADS@firstpage}\ADS@authorinitial}
\newcommand{\ADS@year}{2001} 
\newcommand{\ADS@journal}{\if@proceedings RMxAC\else RMxAA\fi}
\newcommand{\ADS@volume}{....} 
%    \end{macrocode}
% These two are just set to default values here, |\ADS@firstpage| is set
% in the |document| environment for included articles, in the
% |frontsection| environment, |\tableofcontents| command,  |index| and
% |abstracts| environments. Also by |\AtBeginDocument| for stand-alone
% articles. 
% 
% God knows how we are going to do |\ADS@authorinitial|!
%    \begin{macrocode}
\newcommand{\ADS@firstpage}{....}
\newcommand{\ADS@authorinitial}{A}
%    \end{macrocode}
% \end{macro}
% \end{macro}
% \end{macro}
% \end{macro}
% \end{macro}
% \end{macro}
% We also have a whole slew of other fields that can be set. These are
% all listed in Table~\ref{tab:ADS}. Most of them still have to be
% defined. 
% \begin{table}[!ht]\raggedleft
% \caption{ADS Bibliographic Records}\label{tab:ADS}
% \begin{tabular}{lp{8em}lp{8em}}\hline
% Code & Description & Macro Name & Comment \\ \hline
% |%R| & Bibliographic Code & |\ADS@BibCode| & required \\
% |%A| & Author list & |\ADS@authors| & required, generate from |\indexauthor| \\
% |%F| & Author affiliation & |\ADS@affil| & format: AA(\dots)
% AB(\dots) \dots \\
% |%J| & Journal Name & |\ADS@journalname| &  \\
% |%D| & Publication Date & |\ADS@fulldate| &  \\
% |%L| & Last Page & |\ADS@lastpage| & generate from label \\
% |%T| & Title & |\ADS@title| & required, generate from |\title| \\
% |%C| & Abstract Copyright & |\ADS@copyright| & set to IA-UNAM \\
% |%U| & URL for electronic document & |\ADS@url| & set to PDF file \\
% |%d| & DOI for article & |\ADS@doi| & set by |\DOI| \\
% |%K| & Keywords & |\ADS@keywords| & set equal to |\@keywords| \\
% |%A| & Abstract Text & |\ADS@abstract| & set equal to |\@abstract| \\
% |%R| & References & |\ADS@references| & generate from the
% bibliography  \\ \hline
% \end{tabular}
% \end{table}
% A lot of these macros are set by other commands, as mentioned in the
% Table, but some need to be set  explicitly here. 
% 
% \begin{macro}{\ADS@doi}
% \changes{4.0}{2019/11/08}{Save DOI for ADS}
% The user/editor command \cs{DOI} is used to set this. 
%    \begin{macrocode}
\newcommand{\ADS@doi}{https://doi.org/10.22201/ia.01851101p.0000.00.00.00}
%    \end{macrocode}
% \end{macro}
% 
% \begin{macro}{\ADS@journalname}
% This
% does not contain any accents at the moment so as not to cause
% problems---this maybe should be changed though. 
%    \begin{macrocode}
\newcommand{\ADS@journalname}{%
  \if@proceedings\@conftitle\space (Eds. \@editors) \fi%
  Revista Mexicana de Astronom\'\i{}a y Astrof\'\i{}sica%
  \if@proceedings\space (Serie de Conferencias)\fi%
  \space Vol. \@volume, pp. \ADS@firstpage--\ADS@lastpage\space (\@year)%
  \space (http://www.astroscu.unam.mx/\~{}rmaa/)
  }
%    \end{macrocode}
% \end{macro}
% \begin{macro}{\ADS@keywords}
% \begin{macro}{\ADS@abstract}
% \begin{macro}{\ADS@title}
% |\ADS@abstract| is straightforward, since we already have the text
% saved as |\@abstract|. We could do the same with |\ADS@keywords| but
% to avoid unwanted formatting we explicitly set it in |\keywords|
% (p.~\pageref{keywords}). |\ADS@title| is just set to |\@rawtitle|,
% which we were already using for another purpose. 
%    \begin{macrocode}
\let\ADS@title=\@rawtitle
\newcommand{\ADS@keywords}{No keywords specified}
\newcommand{\ADS@abstract}{No abstract specified}
%    \end{macrocode}
% \end{macro}\end{macro}\end{macro}
% \begin{macro}{\ADS@authors}
% |\ADS@authors| wants to be in ``Other, A. N.'' format, but with
% semi-colons between authors. The most natural way to do this seems to
% be with |\indexauthor| (p.~\pageref{indexauthor}). We append each
% author to this macro, which means we must explicitly clear it for
% each new article (done in |\article|). .   
%    \begin{macrocode}
\newcommand{\ADS@authors}{}
%    \end{macrocode}
% \end{macro}
% \begin{macro}{\ADS@fulldate}
% \begin{macro}{\ADS@lastpage}
% \begin{macro}{\ADS@copyright}
% \begin{macro}{\ADS@website}
% \begin{macro}{\ADS@url}
% \begin{macro}{\SetWebsite}
% \begin{macro}{\SetMonthNumeric}
% Other sundry data, together with user commands for bits that can be
% set by the Editors. The |\ADS@url| is constructed from the website
% (currently |http://www.astrosmo.unam.mx/rmaa/|,
% the journal descriptor (e.g., |RMxAC..10/|), followed by the
% path to the PDF file (e.g., |PDF/RMxAC..10_HEW1.pdf|). 
% 
%    \begin{macrocode}
\newcommand{\ADS@website}{http://www.astrosmo.unam.mx/rmaa/}
\newcommand{\ADS@url}
{\ADS@website\ADS@journal\ADS@volume\ADS@number/PDF/\ADS@journal\ADS@volume\ADS@number\_\articlename.pdf}
\newcommand{\ADS@number}{}
\newcommand{\ADS@fulldate}{00/\ADS@year} 
\newcommand{\ADS@lastpage}{00} 
\newcommand{\ADS@copyright}{(C) \ADS@year; Instituto de Astronom\'\i{}a,
  Universidad Nacional Aut\'onoma de M\'exico} 
\newcommand{\SetNumber}[1]{\renewcommand{\ADS@number}{-#1}}
\newcommand{\SetWebsite}[1]{\renewcommand{\ADS@website}{#1}}
\newcommand{\SetMonthNumeric}[1]{\renewcommand{\ADS@fulldate}{#1/\ADS@year}}
%    \end{macrocode}
% \end{macro}\end{macro}\end{macro}\end{macro}\end{macro}\end{macro}\end{macro}
% \begin{macro}{\ADS@affil}
% \begin{macro}{\ADS@references}
% These two are somewhat more problematical. Affiliation currently
% only works for the first author (see \dots) pending redesign of the
% user interface. References are now done by wrapping the
% |thebibliography| environment to open another input stream on the
% same file and search through for the |\bibitems|. 
%    \begin{macrocode}
\newcommand{\ADS@references}{}
\newcommand{\ADS@affil}{}
%    \end{macrocode}
% \end{macro}\end{macro}
% 
% Define two macros that expand to the strings
% |\begin{thebibliography}| and |\end{thebibliography}|. Note that the
% idea for this and much else is inspired by the code for the docstrip
% package. Also define |\StartBibX| in case authors use the more
% standard |\begin{thebibliography}{}| syntax. 
%    \begin{macrocode}
\edef\EndBib{\expandafter\@gobble\string\\end\string{thebibliography\string}}
\edef\StartBib{\expandafter\@gobble\string\\begin\string{thebibliography\string}}
\edef\StartBibX{\expandafter\@gobble\string\\begin\string{thebibliography\string}\string{\string}}
%    \end{macrocode}
% Define two new logical variables. |\ifAnotherline| is true if we
% want to read another line from the input stream. It is set to false
% when we encounter the |\end{thebibliography}|. |\ifInBib| is set to
% true when we frind the |\begin{thebibliography}| and tells us to
% start saving the following lines. 
%    \begin{macrocode}
\newif\ifAnotherLine
\newif\ifInBib
%    \end{macrocode}
% \begin{macro}{makeOther}
% We need a macro to change the catcodes of lots of things to other. 
%    \begin{macrocode}
\def\makeOther#1{\catcode`#1=12\relax}
%    \end{macrocode}
% \end{macro}
% \begin{macro}{\RereadBiblio}
% Define a macro to do all the fiddly stuff. This is used at the
% beginning of |thebibliography| environment (Sec~\ref{sec:biblio}).
%    \begin{macrocode}
\newread\tf@biblio
\newcommand{\RereadBiblio}{%
%    \end{macrocode}
% Open the current file again for reading
%    \begin{macrocode}
  \begingroup
  \immediate\openin\tf@biblio=\articlename
%    \end{macrocode}
% Everything in a group to be safe. Initialize logicals. Black magic
% to change
% catcodes and avoid blank lines being interpreted as |\par|. 
%    \begin{macrocode}
  \AnotherLinetrue\InBibfalse
  \makeOther\ \makeOther\\\makeOther\$%
  \makeOther\#\makeOther\^\makeOther\^^K%
  \makeOther\_\makeOther\^^A\makeOther\%%
  \makeOther\~\makeOther\{\makeOther\}\makeOther\&%
  \endlinechar-1\relax
%    \end{macrocode}
% Start looping through the file. Read each line to the macro
% |\Line|. If we have got to EOF or end of biblio then set
% |\AnotherLinefalse|.
%    \begin{macrocode}
  \loop
  \read\tf@biblio to \Line
  \ifx\Line\EndBib\AnotherLinefalse\InBibfalse\fi
  \ifeof\tf@biblio \AnotherLinefalse\fi
%    \end{macrocode}
% If we are in the biblio and the line is not blank then append it to
% the |\ADS@references| macro together with a newline. 
%    \begin{macrocode}
  \ifInBib
  \ifx\Line\empty\else\xdef\ADS@references{\ADS@references\Line^^J}\fi
  \fi
%    \end{macrocode}
% Check if the line we just saw was the start of biblio, then repeat
% the loop if necessary. 
%    \begin{macrocode}
  \ifx\Line\StartBib\InBibtrue\fi
  \ifx\Line\StartBibX\InBibtrue\fi
  \ifAnotherLine\repeat
  \immediate\closein\tf@biblio
  \endgroup
  }
%    \end{macrocode}
% \end{macro}
% 
% \begin{macro}{\PadToFourWithDots}
% Utility function for left-padding a number to 4 characters with
% periods. Used in 
% constructing |\ADS@BibCode|. This is a very simple algorithm and
% \emph{requires} that the argument actually be a number---if it is
% not then dire consequences will result.  
%    \begin{macrocode}
\newcommand{\PadToFourWithDots}[1]{%
  \ifnum #1 > 999 #1%
  \else\ifnum #1 > 99 .#1%
  \else\ifnum #1 > 9 ..#1%
  \else ...#1%
  \fi\fi\fi}
%%  \ifnum \count0 > 999 #1%
%%  \else\ifnum \count0 > 99 .#1%
%%  \else\ifnum \count0 > 9 ..#1%
%%  \else ...#1%
%%  \fi\fi\fi}
%    \end{macrocode}
% \end{macro}
%
% \begin{macro}{\AuthorADS}
% \begin{macro}{\ADS@affiliations}
% New mechanism for setting the affiliations for ADS. The author labels (AA, AB, \dots) are generated using two nested counters.  We use the |\robustify| command from \texttt{etoolbox} so that accented characters in the affiliation addresses do not blow up.  The list of accents may need to be expanded. 
%    \begin{macrocode}
\newcounter{AUTHOR}
\newcounter{author}[AUTHOR]
\stepcounter{AUTHOR}
\newcommand\ADS@affiliations{}
\robustify{\'}\robustify{\i}\robustify{\~}\robustify{\`}
\def\AuthorADS#1#2{%
  \begingroup
  \def\ORCIDiD##1{<ID system="ORCID">##1</ID>}
  \def\Email##1{<EMAIL>##1</EMAIL>}
  \def\Affil##1{##1}
  \stepcounter{author}
  \ifthenelse{\value{author} > 26}{\stepcounter{AUTHOR}\stepcounter{author}}{}
  \ifthenelse{\equal{\ADS@affiliations}{\@rmaaEmpty}}{%
    \xdef\ADS@affiliations{\Alph{AUTHOR}\Alph{author}(#2)}
  }{%
    \xdef\ADS@affiliations{\ADS@affiliations, \Alph{AUTHOR}\Alph{author}(#2)}
  }%
  \endgroup
}
%    \end{macrocode}
% \end{macro}
% \end{macro}
% %
% \begin{macro}{\ADS@dumpdata}
% \changes{4.0}{2019/11/13}{Write all author affiliations for ADS}
% \begin{macro}{\ADS@write}
% \begin{macro}{\ADS@openfile}
% |\ADS@dumpdata| is the macro that actually does the job of writing out all the
% data to an auxiliary file, using |\ADS@write| to write each
% field. The file is set up with |\ADS@openfile|, which currently is
% done in the TOC of a book.  The new multi-author affiliations are used if |\ADS@affiliations| has been set, otherwise it falls back to the single-author version |\ADS@affil|. 
%    \begin{macrocode}
\newcommand{\ADSfield}[2]{\relax}
\newcommand{\ADSvfield}[2]{\relax}
%%\renewcommand{\ADSfield}[12]{\%#1 #2}
\newcommand{\ADS@dumpdata}
{ 
  \ADS@write{\string\begin{ADSarticle}}
  \ADS@WRITE{\string\ADSfield{R}{\ADS@BibCode}}
  \ADS@write{\string\ADSfield{A}{\ADS@authors}}
  \ifthenelse{\equal{\ADS@affiliations}{\@rmaaEmpty}}{%
    {\ADS@sanitize\ADS@write{\string\ADSfield{F}{AA(\ADS@affil)}}}
  }{%
    {\ADS@sanitize\ADS@write{\string\ADSfield{F}{\ADS@affiliations}}}
  }
  \ADS@write{\string\ADSfield{J}{\ADS@journalname}}
  \ADS@write{\string\ADSfield{D}{\ADS@fulldate}}
  \ADS@write{\string\ADSfield{L}{\ADS@lastpage}}
  {\ADS@sanitize\ADS@write{\string\ADSfield{T}{\ADS@title}}}
  \ADS@write{\string\ADSfield{C}{\ADS@copyright}}
  \ADS@write{\string\ADSfield{d}{\ADS@doi}}
  \ADS@write{\string\ADSfield{U}{\ADS@url}}
  \ADS@write{\string\ADSfield{K}{\ADS@keywords}}
  {\ADS@sanitize\ADS@write{\string\ADSfield{B}{\ADS@abstract}}}
  \ADS@write{\string\ADSfield{Z}{\ADS@references}}
  \ADS@write{\string\end{ADSarticle}}
  }
%%\newcommand{\ADS@write}[1]{\if@book\addtocontents{ads}{#1}\fi}
%%\newcommand{\ADS@openfile}{\@starttoc{ads}}
\newcommand{\ADS@write}[1]{\set@display@protect\immediate\write\tf@ads{#1}}
\newcommand{\ADS@WRITE}[1]{\immediate\write\tf@ads{#1}}
\newcommand{\ADS@openfile}{%
  \newwrite\tf@ads%
  \immediate\openout\tf@ads\jobname.ads\relax
  \ADS@WRITE{\ADS@preamble}}
\newcommand{\ADS@closefile}{\ADS@WRITE{\ADS@postamble}}
\def\ADS@preamble{%
  %%  \string\documentclass{article}
  %%  \string\begin{document}
  }
\def\ADS@postamble{%
  %%  \string\end{document}
  \string%\string% Local Variables:^J 
  \string%\string% mode: LaTeX^J
  \string%\string% mode: Font-lock^J
  \string%\string% End:^J
  }
%    \end{macrocode}
% \end{macro}
% \end{macro}
% \end{macro}
% 
% \subsection{Automatic calculation of the page range}
% \label{sec:page-range}
% Another issue that has cost me a lot of grief. Still, it finally
% seems to be sorted now. All hail Jeff Goldberg. 
%
% Cunning trick to get the page numbers from |\pageref|.  Put a label
% on the first and last page of the document using |\AtBeginDocument|
% and |\AtEndDocument|. .  This way won't work for the book option.
% In that case, we do it by defining |\@endofarticlehook| and
% |\@begofarticlehook| as parallels to |\AtEndDocument| and
% |\AtBeginDocument|.  Note that we have to defer adding
% |\@begofarticlehook| to |\AtBeginDocument| to make sure it comes
% after the redefinitions of |\ref| and |\label|.
%    \begin{macrocode}
\if@book\AtEndDocument{\ADS@closefile}\else
 %\AtBeginDocument{\@begofarticlehook}
\AtEndDocument{\@endofarticlehook}
\fi
%    \end{macrocode}
% \begin{macro}{\@begofarticlehook}
% Since we are getting more and more things that need to be done at
% the beginning of every article, frontsection, etc we define a sister
% hook to 
% |\@endofarticlehook|. This hook does the following: 
% \begin{enumerate}
% \item Labels the first page
% \item Sets |\ADS@firstpage|
% \item runs and then clears |\indexauthor@deferlist|
% \item resets the keyword counter
% \end{enumerate}
% It is executed in the following circumstances: 
% \begin{enumerate}
% \item At |\begin{document}| for standalone files
% \item At |\begin{document}| for each included file
% \item At beginning of front sections 
% \item At beginning of TOC
% \item At beginning of Author Index
% \item At beginning of abstracts
% \end{enumerate}
% For some reason, this is not always within the right group, so we
% use |\xdef| just to be sure. 
%    \begin{macrocode}
\newcommand{\@begofarticlehook}{%
  \if@frontmatter\frontlabel{RMAAFirstPage}\else\label{RMAAFirstPage}\fi%
  \xdef\ADS@firstpage{\@arabic{\c@page}}%
  \if@book\else\makeindex\fi % needed so default index is defined in stand-alone articles
%%  \indexauthor@deferlist
%%  \xdef\indexauthor@deferlist{}
  \@inpreamblefalse
  \setcounter{kwdcounter}{0}
  }
%    \end{macrocode}
% \end{macro}
% \begin{macro}{\@endofarticlehook}
% This is a hook to be run at the end of each article
% currently, it just labels the last page and writes out all the ADS
% data to an auxiliary file.  
% It should not produce any printed output. 
% It is modelled on the |\AtEndDocument| command, to which it reduces
% in the case of a single article. 
%    \begin{macrocode}
\newcommand{\@endofarticlehook}{%
  \if@debug{\ClassWarning{rmaa}{***Doing Hook***}}\fi
  \if@debug{\ClassWarning{rmaa}{ADSabstract is: \ADS@abstract}}\fi
  \clearpage\rmaa@lastpagelabel\if@book\ADS@dumpdata\fi
  \ifx\NotesList\@rmaaEmpty\else\OutputNotes\fi
  \if@debug{\ClassWarning{rmaa}{Notes list is: \NotesList}}\fi
}
%    \end{macrocode}
% \begin{macro}{\rmaa@lastpagelabel}
% New version of last page label command copied from Jeff Goldberg's
% lastpage.sty package - we can't just use |\label| here in case the
% output routine for the last page has already been called. Instead,
% we first explicitly clear the page, step back the page counter,
% explicitly write to the .aux file, then restep the page counter
% again.
%
% But this will conflict with the hyperref package, which does labels differently.  Compensate for that by using the five-argument form.
% 
% Also set the |\ADS@LastPage| macro to the appropriate number. 
%    \begin{macrocode}
\newcommand{\rmaa@lastpagelabel}{%
  \addtocounter{page}{-1}%
  \immediate\write\@auxout{\string
    \newlabel{\articlename:RMAALastPage}{{}{\thepage}{}{}{}}}%
  \edef\ADS@lastpage{\thepage}%
  \addtocounter{page}{1}}
%    \end{macrocode}
% \end{macro}
% \end{macro}

% \subsection{Deal with explicit change in number of columns}
% \label{sec:columns}
% Sometimes, when formatting a book, we need to switch between
% one-column articles and two-column articles. For this to work, we
% need to wrap the standard \LaTeX{} commands to set the |\@rmtwocol|
% flag appropriately and also reset the package default values of all
% scalable lengths. .  
%    \begin{macrocode}
\let\latexonecolumn=\onecolumn
\let\latextwocolumn=\twocolumn
\renewcommand{\onecolumn}{\@rmtwocolfalse\@package@default@lengths\latexonecolumn}
\renewcommand{\twocolumn}{\@rmtwocoltrue\@package@default@lengths\latextwocolumn}
%    \end{macrocode}
% 
% \subsection{The actual formatting of the title}
% \label{sec:title-format}

% Note, we have to redefine the |\maketitle| command to use the
% numbered footnotes. This means we needn't use the atbeginend stuff. 
%    \begin{macrocode}
\renewcommand{\maketitle}{\par
  \begingroup
%    \end{macrocode}
% Don't switch to marked footnotes. 
% This is consistent with existing style but I'm not sure I like it
%    \begin{macrocode}
  %%  \renewcommand\thefootnote{\fnsymbol{footnote}}%
%%  \def\@makefnmark{\hbox to\z@{$\m@th^{\@thefnmark}$\hss}}%
  \def\@makefnmark{\textsuperscript{\@thefnmark}\hspace{-0.2em}}%
  \long\def\@makefntext##1{\parindent 1em\noindent
    \hbox to1.8em{\hss$\m@th^{\@thefnmark}$}##1}%
  \if@rmtwocol
  { \if@debug{\ClassWarning{rmaa}{TWO COLUMN}}\fi
    \twocolumn[\@maketitle]%
    }
  \else
  { \if@debug{\ClassWarning{rmaa}{ONE COLUMN}}\fi
    \newpage
    \global\@topnum\z@   % Prevents figures from going at top of page.
    \if@manuscript \@MSmaketitle \else \@maketitle \fi
    }
  \fi
  \if@manuscript\else\thispagestyle{rmaafirstpage}\@thanks\fi
  \ifx\NotesList\@rmaaEmpty\else{%
    \NoteMark{1--\arabic{notecounter}}\vspace*{-2\baselineskip}%
  }\fi%
  \endgroup
%    \end{macrocode}
% Don't reset footnote counter (not sure whether this is good thing
% or not).
%    \begin{macrocode}
  %%  \setcounter{footnote}{0}%
  \let\thanks\relax\let\affil\relax\let\altaffil\relax
  \let\maketitle\relax\let\@maketitle\relax
  \gdef\@thanks{}\gdef\@author{}\gdef\@title{}}
%    \end{macrocode}

% \begin{macro}{\@kwhwidth}
% \begin{macro}{\@kwtwidth}
% Redefine \cs{@maketitle} to include the abstract, resumen and
% keywords. Note that in the manuscripty style we use
% \cs{@MSmaketitle} instead.  
% \changes{3.16}{2003/04/07}{
%  Added extra vertical space between authors and accepted date.}
% \changes{3.26}{2011/08/15}{Add extra label to review articles}
%    \begin{macrocode}
\newlength{\@kwhwidth}\newlength{\@kwtwidth}
\newcommand\reviewname{Review}
\renewcommand{\@maketitle}{%
  \newpage
  \null
  \if@reviewarticle
  \smallskip
  {\large\bfseries\itshape\reviewname}%
  \bigskip
  \else
  \vskip\@before@title@skip%
  \fi
  \begin{center}%
    {\@title \par}%
    \vskip\@before@author@skip%
    \parbox{\rmaa@AuthorWidthFrac\textwidth}{\centering
      \lineskip\@between@author@skip%
      \@author
      \par}
    %% this should be replaced with `Submitted to...' for preprints
    \if@printdates
    \vskip\@before@dates@skip%
    \if@preprint
    \rmaa@ReceivedStyle{Submitted to RMxAA \@received}
    \else
    \rmaa@ReceivedStyle{Received \@received; accepted \@accepted}
    \fi
    \fi
    \if@manuscript
    \rmaa@ReceivedStyle{Submitted to RMxAA \@received}
    \fi
    \if@guide
    \rmaa@ReceivedStyle{Version \rmaa@version, \rmaa@date}
    \fi
    \vskip\@before@abstract@skip%
    \if@abstractpresent
    \parbox{\rmaa@AbstractWidthFrac\textwidth}{%
      \setlength{\parindent}{0.9cm}%
      \centerline{\rmaa@AbstractHeaderStyle{Abstract}}\par%
      \vskip\@after@abstracthead@skip%
      \if@newstyle\noindent\fi\@abstract
      \vskip\@between@abstract@skip%
      \centerline{\rmaa@AbstractHeaderStyle{Resumen}}\par%
      \vskip\@after@abstracthead@skip%
      \if@newstyle\noindent\fi\@resumen
      \if@setresumen\else\@abstract\fi % append english abstract if
                                       % spanish one not set
      \if@guide\else            %do not use keywords for guide
      \vskip\@before@kwd@skip%
      % Format the Key Words with hanging indentation:
      % They are put in a \parbox with width equal to the line width
      % minus the width of the header. 
      { % Widths of keyword
        \settowidth{\@kwhwidth}{\@keywordheader~}   % header and
        \setlength{\@kwtwidth}{\linewidth}          % keyword text
        \addtolength{\@kwtwidth}{-\@kwhwidth}       %
        \noindent\@keywordheader~\parbox[t]{\@kwtwidth}{%
          \rmaa@KeywordStyle{\boldmath\uppercase{\@keywords}}}}%
      \fi
      }
      \fi
  \end{center}%
  \vspace*{\afterkwdskip}
  \par
  }
\newlength{\afterkwdskip}\setlength{\afterkwdskip}{\z@}
%    \end{macrocode}
% \end{macro}
% \end{macro}
%
%
% \begin{macro}{\reviewarticle}
% \begin{macro}{\reviewtoc}
% \begin{macro}{\reviewlof}
%   User command to specify a review article. Must be used in
%   preamble. Main effect is to place the text ``Review'' before the
%   article title.  Also add command for setting a table of contents
%   and a list of figures at the beginning of the article.
%    \begin{macrocode}
\newcommand{\reviewarticle}{\@reviewarticletrue}
\newcommand\reviewtoc{%
  \section*{CONTENTS}%
  \@starttoc{rtc}%
}
\newcommand\reviewlof{%
  \section*{LIST OF FIGURES}%
  \@starttoc{lof}%
}
%    \end{macrocode}
% \end{macro}
% \end{macro}
% \end{macro}
% \begin{macro}{\l@section}
%    \begin{macrocode}
\renewcommand*\l@section[2]{%
  \ifnum \c@tocdepth >\z@
    \addpenalty\@secpenalty
    \addvspace{0.8em \@plus\p@}%
    \setlength\@tempdima{1.5em}%
    \begingroup
      \parindent \z@ \rightskip \@pnumwidth
      \parfillskip -\@pnumwidth
      \leavevmode %\bfseries
      \advance\leftskip\@tempdima
      \hskip -\leftskip
      #1\nobreak\hfil \nobreak\hb@xt@\@pnumwidth{\hss #2}\par
    \endgroup
  \fi}
%    \end{macrocode}
% \end{macro}
%
% \begin{macro}{\@MSmaketitle}
% A special version of |\@maketitle| for the \texttt{manuscript}
% option. This can deal with a multi-page abstract/resumen. Not very
% well, mind. 
%    \begin{macrocode}
\newcommand{\@MSmaketitle}{%
  \if@debug\ClassInfo{rmaa}{Using special MS maketitle command}\fi
  \newpage
  {
    \thispagestyle{rmaafirstpage}
%     \setlength{\parindent}{0pt}
    \null\vskip 2ex
    \begin{center}
      \@title\par
      \@author\par
      \rmaa@ReceivedStyle{%
        \ifx\@received\@rmaaEmpty
          Draft version: \@date
        \else
          Received: \@received
          \ifx\@accepted\@rmaaEmpty
          \else
            ; Accepted: \@accepted
          \fi
        \fi}\par
    \end{center}
    \@thanks
    \centerline{\rmaa@AbstractHeaderStyle{Resumen}}
    \noindent\@resumen\par\vskip 2ex
    \centerline{\rmaa@AbstractHeaderStyle{Abstract}}
    \noindent\@abstract\par\vskip 2ex
    \noindent\@keywordheader~\rmaa@KeywordStyle{\@keywords}\par
%     \@keywordheader~\rmaa@KeywordStyle{\uppercase{\@keywords}}\par
  }%
  }
%    \end{macrocode}
% \end{macro}

% \subsection{The sectioning commands}
% \label{sec:section-commands}

% \begin{macro}{\section}
% \begin{macro}{\subsection}
% \begin{macro}{\subsubsection}
% \begin{macro}{\paragraph}
% \begin{macro}{\subparagraph}
% Redefine all the standard commands to use these lengths. 
%    \begin{macrocode}
\renewcommand\section{\@startsection {section}{1}{\z@}%
                                   {\abovesectionskip}%
                                   {\belowsectionskip}%
                                   %% how do we get upper case here
                                   {\centering\reset@font\normalsize}}
\renewcommand\subsection{\@startsection{subsection}{2}{\z@}%
                                     {\abovesectionskip}%
                                     {\belowsectionskip}%
                                     {\centering\reset@font\normalsize\itshape}}
\renewcommand\subsubsection{\@startsection{subsubsection}{3}{\z@}%
                                     {\abovesectionskip}%
                                     {\belowsectionskip}%
                                     {\centering\reset@font\normalsize\itshape}}
\renewcommand\paragraph{\@startsection{paragraph}{4}{\z@}%
                                    {\aboveparaskip}%
                                    {\belowparaskip}%
                                    {\reset@font\normalsize\bfseries}}
\renewcommand\subparagraph{\@startsection{subparagraph}{5}{\parindent}%
                                      {\aboveparaskip}%
                                      {\belowparaskip}%
                                      {\reset@font\normalsize\bfseries}}
%    \end{macrocode}
% \end{macro}
% \end{macro}
% \end{macro}
% \end{macro}
% \end{macro}

% \begin{macro}{secnumdepth}
% Make sure that we number the subsubsections---this is done by
% default in the |article| class but not the |book| class, which is
% why we must set it explicitly here.
%    \begin{macrocode}
\setcounter{secnumdepth}{3}
%    \end{macrocode}
% \end{macro}
% \changes{2.2beta10}{2001/11/07}{Fixed \cs{secnumdepth} for \cs{book}
% style.}

% \begin{macro}{\thesubsection}
% \begin{macro}{\thesubsubsection}
% For subsection and subsubsection, we want the numbers to be roman
% but the title to be italic. This seems the easiest way to go about it.
%    \begin{macrocode}
\renewcommand{\thesubsection}{{\upshape\thesection.\@arabic\c@subsection}}
\renewcommand{\thesubsubsection}{\thesubsection{\upshape.\@arabic\c@subsubsection}}
%    \end{macrocode}
% \end{macro}
% \end{macro}

% \begin{macro}{\@seccntformat}
% Add ``.'' to the number here so we can remove it from |\thesection|,
% etc. Reduce gap between number and title WJH 27 Jan 1995
%    \begin{macrocode}
\def\@seccntformat#1{\csname the#1\endcsname .\hskip 0.4em}
%    \end{macrocode}
% \end{macro}

% \begin{macro}{\@sect}
% I didn't really want to mess with the |\@sect| command, but it was
% the only way I could manage to force upper case. Other than the
% commented parts, it is identical to the standard version. 
%    \begin{macrocode}
\def\@sect#1#2#3#4#5#6[#7]#8{
  \ifnum #2>\c@secnumdepth
     \let\@svsec\@empty
  \else
     \refstepcounter{#1}%
     \let\@@protect\protect
     \def\protect{\noexpand\protect\noexpand}%
     \edef\@svsec{\@seccntformat{#1}}%
     \let\protect\@@protect
  \fi
%    \end{macrocode}
% This is the bit I changed (27 Jan 1995). Use upper case for
% |\section|, use mixed case for the rest. 
%    \begin{macrocode}
  \ifnum #2=1                   
     \def\sec@upcase##1{\uppercase{##1}}%
     \def\sec@shape{\upshape}%
  \else                         
     \def\sec@upcase##1{\relax{##1}}%
     \def\sec@shape{\itshape}%
  \fi
  \@tempskipa #5\relax
  \ifdim \@tempskipa>\z@
     \begingroup #6\relax
          \@hangfrom{\hskip #3\relax\@svsec}%
                    {\interlinepenalty \@M \sec@upcase{#8}\par}%
     \endgroup
     \csname #1mark\endcsname{#7}%
%    \end{macrocode}
% Only write TOC entries if doing the Guide style. 
%    \begin{macrocode}
     \if@guide 
        \addcontentsline{toc}{#1}{%
          \ifnum #2>\c@secnumdepth \else
             \protect\numberline{\csname the#1\endcsname}%
          \fi
          #7}%
     \fi
%    \end{macrocode}
% Also allow per-article TOC for review papers
%    \begin{macrocode}
     \if@reviewarticle
        \ifnum #2>2 \else
           \addcontentsline{rtc}{#1}{%
             \protect\numberline{\csname the#1\endcsname}%
             \sec@upcase{\sec@shape #7}}%
        \fi
     \fi
  \else
     \def\@svsechd{#6\hskip #3\relax  %% \relax added 2 May 90
       \@svsec #8\csname #1mark\endcsname{#7}%
       \if@guide
          \addcontentsline{toc}{#1}{%
            \ifnum #2>\c@secnumdepth \else
               \protect\numberline{\csname the#1\endcsname}%
            \fi
            #7}%
       \fi
       }%
  \fi
  \@xsect{#5}}
%    \end{macrocode}
% \end{macro}


% \subsection{Appendices}
% \label{sec:appendices}
% Make sure that appendices in individual articles work properly. We
% copy the definition from the article class (the book one doesn't do
% what we want).  
% \begin{environment}{appendices}
%   This uses an environment |appendices| inside of which each
%   |\section| will become a separate appendix.  Note that the
%   appendices must come before the references (otherwise the full
%   addresses mechanism will not work properly). We make the equation
%   numbering be subordinate to the appendix (equations A12, B6,
%   etc). Finally, we have a medskip followed by APPENDICES to
%   introduce them. \changes{3.17}{2002/04/08}{ Added appendices environment
%   that does the right thing.}
%    \begin{macrocode}
\newenvironment{appendices}{\par
  \begingroup
  \setcounter{section}{0}%
  \setcounter{subsection}{0}%
  \renewcommand\thesection{\@Alph\c@section}
  \renewcommand\theequation{\thesection\arabic{equation}}
  \medskip
  \section*{APPENDICES}}
{\endgroup}
%    \end{macrocode}
% \end{environment}
% \begin{environment}{appendix}
%   If there is only one appendix, then use |appendix| environment
%   instead.  \changes{4.0}{2019/11/12}{ Added
%   appendix environment for case where there is only one section.}
%    \begin{macrocode}
\renewenvironment{appendix}{\par
  \begingroup
  \setcounter{section}{0}%
  \setcounter{subsection}{0}%
  \renewcommand\thesection{\@Alph\c@section}
  \renewcommand\theequation{\thesection\arabic{equation}}
  \medskip
  \section*{APPENDIX}}
{\endgroup}
%    \end{macrocode}
% \end{environment}

% \subsection{Floating thingies}
% \label{sec:floats}   

%% 
%    \begin{macrocode}
\renewcommand{\thetable}{\@arabic\c@table}
\renewcommand{\thefigure}{\@arabic\c@figure}
\renewcommand{\figurename}{Fig.}
\renewcommand{\tablename}{TABLE}
%\renewcommand{\fnum@table}{\MakeUpperCase{\tablename}~\thetable}
%    \end{macrocode}
% For some reason, I have to write `TABLE' explicitly here
% If I use |\tablename|, it gets converted to mixed case---God knows
% why. 
% |\figurename| is also being ignored. 
%    \begin{macrocode}
\renewcommand{\fnum@table}{TABLE~\thetable}
\renewcommand{\fnum@figure}{Fig.~\thefigure}
%    \end{macrocode}

% Do captions in 9 pt face in the main journal style and in new-style
% anything. 
%    \begin{macrocode}
\long\def\@makefigcap#1#2{%
  \if@proceedings\else\small\fi
  \if@newstyle\small\fi 
  \vskip\abovecaptionskip
  \sbox\@tempboxa{#1. #2}%
  \ifdim \wd\@tempboxa >\hsize
    #1. #2 \par
  \else
    \global \@minipagefalse
    \hb@xt@\hsize{\hfil\box\@tempboxa\hfil}%
  \fi
  \vskip\belowcaptionskip}
%    \end{macrocode}

% Make sure table caption is upper case. 
%    \begin{macrocode}
\long\def\@maketabcap#1#2{%
  \hbox to\hsize{\hfil #1 \hfil}
  \vskip 6pt 
  \centering\MakeUpperCase{#2}\par
  \vskip 6pt}
%    \end{macrocode}

%    \begin{macrocode}
\newlength{\@rottabcapwidth}
\long\def\@makerottabcap#1#2{%
  \setbox\@tempboxa\hbox{\MakeUpperCase{#2}}%
  \ifdim\wd\@tempboxa >.8\vsize
     \setlength\@rottabcapwidth{.8\vsize}
  \else
     \setlength\@rottabcapwidth{\wd\@tempboxa}
  \fi
  \rotatebox{90}{%
    \begin{minipage}{\vsize}\centering
      \begin{minipage}{\@rottabcapwidth}\centering
        #1\\ \vskip 6pt \MakeUpperCase{#2}
      \end{minipage}%
    \end{minipage}
    }%
  \hskip 6pt
  }

\newenvironment{Sideways}
    {\begin{sideways}\begin{minipage}{\vsize}\centering}
    {\end{minipage}\end{sideways}}

\newenvironment{Sidewaystable}
{%
  \let\oldcaption=\caption%
  \def\caption{
    \def\@captype{table}\renewcommand{\@makecaption}{\@maketabcap}
    \oldcaption}
  \begin{sidewaystable}\centering}
  {\end{sidewaystable}\let\caption=\oldcaption}
\newenvironment{Sidewaystable*}
{%
  \let\oldcaption=\caption%
  \def\caption{
    \def\@captype{table}\renewcommand{\@makecaption}{\@maketabcap}
    \oldcaption}
  \begin{sidewaystable*}\centering}
  {\end{sidewaystable*}\let\caption=\oldcaption}
%    \end{macrocode}

% Not sure why we need to stick these lines in. 
%    \begin{macrocode}
\renewenvironment{figure}
                 {\renewcommand{\@makecaption}{\@makefigcap}\@float{figure}}
                 {\end@float}
\renewenvironment{figure*}
                 {\renewcommand{\@makecaption}{\@makefigcap}\@dblfloat{figure}}
                 {\end@dblfloat}
\renewenvironment{table}
                 {\renewcommand{\@makecaption}{\@maketabcap}\@float{table}}
                 {\end@float}
\renewenvironment{table*}
                 {\renewcommand{\@makecaption}{\@maketabcap}\@dblfloat{table}}
                 {\end@dblfloat}
%    \end{macrocode}

% Redefine |\@caption| so that we don't bother with the list of
% figures. This saves us having to worry about fragile commands in the
% captions. 
%    \begin{macrocode}
\long\def\@caption#1[#2]#3{%
  \par
%%  \addcontentsline{\csname ext@#1\endcsname}{#1}%
%%    {\protect\numberline{\csname the#1\endcsname}{\ignorespaces #2}}%
  \begingroup
    \@parboxrestore
    \if@minipage
      \@setminipage
    \fi
    \normalsize
    \@makecaption{\csname fnum@#1\endcsname}{\ignorespaces #3}\par
  \endgroup}
%    \end{macrocode}

% Make tables not quite so cramped. Allow |\hline\hline| to give a
% poor man's thick rule. 
%    \begin{macrocode}
\renewcommand{\arraystretch}{1.15}
\setlength{\doublerulesep}{0.4pt}
%    \end{macrocode}

% \begin{macro}{\tablecols}
% \begin{macro}{\@tablecols}
% \begin{macro}{\tabcaptioncontinued}
% |\tablecols| is a user command for setting the number of columns. It
% is used by commands that try and span the entire table (I think
% there must be a better way of acheiving
% this). |\tabcaptioncontinued| is 
% a user command for setting the caption of a continued table on the
% next page. 
%    \begin{macrocode}
\newcommand{\tablecols}[1]{\renewcommand{\@tablecols}{#1}}
\newcommand{\@tablecols}{1}
\newcommand{\tabcaptioncontinued}{%
  \multicolumn{\@tablecols}{c}{\addtocounter{table}{-1}TABLE~\thetable\ (CONTINUED)} \\ \\
  }
%    \end{macrocode}
% \end{macro}
% \end{macro}
% \end{macro}

% \subsubsection{Footnotes in tables}
% \label{sec:table-footnotes}
% \begin{macro}{\tablenotemark}
% \begin{macro}{\tablenote}
% \begin{macro}{\tablenoteskip}
% Old user commands |\tablenotemark{MARK}|,
% |\tablenote{MARK}{TEXT}|---now deprecated. 
%    \begin{macrocode}
\newcommand\tablenotemark[1]{$^\mathrm{#1}$} 
\newcommand\tablenote[2]{\if@proceedings\else\small\fi ${}^\mathrm{#1}$ #2} 
\newcommand\tablenoteskip{\\[0.5\baselineskip]} 
%    \end{macrocode}
% \end{macro}
% \end{macro}
% \end{macro}

% \begin{macro}{\tabnotemark}
% \begin{macro}{\tabnotetext}
% \begin{macro}{\tabnotetextTWO}
% Alternative way of doing footnotes in tables. 
%    \begin{macrocode}
\newcommand{\tabnotemark}[1]{\textsuperscript{\lowercase{#1}}}
%\newcommand{\tabnotetextTWO}[4]{\tabnotetext{#1}{#2}\tabnotetext{#3}{#4}}
\newcommand{\tabnotetextTWO}[4]{%
  \multicolumn{\@tablecols}{l}{%
    \parbox{0.5\tabnotewidth}{\small\textsuperscript{#1}#2}%
    \parbox{0.5\tabnotewidth}{\small\textsuperscript{#3}#4}%
    }%
  \\}
\newcommand{\tabnotetext}[2]{%
  \multicolumn{\@tablecols}{@{}p{\tabnotewidth}@{}}{\small\textsuperscript{#1}#2}\\}
\newlength{\tabnotewidth}
%    \end{macrocode}
% \end{macro}
% \end{macro}
% \end{macro}


% \begin{macro}{changemargin}
%    \begin{macrocode}
% Command to allow margins to be changed on the fly. Useful to allow
% tables and Figures to protrude into the margins. 
%
% \changes{3.24}{2007/09/14}{Added \texttt{changemargin} environment
% (copied from TeX FAQ).}
\newenvironment{changemargin}[2]{% 
  % From TeX FAQ http://www.tex.ac.uk/cgi-bin/texfaq2html?label=chngmargonfly
  \begin{list}{}{%
    \setlength{\topsep}{0pt}%
    \setlength{\leftmargin}{#1}%
    \setlength{\rightmargin}{#2}%
    \setlength{\listparindent}{\parindent}%
    \setlength{\itemindent}{\parindent}%
    \setlength{\parsep}{\parskip}%
  }%
  \item[]}{\end{list}}
%    \end{macrocode}
% \end{macro}

% \subsection{The bibliography and references}
% \label{sec:biblio}
% \begin{macro}{\beforebiblioskip}
% \begin{macro}{\afterbiblioskip}
% I'm not entirely sure why I introduced these two commands, since
% directly putting |\vspace| before and after
% |\begin{thebibliography}| seems to work well enough. 
%    \begin{macrocode}
\newlength{\beforebiblioskip}
\setlength{\beforebiblioskip}{0pt}
\newlength{\afterbiblioskip}
\setlength{\afterbiblioskip}{0pt}
%    \end{macrocode}
% \end{macro}
% \end{macro}
% \begin{macro}{thebibliography}
% The environment itself is done as starred section with extra
% vertical space after heading. This space is currently zero. 
% In the book style we do the cunning trick with rereading the file in
% order to get the references into the |\ADS@references| macro.
% 
% \changes{3.05}{2003/01/02}{Removed the \cs{beforebiblioskip} and
% \cs{afterbiblioskip} commands, since they cause unwanted vertical
% space in some instances.}  
%    \begin{macrocode}
\renewenvironment{thebibliography}{%
  \if@book\RereadBiblio\fi
%%  \vspace*{\beforebiblioskip}
  \section*{REFERENCES} \frenchspacing
%%  \vspace*{\afterbiblioskip}
  \small
%    \end{macrocode} 
% The references themselves are set as a list in order 
% to get the hanging indentation. They are done in a smaller font and
% with french spacing so we needn't worry about the distinction
% between sentence-ending periods and abbreviations. 
%    \begin{macrocode} 
  \list{\null}{\leftmargin 1.5em\labelwidth\z@\labelsep\z@\itemindent -1.5em
    \itemsep\z@\parsep\z@\usecounter{enumi}}
  \sloppy\clubpenalty4000\widowpenalty4000
  \sfcode`\.=1000\relax\@biblisttrue}
{\endlist\@biblistfalse\normalsize\@outputaddressesifnecessary}
%    \end{macrocode}
% \end{macro}
% \subsection{Uniquification of cross-references}
% \label{sec:uniquification}
% 
% \begin{macro}{\articlename}
% In order that the |\label{}| command in individual articles be
% guaranteed to give a unique anchor, we append the article name (held
% in the |\articlename| command) to each key. This requires
% redefinition of the following commands: |\cite|, |\nocite|, |\ref|,
% |\pageref|, |\label|. 
% \begin{macro}{\oldcite}
% \begin{macro}{\oldnocite}
% \begin{macro}{\oldref}
% \begin{macro}{\oldpageref}
% \begin{macro}{\oldlabel}
% The original versions of these commands are
% saved as |\oldcite|, |\oldnocite|, etc. 
%
% The uniquifying of cross-references is now always done, so 
% we don't get multiply-defined labels in the single article case
%    \begin{macrocode}
\newcommand{\articlename}{anon}
\AtBeginDocument{
  \let\oldcite=\cite 
  \let\oldnocite=\nocite 
  \let\oldref=\ref 
  \let\oldpageref=\pageref
  \let\oldlabel=\label
}
%    \end{macrocode}
% \end{macro}
% \end{macro}
% \end{macro}
% \end{macro}
% \end{macro}
% \end{macro}
%
% Uniquify all cross-references by prepending the article name.  This
% now has to go into |\AtBeginDocument| so it is working on the new
% versions that hyperref.sty has set up.  Then we need to put the
% |\@begofarticlehook| afterwards, so that the first page label is using
% the new definitions.  If amsmath is loaded, then we also need to patch
% the special label command that they use. \changes{4.0}{2019/11/18}{Fix labels for amsmath display equations, in case that package has been loaded.}
%    \begin{macrocode}
\AtBeginDocument{
  \renewcommand{\ref}[1]{\oldref{\articlename:#1}%
    \if@debug{\ClassInfo{rmaa}{Reference to \articlename:#1}}\fi}
  \renewcommand{\pageref}[1]{\oldpageref{\articlename:#1}%
    \if@debug{\ClassInfo{rmaa}{Page reference to \articlename:#1}}\fi}
  \renewcommand{\label}[1]{\oldlabel{\articlename:#1}%
    \if@debug{\ClassInfo{rmaa}{Label for \articlename:#1}}\fi}%
  \ifdef{\label@in@display}{%
    \let\oldlabel@in@display=\label@in@display
    \renewcommand{\label@in@display}[1]{\oldlabel@in@display{\articlename:#1}%
      \if@debug{\ClassInfo{rmaa}{AMSTeX Label for \articlename:#1}}\fi}%
  }{}
  \@begofarticlehook
}%
%    \end{macrocode}
% Citations are more complicated since they can contain a list of
% keys. These are really only here for use with the rmaacite package,
% which is utterly deprecated (02 Jan 2003) now that I have the natbib
% code working. The command \cs{cite} should \emph{not} be used with
% natbib: use \cs{citep}, \cs{citet}, etc.\@ instead. 
%    \begin{macrocode}
\renewcommand{\nocite}[1]{\@bsphack
  \@for\@citeb:=#1\do{%
    \edef\@citeb{\expandafter\@firstofone\@citeb:\articlename}%
    \if@filesw\immediate\write\@auxout{\string\citation{\@citeb}}\fi
    \@ifundefined{b@\@citeb:\articlename}{\G@refundefinedtrue
      \@latex@warning{Citation `\@citeb' undefined}}{}}%
  \@esphack}
\renewcommand{\cite}[1]{\oldcite{#1:\articlename}}
%   \end{macrocode}


% Horrible fiddling with the bibliography citation mechanism. 
% This has to be globally available in case rmaacite puts it in the
% \texttt{.aux} file. 
%    \begin{macrocode}
\def\lbibcite#1#2#3{\global\@namedef{b@#1}{#2}
                    \global\@namedef{by@#1}{#3}}
%    \end{macrocode}

% \begin{macro}{\oldbibitem}
% \begin{macro}{\@rmaaBibArg}
% \begin{macro}{\@rmaaSpace}
% \begin{macro}{rmaabib}
% Of course, |\bibitem| needs to get the uniquification treatment as
% well. While we are at it, we redefine it to use consecutive integers
% for the key if the user leaves it blank (using the counter
% |rmaabib|). This allows one to put 
% |\bibitem{}| for all your references without getting annoying
% messages about multiply-defined labels.
% 
% We also take the opportunity to insert the optional argument to
% \cs{bibitem} if it is missing in the form Author(DATE) so that natbib
% doesn't try to do numerical citations. Natbib already provides a
% hook \cs{@extra@b@citeb} to help uniquify the citation key, so we
% just plug into that. However, it puts it \emph{after} the original
% key, not before as I had had it. Easiest solution is to change the
% way that I do it. 02 Jan 2003. Note that the \cs{@extra@b@citeb}
% does not get written to the .aux  file in the \cs{citation} command
% but it does work. 
% \changes{3.14}{06 Feb 2003}{Reinstalled old definition of \cs{bibitem}
%  for the case where the nonatbib option is requested, so that the
%  rmaacite package will work. This means that you can't use rmaacite
%  and natbib in the same volume.}
%    \begin{macrocode}
\if@natbib\def\@extra@b@citeb{:\articlename}\fi
%%\else
\let\oldbibitem\bibitem
\newcommand{\@rmaaBibArg}{}
\newcommand{\@rmaaSpace}{ }
\newcounter{rmaabib}
\if@natbib
\renewcommand{\bibitem}[2][ABC(2003)]{%
  \renewcommand{\@rmaaBibArg}{#2}
  \ifx\@rmaaBibArg\@rmaaEmpty
     \addtocounter{rmaabib}{1}
     \oldbibitem[#1]{\thermaabib:\articlename}
  \else\ifx\@rmaaBibArg\@rmaaSpace
     \addtocounter{rmaabib}{1}
     \oldbibitem[#1]{\thermaabib:\articlename}
  \else 
     \oldbibitem[#1]{#2:\articlename}
  \fi\fi
  \renewcommand{\@rmaaBibArg}{}
  }
\else
\renewcommand{\bibitem}[1]{%
  \renewcommand{\@rmaaBibArg}{#1}
  \ifx\@rmaaBibArg\@rmaaEmpty
     \addtocounter{rmaabib}{1}
     \oldbibitem{\articlename:\thermaabib}
  \else\ifx\@rmaaBibArg\@rmaaSpace
     \addtocounter{rmaabib}{1}
     \oldbibitem{\articlename:\thermaabib}
  \else
     \oldbibitem{\articlename:#1}
  \fi\fi
  \renewcommand{\@rmaaBibArg}{}
  }
\fi
%%\fi
%    \end{macrocode}
% \end{macro}
% \end{macro}
% \end{macro}
% \end{macro}

% \begin{macro}{oldpcite}
% \begin{macro}{oldscite}
% Then, we need to repeat the above scheme, but for the |rmaacite|
% definition. 
%    \begin{macrocode}
\newcommand{\@redefinebibitem}{%
  \let\oldbibitem\bibitem
  \let\oldcite\cite
  \let\oldpcite\pcite
  \let\oldscite\scite
  \def\bibitem[##1<##2>]##3{%
    \renewcommand{\@rmaaBibArg}{##3}
    \ifx\@rmaaBibArg\@rmaaEmpty
       \addtocounter{rmaabib}{1}
       \oldbibitem[##1<##2>]{\thermaabib:\articlename}
    \else\ifx\@rmaaBibArg\@rmaaSpace
       \addtocounter{rmaabib}{1}
       \oldbibitem[##1<##2>]{\thermaabib:\articlename}
    \else 
       \oldbibitem[##1<##2>]{##3:\articlename}
    % trust the author here
%%       \oldbibitem[##1<##2>]{##3}
    \fi\fi
    }
  }
%    \end{macrocode}
% \end{macro}
% \end{macro}


% \subsection{HTML Table of Contents}
% \label{sec:hyper-TOC}
% \begin{macro}{\HyperTOCEntry}
% \begin{macro}{\HyperTOCpart}
% \begin{macro}{\HyperTOCmeeting}
% \begin{macro}{\HyperTOCeditors}
% \begin{macro}{\HyperTOCpreamble}
% These commands are used to write information to the \texttt{.htoc}
% file, which is used in generating a Web Page of the Table of
% Contents. Since they only ever appear preceded by |\protect| in the
% argument to an |\addtocontents| command,  they should never need to
% be expanded by the macros. However, we define them here just in case
% \dots
%    \begin{macrocode}
\newcommand{\HyperTOCEntry}[3]{\relax}
\newcommand{\HyperTOCpart}[1]{\relax}
\newcommand{\HyperTOCmeeting}[3]{\relax}
\newcommand{\HyperTOCeditors}[1]{\relax}
\newcommand{\HyperTOCpreamble}{\relax}
%    \end{macrocode}
% \end{macro}
% \end{macro}
% \end{macro}
% \end{macro}
% \end{macro}
%

% \subsection{Formatting the poster abstracts} 
% These can now be used without using the book option.
%    \begin{macrocode}

\newcommand{\@makeminititle}{%
  \newpage
  \null 
  \vskip 3em%
  \begin{center}%
    {\@title \par}%
  \end{center}
  \vskip 3em%
  }

\newif\if@abstractsads
\@abstractsadsfalse
\newcommand\abstractsADS{\@abstractsadstrue}

\if@book\RequirePackage{multicol}\fi
\newenvironment{abstracts}[1][Abstracts of Contributed Papers]{%
  %% First, make sure we start on an odd page (06 Jun 2000)
  \if@book\clearpage\thispagestyle{empty}\cleardoublepage\fi
  %%
  \renewcommand{\articlename}{ABSTRACTS}
  \title{#1}\shortauthor{Abstracts}\shorttitle{Abstracts}
  \listofauthors{}
  \if@book
  \addcontentsline{toc}{chapter}{\protect\@fiddleauthors{\@TOCtitle}{\rmaaTOCauthors}}%
  \addtocontents{htc}{\protect\HyperTOCEntry{\@rawtitle}{\@rawauthors}{\articlename}}
  \fi
  \pagestyle{rmaa} % this now seems necessary (06 Jun 2000)
  \thispagestyle{rmaafirstpage}
%%  \twocolumn[\@makeminititle]
  \renewcommand{\abstract}[1]{%
    \renewcommand\ADS@abstract{##1}%
    \renewcommand{\@abstract}{##1}\@writeabstract%
  }
  \renewcommand{\title}[2][]{%
    \gdef\ADS@authors{}%
    \gdef\ADS@affil{}%
    \gdef\ADS@references{}%
    \gdef\ADS@keywords{}%
    \renewcommand{\@rmaaArg}{##1}
    \ifx\@rmaaArg\@rmaaEmpty
       \renewcommand{\@TOCaftertitle}{}
    \else
       \renewcommand{\@TOCaftertitle}{ (\@rmaaArg)}
    \fi
    \renewcommand{\@title}{\rmaa@TOCTitleStyle{##2}}%
    \renewcommand{\@TOCtitle}{\rmaa@TOCTitleStyle{##2}\@TOCaftertitle}%
    \renewcommand{\rmaaIndexTitle}{\rmaa@IndexTitleStyle{##2\@TOCaftertitle.}}%
    \renewcommand\ADS@title{##2}%
    \xdef\ADS@firstpage{\@arabic{\c@page}}%
    }
  \renewcommand{\author}[1]{\let\@thanks\@empty\renewcommand\@author{##1} }
  \renewcommand{\footnotetext}[2][XX]{\par\noindent\textsuperscript{##1}~##2}
  \renewcommand{\footnoterule}{\kern-3\p@ \hrule \@width 3em \kern 2.6\p@}
  \if@book\@begofarticlehook\fi
  \if@book\begin{multicols}{2}[\@makeminititle]\else\@makeminititle\fi
  \begingroup
  \setcounter{footnote}{0}
  }% and at the end of the environment....
  {\endgroup\if@book\end{multicols}\listofabstracts\label{RMAALastPage}\clearpage\fi}
%% auxiliary commands used:
%% \@writeabstract does all the actual formatting 
\newcommand{\@writeabstract}{%
  \if@book%
  \addcontentsline{loa}{abstract}{\protect\@fiddleauthors{\@TOCtitle}{\rmaaTOCauthors}}%
  \fi%
  \begin{center}\@title\par\@author\end{center}
  \@abstract
  \par\noindent\footnoterule\nobreak\footnotesize\@thanks
  \par\noindent\hrulefill\goodbreak
  \vskip 1cm \@plus0.5cm
  %% write ADS entry
  \if@book\if@abstractsads\ADS@dumpdata\fi\fi%
  \endgroup % this pair is to make each abstract a self-contained group
  \begingroup
  \setcounter{footnote}{0}
  }

\if@book
%% List of abstracts (stuff written to .loa file)
\newcommand{\listofabstracts}{%
  \begin{multicols}{2}[%
    \vskip 1.5cm \vfill % leave at least 1.5 cm and go to bottom of page if possible
    \centerline{\MakeUpperCase{List of Abstracts}}
    \vskip 1.5cm \@plus0.5cm
    ][6cm] % this makes it start a new page if there is not at least 6cm of space
    \renewcommand{\@fiddleauthors}[2]{##1\hskip1em plus2em minus0.4em\penalty-100  ##2}
    \parindent\z@
%%    \parskip.5\p@ \@plus .7\p@\relax 
    \parskip3\p@ \@plus 1\p@ \@minus 1\p@\relax % allow quite a lot of
                                % leeway here to avoid breaks inside
                                % entries 
    \renewcommand{\\}{\relax} % temporarily turn off explicit line breaks just in case
    \renewcommand{\altaffilmark}[1]{\relax} % ditto with this
    \@starttoc{loa}
  \end{multicols}
  }
\newcommand*{\l@abstract}[2]{%
  \par\penalty-1000
  #1\hskip 1em\penalty0\null\nobreak\hfill#2}
\fi

%    \end{macrocode}
 
% \subsection{Code for typesetting an entire volume of conference
% proceedings} \label{sec:whole-volume}
%    \begin{macrocode}

\if@book
%% environment for including the individual articles
\newcommand{\@fiddleauthors}[2]{\relax}
\renewcommand{\theequation}{\@arabic\c@equation}
\newenvironment{articles}%
{ \begingroup % make it a group so that author definitions do not
              % spill over
  \gdef\ADS@authors{}
  \pagestyle{rmaa}
  \renewenvironment{document}{%
    \@begofarticlehook
    \addtocounter{chapter}{1}%
    \addcontentsline{toc}{chapter}{%
      \protect\@fiddleauthors{\@TOCtitle}{\rmaaTOCauthors}}%
    \addtocontents{htc}{\protect\HyperTOCEntry{\@rawtitle}{\@rawauthors}{\articlename}}
    }{\relax}
  \renewcommand{\documentclass}[2][foo]{\relax} 
  \renewcommand{\usepackage}[2][foo]{%
    \def\@packagename{##2}\def\@rmaacitename{rmaacite}%
    \ifx\@packagename\@rmaacitename
    \@rmcitetrue\ClassInfo{rmaa}{This paper uses rmaacite}
    \makeatletter\input{rmaacite.sty}\makeatother % read it in here so that all changes
    \else                                         % are local to paper
    \if@debug\ClassWarning{rmaa}{Package \@packagename\space was requested by \articlename\space}\fi
    \fi} 
  %% Don't allow individual articles to mess with the pages, volume or year
  \renewcommand{\SetFirstPage}[1]{\relax}
  \renewcommand{\SetYear}[1]{\relax}
  \renewcommand{\SetVolume}[1]{\relax}
  \renewcommand{\SetConfTitle}[1]{\relax}
  %% Don't allow individual articles to change these
  \renewcommand{\nonstopmode}{\relax}
  \renewcommand{\scrollmode}{\relax}
  \renewcommand{\outputplates}{\relax}
  %% XXXX need to check if using rmaacite and act accordingly
  %% avoid multiple declarations of the same length
  \let\oldnewlength=\newlength
  \renewcommand{\newlength}[1]{%
    \edef\reserved@a{\expandafter\@gobble\string ##1}%
    \@ifundefined{\reserved@a}{%
      \oldnewlength{##1}}{%
      \if@debug{\ClassInfo{rmaa}{\reserved@a\space already defined: doing nothing}}\fi%
      \relax}}%
  \renewcommand\newcommand{\providecommand}
  }%
{\endgroup\onecolumn}
     

%    \end{macrocode}
% Command to be used in `articles' environment for including an
% article
%    \begin{macrocode}
\newcommand{\article}[2][]{%
  % optional first argument is extra text to put in brackets after title in TOC
  % e.g. ``Invited Talk'' or ``Poster''
  \renewcommand{\articlename}{#2}
  % reset loads of counters
  \setcounter{footnote}{0}
  \setcounter{figure}{0}
  \setcounter{table}{0}
  \setcounter{section}{0}
  \setcounter{equation}{0}
  \setcounter{rmaabib}{0}
  \setcounter{author}{0}
  \setcounter{AUTHOR}{1}
  \gdef\ADS@authors{}%
  \gdef\ADS@affil{}%
  \gdef\ADS@affiliations{}%
  \gdef\ADS@references{}%
  \@firstauthortrue
  \@firstaffiltrue
  \@firstaltaffiltrue
  \@firstkwdtrue
  \@rmcitefalse
  \if@suppressfulladdresses\else\@needaddressestrue\fi
  \renewcommand{\@rmaaArg}{#1}
  \ifx\@rmaaArg\@rmaaEmpty
     \renewcommand{\@TOCaftertitle}{}
  \else
     \renewcommand{\@TOCaftertitle}{ (#1)}
  \fi
%    \end{macrocode}
% In the main journal articles must start on an odd page. 
%    \begin{macrocode}
  \if@journal\clearpage\thispagestyle{empty}\cleardoublepage\fi
%% Now input the file, then run the end-of-article hook (reinstated 28
  %% Sep 2001)
  {\input{\@articledir/\articlename.tex}\@endofarticlehook}}
%    \end{macrocode}
% Command to separate articles into sessions. Write an entry in the
% TOC, makes sure we start on an odd page (added 06 Jun 2000), then do
% a title page if \cs{SessionTitlePages} was called. 
% \changes{3.12}{04 Feb 2003}{Added option to have title page for each
% session. Turn it on with \cs{SessionTitlePages}.}
%    \begin{macrocode}
\newcommand{\session}[1]{
  %% this does nothing except write an entry in the TOC
  \addcontentsline{toc}{part}{\rmaa@TOCPartStyle{#1}}
  \addtocontents{htc}{\protect\HyperTOCpart{#1}}
  %% and also makes sure we start on an odd page (added 06 Jun 2000)
  \clearpage\thispagestyle{empty}\cleardoublepage
  \stepcounter{part}
  \if@sessiontitlepages{%
    \onecolumn
    \thispagestyle{empty}
    \renewcommand{\articlename}{Section\thepart TitlePage}
    \label{RMAAFirstPage}
    \make@session@title{#1}
    \label{RMAALastPage}
    \clearpage\thispagestyle{empty}\cleardoublepage
  }\fi
  }
\newcommand{\make@session@title}[1]{%
  {%
    \setlength\parindent{0pt}
    \vspace*{\stretch{1}}
    \rule{\linewidth}{1mm}
    \begin{flushright}\Huge\bfseries\boldmath
      #1
    \end{flushright}
    \rule{\linewidth}{1mm}
    \vspace*{\stretch{2}}
  }
}
%    \end{macrocode}
% Command to turn on separate title pages for each session. 
%    \begin{macrocode}
\newcommand{\SessionTitlePages}{\@sessiontitlepagestrue}
\newcommand{\BoldSessionTitles}{%
  \renewcommand{\rmaa@TOCPartStyle}[1]{\textbf{\rmaa@TOCTitleStyle{##1}}}
  }
\newcommand{\ArticleDir}[1]{\renewcommand{\@articledir}{#1/}\graphicspath{{1/}}}
\newcommand{\@articledir}{.}
%\newcommand{\ps@rmaafront}{\ps@rmaa}% page style for front matter (same as ever, actually)
\newcommand{\ps@rmaafront}{% page style for front matter:  page # in each foot
  \renewcommand\@oddfoot{\hfil{\thepage}}%
  \renewcommand\@evenfoot{{\thepage}\hfil}%
  % shorttitle in both heads
  \renewcommand\@evenhead{\hfil\@shorttitle\hfil}%
  \renewcommand\@oddhead{\@evenhead}
  }
\newcommand{\ps@rmaafrontfirstpage}{% page style for first page of front matter
  % page # in each foot, on outside edge
  \renewcommand\@oddfoot{\hfil\thepage}%
  \renewcommand\@evenfoot{\thepage\hfil}%
  % nothing in the even heads
  \renewcommand\@evenhead{\relax}%
  \renewcommand\@oddhead{\relax}%
  }
%% slight modification to the standard \frontmatter and \mainmatter
%% commands, just to set the \pagestyle
\renewcommand\frontmatter{%
  \latexonecolumn
  \cleardoublepage
  \@frontmattertrue
  \@mainmatterfalse
  \pagenumbering{roman}
  \pagestyle{rmaafront}
  \addtocounter{page}{2} % leave space for the title page
  }
\renewcommand\mainmatter{%
  \pagestyle{empty}
  \if@rmtwocol\twocolumn\fi
  \cleardoublepage
  \@frontmatterfalse
  \@mainmattertrue
  \pagenumbering{arabic}
  \pagestyle{rmaa}
  }
\renewcommand\backmatter{%
  \clearpage\pagestyle{empty}\cleardoublepage
  \@mainmatterfalse
  %% put extra space in the TOC
  \addtocontents{toc}{\bigskip}
  }
%    \end{macrocode}
%
% \subsection{Front sections}\label{sec:front-sections}
% \begin{macro}{\notitle}
% \begin{macro}{\shortname}
% \begin{macro}{frontsection}
% The environment \texttt{frontsection} is for doing a section in the
% front matter (preface, list of participants, or whatever). It has 
% one optional argument:
% \begin{enumerate}
% \item[\#1] Short title (no embedded spaces) for
% use in making PS files. E.g., ``list\_of\_participants''.   
% \end{enumerate}
% and two mandatory arguments: 
% \begin{enumerate}
% \item[\#2] Title of section (e.g., `Preface')
% \item[\#3] Commands to be evaluated before processing the contents
% of the environment. This is mainly useful for putting in the things
% that would go in the preamble of a normal article  (e.g.,
% |\listofauthors|, |\indexauthor|, |\notitle|) 
% \end{enumerate}
%    \begin{macrocode}
\newcommand{\notitle}{\@titlevisiblefalse} \newcommand{\shortname}{}
\newenvironment{frontsection}[3][]{%
  \begingroup
  \renewcommand{\shortname}{#1}
  \ifx\shortname\@rmaaEmpty
     \renewcommand{\articlename}{#2}
  \else
     \renewcommand{\articlename}{\shortname}
  \fi
  \@begofarticlehook
  \@titlevisibletrue % default
  \thispagestyle{rmaafrontfirstpage}
  \listofauthors{} % blank out this in case it is not specified
  \title{{#2}}\shorttitle{#2}\shortauthor{#2}
  #3 % evaluate the command given as second mandatory arg
  \addcontentsline{toc}{chapter}{\protect\@fiddleauthors{\@TOCtitle}{\rmaaTOCauthors}}%
  \addtocontents{htc}{\protect\HyperTOCEntry{\@rawtitle}{\@rawauthors}{\articlename}}
  \if@titlevisible
     \vspace*{1cm}\centerline{\rmaa@FrontSecStyle{\@title}}\vspace*{1cm}
  \fi
  \setlength{\parskip}{3ex}
  }{\frontlabel{RMAALastPage}\clearpage\endgroup}
%    \end{macrocode}
% \end{macro}
% \end{macro}
% \end{macro}
%    \begin{macrocode}
%% Definition of the command \frontlabel, which rights out something like `F4' to
%% the .aux file, instead of `\textlatin  {i}'. This is so it is easy to post-process
%% the file for making PS files of individual sections
\newcommand{\frontlabel}[1]{%
  \@bsphack
  \protected@write\@auxout{}%
     {\string\newlabel{\articlename:#1}{{}{F\@arabic{\c@page}}}}
  \@esphack
  }
%% User command to say X, Y, Z,\\ Editors of the Proceedings
\newcommand{\EditorSignature}{\vspace*{\parskip}\rightline{\@editors}%
  \rightline{Editors of the proceedings}}
%% Fiddle with \thesection, etc. so that they don't include the chapter number
%% Basically go back to the `article' definitions of these
\renewcommand{\thesection}{\@arabic\c@section}
\renewcommand{\thefigure}{\@arabic\c@figure}
\renewcommand{\thetable}{\@arabic\c@table}

\newcommand{\@oldargs}{}
\newcommand{\@newargs}{}
%% add TOC entry for each paper
%\renewcommand{\title}[1]%
%{\renewcommand{\@title}{\rmaa@TitleStyle{#1}}%
%\addtocounter{chapter}{1}\addcontentsline{toc}{chapter}{#1}%
%}

%% Formatting of the table of contents
%% dimensions of various parts
\newlength{\TOCtextwidth}\setlength{\TOCtextwidth}{12.65cm}
\newlength{\TOCpagenumwidth}\setlength{\TOCpagenumwidth}{1.3cm}
\newlength{\TOCentrywidth}\setlength{\TOCentrywidth}{\TOCtextwidth}
\addtolength{\TOCentrywidth}{-\TOCpagenumwidth} 
\newlength{\TOCextramargin}\setlength{\TOCextramargin}{1.25cm}
\newcommand{\TOCbox}[1]{\hspace{\TOCextramargin}\parbox{\TOCtextwidth}{#1}\par}
\newlength{\TOCextramarginx}\setlength{\TOCextramarginx}{\TOCextramargin}\addtolength{\TOCextramarginx}{4.3cm}
\newlength{\TOCtextwidthx}\setlength{\TOCtextwidthx}{\TOCtextwidth}\addtolength{\TOCtextwidthx}{-4.3cm}
\newcommand{\TOCboxx}[1]{\hspace{\TOCextramarginx}\parbox{\TOCtextwidthx}{#1}\par}

\newcommand{\RMAAbannerimage}{rmsc}
\renewcommand\tableofcontents{%
  \renewcommand{\articlename}{CONTENTS}
  \begingroup
  \@begofarticlehook
  \if@twocolumn
     \@restonecoltrue\onecolumn
  \else
     \@restonecolfalse
  \fi
  %% 9pt for journal style
  \if@proceedings\else\small\fi
  \if@newstyle\small\fi %% same for new style
  \if@proceedings\thispagestyle{rmaafrontfirstpage}\else\thispagestyle{empty}\fi
  \shortauthor{Contents}\shorttitle{Contents}
  \vspace*{18mm}
  \if@printbanner
  %% put in the banner header (need to get better image of this)
     \TOCbox{\hspace*{-0.15cm}\includegraphics[width=0.9\TOCtextwidth]{\RMAAbannerimage}}
     \vspace*{8mm}
  %% volume and publication date (should really be in funny font)
     \TOCboxx{\@volumen\\[2mm] \@mes\ \@anyo}
  \else
  %% put in the banner header (need to get better image of this)
     \TOCbox{\hspace*{-0.15cm}\includegraphics[width=0.9\TOCtextwidth]{\RMAAbannerimage}}
     \vspace*{12.5mm}
  %% volume and publication date (should really be in funny font)
     \TOCboxx{~\\[2mm] ~}
  \fi
  \vspace*{16mm}
  %% This is the command that fiddles with the TOC entry so that the
  %% title is flushleft but the authors are flushright
  \renewcommand{\@fiddleauthors}[2]{%
    % first put the paper title
    \renewcommand{\\}{\relax} % temporarily turn off explicit line breaks just in case
    \renewcommand{\altaffilmark}[1]{\relax} % ditto with this
    \renewcommand{\textsuperscript}[1]{\relax} % ditto with this 
    ##1
    % now save the author list in a box
    \sbox\@tempboxa{##2}%
    % check that box is shorter than the total width available for the entry
    \ifdim \wd\@tempboxa <\TOCentrywidth
       % if so, then put it right here, flush against right margin
       \penalty10000\hskip1em plus10em\penalty0\hbox{}\nobreak\hfill\box\@tempboxa
       % Note that we force a space of at least 1em between title and authorlist, 
       % stretchable to 11em since we may have very short titles but still with an 
       % authorlist too long to go on the same line.  The only linebreak allowed is 
       % directly before the authorlist.
    \else
       % if it is too long, we set the original authorlist text in its own parbox
       \par\vskip1ex\parbox[b]{\TOCentrywidth}{\raggedleft ##2}
    \fi
    }
  \@starttoc{toc}%
  \if@restonecol\twocolumn\fi
  \frontlabel{RMAALastPage}
  \clearpage
  \thispagestyle{empty}
  \cleardoublepage % make sure we start on an odd page
  \endgroup

  %%% ALSO, do the HTML TOC
  %% The file that will be used to make the HTML TOC
  \@starttoc{htc}
  %% All the stuff in this file does nothing in the book style
  \addtocontents{htc}{\protect\HyperTOCmeeting{\@conftitle}{\@reunionplace}{\@reuniondate}}
  \addtocontents{htc}{\protect\HyperTOCeditors{\@editors}}
  \addtocontents{htc}{\protect\HyperTOCpreamble}
  \addtocontents{htc}{\protect\HyperTOCEntry{Table of Contents}{}{\articlename}}

  %% and the ADS file
  \ADS@openfile
  }
%% This is now much simpler
\renewcommand\l@chapter[2]{%TOC entry for articles
  \ifnum \c@tocdepth >\m@ne
    \vskip 1.0em \@plus\p@
    \TOCbox{%
      \parbox[b]{\TOCentrywidth}{#1}\parbox[b]{\TOCpagenumwidth}{\hfill#2}
    }
  \fi} 
\renewcommand\l@part[2]{%TOC entry for sessions: centered, no page number
  \ifnum \c@tocdepth >\m@ne
    \vskip 1.5em \@plus1.5em \goodbreak
    \TOCbox{%
      \parbox[b]{\TOCentrywidth}{\hfil#1\hfil}\hfill
    }
  \fi\nobreak} 

%% Formatting theindex
%% strongly encourage page breaks between index entries
\renewcommand{\@idxitem}{\par\penalty-1000\hangindent.45cm} 
\renewcommand{\indexspace}{\par\penalty-1000\vskip 10\p@ \@plus15\p@ \@minus3\p@\relax}
\renewenvironment{theindex}{%
  %% first, make sure we start on an odd page (06 Jun 2000)
  \clearpage\thispagestyle{empty}\cleardoublepage
  %% this looks a bit like an article but is two columns and without abstract, etc
  \renewcommand{\articlename}{INDEX}
  \title{Author Index}\shortauthor{Author Index}\shorttitle{Author Index}
  \listofauthors{}
  \addcontentsline{toc}{chapter}{\protect\@fiddleauthors{\@TOCtitle}{\rmaaTOCauthors}}%
  \addtocontents{htc}{\protect\HyperTOCEntry{\@rawtitle}{\@rawauthors}{\articlename}}
  \pagestyle{rmaa}
  \thispagestyle{rmaafirstpage}
  \@begofarticlehook
  %% 9pt for journal style
  \if@proceedings\else\small\fi
  \if@newstyle\small\fi
  \twocolumn[\@makeminititle]
  %% vertical spacing (see also definition of \@idxitem and \indexspace
  \parskip\z@ \@plus .6\p@\relax % allow quite a lot of leaway here to avoid breaks inside entries
  %% horizontal spacing, breaking and glue parameters (see also the stuff in the definition 
  %% of \indexauthor
  \parindent\z@ 
  \exhyphenpenalty10000 % don't break on explicit hyphens, it always looks bad
  \rightskip=0pt plus.2em % allow a tiny bit of leeway at the ends of lines
  \spaceskip.3333em plus.5em minus.15em % a lot of leeway between words
  \xspaceskip.5em plus.5em minus.4em % this seems to have no effect
  \let\item\@idxitem
  \renewcommand{\\}{\relax} % temporarily turn off explicit line breaks just in case
  \renewcommand{\altaffilmark}[1]{\relax} % ditto with this
  \renewcommand{\textsuperscript}[1]{\relax} % ditto with this
  }{\label{RMAALastPage}\clearpage} % clearpage necessary so that we know about \shortauthor



%% end definitions for book styles
\fi

%    \end{macrocode}
% In order to do away with plates, we sometimes want a figure to
% appear in color only in the online version. The following command
% can bve used in a figure caption to typeset a note to this effect. 
%    \begin{macrocode}
\newcommand\ColorOnlineOnly{%
  \textsc{Note: this figure is available in color in the electronic
    version of this article, obtainable from}
  \texttt{http://\discretionary{}{}{}www.\discretionary{}{}{}astroscu.\discretionary{}{}{}unam.\discretionary{}{}{}mx/\discretionary{}{}{}\~{}rmaa/}.}
%% Plates
%% Define 2 counters: plate is the number in the whole document, plateloop is used 
%% internally by \outputplates
\newcounter{plate}
\newcounter{plateloop}
%% Currently limited to 15 plates in document
%% when we do a label, we need to write to \newlabel entries to the .aux file, one 
%% that says RMAAPlateX so that it can be referred to in \ouputplates, and one that
%% says articlename:RMAAPlateX so that the dvips program can find it and asociate it
%% with the right article
\newcommand{\@platelabeli}{\oldlabel{RMAAPlate1}\label{RMAAPlate1}}
\newcommand{\@platelabelii}{\oldlabel{RMAAPlate2}\label{RMAAPlate2}}
\newcommand{\@platelabeliii}{\oldlabel{RMAAPlate3}\label{RMAAPlate3}}
\newcommand{\@platelabeliv}{\oldlabel{RMAAPlate4}\label{RMAAPlate4}}
\newcommand{\@platelabelv}{\oldlabel{RMAAPlate5}\label{RMAAPlate5}}
\newcommand{\@platelabelvi}{\oldlabel{RMAAPlate6}\label{RMAAPlate6}}
\newcommand{\@platelabelvii}{\oldlabel{RMAAPlate7}\label{RMAAPlate7}}
\newcommand{\@platelabelviii}{\oldlabel{RMAAPlate8}\label{RMAAPlate8}}
\newcommand{\@platelabelix}{\oldlabel{RMAAPlate9}\label{RMAAPlate9}}
\newcommand{\@platelabelx}{\oldlabel{RMAAPlate10}\label{RMAAPlate10}}
\newcommand{\@platelabelxi}{\oldlabel{RMAAPlate11}\label{RMAAPlate11}}
\newcommand{\@platelabelxii}{\oldlabel{RMAAPlate12}\label{RMAAPlate12}}
\newcommand{\@platelabelxiii}{\oldlabel{RMAAPlate13}\label{RMAAPlate13}}
\newcommand{\@platelabelxiv}{\oldlabel{RMAAPlate14}\label{RMAAPlate14}}
\newcommand{\@platelabelxv}{\oldlabel{RMAAPlate15}\label{RMAAPlate15}}
\newcommand{\@outputplatelabeli}{\oldlabel{RMAAOutputPlate1}}
\newcommand{\@outputplatelabelii}{\oldlabel{RMAAOutputPlate2}}
\newcommand{\@outputplatelabeliii}{\oldlabel{RMAAOutputPlate3}}
\newcommand{\@outputplatelabeliv}{\oldlabel{RMAAOutputPlate4}}
\newcommand{\@outputplatelabelv}{\oldlabel{RMAAOutputPlate5}}
\newcommand{\@outputplatelabelvi}{\oldlabel{RMAAOutputPlate6}}
\newcommand{\@outputplatelabelvii}{\oldlabel{RMAAOutputPlate7}}
\newcommand{\@outputplatelabelviii}{\oldlabel{RMAAOutputPlate8}}
\newcommand{\@outputplatelabelix}{\oldlabel{RMAAOutputPlate9}}
\newcommand{\@outputplatelabelx}{\oldlabel{RMAAOutputPlate10}}
\newcommand{\@outputplatelabelxi}{\oldlabel{RMAAOutputPlate11}}
\newcommand{\@outputplatelabelxii}{\oldlabel{RMAAOutputPlate12}}
\newcommand{\@outputplatelabelxiii}{\oldlabel{RMAAOutputPlate13}}
\newcommand{\@outputplatelabelxiv}{\oldlabel{RMAAOutputPlate14}}
\newcommand{\@outputplatelabelxv}{\oldlabel{RMAAOutputPlate15}}
\newcommand{\@platerefi}{\oldpageref{RMAAPlate1}}
\newcommand{\@platerefii}{\oldpageref{RMAAPlate2}}
\newcommand{\@platerefiii}{\oldpageref{RMAAPlate3}}
\newcommand{\@platerefiv}{\oldpageref{RMAAPlate4}}
\newcommand{\@platerefv}{\oldpageref{RMAAPlate5}}
\newcommand{\@platerefvi}{\oldpageref{RMAAPlate6}}
\newcommand{\@platerefvii}{\oldpageref{RMAAPlate7}}
\newcommand{\@platerefviii}{\oldpageref{RMAAPlate8}}
\newcommand{\@platerefix}{\oldpageref{RMAAPlate9}}
\newcommand{\@platerefx}{\oldpageref{RMAAPlate10}}
\newcommand{\@platerefxi}{\oldpageref{RMAAPlate11}}
\newcommand{\@platerefxii}{\oldpageref{RMAAPlate12}}
\newcommand{\@platerefxiii}{\oldpageref{RMAAPlate13}}
\newcommand{\@platerefxiv}{\oldpageref{RMAAPlate14}}
\newcommand{\@platerefxv}{\oldpageref{RMAAPlate15}}
\newcounter{figplatei}
\newcounter{figplateii}
\newcounter{figplateiii}
\newcounter{figplateiv}
\newcounter{figplatev}
\newcounter{figplatevi}
\newcounter{figplatevii}
\newcounter{figplateviii}
\newcounter{figplateix}
\newcounter{figplatex}
\newcounter{figplatexi}
\newcounter{figplatexii}
\newcounter{figplatexiii}
\newcounter{figplatexiv}
\newcounter{figplatexv}
%% Printed output should use the counter `plate'
\renewcommand{\theplate}{\@arabic\c@plate}
\renewcommand{\theplateloop}{\@arabic\c@plateloop}
%% Command to output all the saved plates in this paper
\newcommand{\outputplates}{%
  \setcounter{plateloop}{0}
  \if@book\else\@endofarticlehook\fi % not entirely sure why, but this
                                     % seems to be needed
  \clearpage % first output the last page of text
  % now set pagestyle for blank even pages and make sure we start on an odd page
  \pagestyle{empty}
  \cleardoublepage
  \onecolumn
  \parindent\z@
  % first put in an extra sheet
  % the main purpose of this is so there is always a blank even page before 
  % each plate for filling in the PS file if the article has an odd number
  % of pages (see splitbook.sh)
  \null\vfil\centerline{PLATES}\vfil\clearpage\pagestyle{rmaaplate}\cleardoublepage
  % output each plate in turn
  \@whilenum\value{plateloop}<\value{plate}\do{%
    \addtocounter{plateloop}{1}
    % restore the figure counter from \@figcounti, \@figcountii, etc
    \setcounter{figure}{\value{figplate\roman{plateloop}}}
    % spew out the plate (in its own group for safety)
    {\csname @plate\roman{plateloop}\endcsname\par}
    }
  \pagestyle{rmaa}
  % reset the page number as if nothing had happened
  \addtocounter{page}{-\value{plate}}
  \addtocounter{page}{-\value{plate}}
  }
\newcommand\figcaption{%
  \def\@captype{figure}
  \renewcommand{\@makecaption}{\@makefigcap}
  \caption}
\newcommand\tabcaption{%
  \def\@captype{table}
  \renewcommand{\@makecaption}{\@maketabcap}
  \caption}

\newcommand{\plate}[2][1]{%
  %% User command to make a plate page: 
  %% Mandatory argument is what goes on the plate - this should not 
  %% include any floating environments (e.g. figure). Instead, \figcaption 
  %% should be used to do the captions. 
  %% Optional argument is the number of figures (i.e. captions) on the plate
  %% page (default is 1). This is necessary for incrementing the figure counter. 
  %%
  %% Note that we have to use a \gdef rather than an \xdef when we 
  %% save the plate output in a command. This means that in the book styles when it 
  %% comes to spewing out the plate, we will have forgotten about any user macros 
  %% that may have been defined. This can be got round by repeating the macros inside the plate.
  %% (Presumably, one would use a \newcommand rather than \renewcommand if they were new.)
  % increment global and local counters
  \addtocounter{plate}{1}
  % write appropriate label to .aux file
  \csname @platelabel\roman{plate}\endcsname
  % define a global macro called \@platei, \@plateii, \@plateiii, \@plateiv, ...
  \expandafter\gdef\csname @plate\roman{plate}\endcsname{%
    \csname @outputplatelabel\roman{plateloop}\endcsname
    \null\vfil #2 \vfil
    \cleardoublepage
    }
  % save the current shortauthor and shorttitle
  % need to be \xdef to take current values of \@shortauthor, etc 
  \expandafter\xdef\csname @shorta\roman{plate}\endcsname{\@shortauthor}
  \expandafter\xdef\csname @shortt\roman{plate}\endcsname{\@shorttitle}
  % save the current value of the figure counter for later use in \outputplates
  \setcounter{figplate\roman{plate}}{\value{figure}}
  % increment the figure counter according to the optional argument
  \addtocounter{figure}{#1}
  }
%% page style for plate pages 
\newcommand{\ps@rmaaplate}{%
  % Plate X (see page Y) in the foot
  \renewcommand\@oddfoot{\hfil\csname @shorta\roman{plateloop}\endcsname\ 
    (see page~\csname @plateref\roman{plateloop}\endcsname)\hfil}%
  \renewcommand\@evenfoot{\relax}%
  % nothing in even heads (blank pages)
  \renewcommand\@evenhead{\relax}
  % short title and page num in the odd heads
  \renewcommand\@oddhead{\hfil\csname @shortt\roman{plateloop}\endcsname%
    \hfil\llap{PLATE~\theplateloop}}
  }





%% Things for the User Guide
\if@guide
%% stolen from doc.sty
\def\MakeShortVerb#1{%
  \expandafter\ifx\csname cc\string#1\endcsname\relax
  \@shortvrbinfo{Made }{#1}%
  \add@special{#1}%
  \expandafter
  \xdef\csname cc\string#1\endcsname{\the\catcode`#1}%
  \begingroup
  \catcode`\~\active  \lccode`\~`#1%
  \lowercase{%
    \global\expandafter\let
    \csname ac\string#1\endcsname~%
    \gdef~{\verb~}}%
  \endgroup
  \global\catcode`#1\active
  \else
  \@shortvrbinfo\@empty{#1 already}%
  \fi}
\def\DeleteShortVerb#1{%
  \expandafter\ifx\csname cc\string#1\endcsname\relax
  \else
  \@shortvrbinfo{Deleted }{#1 as}%
  \rem@special{#1}%
  \global\catcode`#1\csname cc\string#1\endcsname
  \global \expandafter\let \csname cc\string#1\endcsname \relax
  \ifnum\catcode`#1=\active
  \begingroup
  \catcode`\~\active   \lccode`\~`#1%
  \lowercase{%
    \global\expandafter\let\expandafter~%
    \csname ac\string#1\endcsname}%
  \endgroup \fi \fi}
\def\@shortvrbinfo#1#2{%
  \PackageInfo{doc}{%
    #1\expandafter\@gobble\string#2 a short reference
    for \string\verb}}
\def\add@special#1{%
  \rem@special{#1}%
  \expandafter\gdef\expandafter\dospecials\expandafter
  {\dospecials \do #1}%
  \expandafter\gdef\expandafter\@sanitize\expandafter
  {\@sanitize \@makeother #1}}
\def\rem@special#1{%
  \def\do##1{%
    \ifnum`#1=`##1 \else \noexpand\do\noexpand##1\fi}%
  \xdef\dospecials{\dospecials}%
  \begingroup
  \def\@makeother##1{%
    \ifnum`#1=`##1 \else \noexpand\@makeother\noexpand##1\fi}%
  \xdef\@sanitize{\@sanitize}%
  \endgroup}
%% define | as shortcut for \verbatim
\MakeShortVerb{\|}
%% comand to specify arguments to LaTeX commands
\newcommand{\rmArg}[1]{%
  \ensuremath{\langle}{\scriptsize \textsf{\MakeUpperCase{#1}}}\ensuremath{\rangle}}
%\newcommand{\rmArg}[1]{\ensuremath{\langle}\textsc{#1}\ensuremath{\rangle}}
\providecommand{\BibTeX}{{\rmfamily B\kern-.05em%
    \textsc{i\kern-.025em b}\kern-.08em%
    T\kern-.1667em\lower.7ex\hbox{E}\kern-.125emX}}
%% Environment for code examples
\newenvironment{Code}{%
  \begin{list}{}{\setlength{\leftmargin}{0.05\columnwidth}} 
  \item[]}{\end{list}}
\fi
%</package>
%    \end{macrocode}

% \section{User commands for astronomical articles}
% \label{sec:user-commands}
% Some random definitions. These go both into the main class file and
% also into |adshtml.sty|.
% 
% Definitions of the journal name in different incarnations.  
%    \begin{macrocode}
%<*package|adshtml-style>
\newcommand{\RMAA}{RMxAA}
\newcommand{\RMAAlong}{Revista Mexicana de Astronom\'{\i}a y Astrof\'{\i}sica}
\newcommand{\RMAASC}{RMxAC}
\newcommand{\RMAASClong}{\RMAA\ (Serie de Conferencias)}
%    \end{macrocode}
% \begin{macro}{\ibidrule}
% This puts a horizontal rule with a period after it. It is used when
% the author list of a reference is identical to the preceding one. 
%    \begin{macrocode}
\newcommand{\ibidrule}{\rule[-0.03cm]{4.7em}{.01cm}.\ }
%    \end{macrocode}
% \end{macro}
% \begin{macro}{\ThisVolume}
% References to other articles in the same volume can be coded as,
% e.g., |\bibitem{} Bloggs, J. \ThisVolume{bloggs}|.
%    \begin{macrocode}
\newcommand{\ThisVolume}[1]{%
  \@year, \if@proceedings\RMAASC\else\RMAA\fi, \@volume,
  \if@book\oldpageref{#1:RMAAFirstPage}\else 000\fi\ (this volume)}
%    \end{macrocode}
% \end{macro}
% Other stuff that may be useful in ref lists.
%    \begin{macrocode}
\newcommand{\etal}{et al.\@}
%    \end{macrocode}
% Various new commands from AAS macros
%    \begin{macrocode}
\def\acknowledgments{\vskip 3ex plus .8ex minus .4ex}
\let\acknowledgements=\acknowledgments                  % second spelling
\let\jnl@style=\rm
\def\ref@jnl#1{{\jnl@style#1}}
\DeclareRobustCommand\aj{\ref@jnl{AJ}}                   
\DeclareRobustCommand\araa{\ref@jnl{ARA\&A}}             
\DeclareRobustCommand\apj{\ref@jnl{ApJ}}                 
\DeclareRobustCommand\apjl{\ref@jnl{ApJ}}                
\DeclareRobustCommand\apjs{\ref@jnl{ApJS}}               
\DeclareRobustCommand\ao{\ref@jnl{Appl.~Opt.}}           
\DeclareRobustCommand\apss{\ref@jnl{Ap\&SS}}             
\DeclareRobustCommand\aap{\ref@jnl{A\&A}}                
\DeclareRobustCommand\aapr{\ref@jnl{A\&A~Rev.}}          
\DeclareRobustCommand\aaps{\ref@jnl{A\&AS}}              
\DeclareRobustCommand\azh{\ref@jnl{AZh}}                 
\DeclareRobustCommand\baas{\ref@jnl{BAAS}}               
\DeclareRobustCommand\jrasc{\ref@jnl{JRASC}}             
\DeclareRobustCommand\memras{\ref@jnl{MmRAS}}            
\DeclareRobustCommand\mnras{\ref@jnl{MNRAS}}             
\DeclareRobustCommand\pra{\ref@jnl{Phys.~Rev.~A}}        
\DeclareRobustCommand\prb{\ref@jnl{Phys.~Rev.~B}}        
\DeclareRobustCommand\prc{\ref@jnl{Phys.~Rev.~C}}        
\DeclareRobustCommand\prd{\ref@jnl{Phys.~Rev.~D}}        
\DeclareRobustCommand\pre{\ref@jnl{Phys.~Rev.~E}}        
\DeclareRobustCommand\prl{\ref@jnl{Phys.~Rev.~Lett.}}    
\DeclareRobustCommand\pasp{\ref@jnl{PASP}}               
\DeclareRobustCommand\pasj{\ref@jnl{PASJ}}               
\DeclareRobustCommand\qjras{\ref@jnl{QJRAS}}             
\DeclareRobustCommand\skytel{\ref@jnl{S\&T}}             
\DeclareRobustCommand\solphys{\ref@jnl{Sol.~Phys.}}      
\DeclareRobustCommand\sovast{\ref@jnl{Soviet~Ast.}}      
\DeclareRobustCommand\ssr{\ref@jnl{Space~Sci.~Rev.}}     
\DeclareRobustCommand\zap{\ref@jnl{ZAp}}                 
\DeclareRobustCommand\nat{\ref@jnl{Nature}}              
\DeclareRobustCommand\iaucirc{\ref@jnl{IAU~Circ.}}
\DeclareRobustCommand\aplett{\ref@jnl{Astrophys.~Lett.}}
\DeclareRobustCommand\apspr{\ref@jnl{Astrophys.~Space~Phys.~Res.}}
\DeclareRobustCommand\bain{\ref@jnl{Bull.~Astron.~Inst.~Netherlands}}
\DeclareRobustCommand\fcp{\ref@jnl{Fund.~Cosmic~Phys.}}
\DeclareRobustCommand\gca{\ref@jnl{Geochim.~Cosmochim.~Acta}}
\DeclareRobustCommand\grl{\ref@jnl{Geophys.~Res.~Lett.}}
\DeclareRobustCommand\jcp{\ref@jnl{J.~Chem.~Phys.}}      
\DeclareRobustCommand\jgr{\ref@jnl{J.~Geophys.~Res.}}    
\DeclareRobustCommand\jqsrt{\ref@jnl{J.~Quant.~Spec.~Radiat.~Transf.}}
\DeclareRobustCommand\memsai{\ref@jnl{Mem.~Soc.~Astron.~Italiana}}
\DeclareRobustCommand\nphysa{\ref@jnl{Nucl.~Phys.~A}}
\DeclareRobustCommand\physrep{\ref@jnl{Phys.~Rep.}}
\DeclareRobustCommand\physscr{\ref@jnl{Phys.~Scr}}
\DeclareRobustCommand\planss{\ref@jnl{Planet.~Space~Sci.}}
\DeclareRobustCommand\procspie{\ref@jnl{Proc.~SPIE}}
\let\astap=\aap
\let\apjlett=\apjl
\let\apjsupp=\apjs
\let\applopt=\ao
\def\phn{\phantom{0}}
\def\phd{\phantom{.}}
\def\phs{\phantom{\ensuremath{-}}}
\def\phm#1{\phantom{#1}}
%    \end{macrocode}
% All these are redone as robust and using more modern commands where
% possible. The robustness is so they will not get expanded if used in
% abstracts and titles when these are passed to the ADS file. 
% \changes{3.16}{2003/04/07}{
%  Added \cs{Color} command for doing photometric colors.}
%    \begin{macrocode}
%%\def\sun{\hbox{$\odot$}}
\DeclareRobustCommand\sun{\ensuremath{\odot}}
%%\def\earth{\hbox{$\oplus$}}
\DeclareRobustCommand\earth{\ensuremath{\oplus}}
%%\def\lesssim{\mathrel{\hbox{\rlap{\hbox{\lower4pt\hbox{$\sim$}}}\hbox{$<$}}}}
%%\def\gtrsim{\mathrel{\hbox{\rlap{\hbox{\lower4pt\hbox{$\sim$}}}\hbox{$>$}}}}
\DeclareRobustCommand\lesssim{\mathrel{\hbox{\smaller\smaller\rlap{\hbox{\lower3pt\hbox{$\sim$}}}\raise1.5pt\hbox{$<$}}}}
\DeclareRobustCommand\gtrsim{\mathrel{\hbox{\smaller\smaller\rlap{\hbox{\lower3pt\hbox{$\sim$}}}\raise1.5pt\hbox{$>$}}}}
%%\def\lta{\mathrel{\spose{\lower 3pt\hbox{$\mathchar"218$}}
%%     \raise 2.0pt\hbox{$\mathchar"13C$}}}
%%\def\gta{\mathrel{\spose{\lower 3pt\hbox{$\mathchar"218$}}
%%     \raise 2.0pt\hbox{$\mathchar"13E$}}}
\DeclareRobustCommand\sq{\hbox{\rlap{$\sqcap$}$\sqcup$}}
%%\def\arcdeg{\hbox{$^\circ$}}
%%\def\arcmin{\hbox{$^\prime$}}
%%\def\arcsec{\hbox{$^{\prime\prime}$}}
\DeclareRobustCommand\arcdeg{\ensuremath{^\circ}}
\DeclareRobustCommand\arcmin{\ensuremath{^\prime}}
\DeclareRobustCommand\arcsec{\ensuremath{^{\prime\prime}}}
%%\def\fd{\hbox{$.\!\!^{\rm d}$}}
%%\def\fh{\hbox{$.\!\!^{\rm h}$}}
%%\def\fm{\hbox{$.\!\!^{\rm m}$}}
%%\def\fs{\hbox{$.\!\!^{\rm s}$}}
%%\def\fdg{\hbox{$.\!\!^\circ$}}
%%\def\farcm{\hbox{$.\mkern-4mu^\prime$}}
%%\def\farcs{\hbox{$.\!\!^{\prime\prime}$}}
%%\def\fp{\hbox{$.\!\!^{\scriptscriptstyle\rm p}$}}
\DeclareRobustCommand\fd{\ensuremath{.\!\!^\mathrm{d}}}
\DeclareRobustCommand\fh{\ensuremath{.\!\!^\mathrm{h}}}
\DeclareRobustCommand\fm{\ensuremath{.\!\!^\mathrm{m}}}
\DeclareRobustCommand\fs{\ensuremath{\mkern0.5mu.\mkern-5mu^\mathrm{s}\mkern-1mu}}
\DeclareRobustCommand\fdg{\ensuremath{.\!\!^\circ}}
\DeclareRobustCommand\farcm{\ensuremath{.\mkern-4mu^\prime}}
\DeclareRobustCommand\farcs{\ensuremath{\mkern1mu.\mkern-6mu^{\prime\prime}\mkern-1.5mu}}
%%\DeclareRobustCommand\farcs{\ensuremath{.\!\!^{\prime\prime}}}
\DeclareRobustCommand\fp{\ensuremath{.\!\!^{\scriptscriptstyle\rm p}}}
%%\def\micron{\hbox{$\mu$m}}
\DeclareRobustCommand\micron{\ensuremath{\mu\mathrm{m}}}
\let\la=\lesssim                        
\let\ga=\gtrsim
%%\def\case#1#2{\hbox{$\frac{#1}{#2}$}}
%%\def\slantfrac#1#2{\hbox{$\,^#1\!/_#2$}}
%%\def\onehalf{\slantfrac{1}{2}}
%%\def\onethird{\slantfrac{1}{3}}
%%\def\twothirds{\slantfrac{2}{3}}
%%\def\onequarter{\slantfrac{1}{4}}
%%\def\threequarters{\slantfrac{3}{4}}
\DeclareRobustCommand\case[2]{\ensuremath{\frac{#1}{#2}}}
\DeclareRobustCommand\slantfrac[2]{\ensuremath{\,^#1\!/_#2}}
\DeclareRobustCommand\onehalf{\slantfrac{1}{2}}
\DeclareRobustCommand\onethird{\slantfrac{1}{3}}
\DeclareRobustCommand\twothirds{\slantfrac{2}{3}}
\DeclareRobustCommand\onequarter{\slantfrac{1}{4}}
\DeclareRobustCommand\threequarters{\slantfrac{3}{4}}
%%\def\ubvr{\hbox{$U\!BV\!R$}}            
%%\def\ub{\hbox{$U\!-\!B$}}               
%%\def\bv{\hbox{$B\!-\!V$}}               
%%\def\vr{\hbox{$V\!-\!R$}}               
%%\def\ur{\hbox{$U\!-\!R$}}
%    \end{macrocode}
% New commands for ``photometric'' colors 
%    \begin{macrocode}
\DeclareRobustCommand\ubvr{\ensuremath{U\!BV\!R}}            
\DeclareRobustCommand\Color[2]{%
  \ensuremath{%
    #1\!\ifx#1V\!\fi-\!#2
  }}
\DeclareRobustCommand\ub{\Color{U}{B}}
\DeclareRobustCommand\bv{\Color{B}{V}}
\DeclareRobustCommand\vr{\Color{V}{R}}
\DeclareRobustCommand\ur{\Color{U}{R}}
%% New version of the \ion command, that actually works! 
%    \end{macrocode}
% \begin{macro}{\ion}
% This version of the |\ion| command allows you to write either
% |\ion{N}{4}| or |\ion{N}{iv}| in order to get N~IV. 
%    \begin{macrocode}
\DeclareRobustCommand\ion[2]{%
  \mbox{#1\kern0.2em%
    \smaller\rmfamily%
    \edef\@tempa{\@car#2\@nil}%
    \ifcat1\@tempa%
    \@Roman{#2}%
    \else%
    \uppercase{#2}%
    \fi}}
\newcounter{IonStage}
\DeclareRobustCommand\plainion[2]{\setcounter{IonStage}{#2}#1 \Roman{IonStage}}
%%\def\ion#1#2{\ensuremath{\mathrm{#1}\;}{\protect\small\rm\@Roman{#2}}}
%    \end{macrocode}
% \end{macro}
%    \begin{macrocode}
\DeclareRobustCommand\nodata{$\cdots$}
%    \end{macrocode}
% These are needed by the |\ion| command. Isn't there an easier way of
% doing this?
%    \begin{macrocode}
\let\@small=\small
\let\@normalsize=\normalsize
\let\@footnotesize=\footnotesize
\let\@scriptsize=\scriptsize
\let\@tiny=\tiny
\let\@large=\large
\let\@Large=\Large
\let\@LARGE=\LARGE
\let\@huge=\huge
\let\@Huge=\Huge
\newcommand{\smaller}{\small}
\renewcommand{\normalsize}{\renewcommand{\smaller}{\small}\@normalsize}
\renewcommand{\small}{\renewcommand{\smaller}{\footnotesize}\@small}
\renewcommand{\footnotesize}{\renewcommand{\smaller}{\scriptsize}\@footnotesize}
\renewcommand{\scriptsize}{\renewcommand{\smaller}{\tiny}\@scriptsize}
\renewcommand{\tiny}{\renewcommand{\smaller}{\tiny}\@tiny}
\renewcommand{\large}{\renewcommand{\smaller}{\normalsize}\@large}  
\renewcommand{\Large}{\renewcommand{\smaller}{\large}\@Large}  
\renewcommand{\LARGE}{\renewcommand{\smaller}{\Large}\@LARGE}  
\renewcommand{\huge}{\renewcommand{\smaller}{\LARGE}\@huge}  
\renewcommand{\Huge}{\renewcommand{\smaller}{\huge}\@Huge}  
%</package|adshtml-style>
%    \end{macrocode}


% \section{Other files}

% \subsection{The driver and make files}
% \label{sec:driver_make}
% These are put here for completeness, although they still need to be
% distributed separately since otherwise the user won't be able to
% unpack this file into its constituent parts.
% 
% The |Makefile| has two principal rules: |docs|, and |macros|. The
% |docs| rule produces a PDF and gzipped PS file of the documentation.
% I still haven't perfected the combination of rules required to get
% it to run \LaTeX{} enough times in all circumstances. The |macros|
% rule runs \LaTeX{} on |rmaa.ins| in order to unpack |rmaa.cls| and
% all the other little files (including the |Makefile| and the
% |rmaa.ins| files themselves\dots 
% 
% First set up lists of various types of file. 
%    \begin{macrocode}
%<*makefile>
DOCFILES=rmaa.pdf
MACRO_LTX_FILES=rmaa.cls rmaacite.sty adshtml.sty
SHELLSCRIPTS=makeads.sh
MACRO_OTHER_FILES=adshtml.perl sample.latex2html-init
MAKEFILE=Makefile
MACROFILES=$(MACRO_LTX_FILES) $(SHELLSCRIPTS) $(MACRO_OTHER_FILES)         
SKELETONFILES=ADS.tex 
%    \end{macrocode}
% Install the files read only into the package directory, specified via a path relative to the source directory. \changes{4.0}{2019/11/14}{Reorganize folder layout so that hardly anything is in the top level.  That way, things are clearer on github.}
%    \begin{macrocode}
INSTALL=install 
INSTALL_DIR=../packages
%    \end{macrocode}
% The default rule is to make everything. Installation must be
% explicitly asked for.
%    \begin{macrocode}
all: docs macros

install: install-docs install-macros

%    \end{macrocode}
% Most files are installed read-only to prevent editing. Skeleton
% files must be writeable. Must be careful not to overwrite the
% Makefile in the parent directory. 
%    \begin{macrocode}
install-macros: macros
        $(INSTALL) -m 444 $(MACRO_LTX_FILES) $(MACRO_OTHER_FILES) $(INSTALL_DIR)
        $(INSTALL) -m 555 $(SHELLSCRIPTS) $(INSTALL_DIR)
        $(INSTALL) -m 644 $(SKELETONFILES) $(INSTALL_DIR)

install-docs: docs
        $(INSTALL) rmaa.pdf $(INSTALL_DIR)/rm-fulldocs.pdf

%    \end{macrocode}
% Make the PDF and gzipped PS files. 
%    \begin{macrocode}
docs: $(DOCFILES)

%    \end{macrocode}
% Make everything else if either the |.dtx| or |.ins| files have
% changed. 
%    \begin{macrocode}
macros: rmaa.dtx rmaa.ins
%    \end{macrocode}
% Ensure that all the macro files can be written to then recreate
% them. The |touch| is there in case any of the files don't exist
% yet. 
%    \begin{macrocode}
        touch $(MACROFILES) $(MAKEFILE)
        chmod u+w $(MACROFILES) $(MAKEFILE) rmaa.ins
        latex "\input{rmaa.ins}"
%    \end{macrocode}
% \TeX{} automatically adds the |.tex| extension to extensionless
% files. We fix this and take the opportunity to convert
% back to TABs all that it has lovingly converted to spaces. 
%    \begin{macrocode}
        cp -f $(MAKEFILE) $(MAKEFILE).BCK && unexpand $(MAKEFILE).tex >| $(MAKEFILE) && rm $(MAKEFILE).tex
%    \end{macrocode}
% Make the shell scripts executable and everything unwriteable (so
% they don't get hand-edited by mistake). However, undo this for the
% files that the user is \emph{supposed} to modify. 
%    \begin{macrocode}
        chmod a+x $(SHELLSCRIPTS)
        chmod a-w $(MACROFILES) $(MAKEFILE)
        chmod u+w $(SKELETONFILES)

%    \end{macrocode}
% We now use PDF\LaTeX{} for this.
%    \begin{macrocode}
rmaa.idx: rmaa.dtx
        @echo "REGENERATING .idx FILE"
        pdflatex "\nonstopmode\input{rmaa.dtx}"

rmaa.aux: rmaa.dtx rmaa.ind rmaa.gls
        @echo "REGENERATING .aux FILE"
        pdflatex "\nonstopmode\input{rmaa.dtx}"

rmaa.ind: rmaa.idx
        @echo "REGENERATING INDEX"
        makeindex -s gind.ist rmaa

rmaa.gls: rmaa.glo
        @echo "REGENERATING CHANGELOG"
        makeindex -s gglo.ist -o rmaa.gls rmaa.glo

rmaa.pdf: rmaa.aux rmaa.ind rmaa.gls
        @echo "REGENERATING .pdf FILE"
        pdflatex "\nonstopmode\input{rmaa.dtx}"

clean: 
        rm -f rmaa.aux rmaa.dvi rmaa.ind rmaa.gls
%</makefile>
%    \end{macrocode}

% The |rmaa.ins| file uses the \textsf{docstrip} package to split out
% all the files and strip off this documentation.
%    \begin{macrocode}
%<*ins>
%%
%% Installation batch file for the RMAA LaTeX package
%%
\def\batchfile{rmaa.ins}
\input docstrip.tex
\preamble

PLEASE DO NOT EDIT THIS FILE! It is generated automatically from
the file rmaa.dtx using the accompanying Makefile and rmaa.ins 
Any hand-edited changes will be lost

A LaTeX document class for the RMxAA

Author: William Henney (w.henney@irya.unam.mx)

Latest version available from: http://www.irya.unam.mx/rmaa

Copyright (C) 1999-2019 Instituto de Astronomia, UNAM, Mexico

Permission is hereby granted, free of charge, to any person obtaining
a copy of this software and associated documentation files (the
"Software"), to deal in the Software without restriction, including
without limitation the rights to use, copy, modify, merge, publish,
distribute, sublicense, and/or sell copies of the Software, and to
permit persons to whom the Software is furnished to do so, subject to
the following conditions:

The above copyright notice and this permission notice shall be
included in all copies or substantial portions of the Software.

THE SOFTWARE IS PROVIDED "AS IS", WITHOUT WARRANTY OF ANY KIND,
EXPRESS OR IMPLIED, INCLUDING BUT NOT LIMITED TO THE WARRANTIES OF
MERCHANTABILITY, FITNESS FOR A PARTICULAR PURPOSE AND
NONINFRINGEMENT. IN NO EVENT SHALL THE AUTHORS OR COPYRIGHT HOLDERS BE
LIABLE FOR ANY CLAIM, DAMAGES OR OTHER LIABILITY, WHETHER IN AN ACTION
OF CONTRACT, TORT OR OTHERWISE, ARISING FROM, OUT OF OR IN CONNECTION
WITH THE SOFTWARE OR THE USE OR OTHER DEALINGS IN THE SOFTWARE.
\endpreamble
%    \end{macrocode}
% Clobber files without asking
%    \begin{macrocode}
\askforoverwritefalse
%    \end{macrocode}
% Make the \LaTeX{} macro files (class and styles)
%    \begin{macrocode}
\generate{% 
  \file{rmaa.cls}{\from{rmaa.dtx}{package}}
  \file{rmaa.ins}{\from{rmaa.dtx}{ins}}
  \file{rmaacite.sty}{\from{rmaa.dtx}{rmaacite}}
  \file{adshtml.sty}{\from{rmaa.dtx}{adshtml-style}}
  \file{ADS.tex}{\from{rmaa.dtx}{ads-container}}
  }
%    \end{macrocode}
% Make the other non-\LaTeX{} files. These are done separately since
% they can't use `\%\%' as a comment character for the
% preamble/postamble (and Perl files can't have the postamble I don't
% think). 
%    \begin{macrocode}
\generate{\nopreamble\nopostamble% 
  \file{Makefile}{\from{rmaa.dtx}{makefile}}
  \file{adshtml.perl}{\from{rmaa.dtx}{adshtml-perl}}
  \file{makeads.sh}{\from{rmaa.dtx}{ads-shell-script}}
  \file{sample.latex2html-init}{\from{rmaa.dtx}{latex2html-init}}
  }

%% Local Variables: 
%% mode: LaTeX
%% mode: Font-lock
%% End:
%</ins>
%    \end{macrocode}

% \subsection{Extended citation syntax}
% Modified from a file called mnrascite.sty
% I no longer remember where I stole the original version of this from
% 
% Alterations to the cite commands to make them do more sensible things.
% N.B. the bibliography style should produce entries which have labels where
% the date is delimited by angle brackets '<>' for these macros to work
% properly. They will not give very sensible error messages if you do not do 
% this properly.
% 
% A new command |\scite{key}| is defined which produces a citation which is 
% suitable for insertion in a sentence. i.e., it does not put brackets around 
% the whole citation, only the date. Also |\pcite{key}| produce output with no
% parentheses whatsoever. 
% 
% Should redefine |\@bibitem| really to have any hope of compatibility
% the command |\lbibcite| is defined to deal with the citations.
%    \begin{macrocode}
%<*rmaacite>
\def\bibitem[#1<#2>]#3{\par\hangindent=1.5em\hangafter=1
  \if@filesw
  {\def\protect##1{\string ##1\space}
    \immediate\write\@auxout{\string\lbibcite{#3}{#1}{#2}}}
  \fi}
\@redefinebibitem % what it says
%    \end{macrocode}
% Note that the |\lbibcite| defines two types of citation one contains the names,
% the other the year
%    \begin{macrocode}
\def\lbibcite#1#2#3{\global\@namedef{b@#1}{#2}
                    \global\@namedef{by@#1}{#3}}

\def\@bcite#1#2{({#1\if@tempswa , #2\fi})}
\def\@pcite#1#2{#1\if@tempswa , #2\fi}
%    \end{macrocode}
% This declaration contains what is put between the name and the date
% \changes{3.11}{2003/01/31}{Moved the \cs{articlename} to after the key for consistency
% with the new natbib stuff.}
%    \begin{macrocode}
\def\@brcite{\ (}
\def\@ebrcite{)}
\def\@fmttop#1{\@fmtcite[\csname #1\endcsname]}
\def\@citex[#1]#2{%
%%  \if@filesw\immediate\write\@auxout{\string\citation{#2}}\fi
%% WJH added 16 Jan 2002
  \@for\@citeb:=#2\do{%
    \edef\@citeb{\expandafter\@firstofone\@citeb:\articlename}%
    \if@filesw\immediate\write\@auxout{\string\citation{\@citeb}}\fi}%
%% to here
  \def\@citea{}%
  \@cite{\@for\@citeb:=#2\do
    {%
      \edef\@citeb{\expandafter\@firstofone\@citeb:\articlename}%
      \@citea\def\@citea{;\ }\@ifundefined
      {b@\@citeb}{{\bf ?}\@warning {Citation `\@citeb' on page \thepage \space undefined}}%
      {\csname b@\@citeb\endcsname\@citebro\csname by@\@citeb\endcsname\@citebrc}}}%
  {#1}}

\def\cite{\@ifnextchar [{\let\@citebro=\space
                          \let\@citebrc=\relax
                          \let\@cite=\@bcite\@tempswatrue\@citex}
                        {\let\@citebro=\space
                          \let\@citebrc=\relax
                          \let\@cite=\@bcite\@tempswafalse\@citex[]}}
\def\pcite{\@ifnextchar [{\let\@citebro=\space
                          \let\@citebrc=\relax
                          \let\@cite=\@pcite\@tempswatrue\@citex}
                        {\let\@citebro=\space
                          \let\@citebrc=\relax
                          \let\@cite=\@pcite\@tempswafalse\@citex[]}}
\def\scite{\@ifnextchar [{\let\@citebro=\@brcite
                          \let\@citebrc=\@ebrcite
                          \let\@cite=\@pcite\@tempswatrue\@citex}
                        {\let\@citebro=\@brcite
                          \let\@citebrc=\@ebrcite
                          \let\@cite=\@pcite\@tempswafalse\@citex[]}}

\renewenvironment{thebibliography}
  {\section*{REFERENCES}\vspace*{\afterbiblioskip}\small
    \if@book\RereadBiblio\fi
    \sloppy\clubpenalty4000\widowpenalty4000
    \sfcode`\.=1000\relax
    \setlength\parindent{\z@}\setlength\parskip{\z@}}
  {\relax\normalsize\@outputaddressesifnecessary} 
%</rmaacite>
%    \end{macrocode}

% \subsection{Auxilliary files for generating ADS data}

% The ADS data are written by \LaTeX{} to a |.ads| file. This file
% then has to be included in the following file (|ADS.tex|),
% which is run through \LaTeX{}2HTML (see below). 

%    \begin{macrocode}
%<*ads-container>
%%
%% Instructions:
%% Edit the \input{} line below to reflect the name of your file. 
\documentclass{article}
\usepackage[latin1]{inputenc}
\usepackage{adshtml}
\newcommand{\ADSfield}[2]
{\%#1  #2\\}
%%\newcommand{\ADSvfield}[2]
%%{\%#1 \begin{verbatim}#2\end{verbatim}}
\newenvironment{ADSarticle}{}{\par\bigskip}
\renewcommand{\bibitem}[2][]{\par}
\setlength{\parindent}{0pt}
\renewcommand{\ThisVolume}[1]{%  
  2001, \RMAASC, 12, \oldpageref{#1:RMAAFirstPage}}  
 

\begin{document}
\makeatletter
\input{rmbooksamp.ads}
\end{document}
%</ads-container>
%    \end{macrocode}

% This is the |adshtml.sty| style file used by |ADS.tex|. It just
% defines a few 
% macros that might be used in abstracts, etc (really it should define
% them all). This now gets all its code from
% Sec~\ref{sec:user-commands}. 
%    \begin{macrocode}
%<*adshtml-style>
%% Nothing else needs to go here now. 
%</adshtml-style>
%    \end{macrocode}
% The init file for latex2html.
%    \begin{macrocode}
%<*latex2html-init>
$PDFTOCAIRO = "pdftocairo";
$PDFLATEX = "pdflatex";
$LATEX = "latex";
$MAX_SPLIT_DEPTH = 0;
$HTML_VERSION = "4.0,math";
$NO_SIMPLE_MATH = 1;
$ADDRESS = "<p><address><a href=\"mailto:rmaa\@astro.unam.mx\">rmaa\@astro.unam.mx</address></p>"; 
$INFO = "<p>This document was created using latex2html and the rmaa latex macros</p>"; 
$NO_NAVIGATION = 1;
1; 
%</latex2html-init>
%    \end{macrocode}
% This is the |adshtml.perl| file that is used by \LaTeX{}2HTML. Some
% bizarre \TeX{}y command sometimes show up so we tell it to ignore
% these. All the astronomy macros that come from AAS are processed in
% \TeX{}. That way they get made to images but with an |ALT| tag that
% has the original code, which is what will be picked up by |lynx|,
% with the result that they get passed intact into the final |.txt|
% file. 
%    \begin{macrocode}
%<*adshtml-perl>
package main;

$dum = &do_cmd_makeatletter();

sub do_cmd_RMAA {
    local($_) = @_;
    "RMxAA" . $_;
}

sub do_cmd_RMAASC {
    local($_) = @_;
    "RMxAASC" . $_;
}

sub do_cmd_etal {
    local($_) = @_;
    "et al\." . $_;
}

sub do_cmd_ibidrule {
    local($_) = @_;
    "<img src=\"foo\" alt=\"\ibidrule\">" . $_;
}


sub do_cmd_bibitem {
    local($_) = @_;
    local($key);
    local($opt, $dummy)=&get_next_optional_argument; 
    $key = &missing_braces unless
        ((s/$next_pair_pr_rx/$key = $2; ''/eo)
        ||(s/$next_pair_rx/$key = $2; ''/eo));
    "<p>" . $_; 
}

&ignore_commands (<<_RAW_ARG_CMDS_);
lowercase
adjustfinalcols
textsuperscript # {}
newline # []
thinmuskip
mskip
smallskip
medskip
bigskip
hspace # {}
hspacestar # {}
vspace # {}
vspacestar # {}
mathrm
_at_mathrm
spacefactor
discretionary
unhbox 
voidb_at_x 
penalty 
_at_M
_at_m
null
IeC
_RAW_ARG_CMDS_

&process_commands_in_tex (<<_RAW_ARG_CMDS_);
sun
earth
lesssim
gtrsim
sq
argdeg
arcmin
arcsec
fd
fh
fm
fs
fdg
farcm
farcs
fp
micron
la
ga
case # {} # {}
slantfrac # {} # {}
onehalf
onethird
twothirds
onequarter
threequarters
ubvr
ub
bv
vr
ur
aj
araa
apj
apjl
apjs
ao
apss
aap
aapr
aaps
azh
baas
jrasc
memras
mnras
pra
prb
prc
prd
pre
prl
pasp
pasj
qjras
skytel
solphys
sovast
ssr
zap
nat
iaucirc
aplett
apspr
bain
fcp
gca
grl
jcp
jgr
jqsrt
memsai
nphysa
physrep
physscr
planss
procspie
_RAW_ARG_CMDS_

1;                              # This must be the last line
%</adshtml-perl> 
%    \end{macrocode}

% This is a shell script |makeads.sh| that generates the |_ads.txt| file from
% |_ads.tex| by using \LaTeX{}2HTML and |lynx|.  \changes{4.0}{2019/11/13}{Make sure that lynx writes utf-8 output instead of latin-1.}
% \changes{2.2beta13}{2002/03/12}{Complete rewrite of |makeads.sh|. Now much 
% improved I think}
% \changes{4.0}{2019/11/13}{Fix the at-sign substitution so that it only eliminates \cs{@} and does not mess with email addresses.}

%    \begin{macrocode}
%<*ads-shell-script>
#! /bin/sh
if [ -z $1 ]; then
    echo "Usage: $0 basename"
    exit
fi
BASENAME=$1
ADSBASENAME=${BASENAME}_ads
AUXFILE=$BASENAME.aux
TMPFILE=tmp.$$


echo "Generating $ADSBASENAME.tex"

# For some reason, must double-up \\ to \\\\ but not \ to \\ - odd!
cat > $ADSBASENAME.tex <<EOF
%% PLEASE DO NOT EDIT THIS FILE!  It is generated automatically from
%% the script makeads.sh.  Any hand-edited changes will be lost
%% 
\documentclass{article}
\usepackage[latin1]{inputenc}
\usepackage{html}
\usepackage{adshtml}
\newcommand{\ADSfield}[2]{\%#1  #2\\\\}
\newenvironment{ADSarticle}{}{\par\bigskip}
\setlength{\parindent}{0pt}
\newcounter{istage}
\renewcommand{\ion}[2]{\setcounter{istage}{#2}#1 \Roman{istage}}
\begin{document}
\makeatletter
EOF
echo >> $ADSBASENAME.tex "\input{$BASENAME.ads2}" 
echo >> $ADSBASENAME.tex "\end{document}"



echo "Generating $ADSBASENAME.sed" 

# Find type of journal
grep '\\documentclass\[.*proceedings.*\]' $BASENAME.tex >/dev/null \
     && JOURNAL='RMxAC' \
     || JOURNAL='RMxAA'

# Find the year and volume number
YEAR=`grep '\\SetYear' $BASENAME.tex | sed 's/.*\\SetYear{\(.*\)}.*/\1/'`
VOLUME=`grep '\\SetVolume' $BASENAME.tex | sed 's/.*\\SetVolume{\(.*\)}.*/\1/'`

# Generate an auxilliary sed script 
cat > $TMPFILE <<EOF
s|\\\\newlabel{\(.*\):RMAAFirstPage}{{.*}{\(.*\)}}|s/\\\\\\\\ThisVolume{\1}/$YEAR, RMxAC, $VOLUME, \2/g|g
EOF

# Now go through the .aux file, looking for \newlabels and use the sed 
# script just generated in order to make another sed script, one that 
# will replace \ThisVolume{LABEL} with the correct reference
grep '^\\newlabel{.*:RMAAFirstPage}' $AUXFILE \
     | sed -f $TMPFILE > $ADSBASENAME.sed && rm $TMPFILE

# Apend a further sed command to get rid of \@, since redefining 
# it in _ads.tex didn't work
echo >> $ADSBASENAME.sed 's/\\\\@//g'

echo "Generating $BASENAME.ads2 from $ADSBASENAME.sed and $BASENAME.ads" 

# Apply this to the .ads file, generating the .ads2 file, which is included
# by the _ads.tex file.
sed -f $ADSBASENAME.sed $BASENAME.ads > $BASENAME.ads2

echo "Generating $ADSBASENAME.txt from $ADSBASENAME.tex and $BASENAME.ads2"

# Now run latex2html on the _ads.tex file
latex2html -no_navigation -info 0 -address 0 -html_version 4.0,math $ADSBASENAME

# Then, convert to utf8 with lynx and tidy it up a bit
lynx --justify=off --width=100 --dump --nolist --display_charset=utf-8 \
     $ADSBASENAME/$ADSBASENAME.html | cut -c4- \
     | sed 's/ibidrule/\\ibidrule/g' > $ADSBASENAME.txt
%</ads-shell-script>
%    \end{macrocode}

% \section{Epilogue}
% Err, that's it \dots

\endinput


% \endinput
% Local Variables:
% mode: doctex
% TeX-master: t
% End:
"

clean: 
        rm -f rmaa.aux rmaa.dvi rmaa.ind rmaa.gls
%</makefile>
%    \end{macrocode}

% The |rmaa.ins| file uses the \textsf{docstrip} package to split out
% all the files and strip off this documentation.
%    \begin{macrocode}
%<*ins>
%%
%% Installation batch file for the RMAA LaTeX package
%%
\def\batchfile{rmaa.ins}
\input docstrip.tex
\preamble

PLEASE DO NOT EDIT THIS FILE! It is generated automatically from
the file rmaa.dtx using the accompanying Makefile and rmaa.ins 
Any hand-edited changes will be lost

A LaTeX document class for the RMxAA

Author: William Henney (w.henney@irya.unam.mx)

Latest version available from: http://www.irya.unam.mx/rmaa

Copyright (C) 1999-2019 Instituto de Astronomia, UNAM, Mexico

Permission is hereby granted, free of charge, to any person obtaining
a copy of this software and associated documentation files (the
"Software"), to deal in the Software without restriction, including
without limitation the rights to use, copy, modify, merge, publish,
distribute, sublicense, and/or sell copies of the Software, and to
permit persons to whom the Software is furnished to do so, subject to
the following conditions:

The above copyright notice and this permission notice shall be
included in all copies or substantial portions of the Software.

THE SOFTWARE IS PROVIDED "AS IS", WITHOUT WARRANTY OF ANY KIND,
EXPRESS OR IMPLIED, INCLUDING BUT NOT LIMITED TO THE WARRANTIES OF
MERCHANTABILITY, FITNESS FOR A PARTICULAR PURPOSE AND
NONINFRINGEMENT. IN NO EVENT SHALL THE AUTHORS OR COPYRIGHT HOLDERS BE
LIABLE FOR ANY CLAIM, DAMAGES OR OTHER LIABILITY, WHETHER IN AN ACTION
OF CONTRACT, TORT OR OTHERWISE, ARISING FROM, OUT OF OR IN CONNECTION
WITH THE SOFTWARE OR THE USE OR OTHER DEALINGS IN THE SOFTWARE.
\endpreamble
%    \end{macrocode}
% Clobber files without asking
%    \begin{macrocode}
\askforoverwritefalse
%    \end{macrocode}
% Make the \LaTeX{} macro files (class and styles)
%    \begin{macrocode}
\generate{% 
  \file{rmaa.cls}{\from{rmaa.dtx}{package}}
  \file{rmaa.ins}{\from{rmaa.dtx}{ins}}
  \file{rmaacite.sty}{\from{rmaa.dtx}{rmaacite}}
  \file{adshtml.sty}{\from{rmaa.dtx}{adshtml-style}}
  \file{ADS.tex}{\from{rmaa.dtx}{ads-container}}
  }
%    \end{macrocode}
% Make the other non-\LaTeX{} files. These are done separately since
% they can't use `\%\%' as a comment character for the
% preamble/postamble (and Perl files can't have the postamble I don't
% think). 
%    \begin{macrocode}
\generate{\nopreamble\nopostamble% 
  \file{Makefile}{\from{rmaa.dtx}{makefile}}
  \file{adshtml.perl}{\from{rmaa.dtx}{adshtml-perl}}
  \file{makeads.sh}{\from{rmaa.dtx}{ads-shell-script}}
  \file{sample.latex2html-init}{\from{rmaa.dtx}{latex2html-init}}
  }

%% Local Variables: 
%% mode: LaTeX
%% mode: Font-lock
%% End:
%</ins>
%    \end{macrocode}

% \subsection{Extended citation syntax}
% Modified from a file called mnrascite.sty
% I no longer remember where I stole the original version of this from
% 
% Alterations to the cite commands to make them do more sensible things.
% N.B. the bibliography style should produce entries which have labels where
% the date is delimited by angle brackets '<>' for these macros to work
% properly. They will not give very sensible error messages if you do not do 
% this properly.
% 
% A new command |\scite{key}| is defined which produces a citation which is 
% suitable for insertion in a sentence. i.e., it does not put brackets around 
% the whole citation, only the date. Also |\pcite{key}| produce output with no
% parentheses whatsoever. 
% 
% Should redefine |\@bibitem| really to have any hope of compatibility
% the command |\lbibcite| is defined to deal with the citations.
%    \begin{macrocode}
%<*rmaacite>
\def\bibitem[#1<#2>]#3{\par\hangindent=1.5em\hangafter=1
  \if@filesw
  {\def\protect##1{\string ##1\space}
    \immediate\write\@auxout{\string\lbibcite{#3}{#1}{#2}}}
  \fi}
\@redefinebibitem % what it says
%    \end{macrocode}
% Note that the |\lbibcite| defines two types of citation one contains the names,
% the other the year
%    \begin{macrocode}
\def\lbibcite#1#2#3{\global\@namedef{b@#1}{#2}
                    \global\@namedef{by@#1}{#3}}

\def\@bcite#1#2{({#1\if@tempswa , #2\fi})}
\def\@pcite#1#2{#1\if@tempswa , #2\fi}
%    \end{macrocode}
% This declaration contains what is put between the name and the date
% \changes{3.11}{2003/01/31}{Moved the \cs{articlename} to after the key for consistency
% with the new natbib stuff.}
%    \begin{macrocode}
\def\@brcite{\ (}
\def\@ebrcite{)}
\def\@fmttop#1{\@fmtcite[\csname #1\endcsname]}
\def\@citex[#1]#2{%
%%  \if@filesw\immediate\write\@auxout{\string\citation{#2}}\fi
%% WJH added 16 Jan 2002
  \@for\@citeb:=#2\do{%
    \edef\@citeb{\expandafter\@firstofone\@citeb:\articlename}%
    \if@filesw\immediate\write\@auxout{\string\citation{\@citeb}}\fi}%
%% to here
  \def\@citea{}%
  \@cite{\@for\@citeb:=#2\do
    {%
      \edef\@citeb{\expandafter\@firstofone\@citeb:\articlename}%
      \@citea\def\@citea{;\ }\@ifundefined
      {b@\@citeb}{{\bf ?}\@warning {Citation `\@citeb' on page \thepage \space undefined}}%
      {\csname b@\@citeb\endcsname\@citebro\csname by@\@citeb\endcsname\@citebrc}}}%
  {#1}}

\def\cite{\@ifnextchar [{\let\@citebro=\space
                          \let\@citebrc=\relax
                          \let\@cite=\@bcite\@tempswatrue\@citex}
                        {\let\@citebro=\space
                          \let\@citebrc=\relax
                          \let\@cite=\@bcite\@tempswafalse\@citex[]}}
\def\pcite{\@ifnextchar [{\let\@citebro=\space
                          \let\@citebrc=\relax
                          \let\@cite=\@pcite\@tempswatrue\@citex}
                        {\let\@citebro=\space
                          \let\@citebrc=\relax
                          \let\@cite=\@pcite\@tempswafalse\@citex[]}}
\def\scite{\@ifnextchar [{\let\@citebro=\@brcite
                          \let\@citebrc=\@ebrcite
                          \let\@cite=\@pcite\@tempswatrue\@citex}
                        {\let\@citebro=\@brcite
                          \let\@citebrc=\@ebrcite
                          \let\@cite=\@pcite\@tempswafalse\@citex[]}}

\renewenvironment{thebibliography}
  {\section*{REFERENCES}\vspace*{\afterbiblioskip}\small
    \if@book\RereadBiblio\fi
    \sloppy\clubpenalty4000\widowpenalty4000
    \sfcode`\.=1000\relax
    \setlength\parindent{\z@}\setlength\parskip{\z@}}
  {\relax\normalsize\@outputaddressesifnecessary} 
%</rmaacite>
%    \end{macrocode}

% \subsection{Auxilliary files for generating ADS data}

% The ADS data are written by \LaTeX{} to a |.ads| file. This file
% then has to be included in the following file (|ADS.tex|),
% which is run through \LaTeX{}2HTML (see below). 

%    \begin{macrocode}
%<*ads-container>
%%
%% Instructions:
%% Edit the \input{} line below to reflect the name of your file. 
\documentclass{article}
\usepackage[latin1]{inputenc}
\usepackage{adshtml}
\newcommand{\ADSfield}[2]
{\%#1  #2\\}
%%\newcommand{\ADSvfield}[2]
%%{\%#1 \begin{verbatim}#2\end{verbatim}}
\newenvironment{ADSarticle}{}{\par\bigskip}
\renewcommand{\bibitem}[2][]{\par}
\setlength{\parindent}{0pt}
\renewcommand{\ThisVolume}[1]{%  
  2001, \RMAASC, 12, \oldpageref{#1:RMAAFirstPage}}  
 

\begin{document}
\makeatletter
%% Demonstration file for formatting an entire volume
\documentclass[book,journal,twocolumn]{rmaa}
%%%%%%%%%%%%%%%%%%%%%%%%%%%%%%%%%%%%%%%%%%%%%%%%%%%%%%%%%%%%%%%%
%%%%%%%%%%%%%%%%            PREAMBLE            %%%%%%%%%%%%%%%%
%%%%%%%%%%%%%%%%%%%%%%%%%%%%%%%%%%%%%%%%%%%%%%%%%%%%%%%%%%%%%%%%

\usepackage{paralist}
\usepackage{psfrag,color}
% Needed for landscape tables
\usepackage{pdflscape}

% Used in rm-minimal-test
\usepackage{kantlipsum}

% publication details of volume
\SetVolume{9}
\SetYear{2000}
\SetMonthNumeric{06}
\SetMes{junio}
% details of meeting
\SetConfTitle{Astrophysical Plasmas: Codes, Models, \& Observations}
\SetReunion{M\'exico, D. F.}{octubre 25--29, 1999}
\SetEditors{S. Jane Arthur, Nancy S. Brickhouse, \& Jos\'e Franco}

% this writes the .idx file that is used to generate the index 
\makeindex

%% This command makes the session headings come out bold in the
%% Table-of-Contents. Comment it out if you don't like it.
\BoldSessionTitles

%% This command makes a title page for each session
%% Comment it out if you don't like it.
\SessionTitlePages

%% Use this command if you want the abstract-only papers to be
%% indexed by ADS
\abstractsADS

%% At the special request of Alberto Carramin~ana, the following is a
%% way to suppress the author adresses at the end of each article
%\suppressfulladdresses

%% Use this command to add Received/accepted dates for each paper
% \printdates


\begin{document}

%%%%%%%%%%%%%%%%%%%%%%%%%%%%%%%%%%%%%%%%%%%%%%%%%%%%%%%%%%%%%%%%b
%%%%%%%%%%%%%%%%          FRONT MATTER          %%%%%%%%%%%%%%%%
%%%%%%%%%%%%%%%%%%%%%%%%%%%%%%%%%%%%%%%%%%%%%%%%%%%%%%%%%%%%%%%%
\frontmatter % start roman numbering

% format the TOC
\renewcommand{\RMAAbannerimage}{rmsc_blank}
\tableofcontents

%% preface, group photo, list of participants go here
%% use the `frontsection' environment (see editors' guide)
\begin{frontsection}{Preface}{%
    \listofauthors{S.~J.~Arthur, N.~S.~Brickhouse, J.~Franco}
    \indexauthor{Arthur, S. J.}
    \indexauthor{Brickhouse, N.}
    \indexauthor{Franco, J.}
    }
  These are the proceedings of a very nice conference that Jane and
  Pepe organised. This is the preface, which explains what a wonderful
  time was had by all and how grateful they are for all the people who
  gave them money or helped out.
  
  This is another paragraph of the preface, in case they have anything
  else to say.  For example, the Scientific Organizing Committee
  consisted of
  \begin{quote}
    J. Franco (Mexico, Chair), D. Cox (USA) and various other notables
  \end{quote}
  while the Local Organizing Committee comprised
  \begin{quote}
    S. J. Arthur (Chair), S. Kurtz, etc.
  \end{quote}

  \EditorSignature

\end{frontsection}

\begin{frontsection}[Photo]{Group Photograph}{\notitle}
\null\vfill
\centerline{Here goes the photograph.}
\vfill
\end{frontsection}

\begin{frontsection}[Participants]{List of Participants}{}
%% Here you put a formatted list of participants
\vfil
\dotfill First page of list of participants \dotfill
\vfil\clearpage

\null\vfil
\dotfill Second page of list of participants \dotfill
\vfil

\end{frontsection}

\begin{frontsection}[XXX]{Another front section}{}

\end{frontsection}

%%%%%%%%%%%%%%%%%%%%%%%%%%%%%%%%%%%%%%%%%%%%%%%%%%%%%%%%%%%%%%%%
%%%%%%%%%%%%%%%%          MAIN MATTER           %%%%%%%%%%%%%%%%
%%%%%%%%%%%%%%%%%%%%%%%%%%%%%%%%%%%%%%%%%%%%%%%%%%%%%%%%%%%%%%%%
\mainmatter % start arabic numbering

\begin{articles}
  %% Main articles go here
  %% Use \article{ARTICLENAME} for each individual article
  %% If you want any text to appear in brackets after the article
  %% title in the TOC and index, then put it as an optional first
  %% argument, e.g.,  \article[Invited Talk]{ARTICLENAME} 

  \session{The Contributed Papers}
  \article{rm-journal-example} % An illustrative article
  \article{rm-onepage} % another article, which has a plate
  \article{rm-minimal-test}
  \session{Conference Summary}
  \article[Invited Review]{rmsummsamp} % An illustrative conference summary

\end{articles}

%%%%%%%%%%%%%%%%%%%%%%%%%%%%%%%%%%%%%%%%%%%%%%%%%%%%%%%%%%%%%%%%
%%%%%%%%%%%%%%%%          BACK MATTER           %%%%%%%%%%%%%%%%
%%%%%%%%%%%%%%%%%%%%%%%%%%%%%%%%%%%%%%%%%%%%%%%%%%%%%%%%%%%%%%%%
\backmatter

%% Abstracts of presentations for which the
%% authors have not submitted a paper
\begin{abstracts}

  \title{An abstract of a poster}
  \author{W. J. Henney\altaffilmark{1}}
  \listofauthors{W.~J.~Henney}
  \indexauthor{Henney, W. J}
  \altaffiltext{1}{Where I am from done with altaffilmark/text}
  \abstract{This is a test abstract.   This is a test abstract.   This
    is a test abstract.   This is a test abstract.}
  

  \title{An abstract of another  poster}
  \author{W.~J.~Henney\altaffil{This is done with altaffil} and
    A.~N.~Other\altaffil{So is this} }
  \listofauthors{W.~J.~Henney \& A.~N.~Other}
  \indexauthor{Henney, W. J}
  \indexauthor{Other, A. N.}
  \abstract{This is a test abstract.   This is a test abstract.   This
    is a test abstract.   This is a test abstract.}

  \title{An abstract of yet another  poster}
  \author{W.~J.~Henney\altaffil{This is done with altaffil} and
    A.~N.~Other\altaffil{So is this} }
  \listofauthors{W.~J.~Henney \& A.~N.~Other}
  \indexauthor{Henney, W. J}
  \indexauthor{Other, A. N.}
  \abstract{This is a test abstract.   This is a test abstract.   This
    is a test abstract.   This is a test abstract.}

  \title{The abstract of a poster with a very long title to see how
    well it gets split up into lines}
  \author{A.~N.~Other\altaffil{The affiliation of A.~N.~Other is also
      very long so that it will go over several lines. His email
      address is anon@nowhere.org}}
  \listofauthors{A.~N.~Other}
  \indexauthor{Other, A. N.}
  \abstract{We will try and make this abstract longer so there is more
    stuff to put on the page and it doesn't look too empty. However, I
    have nothing whatsoever to say so will have to stop here. Let's
    repeat that a few times. We will try and make this abstract longer
    so there is more stuff to put on the page and it doesn't look too
    empty. However, I have nothing whatsoever to say so will have to
    stop here. Let's repeat that a few times. We will try and make
    this abstract longer so there is more stuff to put on the page and
    it doesn't look too empty. However, I have nothing whatsoever to
    say so will have to stop here. Let's repeat that a few times. We
    will try and make this abstract longer so there is more stuff to
    put on the page and it doesn't look too empty. However, I have
    nothing whatsoever to say so will have to stop here. Let's repeat
    that a few times.}
    
  \title[Comment to go in TOC]{The abstract of a second poster with a
    very long title to see how well it gets split up into lines}
  \author{A.~N.~Other\altaffil{The affiliation of A.~N.~Other is also
      very long so that it will go over several lines. His email
      address is anon@nowhere.org}}
  \listofauthors{A.~N.~Other}
  \indexauthor{Other, A. N.}
  \abstract{We will try and make this abstract longer so there is more
    stuff to put on the page and it doesn't look too empty. However, I
    have nothing whatsoever to say so will have to stop here. Let's
    repeat that a few times. We will try and make this abstract longer
    so there is more stuff to put on the page and it doesn't look too
    empty. However, I have nothing whatsoever to say so will have to
    stop here. Let's repeat that a few times. We will try and make
    this abstract longer so there is more stuff to put on the page and
    it doesn't look too empty. However, I have nothing whatsoever to
    say so will have to stop here. Let's repeat that a few times. We
    will try and make this abstract longer so there is more stuff to
    put on the page and it doesn't look too empty. However, I have
    nothing whatsoever to say so will have to stop here. Let's repeat
    that a few times.}
    
  \title{The abstract of a third poster with a very long title to see how
    well it gets split up into lines}
  \author{A.~N.~Other\altaffil{The affiliation of A.~N.~Other is also
      very long so that it will go over several lines. His email
      address is anon@nowhere.org}}
  \listofauthors{A.~N.~Other}
  \indexauthor{Other, A. N.}
  \abstract{We will try and make this abstract longer so there is more
    stuff to put on the page and it doesn't look too empty. However, I
    have nothing whatsoever to say so will have to stop here. Let's
    repeat that a few times. We will try and make this abstract longer
    so there is more stuff to put on the page and it doesn't look too
    empty. However, I have nothing whatsoever to say so will have to
    stop here. Let's repeat that a few times. We will try and make
    this abstract longer so there is more stuff to put on the page and
    it doesn't look too empty. However, I have nothing whatsoever to
    say so will have to stop here. Let's repeat that a few times. We
    will try and make this abstract longer so there is more stuff to
    put on the page and it doesn't look too empty. However, I have
    nothing whatsoever to say so will have to stop here. Let's repeat
    that a few times.}
    
  \title{The abstract of a fourth poster with a very long title to see how
    well it gets split up into lines}
  \author{A.~N.~Other\altaffil{The affiliation of A.~N.~Other is also
      very long so that it will go over several lines. His email
      address is anon@nowhere.org}}
  \listofauthors{A.~N.~Other}
  \indexauthor{Other, A. N.}
  \abstract{We will try and make this abstract longer so there is more
    stuff to put on the page and it doesn't look too empty. However, I
    have nothing whatsoever to say so will have to stop here. Let's
    repeat that a few times. We will try and make this abstract longer
    so there is more stuff to put on the page and it doesn't look too
    empty. However, I have nothing whatsoever to say so will have to
    stop here. Let's repeat that a few times. We will try and make
    this abstract longer so there is more stuff to put on the page and
    it doesn't look too empty. However, I have nothing whatsoever to
    say so will have to stop here. Let's repeat that a few times. We
    will try and make this abstract longer so there is more stuff to
    put on the page and it doesn't look too empty. However, I have
    nothing whatsoever to say so will have to stop here. Let's repeat
    that a few times.}
    
  \title{The abstract of a fifth poster with a very long title to see how
    well it gets split up into lines}
  \author{A.~N.~Other\altaffil{The affiliation of A.~N.~Other is also
      very long so that it will go over several lines. His email
      address is anon@nowhere.org}}
  \listofauthors{A.~N.~Other}
  \indexauthor{Other, A. N.}
  \abstract{We will try and make this abstract longer so there is more
    stuff to put on the page and it doesn't look too empty. However, I
    have nothing whatsoever to say so will have to stop here. Let's
    repeat that a few times. We will try and make this abstract longer
    so there is more stuff to put on the page and it doesn't look too
    empty. However, I have nothing whatsoever to say so will have to
    stop here. Let's repeat that a few times. We will try and make
    this abstract longer so there is more stuff to put on the page and
    it doesn't look too empty. However, I have nothing whatsoever to
    say so will have to stop here. Let's repeat that a few times. We
    will try and make this abstract longer so there is more stuff to
    put on the page and it doesn't look too empty. However, I have
    nothing whatsoever to say so will have to stop here. Let's repeat
    that a few times.}
    
  \title{The abstract of a sixth poster with a very long title to see how
    well it gets split up into lines}
  \author{A.~N.~Other\altaffil{The affiliation of A.~N.~Other is also
      very long so that it will go over several lines. His email
      address is anon@nowhere.org}}
  \listofauthors{A.~N.~Other}
  \indexauthor{Other, A. N.}
  \abstract{We will try and make this abstract longer so there is more
    stuff to put on the page and it doesn't look too empty. However, I
    have nothing whatsoever to say so will have to stop here. Let's
    repeat that a few times. We will try and make this abstract longer
    so there is more stuff to put on the page and it doesn't look too
    empty. However, I have nothing whatsoever to say so will have to
    stop here. Let's repeat that a few times. We will try and make
    this abstract longer so there is more stuff to put on the page and
    it doesn't look too empty. However, I have nothing whatsoever to
    say so will have to stop here. Let's repeat that a few times. We
    will try and make this abstract longer so there is more stuff to
    put on the page and it doesn't look too empty. However, I have
    nothing whatsoever to say so will have to stop here. Let's repeat
    that a few times.}
    
  \title{The abstract of a seventh poster with a very long title to see how
    well it gets split up into lines}
  \author{A.~N.~Other\altaffil{The affiliation of A.~N.~Other is also
      very long so that it will go over several lines. His email
      address is anon@nowhere.org}}
  \listofauthors{A.~N.~Other}
  \indexauthor{Other, A. N.}
  \abstract{We will try and make this abstract longer so there is more
    stuff to put on the page and it doesn't look too empty. However, I
    have nothing whatsoever to say so will have to stop here. Let's
    repeat that a few times. We will try and make this abstract longer
    so there is more stuff to put on the page and it doesn't look too
    empty. However, I have nothing whatsoever to say so will have to
    stop here. Let's repeat that a few times. We will try and make
    this abstract longer so there is more stuff to put on the page and
    it doesn't look too empty. However, I have nothing whatsoever to
    say so will have to stop here. Let's repeat that a few times. We
    will try and make this abstract longer so there is more stuff to
    put on the page and it doesn't look too empty. However, I have
    nothing whatsoever to say so will have to stop here. Let's repeat
    that a few times.}
  
  \title{An abstract to go in the list of abstracts, but not in the
    index of authors}
  \author{I.~A.~M. Hidden\altaffil{Puerto Escondido}}
  \listofauthors{I.~A.~M. Hidden}
  \abstract{The idea is that by omitting the
    \textbackslash{}indexauthor command, we can avoid an entry in the
    author index.}

\end{abstracts}


%% The Index: after latexing twice, run `makeindex -s rmaa' on this
%% file, then latex twice more
\printindex

%% if you rename this file to something other than `rmbooksamp', then
%% you must change the above line accordingly

%%
%% Plates: If you have any, then the following line will print
%% them out:

\outputplates

\end{document}

%%% Local Variables:
%%% mode: latex
%%% TeX-master: t
%%% End:

\end{document}
%</ads-container>
%    \end{macrocode}

% This is the |adshtml.sty| style file used by |ADS.tex|. It just
% defines a few 
% macros that might be used in abstracts, etc (really it should define
% them all). This now gets all its code from
% Sec~\ref{sec:user-commands}. 
%    \begin{macrocode}
%<*adshtml-style>
%% Nothing else needs to go here now. 
%</adshtml-style>
%    \end{macrocode}
% The init file for latex2html.
%    \begin{macrocode}
%<*latex2html-init>
$PDFTOCAIRO = "pdftocairo";
$PDFLATEX = "pdflatex";
$LATEX = "latex";
$MAX_SPLIT_DEPTH = 0;
$HTML_VERSION = "4.0,math";
$NO_SIMPLE_MATH = 1;
$ADDRESS = "<p><address><a href=\"mailto:rmaa\@astro.unam.mx\">rmaa\@astro.unam.mx</address></p>"; 
$INFO = "<p>This document was created using latex2html and the rmaa latex macros</p>"; 
$NO_NAVIGATION = 1;
1; 
%</latex2html-init>
%    \end{macrocode}
% This is the |adshtml.perl| file that is used by \LaTeX{}2HTML. Some
% bizarre \TeX{}y command sometimes show up so we tell it to ignore
% these. All the astronomy macros that come from AAS are processed in
% \TeX{}. That way they get made to images but with an |ALT| tag that
% has the original code, which is what will be picked up by |lynx|,
% with the result that they get passed intact into the final |.txt|
% file. 
%    \begin{macrocode}
%<*adshtml-perl>
package main;

$dum = &do_cmd_makeatletter();

sub do_cmd_RMAA {
    local($_) = @_;
    "RMxAA" . $_;
}

sub do_cmd_RMAASC {
    local($_) = @_;
    "RMxAASC" . $_;
}

sub do_cmd_etal {
    local($_) = @_;
    "et al\." . $_;
}

sub do_cmd_ibidrule {
    local($_) = @_;
    "<img src=\"foo\" alt=\"\ibidrule\">" . $_;
}


sub do_cmd_bibitem {
    local($_) = @_;
    local($key);
    local($opt, $dummy)=&get_next_optional_argument; 
    $key = &missing_braces unless
        ((s/$next_pair_pr_rx/$key = $2; ''/eo)
        ||(s/$next_pair_rx/$key = $2; ''/eo));
    "<p>" . $_; 
}

&ignore_commands (<<_RAW_ARG_CMDS_);
lowercase
adjustfinalcols
textsuperscript # {}
newline # []
thinmuskip
mskip
smallskip
medskip
bigskip
hspace # {}
hspacestar # {}
vspace # {}
vspacestar # {}
mathrm
_at_mathrm
spacefactor
discretionary
unhbox 
voidb_at_x 
penalty 
_at_M
_at_m
null
IeC
_RAW_ARG_CMDS_

&process_commands_in_tex (<<_RAW_ARG_CMDS_);
sun
earth
lesssim
gtrsim
sq
argdeg
arcmin
arcsec
fd
fh
fm
fs
fdg
farcm
farcs
fp
micron
la
ga
case # {} # {}
slantfrac # {} # {}
onehalf
onethird
twothirds
onequarter
threequarters
ubvr
ub
bv
vr
ur
aj
araa
apj
apjl
apjs
ao
apss
aap
aapr
aaps
azh
baas
jrasc
memras
mnras
pra
prb
prc
prd
pre
prl
pasp
pasj
qjras
skytel
solphys
sovast
ssr
zap
nat
iaucirc
aplett
apspr
bain
fcp
gca
grl
jcp
jgr
jqsrt
memsai
nphysa
physrep
physscr
planss
procspie
_RAW_ARG_CMDS_

1;                              # This must be the last line
%</adshtml-perl> 
%    \end{macrocode}

% This is a shell script |makeads.sh| that generates the |_ads.txt| file from
% |_ads.tex| by using \LaTeX{}2HTML and |lynx|.  \changes{4.0}{2019/11/13}{Make sure that lynx writes utf-8 output instead of latin-1.}
% \changes{2.2beta13}{2002/03/12}{Complete rewrite of |makeads.sh|. Now much 
% improved I think}
% \changes{4.0}{2019/11/13}{Fix the at-sign substitution so that it only eliminates \cs{@} and does not mess with email addresses.}

%    \begin{macrocode}
%<*ads-shell-script>
#! /bin/sh
if [ -z $1 ]; then
    echo "Usage: $0 basename"
    exit
fi
BASENAME=$1
ADSBASENAME=${BASENAME}_ads
AUXFILE=$BASENAME.aux
TMPFILE=tmp.$$


echo "Generating $ADSBASENAME.tex"

# For some reason, must double-up \\ to \\\\ but not \ to \\ - odd!
cat > $ADSBASENAME.tex <<EOF
%% PLEASE DO NOT EDIT THIS FILE!  It is generated automatically from
%% the script makeads.sh.  Any hand-edited changes will be lost
%% 
\documentclass{article}
\usepackage[latin1]{inputenc}
\usepackage{html}
\usepackage{adshtml}
\newcommand{\ADSfield}[2]{\%#1  #2\\\\}
\newenvironment{ADSarticle}{}{\par\bigskip}
\setlength{\parindent}{0pt}
\newcounter{istage}
\renewcommand{\ion}[2]{\setcounter{istage}{#2}#1 \Roman{istage}}
\begin{document}
\makeatletter
EOF
echo >> $ADSBASENAME.tex "\input{$BASENAME.ads2}" 
echo >> $ADSBASENAME.tex "\end{document}"



echo "Generating $ADSBASENAME.sed" 

# Find type of journal
grep '\\documentclass\[.*proceedings.*\]' $BASENAME.tex >/dev/null \
     && JOURNAL='RMxAC' \
     || JOURNAL='RMxAA'

# Find the year and volume number
YEAR=`grep '\\SetYear' $BASENAME.tex | sed 's/.*\\SetYear{\(.*\)}.*/\1/'`
VOLUME=`grep '\\SetVolume' $BASENAME.tex | sed 's/.*\\SetVolume{\(.*\)}.*/\1/'`

# Generate an auxilliary sed script 
cat > $TMPFILE <<EOF
s|\\\\newlabel{\(.*\):RMAAFirstPage}{{.*}{\(.*\)}}|s/\\\\\\\\ThisVolume{\1}/$YEAR, RMxAC, $VOLUME, \2/g|g
EOF

# Now go through the .aux file, looking for \newlabels and use the sed 
# script just generated in order to make another sed script, one that 
# will replace \ThisVolume{LABEL} with the correct reference
grep '^\\newlabel{.*:RMAAFirstPage}' $AUXFILE \
     | sed -f $TMPFILE > $ADSBASENAME.sed && rm $TMPFILE

# Apend a further sed command to get rid of \@, since redefining 
# it in _ads.tex didn't work
echo >> $ADSBASENAME.sed 's/\\\\@//g'

echo "Generating $BASENAME.ads2 from $ADSBASENAME.sed and $BASENAME.ads" 

# Apply this to the .ads file, generating the .ads2 file, which is included
# by the _ads.tex file.
sed -f $ADSBASENAME.sed $BASENAME.ads > $BASENAME.ads2

echo "Generating $ADSBASENAME.txt from $ADSBASENAME.tex and $BASENAME.ads2"

# Now run latex2html on the _ads.tex file
latex2html -no_navigation -info 0 -address 0 -html_version 4.0,math $ADSBASENAME

# Then, convert to utf8 with lynx and tidy it up a bit
lynx --justify=off --width=100 --dump --nolist --display_charset=utf-8 \
     $ADSBASENAME/$ADSBASENAME.html | cut -c4- \
     | sed 's/ibidrule/\\ibidrule/g' > $ADSBASENAME.txt
%</ads-shell-script>
%    \end{macrocode}

% \section{Epilogue}
% Err, that's it \dots

\endinput


% \endinput
% Local Variables:
% mode: doctex
% TeX-master: t
% End:
"

rmaa.ind: rmaa.idx
        @echo "REGENERATING INDEX"
        makeindex -s gind.ist rmaa

rmaa.gls: rmaa.glo
        @echo "REGENERATING CHANGELOG"
        makeindex -s gglo.ist -o rmaa.gls rmaa.glo

rmaa.pdf: rmaa.aux rmaa.ind rmaa.gls
        @echo "REGENERATING .pdf FILE"
        pdflatex "\nonstopmode% \def\fileversion{4.0} \def\filedate{2019/11/12}
% \MakeShortVerb{\|}
% \title{A \LaTeX{} document class for the Revista Mexicana de
% Astronom\'\i a y Astrof\'\i sica\thanks{This file has version
% \fileversion{} last revised \filedate{}.}}
% \author{William J. Henney\thanks{w.henney@irya.unam.mx}}
% \date{\filedate}
% \maketitle
% \paragraph{Abstract}
% This document describes the \texttt{rmaa} \LaTeX{}
% document class for the journal ``Revista Mexicana de
% Astronom\'\i a y Astrof\'\i sica''. The macros can be used for 
% \begin{enumerate}
% \item Preparing author manuscripts for the journal (including
% conference proceedings to be published in the ``Serie de
% Conferencias'')
% \item Typesetting an entire volume of the journal by inputting the
% individual article files from a container file. Facilities are
% included for automatic generation of TOC and Author Index. 
% \end{enumerate}
% 
% \tableofcontents
% \section{The documentation driver file}
% This is the part that generates the documentation that you may or
% may not be reading, which is done with the help of the Makefile
% discussed in section~\ref{sec:driver_make}. 
%    \begin{macrocode}
%<*driver>
\documentclass{ltxdoc}
\usepackage{hyperref}
\usepackage{palatino}
\EnableCrossrefs \CodelineIndex \RecordChanges
\setlength{\parindent}{0pt}
\def\MacroFont{\footnotesize\tt}
\begin{document} \DocInput{rmaa.dtx} \PrintIndex \PrintChanges \end{document}
%</driver>
%    \end{macrocode}

% \section{The implementation}
% 
% \changes{3.24}{2007/09/14}{Added user access to the macro version
% and date via \cs{rmVersion} and \cs{rmDate}}
% 
%    \begin{macrocode}
%<*package>
\newcommand{\rmaa@version}{3.26}
\newcommand{\rmaa@date}{2011/08/15}
\newcommand\rmVersion{\rmaa@version}
\newcommand\rmDate{\rmaa@date}
\NeedsTeXFormat{LaTeX2e}[1995/06/01]
%    \end{macrocode}
% \changes{4.0}{2019/11/02}{Change abstract to appear before
% resumen. Change RevMexAA to RMxAA for consistency with ADS.}
%
% \changes{3.22}{2007/03/22}{Allow proceedings articles to have
%   received/accepted dates} \changes{3.21}{2006/05/15}{Changes for
%   Silvia - abstracts are now shipped out to ADS.}
% \changes{3.19}{2006/01/19}{Changes for Silvia - remove adresses from
%   end, new abstract style} \changes{2.2beta}{2001/10/30}{ Fixed bug
%   whereby twocolumn option was not passed through to book
%   class. Finally sorted out the ``label on last page'' issue. }
% \changes{2.1}{2001/01/23}{ Fixed some inconsistencies in
%   \cs{nocite} and \cs{bibitem} uniqification Fixed the definition of
%   \cs{ion} so it will work in math mode and in contexts where the
%   typesize is different from normal (e.g., title, fig captions).
%   Modified the definition of \cs{@caption}, so we don't write any
%   list-of-figures data (we never use it anyway, and it just causes
%   problems with fragile commands in fig captions). }
% \changes{2.0}{2000/11/30}{ Added `newstyle' option, which should be
%   used in all new proceedings volumes (wider abstract, smaller
%   typeface in fig captions, etc.) Added commands for formatting
%   table captions. New demonstation document rmpostsamp.tex for
%   1-page posters. This will be the last release under a PRI
%   government (for a while...)}  \changes{1.9}{2000/06/06}{
%   Corrected position of page number in footer of articles that open
%   on even pages. Ensured that first article in each session opens on
%   an odd page, as do list of abstracts and author index. Added extra
%   space in TOC before the backmatter. Added command for editors,
%   \cs{BoldSessionTitles}, that makes produces bold titles for
%   session in the TOC.}  \changes{1.8}{2000/05/16}{ Fixed bug in
%   last page label with \cs{suppressfulladdresses} }
% \changes{1.7}{2000/05/11}{ Fixed more little things. Page labels
%   now work, even if you have no bibliography. You can output the
%   full addresses by hand with the \cs{outputfulladdresses}
%   command. Editors' Guide updated but still not finished. Revamped
%   splitbook.sh now produces PS file of each article for web page and
%   also an EPS file of each page for sending to the printers.}
% \changes{1.6}{2000/04/15}{ More cleaning up: full stops after
%   titles in index, running headers and footers in frontmatter,
%   fontsize in title of index, etc} \changes{1.5}{2000/04/15}{
%   Eliminated superscripts from titles in TOC and index}
% \changes{1.4}{2000/04/14}{ Added copyright and license info (GPL)}
% \changes{1.3}{2000/03/01}{ Corrected \cs{figcaption} to do the
%   caption formatting right Corrected font size of sec headings in
%   front matter} \changes{1.2}{2000/02/16}{ Added \cs{plate}
%   command} \changes{1.1}{2000/02/14}{ Added manuscript option,
%   corrected bugs} \changes{1.0}{2000/02/08}{ Now does a whole book
%   - lots of stuff tidied up} \changes{0.7}{1999/10/04}{ Increased
%   spacing between sections in single column styles. Suppressed
%   floats at bottom of pages (didn't work though) }
% \changes{0.6}{1999/09/24}{ Tweaked \cs{adjustfinalcols} (still not
%   perfect) Added proceedings option Added onecolumn/twocolumn
%   options (orthogonal to journal/proceedings) }
% \changes{0.5}{1999/09/23}{ Added \cs{fulladdresses} and
%   \cs{adjustfinalcols} (just about working)}
% \changes{0.4}{1999/09/22}{ Minor stylistic changes}
% \changes{0.3}{1999/06/28}{ Better mechanism for alternative
%   affiliations Added `guide' option for user guide}
% \changes{0.2}{1999/05/27}{ All basic functionality present, I
%   think} \changes{0.1}{1999/05/26}{ Just about working - still
%   missing Fig/Tab/Refs} \changes{0.0}{1999/05/21}{ So bored today,
%   decided to rewrite RMAA LaTeX macros from scratch}
%    \begin{macrocode}
\ProvidesClass{rmaa}[\rmaa@date\space v\rmaa@version\space RMxAA article
class]
%    \end{macrocode}
% \subsection{Option Handling}
% \label{sec:options}
% Declarations that are needed in options
% \changes{3.05}{2003/10/22}{Added natbib option}
% \changes{3.06}{2003/01/03}{Natbib option turned on permanently}
% \changes{3.07}{2003/01/11}{Added proof option}
% \changes{3.14}{2003/01/11}{Option to turn off natbib}
% \changes{3.18}{2003/10/22}{trying out doctex mode} 
% \changes{3.23}{2007/09/09}{Manuscript is now on by default.}
% \changes{3.26}{2011/08/15}{New option for review articles}
%    \begin{macrocode}
\newif\if@reviewarticle
\@reviewarticlefalse
\newif\if@proof
\@prooffalse
\newif\if@journal
\@journaltrue
\newif\if@printdates
\@printdatesfalse
\newif\if@proceedings
\@proceedingsfalse
\newif\if@manuscript
\@manuscripttrue
\newif\if@bigfont
\@bigfonttrue
\newif\if@guide
\@guidefalse
\newif\if@preprint
\@preprintfalse
\newif\if@rmtwocol
\@rmtwocoltrue
\newif\if@biblist
\@biblistfalse
\newif\if@book
\@bookfalse
\newif\if@natbib
\@natbibtrue
\newif\if@rmcite
\@rmcitefalse
\newif\if@titlevisible
\@titlevisibletrue
\newif\if@abstractpresent
\@abstractpresentfalse
\newif\if@printbanner
\@printbannerfalse
\newif\if@suppressfulladdresses
\@suppressfulladdressesfalse
\newif\if@newstyle
\@newstylefalse
\newif\if@debug
\@debugfalse
\newif\if@boldkeywords
\@boldkeywordsfalse
%    \end{macrocode}
% The option handling---mainly setting various booleans.
% \changes{2.2beta11}{2001/11/26}{Made twocol and
% newstyle the default for proceedings.}
% \changes{3.23}{2007/09/09}{Manuscript is now 10 pt instead of 12pt}
% \changes{3.23}{2007/09/10}{Full addresses are now automatically
% suppressed for proceedings style.}
% \changes{3.24}{2007/09/14}{New option \texttt{boldkeywords} that
% does the keywords in the old ugly way. Now off by default.}
%    \begin{macrocode}
\DeclareOption{manuscript}{\typeout{Manuscript for referee} 
  \@manuscripttrue\@journalfalse\@rmtwocolfalse
  \PassOptionsToClass{10pt}{article}}
\DeclareOption{bigfont}{\typeout{Big font manuscript requested by Gary
  Ferland} 
  \@manuscripttrue\@journalfalse\@rmtwocolfalse\@bigfonttrue
  \PassOptionsToClass{11pt}{article}}
\DeclareOption{journal}{\typeout{Main journal option}%
  \@journaltrue\@rmtwocoltrue\@printdatestrue\@manuscriptfalse
  \PassOptionsToClass{10pt}{article}}
\DeclareOption{guide}{\typeout{User guide option}%
  \@guidetrue\@journalfalse\@rmtwocoltrue\@manuscriptfalse
  \PassOptionsToClass{10pt}{article}}
\DeclareOption{preprint}{\typeout{Preprint option}%
  \@preprinttrue\@manuscriptfalse\PassOptionsToClass{10pt}{article}}
\DeclareOption{proof}{\typeout{Author proof option}%
  \@prooftrue\@manuscriptfalse}
\DeclareOption{proceedings}{\typeout{Proceedings option}%
  \@proceedingstrue\@journalfalse\@rmtwocoltrue\@newstyletrue
  \@suppressfulladdressestrue
  \@printdatesfalse\@manuscriptfalse\PassOptionsToClass{10pt}{article}} 
\DeclareOption{printdates}{\typeout{Force printing of received/accepted dates}%
  \@printdatestrue} 
\DeclareOption{book}{\typeout{Whole issue option}%
  \@booktrue} 
\DeclareOption{nonatbib}{\typeout{Turn off natbib support (required
    for compatibility with old rmaacite package).}%
  \@natbibfalse} 
\DeclareOption{newstyle}{\typeout{Use new layout style}%
  \@newstyletrue}
\DeclareOption{oldstyle}{\typeout{Use old layout style}%
  \@newstylefalse}
\DeclareOption{debug}{\typeout{Loads of annoying messages}%
  \@debugtrue}
%    \end{macrocode}
% These two need to be declared explicitly so that |\ExecuteOptions|
% actually does something.
%    \begin{macrocode}
\DeclareOption{twocolumn}{\@rmtwocoltrue}
\DeclareOption{onecolumn}{\@rmtwocolfalse}
\DeclareOption{twoside}{\PassOptionsToClass{\CurrentOption}{article}}
\DeclareOption{boldkeywords}{\@boldkeywordstrue}
%    \end{macrocode}
% Bold, upper-case keywords used to be the default. Now they must be
% turned on explicitly with the \texttt{boldkeywords} option. 
%    \begin{macrocode}
\DeclareOption{boldkeywords}{\@boldkeywordstrue}
%    \end{macrocode}
% Pass any other option along to the \texttt{article} class.
%    \begin{macrocode}
\DeclareOption*{\PassOptionsToClass{\CurrentOption}{article}}
%    \end{macrocode}
% Default options are to be two-sided and based on the standard
% \LaTeX{} \texttt{article} class,
% unless we are processing the entire book, in which case we base
% ourselves on the standard \texttt{book} class. Default is to use
% manuscript option unless otherwise specified. 
%    \begin{macrocode}
\ExecuteOptions{twoside,manuscript} \ProcessOptions\relax
\if@book
\if@rmtwocol\LoadClass[twocolumn]{book}\else\LoadClass{book}\fi
\newcommand{\abstract}{\relax}
\newcommand{\@maketitle}{\relax}
\else
\if@rmtwocol\LoadClass[twocolumn]{article}\else\LoadClass{article}\fi
\fi
%    \end{macrocode}
% Attempt to use natbib for citations 02 Jan 2003. 
% \changes{3.05}{2003/01/02}{If the natbib option is used, then set
% things up so that natbib's version of the bibliography is identical
% to the standard one.}
% \changes{3.09}{2003/01/23}{Explicitly require a recent version of
% natbib, so that we get a decent error message if the package is too
% old.} 
% \changes{3.23}{2007/09/09}{Now allow Bibtex to work properly} 
% \changes{3.25}{2009/10/01}{Change \cs{newcommand} tp \cs{def} when setting \cs{bibfont}. This allows it to work with all versions of natbib.} 
% \changes{3.25}{2010/02/06}{Updated copyright year.} 
% \changes{3.25}{2010/02/06}{Updated email and web addresses to \texttt{crya.unam.mx}.} 
%    \begin{macrocode}
\if@natbib
\RequirePackage{natbib}[1999/02/25]
\bibpunct{(}{)}{;}{a}{}{,}
\def\bibfont{\small}
\renewcommand\bibsection{\section*{REFERENCES}}
\setlength\bibsep{0pt}
\setlength\bibhang{1.5em}
\bibliographystyle{rmaa}
\fi
%    \end{macrocode}
% We want the \texttt{graphicx} package for the extended syntax of
% |\includegraphics|. We want the \texttt{ifthen} package for using in
% these macros (although we don't much yet). We explicitly load
% \texttt{babel} since some \LaTeX{} distributions messed up the
% hyphenation if we didn't.  
%    \begin{macrocode}
\RequirePackage{graphicx}
\RequirePackage{ifthen}
\IfFileExists{babel.sty}{%
%   \let\@restoreindent=\@afterindentfalse
  \RequirePackage[english]{babel}
%   \@restoreindent
  %% reinstate the definition of \@roman so as not to break .idx files
  \@ifundefined{org@@roman}{\relax}{\renewcommand{\@roman}[1]{\org@@roman{##1}}}
  }{%
  \typeout{Package babel.sty not found - don't worry about this}}
%    \end{macrocode}
% Try to load the \texttt{dcolumn} and \texttt{booktabs} packages if
% they can be found. These can produce nicer-looking tables. 
%    \begin{macrocode}
\IfFileExists{dcolumn.sty}{%
  \RequirePackage{dcolumn}
  }{%
  \typeout{Package dcolumn.sty not found - you will not be able to
    align columns on decimal points in tables}}
\IfFileExists{booktabs.sty}{%
  \RequirePackage{booktabs}
  }{%
  \typeout{Package booktabs.sty not found - your tables will not look
    as pretty as they ought too.}
  }
\IfFileExists{makeidx.sty}{%
  \RequirePackage{makeidx}
}{%
  \typeout{Package makeidx.sty not found - Author Index may not work.}
  }
%% commented out for moment since requires more modern version of LaTeX
%% \RequirePackage[latin1]{inputenc}
%    \end{macrocode}
% \changes{4.0}{2019/11/13}{Make \texttt{hyperref} and
% \texttt{etoolbox} into hard requirements, since we now use them
% extensively.}  We use hyperref to get internal and external links.
% But I don't like the colored boxes so I turn those off.  We use
% etoolbox for patching and robustifying at the moment.  In the
% future, I want to port a lot of other stuff over to it.  For
% instance, replace the ifthen package and the TeX ifs too.
%
% We target the latest version (1.7) of PDF, rather than the default
% of 1.5.  This reduces the chance of getting warnings from included
% graphics.  This is not strictly related to hyperref, but we put it
% here anyway.
%    \begin{macrocode}
\RequirePackage{etoolbox}
\RequirePackage{hyperref}
\hypersetup{hidelinks}
\pdfminorversion=7
%    \end{macrocode}
% We patch the \cs{@combinedblfloats} following a suggestion by Heiko
% Oberdiek at \url{https://tex.stackexchange.com/a/249743/3843}, which
% avoids problems when hyperlinks are split across
% pages. \changes{4.0}{2019/11/12}{Fix for issue where breaking
% hyperlink across pages gave an error.}
%    \begin{macrocode}
\urlstyle{same}
\newcommand\email{\begingroup \urlstyle{rm}\Url}
\patchcmd\@combinedblfloats{\box\@outputbox}{\unvbox\@outputbox}{}{%
  \errmessage{\noexpand\@combinedblfloats could not be patched}%
}%
%    \end{macrocode}
% \subsection{Other setup stuff}\label{sec:setup}
% Some declarations: 
%    \begin{macrocode}
\newif\if@needaddresses
\@needaddressestrue
\newif\if@frontmatter
\@frontmatterfalse
\newif\if@inpreamble
\@inpreambletrue
\newif\if@firstauthor
\@firstauthortrue
\newif\if@firstaffil
\@firstaffiltrue
\newif\if@firstaltaffil
\@firstaltaffiltrue
\newif\if@firstkwd
\@firstkwdtrue
\newif\if@noexpandADS
\@noexpandADSfalse
\newif\if@sessiontitlepages
\@sessiontitlepagesfalse
%    \end{macrocode}
% \begin{macro}{\@rmaaEmpty}
% \begin{macro}{\@rmaaArg}
% These are used in various commands to check if arg was left blank
%    \begin{macrocode}
\newcommand{\@rmaaEmpty}{}
\newcommand{\@rmaaArg}{}
%    \end{macrocode}  
% \end{macro}
% \end{macro}
%
% \subsection{Compatibility with old \LaTeX{} versions}
% \label{sec:compat}
% This package is based on the 1998/12/01 release of \LaTeX{}.
% Older releases may not define all the commands that we need.
% Luckily, we can use |\providecommand| to define these if necessary. 
% The following commands are not defined in the 1994/06/01 release. 
%
% \begin{macro}{\MakeUpperCase}
% \begin{macro}{\textsuperscript}
% \begin{macro}{\textbackslash}
% User commands:
%    \begin{macrocode}
\providecommand{\MakeUpperCase}[1]{\uppercase{#1}}
\providecommand{\textsuperscript}[1]{\ensuremath{^\textrm{#1}}}
\providecommand{\textbackslash}{\ensuremath{\backslash}}
%    \end{macrocode}  
% \end{macro}
% \end{macro}
% \end{macro}
% \begin{macro}{\protected@xdef}
% \begin{macro}{\restore@protect}
% \begin{macro}{\@unexpandable@protect}
% \begin{macro}{\hb@xt@}
% \LaTeX{} internal commands:
%    \begin{macrocode}
\providecommand{\protected@xdef}{%
        \let\@@protect\protect
        \let\protect\@unexpandable@protect
        \afterassignment\restore@protect
        \xdef}
\providecommand{\restore@protect}{\let\protect\@@protect}
\providecommand{\@unexpandable@protect}{\noexpand\protect\noexpand}
\providecommand{\hb@xt@}{\hbox to}
%    \end{macrocode}
% \end{macro}
% \end{macro}
% \end{macro}
% \end{macro}
%
% \subsection{Page layout}\label{sec:page-layout}
%
% The dimensions of the text were defined by Elizabeth. Text width is
% 169\,mm with equal margins on the two sides: 
%    \begin{macrocode}
\setlength{\textwidth}{169mm}
\setlength{\@tempdima}{\paperwidth}
\addtolength{\@tempdima}{-\textwidth}
\setlength{\oddsidemargin}{0.5\@tempdima}
\addtolength{\oddsidemargin}{-1in}
\setlength{\evensidemargin}{\oddsidemargin}
%    \end{macrocode}
% Text height is 225\,mm, top margin is 14\,mm. 
%    \begin{macrocode}
\setlength{\textheight}{225mm} 
\setlength{\topmargin}{-1in} 
\addtolength{\topmargin}{14mm}
%    \end{macrocode}
% Space between columns in two-column styles, header, footer,
% paragraph indentation:
%    \begin{macrocode}
\setlength{\columnsep}{9.5mm}
\setlength{\headsep}{0.7cm}
\setlength{\footskip}{0.75cm}
\setlength{\parindent}{1.5em}
%    \end{macrocode}
% 
% \begin{macro}{\@set@scalable@length}
% Try out a new scheme for dealing with various spacing dimensions and
% allowing the user to easily modify them. The command
% \begin{quote}
% |\set@scalable@length|\marg{name}\marg{defaultlength}\marg{stretch}\marg{shrink}
% \end{quote}
% does a 
% |\setlength|\marg{\cs{name}}\marg{defaultlength plus stretch minus shrink}
% and also defines the user command
% \textbackslash\meta{name}|Scale|\marg{factor} that can 
% be used to rescale length by \meta{factor}. Additionally, now
% defines the package command
% \textbackslash\meta{name}|@set@scale|\marg{FACTOR} 
% that also rescales the length by \meta{FACTOR}, so that the length
% becomes \meta{factor} $\times$ \meta{FACTOR} $\times$
% \meta{defaultlength}.   
%    \begin{macrocode}
\newcommand\@set@scalable@length[4]{%
  %% arg #1: NAME of length without leading `\'
  %%
  %% arg #2: default VALUE  
  \expandafter\newlength\csname @#1@default\endcsname
  \expandafter\setlength\csname @#1@default\endcsname{#2}%
  \expandafter\newlength\csname @@#1@@default\endcsname
  \expandafter\setlength\csname @@#1@@default\endcsname{#2}%
  %% arg #3: default STRETCH  
  \expandafter\newlength\csname @#1@stretch\endcsname
  \expandafter\setlength\csname @#1@stretch\endcsname{#3}%
  \expandafter\newlength\csname @@#1@@stretch\endcsname
  \expandafter\setlength\csname @@#1@@stretch\endcsname{#3}%
  %% arg #4: default SHRINK  
  \expandafter\newlength\csname @#1@shrink\endcsname
  \expandafter\setlength\csname @#1@shrink\endcsname{#4}%
  \expandafter\newlength\csname @@#1@@shrink\endcsname
  \expandafter\setlength\csname @@#1@@shrink\endcsname{#4}%
  %%
  %% User-defined scale \@NAME@scale is unity by default
  \expandafter\gdef\csname @#1@scale\endcsname{1}%
  %% Package-defined scale \@NAME@SCALE is unity by default
  \expandafter\gdef\csname @#1@SCALE\endcsname{1}%
  %%
  %% internal command \@reset@NAME should expand to 
  %% \setlength\@@NAME@@default{\@NAME@SCALE\@NAME@default}
  %% \setlength\NAME{\@NAME@scale\@@NAME@@default}
  %% in other words: scale x SCALE x default 
  \expandafter\gdef\csname @reset@#1\endcsname{%
    \expandafter\setlength\csname @@#1@@default\endcsname{%
      \csname @#1@SCALE\endcsname \csname @#1@default\endcsname
      }
    \expandafter\setlength\csname @@#1@@stretch\endcsname{%
      \csname @#1@SCALE\endcsname \csname @#1@stretch\endcsname
      }
    \expandafter\setlength\csname @@#1@@shrink\endcsname{%
      \csname @#1@SCALE\endcsname \csname @#1@shrink\endcsname
      }
    \expandafter\setlength\csname #1\endcsname{%
      \csname @#1@scale\endcsname \csname @@#1@@default\endcsname
      \@plus
      \csname @#1@scale\endcsname \csname @@#1@@stretch\endcsname
      \@minus
      \csname @#1@scale\endcsname \csname @@#1@@shrink\endcsname
      }}
  %%
  %% User command \NAMEScale should expand to 
  %% {\renewcommand\@NAME@scale{#1}\@reset@NAME}
  \expandafter\gdef\csname #1Scale\endcsname##1{%
    \expandafter\gdef\csname @#1@scale\endcsname{##1}%
    \csname @reset@#1\endcsname}
  %%
  %% Package command \NAME@set@SCALE should expand to 
  %% {\renewcommand\@NAME@SCALE{#1}\@reset@NAME}
  \expandafter\gdef\csname #1@set@SCALE\endcsname##1{%
    \expandafter\gdef\csname @#1@SCALE\endcsname{##1}%
    \csname @reset@#1\endcsname}
  %%
  %% initialize at default values with \@reset@NAME
  \csname @reset@#1\endcsname
  }
%    \end{macrocode}
% \end{macro}
% \begin{macro}{\parskip}
% First use of this new infrastructure: set interparagraph space to be
% zero but with a small amount of stretchiness.
%    \begin{macrocode}
%   \ifthenelse{\boolean{@manuscript}}{%
%     \@set@scalable@length{parskip}{0.5\baselineskip}{0.3\p@}{0.1\p@}
%   }{
    \@set@scalable@length{parskip}{0\p@}{0.3\p@}{0\p@}
%   }
%    \end{macrocode}
% \end{macro}
% 
% \begin{macro}{\abovesectionskip}
% \changes{4.0}{2019/11/12}{Use negative values for skip so that first paragraph after section headings is not indented.  This is something that I have been trying to fix for years!  It turns out that the anser was in the documentation for the standard command \cs{@startsection.}}
% \begin{macro}{\belowsectionskip}
% \begin{macro}{\aboveparaskip}
% \begin{macro}{\belowparaskip}
% Define some utility rubber lengths for use in
% section~\ref{sec:section-commands}. 
%    \begin{macrocode}
\newlength{\abovesectionskip}
\newlength{\belowsectionskip}
\newlength{\aboveparaskip}
\newlength{\belowparaskip}
%    \end{macrocode}
% The default values for these lengths are those traditionally used
% for the onecolumn style. 
%    \begin{macrocode}
\@set@scalable@length{abovesectionskip}{-3.5ex}{-4ex}{-2ex}
\@set@scalable@length{belowsectionskip}{2.3ex}{0ex}{2ex}
\@set@scalable@length{aboveparaskip}{3.25ex}{2ex}{0.5ex}
\@set@scalable@length{belowparaskip}{-1em}{\z@}{\z@}
%    \end{macrocode}
% \end{macro}\end{macro}\end{macro}\end{macro}
% We also make scalable some of the standard float parameters. These
% are the same defaults as in the standard article classes, but
% slightly more stretchy. For the twocolumn style, they are currently
% halved, see below. Note that I have also lumped |\abovecaptionskip|
% in with these, since it plays a similar role. There is also a
% |\belowcaptionskip| but it is set to zero in the standard classes. 
%    \begin{macrocode}
\@set@scalable@length{floatsep}{12\p@}{4\p@}{4\p@}
\@set@scalable@length{textfloatsep}{20\p@}{4\p@}{8\p@}
\@set@scalable@length{intextsep}{12\p@}{4\p@}{4\p@}
\@set@scalable@length{dblfloatsep}{12\p@}{4\p@}{4\p@}
\@set@scalable@length{dbltextfloatsep}{20\p@}{4\p@}{8\p@}
\@set@scalable@length{abovecaptionskip}{10\p@}{\z@}{\z@}
%    \end{macrocode}
% \begin{macro}{\@before@title@skip}         
% \begin{macro}{\@before@author@skip}        
% \begin{macro}{\@between@author@skip}       
% \begin{macro}{\@before@abstract@skip}      
% \begin{macro}{\@after@abstracthead@skip}   
% \begin{macro}{\@between@abstract@skip}     
% \begin{macro}{\@before@kwd@skip}           
% \begin{macro}{\@before@dates@skip}           
% Spacing parameters used in |\@maketitle|. These are not changed by
% default except by the user, so they don't need to go into
% |\@package@default@lengths|.  
%    \begin{macrocode}
\newlength{\@before@title@skip}         
\newlength{\@before@author@skip}        
\newlength{\@between@author@skip}       
\newlength{\@before@abstract@skip}      
\newlength{\@after@abstracthead@skip}   
\newlength{\@between@abstract@skip}     
\newlength{\@before@kwd@skip}           
\newlength{\@before@dates@skip}           
\@set@scalable@length{@before@title@skip}{2em}{\z@}{\z@}
\@set@scalable@length{@before@author@skip}{1.5em}{\z@}{\z@}
\@set@scalable@length{@between@author@skip}{0.5em}{\z@}{\z@}
\@set@scalable@length{@before@abstract@skip}{3ex}{1ex}{0.5ex}
\@set@scalable@length{@after@abstracthead@skip}{1ex}{\z@}{\z@}
\@set@scalable@length{@between@abstract@skip}{2ex}{1ex}{0.5ex}
\@set@scalable@length{@before@kwd@skip}{1ex}{0.5ex}{0.2ex}
\@set@scalable@length{@before@dates@skip}{2ex}{0.5ex}{0.2ex}
%    \end{macrocode}
% \end{macro}\end{macro}\end{macro}\end{macro}
% \end{macro}\end{macro}\end{macro}\end{macro}
% 
% Spacing above and below displayed equations also wants to be done
% like this, but it is more complicated since the standard article
% classes make these change with the fontsize. Something to do later,
% maybe. 
% 
%    \begin{macrocode}
%    \end{macrocode}
% \begin{macro}{\@package@default@lengths}
% The value actually used for these lengths is controlled by the
% \meta{SCALE} parameter, which depends on the style option that is in
% effect. This is set using the |\@package@default@lengths| command,
% which must be called every time the style changes (e.g., changing
% from onecolumn to twocolumn). 
%    \begin{macrocode}
\newcommand{\@package@default@lengths}{%
  \ifthenelse{\boolean{@rmtwocol}}{%
    %% twocolumn mode
    \abovesectionskip@set@SCALE{0.5}
    \belowsectionskip@set@SCALE{0.5}
    \aboveparaskip@set@SCALE{0.5}
    \belowparaskip@set@SCALE{1}
    \floatsep@set@SCALE{0.5}
    \textfloatsep@set@SCALE{0.5}
    \intextsep@set@SCALE{0.5}
    \dblfloatsep@set@SCALE{0.5}
    \dbltextfloatsep@set@SCALE{0.5}
    }{%
    %% onecolumn mode
    \abovesectionskip@set@SCALE{1}
    \belowsectionskip@set@SCALE{1}
    \aboveparaskip@set@SCALE{1}
    \belowparaskip@set@SCALE{1}
    \floatsep@set@SCALE{1}
    \textfloatsep@set@SCALE{1}
    \intextsep@set@SCALE{1}
    \dblfloatsep@set@SCALE{1}
    \dbltextfloatsep@set@SCALE{1}
    }%
  }
\@package@default@lengths
%    \end{macrocode}
% \end{macro}
% \begin{macro}{\RescaleLengths}
% \begin{macro}{\RescaleSecLengths}
% \begin{macro}{\RescaleFloatLengths}
% \begin{macro}{\RescaleTitleLengths}
% User commands that will fiddle with all the scalable lengths. 
%    \begin{macrocode}
\newcommand{\RescaleLengths}[1]{%
  \RescaleSecLengths{#1}\RescaleFloatLengths{#1}\RescaleTitleLengths{#1}}
\newcommand{\RescaleSecLengths}[1]{%
%% sections
  \abovesectionskipScale{#1}
  \belowsectionskipScale{#1}
  \aboveparaskipScale{#1}
  }
\newcommand{\RescaleFloatLengths}[1]{%
%% floats
  \floatsepScale{#1}
  \textfloatsepScale{#1}
  \intextsepScale{#1}
  \dblfloatsepScale{#1}
  \dbltextfloatsepScale{#1}
  \abovecaptionskipScale{#1}
  }
\newcommand{\RescaleTitleLengths}[1]{%
%% maketitle
  \@before@title@skipScale{#1}           
  \@before@author@skipScale{#1}        
  \@between@author@skipScale{#1}       
  \@before@abstract@skipScale{#1}      
  \@after@abstracthead@skipScale{#1}   
  \@between@abstract@skipScale{#1}     
  \@before@kwd@skipScale{#1}           
  \@before@dates@skipScale{#1}           
  }
%    \end{macrocode}
% \end{macro}\end{macro}\end{macro}\end{macro}
% \changes{2.2beta14}{2002/03/17}{Finished implementation of new
% infrastructure for scalable lengths. Converted the above/below
% section skips to use it. Also, all the floatsep parameters. Also,
% all the vertical lengths that are used by |\@maketitle|. }

% \begin{macro}{\MSextramargin}
% For the \texttt{manuscript} option we have an extra 0.3\,inch
% vertical margin and 1\,inch horizontal margin for marginal notes (by
% hand or with |\marginpar|. 
%
% \changes{3.23}{2007/09/09}{Reduce baselinestretch in
%   \texttt{manuscript} style to 1.2. No longer use raggedright for
%   manuscript.}
%
% \changes{3.24}{2007/09/14}{Further reduction of baselinestretch in
%   \texttt{manuscript} style to 1.1 since we now have wider margins.}
%
% \changes{3.24}{2007/09/14}{Split \texttt{manuscript}'s extra margin
%   into separate horizontal and vertical parts.}
% 
%    \begin{macrocode}
\newlength{\MSextraHmargin}\newlength{\MSextraVmargin}
\setlength{\MSextraHmargin}{1in} 
\setlength{\MSextraVmargin}{0.3in} 
\if@manuscript
   \renewcommand{\baselinestretch}{1.1}
   \addtolength{\textwidth}{-2\MSextraHmargin}
   \addtolength{\oddsidemargin}{\MSextraHmargin}
   \setlength{\evensidemargin}{\oddsidemargin}
   \addtolength{\textheight}{-2\MSextraVmargin}
   \addtolength{\topmargin}{\MSextraVmargin}
%    \raggedright
\fi
%    \end{macrocode}
% \end{macro}
% Float parameters: don't allow bottom floats by default. 
%    \begin{macrocode}
\renewcommand{\fps@figure}{tp}
\renewcommand{\fps@table}{tp}
%    \end{macrocode}
%
% \begin{macro}{\rmaa@AuthorWidthFrac}
% \begin{macro}{\rmaa@AbstractWidthFrac}
% These parameters control what fraction of the total text width can
% be used by the list of authors and the abstract/resumen,
% respectively.  
%    \begin{macrocode}
\newcommand{\rmaa@AuthorWidthFrac}{1.0}
\newcommand{\rmaa@AbstractWidthFrac}{0.75}
%    \end{macrocode}
% Use a wide abstract to save space in the new style.
%    \begin{macrocode}
\if@newstyle\renewcommand{\rmaa@AbstractWidthFrac}{1.0}\fi
%    \end{macrocode}
% \end{macro}
% \end{macro}
%
% Some things to consider in future versions. Trying to avoid too many
% widows, etc. Apparently, the following will work best with
% \cs{flushbottom}. 
%    \begin{macrocode}
%%\widowpenalty=10000
%%\clubpenalty=10000
%    \end{macrocode}
% Also, should perhaps try getting rid of underfull/overfull box warnings. 
% \changes{3.10}{2003/01/24}{Set large values for \cs{hbadness},
% \cs{vbadness}, \cs{hfuzz}, \cs{vfuzz}. To reduce amount of stuff
% that gets written to the terminal every time you compile something.}
%    \begin{macrocode}
\hbadness=10000
\vbadness=10000
\setlength\hfuzz{\maxdimen}
\setlength\vfuzz{\maxdimen}
%    \end{macrocode}
% \subsection{Various text styles for different contexts}
% \label{sec:text-styles} 
% \begin{macro}{\rmaa@AffilStyle}
% \begin{macro}{\rmaa@PagesStyle}
% \begin{macro}{\rmaa@YearStyle}
% \begin{macro}{\rmaa@VolumeStyle}
% \begin{macro}{\rmaa@VolumenStyle}
% \begin{macro}{\rmaa@JournalNameStyle}
% \begin{macro}{\rmaa@ReceivedStyle}
% \begin{macro}{\rmaa@AbstractHeaderStyle}
% \begin{macro}{\rmaa@JournalNameStyle}
% Style commands for the different fields used by |\maketitle| in
% constructing the title page of the article. These have the form
% |\rmaa@XYZStyle{TEXT}|. I'm not sure these were really a good idea,
% but never mind \dots 
%    \begin{macrocode}
\newcommand{\rmaa@AffilStyle}[1]{{\normalsize\rmfamily#1}}
\newcommand{\rmaa@PagesStyle}[1]{{\normalsize\rmfamily#1}}
\newcommand{\rmaa@YearStyle}[1]{{\normalsize\rmfamily#1}}
\newcommand{\rmaa@VolumeStyle}[1]{{\normalsize\rmfamily\bfseries#1}}
\newcommand{\rmaa@VolumenStyle}[1]{{\Large\sffamily\bfseries#1}}
\newcommand{\rmaa@JournalNameStyle}[1]{{\normalsize\itshape#1}}
\newcommand{\rmaa@ReceivedStyle}[1]{{\normalsize\rmfamily\itshape#1}}
\newcommand{\rmaa@AbstractHeaderStyle}[1]{%
  {\normalsize\rmfamily\MakeUpperCase{#1}}}
\newcommand{\rmaa@TitleStyle}[1]{%
  {\large\bfseries\rmfamily\boldmath\MakeUpperCase{#1}\par}}
%    \end{macrocode}
% \end{macro}
% \end{macro}
% \end{macro}
% \end{macro}
% \end{macro}
% \end{macro}
% \end{macro}
% \end{macro}
% \end{macro}
%
% \begin{macro}{\rmaa@FrontSecStyle}
% Headings of ``front sections'' (see section~\ref{sec:front-sections})
%    \begin{macrocode}
\newcommand{\rmaa@FrontSecStyle}[1]{%
  {\normalsize\bfseries\rmfamily\MakeUpperCase{#1}}}
%    \end{macrocode}
% \end{macro}
%
% \begin{macro}{\rmaa@TOCTitleStyle}
% \begin{macro}{\rmaa@TOCPartStyle}
% \begin{macro}{\rmaa@IndexTitleStyle}
% \begin{macro}{\rmaaTOCAuthorStyle}
% \begin{macro}{\rmaaIndexAuthorStyle}
% Commands to set style of TOC and index entries.
% Rewritten to avoid unnecessary font commands in the .toc and .ind files
%    \begin{macrocode}
\newcommand{\rmaa@TOCTitleStyle}[1]{\MakeUpperCase{#1}}
\newcommand{\rmaa@TOCPartStyle}[1]{\rmaa@TOCTitleStyle{#1}}
\newcommand{\rmaa@IndexTitleStyle}[1]{#1}
\newcommand{\rmaaTOCAuthorStyle}[1]{%
  {\itshape#1}}
\newcommand{\rmaaIndexAuthorStyle}[1]{%
  \protect\mbox{\bfseries#1}}
%    \end{macrocode}
% \end{macro}
% \end{macro}
% \end{macro}
% \end{macro}
% \end{macro}
%
% \begin{macro}{\rmaa@KeywordStyle}
% \begin{macro}{\rmaa@KeywordHeaderStyle}
% The style of the keywords depends on whether we are using the
% \texttt{newstyle} option or not. If we are, then they are small. 
%    \begin{macrocode}
\ifthenelse{\boolean{@newstyle}}{
  \newcommand{\rmaa@KeywordStyle}[1]{%
    {\small%
      \if@boldkeywords 
      \bfseries\MakeUpperCase{\raggedright #1}
      \else
      \raggedright #1
      \fi
    }}
  \newcommand{\rmaa@KeywordHeaderStyle}[1]{%
    {\small\itshape#1}}}
{
  \newcommand{\rmaa@KeywordStyle}[1]{%
    {\normalsize%
      \if@boldkeywords\bfseries\fi\raggedright #1}}
  \newcommand{\rmaa@KeywordHeaderStyle}[1]{%
    {\normalsize\itshape#1}}}
%    \end{macrocode}
% \end{macro}
% \end{macro}
%
% \subsection{Page styles}\label{sec:page-styles}
% \begin{macro}{\ps@rmaa}
% The running headers for normal pages:
%    \begin{macrocode}
\newcommand{\ps@rmaa}{%
%    \end{macrocode}
% nothing in either foot; 
%    \begin{macrocode}
  \renewcommand\@oddfoot{\relax}%
  \renewcommand\@evenfoot{\relax}%
%    \end{macrocode}
% page number and short author list in the even heads; 
%    \begin{macrocode}
  \renewcommand\@evenhead{\rlap{\thepage}\hfil\@shortauthor\hfil}%
%    \end{macrocode}
% short title and page number in the odd heads.
%    \begin{macrocode}
  \renewcommand\@oddhead{\hfil\@shorttitle\hfil\llap{\thepage}}
  }
%    \end{macrocode}
% \begin{macro}{\ps@rmaafirstpage}
% The header of the first page of each
% article:
%    \begin{macrocode}
\newcommand{\ps@rmaafirstpage}{%
  \renewcommand\@oddfoot{\hfil\thepage}%
%    \end{macrocode}
% footer corrected (06 Jun 2000) to have page on outside edge;
%    \begin{macrocode}
  \renewcommand\@evenfoot{\thepage\hfil}%
%    \end{macrocode}
% Header has journal name, volume, and page numbers for a regular
% article, or similar for special cases. 
% \changes{4.0}{2019/11/03}{Change to three-line header with copyright info and DOI.}
%    \begin{macrocode}
\if@preprint
\renewcommand\@oddhead{%
  \if@proceedings 
  \textsc{To appear in ``\@conftitle{} (\rmaa@YearStyle{\@year})''}\hfil
  \textit{RMxAC}
  \else
  \textsc{To appear in} \@journalname\hfil
  \fi}\fi%
\if@guide
\renewcommand\@oddhead{\textsc{Documentation} (\rmaa@YearStyle{\@year})\hfil}
\fi%
\if@manuscript
\renewcommand\@oddhead{Manuscript for \@journalname{} (\rmaa@YearStyle{\@year})\hfil}
\fi%
\if@journal
\renewcommand\@oddhead{%
  \parbox{\textwidth}{
    \@journalname, \rmaa@VolumeStyle{\@volume},
    \rmaa@PagesStyle{\@artpages}\
    (\rmaa@YearStyle{\@year})\hfil\\[1pt]
    {\small\textcopyright{} \@year:
      Instituto de Astronom\'\i{}a,
      Universidad Nacional Aut\'onoma de M\'exico}\hfil\\[1pt]
    {\small\@DOI@url}\hfil
  }%
}\fi%
  \renewcommand\@evenhead{\@oddhead}}
%    \end{macrocode}
% By default, use the |\ps@rmaa| page style.
%    \begin{macrocode}
\pagestyle{rmaa}
%    \end{macrocode}
% \end{macro}
% \end{macro}
%
% \subsection{Preamble commands}\label{sec:preamble} 
% Everything you need for setting the title page, TOC and index
% entries, etc.
% \subsubsection{Commands for the title page}
% \label{sec:title-commands}
% \begin{macro}{\resumen}
% \begin{macro}{\@resumen}
% \begin{macro}{\abstract}
% \begin{macro}{\@abstract}
% User commands for setting the resumen and abstract. 
% These should be used in the document header. The internal commands
% |\@resumen| and |\@abstract| are used to hold the actual text
% pending output via the |\maketitle| command. 
%    \begin{macrocode}
\newif\if@setresumen
\newcommand{\resumen}[1]{%
%%  \renewcommand{\@resumen}{{\selectlanguage{spanish}#1}}%
  \renewcommand{\@resumen}{{\ignorespaces#1}}%
  \@setresumentrue}
\renewcommand{\abstract}[1]{%
  \renewcommand{\@abstract}{{\ignorespaces#1}}
  \renewcommand\ADS@abstract{#1}
%% Not a good idea, since then author macros are not expanded and
%% cause errors later. However, by not doing it we are vulnerable to
%% fragile commands.....
  \if@noexpandADS\@onelevel@sanitize{\ADS@abstract}\fi
  \@abstractpresenttrue
  }
\newcommand{\ADSnoexpand}{\@noexpandADStrue}
%    \end{macrocode}
% \begin{macro}{\boldabstract}
% This is the version for doing 1-page conference proceedings. Still
% save the text for ADS but then write it out immediately as a bold
% paragraph. 
%    \begin{macrocode}
\newcommand{\boldabstract}[1]{%
  \gdef\ADS@abstract{#1}%
  \renewcommand{\@abstract}{#1}%
  \if@noexpandADS\@onelevel@sanitize{\ADS@abstract}\fi%
  \noindent\textbf{\boldmath #1}
  }
%    \end{macrocode}
% \end{macro}
% Default value in each case is to print out a warning if not set
%    \begin{macrocode}
  \newcommand{\@resumen}{%
    \@setresumenfalse \ClassWarningNoLine{rmaa}{No Spanish abstract
      (resumen) specified - if you do not speak Spanish, this will be
      translated by the editors}%
    \textbf{Favor de proporcionar un resumen en espa\~nol. If you are
      unable to translate your abstract into Spanish, the editors will
      do it for you.} }
  \newcommand{\@abstract}{\ClassWarningNoLine{rmaa}{No abstract
      specified}}
%    \end{macrocode}
% \end{macro}
% \end{macro}
% \end{macro}
% \end{macro}
%
% \begin{macro}{\title}
% \begin{macro}{\@rawtitle}
% \begin{macro}{\@TOCtitle}
% \begin{macro}{\rmaaIndexTitle}
% Similarly, |\title| is the preamble command to set the title of the
% article. This is used in various places, typeset in different
% styles, hence the use of the internal |\@rawtitle|, |\@title|,
% |\rmaaIndexTitle|, and |\@TOCtitle| commands, each of which uses the
% relevant style from section~\ref{sec:text-styles}
%    \begin{macrocode}
\newcommand{\@rawtitle}{Raw Title}
\newcommand{\@TOCtitle}{\rmaa@TOCTitleStyle{Title for TOC}}
\newcommand{\rmaaIndexTitle}{\rmaa@IndexTitleStyle{Title for index}.}
\newcommand{\@TOCaftertitle}{}
\renewcommand{\title}[1]{%
%%  {\ADS@sanitize\xdef\@rawtitle{#1}}
  \renewcommand\@rawtitle{\noexpand#1}
  \let\ADS@title=\@rawtitle
  \renewcommand{\@title}{\rmaa@TitleStyle{#1}}%
  \renewcommand{\@TOCtitle}{\rmaa@TOCTitleStyle{#1}\@TOCaftertitle}%
  \renewcommand{\rmaaIndexTitle}{\rmaa@IndexTitleStyle{#1\@TOCaftertitle}.}%
  }
%    \end{macrocode}
% \end{macro}
% \end{macro}
% \end{macro}
% \end{macro}
% 
% All the weird and whacky ways of indicating the author
% affiliations. 
% \begin{macro}{\affil}
% First, the way that sets it directly underneath the authors name.
% Note that |\affil| trumps all other affiliations for the purposes of
% getting into the ADS stuff (we use |\gdef|).   
% \dots 
%    \begin{macrocode}
\newcommand{\affil}[1]{%
  \\[1.5ex]\rmaa@AffilStyle{#1}\\[1.5ex]%
  \if@firstaffil\gdef\ADS@affil{#1}\@firstaffilfalse\fi%
  }
\renewcommand{\and}{and\\[1.5ex]}
%    \end{macrocode}
% \begin{macro}{\altaffil}
% Alternative affiliations are footnoted. 
%    \begin{macrocode}
\newcommand{\altaffil}[1]{%
  \if@firstaltaffil\gdef\ADS@affil{#1}\@firstaltaffilfalse\fi%
  \thanks{#1}}
%    \end{macrocode}
% \begin{macro}{\altaffilmark}
% \begin{macro}{\altaffiltext}
% Alternative-alternative mechanism for affiliations---should not be
% mixed with |\altaffil| (unless you know what you are doing) or there
% will be problems with numbering.
%    \begin{macrocode}
\newcommand{\altaffilmark}[1]{\textsuperscript{#1}}
\newcommand{\altaffiltext}[2]{%
  \if@firstaltaffil\gdef\ADS@affil{#2}\@firstaltaffilfalse\fi
  \stepcounter{footnote}%
  \protected@xdef\@thanks{\@thanks
    \protect\footnotetext[#1]{#2}}} 
%    \end{macrocode}
% \end{macro}
% \end{macro}
% \end{macro}
% \end{macro}
%
% \begin{macro}{\keywords}
% \begin{macro}{\@keywords}
% \begin{macro}{\@keywordheader}
% \label{keywords}Sundry minor preamble commands. All
% self-explanatory, I hope. |\keywords| is the old user interface
% (deprecated). New articles should use |\addkeyword| once for each
% keyword. That way the ADS record comes out right too (comma rather
% than em-dash separator).
% \changes{2.2beta12}{2001/12/07}{\cs{addkeyword} New user interface
%   for keywords.} 
% \changes{2.2beta13}{2002/03/11}{\cs{addkeyword} Improved behaviour
%   with .ads file. Use
%   \cs{ADS@sanitize} to change definitions of `~' and \cs{ion}.}
% \changes{3.23}{2007/09/09}{\cs{addkeyword} now signals a compilation
% error if the author tries to add more than 6 keywords.}
% \changes{3.24}{2007/09/14}{kwdcounter is now reset in
% \cs{@begofarticlehook} to prevent error in book style} 
%    \begin{macrocode}
\newcommand{\@kwdsep}{ --- }
\newcommand{\ADS@kwdsep}{, }
\newcounter{kwdcounter}\setcounter{kwdcounter}{0}
\newcommand{\addkeyword}[1]{%
  \addtocounter{kwdcounter}{1}
  \ifthenelse{\thekwdcounter > 6}{
    \ClassError{rmaa}{Too many keywords. A maximum of 6 keywords is allowed.}{}
  }{}
  \if@firstkwd%
  \renewcommand\@keywords{\if@boldkeywords\uppercase{#1}\else #1\fi}
  {\ADS@sanitize\def\ion{\plainion}\gdef\ADS@keywords{#1}}
  \@firstkwdfalse
  \else%
  \edef\@keywords{\@keywords\@kwdsep \if@boldkeywords\uppercase{#1}\else #1\fi}%
  {\ADS@sanitize\def\ion{\plainion}\xdef\ADS@keywords{\ADS@keywords\ADS@kwdsep #1}}%
  \fi}
\newcommand{\ADS@sanitize}{%
  \def\mathrm{}
  \def\@{\ }
  \def~{\ }
  \def\langle{<}
  \def\rangle{>}
  \def\log{log }
  \def\sin{sin }
  \def\cos{cos }
  \def\tan{tan }
  \def\fd{^d.}
  \def\fh{^h.}
  \def\fm{^m.}
  \def\fs{^s.}
  \def\fdg{�.}
  \def\farcm{'.}
  \def\farcs{''.}
  \def\arcdeg{�}
  \def\,{ }
  }
\newcommand{\keywords}[1]{%
  \renewcommand{\ADS@keywords}{#1}
  \renewcommand{\@keywords}{\uppercase{#1}}
  }
\newcommand{\@keywordheader}{\rmaa@KeywordHeaderStyle{Key Words:}}
\newcommand{\@keywords}{%
  \ClassWarningNoLine{rmaa}{No keywords specified}%
  Keywords go here
  }
%    \end{macrocode}
% \begin{macro}{\shorttitle}
% \begin{macro}{\shortauthor}
% \begin{macro}{\@shorttitle}
% \begin{macro}{\@shortauthor}
% Commands for setting the running headers.
%    \begin{macrocode}
\newcommand{\shorttitle}[1]{%
  \renewcommand\@shorttitle{\MakeUpperCase{#1}}}
\newcommand{\shortauthor}[1]{%
  \renewcommand\@shortauthor{\MakeUpperCase{#1}}}
\newcommand{\@shorttitle}{SHORT TITLE GOES HERE}
\newcommand{\@shortauthor}{SHORT AUTHOR GOES HERE}
%    \end{macrocode}
% \end{macro}
% \end{macro}
% \end{macro}
% \end{macro}
% \end{macro}
% \end{macro}
% \end{macro}


% \subsubsection{The full addresses}
% \label{sec:full-addresses}
% \begin{macro}{\fulladdresses}
% \begin{macro}{\@fulladdresses}
% \begin{macro}{\@fulladdheight}
% \begin{macro}{\@fulladddepth}
% \begin{macro}{\@fulladdbox}
% Aaghh, the full addresses! A perennial source of problems. 
%    \begin{macrocode}
\newlength{\@fulladdheight}
\newlength{\@fulladddepth}
\newsavebox{\@fulladdbox}
%% \fulladdresses command modified (11 May 2000) so that we can put
%% \fulladdresses{} in the article to suppress the output of the
%% addresses
\newcommand{\@FinalFigure}{\relax}
\newcommand{\FinalFigure}[1]{\renewcommand{\@FinalFigure}{#1}}
\newcommand{\fulladdresses}[1]{%
  \renewcommand{\@rmaaArg}{#1}
  \ifx\@rmaaArg\@rmaaEmpty
     \renewcommand{\@fulladdresses}{}
  \else
     \renewcommand{\@fulladdresses}{
       \normalsize%
       \ifx\@FinalFigure\relax\else\vfill\@FinalFigure\vfill\fi
%%       \vspace*{1ex}
       %% do addresses as a list (copied from thebibliography environment) 
       \list{\null}{\leftmargin 1.5em\topsep\z@\labelwidth\z@\labelsep\z@\itemindent -1.5em
         \itemsep\z@\parsep\z@\usecounter{enumi}}
       \sloppy\clubpenalty4000\widowpenalty4000
       \sfcode`\.=1000\relax
       #1
       \endlist
       }
  \fi
  }
\newcommand{\@fulladdresses}{Full addresses go here}
%    \end{macrocode}
% \begin{macro}{\suppressfulladdresses}
% New user-level command for suppressing full addresses should you so
% desire. Currently used in the single-page-poster style. 
%    \begin{macrocode}
\newcommand{\suppressfulladdresses}{\@suppressfulladdressestrue\@needaddressesfalse}
%    \end{macrocode}
% \end{macro}
% \end{macro}
% \end{macro}
% \end{macro}
% \end{macro}
% \end{macro}

% \begin{macro}{\printdates}
% \begin{macro}{\noprintdates}
% Command to activate the printing out of the received/accepted
% dates. This can also be done using the \texttt{printdates} option,
% which is the default for \texttt{journal}, but not for \texttt{proceedings}.
%    \begin{macrocode}
\newcommand{\printdates}{\@printdatestrue}
\newcommand{\noprintdates}{\@printdatesfalse}
%    \end{macrocode}
% \end{macro}
% \end{macro}
% \begin{macro}{\@outputaddressesifnecessary}
% Command to output the addresses if we haven't already done so
% This is called by: 
% \begin{enumerate}
% \item The end of |thebibliography| environment
% \item The |\outputfulladdresses| command
% \end{enumerate}
%    \begin{macrocode}
\newcommand{\@outputaddressesifnecessary}{%
  \ClassInfo{rmaa}{Calling @outputaddressesifnecessary}
  \ifthenelse{\boolean{@needaddresses}}{%
    %% If in one column style, or in the first column of a two column
    %% style then we can try to put it on this page 
    \ifthenelse{\boolean{@firstcolumn} \or \not\boolean{@rmtwocol}}{%
      \@outputaddresses}{%
      %% otherwise give warning it will be on a separate page
      \ClassWarning{rmaa}{Full addresses will go on a separate
        page (please see User Guide for how to remedy this)}}%
    %% If we haven't managed yet, put it on its own page
    \ifthenelse{\boolean{@needaddresses}}{%
      \if@debug{\ClassWarning{rmaa}{In this silly bit!}}\fi
      \clearpage\vfill\@outputaddresses}{}%
    }{
    %% If we don't need them, we still might need to do the eoa hook 
    \ifthenelse{\boolean{@suppressfulladdresses}}{%
      \if@debug{\ClassWarning{rmaa}{XXX}}\fi
%%      \if@book\else\@endofarticlehook\null\fi
      }{}
    }%
  }  
%    \end{macrocode}
% \end{macro}

% \begin{macro}{\outputfulladdresses}
% A user interface to the above command, in case there is no biblio
% (11 May 2000) 
%    \begin{macrocode}
\newcommand{\outputfulladdresses}{\@outputaddressesifnecessary}
%    \end{macrocode}
% \end{macro}

% \begin{macro}{\adjustfinalcols}
% This command starts a new column and should be 
% used to ensure that the columns balance on the last page. It also
% outputs the full 
% author addresses to ensure that this occurs in the left hand
% column. 
% The |\par| command is so that the current line will be taken account
% of when working out |\pagetotal|
%    \begin{macrocode}
\newcommand{\adjustfinalcols}{\if@twocolumn\goodbreak\@outputaddresses\newpage\noindent\fi}
%    \end{macrocode}
% \end{macro}

% \begin{macro}{\@outputaddresses}
% Command to do the work of outputting the addresses. 
% This can be called by 
% \begin{enumerate}
% \item an explicit |\adjustfinalcols|
% \item via |\@outputaddressesifnecessary|, q.v.
% \end{enumerate}
%
%    \begin{macrocode}
\newcommand{\@outputaddresses}{%
  % Allow option of suppressing the full addresses, in which case
  % we just run the hook if we are not in a book style
  \if@suppressfulladdresses
%%     \if@book\else\@endofarticlehook\null\fi
  \else
  %% save the full addresses in a box for output later
  \sbox{\@fulladdbox}{\parbox{\textwidth}{\@fulladdresses}}
  % calculate height of address box  
  \settoheight{\@fulladdheight}{\usebox{\@fulladdbox}}
  \settodepth{\@fulladddepth}{\usebox{\@fulladdbox}}
  \addtolength{\@fulladdheight}{\@fulladddepth}
  \if@debug
  { \typeout{Address total height: \the\@fulladdheight}
    \typeout{Text height: \the\pagetotal}
    \typeout{Goal height: \the\pagegoal}
    }
  \fi
  % reuse \@fulladdheight to mean the total height of text+addresses
  \addtolength{\@fulladdheight}{\pagetotal}
%%  \addtolength{\@fulladdheight}{-\baselineskip}
  \if@debug\typeout{Adjusted text + addresses height: \the\@fulladdheight}\fi
  % check it will fit on current page
  \ifdim\@fulladdheight<\pagegoal
     %% must shift it to the left
     %% to compensate for the hanging indentation in the reference list. 
     %% NOTE: this is not foolproof!
     %% NB. This will give overfull hbox warnings, which is a bit ugly.
     \vfill\noindent\if@biblist\hspace*{-\leftmargin}\fi
     \usebox{\@fulladdbox}
     % Set this to show we've done it
     \@needaddressesfalse
     % and assume we are at the end of the article
     % so that we do the EOA hook if not in the book style
%%     \if@book\else\@endofarticlehook\fi
  \else
     % if not, just do nothing. With luck, this function will get 
     % called again on the next page
     \if@debug{\typeout{Holding over full addresses to next page}}\fi
  \fi\fi
  }
%    \end{macrocode}
% \end{macro}

% \subsubsection{Commands for editors}
% \label{sec:editor-commands}
% Some user/editor commands for doing the TOC and index in the book
% option. 
% They are defined here rather in the book-specific part since they
% need to be always available in order to avoid breaking documents. 
% \begin{macro}{\listofauthors}
% \begin{macro}{\@rawauthors}
% \begin{macro}{\rmaaTOCauthors}
% The command |\listofauthors| formats an author list for use in the
% TOC. It should contain the same names as in the |\author| command,
% but without the affiliation garbage and using |\&| instead of
% |and|. 
%    \begin{macrocode}
\newcommand{\@rawauthors}{Raw Authors}
\newcommand{\rmaaTOCauthors}{\noexpand\rmaaTOCAuthorStyle{Author list for TOC}}
\newcommand{\listofauthors}[1]{%
  \renewcommand{\@rawauthors}{#1}
  \renewcommand\rmaaTOCauthors{\noexpand\rmaaTOCAuthorStyle{#1}}}
%    \end{macrocode}
% \begin{macro}{\indexauthor}\label{indexauthor}
% The command |\indexauthor| produces an index entry for each
% author. 
% This command must come \emph{after} the
% |\title| and the |\listofauthors| commands. Also, there should be one 
% |\indexauthor| command for \emph{each} author. We also add the
% author onto the end of the |\ADS@authors| macro and construct the
% |\ADS@authorinitial| by a most cunning method (see below), but only
% for the first author. 
% 
% No, No, NO! This was all wrong. Backed out completely: 
% \begin{quote}
% Completely rewritten (19 Dec 2002) since we now use the index class, so 
% as to be able to alphabetize accented characters. A side effect of this is
% that |\index| no longer expands out the argument after the |@|, so we 
% instead add all the |\index| commands to a list (|\indexauthor@deferlist|), 
% which is executed in the |\@begofarticlehook|. 
% 
% Note the stuff with |\toks0| is copied from Exercise 20.15 of the \TeX{}book---it is
% to make sure the latin1-accented name gets passed into the idx
% file. 
% \end{quote}
% We now go back to the old way of doing things, using an optional
% argument to the \cs{indexauthor} command that has the name without
% accents. 
% \changes{3.08}{2003/01/16}{Remerged changes from 3.06 that got
% dropped from 3.07 after a synchronization disaster. Basically, the
% accentless optional argument to \cs{indexauthor}. Must add the macro
% directories to my unison set-up to stop this sort of thing from happening.}
%    \begin{macrocode}
\newcommand\@isortkey{}
\newcommand\@emptykey{empty}
\newcommand\indexauthor[2][empty]{
  \renewcommand\@isortkey{#1}
  \ifx\@isortkey\@emptykey \renewcommand\@isortkey{#2}\fi
  \index{\@isortkey @{{\rmaaIndexAuthorStyle{#2}}\hskip1em plus1em minus0.4em
      \rmaaIndexTitle\hskip1em plus2em minus0.4em\penalty-100
      \rmaaTOCauthors}}
  \indexauthorADS{#2}
}
%%\def\indexauthor@deferlist{\relax}
%%\long\def\indexauthor#1{
%%  \def\@@name{#1} \toks0=\expandafter{\@@name}
%%  \toks1=\expandafter{\indexauthor@deferlist}
%%  \protected@xdef\indexauthor@deferlist{%
%%    \the\toks1
%%    \noexpand\index{\the\toks0@{\noexpand\rmaaIndexAuthorStyle{\the\toks0}\noexpand\hskip1em plus1em minus0.4em 
%%        \rmaaIndexTitle\noexpand\hskip1em plus2em minus0.4em\noexpand\penalty-100
%%        \rmaaTOCauthors}}
%%  }
%%  \indexauthorADS{#1}
%%  }
%%\xdef\indexauthor@deferlist{}
\newcommand\indexauthorADS[1]{%
  \if@firstauthor\protected@xdef\ADS@authorinitial{\firstletter{#1}}\fi
  \protected@xdef\ADS@authors{\ADS@authors\if@firstauthor\else ; \fi
    #1}
  \@firstauthorfalse
}
%    \end{macrocode}
% Note that the first |#1| sets the key for the index, making sure it
% actually starts with  
% first letter of the surname. Meanwhile, the |@{...}| expression,
% which starts with some  
% control sequences is what is written to the |.ind| file. 
% See the makeindex man page for details. 
% \end{macro}
% \end{macro}
% \end{macro}
% \end{macro}

% \begin{macro}{\KillAllAccents}
% \begin{macro}{\firstletter}
% These are utility functions used in the above. |lett| is a counter
% indicating the current letter in the string. |\firstletter| just
% prints the first letter of its argument.
%    \begin{macrocode}
\newcommand{\KillAllAccents}{\def\'{}\def\`{}\def\~{}\def\.{}\def\i{i}}
\newcommand\firstletter[1]{\@car #1\@nil}
%    \end{macrocode}
% \end{macro}
% \end{macro}
% \begin{macro}{\SetVolume}
% \begin{macro}{\SetMSnumber}
% \begin{macro}{\SetFirstPage}
% \begin{macro}{\SetYear}
% \begin{macro}{\ReceivedDate}
% \begin{macro}{\AcceptedDate}
% These are for the editors to set the volume number, etc. 
%    \begin{macrocode}
\newcommand{\SetVolume}[1]{%
  \renewcommand{\@volume}{#1}
  \edef\ADS@volume{\PadToFourWithDots{#1}}
  \renewcommand{\@volumen}{\rmaa@VolumenStyle{volumen #1}}%
  }
\newcommand{\SetMSnumber}[1]{\renewcommand{\@MSnum}{#1}}
%    \end{macrocode}
% Note that there is no real reason to use |\SetFirstPage| since all
% the page numbering is done automatically in the book style. 
%    \begin{macrocode}
\newcommand{\SetFirstPage}[1]{\setcounter{page}{#1}}
\newcommand{\SetYear}[1]{%
  \renewcommand{\@year}{#1}
  \renewcommand{\ADS@year}{#1}
  \renewcommand{\@anyo}{\rmaa@VolumenStyle{#1}}
  }
\newcommand{\ReceivedDate}[1]{\renewcommand{\@received}{#1}}
\newcommand{\AcceptedDate}[1]{\renewcommand{\@accepted}{#1}}
%    \end{macrocode}
%
% \begin{macro}{\DOI}
% User/editor command to enter the DOI, which will be typeset as a link.
% \changes{4.0}{2019/11/02}{New user command \cs{DOI} to typeset
% DOI in the header of first page and add it to ADS info.}
%    \begin{macrocode}
\newcommand\DOI[1]{%
  \def\@DOI@url{\url{https://doi.org/10.22201/ia.01851101p.#1}}%
  \def\ADS@doi{https://doi.org/10.22201/ia.01851101p.#1}%
}%
%    \end{macrocode}
% \end{macro}
%
% \begin{macro}{\ORCIDiD}
% \changes{4.0}{2019/11/12}{Add mechanism for giving ORCID iD links}
% \begin{macro}{\ORCIDiDfull}
% \begin{macro}{\@orcid@icon}
% \begin{macro}{\@orcid@iconbig}
% Commands to add icon links to ORCID iDs of authors. For example
% \begin{quote}
% |\ORCIDiD|\marg{0000-0000-0000-0000}
% \end{quote}
% to be used after each author name in the |\author|
% command and in the |\fulladdresses|.  Alternatively,
% |\OORCIDiDfull| will typeset the icon plus the full text of the
% link, but we don't use that at present.  The icon image is stored in |ORCIDiD_iconvector.pdf|, which needs to be on the input path.  In principle, we could make EPS versions if necessary.
%    \begin{macrocode}
\newcommand\@orcid@icon{\includegraphics[height=1.8ex]{ORCIDiD_iconvector}}
\newcommand\@orcid@iconbig{\includegraphics[height=2ex]{ORCIDiD_iconvector}}
\newcommand\ORCIDiD[1]{\href{https://orcid.org/#1}{\,\@orcid@icon}}
\newcommand\ORCIDiDfull[1]{%
  \href{https://orcid.org/#1}{\@orcid@iconbig}\,\url{https://orcid.org/#1}}
%    \end{macrocode}
% \end{macro}
% \end{macro}
% \end{macro}
% \end{macro}
%
% \begin{macro}{\@volume}
% \begin{macro}{\@volumen}
% \begin{macro}{\@artpages}
% \begin{macro}{\@year}
% \begin{macro}{\@anyo}
% \begin{macro}{\@received}
% \begin{macro}{\@accepted}
% \begin{macro}{\@MSnum}
% Default values for the containers used by above commands, plus a
% few more useful ones. 
%    \begin{macrocode}
\newcommand{\@volume}{00}
\newcommand{\@volumen}{\rmaa@VolumenStyle{00}}
\newcommand{\@artpages}{%
  \oldpageref{\articlename:RMAAFirstPage}--\oldpageref{\articlename:RMAALastPage}}
\newcommand{\@year}{\the\year}
\newcommand{\@anyo}{\rmaa@VolumenStyle{2007}}
\newcommand{\@journalname}{\rmaa@JournalNameStyle{%
    \if@proceedings\RMAASClong\else\RMAAlong\fi}}
\newcommand{\@received}{}
\newcommand{\@accepted}{}
\newcommand{\@MSnum}{0000}
\newcommand{\@DOI@prefix}{https://doi.org/10.22201/ia.01851101p}
\def\@DOI@url{\url{https://doi.org/10.22201/ia.01851101p.2099.00.00.00}}
%    \end{macrocode}
% \end{macro}
% \end{macro}
% \end{macro}
% \end{macro}
% \end{macro}
% \end{macro}
% \end{macro}
% \end{macro}
% \end{macro}
% \end{macro}
% \end{macro}
% \end{macro}
% \end{macro}
% \end{macro}
% 
% The following commands are most useful in the \texttt{book}
% style. However, the |\SetConfTitle| command can now be used in
% individual articles to set the header for the \texttt{preprint}
% style. 
%    \begin{macrocode}
\newcommand{\@reunionplace}{Lugar de la reuni\'on}
\newcommand{\@reuniondate}{Fecha de la reuni\'on}
\newcommand{\@conftitle}{Conference Title}
\newcommand{\SetReunion}[2]{\renewcommand{\@reunionplace}{#1}\renewcommand{\@reuniondate}{#2}}
\newcommand{\SetConfTitle}[1]{\renewcommand{\@conftitle}{#1}}
\newcommand{\@editors}{Juan Ram\'on de la Fuente and ``El Mosh''}
\newcommand{\SetEditors}[1]{\renewcommand{\@editors}{#1}}
\newcommand{\@editores}{Juan Ram\'on de la Fuente y ``El Mosh''}
\newcommand{\SetEditores}[1]{\renewcommand{\@editores}{#1}}
\newcommand{\@mes}{mes de la publicaci\'on}
\newcommand{\SetMes}[1]{\renewcommand{\@mes}{\rmaa@VolumenStyle{#1}}}
%    \end{macrocode}

% \subsection{Notes for Author Proofs and Communication between Editors}
% \label{sec:notes}
% We introduce a command \cs{Note} that should be used to describe
% any change made to the manuscript. These are a bit tricky to get
% right. In general we use \cs{marginpar} but this only works in outer
% paragraph mode, so we just bung the box where we are if we are in a
% figure caption or something. If we are in the preamble, then we just
% use a \marginpar as part of the \maketitle command.
%    \begin{macrocode}
\newif\if@arenotes
\@arenotesfalse
\newcounter{notecounter}
\setcounter{notecounter}{0}
\newcommand\NotesList{}
\newcommand\Note[1]{%
  \if@proof{%
    \@arenotestrue%
    \stepcounter{notecounter}%
    \if@inpreamble\else{%
      \NoteMark{\arabic{notecounter}}%
    }\fi
    \protected@xdef\NotesList{\NotesList 
      \protect\item #1}%
  }\fi%
}
\newcommand\Changed[2]{Changed ``#1'' to ``#2''.}
\newcommand\NoteMark[1]{%
  \ifinner
  \raisebox{1ex}[0pt][0pt]{\makebox[0pt][l]{\framebox{\footnotesize\bfseries
        N#1}}}%
  \else
  \marginpar{\framebox{\footnotesize\bfseries N#1}}%
  \fi
}
\newcommand\RefNote[1]{\if@proof\NoteMark{\ref{#1}}\relax\fi}
\newcommand\NoteLabel{\protect\label}
\newcommand\OutputNotes{%
  \bgroup
  \onecolumn
  \section*{EDITORIAL NOTES}
  \begin{enumerate}
    \NotesList
  \end{enumerate}
  \egroup
}
\DeclareRobustCommand\cs[1]{\texttt{\char`\\#1}}

%    \end{macrocode}
% \subsection{Bibliographic Records for ADS}
% \label{sec:ADS}
% \changes{2.2beta11}{2001/11/27}{Just about finished all the ADS
% stuff now. Except for the references, which I have given up on for
% the time being.}
% \begin{macro}{\ADS@BibCode}
% \begin{macro}{\ADS@year}
% \begin{macro}{\ADS@journal}
% \begin{macro}{\ADS@volume}
% \begin{macro}{\ADS@firstpage}
% \begin{macro}{\ADS@authorinitial}
% We use the |\ADS@BibCode| command to hold the bibliographic code for
% ADS. It is 
% used in the |%R| 
% field of the bibliographic record for each article. 
% The format is 
% \begin{quote}
% \texttt{YYYYJJJJJVVVVMPPPPA}  
% \end{quote}
% where \texttt{YYYY} is the year, \texttt{JJJJJ} is the journal code,
% which is \texttt{RMxAA} or \texttt{RMxAC}, \texttt{VVVV} is the
% volume number (right-justified), \texttt{M} is not used by us,
% \texttt{PPPP} is the first page number (right-justified), and
% \texttt{A} is the initial of the first author's surname. The volume
% and page fields are padded with periods (`\texttt{.}'), which causes
% us some fun (this is done with the macro |\PadToFourWithDots| described
% below) . 
% 
% All these individual bits are held in the commands |\ADS@year|,
% |\ADS@journal|, |\ADS@volume|, |\ADS@firstpage|,
% |\ADS@authorinitial|. 
%    \begin{macrocode}
\newcommand{\SetADSBibCode}[1]{\renewcommand{\ADS@BibCode}{#1}}
\newcommand{\ADS@BibCode}
{\ADS@year\ADS@journal\ADS@volume.\PadToFourWithDots{\ADS@firstpage}\ADS@authorinitial}
\newcommand{\ADS@year}{2001} 
\newcommand{\ADS@journal}{\if@proceedings RMxAC\else RMxAA\fi}
\newcommand{\ADS@volume}{....} 
%    \end{macrocode}
% These two are just set to default values here, |\ADS@firstpage| is set
% in the |document| environment for included articles, in the
% |frontsection| environment, |\tableofcontents| command,  |index| and
% |abstracts| environments. Also by |\AtBeginDocument| for stand-alone
% articles. 
% 
% God knows how we are going to do |\ADS@authorinitial|!
%    \begin{macrocode}
\newcommand{\ADS@firstpage}{....}
\newcommand{\ADS@authorinitial}{A}
%    \end{macrocode}
% \end{macro}
% \end{macro}
% \end{macro}
% \end{macro}
% \end{macro}
% \end{macro}
% We also have a whole slew of other fields that can be set. These are
% all listed in Table~\ref{tab:ADS}. Most of them still have to be
% defined. 
% \begin{table}[!ht]\raggedleft
% \caption{ADS Bibliographic Records}\label{tab:ADS}
% \begin{tabular}{lp{8em}lp{8em}}\hline
% Code & Description & Macro Name & Comment \\ \hline
% |%R| & Bibliographic Code & |\ADS@BibCode| & required \\
% |%A| & Author list & |\ADS@authors| & required, generate from |\indexauthor| \\
% |%F| & Author affiliation & |\ADS@affil| & format: AA(\dots)
% AB(\dots) \dots \\
% |%J| & Journal Name & |\ADS@journalname| &  \\
% |%D| & Publication Date & |\ADS@fulldate| &  \\
% |%L| & Last Page & |\ADS@lastpage| & generate from label \\
% |%T| & Title & |\ADS@title| & required, generate from |\title| \\
% |%C| & Abstract Copyright & |\ADS@copyright| & set to IA-UNAM \\
% |%U| & URL for electronic document & |\ADS@url| & set to PDF file \\
% |%d| & DOI for article & |\ADS@doi| & set by |\DOI| \\
% |%K| & Keywords & |\ADS@keywords| & set equal to |\@keywords| \\
% |%A| & Abstract Text & |\ADS@abstract| & set equal to |\@abstract| \\
% |%R| & References & |\ADS@references| & generate from the
% bibliography  \\ \hline
% \end{tabular}
% \end{table}
% A lot of these macros are set by other commands, as mentioned in the
% Table, but some need to be set  explicitly here. 
% 
% \begin{macro}{\ADS@doi}
% \changes{4.0}{2019/11/08}{Save DOI for ADS}
% The user/editor command \cs{DOI} is used to set this. 
%    \begin{macrocode}
\newcommand{\ADS@doi}{https://doi.org/10.22201/ia.01851101p.0000.00.00.00}
%    \end{macrocode}
% \end{macro}
% 
% \begin{macro}{\ADS@journalname}
% This
% does not contain any accents at the moment so as not to cause
% problems---this maybe should be changed though. 
%    \begin{macrocode}
\newcommand{\ADS@journalname}{%
  \if@proceedings\@conftitle\space (Eds. \@editors) \fi%
  Revista Mexicana de Astronom\'\i{}a y Astrof\'\i{}sica%
  \if@proceedings\space (Serie de Conferencias)\fi%
  \space Vol. \@volume, pp. \ADS@firstpage--\ADS@lastpage\space (\@year)%
  \space (http://www.astroscu.unam.mx/\~{}rmaa/)
  }
%    \end{macrocode}
% \end{macro}
% \begin{macro}{\ADS@keywords}
% \begin{macro}{\ADS@abstract}
% \begin{macro}{\ADS@title}
% |\ADS@abstract| is straightforward, since we already have the text
% saved as |\@abstract|. We could do the same with |\ADS@keywords| but
% to avoid unwanted formatting we explicitly set it in |\keywords|
% (p.~\pageref{keywords}). |\ADS@title| is just set to |\@rawtitle|,
% which we were already using for another purpose. 
%    \begin{macrocode}
\let\ADS@title=\@rawtitle
\newcommand{\ADS@keywords}{No keywords specified}
\newcommand{\ADS@abstract}{No abstract specified}
%    \end{macrocode}
% \end{macro}\end{macro}\end{macro}
% \begin{macro}{\ADS@authors}
% |\ADS@authors| wants to be in ``Other, A. N.'' format, but with
% semi-colons between authors. The most natural way to do this seems to
% be with |\indexauthor| (p.~\pageref{indexauthor}). We append each
% author to this macro, which means we must explicitly clear it for
% each new article (done in |\article|). .   
%    \begin{macrocode}
\newcommand{\ADS@authors}{}
%    \end{macrocode}
% \end{macro}
% \begin{macro}{\ADS@fulldate}
% \begin{macro}{\ADS@lastpage}
% \begin{macro}{\ADS@copyright}
% \begin{macro}{\ADS@website}
% \begin{macro}{\ADS@url}
% \begin{macro}{\SetWebsite}
% \begin{macro}{\SetMonthNumeric}
% Other sundry data, together with user commands for bits that can be
% set by the Editors. The |\ADS@url| is constructed from the website
% (currently |http://www.astrosmo.unam.mx/rmaa/|,
% the journal descriptor (e.g., |RMxAC..10/|), followed by the
% path to the PDF file (e.g., |PDF/RMxAC..10_HEW1.pdf|). 
% 
%    \begin{macrocode}
\newcommand{\ADS@website}{http://www.astrosmo.unam.mx/rmaa/}
\newcommand{\ADS@url}
{\ADS@website\ADS@journal\ADS@volume\ADS@number/PDF/\ADS@journal\ADS@volume\ADS@number\_\articlename.pdf}
\newcommand{\ADS@number}{}
\newcommand{\ADS@fulldate}{00/\ADS@year} 
\newcommand{\ADS@lastpage}{00} 
\newcommand{\ADS@copyright}{(C) \ADS@year; Instituto de Astronom\'\i{}a,
  Universidad Nacional Aut\'onoma de M\'exico} 
\newcommand{\SetNumber}[1]{\renewcommand{\ADS@number}{-#1}}
\newcommand{\SetWebsite}[1]{\renewcommand{\ADS@website}{#1}}
\newcommand{\SetMonthNumeric}[1]{\renewcommand{\ADS@fulldate}{#1/\ADS@year}}
%    \end{macrocode}
% \end{macro}\end{macro}\end{macro}\end{macro}\end{macro}\end{macro}\end{macro}
% \begin{macro}{\ADS@affil}
% \begin{macro}{\ADS@references}
% These two are somewhat more problematical. Affiliation currently
% only works for the first author (see \dots) pending redesign of the
% user interface. References are now done by wrapping the
% |thebibliography| environment to open another input stream on the
% same file and search through for the |\bibitems|. 
%    \begin{macrocode}
\newcommand{\ADS@references}{}
\newcommand{\ADS@affil}{}
%    \end{macrocode}
% \end{macro}\end{macro}
% 
% Define two macros that expand to the strings
% |\begin{thebibliography}| and |\end{thebibliography}|. Note that the
% idea for this and much else is inspired by the code for the docstrip
% package. Also define |\StartBibX| in case authors use the more
% standard |\begin{thebibliography}{}| syntax. 
%    \begin{macrocode}
\edef\EndBib{\expandafter\@gobble\string\\end\string{thebibliography\string}}
\edef\StartBib{\expandafter\@gobble\string\\begin\string{thebibliography\string}}
\edef\StartBibX{\expandafter\@gobble\string\\begin\string{thebibliography\string}\string{\string}}
%    \end{macrocode}
% Define two new logical variables. |\ifAnotherline| is true if we
% want to read another line from the input stream. It is set to false
% when we encounter the |\end{thebibliography}|. |\ifInBib| is set to
% true when we frind the |\begin{thebibliography}| and tells us to
% start saving the following lines. 
%    \begin{macrocode}
\newif\ifAnotherLine
\newif\ifInBib
%    \end{macrocode}
% \begin{macro}{makeOther}
% We need a macro to change the catcodes of lots of things to other. 
%    \begin{macrocode}
\def\makeOther#1{\catcode`#1=12\relax}
%    \end{macrocode}
% \end{macro}
% \begin{macro}{\RereadBiblio}
% Define a macro to do all the fiddly stuff. This is used at the
% beginning of |thebibliography| environment (Sec~\ref{sec:biblio}).
%    \begin{macrocode}
\newread\tf@biblio
\newcommand{\RereadBiblio}{%
%    \end{macrocode}
% Open the current file again for reading
%    \begin{macrocode}
  \begingroup
  \immediate\openin\tf@biblio=\articlename
%    \end{macrocode}
% Everything in a group to be safe. Initialize logicals. Black magic
% to change
% catcodes and avoid blank lines being interpreted as |\par|. 
%    \begin{macrocode}
  \AnotherLinetrue\InBibfalse
  \makeOther\ \makeOther\\\makeOther\$%
  \makeOther\#\makeOther\^\makeOther\^^K%
  \makeOther\_\makeOther\^^A\makeOther\%%
  \makeOther\~\makeOther\{\makeOther\}\makeOther\&%
  \endlinechar-1\relax
%    \end{macrocode}
% Start looping through the file. Read each line to the macro
% |\Line|. If we have got to EOF or end of biblio then set
% |\AnotherLinefalse|.
%    \begin{macrocode}
  \loop
  \read\tf@biblio to \Line
  \ifx\Line\EndBib\AnotherLinefalse\InBibfalse\fi
  \ifeof\tf@biblio \AnotherLinefalse\fi
%    \end{macrocode}
% If we are in the biblio and the line is not blank then append it to
% the |\ADS@references| macro together with a newline. 
%    \begin{macrocode}
  \ifInBib
  \ifx\Line\empty\else\xdef\ADS@references{\ADS@references\Line^^J}\fi
  \fi
%    \end{macrocode}
% Check if the line we just saw was the start of biblio, then repeat
% the loop if necessary. 
%    \begin{macrocode}
  \ifx\Line\StartBib\InBibtrue\fi
  \ifx\Line\StartBibX\InBibtrue\fi
  \ifAnotherLine\repeat
  \immediate\closein\tf@biblio
  \endgroup
  }
%    \end{macrocode}
% \end{macro}
% 
% \begin{macro}{\PadToFourWithDots}
% Utility function for left-padding a number to 4 characters with
% periods. Used in 
% constructing |\ADS@BibCode|. This is a very simple algorithm and
% \emph{requires} that the argument actually be a number---if it is
% not then dire consequences will result.  
%    \begin{macrocode}
\newcommand{\PadToFourWithDots}[1]{%
  \ifnum #1 > 999 #1%
  \else\ifnum #1 > 99 .#1%
  \else\ifnum #1 > 9 ..#1%
  \else ...#1%
  \fi\fi\fi}
%%  \ifnum \count0 > 999 #1%
%%  \else\ifnum \count0 > 99 .#1%
%%  \else\ifnum \count0 > 9 ..#1%
%%  \else ...#1%
%%  \fi\fi\fi}
%    \end{macrocode}
% \end{macro}
%
% \begin{macro}{\AuthorADS}
% \begin{macro}{\ADS@affiliations}
% New mechanism for setting the affiliations for ADS. The author labels (AA, AB, \dots) are generated using two nested counters.  We use the |\robustify| command from \texttt{etoolbox} so that accented characters in the affiliation addresses do not blow up.  The list of accents may need to be expanded. 
%    \begin{macrocode}
\newcounter{AUTHOR}
\newcounter{author}[AUTHOR]
\stepcounter{AUTHOR}
\newcommand\ADS@affiliations{}
\robustify{\'}\robustify{\i}\robustify{\~}\robustify{\`}
\def\AuthorADS#1#2{%
  \begingroup
  \def\ORCIDiD##1{<ID system="ORCID">##1</ID>}
  \def\Email##1{<EMAIL>##1</EMAIL>}
  \def\Affil##1{##1}
  \stepcounter{author}
  \ifthenelse{\value{author} > 26}{\stepcounter{AUTHOR}\stepcounter{author}}{}
  \ifthenelse{\equal{\ADS@affiliations}{\@rmaaEmpty}}{%
    \xdef\ADS@affiliations{\Alph{AUTHOR}\Alph{author}(#2)}
  }{%
    \xdef\ADS@affiliations{\ADS@affiliations, \Alph{AUTHOR}\Alph{author}(#2)}
  }%
  \endgroup
}
%    \end{macrocode}
% \end{macro}
% \end{macro}
% %
% \begin{macro}{\ADS@dumpdata}
% \changes{4.0}{2019/11/13}{Write all author affiliations for ADS}
% \begin{macro}{\ADS@write}
% \begin{macro}{\ADS@openfile}
% |\ADS@dumpdata| is the macro that actually does the job of writing out all the
% data to an auxiliary file, using |\ADS@write| to write each
% field. The file is set up with |\ADS@openfile|, which currently is
% done in the TOC of a book.  The new multi-author affiliations are used if |\ADS@affiliations| has been set, otherwise it falls back to the single-author version |\ADS@affil|. 
%    \begin{macrocode}
\newcommand{\ADSfield}[2]{\relax}
\newcommand{\ADSvfield}[2]{\relax}
%%\renewcommand{\ADSfield}[12]{\%#1 #2}
\newcommand{\ADS@dumpdata}
{ 
  \ADS@write{\string\begin{ADSarticle}}
  \ADS@WRITE{\string\ADSfield{R}{\ADS@BibCode}}
  \ADS@write{\string\ADSfield{A}{\ADS@authors}}
  \ifthenelse{\equal{\ADS@affiliations}{\@rmaaEmpty}}{%
    {\ADS@sanitize\ADS@write{\string\ADSfield{F}{AA(\ADS@affil)}}}
  }{%
    {\ADS@sanitize\ADS@write{\string\ADSfield{F}{\ADS@affiliations}}}
  }
  \ADS@write{\string\ADSfield{J}{\ADS@journalname}}
  \ADS@write{\string\ADSfield{D}{\ADS@fulldate}}
  \ADS@write{\string\ADSfield{L}{\ADS@lastpage}}
  {\ADS@sanitize\ADS@write{\string\ADSfield{T}{\ADS@title}}}
  \ADS@write{\string\ADSfield{C}{\ADS@copyright}}
  \ADS@write{\string\ADSfield{d}{\ADS@doi}}
  \ADS@write{\string\ADSfield{U}{\ADS@url}}
  \ADS@write{\string\ADSfield{K}{\ADS@keywords}}
  {\ADS@sanitize\ADS@write{\string\ADSfield{B}{\ADS@abstract}}}
  \ADS@write{\string\ADSfield{Z}{\ADS@references}}
  \ADS@write{\string\end{ADSarticle}}
  }
%%\newcommand{\ADS@write}[1]{\if@book\addtocontents{ads}{#1}\fi}
%%\newcommand{\ADS@openfile}{\@starttoc{ads}}
\newcommand{\ADS@write}[1]{\set@display@protect\immediate\write\tf@ads{#1}}
\newcommand{\ADS@WRITE}[1]{\immediate\write\tf@ads{#1}}
\newcommand{\ADS@openfile}{%
  \newwrite\tf@ads%
  \immediate\openout\tf@ads\jobname.ads\relax
  \ADS@WRITE{\ADS@preamble}}
\newcommand{\ADS@closefile}{\ADS@WRITE{\ADS@postamble}}
\def\ADS@preamble{%
  %%  \string\documentclass{article}
  %%  \string\begin{document}
  }
\def\ADS@postamble{%
  %%  \string\end{document}
  \string%\string% Local Variables:^J 
  \string%\string% mode: LaTeX^J
  \string%\string% mode: Font-lock^J
  \string%\string% End:^J
  }
%    \end{macrocode}
% \end{macro}
% \end{macro}
% \end{macro}
% 
% \subsection{Automatic calculation of the page range}
% \label{sec:page-range}
% Another issue that has cost me a lot of grief. Still, it finally
% seems to be sorted now. All hail Jeff Goldberg. 
%
% Cunning trick to get the page numbers from |\pageref|.  Put a label
% on the first and last page of the document using |\AtBeginDocument|
% and |\AtEndDocument|. .  This way won't work for the book option.
% In that case, we do it by defining |\@endofarticlehook| and
% |\@begofarticlehook| as parallels to |\AtEndDocument| and
% |\AtBeginDocument|.  Note that we have to defer adding
% |\@begofarticlehook| to |\AtBeginDocument| to make sure it comes
% after the redefinitions of |\ref| and |\label|.
%    \begin{macrocode}
\if@book\AtEndDocument{\ADS@closefile}\else
 %\AtBeginDocument{\@begofarticlehook}
\AtEndDocument{\@endofarticlehook}
\fi
%    \end{macrocode}
% \begin{macro}{\@begofarticlehook}
% Since we are getting more and more things that need to be done at
% the beginning of every article, frontsection, etc we define a sister
% hook to 
% |\@endofarticlehook|. This hook does the following: 
% \begin{enumerate}
% \item Labels the first page
% \item Sets |\ADS@firstpage|
% \item runs and then clears |\indexauthor@deferlist|
% \item resets the keyword counter
% \end{enumerate}
% It is executed in the following circumstances: 
% \begin{enumerate}
% \item At |\begin{document}| for standalone files
% \item At |\begin{document}| for each included file
% \item At beginning of front sections 
% \item At beginning of TOC
% \item At beginning of Author Index
% \item At beginning of abstracts
% \end{enumerate}
% For some reason, this is not always within the right group, so we
% use |\xdef| just to be sure. 
%    \begin{macrocode}
\newcommand{\@begofarticlehook}{%
  \if@frontmatter\frontlabel{RMAAFirstPage}\else\label{RMAAFirstPage}\fi%
  \xdef\ADS@firstpage{\@arabic{\c@page}}%
  \if@book\else\makeindex\fi % needed so default index is defined in stand-alone articles
%%  \indexauthor@deferlist
%%  \xdef\indexauthor@deferlist{}
  \@inpreamblefalse
  \setcounter{kwdcounter}{0}
  }
%    \end{macrocode}
% \end{macro}
% \begin{macro}{\@endofarticlehook}
% This is a hook to be run at the end of each article
% currently, it just labels the last page and writes out all the ADS
% data to an auxiliary file.  
% It should not produce any printed output. 
% It is modelled on the |\AtEndDocument| command, to which it reduces
% in the case of a single article. 
%    \begin{macrocode}
\newcommand{\@endofarticlehook}{%
  \if@debug{\ClassWarning{rmaa}{***Doing Hook***}}\fi
  \if@debug{\ClassWarning{rmaa}{ADSabstract is: \ADS@abstract}}\fi
  \clearpage\rmaa@lastpagelabel\if@book\ADS@dumpdata\fi
  \ifx\NotesList\@rmaaEmpty\else\OutputNotes\fi
  \if@debug{\ClassWarning{rmaa}{Notes list is: \NotesList}}\fi
}
%    \end{macrocode}
% \begin{macro}{\rmaa@lastpagelabel}
% New version of last page label command copied from Jeff Goldberg's
% lastpage.sty package - we can't just use |\label| here in case the
% output routine for the last page has already been called. Instead,
% we first explicitly clear the page, step back the page counter,
% explicitly write to the .aux file, then restep the page counter
% again.
%
% But this will conflict with the hyperref package, which does labels differently.  Compensate for that by using the five-argument form.
% 
% Also set the |\ADS@LastPage| macro to the appropriate number. 
%    \begin{macrocode}
\newcommand{\rmaa@lastpagelabel}{%
  \addtocounter{page}{-1}%
  \immediate\write\@auxout{\string
    \newlabel{\articlename:RMAALastPage}{{}{\thepage}{}{}{}}}%
  \edef\ADS@lastpage{\thepage}%
  \addtocounter{page}{1}}
%    \end{macrocode}
% \end{macro}
% \end{macro}

% \subsection{Deal with explicit change in number of columns}
% \label{sec:columns}
% Sometimes, when formatting a book, we need to switch between
% one-column articles and two-column articles. For this to work, we
% need to wrap the standard \LaTeX{} commands to set the |\@rmtwocol|
% flag appropriately and also reset the package default values of all
% scalable lengths. .  
%    \begin{macrocode}
\let\latexonecolumn=\onecolumn
\let\latextwocolumn=\twocolumn
\renewcommand{\onecolumn}{\@rmtwocolfalse\@package@default@lengths\latexonecolumn}
\renewcommand{\twocolumn}{\@rmtwocoltrue\@package@default@lengths\latextwocolumn}
%    \end{macrocode}
% 
% \subsection{The actual formatting of the title}
% \label{sec:title-format}

% Note, we have to redefine the |\maketitle| command to use the
% numbered footnotes. This means we needn't use the atbeginend stuff. 
%    \begin{macrocode}
\renewcommand{\maketitle}{\par
  \begingroup
%    \end{macrocode}
% Don't switch to marked footnotes. 
% This is consistent with existing style but I'm not sure I like it
%    \begin{macrocode}
  %%  \renewcommand\thefootnote{\fnsymbol{footnote}}%
%%  \def\@makefnmark{\hbox to\z@{$\m@th^{\@thefnmark}$\hss}}%
  \def\@makefnmark{\textsuperscript{\@thefnmark}\hspace{-0.2em}}%
  \long\def\@makefntext##1{\parindent 1em\noindent
    \hbox to1.8em{\hss$\m@th^{\@thefnmark}$}##1}%
  \if@rmtwocol
  { \if@debug{\ClassWarning{rmaa}{TWO COLUMN}}\fi
    \twocolumn[\@maketitle]%
    }
  \else
  { \if@debug{\ClassWarning{rmaa}{ONE COLUMN}}\fi
    \newpage
    \global\@topnum\z@   % Prevents figures from going at top of page.
    \if@manuscript \@MSmaketitle \else \@maketitle \fi
    }
  \fi
  \if@manuscript\else\thispagestyle{rmaafirstpage}\@thanks\fi
  \ifx\NotesList\@rmaaEmpty\else{%
    \NoteMark{1--\arabic{notecounter}}\vspace*{-2\baselineskip}%
  }\fi%
  \endgroup
%    \end{macrocode}
% Don't reset footnote counter (not sure whether this is good thing
% or not).
%    \begin{macrocode}
  %%  \setcounter{footnote}{0}%
  \let\thanks\relax\let\affil\relax\let\altaffil\relax
  \let\maketitle\relax\let\@maketitle\relax
  \gdef\@thanks{}\gdef\@author{}\gdef\@title{}}
%    \end{macrocode}

% \begin{macro}{\@kwhwidth}
% \begin{macro}{\@kwtwidth}
% Redefine \cs{@maketitle} to include the abstract, resumen and
% keywords. Note that in the manuscripty style we use
% \cs{@MSmaketitle} instead.  
% \changes{3.16}{2003/04/07}{
%  Added extra vertical space between authors and accepted date.}
% \changes{3.26}{2011/08/15}{Add extra label to review articles}
%    \begin{macrocode}
\newlength{\@kwhwidth}\newlength{\@kwtwidth}
\newcommand\reviewname{Review}
\renewcommand{\@maketitle}{%
  \newpage
  \null
  \if@reviewarticle
  \smallskip
  {\large\bfseries\itshape\reviewname}%
  \bigskip
  \else
  \vskip\@before@title@skip%
  \fi
  \begin{center}%
    {\@title \par}%
    \vskip\@before@author@skip%
    \parbox{\rmaa@AuthorWidthFrac\textwidth}{\centering
      \lineskip\@between@author@skip%
      \@author
      \par}
    %% this should be replaced with `Submitted to...' for preprints
    \if@printdates
    \vskip\@before@dates@skip%
    \if@preprint
    \rmaa@ReceivedStyle{Submitted to RMxAA \@received}
    \else
    \rmaa@ReceivedStyle{Received \@received; accepted \@accepted}
    \fi
    \fi
    \if@manuscript
    \rmaa@ReceivedStyle{Submitted to RMxAA \@received}
    \fi
    \if@guide
    \rmaa@ReceivedStyle{Version \rmaa@version, \rmaa@date}
    \fi
    \vskip\@before@abstract@skip%
    \if@abstractpresent
    \parbox{\rmaa@AbstractWidthFrac\textwidth}{%
      \setlength{\parindent}{0.9cm}%
      \centerline{\rmaa@AbstractHeaderStyle{Abstract}}\par%
      \vskip\@after@abstracthead@skip%
      \if@newstyle\noindent\fi\@abstract
      \vskip\@between@abstract@skip%
      \centerline{\rmaa@AbstractHeaderStyle{Resumen}}\par%
      \vskip\@after@abstracthead@skip%
      \if@newstyle\noindent\fi\@resumen
      \if@setresumen\else\@abstract\fi % append english abstract if
                                       % spanish one not set
      \if@guide\else            %do not use keywords for guide
      \vskip\@before@kwd@skip%
      % Format the Key Words with hanging indentation:
      % They are put in a \parbox with width equal to the line width
      % minus the width of the header. 
      { % Widths of keyword
        \settowidth{\@kwhwidth}{\@keywordheader~}   % header and
        \setlength{\@kwtwidth}{\linewidth}          % keyword text
        \addtolength{\@kwtwidth}{-\@kwhwidth}       %
        \noindent\@keywordheader~\parbox[t]{\@kwtwidth}{%
          \rmaa@KeywordStyle{\boldmath\uppercase{\@keywords}}}}%
      \fi
      }
      \fi
  \end{center}%
  \vspace*{\afterkwdskip}
  \par
  }
\newlength{\afterkwdskip}\setlength{\afterkwdskip}{\z@}
%    \end{macrocode}
% \end{macro}
% \end{macro}
%
%
% \begin{macro}{\reviewarticle}
% \begin{macro}{\reviewtoc}
% \begin{macro}{\reviewlof}
%   User command to specify a review article. Must be used in
%   preamble. Main effect is to place the text ``Review'' before the
%   article title.  Also add command for setting a table of contents
%   and a list of figures at the beginning of the article.
%    \begin{macrocode}
\newcommand{\reviewarticle}{\@reviewarticletrue}
\newcommand\reviewtoc{%
  \section*{CONTENTS}%
  \@starttoc{rtc}%
}
\newcommand\reviewlof{%
  \section*{LIST OF FIGURES}%
  \@starttoc{lof}%
}
%    \end{macrocode}
% \end{macro}
% \end{macro}
% \end{macro}
% \begin{macro}{\l@section}
%    \begin{macrocode}
\renewcommand*\l@section[2]{%
  \ifnum \c@tocdepth >\z@
    \addpenalty\@secpenalty
    \addvspace{0.8em \@plus\p@}%
    \setlength\@tempdima{1.5em}%
    \begingroup
      \parindent \z@ \rightskip \@pnumwidth
      \parfillskip -\@pnumwidth
      \leavevmode %\bfseries
      \advance\leftskip\@tempdima
      \hskip -\leftskip
      #1\nobreak\hfil \nobreak\hb@xt@\@pnumwidth{\hss #2}\par
    \endgroup
  \fi}
%    \end{macrocode}
% \end{macro}
%
% \begin{macro}{\@MSmaketitle}
% A special version of |\@maketitle| for the \texttt{manuscript}
% option. This can deal with a multi-page abstract/resumen. Not very
% well, mind. 
%    \begin{macrocode}
\newcommand{\@MSmaketitle}{%
  \if@debug\ClassInfo{rmaa}{Using special MS maketitle command}\fi
  \newpage
  {
    \thispagestyle{rmaafirstpage}
%     \setlength{\parindent}{0pt}
    \null\vskip 2ex
    \begin{center}
      \@title\par
      \@author\par
      \rmaa@ReceivedStyle{%
        \ifx\@received\@rmaaEmpty
          Draft version: \@date
        \else
          Received: \@received
          \ifx\@accepted\@rmaaEmpty
          \else
            ; Accepted: \@accepted
          \fi
        \fi}\par
    \end{center}
    \@thanks
    \centerline{\rmaa@AbstractHeaderStyle{Resumen}}
    \noindent\@resumen\par\vskip 2ex
    \centerline{\rmaa@AbstractHeaderStyle{Abstract}}
    \noindent\@abstract\par\vskip 2ex
    \noindent\@keywordheader~\rmaa@KeywordStyle{\@keywords}\par
%     \@keywordheader~\rmaa@KeywordStyle{\uppercase{\@keywords}}\par
  }%
  }
%    \end{macrocode}
% \end{macro}

% \subsection{The sectioning commands}
% \label{sec:section-commands}

% \begin{macro}{\section}
% \begin{macro}{\subsection}
% \begin{macro}{\subsubsection}
% \begin{macro}{\paragraph}
% \begin{macro}{\subparagraph}
% Redefine all the standard commands to use these lengths. 
%    \begin{macrocode}
\renewcommand\section{\@startsection {section}{1}{\z@}%
                                   {\abovesectionskip}%
                                   {\belowsectionskip}%
                                   %% how do we get upper case here
                                   {\centering\reset@font\normalsize}}
\renewcommand\subsection{\@startsection{subsection}{2}{\z@}%
                                     {\abovesectionskip}%
                                     {\belowsectionskip}%
                                     {\centering\reset@font\normalsize\itshape}}
\renewcommand\subsubsection{\@startsection{subsubsection}{3}{\z@}%
                                     {\abovesectionskip}%
                                     {\belowsectionskip}%
                                     {\centering\reset@font\normalsize\itshape}}
\renewcommand\paragraph{\@startsection{paragraph}{4}{\z@}%
                                    {\aboveparaskip}%
                                    {\belowparaskip}%
                                    {\reset@font\normalsize\bfseries}}
\renewcommand\subparagraph{\@startsection{subparagraph}{5}{\parindent}%
                                      {\aboveparaskip}%
                                      {\belowparaskip}%
                                      {\reset@font\normalsize\bfseries}}
%    \end{macrocode}
% \end{macro}
% \end{macro}
% \end{macro}
% \end{macro}
% \end{macro}

% \begin{macro}{secnumdepth}
% Make sure that we number the subsubsections---this is done by
% default in the |article| class but not the |book| class, which is
% why we must set it explicitly here.
%    \begin{macrocode}
\setcounter{secnumdepth}{3}
%    \end{macrocode}
% \end{macro}
% \changes{2.2beta10}{2001/11/07}{Fixed \cs{secnumdepth} for \cs{book}
% style.}

% \begin{macro}{\thesubsection}
% \begin{macro}{\thesubsubsection}
% For subsection and subsubsection, we want the numbers to be roman
% but the title to be italic. This seems the easiest way to go about it.
%    \begin{macrocode}
\renewcommand{\thesubsection}{{\upshape\thesection.\@arabic\c@subsection}}
\renewcommand{\thesubsubsection}{\thesubsection{\upshape.\@arabic\c@subsubsection}}
%    \end{macrocode}
% \end{macro}
% \end{macro}

% \begin{macro}{\@seccntformat}
% Add ``.'' to the number here so we can remove it from |\thesection|,
% etc. Reduce gap between number and title WJH 27 Jan 1995
%    \begin{macrocode}
\def\@seccntformat#1{\csname the#1\endcsname .\hskip 0.4em}
%    \end{macrocode}
% \end{macro}

% \begin{macro}{\@sect}
% I didn't really want to mess with the |\@sect| command, but it was
% the only way I could manage to force upper case. Other than the
% commented parts, it is identical to the standard version. 
%    \begin{macrocode}
\def\@sect#1#2#3#4#5#6[#7]#8{
  \ifnum #2>\c@secnumdepth
     \let\@svsec\@empty
  \else
     \refstepcounter{#1}%
     \let\@@protect\protect
     \def\protect{\noexpand\protect\noexpand}%
     \edef\@svsec{\@seccntformat{#1}}%
     \let\protect\@@protect
  \fi
%    \end{macrocode}
% This is the bit I changed (27 Jan 1995). Use upper case for
% |\section|, use mixed case for the rest. 
%    \begin{macrocode}
  \ifnum #2=1                   
     \def\sec@upcase##1{\uppercase{##1}}%
     \def\sec@shape{\upshape}%
  \else                         
     \def\sec@upcase##1{\relax{##1}}%
     \def\sec@shape{\itshape}%
  \fi
  \@tempskipa #5\relax
  \ifdim \@tempskipa>\z@
     \begingroup #6\relax
          \@hangfrom{\hskip #3\relax\@svsec}%
                    {\interlinepenalty \@M \sec@upcase{#8}\par}%
     \endgroup
     \csname #1mark\endcsname{#7}%
%    \end{macrocode}
% Only write TOC entries if doing the Guide style. 
%    \begin{macrocode}
     \if@guide 
        \addcontentsline{toc}{#1}{%
          \ifnum #2>\c@secnumdepth \else
             \protect\numberline{\csname the#1\endcsname}%
          \fi
          #7}%
     \fi
%    \end{macrocode}
% Also allow per-article TOC for review papers
%    \begin{macrocode}
     \if@reviewarticle
        \ifnum #2>2 \else
           \addcontentsline{rtc}{#1}{%
             \protect\numberline{\csname the#1\endcsname}%
             \sec@upcase{\sec@shape #7}}%
        \fi
     \fi
  \else
     \def\@svsechd{#6\hskip #3\relax  %% \relax added 2 May 90
       \@svsec #8\csname #1mark\endcsname{#7}%
       \if@guide
          \addcontentsline{toc}{#1}{%
            \ifnum #2>\c@secnumdepth \else
               \protect\numberline{\csname the#1\endcsname}%
            \fi
            #7}%
       \fi
       }%
  \fi
  \@xsect{#5}}
%    \end{macrocode}
% \end{macro}


% \subsection{Appendices}
% \label{sec:appendices}
% Make sure that appendices in individual articles work properly. We
% copy the definition from the article class (the book one doesn't do
% what we want).  
% \begin{environment}{appendices}
%   This uses an environment |appendices| inside of which each
%   |\section| will become a separate appendix.  Note that the
%   appendices must come before the references (otherwise the full
%   addresses mechanism will not work properly). We make the equation
%   numbering be subordinate to the appendix (equations A12, B6,
%   etc). Finally, we have a medskip followed by APPENDICES to
%   introduce them. \changes{3.17}{2002/04/08}{ Added appendices environment
%   that does the right thing.}
%    \begin{macrocode}
\newenvironment{appendices}{\par
  \begingroup
  \setcounter{section}{0}%
  \setcounter{subsection}{0}%
  \renewcommand\thesection{\@Alph\c@section}
  \renewcommand\theequation{\thesection\arabic{equation}}
  \medskip
  \section*{APPENDICES}}
{\endgroup}
%    \end{macrocode}
% \end{environment}
% \begin{environment}{appendix}
%   If there is only one appendix, then use |appendix| environment
%   instead.  \changes{4.0}{2019/11/12}{ Added
%   appendix environment for case where there is only one section.}
%    \begin{macrocode}
\renewenvironment{appendix}{\par
  \begingroup
  \setcounter{section}{0}%
  \setcounter{subsection}{0}%
  \renewcommand\thesection{\@Alph\c@section}
  \renewcommand\theequation{\thesection\arabic{equation}}
  \medskip
  \section*{APPENDIX}}
{\endgroup}
%    \end{macrocode}
% \end{environment}

% \subsection{Floating thingies}
% \label{sec:floats}   

%% 
%    \begin{macrocode}
\renewcommand{\thetable}{\@arabic\c@table}
\renewcommand{\thefigure}{\@arabic\c@figure}
\renewcommand{\figurename}{Fig.}
\renewcommand{\tablename}{TABLE}
%\renewcommand{\fnum@table}{\MakeUpperCase{\tablename}~\thetable}
%    \end{macrocode}
% For some reason, I have to write `TABLE' explicitly here
% If I use |\tablename|, it gets converted to mixed case---God knows
% why. 
% |\figurename| is also being ignored. 
%    \begin{macrocode}
\renewcommand{\fnum@table}{TABLE~\thetable}
\renewcommand{\fnum@figure}{Fig.~\thefigure}
%    \end{macrocode}

% Do captions in 9 pt face in the main journal style and in new-style
% anything. 
%    \begin{macrocode}
\long\def\@makefigcap#1#2{%
  \if@proceedings\else\small\fi
  \if@newstyle\small\fi 
  \vskip\abovecaptionskip
  \sbox\@tempboxa{#1. #2}%
  \ifdim \wd\@tempboxa >\hsize
    #1. #2 \par
  \else
    \global \@minipagefalse
    \hb@xt@\hsize{\hfil\box\@tempboxa\hfil}%
  \fi
  \vskip\belowcaptionskip}
%    \end{macrocode}

% Make sure table caption is upper case. 
%    \begin{macrocode}
\long\def\@maketabcap#1#2{%
  \hbox to\hsize{\hfil #1 \hfil}
  \vskip 6pt 
  \centering\MakeUpperCase{#2}\par
  \vskip 6pt}
%    \end{macrocode}

%    \begin{macrocode}
\newlength{\@rottabcapwidth}
\long\def\@makerottabcap#1#2{%
  \setbox\@tempboxa\hbox{\MakeUpperCase{#2}}%
  \ifdim\wd\@tempboxa >.8\vsize
     \setlength\@rottabcapwidth{.8\vsize}
  \else
     \setlength\@rottabcapwidth{\wd\@tempboxa}
  \fi
  \rotatebox{90}{%
    \begin{minipage}{\vsize}\centering
      \begin{minipage}{\@rottabcapwidth}\centering
        #1\\ \vskip 6pt \MakeUpperCase{#2}
      \end{minipage}%
    \end{minipage}
    }%
  \hskip 6pt
  }

\newenvironment{Sideways}
    {\begin{sideways}\begin{minipage}{\vsize}\centering}
    {\end{minipage}\end{sideways}}

\newenvironment{Sidewaystable}
{%
  \let\oldcaption=\caption%
  \def\caption{
    \def\@captype{table}\renewcommand{\@makecaption}{\@maketabcap}
    \oldcaption}
  \begin{sidewaystable}\centering}
  {\end{sidewaystable}\let\caption=\oldcaption}
\newenvironment{Sidewaystable*}
{%
  \let\oldcaption=\caption%
  \def\caption{
    \def\@captype{table}\renewcommand{\@makecaption}{\@maketabcap}
    \oldcaption}
  \begin{sidewaystable*}\centering}
  {\end{sidewaystable*}\let\caption=\oldcaption}
%    \end{macrocode}

% Not sure why we need to stick these lines in. 
%    \begin{macrocode}
\renewenvironment{figure}
                 {\renewcommand{\@makecaption}{\@makefigcap}\@float{figure}}
                 {\end@float}
\renewenvironment{figure*}
                 {\renewcommand{\@makecaption}{\@makefigcap}\@dblfloat{figure}}
                 {\end@dblfloat}
\renewenvironment{table}
                 {\renewcommand{\@makecaption}{\@maketabcap}\@float{table}}
                 {\end@float}
\renewenvironment{table*}
                 {\renewcommand{\@makecaption}{\@maketabcap}\@dblfloat{table}}
                 {\end@dblfloat}
%    \end{macrocode}

% Redefine |\@caption| so that we don't bother with the list of
% figures. This saves us having to worry about fragile commands in the
% captions. 
%    \begin{macrocode}
\long\def\@caption#1[#2]#3{%
  \par
%%  \addcontentsline{\csname ext@#1\endcsname}{#1}%
%%    {\protect\numberline{\csname the#1\endcsname}{\ignorespaces #2}}%
  \begingroup
    \@parboxrestore
    \if@minipage
      \@setminipage
    \fi
    \normalsize
    \@makecaption{\csname fnum@#1\endcsname}{\ignorespaces #3}\par
  \endgroup}
%    \end{macrocode}

% Make tables not quite so cramped. Allow |\hline\hline| to give a
% poor man's thick rule. 
%    \begin{macrocode}
\renewcommand{\arraystretch}{1.15}
\setlength{\doublerulesep}{0.4pt}
%    \end{macrocode}

% \begin{macro}{\tablecols}
% \begin{macro}{\@tablecols}
% \begin{macro}{\tabcaptioncontinued}
% |\tablecols| is a user command for setting the number of columns. It
% is used by commands that try and span the entire table (I think
% there must be a better way of acheiving
% this). |\tabcaptioncontinued| is 
% a user command for setting the caption of a continued table on the
% next page. 
%    \begin{macrocode}
\newcommand{\tablecols}[1]{\renewcommand{\@tablecols}{#1}}
\newcommand{\@tablecols}{1}
\newcommand{\tabcaptioncontinued}{%
  \multicolumn{\@tablecols}{c}{\addtocounter{table}{-1}TABLE~\thetable\ (CONTINUED)} \\ \\
  }
%    \end{macrocode}
% \end{macro}
% \end{macro}
% \end{macro}

% \subsubsection{Footnotes in tables}
% \label{sec:table-footnotes}
% \begin{macro}{\tablenotemark}
% \begin{macro}{\tablenote}
% \begin{macro}{\tablenoteskip}
% Old user commands |\tablenotemark{MARK}|,
% |\tablenote{MARK}{TEXT}|---now deprecated. 
%    \begin{macrocode}
\newcommand\tablenotemark[1]{$^\mathrm{#1}$} 
\newcommand\tablenote[2]{\if@proceedings\else\small\fi ${}^\mathrm{#1}$ #2} 
\newcommand\tablenoteskip{\\[0.5\baselineskip]} 
%    \end{macrocode}
% \end{macro}
% \end{macro}
% \end{macro}

% \begin{macro}{\tabnotemark}
% \begin{macro}{\tabnotetext}
% \begin{macro}{\tabnotetextTWO}
% Alternative way of doing footnotes in tables. 
%    \begin{macrocode}
\newcommand{\tabnotemark}[1]{\textsuperscript{\lowercase{#1}}}
%\newcommand{\tabnotetextTWO}[4]{\tabnotetext{#1}{#2}\tabnotetext{#3}{#4}}
\newcommand{\tabnotetextTWO}[4]{%
  \multicolumn{\@tablecols}{l}{%
    \parbox{0.5\tabnotewidth}{\small\textsuperscript{#1}#2}%
    \parbox{0.5\tabnotewidth}{\small\textsuperscript{#3}#4}%
    }%
  \\}
\newcommand{\tabnotetext}[2]{%
  \multicolumn{\@tablecols}{@{}p{\tabnotewidth}@{}}{\small\textsuperscript{#1}#2}\\}
\newlength{\tabnotewidth}
%    \end{macrocode}
% \end{macro}
% \end{macro}
% \end{macro}


% \begin{macro}{changemargin}
%    \begin{macrocode}
% Command to allow margins to be changed on the fly. Useful to allow
% tables and Figures to protrude into the margins. 
%
% \changes{3.24}{2007/09/14}{Added \texttt{changemargin} environment
% (copied from TeX FAQ).}
\newenvironment{changemargin}[2]{% 
  % From TeX FAQ http://www.tex.ac.uk/cgi-bin/texfaq2html?label=chngmargonfly
  \begin{list}{}{%
    \setlength{\topsep}{0pt}%
    \setlength{\leftmargin}{#1}%
    \setlength{\rightmargin}{#2}%
    \setlength{\listparindent}{\parindent}%
    \setlength{\itemindent}{\parindent}%
    \setlength{\parsep}{\parskip}%
  }%
  \item[]}{\end{list}}
%    \end{macrocode}
% \end{macro}

% \subsection{The bibliography and references}
% \label{sec:biblio}
% \begin{macro}{\beforebiblioskip}
% \begin{macro}{\afterbiblioskip}
% I'm not entirely sure why I introduced these two commands, since
% directly putting |\vspace| before and after
% |\begin{thebibliography}| seems to work well enough. 
%    \begin{macrocode}
\newlength{\beforebiblioskip}
\setlength{\beforebiblioskip}{0pt}
\newlength{\afterbiblioskip}
\setlength{\afterbiblioskip}{0pt}
%    \end{macrocode}
% \end{macro}
% \end{macro}
% \begin{macro}{thebibliography}
% The environment itself is done as starred section with extra
% vertical space after heading. This space is currently zero. 
% In the book style we do the cunning trick with rereading the file in
% order to get the references into the |\ADS@references| macro.
% 
% \changes{3.05}{2003/01/02}{Removed the \cs{beforebiblioskip} and
% \cs{afterbiblioskip} commands, since they cause unwanted vertical
% space in some instances.}  
%    \begin{macrocode}
\renewenvironment{thebibliography}{%
  \if@book\RereadBiblio\fi
%%  \vspace*{\beforebiblioskip}
  \section*{REFERENCES} \frenchspacing
%%  \vspace*{\afterbiblioskip}
  \small
%    \end{macrocode} 
% The references themselves are set as a list in order 
% to get the hanging indentation. They are done in a smaller font and
% with french spacing so we needn't worry about the distinction
% between sentence-ending periods and abbreviations. 
%    \begin{macrocode} 
  \list{\null}{\leftmargin 1.5em\labelwidth\z@\labelsep\z@\itemindent -1.5em
    \itemsep\z@\parsep\z@\usecounter{enumi}}
  \sloppy\clubpenalty4000\widowpenalty4000
  \sfcode`\.=1000\relax\@biblisttrue}
{\endlist\@biblistfalse\normalsize\@outputaddressesifnecessary}
%    \end{macrocode}
% \end{macro}
% \subsection{Uniquification of cross-references}
% \label{sec:uniquification}
% 
% \begin{macro}{\articlename}
% In order that the |\label{}| command in individual articles be
% guaranteed to give a unique anchor, we append the article name (held
% in the |\articlename| command) to each key. This requires
% redefinition of the following commands: |\cite|, |\nocite|, |\ref|,
% |\pageref|, |\label|. 
% \begin{macro}{\oldcite}
% \begin{macro}{\oldnocite}
% \begin{macro}{\oldref}
% \begin{macro}{\oldpageref}
% \begin{macro}{\oldlabel}
% The original versions of these commands are
% saved as |\oldcite|, |\oldnocite|, etc. 
%
% The uniquifying of cross-references is now always done, so 
% we don't get multiply-defined labels in the single article case
%    \begin{macrocode}
\newcommand{\articlename}{anon}
\AtBeginDocument{
  \let\oldcite=\cite 
  \let\oldnocite=\nocite 
  \let\oldref=\ref 
  \let\oldpageref=\pageref
  \let\oldlabel=\label
}
%    \end{macrocode}
% \end{macro}
% \end{macro}
% \end{macro}
% \end{macro}
% \end{macro}
% \end{macro}
%
% Uniquify all cross-references by prepending the article name.  This
% now has to go into |\AtBeginDocument| so it is working on the new
% versions that hyperref.sty has set up.  Then we need to put the
% |\@begofarticlehook| afterwards, so that the first page label is using
% the new definitions.  If amsmath is loaded, then we also need to patch
% the special label command that they use. \changes{4.0}{2019/11/18}{Fix labels for amsmath display equations, in case that package has been loaded.}
%    \begin{macrocode}
\AtBeginDocument{
  \renewcommand{\ref}[1]{\oldref{\articlename:#1}%
    \if@debug{\ClassInfo{rmaa}{Reference to \articlename:#1}}\fi}
  \renewcommand{\pageref}[1]{\oldpageref{\articlename:#1}%
    \if@debug{\ClassInfo{rmaa}{Page reference to \articlename:#1}}\fi}
  \renewcommand{\label}[1]{\oldlabel{\articlename:#1}%
    \if@debug{\ClassInfo{rmaa}{Label for \articlename:#1}}\fi}%
  \ifdef{\label@in@display}{%
    \let\oldlabel@in@display=\label@in@display
    \renewcommand{\label@in@display}[1]{\oldlabel@in@display{\articlename:#1}%
      \if@debug{\ClassInfo{rmaa}{AMSTeX Label for \articlename:#1}}\fi}%
  }{}
  \@begofarticlehook
}%
%    \end{macrocode}
% Citations are more complicated since they can contain a list of
% keys. These are really only here for use with the rmaacite package,
% which is utterly deprecated (02 Jan 2003) now that I have the natbib
% code working. The command \cs{cite} should \emph{not} be used with
% natbib: use \cs{citep}, \cs{citet}, etc.\@ instead. 
%    \begin{macrocode}
\renewcommand{\nocite}[1]{\@bsphack
  \@for\@citeb:=#1\do{%
    \edef\@citeb{\expandafter\@firstofone\@citeb:\articlename}%
    \if@filesw\immediate\write\@auxout{\string\citation{\@citeb}}\fi
    \@ifundefined{b@\@citeb:\articlename}{\G@refundefinedtrue
      \@latex@warning{Citation `\@citeb' undefined}}{}}%
  \@esphack}
\renewcommand{\cite}[1]{\oldcite{#1:\articlename}}
%   \end{macrocode}


% Horrible fiddling with the bibliography citation mechanism. 
% This has to be globally available in case rmaacite puts it in the
% \texttt{.aux} file. 
%    \begin{macrocode}
\def\lbibcite#1#2#3{\global\@namedef{b@#1}{#2}
                    \global\@namedef{by@#1}{#3}}
%    \end{macrocode}

% \begin{macro}{\oldbibitem}
% \begin{macro}{\@rmaaBibArg}
% \begin{macro}{\@rmaaSpace}
% \begin{macro}{rmaabib}
% Of course, |\bibitem| needs to get the uniquification treatment as
% well. While we are at it, we redefine it to use consecutive integers
% for the key if the user leaves it blank (using the counter
% |rmaabib|). This allows one to put 
% |\bibitem{}| for all your references without getting annoying
% messages about multiply-defined labels.
% 
% We also take the opportunity to insert the optional argument to
% \cs{bibitem} if it is missing in the form Author(DATE) so that natbib
% doesn't try to do numerical citations. Natbib already provides a
% hook \cs{@extra@b@citeb} to help uniquify the citation key, so we
% just plug into that. However, it puts it \emph{after} the original
% key, not before as I had had it. Easiest solution is to change the
% way that I do it. 02 Jan 2003. Note that the \cs{@extra@b@citeb}
% does not get written to the .aux  file in the \cs{citation} command
% but it does work. 
% \changes{3.14}{06 Feb 2003}{Reinstalled old definition of \cs{bibitem}
%  for the case where the nonatbib option is requested, so that the
%  rmaacite package will work. This means that you can't use rmaacite
%  and natbib in the same volume.}
%    \begin{macrocode}
\if@natbib\def\@extra@b@citeb{:\articlename}\fi
%%\else
\let\oldbibitem\bibitem
\newcommand{\@rmaaBibArg}{}
\newcommand{\@rmaaSpace}{ }
\newcounter{rmaabib}
\if@natbib
\renewcommand{\bibitem}[2][ABC(2003)]{%
  \renewcommand{\@rmaaBibArg}{#2}
  \ifx\@rmaaBibArg\@rmaaEmpty
     \addtocounter{rmaabib}{1}
     \oldbibitem[#1]{\thermaabib:\articlename}
  \else\ifx\@rmaaBibArg\@rmaaSpace
     \addtocounter{rmaabib}{1}
     \oldbibitem[#1]{\thermaabib:\articlename}
  \else 
     \oldbibitem[#1]{#2:\articlename}
  \fi\fi
  \renewcommand{\@rmaaBibArg}{}
  }
\else
\renewcommand{\bibitem}[1]{%
  \renewcommand{\@rmaaBibArg}{#1}
  \ifx\@rmaaBibArg\@rmaaEmpty
     \addtocounter{rmaabib}{1}
     \oldbibitem{\articlename:\thermaabib}
  \else\ifx\@rmaaBibArg\@rmaaSpace
     \addtocounter{rmaabib}{1}
     \oldbibitem{\articlename:\thermaabib}
  \else
     \oldbibitem{\articlename:#1}
  \fi\fi
  \renewcommand{\@rmaaBibArg}{}
  }
\fi
%%\fi
%    \end{macrocode}
% \end{macro}
% \end{macro}
% \end{macro}
% \end{macro}

% \begin{macro}{oldpcite}
% \begin{macro}{oldscite}
% Then, we need to repeat the above scheme, but for the |rmaacite|
% definition. 
%    \begin{macrocode}
\newcommand{\@redefinebibitem}{%
  \let\oldbibitem\bibitem
  \let\oldcite\cite
  \let\oldpcite\pcite
  \let\oldscite\scite
  \def\bibitem[##1<##2>]##3{%
    \renewcommand{\@rmaaBibArg}{##3}
    \ifx\@rmaaBibArg\@rmaaEmpty
       \addtocounter{rmaabib}{1}
       \oldbibitem[##1<##2>]{\thermaabib:\articlename}
    \else\ifx\@rmaaBibArg\@rmaaSpace
       \addtocounter{rmaabib}{1}
       \oldbibitem[##1<##2>]{\thermaabib:\articlename}
    \else 
       \oldbibitem[##1<##2>]{##3:\articlename}
    % trust the author here
%%       \oldbibitem[##1<##2>]{##3}
    \fi\fi
    }
  }
%    \end{macrocode}
% \end{macro}
% \end{macro}


% \subsection{HTML Table of Contents}
% \label{sec:hyper-TOC}
% \begin{macro}{\HyperTOCEntry}
% \begin{macro}{\HyperTOCpart}
% \begin{macro}{\HyperTOCmeeting}
% \begin{macro}{\HyperTOCeditors}
% \begin{macro}{\HyperTOCpreamble}
% These commands are used to write information to the \texttt{.htoc}
% file, which is used in generating a Web Page of the Table of
% Contents. Since they only ever appear preceded by |\protect| in the
% argument to an |\addtocontents| command,  they should never need to
% be expanded by the macros. However, we define them here just in case
% \dots
%    \begin{macrocode}
\newcommand{\HyperTOCEntry}[3]{\relax}
\newcommand{\HyperTOCpart}[1]{\relax}
\newcommand{\HyperTOCmeeting}[3]{\relax}
\newcommand{\HyperTOCeditors}[1]{\relax}
\newcommand{\HyperTOCpreamble}{\relax}
%    \end{macrocode}
% \end{macro}
% \end{macro}
% \end{macro}
% \end{macro}
% \end{macro}
%

% \subsection{Formatting the poster abstracts} 
% These can now be used without using the book option.
%    \begin{macrocode}

\newcommand{\@makeminititle}{%
  \newpage
  \null 
  \vskip 3em%
  \begin{center}%
    {\@title \par}%
  \end{center}
  \vskip 3em%
  }

\newif\if@abstractsads
\@abstractsadsfalse
\newcommand\abstractsADS{\@abstractsadstrue}

\if@book\RequirePackage{multicol}\fi
\newenvironment{abstracts}[1][Abstracts of Contributed Papers]{%
  %% First, make sure we start on an odd page (06 Jun 2000)
  \if@book\clearpage\thispagestyle{empty}\cleardoublepage\fi
  %%
  \renewcommand{\articlename}{ABSTRACTS}
  \title{#1}\shortauthor{Abstracts}\shorttitle{Abstracts}
  \listofauthors{}
  \if@book
  \addcontentsline{toc}{chapter}{\protect\@fiddleauthors{\@TOCtitle}{\rmaaTOCauthors}}%
  \addtocontents{htc}{\protect\HyperTOCEntry{\@rawtitle}{\@rawauthors}{\articlename}}
  \fi
  \pagestyle{rmaa} % this now seems necessary (06 Jun 2000)
  \thispagestyle{rmaafirstpage}
%%  \twocolumn[\@makeminititle]
  \renewcommand{\abstract}[1]{%
    \renewcommand\ADS@abstract{##1}%
    \renewcommand{\@abstract}{##1}\@writeabstract%
  }
  \renewcommand{\title}[2][]{%
    \gdef\ADS@authors{}%
    \gdef\ADS@affil{}%
    \gdef\ADS@references{}%
    \gdef\ADS@keywords{}%
    \renewcommand{\@rmaaArg}{##1}
    \ifx\@rmaaArg\@rmaaEmpty
       \renewcommand{\@TOCaftertitle}{}
    \else
       \renewcommand{\@TOCaftertitle}{ (\@rmaaArg)}
    \fi
    \renewcommand{\@title}{\rmaa@TOCTitleStyle{##2}}%
    \renewcommand{\@TOCtitle}{\rmaa@TOCTitleStyle{##2}\@TOCaftertitle}%
    \renewcommand{\rmaaIndexTitle}{\rmaa@IndexTitleStyle{##2\@TOCaftertitle.}}%
    \renewcommand\ADS@title{##2}%
    \xdef\ADS@firstpage{\@arabic{\c@page}}%
    }
  \renewcommand{\author}[1]{\let\@thanks\@empty\renewcommand\@author{##1} }
  \renewcommand{\footnotetext}[2][XX]{\par\noindent\textsuperscript{##1}~##2}
  \renewcommand{\footnoterule}{\kern-3\p@ \hrule \@width 3em \kern 2.6\p@}
  \if@book\@begofarticlehook\fi
  \if@book\begin{multicols}{2}[\@makeminititle]\else\@makeminititle\fi
  \begingroup
  \setcounter{footnote}{0}
  }% and at the end of the environment....
  {\endgroup\if@book\end{multicols}\listofabstracts\label{RMAALastPage}\clearpage\fi}
%% auxiliary commands used:
%% \@writeabstract does all the actual formatting 
\newcommand{\@writeabstract}{%
  \if@book%
  \addcontentsline{loa}{abstract}{\protect\@fiddleauthors{\@TOCtitle}{\rmaaTOCauthors}}%
  \fi%
  \begin{center}\@title\par\@author\end{center}
  \@abstract
  \par\noindent\footnoterule\nobreak\footnotesize\@thanks
  \par\noindent\hrulefill\goodbreak
  \vskip 1cm \@plus0.5cm
  %% write ADS entry
  \if@book\if@abstractsads\ADS@dumpdata\fi\fi%
  \endgroup % this pair is to make each abstract a self-contained group
  \begingroup
  \setcounter{footnote}{0}
  }

\if@book
%% List of abstracts (stuff written to .loa file)
\newcommand{\listofabstracts}{%
  \begin{multicols}{2}[%
    \vskip 1.5cm \vfill % leave at least 1.5 cm and go to bottom of page if possible
    \centerline{\MakeUpperCase{List of Abstracts}}
    \vskip 1.5cm \@plus0.5cm
    ][6cm] % this makes it start a new page if there is not at least 6cm of space
    \renewcommand{\@fiddleauthors}[2]{##1\hskip1em plus2em minus0.4em\penalty-100  ##2}
    \parindent\z@
%%    \parskip.5\p@ \@plus .7\p@\relax 
    \parskip3\p@ \@plus 1\p@ \@minus 1\p@\relax % allow quite a lot of
                                % leeway here to avoid breaks inside
                                % entries 
    \renewcommand{\\}{\relax} % temporarily turn off explicit line breaks just in case
    \renewcommand{\altaffilmark}[1]{\relax} % ditto with this
    \@starttoc{loa}
  \end{multicols}
  }
\newcommand*{\l@abstract}[2]{%
  \par\penalty-1000
  #1\hskip 1em\penalty0\null\nobreak\hfill#2}
\fi

%    \end{macrocode}
 
% \subsection{Code for typesetting an entire volume of conference
% proceedings} \label{sec:whole-volume}
%    \begin{macrocode}

\if@book
%% environment for including the individual articles
\newcommand{\@fiddleauthors}[2]{\relax}
\renewcommand{\theequation}{\@arabic\c@equation}
\newenvironment{articles}%
{ \begingroup % make it a group so that author definitions do not
              % spill over
  \gdef\ADS@authors{}
  \pagestyle{rmaa}
  \renewenvironment{document}{%
    \@begofarticlehook
    \addtocounter{chapter}{1}%
    \addcontentsline{toc}{chapter}{%
      \protect\@fiddleauthors{\@TOCtitle}{\rmaaTOCauthors}}%
    \addtocontents{htc}{\protect\HyperTOCEntry{\@rawtitle}{\@rawauthors}{\articlename}}
    }{\relax}
  \renewcommand{\documentclass}[2][foo]{\relax} 
  \renewcommand{\usepackage}[2][foo]{%
    \def\@packagename{##2}\def\@rmaacitename{rmaacite}%
    \ifx\@packagename\@rmaacitename
    \@rmcitetrue\ClassInfo{rmaa}{This paper uses rmaacite}
    \makeatletter%%
%% This is file `rmaacite.sty',
%% generated with the docstrip utility.
%%
%% The original source files were:
%%
%% rmaa.dtx  (with options: `rmaacite')
%% 
%% PLEASE DO NOT EDIT THIS FILE! It is generated automatically from
%% the file rmaa.dtx using the accompanying Makefile and rmaa.ins
%% Any hand-edited changes will be lost
%% 
%% A LaTeX document class for the RMxAA
%% 
%% Author: William Henney (w.henney@crya.unam.mx)
%% 
%% Latest version available from: http://www.crya.unam.mx/rmaa
%% 
%% Copyright (C) 1999-2010 Instituto de Astronomia, UNAM, Mexico
%% 
%% This program is free software; you can redistribute it and/or
%% modify it under the terms of the GNU General Public License
%% as published by the Free Software Foundation; either version 2
%% of the License, or (at your option) any later version.
%% 
%% This program is distributed in the hope that it will be useful,
%% but WITHOUT ANY WARRANTY; without even the implied warranty of
%% MERCHANTABILITY or FITNESS FOR A PARTICULAR PURPOSE.  See the
%% GNU General Public License for more details.
%% 
%% You should have received a copy of the GNU General Public License
%% along with this program; if not, write to the Free Software
%% Foundation, Inc., 59 Temple Place - Suite 330, Boston, MA  02111-1307, USA.






\def\bibitem[#1<#2>]#3{\par\hangindent=1.5em\hangafter=1
  \if@filesw
  {\def\protect##1{\string ##1\space}
    \immediate\write\@auxout{\string\lbibcite{#3}{#1}{#2}}}
  \fi}
\@redefinebibitem % what it says
\def\lbibcite#1#2#3{\global\@namedef{b@#1}{#2}
                    \global\@namedef{by@#1}{#3}}

\def\@bcite#1#2{({#1\if@tempswa , #2\fi})}
\def\@pcite#1#2{#1\if@tempswa , #2\fi}
\def\@brcite{\ (}
\def\@ebrcite{)}
\def\@fmttop#1{\@fmtcite[\csname #1\endcsname]}
\def\@citex[#1]#2{%
%%  \if@filesw\immediate\write\@auxout{\string\citation{#2}}\fi
%% WJH added 16 Jan 2002
  \@for\@citeb:=#2\do{%
    \edef\@citeb{\expandafter\@firstofone\@citeb:\articlename}%
    \if@filesw\immediate\write\@auxout{\string\citation{\@citeb}}\fi}%
%% to here
  \def\@citea{}%
  \@cite{\@for\@citeb:=#2\do
    {%
      \edef\@citeb{\expandafter\@firstofone\@citeb:\articlename}%
      \@citea\def\@citea{;\ }\@ifundefined
      {b@\@citeb}{{\bf ?}\@warning {Citation `\@citeb' on page \thepage \space undefined}}%
      {\csname b@\@citeb\endcsname\@citebro\csname by@\@citeb\endcsname\@citebrc}}}%
  {#1}}

\def\cite{\@ifnextchar [{\let\@citebro=\space
                          \let\@citebrc=\relax
                          \let\@cite=\@bcite\@tempswatrue\@citex}
                        {\let\@citebro=\space
                          \let\@citebrc=\relax
                          \let\@cite=\@bcite\@tempswafalse\@citex[]}}
\def\pcite{\@ifnextchar [{\let\@citebro=\space
                          \let\@citebrc=\relax
                          \let\@cite=\@pcite\@tempswatrue\@citex}
                        {\let\@citebro=\space
                          \let\@citebrc=\relax
                          \let\@cite=\@pcite\@tempswafalse\@citex[]}}
\def\scite{\@ifnextchar [{\let\@citebro=\@brcite
                          \let\@citebrc=\@ebrcite
                          \let\@cite=\@pcite\@tempswatrue\@citex}
                        {\let\@citebro=\@brcite
                          \let\@citebrc=\@ebrcite
                          \let\@cite=\@pcite\@tempswafalse\@citex[]}}

\renewenvironment{thebibliography}
  {\section*{REFERENCES}\vspace*{\afterbiblioskip}\small
    \if@book\RereadBiblio\fi
    \sloppy\clubpenalty4000\widowpenalty4000
    \sfcode`\.=1000\relax
    \setlength\parindent{\z@}\setlength\parskip{\z@}}
  {\relax\normalsize\@outputaddressesifnecessary}








\endinput
%%
%% End of file `rmaacite.sty'.
\makeatother % read it in here so that all changes
    \else                                         % are local to paper
    \if@debug\ClassWarning{rmaa}{Package \@packagename\space was requested by \articlename\space}\fi
    \fi} 
  %% Don't allow individual articles to mess with the pages, volume or year
  \renewcommand{\SetFirstPage}[1]{\relax}
  \renewcommand{\SetYear}[1]{\relax}
  \renewcommand{\SetVolume}[1]{\relax}
  \renewcommand{\SetConfTitle}[1]{\relax}
  %% Don't allow individual articles to change these
  \renewcommand{\nonstopmode}{\relax}
  \renewcommand{\scrollmode}{\relax}
  \renewcommand{\outputplates}{\relax}
  %% XXXX need to check if using rmaacite and act accordingly
  %% avoid multiple declarations of the same length
  \let\oldnewlength=\newlength
  \renewcommand{\newlength}[1]{%
    \edef\reserved@a{\expandafter\@gobble\string ##1}%
    \@ifundefined{\reserved@a}{%
      \oldnewlength{##1}}{%
      \if@debug{\ClassInfo{rmaa}{\reserved@a\space already defined: doing nothing}}\fi%
      \relax}}%
  \renewcommand\newcommand{\providecommand}
  }%
{\endgroup\onecolumn}
     

%    \end{macrocode}
% Command to be used in `articles' environment for including an
% article
%    \begin{macrocode}
\newcommand{\article}[2][]{%
  % optional first argument is extra text to put in brackets after title in TOC
  % e.g. ``Invited Talk'' or ``Poster''
  \renewcommand{\articlename}{#2}
  % reset loads of counters
  \setcounter{footnote}{0}
  \setcounter{figure}{0}
  \setcounter{table}{0}
  \setcounter{section}{0}
  \setcounter{equation}{0}
  \setcounter{rmaabib}{0}
  \setcounter{author}{0}
  \setcounter{AUTHOR}{1}
  \gdef\ADS@authors{}%
  \gdef\ADS@affil{}%
  \gdef\ADS@affiliations{}%
  \gdef\ADS@references{}%
  \@firstauthortrue
  \@firstaffiltrue
  \@firstaltaffiltrue
  \@firstkwdtrue
  \@rmcitefalse
  \if@suppressfulladdresses\else\@needaddressestrue\fi
  \renewcommand{\@rmaaArg}{#1}
  \ifx\@rmaaArg\@rmaaEmpty
     \renewcommand{\@TOCaftertitle}{}
  \else
     \renewcommand{\@TOCaftertitle}{ (#1)}
  \fi
%    \end{macrocode}
% In the main journal articles must start on an odd page. 
%    \begin{macrocode}
  \if@journal\clearpage\thispagestyle{empty}\cleardoublepage\fi
%% Now input the file, then run the end-of-article hook (reinstated 28
  %% Sep 2001)
  {\input{\@articledir/\articlename.tex}\@endofarticlehook}}
%    \end{macrocode}
% Command to separate articles into sessions. Write an entry in the
% TOC, makes sure we start on an odd page (added 06 Jun 2000), then do
% a title page if \cs{SessionTitlePages} was called. 
% \changes{3.12}{04 Feb 2003}{Added option to have title page for each
% session. Turn it on with \cs{SessionTitlePages}.}
%    \begin{macrocode}
\newcommand{\session}[1]{
  %% this does nothing except write an entry in the TOC
  \addcontentsline{toc}{part}{\rmaa@TOCPartStyle{#1}}
  \addtocontents{htc}{\protect\HyperTOCpart{#1}}
  %% and also makes sure we start on an odd page (added 06 Jun 2000)
  \clearpage\thispagestyle{empty}\cleardoublepage
  \stepcounter{part}
  \if@sessiontitlepages{%
    \onecolumn
    \thispagestyle{empty}
    \renewcommand{\articlename}{Section\thepart TitlePage}
    \label{RMAAFirstPage}
    \make@session@title{#1}
    \label{RMAALastPage}
    \clearpage\thispagestyle{empty}\cleardoublepage
  }\fi
  }
\newcommand{\make@session@title}[1]{%
  {%
    \setlength\parindent{0pt}
    \vspace*{\stretch{1}}
    \rule{\linewidth}{1mm}
    \begin{flushright}\Huge\bfseries\boldmath
      #1
    \end{flushright}
    \rule{\linewidth}{1mm}
    \vspace*{\stretch{2}}
  }
}
%    \end{macrocode}
% Command to turn on separate title pages for each session. 
%    \begin{macrocode}
\newcommand{\SessionTitlePages}{\@sessiontitlepagestrue}
\newcommand{\BoldSessionTitles}{%
  \renewcommand{\rmaa@TOCPartStyle}[1]{\textbf{\rmaa@TOCTitleStyle{##1}}}
  }
\newcommand{\ArticleDir}[1]{\renewcommand{\@articledir}{#1/}\graphicspath{{1/}}}
\newcommand{\@articledir}{.}
%\newcommand{\ps@rmaafront}{\ps@rmaa}% page style for front matter (same as ever, actually)
\newcommand{\ps@rmaafront}{% page style for front matter:  page # in each foot
  \renewcommand\@oddfoot{\hfil{\thepage}}%
  \renewcommand\@evenfoot{{\thepage}\hfil}%
  % shorttitle in both heads
  \renewcommand\@evenhead{\hfil\@shorttitle\hfil}%
  \renewcommand\@oddhead{\@evenhead}
  }
\newcommand{\ps@rmaafrontfirstpage}{% page style for first page of front matter
  % page # in each foot, on outside edge
  \renewcommand\@oddfoot{\hfil\thepage}%
  \renewcommand\@evenfoot{\thepage\hfil}%
  % nothing in the even heads
  \renewcommand\@evenhead{\relax}%
  \renewcommand\@oddhead{\relax}%
  }
%% slight modification to the standard \frontmatter and \mainmatter
%% commands, just to set the \pagestyle
\renewcommand\frontmatter{%
  \latexonecolumn
  \cleardoublepage
  \@frontmattertrue
  \@mainmatterfalse
  \pagenumbering{roman}
  \pagestyle{rmaafront}
  \addtocounter{page}{2} % leave space for the title page
  }
\renewcommand\mainmatter{%
  \pagestyle{empty}
  \if@rmtwocol\twocolumn\fi
  \cleardoublepage
  \@frontmatterfalse
  \@mainmattertrue
  \pagenumbering{arabic}
  \pagestyle{rmaa}
  }
\renewcommand\backmatter{%
  \clearpage\pagestyle{empty}\cleardoublepage
  \@mainmatterfalse
  %% put extra space in the TOC
  \addtocontents{toc}{\bigskip}
  }
%    \end{macrocode}
%
% \subsection{Front sections}\label{sec:front-sections}
% \begin{macro}{\notitle}
% \begin{macro}{\shortname}
% \begin{macro}{frontsection}
% The environment \texttt{frontsection} is for doing a section in the
% front matter (preface, list of participants, or whatever). It has 
% one optional argument:
% \begin{enumerate}
% \item[\#1] Short title (no embedded spaces) for
% use in making PS files. E.g., ``list\_of\_participants''.   
% \end{enumerate}
% and two mandatory arguments: 
% \begin{enumerate}
% \item[\#2] Title of section (e.g., `Preface')
% \item[\#3] Commands to be evaluated before processing the contents
% of the environment. This is mainly useful for putting in the things
% that would go in the preamble of a normal article  (e.g.,
% |\listofauthors|, |\indexauthor|, |\notitle|) 
% \end{enumerate}
%    \begin{macrocode}
\newcommand{\notitle}{\@titlevisiblefalse} \newcommand{\shortname}{}
\newenvironment{frontsection}[3][]{%
  \begingroup
  \renewcommand{\shortname}{#1}
  \ifx\shortname\@rmaaEmpty
     \renewcommand{\articlename}{#2}
  \else
     \renewcommand{\articlename}{\shortname}
  \fi
  \@begofarticlehook
  \@titlevisibletrue % default
  \thispagestyle{rmaafrontfirstpage}
  \listofauthors{} % blank out this in case it is not specified
  \title{{#2}}\shorttitle{#2}\shortauthor{#2}
  #3 % evaluate the command given as second mandatory arg
  \addcontentsline{toc}{chapter}{\protect\@fiddleauthors{\@TOCtitle}{\rmaaTOCauthors}}%
  \addtocontents{htc}{\protect\HyperTOCEntry{\@rawtitle}{\@rawauthors}{\articlename}}
  \if@titlevisible
     \vspace*{1cm}\centerline{\rmaa@FrontSecStyle{\@title}}\vspace*{1cm}
  \fi
  \setlength{\parskip}{3ex}
  }{\frontlabel{RMAALastPage}\clearpage\endgroup}
%    \end{macrocode}
% \end{macro}
% \end{macro}
% \end{macro}
%    \begin{macrocode}
%% Definition of the command \frontlabel, which rights out something like `F4' to
%% the .aux file, instead of `\textlatin  {i}'. This is so it is easy to post-process
%% the file for making PS files of individual sections
\newcommand{\frontlabel}[1]{%
  \@bsphack
  \protected@write\@auxout{}%
     {\string\newlabel{\articlename:#1}{{}{F\@arabic{\c@page}}}}
  \@esphack
  }
%% User command to say X, Y, Z,\\ Editors of the Proceedings
\newcommand{\EditorSignature}{\vspace*{\parskip}\rightline{\@editors}%
  \rightline{Editors of the proceedings}}
%% Fiddle with \thesection, etc. so that they don't include the chapter number
%% Basically go back to the `article' definitions of these
\renewcommand{\thesection}{\@arabic\c@section}
\renewcommand{\thefigure}{\@arabic\c@figure}
\renewcommand{\thetable}{\@arabic\c@table}

\newcommand{\@oldargs}{}
\newcommand{\@newargs}{}
%% add TOC entry for each paper
%\renewcommand{\title}[1]%
%{\renewcommand{\@title}{\rmaa@TitleStyle{#1}}%
%\addtocounter{chapter}{1}\addcontentsline{toc}{chapter}{#1}%
%}

%% Formatting of the table of contents
%% dimensions of various parts
\newlength{\TOCtextwidth}\setlength{\TOCtextwidth}{12.65cm}
\newlength{\TOCpagenumwidth}\setlength{\TOCpagenumwidth}{1.3cm}
\newlength{\TOCentrywidth}\setlength{\TOCentrywidth}{\TOCtextwidth}
\addtolength{\TOCentrywidth}{-\TOCpagenumwidth} 
\newlength{\TOCextramargin}\setlength{\TOCextramargin}{1.25cm}
\newcommand{\TOCbox}[1]{\hspace{\TOCextramargin}\parbox{\TOCtextwidth}{#1}\par}
\newlength{\TOCextramarginx}\setlength{\TOCextramarginx}{\TOCextramargin}\addtolength{\TOCextramarginx}{4.3cm}
\newlength{\TOCtextwidthx}\setlength{\TOCtextwidthx}{\TOCtextwidth}\addtolength{\TOCtextwidthx}{-4.3cm}
\newcommand{\TOCboxx}[1]{\hspace{\TOCextramarginx}\parbox{\TOCtextwidthx}{#1}\par}

\newcommand{\RMAAbannerimage}{rmsc}
\renewcommand\tableofcontents{%
  \renewcommand{\articlename}{CONTENTS}
  \begingroup
  \@begofarticlehook
  \if@twocolumn
     \@restonecoltrue\onecolumn
  \else
     \@restonecolfalse
  \fi
  %% 9pt for journal style
  \if@proceedings\else\small\fi
  \if@newstyle\small\fi %% same for new style
  \if@proceedings\thispagestyle{rmaafrontfirstpage}\else\thispagestyle{empty}\fi
  \shortauthor{Contents}\shorttitle{Contents}
  \vspace*{18mm}
  \if@printbanner
  %% put in the banner header (need to get better image of this)
     \TOCbox{\hspace*{-0.15cm}\includegraphics[width=0.9\TOCtextwidth]{\RMAAbannerimage}}
     \vspace*{8mm}
  %% volume and publication date (should really be in funny font)
     \TOCboxx{\@volumen\\[2mm] \@mes\ \@anyo}
  \else
  %% put in the banner header (need to get better image of this)
     \TOCbox{\hspace*{-0.15cm}\includegraphics[width=0.9\TOCtextwidth]{\RMAAbannerimage}}
     \vspace*{12.5mm}
  %% volume and publication date (should really be in funny font)
     \TOCboxx{~\\[2mm] ~}
  \fi
  \vspace*{16mm}
  %% This is the command that fiddles with the TOC entry so that the
  %% title is flushleft but the authors are flushright
  \renewcommand{\@fiddleauthors}[2]{%
    % first put the paper title
    \renewcommand{\\}{\relax} % temporarily turn off explicit line breaks just in case
    \renewcommand{\altaffilmark}[1]{\relax} % ditto with this
    \renewcommand{\textsuperscript}[1]{\relax} % ditto with this 
    ##1
    % now save the author list in a box
    \sbox\@tempboxa{##2}%
    % check that box is shorter than the total width available for the entry
    \ifdim \wd\@tempboxa <\TOCentrywidth
       % if so, then put it right here, flush against right margin
       \penalty10000\hskip1em plus10em\penalty0\hbox{}\nobreak\hfill\box\@tempboxa
       % Note that we force a space of at least 1em between title and authorlist, 
       % stretchable to 11em since we may have very short titles but still with an 
       % authorlist too long to go on the same line.  The only linebreak allowed is 
       % directly before the authorlist.
    \else
       % if it is too long, we set the original authorlist text in its own parbox
       \par\vskip1ex\parbox[b]{\TOCentrywidth}{\raggedleft ##2}
    \fi
    }
  \@starttoc{toc}%
  \if@restonecol\twocolumn\fi
  \frontlabel{RMAALastPage}
  \clearpage
  \thispagestyle{empty}
  \cleardoublepage % make sure we start on an odd page
  \endgroup

  %%% ALSO, do the HTML TOC
  %% The file that will be used to make the HTML TOC
  \@starttoc{htc}
  %% All the stuff in this file does nothing in the book style
  \addtocontents{htc}{\protect\HyperTOCmeeting{\@conftitle}{\@reunionplace}{\@reuniondate}}
  \addtocontents{htc}{\protect\HyperTOCeditors{\@editors}}
  \addtocontents{htc}{\protect\HyperTOCpreamble}
  \addtocontents{htc}{\protect\HyperTOCEntry{Table of Contents}{}{\articlename}}

  %% and the ADS file
  \ADS@openfile
  }
%% This is now much simpler
\renewcommand\l@chapter[2]{%TOC entry for articles
  \ifnum \c@tocdepth >\m@ne
    \vskip 1.0em \@plus\p@
    \TOCbox{%
      \parbox[b]{\TOCentrywidth}{#1}\parbox[b]{\TOCpagenumwidth}{\hfill#2}
    }
  \fi} 
\renewcommand\l@part[2]{%TOC entry for sessions: centered, no page number
  \ifnum \c@tocdepth >\m@ne
    \vskip 1.5em \@plus1.5em \goodbreak
    \TOCbox{%
      \parbox[b]{\TOCentrywidth}{\hfil#1\hfil}\hfill
    }
  \fi\nobreak} 

%% Formatting theindex
%% strongly encourage page breaks between index entries
\renewcommand{\@idxitem}{\par\penalty-1000\hangindent.45cm} 
\renewcommand{\indexspace}{\par\penalty-1000\vskip 10\p@ \@plus15\p@ \@minus3\p@\relax}
\renewenvironment{theindex}{%
  %% first, make sure we start on an odd page (06 Jun 2000)
  \clearpage\thispagestyle{empty}\cleardoublepage
  %% this looks a bit like an article but is two columns and without abstract, etc
  \renewcommand{\articlename}{INDEX}
  \title{Author Index}\shortauthor{Author Index}\shorttitle{Author Index}
  \listofauthors{}
  \addcontentsline{toc}{chapter}{\protect\@fiddleauthors{\@TOCtitle}{\rmaaTOCauthors}}%
  \addtocontents{htc}{\protect\HyperTOCEntry{\@rawtitle}{\@rawauthors}{\articlename}}
  \pagestyle{rmaa}
  \thispagestyle{rmaafirstpage}
  \@begofarticlehook
  %% 9pt for journal style
  \if@proceedings\else\small\fi
  \if@newstyle\small\fi
  \twocolumn[\@makeminititle]
  %% vertical spacing (see also definition of \@idxitem and \indexspace
  \parskip\z@ \@plus .6\p@\relax % allow quite a lot of leaway here to avoid breaks inside entries
  %% horizontal spacing, breaking and glue parameters (see also the stuff in the definition 
  %% of \indexauthor
  \parindent\z@ 
  \exhyphenpenalty10000 % don't break on explicit hyphens, it always looks bad
  \rightskip=0pt plus.2em % allow a tiny bit of leeway at the ends of lines
  \spaceskip.3333em plus.5em minus.15em % a lot of leeway between words
  \xspaceskip.5em plus.5em minus.4em % this seems to have no effect
  \let\item\@idxitem
  \renewcommand{\\}{\relax} % temporarily turn off explicit line breaks just in case
  \renewcommand{\altaffilmark}[1]{\relax} % ditto with this
  \renewcommand{\textsuperscript}[1]{\relax} % ditto with this
  }{\label{RMAALastPage}\clearpage} % clearpage necessary so that we know about \shortauthor



%% end definitions for book styles
\fi

%    \end{macrocode}
% In order to do away with plates, we sometimes want a figure to
% appear in color only in the online version. The following command
% can bve used in a figure caption to typeset a note to this effect. 
%    \begin{macrocode}
\newcommand\ColorOnlineOnly{%
  \textsc{Note: this figure is available in color in the electronic
    version of this article, obtainable from}
  \texttt{http://\discretionary{}{}{}www.\discretionary{}{}{}astroscu.\discretionary{}{}{}unam.\discretionary{}{}{}mx/\discretionary{}{}{}\~{}rmaa/}.}
%% Plates
%% Define 2 counters: plate is the number in the whole document, plateloop is used 
%% internally by \outputplates
\newcounter{plate}
\newcounter{plateloop}
%% Currently limited to 15 plates in document
%% when we do a label, we need to write to \newlabel entries to the .aux file, one 
%% that says RMAAPlateX so that it can be referred to in \ouputplates, and one that
%% says articlename:RMAAPlateX so that the dvips program can find it and asociate it
%% with the right article
\newcommand{\@platelabeli}{\oldlabel{RMAAPlate1}\label{RMAAPlate1}}
\newcommand{\@platelabelii}{\oldlabel{RMAAPlate2}\label{RMAAPlate2}}
\newcommand{\@platelabeliii}{\oldlabel{RMAAPlate3}\label{RMAAPlate3}}
\newcommand{\@platelabeliv}{\oldlabel{RMAAPlate4}\label{RMAAPlate4}}
\newcommand{\@platelabelv}{\oldlabel{RMAAPlate5}\label{RMAAPlate5}}
\newcommand{\@platelabelvi}{\oldlabel{RMAAPlate6}\label{RMAAPlate6}}
\newcommand{\@platelabelvii}{\oldlabel{RMAAPlate7}\label{RMAAPlate7}}
\newcommand{\@platelabelviii}{\oldlabel{RMAAPlate8}\label{RMAAPlate8}}
\newcommand{\@platelabelix}{\oldlabel{RMAAPlate9}\label{RMAAPlate9}}
\newcommand{\@platelabelx}{\oldlabel{RMAAPlate10}\label{RMAAPlate10}}
\newcommand{\@platelabelxi}{\oldlabel{RMAAPlate11}\label{RMAAPlate11}}
\newcommand{\@platelabelxii}{\oldlabel{RMAAPlate12}\label{RMAAPlate12}}
\newcommand{\@platelabelxiii}{\oldlabel{RMAAPlate13}\label{RMAAPlate13}}
\newcommand{\@platelabelxiv}{\oldlabel{RMAAPlate14}\label{RMAAPlate14}}
\newcommand{\@platelabelxv}{\oldlabel{RMAAPlate15}\label{RMAAPlate15}}
\newcommand{\@outputplatelabeli}{\oldlabel{RMAAOutputPlate1}}
\newcommand{\@outputplatelabelii}{\oldlabel{RMAAOutputPlate2}}
\newcommand{\@outputplatelabeliii}{\oldlabel{RMAAOutputPlate3}}
\newcommand{\@outputplatelabeliv}{\oldlabel{RMAAOutputPlate4}}
\newcommand{\@outputplatelabelv}{\oldlabel{RMAAOutputPlate5}}
\newcommand{\@outputplatelabelvi}{\oldlabel{RMAAOutputPlate6}}
\newcommand{\@outputplatelabelvii}{\oldlabel{RMAAOutputPlate7}}
\newcommand{\@outputplatelabelviii}{\oldlabel{RMAAOutputPlate8}}
\newcommand{\@outputplatelabelix}{\oldlabel{RMAAOutputPlate9}}
\newcommand{\@outputplatelabelx}{\oldlabel{RMAAOutputPlate10}}
\newcommand{\@outputplatelabelxi}{\oldlabel{RMAAOutputPlate11}}
\newcommand{\@outputplatelabelxii}{\oldlabel{RMAAOutputPlate12}}
\newcommand{\@outputplatelabelxiii}{\oldlabel{RMAAOutputPlate13}}
\newcommand{\@outputplatelabelxiv}{\oldlabel{RMAAOutputPlate14}}
\newcommand{\@outputplatelabelxv}{\oldlabel{RMAAOutputPlate15}}
\newcommand{\@platerefi}{\oldpageref{RMAAPlate1}}
\newcommand{\@platerefii}{\oldpageref{RMAAPlate2}}
\newcommand{\@platerefiii}{\oldpageref{RMAAPlate3}}
\newcommand{\@platerefiv}{\oldpageref{RMAAPlate4}}
\newcommand{\@platerefv}{\oldpageref{RMAAPlate5}}
\newcommand{\@platerefvi}{\oldpageref{RMAAPlate6}}
\newcommand{\@platerefvii}{\oldpageref{RMAAPlate7}}
\newcommand{\@platerefviii}{\oldpageref{RMAAPlate8}}
\newcommand{\@platerefix}{\oldpageref{RMAAPlate9}}
\newcommand{\@platerefx}{\oldpageref{RMAAPlate10}}
\newcommand{\@platerefxi}{\oldpageref{RMAAPlate11}}
\newcommand{\@platerefxii}{\oldpageref{RMAAPlate12}}
\newcommand{\@platerefxiii}{\oldpageref{RMAAPlate13}}
\newcommand{\@platerefxiv}{\oldpageref{RMAAPlate14}}
\newcommand{\@platerefxv}{\oldpageref{RMAAPlate15}}
\newcounter{figplatei}
\newcounter{figplateii}
\newcounter{figplateiii}
\newcounter{figplateiv}
\newcounter{figplatev}
\newcounter{figplatevi}
\newcounter{figplatevii}
\newcounter{figplateviii}
\newcounter{figplateix}
\newcounter{figplatex}
\newcounter{figplatexi}
\newcounter{figplatexii}
\newcounter{figplatexiii}
\newcounter{figplatexiv}
\newcounter{figplatexv}
%% Printed output should use the counter `plate'
\renewcommand{\theplate}{\@arabic\c@plate}
\renewcommand{\theplateloop}{\@arabic\c@plateloop}
%% Command to output all the saved plates in this paper
\newcommand{\outputplates}{%
  \setcounter{plateloop}{0}
  \if@book\else\@endofarticlehook\fi % not entirely sure why, but this
                                     % seems to be needed
  \clearpage % first output the last page of text
  % now set pagestyle for blank even pages and make sure we start on an odd page
  \pagestyle{empty}
  \cleardoublepage
  \onecolumn
  \parindent\z@
  % first put in an extra sheet
  % the main purpose of this is so there is always a blank even page before 
  % each plate for filling in the PS file if the article has an odd number
  % of pages (see splitbook.sh)
  \null\vfil\centerline{PLATES}\vfil\clearpage\pagestyle{rmaaplate}\cleardoublepage
  % output each plate in turn
  \@whilenum\value{plateloop}<\value{plate}\do{%
    \addtocounter{plateloop}{1}
    % restore the figure counter from \@figcounti, \@figcountii, etc
    \setcounter{figure}{\value{figplate\roman{plateloop}}}
    % spew out the plate (in its own group for safety)
    {\csname @plate\roman{plateloop}\endcsname\par}
    }
  \pagestyle{rmaa}
  % reset the page number as if nothing had happened
  \addtocounter{page}{-\value{plate}}
  \addtocounter{page}{-\value{plate}}
  }
\newcommand\figcaption{%
  \def\@captype{figure}
  \renewcommand{\@makecaption}{\@makefigcap}
  \caption}
\newcommand\tabcaption{%
  \def\@captype{table}
  \renewcommand{\@makecaption}{\@maketabcap}
  \caption}

\newcommand{\plate}[2][1]{%
  %% User command to make a plate page: 
  %% Mandatory argument is what goes on the plate - this should not 
  %% include any floating environments (e.g. figure). Instead, \figcaption 
  %% should be used to do the captions. 
  %% Optional argument is the number of figures (i.e. captions) on the plate
  %% page (default is 1). This is necessary for incrementing the figure counter. 
  %%
  %% Note that we have to use a \gdef rather than an \xdef when we 
  %% save the plate output in a command. This means that in the book styles when it 
  %% comes to spewing out the plate, we will have forgotten about any user macros 
  %% that may have been defined. This can be got round by repeating the macros inside the plate.
  %% (Presumably, one would use a \newcommand rather than \renewcommand if they were new.)
  % increment global and local counters
  \addtocounter{plate}{1}
  % write appropriate label to .aux file
  \csname @platelabel\roman{plate}\endcsname
  % define a global macro called \@platei, \@plateii, \@plateiii, \@plateiv, ...
  \expandafter\gdef\csname @plate\roman{plate}\endcsname{%
    \csname @outputplatelabel\roman{plateloop}\endcsname
    \null\vfil #2 \vfil
    \cleardoublepage
    }
  % save the current shortauthor and shorttitle
  % need to be \xdef to take current values of \@shortauthor, etc 
  \expandafter\xdef\csname @shorta\roman{plate}\endcsname{\@shortauthor}
  \expandafter\xdef\csname @shortt\roman{plate}\endcsname{\@shorttitle}
  % save the current value of the figure counter for later use in \outputplates
  \setcounter{figplate\roman{plate}}{\value{figure}}
  % increment the figure counter according to the optional argument
  \addtocounter{figure}{#1}
  }
%% page style for plate pages 
\newcommand{\ps@rmaaplate}{%
  % Plate X (see page Y) in the foot
  \renewcommand\@oddfoot{\hfil\csname @shorta\roman{plateloop}\endcsname\ 
    (see page~\csname @plateref\roman{plateloop}\endcsname)\hfil}%
  \renewcommand\@evenfoot{\relax}%
  % nothing in even heads (blank pages)
  \renewcommand\@evenhead{\relax}
  % short title and page num in the odd heads
  \renewcommand\@oddhead{\hfil\csname @shortt\roman{plateloop}\endcsname%
    \hfil\llap{PLATE~\theplateloop}}
  }





%% Things for the User Guide
\if@guide
%% stolen from doc.sty
\def\MakeShortVerb#1{%
  \expandafter\ifx\csname cc\string#1\endcsname\relax
  \@shortvrbinfo{Made }{#1}%
  \add@special{#1}%
  \expandafter
  \xdef\csname cc\string#1\endcsname{\the\catcode`#1}%
  \begingroup
  \catcode`\~\active  \lccode`\~`#1%
  \lowercase{%
    \global\expandafter\let
    \csname ac\string#1\endcsname~%
    \gdef~{\verb~}}%
  \endgroup
  \global\catcode`#1\active
  \else
  \@shortvrbinfo\@empty{#1 already}%
  \fi}
\def\DeleteShortVerb#1{%
  \expandafter\ifx\csname cc\string#1\endcsname\relax
  \else
  \@shortvrbinfo{Deleted }{#1 as}%
  \rem@special{#1}%
  \global\catcode`#1\csname cc\string#1\endcsname
  \global \expandafter\let \csname cc\string#1\endcsname \relax
  \ifnum\catcode`#1=\active
  \begingroup
  \catcode`\~\active   \lccode`\~`#1%
  \lowercase{%
    \global\expandafter\let\expandafter~%
    \csname ac\string#1\endcsname}%
  \endgroup \fi \fi}
\def\@shortvrbinfo#1#2{%
  \PackageInfo{doc}{%
    #1\expandafter\@gobble\string#2 a short reference
    for \string\verb}}
\def\add@special#1{%
  \rem@special{#1}%
  \expandafter\gdef\expandafter\dospecials\expandafter
  {\dospecials \do #1}%
  \expandafter\gdef\expandafter\@sanitize\expandafter
  {\@sanitize \@makeother #1}}
\def\rem@special#1{%
  \def\do##1{%
    \ifnum`#1=`##1 \else \noexpand\do\noexpand##1\fi}%
  \xdef\dospecials{\dospecials}%
  \begingroup
  \def\@makeother##1{%
    \ifnum`#1=`##1 \else \noexpand\@makeother\noexpand##1\fi}%
  \xdef\@sanitize{\@sanitize}%
  \endgroup}
%% define | as shortcut for \verbatim
\MakeShortVerb{\|}
%% comand to specify arguments to LaTeX commands
\newcommand{\rmArg}[1]{%
  \ensuremath{\langle}{\scriptsize \textsf{\MakeUpperCase{#1}}}\ensuremath{\rangle}}
%\newcommand{\rmArg}[1]{\ensuremath{\langle}\textsc{#1}\ensuremath{\rangle}}
\providecommand{\BibTeX}{{\rmfamily B\kern-.05em%
    \textsc{i\kern-.025em b}\kern-.08em%
    T\kern-.1667em\lower.7ex\hbox{E}\kern-.125emX}}
%% Environment for code examples
\newenvironment{Code}{%
  \begin{list}{}{\setlength{\leftmargin}{0.05\columnwidth}} 
  \item[]}{\end{list}}
\fi
%</package>
%    \end{macrocode}

% \section{User commands for astronomical articles}
% \label{sec:user-commands}
% Some random definitions. These go both into the main class file and
% also into |adshtml.sty|.
% 
% Definitions of the journal name in different incarnations.  
%    \begin{macrocode}
%<*package|adshtml-style>
\newcommand{\RMAA}{RMxAA}
\newcommand{\RMAAlong}{Revista Mexicana de Astronom\'{\i}a y Astrof\'{\i}sica}
\newcommand{\RMAASC}{RMxAC}
\newcommand{\RMAASClong}{\RMAA\ (Serie de Conferencias)}
%    \end{macrocode}
% \begin{macro}{\ibidrule}
% This puts a horizontal rule with a period after it. It is used when
% the author list of a reference is identical to the preceding one. 
%    \begin{macrocode}
\newcommand{\ibidrule}{\rule[-0.03cm]{4.7em}{.01cm}.\ }
%    \end{macrocode}
% \end{macro}
% \begin{macro}{\ThisVolume}
% References to other articles in the same volume can be coded as,
% e.g., |\bibitem{} Bloggs, J. \ThisVolume{bloggs}|.
%    \begin{macrocode}
\newcommand{\ThisVolume}[1]{%
  \@year, \if@proceedings\RMAASC\else\RMAA\fi, \@volume,
  \if@book\oldpageref{#1:RMAAFirstPage}\else 000\fi\ (this volume)}
%    \end{macrocode}
% \end{macro}
% Other stuff that may be useful in ref lists.
%    \begin{macrocode}
\newcommand{\etal}{et al.\@}
%    \end{macrocode}
% Various new commands from AAS macros
%    \begin{macrocode}
\def\acknowledgments{\vskip 3ex plus .8ex minus .4ex}
\let\acknowledgements=\acknowledgments                  % second spelling
\let\jnl@style=\rm
\def\ref@jnl#1{{\jnl@style#1}}
\DeclareRobustCommand\aj{\ref@jnl{AJ}}                   
\DeclareRobustCommand\araa{\ref@jnl{ARA\&A}}             
\DeclareRobustCommand\apj{\ref@jnl{ApJ}}                 
\DeclareRobustCommand\apjl{\ref@jnl{ApJ}}                
\DeclareRobustCommand\apjs{\ref@jnl{ApJS}}               
\DeclareRobustCommand\ao{\ref@jnl{Appl.~Opt.}}           
\DeclareRobustCommand\apss{\ref@jnl{Ap\&SS}}             
\DeclareRobustCommand\aap{\ref@jnl{A\&A}}                
\DeclareRobustCommand\aapr{\ref@jnl{A\&A~Rev.}}          
\DeclareRobustCommand\aaps{\ref@jnl{A\&AS}}              
\DeclareRobustCommand\azh{\ref@jnl{AZh}}                 
\DeclareRobustCommand\baas{\ref@jnl{BAAS}}               
\DeclareRobustCommand\jrasc{\ref@jnl{JRASC}}             
\DeclareRobustCommand\memras{\ref@jnl{MmRAS}}            
\DeclareRobustCommand\mnras{\ref@jnl{MNRAS}}             
\DeclareRobustCommand\pra{\ref@jnl{Phys.~Rev.~A}}        
\DeclareRobustCommand\prb{\ref@jnl{Phys.~Rev.~B}}        
\DeclareRobustCommand\prc{\ref@jnl{Phys.~Rev.~C}}        
\DeclareRobustCommand\prd{\ref@jnl{Phys.~Rev.~D}}        
\DeclareRobustCommand\pre{\ref@jnl{Phys.~Rev.~E}}        
\DeclareRobustCommand\prl{\ref@jnl{Phys.~Rev.~Lett.}}    
\DeclareRobustCommand\pasp{\ref@jnl{PASP}}               
\DeclareRobustCommand\pasj{\ref@jnl{PASJ}}               
\DeclareRobustCommand\qjras{\ref@jnl{QJRAS}}             
\DeclareRobustCommand\skytel{\ref@jnl{S\&T}}             
\DeclareRobustCommand\solphys{\ref@jnl{Sol.~Phys.}}      
\DeclareRobustCommand\sovast{\ref@jnl{Soviet~Ast.}}      
\DeclareRobustCommand\ssr{\ref@jnl{Space~Sci.~Rev.}}     
\DeclareRobustCommand\zap{\ref@jnl{ZAp}}                 
\DeclareRobustCommand\nat{\ref@jnl{Nature}}              
\DeclareRobustCommand\iaucirc{\ref@jnl{IAU~Circ.}}
\DeclareRobustCommand\aplett{\ref@jnl{Astrophys.~Lett.}}
\DeclareRobustCommand\apspr{\ref@jnl{Astrophys.~Space~Phys.~Res.}}
\DeclareRobustCommand\bain{\ref@jnl{Bull.~Astron.~Inst.~Netherlands}}
\DeclareRobustCommand\fcp{\ref@jnl{Fund.~Cosmic~Phys.}}
\DeclareRobustCommand\gca{\ref@jnl{Geochim.~Cosmochim.~Acta}}
\DeclareRobustCommand\grl{\ref@jnl{Geophys.~Res.~Lett.}}
\DeclareRobustCommand\jcp{\ref@jnl{J.~Chem.~Phys.}}      
\DeclareRobustCommand\jgr{\ref@jnl{J.~Geophys.~Res.}}    
\DeclareRobustCommand\jqsrt{\ref@jnl{J.~Quant.~Spec.~Radiat.~Transf.}}
\DeclareRobustCommand\memsai{\ref@jnl{Mem.~Soc.~Astron.~Italiana}}
\DeclareRobustCommand\nphysa{\ref@jnl{Nucl.~Phys.~A}}
\DeclareRobustCommand\physrep{\ref@jnl{Phys.~Rep.}}
\DeclareRobustCommand\physscr{\ref@jnl{Phys.~Scr}}
\DeclareRobustCommand\planss{\ref@jnl{Planet.~Space~Sci.}}
\DeclareRobustCommand\procspie{\ref@jnl{Proc.~SPIE}}
\let\astap=\aap
\let\apjlett=\apjl
\let\apjsupp=\apjs
\let\applopt=\ao
\def\phn{\phantom{0}}
\def\phd{\phantom{.}}
\def\phs{\phantom{\ensuremath{-}}}
\def\phm#1{\phantom{#1}}
%    \end{macrocode}
% All these are redone as robust and using more modern commands where
% possible. The robustness is so they will not get expanded if used in
% abstracts and titles when these are passed to the ADS file. 
% \changes{3.16}{2003/04/07}{
%  Added \cs{Color} command for doing photometric colors.}
%    \begin{macrocode}
%%\def\sun{\hbox{$\odot$}}
\DeclareRobustCommand\sun{\ensuremath{\odot}}
%%\def\earth{\hbox{$\oplus$}}
\DeclareRobustCommand\earth{\ensuremath{\oplus}}
%%\def\lesssim{\mathrel{\hbox{\rlap{\hbox{\lower4pt\hbox{$\sim$}}}\hbox{$<$}}}}
%%\def\gtrsim{\mathrel{\hbox{\rlap{\hbox{\lower4pt\hbox{$\sim$}}}\hbox{$>$}}}}
\DeclareRobustCommand\lesssim{\mathrel{\hbox{\smaller\smaller\rlap{\hbox{\lower3pt\hbox{$\sim$}}}\raise1.5pt\hbox{$<$}}}}
\DeclareRobustCommand\gtrsim{\mathrel{\hbox{\smaller\smaller\rlap{\hbox{\lower3pt\hbox{$\sim$}}}\raise1.5pt\hbox{$>$}}}}
%%\def\lta{\mathrel{\spose{\lower 3pt\hbox{$\mathchar"218$}}
%%     \raise 2.0pt\hbox{$\mathchar"13C$}}}
%%\def\gta{\mathrel{\spose{\lower 3pt\hbox{$\mathchar"218$}}
%%     \raise 2.0pt\hbox{$\mathchar"13E$}}}
\DeclareRobustCommand\sq{\hbox{\rlap{$\sqcap$}$\sqcup$}}
%%\def\arcdeg{\hbox{$^\circ$}}
%%\def\arcmin{\hbox{$^\prime$}}
%%\def\arcsec{\hbox{$^{\prime\prime}$}}
\DeclareRobustCommand\arcdeg{\ensuremath{^\circ}}
\DeclareRobustCommand\arcmin{\ensuremath{^\prime}}
\DeclareRobustCommand\arcsec{\ensuremath{^{\prime\prime}}}
%%\def\fd{\hbox{$.\!\!^{\rm d}$}}
%%\def\fh{\hbox{$.\!\!^{\rm h}$}}
%%\def\fm{\hbox{$.\!\!^{\rm m}$}}
%%\def\fs{\hbox{$.\!\!^{\rm s}$}}
%%\def\fdg{\hbox{$.\!\!^\circ$}}
%%\def\farcm{\hbox{$.\mkern-4mu^\prime$}}
%%\def\farcs{\hbox{$.\!\!^{\prime\prime}$}}
%%\def\fp{\hbox{$.\!\!^{\scriptscriptstyle\rm p}$}}
\DeclareRobustCommand\fd{\ensuremath{.\!\!^\mathrm{d}}}
\DeclareRobustCommand\fh{\ensuremath{.\!\!^\mathrm{h}}}
\DeclareRobustCommand\fm{\ensuremath{.\!\!^\mathrm{m}}}
\DeclareRobustCommand\fs{\ensuremath{\mkern0.5mu.\mkern-5mu^\mathrm{s}\mkern-1mu}}
\DeclareRobustCommand\fdg{\ensuremath{.\!\!^\circ}}
\DeclareRobustCommand\farcm{\ensuremath{.\mkern-4mu^\prime}}
\DeclareRobustCommand\farcs{\ensuremath{\mkern1mu.\mkern-6mu^{\prime\prime}\mkern-1.5mu}}
%%\DeclareRobustCommand\farcs{\ensuremath{.\!\!^{\prime\prime}}}
\DeclareRobustCommand\fp{\ensuremath{.\!\!^{\scriptscriptstyle\rm p}}}
%%\def\micron{\hbox{$\mu$m}}
\DeclareRobustCommand\micron{\ensuremath{\mu\mathrm{m}}}
\let\la=\lesssim                        
\let\ga=\gtrsim
%%\def\case#1#2{\hbox{$\frac{#1}{#2}$}}
%%\def\slantfrac#1#2{\hbox{$\,^#1\!/_#2$}}
%%\def\onehalf{\slantfrac{1}{2}}
%%\def\onethird{\slantfrac{1}{3}}
%%\def\twothirds{\slantfrac{2}{3}}
%%\def\onequarter{\slantfrac{1}{4}}
%%\def\threequarters{\slantfrac{3}{4}}
\DeclareRobustCommand\case[2]{\ensuremath{\frac{#1}{#2}}}
\DeclareRobustCommand\slantfrac[2]{\ensuremath{\,^#1\!/_#2}}
\DeclareRobustCommand\onehalf{\slantfrac{1}{2}}
\DeclareRobustCommand\onethird{\slantfrac{1}{3}}
\DeclareRobustCommand\twothirds{\slantfrac{2}{3}}
\DeclareRobustCommand\onequarter{\slantfrac{1}{4}}
\DeclareRobustCommand\threequarters{\slantfrac{3}{4}}
%%\def\ubvr{\hbox{$U\!BV\!R$}}            
%%\def\ub{\hbox{$U\!-\!B$}}               
%%\def\bv{\hbox{$B\!-\!V$}}               
%%\def\vr{\hbox{$V\!-\!R$}}               
%%\def\ur{\hbox{$U\!-\!R$}}
%    \end{macrocode}
% New commands for ``photometric'' colors 
%    \begin{macrocode}
\DeclareRobustCommand\ubvr{\ensuremath{U\!BV\!R}}            
\DeclareRobustCommand\Color[2]{%
  \ensuremath{%
    #1\!\ifx#1V\!\fi-\!#2
  }}
\DeclareRobustCommand\ub{\Color{U}{B}}
\DeclareRobustCommand\bv{\Color{B}{V}}
\DeclareRobustCommand\vr{\Color{V}{R}}
\DeclareRobustCommand\ur{\Color{U}{R}}
%% New version of the \ion command, that actually works! 
%    \end{macrocode}
% \begin{macro}{\ion}
% This version of the |\ion| command allows you to write either
% |\ion{N}{4}| or |\ion{N}{iv}| in order to get N~IV. 
%    \begin{macrocode}
\DeclareRobustCommand\ion[2]{%
  \mbox{#1\kern0.2em%
    \smaller\rmfamily%
    \edef\@tempa{\@car#2\@nil}%
    \ifcat1\@tempa%
    \@Roman{#2}%
    \else%
    \uppercase{#2}%
    \fi}}
\newcounter{IonStage}
\DeclareRobustCommand\plainion[2]{\setcounter{IonStage}{#2}#1 \Roman{IonStage}}
%%\def\ion#1#2{\ensuremath{\mathrm{#1}\;}{\protect\small\rm\@Roman{#2}}}
%    \end{macrocode}
% \end{macro}
%    \begin{macrocode}
\DeclareRobustCommand\nodata{$\cdots$}
%    \end{macrocode}
% These are needed by the |\ion| command. Isn't there an easier way of
% doing this?
%    \begin{macrocode}
\let\@small=\small
\let\@normalsize=\normalsize
\let\@footnotesize=\footnotesize
\let\@scriptsize=\scriptsize
\let\@tiny=\tiny
\let\@large=\large
\let\@Large=\Large
\let\@LARGE=\LARGE
\let\@huge=\huge
\let\@Huge=\Huge
\newcommand{\smaller}{\small}
\renewcommand{\normalsize}{\renewcommand{\smaller}{\small}\@normalsize}
\renewcommand{\small}{\renewcommand{\smaller}{\footnotesize}\@small}
\renewcommand{\footnotesize}{\renewcommand{\smaller}{\scriptsize}\@footnotesize}
\renewcommand{\scriptsize}{\renewcommand{\smaller}{\tiny}\@scriptsize}
\renewcommand{\tiny}{\renewcommand{\smaller}{\tiny}\@tiny}
\renewcommand{\large}{\renewcommand{\smaller}{\normalsize}\@large}  
\renewcommand{\Large}{\renewcommand{\smaller}{\large}\@Large}  
\renewcommand{\LARGE}{\renewcommand{\smaller}{\Large}\@LARGE}  
\renewcommand{\huge}{\renewcommand{\smaller}{\LARGE}\@huge}  
\renewcommand{\Huge}{\renewcommand{\smaller}{\huge}\@Huge}  
%</package|adshtml-style>
%    \end{macrocode}


% \section{Other files}

% \subsection{The driver and make files}
% \label{sec:driver_make}
% These are put here for completeness, although they still need to be
% distributed separately since otherwise the user won't be able to
% unpack this file into its constituent parts.
% 
% The |Makefile| has two principal rules: |docs|, and |macros|. The
% |docs| rule produces a PDF and gzipped PS file of the documentation.
% I still haven't perfected the combination of rules required to get
% it to run \LaTeX{} enough times in all circumstances. The |macros|
% rule runs \LaTeX{} on |rmaa.ins| in order to unpack |rmaa.cls| and
% all the other little files (including the |Makefile| and the
% |rmaa.ins| files themselves\dots 
% 
% First set up lists of various types of file. 
%    \begin{macrocode}
%<*makefile>
DOCFILES=rmaa.pdf
MACRO_LTX_FILES=rmaa.cls rmaacite.sty adshtml.sty
SHELLSCRIPTS=makeads.sh
MACRO_OTHER_FILES=adshtml.perl sample.latex2html-init
MAKEFILE=Makefile
MACROFILES=$(MACRO_LTX_FILES) $(SHELLSCRIPTS) $(MACRO_OTHER_FILES)         
SKELETONFILES=ADS.tex 
%    \end{macrocode}
% Install the files read only into the package directory, specified via a path relative to the source directory. \changes{4.0}{2019/11/14}{Reorganize folder layout so that hardly anything is in the top level.  That way, things are clearer on github.}
%    \begin{macrocode}
INSTALL=install 
INSTALL_DIR=../packages
%    \end{macrocode}
% The default rule is to make everything. Installation must be
% explicitly asked for.
%    \begin{macrocode}
all: docs macros

install: install-docs install-macros

%    \end{macrocode}
% Most files are installed read-only to prevent editing. Skeleton
% files must be writeable. Must be careful not to overwrite the
% Makefile in the parent directory. 
%    \begin{macrocode}
install-macros: macros
        $(INSTALL) -m 444 $(MACRO_LTX_FILES) $(MACRO_OTHER_FILES) $(INSTALL_DIR)
        $(INSTALL) -m 555 $(SHELLSCRIPTS) $(INSTALL_DIR)
        $(INSTALL) -m 644 $(SKELETONFILES) $(INSTALL_DIR)

install-docs: docs
        $(INSTALL) rmaa.pdf $(INSTALL_DIR)/rm-fulldocs.pdf

%    \end{macrocode}
% Make the PDF and gzipped PS files. 
%    \begin{macrocode}
docs: $(DOCFILES)

%    \end{macrocode}
% Make everything else if either the |.dtx| or |.ins| files have
% changed. 
%    \begin{macrocode}
macros: rmaa.dtx rmaa.ins
%    \end{macrocode}
% Ensure that all the macro files can be written to then recreate
% them. The |touch| is there in case any of the files don't exist
% yet. 
%    \begin{macrocode}
        touch $(MACROFILES) $(MAKEFILE)
        chmod u+w $(MACROFILES) $(MAKEFILE) rmaa.ins
        latex "%%
%% This is file `rmaa.ins',
%% generated with the docstrip utility.
%%
%% The original source files were:
%%
%% rmaa.dtx  (with options: `ins')
%% 
%% PLEASE DO NOT EDIT THIS FILE! It is generated automatically from
%% the file rmaa.dtx using the accompanying Makefile and rmaa.ins
%% Any hand-edited changes will be lost
%% 
%% A LaTeX document class for the RMxAA
%% 
%% Author: William Henney (w.henney@crya.unam.mx)
%% 
%% Latest version available from: http://www.crya.unam.mx/rmaa
%% 
%% Copyright (C) 1999-2010 Instituto de Astronomia, UNAM, Mexico
%% 
%% This program is free software; you can redistribute it and/or
%% modify it under the terms of the GNU General Public License
%% as published by the Free Software Foundation; either version 2
%% of the License, or (at your option) any later version.
%% 
%% This program is distributed in the hope that it will be useful,
%% but WITHOUT ANY WARRANTY; without even the implied warranty of
%% MERCHANTABILITY or FITNESS FOR A PARTICULAR PURPOSE.  See the
%% GNU General Public License for more details.
%% 
%% You should have received a copy of the GNU General Public License
%% along with this program; if not, write to the Free Software
%% Foundation, Inc., 59 Temple Place - Suite 330, Boston, MA  02111-1307, USA.





%%
%% Installation batch file for the RMAA LaTeX package
%%
\def\batchfile{rmaa.ins}
\input docstrip.tex
\preamble

PLEASE DO NOT EDIT THIS FILE! It is generated automatically from
the file rmaa.dtx using the accompanying Makefile and rmaa.ins
Any hand-edited changes will be lost

A LaTeX document class for the RMxAA

Author: William Henney (w.henney@crya.unam.mx)

Latest version available from: http://www.crya.unam.mx/rmaa

Copyright (C) 1999-2010 Instituto de Astronomia, UNAM, Mexico

This program is free software; you can redistribute it and/or
modify it under the terms of the GNU General Public License
as published by the Free Software Foundation; either version 2
of the License, or (at your option) any later version.

This program is distributed in the hope that it will be useful,
but WITHOUT ANY WARRANTY; without even the implied warranty of
MERCHANTABILITY or FITNESS FOR A PARTICULAR PURPOSE.  See the
GNU General Public License for more details.

You should have received a copy of the GNU General Public License
along with this program; if not, write to the Free Software
Foundation, Inc., 59 Temple Place - Suite 330, Boston, MA  02111-1307, USA.
\endpreamble
\askforoverwritefalse
\generate{%
  \file{rmaa.cls}{\from{rmaa.dtx}{package}}
  \file{rmaa.ins}{\from{rmaa.dtx}{ins}}
  \file{rmaacite.sty}{\from{rmaa.dtx}{rmaacite}}
  \file{adshtml.sty}{\from{rmaa.dtx}{adshtml-style}}
  \file{ADS.tex}{\from{rmaa.dtx}{ads-container}}
  }
\generate{\nopreamble\nopostamble%
  \file{Makefile}{\from{rmaa.dtx}{makefile}}
  \file{adshtml.perl}{\from{rmaa.dtx}{adshtml-perl}}
  \file{makeads.sh}{\from{rmaa.dtx}{ads-shell-script}}
  \file{sample.latex2html-init}{\from{rmaa.dtx}{latex2html-init}}
  }

%% Local Variables:
%% mode: LaTeX
%% mode: Font-lock
%% End:









\endinput
%%
%% End of file `rmaa.ins'.
"
%    \end{macrocode}
% \TeX{} automatically adds the |.tex| extension to extensionless
% files. We fix this and take the opportunity to convert
% back to TABs all that it has lovingly converted to spaces. 
%    \begin{macrocode}
        cp -f $(MAKEFILE) $(MAKEFILE).BCK && unexpand $(MAKEFILE).tex >| $(MAKEFILE) && rm $(MAKEFILE).tex
%    \end{macrocode}
% Make the shell scripts executable and everything unwriteable (so
% they don't get hand-edited by mistake). However, undo this for the
% files that the user is \emph{supposed} to modify. 
%    \begin{macrocode}
        chmod a+x $(SHELLSCRIPTS)
        chmod a-w $(MACROFILES) $(MAKEFILE)
        chmod u+w $(SKELETONFILES)

%    \end{macrocode}
% We now use PDF\LaTeX{} for this.
%    \begin{macrocode}
rmaa.idx: rmaa.dtx
        @echo "REGENERATING .idx FILE"
        pdflatex "\nonstopmode% \def\fileversion{4.0} \def\filedate{2019/11/12}
% \MakeShortVerb{\|}
% \title{A \LaTeX{} document class for the Revista Mexicana de
% Astronom\'\i a y Astrof\'\i sica\thanks{This file has version
% \fileversion{} last revised \filedate{}.}}
% \author{William J. Henney\thanks{w.henney@irya.unam.mx}}
% \date{\filedate}
% \maketitle
% \paragraph{Abstract}
% This document describes the \texttt{rmaa} \LaTeX{}
% document class for the journal ``Revista Mexicana de
% Astronom\'\i a y Astrof\'\i sica''. The macros can be used for 
% \begin{enumerate}
% \item Preparing author manuscripts for the journal (including
% conference proceedings to be published in the ``Serie de
% Conferencias'')
% \item Typesetting an entire volume of the journal by inputting the
% individual article files from a container file. Facilities are
% included for automatic generation of TOC and Author Index. 
% \end{enumerate}
% 
% \tableofcontents
% \section{The documentation driver file}
% This is the part that generates the documentation that you may or
% may not be reading, which is done with the help of the Makefile
% discussed in section~\ref{sec:driver_make}. 
%    \begin{macrocode}
%<*driver>
\documentclass{ltxdoc}
\usepackage{hyperref}
\usepackage{palatino}
\EnableCrossrefs \CodelineIndex \RecordChanges
\setlength{\parindent}{0pt}
\def\MacroFont{\footnotesize\tt}
\begin{document} \DocInput{rmaa.dtx} \PrintIndex \PrintChanges \end{document}
%</driver>
%    \end{macrocode}

% \section{The implementation}
% 
% \changes{3.24}{2007/09/14}{Added user access to the macro version
% and date via \cs{rmVersion} and \cs{rmDate}}
% 
%    \begin{macrocode}
%<*package>
\newcommand{\rmaa@version}{3.26}
\newcommand{\rmaa@date}{2011/08/15}
\newcommand\rmVersion{\rmaa@version}
\newcommand\rmDate{\rmaa@date}
\NeedsTeXFormat{LaTeX2e}[1995/06/01]
%    \end{macrocode}
% \changes{4.0}{2019/11/02}{Change abstract to appear before
% resumen. Change RevMexAA to RMxAA for consistency with ADS.}
%
% \changes{3.22}{2007/03/22}{Allow proceedings articles to have
%   received/accepted dates} \changes{3.21}{2006/05/15}{Changes for
%   Silvia - abstracts are now shipped out to ADS.}
% \changes{3.19}{2006/01/19}{Changes for Silvia - remove adresses from
%   end, new abstract style} \changes{2.2beta}{2001/10/30}{ Fixed bug
%   whereby twocolumn option was not passed through to book
%   class. Finally sorted out the ``label on last page'' issue. }
% \changes{2.1}{2001/01/23}{ Fixed some inconsistencies in
%   \cs{nocite} and \cs{bibitem} uniqification Fixed the definition of
%   \cs{ion} so it will work in math mode and in contexts where the
%   typesize is different from normal (e.g., title, fig captions).
%   Modified the definition of \cs{@caption}, so we don't write any
%   list-of-figures data (we never use it anyway, and it just causes
%   problems with fragile commands in fig captions). }
% \changes{2.0}{2000/11/30}{ Added `newstyle' option, which should be
%   used in all new proceedings volumes (wider abstract, smaller
%   typeface in fig captions, etc.) Added commands for formatting
%   table captions. New demonstation document rmpostsamp.tex for
%   1-page posters. This will be the last release under a PRI
%   government (for a while...)}  \changes{1.9}{2000/06/06}{
%   Corrected position of page number in footer of articles that open
%   on even pages. Ensured that first article in each session opens on
%   an odd page, as do list of abstracts and author index. Added extra
%   space in TOC before the backmatter. Added command for editors,
%   \cs{BoldSessionTitles}, that makes produces bold titles for
%   session in the TOC.}  \changes{1.8}{2000/05/16}{ Fixed bug in
%   last page label with \cs{suppressfulladdresses} }
% \changes{1.7}{2000/05/11}{ Fixed more little things. Page labels
%   now work, even if you have no bibliography. You can output the
%   full addresses by hand with the \cs{outputfulladdresses}
%   command. Editors' Guide updated but still not finished. Revamped
%   splitbook.sh now produces PS file of each article for web page and
%   also an EPS file of each page for sending to the printers.}
% \changes{1.6}{2000/04/15}{ More cleaning up: full stops after
%   titles in index, running headers and footers in frontmatter,
%   fontsize in title of index, etc} \changes{1.5}{2000/04/15}{
%   Eliminated superscripts from titles in TOC and index}
% \changes{1.4}{2000/04/14}{ Added copyright and license info (GPL)}
% \changes{1.3}{2000/03/01}{ Corrected \cs{figcaption} to do the
%   caption formatting right Corrected font size of sec headings in
%   front matter} \changes{1.2}{2000/02/16}{ Added \cs{plate}
%   command} \changes{1.1}{2000/02/14}{ Added manuscript option,
%   corrected bugs} \changes{1.0}{2000/02/08}{ Now does a whole book
%   - lots of stuff tidied up} \changes{0.7}{1999/10/04}{ Increased
%   spacing between sections in single column styles. Suppressed
%   floats at bottom of pages (didn't work though) }
% \changes{0.6}{1999/09/24}{ Tweaked \cs{adjustfinalcols} (still not
%   perfect) Added proceedings option Added onecolumn/twocolumn
%   options (orthogonal to journal/proceedings) }
% \changes{0.5}{1999/09/23}{ Added \cs{fulladdresses} and
%   \cs{adjustfinalcols} (just about working)}
% \changes{0.4}{1999/09/22}{ Minor stylistic changes}
% \changes{0.3}{1999/06/28}{ Better mechanism for alternative
%   affiliations Added `guide' option for user guide}
% \changes{0.2}{1999/05/27}{ All basic functionality present, I
%   think} \changes{0.1}{1999/05/26}{ Just about working - still
%   missing Fig/Tab/Refs} \changes{0.0}{1999/05/21}{ So bored today,
%   decided to rewrite RMAA LaTeX macros from scratch}
%    \begin{macrocode}
\ProvidesClass{rmaa}[\rmaa@date\space v\rmaa@version\space RMxAA article
class]
%    \end{macrocode}
% \subsection{Option Handling}
% \label{sec:options}
% Declarations that are needed in options
% \changes{3.05}{2003/10/22}{Added natbib option}
% \changes{3.06}{2003/01/03}{Natbib option turned on permanently}
% \changes{3.07}{2003/01/11}{Added proof option}
% \changes{3.14}{2003/01/11}{Option to turn off natbib}
% \changes{3.18}{2003/10/22}{trying out doctex mode} 
% \changes{3.23}{2007/09/09}{Manuscript is now on by default.}
% \changes{3.26}{2011/08/15}{New option for review articles}
%    \begin{macrocode}
\newif\if@reviewarticle
\@reviewarticlefalse
\newif\if@proof
\@prooffalse
\newif\if@journal
\@journaltrue
\newif\if@printdates
\@printdatesfalse
\newif\if@proceedings
\@proceedingsfalse
\newif\if@manuscript
\@manuscripttrue
\newif\if@bigfont
\@bigfonttrue
\newif\if@guide
\@guidefalse
\newif\if@preprint
\@preprintfalse
\newif\if@rmtwocol
\@rmtwocoltrue
\newif\if@biblist
\@biblistfalse
\newif\if@book
\@bookfalse
\newif\if@natbib
\@natbibtrue
\newif\if@rmcite
\@rmcitefalse
\newif\if@titlevisible
\@titlevisibletrue
\newif\if@abstractpresent
\@abstractpresentfalse
\newif\if@printbanner
\@printbannerfalse
\newif\if@suppressfulladdresses
\@suppressfulladdressesfalse
\newif\if@newstyle
\@newstylefalse
\newif\if@debug
\@debugfalse
\newif\if@boldkeywords
\@boldkeywordsfalse
%    \end{macrocode}
% The option handling---mainly setting various booleans.
% \changes{2.2beta11}{2001/11/26}{Made twocol and
% newstyle the default for proceedings.}
% \changes{3.23}{2007/09/09}{Manuscript is now 10 pt instead of 12pt}
% \changes{3.23}{2007/09/10}{Full addresses are now automatically
% suppressed for proceedings style.}
% \changes{3.24}{2007/09/14}{New option \texttt{boldkeywords} that
% does the keywords in the old ugly way. Now off by default.}
%    \begin{macrocode}
\DeclareOption{manuscript}{\typeout{Manuscript for referee} 
  \@manuscripttrue\@journalfalse\@rmtwocolfalse
  \PassOptionsToClass{10pt}{article}}
\DeclareOption{bigfont}{\typeout{Big font manuscript requested by Gary
  Ferland} 
  \@manuscripttrue\@journalfalse\@rmtwocolfalse\@bigfonttrue
  \PassOptionsToClass{11pt}{article}}
\DeclareOption{journal}{\typeout{Main journal option}%
  \@journaltrue\@rmtwocoltrue\@printdatestrue\@manuscriptfalse
  \PassOptionsToClass{10pt}{article}}
\DeclareOption{guide}{\typeout{User guide option}%
  \@guidetrue\@journalfalse\@rmtwocoltrue\@manuscriptfalse
  \PassOptionsToClass{10pt}{article}}
\DeclareOption{preprint}{\typeout{Preprint option}%
  \@preprinttrue\@manuscriptfalse\PassOptionsToClass{10pt}{article}}
\DeclareOption{proof}{\typeout{Author proof option}%
  \@prooftrue\@manuscriptfalse}
\DeclareOption{proceedings}{\typeout{Proceedings option}%
  \@proceedingstrue\@journalfalse\@rmtwocoltrue\@newstyletrue
  \@suppressfulladdressestrue
  \@printdatesfalse\@manuscriptfalse\PassOptionsToClass{10pt}{article}} 
\DeclareOption{printdates}{\typeout{Force printing of received/accepted dates}%
  \@printdatestrue} 
\DeclareOption{book}{\typeout{Whole issue option}%
  \@booktrue} 
\DeclareOption{nonatbib}{\typeout{Turn off natbib support (required
    for compatibility with old rmaacite package).}%
  \@natbibfalse} 
\DeclareOption{newstyle}{\typeout{Use new layout style}%
  \@newstyletrue}
\DeclareOption{oldstyle}{\typeout{Use old layout style}%
  \@newstylefalse}
\DeclareOption{debug}{\typeout{Loads of annoying messages}%
  \@debugtrue}
%    \end{macrocode}
% These two need to be declared explicitly so that |\ExecuteOptions|
% actually does something.
%    \begin{macrocode}
\DeclareOption{twocolumn}{\@rmtwocoltrue}
\DeclareOption{onecolumn}{\@rmtwocolfalse}
\DeclareOption{twoside}{\PassOptionsToClass{\CurrentOption}{article}}
\DeclareOption{boldkeywords}{\@boldkeywordstrue}
%    \end{macrocode}
% Bold, upper-case keywords used to be the default. Now they must be
% turned on explicitly with the \texttt{boldkeywords} option. 
%    \begin{macrocode}
\DeclareOption{boldkeywords}{\@boldkeywordstrue}
%    \end{macrocode}
% Pass any other option along to the \texttt{article} class.
%    \begin{macrocode}
\DeclareOption*{\PassOptionsToClass{\CurrentOption}{article}}
%    \end{macrocode}
% Default options are to be two-sided and based on the standard
% \LaTeX{} \texttt{article} class,
% unless we are processing the entire book, in which case we base
% ourselves on the standard \texttt{book} class. Default is to use
% manuscript option unless otherwise specified. 
%    \begin{macrocode}
\ExecuteOptions{twoside,manuscript} \ProcessOptions\relax
\if@book
\if@rmtwocol\LoadClass[twocolumn]{book}\else\LoadClass{book}\fi
\newcommand{\abstract}{\relax}
\newcommand{\@maketitle}{\relax}
\else
\if@rmtwocol\LoadClass[twocolumn]{article}\else\LoadClass{article}\fi
\fi
%    \end{macrocode}
% Attempt to use natbib for citations 02 Jan 2003. 
% \changes{3.05}{2003/01/02}{If the natbib option is used, then set
% things up so that natbib's version of the bibliography is identical
% to the standard one.}
% \changes{3.09}{2003/01/23}{Explicitly require a recent version of
% natbib, so that we get a decent error message if the package is too
% old.} 
% \changes{3.23}{2007/09/09}{Now allow Bibtex to work properly} 
% \changes{3.25}{2009/10/01}{Change \cs{newcommand} tp \cs{def} when setting \cs{bibfont}. This allows it to work with all versions of natbib.} 
% \changes{3.25}{2010/02/06}{Updated copyright year.} 
% \changes{3.25}{2010/02/06}{Updated email and web addresses to \texttt{crya.unam.mx}.} 
%    \begin{macrocode}
\if@natbib
\RequirePackage{natbib}[1999/02/25]
\bibpunct{(}{)}{;}{a}{}{,}
\def\bibfont{\small}
\renewcommand\bibsection{\section*{REFERENCES}}
\setlength\bibsep{0pt}
\setlength\bibhang{1.5em}
\bibliographystyle{rmaa}
\fi
%    \end{macrocode}
% We want the \texttt{graphicx} package for the extended syntax of
% |\includegraphics|. We want the \texttt{ifthen} package for using in
% these macros (although we don't much yet). We explicitly load
% \texttt{babel} since some \LaTeX{} distributions messed up the
% hyphenation if we didn't.  
%    \begin{macrocode}
\RequirePackage{graphicx}
\RequirePackage{ifthen}
\IfFileExists{babel.sty}{%
%   \let\@restoreindent=\@afterindentfalse
  \RequirePackage[english]{babel}
%   \@restoreindent
  %% reinstate the definition of \@roman so as not to break .idx files
  \@ifundefined{org@@roman}{\relax}{\renewcommand{\@roman}[1]{\org@@roman{##1}}}
  }{%
  \typeout{Package babel.sty not found - don't worry about this}}
%    \end{macrocode}
% Try to load the \texttt{dcolumn} and \texttt{booktabs} packages if
% they can be found. These can produce nicer-looking tables. 
%    \begin{macrocode}
\IfFileExists{dcolumn.sty}{%
  \RequirePackage{dcolumn}
  }{%
  \typeout{Package dcolumn.sty not found - you will not be able to
    align columns on decimal points in tables}}
\IfFileExists{booktabs.sty}{%
  \RequirePackage{booktabs}
  }{%
  \typeout{Package booktabs.sty not found - your tables will not look
    as pretty as they ought too.}
  }
\IfFileExists{makeidx.sty}{%
  \RequirePackage{makeidx}
}{%
  \typeout{Package makeidx.sty not found - Author Index may not work.}
  }
%% commented out for moment since requires more modern version of LaTeX
%% \RequirePackage[latin1]{inputenc}
%    \end{macrocode}
% \changes{4.0}{2019/11/13}{Make \texttt{hyperref} and
% \texttt{etoolbox} into hard requirements, since we now use them
% extensively.}  We use hyperref to get internal and external links.
% But I don't like the colored boxes so I turn those off.  We use
% etoolbox for patching and robustifying at the moment.  In the
% future, I want to port a lot of other stuff over to it.  For
% instance, replace the ifthen package and the TeX ifs too.
%
% We target the latest version (1.7) of PDF, rather than the default
% of 1.5.  This reduces the chance of getting warnings from included
% graphics.  This is not strictly related to hyperref, but we put it
% here anyway.
%    \begin{macrocode}
\RequirePackage{etoolbox}
\RequirePackage{hyperref}
\hypersetup{hidelinks}
\pdfminorversion=7
%    \end{macrocode}
% We patch the \cs{@combinedblfloats} following a suggestion by Heiko
% Oberdiek at \url{https://tex.stackexchange.com/a/249743/3843}, which
% avoids problems when hyperlinks are split across
% pages. \changes{4.0}{2019/11/12}{Fix for issue where breaking
% hyperlink across pages gave an error.}
%    \begin{macrocode}
\urlstyle{same}
\newcommand\email{\begingroup \urlstyle{rm}\Url}
\patchcmd\@combinedblfloats{\box\@outputbox}{\unvbox\@outputbox}{}{%
  \errmessage{\noexpand\@combinedblfloats could not be patched}%
}%
%    \end{macrocode}
% \subsection{Other setup stuff}\label{sec:setup}
% Some declarations: 
%    \begin{macrocode}
\newif\if@needaddresses
\@needaddressestrue
\newif\if@frontmatter
\@frontmatterfalse
\newif\if@inpreamble
\@inpreambletrue
\newif\if@firstauthor
\@firstauthortrue
\newif\if@firstaffil
\@firstaffiltrue
\newif\if@firstaltaffil
\@firstaltaffiltrue
\newif\if@firstkwd
\@firstkwdtrue
\newif\if@noexpandADS
\@noexpandADSfalse
\newif\if@sessiontitlepages
\@sessiontitlepagesfalse
%    \end{macrocode}
% \begin{macro}{\@rmaaEmpty}
% \begin{macro}{\@rmaaArg}
% These are used in various commands to check if arg was left blank
%    \begin{macrocode}
\newcommand{\@rmaaEmpty}{}
\newcommand{\@rmaaArg}{}
%    \end{macrocode}  
% \end{macro}
% \end{macro}
%
% \subsection{Compatibility with old \LaTeX{} versions}
% \label{sec:compat}
% This package is based on the 1998/12/01 release of \LaTeX{}.
% Older releases may not define all the commands that we need.
% Luckily, we can use |\providecommand| to define these if necessary. 
% The following commands are not defined in the 1994/06/01 release. 
%
% \begin{macro}{\MakeUpperCase}
% \begin{macro}{\textsuperscript}
% \begin{macro}{\textbackslash}
% User commands:
%    \begin{macrocode}
\providecommand{\MakeUpperCase}[1]{\uppercase{#1}}
\providecommand{\textsuperscript}[1]{\ensuremath{^\textrm{#1}}}
\providecommand{\textbackslash}{\ensuremath{\backslash}}
%    \end{macrocode}  
% \end{macro}
% \end{macro}
% \end{macro}
% \begin{macro}{\protected@xdef}
% \begin{macro}{\restore@protect}
% \begin{macro}{\@unexpandable@protect}
% \begin{macro}{\hb@xt@}
% \LaTeX{} internal commands:
%    \begin{macrocode}
\providecommand{\protected@xdef}{%
        \let\@@protect\protect
        \let\protect\@unexpandable@protect
        \afterassignment\restore@protect
        \xdef}
\providecommand{\restore@protect}{\let\protect\@@protect}
\providecommand{\@unexpandable@protect}{\noexpand\protect\noexpand}
\providecommand{\hb@xt@}{\hbox to}
%    \end{macrocode}
% \end{macro}
% \end{macro}
% \end{macro}
% \end{macro}
%
% \subsection{Page layout}\label{sec:page-layout}
%
% The dimensions of the text were defined by Elizabeth. Text width is
% 169\,mm with equal margins on the two sides: 
%    \begin{macrocode}
\setlength{\textwidth}{169mm}
\setlength{\@tempdima}{\paperwidth}
\addtolength{\@tempdima}{-\textwidth}
\setlength{\oddsidemargin}{0.5\@tempdima}
\addtolength{\oddsidemargin}{-1in}
\setlength{\evensidemargin}{\oddsidemargin}
%    \end{macrocode}
% Text height is 225\,mm, top margin is 14\,mm. 
%    \begin{macrocode}
\setlength{\textheight}{225mm} 
\setlength{\topmargin}{-1in} 
\addtolength{\topmargin}{14mm}
%    \end{macrocode}
% Space between columns in two-column styles, header, footer,
% paragraph indentation:
%    \begin{macrocode}
\setlength{\columnsep}{9.5mm}
\setlength{\headsep}{0.7cm}
\setlength{\footskip}{0.75cm}
\setlength{\parindent}{1.5em}
%    \end{macrocode}
% 
% \begin{macro}{\@set@scalable@length}
% Try out a new scheme for dealing with various spacing dimensions and
% allowing the user to easily modify them. The command
% \begin{quote}
% |\set@scalable@length|\marg{name}\marg{defaultlength}\marg{stretch}\marg{shrink}
% \end{quote}
% does a 
% |\setlength|\marg{\cs{name}}\marg{defaultlength plus stretch minus shrink}
% and also defines the user command
% \textbackslash\meta{name}|Scale|\marg{factor} that can 
% be used to rescale length by \meta{factor}. Additionally, now
% defines the package command
% \textbackslash\meta{name}|@set@scale|\marg{FACTOR} 
% that also rescales the length by \meta{FACTOR}, so that the length
% becomes \meta{factor} $\times$ \meta{FACTOR} $\times$
% \meta{defaultlength}.   
%    \begin{macrocode}
\newcommand\@set@scalable@length[4]{%
  %% arg #1: NAME of length without leading `\'
  %%
  %% arg #2: default VALUE  
  \expandafter\newlength\csname @#1@default\endcsname
  \expandafter\setlength\csname @#1@default\endcsname{#2}%
  \expandafter\newlength\csname @@#1@@default\endcsname
  \expandafter\setlength\csname @@#1@@default\endcsname{#2}%
  %% arg #3: default STRETCH  
  \expandafter\newlength\csname @#1@stretch\endcsname
  \expandafter\setlength\csname @#1@stretch\endcsname{#3}%
  \expandafter\newlength\csname @@#1@@stretch\endcsname
  \expandafter\setlength\csname @@#1@@stretch\endcsname{#3}%
  %% arg #4: default SHRINK  
  \expandafter\newlength\csname @#1@shrink\endcsname
  \expandafter\setlength\csname @#1@shrink\endcsname{#4}%
  \expandafter\newlength\csname @@#1@@shrink\endcsname
  \expandafter\setlength\csname @@#1@@shrink\endcsname{#4}%
  %%
  %% User-defined scale \@NAME@scale is unity by default
  \expandafter\gdef\csname @#1@scale\endcsname{1}%
  %% Package-defined scale \@NAME@SCALE is unity by default
  \expandafter\gdef\csname @#1@SCALE\endcsname{1}%
  %%
  %% internal command \@reset@NAME should expand to 
  %% \setlength\@@NAME@@default{\@NAME@SCALE\@NAME@default}
  %% \setlength\NAME{\@NAME@scale\@@NAME@@default}
  %% in other words: scale x SCALE x default 
  \expandafter\gdef\csname @reset@#1\endcsname{%
    \expandafter\setlength\csname @@#1@@default\endcsname{%
      \csname @#1@SCALE\endcsname \csname @#1@default\endcsname
      }
    \expandafter\setlength\csname @@#1@@stretch\endcsname{%
      \csname @#1@SCALE\endcsname \csname @#1@stretch\endcsname
      }
    \expandafter\setlength\csname @@#1@@shrink\endcsname{%
      \csname @#1@SCALE\endcsname \csname @#1@shrink\endcsname
      }
    \expandafter\setlength\csname #1\endcsname{%
      \csname @#1@scale\endcsname \csname @@#1@@default\endcsname
      \@plus
      \csname @#1@scale\endcsname \csname @@#1@@stretch\endcsname
      \@minus
      \csname @#1@scale\endcsname \csname @@#1@@shrink\endcsname
      }}
  %%
  %% User command \NAMEScale should expand to 
  %% {\renewcommand\@NAME@scale{#1}\@reset@NAME}
  \expandafter\gdef\csname #1Scale\endcsname##1{%
    \expandafter\gdef\csname @#1@scale\endcsname{##1}%
    \csname @reset@#1\endcsname}
  %%
  %% Package command \NAME@set@SCALE should expand to 
  %% {\renewcommand\@NAME@SCALE{#1}\@reset@NAME}
  \expandafter\gdef\csname #1@set@SCALE\endcsname##1{%
    \expandafter\gdef\csname @#1@SCALE\endcsname{##1}%
    \csname @reset@#1\endcsname}
  %%
  %% initialize at default values with \@reset@NAME
  \csname @reset@#1\endcsname
  }
%    \end{macrocode}
% \end{macro}
% \begin{macro}{\parskip}
% First use of this new infrastructure: set interparagraph space to be
% zero but with a small amount of stretchiness.
%    \begin{macrocode}
%   \ifthenelse{\boolean{@manuscript}}{%
%     \@set@scalable@length{parskip}{0.5\baselineskip}{0.3\p@}{0.1\p@}
%   }{
    \@set@scalable@length{parskip}{0\p@}{0.3\p@}{0\p@}
%   }
%    \end{macrocode}
% \end{macro}
% 
% \begin{macro}{\abovesectionskip}
% \changes{4.0}{2019/11/12}{Use negative values for skip so that first paragraph after section headings is not indented.  This is something that I have been trying to fix for years!  It turns out that the anser was in the documentation for the standard command \cs{@startsection.}}
% \begin{macro}{\belowsectionskip}
% \begin{macro}{\aboveparaskip}
% \begin{macro}{\belowparaskip}
% Define some utility rubber lengths for use in
% section~\ref{sec:section-commands}. 
%    \begin{macrocode}
\newlength{\abovesectionskip}
\newlength{\belowsectionskip}
\newlength{\aboveparaskip}
\newlength{\belowparaskip}
%    \end{macrocode}
% The default values for these lengths are those traditionally used
% for the onecolumn style. 
%    \begin{macrocode}
\@set@scalable@length{abovesectionskip}{-3.5ex}{-4ex}{-2ex}
\@set@scalable@length{belowsectionskip}{2.3ex}{0ex}{2ex}
\@set@scalable@length{aboveparaskip}{3.25ex}{2ex}{0.5ex}
\@set@scalable@length{belowparaskip}{-1em}{\z@}{\z@}
%    \end{macrocode}
% \end{macro}\end{macro}\end{macro}\end{macro}
% We also make scalable some of the standard float parameters. These
% are the same defaults as in the standard article classes, but
% slightly more stretchy. For the twocolumn style, they are currently
% halved, see below. Note that I have also lumped |\abovecaptionskip|
% in with these, since it plays a similar role. There is also a
% |\belowcaptionskip| but it is set to zero in the standard classes. 
%    \begin{macrocode}
\@set@scalable@length{floatsep}{12\p@}{4\p@}{4\p@}
\@set@scalable@length{textfloatsep}{20\p@}{4\p@}{8\p@}
\@set@scalable@length{intextsep}{12\p@}{4\p@}{4\p@}
\@set@scalable@length{dblfloatsep}{12\p@}{4\p@}{4\p@}
\@set@scalable@length{dbltextfloatsep}{20\p@}{4\p@}{8\p@}
\@set@scalable@length{abovecaptionskip}{10\p@}{\z@}{\z@}
%    \end{macrocode}
% \begin{macro}{\@before@title@skip}         
% \begin{macro}{\@before@author@skip}        
% \begin{macro}{\@between@author@skip}       
% \begin{macro}{\@before@abstract@skip}      
% \begin{macro}{\@after@abstracthead@skip}   
% \begin{macro}{\@between@abstract@skip}     
% \begin{macro}{\@before@kwd@skip}           
% \begin{macro}{\@before@dates@skip}           
% Spacing parameters used in |\@maketitle|. These are not changed by
% default except by the user, so they don't need to go into
% |\@package@default@lengths|.  
%    \begin{macrocode}
\newlength{\@before@title@skip}         
\newlength{\@before@author@skip}        
\newlength{\@between@author@skip}       
\newlength{\@before@abstract@skip}      
\newlength{\@after@abstracthead@skip}   
\newlength{\@between@abstract@skip}     
\newlength{\@before@kwd@skip}           
\newlength{\@before@dates@skip}           
\@set@scalable@length{@before@title@skip}{2em}{\z@}{\z@}
\@set@scalable@length{@before@author@skip}{1.5em}{\z@}{\z@}
\@set@scalable@length{@between@author@skip}{0.5em}{\z@}{\z@}
\@set@scalable@length{@before@abstract@skip}{3ex}{1ex}{0.5ex}
\@set@scalable@length{@after@abstracthead@skip}{1ex}{\z@}{\z@}
\@set@scalable@length{@between@abstract@skip}{2ex}{1ex}{0.5ex}
\@set@scalable@length{@before@kwd@skip}{1ex}{0.5ex}{0.2ex}
\@set@scalable@length{@before@dates@skip}{2ex}{0.5ex}{0.2ex}
%    \end{macrocode}
% \end{macro}\end{macro}\end{macro}\end{macro}
% \end{macro}\end{macro}\end{macro}\end{macro}
% 
% Spacing above and below displayed equations also wants to be done
% like this, but it is more complicated since the standard article
% classes make these change with the fontsize. Something to do later,
% maybe. 
% 
%    \begin{macrocode}
%    \end{macrocode}
% \begin{macro}{\@package@default@lengths}
% The value actually used for these lengths is controlled by the
% \meta{SCALE} parameter, which depends on the style option that is in
% effect. This is set using the |\@package@default@lengths| command,
% which must be called every time the style changes (e.g., changing
% from onecolumn to twocolumn). 
%    \begin{macrocode}
\newcommand{\@package@default@lengths}{%
  \ifthenelse{\boolean{@rmtwocol}}{%
    %% twocolumn mode
    \abovesectionskip@set@SCALE{0.5}
    \belowsectionskip@set@SCALE{0.5}
    \aboveparaskip@set@SCALE{0.5}
    \belowparaskip@set@SCALE{1}
    \floatsep@set@SCALE{0.5}
    \textfloatsep@set@SCALE{0.5}
    \intextsep@set@SCALE{0.5}
    \dblfloatsep@set@SCALE{0.5}
    \dbltextfloatsep@set@SCALE{0.5}
    }{%
    %% onecolumn mode
    \abovesectionskip@set@SCALE{1}
    \belowsectionskip@set@SCALE{1}
    \aboveparaskip@set@SCALE{1}
    \belowparaskip@set@SCALE{1}
    \floatsep@set@SCALE{1}
    \textfloatsep@set@SCALE{1}
    \intextsep@set@SCALE{1}
    \dblfloatsep@set@SCALE{1}
    \dbltextfloatsep@set@SCALE{1}
    }%
  }
\@package@default@lengths
%    \end{macrocode}
% \end{macro}
% \begin{macro}{\RescaleLengths}
% \begin{macro}{\RescaleSecLengths}
% \begin{macro}{\RescaleFloatLengths}
% \begin{macro}{\RescaleTitleLengths}
% User commands that will fiddle with all the scalable lengths. 
%    \begin{macrocode}
\newcommand{\RescaleLengths}[1]{%
  \RescaleSecLengths{#1}\RescaleFloatLengths{#1}\RescaleTitleLengths{#1}}
\newcommand{\RescaleSecLengths}[1]{%
%% sections
  \abovesectionskipScale{#1}
  \belowsectionskipScale{#1}
  \aboveparaskipScale{#1}
  }
\newcommand{\RescaleFloatLengths}[1]{%
%% floats
  \floatsepScale{#1}
  \textfloatsepScale{#1}
  \intextsepScale{#1}
  \dblfloatsepScale{#1}
  \dbltextfloatsepScale{#1}
  \abovecaptionskipScale{#1}
  }
\newcommand{\RescaleTitleLengths}[1]{%
%% maketitle
  \@before@title@skipScale{#1}           
  \@before@author@skipScale{#1}        
  \@between@author@skipScale{#1}       
  \@before@abstract@skipScale{#1}      
  \@after@abstracthead@skipScale{#1}   
  \@between@abstract@skipScale{#1}     
  \@before@kwd@skipScale{#1}           
  \@before@dates@skipScale{#1}           
  }
%    \end{macrocode}
% \end{macro}\end{macro}\end{macro}\end{macro}
% \changes{2.2beta14}{2002/03/17}{Finished implementation of new
% infrastructure for scalable lengths. Converted the above/below
% section skips to use it. Also, all the floatsep parameters. Also,
% all the vertical lengths that are used by |\@maketitle|. }

% \begin{macro}{\MSextramargin}
% For the \texttt{manuscript} option we have an extra 0.3\,inch
% vertical margin and 1\,inch horizontal margin for marginal notes (by
% hand or with |\marginpar|. 
%
% \changes{3.23}{2007/09/09}{Reduce baselinestretch in
%   \texttt{manuscript} style to 1.2. No longer use raggedright for
%   manuscript.}
%
% \changes{3.24}{2007/09/14}{Further reduction of baselinestretch in
%   \texttt{manuscript} style to 1.1 since we now have wider margins.}
%
% \changes{3.24}{2007/09/14}{Split \texttt{manuscript}'s extra margin
%   into separate horizontal and vertical parts.}
% 
%    \begin{macrocode}
\newlength{\MSextraHmargin}\newlength{\MSextraVmargin}
\setlength{\MSextraHmargin}{1in} 
\setlength{\MSextraVmargin}{0.3in} 
\if@manuscript
   \renewcommand{\baselinestretch}{1.1}
   \addtolength{\textwidth}{-2\MSextraHmargin}
   \addtolength{\oddsidemargin}{\MSextraHmargin}
   \setlength{\evensidemargin}{\oddsidemargin}
   \addtolength{\textheight}{-2\MSextraVmargin}
   \addtolength{\topmargin}{\MSextraVmargin}
%    \raggedright
\fi
%    \end{macrocode}
% \end{macro}
% Float parameters: don't allow bottom floats by default. 
%    \begin{macrocode}
\renewcommand{\fps@figure}{tp}
\renewcommand{\fps@table}{tp}
%    \end{macrocode}
%
% \begin{macro}{\rmaa@AuthorWidthFrac}
% \begin{macro}{\rmaa@AbstractWidthFrac}
% These parameters control what fraction of the total text width can
% be used by the list of authors and the abstract/resumen,
% respectively.  
%    \begin{macrocode}
\newcommand{\rmaa@AuthorWidthFrac}{1.0}
\newcommand{\rmaa@AbstractWidthFrac}{0.75}
%    \end{macrocode}
% Use a wide abstract to save space in the new style.
%    \begin{macrocode}
\if@newstyle\renewcommand{\rmaa@AbstractWidthFrac}{1.0}\fi
%    \end{macrocode}
% \end{macro}
% \end{macro}
%
% Some things to consider in future versions. Trying to avoid too many
% widows, etc. Apparently, the following will work best with
% \cs{flushbottom}. 
%    \begin{macrocode}
%%\widowpenalty=10000
%%\clubpenalty=10000
%    \end{macrocode}
% Also, should perhaps try getting rid of underfull/overfull box warnings. 
% \changes{3.10}{2003/01/24}{Set large values for \cs{hbadness},
% \cs{vbadness}, \cs{hfuzz}, \cs{vfuzz}. To reduce amount of stuff
% that gets written to the terminal every time you compile something.}
%    \begin{macrocode}
\hbadness=10000
\vbadness=10000
\setlength\hfuzz{\maxdimen}
\setlength\vfuzz{\maxdimen}
%    \end{macrocode}
% \subsection{Various text styles for different contexts}
% \label{sec:text-styles} 
% \begin{macro}{\rmaa@AffilStyle}
% \begin{macro}{\rmaa@PagesStyle}
% \begin{macro}{\rmaa@YearStyle}
% \begin{macro}{\rmaa@VolumeStyle}
% \begin{macro}{\rmaa@VolumenStyle}
% \begin{macro}{\rmaa@JournalNameStyle}
% \begin{macro}{\rmaa@ReceivedStyle}
% \begin{macro}{\rmaa@AbstractHeaderStyle}
% \begin{macro}{\rmaa@JournalNameStyle}
% Style commands for the different fields used by |\maketitle| in
% constructing the title page of the article. These have the form
% |\rmaa@XYZStyle{TEXT}|. I'm not sure these were really a good idea,
% but never mind \dots 
%    \begin{macrocode}
\newcommand{\rmaa@AffilStyle}[1]{{\normalsize\rmfamily#1}}
\newcommand{\rmaa@PagesStyle}[1]{{\normalsize\rmfamily#1}}
\newcommand{\rmaa@YearStyle}[1]{{\normalsize\rmfamily#1}}
\newcommand{\rmaa@VolumeStyle}[1]{{\normalsize\rmfamily\bfseries#1}}
\newcommand{\rmaa@VolumenStyle}[1]{{\Large\sffamily\bfseries#1}}
\newcommand{\rmaa@JournalNameStyle}[1]{{\normalsize\itshape#1}}
\newcommand{\rmaa@ReceivedStyle}[1]{{\normalsize\rmfamily\itshape#1}}
\newcommand{\rmaa@AbstractHeaderStyle}[1]{%
  {\normalsize\rmfamily\MakeUpperCase{#1}}}
\newcommand{\rmaa@TitleStyle}[1]{%
  {\large\bfseries\rmfamily\boldmath\MakeUpperCase{#1}\par}}
%    \end{macrocode}
% \end{macro}
% \end{macro}
% \end{macro}
% \end{macro}
% \end{macro}
% \end{macro}
% \end{macro}
% \end{macro}
% \end{macro}
%
% \begin{macro}{\rmaa@FrontSecStyle}
% Headings of ``front sections'' (see section~\ref{sec:front-sections})
%    \begin{macrocode}
\newcommand{\rmaa@FrontSecStyle}[1]{%
  {\normalsize\bfseries\rmfamily\MakeUpperCase{#1}}}
%    \end{macrocode}
% \end{macro}
%
% \begin{macro}{\rmaa@TOCTitleStyle}
% \begin{macro}{\rmaa@TOCPartStyle}
% \begin{macro}{\rmaa@IndexTitleStyle}
% \begin{macro}{\rmaaTOCAuthorStyle}
% \begin{macro}{\rmaaIndexAuthorStyle}
% Commands to set style of TOC and index entries.
% Rewritten to avoid unnecessary font commands in the .toc and .ind files
%    \begin{macrocode}
\newcommand{\rmaa@TOCTitleStyle}[1]{\MakeUpperCase{#1}}
\newcommand{\rmaa@TOCPartStyle}[1]{\rmaa@TOCTitleStyle{#1}}
\newcommand{\rmaa@IndexTitleStyle}[1]{#1}
\newcommand{\rmaaTOCAuthorStyle}[1]{%
  {\itshape#1}}
\newcommand{\rmaaIndexAuthorStyle}[1]{%
  \protect\mbox{\bfseries#1}}
%    \end{macrocode}
% \end{macro}
% \end{macro}
% \end{macro}
% \end{macro}
% \end{macro}
%
% \begin{macro}{\rmaa@KeywordStyle}
% \begin{macro}{\rmaa@KeywordHeaderStyle}
% The style of the keywords depends on whether we are using the
% \texttt{newstyle} option or not. If we are, then they are small. 
%    \begin{macrocode}
\ifthenelse{\boolean{@newstyle}}{
  \newcommand{\rmaa@KeywordStyle}[1]{%
    {\small%
      \if@boldkeywords 
      \bfseries\MakeUpperCase{\raggedright #1}
      \else
      \raggedright #1
      \fi
    }}
  \newcommand{\rmaa@KeywordHeaderStyle}[1]{%
    {\small\itshape#1}}}
{
  \newcommand{\rmaa@KeywordStyle}[1]{%
    {\normalsize%
      \if@boldkeywords\bfseries\fi\raggedright #1}}
  \newcommand{\rmaa@KeywordHeaderStyle}[1]{%
    {\normalsize\itshape#1}}}
%    \end{macrocode}
% \end{macro}
% \end{macro}
%
% \subsection{Page styles}\label{sec:page-styles}
% \begin{macro}{\ps@rmaa}
% The running headers for normal pages:
%    \begin{macrocode}
\newcommand{\ps@rmaa}{%
%    \end{macrocode}
% nothing in either foot; 
%    \begin{macrocode}
  \renewcommand\@oddfoot{\relax}%
  \renewcommand\@evenfoot{\relax}%
%    \end{macrocode}
% page number and short author list in the even heads; 
%    \begin{macrocode}
  \renewcommand\@evenhead{\rlap{\thepage}\hfil\@shortauthor\hfil}%
%    \end{macrocode}
% short title and page number in the odd heads.
%    \begin{macrocode}
  \renewcommand\@oddhead{\hfil\@shorttitle\hfil\llap{\thepage}}
  }
%    \end{macrocode}
% \begin{macro}{\ps@rmaafirstpage}
% The header of the first page of each
% article:
%    \begin{macrocode}
\newcommand{\ps@rmaafirstpage}{%
  \renewcommand\@oddfoot{\hfil\thepage}%
%    \end{macrocode}
% footer corrected (06 Jun 2000) to have page on outside edge;
%    \begin{macrocode}
  \renewcommand\@evenfoot{\thepage\hfil}%
%    \end{macrocode}
% Header has journal name, volume, and page numbers for a regular
% article, or similar for special cases. 
% \changes{4.0}{2019/11/03}{Change to three-line header with copyright info and DOI.}
%    \begin{macrocode}
\if@preprint
\renewcommand\@oddhead{%
  \if@proceedings 
  \textsc{To appear in ``\@conftitle{} (\rmaa@YearStyle{\@year})''}\hfil
  \textit{RMxAC}
  \else
  \textsc{To appear in} \@journalname\hfil
  \fi}\fi%
\if@guide
\renewcommand\@oddhead{\textsc{Documentation} (\rmaa@YearStyle{\@year})\hfil}
\fi%
\if@manuscript
\renewcommand\@oddhead{Manuscript for \@journalname{} (\rmaa@YearStyle{\@year})\hfil}
\fi%
\if@journal
\renewcommand\@oddhead{%
  \parbox{\textwidth}{
    \@journalname, \rmaa@VolumeStyle{\@volume},
    \rmaa@PagesStyle{\@artpages}\
    (\rmaa@YearStyle{\@year})\hfil\\[1pt]
    {\small\textcopyright{} \@year:
      Instituto de Astronom\'\i{}a,
      Universidad Nacional Aut\'onoma de M\'exico}\hfil\\[1pt]
    {\small\@DOI@url}\hfil
  }%
}\fi%
  \renewcommand\@evenhead{\@oddhead}}
%    \end{macrocode}
% By default, use the |\ps@rmaa| page style.
%    \begin{macrocode}
\pagestyle{rmaa}
%    \end{macrocode}
% \end{macro}
% \end{macro}
%
% \subsection{Preamble commands}\label{sec:preamble} 
% Everything you need for setting the title page, TOC and index
% entries, etc.
% \subsubsection{Commands for the title page}
% \label{sec:title-commands}
% \begin{macro}{\resumen}
% \begin{macro}{\@resumen}
% \begin{macro}{\abstract}
% \begin{macro}{\@abstract}
% User commands for setting the resumen and abstract. 
% These should be used in the document header. The internal commands
% |\@resumen| and |\@abstract| are used to hold the actual text
% pending output via the |\maketitle| command. 
%    \begin{macrocode}
\newif\if@setresumen
\newcommand{\resumen}[1]{%
%%  \renewcommand{\@resumen}{{\selectlanguage{spanish}#1}}%
  \renewcommand{\@resumen}{{\ignorespaces#1}}%
  \@setresumentrue}
\renewcommand{\abstract}[1]{%
  \renewcommand{\@abstract}{{\ignorespaces#1}}
  \renewcommand\ADS@abstract{#1}
%% Not a good idea, since then author macros are not expanded and
%% cause errors later. However, by not doing it we are vulnerable to
%% fragile commands.....
  \if@noexpandADS\@onelevel@sanitize{\ADS@abstract}\fi
  \@abstractpresenttrue
  }
\newcommand{\ADSnoexpand}{\@noexpandADStrue}
%    \end{macrocode}
% \begin{macro}{\boldabstract}
% This is the version for doing 1-page conference proceedings. Still
% save the text for ADS but then write it out immediately as a bold
% paragraph. 
%    \begin{macrocode}
\newcommand{\boldabstract}[1]{%
  \gdef\ADS@abstract{#1}%
  \renewcommand{\@abstract}{#1}%
  \if@noexpandADS\@onelevel@sanitize{\ADS@abstract}\fi%
  \noindent\textbf{\boldmath #1}
  }
%    \end{macrocode}
% \end{macro}
% Default value in each case is to print out a warning if not set
%    \begin{macrocode}
  \newcommand{\@resumen}{%
    \@setresumenfalse \ClassWarningNoLine{rmaa}{No Spanish abstract
      (resumen) specified - if you do not speak Spanish, this will be
      translated by the editors}%
    \textbf{Favor de proporcionar un resumen en espa\~nol. If you are
      unable to translate your abstract into Spanish, the editors will
      do it for you.} }
  \newcommand{\@abstract}{\ClassWarningNoLine{rmaa}{No abstract
      specified}}
%    \end{macrocode}
% \end{macro}
% \end{macro}
% \end{macro}
% \end{macro}
%
% \begin{macro}{\title}
% \begin{macro}{\@rawtitle}
% \begin{macro}{\@TOCtitle}
% \begin{macro}{\rmaaIndexTitle}
% Similarly, |\title| is the preamble command to set the title of the
% article. This is used in various places, typeset in different
% styles, hence the use of the internal |\@rawtitle|, |\@title|,
% |\rmaaIndexTitle|, and |\@TOCtitle| commands, each of which uses the
% relevant style from section~\ref{sec:text-styles}
%    \begin{macrocode}
\newcommand{\@rawtitle}{Raw Title}
\newcommand{\@TOCtitle}{\rmaa@TOCTitleStyle{Title for TOC}}
\newcommand{\rmaaIndexTitle}{\rmaa@IndexTitleStyle{Title for index}.}
\newcommand{\@TOCaftertitle}{}
\renewcommand{\title}[1]{%
%%  {\ADS@sanitize\xdef\@rawtitle{#1}}
  \renewcommand\@rawtitle{\noexpand#1}
  \let\ADS@title=\@rawtitle
  \renewcommand{\@title}{\rmaa@TitleStyle{#1}}%
  \renewcommand{\@TOCtitle}{\rmaa@TOCTitleStyle{#1}\@TOCaftertitle}%
  \renewcommand{\rmaaIndexTitle}{\rmaa@IndexTitleStyle{#1\@TOCaftertitle}.}%
  }
%    \end{macrocode}
% \end{macro}
% \end{macro}
% \end{macro}
% \end{macro}
% 
% All the weird and whacky ways of indicating the author
% affiliations. 
% \begin{macro}{\affil}
% First, the way that sets it directly underneath the authors name.
% Note that |\affil| trumps all other affiliations for the purposes of
% getting into the ADS stuff (we use |\gdef|).   
% \dots 
%    \begin{macrocode}
\newcommand{\affil}[1]{%
  \\[1.5ex]\rmaa@AffilStyle{#1}\\[1.5ex]%
  \if@firstaffil\gdef\ADS@affil{#1}\@firstaffilfalse\fi%
  }
\renewcommand{\and}{and\\[1.5ex]}
%    \end{macrocode}
% \begin{macro}{\altaffil}
% Alternative affiliations are footnoted. 
%    \begin{macrocode}
\newcommand{\altaffil}[1]{%
  \if@firstaltaffil\gdef\ADS@affil{#1}\@firstaltaffilfalse\fi%
  \thanks{#1}}
%    \end{macrocode}
% \begin{macro}{\altaffilmark}
% \begin{macro}{\altaffiltext}
% Alternative-alternative mechanism for affiliations---should not be
% mixed with |\altaffil| (unless you know what you are doing) or there
% will be problems with numbering.
%    \begin{macrocode}
\newcommand{\altaffilmark}[1]{\textsuperscript{#1}}
\newcommand{\altaffiltext}[2]{%
  \if@firstaltaffil\gdef\ADS@affil{#2}\@firstaltaffilfalse\fi
  \stepcounter{footnote}%
  \protected@xdef\@thanks{\@thanks
    \protect\footnotetext[#1]{#2}}} 
%    \end{macrocode}
% \end{macro}
% \end{macro}
% \end{macro}
% \end{macro}
%
% \begin{macro}{\keywords}
% \begin{macro}{\@keywords}
% \begin{macro}{\@keywordheader}
% \label{keywords}Sundry minor preamble commands. All
% self-explanatory, I hope. |\keywords| is the old user interface
% (deprecated). New articles should use |\addkeyword| once for each
% keyword. That way the ADS record comes out right too (comma rather
% than em-dash separator).
% \changes{2.2beta12}{2001/12/07}{\cs{addkeyword} New user interface
%   for keywords.} 
% \changes{2.2beta13}{2002/03/11}{\cs{addkeyword} Improved behaviour
%   with .ads file. Use
%   \cs{ADS@sanitize} to change definitions of `~' and \cs{ion}.}
% \changes{3.23}{2007/09/09}{\cs{addkeyword} now signals a compilation
% error if the author tries to add more than 6 keywords.}
% \changes{3.24}{2007/09/14}{kwdcounter is now reset in
% \cs{@begofarticlehook} to prevent error in book style} 
%    \begin{macrocode}
\newcommand{\@kwdsep}{ --- }
\newcommand{\ADS@kwdsep}{, }
\newcounter{kwdcounter}\setcounter{kwdcounter}{0}
\newcommand{\addkeyword}[1]{%
  \addtocounter{kwdcounter}{1}
  \ifthenelse{\thekwdcounter > 6}{
    \ClassError{rmaa}{Too many keywords. A maximum of 6 keywords is allowed.}{}
  }{}
  \if@firstkwd%
  \renewcommand\@keywords{\if@boldkeywords\uppercase{#1}\else #1\fi}
  {\ADS@sanitize\def\ion{\plainion}\gdef\ADS@keywords{#1}}
  \@firstkwdfalse
  \else%
  \edef\@keywords{\@keywords\@kwdsep \if@boldkeywords\uppercase{#1}\else #1\fi}%
  {\ADS@sanitize\def\ion{\plainion}\xdef\ADS@keywords{\ADS@keywords\ADS@kwdsep #1}}%
  \fi}
\newcommand{\ADS@sanitize}{%
  \def\mathrm{}
  \def\@{\ }
  \def~{\ }
  \def\langle{<}
  \def\rangle{>}
  \def\log{log }
  \def\sin{sin }
  \def\cos{cos }
  \def\tan{tan }
  \def\fd{^d.}
  \def\fh{^h.}
  \def\fm{^m.}
  \def\fs{^s.}
  \def\fdg{�.}
  \def\farcm{'.}
  \def\farcs{''.}
  \def\arcdeg{�}
  \def\,{ }
  }
\newcommand{\keywords}[1]{%
  \renewcommand{\ADS@keywords}{#1}
  \renewcommand{\@keywords}{\uppercase{#1}}
  }
\newcommand{\@keywordheader}{\rmaa@KeywordHeaderStyle{Key Words:}}
\newcommand{\@keywords}{%
  \ClassWarningNoLine{rmaa}{No keywords specified}%
  Keywords go here
  }
%    \end{macrocode}
% \begin{macro}{\shorttitle}
% \begin{macro}{\shortauthor}
% \begin{macro}{\@shorttitle}
% \begin{macro}{\@shortauthor}
% Commands for setting the running headers.
%    \begin{macrocode}
\newcommand{\shorttitle}[1]{%
  \renewcommand\@shorttitle{\MakeUpperCase{#1}}}
\newcommand{\shortauthor}[1]{%
  \renewcommand\@shortauthor{\MakeUpperCase{#1}}}
\newcommand{\@shorttitle}{SHORT TITLE GOES HERE}
\newcommand{\@shortauthor}{SHORT AUTHOR GOES HERE}
%    \end{macrocode}
% \end{macro}
% \end{macro}
% \end{macro}
% \end{macro}
% \end{macro}
% \end{macro}
% \end{macro}


% \subsubsection{The full addresses}
% \label{sec:full-addresses}
% \begin{macro}{\fulladdresses}
% \begin{macro}{\@fulladdresses}
% \begin{macro}{\@fulladdheight}
% \begin{macro}{\@fulladddepth}
% \begin{macro}{\@fulladdbox}
% Aaghh, the full addresses! A perennial source of problems. 
%    \begin{macrocode}
\newlength{\@fulladdheight}
\newlength{\@fulladddepth}
\newsavebox{\@fulladdbox}
%% \fulladdresses command modified (11 May 2000) so that we can put
%% \fulladdresses{} in the article to suppress the output of the
%% addresses
\newcommand{\@FinalFigure}{\relax}
\newcommand{\FinalFigure}[1]{\renewcommand{\@FinalFigure}{#1}}
\newcommand{\fulladdresses}[1]{%
  \renewcommand{\@rmaaArg}{#1}
  \ifx\@rmaaArg\@rmaaEmpty
     \renewcommand{\@fulladdresses}{}
  \else
     \renewcommand{\@fulladdresses}{
       \normalsize%
       \ifx\@FinalFigure\relax\else\vfill\@FinalFigure\vfill\fi
%%       \vspace*{1ex}
       %% do addresses as a list (copied from thebibliography environment) 
       \list{\null}{\leftmargin 1.5em\topsep\z@\labelwidth\z@\labelsep\z@\itemindent -1.5em
         \itemsep\z@\parsep\z@\usecounter{enumi}}
       \sloppy\clubpenalty4000\widowpenalty4000
       \sfcode`\.=1000\relax
       #1
       \endlist
       }
  \fi
  }
\newcommand{\@fulladdresses}{Full addresses go here}
%    \end{macrocode}
% \begin{macro}{\suppressfulladdresses}
% New user-level command for suppressing full addresses should you so
% desire. Currently used in the single-page-poster style. 
%    \begin{macrocode}
\newcommand{\suppressfulladdresses}{\@suppressfulladdressestrue\@needaddressesfalse}
%    \end{macrocode}
% \end{macro}
% \end{macro}
% \end{macro}
% \end{macro}
% \end{macro}
% \end{macro}

% \begin{macro}{\printdates}
% \begin{macro}{\noprintdates}
% Command to activate the printing out of the received/accepted
% dates. This can also be done using the \texttt{printdates} option,
% which is the default for \texttt{journal}, but not for \texttt{proceedings}.
%    \begin{macrocode}
\newcommand{\printdates}{\@printdatestrue}
\newcommand{\noprintdates}{\@printdatesfalse}
%    \end{macrocode}
% \end{macro}
% \end{macro}
% \begin{macro}{\@outputaddressesifnecessary}
% Command to output the addresses if we haven't already done so
% This is called by: 
% \begin{enumerate}
% \item The end of |thebibliography| environment
% \item The |\outputfulladdresses| command
% \end{enumerate}
%    \begin{macrocode}
\newcommand{\@outputaddressesifnecessary}{%
  \ClassInfo{rmaa}{Calling @outputaddressesifnecessary}
  \ifthenelse{\boolean{@needaddresses}}{%
    %% If in one column style, or in the first column of a two column
    %% style then we can try to put it on this page 
    \ifthenelse{\boolean{@firstcolumn} \or \not\boolean{@rmtwocol}}{%
      \@outputaddresses}{%
      %% otherwise give warning it will be on a separate page
      \ClassWarning{rmaa}{Full addresses will go on a separate
        page (please see User Guide for how to remedy this)}}%
    %% If we haven't managed yet, put it on its own page
    \ifthenelse{\boolean{@needaddresses}}{%
      \if@debug{\ClassWarning{rmaa}{In this silly bit!}}\fi
      \clearpage\vfill\@outputaddresses}{}%
    }{
    %% If we don't need them, we still might need to do the eoa hook 
    \ifthenelse{\boolean{@suppressfulladdresses}}{%
      \if@debug{\ClassWarning{rmaa}{XXX}}\fi
%%      \if@book\else\@endofarticlehook\null\fi
      }{}
    }%
  }  
%    \end{macrocode}
% \end{macro}

% \begin{macro}{\outputfulladdresses}
% A user interface to the above command, in case there is no biblio
% (11 May 2000) 
%    \begin{macrocode}
\newcommand{\outputfulladdresses}{\@outputaddressesifnecessary}
%    \end{macrocode}
% \end{macro}

% \begin{macro}{\adjustfinalcols}
% This command starts a new column and should be 
% used to ensure that the columns balance on the last page. It also
% outputs the full 
% author addresses to ensure that this occurs in the left hand
% column. 
% The |\par| command is so that the current line will be taken account
% of when working out |\pagetotal|
%    \begin{macrocode}
\newcommand{\adjustfinalcols}{\if@twocolumn\goodbreak\@outputaddresses\newpage\noindent\fi}
%    \end{macrocode}
% \end{macro}

% \begin{macro}{\@outputaddresses}
% Command to do the work of outputting the addresses. 
% This can be called by 
% \begin{enumerate}
% \item an explicit |\adjustfinalcols|
% \item via |\@outputaddressesifnecessary|, q.v.
% \end{enumerate}
%
%    \begin{macrocode}
\newcommand{\@outputaddresses}{%
  % Allow option of suppressing the full addresses, in which case
  % we just run the hook if we are not in a book style
  \if@suppressfulladdresses
%%     \if@book\else\@endofarticlehook\null\fi
  \else
  %% save the full addresses in a box for output later
  \sbox{\@fulladdbox}{\parbox{\textwidth}{\@fulladdresses}}
  % calculate height of address box  
  \settoheight{\@fulladdheight}{\usebox{\@fulladdbox}}
  \settodepth{\@fulladddepth}{\usebox{\@fulladdbox}}
  \addtolength{\@fulladdheight}{\@fulladddepth}
  \if@debug
  { \typeout{Address total height: \the\@fulladdheight}
    \typeout{Text height: \the\pagetotal}
    \typeout{Goal height: \the\pagegoal}
    }
  \fi
  % reuse \@fulladdheight to mean the total height of text+addresses
  \addtolength{\@fulladdheight}{\pagetotal}
%%  \addtolength{\@fulladdheight}{-\baselineskip}
  \if@debug\typeout{Adjusted text + addresses height: \the\@fulladdheight}\fi
  % check it will fit on current page
  \ifdim\@fulladdheight<\pagegoal
     %% must shift it to the left
     %% to compensate for the hanging indentation in the reference list. 
     %% NOTE: this is not foolproof!
     %% NB. This will give overfull hbox warnings, which is a bit ugly.
     \vfill\noindent\if@biblist\hspace*{-\leftmargin}\fi
     \usebox{\@fulladdbox}
     % Set this to show we've done it
     \@needaddressesfalse
     % and assume we are at the end of the article
     % so that we do the EOA hook if not in the book style
%%     \if@book\else\@endofarticlehook\fi
  \else
     % if not, just do nothing. With luck, this function will get 
     % called again on the next page
     \if@debug{\typeout{Holding over full addresses to next page}}\fi
  \fi\fi
  }
%    \end{macrocode}
% \end{macro}

% \subsubsection{Commands for editors}
% \label{sec:editor-commands}
% Some user/editor commands for doing the TOC and index in the book
% option. 
% They are defined here rather in the book-specific part since they
% need to be always available in order to avoid breaking documents. 
% \begin{macro}{\listofauthors}
% \begin{macro}{\@rawauthors}
% \begin{macro}{\rmaaTOCauthors}
% The command |\listofauthors| formats an author list for use in the
% TOC. It should contain the same names as in the |\author| command,
% but without the affiliation garbage and using |\&| instead of
% |and|. 
%    \begin{macrocode}
\newcommand{\@rawauthors}{Raw Authors}
\newcommand{\rmaaTOCauthors}{\noexpand\rmaaTOCAuthorStyle{Author list for TOC}}
\newcommand{\listofauthors}[1]{%
  \renewcommand{\@rawauthors}{#1}
  \renewcommand\rmaaTOCauthors{\noexpand\rmaaTOCAuthorStyle{#1}}}
%    \end{macrocode}
% \begin{macro}{\indexauthor}\label{indexauthor}
% The command |\indexauthor| produces an index entry for each
% author. 
% This command must come \emph{after} the
% |\title| and the |\listofauthors| commands. Also, there should be one 
% |\indexauthor| command for \emph{each} author. We also add the
% author onto the end of the |\ADS@authors| macro and construct the
% |\ADS@authorinitial| by a most cunning method (see below), but only
% for the first author. 
% 
% No, No, NO! This was all wrong. Backed out completely: 
% \begin{quote}
% Completely rewritten (19 Dec 2002) since we now use the index class, so 
% as to be able to alphabetize accented characters. A side effect of this is
% that |\index| no longer expands out the argument after the |@|, so we 
% instead add all the |\index| commands to a list (|\indexauthor@deferlist|), 
% which is executed in the |\@begofarticlehook|. 
% 
% Note the stuff with |\toks0| is copied from Exercise 20.15 of the \TeX{}book---it is
% to make sure the latin1-accented name gets passed into the idx
% file. 
% \end{quote}
% We now go back to the old way of doing things, using an optional
% argument to the \cs{indexauthor} command that has the name without
% accents. 
% \changes{3.08}{2003/01/16}{Remerged changes from 3.06 that got
% dropped from 3.07 after a synchronization disaster. Basically, the
% accentless optional argument to \cs{indexauthor}. Must add the macro
% directories to my unison set-up to stop this sort of thing from happening.}
%    \begin{macrocode}
\newcommand\@isortkey{}
\newcommand\@emptykey{empty}
\newcommand\indexauthor[2][empty]{
  \renewcommand\@isortkey{#1}
  \ifx\@isortkey\@emptykey \renewcommand\@isortkey{#2}\fi
  \index{\@isortkey @{{\rmaaIndexAuthorStyle{#2}}\hskip1em plus1em minus0.4em
      \rmaaIndexTitle\hskip1em plus2em minus0.4em\penalty-100
      \rmaaTOCauthors}}
  \indexauthorADS{#2}
}
%%\def\indexauthor@deferlist{\relax}
%%\long\def\indexauthor#1{
%%  \def\@@name{#1} \toks0=\expandafter{\@@name}
%%  \toks1=\expandafter{\indexauthor@deferlist}
%%  \protected@xdef\indexauthor@deferlist{%
%%    \the\toks1
%%    \noexpand\index{\the\toks0@{\noexpand\rmaaIndexAuthorStyle{\the\toks0}\noexpand\hskip1em plus1em minus0.4em 
%%        \rmaaIndexTitle\noexpand\hskip1em plus2em minus0.4em\noexpand\penalty-100
%%        \rmaaTOCauthors}}
%%  }
%%  \indexauthorADS{#1}
%%  }
%%\xdef\indexauthor@deferlist{}
\newcommand\indexauthorADS[1]{%
  \if@firstauthor\protected@xdef\ADS@authorinitial{\firstletter{#1}}\fi
  \protected@xdef\ADS@authors{\ADS@authors\if@firstauthor\else ; \fi
    #1}
  \@firstauthorfalse
}
%    \end{macrocode}
% Note that the first |#1| sets the key for the index, making sure it
% actually starts with  
% first letter of the surname. Meanwhile, the |@{...}| expression,
% which starts with some  
% control sequences is what is written to the |.ind| file. 
% See the makeindex man page for details. 
% \end{macro}
% \end{macro}
% \end{macro}
% \end{macro}

% \begin{macro}{\KillAllAccents}
% \begin{macro}{\firstletter}
% These are utility functions used in the above. |lett| is a counter
% indicating the current letter in the string. |\firstletter| just
% prints the first letter of its argument.
%    \begin{macrocode}
\newcommand{\KillAllAccents}{\def\'{}\def\`{}\def\~{}\def\.{}\def\i{i}}
\newcommand\firstletter[1]{\@car #1\@nil}
%    \end{macrocode}
% \end{macro}
% \end{macro}
% \begin{macro}{\SetVolume}
% \begin{macro}{\SetMSnumber}
% \begin{macro}{\SetFirstPage}
% \begin{macro}{\SetYear}
% \begin{macro}{\ReceivedDate}
% \begin{macro}{\AcceptedDate}
% These are for the editors to set the volume number, etc. 
%    \begin{macrocode}
\newcommand{\SetVolume}[1]{%
  \renewcommand{\@volume}{#1}
  \edef\ADS@volume{\PadToFourWithDots{#1}}
  \renewcommand{\@volumen}{\rmaa@VolumenStyle{volumen #1}}%
  }
\newcommand{\SetMSnumber}[1]{\renewcommand{\@MSnum}{#1}}
%    \end{macrocode}
% Note that there is no real reason to use |\SetFirstPage| since all
% the page numbering is done automatically in the book style. 
%    \begin{macrocode}
\newcommand{\SetFirstPage}[1]{\setcounter{page}{#1}}
\newcommand{\SetYear}[1]{%
  \renewcommand{\@year}{#1}
  \renewcommand{\ADS@year}{#1}
  \renewcommand{\@anyo}{\rmaa@VolumenStyle{#1}}
  }
\newcommand{\ReceivedDate}[1]{\renewcommand{\@received}{#1}}
\newcommand{\AcceptedDate}[1]{\renewcommand{\@accepted}{#1}}
%    \end{macrocode}
%
% \begin{macro}{\DOI}
% User/editor command to enter the DOI, which will be typeset as a link.
% \changes{4.0}{2019/11/02}{New user command \cs{DOI} to typeset
% DOI in the header of first page and add it to ADS info.}
%    \begin{macrocode}
\newcommand\DOI[1]{%
  \def\@DOI@url{\url{https://doi.org/10.22201/ia.01851101p.#1}}%
  \def\ADS@doi{https://doi.org/10.22201/ia.01851101p.#1}%
}%
%    \end{macrocode}
% \end{macro}
%
% \begin{macro}{\ORCIDiD}
% \changes{4.0}{2019/11/12}{Add mechanism for giving ORCID iD links}
% \begin{macro}{\ORCIDiDfull}
% \begin{macro}{\@orcid@icon}
% \begin{macro}{\@orcid@iconbig}
% Commands to add icon links to ORCID iDs of authors. For example
% \begin{quote}
% |\ORCIDiD|\marg{0000-0000-0000-0000}
% \end{quote}
% to be used after each author name in the |\author|
% command and in the |\fulladdresses|.  Alternatively,
% |\OORCIDiDfull| will typeset the icon plus the full text of the
% link, but we don't use that at present.  The icon image is stored in |ORCIDiD_iconvector.pdf|, which needs to be on the input path.  In principle, we could make EPS versions if necessary.
%    \begin{macrocode}
\newcommand\@orcid@icon{\includegraphics[height=1.8ex]{ORCIDiD_iconvector}}
\newcommand\@orcid@iconbig{\includegraphics[height=2ex]{ORCIDiD_iconvector}}
\newcommand\ORCIDiD[1]{\href{https://orcid.org/#1}{\,\@orcid@icon}}
\newcommand\ORCIDiDfull[1]{%
  \href{https://orcid.org/#1}{\@orcid@iconbig}\,\url{https://orcid.org/#1}}
%    \end{macrocode}
% \end{macro}
% \end{macro}
% \end{macro}
% \end{macro}
%
% \begin{macro}{\@volume}
% \begin{macro}{\@volumen}
% \begin{macro}{\@artpages}
% \begin{macro}{\@year}
% \begin{macro}{\@anyo}
% \begin{macro}{\@received}
% \begin{macro}{\@accepted}
% \begin{macro}{\@MSnum}
% Default values for the containers used by above commands, plus a
% few more useful ones. 
%    \begin{macrocode}
\newcommand{\@volume}{00}
\newcommand{\@volumen}{\rmaa@VolumenStyle{00}}
\newcommand{\@artpages}{%
  \oldpageref{\articlename:RMAAFirstPage}--\oldpageref{\articlename:RMAALastPage}}
\newcommand{\@year}{\the\year}
\newcommand{\@anyo}{\rmaa@VolumenStyle{2007}}
\newcommand{\@journalname}{\rmaa@JournalNameStyle{%
    \if@proceedings\RMAASClong\else\RMAAlong\fi}}
\newcommand{\@received}{}
\newcommand{\@accepted}{}
\newcommand{\@MSnum}{0000}
\newcommand{\@DOI@prefix}{https://doi.org/10.22201/ia.01851101p}
\def\@DOI@url{\url{https://doi.org/10.22201/ia.01851101p.2099.00.00.00}}
%    \end{macrocode}
% \end{macro}
% \end{macro}
% \end{macro}
% \end{macro}
% \end{macro}
% \end{macro}
% \end{macro}
% \end{macro}
% \end{macro}
% \end{macro}
% \end{macro}
% \end{macro}
% \end{macro}
% \end{macro}
% 
% The following commands are most useful in the \texttt{book}
% style. However, the |\SetConfTitle| command can now be used in
% individual articles to set the header for the \texttt{preprint}
% style. 
%    \begin{macrocode}
\newcommand{\@reunionplace}{Lugar de la reuni\'on}
\newcommand{\@reuniondate}{Fecha de la reuni\'on}
\newcommand{\@conftitle}{Conference Title}
\newcommand{\SetReunion}[2]{\renewcommand{\@reunionplace}{#1}\renewcommand{\@reuniondate}{#2}}
\newcommand{\SetConfTitle}[1]{\renewcommand{\@conftitle}{#1}}
\newcommand{\@editors}{Juan Ram\'on de la Fuente and ``El Mosh''}
\newcommand{\SetEditors}[1]{\renewcommand{\@editors}{#1}}
\newcommand{\@editores}{Juan Ram\'on de la Fuente y ``El Mosh''}
\newcommand{\SetEditores}[1]{\renewcommand{\@editores}{#1}}
\newcommand{\@mes}{mes de la publicaci\'on}
\newcommand{\SetMes}[1]{\renewcommand{\@mes}{\rmaa@VolumenStyle{#1}}}
%    \end{macrocode}

% \subsection{Notes for Author Proofs and Communication between Editors}
% \label{sec:notes}
% We introduce a command \cs{Note} that should be used to describe
% any change made to the manuscript. These are a bit tricky to get
% right. In general we use \cs{marginpar} but this only works in outer
% paragraph mode, so we just bung the box where we are if we are in a
% figure caption or something. If we are in the preamble, then we just
% use a \marginpar as part of the \maketitle command.
%    \begin{macrocode}
\newif\if@arenotes
\@arenotesfalse
\newcounter{notecounter}
\setcounter{notecounter}{0}
\newcommand\NotesList{}
\newcommand\Note[1]{%
  \if@proof{%
    \@arenotestrue%
    \stepcounter{notecounter}%
    \if@inpreamble\else{%
      \NoteMark{\arabic{notecounter}}%
    }\fi
    \protected@xdef\NotesList{\NotesList 
      \protect\item #1}%
  }\fi%
}
\newcommand\Changed[2]{Changed ``#1'' to ``#2''.}
\newcommand\NoteMark[1]{%
  \ifinner
  \raisebox{1ex}[0pt][0pt]{\makebox[0pt][l]{\framebox{\footnotesize\bfseries
        N#1}}}%
  \else
  \marginpar{\framebox{\footnotesize\bfseries N#1}}%
  \fi
}
\newcommand\RefNote[1]{\if@proof\NoteMark{\ref{#1}}\relax\fi}
\newcommand\NoteLabel{\protect\label}
\newcommand\OutputNotes{%
  \bgroup
  \onecolumn
  \section*{EDITORIAL NOTES}
  \begin{enumerate}
    \NotesList
  \end{enumerate}
  \egroup
}
\DeclareRobustCommand\cs[1]{\texttt{\char`\\#1}}

%    \end{macrocode}
% \subsection{Bibliographic Records for ADS}
% \label{sec:ADS}
% \changes{2.2beta11}{2001/11/27}{Just about finished all the ADS
% stuff now. Except for the references, which I have given up on for
% the time being.}
% \begin{macro}{\ADS@BibCode}
% \begin{macro}{\ADS@year}
% \begin{macro}{\ADS@journal}
% \begin{macro}{\ADS@volume}
% \begin{macro}{\ADS@firstpage}
% \begin{macro}{\ADS@authorinitial}
% We use the |\ADS@BibCode| command to hold the bibliographic code for
% ADS. It is 
% used in the |%R| 
% field of the bibliographic record for each article. 
% The format is 
% \begin{quote}
% \texttt{YYYYJJJJJVVVVMPPPPA}  
% \end{quote}
% where \texttt{YYYY} is the year, \texttt{JJJJJ} is the journal code,
% which is \texttt{RMxAA} or \texttt{RMxAC}, \texttt{VVVV} is the
% volume number (right-justified), \texttt{M} is not used by us,
% \texttt{PPPP} is the first page number (right-justified), and
% \texttt{A} is the initial of the first author's surname. The volume
% and page fields are padded with periods (`\texttt{.}'), which causes
% us some fun (this is done with the macro |\PadToFourWithDots| described
% below) . 
% 
% All these individual bits are held in the commands |\ADS@year|,
% |\ADS@journal|, |\ADS@volume|, |\ADS@firstpage|,
% |\ADS@authorinitial|. 
%    \begin{macrocode}
\newcommand{\SetADSBibCode}[1]{\renewcommand{\ADS@BibCode}{#1}}
\newcommand{\ADS@BibCode}
{\ADS@year\ADS@journal\ADS@volume.\PadToFourWithDots{\ADS@firstpage}\ADS@authorinitial}
\newcommand{\ADS@year}{2001} 
\newcommand{\ADS@journal}{\if@proceedings RMxAC\else RMxAA\fi}
\newcommand{\ADS@volume}{....} 
%    \end{macrocode}
% These two are just set to default values here, |\ADS@firstpage| is set
% in the |document| environment for included articles, in the
% |frontsection| environment, |\tableofcontents| command,  |index| and
% |abstracts| environments. Also by |\AtBeginDocument| for stand-alone
% articles. 
% 
% God knows how we are going to do |\ADS@authorinitial|!
%    \begin{macrocode}
\newcommand{\ADS@firstpage}{....}
\newcommand{\ADS@authorinitial}{A}
%    \end{macrocode}
% \end{macro}
% \end{macro}
% \end{macro}
% \end{macro}
% \end{macro}
% \end{macro}
% We also have a whole slew of other fields that can be set. These are
% all listed in Table~\ref{tab:ADS}. Most of them still have to be
% defined. 
% \begin{table}[!ht]\raggedleft
% \caption{ADS Bibliographic Records}\label{tab:ADS}
% \begin{tabular}{lp{8em}lp{8em}}\hline
% Code & Description & Macro Name & Comment \\ \hline
% |%R| & Bibliographic Code & |\ADS@BibCode| & required \\
% |%A| & Author list & |\ADS@authors| & required, generate from |\indexauthor| \\
% |%F| & Author affiliation & |\ADS@affil| & format: AA(\dots)
% AB(\dots) \dots \\
% |%J| & Journal Name & |\ADS@journalname| &  \\
% |%D| & Publication Date & |\ADS@fulldate| &  \\
% |%L| & Last Page & |\ADS@lastpage| & generate from label \\
% |%T| & Title & |\ADS@title| & required, generate from |\title| \\
% |%C| & Abstract Copyright & |\ADS@copyright| & set to IA-UNAM \\
% |%U| & URL for electronic document & |\ADS@url| & set to PDF file \\
% |%d| & DOI for article & |\ADS@doi| & set by |\DOI| \\
% |%K| & Keywords & |\ADS@keywords| & set equal to |\@keywords| \\
% |%A| & Abstract Text & |\ADS@abstract| & set equal to |\@abstract| \\
% |%R| & References & |\ADS@references| & generate from the
% bibliography  \\ \hline
% \end{tabular}
% \end{table}
% A lot of these macros are set by other commands, as mentioned in the
% Table, but some need to be set  explicitly here. 
% 
% \begin{macro}{\ADS@doi}
% \changes{4.0}{2019/11/08}{Save DOI for ADS}
% The user/editor command \cs{DOI} is used to set this. 
%    \begin{macrocode}
\newcommand{\ADS@doi}{https://doi.org/10.22201/ia.01851101p.0000.00.00.00}
%    \end{macrocode}
% \end{macro}
% 
% \begin{macro}{\ADS@journalname}
% This
% does not contain any accents at the moment so as not to cause
% problems---this maybe should be changed though. 
%    \begin{macrocode}
\newcommand{\ADS@journalname}{%
  \if@proceedings\@conftitle\space (Eds. \@editors) \fi%
  Revista Mexicana de Astronom\'\i{}a y Astrof\'\i{}sica%
  \if@proceedings\space (Serie de Conferencias)\fi%
  \space Vol. \@volume, pp. \ADS@firstpage--\ADS@lastpage\space (\@year)%
  \space (http://www.astroscu.unam.mx/\~{}rmaa/)
  }
%    \end{macrocode}
% \end{macro}
% \begin{macro}{\ADS@keywords}
% \begin{macro}{\ADS@abstract}
% \begin{macro}{\ADS@title}
% |\ADS@abstract| is straightforward, since we already have the text
% saved as |\@abstract|. We could do the same with |\ADS@keywords| but
% to avoid unwanted formatting we explicitly set it in |\keywords|
% (p.~\pageref{keywords}). |\ADS@title| is just set to |\@rawtitle|,
% which we were already using for another purpose. 
%    \begin{macrocode}
\let\ADS@title=\@rawtitle
\newcommand{\ADS@keywords}{No keywords specified}
\newcommand{\ADS@abstract}{No abstract specified}
%    \end{macrocode}
% \end{macro}\end{macro}\end{macro}
% \begin{macro}{\ADS@authors}
% |\ADS@authors| wants to be in ``Other, A. N.'' format, but with
% semi-colons between authors. The most natural way to do this seems to
% be with |\indexauthor| (p.~\pageref{indexauthor}). We append each
% author to this macro, which means we must explicitly clear it for
% each new article (done in |\article|). .   
%    \begin{macrocode}
\newcommand{\ADS@authors}{}
%    \end{macrocode}
% \end{macro}
% \begin{macro}{\ADS@fulldate}
% \begin{macro}{\ADS@lastpage}
% \begin{macro}{\ADS@copyright}
% \begin{macro}{\ADS@website}
% \begin{macro}{\ADS@url}
% \begin{macro}{\SetWebsite}
% \begin{macro}{\SetMonthNumeric}
% Other sundry data, together with user commands for bits that can be
% set by the Editors. The |\ADS@url| is constructed from the website
% (currently |http://www.astrosmo.unam.mx/rmaa/|,
% the journal descriptor (e.g., |RMxAC..10/|), followed by the
% path to the PDF file (e.g., |PDF/RMxAC..10_HEW1.pdf|). 
% 
%    \begin{macrocode}
\newcommand{\ADS@website}{http://www.astrosmo.unam.mx/rmaa/}
\newcommand{\ADS@url}
{\ADS@website\ADS@journal\ADS@volume\ADS@number/PDF/\ADS@journal\ADS@volume\ADS@number\_\articlename.pdf}
\newcommand{\ADS@number}{}
\newcommand{\ADS@fulldate}{00/\ADS@year} 
\newcommand{\ADS@lastpage}{00} 
\newcommand{\ADS@copyright}{(C) \ADS@year; Instituto de Astronom\'\i{}a,
  Universidad Nacional Aut\'onoma de M\'exico} 
\newcommand{\SetNumber}[1]{\renewcommand{\ADS@number}{-#1}}
\newcommand{\SetWebsite}[1]{\renewcommand{\ADS@website}{#1}}
\newcommand{\SetMonthNumeric}[1]{\renewcommand{\ADS@fulldate}{#1/\ADS@year}}
%    \end{macrocode}
% \end{macro}\end{macro}\end{macro}\end{macro}\end{macro}\end{macro}\end{macro}
% \begin{macro}{\ADS@affil}
% \begin{macro}{\ADS@references}
% These two are somewhat more problematical. Affiliation currently
% only works for the first author (see \dots) pending redesign of the
% user interface. References are now done by wrapping the
% |thebibliography| environment to open another input stream on the
% same file and search through for the |\bibitems|. 
%    \begin{macrocode}
\newcommand{\ADS@references}{}
\newcommand{\ADS@affil}{}
%    \end{macrocode}
% \end{macro}\end{macro}
% 
% Define two macros that expand to the strings
% |\begin{thebibliography}| and |\end{thebibliography}|. Note that the
% idea for this and much else is inspired by the code for the docstrip
% package. Also define |\StartBibX| in case authors use the more
% standard |\begin{thebibliography}{}| syntax. 
%    \begin{macrocode}
\edef\EndBib{\expandafter\@gobble\string\\end\string{thebibliography\string}}
\edef\StartBib{\expandafter\@gobble\string\\begin\string{thebibliography\string}}
\edef\StartBibX{\expandafter\@gobble\string\\begin\string{thebibliography\string}\string{\string}}
%    \end{macrocode}
% Define two new logical variables. |\ifAnotherline| is true if we
% want to read another line from the input stream. It is set to false
% when we encounter the |\end{thebibliography}|. |\ifInBib| is set to
% true when we frind the |\begin{thebibliography}| and tells us to
% start saving the following lines. 
%    \begin{macrocode}
\newif\ifAnotherLine
\newif\ifInBib
%    \end{macrocode}
% \begin{macro}{makeOther}
% We need a macro to change the catcodes of lots of things to other. 
%    \begin{macrocode}
\def\makeOther#1{\catcode`#1=12\relax}
%    \end{macrocode}
% \end{macro}
% \begin{macro}{\RereadBiblio}
% Define a macro to do all the fiddly stuff. This is used at the
% beginning of |thebibliography| environment (Sec~\ref{sec:biblio}).
%    \begin{macrocode}
\newread\tf@biblio
\newcommand{\RereadBiblio}{%
%    \end{macrocode}
% Open the current file again for reading
%    \begin{macrocode}
  \begingroup
  \immediate\openin\tf@biblio=\articlename
%    \end{macrocode}
% Everything in a group to be safe. Initialize logicals. Black magic
% to change
% catcodes and avoid blank lines being interpreted as |\par|. 
%    \begin{macrocode}
  \AnotherLinetrue\InBibfalse
  \makeOther\ \makeOther\\\makeOther\$%
  \makeOther\#\makeOther\^\makeOther\^^K%
  \makeOther\_\makeOther\^^A\makeOther\%%
  \makeOther\~\makeOther\{\makeOther\}\makeOther\&%
  \endlinechar-1\relax
%    \end{macrocode}
% Start looping through the file. Read each line to the macro
% |\Line|. If we have got to EOF or end of biblio then set
% |\AnotherLinefalse|.
%    \begin{macrocode}
  \loop
  \read\tf@biblio to \Line
  \ifx\Line\EndBib\AnotherLinefalse\InBibfalse\fi
  \ifeof\tf@biblio \AnotherLinefalse\fi
%    \end{macrocode}
% If we are in the biblio and the line is not blank then append it to
% the |\ADS@references| macro together with a newline. 
%    \begin{macrocode}
  \ifInBib
  \ifx\Line\empty\else\xdef\ADS@references{\ADS@references\Line^^J}\fi
  \fi
%    \end{macrocode}
% Check if the line we just saw was the start of biblio, then repeat
% the loop if necessary. 
%    \begin{macrocode}
  \ifx\Line\StartBib\InBibtrue\fi
  \ifx\Line\StartBibX\InBibtrue\fi
  \ifAnotherLine\repeat
  \immediate\closein\tf@biblio
  \endgroup
  }
%    \end{macrocode}
% \end{macro}
% 
% \begin{macro}{\PadToFourWithDots}
% Utility function for left-padding a number to 4 characters with
% periods. Used in 
% constructing |\ADS@BibCode|. This is a very simple algorithm and
% \emph{requires} that the argument actually be a number---if it is
% not then dire consequences will result.  
%    \begin{macrocode}
\newcommand{\PadToFourWithDots}[1]{%
  \ifnum #1 > 999 #1%
  \else\ifnum #1 > 99 .#1%
  \else\ifnum #1 > 9 ..#1%
  \else ...#1%
  \fi\fi\fi}
%%  \ifnum \count0 > 999 #1%
%%  \else\ifnum \count0 > 99 .#1%
%%  \else\ifnum \count0 > 9 ..#1%
%%  \else ...#1%
%%  \fi\fi\fi}
%    \end{macrocode}
% \end{macro}
%
% \begin{macro}{\AuthorADS}
% \begin{macro}{\ADS@affiliations}
% New mechanism for setting the affiliations for ADS. The author labels (AA, AB, \dots) are generated using two nested counters.  We use the |\robustify| command from \texttt{etoolbox} so that accented characters in the affiliation addresses do not blow up.  The list of accents may need to be expanded. 
%    \begin{macrocode}
\newcounter{AUTHOR}
\newcounter{author}[AUTHOR]
\stepcounter{AUTHOR}
\newcommand\ADS@affiliations{}
\robustify{\'}\robustify{\i}\robustify{\~}\robustify{\`}
\def\AuthorADS#1#2{%
  \begingroup
  \def\ORCIDiD##1{<ID system="ORCID">##1</ID>}
  \def\Email##1{<EMAIL>##1</EMAIL>}
  \def\Affil##1{##1}
  \stepcounter{author}
  \ifthenelse{\value{author} > 26}{\stepcounter{AUTHOR}\stepcounter{author}}{}
  \ifthenelse{\equal{\ADS@affiliations}{\@rmaaEmpty}}{%
    \xdef\ADS@affiliations{\Alph{AUTHOR}\Alph{author}(#2)}
  }{%
    \xdef\ADS@affiliations{\ADS@affiliations, \Alph{AUTHOR}\Alph{author}(#2)}
  }%
  \endgroup
}
%    \end{macrocode}
% \end{macro}
% \end{macro}
% %
% \begin{macro}{\ADS@dumpdata}
% \changes{4.0}{2019/11/13}{Write all author affiliations for ADS}
% \begin{macro}{\ADS@write}
% \begin{macro}{\ADS@openfile}
% |\ADS@dumpdata| is the macro that actually does the job of writing out all the
% data to an auxiliary file, using |\ADS@write| to write each
% field. The file is set up with |\ADS@openfile|, which currently is
% done in the TOC of a book.  The new multi-author affiliations are used if |\ADS@affiliations| has been set, otherwise it falls back to the single-author version |\ADS@affil|. 
%    \begin{macrocode}
\newcommand{\ADSfield}[2]{\relax}
\newcommand{\ADSvfield}[2]{\relax}
%%\renewcommand{\ADSfield}[12]{\%#1 #2}
\newcommand{\ADS@dumpdata}
{ 
  \ADS@write{\string\begin{ADSarticle}}
  \ADS@WRITE{\string\ADSfield{R}{\ADS@BibCode}}
  \ADS@write{\string\ADSfield{A}{\ADS@authors}}
  \ifthenelse{\equal{\ADS@affiliations}{\@rmaaEmpty}}{%
    {\ADS@sanitize\ADS@write{\string\ADSfield{F}{AA(\ADS@affil)}}}
  }{%
    {\ADS@sanitize\ADS@write{\string\ADSfield{F}{\ADS@affiliations}}}
  }
  \ADS@write{\string\ADSfield{J}{\ADS@journalname}}
  \ADS@write{\string\ADSfield{D}{\ADS@fulldate}}
  \ADS@write{\string\ADSfield{L}{\ADS@lastpage}}
  {\ADS@sanitize\ADS@write{\string\ADSfield{T}{\ADS@title}}}
  \ADS@write{\string\ADSfield{C}{\ADS@copyright}}
  \ADS@write{\string\ADSfield{d}{\ADS@doi}}
  \ADS@write{\string\ADSfield{U}{\ADS@url}}
  \ADS@write{\string\ADSfield{K}{\ADS@keywords}}
  {\ADS@sanitize\ADS@write{\string\ADSfield{B}{\ADS@abstract}}}
  \ADS@write{\string\ADSfield{Z}{\ADS@references}}
  \ADS@write{\string\end{ADSarticle}}
  }
%%\newcommand{\ADS@write}[1]{\if@book\addtocontents{ads}{#1}\fi}
%%\newcommand{\ADS@openfile}{\@starttoc{ads}}
\newcommand{\ADS@write}[1]{\set@display@protect\immediate\write\tf@ads{#1}}
\newcommand{\ADS@WRITE}[1]{\immediate\write\tf@ads{#1}}
\newcommand{\ADS@openfile}{%
  \newwrite\tf@ads%
  \immediate\openout\tf@ads\jobname.ads\relax
  \ADS@WRITE{\ADS@preamble}}
\newcommand{\ADS@closefile}{\ADS@WRITE{\ADS@postamble}}
\def\ADS@preamble{%
  %%  \string\documentclass{article}
  %%  \string\begin{document}
  }
\def\ADS@postamble{%
  %%  \string\end{document}
  \string%\string% Local Variables:^J 
  \string%\string% mode: LaTeX^J
  \string%\string% mode: Font-lock^J
  \string%\string% End:^J
  }
%    \end{macrocode}
% \end{macro}
% \end{macro}
% \end{macro}
% 
% \subsection{Automatic calculation of the page range}
% \label{sec:page-range}
% Another issue that has cost me a lot of grief. Still, it finally
% seems to be sorted now. All hail Jeff Goldberg. 
%
% Cunning trick to get the page numbers from |\pageref|.  Put a label
% on the first and last page of the document using |\AtBeginDocument|
% and |\AtEndDocument|. .  This way won't work for the book option.
% In that case, we do it by defining |\@endofarticlehook| and
% |\@begofarticlehook| as parallels to |\AtEndDocument| and
% |\AtBeginDocument|.  Note that we have to defer adding
% |\@begofarticlehook| to |\AtBeginDocument| to make sure it comes
% after the redefinitions of |\ref| and |\label|.
%    \begin{macrocode}
\if@book\AtEndDocument{\ADS@closefile}\else
 %\AtBeginDocument{\@begofarticlehook}
\AtEndDocument{\@endofarticlehook}
\fi
%    \end{macrocode}
% \begin{macro}{\@begofarticlehook}
% Since we are getting more and more things that need to be done at
% the beginning of every article, frontsection, etc we define a sister
% hook to 
% |\@endofarticlehook|. This hook does the following: 
% \begin{enumerate}
% \item Labels the first page
% \item Sets |\ADS@firstpage|
% \item runs and then clears |\indexauthor@deferlist|
% \item resets the keyword counter
% \end{enumerate}
% It is executed in the following circumstances: 
% \begin{enumerate}
% \item At |\begin{document}| for standalone files
% \item At |\begin{document}| for each included file
% \item At beginning of front sections 
% \item At beginning of TOC
% \item At beginning of Author Index
% \item At beginning of abstracts
% \end{enumerate}
% For some reason, this is not always within the right group, so we
% use |\xdef| just to be sure. 
%    \begin{macrocode}
\newcommand{\@begofarticlehook}{%
  \if@frontmatter\frontlabel{RMAAFirstPage}\else\label{RMAAFirstPage}\fi%
  \xdef\ADS@firstpage{\@arabic{\c@page}}%
  \if@book\else\makeindex\fi % needed so default index is defined in stand-alone articles
%%  \indexauthor@deferlist
%%  \xdef\indexauthor@deferlist{}
  \@inpreamblefalse
  \setcounter{kwdcounter}{0}
  }
%    \end{macrocode}
% \end{macro}
% \begin{macro}{\@endofarticlehook}
% This is a hook to be run at the end of each article
% currently, it just labels the last page and writes out all the ADS
% data to an auxiliary file.  
% It should not produce any printed output. 
% It is modelled on the |\AtEndDocument| command, to which it reduces
% in the case of a single article. 
%    \begin{macrocode}
\newcommand{\@endofarticlehook}{%
  \if@debug{\ClassWarning{rmaa}{***Doing Hook***}}\fi
  \if@debug{\ClassWarning{rmaa}{ADSabstract is: \ADS@abstract}}\fi
  \clearpage\rmaa@lastpagelabel\if@book\ADS@dumpdata\fi
  \ifx\NotesList\@rmaaEmpty\else\OutputNotes\fi
  \if@debug{\ClassWarning{rmaa}{Notes list is: \NotesList}}\fi
}
%    \end{macrocode}
% \begin{macro}{\rmaa@lastpagelabel}
% New version of last page label command copied from Jeff Goldberg's
% lastpage.sty package - we can't just use |\label| here in case the
% output routine for the last page has already been called. Instead,
% we first explicitly clear the page, step back the page counter,
% explicitly write to the .aux file, then restep the page counter
% again.
%
% But this will conflict with the hyperref package, which does labels differently.  Compensate for that by using the five-argument form.
% 
% Also set the |\ADS@LastPage| macro to the appropriate number. 
%    \begin{macrocode}
\newcommand{\rmaa@lastpagelabel}{%
  \addtocounter{page}{-1}%
  \immediate\write\@auxout{\string
    \newlabel{\articlename:RMAALastPage}{{}{\thepage}{}{}{}}}%
  \edef\ADS@lastpage{\thepage}%
  \addtocounter{page}{1}}
%    \end{macrocode}
% \end{macro}
% \end{macro}

% \subsection{Deal with explicit change in number of columns}
% \label{sec:columns}
% Sometimes, when formatting a book, we need to switch between
% one-column articles and two-column articles. For this to work, we
% need to wrap the standard \LaTeX{} commands to set the |\@rmtwocol|
% flag appropriately and also reset the package default values of all
% scalable lengths. .  
%    \begin{macrocode}
\let\latexonecolumn=\onecolumn
\let\latextwocolumn=\twocolumn
\renewcommand{\onecolumn}{\@rmtwocolfalse\@package@default@lengths\latexonecolumn}
\renewcommand{\twocolumn}{\@rmtwocoltrue\@package@default@lengths\latextwocolumn}
%    \end{macrocode}
% 
% \subsection{The actual formatting of the title}
% \label{sec:title-format}

% Note, we have to redefine the |\maketitle| command to use the
% numbered footnotes. This means we needn't use the atbeginend stuff. 
%    \begin{macrocode}
\renewcommand{\maketitle}{\par
  \begingroup
%    \end{macrocode}
% Don't switch to marked footnotes. 
% This is consistent with existing style but I'm not sure I like it
%    \begin{macrocode}
  %%  \renewcommand\thefootnote{\fnsymbol{footnote}}%
%%  \def\@makefnmark{\hbox to\z@{$\m@th^{\@thefnmark}$\hss}}%
  \def\@makefnmark{\textsuperscript{\@thefnmark}\hspace{-0.2em}}%
  \long\def\@makefntext##1{\parindent 1em\noindent
    \hbox to1.8em{\hss$\m@th^{\@thefnmark}$}##1}%
  \if@rmtwocol
  { \if@debug{\ClassWarning{rmaa}{TWO COLUMN}}\fi
    \twocolumn[\@maketitle]%
    }
  \else
  { \if@debug{\ClassWarning{rmaa}{ONE COLUMN}}\fi
    \newpage
    \global\@topnum\z@   % Prevents figures from going at top of page.
    \if@manuscript \@MSmaketitle \else \@maketitle \fi
    }
  \fi
  \if@manuscript\else\thispagestyle{rmaafirstpage}\@thanks\fi
  \ifx\NotesList\@rmaaEmpty\else{%
    \NoteMark{1--\arabic{notecounter}}\vspace*{-2\baselineskip}%
  }\fi%
  \endgroup
%    \end{macrocode}
% Don't reset footnote counter (not sure whether this is good thing
% or not).
%    \begin{macrocode}
  %%  \setcounter{footnote}{0}%
  \let\thanks\relax\let\affil\relax\let\altaffil\relax
  \let\maketitle\relax\let\@maketitle\relax
  \gdef\@thanks{}\gdef\@author{}\gdef\@title{}}
%    \end{macrocode}

% \begin{macro}{\@kwhwidth}
% \begin{macro}{\@kwtwidth}
% Redefine \cs{@maketitle} to include the abstract, resumen and
% keywords. Note that in the manuscripty style we use
% \cs{@MSmaketitle} instead.  
% \changes{3.16}{2003/04/07}{
%  Added extra vertical space between authors and accepted date.}
% \changes{3.26}{2011/08/15}{Add extra label to review articles}
%    \begin{macrocode}
\newlength{\@kwhwidth}\newlength{\@kwtwidth}
\newcommand\reviewname{Review}
\renewcommand{\@maketitle}{%
  \newpage
  \null
  \if@reviewarticle
  \smallskip
  {\large\bfseries\itshape\reviewname}%
  \bigskip
  \else
  \vskip\@before@title@skip%
  \fi
  \begin{center}%
    {\@title \par}%
    \vskip\@before@author@skip%
    \parbox{\rmaa@AuthorWidthFrac\textwidth}{\centering
      \lineskip\@between@author@skip%
      \@author
      \par}
    %% this should be replaced with `Submitted to...' for preprints
    \if@printdates
    \vskip\@before@dates@skip%
    \if@preprint
    \rmaa@ReceivedStyle{Submitted to RMxAA \@received}
    \else
    \rmaa@ReceivedStyle{Received \@received; accepted \@accepted}
    \fi
    \fi
    \if@manuscript
    \rmaa@ReceivedStyle{Submitted to RMxAA \@received}
    \fi
    \if@guide
    \rmaa@ReceivedStyle{Version \rmaa@version, \rmaa@date}
    \fi
    \vskip\@before@abstract@skip%
    \if@abstractpresent
    \parbox{\rmaa@AbstractWidthFrac\textwidth}{%
      \setlength{\parindent}{0.9cm}%
      \centerline{\rmaa@AbstractHeaderStyle{Abstract}}\par%
      \vskip\@after@abstracthead@skip%
      \if@newstyle\noindent\fi\@abstract
      \vskip\@between@abstract@skip%
      \centerline{\rmaa@AbstractHeaderStyle{Resumen}}\par%
      \vskip\@after@abstracthead@skip%
      \if@newstyle\noindent\fi\@resumen
      \if@setresumen\else\@abstract\fi % append english abstract if
                                       % spanish one not set
      \if@guide\else            %do not use keywords for guide
      \vskip\@before@kwd@skip%
      % Format the Key Words with hanging indentation:
      % They are put in a \parbox with width equal to the line width
      % minus the width of the header. 
      { % Widths of keyword
        \settowidth{\@kwhwidth}{\@keywordheader~}   % header and
        \setlength{\@kwtwidth}{\linewidth}          % keyword text
        \addtolength{\@kwtwidth}{-\@kwhwidth}       %
        \noindent\@keywordheader~\parbox[t]{\@kwtwidth}{%
          \rmaa@KeywordStyle{\boldmath\uppercase{\@keywords}}}}%
      \fi
      }
      \fi
  \end{center}%
  \vspace*{\afterkwdskip}
  \par
  }
\newlength{\afterkwdskip}\setlength{\afterkwdskip}{\z@}
%    \end{macrocode}
% \end{macro}
% \end{macro}
%
%
% \begin{macro}{\reviewarticle}
% \begin{macro}{\reviewtoc}
% \begin{macro}{\reviewlof}
%   User command to specify a review article. Must be used in
%   preamble. Main effect is to place the text ``Review'' before the
%   article title.  Also add command for setting a table of contents
%   and a list of figures at the beginning of the article.
%    \begin{macrocode}
\newcommand{\reviewarticle}{\@reviewarticletrue}
\newcommand\reviewtoc{%
  \section*{CONTENTS}%
  \@starttoc{rtc}%
}
\newcommand\reviewlof{%
  \section*{LIST OF FIGURES}%
  \@starttoc{lof}%
}
%    \end{macrocode}
% \end{macro}
% \end{macro}
% \end{macro}
% \begin{macro}{\l@section}
%    \begin{macrocode}
\renewcommand*\l@section[2]{%
  \ifnum \c@tocdepth >\z@
    \addpenalty\@secpenalty
    \addvspace{0.8em \@plus\p@}%
    \setlength\@tempdima{1.5em}%
    \begingroup
      \parindent \z@ \rightskip \@pnumwidth
      \parfillskip -\@pnumwidth
      \leavevmode %\bfseries
      \advance\leftskip\@tempdima
      \hskip -\leftskip
      #1\nobreak\hfil \nobreak\hb@xt@\@pnumwidth{\hss #2}\par
    \endgroup
  \fi}
%    \end{macrocode}
% \end{macro}
%
% \begin{macro}{\@MSmaketitle}
% A special version of |\@maketitle| for the \texttt{manuscript}
% option. This can deal with a multi-page abstract/resumen. Not very
% well, mind. 
%    \begin{macrocode}
\newcommand{\@MSmaketitle}{%
  \if@debug\ClassInfo{rmaa}{Using special MS maketitle command}\fi
  \newpage
  {
    \thispagestyle{rmaafirstpage}
%     \setlength{\parindent}{0pt}
    \null\vskip 2ex
    \begin{center}
      \@title\par
      \@author\par
      \rmaa@ReceivedStyle{%
        \ifx\@received\@rmaaEmpty
          Draft version: \@date
        \else
          Received: \@received
          \ifx\@accepted\@rmaaEmpty
          \else
            ; Accepted: \@accepted
          \fi
        \fi}\par
    \end{center}
    \@thanks
    \centerline{\rmaa@AbstractHeaderStyle{Resumen}}
    \noindent\@resumen\par\vskip 2ex
    \centerline{\rmaa@AbstractHeaderStyle{Abstract}}
    \noindent\@abstract\par\vskip 2ex
    \noindent\@keywordheader~\rmaa@KeywordStyle{\@keywords}\par
%     \@keywordheader~\rmaa@KeywordStyle{\uppercase{\@keywords}}\par
  }%
  }
%    \end{macrocode}
% \end{macro}

% \subsection{The sectioning commands}
% \label{sec:section-commands}

% \begin{macro}{\section}
% \begin{macro}{\subsection}
% \begin{macro}{\subsubsection}
% \begin{macro}{\paragraph}
% \begin{macro}{\subparagraph}
% Redefine all the standard commands to use these lengths. 
%    \begin{macrocode}
\renewcommand\section{\@startsection {section}{1}{\z@}%
                                   {\abovesectionskip}%
                                   {\belowsectionskip}%
                                   %% how do we get upper case here
                                   {\centering\reset@font\normalsize}}
\renewcommand\subsection{\@startsection{subsection}{2}{\z@}%
                                     {\abovesectionskip}%
                                     {\belowsectionskip}%
                                     {\centering\reset@font\normalsize\itshape}}
\renewcommand\subsubsection{\@startsection{subsubsection}{3}{\z@}%
                                     {\abovesectionskip}%
                                     {\belowsectionskip}%
                                     {\centering\reset@font\normalsize\itshape}}
\renewcommand\paragraph{\@startsection{paragraph}{4}{\z@}%
                                    {\aboveparaskip}%
                                    {\belowparaskip}%
                                    {\reset@font\normalsize\bfseries}}
\renewcommand\subparagraph{\@startsection{subparagraph}{5}{\parindent}%
                                      {\aboveparaskip}%
                                      {\belowparaskip}%
                                      {\reset@font\normalsize\bfseries}}
%    \end{macrocode}
% \end{macro}
% \end{macro}
% \end{macro}
% \end{macro}
% \end{macro}

% \begin{macro}{secnumdepth}
% Make sure that we number the subsubsections---this is done by
% default in the |article| class but not the |book| class, which is
% why we must set it explicitly here.
%    \begin{macrocode}
\setcounter{secnumdepth}{3}
%    \end{macrocode}
% \end{macro}
% \changes{2.2beta10}{2001/11/07}{Fixed \cs{secnumdepth} for \cs{book}
% style.}

% \begin{macro}{\thesubsection}
% \begin{macro}{\thesubsubsection}
% For subsection and subsubsection, we want the numbers to be roman
% but the title to be italic. This seems the easiest way to go about it.
%    \begin{macrocode}
\renewcommand{\thesubsection}{{\upshape\thesection.\@arabic\c@subsection}}
\renewcommand{\thesubsubsection}{\thesubsection{\upshape.\@arabic\c@subsubsection}}
%    \end{macrocode}
% \end{macro}
% \end{macro}

% \begin{macro}{\@seccntformat}
% Add ``.'' to the number here so we can remove it from |\thesection|,
% etc. Reduce gap between number and title WJH 27 Jan 1995
%    \begin{macrocode}
\def\@seccntformat#1{\csname the#1\endcsname .\hskip 0.4em}
%    \end{macrocode}
% \end{macro}

% \begin{macro}{\@sect}
% I didn't really want to mess with the |\@sect| command, but it was
% the only way I could manage to force upper case. Other than the
% commented parts, it is identical to the standard version. 
%    \begin{macrocode}
\def\@sect#1#2#3#4#5#6[#7]#8{
  \ifnum #2>\c@secnumdepth
     \let\@svsec\@empty
  \else
     \refstepcounter{#1}%
     \let\@@protect\protect
     \def\protect{\noexpand\protect\noexpand}%
     \edef\@svsec{\@seccntformat{#1}}%
     \let\protect\@@protect
  \fi
%    \end{macrocode}
% This is the bit I changed (27 Jan 1995). Use upper case for
% |\section|, use mixed case for the rest. 
%    \begin{macrocode}
  \ifnum #2=1                   
     \def\sec@upcase##1{\uppercase{##1}}%
     \def\sec@shape{\upshape}%
  \else                         
     \def\sec@upcase##1{\relax{##1}}%
     \def\sec@shape{\itshape}%
  \fi
  \@tempskipa #5\relax
  \ifdim \@tempskipa>\z@
     \begingroup #6\relax
          \@hangfrom{\hskip #3\relax\@svsec}%
                    {\interlinepenalty \@M \sec@upcase{#8}\par}%
     \endgroup
     \csname #1mark\endcsname{#7}%
%    \end{macrocode}
% Only write TOC entries if doing the Guide style. 
%    \begin{macrocode}
     \if@guide 
        \addcontentsline{toc}{#1}{%
          \ifnum #2>\c@secnumdepth \else
             \protect\numberline{\csname the#1\endcsname}%
          \fi
          #7}%
     \fi
%    \end{macrocode}
% Also allow per-article TOC for review papers
%    \begin{macrocode}
     \if@reviewarticle
        \ifnum #2>2 \else
           \addcontentsline{rtc}{#1}{%
             \protect\numberline{\csname the#1\endcsname}%
             \sec@upcase{\sec@shape #7}}%
        \fi
     \fi
  \else
     \def\@svsechd{#6\hskip #3\relax  %% \relax added 2 May 90
       \@svsec #8\csname #1mark\endcsname{#7}%
       \if@guide
          \addcontentsline{toc}{#1}{%
            \ifnum #2>\c@secnumdepth \else
               \protect\numberline{\csname the#1\endcsname}%
            \fi
            #7}%
       \fi
       }%
  \fi
  \@xsect{#5}}
%    \end{macrocode}
% \end{macro}


% \subsection{Appendices}
% \label{sec:appendices}
% Make sure that appendices in individual articles work properly. We
% copy the definition from the article class (the book one doesn't do
% what we want).  
% \begin{environment}{appendices}
%   This uses an environment |appendices| inside of which each
%   |\section| will become a separate appendix.  Note that the
%   appendices must come before the references (otherwise the full
%   addresses mechanism will not work properly). We make the equation
%   numbering be subordinate to the appendix (equations A12, B6,
%   etc). Finally, we have a medskip followed by APPENDICES to
%   introduce them. \changes{3.17}{2002/04/08}{ Added appendices environment
%   that does the right thing.}
%    \begin{macrocode}
\newenvironment{appendices}{\par
  \begingroup
  \setcounter{section}{0}%
  \setcounter{subsection}{0}%
  \renewcommand\thesection{\@Alph\c@section}
  \renewcommand\theequation{\thesection\arabic{equation}}
  \medskip
  \section*{APPENDICES}}
{\endgroup}
%    \end{macrocode}
% \end{environment}
% \begin{environment}{appendix}
%   If there is only one appendix, then use |appendix| environment
%   instead.  \changes{4.0}{2019/11/12}{ Added
%   appendix environment for case where there is only one section.}
%    \begin{macrocode}
\renewenvironment{appendix}{\par
  \begingroup
  \setcounter{section}{0}%
  \setcounter{subsection}{0}%
  \renewcommand\thesection{\@Alph\c@section}
  \renewcommand\theequation{\thesection\arabic{equation}}
  \medskip
  \section*{APPENDIX}}
{\endgroup}
%    \end{macrocode}
% \end{environment}

% \subsection{Floating thingies}
% \label{sec:floats}   

%% 
%    \begin{macrocode}
\renewcommand{\thetable}{\@arabic\c@table}
\renewcommand{\thefigure}{\@arabic\c@figure}
\renewcommand{\figurename}{Fig.}
\renewcommand{\tablename}{TABLE}
%\renewcommand{\fnum@table}{\MakeUpperCase{\tablename}~\thetable}
%    \end{macrocode}
% For some reason, I have to write `TABLE' explicitly here
% If I use |\tablename|, it gets converted to mixed case---God knows
% why. 
% |\figurename| is also being ignored. 
%    \begin{macrocode}
\renewcommand{\fnum@table}{TABLE~\thetable}
\renewcommand{\fnum@figure}{Fig.~\thefigure}
%    \end{macrocode}

% Do captions in 9 pt face in the main journal style and in new-style
% anything. 
%    \begin{macrocode}
\long\def\@makefigcap#1#2{%
  \if@proceedings\else\small\fi
  \if@newstyle\small\fi 
  \vskip\abovecaptionskip
  \sbox\@tempboxa{#1. #2}%
  \ifdim \wd\@tempboxa >\hsize
    #1. #2 \par
  \else
    \global \@minipagefalse
    \hb@xt@\hsize{\hfil\box\@tempboxa\hfil}%
  \fi
  \vskip\belowcaptionskip}
%    \end{macrocode}

% Make sure table caption is upper case. 
%    \begin{macrocode}
\long\def\@maketabcap#1#2{%
  \hbox to\hsize{\hfil #1 \hfil}
  \vskip 6pt 
  \centering\MakeUpperCase{#2}\par
  \vskip 6pt}
%    \end{macrocode}

%    \begin{macrocode}
\newlength{\@rottabcapwidth}
\long\def\@makerottabcap#1#2{%
  \setbox\@tempboxa\hbox{\MakeUpperCase{#2}}%
  \ifdim\wd\@tempboxa >.8\vsize
     \setlength\@rottabcapwidth{.8\vsize}
  \else
     \setlength\@rottabcapwidth{\wd\@tempboxa}
  \fi
  \rotatebox{90}{%
    \begin{minipage}{\vsize}\centering
      \begin{minipage}{\@rottabcapwidth}\centering
        #1\\ \vskip 6pt \MakeUpperCase{#2}
      \end{minipage}%
    \end{minipage}
    }%
  \hskip 6pt
  }

\newenvironment{Sideways}
    {\begin{sideways}\begin{minipage}{\vsize}\centering}
    {\end{minipage}\end{sideways}}

\newenvironment{Sidewaystable}
{%
  \let\oldcaption=\caption%
  \def\caption{
    \def\@captype{table}\renewcommand{\@makecaption}{\@maketabcap}
    \oldcaption}
  \begin{sidewaystable}\centering}
  {\end{sidewaystable}\let\caption=\oldcaption}
\newenvironment{Sidewaystable*}
{%
  \let\oldcaption=\caption%
  \def\caption{
    \def\@captype{table}\renewcommand{\@makecaption}{\@maketabcap}
    \oldcaption}
  \begin{sidewaystable*}\centering}
  {\end{sidewaystable*}\let\caption=\oldcaption}
%    \end{macrocode}

% Not sure why we need to stick these lines in. 
%    \begin{macrocode}
\renewenvironment{figure}
                 {\renewcommand{\@makecaption}{\@makefigcap}\@float{figure}}
                 {\end@float}
\renewenvironment{figure*}
                 {\renewcommand{\@makecaption}{\@makefigcap}\@dblfloat{figure}}
                 {\end@dblfloat}
\renewenvironment{table}
                 {\renewcommand{\@makecaption}{\@maketabcap}\@float{table}}
                 {\end@float}
\renewenvironment{table*}
                 {\renewcommand{\@makecaption}{\@maketabcap}\@dblfloat{table}}
                 {\end@dblfloat}
%    \end{macrocode}

% Redefine |\@caption| so that we don't bother with the list of
% figures. This saves us having to worry about fragile commands in the
% captions. 
%    \begin{macrocode}
\long\def\@caption#1[#2]#3{%
  \par
%%  \addcontentsline{\csname ext@#1\endcsname}{#1}%
%%    {\protect\numberline{\csname the#1\endcsname}{\ignorespaces #2}}%
  \begingroup
    \@parboxrestore
    \if@minipage
      \@setminipage
    \fi
    \normalsize
    \@makecaption{\csname fnum@#1\endcsname}{\ignorespaces #3}\par
  \endgroup}
%    \end{macrocode}

% Make tables not quite so cramped. Allow |\hline\hline| to give a
% poor man's thick rule. 
%    \begin{macrocode}
\renewcommand{\arraystretch}{1.15}
\setlength{\doublerulesep}{0.4pt}
%    \end{macrocode}

% \begin{macro}{\tablecols}
% \begin{macro}{\@tablecols}
% \begin{macro}{\tabcaptioncontinued}
% |\tablecols| is a user command for setting the number of columns. It
% is used by commands that try and span the entire table (I think
% there must be a better way of acheiving
% this). |\tabcaptioncontinued| is 
% a user command for setting the caption of a continued table on the
% next page. 
%    \begin{macrocode}
\newcommand{\tablecols}[1]{\renewcommand{\@tablecols}{#1}}
\newcommand{\@tablecols}{1}
\newcommand{\tabcaptioncontinued}{%
  \multicolumn{\@tablecols}{c}{\addtocounter{table}{-1}TABLE~\thetable\ (CONTINUED)} \\ \\
  }
%    \end{macrocode}
% \end{macro}
% \end{macro}
% \end{macro}

% \subsubsection{Footnotes in tables}
% \label{sec:table-footnotes}
% \begin{macro}{\tablenotemark}
% \begin{macro}{\tablenote}
% \begin{macro}{\tablenoteskip}
% Old user commands |\tablenotemark{MARK}|,
% |\tablenote{MARK}{TEXT}|---now deprecated. 
%    \begin{macrocode}
\newcommand\tablenotemark[1]{$^\mathrm{#1}$} 
\newcommand\tablenote[2]{\if@proceedings\else\small\fi ${}^\mathrm{#1}$ #2} 
\newcommand\tablenoteskip{\\[0.5\baselineskip]} 
%    \end{macrocode}
% \end{macro}
% \end{macro}
% \end{macro}

% \begin{macro}{\tabnotemark}
% \begin{macro}{\tabnotetext}
% \begin{macro}{\tabnotetextTWO}
% Alternative way of doing footnotes in tables. 
%    \begin{macrocode}
\newcommand{\tabnotemark}[1]{\textsuperscript{\lowercase{#1}}}
%\newcommand{\tabnotetextTWO}[4]{\tabnotetext{#1}{#2}\tabnotetext{#3}{#4}}
\newcommand{\tabnotetextTWO}[4]{%
  \multicolumn{\@tablecols}{l}{%
    \parbox{0.5\tabnotewidth}{\small\textsuperscript{#1}#2}%
    \parbox{0.5\tabnotewidth}{\small\textsuperscript{#3}#4}%
    }%
  \\}
\newcommand{\tabnotetext}[2]{%
  \multicolumn{\@tablecols}{@{}p{\tabnotewidth}@{}}{\small\textsuperscript{#1}#2}\\}
\newlength{\tabnotewidth}
%    \end{macrocode}
% \end{macro}
% \end{macro}
% \end{macro}


% \begin{macro}{changemargin}
%    \begin{macrocode}
% Command to allow margins to be changed on the fly. Useful to allow
% tables and Figures to protrude into the margins. 
%
% \changes{3.24}{2007/09/14}{Added \texttt{changemargin} environment
% (copied from TeX FAQ).}
\newenvironment{changemargin}[2]{% 
  % From TeX FAQ http://www.tex.ac.uk/cgi-bin/texfaq2html?label=chngmargonfly
  \begin{list}{}{%
    \setlength{\topsep}{0pt}%
    \setlength{\leftmargin}{#1}%
    \setlength{\rightmargin}{#2}%
    \setlength{\listparindent}{\parindent}%
    \setlength{\itemindent}{\parindent}%
    \setlength{\parsep}{\parskip}%
  }%
  \item[]}{\end{list}}
%    \end{macrocode}
% \end{macro}

% \subsection{The bibliography and references}
% \label{sec:biblio}
% \begin{macro}{\beforebiblioskip}
% \begin{macro}{\afterbiblioskip}
% I'm not entirely sure why I introduced these two commands, since
% directly putting |\vspace| before and after
% |\begin{thebibliography}| seems to work well enough. 
%    \begin{macrocode}
\newlength{\beforebiblioskip}
\setlength{\beforebiblioskip}{0pt}
\newlength{\afterbiblioskip}
\setlength{\afterbiblioskip}{0pt}
%    \end{macrocode}
% \end{macro}
% \end{macro}
% \begin{macro}{thebibliography}
% The environment itself is done as starred section with extra
% vertical space after heading. This space is currently zero. 
% In the book style we do the cunning trick with rereading the file in
% order to get the references into the |\ADS@references| macro.
% 
% \changes{3.05}{2003/01/02}{Removed the \cs{beforebiblioskip} and
% \cs{afterbiblioskip} commands, since they cause unwanted vertical
% space in some instances.}  
%    \begin{macrocode}
\renewenvironment{thebibliography}{%
  \if@book\RereadBiblio\fi
%%  \vspace*{\beforebiblioskip}
  \section*{REFERENCES} \frenchspacing
%%  \vspace*{\afterbiblioskip}
  \small
%    \end{macrocode} 
% The references themselves are set as a list in order 
% to get the hanging indentation. They are done in a smaller font and
% with french spacing so we needn't worry about the distinction
% between sentence-ending periods and abbreviations. 
%    \begin{macrocode} 
  \list{\null}{\leftmargin 1.5em\labelwidth\z@\labelsep\z@\itemindent -1.5em
    \itemsep\z@\parsep\z@\usecounter{enumi}}
  \sloppy\clubpenalty4000\widowpenalty4000
  \sfcode`\.=1000\relax\@biblisttrue}
{\endlist\@biblistfalse\normalsize\@outputaddressesifnecessary}
%    \end{macrocode}
% \end{macro}
% \subsection{Uniquification of cross-references}
% \label{sec:uniquification}
% 
% \begin{macro}{\articlename}
% In order that the |\label{}| command in individual articles be
% guaranteed to give a unique anchor, we append the article name (held
% in the |\articlename| command) to each key. This requires
% redefinition of the following commands: |\cite|, |\nocite|, |\ref|,
% |\pageref|, |\label|. 
% \begin{macro}{\oldcite}
% \begin{macro}{\oldnocite}
% \begin{macro}{\oldref}
% \begin{macro}{\oldpageref}
% \begin{macro}{\oldlabel}
% The original versions of these commands are
% saved as |\oldcite|, |\oldnocite|, etc. 
%
% The uniquifying of cross-references is now always done, so 
% we don't get multiply-defined labels in the single article case
%    \begin{macrocode}
\newcommand{\articlename}{anon}
\AtBeginDocument{
  \let\oldcite=\cite 
  \let\oldnocite=\nocite 
  \let\oldref=\ref 
  \let\oldpageref=\pageref
  \let\oldlabel=\label
}
%    \end{macrocode}
% \end{macro}
% \end{macro}
% \end{macro}
% \end{macro}
% \end{macro}
% \end{macro}
%
% Uniquify all cross-references by prepending the article name.  This
% now has to go into |\AtBeginDocument| so it is working on the new
% versions that hyperref.sty has set up.  Then we need to put the
% |\@begofarticlehook| afterwards, so that the first page label is using
% the new definitions.  If amsmath is loaded, then we also need to patch
% the special label command that they use. \changes{4.0}{2019/11/18}{Fix labels for amsmath display equations, in case that package has been loaded.}
%    \begin{macrocode}
\AtBeginDocument{
  \renewcommand{\ref}[1]{\oldref{\articlename:#1}%
    \if@debug{\ClassInfo{rmaa}{Reference to \articlename:#1}}\fi}
  \renewcommand{\pageref}[1]{\oldpageref{\articlename:#1}%
    \if@debug{\ClassInfo{rmaa}{Page reference to \articlename:#1}}\fi}
  \renewcommand{\label}[1]{\oldlabel{\articlename:#1}%
    \if@debug{\ClassInfo{rmaa}{Label for \articlename:#1}}\fi}%
  \ifdef{\label@in@display}{%
    \let\oldlabel@in@display=\label@in@display
    \renewcommand{\label@in@display}[1]{\oldlabel@in@display{\articlename:#1}%
      \if@debug{\ClassInfo{rmaa}{AMSTeX Label for \articlename:#1}}\fi}%
  }{}
  \@begofarticlehook
}%
%    \end{macrocode}
% Citations are more complicated since they can contain a list of
% keys. These are really only here for use with the rmaacite package,
% which is utterly deprecated (02 Jan 2003) now that I have the natbib
% code working. The command \cs{cite} should \emph{not} be used with
% natbib: use \cs{citep}, \cs{citet}, etc.\@ instead. 
%    \begin{macrocode}
\renewcommand{\nocite}[1]{\@bsphack
  \@for\@citeb:=#1\do{%
    \edef\@citeb{\expandafter\@firstofone\@citeb:\articlename}%
    \if@filesw\immediate\write\@auxout{\string\citation{\@citeb}}\fi
    \@ifundefined{b@\@citeb:\articlename}{\G@refundefinedtrue
      \@latex@warning{Citation `\@citeb' undefined}}{}}%
  \@esphack}
\renewcommand{\cite}[1]{\oldcite{#1:\articlename}}
%   \end{macrocode}


% Horrible fiddling with the bibliography citation mechanism. 
% This has to be globally available in case rmaacite puts it in the
% \texttt{.aux} file. 
%    \begin{macrocode}
\def\lbibcite#1#2#3{\global\@namedef{b@#1}{#2}
                    \global\@namedef{by@#1}{#3}}
%    \end{macrocode}

% \begin{macro}{\oldbibitem}
% \begin{macro}{\@rmaaBibArg}
% \begin{macro}{\@rmaaSpace}
% \begin{macro}{rmaabib}
% Of course, |\bibitem| needs to get the uniquification treatment as
% well. While we are at it, we redefine it to use consecutive integers
% for the key if the user leaves it blank (using the counter
% |rmaabib|). This allows one to put 
% |\bibitem{}| for all your references without getting annoying
% messages about multiply-defined labels.
% 
% We also take the opportunity to insert the optional argument to
% \cs{bibitem} if it is missing in the form Author(DATE) so that natbib
% doesn't try to do numerical citations. Natbib already provides a
% hook \cs{@extra@b@citeb} to help uniquify the citation key, so we
% just plug into that. However, it puts it \emph{after} the original
% key, not before as I had had it. Easiest solution is to change the
% way that I do it. 02 Jan 2003. Note that the \cs{@extra@b@citeb}
% does not get written to the .aux  file in the \cs{citation} command
% but it does work. 
% \changes{3.14}{06 Feb 2003}{Reinstalled old definition of \cs{bibitem}
%  for the case where the nonatbib option is requested, so that the
%  rmaacite package will work. This means that you can't use rmaacite
%  and natbib in the same volume.}
%    \begin{macrocode}
\if@natbib\def\@extra@b@citeb{:\articlename}\fi
%%\else
\let\oldbibitem\bibitem
\newcommand{\@rmaaBibArg}{}
\newcommand{\@rmaaSpace}{ }
\newcounter{rmaabib}
\if@natbib
\renewcommand{\bibitem}[2][ABC(2003)]{%
  \renewcommand{\@rmaaBibArg}{#2}
  \ifx\@rmaaBibArg\@rmaaEmpty
     \addtocounter{rmaabib}{1}
     \oldbibitem[#1]{\thermaabib:\articlename}
  \else\ifx\@rmaaBibArg\@rmaaSpace
     \addtocounter{rmaabib}{1}
     \oldbibitem[#1]{\thermaabib:\articlename}
  \else 
     \oldbibitem[#1]{#2:\articlename}
  \fi\fi
  \renewcommand{\@rmaaBibArg}{}
  }
\else
\renewcommand{\bibitem}[1]{%
  \renewcommand{\@rmaaBibArg}{#1}
  \ifx\@rmaaBibArg\@rmaaEmpty
     \addtocounter{rmaabib}{1}
     \oldbibitem{\articlename:\thermaabib}
  \else\ifx\@rmaaBibArg\@rmaaSpace
     \addtocounter{rmaabib}{1}
     \oldbibitem{\articlename:\thermaabib}
  \else
     \oldbibitem{\articlename:#1}
  \fi\fi
  \renewcommand{\@rmaaBibArg}{}
  }
\fi
%%\fi
%    \end{macrocode}
% \end{macro}
% \end{macro}
% \end{macro}
% \end{macro}

% \begin{macro}{oldpcite}
% \begin{macro}{oldscite}
% Then, we need to repeat the above scheme, but for the |rmaacite|
% definition. 
%    \begin{macrocode}
\newcommand{\@redefinebibitem}{%
  \let\oldbibitem\bibitem
  \let\oldcite\cite
  \let\oldpcite\pcite
  \let\oldscite\scite
  \def\bibitem[##1<##2>]##3{%
    \renewcommand{\@rmaaBibArg}{##3}
    \ifx\@rmaaBibArg\@rmaaEmpty
       \addtocounter{rmaabib}{1}
       \oldbibitem[##1<##2>]{\thermaabib:\articlename}
    \else\ifx\@rmaaBibArg\@rmaaSpace
       \addtocounter{rmaabib}{1}
       \oldbibitem[##1<##2>]{\thermaabib:\articlename}
    \else 
       \oldbibitem[##1<##2>]{##3:\articlename}
    % trust the author here
%%       \oldbibitem[##1<##2>]{##3}
    \fi\fi
    }
  }
%    \end{macrocode}
% \end{macro}
% \end{macro}


% \subsection{HTML Table of Contents}
% \label{sec:hyper-TOC}
% \begin{macro}{\HyperTOCEntry}
% \begin{macro}{\HyperTOCpart}
% \begin{macro}{\HyperTOCmeeting}
% \begin{macro}{\HyperTOCeditors}
% \begin{macro}{\HyperTOCpreamble}
% These commands are used to write information to the \texttt{.htoc}
% file, which is used in generating a Web Page of the Table of
% Contents. Since they only ever appear preceded by |\protect| in the
% argument to an |\addtocontents| command,  they should never need to
% be expanded by the macros. However, we define them here just in case
% \dots
%    \begin{macrocode}
\newcommand{\HyperTOCEntry}[3]{\relax}
\newcommand{\HyperTOCpart}[1]{\relax}
\newcommand{\HyperTOCmeeting}[3]{\relax}
\newcommand{\HyperTOCeditors}[1]{\relax}
\newcommand{\HyperTOCpreamble}{\relax}
%    \end{macrocode}
% \end{macro}
% \end{macro}
% \end{macro}
% \end{macro}
% \end{macro}
%

% \subsection{Formatting the poster abstracts} 
% These can now be used without using the book option.
%    \begin{macrocode}

\newcommand{\@makeminititle}{%
  \newpage
  \null 
  \vskip 3em%
  \begin{center}%
    {\@title \par}%
  \end{center}
  \vskip 3em%
  }

\newif\if@abstractsads
\@abstractsadsfalse
\newcommand\abstractsADS{\@abstractsadstrue}

\if@book\RequirePackage{multicol}\fi
\newenvironment{abstracts}[1][Abstracts of Contributed Papers]{%
  %% First, make sure we start on an odd page (06 Jun 2000)
  \if@book\clearpage\thispagestyle{empty}\cleardoublepage\fi
  %%
  \renewcommand{\articlename}{ABSTRACTS}
  \title{#1}\shortauthor{Abstracts}\shorttitle{Abstracts}
  \listofauthors{}
  \if@book
  \addcontentsline{toc}{chapter}{\protect\@fiddleauthors{\@TOCtitle}{\rmaaTOCauthors}}%
  \addtocontents{htc}{\protect\HyperTOCEntry{\@rawtitle}{\@rawauthors}{\articlename}}
  \fi
  \pagestyle{rmaa} % this now seems necessary (06 Jun 2000)
  \thispagestyle{rmaafirstpage}
%%  \twocolumn[\@makeminititle]
  \renewcommand{\abstract}[1]{%
    \renewcommand\ADS@abstract{##1}%
    \renewcommand{\@abstract}{##1}\@writeabstract%
  }
  \renewcommand{\title}[2][]{%
    \gdef\ADS@authors{}%
    \gdef\ADS@affil{}%
    \gdef\ADS@references{}%
    \gdef\ADS@keywords{}%
    \renewcommand{\@rmaaArg}{##1}
    \ifx\@rmaaArg\@rmaaEmpty
       \renewcommand{\@TOCaftertitle}{}
    \else
       \renewcommand{\@TOCaftertitle}{ (\@rmaaArg)}
    \fi
    \renewcommand{\@title}{\rmaa@TOCTitleStyle{##2}}%
    \renewcommand{\@TOCtitle}{\rmaa@TOCTitleStyle{##2}\@TOCaftertitle}%
    \renewcommand{\rmaaIndexTitle}{\rmaa@IndexTitleStyle{##2\@TOCaftertitle.}}%
    \renewcommand\ADS@title{##2}%
    \xdef\ADS@firstpage{\@arabic{\c@page}}%
    }
  \renewcommand{\author}[1]{\let\@thanks\@empty\renewcommand\@author{##1} }
  \renewcommand{\footnotetext}[2][XX]{\par\noindent\textsuperscript{##1}~##2}
  \renewcommand{\footnoterule}{\kern-3\p@ \hrule \@width 3em \kern 2.6\p@}
  \if@book\@begofarticlehook\fi
  \if@book\begin{multicols}{2}[\@makeminititle]\else\@makeminititle\fi
  \begingroup
  \setcounter{footnote}{0}
  }% and at the end of the environment....
  {\endgroup\if@book\end{multicols}\listofabstracts\label{RMAALastPage}\clearpage\fi}
%% auxiliary commands used:
%% \@writeabstract does all the actual formatting 
\newcommand{\@writeabstract}{%
  \if@book%
  \addcontentsline{loa}{abstract}{\protect\@fiddleauthors{\@TOCtitle}{\rmaaTOCauthors}}%
  \fi%
  \begin{center}\@title\par\@author\end{center}
  \@abstract
  \par\noindent\footnoterule\nobreak\footnotesize\@thanks
  \par\noindent\hrulefill\goodbreak
  \vskip 1cm \@plus0.5cm
  %% write ADS entry
  \if@book\if@abstractsads\ADS@dumpdata\fi\fi%
  \endgroup % this pair is to make each abstract a self-contained group
  \begingroup
  \setcounter{footnote}{0}
  }

\if@book
%% List of abstracts (stuff written to .loa file)
\newcommand{\listofabstracts}{%
  \begin{multicols}{2}[%
    \vskip 1.5cm \vfill % leave at least 1.5 cm and go to bottom of page if possible
    \centerline{\MakeUpperCase{List of Abstracts}}
    \vskip 1.5cm \@plus0.5cm
    ][6cm] % this makes it start a new page if there is not at least 6cm of space
    \renewcommand{\@fiddleauthors}[2]{##1\hskip1em plus2em minus0.4em\penalty-100  ##2}
    \parindent\z@
%%    \parskip.5\p@ \@plus .7\p@\relax 
    \parskip3\p@ \@plus 1\p@ \@minus 1\p@\relax % allow quite a lot of
                                % leeway here to avoid breaks inside
                                % entries 
    \renewcommand{\\}{\relax} % temporarily turn off explicit line breaks just in case
    \renewcommand{\altaffilmark}[1]{\relax} % ditto with this
    \@starttoc{loa}
  \end{multicols}
  }
\newcommand*{\l@abstract}[2]{%
  \par\penalty-1000
  #1\hskip 1em\penalty0\null\nobreak\hfill#2}
\fi

%    \end{macrocode}
 
% \subsection{Code for typesetting an entire volume of conference
% proceedings} \label{sec:whole-volume}
%    \begin{macrocode}

\if@book
%% environment for including the individual articles
\newcommand{\@fiddleauthors}[2]{\relax}
\renewcommand{\theequation}{\@arabic\c@equation}
\newenvironment{articles}%
{ \begingroup % make it a group so that author definitions do not
              % spill over
  \gdef\ADS@authors{}
  \pagestyle{rmaa}
  \renewenvironment{document}{%
    \@begofarticlehook
    \addtocounter{chapter}{1}%
    \addcontentsline{toc}{chapter}{%
      \protect\@fiddleauthors{\@TOCtitle}{\rmaaTOCauthors}}%
    \addtocontents{htc}{\protect\HyperTOCEntry{\@rawtitle}{\@rawauthors}{\articlename}}
    }{\relax}
  \renewcommand{\documentclass}[2][foo]{\relax} 
  \renewcommand{\usepackage}[2][foo]{%
    \def\@packagename{##2}\def\@rmaacitename{rmaacite}%
    \ifx\@packagename\@rmaacitename
    \@rmcitetrue\ClassInfo{rmaa}{This paper uses rmaacite}
    \makeatletter\input{rmaacite.sty}\makeatother % read it in here so that all changes
    \else                                         % are local to paper
    \if@debug\ClassWarning{rmaa}{Package \@packagename\space was requested by \articlename\space}\fi
    \fi} 
  %% Don't allow individual articles to mess with the pages, volume or year
  \renewcommand{\SetFirstPage}[1]{\relax}
  \renewcommand{\SetYear}[1]{\relax}
  \renewcommand{\SetVolume}[1]{\relax}
  \renewcommand{\SetConfTitle}[1]{\relax}
  %% Don't allow individual articles to change these
  \renewcommand{\nonstopmode}{\relax}
  \renewcommand{\scrollmode}{\relax}
  \renewcommand{\outputplates}{\relax}
  %% XXXX need to check if using rmaacite and act accordingly
  %% avoid multiple declarations of the same length
  \let\oldnewlength=\newlength
  \renewcommand{\newlength}[1]{%
    \edef\reserved@a{\expandafter\@gobble\string ##1}%
    \@ifundefined{\reserved@a}{%
      \oldnewlength{##1}}{%
      \if@debug{\ClassInfo{rmaa}{\reserved@a\space already defined: doing nothing}}\fi%
      \relax}}%
  \renewcommand\newcommand{\providecommand}
  }%
{\endgroup\onecolumn}
     

%    \end{macrocode}
% Command to be used in `articles' environment for including an
% article
%    \begin{macrocode}
\newcommand{\article}[2][]{%
  % optional first argument is extra text to put in brackets after title in TOC
  % e.g. ``Invited Talk'' or ``Poster''
  \renewcommand{\articlename}{#2}
  % reset loads of counters
  \setcounter{footnote}{0}
  \setcounter{figure}{0}
  \setcounter{table}{0}
  \setcounter{section}{0}
  \setcounter{equation}{0}
  \setcounter{rmaabib}{0}
  \setcounter{author}{0}
  \setcounter{AUTHOR}{1}
  \gdef\ADS@authors{}%
  \gdef\ADS@affil{}%
  \gdef\ADS@affiliations{}%
  \gdef\ADS@references{}%
  \@firstauthortrue
  \@firstaffiltrue
  \@firstaltaffiltrue
  \@firstkwdtrue
  \@rmcitefalse
  \if@suppressfulladdresses\else\@needaddressestrue\fi
  \renewcommand{\@rmaaArg}{#1}
  \ifx\@rmaaArg\@rmaaEmpty
     \renewcommand{\@TOCaftertitle}{}
  \else
     \renewcommand{\@TOCaftertitle}{ (#1)}
  \fi
%    \end{macrocode}
% In the main journal articles must start on an odd page. 
%    \begin{macrocode}
  \if@journal\clearpage\thispagestyle{empty}\cleardoublepage\fi
%% Now input the file, then run the end-of-article hook (reinstated 28
  %% Sep 2001)
  {\input{\@articledir/\articlename.tex}\@endofarticlehook}}
%    \end{macrocode}
% Command to separate articles into sessions. Write an entry in the
% TOC, makes sure we start on an odd page (added 06 Jun 2000), then do
% a title page if \cs{SessionTitlePages} was called. 
% \changes{3.12}{04 Feb 2003}{Added option to have title page for each
% session. Turn it on with \cs{SessionTitlePages}.}
%    \begin{macrocode}
\newcommand{\session}[1]{
  %% this does nothing except write an entry in the TOC
  \addcontentsline{toc}{part}{\rmaa@TOCPartStyle{#1}}
  \addtocontents{htc}{\protect\HyperTOCpart{#1}}
  %% and also makes sure we start on an odd page (added 06 Jun 2000)
  \clearpage\thispagestyle{empty}\cleardoublepage
  \stepcounter{part}
  \if@sessiontitlepages{%
    \onecolumn
    \thispagestyle{empty}
    \renewcommand{\articlename}{Section\thepart TitlePage}
    \label{RMAAFirstPage}
    \make@session@title{#1}
    \label{RMAALastPage}
    \clearpage\thispagestyle{empty}\cleardoublepage
  }\fi
  }
\newcommand{\make@session@title}[1]{%
  {%
    \setlength\parindent{0pt}
    \vspace*{\stretch{1}}
    \rule{\linewidth}{1mm}
    \begin{flushright}\Huge\bfseries\boldmath
      #1
    \end{flushright}
    \rule{\linewidth}{1mm}
    \vspace*{\stretch{2}}
  }
}
%    \end{macrocode}
% Command to turn on separate title pages for each session. 
%    \begin{macrocode}
\newcommand{\SessionTitlePages}{\@sessiontitlepagestrue}
\newcommand{\BoldSessionTitles}{%
  \renewcommand{\rmaa@TOCPartStyle}[1]{\textbf{\rmaa@TOCTitleStyle{##1}}}
  }
\newcommand{\ArticleDir}[1]{\renewcommand{\@articledir}{#1/}\graphicspath{{1/}}}
\newcommand{\@articledir}{.}
%\newcommand{\ps@rmaafront}{\ps@rmaa}% page style for front matter (same as ever, actually)
\newcommand{\ps@rmaafront}{% page style for front matter:  page # in each foot
  \renewcommand\@oddfoot{\hfil{\thepage}}%
  \renewcommand\@evenfoot{{\thepage}\hfil}%
  % shorttitle in both heads
  \renewcommand\@evenhead{\hfil\@shorttitle\hfil}%
  \renewcommand\@oddhead{\@evenhead}
  }
\newcommand{\ps@rmaafrontfirstpage}{% page style for first page of front matter
  % page # in each foot, on outside edge
  \renewcommand\@oddfoot{\hfil\thepage}%
  \renewcommand\@evenfoot{\thepage\hfil}%
  % nothing in the even heads
  \renewcommand\@evenhead{\relax}%
  \renewcommand\@oddhead{\relax}%
  }
%% slight modification to the standard \frontmatter and \mainmatter
%% commands, just to set the \pagestyle
\renewcommand\frontmatter{%
  \latexonecolumn
  \cleardoublepage
  \@frontmattertrue
  \@mainmatterfalse
  \pagenumbering{roman}
  \pagestyle{rmaafront}
  \addtocounter{page}{2} % leave space for the title page
  }
\renewcommand\mainmatter{%
  \pagestyle{empty}
  \if@rmtwocol\twocolumn\fi
  \cleardoublepage
  \@frontmatterfalse
  \@mainmattertrue
  \pagenumbering{arabic}
  \pagestyle{rmaa}
  }
\renewcommand\backmatter{%
  \clearpage\pagestyle{empty}\cleardoublepage
  \@mainmatterfalse
  %% put extra space in the TOC
  \addtocontents{toc}{\bigskip}
  }
%    \end{macrocode}
%
% \subsection{Front sections}\label{sec:front-sections}
% \begin{macro}{\notitle}
% \begin{macro}{\shortname}
% \begin{macro}{frontsection}
% The environment \texttt{frontsection} is for doing a section in the
% front matter (preface, list of participants, or whatever). It has 
% one optional argument:
% \begin{enumerate}
% \item[\#1] Short title (no embedded spaces) for
% use in making PS files. E.g., ``list\_of\_participants''.   
% \end{enumerate}
% and two mandatory arguments: 
% \begin{enumerate}
% \item[\#2] Title of section (e.g., `Preface')
% \item[\#3] Commands to be evaluated before processing the contents
% of the environment. This is mainly useful for putting in the things
% that would go in the preamble of a normal article  (e.g.,
% |\listofauthors|, |\indexauthor|, |\notitle|) 
% \end{enumerate}
%    \begin{macrocode}
\newcommand{\notitle}{\@titlevisiblefalse} \newcommand{\shortname}{}
\newenvironment{frontsection}[3][]{%
  \begingroup
  \renewcommand{\shortname}{#1}
  \ifx\shortname\@rmaaEmpty
     \renewcommand{\articlename}{#2}
  \else
     \renewcommand{\articlename}{\shortname}
  \fi
  \@begofarticlehook
  \@titlevisibletrue % default
  \thispagestyle{rmaafrontfirstpage}
  \listofauthors{} % blank out this in case it is not specified
  \title{{#2}}\shorttitle{#2}\shortauthor{#2}
  #3 % evaluate the command given as second mandatory arg
  \addcontentsline{toc}{chapter}{\protect\@fiddleauthors{\@TOCtitle}{\rmaaTOCauthors}}%
  \addtocontents{htc}{\protect\HyperTOCEntry{\@rawtitle}{\@rawauthors}{\articlename}}
  \if@titlevisible
     \vspace*{1cm}\centerline{\rmaa@FrontSecStyle{\@title}}\vspace*{1cm}
  \fi
  \setlength{\parskip}{3ex}
  }{\frontlabel{RMAALastPage}\clearpage\endgroup}
%    \end{macrocode}
% \end{macro}
% \end{macro}
% \end{macro}
%    \begin{macrocode}
%% Definition of the command \frontlabel, which rights out something like `F4' to
%% the .aux file, instead of `\textlatin  {i}'. This is so it is easy to post-process
%% the file for making PS files of individual sections
\newcommand{\frontlabel}[1]{%
  \@bsphack
  \protected@write\@auxout{}%
     {\string\newlabel{\articlename:#1}{{}{F\@arabic{\c@page}}}}
  \@esphack
  }
%% User command to say X, Y, Z,\\ Editors of the Proceedings
\newcommand{\EditorSignature}{\vspace*{\parskip}\rightline{\@editors}%
  \rightline{Editors of the proceedings}}
%% Fiddle with \thesection, etc. so that they don't include the chapter number
%% Basically go back to the `article' definitions of these
\renewcommand{\thesection}{\@arabic\c@section}
\renewcommand{\thefigure}{\@arabic\c@figure}
\renewcommand{\thetable}{\@arabic\c@table}

\newcommand{\@oldargs}{}
\newcommand{\@newargs}{}
%% add TOC entry for each paper
%\renewcommand{\title}[1]%
%{\renewcommand{\@title}{\rmaa@TitleStyle{#1}}%
%\addtocounter{chapter}{1}\addcontentsline{toc}{chapter}{#1}%
%}

%% Formatting of the table of contents
%% dimensions of various parts
\newlength{\TOCtextwidth}\setlength{\TOCtextwidth}{12.65cm}
\newlength{\TOCpagenumwidth}\setlength{\TOCpagenumwidth}{1.3cm}
\newlength{\TOCentrywidth}\setlength{\TOCentrywidth}{\TOCtextwidth}
\addtolength{\TOCentrywidth}{-\TOCpagenumwidth} 
\newlength{\TOCextramargin}\setlength{\TOCextramargin}{1.25cm}
\newcommand{\TOCbox}[1]{\hspace{\TOCextramargin}\parbox{\TOCtextwidth}{#1}\par}
\newlength{\TOCextramarginx}\setlength{\TOCextramarginx}{\TOCextramargin}\addtolength{\TOCextramarginx}{4.3cm}
\newlength{\TOCtextwidthx}\setlength{\TOCtextwidthx}{\TOCtextwidth}\addtolength{\TOCtextwidthx}{-4.3cm}
\newcommand{\TOCboxx}[1]{\hspace{\TOCextramarginx}\parbox{\TOCtextwidthx}{#1}\par}

\newcommand{\RMAAbannerimage}{rmsc}
\renewcommand\tableofcontents{%
  \renewcommand{\articlename}{CONTENTS}
  \begingroup
  \@begofarticlehook
  \if@twocolumn
     \@restonecoltrue\onecolumn
  \else
     \@restonecolfalse
  \fi
  %% 9pt for journal style
  \if@proceedings\else\small\fi
  \if@newstyle\small\fi %% same for new style
  \if@proceedings\thispagestyle{rmaafrontfirstpage}\else\thispagestyle{empty}\fi
  \shortauthor{Contents}\shorttitle{Contents}
  \vspace*{18mm}
  \if@printbanner
  %% put in the banner header (need to get better image of this)
     \TOCbox{\hspace*{-0.15cm}\includegraphics[width=0.9\TOCtextwidth]{\RMAAbannerimage}}
     \vspace*{8mm}
  %% volume and publication date (should really be in funny font)
     \TOCboxx{\@volumen\\[2mm] \@mes\ \@anyo}
  \else
  %% put in the banner header (need to get better image of this)
     \TOCbox{\hspace*{-0.15cm}\includegraphics[width=0.9\TOCtextwidth]{\RMAAbannerimage}}
     \vspace*{12.5mm}
  %% volume and publication date (should really be in funny font)
     \TOCboxx{~\\[2mm] ~}
  \fi
  \vspace*{16mm}
  %% This is the command that fiddles with the TOC entry so that the
  %% title is flushleft but the authors are flushright
  \renewcommand{\@fiddleauthors}[2]{%
    % first put the paper title
    \renewcommand{\\}{\relax} % temporarily turn off explicit line breaks just in case
    \renewcommand{\altaffilmark}[1]{\relax} % ditto with this
    \renewcommand{\textsuperscript}[1]{\relax} % ditto with this 
    ##1
    % now save the author list in a box
    \sbox\@tempboxa{##2}%
    % check that box is shorter than the total width available for the entry
    \ifdim \wd\@tempboxa <\TOCentrywidth
       % if so, then put it right here, flush against right margin
       \penalty10000\hskip1em plus10em\penalty0\hbox{}\nobreak\hfill\box\@tempboxa
       % Note that we force a space of at least 1em between title and authorlist, 
       % stretchable to 11em since we may have very short titles but still with an 
       % authorlist too long to go on the same line.  The only linebreak allowed is 
       % directly before the authorlist.
    \else
       % if it is too long, we set the original authorlist text in its own parbox
       \par\vskip1ex\parbox[b]{\TOCentrywidth}{\raggedleft ##2}
    \fi
    }
  \@starttoc{toc}%
  \if@restonecol\twocolumn\fi
  \frontlabel{RMAALastPage}
  \clearpage
  \thispagestyle{empty}
  \cleardoublepage % make sure we start on an odd page
  \endgroup

  %%% ALSO, do the HTML TOC
  %% The file that will be used to make the HTML TOC
  \@starttoc{htc}
  %% All the stuff in this file does nothing in the book style
  \addtocontents{htc}{\protect\HyperTOCmeeting{\@conftitle}{\@reunionplace}{\@reuniondate}}
  \addtocontents{htc}{\protect\HyperTOCeditors{\@editors}}
  \addtocontents{htc}{\protect\HyperTOCpreamble}
  \addtocontents{htc}{\protect\HyperTOCEntry{Table of Contents}{}{\articlename}}

  %% and the ADS file
  \ADS@openfile
  }
%% This is now much simpler
\renewcommand\l@chapter[2]{%TOC entry for articles
  \ifnum \c@tocdepth >\m@ne
    \vskip 1.0em \@plus\p@
    \TOCbox{%
      \parbox[b]{\TOCentrywidth}{#1}\parbox[b]{\TOCpagenumwidth}{\hfill#2}
    }
  \fi} 
\renewcommand\l@part[2]{%TOC entry for sessions: centered, no page number
  \ifnum \c@tocdepth >\m@ne
    \vskip 1.5em \@plus1.5em \goodbreak
    \TOCbox{%
      \parbox[b]{\TOCentrywidth}{\hfil#1\hfil}\hfill
    }
  \fi\nobreak} 

%% Formatting theindex
%% strongly encourage page breaks between index entries
\renewcommand{\@idxitem}{\par\penalty-1000\hangindent.45cm} 
\renewcommand{\indexspace}{\par\penalty-1000\vskip 10\p@ \@plus15\p@ \@minus3\p@\relax}
\renewenvironment{theindex}{%
  %% first, make sure we start on an odd page (06 Jun 2000)
  \clearpage\thispagestyle{empty}\cleardoublepage
  %% this looks a bit like an article but is two columns and without abstract, etc
  \renewcommand{\articlename}{INDEX}
  \title{Author Index}\shortauthor{Author Index}\shorttitle{Author Index}
  \listofauthors{}
  \addcontentsline{toc}{chapter}{\protect\@fiddleauthors{\@TOCtitle}{\rmaaTOCauthors}}%
  \addtocontents{htc}{\protect\HyperTOCEntry{\@rawtitle}{\@rawauthors}{\articlename}}
  \pagestyle{rmaa}
  \thispagestyle{rmaafirstpage}
  \@begofarticlehook
  %% 9pt for journal style
  \if@proceedings\else\small\fi
  \if@newstyle\small\fi
  \twocolumn[\@makeminititle]
  %% vertical spacing (see also definition of \@idxitem and \indexspace
  \parskip\z@ \@plus .6\p@\relax % allow quite a lot of leaway here to avoid breaks inside entries
  %% horizontal spacing, breaking and glue parameters (see also the stuff in the definition 
  %% of \indexauthor
  \parindent\z@ 
  \exhyphenpenalty10000 % don't break on explicit hyphens, it always looks bad
  \rightskip=0pt plus.2em % allow a tiny bit of leeway at the ends of lines
  \spaceskip.3333em plus.5em minus.15em % a lot of leeway between words
  \xspaceskip.5em plus.5em minus.4em % this seems to have no effect
  \let\item\@idxitem
  \renewcommand{\\}{\relax} % temporarily turn off explicit line breaks just in case
  \renewcommand{\altaffilmark}[1]{\relax} % ditto with this
  \renewcommand{\textsuperscript}[1]{\relax} % ditto with this
  }{\label{RMAALastPage}\clearpage} % clearpage necessary so that we know about \shortauthor



%% end definitions for book styles
\fi

%    \end{macrocode}
% In order to do away with plates, we sometimes want a figure to
% appear in color only in the online version. The following command
% can bve used in a figure caption to typeset a note to this effect. 
%    \begin{macrocode}
\newcommand\ColorOnlineOnly{%
  \textsc{Note: this figure is available in color in the electronic
    version of this article, obtainable from}
  \texttt{http://\discretionary{}{}{}www.\discretionary{}{}{}astroscu.\discretionary{}{}{}unam.\discretionary{}{}{}mx/\discretionary{}{}{}\~{}rmaa/}.}
%% Plates
%% Define 2 counters: plate is the number in the whole document, plateloop is used 
%% internally by \outputplates
\newcounter{plate}
\newcounter{plateloop}
%% Currently limited to 15 plates in document
%% when we do a label, we need to write to \newlabel entries to the .aux file, one 
%% that says RMAAPlateX so that it can be referred to in \ouputplates, and one that
%% says articlename:RMAAPlateX so that the dvips program can find it and asociate it
%% with the right article
\newcommand{\@platelabeli}{\oldlabel{RMAAPlate1}\label{RMAAPlate1}}
\newcommand{\@platelabelii}{\oldlabel{RMAAPlate2}\label{RMAAPlate2}}
\newcommand{\@platelabeliii}{\oldlabel{RMAAPlate3}\label{RMAAPlate3}}
\newcommand{\@platelabeliv}{\oldlabel{RMAAPlate4}\label{RMAAPlate4}}
\newcommand{\@platelabelv}{\oldlabel{RMAAPlate5}\label{RMAAPlate5}}
\newcommand{\@platelabelvi}{\oldlabel{RMAAPlate6}\label{RMAAPlate6}}
\newcommand{\@platelabelvii}{\oldlabel{RMAAPlate7}\label{RMAAPlate7}}
\newcommand{\@platelabelviii}{\oldlabel{RMAAPlate8}\label{RMAAPlate8}}
\newcommand{\@platelabelix}{\oldlabel{RMAAPlate9}\label{RMAAPlate9}}
\newcommand{\@platelabelx}{\oldlabel{RMAAPlate10}\label{RMAAPlate10}}
\newcommand{\@platelabelxi}{\oldlabel{RMAAPlate11}\label{RMAAPlate11}}
\newcommand{\@platelabelxii}{\oldlabel{RMAAPlate12}\label{RMAAPlate12}}
\newcommand{\@platelabelxiii}{\oldlabel{RMAAPlate13}\label{RMAAPlate13}}
\newcommand{\@platelabelxiv}{\oldlabel{RMAAPlate14}\label{RMAAPlate14}}
\newcommand{\@platelabelxv}{\oldlabel{RMAAPlate15}\label{RMAAPlate15}}
\newcommand{\@outputplatelabeli}{\oldlabel{RMAAOutputPlate1}}
\newcommand{\@outputplatelabelii}{\oldlabel{RMAAOutputPlate2}}
\newcommand{\@outputplatelabeliii}{\oldlabel{RMAAOutputPlate3}}
\newcommand{\@outputplatelabeliv}{\oldlabel{RMAAOutputPlate4}}
\newcommand{\@outputplatelabelv}{\oldlabel{RMAAOutputPlate5}}
\newcommand{\@outputplatelabelvi}{\oldlabel{RMAAOutputPlate6}}
\newcommand{\@outputplatelabelvii}{\oldlabel{RMAAOutputPlate7}}
\newcommand{\@outputplatelabelviii}{\oldlabel{RMAAOutputPlate8}}
\newcommand{\@outputplatelabelix}{\oldlabel{RMAAOutputPlate9}}
\newcommand{\@outputplatelabelx}{\oldlabel{RMAAOutputPlate10}}
\newcommand{\@outputplatelabelxi}{\oldlabel{RMAAOutputPlate11}}
\newcommand{\@outputplatelabelxii}{\oldlabel{RMAAOutputPlate12}}
\newcommand{\@outputplatelabelxiii}{\oldlabel{RMAAOutputPlate13}}
\newcommand{\@outputplatelabelxiv}{\oldlabel{RMAAOutputPlate14}}
\newcommand{\@outputplatelabelxv}{\oldlabel{RMAAOutputPlate15}}
\newcommand{\@platerefi}{\oldpageref{RMAAPlate1}}
\newcommand{\@platerefii}{\oldpageref{RMAAPlate2}}
\newcommand{\@platerefiii}{\oldpageref{RMAAPlate3}}
\newcommand{\@platerefiv}{\oldpageref{RMAAPlate4}}
\newcommand{\@platerefv}{\oldpageref{RMAAPlate5}}
\newcommand{\@platerefvi}{\oldpageref{RMAAPlate6}}
\newcommand{\@platerefvii}{\oldpageref{RMAAPlate7}}
\newcommand{\@platerefviii}{\oldpageref{RMAAPlate8}}
\newcommand{\@platerefix}{\oldpageref{RMAAPlate9}}
\newcommand{\@platerefx}{\oldpageref{RMAAPlate10}}
\newcommand{\@platerefxi}{\oldpageref{RMAAPlate11}}
\newcommand{\@platerefxii}{\oldpageref{RMAAPlate12}}
\newcommand{\@platerefxiii}{\oldpageref{RMAAPlate13}}
\newcommand{\@platerefxiv}{\oldpageref{RMAAPlate14}}
\newcommand{\@platerefxv}{\oldpageref{RMAAPlate15}}
\newcounter{figplatei}
\newcounter{figplateii}
\newcounter{figplateiii}
\newcounter{figplateiv}
\newcounter{figplatev}
\newcounter{figplatevi}
\newcounter{figplatevii}
\newcounter{figplateviii}
\newcounter{figplateix}
\newcounter{figplatex}
\newcounter{figplatexi}
\newcounter{figplatexii}
\newcounter{figplatexiii}
\newcounter{figplatexiv}
\newcounter{figplatexv}
%% Printed output should use the counter `plate'
\renewcommand{\theplate}{\@arabic\c@plate}
\renewcommand{\theplateloop}{\@arabic\c@plateloop}
%% Command to output all the saved plates in this paper
\newcommand{\outputplates}{%
  \setcounter{plateloop}{0}
  \if@book\else\@endofarticlehook\fi % not entirely sure why, but this
                                     % seems to be needed
  \clearpage % first output the last page of text
  % now set pagestyle for blank even pages and make sure we start on an odd page
  \pagestyle{empty}
  \cleardoublepage
  \onecolumn
  \parindent\z@
  % first put in an extra sheet
  % the main purpose of this is so there is always a blank even page before 
  % each plate for filling in the PS file if the article has an odd number
  % of pages (see splitbook.sh)
  \null\vfil\centerline{PLATES}\vfil\clearpage\pagestyle{rmaaplate}\cleardoublepage
  % output each plate in turn
  \@whilenum\value{plateloop}<\value{plate}\do{%
    \addtocounter{plateloop}{1}
    % restore the figure counter from \@figcounti, \@figcountii, etc
    \setcounter{figure}{\value{figplate\roman{plateloop}}}
    % spew out the plate (in its own group for safety)
    {\csname @plate\roman{plateloop}\endcsname\par}
    }
  \pagestyle{rmaa}
  % reset the page number as if nothing had happened
  \addtocounter{page}{-\value{plate}}
  \addtocounter{page}{-\value{plate}}
  }
\newcommand\figcaption{%
  \def\@captype{figure}
  \renewcommand{\@makecaption}{\@makefigcap}
  \caption}
\newcommand\tabcaption{%
  \def\@captype{table}
  \renewcommand{\@makecaption}{\@maketabcap}
  \caption}

\newcommand{\plate}[2][1]{%
  %% User command to make a plate page: 
  %% Mandatory argument is what goes on the plate - this should not 
  %% include any floating environments (e.g. figure). Instead, \figcaption 
  %% should be used to do the captions. 
  %% Optional argument is the number of figures (i.e. captions) on the plate
  %% page (default is 1). This is necessary for incrementing the figure counter. 
  %%
  %% Note that we have to use a \gdef rather than an \xdef when we 
  %% save the plate output in a command. This means that in the book styles when it 
  %% comes to spewing out the plate, we will have forgotten about any user macros 
  %% that may have been defined. This can be got round by repeating the macros inside the plate.
  %% (Presumably, one would use a \newcommand rather than \renewcommand if they were new.)
  % increment global and local counters
  \addtocounter{plate}{1}
  % write appropriate label to .aux file
  \csname @platelabel\roman{plate}\endcsname
  % define a global macro called \@platei, \@plateii, \@plateiii, \@plateiv, ...
  \expandafter\gdef\csname @plate\roman{plate}\endcsname{%
    \csname @outputplatelabel\roman{plateloop}\endcsname
    \null\vfil #2 \vfil
    \cleardoublepage
    }
  % save the current shortauthor and shorttitle
  % need to be \xdef to take current values of \@shortauthor, etc 
  \expandafter\xdef\csname @shorta\roman{plate}\endcsname{\@shortauthor}
  \expandafter\xdef\csname @shortt\roman{plate}\endcsname{\@shorttitle}
  % save the current value of the figure counter for later use in \outputplates
  \setcounter{figplate\roman{plate}}{\value{figure}}
  % increment the figure counter according to the optional argument
  \addtocounter{figure}{#1}
  }
%% page style for plate pages 
\newcommand{\ps@rmaaplate}{%
  % Plate X (see page Y) in the foot
  \renewcommand\@oddfoot{\hfil\csname @shorta\roman{plateloop}\endcsname\ 
    (see page~\csname @plateref\roman{plateloop}\endcsname)\hfil}%
  \renewcommand\@evenfoot{\relax}%
  % nothing in even heads (blank pages)
  \renewcommand\@evenhead{\relax}
  % short title and page num in the odd heads
  \renewcommand\@oddhead{\hfil\csname @shortt\roman{plateloop}\endcsname%
    \hfil\llap{PLATE~\theplateloop}}
  }





%% Things for the User Guide
\if@guide
%% stolen from doc.sty
\def\MakeShortVerb#1{%
  \expandafter\ifx\csname cc\string#1\endcsname\relax
  \@shortvrbinfo{Made }{#1}%
  \add@special{#1}%
  \expandafter
  \xdef\csname cc\string#1\endcsname{\the\catcode`#1}%
  \begingroup
  \catcode`\~\active  \lccode`\~`#1%
  \lowercase{%
    \global\expandafter\let
    \csname ac\string#1\endcsname~%
    \gdef~{\verb~}}%
  \endgroup
  \global\catcode`#1\active
  \else
  \@shortvrbinfo\@empty{#1 already}%
  \fi}
\def\DeleteShortVerb#1{%
  \expandafter\ifx\csname cc\string#1\endcsname\relax
  \else
  \@shortvrbinfo{Deleted }{#1 as}%
  \rem@special{#1}%
  \global\catcode`#1\csname cc\string#1\endcsname
  \global \expandafter\let \csname cc\string#1\endcsname \relax
  \ifnum\catcode`#1=\active
  \begingroup
  \catcode`\~\active   \lccode`\~`#1%
  \lowercase{%
    \global\expandafter\let\expandafter~%
    \csname ac\string#1\endcsname}%
  \endgroup \fi \fi}
\def\@shortvrbinfo#1#2{%
  \PackageInfo{doc}{%
    #1\expandafter\@gobble\string#2 a short reference
    for \string\verb}}
\def\add@special#1{%
  \rem@special{#1}%
  \expandafter\gdef\expandafter\dospecials\expandafter
  {\dospecials \do #1}%
  \expandafter\gdef\expandafter\@sanitize\expandafter
  {\@sanitize \@makeother #1}}
\def\rem@special#1{%
  \def\do##1{%
    \ifnum`#1=`##1 \else \noexpand\do\noexpand##1\fi}%
  \xdef\dospecials{\dospecials}%
  \begingroup
  \def\@makeother##1{%
    \ifnum`#1=`##1 \else \noexpand\@makeother\noexpand##1\fi}%
  \xdef\@sanitize{\@sanitize}%
  \endgroup}
%% define | as shortcut for \verbatim
\MakeShortVerb{\|}
%% comand to specify arguments to LaTeX commands
\newcommand{\rmArg}[1]{%
  \ensuremath{\langle}{\scriptsize \textsf{\MakeUpperCase{#1}}}\ensuremath{\rangle}}
%\newcommand{\rmArg}[1]{\ensuremath{\langle}\textsc{#1}\ensuremath{\rangle}}
\providecommand{\BibTeX}{{\rmfamily B\kern-.05em%
    \textsc{i\kern-.025em b}\kern-.08em%
    T\kern-.1667em\lower.7ex\hbox{E}\kern-.125emX}}
%% Environment for code examples
\newenvironment{Code}{%
  \begin{list}{}{\setlength{\leftmargin}{0.05\columnwidth}} 
  \item[]}{\end{list}}
\fi
%</package>
%    \end{macrocode}

% \section{User commands for astronomical articles}
% \label{sec:user-commands}
% Some random definitions. These go both into the main class file and
% also into |adshtml.sty|.
% 
% Definitions of the journal name in different incarnations.  
%    \begin{macrocode}
%<*package|adshtml-style>
\newcommand{\RMAA}{RMxAA}
\newcommand{\RMAAlong}{Revista Mexicana de Astronom\'{\i}a y Astrof\'{\i}sica}
\newcommand{\RMAASC}{RMxAC}
\newcommand{\RMAASClong}{\RMAA\ (Serie de Conferencias)}
%    \end{macrocode}
% \begin{macro}{\ibidrule}
% This puts a horizontal rule with a period after it. It is used when
% the author list of a reference is identical to the preceding one. 
%    \begin{macrocode}
\newcommand{\ibidrule}{\rule[-0.03cm]{4.7em}{.01cm}.\ }
%    \end{macrocode}
% \end{macro}
% \begin{macro}{\ThisVolume}
% References to other articles in the same volume can be coded as,
% e.g., |\bibitem{} Bloggs, J. \ThisVolume{bloggs}|.
%    \begin{macrocode}
\newcommand{\ThisVolume}[1]{%
  \@year, \if@proceedings\RMAASC\else\RMAA\fi, \@volume,
  \if@book\oldpageref{#1:RMAAFirstPage}\else 000\fi\ (this volume)}
%    \end{macrocode}
% \end{macro}
% Other stuff that may be useful in ref lists.
%    \begin{macrocode}
\newcommand{\etal}{et al.\@}
%    \end{macrocode}
% Various new commands from AAS macros
%    \begin{macrocode}
\def\acknowledgments{\vskip 3ex plus .8ex minus .4ex}
\let\acknowledgements=\acknowledgments                  % second spelling
\let\jnl@style=\rm
\def\ref@jnl#1{{\jnl@style#1}}
\DeclareRobustCommand\aj{\ref@jnl{AJ}}                   
\DeclareRobustCommand\araa{\ref@jnl{ARA\&A}}             
\DeclareRobustCommand\apj{\ref@jnl{ApJ}}                 
\DeclareRobustCommand\apjl{\ref@jnl{ApJ}}                
\DeclareRobustCommand\apjs{\ref@jnl{ApJS}}               
\DeclareRobustCommand\ao{\ref@jnl{Appl.~Opt.}}           
\DeclareRobustCommand\apss{\ref@jnl{Ap\&SS}}             
\DeclareRobustCommand\aap{\ref@jnl{A\&A}}                
\DeclareRobustCommand\aapr{\ref@jnl{A\&A~Rev.}}          
\DeclareRobustCommand\aaps{\ref@jnl{A\&AS}}              
\DeclareRobustCommand\azh{\ref@jnl{AZh}}                 
\DeclareRobustCommand\baas{\ref@jnl{BAAS}}               
\DeclareRobustCommand\jrasc{\ref@jnl{JRASC}}             
\DeclareRobustCommand\memras{\ref@jnl{MmRAS}}            
\DeclareRobustCommand\mnras{\ref@jnl{MNRAS}}             
\DeclareRobustCommand\pra{\ref@jnl{Phys.~Rev.~A}}        
\DeclareRobustCommand\prb{\ref@jnl{Phys.~Rev.~B}}        
\DeclareRobustCommand\prc{\ref@jnl{Phys.~Rev.~C}}        
\DeclareRobustCommand\prd{\ref@jnl{Phys.~Rev.~D}}        
\DeclareRobustCommand\pre{\ref@jnl{Phys.~Rev.~E}}        
\DeclareRobustCommand\prl{\ref@jnl{Phys.~Rev.~Lett.}}    
\DeclareRobustCommand\pasp{\ref@jnl{PASP}}               
\DeclareRobustCommand\pasj{\ref@jnl{PASJ}}               
\DeclareRobustCommand\qjras{\ref@jnl{QJRAS}}             
\DeclareRobustCommand\skytel{\ref@jnl{S\&T}}             
\DeclareRobustCommand\solphys{\ref@jnl{Sol.~Phys.}}      
\DeclareRobustCommand\sovast{\ref@jnl{Soviet~Ast.}}      
\DeclareRobustCommand\ssr{\ref@jnl{Space~Sci.~Rev.}}     
\DeclareRobustCommand\zap{\ref@jnl{ZAp}}                 
\DeclareRobustCommand\nat{\ref@jnl{Nature}}              
\DeclareRobustCommand\iaucirc{\ref@jnl{IAU~Circ.}}
\DeclareRobustCommand\aplett{\ref@jnl{Astrophys.~Lett.}}
\DeclareRobustCommand\apspr{\ref@jnl{Astrophys.~Space~Phys.~Res.}}
\DeclareRobustCommand\bain{\ref@jnl{Bull.~Astron.~Inst.~Netherlands}}
\DeclareRobustCommand\fcp{\ref@jnl{Fund.~Cosmic~Phys.}}
\DeclareRobustCommand\gca{\ref@jnl{Geochim.~Cosmochim.~Acta}}
\DeclareRobustCommand\grl{\ref@jnl{Geophys.~Res.~Lett.}}
\DeclareRobustCommand\jcp{\ref@jnl{J.~Chem.~Phys.}}      
\DeclareRobustCommand\jgr{\ref@jnl{J.~Geophys.~Res.}}    
\DeclareRobustCommand\jqsrt{\ref@jnl{J.~Quant.~Spec.~Radiat.~Transf.}}
\DeclareRobustCommand\memsai{\ref@jnl{Mem.~Soc.~Astron.~Italiana}}
\DeclareRobustCommand\nphysa{\ref@jnl{Nucl.~Phys.~A}}
\DeclareRobustCommand\physrep{\ref@jnl{Phys.~Rep.}}
\DeclareRobustCommand\physscr{\ref@jnl{Phys.~Scr}}
\DeclareRobustCommand\planss{\ref@jnl{Planet.~Space~Sci.}}
\DeclareRobustCommand\procspie{\ref@jnl{Proc.~SPIE}}
\let\astap=\aap
\let\apjlett=\apjl
\let\apjsupp=\apjs
\let\applopt=\ao
\def\phn{\phantom{0}}
\def\phd{\phantom{.}}
\def\phs{\phantom{\ensuremath{-}}}
\def\phm#1{\phantom{#1}}
%    \end{macrocode}
% All these are redone as robust and using more modern commands where
% possible. The robustness is so they will not get expanded if used in
% abstracts and titles when these are passed to the ADS file. 
% \changes{3.16}{2003/04/07}{
%  Added \cs{Color} command for doing photometric colors.}
%    \begin{macrocode}
%%\def\sun{\hbox{$\odot$}}
\DeclareRobustCommand\sun{\ensuremath{\odot}}
%%\def\earth{\hbox{$\oplus$}}
\DeclareRobustCommand\earth{\ensuremath{\oplus}}
%%\def\lesssim{\mathrel{\hbox{\rlap{\hbox{\lower4pt\hbox{$\sim$}}}\hbox{$<$}}}}
%%\def\gtrsim{\mathrel{\hbox{\rlap{\hbox{\lower4pt\hbox{$\sim$}}}\hbox{$>$}}}}
\DeclareRobustCommand\lesssim{\mathrel{\hbox{\smaller\smaller\rlap{\hbox{\lower3pt\hbox{$\sim$}}}\raise1.5pt\hbox{$<$}}}}
\DeclareRobustCommand\gtrsim{\mathrel{\hbox{\smaller\smaller\rlap{\hbox{\lower3pt\hbox{$\sim$}}}\raise1.5pt\hbox{$>$}}}}
%%\def\lta{\mathrel{\spose{\lower 3pt\hbox{$\mathchar"218$}}
%%     \raise 2.0pt\hbox{$\mathchar"13C$}}}
%%\def\gta{\mathrel{\spose{\lower 3pt\hbox{$\mathchar"218$}}
%%     \raise 2.0pt\hbox{$\mathchar"13E$}}}
\DeclareRobustCommand\sq{\hbox{\rlap{$\sqcap$}$\sqcup$}}
%%\def\arcdeg{\hbox{$^\circ$}}
%%\def\arcmin{\hbox{$^\prime$}}
%%\def\arcsec{\hbox{$^{\prime\prime}$}}
\DeclareRobustCommand\arcdeg{\ensuremath{^\circ}}
\DeclareRobustCommand\arcmin{\ensuremath{^\prime}}
\DeclareRobustCommand\arcsec{\ensuremath{^{\prime\prime}}}
%%\def\fd{\hbox{$.\!\!^{\rm d}$}}
%%\def\fh{\hbox{$.\!\!^{\rm h}$}}
%%\def\fm{\hbox{$.\!\!^{\rm m}$}}
%%\def\fs{\hbox{$.\!\!^{\rm s}$}}
%%\def\fdg{\hbox{$.\!\!^\circ$}}
%%\def\farcm{\hbox{$.\mkern-4mu^\prime$}}
%%\def\farcs{\hbox{$.\!\!^{\prime\prime}$}}
%%\def\fp{\hbox{$.\!\!^{\scriptscriptstyle\rm p}$}}
\DeclareRobustCommand\fd{\ensuremath{.\!\!^\mathrm{d}}}
\DeclareRobustCommand\fh{\ensuremath{.\!\!^\mathrm{h}}}
\DeclareRobustCommand\fm{\ensuremath{.\!\!^\mathrm{m}}}
\DeclareRobustCommand\fs{\ensuremath{\mkern0.5mu.\mkern-5mu^\mathrm{s}\mkern-1mu}}
\DeclareRobustCommand\fdg{\ensuremath{.\!\!^\circ}}
\DeclareRobustCommand\farcm{\ensuremath{.\mkern-4mu^\prime}}
\DeclareRobustCommand\farcs{\ensuremath{\mkern1mu.\mkern-6mu^{\prime\prime}\mkern-1.5mu}}
%%\DeclareRobustCommand\farcs{\ensuremath{.\!\!^{\prime\prime}}}
\DeclareRobustCommand\fp{\ensuremath{.\!\!^{\scriptscriptstyle\rm p}}}
%%\def\micron{\hbox{$\mu$m}}
\DeclareRobustCommand\micron{\ensuremath{\mu\mathrm{m}}}
\let\la=\lesssim                        
\let\ga=\gtrsim
%%\def\case#1#2{\hbox{$\frac{#1}{#2}$}}
%%\def\slantfrac#1#2{\hbox{$\,^#1\!/_#2$}}
%%\def\onehalf{\slantfrac{1}{2}}
%%\def\onethird{\slantfrac{1}{3}}
%%\def\twothirds{\slantfrac{2}{3}}
%%\def\onequarter{\slantfrac{1}{4}}
%%\def\threequarters{\slantfrac{3}{4}}
\DeclareRobustCommand\case[2]{\ensuremath{\frac{#1}{#2}}}
\DeclareRobustCommand\slantfrac[2]{\ensuremath{\,^#1\!/_#2}}
\DeclareRobustCommand\onehalf{\slantfrac{1}{2}}
\DeclareRobustCommand\onethird{\slantfrac{1}{3}}
\DeclareRobustCommand\twothirds{\slantfrac{2}{3}}
\DeclareRobustCommand\onequarter{\slantfrac{1}{4}}
\DeclareRobustCommand\threequarters{\slantfrac{3}{4}}
%%\def\ubvr{\hbox{$U\!BV\!R$}}            
%%\def\ub{\hbox{$U\!-\!B$}}               
%%\def\bv{\hbox{$B\!-\!V$}}               
%%\def\vr{\hbox{$V\!-\!R$}}               
%%\def\ur{\hbox{$U\!-\!R$}}
%    \end{macrocode}
% New commands for ``photometric'' colors 
%    \begin{macrocode}
\DeclareRobustCommand\ubvr{\ensuremath{U\!BV\!R}}            
\DeclareRobustCommand\Color[2]{%
  \ensuremath{%
    #1\!\ifx#1V\!\fi-\!#2
  }}
\DeclareRobustCommand\ub{\Color{U}{B}}
\DeclareRobustCommand\bv{\Color{B}{V}}
\DeclareRobustCommand\vr{\Color{V}{R}}
\DeclareRobustCommand\ur{\Color{U}{R}}
%% New version of the \ion command, that actually works! 
%    \end{macrocode}
% \begin{macro}{\ion}
% This version of the |\ion| command allows you to write either
% |\ion{N}{4}| or |\ion{N}{iv}| in order to get N~IV. 
%    \begin{macrocode}
\DeclareRobustCommand\ion[2]{%
  \mbox{#1\kern0.2em%
    \smaller\rmfamily%
    \edef\@tempa{\@car#2\@nil}%
    \ifcat1\@tempa%
    \@Roman{#2}%
    \else%
    \uppercase{#2}%
    \fi}}
\newcounter{IonStage}
\DeclareRobustCommand\plainion[2]{\setcounter{IonStage}{#2}#1 \Roman{IonStage}}
%%\def\ion#1#2{\ensuremath{\mathrm{#1}\;}{\protect\small\rm\@Roman{#2}}}
%    \end{macrocode}
% \end{macro}
%    \begin{macrocode}
\DeclareRobustCommand\nodata{$\cdots$}
%    \end{macrocode}
% These are needed by the |\ion| command. Isn't there an easier way of
% doing this?
%    \begin{macrocode}
\let\@small=\small
\let\@normalsize=\normalsize
\let\@footnotesize=\footnotesize
\let\@scriptsize=\scriptsize
\let\@tiny=\tiny
\let\@large=\large
\let\@Large=\Large
\let\@LARGE=\LARGE
\let\@huge=\huge
\let\@Huge=\Huge
\newcommand{\smaller}{\small}
\renewcommand{\normalsize}{\renewcommand{\smaller}{\small}\@normalsize}
\renewcommand{\small}{\renewcommand{\smaller}{\footnotesize}\@small}
\renewcommand{\footnotesize}{\renewcommand{\smaller}{\scriptsize}\@footnotesize}
\renewcommand{\scriptsize}{\renewcommand{\smaller}{\tiny}\@scriptsize}
\renewcommand{\tiny}{\renewcommand{\smaller}{\tiny}\@tiny}
\renewcommand{\large}{\renewcommand{\smaller}{\normalsize}\@large}  
\renewcommand{\Large}{\renewcommand{\smaller}{\large}\@Large}  
\renewcommand{\LARGE}{\renewcommand{\smaller}{\Large}\@LARGE}  
\renewcommand{\huge}{\renewcommand{\smaller}{\LARGE}\@huge}  
\renewcommand{\Huge}{\renewcommand{\smaller}{\huge}\@Huge}  
%</package|adshtml-style>
%    \end{macrocode}


% \section{Other files}

% \subsection{The driver and make files}
% \label{sec:driver_make}
% These are put here for completeness, although they still need to be
% distributed separately since otherwise the user won't be able to
% unpack this file into its constituent parts.
% 
% The |Makefile| has two principal rules: |docs|, and |macros|. The
% |docs| rule produces a PDF and gzipped PS file of the documentation.
% I still haven't perfected the combination of rules required to get
% it to run \LaTeX{} enough times in all circumstances. The |macros|
% rule runs \LaTeX{} on |rmaa.ins| in order to unpack |rmaa.cls| and
% all the other little files (including the |Makefile| and the
% |rmaa.ins| files themselves\dots 
% 
% First set up lists of various types of file. 
%    \begin{macrocode}
%<*makefile>
DOCFILES=rmaa.pdf
MACRO_LTX_FILES=rmaa.cls rmaacite.sty adshtml.sty
SHELLSCRIPTS=makeads.sh
MACRO_OTHER_FILES=adshtml.perl sample.latex2html-init
MAKEFILE=Makefile
MACROFILES=$(MACRO_LTX_FILES) $(SHELLSCRIPTS) $(MACRO_OTHER_FILES)         
SKELETONFILES=ADS.tex 
%    \end{macrocode}
% Install the files read only into the package directory, specified via a path relative to the source directory. \changes{4.0}{2019/11/14}{Reorganize folder layout so that hardly anything is in the top level.  That way, things are clearer on github.}
%    \begin{macrocode}
INSTALL=install 
INSTALL_DIR=../packages
%    \end{macrocode}
% The default rule is to make everything. Installation must be
% explicitly asked for.
%    \begin{macrocode}
all: docs macros

install: install-docs install-macros

%    \end{macrocode}
% Most files are installed read-only to prevent editing. Skeleton
% files must be writeable. Must be careful not to overwrite the
% Makefile in the parent directory. 
%    \begin{macrocode}
install-macros: macros
        $(INSTALL) -m 444 $(MACRO_LTX_FILES) $(MACRO_OTHER_FILES) $(INSTALL_DIR)
        $(INSTALL) -m 555 $(SHELLSCRIPTS) $(INSTALL_DIR)
        $(INSTALL) -m 644 $(SKELETONFILES) $(INSTALL_DIR)

install-docs: docs
        $(INSTALL) rmaa.pdf $(INSTALL_DIR)/rm-fulldocs.pdf

%    \end{macrocode}
% Make the PDF and gzipped PS files. 
%    \begin{macrocode}
docs: $(DOCFILES)

%    \end{macrocode}
% Make everything else if either the |.dtx| or |.ins| files have
% changed. 
%    \begin{macrocode}
macros: rmaa.dtx rmaa.ins
%    \end{macrocode}
% Ensure that all the macro files can be written to then recreate
% them. The |touch| is there in case any of the files don't exist
% yet. 
%    \begin{macrocode}
        touch $(MACROFILES) $(MAKEFILE)
        chmod u+w $(MACROFILES) $(MAKEFILE) rmaa.ins
        latex "\input{rmaa.ins}"
%    \end{macrocode}
% \TeX{} automatically adds the |.tex| extension to extensionless
% files. We fix this and take the opportunity to convert
% back to TABs all that it has lovingly converted to spaces. 
%    \begin{macrocode}
        cp -f $(MAKEFILE) $(MAKEFILE).BCK && unexpand $(MAKEFILE).tex >| $(MAKEFILE) && rm $(MAKEFILE).tex
%    \end{macrocode}
% Make the shell scripts executable and everything unwriteable (so
% they don't get hand-edited by mistake). However, undo this for the
% files that the user is \emph{supposed} to modify. 
%    \begin{macrocode}
        chmod a+x $(SHELLSCRIPTS)
        chmod a-w $(MACROFILES) $(MAKEFILE)
        chmod u+w $(SKELETONFILES)

%    \end{macrocode}
% We now use PDF\LaTeX{} for this.
%    \begin{macrocode}
rmaa.idx: rmaa.dtx
        @echo "REGENERATING .idx FILE"
        pdflatex "\nonstopmode\input{rmaa.dtx}"

rmaa.aux: rmaa.dtx rmaa.ind rmaa.gls
        @echo "REGENERATING .aux FILE"
        pdflatex "\nonstopmode\input{rmaa.dtx}"

rmaa.ind: rmaa.idx
        @echo "REGENERATING INDEX"
        makeindex -s gind.ist rmaa

rmaa.gls: rmaa.glo
        @echo "REGENERATING CHANGELOG"
        makeindex -s gglo.ist -o rmaa.gls rmaa.glo

rmaa.pdf: rmaa.aux rmaa.ind rmaa.gls
        @echo "REGENERATING .pdf FILE"
        pdflatex "\nonstopmode\input{rmaa.dtx}"

clean: 
        rm -f rmaa.aux rmaa.dvi rmaa.ind rmaa.gls
%</makefile>
%    \end{macrocode}

% The |rmaa.ins| file uses the \textsf{docstrip} package to split out
% all the files and strip off this documentation.
%    \begin{macrocode}
%<*ins>
%%
%% Installation batch file for the RMAA LaTeX package
%%
\def\batchfile{rmaa.ins}
\input docstrip.tex
\preamble

PLEASE DO NOT EDIT THIS FILE! It is generated automatically from
the file rmaa.dtx using the accompanying Makefile and rmaa.ins 
Any hand-edited changes will be lost

A LaTeX document class for the RMxAA

Author: William Henney (w.henney@irya.unam.mx)

Latest version available from: http://www.irya.unam.mx/rmaa

Copyright (C) 1999-2019 Instituto de Astronomia, UNAM, Mexico

Permission is hereby granted, free of charge, to any person obtaining
a copy of this software and associated documentation files (the
"Software"), to deal in the Software without restriction, including
without limitation the rights to use, copy, modify, merge, publish,
distribute, sublicense, and/or sell copies of the Software, and to
permit persons to whom the Software is furnished to do so, subject to
the following conditions:

The above copyright notice and this permission notice shall be
included in all copies or substantial portions of the Software.

THE SOFTWARE IS PROVIDED "AS IS", WITHOUT WARRANTY OF ANY KIND,
EXPRESS OR IMPLIED, INCLUDING BUT NOT LIMITED TO THE WARRANTIES OF
MERCHANTABILITY, FITNESS FOR A PARTICULAR PURPOSE AND
NONINFRINGEMENT. IN NO EVENT SHALL THE AUTHORS OR COPYRIGHT HOLDERS BE
LIABLE FOR ANY CLAIM, DAMAGES OR OTHER LIABILITY, WHETHER IN AN ACTION
OF CONTRACT, TORT OR OTHERWISE, ARISING FROM, OUT OF OR IN CONNECTION
WITH THE SOFTWARE OR THE USE OR OTHER DEALINGS IN THE SOFTWARE.
\endpreamble
%    \end{macrocode}
% Clobber files without asking
%    \begin{macrocode}
\askforoverwritefalse
%    \end{macrocode}
% Make the \LaTeX{} macro files (class and styles)
%    \begin{macrocode}
\generate{% 
  \file{rmaa.cls}{\from{rmaa.dtx}{package}}
  \file{rmaa.ins}{\from{rmaa.dtx}{ins}}
  \file{rmaacite.sty}{\from{rmaa.dtx}{rmaacite}}
  \file{adshtml.sty}{\from{rmaa.dtx}{adshtml-style}}
  \file{ADS.tex}{\from{rmaa.dtx}{ads-container}}
  }
%    \end{macrocode}
% Make the other non-\LaTeX{} files. These are done separately since
% they can't use `\%\%' as a comment character for the
% preamble/postamble (and Perl files can't have the postamble I don't
% think). 
%    \begin{macrocode}
\generate{\nopreamble\nopostamble% 
  \file{Makefile}{\from{rmaa.dtx}{makefile}}
  \file{adshtml.perl}{\from{rmaa.dtx}{adshtml-perl}}
  \file{makeads.sh}{\from{rmaa.dtx}{ads-shell-script}}
  \file{sample.latex2html-init}{\from{rmaa.dtx}{latex2html-init}}
  }

%% Local Variables: 
%% mode: LaTeX
%% mode: Font-lock
%% End:
%</ins>
%    \end{macrocode}

% \subsection{Extended citation syntax}
% Modified from a file called mnrascite.sty
% I no longer remember where I stole the original version of this from
% 
% Alterations to the cite commands to make them do more sensible things.
% N.B. the bibliography style should produce entries which have labels where
% the date is delimited by angle brackets '<>' for these macros to work
% properly. They will not give very sensible error messages if you do not do 
% this properly.
% 
% A new command |\scite{key}| is defined which produces a citation which is 
% suitable for insertion in a sentence. i.e., it does not put brackets around 
% the whole citation, only the date. Also |\pcite{key}| produce output with no
% parentheses whatsoever. 
% 
% Should redefine |\@bibitem| really to have any hope of compatibility
% the command |\lbibcite| is defined to deal with the citations.
%    \begin{macrocode}
%<*rmaacite>
\def\bibitem[#1<#2>]#3{\par\hangindent=1.5em\hangafter=1
  \if@filesw
  {\def\protect##1{\string ##1\space}
    \immediate\write\@auxout{\string\lbibcite{#3}{#1}{#2}}}
  \fi}
\@redefinebibitem % what it says
%    \end{macrocode}
% Note that the |\lbibcite| defines two types of citation one contains the names,
% the other the year
%    \begin{macrocode}
\def\lbibcite#1#2#3{\global\@namedef{b@#1}{#2}
                    \global\@namedef{by@#1}{#3}}

\def\@bcite#1#2{({#1\if@tempswa , #2\fi})}
\def\@pcite#1#2{#1\if@tempswa , #2\fi}
%    \end{macrocode}
% This declaration contains what is put between the name and the date
% \changes{3.11}{2003/01/31}{Moved the \cs{articlename} to after the key for consistency
% with the new natbib stuff.}
%    \begin{macrocode}
\def\@brcite{\ (}
\def\@ebrcite{)}
\def\@fmttop#1{\@fmtcite[\csname #1\endcsname]}
\def\@citex[#1]#2{%
%%  \if@filesw\immediate\write\@auxout{\string\citation{#2}}\fi
%% WJH added 16 Jan 2002
  \@for\@citeb:=#2\do{%
    \edef\@citeb{\expandafter\@firstofone\@citeb:\articlename}%
    \if@filesw\immediate\write\@auxout{\string\citation{\@citeb}}\fi}%
%% to here
  \def\@citea{}%
  \@cite{\@for\@citeb:=#2\do
    {%
      \edef\@citeb{\expandafter\@firstofone\@citeb:\articlename}%
      \@citea\def\@citea{;\ }\@ifundefined
      {b@\@citeb}{{\bf ?}\@warning {Citation `\@citeb' on page \thepage \space undefined}}%
      {\csname b@\@citeb\endcsname\@citebro\csname by@\@citeb\endcsname\@citebrc}}}%
  {#1}}

\def\cite{\@ifnextchar [{\let\@citebro=\space
                          \let\@citebrc=\relax
                          \let\@cite=\@bcite\@tempswatrue\@citex}
                        {\let\@citebro=\space
                          \let\@citebrc=\relax
                          \let\@cite=\@bcite\@tempswafalse\@citex[]}}
\def\pcite{\@ifnextchar [{\let\@citebro=\space
                          \let\@citebrc=\relax
                          \let\@cite=\@pcite\@tempswatrue\@citex}
                        {\let\@citebro=\space
                          \let\@citebrc=\relax
                          \let\@cite=\@pcite\@tempswafalse\@citex[]}}
\def\scite{\@ifnextchar [{\let\@citebro=\@brcite
                          \let\@citebrc=\@ebrcite
                          \let\@cite=\@pcite\@tempswatrue\@citex}
                        {\let\@citebro=\@brcite
                          \let\@citebrc=\@ebrcite
                          \let\@cite=\@pcite\@tempswafalse\@citex[]}}

\renewenvironment{thebibliography}
  {\section*{REFERENCES}\vspace*{\afterbiblioskip}\small
    \if@book\RereadBiblio\fi
    \sloppy\clubpenalty4000\widowpenalty4000
    \sfcode`\.=1000\relax
    \setlength\parindent{\z@}\setlength\parskip{\z@}}
  {\relax\normalsize\@outputaddressesifnecessary} 
%</rmaacite>
%    \end{macrocode}

% \subsection{Auxilliary files for generating ADS data}

% The ADS data are written by \LaTeX{} to a |.ads| file. This file
% then has to be included in the following file (|ADS.tex|),
% which is run through \LaTeX{}2HTML (see below). 

%    \begin{macrocode}
%<*ads-container>
%%
%% Instructions:
%% Edit the \input{} line below to reflect the name of your file. 
\documentclass{article}
\usepackage[latin1]{inputenc}
\usepackage{adshtml}
\newcommand{\ADSfield}[2]
{\%#1  #2\\}
%%\newcommand{\ADSvfield}[2]
%%{\%#1 \begin{verbatim}#2\end{verbatim}}
\newenvironment{ADSarticle}{}{\par\bigskip}
\renewcommand{\bibitem}[2][]{\par}
\setlength{\parindent}{0pt}
\renewcommand{\ThisVolume}[1]{%  
  2001, \RMAASC, 12, \oldpageref{#1:RMAAFirstPage}}  
 

\begin{document}
\makeatletter
\input{rmbooksamp.ads}
\end{document}
%</ads-container>
%    \end{macrocode}

% This is the |adshtml.sty| style file used by |ADS.tex|. It just
% defines a few 
% macros that might be used in abstracts, etc (really it should define
% them all). This now gets all its code from
% Sec~\ref{sec:user-commands}. 
%    \begin{macrocode}
%<*adshtml-style>
%% Nothing else needs to go here now. 
%</adshtml-style>
%    \end{macrocode}
% The init file for latex2html.
%    \begin{macrocode}
%<*latex2html-init>
$PDFTOCAIRO = "pdftocairo";
$PDFLATEX = "pdflatex";
$LATEX = "latex";
$MAX_SPLIT_DEPTH = 0;
$HTML_VERSION = "4.0,math";
$NO_SIMPLE_MATH = 1;
$ADDRESS = "<p><address><a href=\"mailto:rmaa\@astro.unam.mx\">rmaa\@astro.unam.mx</address></p>"; 
$INFO = "<p>This document was created using latex2html and the rmaa latex macros</p>"; 
$NO_NAVIGATION = 1;
1; 
%</latex2html-init>
%    \end{macrocode}
% This is the |adshtml.perl| file that is used by \LaTeX{}2HTML. Some
% bizarre \TeX{}y command sometimes show up so we tell it to ignore
% these. All the astronomy macros that come from AAS are processed in
% \TeX{}. That way they get made to images but with an |ALT| tag that
% has the original code, which is what will be picked up by |lynx|,
% with the result that they get passed intact into the final |.txt|
% file. 
%    \begin{macrocode}
%<*adshtml-perl>
package main;

$dum = &do_cmd_makeatletter();

sub do_cmd_RMAA {
    local($_) = @_;
    "RMxAA" . $_;
}

sub do_cmd_RMAASC {
    local($_) = @_;
    "RMxAASC" . $_;
}

sub do_cmd_etal {
    local($_) = @_;
    "et al\." . $_;
}

sub do_cmd_ibidrule {
    local($_) = @_;
    "<img src=\"foo\" alt=\"\ibidrule\">" . $_;
}


sub do_cmd_bibitem {
    local($_) = @_;
    local($key);
    local($opt, $dummy)=&get_next_optional_argument; 
    $key = &missing_braces unless
        ((s/$next_pair_pr_rx/$key = $2; ''/eo)
        ||(s/$next_pair_rx/$key = $2; ''/eo));
    "<p>" . $_; 
}

&ignore_commands (<<_RAW_ARG_CMDS_);
lowercase
adjustfinalcols
textsuperscript # {}
newline # []
thinmuskip
mskip
smallskip
medskip
bigskip
hspace # {}
hspacestar # {}
vspace # {}
vspacestar # {}
mathrm
_at_mathrm
spacefactor
discretionary
unhbox 
voidb_at_x 
penalty 
_at_M
_at_m
null
IeC
_RAW_ARG_CMDS_

&process_commands_in_tex (<<_RAW_ARG_CMDS_);
sun
earth
lesssim
gtrsim
sq
argdeg
arcmin
arcsec
fd
fh
fm
fs
fdg
farcm
farcs
fp
micron
la
ga
case # {} # {}
slantfrac # {} # {}
onehalf
onethird
twothirds
onequarter
threequarters
ubvr
ub
bv
vr
ur
aj
araa
apj
apjl
apjs
ao
apss
aap
aapr
aaps
azh
baas
jrasc
memras
mnras
pra
prb
prc
prd
pre
prl
pasp
pasj
qjras
skytel
solphys
sovast
ssr
zap
nat
iaucirc
aplett
apspr
bain
fcp
gca
grl
jcp
jgr
jqsrt
memsai
nphysa
physrep
physscr
planss
procspie
_RAW_ARG_CMDS_

1;                              # This must be the last line
%</adshtml-perl> 
%    \end{macrocode}

% This is a shell script |makeads.sh| that generates the |_ads.txt| file from
% |_ads.tex| by using \LaTeX{}2HTML and |lynx|.  \changes{4.0}{2019/11/13}{Make sure that lynx writes utf-8 output instead of latin-1.}
% \changes{2.2beta13}{2002/03/12}{Complete rewrite of |makeads.sh|. Now much 
% improved I think}
% \changes{4.0}{2019/11/13}{Fix the at-sign substitution so that it only eliminates \cs{@} and does not mess with email addresses.}

%    \begin{macrocode}
%<*ads-shell-script>
#! /bin/sh
if [ -z $1 ]; then
    echo "Usage: $0 basename"
    exit
fi
BASENAME=$1
ADSBASENAME=${BASENAME}_ads
AUXFILE=$BASENAME.aux
TMPFILE=tmp.$$


echo "Generating $ADSBASENAME.tex"

# For some reason, must double-up \\ to \\\\ but not \ to \\ - odd!
cat > $ADSBASENAME.tex <<EOF
%% PLEASE DO NOT EDIT THIS FILE!  It is generated automatically from
%% the script makeads.sh.  Any hand-edited changes will be lost
%% 
\documentclass{article}
\usepackage[latin1]{inputenc}
\usepackage{html}
\usepackage{adshtml}
\newcommand{\ADSfield}[2]{\%#1  #2\\\\}
\newenvironment{ADSarticle}{}{\par\bigskip}
\setlength{\parindent}{0pt}
\newcounter{istage}
\renewcommand{\ion}[2]{\setcounter{istage}{#2}#1 \Roman{istage}}
\begin{document}
\makeatletter
EOF
echo >> $ADSBASENAME.tex "\input{$BASENAME.ads2}" 
echo >> $ADSBASENAME.tex "\end{document}"



echo "Generating $ADSBASENAME.sed" 

# Find type of journal
grep '\\documentclass\[.*proceedings.*\]' $BASENAME.tex >/dev/null \
     && JOURNAL='RMxAC' \
     || JOURNAL='RMxAA'

# Find the year and volume number
YEAR=`grep '\\SetYear' $BASENAME.tex | sed 's/.*\\SetYear{\(.*\)}.*/\1/'`
VOLUME=`grep '\\SetVolume' $BASENAME.tex | sed 's/.*\\SetVolume{\(.*\)}.*/\1/'`

# Generate an auxilliary sed script 
cat > $TMPFILE <<EOF
s|\\\\newlabel{\(.*\):RMAAFirstPage}{{.*}{\(.*\)}}|s/\\\\\\\\ThisVolume{\1}/$YEAR, RMxAC, $VOLUME, \2/g|g
EOF

# Now go through the .aux file, looking for \newlabels and use the sed 
# script just generated in order to make another sed script, one that 
# will replace \ThisVolume{LABEL} with the correct reference
grep '^\\newlabel{.*:RMAAFirstPage}' $AUXFILE \
     | sed -f $TMPFILE > $ADSBASENAME.sed && rm $TMPFILE

# Apend a further sed command to get rid of \@, since redefining 
# it in _ads.tex didn't work
echo >> $ADSBASENAME.sed 's/\\\\@//g'

echo "Generating $BASENAME.ads2 from $ADSBASENAME.sed and $BASENAME.ads" 

# Apply this to the .ads file, generating the .ads2 file, which is included
# by the _ads.tex file.
sed -f $ADSBASENAME.sed $BASENAME.ads > $BASENAME.ads2

echo "Generating $ADSBASENAME.txt from $ADSBASENAME.tex and $BASENAME.ads2"

# Now run latex2html on the _ads.tex file
latex2html -no_navigation -info 0 -address 0 -html_version 4.0,math $ADSBASENAME

# Then, convert to utf8 with lynx and tidy it up a bit
lynx --justify=off --width=100 --dump --nolist --display_charset=utf-8 \
     $ADSBASENAME/$ADSBASENAME.html | cut -c4- \
     | sed 's/ibidrule/\\ibidrule/g' > $ADSBASENAME.txt
%</ads-shell-script>
%    \end{macrocode}

% \section{Epilogue}
% Err, that's it \dots

\endinput


% \endinput
% Local Variables:
% mode: doctex
% TeX-master: t
% End:
"

rmaa.aux: rmaa.dtx rmaa.ind rmaa.gls
        @echo "REGENERATING .aux FILE"
        pdflatex "\nonstopmode% \def\fileversion{4.0} \def\filedate{2019/11/12}
% \MakeShortVerb{\|}
% \title{A \LaTeX{} document class for the Revista Mexicana de
% Astronom\'\i a y Astrof\'\i sica\thanks{This file has version
% \fileversion{} last revised \filedate{}.}}
% \author{William J. Henney\thanks{w.henney@irya.unam.mx}}
% \date{\filedate}
% \maketitle
% \paragraph{Abstract}
% This document describes the \texttt{rmaa} \LaTeX{}
% document class for the journal ``Revista Mexicana de
% Astronom\'\i a y Astrof\'\i sica''. The macros can be used for 
% \begin{enumerate}
% \item Preparing author manuscripts for the journal (including
% conference proceedings to be published in the ``Serie de
% Conferencias'')
% \item Typesetting an entire volume of the journal by inputting the
% individual article files from a container file. Facilities are
% included for automatic generation of TOC and Author Index. 
% \end{enumerate}
% 
% \tableofcontents
% \section{The documentation driver file}
% This is the part that generates the documentation that you may or
% may not be reading, which is done with the help of the Makefile
% discussed in section~\ref{sec:driver_make}. 
%    \begin{macrocode}
%<*driver>
\documentclass{ltxdoc}
\usepackage{hyperref}
\usepackage{palatino}
\EnableCrossrefs \CodelineIndex \RecordChanges
\setlength{\parindent}{0pt}
\def\MacroFont{\footnotesize\tt}
\begin{document} \DocInput{rmaa.dtx} \PrintIndex \PrintChanges \end{document}
%</driver>
%    \end{macrocode}

% \section{The implementation}
% 
% \changes{3.24}{2007/09/14}{Added user access to the macro version
% and date via \cs{rmVersion} and \cs{rmDate}}
% 
%    \begin{macrocode}
%<*package>
\newcommand{\rmaa@version}{3.26}
\newcommand{\rmaa@date}{2011/08/15}
\newcommand\rmVersion{\rmaa@version}
\newcommand\rmDate{\rmaa@date}
\NeedsTeXFormat{LaTeX2e}[1995/06/01]
%    \end{macrocode}
% \changes{4.0}{2019/11/02}{Change abstract to appear before
% resumen. Change RevMexAA to RMxAA for consistency with ADS.}
%
% \changes{3.22}{2007/03/22}{Allow proceedings articles to have
%   received/accepted dates} \changes{3.21}{2006/05/15}{Changes for
%   Silvia - abstracts are now shipped out to ADS.}
% \changes{3.19}{2006/01/19}{Changes for Silvia - remove adresses from
%   end, new abstract style} \changes{2.2beta}{2001/10/30}{ Fixed bug
%   whereby twocolumn option was not passed through to book
%   class. Finally sorted out the ``label on last page'' issue. }
% \changes{2.1}{2001/01/23}{ Fixed some inconsistencies in
%   \cs{nocite} and \cs{bibitem} uniqification Fixed the definition of
%   \cs{ion} so it will work in math mode and in contexts where the
%   typesize is different from normal (e.g., title, fig captions).
%   Modified the definition of \cs{@caption}, so we don't write any
%   list-of-figures data (we never use it anyway, and it just causes
%   problems with fragile commands in fig captions). }
% \changes{2.0}{2000/11/30}{ Added `newstyle' option, which should be
%   used in all new proceedings volumes (wider abstract, smaller
%   typeface in fig captions, etc.) Added commands for formatting
%   table captions. New demonstation document rmpostsamp.tex for
%   1-page posters. This will be the last release under a PRI
%   government (for a while...)}  \changes{1.9}{2000/06/06}{
%   Corrected position of page number in footer of articles that open
%   on even pages. Ensured that first article in each session opens on
%   an odd page, as do list of abstracts and author index. Added extra
%   space in TOC before the backmatter. Added command for editors,
%   \cs{BoldSessionTitles}, that makes produces bold titles for
%   session in the TOC.}  \changes{1.8}{2000/05/16}{ Fixed bug in
%   last page label with \cs{suppressfulladdresses} }
% \changes{1.7}{2000/05/11}{ Fixed more little things. Page labels
%   now work, even if you have no bibliography. You can output the
%   full addresses by hand with the \cs{outputfulladdresses}
%   command. Editors' Guide updated but still not finished. Revamped
%   splitbook.sh now produces PS file of each article for web page and
%   also an EPS file of each page for sending to the printers.}
% \changes{1.6}{2000/04/15}{ More cleaning up: full stops after
%   titles in index, running headers and footers in frontmatter,
%   fontsize in title of index, etc} \changes{1.5}{2000/04/15}{
%   Eliminated superscripts from titles in TOC and index}
% \changes{1.4}{2000/04/14}{ Added copyright and license info (GPL)}
% \changes{1.3}{2000/03/01}{ Corrected \cs{figcaption} to do the
%   caption formatting right Corrected font size of sec headings in
%   front matter} \changes{1.2}{2000/02/16}{ Added \cs{plate}
%   command} \changes{1.1}{2000/02/14}{ Added manuscript option,
%   corrected bugs} \changes{1.0}{2000/02/08}{ Now does a whole book
%   - lots of stuff tidied up} \changes{0.7}{1999/10/04}{ Increased
%   spacing between sections in single column styles. Suppressed
%   floats at bottom of pages (didn't work though) }
% \changes{0.6}{1999/09/24}{ Tweaked \cs{adjustfinalcols} (still not
%   perfect) Added proceedings option Added onecolumn/twocolumn
%   options (orthogonal to journal/proceedings) }
% \changes{0.5}{1999/09/23}{ Added \cs{fulladdresses} and
%   \cs{adjustfinalcols} (just about working)}
% \changes{0.4}{1999/09/22}{ Minor stylistic changes}
% \changes{0.3}{1999/06/28}{ Better mechanism for alternative
%   affiliations Added `guide' option for user guide}
% \changes{0.2}{1999/05/27}{ All basic functionality present, I
%   think} \changes{0.1}{1999/05/26}{ Just about working - still
%   missing Fig/Tab/Refs} \changes{0.0}{1999/05/21}{ So bored today,
%   decided to rewrite RMAA LaTeX macros from scratch}
%    \begin{macrocode}
\ProvidesClass{rmaa}[\rmaa@date\space v\rmaa@version\space RMxAA article
class]
%    \end{macrocode}
% \subsection{Option Handling}
% \label{sec:options}
% Declarations that are needed in options
% \changes{3.05}{2003/10/22}{Added natbib option}
% \changes{3.06}{2003/01/03}{Natbib option turned on permanently}
% \changes{3.07}{2003/01/11}{Added proof option}
% \changes{3.14}{2003/01/11}{Option to turn off natbib}
% \changes{3.18}{2003/10/22}{trying out doctex mode} 
% \changes{3.23}{2007/09/09}{Manuscript is now on by default.}
% \changes{3.26}{2011/08/15}{New option for review articles}
%    \begin{macrocode}
\newif\if@reviewarticle
\@reviewarticlefalse
\newif\if@proof
\@prooffalse
\newif\if@journal
\@journaltrue
\newif\if@printdates
\@printdatesfalse
\newif\if@proceedings
\@proceedingsfalse
\newif\if@manuscript
\@manuscripttrue
\newif\if@bigfont
\@bigfonttrue
\newif\if@guide
\@guidefalse
\newif\if@preprint
\@preprintfalse
\newif\if@rmtwocol
\@rmtwocoltrue
\newif\if@biblist
\@biblistfalse
\newif\if@book
\@bookfalse
\newif\if@natbib
\@natbibtrue
\newif\if@rmcite
\@rmcitefalse
\newif\if@titlevisible
\@titlevisibletrue
\newif\if@abstractpresent
\@abstractpresentfalse
\newif\if@printbanner
\@printbannerfalse
\newif\if@suppressfulladdresses
\@suppressfulladdressesfalse
\newif\if@newstyle
\@newstylefalse
\newif\if@debug
\@debugfalse
\newif\if@boldkeywords
\@boldkeywordsfalse
%    \end{macrocode}
% The option handling---mainly setting various booleans.
% \changes{2.2beta11}{2001/11/26}{Made twocol and
% newstyle the default for proceedings.}
% \changes{3.23}{2007/09/09}{Manuscript is now 10 pt instead of 12pt}
% \changes{3.23}{2007/09/10}{Full addresses are now automatically
% suppressed for proceedings style.}
% \changes{3.24}{2007/09/14}{New option \texttt{boldkeywords} that
% does the keywords in the old ugly way. Now off by default.}
%    \begin{macrocode}
\DeclareOption{manuscript}{\typeout{Manuscript for referee} 
  \@manuscripttrue\@journalfalse\@rmtwocolfalse
  \PassOptionsToClass{10pt}{article}}
\DeclareOption{bigfont}{\typeout{Big font manuscript requested by Gary
  Ferland} 
  \@manuscripttrue\@journalfalse\@rmtwocolfalse\@bigfonttrue
  \PassOptionsToClass{11pt}{article}}
\DeclareOption{journal}{\typeout{Main journal option}%
  \@journaltrue\@rmtwocoltrue\@printdatestrue\@manuscriptfalse
  \PassOptionsToClass{10pt}{article}}
\DeclareOption{guide}{\typeout{User guide option}%
  \@guidetrue\@journalfalse\@rmtwocoltrue\@manuscriptfalse
  \PassOptionsToClass{10pt}{article}}
\DeclareOption{preprint}{\typeout{Preprint option}%
  \@preprinttrue\@manuscriptfalse\PassOptionsToClass{10pt}{article}}
\DeclareOption{proof}{\typeout{Author proof option}%
  \@prooftrue\@manuscriptfalse}
\DeclareOption{proceedings}{\typeout{Proceedings option}%
  \@proceedingstrue\@journalfalse\@rmtwocoltrue\@newstyletrue
  \@suppressfulladdressestrue
  \@printdatesfalse\@manuscriptfalse\PassOptionsToClass{10pt}{article}} 
\DeclareOption{printdates}{\typeout{Force printing of received/accepted dates}%
  \@printdatestrue} 
\DeclareOption{book}{\typeout{Whole issue option}%
  \@booktrue} 
\DeclareOption{nonatbib}{\typeout{Turn off natbib support (required
    for compatibility with old rmaacite package).}%
  \@natbibfalse} 
\DeclareOption{newstyle}{\typeout{Use new layout style}%
  \@newstyletrue}
\DeclareOption{oldstyle}{\typeout{Use old layout style}%
  \@newstylefalse}
\DeclareOption{debug}{\typeout{Loads of annoying messages}%
  \@debugtrue}
%    \end{macrocode}
% These two need to be declared explicitly so that |\ExecuteOptions|
% actually does something.
%    \begin{macrocode}
\DeclareOption{twocolumn}{\@rmtwocoltrue}
\DeclareOption{onecolumn}{\@rmtwocolfalse}
\DeclareOption{twoside}{\PassOptionsToClass{\CurrentOption}{article}}
\DeclareOption{boldkeywords}{\@boldkeywordstrue}
%    \end{macrocode}
% Bold, upper-case keywords used to be the default. Now they must be
% turned on explicitly with the \texttt{boldkeywords} option. 
%    \begin{macrocode}
\DeclareOption{boldkeywords}{\@boldkeywordstrue}
%    \end{macrocode}
% Pass any other option along to the \texttt{article} class.
%    \begin{macrocode}
\DeclareOption*{\PassOptionsToClass{\CurrentOption}{article}}
%    \end{macrocode}
% Default options are to be two-sided and based on the standard
% \LaTeX{} \texttt{article} class,
% unless we are processing the entire book, in which case we base
% ourselves on the standard \texttt{book} class. Default is to use
% manuscript option unless otherwise specified. 
%    \begin{macrocode}
\ExecuteOptions{twoside,manuscript} \ProcessOptions\relax
\if@book
\if@rmtwocol\LoadClass[twocolumn]{book}\else\LoadClass{book}\fi
\newcommand{\abstract}{\relax}
\newcommand{\@maketitle}{\relax}
\else
\if@rmtwocol\LoadClass[twocolumn]{article}\else\LoadClass{article}\fi
\fi
%    \end{macrocode}
% Attempt to use natbib for citations 02 Jan 2003. 
% \changes{3.05}{2003/01/02}{If the natbib option is used, then set
% things up so that natbib's version of the bibliography is identical
% to the standard one.}
% \changes{3.09}{2003/01/23}{Explicitly require a recent version of
% natbib, so that we get a decent error message if the package is too
% old.} 
% \changes{3.23}{2007/09/09}{Now allow Bibtex to work properly} 
% \changes{3.25}{2009/10/01}{Change \cs{newcommand} tp \cs{def} when setting \cs{bibfont}. This allows it to work with all versions of natbib.} 
% \changes{3.25}{2010/02/06}{Updated copyright year.} 
% \changes{3.25}{2010/02/06}{Updated email and web addresses to \texttt{crya.unam.mx}.} 
%    \begin{macrocode}
\if@natbib
\RequirePackage{natbib}[1999/02/25]
\bibpunct{(}{)}{;}{a}{}{,}
\def\bibfont{\small}
\renewcommand\bibsection{\section*{REFERENCES}}
\setlength\bibsep{0pt}
\setlength\bibhang{1.5em}
\bibliographystyle{rmaa}
\fi
%    \end{macrocode}
% We want the \texttt{graphicx} package for the extended syntax of
% |\includegraphics|. We want the \texttt{ifthen} package for using in
% these macros (although we don't much yet). We explicitly load
% \texttt{babel} since some \LaTeX{} distributions messed up the
% hyphenation if we didn't.  
%    \begin{macrocode}
\RequirePackage{graphicx}
\RequirePackage{ifthen}
\IfFileExists{babel.sty}{%
%   \let\@restoreindent=\@afterindentfalse
  \RequirePackage[english]{babel}
%   \@restoreindent
  %% reinstate the definition of \@roman so as not to break .idx files
  \@ifundefined{org@@roman}{\relax}{\renewcommand{\@roman}[1]{\org@@roman{##1}}}
  }{%
  \typeout{Package babel.sty not found - don't worry about this}}
%    \end{macrocode}
% Try to load the \texttt{dcolumn} and \texttt{booktabs} packages if
% they can be found. These can produce nicer-looking tables. 
%    \begin{macrocode}
\IfFileExists{dcolumn.sty}{%
  \RequirePackage{dcolumn}
  }{%
  \typeout{Package dcolumn.sty not found - you will not be able to
    align columns on decimal points in tables}}
\IfFileExists{booktabs.sty}{%
  \RequirePackage{booktabs}
  }{%
  \typeout{Package booktabs.sty not found - your tables will not look
    as pretty as they ought too.}
  }
\IfFileExists{makeidx.sty}{%
  \RequirePackage{makeidx}
}{%
  \typeout{Package makeidx.sty not found - Author Index may not work.}
  }
%% commented out for moment since requires more modern version of LaTeX
%% \RequirePackage[latin1]{inputenc}
%    \end{macrocode}
% \changes{4.0}{2019/11/13}{Make \texttt{hyperref} and
% \texttt{etoolbox} into hard requirements, since we now use them
% extensively.}  We use hyperref to get internal and external links.
% But I don't like the colored boxes so I turn those off.  We use
% etoolbox for patching and robustifying at the moment.  In the
% future, I want to port a lot of other stuff over to it.  For
% instance, replace the ifthen package and the TeX ifs too.
%
% We target the latest version (1.7) of PDF, rather than the default
% of 1.5.  This reduces the chance of getting warnings from included
% graphics.  This is not strictly related to hyperref, but we put it
% here anyway.
%    \begin{macrocode}
\RequirePackage{etoolbox}
\RequirePackage{hyperref}
\hypersetup{hidelinks}
\pdfminorversion=7
%    \end{macrocode}
% We patch the \cs{@combinedblfloats} following a suggestion by Heiko
% Oberdiek at \url{https://tex.stackexchange.com/a/249743/3843}, which
% avoids problems when hyperlinks are split across
% pages. \changes{4.0}{2019/11/12}{Fix for issue where breaking
% hyperlink across pages gave an error.}
%    \begin{macrocode}
\urlstyle{same}
\newcommand\email{\begingroup \urlstyle{rm}\Url}
\patchcmd\@combinedblfloats{\box\@outputbox}{\unvbox\@outputbox}{}{%
  \errmessage{\noexpand\@combinedblfloats could not be patched}%
}%
%    \end{macrocode}
% \subsection{Other setup stuff}\label{sec:setup}
% Some declarations: 
%    \begin{macrocode}
\newif\if@needaddresses
\@needaddressestrue
\newif\if@frontmatter
\@frontmatterfalse
\newif\if@inpreamble
\@inpreambletrue
\newif\if@firstauthor
\@firstauthortrue
\newif\if@firstaffil
\@firstaffiltrue
\newif\if@firstaltaffil
\@firstaltaffiltrue
\newif\if@firstkwd
\@firstkwdtrue
\newif\if@noexpandADS
\@noexpandADSfalse
\newif\if@sessiontitlepages
\@sessiontitlepagesfalse
%    \end{macrocode}
% \begin{macro}{\@rmaaEmpty}
% \begin{macro}{\@rmaaArg}
% These are used in various commands to check if arg was left blank
%    \begin{macrocode}
\newcommand{\@rmaaEmpty}{}
\newcommand{\@rmaaArg}{}
%    \end{macrocode}  
% \end{macro}
% \end{macro}
%
% \subsection{Compatibility with old \LaTeX{} versions}
% \label{sec:compat}
% This package is based on the 1998/12/01 release of \LaTeX{}.
% Older releases may not define all the commands that we need.
% Luckily, we can use |\providecommand| to define these if necessary. 
% The following commands are not defined in the 1994/06/01 release. 
%
% \begin{macro}{\MakeUpperCase}
% \begin{macro}{\textsuperscript}
% \begin{macro}{\textbackslash}
% User commands:
%    \begin{macrocode}
\providecommand{\MakeUpperCase}[1]{\uppercase{#1}}
\providecommand{\textsuperscript}[1]{\ensuremath{^\textrm{#1}}}
\providecommand{\textbackslash}{\ensuremath{\backslash}}
%    \end{macrocode}  
% \end{macro}
% \end{macro}
% \end{macro}
% \begin{macro}{\protected@xdef}
% \begin{macro}{\restore@protect}
% \begin{macro}{\@unexpandable@protect}
% \begin{macro}{\hb@xt@}
% \LaTeX{} internal commands:
%    \begin{macrocode}
\providecommand{\protected@xdef}{%
        \let\@@protect\protect
        \let\protect\@unexpandable@protect
        \afterassignment\restore@protect
        \xdef}
\providecommand{\restore@protect}{\let\protect\@@protect}
\providecommand{\@unexpandable@protect}{\noexpand\protect\noexpand}
\providecommand{\hb@xt@}{\hbox to}
%    \end{macrocode}
% \end{macro}
% \end{macro}
% \end{macro}
% \end{macro}
%
% \subsection{Page layout}\label{sec:page-layout}
%
% The dimensions of the text were defined by Elizabeth. Text width is
% 169\,mm with equal margins on the two sides: 
%    \begin{macrocode}
\setlength{\textwidth}{169mm}
\setlength{\@tempdima}{\paperwidth}
\addtolength{\@tempdima}{-\textwidth}
\setlength{\oddsidemargin}{0.5\@tempdima}
\addtolength{\oddsidemargin}{-1in}
\setlength{\evensidemargin}{\oddsidemargin}
%    \end{macrocode}
% Text height is 225\,mm, top margin is 14\,mm. 
%    \begin{macrocode}
\setlength{\textheight}{225mm} 
\setlength{\topmargin}{-1in} 
\addtolength{\topmargin}{14mm}
%    \end{macrocode}
% Space between columns in two-column styles, header, footer,
% paragraph indentation:
%    \begin{macrocode}
\setlength{\columnsep}{9.5mm}
\setlength{\headsep}{0.7cm}
\setlength{\footskip}{0.75cm}
\setlength{\parindent}{1.5em}
%    \end{macrocode}
% 
% \begin{macro}{\@set@scalable@length}
% Try out a new scheme for dealing with various spacing dimensions and
% allowing the user to easily modify them. The command
% \begin{quote}
% |\set@scalable@length|\marg{name}\marg{defaultlength}\marg{stretch}\marg{shrink}
% \end{quote}
% does a 
% |\setlength|\marg{\cs{name}}\marg{defaultlength plus stretch minus shrink}
% and also defines the user command
% \textbackslash\meta{name}|Scale|\marg{factor} that can 
% be used to rescale length by \meta{factor}. Additionally, now
% defines the package command
% \textbackslash\meta{name}|@set@scale|\marg{FACTOR} 
% that also rescales the length by \meta{FACTOR}, so that the length
% becomes \meta{factor} $\times$ \meta{FACTOR} $\times$
% \meta{defaultlength}.   
%    \begin{macrocode}
\newcommand\@set@scalable@length[4]{%
  %% arg #1: NAME of length without leading `\'
  %%
  %% arg #2: default VALUE  
  \expandafter\newlength\csname @#1@default\endcsname
  \expandafter\setlength\csname @#1@default\endcsname{#2}%
  \expandafter\newlength\csname @@#1@@default\endcsname
  \expandafter\setlength\csname @@#1@@default\endcsname{#2}%
  %% arg #3: default STRETCH  
  \expandafter\newlength\csname @#1@stretch\endcsname
  \expandafter\setlength\csname @#1@stretch\endcsname{#3}%
  \expandafter\newlength\csname @@#1@@stretch\endcsname
  \expandafter\setlength\csname @@#1@@stretch\endcsname{#3}%
  %% arg #4: default SHRINK  
  \expandafter\newlength\csname @#1@shrink\endcsname
  \expandafter\setlength\csname @#1@shrink\endcsname{#4}%
  \expandafter\newlength\csname @@#1@@shrink\endcsname
  \expandafter\setlength\csname @@#1@@shrink\endcsname{#4}%
  %%
  %% User-defined scale \@NAME@scale is unity by default
  \expandafter\gdef\csname @#1@scale\endcsname{1}%
  %% Package-defined scale \@NAME@SCALE is unity by default
  \expandafter\gdef\csname @#1@SCALE\endcsname{1}%
  %%
  %% internal command \@reset@NAME should expand to 
  %% \setlength\@@NAME@@default{\@NAME@SCALE\@NAME@default}
  %% \setlength\NAME{\@NAME@scale\@@NAME@@default}
  %% in other words: scale x SCALE x default 
  \expandafter\gdef\csname @reset@#1\endcsname{%
    \expandafter\setlength\csname @@#1@@default\endcsname{%
      \csname @#1@SCALE\endcsname \csname @#1@default\endcsname
      }
    \expandafter\setlength\csname @@#1@@stretch\endcsname{%
      \csname @#1@SCALE\endcsname \csname @#1@stretch\endcsname
      }
    \expandafter\setlength\csname @@#1@@shrink\endcsname{%
      \csname @#1@SCALE\endcsname \csname @#1@shrink\endcsname
      }
    \expandafter\setlength\csname #1\endcsname{%
      \csname @#1@scale\endcsname \csname @@#1@@default\endcsname
      \@plus
      \csname @#1@scale\endcsname \csname @@#1@@stretch\endcsname
      \@minus
      \csname @#1@scale\endcsname \csname @@#1@@shrink\endcsname
      }}
  %%
  %% User command \NAMEScale should expand to 
  %% {\renewcommand\@NAME@scale{#1}\@reset@NAME}
  \expandafter\gdef\csname #1Scale\endcsname##1{%
    \expandafter\gdef\csname @#1@scale\endcsname{##1}%
    \csname @reset@#1\endcsname}
  %%
  %% Package command \NAME@set@SCALE should expand to 
  %% {\renewcommand\@NAME@SCALE{#1}\@reset@NAME}
  \expandafter\gdef\csname #1@set@SCALE\endcsname##1{%
    \expandafter\gdef\csname @#1@SCALE\endcsname{##1}%
    \csname @reset@#1\endcsname}
  %%
  %% initialize at default values with \@reset@NAME
  \csname @reset@#1\endcsname
  }
%    \end{macrocode}
% \end{macro}
% \begin{macro}{\parskip}
% First use of this new infrastructure: set interparagraph space to be
% zero but with a small amount of stretchiness.
%    \begin{macrocode}
%   \ifthenelse{\boolean{@manuscript}}{%
%     \@set@scalable@length{parskip}{0.5\baselineskip}{0.3\p@}{0.1\p@}
%   }{
    \@set@scalable@length{parskip}{0\p@}{0.3\p@}{0\p@}
%   }
%    \end{macrocode}
% \end{macro}
% 
% \begin{macro}{\abovesectionskip}
% \changes{4.0}{2019/11/12}{Use negative values for skip so that first paragraph after section headings is not indented.  This is something that I have been trying to fix for years!  It turns out that the anser was in the documentation for the standard command \cs{@startsection.}}
% \begin{macro}{\belowsectionskip}
% \begin{macro}{\aboveparaskip}
% \begin{macro}{\belowparaskip}
% Define some utility rubber lengths for use in
% section~\ref{sec:section-commands}. 
%    \begin{macrocode}
\newlength{\abovesectionskip}
\newlength{\belowsectionskip}
\newlength{\aboveparaskip}
\newlength{\belowparaskip}
%    \end{macrocode}
% The default values for these lengths are those traditionally used
% for the onecolumn style. 
%    \begin{macrocode}
\@set@scalable@length{abovesectionskip}{-3.5ex}{-4ex}{-2ex}
\@set@scalable@length{belowsectionskip}{2.3ex}{0ex}{2ex}
\@set@scalable@length{aboveparaskip}{3.25ex}{2ex}{0.5ex}
\@set@scalable@length{belowparaskip}{-1em}{\z@}{\z@}
%    \end{macrocode}
% \end{macro}\end{macro}\end{macro}\end{macro}
% We also make scalable some of the standard float parameters. These
% are the same defaults as in the standard article classes, but
% slightly more stretchy. For the twocolumn style, they are currently
% halved, see below. Note that I have also lumped |\abovecaptionskip|
% in with these, since it plays a similar role. There is also a
% |\belowcaptionskip| but it is set to zero in the standard classes. 
%    \begin{macrocode}
\@set@scalable@length{floatsep}{12\p@}{4\p@}{4\p@}
\@set@scalable@length{textfloatsep}{20\p@}{4\p@}{8\p@}
\@set@scalable@length{intextsep}{12\p@}{4\p@}{4\p@}
\@set@scalable@length{dblfloatsep}{12\p@}{4\p@}{4\p@}
\@set@scalable@length{dbltextfloatsep}{20\p@}{4\p@}{8\p@}
\@set@scalable@length{abovecaptionskip}{10\p@}{\z@}{\z@}
%    \end{macrocode}
% \begin{macro}{\@before@title@skip}         
% \begin{macro}{\@before@author@skip}        
% \begin{macro}{\@between@author@skip}       
% \begin{macro}{\@before@abstract@skip}      
% \begin{macro}{\@after@abstracthead@skip}   
% \begin{macro}{\@between@abstract@skip}     
% \begin{macro}{\@before@kwd@skip}           
% \begin{macro}{\@before@dates@skip}           
% Spacing parameters used in |\@maketitle|. These are not changed by
% default except by the user, so they don't need to go into
% |\@package@default@lengths|.  
%    \begin{macrocode}
\newlength{\@before@title@skip}         
\newlength{\@before@author@skip}        
\newlength{\@between@author@skip}       
\newlength{\@before@abstract@skip}      
\newlength{\@after@abstracthead@skip}   
\newlength{\@between@abstract@skip}     
\newlength{\@before@kwd@skip}           
\newlength{\@before@dates@skip}           
\@set@scalable@length{@before@title@skip}{2em}{\z@}{\z@}
\@set@scalable@length{@before@author@skip}{1.5em}{\z@}{\z@}
\@set@scalable@length{@between@author@skip}{0.5em}{\z@}{\z@}
\@set@scalable@length{@before@abstract@skip}{3ex}{1ex}{0.5ex}
\@set@scalable@length{@after@abstracthead@skip}{1ex}{\z@}{\z@}
\@set@scalable@length{@between@abstract@skip}{2ex}{1ex}{0.5ex}
\@set@scalable@length{@before@kwd@skip}{1ex}{0.5ex}{0.2ex}
\@set@scalable@length{@before@dates@skip}{2ex}{0.5ex}{0.2ex}
%    \end{macrocode}
% \end{macro}\end{macro}\end{macro}\end{macro}
% \end{macro}\end{macro}\end{macro}\end{macro}
% 
% Spacing above and below displayed equations also wants to be done
% like this, but it is more complicated since the standard article
% classes make these change with the fontsize. Something to do later,
% maybe. 
% 
%    \begin{macrocode}
%    \end{macrocode}
% \begin{macro}{\@package@default@lengths}
% The value actually used for these lengths is controlled by the
% \meta{SCALE} parameter, which depends on the style option that is in
% effect. This is set using the |\@package@default@lengths| command,
% which must be called every time the style changes (e.g., changing
% from onecolumn to twocolumn). 
%    \begin{macrocode}
\newcommand{\@package@default@lengths}{%
  \ifthenelse{\boolean{@rmtwocol}}{%
    %% twocolumn mode
    \abovesectionskip@set@SCALE{0.5}
    \belowsectionskip@set@SCALE{0.5}
    \aboveparaskip@set@SCALE{0.5}
    \belowparaskip@set@SCALE{1}
    \floatsep@set@SCALE{0.5}
    \textfloatsep@set@SCALE{0.5}
    \intextsep@set@SCALE{0.5}
    \dblfloatsep@set@SCALE{0.5}
    \dbltextfloatsep@set@SCALE{0.5}
    }{%
    %% onecolumn mode
    \abovesectionskip@set@SCALE{1}
    \belowsectionskip@set@SCALE{1}
    \aboveparaskip@set@SCALE{1}
    \belowparaskip@set@SCALE{1}
    \floatsep@set@SCALE{1}
    \textfloatsep@set@SCALE{1}
    \intextsep@set@SCALE{1}
    \dblfloatsep@set@SCALE{1}
    \dbltextfloatsep@set@SCALE{1}
    }%
  }
\@package@default@lengths
%    \end{macrocode}
% \end{macro}
% \begin{macro}{\RescaleLengths}
% \begin{macro}{\RescaleSecLengths}
% \begin{macro}{\RescaleFloatLengths}
% \begin{macro}{\RescaleTitleLengths}
% User commands that will fiddle with all the scalable lengths. 
%    \begin{macrocode}
\newcommand{\RescaleLengths}[1]{%
  \RescaleSecLengths{#1}\RescaleFloatLengths{#1}\RescaleTitleLengths{#1}}
\newcommand{\RescaleSecLengths}[1]{%
%% sections
  \abovesectionskipScale{#1}
  \belowsectionskipScale{#1}
  \aboveparaskipScale{#1}
  }
\newcommand{\RescaleFloatLengths}[1]{%
%% floats
  \floatsepScale{#1}
  \textfloatsepScale{#1}
  \intextsepScale{#1}
  \dblfloatsepScale{#1}
  \dbltextfloatsepScale{#1}
  \abovecaptionskipScale{#1}
  }
\newcommand{\RescaleTitleLengths}[1]{%
%% maketitle
  \@before@title@skipScale{#1}           
  \@before@author@skipScale{#1}        
  \@between@author@skipScale{#1}       
  \@before@abstract@skipScale{#1}      
  \@after@abstracthead@skipScale{#1}   
  \@between@abstract@skipScale{#1}     
  \@before@kwd@skipScale{#1}           
  \@before@dates@skipScale{#1}           
  }
%    \end{macrocode}
% \end{macro}\end{macro}\end{macro}\end{macro}
% \changes{2.2beta14}{2002/03/17}{Finished implementation of new
% infrastructure for scalable lengths. Converted the above/below
% section skips to use it. Also, all the floatsep parameters. Also,
% all the vertical lengths that are used by |\@maketitle|. }

% \begin{macro}{\MSextramargin}
% For the \texttt{manuscript} option we have an extra 0.3\,inch
% vertical margin and 1\,inch horizontal margin for marginal notes (by
% hand or with |\marginpar|. 
%
% \changes{3.23}{2007/09/09}{Reduce baselinestretch in
%   \texttt{manuscript} style to 1.2. No longer use raggedright for
%   manuscript.}
%
% \changes{3.24}{2007/09/14}{Further reduction of baselinestretch in
%   \texttt{manuscript} style to 1.1 since we now have wider margins.}
%
% \changes{3.24}{2007/09/14}{Split \texttt{manuscript}'s extra margin
%   into separate horizontal and vertical parts.}
% 
%    \begin{macrocode}
\newlength{\MSextraHmargin}\newlength{\MSextraVmargin}
\setlength{\MSextraHmargin}{1in} 
\setlength{\MSextraVmargin}{0.3in} 
\if@manuscript
   \renewcommand{\baselinestretch}{1.1}
   \addtolength{\textwidth}{-2\MSextraHmargin}
   \addtolength{\oddsidemargin}{\MSextraHmargin}
   \setlength{\evensidemargin}{\oddsidemargin}
   \addtolength{\textheight}{-2\MSextraVmargin}
   \addtolength{\topmargin}{\MSextraVmargin}
%    \raggedright
\fi
%    \end{macrocode}
% \end{macro}
% Float parameters: don't allow bottom floats by default. 
%    \begin{macrocode}
\renewcommand{\fps@figure}{tp}
\renewcommand{\fps@table}{tp}
%    \end{macrocode}
%
% \begin{macro}{\rmaa@AuthorWidthFrac}
% \begin{macro}{\rmaa@AbstractWidthFrac}
% These parameters control what fraction of the total text width can
% be used by the list of authors and the abstract/resumen,
% respectively.  
%    \begin{macrocode}
\newcommand{\rmaa@AuthorWidthFrac}{1.0}
\newcommand{\rmaa@AbstractWidthFrac}{0.75}
%    \end{macrocode}
% Use a wide abstract to save space in the new style.
%    \begin{macrocode}
\if@newstyle\renewcommand{\rmaa@AbstractWidthFrac}{1.0}\fi
%    \end{macrocode}
% \end{macro}
% \end{macro}
%
% Some things to consider in future versions. Trying to avoid too many
% widows, etc. Apparently, the following will work best with
% \cs{flushbottom}. 
%    \begin{macrocode}
%%\widowpenalty=10000
%%\clubpenalty=10000
%    \end{macrocode}
% Also, should perhaps try getting rid of underfull/overfull box warnings. 
% \changes{3.10}{2003/01/24}{Set large values for \cs{hbadness},
% \cs{vbadness}, \cs{hfuzz}, \cs{vfuzz}. To reduce amount of stuff
% that gets written to the terminal every time you compile something.}
%    \begin{macrocode}
\hbadness=10000
\vbadness=10000
\setlength\hfuzz{\maxdimen}
\setlength\vfuzz{\maxdimen}
%    \end{macrocode}
% \subsection{Various text styles for different contexts}
% \label{sec:text-styles} 
% \begin{macro}{\rmaa@AffilStyle}
% \begin{macro}{\rmaa@PagesStyle}
% \begin{macro}{\rmaa@YearStyle}
% \begin{macro}{\rmaa@VolumeStyle}
% \begin{macro}{\rmaa@VolumenStyle}
% \begin{macro}{\rmaa@JournalNameStyle}
% \begin{macro}{\rmaa@ReceivedStyle}
% \begin{macro}{\rmaa@AbstractHeaderStyle}
% \begin{macro}{\rmaa@JournalNameStyle}
% Style commands for the different fields used by |\maketitle| in
% constructing the title page of the article. These have the form
% |\rmaa@XYZStyle{TEXT}|. I'm not sure these were really a good idea,
% but never mind \dots 
%    \begin{macrocode}
\newcommand{\rmaa@AffilStyle}[1]{{\normalsize\rmfamily#1}}
\newcommand{\rmaa@PagesStyle}[1]{{\normalsize\rmfamily#1}}
\newcommand{\rmaa@YearStyle}[1]{{\normalsize\rmfamily#1}}
\newcommand{\rmaa@VolumeStyle}[1]{{\normalsize\rmfamily\bfseries#1}}
\newcommand{\rmaa@VolumenStyle}[1]{{\Large\sffamily\bfseries#1}}
\newcommand{\rmaa@JournalNameStyle}[1]{{\normalsize\itshape#1}}
\newcommand{\rmaa@ReceivedStyle}[1]{{\normalsize\rmfamily\itshape#1}}
\newcommand{\rmaa@AbstractHeaderStyle}[1]{%
  {\normalsize\rmfamily\MakeUpperCase{#1}}}
\newcommand{\rmaa@TitleStyle}[1]{%
  {\large\bfseries\rmfamily\boldmath\MakeUpperCase{#1}\par}}
%    \end{macrocode}
% \end{macro}
% \end{macro}
% \end{macro}
% \end{macro}
% \end{macro}
% \end{macro}
% \end{macro}
% \end{macro}
% \end{macro}
%
% \begin{macro}{\rmaa@FrontSecStyle}
% Headings of ``front sections'' (see section~\ref{sec:front-sections})
%    \begin{macrocode}
\newcommand{\rmaa@FrontSecStyle}[1]{%
  {\normalsize\bfseries\rmfamily\MakeUpperCase{#1}}}
%    \end{macrocode}
% \end{macro}
%
% \begin{macro}{\rmaa@TOCTitleStyle}
% \begin{macro}{\rmaa@TOCPartStyle}
% \begin{macro}{\rmaa@IndexTitleStyle}
% \begin{macro}{\rmaaTOCAuthorStyle}
% \begin{macro}{\rmaaIndexAuthorStyle}
% Commands to set style of TOC and index entries.
% Rewritten to avoid unnecessary font commands in the .toc and .ind files
%    \begin{macrocode}
\newcommand{\rmaa@TOCTitleStyle}[1]{\MakeUpperCase{#1}}
\newcommand{\rmaa@TOCPartStyle}[1]{\rmaa@TOCTitleStyle{#1}}
\newcommand{\rmaa@IndexTitleStyle}[1]{#1}
\newcommand{\rmaaTOCAuthorStyle}[1]{%
  {\itshape#1}}
\newcommand{\rmaaIndexAuthorStyle}[1]{%
  \protect\mbox{\bfseries#1}}
%    \end{macrocode}
% \end{macro}
% \end{macro}
% \end{macro}
% \end{macro}
% \end{macro}
%
% \begin{macro}{\rmaa@KeywordStyle}
% \begin{macro}{\rmaa@KeywordHeaderStyle}
% The style of the keywords depends on whether we are using the
% \texttt{newstyle} option or not. If we are, then they are small. 
%    \begin{macrocode}
\ifthenelse{\boolean{@newstyle}}{
  \newcommand{\rmaa@KeywordStyle}[1]{%
    {\small%
      \if@boldkeywords 
      \bfseries\MakeUpperCase{\raggedright #1}
      \else
      \raggedright #1
      \fi
    }}
  \newcommand{\rmaa@KeywordHeaderStyle}[1]{%
    {\small\itshape#1}}}
{
  \newcommand{\rmaa@KeywordStyle}[1]{%
    {\normalsize%
      \if@boldkeywords\bfseries\fi\raggedright #1}}
  \newcommand{\rmaa@KeywordHeaderStyle}[1]{%
    {\normalsize\itshape#1}}}
%    \end{macrocode}
% \end{macro}
% \end{macro}
%
% \subsection{Page styles}\label{sec:page-styles}
% \begin{macro}{\ps@rmaa}
% The running headers for normal pages:
%    \begin{macrocode}
\newcommand{\ps@rmaa}{%
%    \end{macrocode}
% nothing in either foot; 
%    \begin{macrocode}
  \renewcommand\@oddfoot{\relax}%
  \renewcommand\@evenfoot{\relax}%
%    \end{macrocode}
% page number and short author list in the even heads; 
%    \begin{macrocode}
  \renewcommand\@evenhead{\rlap{\thepage}\hfil\@shortauthor\hfil}%
%    \end{macrocode}
% short title and page number in the odd heads.
%    \begin{macrocode}
  \renewcommand\@oddhead{\hfil\@shorttitle\hfil\llap{\thepage}}
  }
%    \end{macrocode}
% \begin{macro}{\ps@rmaafirstpage}
% The header of the first page of each
% article:
%    \begin{macrocode}
\newcommand{\ps@rmaafirstpage}{%
  \renewcommand\@oddfoot{\hfil\thepage}%
%    \end{macrocode}
% footer corrected (06 Jun 2000) to have page on outside edge;
%    \begin{macrocode}
  \renewcommand\@evenfoot{\thepage\hfil}%
%    \end{macrocode}
% Header has journal name, volume, and page numbers for a regular
% article, or similar for special cases. 
% \changes{4.0}{2019/11/03}{Change to three-line header with copyright info and DOI.}
%    \begin{macrocode}
\if@preprint
\renewcommand\@oddhead{%
  \if@proceedings 
  \textsc{To appear in ``\@conftitle{} (\rmaa@YearStyle{\@year})''}\hfil
  \textit{RMxAC}
  \else
  \textsc{To appear in} \@journalname\hfil
  \fi}\fi%
\if@guide
\renewcommand\@oddhead{\textsc{Documentation} (\rmaa@YearStyle{\@year})\hfil}
\fi%
\if@manuscript
\renewcommand\@oddhead{Manuscript for \@journalname{} (\rmaa@YearStyle{\@year})\hfil}
\fi%
\if@journal
\renewcommand\@oddhead{%
  \parbox{\textwidth}{
    \@journalname, \rmaa@VolumeStyle{\@volume},
    \rmaa@PagesStyle{\@artpages}\
    (\rmaa@YearStyle{\@year})\hfil\\[1pt]
    {\small\textcopyright{} \@year:
      Instituto de Astronom\'\i{}a,
      Universidad Nacional Aut\'onoma de M\'exico}\hfil\\[1pt]
    {\small\@DOI@url}\hfil
  }%
}\fi%
  \renewcommand\@evenhead{\@oddhead}}
%    \end{macrocode}
% By default, use the |\ps@rmaa| page style.
%    \begin{macrocode}
\pagestyle{rmaa}
%    \end{macrocode}
% \end{macro}
% \end{macro}
%
% \subsection{Preamble commands}\label{sec:preamble} 
% Everything you need for setting the title page, TOC and index
% entries, etc.
% \subsubsection{Commands for the title page}
% \label{sec:title-commands}
% \begin{macro}{\resumen}
% \begin{macro}{\@resumen}
% \begin{macro}{\abstract}
% \begin{macro}{\@abstract}
% User commands for setting the resumen and abstract. 
% These should be used in the document header. The internal commands
% |\@resumen| and |\@abstract| are used to hold the actual text
% pending output via the |\maketitle| command. 
%    \begin{macrocode}
\newif\if@setresumen
\newcommand{\resumen}[1]{%
%%  \renewcommand{\@resumen}{{\selectlanguage{spanish}#1}}%
  \renewcommand{\@resumen}{{\ignorespaces#1}}%
  \@setresumentrue}
\renewcommand{\abstract}[1]{%
  \renewcommand{\@abstract}{{\ignorespaces#1}}
  \renewcommand\ADS@abstract{#1}
%% Not a good idea, since then author macros are not expanded and
%% cause errors later. However, by not doing it we are vulnerable to
%% fragile commands.....
  \if@noexpandADS\@onelevel@sanitize{\ADS@abstract}\fi
  \@abstractpresenttrue
  }
\newcommand{\ADSnoexpand}{\@noexpandADStrue}
%    \end{macrocode}
% \begin{macro}{\boldabstract}
% This is the version for doing 1-page conference proceedings. Still
% save the text for ADS but then write it out immediately as a bold
% paragraph. 
%    \begin{macrocode}
\newcommand{\boldabstract}[1]{%
  \gdef\ADS@abstract{#1}%
  \renewcommand{\@abstract}{#1}%
  \if@noexpandADS\@onelevel@sanitize{\ADS@abstract}\fi%
  \noindent\textbf{\boldmath #1}
  }
%    \end{macrocode}
% \end{macro}
% Default value in each case is to print out a warning if not set
%    \begin{macrocode}
  \newcommand{\@resumen}{%
    \@setresumenfalse \ClassWarningNoLine{rmaa}{No Spanish abstract
      (resumen) specified - if you do not speak Spanish, this will be
      translated by the editors}%
    \textbf{Favor de proporcionar un resumen en espa\~nol. If you are
      unable to translate your abstract into Spanish, the editors will
      do it for you.} }
  \newcommand{\@abstract}{\ClassWarningNoLine{rmaa}{No abstract
      specified}}
%    \end{macrocode}
% \end{macro}
% \end{macro}
% \end{macro}
% \end{macro}
%
% \begin{macro}{\title}
% \begin{macro}{\@rawtitle}
% \begin{macro}{\@TOCtitle}
% \begin{macro}{\rmaaIndexTitle}
% Similarly, |\title| is the preamble command to set the title of the
% article. This is used in various places, typeset in different
% styles, hence the use of the internal |\@rawtitle|, |\@title|,
% |\rmaaIndexTitle|, and |\@TOCtitle| commands, each of which uses the
% relevant style from section~\ref{sec:text-styles}
%    \begin{macrocode}
\newcommand{\@rawtitle}{Raw Title}
\newcommand{\@TOCtitle}{\rmaa@TOCTitleStyle{Title for TOC}}
\newcommand{\rmaaIndexTitle}{\rmaa@IndexTitleStyle{Title for index}.}
\newcommand{\@TOCaftertitle}{}
\renewcommand{\title}[1]{%
%%  {\ADS@sanitize\xdef\@rawtitle{#1}}
  \renewcommand\@rawtitle{\noexpand#1}
  \let\ADS@title=\@rawtitle
  \renewcommand{\@title}{\rmaa@TitleStyle{#1}}%
  \renewcommand{\@TOCtitle}{\rmaa@TOCTitleStyle{#1}\@TOCaftertitle}%
  \renewcommand{\rmaaIndexTitle}{\rmaa@IndexTitleStyle{#1\@TOCaftertitle}.}%
  }
%    \end{macrocode}
% \end{macro}
% \end{macro}
% \end{macro}
% \end{macro}
% 
% All the weird and whacky ways of indicating the author
% affiliations. 
% \begin{macro}{\affil}
% First, the way that sets it directly underneath the authors name.
% Note that |\affil| trumps all other affiliations for the purposes of
% getting into the ADS stuff (we use |\gdef|).   
% \dots 
%    \begin{macrocode}
\newcommand{\affil}[1]{%
  \\[1.5ex]\rmaa@AffilStyle{#1}\\[1.5ex]%
  \if@firstaffil\gdef\ADS@affil{#1}\@firstaffilfalse\fi%
  }
\renewcommand{\and}{and\\[1.5ex]}
%    \end{macrocode}
% \begin{macro}{\altaffil}
% Alternative affiliations are footnoted. 
%    \begin{macrocode}
\newcommand{\altaffil}[1]{%
  \if@firstaltaffil\gdef\ADS@affil{#1}\@firstaltaffilfalse\fi%
  \thanks{#1}}
%    \end{macrocode}
% \begin{macro}{\altaffilmark}
% \begin{macro}{\altaffiltext}
% Alternative-alternative mechanism for affiliations---should not be
% mixed with |\altaffil| (unless you know what you are doing) or there
% will be problems with numbering.
%    \begin{macrocode}
\newcommand{\altaffilmark}[1]{\textsuperscript{#1}}
\newcommand{\altaffiltext}[2]{%
  \if@firstaltaffil\gdef\ADS@affil{#2}\@firstaltaffilfalse\fi
  \stepcounter{footnote}%
  \protected@xdef\@thanks{\@thanks
    \protect\footnotetext[#1]{#2}}} 
%    \end{macrocode}
% \end{macro}
% \end{macro}
% \end{macro}
% \end{macro}
%
% \begin{macro}{\keywords}
% \begin{macro}{\@keywords}
% \begin{macro}{\@keywordheader}
% \label{keywords}Sundry minor preamble commands. All
% self-explanatory, I hope. |\keywords| is the old user interface
% (deprecated). New articles should use |\addkeyword| once for each
% keyword. That way the ADS record comes out right too (comma rather
% than em-dash separator).
% \changes{2.2beta12}{2001/12/07}{\cs{addkeyword} New user interface
%   for keywords.} 
% \changes{2.2beta13}{2002/03/11}{\cs{addkeyword} Improved behaviour
%   with .ads file. Use
%   \cs{ADS@sanitize} to change definitions of `~' and \cs{ion}.}
% \changes{3.23}{2007/09/09}{\cs{addkeyword} now signals a compilation
% error if the author tries to add more than 6 keywords.}
% \changes{3.24}{2007/09/14}{kwdcounter is now reset in
% \cs{@begofarticlehook} to prevent error in book style} 
%    \begin{macrocode}
\newcommand{\@kwdsep}{ --- }
\newcommand{\ADS@kwdsep}{, }
\newcounter{kwdcounter}\setcounter{kwdcounter}{0}
\newcommand{\addkeyword}[1]{%
  \addtocounter{kwdcounter}{1}
  \ifthenelse{\thekwdcounter > 6}{
    \ClassError{rmaa}{Too many keywords. A maximum of 6 keywords is allowed.}{}
  }{}
  \if@firstkwd%
  \renewcommand\@keywords{\if@boldkeywords\uppercase{#1}\else #1\fi}
  {\ADS@sanitize\def\ion{\plainion}\gdef\ADS@keywords{#1}}
  \@firstkwdfalse
  \else%
  \edef\@keywords{\@keywords\@kwdsep \if@boldkeywords\uppercase{#1}\else #1\fi}%
  {\ADS@sanitize\def\ion{\plainion}\xdef\ADS@keywords{\ADS@keywords\ADS@kwdsep #1}}%
  \fi}
\newcommand{\ADS@sanitize}{%
  \def\mathrm{}
  \def\@{\ }
  \def~{\ }
  \def\langle{<}
  \def\rangle{>}
  \def\log{log }
  \def\sin{sin }
  \def\cos{cos }
  \def\tan{tan }
  \def\fd{^d.}
  \def\fh{^h.}
  \def\fm{^m.}
  \def\fs{^s.}
  \def\fdg{�.}
  \def\farcm{'.}
  \def\farcs{''.}
  \def\arcdeg{�}
  \def\,{ }
  }
\newcommand{\keywords}[1]{%
  \renewcommand{\ADS@keywords}{#1}
  \renewcommand{\@keywords}{\uppercase{#1}}
  }
\newcommand{\@keywordheader}{\rmaa@KeywordHeaderStyle{Key Words:}}
\newcommand{\@keywords}{%
  \ClassWarningNoLine{rmaa}{No keywords specified}%
  Keywords go here
  }
%    \end{macrocode}
% \begin{macro}{\shorttitle}
% \begin{macro}{\shortauthor}
% \begin{macro}{\@shorttitle}
% \begin{macro}{\@shortauthor}
% Commands for setting the running headers.
%    \begin{macrocode}
\newcommand{\shorttitle}[1]{%
  \renewcommand\@shorttitle{\MakeUpperCase{#1}}}
\newcommand{\shortauthor}[1]{%
  \renewcommand\@shortauthor{\MakeUpperCase{#1}}}
\newcommand{\@shorttitle}{SHORT TITLE GOES HERE}
\newcommand{\@shortauthor}{SHORT AUTHOR GOES HERE}
%    \end{macrocode}
% \end{macro}
% \end{macro}
% \end{macro}
% \end{macro}
% \end{macro}
% \end{macro}
% \end{macro}


% \subsubsection{The full addresses}
% \label{sec:full-addresses}
% \begin{macro}{\fulladdresses}
% \begin{macro}{\@fulladdresses}
% \begin{macro}{\@fulladdheight}
% \begin{macro}{\@fulladddepth}
% \begin{macro}{\@fulladdbox}
% Aaghh, the full addresses! A perennial source of problems. 
%    \begin{macrocode}
\newlength{\@fulladdheight}
\newlength{\@fulladddepth}
\newsavebox{\@fulladdbox}
%% \fulladdresses command modified (11 May 2000) so that we can put
%% \fulladdresses{} in the article to suppress the output of the
%% addresses
\newcommand{\@FinalFigure}{\relax}
\newcommand{\FinalFigure}[1]{\renewcommand{\@FinalFigure}{#1}}
\newcommand{\fulladdresses}[1]{%
  \renewcommand{\@rmaaArg}{#1}
  \ifx\@rmaaArg\@rmaaEmpty
     \renewcommand{\@fulladdresses}{}
  \else
     \renewcommand{\@fulladdresses}{
       \normalsize%
       \ifx\@FinalFigure\relax\else\vfill\@FinalFigure\vfill\fi
%%       \vspace*{1ex}
       %% do addresses as a list (copied from thebibliography environment) 
       \list{\null}{\leftmargin 1.5em\topsep\z@\labelwidth\z@\labelsep\z@\itemindent -1.5em
         \itemsep\z@\parsep\z@\usecounter{enumi}}
       \sloppy\clubpenalty4000\widowpenalty4000
       \sfcode`\.=1000\relax
       #1
       \endlist
       }
  \fi
  }
\newcommand{\@fulladdresses}{Full addresses go here}
%    \end{macrocode}
% \begin{macro}{\suppressfulladdresses}
% New user-level command for suppressing full addresses should you so
% desire. Currently used in the single-page-poster style. 
%    \begin{macrocode}
\newcommand{\suppressfulladdresses}{\@suppressfulladdressestrue\@needaddressesfalse}
%    \end{macrocode}
% \end{macro}
% \end{macro}
% \end{macro}
% \end{macro}
% \end{macro}
% \end{macro}

% \begin{macro}{\printdates}
% \begin{macro}{\noprintdates}
% Command to activate the printing out of the received/accepted
% dates. This can also be done using the \texttt{printdates} option,
% which is the default for \texttt{journal}, but not for \texttt{proceedings}.
%    \begin{macrocode}
\newcommand{\printdates}{\@printdatestrue}
\newcommand{\noprintdates}{\@printdatesfalse}
%    \end{macrocode}
% \end{macro}
% \end{macro}
% \begin{macro}{\@outputaddressesifnecessary}
% Command to output the addresses if we haven't already done so
% This is called by: 
% \begin{enumerate}
% \item The end of |thebibliography| environment
% \item The |\outputfulladdresses| command
% \end{enumerate}
%    \begin{macrocode}
\newcommand{\@outputaddressesifnecessary}{%
  \ClassInfo{rmaa}{Calling @outputaddressesifnecessary}
  \ifthenelse{\boolean{@needaddresses}}{%
    %% If in one column style, or in the first column of a two column
    %% style then we can try to put it on this page 
    \ifthenelse{\boolean{@firstcolumn} \or \not\boolean{@rmtwocol}}{%
      \@outputaddresses}{%
      %% otherwise give warning it will be on a separate page
      \ClassWarning{rmaa}{Full addresses will go on a separate
        page (please see User Guide for how to remedy this)}}%
    %% If we haven't managed yet, put it on its own page
    \ifthenelse{\boolean{@needaddresses}}{%
      \if@debug{\ClassWarning{rmaa}{In this silly bit!}}\fi
      \clearpage\vfill\@outputaddresses}{}%
    }{
    %% If we don't need them, we still might need to do the eoa hook 
    \ifthenelse{\boolean{@suppressfulladdresses}}{%
      \if@debug{\ClassWarning{rmaa}{XXX}}\fi
%%      \if@book\else\@endofarticlehook\null\fi
      }{}
    }%
  }  
%    \end{macrocode}
% \end{macro}

% \begin{macro}{\outputfulladdresses}
% A user interface to the above command, in case there is no biblio
% (11 May 2000) 
%    \begin{macrocode}
\newcommand{\outputfulladdresses}{\@outputaddressesifnecessary}
%    \end{macrocode}
% \end{macro}

% \begin{macro}{\adjustfinalcols}
% This command starts a new column and should be 
% used to ensure that the columns balance on the last page. It also
% outputs the full 
% author addresses to ensure that this occurs in the left hand
% column. 
% The |\par| command is so that the current line will be taken account
% of when working out |\pagetotal|
%    \begin{macrocode}
\newcommand{\adjustfinalcols}{\if@twocolumn\goodbreak\@outputaddresses\newpage\noindent\fi}
%    \end{macrocode}
% \end{macro}

% \begin{macro}{\@outputaddresses}
% Command to do the work of outputting the addresses. 
% This can be called by 
% \begin{enumerate}
% \item an explicit |\adjustfinalcols|
% \item via |\@outputaddressesifnecessary|, q.v.
% \end{enumerate}
%
%    \begin{macrocode}
\newcommand{\@outputaddresses}{%
  % Allow option of suppressing the full addresses, in which case
  % we just run the hook if we are not in a book style
  \if@suppressfulladdresses
%%     \if@book\else\@endofarticlehook\null\fi
  \else
  %% save the full addresses in a box for output later
  \sbox{\@fulladdbox}{\parbox{\textwidth}{\@fulladdresses}}
  % calculate height of address box  
  \settoheight{\@fulladdheight}{\usebox{\@fulladdbox}}
  \settodepth{\@fulladddepth}{\usebox{\@fulladdbox}}
  \addtolength{\@fulladdheight}{\@fulladddepth}
  \if@debug
  { \typeout{Address total height: \the\@fulladdheight}
    \typeout{Text height: \the\pagetotal}
    \typeout{Goal height: \the\pagegoal}
    }
  \fi
  % reuse \@fulladdheight to mean the total height of text+addresses
  \addtolength{\@fulladdheight}{\pagetotal}
%%  \addtolength{\@fulladdheight}{-\baselineskip}
  \if@debug\typeout{Adjusted text + addresses height: \the\@fulladdheight}\fi
  % check it will fit on current page
  \ifdim\@fulladdheight<\pagegoal
     %% must shift it to the left
     %% to compensate for the hanging indentation in the reference list. 
     %% NOTE: this is not foolproof!
     %% NB. This will give overfull hbox warnings, which is a bit ugly.
     \vfill\noindent\if@biblist\hspace*{-\leftmargin}\fi
     \usebox{\@fulladdbox}
     % Set this to show we've done it
     \@needaddressesfalse
     % and assume we are at the end of the article
     % so that we do the EOA hook if not in the book style
%%     \if@book\else\@endofarticlehook\fi
  \else
     % if not, just do nothing. With luck, this function will get 
     % called again on the next page
     \if@debug{\typeout{Holding over full addresses to next page}}\fi
  \fi\fi
  }
%    \end{macrocode}
% \end{macro}

% \subsubsection{Commands for editors}
% \label{sec:editor-commands}
% Some user/editor commands for doing the TOC and index in the book
% option. 
% They are defined here rather in the book-specific part since they
% need to be always available in order to avoid breaking documents. 
% \begin{macro}{\listofauthors}
% \begin{macro}{\@rawauthors}
% \begin{macro}{\rmaaTOCauthors}
% The command |\listofauthors| formats an author list for use in the
% TOC. It should contain the same names as in the |\author| command,
% but without the affiliation garbage and using |\&| instead of
% |and|. 
%    \begin{macrocode}
\newcommand{\@rawauthors}{Raw Authors}
\newcommand{\rmaaTOCauthors}{\noexpand\rmaaTOCAuthorStyle{Author list for TOC}}
\newcommand{\listofauthors}[1]{%
  \renewcommand{\@rawauthors}{#1}
  \renewcommand\rmaaTOCauthors{\noexpand\rmaaTOCAuthorStyle{#1}}}
%    \end{macrocode}
% \begin{macro}{\indexauthor}\label{indexauthor}
% The command |\indexauthor| produces an index entry for each
% author. 
% This command must come \emph{after} the
% |\title| and the |\listofauthors| commands. Also, there should be one 
% |\indexauthor| command for \emph{each} author. We also add the
% author onto the end of the |\ADS@authors| macro and construct the
% |\ADS@authorinitial| by a most cunning method (see below), but only
% for the first author. 
% 
% No, No, NO! This was all wrong. Backed out completely: 
% \begin{quote}
% Completely rewritten (19 Dec 2002) since we now use the index class, so 
% as to be able to alphabetize accented characters. A side effect of this is
% that |\index| no longer expands out the argument after the |@|, so we 
% instead add all the |\index| commands to a list (|\indexauthor@deferlist|), 
% which is executed in the |\@begofarticlehook|. 
% 
% Note the stuff with |\toks0| is copied from Exercise 20.15 of the \TeX{}book---it is
% to make sure the latin1-accented name gets passed into the idx
% file. 
% \end{quote}
% We now go back to the old way of doing things, using an optional
% argument to the \cs{indexauthor} command that has the name without
% accents. 
% \changes{3.08}{2003/01/16}{Remerged changes from 3.06 that got
% dropped from 3.07 after a synchronization disaster. Basically, the
% accentless optional argument to \cs{indexauthor}. Must add the macro
% directories to my unison set-up to stop this sort of thing from happening.}
%    \begin{macrocode}
\newcommand\@isortkey{}
\newcommand\@emptykey{empty}
\newcommand\indexauthor[2][empty]{
  \renewcommand\@isortkey{#1}
  \ifx\@isortkey\@emptykey \renewcommand\@isortkey{#2}\fi
  \index{\@isortkey @{{\rmaaIndexAuthorStyle{#2}}\hskip1em plus1em minus0.4em
      \rmaaIndexTitle\hskip1em plus2em minus0.4em\penalty-100
      \rmaaTOCauthors}}
  \indexauthorADS{#2}
}
%%\def\indexauthor@deferlist{\relax}
%%\long\def\indexauthor#1{
%%  \def\@@name{#1} \toks0=\expandafter{\@@name}
%%  \toks1=\expandafter{\indexauthor@deferlist}
%%  \protected@xdef\indexauthor@deferlist{%
%%    \the\toks1
%%    \noexpand\index{\the\toks0@{\noexpand\rmaaIndexAuthorStyle{\the\toks0}\noexpand\hskip1em plus1em minus0.4em 
%%        \rmaaIndexTitle\noexpand\hskip1em plus2em minus0.4em\noexpand\penalty-100
%%        \rmaaTOCauthors}}
%%  }
%%  \indexauthorADS{#1}
%%  }
%%\xdef\indexauthor@deferlist{}
\newcommand\indexauthorADS[1]{%
  \if@firstauthor\protected@xdef\ADS@authorinitial{\firstletter{#1}}\fi
  \protected@xdef\ADS@authors{\ADS@authors\if@firstauthor\else ; \fi
    #1}
  \@firstauthorfalse
}
%    \end{macrocode}
% Note that the first |#1| sets the key for the index, making sure it
% actually starts with  
% first letter of the surname. Meanwhile, the |@{...}| expression,
% which starts with some  
% control sequences is what is written to the |.ind| file. 
% See the makeindex man page for details. 
% \end{macro}
% \end{macro}
% \end{macro}
% \end{macro}

% \begin{macro}{\KillAllAccents}
% \begin{macro}{\firstletter}
% These are utility functions used in the above. |lett| is a counter
% indicating the current letter in the string. |\firstletter| just
% prints the first letter of its argument.
%    \begin{macrocode}
\newcommand{\KillAllAccents}{\def\'{}\def\`{}\def\~{}\def\.{}\def\i{i}}
\newcommand\firstletter[1]{\@car #1\@nil}
%    \end{macrocode}
% \end{macro}
% \end{macro}
% \begin{macro}{\SetVolume}
% \begin{macro}{\SetMSnumber}
% \begin{macro}{\SetFirstPage}
% \begin{macro}{\SetYear}
% \begin{macro}{\ReceivedDate}
% \begin{macro}{\AcceptedDate}
% These are for the editors to set the volume number, etc. 
%    \begin{macrocode}
\newcommand{\SetVolume}[1]{%
  \renewcommand{\@volume}{#1}
  \edef\ADS@volume{\PadToFourWithDots{#1}}
  \renewcommand{\@volumen}{\rmaa@VolumenStyle{volumen #1}}%
  }
\newcommand{\SetMSnumber}[1]{\renewcommand{\@MSnum}{#1}}
%    \end{macrocode}
% Note that there is no real reason to use |\SetFirstPage| since all
% the page numbering is done automatically in the book style. 
%    \begin{macrocode}
\newcommand{\SetFirstPage}[1]{\setcounter{page}{#1}}
\newcommand{\SetYear}[1]{%
  \renewcommand{\@year}{#1}
  \renewcommand{\ADS@year}{#1}
  \renewcommand{\@anyo}{\rmaa@VolumenStyle{#1}}
  }
\newcommand{\ReceivedDate}[1]{\renewcommand{\@received}{#1}}
\newcommand{\AcceptedDate}[1]{\renewcommand{\@accepted}{#1}}
%    \end{macrocode}
%
% \begin{macro}{\DOI}
% User/editor command to enter the DOI, which will be typeset as a link.
% \changes{4.0}{2019/11/02}{New user command \cs{DOI} to typeset
% DOI in the header of first page and add it to ADS info.}
%    \begin{macrocode}
\newcommand\DOI[1]{%
  \def\@DOI@url{\url{https://doi.org/10.22201/ia.01851101p.#1}}%
  \def\ADS@doi{https://doi.org/10.22201/ia.01851101p.#1}%
}%
%    \end{macrocode}
% \end{macro}
%
% \begin{macro}{\ORCIDiD}
% \changes{4.0}{2019/11/12}{Add mechanism for giving ORCID iD links}
% \begin{macro}{\ORCIDiDfull}
% \begin{macro}{\@orcid@icon}
% \begin{macro}{\@orcid@iconbig}
% Commands to add icon links to ORCID iDs of authors. For example
% \begin{quote}
% |\ORCIDiD|\marg{0000-0000-0000-0000}
% \end{quote}
% to be used after each author name in the |\author|
% command and in the |\fulladdresses|.  Alternatively,
% |\OORCIDiDfull| will typeset the icon plus the full text of the
% link, but we don't use that at present.  The icon image is stored in |ORCIDiD_iconvector.pdf|, which needs to be on the input path.  In principle, we could make EPS versions if necessary.
%    \begin{macrocode}
\newcommand\@orcid@icon{\includegraphics[height=1.8ex]{ORCIDiD_iconvector}}
\newcommand\@orcid@iconbig{\includegraphics[height=2ex]{ORCIDiD_iconvector}}
\newcommand\ORCIDiD[1]{\href{https://orcid.org/#1}{\,\@orcid@icon}}
\newcommand\ORCIDiDfull[1]{%
  \href{https://orcid.org/#1}{\@orcid@iconbig}\,\url{https://orcid.org/#1}}
%    \end{macrocode}
% \end{macro}
% \end{macro}
% \end{macro}
% \end{macro}
%
% \begin{macro}{\@volume}
% \begin{macro}{\@volumen}
% \begin{macro}{\@artpages}
% \begin{macro}{\@year}
% \begin{macro}{\@anyo}
% \begin{macro}{\@received}
% \begin{macro}{\@accepted}
% \begin{macro}{\@MSnum}
% Default values for the containers used by above commands, plus a
% few more useful ones. 
%    \begin{macrocode}
\newcommand{\@volume}{00}
\newcommand{\@volumen}{\rmaa@VolumenStyle{00}}
\newcommand{\@artpages}{%
  \oldpageref{\articlename:RMAAFirstPage}--\oldpageref{\articlename:RMAALastPage}}
\newcommand{\@year}{\the\year}
\newcommand{\@anyo}{\rmaa@VolumenStyle{2007}}
\newcommand{\@journalname}{\rmaa@JournalNameStyle{%
    \if@proceedings\RMAASClong\else\RMAAlong\fi}}
\newcommand{\@received}{}
\newcommand{\@accepted}{}
\newcommand{\@MSnum}{0000}
\newcommand{\@DOI@prefix}{https://doi.org/10.22201/ia.01851101p}
\def\@DOI@url{\url{https://doi.org/10.22201/ia.01851101p.2099.00.00.00}}
%    \end{macrocode}
% \end{macro}
% \end{macro}
% \end{macro}
% \end{macro}
% \end{macro}
% \end{macro}
% \end{macro}
% \end{macro}
% \end{macro}
% \end{macro}
% \end{macro}
% \end{macro}
% \end{macro}
% \end{macro}
% 
% The following commands are most useful in the \texttt{book}
% style. However, the |\SetConfTitle| command can now be used in
% individual articles to set the header for the \texttt{preprint}
% style. 
%    \begin{macrocode}
\newcommand{\@reunionplace}{Lugar de la reuni\'on}
\newcommand{\@reuniondate}{Fecha de la reuni\'on}
\newcommand{\@conftitle}{Conference Title}
\newcommand{\SetReunion}[2]{\renewcommand{\@reunionplace}{#1}\renewcommand{\@reuniondate}{#2}}
\newcommand{\SetConfTitle}[1]{\renewcommand{\@conftitle}{#1}}
\newcommand{\@editors}{Juan Ram\'on de la Fuente and ``El Mosh''}
\newcommand{\SetEditors}[1]{\renewcommand{\@editors}{#1}}
\newcommand{\@editores}{Juan Ram\'on de la Fuente y ``El Mosh''}
\newcommand{\SetEditores}[1]{\renewcommand{\@editores}{#1}}
\newcommand{\@mes}{mes de la publicaci\'on}
\newcommand{\SetMes}[1]{\renewcommand{\@mes}{\rmaa@VolumenStyle{#1}}}
%    \end{macrocode}

% \subsection{Notes for Author Proofs and Communication between Editors}
% \label{sec:notes}
% We introduce a command \cs{Note} that should be used to describe
% any change made to the manuscript. These are a bit tricky to get
% right. In general we use \cs{marginpar} but this only works in outer
% paragraph mode, so we just bung the box where we are if we are in a
% figure caption or something. If we are in the preamble, then we just
% use a \marginpar as part of the \maketitle command.
%    \begin{macrocode}
\newif\if@arenotes
\@arenotesfalse
\newcounter{notecounter}
\setcounter{notecounter}{0}
\newcommand\NotesList{}
\newcommand\Note[1]{%
  \if@proof{%
    \@arenotestrue%
    \stepcounter{notecounter}%
    \if@inpreamble\else{%
      \NoteMark{\arabic{notecounter}}%
    }\fi
    \protected@xdef\NotesList{\NotesList 
      \protect\item #1}%
  }\fi%
}
\newcommand\Changed[2]{Changed ``#1'' to ``#2''.}
\newcommand\NoteMark[1]{%
  \ifinner
  \raisebox{1ex}[0pt][0pt]{\makebox[0pt][l]{\framebox{\footnotesize\bfseries
        N#1}}}%
  \else
  \marginpar{\framebox{\footnotesize\bfseries N#1}}%
  \fi
}
\newcommand\RefNote[1]{\if@proof\NoteMark{\ref{#1}}\relax\fi}
\newcommand\NoteLabel{\protect\label}
\newcommand\OutputNotes{%
  \bgroup
  \onecolumn
  \section*{EDITORIAL NOTES}
  \begin{enumerate}
    \NotesList
  \end{enumerate}
  \egroup
}
\DeclareRobustCommand\cs[1]{\texttt{\char`\\#1}}

%    \end{macrocode}
% \subsection{Bibliographic Records for ADS}
% \label{sec:ADS}
% \changes{2.2beta11}{2001/11/27}{Just about finished all the ADS
% stuff now. Except for the references, which I have given up on for
% the time being.}
% \begin{macro}{\ADS@BibCode}
% \begin{macro}{\ADS@year}
% \begin{macro}{\ADS@journal}
% \begin{macro}{\ADS@volume}
% \begin{macro}{\ADS@firstpage}
% \begin{macro}{\ADS@authorinitial}
% We use the |\ADS@BibCode| command to hold the bibliographic code for
% ADS. It is 
% used in the |%R| 
% field of the bibliographic record for each article. 
% The format is 
% \begin{quote}
% \texttt{YYYYJJJJJVVVVMPPPPA}  
% \end{quote}
% where \texttt{YYYY} is the year, \texttt{JJJJJ} is the journal code,
% which is \texttt{RMxAA} or \texttt{RMxAC}, \texttt{VVVV} is the
% volume number (right-justified), \texttt{M} is not used by us,
% \texttt{PPPP} is the first page number (right-justified), and
% \texttt{A} is the initial of the first author's surname. The volume
% and page fields are padded with periods (`\texttt{.}'), which causes
% us some fun (this is done with the macro |\PadToFourWithDots| described
% below) . 
% 
% All these individual bits are held in the commands |\ADS@year|,
% |\ADS@journal|, |\ADS@volume|, |\ADS@firstpage|,
% |\ADS@authorinitial|. 
%    \begin{macrocode}
\newcommand{\SetADSBibCode}[1]{\renewcommand{\ADS@BibCode}{#1}}
\newcommand{\ADS@BibCode}
{\ADS@year\ADS@journal\ADS@volume.\PadToFourWithDots{\ADS@firstpage}\ADS@authorinitial}
\newcommand{\ADS@year}{2001} 
\newcommand{\ADS@journal}{\if@proceedings RMxAC\else RMxAA\fi}
\newcommand{\ADS@volume}{....} 
%    \end{macrocode}
% These two are just set to default values here, |\ADS@firstpage| is set
% in the |document| environment for included articles, in the
% |frontsection| environment, |\tableofcontents| command,  |index| and
% |abstracts| environments. Also by |\AtBeginDocument| for stand-alone
% articles. 
% 
% God knows how we are going to do |\ADS@authorinitial|!
%    \begin{macrocode}
\newcommand{\ADS@firstpage}{....}
\newcommand{\ADS@authorinitial}{A}
%    \end{macrocode}
% \end{macro}
% \end{macro}
% \end{macro}
% \end{macro}
% \end{macro}
% \end{macro}
% We also have a whole slew of other fields that can be set. These are
% all listed in Table~\ref{tab:ADS}. Most of them still have to be
% defined. 
% \begin{table}[!ht]\raggedleft
% \caption{ADS Bibliographic Records}\label{tab:ADS}
% \begin{tabular}{lp{8em}lp{8em}}\hline
% Code & Description & Macro Name & Comment \\ \hline
% |%R| & Bibliographic Code & |\ADS@BibCode| & required \\
% |%A| & Author list & |\ADS@authors| & required, generate from |\indexauthor| \\
% |%F| & Author affiliation & |\ADS@affil| & format: AA(\dots)
% AB(\dots) \dots \\
% |%J| & Journal Name & |\ADS@journalname| &  \\
% |%D| & Publication Date & |\ADS@fulldate| &  \\
% |%L| & Last Page & |\ADS@lastpage| & generate from label \\
% |%T| & Title & |\ADS@title| & required, generate from |\title| \\
% |%C| & Abstract Copyright & |\ADS@copyright| & set to IA-UNAM \\
% |%U| & URL for electronic document & |\ADS@url| & set to PDF file \\
% |%d| & DOI for article & |\ADS@doi| & set by |\DOI| \\
% |%K| & Keywords & |\ADS@keywords| & set equal to |\@keywords| \\
% |%A| & Abstract Text & |\ADS@abstract| & set equal to |\@abstract| \\
% |%R| & References & |\ADS@references| & generate from the
% bibliography  \\ \hline
% \end{tabular}
% \end{table}
% A lot of these macros are set by other commands, as mentioned in the
% Table, but some need to be set  explicitly here. 
% 
% \begin{macro}{\ADS@doi}
% \changes{4.0}{2019/11/08}{Save DOI for ADS}
% The user/editor command \cs{DOI} is used to set this. 
%    \begin{macrocode}
\newcommand{\ADS@doi}{https://doi.org/10.22201/ia.01851101p.0000.00.00.00}
%    \end{macrocode}
% \end{macro}
% 
% \begin{macro}{\ADS@journalname}
% This
% does not contain any accents at the moment so as not to cause
% problems---this maybe should be changed though. 
%    \begin{macrocode}
\newcommand{\ADS@journalname}{%
  \if@proceedings\@conftitle\space (Eds. \@editors) \fi%
  Revista Mexicana de Astronom\'\i{}a y Astrof\'\i{}sica%
  \if@proceedings\space (Serie de Conferencias)\fi%
  \space Vol. \@volume, pp. \ADS@firstpage--\ADS@lastpage\space (\@year)%
  \space (http://www.astroscu.unam.mx/\~{}rmaa/)
  }
%    \end{macrocode}
% \end{macro}
% \begin{macro}{\ADS@keywords}
% \begin{macro}{\ADS@abstract}
% \begin{macro}{\ADS@title}
% |\ADS@abstract| is straightforward, since we already have the text
% saved as |\@abstract|. We could do the same with |\ADS@keywords| but
% to avoid unwanted formatting we explicitly set it in |\keywords|
% (p.~\pageref{keywords}). |\ADS@title| is just set to |\@rawtitle|,
% which we were already using for another purpose. 
%    \begin{macrocode}
\let\ADS@title=\@rawtitle
\newcommand{\ADS@keywords}{No keywords specified}
\newcommand{\ADS@abstract}{No abstract specified}
%    \end{macrocode}
% \end{macro}\end{macro}\end{macro}
% \begin{macro}{\ADS@authors}
% |\ADS@authors| wants to be in ``Other, A. N.'' format, but with
% semi-colons between authors. The most natural way to do this seems to
% be with |\indexauthor| (p.~\pageref{indexauthor}). We append each
% author to this macro, which means we must explicitly clear it for
% each new article (done in |\article|). .   
%    \begin{macrocode}
\newcommand{\ADS@authors}{}
%    \end{macrocode}
% \end{macro}
% \begin{macro}{\ADS@fulldate}
% \begin{macro}{\ADS@lastpage}
% \begin{macro}{\ADS@copyright}
% \begin{macro}{\ADS@website}
% \begin{macro}{\ADS@url}
% \begin{macro}{\SetWebsite}
% \begin{macro}{\SetMonthNumeric}
% Other sundry data, together with user commands for bits that can be
% set by the Editors. The |\ADS@url| is constructed from the website
% (currently |http://www.astrosmo.unam.mx/rmaa/|,
% the journal descriptor (e.g., |RMxAC..10/|), followed by the
% path to the PDF file (e.g., |PDF/RMxAC..10_HEW1.pdf|). 
% 
%    \begin{macrocode}
\newcommand{\ADS@website}{http://www.astrosmo.unam.mx/rmaa/}
\newcommand{\ADS@url}
{\ADS@website\ADS@journal\ADS@volume\ADS@number/PDF/\ADS@journal\ADS@volume\ADS@number\_\articlename.pdf}
\newcommand{\ADS@number}{}
\newcommand{\ADS@fulldate}{00/\ADS@year} 
\newcommand{\ADS@lastpage}{00} 
\newcommand{\ADS@copyright}{(C) \ADS@year; Instituto de Astronom\'\i{}a,
  Universidad Nacional Aut\'onoma de M\'exico} 
\newcommand{\SetNumber}[1]{\renewcommand{\ADS@number}{-#1}}
\newcommand{\SetWebsite}[1]{\renewcommand{\ADS@website}{#1}}
\newcommand{\SetMonthNumeric}[1]{\renewcommand{\ADS@fulldate}{#1/\ADS@year}}
%    \end{macrocode}
% \end{macro}\end{macro}\end{macro}\end{macro}\end{macro}\end{macro}\end{macro}
% \begin{macro}{\ADS@affil}
% \begin{macro}{\ADS@references}
% These two are somewhat more problematical. Affiliation currently
% only works for the first author (see \dots) pending redesign of the
% user interface. References are now done by wrapping the
% |thebibliography| environment to open another input stream on the
% same file and search through for the |\bibitems|. 
%    \begin{macrocode}
\newcommand{\ADS@references}{}
\newcommand{\ADS@affil}{}
%    \end{macrocode}
% \end{macro}\end{macro}
% 
% Define two macros that expand to the strings
% |\begin{thebibliography}| and |\end{thebibliography}|. Note that the
% idea for this and much else is inspired by the code for the docstrip
% package. Also define |\StartBibX| in case authors use the more
% standard |\begin{thebibliography}{}| syntax. 
%    \begin{macrocode}
\edef\EndBib{\expandafter\@gobble\string\\end\string{thebibliography\string}}
\edef\StartBib{\expandafter\@gobble\string\\begin\string{thebibliography\string}}
\edef\StartBibX{\expandafter\@gobble\string\\begin\string{thebibliography\string}\string{\string}}
%    \end{macrocode}
% Define two new logical variables. |\ifAnotherline| is true if we
% want to read another line from the input stream. It is set to false
% when we encounter the |\end{thebibliography}|. |\ifInBib| is set to
% true when we frind the |\begin{thebibliography}| and tells us to
% start saving the following lines. 
%    \begin{macrocode}
\newif\ifAnotherLine
\newif\ifInBib
%    \end{macrocode}
% \begin{macro}{makeOther}
% We need a macro to change the catcodes of lots of things to other. 
%    \begin{macrocode}
\def\makeOther#1{\catcode`#1=12\relax}
%    \end{macrocode}
% \end{macro}
% \begin{macro}{\RereadBiblio}
% Define a macro to do all the fiddly stuff. This is used at the
% beginning of |thebibliography| environment (Sec~\ref{sec:biblio}).
%    \begin{macrocode}
\newread\tf@biblio
\newcommand{\RereadBiblio}{%
%    \end{macrocode}
% Open the current file again for reading
%    \begin{macrocode}
  \begingroup
  \immediate\openin\tf@biblio=\articlename
%    \end{macrocode}
% Everything in a group to be safe. Initialize logicals. Black magic
% to change
% catcodes and avoid blank lines being interpreted as |\par|. 
%    \begin{macrocode}
  \AnotherLinetrue\InBibfalse
  \makeOther\ \makeOther\\\makeOther\$%
  \makeOther\#\makeOther\^\makeOther\^^K%
  \makeOther\_\makeOther\^^A\makeOther\%%
  \makeOther\~\makeOther\{\makeOther\}\makeOther\&%
  \endlinechar-1\relax
%    \end{macrocode}
% Start looping through the file. Read each line to the macro
% |\Line|. If we have got to EOF or end of biblio then set
% |\AnotherLinefalse|.
%    \begin{macrocode}
  \loop
  \read\tf@biblio to \Line
  \ifx\Line\EndBib\AnotherLinefalse\InBibfalse\fi
  \ifeof\tf@biblio \AnotherLinefalse\fi
%    \end{macrocode}
% If we are in the biblio and the line is not blank then append it to
% the |\ADS@references| macro together with a newline. 
%    \begin{macrocode}
  \ifInBib
  \ifx\Line\empty\else\xdef\ADS@references{\ADS@references\Line^^J}\fi
  \fi
%    \end{macrocode}
% Check if the line we just saw was the start of biblio, then repeat
% the loop if necessary. 
%    \begin{macrocode}
  \ifx\Line\StartBib\InBibtrue\fi
  \ifx\Line\StartBibX\InBibtrue\fi
  \ifAnotherLine\repeat
  \immediate\closein\tf@biblio
  \endgroup
  }
%    \end{macrocode}
% \end{macro}
% 
% \begin{macro}{\PadToFourWithDots}
% Utility function for left-padding a number to 4 characters with
% periods. Used in 
% constructing |\ADS@BibCode|. This is a very simple algorithm and
% \emph{requires} that the argument actually be a number---if it is
% not then dire consequences will result.  
%    \begin{macrocode}
\newcommand{\PadToFourWithDots}[1]{%
  \ifnum #1 > 999 #1%
  \else\ifnum #1 > 99 .#1%
  \else\ifnum #1 > 9 ..#1%
  \else ...#1%
  \fi\fi\fi}
%%  \ifnum \count0 > 999 #1%
%%  \else\ifnum \count0 > 99 .#1%
%%  \else\ifnum \count0 > 9 ..#1%
%%  \else ...#1%
%%  \fi\fi\fi}
%    \end{macrocode}
% \end{macro}
%
% \begin{macro}{\AuthorADS}
% \begin{macro}{\ADS@affiliations}
% New mechanism for setting the affiliations for ADS. The author labels (AA, AB, \dots) are generated using two nested counters.  We use the |\robustify| command from \texttt{etoolbox} so that accented characters in the affiliation addresses do not blow up.  The list of accents may need to be expanded. 
%    \begin{macrocode}
\newcounter{AUTHOR}
\newcounter{author}[AUTHOR]
\stepcounter{AUTHOR}
\newcommand\ADS@affiliations{}
\robustify{\'}\robustify{\i}\robustify{\~}\robustify{\`}
\def\AuthorADS#1#2{%
  \begingroup
  \def\ORCIDiD##1{<ID system="ORCID">##1</ID>}
  \def\Email##1{<EMAIL>##1</EMAIL>}
  \def\Affil##1{##1}
  \stepcounter{author}
  \ifthenelse{\value{author} > 26}{\stepcounter{AUTHOR}\stepcounter{author}}{}
  \ifthenelse{\equal{\ADS@affiliations}{\@rmaaEmpty}}{%
    \xdef\ADS@affiliations{\Alph{AUTHOR}\Alph{author}(#2)}
  }{%
    \xdef\ADS@affiliations{\ADS@affiliations, \Alph{AUTHOR}\Alph{author}(#2)}
  }%
  \endgroup
}
%    \end{macrocode}
% \end{macro}
% \end{macro}
% %
% \begin{macro}{\ADS@dumpdata}
% \changes{4.0}{2019/11/13}{Write all author affiliations for ADS}
% \begin{macro}{\ADS@write}
% \begin{macro}{\ADS@openfile}
% |\ADS@dumpdata| is the macro that actually does the job of writing out all the
% data to an auxiliary file, using |\ADS@write| to write each
% field. The file is set up with |\ADS@openfile|, which currently is
% done in the TOC of a book.  The new multi-author affiliations are used if |\ADS@affiliations| has been set, otherwise it falls back to the single-author version |\ADS@affil|. 
%    \begin{macrocode}
\newcommand{\ADSfield}[2]{\relax}
\newcommand{\ADSvfield}[2]{\relax}
%%\renewcommand{\ADSfield}[12]{\%#1 #2}
\newcommand{\ADS@dumpdata}
{ 
  \ADS@write{\string\begin{ADSarticle}}
  \ADS@WRITE{\string\ADSfield{R}{\ADS@BibCode}}
  \ADS@write{\string\ADSfield{A}{\ADS@authors}}
  \ifthenelse{\equal{\ADS@affiliations}{\@rmaaEmpty}}{%
    {\ADS@sanitize\ADS@write{\string\ADSfield{F}{AA(\ADS@affil)}}}
  }{%
    {\ADS@sanitize\ADS@write{\string\ADSfield{F}{\ADS@affiliations}}}
  }
  \ADS@write{\string\ADSfield{J}{\ADS@journalname}}
  \ADS@write{\string\ADSfield{D}{\ADS@fulldate}}
  \ADS@write{\string\ADSfield{L}{\ADS@lastpage}}
  {\ADS@sanitize\ADS@write{\string\ADSfield{T}{\ADS@title}}}
  \ADS@write{\string\ADSfield{C}{\ADS@copyright}}
  \ADS@write{\string\ADSfield{d}{\ADS@doi}}
  \ADS@write{\string\ADSfield{U}{\ADS@url}}
  \ADS@write{\string\ADSfield{K}{\ADS@keywords}}
  {\ADS@sanitize\ADS@write{\string\ADSfield{B}{\ADS@abstract}}}
  \ADS@write{\string\ADSfield{Z}{\ADS@references}}
  \ADS@write{\string\end{ADSarticle}}
  }
%%\newcommand{\ADS@write}[1]{\if@book\addtocontents{ads}{#1}\fi}
%%\newcommand{\ADS@openfile}{\@starttoc{ads}}
\newcommand{\ADS@write}[1]{\set@display@protect\immediate\write\tf@ads{#1}}
\newcommand{\ADS@WRITE}[1]{\immediate\write\tf@ads{#1}}
\newcommand{\ADS@openfile}{%
  \newwrite\tf@ads%
  \immediate\openout\tf@ads\jobname.ads\relax
  \ADS@WRITE{\ADS@preamble}}
\newcommand{\ADS@closefile}{\ADS@WRITE{\ADS@postamble}}
\def\ADS@preamble{%
  %%  \string\documentclass{article}
  %%  \string\begin{document}
  }
\def\ADS@postamble{%
  %%  \string\end{document}
  \string%\string% Local Variables:^J 
  \string%\string% mode: LaTeX^J
  \string%\string% mode: Font-lock^J
  \string%\string% End:^J
  }
%    \end{macrocode}
% \end{macro}
% \end{macro}
% \end{macro}
% 
% \subsection{Automatic calculation of the page range}
% \label{sec:page-range}
% Another issue that has cost me a lot of grief. Still, it finally
% seems to be sorted now. All hail Jeff Goldberg. 
%
% Cunning trick to get the page numbers from |\pageref|.  Put a label
% on the first and last page of the document using |\AtBeginDocument|
% and |\AtEndDocument|. .  This way won't work for the book option.
% In that case, we do it by defining |\@endofarticlehook| and
% |\@begofarticlehook| as parallels to |\AtEndDocument| and
% |\AtBeginDocument|.  Note that we have to defer adding
% |\@begofarticlehook| to |\AtBeginDocument| to make sure it comes
% after the redefinitions of |\ref| and |\label|.
%    \begin{macrocode}
\if@book\AtEndDocument{\ADS@closefile}\else
 %\AtBeginDocument{\@begofarticlehook}
\AtEndDocument{\@endofarticlehook}
\fi
%    \end{macrocode}
% \begin{macro}{\@begofarticlehook}
% Since we are getting more and more things that need to be done at
% the beginning of every article, frontsection, etc we define a sister
% hook to 
% |\@endofarticlehook|. This hook does the following: 
% \begin{enumerate}
% \item Labels the first page
% \item Sets |\ADS@firstpage|
% \item runs and then clears |\indexauthor@deferlist|
% \item resets the keyword counter
% \end{enumerate}
% It is executed in the following circumstances: 
% \begin{enumerate}
% \item At |\begin{document}| for standalone files
% \item At |\begin{document}| for each included file
% \item At beginning of front sections 
% \item At beginning of TOC
% \item At beginning of Author Index
% \item At beginning of abstracts
% \end{enumerate}
% For some reason, this is not always within the right group, so we
% use |\xdef| just to be sure. 
%    \begin{macrocode}
\newcommand{\@begofarticlehook}{%
  \if@frontmatter\frontlabel{RMAAFirstPage}\else\label{RMAAFirstPage}\fi%
  \xdef\ADS@firstpage{\@arabic{\c@page}}%
  \if@book\else\makeindex\fi % needed so default index is defined in stand-alone articles
%%  \indexauthor@deferlist
%%  \xdef\indexauthor@deferlist{}
  \@inpreamblefalse
  \setcounter{kwdcounter}{0}
  }
%    \end{macrocode}
% \end{macro}
% \begin{macro}{\@endofarticlehook}
% This is a hook to be run at the end of each article
% currently, it just labels the last page and writes out all the ADS
% data to an auxiliary file.  
% It should not produce any printed output. 
% It is modelled on the |\AtEndDocument| command, to which it reduces
% in the case of a single article. 
%    \begin{macrocode}
\newcommand{\@endofarticlehook}{%
  \if@debug{\ClassWarning{rmaa}{***Doing Hook***}}\fi
  \if@debug{\ClassWarning{rmaa}{ADSabstract is: \ADS@abstract}}\fi
  \clearpage\rmaa@lastpagelabel\if@book\ADS@dumpdata\fi
  \ifx\NotesList\@rmaaEmpty\else\OutputNotes\fi
  \if@debug{\ClassWarning{rmaa}{Notes list is: \NotesList}}\fi
}
%    \end{macrocode}
% \begin{macro}{\rmaa@lastpagelabel}
% New version of last page label command copied from Jeff Goldberg's
% lastpage.sty package - we can't just use |\label| here in case the
% output routine for the last page has already been called. Instead,
% we first explicitly clear the page, step back the page counter,
% explicitly write to the .aux file, then restep the page counter
% again.
%
% But this will conflict with the hyperref package, which does labels differently.  Compensate for that by using the five-argument form.
% 
% Also set the |\ADS@LastPage| macro to the appropriate number. 
%    \begin{macrocode}
\newcommand{\rmaa@lastpagelabel}{%
  \addtocounter{page}{-1}%
  \immediate\write\@auxout{\string
    \newlabel{\articlename:RMAALastPage}{{}{\thepage}{}{}{}}}%
  \edef\ADS@lastpage{\thepage}%
  \addtocounter{page}{1}}
%    \end{macrocode}
% \end{macro}
% \end{macro}

% \subsection{Deal with explicit change in number of columns}
% \label{sec:columns}
% Sometimes, when formatting a book, we need to switch between
% one-column articles and two-column articles. For this to work, we
% need to wrap the standard \LaTeX{} commands to set the |\@rmtwocol|
% flag appropriately and also reset the package default values of all
% scalable lengths. .  
%    \begin{macrocode}
\let\latexonecolumn=\onecolumn
\let\latextwocolumn=\twocolumn
\renewcommand{\onecolumn}{\@rmtwocolfalse\@package@default@lengths\latexonecolumn}
\renewcommand{\twocolumn}{\@rmtwocoltrue\@package@default@lengths\latextwocolumn}
%    \end{macrocode}
% 
% \subsection{The actual formatting of the title}
% \label{sec:title-format}

% Note, we have to redefine the |\maketitle| command to use the
% numbered footnotes. This means we needn't use the atbeginend stuff. 
%    \begin{macrocode}
\renewcommand{\maketitle}{\par
  \begingroup
%    \end{macrocode}
% Don't switch to marked footnotes. 
% This is consistent with existing style but I'm not sure I like it
%    \begin{macrocode}
  %%  \renewcommand\thefootnote{\fnsymbol{footnote}}%
%%  \def\@makefnmark{\hbox to\z@{$\m@th^{\@thefnmark}$\hss}}%
  \def\@makefnmark{\textsuperscript{\@thefnmark}\hspace{-0.2em}}%
  \long\def\@makefntext##1{\parindent 1em\noindent
    \hbox to1.8em{\hss$\m@th^{\@thefnmark}$}##1}%
  \if@rmtwocol
  { \if@debug{\ClassWarning{rmaa}{TWO COLUMN}}\fi
    \twocolumn[\@maketitle]%
    }
  \else
  { \if@debug{\ClassWarning{rmaa}{ONE COLUMN}}\fi
    \newpage
    \global\@topnum\z@   % Prevents figures from going at top of page.
    \if@manuscript \@MSmaketitle \else \@maketitle \fi
    }
  \fi
  \if@manuscript\else\thispagestyle{rmaafirstpage}\@thanks\fi
  \ifx\NotesList\@rmaaEmpty\else{%
    \NoteMark{1--\arabic{notecounter}}\vspace*{-2\baselineskip}%
  }\fi%
  \endgroup
%    \end{macrocode}
% Don't reset footnote counter (not sure whether this is good thing
% or not).
%    \begin{macrocode}
  %%  \setcounter{footnote}{0}%
  \let\thanks\relax\let\affil\relax\let\altaffil\relax
  \let\maketitle\relax\let\@maketitle\relax
  \gdef\@thanks{}\gdef\@author{}\gdef\@title{}}
%    \end{macrocode}

% \begin{macro}{\@kwhwidth}
% \begin{macro}{\@kwtwidth}
% Redefine \cs{@maketitle} to include the abstract, resumen and
% keywords. Note that in the manuscripty style we use
% \cs{@MSmaketitle} instead.  
% \changes{3.16}{2003/04/07}{
%  Added extra vertical space between authors and accepted date.}
% \changes{3.26}{2011/08/15}{Add extra label to review articles}
%    \begin{macrocode}
\newlength{\@kwhwidth}\newlength{\@kwtwidth}
\newcommand\reviewname{Review}
\renewcommand{\@maketitle}{%
  \newpage
  \null
  \if@reviewarticle
  \smallskip
  {\large\bfseries\itshape\reviewname}%
  \bigskip
  \else
  \vskip\@before@title@skip%
  \fi
  \begin{center}%
    {\@title \par}%
    \vskip\@before@author@skip%
    \parbox{\rmaa@AuthorWidthFrac\textwidth}{\centering
      \lineskip\@between@author@skip%
      \@author
      \par}
    %% this should be replaced with `Submitted to...' for preprints
    \if@printdates
    \vskip\@before@dates@skip%
    \if@preprint
    \rmaa@ReceivedStyle{Submitted to RMxAA \@received}
    \else
    \rmaa@ReceivedStyle{Received \@received; accepted \@accepted}
    \fi
    \fi
    \if@manuscript
    \rmaa@ReceivedStyle{Submitted to RMxAA \@received}
    \fi
    \if@guide
    \rmaa@ReceivedStyle{Version \rmaa@version, \rmaa@date}
    \fi
    \vskip\@before@abstract@skip%
    \if@abstractpresent
    \parbox{\rmaa@AbstractWidthFrac\textwidth}{%
      \setlength{\parindent}{0.9cm}%
      \centerline{\rmaa@AbstractHeaderStyle{Abstract}}\par%
      \vskip\@after@abstracthead@skip%
      \if@newstyle\noindent\fi\@abstract
      \vskip\@between@abstract@skip%
      \centerline{\rmaa@AbstractHeaderStyle{Resumen}}\par%
      \vskip\@after@abstracthead@skip%
      \if@newstyle\noindent\fi\@resumen
      \if@setresumen\else\@abstract\fi % append english abstract if
                                       % spanish one not set
      \if@guide\else            %do not use keywords for guide
      \vskip\@before@kwd@skip%
      % Format the Key Words with hanging indentation:
      % They are put in a \parbox with width equal to the line width
      % minus the width of the header. 
      { % Widths of keyword
        \settowidth{\@kwhwidth}{\@keywordheader~}   % header and
        \setlength{\@kwtwidth}{\linewidth}          % keyword text
        \addtolength{\@kwtwidth}{-\@kwhwidth}       %
        \noindent\@keywordheader~\parbox[t]{\@kwtwidth}{%
          \rmaa@KeywordStyle{\boldmath\uppercase{\@keywords}}}}%
      \fi
      }
      \fi
  \end{center}%
  \vspace*{\afterkwdskip}
  \par
  }
\newlength{\afterkwdskip}\setlength{\afterkwdskip}{\z@}
%    \end{macrocode}
% \end{macro}
% \end{macro}
%
%
% \begin{macro}{\reviewarticle}
% \begin{macro}{\reviewtoc}
% \begin{macro}{\reviewlof}
%   User command to specify a review article. Must be used in
%   preamble. Main effect is to place the text ``Review'' before the
%   article title.  Also add command for setting a table of contents
%   and a list of figures at the beginning of the article.
%    \begin{macrocode}
\newcommand{\reviewarticle}{\@reviewarticletrue}
\newcommand\reviewtoc{%
  \section*{CONTENTS}%
  \@starttoc{rtc}%
}
\newcommand\reviewlof{%
  \section*{LIST OF FIGURES}%
  \@starttoc{lof}%
}
%    \end{macrocode}
% \end{macro}
% \end{macro}
% \end{macro}
% \begin{macro}{\l@section}
%    \begin{macrocode}
\renewcommand*\l@section[2]{%
  \ifnum \c@tocdepth >\z@
    \addpenalty\@secpenalty
    \addvspace{0.8em \@plus\p@}%
    \setlength\@tempdima{1.5em}%
    \begingroup
      \parindent \z@ \rightskip \@pnumwidth
      \parfillskip -\@pnumwidth
      \leavevmode %\bfseries
      \advance\leftskip\@tempdima
      \hskip -\leftskip
      #1\nobreak\hfil \nobreak\hb@xt@\@pnumwidth{\hss #2}\par
    \endgroup
  \fi}
%    \end{macrocode}
% \end{macro}
%
% \begin{macro}{\@MSmaketitle}
% A special version of |\@maketitle| for the \texttt{manuscript}
% option. This can deal with a multi-page abstract/resumen. Not very
% well, mind. 
%    \begin{macrocode}
\newcommand{\@MSmaketitle}{%
  \if@debug\ClassInfo{rmaa}{Using special MS maketitle command}\fi
  \newpage
  {
    \thispagestyle{rmaafirstpage}
%     \setlength{\parindent}{0pt}
    \null\vskip 2ex
    \begin{center}
      \@title\par
      \@author\par
      \rmaa@ReceivedStyle{%
        \ifx\@received\@rmaaEmpty
          Draft version: \@date
        \else
          Received: \@received
          \ifx\@accepted\@rmaaEmpty
          \else
            ; Accepted: \@accepted
          \fi
        \fi}\par
    \end{center}
    \@thanks
    \centerline{\rmaa@AbstractHeaderStyle{Resumen}}
    \noindent\@resumen\par\vskip 2ex
    \centerline{\rmaa@AbstractHeaderStyle{Abstract}}
    \noindent\@abstract\par\vskip 2ex
    \noindent\@keywordheader~\rmaa@KeywordStyle{\@keywords}\par
%     \@keywordheader~\rmaa@KeywordStyle{\uppercase{\@keywords}}\par
  }%
  }
%    \end{macrocode}
% \end{macro}

% \subsection{The sectioning commands}
% \label{sec:section-commands}

% \begin{macro}{\section}
% \begin{macro}{\subsection}
% \begin{macro}{\subsubsection}
% \begin{macro}{\paragraph}
% \begin{macro}{\subparagraph}
% Redefine all the standard commands to use these lengths. 
%    \begin{macrocode}
\renewcommand\section{\@startsection {section}{1}{\z@}%
                                   {\abovesectionskip}%
                                   {\belowsectionskip}%
                                   %% how do we get upper case here
                                   {\centering\reset@font\normalsize}}
\renewcommand\subsection{\@startsection{subsection}{2}{\z@}%
                                     {\abovesectionskip}%
                                     {\belowsectionskip}%
                                     {\centering\reset@font\normalsize\itshape}}
\renewcommand\subsubsection{\@startsection{subsubsection}{3}{\z@}%
                                     {\abovesectionskip}%
                                     {\belowsectionskip}%
                                     {\centering\reset@font\normalsize\itshape}}
\renewcommand\paragraph{\@startsection{paragraph}{4}{\z@}%
                                    {\aboveparaskip}%
                                    {\belowparaskip}%
                                    {\reset@font\normalsize\bfseries}}
\renewcommand\subparagraph{\@startsection{subparagraph}{5}{\parindent}%
                                      {\aboveparaskip}%
                                      {\belowparaskip}%
                                      {\reset@font\normalsize\bfseries}}
%    \end{macrocode}
% \end{macro}
% \end{macro}
% \end{macro}
% \end{macro}
% \end{macro}

% \begin{macro}{secnumdepth}
% Make sure that we number the subsubsections---this is done by
% default in the |article| class but not the |book| class, which is
% why we must set it explicitly here.
%    \begin{macrocode}
\setcounter{secnumdepth}{3}
%    \end{macrocode}
% \end{macro}
% \changes{2.2beta10}{2001/11/07}{Fixed \cs{secnumdepth} for \cs{book}
% style.}

% \begin{macro}{\thesubsection}
% \begin{macro}{\thesubsubsection}
% For subsection and subsubsection, we want the numbers to be roman
% but the title to be italic. This seems the easiest way to go about it.
%    \begin{macrocode}
\renewcommand{\thesubsection}{{\upshape\thesection.\@arabic\c@subsection}}
\renewcommand{\thesubsubsection}{\thesubsection{\upshape.\@arabic\c@subsubsection}}
%    \end{macrocode}
% \end{macro}
% \end{macro}

% \begin{macro}{\@seccntformat}
% Add ``.'' to the number here so we can remove it from |\thesection|,
% etc. Reduce gap between number and title WJH 27 Jan 1995
%    \begin{macrocode}
\def\@seccntformat#1{\csname the#1\endcsname .\hskip 0.4em}
%    \end{macrocode}
% \end{macro}

% \begin{macro}{\@sect}
% I didn't really want to mess with the |\@sect| command, but it was
% the only way I could manage to force upper case. Other than the
% commented parts, it is identical to the standard version. 
%    \begin{macrocode}
\def\@sect#1#2#3#4#5#6[#7]#8{
  \ifnum #2>\c@secnumdepth
     \let\@svsec\@empty
  \else
     \refstepcounter{#1}%
     \let\@@protect\protect
     \def\protect{\noexpand\protect\noexpand}%
     \edef\@svsec{\@seccntformat{#1}}%
     \let\protect\@@protect
  \fi
%    \end{macrocode}
% This is the bit I changed (27 Jan 1995). Use upper case for
% |\section|, use mixed case for the rest. 
%    \begin{macrocode}
  \ifnum #2=1                   
     \def\sec@upcase##1{\uppercase{##1}}%
     \def\sec@shape{\upshape}%
  \else                         
     \def\sec@upcase##1{\relax{##1}}%
     \def\sec@shape{\itshape}%
  \fi
  \@tempskipa #5\relax
  \ifdim \@tempskipa>\z@
     \begingroup #6\relax
          \@hangfrom{\hskip #3\relax\@svsec}%
                    {\interlinepenalty \@M \sec@upcase{#8}\par}%
     \endgroup
     \csname #1mark\endcsname{#7}%
%    \end{macrocode}
% Only write TOC entries if doing the Guide style. 
%    \begin{macrocode}
     \if@guide 
        \addcontentsline{toc}{#1}{%
          \ifnum #2>\c@secnumdepth \else
             \protect\numberline{\csname the#1\endcsname}%
          \fi
          #7}%
     \fi
%    \end{macrocode}
% Also allow per-article TOC for review papers
%    \begin{macrocode}
     \if@reviewarticle
        \ifnum #2>2 \else
           \addcontentsline{rtc}{#1}{%
             \protect\numberline{\csname the#1\endcsname}%
             \sec@upcase{\sec@shape #7}}%
        \fi
     \fi
  \else
     \def\@svsechd{#6\hskip #3\relax  %% \relax added 2 May 90
       \@svsec #8\csname #1mark\endcsname{#7}%
       \if@guide
          \addcontentsline{toc}{#1}{%
            \ifnum #2>\c@secnumdepth \else
               \protect\numberline{\csname the#1\endcsname}%
            \fi
            #7}%
       \fi
       }%
  \fi
  \@xsect{#5}}
%    \end{macrocode}
% \end{macro}


% \subsection{Appendices}
% \label{sec:appendices}
% Make sure that appendices in individual articles work properly. We
% copy the definition from the article class (the book one doesn't do
% what we want).  
% \begin{environment}{appendices}
%   This uses an environment |appendices| inside of which each
%   |\section| will become a separate appendix.  Note that the
%   appendices must come before the references (otherwise the full
%   addresses mechanism will not work properly). We make the equation
%   numbering be subordinate to the appendix (equations A12, B6,
%   etc). Finally, we have a medskip followed by APPENDICES to
%   introduce them. \changes{3.17}{2002/04/08}{ Added appendices environment
%   that does the right thing.}
%    \begin{macrocode}
\newenvironment{appendices}{\par
  \begingroup
  \setcounter{section}{0}%
  \setcounter{subsection}{0}%
  \renewcommand\thesection{\@Alph\c@section}
  \renewcommand\theequation{\thesection\arabic{equation}}
  \medskip
  \section*{APPENDICES}}
{\endgroup}
%    \end{macrocode}
% \end{environment}
% \begin{environment}{appendix}
%   If there is only one appendix, then use |appendix| environment
%   instead.  \changes{4.0}{2019/11/12}{ Added
%   appendix environment for case where there is only one section.}
%    \begin{macrocode}
\renewenvironment{appendix}{\par
  \begingroup
  \setcounter{section}{0}%
  \setcounter{subsection}{0}%
  \renewcommand\thesection{\@Alph\c@section}
  \renewcommand\theequation{\thesection\arabic{equation}}
  \medskip
  \section*{APPENDIX}}
{\endgroup}
%    \end{macrocode}
% \end{environment}

% \subsection{Floating thingies}
% \label{sec:floats}   

%% 
%    \begin{macrocode}
\renewcommand{\thetable}{\@arabic\c@table}
\renewcommand{\thefigure}{\@arabic\c@figure}
\renewcommand{\figurename}{Fig.}
\renewcommand{\tablename}{TABLE}
%\renewcommand{\fnum@table}{\MakeUpperCase{\tablename}~\thetable}
%    \end{macrocode}
% For some reason, I have to write `TABLE' explicitly here
% If I use |\tablename|, it gets converted to mixed case---God knows
% why. 
% |\figurename| is also being ignored. 
%    \begin{macrocode}
\renewcommand{\fnum@table}{TABLE~\thetable}
\renewcommand{\fnum@figure}{Fig.~\thefigure}
%    \end{macrocode}

% Do captions in 9 pt face in the main journal style and in new-style
% anything. 
%    \begin{macrocode}
\long\def\@makefigcap#1#2{%
  \if@proceedings\else\small\fi
  \if@newstyle\small\fi 
  \vskip\abovecaptionskip
  \sbox\@tempboxa{#1. #2}%
  \ifdim \wd\@tempboxa >\hsize
    #1. #2 \par
  \else
    \global \@minipagefalse
    \hb@xt@\hsize{\hfil\box\@tempboxa\hfil}%
  \fi
  \vskip\belowcaptionskip}
%    \end{macrocode}

% Make sure table caption is upper case. 
%    \begin{macrocode}
\long\def\@maketabcap#1#2{%
  \hbox to\hsize{\hfil #1 \hfil}
  \vskip 6pt 
  \centering\MakeUpperCase{#2}\par
  \vskip 6pt}
%    \end{macrocode}

%    \begin{macrocode}
\newlength{\@rottabcapwidth}
\long\def\@makerottabcap#1#2{%
  \setbox\@tempboxa\hbox{\MakeUpperCase{#2}}%
  \ifdim\wd\@tempboxa >.8\vsize
     \setlength\@rottabcapwidth{.8\vsize}
  \else
     \setlength\@rottabcapwidth{\wd\@tempboxa}
  \fi
  \rotatebox{90}{%
    \begin{minipage}{\vsize}\centering
      \begin{minipage}{\@rottabcapwidth}\centering
        #1\\ \vskip 6pt \MakeUpperCase{#2}
      \end{minipage}%
    \end{minipage}
    }%
  \hskip 6pt
  }

\newenvironment{Sideways}
    {\begin{sideways}\begin{minipage}{\vsize}\centering}
    {\end{minipage}\end{sideways}}

\newenvironment{Sidewaystable}
{%
  \let\oldcaption=\caption%
  \def\caption{
    \def\@captype{table}\renewcommand{\@makecaption}{\@maketabcap}
    \oldcaption}
  \begin{sidewaystable}\centering}
  {\end{sidewaystable}\let\caption=\oldcaption}
\newenvironment{Sidewaystable*}
{%
  \let\oldcaption=\caption%
  \def\caption{
    \def\@captype{table}\renewcommand{\@makecaption}{\@maketabcap}
    \oldcaption}
  \begin{sidewaystable*}\centering}
  {\end{sidewaystable*}\let\caption=\oldcaption}
%    \end{macrocode}

% Not sure why we need to stick these lines in. 
%    \begin{macrocode}
\renewenvironment{figure}
                 {\renewcommand{\@makecaption}{\@makefigcap}\@float{figure}}
                 {\end@float}
\renewenvironment{figure*}
                 {\renewcommand{\@makecaption}{\@makefigcap}\@dblfloat{figure}}
                 {\end@dblfloat}
\renewenvironment{table}
                 {\renewcommand{\@makecaption}{\@maketabcap}\@float{table}}
                 {\end@float}
\renewenvironment{table*}
                 {\renewcommand{\@makecaption}{\@maketabcap}\@dblfloat{table}}
                 {\end@dblfloat}
%    \end{macrocode}

% Redefine |\@caption| so that we don't bother with the list of
% figures. This saves us having to worry about fragile commands in the
% captions. 
%    \begin{macrocode}
\long\def\@caption#1[#2]#3{%
  \par
%%  \addcontentsline{\csname ext@#1\endcsname}{#1}%
%%    {\protect\numberline{\csname the#1\endcsname}{\ignorespaces #2}}%
  \begingroup
    \@parboxrestore
    \if@minipage
      \@setminipage
    \fi
    \normalsize
    \@makecaption{\csname fnum@#1\endcsname}{\ignorespaces #3}\par
  \endgroup}
%    \end{macrocode}

% Make tables not quite so cramped. Allow |\hline\hline| to give a
% poor man's thick rule. 
%    \begin{macrocode}
\renewcommand{\arraystretch}{1.15}
\setlength{\doublerulesep}{0.4pt}
%    \end{macrocode}

% \begin{macro}{\tablecols}
% \begin{macro}{\@tablecols}
% \begin{macro}{\tabcaptioncontinued}
% |\tablecols| is a user command for setting the number of columns. It
% is used by commands that try and span the entire table (I think
% there must be a better way of acheiving
% this). |\tabcaptioncontinued| is 
% a user command for setting the caption of a continued table on the
% next page. 
%    \begin{macrocode}
\newcommand{\tablecols}[1]{\renewcommand{\@tablecols}{#1}}
\newcommand{\@tablecols}{1}
\newcommand{\tabcaptioncontinued}{%
  \multicolumn{\@tablecols}{c}{\addtocounter{table}{-1}TABLE~\thetable\ (CONTINUED)} \\ \\
  }
%    \end{macrocode}
% \end{macro}
% \end{macro}
% \end{macro}

% \subsubsection{Footnotes in tables}
% \label{sec:table-footnotes}
% \begin{macro}{\tablenotemark}
% \begin{macro}{\tablenote}
% \begin{macro}{\tablenoteskip}
% Old user commands |\tablenotemark{MARK}|,
% |\tablenote{MARK}{TEXT}|---now deprecated. 
%    \begin{macrocode}
\newcommand\tablenotemark[1]{$^\mathrm{#1}$} 
\newcommand\tablenote[2]{\if@proceedings\else\small\fi ${}^\mathrm{#1}$ #2} 
\newcommand\tablenoteskip{\\[0.5\baselineskip]} 
%    \end{macrocode}
% \end{macro}
% \end{macro}
% \end{macro}

% \begin{macro}{\tabnotemark}
% \begin{macro}{\tabnotetext}
% \begin{macro}{\tabnotetextTWO}
% Alternative way of doing footnotes in tables. 
%    \begin{macrocode}
\newcommand{\tabnotemark}[1]{\textsuperscript{\lowercase{#1}}}
%\newcommand{\tabnotetextTWO}[4]{\tabnotetext{#1}{#2}\tabnotetext{#3}{#4}}
\newcommand{\tabnotetextTWO}[4]{%
  \multicolumn{\@tablecols}{l}{%
    \parbox{0.5\tabnotewidth}{\small\textsuperscript{#1}#2}%
    \parbox{0.5\tabnotewidth}{\small\textsuperscript{#3}#4}%
    }%
  \\}
\newcommand{\tabnotetext}[2]{%
  \multicolumn{\@tablecols}{@{}p{\tabnotewidth}@{}}{\small\textsuperscript{#1}#2}\\}
\newlength{\tabnotewidth}
%    \end{macrocode}
% \end{macro}
% \end{macro}
% \end{macro}


% \begin{macro}{changemargin}
%    \begin{macrocode}
% Command to allow margins to be changed on the fly. Useful to allow
% tables and Figures to protrude into the margins. 
%
% \changes{3.24}{2007/09/14}{Added \texttt{changemargin} environment
% (copied from TeX FAQ).}
\newenvironment{changemargin}[2]{% 
  % From TeX FAQ http://www.tex.ac.uk/cgi-bin/texfaq2html?label=chngmargonfly
  \begin{list}{}{%
    \setlength{\topsep}{0pt}%
    \setlength{\leftmargin}{#1}%
    \setlength{\rightmargin}{#2}%
    \setlength{\listparindent}{\parindent}%
    \setlength{\itemindent}{\parindent}%
    \setlength{\parsep}{\parskip}%
  }%
  \item[]}{\end{list}}
%    \end{macrocode}
% \end{macro}

% \subsection{The bibliography and references}
% \label{sec:biblio}
% \begin{macro}{\beforebiblioskip}
% \begin{macro}{\afterbiblioskip}
% I'm not entirely sure why I introduced these two commands, since
% directly putting |\vspace| before and after
% |\begin{thebibliography}| seems to work well enough. 
%    \begin{macrocode}
\newlength{\beforebiblioskip}
\setlength{\beforebiblioskip}{0pt}
\newlength{\afterbiblioskip}
\setlength{\afterbiblioskip}{0pt}
%    \end{macrocode}
% \end{macro}
% \end{macro}
% \begin{macro}{thebibliography}
% The environment itself is done as starred section with extra
% vertical space after heading. This space is currently zero. 
% In the book style we do the cunning trick with rereading the file in
% order to get the references into the |\ADS@references| macro.
% 
% \changes{3.05}{2003/01/02}{Removed the \cs{beforebiblioskip} and
% \cs{afterbiblioskip} commands, since they cause unwanted vertical
% space in some instances.}  
%    \begin{macrocode}
\renewenvironment{thebibliography}{%
  \if@book\RereadBiblio\fi
%%  \vspace*{\beforebiblioskip}
  \section*{REFERENCES} \frenchspacing
%%  \vspace*{\afterbiblioskip}
  \small
%    \end{macrocode} 
% The references themselves are set as a list in order 
% to get the hanging indentation. They are done in a smaller font and
% with french spacing so we needn't worry about the distinction
% between sentence-ending periods and abbreviations. 
%    \begin{macrocode} 
  \list{\null}{\leftmargin 1.5em\labelwidth\z@\labelsep\z@\itemindent -1.5em
    \itemsep\z@\parsep\z@\usecounter{enumi}}
  \sloppy\clubpenalty4000\widowpenalty4000
  \sfcode`\.=1000\relax\@biblisttrue}
{\endlist\@biblistfalse\normalsize\@outputaddressesifnecessary}
%    \end{macrocode}
% \end{macro}
% \subsection{Uniquification of cross-references}
% \label{sec:uniquification}
% 
% \begin{macro}{\articlename}
% In order that the |\label{}| command in individual articles be
% guaranteed to give a unique anchor, we append the article name (held
% in the |\articlename| command) to each key. This requires
% redefinition of the following commands: |\cite|, |\nocite|, |\ref|,
% |\pageref|, |\label|. 
% \begin{macro}{\oldcite}
% \begin{macro}{\oldnocite}
% \begin{macro}{\oldref}
% \begin{macro}{\oldpageref}
% \begin{macro}{\oldlabel}
% The original versions of these commands are
% saved as |\oldcite|, |\oldnocite|, etc. 
%
% The uniquifying of cross-references is now always done, so 
% we don't get multiply-defined labels in the single article case
%    \begin{macrocode}
\newcommand{\articlename}{anon}
\AtBeginDocument{
  \let\oldcite=\cite 
  \let\oldnocite=\nocite 
  \let\oldref=\ref 
  \let\oldpageref=\pageref
  \let\oldlabel=\label
}
%    \end{macrocode}
% \end{macro}
% \end{macro}
% \end{macro}
% \end{macro}
% \end{macro}
% \end{macro}
%
% Uniquify all cross-references by prepending the article name.  This
% now has to go into |\AtBeginDocument| so it is working on the new
% versions that hyperref.sty has set up.  Then we need to put the
% |\@begofarticlehook| afterwards, so that the first page label is using
% the new definitions.  If amsmath is loaded, then we also need to patch
% the special label command that they use. \changes{4.0}{2019/11/18}{Fix labels for amsmath display equations, in case that package has been loaded.}
%    \begin{macrocode}
\AtBeginDocument{
  \renewcommand{\ref}[1]{\oldref{\articlename:#1}%
    \if@debug{\ClassInfo{rmaa}{Reference to \articlename:#1}}\fi}
  \renewcommand{\pageref}[1]{\oldpageref{\articlename:#1}%
    \if@debug{\ClassInfo{rmaa}{Page reference to \articlename:#1}}\fi}
  \renewcommand{\label}[1]{\oldlabel{\articlename:#1}%
    \if@debug{\ClassInfo{rmaa}{Label for \articlename:#1}}\fi}%
  \ifdef{\label@in@display}{%
    \let\oldlabel@in@display=\label@in@display
    \renewcommand{\label@in@display}[1]{\oldlabel@in@display{\articlename:#1}%
      \if@debug{\ClassInfo{rmaa}{AMSTeX Label for \articlename:#1}}\fi}%
  }{}
  \@begofarticlehook
}%
%    \end{macrocode}
% Citations are more complicated since they can contain a list of
% keys. These are really only here for use with the rmaacite package,
% which is utterly deprecated (02 Jan 2003) now that I have the natbib
% code working. The command \cs{cite} should \emph{not} be used with
% natbib: use \cs{citep}, \cs{citet}, etc.\@ instead. 
%    \begin{macrocode}
\renewcommand{\nocite}[1]{\@bsphack
  \@for\@citeb:=#1\do{%
    \edef\@citeb{\expandafter\@firstofone\@citeb:\articlename}%
    \if@filesw\immediate\write\@auxout{\string\citation{\@citeb}}\fi
    \@ifundefined{b@\@citeb:\articlename}{\G@refundefinedtrue
      \@latex@warning{Citation `\@citeb' undefined}}{}}%
  \@esphack}
\renewcommand{\cite}[1]{\oldcite{#1:\articlename}}
%   \end{macrocode}


% Horrible fiddling with the bibliography citation mechanism. 
% This has to be globally available in case rmaacite puts it in the
% \texttt{.aux} file. 
%    \begin{macrocode}
\def\lbibcite#1#2#3{\global\@namedef{b@#1}{#2}
                    \global\@namedef{by@#1}{#3}}
%    \end{macrocode}

% \begin{macro}{\oldbibitem}
% \begin{macro}{\@rmaaBibArg}
% \begin{macro}{\@rmaaSpace}
% \begin{macro}{rmaabib}
% Of course, |\bibitem| needs to get the uniquification treatment as
% well. While we are at it, we redefine it to use consecutive integers
% for the key if the user leaves it blank (using the counter
% |rmaabib|). This allows one to put 
% |\bibitem{}| for all your references without getting annoying
% messages about multiply-defined labels.
% 
% We also take the opportunity to insert the optional argument to
% \cs{bibitem} if it is missing in the form Author(DATE) so that natbib
% doesn't try to do numerical citations. Natbib already provides a
% hook \cs{@extra@b@citeb} to help uniquify the citation key, so we
% just plug into that. However, it puts it \emph{after} the original
% key, not before as I had had it. Easiest solution is to change the
% way that I do it. 02 Jan 2003. Note that the \cs{@extra@b@citeb}
% does not get written to the .aux  file in the \cs{citation} command
% but it does work. 
% \changes{3.14}{06 Feb 2003}{Reinstalled old definition of \cs{bibitem}
%  for the case where the nonatbib option is requested, so that the
%  rmaacite package will work. This means that you can't use rmaacite
%  and natbib in the same volume.}
%    \begin{macrocode}
\if@natbib\def\@extra@b@citeb{:\articlename}\fi
%%\else
\let\oldbibitem\bibitem
\newcommand{\@rmaaBibArg}{}
\newcommand{\@rmaaSpace}{ }
\newcounter{rmaabib}
\if@natbib
\renewcommand{\bibitem}[2][ABC(2003)]{%
  \renewcommand{\@rmaaBibArg}{#2}
  \ifx\@rmaaBibArg\@rmaaEmpty
     \addtocounter{rmaabib}{1}
     \oldbibitem[#1]{\thermaabib:\articlename}
  \else\ifx\@rmaaBibArg\@rmaaSpace
     \addtocounter{rmaabib}{1}
     \oldbibitem[#1]{\thermaabib:\articlename}
  \else 
     \oldbibitem[#1]{#2:\articlename}
  \fi\fi
  \renewcommand{\@rmaaBibArg}{}
  }
\else
\renewcommand{\bibitem}[1]{%
  \renewcommand{\@rmaaBibArg}{#1}
  \ifx\@rmaaBibArg\@rmaaEmpty
     \addtocounter{rmaabib}{1}
     \oldbibitem{\articlename:\thermaabib}
  \else\ifx\@rmaaBibArg\@rmaaSpace
     \addtocounter{rmaabib}{1}
     \oldbibitem{\articlename:\thermaabib}
  \else
     \oldbibitem{\articlename:#1}
  \fi\fi
  \renewcommand{\@rmaaBibArg}{}
  }
\fi
%%\fi
%    \end{macrocode}
% \end{macro}
% \end{macro}
% \end{macro}
% \end{macro}

% \begin{macro}{oldpcite}
% \begin{macro}{oldscite}
% Then, we need to repeat the above scheme, but for the |rmaacite|
% definition. 
%    \begin{macrocode}
\newcommand{\@redefinebibitem}{%
  \let\oldbibitem\bibitem
  \let\oldcite\cite
  \let\oldpcite\pcite
  \let\oldscite\scite
  \def\bibitem[##1<##2>]##3{%
    \renewcommand{\@rmaaBibArg}{##3}
    \ifx\@rmaaBibArg\@rmaaEmpty
       \addtocounter{rmaabib}{1}
       \oldbibitem[##1<##2>]{\thermaabib:\articlename}
    \else\ifx\@rmaaBibArg\@rmaaSpace
       \addtocounter{rmaabib}{1}
       \oldbibitem[##1<##2>]{\thermaabib:\articlename}
    \else 
       \oldbibitem[##1<##2>]{##3:\articlename}
    % trust the author here
%%       \oldbibitem[##1<##2>]{##3}
    \fi\fi
    }
  }
%    \end{macrocode}
% \end{macro}
% \end{macro}


% \subsection{HTML Table of Contents}
% \label{sec:hyper-TOC}
% \begin{macro}{\HyperTOCEntry}
% \begin{macro}{\HyperTOCpart}
% \begin{macro}{\HyperTOCmeeting}
% \begin{macro}{\HyperTOCeditors}
% \begin{macro}{\HyperTOCpreamble}
% These commands are used to write information to the \texttt{.htoc}
% file, which is used in generating a Web Page of the Table of
% Contents. Since they only ever appear preceded by |\protect| in the
% argument to an |\addtocontents| command,  they should never need to
% be expanded by the macros. However, we define them here just in case
% \dots
%    \begin{macrocode}
\newcommand{\HyperTOCEntry}[3]{\relax}
\newcommand{\HyperTOCpart}[1]{\relax}
\newcommand{\HyperTOCmeeting}[3]{\relax}
\newcommand{\HyperTOCeditors}[1]{\relax}
\newcommand{\HyperTOCpreamble}{\relax}
%    \end{macrocode}
% \end{macro}
% \end{macro}
% \end{macro}
% \end{macro}
% \end{macro}
%

% \subsection{Formatting the poster abstracts} 
% These can now be used without using the book option.
%    \begin{macrocode}

\newcommand{\@makeminititle}{%
  \newpage
  \null 
  \vskip 3em%
  \begin{center}%
    {\@title \par}%
  \end{center}
  \vskip 3em%
  }

\newif\if@abstractsads
\@abstractsadsfalse
\newcommand\abstractsADS{\@abstractsadstrue}

\if@book\RequirePackage{multicol}\fi
\newenvironment{abstracts}[1][Abstracts of Contributed Papers]{%
  %% First, make sure we start on an odd page (06 Jun 2000)
  \if@book\clearpage\thispagestyle{empty}\cleardoublepage\fi
  %%
  \renewcommand{\articlename}{ABSTRACTS}
  \title{#1}\shortauthor{Abstracts}\shorttitle{Abstracts}
  \listofauthors{}
  \if@book
  \addcontentsline{toc}{chapter}{\protect\@fiddleauthors{\@TOCtitle}{\rmaaTOCauthors}}%
  \addtocontents{htc}{\protect\HyperTOCEntry{\@rawtitle}{\@rawauthors}{\articlename}}
  \fi
  \pagestyle{rmaa} % this now seems necessary (06 Jun 2000)
  \thispagestyle{rmaafirstpage}
%%  \twocolumn[\@makeminititle]
  \renewcommand{\abstract}[1]{%
    \renewcommand\ADS@abstract{##1}%
    \renewcommand{\@abstract}{##1}\@writeabstract%
  }
  \renewcommand{\title}[2][]{%
    \gdef\ADS@authors{}%
    \gdef\ADS@affil{}%
    \gdef\ADS@references{}%
    \gdef\ADS@keywords{}%
    \renewcommand{\@rmaaArg}{##1}
    \ifx\@rmaaArg\@rmaaEmpty
       \renewcommand{\@TOCaftertitle}{}
    \else
       \renewcommand{\@TOCaftertitle}{ (\@rmaaArg)}
    \fi
    \renewcommand{\@title}{\rmaa@TOCTitleStyle{##2}}%
    \renewcommand{\@TOCtitle}{\rmaa@TOCTitleStyle{##2}\@TOCaftertitle}%
    \renewcommand{\rmaaIndexTitle}{\rmaa@IndexTitleStyle{##2\@TOCaftertitle.}}%
    \renewcommand\ADS@title{##2}%
    \xdef\ADS@firstpage{\@arabic{\c@page}}%
    }
  \renewcommand{\author}[1]{\let\@thanks\@empty\renewcommand\@author{##1} }
  \renewcommand{\footnotetext}[2][XX]{\par\noindent\textsuperscript{##1}~##2}
  \renewcommand{\footnoterule}{\kern-3\p@ \hrule \@width 3em \kern 2.6\p@}
  \if@book\@begofarticlehook\fi
  \if@book\begin{multicols}{2}[\@makeminititle]\else\@makeminititle\fi
  \begingroup
  \setcounter{footnote}{0}
  }% and at the end of the environment....
  {\endgroup\if@book\end{multicols}\listofabstracts\label{RMAALastPage}\clearpage\fi}
%% auxiliary commands used:
%% \@writeabstract does all the actual formatting 
\newcommand{\@writeabstract}{%
  \if@book%
  \addcontentsline{loa}{abstract}{\protect\@fiddleauthors{\@TOCtitle}{\rmaaTOCauthors}}%
  \fi%
  \begin{center}\@title\par\@author\end{center}
  \@abstract
  \par\noindent\footnoterule\nobreak\footnotesize\@thanks
  \par\noindent\hrulefill\goodbreak
  \vskip 1cm \@plus0.5cm
  %% write ADS entry
  \if@book\if@abstractsads\ADS@dumpdata\fi\fi%
  \endgroup % this pair is to make each abstract a self-contained group
  \begingroup
  \setcounter{footnote}{0}
  }

\if@book
%% List of abstracts (stuff written to .loa file)
\newcommand{\listofabstracts}{%
  \begin{multicols}{2}[%
    \vskip 1.5cm \vfill % leave at least 1.5 cm and go to bottom of page if possible
    \centerline{\MakeUpperCase{List of Abstracts}}
    \vskip 1.5cm \@plus0.5cm
    ][6cm] % this makes it start a new page if there is not at least 6cm of space
    \renewcommand{\@fiddleauthors}[2]{##1\hskip1em plus2em minus0.4em\penalty-100  ##2}
    \parindent\z@
%%    \parskip.5\p@ \@plus .7\p@\relax 
    \parskip3\p@ \@plus 1\p@ \@minus 1\p@\relax % allow quite a lot of
                                % leeway here to avoid breaks inside
                                % entries 
    \renewcommand{\\}{\relax} % temporarily turn off explicit line breaks just in case
    \renewcommand{\altaffilmark}[1]{\relax} % ditto with this
    \@starttoc{loa}
  \end{multicols}
  }
\newcommand*{\l@abstract}[2]{%
  \par\penalty-1000
  #1\hskip 1em\penalty0\null\nobreak\hfill#2}
\fi

%    \end{macrocode}
 
% \subsection{Code for typesetting an entire volume of conference
% proceedings} \label{sec:whole-volume}
%    \begin{macrocode}

\if@book
%% environment for including the individual articles
\newcommand{\@fiddleauthors}[2]{\relax}
\renewcommand{\theequation}{\@arabic\c@equation}
\newenvironment{articles}%
{ \begingroup % make it a group so that author definitions do not
              % spill over
  \gdef\ADS@authors{}
  \pagestyle{rmaa}
  \renewenvironment{document}{%
    \@begofarticlehook
    \addtocounter{chapter}{1}%
    \addcontentsline{toc}{chapter}{%
      \protect\@fiddleauthors{\@TOCtitle}{\rmaaTOCauthors}}%
    \addtocontents{htc}{\protect\HyperTOCEntry{\@rawtitle}{\@rawauthors}{\articlename}}
    }{\relax}
  \renewcommand{\documentclass}[2][foo]{\relax} 
  \renewcommand{\usepackage}[2][foo]{%
    \def\@packagename{##2}\def\@rmaacitename{rmaacite}%
    \ifx\@packagename\@rmaacitename
    \@rmcitetrue\ClassInfo{rmaa}{This paper uses rmaacite}
    \makeatletter\input{rmaacite.sty}\makeatother % read it in here so that all changes
    \else                                         % are local to paper
    \if@debug\ClassWarning{rmaa}{Package \@packagename\space was requested by \articlename\space}\fi
    \fi} 
  %% Don't allow individual articles to mess with the pages, volume or year
  \renewcommand{\SetFirstPage}[1]{\relax}
  \renewcommand{\SetYear}[1]{\relax}
  \renewcommand{\SetVolume}[1]{\relax}
  \renewcommand{\SetConfTitle}[1]{\relax}
  %% Don't allow individual articles to change these
  \renewcommand{\nonstopmode}{\relax}
  \renewcommand{\scrollmode}{\relax}
  \renewcommand{\outputplates}{\relax}
  %% XXXX need to check if using rmaacite and act accordingly
  %% avoid multiple declarations of the same length
  \let\oldnewlength=\newlength
  \renewcommand{\newlength}[1]{%
    \edef\reserved@a{\expandafter\@gobble\string ##1}%
    \@ifundefined{\reserved@a}{%
      \oldnewlength{##1}}{%
      \if@debug{\ClassInfo{rmaa}{\reserved@a\space already defined: doing nothing}}\fi%
      \relax}}%
  \renewcommand\newcommand{\providecommand}
  }%
{\endgroup\onecolumn}
     

%    \end{macrocode}
% Command to be used in `articles' environment for including an
% article
%    \begin{macrocode}
\newcommand{\article}[2][]{%
  % optional first argument is extra text to put in brackets after title in TOC
  % e.g. ``Invited Talk'' or ``Poster''
  \renewcommand{\articlename}{#2}
  % reset loads of counters
  \setcounter{footnote}{0}
  \setcounter{figure}{0}
  \setcounter{table}{0}
  \setcounter{section}{0}
  \setcounter{equation}{0}
  \setcounter{rmaabib}{0}
  \setcounter{author}{0}
  \setcounter{AUTHOR}{1}
  \gdef\ADS@authors{}%
  \gdef\ADS@affil{}%
  \gdef\ADS@affiliations{}%
  \gdef\ADS@references{}%
  \@firstauthortrue
  \@firstaffiltrue
  \@firstaltaffiltrue
  \@firstkwdtrue
  \@rmcitefalse
  \if@suppressfulladdresses\else\@needaddressestrue\fi
  \renewcommand{\@rmaaArg}{#1}
  \ifx\@rmaaArg\@rmaaEmpty
     \renewcommand{\@TOCaftertitle}{}
  \else
     \renewcommand{\@TOCaftertitle}{ (#1)}
  \fi
%    \end{macrocode}
% In the main journal articles must start on an odd page. 
%    \begin{macrocode}
  \if@journal\clearpage\thispagestyle{empty}\cleardoublepage\fi
%% Now input the file, then run the end-of-article hook (reinstated 28
  %% Sep 2001)
  {\input{\@articledir/\articlename.tex}\@endofarticlehook}}
%    \end{macrocode}
% Command to separate articles into sessions. Write an entry in the
% TOC, makes sure we start on an odd page (added 06 Jun 2000), then do
% a title page if \cs{SessionTitlePages} was called. 
% \changes{3.12}{04 Feb 2003}{Added option to have title page for each
% session. Turn it on with \cs{SessionTitlePages}.}
%    \begin{macrocode}
\newcommand{\session}[1]{
  %% this does nothing except write an entry in the TOC
  \addcontentsline{toc}{part}{\rmaa@TOCPartStyle{#1}}
  \addtocontents{htc}{\protect\HyperTOCpart{#1}}
  %% and also makes sure we start on an odd page (added 06 Jun 2000)
  \clearpage\thispagestyle{empty}\cleardoublepage
  \stepcounter{part}
  \if@sessiontitlepages{%
    \onecolumn
    \thispagestyle{empty}
    \renewcommand{\articlename}{Section\thepart TitlePage}
    \label{RMAAFirstPage}
    \make@session@title{#1}
    \label{RMAALastPage}
    \clearpage\thispagestyle{empty}\cleardoublepage
  }\fi
  }
\newcommand{\make@session@title}[1]{%
  {%
    \setlength\parindent{0pt}
    \vspace*{\stretch{1}}
    \rule{\linewidth}{1mm}
    \begin{flushright}\Huge\bfseries\boldmath
      #1
    \end{flushright}
    \rule{\linewidth}{1mm}
    \vspace*{\stretch{2}}
  }
}
%    \end{macrocode}
% Command to turn on separate title pages for each session. 
%    \begin{macrocode}
\newcommand{\SessionTitlePages}{\@sessiontitlepagestrue}
\newcommand{\BoldSessionTitles}{%
  \renewcommand{\rmaa@TOCPartStyle}[1]{\textbf{\rmaa@TOCTitleStyle{##1}}}
  }
\newcommand{\ArticleDir}[1]{\renewcommand{\@articledir}{#1/}\graphicspath{{1/}}}
\newcommand{\@articledir}{.}
%\newcommand{\ps@rmaafront}{\ps@rmaa}% page style for front matter (same as ever, actually)
\newcommand{\ps@rmaafront}{% page style for front matter:  page # in each foot
  \renewcommand\@oddfoot{\hfil{\thepage}}%
  \renewcommand\@evenfoot{{\thepage}\hfil}%
  % shorttitle in both heads
  \renewcommand\@evenhead{\hfil\@shorttitle\hfil}%
  \renewcommand\@oddhead{\@evenhead}
  }
\newcommand{\ps@rmaafrontfirstpage}{% page style for first page of front matter
  % page # in each foot, on outside edge
  \renewcommand\@oddfoot{\hfil\thepage}%
  \renewcommand\@evenfoot{\thepage\hfil}%
  % nothing in the even heads
  \renewcommand\@evenhead{\relax}%
  \renewcommand\@oddhead{\relax}%
  }
%% slight modification to the standard \frontmatter and \mainmatter
%% commands, just to set the \pagestyle
\renewcommand\frontmatter{%
  \latexonecolumn
  \cleardoublepage
  \@frontmattertrue
  \@mainmatterfalse
  \pagenumbering{roman}
  \pagestyle{rmaafront}
  \addtocounter{page}{2} % leave space for the title page
  }
\renewcommand\mainmatter{%
  \pagestyle{empty}
  \if@rmtwocol\twocolumn\fi
  \cleardoublepage
  \@frontmatterfalse
  \@mainmattertrue
  \pagenumbering{arabic}
  \pagestyle{rmaa}
  }
\renewcommand\backmatter{%
  \clearpage\pagestyle{empty}\cleardoublepage
  \@mainmatterfalse
  %% put extra space in the TOC
  \addtocontents{toc}{\bigskip}
  }
%    \end{macrocode}
%
% \subsection{Front sections}\label{sec:front-sections}
% \begin{macro}{\notitle}
% \begin{macro}{\shortname}
% \begin{macro}{frontsection}
% The environment \texttt{frontsection} is for doing a section in the
% front matter (preface, list of participants, or whatever). It has 
% one optional argument:
% \begin{enumerate}
% \item[\#1] Short title (no embedded spaces) for
% use in making PS files. E.g., ``list\_of\_participants''.   
% \end{enumerate}
% and two mandatory arguments: 
% \begin{enumerate}
% \item[\#2] Title of section (e.g., `Preface')
% \item[\#3] Commands to be evaluated before processing the contents
% of the environment. This is mainly useful for putting in the things
% that would go in the preamble of a normal article  (e.g.,
% |\listofauthors|, |\indexauthor|, |\notitle|) 
% \end{enumerate}
%    \begin{macrocode}
\newcommand{\notitle}{\@titlevisiblefalse} \newcommand{\shortname}{}
\newenvironment{frontsection}[3][]{%
  \begingroup
  \renewcommand{\shortname}{#1}
  \ifx\shortname\@rmaaEmpty
     \renewcommand{\articlename}{#2}
  \else
     \renewcommand{\articlename}{\shortname}
  \fi
  \@begofarticlehook
  \@titlevisibletrue % default
  \thispagestyle{rmaafrontfirstpage}
  \listofauthors{} % blank out this in case it is not specified
  \title{{#2}}\shorttitle{#2}\shortauthor{#2}
  #3 % evaluate the command given as second mandatory arg
  \addcontentsline{toc}{chapter}{\protect\@fiddleauthors{\@TOCtitle}{\rmaaTOCauthors}}%
  \addtocontents{htc}{\protect\HyperTOCEntry{\@rawtitle}{\@rawauthors}{\articlename}}
  \if@titlevisible
     \vspace*{1cm}\centerline{\rmaa@FrontSecStyle{\@title}}\vspace*{1cm}
  \fi
  \setlength{\parskip}{3ex}
  }{\frontlabel{RMAALastPage}\clearpage\endgroup}
%    \end{macrocode}
% \end{macro}
% \end{macro}
% \end{macro}
%    \begin{macrocode}
%% Definition of the command \frontlabel, which rights out something like `F4' to
%% the .aux file, instead of `\textlatin  {i}'. This is so it is easy to post-process
%% the file for making PS files of individual sections
\newcommand{\frontlabel}[1]{%
  \@bsphack
  \protected@write\@auxout{}%
     {\string\newlabel{\articlename:#1}{{}{F\@arabic{\c@page}}}}
  \@esphack
  }
%% User command to say X, Y, Z,\\ Editors of the Proceedings
\newcommand{\EditorSignature}{\vspace*{\parskip}\rightline{\@editors}%
  \rightline{Editors of the proceedings}}
%% Fiddle with \thesection, etc. so that they don't include the chapter number
%% Basically go back to the `article' definitions of these
\renewcommand{\thesection}{\@arabic\c@section}
\renewcommand{\thefigure}{\@arabic\c@figure}
\renewcommand{\thetable}{\@arabic\c@table}

\newcommand{\@oldargs}{}
\newcommand{\@newargs}{}
%% add TOC entry for each paper
%\renewcommand{\title}[1]%
%{\renewcommand{\@title}{\rmaa@TitleStyle{#1}}%
%\addtocounter{chapter}{1}\addcontentsline{toc}{chapter}{#1}%
%}

%% Formatting of the table of contents
%% dimensions of various parts
\newlength{\TOCtextwidth}\setlength{\TOCtextwidth}{12.65cm}
\newlength{\TOCpagenumwidth}\setlength{\TOCpagenumwidth}{1.3cm}
\newlength{\TOCentrywidth}\setlength{\TOCentrywidth}{\TOCtextwidth}
\addtolength{\TOCentrywidth}{-\TOCpagenumwidth} 
\newlength{\TOCextramargin}\setlength{\TOCextramargin}{1.25cm}
\newcommand{\TOCbox}[1]{\hspace{\TOCextramargin}\parbox{\TOCtextwidth}{#1}\par}
\newlength{\TOCextramarginx}\setlength{\TOCextramarginx}{\TOCextramargin}\addtolength{\TOCextramarginx}{4.3cm}
\newlength{\TOCtextwidthx}\setlength{\TOCtextwidthx}{\TOCtextwidth}\addtolength{\TOCtextwidthx}{-4.3cm}
\newcommand{\TOCboxx}[1]{\hspace{\TOCextramarginx}\parbox{\TOCtextwidthx}{#1}\par}

\newcommand{\RMAAbannerimage}{rmsc}
\renewcommand\tableofcontents{%
  \renewcommand{\articlename}{CONTENTS}
  \begingroup
  \@begofarticlehook
  \if@twocolumn
     \@restonecoltrue\onecolumn
  \else
     \@restonecolfalse
  \fi
  %% 9pt for journal style
  \if@proceedings\else\small\fi
  \if@newstyle\small\fi %% same for new style
  \if@proceedings\thispagestyle{rmaafrontfirstpage}\else\thispagestyle{empty}\fi
  \shortauthor{Contents}\shorttitle{Contents}
  \vspace*{18mm}
  \if@printbanner
  %% put in the banner header (need to get better image of this)
     \TOCbox{\hspace*{-0.15cm}\includegraphics[width=0.9\TOCtextwidth]{\RMAAbannerimage}}
     \vspace*{8mm}
  %% volume and publication date (should really be in funny font)
     \TOCboxx{\@volumen\\[2mm] \@mes\ \@anyo}
  \else
  %% put in the banner header (need to get better image of this)
     \TOCbox{\hspace*{-0.15cm}\includegraphics[width=0.9\TOCtextwidth]{\RMAAbannerimage}}
     \vspace*{12.5mm}
  %% volume and publication date (should really be in funny font)
     \TOCboxx{~\\[2mm] ~}
  \fi
  \vspace*{16mm}
  %% This is the command that fiddles with the TOC entry so that the
  %% title is flushleft but the authors are flushright
  \renewcommand{\@fiddleauthors}[2]{%
    % first put the paper title
    \renewcommand{\\}{\relax} % temporarily turn off explicit line breaks just in case
    \renewcommand{\altaffilmark}[1]{\relax} % ditto with this
    \renewcommand{\textsuperscript}[1]{\relax} % ditto with this 
    ##1
    % now save the author list in a box
    \sbox\@tempboxa{##2}%
    % check that box is shorter than the total width available for the entry
    \ifdim \wd\@tempboxa <\TOCentrywidth
       % if so, then put it right here, flush against right margin
       \penalty10000\hskip1em plus10em\penalty0\hbox{}\nobreak\hfill\box\@tempboxa
       % Note that we force a space of at least 1em between title and authorlist, 
       % stretchable to 11em since we may have very short titles but still with an 
       % authorlist too long to go on the same line.  The only linebreak allowed is 
       % directly before the authorlist.
    \else
       % if it is too long, we set the original authorlist text in its own parbox
       \par\vskip1ex\parbox[b]{\TOCentrywidth}{\raggedleft ##2}
    \fi
    }
  \@starttoc{toc}%
  \if@restonecol\twocolumn\fi
  \frontlabel{RMAALastPage}
  \clearpage
  \thispagestyle{empty}
  \cleardoublepage % make sure we start on an odd page
  \endgroup

  %%% ALSO, do the HTML TOC
  %% The file that will be used to make the HTML TOC
  \@starttoc{htc}
  %% All the stuff in this file does nothing in the book style
  \addtocontents{htc}{\protect\HyperTOCmeeting{\@conftitle}{\@reunionplace}{\@reuniondate}}
  \addtocontents{htc}{\protect\HyperTOCeditors{\@editors}}
  \addtocontents{htc}{\protect\HyperTOCpreamble}
  \addtocontents{htc}{\protect\HyperTOCEntry{Table of Contents}{}{\articlename}}

  %% and the ADS file
  \ADS@openfile
  }
%% This is now much simpler
\renewcommand\l@chapter[2]{%TOC entry for articles
  \ifnum \c@tocdepth >\m@ne
    \vskip 1.0em \@plus\p@
    \TOCbox{%
      \parbox[b]{\TOCentrywidth}{#1}\parbox[b]{\TOCpagenumwidth}{\hfill#2}
    }
  \fi} 
\renewcommand\l@part[2]{%TOC entry for sessions: centered, no page number
  \ifnum \c@tocdepth >\m@ne
    \vskip 1.5em \@plus1.5em \goodbreak
    \TOCbox{%
      \parbox[b]{\TOCentrywidth}{\hfil#1\hfil}\hfill
    }
  \fi\nobreak} 

%% Formatting theindex
%% strongly encourage page breaks between index entries
\renewcommand{\@idxitem}{\par\penalty-1000\hangindent.45cm} 
\renewcommand{\indexspace}{\par\penalty-1000\vskip 10\p@ \@plus15\p@ \@minus3\p@\relax}
\renewenvironment{theindex}{%
  %% first, make sure we start on an odd page (06 Jun 2000)
  \clearpage\thispagestyle{empty}\cleardoublepage
  %% this looks a bit like an article but is two columns and without abstract, etc
  \renewcommand{\articlename}{INDEX}
  \title{Author Index}\shortauthor{Author Index}\shorttitle{Author Index}
  \listofauthors{}
  \addcontentsline{toc}{chapter}{\protect\@fiddleauthors{\@TOCtitle}{\rmaaTOCauthors}}%
  \addtocontents{htc}{\protect\HyperTOCEntry{\@rawtitle}{\@rawauthors}{\articlename}}
  \pagestyle{rmaa}
  \thispagestyle{rmaafirstpage}
  \@begofarticlehook
  %% 9pt for journal style
  \if@proceedings\else\small\fi
  \if@newstyle\small\fi
  \twocolumn[\@makeminititle]
  %% vertical spacing (see also definition of \@idxitem and \indexspace
  \parskip\z@ \@plus .6\p@\relax % allow quite a lot of leaway here to avoid breaks inside entries
  %% horizontal spacing, breaking and glue parameters (see also the stuff in the definition 
  %% of \indexauthor
  \parindent\z@ 
  \exhyphenpenalty10000 % don't break on explicit hyphens, it always looks bad
  \rightskip=0pt plus.2em % allow a tiny bit of leeway at the ends of lines
  \spaceskip.3333em plus.5em minus.15em % a lot of leeway between words
  \xspaceskip.5em plus.5em minus.4em % this seems to have no effect
  \let\item\@idxitem
  \renewcommand{\\}{\relax} % temporarily turn off explicit line breaks just in case
  \renewcommand{\altaffilmark}[1]{\relax} % ditto with this
  \renewcommand{\textsuperscript}[1]{\relax} % ditto with this
  }{\label{RMAALastPage}\clearpage} % clearpage necessary so that we know about \shortauthor



%% end definitions for book styles
\fi

%    \end{macrocode}
% In order to do away with plates, we sometimes want a figure to
% appear in color only in the online version. The following command
% can bve used in a figure caption to typeset a note to this effect. 
%    \begin{macrocode}
\newcommand\ColorOnlineOnly{%
  \textsc{Note: this figure is available in color in the electronic
    version of this article, obtainable from}
  \texttt{http://\discretionary{}{}{}www.\discretionary{}{}{}astroscu.\discretionary{}{}{}unam.\discretionary{}{}{}mx/\discretionary{}{}{}\~{}rmaa/}.}
%% Plates
%% Define 2 counters: plate is the number in the whole document, plateloop is used 
%% internally by \outputplates
\newcounter{plate}
\newcounter{plateloop}
%% Currently limited to 15 plates in document
%% when we do a label, we need to write to \newlabel entries to the .aux file, one 
%% that says RMAAPlateX so that it can be referred to in \ouputplates, and one that
%% says articlename:RMAAPlateX so that the dvips program can find it and asociate it
%% with the right article
\newcommand{\@platelabeli}{\oldlabel{RMAAPlate1}\label{RMAAPlate1}}
\newcommand{\@platelabelii}{\oldlabel{RMAAPlate2}\label{RMAAPlate2}}
\newcommand{\@platelabeliii}{\oldlabel{RMAAPlate3}\label{RMAAPlate3}}
\newcommand{\@platelabeliv}{\oldlabel{RMAAPlate4}\label{RMAAPlate4}}
\newcommand{\@platelabelv}{\oldlabel{RMAAPlate5}\label{RMAAPlate5}}
\newcommand{\@platelabelvi}{\oldlabel{RMAAPlate6}\label{RMAAPlate6}}
\newcommand{\@platelabelvii}{\oldlabel{RMAAPlate7}\label{RMAAPlate7}}
\newcommand{\@platelabelviii}{\oldlabel{RMAAPlate8}\label{RMAAPlate8}}
\newcommand{\@platelabelix}{\oldlabel{RMAAPlate9}\label{RMAAPlate9}}
\newcommand{\@platelabelx}{\oldlabel{RMAAPlate10}\label{RMAAPlate10}}
\newcommand{\@platelabelxi}{\oldlabel{RMAAPlate11}\label{RMAAPlate11}}
\newcommand{\@platelabelxii}{\oldlabel{RMAAPlate12}\label{RMAAPlate12}}
\newcommand{\@platelabelxiii}{\oldlabel{RMAAPlate13}\label{RMAAPlate13}}
\newcommand{\@platelabelxiv}{\oldlabel{RMAAPlate14}\label{RMAAPlate14}}
\newcommand{\@platelabelxv}{\oldlabel{RMAAPlate15}\label{RMAAPlate15}}
\newcommand{\@outputplatelabeli}{\oldlabel{RMAAOutputPlate1}}
\newcommand{\@outputplatelabelii}{\oldlabel{RMAAOutputPlate2}}
\newcommand{\@outputplatelabeliii}{\oldlabel{RMAAOutputPlate3}}
\newcommand{\@outputplatelabeliv}{\oldlabel{RMAAOutputPlate4}}
\newcommand{\@outputplatelabelv}{\oldlabel{RMAAOutputPlate5}}
\newcommand{\@outputplatelabelvi}{\oldlabel{RMAAOutputPlate6}}
\newcommand{\@outputplatelabelvii}{\oldlabel{RMAAOutputPlate7}}
\newcommand{\@outputplatelabelviii}{\oldlabel{RMAAOutputPlate8}}
\newcommand{\@outputplatelabelix}{\oldlabel{RMAAOutputPlate9}}
\newcommand{\@outputplatelabelx}{\oldlabel{RMAAOutputPlate10}}
\newcommand{\@outputplatelabelxi}{\oldlabel{RMAAOutputPlate11}}
\newcommand{\@outputplatelabelxii}{\oldlabel{RMAAOutputPlate12}}
\newcommand{\@outputplatelabelxiii}{\oldlabel{RMAAOutputPlate13}}
\newcommand{\@outputplatelabelxiv}{\oldlabel{RMAAOutputPlate14}}
\newcommand{\@outputplatelabelxv}{\oldlabel{RMAAOutputPlate15}}
\newcommand{\@platerefi}{\oldpageref{RMAAPlate1}}
\newcommand{\@platerefii}{\oldpageref{RMAAPlate2}}
\newcommand{\@platerefiii}{\oldpageref{RMAAPlate3}}
\newcommand{\@platerefiv}{\oldpageref{RMAAPlate4}}
\newcommand{\@platerefv}{\oldpageref{RMAAPlate5}}
\newcommand{\@platerefvi}{\oldpageref{RMAAPlate6}}
\newcommand{\@platerefvii}{\oldpageref{RMAAPlate7}}
\newcommand{\@platerefviii}{\oldpageref{RMAAPlate8}}
\newcommand{\@platerefix}{\oldpageref{RMAAPlate9}}
\newcommand{\@platerefx}{\oldpageref{RMAAPlate10}}
\newcommand{\@platerefxi}{\oldpageref{RMAAPlate11}}
\newcommand{\@platerefxii}{\oldpageref{RMAAPlate12}}
\newcommand{\@platerefxiii}{\oldpageref{RMAAPlate13}}
\newcommand{\@platerefxiv}{\oldpageref{RMAAPlate14}}
\newcommand{\@platerefxv}{\oldpageref{RMAAPlate15}}
\newcounter{figplatei}
\newcounter{figplateii}
\newcounter{figplateiii}
\newcounter{figplateiv}
\newcounter{figplatev}
\newcounter{figplatevi}
\newcounter{figplatevii}
\newcounter{figplateviii}
\newcounter{figplateix}
\newcounter{figplatex}
\newcounter{figplatexi}
\newcounter{figplatexii}
\newcounter{figplatexiii}
\newcounter{figplatexiv}
\newcounter{figplatexv}
%% Printed output should use the counter `plate'
\renewcommand{\theplate}{\@arabic\c@plate}
\renewcommand{\theplateloop}{\@arabic\c@plateloop}
%% Command to output all the saved plates in this paper
\newcommand{\outputplates}{%
  \setcounter{plateloop}{0}
  \if@book\else\@endofarticlehook\fi % not entirely sure why, but this
                                     % seems to be needed
  \clearpage % first output the last page of text
  % now set pagestyle for blank even pages and make sure we start on an odd page
  \pagestyle{empty}
  \cleardoublepage
  \onecolumn
  \parindent\z@
  % first put in an extra sheet
  % the main purpose of this is so there is always a blank even page before 
  % each plate for filling in the PS file if the article has an odd number
  % of pages (see splitbook.sh)
  \null\vfil\centerline{PLATES}\vfil\clearpage\pagestyle{rmaaplate}\cleardoublepage
  % output each plate in turn
  \@whilenum\value{plateloop}<\value{plate}\do{%
    \addtocounter{plateloop}{1}
    % restore the figure counter from \@figcounti, \@figcountii, etc
    \setcounter{figure}{\value{figplate\roman{plateloop}}}
    % spew out the plate (in its own group for safety)
    {\csname @plate\roman{plateloop}\endcsname\par}
    }
  \pagestyle{rmaa}
  % reset the page number as if nothing had happened
  \addtocounter{page}{-\value{plate}}
  \addtocounter{page}{-\value{plate}}
  }
\newcommand\figcaption{%
  \def\@captype{figure}
  \renewcommand{\@makecaption}{\@makefigcap}
  \caption}
\newcommand\tabcaption{%
  \def\@captype{table}
  \renewcommand{\@makecaption}{\@maketabcap}
  \caption}

\newcommand{\plate}[2][1]{%
  %% User command to make a plate page: 
  %% Mandatory argument is what goes on the plate - this should not 
  %% include any floating environments (e.g. figure). Instead, \figcaption 
  %% should be used to do the captions. 
  %% Optional argument is the number of figures (i.e. captions) on the plate
  %% page (default is 1). This is necessary for incrementing the figure counter. 
  %%
  %% Note that we have to use a \gdef rather than an \xdef when we 
  %% save the plate output in a command. This means that in the book styles when it 
  %% comes to spewing out the plate, we will have forgotten about any user macros 
  %% that may have been defined. This can be got round by repeating the macros inside the plate.
  %% (Presumably, one would use a \newcommand rather than \renewcommand if they were new.)
  % increment global and local counters
  \addtocounter{plate}{1}
  % write appropriate label to .aux file
  \csname @platelabel\roman{plate}\endcsname
  % define a global macro called \@platei, \@plateii, \@plateiii, \@plateiv, ...
  \expandafter\gdef\csname @plate\roman{plate}\endcsname{%
    \csname @outputplatelabel\roman{plateloop}\endcsname
    \null\vfil #2 \vfil
    \cleardoublepage
    }
  % save the current shortauthor and shorttitle
  % need to be \xdef to take current values of \@shortauthor, etc 
  \expandafter\xdef\csname @shorta\roman{plate}\endcsname{\@shortauthor}
  \expandafter\xdef\csname @shortt\roman{plate}\endcsname{\@shorttitle}
  % save the current value of the figure counter for later use in \outputplates
  \setcounter{figplate\roman{plate}}{\value{figure}}
  % increment the figure counter according to the optional argument
  \addtocounter{figure}{#1}
  }
%% page style for plate pages 
\newcommand{\ps@rmaaplate}{%
  % Plate X (see page Y) in the foot
  \renewcommand\@oddfoot{\hfil\csname @shorta\roman{plateloop}\endcsname\ 
    (see page~\csname @plateref\roman{plateloop}\endcsname)\hfil}%
  \renewcommand\@evenfoot{\relax}%
  % nothing in even heads (blank pages)
  \renewcommand\@evenhead{\relax}
  % short title and page num in the odd heads
  \renewcommand\@oddhead{\hfil\csname @shortt\roman{plateloop}\endcsname%
    \hfil\llap{PLATE~\theplateloop}}
  }





%% Things for the User Guide
\if@guide
%% stolen from doc.sty
\def\MakeShortVerb#1{%
  \expandafter\ifx\csname cc\string#1\endcsname\relax
  \@shortvrbinfo{Made }{#1}%
  \add@special{#1}%
  \expandafter
  \xdef\csname cc\string#1\endcsname{\the\catcode`#1}%
  \begingroup
  \catcode`\~\active  \lccode`\~`#1%
  \lowercase{%
    \global\expandafter\let
    \csname ac\string#1\endcsname~%
    \gdef~{\verb~}}%
  \endgroup
  \global\catcode`#1\active
  \else
  \@shortvrbinfo\@empty{#1 already}%
  \fi}
\def\DeleteShortVerb#1{%
  \expandafter\ifx\csname cc\string#1\endcsname\relax
  \else
  \@shortvrbinfo{Deleted }{#1 as}%
  \rem@special{#1}%
  \global\catcode`#1\csname cc\string#1\endcsname
  \global \expandafter\let \csname cc\string#1\endcsname \relax
  \ifnum\catcode`#1=\active
  \begingroup
  \catcode`\~\active   \lccode`\~`#1%
  \lowercase{%
    \global\expandafter\let\expandafter~%
    \csname ac\string#1\endcsname}%
  \endgroup \fi \fi}
\def\@shortvrbinfo#1#2{%
  \PackageInfo{doc}{%
    #1\expandafter\@gobble\string#2 a short reference
    for \string\verb}}
\def\add@special#1{%
  \rem@special{#1}%
  \expandafter\gdef\expandafter\dospecials\expandafter
  {\dospecials \do #1}%
  \expandafter\gdef\expandafter\@sanitize\expandafter
  {\@sanitize \@makeother #1}}
\def\rem@special#1{%
  \def\do##1{%
    \ifnum`#1=`##1 \else \noexpand\do\noexpand##1\fi}%
  \xdef\dospecials{\dospecials}%
  \begingroup
  \def\@makeother##1{%
    \ifnum`#1=`##1 \else \noexpand\@makeother\noexpand##1\fi}%
  \xdef\@sanitize{\@sanitize}%
  \endgroup}
%% define | as shortcut for \verbatim
\MakeShortVerb{\|}
%% comand to specify arguments to LaTeX commands
\newcommand{\rmArg}[1]{%
  \ensuremath{\langle}{\scriptsize \textsf{\MakeUpperCase{#1}}}\ensuremath{\rangle}}
%\newcommand{\rmArg}[1]{\ensuremath{\langle}\textsc{#1}\ensuremath{\rangle}}
\providecommand{\BibTeX}{{\rmfamily B\kern-.05em%
    \textsc{i\kern-.025em b}\kern-.08em%
    T\kern-.1667em\lower.7ex\hbox{E}\kern-.125emX}}
%% Environment for code examples
\newenvironment{Code}{%
  \begin{list}{}{\setlength{\leftmargin}{0.05\columnwidth}} 
  \item[]}{\end{list}}
\fi
%</package>
%    \end{macrocode}

% \section{User commands for astronomical articles}
% \label{sec:user-commands}
% Some random definitions. These go both into the main class file and
% also into |adshtml.sty|.
% 
% Definitions of the journal name in different incarnations.  
%    \begin{macrocode}
%<*package|adshtml-style>
\newcommand{\RMAA}{RMxAA}
\newcommand{\RMAAlong}{Revista Mexicana de Astronom\'{\i}a y Astrof\'{\i}sica}
\newcommand{\RMAASC}{RMxAC}
\newcommand{\RMAASClong}{\RMAA\ (Serie de Conferencias)}
%    \end{macrocode}
% \begin{macro}{\ibidrule}
% This puts a horizontal rule with a period after it. It is used when
% the author list of a reference is identical to the preceding one. 
%    \begin{macrocode}
\newcommand{\ibidrule}{\rule[-0.03cm]{4.7em}{.01cm}.\ }
%    \end{macrocode}
% \end{macro}
% \begin{macro}{\ThisVolume}
% References to other articles in the same volume can be coded as,
% e.g., |\bibitem{} Bloggs, J. \ThisVolume{bloggs}|.
%    \begin{macrocode}
\newcommand{\ThisVolume}[1]{%
  \@year, \if@proceedings\RMAASC\else\RMAA\fi, \@volume,
  \if@book\oldpageref{#1:RMAAFirstPage}\else 000\fi\ (this volume)}
%    \end{macrocode}
% \end{macro}
% Other stuff that may be useful in ref lists.
%    \begin{macrocode}
\newcommand{\etal}{et al.\@}
%    \end{macrocode}
% Various new commands from AAS macros
%    \begin{macrocode}
\def\acknowledgments{\vskip 3ex plus .8ex minus .4ex}
\let\acknowledgements=\acknowledgments                  % second spelling
\let\jnl@style=\rm
\def\ref@jnl#1{{\jnl@style#1}}
\DeclareRobustCommand\aj{\ref@jnl{AJ}}                   
\DeclareRobustCommand\araa{\ref@jnl{ARA\&A}}             
\DeclareRobustCommand\apj{\ref@jnl{ApJ}}                 
\DeclareRobustCommand\apjl{\ref@jnl{ApJ}}                
\DeclareRobustCommand\apjs{\ref@jnl{ApJS}}               
\DeclareRobustCommand\ao{\ref@jnl{Appl.~Opt.}}           
\DeclareRobustCommand\apss{\ref@jnl{Ap\&SS}}             
\DeclareRobustCommand\aap{\ref@jnl{A\&A}}                
\DeclareRobustCommand\aapr{\ref@jnl{A\&A~Rev.}}          
\DeclareRobustCommand\aaps{\ref@jnl{A\&AS}}              
\DeclareRobustCommand\azh{\ref@jnl{AZh}}                 
\DeclareRobustCommand\baas{\ref@jnl{BAAS}}               
\DeclareRobustCommand\jrasc{\ref@jnl{JRASC}}             
\DeclareRobustCommand\memras{\ref@jnl{MmRAS}}            
\DeclareRobustCommand\mnras{\ref@jnl{MNRAS}}             
\DeclareRobustCommand\pra{\ref@jnl{Phys.~Rev.~A}}        
\DeclareRobustCommand\prb{\ref@jnl{Phys.~Rev.~B}}        
\DeclareRobustCommand\prc{\ref@jnl{Phys.~Rev.~C}}        
\DeclareRobustCommand\prd{\ref@jnl{Phys.~Rev.~D}}        
\DeclareRobustCommand\pre{\ref@jnl{Phys.~Rev.~E}}        
\DeclareRobustCommand\prl{\ref@jnl{Phys.~Rev.~Lett.}}    
\DeclareRobustCommand\pasp{\ref@jnl{PASP}}               
\DeclareRobustCommand\pasj{\ref@jnl{PASJ}}               
\DeclareRobustCommand\qjras{\ref@jnl{QJRAS}}             
\DeclareRobustCommand\skytel{\ref@jnl{S\&T}}             
\DeclareRobustCommand\solphys{\ref@jnl{Sol.~Phys.}}      
\DeclareRobustCommand\sovast{\ref@jnl{Soviet~Ast.}}      
\DeclareRobustCommand\ssr{\ref@jnl{Space~Sci.~Rev.}}     
\DeclareRobustCommand\zap{\ref@jnl{ZAp}}                 
\DeclareRobustCommand\nat{\ref@jnl{Nature}}              
\DeclareRobustCommand\iaucirc{\ref@jnl{IAU~Circ.}}
\DeclareRobustCommand\aplett{\ref@jnl{Astrophys.~Lett.}}
\DeclareRobustCommand\apspr{\ref@jnl{Astrophys.~Space~Phys.~Res.}}
\DeclareRobustCommand\bain{\ref@jnl{Bull.~Astron.~Inst.~Netherlands}}
\DeclareRobustCommand\fcp{\ref@jnl{Fund.~Cosmic~Phys.}}
\DeclareRobustCommand\gca{\ref@jnl{Geochim.~Cosmochim.~Acta}}
\DeclareRobustCommand\grl{\ref@jnl{Geophys.~Res.~Lett.}}
\DeclareRobustCommand\jcp{\ref@jnl{J.~Chem.~Phys.}}      
\DeclareRobustCommand\jgr{\ref@jnl{J.~Geophys.~Res.}}    
\DeclareRobustCommand\jqsrt{\ref@jnl{J.~Quant.~Spec.~Radiat.~Transf.}}
\DeclareRobustCommand\memsai{\ref@jnl{Mem.~Soc.~Astron.~Italiana}}
\DeclareRobustCommand\nphysa{\ref@jnl{Nucl.~Phys.~A}}
\DeclareRobustCommand\physrep{\ref@jnl{Phys.~Rep.}}
\DeclareRobustCommand\physscr{\ref@jnl{Phys.~Scr}}
\DeclareRobustCommand\planss{\ref@jnl{Planet.~Space~Sci.}}
\DeclareRobustCommand\procspie{\ref@jnl{Proc.~SPIE}}
\let\astap=\aap
\let\apjlett=\apjl
\let\apjsupp=\apjs
\let\applopt=\ao
\def\phn{\phantom{0}}
\def\phd{\phantom{.}}
\def\phs{\phantom{\ensuremath{-}}}
\def\phm#1{\phantom{#1}}
%    \end{macrocode}
% All these are redone as robust and using more modern commands where
% possible. The robustness is so they will not get expanded if used in
% abstracts and titles when these are passed to the ADS file. 
% \changes{3.16}{2003/04/07}{
%  Added \cs{Color} command for doing photometric colors.}
%    \begin{macrocode}
%%\def\sun{\hbox{$\odot$}}
\DeclareRobustCommand\sun{\ensuremath{\odot}}
%%\def\earth{\hbox{$\oplus$}}
\DeclareRobustCommand\earth{\ensuremath{\oplus}}
%%\def\lesssim{\mathrel{\hbox{\rlap{\hbox{\lower4pt\hbox{$\sim$}}}\hbox{$<$}}}}
%%\def\gtrsim{\mathrel{\hbox{\rlap{\hbox{\lower4pt\hbox{$\sim$}}}\hbox{$>$}}}}
\DeclareRobustCommand\lesssim{\mathrel{\hbox{\smaller\smaller\rlap{\hbox{\lower3pt\hbox{$\sim$}}}\raise1.5pt\hbox{$<$}}}}
\DeclareRobustCommand\gtrsim{\mathrel{\hbox{\smaller\smaller\rlap{\hbox{\lower3pt\hbox{$\sim$}}}\raise1.5pt\hbox{$>$}}}}
%%\def\lta{\mathrel{\spose{\lower 3pt\hbox{$\mathchar"218$}}
%%     \raise 2.0pt\hbox{$\mathchar"13C$}}}
%%\def\gta{\mathrel{\spose{\lower 3pt\hbox{$\mathchar"218$}}
%%     \raise 2.0pt\hbox{$\mathchar"13E$}}}
\DeclareRobustCommand\sq{\hbox{\rlap{$\sqcap$}$\sqcup$}}
%%\def\arcdeg{\hbox{$^\circ$}}
%%\def\arcmin{\hbox{$^\prime$}}
%%\def\arcsec{\hbox{$^{\prime\prime}$}}
\DeclareRobustCommand\arcdeg{\ensuremath{^\circ}}
\DeclareRobustCommand\arcmin{\ensuremath{^\prime}}
\DeclareRobustCommand\arcsec{\ensuremath{^{\prime\prime}}}
%%\def\fd{\hbox{$.\!\!^{\rm d}$}}
%%\def\fh{\hbox{$.\!\!^{\rm h}$}}
%%\def\fm{\hbox{$.\!\!^{\rm m}$}}
%%\def\fs{\hbox{$.\!\!^{\rm s}$}}
%%\def\fdg{\hbox{$.\!\!^\circ$}}
%%\def\farcm{\hbox{$.\mkern-4mu^\prime$}}
%%\def\farcs{\hbox{$.\!\!^{\prime\prime}$}}
%%\def\fp{\hbox{$.\!\!^{\scriptscriptstyle\rm p}$}}
\DeclareRobustCommand\fd{\ensuremath{.\!\!^\mathrm{d}}}
\DeclareRobustCommand\fh{\ensuremath{.\!\!^\mathrm{h}}}
\DeclareRobustCommand\fm{\ensuremath{.\!\!^\mathrm{m}}}
\DeclareRobustCommand\fs{\ensuremath{\mkern0.5mu.\mkern-5mu^\mathrm{s}\mkern-1mu}}
\DeclareRobustCommand\fdg{\ensuremath{.\!\!^\circ}}
\DeclareRobustCommand\farcm{\ensuremath{.\mkern-4mu^\prime}}
\DeclareRobustCommand\farcs{\ensuremath{\mkern1mu.\mkern-6mu^{\prime\prime}\mkern-1.5mu}}
%%\DeclareRobustCommand\farcs{\ensuremath{.\!\!^{\prime\prime}}}
\DeclareRobustCommand\fp{\ensuremath{.\!\!^{\scriptscriptstyle\rm p}}}
%%\def\micron{\hbox{$\mu$m}}
\DeclareRobustCommand\micron{\ensuremath{\mu\mathrm{m}}}
\let\la=\lesssim                        
\let\ga=\gtrsim
%%\def\case#1#2{\hbox{$\frac{#1}{#2}$}}
%%\def\slantfrac#1#2{\hbox{$\,^#1\!/_#2$}}
%%\def\onehalf{\slantfrac{1}{2}}
%%\def\onethird{\slantfrac{1}{3}}
%%\def\twothirds{\slantfrac{2}{3}}
%%\def\onequarter{\slantfrac{1}{4}}
%%\def\threequarters{\slantfrac{3}{4}}
\DeclareRobustCommand\case[2]{\ensuremath{\frac{#1}{#2}}}
\DeclareRobustCommand\slantfrac[2]{\ensuremath{\,^#1\!/_#2}}
\DeclareRobustCommand\onehalf{\slantfrac{1}{2}}
\DeclareRobustCommand\onethird{\slantfrac{1}{3}}
\DeclareRobustCommand\twothirds{\slantfrac{2}{3}}
\DeclareRobustCommand\onequarter{\slantfrac{1}{4}}
\DeclareRobustCommand\threequarters{\slantfrac{3}{4}}
%%\def\ubvr{\hbox{$U\!BV\!R$}}            
%%\def\ub{\hbox{$U\!-\!B$}}               
%%\def\bv{\hbox{$B\!-\!V$}}               
%%\def\vr{\hbox{$V\!-\!R$}}               
%%\def\ur{\hbox{$U\!-\!R$}}
%    \end{macrocode}
% New commands for ``photometric'' colors 
%    \begin{macrocode}
\DeclareRobustCommand\ubvr{\ensuremath{U\!BV\!R}}            
\DeclareRobustCommand\Color[2]{%
  \ensuremath{%
    #1\!\ifx#1V\!\fi-\!#2
  }}
\DeclareRobustCommand\ub{\Color{U}{B}}
\DeclareRobustCommand\bv{\Color{B}{V}}
\DeclareRobustCommand\vr{\Color{V}{R}}
\DeclareRobustCommand\ur{\Color{U}{R}}
%% New version of the \ion command, that actually works! 
%    \end{macrocode}
% \begin{macro}{\ion}
% This version of the |\ion| command allows you to write either
% |\ion{N}{4}| or |\ion{N}{iv}| in order to get N~IV. 
%    \begin{macrocode}
\DeclareRobustCommand\ion[2]{%
  \mbox{#1\kern0.2em%
    \smaller\rmfamily%
    \edef\@tempa{\@car#2\@nil}%
    \ifcat1\@tempa%
    \@Roman{#2}%
    \else%
    \uppercase{#2}%
    \fi}}
\newcounter{IonStage}
\DeclareRobustCommand\plainion[2]{\setcounter{IonStage}{#2}#1 \Roman{IonStage}}
%%\def\ion#1#2{\ensuremath{\mathrm{#1}\;}{\protect\small\rm\@Roman{#2}}}
%    \end{macrocode}
% \end{macro}
%    \begin{macrocode}
\DeclareRobustCommand\nodata{$\cdots$}
%    \end{macrocode}
% These are needed by the |\ion| command. Isn't there an easier way of
% doing this?
%    \begin{macrocode}
\let\@small=\small
\let\@normalsize=\normalsize
\let\@footnotesize=\footnotesize
\let\@scriptsize=\scriptsize
\let\@tiny=\tiny
\let\@large=\large
\let\@Large=\Large
\let\@LARGE=\LARGE
\let\@huge=\huge
\let\@Huge=\Huge
\newcommand{\smaller}{\small}
\renewcommand{\normalsize}{\renewcommand{\smaller}{\small}\@normalsize}
\renewcommand{\small}{\renewcommand{\smaller}{\footnotesize}\@small}
\renewcommand{\footnotesize}{\renewcommand{\smaller}{\scriptsize}\@footnotesize}
\renewcommand{\scriptsize}{\renewcommand{\smaller}{\tiny}\@scriptsize}
\renewcommand{\tiny}{\renewcommand{\smaller}{\tiny}\@tiny}
\renewcommand{\large}{\renewcommand{\smaller}{\normalsize}\@large}  
\renewcommand{\Large}{\renewcommand{\smaller}{\large}\@Large}  
\renewcommand{\LARGE}{\renewcommand{\smaller}{\Large}\@LARGE}  
\renewcommand{\huge}{\renewcommand{\smaller}{\LARGE}\@huge}  
\renewcommand{\Huge}{\renewcommand{\smaller}{\huge}\@Huge}  
%</package|adshtml-style>
%    \end{macrocode}


% \section{Other files}

% \subsection{The driver and make files}
% \label{sec:driver_make}
% These are put here for completeness, although they still need to be
% distributed separately since otherwise the user won't be able to
% unpack this file into its constituent parts.
% 
% The |Makefile| has two principal rules: |docs|, and |macros|. The
% |docs| rule produces a PDF and gzipped PS file of the documentation.
% I still haven't perfected the combination of rules required to get
% it to run \LaTeX{} enough times in all circumstances. The |macros|
% rule runs \LaTeX{} on |rmaa.ins| in order to unpack |rmaa.cls| and
% all the other little files (including the |Makefile| and the
% |rmaa.ins| files themselves\dots 
% 
% First set up lists of various types of file. 
%    \begin{macrocode}
%<*makefile>
DOCFILES=rmaa.pdf
MACRO_LTX_FILES=rmaa.cls rmaacite.sty adshtml.sty
SHELLSCRIPTS=makeads.sh
MACRO_OTHER_FILES=adshtml.perl sample.latex2html-init
MAKEFILE=Makefile
MACROFILES=$(MACRO_LTX_FILES) $(SHELLSCRIPTS) $(MACRO_OTHER_FILES)         
SKELETONFILES=ADS.tex 
%    \end{macrocode}
% Install the files read only into the package directory, specified via a path relative to the source directory. \changes{4.0}{2019/11/14}{Reorganize folder layout so that hardly anything is in the top level.  That way, things are clearer on github.}
%    \begin{macrocode}
INSTALL=install 
INSTALL_DIR=../packages
%    \end{macrocode}
% The default rule is to make everything. Installation must be
% explicitly asked for.
%    \begin{macrocode}
all: docs macros

install: install-docs install-macros

%    \end{macrocode}
% Most files are installed read-only to prevent editing. Skeleton
% files must be writeable. Must be careful not to overwrite the
% Makefile in the parent directory. 
%    \begin{macrocode}
install-macros: macros
        $(INSTALL) -m 444 $(MACRO_LTX_FILES) $(MACRO_OTHER_FILES) $(INSTALL_DIR)
        $(INSTALL) -m 555 $(SHELLSCRIPTS) $(INSTALL_DIR)
        $(INSTALL) -m 644 $(SKELETONFILES) $(INSTALL_DIR)

install-docs: docs
        $(INSTALL) rmaa.pdf $(INSTALL_DIR)/rm-fulldocs.pdf

%    \end{macrocode}
% Make the PDF and gzipped PS files. 
%    \begin{macrocode}
docs: $(DOCFILES)

%    \end{macrocode}
% Make everything else if either the |.dtx| or |.ins| files have
% changed. 
%    \begin{macrocode}
macros: rmaa.dtx rmaa.ins
%    \end{macrocode}
% Ensure that all the macro files can be written to then recreate
% them. The |touch| is there in case any of the files don't exist
% yet. 
%    \begin{macrocode}
        touch $(MACROFILES) $(MAKEFILE)
        chmod u+w $(MACROFILES) $(MAKEFILE) rmaa.ins
        latex "\input{rmaa.ins}"
%    \end{macrocode}
% \TeX{} automatically adds the |.tex| extension to extensionless
% files. We fix this and take the opportunity to convert
% back to TABs all that it has lovingly converted to spaces. 
%    \begin{macrocode}
        cp -f $(MAKEFILE) $(MAKEFILE).BCK && unexpand $(MAKEFILE).tex >| $(MAKEFILE) && rm $(MAKEFILE).tex
%    \end{macrocode}
% Make the shell scripts executable and everything unwriteable (so
% they don't get hand-edited by mistake). However, undo this for the
% files that the user is \emph{supposed} to modify. 
%    \begin{macrocode}
        chmod a+x $(SHELLSCRIPTS)
        chmod a-w $(MACROFILES) $(MAKEFILE)
        chmod u+w $(SKELETONFILES)

%    \end{macrocode}
% We now use PDF\LaTeX{} for this.
%    \begin{macrocode}
rmaa.idx: rmaa.dtx
        @echo "REGENERATING .idx FILE"
        pdflatex "\nonstopmode\input{rmaa.dtx}"

rmaa.aux: rmaa.dtx rmaa.ind rmaa.gls
        @echo "REGENERATING .aux FILE"
        pdflatex "\nonstopmode\input{rmaa.dtx}"

rmaa.ind: rmaa.idx
        @echo "REGENERATING INDEX"
        makeindex -s gind.ist rmaa

rmaa.gls: rmaa.glo
        @echo "REGENERATING CHANGELOG"
        makeindex -s gglo.ist -o rmaa.gls rmaa.glo

rmaa.pdf: rmaa.aux rmaa.ind rmaa.gls
        @echo "REGENERATING .pdf FILE"
        pdflatex "\nonstopmode\input{rmaa.dtx}"

clean: 
        rm -f rmaa.aux rmaa.dvi rmaa.ind rmaa.gls
%</makefile>
%    \end{macrocode}

% The |rmaa.ins| file uses the \textsf{docstrip} package to split out
% all the files and strip off this documentation.
%    \begin{macrocode}
%<*ins>
%%
%% Installation batch file for the RMAA LaTeX package
%%
\def\batchfile{rmaa.ins}
\input docstrip.tex
\preamble

PLEASE DO NOT EDIT THIS FILE! It is generated automatically from
the file rmaa.dtx using the accompanying Makefile and rmaa.ins 
Any hand-edited changes will be lost

A LaTeX document class for the RMxAA

Author: William Henney (w.henney@irya.unam.mx)

Latest version available from: http://www.irya.unam.mx/rmaa

Copyright (C) 1999-2019 Instituto de Astronomia, UNAM, Mexico

Permission is hereby granted, free of charge, to any person obtaining
a copy of this software and associated documentation files (the
"Software"), to deal in the Software without restriction, including
without limitation the rights to use, copy, modify, merge, publish,
distribute, sublicense, and/or sell copies of the Software, and to
permit persons to whom the Software is furnished to do so, subject to
the following conditions:

The above copyright notice and this permission notice shall be
included in all copies or substantial portions of the Software.

THE SOFTWARE IS PROVIDED "AS IS", WITHOUT WARRANTY OF ANY KIND,
EXPRESS OR IMPLIED, INCLUDING BUT NOT LIMITED TO THE WARRANTIES OF
MERCHANTABILITY, FITNESS FOR A PARTICULAR PURPOSE AND
NONINFRINGEMENT. IN NO EVENT SHALL THE AUTHORS OR COPYRIGHT HOLDERS BE
LIABLE FOR ANY CLAIM, DAMAGES OR OTHER LIABILITY, WHETHER IN AN ACTION
OF CONTRACT, TORT OR OTHERWISE, ARISING FROM, OUT OF OR IN CONNECTION
WITH THE SOFTWARE OR THE USE OR OTHER DEALINGS IN THE SOFTWARE.
\endpreamble
%    \end{macrocode}
% Clobber files without asking
%    \begin{macrocode}
\askforoverwritefalse
%    \end{macrocode}
% Make the \LaTeX{} macro files (class and styles)
%    \begin{macrocode}
\generate{% 
  \file{rmaa.cls}{\from{rmaa.dtx}{package}}
  \file{rmaa.ins}{\from{rmaa.dtx}{ins}}
  \file{rmaacite.sty}{\from{rmaa.dtx}{rmaacite}}
  \file{adshtml.sty}{\from{rmaa.dtx}{adshtml-style}}
  \file{ADS.tex}{\from{rmaa.dtx}{ads-container}}
  }
%    \end{macrocode}
% Make the other non-\LaTeX{} files. These are done separately since
% they can't use `\%\%' as a comment character for the
% preamble/postamble (and Perl files can't have the postamble I don't
% think). 
%    \begin{macrocode}
\generate{\nopreamble\nopostamble% 
  \file{Makefile}{\from{rmaa.dtx}{makefile}}
  \file{adshtml.perl}{\from{rmaa.dtx}{adshtml-perl}}
  \file{makeads.sh}{\from{rmaa.dtx}{ads-shell-script}}
  \file{sample.latex2html-init}{\from{rmaa.dtx}{latex2html-init}}
  }

%% Local Variables: 
%% mode: LaTeX
%% mode: Font-lock
%% End:
%</ins>
%    \end{macrocode}

% \subsection{Extended citation syntax}
% Modified from a file called mnrascite.sty
% I no longer remember where I stole the original version of this from
% 
% Alterations to the cite commands to make them do more sensible things.
% N.B. the bibliography style should produce entries which have labels where
% the date is delimited by angle brackets '<>' for these macros to work
% properly. They will not give very sensible error messages if you do not do 
% this properly.
% 
% A new command |\scite{key}| is defined which produces a citation which is 
% suitable for insertion in a sentence. i.e., it does not put brackets around 
% the whole citation, only the date. Also |\pcite{key}| produce output with no
% parentheses whatsoever. 
% 
% Should redefine |\@bibitem| really to have any hope of compatibility
% the command |\lbibcite| is defined to deal with the citations.
%    \begin{macrocode}
%<*rmaacite>
\def\bibitem[#1<#2>]#3{\par\hangindent=1.5em\hangafter=1
  \if@filesw
  {\def\protect##1{\string ##1\space}
    \immediate\write\@auxout{\string\lbibcite{#3}{#1}{#2}}}
  \fi}
\@redefinebibitem % what it says
%    \end{macrocode}
% Note that the |\lbibcite| defines two types of citation one contains the names,
% the other the year
%    \begin{macrocode}
\def\lbibcite#1#2#3{\global\@namedef{b@#1}{#2}
                    \global\@namedef{by@#1}{#3}}

\def\@bcite#1#2{({#1\if@tempswa , #2\fi})}
\def\@pcite#1#2{#1\if@tempswa , #2\fi}
%    \end{macrocode}
% This declaration contains what is put between the name and the date
% \changes{3.11}{2003/01/31}{Moved the \cs{articlename} to after the key for consistency
% with the new natbib stuff.}
%    \begin{macrocode}
\def\@brcite{\ (}
\def\@ebrcite{)}
\def\@fmttop#1{\@fmtcite[\csname #1\endcsname]}
\def\@citex[#1]#2{%
%%  \if@filesw\immediate\write\@auxout{\string\citation{#2}}\fi
%% WJH added 16 Jan 2002
  \@for\@citeb:=#2\do{%
    \edef\@citeb{\expandafter\@firstofone\@citeb:\articlename}%
    \if@filesw\immediate\write\@auxout{\string\citation{\@citeb}}\fi}%
%% to here
  \def\@citea{}%
  \@cite{\@for\@citeb:=#2\do
    {%
      \edef\@citeb{\expandafter\@firstofone\@citeb:\articlename}%
      \@citea\def\@citea{;\ }\@ifundefined
      {b@\@citeb}{{\bf ?}\@warning {Citation `\@citeb' on page \thepage \space undefined}}%
      {\csname b@\@citeb\endcsname\@citebro\csname by@\@citeb\endcsname\@citebrc}}}%
  {#1}}

\def\cite{\@ifnextchar [{\let\@citebro=\space
                          \let\@citebrc=\relax
                          \let\@cite=\@bcite\@tempswatrue\@citex}
                        {\let\@citebro=\space
                          \let\@citebrc=\relax
                          \let\@cite=\@bcite\@tempswafalse\@citex[]}}
\def\pcite{\@ifnextchar [{\let\@citebro=\space
                          \let\@citebrc=\relax
                          \let\@cite=\@pcite\@tempswatrue\@citex}
                        {\let\@citebro=\space
                          \let\@citebrc=\relax
                          \let\@cite=\@pcite\@tempswafalse\@citex[]}}
\def\scite{\@ifnextchar [{\let\@citebro=\@brcite
                          \let\@citebrc=\@ebrcite
                          \let\@cite=\@pcite\@tempswatrue\@citex}
                        {\let\@citebro=\@brcite
                          \let\@citebrc=\@ebrcite
                          \let\@cite=\@pcite\@tempswafalse\@citex[]}}

\renewenvironment{thebibliography}
  {\section*{REFERENCES}\vspace*{\afterbiblioskip}\small
    \if@book\RereadBiblio\fi
    \sloppy\clubpenalty4000\widowpenalty4000
    \sfcode`\.=1000\relax
    \setlength\parindent{\z@}\setlength\parskip{\z@}}
  {\relax\normalsize\@outputaddressesifnecessary} 
%</rmaacite>
%    \end{macrocode}

% \subsection{Auxilliary files for generating ADS data}

% The ADS data are written by \LaTeX{} to a |.ads| file. This file
% then has to be included in the following file (|ADS.tex|),
% which is run through \LaTeX{}2HTML (see below). 

%    \begin{macrocode}
%<*ads-container>
%%
%% Instructions:
%% Edit the \input{} line below to reflect the name of your file. 
\documentclass{article}
\usepackage[latin1]{inputenc}
\usepackage{adshtml}
\newcommand{\ADSfield}[2]
{\%#1  #2\\}
%%\newcommand{\ADSvfield}[2]
%%{\%#1 \begin{verbatim}#2\end{verbatim}}
\newenvironment{ADSarticle}{}{\par\bigskip}
\renewcommand{\bibitem}[2][]{\par}
\setlength{\parindent}{0pt}
\renewcommand{\ThisVolume}[1]{%  
  2001, \RMAASC, 12, \oldpageref{#1:RMAAFirstPage}}  
 

\begin{document}
\makeatletter
\input{rmbooksamp.ads}
\end{document}
%</ads-container>
%    \end{macrocode}

% This is the |adshtml.sty| style file used by |ADS.tex|. It just
% defines a few 
% macros that might be used in abstracts, etc (really it should define
% them all). This now gets all its code from
% Sec~\ref{sec:user-commands}. 
%    \begin{macrocode}
%<*adshtml-style>
%% Nothing else needs to go here now. 
%</adshtml-style>
%    \end{macrocode}
% The init file for latex2html.
%    \begin{macrocode}
%<*latex2html-init>
$PDFTOCAIRO = "pdftocairo";
$PDFLATEX = "pdflatex";
$LATEX = "latex";
$MAX_SPLIT_DEPTH = 0;
$HTML_VERSION = "4.0,math";
$NO_SIMPLE_MATH = 1;
$ADDRESS = "<p><address><a href=\"mailto:rmaa\@astro.unam.mx\">rmaa\@astro.unam.mx</address></p>"; 
$INFO = "<p>This document was created using latex2html and the rmaa latex macros</p>"; 
$NO_NAVIGATION = 1;
1; 
%</latex2html-init>
%    \end{macrocode}
% This is the |adshtml.perl| file that is used by \LaTeX{}2HTML. Some
% bizarre \TeX{}y command sometimes show up so we tell it to ignore
% these. All the astronomy macros that come from AAS are processed in
% \TeX{}. That way they get made to images but with an |ALT| tag that
% has the original code, which is what will be picked up by |lynx|,
% with the result that they get passed intact into the final |.txt|
% file. 
%    \begin{macrocode}
%<*adshtml-perl>
package main;

$dum = &do_cmd_makeatletter();

sub do_cmd_RMAA {
    local($_) = @_;
    "RMxAA" . $_;
}

sub do_cmd_RMAASC {
    local($_) = @_;
    "RMxAASC" . $_;
}

sub do_cmd_etal {
    local($_) = @_;
    "et al\." . $_;
}

sub do_cmd_ibidrule {
    local($_) = @_;
    "<img src=\"foo\" alt=\"\ibidrule\">" . $_;
}


sub do_cmd_bibitem {
    local($_) = @_;
    local($key);
    local($opt, $dummy)=&get_next_optional_argument; 
    $key = &missing_braces unless
        ((s/$next_pair_pr_rx/$key = $2; ''/eo)
        ||(s/$next_pair_rx/$key = $2; ''/eo));
    "<p>" . $_; 
}

&ignore_commands (<<_RAW_ARG_CMDS_);
lowercase
adjustfinalcols
textsuperscript # {}
newline # []
thinmuskip
mskip
smallskip
medskip
bigskip
hspace # {}
hspacestar # {}
vspace # {}
vspacestar # {}
mathrm
_at_mathrm
spacefactor
discretionary
unhbox 
voidb_at_x 
penalty 
_at_M
_at_m
null
IeC
_RAW_ARG_CMDS_

&process_commands_in_tex (<<_RAW_ARG_CMDS_);
sun
earth
lesssim
gtrsim
sq
argdeg
arcmin
arcsec
fd
fh
fm
fs
fdg
farcm
farcs
fp
micron
la
ga
case # {} # {}
slantfrac # {} # {}
onehalf
onethird
twothirds
onequarter
threequarters
ubvr
ub
bv
vr
ur
aj
araa
apj
apjl
apjs
ao
apss
aap
aapr
aaps
azh
baas
jrasc
memras
mnras
pra
prb
prc
prd
pre
prl
pasp
pasj
qjras
skytel
solphys
sovast
ssr
zap
nat
iaucirc
aplett
apspr
bain
fcp
gca
grl
jcp
jgr
jqsrt
memsai
nphysa
physrep
physscr
planss
procspie
_RAW_ARG_CMDS_

1;                              # This must be the last line
%</adshtml-perl> 
%    \end{macrocode}

% This is a shell script |makeads.sh| that generates the |_ads.txt| file from
% |_ads.tex| by using \LaTeX{}2HTML and |lynx|.  \changes{4.0}{2019/11/13}{Make sure that lynx writes utf-8 output instead of latin-1.}
% \changes{2.2beta13}{2002/03/12}{Complete rewrite of |makeads.sh|. Now much 
% improved I think}
% \changes{4.0}{2019/11/13}{Fix the at-sign substitution so that it only eliminates \cs{@} and does not mess with email addresses.}

%    \begin{macrocode}
%<*ads-shell-script>
#! /bin/sh
if [ -z $1 ]; then
    echo "Usage: $0 basename"
    exit
fi
BASENAME=$1
ADSBASENAME=${BASENAME}_ads
AUXFILE=$BASENAME.aux
TMPFILE=tmp.$$


echo "Generating $ADSBASENAME.tex"

# For some reason, must double-up \\ to \\\\ but not \ to \\ - odd!
cat > $ADSBASENAME.tex <<EOF
%% PLEASE DO NOT EDIT THIS FILE!  It is generated automatically from
%% the script makeads.sh.  Any hand-edited changes will be lost
%% 
\documentclass{article}
\usepackage[latin1]{inputenc}
\usepackage{html}
\usepackage{adshtml}
\newcommand{\ADSfield}[2]{\%#1  #2\\\\}
\newenvironment{ADSarticle}{}{\par\bigskip}
\setlength{\parindent}{0pt}
\newcounter{istage}
\renewcommand{\ion}[2]{\setcounter{istage}{#2}#1 \Roman{istage}}
\begin{document}
\makeatletter
EOF
echo >> $ADSBASENAME.tex "\input{$BASENAME.ads2}" 
echo >> $ADSBASENAME.tex "\end{document}"



echo "Generating $ADSBASENAME.sed" 

# Find type of journal
grep '\\documentclass\[.*proceedings.*\]' $BASENAME.tex >/dev/null \
     && JOURNAL='RMxAC' \
     || JOURNAL='RMxAA'

# Find the year and volume number
YEAR=`grep '\\SetYear' $BASENAME.tex | sed 's/.*\\SetYear{\(.*\)}.*/\1/'`
VOLUME=`grep '\\SetVolume' $BASENAME.tex | sed 's/.*\\SetVolume{\(.*\)}.*/\1/'`

# Generate an auxilliary sed script 
cat > $TMPFILE <<EOF
s|\\\\newlabel{\(.*\):RMAAFirstPage}{{.*}{\(.*\)}}|s/\\\\\\\\ThisVolume{\1}/$YEAR, RMxAC, $VOLUME, \2/g|g
EOF

# Now go through the .aux file, looking for \newlabels and use the sed 
# script just generated in order to make another sed script, one that 
# will replace \ThisVolume{LABEL} with the correct reference
grep '^\\newlabel{.*:RMAAFirstPage}' $AUXFILE \
     | sed -f $TMPFILE > $ADSBASENAME.sed && rm $TMPFILE

# Apend a further sed command to get rid of \@, since redefining 
# it in _ads.tex didn't work
echo >> $ADSBASENAME.sed 's/\\\\@//g'

echo "Generating $BASENAME.ads2 from $ADSBASENAME.sed and $BASENAME.ads" 

# Apply this to the .ads file, generating the .ads2 file, which is included
# by the _ads.tex file.
sed -f $ADSBASENAME.sed $BASENAME.ads > $BASENAME.ads2

echo "Generating $ADSBASENAME.txt from $ADSBASENAME.tex and $BASENAME.ads2"

# Now run latex2html on the _ads.tex file
latex2html -no_navigation -info 0 -address 0 -html_version 4.0,math $ADSBASENAME

# Then, convert to utf8 with lynx and tidy it up a bit
lynx --justify=off --width=100 --dump --nolist --display_charset=utf-8 \
     $ADSBASENAME/$ADSBASENAME.html | cut -c4- \
     | sed 's/ibidrule/\\ibidrule/g' > $ADSBASENAME.txt
%</ads-shell-script>
%    \end{macrocode}

% \section{Epilogue}
% Err, that's it \dots

\endinput


% \endinput
% Local Variables:
% mode: doctex
% TeX-master: t
% End:
"

rmaa.ind: rmaa.idx
        @echo "REGENERATING INDEX"
        makeindex -s gind.ist rmaa

rmaa.gls: rmaa.glo
        @echo "REGENERATING CHANGELOG"
        makeindex -s gglo.ist -o rmaa.gls rmaa.glo

rmaa.pdf: rmaa.aux rmaa.ind rmaa.gls
        @echo "REGENERATING .pdf FILE"
        pdflatex "\nonstopmode% \def\fileversion{4.0} \def\filedate{2019/11/12}
% \MakeShortVerb{\|}
% \title{A \LaTeX{} document class for the Revista Mexicana de
% Astronom\'\i a y Astrof\'\i sica\thanks{This file has version
% \fileversion{} last revised \filedate{}.}}
% \author{William J. Henney\thanks{w.henney@irya.unam.mx}}
% \date{\filedate}
% \maketitle
% \paragraph{Abstract}
% This document describes the \texttt{rmaa} \LaTeX{}
% document class for the journal ``Revista Mexicana de
% Astronom\'\i a y Astrof\'\i sica''. The macros can be used for 
% \begin{enumerate}
% \item Preparing author manuscripts for the journal (including
% conference proceedings to be published in the ``Serie de
% Conferencias'')
% \item Typesetting an entire volume of the journal by inputting the
% individual article files from a container file. Facilities are
% included for automatic generation of TOC and Author Index. 
% \end{enumerate}
% 
% \tableofcontents
% \section{The documentation driver file}
% This is the part that generates the documentation that you may or
% may not be reading, which is done with the help of the Makefile
% discussed in section~\ref{sec:driver_make}. 
%    \begin{macrocode}
%<*driver>
\documentclass{ltxdoc}
\usepackage{hyperref}
\usepackage{palatino}
\EnableCrossrefs \CodelineIndex \RecordChanges
\setlength{\parindent}{0pt}
\def\MacroFont{\footnotesize\tt}
\begin{document} \DocInput{rmaa.dtx} \PrintIndex \PrintChanges \end{document}
%</driver>
%    \end{macrocode}

% \section{The implementation}
% 
% \changes{3.24}{2007/09/14}{Added user access to the macro version
% and date via \cs{rmVersion} and \cs{rmDate}}
% 
%    \begin{macrocode}
%<*package>
\newcommand{\rmaa@version}{3.26}
\newcommand{\rmaa@date}{2011/08/15}
\newcommand\rmVersion{\rmaa@version}
\newcommand\rmDate{\rmaa@date}
\NeedsTeXFormat{LaTeX2e}[1995/06/01]
%    \end{macrocode}
% \changes{4.0}{2019/11/02}{Change abstract to appear before
% resumen. Change RevMexAA to RMxAA for consistency with ADS.}
%
% \changes{3.22}{2007/03/22}{Allow proceedings articles to have
%   received/accepted dates} \changes{3.21}{2006/05/15}{Changes for
%   Silvia - abstracts are now shipped out to ADS.}
% \changes{3.19}{2006/01/19}{Changes for Silvia - remove adresses from
%   end, new abstract style} \changes{2.2beta}{2001/10/30}{ Fixed bug
%   whereby twocolumn option was not passed through to book
%   class. Finally sorted out the ``label on last page'' issue. }
% \changes{2.1}{2001/01/23}{ Fixed some inconsistencies in
%   \cs{nocite} and \cs{bibitem} uniqification Fixed the definition of
%   \cs{ion} so it will work in math mode and in contexts where the
%   typesize is different from normal (e.g., title, fig captions).
%   Modified the definition of \cs{@caption}, so we don't write any
%   list-of-figures data (we never use it anyway, and it just causes
%   problems with fragile commands in fig captions). }
% \changes{2.0}{2000/11/30}{ Added `newstyle' option, which should be
%   used in all new proceedings volumes (wider abstract, smaller
%   typeface in fig captions, etc.) Added commands for formatting
%   table captions. New demonstation document rmpostsamp.tex for
%   1-page posters. This will be the last release under a PRI
%   government (for a while...)}  \changes{1.9}{2000/06/06}{
%   Corrected position of page number in footer of articles that open
%   on even pages. Ensured that first article in each session opens on
%   an odd page, as do list of abstracts and author index. Added extra
%   space in TOC before the backmatter. Added command for editors,
%   \cs{BoldSessionTitles}, that makes produces bold titles for
%   session in the TOC.}  \changes{1.8}{2000/05/16}{ Fixed bug in
%   last page label with \cs{suppressfulladdresses} }
% \changes{1.7}{2000/05/11}{ Fixed more little things. Page labels
%   now work, even if you have no bibliography. You can output the
%   full addresses by hand with the \cs{outputfulladdresses}
%   command. Editors' Guide updated but still not finished. Revamped
%   splitbook.sh now produces PS file of each article for web page and
%   also an EPS file of each page for sending to the printers.}
% \changes{1.6}{2000/04/15}{ More cleaning up: full stops after
%   titles in index, running headers and footers in frontmatter,
%   fontsize in title of index, etc} \changes{1.5}{2000/04/15}{
%   Eliminated superscripts from titles in TOC and index}
% \changes{1.4}{2000/04/14}{ Added copyright and license info (GPL)}
% \changes{1.3}{2000/03/01}{ Corrected \cs{figcaption} to do the
%   caption formatting right Corrected font size of sec headings in
%   front matter} \changes{1.2}{2000/02/16}{ Added \cs{plate}
%   command} \changes{1.1}{2000/02/14}{ Added manuscript option,
%   corrected bugs} \changes{1.0}{2000/02/08}{ Now does a whole book
%   - lots of stuff tidied up} \changes{0.7}{1999/10/04}{ Increased
%   spacing between sections in single column styles. Suppressed
%   floats at bottom of pages (didn't work though) }
% \changes{0.6}{1999/09/24}{ Tweaked \cs{adjustfinalcols} (still not
%   perfect) Added proceedings option Added onecolumn/twocolumn
%   options (orthogonal to journal/proceedings) }
% \changes{0.5}{1999/09/23}{ Added \cs{fulladdresses} and
%   \cs{adjustfinalcols} (just about working)}
% \changes{0.4}{1999/09/22}{ Minor stylistic changes}
% \changes{0.3}{1999/06/28}{ Better mechanism for alternative
%   affiliations Added `guide' option for user guide}
% \changes{0.2}{1999/05/27}{ All basic functionality present, I
%   think} \changes{0.1}{1999/05/26}{ Just about working - still
%   missing Fig/Tab/Refs} \changes{0.0}{1999/05/21}{ So bored today,
%   decided to rewrite RMAA LaTeX macros from scratch}
%    \begin{macrocode}
\ProvidesClass{rmaa}[\rmaa@date\space v\rmaa@version\space RMxAA article
class]
%    \end{macrocode}
% \subsection{Option Handling}
% \label{sec:options}
% Declarations that are needed in options
% \changes{3.05}{2003/10/22}{Added natbib option}
% \changes{3.06}{2003/01/03}{Natbib option turned on permanently}
% \changes{3.07}{2003/01/11}{Added proof option}
% \changes{3.14}{2003/01/11}{Option to turn off natbib}
% \changes{3.18}{2003/10/22}{trying out doctex mode} 
% \changes{3.23}{2007/09/09}{Manuscript is now on by default.}
% \changes{3.26}{2011/08/15}{New option for review articles}
%    \begin{macrocode}
\newif\if@reviewarticle
\@reviewarticlefalse
\newif\if@proof
\@prooffalse
\newif\if@journal
\@journaltrue
\newif\if@printdates
\@printdatesfalse
\newif\if@proceedings
\@proceedingsfalse
\newif\if@manuscript
\@manuscripttrue
\newif\if@bigfont
\@bigfonttrue
\newif\if@guide
\@guidefalse
\newif\if@preprint
\@preprintfalse
\newif\if@rmtwocol
\@rmtwocoltrue
\newif\if@biblist
\@biblistfalse
\newif\if@book
\@bookfalse
\newif\if@natbib
\@natbibtrue
\newif\if@rmcite
\@rmcitefalse
\newif\if@titlevisible
\@titlevisibletrue
\newif\if@abstractpresent
\@abstractpresentfalse
\newif\if@printbanner
\@printbannerfalse
\newif\if@suppressfulladdresses
\@suppressfulladdressesfalse
\newif\if@newstyle
\@newstylefalse
\newif\if@debug
\@debugfalse
\newif\if@boldkeywords
\@boldkeywordsfalse
%    \end{macrocode}
% The option handling---mainly setting various booleans.
% \changes{2.2beta11}{2001/11/26}{Made twocol and
% newstyle the default for proceedings.}
% \changes{3.23}{2007/09/09}{Manuscript is now 10 pt instead of 12pt}
% \changes{3.23}{2007/09/10}{Full addresses are now automatically
% suppressed for proceedings style.}
% \changes{3.24}{2007/09/14}{New option \texttt{boldkeywords} that
% does the keywords in the old ugly way. Now off by default.}
%    \begin{macrocode}
\DeclareOption{manuscript}{\typeout{Manuscript for referee} 
  \@manuscripttrue\@journalfalse\@rmtwocolfalse
  \PassOptionsToClass{10pt}{article}}
\DeclareOption{bigfont}{\typeout{Big font manuscript requested by Gary
  Ferland} 
  \@manuscripttrue\@journalfalse\@rmtwocolfalse\@bigfonttrue
  \PassOptionsToClass{11pt}{article}}
\DeclareOption{journal}{\typeout{Main journal option}%
  \@journaltrue\@rmtwocoltrue\@printdatestrue\@manuscriptfalse
  \PassOptionsToClass{10pt}{article}}
\DeclareOption{guide}{\typeout{User guide option}%
  \@guidetrue\@journalfalse\@rmtwocoltrue\@manuscriptfalse
  \PassOptionsToClass{10pt}{article}}
\DeclareOption{preprint}{\typeout{Preprint option}%
  \@preprinttrue\@manuscriptfalse\PassOptionsToClass{10pt}{article}}
\DeclareOption{proof}{\typeout{Author proof option}%
  \@prooftrue\@manuscriptfalse}
\DeclareOption{proceedings}{\typeout{Proceedings option}%
  \@proceedingstrue\@journalfalse\@rmtwocoltrue\@newstyletrue
  \@suppressfulladdressestrue
  \@printdatesfalse\@manuscriptfalse\PassOptionsToClass{10pt}{article}} 
\DeclareOption{printdates}{\typeout{Force printing of received/accepted dates}%
  \@printdatestrue} 
\DeclareOption{book}{\typeout{Whole issue option}%
  \@booktrue} 
\DeclareOption{nonatbib}{\typeout{Turn off natbib support (required
    for compatibility with old rmaacite package).}%
  \@natbibfalse} 
\DeclareOption{newstyle}{\typeout{Use new layout style}%
  \@newstyletrue}
\DeclareOption{oldstyle}{\typeout{Use old layout style}%
  \@newstylefalse}
\DeclareOption{debug}{\typeout{Loads of annoying messages}%
  \@debugtrue}
%    \end{macrocode}
% These two need to be declared explicitly so that |\ExecuteOptions|
% actually does something.
%    \begin{macrocode}
\DeclareOption{twocolumn}{\@rmtwocoltrue}
\DeclareOption{onecolumn}{\@rmtwocolfalse}
\DeclareOption{twoside}{\PassOptionsToClass{\CurrentOption}{article}}
\DeclareOption{boldkeywords}{\@boldkeywordstrue}
%    \end{macrocode}
% Bold, upper-case keywords used to be the default. Now they must be
% turned on explicitly with the \texttt{boldkeywords} option. 
%    \begin{macrocode}
\DeclareOption{boldkeywords}{\@boldkeywordstrue}
%    \end{macrocode}
% Pass any other option along to the \texttt{article} class.
%    \begin{macrocode}
\DeclareOption*{\PassOptionsToClass{\CurrentOption}{article}}
%    \end{macrocode}
% Default options are to be two-sided and based on the standard
% \LaTeX{} \texttt{article} class,
% unless we are processing the entire book, in which case we base
% ourselves on the standard \texttt{book} class. Default is to use
% manuscript option unless otherwise specified. 
%    \begin{macrocode}
\ExecuteOptions{twoside,manuscript} \ProcessOptions\relax
\if@book
\if@rmtwocol\LoadClass[twocolumn]{book}\else\LoadClass{book}\fi
\newcommand{\abstract}{\relax}
\newcommand{\@maketitle}{\relax}
\else
\if@rmtwocol\LoadClass[twocolumn]{article}\else\LoadClass{article}\fi
\fi
%    \end{macrocode}
% Attempt to use natbib for citations 02 Jan 2003. 
% \changes{3.05}{2003/01/02}{If the natbib option is used, then set
% things up so that natbib's version of the bibliography is identical
% to the standard one.}
% \changes{3.09}{2003/01/23}{Explicitly require a recent version of
% natbib, so that we get a decent error message if the package is too
% old.} 
% \changes{3.23}{2007/09/09}{Now allow Bibtex to work properly} 
% \changes{3.25}{2009/10/01}{Change \cs{newcommand} tp \cs{def} when setting \cs{bibfont}. This allows it to work with all versions of natbib.} 
% \changes{3.25}{2010/02/06}{Updated copyright year.} 
% \changes{3.25}{2010/02/06}{Updated email and web addresses to \texttt{crya.unam.mx}.} 
%    \begin{macrocode}
\if@natbib
\RequirePackage{natbib}[1999/02/25]
\bibpunct{(}{)}{;}{a}{}{,}
\def\bibfont{\small}
\renewcommand\bibsection{\section*{REFERENCES}}
\setlength\bibsep{0pt}
\setlength\bibhang{1.5em}
\bibliographystyle{rmaa}
\fi
%    \end{macrocode}
% We want the \texttt{graphicx} package for the extended syntax of
% |\includegraphics|. We want the \texttt{ifthen} package for using in
% these macros (although we don't much yet). We explicitly load
% \texttt{babel} since some \LaTeX{} distributions messed up the
% hyphenation if we didn't.  
%    \begin{macrocode}
\RequirePackage{graphicx}
\RequirePackage{ifthen}
\IfFileExists{babel.sty}{%
%   \let\@restoreindent=\@afterindentfalse
  \RequirePackage[english]{babel}
%   \@restoreindent
  %% reinstate the definition of \@roman so as not to break .idx files
  \@ifundefined{org@@roman}{\relax}{\renewcommand{\@roman}[1]{\org@@roman{##1}}}
  }{%
  \typeout{Package babel.sty not found - don't worry about this}}
%    \end{macrocode}
% Try to load the \texttt{dcolumn} and \texttt{booktabs} packages if
% they can be found. These can produce nicer-looking tables. 
%    \begin{macrocode}
\IfFileExists{dcolumn.sty}{%
  \RequirePackage{dcolumn}
  }{%
  \typeout{Package dcolumn.sty not found - you will not be able to
    align columns on decimal points in tables}}
\IfFileExists{booktabs.sty}{%
  \RequirePackage{booktabs}
  }{%
  \typeout{Package booktabs.sty not found - your tables will not look
    as pretty as they ought too.}
  }
\IfFileExists{makeidx.sty}{%
  \RequirePackage{makeidx}
}{%
  \typeout{Package makeidx.sty not found - Author Index may not work.}
  }
%% commented out for moment since requires more modern version of LaTeX
%% \RequirePackage[latin1]{inputenc}
%    \end{macrocode}
% \changes{4.0}{2019/11/13}{Make \texttt{hyperref} and
% \texttt{etoolbox} into hard requirements, since we now use them
% extensively.}  We use hyperref to get internal and external links.
% But I don't like the colored boxes so I turn those off.  We use
% etoolbox for patching and robustifying at the moment.  In the
% future, I want to port a lot of other stuff over to it.  For
% instance, replace the ifthen package and the TeX ifs too.
%
% We target the latest version (1.7) of PDF, rather than the default
% of 1.5.  This reduces the chance of getting warnings from included
% graphics.  This is not strictly related to hyperref, but we put it
% here anyway.
%    \begin{macrocode}
\RequirePackage{etoolbox}
\RequirePackage{hyperref}
\hypersetup{hidelinks}
\pdfminorversion=7
%    \end{macrocode}
% We patch the \cs{@combinedblfloats} following a suggestion by Heiko
% Oberdiek at \url{https://tex.stackexchange.com/a/249743/3843}, which
% avoids problems when hyperlinks are split across
% pages. \changes{4.0}{2019/11/12}{Fix for issue where breaking
% hyperlink across pages gave an error.}
%    \begin{macrocode}
\urlstyle{same}
\newcommand\email{\begingroup \urlstyle{rm}\Url}
\patchcmd\@combinedblfloats{\box\@outputbox}{\unvbox\@outputbox}{}{%
  \errmessage{\noexpand\@combinedblfloats could not be patched}%
}%
%    \end{macrocode}
% \subsection{Other setup stuff}\label{sec:setup}
% Some declarations: 
%    \begin{macrocode}
\newif\if@needaddresses
\@needaddressestrue
\newif\if@frontmatter
\@frontmatterfalse
\newif\if@inpreamble
\@inpreambletrue
\newif\if@firstauthor
\@firstauthortrue
\newif\if@firstaffil
\@firstaffiltrue
\newif\if@firstaltaffil
\@firstaltaffiltrue
\newif\if@firstkwd
\@firstkwdtrue
\newif\if@noexpandADS
\@noexpandADSfalse
\newif\if@sessiontitlepages
\@sessiontitlepagesfalse
%    \end{macrocode}
% \begin{macro}{\@rmaaEmpty}
% \begin{macro}{\@rmaaArg}
% These are used in various commands to check if arg was left blank
%    \begin{macrocode}
\newcommand{\@rmaaEmpty}{}
\newcommand{\@rmaaArg}{}
%    \end{macrocode}  
% \end{macro}
% \end{macro}
%
% \subsection{Compatibility with old \LaTeX{} versions}
% \label{sec:compat}
% This package is based on the 1998/12/01 release of \LaTeX{}.
% Older releases may not define all the commands that we need.
% Luckily, we can use |\providecommand| to define these if necessary. 
% The following commands are not defined in the 1994/06/01 release. 
%
% \begin{macro}{\MakeUpperCase}
% \begin{macro}{\textsuperscript}
% \begin{macro}{\textbackslash}
% User commands:
%    \begin{macrocode}
\providecommand{\MakeUpperCase}[1]{\uppercase{#1}}
\providecommand{\textsuperscript}[1]{\ensuremath{^\textrm{#1}}}
\providecommand{\textbackslash}{\ensuremath{\backslash}}
%    \end{macrocode}  
% \end{macro}
% \end{macro}
% \end{macro}
% \begin{macro}{\protected@xdef}
% \begin{macro}{\restore@protect}
% \begin{macro}{\@unexpandable@protect}
% \begin{macro}{\hb@xt@}
% \LaTeX{} internal commands:
%    \begin{macrocode}
\providecommand{\protected@xdef}{%
        \let\@@protect\protect
        \let\protect\@unexpandable@protect
        \afterassignment\restore@protect
        \xdef}
\providecommand{\restore@protect}{\let\protect\@@protect}
\providecommand{\@unexpandable@protect}{\noexpand\protect\noexpand}
\providecommand{\hb@xt@}{\hbox to}
%    \end{macrocode}
% \end{macro}
% \end{macro}
% \end{macro}
% \end{macro}
%
% \subsection{Page layout}\label{sec:page-layout}
%
% The dimensions of the text were defined by Elizabeth. Text width is
% 169\,mm with equal margins on the two sides: 
%    \begin{macrocode}
\setlength{\textwidth}{169mm}
\setlength{\@tempdima}{\paperwidth}
\addtolength{\@tempdima}{-\textwidth}
\setlength{\oddsidemargin}{0.5\@tempdima}
\addtolength{\oddsidemargin}{-1in}
\setlength{\evensidemargin}{\oddsidemargin}
%    \end{macrocode}
% Text height is 225\,mm, top margin is 14\,mm. 
%    \begin{macrocode}
\setlength{\textheight}{225mm} 
\setlength{\topmargin}{-1in} 
\addtolength{\topmargin}{14mm}
%    \end{macrocode}
% Space between columns in two-column styles, header, footer,
% paragraph indentation:
%    \begin{macrocode}
\setlength{\columnsep}{9.5mm}
\setlength{\headsep}{0.7cm}
\setlength{\footskip}{0.75cm}
\setlength{\parindent}{1.5em}
%    \end{macrocode}
% 
% \begin{macro}{\@set@scalable@length}
% Try out a new scheme for dealing with various spacing dimensions and
% allowing the user to easily modify them. The command
% \begin{quote}
% |\set@scalable@length|\marg{name}\marg{defaultlength}\marg{stretch}\marg{shrink}
% \end{quote}
% does a 
% |\setlength|\marg{\cs{name}}\marg{defaultlength plus stretch minus shrink}
% and also defines the user command
% \textbackslash\meta{name}|Scale|\marg{factor} that can 
% be used to rescale length by \meta{factor}. Additionally, now
% defines the package command
% \textbackslash\meta{name}|@set@scale|\marg{FACTOR} 
% that also rescales the length by \meta{FACTOR}, so that the length
% becomes \meta{factor} $\times$ \meta{FACTOR} $\times$
% \meta{defaultlength}.   
%    \begin{macrocode}
\newcommand\@set@scalable@length[4]{%
  %% arg #1: NAME of length without leading `\'
  %%
  %% arg #2: default VALUE  
  \expandafter\newlength\csname @#1@default\endcsname
  \expandafter\setlength\csname @#1@default\endcsname{#2}%
  \expandafter\newlength\csname @@#1@@default\endcsname
  \expandafter\setlength\csname @@#1@@default\endcsname{#2}%
  %% arg #3: default STRETCH  
  \expandafter\newlength\csname @#1@stretch\endcsname
  \expandafter\setlength\csname @#1@stretch\endcsname{#3}%
  \expandafter\newlength\csname @@#1@@stretch\endcsname
  \expandafter\setlength\csname @@#1@@stretch\endcsname{#3}%
  %% arg #4: default SHRINK  
  \expandafter\newlength\csname @#1@shrink\endcsname
  \expandafter\setlength\csname @#1@shrink\endcsname{#4}%
  \expandafter\newlength\csname @@#1@@shrink\endcsname
  \expandafter\setlength\csname @@#1@@shrink\endcsname{#4}%
  %%
  %% User-defined scale \@NAME@scale is unity by default
  \expandafter\gdef\csname @#1@scale\endcsname{1}%
  %% Package-defined scale \@NAME@SCALE is unity by default
  \expandafter\gdef\csname @#1@SCALE\endcsname{1}%
  %%
  %% internal command \@reset@NAME should expand to 
  %% \setlength\@@NAME@@default{\@NAME@SCALE\@NAME@default}
  %% \setlength\NAME{\@NAME@scale\@@NAME@@default}
  %% in other words: scale x SCALE x default 
  \expandafter\gdef\csname @reset@#1\endcsname{%
    \expandafter\setlength\csname @@#1@@default\endcsname{%
      \csname @#1@SCALE\endcsname \csname @#1@default\endcsname
      }
    \expandafter\setlength\csname @@#1@@stretch\endcsname{%
      \csname @#1@SCALE\endcsname \csname @#1@stretch\endcsname
      }
    \expandafter\setlength\csname @@#1@@shrink\endcsname{%
      \csname @#1@SCALE\endcsname \csname @#1@shrink\endcsname
      }
    \expandafter\setlength\csname #1\endcsname{%
      \csname @#1@scale\endcsname \csname @@#1@@default\endcsname
      \@plus
      \csname @#1@scale\endcsname \csname @@#1@@stretch\endcsname
      \@minus
      \csname @#1@scale\endcsname \csname @@#1@@shrink\endcsname
      }}
  %%
  %% User command \NAMEScale should expand to 
  %% {\renewcommand\@NAME@scale{#1}\@reset@NAME}
  \expandafter\gdef\csname #1Scale\endcsname##1{%
    \expandafter\gdef\csname @#1@scale\endcsname{##1}%
    \csname @reset@#1\endcsname}
  %%
  %% Package command \NAME@set@SCALE should expand to 
  %% {\renewcommand\@NAME@SCALE{#1}\@reset@NAME}
  \expandafter\gdef\csname #1@set@SCALE\endcsname##1{%
    \expandafter\gdef\csname @#1@SCALE\endcsname{##1}%
    \csname @reset@#1\endcsname}
  %%
  %% initialize at default values with \@reset@NAME
  \csname @reset@#1\endcsname
  }
%    \end{macrocode}
% \end{macro}
% \begin{macro}{\parskip}
% First use of this new infrastructure: set interparagraph space to be
% zero but with a small amount of stretchiness.
%    \begin{macrocode}
%   \ifthenelse{\boolean{@manuscript}}{%
%     \@set@scalable@length{parskip}{0.5\baselineskip}{0.3\p@}{0.1\p@}
%   }{
    \@set@scalable@length{parskip}{0\p@}{0.3\p@}{0\p@}
%   }
%    \end{macrocode}
% \end{macro}
% 
% \begin{macro}{\abovesectionskip}
% \changes{4.0}{2019/11/12}{Use negative values for skip so that first paragraph after section headings is not indented.  This is something that I have been trying to fix for years!  It turns out that the anser was in the documentation for the standard command \cs{@startsection.}}
% \begin{macro}{\belowsectionskip}
% \begin{macro}{\aboveparaskip}
% \begin{macro}{\belowparaskip}
% Define some utility rubber lengths for use in
% section~\ref{sec:section-commands}. 
%    \begin{macrocode}
\newlength{\abovesectionskip}
\newlength{\belowsectionskip}
\newlength{\aboveparaskip}
\newlength{\belowparaskip}
%    \end{macrocode}
% The default values for these lengths are those traditionally used
% for the onecolumn style. 
%    \begin{macrocode}
\@set@scalable@length{abovesectionskip}{-3.5ex}{-4ex}{-2ex}
\@set@scalable@length{belowsectionskip}{2.3ex}{0ex}{2ex}
\@set@scalable@length{aboveparaskip}{3.25ex}{2ex}{0.5ex}
\@set@scalable@length{belowparaskip}{-1em}{\z@}{\z@}
%    \end{macrocode}
% \end{macro}\end{macro}\end{macro}\end{macro}
% We also make scalable some of the standard float parameters. These
% are the same defaults as in the standard article classes, but
% slightly more stretchy. For the twocolumn style, they are currently
% halved, see below. Note that I have also lumped |\abovecaptionskip|
% in with these, since it plays a similar role. There is also a
% |\belowcaptionskip| but it is set to zero in the standard classes. 
%    \begin{macrocode}
\@set@scalable@length{floatsep}{12\p@}{4\p@}{4\p@}
\@set@scalable@length{textfloatsep}{20\p@}{4\p@}{8\p@}
\@set@scalable@length{intextsep}{12\p@}{4\p@}{4\p@}
\@set@scalable@length{dblfloatsep}{12\p@}{4\p@}{4\p@}
\@set@scalable@length{dbltextfloatsep}{20\p@}{4\p@}{8\p@}
\@set@scalable@length{abovecaptionskip}{10\p@}{\z@}{\z@}
%    \end{macrocode}
% \begin{macro}{\@before@title@skip}         
% \begin{macro}{\@before@author@skip}        
% \begin{macro}{\@between@author@skip}       
% \begin{macro}{\@before@abstract@skip}      
% \begin{macro}{\@after@abstracthead@skip}   
% \begin{macro}{\@between@abstract@skip}     
% \begin{macro}{\@before@kwd@skip}           
% \begin{macro}{\@before@dates@skip}           
% Spacing parameters used in |\@maketitle|. These are not changed by
% default except by the user, so they don't need to go into
% |\@package@default@lengths|.  
%    \begin{macrocode}
\newlength{\@before@title@skip}         
\newlength{\@before@author@skip}        
\newlength{\@between@author@skip}       
\newlength{\@before@abstract@skip}      
\newlength{\@after@abstracthead@skip}   
\newlength{\@between@abstract@skip}     
\newlength{\@before@kwd@skip}           
\newlength{\@before@dates@skip}           
\@set@scalable@length{@before@title@skip}{2em}{\z@}{\z@}
\@set@scalable@length{@before@author@skip}{1.5em}{\z@}{\z@}
\@set@scalable@length{@between@author@skip}{0.5em}{\z@}{\z@}
\@set@scalable@length{@before@abstract@skip}{3ex}{1ex}{0.5ex}
\@set@scalable@length{@after@abstracthead@skip}{1ex}{\z@}{\z@}
\@set@scalable@length{@between@abstract@skip}{2ex}{1ex}{0.5ex}
\@set@scalable@length{@before@kwd@skip}{1ex}{0.5ex}{0.2ex}
\@set@scalable@length{@before@dates@skip}{2ex}{0.5ex}{0.2ex}
%    \end{macrocode}
% \end{macro}\end{macro}\end{macro}\end{macro}
% \end{macro}\end{macro}\end{macro}\end{macro}
% 
% Spacing above and below displayed equations also wants to be done
% like this, but it is more complicated since the standard article
% classes make these change with the fontsize. Something to do later,
% maybe. 
% 
%    \begin{macrocode}
%    \end{macrocode}
% \begin{macro}{\@package@default@lengths}
% The value actually used for these lengths is controlled by the
% \meta{SCALE} parameter, which depends on the style option that is in
% effect. This is set using the |\@package@default@lengths| command,
% which must be called every time the style changes (e.g., changing
% from onecolumn to twocolumn). 
%    \begin{macrocode}
\newcommand{\@package@default@lengths}{%
  \ifthenelse{\boolean{@rmtwocol}}{%
    %% twocolumn mode
    \abovesectionskip@set@SCALE{0.5}
    \belowsectionskip@set@SCALE{0.5}
    \aboveparaskip@set@SCALE{0.5}
    \belowparaskip@set@SCALE{1}
    \floatsep@set@SCALE{0.5}
    \textfloatsep@set@SCALE{0.5}
    \intextsep@set@SCALE{0.5}
    \dblfloatsep@set@SCALE{0.5}
    \dbltextfloatsep@set@SCALE{0.5}
    }{%
    %% onecolumn mode
    \abovesectionskip@set@SCALE{1}
    \belowsectionskip@set@SCALE{1}
    \aboveparaskip@set@SCALE{1}
    \belowparaskip@set@SCALE{1}
    \floatsep@set@SCALE{1}
    \textfloatsep@set@SCALE{1}
    \intextsep@set@SCALE{1}
    \dblfloatsep@set@SCALE{1}
    \dbltextfloatsep@set@SCALE{1}
    }%
  }
\@package@default@lengths
%    \end{macrocode}
% \end{macro}
% \begin{macro}{\RescaleLengths}
% \begin{macro}{\RescaleSecLengths}
% \begin{macro}{\RescaleFloatLengths}
% \begin{macro}{\RescaleTitleLengths}
% User commands that will fiddle with all the scalable lengths. 
%    \begin{macrocode}
\newcommand{\RescaleLengths}[1]{%
  \RescaleSecLengths{#1}\RescaleFloatLengths{#1}\RescaleTitleLengths{#1}}
\newcommand{\RescaleSecLengths}[1]{%
%% sections
  \abovesectionskipScale{#1}
  \belowsectionskipScale{#1}
  \aboveparaskipScale{#1}
  }
\newcommand{\RescaleFloatLengths}[1]{%
%% floats
  \floatsepScale{#1}
  \textfloatsepScale{#1}
  \intextsepScale{#1}
  \dblfloatsepScale{#1}
  \dbltextfloatsepScale{#1}
  \abovecaptionskipScale{#1}
  }
\newcommand{\RescaleTitleLengths}[1]{%
%% maketitle
  \@before@title@skipScale{#1}           
  \@before@author@skipScale{#1}        
  \@between@author@skipScale{#1}       
  \@before@abstract@skipScale{#1}      
  \@after@abstracthead@skipScale{#1}   
  \@between@abstract@skipScale{#1}     
  \@before@kwd@skipScale{#1}           
  \@before@dates@skipScale{#1}           
  }
%    \end{macrocode}
% \end{macro}\end{macro}\end{macro}\end{macro}
% \changes{2.2beta14}{2002/03/17}{Finished implementation of new
% infrastructure for scalable lengths. Converted the above/below
% section skips to use it. Also, all the floatsep parameters. Also,
% all the vertical lengths that are used by |\@maketitle|. }

% \begin{macro}{\MSextramargin}
% For the \texttt{manuscript} option we have an extra 0.3\,inch
% vertical margin and 1\,inch horizontal margin for marginal notes (by
% hand or with |\marginpar|. 
%
% \changes{3.23}{2007/09/09}{Reduce baselinestretch in
%   \texttt{manuscript} style to 1.2. No longer use raggedright for
%   manuscript.}
%
% \changes{3.24}{2007/09/14}{Further reduction of baselinestretch in
%   \texttt{manuscript} style to 1.1 since we now have wider margins.}
%
% \changes{3.24}{2007/09/14}{Split \texttt{manuscript}'s extra margin
%   into separate horizontal and vertical parts.}
% 
%    \begin{macrocode}
\newlength{\MSextraHmargin}\newlength{\MSextraVmargin}
\setlength{\MSextraHmargin}{1in} 
\setlength{\MSextraVmargin}{0.3in} 
\if@manuscript
   \renewcommand{\baselinestretch}{1.1}
   \addtolength{\textwidth}{-2\MSextraHmargin}
   \addtolength{\oddsidemargin}{\MSextraHmargin}
   \setlength{\evensidemargin}{\oddsidemargin}
   \addtolength{\textheight}{-2\MSextraVmargin}
   \addtolength{\topmargin}{\MSextraVmargin}
%    \raggedright
\fi
%    \end{macrocode}
% \end{macro}
% Float parameters: don't allow bottom floats by default. 
%    \begin{macrocode}
\renewcommand{\fps@figure}{tp}
\renewcommand{\fps@table}{tp}
%    \end{macrocode}
%
% \begin{macro}{\rmaa@AuthorWidthFrac}
% \begin{macro}{\rmaa@AbstractWidthFrac}
% These parameters control what fraction of the total text width can
% be used by the list of authors and the abstract/resumen,
% respectively.  
%    \begin{macrocode}
\newcommand{\rmaa@AuthorWidthFrac}{1.0}
\newcommand{\rmaa@AbstractWidthFrac}{0.75}
%    \end{macrocode}
% Use a wide abstract to save space in the new style.
%    \begin{macrocode}
\if@newstyle\renewcommand{\rmaa@AbstractWidthFrac}{1.0}\fi
%    \end{macrocode}
% \end{macro}
% \end{macro}
%
% Some things to consider in future versions. Trying to avoid too many
% widows, etc. Apparently, the following will work best with
% \cs{flushbottom}. 
%    \begin{macrocode}
%%\widowpenalty=10000
%%\clubpenalty=10000
%    \end{macrocode}
% Also, should perhaps try getting rid of underfull/overfull box warnings. 
% \changes{3.10}{2003/01/24}{Set large values for \cs{hbadness},
% \cs{vbadness}, \cs{hfuzz}, \cs{vfuzz}. To reduce amount of stuff
% that gets written to the terminal every time you compile something.}
%    \begin{macrocode}
\hbadness=10000
\vbadness=10000
\setlength\hfuzz{\maxdimen}
\setlength\vfuzz{\maxdimen}
%    \end{macrocode}
% \subsection{Various text styles for different contexts}
% \label{sec:text-styles} 
% \begin{macro}{\rmaa@AffilStyle}
% \begin{macro}{\rmaa@PagesStyle}
% \begin{macro}{\rmaa@YearStyle}
% \begin{macro}{\rmaa@VolumeStyle}
% \begin{macro}{\rmaa@VolumenStyle}
% \begin{macro}{\rmaa@JournalNameStyle}
% \begin{macro}{\rmaa@ReceivedStyle}
% \begin{macro}{\rmaa@AbstractHeaderStyle}
% \begin{macro}{\rmaa@JournalNameStyle}
% Style commands for the different fields used by |\maketitle| in
% constructing the title page of the article. These have the form
% |\rmaa@XYZStyle{TEXT}|. I'm not sure these were really a good idea,
% but never mind \dots 
%    \begin{macrocode}
\newcommand{\rmaa@AffilStyle}[1]{{\normalsize\rmfamily#1}}
\newcommand{\rmaa@PagesStyle}[1]{{\normalsize\rmfamily#1}}
\newcommand{\rmaa@YearStyle}[1]{{\normalsize\rmfamily#1}}
\newcommand{\rmaa@VolumeStyle}[1]{{\normalsize\rmfamily\bfseries#1}}
\newcommand{\rmaa@VolumenStyle}[1]{{\Large\sffamily\bfseries#1}}
\newcommand{\rmaa@JournalNameStyle}[1]{{\normalsize\itshape#1}}
\newcommand{\rmaa@ReceivedStyle}[1]{{\normalsize\rmfamily\itshape#1}}
\newcommand{\rmaa@AbstractHeaderStyle}[1]{%
  {\normalsize\rmfamily\MakeUpperCase{#1}}}
\newcommand{\rmaa@TitleStyle}[1]{%
  {\large\bfseries\rmfamily\boldmath\MakeUpperCase{#1}\par}}
%    \end{macrocode}
% \end{macro}
% \end{macro}
% \end{macro}
% \end{macro}
% \end{macro}
% \end{macro}
% \end{macro}
% \end{macro}
% \end{macro}
%
% \begin{macro}{\rmaa@FrontSecStyle}
% Headings of ``front sections'' (see section~\ref{sec:front-sections})
%    \begin{macrocode}
\newcommand{\rmaa@FrontSecStyle}[1]{%
  {\normalsize\bfseries\rmfamily\MakeUpperCase{#1}}}
%    \end{macrocode}
% \end{macro}
%
% \begin{macro}{\rmaa@TOCTitleStyle}
% \begin{macro}{\rmaa@TOCPartStyle}
% \begin{macro}{\rmaa@IndexTitleStyle}
% \begin{macro}{\rmaaTOCAuthorStyle}
% \begin{macro}{\rmaaIndexAuthorStyle}
% Commands to set style of TOC and index entries.
% Rewritten to avoid unnecessary font commands in the .toc and .ind files
%    \begin{macrocode}
\newcommand{\rmaa@TOCTitleStyle}[1]{\MakeUpperCase{#1}}
\newcommand{\rmaa@TOCPartStyle}[1]{\rmaa@TOCTitleStyle{#1}}
\newcommand{\rmaa@IndexTitleStyle}[1]{#1}
\newcommand{\rmaaTOCAuthorStyle}[1]{%
  {\itshape#1}}
\newcommand{\rmaaIndexAuthorStyle}[1]{%
  \protect\mbox{\bfseries#1}}
%    \end{macrocode}
% \end{macro}
% \end{macro}
% \end{macro}
% \end{macro}
% \end{macro}
%
% \begin{macro}{\rmaa@KeywordStyle}
% \begin{macro}{\rmaa@KeywordHeaderStyle}
% The style of the keywords depends on whether we are using the
% \texttt{newstyle} option or not. If we are, then they are small. 
%    \begin{macrocode}
\ifthenelse{\boolean{@newstyle}}{
  \newcommand{\rmaa@KeywordStyle}[1]{%
    {\small%
      \if@boldkeywords 
      \bfseries\MakeUpperCase{\raggedright #1}
      \else
      \raggedright #1
      \fi
    }}
  \newcommand{\rmaa@KeywordHeaderStyle}[1]{%
    {\small\itshape#1}}}
{
  \newcommand{\rmaa@KeywordStyle}[1]{%
    {\normalsize%
      \if@boldkeywords\bfseries\fi\raggedright #1}}
  \newcommand{\rmaa@KeywordHeaderStyle}[1]{%
    {\normalsize\itshape#1}}}
%    \end{macrocode}
% \end{macro}
% \end{macro}
%
% \subsection{Page styles}\label{sec:page-styles}
% \begin{macro}{\ps@rmaa}
% The running headers for normal pages:
%    \begin{macrocode}
\newcommand{\ps@rmaa}{%
%    \end{macrocode}
% nothing in either foot; 
%    \begin{macrocode}
  \renewcommand\@oddfoot{\relax}%
  \renewcommand\@evenfoot{\relax}%
%    \end{macrocode}
% page number and short author list in the even heads; 
%    \begin{macrocode}
  \renewcommand\@evenhead{\rlap{\thepage}\hfil\@shortauthor\hfil}%
%    \end{macrocode}
% short title and page number in the odd heads.
%    \begin{macrocode}
  \renewcommand\@oddhead{\hfil\@shorttitle\hfil\llap{\thepage}}
  }
%    \end{macrocode}
% \begin{macro}{\ps@rmaafirstpage}
% The header of the first page of each
% article:
%    \begin{macrocode}
\newcommand{\ps@rmaafirstpage}{%
  \renewcommand\@oddfoot{\hfil\thepage}%
%    \end{macrocode}
% footer corrected (06 Jun 2000) to have page on outside edge;
%    \begin{macrocode}
  \renewcommand\@evenfoot{\thepage\hfil}%
%    \end{macrocode}
% Header has journal name, volume, and page numbers for a regular
% article, or similar for special cases. 
% \changes{4.0}{2019/11/03}{Change to three-line header with copyright info and DOI.}
%    \begin{macrocode}
\if@preprint
\renewcommand\@oddhead{%
  \if@proceedings 
  \textsc{To appear in ``\@conftitle{} (\rmaa@YearStyle{\@year})''}\hfil
  \textit{RMxAC}
  \else
  \textsc{To appear in} \@journalname\hfil
  \fi}\fi%
\if@guide
\renewcommand\@oddhead{\textsc{Documentation} (\rmaa@YearStyle{\@year})\hfil}
\fi%
\if@manuscript
\renewcommand\@oddhead{Manuscript for \@journalname{} (\rmaa@YearStyle{\@year})\hfil}
\fi%
\if@journal
\renewcommand\@oddhead{%
  \parbox{\textwidth}{
    \@journalname, \rmaa@VolumeStyle{\@volume},
    \rmaa@PagesStyle{\@artpages}\
    (\rmaa@YearStyle{\@year})\hfil\\[1pt]
    {\small\textcopyright{} \@year:
      Instituto de Astronom\'\i{}a,
      Universidad Nacional Aut\'onoma de M\'exico}\hfil\\[1pt]
    {\small\@DOI@url}\hfil
  }%
}\fi%
  \renewcommand\@evenhead{\@oddhead}}
%    \end{macrocode}
% By default, use the |\ps@rmaa| page style.
%    \begin{macrocode}
\pagestyle{rmaa}
%    \end{macrocode}
% \end{macro}
% \end{macro}
%
% \subsection{Preamble commands}\label{sec:preamble} 
% Everything you need for setting the title page, TOC and index
% entries, etc.
% \subsubsection{Commands for the title page}
% \label{sec:title-commands}
% \begin{macro}{\resumen}
% \begin{macro}{\@resumen}
% \begin{macro}{\abstract}
% \begin{macro}{\@abstract}
% User commands for setting the resumen and abstract. 
% These should be used in the document header. The internal commands
% |\@resumen| and |\@abstract| are used to hold the actual text
% pending output via the |\maketitle| command. 
%    \begin{macrocode}
\newif\if@setresumen
\newcommand{\resumen}[1]{%
%%  \renewcommand{\@resumen}{{\selectlanguage{spanish}#1}}%
  \renewcommand{\@resumen}{{\ignorespaces#1}}%
  \@setresumentrue}
\renewcommand{\abstract}[1]{%
  \renewcommand{\@abstract}{{\ignorespaces#1}}
  \renewcommand\ADS@abstract{#1}
%% Not a good idea, since then author macros are not expanded and
%% cause errors later. However, by not doing it we are vulnerable to
%% fragile commands.....
  \if@noexpandADS\@onelevel@sanitize{\ADS@abstract}\fi
  \@abstractpresenttrue
  }
\newcommand{\ADSnoexpand}{\@noexpandADStrue}
%    \end{macrocode}
% \begin{macro}{\boldabstract}
% This is the version for doing 1-page conference proceedings. Still
% save the text for ADS but then write it out immediately as a bold
% paragraph. 
%    \begin{macrocode}
\newcommand{\boldabstract}[1]{%
  \gdef\ADS@abstract{#1}%
  \renewcommand{\@abstract}{#1}%
  \if@noexpandADS\@onelevel@sanitize{\ADS@abstract}\fi%
  \noindent\textbf{\boldmath #1}
  }
%    \end{macrocode}
% \end{macro}
% Default value in each case is to print out a warning if not set
%    \begin{macrocode}
  \newcommand{\@resumen}{%
    \@setresumenfalse \ClassWarningNoLine{rmaa}{No Spanish abstract
      (resumen) specified - if you do not speak Spanish, this will be
      translated by the editors}%
    \textbf{Favor de proporcionar un resumen en espa\~nol. If you are
      unable to translate your abstract into Spanish, the editors will
      do it for you.} }
  \newcommand{\@abstract}{\ClassWarningNoLine{rmaa}{No abstract
      specified}}
%    \end{macrocode}
% \end{macro}
% \end{macro}
% \end{macro}
% \end{macro}
%
% \begin{macro}{\title}
% \begin{macro}{\@rawtitle}
% \begin{macro}{\@TOCtitle}
% \begin{macro}{\rmaaIndexTitle}
% Similarly, |\title| is the preamble command to set the title of the
% article. This is used in various places, typeset in different
% styles, hence the use of the internal |\@rawtitle|, |\@title|,
% |\rmaaIndexTitle|, and |\@TOCtitle| commands, each of which uses the
% relevant style from section~\ref{sec:text-styles}
%    \begin{macrocode}
\newcommand{\@rawtitle}{Raw Title}
\newcommand{\@TOCtitle}{\rmaa@TOCTitleStyle{Title for TOC}}
\newcommand{\rmaaIndexTitle}{\rmaa@IndexTitleStyle{Title for index}.}
\newcommand{\@TOCaftertitle}{}
\renewcommand{\title}[1]{%
%%  {\ADS@sanitize\xdef\@rawtitle{#1}}
  \renewcommand\@rawtitle{\noexpand#1}
  \let\ADS@title=\@rawtitle
  \renewcommand{\@title}{\rmaa@TitleStyle{#1}}%
  \renewcommand{\@TOCtitle}{\rmaa@TOCTitleStyle{#1}\@TOCaftertitle}%
  \renewcommand{\rmaaIndexTitle}{\rmaa@IndexTitleStyle{#1\@TOCaftertitle}.}%
  }
%    \end{macrocode}
% \end{macro}
% \end{macro}
% \end{macro}
% \end{macro}
% 
% All the weird and whacky ways of indicating the author
% affiliations. 
% \begin{macro}{\affil}
% First, the way that sets it directly underneath the authors name.
% Note that |\affil| trumps all other affiliations for the purposes of
% getting into the ADS stuff (we use |\gdef|).   
% \dots 
%    \begin{macrocode}
\newcommand{\affil}[1]{%
  \\[1.5ex]\rmaa@AffilStyle{#1}\\[1.5ex]%
  \if@firstaffil\gdef\ADS@affil{#1}\@firstaffilfalse\fi%
  }
\renewcommand{\and}{and\\[1.5ex]}
%    \end{macrocode}
% \begin{macro}{\altaffil}
% Alternative affiliations are footnoted. 
%    \begin{macrocode}
\newcommand{\altaffil}[1]{%
  \if@firstaltaffil\gdef\ADS@affil{#1}\@firstaltaffilfalse\fi%
  \thanks{#1}}
%    \end{macrocode}
% \begin{macro}{\altaffilmark}
% \begin{macro}{\altaffiltext}
% Alternative-alternative mechanism for affiliations---should not be
% mixed with |\altaffil| (unless you know what you are doing) or there
% will be problems with numbering.
%    \begin{macrocode}
\newcommand{\altaffilmark}[1]{\textsuperscript{#1}}
\newcommand{\altaffiltext}[2]{%
  \if@firstaltaffil\gdef\ADS@affil{#2}\@firstaltaffilfalse\fi
  \stepcounter{footnote}%
  \protected@xdef\@thanks{\@thanks
    \protect\footnotetext[#1]{#2}}} 
%    \end{macrocode}
% \end{macro}
% \end{macro}
% \end{macro}
% \end{macro}
%
% \begin{macro}{\keywords}
% \begin{macro}{\@keywords}
% \begin{macro}{\@keywordheader}
% \label{keywords}Sundry minor preamble commands. All
% self-explanatory, I hope. |\keywords| is the old user interface
% (deprecated). New articles should use |\addkeyword| once for each
% keyword. That way the ADS record comes out right too (comma rather
% than em-dash separator).
% \changes{2.2beta12}{2001/12/07}{\cs{addkeyword} New user interface
%   for keywords.} 
% \changes{2.2beta13}{2002/03/11}{\cs{addkeyword} Improved behaviour
%   with .ads file. Use
%   \cs{ADS@sanitize} to change definitions of `~' and \cs{ion}.}
% \changes{3.23}{2007/09/09}{\cs{addkeyword} now signals a compilation
% error if the author tries to add more than 6 keywords.}
% \changes{3.24}{2007/09/14}{kwdcounter is now reset in
% \cs{@begofarticlehook} to prevent error in book style} 
%    \begin{macrocode}
\newcommand{\@kwdsep}{ --- }
\newcommand{\ADS@kwdsep}{, }
\newcounter{kwdcounter}\setcounter{kwdcounter}{0}
\newcommand{\addkeyword}[1]{%
  \addtocounter{kwdcounter}{1}
  \ifthenelse{\thekwdcounter > 6}{
    \ClassError{rmaa}{Too many keywords. A maximum of 6 keywords is allowed.}{}
  }{}
  \if@firstkwd%
  \renewcommand\@keywords{\if@boldkeywords\uppercase{#1}\else #1\fi}
  {\ADS@sanitize\def\ion{\plainion}\gdef\ADS@keywords{#1}}
  \@firstkwdfalse
  \else%
  \edef\@keywords{\@keywords\@kwdsep \if@boldkeywords\uppercase{#1}\else #1\fi}%
  {\ADS@sanitize\def\ion{\plainion}\xdef\ADS@keywords{\ADS@keywords\ADS@kwdsep #1}}%
  \fi}
\newcommand{\ADS@sanitize}{%
  \def\mathrm{}
  \def\@{\ }
  \def~{\ }
  \def\langle{<}
  \def\rangle{>}
  \def\log{log }
  \def\sin{sin }
  \def\cos{cos }
  \def\tan{tan }
  \def\fd{^d.}
  \def\fh{^h.}
  \def\fm{^m.}
  \def\fs{^s.}
  \def\fdg{�.}
  \def\farcm{'.}
  \def\farcs{''.}
  \def\arcdeg{�}
  \def\,{ }
  }
\newcommand{\keywords}[1]{%
  \renewcommand{\ADS@keywords}{#1}
  \renewcommand{\@keywords}{\uppercase{#1}}
  }
\newcommand{\@keywordheader}{\rmaa@KeywordHeaderStyle{Key Words:}}
\newcommand{\@keywords}{%
  \ClassWarningNoLine{rmaa}{No keywords specified}%
  Keywords go here
  }
%    \end{macrocode}
% \begin{macro}{\shorttitle}
% \begin{macro}{\shortauthor}
% \begin{macro}{\@shorttitle}
% \begin{macro}{\@shortauthor}
% Commands for setting the running headers.
%    \begin{macrocode}
\newcommand{\shorttitle}[1]{%
  \renewcommand\@shorttitle{\MakeUpperCase{#1}}}
\newcommand{\shortauthor}[1]{%
  \renewcommand\@shortauthor{\MakeUpperCase{#1}}}
\newcommand{\@shorttitle}{SHORT TITLE GOES HERE}
\newcommand{\@shortauthor}{SHORT AUTHOR GOES HERE}
%    \end{macrocode}
% \end{macro}
% \end{macro}
% \end{macro}
% \end{macro}
% \end{macro}
% \end{macro}
% \end{macro}


% \subsubsection{The full addresses}
% \label{sec:full-addresses}
% \begin{macro}{\fulladdresses}
% \begin{macro}{\@fulladdresses}
% \begin{macro}{\@fulladdheight}
% \begin{macro}{\@fulladddepth}
% \begin{macro}{\@fulladdbox}
% Aaghh, the full addresses! A perennial source of problems. 
%    \begin{macrocode}
\newlength{\@fulladdheight}
\newlength{\@fulladddepth}
\newsavebox{\@fulladdbox}
%% \fulladdresses command modified (11 May 2000) so that we can put
%% \fulladdresses{} in the article to suppress the output of the
%% addresses
\newcommand{\@FinalFigure}{\relax}
\newcommand{\FinalFigure}[1]{\renewcommand{\@FinalFigure}{#1}}
\newcommand{\fulladdresses}[1]{%
  \renewcommand{\@rmaaArg}{#1}
  \ifx\@rmaaArg\@rmaaEmpty
     \renewcommand{\@fulladdresses}{}
  \else
     \renewcommand{\@fulladdresses}{
       \normalsize%
       \ifx\@FinalFigure\relax\else\vfill\@FinalFigure\vfill\fi
%%       \vspace*{1ex}
       %% do addresses as a list (copied from thebibliography environment) 
       \list{\null}{\leftmargin 1.5em\topsep\z@\labelwidth\z@\labelsep\z@\itemindent -1.5em
         \itemsep\z@\parsep\z@\usecounter{enumi}}
       \sloppy\clubpenalty4000\widowpenalty4000
       \sfcode`\.=1000\relax
       #1
       \endlist
       }
  \fi
  }
\newcommand{\@fulladdresses}{Full addresses go here}
%    \end{macrocode}
% \begin{macro}{\suppressfulladdresses}
% New user-level command for suppressing full addresses should you so
% desire. Currently used in the single-page-poster style. 
%    \begin{macrocode}
\newcommand{\suppressfulladdresses}{\@suppressfulladdressestrue\@needaddressesfalse}
%    \end{macrocode}
% \end{macro}
% \end{macro}
% \end{macro}
% \end{macro}
% \end{macro}
% \end{macro}

% \begin{macro}{\printdates}
% \begin{macro}{\noprintdates}
% Command to activate the printing out of the received/accepted
% dates. This can also be done using the \texttt{printdates} option,
% which is the default for \texttt{journal}, but not for \texttt{proceedings}.
%    \begin{macrocode}
\newcommand{\printdates}{\@printdatestrue}
\newcommand{\noprintdates}{\@printdatesfalse}
%    \end{macrocode}
% \end{macro}
% \end{macro}
% \begin{macro}{\@outputaddressesifnecessary}
% Command to output the addresses if we haven't already done so
% This is called by: 
% \begin{enumerate}
% \item The end of |thebibliography| environment
% \item The |\outputfulladdresses| command
% \end{enumerate}
%    \begin{macrocode}
\newcommand{\@outputaddressesifnecessary}{%
  \ClassInfo{rmaa}{Calling @outputaddressesifnecessary}
  \ifthenelse{\boolean{@needaddresses}}{%
    %% If in one column style, or in the first column of a two column
    %% style then we can try to put it on this page 
    \ifthenelse{\boolean{@firstcolumn} \or \not\boolean{@rmtwocol}}{%
      \@outputaddresses}{%
      %% otherwise give warning it will be on a separate page
      \ClassWarning{rmaa}{Full addresses will go on a separate
        page (please see User Guide for how to remedy this)}}%
    %% If we haven't managed yet, put it on its own page
    \ifthenelse{\boolean{@needaddresses}}{%
      \if@debug{\ClassWarning{rmaa}{In this silly bit!}}\fi
      \clearpage\vfill\@outputaddresses}{}%
    }{
    %% If we don't need them, we still might need to do the eoa hook 
    \ifthenelse{\boolean{@suppressfulladdresses}}{%
      \if@debug{\ClassWarning{rmaa}{XXX}}\fi
%%      \if@book\else\@endofarticlehook\null\fi
      }{}
    }%
  }  
%    \end{macrocode}
% \end{macro}

% \begin{macro}{\outputfulladdresses}
% A user interface to the above command, in case there is no biblio
% (11 May 2000) 
%    \begin{macrocode}
\newcommand{\outputfulladdresses}{\@outputaddressesifnecessary}
%    \end{macrocode}
% \end{macro}

% \begin{macro}{\adjustfinalcols}
% This command starts a new column and should be 
% used to ensure that the columns balance on the last page. It also
% outputs the full 
% author addresses to ensure that this occurs in the left hand
% column. 
% The |\par| command is so that the current line will be taken account
% of when working out |\pagetotal|
%    \begin{macrocode}
\newcommand{\adjustfinalcols}{\if@twocolumn\goodbreak\@outputaddresses\newpage\noindent\fi}
%    \end{macrocode}
% \end{macro}

% \begin{macro}{\@outputaddresses}
% Command to do the work of outputting the addresses. 
% This can be called by 
% \begin{enumerate}
% \item an explicit |\adjustfinalcols|
% \item via |\@outputaddressesifnecessary|, q.v.
% \end{enumerate}
%
%    \begin{macrocode}
\newcommand{\@outputaddresses}{%
  % Allow option of suppressing the full addresses, in which case
  % we just run the hook if we are not in a book style
  \if@suppressfulladdresses
%%     \if@book\else\@endofarticlehook\null\fi
  \else
  %% save the full addresses in a box for output later
  \sbox{\@fulladdbox}{\parbox{\textwidth}{\@fulladdresses}}
  % calculate height of address box  
  \settoheight{\@fulladdheight}{\usebox{\@fulladdbox}}
  \settodepth{\@fulladddepth}{\usebox{\@fulladdbox}}
  \addtolength{\@fulladdheight}{\@fulladddepth}
  \if@debug
  { \typeout{Address total height: \the\@fulladdheight}
    \typeout{Text height: \the\pagetotal}
    \typeout{Goal height: \the\pagegoal}
    }
  \fi
  % reuse \@fulladdheight to mean the total height of text+addresses
  \addtolength{\@fulladdheight}{\pagetotal}
%%  \addtolength{\@fulladdheight}{-\baselineskip}
  \if@debug\typeout{Adjusted text + addresses height: \the\@fulladdheight}\fi
  % check it will fit on current page
  \ifdim\@fulladdheight<\pagegoal
     %% must shift it to the left
     %% to compensate for the hanging indentation in the reference list. 
     %% NOTE: this is not foolproof!
     %% NB. This will give overfull hbox warnings, which is a bit ugly.
     \vfill\noindent\if@biblist\hspace*{-\leftmargin}\fi
     \usebox{\@fulladdbox}
     % Set this to show we've done it
     \@needaddressesfalse
     % and assume we are at the end of the article
     % so that we do the EOA hook if not in the book style
%%     \if@book\else\@endofarticlehook\fi
  \else
     % if not, just do nothing. With luck, this function will get 
     % called again on the next page
     \if@debug{\typeout{Holding over full addresses to next page}}\fi
  \fi\fi
  }
%    \end{macrocode}
% \end{macro}

% \subsubsection{Commands for editors}
% \label{sec:editor-commands}
% Some user/editor commands for doing the TOC and index in the book
% option. 
% They are defined here rather in the book-specific part since they
% need to be always available in order to avoid breaking documents. 
% \begin{macro}{\listofauthors}
% \begin{macro}{\@rawauthors}
% \begin{macro}{\rmaaTOCauthors}
% The command |\listofauthors| formats an author list for use in the
% TOC. It should contain the same names as in the |\author| command,
% but without the affiliation garbage and using |\&| instead of
% |and|. 
%    \begin{macrocode}
\newcommand{\@rawauthors}{Raw Authors}
\newcommand{\rmaaTOCauthors}{\noexpand\rmaaTOCAuthorStyle{Author list for TOC}}
\newcommand{\listofauthors}[1]{%
  \renewcommand{\@rawauthors}{#1}
  \renewcommand\rmaaTOCauthors{\noexpand\rmaaTOCAuthorStyle{#1}}}
%    \end{macrocode}
% \begin{macro}{\indexauthor}\label{indexauthor}
% The command |\indexauthor| produces an index entry for each
% author. 
% This command must come \emph{after} the
% |\title| and the |\listofauthors| commands. Also, there should be one 
% |\indexauthor| command for \emph{each} author. We also add the
% author onto the end of the |\ADS@authors| macro and construct the
% |\ADS@authorinitial| by a most cunning method (see below), but only
% for the first author. 
% 
% No, No, NO! This was all wrong. Backed out completely: 
% \begin{quote}
% Completely rewritten (19 Dec 2002) since we now use the index class, so 
% as to be able to alphabetize accented characters. A side effect of this is
% that |\index| no longer expands out the argument after the |@|, so we 
% instead add all the |\index| commands to a list (|\indexauthor@deferlist|), 
% which is executed in the |\@begofarticlehook|. 
% 
% Note the stuff with |\toks0| is copied from Exercise 20.15 of the \TeX{}book---it is
% to make sure the latin1-accented name gets passed into the idx
% file. 
% \end{quote}
% We now go back to the old way of doing things, using an optional
% argument to the \cs{indexauthor} command that has the name without
% accents. 
% \changes{3.08}{2003/01/16}{Remerged changes from 3.06 that got
% dropped from 3.07 after a synchronization disaster. Basically, the
% accentless optional argument to \cs{indexauthor}. Must add the macro
% directories to my unison set-up to stop this sort of thing from happening.}
%    \begin{macrocode}
\newcommand\@isortkey{}
\newcommand\@emptykey{empty}
\newcommand\indexauthor[2][empty]{
  \renewcommand\@isortkey{#1}
  \ifx\@isortkey\@emptykey \renewcommand\@isortkey{#2}\fi
  \index{\@isortkey @{{\rmaaIndexAuthorStyle{#2}}\hskip1em plus1em minus0.4em
      \rmaaIndexTitle\hskip1em plus2em minus0.4em\penalty-100
      \rmaaTOCauthors}}
  \indexauthorADS{#2}
}
%%\def\indexauthor@deferlist{\relax}
%%\long\def\indexauthor#1{
%%  \def\@@name{#1} \toks0=\expandafter{\@@name}
%%  \toks1=\expandafter{\indexauthor@deferlist}
%%  \protected@xdef\indexauthor@deferlist{%
%%    \the\toks1
%%    \noexpand\index{\the\toks0@{\noexpand\rmaaIndexAuthorStyle{\the\toks0}\noexpand\hskip1em plus1em minus0.4em 
%%        \rmaaIndexTitle\noexpand\hskip1em plus2em minus0.4em\noexpand\penalty-100
%%        \rmaaTOCauthors}}
%%  }
%%  \indexauthorADS{#1}
%%  }
%%\xdef\indexauthor@deferlist{}
\newcommand\indexauthorADS[1]{%
  \if@firstauthor\protected@xdef\ADS@authorinitial{\firstletter{#1}}\fi
  \protected@xdef\ADS@authors{\ADS@authors\if@firstauthor\else ; \fi
    #1}
  \@firstauthorfalse
}
%    \end{macrocode}
% Note that the first |#1| sets the key for the index, making sure it
% actually starts with  
% first letter of the surname. Meanwhile, the |@{...}| expression,
% which starts with some  
% control sequences is what is written to the |.ind| file. 
% See the makeindex man page for details. 
% \end{macro}
% \end{macro}
% \end{macro}
% \end{macro}

% \begin{macro}{\KillAllAccents}
% \begin{macro}{\firstletter}
% These are utility functions used in the above. |lett| is a counter
% indicating the current letter in the string. |\firstletter| just
% prints the first letter of its argument.
%    \begin{macrocode}
\newcommand{\KillAllAccents}{\def\'{}\def\`{}\def\~{}\def\.{}\def\i{i}}
\newcommand\firstletter[1]{\@car #1\@nil}
%    \end{macrocode}
% \end{macro}
% \end{macro}
% \begin{macro}{\SetVolume}
% \begin{macro}{\SetMSnumber}
% \begin{macro}{\SetFirstPage}
% \begin{macro}{\SetYear}
% \begin{macro}{\ReceivedDate}
% \begin{macro}{\AcceptedDate}
% These are for the editors to set the volume number, etc. 
%    \begin{macrocode}
\newcommand{\SetVolume}[1]{%
  \renewcommand{\@volume}{#1}
  \edef\ADS@volume{\PadToFourWithDots{#1}}
  \renewcommand{\@volumen}{\rmaa@VolumenStyle{volumen #1}}%
  }
\newcommand{\SetMSnumber}[1]{\renewcommand{\@MSnum}{#1}}
%    \end{macrocode}
% Note that there is no real reason to use |\SetFirstPage| since all
% the page numbering is done automatically in the book style. 
%    \begin{macrocode}
\newcommand{\SetFirstPage}[1]{\setcounter{page}{#1}}
\newcommand{\SetYear}[1]{%
  \renewcommand{\@year}{#1}
  \renewcommand{\ADS@year}{#1}
  \renewcommand{\@anyo}{\rmaa@VolumenStyle{#1}}
  }
\newcommand{\ReceivedDate}[1]{\renewcommand{\@received}{#1}}
\newcommand{\AcceptedDate}[1]{\renewcommand{\@accepted}{#1}}
%    \end{macrocode}
%
% \begin{macro}{\DOI}
% User/editor command to enter the DOI, which will be typeset as a link.
% \changes{4.0}{2019/11/02}{New user command \cs{DOI} to typeset
% DOI in the header of first page and add it to ADS info.}
%    \begin{macrocode}
\newcommand\DOI[1]{%
  \def\@DOI@url{\url{https://doi.org/10.22201/ia.01851101p.#1}}%
  \def\ADS@doi{https://doi.org/10.22201/ia.01851101p.#1}%
}%
%    \end{macrocode}
% \end{macro}
%
% \begin{macro}{\ORCIDiD}
% \changes{4.0}{2019/11/12}{Add mechanism for giving ORCID iD links}
% \begin{macro}{\ORCIDiDfull}
% \begin{macro}{\@orcid@icon}
% \begin{macro}{\@orcid@iconbig}
% Commands to add icon links to ORCID iDs of authors. For example
% \begin{quote}
% |\ORCIDiD|\marg{0000-0000-0000-0000}
% \end{quote}
% to be used after each author name in the |\author|
% command and in the |\fulladdresses|.  Alternatively,
% |\OORCIDiDfull| will typeset the icon plus the full text of the
% link, but we don't use that at present.  The icon image is stored in |ORCIDiD_iconvector.pdf|, which needs to be on the input path.  In principle, we could make EPS versions if necessary.
%    \begin{macrocode}
\newcommand\@orcid@icon{\includegraphics[height=1.8ex]{ORCIDiD_iconvector}}
\newcommand\@orcid@iconbig{\includegraphics[height=2ex]{ORCIDiD_iconvector}}
\newcommand\ORCIDiD[1]{\href{https://orcid.org/#1}{\,\@orcid@icon}}
\newcommand\ORCIDiDfull[1]{%
  \href{https://orcid.org/#1}{\@orcid@iconbig}\,\url{https://orcid.org/#1}}
%    \end{macrocode}
% \end{macro}
% \end{macro}
% \end{macro}
% \end{macro}
%
% \begin{macro}{\@volume}
% \begin{macro}{\@volumen}
% \begin{macro}{\@artpages}
% \begin{macro}{\@year}
% \begin{macro}{\@anyo}
% \begin{macro}{\@received}
% \begin{macro}{\@accepted}
% \begin{macro}{\@MSnum}
% Default values for the containers used by above commands, plus a
% few more useful ones. 
%    \begin{macrocode}
\newcommand{\@volume}{00}
\newcommand{\@volumen}{\rmaa@VolumenStyle{00}}
\newcommand{\@artpages}{%
  \oldpageref{\articlename:RMAAFirstPage}--\oldpageref{\articlename:RMAALastPage}}
\newcommand{\@year}{\the\year}
\newcommand{\@anyo}{\rmaa@VolumenStyle{2007}}
\newcommand{\@journalname}{\rmaa@JournalNameStyle{%
    \if@proceedings\RMAASClong\else\RMAAlong\fi}}
\newcommand{\@received}{}
\newcommand{\@accepted}{}
\newcommand{\@MSnum}{0000}
\newcommand{\@DOI@prefix}{https://doi.org/10.22201/ia.01851101p}
\def\@DOI@url{\url{https://doi.org/10.22201/ia.01851101p.2099.00.00.00}}
%    \end{macrocode}
% \end{macro}
% \end{macro}
% \end{macro}
% \end{macro}
% \end{macro}
% \end{macro}
% \end{macro}
% \end{macro}
% \end{macro}
% \end{macro}
% \end{macro}
% \end{macro}
% \end{macro}
% \end{macro}
% 
% The following commands are most useful in the \texttt{book}
% style. However, the |\SetConfTitle| command can now be used in
% individual articles to set the header for the \texttt{preprint}
% style. 
%    \begin{macrocode}
\newcommand{\@reunionplace}{Lugar de la reuni\'on}
\newcommand{\@reuniondate}{Fecha de la reuni\'on}
\newcommand{\@conftitle}{Conference Title}
\newcommand{\SetReunion}[2]{\renewcommand{\@reunionplace}{#1}\renewcommand{\@reuniondate}{#2}}
\newcommand{\SetConfTitle}[1]{\renewcommand{\@conftitle}{#1}}
\newcommand{\@editors}{Juan Ram\'on de la Fuente and ``El Mosh''}
\newcommand{\SetEditors}[1]{\renewcommand{\@editors}{#1}}
\newcommand{\@editores}{Juan Ram\'on de la Fuente y ``El Mosh''}
\newcommand{\SetEditores}[1]{\renewcommand{\@editores}{#1}}
\newcommand{\@mes}{mes de la publicaci\'on}
\newcommand{\SetMes}[1]{\renewcommand{\@mes}{\rmaa@VolumenStyle{#1}}}
%    \end{macrocode}

% \subsection{Notes for Author Proofs and Communication between Editors}
% \label{sec:notes}
% We introduce a command \cs{Note} that should be used to describe
% any change made to the manuscript. These are a bit tricky to get
% right. In general we use \cs{marginpar} but this only works in outer
% paragraph mode, so we just bung the box where we are if we are in a
% figure caption or something. If we are in the preamble, then we just
% use a \marginpar as part of the \maketitle command.
%    \begin{macrocode}
\newif\if@arenotes
\@arenotesfalse
\newcounter{notecounter}
\setcounter{notecounter}{0}
\newcommand\NotesList{}
\newcommand\Note[1]{%
  \if@proof{%
    \@arenotestrue%
    \stepcounter{notecounter}%
    \if@inpreamble\else{%
      \NoteMark{\arabic{notecounter}}%
    }\fi
    \protected@xdef\NotesList{\NotesList 
      \protect\item #1}%
  }\fi%
}
\newcommand\Changed[2]{Changed ``#1'' to ``#2''.}
\newcommand\NoteMark[1]{%
  \ifinner
  \raisebox{1ex}[0pt][0pt]{\makebox[0pt][l]{\framebox{\footnotesize\bfseries
        N#1}}}%
  \else
  \marginpar{\framebox{\footnotesize\bfseries N#1}}%
  \fi
}
\newcommand\RefNote[1]{\if@proof\NoteMark{\ref{#1}}\relax\fi}
\newcommand\NoteLabel{\protect\label}
\newcommand\OutputNotes{%
  \bgroup
  \onecolumn
  \section*{EDITORIAL NOTES}
  \begin{enumerate}
    \NotesList
  \end{enumerate}
  \egroup
}
\DeclareRobustCommand\cs[1]{\texttt{\char`\\#1}}

%    \end{macrocode}
% \subsection{Bibliographic Records for ADS}
% \label{sec:ADS}
% \changes{2.2beta11}{2001/11/27}{Just about finished all the ADS
% stuff now. Except for the references, which I have given up on for
% the time being.}
% \begin{macro}{\ADS@BibCode}
% \begin{macro}{\ADS@year}
% \begin{macro}{\ADS@journal}
% \begin{macro}{\ADS@volume}
% \begin{macro}{\ADS@firstpage}
% \begin{macro}{\ADS@authorinitial}
% We use the |\ADS@BibCode| command to hold the bibliographic code for
% ADS. It is 
% used in the |%R| 
% field of the bibliographic record for each article. 
% The format is 
% \begin{quote}
% \texttt{YYYYJJJJJVVVVMPPPPA}  
% \end{quote}
% where \texttt{YYYY} is the year, \texttt{JJJJJ} is the journal code,
% which is \texttt{RMxAA} or \texttt{RMxAC}, \texttt{VVVV} is the
% volume number (right-justified), \texttt{M} is not used by us,
% \texttt{PPPP} is the first page number (right-justified), and
% \texttt{A} is the initial of the first author's surname. The volume
% and page fields are padded with periods (`\texttt{.}'), which causes
% us some fun (this is done with the macro |\PadToFourWithDots| described
% below) . 
% 
% All these individual bits are held in the commands |\ADS@year|,
% |\ADS@journal|, |\ADS@volume|, |\ADS@firstpage|,
% |\ADS@authorinitial|. 
%    \begin{macrocode}
\newcommand{\SetADSBibCode}[1]{\renewcommand{\ADS@BibCode}{#1}}
\newcommand{\ADS@BibCode}
{\ADS@year\ADS@journal\ADS@volume.\PadToFourWithDots{\ADS@firstpage}\ADS@authorinitial}
\newcommand{\ADS@year}{2001} 
\newcommand{\ADS@journal}{\if@proceedings RMxAC\else RMxAA\fi}
\newcommand{\ADS@volume}{....} 
%    \end{macrocode}
% These two are just set to default values here, |\ADS@firstpage| is set
% in the |document| environment for included articles, in the
% |frontsection| environment, |\tableofcontents| command,  |index| and
% |abstracts| environments. Also by |\AtBeginDocument| for stand-alone
% articles. 
% 
% God knows how we are going to do |\ADS@authorinitial|!
%    \begin{macrocode}
\newcommand{\ADS@firstpage}{....}
\newcommand{\ADS@authorinitial}{A}
%    \end{macrocode}
% \end{macro}
% \end{macro}
% \end{macro}
% \end{macro}
% \end{macro}
% \end{macro}
% We also have a whole slew of other fields that can be set. These are
% all listed in Table~\ref{tab:ADS}. Most of them still have to be
% defined. 
% \begin{table}[!ht]\raggedleft
% \caption{ADS Bibliographic Records}\label{tab:ADS}
% \begin{tabular}{lp{8em}lp{8em}}\hline
% Code & Description & Macro Name & Comment \\ \hline
% |%R| & Bibliographic Code & |\ADS@BibCode| & required \\
% |%A| & Author list & |\ADS@authors| & required, generate from |\indexauthor| \\
% |%F| & Author affiliation & |\ADS@affil| & format: AA(\dots)
% AB(\dots) \dots \\
% |%J| & Journal Name & |\ADS@journalname| &  \\
% |%D| & Publication Date & |\ADS@fulldate| &  \\
% |%L| & Last Page & |\ADS@lastpage| & generate from label \\
% |%T| & Title & |\ADS@title| & required, generate from |\title| \\
% |%C| & Abstract Copyright & |\ADS@copyright| & set to IA-UNAM \\
% |%U| & URL for electronic document & |\ADS@url| & set to PDF file \\
% |%d| & DOI for article & |\ADS@doi| & set by |\DOI| \\
% |%K| & Keywords & |\ADS@keywords| & set equal to |\@keywords| \\
% |%A| & Abstract Text & |\ADS@abstract| & set equal to |\@abstract| \\
% |%R| & References & |\ADS@references| & generate from the
% bibliography  \\ \hline
% \end{tabular}
% \end{table}
% A lot of these macros are set by other commands, as mentioned in the
% Table, but some need to be set  explicitly here. 
% 
% \begin{macro}{\ADS@doi}
% \changes{4.0}{2019/11/08}{Save DOI for ADS}
% The user/editor command \cs{DOI} is used to set this. 
%    \begin{macrocode}
\newcommand{\ADS@doi}{https://doi.org/10.22201/ia.01851101p.0000.00.00.00}
%    \end{macrocode}
% \end{macro}
% 
% \begin{macro}{\ADS@journalname}
% This
% does not contain any accents at the moment so as not to cause
% problems---this maybe should be changed though. 
%    \begin{macrocode}
\newcommand{\ADS@journalname}{%
  \if@proceedings\@conftitle\space (Eds. \@editors) \fi%
  Revista Mexicana de Astronom\'\i{}a y Astrof\'\i{}sica%
  \if@proceedings\space (Serie de Conferencias)\fi%
  \space Vol. \@volume, pp. \ADS@firstpage--\ADS@lastpage\space (\@year)%
  \space (http://www.astroscu.unam.mx/\~{}rmaa/)
  }
%    \end{macrocode}
% \end{macro}
% \begin{macro}{\ADS@keywords}
% \begin{macro}{\ADS@abstract}
% \begin{macro}{\ADS@title}
% |\ADS@abstract| is straightforward, since we already have the text
% saved as |\@abstract|. We could do the same with |\ADS@keywords| but
% to avoid unwanted formatting we explicitly set it in |\keywords|
% (p.~\pageref{keywords}). |\ADS@title| is just set to |\@rawtitle|,
% which we were already using for another purpose. 
%    \begin{macrocode}
\let\ADS@title=\@rawtitle
\newcommand{\ADS@keywords}{No keywords specified}
\newcommand{\ADS@abstract}{No abstract specified}
%    \end{macrocode}
% \end{macro}\end{macro}\end{macro}
% \begin{macro}{\ADS@authors}
% |\ADS@authors| wants to be in ``Other, A. N.'' format, but with
% semi-colons between authors. The most natural way to do this seems to
% be with |\indexauthor| (p.~\pageref{indexauthor}). We append each
% author to this macro, which means we must explicitly clear it for
% each new article (done in |\article|). .   
%    \begin{macrocode}
\newcommand{\ADS@authors}{}
%    \end{macrocode}
% \end{macro}
% \begin{macro}{\ADS@fulldate}
% \begin{macro}{\ADS@lastpage}
% \begin{macro}{\ADS@copyright}
% \begin{macro}{\ADS@website}
% \begin{macro}{\ADS@url}
% \begin{macro}{\SetWebsite}
% \begin{macro}{\SetMonthNumeric}
% Other sundry data, together with user commands for bits that can be
% set by the Editors. The |\ADS@url| is constructed from the website
% (currently |http://www.astrosmo.unam.mx/rmaa/|,
% the journal descriptor (e.g., |RMxAC..10/|), followed by the
% path to the PDF file (e.g., |PDF/RMxAC..10_HEW1.pdf|). 
% 
%    \begin{macrocode}
\newcommand{\ADS@website}{http://www.astrosmo.unam.mx/rmaa/}
\newcommand{\ADS@url}
{\ADS@website\ADS@journal\ADS@volume\ADS@number/PDF/\ADS@journal\ADS@volume\ADS@number\_\articlename.pdf}
\newcommand{\ADS@number}{}
\newcommand{\ADS@fulldate}{00/\ADS@year} 
\newcommand{\ADS@lastpage}{00} 
\newcommand{\ADS@copyright}{(C) \ADS@year; Instituto de Astronom\'\i{}a,
  Universidad Nacional Aut\'onoma de M\'exico} 
\newcommand{\SetNumber}[1]{\renewcommand{\ADS@number}{-#1}}
\newcommand{\SetWebsite}[1]{\renewcommand{\ADS@website}{#1}}
\newcommand{\SetMonthNumeric}[1]{\renewcommand{\ADS@fulldate}{#1/\ADS@year}}
%    \end{macrocode}
% \end{macro}\end{macro}\end{macro}\end{macro}\end{macro}\end{macro}\end{macro}
% \begin{macro}{\ADS@affil}
% \begin{macro}{\ADS@references}
% These two are somewhat more problematical. Affiliation currently
% only works for the first author (see \dots) pending redesign of the
% user interface. References are now done by wrapping the
% |thebibliography| environment to open another input stream on the
% same file and search through for the |\bibitems|. 
%    \begin{macrocode}
\newcommand{\ADS@references}{}
\newcommand{\ADS@affil}{}
%    \end{macrocode}
% \end{macro}\end{macro}
% 
% Define two macros that expand to the strings
% |\begin{thebibliography}| and |\end{thebibliography}|. Note that the
% idea for this and much else is inspired by the code for the docstrip
% package. Also define |\StartBibX| in case authors use the more
% standard |\begin{thebibliography}{}| syntax. 
%    \begin{macrocode}
\edef\EndBib{\expandafter\@gobble\string\\end\string{thebibliography\string}}
\edef\StartBib{\expandafter\@gobble\string\\begin\string{thebibliography\string}}
\edef\StartBibX{\expandafter\@gobble\string\\begin\string{thebibliography\string}\string{\string}}
%    \end{macrocode}
% Define two new logical variables. |\ifAnotherline| is true if we
% want to read another line from the input stream. It is set to false
% when we encounter the |\end{thebibliography}|. |\ifInBib| is set to
% true when we frind the |\begin{thebibliography}| and tells us to
% start saving the following lines. 
%    \begin{macrocode}
\newif\ifAnotherLine
\newif\ifInBib
%    \end{macrocode}
% \begin{macro}{makeOther}
% We need a macro to change the catcodes of lots of things to other. 
%    \begin{macrocode}
\def\makeOther#1{\catcode`#1=12\relax}
%    \end{macrocode}
% \end{macro}
% \begin{macro}{\RereadBiblio}
% Define a macro to do all the fiddly stuff. This is used at the
% beginning of |thebibliography| environment (Sec~\ref{sec:biblio}).
%    \begin{macrocode}
\newread\tf@biblio
\newcommand{\RereadBiblio}{%
%    \end{macrocode}
% Open the current file again for reading
%    \begin{macrocode}
  \begingroup
  \immediate\openin\tf@biblio=\articlename
%    \end{macrocode}
% Everything in a group to be safe. Initialize logicals. Black magic
% to change
% catcodes and avoid blank lines being interpreted as |\par|. 
%    \begin{macrocode}
  \AnotherLinetrue\InBibfalse
  \makeOther\ \makeOther\\\makeOther\$%
  \makeOther\#\makeOther\^\makeOther\^^K%
  \makeOther\_\makeOther\^^A\makeOther\%%
  \makeOther\~\makeOther\{\makeOther\}\makeOther\&%
  \endlinechar-1\relax
%    \end{macrocode}
% Start looping through the file. Read each line to the macro
% |\Line|. If we have got to EOF or end of biblio then set
% |\AnotherLinefalse|.
%    \begin{macrocode}
  \loop
  \read\tf@biblio to \Line
  \ifx\Line\EndBib\AnotherLinefalse\InBibfalse\fi
  \ifeof\tf@biblio \AnotherLinefalse\fi
%    \end{macrocode}
% If we are in the biblio and the line is not blank then append it to
% the |\ADS@references| macro together with a newline. 
%    \begin{macrocode}
  \ifInBib
  \ifx\Line\empty\else\xdef\ADS@references{\ADS@references\Line^^J}\fi
  \fi
%    \end{macrocode}
% Check if the line we just saw was the start of biblio, then repeat
% the loop if necessary. 
%    \begin{macrocode}
  \ifx\Line\StartBib\InBibtrue\fi
  \ifx\Line\StartBibX\InBibtrue\fi
  \ifAnotherLine\repeat
  \immediate\closein\tf@biblio
  \endgroup
  }
%    \end{macrocode}
% \end{macro}
% 
% \begin{macro}{\PadToFourWithDots}
% Utility function for left-padding a number to 4 characters with
% periods. Used in 
% constructing |\ADS@BibCode|. This is a very simple algorithm and
% \emph{requires} that the argument actually be a number---if it is
% not then dire consequences will result.  
%    \begin{macrocode}
\newcommand{\PadToFourWithDots}[1]{%
  \ifnum #1 > 999 #1%
  \else\ifnum #1 > 99 .#1%
  \else\ifnum #1 > 9 ..#1%
  \else ...#1%
  \fi\fi\fi}
%%  \ifnum \count0 > 999 #1%
%%  \else\ifnum \count0 > 99 .#1%
%%  \else\ifnum \count0 > 9 ..#1%
%%  \else ...#1%
%%  \fi\fi\fi}
%    \end{macrocode}
% \end{macro}
%
% \begin{macro}{\AuthorADS}
% \begin{macro}{\ADS@affiliations}
% New mechanism for setting the affiliations for ADS. The author labels (AA, AB, \dots) are generated using two nested counters.  We use the |\robustify| command from \texttt{etoolbox} so that accented characters in the affiliation addresses do not blow up.  The list of accents may need to be expanded. 
%    \begin{macrocode}
\newcounter{AUTHOR}
\newcounter{author}[AUTHOR]
\stepcounter{AUTHOR}
\newcommand\ADS@affiliations{}
\robustify{\'}\robustify{\i}\robustify{\~}\robustify{\`}
\def\AuthorADS#1#2{%
  \begingroup
  \def\ORCIDiD##1{<ID system="ORCID">##1</ID>}
  \def\Email##1{<EMAIL>##1</EMAIL>}
  \def\Affil##1{##1}
  \stepcounter{author}
  \ifthenelse{\value{author} > 26}{\stepcounter{AUTHOR}\stepcounter{author}}{}
  \ifthenelse{\equal{\ADS@affiliations}{\@rmaaEmpty}}{%
    \xdef\ADS@affiliations{\Alph{AUTHOR}\Alph{author}(#2)}
  }{%
    \xdef\ADS@affiliations{\ADS@affiliations, \Alph{AUTHOR}\Alph{author}(#2)}
  }%
  \endgroup
}
%    \end{macrocode}
% \end{macro}
% \end{macro}
% %
% \begin{macro}{\ADS@dumpdata}
% \changes{4.0}{2019/11/13}{Write all author affiliations for ADS}
% \begin{macro}{\ADS@write}
% \begin{macro}{\ADS@openfile}
% |\ADS@dumpdata| is the macro that actually does the job of writing out all the
% data to an auxiliary file, using |\ADS@write| to write each
% field. The file is set up with |\ADS@openfile|, which currently is
% done in the TOC of a book.  The new multi-author affiliations are used if |\ADS@affiliations| has been set, otherwise it falls back to the single-author version |\ADS@affil|. 
%    \begin{macrocode}
\newcommand{\ADSfield}[2]{\relax}
\newcommand{\ADSvfield}[2]{\relax}
%%\renewcommand{\ADSfield}[12]{\%#1 #2}
\newcommand{\ADS@dumpdata}
{ 
  \ADS@write{\string\begin{ADSarticle}}
  \ADS@WRITE{\string\ADSfield{R}{\ADS@BibCode}}
  \ADS@write{\string\ADSfield{A}{\ADS@authors}}
  \ifthenelse{\equal{\ADS@affiliations}{\@rmaaEmpty}}{%
    {\ADS@sanitize\ADS@write{\string\ADSfield{F}{AA(\ADS@affil)}}}
  }{%
    {\ADS@sanitize\ADS@write{\string\ADSfield{F}{\ADS@affiliations}}}
  }
  \ADS@write{\string\ADSfield{J}{\ADS@journalname}}
  \ADS@write{\string\ADSfield{D}{\ADS@fulldate}}
  \ADS@write{\string\ADSfield{L}{\ADS@lastpage}}
  {\ADS@sanitize\ADS@write{\string\ADSfield{T}{\ADS@title}}}
  \ADS@write{\string\ADSfield{C}{\ADS@copyright}}
  \ADS@write{\string\ADSfield{d}{\ADS@doi}}
  \ADS@write{\string\ADSfield{U}{\ADS@url}}
  \ADS@write{\string\ADSfield{K}{\ADS@keywords}}
  {\ADS@sanitize\ADS@write{\string\ADSfield{B}{\ADS@abstract}}}
  \ADS@write{\string\ADSfield{Z}{\ADS@references}}
  \ADS@write{\string\end{ADSarticle}}
  }
%%\newcommand{\ADS@write}[1]{\if@book\addtocontents{ads}{#1}\fi}
%%\newcommand{\ADS@openfile}{\@starttoc{ads}}
\newcommand{\ADS@write}[1]{\set@display@protect\immediate\write\tf@ads{#1}}
\newcommand{\ADS@WRITE}[1]{\immediate\write\tf@ads{#1}}
\newcommand{\ADS@openfile}{%
  \newwrite\tf@ads%
  \immediate\openout\tf@ads\jobname.ads\relax
  \ADS@WRITE{\ADS@preamble}}
\newcommand{\ADS@closefile}{\ADS@WRITE{\ADS@postamble}}
\def\ADS@preamble{%
  %%  \string\documentclass{article}
  %%  \string\begin{document}
  }
\def\ADS@postamble{%
  %%  \string\end{document}
  \string%\string% Local Variables:^J 
  \string%\string% mode: LaTeX^J
  \string%\string% mode: Font-lock^J
  \string%\string% End:^J
  }
%    \end{macrocode}
% \end{macro}
% \end{macro}
% \end{macro}
% 
% \subsection{Automatic calculation of the page range}
% \label{sec:page-range}
% Another issue that has cost me a lot of grief. Still, it finally
% seems to be sorted now. All hail Jeff Goldberg. 
%
% Cunning trick to get the page numbers from |\pageref|.  Put a label
% on the first and last page of the document using |\AtBeginDocument|
% and |\AtEndDocument|. .  This way won't work for the book option.
% In that case, we do it by defining |\@endofarticlehook| and
% |\@begofarticlehook| as parallels to |\AtEndDocument| and
% |\AtBeginDocument|.  Note that we have to defer adding
% |\@begofarticlehook| to |\AtBeginDocument| to make sure it comes
% after the redefinitions of |\ref| and |\label|.
%    \begin{macrocode}
\if@book\AtEndDocument{\ADS@closefile}\else
 %\AtBeginDocument{\@begofarticlehook}
\AtEndDocument{\@endofarticlehook}
\fi
%    \end{macrocode}
% \begin{macro}{\@begofarticlehook}
% Since we are getting more and more things that need to be done at
% the beginning of every article, frontsection, etc we define a sister
% hook to 
% |\@endofarticlehook|. This hook does the following: 
% \begin{enumerate}
% \item Labels the first page
% \item Sets |\ADS@firstpage|
% \item runs and then clears |\indexauthor@deferlist|
% \item resets the keyword counter
% \end{enumerate}
% It is executed in the following circumstances: 
% \begin{enumerate}
% \item At |\begin{document}| for standalone files
% \item At |\begin{document}| for each included file
% \item At beginning of front sections 
% \item At beginning of TOC
% \item At beginning of Author Index
% \item At beginning of abstracts
% \end{enumerate}
% For some reason, this is not always within the right group, so we
% use |\xdef| just to be sure. 
%    \begin{macrocode}
\newcommand{\@begofarticlehook}{%
  \if@frontmatter\frontlabel{RMAAFirstPage}\else\label{RMAAFirstPage}\fi%
  \xdef\ADS@firstpage{\@arabic{\c@page}}%
  \if@book\else\makeindex\fi % needed so default index is defined in stand-alone articles
%%  \indexauthor@deferlist
%%  \xdef\indexauthor@deferlist{}
  \@inpreamblefalse
  \setcounter{kwdcounter}{0}
  }
%    \end{macrocode}
% \end{macro}
% \begin{macro}{\@endofarticlehook}
% This is a hook to be run at the end of each article
% currently, it just labels the last page and writes out all the ADS
% data to an auxiliary file.  
% It should not produce any printed output. 
% It is modelled on the |\AtEndDocument| command, to which it reduces
% in the case of a single article. 
%    \begin{macrocode}
\newcommand{\@endofarticlehook}{%
  \if@debug{\ClassWarning{rmaa}{***Doing Hook***}}\fi
  \if@debug{\ClassWarning{rmaa}{ADSabstract is: \ADS@abstract}}\fi
  \clearpage\rmaa@lastpagelabel\if@book\ADS@dumpdata\fi
  \ifx\NotesList\@rmaaEmpty\else\OutputNotes\fi
  \if@debug{\ClassWarning{rmaa}{Notes list is: \NotesList}}\fi
}
%    \end{macrocode}
% \begin{macro}{\rmaa@lastpagelabel}
% New version of last page label command copied from Jeff Goldberg's
% lastpage.sty package - we can't just use |\label| here in case the
% output routine for the last page has already been called. Instead,
% we first explicitly clear the page, step back the page counter,
% explicitly write to the .aux file, then restep the page counter
% again.
%
% But this will conflict with the hyperref package, which does labels differently.  Compensate for that by using the five-argument form.
% 
% Also set the |\ADS@LastPage| macro to the appropriate number. 
%    \begin{macrocode}
\newcommand{\rmaa@lastpagelabel}{%
  \addtocounter{page}{-1}%
  \immediate\write\@auxout{\string
    \newlabel{\articlename:RMAALastPage}{{}{\thepage}{}{}{}}}%
  \edef\ADS@lastpage{\thepage}%
  \addtocounter{page}{1}}
%    \end{macrocode}
% \end{macro}
% \end{macro}

% \subsection{Deal with explicit change in number of columns}
% \label{sec:columns}
% Sometimes, when formatting a book, we need to switch between
% one-column articles and two-column articles. For this to work, we
% need to wrap the standard \LaTeX{} commands to set the |\@rmtwocol|
% flag appropriately and also reset the package default values of all
% scalable lengths. .  
%    \begin{macrocode}
\let\latexonecolumn=\onecolumn
\let\latextwocolumn=\twocolumn
\renewcommand{\onecolumn}{\@rmtwocolfalse\@package@default@lengths\latexonecolumn}
\renewcommand{\twocolumn}{\@rmtwocoltrue\@package@default@lengths\latextwocolumn}
%    \end{macrocode}
% 
% \subsection{The actual formatting of the title}
% \label{sec:title-format}

% Note, we have to redefine the |\maketitle| command to use the
% numbered footnotes. This means we needn't use the atbeginend stuff. 
%    \begin{macrocode}
\renewcommand{\maketitle}{\par
  \begingroup
%    \end{macrocode}
% Don't switch to marked footnotes. 
% This is consistent with existing style but I'm not sure I like it
%    \begin{macrocode}
  %%  \renewcommand\thefootnote{\fnsymbol{footnote}}%
%%  \def\@makefnmark{\hbox to\z@{$\m@th^{\@thefnmark}$\hss}}%
  \def\@makefnmark{\textsuperscript{\@thefnmark}\hspace{-0.2em}}%
  \long\def\@makefntext##1{\parindent 1em\noindent
    \hbox to1.8em{\hss$\m@th^{\@thefnmark}$}##1}%
  \if@rmtwocol
  { \if@debug{\ClassWarning{rmaa}{TWO COLUMN}}\fi
    \twocolumn[\@maketitle]%
    }
  \else
  { \if@debug{\ClassWarning{rmaa}{ONE COLUMN}}\fi
    \newpage
    \global\@topnum\z@   % Prevents figures from going at top of page.
    \if@manuscript \@MSmaketitle \else \@maketitle \fi
    }
  \fi
  \if@manuscript\else\thispagestyle{rmaafirstpage}\@thanks\fi
  \ifx\NotesList\@rmaaEmpty\else{%
    \NoteMark{1--\arabic{notecounter}}\vspace*{-2\baselineskip}%
  }\fi%
  \endgroup
%    \end{macrocode}
% Don't reset footnote counter (not sure whether this is good thing
% or not).
%    \begin{macrocode}
  %%  \setcounter{footnote}{0}%
  \let\thanks\relax\let\affil\relax\let\altaffil\relax
  \let\maketitle\relax\let\@maketitle\relax
  \gdef\@thanks{}\gdef\@author{}\gdef\@title{}}
%    \end{macrocode}

% \begin{macro}{\@kwhwidth}
% \begin{macro}{\@kwtwidth}
% Redefine \cs{@maketitle} to include the abstract, resumen and
% keywords. Note that in the manuscripty style we use
% \cs{@MSmaketitle} instead.  
% \changes{3.16}{2003/04/07}{
%  Added extra vertical space between authors and accepted date.}
% \changes{3.26}{2011/08/15}{Add extra label to review articles}
%    \begin{macrocode}
\newlength{\@kwhwidth}\newlength{\@kwtwidth}
\newcommand\reviewname{Review}
\renewcommand{\@maketitle}{%
  \newpage
  \null
  \if@reviewarticle
  \smallskip
  {\large\bfseries\itshape\reviewname}%
  \bigskip
  \else
  \vskip\@before@title@skip%
  \fi
  \begin{center}%
    {\@title \par}%
    \vskip\@before@author@skip%
    \parbox{\rmaa@AuthorWidthFrac\textwidth}{\centering
      \lineskip\@between@author@skip%
      \@author
      \par}
    %% this should be replaced with `Submitted to...' for preprints
    \if@printdates
    \vskip\@before@dates@skip%
    \if@preprint
    \rmaa@ReceivedStyle{Submitted to RMxAA \@received}
    \else
    \rmaa@ReceivedStyle{Received \@received; accepted \@accepted}
    \fi
    \fi
    \if@manuscript
    \rmaa@ReceivedStyle{Submitted to RMxAA \@received}
    \fi
    \if@guide
    \rmaa@ReceivedStyle{Version \rmaa@version, \rmaa@date}
    \fi
    \vskip\@before@abstract@skip%
    \if@abstractpresent
    \parbox{\rmaa@AbstractWidthFrac\textwidth}{%
      \setlength{\parindent}{0.9cm}%
      \centerline{\rmaa@AbstractHeaderStyle{Abstract}}\par%
      \vskip\@after@abstracthead@skip%
      \if@newstyle\noindent\fi\@abstract
      \vskip\@between@abstract@skip%
      \centerline{\rmaa@AbstractHeaderStyle{Resumen}}\par%
      \vskip\@after@abstracthead@skip%
      \if@newstyle\noindent\fi\@resumen
      \if@setresumen\else\@abstract\fi % append english abstract if
                                       % spanish one not set
      \if@guide\else            %do not use keywords for guide
      \vskip\@before@kwd@skip%
      % Format the Key Words with hanging indentation:
      % They are put in a \parbox with width equal to the line width
      % minus the width of the header. 
      { % Widths of keyword
        \settowidth{\@kwhwidth}{\@keywordheader~}   % header and
        \setlength{\@kwtwidth}{\linewidth}          % keyword text
        \addtolength{\@kwtwidth}{-\@kwhwidth}       %
        \noindent\@keywordheader~\parbox[t]{\@kwtwidth}{%
          \rmaa@KeywordStyle{\boldmath\uppercase{\@keywords}}}}%
      \fi
      }
      \fi
  \end{center}%
  \vspace*{\afterkwdskip}
  \par
  }
\newlength{\afterkwdskip}\setlength{\afterkwdskip}{\z@}
%    \end{macrocode}
% \end{macro}
% \end{macro}
%
%
% \begin{macro}{\reviewarticle}
% \begin{macro}{\reviewtoc}
% \begin{macro}{\reviewlof}
%   User command to specify a review article. Must be used in
%   preamble. Main effect is to place the text ``Review'' before the
%   article title.  Also add command for setting a table of contents
%   and a list of figures at the beginning of the article.
%    \begin{macrocode}
\newcommand{\reviewarticle}{\@reviewarticletrue}
\newcommand\reviewtoc{%
  \section*{CONTENTS}%
  \@starttoc{rtc}%
}
\newcommand\reviewlof{%
  \section*{LIST OF FIGURES}%
  \@starttoc{lof}%
}
%    \end{macrocode}
% \end{macro}
% \end{macro}
% \end{macro}
% \begin{macro}{\l@section}
%    \begin{macrocode}
\renewcommand*\l@section[2]{%
  \ifnum \c@tocdepth >\z@
    \addpenalty\@secpenalty
    \addvspace{0.8em \@plus\p@}%
    \setlength\@tempdima{1.5em}%
    \begingroup
      \parindent \z@ \rightskip \@pnumwidth
      \parfillskip -\@pnumwidth
      \leavevmode %\bfseries
      \advance\leftskip\@tempdima
      \hskip -\leftskip
      #1\nobreak\hfil \nobreak\hb@xt@\@pnumwidth{\hss #2}\par
    \endgroup
  \fi}
%    \end{macrocode}
% \end{macro}
%
% \begin{macro}{\@MSmaketitle}
% A special version of |\@maketitle| for the \texttt{manuscript}
% option. This can deal with a multi-page abstract/resumen. Not very
% well, mind. 
%    \begin{macrocode}
\newcommand{\@MSmaketitle}{%
  \if@debug\ClassInfo{rmaa}{Using special MS maketitle command}\fi
  \newpage
  {
    \thispagestyle{rmaafirstpage}
%     \setlength{\parindent}{0pt}
    \null\vskip 2ex
    \begin{center}
      \@title\par
      \@author\par
      \rmaa@ReceivedStyle{%
        \ifx\@received\@rmaaEmpty
          Draft version: \@date
        \else
          Received: \@received
          \ifx\@accepted\@rmaaEmpty
          \else
            ; Accepted: \@accepted
          \fi
        \fi}\par
    \end{center}
    \@thanks
    \centerline{\rmaa@AbstractHeaderStyle{Resumen}}
    \noindent\@resumen\par\vskip 2ex
    \centerline{\rmaa@AbstractHeaderStyle{Abstract}}
    \noindent\@abstract\par\vskip 2ex
    \noindent\@keywordheader~\rmaa@KeywordStyle{\@keywords}\par
%     \@keywordheader~\rmaa@KeywordStyle{\uppercase{\@keywords}}\par
  }%
  }
%    \end{macrocode}
% \end{macro}

% \subsection{The sectioning commands}
% \label{sec:section-commands}

% \begin{macro}{\section}
% \begin{macro}{\subsection}
% \begin{macro}{\subsubsection}
% \begin{macro}{\paragraph}
% \begin{macro}{\subparagraph}
% Redefine all the standard commands to use these lengths. 
%    \begin{macrocode}
\renewcommand\section{\@startsection {section}{1}{\z@}%
                                   {\abovesectionskip}%
                                   {\belowsectionskip}%
                                   %% how do we get upper case here
                                   {\centering\reset@font\normalsize}}
\renewcommand\subsection{\@startsection{subsection}{2}{\z@}%
                                     {\abovesectionskip}%
                                     {\belowsectionskip}%
                                     {\centering\reset@font\normalsize\itshape}}
\renewcommand\subsubsection{\@startsection{subsubsection}{3}{\z@}%
                                     {\abovesectionskip}%
                                     {\belowsectionskip}%
                                     {\centering\reset@font\normalsize\itshape}}
\renewcommand\paragraph{\@startsection{paragraph}{4}{\z@}%
                                    {\aboveparaskip}%
                                    {\belowparaskip}%
                                    {\reset@font\normalsize\bfseries}}
\renewcommand\subparagraph{\@startsection{subparagraph}{5}{\parindent}%
                                      {\aboveparaskip}%
                                      {\belowparaskip}%
                                      {\reset@font\normalsize\bfseries}}
%    \end{macrocode}
% \end{macro}
% \end{macro}
% \end{macro}
% \end{macro}
% \end{macro}

% \begin{macro}{secnumdepth}
% Make sure that we number the subsubsections---this is done by
% default in the |article| class but not the |book| class, which is
% why we must set it explicitly here.
%    \begin{macrocode}
\setcounter{secnumdepth}{3}
%    \end{macrocode}
% \end{macro}
% \changes{2.2beta10}{2001/11/07}{Fixed \cs{secnumdepth} for \cs{book}
% style.}

% \begin{macro}{\thesubsection}
% \begin{macro}{\thesubsubsection}
% For subsection and subsubsection, we want the numbers to be roman
% but the title to be italic. This seems the easiest way to go about it.
%    \begin{macrocode}
\renewcommand{\thesubsection}{{\upshape\thesection.\@arabic\c@subsection}}
\renewcommand{\thesubsubsection}{\thesubsection{\upshape.\@arabic\c@subsubsection}}
%    \end{macrocode}
% \end{macro}
% \end{macro}

% \begin{macro}{\@seccntformat}
% Add ``.'' to the number here so we can remove it from |\thesection|,
% etc. Reduce gap between number and title WJH 27 Jan 1995
%    \begin{macrocode}
\def\@seccntformat#1{\csname the#1\endcsname .\hskip 0.4em}
%    \end{macrocode}
% \end{macro}

% \begin{macro}{\@sect}
% I didn't really want to mess with the |\@sect| command, but it was
% the only way I could manage to force upper case. Other than the
% commented parts, it is identical to the standard version. 
%    \begin{macrocode}
\def\@sect#1#2#3#4#5#6[#7]#8{
  \ifnum #2>\c@secnumdepth
     \let\@svsec\@empty
  \else
     \refstepcounter{#1}%
     \let\@@protect\protect
     \def\protect{\noexpand\protect\noexpand}%
     \edef\@svsec{\@seccntformat{#1}}%
     \let\protect\@@protect
  \fi
%    \end{macrocode}
% This is the bit I changed (27 Jan 1995). Use upper case for
% |\section|, use mixed case for the rest. 
%    \begin{macrocode}
  \ifnum #2=1                   
     \def\sec@upcase##1{\uppercase{##1}}%
     \def\sec@shape{\upshape}%
  \else                         
     \def\sec@upcase##1{\relax{##1}}%
     \def\sec@shape{\itshape}%
  \fi
  \@tempskipa #5\relax
  \ifdim \@tempskipa>\z@
     \begingroup #6\relax
          \@hangfrom{\hskip #3\relax\@svsec}%
                    {\interlinepenalty \@M \sec@upcase{#8}\par}%
     \endgroup
     \csname #1mark\endcsname{#7}%
%    \end{macrocode}
% Only write TOC entries if doing the Guide style. 
%    \begin{macrocode}
     \if@guide 
        \addcontentsline{toc}{#1}{%
          \ifnum #2>\c@secnumdepth \else
             \protect\numberline{\csname the#1\endcsname}%
          \fi
          #7}%
     \fi
%    \end{macrocode}
% Also allow per-article TOC for review papers
%    \begin{macrocode}
     \if@reviewarticle
        \ifnum #2>2 \else
           \addcontentsline{rtc}{#1}{%
             \protect\numberline{\csname the#1\endcsname}%
             \sec@upcase{\sec@shape #7}}%
        \fi
     \fi
  \else
     \def\@svsechd{#6\hskip #3\relax  %% \relax added 2 May 90
       \@svsec #8\csname #1mark\endcsname{#7}%
       \if@guide
          \addcontentsline{toc}{#1}{%
            \ifnum #2>\c@secnumdepth \else
               \protect\numberline{\csname the#1\endcsname}%
            \fi
            #7}%
       \fi
       }%
  \fi
  \@xsect{#5}}
%    \end{macrocode}
% \end{macro}


% \subsection{Appendices}
% \label{sec:appendices}
% Make sure that appendices in individual articles work properly. We
% copy the definition from the article class (the book one doesn't do
% what we want).  
% \begin{environment}{appendices}
%   This uses an environment |appendices| inside of which each
%   |\section| will become a separate appendix.  Note that the
%   appendices must come before the references (otherwise the full
%   addresses mechanism will not work properly). We make the equation
%   numbering be subordinate to the appendix (equations A12, B6,
%   etc). Finally, we have a medskip followed by APPENDICES to
%   introduce them. \changes{3.17}{2002/04/08}{ Added appendices environment
%   that does the right thing.}
%    \begin{macrocode}
\newenvironment{appendices}{\par
  \begingroup
  \setcounter{section}{0}%
  \setcounter{subsection}{0}%
  \renewcommand\thesection{\@Alph\c@section}
  \renewcommand\theequation{\thesection\arabic{equation}}
  \medskip
  \section*{APPENDICES}}
{\endgroup}
%    \end{macrocode}
% \end{environment}
% \begin{environment}{appendix}
%   If there is only one appendix, then use |appendix| environment
%   instead.  \changes{4.0}{2019/11/12}{ Added
%   appendix environment for case where there is only one section.}
%    \begin{macrocode}
\renewenvironment{appendix}{\par
  \begingroup
  \setcounter{section}{0}%
  \setcounter{subsection}{0}%
  \renewcommand\thesection{\@Alph\c@section}
  \renewcommand\theequation{\thesection\arabic{equation}}
  \medskip
  \section*{APPENDIX}}
{\endgroup}
%    \end{macrocode}
% \end{environment}

% \subsection{Floating thingies}
% \label{sec:floats}   

%% 
%    \begin{macrocode}
\renewcommand{\thetable}{\@arabic\c@table}
\renewcommand{\thefigure}{\@arabic\c@figure}
\renewcommand{\figurename}{Fig.}
\renewcommand{\tablename}{TABLE}
%\renewcommand{\fnum@table}{\MakeUpperCase{\tablename}~\thetable}
%    \end{macrocode}
% For some reason, I have to write `TABLE' explicitly here
% If I use |\tablename|, it gets converted to mixed case---God knows
% why. 
% |\figurename| is also being ignored. 
%    \begin{macrocode}
\renewcommand{\fnum@table}{TABLE~\thetable}
\renewcommand{\fnum@figure}{Fig.~\thefigure}
%    \end{macrocode}

% Do captions in 9 pt face in the main journal style and in new-style
% anything. 
%    \begin{macrocode}
\long\def\@makefigcap#1#2{%
  \if@proceedings\else\small\fi
  \if@newstyle\small\fi 
  \vskip\abovecaptionskip
  \sbox\@tempboxa{#1. #2}%
  \ifdim \wd\@tempboxa >\hsize
    #1. #2 \par
  \else
    \global \@minipagefalse
    \hb@xt@\hsize{\hfil\box\@tempboxa\hfil}%
  \fi
  \vskip\belowcaptionskip}
%    \end{macrocode}

% Make sure table caption is upper case. 
%    \begin{macrocode}
\long\def\@maketabcap#1#2{%
  \hbox to\hsize{\hfil #1 \hfil}
  \vskip 6pt 
  \centering\MakeUpperCase{#2}\par
  \vskip 6pt}
%    \end{macrocode}

%    \begin{macrocode}
\newlength{\@rottabcapwidth}
\long\def\@makerottabcap#1#2{%
  \setbox\@tempboxa\hbox{\MakeUpperCase{#2}}%
  \ifdim\wd\@tempboxa >.8\vsize
     \setlength\@rottabcapwidth{.8\vsize}
  \else
     \setlength\@rottabcapwidth{\wd\@tempboxa}
  \fi
  \rotatebox{90}{%
    \begin{minipage}{\vsize}\centering
      \begin{minipage}{\@rottabcapwidth}\centering
        #1\\ \vskip 6pt \MakeUpperCase{#2}
      \end{minipage}%
    \end{minipage}
    }%
  \hskip 6pt
  }

\newenvironment{Sideways}
    {\begin{sideways}\begin{minipage}{\vsize}\centering}
    {\end{minipage}\end{sideways}}

\newenvironment{Sidewaystable}
{%
  \let\oldcaption=\caption%
  \def\caption{
    \def\@captype{table}\renewcommand{\@makecaption}{\@maketabcap}
    \oldcaption}
  \begin{sidewaystable}\centering}
  {\end{sidewaystable}\let\caption=\oldcaption}
\newenvironment{Sidewaystable*}
{%
  \let\oldcaption=\caption%
  \def\caption{
    \def\@captype{table}\renewcommand{\@makecaption}{\@maketabcap}
    \oldcaption}
  \begin{sidewaystable*}\centering}
  {\end{sidewaystable*}\let\caption=\oldcaption}
%    \end{macrocode}

% Not sure why we need to stick these lines in. 
%    \begin{macrocode}
\renewenvironment{figure}
                 {\renewcommand{\@makecaption}{\@makefigcap}\@float{figure}}
                 {\end@float}
\renewenvironment{figure*}
                 {\renewcommand{\@makecaption}{\@makefigcap}\@dblfloat{figure}}
                 {\end@dblfloat}
\renewenvironment{table}
                 {\renewcommand{\@makecaption}{\@maketabcap}\@float{table}}
                 {\end@float}
\renewenvironment{table*}
                 {\renewcommand{\@makecaption}{\@maketabcap}\@dblfloat{table}}
                 {\end@dblfloat}
%    \end{macrocode}

% Redefine |\@caption| so that we don't bother with the list of
% figures. This saves us having to worry about fragile commands in the
% captions. 
%    \begin{macrocode}
\long\def\@caption#1[#2]#3{%
  \par
%%  \addcontentsline{\csname ext@#1\endcsname}{#1}%
%%    {\protect\numberline{\csname the#1\endcsname}{\ignorespaces #2}}%
  \begingroup
    \@parboxrestore
    \if@minipage
      \@setminipage
    \fi
    \normalsize
    \@makecaption{\csname fnum@#1\endcsname}{\ignorespaces #3}\par
  \endgroup}
%    \end{macrocode}

% Make tables not quite so cramped. Allow |\hline\hline| to give a
% poor man's thick rule. 
%    \begin{macrocode}
\renewcommand{\arraystretch}{1.15}
\setlength{\doublerulesep}{0.4pt}
%    \end{macrocode}

% \begin{macro}{\tablecols}
% \begin{macro}{\@tablecols}
% \begin{macro}{\tabcaptioncontinued}
% |\tablecols| is a user command for setting the number of columns. It
% is used by commands that try and span the entire table (I think
% there must be a better way of acheiving
% this). |\tabcaptioncontinued| is 
% a user command for setting the caption of a continued table on the
% next page. 
%    \begin{macrocode}
\newcommand{\tablecols}[1]{\renewcommand{\@tablecols}{#1}}
\newcommand{\@tablecols}{1}
\newcommand{\tabcaptioncontinued}{%
  \multicolumn{\@tablecols}{c}{\addtocounter{table}{-1}TABLE~\thetable\ (CONTINUED)} \\ \\
  }
%    \end{macrocode}
% \end{macro}
% \end{macro}
% \end{macro}

% \subsubsection{Footnotes in tables}
% \label{sec:table-footnotes}
% \begin{macro}{\tablenotemark}
% \begin{macro}{\tablenote}
% \begin{macro}{\tablenoteskip}
% Old user commands |\tablenotemark{MARK}|,
% |\tablenote{MARK}{TEXT}|---now deprecated. 
%    \begin{macrocode}
\newcommand\tablenotemark[1]{$^\mathrm{#1}$} 
\newcommand\tablenote[2]{\if@proceedings\else\small\fi ${}^\mathrm{#1}$ #2} 
\newcommand\tablenoteskip{\\[0.5\baselineskip]} 
%    \end{macrocode}
% \end{macro}
% \end{macro}
% \end{macro}

% \begin{macro}{\tabnotemark}
% \begin{macro}{\tabnotetext}
% \begin{macro}{\tabnotetextTWO}
% Alternative way of doing footnotes in tables. 
%    \begin{macrocode}
\newcommand{\tabnotemark}[1]{\textsuperscript{\lowercase{#1}}}
%\newcommand{\tabnotetextTWO}[4]{\tabnotetext{#1}{#2}\tabnotetext{#3}{#4}}
\newcommand{\tabnotetextTWO}[4]{%
  \multicolumn{\@tablecols}{l}{%
    \parbox{0.5\tabnotewidth}{\small\textsuperscript{#1}#2}%
    \parbox{0.5\tabnotewidth}{\small\textsuperscript{#3}#4}%
    }%
  \\}
\newcommand{\tabnotetext}[2]{%
  \multicolumn{\@tablecols}{@{}p{\tabnotewidth}@{}}{\small\textsuperscript{#1}#2}\\}
\newlength{\tabnotewidth}
%    \end{macrocode}
% \end{macro}
% \end{macro}
% \end{macro}


% \begin{macro}{changemargin}
%    \begin{macrocode}
% Command to allow margins to be changed on the fly. Useful to allow
% tables and Figures to protrude into the margins. 
%
% \changes{3.24}{2007/09/14}{Added \texttt{changemargin} environment
% (copied from TeX FAQ).}
\newenvironment{changemargin}[2]{% 
  % From TeX FAQ http://www.tex.ac.uk/cgi-bin/texfaq2html?label=chngmargonfly
  \begin{list}{}{%
    \setlength{\topsep}{0pt}%
    \setlength{\leftmargin}{#1}%
    \setlength{\rightmargin}{#2}%
    \setlength{\listparindent}{\parindent}%
    \setlength{\itemindent}{\parindent}%
    \setlength{\parsep}{\parskip}%
  }%
  \item[]}{\end{list}}
%    \end{macrocode}
% \end{macro}

% \subsection{The bibliography and references}
% \label{sec:biblio}
% \begin{macro}{\beforebiblioskip}
% \begin{macro}{\afterbiblioskip}
% I'm not entirely sure why I introduced these two commands, since
% directly putting |\vspace| before and after
% |\begin{thebibliography}| seems to work well enough. 
%    \begin{macrocode}
\newlength{\beforebiblioskip}
\setlength{\beforebiblioskip}{0pt}
\newlength{\afterbiblioskip}
\setlength{\afterbiblioskip}{0pt}
%    \end{macrocode}
% \end{macro}
% \end{macro}
% \begin{macro}{thebibliography}
% The environment itself is done as starred section with extra
% vertical space after heading. This space is currently zero. 
% In the book style we do the cunning trick with rereading the file in
% order to get the references into the |\ADS@references| macro.
% 
% \changes{3.05}{2003/01/02}{Removed the \cs{beforebiblioskip} and
% \cs{afterbiblioskip} commands, since they cause unwanted vertical
% space in some instances.}  
%    \begin{macrocode}
\renewenvironment{thebibliography}{%
  \if@book\RereadBiblio\fi
%%  \vspace*{\beforebiblioskip}
  \section*{REFERENCES} \frenchspacing
%%  \vspace*{\afterbiblioskip}
  \small
%    \end{macrocode} 
% The references themselves are set as a list in order 
% to get the hanging indentation. They are done in a smaller font and
% with french spacing so we needn't worry about the distinction
% between sentence-ending periods and abbreviations. 
%    \begin{macrocode} 
  \list{\null}{\leftmargin 1.5em\labelwidth\z@\labelsep\z@\itemindent -1.5em
    \itemsep\z@\parsep\z@\usecounter{enumi}}
  \sloppy\clubpenalty4000\widowpenalty4000
  \sfcode`\.=1000\relax\@biblisttrue}
{\endlist\@biblistfalse\normalsize\@outputaddressesifnecessary}
%    \end{macrocode}
% \end{macro}
% \subsection{Uniquification of cross-references}
% \label{sec:uniquification}
% 
% \begin{macro}{\articlename}
% In order that the |\label{}| command in individual articles be
% guaranteed to give a unique anchor, we append the article name (held
% in the |\articlename| command) to each key. This requires
% redefinition of the following commands: |\cite|, |\nocite|, |\ref|,
% |\pageref|, |\label|. 
% \begin{macro}{\oldcite}
% \begin{macro}{\oldnocite}
% \begin{macro}{\oldref}
% \begin{macro}{\oldpageref}
% \begin{macro}{\oldlabel}
% The original versions of these commands are
% saved as |\oldcite|, |\oldnocite|, etc. 
%
% The uniquifying of cross-references is now always done, so 
% we don't get multiply-defined labels in the single article case
%    \begin{macrocode}
\newcommand{\articlename}{anon}
\AtBeginDocument{
  \let\oldcite=\cite 
  \let\oldnocite=\nocite 
  \let\oldref=\ref 
  \let\oldpageref=\pageref
  \let\oldlabel=\label
}
%    \end{macrocode}
% \end{macro}
% \end{macro}
% \end{macro}
% \end{macro}
% \end{macro}
% \end{macro}
%
% Uniquify all cross-references by prepending the article name.  This
% now has to go into |\AtBeginDocument| so it is working on the new
% versions that hyperref.sty has set up.  Then we need to put the
% |\@begofarticlehook| afterwards, so that the first page label is using
% the new definitions.  If amsmath is loaded, then we also need to patch
% the special label command that they use. \changes{4.0}{2019/11/18}{Fix labels for amsmath display equations, in case that package has been loaded.}
%    \begin{macrocode}
\AtBeginDocument{
  \renewcommand{\ref}[1]{\oldref{\articlename:#1}%
    \if@debug{\ClassInfo{rmaa}{Reference to \articlename:#1}}\fi}
  \renewcommand{\pageref}[1]{\oldpageref{\articlename:#1}%
    \if@debug{\ClassInfo{rmaa}{Page reference to \articlename:#1}}\fi}
  \renewcommand{\label}[1]{\oldlabel{\articlename:#1}%
    \if@debug{\ClassInfo{rmaa}{Label for \articlename:#1}}\fi}%
  \ifdef{\label@in@display}{%
    \let\oldlabel@in@display=\label@in@display
    \renewcommand{\label@in@display}[1]{\oldlabel@in@display{\articlename:#1}%
      \if@debug{\ClassInfo{rmaa}{AMSTeX Label for \articlename:#1}}\fi}%
  }{}
  \@begofarticlehook
}%
%    \end{macrocode}
% Citations are more complicated since they can contain a list of
% keys. These are really only here for use with the rmaacite package,
% which is utterly deprecated (02 Jan 2003) now that I have the natbib
% code working. The command \cs{cite} should \emph{not} be used with
% natbib: use \cs{citep}, \cs{citet}, etc.\@ instead. 
%    \begin{macrocode}
\renewcommand{\nocite}[1]{\@bsphack
  \@for\@citeb:=#1\do{%
    \edef\@citeb{\expandafter\@firstofone\@citeb:\articlename}%
    \if@filesw\immediate\write\@auxout{\string\citation{\@citeb}}\fi
    \@ifundefined{b@\@citeb:\articlename}{\G@refundefinedtrue
      \@latex@warning{Citation `\@citeb' undefined}}{}}%
  \@esphack}
\renewcommand{\cite}[1]{\oldcite{#1:\articlename}}
%   \end{macrocode}


% Horrible fiddling with the bibliography citation mechanism. 
% This has to be globally available in case rmaacite puts it in the
% \texttt{.aux} file. 
%    \begin{macrocode}
\def\lbibcite#1#2#3{\global\@namedef{b@#1}{#2}
                    \global\@namedef{by@#1}{#3}}
%    \end{macrocode}

% \begin{macro}{\oldbibitem}
% \begin{macro}{\@rmaaBibArg}
% \begin{macro}{\@rmaaSpace}
% \begin{macro}{rmaabib}
% Of course, |\bibitem| needs to get the uniquification treatment as
% well. While we are at it, we redefine it to use consecutive integers
% for the key if the user leaves it blank (using the counter
% |rmaabib|). This allows one to put 
% |\bibitem{}| for all your references without getting annoying
% messages about multiply-defined labels.
% 
% We also take the opportunity to insert the optional argument to
% \cs{bibitem} if it is missing in the form Author(DATE) so that natbib
% doesn't try to do numerical citations. Natbib already provides a
% hook \cs{@extra@b@citeb} to help uniquify the citation key, so we
% just plug into that. However, it puts it \emph{after} the original
% key, not before as I had had it. Easiest solution is to change the
% way that I do it. 02 Jan 2003. Note that the \cs{@extra@b@citeb}
% does not get written to the .aux  file in the \cs{citation} command
% but it does work. 
% \changes{3.14}{06 Feb 2003}{Reinstalled old definition of \cs{bibitem}
%  for the case where the nonatbib option is requested, so that the
%  rmaacite package will work. This means that you can't use rmaacite
%  and natbib in the same volume.}
%    \begin{macrocode}
\if@natbib\def\@extra@b@citeb{:\articlename}\fi
%%\else
\let\oldbibitem\bibitem
\newcommand{\@rmaaBibArg}{}
\newcommand{\@rmaaSpace}{ }
\newcounter{rmaabib}
\if@natbib
\renewcommand{\bibitem}[2][ABC(2003)]{%
  \renewcommand{\@rmaaBibArg}{#2}
  \ifx\@rmaaBibArg\@rmaaEmpty
     \addtocounter{rmaabib}{1}
     \oldbibitem[#1]{\thermaabib:\articlename}
  \else\ifx\@rmaaBibArg\@rmaaSpace
     \addtocounter{rmaabib}{1}
     \oldbibitem[#1]{\thermaabib:\articlename}
  \else 
     \oldbibitem[#1]{#2:\articlename}
  \fi\fi
  \renewcommand{\@rmaaBibArg}{}
  }
\else
\renewcommand{\bibitem}[1]{%
  \renewcommand{\@rmaaBibArg}{#1}
  \ifx\@rmaaBibArg\@rmaaEmpty
     \addtocounter{rmaabib}{1}
     \oldbibitem{\articlename:\thermaabib}
  \else\ifx\@rmaaBibArg\@rmaaSpace
     \addtocounter{rmaabib}{1}
     \oldbibitem{\articlename:\thermaabib}
  \else
     \oldbibitem{\articlename:#1}
  \fi\fi
  \renewcommand{\@rmaaBibArg}{}
  }
\fi
%%\fi
%    \end{macrocode}
% \end{macro}
% \end{macro}
% \end{macro}
% \end{macro}

% \begin{macro}{oldpcite}
% \begin{macro}{oldscite}
% Then, we need to repeat the above scheme, but for the |rmaacite|
% definition. 
%    \begin{macrocode}
\newcommand{\@redefinebibitem}{%
  \let\oldbibitem\bibitem
  \let\oldcite\cite
  \let\oldpcite\pcite
  \let\oldscite\scite
  \def\bibitem[##1<##2>]##3{%
    \renewcommand{\@rmaaBibArg}{##3}
    \ifx\@rmaaBibArg\@rmaaEmpty
       \addtocounter{rmaabib}{1}
       \oldbibitem[##1<##2>]{\thermaabib:\articlename}
    \else\ifx\@rmaaBibArg\@rmaaSpace
       \addtocounter{rmaabib}{1}
       \oldbibitem[##1<##2>]{\thermaabib:\articlename}
    \else 
       \oldbibitem[##1<##2>]{##3:\articlename}
    % trust the author here
%%       \oldbibitem[##1<##2>]{##3}
    \fi\fi
    }
  }
%    \end{macrocode}
% \end{macro}
% \end{macro}


% \subsection{HTML Table of Contents}
% \label{sec:hyper-TOC}
% \begin{macro}{\HyperTOCEntry}
% \begin{macro}{\HyperTOCpart}
% \begin{macro}{\HyperTOCmeeting}
% \begin{macro}{\HyperTOCeditors}
% \begin{macro}{\HyperTOCpreamble}
% These commands are used to write information to the \texttt{.htoc}
% file, which is used in generating a Web Page of the Table of
% Contents. Since they only ever appear preceded by |\protect| in the
% argument to an |\addtocontents| command,  they should never need to
% be expanded by the macros. However, we define them here just in case
% \dots
%    \begin{macrocode}
\newcommand{\HyperTOCEntry}[3]{\relax}
\newcommand{\HyperTOCpart}[1]{\relax}
\newcommand{\HyperTOCmeeting}[3]{\relax}
\newcommand{\HyperTOCeditors}[1]{\relax}
\newcommand{\HyperTOCpreamble}{\relax}
%    \end{macrocode}
% \end{macro}
% \end{macro}
% \end{macro}
% \end{macro}
% \end{macro}
%

% \subsection{Formatting the poster abstracts} 
% These can now be used without using the book option.
%    \begin{macrocode}

\newcommand{\@makeminititle}{%
  \newpage
  \null 
  \vskip 3em%
  \begin{center}%
    {\@title \par}%
  \end{center}
  \vskip 3em%
  }

\newif\if@abstractsads
\@abstractsadsfalse
\newcommand\abstractsADS{\@abstractsadstrue}

\if@book\RequirePackage{multicol}\fi
\newenvironment{abstracts}[1][Abstracts of Contributed Papers]{%
  %% First, make sure we start on an odd page (06 Jun 2000)
  \if@book\clearpage\thispagestyle{empty}\cleardoublepage\fi
  %%
  \renewcommand{\articlename}{ABSTRACTS}
  \title{#1}\shortauthor{Abstracts}\shorttitle{Abstracts}
  \listofauthors{}
  \if@book
  \addcontentsline{toc}{chapter}{\protect\@fiddleauthors{\@TOCtitle}{\rmaaTOCauthors}}%
  \addtocontents{htc}{\protect\HyperTOCEntry{\@rawtitle}{\@rawauthors}{\articlename}}
  \fi
  \pagestyle{rmaa} % this now seems necessary (06 Jun 2000)
  \thispagestyle{rmaafirstpage}
%%  \twocolumn[\@makeminititle]
  \renewcommand{\abstract}[1]{%
    \renewcommand\ADS@abstract{##1}%
    \renewcommand{\@abstract}{##1}\@writeabstract%
  }
  \renewcommand{\title}[2][]{%
    \gdef\ADS@authors{}%
    \gdef\ADS@affil{}%
    \gdef\ADS@references{}%
    \gdef\ADS@keywords{}%
    \renewcommand{\@rmaaArg}{##1}
    \ifx\@rmaaArg\@rmaaEmpty
       \renewcommand{\@TOCaftertitle}{}
    \else
       \renewcommand{\@TOCaftertitle}{ (\@rmaaArg)}
    \fi
    \renewcommand{\@title}{\rmaa@TOCTitleStyle{##2}}%
    \renewcommand{\@TOCtitle}{\rmaa@TOCTitleStyle{##2}\@TOCaftertitle}%
    \renewcommand{\rmaaIndexTitle}{\rmaa@IndexTitleStyle{##2\@TOCaftertitle.}}%
    \renewcommand\ADS@title{##2}%
    \xdef\ADS@firstpage{\@arabic{\c@page}}%
    }
  \renewcommand{\author}[1]{\let\@thanks\@empty\renewcommand\@author{##1} }
  \renewcommand{\footnotetext}[2][XX]{\par\noindent\textsuperscript{##1}~##2}
  \renewcommand{\footnoterule}{\kern-3\p@ \hrule \@width 3em \kern 2.6\p@}
  \if@book\@begofarticlehook\fi
  \if@book\begin{multicols}{2}[\@makeminititle]\else\@makeminititle\fi
  \begingroup
  \setcounter{footnote}{0}
  }% and at the end of the environment....
  {\endgroup\if@book\end{multicols}\listofabstracts\label{RMAALastPage}\clearpage\fi}
%% auxiliary commands used:
%% \@writeabstract does all the actual formatting 
\newcommand{\@writeabstract}{%
  \if@book%
  \addcontentsline{loa}{abstract}{\protect\@fiddleauthors{\@TOCtitle}{\rmaaTOCauthors}}%
  \fi%
  \begin{center}\@title\par\@author\end{center}
  \@abstract
  \par\noindent\footnoterule\nobreak\footnotesize\@thanks
  \par\noindent\hrulefill\goodbreak
  \vskip 1cm \@plus0.5cm
  %% write ADS entry
  \if@book\if@abstractsads\ADS@dumpdata\fi\fi%
  \endgroup % this pair is to make each abstract a self-contained group
  \begingroup
  \setcounter{footnote}{0}
  }

\if@book
%% List of abstracts (stuff written to .loa file)
\newcommand{\listofabstracts}{%
  \begin{multicols}{2}[%
    \vskip 1.5cm \vfill % leave at least 1.5 cm and go to bottom of page if possible
    \centerline{\MakeUpperCase{List of Abstracts}}
    \vskip 1.5cm \@plus0.5cm
    ][6cm] % this makes it start a new page if there is not at least 6cm of space
    \renewcommand{\@fiddleauthors}[2]{##1\hskip1em plus2em minus0.4em\penalty-100  ##2}
    \parindent\z@
%%    \parskip.5\p@ \@plus .7\p@\relax 
    \parskip3\p@ \@plus 1\p@ \@minus 1\p@\relax % allow quite a lot of
                                % leeway here to avoid breaks inside
                                % entries 
    \renewcommand{\\}{\relax} % temporarily turn off explicit line breaks just in case
    \renewcommand{\altaffilmark}[1]{\relax} % ditto with this
    \@starttoc{loa}
  \end{multicols}
  }
\newcommand*{\l@abstract}[2]{%
  \par\penalty-1000
  #1\hskip 1em\penalty0\null\nobreak\hfill#2}
\fi

%    \end{macrocode}
 
% \subsection{Code for typesetting an entire volume of conference
% proceedings} \label{sec:whole-volume}
%    \begin{macrocode}

\if@book
%% environment for including the individual articles
\newcommand{\@fiddleauthors}[2]{\relax}
\renewcommand{\theequation}{\@arabic\c@equation}
\newenvironment{articles}%
{ \begingroup % make it a group so that author definitions do not
              % spill over
  \gdef\ADS@authors{}
  \pagestyle{rmaa}
  \renewenvironment{document}{%
    \@begofarticlehook
    \addtocounter{chapter}{1}%
    \addcontentsline{toc}{chapter}{%
      \protect\@fiddleauthors{\@TOCtitle}{\rmaaTOCauthors}}%
    \addtocontents{htc}{\protect\HyperTOCEntry{\@rawtitle}{\@rawauthors}{\articlename}}
    }{\relax}
  \renewcommand{\documentclass}[2][foo]{\relax} 
  \renewcommand{\usepackage}[2][foo]{%
    \def\@packagename{##2}\def\@rmaacitename{rmaacite}%
    \ifx\@packagename\@rmaacitename
    \@rmcitetrue\ClassInfo{rmaa}{This paper uses rmaacite}
    \makeatletter\input{rmaacite.sty}\makeatother % read it in here so that all changes
    \else                                         % are local to paper
    \if@debug\ClassWarning{rmaa}{Package \@packagename\space was requested by \articlename\space}\fi
    \fi} 
  %% Don't allow individual articles to mess with the pages, volume or year
  \renewcommand{\SetFirstPage}[1]{\relax}
  \renewcommand{\SetYear}[1]{\relax}
  \renewcommand{\SetVolume}[1]{\relax}
  \renewcommand{\SetConfTitle}[1]{\relax}
  %% Don't allow individual articles to change these
  \renewcommand{\nonstopmode}{\relax}
  \renewcommand{\scrollmode}{\relax}
  \renewcommand{\outputplates}{\relax}
  %% XXXX need to check if using rmaacite and act accordingly
  %% avoid multiple declarations of the same length
  \let\oldnewlength=\newlength
  \renewcommand{\newlength}[1]{%
    \edef\reserved@a{\expandafter\@gobble\string ##1}%
    \@ifundefined{\reserved@a}{%
      \oldnewlength{##1}}{%
      \if@debug{\ClassInfo{rmaa}{\reserved@a\space already defined: doing nothing}}\fi%
      \relax}}%
  \renewcommand\newcommand{\providecommand}
  }%
{\endgroup\onecolumn}
     

%    \end{macrocode}
% Command to be used in `articles' environment for including an
% article
%    \begin{macrocode}
\newcommand{\article}[2][]{%
  % optional first argument is extra text to put in brackets after title in TOC
  % e.g. ``Invited Talk'' or ``Poster''
  \renewcommand{\articlename}{#2}
  % reset loads of counters
  \setcounter{footnote}{0}
  \setcounter{figure}{0}
  \setcounter{table}{0}
  \setcounter{section}{0}
  \setcounter{equation}{0}
  \setcounter{rmaabib}{0}
  \setcounter{author}{0}
  \setcounter{AUTHOR}{1}
  \gdef\ADS@authors{}%
  \gdef\ADS@affil{}%
  \gdef\ADS@affiliations{}%
  \gdef\ADS@references{}%
  \@firstauthortrue
  \@firstaffiltrue
  \@firstaltaffiltrue
  \@firstkwdtrue
  \@rmcitefalse
  \if@suppressfulladdresses\else\@needaddressestrue\fi
  \renewcommand{\@rmaaArg}{#1}
  \ifx\@rmaaArg\@rmaaEmpty
     \renewcommand{\@TOCaftertitle}{}
  \else
     \renewcommand{\@TOCaftertitle}{ (#1)}
  \fi
%    \end{macrocode}
% In the main journal articles must start on an odd page. 
%    \begin{macrocode}
  \if@journal\clearpage\thispagestyle{empty}\cleardoublepage\fi
%% Now input the file, then run the end-of-article hook (reinstated 28
  %% Sep 2001)
  {\input{\@articledir/\articlename.tex}\@endofarticlehook}}
%    \end{macrocode}
% Command to separate articles into sessions. Write an entry in the
% TOC, makes sure we start on an odd page (added 06 Jun 2000), then do
% a title page if \cs{SessionTitlePages} was called. 
% \changes{3.12}{04 Feb 2003}{Added option to have title page for each
% session. Turn it on with \cs{SessionTitlePages}.}
%    \begin{macrocode}
\newcommand{\session}[1]{
  %% this does nothing except write an entry in the TOC
  \addcontentsline{toc}{part}{\rmaa@TOCPartStyle{#1}}
  \addtocontents{htc}{\protect\HyperTOCpart{#1}}
  %% and also makes sure we start on an odd page (added 06 Jun 2000)
  \clearpage\thispagestyle{empty}\cleardoublepage
  \stepcounter{part}
  \if@sessiontitlepages{%
    \onecolumn
    \thispagestyle{empty}
    \renewcommand{\articlename}{Section\thepart TitlePage}
    \label{RMAAFirstPage}
    \make@session@title{#1}
    \label{RMAALastPage}
    \clearpage\thispagestyle{empty}\cleardoublepage
  }\fi
  }
\newcommand{\make@session@title}[1]{%
  {%
    \setlength\parindent{0pt}
    \vspace*{\stretch{1}}
    \rule{\linewidth}{1mm}
    \begin{flushright}\Huge\bfseries\boldmath
      #1
    \end{flushright}
    \rule{\linewidth}{1mm}
    \vspace*{\stretch{2}}
  }
}
%    \end{macrocode}
% Command to turn on separate title pages for each session. 
%    \begin{macrocode}
\newcommand{\SessionTitlePages}{\@sessiontitlepagestrue}
\newcommand{\BoldSessionTitles}{%
  \renewcommand{\rmaa@TOCPartStyle}[1]{\textbf{\rmaa@TOCTitleStyle{##1}}}
  }
\newcommand{\ArticleDir}[1]{\renewcommand{\@articledir}{#1/}\graphicspath{{1/}}}
\newcommand{\@articledir}{.}
%\newcommand{\ps@rmaafront}{\ps@rmaa}% page style for front matter (same as ever, actually)
\newcommand{\ps@rmaafront}{% page style for front matter:  page # in each foot
  \renewcommand\@oddfoot{\hfil{\thepage}}%
  \renewcommand\@evenfoot{{\thepage}\hfil}%
  % shorttitle in both heads
  \renewcommand\@evenhead{\hfil\@shorttitle\hfil}%
  \renewcommand\@oddhead{\@evenhead}
  }
\newcommand{\ps@rmaafrontfirstpage}{% page style for first page of front matter
  % page # in each foot, on outside edge
  \renewcommand\@oddfoot{\hfil\thepage}%
  \renewcommand\@evenfoot{\thepage\hfil}%
  % nothing in the even heads
  \renewcommand\@evenhead{\relax}%
  \renewcommand\@oddhead{\relax}%
  }
%% slight modification to the standard \frontmatter and \mainmatter
%% commands, just to set the \pagestyle
\renewcommand\frontmatter{%
  \latexonecolumn
  \cleardoublepage
  \@frontmattertrue
  \@mainmatterfalse
  \pagenumbering{roman}
  \pagestyle{rmaafront}
  \addtocounter{page}{2} % leave space for the title page
  }
\renewcommand\mainmatter{%
  \pagestyle{empty}
  \if@rmtwocol\twocolumn\fi
  \cleardoublepage
  \@frontmatterfalse
  \@mainmattertrue
  \pagenumbering{arabic}
  \pagestyle{rmaa}
  }
\renewcommand\backmatter{%
  \clearpage\pagestyle{empty}\cleardoublepage
  \@mainmatterfalse
  %% put extra space in the TOC
  \addtocontents{toc}{\bigskip}
  }
%    \end{macrocode}
%
% \subsection{Front sections}\label{sec:front-sections}
% \begin{macro}{\notitle}
% \begin{macro}{\shortname}
% \begin{macro}{frontsection}
% The environment \texttt{frontsection} is for doing a section in the
% front matter (preface, list of participants, or whatever). It has 
% one optional argument:
% \begin{enumerate}
% \item[\#1] Short title (no embedded spaces) for
% use in making PS files. E.g., ``list\_of\_participants''.   
% \end{enumerate}
% and two mandatory arguments: 
% \begin{enumerate}
% \item[\#2] Title of section (e.g., `Preface')
% \item[\#3] Commands to be evaluated before processing the contents
% of the environment. This is mainly useful for putting in the things
% that would go in the preamble of a normal article  (e.g.,
% |\listofauthors|, |\indexauthor|, |\notitle|) 
% \end{enumerate}
%    \begin{macrocode}
\newcommand{\notitle}{\@titlevisiblefalse} \newcommand{\shortname}{}
\newenvironment{frontsection}[3][]{%
  \begingroup
  \renewcommand{\shortname}{#1}
  \ifx\shortname\@rmaaEmpty
     \renewcommand{\articlename}{#2}
  \else
     \renewcommand{\articlename}{\shortname}
  \fi
  \@begofarticlehook
  \@titlevisibletrue % default
  \thispagestyle{rmaafrontfirstpage}
  \listofauthors{} % blank out this in case it is not specified
  \title{{#2}}\shorttitle{#2}\shortauthor{#2}
  #3 % evaluate the command given as second mandatory arg
  \addcontentsline{toc}{chapter}{\protect\@fiddleauthors{\@TOCtitle}{\rmaaTOCauthors}}%
  \addtocontents{htc}{\protect\HyperTOCEntry{\@rawtitle}{\@rawauthors}{\articlename}}
  \if@titlevisible
     \vspace*{1cm}\centerline{\rmaa@FrontSecStyle{\@title}}\vspace*{1cm}
  \fi
  \setlength{\parskip}{3ex}
  }{\frontlabel{RMAALastPage}\clearpage\endgroup}
%    \end{macrocode}
% \end{macro}
% \end{macro}
% \end{macro}
%    \begin{macrocode}
%% Definition of the command \frontlabel, which rights out something like `F4' to
%% the .aux file, instead of `\textlatin  {i}'. This is so it is easy to post-process
%% the file for making PS files of individual sections
\newcommand{\frontlabel}[1]{%
  \@bsphack
  \protected@write\@auxout{}%
     {\string\newlabel{\articlename:#1}{{}{F\@arabic{\c@page}}}}
  \@esphack
  }
%% User command to say X, Y, Z,\\ Editors of the Proceedings
\newcommand{\EditorSignature}{\vspace*{\parskip}\rightline{\@editors}%
  \rightline{Editors of the proceedings}}
%% Fiddle with \thesection, etc. so that they don't include the chapter number
%% Basically go back to the `article' definitions of these
\renewcommand{\thesection}{\@arabic\c@section}
\renewcommand{\thefigure}{\@arabic\c@figure}
\renewcommand{\thetable}{\@arabic\c@table}

\newcommand{\@oldargs}{}
\newcommand{\@newargs}{}
%% add TOC entry for each paper
%\renewcommand{\title}[1]%
%{\renewcommand{\@title}{\rmaa@TitleStyle{#1}}%
%\addtocounter{chapter}{1}\addcontentsline{toc}{chapter}{#1}%
%}

%% Formatting of the table of contents
%% dimensions of various parts
\newlength{\TOCtextwidth}\setlength{\TOCtextwidth}{12.65cm}
\newlength{\TOCpagenumwidth}\setlength{\TOCpagenumwidth}{1.3cm}
\newlength{\TOCentrywidth}\setlength{\TOCentrywidth}{\TOCtextwidth}
\addtolength{\TOCentrywidth}{-\TOCpagenumwidth} 
\newlength{\TOCextramargin}\setlength{\TOCextramargin}{1.25cm}
\newcommand{\TOCbox}[1]{\hspace{\TOCextramargin}\parbox{\TOCtextwidth}{#1}\par}
\newlength{\TOCextramarginx}\setlength{\TOCextramarginx}{\TOCextramargin}\addtolength{\TOCextramarginx}{4.3cm}
\newlength{\TOCtextwidthx}\setlength{\TOCtextwidthx}{\TOCtextwidth}\addtolength{\TOCtextwidthx}{-4.3cm}
\newcommand{\TOCboxx}[1]{\hspace{\TOCextramarginx}\parbox{\TOCtextwidthx}{#1}\par}

\newcommand{\RMAAbannerimage}{rmsc}
\renewcommand\tableofcontents{%
  \renewcommand{\articlename}{CONTENTS}
  \begingroup
  \@begofarticlehook
  \if@twocolumn
     \@restonecoltrue\onecolumn
  \else
     \@restonecolfalse
  \fi
  %% 9pt for journal style
  \if@proceedings\else\small\fi
  \if@newstyle\small\fi %% same for new style
  \if@proceedings\thispagestyle{rmaafrontfirstpage}\else\thispagestyle{empty}\fi
  \shortauthor{Contents}\shorttitle{Contents}
  \vspace*{18mm}
  \if@printbanner
  %% put in the banner header (need to get better image of this)
     \TOCbox{\hspace*{-0.15cm}\includegraphics[width=0.9\TOCtextwidth]{\RMAAbannerimage}}
     \vspace*{8mm}
  %% volume and publication date (should really be in funny font)
     \TOCboxx{\@volumen\\[2mm] \@mes\ \@anyo}
  \else
  %% put in the banner header (need to get better image of this)
     \TOCbox{\hspace*{-0.15cm}\includegraphics[width=0.9\TOCtextwidth]{\RMAAbannerimage}}
     \vspace*{12.5mm}
  %% volume and publication date (should really be in funny font)
     \TOCboxx{~\\[2mm] ~}
  \fi
  \vspace*{16mm}
  %% This is the command that fiddles with the TOC entry so that the
  %% title is flushleft but the authors are flushright
  \renewcommand{\@fiddleauthors}[2]{%
    % first put the paper title
    \renewcommand{\\}{\relax} % temporarily turn off explicit line breaks just in case
    \renewcommand{\altaffilmark}[1]{\relax} % ditto with this
    \renewcommand{\textsuperscript}[1]{\relax} % ditto with this 
    ##1
    % now save the author list in a box
    \sbox\@tempboxa{##2}%
    % check that box is shorter than the total width available for the entry
    \ifdim \wd\@tempboxa <\TOCentrywidth
       % if so, then put it right here, flush against right margin
       \penalty10000\hskip1em plus10em\penalty0\hbox{}\nobreak\hfill\box\@tempboxa
       % Note that we force a space of at least 1em between title and authorlist, 
       % stretchable to 11em since we may have very short titles but still with an 
       % authorlist too long to go on the same line.  The only linebreak allowed is 
       % directly before the authorlist.
    \else
       % if it is too long, we set the original authorlist text in its own parbox
       \par\vskip1ex\parbox[b]{\TOCentrywidth}{\raggedleft ##2}
    \fi
    }
  \@starttoc{toc}%
  \if@restonecol\twocolumn\fi
  \frontlabel{RMAALastPage}
  \clearpage
  \thispagestyle{empty}
  \cleardoublepage % make sure we start on an odd page
  \endgroup

  %%% ALSO, do the HTML TOC
  %% The file that will be used to make the HTML TOC
  \@starttoc{htc}
  %% All the stuff in this file does nothing in the book style
  \addtocontents{htc}{\protect\HyperTOCmeeting{\@conftitle}{\@reunionplace}{\@reuniondate}}
  \addtocontents{htc}{\protect\HyperTOCeditors{\@editors}}
  \addtocontents{htc}{\protect\HyperTOCpreamble}
  \addtocontents{htc}{\protect\HyperTOCEntry{Table of Contents}{}{\articlename}}

  %% and the ADS file
  \ADS@openfile
  }
%% This is now much simpler
\renewcommand\l@chapter[2]{%TOC entry for articles
  \ifnum \c@tocdepth >\m@ne
    \vskip 1.0em \@plus\p@
    \TOCbox{%
      \parbox[b]{\TOCentrywidth}{#1}\parbox[b]{\TOCpagenumwidth}{\hfill#2}
    }
  \fi} 
\renewcommand\l@part[2]{%TOC entry for sessions: centered, no page number
  \ifnum \c@tocdepth >\m@ne
    \vskip 1.5em \@plus1.5em \goodbreak
    \TOCbox{%
      \parbox[b]{\TOCentrywidth}{\hfil#1\hfil}\hfill
    }
  \fi\nobreak} 

%% Formatting theindex
%% strongly encourage page breaks between index entries
\renewcommand{\@idxitem}{\par\penalty-1000\hangindent.45cm} 
\renewcommand{\indexspace}{\par\penalty-1000\vskip 10\p@ \@plus15\p@ \@minus3\p@\relax}
\renewenvironment{theindex}{%
  %% first, make sure we start on an odd page (06 Jun 2000)
  \clearpage\thispagestyle{empty}\cleardoublepage
  %% this looks a bit like an article but is two columns and without abstract, etc
  \renewcommand{\articlename}{INDEX}
  \title{Author Index}\shortauthor{Author Index}\shorttitle{Author Index}
  \listofauthors{}
  \addcontentsline{toc}{chapter}{\protect\@fiddleauthors{\@TOCtitle}{\rmaaTOCauthors}}%
  \addtocontents{htc}{\protect\HyperTOCEntry{\@rawtitle}{\@rawauthors}{\articlename}}
  \pagestyle{rmaa}
  \thispagestyle{rmaafirstpage}
  \@begofarticlehook
  %% 9pt for journal style
  \if@proceedings\else\small\fi
  \if@newstyle\small\fi
  \twocolumn[\@makeminititle]
  %% vertical spacing (see also definition of \@idxitem and \indexspace
  \parskip\z@ \@plus .6\p@\relax % allow quite a lot of leaway here to avoid breaks inside entries
  %% horizontal spacing, breaking and glue parameters (see also the stuff in the definition 
  %% of \indexauthor
  \parindent\z@ 
  \exhyphenpenalty10000 % don't break on explicit hyphens, it always looks bad
  \rightskip=0pt plus.2em % allow a tiny bit of leeway at the ends of lines
  \spaceskip.3333em plus.5em minus.15em % a lot of leeway between words
  \xspaceskip.5em plus.5em minus.4em % this seems to have no effect
  \let\item\@idxitem
  \renewcommand{\\}{\relax} % temporarily turn off explicit line breaks just in case
  \renewcommand{\altaffilmark}[1]{\relax} % ditto with this
  \renewcommand{\textsuperscript}[1]{\relax} % ditto with this
  }{\label{RMAALastPage}\clearpage} % clearpage necessary so that we know about \shortauthor



%% end definitions for book styles
\fi

%    \end{macrocode}
% In order to do away with plates, we sometimes want a figure to
% appear in color only in the online version. The following command
% can bve used in a figure caption to typeset a note to this effect. 
%    \begin{macrocode}
\newcommand\ColorOnlineOnly{%
  \textsc{Note: this figure is available in color in the electronic
    version of this article, obtainable from}
  \texttt{http://\discretionary{}{}{}www.\discretionary{}{}{}astroscu.\discretionary{}{}{}unam.\discretionary{}{}{}mx/\discretionary{}{}{}\~{}rmaa/}.}
%% Plates
%% Define 2 counters: plate is the number in the whole document, plateloop is used 
%% internally by \outputplates
\newcounter{plate}
\newcounter{plateloop}
%% Currently limited to 15 plates in document
%% when we do a label, we need to write to \newlabel entries to the .aux file, one 
%% that says RMAAPlateX so that it can be referred to in \ouputplates, and one that
%% says articlename:RMAAPlateX so that the dvips program can find it and asociate it
%% with the right article
\newcommand{\@platelabeli}{\oldlabel{RMAAPlate1}\label{RMAAPlate1}}
\newcommand{\@platelabelii}{\oldlabel{RMAAPlate2}\label{RMAAPlate2}}
\newcommand{\@platelabeliii}{\oldlabel{RMAAPlate3}\label{RMAAPlate3}}
\newcommand{\@platelabeliv}{\oldlabel{RMAAPlate4}\label{RMAAPlate4}}
\newcommand{\@platelabelv}{\oldlabel{RMAAPlate5}\label{RMAAPlate5}}
\newcommand{\@platelabelvi}{\oldlabel{RMAAPlate6}\label{RMAAPlate6}}
\newcommand{\@platelabelvii}{\oldlabel{RMAAPlate7}\label{RMAAPlate7}}
\newcommand{\@platelabelviii}{\oldlabel{RMAAPlate8}\label{RMAAPlate8}}
\newcommand{\@platelabelix}{\oldlabel{RMAAPlate9}\label{RMAAPlate9}}
\newcommand{\@platelabelx}{\oldlabel{RMAAPlate10}\label{RMAAPlate10}}
\newcommand{\@platelabelxi}{\oldlabel{RMAAPlate11}\label{RMAAPlate11}}
\newcommand{\@platelabelxii}{\oldlabel{RMAAPlate12}\label{RMAAPlate12}}
\newcommand{\@platelabelxiii}{\oldlabel{RMAAPlate13}\label{RMAAPlate13}}
\newcommand{\@platelabelxiv}{\oldlabel{RMAAPlate14}\label{RMAAPlate14}}
\newcommand{\@platelabelxv}{\oldlabel{RMAAPlate15}\label{RMAAPlate15}}
\newcommand{\@outputplatelabeli}{\oldlabel{RMAAOutputPlate1}}
\newcommand{\@outputplatelabelii}{\oldlabel{RMAAOutputPlate2}}
\newcommand{\@outputplatelabeliii}{\oldlabel{RMAAOutputPlate3}}
\newcommand{\@outputplatelabeliv}{\oldlabel{RMAAOutputPlate4}}
\newcommand{\@outputplatelabelv}{\oldlabel{RMAAOutputPlate5}}
\newcommand{\@outputplatelabelvi}{\oldlabel{RMAAOutputPlate6}}
\newcommand{\@outputplatelabelvii}{\oldlabel{RMAAOutputPlate7}}
\newcommand{\@outputplatelabelviii}{\oldlabel{RMAAOutputPlate8}}
\newcommand{\@outputplatelabelix}{\oldlabel{RMAAOutputPlate9}}
\newcommand{\@outputplatelabelx}{\oldlabel{RMAAOutputPlate10}}
\newcommand{\@outputplatelabelxi}{\oldlabel{RMAAOutputPlate11}}
\newcommand{\@outputplatelabelxii}{\oldlabel{RMAAOutputPlate12}}
\newcommand{\@outputplatelabelxiii}{\oldlabel{RMAAOutputPlate13}}
\newcommand{\@outputplatelabelxiv}{\oldlabel{RMAAOutputPlate14}}
\newcommand{\@outputplatelabelxv}{\oldlabel{RMAAOutputPlate15}}
\newcommand{\@platerefi}{\oldpageref{RMAAPlate1}}
\newcommand{\@platerefii}{\oldpageref{RMAAPlate2}}
\newcommand{\@platerefiii}{\oldpageref{RMAAPlate3}}
\newcommand{\@platerefiv}{\oldpageref{RMAAPlate4}}
\newcommand{\@platerefv}{\oldpageref{RMAAPlate5}}
\newcommand{\@platerefvi}{\oldpageref{RMAAPlate6}}
\newcommand{\@platerefvii}{\oldpageref{RMAAPlate7}}
\newcommand{\@platerefviii}{\oldpageref{RMAAPlate8}}
\newcommand{\@platerefix}{\oldpageref{RMAAPlate9}}
\newcommand{\@platerefx}{\oldpageref{RMAAPlate10}}
\newcommand{\@platerefxi}{\oldpageref{RMAAPlate11}}
\newcommand{\@platerefxii}{\oldpageref{RMAAPlate12}}
\newcommand{\@platerefxiii}{\oldpageref{RMAAPlate13}}
\newcommand{\@platerefxiv}{\oldpageref{RMAAPlate14}}
\newcommand{\@platerefxv}{\oldpageref{RMAAPlate15}}
\newcounter{figplatei}
\newcounter{figplateii}
\newcounter{figplateiii}
\newcounter{figplateiv}
\newcounter{figplatev}
\newcounter{figplatevi}
\newcounter{figplatevii}
\newcounter{figplateviii}
\newcounter{figplateix}
\newcounter{figplatex}
\newcounter{figplatexi}
\newcounter{figplatexii}
\newcounter{figplatexiii}
\newcounter{figplatexiv}
\newcounter{figplatexv}
%% Printed output should use the counter `plate'
\renewcommand{\theplate}{\@arabic\c@plate}
\renewcommand{\theplateloop}{\@arabic\c@plateloop}
%% Command to output all the saved plates in this paper
\newcommand{\outputplates}{%
  \setcounter{plateloop}{0}
  \if@book\else\@endofarticlehook\fi % not entirely sure why, but this
                                     % seems to be needed
  \clearpage % first output the last page of text
  % now set pagestyle for blank even pages and make sure we start on an odd page
  \pagestyle{empty}
  \cleardoublepage
  \onecolumn
  \parindent\z@
  % first put in an extra sheet
  % the main purpose of this is so there is always a blank even page before 
  % each plate for filling in the PS file if the article has an odd number
  % of pages (see splitbook.sh)
  \null\vfil\centerline{PLATES}\vfil\clearpage\pagestyle{rmaaplate}\cleardoublepage
  % output each plate in turn
  \@whilenum\value{plateloop}<\value{plate}\do{%
    \addtocounter{plateloop}{1}
    % restore the figure counter from \@figcounti, \@figcountii, etc
    \setcounter{figure}{\value{figplate\roman{plateloop}}}
    % spew out the plate (in its own group for safety)
    {\csname @plate\roman{plateloop}\endcsname\par}
    }
  \pagestyle{rmaa}
  % reset the page number as if nothing had happened
  \addtocounter{page}{-\value{plate}}
  \addtocounter{page}{-\value{plate}}
  }
\newcommand\figcaption{%
  \def\@captype{figure}
  \renewcommand{\@makecaption}{\@makefigcap}
  \caption}
\newcommand\tabcaption{%
  \def\@captype{table}
  \renewcommand{\@makecaption}{\@maketabcap}
  \caption}

\newcommand{\plate}[2][1]{%
  %% User command to make a plate page: 
  %% Mandatory argument is what goes on the plate - this should not 
  %% include any floating environments (e.g. figure). Instead, \figcaption 
  %% should be used to do the captions. 
  %% Optional argument is the number of figures (i.e. captions) on the plate
  %% page (default is 1). This is necessary for incrementing the figure counter. 
  %%
  %% Note that we have to use a \gdef rather than an \xdef when we 
  %% save the plate output in a command. This means that in the book styles when it 
  %% comes to spewing out the plate, we will have forgotten about any user macros 
  %% that may have been defined. This can be got round by repeating the macros inside the plate.
  %% (Presumably, one would use a \newcommand rather than \renewcommand if they were new.)
  % increment global and local counters
  \addtocounter{plate}{1}
  % write appropriate label to .aux file
  \csname @platelabel\roman{plate}\endcsname
  % define a global macro called \@platei, \@plateii, \@plateiii, \@plateiv, ...
  \expandafter\gdef\csname @plate\roman{plate}\endcsname{%
    \csname @outputplatelabel\roman{plateloop}\endcsname
    \null\vfil #2 \vfil
    \cleardoublepage
    }
  % save the current shortauthor and shorttitle
  % need to be \xdef to take current values of \@shortauthor, etc 
  \expandafter\xdef\csname @shorta\roman{plate}\endcsname{\@shortauthor}
  \expandafter\xdef\csname @shortt\roman{plate}\endcsname{\@shorttitle}
  % save the current value of the figure counter for later use in \outputplates
  \setcounter{figplate\roman{plate}}{\value{figure}}
  % increment the figure counter according to the optional argument
  \addtocounter{figure}{#1}
  }
%% page style for plate pages 
\newcommand{\ps@rmaaplate}{%
  % Plate X (see page Y) in the foot
  \renewcommand\@oddfoot{\hfil\csname @shorta\roman{plateloop}\endcsname\ 
    (see page~\csname @plateref\roman{plateloop}\endcsname)\hfil}%
  \renewcommand\@evenfoot{\relax}%
  % nothing in even heads (blank pages)
  \renewcommand\@evenhead{\relax}
  % short title and page num in the odd heads
  \renewcommand\@oddhead{\hfil\csname @shortt\roman{plateloop}\endcsname%
    \hfil\llap{PLATE~\theplateloop}}
  }





%% Things for the User Guide
\if@guide
%% stolen from doc.sty
\def\MakeShortVerb#1{%
  \expandafter\ifx\csname cc\string#1\endcsname\relax
  \@shortvrbinfo{Made }{#1}%
  \add@special{#1}%
  \expandafter
  \xdef\csname cc\string#1\endcsname{\the\catcode`#1}%
  \begingroup
  \catcode`\~\active  \lccode`\~`#1%
  \lowercase{%
    \global\expandafter\let
    \csname ac\string#1\endcsname~%
    \gdef~{\verb~}}%
  \endgroup
  \global\catcode`#1\active
  \else
  \@shortvrbinfo\@empty{#1 already}%
  \fi}
\def\DeleteShortVerb#1{%
  \expandafter\ifx\csname cc\string#1\endcsname\relax
  \else
  \@shortvrbinfo{Deleted }{#1 as}%
  \rem@special{#1}%
  \global\catcode`#1\csname cc\string#1\endcsname
  \global \expandafter\let \csname cc\string#1\endcsname \relax
  \ifnum\catcode`#1=\active
  \begingroup
  \catcode`\~\active   \lccode`\~`#1%
  \lowercase{%
    \global\expandafter\let\expandafter~%
    \csname ac\string#1\endcsname}%
  \endgroup \fi \fi}
\def\@shortvrbinfo#1#2{%
  \PackageInfo{doc}{%
    #1\expandafter\@gobble\string#2 a short reference
    for \string\verb}}
\def\add@special#1{%
  \rem@special{#1}%
  \expandafter\gdef\expandafter\dospecials\expandafter
  {\dospecials \do #1}%
  \expandafter\gdef\expandafter\@sanitize\expandafter
  {\@sanitize \@makeother #1}}
\def\rem@special#1{%
  \def\do##1{%
    \ifnum`#1=`##1 \else \noexpand\do\noexpand##1\fi}%
  \xdef\dospecials{\dospecials}%
  \begingroup
  \def\@makeother##1{%
    \ifnum`#1=`##1 \else \noexpand\@makeother\noexpand##1\fi}%
  \xdef\@sanitize{\@sanitize}%
  \endgroup}
%% define | as shortcut for \verbatim
\MakeShortVerb{\|}
%% comand to specify arguments to LaTeX commands
\newcommand{\rmArg}[1]{%
  \ensuremath{\langle}{\scriptsize \textsf{\MakeUpperCase{#1}}}\ensuremath{\rangle}}
%\newcommand{\rmArg}[1]{\ensuremath{\langle}\textsc{#1}\ensuremath{\rangle}}
\providecommand{\BibTeX}{{\rmfamily B\kern-.05em%
    \textsc{i\kern-.025em b}\kern-.08em%
    T\kern-.1667em\lower.7ex\hbox{E}\kern-.125emX}}
%% Environment for code examples
\newenvironment{Code}{%
  \begin{list}{}{\setlength{\leftmargin}{0.05\columnwidth}} 
  \item[]}{\end{list}}
\fi
%</package>
%    \end{macrocode}

% \section{User commands for astronomical articles}
% \label{sec:user-commands}
% Some random definitions. These go both into the main class file and
% also into |adshtml.sty|.
% 
% Definitions of the journal name in different incarnations.  
%    \begin{macrocode}
%<*package|adshtml-style>
\newcommand{\RMAA}{RMxAA}
\newcommand{\RMAAlong}{Revista Mexicana de Astronom\'{\i}a y Astrof\'{\i}sica}
\newcommand{\RMAASC}{RMxAC}
\newcommand{\RMAASClong}{\RMAA\ (Serie de Conferencias)}
%    \end{macrocode}
% \begin{macro}{\ibidrule}
% This puts a horizontal rule with a period after it. It is used when
% the author list of a reference is identical to the preceding one. 
%    \begin{macrocode}
\newcommand{\ibidrule}{\rule[-0.03cm]{4.7em}{.01cm}.\ }
%    \end{macrocode}
% \end{macro}
% \begin{macro}{\ThisVolume}
% References to other articles in the same volume can be coded as,
% e.g., |\bibitem{} Bloggs, J. \ThisVolume{bloggs}|.
%    \begin{macrocode}
\newcommand{\ThisVolume}[1]{%
  \@year, \if@proceedings\RMAASC\else\RMAA\fi, \@volume,
  \if@book\oldpageref{#1:RMAAFirstPage}\else 000\fi\ (this volume)}
%    \end{macrocode}
% \end{macro}
% Other stuff that may be useful in ref lists.
%    \begin{macrocode}
\newcommand{\etal}{et al.\@}
%    \end{macrocode}
% Various new commands from AAS macros
%    \begin{macrocode}
\def\acknowledgments{\vskip 3ex plus .8ex minus .4ex}
\let\acknowledgements=\acknowledgments                  % second spelling
\let\jnl@style=\rm
\def\ref@jnl#1{{\jnl@style#1}}
\DeclareRobustCommand\aj{\ref@jnl{AJ}}                   
\DeclareRobustCommand\araa{\ref@jnl{ARA\&A}}             
\DeclareRobustCommand\apj{\ref@jnl{ApJ}}                 
\DeclareRobustCommand\apjl{\ref@jnl{ApJ}}                
\DeclareRobustCommand\apjs{\ref@jnl{ApJS}}               
\DeclareRobustCommand\ao{\ref@jnl{Appl.~Opt.}}           
\DeclareRobustCommand\apss{\ref@jnl{Ap\&SS}}             
\DeclareRobustCommand\aap{\ref@jnl{A\&A}}                
\DeclareRobustCommand\aapr{\ref@jnl{A\&A~Rev.}}          
\DeclareRobustCommand\aaps{\ref@jnl{A\&AS}}              
\DeclareRobustCommand\azh{\ref@jnl{AZh}}                 
\DeclareRobustCommand\baas{\ref@jnl{BAAS}}               
\DeclareRobustCommand\jrasc{\ref@jnl{JRASC}}             
\DeclareRobustCommand\memras{\ref@jnl{MmRAS}}            
\DeclareRobustCommand\mnras{\ref@jnl{MNRAS}}             
\DeclareRobustCommand\pra{\ref@jnl{Phys.~Rev.~A}}        
\DeclareRobustCommand\prb{\ref@jnl{Phys.~Rev.~B}}        
\DeclareRobustCommand\prc{\ref@jnl{Phys.~Rev.~C}}        
\DeclareRobustCommand\prd{\ref@jnl{Phys.~Rev.~D}}        
\DeclareRobustCommand\pre{\ref@jnl{Phys.~Rev.~E}}        
\DeclareRobustCommand\prl{\ref@jnl{Phys.~Rev.~Lett.}}    
\DeclareRobustCommand\pasp{\ref@jnl{PASP}}               
\DeclareRobustCommand\pasj{\ref@jnl{PASJ}}               
\DeclareRobustCommand\qjras{\ref@jnl{QJRAS}}             
\DeclareRobustCommand\skytel{\ref@jnl{S\&T}}             
\DeclareRobustCommand\solphys{\ref@jnl{Sol.~Phys.}}      
\DeclareRobustCommand\sovast{\ref@jnl{Soviet~Ast.}}      
\DeclareRobustCommand\ssr{\ref@jnl{Space~Sci.~Rev.}}     
\DeclareRobustCommand\zap{\ref@jnl{ZAp}}                 
\DeclareRobustCommand\nat{\ref@jnl{Nature}}              
\DeclareRobustCommand\iaucirc{\ref@jnl{IAU~Circ.}}
\DeclareRobustCommand\aplett{\ref@jnl{Astrophys.~Lett.}}
\DeclareRobustCommand\apspr{\ref@jnl{Astrophys.~Space~Phys.~Res.}}
\DeclareRobustCommand\bain{\ref@jnl{Bull.~Astron.~Inst.~Netherlands}}
\DeclareRobustCommand\fcp{\ref@jnl{Fund.~Cosmic~Phys.}}
\DeclareRobustCommand\gca{\ref@jnl{Geochim.~Cosmochim.~Acta}}
\DeclareRobustCommand\grl{\ref@jnl{Geophys.~Res.~Lett.}}
\DeclareRobustCommand\jcp{\ref@jnl{J.~Chem.~Phys.}}      
\DeclareRobustCommand\jgr{\ref@jnl{J.~Geophys.~Res.}}    
\DeclareRobustCommand\jqsrt{\ref@jnl{J.~Quant.~Spec.~Radiat.~Transf.}}
\DeclareRobustCommand\memsai{\ref@jnl{Mem.~Soc.~Astron.~Italiana}}
\DeclareRobustCommand\nphysa{\ref@jnl{Nucl.~Phys.~A}}
\DeclareRobustCommand\physrep{\ref@jnl{Phys.~Rep.}}
\DeclareRobustCommand\physscr{\ref@jnl{Phys.~Scr}}
\DeclareRobustCommand\planss{\ref@jnl{Planet.~Space~Sci.}}
\DeclareRobustCommand\procspie{\ref@jnl{Proc.~SPIE}}
\let\astap=\aap
\let\apjlett=\apjl
\let\apjsupp=\apjs
\let\applopt=\ao
\def\phn{\phantom{0}}
\def\phd{\phantom{.}}
\def\phs{\phantom{\ensuremath{-}}}
\def\phm#1{\phantom{#1}}
%    \end{macrocode}
% All these are redone as robust and using more modern commands where
% possible. The robustness is so they will not get expanded if used in
% abstracts and titles when these are passed to the ADS file. 
% \changes{3.16}{2003/04/07}{
%  Added \cs{Color} command for doing photometric colors.}
%    \begin{macrocode}
%%\def\sun{\hbox{$\odot$}}
\DeclareRobustCommand\sun{\ensuremath{\odot}}
%%\def\earth{\hbox{$\oplus$}}
\DeclareRobustCommand\earth{\ensuremath{\oplus}}
%%\def\lesssim{\mathrel{\hbox{\rlap{\hbox{\lower4pt\hbox{$\sim$}}}\hbox{$<$}}}}
%%\def\gtrsim{\mathrel{\hbox{\rlap{\hbox{\lower4pt\hbox{$\sim$}}}\hbox{$>$}}}}
\DeclareRobustCommand\lesssim{\mathrel{\hbox{\smaller\smaller\rlap{\hbox{\lower3pt\hbox{$\sim$}}}\raise1.5pt\hbox{$<$}}}}
\DeclareRobustCommand\gtrsim{\mathrel{\hbox{\smaller\smaller\rlap{\hbox{\lower3pt\hbox{$\sim$}}}\raise1.5pt\hbox{$>$}}}}
%%\def\lta{\mathrel{\spose{\lower 3pt\hbox{$\mathchar"218$}}
%%     \raise 2.0pt\hbox{$\mathchar"13C$}}}
%%\def\gta{\mathrel{\spose{\lower 3pt\hbox{$\mathchar"218$}}
%%     \raise 2.0pt\hbox{$\mathchar"13E$}}}
\DeclareRobustCommand\sq{\hbox{\rlap{$\sqcap$}$\sqcup$}}
%%\def\arcdeg{\hbox{$^\circ$}}
%%\def\arcmin{\hbox{$^\prime$}}
%%\def\arcsec{\hbox{$^{\prime\prime}$}}
\DeclareRobustCommand\arcdeg{\ensuremath{^\circ}}
\DeclareRobustCommand\arcmin{\ensuremath{^\prime}}
\DeclareRobustCommand\arcsec{\ensuremath{^{\prime\prime}}}
%%\def\fd{\hbox{$.\!\!^{\rm d}$}}
%%\def\fh{\hbox{$.\!\!^{\rm h}$}}
%%\def\fm{\hbox{$.\!\!^{\rm m}$}}
%%\def\fs{\hbox{$.\!\!^{\rm s}$}}
%%\def\fdg{\hbox{$.\!\!^\circ$}}
%%\def\farcm{\hbox{$.\mkern-4mu^\prime$}}
%%\def\farcs{\hbox{$.\!\!^{\prime\prime}$}}
%%\def\fp{\hbox{$.\!\!^{\scriptscriptstyle\rm p}$}}
\DeclareRobustCommand\fd{\ensuremath{.\!\!^\mathrm{d}}}
\DeclareRobustCommand\fh{\ensuremath{.\!\!^\mathrm{h}}}
\DeclareRobustCommand\fm{\ensuremath{.\!\!^\mathrm{m}}}
\DeclareRobustCommand\fs{\ensuremath{\mkern0.5mu.\mkern-5mu^\mathrm{s}\mkern-1mu}}
\DeclareRobustCommand\fdg{\ensuremath{.\!\!^\circ}}
\DeclareRobustCommand\farcm{\ensuremath{.\mkern-4mu^\prime}}
\DeclareRobustCommand\farcs{\ensuremath{\mkern1mu.\mkern-6mu^{\prime\prime}\mkern-1.5mu}}
%%\DeclareRobustCommand\farcs{\ensuremath{.\!\!^{\prime\prime}}}
\DeclareRobustCommand\fp{\ensuremath{.\!\!^{\scriptscriptstyle\rm p}}}
%%\def\micron{\hbox{$\mu$m}}
\DeclareRobustCommand\micron{\ensuremath{\mu\mathrm{m}}}
\let\la=\lesssim                        
\let\ga=\gtrsim
%%\def\case#1#2{\hbox{$\frac{#1}{#2}$}}
%%\def\slantfrac#1#2{\hbox{$\,^#1\!/_#2$}}
%%\def\onehalf{\slantfrac{1}{2}}
%%\def\onethird{\slantfrac{1}{3}}
%%\def\twothirds{\slantfrac{2}{3}}
%%\def\onequarter{\slantfrac{1}{4}}
%%\def\threequarters{\slantfrac{3}{4}}
\DeclareRobustCommand\case[2]{\ensuremath{\frac{#1}{#2}}}
\DeclareRobustCommand\slantfrac[2]{\ensuremath{\,^#1\!/_#2}}
\DeclareRobustCommand\onehalf{\slantfrac{1}{2}}
\DeclareRobustCommand\onethird{\slantfrac{1}{3}}
\DeclareRobustCommand\twothirds{\slantfrac{2}{3}}
\DeclareRobustCommand\onequarter{\slantfrac{1}{4}}
\DeclareRobustCommand\threequarters{\slantfrac{3}{4}}
%%\def\ubvr{\hbox{$U\!BV\!R$}}            
%%\def\ub{\hbox{$U\!-\!B$}}               
%%\def\bv{\hbox{$B\!-\!V$}}               
%%\def\vr{\hbox{$V\!-\!R$}}               
%%\def\ur{\hbox{$U\!-\!R$}}
%    \end{macrocode}
% New commands for ``photometric'' colors 
%    \begin{macrocode}
\DeclareRobustCommand\ubvr{\ensuremath{U\!BV\!R}}            
\DeclareRobustCommand\Color[2]{%
  \ensuremath{%
    #1\!\ifx#1V\!\fi-\!#2
  }}
\DeclareRobustCommand\ub{\Color{U}{B}}
\DeclareRobustCommand\bv{\Color{B}{V}}
\DeclareRobustCommand\vr{\Color{V}{R}}
\DeclareRobustCommand\ur{\Color{U}{R}}
%% New version of the \ion command, that actually works! 
%    \end{macrocode}
% \begin{macro}{\ion}
% This version of the |\ion| command allows you to write either
% |\ion{N}{4}| or |\ion{N}{iv}| in order to get N~IV. 
%    \begin{macrocode}
\DeclareRobustCommand\ion[2]{%
  \mbox{#1\kern0.2em%
    \smaller\rmfamily%
    \edef\@tempa{\@car#2\@nil}%
    \ifcat1\@tempa%
    \@Roman{#2}%
    \else%
    \uppercase{#2}%
    \fi}}
\newcounter{IonStage}
\DeclareRobustCommand\plainion[2]{\setcounter{IonStage}{#2}#1 \Roman{IonStage}}
%%\def\ion#1#2{\ensuremath{\mathrm{#1}\;}{\protect\small\rm\@Roman{#2}}}
%    \end{macrocode}
% \end{macro}
%    \begin{macrocode}
\DeclareRobustCommand\nodata{$\cdots$}
%    \end{macrocode}
% These are needed by the |\ion| command. Isn't there an easier way of
% doing this?
%    \begin{macrocode}
\let\@small=\small
\let\@normalsize=\normalsize
\let\@footnotesize=\footnotesize
\let\@scriptsize=\scriptsize
\let\@tiny=\tiny
\let\@large=\large
\let\@Large=\Large
\let\@LARGE=\LARGE
\let\@huge=\huge
\let\@Huge=\Huge
\newcommand{\smaller}{\small}
\renewcommand{\normalsize}{\renewcommand{\smaller}{\small}\@normalsize}
\renewcommand{\small}{\renewcommand{\smaller}{\footnotesize}\@small}
\renewcommand{\footnotesize}{\renewcommand{\smaller}{\scriptsize}\@footnotesize}
\renewcommand{\scriptsize}{\renewcommand{\smaller}{\tiny}\@scriptsize}
\renewcommand{\tiny}{\renewcommand{\smaller}{\tiny}\@tiny}
\renewcommand{\large}{\renewcommand{\smaller}{\normalsize}\@large}  
\renewcommand{\Large}{\renewcommand{\smaller}{\large}\@Large}  
\renewcommand{\LARGE}{\renewcommand{\smaller}{\Large}\@LARGE}  
\renewcommand{\huge}{\renewcommand{\smaller}{\LARGE}\@huge}  
\renewcommand{\Huge}{\renewcommand{\smaller}{\huge}\@Huge}  
%</package|adshtml-style>
%    \end{macrocode}


% \section{Other files}

% \subsection{The driver and make files}
% \label{sec:driver_make}
% These are put here for completeness, although they still need to be
% distributed separately since otherwise the user won't be able to
% unpack this file into its constituent parts.
% 
% The |Makefile| has two principal rules: |docs|, and |macros|. The
% |docs| rule produces a PDF and gzipped PS file of the documentation.
% I still haven't perfected the combination of rules required to get
% it to run \LaTeX{} enough times in all circumstances. The |macros|
% rule runs \LaTeX{} on |rmaa.ins| in order to unpack |rmaa.cls| and
% all the other little files (including the |Makefile| and the
% |rmaa.ins| files themselves\dots 
% 
% First set up lists of various types of file. 
%    \begin{macrocode}
%<*makefile>
DOCFILES=rmaa.pdf
MACRO_LTX_FILES=rmaa.cls rmaacite.sty adshtml.sty
SHELLSCRIPTS=makeads.sh
MACRO_OTHER_FILES=adshtml.perl sample.latex2html-init
MAKEFILE=Makefile
MACROFILES=$(MACRO_LTX_FILES) $(SHELLSCRIPTS) $(MACRO_OTHER_FILES)         
SKELETONFILES=ADS.tex 
%    \end{macrocode}
% Install the files read only into the package directory, specified via a path relative to the source directory. \changes{4.0}{2019/11/14}{Reorganize folder layout so that hardly anything is in the top level.  That way, things are clearer on github.}
%    \begin{macrocode}
INSTALL=install 
INSTALL_DIR=../packages
%    \end{macrocode}
% The default rule is to make everything. Installation must be
% explicitly asked for.
%    \begin{macrocode}
all: docs macros

install: install-docs install-macros

%    \end{macrocode}
% Most files are installed read-only to prevent editing. Skeleton
% files must be writeable. Must be careful not to overwrite the
% Makefile in the parent directory. 
%    \begin{macrocode}
install-macros: macros
        $(INSTALL) -m 444 $(MACRO_LTX_FILES) $(MACRO_OTHER_FILES) $(INSTALL_DIR)
        $(INSTALL) -m 555 $(SHELLSCRIPTS) $(INSTALL_DIR)
        $(INSTALL) -m 644 $(SKELETONFILES) $(INSTALL_DIR)

install-docs: docs
        $(INSTALL) rmaa.pdf $(INSTALL_DIR)/rm-fulldocs.pdf

%    \end{macrocode}
% Make the PDF and gzipped PS files. 
%    \begin{macrocode}
docs: $(DOCFILES)

%    \end{macrocode}
% Make everything else if either the |.dtx| or |.ins| files have
% changed. 
%    \begin{macrocode}
macros: rmaa.dtx rmaa.ins
%    \end{macrocode}
% Ensure that all the macro files can be written to then recreate
% them. The |touch| is there in case any of the files don't exist
% yet. 
%    \begin{macrocode}
        touch $(MACROFILES) $(MAKEFILE)
        chmod u+w $(MACROFILES) $(MAKEFILE) rmaa.ins
        latex "\input{rmaa.ins}"
%    \end{macrocode}
% \TeX{} automatically adds the |.tex| extension to extensionless
% files. We fix this and take the opportunity to convert
% back to TABs all that it has lovingly converted to spaces. 
%    \begin{macrocode}
        cp -f $(MAKEFILE) $(MAKEFILE).BCK && unexpand $(MAKEFILE).tex >| $(MAKEFILE) && rm $(MAKEFILE).tex
%    \end{macrocode}
% Make the shell scripts executable and everything unwriteable (so
% they don't get hand-edited by mistake). However, undo this for the
% files that the user is \emph{supposed} to modify. 
%    \begin{macrocode}
        chmod a+x $(SHELLSCRIPTS)
        chmod a-w $(MACROFILES) $(MAKEFILE)
        chmod u+w $(SKELETONFILES)

%    \end{macrocode}
% We now use PDF\LaTeX{} for this.
%    \begin{macrocode}
rmaa.idx: rmaa.dtx
        @echo "REGENERATING .idx FILE"
        pdflatex "\nonstopmode\input{rmaa.dtx}"

rmaa.aux: rmaa.dtx rmaa.ind rmaa.gls
        @echo "REGENERATING .aux FILE"
        pdflatex "\nonstopmode\input{rmaa.dtx}"

rmaa.ind: rmaa.idx
        @echo "REGENERATING INDEX"
        makeindex -s gind.ist rmaa

rmaa.gls: rmaa.glo
        @echo "REGENERATING CHANGELOG"
        makeindex -s gglo.ist -o rmaa.gls rmaa.glo

rmaa.pdf: rmaa.aux rmaa.ind rmaa.gls
        @echo "REGENERATING .pdf FILE"
        pdflatex "\nonstopmode\input{rmaa.dtx}"

clean: 
        rm -f rmaa.aux rmaa.dvi rmaa.ind rmaa.gls
%</makefile>
%    \end{macrocode}

% The |rmaa.ins| file uses the \textsf{docstrip} package to split out
% all the files and strip off this documentation.
%    \begin{macrocode}
%<*ins>
%%
%% Installation batch file for the RMAA LaTeX package
%%
\def\batchfile{rmaa.ins}
\input docstrip.tex
\preamble

PLEASE DO NOT EDIT THIS FILE! It is generated automatically from
the file rmaa.dtx using the accompanying Makefile and rmaa.ins 
Any hand-edited changes will be lost

A LaTeX document class for the RMxAA

Author: William Henney (w.henney@irya.unam.mx)

Latest version available from: http://www.irya.unam.mx/rmaa

Copyright (C) 1999-2019 Instituto de Astronomia, UNAM, Mexico

Permission is hereby granted, free of charge, to any person obtaining
a copy of this software and associated documentation files (the
"Software"), to deal in the Software without restriction, including
without limitation the rights to use, copy, modify, merge, publish,
distribute, sublicense, and/or sell copies of the Software, and to
permit persons to whom the Software is furnished to do so, subject to
the following conditions:

The above copyright notice and this permission notice shall be
included in all copies or substantial portions of the Software.

THE SOFTWARE IS PROVIDED "AS IS", WITHOUT WARRANTY OF ANY KIND,
EXPRESS OR IMPLIED, INCLUDING BUT NOT LIMITED TO THE WARRANTIES OF
MERCHANTABILITY, FITNESS FOR A PARTICULAR PURPOSE AND
NONINFRINGEMENT. IN NO EVENT SHALL THE AUTHORS OR COPYRIGHT HOLDERS BE
LIABLE FOR ANY CLAIM, DAMAGES OR OTHER LIABILITY, WHETHER IN AN ACTION
OF CONTRACT, TORT OR OTHERWISE, ARISING FROM, OUT OF OR IN CONNECTION
WITH THE SOFTWARE OR THE USE OR OTHER DEALINGS IN THE SOFTWARE.
\endpreamble
%    \end{macrocode}
% Clobber files without asking
%    \begin{macrocode}
\askforoverwritefalse
%    \end{macrocode}
% Make the \LaTeX{} macro files (class and styles)
%    \begin{macrocode}
\generate{% 
  \file{rmaa.cls}{\from{rmaa.dtx}{package}}
  \file{rmaa.ins}{\from{rmaa.dtx}{ins}}
  \file{rmaacite.sty}{\from{rmaa.dtx}{rmaacite}}
  \file{adshtml.sty}{\from{rmaa.dtx}{adshtml-style}}
  \file{ADS.tex}{\from{rmaa.dtx}{ads-container}}
  }
%    \end{macrocode}
% Make the other non-\LaTeX{} files. These are done separately since
% they can't use `\%\%' as a comment character for the
% preamble/postamble (and Perl files can't have the postamble I don't
% think). 
%    \begin{macrocode}
\generate{\nopreamble\nopostamble% 
  \file{Makefile}{\from{rmaa.dtx}{makefile}}
  \file{adshtml.perl}{\from{rmaa.dtx}{adshtml-perl}}
  \file{makeads.sh}{\from{rmaa.dtx}{ads-shell-script}}
  \file{sample.latex2html-init}{\from{rmaa.dtx}{latex2html-init}}
  }

%% Local Variables: 
%% mode: LaTeX
%% mode: Font-lock
%% End:
%</ins>
%    \end{macrocode}

% \subsection{Extended citation syntax}
% Modified from a file called mnrascite.sty
% I no longer remember where I stole the original version of this from
% 
% Alterations to the cite commands to make them do more sensible things.
% N.B. the bibliography style should produce entries which have labels where
% the date is delimited by angle brackets '<>' for these macros to work
% properly. They will not give very sensible error messages if you do not do 
% this properly.
% 
% A new command |\scite{key}| is defined which produces a citation which is 
% suitable for insertion in a sentence. i.e., it does not put brackets around 
% the whole citation, only the date. Also |\pcite{key}| produce output with no
% parentheses whatsoever. 
% 
% Should redefine |\@bibitem| really to have any hope of compatibility
% the command |\lbibcite| is defined to deal with the citations.
%    \begin{macrocode}
%<*rmaacite>
\def\bibitem[#1<#2>]#3{\par\hangindent=1.5em\hangafter=1
  \if@filesw
  {\def\protect##1{\string ##1\space}
    \immediate\write\@auxout{\string\lbibcite{#3}{#1}{#2}}}
  \fi}
\@redefinebibitem % what it says
%    \end{macrocode}
% Note that the |\lbibcite| defines two types of citation one contains the names,
% the other the year
%    \begin{macrocode}
\def\lbibcite#1#2#3{\global\@namedef{b@#1}{#2}
                    \global\@namedef{by@#1}{#3}}

\def\@bcite#1#2{({#1\if@tempswa , #2\fi})}
\def\@pcite#1#2{#1\if@tempswa , #2\fi}
%    \end{macrocode}
% This declaration contains what is put between the name and the date
% \changes{3.11}{2003/01/31}{Moved the \cs{articlename} to after the key for consistency
% with the new natbib stuff.}
%    \begin{macrocode}
\def\@brcite{\ (}
\def\@ebrcite{)}
\def\@fmttop#1{\@fmtcite[\csname #1\endcsname]}
\def\@citex[#1]#2{%
%%  \if@filesw\immediate\write\@auxout{\string\citation{#2}}\fi
%% WJH added 16 Jan 2002
  \@for\@citeb:=#2\do{%
    \edef\@citeb{\expandafter\@firstofone\@citeb:\articlename}%
    \if@filesw\immediate\write\@auxout{\string\citation{\@citeb}}\fi}%
%% to here
  \def\@citea{}%
  \@cite{\@for\@citeb:=#2\do
    {%
      \edef\@citeb{\expandafter\@firstofone\@citeb:\articlename}%
      \@citea\def\@citea{;\ }\@ifundefined
      {b@\@citeb}{{\bf ?}\@warning {Citation `\@citeb' on page \thepage \space undefined}}%
      {\csname b@\@citeb\endcsname\@citebro\csname by@\@citeb\endcsname\@citebrc}}}%
  {#1}}

\def\cite{\@ifnextchar [{\let\@citebro=\space
                          \let\@citebrc=\relax
                          \let\@cite=\@bcite\@tempswatrue\@citex}
                        {\let\@citebro=\space
                          \let\@citebrc=\relax
                          \let\@cite=\@bcite\@tempswafalse\@citex[]}}
\def\pcite{\@ifnextchar [{\let\@citebro=\space
                          \let\@citebrc=\relax
                          \let\@cite=\@pcite\@tempswatrue\@citex}
                        {\let\@citebro=\space
                          \let\@citebrc=\relax
                          \let\@cite=\@pcite\@tempswafalse\@citex[]}}
\def\scite{\@ifnextchar [{\let\@citebro=\@brcite
                          \let\@citebrc=\@ebrcite
                          \let\@cite=\@pcite\@tempswatrue\@citex}
                        {\let\@citebro=\@brcite
                          \let\@citebrc=\@ebrcite
                          \let\@cite=\@pcite\@tempswafalse\@citex[]}}

\renewenvironment{thebibliography}
  {\section*{REFERENCES}\vspace*{\afterbiblioskip}\small
    \if@book\RereadBiblio\fi
    \sloppy\clubpenalty4000\widowpenalty4000
    \sfcode`\.=1000\relax
    \setlength\parindent{\z@}\setlength\parskip{\z@}}
  {\relax\normalsize\@outputaddressesifnecessary} 
%</rmaacite>
%    \end{macrocode}

% \subsection{Auxilliary files for generating ADS data}

% The ADS data are written by \LaTeX{} to a |.ads| file. This file
% then has to be included in the following file (|ADS.tex|),
% which is run through \LaTeX{}2HTML (see below). 

%    \begin{macrocode}
%<*ads-container>
%%
%% Instructions:
%% Edit the \input{} line below to reflect the name of your file. 
\documentclass{article}
\usepackage[latin1]{inputenc}
\usepackage{adshtml}
\newcommand{\ADSfield}[2]
{\%#1  #2\\}
%%\newcommand{\ADSvfield}[2]
%%{\%#1 \begin{verbatim}#2\end{verbatim}}
\newenvironment{ADSarticle}{}{\par\bigskip}
\renewcommand{\bibitem}[2][]{\par}
\setlength{\parindent}{0pt}
\renewcommand{\ThisVolume}[1]{%  
  2001, \RMAASC, 12, \oldpageref{#1:RMAAFirstPage}}  
 

\begin{document}
\makeatletter
\input{rmbooksamp.ads}
\end{document}
%</ads-container>
%    \end{macrocode}

% This is the |adshtml.sty| style file used by |ADS.tex|. It just
% defines a few 
% macros that might be used in abstracts, etc (really it should define
% them all). This now gets all its code from
% Sec~\ref{sec:user-commands}. 
%    \begin{macrocode}
%<*adshtml-style>
%% Nothing else needs to go here now. 
%</adshtml-style>
%    \end{macrocode}
% The init file for latex2html.
%    \begin{macrocode}
%<*latex2html-init>
$PDFTOCAIRO = "pdftocairo";
$PDFLATEX = "pdflatex";
$LATEX = "latex";
$MAX_SPLIT_DEPTH = 0;
$HTML_VERSION = "4.0,math";
$NO_SIMPLE_MATH = 1;
$ADDRESS = "<p><address><a href=\"mailto:rmaa\@astro.unam.mx\">rmaa\@astro.unam.mx</address></p>"; 
$INFO = "<p>This document was created using latex2html and the rmaa latex macros</p>"; 
$NO_NAVIGATION = 1;
1; 
%</latex2html-init>
%    \end{macrocode}
% This is the |adshtml.perl| file that is used by \LaTeX{}2HTML. Some
% bizarre \TeX{}y command sometimes show up so we tell it to ignore
% these. All the astronomy macros that come from AAS are processed in
% \TeX{}. That way they get made to images but with an |ALT| tag that
% has the original code, which is what will be picked up by |lynx|,
% with the result that they get passed intact into the final |.txt|
% file. 
%    \begin{macrocode}
%<*adshtml-perl>
package main;

$dum = &do_cmd_makeatletter();

sub do_cmd_RMAA {
    local($_) = @_;
    "RMxAA" . $_;
}

sub do_cmd_RMAASC {
    local($_) = @_;
    "RMxAASC" . $_;
}

sub do_cmd_etal {
    local($_) = @_;
    "et al\." . $_;
}

sub do_cmd_ibidrule {
    local($_) = @_;
    "<img src=\"foo\" alt=\"\ibidrule\">" . $_;
}


sub do_cmd_bibitem {
    local($_) = @_;
    local($key);
    local($opt, $dummy)=&get_next_optional_argument; 
    $key = &missing_braces unless
        ((s/$next_pair_pr_rx/$key = $2; ''/eo)
        ||(s/$next_pair_rx/$key = $2; ''/eo));
    "<p>" . $_; 
}

&ignore_commands (<<_RAW_ARG_CMDS_);
lowercase
adjustfinalcols
textsuperscript # {}
newline # []
thinmuskip
mskip
smallskip
medskip
bigskip
hspace # {}
hspacestar # {}
vspace # {}
vspacestar # {}
mathrm
_at_mathrm
spacefactor
discretionary
unhbox 
voidb_at_x 
penalty 
_at_M
_at_m
null
IeC
_RAW_ARG_CMDS_

&process_commands_in_tex (<<_RAW_ARG_CMDS_);
sun
earth
lesssim
gtrsim
sq
argdeg
arcmin
arcsec
fd
fh
fm
fs
fdg
farcm
farcs
fp
micron
la
ga
case # {} # {}
slantfrac # {} # {}
onehalf
onethird
twothirds
onequarter
threequarters
ubvr
ub
bv
vr
ur
aj
araa
apj
apjl
apjs
ao
apss
aap
aapr
aaps
azh
baas
jrasc
memras
mnras
pra
prb
prc
prd
pre
prl
pasp
pasj
qjras
skytel
solphys
sovast
ssr
zap
nat
iaucirc
aplett
apspr
bain
fcp
gca
grl
jcp
jgr
jqsrt
memsai
nphysa
physrep
physscr
planss
procspie
_RAW_ARG_CMDS_

1;                              # This must be the last line
%</adshtml-perl> 
%    \end{macrocode}

% This is a shell script |makeads.sh| that generates the |_ads.txt| file from
% |_ads.tex| by using \LaTeX{}2HTML and |lynx|.  \changes{4.0}{2019/11/13}{Make sure that lynx writes utf-8 output instead of latin-1.}
% \changes{2.2beta13}{2002/03/12}{Complete rewrite of |makeads.sh|. Now much 
% improved I think}
% \changes{4.0}{2019/11/13}{Fix the at-sign substitution so that it only eliminates \cs{@} and does not mess with email addresses.}

%    \begin{macrocode}
%<*ads-shell-script>
#! /bin/sh
if [ -z $1 ]; then
    echo "Usage: $0 basename"
    exit
fi
BASENAME=$1
ADSBASENAME=${BASENAME}_ads
AUXFILE=$BASENAME.aux
TMPFILE=tmp.$$


echo "Generating $ADSBASENAME.tex"

# For some reason, must double-up \\ to \\\\ but not \ to \\ - odd!
cat > $ADSBASENAME.tex <<EOF
%% PLEASE DO NOT EDIT THIS FILE!  It is generated automatically from
%% the script makeads.sh.  Any hand-edited changes will be lost
%% 
\documentclass{article}
\usepackage[latin1]{inputenc}
\usepackage{html}
\usepackage{adshtml}
\newcommand{\ADSfield}[2]{\%#1  #2\\\\}
\newenvironment{ADSarticle}{}{\par\bigskip}
\setlength{\parindent}{0pt}
\newcounter{istage}
\renewcommand{\ion}[2]{\setcounter{istage}{#2}#1 \Roman{istage}}
\begin{document}
\makeatletter
EOF
echo >> $ADSBASENAME.tex "\input{$BASENAME.ads2}" 
echo >> $ADSBASENAME.tex "\end{document}"



echo "Generating $ADSBASENAME.sed" 

# Find type of journal
grep '\\documentclass\[.*proceedings.*\]' $BASENAME.tex >/dev/null \
     && JOURNAL='RMxAC' \
     || JOURNAL='RMxAA'

# Find the year and volume number
YEAR=`grep '\\SetYear' $BASENAME.tex | sed 's/.*\\SetYear{\(.*\)}.*/\1/'`
VOLUME=`grep '\\SetVolume' $BASENAME.tex | sed 's/.*\\SetVolume{\(.*\)}.*/\1/'`

# Generate an auxilliary sed script 
cat > $TMPFILE <<EOF
s|\\\\newlabel{\(.*\):RMAAFirstPage}{{.*}{\(.*\)}}|s/\\\\\\\\ThisVolume{\1}/$YEAR, RMxAC, $VOLUME, \2/g|g
EOF

# Now go through the .aux file, looking for \newlabels and use the sed 
# script just generated in order to make another sed script, one that 
# will replace \ThisVolume{LABEL} with the correct reference
grep '^\\newlabel{.*:RMAAFirstPage}' $AUXFILE \
     | sed -f $TMPFILE > $ADSBASENAME.sed && rm $TMPFILE

# Apend a further sed command to get rid of \@, since redefining 
# it in _ads.tex didn't work
echo >> $ADSBASENAME.sed 's/\\\\@//g'

echo "Generating $BASENAME.ads2 from $ADSBASENAME.sed and $BASENAME.ads" 

# Apply this to the .ads file, generating the .ads2 file, which is included
# by the _ads.tex file.
sed -f $ADSBASENAME.sed $BASENAME.ads > $BASENAME.ads2

echo "Generating $ADSBASENAME.txt from $ADSBASENAME.tex and $BASENAME.ads2"

# Now run latex2html on the _ads.tex file
latex2html -no_navigation -info 0 -address 0 -html_version 4.0,math $ADSBASENAME

# Then, convert to utf8 with lynx and tidy it up a bit
lynx --justify=off --width=100 --dump --nolist --display_charset=utf-8 \
     $ADSBASENAME/$ADSBASENAME.html | cut -c4- \
     | sed 's/ibidrule/\\ibidrule/g' > $ADSBASENAME.txt
%</ads-shell-script>
%    \end{macrocode}

% \section{Epilogue}
% Err, that's it \dots

\endinput


% \endinput
% Local Variables:
% mode: doctex
% TeX-master: t
% End:
"

clean: 
        rm -f rmaa.aux rmaa.dvi rmaa.ind rmaa.gls
%</makefile>
%    \end{macrocode}

% The |rmaa.ins| file uses the \textsf{docstrip} package to split out
% all the files and strip off this documentation.
%    \begin{macrocode}
%<*ins>
%%
%% Installation batch file for the RMAA LaTeX package
%%
\def\batchfile{rmaa.ins}
\input docstrip.tex
\preamble

PLEASE DO NOT EDIT THIS FILE! It is generated automatically from
the file rmaa.dtx using the accompanying Makefile and rmaa.ins 
Any hand-edited changes will be lost

A LaTeX document class for the RMxAA

Author: William Henney (w.henney@irya.unam.mx)

Latest version available from: http://www.irya.unam.mx/rmaa

Copyright (C) 1999-2019 Instituto de Astronomia, UNAM, Mexico

Permission is hereby granted, free of charge, to any person obtaining
a copy of this software and associated documentation files (the
"Software"), to deal in the Software without restriction, including
without limitation the rights to use, copy, modify, merge, publish,
distribute, sublicense, and/or sell copies of the Software, and to
permit persons to whom the Software is furnished to do so, subject to
the following conditions:

The above copyright notice and this permission notice shall be
included in all copies or substantial portions of the Software.

THE SOFTWARE IS PROVIDED "AS IS", WITHOUT WARRANTY OF ANY KIND,
EXPRESS OR IMPLIED, INCLUDING BUT NOT LIMITED TO THE WARRANTIES OF
MERCHANTABILITY, FITNESS FOR A PARTICULAR PURPOSE AND
NONINFRINGEMENT. IN NO EVENT SHALL THE AUTHORS OR COPYRIGHT HOLDERS BE
LIABLE FOR ANY CLAIM, DAMAGES OR OTHER LIABILITY, WHETHER IN AN ACTION
OF CONTRACT, TORT OR OTHERWISE, ARISING FROM, OUT OF OR IN CONNECTION
WITH THE SOFTWARE OR THE USE OR OTHER DEALINGS IN THE SOFTWARE.
\endpreamble
%    \end{macrocode}
% Clobber files without asking
%    \begin{macrocode}
\askforoverwritefalse
%    \end{macrocode}
% Make the \LaTeX{} macro files (class and styles)
%    \begin{macrocode}
\generate{% 
  \file{rmaa.cls}{\from{rmaa.dtx}{package}}
  \file{rmaa.ins}{\from{rmaa.dtx}{ins}}
  \file{rmaacite.sty}{\from{rmaa.dtx}{rmaacite}}
  \file{adshtml.sty}{\from{rmaa.dtx}{adshtml-style}}
  \file{ADS.tex}{\from{rmaa.dtx}{ads-container}}
  }
%    \end{macrocode}
% Make the other non-\LaTeX{} files. These are done separately since
% they can't use `\%\%' as a comment character for the
% preamble/postamble (and Perl files can't have the postamble I don't
% think). 
%    \begin{macrocode}
\generate{\nopreamble\nopostamble% 
  \file{Makefile}{\from{rmaa.dtx}{makefile}}
  \file{adshtml.perl}{\from{rmaa.dtx}{adshtml-perl}}
  \file{makeads.sh}{\from{rmaa.dtx}{ads-shell-script}}
  \file{sample.latex2html-init}{\from{rmaa.dtx}{latex2html-init}}
  }

%% Local Variables: 
%% mode: LaTeX
%% mode: Font-lock
%% End:
%</ins>
%    \end{macrocode}

% \subsection{Extended citation syntax}
% Modified from a file called mnrascite.sty
% I no longer remember where I stole the original version of this from
% 
% Alterations to the cite commands to make them do more sensible things.
% N.B. the bibliography style should produce entries which have labels where
% the date is delimited by angle brackets '<>' for these macros to work
% properly. They will not give very sensible error messages if you do not do 
% this properly.
% 
% A new command |\scite{key}| is defined which produces a citation which is 
% suitable for insertion in a sentence. i.e., it does not put brackets around 
% the whole citation, only the date. Also |\pcite{key}| produce output with no
% parentheses whatsoever. 
% 
% Should redefine |\@bibitem| really to have any hope of compatibility
% the command |\lbibcite| is defined to deal with the citations.
%    \begin{macrocode}
%<*rmaacite>
\def\bibitem[#1<#2>]#3{\par\hangindent=1.5em\hangafter=1
  \if@filesw
  {\def\protect##1{\string ##1\space}
    \immediate\write\@auxout{\string\lbibcite{#3}{#1}{#2}}}
  \fi}
\@redefinebibitem % what it says
%    \end{macrocode}
% Note that the |\lbibcite| defines two types of citation one contains the names,
% the other the year
%    \begin{macrocode}
\def\lbibcite#1#2#3{\global\@namedef{b@#1}{#2}
                    \global\@namedef{by@#1}{#3}}

\def\@bcite#1#2{({#1\if@tempswa , #2\fi})}
\def\@pcite#1#2{#1\if@tempswa , #2\fi}
%    \end{macrocode}
% This declaration contains what is put between the name and the date
% \changes{3.11}{2003/01/31}{Moved the \cs{articlename} to after the key for consistency
% with the new natbib stuff.}
%    \begin{macrocode}
\def\@brcite{\ (}
\def\@ebrcite{)}
\def\@fmttop#1{\@fmtcite[\csname #1\endcsname]}
\def\@citex[#1]#2{%
%%  \if@filesw\immediate\write\@auxout{\string\citation{#2}}\fi
%% WJH added 16 Jan 2002
  \@for\@citeb:=#2\do{%
    \edef\@citeb{\expandafter\@firstofone\@citeb:\articlename}%
    \if@filesw\immediate\write\@auxout{\string\citation{\@citeb}}\fi}%
%% to here
  \def\@citea{}%
  \@cite{\@for\@citeb:=#2\do
    {%
      \edef\@citeb{\expandafter\@firstofone\@citeb:\articlename}%
      \@citea\def\@citea{;\ }\@ifundefined
      {b@\@citeb}{{\bf ?}\@warning {Citation `\@citeb' on page \thepage \space undefined}}%
      {\csname b@\@citeb\endcsname\@citebro\csname by@\@citeb\endcsname\@citebrc}}}%
  {#1}}

\def\cite{\@ifnextchar [{\let\@citebro=\space
                          \let\@citebrc=\relax
                          \let\@cite=\@bcite\@tempswatrue\@citex}
                        {\let\@citebro=\space
                          \let\@citebrc=\relax
                          \let\@cite=\@bcite\@tempswafalse\@citex[]}}
\def\pcite{\@ifnextchar [{\let\@citebro=\space
                          \let\@citebrc=\relax
                          \let\@cite=\@pcite\@tempswatrue\@citex}
                        {\let\@citebro=\space
                          \let\@citebrc=\relax
                          \let\@cite=\@pcite\@tempswafalse\@citex[]}}
\def\scite{\@ifnextchar [{\let\@citebro=\@brcite
                          \let\@citebrc=\@ebrcite
                          \let\@cite=\@pcite\@tempswatrue\@citex}
                        {\let\@citebro=\@brcite
                          \let\@citebrc=\@ebrcite
                          \let\@cite=\@pcite\@tempswafalse\@citex[]}}

\renewenvironment{thebibliography}
  {\section*{REFERENCES}\vspace*{\afterbiblioskip}\small
    \if@book\RereadBiblio\fi
    \sloppy\clubpenalty4000\widowpenalty4000
    \sfcode`\.=1000\relax
    \setlength\parindent{\z@}\setlength\parskip{\z@}}
  {\relax\normalsize\@outputaddressesifnecessary} 
%</rmaacite>
%    \end{macrocode}

% \subsection{Auxilliary files for generating ADS data}

% The ADS data are written by \LaTeX{} to a |.ads| file. This file
% then has to be included in the following file (|ADS.tex|),
% which is run through \LaTeX{}2HTML (see below). 

%    \begin{macrocode}
%<*ads-container>
%%
%% Instructions:
%% Edit the \input{} line below to reflect the name of your file. 
\documentclass{article}
\usepackage[latin1]{inputenc}
\usepackage{adshtml}
\newcommand{\ADSfield}[2]
{\%#1  #2\\}
%%\newcommand{\ADSvfield}[2]
%%{\%#1 \begin{verbatim}#2\end{verbatim}}
\newenvironment{ADSarticle}{}{\par\bigskip}
\renewcommand{\bibitem}[2][]{\par}
\setlength{\parindent}{0pt}
\renewcommand{\ThisVolume}[1]{%  
  2001, \RMAASC, 12, \oldpageref{#1:RMAAFirstPage}}  
 

\begin{document}
\makeatletter
%% Demonstration file for formatting an entire volume
\documentclass[book,journal,twocolumn]{rmaa}
%%%%%%%%%%%%%%%%%%%%%%%%%%%%%%%%%%%%%%%%%%%%%%%%%%%%%%%%%%%%%%%%
%%%%%%%%%%%%%%%%            PREAMBLE            %%%%%%%%%%%%%%%%
%%%%%%%%%%%%%%%%%%%%%%%%%%%%%%%%%%%%%%%%%%%%%%%%%%%%%%%%%%%%%%%%

\usepackage{paralist}
\usepackage{psfrag,color}
% Needed for landscape tables
\usepackage{pdflscape}

% Used in rm-minimal-test
\usepackage{kantlipsum}

% publication details of volume
\SetVolume{9}
\SetYear{2000}
\SetMonthNumeric{06}
\SetMes{junio}
% details of meeting
\SetConfTitle{Astrophysical Plasmas: Codes, Models, \& Observations}
\SetReunion{M\'exico, D. F.}{octubre 25--29, 1999}
\SetEditors{S. Jane Arthur, Nancy S. Brickhouse, \& Jos\'e Franco}

% this writes the .idx file that is used to generate the index 
\makeindex

%% This command makes the session headings come out bold in the
%% Table-of-Contents. Comment it out if you don't like it.
\BoldSessionTitles

%% This command makes a title page for each session
%% Comment it out if you don't like it.
\SessionTitlePages

%% Use this command if you want the abstract-only papers to be
%% indexed by ADS
\abstractsADS

%% At the special request of Alberto Carramin~ana, the following is a
%% way to suppress the author adresses at the end of each article
%\suppressfulladdresses

%% Use this command to add Received/accepted dates for each paper
% \printdates


\begin{document}

%%%%%%%%%%%%%%%%%%%%%%%%%%%%%%%%%%%%%%%%%%%%%%%%%%%%%%%%%%%%%%%%b
%%%%%%%%%%%%%%%%          FRONT MATTER          %%%%%%%%%%%%%%%%
%%%%%%%%%%%%%%%%%%%%%%%%%%%%%%%%%%%%%%%%%%%%%%%%%%%%%%%%%%%%%%%%
\frontmatter % start roman numbering

% format the TOC
\renewcommand{\RMAAbannerimage}{rmsc_blank}
\tableofcontents

%% preface, group photo, list of participants go here
%% use the `frontsection' environment (see editors' guide)
\begin{frontsection}{Preface}{%
    \listofauthors{S.~J.~Arthur, N.~S.~Brickhouse, J.~Franco}
    \indexauthor{Arthur, S. J.}
    \indexauthor{Brickhouse, N.}
    \indexauthor{Franco, J.}
    }
  These are the proceedings of a very nice conference that Jane and
  Pepe organised. This is the preface, which explains what a wonderful
  time was had by all and how grateful they are for all the people who
  gave them money or helped out.
  
  This is another paragraph of the preface, in case they have anything
  else to say.  For example, the Scientific Organizing Committee
  consisted of
  \begin{quote}
    J. Franco (Mexico, Chair), D. Cox (USA) and various other notables
  \end{quote}
  while the Local Organizing Committee comprised
  \begin{quote}
    S. J. Arthur (Chair), S. Kurtz, etc.
  \end{quote}

  \EditorSignature

\end{frontsection}

\begin{frontsection}[Photo]{Group Photograph}{\notitle}
\null\vfill
\centerline{Here goes the photograph.}
\vfill
\end{frontsection}

\begin{frontsection}[Participants]{List of Participants}{}
%% Here you put a formatted list of participants
\vfil
\dotfill First page of list of participants \dotfill
\vfil\clearpage

\null\vfil
\dotfill Second page of list of participants \dotfill
\vfil

\end{frontsection}

\begin{frontsection}[XXX]{Another front section}{}

\end{frontsection}

%%%%%%%%%%%%%%%%%%%%%%%%%%%%%%%%%%%%%%%%%%%%%%%%%%%%%%%%%%%%%%%%
%%%%%%%%%%%%%%%%          MAIN MATTER           %%%%%%%%%%%%%%%%
%%%%%%%%%%%%%%%%%%%%%%%%%%%%%%%%%%%%%%%%%%%%%%%%%%%%%%%%%%%%%%%%
\mainmatter % start arabic numbering

\begin{articles}
  %% Main articles go here
  %% Use \article{ARTICLENAME} for each individual article
  %% If you want any text to appear in brackets after the article
  %% title in the TOC and index, then put it as an optional first
  %% argument, e.g.,  \article[Invited Talk]{ARTICLENAME} 

  \session{The Contributed Papers}
  \article{rm-journal-example} % An illustrative article
  \article{rm-onepage} % another article, which has a plate
  \article{rm-minimal-test}
  \session{Conference Summary}
  \article[Invited Review]{rmsummsamp} % An illustrative conference summary

\end{articles}

%%%%%%%%%%%%%%%%%%%%%%%%%%%%%%%%%%%%%%%%%%%%%%%%%%%%%%%%%%%%%%%%
%%%%%%%%%%%%%%%%          BACK MATTER           %%%%%%%%%%%%%%%%
%%%%%%%%%%%%%%%%%%%%%%%%%%%%%%%%%%%%%%%%%%%%%%%%%%%%%%%%%%%%%%%%
\backmatter

%% Abstracts of presentations for which the
%% authors have not submitted a paper
\begin{abstracts}

  \title{An abstract of a poster}
  \author{W. J. Henney\altaffilmark{1}}
  \listofauthors{W.~J.~Henney}
  \indexauthor{Henney, W. J}
  \altaffiltext{1}{Where I am from done with altaffilmark/text}
  \abstract{This is a test abstract.   This is a test abstract.   This
    is a test abstract.   This is a test abstract.}
  

  \title{An abstract of another  poster}
  \author{W.~J.~Henney\altaffil{This is done with altaffil} and
    A.~N.~Other\altaffil{So is this} }
  \listofauthors{W.~J.~Henney \& A.~N.~Other}
  \indexauthor{Henney, W. J}
  \indexauthor{Other, A. N.}
  \abstract{This is a test abstract.   This is a test abstract.   This
    is a test abstract.   This is a test abstract.}

  \title{An abstract of yet another  poster}
  \author{W.~J.~Henney\altaffil{This is done with altaffil} and
    A.~N.~Other\altaffil{So is this} }
  \listofauthors{W.~J.~Henney \& A.~N.~Other}
  \indexauthor{Henney, W. J}
  \indexauthor{Other, A. N.}
  \abstract{This is a test abstract.   This is a test abstract.   This
    is a test abstract.   This is a test abstract.}

  \title{The abstract of a poster with a very long title to see how
    well it gets split up into lines}
  \author{A.~N.~Other\altaffil{The affiliation of A.~N.~Other is also
      very long so that it will go over several lines. His email
      address is anon@nowhere.org}}
  \listofauthors{A.~N.~Other}
  \indexauthor{Other, A. N.}
  \abstract{We will try and make this abstract longer so there is more
    stuff to put on the page and it doesn't look too empty. However, I
    have nothing whatsoever to say so will have to stop here. Let's
    repeat that a few times. We will try and make this abstract longer
    so there is more stuff to put on the page and it doesn't look too
    empty. However, I have nothing whatsoever to say so will have to
    stop here. Let's repeat that a few times. We will try and make
    this abstract longer so there is more stuff to put on the page and
    it doesn't look too empty. However, I have nothing whatsoever to
    say so will have to stop here. Let's repeat that a few times. We
    will try and make this abstract longer so there is more stuff to
    put on the page and it doesn't look too empty. However, I have
    nothing whatsoever to say so will have to stop here. Let's repeat
    that a few times.}
    
  \title[Comment to go in TOC]{The abstract of a second poster with a
    very long title to see how well it gets split up into lines}
  \author{A.~N.~Other\altaffil{The affiliation of A.~N.~Other is also
      very long so that it will go over several lines. His email
      address is anon@nowhere.org}}
  \listofauthors{A.~N.~Other}
  \indexauthor{Other, A. N.}
  \abstract{We will try and make this abstract longer so there is more
    stuff to put on the page and it doesn't look too empty. However, I
    have nothing whatsoever to say so will have to stop here. Let's
    repeat that a few times. We will try and make this abstract longer
    so there is more stuff to put on the page and it doesn't look too
    empty. However, I have nothing whatsoever to say so will have to
    stop here. Let's repeat that a few times. We will try and make
    this abstract longer so there is more stuff to put on the page and
    it doesn't look too empty. However, I have nothing whatsoever to
    say so will have to stop here. Let's repeat that a few times. We
    will try and make this abstract longer so there is more stuff to
    put on the page and it doesn't look too empty. However, I have
    nothing whatsoever to say so will have to stop here. Let's repeat
    that a few times.}
    
  \title{The abstract of a third poster with a very long title to see how
    well it gets split up into lines}
  \author{A.~N.~Other\altaffil{The affiliation of A.~N.~Other is also
      very long so that it will go over several lines. His email
      address is anon@nowhere.org}}
  \listofauthors{A.~N.~Other}
  \indexauthor{Other, A. N.}
  \abstract{We will try and make this abstract longer so there is more
    stuff to put on the page and it doesn't look too empty. However, I
    have nothing whatsoever to say so will have to stop here. Let's
    repeat that a few times. We will try and make this abstract longer
    so there is more stuff to put on the page and it doesn't look too
    empty. However, I have nothing whatsoever to say so will have to
    stop here. Let's repeat that a few times. We will try and make
    this abstract longer so there is more stuff to put on the page and
    it doesn't look too empty. However, I have nothing whatsoever to
    say so will have to stop here. Let's repeat that a few times. We
    will try and make this abstract longer so there is more stuff to
    put on the page and it doesn't look too empty. However, I have
    nothing whatsoever to say so will have to stop here. Let's repeat
    that a few times.}
    
  \title{The abstract of a fourth poster with a very long title to see how
    well it gets split up into lines}
  \author{A.~N.~Other\altaffil{The affiliation of A.~N.~Other is also
      very long so that it will go over several lines. His email
      address is anon@nowhere.org}}
  \listofauthors{A.~N.~Other}
  \indexauthor{Other, A. N.}
  \abstract{We will try and make this abstract longer so there is more
    stuff to put on the page and it doesn't look too empty. However, I
    have nothing whatsoever to say so will have to stop here. Let's
    repeat that a few times. We will try and make this abstract longer
    so there is more stuff to put on the page and it doesn't look too
    empty. However, I have nothing whatsoever to say so will have to
    stop here. Let's repeat that a few times. We will try and make
    this abstract longer so there is more stuff to put on the page and
    it doesn't look too empty. However, I have nothing whatsoever to
    say so will have to stop here. Let's repeat that a few times. We
    will try and make this abstract longer so there is more stuff to
    put on the page and it doesn't look too empty. However, I have
    nothing whatsoever to say so will have to stop here. Let's repeat
    that a few times.}
    
  \title{The abstract of a fifth poster with a very long title to see how
    well it gets split up into lines}
  \author{A.~N.~Other\altaffil{The affiliation of A.~N.~Other is also
      very long so that it will go over several lines. His email
      address is anon@nowhere.org}}
  \listofauthors{A.~N.~Other}
  \indexauthor{Other, A. N.}
  \abstract{We will try and make this abstract longer so there is more
    stuff to put on the page and it doesn't look too empty. However, I
    have nothing whatsoever to say so will have to stop here. Let's
    repeat that a few times. We will try and make this abstract longer
    so there is more stuff to put on the page and it doesn't look too
    empty. However, I have nothing whatsoever to say so will have to
    stop here. Let's repeat that a few times. We will try and make
    this abstract longer so there is more stuff to put on the page and
    it doesn't look too empty. However, I have nothing whatsoever to
    say so will have to stop here. Let's repeat that a few times. We
    will try and make this abstract longer so there is more stuff to
    put on the page and it doesn't look too empty. However, I have
    nothing whatsoever to say so will have to stop here. Let's repeat
    that a few times.}
    
  \title{The abstract of a sixth poster with a very long title to see how
    well it gets split up into lines}
  \author{A.~N.~Other\altaffil{The affiliation of A.~N.~Other is also
      very long so that it will go over several lines. His email
      address is anon@nowhere.org}}
  \listofauthors{A.~N.~Other}
  \indexauthor{Other, A. N.}
  \abstract{We will try and make this abstract longer so there is more
    stuff to put on the page and it doesn't look too empty. However, I
    have nothing whatsoever to say so will have to stop here. Let's
    repeat that a few times. We will try and make this abstract longer
    so there is more stuff to put on the page and it doesn't look too
    empty. However, I have nothing whatsoever to say so will have to
    stop here. Let's repeat that a few times. We will try and make
    this abstract longer so there is more stuff to put on the page and
    it doesn't look too empty. However, I have nothing whatsoever to
    say so will have to stop here. Let's repeat that a few times. We
    will try and make this abstract longer so there is more stuff to
    put on the page and it doesn't look too empty. However, I have
    nothing whatsoever to say so will have to stop here. Let's repeat
    that a few times.}
    
  \title{The abstract of a seventh poster with a very long title to see how
    well it gets split up into lines}
  \author{A.~N.~Other\altaffil{The affiliation of A.~N.~Other is also
      very long so that it will go over several lines. His email
      address is anon@nowhere.org}}
  \listofauthors{A.~N.~Other}
  \indexauthor{Other, A. N.}
  \abstract{We will try and make this abstract longer so there is more
    stuff to put on the page and it doesn't look too empty. However, I
    have nothing whatsoever to say so will have to stop here. Let's
    repeat that a few times. We will try and make this abstract longer
    so there is more stuff to put on the page and it doesn't look too
    empty. However, I have nothing whatsoever to say so will have to
    stop here. Let's repeat that a few times. We will try and make
    this abstract longer so there is more stuff to put on the page and
    it doesn't look too empty. However, I have nothing whatsoever to
    say so will have to stop here. Let's repeat that a few times. We
    will try and make this abstract longer so there is more stuff to
    put on the page and it doesn't look too empty. However, I have
    nothing whatsoever to say so will have to stop here. Let's repeat
    that a few times.}
  
  \title{An abstract to go in the list of abstracts, but not in the
    index of authors}
  \author{I.~A.~M. Hidden\altaffil{Puerto Escondido}}
  \listofauthors{I.~A.~M. Hidden}
  \abstract{The idea is that by omitting the
    \textbackslash{}indexauthor command, we can avoid an entry in the
    author index.}

\end{abstracts}


%% The Index: after latexing twice, run `makeindex -s rmaa' on this
%% file, then latex twice more
\printindex

%% if you rename this file to something other than `rmbooksamp', then
%% you must change the above line accordingly

%%
%% Plates: If you have any, then the following line will print
%% them out:

\outputplates

\end{document}

%%% Local Variables:
%%% mode: latex
%%% TeX-master: t
%%% End:

\end{document}
%</ads-container>
%    \end{macrocode}

% This is the |adshtml.sty| style file used by |ADS.tex|. It just
% defines a few 
% macros that might be used in abstracts, etc (really it should define
% them all). This now gets all its code from
% Sec~\ref{sec:user-commands}. 
%    \begin{macrocode}
%<*adshtml-style>
%% Nothing else needs to go here now. 
%</adshtml-style>
%    \end{macrocode}
% The init file for latex2html.
%    \begin{macrocode}
%<*latex2html-init>
$PDFTOCAIRO = "pdftocairo";
$PDFLATEX = "pdflatex";
$LATEX = "latex";
$MAX_SPLIT_DEPTH = 0;
$HTML_VERSION = "4.0,math";
$NO_SIMPLE_MATH = 1;
$ADDRESS = "<p><address><a href=\"mailto:rmaa\@astro.unam.mx\">rmaa\@astro.unam.mx</address></p>"; 
$INFO = "<p>This document was created using latex2html and the rmaa latex macros</p>"; 
$NO_NAVIGATION = 1;
1; 
%</latex2html-init>
%    \end{macrocode}
% This is the |adshtml.perl| file that is used by \LaTeX{}2HTML. Some
% bizarre \TeX{}y command sometimes show up so we tell it to ignore
% these. All the astronomy macros that come from AAS are processed in
% \TeX{}. That way they get made to images but with an |ALT| tag that
% has the original code, which is what will be picked up by |lynx|,
% with the result that they get passed intact into the final |.txt|
% file. 
%    \begin{macrocode}
%<*adshtml-perl>
package main;

$dum = &do_cmd_makeatletter();

sub do_cmd_RMAA {
    local($_) = @_;
    "RMxAA" . $_;
}

sub do_cmd_RMAASC {
    local($_) = @_;
    "RMxAASC" . $_;
}

sub do_cmd_etal {
    local($_) = @_;
    "et al\." . $_;
}

sub do_cmd_ibidrule {
    local($_) = @_;
    "<img src=\"foo\" alt=\"\ibidrule\">" . $_;
}


sub do_cmd_bibitem {
    local($_) = @_;
    local($key);
    local($opt, $dummy)=&get_next_optional_argument; 
    $key = &missing_braces unless
        ((s/$next_pair_pr_rx/$key = $2; ''/eo)
        ||(s/$next_pair_rx/$key = $2; ''/eo));
    "<p>" . $_; 
}

&ignore_commands (<<_RAW_ARG_CMDS_);
lowercase
adjustfinalcols
textsuperscript # {}
newline # []
thinmuskip
mskip
smallskip
medskip
bigskip
hspace # {}
hspacestar # {}
vspace # {}
vspacestar # {}
mathrm
_at_mathrm
spacefactor
discretionary
unhbox 
voidb_at_x 
penalty 
_at_M
_at_m
null
IeC
_RAW_ARG_CMDS_

&process_commands_in_tex (<<_RAW_ARG_CMDS_);
sun
earth
lesssim
gtrsim
sq
argdeg
arcmin
arcsec
fd
fh
fm
fs
fdg
farcm
farcs
fp
micron
la
ga
case # {} # {}
slantfrac # {} # {}
onehalf
onethird
twothirds
onequarter
threequarters
ubvr
ub
bv
vr
ur
aj
araa
apj
apjl
apjs
ao
apss
aap
aapr
aaps
azh
baas
jrasc
memras
mnras
pra
prb
prc
prd
pre
prl
pasp
pasj
qjras
skytel
solphys
sovast
ssr
zap
nat
iaucirc
aplett
apspr
bain
fcp
gca
grl
jcp
jgr
jqsrt
memsai
nphysa
physrep
physscr
planss
procspie
_RAW_ARG_CMDS_

1;                              # This must be the last line
%</adshtml-perl> 
%    \end{macrocode}

% This is a shell script |makeads.sh| that generates the |_ads.txt| file from
% |_ads.tex| by using \LaTeX{}2HTML and |lynx|.  \changes{4.0}{2019/11/13}{Make sure that lynx writes utf-8 output instead of latin-1.}
% \changes{2.2beta13}{2002/03/12}{Complete rewrite of |makeads.sh|. Now much 
% improved I think}
% \changes{4.0}{2019/11/13}{Fix the at-sign substitution so that it only eliminates \cs{@} and does not mess with email addresses.}

%    \begin{macrocode}
%<*ads-shell-script>
#! /bin/sh
if [ -z $1 ]; then
    echo "Usage: $0 basename"
    exit
fi
BASENAME=$1
ADSBASENAME=${BASENAME}_ads
AUXFILE=$BASENAME.aux
TMPFILE=tmp.$$


echo "Generating $ADSBASENAME.tex"

# For some reason, must double-up \\ to \\\\ but not \ to \\ - odd!
cat > $ADSBASENAME.tex <<EOF
%% PLEASE DO NOT EDIT THIS FILE!  It is generated automatically from
%% the script makeads.sh.  Any hand-edited changes will be lost
%% 
\documentclass{article}
\usepackage[latin1]{inputenc}
\usepackage{html}
\usepackage{adshtml}
\newcommand{\ADSfield}[2]{\%#1  #2\\\\}
\newenvironment{ADSarticle}{}{\par\bigskip}
\setlength{\parindent}{0pt}
\newcounter{istage}
\renewcommand{\ion}[2]{\setcounter{istage}{#2}#1 \Roman{istage}}
\begin{document}
\makeatletter
EOF
echo >> $ADSBASENAME.tex "\input{$BASENAME.ads2}" 
echo >> $ADSBASENAME.tex "\end{document}"



echo "Generating $ADSBASENAME.sed" 

# Find type of journal
grep '\\documentclass\[.*proceedings.*\]' $BASENAME.tex >/dev/null \
     && JOURNAL='RMxAC' \
     || JOURNAL='RMxAA'

# Find the year and volume number
YEAR=`grep '\\SetYear' $BASENAME.tex | sed 's/.*\\SetYear{\(.*\)}.*/\1/'`
VOLUME=`grep '\\SetVolume' $BASENAME.tex | sed 's/.*\\SetVolume{\(.*\)}.*/\1/'`

# Generate an auxilliary sed script 
cat > $TMPFILE <<EOF
s|\\\\newlabel{\(.*\):RMAAFirstPage}{{.*}{\(.*\)}}|s/\\\\\\\\ThisVolume{\1}/$YEAR, RMxAC, $VOLUME, \2/g|g
EOF

# Now go through the .aux file, looking for \newlabels and use the sed 
# script just generated in order to make another sed script, one that 
# will replace \ThisVolume{LABEL} with the correct reference
grep '^\\newlabel{.*:RMAAFirstPage}' $AUXFILE \
     | sed -f $TMPFILE > $ADSBASENAME.sed && rm $TMPFILE

# Apend a further sed command to get rid of \@, since redefining 
# it in _ads.tex didn't work
echo >> $ADSBASENAME.sed 's/\\\\@//g'

echo "Generating $BASENAME.ads2 from $ADSBASENAME.sed and $BASENAME.ads" 

# Apply this to the .ads file, generating the .ads2 file, which is included
# by the _ads.tex file.
sed -f $ADSBASENAME.sed $BASENAME.ads > $BASENAME.ads2

echo "Generating $ADSBASENAME.txt from $ADSBASENAME.tex and $BASENAME.ads2"

# Now run latex2html on the _ads.tex file
latex2html -no_navigation -info 0 -address 0 -html_version 4.0,math $ADSBASENAME

# Then, convert to utf8 with lynx and tidy it up a bit
lynx --justify=off --width=100 --dump --nolist --display_charset=utf-8 \
     $ADSBASENAME/$ADSBASENAME.html | cut -c4- \
     | sed 's/ibidrule/\\ibidrule/g' > $ADSBASENAME.txt
%</ads-shell-script>
%    \end{macrocode}

% \section{Epilogue}
% Err, that's it \dots

\endinput


% \endinput
% Local Variables:
% mode: doctex
% TeX-master: t
% End:
"

clean: 
        rm -f rmaa.aux rmaa.dvi rmaa.ind rmaa.gls
%</makefile>
%    \end{macrocode}

% The |rmaa.ins| file uses the \textsf{docstrip} package to split out
% all the files and strip off this documentation.
%    \begin{macrocode}
%<*ins>
%%
%% Installation batch file for the RMAA LaTeX package
%%
\def\batchfile{rmaa.ins}
\input docstrip.tex
\preamble

PLEASE DO NOT EDIT THIS FILE! It is generated automatically from
the file rmaa.dtx using the accompanying Makefile and rmaa.ins 
Any hand-edited changes will be lost

A LaTeX document class for the RMxAA

Author: William Henney (w.henney@irya.unam.mx)

Latest version available from: http://www.irya.unam.mx/rmaa

Copyright (C) 1999-2019 Instituto de Astronomia, UNAM, Mexico

Permission is hereby granted, free of charge, to any person obtaining
a copy of this software and associated documentation files (the
"Software"), to deal in the Software without restriction, including
without limitation the rights to use, copy, modify, merge, publish,
distribute, sublicense, and/or sell copies of the Software, and to
permit persons to whom the Software is furnished to do so, subject to
the following conditions:

The above copyright notice and this permission notice shall be
included in all copies or substantial portions of the Software.

THE SOFTWARE IS PROVIDED "AS IS", WITHOUT WARRANTY OF ANY KIND,
EXPRESS OR IMPLIED, INCLUDING BUT NOT LIMITED TO THE WARRANTIES OF
MERCHANTABILITY, FITNESS FOR A PARTICULAR PURPOSE AND
NONINFRINGEMENT. IN NO EVENT SHALL THE AUTHORS OR COPYRIGHT HOLDERS BE
LIABLE FOR ANY CLAIM, DAMAGES OR OTHER LIABILITY, WHETHER IN AN ACTION
OF CONTRACT, TORT OR OTHERWISE, ARISING FROM, OUT OF OR IN CONNECTION
WITH THE SOFTWARE OR THE USE OR OTHER DEALINGS IN THE SOFTWARE.
\endpreamble
%    \end{macrocode}
% Clobber files without asking
%    \begin{macrocode}
\askforoverwritefalse
%    \end{macrocode}
% Make the \LaTeX{} macro files (class and styles)
%    \begin{macrocode}
\generate{% 
  \file{rmaa.cls}{\from{rmaa.dtx}{package}}
  \file{rmaa.ins}{\from{rmaa.dtx}{ins}}
  \file{rmaacite.sty}{\from{rmaa.dtx}{rmaacite}}
  \file{adshtml.sty}{\from{rmaa.dtx}{adshtml-style}}
  \file{ADS.tex}{\from{rmaa.dtx}{ads-container}}
  }
%    \end{macrocode}
% Make the other non-\LaTeX{} files. These are done separately since
% they can't use `\%\%' as a comment character for the
% preamble/postamble (and Perl files can't have the postamble I don't
% think). 
%    \begin{macrocode}
\generate{\nopreamble\nopostamble% 
  \file{Makefile}{\from{rmaa.dtx}{makefile}}
  \file{adshtml.perl}{\from{rmaa.dtx}{adshtml-perl}}
  \file{makeads.sh}{\from{rmaa.dtx}{ads-shell-script}}
  \file{sample.latex2html-init}{\from{rmaa.dtx}{latex2html-init}}
  }

%% Local Variables: 
%% mode: LaTeX
%% mode: Font-lock
%% End:
%</ins>
%    \end{macrocode}

% \subsection{Extended citation syntax}
% Modified from a file called mnrascite.sty
% I no longer remember where I stole the original version of this from
% 
% Alterations to the cite commands to make them do more sensible things.
% N.B. the bibliography style should produce entries which have labels where
% the date is delimited by angle brackets '<>' for these macros to work
% properly. They will not give very sensible error messages if you do not do 
% this properly.
% 
% A new command |\scite{key}| is defined which produces a citation which is 
% suitable for insertion in a sentence. i.e., it does not put brackets around 
% the whole citation, only the date. Also |\pcite{key}| produce output with no
% parentheses whatsoever. 
% 
% Should redefine |\@bibitem| really to have any hope of compatibility
% the command |\lbibcite| is defined to deal with the citations.
%    \begin{macrocode}
%<*rmaacite>
\def\bibitem[#1<#2>]#3{\par\hangindent=1.5em\hangafter=1
  \if@filesw
  {\def\protect##1{\string ##1\space}
    \immediate\write\@auxout{\string\lbibcite{#3}{#1}{#2}}}
  \fi}
\@redefinebibitem % what it says
%    \end{macrocode}
% Note that the |\lbibcite| defines two types of citation one contains the names,
% the other the year
%    \begin{macrocode}
\def\lbibcite#1#2#3{\global\@namedef{b@#1}{#2}
                    \global\@namedef{by@#1}{#3}}

\def\@bcite#1#2{({#1\if@tempswa , #2\fi})}
\def\@pcite#1#2{#1\if@tempswa , #2\fi}
%    \end{macrocode}
% This declaration contains what is put between the name and the date
% \changes{3.11}{2003/01/31}{Moved the \cs{articlename} to after the key for consistency
% with the new natbib stuff.}
%    \begin{macrocode}
\def\@brcite{\ (}
\def\@ebrcite{)}
\def\@fmttop#1{\@fmtcite[\csname #1\endcsname]}
\def\@citex[#1]#2{%
%%  \if@filesw\immediate\write\@auxout{\string\citation{#2}}\fi
%% WJH added 16 Jan 2002
  \@for\@citeb:=#2\do{%
    \edef\@citeb{\expandafter\@firstofone\@citeb:\articlename}%
    \if@filesw\immediate\write\@auxout{\string\citation{\@citeb}}\fi}%
%% to here
  \def\@citea{}%
  \@cite{\@for\@citeb:=#2\do
    {%
      \edef\@citeb{\expandafter\@firstofone\@citeb:\articlename}%
      \@citea\def\@citea{;\ }\@ifundefined
      {b@\@citeb}{{\bf ?}\@warning {Citation `\@citeb' on page \thepage \space undefined}}%
      {\csname b@\@citeb\endcsname\@citebro\csname by@\@citeb\endcsname\@citebrc}}}%
  {#1}}

\def\cite{\@ifnextchar [{\let\@citebro=\space
                          \let\@citebrc=\relax
                          \let\@cite=\@bcite\@tempswatrue\@citex}
                        {\let\@citebro=\space
                          \let\@citebrc=\relax
                          \let\@cite=\@bcite\@tempswafalse\@citex[]}}
\def\pcite{\@ifnextchar [{\let\@citebro=\space
                          \let\@citebrc=\relax
                          \let\@cite=\@pcite\@tempswatrue\@citex}
                        {\let\@citebro=\space
                          \let\@citebrc=\relax
                          \let\@cite=\@pcite\@tempswafalse\@citex[]}}
\def\scite{\@ifnextchar [{\let\@citebro=\@brcite
                          \let\@citebrc=\@ebrcite
                          \let\@cite=\@pcite\@tempswatrue\@citex}
                        {\let\@citebro=\@brcite
                          \let\@citebrc=\@ebrcite
                          \let\@cite=\@pcite\@tempswafalse\@citex[]}}

\renewenvironment{thebibliography}
  {\section*{REFERENCES}\vspace*{\afterbiblioskip}\small
    \if@book\RereadBiblio\fi
    \sloppy\clubpenalty4000\widowpenalty4000
    \sfcode`\.=1000\relax
    \setlength\parindent{\z@}\setlength\parskip{\z@}}
  {\relax\normalsize\@outputaddressesifnecessary} 
%</rmaacite>
%    \end{macrocode}

% \subsection{Auxilliary files for generating ADS data}

% The ADS data are written by \LaTeX{} to a |.ads| file. This file
% then has to be included in the following file (|ADS.tex|),
% which is run through \LaTeX{}2HTML (see below). 

%    \begin{macrocode}
%<*ads-container>
%%
%% Instructions:
%% Edit the \input{} line below to reflect the name of your file. 
\documentclass{article}
\usepackage[latin1]{inputenc}
\usepackage{adshtml}
\newcommand{\ADSfield}[2]
{\%#1  #2\\}
%%\newcommand{\ADSvfield}[2]
%%{\%#1 \begin{verbatim}#2\end{verbatim}}
\newenvironment{ADSarticle}{}{\par\bigskip}
\renewcommand{\bibitem}[2][]{\par}
\setlength{\parindent}{0pt}
\renewcommand{\ThisVolume}[1]{%  
  2001, \RMAASC, 12, \oldpageref{#1:RMAAFirstPage}}  
 

\begin{document}
\makeatletter
%% Demonstration file for formatting an entire volume
\documentclass[book,journal,twocolumn]{rmaa}
%%%%%%%%%%%%%%%%%%%%%%%%%%%%%%%%%%%%%%%%%%%%%%%%%%%%%%%%%%%%%%%%
%%%%%%%%%%%%%%%%            PREAMBLE            %%%%%%%%%%%%%%%%
%%%%%%%%%%%%%%%%%%%%%%%%%%%%%%%%%%%%%%%%%%%%%%%%%%%%%%%%%%%%%%%%

\usepackage{paralist}
\usepackage{psfrag,color}
% Needed for landscape tables
\usepackage{pdflscape}

% Used in rm-minimal-test
\usepackage{kantlipsum}

% publication details of volume
\SetVolume{9}
\SetYear{2000}
\SetMonthNumeric{06}
\SetMes{junio}
% details of meeting
\SetConfTitle{Astrophysical Plasmas: Codes, Models, \& Observations}
\SetReunion{M\'exico, D. F.}{octubre 25--29, 1999}
\SetEditors{S. Jane Arthur, Nancy S. Brickhouse, \& Jos\'e Franco}

% this writes the .idx file that is used to generate the index 
\makeindex

%% This command makes the session headings come out bold in the
%% Table-of-Contents. Comment it out if you don't like it.
\BoldSessionTitles

%% This command makes a title page for each session
%% Comment it out if you don't like it.
\SessionTitlePages

%% Use this command if you want the abstract-only papers to be
%% indexed by ADS
\abstractsADS

%% At the special request of Alberto Carramin~ana, the following is a
%% way to suppress the author adresses at the end of each article
%\suppressfulladdresses

%% Use this command to add Received/accepted dates for each paper
% \printdates


\begin{document}

%%%%%%%%%%%%%%%%%%%%%%%%%%%%%%%%%%%%%%%%%%%%%%%%%%%%%%%%%%%%%%%%b
%%%%%%%%%%%%%%%%          FRONT MATTER          %%%%%%%%%%%%%%%%
%%%%%%%%%%%%%%%%%%%%%%%%%%%%%%%%%%%%%%%%%%%%%%%%%%%%%%%%%%%%%%%%
\frontmatter % start roman numbering

% format the TOC
\renewcommand{\RMAAbannerimage}{rmsc_blank}
\tableofcontents

%% preface, group photo, list of participants go here
%% use the `frontsection' environment (see editors' guide)
\begin{frontsection}{Preface}{%
    \listofauthors{S.~J.~Arthur, N.~S.~Brickhouse, J.~Franco}
    \indexauthor{Arthur, S. J.}
    \indexauthor{Brickhouse, N.}
    \indexauthor{Franco, J.}
    }
  These are the proceedings of a very nice conference that Jane and
  Pepe organised. This is the preface, which explains what a wonderful
  time was had by all and how grateful they are for all the people who
  gave them money or helped out.
  
  This is another paragraph of the preface, in case they have anything
  else to say.  For example, the Scientific Organizing Committee
  consisted of
  \begin{quote}
    J. Franco (Mexico, Chair), D. Cox (USA) and various other notables
  \end{quote}
  while the Local Organizing Committee comprised
  \begin{quote}
    S. J. Arthur (Chair), S. Kurtz, etc.
  \end{quote}

  \EditorSignature

\end{frontsection}

\begin{frontsection}[Photo]{Group Photograph}{\notitle}
\null\vfill
\centerline{Here goes the photograph.}
\vfill
\end{frontsection}

\begin{frontsection}[Participants]{List of Participants}{}
%% Here you put a formatted list of participants
\vfil
\dotfill First page of list of participants \dotfill
\vfil\clearpage

\null\vfil
\dotfill Second page of list of participants \dotfill
\vfil

\end{frontsection}

\begin{frontsection}[XXX]{Another front section}{}

\end{frontsection}

%%%%%%%%%%%%%%%%%%%%%%%%%%%%%%%%%%%%%%%%%%%%%%%%%%%%%%%%%%%%%%%%
%%%%%%%%%%%%%%%%          MAIN MATTER           %%%%%%%%%%%%%%%%
%%%%%%%%%%%%%%%%%%%%%%%%%%%%%%%%%%%%%%%%%%%%%%%%%%%%%%%%%%%%%%%%
\mainmatter % start arabic numbering

\begin{articles}
  %% Main articles go here
  %% Use \article{ARTICLENAME} for each individual article
  %% If you want any text to appear in brackets after the article
  %% title in the TOC and index, then put it as an optional first
  %% argument, e.g.,  \article[Invited Talk]{ARTICLENAME} 

  \session{The Contributed Papers}
  \article{rm-journal-example} % An illustrative article
  \article{rm-onepage} % another article, which has a plate
  \article{rm-minimal-test}
  \session{Conference Summary}
  \article[Invited Review]{rmsummsamp} % An illustrative conference summary

\end{articles}

%%%%%%%%%%%%%%%%%%%%%%%%%%%%%%%%%%%%%%%%%%%%%%%%%%%%%%%%%%%%%%%%
%%%%%%%%%%%%%%%%          BACK MATTER           %%%%%%%%%%%%%%%%
%%%%%%%%%%%%%%%%%%%%%%%%%%%%%%%%%%%%%%%%%%%%%%%%%%%%%%%%%%%%%%%%
\backmatter

%% Abstracts of presentations for which the
%% authors have not submitted a paper
\begin{abstracts}

  \title{An abstract of a poster}
  \author{W. J. Henney\altaffilmark{1}}
  \listofauthors{W.~J.~Henney}
  \indexauthor{Henney, W. J}
  \altaffiltext{1}{Where I am from done with altaffilmark/text}
  \abstract{This is a test abstract.   This is a test abstract.   This
    is a test abstract.   This is a test abstract.}
  

  \title{An abstract of another  poster}
  \author{W.~J.~Henney\altaffil{This is done with altaffil} and
    A.~N.~Other\altaffil{So is this} }
  \listofauthors{W.~J.~Henney \& A.~N.~Other}
  \indexauthor{Henney, W. J}
  \indexauthor{Other, A. N.}
  \abstract{This is a test abstract.   This is a test abstract.   This
    is a test abstract.   This is a test abstract.}

  \title{An abstract of yet another  poster}
  \author{W.~J.~Henney\altaffil{This is done with altaffil} and
    A.~N.~Other\altaffil{So is this} }
  \listofauthors{W.~J.~Henney \& A.~N.~Other}
  \indexauthor{Henney, W. J}
  \indexauthor{Other, A. N.}
  \abstract{This is a test abstract.   This is a test abstract.   This
    is a test abstract.   This is a test abstract.}

  \title{The abstract of a poster with a very long title to see how
    well it gets split up into lines}
  \author{A.~N.~Other\altaffil{The affiliation of A.~N.~Other is also
      very long so that it will go over several lines. His email
      address is anon@nowhere.org}}
  \listofauthors{A.~N.~Other}
  \indexauthor{Other, A. N.}
  \abstract{We will try and make this abstract longer so there is more
    stuff to put on the page and it doesn't look too empty. However, I
    have nothing whatsoever to say so will have to stop here. Let's
    repeat that a few times. We will try and make this abstract longer
    so there is more stuff to put on the page and it doesn't look too
    empty. However, I have nothing whatsoever to say so will have to
    stop here. Let's repeat that a few times. We will try and make
    this abstract longer so there is more stuff to put on the page and
    it doesn't look too empty. However, I have nothing whatsoever to
    say so will have to stop here. Let's repeat that a few times. We
    will try and make this abstract longer so there is more stuff to
    put on the page and it doesn't look too empty. However, I have
    nothing whatsoever to say so will have to stop here. Let's repeat
    that a few times.}
    
  \title[Comment to go in TOC]{The abstract of a second poster with a
    very long title to see how well it gets split up into lines}
  \author{A.~N.~Other\altaffil{The affiliation of A.~N.~Other is also
      very long so that it will go over several lines. His email
      address is anon@nowhere.org}}
  \listofauthors{A.~N.~Other}
  \indexauthor{Other, A. N.}
  \abstract{We will try and make this abstract longer so there is more
    stuff to put on the page and it doesn't look too empty. However, I
    have nothing whatsoever to say so will have to stop here. Let's
    repeat that a few times. We will try and make this abstract longer
    so there is more stuff to put on the page and it doesn't look too
    empty. However, I have nothing whatsoever to say so will have to
    stop here. Let's repeat that a few times. We will try and make
    this abstract longer so there is more stuff to put on the page and
    it doesn't look too empty. However, I have nothing whatsoever to
    say so will have to stop here. Let's repeat that a few times. We
    will try and make this abstract longer so there is more stuff to
    put on the page and it doesn't look too empty. However, I have
    nothing whatsoever to say so will have to stop here. Let's repeat
    that a few times.}
    
  \title{The abstract of a third poster with a very long title to see how
    well it gets split up into lines}
  \author{A.~N.~Other\altaffil{The affiliation of A.~N.~Other is also
      very long so that it will go over several lines. His email
      address is anon@nowhere.org}}
  \listofauthors{A.~N.~Other}
  \indexauthor{Other, A. N.}
  \abstract{We will try and make this abstract longer so there is more
    stuff to put on the page and it doesn't look too empty. However, I
    have nothing whatsoever to say so will have to stop here. Let's
    repeat that a few times. We will try and make this abstract longer
    so there is more stuff to put on the page and it doesn't look too
    empty. However, I have nothing whatsoever to say so will have to
    stop here. Let's repeat that a few times. We will try and make
    this abstract longer so there is more stuff to put on the page and
    it doesn't look too empty. However, I have nothing whatsoever to
    say so will have to stop here. Let's repeat that a few times. We
    will try and make this abstract longer so there is more stuff to
    put on the page and it doesn't look too empty. However, I have
    nothing whatsoever to say so will have to stop here. Let's repeat
    that a few times.}
    
  \title{The abstract of a fourth poster with a very long title to see how
    well it gets split up into lines}
  \author{A.~N.~Other\altaffil{The affiliation of A.~N.~Other is also
      very long so that it will go over several lines. His email
      address is anon@nowhere.org}}
  \listofauthors{A.~N.~Other}
  \indexauthor{Other, A. N.}
  \abstract{We will try and make this abstract longer so there is more
    stuff to put on the page and it doesn't look too empty. However, I
    have nothing whatsoever to say so will have to stop here. Let's
    repeat that a few times. We will try and make this abstract longer
    so there is more stuff to put on the page and it doesn't look too
    empty. However, I have nothing whatsoever to say so will have to
    stop here. Let's repeat that a few times. We will try and make
    this abstract longer so there is more stuff to put on the page and
    it doesn't look too empty. However, I have nothing whatsoever to
    say so will have to stop here. Let's repeat that a few times. We
    will try and make this abstract longer so there is more stuff to
    put on the page and it doesn't look too empty. However, I have
    nothing whatsoever to say so will have to stop here. Let's repeat
    that a few times.}
    
  \title{The abstract of a fifth poster with a very long title to see how
    well it gets split up into lines}
  \author{A.~N.~Other\altaffil{The affiliation of A.~N.~Other is also
      very long so that it will go over several lines. His email
      address is anon@nowhere.org}}
  \listofauthors{A.~N.~Other}
  \indexauthor{Other, A. N.}
  \abstract{We will try and make this abstract longer so there is more
    stuff to put on the page and it doesn't look too empty. However, I
    have nothing whatsoever to say so will have to stop here. Let's
    repeat that a few times. We will try and make this abstract longer
    so there is more stuff to put on the page and it doesn't look too
    empty. However, I have nothing whatsoever to say so will have to
    stop here. Let's repeat that a few times. We will try and make
    this abstract longer so there is more stuff to put on the page and
    it doesn't look too empty. However, I have nothing whatsoever to
    say so will have to stop here. Let's repeat that a few times. We
    will try and make this abstract longer so there is more stuff to
    put on the page and it doesn't look too empty. However, I have
    nothing whatsoever to say so will have to stop here. Let's repeat
    that a few times.}
    
  \title{The abstract of a sixth poster with a very long title to see how
    well it gets split up into lines}
  \author{A.~N.~Other\altaffil{The affiliation of A.~N.~Other is also
      very long so that it will go over several lines. His email
      address is anon@nowhere.org}}
  \listofauthors{A.~N.~Other}
  \indexauthor{Other, A. N.}
  \abstract{We will try and make this abstract longer so there is more
    stuff to put on the page and it doesn't look too empty. However, I
    have nothing whatsoever to say so will have to stop here. Let's
    repeat that a few times. We will try and make this abstract longer
    so there is more stuff to put on the page and it doesn't look too
    empty. However, I have nothing whatsoever to say so will have to
    stop here. Let's repeat that a few times. We will try and make
    this abstract longer so there is more stuff to put on the page and
    it doesn't look too empty. However, I have nothing whatsoever to
    say so will have to stop here. Let's repeat that a few times. We
    will try and make this abstract longer so there is more stuff to
    put on the page and it doesn't look too empty. However, I have
    nothing whatsoever to say so will have to stop here. Let's repeat
    that a few times.}
    
  \title{The abstract of a seventh poster with a very long title to see how
    well it gets split up into lines}
  \author{A.~N.~Other\altaffil{The affiliation of A.~N.~Other is also
      very long so that it will go over several lines. His email
      address is anon@nowhere.org}}
  \listofauthors{A.~N.~Other}
  \indexauthor{Other, A. N.}
  \abstract{We will try and make this abstract longer so there is more
    stuff to put on the page and it doesn't look too empty. However, I
    have nothing whatsoever to say so will have to stop here. Let's
    repeat that a few times. We will try and make this abstract longer
    so there is more stuff to put on the page and it doesn't look too
    empty. However, I have nothing whatsoever to say so will have to
    stop here. Let's repeat that a few times. We will try and make
    this abstract longer so there is more stuff to put on the page and
    it doesn't look too empty. However, I have nothing whatsoever to
    say so will have to stop here. Let's repeat that a few times. We
    will try and make this abstract longer so there is more stuff to
    put on the page and it doesn't look too empty. However, I have
    nothing whatsoever to say so will have to stop here. Let's repeat
    that a few times.}
  
  \title{An abstract to go in the list of abstracts, but not in the
    index of authors}
  \author{I.~A.~M. Hidden\altaffil{Puerto Escondido}}
  \listofauthors{I.~A.~M. Hidden}
  \abstract{The idea is that by omitting the
    \textbackslash{}indexauthor command, we can avoid an entry in the
    author index.}

\end{abstracts}


%% The Index: after latexing twice, run `makeindex -s rmaa' on this
%% file, then latex twice more
\printindex

%% if you rename this file to something other than `rmbooksamp', then
%% you must change the above line accordingly

%%
%% Plates: If you have any, then the following line will print
%% them out:

\outputplates

\end{document}

%%% Local Variables:
%%% mode: latex
%%% TeX-master: t
%%% End:

\end{document}
%</ads-container>
%    \end{macrocode}

% This is the |adshtml.sty| style file used by |ADS.tex|. It just
% defines a few 
% macros that might be used in abstracts, etc (really it should define
% them all). This now gets all its code from
% Sec~\ref{sec:user-commands}. 
%    \begin{macrocode}
%<*adshtml-style>
%% Nothing else needs to go here now. 
%</adshtml-style>
%    \end{macrocode}
% The init file for latex2html.
%    \begin{macrocode}
%<*latex2html-init>
$PDFTOCAIRO = "pdftocairo";
$PDFLATEX = "pdflatex";
$LATEX = "latex";
$MAX_SPLIT_DEPTH = 0;
$HTML_VERSION = "4.0,math";
$NO_SIMPLE_MATH = 1;
$ADDRESS = "<p><address><a href=\"mailto:rmaa\@astro.unam.mx\">rmaa\@astro.unam.mx</address></p>"; 
$INFO = "<p>This document was created using latex2html and the rmaa latex macros</p>"; 
$NO_NAVIGATION = 1;
1; 
%</latex2html-init>
%    \end{macrocode}
% This is the |adshtml.perl| file that is used by \LaTeX{}2HTML. Some
% bizarre \TeX{}y command sometimes show up so we tell it to ignore
% these. All the astronomy macros that come from AAS are processed in
% \TeX{}. That way they get made to images but with an |ALT| tag that
% has the original code, which is what will be picked up by |lynx|,
% with the result that they get passed intact into the final |.txt|
% file. 
%    \begin{macrocode}
%<*adshtml-perl>
package main;

$dum = &do_cmd_makeatletter();

sub do_cmd_RMAA {
    local($_) = @_;
    "RMxAA" . $_;
}

sub do_cmd_RMAASC {
    local($_) = @_;
    "RMxAASC" . $_;
}

sub do_cmd_etal {
    local($_) = @_;
    "et al\." . $_;
}

sub do_cmd_ibidrule {
    local($_) = @_;
    "<img src=\"foo\" alt=\"\ibidrule\">" . $_;
}


sub do_cmd_bibitem {
    local($_) = @_;
    local($key);
    local($opt, $dummy)=&get_next_optional_argument; 
    $key = &missing_braces unless
        ((s/$next_pair_pr_rx/$key = $2; ''/eo)
        ||(s/$next_pair_rx/$key = $2; ''/eo));
    "<p>" . $_; 
}

&ignore_commands (<<_RAW_ARG_CMDS_);
lowercase
adjustfinalcols
textsuperscript # {}
newline # []
thinmuskip
mskip
smallskip
medskip
bigskip
hspace # {}
hspacestar # {}
vspace # {}
vspacestar # {}
mathrm
_at_mathrm
spacefactor
discretionary
unhbox 
voidb_at_x 
penalty 
_at_M
_at_m
null
IeC
_RAW_ARG_CMDS_

&process_commands_in_tex (<<_RAW_ARG_CMDS_);
sun
earth
lesssim
gtrsim
sq
argdeg
arcmin
arcsec
fd
fh
fm
fs
fdg
farcm
farcs
fp
micron
la
ga
case # {} # {}
slantfrac # {} # {}
onehalf
onethird
twothirds
onequarter
threequarters
ubvr
ub
bv
vr
ur
aj
araa
apj
apjl
apjs
ao
apss
aap
aapr
aaps
azh
baas
jrasc
memras
mnras
pra
prb
prc
prd
pre
prl
pasp
pasj
qjras
skytel
solphys
sovast
ssr
zap
nat
iaucirc
aplett
apspr
bain
fcp
gca
grl
jcp
jgr
jqsrt
memsai
nphysa
physrep
physscr
planss
procspie
_RAW_ARG_CMDS_

1;                              # This must be the last line
%</adshtml-perl> 
%    \end{macrocode}

% This is a shell script |makeads.sh| that generates the |_ads.txt| file from
% |_ads.tex| by using \LaTeX{}2HTML and |lynx|.  \changes{4.0}{2019/11/13}{Make sure that lynx writes utf-8 output instead of latin-1.}
% \changes{2.2beta13}{2002/03/12}{Complete rewrite of |makeads.sh|. Now much 
% improved I think}
% \changes{4.0}{2019/11/13}{Fix the at-sign substitution so that it only eliminates \cs{@} and does not mess with email addresses.}

%    \begin{macrocode}
%<*ads-shell-script>
#! /bin/sh
if [ -z $1 ]; then
    echo "Usage: $0 basename"
    exit
fi
BASENAME=$1
ADSBASENAME=${BASENAME}_ads
AUXFILE=$BASENAME.aux
TMPFILE=tmp.$$


echo "Generating $ADSBASENAME.tex"

# For some reason, must double-up \\ to \\\\ but not \ to \\ - odd!
cat > $ADSBASENAME.tex <<EOF
%% PLEASE DO NOT EDIT THIS FILE!  It is generated automatically from
%% the script makeads.sh.  Any hand-edited changes will be lost
%% 
\documentclass{article}
\usepackage[latin1]{inputenc}
\usepackage{html}
\usepackage{adshtml}
\newcommand{\ADSfield}[2]{\%#1  #2\\\\}
\newenvironment{ADSarticle}{}{\par\bigskip}
\setlength{\parindent}{0pt}
\newcounter{istage}
\renewcommand{\ion}[2]{\setcounter{istage}{#2}#1 \Roman{istage}}
\begin{document}
\makeatletter
EOF
echo >> $ADSBASENAME.tex "\input{$BASENAME.ads2}" 
echo >> $ADSBASENAME.tex "\end{document}"



echo "Generating $ADSBASENAME.sed" 

# Find type of journal
grep '\\documentclass\[.*proceedings.*\]' $BASENAME.tex >/dev/null \
     && JOURNAL='RMxAC' \
     || JOURNAL='RMxAA'

# Find the year and volume number
YEAR=`grep '\\SetYear' $BASENAME.tex | sed 's/.*\\SetYear{\(.*\)}.*/\1/'`
VOLUME=`grep '\\SetVolume' $BASENAME.tex | sed 's/.*\\SetVolume{\(.*\)}.*/\1/'`

# Generate an auxilliary sed script 
cat > $TMPFILE <<EOF
s|\\\\newlabel{\(.*\):RMAAFirstPage}{{.*}{\(.*\)}}|s/\\\\\\\\ThisVolume{\1}/$YEAR, RMxAC, $VOLUME, \2/g|g
EOF

# Now go through the .aux file, looking for \newlabels and use the sed 
# script just generated in order to make another sed script, one that 
# will replace \ThisVolume{LABEL} with the correct reference
grep '^\\newlabel{.*:RMAAFirstPage}' $AUXFILE \
     | sed -f $TMPFILE > $ADSBASENAME.sed && rm $TMPFILE

# Apend a further sed command to get rid of \@, since redefining 
# it in _ads.tex didn't work
echo >> $ADSBASENAME.sed 's/\\\\@//g'

echo "Generating $BASENAME.ads2 from $ADSBASENAME.sed and $BASENAME.ads" 

# Apply this to the .ads file, generating the .ads2 file, which is included
# by the _ads.tex file.
sed -f $ADSBASENAME.sed $BASENAME.ads > $BASENAME.ads2

echo "Generating $ADSBASENAME.txt from $ADSBASENAME.tex and $BASENAME.ads2"

# Now run latex2html on the _ads.tex file
latex2html -no_navigation -info 0 -address 0 -html_version 4.0,math $ADSBASENAME

# Then, convert to utf8 with lynx and tidy it up a bit
lynx --justify=off --width=100 --dump --nolist --display_charset=utf-8 \
     $ADSBASENAME/$ADSBASENAME.html | cut -c4- \
     | sed 's/ibidrule/\\ibidrule/g' > $ADSBASENAME.txt
%</ads-shell-script>
%    \end{macrocode}

% \section{Epilogue}
% Err, that's it \dots

\endinput


% \endinput
% Local Variables:
% mode: doctex
% TeX-master: t
% End:
"

rmaa.aux: rmaa.dtx rmaa.ind rmaa.gls
        @echo "REGENERATING .aux FILE"
        pdflatex "\nonstopmode% \def\fileversion{4.0} \def\filedate{2019/11/12}
% \MakeShortVerb{\|}
% \title{A \LaTeX{} document class for the Revista Mexicana de
% Astronom\'\i a y Astrof\'\i sica\thanks{This file has version
% \fileversion{} last revised \filedate{}.}}
% \author{William J. Henney\thanks{w.henney@irya.unam.mx}}
% \date{\filedate}
% \maketitle
% \paragraph{Abstract}
% This document describes the \texttt{rmaa} \LaTeX{}
% document class for the journal ``Revista Mexicana de
% Astronom\'\i a y Astrof\'\i sica''. The macros can be used for 
% \begin{enumerate}
% \item Preparing author manuscripts for the journal (including
% conference proceedings to be published in the ``Serie de
% Conferencias'')
% \item Typesetting an entire volume of the journal by inputting the
% individual article files from a container file. Facilities are
% included for automatic generation of TOC and Author Index. 
% \end{enumerate}
% 
% \tableofcontents
% \section{The documentation driver file}
% This is the part that generates the documentation that you may or
% may not be reading, which is done with the help of the Makefile
% discussed in section~\ref{sec:driver_make}. 
%    \begin{macrocode}
%<*driver>
\documentclass{ltxdoc}
\usepackage{hyperref}
\usepackage{palatino}
\EnableCrossrefs \CodelineIndex \RecordChanges
\setlength{\parindent}{0pt}
\def\MacroFont{\footnotesize\tt}
\begin{document} \DocInput{rmaa.dtx} \PrintIndex \PrintChanges \end{document}
%</driver>
%    \end{macrocode}

% \section{The implementation}
% 
% \changes{3.24}{2007/09/14}{Added user access to the macro version
% and date via \cs{rmVersion} and \cs{rmDate}}
% 
%    \begin{macrocode}
%<*package>
\newcommand{\rmaa@version}{3.26}
\newcommand{\rmaa@date}{2011/08/15}
\newcommand\rmVersion{\rmaa@version}
\newcommand\rmDate{\rmaa@date}
\NeedsTeXFormat{LaTeX2e}[1995/06/01]
%    \end{macrocode}
% \changes{4.0}{2019/11/02}{Change abstract to appear before
% resumen. Change RevMexAA to RMxAA for consistency with ADS.}
%
% \changes{3.22}{2007/03/22}{Allow proceedings articles to have
%   received/accepted dates} \changes{3.21}{2006/05/15}{Changes for
%   Silvia - abstracts are now shipped out to ADS.}
% \changes{3.19}{2006/01/19}{Changes for Silvia - remove adresses from
%   end, new abstract style} \changes{2.2beta}{2001/10/30}{ Fixed bug
%   whereby twocolumn option was not passed through to book
%   class. Finally sorted out the ``label on last page'' issue. }
% \changes{2.1}{2001/01/23}{ Fixed some inconsistencies in
%   \cs{nocite} and \cs{bibitem} uniqification Fixed the definition of
%   \cs{ion} so it will work in math mode and in contexts where the
%   typesize is different from normal (e.g., title, fig captions).
%   Modified the definition of \cs{@caption}, so we don't write any
%   list-of-figures data (we never use it anyway, and it just causes
%   problems with fragile commands in fig captions). }
% \changes{2.0}{2000/11/30}{ Added `newstyle' option, which should be
%   used in all new proceedings volumes (wider abstract, smaller
%   typeface in fig captions, etc.) Added commands for formatting
%   table captions. New demonstation document rmpostsamp.tex for
%   1-page posters. This will be the last release under a PRI
%   government (for a while...)}  \changes{1.9}{2000/06/06}{
%   Corrected position of page number in footer of articles that open
%   on even pages. Ensured that first article in each session opens on
%   an odd page, as do list of abstracts and author index. Added extra
%   space in TOC before the backmatter. Added command for editors,
%   \cs{BoldSessionTitles}, that makes produces bold titles for
%   session in the TOC.}  \changes{1.8}{2000/05/16}{ Fixed bug in
%   last page label with \cs{suppressfulladdresses} }
% \changes{1.7}{2000/05/11}{ Fixed more little things. Page labels
%   now work, even if you have no bibliography. You can output the
%   full addresses by hand with the \cs{outputfulladdresses}
%   command. Editors' Guide updated but still not finished. Revamped
%   splitbook.sh now produces PS file of each article for web page and
%   also an EPS file of each page for sending to the printers.}
% \changes{1.6}{2000/04/15}{ More cleaning up: full stops after
%   titles in index, running headers and footers in frontmatter,
%   fontsize in title of index, etc} \changes{1.5}{2000/04/15}{
%   Eliminated superscripts from titles in TOC and index}
% \changes{1.4}{2000/04/14}{ Added copyright and license info (GPL)}
% \changes{1.3}{2000/03/01}{ Corrected \cs{figcaption} to do the
%   caption formatting right Corrected font size of sec headings in
%   front matter} \changes{1.2}{2000/02/16}{ Added \cs{plate}
%   command} \changes{1.1}{2000/02/14}{ Added manuscript option,
%   corrected bugs} \changes{1.0}{2000/02/08}{ Now does a whole book
%   - lots of stuff tidied up} \changes{0.7}{1999/10/04}{ Increased
%   spacing between sections in single column styles. Suppressed
%   floats at bottom of pages (didn't work though) }
% \changes{0.6}{1999/09/24}{ Tweaked \cs{adjustfinalcols} (still not
%   perfect) Added proceedings option Added onecolumn/twocolumn
%   options (orthogonal to journal/proceedings) }
% \changes{0.5}{1999/09/23}{ Added \cs{fulladdresses} and
%   \cs{adjustfinalcols} (just about working)}
% \changes{0.4}{1999/09/22}{ Minor stylistic changes}
% \changes{0.3}{1999/06/28}{ Better mechanism for alternative
%   affiliations Added `guide' option for user guide}
% \changes{0.2}{1999/05/27}{ All basic functionality present, I
%   think} \changes{0.1}{1999/05/26}{ Just about working - still
%   missing Fig/Tab/Refs} \changes{0.0}{1999/05/21}{ So bored today,
%   decided to rewrite RMAA LaTeX macros from scratch}
%    \begin{macrocode}
\ProvidesClass{rmaa}[\rmaa@date\space v\rmaa@version\space RMxAA article
class]
%    \end{macrocode}
% \subsection{Option Handling}
% \label{sec:options}
% Declarations that are needed in options
% \changes{3.05}{2003/10/22}{Added natbib option}
% \changes{3.06}{2003/01/03}{Natbib option turned on permanently}
% \changes{3.07}{2003/01/11}{Added proof option}
% \changes{3.14}{2003/01/11}{Option to turn off natbib}
% \changes{3.18}{2003/10/22}{trying out doctex mode} 
% \changes{3.23}{2007/09/09}{Manuscript is now on by default.}
% \changes{3.26}{2011/08/15}{New option for review articles}
%    \begin{macrocode}
\newif\if@reviewarticle
\@reviewarticlefalse
\newif\if@proof
\@prooffalse
\newif\if@journal
\@journaltrue
\newif\if@printdates
\@printdatesfalse
\newif\if@proceedings
\@proceedingsfalse
\newif\if@manuscript
\@manuscripttrue
\newif\if@bigfont
\@bigfonttrue
\newif\if@guide
\@guidefalse
\newif\if@preprint
\@preprintfalse
\newif\if@rmtwocol
\@rmtwocoltrue
\newif\if@biblist
\@biblistfalse
\newif\if@book
\@bookfalse
\newif\if@natbib
\@natbibtrue
\newif\if@rmcite
\@rmcitefalse
\newif\if@titlevisible
\@titlevisibletrue
\newif\if@abstractpresent
\@abstractpresentfalse
\newif\if@printbanner
\@printbannerfalse
\newif\if@suppressfulladdresses
\@suppressfulladdressesfalse
\newif\if@newstyle
\@newstylefalse
\newif\if@debug
\@debugfalse
\newif\if@boldkeywords
\@boldkeywordsfalse
%    \end{macrocode}
% The option handling---mainly setting various booleans.
% \changes{2.2beta11}{2001/11/26}{Made twocol and
% newstyle the default for proceedings.}
% \changes{3.23}{2007/09/09}{Manuscript is now 10 pt instead of 12pt}
% \changes{3.23}{2007/09/10}{Full addresses are now automatically
% suppressed for proceedings style.}
% \changes{3.24}{2007/09/14}{New option \texttt{boldkeywords} that
% does the keywords in the old ugly way. Now off by default.}
%    \begin{macrocode}
\DeclareOption{manuscript}{\typeout{Manuscript for referee} 
  \@manuscripttrue\@journalfalse\@rmtwocolfalse
  \PassOptionsToClass{10pt}{article}}
\DeclareOption{bigfont}{\typeout{Big font manuscript requested by Gary
  Ferland} 
  \@manuscripttrue\@journalfalse\@rmtwocolfalse\@bigfonttrue
  \PassOptionsToClass{11pt}{article}}
\DeclareOption{journal}{\typeout{Main journal option}%
  \@journaltrue\@rmtwocoltrue\@printdatestrue\@manuscriptfalse
  \PassOptionsToClass{10pt}{article}}
\DeclareOption{guide}{\typeout{User guide option}%
  \@guidetrue\@journalfalse\@rmtwocoltrue\@manuscriptfalse
  \PassOptionsToClass{10pt}{article}}
\DeclareOption{preprint}{\typeout{Preprint option}%
  \@preprinttrue\@manuscriptfalse\PassOptionsToClass{10pt}{article}}
\DeclareOption{proof}{\typeout{Author proof option}%
  \@prooftrue\@manuscriptfalse}
\DeclareOption{proceedings}{\typeout{Proceedings option}%
  \@proceedingstrue\@journalfalse\@rmtwocoltrue\@newstyletrue
  \@suppressfulladdressestrue
  \@printdatesfalse\@manuscriptfalse\PassOptionsToClass{10pt}{article}} 
\DeclareOption{printdates}{\typeout{Force printing of received/accepted dates}%
  \@printdatestrue} 
\DeclareOption{book}{\typeout{Whole issue option}%
  \@booktrue} 
\DeclareOption{nonatbib}{\typeout{Turn off natbib support (required
    for compatibility with old rmaacite package).}%
  \@natbibfalse} 
\DeclareOption{newstyle}{\typeout{Use new layout style}%
  \@newstyletrue}
\DeclareOption{oldstyle}{\typeout{Use old layout style}%
  \@newstylefalse}
\DeclareOption{debug}{\typeout{Loads of annoying messages}%
  \@debugtrue}
%    \end{macrocode}
% These two need to be declared explicitly so that |\ExecuteOptions|
% actually does something.
%    \begin{macrocode}
\DeclareOption{twocolumn}{\@rmtwocoltrue}
\DeclareOption{onecolumn}{\@rmtwocolfalse}
\DeclareOption{twoside}{\PassOptionsToClass{\CurrentOption}{article}}
\DeclareOption{boldkeywords}{\@boldkeywordstrue}
%    \end{macrocode}
% Bold, upper-case keywords used to be the default. Now they must be
% turned on explicitly with the \texttt{boldkeywords} option. 
%    \begin{macrocode}
\DeclareOption{boldkeywords}{\@boldkeywordstrue}
%    \end{macrocode}
% Pass any other option along to the \texttt{article} class.
%    \begin{macrocode}
\DeclareOption*{\PassOptionsToClass{\CurrentOption}{article}}
%    \end{macrocode}
% Default options are to be two-sided and based on the standard
% \LaTeX{} \texttt{article} class,
% unless we are processing the entire book, in which case we base
% ourselves on the standard \texttt{book} class. Default is to use
% manuscript option unless otherwise specified. 
%    \begin{macrocode}
\ExecuteOptions{twoside,manuscript} \ProcessOptions\relax
\if@book
\if@rmtwocol\LoadClass[twocolumn]{book}\else\LoadClass{book}\fi
\newcommand{\abstract}{\relax}
\newcommand{\@maketitle}{\relax}
\else
\if@rmtwocol\LoadClass[twocolumn]{article}\else\LoadClass{article}\fi
\fi
%    \end{macrocode}
% Attempt to use natbib for citations 02 Jan 2003. 
% \changes{3.05}{2003/01/02}{If the natbib option is used, then set
% things up so that natbib's version of the bibliography is identical
% to the standard one.}
% \changes{3.09}{2003/01/23}{Explicitly require a recent version of
% natbib, so that we get a decent error message if the package is too
% old.} 
% \changes{3.23}{2007/09/09}{Now allow Bibtex to work properly} 
% \changes{3.25}{2009/10/01}{Change \cs{newcommand} tp \cs{def} when setting \cs{bibfont}. This allows it to work with all versions of natbib.} 
% \changes{3.25}{2010/02/06}{Updated copyright year.} 
% \changes{3.25}{2010/02/06}{Updated email and web addresses to \texttt{crya.unam.mx}.} 
%    \begin{macrocode}
\if@natbib
\RequirePackage{natbib}[1999/02/25]
\bibpunct{(}{)}{;}{a}{}{,}
\def\bibfont{\small}
\renewcommand\bibsection{\section*{REFERENCES}}
\setlength\bibsep{0pt}
\setlength\bibhang{1.5em}
\bibliographystyle{rmaa}
\fi
%    \end{macrocode}
% We want the \texttt{graphicx} package for the extended syntax of
% |\includegraphics|. We want the \texttt{ifthen} package for using in
% these macros (although we don't much yet). We explicitly load
% \texttt{babel} since some \LaTeX{} distributions messed up the
% hyphenation if we didn't.  
%    \begin{macrocode}
\RequirePackage{graphicx}
\RequirePackage{ifthen}
\IfFileExists{babel.sty}{%
%   \let\@restoreindent=\@afterindentfalse
  \RequirePackage[english]{babel}
%   \@restoreindent
  %% reinstate the definition of \@roman so as not to break .idx files
  \@ifundefined{org@@roman}{\relax}{\renewcommand{\@roman}[1]{\org@@roman{##1}}}
  }{%
  \typeout{Package babel.sty not found - don't worry about this}}
%    \end{macrocode}
% Try to load the \texttt{dcolumn} and \texttt{booktabs} packages if
% they can be found. These can produce nicer-looking tables. 
%    \begin{macrocode}
\IfFileExists{dcolumn.sty}{%
  \RequirePackage{dcolumn}
  }{%
  \typeout{Package dcolumn.sty not found - you will not be able to
    align columns on decimal points in tables}}
\IfFileExists{booktabs.sty}{%
  \RequirePackage{booktabs}
  }{%
  \typeout{Package booktabs.sty not found - your tables will not look
    as pretty as they ought too.}
  }
\IfFileExists{makeidx.sty}{%
  \RequirePackage{makeidx}
}{%
  \typeout{Package makeidx.sty not found - Author Index may not work.}
  }
%% commented out for moment since requires more modern version of LaTeX
%% \RequirePackage[latin1]{inputenc}
%    \end{macrocode}
% \changes{4.0}{2019/11/13}{Make \texttt{hyperref} and
% \texttt{etoolbox} into hard requirements, since we now use them
% extensively.}  We use hyperref to get internal and external links.
% But I don't like the colored boxes so I turn those off.  We use
% etoolbox for patching and robustifying at the moment.  In the
% future, I want to port a lot of other stuff over to it.  For
% instance, replace the ifthen package and the TeX ifs too.
%
% We target the latest version (1.7) of PDF, rather than the default
% of 1.5.  This reduces the chance of getting warnings from included
% graphics.  This is not strictly related to hyperref, but we put it
% here anyway.
%    \begin{macrocode}
\RequirePackage{etoolbox}
\RequirePackage{hyperref}
\hypersetup{hidelinks}
\pdfminorversion=7
%    \end{macrocode}
% We patch the \cs{@combinedblfloats} following a suggestion by Heiko
% Oberdiek at \url{https://tex.stackexchange.com/a/249743/3843}, which
% avoids problems when hyperlinks are split across
% pages. \changes{4.0}{2019/11/12}{Fix for issue where breaking
% hyperlink across pages gave an error.}
%    \begin{macrocode}
\urlstyle{same}
\newcommand\email{\begingroup \urlstyle{rm}\Url}
\patchcmd\@combinedblfloats{\box\@outputbox}{\unvbox\@outputbox}{}{%
  \errmessage{\noexpand\@combinedblfloats could not be patched}%
}%
%    \end{macrocode}
% \subsection{Other setup stuff}\label{sec:setup}
% Some declarations: 
%    \begin{macrocode}
\newif\if@needaddresses
\@needaddressestrue
\newif\if@frontmatter
\@frontmatterfalse
\newif\if@inpreamble
\@inpreambletrue
\newif\if@firstauthor
\@firstauthortrue
\newif\if@firstaffil
\@firstaffiltrue
\newif\if@firstaltaffil
\@firstaltaffiltrue
\newif\if@firstkwd
\@firstkwdtrue
\newif\if@noexpandADS
\@noexpandADSfalse
\newif\if@sessiontitlepages
\@sessiontitlepagesfalse
%    \end{macrocode}
% \begin{macro}{\@rmaaEmpty}
% \begin{macro}{\@rmaaArg}
% These are used in various commands to check if arg was left blank
%    \begin{macrocode}
\newcommand{\@rmaaEmpty}{}
\newcommand{\@rmaaArg}{}
%    \end{macrocode}  
% \end{macro}
% \end{macro}
%
% \subsection{Compatibility with old \LaTeX{} versions}
% \label{sec:compat}
% This package is based on the 1998/12/01 release of \LaTeX{}.
% Older releases may not define all the commands that we need.
% Luckily, we can use |\providecommand| to define these if necessary. 
% The following commands are not defined in the 1994/06/01 release. 
%
% \begin{macro}{\MakeUpperCase}
% \begin{macro}{\textsuperscript}
% \begin{macro}{\textbackslash}
% User commands:
%    \begin{macrocode}
\providecommand{\MakeUpperCase}[1]{\uppercase{#1}}
\providecommand{\textsuperscript}[1]{\ensuremath{^\textrm{#1}}}
\providecommand{\textbackslash}{\ensuremath{\backslash}}
%    \end{macrocode}  
% \end{macro}
% \end{macro}
% \end{macro}
% \begin{macro}{\protected@xdef}
% \begin{macro}{\restore@protect}
% \begin{macro}{\@unexpandable@protect}
% \begin{macro}{\hb@xt@}
% \LaTeX{} internal commands:
%    \begin{macrocode}
\providecommand{\protected@xdef}{%
        \let\@@protect\protect
        \let\protect\@unexpandable@protect
        \afterassignment\restore@protect
        \xdef}
\providecommand{\restore@protect}{\let\protect\@@protect}
\providecommand{\@unexpandable@protect}{\noexpand\protect\noexpand}
\providecommand{\hb@xt@}{\hbox to}
%    \end{macrocode}
% \end{macro}
% \end{macro}
% \end{macro}
% \end{macro}
%
% \subsection{Page layout}\label{sec:page-layout}
%
% The dimensions of the text were defined by Elizabeth. Text width is
% 169\,mm with equal margins on the two sides: 
%    \begin{macrocode}
\setlength{\textwidth}{169mm}
\setlength{\@tempdima}{\paperwidth}
\addtolength{\@tempdima}{-\textwidth}
\setlength{\oddsidemargin}{0.5\@tempdima}
\addtolength{\oddsidemargin}{-1in}
\setlength{\evensidemargin}{\oddsidemargin}
%    \end{macrocode}
% Text height is 225\,mm, top margin is 14\,mm. 
%    \begin{macrocode}
\setlength{\textheight}{225mm} 
\setlength{\topmargin}{-1in} 
\addtolength{\topmargin}{14mm}
%    \end{macrocode}
% Space between columns in two-column styles, header, footer,
% paragraph indentation:
%    \begin{macrocode}
\setlength{\columnsep}{9.5mm}
\setlength{\headsep}{0.7cm}
\setlength{\footskip}{0.75cm}
\setlength{\parindent}{1.5em}
%    \end{macrocode}
% 
% \begin{macro}{\@set@scalable@length}
% Try out a new scheme for dealing with various spacing dimensions and
% allowing the user to easily modify them. The command
% \begin{quote}
% |\set@scalable@length|\marg{name}\marg{defaultlength}\marg{stretch}\marg{shrink}
% \end{quote}
% does a 
% |\setlength|\marg{\cs{name}}\marg{defaultlength plus stretch minus shrink}
% and also defines the user command
% \textbackslash\meta{name}|Scale|\marg{factor} that can 
% be used to rescale length by \meta{factor}. Additionally, now
% defines the package command
% \textbackslash\meta{name}|@set@scale|\marg{FACTOR} 
% that also rescales the length by \meta{FACTOR}, so that the length
% becomes \meta{factor} $\times$ \meta{FACTOR} $\times$
% \meta{defaultlength}.   
%    \begin{macrocode}
\newcommand\@set@scalable@length[4]{%
  %% arg #1: NAME of length without leading `\'
  %%
  %% arg #2: default VALUE  
  \expandafter\newlength\csname @#1@default\endcsname
  \expandafter\setlength\csname @#1@default\endcsname{#2}%
  \expandafter\newlength\csname @@#1@@default\endcsname
  \expandafter\setlength\csname @@#1@@default\endcsname{#2}%
  %% arg #3: default STRETCH  
  \expandafter\newlength\csname @#1@stretch\endcsname
  \expandafter\setlength\csname @#1@stretch\endcsname{#3}%
  \expandafter\newlength\csname @@#1@@stretch\endcsname
  \expandafter\setlength\csname @@#1@@stretch\endcsname{#3}%
  %% arg #4: default SHRINK  
  \expandafter\newlength\csname @#1@shrink\endcsname
  \expandafter\setlength\csname @#1@shrink\endcsname{#4}%
  \expandafter\newlength\csname @@#1@@shrink\endcsname
  \expandafter\setlength\csname @@#1@@shrink\endcsname{#4}%
  %%
  %% User-defined scale \@NAME@scale is unity by default
  \expandafter\gdef\csname @#1@scale\endcsname{1}%
  %% Package-defined scale \@NAME@SCALE is unity by default
  \expandafter\gdef\csname @#1@SCALE\endcsname{1}%
  %%
  %% internal command \@reset@NAME should expand to 
  %% \setlength\@@NAME@@default{\@NAME@SCALE\@NAME@default}
  %% \setlength\NAME{\@NAME@scale\@@NAME@@default}
  %% in other words: scale x SCALE x default 
  \expandafter\gdef\csname @reset@#1\endcsname{%
    \expandafter\setlength\csname @@#1@@default\endcsname{%
      \csname @#1@SCALE\endcsname \csname @#1@default\endcsname
      }
    \expandafter\setlength\csname @@#1@@stretch\endcsname{%
      \csname @#1@SCALE\endcsname \csname @#1@stretch\endcsname
      }
    \expandafter\setlength\csname @@#1@@shrink\endcsname{%
      \csname @#1@SCALE\endcsname \csname @#1@shrink\endcsname
      }
    \expandafter\setlength\csname #1\endcsname{%
      \csname @#1@scale\endcsname \csname @@#1@@default\endcsname
      \@plus
      \csname @#1@scale\endcsname \csname @@#1@@stretch\endcsname
      \@minus
      \csname @#1@scale\endcsname \csname @@#1@@shrink\endcsname
      }}
  %%
  %% User command \NAMEScale should expand to 
  %% {\renewcommand\@NAME@scale{#1}\@reset@NAME}
  \expandafter\gdef\csname #1Scale\endcsname##1{%
    \expandafter\gdef\csname @#1@scale\endcsname{##1}%
    \csname @reset@#1\endcsname}
  %%
  %% Package command \NAME@set@SCALE should expand to 
  %% {\renewcommand\@NAME@SCALE{#1}\@reset@NAME}
  \expandafter\gdef\csname #1@set@SCALE\endcsname##1{%
    \expandafter\gdef\csname @#1@SCALE\endcsname{##1}%
    \csname @reset@#1\endcsname}
  %%
  %% initialize at default values with \@reset@NAME
  \csname @reset@#1\endcsname
  }
%    \end{macrocode}
% \end{macro}
% \begin{macro}{\parskip}
% First use of this new infrastructure: set interparagraph space to be
% zero but with a small amount of stretchiness.
%    \begin{macrocode}
%   \ifthenelse{\boolean{@manuscript}}{%
%     \@set@scalable@length{parskip}{0.5\baselineskip}{0.3\p@}{0.1\p@}
%   }{
    \@set@scalable@length{parskip}{0\p@}{0.3\p@}{0\p@}
%   }
%    \end{macrocode}
% \end{macro}
% 
% \begin{macro}{\abovesectionskip}
% \changes{4.0}{2019/11/12}{Use negative values for skip so that first paragraph after section headings is not indented.  This is something that I have been trying to fix for years!  It turns out that the anser was in the documentation for the standard command \cs{@startsection.}}
% \begin{macro}{\belowsectionskip}
% \begin{macro}{\aboveparaskip}
% \begin{macro}{\belowparaskip}
% Define some utility rubber lengths for use in
% section~\ref{sec:section-commands}. 
%    \begin{macrocode}
\newlength{\abovesectionskip}
\newlength{\belowsectionskip}
\newlength{\aboveparaskip}
\newlength{\belowparaskip}
%    \end{macrocode}
% The default values for these lengths are those traditionally used
% for the onecolumn style. 
%    \begin{macrocode}
\@set@scalable@length{abovesectionskip}{-3.5ex}{-4ex}{-2ex}
\@set@scalable@length{belowsectionskip}{2.3ex}{0ex}{2ex}
\@set@scalable@length{aboveparaskip}{3.25ex}{2ex}{0.5ex}
\@set@scalable@length{belowparaskip}{-1em}{\z@}{\z@}
%    \end{macrocode}
% \end{macro}\end{macro}\end{macro}\end{macro}
% We also make scalable some of the standard float parameters. These
% are the same defaults as in the standard article classes, but
% slightly more stretchy. For the twocolumn style, they are currently
% halved, see below. Note that I have also lumped |\abovecaptionskip|
% in with these, since it plays a similar role. There is also a
% |\belowcaptionskip| but it is set to zero in the standard classes. 
%    \begin{macrocode}
\@set@scalable@length{floatsep}{12\p@}{4\p@}{4\p@}
\@set@scalable@length{textfloatsep}{20\p@}{4\p@}{8\p@}
\@set@scalable@length{intextsep}{12\p@}{4\p@}{4\p@}
\@set@scalable@length{dblfloatsep}{12\p@}{4\p@}{4\p@}
\@set@scalable@length{dbltextfloatsep}{20\p@}{4\p@}{8\p@}
\@set@scalable@length{abovecaptionskip}{10\p@}{\z@}{\z@}
%    \end{macrocode}
% \begin{macro}{\@before@title@skip}         
% \begin{macro}{\@before@author@skip}        
% \begin{macro}{\@between@author@skip}       
% \begin{macro}{\@before@abstract@skip}      
% \begin{macro}{\@after@abstracthead@skip}   
% \begin{macro}{\@between@abstract@skip}     
% \begin{macro}{\@before@kwd@skip}           
% \begin{macro}{\@before@dates@skip}           
% Spacing parameters used in |\@maketitle|. These are not changed by
% default except by the user, so they don't need to go into
% |\@package@default@lengths|.  
%    \begin{macrocode}
\newlength{\@before@title@skip}         
\newlength{\@before@author@skip}        
\newlength{\@between@author@skip}       
\newlength{\@before@abstract@skip}      
\newlength{\@after@abstracthead@skip}   
\newlength{\@between@abstract@skip}     
\newlength{\@before@kwd@skip}           
\newlength{\@before@dates@skip}           
\@set@scalable@length{@before@title@skip}{2em}{\z@}{\z@}
\@set@scalable@length{@before@author@skip}{1.5em}{\z@}{\z@}
\@set@scalable@length{@between@author@skip}{0.5em}{\z@}{\z@}
\@set@scalable@length{@before@abstract@skip}{3ex}{1ex}{0.5ex}
\@set@scalable@length{@after@abstracthead@skip}{1ex}{\z@}{\z@}
\@set@scalable@length{@between@abstract@skip}{2ex}{1ex}{0.5ex}
\@set@scalable@length{@before@kwd@skip}{1ex}{0.5ex}{0.2ex}
\@set@scalable@length{@before@dates@skip}{2ex}{0.5ex}{0.2ex}
%    \end{macrocode}
% \end{macro}\end{macro}\end{macro}\end{macro}
% \end{macro}\end{macro}\end{macro}\end{macro}
% 
% Spacing above and below displayed equations also wants to be done
% like this, but it is more complicated since the standard article
% classes make these change with the fontsize. Something to do later,
% maybe. 
% 
%    \begin{macrocode}
%    \end{macrocode}
% \begin{macro}{\@package@default@lengths}
% The value actually used for these lengths is controlled by the
% \meta{SCALE} parameter, which depends on the style option that is in
% effect. This is set using the |\@package@default@lengths| command,
% which must be called every time the style changes (e.g., changing
% from onecolumn to twocolumn). 
%    \begin{macrocode}
\newcommand{\@package@default@lengths}{%
  \ifthenelse{\boolean{@rmtwocol}}{%
    %% twocolumn mode
    \abovesectionskip@set@SCALE{0.5}
    \belowsectionskip@set@SCALE{0.5}
    \aboveparaskip@set@SCALE{0.5}
    \belowparaskip@set@SCALE{1}
    \floatsep@set@SCALE{0.5}
    \textfloatsep@set@SCALE{0.5}
    \intextsep@set@SCALE{0.5}
    \dblfloatsep@set@SCALE{0.5}
    \dbltextfloatsep@set@SCALE{0.5}
    }{%
    %% onecolumn mode
    \abovesectionskip@set@SCALE{1}
    \belowsectionskip@set@SCALE{1}
    \aboveparaskip@set@SCALE{1}
    \belowparaskip@set@SCALE{1}
    \floatsep@set@SCALE{1}
    \textfloatsep@set@SCALE{1}
    \intextsep@set@SCALE{1}
    \dblfloatsep@set@SCALE{1}
    \dbltextfloatsep@set@SCALE{1}
    }%
  }
\@package@default@lengths
%    \end{macrocode}
% \end{macro}
% \begin{macro}{\RescaleLengths}
% \begin{macro}{\RescaleSecLengths}
% \begin{macro}{\RescaleFloatLengths}
% \begin{macro}{\RescaleTitleLengths}
% User commands that will fiddle with all the scalable lengths. 
%    \begin{macrocode}
\newcommand{\RescaleLengths}[1]{%
  \RescaleSecLengths{#1}\RescaleFloatLengths{#1}\RescaleTitleLengths{#1}}
\newcommand{\RescaleSecLengths}[1]{%
%% sections
  \abovesectionskipScale{#1}
  \belowsectionskipScale{#1}
  \aboveparaskipScale{#1}
  }
\newcommand{\RescaleFloatLengths}[1]{%
%% floats
  \floatsepScale{#1}
  \textfloatsepScale{#1}
  \intextsepScale{#1}
  \dblfloatsepScale{#1}
  \dbltextfloatsepScale{#1}
  \abovecaptionskipScale{#1}
  }
\newcommand{\RescaleTitleLengths}[1]{%
%% maketitle
  \@before@title@skipScale{#1}           
  \@before@author@skipScale{#1}        
  \@between@author@skipScale{#1}       
  \@before@abstract@skipScale{#1}      
  \@after@abstracthead@skipScale{#1}   
  \@between@abstract@skipScale{#1}     
  \@before@kwd@skipScale{#1}           
  \@before@dates@skipScale{#1}           
  }
%    \end{macrocode}
% \end{macro}\end{macro}\end{macro}\end{macro}
% \changes{2.2beta14}{2002/03/17}{Finished implementation of new
% infrastructure for scalable lengths. Converted the above/below
% section skips to use it. Also, all the floatsep parameters. Also,
% all the vertical lengths that are used by |\@maketitle|. }

% \begin{macro}{\MSextramargin}
% For the \texttt{manuscript} option we have an extra 0.3\,inch
% vertical margin and 1\,inch horizontal margin for marginal notes (by
% hand or with |\marginpar|. 
%
% \changes{3.23}{2007/09/09}{Reduce baselinestretch in
%   \texttt{manuscript} style to 1.2. No longer use raggedright for
%   manuscript.}
%
% \changes{3.24}{2007/09/14}{Further reduction of baselinestretch in
%   \texttt{manuscript} style to 1.1 since we now have wider margins.}
%
% \changes{3.24}{2007/09/14}{Split \texttt{manuscript}'s extra margin
%   into separate horizontal and vertical parts.}
% 
%    \begin{macrocode}
\newlength{\MSextraHmargin}\newlength{\MSextraVmargin}
\setlength{\MSextraHmargin}{1in} 
\setlength{\MSextraVmargin}{0.3in} 
\if@manuscript
   \renewcommand{\baselinestretch}{1.1}
   \addtolength{\textwidth}{-2\MSextraHmargin}
   \addtolength{\oddsidemargin}{\MSextraHmargin}
   \setlength{\evensidemargin}{\oddsidemargin}
   \addtolength{\textheight}{-2\MSextraVmargin}
   \addtolength{\topmargin}{\MSextraVmargin}
%    \raggedright
\fi
%    \end{macrocode}
% \end{macro}
% Float parameters: don't allow bottom floats by default. 
%    \begin{macrocode}
\renewcommand{\fps@figure}{tp}
\renewcommand{\fps@table}{tp}
%    \end{macrocode}
%
% \begin{macro}{\rmaa@AuthorWidthFrac}
% \begin{macro}{\rmaa@AbstractWidthFrac}
% These parameters control what fraction of the total text width can
% be used by the list of authors and the abstract/resumen,
% respectively.  
%    \begin{macrocode}
\newcommand{\rmaa@AuthorWidthFrac}{1.0}
\newcommand{\rmaa@AbstractWidthFrac}{0.75}
%    \end{macrocode}
% Use a wide abstract to save space in the new style.
%    \begin{macrocode}
\if@newstyle\renewcommand{\rmaa@AbstractWidthFrac}{1.0}\fi
%    \end{macrocode}
% \end{macro}
% \end{macro}
%
% Some things to consider in future versions. Trying to avoid too many
% widows, etc. Apparently, the following will work best with
% \cs{flushbottom}. 
%    \begin{macrocode}
%%\widowpenalty=10000
%%\clubpenalty=10000
%    \end{macrocode}
% Also, should perhaps try getting rid of underfull/overfull box warnings. 
% \changes{3.10}{2003/01/24}{Set large values for \cs{hbadness},
% \cs{vbadness}, \cs{hfuzz}, \cs{vfuzz}. To reduce amount of stuff
% that gets written to the terminal every time you compile something.}
%    \begin{macrocode}
\hbadness=10000
\vbadness=10000
\setlength\hfuzz{\maxdimen}
\setlength\vfuzz{\maxdimen}
%    \end{macrocode}
% \subsection{Various text styles for different contexts}
% \label{sec:text-styles} 
% \begin{macro}{\rmaa@AffilStyle}
% \begin{macro}{\rmaa@PagesStyle}
% \begin{macro}{\rmaa@YearStyle}
% \begin{macro}{\rmaa@VolumeStyle}
% \begin{macro}{\rmaa@VolumenStyle}
% \begin{macro}{\rmaa@JournalNameStyle}
% \begin{macro}{\rmaa@ReceivedStyle}
% \begin{macro}{\rmaa@AbstractHeaderStyle}
% \begin{macro}{\rmaa@JournalNameStyle}
% Style commands for the different fields used by |\maketitle| in
% constructing the title page of the article. These have the form
% |\rmaa@XYZStyle{TEXT}|. I'm not sure these were really a good idea,
% but never mind \dots 
%    \begin{macrocode}
\newcommand{\rmaa@AffilStyle}[1]{{\normalsize\rmfamily#1}}
\newcommand{\rmaa@PagesStyle}[1]{{\normalsize\rmfamily#1}}
\newcommand{\rmaa@YearStyle}[1]{{\normalsize\rmfamily#1}}
\newcommand{\rmaa@VolumeStyle}[1]{{\normalsize\rmfamily\bfseries#1}}
\newcommand{\rmaa@VolumenStyle}[1]{{\Large\sffamily\bfseries#1}}
\newcommand{\rmaa@JournalNameStyle}[1]{{\normalsize\itshape#1}}
\newcommand{\rmaa@ReceivedStyle}[1]{{\normalsize\rmfamily\itshape#1}}
\newcommand{\rmaa@AbstractHeaderStyle}[1]{%
  {\normalsize\rmfamily\MakeUpperCase{#1}}}
\newcommand{\rmaa@TitleStyle}[1]{%
  {\large\bfseries\rmfamily\boldmath\MakeUpperCase{#1}\par}}
%    \end{macrocode}
% \end{macro}
% \end{macro}
% \end{macro}
% \end{macro}
% \end{macro}
% \end{macro}
% \end{macro}
% \end{macro}
% \end{macro}
%
% \begin{macro}{\rmaa@FrontSecStyle}
% Headings of ``front sections'' (see section~\ref{sec:front-sections})
%    \begin{macrocode}
\newcommand{\rmaa@FrontSecStyle}[1]{%
  {\normalsize\bfseries\rmfamily\MakeUpperCase{#1}}}
%    \end{macrocode}
% \end{macro}
%
% \begin{macro}{\rmaa@TOCTitleStyle}
% \begin{macro}{\rmaa@TOCPartStyle}
% \begin{macro}{\rmaa@IndexTitleStyle}
% \begin{macro}{\rmaaTOCAuthorStyle}
% \begin{macro}{\rmaaIndexAuthorStyle}
% Commands to set style of TOC and index entries.
% Rewritten to avoid unnecessary font commands in the .toc and .ind files
%    \begin{macrocode}
\newcommand{\rmaa@TOCTitleStyle}[1]{\MakeUpperCase{#1}}
\newcommand{\rmaa@TOCPartStyle}[1]{\rmaa@TOCTitleStyle{#1}}
\newcommand{\rmaa@IndexTitleStyle}[1]{#1}
\newcommand{\rmaaTOCAuthorStyle}[1]{%
  {\itshape#1}}
\newcommand{\rmaaIndexAuthorStyle}[1]{%
  \protect\mbox{\bfseries#1}}
%    \end{macrocode}
% \end{macro}
% \end{macro}
% \end{macro}
% \end{macro}
% \end{macro}
%
% \begin{macro}{\rmaa@KeywordStyle}
% \begin{macro}{\rmaa@KeywordHeaderStyle}
% The style of the keywords depends on whether we are using the
% \texttt{newstyle} option or not. If we are, then they are small. 
%    \begin{macrocode}
\ifthenelse{\boolean{@newstyle}}{
  \newcommand{\rmaa@KeywordStyle}[1]{%
    {\small%
      \if@boldkeywords 
      \bfseries\MakeUpperCase{\raggedright #1}
      \else
      \raggedright #1
      \fi
    }}
  \newcommand{\rmaa@KeywordHeaderStyle}[1]{%
    {\small\itshape#1}}}
{
  \newcommand{\rmaa@KeywordStyle}[1]{%
    {\normalsize%
      \if@boldkeywords\bfseries\fi\raggedright #1}}
  \newcommand{\rmaa@KeywordHeaderStyle}[1]{%
    {\normalsize\itshape#1}}}
%    \end{macrocode}
% \end{macro}
% \end{macro}
%
% \subsection{Page styles}\label{sec:page-styles}
% \begin{macro}{\ps@rmaa}
% The running headers for normal pages:
%    \begin{macrocode}
\newcommand{\ps@rmaa}{%
%    \end{macrocode}
% nothing in either foot; 
%    \begin{macrocode}
  \renewcommand\@oddfoot{\relax}%
  \renewcommand\@evenfoot{\relax}%
%    \end{macrocode}
% page number and short author list in the even heads; 
%    \begin{macrocode}
  \renewcommand\@evenhead{\rlap{\thepage}\hfil\@shortauthor\hfil}%
%    \end{macrocode}
% short title and page number in the odd heads.
%    \begin{macrocode}
  \renewcommand\@oddhead{\hfil\@shorttitle\hfil\llap{\thepage}}
  }
%    \end{macrocode}
% \begin{macro}{\ps@rmaafirstpage}
% The header of the first page of each
% article:
%    \begin{macrocode}
\newcommand{\ps@rmaafirstpage}{%
  \renewcommand\@oddfoot{\hfil\thepage}%
%    \end{macrocode}
% footer corrected (06 Jun 2000) to have page on outside edge;
%    \begin{macrocode}
  \renewcommand\@evenfoot{\thepage\hfil}%
%    \end{macrocode}
% Header has journal name, volume, and page numbers for a regular
% article, or similar for special cases. 
% \changes{4.0}{2019/11/03}{Change to three-line header with copyright info and DOI.}
%    \begin{macrocode}
\if@preprint
\renewcommand\@oddhead{%
  \if@proceedings 
  \textsc{To appear in ``\@conftitle{} (\rmaa@YearStyle{\@year})''}\hfil
  \textit{RMxAC}
  \else
  \textsc{To appear in} \@journalname\hfil
  \fi}\fi%
\if@guide
\renewcommand\@oddhead{\textsc{Documentation} (\rmaa@YearStyle{\@year})\hfil}
\fi%
\if@manuscript
\renewcommand\@oddhead{Manuscript for \@journalname{} (\rmaa@YearStyle{\@year})\hfil}
\fi%
\if@journal
\renewcommand\@oddhead{%
  \parbox{\textwidth}{
    \@journalname, \rmaa@VolumeStyle{\@volume},
    \rmaa@PagesStyle{\@artpages}\
    (\rmaa@YearStyle{\@year})\hfil\\[1pt]
    {\small\textcopyright{} \@year:
      Instituto de Astronom\'\i{}a,
      Universidad Nacional Aut\'onoma de M\'exico}\hfil\\[1pt]
    {\small\@DOI@url}\hfil
  }%
}\fi%
  \renewcommand\@evenhead{\@oddhead}}
%    \end{macrocode}
% By default, use the |\ps@rmaa| page style.
%    \begin{macrocode}
\pagestyle{rmaa}
%    \end{macrocode}
% \end{macro}
% \end{macro}
%
% \subsection{Preamble commands}\label{sec:preamble} 
% Everything you need for setting the title page, TOC and index
% entries, etc.
% \subsubsection{Commands for the title page}
% \label{sec:title-commands}
% \begin{macro}{\resumen}
% \begin{macro}{\@resumen}
% \begin{macro}{\abstract}
% \begin{macro}{\@abstract}
% User commands for setting the resumen and abstract. 
% These should be used in the document header. The internal commands
% |\@resumen| and |\@abstract| are used to hold the actual text
% pending output via the |\maketitle| command. 
%    \begin{macrocode}
\newif\if@setresumen
\newcommand{\resumen}[1]{%
%%  \renewcommand{\@resumen}{{\selectlanguage{spanish}#1}}%
  \renewcommand{\@resumen}{{\ignorespaces#1}}%
  \@setresumentrue}
\renewcommand{\abstract}[1]{%
  \renewcommand{\@abstract}{{\ignorespaces#1}}
  \renewcommand\ADS@abstract{#1}
%% Not a good idea, since then author macros are not expanded and
%% cause errors later. However, by not doing it we are vulnerable to
%% fragile commands.....
  \if@noexpandADS\@onelevel@sanitize{\ADS@abstract}\fi
  \@abstractpresenttrue
  }
\newcommand{\ADSnoexpand}{\@noexpandADStrue}
%    \end{macrocode}
% \begin{macro}{\boldabstract}
% This is the version for doing 1-page conference proceedings. Still
% save the text for ADS but then write it out immediately as a bold
% paragraph. 
%    \begin{macrocode}
\newcommand{\boldabstract}[1]{%
  \gdef\ADS@abstract{#1}%
  \renewcommand{\@abstract}{#1}%
  \if@noexpandADS\@onelevel@sanitize{\ADS@abstract}\fi%
  \noindent\textbf{\boldmath #1}
  }
%    \end{macrocode}
% \end{macro}
% Default value in each case is to print out a warning if not set
%    \begin{macrocode}
  \newcommand{\@resumen}{%
    \@setresumenfalse \ClassWarningNoLine{rmaa}{No Spanish abstract
      (resumen) specified - if you do not speak Spanish, this will be
      translated by the editors}%
    \textbf{Favor de proporcionar un resumen en espa\~nol. If you are
      unable to translate your abstract into Spanish, the editors will
      do it for you.} }
  \newcommand{\@abstract}{\ClassWarningNoLine{rmaa}{No abstract
      specified}}
%    \end{macrocode}
% \end{macro}
% \end{macro}
% \end{macro}
% \end{macro}
%
% \begin{macro}{\title}
% \begin{macro}{\@rawtitle}
% \begin{macro}{\@TOCtitle}
% \begin{macro}{\rmaaIndexTitle}
% Similarly, |\title| is the preamble command to set the title of the
% article. This is used in various places, typeset in different
% styles, hence the use of the internal |\@rawtitle|, |\@title|,
% |\rmaaIndexTitle|, and |\@TOCtitle| commands, each of which uses the
% relevant style from section~\ref{sec:text-styles}
%    \begin{macrocode}
\newcommand{\@rawtitle}{Raw Title}
\newcommand{\@TOCtitle}{\rmaa@TOCTitleStyle{Title for TOC}}
\newcommand{\rmaaIndexTitle}{\rmaa@IndexTitleStyle{Title for index}.}
\newcommand{\@TOCaftertitle}{}
\renewcommand{\title}[1]{%
%%  {\ADS@sanitize\xdef\@rawtitle{#1}}
  \renewcommand\@rawtitle{\noexpand#1}
  \let\ADS@title=\@rawtitle
  \renewcommand{\@title}{\rmaa@TitleStyle{#1}}%
  \renewcommand{\@TOCtitle}{\rmaa@TOCTitleStyle{#1}\@TOCaftertitle}%
  \renewcommand{\rmaaIndexTitle}{\rmaa@IndexTitleStyle{#1\@TOCaftertitle}.}%
  }
%    \end{macrocode}
% \end{macro}
% \end{macro}
% \end{macro}
% \end{macro}
% 
% All the weird and whacky ways of indicating the author
% affiliations. 
% \begin{macro}{\affil}
% First, the way that sets it directly underneath the authors name.
% Note that |\affil| trumps all other affiliations for the purposes of
% getting into the ADS stuff (we use |\gdef|).   
% \dots 
%    \begin{macrocode}
\newcommand{\affil}[1]{%
  \\[1.5ex]\rmaa@AffilStyle{#1}\\[1.5ex]%
  \if@firstaffil\gdef\ADS@affil{#1}\@firstaffilfalse\fi%
  }
\renewcommand{\and}{and\\[1.5ex]}
%    \end{macrocode}
% \begin{macro}{\altaffil}
% Alternative affiliations are footnoted. 
%    \begin{macrocode}
\newcommand{\altaffil}[1]{%
  \if@firstaltaffil\gdef\ADS@affil{#1}\@firstaltaffilfalse\fi%
  \thanks{#1}}
%    \end{macrocode}
% \begin{macro}{\altaffilmark}
% \begin{macro}{\altaffiltext}
% Alternative-alternative mechanism for affiliations---should not be
% mixed with |\altaffil| (unless you know what you are doing) or there
% will be problems with numbering.
%    \begin{macrocode}
\newcommand{\altaffilmark}[1]{\textsuperscript{#1}}
\newcommand{\altaffiltext}[2]{%
  \if@firstaltaffil\gdef\ADS@affil{#2}\@firstaltaffilfalse\fi
  \stepcounter{footnote}%
  \protected@xdef\@thanks{\@thanks
    \protect\footnotetext[#1]{#2}}} 
%    \end{macrocode}
% \end{macro}
% \end{macro}
% \end{macro}
% \end{macro}
%
% \begin{macro}{\keywords}
% \begin{macro}{\@keywords}
% \begin{macro}{\@keywordheader}
% \label{keywords}Sundry minor preamble commands. All
% self-explanatory, I hope. |\keywords| is the old user interface
% (deprecated). New articles should use |\addkeyword| once for each
% keyword. That way the ADS record comes out right too (comma rather
% than em-dash separator).
% \changes{2.2beta12}{2001/12/07}{\cs{addkeyword} New user interface
%   for keywords.} 
% \changes{2.2beta13}{2002/03/11}{\cs{addkeyword} Improved behaviour
%   with .ads file. Use
%   \cs{ADS@sanitize} to change definitions of `~' and \cs{ion}.}
% \changes{3.23}{2007/09/09}{\cs{addkeyword} now signals a compilation
% error if the author tries to add more than 6 keywords.}
% \changes{3.24}{2007/09/14}{kwdcounter is now reset in
% \cs{@begofarticlehook} to prevent error in book style} 
%    \begin{macrocode}
\newcommand{\@kwdsep}{ --- }
\newcommand{\ADS@kwdsep}{, }
\newcounter{kwdcounter}\setcounter{kwdcounter}{0}
\newcommand{\addkeyword}[1]{%
  \addtocounter{kwdcounter}{1}
  \ifthenelse{\thekwdcounter > 6}{
    \ClassError{rmaa}{Too many keywords. A maximum of 6 keywords is allowed.}{}
  }{}
  \if@firstkwd%
  \renewcommand\@keywords{\if@boldkeywords\uppercase{#1}\else #1\fi}
  {\ADS@sanitize\def\ion{\plainion}\gdef\ADS@keywords{#1}}
  \@firstkwdfalse
  \else%
  \edef\@keywords{\@keywords\@kwdsep \if@boldkeywords\uppercase{#1}\else #1\fi}%
  {\ADS@sanitize\def\ion{\plainion}\xdef\ADS@keywords{\ADS@keywords\ADS@kwdsep #1}}%
  \fi}
\newcommand{\ADS@sanitize}{%
  \def\mathrm{}
  \def\@{\ }
  \def~{\ }
  \def\langle{<}
  \def\rangle{>}
  \def\log{log }
  \def\sin{sin }
  \def\cos{cos }
  \def\tan{tan }
  \def\fd{^d.}
  \def\fh{^h.}
  \def\fm{^m.}
  \def\fs{^s.}
  \def\fdg{�.}
  \def\farcm{'.}
  \def\farcs{''.}
  \def\arcdeg{�}
  \def\,{ }
  }
\newcommand{\keywords}[1]{%
  \renewcommand{\ADS@keywords}{#1}
  \renewcommand{\@keywords}{\uppercase{#1}}
  }
\newcommand{\@keywordheader}{\rmaa@KeywordHeaderStyle{Key Words:}}
\newcommand{\@keywords}{%
  \ClassWarningNoLine{rmaa}{No keywords specified}%
  Keywords go here
  }
%    \end{macrocode}
% \begin{macro}{\shorttitle}
% \begin{macro}{\shortauthor}
% \begin{macro}{\@shorttitle}
% \begin{macro}{\@shortauthor}
% Commands for setting the running headers.
%    \begin{macrocode}
\newcommand{\shorttitle}[1]{%
  \renewcommand\@shorttitle{\MakeUpperCase{#1}}}
\newcommand{\shortauthor}[1]{%
  \renewcommand\@shortauthor{\MakeUpperCase{#1}}}
\newcommand{\@shorttitle}{SHORT TITLE GOES HERE}
\newcommand{\@shortauthor}{SHORT AUTHOR GOES HERE}
%    \end{macrocode}
% \end{macro}
% \end{macro}
% \end{macro}
% \end{macro}
% \end{macro}
% \end{macro}
% \end{macro}


% \subsubsection{The full addresses}
% \label{sec:full-addresses}
% \begin{macro}{\fulladdresses}
% \begin{macro}{\@fulladdresses}
% \begin{macro}{\@fulladdheight}
% \begin{macro}{\@fulladddepth}
% \begin{macro}{\@fulladdbox}
% Aaghh, the full addresses! A perennial source of problems. 
%    \begin{macrocode}
\newlength{\@fulladdheight}
\newlength{\@fulladddepth}
\newsavebox{\@fulladdbox}
%% \fulladdresses command modified (11 May 2000) so that we can put
%% \fulladdresses{} in the article to suppress the output of the
%% addresses
\newcommand{\@FinalFigure}{\relax}
\newcommand{\FinalFigure}[1]{\renewcommand{\@FinalFigure}{#1}}
\newcommand{\fulladdresses}[1]{%
  \renewcommand{\@rmaaArg}{#1}
  \ifx\@rmaaArg\@rmaaEmpty
     \renewcommand{\@fulladdresses}{}
  \else
     \renewcommand{\@fulladdresses}{
       \normalsize%
       \ifx\@FinalFigure\relax\else\vfill\@FinalFigure\vfill\fi
%%       \vspace*{1ex}
       %% do addresses as a list (copied from thebibliography environment) 
       \list{\null}{\leftmargin 1.5em\topsep\z@\labelwidth\z@\labelsep\z@\itemindent -1.5em
         \itemsep\z@\parsep\z@\usecounter{enumi}}
       \sloppy\clubpenalty4000\widowpenalty4000
       \sfcode`\.=1000\relax
       #1
       \endlist
       }
  \fi
  }
\newcommand{\@fulladdresses}{Full addresses go here}
%    \end{macrocode}
% \begin{macro}{\suppressfulladdresses}
% New user-level command for suppressing full addresses should you so
% desire. Currently used in the single-page-poster style. 
%    \begin{macrocode}
\newcommand{\suppressfulladdresses}{\@suppressfulladdressestrue\@needaddressesfalse}
%    \end{macrocode}
% \end{macro}
% \end{macro}
% \end{macro}
% \end{macro}
% \end{macro}
% \end{macro}

% \begin{macro}{\printdates}
% \begin{macro}{\noprintdates}
% Command to activate the printing out of the received/accepted
% dates. This can also be done using the \texttt{printdates} option,
% which is the default for \texttt{journal}, but not for \texttt{proceedings}.
%    \begin{macrocode}
\newcommand{\printdates}{\@printdatestrue}
\newcommand{\noprintdates}{\@printdatesfalse}
%    \end{macrocode}
% \end{macro}
% \end{macro}
% \begin{macro}{\@outputaddressesifnecessary}
% Command to output the addresses if we haven't already done so
% This is called by: 
% \begin{enumerate}
% \item The end of |thebibliography| environment
% \item The |\outputfulladdresses| command
% \end{enumerate}
%    \begin{macrocode}
\newcommand{\@outputaddressesifnecessary}{%
  \ClassInfo{rmaa}{Calling @outputaddressesifnecessary}
  \ifthenelse{\boolean{@needaddresses}}{%
    %% If in one column style, or in the first column of a two column
    %% style then we can try to put it on this page 
    \ifthenelse{\boolean{@firstcolumn} \or \not\boolean{@rmtwocol}}{%
      \@outputaddresses}{%
      %% otherwise give warning it will be on a separate page
      \ClassWarning{rmaa}{Full addresses will go on a separate
        page (please see User Guide for how to remedy this)}}%
    %% If we haven't managed yet, put it on its own page
    \ifthenelse{\boolean{@needaddresses}}{%
      \if@debug{\ClassWarning{rmaa}{In this silly bit!}}\fi
      \clearpage\vfill\@outputaddresses}{}%
    }{
    %% If we don't need them, we still might need to do the eoa hook 
    \ifthenelse{\boolean{@suppressfulladdresses}}{%
      \if@debug{\ClassWarning{rmaa}{XXX}}\fi
%%      \if@book\else\@endofarticlehook\null\fi
      }{}
    }%
  }  
%    \end{macrocode}
% \end{macro}

% \begin{macro}{\outputfulladdresses}
% A user interface to the above command, in case there is no biblio
% (11 May 2000) 
%    \begin{macrocode}
\newcommand{\outputfulladdresses}{\@outputaddressesifnecessary}
%    \end{macrocode}
% \end{macro}

% \begin{macro}{\adjustfinalcols}
% This command starts a new column and should be 
% used to ensure that the columns balance on the last page. It also
% outputs the full 
% author addresses to ensure that this occurs in the left hand
% column. 
% The |\par| command is so that the current line will be taken account
% of when working out |\pagetotal|
%    \begin{macrocode}
\newcommand{\adjustfinalcols}{\if@twocolumn\goodbreak\@outputaddresses\newpage\noindent\fi}
%    \end{macrocode}
% \end{macro}

% \begin{macro}{\@outputaddresses}
% Command to do the work of outputting the addresses. 
% This can be called by 
% \begin{enumerate}
% \item an explicit |\adjustfinalcols|
% \item via |\@outputaddressesifnecessary|, q.v.
% \end{enumerate}
%
%    \begin{macrocode}
\newcommand{\@outputaddresses}{%
  % Allow option of suppressing the full addresses, in which case
  % we just run the hook if we are not in a book style
  \if@suppressfulladdresses
%%     \if@book\else\@endofarticlehook\null\fi
  \else
  %% save the full addresses in a box for output later
  \sbox{\@fulladdbox}{\parbox{\textwidth}{\@fulladdresses}}
  % calculate height of address box  
  \settoheight{\@fulladdheight}{\usebox{\@fulladdbox}}
  \settodepth{\@fulladddepth}{\usebox{\@fulladdbox}}
  \addtolength{\@fulladdheight}{\@fulladddepth}
  \if@debug
  { \typeout{Address total height: \the\@fulladdheight}
    \typeout{Text height: \the\pagetotal}
    \typeout{Goal height: \the\pagegoal}
    }
  \fi
  % reuse \@fulladdheight to mean the total height of text+addresses
  \addtolength{\@fulladdheight}{\pagetotal}
%%  \addtolength{\@fulladdheight}{-\baselineskip}
  \if@debug\typeout{Adjusted text + addresses height: \the\@fulladdheight}\fi
  % check it will fit on current page
  \ifdim\@fulladdheight<\pagegoal
     %% must shift it to the left
     %% to compensate for the hanging indentation in the reference list. 
     %% NOTE: this is not foolproof!
     %% NB. This will give overfull hbox warnings, which is a bit ugly.
     \vfill\noindent\if@biblist\hspace*{-\leftmargin}\fi
     \usebox{\@fulladdbox}
     % Set this to show we've done it
     \@needaddressesfalse
     % and assume we are at the end of the article
     % so that we do the EOA hook if not in the book style
%%     \if@book\else\@endofarticlehook\fi
  \else
     % if not, just do nothing. With luck, this function will get 
     % called again on the next page
     \if@debug{\typeout{Holding over full addresses to next page}}\fi
  \fi\fi
  }
%    \end{macrocode}
% \end{macro}

% \subsubsection{Commands for editors}
% \label{sec:editor-commands}
% Some user/editor commands for doing the TOC and index in the book
% option. 
% They are defined here rather in the book-specific part since they
% need to be always available in order to avoid breaking documents. 
% \begin{macro}{\listofauthors}
% \begin{macro}{\@rawauthors}
% \begin{macro}{\rmaaTOCauthors}
% The command |\listofauthors| formats an author list for use in the
% TOC. It should contain the same names as in the |\author| command,
% but without the affiliation garbage and using |\&| instead of
% |and|. 
%    \begin{macrocode}
\newcommand{\@rawauthors}{Raw Authors}
\newcommand{\rmaaTOCauthors}{\noexpand\rmaaTOCAuthorStyle{Author list for TOC}}
\newcommand{\listofauthors}[1]{%
  \renewcommand{\@rawauthors}{#1}
  \renewcommand\rmaaTOCauthors{\noexpand\rmaaTOCAuthorStyle{#1}}}
%    \end{macrocode}
% \begin{macro}{\indexauthor}\label{indexauthor}
% The command |\indexauthor| produces an index entry for each
% author. 
% This command must come \emph{after} the
% |\title| and the |\listofauthors| commands. Also, there should be one 
% |\indexauthor| command for \emph{each} author. We also add the
% author onto the end of the |\ADS@authors| macro and construct the
% |\ADS@authorinitial| by a most cunning method (see below), but only
% for the first author. 
% 
% No, No, NO! This was all wrong. Backed out completely: 
% \begin{quote}
% Completely rewritten (19 Dec 2002) since we now use the index class, so 
% as to be able to alphabetize accented characters. A side effect of this is
% that |\index| no longer expands out the argument after the |@|, so we 
% instead add all the |\index| commands to a list (|\indexauthor@deferlist|), 
% which is executed in the |\@begofarticlehook|. 
% 
% Note the stuff with |\toks0| is copied from Exercise 20.15 of the \TeX{}book---it is
% to make sure the latin1-accented name gets passed into the idx
% file. 
% \end{quote}
% We now go back to the old way of doing things, using an optional
% argument to the \cs{indexauthor} command that has the name without
% accents. 
% \changes{3.08}{2003/01/16}{Remerged changes from 3.06 that got
% dropped from 3.07 after a synchronization disaster. Basically, the
% accentless optional argument to \cs{indexauthor}. Must add the macro
% directories to my unison set-up to stop this sort of thing from happening.}
%    \begin{macrocode}
\newcommand\@isortkey{}
\newcommand\@emptykey{empty}
\newcommand\indexauthor[2][empty]{
  \renewcommand\@isortkey{#1}
  \ifx\@isortkey\@emptykey \renewcommand\@isortkey{#2}\fi
  \index{\@isortkey @{{\rmaaIndexAuthorStyle{#2}}\hskip1em plus1em minus0.4em
      \rmaaIndexTitle\hskip1em plus2em minus0.4em\penalty-100
      \rmaaTOCauthors}}
  \indexauthorADS{#2}
}
%%\def\indexauthor@deferlist{\relax}
%%\long\def\indexauthor#1{
%%  \def\@@name{#1} \toks0=\expandafter{\@@name}
%%  \toks1=\expandafter{\indexauthor@deferlist}
%%  \protected@xdef\indexauthor@deferlist{%
%%    \the\toks1
%%    \noexpand\index{\the\toks0@{\noexpand\rmaaIndexAuthorStyle{\the\toks0}\noexpand\hskip1em plus1em minus0.4em 
%%        \rmaaIndexTitle\noexpand\hskip1em plus2em minus0.4em\noexpand\penalty-100
%%        \rmaaTOCauthors}}
%%  }
%%  \indexauthorADS{#1}
%%  }
%%\xdef\indexauthor@deferlist{}
\newcommand\indexauthorADS[1]{%
  \if@firstauthor\protected@xdef\ADS@authorinitial{\firstletter{#1}}\fi
  \protected@xdef\ADS@authors{\ADS@authors\if@firstauthor\else ; \fi
    #1}
  \@firstauthorfalse
}
%    \end{macrocode}
% Note that the first |#1| sets the key for the index, making sure it
% actually starts with  
% first letter of the surname. Meanwhile, the |@{...}| expression,
% which starts with some  
% control sequences is what is written to the |.ind| file. 
% See the makeindex man page for details. 
% \end{macro}
% \end{macro}
% \end{macro}
% \end{macro}

% \begin{macro}{\KillAllAccents}
% \begin{macro}{\firstletter}
% These are utility functions used in the above. |lett| is a counter
% indicating the current letter in the string. |\firstletter| just
% prints the first letter of its argument.
%    \begin{macrocode}
\newcommand{\KillAllAccents}{\def\'{}\def\`{}\def\~{}\def\.{}\def\i{i}}
\newcommand\firstletter[1]{\@car #1\@nil}
%    \end{macrocode}
% \end{macro}
% \end{macro}
% \begin{macro}{\SetVolume}
% \begin{macro}{\SetMSnumber}
% \begin{macro}{\SetFirstPage}
% \begin{macro}{\SetYear}
% \begin{macro}{\ReceivedDate}
% \begin{macro}{\AcceptedDate}
% These are for the editors to set the volume number, etc. 
%    \begin{macrocode}
\newcommand{\SetVolume}[1]{%
  \renewcommand{\@volume}{#1}
  \edef\ADS@volume{\PadToFourWithDots{#1}}
  \renewcommand{\@volumen}{\rmaa@VolumenStyle{volumen #1}}%
  }
\newcommand{\SetMSnumber}[1]{\renewcommand{\@MSnum}{#1}}
%    \end{macrocode}
% Note that there is no real reason to use |\SetFirstPage| since all
% the page numbering is done automatically in the book style. 
%    \begin{macrocode}
\newcommand{\SetFirstPage}[1]{\setcounter{page}{#1}}
\newcommand{\SetYear}[1]{%
  \renewcommand{\@year}{#1}
  \renewcommand{\ADS@year}{#1}
  \renewcommand{\@anyo}{\rmaa@VolumenStyle{#1}}
  }
\newcommand{\ReceivedDate}[1]{\renewcommand{\@received}{#1}}
\newcommand{\AcceptedDate}[1]{\renewcommand{\@accepted}{#1}}
%    \end{macrocode}
%
% \begin{macro}{\DOI}
% User/editor command to enter the DOI, which will be typeset as a link.
% \changes{4.0}{2019/11/02}{New user command \cs{DOI} to typeset
% DOI in the header of first page and add it to ADS info.}
%    \begin{macrocode}
\newcommand\DOI[1]{%
  \def\@DOI@url{\url{https://doi.org/10.22201/ia.01851101p.#1}}%
  \def\ADS@doi{https://doi.org/10.22201/ia.01851101p.#1}%
}%
%    \end{macrocode}
% \end{macro}
%
% \begin{macro}{\ORCIDiD}
% \changes{4.0}{2019/11/12}{Add mechanism for giving ORCID iD links}
% \begin{macro}{\ORCIDiDfull}
% \begin{macro}{\@orcid@icon}
% \begin{macro}{\@orcid@iconbig}
% Commands to add icon links to ORCID iDs of authors. For example
% \begin{quote}
% |\ORCIDiD|\marg{0000-0000-0000-0000}
% \end{quote}
% to be used after each author name in the |\author|
% command and in the |\fulladdresses|.  Alternatively,
% |\OORCIDiDfull| will typeset the icon plus the full text of the
% link, but we don't use that at present.  The icon image is stored in |ORCIDiD_iconvector.pdf|, which needs to be on the input path.  In principle, we could make EPS versions if necessary.
%    \begin{macrocode}
\newcommand\@orcid@icon{\includegraphics[height=1.8ex]{ORCIDiD_iconvector}}
\newcommand\@orcid@iconbig{\includegraphics[height=2ex]{ORCIDiD_iconvector}}
\newcommand\ORCIDiD[1]{\href{https://orcid.org/#1}{\,\@orcid@icon}}
\newcommand\ORCIDiDfull[1]{%
  \href{https://orcid.org/#1}{\@orcid@iconbig}\,\url{https://orcid.org/#1}}
%    \end{macrocode}
% \end{macro}
% \end{macro}
% \end{macro}
% \end{macro}
%
% \begin{macro}{\@volume}
% \begin{macro}{\@volumen}
% \begin{macro}{\@artpages}
% \begin{macro}{\@year}
% \begin{macro}{\@anyo}
% \begin{macro}{\@received}
% \begin{macro}{\@accepted}
% \begin{macro}{\@MSnum}
% Default values for the containers used by above commands, plus a
% few more useful ones. 
%    \begin{macrocode}
\newcommand{\@volume}{00}
\newcommand{\@volumen}{\rmaa@VolumenStyle{00}}
\newcommand{\@artpages}{%
  \oldpageref{\articlename:RMAAFirstPage}--\oldpageref{\articlename:RMAALastPage}}
\newcommand{\@year}{\the\year}
\newcommand{\@anyo}{\rmaa@VolumenStyle{2007}}
\newcommand{\@journalname}{\rmaa@JournalNameStyle{%
    \if@proceedings\RMAASClong\else\RMAAlong\fi}}
\newcommand{\@received}{}
\newcommand{\@accepted}{}
\newcommand{\@MSnum}{0000}
\newcommand{\@DOI@prefix}{https://doi.org/10.22201/ia.01851101p}
\def\@DOI@url{\url{https://doi.org/10.22201/ia.01851101p.2099.00.00.00}}
%    \end{macrocode}
% \end{macro}
% \end{macro}
% \end{macro}
% \end{macro}
% \end{macro}
% \end{macro}
% \end{macro}
% \end{macro}
% \end{macro}
% \end{macro}
% \end{macro}
% \end{macro}
% \end{macro}
% \end{macro}
% 
% The following commands are most useful in the \texttt{book}
% style. However, the |\SetConfTitle| command can now be used in
% individual articles to set the header for the \texttt{preprint}
% style. 
%    \begin{macrocode}
\newcommand{\@reunionplace}{Lugar de la reuni\'on}
\newcommand{\@reuniondate}{Fecha de la reuni\'on}
\newcommand{\@conftitle}{Conference Title}
\newcommand{\SetReunion}[2]{\renewcommand{\@reunionplace}{#1}\renewcommand{\@reuniondate}{#2}}
\newcommand{\SetConfTitle}[1]{\renewcommand{\@conftitle}{#1}}
\newcommand{\@editors}{Juan Ram\'on de la Fuente and ``El Mosh''}
\newcommand{\SetEditors}[1]{\renewcommand{\@editors}{#1}}
\newcommand{\@editores}{Juan Ram\'on de la Fuente y ``El Mosh''}
\newcommand{\SetEditores}[1]{\renewcommand{\@editores}{#1}}
\newcommand{\@mes}{mes de la publicaci\'on}
\newcommand{\SetMes}[1]{\renewcommand{\@mes}{\rmaa@VolumenStyle{#1}}}
%    \end{macrocode}

% \subsection{Notes for Author Proofs and Communication between Editors}
% \label{sec:notes}
% We introduce a command \cs{Note} that should be used to describe
% any change made to the manuscript. These are a bit tricky to get
% right. In general we use \cs{marginpar} but this only works in outer
% paragraph mode, so we just bung the box where we are if we are in a
% figure caption or something. If we are in the preamble, then we just
% use a \marginpar as part of the \maketitle command.
%    \begin{macrocode}
\newif\if@arenotes
\@arenotesfalse
\newcounter{notecounter}
\setcounter{notecounter}{0}
\newcommand\NotesList{}
\newcommand\Note[1]{%
  \if@proof{%
    \@arenotestrue%
    \stepcounter{notecounter}%
    \if@inpreamble\else{%
      \NoteMark{\arabic{notecounter}}%
    }\fi
    \protected@xdef\NotesList{\NotesList 
      \protect\item #1}%
  }\fi%
}
\newcommand\Changed[2]{Changed ``#1'' to ``#2''.}
\newcommand\NoteMark[1]{%
  \ifinner
  \raisebox{1ex}[0pt][0pt]{\makebox[0pt][l]{\framebox{\footnotesize\bfseries
        N#1}}}%
  \else
  \marginpar{\framebox{\footnotesize\bfseries N#1}}%
  \fi
}
\newcommand\RefNote[1]{\if@proof\NoteMark{\ref{#1}}\relax\fi}
\newcommand\NoteLabel{\protect\label}
\newcommand\OutputNotes{%
  \bgroup
  \onecolumn
  \section*{EDITORIAL NOTES}
  \begin{enumerate}
    \NotesList
  \end{enumerate}
  \egroup
}
\DeclareRobustCommand\cs[1]{\texttt{\char`\\#1}}

%    \end{macrocode}
% \subsection{Bibliographic Records for ADS}
% \label{sec:ADS}
% \changes{2.2beta11}{2001/11/27}{Just about finished all the ADS
% stuff now. Except for the references, which I have given up on for
% the time being.}
% \begin{macro}{\ADS@BibCode}
% \begin{macro}{\ADS@year}
% \begin{macro}{\ADS@journal}
% \begin{macro}{\ADS@volume}
% \begin{macro}{\ADS@firstpage}
% \begin{macro}{\ADS@authorinitial}
% We use the |\ADS@BibCode| command to hold the bibliographic code for
% ADS. It is 
% used in the |%R| 
% field of the bibliographic record for each article. 
% The format is 
% \begin{quote}
% \texttt{YYYYJJJJJVVVVMPPPPA}  
% \end{quote}
% where \texttt{YYYY} is the year, \texttt{JJJJJ} is the journal code,
% which is \texttt{RMxAA} or \texttt{RMxAC}, \texttt{VVVV} is the
% volume number (right-justified), \texttt{M} is not used by us,
% \texttt{PPPP} is the first page number (right-justified), and
% \texttt{A} is the initial of the first author's surname. The volume
% and page fields are padded with periods (`\texttt{.}'), which causes
% us some fun (this is done with the macro |\PadToFourWithDots| described
% below) . 
% 
% All these individual bits are held in the commands |\ADS@year|,
% |\ADS@journal|, |\ADS@volume|, |\ADS@firstpage|,
% |\ADS@authorinitial|. 
%    \begin{macrocode}
\newcommand{\SetADSBibCode}[1]{\renewcommand{\ADS@BibCode}{#1}}
\newcommand{\ADS@BibCode}
{\ADS@year\ADS@journal\ADS@volume.\PadToFourWithDots{\ADS@firstpage}\ADS@authorinitial}
\newcommand{\ADS@year}{2001} 
\newcommand{\ADS@journal}{\if@proceedings RMxAC\else RMxAA\fi}
\newcommand{\ADS@volume}{....} 
%    \end{macrocode}
% These two are just set to default values here, |\ADS@firstpage| is set
% in the |document| environment for included articles, in the
% |frontsection| environment, |\tableofcontents| command,  |index| and
% |abstracts| environments. Also by |\AtBeginDocument| for stand-alone
% articles. 
% 
% God knows how we are going to do |\ADS@authorinitial|!
%    \begin{macrocode}
\newcommand{\ADS@firstpage}{....}
\newcommand{\ADS@authorinitial}{A}
%    \end{macrocode}
% \end{macro}
% \end{macro}
% \end{macro}
% \end{macro}
% \end{macro}
% \end{macro}
% We also have a whole slew of other fields that can be set. These are
% all listed in Table~\ref{tab:ADS}. Most of them still have to be
% defined. 
% \begin{table}[!ht]\raggedleft
% \caption{ADS Bibliographic Records}\label{tab:ADS}
% \begin{tabular}{lp{8em}lp{8em}}\hline
% Code & Description & Macro Name & Comment \\ \hline
% |%R| & Bibliographic Code & |\ADS@BibCode| & required \\
% |%A| & Author list & |\ADS@authors| & required, generate from |\indexauthor| \\
% |%F| & Author affiliation & |\ADS@affil| & format: AA(\dots)
% AB(\dots) \dots \\
% |%J| & Journal Name & |\ADS@journalname| &  \\
% |%D| & Publication Date & |\ADS@fulldate| &  \\
% |%L| & Last Page & |\ADS@lastpage| & generate from label \\
% |%T| & Title & |\ADS@title| & required, generate from |\title| \\
% |%C| & Abstract Copyright & |\ADS@copyright| & set to IA-UNAM \\
% |%U| & URL for electronic document & |\ADS@url| & set to PDF file \\
% |%d| & DOI for article & |\ADS@doi| & set by |\DOI| \\
% |%K| & Keywords & |\ADS@keywords| & set equal to |\@keywords| \\
% |%A| & Abstract Text & |\ADS@abstract| & set equal to |\@abstract| \\
% |%R| & References & |\ADS@references| & generate from the
% bibliography  \\ \hline
% \end{tabular}
% \end{table}
% A lot of these macros are set by other commands, as mentioned in the
% Table, but some need to be set  explicitly here. 
% 
% \begin{macro}{\ADS@doi}
% \changes{4.0}{2019/11/08}{Save DOI for ADS}
% The user/editor command \cs{DOI} is used to set this. 
%    \begin{macrocode}
\newcommand{\ADS@doi}{https://doi.org/10.22201/ia.01851101p.0000.00.00.00}
%    \end{macrocode}
% \end{macro}
% 
% \begin{macro}{\ADS@journalname}
% This
% does not contain any accents at the moment so as not to cause
% problems---this maybe should be changed though. 
%    \begin{macrocode}
\newcommand{\ADS@journalname}{%
  \if@proceedings\@conftitle\space (Eds. \@editors) \fi%
  Revista Mexicana de Astronom\'\i{}a y Astrof\'\i{}sica%
  \if@proceedings\space (Serie de Conferencias)\fi%
  \space Vol. \@volume, pp. \ADS@firstpage--\ADS@lastpage\space (\@year)%
  \space (http://www.astroscu.unam.mx/\~{}rmaa/)
  }
%    \end{macrocode}
% \end{macro}
% \begin{macro}{\ADS@keywords}
% \begin{macro}{\ADS@abstract}
% \begin{macro}{\ADS@title}
% |\ADS@abstract| is straightforward, since we already have the text
% saved as |\@abstract|. We could do the same with |\ADS@keywords| but
% to avoid unwanted formatting we explicitly set it in |\keywords|
% (p.~\pageref{keywords}). |\ADS@title| is just set to |\@rawtitle|,
% which we were already using for another purpose. 
%    \begin{macrocode}
\let\ADS@title=\@rawtitle
\newcommand{\ADS@keywords}{No keywords specified}
\newcommand{\ADS@abstract}{No abstract specified}
%    \end{macrocode}
% \end{macro}\end{macro}\end{macro}
% \begin{macro}{\ADS@authors}
% |\ADS@authors| wants to be in ``Other, A. N.'' format, but with
% semi-colons between authors. The most natural way to do this seems to
% be with |\indexauthor| (p.~\pageref{indexauthor}). We append each
% author to this macro, which means we must explicitly clear it for
% each new article (done in |\article|). .   
%    \begin{macrocode}
\newcommand{\ADS@authors}{}
%    \end{macrocode}
% \end{macro}
% \begin{macro}{\ADS@fulldate}
% \begin{macro}{\ADS@lastpage}
% \begin{macro}{\ADS@copyright}
% \begin{macro}{\ADS@website}
% \begin{macro}{\ADS@url}
% \begin{macro}{\SetWebsite}
% \begin{macro}{\SetMonthNumeric}
% Other sundry data, together with user commands for bits that can be
% set by the Editors. The |\ADS@url| is constructed from the website
% (currently |http://www.astrosmo.unam.mx/rmaa/|,
% the journal descriptor (e.g., |RMxAC..10/|), followed by the
% path to the PDF file (e.g., |PDF/RMxAC..10_HEW1.pdf|). 
% 
%    \begin{macrocode}
\newcommand{\ADS@website}{http://www.astrosmo.unam.mx/rmaa/}
\newcommand{\ADS@url}
{\ADS@website\ADS@journal\ADS@volume\ADS@number/PDF/\ADS@journal\ADS@volume\ADS@number\_\articlename.pdf}
\newcommand{\ADS@number}{}
\newcommand{\ADS@fulldate}{00/\ADS@year} 
\newcommand{\ADS@lastpage}{00} 
\newcommand{\ADS@copyright}{(C) \ADS@year; Instituto de Astronom\'\i{}a,
  Universidad Nacional Aut\'onoma de M\'exico} 
\newcommand{\SetNumber}[1]{\renewcommand{\ADS@number}{-#1}}
\newcommand{\SetWebsite}[1]{\renewcommand{\ADS@website}{#1}}
\newcommand{\SetMonthNumeric}[1]{\renewcommand{\ADS@fulldate}{#1/\ADS@year}}
%    \end{macrocode}
% \end{macro}\end{macro}\end{macro}\end{macro}\end{macro}\end{macro}\end{macro}
% \begin{macro}{\ADS@affil}
% \begin{macro}{\ADS@references}
% These two are somewhat more problematical. Affiliation currently
% only works for the first author (see \dots) pending redesign of the
% user interface. References are now done by wrapping the
% |thebibliography| environment to open another input stream on the
% same file and search through for the |\bibitems|. 
%    \begin{macrocode}
\newcommand{\ADS@references}{}
\newcommand{\ADS@affil}{}
%    \end{macrocode}
% \end{macro}\end{macro}
% 
% Define two macros that expand to the strings
% |\begin{thebibliography}| and |\end{thebibliography}|. Note that the
% idea for this and much else is inspired by the code for the docstrip
% package. Also define |\StartBibX| in case authors use the more
% standard |\begin{thebibliography}{}| syntax. 
%    \begin{macrocode}
\edef\EndBib{\expandafter\@gobble\string\\end\string{thebibliography\string}}
\edef\StartBib{\expandafter\@gobble\string\\begin\string{thebibliography\string}}
\edef\StartBibX{\expandafter\@gobble\string\\begin\string{thebibliography\string}\string{\string}}
%    \end{macrocode}
% Define two new logical variables. |\ifAnotherline| is true if we
% want to read another line from the input stream. It is set to false
% when we encounter the |\end{thebibliography}|. |\ifInBib| is set to
% true when we frind the |\begin{thebibliography}| and tells us to
% start saving the following lines. 
%    \begin{macrocode}
\newif\ifAnotherLine
\newif\ifInBib
%    \end{macrocode}
% \begin{macro}{makeOther}
% We need a macro to change the catcodes of lots of things to other. 
%    \begin{macrocode}
\def\makeOther#1{\catcode`#1=12\relax}
%    \end{macrocode}
% \end{macro}
% \begin{macro}{\RereadBiblio}
% Define a macro to do all the fiddly stuff. This is used at the
% beginning of |thebibliography| environment (Sec~\ref{sec:biblio}).
%    \begin{macrocode}
\newread\tf@biblio
\newcommand{\RereadBiblio}{%
%    \end{macrocode}
% Open the current file again for reading
%    \begin{macrocode}
  \begingroup
  \immediate\openin\tf@biblio=\articlename
%    \end{macrocode}
% Everything in a group to be safe. Initialize logicals. Black magic
% to change
% catcodes and avoid blank lines being interpreted as |\par|. 
%    \begin{macrocode}
  \AnotherLinetrue\InBibfalse
  \makeOther\ \makeOther\\\makeOther\$%
  \makeOther\#\makeOther\^\makeOther\^^K%
  \makeOther\_\makeOther\^^A\makeOther\%%
  \makeOther\~\makeOther\{\makeOther\}\makeOther\&%
  \endlinechar-1\relax
%    \end{macrocode}
% Start looping through the file. Read each line to the macro
% |\Line|. If we have got to EOF or end of biblio then set
% |\AnotherLinefalse|.
%    \begin{macrocode}
  \loop
  \read\tf@biblio to \Line
  \ifx\Line\EndBib\AnotherLinefalse\InBibfalse\fi
  \ifeof\tf@biblio \AnotherLinefalse\fi
%    \end{macrocode}
% If we are in the biblio and the line is not blank then append it to
% the |\ADS@references| macro together with a newline. 
%    \begin{macrocode}
  \ifInBib
  \ifx\Line\empty\else\xdef\ADS@references{\ADS@references\Line^^J}\fi
  \fi
%    \end{macrocode}
% Check if the line we just saw was the start of biblio, then repeat
% the loop if necessary. 
%    \begin{macrocode}
  \ifx\Line\StartBib\InBibtrue\fi
  \ifx\Line\StartBibX\InBibtrue\fi
  \ifAnotherLine\repeat
  \immediate\closein\tf@biblio
  \endgroup
  }
%    \end{macrocode}
% \end{macro}
% 
% \begin{macro}{\PadToFourWithDots}
% Utility function for left-padding a number to 4 characters with
% periods. Used in 
% constructing |\ADS@BibCode|. This is a very simple algorithm and
% \emph{requires} that the argument actually be a number---if it is
% not then dire consequences will result.  
%    \begin{macrocode}
\newcommand{\PadToFourWithDots}[1]{%
  \ifnum #1 > 999 #1%
  \else\ifnum #1 > 99 .#1%
  \else\ifnum #1 > 9 ..#1%
  \else ...#1%
  \fi\fi\fi}
%%  \ifnum \count0 > 999 #1%
%%  \else\ifnum \count0 > 99 .#1%
%%  \else\ifnum \count0 > 9 ..#1%
%%  \else ...#1%
%%  \fi\fi\fi}
%    \end{macrocode}
% \end{macro}
%
% \begin{macro}{\AuthorADS}
% \begin{macro}{\ADS@affiliations}
% New mechanism for setting the affiliations for ADS. The author labels (AA, AB, \dots) are generated using two nested counters.  We use the |\robustify| command from \texttt{etoolbox} so that accented characters in the affiliation addresses do not blow up.  The list of accents may need to be expanded. 
%    \begin{macrocode}
\newcounter{AUTHOR}
\newcounter{author}[AUTHOR]
\stepcounter{AUTHOR}
\newcommand\ADS@affiliations{}
\robustify{\'}\robustify{\i}\robustify{\~}\robustify{\`}
\def\AuthorADS#1#2{%
  \begingroup
  \def\ORCIDiD##1{<ID system="ORCID">##1</ID>}
  \def\Email##1{<EMAIL>##1</EMAIL>}
  \def\Affil##1{##1}
  \stepcounter{author}
  \ifthenelse{\value{author} > 26}{\stepcounter{AUTHOR}\stepcounter{author}}{}
  \ifthenelse{\equal{\ADS@affiliations}{\@rmaaEmpty}}{%
    \xdef\ADS@affiliations{\Alph{AUTHOR}\Alph{author}(#2)}
  }{%
    \xdef\ADS@affiliations{\ADS@affiliations, \Alph{AUTHOR}\Alph{author}(#2)}
  }%
  \endgroup
}
%    \end{macrocode}
% \end{macro}
% \end{macro}
% %
% \begin{macro}{\ADS@dumpdata}
% \changes{4.0}{2019/11/13}{Write all author affiliations for ADS}
% \begin{macro}{\ADS@write}
% \begin{macro}{\ADS@openfile}
% |\ADS@dumpdata| is the macro that actually does the job of writing out all the
% data to an auxiliary file, using |\ADS@write| to write each
% field. The file is set up with |\ADS@openfile|, which currently is
% done in the TOC of a book.  The new multi-author affiliations are used if |\ADS@affiliations| has been set, otherwise it falls back to the single-author version |\ADS@affil|. 
%    \begin{macrocode}
\newcommand{\ADSfield}[2]{\relax}
\newcommand{\ADSvfield}[2]{\relax}
%%\renewcommand{\ADSfield}[12]{\%#1 #2}
\newcommand{\ADS@dumpdata}
{ 
  \ADS@write{\string\begin{ADSarticle}}
  \ADS@WRITE{\string\ADSfield{R}{\ADS@BibCode}}
  \ADS@write{\string\ADSfield{A}{\ADS@authors}}
  \ifthenelse{\equal{\ADS@affiliations}{\@rmaaEmpty}}{%
    {\ADS@sanitize\ADS@write{\string\ADSfield{F}{AA(\ADS@affil)}}}
  }{%
    {\ADS@sanitize\ADS@write{\string\ADSfield{F}{\ADS@affiliations}}}
  }
  \ADS@write{\string\ADSfield{J}{\ADS@journalname}}
  \ADS@write{\string\ADSfield{D}{\ADS@fulldate}}
  \ADS@write{\string\ADSfield{L}{\ADS@lastpage}}
  {\ADS@sanitize\ADS@write{\string\ADSfield{T}{\ADS@title}}}
  \ADS@write{\string\ADSfield{C}{\ADS@copyright}}
  \ADS@write{\string\ADSfield{d}{\ADS@doi}}
  \ADS@write{\string\ADSfield{U}{\ADS@url}}
  \ADS@write{\string\ADSfield{K}{\ADS@keywords}}
  {\ADS@sanitize\ADS@write{\string\ADSfield{B}{\ADS@abstract}}}
  \ADS@write{\string\ADSfield{Z}{\ADS@references}}
  \ADS@write{\string\end{ADSarticle}}
  }
%%\newcommand{\ADS@write}[1]{\if@book\addtocontents{ads}{#1}\fi}
%%\newcommand{\ADS@openfile}{\@starttoc{ads}}
\newcommand{\ADS@write}[1]{\set@display@protect\immediate\write\tf@ads{#1}}
\newcommand{\ADS@WRITE}[1]{\immediate\write\tf@ads{#1}}
\newcommand{\ADS@openfile}{%
  \newwrite\tf@ads%
  \immediate\openout\tf@ads\jobname.ads\relax
  \ADS@WRITE{\ADS@preamble}}
\newcommand{\ADS@closefile}{\ADS@WRITE{\ADS@postamble}}
\def\ADS@preamble{%
  %%  \string\documentclass{article}
  %%  \string\begin{document}
  }
\def\ADS@postamble{%
  %%  \string\end{document}
  \string%\string% Local Variables:^J 
  \string%\string% mode: LaTeX^J
  \string%\string% mode: Font-lock^J
  \string%\string% End:^J
  }
%    \end{macrocode}
% \end{macro}
% \end{macro}
% \end{macro}
% 
% \subsection{Automatic calculation of the page range}
% \label{sec:page-range}
% Another issue that has cost me a lot of grief. Still, it finally
% seems to be sorted now. All hail Jeff Goldberg. 
%
% Cunning trick to get the page numbers from |\pageref|.  Put a label
% on the first and last page of the document using |\AtBeginDocument|
% and |\AtEndDocument|. .  This way won't work for the book option.
% In that case, we do it by defining |\@endofarticlehook| and
% |\@begofarticlehook| as parallels to |\AtEndDocument| and
% |\AtBeginDocument|.  Note that we have to defer adding
% |\@begofarticlehook| to |\AtBeginDocument| to make sure it comes
% after the redefinitions of |\ref| and |\label|.
%    \begin{macrocode}
\if@book\AtEndDocument{\ADS@closefile}\else
 %\AtBeginDocument{\@begofarticlehook}
\AtEndDocument{\@endofarticlehook}
\fi
%    \end{macrocode}
% \begin{macro}{\@begofarticlehook}
% Since we are getting more and more things that need to be done at
% the beginning of every article, frontsection, etc we define a sister
% hook to 
% |\@endofarticlehook|. This hook does the following: 
% \begin{enumerate}
% \item Labels the first page
% \item Sets |\ADS@firstpage|
% \item runs and then clears |\indexauthor@deferlist|
% \item resets the keyword counter
% \end{enumerate}
% It is executed in the following circumstances: 
% \begin{enumerate}
% \item At |\begin{document}| for standalone files
% \item At |\begin{document}| for each included file
% \item At beginning of front sections 
% \item At beginning of TOC
% \item At beginning of Author Index
% \item At beginning of abstracts
% \end{enumerate}
% For some reason, this is not always within the right group, so we
% use |\xdef| just to be sure. 
%    \begin{macrocode}
\newcommand{\@begofarticlehook}{%
  \if@frontmatter\frontlabel{RMAAFirstPage}\else\label{RMAAFirstPage}\fi%
  \xdef\ADS@firstpage{\@arabic{\c@page}}%
  \if@book\else\makeindex\fi % needed so default index is defined in stand-alone articles
%%  \indexauthor@deferlist
%%  \xdef\indexauthor@deferlist{}
  \@inpreamblefalse
  \setcounter{kwdcounter}{0}
  }
%    \end{macrocode}
% \end{macro}
% \begin{macro}{\@endofarticlehook}
% This is a hook to be run at the end of each article
% currently, it just labels the last page and writes out all the ADS
% data to an auxiliary file.  
% It should not produce any printed output. 
% It is modelled on the |\AtEndDocument| command, to which it reduces
% in the case of a single article. 
%    \begin{macrocode}
\newcommand{\@endofarticlehook}{%
  \if@debug{\ClassWarning{rmaa}{***Doing Hook***}}\fi
  \if@debug{\ClassWarning{rmaa}{ADSabstract is: \ADS@abstract}}\fi
  \clearpage\rmaa@lastpagelabel\if@book\ADS@dumpdata\fi
  \ifx\NotesList\@rmaaEmpty\else\OutputNotes\fi
  \if@debug{\ClassWarning{rmaa}{Notes list is: \NotesList}}\fi
}
%    \end{macrocode}
% \begin{macro}{\rmaa@lastpagelabel}
% New version of last page label command copied from Jeff Goldberg's
% lastpage.sty package - we can't just use |\label| here in case the
% output routine for the last page has already been called. Instead,
% we first explicitly clear the page, step back the page counter,
% explicitly write to the .aux file, then restep the page counter
% again.
%
% But this will conflict with the hyperref package, which does labels differently.  Compensate for that by using the five-argument form.
% 
% Also set the |\ADS@LastPage| macro to the appropriate number. 
%    \begin{macrocode}
\newcommand{\rmaa@lastpagelabel}{%
  \addtocounter{page}{-1}%
  \immediate\write\@auxout{\string
    \newlabel{\articlename:RMAALastPage}{{}{\thepage}{}{}{}}}%
  \edef\ADS@lastpage{\thepage}%
  \addtocounter{page}{1}}
%    \end{macrocode}
% \end{macro}
% \end{macro}

% \subsection{Deal with explicit change in number of columns}
% \label{sec:columns}
% Sometimes, when formatting a book, we need to switch between
% one-column articles and two-column articles. For this to work, we
% need to wrap the standard \LaTeX{} commands to set the |\@rmtwocol|
% flag appropriately and also reset the package default values of all
% scalable lengths. .  
%    \begin{macrocode}
\let\latexonecolumn=\onecolumn
\let\latextwocolumn=\twocolumn
\renewcommand{\onecolumn}{\@rmtwocolfalse\@package@default@lengths\latexonecolumn}
\renewcommand{\twocolumn}{\@rmtwocoltrue\@package@default@lengths\latextwocolumn}
%    \end{macrocode}
% 
% \subsection{The actual formatting of the title}
% \label{sec:title-format}

% Note, we have to redefine the |\maketitle| command to use the
% numbered footnotes. This means we needn't use the atbeginend stuff. 
%    \begin{macrocode}
\renewcommand{\maketitle}{\par
  \begingroup
%    \end{macrocode}
% Don't switch to marked footnotes. 
% This is consistent with existing style but I'm not sure I like it
%    \begin{macrocode}
  %%  \renewcommand\thefootnote{\fnsymbol{footnote}}%
%%  \def\@makefnmark{\hbox to\z@{$\m@th^{\@thefnmark}$\hss}}%
  \def\@makefnmark{\textsuperscript{\@thefnmark}\hspace{-0.2em}}%
  \long\def\@makefntext##1{\parindent 1em\noindent
    \hbox to1.8em{\hss$\m@th^{\@thefnmark}$}##1}%
  \if@rmtwocol
  { \if@debug{\ClassWarning{rmaa}{TWO COLUMN}}\fi
    \twocolumn[\@maketitle]%
    }
  \else
  { \if@debug{\ClassWarning{rmaa}{ONE COLUMN}}\fi
    \newpage
    \global\@topnum\z@   % Prevents figures from going at top of page.
    \if@manuscript \@MSmaketitle \else \@maketitle \fi
    }
  \fi
  \if@manuscript\else\thispagestyle{rmaafirstpage}\@thanks\fi
  \ifx\NotesList\@rmaaEmpty\else{%
    \NoteMark{1--\arabic{notecounter}}\vspace*{-2\baselineskip}%
  }\fi%
  \endgroup
%    \end{macrocode}
% Don't reset footnote counter (not sure whether this is good thing
% or not).
%    \begin{macrocode}
  %%  \setcounter{footnote}{0}%
  \let\thanks\relax\let\affil\relax\let\altaffil\relax
  \let\maketitle\relax\let\@maketitle\relax
  \gdef\@thanks{}\gdef\@author{}\gdef\@title{}}
%    \end{macrocode}

% \begin{macro}{\@kwhwidth}
% \begin{macro}{\@kwtwidth}
% Redefine \cs{@maketitle} to include the abstract, resumen and
% keywords. Note that in the manuscripty style we use
% \cs{@MSmaketitle} instead.  
% \changes{3.16}{2003/04/07}{
%  Added extra vertical space between authors and accepted date.}
% \changes{3.26}{2011/08/15}{Add extra label to review articles}
%    \begin{macrocode}
\newlength{\@kwhwidth}\newlength{\@kwtwidth}
\newcommand\reviewname{Review}
\renewcommand{\@maketitle}{%
  \newpage
  \null
  \if@reviewarticle
  \smallskip
  {\large\bfseries\itshape\reviewname}%
  \bigskip
  \else
  \vskip\@before@title@skip%
  \fi
  \begin{center}%
    {\@title \par}%
    \vskip\@before@author@skip%
    \parbox{\rmaa@AuthorWidthFrac\textwidth}{\centering
      \lineskip\@between@author@skip%
      \@author
      \par}
    %% this should be replaced with `Submitted to...' for preprints
    \if@printdates
    \vskip\@before@dates@skip%
    \if@preprint
    \rmaa@ReceivedStyle{Submitted to RMxAA \@received}
    \else
    \rmaa@ReceivedStyle{Received \@received; accepted \@accepted}
    \fi
    \fi
    \if@manuscript
    \rmaa@ReceivedStyle{Submitted to RMxAA \@received}
    \fi
    \if@guide
    \rmaa@ReceivedStyle{Version \rmaa@version, \rmaa@date}
    \fi
    \vskip\@before@abstract@skip%
    \if@abstractpresent
    \parbox{\rmaa@AbstractWidthFrac\textwidth}{%
      \setlength{\parindent}{0.9cm}%
      \centerline{\rmaa@AbstractHeaderStyle{Abstract}}\par%
      \vskip\@after@abstracthead@skip%
      \if@newstyle\noindent\fi\@abstract
      \vskip\@between@abstract@skip%
      \centerline{\rmaa@AbstractHeaderStyle{Resumen}}\par%
      \vskip\@after@abstracthead@skip%
      \if@newstyle\noindent\fi\@resumen
      \if@setresumen\else\@abstract\fi % append english abstract if
                                       % spanish one not set
      \if@guide\else            %do not use keywords for guide
      \vskip\@before@kwd@skip%
      % Format the Key Words with hanging indentation:
      % They are put in a \parbox with width equal to the line width
      % minus the width of the header. 
      { % Widths of keyword
        \settowidth{\@kwhwidth}{\@keywordheader~}   % header and
        \setlength{\@kwtwidth}{\linewidth}          % keyword text
        \addtolength{\@kwtwidth}{-\@kwhwidth}       %
        \noindent\@keywordheader~\parbox[t]{\@kwtwidth}{%
          \rmaa@KeywordStyle{\boldmath\uppercase{\@keywords}}}}%
      \fi
      }
      \fi
  \end{center}%
  \vspace*{\afterkwdskip}
  \par
  }
\newlength{\afterkwdskip}\setlength{\afterkwdskip}{\z@}
%    \end{macrocode}
% \end{macro}
% \end{macro}
%
%
% \begin{macro}{\reviewarticle}
% \begin{macro}{\reviewtoc}
% \begin{macro}{\reviewlof}
%   User command to specify a review article. Must be used in
%   preamble. Main effect is to place the text ``Review'' before the
%   article title.  Also add command for setting a table of contents
%   and a list of figures at the beginning of the article.
%    \begin{macrocode}
\newcommand{\reviewarticle}{\@reviewarticletrue}
\newcommand\reviewtoc{%
  \section*{CONTENTS}%
  \@starttoc{rtc}%
}
\newcommand\reviewlof{%
  \section*{LIST OF FIGURES}%
  \@starttoc{lof}%
}
%    \end{macrocode}
% \end{macro}
% \end{macro}
% \end{macro}
% \begin{macro}{\l@section}
%    \begin{macrocode}
\renewcommand*\l@section[2]{%
  \ifnum \c@tocdepth >\z@
    \addpenalty\@secpenalty
    \addvspace{0.8em \@plus\p@}%
    \setlength\@tempdima{1.5em}%
    \begingroup
      \parindent \z@ \rightskip \@pnumwidth
      \parfillskip -\@pnumwidth
      \leavevmode %\bfseries
      \advance\leftskip\@tempdima
      \hskip -\leftskip
      #1\nobreak\hfil \nobreak\hb@xt@\@pnumwidth{\hss #2}\par
    \endgroup
  \fi}
%    \end{macrocode}
% \end{macro}
%
% \begin{macro}{\@MSmaketitle}
% A special version of |\@maketitle| for the \texttt{manuscript}
% option. This can deal with a multi-page abstract/resumen. Not very
% well, mind. 
%    \begin{macrocode}
\newcommand{\@MSmaketitle}{%
  \if@debug\ClassInfo{rmaa}{Using special MS maketitle command}\fi
  \newpage
  {
    \thispagestyle{rmaafirstpage}
%     \setlength{\parindent}{0pt}
    \null\vskip 2ex
    \begin{center}
      \@title\par
      \@author\par
      \rmaa@ReceivedStyle{%
        \ifx\@received\@rmaaEmpty
          Draft version: \@date
        \else
          Received: \@received
          \ifx\@accepted\@rmaaEmpty
          \else
            ; Accepted: \@accepted
          \fi
        \fi}\par
    \end{center}
    \@thanks
    \centerline{\rmaa@AbstractHeaderStyle{Resumen}}
    \noindent\@resumen\par\vskip 2ex
    \centerline{\rmaa@AbstractHeaderStyle{Abstract}}
    \noindent\@abstract\par\vskip 2ex
    \noindent\@keywordheader~\rmaa@KeywordStyle{\@keywords}\par
%     \@keywordheader~\rmaa@KeywordStyle{\uppercase{\@keywords}}\par
  }%
  }
%    \end{macrocode}
% \end{macro}

% \subsection{The sectioning commands}
% \label{sec:section-commands}

% \begin{macro}{\section}
% \begin{macro}{\subsection}
% \begin{macro}{\subsubsection}
% \begin{macro}{\paragraph}
% \begin{macro}{\subparagraph}
% Redefine all the standard commands to use these lengths. 
%    \begin{macrocode}
\renewcommand\section{\@startsection {section}{1}{\z@}%
                                   {\abovesectionskip}%
                                   {\belowsectionskip}%
                                   %% how do we get upper case here
                                   {\centering\reset@font\normalsize}}
\renewcommand\subsection{\@startsection{subsection}{2}{\z@}%
                                     {\abovesectionskip}%
                                     {\belowsectionskip}%
                                     {\centering\reset@font\normalsize\itshape}}
\renewcommand\subsubsection{\@startsection{subsubsection}{3}{\z@}%
                                     {\abovesectionskip}%
                                     {\belowsectionskip}%
                                     {\centering\reset@font\normalsize\itshape}}
\renewcommand\paragraph{\@startsection{paragraph}{4}{\z@}%
                                    {\aboveparaskip}%
                                    {\belowparaskip}%
                                    {\reset@font\normalsize\bfseries}}
\renewcommand\subparagraph{\@startsection{subparagraph}{5}{\parindent}%
                                      {\aboveparaskip}%
                                      {\belowparaskip}%
                                      {\reset@font\normalsize\bfseries}}
%    \end{macrocode}
% \end{macro}
% \end{macro}
% \end{macro}
% \end{macro}
% \end{macro}

% \begin{macro}{secnumdepth}
% Make sure that we number the subsubsections---this is done by
% default in the |article| class but not the |book| class, which is
% why we must set it explicitly here.
%    \begin{macrocode}
\setcounter{secnumdepth}{3}
%    \end{macrocode}
% \end{macro}
% \changes{2.2beta10}{2001/11/07}{Fixed \cs{secnumdepth} for \cs{book}
% style.}

% \begin{macro}{\thesubsection}
% \begin{macro}{\thesubsubsection}
% For subsection and subsubsection, we want the numbers to be roman
% but the title to be italic. This seems the easiest way to go about it.
%    \begin{macrocode}
\renewcommand{\thesubsection}{{\upshape\thesection.\@arabic\c@subsection}}
\renewcommand{\thesubsubsection}{\thesubsection{\upshape.\@arabic\c@subsubsection}}
%    \end{macrocode}
% \end{macro}
% \end{macro}

% \begin{macro}{\@seccntformat}
% Add ``.'' to the number here so we can remove it from |\thesection|,
% etc. Reduce gap between number and title WJH 27 Jan 1995
%    \begin{macrocode}
\def\@seccntformat#1{\csname the#1\endcsname .\hskip 0.4em}
%    \end{macrocode}
% \end{macro}

% \begin{macro}{\@sect}
% I didn't really want to mess with the |\@sect| command, but it was
% the only way I could manage to force upper case. Other than the
% commented parts, it is identical to the standard version. 
%    \begin{macrocode}
\def\@sect#1#2#3#4#5#6[#7]#8{
  \ifnum #2>\c@secnumdepth
     \let\@svsec\@empty
  \else
     \refstepcounter{#1}%
     \let\@@protect\protect
     \def\protect{\noexpand\protect\noexpand}%
     \edef\@svsec{\@seccntformat{#1}}%
     \let\protect\@@protect
  \fi
%    \end{macrocode}
% This is the bit I changed (27 Jan 1995). Use upper case for
% |\section|, use mixed case for the rest. 
%    \begin{macrocode}
  \ifnum #2=1                   
     \def\sec@upcase##1{\uppercase{##1}}%
     \def\sec@shape{\upshape}%
  \else                         
     \def\sec@upcase##1{\relax{##1}}%
     \def\sec@shape{\itshape}%
  \fi
  \@tempskipa #5\relax
  \ifdim \@tempskipa>\z@
     \begingroup #6\relax
          \@hangfrom{\hskip #3\relax\@svsec}%
                    {\interlinepenalty \@M \sec@upcase{#8}\par}%
     \endgroup
     \csname #1mark\endcsname{#7}%
%    \end{macrocode}
% Only write TOC entries if doing the Guide style. 
%    \begin{macrocode}
     \if@guide 
        \addcontentsline{toc}{#1}{%
          \ifnum #2>\c@secnumdepth \else
             \protect\numberline{\csname the#1\endcsname}%
          \fi
          #7}%
     \fi
%    \end{macrocode}
% Also allow per-article TOC for review papers
%    \begin{macrocode}
     \if@reviewarticle
        \ifnum #2>2 \else
           \addcontentsline{rtc}{#1}{%
             \protect\numberline{\csname the#1\endcsname}%
             \sec@upcase{\sec@shape #7}}%
        \fi
     \fi
  \else
     \def\@svsechd{#6\hskip #3\relax  %% \relax added 2 May 90
       \@svsec #8\csname #1mark\endcsname{#7}%
       \if@guide
          \addcontentsline{toc}{#1}{%
            \ifnum #2>\c@secnumdepth \else
               \protect\numberline{\csname the#1\endcsname}%
            \fi
            #7}%
       \fi
       }%
  \fi
  \@xsect{#5}}
%    \end{macrocode}
% \end{macro}


% \subsection{Appendices}
% \label{sec:appendices}
% Make sure that appendices in individual articles work properly. We
% copy the definition from the article class (the book one doesn't do
% what we want).  
% \begin{environment}{appendices}
%   This uses an environment |appendices| inside of which each
%   |\section| will become a separate appendix.  Note that the
%   appendices must come before the references (otherwise the full
%   addresses mechanism will not work properly). We make the equation
%   numbering be subordinate to the appendix (equations A12, B6,
%   etc). Finally, we have a medskip followed by APPENDICES to
%   introduce them. \changes{3.17}{2002/04/08}{ Added appendices environment
%   that does the right thing.}
%    \begin{macrocode}
\newenvironment{appendices}{\par
  \begingroup
  \setcounter{section}{0}%
  \setcounter{subsection}{0}%
  \renewcommand\thesection{\@Alph\c@section}
  \renewcommand\theequation{\thesection\arabic{equation}}
  \medskip
  \section*{APPENDICES}}
{\endgroup}
%    \end{macrocode}
% \end{environment}
% \begin{environment}{appendix}
%   If there is only one appendix, then use |appendix| environment
%   instead.  \changes{4.0}{2019/11/12}{ Added
%   appendix environment for case where there is only one section.}
%    \begin{macrocode}
\renewenvironment{appendix}{\par
  \begingroup
  \setcounter{section}{0}%
  \setcounter{subsection}{0}%
  \renewcommand\thesection{\@Alph\c@section}
  \renewcommand\theequation{\thesection\arabic{equation}}
  \medskip
  \section*{APPENDIX}}
{\endgroup}
%    \end{macrocode}
% \end{environment}

% \subsection{Floating thingies}
% \label{sec:floats}   

%% 
%    \begin{macrocode}
\renewcommand{\thetable}{\@arabic\c@table}
\renewcommand{\thefigure}{\@arabic\c@figure}
\renewcommand{\figurename}{Fig.}
\renewcommand{\tablename}{TABLE}
%\renewcommand{\fnum@table}{\MakeUpperCase{\tablename}~\thetable}
%    \end{macrocode}
% For some reason, I have to write `TABLE' explicitly here
% If I use |\tablename|, it gets converted to mixed case---God knows
% why. 
% |\figurename| is also being ignored. 
%    \begin{macrocode}
\renewcommand{\fnum@table}{TABLE~\thetable}
\renewcommand{\fnum@figure}{Fig.~\thefigure}
%    \end{macrocode}

% Do captions in 9 pt face in the main journal style and in new-style
% anything. 
%    \begin{macrocode}
\long\def\@makefigcap#1#2{%
  \if@proceedings\else\small\fi
  \if@newstyle\small\fi 
  \vskip\abovecaptionskip
  \sbox\@tempboxa{#1. #2}%
  \ifdim \wd\@tempboxa >\hsize
    #1. #2 \par
  \else
    \global \@minipagefalse
    \hb@xt@\hsize{\hfil\box\@tempboxa\hfil}%
  \fi
  \vskip\belowcaptionskip}
%    \end{macrocode}

% Make sure table caption is upper case. 
%    \begin{macrocode}
\long\def\@maketabcap#1#2{%
  \hbox to\hsize{\hfil #1 \hfil}
  \vskip 6pt 
  \centering\MakeUpperCase{#2}\par
  \vskip 6pt}
%    \end{macrocode}

%    \begin{macrocode}
\newlength{\@rottabcapwidth}
\long\def\@makerottabcap#1#2{%
  \setbox\@tempboxa\hbox{\MakeUpperCase{#2}}%
  \ifdim\wd\@tempboxa >.8\vsize
     \setlength\@rottabcapwidth{.8\vsize}
  \else
     \setlength\@rottabcapwidth{\wd\@tempboxa}
  \fi
  \rotatebox{90}{%
    \begin{minipage}{\vsize}\centering
      \begin{minipage}{\@rottabcapwidth}\centering
        #1\\ \vskip 6pt \MakeUpperCase{#2}
      \end{minipage}%
    \end{minipage}
    }%
  \hskip 6pt
  }

\newenvironment{Sideways}
    {\begin{sideways}\begin{minipage}{\vsize}\centering}
    {\end{minipage}\end{sideways}}

\newenvironment{Sidewaystable}
{%
  \let\oldcaption=\caption%
  \def\caption{
    \def\@captype{table}\renewcommand{\@makecaption}{\@maketabcap}
    \oldcaption}
  \begin{sidewaystable}\centering}
  {\end{sidewaystable}\let\caption=\oldcaption}
\newenvironment{Sidewaystable*}
{%
  \let\oldcaption=\caption%
  \def\caption{
    \def\@captype{table}\renewcommand{\@makecaption}{\@maketabcap}
    \oldcaption}
  \begin{sidewaystable*}\centering}
  {\end{sidewaystable*}\let\caption=\oldcaption}
%    \end{macrocode}

% Not sure why we need to stick these lines in. 
%    \begin{macrocode}
\renewenvironment{figure}
                 {\renewcommand{\@makecaption}{\@makefigcap}\@float{figure}}
                 {\end@float}
\renewenvironment{figure*}
                 {\renewcommand{\@makecaption}{\@makefigcap}\@dblfloat{figure}}
                 {\end@dblfloat}
\renewenvironment{table}
                 {\renewcommand{\@makecaption}{\@maketabcap}\@float{table}}
                 {\end@float}
\renewenvironment{table*}
                 {\renewcommand{\@makecaption}{\@maketabcap}\@dblfloat{table}}
                 {\end@dblfloat}
%    \end{macrocode}

% Redefine |\@caption| so that we don't bother with the list of
% figures. This saves us having to worry about fragile commands in the
% captions. 
%    \begin{macrocode}
\long\def\@caption#1[#2]#3{%
  \par
%%  \addcontentsline{\csname ext@#1\endcsname}{#1}%
%%    {\protect\numberline{\csname the#1\endcsname}{\ignorespaces #2}}%
  \begingroup
    \@parboxrestore
    \if@minipage
      \@setminipage
    \fi
    \normalsize
    \@makecaption{\csname fnum@#1\endcsname}{\ignorespaces #3}\par
  \endgroup}
%    \end{macrocode}

% Make tables not quite so cramped. Allow |\hline\hline| to give a
% poor man's thick rule. 
%    \begin{macrocode}
\renewcommand{\arraystretch}{1.15}
\setlength{\doublerulesep}{0.4pt}
%    \end{macrocode}

% \begin{macro}{\tablecols}
% \begin{macro}{\@tablecols}
% \begin{macro}{\tabcaptioncontinued}
% |\tablecols| is a user command for setting the number of columns. It
% is used by commands that try and span the entire table (I think
% there must be a better way of acheiving
% this). |\tabcaptioncontinued| is 
% a user command for setting the caption of a continued table on the
% next page. 
%    \begin{macrocode}
\newcommand{\tablecols}[1]{\renewcommand{\@tablecols}{#1}}
\newcommand{\@tablecols}{1}
\newcommand{\tabcaptioncontinued}{%
  \multicolumn{\@tablecols}{c}{\addtocounter{table}{-1}TABLE~\thetable\ (CONTINUED)} \\ \\
  }
%    \end{macrocode}
% \end{macro}
% \end{macro}
% \end{macro}

% \subsubsection{Footnotes in tables}
% \label{sec:table-footnotes}
% \begin{macro}{\tablenotemark}
% \begin{macro}{\tablenote}
% \begin{macro}{\tablenoteskip}
% Old user commands |\tablenotemark{MARK}|,
% |\tablenote{MARK}{TEXT}|---now deprecated. 
%    \begin{macrocode}
\newcommand\tablenotemark[1]{$^\mathrm{#1}$} 
\newcommand\tablenote[2]{\if@proceedings\else\small\fi ${}^\mathrm{#1}$ #2} 
\newcommand\tablenoteskip{\\[0.5\baselineskip]} 
%    \end{macrocode}
% \end{macro}
% \end{macro}
% \end{macro}

% \begin{macro}{\tabnotemark}
% \begin{macro}{\tabnotetext}
% \begin{macro}{\tabnotetextTWO}
% Alternative way of doing footnotes in tables. 
%    \begin{macrocode}
\newcommand{\tabnotemark}[1]{\textsuperscript{\lowercase{#1}}}
%\newcommand{\tabnotetextTWO}[4]{\tabnotetext{#1}{#2}\tabnotetext{#3}{#4}}
\newcommand{\tabnotetextTWO}[4]{%
  \multicolumn{\@tablecols}{l}{%
    \parbox{0.5\tabnotewidth}{\small\textsuperscript{#1}#2}%
    \parbox{0.5\tabnotewidth}{\small\textsuperscript{#3}#4}%
    }%
  \\}
\newcommand{\tabnotetext}[2]{%
  \multicolumn{\@tablecols}{@{}p{\tabnotewidth}@{}}{\small\textsuperscript{#1}#2}\\}
\newlength{\tabnotewidth}
%    \end{macrocode}
% \end{macro}
% \end{macro}
% \end{macro}


% \begin{macro}{changemargin}
%    \begin{macrocode}
% Command to allow margins to be changed on the fly. Useful to allow
% tables and Figures to protrude into the margins. 
%
% \changes{3.24}{2007/09/14}{Added \texttt{changemargin} environment
% (copied from TeX FAQ).}
\newenvironment{changemargin}[2]{% 
  % From TeX FAQ http://www.tex.ac.uk/cgi-bin/texfaq2html?label=chngmargonfly
  \begin{list}{}{%
    \setlength{\topsep}{0pt}%
    \setlength{\leftmargin}{#1}%
    \setlength{\rightmargin}{#2}%
    \setlength{\listparindent}{\parindent}%
    \setlength{\itemindent}{\parindent}%
    \setlength{\parsep}{\parskip}%
  }%
  \item[]}{\end{list}}
%    \end{macrocode}
% \end{macro}

% \subsection{The bibliography and references}
% \label{sec:biblio}
% \begin{macro}{\beforebiblioskip}
% \begin{macro}{\afterbiblioskip}
% I'm not entirely sure why I introduced these two commands, since
% directly putting |\vspace| before and after
% |\begin{thebibliography}| seems to work well enough. 
%    \begin{macrocode}
\newlength{\beforebiblioskip}
\setlength{\beforebiblioskip}{0pt}
\newlength{\afterbiblioskip}
\setlength{\afterbiblioskip}{0pt}
%    \end{macrocode}
% \end{macro}
% \end{macro}
% \begin{macro}{thebibliography}
% The environment itself is done as starred section with extra
% vertical space after heading. This space is currently zero. 
% In the book style we do the cunning trick with rereading the file in
% order to get the references into the |\ADS@references| macro.
% 
% \changes{3.05}{2003/01/02}{Removed the \cs{beforebiblioskip} and
% \cs{afterbiblioskip} commands, since they cause unwanted vertical
% space in some instances.}  
%    \begin{macrocode}
\renewenvironment{thebibliography}{%
  \if@book\RereadBiblio\fi
%%  \vspace*{\beforebiblioskip}
  \section*{REFERENCES} \frenchspacing
%%  \vspace*{\afterbiblioskip}
  \small
%    \end{macrocode} 
% The references themselves are set as a list in order 
% to get the hanging indentation. They are done in a smaller font and
% with french spacing so we needn't worry about the distinction
% between sentence-ending periods and abbreviations. 
%    \begin{macrocode} 
  \list{\null}{\leftmargin 1.5em\labelwidth\z@\labelsep\z@\itemindent -1.5em
    \itemsep\z@\parsep\z@\usecounter{enumi}}
  \sloppy\clubpenalty4000\widowpenalty4000
  \sfcode`\.=1000\relax\@biblisttrue}
{\endlist\@biblistfalse\normalsize\@outputaddressesifnecessary}
%    \end{macrocode}
% \end{macro}
% \subsection{Uniquification of cross-references}
% \label{sec:uniquification}
% 
% \begin{macro}{\articlename}
% In order that the |\label{}| command in individual articles be
% guaranteed to give a unique anchor, we append the article name (held
% in the |\articlename| command) to each key. This requires
% redefinition of the following commands: |\cite|, |\nocite|, |\ref|,
% |\pageref|, |\label|. 
% \begin{macro}{\oldcite}
% \begin{macro}{\oldnocite}
% \begin{macro}{\oldref}
% \begin{macro}{\oldpageref}
% \begin{macro}{\oldlabel}
% The original versions of these commands are
% saved as |\oldcite|, |\oldnocite|, etc. 
%
% The uniquifying of cross-references is now always done, so 
% we don't get multiply-defined labels in the single article case
%    \begin{macrocode}
\newcommand{\articlename}{anon}
\AtBeginDocument{
  \let\oldcite=\cite 
  \let\oldnocite=\nocite 
  \let\oldref=\ref 
  \let\oldpageref=\pageref
  \let\oldlabel=\label
}
%    \end{macrocode}
% \end{macro}
% \end{macro}
% \end{macro}
% \end{macro}
% \end{macro}
% \end{macro}
%
% Uniquify all cross-references by prepending the article name.  This
% now has to go into |\AtBeginDocument| so it is working on the new
% versions that hyperref.sty has set up.  Then we need to put the
% |\@begofarticlehook| afterwards, so that the first page label is using
% the new definitions.  If amsmath is loaded, then we also need to patch
% the special label command that they use. \changes{4.0}{2019/11/18}{Fix labels for amsmath display equations, in case that package has been loaded.}
%    \begin{macrocode}
\AtBeginDocument{
  \renewcommand{\ref}[1]{\oldref{\articlename:#1}%
    \if@debug{\ClassInfo{rmaa}{Reference to \articlename:#1}}\fi}
  \renewcommand{\pageref}[1]{\oldpageref{\articlename:#1}%
    \if@debug{\ClassInfo{rmaa}{Page reference to \articlename:#1}}\fi}
  \renewcommand{\label}[1]{\oldlabel{\articlename:#1}%
    \if@debug{\ClassInfo{rmaa}{Label for \articlename:#1}}\fi}%
  \ifdef{\label@in@display}{%
    \let\oldlabel@in@display=\label@in@display
    \renewcommand{\label@in@display}[1]{\oldlabel@in@display{\articlename:#1}%
      \if@debug{\ClassInfo{rmaa}{AMSTeX Label for \articlename:#1}}\fi}%
  }{}
  \@begofarticlehook
}%
%    \end{macrocode}
% Citations are more complicated since they can contain a list of
% keys. These are really only here for use with the rmaacite package,
% which is utterly deprecated (02 Jan 2003) now that I have the natbib
% code working. The command \cs{cite} should \emph{not} be used with
% natbib: use \cs{citep}, \cs{citet}, etc.\@ instead. 
%    \begin{macrocode}
\renewcommand{\nocite}[1]{\@bsphack
  \@for\@citeb:=#1\do{%
    \edef\@citeb{\expandafter\@firstofone\@citeb:\articlename}%
    \if@filesw\immediate\write\@auxout{\string\citation{\@citeb}}\fi
    \@ifundefined{b@\@citeb:\articlename}{\G@refundefinedtrue
      \@latex@warning{Citation `\@citeb' undefined}}{}}%
  \@esphack}
\renewcommand{\cite}[1]{\oldcite{#1:\articlename}}
%   \end{macrocode}


% Horrible fiddling with the bibliography citation mechanism. 
% This has to be globally available in case rmaacite puts it in the
% \texttt{.aux} file. 
%    \begin{macrocode}
\def\lbibcite#1#2#3{\global\@namedef{b@#1}{#2}
                    \global\@namedef{by@#1}{#3}}
%    \end{macrocode}

% \begin{macro}{\oldbibitem}
% \begin{macro}{\@rmaaBibArg}
% \begin{macro}{\@rmaaSpace}
% \begin{macro}{rmaabib}
% Of course, |\bibitem| needs to get the uniquification treatment as
% well. While we are at it, we redefine it to use consecutive integers
% for the key if the user leaves it blank (using the counter
% |rmaabib|). This allows one to put 
% |\bibitem{}| for all your references without getting annoying
% messages about multiply-defined labels.
% 
% We also take the opportunity to insert the optional argument to
% \cs{bibitem} if it is missing in the form Author(DATE) so that natbib
% doesn't try to do numerical citations. Natbib already provides a
% hook \cs{@extra@b@citeb} to help uniquify the citation key, so we
% just plug into that. However, it puts it \emph{after} the original
% key, not before as I had had it. Easiest solution is to change the
% way that I do it. 02 Jan 2003. Note that the \cs{@extra@b@citeb}
% does not get written to the .aux  file in the \cs{citation} command
% but it does work. 
% \changes{3.14}{06 Feb 2003}{Reinstalled old definition of \cs{bibitem}
%  for the case where the nonatbib option is requested, so that the
%  rmaacite package will work. This means that you can't use rmaacite
%  and natbib in the same volume.}
%    \begin{macrocode}
\if@natbib\def\@extra@b@citeb{:\articlename}\fi
%%\else
\let\oldbibitem\bibitem
\newcommand{\@rmaaBibArg}{}
\newcommand{\@rmaaSpace}{ }
\newcounter{rmaabib}
\if@natbib
\renewcommand{\bibitem}[2][ABC(2003)]{%
  \renewcommand{\@rmaaBibArg}{#2}
  \ifx\@rmaaBibArg\@rmaaEmpty
     \addtocounter{rmaabib}{1}
     \oldbibitem[#1]{\thermaabib:\articlename}
  \else\ifx\@rmaaBibArg\@rmaaSpace
     \addtocounter{rmaabib}{1}
     \oldbibitem[#1]{\thermaabib:\articlename}
  \else 
     \oldbibitem[#1]{#2:\articlename}
  \fi\fi
  \renewcommand{\@rmaaBibArg}{}
  }
\else
\renewcommand{\bibitem}[1]{%
  \renewcommand{\@rmaaBibArg}{#1}
  \ifx\@rmaaBibArg\@rmaaEmpty
     \addtocounter{rmaabib}{1}
     \oldbibitem{\articlename:\thermaabib}
  \else\ifx\@rmaaBibArg\@rmaaSpace
     \addtocounter{rmaabib}{1}
     \oldbibitem{\articlename:\thermaabib}
  \else
     \oldbibitem{\articlename:#1}
  \fi\fi
  \renewcommand{\@rmaaBibArg}{}
  }
\fi
%%\fi
%    \end{macrocode}
% \end{macro}
% \end{macro}
% \end{macro}
% \end{macro}

% \begin{macro}{oldpcite}
% \begin{macro}{oldscite}
% Then, we need to repeat the above scheme, but for the |rmaacite|
% definition. 
%    \begin{macrocode}
\newcommand{\@redefinebibitem}{%
  \let\oldbibitem\bibitem
  \let\oldcite\cite
  \let\oldpcite\pcite
  \let\oldscite\scite
  \def\bibitem[##1<##2>]##3{%
    \renewcommand{\@rmaaBibArg}{##3}
    \ifx\@rmaaBibArg\@rmaaEmpty
       \addtocounter{rmaabib}{1}
       \oldbibitem[##1<##2>]{\thermaabib:\articlename}
    \else\ifx\@rmaaBibArg\@rmaaSpace
       \addtocounter{rmaabib}{1}
       \oldbibitem[##1<##2>]{\thermaabib:\articlename}
    \else 
       \oldbibitem[##1<##2>]{##3:\articlename}
    % trust the author here
%%       \oldbibitem[##1<##2>]{##3}
    \fi\fi
    }
  }
%    \end{macrocode}
% \end{macro}
% \end{macro}


% \subsection{HTML Table of Contents}
% \label{sec:hyper-TOC}
% \begin{macro}{\HyperTOCEntry}
% \begin{macro}{\HyperTOCpart}
% \begin{macro}{\HyperTOCmeeting}
% \begin{macro}{\HyperTOCeditors}
% \begin{macro}{\HyperTOCpreamble}
% These commands are used to write information to the \texttt{.htoc}
% file, which is used in generating a Web Page of the Table of
% Contents. Since they only ever appear preceded by |\protect| in the
% argument to an |\addtocontents| command,  they should never need to
% be expanded by the macros. However, we define them here just in case
% \dots
%    \begin{macrocode}
\newcommand{\HyperTOCEntry}[3]{\relax}
\newcommand{\HyperTOCpart}[1]{\relax}
\newcommand{\HyperTOCmeeting}[3]{\relax}
\newcommand{\HyperTOCeditors}[1]{\relax}
\newcommand{\HyperTOCpreamble}{\relax}
%    \end{macrocode}
% \end{macro}
% \end{macro}
% \end{macro}
% \end{macro}
% \end{macro}
%

% \subsection{Formatting the poster abstracts} 
% These can now be used without using the book option.
%    \begin{macrocode}

\newcommand{\@makeminititle}{%
  \newpage
  \null 
  \vskip 3em%
  \begin{center}%
    {\@title \par}%
  \end{center}
  \vskip 3em%
  }

\newif\if@abstractsads
\@abstractsadsfalse
\newcommand\abstractsADS{\@abstractsadstrue}

\if@book\RequirePackage{multicol}\fi
\newenvironment{abstracts}[1][Abstracts of Contributed Papers]{%
  %% First, make sure we start on an odd page (06 Jun 2000)
  \if@book\clearpage\thispagestyle{empty}\cleardoublepage\fi
  %%
  \renewcommand{\articlename}{ABSTRACTS}
  \title{#1}\shortauthor{Abstracts}\shorttitle{Abstracts}
  \listofauthors{}
  \if@book
  \addcontentsline{toc}{chapter}{\protect\@fiddleauthors{\@TOCtitle}{\rmaaTOCauthors}}%
  \addtocontents{htc}{\protect\HyperTOCEntry{\@rawtitle}{\@rawauthors}{\articlename}}
  \fi
  \pagestyle{rmaa} % this now seems necessary (06 Jun 2000)
  \thispagestyle{rmaafirstpage}
%%  \twocolumn[\@makeminititle]
  \renewcommand{\abstract}[1]{%
    \renewcommand\ADS@abstract{##1}%
    \renewcommand{\@abstract}{##1}\@writeabstract%
  }
  \renewcommand{\title}[2][]{%
    \gdef\ADS@authors{}%
    \gdef\ADS@affil{}%
    \gdef\ADS@references{}%
    \gdef\ADS@keywords{}%
    \renewcommand{\@rmaaArg}{##1}
    \ifx\@rmaaArg\@rmaaEmpty
       \renewcommand{\@TOCaftertitle}{}
    \else
       \renewcommand{\@TOCaftertitle}{ (\@rmaaArg)}
    \fi
    \renewcommand{\@title}{\rmaa@TOCTitleStyle{##2}}%
    \renewcommand{\@TOCtitle}{\rmaa@TOCTitleStyle{##2}\@TOCaftertitle}%
    \renewcommand{\rmaaIndexTitle}{\rmaa@IndexTitleStyle{##2\@TOCaftertitle.}}%
    \renewcommand\ADS@title{##2}%
    \xdef\ADS@firstpage{\@arabic{\c@page}}%
    }
  \renewcommand{\author}[1]{\let\@thanks\@empty\renewcommand\@author{##1} }
  \renewcommand{\footnotetext}[2][XX]{\par\noindent\textsuperscript{##1}~##2}
  \renewcommand{\footnoterule}{\kern-3\p@ \hrule \@width 3em \kern 2.6\p@}
  \if@book\@begofarticlehook\fi
  \if@book\begin{multicols}{2}[\@makeminititle]\else\@makeminititle\fi
  \begingroup
  \setcounter{footnote}{0}
  }% and at the end of the environment....
  {\endgroup\if@book\end{multicols}\listofabstracts\label{RMAALastPage}\clearpage\fi}
%% auxiliary commands used:
%% \@writeabstract does all the actual formatting 
\newcommand{\@writeabstract}{%
  \if@book%
  \addcontentsline{loa}{abstract}{\protect\@fiddleauthors{\@TOCtitle}{\rmaaTOCauthors}}%
  \fi%
  \begin{center}\@title\par\@author\end{center}
  \@abstract
  \par\noindent\footnoterule\nobreak\footnotesize\@thanks
  \par\noindent\hrulefill\goodbreak
  \vskip 1cm \@plus0.5cm
  %% write ADS entry
  \if@book\if@abstractsads\ADS@dumpdata\fi\fi%
  \endgroup % this pair is to make each abstract a self-contained group
  \begingroup
  \setcounter{footnote}{0}
  }

\if@book
%% List of abstracts (stuff written to .loa file)
\newcommand{\listofabstracts}{%
  \begin{multicols}{2}[%
    \vskip 1.5cm \vfill % leave at least 1.5 cm and go to bottom of page if possible
    \centerline{\MakeUpperCase{List of Abstracts}}
    \vskip 1.5cm \@plus0.5cm
    ][6cm] % this makes it start a new page if there is not at least 6cm of space
    \renewcommand{\@fiddleauthors}[2]{##1\hskip1em plus2em minus0.4em\penalty-100  ##2}
    \parindent\z@
%%    \parskip.5\p@ \@plus .7\p@\relax 
    \parskip3\p@ \@plus 1\p@ \@minus 1\p@\relax % allow quite a lot of
                                % leeway here to avoid breaks inside
                                % entries 
    \renewcommand{\\}{\relax} % temporarily turn off explicit line breaks just in case
    \renewcommand{\altaffilmark}[1]{\relax} % ditto with this
    \@starttoc{loa}
  \end{multicols}
  }
\newcommand*{\l@abstract}[2]{%
  \par\penalty-1000
  #1\hskip 1em\penalty0\null\nobreak\hfill#2}
\fi

%    \end{macrocode}
 
% \subsection{Code for typesetting an entire volume of conference
% proceedings} \label{sec:whole-volume}
%    \begin{macrocode}

\if@book
%% environment for including the individual articles
\newcommand{\@fiddleauthors}[2]{\relax}
\renewcommand{\theequation}{\@arabic\c@equation}
\newenvironment{articles}%
{ \begingroup % make it a group so that author definitions do not
              % spill over
  \gdef\ADS@authors{}
  \pagestyle{rmaa}
  \renewenvironment{document}{%
    \@begofarticlehook
    \addtocounter{chapter}{1}%
    \addcontentsline{toc}{chapter}{%
      \protect\@fiddleauthors{\@TOCtitle}{\rmaaTOCauthors}}%
    \addtocontents{htc}{\protect\HyperTOCEntry{\@rawtitle}{\@rawauthors}{\articlename}}
    }{\relax}
  \renewcommand{\documentclass}[2][foo]{\relax} 
  \renewcommand{\usepackage}[2][foo]{%
    \def\@packagename{##2}\def\@rmaacitename{rmaacite}%
    \ifx\@packagename\@rmaacitename
    \@rmcitetrue\ClassInfo{rmaa}{This paper uses rmaacite}
    \makeatletter%%
%% This is file `rmaacite.sty',
%% generated with the docstrip utility.
%%
%% The original source files were:
%%
%% rmaa.dtx  (with options: `rmaacite')
%% 
%% PLEASE DO NOT EDIT THIS FILE! It is generated automatically from
%% the file rmaa.dtx using the accompanying Makefile and rmaa.ins
%% Any hand-edited changes will be lost
%% 
%% A LaTeX document class for the RMxAA
%% 
%% Author: William Henney (w.henney@crya.unam.mx)
%% 
%% Latest version available from: http://www.crya.unam.mx/rmaa
%% 
%% Copyright (C) 1999-2010 Instituto de Astronomia, UNAM, Mexico
%% 
%% This program is free software; you can redistribute it and/or
%% modify it under the terms of the GNU General Public License
%% as published by the Free Software Foundation; either version 2
%% of the License, or (at your option) any later version.
%% 
%% This program is distributed in the hope that it will be useful,
%% but WITHOUT ANY WARRANTY; without even the implied warranty of
%% MERCHANTABILITY or FITNESS FOR A PARTICULAR PURPOSE.  See the
%% GNU General Public License for more details.
%% 
%% You should have received a copy of the GNU General Public License
%% along with this program; if not, write to the Free Software
%% Foundation, Inc., 59 Temple Place - Suite 330, Boston, MA  02111-1307, USA.






\def\bibitem[#1<#2>]#3{\par\hangindent=1.5em\hangafter=1
  \if@filesw
  {\def\protect##1{\string ##1\space}
    \immediate\write\@auxout{\string\lbibcite{#3}{#1}{#2}}}
  \fi}
\@redefinebibitem % what it says
\def\lbibcite#1#2#3{\global\@namedef{b@#1}{#2}
                    \global\@namedef{by@#1}{#3}}

\def\@bcite#1#2{({#1\if@tempswa , #2\fi})}
\def\@pcite#1#2{#1\if@tempswa , #2\fi}
\def\@brcite{\ (}
\def\@ebrcite{)}
\def\@fmttop#1{\@fmtcite[\csname #1\endcsname]}
\def\@citex[#1]#2{%
%%  \if@filesw\immediate\write\@auxout{\string\citation{#2}}\fi
%% WJH added 16 Jan 2002
  \@for\@citeb:=#2\do{%
    \edef\@citeb{\expandafter\@firstofone\@citeb:\articlename}%
    \if@filesw\immediate\write\@auxout{\string\citation{\@citeb}}\fi}%
%% to here
  \def\@citea{}%
  \@cite{\@for\@citeb:=#2\do
    {%
      \edef\@citeb{\expandafter\@firstofone\@citeb:\articlename}%
      \@citea\def\@citea{;\ }\@ifundefined
      {b@\@citeb}{{\bf ?}\@warning {Citation `\@citeb' on page \thepage \space undefined}}%
      {\csname b@\@citeb\endcsname\@citebro\csname by@\@citeb\endcsname\@citebrc}}}%
  {#1}}

\def\cite{\@ifnextchar [{\let\@citebro=\space
                          \let\@citebrc=\relax
                          \let\@cite=\@bcite\@tempswatrue\@citex}
                        {\let\@citebro=\space
                          \let\@citebrc=\relax
                          \let\@cite=\@bcite\@tempswafalse\@citex[]}}
\def\pcite{\@ifnextchar [{\let\@citebro=\space
                          \let\@citebrc=\relax
                          \let\@cite=\@pcite\@tempswatrue\@citex}
                        {\let\@citebro=\space
                          \let\@citebrc=\relax
                          \let\@cite=\@pcite\@tempswafalse\@citex[]}}
\def\scite{\@ifnextchar [{\let\@citebro=\@brcite
                          \let\@citebrc=\@ebrcite
                          \let\@cite=\@pcite\@tempswatrue\@citex}
                        {\let\@citebro=\@brcite
                          \let\@citebrc=\@ebrcite
                          \let\@cite=\@pcite\@tempswafalse\@citex[]}}

\renewenvironment{thebibliography}
  {\section*{REFERENCES}\vspace*{\afterbiblioskip}\small
    \if@book\RereadBiblio\fi
    \sloppy\clubpenalty4000\widowpenalty4000
    \sfcode`\.=1000\relax
    \setlength\parindent{\z@}\setlength\parskip{\z@}}
  {\relax\normalsize\@outputaddressesifnecessary}








\endinput
%%
%% End of file `rmaacite.sty'.
\makeatother % read it in here so that all changes
    \else                                         % are local to paper
    \if@debug\ClassWarning{rmaa}{Package \@packagename\space was requested by \articlename\space}\fi
    \fi} 
  %% Don't allow individual articles to mess with the pages, volume or year
  \renewcommand{\SetFirstPage}[1]{\relax}
  \renewcommand{\SetYear}[1]{\relax}
  \renewcommand{\SetVolume}[1]{\relax}
  \renewcommand{\SetConfTitle}[1]{\relax}
  %% Don't allow individual articles to change these
  \renewcommand{\nonstopmode}{\relax}
  \renewcommand{\scrollmode}{\relax}
  \renewcommand{\outputplates}{\relax}
  %% XXXX need to check if using rmaacite and act accordingly
  %% avoid multiple declarations of the same length
  \let\oldnewlength=\newlength
  \renewcommand{\newlength}[1]{%
    \edef\reserved@a{\expandafter\@gobble\string ##1}%
    \@ifundefined{\reserved@a}{%
      \oldnewlength{##1}}{%
      \if@debug{\ClassInfo{rmaa}{\reserved@a\space already defined: doing nothing}}\fi%
      \relax}}%
  \renewcommand\newcommand{\providecommand}
  }%
{\endgroup\onecolumn}
     

%    \end{macrocode}
% Command to be used in `articles' environment for including an
% article
%    \begin{macrocode}
\newcommand{\article}[2][]{%
  % optional first argument is extra text to put in brackets after title in TOC
  % e.g. ``Invited Talk'' or ``Poster''
  \renewcommand{\articlename}{#2}
  % reset loads of counters
  \setcounter{footnote}{0}
  \setcounter{figure}{0}
  \setcounter{table}{0}
  \setcounter{section}{0}
  \setcounter{equation}{0}
  \setcounter{rmaabib}{0}
  \setcounter{author}{0}
  \setcounter{AUTHOR}{1}
  \gdef\ADS@authors{}%
  \gdef\ADS@affil{}%
  \gdef\ADS@affiliations{}%
  \gdef\ADS@references{}%
  \@firstauthortrue
  \@firstaffiltrue
  \@firstaltaffiltrue
  \@firstkwdtrue
  \@rmcitefalse
  \if@suppressfulladdresses\else\@needaddressestrue\fi
  \renewcommand{\@rmaaArg}{#1}
  \ifx\@rmaaArg\@rmaaEmpty
     \renewcommand{\@TOCaftertitle}{}
  \else
     \renewcommand{\@TOCaftertitle}{ (#1)}
  \fi
%    \end{macrocode}
% In the main journal articles must start on an odd page. 
%    \begin{macrocode}
  \if@journal\clearpage\thispagestyle{empty}\cleardoublepage\fi
%% Now input the file, then run the end-of-article hook (reinstated 28
  %% Sep 2001)
  {\input{\@articledir/\articlename.tex}\@endofarticlehook}}
%    \end{macrocode}
% Command to separate articles into sessions. Write an entry in the
% TOC, makes sure we start on an odd page (added 06 Jun 2000), then do
% a title page if \cs{SessionTitlePages} was called. 
% \changes{3.12}{04 Feb 2003}{Added option to have title page for each
% session. Turn it on with \cs{SessionTitlePages}.}
%    \begin{macrocode}
\newcommand{\session}[1]{
  %% this does nothing except write an entry in the TOC
  \addcontentsline{toc}{part}{\rmaa@TOCPartStyle{#1}}
  \addtocontents{htc}{\protect\HyperTOCpart{#1}}
  %% and also makes sure we start on an odd page (added 06 Jun 2000)
  \clearpage\thispagestyle{empty}\cleardoublepage
  \stepcounter{part}
  \if@sessiontitlepages{%
    \onecolumn
    \thispagestyle{empty}
    \renewcommand{\articlename}{Section\thepart TitlePage}
    \label{RMAAFirstPage}
    \make@session@title{#1}
    \label{RMAALastPage}
    \clearpage\thispagestyle{empty}\cleardoublepage
  }\fi
  }
\newcommand{\make@session@title}[1]{%
  {%
    \setlength\parindent{0pt}
    \vspace*{\stretch{1}}
    \rule{\linewidth}{1mm}
    \begin{flushright}\Huge\bfseries\boldmath
      #1
    \end{flushright}
    \rule{\linewidth}{1mm}
    \vspace*{\stretch{2}}
  }
}
%    \end{macrocode}
% Command to turn on separate title pages for each session. 
%    \begin{macrocode}
\newcommand{\SessionTitlePages}{\@sessiontitlepagestrue}
\newcommand{\BoldSessionTitles}{%
  \renewcommand{\rmaa@TOCPartStyle}[1]{\textbf{\rmaa@TOCTitleStyle{##1}}}
  }
\newcommand{\ArticleDir}[1]{\renewcommand{\@articledir}{#1/}\graphicspath{{1/}}}
\newcommand{\@articledir}{.}
%\newcommand{\ps@rmaafront}{\ps@rmaa}% page style for front matter (same as ever, actually)
\newcommand{\ps@rmaafront}{% page style for front matter:  page # in each foot
  \renewcommand\@oddfoot{\hfil{\thepage}}%
  \renewcommand\@evenfoot{{\thepage}\hfil}%
  % shorttitle in both heads
  \renewcommand\@evenhead{\hfil\@shorttitle\hfil}%
  \renewcommand\@oddhead{\@evenhead}
  }
\newcommand{\ps@rmaafrontfirstpage}{% page style for first page of front matter
  % page # in each foot, on outside edge
  \renewcommand\@oddfoot{\hfil\thepage}%
  \renewcommand\@evenfoot{\thepage\hfil}%
  % nothing in the even heads
  \renewcommand\@evenhead{\relax}%
  \renewcommand\@oddhead{\relax}%
  }
%% slight modification to the standard \frontmatter and \mainmatter
%% commands, just to set the \pagestyle
\renewcommand\frontmatter{%
  \latexonecolumn
  \cleardoublepage
  \@frontmattertrue
  \@mainmatterfalse
  \pagenumbering{roman}
  \pagestyle{rmaafront}
  \addtocounter{page}{2} % leave space for the title page
  }
\renewcommand\mainmatter{%
  \pagestyle{empty}
  \if@rmtwocol\twocolumn\fi
  \cleardoublepage
  \@frontmatterfalse
  \@mainmattertrue
  \pagenumbering{arabic}
  \pagestyle{rmaa}
  }
\renewcommand\backmatter{%
  \clearpage\pagestyle{empty}\cleardoublepage
  \@mainmatterfalse
  %% put extra space in the TOC
  \addtocontents{toc}{\bigskip}
  }
%    \end{macrocode}
%
% \subsection{Front sections}\label{sec:front-sections}
% \begin{macro}{\notitle}
% \begin{macro}{\shortname}
% \begin{macro}{frontsection}
% The environment \texttt{frontsection} is for doing a section in the
% front matter (preface, list of participants, or whatever). It has 
% one optional argument:
% \begin{enumerate}
% \item[\#1] Short title (no embedded spaces) for
% use in making PS files. E.g., ``list\_of\_participants''.   
% \end{enumerate}
% and two mandatory arguments: 
% \begin{enumerate}
% \item[\#2] Title of section (e.g., `Preface')
% \item[\#3] Commands to be evaluated before processing the contents
% of the environment. This is mainly useful for putting in the things
% that would go in the preamble of a normal article  (e.g.,
% |\listofauthors|, |\indexauthor|, |\notitle|) 
% \end{enumerate}
%    \begin{macrocode}
\newcommand{\notitle}{\@titlevisiblefalse} \newcommand{\shortname}{}
\newenvironment{frontsection}[3][]{%
  \begingroup
  \renewcommand{\shortname}{#1}
  \ifx\shortname\@rmaaEmpty
     \renewcommand{\articlename}{#2}
  \else
     \renewcommand{\articlename}{\shortname}
  \fi
  \@begofarticlehook
  \@titlevisibletrue % default
  \thispagestyle{rmaafrontfirstpage}
  \listofauthors{} % blank out this in case it is not specified
  \title{{#2}}\shorttitle{#2}\shortauthor{#2}
  #3 % evaluate the command given as second mandatory arg
  \addcontentsline{toc}{chapter}{\protect\@fiddleauthors{\@TOCtitle}{\rmaaTOCauthors}}%
  \addtocontents{htc}{\protect\HyperTOCEntry{\@rawtitle}{\@rawauthors}{\articlename}}
  \if@titlevisible
     \vspace*{1cm}\centerline{\rmaa@FrontSecStyle{\@title}}\vspace*{1cm}
  \fi
  \setlength{\parskip}{3ex}
  }{\frontlabel{RMAALastPage}\clearpage\endgroup}
%    \end{macrocode}
% \end{macro}
% \end{macro}
% \end{macro}
%    \begin{macrocode}
%% Definition of the command \frontlabel, which rights out something like `F4' to
%% the .aux file, instead of `\textlatin  {i}'. This is so it is easy to post-process
%% the file for making PS files of individual sections
\newcommand{\frontlabel}[1]{%
  \@bsphack
  \protected@write\@auxout{}%
     {\string\newlabel{\articlename:#1}{{}{F\@arabic{\c@page}}}}
  \@esphack
  }
%% User command to say X, Y, Z,\\ Editors of the Proceedings
\newcommand{\EditorSignature}{\vspace*{\parskip}\rightline{\@editors}%
  \rightline{Editors of the proceedings}}
%% Fiddle with \thesection, etc. so that they don't include the chapter number
%% Basically go back to the `article' definitions of these
\renewcommand{\thesection}{\@arabic\c@section}
\renewcommand{\thefigure}{\@arabic\c@figure}
\renewcommand{\thetable}{\@arabic\c@table}

\newcommand{\@oldargs}{}
\newcommand{\@newargs}{}
%% add TOC entry for each paper
%\renewcommand{\title}[1]%
%{\renewcommand{\@title}{\rmaa@TitleStyle{#1}}%
%\addtocounter{chapter}{1}\addcontentsline{toc}{chapter}{#1}%
%}

%% Formatting of the table of contents
%% dimensions of various parts
\newlength{\TOCtextwidth}\setlength{\TOCtextwidth}{12.65cm}
\newlength{\TOCpagenumwidth}\setlength{\TOCpagenumwidth}{1.3cm}
\newlength{\TOCentrywidth}\setlength{\TOCentrywidth}{\TOCtextwidth}
\addtolength{\TOCentrywidth}{-\TOCpagenumwidth} 
\newlength{\TOCextramargin}\setlength{\TOCextramargin}{1.25cm}
\newcommand{\TOCbox}[1]{\hspace{\TOCextramargin}\parbox{\TOCtextwidth}{#1}\par}
\newlength{\TOCextramarginx}\setlength{\TOCextramarginx}{\TOCextramargin}\addtolength{\TOCextramarginx}{4.3cm}
\newlength{\TOCtextwidthx}\setlength{\TOCtextwidthx}{\TOCtextwidth}\addtolength{\TOCtextwidthx}{-4.3cm}
\newcommand{\TOCboxx}[1]{\hspace{\TOCextramarginx}\parbox{\TOCtextwidthx}{#1}\par}

\newcommand{\RMAAbannerimage}{rmsc}
\renewcommand\tableofcontents{%
  \renewcommand{\articlename}{CONTENTS}
  \begingroup
  \@begofarticlehook
  \if@twocolumn
     \@restonecoltrue\onecolumn
  \else
     \@restonecolfalse
  \fi
  %% 9pt for journal style
  \if@proceedings\else\small\fi
  \if@newstyle\small\fi %% same for new style
  \if@proceedings\thispagestyle{rmaafrontfirstpage}\else\thispagestyle{empty}\fi
  \shortauthor{Contents}\shorttitle{Contents}
  \vspace*{18mm}
  \if@printbanner
  %% put in the banner header (need to get better image of this)
     \TOCbox{\hspace*{-0.15cm}\includegraphics[width=0.9\TOCtextwidth]{\RMAAbannerimage}}
     \vspace*{8mm}
  %% volume and publication date (should really be in funny font)
     \TOCboxx{\@volumen\\[2mm] \@mes\ \@anyo}
  \else
  %% put in the banner header (need to get better image of this)
     \TOCbox{\hspace*{-0.15cm}\includegraphics[width=0.9\TOCtextwidth]{\RMAAbannerimage}}
     \vspace*{12.5mm}
  %% volume and publication date (should really be in funny font)
     \TOCboxx{~\\[2mm] ~}
  \fi
  \vspace*{16mm}
  %% This is the command that fiddles with the TOC entry so that the
  %% title is flushleft but the authors are flushright
  \renewcommand{\@fiddleauthors}[2]{%
    % first put the paper title
    \renewcommand{\\}{\relax} % temporarily turn off explicit line breaks just in case
    \renewcommand{\altaffilmark}[1]{\relax} % ditto with this
    \renewcommand{\textsuperscript}[1]{\relax} % ditto with this 
    ##1
    % now save the author list in a box
    \sbox\@tempboxa{##2}%
    % check that box is shorter than the total width available for the entry
    \ifdim \wd\@tempboxa <\TOCentrywidth
       % if so, then put it right here, flush against right margin
       \penalty10000\hskip1em plus10em\penalty0\hbox{}\nobreak\hfill\box\@tempboxa
       % Note that we force a space of at least 1em between title and authorlist, 
       % stretchable to 11em since we may have very short titles but still with an 
       % authorlist too long to go on the same line.  The only linebreak allowed is 
       % directly before the authorlist.
    \else
       % if it is too long, we set the original authorlist text in its own parbox
       \par\vskip1ex\parbox[b]{\TOCentrywidth}{\raggedleft ##2}
    \fi
    }
  \@starttoc{toc}%
  \if@restonecol\twocolumn\fi
  \frontlabel{RMAALastPage}
  \clearpage
  \thispagestyle{empty}
  \cleardoublepage % make sure we start on an odd page
  \endgroup

  %%% ALSO, do the HTML TOC
  %% The file that will be used to make the HTML TOC
  \@starttoc{htc}
  %% All the stuff in this file does nothing in the book style
  \addtocontents{htc}{\protect\HyperTOCmeeting{\@conftitle}{\@reunionplace}{\@reuniondate}}
  \addtocontents{htc}{\protect\HyperTOCeditors{\@editors}}
  \addtocontents{htc}{\protect\HyperTOCpreamble}
  \addtocontents{htc}{\protect\HyperTOCEntry{Table of Contents}{}{\articlename}}

  %% and the ADS file
  \ADS@openfile
  }
%% This is now much simpler
\renewcommand\l@chapter[2]{%TOC entry for articles
  \ifnum \c@tocdepth >\m@ne
    \vskip 1.0em \@plus\p@
    \TOCbox{%
      \parbox[b]{\TOCentrywidth}{#1}\parbox[b]{\TOCpagenumwidth}{\hfill#2}
    }
  \fi} 
\renewcommand\l@part[2]{%TOC entry for sessions: centered, no page number
  \ifnum \c@tocdepth >\m@ne
    \vskip 1.5em \@plus1.5em \goodbreak
    \TOCbox{%
      \parbox[b]{\TOCentrywidth}{\hfil#1\hfil}\hfill
    }
  \fi\nobreak} 

%% Formatting theindex
%% strongly encourage page breaks between index entries
\renewcommand{\@idxitem}{\par\penalty-1000\hangindent.45cm} 
\renewcommand{\indexspace}{\par\penalty-1000\vskip 10\p@ \@plus15\p@ \@minus3\p@\relax}
\renewenvironment{theindex}{%
  %% first, make sure we start on an odd page (06 Jun 2000)
  \clearpage\thispagestyle{empty}\cleardoublepage
  %% this looks a bit like an article but is two columns and without abstract, etc
  \renewcommand{\articlename}{INDEX}
  \title{Author Index}\shortauthor{Author Index}\shorttitle{Author Index}
  \listofauthors{}
  \addcontentsline{toc}{chapter}{\protect\@fiddleauthors{\@TOCtitle}{\rmaaTOCauthors}}%
  \addtocontents{htc}{\protect\HyperTOCEntry{\@rawtitle}{\@rawauthors}{\articlename}}
  \pagestyle{rmaa}
  \thispagestyle{rmaafirstpage}
  \@begofarticlehook
  %% 9pt for journal style
  \if@proceedings\else\small\fi
  \if@newstyle\small\fi
  \twocolumn[\@makeminititle]
  %% vertical spacing (see also definition of \@idxitem and \indexspace
  \parskip\z@ \@plus .6\p@\relax % allow quite a lot of leaway here to avoid breaks inside entries
  %% horizontal spacing, breaking and glue parameters (see also the stuff in the definition 
  %% of \indexauthor
  \parindent\z@ 
  \exhyphenpenalty10000 % don't break on explicit hyphens, it always looks bad
  \rightskip=0pt plus.2em % allow a tiny bit of leeway at the ends of lines
  \spaceskip.3333em plus.5em minus.15em % a lot of leeway between words
  \xspaceskip.5em plus.5em minus.4em % this seems to have no effect
  \let\item\@idxitem
  \renewcommand{\\}{\relax} % temporarily turn off explicit line breaks just in case
  \renewcommand{\altaffilmark}[1]{\relax} % ditto with this
  \renewcommand{\textsuperscript}[1]{\relax} % ditto with this
  }{\label{RMAALastPage}\clearpage} % clearpage necessary so that we know about \shortauthor



%% end definitions for book styles
\fi

%    \end{macrocode}
% In order to do away with plates, we sometimes want a figure to
% appear in color only in the online version. The following command
% can bve used in a figure caption to typeset a note to this effect. 
%    \begin{macrocode}
\newcommand\ColorOnlineOnly{%
  \textsc{Note: this figure is available in color in the electronic
    version of this article, obtainable from}
  \texttt{http://\discretionary{}{}{}www.\discretionary{}{}{}astroscu.\discretionary{}{}{}unam.\discretionary{}{}{}mx/\discretionary{}{}{}\~{}rmaa/}.}
%% Plates
%% Define 2 counters: plate is the number in the whole document, plateloop is used 
%% internally by \outputplates
\newcounter{plate}
\newcounter{plateloop}
%% Currently limited to 15 plates in document
%% when we do a label, we need to write to \newlabel entries to the .aux file, one 
%% that says RMAAPlateX so that it can be referred to in \ouputplates, and one that
%% says articlename:RMAAPlateX so that the dvips program can find it and asociate it
%% with the right article
\newcommand{\@platelabeli}{\oldlabel{RMAAPlate1}\label{RMAAPlate1}}
\newcommand{\@platelabelii}{\oldlabel{RMAAPlate2}\label{RMAAPlate2}}
\newcommand{\@platelabeliii}{\oldlabel{RMAAPlate3}\label{RMAAPlate3}}
\newcommand{\@platelabeliv}{\oldlabel{RMAAPlate4}\label{RMAAPlate4}}
\newcommand{\@platelabelv}{\oldlabel{RMAAPlate5}\label{RMAAPlate5}}
\newcommand{\@platelabelvi}{\oldlabel{RMAAPlate6}\label{RMAAPlate6}}
\newcommand{\@platelabelvii}{\oldlabel{RMAAPlate7}\label{RMAAPlate7}}
\newcommand{\@platelabelviii}{\oldlabel{RMAAPlate8}\label{RMAAPlate8}}
\newcommand{\@platelabelix}{\oldlabel{RMAAPlate9}\label{RMAAPlate9}}
\newcommand{\@platelabelx}{\oldlabel{RMAAPlate10}\label{RMAAPlate10}}
\newcommand{\@platelabelxi}{\oldlabel{RMAAPlate11}\label{RMAAPlate11}}
\newcommand{\@platelabelxii}{\oldlabel{RMAAPlate12}\label{RMAAPlate12}}
\newcommand{\@platelabelxiii}{\oldlabel{RMAAPlate13}\label{RMAAPlate13}}
\newcommand{\@platelabelxiv}{\oldlabel{RMAAPlate14}\label{RMAAPlate14}}
\newcommand{\@platelabelxv}{\oldlabel{RMAAPlate15}\label{RMAAPlate15}}
\newcommand{\@outputplatelabeli}{\oldlabel{RMAAOutputPlate1}}
\newcommand{\@outputplatelabelii}{\oldlabel{RMAAOutputPlate2}}
\newcommand{\@outputplatelabeliii}{\oldlabel{RMAAOutputPlate3}}
\newcommand{\@outputplatelabeliv}{\oldlabel{RMAAOutputPlate4}}
\newcommand{\@outputplatelabelv}{\oldlabel{RMAAOutputPlate5}}
\newcommand{\@outputplatelabelvi}{\oldlabel{RMAAOutputPlate6}}
\newcommand{\@outputplatelabelvii}{\oldlabel{RMAAOutputPlate7}}
\newcommand{\@outputplatelabelviii}{\oldlabel{RMAAOutputPlate8}}
\newcommand{\@outputplatelabelix}{\oldlabel{RMAAOutputPlate9}}
\newcommand{\@outputplatelabelx}{\oldlabel{RMAAOutputPlate10}}
\newcommand{\@outputplatelabelxi}{\oldlabel{RMAAOutputPlate11}}
\newcommand{\@outputplatelabelxii}{\oldlabel{RMAAOutputPlate12}}
\newcommand{\@outputplatelabelxiii}{\oldlabel{RMAAOutputPlate13}}
\newcommand{\@outputplatelabelxiv}{\oldlabel{RMAAOutputPlate14}}
\newcommand{\@outputplatelabelxv}{\oldlabel{RMAAOutputPlate15}}
\newcommand{\@platerefi}{\oldpageref{RMAAPlate1}}
\newcommand{\@platerefii}{\oldpageref{RMAAPlate2}}
\newcommand{\@platerefiii}{\oldpageref{RMAAPlate3}}
\newcommand{\@platerefiv}{\oldpageref{RMAAPlate4}}
\newcommand{\@platerefv}{\oldpageref{RMAAPlate5}}
\newcommand{\@platerefvi}{\oldpageref{RMAAPlate6}}
\newcommand{\@platerefvii}{\oldpageref{RMAAPlate7}}
\newcommand{\@platerefviii}{\oldpageref{RMAAPlate8}}
\newcommand{\@platerefix}{\oldpageref{RMAAPlate9}}
\newcommand{\@platerefx}{\oldpageref{RMAAPlate10}}
\newcommand{\@platerefxi}{\oldpageref{RMAAPlate11}}
\newcommand{\@platerefxii}{\oldpageref{RMAAPlate12}}
\newcommand{\@platerefxiii}{\oldpageref{RMAAPlate13}}
\newcommand{\@platerefxiv}{\oldpageref{RMAAPlate14}}
\newcommand{\@platerefxv}{\oldpageref{RMAAPlate15}}
\newcounter{figplatei}
\newcounter{figplateii}
\newcounter{figplateiii}
\newcounter{figplateiv}
\newcounter{figplatev}
\newcounter{figplatevi}
\newcounter{figplatevii}
\newcounter{figplateviii}
\newcounter{figplateix}
\newcounter{figplatex}
\newcounter{figplatexi}
\newcounter{figplatexii}
\newcounter{figplatexiii}
\newcounter{figplatexiv}
\newcounter{figplatexv}
%% Printed output should use the counter `plate'
\renewcommand{\theplate}{\@arabic\c@plate}
\renewcommand{\theplateloop}{\@arabic\c@plateloop}
%% Command to output all the saved plates in this paper
\newcommand{\outputplates}{%
  \setcounter{plateloop}{0}
  \if@book\else\@endofarticlehook\fi % not entirely sure why, but this
                                     % seems to be needed
  \clearpage % first output the last page of text
  % now set pagestyle for blank even pages and make sure we start on an odd page
  \pagestyle{empty}
  \cleardoublepage
  \onecolumn
  \parindent\z@
  % first put in an extra sheet
  % the main purpose of this is so there is always a blank even page before 
  % each plate for filling in the PS file if the article has an odd number
  % of pages (see splitbook.sh)
  \null\vfil\centerline{PLATES}\vfil\clearpage\pagestyle{rmaaplate}\cleardoublepage
  % output each plate in turn
  \@whilenum\value{plateloop}<\value{plate}\do{%
    \addtocounter{plateloop}{1}
    % restore the figure counter from \@figcounti, \@figcountii, etc
    \setcounter{figure}{\value{figplate\roman{plateloop}}}
    % spew out the plate (in its own group for safety)
    {\csname @plate\roman{plateloop}\endcsname\par}
    }
  \pagestyle{rmaa}
  % reset the page number as if nothing had happened
  \addtocounter{page}{-\value{plate}}
  \addtocounter{page}{-\value{plate}}
  }
\newcommand\figcaption{%
  \def\@captype{figure}
  \renewcommand{\@makecaption}{\@makefigcap}
  \caption}
\newcommand\tabcaption{%
  \def\@captype{table}
  \renewcommand{\@makecaption}{\@maketabcap}
  \caption}

\newcommand{\plate}[2][1]{%
  %% User command to make a plate page: 
  %% Mandatory argument is what goes on the plate - this should not 
  %% include any floating environments (e.g. figure). Instead, \figcaption 
  %% should be used to do the captions. 
  %% Optional argument is the number of figures (i.e. captions) on the plate
  %% page (default is 1). This is necessary for incrementing the figure counter. 
  %%
  %% Note that we have to use a \gdef rather than an \xdef when we 
  %% save the plate output in a command. This means that in the book styles when it 
  %% comes to spewing out the plate, we will have forgotten about any user macros 
  %% that may have been defined. This can be got round by repeating the macros inside the plate.
  %% (Presumably, one would use a \newcommand rather than \renewcommand if they were new.)
  % increment global and local counters
  \addtocounter{plate}{1}
  % write appropriate label to .aux file
  \csname @platelabel\roman{plate}\endcsname
  % define a global macro called \@platei, \@plateii, \@plateiii, \@plateiv, ...
  \expandafter\gdef\csname @plate\roman{plate}\endcsname{%
    \csname @outputplatelabel\roman{plateloop}\endcsname
    \null\vfil #2 \vfil
    \cleardoublepage
    }
  % save the current shortauthor and shorttitle
  % need to be \xdef to take current values of \@shortauthor, etc 
  \expandafter\xdef\csname @shorta\roman{plate}\endcsname{\@shortauthor}
  \expandafter\xdef\csname @shortt\roman{plate}\endcsname{\@shorttitle}
  % save the current value of the figure counter for later use in \outputplates
  \setcounter{figplate\roman{plate}}{\value{figure}}
  % increment the figure counter according to the optional argument
  \addtocounter{figure}{#1}
  }
%% page style for plate pages 
\newcommand{\ps@rmaaplate}{%
  % Plate X (see page Y) in the foot
  \renewcommand\@oddfoot{\hfil\csname @shorta\roman{plateloop}\endcsname\ 
    (see page~\csname @plateref\roman{plateloop}\endcsname)\hfil}%
  \renewcommand\@evenfoot{\relax}%
  % nothing in even heads (blank pages)
  \renewcommand\@evenhead{\relax}
  % short title and page num in the odd heads
  \renewcommand\@oddhead{\hfil\csname @shortt\roman{plateloop}\endcsname%
    \hfil\llap{PLATE~\theplateloop}}
  }





%% Things for the User Guide
\if@guide
%% stolen from doc.sty
\def\MakeShortVerb#1{%
  \expandafter\ifx\csname cc\string#1\endcsname\relax
  \@shortvrbinfo{Made }{#1}%
  \add@special{#1}%
  \expandafter
  \xdef\csname cc\string#1\endcsname{\the\catcode`#1}%
  \begingroup
  \catcode`\~\active  \lccode`\~`#1%
  \lowercase{%
    \global\expandafter\let
    \csname ac\string#1\endcsname~%
    \gdef~{\verb~}}%
  \endgroup
  \global\catcode`#1\active
  \else
  \@shortvrbinfo\@empty{#1 already}%
  \fi}
\def\DeleteShortVerb#1{%
  \expandafter\ifx\csname cc\string#1\endcsname\relax
  \else
  \@shortvrbinfo{Deleted }{#1 as}%
  \rem@special{#1}%
  \global\catcode`#1\csname cc\string#1\endcsname
  \global \expandafter\let \csname cc\string#1\endcsname \relax
  \ifnum\catcode`#1=\active
  \begingroup
  \catcode`\~\active   \lccode`\~`#1%
  \lowercase{%
    \global\expandafter\let\expandafter~%
    \csname ac\string#1\endcsname}%
  \endgroup \fi \fi}
\def\@shortvrbinfo#1#2{%
  \PackageInfo{doc}{%
    #1\expandafter\@gobble\string#2 a short reference
    for \string\verb}}
\def\add@special#1{%
  \rem@special{#1}%
  \expandafter\gdef\expandafter\dospecials\expandafter
  {\dospecials \do #1}%
  \expandafter\gdef\expandafter\@sanitize\expandafter
  {\@sanitize \@makeother #1}}
\def\rem@special#1{%
  \def\do##1{%
    \ifnum`#1=`##1 \else \noexpand\do\noexpand##1\fi}%
  \xdef\dospecials{\dospecials}%
  \begingroup
  \def\@makeother##1{%
    \ifnum`#1=`##1 \else \noexpand\@makeother\noexpand##1\fi}%
  \xdef\@sanitize{\@sanitize}%
  \endgroup}
%% define | as shortcut for \verbatim
\MakeShortVerb{\|}
%% comand to specify arguments to LaTeX commands
\newcommand{\rmArg}[1]{%
  \ensuremath{\langle}{\scriptsize \textsf{\MakeUpperCase{#1}}}\ensuremath{\rangle}}
%\newcommand{\rmArg}[1]{\ensuremath{\langle}\textsc{#1}\ensuremath{\rangle}}
\providecommand{\BibTeX}{{\rmfamily B\kern-.05em%
    \textsc{i\kern-.025em b}\kern-.08em%
    T\kern-.1667em\lower.7ex\hbox{E}\kern-.125emX}}
%% Environment for code examples
\newenvironment{Code}{%
  \begin{list}{}{\setlength{\leftmargin}{0.05\columnwidth}} 
  \item[]}{\end{list}}
\fi
%</package>
%    \end{macrocode}

% \section{User commands for astronomical articles}
% \label{sec:user-commands}
% Some random definitions. These go both into the main class file and
% also into |adshtml.sty|.
% 
% Definitions of the journal name in different incarnations.  
%    \begin{macrocode}
%<*package|adshtml-style>
\newcommand{\RMAA}{RMxAA}
\newcommand{\RMAAlong}{Revista Mexicana de Astronom\'{\i}a y Astrof\'{\i}sica}
\newcommand{\RMAASC}{RMxAC}
\newcommand{\RMAASClong}{\RMAA\ (Serie de Conferencias)}
%    \end{macrocode}
% \begin{macro}{\ibidrule}
% This puts a horizontal rule with a period after it. It is used when
% the author list of a reference is identical to the preceding one. 
%    \begin{macrocode}
\newcommand{\ibidrule}{\rule[-0.03cm]{4.7em}{.01cm}.\ }
%    \end{macrocode}
% \end{macro}
% \begin{macro}{\ThisVolume}
% References to other articles in the same volume can be coded as,
% e.g., |\bibitem{} Bloggs, J. \ThisVolume{bloggs}|.
%    \begin{macrocode}
\newcommand{\ThisVolume}[1]{%
  \@year, \if@proceedings\RMAASC\else\RMAA\fi, \@volume,
  \if@book\oldpageref{#1:RMAAFirstPage}\else 000\fi\ (this volume)}
%    \end{macrocode}
% \end{macro}
% Other stuff that may be useful in ref lists.
%    \begin{macrocode}
\newcommand{\etal}{et al.\@}
%    \end{macrocode}
% Various new commands from AAS macros
%    \begin{macrocode}
\def\acknowledgments{\vskip 3ex plus .8ex minus .4ex}
\let\acknowledgements=\acknowledgments                  % second spelling
\let\jnl@style=\rm
\def\ref@jnl#1{{\jnl@style#1}}
\DeclareRobustCommand\aj{\ref@jnl{AJ}}                   
\DeclareRobustCommand\araa{\ref@jnl{ARA\&A}}             
\DeclareRobustCommand\apj{\ref@jnl{ApJ}}                 
\DeclareRobustCommand\apjl{\ref@jnl{ApJ}}                
\DeclareRobustCommand\apjs{\ref@jnl{ApJS}}               
\DeclareRobustCommand\ao{\ref@jnl{Appl.~Opt.}}           
\DeclareRobustCommand\apss{\ref@jnl{Ap\&SS}}             
\DeclareRobustCommand\aap{\ref@jnl{A\&A}}                
\DeclareRobustCommand\aapr{\ref@jnl{A\&A~Rev.}}          
\DeclareRobustCommand\aaps{\ref@jnl{A\&AS}}              
\DeclareRobustCommand\azh{\ref@jnl{AZh}}                 
\DeclareRobustCommand\baas{\ref@jnl{BAAS}}               
\DeclareRobustCommand\jrasc{\ref@jnl{JRASC}}             
\DeclareRobustCommand\memras{\ref@jnl{MmRAS}}            
\DeclareRobustCommand\mnras{\ref@jnl{MNRAS}}             
\DeclareRobustCommand\pra{\ref@jnl{Phys.~Rev.~A}}        
\DeclareRobustCommand\prb{\ref@jnl{Phys.~Rev.~B}}        
\DeclareRobustCommand\prc{\ref@jnl{Phys.~Rev.~C}}        
\DeclareRobustCommand\prd{\ref@jnl{Phys.~Rev.~D}}        
\DeclareRobustCommand\pre{\ref@jnl{Phys.~Rev.~E}}        
\DeclareRobustCommand\prl{\ref@jnl{Phys.~Rev.~Lett.}}    
\DeclareRobustCommand\pasp{\ref@jnl{PASP}}               
\DeclareRobustCommand\pasj{\ref@jnl{PASJ}}               
\DeclareRobustCommand\qjras{\ref@jnl{QJRAS}}             
\DeclareRobustCommand\skytel{\ref@jnl{S\&T}}             
\DeclareRobustCommand\solphys{\ref@jnl{Sol.~Phys.}}      
\DeclareRobustCommand\sovast{\ref@jnl{Soviet~Ast.}}      
\DeclareRobustCommand\ssr{\ref@jnl{Space~Sci.~Rev.}}     
\DeclareRobustCommand\zap{\ref@jnl{ZAp}}                 
\DeclareRobustCommand\nat{\ref@jnl{Nature}}              
\DeclareRobustCommand\iaucirc{\ref@jnl{IAU~Circ.}}
\DeclareRobustCommand\aplett{\ref@jnl{Astrophys.~Lett.}}
\DeclareRobustCommand\apspr{\ref@jnl{Astrophys.~Space~Phys.~Res.}}
\DeclareRobustCommand\bain{\ref@jnl{Bull.~Astron.~Inst.~Netherlands}}
\DeclareRobustCommand\fcp{\ref@jnl{Fund.~Cosmic~Phys.}}
\DeclareRobustCommand\gca{\ref@jnl{Geochim.~Cosmochim.~Acta}}
\DeclareRobustCommand\grl{\ref@jnl{Geophys.~Res.~Lett.}}
\DeclareRobustCommand\jcp{\ref@jnl{J.~Chem.~Phys.}}      
\DeclareRobustCommand\jgr{\ref@jnl{J.~Geophys.~Res.}}    
\DeclareRobustCommand\jqsrt{\ref@jnl{J.~Quant.~Spec.~Radiat.~Transf.}}
\DeclareRobustCommand\memsai{\ref@jnl{Mem.~Soc.~Astron.~Italiana}}
\DeclareRobustCommand\nphysa{\ref@jnl{Nucl.~Phys.~A}}
\DeclareRobustCommand\physrep{\ref@jnl{Phys.~Rep.}}
\DeclareRobustCommand\physscr{\ref@jnl{Phys.~Scr}}
\DeclareRobustCommand\planss{\ref@jnl{Planet.~Space~Sci.}}
\DeclareRobustCommand\procspie{\ref@jnl{Proc.~SPIE}}
\let\astap=\aap
\let\apjlett=\apjl
\let\apjsupp=\apjs
\let\applopt=\ao
\def\phn{\phantom{0}}
\def\phd{\phantom{.}}
\def\phs{\phantom{\ensuremath{-}}}
\def\phm#1{\phantom{#1}}
%    \end{macrocode}
% All these are redone as robust and using more modern commands where
% possible. The robustness is so they will not get expanded if used in
% abstracts and titles when these are passed to the ADS file. 
% \changes{3.16}{2003/04/07}{
%  Added \cs{Color} command for doing photometric colors.}
%    \begin{macrocode}
%%\def\sun{\hbox{$\odot$}}
\DeclareRobustCommand\sun{\ensuremath{\odot}}
%%\def\earth{\hbox{$\oplus$}}
\DeclareRobustCommand\earth{\ensuremath{\oplus}}
%%\def\lesssim{\mathrel{\hbox{\rlap{\hbox{\lower4pt\hbox{$\sim$}}}\hbox{$<$}}}}
%%\def\gtrsim{\mathrel{\hbox{\rlap{\hbox{\lower4pt\hbox{$\sim$}}}\hbox{$>$}}}}
\DeclareRobustCommand\lesssim{\mathrel{\hbox{\smaller\smaller\rlap{\hbox{\lower3pt\hbox{$\sim$}}}\raise1.5pt\hbox{$<$}}}}
\DeclareRobustCommand\gtrsim{\mathrel{\hbox{\smaller\smaller\rlap{\hbox{\lower3pt\hbox{$\sim$}}}\raise1.5pt\hbox{$>$}}}}
%%\def\lta{\mathrel{\spose{\lower 3pt\hbox{$\mathchar"218$}}
%%     \raise 2.0pt\hbox{$\mathchar"13C$}}}
%%\def\gta{\mathrel{\spose{\lower 3pt\hbox{$\mathchar"218$}}
%%     \raise 2.0pt\hbox{$\mathchar"13E$}}}
\DeclareRobustCommand\sq{\hbox{\rlap{$\sqcap$}$\sqcup$}}
%%\def\arcdeg{\hbox{$^\circ$}}
%%\def\arcmin{\hbox{$^\prime$}}
%%\def\arcsec{\hbox{$^{\prime\prime}$}}
\DeclareRobustCommand\arcdeg{\ensuremath{^\circ}}
\DeclareRobustCommand\arcmin{\ensuremath{^\prime}}
\DeclareRobustCommand\arcsec{\ensuremath{^{\prime\prime}}}
%%\def\fd{\hbox{$.\!\!^{\rm d}$}}
%%\def\fh{\hbox{$.\!\!^{\rm h}$}}
%%\def\fm{\hbox{$.\!\!^{\rm m}$}}
%%\def\fs{\hbox{$.\!\!^{\rm s}$}}
%%\def\fdg{\hbox{$.\!\!^\circ$}}
%%\def\farcm{\hbox{$.\mkern-4mu^\prime$}}
%%\def\farcs{\hbox{$.\!\!^{\prime\prime}$}}
%%\def\fp{\hbox{$.\!\!^{\scriptscriptstyle\rm p}$}}
\DeclareRobustCommand\fd{\ensuremath{.\!\!^\mathrm{d}}}
\DeclareRobustCommand\fh{\ensuremath{.\!\!^\mathrm{h}}}
\DeclareRobustCommand\fm{\ensuremath{.\!\!^\mathrm{m}}}
\DeclareRobustCommand\fs{\ensuremath{\mkern0.5mu.\mkern-5mu^\mathrm{s}\mkern-1mu}}
\DeclareRobustCommand\fdg{\ensuremath{.\!\!^\circ}}
\DeclareRobustCommand\farcm{\ensuremath{.\mkern-4mu^\prime}}
\DeclareRobustCommand\farcs{\ensuremath{\mkern1mu.\mkern-6mu^{\prime\prime}\mkern-1.5mu}}
%%\DeclareRobustCommand\farcs{\ensuremath{.\!\!^{\prime\prime}}}
\DeclareRobustCommand\fp{\ensuremath{.\!\!^{\scriptscriptstyle\rm p}}}
%%\def\micron{\hbox{$\mu$m}}
\DeclareRobustCommand\micron{\ensuremath{\mu\mathrm{m}}}
\let\la=\lesssim                        
\let\ga=\gtrsim
%%\def\case#1#2{\hbox{$\frac{#1}{#2}$}}
%%\def\slantfrac#1#2{\hbox{$\,^#1\!/_#2$}}
%%\def\onehalf{\slantfrac{1}{2}}
%%\def\onethird{\slantfrac{1}{3}}
%%\def\twothirds{\slantfrac{2}{3}}
%%\def\onequarter{\slantfrac{1}{4}}
%%\def\threequarters{\slantfrac{3}{4}}
\DeclareRobustCommand\case[2]{\ensuremath{\frac{#1}{#2}}}
\DeclareRobustCommand\slantfrac[2]{\ensuremath{\,^#1\!/_#2}}
\DeclareRobustCommand\onehalf{\slantfrac{1}{2}}
\DeclareRobustCommand\onethird{\slantfrac{1}{3}}
\DeclareRobustCommand\twothirds{\slantfrac{2}{3}}
\DeclareRobustCommand\onequarter{\slantfrac{1}{4}}
\DeclareRobustCommand\threequarters{\slantfrac{3}{4}}
%%\def\ubvr{\hbox{$U\!BV\!R$}}            
%%\def\ub{\hbox{$U\!-\!B$}}               
%%\def\bv{\hbox{$B\!-\!V$}}               
%%\def\vr{\hbox{$V\!-\!R$}}               
%%\def\ur{\hbox{$U\!-\!R$}}
%    \end{macrocode}
% New commands for ``photometric'' colors 
%    \begin{macrocode}
\DeclareRobustCommand\ubvr{\ensuremath{U\!BV\!R}}            
\DeclareRobustCommand\Color[2]{%
  \ensuremath{%
    #1\!\ifx#1V\!\fi-\!#2
  }}
\DeclareRobustCommand\ub{\Color{U}{B}}
\DeclareRobustCommand\bv{\Color{B}{V}}
\DeclareRobustCommand\vr{\Color{V}{R}}
\DeclareRobustCommand\ur{\Color{U}{R}}
%% New version of the \ion command, that actually works! 
%    \end{macrocode}
% \begin{macro}{\ion}
% This version of the |\ion| command allows you to write either
% |\ion{N}{4}| or |\ion{N}{iv}| in order to get N~IV. 
%    \begin{macrocode}
\DeclareRobustCommand\ion[2]{%
  \mbox{#1\kern0.2em%
    \smaller\rmfamily%
    \edef\@tempa{\@car#2\@nil}%
    \ifcat1\@tempa%
    \@Roman{#2}%
    \else%
    \uppercase{#2}%
    \fi}}
\newcounter{IonStage}
\DeclareRobustCommand\plainion[2]{\setcounter{IonStage}{#2}#1 \Roman{IonStage}}
%%\def\ion#1#2{\ensuremath{\mathrm{#1}\;}{\protect\small\rm\@Roman{#2}}}
%    \end{macrocode}
% \end{macro}
%    \begin{macrocode}
\DeclareRobustCommand\nodata{$\cdots$}
%    \end{macrocode}
% These are needed by the |\ion| command. Isn't there an easier way of
% doing this?
%    \begin{macrocode}
\let\@small=\small
\let\@normalsize=\normalsize
\let\@footnotesize=\footnotesize
\let\@scriptsize=\scriptsize
\let\@tiny=\tiny
\let\@large=\large
\let\@Large=\Large
\let\@LARGE=\LARGE
\let\@huge=\huge
\let\@Huge=\Huge
\newcommand{\smaller}{\small}
\renewcommand{\normalsize}{\renewcommand{\smaller}{\small}\@normalsize}
\renewcommand{\small}{\renewcommand{\smaller}{\footnotesize}\@small}
\renewcommand{\footnotesize}{\renewcommand{\smaller}{\scriptsize}\@footnotesize}
\renewcommand{\scriptsize}{\renewcommand{\smaller}{\tiny}\@scriptsize}
\renewcommand{\tiny}{\renewcommand{\smaller}{\tiny}\@tiny}
\renewcommand{\large}{\renewcommand{\smaller}{\normalsize}\@large}  
\renewcommand{\Large}{\renewcommand{\smaller}{\large}\@Large}  
\renewcommand{\LARGE}{\renewcommand{\smaller}{\Large}\@LARGE}  
\renewcommand{\huge}{\renewcommand{\smaller}{\LARGE}\@huge}  
\renewcommand{\Huge}{\renewcommand{\smaller}{\huge}\@Huge}  
%</package|adshtml-style>
%    \end{macrocode}


% \section{Other files}

% \subsection{The driver and make files}
% \label{sec:driver_make}
% These are put here for completeness, although they still need to be
% distributed separately since otherwise the user won't be able to
% unpack this file into its constituent parts.
% 
% The |Makefile| has two principal rules: |docs|, and |macros|. The
% |docs| rule produces a PDF and gzipped PS file of the documentation.
% I still haven't perfected the combination of rules required to get
% it to run \LaTeX{} enough times in all circumstances. The |macros|
% rule runs \LaTeX{} on |rmaa.ins| in order to unpack |rmaa.cls| and
% all the other little files (including the |Makefile| and the
% |rmaa.ins| files themselves\dots 
% 
% First set up lists of various types of file. 
%    \begin{macrocode}
%<*makefile>
DOCFILES=rmaa.pdf
MACRO_LTX_FILES=rmaa.cls rmaacite.sty adshtml.sty
SHELLSCRIPTS=makeads.sh
MACRO_OTHER_FILES=adshtml.perl sample.latex2html-init
MAKEFILE=Makefile
MACROFILES=$(MACRO_LTX_FILES) $(SHELLSCRIPTS) $(MACRO_OTHER_FILES)         
SKELETONFILES=ADS.tex 
%    \end{macrocode}
% Install the files read only into the package directory, specified via a path relative to the source directory. \changes{4.0}{2019/11/14}{Reorganize folder layout so that hardly anything is in the top level.  That way, things are clearer on github.}
%    \begin{macrocode}
INSTALL=install 
INSTALL_DIR=../packages
%    \end{macrocode}
% The default rule is to make everything. Installation must be
% explicitly asked for.
%    \begin{macrocode}
all: docs macros

install: install-docs install-macros

%    \end{macrocode}
% Most files are installed read-only to prevent editing. Skeleton
% files must be writeable. Must be careful not to overwrite the
% Makefile in the parent directory. 
%    \begin{macrocode}
install-macros: macros
        $(INSTALL) -m 444 $(MACRO_LTX_FILES) $(MACRO_OTHER_FILES) $(INSTALL_DIR)
        $(INSTALL) -m 555 $(SHELLSCRIPTS) $(INSTALL_DIR)
        $(INSTALL) -m 644 $(SKELETONFILES) $(INSTALL_DIR)

install-docs: docs
        $(INSTALL) rmaa.pdf $(INSTALL_DIR)/rm-fulldocs.pdf

%    \end{macrocode}
% Make the PDF and gzipped PS files. 
%    \begin{macrocode}
docs: $(DOCFILES)

%    \end{macrocode}
% Make everything else if either the |.dtx| or |.ins| files have
% changed. 
%    \begin{macrocode}
macros: rmaa.dtx rmaa.ins
%    \end{macrocode}
% Ensure that all the macro files can be written to then recreate
% them. The |touch| is there in case any of the files don't exist
% yet. 
%    \begin{macrocode}
        touch $(MACROFILES) $(MAKEFILE)
        chmod u+w $(MACROFILES) $(MAKEFILE) rmaa.ins
        latex "%%
%% This is file `rmaa.ins',
%% generated with the docstrip utility.
%%
%% The original source files were:
%%
%% rmaa.dtx  (with options: `ins')
%% 
%% PLEASE DO NOT EDIT THIS FILE! It is generated automatically from
%% the file rmaa.dtx using the accompanying Makefile and rmaa.ins
%% Any hand-edited changes will be lost
%% 
%% A LaTeX document class for the RMxAA
%% 
%% Author: William Henney (w.henney@crya.unam.mx)
%% 
%% Latest version available from: http://www.crya.unam.mx/rmaa
%% 
%% Copyright (C) 1999-2010 Instituto de Astronomia, UNAM, Mexico
%% 
%% This program is free software; you can redistribute it and/or
%% modify it under the terms of the GNU General Public License
%% as published by the Free Software Foundation; either version 2
%% of the License, or (at your option) any later version.
%% 
%% This program is distributed in the hope that it will be useful,
%% but WITHOUT ANY WARRANTY; without even the implied warranty of
%% MERCHANTABILITY or FITNESS FOR A PARTICULAR PURPOSE.  See the
%% GNU General Public License for more details.
%% 
%% You should have received a copy of the GNU General Public License
%% along with this program; if not, write to the Free Software
%% Foundation, Inc., 59 Temple Place - Suite 330, Boston, MA  02111-1307, USA.





%%
%% Installation batch file for the RMAA LaTeX package
%%
\def\batchfile{rmaa.ins}
\input docstrip.tex
\preamble

PLEASE DO NOT EDIT THIS FILE! It is generated automatically from
the file rmaa.dtx using the accompanying Makefile and rmaa.ins
Any hand-edited changes will be lost

A LaTeX document class for the RMxAA

Author: William Henney (w.henney@crya.unam.mx)

Latest version available from: http://www.crya.unam.mx/rmaa

Copyright (C) 1999-2010 Instituto de Astronomia, UNAM, Mexico

This program is free software; you can redistribute it and/or
modify it under the terms of the GNU General Public License
as published by the Free Software Foundation; either version 2
of the License, or (at your option) any later version.

This program is distributed in the hope that it will be useful,
but WITHOUT ANY WARRANTY; without even the implied warranty of
MERCHANTABILITY or FITNESS FOR A PARTICULAR PURPOSE.  See the
GNU General Public License for more details.

You should have received a copy of the GNU General Public License
along with this program; if not, write to the Free Software
Foundation, Inc., 59 Temple Place - Suite 330, Boston, MA  02111-1307, USA.
\endpreamble
\askforoverwritefalse
\generate{%
  \file{rmaa.cls}{\from{rmaa.dtx}{package}}
  \file{rmaa.ins}{\from{rmaa.dtx}{ins}}
  \file{rmaacite.sty}{\from{rmaa.dtx}{rmaacite}}
  \file{adshtml.sty}{\from{rmaa.dtx}{adshtml-style}}
  \file{ADS.tex}{\from{rmaa.dtx}{ads-container}}
  }
\generate{\nopreamble\nopostamble%
  \file{Makefile}{\from{rmaa.dtx}{makefile}}
  \file{adshtml.perl}{\from{rmaa.dtx}{adshtml-perl}}
  \file{makeads.sh}{\from{rmaa.dtx}{ads-shell-script}}
  \file{sample.latex2html-init}{\from{rmaa.dtx}{latex2html-init}}
  }

%% Local Variables:
%% mode: LaTeX
%% mode: Font-lock
%% End:









\endinput
%%
%% End of file `rmaa.ins'.
"
%    \end{macrocode}
% \TeX{} automatically adds the |.tex| extension to extensionless
% files. We fix this and take the opportunity to convert
% back to TABs all that it has lovingly converted to spaces. 
%    \begin{macrocode}
        cp -f $(MAKEFILE) $(MAKEFILE).BCK && unexpand $(MAKEFILE).tex >| $(MAKEFILE) && rm $(MAKEFILE).tex
%    \end{macrocode}
% Make the shell scripts executable and everything unwriteable (so
% they don't get hand-edited by mistake). However, undo this for the
% files that the user is \emph{supposed} to modify. 
%    \begin{macrocode}
        chmod a+x $(SHELLSCRIPTS)
        chmod a-w $(MACROFILES) $(MAKEFILE)
        chmod u+w $(SKELETONFILES)

%    \end{macrocode}
% We now use PDF\LaTeX{} for this.
%    \begin{macrocode}
rmaa.idx: rmaa.dtx
        @echo "REGENERATING .idx FILE"
        pdflatex "\nonstopmode% \def\fileversion{4.0} \def\filedate{2019/11/12}
% \MakeShortVerb{\|}
% \title{A \LaTeX{} document class for the Revista Mexicana de
% Astronom\'\i a y Astrof\'\i sica\thanks{This file has version
% \fileversion{} last revised \filedate{}.}}
% \author{William J. Henney\thanks{w.henney@irya.unam.mx}}
% \date{\filedate}
% \maketitle
% \paragraph{Abstract}
% This document describes the \texttt{rmaa} \LaTeX{}
% document class for the journal ``Revista Mexicana de
% Astronom\'\i a y Astrof\'\i sica''. The macros can be used for 
% \begin{enumerate}
% \item Preparing author manuscripts for the journal (including
% conference proceedings to be published in the ``Serie de
% Conferencias'')
% \item Typesetting an entire volume of the journal by inputting the
% individual article files from a container file. Facilities are
% included for automatic generation of TOC and Author Index. 
% \end{enumerate}
% 
% \tableofcontents
% \section{The documentation driver file}
% This is the part that generates the documentation that you may or
% may not be reading, which is done with the help of the Makefile
% discussed in section~\ref{sec:driver_make}. 
%    \begin{macrocode}
%<*driver>
\documentclass{ltxdoc}
\usepackage{hyperref}
\usepackage{palatino}
\EnableCrossrefs \CodelineIndex \RecordChanges
\setlength{\parindent}{0pt}
\def\MacroFont{\footnotesize\tt}
\begin{document} \DocInput{rmaa.dtx} \PrintIndex \PrintChanges \end{document}
%</driver>
%    \end{macrocode}

% \section{The implementation}
% 
% \changes{3.24}{2007/09/14}{Added user access to the macro version
% and date via \cs{rmVersion} and \cs{rmDate}}
% 
%    \begin{macrocode}
%<*package>
\newcommand{\rmaa@version}{3.26}
\newcommand{\rmaa@date}{2011/08/15}
\newcommand\rmVersion{\rmaa@version}
\newcommand\rmDate{\rmaa@date}
\NeedsTeXFormat{LaTeX2e}[1995/06/01]
%    \end{macrocode}
% \changes{4.0}{2019/11/02}{Change abstract to appear before
% resumen. Change RevMexAA to RMxAA for consistency with ADS.}
%
% \changes{3.22}{2007/03/22}{Allow proceedings articles to have
%   received/accepted dates} \changes{3.21}{2006/05/15}{Changes for
%   Silvia - abstracts are now shipped out to ADS.}
% \changes{3.19}{2006/01/19}{Changes for Silvia - remove adresses from
%   end, new abstract style} \changes{2.2beta}{2001/10/30}{ Fixed bug
%   whereby twocolumn option was not passed through to book
%   class. Finally sorted out the ``label on last page'' issue. }
% \changes{2.1}{2001/01/23}{ Fixed some inconsistencies in
%   \cs{nocite} and \cs{bibitem} uniqification Fixed the definition of
%   \cs{ion} so it will work in math mode and in contexts where the
%   typesize is different from normal (e.g., title, fig captions).
%   Modified the definition of \cs{@caption}, so we don't write any
%   list-of-figures data (we never use it anyway, and it just causes
%   problems with fragile commands in fig captions). }
% \changes{2.0}{2000/11/30}{ Added `newstyle' option, which should be
%   used in all new proceedings volumes (wider abstract, smaller
%   typeface in fig captions, etc.) Added commands for formatting
%   table captions. New demonstation document rmpostsamp.tex for
%   1-page posters. This will be the last release under a PRI
%   government (for a while...)}  \changes{1.9}{2000/06/06}{
%   Corrected position of page number in footer of articles that open
%   on even pages. Ensured that first article in each session opens on
%   an odd page, as do list of abstracts and author index. Added extra
%   space in TOC before the backmatter. Added command for editors,
%   \cs{BoldSessionTitles}, that makes produces bold titles for
%   session in the TOC.}  \changes{1.8}{2000/05/16}{ Fixed bug in
%   last page label with \cs{suppressfulladdresses} }
% \changes{1.7}{2000/05/11}{ Fixed more little things. Page labels
%   now work, even if you have no bibliography. You can output the
%   full addresses by hand with the \cs{outputfulladdresses}
%   command. Editors' Guide updated but still not finished. Revamped
%   splitbook.sh now produces PS file of each article for web page and
%   also an EPS file of each page for sending to the printers.}
% \changes{1.6}{2000/04/15}{ More cleaning up: full stops after
%   titles in index, running headers and footers in frontmatter,
%   fontsize in title of index, etc} \changes{1.5}{2000/04/15}{
%   Eliminated superscripts from titles in TOC and index}
% \changes{1.4}{2000/04/14}{ Added copyright and license info (GPL)}
% \changes{1.3}{2000/03/01}{ Corrected \cs{figcaption} to do the
%   caption formatting right Corrected font size of sec headings in
%   front matter} \changes{1.2}{2000/02/16}{ Added \cs{plate}
%   command} \changes{1.1}{2000/02/14}{ Added manuscript option,
%   corrected bugs} \changes{1.0}{2000/02/08}{ Now does a whole book
%   - lots of stuff tidied up} \changes{0.7}{1999/10/04}{ Increased
%   spacing between sections in single column styles. Suppressed
%   floats at bottom of pages (didn't work though) }
% \changes{0.6}{1999/09/24}{ Tweaked \cs{adjustfinalcols} (still not
%   perfect) Added proceedings option Added onecolumn/twocolumn
%   options (orthogonal to journal/proceedings) }
% \changes{0.5}{1999/09/23}{ Added \cs{fulladdresses} and
%   \cs{adjustfinalcols} (just about working)}
% \changes{0.4}{1999/09/22}{ Minor stylistic changes}
% \changes{0.3}{1999/06/28}{ Better mechanism for alternative
%   affiliations Added `guide' option for user guide}
% \changes{0.2}{1999/05/27}{ All basic functionality present, I
%   think} \changes{0.1}{1999/05/26}{ Just about working - still
%   missing Fig/Tab/Refs} \changes{0.0}{1999/05/21}{ So bored today,
%   decided to rewrite RMAA LaTeX macros from scratch}
%    \begin{macrocode}
\ProvidesClass{rmaa}[\rmaa@date\space v\rmaa@version\space RMxAA article
class]
%    \end{macrocode}
% \subsection{Option Handling}
% \label{sec:options}
% Declarations that are needed in options
% \changes{3.05}{2003/10/22}{Added natbib option}
% \changes{3.06}{2003/01/03}{Natbib option turned on permanently}
% \changes{3.07}{2003/01/11}{Added proof option}
% \changes{3.14}{2003/01/11}{Option to turn off natbib}
% \changes{3.18}{2003/10/22}{trying out doctex mode} 
% \changes{3.23}{2007/09/09}{Manuscript is now on by default.}
% \changes{3.26}{2011/08/15}{New option for review articles}
%    \begin{macrocode}
\newif\if@reviewarticle
\@reviewarticlefalse
\newif\if@proof
\@prooffalse
\newif\if@journal
\@journaltrue
\newif\if@printdates
\@printdatesfalse
\newif\if@proceedings
\@proceedingsfalse
\newif\if@manuscript
\@manuscripttrue
\newif\if@bigfont
\@bigfonttrue
\newif\if@guide
\@guidefalse
\newif\if@preprint
\@preprintfalse
\newif\if@rmtwocol
\@rmtwocoltrue
\newif\if@biblist
\@biblistfalse
\newif\if@book
\@bookfalse
\newif\if@natbib
\@natbibtrue
\newif\if@rmcite
\@rmcitefalse
\newif\if@titlevisible
\@titlevisibletrue
\newif\if@abstractpresent
\@abstractpresentfalse
\newif\if@printbanner
\@printbannerfalse
\newif\if@suppressfulladdresses
\@suppressfulladdressesfalse
\newif\if@newstyle
\@newstylefalse
\newif\if@debug
\@debugfalse
\newif\if@boldkeywords
\@boldkeywordsfalse
%    \end{macrocode}
% The option handling---mainly setting various booleans.
% \changes{2.2beta11}{2001/11/26}{Made twocol and
% newstyle the default for proceedings.}
% \changes{3.23}{2007/09/09}{Manuscript is now 10 pt instead of 12pt}
% \changes{3.23}{2007/09/10}{Full addresses are now automatically
% suppressed for proceedings style.}
% \changes{3.24}{2007/09/14}{New option \texttt{boldkeywords} that
% does the keywords in the old ugly way. Now off by default.}
%    \begin{macrocode}
\DeclareOption{manuscript}{\typeout{Manuscript for referee} 
  \@manuscripttrue\@journalfalse\@rmtwocolfalse
  \PassOptionsToClass{10pt}{article}}
\DeclareOption{bigfont}{\typeout{Big font manuscript requested by Gary
  Ferland} 
  \@manuscripttrue\@journalfalse\@rmtwocolfalse\@bigfonttrue
  \PassOptionsToClass{11pt}{article}}
\DeclareOption{journal}{\typeout{Main journal option}%
  \@journaltrue\@rmtwocoltrue\@printdatestrue\@manuscriptfalse
  \PassOptionsToClass{10pt}{article}}
\DeclareOption{guide}{\typeout{User guide option}%
  \@guidetrue\@journalfalse\@rmtwocoltrue\@manuscriptfalse
  \PassOptionsToClass{10pt}{article}}
\DeclareOption{preprint}{\typeout{Preprint option}%
  \@preprinttrue\@manuscriptfalse\PassOptionsToClass{10pt}{article}}
\DeclareOption{proof}{\typeout{Author proof option}%
  \@prooftrue\@manuscriptfalse}
\DeclareOption{proceedings}{\typeout{Proceedings option}%
  \@proceedingstrue\@journalfalse\@rmtwocoltrue\@newstyletrue
  \@suppressfulladdressestrue
  \@printdatesfalse\@manuscriptfalse\PassOptionsToClass{10pt}{article}} 
\DeclareOption{printdates}{\typeout{Force printing of received/accepted dates}%
  \@printdatestrue} 
\DeclareOption{book}{\typeout{Whole issue option}%
  \@booktrue} 
\DeclareOption{nonatbib}{\typeout{Turn off natbib support (required
    for compatibility with old rmaacite package).}%
  \@natbibfalse} 
\DeclareOption{newstyle}{\typeout{Use new layout style}%
  \@newstyletrue}
\DeclareOption{oldstyle}{\typeout{Use old layout style}%
  \@newstylefalse}
\DeclareOption{debug}{\typeout{Loads of annoying messages}%
  \@debugtrue}
%    \end{macrocode}
% These two need to be declared explicitly so that |\ExecuteOptions|
% actually does something.
%    \begin{macrocode}
\DeclareOption{twocolumn}{\@rmtwocoltrue}
\DeclareOption{onecolumn}{\@rmtwocolfalse}
\DeclareOption{twoside}{\PassOptionsToClass{\CurrentOption}{article}}
\DeclareOption{boldkeywords}{\@boldkeywordstrue}
%    \end{macrocode}
% Bold, upper-case keywords used to be the default. Now they must be
% turned on explicitly with the \texttt{boldkeywords} option. 
%    \begin{macrocode}
\DeclareOption{boldkeywords}{\@boldkeywordstrue}
%    \end{macrocode}
% Pass any other option along to the \texttt{article} class.
%    \begin{macrocode}
\DeclareOption*{\PassOptionsToClass{\CurrentOption}{article}}
%    \end{macrocode}
% Default options are to be two-sided and based on the standard
% \LaTeX{} \texttt{article} class,
% unless we are processing the entire book, in which case we base
% ourselves on the standard \texttt{book} class. Default is to use
% manuscript option unless otherwise specified. 
%    \begin{macrocode}
\ExecuteOptions{twoside,manuscript} \ProcessOptions\relax
\if@book
\if@rmtwocol\LoadClass[twocolumn]{book}\else\LoadClass{book}\fi
\newcommand{\abstract}{\relax}
\newcommand{\@maketitle}{\relax}
\else
\if@rmtwocol\LoadClass[twocolumn]{article}\else\LoadClass{article}\fi
\fi
%    \end{macrocode}
% Attempt to use natbib for citations 02 Jan 2003. 
% \changes{3.05}{2003/01/02}{If the natbib option is used, then set
% things up so that natbib's version of the bibliography is identical
% to the standard one.}
% \changes{3.09}{2003/01/23}{Explicitly require a recent version of
% natbib, so that we get a decent error message if the package is too
% old.} 
% \changes{3.23}{2007/09/09}{Now allow Bibtex to work properly} 
% \changes{3.25}{2009/10/01}{Change \cs{newcommand} tp \cs{def} when setting \cs{bibfont}. This allows it to work with all versions of natbib.} 
% \changes{3.25}{2010/02/06}{Updated copyright year.} 
% \changes{3.25}{2010/02/06}{Updated email and web addresses to \texttt{crya.unam.mx}.} 
%    \begin{macrocode}
\if@natbib
\RequirePackage{natbib}[1999/02/25]
\bibpunct{(}{)}{;}{a}{}{,}
\def\bibfont{\small}
\renewcommand\bibsection{\section*{REFERENCES}}
\setlength\bibsep{0pt}
\setlength\bibhang{1.5em}
\bibliographystyle{rmaa}
\fi
%    \end{macrocode}
% We want the \texttt{graphicx} package for the extended syntax of
% |\includegraphics|. We want the \texttt{ifthen} package for using in
% these macros (although we don't much yet). We explicitly load
% \texttt{babel} since some \LaTeX{} distributions messed up the
% hyphenation if we didn't.  
%    \begin{macrocode}
\RequirePackage{graphicx}
\RequirePackage{ifthen}
\IfFileExists{babel.sty}{%
%   \let\@restoreindent=\@afterindentfalse
  \RequirePackage[english]{babel}
%   \@restoreindent
  %% reinstate the definition of \@roman so as not to break .idx files
  \@ifundefined{org@@roman}{\relax}{\renewcommand{\@roman}[1]{\org@@roman{##1}}}
  }{%
  \typeout{Package babel.sty not found - don't worry about this}}
%    \end{macrocode}
% Try to load the \texttt{dcolumn} and \texttt{booktabs} packages if
% they can be found. These can produce nicer-looking tables. 
%    \begin{macrocode}
\IfFileExists{dcolumn.sty}{%
  \RequirePackage{dcolumn}
  }{%
  \typeout{Package dcolumn.sty not found - you will not be able to
    align columns on decimal points in tables}}
\IfFileExists{booktabs.sty}{%
  \RequirePackage{booktabs}
  }{%
  \typeout{Package booktabs.sty not found - your tables will not look
    as pretty as they ought too.}
  }
\IfFileExists{makeidx.sty}{%
  \RequirePackage{makeidx}
}{%
  \typeout{Package makeidx.sty not found - Author Index may not work.}
  }
%% commented out for moment since requires more modern version of LaTeX
%% \RequirePackage[latin1]{inputenc}
%    \end{macrocode}
% \changes{4.0}{2019/11/13}{Make \texttt{hyperref} and
% \texttt{etoolbox} into hard requirements, since we now use them
% extensively.}  We use hyperref to get internal and external links.
% But I don't like the colored boxes so I turn those off.  We use
% etoolbox for patching and robustifying at the moment.  In the
% future, I want to port a lot of other stuff over to it.  For
% instance, replace the ifthen package and the TeX ifs too.
%
% We target the latest version (1.7) of PDF, rather than the default
% of 1.5.  This reduces the chance of getting warnings from included
% graphics.  This is not strictly related to hyperref, but we put it
% here anyway.
%    \begin{macrocode}
\RequirePackage{etoolbox}
\RequirePackage{hyperref}
\hypersetup{hidelinks}
\pdfminorversion=7
%    \end{macrocode}
% We patch the \cs{@combinedblfloats} following a suggestion by Heiko
% Oberdiek at \url{https://tex.stackexchange.com/a/249743/3843}, which
% avoids problems when hyperlinks are split across
% pages. \changes{4.0}{2019/11/12}{Fix for issue where breaking
% hyperlink across pages gave an error.}
%    \begin{macrocode}
\urlstyle{same}
\newcommand\email{\begingroup \urlstyle{rm}\Url}
\patchcmd\@combinedblfloats{\box\@outputbox}{\unvbox\@outputbox}{}{%
  \errmessage{\noexpand\@combinedblfloats could not be patched}%
}%
%    \end{macrocode}
% \subsection{Other setup stuff}\label{sec:setup}
% Some declarations: 
%    \begin{macrocode}
\newif\if@needaddresses
\@needaddressestrue
\newif\if@frontmatter
\@frontmatterfalse
\newif\if@inpreamble
\@inpreambletrue
\newif\if@firstauthor
\@firstauthortrue
\newif\if@firstaffil
\@firstaffiltrue
\newif\if@firstaltaffil
\@firstaltaffiltrue
\newif\if@firstkwd
\@firstkwdtrue
\newif\if@noexpandADS
\@noexpandADSfalse
\newif\if@sessiontitlepages
\@sessiontitlepagesfalse
%    \end{macrocode}
% \begin{macro}{\@rmaaEmpty}
% \begin{macro}{\@rmaaArg}
% These are used in various commands to check if arg was left blank
%    \begin{macrocode}
\newcommand{\@rmaaEmpty}{}
\newcommand{\@rmaaArg}{}
%    \end{macrocode}  
% \end{macro}
% \end{macro}
%
% \subsection{Compatibility with old \LaTeX{} versions}
% \label{sec:compat}
% This package is based on the 1998/12/01 release of \LaTeX{}.
% Older releases may not define all the commands that we need.
% Luckily, we can use |\providecommand| to define these if necessary. 
% The following commands are not defined in the 1994/06/01 release. 
%
% \begin{macro}{\MakeUpperCase}
% \begin{macro}{\textsuperscript}
% \begin{macro}{\textbackslash}
% User commands:
%    \begin{macrocode}
\providecommand{\MakeUpperCase}[1]{\uppercase{#1}}
\providecommand{\textsuperscript}[1]{\ensuremath{^\textrm{#1}}}
\providecommand{\textbackslash}{\ensuremath{\backslash}}
%    \end{macrocode}  
% \end{macro}
% \end{macro}
% \end{macro}
% \begin{macro}{\protected@xdef}
% \begin{macro}{\restore@protect}
% \begin{macro}{\@unexpandable@protect}
% \begin{macro}{\hb@xt@}
% \LaTeX{} internal commands:
%    \begin{macrocode}
\providecommand{\protected@xdef}{%
        \let\@@protect\protect
        \let\protect\@unexpandable@protect
        \afterassignment\restore@protect
        \xdef}
\providecommand{\restore@protect}{\let\protect\@@protect}
\providecommand{\@unexpandable@protect}{\noexpand\protect\noexpand}
\providecommand{\hb@xt@}{\hbox to}
%    \end{macrocode}
% \end{macro}
% \end{macro}
% \end{macro}
% \end{macro}
%
% \subsection{Page layout}\label{sec:page-layout}
%
% The dimensions of the text were defined by Elizabeth. Text width is
% 169\,mm with equal margins on the two sides: 
%    \begin{macrocode}
\setlength{\textwidth}{169mm}
\setlength{\@tempdima}{\paperwidth}
\addtolength{\@tempdima}{-\textwidth}
\setlength{\oddsidemargin}{0.5\@tempdima}
\addtolength{\oddsidemargin}{-1in}
\setlength{\evensidemargin}{\oddsidemargin}
%    \end{macrocode}
% Text height is 225\,mm, top margin is 14\,mm. 
%    \begin{macrocode}
\setlength{\textheight}{225mm} 
\setlength{\topmargin}{-1in} 
\addtolength{\topmargin}{14mm}
%    \end{macrocode}
% Space between columns in two-column styles, header, footer,
% paragraph indentation:
%    \begin{macrocode}
\setlength{\columnsep}{9.5mm}
\setlength{\headsep}{0.7cm}
\setlength{\footskip}{0.75cm}
\setlength{\parindent}{1.5em}
%    \end{macrocode}
% 
% \begin{macro}{\@set@scalable@length}
% Try out a new scheme for dealing with various spacing dimensions and
% allowing the user to easily modify them. The command
% \begin{quote}
% |\set@scalable@length|\marg{name}\marg{defaultlength}\marg{stretch}\marg{shrink}
% \end{quote}
% does a 
% |\setlength|\marg{\cs{name}}\marg{defaultlength plus stretch minus shrink}
% and also defines the user command
% \textbackslash\meta{name}|Scale|\marg{factor} that can 
% be used to rescale length by \meta{factor}. Additionally, now
% defines the package command
% \textbackslash\meta{name}|@set@scale|\marg{FACTOR} 
% that also rescales the length by \meta{FACTOR}, so that the length
% becomes \meta{factor} $\times$ \meta{FACTOR} $\times$
% \meta{defaultlength}.   
%    \begin{macrocode}
\newcommand\@set@scalable@length[4]{%
  %% arg #1: NAME of length without leading `\'
  %%
  %% arg #2: default VALUE  
  \expandafter\newlength\csname @#1@default\endcsname
  \expandafter\setlength\csname @#1@default\endcsname{#2}%
  \expandafter\newlength\csname @@#1@@default\endcsname
  \expandafter\setlength\csname @@#1@@default\endcsname{#2}%
  %% arg #3: default STRETCH  
  \expandafter\newlength\csname @#1@stretch\endcsname
  \expandafter\setlength\csname @#1@stretch\endcsname{#3}%
  \expandafter\newlength\csname @@#1@@stretch\endcsname
  \expandafter\setlength\csname @@#1@@stretch\endcsname{#3}%
  %% arg #4: default SHRINK  
  \expandafter\newlength\csname @#1@shrink\endcsname
  \expandafter\setlength\csname @#1@shrink\endcsname{#4}%
  \expandafter\newlength\csname @@#1@@shrink\endcsname
  \expandafter\setlength\csname @@#1@@shrink\endcsname{#4}%
  %%
  %% User-defined scale \@NAME@scale is unity by default
  \expandafter\gdef\csname @#1@scale\endcsname{1}%
  %% Package-defined scale \@NAME@SCALE is unity by default
  \expandafter\gdef\csname @#1@SCALE\endcsname{1}%
  %%
  %% internal command \@reset@NAME should expand to 
  %% \setlength\@@NAME@@default{\@NAME@SCALE\@NAME@default}
  %% \setlength\NAME{\@NAME@scale\@@NAME@@default}
  %% in other words: scale x SCALE x default 
  \expandafter\gdef\csname @reset@#1\endcsname{%
    \expandafter\setlength\csname @@#1@@default\endcsname{%
      \csname @#1@SCALE\endcsname \csname @#1@default\endcsname
      }
    \expandafter\setlength\csname @@#1@@stretch\endcsname{%
      \csname @#1@SCALE\endcsname \csname @#1@stretch\endcsname
      }
    \expandafter\setlength\csname @@#1@@shrink\endcsname{%
      \csname @#1@SCALE\endcsname \csname @#1@shrink\endcsname
      }
    \expandafter\setlength\csname #1\endcsname{%
      \csname @#1@scale\endcsname \csname @@#1@@default\endcsname
      \@plus
      \csname @#1@scale\endcsname \csname @@#1@@stretch\endcsname
      \@minus
      \csname @#1@scale\endcsname \csname @@#1@@shrink\endcsname
      }}
  %%
  %% User command \NAMEScale should expand to 
  %% {\renewcommand\@NAME@scale{#1}\@reset@NAME}
  \expandafter\gdef\csname #1Scale\endcsname##1{%
    \expandafter\gdef\csname @#1@scale\endcsname{##1}%
    \csname @reset@#1\endcsname}
  %%
  %% Package command \NAME@set@SCALE should expand to 
  %% {\renewcommand\@NAME@SCALE{#1}\@reset@NAME}
  \expandafter\gdef\csname #1@set@SCALE\endcsname##1{%
    \expandafter\gdef\csname @#1@SCALE\endcsname{##1}%
    \csname @reset@#1\endcsname}
  %%
  %% initialize at default values with \@reset@NAME
  \csname @reset@#1\endcsname
  }
%    \end{macrocode}
% \end{macro}
% \begin{macro}{\parskip}
% First use of this new infrastructure: set interparagraph space to be
% zero but with a small amount of stretchiness.
%    \begin{macrocode}
%   \ifthenelse{\boolean{@manuscript}}{%
%     \@set@scalable@length{parskip}{0.5\baselineskip}{0.3\p@}{0.1\p@}
%   }{
    \@set@scalable@length{parskip}{0\p@}{0.3\p@}{0\p@}
%   }
%    \end{macrocode}
% \end{macro}
% 
% \begin{macro}{\abovesectionskip}
% \changes{4.0}{2019/11/12}{Use negative values for skip so that first paragraph after section headings is not indented.  This is something that I have been trying to fix for years!  It turns out that the anser was in the documentation for the standard command \cs{@startsection.}}
% \begin{macro}{\belowsectionskip}
% \begin{macro}{\aboveparaskip}
% \begin{macro}{\belowparaskip}
% Define some utility rubber lengths for use in
% section~\ref{sec:section-commands}. 
%    \begin{macrocode}
\newlength{\abovesectionskip}
\newlength{\belowsectionskip}
\newlength{\aboveparaskip}
\newlength{\belowparaskip}
%    \end{macrocode}
% The default values for these lengths are those traditionally used
% for the onecolumn style. 
%    \begin{macrocode}
\@set@scalable@length{abovesectionskip}{-3.5ex}{-4ex}{-2ex}
\@set@scalable@length{belowsectionskip}{2.3ex}{0ex}{2ex}
\@set@scalable@length{aboveparaskip}{3.25ex}{2ex}{0.5ex}
\@set@scalable@length{belowparaskip}{-1em}{\z@}{\z@}
%    \end{macrocode}
% \end{macro}\end{macro}\end{macro}\end{macro}
% We also make scalable some of the standard float parameters. These
% are the same defaults as in the standard article classes, but
% slightly more stretchy. For the twocolumn style, they are currently
% halved, see below. Note that I have also lumped |\abovecaptionskip|
% in with these, since it plays a similar role. There is also a
% |\belowcaptionskip| but it is set to zero in the standard classes. 
%    \begin{macrocode}
\@set@scalable@length{floatsep}{12\p@}{4\p@}{4\p@}
\@set@scalable@length{textfloatsep}{20\p@}{4\p@}{8\p@}
\@set@scalable@length{intextsep}{12\p@}{4\p@}{4\p@}
\@set@scalable@length{dblfloatsep}{12\p@}{4\p@}{4\p@}
\@set@scalable@length{dbltextfloatsep}{20\p@}{4\p@}{8\p@}
\@set@scalable@length{abovecaptionskip}{10\p@}{\z@}{\z@}
%    \end{macrocode}
% \begin{macro}{\@before@title@skip}         
% \begin{macro}{\@before@author@skip}        
% \begin{macro}{\@between@author@skip}       
% \begin{macro}{\@before@abstract@skip}      
% \begin{macro}{\@after@abstracthead@skip}   
% \begin{macro}{\@between@abstract@skip}     
% \begin{macro}{\@before@kwd@skip}           
% \begin{macro}{\@before@dates@skip}           
% Spacing parameters used in |\@maketitle|. These are not changed by
% default except by the user, so they don't need to go into
% |\@package@default@lengths|.  
%    \begin{macrocode}
\newlength{\@before@title@skip}         
\newlength{\@before@author@skip}        
\newlength{\@between@author@skip}       
\newlength{\@before@abstract@skip}      
\newlength{\@after@abstracthead@skip}   
\newlength{\@between@abstract@skip}     
\newlength{\@before@kwd@skip}           
\newlength{\@before@dates@skip}           
\@set@scalable@length{@before@title@skip}{2em}{\z@}{\z@}
\@set@scalable@length{@before@author@skip}{1.5em}{\z@}{\z@}
\@set@scalable@length{@between@author@skip}{0.5em}{\z@}{\z@}
\@set@scalable@length{@before@abstract@skip}{3ex}{1ex}{0.5ex}
\@set@scalable@length{@after@abstracthead@skip}{1ex}{\z@}{\z@}
\@set@scalable@length{@between@abstract@skip}{2ex}{1ex}{0.5ex}
\@set@scalable@length{@before@kwd@skip}{1ex}{0.5ex}{0.2ex}
\@set@scalable@length{@before@dates@skip}{2ex}{0.5ex}{0.2ex}
%    \end{macrocode}
% \end{macro}\end{macro}\end{macro}\end{macro}
% \end{macro}\end{macro}\end{macro}\end{macro}
% 
% Spacing above and below displayed equations also wants to be done
% like this, but it is more complicated since the standard article
% classes make these change with the fontsize. Something to do later,
% maybe. 
% 
%    \begin{macrocode}
%    \end{macrocode}
% \begin{macro}{\@package@default@lengths}
% The value actually used for these lengths is controlled by the
% \meta{SCALE} parameter, which depends on the style option that is in
% effect. This is set using the |\@package@default@lengths| command,
% which must be called every time the style changes (e.g., changing
% from onecolumn to twocolumn). 
%    \begin{macrocode}
\newcommand{\@package@default@lengths}{%
  \ifthenelse{\boolean{@rmtwocol}}{%
    %% twocolumn mode
    \abovesectionskip@set@SCALE{0.5}
    \belowsectionskip@set@SCALE{0.5}
    \aboveparaskip@set@SCALE{0.5}
    \belowparaskip@set@SCALE{1}
    \floatsep@set@SCALE{0.5}
    \textfloatsep@set@SCALE{0.5}
    \intextsep@set@SCALE{0.5}
    \dblfloatsep@set@SCALE{0.5}
    \dbltextfloatsep@set@SCALE{0.5}
    }{%
    %% onecolumn mode
    \abovesectionskip@set@SCALE{1}
    \belowsectionskip@set@SCALE{1}
    \aboveparaskip@set@SCALE{1}
    \belowparaskip@set@SCALE{1}
    \floatsep@set@SCALE{1}
    \textfloatsep@set@SCALE{1}
    \intextsep@set@SCALE{1}
    \dblfloatsep@set@SCALE{1}
    \dbltextfloatsep@set@SCALE{1}
    }%
  }
\@package@default@lengths
%    \end{macrocode}
% \end{macro}
% \begin{macro}{\RescaleLengths}
% \begin{macro}{\RescaleSecLengths}
% \begin{macro}{\RescaleFloatLengths}
% \begin{macro}{\RescaleTitleLengths}
% User commands that will fiddle with all the scalable lengths. 
%    \begin{macrocode}
\newcommand{\RescaleLengths}[1]{%
  \RescaleSecLengths{#1}\RescaleFloatLengths{#1}\RescaleTitleLengths{#1}}
\newcommand{\RescaleSecLengths}[1]{%
%% sections
  \abovesectionskipScale{#1}
  \belowsectionskipScale{#1}
  \aboveparaskipScale{#1}
  }
\newcommand{\RescaleFloatLengths}[1]{%
%% floats
  \floatsepScale{#1}
  \textfloatsepScale{#1}
  \intextsepScale{#1}
  \dblfloatsepScale{#1}
  \dbltextfloatsepScale{#1}
  \abovecaptionskipScale{#1}
  }
\newcommand{\RescaleTitleLengths}[1]{%
%% maketitle
  \@before@title@skipScale{#1}           
  \@before@author@skipScale{#1}        
  \@between@author@skipScale{#1}       
  \@before@abstract@skipScale{#1}      
  \@after@abstracthead@skipScale{#1}   
  \@between@abstract@skipScale{#1}     
  \@before@kwd@skipScale{#1}           
  \@before@dates@skipScale{#1}           
  }
%    \end{macrocode}
% \end{macro}\end{macro}\end{macro}\end{macro}
% \changes{2.2beta14}{2002/03/17}{Finished implementation of new
% infrastructure for scalable lengths. Converted the above/below
% section skips to use it. Also, all the floatsep parameters. Also,
% all the vertical lengths that are used by |\@maketitle|. }

% \begin{macro}{\MSextramargin}
% For the \texttt{manuscript} option we have an extra 0.3\,inch
% vertical margin and 1\,inch horizontal margin for marginal notes (by
% hand or with |\marginpar|. 
%
% \changes{3.23}{2007/09/09}{Reduce baselinestretch in
%   \texttt{manuscript} style to 1.2. No longer use raggedright for
%   manuscript.}
%
% \changes{3.24}{2007/09/14}{Further reduction of baselinestretch in
%   \texttt{manuscript} style to 1.1 since we now have wider margins.}
%
% \changes{3.24}{2007/09/14}{Split \texttt{manuscript}'s extra margin
%   into separate horizontal and vertical parts.}
% 
%    \begin{macrocode}
\newlength{\MSextraHmargin}\newlength{\MSextraVmargin}
\setlength{\MSextraHmargin}{1in} 
\setlength{\MSextraVmargin}{0.3in} 
\if@manuscript
   \renewcommand{\baselinestretch}{1.1}
   \addtolength{\textwidth}{-2\MSextraHmargin}
   \addtolength{\oddsidemargin}{\MSextraHmargin}
   \setlength{\evensidemargin}{\oddsidemargin}
   \addtolength{\textheight}{-2\MSextraVmargin}
   \addtolength{\topmargin}{\MSextraVmargin}
%    \raggedright
\fi
%    \end{macrocode}
% \end{macro}
% Float parameters: don't allow bottom floats by default. 
%    \begin{macrocode}
\renewcommand{\fps@figure}{tp}
\renewcommand{\fps@table}{tp}
%    \end{macrocode}
%
% \begin{macro}{\rmaa@AuthorWidthFrac}
% \begin{macro}{\rmaa@AbstractWidthFrac}
% These parameters control what fraction of the total text width can
% be used by the list of authors and the abstract/resumen,
% respectively.  
%    \begin{macrocode}
\newcommand{\rmaa@AuthorWidthFrac}{1.0}
\newcommand{\rmaa@AbstractWidthFrac}{0.75}
%    \end{macrocode}
% Use a wide abstract to save space in the new style.
%    \begin{macrocode}
\if@newstyle\renewcommand{\rmaa@AbstractWidthFrac}{1.0}\fi
%    \end{macrocode}
% \end{macro}
% \end{macro}
%
% Some things to consider in future versions. Trying to avoid too many
% widows, etc. Apparently, the following will work best with
% \cs{flushbottom}. 
%    \begin{macrocode}
%%\widowpenalty=10000
%%\clubpenalty=10000
%    \end{macrocode}
% Also, should perhaps try getting rid of underfull/overfull box warnings. 
% \changes{3.10}{2003/01/24}{Set large values for \cs{hbadness},
% \cs{vbadness}, \cs{hfuzz}, \cs{vfuzz}. To reduce amount of stuff
% that gets written to the terminal every time you compile something.}
%    \begin{macrocode}
\hbadness=10000
\vbadness=10000
\setlength\hfuzz{\maxdimen}
\setlength\vfuzz{\maxdimen}
%    \end{macrocode}
% \subsection{Various text styles for different contexts}
% \label{sec:text-styles} 
% \begin{macro}{\rmaa@AffilStyle}
% \begin{macro}{\rmaa@PagesStyle}
% \begin{macro}{\rmaa@YearStyle}
% \begin{macro}{\rmaa@VolumeStyle}
% \begin{macro}{\rmaa@VolumenStyle}
% \begin{macro}{\rmaa@JournalNameStyle}
% \begin{macro}{\rmaa@ReceivedStyle}
% \begin{macro}{\rmaa@AbstractHeaderStyle}
% \begin{macro}{\rmaa@JournalNameStyle}
% Style commands for the different fields used by |\maketitle| in
% constructing the title page of the article. These have the form
% |\rmaa@XYZStyle{TEXT}|. I'm not sure these were really a good idea,
% but never mind \dots 
%    \begin{macrocode}
\newcommand{\rmaa@AffilStyle}[1]{{\normalsize\rmfamily#1}}
\newcommand{\rmaa@PagesStyle}[1]{{\normalsize\rmfamily#1}}
\newcommand{\rmaa@YearStyle}[1]{{\normalsize\rmfamily#1}}
\newcommand{\rmaa@VolumeStyle}[1]{{\normalsize\rmfamily\bfseries#1}}
\newcommand{\rmaa@VolumenStyle}[1]{{\Large\sffamily\bfseries#1}}
\newcommand{\rmaa@JournalNameStyle}[1]{{\normalsize\itshape#1}}
\newcommand{\rmaa@ReceivedStyle}[1]{{\normalsize\rmfamily\itshape#1}}
\newcommand{\rmaa@AbstractHeaderStyle}[1]{%
  {\normalsize\rmfamily\MakeUpperCase{#1}}}
\newcommand{\rmaa@TitleStyle}[1]{%
  {\large\bfseries\rmfamily\boldmath\MakeUpperCase{#1}\par}}
%    \end{macrocode}
% \end{macro}
% \end{macro}
% \end{macro}
% \end{macro}
% \end{macro}
% \end{macro}
% \end{macro}
% \end{macro}
% \end{macro}
%
% \begin{macro}{\rmaa@FrontSecStyle}
% Headings of ``front sections'' (see section~\ref{sec:front-sections})
%    \begin{macrocode}
\newcommand{\rmaa@FrontSecStyle}[1]{%
  {\normalsize\bfseries\rmfamily\MakeUpperCase{#1}}}
%    \end{macrocode}
% \end{macro}
%
% \begin{macro}{\rmaa@TOCTitleStyle}
% \begin{macro}{\rmaa@TOCPartStyle}
% \begin{macro}{\rmaa@IndexTitleStyle}
% \begin{macro}{\rmaaTOCAuthorStyle}
% \begin{macro}{\rmaaIndexAuthorStyle}
% Commands to set style of TOC and index entries.
% Rewritten to avoid unnecessary font commands in the .toc and .ind files
%    \begin{macrocode}
\newcommand{\rmaa@TOCTitleStyle}[1]{\MakeUpperCase{#1}}
\newcommand{\rmaa@TOCPartStyle}[1]{\rmaa@TOCTitleStyle{#1}}
\newcommand{\rmaa@IndexTitleStyle}[1]{#1}
\newcommand{\rmaaTOCAuthorStyle}[1]{%
  {\itshape#1}}
\newcommand{\rmaaIndexAuthorStyle}[1]{%
  \protect\mbox{\bfseries#1}}
%    \end{macrocode}
% \end{macro}
% \end{macro}
% \end{macro}
% \end{macro}
% \end{macro}
%
% \begin{macro}{\rmaa@KeywordStyle}
% \begin{macro}{\rmaa@KeywordHeaderStyle}
% The style of the keywords depends on whether we are using the
% \texttt{newstyle} option or not. If we are, then they are small. 
%    \begin{macrocode}
\ifthenelse{\boolean{@newstyle}}{
  \newcommand{\rmaa@KeywordStyle}[1]{%
    {\small%
      \if@boldkeywords 
      \bfseries\MakeUpperCase{\raggedright #1}
      \else
      \raggedright #1
      \fi
    }}
  \newcommand{\rmaa@KeywordHeaderStyle}[1]{%
    {\small\itshape#1}}}
{
  \newcommand{\rmaa@KeywordStyle}[1]{%
    {\normalsize%
      \if@boldkeywords\bfseries\fi\raggedright #1}}
  \newcommand{\rmaa@KeywordHeaderStyle}[1]{%
    {\normalsize\itshape#1}}}
%    \end{macrocode}
% \end{macro}
% \end{macro}
%
% \subsection{Page styles}\label{sec:page-styles}
% \begin{macro}{\ps@rmaa}
% The running headers for normal pages:
%    \begin{macrocode}
\newcommand{\ps@rmaa}{%
%    \end{macrocode}
% nothing in either foot; 
%    \begin{macrocode}
  \renewcommand\@oddfoot{\relax}%
  \renewcommand\@evenfoot{\relax}%
%    \end{macrocode}
% page number and short author list in the even heads; 
%    \begin{macrocode}
  \renewcommand\@evenhead{\rlap{\thepage}\hfil\@shortauthor\hfil}%
%    \end{macrocode}
% short title and page number in the odd heads.
%    \begin{macrocode}
  \renewcommand\@oddhead{\hfil\@shorttitle\hfil\llap{\thepage}}
  }
%    \end{macrocode}
% \begin{macro}{\ps@rmaafirstpage}
% The header of the first page of each
% article:
%    \begin{macrocode}
\newcommand{\ps@rmaafirstpage}{%
  \renewcommand\@oddfoot{\hfil\thepage}%
%    \end{macrocode}
% footer corrected (06 Jun 2000) to have page on outside edge;
%    \begin{macrocode}
  \renewcommand\@evenfoot{\thepage\hfil}%
%    \end{macrocode}
% Header has journal name, volume, and page numbers for a regular
% article, or similar for special cases. 
% \changes{4.0}{2019/11/03}{Change to three-line header with copyright info and DOI.}
%    \begin{macrocode}
\if@preprint
\renewcommand\@oddhead{%
  \if@proceedings 
  \textsc{To appear in ``\@conftitle{} (\rmaa@YearStyle{\@year})''}\hfil
  \textit{RMxAC}
  \else
  \textsc{To appear in} \@journalname\hfil
  \fi}\fi%
\if@guide
\renewcommand\@oddhead{\textsc{Documentation} (\rmaa@YearStyle{\@year})\hfil}
\fi%
\if@manuscript
\renewcommand\@oddhead{Manuscript for \@journalname{} (\rmaa@YearStyle{\@year})\hfil}
\fi%
\if@journal
\renewcommand\@oddhead{%
  \parbox{\textwidth}{
    \@journalname, \rmaa@VolumeStyle{\@volume},
    \rmaa@PagesStyle{\@artpages}\
    (\rmaa@YearStyle{\@year})\hfil\\[1pt]
    {\small\textcopyright{} \@year:
      Instituto de Astronom\'\i{}a,
      Universidad Nacional Aut\'onoma de M\'exico}\hfil\\[1pt]
    {\small\@DOI@url}\hfil
  }%
}\fi%
  \renewcommand\@evenhead{\@oddhead}}
%    \end{macrocode}
% By default, use the |\ps@rmaa| page style.
%    \begin{macrocode}
\pagestyle{rmaa}
%    \end{macrocode}
% \end{macro}
% \end{macro}
%
% \subsection{Preamble commands}\label{sec:preamble} 
% Everything you need for setting the title page, TOC and index
% entries, etc.
% \subsubsection{Commands for the title page}
% \label{sec:title-commands}
% \begin{macro}{\resumen}
% \begin{macro}{\@resumen}
% \begin{macro}{\abstract}
% \begin{macro}{\@abstract}
% User commands for setting the resumen and abstract. 
% These should be used in the document header. The internal commands
% |\@resumen| and |\@abstract| are used to hold the actual text
% pending output via the |\maketitle| command. 
%    \begin{macrocode}
\newif\if@setresumen
\newcommand{\resumen}[1]{%
%%  \renewcommand{\@resumen}{{\selectlanguage{spanish}#1}}%
  \renewcommand{\@resumen}{{\ignorespaces#1}}%
  \@setresumentrue}
\renewcommand{\abstract}[1]{%
  \renewcommand{\@abstract}{{\ignorespaces#1}}
  \renewcommand\ADS@abstract{#1}
%% Not a good idea, since then author macros are not expanded and
%% cause errors later. However, by not doing it we are vulnerable to
%% fragile commands.....
  \if@noexpandADS\@onelevel@sanitize{\ADS@abstract}\fi
  \@abstractpresenttrue
  }
\newcommand{\ADSnoexpand}{\@noexpandADStrue}
%    \end{macrocode}
% \begin{macro}{\boldabstract}
% This is the version for doing 1-page conference proceedings. Still
% save the text for ADS but then write it out immediately as a bold
% paragraph. 
%    \begin{macrocode}
\newcommand{\boldabstract}[1]{%
  \gdef\ADS@abstract{#1}%
  \renewcommand{\@abstract}{#1}%
  \if@noexpandADS\@onelevel@sanitize{\ADS@abstract}\fi%
  \noindent\textbf{\boldmath #1}
  }
%    \end{macrocode}
% \end{macro}
% Default value in each case is to print out a warning if not set
%    \begin{macrocode}
  \newcommand{\@resumen}{%
    \@setresumenfalse \ClassWarningNoLine{rmaa}{No Spanish abstract
      (resumen) specified - if you do not speak Spanish, this will be
      translated by the editors}%
    \textbf{Favor de proporcionar un resumen en espa\~nol. If you are
      unable to translate your abstract into Spanish, the editors will
      do it for you.} }
  \newcommand{\@abstract}{\ClassWarningNoLine{rmaa}{No abstract
      specified}}
%    \end{macrocode}
% \end{macro}
% \end{macro}
% \end{macro}
% \end{macro}
%
% \begin{macro}{\title}
% \begin{macro}{\@rawtitle}
% \begin{macro}{\@TOCtitle}
% \begin{macro}{\rmaaIndexTitle}
% Similarly, |\title| is the preamble command to set the title of the
% article. This is used in various places, typeset in different
% styles, hence the use of the internal |\@rawtitle|, |\@title|,
% |\rmaaIndexTitle|, and |\@TOCtitle| commands, each of which uses the
% relevant style from section~\ref{sec:text-styles}
%    \begin{macrocode}
\newcommand{\@rawtitle}{Raw Title}
\newcommand{\@TOCtitle}{\rmaa@TOCTitleStyle{Title for TOC}}
\newcommand{\rmaaIndexTitle}{\rmaa@IndexTitleStyle{Title for index}.}
\newcommand{\@TOCaftertitle}{}
\renewcommand{\title}[1]{%
%%  {\ADS@sanitize\xdef\@rawtitle{#1}}
  \renewcommand\@rawtitle{\noexpand#1}
  \let\ADS@title=\@rawtitle
  \renewcommand{\@title}{\rmaa@TitleStyle{#1}}%
  \renewcommand{\@TOCtitle}{\rmaa@TOCTitleStyle{#1}\@TOCaftertitle}%
  \renewcommand{\rmaaIndexTitle}{\rmaa@IndexTitleStyle{#1\@TOCaftertitle}.}%
  }
%    \end{macrocode}
% \end{macro}
% \end{macro}
% \end{macro}
% \end{macro}
% 
% All the weird and whacky ways of indicating the author
% affiliations. 
% \begin{macro}{\affil}
% First, the way that sets it directly underneath the authors name.
% Note that |\affil| trumps all other affiliations for the purposes of
% getting into the ADS stuff (we use |\gdef|).   
% \dots 
%    \begin{macrocode}
\newcommand{\affil}[1]{%
  \\[1.5ex]\rmaa@AffilStyle{#1}\\[1.5ex]%
  \if@firstaffil\gdef\ADS@affil{#1}\@firstaffilfalse\fi%
  }
\renewcommand{\and}{and\\[1.5ex]}
%    \end{macrocode}
% \begin{macro}{\altaffil}
% Alternative affiliations are footnoted. 
%    \begin{macrocode}
\newcommand{\altaffil}[1]{%
  \if@firstaltaffil\gdef\ADS@affil{#1}\@firstaltaffilfalse\fi%
  \thanks{#1}}
%    \end{macrocode}
% \begin{macro}{\altaffilmark}
% \begin{macro}{\altaffiltext}
% Alternative-alternative mechanism for affiliations---should not be
% mixed with |\altaffil| (unless you know what you are doing) or there
% will be problems with numbering.
%    \begin{macrocode}
\newcommand{\altaffilmark}[1]{\textsuperscript{#1}}
\newcommand{\altaffiltext}[2]{%
  \if@firstaltaffil\gdef\ADS@affil{#2}\@firstaltaffilfalse\fi
  \stepcounter{footnote}%
  \protected@xdef\@thanks{\@thanks
    \protect\footnotetext[#1]{#2}}} 
%    \end{macrocode}
% \end{macro}
% \end{macro}
% \end{macro}
% \end{macro}
%
% \begin{macro}{\keywords}
% \begin{macro}{\@keywords}
% \begin{macro}{\@keywordheader}
% \label{keywords}Sundry minor preamble commands. All
% self-explanatory, I hope. |\keywords| is the old user interface
% (deprecated). New articles should use |\addkeyword| once for each
% keyword. That way the ADS record comes out right too (comma rather
% than em-dash separator).
% \changes{2.2beta12}{2001/12/07}{\cs{addkeyword} New user interface
%   for keywords.} 
% \changes{2.2beta13}{2002/03/11}{\cs{addkeyword} Improved behaviour
%   with .ads file. Use
%   \cs{ADS@sanitize} to change definitions of `~' and \cs{ion}.}
% \changes{3.23}{2007/09/09}{\cs{addkeyword} now signals a compilation
% error if the author tries to add more than 6 keywords.}
% \changes{3.24}{2007/09/14}{kwdcounter is now reset in
% \cs{@begofarticlehook} to prevent error in book style} 
%    \begin{macrocode}
\newcommand{\@kwdsep}{ --- }
\newcommand{\ADS@kwdsep}{, }
\newcounter{kwdcounter}\setcounter{kwdcounter}{0}
\newcommand{\addkeyword}[1]{%
  \addtocounter{kwdcounter}{1}
  \ifthenelse{\thekwdcounter > 6}{
    \ClassError{rmaa}{Too many keywords. A maximum of 6 keywords is allowed.}{}
  }{}
  \if@firstkwd%
  \renewcommand\@keywords{\if@boldkeywords\uppercase{#1}\else #1\fi}
  {\ADS@sanitize\def\ion{\plainion}\gdef\ADS@keywords{#1}}
  \@firstkwdfalse
  \else%
  \edef\@keywords{\@keywords\@kwdsep \if@boldkeywords\uppercase{#1}\else #1\fi}%
  {\ADS@sanitize\def\ion{\plainion}\xdef\ADS@keywords{\ADS@keywords\ADS@kwdsep #1}}%
  \fi}
\newcommand{\ADS@sanitize}{%
  \def\mathrm{}
  \def\@{\ }
  \def~{\ }
  \def\langle{<}
  \def\rangle{>}
  \def\log{log }
  \def\sin{sin }
  \def\cos{cos }
  \def\tan{tan }
  \def\fd{^d.}
  \def\fh{^h.}
  \def\fm{^m.}
  \def\fs{^s.}
  \def\fdg{�.}
  \def\farcm{'.}
  \def\farcs{''.}
  \def\arcdeg{�}
  \def\,{ }
  }
\newcommand{\keywords}[1]{%
  \renewcommand{\ADS@keywords}{#1}
  \renewcommand{\@keywords}{\uppercase{#1}}
  }
\newcommand{\@keywordheader}{\rmaa@KeywordHeaderStyle{Key Words:}}
\newcommand{\@keywords}{%
  \ClassWarningNoLine{rmaa}{No keywords specified}%
  Keywords go here
  }
%    \end{macrocode}
% \begin{macro}{\shorttitle}
% \begin{macro}{\shortauthor}
% \begin{macro}{\@shorttitle}
% \begin{macro}{\@shortauthor}
% Commands for setting the running headers.
%    \begin{macrocode}
\newcommand{\shorttitle}[1]{%
  \renewcommand\@shorttitle{\MakeUpperCase{#1}}}
\newcommand{\shortauthor}[1]{%
  \renewcommand\@shortauthor{\MakeUpperCase{#1}}}
\newcommand{\@shorttitle}{SHORT TITLE GOES HERE}
\newcommand{\@shortauthor}{SHORT AUTHOR GOES HERE}
%    \end{macrocode}
% \end{macro}
% \end{macro}
% \end{macro}
% \end{macro}
% \end{macro}
% \end{macro}
% \end{macro}


% \subsubsection{The full addresses}
% \label{sec:full-addresses}
% \begin{macro}{\fulladdresses}
% \begin{macro}{\@fulladdresses}
% \begin{macro}{\@fulladdheight}
% \begin{macro}{\@fulladddepth}
% \begin{macro}{\@fulladdbox}
% Aaghh, the full addresses! A perennial source of problems. 
%    \begin{macrocode}
\newlength{\@fulladdheight}
\newlength{\@fulladddepth}
\newsavebox{\@fulladdbox}
%% \fulladdresses command modified (11 May 2000) so that we can put
%% \fulladdresses{} in the article to suppress the output of the
%% addresses
\newcommand{\@FinalFigure}{\relax}
\newcommand{\FinalFigure}[1]{\renewcommand{\@FinalFigure}{#1}}
\newcommand{\fulladdresses}[1]{%
  \renewcommand{\@rmaaArg}{#1}
  \ifx\@rmaaArg\@rmaaEmpty
     \renewcommand{\@fulladdresses}{}
  \else
     \renewcommand{\@fulladdresses}{
       \normalsize%
       \ifx\@FinalFigure\relax\else\vfill\@FinalFigure\vfill\fi
%%       \vspace*{1ex}
       %% do addresses as a list (copied from thebibliography environment) 
       \list{\null}{\leftmargin 1.5em\topsep\z@\labelwidth\z@\labelsep\z@\itemindent -1.5em
         \itemsep\z@\parsep\z@\usecounter{enumi}}
       \sloppy\clubpenalty4000\widowpenalty4000
       \sfcode`\.=1000\relax
       #1
       \endlist
       }
  \fi
  }
\newcommand{\@fulladdresses}{Full addresses go here}
%    \end{macrocode}
% \begin{macro}{\suppressfulladdresses}
% New user-level command for suppressing full addresses should you so
% desire. Currently used in the single-page-poster style. 
%    \begin{macrocode}
\newcommand{\suppressfulladdresses}{\@suppressfulladdressestrue\@needaddressesfalse}
%    \end{macrocode}
% \end{macro}
% \end{macro}
% \end{macro}
% \end{macro}
% \end{macro}
% \end{macro}

% \begin{macro}{\printdates}
% \begin{macro}{\noprintdates}
% Command to activate the printing out of the received/accepted
% dates. This can also be done using the \texttt{printdates} option,
% which is the default for \texttt{journal}, but not for \texttt{proceedings}.
%    \begin{macrocode}
\newcommand{\printdates}{\@printdatestrue}
\newcommand{\noprintdates}{\@printdatesfalse}
%    \end{macrocode}
% \end{macro}
% \end{macro}
% \begin{macro}{\@outputaddressesifnecessary}
% Command to output the addresses if we haven't already done so
% This is called by: 
% \begin{enumerate}
% \item The end of |thebibliography| environment
% \item The |\outputfulladdresses| command
% \end{enumerate}
%    \begin{macrocode}
\newcommand{\@outputaddressesifnecessary}{%
  \ClassInfo{rmaa}{Calling @outputaddressesifnecessary}
  \ifthenelse{\boolean{@needaddresses}}{%
    %% If in one column style, or in the first column of a two column
    %% style then we can try to put it on this page 
    \ifthenelse{\boolean{@firstcolumn} \or \not\boolean{@rmtwocol}}{%
      \@outputaddresses}{%
      %% otherwise give warning it will be on a separate page
      \ClassWarning{rmaa}{Full addresses will go on a separate
        page (please see User Guide for how to remedy this)}}%
    %% If we haven't managed yet, put it on its own page
    \ifthenelse{\boolean{@needaddresses}}{%
      \if@debug{\ClassWarning{rmaa}{In this silly bit!}}\fi
      \clearpage\vfill\@outputaddresses}{}%
    }{
    %% If we don't need them, we still might need to do the eoa hook 
    \ifthenelse{\boolean{@suppressfulladdresses}}{%
      \if@debug{\ClassWarning{rmaa}{XXX}}\fi
%%      \if@book\else\@endofarticlehook\null\fi
      }{}
    }%
  }  
%    \end{macrocode}
% \end{macro}

% \begin{macro}{\outputfulladdresses}
% A user interface to the above command, in case there is no biblio
% (11 May 2000) 
%    \begin{macrocode}
\newcommand{\outputfulladdresses}{\@outputaddressesifnecessary}
%    \end{macrocode}
% \end{macro}

% \begin{macro}{\adjustfinalcols}
% This command starts a new column and should be 
% used to ensure that the columns balance on the last page. It also
% outputs the full 
% author addresses to ensure that this occurs in the left hand
% column. 
% The |\par| command is so that the current line will be taken account
% of when working out |\pagetotal|
%    \begin{macrocode}
\newcommand{\adjustfinalcols}{\if@twocolumn\goodbreak\@outputaddresses\newpage\noindent\fi}
%    \end{macrocode}
% \end{macro}

% \begin{macro}{\@outputaddresses}
% Command to do the work of outputting the addresses. 
% This can be called by 
% \begin{enumerate}
% \item an explicit |\adjustfinalcols|
% \item via |\@outputaddressesifnecessary|, q.v.
% \end{enumerate}
%
%    \begin{macrocode}
\newcommand{\@outputaddresses}{%
  % Allow option of suppressing the full addresses, in which case
  % we just run the hook if we are not in a book style
  \if@suppressfulladdresses
%%     \if@book\else\@endofarticlehook\null\fi
  \else
  %% save the full addresses in a box for output later
  \sbox{\@fulladdbox}{\parbox{\textwidth}{\@fulladdresses}}
  % calculate height of address box  
  \settoheight{\@fulladdheight}{\usebox{\@fulladdbox}}
  \settodepth{\@fulladddepth}{\usebox{\@fulladdbox}}
  \addtolength{\@fulladdheight}{\@fulladddepth}
  \if@debug
  { \typeout{Address total height: \the\@fulladdheight}
    \typeout{Text height: \the\pagetotal}
    \typeout{Goal height: \the\pagegoal}
    }
  \fi
  % reuse \@fulladdheight to mean the total height of text+addresses
  \addtolength{\@fulladdheight}{\pagetotal}
%%  \addtolength{\@fulladdheight}{-\baselineskip}
  \if@debug\typeout{Adjusted text + addresses height: \the\@fulladdheight}\fi
  % check it will fit on current page
  \ifdim\@fulladdheight<\pagegoal
     %% must shift it to the left
     %% to compensate for the hanging indentation in the reference list. 
     %% NOTE: this is not foolproof!
     %% NB. This will give overfull hbox warnings, which is a bit ugly.
     \vfill\noindent\if@biblist\hspace*{-\leftmargin}\fi
     \usebox{\@fulladdbox}
     % Set this to show we've done it
     \@needaddressesfalse
     % and assume we are at the end of the article
     % so that we do the EOA hook if not in the book style
%%     \if@book\else\@endofarticlehook\fi
  \else
     % if not, just do nothing. With luck, this function will get 
     % called again on the next page
     \if@debug{\typeout{Holding over full addresses to next page}}\fi
  \fi\fi
  }
%    \end{macrocode}
% \end{macro}

% \subsubsection{Commands for editors}
% \label{sec:editor-commands}
% Some user/editor commands for doing the TOC and index in the book
% option. 
% They are defined here rather in the book-specific part since they
% need to be always available in order to avoid breaking documents. 
% \begin{macro}{\listofauthors}
% \begin{macro}{\@rawauthors}
% \begin{macro}{\rmaaTOCauthors}
% The command |\listofauthors| formats an author list for use in the
% TOC. It should contain the same names as in the |\author| command,
% but without the affiliation garbage and using |\&| instead of
% |and|. 
%    \begin{macrocode}
\newcommand{\@rawauthors}{Raw Authors}
\newcommand{\rmaaTOCauthors}{\noexpand\rmaaTOCAuthorStyle{Author list for TOC}}
\newcommand{\listofauthors}[1]{%
  \renewcommand{\@rawauthors}{#1}
  \renewcommand\rmaaTOCauthors{\noexpand\rmaaTOCAuthorStyle{#1}}}
%    \end{macrocode}
% \begin{macro}{\indexauthor}\label{indexauthor}
% The command |\indexauthor| produces an index entry for each
% author. 
% This command must come \emph{after} the
% |\title| and the |\listofauthors| commands. Also, there should be one 
% |\indexauthor| command for \emph{each} author. We also add the
% author onto the end of the |\ADS@authors| macro and construct the
% |\ADS@authorinitial| by a most cunning method (see below), but only
% for the first author. 
% 
% No, No, NO! This was all wrong. Backed out completely: 
% \begin{quote}
% Completely rewritten (19 Dec 2002) since we now use the index class, so 
% as to be able to alphabetize accented characters. A side effect of this is
% that |\index| no longer expands out the argument after the |@|, so we 
% instead add all the |\index| commands to a list (|\indexauthor@deferlist|), 
% which is executed in the |\@begofarticlehook|. 
% 
% Note the stuff with |\toks0| is copied from Exercise 20.15 of the \TeX{}book---it is
% to make sure the latin1-accented name gets passed into the idx
% file. 
% \end{quote}
% We now go back to the old way of doing things, using an optional
% argument to the \cs{indexauthor} command that has the name without
% accents. 
% \changes{3.08}{2003/01/16}{Remerged changes from 3.06 that got
% dropped from 3.07 after a synchronization disaster. Basically, the
% accentless optional argument to \cs{indexauthor}. Must add the macro
% directories to my unison set-up to stop this sort of thing from happening.}
%    \begin{macrocode}
\newcommand\@isortkey{}
\newcommand\@emptykey{empty}
\newcommand\indexauthor[2][empty]{
  \renewcommand\@isortkey{#1}
  \ifx\@isortkey\@emptykey \renewcommand\@isortkey{#2}\fi
  \index{\@isortkey @{{\rmaaIndexAuthorStyle{#2}}\hskip1em plus1em minus0.4em
      \rmaaIndexTitle\hskip1em plus2em minus0.4em\penalty-100
      \rmaaTOCauthors}}
  \indexauthorADS{#2}
}
%%\def\indexauthor@deferlist{\relax}
%%\long\def\indexauthor#1{
%%  \def\@@name{#1} \toks0=\expandafter{\@@name}
%%  \toks1=\expandafter{\indexauthor@deferlist}
%%  \protected@xdef\indexauthor@deferlist{%
%%    \the\toks1
%%    \noexpand\index{\the\toks0@{\noexpand\rmaaIndexAuthorStyle{\the\toks0}\noexpand\hskip1em plus1em minus0.4em 
%%        \rmaaIndexTitle\noexpand\hskip1em plus2em minus0.4em\noexpand\penalty-100
%%        \rmaaTOCauthors}}
%%  }
%%  \indexauthorADS{#1}
%%  }
%%\xdef\indexauthor@deferlist{}
\newcommand\indexauthorADS[1]{%
  \if@firstauthor\protected@xdef\ADS@authorinitial{\firstletter{#1}}\fi
  \protected@xdef\ADS@authors{\ADS@authors\if@firstauthor\else ; \fi
    #1}
  \@firstauthorfalse
}
%    \end{macrocode}
% Note that the first |#1| sets the key for the index, making sure it
% actually starts with  
% first letter of the surname. Meanwhile, the |@{...}| expression,
% which starts with some  
% control sequences is what is written to the |.ind| file. 
% See the makeindex man page for details. 
% \end{macro}
% \end{macro}
% \end{macro}
% \end{macro}

% \begin{macro}{\KillAllAccents}
% \begin{macro}{\firstletter}
% These are utility functions used in the above. |lett| is a counter
% indicating the current letter in the string. |\firstletter| just
% prints the first letter of its argument.
%    \begin{macrocode}
\newcommand{\KillAllAccents}{\def\'{}\def\`{}\def\~{}\def\.{}\def\i{i}}
\newcommand\firstletter[1]{\@car #1\@nil}
%    \end{macrocode}
% \end{macro}
% \end{macro}
% \begin{macro}{\SetVolume}
% \begin{macro}{\SetMSnumber}
% \begin{macro}{\SetFirstPage}
% \begin{macro}{\SetYear}
% \begin{macro}{\ReceivedDate}
% \begin{macro}{\AcceptedDate}
% These are for the editors to set the volume number, etc. 
%    \begin{macrocode}
\newcommand{\SetVolume}[1]{%
  \renewcommand{\@volume}{#1}
  \edef\ADS@volume{\PadToFourWithDots{#1}}
  \renewcommand{\@volumen}{\rmaa@VolumenStyle{volumen #1}}%
  }
\newcommand{\SetMSnumber}[1]{\renewcommand{\@MSnum}{#1}}
%    \end{macrocode}
% Note that there is no real reason to use |\SetFirstPage| since all
% the page numbering is done automatically in the book style. 
%    \begin{macrocode}
\newcommand{\SetFirstPage}[1]{\setcounter{page}{#1}}
\newcommand{\SetYear}[1]{%
  \renewcommand{\@year}{#1}
  \renewcommand{\ADS@year}{#1}
  \renewcommand{\@anyo}{\rmaa@VolumenStyle{#1}}
  }
\newcommand{\ReceivedDate}[1]{\renewcommand{\@received}{#1}}
\newcommand{\AcceptedDate}[1]{\renewcommand{\@accepted}{#1}}
%    \end{macrocode}
%
% \begin{macro}{\DOI}
% User/editor command to enter the DOI, which will be typeset as a link.
% \changes{4.0}{2019/11/02}{New user command \cs{DOI} to typeset
% DOI in the header of first page and add it to ADS info.}
%    \begin{macrocode}
\newcommand\DOI[1]{%
  \def\@DOI@url{\url{https://doi.org/10.22201/ia.01851101p.#1}}%
  \def\ADS@doi{https://doi.org/10.22201/ia.01851101p.#1}%
}%
%    \end{macrocode}
% \end{macro}
%
% \begin{macro}{\ORCIDiD}
% \changes{4.0}{2019/11/12}{Add mechanism for giving ORCID iD links}
% \begin{macro}{\ORCIDiDfull}
% \begin{macro}{\@orcid@icon}
% \begin{macro}{\@orcid@iconbig}
% Commands to add icon links to ORCID iDs of authors. For example
% \begin{quote}
% |\ORCIDiD|\marg{0000-0000-0000-0000}
% \end{quote}
% to be used after each author name in the |\author|
% command and in the |\fulladdresses|.  Alternatively,
% |\OORCIDiDfull| will typeset the icon plus the full text of the
% link, but we don't use that at present.  The icon image is stored in |ORCIDiD_iconvector.pdf|, which needs to be on the input path.  In principle, we could make EPS versions if necessary.
%    \begin{macrocode}
\newcommand\@orcid@icon{\includegraphics[height=1.8ex]{ORCIDiD_iconvector}}
\newcommand\@orcid@iconbig{\includegraphics[height=2ex]{ORCIDiD_iconvector}}
\newcommand\ORCIDiD[1]{\href{https://orcid.org/#1}{\,\@orcid@icon}}
\newcommand\ORCIDiDfull[1]{%
  \href{https://orcid.org/#1}{\@orcid@iconbig}\,\url{https://orcid.org/#1}}
%    \end{macrocode}
% \end{macro}
% \end{macro}
% \end{macro}
% \end{macro}
%
% \begin{macro}{\@volume}
% \begin{macro}{\@volumen}
% \begin{macro}{\@artpages}
% \begin{macro}{\@year}
% \begin{macro}{\@anyo}
% \begin{macro}{\@received}
% \begin{macro}{\@accepted}
% \begin{macro}{\@MSnum}
% Default values for the containers used by above commands, plus a
% few more useful ones. 
%    \begin{macrocode}
\newcommand{\@volume}{00}
\newcommand{\@volumen}{\rmaa@VolumenStyle{00}}
\newcommand{\@artpages}{%
  \oldpageref{\articlename:RMAAFirstPage}--\oldpageref{\articlename:RMAALastPage}}
\newcommand{\@year}{\the\year}
\newcommand{\@anyo}{\rmaa@VolumenStyle{2007}}
\newcommand{\@journalname}{\rmaa@JournalNameStyle{%
    \if@proceedings\RMAASClong\else\RMAAlong\fi}}
\newcommand{\@received}{}
\newcommand{\@accepted}{}
\newcommand{\@MSnum}{0000}
\newcommand{\@DOI@prefix}{https://doi.org/10.22201/ia.01851101p}
\def\@DOI@url{\url{https://doi.org/10.22201/ia.01851101p.2099.00.00.00}}
%    \end{macrocode}
% \end{macro}
% \end{macro}
% \end{macro}
% \end{macro}
% \end{macro}
% \end{macro}
% \end{macro}
% \end{macro}
% \end{macro}
% \end{macro}
% \end{macro}
% \end{macro}
% \end{macro}
% \end{macro}
% 
% The following commands are most useful in the \texttt{book}
% style. However, the |\SetConfTitle| command can now be used in
% individual articles to set the header for the \texttt{preprint}
% style. 
%    \begin{macrocode}
\newcommand{\@reunionplace}{Lugar de la reuni\'on}
\newcommand{\@reuniondate}{Fecha de la reuni\'on}
\newcommand{\@conftitle}{Conference Title}
\newcommand{\SetReunion}[2]{\renewcommand{\@reunionplace}{#1}\renewcommand{\@reuniondate}{#2}}
\newcommand{\SetConfTitle}[1]{\renewcommand{\@conftitle}{#1}}
\newcommand{\@editors}{Juan Ram\'on de la Fuente and ``El Mosh''}
\newcommand{\SetEditors}[1]{\renewcommand{\@editors}{#1}}
\newcommand{\@editores}{Juan Ram\'on de la Fuente y ``El Mosh''}
\newcommand{\SetEditores}[1]{\renewcommand{\@editores}{#1}}
\newcommand{\@mes}{mes de la publicaci\'on}
\newcommand{\SetMes}[1]{\renewcommand{\@mes}{\rmaa@VolumenStyle{#1}}}
%    \end{macrocode}

% \subsection{Notes for Author Proofs and Communication between Editors}
% \label{sec:notes}
% We introduce a command \cs{Note} that should be used to describe
% any change made to the manuscript. These are a bit tricky to get
% right. In general we use \cs{marginpar} but this only works in outer
% paragraph mode, so we just bung the box where we are if we are in a
% figure caption or something. If we are in the preamble, then we just
% use a \marginpar as part of the \maketitle command.
%    \begin{macrocode}
\newif\if@arenotes
\@arenotesfalse
\newcounter{notecounter}
\setcounter{notecounter}{0}
\newcommand\NotesList{}
\newcommand\Note[1]{%
  \if@proof{%
    \@arenotestrue%
    \stepcounter{notecounter}%
    \if@inpreamble\else{%
      \NoteMark{\arabic{notecounter}}%
    }\fi
    \protected@xdef\NotesList{\NotesList 
      \protect\item #1}%
  }\fi%
}
\newcommand\Changed[2]{Changed ``#1'' to ``#2''.}
\newcommand\NoteMark[1]{%
  \ifinner
  \raisebox{1ex}[0pt][0pt]{\makebox[0pt][l]{\framebox{\footnotesize\bfseries
        N#1}}}%
  \else
  \marginpar{\framebox{\footnotesize\bfseries N#1}}%
  \fi
}
\newcommand\RefNote[1]{\if@proof\NoteMark{\ref{#1}}\relax\fi}
\newcommand\NoteLabel{\protect\label}
\newcommand\OutputNotes{%
  \bgroup
  \onecolumn
  \section*{EDITORIAL NOTES}
  \begin{enumerate}
    \NotesList
  \end{enumerate}
  \egroup
}
\DeclareRobustCommand\cs[1]{\texttt{\char`\\#1}}

%    \end{macrocode}
% \subsection{Bibliographic Records for ADS}
% \label{sec:ADS}
% \changes{2.2beta11}{2001/11/27}{Just about finished all the ADS
% stuff now. Except for the references, which I have given up on for
% the time being.}
% \begin{macro}{\ADS@BibCode}
% \begin{macro}{\ADS@year}
% \begin{macro}{\ADS@journal}
% \begin{macro}{\ADS@volume}
% \begin{macro}{\ADS@firstpage}
% \begin{macro}{\ADS@authorinitial}
% We use the |\ADS@BibCode| command to hold the bibliographic code for
% ADS. It is 
% used in the |%R| 
% field of the bibliographic record for each article. 
% The format is 
% \begin{quote}
% \texttt{YYYYJJJJJVVVVMPPPPA}  
% \end{quote}
% where \texttt{YYYY} is the year, \texttt{JJJJJ} is the journal code,
% which is \texttt{RMxAA} or \texttt{RMxAC}, \texttt{VVVV} is the
% volume number (right-justified), \texttt{M} is not used by us,
% \texttt{PPPP} is the first page number (right-justified), and
% \texttt{A} is the initial of the first author's surname. The volume
% and page fields are padded with periods (`\texttt{.}'), which causes
% us some fun (this is done with the macro |\PadToFourWithDots| described
% below) . 
% 
% All these individual bits are held in the commands |\ADS@year|,
% |\ADS@journal|, |\ADS@volume|, |\ADS@firstpage|,
% |\ADS@authorinitial|. 
%    \begin{macrocode}
\newcommand{\SetADSBibCode}[1]{\renewcommand{\ADS@BibCode}{#1}}
\newcommand{\ADS@BibCode}
{\ADS@year\ADS@journal\ADS@volume.\PadToFourWithDots{\ADS@firstpage}\ADS@authorinitial}
\newcommand{\ADS@year}{2001} 
\newcommand{\ADS@journal}{\if@proceedings RMxAC\else RMxAA\fi}
\newcommand{\ADS@volume}{....} 
%    \end{macrocode}
% These two are just set to default values here, |\ADS@firstpage| is set
% in the |document| environment for included articles, in the
% |frontsection| environment, |\tableofcontents| command,  |index| and
% |abstracts| environments. Also by |\AtBeginDocument| for stand-alone
% articles. 
% 
% God knows how we are going to do |\ADS@authorinitial|!
%    \begin{macrocode}
\newcommand{\ADS@firstpage}{....}
\newcommand{\ADS@authorinitial}{A}
%    \end{macrocode}
% \end{macro}
% \end{macro}
% \end{macro}
% \end{macro}
% \end{macro}
% \end{macro}
% We also have a whole slew of other fields that can be set. These are
% all listed in Table~\ref{tab:ADS}. Most of them still have to be
% defined. 
% \begin{table}[!ht]\raggedleft
% \caption{ADS Bibliographic Records}\label{tab:ADS}
% \begin{tabular}{lp{8em}lp{8em}}\hline
% Code & Description & Macro Name & Comment \\ \hline
% |%R| & Bibliographic Code & |\ADS@BibCode| & required \\
% |%A| & Author list & |\ADS@authors| & required, generate from |\indexauthor| \\
% |%F| & Author affiliation & |\ADS@affil| & format: AA(\dots)
% AB(\dots) \dots \\
% |%J| & Journal Name & |\ADS@journalname| &  \\
% |%D| & Publication Date & |\ADS@fulldate| &  \\
% |%L| & Last Page & |\ADS@lastpage| & generate from label \\
% |%T| & Title & |\ADS@title| & required, generate from |\title| \\
% |%C| & Abstract Copyright & |\ADS@copyright| & set to IA-UNAM \\
% |%U| & URL for electronic document & |\ADS@url| & set to PDF file \\
% |%d| & DOI for article & |\ADS@doi| & set by |\DOI| \\
% |%K| & Keywords & |\ADS@keywords| & set equal to |\@keywords| \\
% |%A| & Abstract Text & |\ADS@abstract| & set equal to |\@abstract| \\
% |%R| & References & |\ADS@references| & generate from the
% bibliography  \\ \hline
% \end{tabular}
% \end{table}
% A lot of these macros are set by other commands, as mentioned in the
% Table, but some need to be set  explicitly here. 
% 
% \begin{macro}{\ADS@doi}
% \changes{4.0}{2019/11/08}{Save DOI for ADS}
% The user/editor command \cs{DOI} is used to set this. 
%    \begin{macrocode}
\newcommand{\ADS@doi}{https://doi.org/10.22201/ia.01851101p.0000.00.00.00}
%    \end{macrocode}
% \end{macro}
% 
% \begin{macro}{\ADS@journalname}
% This
% does not contain any accents at the moment so as not to cause
% problems---this maybe should be changed though. 
%    \begin{macrocode}
\newcommand{\ADS@journalname}{%
  \if@proceedings\@conftitle\space (Eds. \@editors) \fi%
  Revista Mexicana de Astronom\'\i{}a y Astrof\'\i{}sica%
  \if@proceedings\space (Serie de Conferencias)\fi%
  \space Vol. \@volume, pp. \ADS@firstpage--\ADS@lastpage\space (\@year)%
  \space (http://www.astroscu.unam.mx/\~{}rmaa/)
  }
%    \end{macrocode}
% \end{macro}
% \begin{macro}{\ADS@keywords}
% \begin{macro}{\ADS@abstract}
% \begin{macro}{\ADS@title}
% |\ADS@abstract| is straightforward, since we already have the text
% saved as |\@abstract|. We could do the same with |\ADS@keywords| but
% to avoid unwanted formatting we explicitly set it in |\keywords|
% (p.~\pageref{keywords}). |\ADS@title| is just set to |\@rawtitle|,
% which we were already using for another purpose. 
%    \begin{macrocode}
\let\ADS@title=\@rawtitle
\newcommand{\ADS@keywords}{No keywords specified}
\newcommand{\ADS@abstract}{No abstract specified}
%    \end{macrocode}
% \end{macro}\end{macro}\end{macro}
% \begin{macro}{\ADS@authors}
% |\ADS@authors| wants to be in ``Other, A. N.'' format, but with
% semi-colons between authors. The most natural way to do this seems to
% be with |\indexauthor| (p.~\pageref{indexauthor}). We append each
% author to this macro, which means we must explicitly clear it for
% each new article (done in |\article|). .   
%    \begin{macrocode}
\newcommand{\ADS@authors}{}
%    \end{macrocode}
% \end{macro}
% \begin{macro}{\ADS@fulldate}
% \begin{macro}{\ADS@lastpage}
% \begin{macro}{\ADS@copyright}
% \begin{macro}{\ADS@website}
% \begin{macro}{\ADS@url}
% \begin{macro}{\SetWebsite}
% \begin{macro}{\SetMonthNumeric}
% Other sundry data, together with user commands for bits that can be
% set by the Editors. The |\ADS@url| is constructed from the website
% (currently |http://www.astrosmo.unam.mx/rmaa/|,
% the journal descriptor (e.g., |RMxAC..10/|), followed by the
% path to the PDF file (e.g., |PDF/RMxAC..10_HEW1.pdf|). 
% 
%    \begin{macrocode}
\newcommand{\ADS@website}{http://www.astrosmo.unam.mx/rmaa/}
\newcommand{\ADS@url}
{\ADS@website\ADS@journal\ADS@volume\ADS@number/PDF/\ADS@journal\ADS@volume\ADS@number\_\articlename.pdf}
\newcommand{\ADS@number}{}
\newcommand{\ADS@fulldate}{00/\ADS@year} 
\newcommand{\ADS@lastpage}{00} 
\newcommand{\ADS@copyright}{(C) \ADS@year; Instituto de Astronom\'\i{}a,
  Universidad Nacional Aut\'onoma de M\'exico} 
\newcommand{\SetNumber}[1]{\renewcommand{\ADS@number}{-#1}}
\newcommand{\SetWebsite}[1]{\renewcommand{\ADS@website}{#1}}
\newcommand{\SetMonthNumeric}[1]{\renewcommand{\ADS@fulldate}{#1/\ADS@year}}
%    \end{macrocode}
% \end{macro}\end{macro}\end{macro}\end{macro}\end{macro}\end{macro}\end{macro}
% \begin{macro}{\ADS@affil}
% \begin{macro}{\ADS@references}
% These two are somewhat more problematical. Affiliation currently
% only works for the first author (see \dots) pending redesign of the
% user interface. References are now done by wrapping the
% |thebibliography| environment to open another input stream on the
% same file and search through for the |\bibitems|. 
%    \begin{macrocode}
\newcommand{\ADS@references}{}
\newcommand{\ADS@affil}{}
%    \end{macrocode}
% \end{macro}\end{macro}
% 
% Define two macros that expand to the strings
% |\begin{thebibliography}| and |\end{thebibliography}|. Note that the
% idea for this and much else is inspired by the code for the docstrip
% package. Also define |\StartBibX| in case authors use the more
% standard |\begin{thebibliography}{}| syntax. 
%    \begin{macrocode}
\edef\EndBib{\expandafter\@gobble\string\\end\string{thebibliography\string}}
\edef\StartBib{\expandafter\@gobble\string\\begin\string{thebibliography\string}}
\edef\StartBibX{\expandafter\@gobble\string\\begin\string{thebibliography\string}\string{\string}}
%    \end{macrocode}
% Define two new logical variables. |\ifAnotherline| is true if we
% want to read another line from the input stream. It is set to false
% when we encounter the |\end{thebibliography}|. |\ifInBib| is set to
% true when we frind the |\begin{thebibliography}| and tells us to
% start saving the following lines. 
%    \begin{macrocode}
\newif\ifAnotherLine
\newif\ifInBib
%    \end{macrocode}
% \begin{macro}{makeOther}
% We need a macro to change the catcodes of lots of things to other. 
%    \begin{macrocode}
\def\makeOther#1{\catcode`#1=12\relax}
%    \end{macrocode}
% \end{macro}
% \begin{macro}{\RereadBiblio}
% Define a macro to do all the fiddly stuff. This is used at the
% beginning of |thebibliography| environment (Sec~\ref{sec:biblio}).
%    \begin{macrocode}
\newread\tf@biblio
\newcommand{\RereadBiblio}{%
%    \end{macrocode}
% Open the current file again for reading
%    \begin{macrocode}
  \begingroup
  \immediate\openin\tf@biblio=\articlename
%    \end{macrocode}
% Everything in a group to be safe. Initialize logicals. Black magic
% to change
% catcodes and avoid blank lines being interpreted as |\par|. 
%    \begin{macrocode}
  \AnotherLinetrue\InBibfalse
  \makeOther\ \makeOther\\\makeOther\$%
  \makeOther\#\makeOther\^\makeOther\^^K%
  \makeOther\_\makeOther\^^A\makeOther\%%
  \makeOther\~\makeOther\{\makeOther\}\makeOther\&%
  \endlinechar-1\relax
%    \end{macrocode}
% Start looping through the file. Read each line to the macro
% |\Line|. If we have got to EOF or end of biblio then set
% |\AnotherLinefalse|.
%    \begin{macrocode}
  \loop
  \read\tf@biblio to \Line
  \ifx\Line\EndBib\AnotherLinefalse\InBibfalse\fi
  \ifeof\tf@biblio \AnotherLinefalse\fi
%    \end{macrocode}
% If we are in the biblio and the line is not blank then append it to
% the |\ADS@references| macro together with a newline. 
%    \begin{macrocode}
  \ifInBib
  \ifx\Line\empty\else\xdef\ADS@references{\ADS@references\Line^^J}\fi
  \fi
%    \end{macrocode}
% Check if the line we just saw was the start of biblio, then repeat
% the loop if necessary. 
%    \begin{macrocode}
  \ifx\Line\StartBib\InBibtrue\fi
  \ifx\Line\StartBibX\InBibtrue\fi
  \ifAnotherLine\repeat
  \immediate\closein\tf@biblio
  \endgroup
  }
%    \end{macrocode}
% \end{macro}
% 
% \begin{macro}{\PadToFourWithDots}
% Utility function for left-padding a number to 4 characters with
% periods. Used in 
% constructing |\ADS@BibCode|. This is a very simple algorithm and
% \emph{requires} that the argument actually be a number---if it is
% not then dire consequences will result.  
%    \begin{macrocode}
\newcommand{\PadToFourWithDots}[1]{%
  \ifnum #1 > 999 #1%
  \else\ifnum #1 > 99 .#1%
  \else\ifnum #1 > 9 ..#1%
  \else ...#1%
  \fi\fi\fi}
%%  \ifnum \count0 > 999 #1%
%%  \else\ifnum \count0 > 99 .#1%
%%  \else\ifnum \count0 > 9 ..#1%
%%  \else ...#1%
%%  \fi\fi\fi}
%    \end{macrocode}
% \end{macro}
%
% \begin{macro}{\AuthorADS}
% \begin{macro}{\ADS@affiliations}
% New mechanism for setting the affiliations for ADS. The author labels (AA, AB, \dots) are generated using two nested counters.  We use the |\robustify| command from \texttt{etoolbox} so that accented characters in the affiliation addresses do not blow up.  The list of accents may need to be expanded. 
%    \begin{macrocode}
\newcounter{AUTHOR}
\newcounter{author}[AUTHOR]
\stepcounter{AUTHOR}
\newcommand\ADS@affiliations{}
\robustify{\'}\robustify{\i}\robustify{\~}\robustify{\`}
\def\AuthorADS#1#2{%
  \begingroup
  \def\ORCIDiD##1{<ID system="ORCID">##1</ID>}
  \def\Email##1{<EMAIL>##1</EMAIL>}
  \def\Affil##1{##1}
  \stepcounter{author}
  \ifthenelse{\value{author} > 26}{\stepcounter{AUTHOR}\stepcounter{author}}{}
  \ifthenelse{\equal{\ADS@affiliations}{\@rmaaEmpty}}{%
    \xdef\ADS@affiliations{\Alph{AUTHOR}\Alph{author}(#2)}
  }{%
    \xdef\ADS@affiliations{\ADS@affiliations, \Alph{AUTHOR}\Alph{author}(#2)}
  }%
  \endgroup
}
%    \end{macrocode}
% \end{macro}
% \end{macro}
% %
% \begin{macro}{\ADS@dumpdata}
% \changes{4.0}{2019/11/13}{Write all author affiliations for ADS}
% \begin{macro}{\ADS@write}
% \begin{macro}{\ADS@openfile}
% |\ADS@dumpdata| is the macro that actually does the job of writing out all the
% data to an auxiliary file, using |\ADS@write| to write each
% field. The file is set up with |\ADS@openfile|, which currently is
% done in the TOC of a book.  The new multi-author affiliations are used if |\ADS@affiliations| has been set, otherwise it falls back to the single-author version |\ADS@affil|. 
%    \begin{macrocode}
\newcommand{\ADSfield}[2]{\relax}
\newcommand{\ADSvfield}[2]{\relax}
%%\renewcommand{\ADSfield}[12]{\%#1 #2}
\newcommand{\ADS@dumpdata}
{ 
  \ADS@write{\string\begin{ADSarticle}}
  \ADS@WRITE{\string\ADSfield{R}{\ADS@BibCode}}
  \ADS@write{\string\ADSfield{A}{\ADS@authors}}
  \ifthenelse{\equal{\ADS@affiliations}{\@rmaaEmpty}}{%
    {\ADS@sanitize\ADS@write{\string\ADSfield{F}{AA(\ADS@affil)}}}
  }{%
    {\ADS@sanitize\ADS@write{\string\ADSfield{F}{\ADS@affiliations}}}
  }
  \ADS@write{\string\ADSfield{J}{\ADS@journalname}}
  \ADS@write{\string\ADSfield{D}{\ADS@fulldate}}
  \ADS@write{\string\ADSfield{L}{\ADS@lastpage}}
  {\ADS@sanitize\ADS@write{\string\ADSfield{T}{\ADS@title}}}
  \ADS@write{\string\ADSfield{C}{\ADS@copyright}}
  \ADS@write{\string\ADSfield{d}{\ADS@doi}}
  \ADS@write{\string\ADSfield{U}{\ADS@url}}
  \ADS@write{\string\ADSfield{K}{\ADS@keywords}}
  {\ADS@sanitize\ADS@write{\string\ADSfield{B}{\ADS@abstract}}}
  \ADS@write{\string\ADSfield{Z}{\ADS@references}}
  \ADS@write{\string\end{ADSarticle}}
  }
%%\newcommand{\ADS@write}[1]{\if@book\addtocontents{ads}{#1}\fi}
%%\newcommand{\ADS@openfile}{\@starttoc{ads}}
\newcommand{\ADS@write}[1]{\set@display@protect\immediate\write\tf@ads{#1}}
\newcommand{\ADS@WRITE}[1]{\immediate\write\tf@ads{#1}}
\newcommand{\ADS@openfile}{%
  \newwrite\tf@ads%
  \immediate\openout\tf@ads\jobname.ads\relax
  \ADS@WRITE{\ADS@preamble}}
\newcommand{\ADS@closefile}{\ADS@WRITE{\ADS@postamble}}
\def\ADS@preamble{%
  %%  \string\documentclass{article}
  %%  \string\begin{document}
  }
\def\ADS@postamble{%
  %%  \string\end{document}
  \string%\string% Local Variables:^J 
  \string%\string% mode: LaTeX^J
  \string%\string% mode: Font-lock^J
  \string%\string% End:^J
  }
%    \end{macrocode}
% \end{macro}
% \end{macro}
% \end{macro}
% 
% \subsection{Automatic calculation of the page range}
% \label{sec:page-range}
% Another issue that has cost me a lot of grief. Still, it finally
% seems to be sorted now. All hail Jeff Goldberg. 
%
% Cunning trick to get the page numbers from |\pageref|.  Put a label
% on the first and last page of the document using |\AtBeginDocument|
% and |\AtEndDocument|. .  This way won't work for the book option.
% In that case, we do it by defining |\@endofarticlehook| and
% |\@begofarticlehook| as parallels to |\AtEndDocument| and
% |\AtBeginDocument|.  Note that we have to defer adding
% |\@begofarticlehook| to |\AtBeginDocument| to make sure it comes
% after the redefinitions of |\ref| and |\label|.
%    \begin{macrocode}
\if@book\AtEndDocument{\ADS@closefile}\else
 %\AtBeginDocument{\@begofarticlehook}
\AtEndDocument{\@endofarticlehook}
\fi
%    \end{macrocode}
% \begin{macro}{\@begofarticlehook}
% Since we are getting more and more things that need to be done at
% the beginning of every article, frontsection, etc we define a sister
% hook to 
% |\@endofarticlehook|. This hook does the following: 
% \begin{enumerate}
% \item Labels the first page
% \item Sets |\ADS@firstpage|
% \item runs and then clears |\indexauthor@deferlist|
% \item resets the keyword counter
% \end{enumerate}
% It is executed in the following circumstances: 
% \begin{enumerate}
% \item At |\begin{document}| for standalone files
% \item At |\begin{document}| for each included file
% \item At beginning of front sections 
% \item At beginning of TOC
% \item At beginning of Author Index
% \item At beginning of abstracts
% \end{enumerate}
% For some reason, this is not always within the right group, so we
% use |\xdef| just to be sure. 
%    \begin{macrocode}
\newcommand{\@begofarticlehook}{%
  \if@frontmatter\frontlabel{RMAAFirstPage}\else\label{RMAAFirstPage}\fi%
  \xdef\ADS@firstpage{\@arabic{\c@page}}%
  \if@book\else\makeindex\fi % needed so default index is defined in stand-alone articles
%%  \indexauthor@deferlist
%%  \xdef\indexauthor@deferlist{}
  \@inpreamblefalse
  \setcounter{kwdcounter}{0}
  }
%    \end{macrocode}
% \end{macro}
% \begin{macro}{\@endofarticlehook}
% This is a hook to be run at the end of each article
% currently, it just labels the last page and writes out all the ADS
% data to an auxiliary file.  
% It should not produce any printed output. 
% It is modelled on the |\AtEndDocument| command, to which it reduces
% in the case of a single article. 
%    \begin{macrocode}
\newcommand{\@endofarticlehook}{%
  \if@debug{\ClassWarning{rmaa}{***Doing Hook***}}\fi
  \if@debug{\ClassWarning{rmaa}{ADSabstract is: \ADS@abstract}}\fi
  \clearpage\rmaa@lastpagelabel\if@book\ADS@dumpdata\fi
  \ifx\NotesList\@rmaaEmpty\else\OutputNotes\fi
  \if@debug{\ClassWarning{rmaa}{Notes list is: \NotesList}}\fi
}
%    \end{macrocode}
% \begin{macro}{\rmaa@lastpagelabel}
% New version of last page label command copied from Jeff Goldberg's
% lastpage.sty package - we can't just use |\label| here in case the
% output routine for the last page has already been called. Instead,
% we first explicitly clear the page, step back the page counter,
% explicitly write to the .aux file, then restep the page counter
% again.
%
% But this will conflict with the hyperref package, which does labels differently.  Compensate for that by using the five-argument form.
% 
% Also set the |\ADS@LastPage| macro to the appropriate number. 
%    \begin{macrocode}
\newcommand{\rmaa@lastpagelabel}{%
  \addtocounter{page}{-1}%
  \immediate\write\@auxout{\string
    \newlabel{\articlename:RMAALastPage}{{}{\thepage}{}{}{}}}%
  \edef\ADS@lastpage{\thepage}%
  \addtocounter{page}{1}}
%    \end{macrocode}
% \end{macro}
% \end{macro}

% \subsection{Deal with explicit change in number of columns}
% \label{sec:columns}
% Sometimes, when formatting a book, we need to switch between
% one-column articles and two-column articles. For this to work, we
% need to wrap the standard \LaTeX{} commands to set the |\@rmtwocol|
% flag appropriately and also reset the package default values of all
% scalable lengths. .  
%    \begin{macrocode}
\let\latexonecolumn=\onecolumn
\let\latextwocolumn=\twocolumn
\renewcommand{\onecolumn}{\@rmtwocolfalse\@package@default@lengths\latexonecolumn}
\renewcommand{\twocolumn}{\@rmtwocoltrue\@package@default@lengths\latextwocolumn}
%    \end{macrocode}
% 
% \subsection{The actual formatting of the title}
% \label{sec:title-format}

% Note, we have to redefine the |\maketitle| command to use the
% numbered footnotes. This means we needn't use the atbeginend stuff. 
%    \begin{macrocode}
\renewcommand{\maketitle}{\par
  \begingroup
%    \end{macrocode}
% Don't switch to marked footnotes. 
% This is consistent with existing style but I'm not sure I like it
%    \begin{macrocode}
  %%  \renewcommand\thefootnote{\fnsymbol{footnote}}%
%%  \def\@makefnmark{\hbox to\z@{$\m@th^{\@thefnmark}$\hss}}%
  \def\@makefnmark{\textsuperscript{\@thefnmark}\hspace{-0.2em}}%
  \long\def\@makefntext##1{\parindent 1em\noindent
    \hbox to1.8em{\hss$\m@th^{\@thefnmark}$}##1}%
  \if@rmtwocol
  { \if@debug{\ClassWarning{rmaa}{TWO COLUMN}}\fi
    \twocolumn[\@maketitle]%
    }
  \else
  { \if@debug{\ClassWarning{rmaa}{ONE COLUMN}}\fi
    \newpage
    \global\@topnum\z@   % Prevents figures from going at top of page.
    \if@manuscript \@MSmaketitle \else \@maketitle \fi
    }
  \fi
  \if@manuscript\else\thispagestyle{rmaafirstpage}\@thanks\fi
  \ifx\NotesList\@rmaaEmpty\else{%
    \NoteMark{1--\arabic{notecounter}}\vspace*{-2\baselineskip}%
  }\fi%
  \endgroup
%    \end{macrocode}
% Don't reset footnote counter (not sure whether this is good thing
% or not).
%    \begin{macrocode}
  %%  \setcounter{footnote}{0}%
  \let\thanks\relax\let\affil\relax\let\altaffil\relax
  \let\maketitle\relax\let\@maketitle\relax
  \gdef\@thanks{}\gdef\@author{}\gdef\@title{}}
%    \end{macrocode}

% \begin{macro}{\@kwhwidth}
% \begin{macro}{\@kwtwidth}
% Redefine \cs{@maketitle} to include the abstract, resumen and
% keywords. Note that in the manuscripty style we use
% \cs{@MSmaketitle} instead.  
% \changes{3.16}{2003/04/07}{
%  Added extra vertical space between authors and accepted date.}
% \changes{3.26}{2011/08/15}{Add extra label to review articles}
%    \begin{macrocode}
\newlength{\@kwhwidth}\newlength{\@kwtwidth}
\newcommand\reviewname{Review}
\renewcommand{\@maketitle}{%
  \newpage
  \null
  \if@reviewarticle
  \smallskip
  {\large\bfseries\itshape\reviewname}%
  \bigskip
  \else
  \vskip\@before@title@skip%
  \fi
  \begin{center}%
    {\@title \par}%
    \vskip\@before@author@skip%
    \parbox{\rmaa@AuthorWidthFrac\textwidth}{\centering
      \lineskip\@between@author@skip%
      \@author
      \par}
    %% this should be replaced with `Submitted to...' for preprints
    \if@printdates
    \vskip\@before@dates@skip%
    \if@preprint
    \rmaa@ReceivedStyle{Submitted to RMxAA \@received}
    \else
    \rmaa@ReceivedStyle{Received \@received; accepted \@accepted}
    \fi
    \fi
    \if@manuscript
    \rmaa@ReceivedStyle{Submitted to RMxAA \@received}
    \fi
    \if@guide
    \rmaa@ReceivedStyle{Version \rmaa@version, \rmaa@date}
    \fi
    \vskip\@before@abstract@skip%
    \if@abstractpresent
    \parbox{\rmaa@AbstractWidthFrac\textwidth}{%
      \setlength{\parindent}{0.9cm}%
      \centerline{\rmaa@AbstractHeaderStyle{Abstract}}\par%
      \vskip\@after@abstracthead@skip%
      \if@newstyle\noindent\fi\@abstract
      \vskip\@between@abstract@skip%
      \centerline{\rmaa@AbstractHeaderStyle{Resumen}}\par%
      \vskip\@after@abstracthead@skip%
      \if@newstyle\noindent\fi\@resumen
      \if@setresumen\else\@abstract\fi % append english abstract if
                                       % spanish one not set
      \if@guide\else            %do not use keywords for guide
      \vskip\@before@kwd@skip%
      % Format the Key Words with hanging indentation:
      % They are put in a \parbox with width equal to the line width
      % minus the width of the header. 
      { % Widths of keyword
        \settowidth{\@kwhwidth}{\@keywordheader~}   % header and
        \setlength{\@kwtwidth}{\linewidth}          % keyword text
        \addtolength{\@kwtwidth}{-\@kwhwidth}       %
        \noindent\@keywordheader~\parbox[t]{\@kwtwidth}{%
          \rmaa@KeywordStyle{\boldmath\uppercase{\@keywords}}}}%
      \fi
      }
      \fi
  \end{center}%
  \vspace*{\afterkwdskip}
  \par
  }
\newlength{\afterkwdskip}\setlength{\afterkwdskip}{\z@}
%    \end{macrocode}
% \end{macro}
% \end{macro}
%
%
% \begin{macro}{\reviewarticle}
% \begin{macro}{\reviewtoc}
% \begin{macro}{\reviewlof}
%   User command to specify a review article. Must be used in
%   preamble. Main effect is to place the text ``Review'' before the
%   article title.  Also add command for setting a table of contents
%   and a list of figures at the beginning of the article.
%    \begin{macrocode}
\newcommand{\reviewarticle}{\@reviewarticletrue}
\newcommand\reviewtoc{%
  \section*{CONTENTS}%
  \@starttoc{rtc}%
}
\newcommand\reviewlof{%
  \section*{LIST OF FIGURES}%
  \@starttoc{lof}%
}
%    \end{macrocode}
% \end{macro}
% \end{macro}
% \end{macro}
% \begin{macro}{\l@section}
%    \begin{macrocode}
\renewcommand*\l@section[2]{%
  \ifnum \c@tocdepth >\z@
    \addpenalty\@secpenalty
    \addvspace{0.8em \@plus\p@}%
    \setlength\@tempdima{1.5em}%
    \begingroup
      \parindent \z@ \rightskip \@pnumwidth
      \parfillskip -\@pnumwidth
      \leavevmode %\bfseries
      \advance\leftskip\@tempdima
      \hskip -\leftskip
      #1\nobreak\hfil \nobreak\hb@xt@\@pnumwidth{\hss #2}\par
    \endgroup
  \fi}
%    \end{macrocode}
% \end{macro}
%
% \begin{macro}{\@MSmaketitle}
% A special version of |\@maketitle| for the \texttt{manuscript}
% option. This can deal with a multi-page abstract/resumen. Not very
% well, mind. 
%    \begin{macrocode}
\newcommand{\@MSmaketitle}{%
  \if@debug\ClassInfo{rmaa}{Using special MS maketitle command}\fi
  \newpage
  {
    \thispagestyle{rmaafirstpage}
%     \setlength{\parindent}{0pt}
    \null\vskip 2ex
    \begin{center}
      \@title\par
      \@author\par
      \rmaa@ReceivedStyle{%
        \ifx\@received\@rmaaEmpty
          Draft version: \@date
        \else
          Received: \@received
          \ifx\@accepted\@rmaaEmpty
          \else
            ; Accepted: \@accepted
          \fi
        \fi}\par
    \end{center}
    \@thanks
    \centerline{\rmaa@AbstractHeaderStyle{Resumen}}
    \noindent\@resumen\par\vskip 2ex
    \centerline{\rmaa@AbstractHeaderStyle{Abstract}}
    \noindent\@abstract\par\vskip 2ex
    \noindent\@keywordheader~\rmaa@KeywordStyle{\@keywords}\par
%     \@keywordheader~\rmaa@KeywordStyle{\uppercase{\@keywords}}\par
  }%
  }
%    \end{macrocode}
% \end{macro}

% \subsection{The sectioning commands}
% \label{sec:section-commands}

% \begin{macro}{\section}
% \begin{macro}{\subsection}
% \begin{macro}{\subsubsection}
% \begin{macro}{\paragraph}
% \begin{macro}{\subparagraph}
% Redefine all the standard commands to use these lengths. 
%    \begin{macrocode}
\renewcommand\section{\@startsection {section}{1}{\z@}%
                                   {\abovesectionskip}%
                                   {\belowsectionskip}%
                                   %% how do we get upper case here
                                   {\centering\reset@font\normalsize}}
\renewcommand\subsection{\@startsection{subsection}{2}{\z@}%
                                     {\abovesectionskip}%
                                     {\belowsectionskip}%
                                     {\centering\reset@font\normalsize\itshape}}
\renewcommand\subsubsection{\@startsection{subsubsection}{3}{\z@}%
                                     {\abovesectionskip}%
                                     {\belowsectionskip}%
                                     {\centering\reset@font\normalsize\itshape}}
\renewcommand\paragraph{\@startsection{paragraph}{4}{\z@}%
                                    {\aboveparaskip}%
                                    {\belowparaskip}%
                                    {\reset@font\normalsize\bfseries}}
\renewcommand\subparagraph{\@startsection{subparagraph}{5}{\parindent}%
                                      {\aboveparaskip}%
                                      {\belowparaskip}%
                                      {\reset@font\normalsize\bfseries}}
%    \end{macrocode}
% \end{macro}
% \end{macro}
% \end{macro}
% \end{macro}
% \end{macro}

% \begin{macro}{secnumdepth}
% Make sure that we number the subsubsections---this is done by
% default in the |article| class but not the |book| class, which is
% why we must set it explicitly here.
%    \begin{macrocode}
\setcounter{secnumdepth}{3}
%    \end{macrocode}
% \end{macro}
% \changes{2.2beta10}{2001/11/07}{Fixed \cs{secnumdepth} for \cs{book}
% style.}

% \begin{macro}{\thesubsection}
% \begin{macro}{\thesubsubsection}
% For subsection and subsubsection, we want the numbers to be roman
% but the title to be italic. This seems the easiest way to go about it.
%    \begin{macrocode}
\renewcommand{\thesubsection}{{\upshape\thesection.\@arabic\c@subsection}}
\renewcommand{\thesubsubsection}{\thesubsection{\upshape.\@arabic\c@subsubsection}}
%    \end{macrocode}
% \end{macro}
% \end{macro}

% \begin{macro}{\@seccntformat}
% Add ``.'' to the number here so we can remove it from |\thesection|,
% etc. Reduce gap between number and title WJH 27 Jan 1995
%    \begin{macrocode}
\def\@seccntformat#1{\csname the#1\endcsname .\hskip 0.4em}
%    \end{macrocode}
% \end{macro}

% \begin{macro}{\@sect}
% I didn't really want to mess with the |\@sect| command, but it was
% the only way I could manage to force upper case. Other than the
% commented parts, it is identical to the standard version. 
%    \begin{macrocode}
\def\@sect#1#2#3#4#5#6[#7]#8{
  \ifnum #2>\c@secnumdepth
     \let\@svsec\@empty
  \else
     \refstepcounter{#1}%
     \let\@@protect\protect
     \def\protect{\noexpand\protect\noexpand}%
     \edef\@svsec{\@seccntformat{#1}}%
     \let\protect\@@protect
  \fi
%    \end{macrocode}
% This is the bit I changed (27 Jan 1995). Use upper case for
% |\section|, use mixed case for the rest. 
%    \begin{macrocode}
  \ifnum #2=1                   
     \def\sec@upcase##1{\uppercase{##1}}%
     \def\sec@shape{\upshape}%
  \else                         
     \def\sec@upcase##1{\relax{##1}}%
     \def\sec@shape{\itshape}%
  \fi
  \@tempskipa #5\relax
  \ifdim \@tempskipa>\z@
     \begingroup #6\relax
          \@hangfrom{\hskip #3\relax\@svsec}%
                    {\interlinepenalty \@M \sec@upcase{#8}\par}%
     \endgroup
     \csname #1mark\endcsname{#7}%
%    \end{macrocode}
% Only write TOC entries if doing the Guide style. 
%    \begin{macrocode}
     \if@guide 
        \addcontentsline{toc}{#1}{%
          \ifnum #2>\c@secnumdepth \else
             \protect\numberline{\csname the#1\endcsname}%
          \fi
          #7}%
     \fi
%    \end{macrocode}
% Also allow per-article TOC for review papers
%    \begin{macrocode}
     \if@reviewarticle
        \ifnum #2>2 \else
           \addcontentsline{rtc}{#1}{%
             \protect\numberline{\csname the#1\endcsname}%
             \sec@upcase{\sec@shape #7}}%
        \fi
     \fi
  \else
     \def\@svsechd{#6\hskip #3\relax  %% \relax added 2 May 90
       \@svsec #8\csname #1mark\endcsname{#7}%
       \if@guide
          \addcontentsline{toc}{#1}{%
            \ifnum #2>\c@secnumdepth \else
               \protect\numberline{\csname the#1\endcsname}%
            \fi
            #7}%
       \fi
       }%
  \fi
  \@xsect{#5}}
%    \end{macrocode}
% \end{macro}


% \subsection{Appendices}
% \label{sec:appendices}
% Make sure that appendices in individual articles work properly. We
% copy the definition from the article class (the book one doesn't do
% what we want).  
% \begin{environment}{appendices}
%   This uses an environment |appendices| inside of which each
%   |\section| will become a separate appendix.  Note that the
%   appendices must come before the references (otherwise the full
%   addresses mechanism will not work properly). We make the equation
%   numbering be subordinate to the appendix (equations A12, B6,
%   etc). Finally, we have a medskip followed by APPENDICES to
%   introduce them. \changes{3.17}{2002/04/08}{ Added appendices environment
%   that does the right thing.}
%    \begin{macrocode}
\newenvironment{appendices}{\par
  \begingroup
  \setcounter{section}{0}%
  \setcounter{subsection}{0}%
  \renewcommand\thesection{\@Alph\c@section}
  \renewcommand\theequation{\thesection\arabic{equation}}
  \medskip
  \section*{APPENDICES}}
{\endgroup}
%    \end{macrocode}
% \end{environment}
% \begin{environment}{appendix}
%   If there is only one appendix, then use |appendix| environment
%   instead.  \changes{4.0}{2019/11/12}{ Added
%   appendix environment for case where there is only one section.}
%    \begin{macrocode}
\renewenvironment{appendix}{\par
  \begingroup
  \setcounter{section}{0}%
  \setcounter{subsection}{0}%
  \renewcommand\thesection{\@Alph\c@section}
  \renewcommand\theequation{\thesection\arabic{equation}}
  \medskip
  \section*{APPENDIX}}
{\endgroup}
%    \end{macrocode}
% \end{environment}

% \subsection{Floating thingies}
% \label{sec:floats}   

%% 
%    \begin{macrocode}
\renewcommand{\thetable}{\@arabic\c@table}
\renewcommand{\thefigure}{\@arabic\c@figure}
\renewcommand{\figurename}{Fig.}
\renewcommand{\tablename}{TABLE}
%\renewcommand{\fnum@table}{\MakeUpperCase{\tablename}~\thetable}
%    \end{macrocode}
% For some reason, I have to write `TABLE' explicitly here
% If I use |\tablename|, it gets converted to mixed case---God knows
% why. 
% |\figurename| is also being ignored. 
%    \begin{macrocode}
\renewcommand{\fnum@table}{TABLE~\thetable}
\renewcommand{\fnum@figure}{Fig.~\thefigure}
%    \end{macrocode}

% Do captions in 9 pt face in the main journal style and in new-style
% anything. 
%    \begin{macrocode}
\long\def\@makefigcap#1#2{%
  \if@proceedings\else\small\fi
  \if@newstyle\small\fi 
  \vskip\abovecaptionskip
  \sbox\@tempboxa{#1. #2}%
  \ifdim \wd\@tempboxa >\hsize
    #1. #2 \par
  \else
    \global \@minipagefalse
    \hb@xt@\hsize{\hfil\box\@tempboxa\hfil}%
  \fi
  \vskip\belowcaptionskip}
%    \end{macrocode}

% Make sure table caption is upper case. 
%    \begin{macrocode}
\long\def\@maketabcap#1#2{%
  \hbox to\hsize{\hfil #1 \hfil}
  \vskip 6pt 
  \centering\MakeUpperCase{#2}\par
  \vskip 6pt}
%    \end{macrocode}

%    \begin{macrocode}
\newlength{\@rottabcapwidth}
\long\def\@makerottabcap#1#2{%
  \setbox\@tempboxa\hbox{\MakeUpperCase{#2}}%
  \ifdim\wd\@tempboxa >.8\vsize
     \setlength\@rottabcapwidth{.8\vsize}
  \else
     \setlength\@rottabcapwidth{\wd\@tempboxa}
  \fi
  \rotatebox{90}{%
    \begin{minipage}{\vsize}\centering
      \begin{minipage}{\@rottabcapwidth}\centering
        #1\\ \vskip 6pt \MakeUpperCase{#2}
      \end{minipage}%
    \end{minipage}
    }%
  \hskip 6pt
  }

\newenvironment{Sideways}
    {\begin{sideways}\begin{minipage}{\vsize}\centering}
    {\end{minipage}\end{sideways}}

\newenvironment{Sidewaystable}
{%
  \let\oldcaption=\caption%
  \def\caption{
    \def\@captype{table}\renewcommand{\@makecaption}{\@maketabcap}
    \oldcaption}
  \begin{sidewaystable}\centering}
  {\end{sidewaystable}\let\caption=\oldcaption}
\newenvironment{Sidewaystable*}
{%
  \let\oldcaption=\caption%
  \def\caption{
    \def\@captype{table}\renewcommand{\@makecaption}{\@maketabcap}
    \oldcaption}
  \begin{sidewaystable*}\centering}
  {\end{sidewaystable*}\let\caption=\oldcaption}
%    \end{macrocode}

% Not sure why we need to stick these lines in. 
%    \begin{macrocode}
\renewenvironment{figure}
                 {\renewcommand{\@makecaption}{\@makefigcap}\@float{figure}}
                 {\end@float}
\renewenvironment{figure*}
                 {\renewcommand{\@makecaption}{\@makefigcap}\@dblfloat{figure}}
                 {\end@dblfloat}
\renewenvironment{table}
                 {\renewcommand{\@makecaption}{\@maketabcap}\@float{table}}
                 {\end@float}
\renewenvironment{table*}
                 {\renewcommand{\@makecaption}{\@maketabcap}\@dblfloat{table}}
                 {\end@dblfloat}
%    \end{macrocode}

% Redefine |\@caption| so that we don't bother with the list of
% figures. This saves us having to worry about fragile commands in the
% captions. 
%    \begin{macrocode}
\long\def\@caption#1[#2]#3{%
  \par
%%  \addcontentsline{\csname ext@#1\endcsname}{#1}%
%%    {\protect\numberline{\csname the#1\endcsname}{\ignorespaces #2}}%
  \begingroup
    \@parboxrestore
    \if@minipage
      \@setminipage
    \fi
    \normalsize
    \@makecaption{\csname fnum@#1\endcsname}{\ignorespaces #3}\par
  \endgroup}
%    \end{macrocode}

% Make tables not quite so cramped. Allow |\hline\hline| to give a
% poor man's thick rule. 
%    \begin{macrocode}
\renewcommand{\arraystretch}{1.15}
\setlength{\doublerulesep}{0.4pt}
%    \end{macrocode}

% \begin{macro}{\tablecols}
% \begin{macro}{\@tablecols}
% \begin{macro}{\tabcaptioncontinued}
% |\tablecols| is a user command for setting the number of columns. It
% is used by commands that try and span the entire table (I think
% there must be a better way of acheiving
% this). |\tabcaptioncontinued| is 
% a user command for setting the caption of a continued table on the
% next page. 
%    \begin{macrocode}
\newcommand{\tablecols}[1]{\renewcommand{\@tablecols}{#1}}
\newcommand{\@tablecols}{1}
\newcommand{\tabcaptioncontinued}{%
  \multicolumn{\@tablecols}{c}{\addtocounter{table}{-1}TABLE~\thetable\ (CONTINUED)} \\ \\
  }
%    \end{macrocode}
% \end{macro}
% \end{macro}
% \end{macro}

% \subsubsection{Footnotes in tables}
% \label{sec:table-footnotes}
% \begin{macro}{\tablenotemark}
% \begin{macro}{\tablenote}
% \begin{macro}{\tablenoteskip}
% Old user commands |\tablenotemark{MARK}|,
% |\tablenote{MARK}{TEXT}|---now deprecated. 
%    \begin{macrocode}
\newcommand\tablenotemark[1]{$^\mathrm{#1}$} 
\newcommand\tablenote[2]{\if@proceedings\else\small\fi ${}^\mathrm{#1}$ #2} 
\newcommand\tablenoteskip{\\[0.5\baselineskip]} 
%    \end{macrocode}
% \end{macro}
% \end{macro}
% \end{macro}

% \begin{macro}{\tabnotemark}
% \begin{macro}{\tabnotetext}
% \begin{macro}{\tabnotetextTWO}
% Alternative way of doing footnotes in tables. 
%    \begin{macrocode}
\newcommand{\tabnotemark}[1]{\textsuperscript{\lowercase{#1}}}
%\newcommand{\tabnotetextTWO}[4]{\tabnotetext{#1}{#2}\tabnotetext{#3}{#4}}
\newcommand{\tabnotetextTWO}[4]{%
  \multicolumn{\@tablecols}{l}{%
    \parbox{0.5\tabnotewidth}{\small\textsuperscript{#1}#2}%
    \parbox{0.5\tabnotewidth}{\small\textsuperscript{#3}#4}%
    }%
  \\}
\newcommand{\tabnotetext}[2]{%
  \multicolumn{\@tablecols}{@{}p{\tabnotewidth}@{}}{\small\textsuperscript{#1}#2}\\}
\newlength{\tabnotewidth}
%    \end{macrocode}
% \end{macro}
% \end{macro}
% \end{macro}


% \begin{macro}{changemargin}
%    \begin{macrocode}
% Command to allow margins to be changed on the fly. Useful to allow
% tables and Figures to protrude into the margins. 
%
% \changes{3.24}{2007/09/14}{Added \texttt{changemargin} environment
% (copied from TeX FAQ).}
\newenvironment{changemargin}[2]{% 
  % From TeX FAQ http://www.tex.ac.uk/cgi-bin/texfaq2html?label=chngmargonfly
  \begin{list}{}{%
    \setlength{\topsep}{0pt}%
    \setlength{\leftmargin}{#1}%
    \setlength{\rightmargin}{#2}%
    \setlength{\listparindent}{\parindent}%
    \setlength{\itemindent}{\parindent}%
    \setlength{\parsep}{\parskip}%
  }%
  \item[]}{\end{list}}
%    \end{macrocode}
% \end{macro}

% \subsection{The bibliography and references}
% \label{sec:biblio}
% \begin{macro}{\beforebiblioskip}
% \begin{macro}{\afterbiblioskip}
% I'm not entirely sure why I introduced these two commands, since
% directly putting |\vspace| before and after
% |\begin{thebibliography}| seems to work well enough. 
%    \begin{macrocode}
\newlength{\beforebiblioskip}
\setlength{\beforebiblioskip}{0pt}
\newlength{\afterbiblioskip}
\setlength{\afterbiblioskip}{0pt}
%    \end{macrocode}
% \end{macro}
% \end{macro}
% \begin{macro}{thebibliography}
% The environment itself is done as starred section with extra
% vertical space after heading. This space is currently zero. 
% In the book style we do the cunning trick with rereading the file in
% order to get the references into the |\ADS@references| macro.
% 
% \changes{3.05}{2003/01/02}{Removed the \cs{beforebiblioskip} and
% \cs{afterbiblioskip} commands, since they cause unwanted vertical
% space in some instances.}  
%    \begin{macrocode}
\renewenvironment{thebibliography}{%
  \if@book\RereadBiblio\fi
%%  \vspace*{\beforebiblioskip}
  \section*{REFERENCES} \frenchspacing
%%  \vspace*{\afterbiblioskip}
  \small
%    \end{macrocode} 
% The references themselves are set as a list in order 
% to get the hanging indentation. They are done in a smaller font and
% with french spacing so we needn't worry about the distinction
% between sentence-ending periods and abbreviations. 
%    \begin{macrocode} 
  \list{\null}{\leftmargin 1.5em\labelwidth\z@\labelsep\z@\itemindent -1.5em
    \itemsep\z@\parsep\z@\usecounter{enumi}}
  \sloppy\clubpenalty4000\widowpenalty4000
  \sfcode`\.=1000\relax\@biblisttrue}
{\endlist\@biblistfalse\normalsize\@outputaddressesifnecessary}
%    \end{macrocode}
% \end{macro}
% \subsection{Uniquification of cross-references}
% \label{sec:uniquification}
% 
% \begin{macro}{\articlename}
% In order that the |\label{}| command in individual articles be
% guaranteed to give a unique anchor, we append the article name (held
% in the |\articlename| command) to each key. This requires
% redefinition of the following commands: |\cite|, |\nocite|, |\ref|,
% |\pageref|, |\label|. 
% \begin{macro}{\oldcite}
% \begin{macro}{\oldnocite}
% \begin{macro}{\oldref}
% \begin{macro}{\oldpageref}
% \begin{macro}{\oldlabel}
% The original versions of these commands are
% saved as |\oldcite|, |\oldnocite|, etc. 
%
% The uniquifying of cross-references is now always done, so 
% we don't get multiply-defined labels in the single article case
%    \begin{macrocode}
\newcommand{\articlename}{anon}
\AtBeginDocument{
  \let\oldcite=\cite 
  \let\oldnocite=\nocite 
  \let\oldref=\ref 
  \let\oldpageref=\pageref
  \let\oldlabel=\label
}
%    \end{macrocode}
% \end{macro}
% \end{macro}
% \end{macro}
% \end{macro}
% \end{macro}
% \end{macro}
%
% Uniquify all cross-references by prepending the article name.  This
% now has to go into |\AtBeginDocument| so it is working on the new
% versions that hyperref.sty has set up.  Then we need to put the
% |\@begofarticlehook| afterwards, so that the first page label is using
% the new definitions.  If amsmath is loaded, then we also need to patch
% the special label command that they use. \changes{4.0}{2019/11/18}{Fix labels for amsmath display equations, in case that package has been loaded.}
%    \begin{macrocode}
\AtBeginDocument{
  \renewcommand{\ref}[1]{\oldref{\articlename:#1}%
    \if@debug{\ClassInfo{rmaa}{Reference to \articlename:#1}}\fi}
  \renewcommand{\pageref}[1]{\oldpageref{\articlename:#1}%
    \if@debug{\ClassInfo{rmaa}{Page reference to \articlename:#1}}\fi}
  \renewcommand{\label}[1]{\oldlabel{\articlename:#1}%
    \if@debug{\ClassInfo{rmaa}{Label for \articlename:#1}}\fi}%
  \ifdef{\label@in@display}{%
    \let\oldlabel@in@display=\label@in@display
    \renewcommand{\label@in@display}[1]{\oldlabel@in@display{\articlename:#1}%
      \if@debug{\ClassInfo{rmaa}{AMSTeX Label for \articlename:#1}}\fi}%
  }{}
  \@begofarticlehook
}%
%    \end{macrocode}
% Citations are more complicated since they can contain a list of
% keys. These are really only here for use with the rmaacite package,
% which is utterly deprecated (02 Jan 2003) now that I have the natbib
% code working. The command \cs{cite} should \emph{not} be used with
% natbib: use \cs{citep}, \cs{citet}, etc.\@ instead. 
%    \begin{macrocode}
\renewcommand{\nocite}[1]{\@bsphack
  \@for\@citeb:=#1\do{%
    \edef\@citeb{\expandafter\@firstofone\@citeb:\articlename}%
    \if@filesw\immediate\write\@auxout{\string\citation{\@citeb}}\fi
    \@ifundefined{b@\@citeb:\articlename}{\G@refundefinedtrue
      \@latex@warning{Citation `\@citeb' undefined}}{}}%
  \@esphack}
\renewcommand{\cite}[1]{\oldcite{#1:\articlename}}
%   \end{macrocode}


% Horrible fiddling with the bibliography citation mechanism. 
% This has to be globally available in case rmaacite puts it in the
% \texttt{.aux} file. 
%    \begin{macrocode}
\def\lbibcite#1#2#3{\global\@namedef{b@#1}{#2}
                    \global\@namedef{by@#1}{#3}}
%    \end{macrocode}

% \begin{macro}{\oldbibitem}
% \begin{macro}{\@rmaaBibArg}
% \begin{macro}{\@rmaaSpace}
% \begin{macro}{rmaabib}
% Of course, |\bibitem| needs to get the uniquification treatment as
% well. While we are at it, we redefine it to use consecutive integers
% for the key if the user leaves it blank (using the counter
% |rmaabib|). This allows one to put 
% |\bibitem{}| for all your references without getting annoying
% messages about multiply-defined labels.
% 
% We also take the opportunity to insert the optional argument to
% \cs{bibitem} if it is missing in the form Author(DATE) so that natbib
% doesn't try to do numerical citations. Natbib already provides a
% hook \cs{@extra@b@citeb} to help uniquify the citation key, so we
% just plug into that. However, it puts it \emph{after} the original
% key, not before as I had had it. Easiest solution is to change the
% way that I do it. 02 Jan 2003. Note that the \cs{@extra@b@citeb}
% does not get written to the .aux  file in the \cs{citation} command
% but it does work. 
% \changes{3.14}{06 Feb 2003}{Reinstalled old definition of \cs{bibitem}
%  for the case where the nonatbib option is requested, so that the
%  rmaacite package will work. This means that you can't use rmaacite
%  and natbib in the same volume.}
%    \begin{macrocode}
\if@natbib\def\@extra@b@citeb{:\articlename}\fi
%%\else
\let\oldbibitem\bibitem
\newcommand{\@rmaaBibArg}{}
\newcommand{\@rmaaSpace}{ }
\newcounter{rmaabib}
\if@natbib
\renewcommand{\bibitem}[2][ABC(2003)]{%
  \renewcommand{\@rmaaBibArg}{#2}
  \ifx\@rmaaBibArg\@rmaaEmpty
     \addtocounter{rmaabib}{1}
     \oldbibitem[#1]{\thermaabib:\articlename}
  \else\ifx\@rmaaBibArg\@rmaaSpace
     \addtocounter{rmaabib}{1}
     \oldbibitem[#1]{\thermaabib:\articlename}
  \else 
     \oldbibitem[#1]{#2:\articlename}
  \fi\fi
  \renewcommand{\@rmaaBibArg}{}
  }
\else
\renewcommand{\bibitem}[1]{%
  \renewcommand{\@rmaaBibArg}{#1}
  \ifx\@rmaaBibArg\@rmaaEmpty
     \addtocounter{rmaabib}{1}
     \oldbibitem{\articlename:\thermaabib}
  \else\ifx\@rmaaBibArg\@rmaaSpace
     \addtocounter{rmaabib}{1}
     \oldbibitem{\articlename:\thermaabib}
  \else
     \oldbibitem{\articlename:#1}
  \fi\fi
  \renewcommand{\@rmaaBibArg}{}
  }
\fi
%%\fi
%    \end{macrocode}
% \end{macro}
% \end{macro}
% \end{macro}
% \end{macro}

% \begin{macro}{oldpcite}
% \begin{macro}{oldscite}
% Then, we need to repeat the above scheme, but for the |rmaacite|
% definition. 
%    \begin{macrocode}
\newcommand{\@redefinebibitem}{%
  \let\oldbibitem\bibitem
  \let\oldcite\cite
  \let\oldpcite\pcite
  \let\oldscite\scite
  \def\bibitem[##1<##2>]##3{%
    \renewcommand{\@rmaaBibArg}{##3}
    \ifx\@rmaaBibArg\@rmaaEmpty
       \addtocounter{rmaabib}{1}
       \oldbibitem[##1<##2>]{\thermaabib:\articlename}
    \else\ifx\@rmaaBibArg\@rmaaSpace
       \addtocounter{rmaabib}{1}
       \oldbibitem[##1<##2>]{\thermaabib:\articlename}
    \else 
       \oldbibitem[##1<##2>]{##3:\articlename}
    % trust the author here
%%       \oldbibitem[##1<##2>]{##3}
    \fi\fi
    }
  }
%    \end{macrocode}
% \end{macro}
% \end{macro}


% \subsection{HTML Table of Contents}
% \label{sec:hyper-TOC}
% \begin{macro}{\HyperTOCEntry}
% \begin{macro}{\HyperTOCpart}
% \begin{macro}{\HyperTOCmeeting}
% \begin{macro}{\HyperTOCeditors}
% \begin{macro}{\HyperTOCpreamble}
% These commands are used to write information to the \texttt{.htoc}
% file, which is used in generating a Web Page of the Table of
% Contents. Since they only ever appear preceded by |\protect| in the
% argument to an |\addtocontents| command,  they should never need to
% be expanded by the macros. However, we define them here just in case
% \dots
%    \begin{macrocode}
\newcommand{\HyperTOCEntry}[3]{\relax}
\newcommand{\HyperTOCpart}[1]{\relax}
\newcommand{\HyperTOCmeeting}[3]{\relax}
\newcommand{\HyperTOCeditors}[1]{\relax}
\newcommand{\HyperTOCpreamble}{\relax}
%    \end{macrocode}
% \end{macro}
% \end{macro}
% \end{macro}
% \end{macro}
% \end{macro}
%

% \subsection{Formatting the poster abstracts} 
% These can now be used without using the book option.
%    \begin{macrocode}

\newcommand{\@makeminititle}{%
  \newpage
  \null 
  \vskip 3em%
  \begin{center}%
    {\@title \par}%
  \end{center}
  \vskip 3em%
  }

\newif\if@abstractsads
\@abstractsadsfalse
\newcommand\abstractsADS{\@abstractsadstrue}

\if@book\RequirePackage{multicol}\fi
\newenvironment{abstracts}[1][Abstracts of Contributed Papers]{%
  %% First, make sure we start on an odd page (06 Jun 2000)
  \if@book\clearpage\thispagestyle{empty}\cleardoublepage\fi
  %%
  \renewcommand{\articlename}{ABSTRACTS}
  \title{#1}\shortauthor{Abstracts}\shorttitle{Abstracts}
  \listofauthors{}
  \if@book
  \addcontentsline{toc}{chapter}{\protect\@fiddleauthors{\@TOCtitle}{\rmaaTOCauthors}}%
  \addtocontents{htc}{\protect\HyperTOCEntry{\@rawtitle}{\@rawauthors}{\articlename}}
  \fi
  \pagestyle{rmaa} % this now seems necessary (06 Jun 2000)
  \thispagestyle{rmaafirstpage}
%%  \twocolumn[\@makeminititle]
  \renewcommand{\abstract}[1]{%
    \renewcommand\ADS@abstract{##1}%
    \renewcommand{\@abstract}{##1}\@writeabstract%
  }
  \renewcommand{\title}[2][]{%
    \gdef\ADS@authors{}%
    \gdef\ADS@affil{}%
    \gdef\ADS@references{}%
    \gdef\ADS@keywords{}%
    \renewcommand{\@rmaaArg}{##1}
    \ifx\@rmaaArg\@rmaaEmpty
       \renewcommand{\@TOCaftertitle}{}
    \else
       \renewcommand{\@TOCaftertitle}{ (\@rmaaArg)}
    \fi
    \renewcommand{\@title}{\rmaa@TOCTitleStyle{##2}}%
    \renewcommand{\@TOCtitle}{\rmaa@TOCTitleStyle{##2}\@TOCaftertitle}%
    \renewcommand{\rmaaIndexTitle}{\rmaa@IndexTitleStyle{##2\@TOCaftertitle.}}%
    \renewcommand\ADS@title{##2}%
    \xdef\ADS@firstpage{\@arabic{\c@page}}%
    }
  \renewcommand{\author}[1]{\let\@thanks\@empty\renewcommand\@author{##1} }
  \renewcommand{\footnotetext}[2][XX]{\par\noindent\textsuperscript{##1}~##2}
  \renewcommand{\footnoterule}{\kern-3\p@ \hrule \@width 3em \kern 2.6\p@}
  \if@book\@begofarticlehook\fi
  \if@book\begin{multicols}{2}[\@makeminititle]\else\@makeminititle\fi
  \begingroup
  \setcounter{footnote}{0}
  }% and at the end of the environment....
  {\endgroup\if@book\end{multicols}\listofabstracts\label{RMAALastPage}\clearpage\fi}
%% auxiliary commands used:
%% \@writeabstract does all the actual formatting 
\newcommand{\@writeabstract}{%
  \if@book%
  \addcontentsline{loa}{abstract}{\protect\@fiddleauthors{\@TOCtitle}{\rmaaTOCauthors}}%
  \fi%
  \begin{center}\@title\par\@author\end{center}
  \@abstract
  \par\noindent\footnoterule\nobreak\footnotesize\@thanks
  \par\noindent\hrulefill\goodbreak
  \vskip 1cm \@plus0.5cm
  %% write ADS entry
  \if@book\if@abstractsads\ADS@dumpdata\fi\fi%
  \endgroup % this pair is to make each abstract a self-contained group
  \begingroup
  \setcounter{footnote}{0}
  }

\if@book
%% List of abstracts (stuff written to .loa file)
\newcommand{\listofabstracts}{%
  \begin{multicols}{2}[%
    \vskip 1.5cm \vfill % leave at least 1.5 cm and go to bottom of page if possible
    \centerline{\MakeUpperCase{List of Abstracts}}
    \vskip 1.5cm \@plus0.5cm
    ][6cm] % this makes it start a new page if there is not at least 6cm of space
    \renewcommand{\@fiddleauthors}[2]{##1\hskip1em plus2em minus0.4em\penalty-100  ##2}
    \parindent\z@
%%    \parskip.5\p@ \@plus .7\p@\relax 
    \parskip3\p@ \@plus 1\p@ \@minus 1\p@\relax % allow quite a lot of
                                % leeway here to avoid breaks inside
                                % entries 
    \renewcommand{\\}{\relax} % temporarily turn off explicit line breaks just in case
    \renewcommand{\altaffilmark}[1]{\relax} % ditto with this
    \@starttoc{loa}
  \end{multicols}
  }
\newcommand*{\l@abstract}[2]{%
  \par\penalty-1000
  #1\hskip 1em\penalty0\null\nobreak\hfill#2}
\fi

%    \end{macrocode}
 
% \subsection{Code for typesetting an entire volume of conference
% proceedings} \label{sec:whole-volume}
%    \begin{macrocode}

\if@book
%% environment for including the individual articles
\newcommand{\@fiddleauthors}[2]{\relax}
\renewcommand{\theequation}{\@arabic\c@equation}
\newenvironment{articles}%
{ \begingroup % make it a group so that author definitions do not
              % spill over
  \gdef\ADS@authors{}
  \pagestyle{rmaa}
  \renewenvironment{document}{%
    \@begofarticlehook
    \addtocounter{chapter}{1}%
    \addcontentsline{toc}{chapter}{%
      \protect\@fiddleauthors{\@TOCtitle}{\rmaaTOCauthors}}%
    \addtocontents{htc}{\protect\HyperTOCEntry{\@rawtitle}{\@rawauthors}{\articlename}}
    }{\relax}
  \renewcommand{\documentclass}[2][foo]{\relax} 
  \renewcommand{\usepackage}[2][foo]{%
    \def\@packagename{##2}\def\@rmaacitename{rmaacite}%
    \ifx\@packagename\@rmaacitename
    \@rmcitetrue\ClassInfo{rmaa}{This paper uses rmaacite}
    \makeatletter%%
%% This is file `rmaacite.sty',
%% generated with the docstrip utility.
%%
%% The original source files were:
%%
%% rmaa.dtx  (with options: `rmaacite')
%% 
%% PLEASE DO NOT EDIT THIS FILE! It is generated automatically from
%% the file rmaa.dtx using the accompanying Makefile and rmaa.ins
%% Any hand-edited changes will be lost
%% 
%% A LaTeX document class for the RMxAA
%% 
%% Author: William Henney (w.henney@crya.unam.mx)
%% 
%% Latest version available from: http://www.crya.unam.mx/rmaa
%% 
%% Copyright (C) 1999-2010 Instituto de Astronomia, UNAM, Mexico
%% 
%% This program is free software; you can redistribute it and/or
%% modify it under the terms of the GNU General Public License
%% as published by the Free Software Foundation; either version 2
%% of the License, or (at your option) any later version.
%% 
%% This program is distributed in the hope that it will be useful,
%% but WITHOUT ANY WARRANTY; without even the implied warranty of
%% MERCHANTABILITY or FITNESS FOR A PARTICULAR PURPOSE.  See the
%% GNU General Public License for more details.
%% 
%% You should have received a copy of the GNU General Public License
%% along with this program; if not, write to the Free Software
%% Foundation, Inc., 59 Temple Place - Suite 330, Boston, MA  02111-1307, USA.






\def\bibitem[#1<#2>]#3{\par\hangindent=1.5em\hangafter=1
  \if@filesw
  {\def\protect##1{\string ##1\space}
    \immediate\write\@auxout{\string\lbibcite{#3}{#1}{#2}}}
  \fi}
\@redefinebibitem % what it says
\def\lbibcite#1#2#3{\global\@namedef{b@#1}{#2}
                    \global\@namedef{by@#1}{#3}}

\def\@bcite#1#2{({#1\if@tempswa , #2\fi})}
\def\@pcite#1#2{#1\if@tempswa , #2\fi}
\def\@brcite{\ (}
\def\@ebrcite{)}
\def\@fmttop#1{\@fmtcite[\csname #1\endcsname]}
\def\@citex[#1]#2{%
%%  \if@filesw\immediate\write\@auxout{\string\citation{#2}}\fi
%% WJH added 16 Jan 2002
  \@for\@citeb:=#2\do{%
    \edef\@citeb{\expandafter\@firstofone\@citeb:\articlename}%
    \if@filesw\immediate\write\@auxout{\string\citation{\@citeb}}\fi}%
%% to here
  \def\@citea{}%
  \@cite{\@for\@citeb:=#2\do
    {%
      \edef\@citeb{\expandafter\@firstofone\@citeb:\articlename}%
      \@citea\def\@citea{;\ }\@ifundefined
      {b@\@citeb}{{\bf ?}\@warning {Citation `\@citeb' on page \thepage \space undefined}}%
      {\csname b@\@citeb\endcsname\@citebro\csname by@\@citeb\endcsname\@citebrc}}}%
  {#1}}

\def\cite{\@ifnextchar [{\let\@citebro=\space
                          \let\@citebrc=\relax
                          \let\@cite=\@bcite\@tempswatrue\@citex}
                        {\let\@citebro=\space
                          \let\@citebrc=\relax
                          \let\@cite=\@bcite\@tempswafalse\@citex[]}}
\def\pcite{\@ifnextchar [{\let\@citebro=\space
                          \let\@citebrc=\relax
                          \let\@cite=\@pcite\@tempswatrue\@citex}
                        {\let\@citebro=\space
                          \let\@citebrc=\relax
                          \let\@cite=\@pcite\@tempswafalse\@citex[]}}
\def\scite{\@ifnextchar [{\let\@citebro=\@brcite
                          \let\@citebrc=\@ebrcite
                          \let\@cite=\@pcite\@tempswatrue\@citex}
                        {\let\@citebro=\@brcite
                          \let\@citebrc=\@ebrcite
                          \let\@cite=\@pcite\@tempswafalse\@citex[]}}

\renewenvironment{thebibliography}
  {\section*{REFERENCES}\vspace*{\afterbiblioskip}\small
    \if@book\RereadBiblio\fi
    \sloppy\clubpenalty4000\widowpenalty4000
    \sfcode`\.=1000\relax
    \setlength\parindent{\z@}\setlength\parskip{\z@}}
  {\relax\normalsize\@outputaddressesifnecessary}








\endinput
%%
%% End of file `rmaacite.sty'.
\makeatother % read it in here so that all changes
    \else                                         % are local to paper
    \if@debug\ClassWarning{rmaa}{Package \@packagename\space was requested by \articlename\space}\fi
    \fi} 
  %% Don't allow individual articles to mess with the pages, volume or year
  \renewcommand{\SetFirstPage}[1]{\relax}
  \renewcommand{\SetYear}[1]{\relax}
  \renewcommand{\SetVolume}[1]{\relax}
  \renewcommand{\SetConfTitle}[1]{\relax}
  %% Don't allow individual articles to change these
  \renewcommand{\nonstopmode}{\relax}
  \renewcommand{\scrollmode}{\relax}
  \renewcommand{\outputplates}{\relax}
  %% XXXX need to check if using rmaacite and act accordingly
  %% avoid multiple declarations of the same length
  \let\oldnewlength=\newlength
  \renewcommand{\newlength}[1]{%
    \edef\reserved@a{\expandafter\@gobble\string ##1}%
    \@ifundefined{\reserved@a}{%
      \oldnewlength{##1}}{%
      \if@debug{\ClassInfo{rmaa}{\reserved@a\space already defined: doing nothing}}\fi%
      \relax}}%
  \renewcommand\newcommand{\providecommand}
  }%
{\endgroup\onecolumn}
     

%    \end{macrocode}
% Command to be used in `articles' environment for including an
% article
%    \begin{macrocode}
\newcommand{\article}[2][]{%
  % optional first argument is extra text to put in brackets after title in TOC
  % e.g. ``Invited Talk'' or ``Poster''
  \renewcommand{\articlename}{#2}
  % reset loads of counters
  \setcounter{footnote}{0}
  \setcounter{figure}{0}
  \setcounter{table}{0}
  \setcounter{section}{0}
  \setcounter{equation}{0}
  \setcounter{rmaabib}{0}
  \setcounter{author}{0}
  \setcounter{AUTHOR}{1}
  \gdef\ADS@authors{}%
  \gdef\ADS@affil{}%
  \gdef\ADS@affiliations{}%
  \gdef\ADS@references{}%
  \@firstauthortrue
  \@firstaffiltrue
  \@firstaltaffiltrue
  \@firstkwdtrue
  \@rmcitefalse
  \if@suppressfulladdresses\else\@needaddressestrue\fi
  \renewcommand{\@rmaaArg}{#1}
  \ifx\@rmaaArg\@rmaaEmpty
     \renewcommand{\@TOCaftertitle}{}
  \else
     \renewcommand{\@TOCaftertitle}{ (#1)}
  \fi
%    \end{macrocode}
% In the main journal articles must start on an odd page. 
%    \begin{macrocode}
  \if@journal\clearpage\thispagestyle{empty}\cleardoublepage\fi
%% Now input the file, then run the end-of-article hook (reinstated 28
  %% Sep 2001)
  {\input{\@articledir/\articlename.tex}\@endofarticlehook}}
%    \end{macrocode}
% Command to separate articles into sessions. Write an entry in the
% TOC, makes sure we start on an odd page (added 06 Jun 2000), then do
% a title page if \cs{SessionTitlePages} was called. 
% \changes{3.12}{04 Feb 2003}{Added option to have title page for each
% session. Turn it on with \cs{SessionTitlePages}.}
%    \begin{macrocode}
\newcommand{\session}[1]{
  %% this does nothing except write an entry in the TOC
  \addcontentsline{toc}{part}{\rmaa@TOCPartStyle{#1}}
  \addtocontents{htc}{\protect\HyperTOCpart{#1}}
  %% and also makes sure we start on an odd page (added 06 Jun 2000)
  \clearpage\thispagestyle{empty}\cleardoublepage
  \stepcounter{part}
  \if@sessiontitlepages{%
    \onecolumn
    \thispagestyle{empty}
    \renewcommand{\articlename}{Section\thepart TitlePage}
    \label{RMAAFirstPage}
    \make@session@title{#1}
    \label{RMAALastPage}
    \clearpage\thispagestyle{empty}\cleardoublepage
  }\fi
  }
\newcommand{\make@session@title}[1]{%
  {%
    \setlength\parindent{0pt}
    \vspace*{\stretch{1}}
    \rule{\linewidth}{1mm}
    \begin{flushright}\Huge\bfseries\boldmath
      #1
    \end{flushright}
    \rule{\linewidth}{1mm}
    \vspace*{\stretch{2}}
  }
}
%    \end{macrocode}
% Command to turn on separate title pages for each session. 
%    \begin{macrocode}
\newcommand{\SessionTitlePages}{\@sessiontitlepagestrue}
\newcommand{\BoldSessionTitles}{%
  \renewcommand{\rmaa@TOCPartStyle}[1]{\textbf{\rmaa@TOCTitleStyle{##1}}}
  }
\newcommand{\ArticleDir}[1]{\renewcommand{\@articledir}{#1/}\graphicspath{{1/}}}
\newcommand{\@articledir}{.}
%\newcommand{\ps@rmaafront}{\ps@rmaa}% page style for front matter (same as ever, actually)
\newcommand{\ps@rmaafront}{% page style for front matter:  page # in each foot
  \renewcommand\@oddfoot{\hfil{\thepage}}%
  \renewcommand\@evenfoot{{\thepage}\hfil}%
  % shorttitle in both heads
  \renewcommand\@evenhead{\hfil\@shorttitle\hfil}%
  \renewcommand\@oddhead{\@evenhead}
  }
\newcommand{\ps@rmaafrontfirstpage}{% page style for first page of front matter
  % page # in each foot, on outside edge
  \renewcommand\@oddfoot{\hfil\thepage}%
  \renewcommand\@evenfoot{\thepage\hfil}%
  % nothing in the even heads
  \renewcommand\@evenhead{\relax}%
  \renewcommand\@oddhead{\relax}%
  }
%% slight modification to the standard \frontmatter and \mainmatter
%% commands, just to set the \pagestyle
\renewcommand\frontmatter{%
  \latexonecolumn
  \cleardoublepage
  \@frontmattertrue
  \@mainmatterfalse
  \pagenumbering{roman}
  \pagestyle{rmaafront}
  \addtocounter{page}{2} % leave space for the title page
  }
\renewcommand\mainmatter{%
  \pagestyle{empty}
  \if@rmtwocol\twocolumn\fi
  \cleardoublepage
  \@frontmatterfalse
  \@mainmattertrue
  \pagenumbering{arabic}
  \pagestyle{rmaa}
  }
\renewcommand\backmatter{%
  \clearpage\pagestyle{empty}\cleardoublepage
  \@mainmatterfalse
  %% put extra space in the TOC
  \addtocontents{toc}{\bigskip}
  }
%    \end{macrocode}
%
% \subsection{Front sections}\label{sec:front-sections}
% \begin{macro}{\notitle}
% \begin{macro}{\shortname}
% \begin{macro}{frontsection}
% The environment \texttt{frontsection} is for doing a section in the
% front matter (preface, list of participants, or whatever). It has 
% one optional argument:
% \begin{enumerate}
% \item[\#1] Short title (no embedded spaces) for
% use in making PS files. E.g., ``list\_of\_participants''.   
% \end{enumerate}
% and two mandatory arguments: 
% \begin{enumerate}
% \item[\#2] Title of section (e.g., `Preface')
% \item[\#3] Commands to be evaluated before processing the contents
% of the environment. This is mainly useful for putting in the things
% that would go in the preamble of a normal article  (e.g.,
% |\listofauthors|, |\indexauthor|, |\notitle|) 
% \end{enumerate}
%    \begin{macrocode}
\newcommand{\notitle}{\@titlevisiblefalse} \newcommand{\shortname}{}
\newenvironment{frontsection}[3][]{%
  \begingroup
  \renewcommand{\shortname}{#1}
  \ifx\shortname\@rmaaEmpty
     \renewcommand{\articlename}{#2}
  \else
     \renewcommand{\articlename}{\shortname}
  \fi
  \@begofarticlehook
  \@titlevisibletrue % default
  \thispagestyle{rmaafrontfirstpage}
  \listofauthors{} % blank out this in case it is not specified
  \title{{#2}}\shorttitle{#2}\shortauthor{#2}
  #3 % evaluate the command given as second mandatory arg
  \addcontentsline{toc}{chapter}{\protect\@fiddleauthors{\@TOCtitle}{\rmaaTOCauthors}}%
  \addtocontents{htc}{\protect\HyperTOCEntry{\@rawtitle}{\@rawauthors}{\articlename}}
  \if@titlevisible
     \vspace*{1cm}\centerline{\rmaa@FrontSecStyle{\@title}}\vspace*{1cm}
  \fi
  \setlength{\parskip}{3ex}
  }{\frontlabel{RMAALastPage}\clearpage\endgroup}
%    \end{macrocode}
% \end{macro}
% \end{macro}
% \end{macro}
%    \begin{macrocode}
%% Definition of the command \frontlabel, which rights out something like `F4' to
%% the .aux file, instead of `\textlatin  {i}'. This is so it is easy to post-process
%% the file for making PS files of individual sections
\newcommand{\frontlabel}[1]{%
  \@bsphack
  \protected@write\@auxout{}%
     {\string\newlabel{\articlename:#1}{{}{F\@arabic{\c@page}}}}
  \@esphack
  }
%% User command to say X, Y, Z,\\ Editors of the Proceedings
\newcommand{\EditorSignature}{\vspace*{\parskip}\rightline{\@editors}%
  \rightline{Editors of the proceedings}}
%% Fiddle with \thesection, etc. so that they don't include the chapter number
%% Basically go back to the `article' definitions of these
\renewcommand{\thesection}{\@arabic\c@section}
\renewcommand{\thefigure}{\@arabic\c@figure}
\renewcommand{\thetable}{\@arabic\c@table}

\newcommand{\@oldargs}{}
\newcommand{\@newargs}{}
%% add TOC entry for each paper
%\renewcommand{\title}[1]%
%{\renewcommand{\@title}{\rmaa@TitleStyle{#1}}%
%\addtocounter{chapter}{1}\addcontentsline{toc}{chapter}{#1}%
%}

%% Formatting of the table of contents
%% dimensions of various parts
\newlength{\TOCtextwidth}\setlength{\TOCtextwidth}{12.65cm}
\newlength{\TOCpagenumwidth}\setlength{\TOCpagenumwidth}{1.3cm}
\newlength{\TOCentrywidth}\setlength{\TOCentrywidth}{\TOCtextwidth}
\addtolength{\TOCentrywidth}{-\TOCpagenumwidth} 
\newlength{\TOCextramargin}\setlength{\TOCextramargin}{1.25cm}
\newcommand{\TOCbox}[1]{\hspace{\TOCextramargin}\parbox{\TOCtextwidth}{#1}\par}
\newlength{\TOCextramarginx}\setlength{\TOCextramarginx}{\TOCextramargin}\addtolength{\TOCextramarginx}{4.3cm}
\newlength{\TOCtextwidthx}\setlength{\TOCtextwidthx}{\TOCtextwidth}\addtolength{\TOCtextwidthx}{-4.3cm}
\newcommand{\TOCboxx}[1]{\hspace{\TOCextramarginx}\parbox{\TOCtextwidthx}{#1}\par}

\newcommand{\RMAAbannerimage}{rmsc}
\renewcommand\tableofcontents{%
  \renewcommand{\articlename}{CONTENTS}
  \begingroup
  \@begofarticlehook
  \if@twocolumn
     \@restonecoltrue\onecolumn
  \else
     \@restonecolfalse
  \fi
  %% 9pt for journal style
  \if@proceedings\else\small\fi
  \if@newstyle\small\fi %% same for new style
  \if@proceedings\thispagestyle{rmaafrontfirstpage}\else\thispagestyle{empty}\fi
  \shortauthor{Contents}\shorttitle{Contents}
  \vspace*{18mm}
  \if@printbanner
  %% put in the banner header (need to get better image of this)
     \TOCbox{\hspace*{-0.15cm}\includegraphics[width=0.9\TOCtextwidth]{\RMAAbannerimage}}
     \vspace*{8mm}
  %% volume and publication date (should really be in funny font)
     \TOCboxx{\@volumen\\[2mm] \@mes\ \@anyo}
  \else
  %% put in the banner header (need to get better image of this)
     \TOCbox{\hspace*{-0.15cm}\includegraphics[width=0.9\TOCtextwidth]{\RMAAbannerimage}}
     \vspace*{12.5mm}
  %% volume and publication date (should really be in funny font)
     \TOCboxx{~\\[2mm] ~}
  \fi
  \vspace*{16mm}
  %% This is the command that fiddles with the TOC entry so that the
  %% title is flushleft but the authors are flushright
  \renewcommand{\@fiddleauthors}[2]{%
    % first put the paper title
    \renewcommand{\\}{\relax} % temporarily turn off explicit line breaks just in case
    \renewcommand{\altaffilmark}[1]{\relax} % ditto with this
    \renewcommand{\textsuperscript}[1]{\relax} % ditto with this 
    ##1
    % now save the author list in a box
    \sbox\@tempboxa{##2}%
    % check that box is shorter than the total width available for the entry
    \ifdim \wd\@tempboxa <\TOCentrywidth
       % if so, then put it right here, flush against right margin
       \penalty10000\hskip1em plus10em\penalty0\hbox{}\nobreak\hfill\box\@tempboxa
       % Note that we force a space of at least 1em between title and authorlist, 
       % stretchable to 11em since we may have very short titles but still with an 
       % authorlist too long to go on the same line.  The only linebreak allowed is 
       % directly before the authorlist.
    \else
       % if it is too long, we set the original authorlist text in its own parbox
       \par\vskip1ex\parbox[b]{\TOCentrywidth}{\raggedleft ##2}
    \fi
    }
  \@starttoc{toc}%
  \if@restonecol\twocolumn\fi
  \frontlabel{RMAALastPage}
  \clearpage
  \thispagestyle{empty}
  \cleardoublepage % make sure we start on an odd page
  \endgroup

  %%% ALSO, do the HTML TOC
  %% The file that will be used to make the HTML TOC
  \@starttoc{htc}
  %% All the stuff in this file does nothing in the book style
  \addtocontents{htc}{\protect\HyperTOCmeeting{\@conftitle}{\@reunionplace}{\@reuniondate}}
  \addtocontents{htc}{\protect\HyperTOCeditors{\@editors}}
  \addtocontents{htc}{\protect\HyperTOCpreamble}
  \addtocontents{htc}{\protect\HyperTOCEntry{Table of Contents}{}{\articlename}}

  %% and the ADS file
  \ADS@openfile
  }
%% This is now much simpler
\renewcommand\l@chapter[2]{%TOC entry for articles
  \ifnum \c@tocdepth >\m@ne
    \vskip 1.0em \@plus\p@
    \TOCbox{%
      \parbox[b]{\TOCentrywidth}{#1}\parbox[b]{\TOCpagenumwidth}{\hfill#2}
    }
  \fi} 
\renewcommand\l@part[2]{%TOC entry for sessions: centered, no page number
  \ifnum \c@tocdepth >\m@ne
    \vskip 1.5em \@plus1.5em \goodbreak
    \TOCbox{%
      \parbox[b]{\TOCentrywidth}{\hfil#1\hfil}\hfill
    }
  \fi\nobreak} 

%% Formatting theindex
%% strongly encourage page breaks between index entries
\renewcommand{\@idxitem}{\par\penalty-1000\hangindent.45cm} 
\renewcommand{\indexspace}{\par\penalty-1000\vskip 10\p@ \@plus15\p@ \@minus3\p@\relax}
\renewenvironment{theindex}{%
  %% first, make sure we start on an odd page (06 Jun 2000)
  \clearpage\thispagestyle{empty}\cleardoublepage
  %% this looks a bit like an article but is two columns and without abstract, etc
  \renewcommand{\articlename}{INDEX}
  \title{Author Index}\shortauthor{Author Index}\shorttitle{Author Index}
  \listofauthors{}
  \addcontentsline{toc}{chapter}{\protect\@fiddleauthors{\@TOCtitle}{\rmaaTOCauthors}}%
  \addtocontents{htc}{\protect\HyperTOCEntry{\@rawtitle}{\@rawauthors}{\articlename}}
  \pagestyle{rmaa}
  \thispagestyle{rmaafirstpage}
  \@begofarticlehook
  %% 9pt for journal style
  \if@proceedings\else\small\fi
  \if@newstyle\small\fi
  \twocolumn[\@makeminititle]
  %% vertical spacing (see also definition of \@idxitem and \indexspace
  \parskip\z@ \@plus .6\p@\relax % allow quite a lot of leaway here to avoid breaks inside entries
  %% horizontal spacing, breaking and glue parameters (see also the stuff in the definition 
  %% of \indexauthor
  \parindent\z@ 
  \exhyphenpenalty10000 % don't break on explicit hyphens, it always looks bad
  \rightskip=0pt plus.2em % allow a tiny bit of leeway at the ends of lines
  \spaceskip.3333em plus.5em minus.15em % a lot of leeway between words
  \xspaceskip.5em plus.5em minus.4em % this seems to have no effect
  \let\item\@idxitem
  \renewcommand{\\}{\relax} % temporarily turn off explicit line breaks just in case
  \renewcommand{\altaffilmark}[1]{\relax} % ditto with this
  \renewcommand{\textsuperscript}[1]{\relax} % ditto with this
  }{\label{RMAALastPage}\clearpage} % clearpage necessary so that we know about \shortauthor



%% end definitions for book styles
\fi

%    \end{macrocode}
% In order to do away with plates, we sometimes want a figure to
% appear in color only in the online version. The following command
% can bve used in a figure caption to typeset a note to this effect. 
%    \begin{macrocode}
\newcommand\ColorOnlineOnly{%
  \textsc{Note: this figure is available in color in the electronic
    version of this article, obtainable from}
  \texttt{http://\discretionary{}{}{}www.\discretionary{}{}{}astroscu.\discretionary{}{}{}unam.\discretionary{}{}{}mx/\discretionary{}{}{}\~{}rmaa/}.}
%% Plates
%% Define 2 counters: plate is the number in the whole document, plateloop is used 
%% internally by \outputplates
\newcounter{plate}
\newcounter{plateloop}
%% Currently limited to 15 plates in document
%% when we do a label, we need to write to \newlabel entries to the .aux file, one 
%% that says RMAAPlateX so that it can be referred to in \ouputplates, and one that
%% says articlename:RMAAPlateX so that the dvips program can find it and asociate it
%% with the right article
\newcommand{\@platelabeli}{\oldlabel{RMAAPlate1}\label{RMAAPlate1}}
\newcommand{\@platelabelii}{\oldlabel{RMAAPlate2}\label{RMAAPlate2}}
\newcommand{\@platelabeliii}{\oldlabel{RMAAPlate3}\label{RMAAPlate3}}
\newcommand{\@platelabeliv}{\oldlabel{RMAAPlate4}\label{RMAAPlate4}}
\newcommand{\@platelabelv}{\oldlabel{RMAAPlate5}\label{RMAAPlate5}}
\newcommand{\@platelabelvi}{\oldlabel{RMAAPlate6}\label{RMAAPlate6}}
\newcommand{\@platelabelvii}{\oldlabel{RMAAPlate7}\label{RMAAPlate7}}
\newcommand{\@platelabelviii}{\oldlabel{RMAAPlate8}\label{RMAAPlate8}}
\newcommand{\@platelabelix}{\oldlabel{RMAAPlate9}\label{RMAAPlate9}}
\newcommand{\@platelabelx}{\oldlabel{RMAAPlate10}\label{RMAAPlate10}}
\newcommand{\@platelabelxi}{\oldlabel{RMAAPlate11}\label{RMAAPlate11}}
\newcommand{\@platelabelxii}{\oldlabel{RMAAPlate12}\label{RMAAPlate12}}
\newcommand{\@platelabelxiii}{\oldlabel{RMAAPlate13}\label{RMAAPlate13}}
\newcommand{\@platelabelxiv}{\oldlabel{RMAAPlate14}\label{RMAAPlate14}}
\newcommand{\@platelabelxv}{\oldlabel{RMAAPlate15}\label{RMAAPlate15}}
\newcommand{\@outputplatelabeli}{\oldlabel{RMAAOutputPlate1}}
\newcommand{\@outputplatelabelii}{\oldlabel{RMAAOutputPlate2}}
\newcommand{\@outputplatelabeliii}{\oldlabel{RMAAOutputPlate3}}
\newcommand{\@outputplatelabeliv}{\oldlabel{RMAAOutputPlate4}}
\newcommand{\@outputplatelabelv}{\oldlabel{RMAAOutputPlate5}}
\newcommand{\@outputplatelabelvi}{\oldlabel{RMAAOutputPlate6}}
\newcommand{\@outputplatelabelvii}{\oldlabel{RMAAOutputPlate7}}
\newcommand{\@outputplatelabelviii}{\oldlabel{RMAAOutputPlate8}}
\newcommand{\@outputplatelabelix}{\oldlabel{RMAAOutputPlate9}}
\newcommand{\@outputplatelabelx}{\oldlabel{RMAAOutputPlate10}}
\newcommand{\@outputplatelabelxi}{\oldlabel{RMAAOutputPlate11}}
\newcommand{\@outputplatelabelxii}{\oldlabel{RMAAOutputPlate12}}
\newcommand{\@outputplatelabelxiii}{\oldlabel{RMAAOutputPlate13}}
\newcommand{\@outputplatelabelxiv}{\oldlabel{RMAAOutputPlate14}}
\newcommand{\@outputplatelabelxv}{\oldlabel{RMAAOutputPlate15}}
\newcommand{\@platerefi}{\oldpageref{RMAAPlate1}}
\newcommand{\@platerefii}{\oldpageref{RMAAPlate2}}
\newcommand{\@platerefiii}{\oldpageref{RMAAPlate3}}
\newcommand{\@platerefiv}{\oldpageref{RMAAPlate4}}
\newcommand{\@platerefv}{\oldpageref{RMAAPlate5}}
\newcommand{\@platerefvi}{\oldpageref{RMAAPlate6}}
\newcommand{\@platerefvii}{\oldpageref{RMAAPlate7}}
\newcommand{\@platerefviii}{\oldpageref{RMAAPlate8}}
\newcommand{\@platerefix}{\oldpageref{RMAAPlate9}}
\newcommand{\@platerefx}{\oldpageref{RMAAPlate10}}
\newcommand{\@platerefxi}{\oldpageref{RMAAPlate11}}
\newcommand{\@platerefxii}{\oldpageref{RMAAPlate12}}
\newcommand{\@platerefxiii}{\oldpageref{RMAAPlate13}}
\newcommand{\@platerefxiv}{\oldpageref{RMAAPlate14}}
\newcommand{\@platerefxv}{\oldpageref{RMAAPlate15}}
\newcounter{figplatei}
\newcounter{figplateii}
\newcounter{figplateiii}
\newcounter{figplateiv}
\newcounter{figplatev}
\newcounter{figplatevi}
\newcounter{figplatevii}
\newcounter{figplateviii}
\newcounter{figplateix}
\newcounter{figplatex}
\newcounter{figplatexi}
\newcounter{figplatexii}
\newcounter{figplatexiii}
\newcounter{figplatexiv}
\newcounter{figplatexv}
%% Printed output should use the counter `plate'
\renewcommand{\theplate}{\@arabic\c@plate}
\renewcommand{\theplateloop}{\@arabic\c@plateloop}
%% Command to output all the saved plates in this paper
\newcommand{\outputplates}{%
  \setcounter{plateloop}{0}
  \if@book\else\@endofarticlehook\fi % not entirely sure why, but this
                                     % seems to be needed
  \clearpage % first output the last page of text
  % now set pagestyle for blank even pages and make sure we start on an odd page
  \pagestyle{empty}
  \cleardoublepage
  \onecolumn
  \parindent\z@
  % first put in an extra sheet
  % the main purpose of this is so there is always a blank even page before 
  % each plate for filling in the PS file if the article has an odd number
  % of pages (see splitbook.sh)
  \null\vfil\centerline{PLATES}\vfil\clearpage\pagestyle{rmaaplate}\cleardoublepage
  % output each plate in turn
  \@whilenum\value{plateloop}<\value{plate}\do{%
    \addtocounter{plateloop}{1}
    % restore the figure counter from \@figcounti, \@figcountii, etc
    \setcounter{figure}{\value{figplate\roman{plateloop}}}
    % spew out the plate (in its own group for safety)
    {\csname @plate\roman{plateloop}\endcsname\par}
    }
  \pagestyle{rmaa}
  % reset the page number as if nothing had happened
  \addtocounter{page}{-\value{plate}}
  \addtocounter{page}{-\value{plate}}
  }
\newcommand\figcaption{%
  \def\@captype{figure}
  \renewcommand{\@makecaption}{\@makefigcap}
  \caption}
\newcommand\tabcaption{%
  \def\@captype{table}
  \renewcommand{\@makecaption}{\@maketabcap}
  \caption}

\newcommand{\plate}[2][1]{%
  %% User command to make a plate page: 
  %% Mandatory argument is what goes on the plate - this should not 
  %% include any floating environments (e.g. figure). Instead, \figcaption 
  %% should be used to do the captions. 
  %% Optional argument is the number of figures (i.e. captions) on the plate
  %% page (default is 1). This is necessary for incrementing the figure counter. 
  %%
  %% Note that we have to use a \gdef rather than an \xdef when we 
  %% save the plate output in a command. This means that in the book styles when it 
  %% comes to spewing out the plate, we will have forgotten about any user macros 
  %% that may have been defined. This can be got round by repeating the macros inside the plate.
  %% (Presumably, one would use a \newcommand rather than \renewcommand if they were new.)
  % increment global and local counters
  \addtocounter{plate}{1}
  % write appropriate label to .aux file
  \csname @platelabel\roman{plate}\endcsname
  % define a global macro called \@platei, \@plateii, \@plateiii, \@plateiv, ...
  \expandafter\gdef\csname @plate\roman{plate}\endcsname{%
    \csname @outputplatelabel\roman{plateloop}\endcsname
    \null\vfil #2 \vfil
    \cleardoublepage
    }
  % save the current shortauthor and shorttitle
  % need to be \xdef to take current values of \@shortauthor, etc 
  \expandafter\xdef\csname @shorta\roman{plate}\endcsname{\@shortauthor}
  \expandafter\xdef\csname @shortt\roman{plate}\endcsname{\@shorttitle}
  % save the current value of the figure counter for later use in \outputplates
  \setcounter{figplate\roman{plate}}{\value{figure}}
  % increment the figure counter according to the optional argument
  \addtocounter{figure}{#1}
  }
%% page style for plate pages 
\newcommand{\ps@rmaaplate}{%
  % Plate X (see page Y) in the foot
  \renewcommand\@oddfoot{\hfil\csname @shorta\roman{plateloop}\endcsname\ 
    (see page~\csname @plateref\roman{plateloop}\endcsname)\hfil}%
  \renewcommand\@evenfoot{\relax}%
  % nothing in even heads (blank pages)
  \renewcommand\@evenhead{\relax}
  % short title and page num in the odd heads
  \renewcommand\@oddhead{\hfil\csname @shortt\roman{plateloop}\endcsname%
    \hfil\llap{PLATE~\theplateloop}}
  }





%% Things for the User Guide
\if@guide
%% stolen from doc.sty
\def\MakeShortVerb#1{%
  \expandafter\ifx\csname cc\string#1\endcsname\relax
  \@shortvrbinfo{Made }{#1}%
  \add@special{#1}%
  \expandafter
  \xdef\csname cc\string#1\endcsname{\the\catcode`#1}%
  \begingroup
  \catcode`\~\active  \lccode`\~`#1%
  \lowercase{%
    \global\expandafter\let
    \csname ac\string#1\endcsname~%
    \gdef~{\verb~}}%
  \endgroup
  \global\catcode`#1\active
  \else
  \@shortvrbinfo\@empty{#1 already}%
  \fi}
\def\DeleteShortVerb#1{%
  \expandafter\ifx\csname cc\string#1\endcsname\relax
  \else
  \@shortvrbinfo{Deleted }{#1 as}%
  \rem@special{#1}%
  \global\catcode`#1\csname cc\string#1\endcsname
  \global \expandafter\let \csname cc\string#1\endcsname \relax
  \ifnum\catcode`#1=\active
  \begingroup
  \catcode`\~\active   \lccode`\~`#1%
  \lowercase{%
    \global\expandafter\let\expandafter~%
    \csname ac\string#1\endcsname}%
  \endgroup \fi \fi}
\def\@shortvrbinfo#1#2{%
  \PackageInfo{doc}{%
    #1\expandafter\@gobble\string#2 a short reference
    for \string\verb}}
\def\add@special#1{%
  \rem@special{#1}%
  \expandafter\gdef\expandafter\dospecials\expandafter
  {\dospecials \do #1}%
  \expandafter\gdef\expandafter\@sanitize\expandafter
  {\@sanitize \@makeother #1}}
\def\rem@special#1{%
  \def\do##1{%
    \ifnum`#1=`##1 \else \noexpand\do\noexpand##1\fi}%
  \xdef\dospecials{\dospecials}%
  \begingroup
  \def\@makeother##1{%
    \ifnum`#1=`##1 \else \noexpand\@makeother\noexpand##1\fi}%
  \xdef\@sanitize{\@sanitize}%
  \endgroup}
%% define | as shortcut for \verbatim
\MakeShortVerb{\|}
%% comand to specify arguments to LaTeX commands
\newcommand{\rmArg}[1]{%
  \ensuremath{\langle}{\scriptsize \textsf{\MakeUpperCase{#1}}}\ensuremath{\rangle}}
%\newcommand{\rmArg}[1]{\ensuremath{\langle}\textsc{#1}\ensuremath{\rangle}}
\providecommand{\BibTeX}{{\rmfamily B\kern-.05em%
    \textsc{i\kern-.025em b}\kern-.08em%
    T\kern-.1667em\lower.7ex\hbox{E}\kern-.125emX}}
%% Environment for code examples
\newenvironment{Code}{%
  \begin{list}{}{\setlength{\leftmargin}{0.05\columnwidth}} 
  \item[]}{\end{list}}
\fi
%</package>
%    \end{macrocode}

% \section{User commands for astronomical articles}
% \label{sec:user-commands}
% Some random definitions. These go both into the main class file and
% also into |adshtml.sty|.
% 
% Definitions of the journal name in different incarnations.  
%    \begin{macrocode}
%<*package|adshtml-style>
\newcommand{\RMAA}{RMxAA}
\newcommand{\RMAAlong}{Revista Mexicana de Astronom\'{\i}a y Astrof\'{\i}sica}
\newcommand{\RMAASC}{RMxAC}
\newcommand{\RMAASClong}{\RMAA\ (Serie de Conferencias)}
%    \end{macrocode}
% \begin{macro}{\ibidrule}
% This puts a horizontal rule with a period after it. It is used when
% the author list of a reference is identical to the preceding one. 
%    \begin{macrocode}
\newcommand{\ibidrule}{\rule[-0.03cm]{4.7em}{.01cm}.\ }
%    \end{macrocode}
% \end{macro}
% \begin{macro}{\ThisVolume}
% References to other articles in the same volume can be coded as,
% e.g., |\bibitem{} Bloggs, J. \ThisVolume{bloggs}|.
%    \begin{macrocode}
\newcommand{\ThisVolume}[1]{%
  \@year, \if@proceedings\RMAASC\else\RMAA\fi, \@volume,
  \if@book\oldpageref{#1:RMAAFirstPage}\else 000\fi\ (this volume)}
%    \end{macrocode}
% \end{macro}
% Other stuff that may be useful in ref lists.
%    \begin{macrocode}
\newcommand{\etal}{et al.\@}
%    \end{macrocode}
% Various new commands from AAS macros
%    \begin{macrocode}
\def\acknowledgments{\vskip 3ex plus .8ex minus .4ex}
\let\acknowledgements=\acknowledgments                  % second spelling
\let\jnl@style=\rm
\def\ref@jnl#1{{\jnl@style#1}}
\DeclareRobustCommand\aj{\ref@jnl{AJ}}                   
\DeclareRobustCommand\araa{\ref@jnl{ARA\&A}}             
\DeclareRobustCommand\apj{\ref@jnl{ApJ}}                 
\DeclareRobustCommand\apjl{\ref@jnl{ApJ}}                
\DeclareRobustCommand\apjs{\ref@jnl{ApJS}}               
\DeclareRobustCommand\ao{\ref@jnl{Appl.~Opt.}}           
\DeclareRobustCommand\apss{\ref@jnl{Ap\&SS}}             
\DeclareRobustCommand\aap{\ref@jnl{A\&A}}                
\DeclareRobustCommand\aapr{\ref@jnl{A\&A~Rev.}}          
\DeclareRobustCommand\aaps{\ref@jnl{A\&AS}}              
\DeclareRobustCommand\azh{\ref@jnl{AZh}}                 
\DeclareRobustCommand\baas{\ref@jnl{BAAS}}               
\DeclareRobustCommand\jrasc{\ref@jnl{JRASC}}             
\DeclareRobustCommand\memras{\ref@jnl{MmRAS}}            
\DeclareRobustCommand\mnras{\ref@jnl{MNRAS}}             
\DeclareRobustCommand\pra{\ref@jnl{Phys.~Rev.~A}}        
\DeclareRobustCommand\prb{\ref@jnl{Phys.~Rev.~B}}        
\DeclareRobustCommand\prc{\ref@jnl{Phys.~Rev.~C}}        
\DeclareRobustCommand\prd{\ref@jnl{Phys.~Rev.~D}}        
\DeclareRobustCommand\pre{\ref@jnl{Phys.~Rev.~E}}        
\DeclareRobustCommand\prl{\ref@jnl{Phys.~Rev.~Lett.}}    
\DeclareRobustCommand\pasp{\ref@jnl{PASP}}               
\DeclareRobustCommand\pasj{\ref@jnl{PASJ}}               
\DeclareRobustCommand\qjras{\ref@jnl{QJRAS}}             
\DeclareRobustCommand\skytel{\ref@jnl{S\&T}}             
\DeclareRobustCommand\solphys{\ref@jnl{Sol.~Phys.}}      
\DeclareRobustCommand\sovast{\ref@jnl{Soviet~Ast.}}      
\DeclareRobustCommand\ssr{\ref@jnl{Space~Sci.~Rev.}}     
\DeclareRobustCommand\zap{\ref@jnl{ZAp}}                 
\DeclareRobustCommand\nat{\ref@jnl{Nature}}              
\DeclareRobustCommand\iaucirc{\ref@jnl{IAU~Circ.}}
\DeclareRobustCommand\aplett{\ref@jnl{Astrophys.~Lett.}}
\DeclareRobustCommand\apspr{\ref@jnl{Astrophys.~Space~Phys.~Res.}}
\DeclareRobustCommand\bain{\ref@jnl{Bull.~Astron.~Inst.~Netherlands}}
\DeclareRobustCommand\fcp{\ref@jnl{Fund.~Cosmic~Phys.}}
\DeclareRobustCommand\gca{\ref@jnl{Geochim.~Cosmochim.~Acta}}
\DeclareRobustCommand\grl{\ref@jnl{Geophys.~Res.~Lett.}}
\DeclareRobustCommand\jcp{\ref@jnl{J.~Chem.~Phys.}}      
\DeclareRobustCommand\jgr{\ref@jnl{J.~Geophys.~Res.}}    
\DeclareRobustCommand\jqsrt{\ref@jnl{J.~Quant.~Spec.~Radiat.~Transf.}}
\DeclareRobustCommand\memsai{\ref@jnl{Mem.~Soc.~Astron.~Italiana}}
\DeclareRobustCommand\nphysa{\ref@jnl{Nucl.~Phys.~A}}
\DeclareRobustCommand\physrep{\ref@jnl{Phys.~Rep.}}
\DeclareRobustCommand\physscr{\ref@jnl{Phys.~Scr}}
\DeclareRobustCommand\planss{\ref@jnl{Planet.~Space~Sci.}}
\DeclareRobustCommand\procspie{\ref@jnl{Proc.~SPIE}}
\let\astap=\aap
\let\apjlett=\apjl
\let\apjsupp=\apjs
\let\applopt=\ao
\def\phn{\phantom{0}}
\def\phd{\phantom{.}}
\def\phs{\phantom{\ensuremath{-}}}
\def\phm#1{\phantom{#1}}
%    \end{macrocode}
% All these are redone as robust and using more modern commands where
% possible. The robustness is so they will not get expanded if used in
% abstracts and titles when these are passed to the ADS file. 
% \changes{3.16}{2003/04/07}{
%  Added \cs{Color} command for doing photometric colors.}
%    \begin{macrocode}
%%\def\sun{\hbox{$\odot$}}
\DeclareRobustCommand\sun{\ensuremath{\odot}}
%%\def\earth{\hbox{$\oplus$}}
\DeclareRobustCommand\earth{\ensuremath{\oplus}}
%%\def\lesssim{\mathrel{\hbox{\rlap{\hbox{\lower4pt\hbox{$\sim$}}}\hbox{$<$}}}}
%%\def\gtrsim{\mathrel{\hbox{\rlap{\hbox{\lower4pt\hbox{$\sim$}}}\hbox{$>$}}}}
\DeclareRobustCommand\lesssim{\mathrel{\hbox{\smaller\smaller\rlap{\hbox{\lower3pt\hbox{$\sim$}}}\raise1.5pt\hbox{$<$}}}}
\DeclareRobustCommand\gtrsim{\mathrel{\hbox{\smaller\smaller\rlap{\hbox{\lower3pt\hbox{$\sim$}}}\raise1.5pt\hbox{$>$}}}}
%%\def\lta{\mathrel{\spose{\lower 3pt\hbox{$\mathchar"218$}}
%%     \raise 2.0pt\hbox{$\mathchar"13C$}}}
%%\def\gta{\mathrel{\spose{\lower 3pt\hbox{$\mathchar"218$}}
%%     \raise 2.0pt\hbox{$\mathchar"13E$}}}
\DeclareRobustCommand\sq{\hbox{\rlap{$\sqcap$}$\sqcup$}}
%%\def\arcdeg{\hbox{$^\circ$}}
%%\def\arcmin{\hbox{$^\prime$}}
%%\def\arcsec{\hbox{$^{\prime\prime}$}}
\DeclareRobustCommand\arcdeg{\ensuremath{^\circ}}
\DeclareRobustCommand\arcmin{\ensuremath{^\prime}}
\DeclareRobustCommand\arcsec{\ensuremath{^{\prime\prime}}}
%%\def\fd{\hbox{$.\!\!^{\rm d}$}}
%%\def\fh{\hbox{$.\!\!^{\rm h}$}}
%%\def\fm{\hbox{$.\!\!^{\rm m}$}}
%%\def\fs{\hbox{$.\!\!^{\rm s}$}}
%%\def\fdg{\hbox{$.\!\!^\circ$}}
%%\def\farcm{\hbox{$.\mkern-4mu^\prime$}}
%%\def\farcs{\hbox{$.\!\!^{\prime\prime}$}}
%%\def\fp{\hbox{$.\!\!^{\scriptscriptstyle\rm p}$}}
\DeclareRobustCommand\fd{\ensuremath{.\!\!^\mathrm{d}}}
\DeclareRobustCommand\fh{\ensuremath{.\!\!^\mathrm{h}}}
\DeclareRobustCommand\fm{\ensuremath{.\!\!^\mathrm{m}}}
\DeclareRobustCommand\fs{\ensuremath{\mkern0.5mu.\mkern-5mu^\mathrm{s}\mkern-1mu}}
\DeclareRobustCommand\fdg{\ensuremath{.\!\!^\circ}}
\DeclareRobustCommand\farcm{\ensuremath{.\mkern-4mu^\prime}}
\DeclareRobustCommand\farcs{\ensuremath{\mkern1mu.\mkern-6mu^{\prime\prime}\mkern-1.5mu}}
%%\DeclareRobustCommand\farcs{\ensuremath{.\!\!^{\prime\prime}}}
\DeclareRobustCommand\fp{\ensuremath{.\!\!^{\scriptscriptstyle\rm p}}}
%%\def\micron{\hbox{$\mu$m}}
\DeclareRobustCommand\micron{\ensuremath{\mu\mathrm{m}}}
\let\la=\lesssim                        
\let\ga=\gtrsim
%%\def\case#1#2{\hbox{$\frac{#1}{#2}$}}
%%\def\slantfrac#1#2{\hbox{$\,^#1\!/_#2$}}
%%\def\onehalf{\slantfrac{1}{2}}
%%\def\onethird{\slantfrac{1}{3}}
%%\def\twothirds{\slantfrac{2}{3}}
%%\def\onequarter{\slantfrac{1}{4}}
%%\def\threequarters{\slantfrac{3}{4}}
\DeclareRobustCommand\case[2]{\ensuremath{\frac{#1}{#2}}}
\DeclareRobustCommand\slantfrac[2]{\ensuremath{\,^#1\!/_#2}}
\DeclareRobustCommand\onehalf{\slantfrac{1}{2}}
\DeclareRobustCommand\onethird{\slantfrac{1}{3}}
\DeclareRobustCommand\twothirds{\slantfrac{2}{3}}
\DeclareRobustCommand\onequarter{\slantfrac{1}{4}}
\DeclareRobustCommand\threequarters{\slantfrac{3}{4}}
%%\def\ubvr{\hbox{$U\!BV\!R$}}            
%%\def\ub{\hbox{$U\!-\!B$}}               
%%\def\bv{\hbox{$B\!-\!V$}}               
%%\def\vr{\hbox{$V\!-\!R$}}               
%%\def\ur{\hbox{$U\!-\!R$}}
%    \end{macrocode}
% New commands for ``photometric'' colors 
%    \begin{macrocode}
\DeclareRobustCommand\ubvr{\ensuremath{U\!BV\!R}}            
\DeclareRobustCommand\Color[2]{%
  \ensuremath{%
    #1\!\ifx#1V\!\fi-\!#2
  }}
\DeclareRobustCommand\ub{\Color{U}{B}}
\DeclareRobustCommand\bv{\Color{B}{V}}
\DeclareRobustCommand\vr{\Color{V}{R}}
\DeclareRobustCommand\ur{\Color{U}{R}}
%% New version of the \ion command, that actually works! 
%    \end{macrocode}
% \begin{macro}{\ion}
% This version of the |\ion| command allows you to write either
% |\ion{N}{4}| or |\ion{N}{iv}| in order to get N~IV. 
%    \begin{macrocode}
\DeclareRobustCommand\ion[2]{%
  \mbox{#1\kern0.2em%
    \smaller\rmfamily%
    \edef\@tempa{\@car#2\@nil}%
    \ifcat1\@tempa%
    \@Roman{#2}%
    \else%
    \uppercase{#2}%
    \fi}}
\newcounter{IonStage}
\DeclareRobustCommand\plainion[2]{\setcounter{IonStage}{#2}#1 \Roman{IonStage}}
%%\def\ion#1#2{\ensuremath{\mathrm{#1}\;}{\protect\small\rm\@Roman{#2}}}
%    \end{macrocode}
% \end{macro}
%    \begin{macrocode}
\DeclareRobustCommand\nodata{$\cdots$}
%    \end{macrocode}
% These are needed by the |\ion| command. Isn't there an easier way of
% doing this?
%    \begin{macrocode}
\let\@small=\small
\let\@normalsize=\normalsize
\let\@footnotesize=\footnotesize
\let\@scriptsize=\scriptsize
\let\@tiny=\tiny
\let\@large=\large
\let\@Large=\Large
\let\@LARGE=\LARGE
\let\@huge=\huge
\let\@Huge=\Huge
\newcommand{\smaller}{\small}
\renewcommand{\normalsize}{\renewcommand{\smaller}{\small}\@normalsize}
\renewcommand{\small}{\renewcommand{\smaller}{\footnotesize}\@small}
\renewcommand{\footnotesize}{\renewcommand{\smaller}{\scriptsize}\@footnotesize}
\renewcommand{\scriptsize}{\renewcommand{\smaller}{\tiny}\@scriptsize}
\renewcommand{\tiny}{\renewcommand{\smaller}{\tiny}\@tiny}
\renewcommand{\large}{\renewcommand{\smaller}{\normalsize}\@large}  
\renewcommand{\Large}{\renewcommand{\smaller}{\large}\@Large}  
\renewcommand{\LARGE}{\renewcommand{\smaller}{\Large}\@LARGE}  
\renewcommand{\huge}{\renewcommand{\smaller}{\LARGE}\@huge}  
\renewcommand{\Huge}{\renewcommand{\smaller}{\huge}\@Huge}  
%</package|adshtml-style>
%    \end{macrocode}


% \section{Other files}

% \subsection{The driver and make files}
% \label{sec:driver_make}
% These are put here for completeness, although they still need to be
% distributed separately since otherwise the user won't be able to
% unpack this file into its constituent parts.
% 
% The |Makefile| has two principal rules: |docs|, and |macros|. The
% |docs| rule produces a PDF and gzipped PS file of the documentation.
% I still haven't perfected the combination of rules required to get
% it to run \LaTeX{} enough times in all circumstances. The |macros|
% rule runs \LaTeX{} on |rmaa.ins| in order to unpack |rmaa.cls| and
% all the other little files (including the |Makefile| and the
% |rmaa.ins| files themselves\dots 
% 
% First set up lists of various types of file. 
%    \begin{macrocode}
%<*makefile>
DOCFILES=rmaa.pdf
MACRO_LTX_FILES=rmaa.cls rmaacite.sty adshtml.sty
SHELLSCRIPTS=makeads.sh
MACRO_OTHER_FILES=adshtml.perl sample.latex2html-init
MAKEFILE=Makefile
MACROFILES=$(MACRO_LTX_FILES) $(SHELLSCRIPTS) $(MACRO_OTHER_FILES)         
SKELETONFILES=ADS.tex 
%    \end{macrocode}
% Install the files read only into the package directory, specified via a path relative to the source directory. \changes{4.0}{2019/11/14}{Reorganize folder layout so that hardly anything is in the top level.  That way, things are clearer on github.}
%    \begin{macrocode}
INSTALL=install 
INSTALL_DIR=../packages
%    \end{macrocode}
% The default rule is to make everything. Installation must be
% explicitly asked for.
%    \begin{macrocode}
all: docs macros

install: install-docs install-macros

%    \end{macrocode}
% Most files are installed read-only to prevent editing. Skeleton
% files must be writeable. Must be careful not to overwrite the
% Makefile in the parent directory. 
%    \begin{macrocode}
install-macros: macros
        $(INSTALL) -m 444 $(MACRO_LTX_FILES) $(MACRO_OTHER_FILES) $(INSTALL_DIR)
        $(INSTALL) -m 555 $(SHELLSCRIPTS) $(INSTALL_DIR)
        $(INSTALL) -m 644 $(SKELETONFILES) $(INSTALL_DIR)

install-docs: docs
        $(INSTALL) rmaa.pdf $(INSTALL_DIR)/rm-fulldocs.pdf

%    \end{macrocode}
% Make the PDF and gzipped PS files. 
%    \begin{macrocode}
docs: $(DOCFILES)

%    \end{macrocode}
% Make everything else if either the |.dtx| or |.ins| files have
% changed. 
%    \begin{macrocode}
macros: rmaa.dtx rmaa.ins
%    \end{macrocode}
% Ensure that all the macro files can be written to then recreate
% them. The |touch| is there in case any of the files don't exist
% yet. 
%    \begin{macrocode}
        touch $(MACROFILES) $(MAKEFILE)
        chmod u+w $(MACROFILES) $(MAKEFILE) rmaa.ins
        latex "%%
%% This is file `rmaa.ins',
%% generated with the docstrip utility.
%%
%% The original source files were:
%%
%% rmaa.dtx  (with options: `ins')
%% 
%% PLEASE DO NOT EDIT THIS FILE! It is generated automatically from
%% the file rmaa.dtx using the accompanying Makefile and rmaa.ins
%% Any hand-edited changes will be lost
%% 
%% A LaTeX document class for the RMxAA
%% 
%% Author: William Henney (w.henney@crya.unam.mx)
%% 
%% Latest version available from: http://www.crya.unam.mx/rmaa
%% 
%% Copyright (C) 1999-2010 Instituto de Astronomia, UNAM, Mexico
%% 
%% This program is free software; you can redistribute it and/or
%% modify it under the terms of the GNU General Public License
%% as published by the Free Software Foundation; either version 2
%% of the License, or (at your option) any later version.
%% 
%% This program is distributed in the hope that it will be useful,
%% but WITHOUT ANY WARRANTY; without even the implied warranty of
%% MERCHANTABILITY or FITNESS FOR A PARTICULAR PURPOSE.  See the
%% GNU General Public License for more details.
%% 
%% You should have received a copy of the GNU General Public License
%% along with this program; if not, write to the Free Software
%% Foundation, Inc., 59 Temple Place - Suite 330, Boston, MA  02111-1307, USA.





%%
%% Installation batch file for the RMAA LaTeX package
%%
\def\batchfile{rmaa.ins}
\input docstrip.tex
\preamble

PLEASE DO NOT EDIT THIS FILE! It is generated automatically from
the file rmaa.dtx using the accompanying Makefile and rmaa.ins
Any hand-edited changes will be lost

A LaTeX document class for the RMxAA

Author: William Henney (w.henney@crya.unam.mx)

Latest version available from: http://www.crya.unam.mx/rmaa

Copyright (C) 1999-2010 Instituto de Astronomia, UNAM, Mexico

This program is free software; you can redistribute it and/or
modify it under the terms of the GNU General Public License
as published by the Free Software Foundation; either version 2
of the License, or (at your option) any later version.

This program is distributed in the hope that it will be useful,
but WITHOUT ANY WARRANTY; without even the implied warranty of
MERCHANTABILITY or FITNESS FOR A PARTICULAR PURPOSE.  See the
GNU General Public License for more details.

You should have received a copy of the GNU General Public License
along with this program; if not, write to the Free Software
Foundation, Inc., 59 Temple Place - Suite 330, Boston, MA  02111-1307, USA.
\endpreamble
\askforoverwritefalse
\generate{%
  \file{rmaa.cls}{\from{rmaa.dtx}{package}}
  \file{rmaa.ins}{\from{rmaa.dtx}{ins}}
  \file{rmaacite.sty}{\from{rmaa.dtx}{rmaacite}}
  \file{adshtml.sty}{\from{rmaa.dtx}{adshtml-style}}
  \file{ADS.tex}{\from{rmaa.dtx}{ads-container}}
  }
\generate{\nopreamble\nopostamble%
  \file{Makefile}{\from{rmaa.dtx}{makefile}}
  \file{adshtml.perl}{\from{rmaa.dtx}{adshtml-perl}}
  \file{makeads.sh}{\from{rmaa.dtx}{ads-shell-script}}
  \file{sample.latex2html-init}{\from{rmaa.dtx}{latex2html-init}}
  }

%% Local Variables:
%% mode: LaTeX
%% mode: Font-lock
%% End:









\endinput
%%
%% End of file `rmaa.ins'.
"
%    \end{macrocode}
% \TeX{} automatically adds the |.tex| extension to extensionless
% files. We fix this and take the opportunity to convert
% back to TABs all that it has lovingly converted to spaces. 
%    \begin{macrocode}
        cp -f $(MAKEFILE) $(MAKEFILE).BCK && unexpand $(MAKEFILE).tex >| $(MAKEFILE) && rm $(MAKEFILE).tex
%    \end{macrocode}
% Make the shell scripts executable and everything unwriteable (so
% they don't get hand-edited by mistake). However, undo this for the
% files that the user is \emph{supposed} to modify. 
%    \begin{macrocode}
        chmod a+x $(SHELLSCRIPTS)
        chmod a-w $(MACROFILES) $(MAKEFILE)
        chmod u+w $(SKELETONFILES)

%    \end{macrocode}
% We now use PDF\LaTeX{} for this.
%    \begin{macrocode}
rmaa.idx: rmaa.dtx
        @echo "REGENERATING .idx FILE"
        pdflatex "\nonstopmode% \def\fileversion{4.0} \def\filedate{2019/11/12}
% \MakeShortVerb{\|}
% \title{A \LaTeX{} document class for the Revista Mexicana de
% Astronom\'\i a y Astrof\'\i sica\thanks{This file has version
% \fileversion{} last revised \filedate{}.}}
% \author{William J. Henney\thanks{w.henney@irya.unam.mx}}
% \date{\filedate}
% \maketitle
% \paragraph{Abstract}
% This document describes the \texttt{rmaa} \LaTeX{}
% document class for the journal ``Revista Mexicana de
% Astronom\'\i a y Astrof\'\i sica''. The macros can be used for 
% \begin{enumerate}
% \item Preparing author manuscripts for the journal (including
% conference proceedings to be published in the ``Serie de
% Conferencias'')
% \item Typesetting an entire volume of the journal by inputting the
% individual article files from a container file. Facilities are
% included for automatic generation of TOC and Author Index. 
% \end{enumerate}
% 
% \tableofcontents
% \section{The documentation driver file}
% This is the part that generates the documentation that you may or
% may not be reading, which is done with the help of the Makefile
% discussed in section~\ref{sec:driver_make}. 
%    \begin{macrocode}
%<*driver>
\documentclass{ltxdoc}
\usepackage{hyperref}
\usepackage{palatino}
\EnableCrossrefs \CodelineIndex \RecordChanges
\setlength{\parindent}{0pt}
\def\MacroFont{\footnotesize\tt}
\begin{document} \DocInput{rmaa.dtx} \PrintIndex \PrintChanges \end{document}
%</driver>
%    \end{macrocode}

% \section{The implementation}
% 
% \changes{3.24}{2007/09/14}{Added user access to the macro version
% and date via \cs{rmVersion} and \cs{rmDate}}
% 
%    \begin{macrocode}
%<*package>
\newcommand{\rmaa@version}{3.26}
\newcommand{\rmaa@date}{2011/08/15}
\newcommand\rmVersion{\rmaa@version}
\newcommand\rmDate{\rmaa@date}
\NeedsTeXFormat{LaTeX2e}[1995/06/01]
%    \end{macrocode}
% \changes{4.0}{2019/11/02}{Change abstract to appear before
% resumen. Change RevMexAA to RMxAA for consistency with ADS.}
%
% \changes{3.22}{2007/03/22}{Allow proceedings articles to have
%   received/accepted dates} \changes{3.21}{2006/05/15}{Changes for
%   Silvia - abstracts are now shipped out to ADS.}
% \changes{3.19}{2006/01/19}{Changes for Silvia - remove adresses from
%   end, new abstract style} \changes{2.2beta}{2001/10/30}{ Fixed bug
%   whereby twocolumn option was not passed through to book
%   class. Finally sorted out the ``label on last page'' issue. }
% \changes{2.1}{2001/01/23}{ Fixed some inconsistencies in
%   \cs{nocite} and \cs{bibitem} uniqification Fixed the definition of
%   \cs{ion} so it will work in math mode and in contexts where the
%   typesize is different from normal (e.g., title, fig captions).
%   Modified the definition of \cs{@caption}, so we don't write any
%   list-of-figures data (we never use it anyway, and it just causes
%   problems with fragile commands in fig captions). }
% \changes{2.0}{2000/11/30}{ Added `newstyle' option, which should be
%   used in all new proceedings volumes (wider abstract, smaller
%   typeface in fig captions, etc.) Added commands for formatting
%   table captions. New demonstation document rmpostsamp.tex for
%   1-page posters. This will be the last release under a PRI
%   government (for a while...)}  \changes{1.9}{2000/06/06}{
%   Corrected position of page number in footer of articles that open
%   on even pages. Ensured that first article in each session opens on
%   an odd page, as do list of abstracts and author index. Added extra
%   space in TOC before the backmatter. Added command for editors,
%   \cs{BoldSessionTitles}, that makes produces bold titles for
%   session in the TOC.}  \changes{1.8}{2000/05/16}{ Fixed bug in
%   last page label with \cs{suppressfulladdresses} }
% \changes{1.7}{2000/05/11}{ Fixed more little things. Page labels
%   now work, even if you have no bibliography. You can output the
%   full addresses by hand with the \cs{outputfulladdresses}
%   command. Editors' Guide updated but still not finished. Revamped
%   splitbook.sh now produces PS file of each article for web page and
%   also an EPS file of each page for sending to the printers.}
% \changes{1.6}{2000/04/15}{ More cleaning up: full stops after
%   titles in index, running headers and footers in frontmatter,
%   fontsize in title of index, etc} \changes{1.5}{2000/04/15}{
%   Eliminated superscripts from titles in TOC and index}
% \changes{1.4}{2000/04/14}{ Added copyright and license info (GPL)}
% \changes{1.3}{2000/03/01}{ Corrected \cs{figcaption} to do the
%   caption formatting right Corrected font size of sec headings in
%   front matter} \changes{1.2}{2000/02/16}{ Added \cs{plate}
%   command} \changes{1.1}{2000/02/14}{ Added manuscript option,
%   corrected bugs} \changes{1.0}{2000/02/08}{ Now does a whole book
%   - lots of stuff tidied up} \changes{0.7}{1999/10/04}{ Increased
%   spacing between sections in single column styles. Suppressed
%   floats at bottom of pages (didn't work though) }
% \changes{0.6}{1999/09/24}{ Tweaked \cs{adjustfinalcols} (still not
%   perfect) Added proceedings option Added onecolumn/twocolumn
%   options (orthogonal to journal/proceedings) }
% \changes{0.5}{1999/09/23}{ Added \cs{fulladdresses} and
%   \cs{adjustfinalcols} (just about working)}
% \changes{0.4}{1999/09/22}{ Minor stylistic changes}
% \changes{0.3}{1999/06/28}{ Better mechanism for alternative
%   affiliations Added `guide' option for user guide}
% \changes{0.2}{1999/05/27}{ All basic functionality present, I
%   think} \changes{0.1}{1999/05/26}{ Just about working - still
%   missing Fig/Tab/Refs} \changes{0.0}{1999/05/21}{ So bored today,
%   decided to rewrite RMAA LaTeX macros from scratch}
%    \begin{macrocode}
\ProvidesClass{rmaa}[\rmaa@date\space v\rmaa@version\space RMxAA article
class]
%    \end{macrocode}
% \subsection{Option Handling}
% \label{sec:options}
% Declarations that are needed in options
% \changes{3.05}{2003/10/22}{Added natbib option}
% \changes{3.06}{2003/01/03}{Natbib option turned on permanently}
% \changes{3.07}{2003/01/11}{Added proof option}
% \changes{3.14}{2003/01/11}{Option to turn off natbib}
% \changes{3.18}{2003/10/22}{trying out doctex mode} 
% \changes{3.23}{2007/09/09}{Manuscript is now on by default.}
% \changes{3.26}{2011/08/15}{New option for review articles}
%    \begin{macrocode}
\newif\if@reviewarticle
\@reviewarticlefalse
\newif\if@proof
\@prooffalse
\newif\if@journal
\@journaltrue
\newif\if@printdates
\@printdatesfalse
\newif\if@proceedings
\@proceedingsfalse
\newif\if@manuscript
\@manuscripttrue
\newif\if@bigfont
\@bigfonttrue
\newif\if@guide
\@guidefalse
\newif\if@preprint
\@preprintfalse
\newif\if@rmtwocol
\@rmtwocoltrue
\newif\if@biblist
\@biblistfalse
\newif\if@book
\@bookfalse
\newif\if@natbib
\@natbibtrue
\newif\if@rmcite
\@rmcitefalse
\newif\if@titlevisible
\@titlevisibletrue
\newif\if@abstractpresent
\@abstractpresentfalse
\newif\if@printbanner
\@printbannerfalse
\newif\if@suppressfulladdresses
\@suppressfulladdressesfalse
\newif\if@newstyle
\@newstylefalse
\newif\if@debug
\@debugfalse
\newif\if@boldkeywords
\@boldkeywordsfalse
%    \end{macrocode}
% The option handling---mainly setting various booleans.
% \changes{2.2beta11}{2001/11/26}{Made twocol and
% newstyle the default for proceedings.}
% \changes{3.23}{2007/09/09}{Manuscript is now 10 pt instead of 12pt}
% \changes{3.23}{2007/09/10}{Full addresses are now automatically
% suppressed for proceedings style.}
% \changes{3.24}{2007/09/14}{New option \texttt{boldkeywords} that
% does the keywords in the old ugly way. Now off by default.}
%    \begin{macrocode}
\DeclareOption{manuscript}{\typeout{Manuscript for referee} 
  \@manuscripttrue\@journalfalse\@rmtwocolfalse
  \PassOptionsToClass{10pt}{article}}
\DeclareOption{bigfont}{\typeout{Big font manuscript requested by Gary
  Ferland} 
  \@manuscripttrue\@journalfalse\@rmtwocolfalse\@bigfonttrue
  \PassOptionsToClass{11pt}{article}}
\DeclareOption{journal}{\typeout{Main journal option}%
  \@journaltrue\@rmtwocoltrue\@printdatestrue\@manuscriptfalse
  \PassOptionsToClass{10pt}{article}}
\DeclareOption{guide}{\typeout{User guide option}%
  \@guidetrue\@journalfalse\@rmtwocoltrue\@manuscriptfalse
  \PassOptionsToClass{10pt}{article}}
\DeclareOption{preprint}{\typeout{Preprint option}%
  \@preprinttrue\@manuscriptfalse\PassOptionsToClass{10pt}{article}}
\DeclareOption{proof}{\typeout{Author proof option}%
  \@prooftrue\@manuscriptfalse}
\DeclareOption{proceedings}{\typeout{Proceedings option}%
  \@proceedingstrue\@journalfalse\@rmtwocoltrue\@newstyletrue
  \@suppressfulladdressestrue
  \@printdatesfalse\@manuscriptfalse\PassOptionsToClass{10pt}{article}} 
\DeclareOption{printdates}{\typeout{Force printing of received/accepted dates}%
  \@printdatestrue} 
\DeclareOption{book}{\typeout{Whole issue option}%
  \@booktrue} 
\DeclareOption{nonatbib}{\typeout{Turn off natbib support (required
    for compatibility with old rmaacite package).}%
  \@natbibfalse} 
\DeclareOption{newstyle}{\typeout{Use new layout style}%
  \@newstyletrue}
\DeclareOption{oldstyle}{\typeout{Use old layout style}%
  \@newstylefalse}
\DeclareOption{debug}{\typeout{Loads of annoying messages}%
  \@debugtrue}
%    \end{macrocode}
% These two need to be declared explicitly so that |\ExecuteOptions|
% actually does something.
%    \begin{macrocode}
\DeclareOption{twocolumn}{\@rmtwocoltrue}
\DeclareOption{onecolumn}{\@rmtwocolfalse}
\DeclareOption{twoside}{\PassOptionsToClass{\CurrentOption}{article}}
\DeclareOption{boldkeywords}{\@boldkeywordstrue}
%    \end{macrocode}
% Bold, upper-case keywords used to be the default. Now they must be
% turned on explicitly with the \texttt{boldkeywords} option. 
%    \begin{macrocode}
\DeclareOption{boldkeywords}{\@boldkeywordstrue}
%    \end{macrocode}
% Pass any other option along to the \texttt{article} class.
%    \begin{macrocode}
\DeclareOption*{\PassOptionsToClass{\CurrentOption}{article}}
%    \end{macrocode}
% Default options are to be two-sided and based on the standard
% \LaTeX{} \texttt{article} class,
% unless we are processing the entire book, in which case we base
% ourselves on the standard \texttt{book} class. Default is to use
% manuscript option unless otherwise specified. 
%    \begin{macrocode}
\ExecuteOptions{twoside,manuscript} \ProcessOptions\relax
\if@book
\if@rmtwocol\LoadClass[twocolumn]{book}\else\LoadClass{book}\fi
\newcommand{\abstract}{\relax}
\newcommand{\@maketitle}{\relax}
\else
\if@rmtwocol\LoadClass[twocolumn]{article}\else\LoadClass{article}\fi
\fi
%    \end{macrocode}
% Attempt to use natbib for citations 02 Jan 2003. 
% \changes{3.05}{2003/01/02}{If the natbib option is used, then set
% things up so that natbib's version of the bibliography is identical
% to the standard one.}
% \changes{3.09}{2003/01/23}{Explicitly require a recent version of
% natbib, so that we get a decent error message if the package is too
% old.} 
% \changes{3.23}{2007/09/09}{Now allow Bibtex to work properly} 
% \changes{3.25}{2009/10/01}{Change \cs{newcommand} tp \cs{def} when setting \cs{bibfont}. This allows it to work with all versions of natbib.} 
% \changes{3.25}{2010/02/06}{Updated copyright year.} 
% \changes{3.25}{2010/02/06}{Updated email and web addresses to \texttt{crya.unam.mx}.} 
%    \begin{macrocode}
\if@natbib
\RequirePackage{natbib}[1999/02/25]
\bibpunct{(}{)}{;}{a}{}{,}
\def\bibfont{\small}
\renewcommand\bibsection{\section*{REFERENCES}}
\setlength\bibsep{0pt}
\setlength\bibhang{1.5em}
\bibliographystyle{rmaa}
\fi
%    \end{macrocode}
% We want the \texttt{graphicx} package for the extended syntax of
% |\includegraphics|. We want the \texttt{ifthen} package for using in
% these macros (although we don't much yet). We explicitly load
% \texttt{babel} since some \LaTeX{} distributions messed up the
% hyphenation if we didn't.  
%    \begin{macrocode}
\RequirePackage{graphicx}
\RequirePackage{ifthen}
\IfFileExists{babel.sty}{%
%   \let\@restoreindent=\@afterindentfalse
  \RequirePackage[english]{babel}
%   \@restoreindent
  %% reinstate the definition of \@roman so as not to break .idx files
  \@ifundefined{org@@roman}{\relax}{\renewcommand{\@roman}[1]{\org@@roman{##1}}}
  }{%
  \typeout{Package babel.sty not found - don't worry about this}}
%    \end{macrocode}
% Try to load the \texttt{dcolumn} and \texttt{booktabs} packages if
% they can be found. These can produce nicer-looking tables. 
%    \begin{macrocode}
\IfFileExists{dcolumn.sty}{%
  \RequirePackage{dcolumn}
  }{%
  \typeout{Package dcolumn.sty not found - you will not be able to
    align columns on decimal points in tables}}
\IfFileExists{booktabs.sty}{%
  \RequirePackage{booktabs}
  }{%
  \typeout{Package booktabs.sty not found - your tables will not look
    as pretty as they ought too.}
  }
\IfFileExists{makeidx.sty}{%
  \RequirePackage{makeidx}
}{%
  \typeout{Package makeidx.sty not found - Author Index may not work.}
  }
%% commented out for moment since requires more modern version of LaTeX
%% \RequirePackage[latin1]{inputenc}
%    \end{macrocode}
% \changes{4.0}{2019/11/13}{Make \texttt{hyperref} and
% \texttt{etoolbox} into hard requirements, since we now use them
% extensively.}  We use hyperref to get internal and external links.
% But I don't like the colored boxes so I turn those off.  We use
% etoolbox for patching and robustifying at the moment.  In the
% future, I want to port a lot of other stuff over to it.  For
% instance, replace the ifthen package and the TeX ifs too.
%
% We target the latest version (1.7) of PDF, rather than the default
% of 1.5.  This reduces the chance of getting warnings from included
% graphics.  This is not strictly related to hyperref, but we put it
% here anyway.
%    \begin{macrocode}
\RequirePackage{etoolbox}
\RequirePackage{hyperref}
\hypersetup{hidelinks}
\pdfminorversion=7
%    \end{macrocode}
% We patch the \cs{@combinedblfloats} following a suggestion by Heiko
% Oberdiek at \url{https://tex.stackexchange.com/a/249743/3843}, which
% avoids problems when hyperlinks are split across
% pages. \changes{4.0}{2019/11/12}{Fix for issue where breaking
% hyperlink across pages gave an error.}
%    \begin{macrocode}
\urlstyle{same}
\newcommand\email{\begingroup \urlstyle{rm}\Url}
\patchcmd\@combinedblfloats{\box\@outputbox}{\unvbox\@outputbox}{}{%
  \errmessage{\noexpand\@combinedblfloats could not be patched}%
}%
%    \end{macrocode}
% \subsection{Other setup stuff}\label{sec:setup}
% Some declarations: 
%    \begin{macrocode}
\newif\if@needaddresses
\@needaddressestrue
\newif\if@frontmatter
\@frontmatterfalse
\newif\if@inpreamble
\@inpreambletrue
\newif\if@firstauthor
\@firstauthortrue
\newif\if@firstaffil
\@firstaffiltrue
\newif\if@firstaltaffil
\@firstaltaffiltrue
\newif\if@firstkwd
\@firstkwdtrue
\newif\if@noexpandADS
\@noexpandADSfalse
\newif\if@sessiontitlepages
\@sessiontitlepagesfalse
%    \end{macrocode}
% \begin{macro}{\@rmaaEmpty}
% \begin{macro}{\@rmaaArg}
% These are used in various commands to check if arg was left blank
%    \begin{macrocode}
\newcommand{\@rmaaEmpty}{}
\newcommand{\@rmaaArg}{}
%    \end{macrocode}  
% \end{macro}
% \end{macro}
%
% \subsection{Compatibility with old \LaTeX{} versions}
% \label{sec:compat}
% This package is based on the 1998/12/01 release of \LaTeX{}.
% Older releases may not define all the commands that we need.
% Luckily, we can use |\providecommand| to define these if necessary. 
% The following commands are not defined in the 1994/06/01 release. 
%
% \begin{macro}{\MakeUpperCase}
% \begin{macro}{\textsuperscript}
% \begin{macro}{\textbackslash}
% User commands:
%    \begin{macrocode}
\providecommand{\MakeUpperCase}[1]{\uppercase{#1}}
\providecommand{\textsuperscript}[1]{\ensuremath{^\textrm{#1}}}
\providecommand{\textbackslash}{\ensuremath{\backslash}}
%    \end{macrocode}  
% \end{macro}
% \end{macro}
% \end{macro}
% \begin{macro}{\protected@xdef}
% \begin{macro}{\restore@protect}
% \begin{macro}{\@unexpandable@protect}
% \begin{macro}{\hb@xt@}
% \LaTeX{} internal commands:
%    \begin{macrocode}
\providecommand{\protected@xdef}{%
        \let\@@protect\protect
        \let\protect\@unexpandable@protect
        \afterassignment\restore@protect
        \xdef}
\providecommand{\restore@protect}{\let\protect\@@protect}
\providecommand{\@unexpandable@protect}{\noexpand\protect\noexpand}
\providecommand{\hb@xt@}{\hbox to}
%    \end{macrocode}
% \end{macro}
% \end{macro}
% \end{macro}
% \end{macro}
%
% \subsection{Page layout}\label{sec:page-layout}
%
% The dimensions of the text were defined by Elizabeth. Text width is
% 169\,mm with equal margins on the two sides: 
%    \begin{macrocode}
\setlength{\textwidth}{169mm}
\setlength{\@tempdima}{\paperwidth}
\addtolength{\@tempdima}{-\textwidth}
\setlength{\oddsidemargin}{0.5\@tempdima}
\addtolength{\oddsidemargin}{-1in}
\setlength{\evensidemargin}{\oddsidemargin}
%    \end{macrocode}
% Text height is 225\,mm, top margin is 14\,mm. 
%    \begin{macrocode}
\setlength{\textheight}{225mm} 
\setlength{\topmargin}{-1in} 
\addtolength{\topmargin}{14mm}
%    \end{macrocode}
% Space between columns in two-column styles, header, footer,
% paragraph indentation:
%    \begin{macrocode}
\setlength{\columnsep}{9.5mm}
\setlength{\headsep}{0.7cm}
\setlength{\footskip}{0.75cm}
\setlength{\parindent}{1.5em}
%    \end{macrocode}
% 
% \begin{macro}{\@set@scalable@length}
% Try out a new scheme for dealing with various spacing dimensions and
% allowing the user to easily modify them. The command
% \begin{quote}
% |\set@scalable@length|\marg{name}\marg{defaultlength}\marg{stretch}\marg{shrink}
% \end{quote}
% does a 
% |\setlength|\marg{\cs{name}}\marg{defaultlength plus stretch minus shrink}
% and also defines the user command
% \textbackslash\meta{name}|Scale|\marg{factor} that can 
% be used to rescale length by \meta{factor}. Additionally, now
% defines the package command
% \textbackslash\meta{name}|@set@scale|\marg{FACTOR} 
% that also rescales the length by \meta{FACTOR}, so that the length
% becomes \meta{factor} $\times$ \meta{FACTOR} $\times$
% \meta{defaultlength}.   
%    \begin{macrocode}
\newcommand\@set@scalable@length[4]{%
  %% arg #1: NAME of length without leading `\'
  %%
  %% arg #2: default VALUE  
  \expandafter\newlength\csname @#1@default\endcsname
  \expandafter\setlength\csname @#1@default\endcsname{#2}%
  \expandafter\newlength\csname @@#1@@default\endcsname
  \expandafter\setlength\csname @@#1@@default\endcsname{#2}%
  %% arg #3: default STRETCH  
  \expandafter\newlength\csname @#1@stretch\endcsname
  \expandafter\setlength\csname @#1@stretch\endcsname{#3}%
  \expandafter\newlength\csname @@#1@@stretch\endcsname
  \expandafter\setlength\csname @@#1@@stretch\endcsname{#3}%
  %% arg #4: default SHRINK  
  \expandafter\newlength\csname @#1@shrink\endcsname
  \expandafter\setlength\csname @#1@shrink\endcsname{#4}%
  \expandafter\newlength\csname @@#1@@shrink\endcsname
  \expandafter\setlength\csname @@#1@@shrink\endcsname{#4}%
  %%
  %% User-defined scale \@NAME@scale is unity by default
  \expandafter\gdef\csname @#1@scale\endcsname{1}%
  %% Package-defined scale \@NAME@SCALE is unity by default
  \expandafter\gdef\csname @#1@SCALE\endcsname{1}%
  %%
  %% internal command \@reset@NAME should expand to 
  %% \setlength\@@NAME@@default{\@NAME@SCALE\@NAME@default}
  %% \setlength\NAME{\@NAME@scale\@@NAME@@default}
  %% in other words: scale x SCALE x default 
  \expandafter\gdef\csname @reset@#1\endcsname{%
    \expandafter\setlength\csname @@#1@@default\endcsname{%
      \csname @#1@SCALE\endcsname \csname @#1@default\endcsname
      }
    \expandafter\setlength\csname @@#1@@stretch\endcsname{%
      \csname @#1@SCALE\endcsname \csname @#1@stretch\endcsname
      }
    \expandafter\setlength\csname @@#1@@shrink\endcsname{%
      \csname @#1@SCALE\endcsname \csname @#1@shrink\endcsname
      }
    \expandafter\setlength\csname #1\endcsname{%
      \csname @#1@scale\endcsname \csname @@#1@@default\endcsname
      \@plus
      \csname @#1@scale\endcsname \csname @@#1@@stretch\endcsname
      \@minus
      \csname @#1@scale\endcsname \csname @@#1@@shrink\endcsname
      }}
  %%
  %% User command \NAMEScale should expand to 
  %% {\renewcommand\@NAME@scale{#1}\@reset@NAME}
  \expandafter\gdef\csname #1Scale\endcsname##1{%
    \expandafter\gdef\csname @#1@scale\endcsname{##1}%
    \csname @reset@#1\endcsname}
  %%
  %% Package command \NAME@set@SCALE should expand to 
  %% {\renewcommand\@NAME@SCALE{#1}\@reset@NAME}
  \expandafter\gdef\csname #1@set@SCALE\endcsname##1{%
    \expandafter\gdef\csname @#1@SCALE\endcsname{##1}%
    \csname @reset@#1\endcsname}
  %%
  %% initialize at default values with \@reset@NAME
  \csname @reset@#1\endcsname
  }
%    \end{macrocode}
% \end{macro}
% \begin{macro}{\parskip}
% First use of this new infrastructure: set interparagraph space to be
% zero but with a small amount of stretchiness.
%    \begin{macrocode}
%   \ifthenelse{\boolean{@manuscript}}{%
%     \@set@scalable@length{parskip}{0.5\baselineskip}{0.3\p@}{0.1\p@}
%   }{
    \@set@scalable@length{parskip}{0\p@}{0.3\p@}{0\p@}
%   }
%    \end{macrocode}
% \end{macro}
% 
% \begin{macro}{\abovesectionskip}
% \changes{4.0}{2019/11/12}{Use negative values for skip so that first paragraph after section headings is not indented.  This is something that I have been trying to fix for years!  It turns out that the anser was in the documentation for the standard command \cs{@startsection.}}
% \begin{macro}{\belowsectionskip}
% \begin{macro}{\aboveparaskip}
% \begin{macro}{\belowparaskip}
% Define some utility rubber lengths for use in
% section~\ref{sec:section-commands}. 
%    \begin{macrocode}
\newlength{\abovesectionskip}
\newlength{\belowsectionskip}
\newlength{\aboveparaskip}
\newlength{\belowparaskip}
%    \end{macrocode}
% The default values for these lengths are those traditionally used
% for the onecolumn style. 
%    \begin{macrocode}
\@set@scalable@length{abovesectionskip}{-3.5ex}{-4ex}{-2ex}
\@set@scalable@length{belowsectionskip}{2.3ex}{0ex}{2ex}
\@set@scalable@length{aboveparaskip}{3.25ex}{2ex}{0.5ex}
\@set@scalable@length{belowparaskip}{-1em}{\z@}{\z@}
%    \end{macrocode}
% \end{macro}\end{macro}\end{macro}\end{macro}
% We also make scalable some of the standard float parameters. These
% are the same defaults as in the standard article classes, but
% slightly more stretchy. For the twocolumn style, they are currently
% halved, see below. Note that I have also lumped |\abovecaptionskip|
% in with these, since it plays a similar role. There is also a
% |\belowcaptionskip| but it is set to zero in the standard classes. 
%    \begin{macrocode}
\@set@scalable@length{floatsep}{12\p@}{4\p@}{4\p@}
\@set@scalable@length{textfloatsep}{20\p@}{4\p@}{8\p@}
\@set@scalable@length{intextsep}{12\p@}{4\p@}{4\p@}
\@set@scalable@length{dblfloatsep}{12\p@}{4\p@}{4\p@}
\@set@scalable@length{dbltextfloatsep}{20\p@}{4\p@}{8\p@}
\@set@scalable@length{abovecaptionskip}{10\p@}{\z@}{\z@}
%    \end{macrocode}
% \begin{macro}{\@before@title@skip}         
% \begin{macro}{\@before@author@skip}        
% \begin{macro}{\@between@author@skip}       
% \begin{macro}{\@before@abstract@skip}      
% \begin{macro}{\@after@abstracthead@skip}   
% \begin{macro}{\@between@abstract@skip}     
% \begin{macro}{\@before@kwd@skip}           
% \begin{macro}{\@before@dates@skip}           
% Spacing parameters used in |\@maketitle|. These are not changed by
% default except by the user, so they don't need to go into
% |\@package@default@lengths|.  
%    \begin{macrocode}
\newlength{\@before@title@skip}         
\newlength{\@before@author@skip}        
\newlength{\@between@author@skip}       
\newlength{\@before@abstract@skip}      
\newlength{\@after@abstracthead@skip}   
\newlength{\@between@abstract@skip}     
\newlength{\@before@kwd@skip}           
\newlength{\@before@dates@skip}           
\@set@scalable@length{@before@title@skip}{2em}{\z@}{\z@}
\@set@scalable@length{@before@author@skip}{1.5em}{\z@}{\z@}
\@set@scalable@length{@between@author@skip}{0.5em}{\z@}{\z@}
\@set@scalable@length{@before@abstract@skip}{3ex}{1ex}{0.5ex}
\@set@scalable@length{@after@abstracthead@skip}{1ex}{\z@}{\z@}
\@set@scalable@length{@between@abstract@skip}{2ex}{1ex}{0.5ex}
\@set@scalable@length{@before@kwd@skip}{1ex}{0.5ex}{0.2ex}
\@set@scalable@length{@before@dates@skip}{2ex}{0.5ex}{0.2ex}
%    \end{macrocode}
% \end{macro}\end{macro}\end{macro}\end{macro}
% \end{macro}\end{macro}\end{macro}\end{macro}
% 
% Spacing above and below displayed equations also wants to be done
% like this, but it is more complicated since the standard article
% classes make these change with the fontsize. Something to do later,
% maybe. 
% 
%    \begin{macrocode}
%    \end{macrocode}
% \begin{macro}{\@package@default@lengths}
% The value actually used for these lengths is controlled by the
% \meta{SCALE} parameter, which depends on the style option that is in
% effect. This is set using the |\@package@default@lengths| command,
% which must be called every time the style changes (e.g., changing
% from onecolumn to twocolumn). 
%    \begin{macrocode}
\newcommand{\@package@default@lengths}{%
  \ifthenelse{\boolean{@rmtwocol}}{%
    %% twocolumn mode
    \abovesectionskip@set@SCALE{0.5}
    \belowsectionskip@set@SCALE{0.5}
    \aboveparaskip@set@SCALE{0.5}
    \belowparaskip@set@SCALE{1}
    \floatsep@set@SCALE{0.5}
    \textfloatsep@set@SCALE{0.5}
    \intextsep@set@SCALE{0.5}
    \dblfloatsep@set@SCALE{0.5}
    \dbltextfloatsep@set@SCALE{0.5}
    }{%
    %% onecolumn mode
    \abovesectionskip@set@SCALE{1}
    \belowsectionskip@set@SCALE{1}
    \aboveparaskip@set@SCALE{1}
    \belowparaskip@set@SCALE{1}
    \floatsep@set@SCALE{1}
    \textfloatsep@set@SCALE{1}
    \intextsep@set@SCALE{1}
    \dblfloatsep@set@SCALE{1}
    \dbltextfloatsep@set@SCALE{1}
    }%
  }
\@package@default@lengths
%    \end{macrocode}
% \end{macro}
% \begin{macro}{\RescaleLengths}
% \begin{macro}{\RescaleSecLengths}
% \begin{macro}{\RescaleFloatLengths}
% \begin{macro}{\RescaleTitleLengths}
% User commands that will fiddle with all the scalable lengths. 
%    \begin{macrocode}
\newcommand{\RescaleLengths}[1]{%
  \RescaleSecLengths{#1}\RescaleFloatLengths{#1}\RescaleTitleLengths{#1}}
\newcommand{\RescaleSecLengths}[1]{%
%% sections
  \abovesectionskipScale{#1}
  \belowsectionskipScale{#1}
  \aboveparaskipScale{#1}
  }
\newcommand{\RescaleFloatLengths}[1]{%
%% floats
  \floatsepScale{#1}
  \textfloatsepScale{#1}
  \intextsepScale{#1}
  \dblfloatsepScale{#1}
  \dbltextfloatsepScale{#1}
  \abovecaptionskipScale{#1}
  }
\newcommand{\RescaleTitleLengths}[1]{%
%% maketitle
  \@before@title@skipScale{#1}           
  \@before@author@skipScale{#1}        
  \@between@author@skipScale{#1}       
  \@before@abstract@skipScale{#1}      
  \@after@abstracthead@skipScale{#1}   
  \@between@abstract@skipScale{#1}     
  \@before@kwd@skipScale{#1}           
  \@before@dates@skipScale{#1}           
  }
%    \end{macrocode}
% \end{macro}\end{macro}\end{macro}\end{macro}
% \changes{2.2beta14}{2002/03/17}{Finished implementation of new
% infrastructure for scalable lengths. Converted the above/below
% section skips to use it. Also, all the floatsep parameters. Also,
% all the vertical lengths that are used by |\@maketitle|. }

% \begin{macro}{\MSextramargin}
% For the \texttt{manuscript} option we have an extra 0.3\,inch
% vertical margin and 1\,inch horizontal margin for marginal notes (by
% hand or with |\marginpar|. 
%
% \changes{3.23}{2007/09/09}{Reduce baselinestretch in
%   \texttt{manuscript} style to 1.2. No longer use raggedright for
%   manuscript.}
%
% \changes{3.24}{2007/09/14}{Further reduction of baselinestretch in
%   \texttt{manuscript} style to 1.1 since we now have wider margins.}
%
% \changes{3.24}{2007/09/14}{Split \texttt{manuscript}'s extra margin
%   into separate horizontal and vertical parts.}
% 
%    \begin{macrocode}
\newlength{\MSextraHmargin}\newlength{\MSextraVmargin}
\setlength{\MSextraHmargin}{1in} 
\setlength{\MSextraVmargin}{0.3in} 
\if@manuscript
   \renewcommand{\baselinestretch}{1.1}
   \addtolength{\textwidth}{-2\MSextraHmargin}
   \addtolength{\oddsidemargin}{\MSextraHmargin}
   \setlength{\evensidemargin}{\oddsidemargin}
   \addtolength{\textheight}{-2\MSextraVmargin}
   \addtolength{\topmargin}{\MSextraVmargin}
%    \raggedright
\fi
%    \end{macrocode}
% \end{macro}
% Float parameters: don't allow bottom floats by default. 
%    \begin{macrocode}
\renewcommand{\fps@figure}{tp}
\renewcommand{\fps@table}{tp}
%    \end{macrocode}
%
% \begin{macro}{\rmaa@AuthorWidthFrac}
% \begin{macro}{\rmaa@AbstractWidthFrac}
% These parameters control what fraction of the total text width can
% be used by the list of authors and the abstract/resumen,
% respectively.  
%    \begin{macrocode}
\newcommand{\rmaa@AuthorWidthFrac}{1.0}
\newcommand{\rmaa@AbstractWidthFrac}{0.75}
%    \end{macrocode}
% Use a wide abstract to save space in the new style.
%    \begin{macrocode}
\if@newstyle\renewcommand{\rmaa@AbstractWidthFrac}{1.0}\fi
%    \end{macrocode}
% \end{macro}
% \end{macro}
%
% Some things to consider in future versions. Trying to avoid too many
% widows, etc. Apparently, the following will work best with
% \cs{flushbottom}. 
%    \begin{macrocode}
%%\widowpenalty=10000
%%\clubpenalty=10000
%    \end{macrocode}
% Also, should perhaps try getting rid of underfull/overfull box warnings. 
% \changes{3.10}{2003/01/24}{Set large values for \cs{hbadness},
% \cs{vbadness}, \cs{hfuzz}, \cs{vfuzz}. To reduce amount of stuff
% that gets written to the terminal every time you compile something.}
%    \begin{macrocode}
\hbadness=10000
\vbadness=10000
\setlength\hfuzz{\maxdimen}
\setlength\vfuzz{\maxdimen}
%    \end{macrocode}
% \subsection{Various text styles for different contexts}
% \label{sec:text-styles} 
% \begin{macro}{\rmaa@AffilStyle}
% \begin{macro}{\rmaa@PagesStyle}
% \begin{macro}{\rmaa@YearStyle}
% \begin{macro}{\rmaa@VolumeStyle}
% \begin{macro}{\rmaa@VolumenStyle}
% \begin{macro}{\rmaa@JournalNameStyle}
% \begin{macro}{\rmaa@ReceivedStyle}
% \begin{macro}{\rmaa@AbstractHeaderStyle}
% \begin{macro}{\rmaa@JournalNameStyle}
% Style commands for the different fields used by |\maketitle| in
% constructing the title page of the article. These have the form
% |\rmaa@XYZStyle{TEXT}|. I'm not sure these were really a good idea,
% but never mind \dots 
%    \begin{macrocode}
\newcommand{\rmaa@AffilStyle}[1]{{\normalsize\rmfamily#1}}
\newcommand{\rmaa@PagesStyle}[1]{{\normalsize\rmfamily#1}}
\newcommand{\rmaa@YearStyle}[1]{{\normalsize\rmfamily#1}}
\newcommand{\rmaa@VolumeStyle}[1]{{\normalsize\rmfamily\bfseries#1}}
\newcommand{\rmaa@VolumenStyle}[1]{{\Large\sffamily\bfseries#1}}
\newcommand{\rmaa@JournalNameStyle}[1]{{\normalsize\itshape#1}}
\newcommand{\rmaa@ReceivedStyle}[1]{{\normalsize\rmfamily\itshape#1}}
\newcommand{\rmaa@AbstractHeaderStyle}[1]{%
  {\normalsize\rmfamily\MakeUpperCase{#1}}}
\newcommand{\rmaa@TitleStyle}[1]{%
  {\large\bfseries\rmfamily\boldmath\MakeUpperCase{#1}\par}}
%    \end{macrocode}
% \end{macro}
% \end{macro}
% \end{macro}
% \end{macro}
% \end{macro}
% \end{macro}
% \end{macro}
% \end{macro}
% \end{macro}
%
% \begin{macro}{\rmaa@FrontSecStyle}
% Headings of ``front sections'' (see section~\ref{sec:front-sections})
%    \begin{macrocode}
\newcommand{\rmaa@FrontSecStyle}[1]{%
  {\normalsize\bfseries\rmfamily\MakeUpperCase{#1}}}
%    \end{macrocode}
% \end{macro}
%
% \begin{macro}{\rmaa@TOCTitleStyle}
% \begin{macro}{\rmaa@TOCPartStyle}
% \begin{macro}{\rmaa@IndexTitleStyle}
% \begin{macro}{\rmaaTOCAuthorStyle}
% \begin{macro}{\rmaaIndexAuthorStyle}
% Commands to set style of TOC and index entries.
% Rewritten to avoid unnecessary font commands in the .toc and .ind files
%    \begin{macrocode}
\newcommand{\rmaa@TOCTitleStyle}[1]{\MakeUpperCase{#1}}
\newcommand{\rmaa@TOCPartStyle}[1]{\rmaa@TOCTitleStyle{#1}}
\newcommand{\rmaa@IndexTitleStyle}[1]{#1}
\newcommand{\rmaaTOCAuthorStyle}[1]{%
  {\itshape#1}}
\newcommand{\rmaaIndexAuthorStyle}[1]{%
  \protect\mbox{\bfseries#1}}
%    \end{macrocode}
% \end{macro}
% \end{macro}
% \end{macro}
% \end{macro}
% \end{macro}
%
% \begin{macro}{\rmaa@KeywordStyle}
% \begin{macro}{\rmaa@KeywordHeaderStyle}
% The style of the keywords depends on whether we are using the
% \texttt{newstyle} option or not. If we are, then they are small. 
%    \begin{macrocode}
\ifthenelse{\boolean{@newstyle}}{
  \newcommand{\rmaa@KeywordStyle}[1]{%
    {\small%
      \if@boldkeywords 
      \bfseries\MakeUpperCase{\raggedright #1}
      \else
      \raggedright #1
      \fi
    }}
  \newcommand{\rmaa@KeywordHeaderStyle}[1]{%
    {\small\itshape#1}}}
{
  \newcommand{\rmaa@KeywordStyle}[1]{%
    {\normalsize%
      \if@boldkeywords\bfseries\fi\raggedright #1}}
  \newcommand{\rmaa@KeywordHeaderStyle}[1]{%
    {\normalsize\itshape#1}}}
%    \end{macrocode}
% \end{macro}
% \end{macro}
%
% \subsection{Page styles}\label{sec:page-styles}
% \begin{macro}{\ps@rmaa}
% The running headers for normal pages:
%    \begin{macrocode}
\newcommand{\ps@rmaa}{%
%    \end{macrocode}
% nothing in either foot; 
%    \begin{macrocode}
  \renewcommand\@oddfoot{\relax}%
  \renewcommand\@evenfoot{\relax}%
%    \end{macrocode}
% page number and short author list in the even heads; 
%    \begin{macrocode}
  \renewcommand\@evenhead{\rlap{\thepage}\hfil\@shortauthor\hfil}%
%    \end{macrocode}
% short title and page number in the odd heads.
%    \begin{macrocode}
  \renewcommand\@oddhead{\hfil\@shorttitle\hfil\llap{\thepage}}
  }
%    \end{macrocode}
% \begin{macro}{\ps@rmaafirstpage}
% The header of the first page of each
% article:
%    \begin{macrocode}
\newcommand{\ps@rmaafirstpage}{%
  \renewcommand\@oddfoot{\hfil\thepage}%
%    \end{macrocode}
% footer corrected (06 Jun 2000) to have page on outside edge;
%    \begin{macrocode}
  \renewcommand\@evenfoot{\thepage\hfil}%
%    \end{macrocode}
% Header has journal name, volume, and page numbers for a regular
% article, or similar for special cases. 
% \changes{4.0}{2019/11/03}{Change to three-line header with copyright info and DOI.}
%    \begin{macrocode}
\if@preprint
\renewcommand\@oddhead{%
  \if@proceedings 
  \textsc{To appear in ``\@conftitle{} (\rmaa@YearStyle{\@year})''}\hfil
  \textit{RMxAC}
  \else
  \textsc{To appear in} \@journalname\hfil
  \fi}\fi%
\if@guide
\renewcommand\@oddhead{\textsc{Documentation} (\rmaa@YearStyle{\@year})\hfil}
\fi%
\if@manuscript
\renewcommand\@oddhead{Manuscript for \@journalname{} (\rmaa@YearStyle{\@year})\hfil}
\fi%
\if@journal
\renewcommand\@oddhead{%
  \parbox{\textwidth}{
    \@journalname, \rmaa@VolumeStyle{\@volume},
    \rmaa@PagesStyle{\@artpages}\
    (\rmaa@YearStyle{\@year})\hfil\\[1pt]
    {\small\textcopyright{} \@year:
      Instituto de Astronom\'\i{}a,
      Universidad Nacional Aut\'onoma de M\'exico}\hfil\\[1pt]
    {\small\@DOI@url}\hfil
  }%
}\fi%
  \renewcommand\@evenhead{\@oddhead}}
%    \end{macrocode}
% By default, use the |\ps@rmaa| page style.
%    \begin{macrocode}
\pagestyle{rmaa}
%    \end{macrocode}
% \end{macro}
% \end{macro}
%
% \subsection{Preamble commands}\label{sec:preamble} 
% Everything you need for setting the title page, TOC and index
% entries, etc.
% \subsubsection{Commands for the title page}
% \label{sec:title-commands}
% \begin{macro}{\resumen}
% \begin{macro}{\@resumen}
% \begin{macro}{\abstract}
% \begin{macro}{\@abstract}
% User commands for setting the resumen and abstract. 
% These should be used in the document header. The internal commands
% |\@resumen| and |\@abstract| are used to hold the actual text
% pending output via the |\maketitle| command. 
%    \begin{macrocode}
\newif\if@setresumen
\newcommand{\resumen}[1]{%
%%  \renewcommand{\@resumen}{{\selectlanguage{spanish}#1}}%
  \renewcommand{\@resumen}{{\ignorespaces#1}}%
  \@setresumentrue}
\renewcommand{\abstract}[1]{%
  \renewcommand{\@abstract}{{\ignorespaces#1}}
  \renewcommand\ADS@abstract{#1}
%% Not a good idea, since then author macros are not expanded and
%% cause errors later. However, by not doing it we are vulnerable to
%% fragile commands.....
  \if@noexpandADS\@onelevel@sanitize{\ADS@abstract}\fi
  \@abstractpresenttrue
  }
\newcommand{\ADSnoexpand}{\@noexpandADStrue}
%    \end{macrocode}
% \begin{macro}{\boldabstract}
% This is the version for doing 1-page conference proceedings. Still
% save the text for ADS but then write it out immediately as a bold
% paragraph. 
%    \begin{macrocode}
\newcommand{\boldabstract}[1]{%
  \gdef\ADS@abstract{#1}%
  \renewcommand{\@abstract}{#1}%
  \if@noexpandADS\@onelevel@sanitize{\ADS@abstract}\fi%
  \noindent\textbf{\boldmath #1}
  }
%    \end{macrocode}
% \end{macro}
% Default value in each case is to print out a warning if not set
%    \begin{macrocode}
  \newcommand{\@resumen}{%
    \@setresumenfalse \ClassWarningNoLine{rmaa}{No Spanish abstract
      (resumen) specified - if you do not speak Spanish, this will be
      translated by the editors}%
    \textbf{Favor de proporcionar un resumen en espa\~nol. If you are
      unable to translate your abstract into Spanish, the editors will
      do it for you.} }
  \newcommand{\@abstract}{\ClassWarningNoLine{rmaa}{No abstract
      specified}}
%    \end{macrocode}
% \end{macro}
% \end{macro}
% \end{macro}
% \end{macro}
%
% \begin{macro}{\title}
% \begin{macro}{\@rawtitle}
% \begin{macro}{\@TOCtitle}
% \begin{macro}{\rmaaIndexTitle}
% Similarly, |\title| is the preamble command to set the title of the
% article. This is used in various places, typeset in different
% styles, hence the use of the internal |\@rawtitle|, |\@title|,
% |\rmaaIndexTitle|, and |\@TOCtitle| commands, each of which uses the
% relevant style from section~\ref{sec:text-styles}
%    \begin{macrocode}
\newcommand{\@rawtitle}{Raw Title}
\newcommand{\@TOCtitle}{\rmaa@TOCTitleStyle{Title for TOC}}
\newcommand{\rmaaIndexTitle}{\rmaa@IndexTitleStyle{Title for index}.}
\newcommand{\@TOCaftertitle}{}
\renewcommand{\title}[1]{%
%%  {\ADS@sanitize\xdef\@rawtitle{#1}}
  \renewcommand\@rawtitle{\noexpand#1}
  \let\ADS@title=\@rawtitle
  \renewcommand{\@title}{\rmaa@TitleStyle{#1}}%
  \renewcommand{\@TOCtitle}{\rmaa@TOCTitleStyle{#1}\@TOCaftertitle}%
  \renewcommand{\rmaaIndexTitle}{\rmaa@IndexTitleStyle{#1\@TOCaftertitle}.}%
  }
%    \end{macrocode}
% \end{macro}
% \end{macro}
% \end{macro}
% \end{macro}
% 
% All the weird and whacky ways of indicating the author
% affiliations. 
% \begin{macro}{\affil}
% First, the way that sets it directly underneath the authors name.
% Note that |\affil| trumps all other affiliations for the purposes of
% getting into the ADS stuff (we use |\gdef|).   
% \dots 
%    \begin{macrocode}
\newcommand{\affil}[1]{%
  \\[1.5ex]\rmaa@AffilStyle{#1}\\[1.5ex]%
  \if@firstaffil\gdef\ADS@affil{#1}\@firstaffilfalse\fi%
  }
\renewcommand{\and}{and\\[1.5ex]}
%    \end{macrocode}
% \begin{macro}{\altaffil}
% Alternative affiliations are footnoted. 
%    \begin{macrocode}
\newcommand{\altaffil}[1]{%
  \if@firstaltaffil\gdef\ADS@affil{#1}\@firstaltaffilfalse\fi%
  \thanks{#1}}
%    \end{macrocode}
% \begin{macro}{\altaffilmark}
% \begin{macro}{\altaffiltext}
% Alternative-alternative mechanism for affiliations---should not be
% mixed with |\altaffil| (unless you know what you are doing) or there
% will be problems with numbering.
%    \begin{macrocode}
\newcommand{\altaffilmark}[1]{\textsuperscript{#1}}
\newcommand{\altaffiltext}[2]{%
  \if@firstaltaffil\gdef\ADS@affil{#2}\@firstaltaffilfalse\fi
  \stepcounter{footnote}%
  \protected@xdef\@thanks{\@thanks
    \protect\footnotetext[#1]{#2}}} 
%    \end{macrocode}
% \end{macro}
% \end{macro}
% \end{macro}
% \end{macro}
%
% \begin{macro}{\keywords}
% \begin{macro}{\@keywords}
% \begin{macro}{\@keywordheader}
% \label{keywords}Sundry minor preamble commands. All
% self-explanatory, I hope. |\keywords| is the old user interface
% (deprecated). New articles should use |\addkeyword| once for each
% keyword. That way the ADS record comes out right too (comma rather
% than em-dash separator).
% \changes{2.2beta12}{2001/12/07}{\cs{addkeyword} New user interface
%   for keywords.} 
% \changes{2.2beta13}{2002/03/11}{\cs{addkeyword} Improved behaviour
%   with .ads file. Use
%   \cs{ADS@sanitize} to change definitions of `~' and \cs{ion}.}
% \changes{3.23}{2007/09/09}{\cs{addkeyword} now signals a compilation
% error if the author tries to add more than 6 keywords.}
% \changes{3.24}{2007/09/14}{kwdcounter is now reset in
% \cs{@begofarticlehook} to prevent error in book style} 
%    \begin{macrocode}
\newcommand{\@kwdsep}{ --- }
\newcommand{\ADS@kwdsep}{, }
\newcounter{kwdcounter}\setcounter{kwdcounter}{0}
\newcommand{\addkeyword}[1]{%
  \addtocounter{kwdcounter}{1}
  \ifthenelse{\thekwdcounter > 6}{
    \ClassError{rmaa}{Too many keywords. A maximum of 6 keywords is allowed.}{}
  }{}
  \if@firstkwd%
  \renewcommand\@keywords{\if@boldkeywords\uppercase{#1}\else #1\fi}
  {\ADS@sanitize\def\ion{\plainion}\gdef\ADS@keywords{#1}}
  \@firstkwdfalse
  \else%
  \edef\@keywords{\@keywords\@kwdsep \if@boldkeywords\uppercase{#1}\else #1\fi}%
  {\ADS@sanitize\def\ion{\plainion}\xdef\ADS@keywords{\ADS@keywords\ADS@kwdsep #1}}%
  \fi}
\newcommand{\ADS@sanitize}{%
  \def\mathrm{}
  \def\@{\ }
  \def~{\ }
  \def\langle{<}
  \def\rangle{>}
  \def\log{log }
  \def\sin{sin }
  \def\cos{cos }
  \def\tan{tan }
  \def\fd{^d.}
  \def\fh{^h.}
  \def\fm{^m.}
  \def\fs{^s.}
  \def\fdg{�.}
  \def\farcm{'.}
  \def\farcs{''.}
  \def\arcdeg{�}
  \def\,{ }
  }
\newcommand{\keywords}[1]{%
  \renewcommand{\ADS@keywords}{#1}
  \renewcommand{\@keywords}{\uppercase{#1}}
  }
\newcommand{\@keywordheader}{\rmaa@KeywordHeaderStyle{Key Words:}}
\newcommand{\@keywords}{%
  \ClassWarningNoLine{rmaa}{No keywords specified}%
  Keywords go here
  }
%    \end{macrocode}
% \begin{macro}{\shorttitle}
% \begin{macro}{\shortauthor}
% \begin{macro}{\@shorttitle}
% \begin{macro}{\@shortauthor}
% Commands for setting the running headers.
%    \begin{macrocode}
\newcommand{\shorttitle}[1]{%
  \renewcommand\@shorttitle{\MakeUpperCase{#1}}}
\newcommand{\shortauthor}[1]{%
  \renewcommand\@shortauthor{\MakeUpperCase{#1}}}
\newcommand{\@shorttitle}{SHORT TITLE GOES HERE}
\newcommand{\@shortauthor}{SHORT AUTHOR GOES HERE}
%    \end{macrocode}
% \end{macro}
% \end{macro}
% \end{macro}
% \end{macro}
% \end{macro}
% \end{macro}
% \end{macro}


% \subsubsection{The full addresses}
% \label{sec:full-addresses}
% \begin{macro}{\fulladdresses}
% \begin{macro}{\@fulladdresses}
% \begin{macro}{\@fulladdheight}
% \begin{macro}{\@fulladddepth}
% \begin{macro}{\@fulladdbox}
% Aaghh, the full addresses! A perennial source of problems. 
%    \begin{macrocode}
\newlength{\@fulladdheight}
\newlength{\@fulladddepth}
\newsavebox{\@fulladdbox}
%% \fulladdresses command modified (11 May 2000) so that we can put
%% \fulladdresses{} in the article to suppress the output of the
%% addresses
\newcommand{\@FinalFigure}{\relax}
\newcommand{\FinalFigure}[1]{\renewcommand{\@FinalFigure}{#1}}
\newcommand{\fulladdresses}[1]{%
  \renewcommand{\@rmaaArg}{#1}
  \ifx\@rmaaArg\@rmaaEmpty
     \renewcommand{\@fulladdresses}{}
  \else
     \renewcommand{\@fulladdresses}{
       \normalsize%
       \ifx\@FinalFigure\relax\else\vfill\@FinalFigure\vfill\fi
%%       \vspace*{1ex}
       %% do addresses as a list (copied from thebibliography environment) 
       \list{\null}{\leftmargin 1.5em\topsep\z@\labelwidth\z@\labelsep\z@\itemindent -1.5em
         \itemsep\z@\parsep\z@\usecounter{enumi}}
       \sloppy\clubpenalty4000\widowpenalty4000
       \sfcode`\.=1000\relax
       #1
       \endlist
       }
  \fi
  }
\newcommand{\@fulladdresses}{Full addresses go here}
%    \end{macrocode}
% \begin{macro}{\suppressfulladdresses}
% New user-level command for suppressing full addresses should you so
% desire. Currently used in the single-page-poster style. 
%    \begin{macrocode}
\newcommand{\suppressfulladdresses}{\@suppressfulladdressestrue\@needaddressesfalse}
%    \end{macrocode}
% \end{macro}
% \end{macro}
% \end{macro}
% \end{macro}
% \end{macro}
% \end{macro}

% \begin{macro}{\printdates}
% \begin{macro}{\noprintdates}
% Command to activate the printing out of the received/accepted
% dates. This can also be done using the \texttt{printdates} option,
% which is the default for \texttt{journal}, but not for \texttt{proceedings}.
%    \begin{macrocode}
\newcommand{\printdates}{\@printdatestrue}
\newcommand{\noprintdates}{\@printdatesfalse}
%    \end{macrocode}
% \end{macro}
% \end{macro}
% \begin{macro}{\@outputaddressesifnecessary}
% Command to output the addresses if we haven't already done so
% This is called by: 
% \begin{enumerate}
% \item The end of |thebibliography| environment
% \item The |\outputfulladdresses| command
% \end{enumerate}
%    \begin{macrocode}
\newcommand{\@outputaddressesifnecessary}{%
  \ClassInfo{rmaa}{Calling @outputaddressesifnecessary}
  \ifthenelse{\boolean{@needaddresses}}{%
    %% If in one column style, or in the first column of a two column
    %% style then we can try to put it on this page 
    \ifthenelse{\boolean{@firstcolumn} \or \not\boolean{@rmtwocol}}{%
      \@outputaddresses}{%
      %% otherwise give warning it will be on a separate page
      \ClassWarning{rmaa}{Full addresses will go on a separate
        page (please see User Guide for how to remedy this)}}%
    %% If we haven't managed yet, put it on its own page
    \ifthenelse{\boolean{@needaddresses}}{%
      \if@debug{\ClassWarning{rmaa}{In this silly bit!}}\fi
      \clearpage\vfill\@outputaddresses}{}%
    }{
    %% If we don't need them, we still might need to do the eoa hook 
    \ifthenelse{\boolean{@suppressfulladdresses}}{%
      \if@debug{\ClassWarning{rmaa}{XXX}}\fi
%%      \if@book\else\@endofarticlehook\null\fi
      }{}
    }%
  }  
%    \end{macrocode}
% \end{macro}

% \begin{macro}{\outputfulladdresses}
% A user interface to the above command, in case there is no biblio
% (11 May 2000) 
%    \begin{macrocode}
\newcommand{\outputfulladdresses}{\@outputaddressesifnecessary}
%    \end{macrocode}
% \end{macro}

% \begin{macro}{\adjustfinalcols}
% This command starts a new column and should be 
% used to ensure that the columns balance on the last page. It also
% outputs the full 
% author addresses to ensure that this occurs in the left hand
% column. 
% The |\par| command is so that the current line will be taken account
% of when working out |\pagetotal|
%    \begin{macrocode}
\newcommand{\adjustfinalcols}{\if@twocolumn\goodbreak\@outputaddresses\newpage\noindent\fi}
%    \end{macrocode}
% \end{macro}

% \begin{macro}{\@outputaddresses}
% Command to do the work of outputting the addresses. 
% This can be called by 
% \begin{enumerate}
% \item an explicit |\adjustfinalcols|
% \item via |\@outputaddressesifnecessary|, q.v.
% \end{enumerate}
%
%    \begin{macrocode}
\newcommand{\@outputaddresses}{%
  % Allow option of suppressing the full addresses, in which case
  % we just run the hook if we are not in a book style
  \if@suppressfulladdresses
%%     \if@book\else\@endofarticlehook\null\fi
  \else
  %% save the full addresses in a box for output later
  \sbox{\@fulladdbox}{\parbox{\textwidth}{\@fulladdresses}}
  % calculate height of address box  
  \settoheight{\@fulladdheight}{\usebox{\@fulladdbox}}
  \settodepth{\@fulladddepth}{\usebox{\@fulladdbox}}
  \addtolength{\@fulladdheight}{\@fulladddepth}
  \if@debug
  { \typeout{Address total height: \the\@fulladdheight}
    \typeout{Text height: \the\pagetotal}
    \typeout{Goal height: \the\pagegoal}
    }
  \fi
  % reuse \@fulladdheight to mean the total height of text+addresses
  \addtolength{\@fulladdheight}{\pagetotal}
%%  \addtolength{\@fulladdheight}{-\baselineskip}
  \if@debug\typeout{Adjusted text + addresses height: \the\@fulladdheight}\fi
  % check it will fit on current page
  \ifdim\@fulladdheight<\pagegoal
     %% must shift it to the left
     %% to compensate for the hanging indentation in the reference list. 
     %% NOTE: this is not foolproof!
     %% NB. This will give overfull hbox warnings, which is a bit ugly.
     \vfill\noindent\if@biblist\hspace*{-\leftmargin}\fi
     \usebox{\@fulladdbox}
     % Set this to show we've done it
     \@needaddressesfalse
     % and assume we are at the end of the article
     % so that we do the EOA hook if not in the book style
%%     \if@book\else\@endofarticlehook\fi
  \else
     % if not, just do nothing. With luck, this function will get 
     % called again on the next page
     \if@debug{\typeout{Holding over full addresses to next page}}\fi
  \fi\fi
  }
%    \end{macrocode}
% \end{macro}

% \subsubsection{Commands for editors}
% \label{sec:editor-commands}
% Some user/editor commands for doing the TOC and index in the book
% option. 
% They are defined here rather in the book-specific part since they
% need to be always available in order to avoid breaking documents. 
% \begin{macro}{\listofauthors}
% \begin{macro}{\@rawauthors}
% \begin{macro}{\rmaaTOCauthors}
% The command |\listofauthors| formats an author list for use in the
% TOC. It should contain the same names as in the |\author| command,
% but without the affiliation garbage and using |\&| instead of
% |and|. 
%    \begin{macrocode}
\newcommand{\@rawauthors}{Raw Authors}
\newcommand{\rmaaTOCauthors}{\noexpand\rmaaTOCAuthorStyle{Author list for TOC}}
\newcommand{\listofauthors}[1]{%
  \renewcommand{\@rawauthors}{#1}
  \renewcommand\rmaaTOCauthors{\noexpand\rmaaTOCAuthorStyle{#1}}}
%    \end{macrocode}
% \begin{macro}{\indexauthor}\label{indexauthor}
% The command |\indexauthor| produces an index entry for each
% author. 
% This command must come \emph{after} the
% |\title| and the |\listofauthors| commands. Also, there should be one 
% |\indexauthor| command for \emph{each} author. We also add the
% author onto the end of the |\ADS@authors| macro and construct the
% |\ADS@authorinitial| by a most cunning method (see below), but only
% for the first author. 
% 
% No, No, NO! This was all wrong. Backed out completely: 
% \begin{quote}
% Completely rewritten (19 Dec 2002) since we now use the index class, so 
% as to be able to alphabetize accented characters. A side effect of this is
% that |\index| no longer expands out the argument after the |@|, so we 
% instead add all the |\index| commands to a list (|\indexauthor@deferlist|), 
% which is executed in the |\@begofarticlehook|. 
% 
% Note the stuff with |\toks0| is copied from Exercise 20.15 of the \TeX{}book---it is
% to make sure the latin1-accented name gets passed into the idx
% file. 
% \end{quote}
% We now go back to the old way of doing things, using an optional
% argument to the \cs{indexauthor} command that has the name without
% accents. 
% \changes{3.08}{2003/01/16}{Remerged changes from 3.06 that got
% dropped from 3.07 after a synchronization disaster. Basically, the
% accentless optional argument to \cs{indexauthor}. Must add the macro
% directories to my unison set-up to stop this sort of thing from happening.}
%    \begin{macrocode}
\newcommand\@isortkey{}
\newcommand\@emptykey{empty}
\newcommand\indexauthor[2][empty]{
  \renewcommand\@isortkey{#1}
  \ifx\@isortkey\@emptykey \renewcommand\@isortkey{#2}\fi
  \index{\@isortkey @{{\rmaaIndexAuthorStyle{#2}}\hskip1em plus1em minus0.4em
      \rmaaIndexTitle\hskip1em plus2em minus0.4em\penalty-100
      \rmaaTOCauthors}}
  \indexauthorADS{#2}
}
%%\def\indexauthor@deferlist{\relax}
%%\long\def\indexauthor#1{
%%  \def\@@name{#1} \toks0=\expandafter{\@@name}
%%  \toks1=\expandafter{\indexauthor@deferlist}
%%  \protected@xdef\indexauthor@deferlist{%
%%    \the\toks1
%%    \noexpand\index{\the\toks0@{\noexpand\rmaaIndexAuthorStyle{\the\toks0}\noexpand\hskip1em plus1em minus0.4em 
%%        \rmaaIndexTitle\noexpand\hskip1em plus2em minus0.4em\noexpand\penalty-100
%%        \rmaaTOCauthors}}
%%  }
%%  \indexauthorADS{#1}
%%  }
%%\xdef\indexauthor@deferlist{}
\newcommand\indexauthorADS[1]{%
  \if@firstauthor\protected@xdef\ADS@authorinitial{\firstletter{#1}}\fi
  \protected@xdef\ADS@authors{\ADS@authors\if@firstauthor\else ; \fi
    #1}
  \@firstauthorfalse
}
%    \end{macrocode}
% Note that the first |#1| sets the key for the index, making sure it
% actually starts with  
% first letter of the surname. Meanwhile, the |@{...}| expression,
% which starts with some  
% control sequences is what is written to the |.ind| file. 
% See the makeindex man page for details. 
% \end{macro}
% \end{macro}
% \end{macro}
% \end{macro}

% \begin{macro}{\KillAllAccents}
% \begin{macro}{\firstletter}
% These are utility functions used in the above. |lett| is a counter
% indicating the current letter in the string. |\firstletter| just
% prints the first letter of its argument.
%    \begin{macrocode}
\newcommand{\KillAllAccents}{\def\'{}\def\`{}\def\~{}\def\.{}\def\i{i}}
\newcommand\firstletter[1]{\@car #1\@nil}
%    \end{macrocode}
% \end{macro}
% \end{macro}
% \begin{macro}{\SetVolume}
% \begin{macro}{\SetMSnumber}
% \begin{macro}{\SetFirstPage}
% \begin{macro}{\SetYear}
% \begin{macro}{\ReceivedDate}
% \begin{macro}{\AcceptedDate}
% These are for the editors to set the volume number, etc. 
%    \begin{macrocode}
\newcommand{\SetVolume}[1]{%
  \renewcommand{\@volume}{#1}
  \edef\ADS@volume{\PadToFourWithDots{#1}}
  \renewcommand{\@volumen}{\rmaa@VolumenStyle{volumen #1}}%
  }
\newcommand{\SetMSnumber}[1]{\renewcommand{\@MSnum}{#1}}
%    \end{macrocode}
% Note that there is no real reason to use |\SetFirstPage| since all
% the page numbering is done automatically in the book style. 
%    \begin{macrocode}
\newcommand{\SetFirstPage}[1]{\setcounter{page}{#1}}
\newcommand{\SetYear}[1]{%
  \renewcommand{\@year}{#1}
  \renewcommand{\ADS@year}{#1}
  \renewcommand{\@anyo}{\rmaa@VolumenStyle{#1}}
  }
\newcommand{\ReceivedDate}[1]{\renewcommand{\@received}{#1}}
\newcommand{\AcceptedDate}[1]{\renewcommand{\@accepted}{#1}}
%    \end{macrocode}
%
% \begin{macro}{\DOI}
% User/editor command to enter the DOI, which will be typeset as a link.
% \changes{4.0}{2019/11/02}{New user command \cs{DOI} to typeset
% DOI in the header of first page and add it to ADS info.}
%    \begin{macrocode}
\newcommand\DOI[1]{%
  \def\@DOI@url{\url{https://doi.org/10.22201/ia.01851101p.#1}}%
  \def\ADS@doi{https://doi.org/10.22201/ia.01851101p.#1}%
}%
%    \end{macrocode}
% \end{macro}
%
% \begin{macro}{\ORCIDiD}
% \changes{4.0}{2019/11/12}{Add mechanism for giving ORCID iD links}
% \begin{macro}{\ORCIDiDfull}
% \begin{macro}{\@orcid@icon}
% \begin{macro}{\@orcid@iconbig}
% Commands to add icon links to ORCID iDs of authors. For example
% \begin{quote}
% |\ORCIDiD|\marg{0000-0000-0000-0000}
% \end{quote}
% to be used after each author name in the |\author|
% command and in the |\fulladdresses|.  Alternatively,
% |\OORCIDiDfull| will typeset the icon plus the full text of the
% link, but we don't use that at present.  The icon image is stored in |ORCIDiD_iconvector.pdf|, which needs to be on the input path.  In principle, we could make EPS versions if necessary.
%    \begin{macrocode}
\newcommand\@orcid@icon{\includegraphics[height=1.8ex]{ORCIDiD_iconvector}}
\newcommand\@orcid@iconbig{\includegraphics[height=2ex]{ORCIDiD_iconvector}}
\newcommand\ORCIDiD[1]{\href{https://orcid.org/#1}{\,\@orcid@icon}}
\newcommand\ORCIDiDfull[1]{%
  \href{https://orcid.org/#1}{\@orcid@iconbig}\,\url{https://orcid.org/#1}}
%    \end{macrocode}
% \end{macro}
% \end{macro}
% \end{macro}
% \end{macro}
%
% \begin{macro}{\@volume}
% \begin{macro}{\@volumen}
% \begin{macro}{\@artpages}
% \begin{macro}{\@year}
% \begin{macro}{\@anyo}
% \begin{macro}{\@received}
% \begin{macro}{\@accepted}
% \begin{macro}{\@MSnum}
% Default values for the containers used by above commands, plus a
% few more useful ones. 
%    \begin{macrocode}
\newcommand{\@volume}{00}
\newcommand{\@volumen}{\rmaa@VolumenStyle{00}}
\newcommand{\@artpages}{%
  \oldpageref{\articlename:RMAAFirstPage}--\oldpageref{\articlename:RMAALastPage}}
\newcommand{\@year}{\the\year}
\newcommand{\@anyo}{\rmaa@VolumenStyle{2007}}
\newcommand{\@journalname}{\rmaa@JournalNameStyle{%
    \if@proceedings\RMAASClong\else\RMAAlong\fi}}
\newcommand{\@received}{}
\newcommand{\@accepted}{}
\newcommand{\@MSnum}{0000}
\newcommand{\@DOI@prefix}{https://doi.org/10.22201/ia.01851101p}
\def\@DOI@url{\url{https://doi.org/10.22201/ia.01851101p.2099.00.00.00}}
%    \end{macrocode}
% \end{macro}
% \end{macro}
% \end{macro}
% \end{macro}
% \end{macro}
% \end{macro}
% \end{macro}
% \end{macro}
% \end{macro}
% \end{macro}
% \end{macro}
% \end{macro}
% \end{macro}
% \end{macro}
% 
% The following commands are most useful in the \texttt{book}
% style. However, the |\SetConfTitle| command can now be used in
% individual articles to set the header for the \texttt{preprint}
% style. 
%    \begin{macrocode}
\newcommand{\@reunionplace}{Lugar de la reuni\'on}
\newcommand{\@reuniondate}{Fecha de la reuni\'on}
\newcommand{\@conftitle}{Conference Title}
\newcommand{\SetReunion}[2]{\renewcommand{\@reunionplace}{#1}\renewcommand{\@reuniondate}{#2}}
\newcommand{\SetConfTitle}[1]{\renewcommand{\@conftitle}{#1}}
\newcommand{\@editors}{Juan Ram\'on de la Fuente and ``El Mosh''}
\newcommand{\SetEditors}[1]{\renewcommand{\@editors}{#1}}
\newcommand{\@editores}{Juan Ram\'on de la Fuente y ``El Mosh''}
\newcommand{\SetEditores}[1]{\renewcommand{\@editores}{#1}}
\newcommand{\@mes}{mes de la publicaci\'on}
\newcommand{\SetMes}[1]{\renewcommand{\@mes}{\rmaa@VolumenStyle{#1}}}
%    \end{macrocode}

% \subsection{Notes for Author Proofs and Communication between Editors}
% \label{sec:notes}
% We introduce a command \cs{Note} that should be used to describe
% any change made to the manuscript. These are a bit tricky to get
% right. In general we use \cs{marginpar} but this only works in outer
% paragraph mode, so we just bung the box where we are if we are in a
% figure caption or something. If we are in the preamble, then we just
% use a \marginpar as part of the \maketitle command.
%    \begin{macrocode}
\newif\if@arenotes
\@arenotesfalse
\newcounter{notecounter}
\setcounter{notecounter}{0}
\newcommand\NotesList{}
\newcommand\Note[1]{%
  \if@proof{%
    \@arenotestrue%
    \stepcounter{notecounter}%
    \if@inpreamble\else{%
      \NoteMark{\arabic{notecounter}}%
    }\fi
    \protected@xdef\NotesList{\NotesList 
      \protect\item #1}%
  }\fi%
}
\newcommand\Changed[2]{Changed ``#1'' to ``#2''.}
\newcommand\NoteMark[1]{%
  \ifinner
  \raisebox{1ex}[0pt][0pt]{\makebox[0pt][l]{\framebox{\footnotesize\bfseries
        N#1}}}%
  \else
  \marginpar{\framebox{\footnotesize\bfseries N#1}}%
  \fi
}
\newcommand\RefNote[1]{\if@proof\NoteMark{\ref{#1}}\relax\fi}
\newcommand\NoteLabel{\protect\label}
\newcommand\OutputNotes{%
  \bgroup
  \onecolumn
  \section*{EDITORIAL NOTES}
  \begin{enumerate}
    \NotesList
  \end{enumerate}
  \egroup
}
\DeclareRobustCommand\cs[1]{\texttt{\char`\\#1}}

%    \end{macrocode}
% \subsection{Bibliographic Records for ADS}
% \label{sec:ADS}
% \changes{2.2beta11}{2001/11/27}{Just about finished all the ADS
% stuff now. Except for the references, which I have given up on for
% the time being.}
% \begin{macro}{\ADS@BibCode}
% \begin{macro}{\ADS@year}
% \begin{macro}{\ADS@journal}
% \begin{macro}{\ADS@volume}
% \begin{macro}{\ADS@firstpage}
% \begin{macro}{\ADS@authorinitial}
% We use the |\ADS@BibCode| command to hold the bibliographic code for
% ADS. It is 
% used in the |%R| 
% field of the bibliographic record for each article. 
% The format is 
% \begin{quote}
% \texttt{YYYYJJJJJVVVVMPPPPA}  
% \end{quote}
% where \texttt{YYYY} is the year, \texttt{JJJJJ} is the journal code,
% which is \texttt{RMxAA} or \texttt{RMxAC}, \texttt{VVVV} is the
% volume number (right-justified), \texttt{M} is not used by us,
% \texttt{PPPP} is the first page number (right-justified), and
% \texttt{A} is the initial of the first author's surname. The volume
% and page fields are padded with periods (`\texttt{.}'), which causes
% us some fun (this is done with the macro |\PadToFourWithDots| described
% below) . 
% 
% All these individual bits are held in the commands |\ADS@year|,
% |\ADS@journal|, |\ADS@volume|, |\ADS@firstpage|,
% |\ADS@authorinitial|. 
%    \begin{macrocode}
\newcommand{\SetADSBibCode}[1]{\renewcommand{\ADS@BibCode}{#1}}
\newcommand{\ADS@BibCode}
{\ADS@year\ADS@journal\ADS@volume.\PadToFourWithDots{\ADS@firstpage}\ADS@authorinitial}
\newcommand{\ADS@year}{2001} 
\newcommand{\ADS@journal}{\if@proceedings RMxAC\else RMxAA\fi}
\newcommand{\ADS@volume}{....} 
%    \end{macrocode}
% These two are just set to default values here, |\ADS@firstpage| is set
% in the |document| environment for included articles, in the
% |frontsection| environment, |\tableofcontents| command,  |index| and
% |abstracts| environments. Also by |\AtBeginDocument| for stand-alone
% articles. 
% 
% God knows how we are going to do |\ADS@authorinitial|!
%    \begin{macrocode}
\newcommand{\ADS@firstpage}{....}
\newcommand{\ADS@authorinitial}{A}
%    \end{macrocode}
% \end{macro}
% \end{macro}
% \end{macro}
% \end{macro}
% \end{macro}
% \end{macro}
% We also have a whole slew of other fields that can be set. These are
% all listed in Table~\ref{tab:ADS}. Most of them still have to be
% defined. 
% \begin{table}[!ht]\raggedleft
% \caption{ADS Bibliographic Records}\label{tab:ADS}
% \begin{tabular}{lp{8em}lp{8em}}\hline
% Code & Description & Macro Name & Comment \\ \hline
% |%R| & Bibliographic Code & |\ADS@BibCode| & required \\
% |%A| & Author list & |\ADS@authors| & required, generate from |\indexauthor| \\
% |%F| & Author affiliation & |\ADS@affil| & format: AA(\dots)
% AB(\dots) \dots \\
% |%J| & Journal Name & |\ADS@journalname| &  \\
% |%D| & Publication Date & |\ADS@fulldate| &  \\
% |%L| & Last Page & |\ADS@lastpage| & generate from label \\
% |%T| & Title & |\ADS@title| & required, generate from |\title| \\
% |%C| & Abstract Copyright & |\ADS@copyright| & set to IA-UNAM \\
% |%U| & URL for electronic document & |\ADS@url| & set to PDF file \\
% |%d| & DOI for article & |\ADS@doi| & set by |\DOI| \\
% |%K| & Keywords & |\ADS@keywords| & set equal to |\@keywords| \\
% |%A| & Abstract Text & |\ADS@abstract| & set equal to |\@abstract| \\
% |%R| & References & |\ADS@references| & generate from the
% bibliography  \\ \hline
% \end{tabular}
% \end{table}
% A lot of these macros are set by other commands, as mentioned in the
% Table, but some need to be set  explicitly here. 
% 
% \begin{macro}{\ADS@doi}
% \changes{4.0}{2019/11/08}{Save DOI for ADS}
% The user/editor command \cs{DOI} is used to set this. 
%    \begin{macrocode}
\newcommand{\ADS@doi}{https://doi.org/10.22201/ia.01851101p.0000.00.00.00}
%    \end{macrocode}
% \end{macro}
% 
% \begin{macro}{\ADS@journalname}
% This
% does not contain any accents at the moment so as not to cause
% problems---this maybe should be changed though. 
%    \begin{macrocode}
\newcommand{\ADS@journalname}{%
  \if@proceedings\@conftitle\space (Eds. \@editors) \fi%
  Revista Mexicana de Astronom\'\i{}a y Astrof\'\i{}sica%
  \if@proceedings\space (Serie de Conferencias)\fi%
  \space Vol. \@volume, pp. \ADS@firstpage--\ADS@lastpage\space (\@year)%
  \space (http://www.astroscu.unam.mx/\~{}rmaa/)
  }
%    \end{macrocode}
% \end{macro}
% \begin{macro}{\ADS@keywords}
% \begin{macro}{\ADS@abstract}
% \begin{macro}{\ADS@title}
% |\ADS@abstract| is straightforward, since we already have the text
% saved as |\@abstract|. We could do the same with |\ADS@keywords| but
% to avoid unwanted formatting we explicitly set it in |\keywords|
% (p.~\pageref{keywords}). |\ADS@title| is just set to |\@rawtitle|,
% which we were already using for another purpose. 
%    \begin{macrocode}
\let\ADS@title=\@rawtitle
\newcommand{\ADS@keywords}{No keywords specified}
\newcommand{\ADS@abstract}{No abstract specified}
%    \end{macrocode}
% \end{macro}\end{macro}\end{macro}
% \begin{macro}{\ADS@authors}
% |\ADS@authors| wants to be in ``Other, A. N.'' format, but with
% semi-colons between authors. The most natural way to do this seems to
% be with |\indexauthor| (p.~\pageref{indexauthor}). We append each
% author to this macro, which means we must explicitly clear it for
% each new article (done in |\article|). .   
%    \begin{macrocode}
\newcommand{\ADS@authors}{}
%    \end{macrocode}
% \end{macro}
% \begin{macro}{\ADS@fulldate}
% \begin{macro}{\ADS@lastpage}
% \begin{macro}{\ADS@copyright}
% \begin{macro}{\ADS@website}
% \begin{macro}{\ADS@url}
% \begin{macro}{\SetWebsite}
% \begin{macro}{\SetMonthNumeric}
% Other sundry data, together with user commands for bits that can be
% set by the Editors. The |\ADS@url| is constructed from the website
% (currently |http://www.astrosmo.unam.mx/rmaa/|,
% the journal descriptor (e.g., |RMxAC..10/|), followed by the
% path to the PDF file (e.g., |PDF/RMxAC..10_HEW1.pdf|). 
% 
%    \begin{macrocode}
\newcommand{\ADS@website}{http://www.astrosmo.unam.mx/rmaa/}
\newcommand{\ADS@url}
{\ADS@website\ADS@journal\ADS@volume\ADS@number/PDF/\ADS@journal\ADS@volume\ADS@number\_\articlename.pdf}
\newcommand{\ADS@number}{}
\newcommand{\ADS@fulldate}{00/\ADS@year} 
\newcommand{\ADS@lastpage}{00} 
\newcommand{\ADS@copyright}{(C) \ADS@year; Instituto de Astronom\'\i{}a,
  Universidad Nacional Aut\'onoma de M\'exico} 
\newcommand{\SetNumber}[1]{\renewcommand{\ADS@number}{-#1}}
\newcommand{\SetWebsite}[1]{\renewcommand{\ADS@website}{#1}}
\newcommand{\SetMonthNumeric}[1]{\renewcommand{\ADS@fulldate}{#1/\ADS@year}}
%    \end{macrocode}
% \end{macro}\end{macro}\end{macro}\end{macro}\end{macro}\end{macro}\end{macro}
% \begin{macro}{\ADS@affil}
% \begin{macro}{\ADS@references}
% These two are somewhat more problematical. Affiliation currently
% only works for the first author (see \dots) pending redesign of the
% user interface. References are now done by wrapping the
% |thebibliography| environment to open another input stream on the
% same file and search through for the |\bibitems|. 
%    \begin{macrocode}
\newcommand{\ADS@references}{}
\newcommand{\ADS@affil}{}
%    \end{macrocode}
% \end{macro}\end{macro}
% 
% Define two macros that expand to the strings
% |\begin{thebibliography}| and |\end{thebibliography}|. Note that the
% idea for this and much else is inspired by the code for the docstrip
% package. Also define |\StartBibX| in case authors use the more
% standard |\begin{thebibliography}{}| syntax. 
%    \begin{macrocode}
\edef\EndBib{\expandafter\@gobble\string\\end\string{thebibliography\string}}
\edef\StartBib{\expandafter\@gobble\string\\begin\string{thebibliography\string}}
\edef\StartBibX{\expandafter\@gobble\string\\begin\string{thebibliography\string}\string{\string}}
%    \end{macrocode}
% Define two new logical variables. |\ifAnotherline| is true if we
% want to read another line from the input stream. It is set to false
% when we encounter the |\end{thebibliography}|. |\ifInBib| is set to
% true when we frind the |\begin{thebibliography}| and tells us to
% start saving the following lines. 
%    \begin{macrocode}
\newif\ifAnotherLine
\newif\ifInBib
%    \end{macrocode}
% \begin{macro}{makeOther}
% We need a macro to change the catcodes of lots of things to other. 
%    \begin{macrocode}
\def\makeOther#1{\catcode`#1=12\relax}
%    \end{macrocode}
% \end{macro}
% \begin{macro}{\RereadBiblio}
% Define a macro to do all the fiddly stuff. This is used at the
% beginning of |thebibliography| environment (Sec~\ref{sec:biblio}).
%    \begin{macrocode}
\newread\tf@biblio
\newcommand{\RereadBiblio}{%
%    \end{macrocode}
% Open the current file again for reading
%    \begin{macrocode}
  \begingroup
  \immediate\openin\tf@biblio=\articlename
%    \end{macrocode}
% Everything in a group to be safe. Initialize logicals. Black magic
% to change
% catcodes and avoid blank lines being interpreted as |\par|. 
%    \begin{macrocode}
  \AnotherLinetrue\InBibfalse
  \makeOther\ \makeOther\\\makeOther\$%
  \makeOther\#\makeOther\^\makeOther\^^K%
  \makeOther\_\makeOther\^^A\makeOther\%%
  \makeOther\~\makeOther\{\makeOther\}\makeOther\&%
  \endlinechar-1\relax
%    \end{macrocode}
% Start looping through the file. Read each line to the macro
% |\Line|. If we have got to EOF or end of biblio then set
% |\AnotherLinefalse|.
%    \begin{macrocode}
  \loop
  \read\tf@biblio to \Line
  \ifx\Line\EndBib\AnotherLinefalse\InBibfalse\fi
  \ifeof\tf@biblio \AnotherLinefalse\fi
%    \end{macrocode}
% If we are in the biblio and the line is not blank then append it to
% the |\ADS@references| macro together with a newline. 
%    \begin{macrocode}
  \ifInBib
  \ifx\Line\empty\else\xdef\ADS@references{\ADS@references\Line^^J}\fi
  \fi
%    \end{macrocode}
% Check if the line we just saw was the start of biblio, then repeat
% the loop if necessary. 
%    \begin{macrocode}
  \ifx\Line\StartBib\InBibtrue\fi
  \ifx\Line\StartBibX\InBibtrue\fi
  \ifAnotherLine\repeat
  \immediate\closein\tf@biblio
  \endgroup
  }
%    \end{macrocode}
% \end{macro}
% 
% \begin{macro}{\PadToFourWithDots}
% Utility function for left-padding a number to 4 characters with
% periods. Used in 
% constructing |\ADS@BibCode|. This is a very simple algorithm and
% \emph{requires} that the argument actually be a number---if it is
% not then dire consequences will result.  
%    \begin{macrocode}
\newcommand{\PadToFourWithDots}[1]{%
  \ifnum #1 > 999 #1%
  \else\ifnum #1 > 99 .#1%
  \else\ifnum #1 > 9 ..#1%
  \else ...#1%
  \fi\fi\fi}
%%  \ifnum \count0 > 999 #1%
%%  \else\ifnum \count0 > 99 .#1%
%%  \else\ifnum \count0 > 9 ..#1%
%%  \else ...#1%
%%  \fi\fi\fi}
%    \end{macrocode}
% \end{macro}
%
% \begin{macro}{\AuthorADS}
% \begin{macro}{\ADS@affiliations}
% New mechanism for setting the affiliations for ADS. The author labels (AA, AB, \dots) are generated using two nested counters.  We use the |\robustify| command from \texttt{etoolbox} so that accented characters in the affiliation addresses do not blow up.  The list of accents may need to be expanded. 
%    \begin{macrocode}
\newcounter{AUTHOR}
\newcounter{author}[AUTHOR]
\stepcounter{AUTHOR}
\newcommand\ADS@affiliations{}
\robustify{\'}\robustify{\i}\robustify{\~}\robustify{\`}
\def\AuthorADS#1#2{%
  \begingroup
  \def\ORCIDiD##1{<ID system="ORCID">##1</ID>}
  \def\Email##1{<EMAIL>##1</EMAIL>}
  \def\Affil##1{##1}
  \stepcounter{author}
  \ifthenelse{\value{author} > 26}{\stepcounter{AUTHOR}\stepcounter{author}}{}
  \ifthenelse{\equal{\ADS@affiliations}{\@rmaaEmpty}}{%
    \xdef\ADS@affiliations{\Alph{AUTHOR}\Alph{author}(#2)}
  }{%
    \xdef\ADS@affiliations{\ADS@affiliations, \Alph{AUTHOR}\Alph{author}(#2)}
  }%
  \endgroup
}
%    \end{macrocode}
% \end{macro}
% \end{macro}
% %
% \begin{macro}{\ADS@dumpdata}
% \changes{4.0}{2019/11/13}{Write all author affiliations for ADS}
% \begin{macro}{\ADS@write}
% \begin{macro}{\ADS@openfile}
% |\ADS@dumpdata| is the macro that actually does the job of writing out all the
% data to an auxiliary file, using |\ADS@write| to write each
% field. The file is set up with |\ADS@openfile|, which currently is
% done in the TOC of a book.  The new multi-author affiliations are used if |\ADS@affiliations| has been set, otherwise it falls back to the single-author version |\ADS@affil|. 
%    \begin{macrocode}
\newcommand{\ADSfield}[2]{\relax}
\newcommand{\ADSvfield}[2]{\relax}
%%\renewcommand{\ADSfield}[12]{\%#1 #2}
\newcommand{\ADS@dumpdata}
{ 
  \ADS@write{\string\begin{ADSarticle}}
  \ADS@WRITE{\string\ADSfield{R}{\ADS@BibCode}}
  \ADS@write{\string\ADSfield{A}{\ADS@authors}}
  \ifthenelse{\equal{\ADS@affiliations}{\@rmaaEmpty}}{%
    {\ADS@sanitize\ADS@write{\string\ADSfield{F}{AA(\ADS@affil)}}}
  }{%
    {\ADS@sanitize\ADS@write{\string\ADSfield{F}{\ADS@affiliations}}}
  }
  \ADS@write{\string\ADSfield{J}{\ADS@journalname}}
  \ADS@write{\string\ADSfield{D}{\ADS@fulldate}}
  \ADS@write{\string\ADSfield{L}{\ADS@lastpage}}
  {\ADS@sanitize\ADS@write{\string\ADSfield{T}{\ADS@title}}}
  \ADS@write{\string\ADSfield{C}{\ADS@copyright}}
  \ADS@write{\string\ADSfield{d}{\ADS@doi}}
  \ADS@write{\string\ADSfield{U}{\ADS@url}}
  \ADS@write{\string\ADSfield{K}{\ADS@keywords}}
  {\ADS@sanitize\ADS@write{\string\ADSfield{B}{\ADS@abstract}}}
  \ADS@write{\string\ADSfield{Z}{\ADS@references}}
  \ADS@write{\string\end{ADSarticle}}
  }
%%\newcommand{\ADS@write}[1]{\if@book\addtocontents{ads}{#1}\fi}
%%\newcommand{\ADS@openfile}{\@starttoc{ads}}
\newcommand{\ADS@write}[1]{\set@display@protect\immediate\write\tf@ads{#1}}
\newcommand{\ADS@WRITE}[1]{\immediate\write\tf@ads{#1}}
\newcommand{\ADS@openfile}{%
  \newwrite\tf@ads%
  \immediate\openout\tf@ads\jobname.ads\relax
  \ADS@WRITE{\ADS@preamble}}
\newcommand{\ADS@closefile}{\ADS@WRITE{\ADS@postamble}}
\def\ADS@preamble{%
  %%  \string\documentclass{article}
  %%  \string\begin{document}
  }
\def\ADS@postamble{%
  %%  \string\end{document}
  \string%\string% Local Variables:^J 
  \string%\string% mode: LaTeX^J
  \string%\string% mode: Font-lock^J
  \string%\string% End:^J
  }
%    \end{macrocode}
% \end{macro}
% \end{macro}
% \end{macro}
% 
% \subsection{Automatic calculation of the page range}
% \label{sec:page-range}
% Another issue that has cost me a lot of grief. Still, it finally
% seems to be sorted now. All hail Jeff Goldberg. 
%
% Cunning trick to get the page numbers from |\pageref|.  Put a label
% on the first and last page of the document using |\AtBeginDocument|
% and |\AtEndDocument|. .  This way won't work for the book option.
% In that case, we do it by defining |\@endofarticlehook| and
% |\@begofarticlehook| as parallels to |\AtEndDocument| and
% |\AtBeginDocument|.  Note that we have to defer adding
% |\@begofarticlehook| to |\AtBeginDocument| to make sure it comes
% after the redefinitions of |\ref| and |\label|.
%    \begin{macrocode}
\if@book\AtEndDocument{\ADS@closefile}\else
 %\AtBeginDocument{\@begofarticlehook}
\AtEndDocument{\@endofarticlehook}
\fi
%    \end{macrocode}
% \begin{macro}{\@begofarticlehook}
% Since we are getting more and more things that need to be done at
% the beginning of every article, frontsection, etc we define a sister
% hook to 
% |\@endofarticlehook|. This hook does the following: 
% \begin{enumerate}
% \item Labels the first page
% \item Sets |\ADS@firstpage|
% \item runs and then clears |\indexauthor@deferlist|
% \item resets the keyword counter
% \end{enumerate}
% It is executed in the following circumstances: 
% \begin{enumerate}
% \item At |\begin{document}| for standalone files
% \item At |\begin{document}| for each included file
% \item At beginning of front sections 
% \item At beginning of TOC
% \item At beginning of Author Index
% \item At beginning of abstracts
% \end{enumerate}
% For some reason, this is not always within the right group, so we
% use |\xdef| just to be sure. 
%    \begin{macrocode}
\newcommand{\@begofarticlehook}{%
  \if@frontmatter\frontlabel{RMAAFirstPage}\else\label{RMAAFirstPage}\fi%
  \xdef\ADS@firstpage{\@arabic{\c@page}}%
  \if@book\else\makeindex\fi % needed so default index is defined in stand-alone articles
%%  \indexauthor@deferlist
%%  \xdef\indexauthor@deferlist{}
  \@inpreamblefalse
  \setcounter{kwdcounter}{0}
  }
%    \end{macrocode}
% \end{macro}
% \begin{macro}{\@endofarticlehook}
% This is a hook to be run at the end of each article
% currently, it just labels the last page and writes out all the ADS
% data to an auxiliary file.  
% It should not produce any printed output. 
% It is modelled on the |\AtEndDocument| command, to which it reduces
% in the case of a single article. 
%    \begin{macrocode}
\newcommand{\@endofarticlehook}{%
  \if@debug{\ClassWarning{rmaa}{***Doing Hook***}}\fi
  \if@debug{\ClassWarning{rmaa}{ADSabstract is: \ADS@abstract}}\fi
  \clearpage\rmaa@lastpagelabel\if@book\ADS@dumpdata\fi
  \ifx\NotesList\@rmaaEmpty\else\OutputNotes\fi
  \if@debug{\ClassWarning{rmaa}{Notes list is: \NotesList}}\fi
}
%    \end{macrocode}
% \begin{macro}{\rmaa@lastpagelabel}
% New version of last page label command copied from Jeff Goldberg's
% lastpage.sty package - we can't just use |\label| here in case the
% output routine for the last page has already been called. Instead,
% we first explicitly clear the page, step back the page counter,
% explicitly write to the .aux file, then restep the page counter
% again.
%
% But this will conflict with the hyperref package, which does labels differently.  Compensate for that by using the five-argument form.
% 
% Also set the |\ADS@LastPage| macro to the appropriate number. 
%    \begin{macrocode}
\newcommand{\rmaa@lastpagelabel}{%
  \addtocounter{page}{-1}%
  \immediate\write\@auxout{\string
    \newlabel{\articlename:RMAALastPage}{{}{\thepage}{}{}{}}}%
  \edef\ADS@lastpage{\thepage}%
  \addtocounter{page}{1}}
%    \end{macrocode}
% \end{macro}
% \end{macro}

% \subsection{Deal with explicit change in number of columns}
% \label{sec:columns}
% Sometimes, when formatting a book, we need to switch between
% one-column articles and two-column articles. For this to work, we
% need to wrap the standard \LaTeX{} commands to set the |\@rmtwocol|
% flag appropriately and also reset the package default values of all
% scalable lengths. .  
%    \begin{macrocode}
\let\latexonecolumn=\onecolumn
\let\latextwocolumn=\twocolumn
\renewcommand{\onecolumn}{\@rmtwocolfalse\@package@default@lengths\latexonecolumn}
\renewcommand{\twocolumn}{\@rmtwocoltrue\@package@default@lengths\latextwocolumn}
%    \end{macrocode}
% 
% \subsection{The actual formatting of the title}
% \label{sec:title-format}

% Note, we have to redefine the |\maketitle| command to use the
% numbered footnotes. This means we needn't use the atbeginend stuff. 
%    \begin{macrocode}
\renewcommand{\maketitle}{\par
  \begingroup
%    \end{macrocode}
% Don't switch to marked footnotes. 
% This is consistent with existing style but I'm not sure I like it
%    \begin{macrocode}
  %%  \renewcommand\thefootnote{\fnsymbol{footnote}}%
%%  \def\@makefnmark{\hbox to\z@{$\m@th^{\@thefnmark}$\hss}}%
  \def\@makefnmark{\textsuperscript{\@thefnmark}\hspace{-0.2em}}%
  \long\def\@makefntext##1{\parindent 1em\noindent
    \hbox to1.8em{\hss$\m@th^{\@thefnmark}$}##1}%
  \if@rmtwocol
  { \if@debug{\ClassWarning{rmaa}{TWO COLUMN}}\fi
    \twocolumn[\@maketitle]%
    }
  \else
  { \if@debug{\ClassWarning{rmaa}{ONE COLUMN}}\fi
    \newpage
    \global\@topnum\z@   % Prevents figures from going at top of page.
    \if@manuscript \@MSmaketitle \else \@maketitle \fi
    }
  \fi
  \if@manuscript\else\thispagestyle{rmaafirstpage}\@thanks\fi
  \ifx\NotesList\@rmaaEmpty\else{%
    \NoteMark{1--\arabic{notecounter}}\vspace*{-2\baselineskip}%
  }\fi%
  \endgroup
%    \end{macrocode}
% Don't reset footnote counter (not sure whether this is good thing
% or not).
%    \begin{macrocode}
  %%  \setcounter{footnote}{0}%
  \let\thanks\relax\let\affil\relax\let\altaffil\relax
  \let\maketitle\relax\let\@maketitle\relax
  \gdef\@thanks{}\gdef\@author{}\gdef\@title{}}
%    \end{macrocode}

% \begin{macro}{\@kwhwidth}
% \begin{macro}{\@kwtwidth}
% Redefine \cs{@maketitle} to include the abstract, resumen and
% keywords. Note that in the manuscripty style we use
% \cs{@MSmaketitle} instead.  
% \changes{3.16}{2003/04/07}{
%  Added extra vertical space between authors and accepted date.}
% \changes{3.26}{2011/08/15}{Add extra label to review articles}
%    \begin{macrocode}
\newlength{\@kwhwidth}\newlength{\@kwtwidth}
\newcommand\reviewname{Review}
\renewcommand{\@maketitle}{%
  \newpage
  \null
  \if@reviewarticle
  \smallskip
  {\large\bfseries\itshape\reviewname}%
  \bigskip
  \else
  \vskip\@before@title@skip%
  \fi
  \begin{center}%
    {\@title \par}%
    \vskip\@before@author@skip%
    \parbox{\rmaa@AuthorWidthFrac\textwidth}{\centering
      \lineskip\@between@author@skip%
      \@author
      \par}
    %% this should be replaced with `Submitted to...' for preprints
    \if@printdates
    \vskip\@before@dates@skip%
    \if@preprint
    \rmaa@ReceivedStyle{Submitted to RMxAA \@received}
    \else
    \rmaa@ReceivedStyle{Received \@received; accepted \@accepted}
    \fi
    \fi
    \if@manuscript
    \rmaa@ReceivedStyle{Submitted to RMxAA \@received}
    \fi
    \if@guide
    \rmaa@ReceivedStyle{Version \rmaa@version, \rmaa@date}
    \fi
    \vskip\@before@abstract@skip%
    \if@abstractpresent
    \parbox{\rmaa@AbstractWidthFrac\textwidth}{%
      \setlength{\parindent}{0.9cm}%
      \centerline{\rmaa@AbstractHeaderStyle{Abstract}}\par%
      \vskip\@after@abstracthead@skip%
      \if@newstyle\noindent\fi\@abstract
      \vskip\@between@abstract@skip%
      \centerline{\rmaa@AbstractHeaderStyle{Resumen}}\par%
      \vskip\@after@abstracthead@skip%
      \if@newstyle\noindent\fi\@resumen
      \if@setresumen\else\@abstract\fi % append english abstract if
                                       % spanish one not set
      \if@guide\else            %do not use keywords for guide
      \vskip\@before@kwd@skip%
      % Format the Key Words with hanging indentation:
      % They are put in a \parbox with width equal to the line width
      % minus the width of the header. 
      { % Widths of keyword
        \settowidth{\@kwhwidth}{\@keywordheader~}   % header and
        \setlength{\@kwtwidth}{\linewidth}          % keyword text
        \addtolength{\@kwtwidth}{-\@kwhwidth}       %
        \noindent\@keywordheader~\parbox[t]{\@kwtwidth}{%
          \rmaa@KeywordStyle{\boldmath\uppercase{\@keywords}}}}%
      \fi
      }
      \fi
  \end{center}%
  \vspace*{\afterkwdskip}
  \par
  }
\newlength{\afterkwdskip}\setlength{\afterkwdskip}{\z@}
%    \end{macrocode}
% \end{macro}
% \end{macro}
%
%
% \begin{macro}{\reviewarticle}
% \begin{macro}{\reviewtoc}
% \begin{macro}{\reviewlof}
%   User command to specify a review article. Must be used in
%   preamble. Main effect is to place the text ``Review'' before the
%   article title.  Also add command for setting a table of contents
%   and a list of figures at the beginning of the article.
%    \begin{macrocode}
\newcommand{\reviewarticle}{\@reviewarticletrue}
\newcommand\reviewtoc{%
  \section*{CONTENTS}%
  \@starttoc{rtc}%
}
\newcommand\reviewlof{%
  \section*{LIST OF FIGURES}%
  \@starttoc{lof}%
}
%    \end{macrocode}
% \end{macro}
% \end{macro}
% \end{macro}
% \begin{macro}{\l@section}
%    \begin{macrocode}
\renewcommand*\l@section[2]{%
  \ifnum \c@tocdepth >\z@
    \addpenalty\@secpenalty
    \addvspace{0.8em \@plus\p@}%
    \setlength\@tempdima{1.5em}%
    \begingroup
      \parindent \z@ \rightskip \@pnumwidth
      \parfillskip -\@pnumwidth
      \leavevmode %\bfseries
      \advance\leftskip\@tempdima
      \hskip -\leftskip
      #1\nobreak\hfil \nobreak\hb@xt@\@pnumwidth{\hss #2}\par
    \endgroup
  \fi}
%    \end{macrocode}
% \end{macro}
%
% \begin{macro}{\@MSmaketitle}
% A special version of |\@maketitle| for the \texttt{manuscript}
% option. This can deal with a multi-page abstract/resumen. Not very
% well, mind. 
%    \begin{macrocode}
\newcommand{\@MSmaketitle}{%
  \if@debug\ClassInfo{rmaa}{Using special MS maketitle command}\fi
  \newpage
  {
    \thispagestyle{rmaafirstpage}
%     \setlength{\parindent}{0pt}
    \null\vskip 2ex
    \begin{center}
      \@title\par
      \@author\par
      \rmaa@ReceivedStyle{%
        \ifx\@received\@rmaaEmpty
          Draft version: \@date
        \else
          Received: \@received
          \ifx\@accepted\@rmaaEmpty
          \else
            ; Accepted: \@accepted
          \fi
        \fi}\par
    \end{center}
    \@thanks
    \centerline{\rmaa@AbstractHeaderStyle{Resumen}}
    \noindent\@resumen\par\vskip 2ex
    \centerline{\rmaa@AbstractHeaderStyle{Abstract}}
    \noindent\@abstract\par\vskip 2ex
    \noindent\@keywordheader~\rmaa@KeywordStyle{\@keywords}\par
%     \@keywordheader~\rmaa@KeywordStyle{\uppercase{\@keywords}}\par
  }%
  }
%    \end{macrocode}
% \end{macro}

% \subsection{The sectioning commands}
% \label{sec:section-commands}

% \begin{macro}{\section}
% \begin{macro}{\subsection}
% \begin{macro}{\subsubsection}
% \begin{macro}{\paragraph}
% \begin{macro}{\subparagraph}
% Redefine all the standard commands to use these lengths. 
%    \begin{macrocode}
\renewcommand\section{\@startsection {section}{1}{\z@}%
                                   {\abovesectionskip}%
                                   {\belowsectionskip}%
                                   %% how do we get upper case here
                                   {\centering\reset@font\normalsize}}
\renewcommand\subsection{\@startsection{subsection}{2}{\z@}%
                                     {\abovesectionskip}%
                                     {\belowsectionskip}%
                                     {\centering\reset@font\normalsize\itshape}}
\renewcommand\subsubsection{\@startsection{subsubsection}{3}{\z@}%
                                     {\abovesectionskip}%
                                     {\belowsectionskip}%
                                     {\centering\reset@font\normalsize\itshape}}
\renewcommand\paragraph{\@startsection{paragraph}{4}{\z@}%
                                    {\aboveparaskip}%
                                    {\belowparaskip}%
                                    {\reset@font\normalsize\bfseries}}
\renewcommand\subparagraph{\@startsection{subparagraph}{5}{\parindent}%
                                      {\aboveparaskip}%
                                      {\belowparaskip}%
                                      {\reset@font\normalsize\bfseries}}
%    \end{macrocode}
% \end{macro}
% \end{macro}
% \end{macro}
% \end{macro}
% \end{macro}

% \begin{macro}{secnumdepth}
% Make sure that we number the subsubsections---this is done by
% default in the |article| class but not the |book| class, which is
% why we must set it explicitly here.
%    \begin{macrocode}
\setcounter{secnumdepth}{3}
%    \end{macrocode}
% \end{macro}
% \changes{2.2beta10}{2001/11/07}{Fixed \cs{secnumdepth} for \cs{book}
% style.}

% \begin{macro}{\thesubsection}
% \begin{macro}{\thesubsubsection}
% For subsection and subsubsection, we want the numbers to be roman
% but the title to be italic. This seems the easiest way to go about it.
%    \begin{macrocode}
\renewcommand{\thesubsection}{{\upshape\thesection.\@arabic\c@subsection}}
\renewcommand{\thesubsubsection}{\thesubsection{\upshape.\@arabic\c@subsubsection}}
%    \end{macrocode}
% \end{macro}
% \end{macro}

% \begin{macro}{\@seccntformat}
% Add ``.'' to the number here so we can remove it from |\thesection|,
% etc. Reduce gap between number and title WJH 27 Jan 1995
%    \begin{macrocode}
\def\@seccntformat#1{\csname the#1\endcsname .\hskip 0.4em}
%    \end{macrocode}
% \end{macro}

% \begin{macro}{\@sect}
% I didn't really want to mess with the |\@sect| command, but it was
% the only way I could manage to force upper case. Other than the
% commented parts, it is identical to the standard version. 
%    \begin{macrocode}
\def\@sect#1#2#3#4#5#6[#7]#8{
  \ifnum #2>\c@secnumdepth
     \let\@svsec\@empty
  \else
     \refstepcounter{#1}%
     \let\@@protect\protect
     \def\protect{\noexpand\protect\noexpand}%
     \edef\@svsec{\@seccntformat{#1}}%
     \let\protect\@@protect
  \fi
%    \end{macrocode}
% This is the bit I changed (27 Jan 1995). Use upper case for
% |\section|, use mixed case for the rest. 
%    \begin{macrocode}
  \ifnum #2=1                   
     \def\sec@upcase##1{\uppercase{##1}}%
     \def\sec@shape{\upshape}%
  \else                         
     \def\sec@upcase##1{\relax{##1}}%
     \def\sec@shape{\itshape}%
  \fi
  \@tempskipa #5\relax
  \ifdim \@tempskipa>\z@
     \begingroup #6\relax
          \@hangfrom{\hskip #3\relax\@svsec}%
                    {\interlinepenalty \@M \sec@upcase{#8}\par}%
     \endgroup
     \csname #1mark\endcsname{#7}%
%    \end{macrocode}
% Only write TOC entries if doing the Guide style. 
%    \begin{macrocode}
     \if@guide 
        \addcontentsline{toc}{#1}{%
          \ifnum #2>\c@secnumdepth \else
             \protect\numberline{\csname the#1\endcsname}%
          \fi
          #7}%
     \fi
%    \end{macrocode}
% Also allow per-article TOC for review papers
%    \begin{macrocode}
     \if@reviewarticle
        \ifnum #2>2 \else
           \addcontentsline{rtc}{#1}{%
             \protect\numberline{\csname the#1\endcsname}%
             \sec@upcase{\sec@shape #7}}%
        \fi
     \fi
  \else
     \def\@svsechd{#6\hskip #3\relax  %% \relax added 2 May 90
       \@svsec #8\csname #1mark\endcsname{#7}%
       \if@guide
          \addcontentsline{toc}{#1}{%
            \ifnum #2>\c@secnumdepth \else
               \protect\numberline{\csname the#1\endcsname}%
            \fi
            #7}%
       \fi
       }%
  \fi
  \@xsect{#5}}
%    \end{macrocode}
% \end{macro}


% \subsection{Appendices}
% \label{sec:appendices}
% Make sure that appendices in individual articles work properly. We
% copy the definition from the article class (the book one doesn't do
% what we want).  
% \begin{environment}{appendices}
%   This uses an environment |appendices| inside of which each
%   |\section| will become a separate appendix.  Note that the
%   appendices must come before the references (otherwise the full
%   addresses mechanism will not work properly). We make the equation
%   numbering be subordinate to the appendix (equations A12, B6,
%   etc). Finally, we have a medskip followed by APPENDICES to
%   introduce them. \changes{3.17}{2002/04/08}{ Added appendices environment
%   that does the right thing.}
%    \begin{macrocode}
\newenvironment{appendices}{\par
  \begingroup
  \setcounter{section}{0}%
  \setcounter{subsection}{0}%
  \renewcommand\thesection{\@Alph\c@section}
  \renewcommand\theequation{\thesection\arabic{equation}}
  \medskip
  \section*{APPENDICES}}
{\endgroup}
%    \end{macrocode}
% \end{environment}
% \begin{environment}{appendix}
%   If there is only one appendix, then use |appendix| environment
%   instead.  \changes{4.0}{2019/11/12}{ Added
%   appendix environment for case where there is only one section.}
%    \begin{macrocode}
\renewenvironment{appendix}{\par
  \begingroup
  \setcounter{section}{0}%
  \setcounter{subsection}{0}%
  \renewcommand\thesection{\@Alph\c@section}
  \renewcommand\theequation{\thesection\arabic{equation}}
  \medskip
  \section*{APPENDIX}}
{\endgroup}
%    \end{macrocode}
% \end{environment}

% \subsection{Floating thingies}
% \label{sec:floats}   

%% 
%    \begin{macrocode}
\renewcommand{\thetable}{\@arabic\c@table}
\renewcommand{\thefigure}{\@arabic\c@figure}
\renewcommand{\figurename}{Fig.}
\renewcommand{\tablename}{TABLE}
%\renewcommand{\fnum@table}{\MakeUpperCase{\tablename}~\thetable}
%    \end{macrocode}
% For some reason, I have to write `TABLE' explicitly here
% If I use |\tablename|, it gets converted to mixed case---God knows
% why. 
% |\figurename| is also being ignored. 
%    \begin{macrocode}
\renewcommand{\fnum@table}{TABLE~\thetable}
\renewcommand{\fnum@figure}{Fig.~\thefigure}
%    \end{macrocode}

% Do captions in 9 pt face in the main journal style and in new-style
% anything. 
%    \begin{macrocode}
\long\def\@makefigcap#1#2{%
  \if@proceedings\else\small\fi
  \if@newstyle\small\fi 
  \vskip\abovecaptionskip
  \sbox\@tempboxa{#1. #2}%
  \ifdim \wd\@tempboxa >\hsize
    #1. #2 \par
  \else
    \global \@minipagefalse
    \hb@xt@\hsize{\hfil\box\@tempboxa\hfil}%
  \fi
  \vskip\belowcaptionskip}
%    \end{macrocode}

% Make sure table caption is upper case. 
%    \begin{macrocode}
\long\def\@maketabcap#1#2{%
  \hbox to\hsize{\hfil #1 \hfil}
  \vskip 6pt 
  \centering\MakeUpperCase{#2}\par
  \vskip 6pt}
%    \end{macrocode}

%    \begin{macrocode}
\newlength{\@rottabcapwidth}
\long\def\@makerottabcap#1#2{%
  \setbox\@tempboxa\hbox{\MakeUpperCase{#2}}%
  \ifdim\wd\@tempboxa >.8\vsize
     \setlength\@rottabcapwidth{.8\vsize}
  \else
     \setlength\@rottabcapwidth{\wd\@tempboxa}
  \fi
  \rotatebox{90}{%
    \begin{minipage}{\vsize}\centering
      \begin{minipage}{\@rottabcapwidth}\centering
        #1\\ \vskip 6pt \MakeUpperCase{#2}
      \end{minipage}%
    \end{minipage}
    }%
  \hskip 6pt
  }

\newenvironment{Sideways}
    {\begin{sideways}\begin{minipage}{\vsize}\centering}
    {\end{minipage}\end{sideways}}

\newenvironment{Sidewaystable}
{%
  \let\oldcaption=\caption%
  \def\caption{
    \def\@captype{table}\renewcommand{\@makecaption}{\@maketabcap}
    \oldcaption}
  \begin{sidewaystable}\centering}
  {\end{sidewaystable}\let\caption=\oldcaption}
\newenvironment{Sidewaystable*}
{%
  \let\oldcaption=\caption%
  \def\caption{
    \def\@captype{table}\renewcommand{\@makecaption}{\@maketabcap}
    \oldcaption}
  \begin{sidewaystable*}\centering}
  {\end{sidewaystable*}\let\caption=\oldcaption}
%    \end{macrocode}

% Not sure why we need to stick these lines in. 
%    \begin{macrocode}
\renewenvironment{figure}
                 {\renewcommand{\@makecaption}{\@makefigcap}\@float{figure}}
                 {\end@float}
\renewenvironment{figure*}
                 {\renewcommand{\@makecaption}{\@makefigcap}\@dblfloat{figure}}
                 {\end@dblfloat}
\renewenvironment{table}
                 {\renewcommand{\@makecaption}{\@maketabcap}\@float{table}}
                 {\end@float}
\renewenvironment{table*}
                 {\renewcommand{\@makecaption}{\@maketabcap}\@dblfloat{table}}
                 {\end@dblfloat}
%    \end{macrocode}

% Redefine |\@caption| so that we don't bother with the list of
% figures. This saves us having to worry about fragile commands in the
% captions. 
%    \begin{macrocode}
\long\def\@caption#1[#2]#3{%
  \par
%%  \addcontentsline{\csname ext@#1\endcsname}{#1}%
%%    {\protect\numberline{\csname the#1\endcsname}{\ignorespaces #2}}%
  \begingroup
    \@parboxrestore
    \if@minipage
      \@setminipage
    \fi
    \normalsize
    \@makecaption{\csname fnum@#1\endcsname}{\ignorespaces #3}\par
  \endgroup}
%    \end{macrocode}

% Make tables not quite so cramped. Allow |\hline\hline| to give a
% poor man's thick rule. 
%    \begin{macrocode}
\renewcommand{\arraystretch}{1.15}
\setlength{\doublerulesep}{0.4pt}
%    \end{macrocode}

% \begin{macro}{\tablecols}
% \begin{macro}{\@tablecols}
% \begin{macro}{\tabcaptioncontinued}
% |\tablecols| is a user command for setting the number of columns. It
% is used by commands that try and span the entire table (I think
% there must be a better way of acheiving
% this). |\tabcaptioncontinued| is 
% a user command for setting the caption of a continued table on the
% next page. 
%    \begin{macrocode}
\newcommand{\tablecols}[1]{\renewcommand{\@tablecols}{#1}}
\newcommand{\@tablecols}{1}
\newcommand{\tabcaptioncontinued}{%
  \multicolumn{\@tablecols}{c}{\addtocounter{table}{-1}TABLE~\thetable\ (CONTINUED)} \\ \\
  }
%    \end{macrocode}
% \end{macro}
% \end{macro}
% \end{macro}

% \subsubsection{Footnotes in tables}
% \label{sec:table-footnotes}
% \begin{macro}{\tablenotemark}
% \begin{macro}{\tablenote}
% \begin{macro}{\tablenoteskip}
% Old user commands |\tablenotemark{MARK}|,
% |\tablenote{MARK}{TEXT}|---now deprecated. 
%    \begin{macrocode}
\newcommand\tablenotemark[1]{$^\mathrm{#1}$} 
\newcommand\tablenote[2]{\if@proceedings\else\small\fi ${}^\mathrm{#1}$ #2} 
\newcommand\tablenoteskip{\\[0.5\baselineskip]} 
%    \end{macrocode}
% \end{macro}
% \end{macro}
% \end{macro}

% \begin{macro}{\tabnotemark}
% \begin{macro}{\tabnotetext}
% \begin{macro}{\tabnotetextTWO}
% Alternative way of doing footnotes in tables. 
%    \begin{macrocode}
\newcommand{\tabnotemark}[1]{\textsuperscript{\lowercase{#1}}}
%\newcommand{\tabnotetextTWO}[4]{\tabnotetext{#1}{#2}\tabnotetext{#3}{#4}}
\newcommand{\tabnotetextTWO}[4]{%
  \multicolumn{\@tablecols}{l}{%
    \parbox{0.5\tabnotewidth}{\small\textsuperscript{#1}#2}%
    \parbox{0.5\tabnotewidth}{\small\textsuperscript{#3}#4}%
    }%
  \\}
\newcommand{\tabnotetext}[2]{%
  \multicolumn{\@tablecols}{@{}p{\tabnotewidth}@{}}{\small\textsuperscript{#1}#2}\\}
\newlength{\tabnotewidth}
%    \end{macrocode}
% \end{macro}
% \end{macro}
% \end{macro}


% \begin{macro}{changemargin}
%    \begin{macrocode}
% Command to allow margins to be changed on the fly. Useful to allow
% tables and Figures to protrude into the margins. 
%
% \changes{3.24}{2007/09/14}{Added \texttt{changemargin} environment
% (copied from TeX FAQ).}
\newenvironment{changemargin}[2]{% 
  % From TeX FAQ http://www.tex.ac.uk/cgi-bin/texfaq2html?label=chngmargonfly
  \begin{list}{}{%
    \setlength{\topsep}{0pt}%
    \setlength{\leftmargin}{#1}%
    \setlength{\rightmargin}{#2}%
    \setlength{\listparindent}{\parindent}%
    \setlength{\itemindent}{\parindent}%
    \setlength{\parsep}{\parskip}%
  }%
  \item[]}{\end{list}}
%    \end{macrocode}
% \end{macro}

% \subsection{The bibliography and references}
% \label{sec:biblio}
% \begin{macro}{\beforebiblioskip}
% \begin{macro}{\afterbiblioskip}
% I'm not entirely sure why I introduced these two commands, since
% directly putting |\vspace| before and after
% |\begin{thebibliography}| seems to work well enough. 
%    \begin{macrocode}
\newlength{\beforebiblioskip}
\setlength{\beforebiblioskip}{0pt}
\newlength{\afterbiblioskip}
\setlength{\afterbiblioskip}{0pt}
%    \end{macrocode}
% \end{macro}
% \end{macro}
% \begin{macro}{thebibliography}
% The environment itself is done as starred section with extra
% vertical space after heading. This space is currently zero. 
% In the book style we do the cunning trick with rereading the file in
% order to get the references into the |\ADS@references| macro.
% 
% \changes{3.05}{2003/01/02}{Removed the \cs{beforebiblioskip} and
% \cs{afterbiblioskip} commands, since they cause unwanted vertical
% space in some instances.}  
%    \begin{macrocode}
\renewenvironment{thebibliography}{%
  \if@book\RereadBiblio\fi
%%  \vspace*{\beforebiblioskip}
  \section*{REFERENCES} \frenchspacing
%%  \vspace*{\afterbiblioskip}
  \small
%    \end{macrocode} 
% The references themselves are set as a list in order 
% to get the hanging indentation. They are done in a smaller font and
% with french spacing so we needn't worry about the distinction
% between sentence-ending periods and abbreviations. 
%    \begin{macrocode} 
  \list{\null}{\leftmargin 1.5em\labelwidth\z@\labelsep\z@\itemindent -1.5em
    \itemsep\z@\parsep\z@\usecounter{enumi}}
  \sloppy\clubpenalty4000\widowpenalty4000
  \sfcode`\.=1000\relax\@biblisttrue}
{\endlist\@biblistfalse\normalsize\@outputaddressesifnecessary}
%    \end{macrocode}
% \end{macro}
% \subsection{Uniquification of cross-references}
% \label{sec:uniquification}
% 
% \begin{macro}{\articlename}
% In order that the |\label{}| command in individual articles be
% guaranteed to give a unique anchor, we append the article name (held
% in the |\articlename| command) to each key. This requires
% redefinition of the following commands: |\cite|, |\nocite|, |\ref|,
% |\pageref|, |\label|. 
% \begin{macro}{\oldcite}
% \begin{macro}{\oldnocite}
% \begin{macro}{\oldref}
% \begin{macro}{\oldpageref}
% \begin{macro}{\oldlabel}
% The original versions of these commands are
% saved as |\oldcite|, |\oldnocite|, etc. 
%
% The uniquifying of cross-references is now always done, so 
% we don't get multiply-defined labels in the single article case
%    \begin{macrocode}
\newcommand{\articlename}{anon}
\AtBeginDocument{
  \let\oldcite=\cite 
  \let\oldnocite=\nocite 
  \let\oldref=\ref 
  \let\oldpageref=\pageref
  \let\oldlabel=\label
}
%    \end{macrocode}
% \end{macro}
% \end{macro}
% \end{macro}
% \end{macro}
% \end{macro}
% \end{macro}
%
% Uniquify all cross-references by prepending the article name.  This
% now has to go into |\AtBeginDocument| so it is working on the new
% versions that hyperref.sty has set up.  Then we need to put the
% |\@begofarticlehook| afterwards, so that the first page label is using
% the new definitions.  If amsmath is loaded, then we also need to patch
% the special label command that they use. \changes{4.0}{2019/11/18}{Fix labels for amsmath display equations, in case that package has been loaded.}
%    \begin{macrocode}
\AtBeginDocument{
  \renewcommand{\ref}[1]{\oldref{\articlename:#1}%
    \if@debug{\ClassInfo{rmaa}{Reference to \articlename:#1}}\fi}
  \renewcommand{\pageref}[1]{\oldpageref{\articlename:#1}%
    \if@debug{\ClassInfo{rmaa}{Page reference to \articlename:#1}}\fi}
  \renewcommand{\label}[1]{\oldlabel{\articlename:#1}%
    \if@debug{\ClassInfo{rmaa}{Label for \articlename:#1}}\fi}%
  \ifdef{\label@in@display}{%
    \let\oldlabel@in@display=\label@in@display
    \renewcommand{\label@in@display}[1]{\oldlabel@in@display{\articlename:#1}%
      \if@debug{\ClassInfo{rmaa}{AMSTeX Label for \articlename:#1}}\fi}%
  }{}
  \@begofarticlehook
}%
%    \end{macrocode}
% Citations are more complicated since they can contain a list of
% keys. These are really only here for use with the rmaacite package,
% which is utterly deprecated (02 Jan 2003) now that I have the natbib
% code working. The command \cs{cite} should \emph{not} be used with
% natbib: use \cs{citep}, \cs{citet}, etc.\@ instead. 
%    \begin{macrocode}
\renewcommand{\nocite}[1]{\@bsphack
  \@for\@citeb:=#1\do{%
    \edef\@citeb{\expandafter\@firstofone\@citeb:\articlename}%
    \if@filesw\immediate\write\@auxout{\string\citation{\@citeb}}\fi
    \@ifundefined{b@\@citeb:\articlename}{\G@refundefinedtrue
      \@latex@warning{Citation `\@citeb' undefined}}{}}%
  \@esphack}
\renewcommand{\cite}[1]{\oldcite{#1:\articlename}}
%   \end{macrocode}


% Horrible fiddling with the bibliography citation mechanism. 
% This has to be globally available in case rmaacite puts it in the
% \texttt{.aux} file. 
%    \begin{macrocode}
\def\lbibcite#1#2#3{\global\@namedef{b@#1}{#2}
                    \global\@namedef{by@#1}{#3}}
%    \end{macrocode}

% \begin{macro}{\oldbibitem}
% \begin{macro}{\@rmaaBibArg}
% \begin{macro}{\@rmaaSpace}
% \begin{macro}{rmaabib}
% Of course, |\bibitem| needs to get the uniquification treatment as
% well. While we are at it, we redefine it to use consecutive integers
% for the key if the user leaves it blank (using the counter
% |rmaabib|). This allows one to put 
% |\bibitem{}| for all your references without getting annoying
% messages about multiply-defined labels.
% 
% We also take the opportunity to insert the optional argument to
% \cs{bibitem} if it is missing in the form Author(DATE) so that natbib
% doesn't try to do numerical citations. Natbib already provides a
% hook \cs{@extra@b@citeb} to help uniquify the citation key, so we
% just plug into that. However, it puts it \emph{after} the original
% key, not before as I had had it. Easiest solution is to change the
% way that I do it. 02 Jan 2003. Note that the \cs{@extra@b@citeb}
% does not get written to the .aux  file in the \cs{citation} command
% but it does work. 
% \changes{3.14}{06 Feb 2003}{Reinstalled old definition of \cs{bibitem}
%  for the case where the nonatbib option is requested, so that the
%  rmaacite package will work. This means that you can't use rmaacite
%  and natbib in the same volume.}
%    \begin{macrocode}
\if@natbib\def\@extra@b@citeb{:\articlename}\fi
%%\else
\let\oldbibitem\bibitem
\newcommand{\@rmaaBibArg}{}
\newcommand{\@rmaaSpace}{ }
\newcounter{rmaabib}
\if@natbib
\renewcommand{\bibitem}[2][ABC(2003)]{%
  \renewcommand{\@rmaaBibArg}{#2}
  \ifx\@rmaaBibArg\@rmaaEmpty
     \addtocounter{rmaabib}{1}
     \oldbibitem[#1]{\thermaabib:\articlename}
  \else\ifx\@rmaaBibArg\@rmaaSpace
     \addtocounter{rmaabib}{1}
     \oldbibitem[#1]{\thermaabib:\articlename}
  \else 
     \oldbibitem[#1]{#2:\articlename}
  \fi\fi
  \renewcommand{\@rmaaBibArg}{}
  }
\else
\renewcommand{\bibitem}[1]{%
  \renewcommand{\@rmaaBibArg}{#1}
  \ifx\@rmaaBibArg\@rmaaEmpty
     \addtocounter{rmaabib}{1}
     \oldbibitem{\articlename:\thermaabib}
  \else\ifx\@rmaaBibArg\@rmaaSpace
     \addtocounter{rmaabib}{1}
     \oldbibitem{\articlename:\thermaabib}
  \else
     \oldbibitem{\articlename:#1}
  \fi\fi
  \renewcommand{\@rmaaBibArg}{}
  }
\fi
%%\fi
%    \end{macrocode}
% \end{macro}
% \end{macro}
% \end{macro}
% \end{macro}

% \begin{macro}{oldpcite}
% \begin{macro}{oldscite}
% Then, we need to repeat the above scheme, but for the |rmaacite|
% definition. 
%    \begin{macrocode}
\newcommand{\@redefinebibitem}{%
  \let\oldbibitem\bibitem
  \let\oldcite\cite
  \let\oldpcite\pcite
  \let\oldscite\scite
  \def\bibitem[##1<##2>]##3{%
    \renewcommand{\@rmaaBibArg}{##3}
    \ifx\@rmaaBibArg\@rmaaEmpty
       \addtocounter{rmaabib}{1}
       \oldbibitem[##1<##2>]{\thermaabib:\articlename}
    \else\ifx\@rmaaBibArg\@rmaaSpace
       \addtocounter{rmaabib}{1}
       \oldbibitem[##1<##2>]{\thermaabib:\articlename}
    \else 
       \oldbibitem[##1<##2>]{##3:\articlename}
    % trust the author here
%%       \oldbibitem[##1<##2>]{##3}
    \fi\fi
    }
  }
%    \end{macrocode}
% \end{macro}
% \end{macro}


% \subsection{HTML Table of Contents}
% \label{sec:hyper-TOC}
% \begin{macro}{\HyperTOCEntry}
% \begin{macro}{\HyperTOCpart}
% \begin{macro}{\HyperTOCmeeting}
% \begin{macro}{\HyperTOCeditors}
% \begin{macro}{\HyperTOCpreamble}
% These commands are used to write information to the \texttt{.htoc}
% file, which is used in generating a Web Page of the Table of
% Contents. Since they only ever appear preceded by |\protect| in the
% argument to an |\addtocontents| command,  they should never need to
% be expanded by the macros. However, we define them here just in case
% \dots
%    \begin{macrocode}
\newcommand{\HyperTOCEntry}[3]{\relax}
\newcommand{\HyperTOCpart}[1]{\relax}
\newcommand{\HyperTOCmeeting}[3]{\relax}
\newcommand{\HyperTOCeditors}[1]{\relax}
\newcommand{\HyperTOCpreamble}{\relax}
%    \end{macrocode}
% \end{macro}
% \end{macro}
% \end{macro}
% \end{macro}
% \end{macro}
%

% \subsection{Formatting the poster abstracts} 
% These can now be used without using the book option.
%    \begin{macrocode}

\newcommand{\@makeminititle}{%
  \newpage
  \null 
  \vskip 3em%
  \begin{center}%
    {\@title \par}%
  \end{center}
  \vskip 3em%
  }

\newif\if@abstractsads
\@abstractsadsfalse
\newcommand\abstractsADS{\@abstractsadstrue}

\if@book\RequirePackage{multicol}\fi
\newenvironment{abstracts}[1][Abstracts of Contributed Papers]{%
  %% First, make sure we start on an odd page (06 Jun 2000)
  \if@book\clearpage\thispagestyle{empty}\cleardoublepage\fi
  %%
  \renewcommand{\articlename}{ABSTRACTS}
  \title{#1}\shortauthor{Abstracts}\shorttitle{Abstracts}
  \listofauthors{}
  \if@book
  \addcontentsline{toc}{chapter}{\protect\@fiddleauthors{\@TOCtitle}{\rmaaTOCauthors}}%
  \addtocontents{htc}{\protect\HyperTOCEntry{\@rawtitle}{\@rawauthors}{\articlename}}
  \fi
  \pagestyle{rmaa} % this now seems necessary (06 Jun 2000)
  \thispagestyle{rmaafirstpage}
%%  \twocolumn[\@makeminititle]
  \renewcommand{\abstract}[1]{%
    \renewcommand\ADS@abstract{##1}%
    \renewcommand{\@abstract}{##1}\@writeabstract%
  }
  \renewcommand{\title}[2][]{%
    \gdef\ADS@authors{}%
    \gdef\ADS@affil{}%
    \gdef\ADS@references{}%
    \gdef\ADS@keywords{}%
    \renewcommand{\@rmaaArg}{##1}
    \ifx\@rmaaArg\@rmaaEmpty
       \renewcommand{\@TOCaftertitle}{}
    \else
       \renewcommand{\@TOCaftertitle}{ (\@rmaaArg)}
    \fi
    \renewcommand{\@title}{\rmaa@TOCTitleStyle{##2}}%
    \renewcommand{\@TOCtitle}{\rmaa@TOCTitleStyle{##2}\@TOCaftertitle}%
    \renewcommand{\rmaaIndexTitle}{\rmaa@IndexTitleStyle{##2\@TOCaftertitle.}}%
    \renewcommand\ADS@title{##2}%
    \xdef\ADS@firstpage{\@arabic{\c@page}}%
    }
  \renewcommand{\author}[1]{\let\@thanks\@empty\renewcommand\@author{##1} }
  \renewcommand{\footnotetext}[2][XX]{\par\noindent\textsuperscript{##1}~##2}
  \renewcommand{\footnoterule}{\kern-3\p@ \hrule \@width 3em \kern 2.6\p@}
  \if@book\@begofarticlehook\fi
  \if@book\begin{multicols}{2}[\@makeminititle]\else\@makeminititle\fi
  \begingroup
  \setcounter{footnote}{0}
  }% and at the end of the environment....
  {\endgroup\if@book\end{multicols}\listofabstracts\label{RMAALastPage}\clearpage\fi}
%% auxiliary commands used:
%% \@writeabstract does all the actual formatting 
\newcommand{\@writeabstract}{%
  \if@book%
  \addcontentsline{loa}{abstract}{\protect\@fiddleauthors{\@TOCtitle}{\rmaaTOCauthors}}%
  \fi%
  \begin{center}\@title\par\@author\end{center}
  \@abstract
  \par\noindent\footnoterule\nobreak\footnotesize\@thanks
  \par\noindent\hrulefill\goodbreak
  \vskip 1cm \@plus0.5cm
  %% write ADS entry
  \if@book\if@abstractsads\ADS@dumpdata\fi\fi%
  \endgroup % this pair is to make each abstract a self-contained group
  \begingroup
  \setcounter{footnote}{0}
  }

\if@book
%% List of abstracts (stuff written to .loa file)
\newcommand{\listofabstracts}{%
  \begin{multicols}{2}[%
    \vskip 1.5cm \vfill % leave at least 1.5 cm and go to bottom of page if possible
    \centerline{\MakeUpperCase{List of Abstracts}}
    \vskip 1.5cm \@plus0.5cm
    ][6cm] % this makes it start a new page if there is not at least 6cm of space
    \renewcommand{\@fiddleauthors}[2]{##1\hskip1em plus2em minus0.4em\penalty-100  ##2}
    \parindent\z@
%%    \parskip.5\p@ \@plus .7\p@\relax 
    \parskip3\p@ \@plus 1\p@ \@minus 1\p@\relax % allow quite a lot of
                                % leeway here to avoid breaks inside
                                % entries 
    \renewcommand{\\}{\relax} % temporarily turn off explicit line breaks just in case
    \renewcommand{\altaffilmark}[1]{\relax} % ditto with this
    \@starttoc{loa}
  \end{multicols}
  }
\newcommand*{\l@abstract}[2]{%
  \par\penalty-1000
  #1\hskip 1em\penalty0\null\nobreak\hfill#2}
\fi

%    \end{macrocode}
 
% \subsection{Code for typesetting an entire volume of conference
% proceedings} \label{sec:whole-volume}
%    \begin{macrocode}

\if@book
%% environment for including the individual articles
\newcommand{\@fiddleauthors}[2]{\relax}
\renewcommand{\theequation}{\@arabic\c@equation}
\newenvironment{articles}%
{ \begingroup % make it a group so that author definitions do not
              % spill over
  \gdef\ADS@authors{}
  \pagestyle{rmaa}
  \renewenvironment{document}{%
    \@begofarticlehook
    \addtocounter{chapter}{1}%
    \addcontentsline{toc}{chapter}{%
      \protect\@fiddleauthors{\@TOCtitle}{\rmaaTOCauthors}}%
    \addtocontents{htc}{\protect\HyperTOCEntry{\@rawtitle}{\@rawauthors}{\articlename}}
    }{\relax}
  \renewcommand{\documentclass}[2][foo]{\relax} 
  \renewcommand{\usepackage}[2][foo]{%
    \def\@packagename{##2}\def\@rmaacitename{rmaacite}%
    \ifx\@packagename\@rmaacitename
    \@rmcitetrue\ClassInfo{rmaa}{This paper uses rmaacite}
    \makeatletter\input{rmaacite.sty}\makeatother % read it in here so that all changes
    \else                                         % are local to paper
    \if@debug\ClassWarning{rmaa}{Package \@packagename\space was requested by \articlename\space}\fi
    \fi} 
  %% Don't allow individual articles to mess with the pages, volume or year
  \renewcommand{\SetFirstPage}[1]{\relax}
  \renewcommand{\SetYear}[1]{\relax}
  \renewcommand{\SetVolume}[1]{\relax}
  \renewcommand{\SetConfTitle}[1]{\relax}
  %% Don't allow individual articles to change these
  \renewcommand{\nonstopmode}{\relax}
  \renewcommand{\scrollmode}{\relax}
  \renewcommand{\outputplates}{\relax}
  %% XXXX need to check if using rmaacite and act accordingly
  %% avoid multiple declarations of the same length
  \let\oldnewlength=\newlength
  \renewcommand{\newlength}[1]{%
    \edef\reserved@a{\expandafter\@gobble\string ##1}%
    \@ifundefined{\reserved@a}{%
      \oldnewlength{##1}}{%
      \if@debug{\ClassInfo{rmaa}{\reserved@a\space already defined: doing nothing}}\fi%
      \relax}}%
  \renewcommand\newcommand{\providecommand}
  }%
{\endgroup\onecolumn}
     

%    \end{macrocode}
% Command to be used in `articles' environment for including an
% article
%    \begin{macrocode}
\newcommand{\article}[2][]{%
  % optional first argument is extra text to put in brackets after title in TOC
  % e.g. ``Invited Talk'' or ``Poster''
  \renewcommand{\articlename}{#2}
  % reset loads of counters
  \setcounter{footnote}{0}
  \setcounter{figure}{0}
  \setcounter{table}{0}
  \setcounter{section}{0}
  \setcounter{equation}{0}
  \setcounter{rmaabib}{0}
  \setcounter{author}{0}
  \setcounter{AUTHOR}{1}
  \gdef\ADS@authors{}%
  \gdef\ADS@affil{}%
  \gdef\ADS@affiliations{}%
  \gdef\ADS@references{}%
  \@firstauthortrue
  \@firstaffiltrue
  \@firstaltaffiltrue
  \@firstkwdtrue
  \@rmcitefalse
  \if@suppressfulladdresses\else\@needaddressestrue\fi
  \renewcommand{\@rmaaArg}{#1}
  \ifx\@rmaaArg\@rmaaEmpty
     \renewcommand{\@TOCaftertitle}{}
  \else
     \renewcommand{\@TOCaftertitle}{ (#1)}
  \fi
%    \end{macrocode}
% In the main journal articles must start on an odd page. 
%    \begin{macrocode}
  \if@journal\clearpage\thispagestyle{empty}\cleardoublepage\fi
%% Now input the file, then run the end-of-article hook (reinstated 28
  %% Sep 2001)
  {\input{\@articledir/\articlename.tex}\@endofarticlehook}}
%    \end{macrocode}
% Command to separate articles into sessions. Write an entry in the
% TOC, makes sure we start on an odd page (added 06 Jun 2000), then do
% a title page if \cs{SessionTitlePages} was called. 
% \changes{3.12}{04 Feb 2003}{Added option to have title page for each
% session. Turn it on with \cs{SessionTitlePages}.}
%    \begin{macrocode}
\newcommand{\session}[1]{
  %% this does nothing except write an entry in the TOC
  \addcontentsline{toc}{part}{\rmaa@TOCPartStyle{#1}}
  \addtocontents{htc}{\protect\HyperTOCpart{#1}}
  %% and also makes sure we start on an odd page (added 06 Jun 2000)
  \clearpage\thispagestyle{empty}\cleardoublepage
  \stepcounter{part}
  \if@sessiontitlepages{%
    \onecolumn
    \thispagestyle{empty}
    \renewcommand{\articlename}{Section\thepart TitlePage}
    \label{RMAAFirstPage}
    \make@session@title{#1}
    \label{RMAALastPage}
    \clearpage\thispagestyle{empty}\cleardoublepage
  }\fi
  }
\newcommand{\make@session@title}[1]{%
  {%
    \setlength\parindent{0pt}
    \vspace*{\stretch{1}}
    \rule{\linewidth}{1mm}
    \begin{flushright}\Huge\bfseries\boldmath
      #1
    \end{flushright}
    \rule{\linewidth}{1mm}
    \vspace*{\stretch{2}}
  }
}
%    \end{macrocode}
% Command to turn on separate title pages for each session. 
%    \begin{macrocode}
\newcommand{\SessionTitlePages}{\@sessiontitlepagestrue}
\newcommand{\BoldSessionTitles}{%
  \renewcommand{\rmaa@TOCPartStyle}[1]{\textbf{\rmaa@TOCTitleStyle{##1}}}
  }
\newcommand{\ArticleDir}[1]{\renewcommand{\@articledir}{#1/}\graphicspath{{1/}}}
\newcommand{\@articledir}{.}
%\newcommand{\ps@rmaafront}{\ps@rmaa}% page style for front matter (same as ever, actually)
\newcommand{\ps@rmaafront}{% page style for front matter:  page # in each foot
  \renewcommand\@oddfoot{\hfil{\thepage}}%
  \renewcommand\@evenfoot{{\thepage}\hfil}%
  % shorttitle in both heads
  \renewcommand\@evenhead{\hfil\@shorttitle\hfil}%
  \renewcommand\@oddhead{\@evenhead}
  }
\newcommand{\ps@rmaafrontfirstpage}{% page style for first page of front matter
  % page # in each foot, on outside edge
  \renewcommand\@oddfoot{\hfil\thepage}%
  \renewcommand\@evenfoot{\thepage\hfil}%
  % nothing in the even heads
  \renewcommand\@evenhead{\relax}%
  \renewcommand\@oddhead{\relax}%
  }
%% slight modification to the standard \frontmatter and \mainmatter
%% commands, just to set the \pagestyle
\renewcommand\frontmatter{%
  \latexonecolumn
  \cleardoublepage
  \@frontmattertrue
  \@mainmatterfalse
  \pagenumbering{roman}
  \pagestyle{rmaafront}
  \addtocounter{page}{2} % leave space for the title page
  }
\renewcommand\mainmatter{%
  \pagestyle{empty}
  \if@rmtwocol\twocolumn\fi
  \cleardoublepage
  \@frontmatterfalse
  \@mainmattertrue
  \pagenumbering{arabic}
  \pagestyle{rmaa}
  }
\renewcommand\backmatter{%
  \clearpage\pagestyle{empty}\cleardoublepage
  \@mainmatterfalse
  %% put extra space in the TOC
  \addtocontents{toc}{\bigskip}
  }
%    \end{macrocode}
%
% \subsection{Front sections}\label{sec:front-sections}
% \begin{macro}{\notitle}
% \begin{macro}{\shortname}
% \begin{macro}{frontsection}
% The environment \texttt{frontsection} is for doing a section in the
% front matter (preface, list of participants, or whatever). It has 
% one optional argument:
% \begin{enumerate}
% \item[\#1] Short title (no embedded spaces) for
% use in making PS files. E.g., ``list\_of\_participants''.   
% \end{enumerate}
% and two mandatory arguments: 
% \begin{enumerate}
% \item[\#2] Title of section (e.g., `Preface')
% \item[\#3] Commands to be evaluated before processing the contents
% of the environment. This is mainly useful for putting in the things
% that would go in the preamble of a normal article  (e.g.,
% |\listofauthors|, |\indexauthor|, |\notitle|) 
% \end{enumerate}
%    \begin{macrocode}
\newcommand{\notitle}{\@titlevisiblefalse} \newcommand{\shortname}{}
\newenvironment{frontsection}[3][]{%
  \begingroup
  \renewcommand{\shortname}{#1}
  \ifx\shortname\@rmaaEmpty
     \renewcommand{\articlename}{#2}
  \else
     \renewcommand{\articlename}{\shortname}
  \fi
  \@begofarticlehook
  \@titlevisibletrue % default
  \thispagestyle{rmaafrontfirstpage}
  \listofauthors{} % blank out this in case it is not specified
  \title{{#2}}\shorttitle{#2}\shortauthor{#2}
  #3 % evaluate the command given as second mandatory arg
  \addcontentsline{toc}{chapter}{\protect\@fiddleauthors{\@TOCtitle}{\rmaaTOCauthors}}%
  \addtocontents{htc}{\protect\HyperTOCEntry{\@rawtitle}{\@rawauthors}{\articlename}}
  \if@titlevisible
     \vspace*{1cm}\centerline{\rmaa@FrontSecStyle{\@title}}\vspace*{1cm}
  \fi
  \setlength{\parskip}{3ex}
  }{\frontlabel{RMAALastPage}\clearpage\endgroup}
%    \end{macrocode}
% \end{macro}
% \end{macro}
% \end{macro}
%    \begin{macrocode}
%% Definition of the command \frontlabel, which rights out something like `F4' to
%% the .aux file, instead of `\textlatin  {i}'. This is so it is easy to post-process
%% the file for making PS files of individual sections
\newcommand{\frontlabel}[1]{%
  \@bsphack
  \protected@write\@auxout{}%
     {\string\newlabel{\articlename:#1}{{}{F\@arabic{\c@page}}}}
  \@esphack
  }
%% User command to say X, Y, Z,\\ Editors of the Proceedings
\newcommand{\EditorSignature}{\vspace*{\parskip}\rightline{\@editors}%
  \rightline{Editors of the proceedings}}
%% Fiddle with \thesection, etc. so that they don't include the chapter number
%% Basically go back to the `article' definitions of these
\renewcommand{\thesection}{\@arabic\c@section}
\renewcommand{\thefigure}{\@arabic\c@figure}
\renewcommand{\thetable}{\@arabic\c@table}

\newcommand{\@oldargs}{}
\newcommand{\@newargs}{}
%% add TOC entry for each paper
%\renewcommand{\title}[1]%
%{\renewcommand{\@title}{\rmaa@TitleStyle{#1}}%
%\addtocounter{chapter}{1}\addcontentsline{toc}{chapter}{#1}%
%}

%% Formatting of the table of contents
%% dimensions of various parts
\newlength{\TOCtextwidth}\setlength{\TOCtextwidth}{12.65cm}
\newlength{\TOCpagenumwidth}\setlength{\TOCpagenumwidth}{1.3cm}
\newlength{\TOCentrywidth}\setlength{\TOCentrywidth}{\TOCtextwidth}
\addtolength{\TOCentrywidth}{-\TOCpagenumwidth} 
\newlength{\TOCextramargin}\setlength{\TOCextramargin}{1.25cm}
\newcommand{\TOCbox}[1]{\hspace{\TOCextramargin}\parbox{\TOCtextwidth}{#1}\par}
\newlength{\TOCextramarginx}\setlength{\TOCextramarginx}{\TOCextramargin}\addtolength{\TOCextramarginx}{4.3cm}
\newlength{\TOCtextwidthx}\setlength{\TOCtextwidthx}{\TOCtextwidth}\addtolength{\TOCtextwidthx}{-4.3cm}
\newcommand{\TOCboxx}[1]{\hspace{\TOCextramarginx}\parbox{\TOCtextwidthx}{#1}\par}

\newcommand{\RMAAbannerimage}{rmsc}
\renewcommand\tableofcontents{%
  \renewcommand{\articlename}{CONTENTS}
  \begingroup
  \@begofarticlehook
  \if@twocolumn
     \@restonecoltrue\onecolumn
  \else
     \@restonecolfalse
  \fi
  %% 9pt for journal style
  \if@proceedings\else\small\fi
  \if@newstyle\small\fi %% same for new style
  \if@proceedings\thispagestyle{rmaafrontfirstpage}\else\thispagestyle{empty}\fi
  \shortauthor{Contents}\shorttitle{Contents}
  \vspace*{18mm}
  \if@printbanner
  %% put in the banner header (need to get better image of this)
     \TOCbox{\hspace*{-0.15cm}\includegraphics[width=0.9\TOCtextwidth]{\RMAAbannerimage}}
     \vspace*{8mm}
  %% volume and publication date (should really be in funny font)
     \TOCboxx{\@volumen\\[2mm] \@mes\ \@anyo}
  \else
  %% put in the banner header (need to get better image of this)
     \TOCbox{\hspace*{-0.15cm}\includegraphics[width=0.9\TOCtextwidth]{\RMAAbannerimage}}
     \vspace*{12.5mm}
  %% volume and publication date (should really be in funny font)
     \TOCboxx{~\\[2mm] ~}
  \fi
  \vspace*{16mm}
  %% This is the command that fiddles with the TOC entry so that the
  %% title is flushleft but the authors are flushright
  \renewcommand{\@fiddleauthors}[2]{%
    % first put the paper title
    \renewcommand{\\}{\relax} % temporarily turn off explicit line breaks just in case
    \renewcommand{\altaffilmark}[1]{\relax} % ditto with this
    \renewcommand{\textsuperscript}[1]{\relax} % ditto with this 
    ##1
    % now save the author list in a box
    \sbox\@tempboxa{##2}%
    % check that box is shorter than the total width available for the entry
    \ifdim \wd\@tempboxa <\TOCentrywidth
       % if so, then put it right here, flush against right margin
       \penalty10000\hskip1em plus10em\penalty0\hbox{}\nobreak\hfill\box\@tempboxa
       % Note that we force a space of at least 1em between title and authorlist, 
       % stretchable to 11em since we may have very short titles but still with an 
       % authorlist too long to go on the same line.  The only linebreak allowed is 
       % directly before the authorlist.
    \else
       % if it is too long, we set the original authorlist text in its own parbox
       \par\vskip1ex\parbox[b]{\TOCentrywidth}{\raggedleft ##2}
    \fi
    }
  \@starttoc{toc}%
  \if@restonecol\twocolumn\fi
  \frontlabel{RMAALastPage}
  \clearpage
  \thispagestyle{empty}
  \cleardoublepage % make sure we start on an odd page
  \endgroup

  %%% ALSO, do the HTML TOC
  %% The file that will be used to make the HTML TOC
  \@starttoc{htc}
  %% All the stuff in this file does nothing in the book style
  \addtocontents{htc}{\protect\HyperTOCmeeting{\@conftitle}{\@reunionplace}{\@reuniondate}}
  \addtocontents{htc}{\protect\HyperTOCeditors{\@editors}}
  \addtocontents{htc}{\protect\HyperTOCpreamble}
  \addtocontents{htc}{\protect\HyperTOCEntry{Table of Contents}{}{\articlename}}

  %% and the ADS file
  \ADS@openfile
  }
%% This is now much simpler
\renewcommand\l@chapter[2]{%TOC entry for articles
  \ifnum \c@tocdepth >\m@ne
    \vskip 1.0em \@plus\p@
    \TOCbox{%
      \parbox[b]{\TOCentrywidth}{#1}\parbox[b]{\TOCpagenumwidth}{\hfill#2}
    }
  \fi} 
\renewcommand\l@part[2]{%TOC entry for sessions: centered, no page number
  \ifnum \c@tocdepth >\m@ne
    \vskip 1.5em \@plus1.5em \goodbreak
    \TOCbox{%
      \parbox[b]{\TOCentrywidth}{\hfil#1\hfil}\hfill
    }
  \fi\nobreak} 

%% Formatting theindex
%% strongly encourage page breaks between index entries
\renewcommand{\@idxitem}{\par\penalty-1000\hangindent.45cm} 
\renewcommand{\indexspace}{\par\penalty-1000\vskip 10\p@ \@plus15\p@ \@minus3\p@\relax}
\renewenvironment{theindex}{%
  %% first, make sure we start on an odd page (06 Jun 2000)
  \clearpage\thispagestyle{empty}\cleardoublepage
  %% this looks a bit like an article but is two columns and without abstract, etc
  \renewcommand{\articlename}{INDEX}
  \title{Author Index}\shortauthor{Author Index}\shorttitle{Author Index}
  \listofauthors{}
  \addcontentsline{toc}{chapter}{\protect\@fiddleauthors{\@TOCtitle}{\rmaaTOCauthors}}%
  \addtocontents{htc}{\protect\HyperTOCEntry{\@rawtitle}{\@rawauthors}{\articlename}}
  \pagestyle{rmaa}
  \thispagestyle{rmaafirstpage}
  \@begofarticlehook
  %% 9pt for journal style
  \if@proceedings\else\small\fi
  \if@newstyle\small\fi
  \twocolumn[\@makeminititle]
  %% vertical spacing (see also definition of \@idxitem and \indexspace
  \parskip\z@ \@plus .6\p@\relax % allow quite a lot of leaway here to avoid breaks inside entries
  %% horizontal spacing, breaking and glue parameters (see also the stuff in the definition 
  %% of \indexauthor
  \parindent\z@ 
  \exhyphenpenalty10000 % don't break on explicit hyphens, it always looks bad
  \rightskip=0pt plus.2em % allow a tiny bit of leeway at the ends of lines
  \spaceskip.3333em plus.5em minus.15em % a lot of leeway between words
  \xspaceskip.5em plus.5em minus.4em % this seems to have no effect
  \let\item\@idxitem
  \renewcommand{\\}{\relax} % temporarily turn off explicit line breaks just in case
  \renewcommand{\altaffilmark}[1]{\relax} % ditto with this
  \renewcommand{\textsuperscript}[1]{\relax} % ditto with this
  }{\label{RMAALastPage}\clearpage} % clearpage necessary so that we know about \shortauthor



%% end definitions for book styles
\fi

%    \end{macrocode}
% In order to do away with plates, we sometimes want a figure to
% appear in color only in the online version. The following command
% can bve used in a figure caption to typeset a note to this effect. 
%    \begin{macrocode}
\newcommand\ColorOnlineOnly{%
  \textsc{Note: this figure is available in color in the electronic
    version of this article, obtainable from}
  \texttt{http://\discretionary{}{}{}www.\discretionary{}{}{}astroscu.\discretionary{}{}{}unam.\discretionary{}{}{}mx/\discretionary{}{}{}\~{}rmaa/}.}
%% Plates
%% Define 2 counters: plate is the number in the whole document, plateloop is used 
%% internally by \outputplates
\newcounter{plate}
\newcounter{plateloop}
%% Currently limited to 15 plates in document
%% when we do a label, we need to write to \newlabel entries to the .aux file, one 
%% that says RMAAPlateX so that it can be referred to in \ouputplates, and one that
%% says articlename:RMAAPlateX so that the dvips program can find it and asociate it
%% with the right article
\newcommand{\@platelabeli}{\oldlabel{RMAAPlate1}\label{RMAAPlate1}}
\newcommand{\@platelabelii}{\oldlabel{RMAAPlate2}\label{RMAAPlate2}}
\newcommand{\@platelabeliii}{\oldlabel{RMAAPlate3}\label{RMAAPlate3}}
\newcommand{\@platelabeliv}{\oldlabel{RMAAPlate4}\label{RMAAPlate4}}
\newcommand{\@platelabelv}{\oldlabel{RMAAPlate5}\label{RMAAPlate5}}
\newcommand{\@platelabelvi}{\oldlabel{RMAAPlate6}\label{RMAAPlate6}}
\newcommand{\@platelabelvii}{\oldlabel{RMAAPlate7}\label{RMAAPlate7}}
\newcommand{\@platelabelviii}{\oldlabel{RMAAPlate8}\label{RMAAPlate8}}
\newcommand{\@platelabelix}{\oldlabel{RMAAPlate9}\label{RMAAPlate9}}
\newcommand{\@platelabelx}{\oldlabel{RMAAPlate10}\label{RMAAPlate10}}
\newcommand{\@platelabelxi}{\oldlabel{RMAAPlate11}\label{RMAAPlate11}}
\newcommand{\@platelabelxii}{\oldlabel{RMAAPlate12}\label{RMAAPlate12}}
\newcommand{\@platelabelxiii}{\oldlabel{RMAAPlate13}\label{RMAAPlate13}}
\newcommand{\@platelabelxiv}{\oldlabel{RMAAPlate14}\label{RMAAPlate14}}
\newcommand{\@platelabelxv}{\oldlabel{RMAAPlate15}\label{RMAAPlate15}}
\newcommand{\@outputplatelabeli}{\oldlabel{RMAAOutputPlate1}}
\newcommand{\@outputplatelabelii}{\oldlabel{RMAAOutputPlate2}}
\newcommand{\@outputplatelabeliii}{\oldlabel{RMAAOutputPlate3}}
\newcommand{\@outputplatelabeliv}{\oldlabel{RMAAOutputPlate4}}
\newcommand{\@outputplatelabelv}{\oldlabel{RMAAOutputPlate5}}
\newcommand{\@outputplatelabelvi}{\oldlabel{RMAAOutputPlate6}}
\newcommand{\@outputplatelabelvii}{\oldlabel{RMAAOutputPlate7}}
\newcommand{\@outputplatelabelviii}{\oldlabel{RMAAOutputPlate8}}
\newcommand{\@outputplatelabelix}{\oldlabel{RMAAOutputPlate9}}
\newcommand{\@outputplatelabelx}{\oldlabel{RMAAOutputPlate10}}
\newcommand{\@outputplatelabelxi}{\oldlabel{RMAAOutputPlate11}}
\newcommand{\@outputplatelabelxii}{\oldlabel{RMAAOutputPlate12}}
\newcommand{\@outputplatelabelxiii}{\oldlabel{RMAAOutputPlate13}}
\newcommand{\@outputplatelabelxiv}{\oldlabel{RMAAOutputPlate14}}
\newcommand{\@outputplatelabelxv}{\oldlabel{RMAAOutputPlate15}}
\newcommand{\@platerefi}{\oldpageref{RMAAPlate1}}
\newcommand{\@platerefii}{\oldpageref{RMAAPlate2}}
\newcommand{\@platerefiii}{\oldpageref{RMAAPlate3}}
\newcommand{\@platerefiv}{\oldpageref{RMAAPlate4}}
\newcommand{\@platerefv}{\oldpageref{RMAAPlate5}}
\newcommand{\@platerefvi}{\oldpageref{RMAAPlate6}}
\newcommand{\@platerefvii}{\oldpageref{RMAAPlate7}}
\newcommand{\@platerefviii}{\oldpageref{RMAAPlate8}}
\newcommand{\@platerefix}{\oldpageref{RMAAPlate9}}
\newcommand{\@platerefx}{\oldpageref{RMAAPlate10}}
\newcommand{\@platerefxi}{\oldpageref{RMAAPlate11}}
\newcommand{\@platerefxii}{\oldpageref{RMAAPlate12}}
\newcommand{\@platerefxiii}{\oldpageref{RMAAPlate13}}
\newcommand{\@platerefxiv}{\oldpageref{RMAAPlate14}}
\newcommand{\@platerefxv}{\oldpageref{RMAAPlate15}}
\newcounter{figplatei}
\newcounter{figplateii}
\newcounter{figplateiii}
\newcounter{figplateiv}
\newcounter{figplatev}
\newcounter{figplatevi}
\newcounter{figplatevii}
\newcounter{figplateviii}
\newcounter{figplateix}
\newcounter{figplatex}
\newcounter{figplatexi}
\newcounter{figplatexii}
\newcounter{figplatexiii}
\newcounter{figplatexiv}
\newcounter{figplatexv}
%% Printed output should use the counter `plate'
\renewcommand{\theplate}{\@arabic\c@plate}
\renewcommand{\theplateloop}{\@arabic\c@plateloop}
%% Command to output all the saved plates in this paper
\newcommand{\outputplates}{%
  \setcounter{plateloop}{0}
  \if@book\else\@endofarticlehook\fi % not entirely sure why, but this
                                     % seems to be needed
  \clearpage % first output the last page of text
  % now set pagestyle for blank even pages and make sure we start on an odd page
  \pagestyle{empty}
  \cleardoublepage
  \onecolumn
  \parindent\z@
  % first put in an extra sheet
  % the main purpose of this is so there is always a blank even page before 
  % each plate for filling in the PS file if the article has an odd number
  % of pages (see splitbook.sh)
  \null\vfil\centerline{PLATES}\vfil\clearpage\pagestyle{rmaaplate}\cleardoublepage
  % output each plate in turn
  \@whilenum\value{plateloop}<\value{plate}\do{%
    \addtocounter{plateloop}{1}
    % restore the figure counter from \@figcounti, \@figcountii, etc
    \setcounter{figure}{\value{figplate\roman{plateloop}}}
    % spew out the plate (in its own group for safety)
    {\csname @plate\roman{plateloop}\endcsname\par}
    }
  \pagestyle{rmaa}
  % reset the page number as if nothing had happened
  \addtocounter{page}{-\value{plate}}
  \addtocounter{page}{-\value{plate}}
  }
\newcommand\figcaption{%
  \def\@captype{figure}
  \renewcommand{\@makecaption}{\@makefigcap}
  \caption}
\newcommand\tabcaption{%
  \def\@captype{table}
  \renewcommand{\@makecaption}{\@maketabcap}
  \caption}

\newcommand{\plate}[2][1]{%
  %% User command to make a plate page: 
  %% Mandatory argument is what goes on the plate - this should not 
  %% include any floating environments (e.g. figure). Instead, \figcaption 
  %% should be used to do the captions. 
  %% Optional argument is the number of figures (i.e. captions) on the plate
  %% page (default is 1). This is necessary for incrementing the figure counter. 
  %%
  %% Note that we have to use a \gdef rather than an \xdef when we 
  %% save the plate output in a command. This means that in the book styles when it 
  %% comes to spewing out the plate, we will have forgotten about any user macros 
  %% that may have been defined. This can be got round by repeating the macros inside the plate.
  %% (Presumably, one would use a \newcommand rather than \renewcommand if they were new.)
  % increment global and local counters
  \addtocounter{plate}{1}
  % write appropriate label to .aux file
  \csname @platelabel\roman{plate}\endcsname
  % define a global macro called \@platei, \@plateii, \@plateiii, \@plateiv, ...
  \expandafter\gdef\csname @plate\roman{plate}\endcsname{%
    \csname @outputplatelabel\roman{plateloop}\endcsname
    \null\vfil #2 \vfil
    \cleardoublepage
    }
  % save the current shortauthor and shorttitle
  % need to be \xdef to take current values of \@shortauthor, etc 
  \expandafter\xdef\csname @shorta\roman{plate}\endcsname{\@shortauthor}
  \expandafter\xdef\csname @shortt\roman{plate}\endcsname{\@shorttitle}
  % save the current value of the figure counter for later use in \outputplates
  \setcounter{figplate\roman{plate}}{\value{figure}}
  % increment the figure counter according to the optional argument
  \addtocounter{figure}{#1}
  }
%% page style for plate pages 
\newcommand{\ps@rmaaplate}{%
  % Plate X (see page Y) in the foot
  \renewcommand\@oddfoot{\hfil\csname @shorta\roman{plateloop}\endcsname\ 
    (see page~\csname @plateref\roman{plateloop}\endcsname)\hfil}%
  \renewcommand\@evenfoot{\relax}%
  % nothing in even heads (blank pages)
  \renewcommand\@evenhead{\relax}
  % short title and page num in the odd heads
  \renewcommand\@oddhead{\hfil\csname @shortt\roman{plateloop}\endcsname%
    \hfil\llap{PLATE~\theplateloop}}
  }





%% Things for the User Guide
\if@guide
%% stolen from doc.sty
\def\MakeShortVerb#1{%
  \expandafter\ifx\csname cc\string#1\endcsname\relax
  \@shortvrbinfo{Made }{#1}%
  \add@special{#1}%
  \expandafter
  \xdef\csname cc\string#1\endcsname{\the\catcode`#1}%
  \begingroup
  \catcode`\~\active  \lccode`\~`#1%
  \lowercase{%
    \global\expandafter\let
    \csname ac\string#1\endcsname~%
    \gdef~{\verb~}}%
  \endgroup
  \global\catcode`#1\active
  \else
  \@shortvrbinfo\@empty{#1 already}%
  \fi}
\def\DeleteShortVerb#1{%
  \expandafter\ifx\csname cc\string#1\endcsname\relax
  \else
  \@shortvrbinfo{Deleted }{#1 as}%
  \rem@special{#1}%
  \global\catcode`#1\csname cc\string#1\endcsname
  \global \expandafter\let \csname cc\string#1\endcsname \relax
  \ifnum\catcode`#1=\active
  \begingroup
  \catcode`\~\active   \lccode`\~`#1%
  \lowercase{%
    \global\expandafter\let\expandafter~%
    \csname ac\string#1\endcsname}%
  \endgroup \fi \fi}
\def\@shortvrbinfo#1#2{%
  \PackageInfo{doc}{%
    #1\expandafter\@gobble\string#2 a short reference
    for \string\verb}}
\def\add@special#1{%
  \rem@special{#1}%
  \expandafter\gdef\expandafter\dospecials\expandafter
  {\dospecials \do #1}%
  \expandafter\gdef\expandafter\@sanitize\expandafter
  {\@sanitize \@makeother #1}}
\def\rem@special#1{%
  \def\do##1{%
    \ifnum`#1=`##1 \else \noexpand\do\noexpand##1\fi}%
  \xdef\dospecials{\dospecials}%
  \begingroup
  \def\@makeother##1{%
    \ifnum`#1=`##1 \else \noexpand\@makeother\noexpand##1\fi}%
  \xdef\@sanitize{\@sanitize}%
  \endgroup}
%% define | as shortcut for \verbatim
\MakeShortVerb{\|}
%% comand to specify arguments to LaTeX commands
\newcommand{\rmArg}[1]{%
  \ensuremath{\langle}{\scriptsize \textsf{\MakeUpperCase{#1}}}\ensuremath{\rangle}}
%\newcommand{\rmArg}[1]{\ensuremath{\langle}\textsc{#1}\ensuremath{\rangle}}
\providecommand{\BibTeX}{{\rmfamily B\kern-.05em%
    \textsc{i\kern-.025em b}\kern-.08em%
    T\kern-.1667em\lower.7ex\hbox{E}\kern-.125emX}}
%% Environment for code examples
\newenvironment{Code}{%
  \begin{list}{}{\setlength{\leftmargin}{0.05\columnwidth}} 
  \item[]}{\end{list}}
\fi
%</package>
%    \end{macrocode}

% \section{User commands for astronomical articles}
% \label{sec:user-commands}
% Some random definitions. These go both into the main class file and
% also into |adshtml.sty|.
% 
% Definitions of the journal name in different incarnations.  
%    \begin{macrocode}
%<*package|adshtml-style>
\newcommand{\RMAA}{RMxAA}
\newcommand{\RMAAlong}{Revista Mexicana de Astronom\'{\i}a y Astrof\'{\i}sica}
\newcommand{\RMAASC}{RMxAC}
\newcommand{\RMAASClong}{\RMAA\ (Serie de Conferencias)}
%    \end{macrocode}
% \begin{macro}{\ibidrule}
% This puts a horizontal rule with a period after it. It is used when
% the author list of a reference is identical to the preceding one. 
%    \begin{macrocode}
\newcommand{\ibidrule}{\rule[-0.03cm]{4.7em}{.01cm}.\ }
%    \end{macrocode}
% \end{macro}
% \begin{macro}{\ThisVolume}
% References to other articles in the same volume can be coded as,
% e.g., |\bibitem{} Bloggs, J. \ThisVolume{bloggs}|.
%    \begin{macrocode}
\newcommand{\ThisVolume}[1]{%
  \@year, \if@proceedings\RMAASC\else\RMAA\fi, \@volume,
  \if@book\oldpageref{#1:RMAAFirstPage}\else 000\fi\ (this volume)}
%    \end{macrocode}
% \end{macro}
% Other stuff that may be useful in ref lists.
%    \begin{macrocode}
\newcommand{\etal}{et al.\@}
%    \end{macrocode}
% Various new commands from AAS macros
%    \begin{macrocode}
\def\acknowledgments{\vskip 3ex plus .8ex minus .4ex}
\let\acknowledgements=\acknowledgments                  % second spelling
\let\jnl@style=\rm
\def\ref@jnl#1{{\jnl@style#1}}
\DeclareRobustCommand\aj{\ref@jnl{AJ}}                   
\DeclareRobustCommand\araa{\ref@jnl{ARA\&A}}             
\DeclareRobustCommand\apj{\ref@jnl{ApJ}}                 
\DeclareRobustCommand\apjl{\ref@jnl{ApJ}}                
\DeclareRobustCommand\apjs{\ref@jnl{ApJS}}               
\DeclareRobustCommand\ao{\ref@jnl{Appl.~Opt.}}           
\DeclareRobustCommand\apss{\ref@jnl{Ap\&SS}}             
\DeclareRobustCommand\aap{\ref@jnl{A\&A}}                
\DeclareRobustCommand\aapr{\ref@jnl{A\&A~Rev.}}          
\DeclareRobustCommand\aaps{\ref@jnl{A\&AS}}              
\DeclareRobustCommand\azh{\ref@jnl{AZh}}                 
\DeclareRobustCommand\baas{\ref@jnl{BAAS}}               
\DeclareRobustCommand\jrasc{\ref@jnl{JRASC}}             
\DeclareRobustCommand\memras{\ref@jnl{MmRAS}}            
\DeclareRobustCommand\mnras{\ref@jnl{MNRAS}}             
\DeclareRobustCommand\pra{\ref@jnl{Phys.~Rev.~A}}        
\DeclareRobustCommand\prb{\ref@jnl{Phys.~Rev.~B}}        
\DeclareRobustCommand\prc{\ref@jnl{Phys.~Rev.~C}}        
\DeclareRobustCommand\prd{\ref@jnl{Phys.~Rev.~D}}        
\DeclareRobustCommand\pre{\ref@jnl{Phys.~Rev.~E}}        
\DeclareRobustCommand\prl{\ref@jnl{Phys.~Rev.~Lett.}}    
\DeclareRobustCommand\pasp{\ref@jnl{PASP}}               
\DeclareRobustCommand\pasj{\ref@jnl{PASJ}}               
\DeclareRobustCommand\qjras{\ref@jnl{QJRAS}}             
\DeclareRobustCommand\skytel{\ref@jnl{S\&T}}             
\DeclareRobustCommand\solphys{\ref@jnl{Sol.~Phys.}}      
\DeclareRobustCommand\sovast{\ref@jnl{Soviet~Ast.}}      
\DeclareRobustCommand\ssr{\ref@jnl{Space~Sci.~Rev.}}     
\DeclareRobustCommand\zap{\ref@jnl{ZAp}}                 
\DeclareRobustCommand\nat{\ref@jnl{Nature}}              
\DeclareRobustCommand\iaucirc{\ref@jnl{IAU~Circ.}}
\DeclareRobustCommand\aplett{\ref@jnl{Astrophys.~Lett.}}
\DeclareRobustCommand\apspr{\ref@jnl{Astrophys.~Space~Phys.~Res.}}
\DeclareRobustCommand\bain{\ref@jnl{Bull.~Astron.~Inst.~Netherlands}}
\DeclareRobustCommand\fcp{\ref@jnl{Fund.~Cosmic~Phys.}}
\DeclareRobustCommand\gca{\ref@jnl{Geochim.~Cosmochim.~Acta}}
\DeclareRobustCommand\grl{\ref@jnl{Geophys.~Res.~Lett.}}
\DeclareRobustCommand\jcp{\ref@jnl{J.~Chem.~Phys.}}      
\DeclareRobustCommand\jgr{\ref@jnl{J.~Geophys.~Res.}}    
\DeclareRobustCommand\jqsrt{\ref@jnl{J.~Quant.~Spec.~Radiat.~Transf.}}
\DeclareRobustCommand\memsai{\ref@jnl{Mem.~Soc.~Astron.~Italiana}}
\DeclareRobustCommand\nphysa{\ref@jnl{Nucl.~Phys.~A}}
\DeclareRobustCommand\physrep{\ref@jnl{Phys.~Rep.}}
\DeclareRobustCommand\physscr{\ref@jnl{Phys.~Scr}}
\DeclareRobustCommand\planss{\ref@jnl{Planet.~Space~Sci.}}
\DeclareRobustCommand\procspie{\ref@jnl{Proc.~SPIE}}
\let\astap=\aap
\let\apjlett=\apjl
\let\apjsupp=\apjs
\let\applopt=\ao
\def\phn{\phantom{0}}
\def\phd{\phantom{.}}
\def\phs{\phantom{\ensuremath{-}}}
\def\phm#1{\phantom{#1}}
%    \end{macrocode}
% All these are redone as robust and using more modern commands where
% possible. The robustness is so they will not get expanded if used in
% abstracts and titles when these are passed to the ADS file. 
% \changes{3.16}{2003/04/07}{
%  Added \cs{Color} command for doing photometric colors.}
%    \begin{macrocode}
%%\def\sun{\hbox{$\odot$}}
\DeclareRobustCommand\sun{\ensuremath{\odot}}
%%\def\earth{\hbox{$\oplus$}}
\DeclareRobustCommand\earth{\ensuremath{\oplus}}
%%\def\lesssim{\mathrel{\hbox{\rlap{\hbox{\lower4pt\hbox{$\sim$}}}\hbox{$<$}}}}
%%\def\gtrsim{\mathrel{\hbox{\rlap{\hbox{\lower4pt\hbox{$\sim$}}}\hbox{$>$}}}}
\DeclareRobustCommand\lesssim{\mathrel{\hbox{\smaller\smaller\rlap{\hbox{\lower3pt\hbox{$\sim$}}}\raise1.5pt\hbox{$<$}}}}
\DeclareRobustCommand\gtrsim{\mathrel{\hbox{\smaller\smaller\rlap{\hbox{\lower3pt\hbox{$\sim$}}}\raise1.5pt\hbox{$>$}}}}
%%\def\lta{\mathrel{\spose{\lower 3pt\hbox{$\mathchar"218$}}
%%     \raise 2.0pt\hbox{$\mathchar"13C$}}}
%%\def\gta{\mathrel{\spose{\lower 3pt\hbox{$\mathchar"218$}}
%%     \raise 2.0pt\hbox{$\mathchar"13E$}}}
\DeclareRobustCommand\sq{\hbox{\rlap{$\sqcap$}$\sqcup$}}
%%\def\arcdeg{\hbox{$^\circ$}}
%%\def\arcmin{\hbox{$^\prime$}}
%%\def\arcsec{\hbox{$^{\prime\prime}$}}
\DeclareRobustCommand\arcdeg{\ensuremath{^\circ}}
\DeclareRobustCommand\arcmin{\ensuremath{^\prime}}
\DeclareRobustCommand\arcsec{\ensuremath{^{\prime\prime}}}
%%\def\fd{\hbox{$.\!\!^{\rm d}$}}
%%\def\fh{\hbox{$.\!\!^{\rm h}$}}
%%\def\fm{\hbox{$.\!\!^{\rm m}$}}
%%\def\fs{\hbox{$.\!\!^{\rm s}$}}
%%\def\fdg{\hbox{$.\!\!^\circ$}}
%%\def\farcm{\hbox{$.\mkern-4mu^\prime$}}
%%\def\farcs{\hbox{$.\!\!^{\prime\prime}$}}
%%\def\fp{\hbox{$.\!\!^{\scriptscriptstyle\rm p}$}}
\DeclareRobustCommand\fd{\ensuremath{.\!\!^\mathrm{d}}}
\DeclareRobustCommand\fh{\ensuremath{.\!\!^\mathrm{h}}}
\DeclareRobustCommand\fm{\ensuremath{.\!\!^\mathrm{m}}}
\DeclareRobustCommand\fs{\ensuremath{\mkern0.5mu.\mkern-5mu^\mathrm{s}\mkern-1mu}}
\DeclareRobustCommand\fdg{\ensuremath{.\!\!^\circ}}
\DeclareRobustCommand\farcm{\ensuremath{.\mkern-4mu^\prime}}
\DeclareRobustCommand\farcs{\ensuremath{\mkern1mu.\mkern-6mu^{\prime\prime}\mkern-1.5mu}}
%%\DeclareRobustCommand\farcs{\ensuremath{.\!\!^{\prime\prime}}}
\DeclareRobustCommand\fp{\ensuremath{.\!\!^{\scriptscriptstyle\rm p}}}
%%\def\micron{\hbox{$\mu$m}}
\DeclareRobustCommand\micron{\ensuremath{\mu\mathrm{m}}}
\let\la=\lesssim                        
\let\ga=\gtrsim
%%\def\case#1#2{\hbox{$\frac{#1}{#2}$}}
%%\def\slantfrac#1#2{\hbox{$\,^#1\!/_#2$}}
%%\def\onehalf{\slantfrac{1}{2}}
%%\def\onethird{\slantfrac{1}{3}}
%%\def\twothirds{\slantfrac{2}{3}}
%%\def\onequarter{\slantfrac{1}{4}}
%%\def\threequarters{\slantfrac{3}{4}}
\DeclareRobustCommand\case[2]{\ensuremath{\frac{#1}{#2}}}
\DeclareRobustCommand\slantfrac[2]{\ensuremath{\,^#1\!/_#2}}
\DeclareRobustCommand\onehalf{\slantfrac{1}{2}}
\DeclareRobustCommand\onethird{\slantfrac{1}{3}}
\DeclareRobustCommand\twothirds{\slantfrac{2}{3}}
\DeclareRobustCommand\onequarter{\slantfrac{1}{4}}
\DeclareRobustCommand\threequarters{\slantfrac{3}{4}}
%%\def\ubvr{\hbox{$U\!BV\!R$}}            
%%\def\ub{\hbox{$U\!-\!B$}}               
%%\def\bv{\hbox{$B\!-\!V$}}               
%%\def\vr{\hbox{$V\!-\!R$}}               
%%\def\ur{\hbox{$U\!-\!R$}}
%    \end{macrocode}
% New commands for ``photometric'' colors 
%    \begin{macrocode}
\DeclareRobustCommand\ubvr{\ensuremath{U\!BV\!R}}            
\DeclareRobustCommand\Color[2]{%
  \ensuremath{%
    #1\!\ifx#1V\!\fi-\!#2
  }}
\DeclareRobustCommand\ub{\Color{U}{B}}
\DeclareRobustCommand\bv{\Color{B}{V}}
\DeclareRobustCommand\vr{\Color{V}{R}}
\DeclareRobustCommand\ur{\Color{U}{R}}
%% New version of the \ion command, that actually works! 
%    \end{macrocode}
% \begin{macro}{\ion}
% This version of the |\ion| command allows you to write either
% |\ion{N}{4}| or |\ion{N}{iv}| in order to get N~IV. 
%    \begin{macrocode}
\DeclareRobustCommand\ion[2]{%
  \mbox{#1\kern0.2em%
    \smaller\rmfamily%
    \edef\@tempa{\@car#2\@nil}%
    \ifcat1\@tempa%
    \@Roman{#2}%
    \else%
    \uppercase{#2}%
    \fi}}
\newcounter{IonStage}
\DeclareRobustCommand\plainion[2]{\setcounter{IonStage}{#2}#1 \Roman{IonStage}}
%%\def\ion#1#2{\ensuremath{\mathrm{#1}\;}{\protect\small\rm\@Roman{#2}}}
%    \end{macrocode}
% \end{macro}
%    \begin{macrocode}
\DeclareRobustCommand\nodata{$\cdots$}
%    \end{macrocode}
% These are needed by the |\ion| command. Isn't there an easier way of
% doing this?
%    \begin{macrocode}
\let\@small=\small
\let\@normalsize=\normalsize
\let\@footnotesize=\footnotesize
\let\@scriptsize=\scriptsize
\let\@tiny=\tiny
\let\@large=\large
\let\@Large=\Large
\let\@LARGE=\LARGE
\let\@huge=\huge
\let\@Huge=\Huge
\newcommand{\smaller}{\small}
\renewcommand{\normalsize}{\renewcommand{\smaller}{\small}\@normalsize}
\renewcommand{\small}{\renewcommand{\smaller}{\footnotesize}\@small}
\renewcommand{\footnotesize}{\renewcommand{\smaller}{\scriptsize}\@footnotesize}
\renewcommand{\scriptsize}{\renewcommand{\smaller}{\tiny}\@scriptsize}
\renewcommand{\tiny}{\renewcommand{\smaller}{\tiny}\@tiny}
\renewcommand{\large}{\renewcommand{\smaller}{\normalsize}\@large}  
\renewcommand{\Large}{\renewcommand{\smaller}{\large}\@Large}  
\renewcommand{\LARGE}{\renewcommand{\smaller}{\Large}\@LARGE}  
\renewcommand{\huge}{\renewcommand{\smaller}{\LARGE}\@huge}  
\renewcommand{\Huge}{\renewcommand{\smaller}{\huge}\@Huge}  
%</package|adshtml-style>
%    \end{macrocode}


% \section{Other files}

% \subsection{The driver and make files}
% \label{sec:driver_make}
% These are put here for completeness, although they still need to be
% distributed separately since otherwise the user won't be able to
% unpack this file into its constituent parts.
% 
% The |Makefile| has two principal rules: |docs|, and |macros|. The
% |docs| rule produces a PDF and gzipped PS file of the documentation.
% I still haven't perfected the combination of rules required to get
% it to run \LaTeX{} enough times in all circumstances. The |macros|
% rule runs \LaTeX{} on |rmaa.ins| in order to unpack |rmaa.cls| and
% all the other little files (including the |Makefile| and the
% |rmaa.ins| files themselves\dots 
% 
% First set up lists of various types of file. 
%    \begin{macrocode}
%<*makefile>
DOCFILES=rmaa.pdf
MACRO_LTX_FILES=rmaa.cls rmaacite.sty adshtml.sty
SHELLSCRIPTS=makeads.sh
MACRO_OTHER_FILES=adshtml.perl sample.latex2html-init
MAKEFILE=Makefile
MACROFILES=$(MACRO_LTX_FILES) $(SHELLSCRIPTS) $(MACRO_OTHER_FILES)         
SKELETONFILES=ADS.tex 
%    \end{macrocode}
% Install the files read only into the package directory, specified via a path relative to the source directory. \changes{4.0}{2019/11/14}{Reorganize folder layout so that hardly anything is in the top level.  That way, things are clearer on github.}
%    \begin{macrocode}
INSTALL=install 
INSTALL_DIR=../packages
%    \end{macrocode}
% The default rule is to make everything. Installation must be
% explicitly asked for.
%    \begin{macrocode}
all: docs macros

install: install-docs install-macros

%    \end{macrocode}
% Most files are installed read-only to prevent editing. Skeleton
% files must be writeable. Must be careful not to overwrite the
% Makefile in the parent directory. 
%    \begin{macrocode}
install-macros: macros
        $(INSTALL) -m 444 $(MACRO_LTX_FILES) $(MACRO_OTHER_FILES) $(INSTALL_DIR)
        $(INSTALL) -m 555 $(SHELLSCRIPTS) $(INSTALL_DIR)
        $(INSTALL) -m 644 $(SKELETONFILES) $(INSTALL_DIR)

install-docs: docs
        $(INSTALL) rmaa.pdf $(INSTALL_DIR)/rm-fulldocs.pdf

%    \end{macrocode}
% Make the PDF and gzipped PS files. 
%    \begin{macrocode}
docs: $(DOCFILES)

%    \end{macrocode}
% Make everything else if either the |.dtx| or |.ins| files have
% changed. 
%    \begin{macrocode}
macros: rmaa.dtx rmaa.ins
%    \end{macrocode}
% Ensure that all the macro files can be written to then recreate
% them. The |touch| is there in case any of the files don't exist
% yet. 
%    \begin{macrocode}
        touch $(MACROFILES) $(MAKEFILE)
        chmod u+w $(MACROFILES) $(MAKEFILE) rmaa.ins
        latex "\input{rmaa.ins}"
%    \end{macrocode}
% \TeX{} automatically adds the |.tex| extension to extensionless
% files. We fix this and take the opportunity to convert
% back to TABs all that it has lovingly converted to spaces. 
%    \begin{macrocode}
        cp -f $(MAKEFILE) $(MAKEFILE).BCK && unexpand $(MAKEFILE).tex >| $(MAKEFILE) && rm $(MAKEFILE).tex
%    \end{macrocode}
% Make the shell scripts executable and everything unwriteable (so
% they don't get hand-edited by mistake). However, undo this for the
% files that the user is \emph{supposed} to modify. 
%    \begin{macrocode}
        chmod a+x $(SHELLSCRIPTS)
        chmod a-w $(MACROFILES) $(MAKEFILE)
        chmod u+w $(SKELETONFILES)

%    \end{macrocode}
% We now use PDF\LaTeX{} for this.
%    \begin{macrocode}
rmaa.idx: rmaa.dtx
        @echo "REGENERATING .idx FILE"
        pdflatex "\nonstopmode\input{rmaa.dtx}"

rmaa.aux: rmaa.dtx rmaa.ind rmaa.gls
        @echo "REGENERATING .aux FILE"
        pdflatex "\nonstopmode\input{rmaa.dtx}"

rmaa.ind: rmaa.idx
        @echo "REGENERATING INDEX"
        makeindex -s gind.ist rmaa

rmaa.gls: rmaa.glo
        @echo "REGENERATING CHANGELOG"
        makeindex -s gglo.ist -o rmaa.gls rmaa.glo

rmaa.pdf: rmaa.aux rmaa.ind rmaa.gls
        @echo "REGENERATING .pdf FILE"
        pdflatex "\nonstopmode\input{rmaa.dtx}"

clean: 
        rm -f rmaa.aux rmaa.dvi rmaa.ind rmaa.gls
%</makefile>
%    \end{macrocode}

% The |rmaa.ins| file uses the \textsf{docstrip} package to split out
% all the files and strip off this documentation.
%    \begin{macrocode}
%<*ins>
%%
%% Installation batch file for the RMAA LaTeX package
%%
\def\batchfile{rmaa.ins}
\input docstrip.tex
\preamble

PLEASE DO NOT EDIT THIS FILE! It is generated automatically from
the file rmaa.dtx using the accompanying Makefile and rmaa.ins 
Any hand-edited changes will be lost

A LaTeX document class for the RMxAA

Author: William Henney (w.henney@irya.unam.mx)

Latest version available from: http://www.irya.unam.mx/rmaa

Copyright (C) 1999-2019 Instituto de Astronomia, UNAM, Mexico

Permission is hereby granted, free of charge, to any person obtaining
a copy of this software and associated documentation files (the
"Software"), to deal in the Software without restriction, including
without limitation the rights to use, copy, modify, merge, publish,
distribute, sublicense, and/or sell copies of the Software, and to
permit persons to whom the Software is furnished to do so, subject to
the following conditions:

The above copyright notice and this permission notice shall be
included in all copies or substantial portions of the Software.

THE SOFTWARE IS PROVIDED "AS IS", WITHOUT WARRANTY OF ANY KIND,
EXPRESS OR IMPLIED, INCLUDING BUT NOT LIMITED TO THE WARRANTIES OF
MERCHANTABILITY, FITNESS FOR A PARTICULAR PURPOSE AND
NONINFRINGEMENT. IN NO EVENT SHALL THE AUTHORS OR COPYRIGHT HOLDERS BE
LIABLE FOR ANY CLAIM, DAMAGES OR OTHER LIABILITY, WHETHER IN AN ACTION
OF CONTRACT, TORT OR OTHERWISE, ARISING FROM, OUT OF OR IN CONNECTION
WITH THE SOFTWARE OR THE USE OR OTHER DEALINGS IN THE SOFTWARE.
\endpreamble
%    \end{macrocode}
% Clobber files without asking
%    \begin{macrocode}
\askforoverwritefalse
%    \end{macrocode}
% Make the \LaTeX{} macro files (class and styles)
%    \begin{macrocode}
\generate{% 
  \file{rmaa.cls}{\from{rmaa.dtx}{package}}
  \file{rmaa.ins}{\from{rmaa.dtx}{ins}}
  \file{rmaacite.sty}{\from{rmaa.dtx}{rmaacite}}
  \file{adshtml.sty}{\from{rmaa.dtx}{adshtml-style}}
  \file{ADS.tex}{\from{rmaa.dtx}{ads-container}}
  }
%    \end{macrocode}
% Make the other non-\LaTeX{} files. These are done separately since
% they can't use `\%\%' as a comment character for the
% preamble/postamble (and Perl files can't have the postamble I don't
% think). 
%    \begin{macrocode}
\generate{\nopreamble\nopostamble% 
  \file{Makefile}{\from{rmaa.dtx}{makefile}}
  \file{adshtml.perl}{\from{rmaa.dtx}{adshtml-perl}}
  \file{makeads.sh}{\from{rmaa.dtx}{ads-shell-script}}
  \file{sample.latex2html-init}{\from{rmaa.dtx}{latex2html-init}}
  }

%% Local Variables: 
%% mode: LaTeX
%% mode: Font-lock
%% End:
%</ins>
%    \end{macrocode}

% \subsection{Extended citation syntax}
% Modified from a file called mnrascite.sty
% I no longer remember where I stole the original version of this from
% 
% Alterations to the cite commands to make them do more sensible things.
% N.B. the bibliography style should produce entries which have labels where
% the date is delimited by angle brackets '<>' for these macros to work
% properly. They will not give very sensible error messages if you do not do 
% this properly.
% 
% A new command |\scite{key}| is defined which produces a citation which is 
% suitable for insertion in a sentence. i.e., it does not put brackets around 
% the whole citation, only the date. Also |\pcite{key}| produce output with no
% parentheses whatsoever. 
% 
% Should redefine |\@bibitem| really to have any hope of compatibility
% the command |\lbibcite| is defined to deal with the citations.
%    \begin{macrocode}
%<*rmaacite>
\def\bibitem[#1<#2>]#3{\par\hangindent=1.5em\hangafter=1
  \if@filesw
  {\def\protect##1{\string ##1\space}
    \immediate\write\@auxout{\string\lbibcite{#3}{#1}{#2}}}
  \fi}
\@redefinebibitem % what it says
%    \end{macrocode}
% Note that the |\lbibcite| defines two types of citation one contains the names,
% the other the year
%    \begin{macrocode}
\def\lbibcite#1#2#3{\global\@namedef{b@#1}{#2}
                    \global\@namedef{by@#1}{#3}}

\def\@bcite#1#2{({#1\if@tempswa , #2\fi})}
\def\@pcite#1#2{#1\if@tempswa , #2\fi}
%    \end{macrocode}
% This declaration contains what is put between the name and the date
% \changes{3.11}{2003/01/31}{Moved the \cs{articlename} to after the key for consistency
% with the new natbib stuff.}
%    \begin{macrocode}
\def\@brcite{\ (}
\def\@ebrcite{)}
\def\@fmttop#1{\@fmtcite[\csname #1\endcsname]}
\def\@citex[#1]#2{%
%%  \if@filesw\immediate\write\@auxout{\string\citation{#2}}\fi
%% WJH added 16 Jan 2002
  \@for\@citeb:=#2\do{%
    \edef\@citeb{\expandafter\@firstofone\@citeb:\articlename}%
    \if@filesw\immediate\write\@auxout{\string\citation{\@citeb}}\fi}%
%% to here
  \def\@citea{}%
  \@cite{\@for\@citeb:=#2\do
    {%
      \edef\@citeb{\expandafter\@firstofone\@citeb:\articlename}%
      \@citea\def\@citea{;\ }\@ifundefined
      {b@\@citeb}{{\bf ?}\@warning {Citation `\@citeb' on page \thepage \space undefined}}%
      {\csname b@\@citeb\endcsname\@citebro\csname by@\@citeb\endcsname\@citebrc}}}%
  {#1}}

\def\cite{\@ifnextchar [{\let\@citebro=\space
                          \let\@citebrc=\relax
                          \let\@cite=\@bcite\@tempswatrue\@citex}
                        {\let\@citebro=\space
                          \let\@citebrc=\relax
                          \let\@cite=\@bcite\@tempswafalse\@citex[]}}
\def\pcite{\@ifnextchar [{\let\@citebro=\space
                          \let\@citebrc=\relax
                          \let\@cite=\@pcite\@tempswatrue\@citex}
                        {\let\@citebro=\space
                          \let\@citebrc=\relax
                          \let\@cite=\@pcite\@tempswafalse\@citex[]}}
\def\scite{\@ifnextchar [{\let\@citebro=\@brcite
                          \let\@citebrc=\@ebrcite
                          \let\@cite=\@pcite\@tempswatrue\@citex}
                        {\let\@citebro=\@brcite
                          \let\@citebrc=\@ebrcite
                          \let\@cite=\@pcite\@tempswafalse\@citex[]}}

\renewenvironment{thebibliography}
  {\section*{REFERENCES}\vspace*{\afterbiblioskip}\small
    \if@book\RereadBiblio\fi
    \sloppy\clubpenalty4000\widowpenalty4000
    \sfcode`\.=1000\relax
    \setlength\parindent{\z@}\setlength\parskip{\z@}}
  {\relax\normalsize\@outputaddressesifnecessary} 
%</rmaacite>
%    \end{macrocode}

% \subsection{Auxilliary files for generating ADS data}

% The ADS data are written by \LaTeX{} to a |.ads| file. This file
% then has to be included in the following file (|ADS.tex|),
% which is run through \LaTeX{}2HTML (see below). 

%    \begin{macrocode}
%<*ads-container>
%%
%% Instructions:
%% Edit the \input{} line below to reflect the name of your file. 
\documentclass{article}
\usepackage[latin1]{inputenc}
\usepackage{adshtml}
\newcommand{\ADSfield}[2]
{\%#1  #2\\}
%%\newcommand{\ADSvfield}[2]
%%{\%#1 \begin{verbatim}#2\end{verbatim}}
\newenvironment{ADSarticle}{}{\par\bigskip}
\renewcommand{\bibitem}[2][]{\par}
\setlength{\parindent}{0pt}
\renewcommand{\ThisVolume}[1]{%  
  2001, \RMAASC, 12, \oldpageref{#1:RMAAFirstPage}}  
 

\begin{document}
\makeatletter
\input{rmbooksamp.ads}
\end{document}
%</ads-container>
%    \end{macrocode}

% This is the |adshtml.sty| style file used by |ADS.tex|. It just
% defines a few 
% macros that might be used in abstracts, etc (really it should define
% them all). This now gets all its code from
% Sec~\ref{sec:user-commands}. 
%    \begin{macrocode}
%<*adshtml-style>
%% Nothing else needs to go here now. 
%</adshtml-style>
%    \end{macrocode}
% The init file for latex2html.
%    \begin{macrocode}
%<*latex2html-init>
$PDFTOCAIRO = "pdftocairo";
$PDFLATEX = "pdflatex";
$LATEX = "latex";
$MAX_SPLIT_DEPTH = 0;
$HTML_VERSION = "4.0,math";
$NO_SIMPLE_MATH = 1;
$ADDRESS = "<p><address><a href=\"mailto:rmaa\@astro.unam.mx\">rmaa\@astro.unam.mx</address></p>"; 
$INFO = "<p>This document was created using latex2html and the rmaa latex macros</p>"; 
$NO_NAVIGATION = 1;
1; 
%</latex2html-init>
%    \end{macrocode}
% This is the |adshtml.perl| file that is used by \LaTeX{}2HTML. Some
% bizarre \TeX{}y command sometimes show up so we tell it to ignore
% these. All the astronomy macros that come from AAS are processed in
% \TeX{}. That way they get made to images but with an |ALT| tag that
% has the original code, which is what will be picked up by |lynx|,
% with the result that they get passed intact into the final |.txt|
% file. 
%    \begin{macrocode}
%<*adshtml-perl>
package main;

$dum = &do_cmd_makeatletter();

sub do_cmd_RMAA {
    local($_) = @_;
    "RMxAA" . $_;
}

sub do_cmd_RMAASC {
    local($_) = @_;
    "RMxAASC" . $_;
}

sub do_cmd_etal {
    local($_) = @_;
    "et al\." . $_;
}

sub do_cmd_ibidrule {
    local($_) = @_;
    "<img src=\"foo\" alt=\"\ibidrule\">" . $_;
}


sub do_cmd_bibitem {
    local($_) = @_;
    local($key);
    local($opt, $dummy)=&get_next_optional_argument; 
    $key = &missing_braces unless
        ((s/$next_pair_pr_rx/$key = $2; ''/eo)
        ||(s/$next_pair_rx/$key = $2; ''/eo));
    "<p>" . $_; 
}

&ignore_commands (<<_RAW_ARG_CMDS_);
lowercase
adjustfinalcols
textsuperscript # {}
newline # []
thinmuskip
mskip
smallskip
medskip
bigskip
hspace # {}
hspacestar # {}
vspace # {}
vspacestar # {}
mathrm
_at_mathrm
spacefactor
discretionary
unhbox 
voidb_at_x 
penalty 
_at_M
_at_m
null
IeC
_RAW_ARG_CMDS_

&process_commands_in_tex (<<_RAW_ARG_CMDS_);
sun
earth
lesssim
gtrsim
sq
argdeg
arcmin
arcsec
fd
fh
fm
fs
fdg
farcm
farcs
fp
micron
la
ga
case # {} # {}
slantfrac # {} # {}
onehalf
onethird
twothirds
onequarter
threequarters
ubvr
ub
bv
vr
ur
aj
araa
apj
apjl
apjs
ao
apss
aap
aapr
aaps
azh
baas
jrasc
memras
mnras
pra
prb
prc
prd
pre
prl
pasp
pasj
qjras
skytel
solphys
sovast
ssr
zap
nat
iaucirc
aplett
apspr
bain
fcp
gca
grl
jcp
jgr
jqsrt
memsai
nphysa
physrep
physscr
planss
procspie
_RAW_ARG_CMDS_

1;                              # This must be the last line
%</adshtml-perl> 
%    \end{macrocode}

% This is a shell script |makeads.sh| that generates the |_ads.txt| file from
% |_ads.tex| by using \LaTeX{}2HTML and |lynx|.  \changes{4.0}{2019/11/13}{Make sure that lynx writes utf-8 output instead of latin-1.}
% \changes{2.2beta13}{2002/03/12}{Complete rewrite of |makeads.sh|. Now much 
% improved I think}
% \changes{4.0}{2019/11/13}{Fix the at-sign substitution so that it only eliminates \cs{@} and does not mess with email addresses.}

%    \begin{macrocode}
%<*ads-shell-script>
#! /bin/sh
if [ -z $1 ]; then
    echo "Usage: $0 basename"
    exit
fi
BASENAME=$1
ADSBASENAME=${BASENAME}_ads
AUXFILE=$BASENAME.aux
TMPFILE=tmp.$$


echo "Generating $ADSBASENAME.tex"

# For some reason, must double-up \\ to \\\\ but not \ to \\ - odd!
cat > $ADSBASENAME.tex <<EOF
%% PLEASE DO NOT EDIT THIS FILE!  It is generated automatically from
%% the script makeads.sh.  Any hand-edited changes will be lost
%% 
\documentclass{article}
\usepackage[latin1]{inputenc}
\usepackage{html}
\usepackage{adshtml}
\newcommand{\ADSfield}[2]{\%#1  #2\\\\}
\newenvironment{ADSarticle}{}{\par\bigskip}
\setlength{\parindent}{0pt}
\newcounter{istage}
\renewcommand{\ion}[2]{\setcounter{istage}{#2}#1 \Roman{istage}}
\begin{document}
\makeatletter
EOF
echo >> $ADSBASENAME.tex "\input{$BASENAME.ads2}" 
echo >> $ADSBASENAME.tex "\end{document}"



echo "Generating $ADSBASENAME.sed" 

# Find type of journal
grep '\\documentclass\[.*proceedings.*\]' $BASENAME.tex >/dev/null \
     && JOURNAL='RMxAC' \
     || JOURNAL='RMxAA'

# Find the year and volume number
YEAR=`grep '\\SetYear' $BASENAME.tex | sed 's/.*\\SetYear{\(.*\)}.*/\1/'`
VOLUME=`grep '\\SetVolume' $BASENAME.tex | sed 's/.*\\SetVolume{\(.*\)}.*/\1/'`

# Generate an auxilliary sed script 
cat > $TMPFILE <<EOF
s|\\\\newlabel{\(.*\):RMAAFirstPage}{{.*}{\(.*\)}}|s/\\\\\\\\ThisVolume{\1}/$YEAR, RMxAC, $VOLUME, \2/g|g
EOF

# Now go through the .aux file, looking for \newlabels and use the sed 
# script just generated in order to make another sed script, one that 
# will replace \ThisVolume{LABEL} with the correct reference
grep '^\\newlabel{.*:RMAAFirstPage}' $AUXFILE \
     | sed -f $TMPFILE > $ADSBASENAME.sed && rm $TMPFILE

# Apend a further sed command to get rid of \@, since redefining 
# it in _ads.tex didn't work
echo >> $ADSBASENAME.sed 's/\\\\@//g'

echo "Generating $BASENAME.ads2 from $ADSBASENAME.sed and $BASENAME.ads" 

# Apply this to the .ads file, generating the .ads2 file, which is included
# by the _ads.tex file.
sed -f $ADSBASENAME.sed $BASENAME.ads > $BASENAME.ads2

echo "Generating $ADSBASENAME.txt from $ADSBASENAME.tex and $BASENAME.ads2"

# Now run latex2html on the _ads.tex file
latex2html -no_navigation -info 0 -address 0 -html_version 4.0,math $ADSBASENAME

# Then, convert to utf8 with lynx and tidy it up a bit
lynx --justify=off --width=100 --dump --nolist --display_charset=utf-8 \
     $ADSBASENAME/$ADSBASENAME.html | cut -c4- \
     | sed 's/ibidrule/\\ibidrule/g' > $ADSBASENAME.txt
%</ads-shell-script>
%    \end{macrocode}

% \section{Epilogue}
% Err, that's it \dots

\endinput


% \endinput
% Local Variables:
% mode: doctex
% TeX-master: t
% End:
"

rmaa.aux: rmaa.dtx rmaa.ind rmaa.gls
        @echo "REGENERATING .aux FILE"
        pdflatex "\nonstopmode% \def\fileversion{4.0} \def\filedate{2019/11/12}
% \MakeShortVerb{\|}
% \title{A \LaTeX{} document class for the Revista Mexicana de
% Astronom\'\i a y Astrof\'\i sica\thanks{This file has version
% \fileversion{} last revised \filedate{}.}}
% \author{William J. Henney\thanks{w.henney@irya.unam.mx}}
% \date{\filedate}
% \maketitle
% \paragraph{Abstract}
% This document describes the \texttt{rmaa} \LaTeX{}
% document class for the journal ``Revista Mexicana de
% Astronom\'\i a y Astrof\'\i sica''. The macros can be used for 
% \begin{enumerate}
% \item Preparing author manuscripts for the journal (including
% conference proceedings to be published in the ``Serie de
% Conferencias'')
% \item Typesetting an entire volume of the journal by inputting the
% individual article files from a container file. Facilities are
% included for automatic generation of TOC and Author Index. 
% \end{enumerate}
% 
% \tableofcontents
% \section{The documentation driver file}
% This is the part that generates the documentation that you may or
% may not be reading, which is done with the help of the Makefile
% discussed in section~\ref{sec:driver_make}. 
%    \begin{macrocode}
%<*driver>
\documentclass{ltxdoc}
\usepackage{hyperref}
\usepackage{palatino}
\EnableCrossrefs \CodelineIndex \RecordChanges
\setlength{\parindent}{0pt}
\def\MacroFont{\footnotesize\tt}
\begin{document} \DocInput{rmaa.dtx} \PrintIndex \PrintChanges \end{document}
%</driver>
%    \end{macrocode}

% \section{The implementation}
% 
% \changes{3.24}{2007/09/14}{Added user access to the macro version
% and date via \cs{rmVersion} and \cs{rmDate}}
% 
%    \begin{macrocode}
%<*package>
\newcommand{\rmaa@version}{3.26}
\newcommand{\rmaa@date}{2011/08/15}
\newcommand\rmVersion{\rmaa@version}
\newcommand\rmDate{\rmaa@date}
\NeedsTeXFormat{LaTeX2e}[1995/06/01]
%    \end{macrocode}
% \changes{4.0}{2019/11/02}{Change abstract to appear before
% resumen. Change RevMexAA to RMxAA for consistency with ADS.}
%
% \changes{3.22}{2007/03/22}{Allow proceedings articles to have
%   received/accepted dates} \changes{3.21}{2006/05/15}{Changes for
%   Silvia - abstracts are now shipped out to ADS.}
% \changes{3.19}{2006/01/19}{Changes for Silvia - remove adresses from
%   end, new abstract style} \changes{2.2beta}{2001/10/30}{ Fixed bug
%   whereby twocolumn option was not passed through to book
%   class. Finally sorted out the ``label on last page'' issue. }
% \changes{2.1}{2001/01/23}{ Fixed some inconsistencies in
%   \cs{nocite} and \cs{bibitem} uniqification Fixed the definition of
%   \cs{ion} so it will work in math mode and in contexts where the
%   typesize is different from normal (e.g., title, fig captions).
%   Modified the definition of \cs{@caption}, so we don't write any
%   list-of-figures data (we never use it anyway, and it just causes
%   problems with fragile commands in fig captions). }
% \changes{2.0}{2000/11/30}{ Added `newstyle' option, which should be
%   used in all new proceedings volumes (wider abstract, smaller
%   typeface in fig captions, etc.) Added commands for formatting
%   table captions. New demonstation document rmpostsamp.tex for
%   1-page posters. This will be the last release under a PRI
%   government (for a while...)}  \changes{1.9}{2000/06/06}{
%   Corrected position of page number in footer of articles that open
%   on even pages. Ensured that first article in each session opens on
%   an odd page, as do list of abstracts and author index. Added extra
%   space in TOC before the backmatter. Added command for editors,
%   \cs{BoldSessionTitles}, that makes produces bold titles for
%   session in the TOC.}  \changes{1.8}{2000/05/16}{ Fixed bug in
%   last page label with \cs{suppressfulladdresses} }
% \changes{1.7}{2000/05/11}{ Fixed more little things. Page labels
%   now work, even if you have no bibliography. You can output the
%   full addresses by hand with the \cs{outputfulladdresses}
%   command. Editors' Guide updated but still not finished. Revamped
%   splitbook.sh now produces PS file of each article for web page and
%   also an EPS file of each page for sending to the printers.}
% \changes{1.6}{2000/04/15}{ More cleaning up: full stops after
%   titles in index, running headers and footers in frontmatter,
%   fontsize in title of index, etc} \changes{1.5}{2000/04/15}{
%   Eliminated superscripts from titles in TOC and index}
% \changes{1.4}{2000/04/14}{ Added copyright and license info (GPL)}
% \changes{1.3}{2000/03/01}{ Corrected \cs{figcaption} to do the
%   caption formatting right Corrected font size of sec headings in
%   front matter} \changes{1.2}{2000/02/16}{ Added \cs{plate}
%   command} \changes{1.1}{2000/02/14}{ Added manuscript option,
%   corrected bugs} \changes{1.0}{2000/02/08}{ Now does a whole book
%   - lots of stuff tidied up} \changes{0.7}{1999/10/04}{ Increased
%   spacing between sections in single column styles. Suppressed
%   floats at bottom of pages (didn't work though) }
% \changes{0.6}{1999/09/24}{ Tweaked \cs{adjustfinalcols} (still not
%   perfect) Added proceedings option Added onecolumn/twocolumn
%   options (orthogonal to journal/proceedings) }
% \changes{0.5}{1999/09/23}{ Added \cs{fulladdresses} and
%   \cs{adjustfinalcols} (just about working)}
% \changes{0.4}{1999/09/22}{ Minor stylistic changes}
% \changes{0.3}{1999/06/28}{ Better mechanism for alternative
%   affiliations Added `guide' option for user guide}
% \changes{0.2}{1999/05/27}{ All basic functionality present, I
%   think} \changes{0.1}{1999/05/26}{ Just about working - still
%   missing Fig/Tab/Refs} \changes{0.0}{1999/05/21}{ So bored today,
%   decided to rewrite RMAA LaTeX macros from scratch}
%    \begin{macrocode}
\ProvidesClass{rmaa}[\rmaa@date\space v\rmaa@version\space RMxAA article
class]
%    \end{macrocode}
% \subsection{Option Handling}
% \label{sec:options}
% Declarations that are needed in options
% \changes{3.05}{2003/10/22}{Added natbib option}
% \changes{3.06}{2003/01/03}{Natbib option turned on permanently}
% \changes{3.07}{2003/01/11}{Added proof option}
% \changes{3.14}{2003/01/11}{Option to turn off natbib}
% \changes{3.18}{2003/10/22}{trying out doctex mode} 
% \changes{3.23}{2007/09/09}{Manuscript is now on by default.}
% \changes{3.26}{2011/08/15}{New option for review articles}
%    \begin{macrocode}
\newif\if@reviewarticle
\@reviewarticlefalse
\newif\if@proof
\@prooffalse
\newif\if@journal
\@journaltrue
\newif\if@printdates
\@printdatesfalse
\newif\if@proceedings
\@proceedingsfalse
\newif\if@manuscript
\@manuscripttrue
\newif\if@bigfont
\@bigfonttrue
\newif\if@guide
\@guidefalse
\newif\if@preprint
\@preprintfalse
\newif\if@rmtwocol
\@rmtwocoltrue
\newif\if@biblist
\@biblistfalse
\newif\if@book
\@bookfalse
\newif\if@natbib
\@natbibtrue
\newif\if@rmcite
\@rmcitefalse
\newif\if@titlevisible
\@titlevisibletrue
\newif\if@abstractpresent
\@abstractpresentfalse
\newif\if@printbanner
\@printbannerfalse
\newif\if@suppressfulladdresses
\@suppressfulladdressesfalse
\newif\if@newstyle
\@newstylefalse
\newif\if@debug
\@debugfalse
\newif\if@boldkeywords
\@boldkeywordsfalse
%    \end{macrocode}
% The option handling---mainly setting various booleans.
% \changes{2.2beta11}{2001/11/26}{Made twocol and
% newstyle the default for proceedings.}
% \changes{3.23}{2007/09/09}{Manuscript is now 10 pt instead of 12pt}
% \changes{3.23}{2007/09/10}{Full addresses are now automatically
% suppressed for proceedings style.}
% \changes{3.24}{2007/09/14}{New option \texttt{boldkeywords} that
% does the keywords in the old ugly way. Now off by default.}
%    \begin{macrocode}
\DeclareOption{manuscript}{\typeout{Manuscript for referee} 
  \@manuscripttrue\@journalfalse\@rmtwocolfalse
  \PassOptionsToClass{10pt}{article}}
\DeclareOption{bigfont}{\typeout{Big font manuscript requested by Gary
  Ferland} 
  \@manuscripttrue\@journalfalse\@rmtwocolfalse\@bigfonttrue
  \PassOptionsToClass{11pt}{article}}
\DeclareOption{journal}{\typeout{Main journal option}%
  \@journaltrue\@rmtwocoltrue\@printdatestrue\@manuscriptfalse
  \PassOptionsToClass{10pt}{article}}
\DeclareOption{guide}{\typeout{User guide option}%
  \@guidetrue\@journalfalse\@rmtwocoltrue\@manuscriptfalse
  \PassOptionsToClass{10pt}{article}}
\DeclareOption{preprint}{\typeout{Preprint option}%
  \@preprinttrue\@manuscriptfalse\PassOptionsToClass{10pt}{article}}
\DeclareOption{proof}{\typeout{Author proof option}%
  \@prooftrue\@manuscriptfalse}
\DeclareOption{proceedings}{\typeout{Proceedings option}%
  \@proceedingstrue\@journalfalse\@rmtwocoltrue\@newstyletrue
  \@suppressfulladdressestrue
  \@printdatesfalse\@manuscriptfalse\PassOptionsToClass{10pt}{article}} 
\DeclareOption{printdates}{\typeout{Force printing of received/accepted dates}%
  \@printdatestrue} 
\DeclareOption{book}{\typeout{Whole issue option}%
  \@booktrue} 
\DeclareOption{nonatbib}{\typeout{Turn off natbib support (required
    for compatibility with old rmaacite package).}%
  \@natbibfalse} 
\DeclareOption{newstyle}{\typeout{Use new layout style}%
  \@newstyletrue}
\DeclareOption{oldstyle}{\typeout{Use old layout style}%
  \@newstylefalse}
\DeclareOption{debug}{\typeout{Loads of annoying messages}%
  \@debugtrue}
%    \end{macrocode}
% These two need to be declared explicitly so that |\ExecuteOptions|
% actually does something.
%    \begin{macrocode}
\DeclareOption{twocolumn}{\@rmtwocoltrue}
\DeclareOption{onecolumn}{\@rmtwocolfalse}
\DeclareOption{twoside}{\PassOptionsToClass{\CurrentOption}{article}}
\DeclareOption{boldkeywords}{\@boldkeywordstrue}
%    \end{macrocode}
% Bold, upper-case keywords used to be the default. Now they must be
% turned on explicitly with the \texttt{boldkeywords} option. 
%    \begin{macrocode}
\DeclareOption{boldkeywords}{\@boldkeywordstrue}
%    \end{macrocode}
% Pass any other option along to the \texttt{article} class.
%    \begin{macrocode}
\DeclareOption*{\PassOptionsToClass{\CurrentOption}{article}}
%    \end{macrocode}
% Default options are to be two-sided and based on the standard
% \LaTeX{} \texttt{article} class,
% unless we are processing the entire book, in which case we base
% ourselves on the standard \texttt{book} class. Default is to use
% manuscript option unless otherwise specified. 
%    \begin{macrocode}
\ExecuteOptions{twoside,manuscript} \ProcessOptions\relax
\if@book
\if@rmtwocol\LoadClass[twocolumn]{book}\else\LoadClass{book}\fi
\newcommand{\abstract}{\relax}
\newcommand{\@maketitle}{\relax}
\else
\if@rmtwocol\LoadClass[twocolumn]{article}\else\LoadClass{article}\fi
\fi
%    \end{macrocode}
% Attempt to use natbib for citations 02 Jan 2003. 
% \changes{3.05}{2003/01/02}{If the natbib option is used, then set
% things up so that natbib's version of the bibliography is identical
% to the standard one.}
% \changes{3.09}{2003/01/23}{Explicitly require a recent version of
% natbib, so that we get a decent error message if the package is too
% old.} 
% \changes{3.23}{2007/09/09}{Now allow Bibtex to work properly} 
% \changes{3.25}{2009/10/01}{Change \cs{newcommand} tp \cs{def} when setting \cs{bibfont}. This allows it to work with all versions of natbib.} 
% \changes{3.25}{2010/02/06}{Updated copyright year.} 
% \changes{3.25}{2010/02/06}{Updated email and web addresses to \texttt{crya.unam.mx}.} 
%    \begin{macrocode}
\if@natbib
\RequirePackage{natbib}[1999/02/25]
\bibpunct{(}{)}{;}{a}{}{,}
\def\bibfont{\small}
\renewcommand\bibsection{\section*{REFERENCES}}
\setlength\bibsep{0pt}
\setlength\bibhang{1.5em}
\bibliographystyle{rmaa}
\fi
%    \end{macrocode}
% We want the \texttt{graphicx} package for the extended syntax of
% |\includegraphics|. We want the \texttt{ifthen} package for using in
% these macros (although we don't much yet). We explicitly load
% \texttt{babel} since some \LaTeX{} distributions messed up the
% hyphenation if we didn't.  
%    \begin{macrocode}
\RequirePackage{graphicx}
\RequirePackage{ifthen}
\IfFileExists{babel.sty}{%
%   \let\@restoreindent=\@afterindentfalse
  \RequirePackage[english]{babel}
%   \@restoreindent
  %% reinstate the definition of \@roman so as not to break .idx files
  \@ifundefined{org@@roman}{\relax}{\renewcommand{\@roman}[1]{\org@@roman{##1}}}
  }{%
  \typeout{Package babel.sty not found - don't worry about this}}
%    \end{macrocode}
% Try to load the \texttt{dcolumn} and \texttt{booktabs} packages if
% they can be found. These can produce nicer-looking tables. 
%    \begin{macrocode}
\IfFileExists{dcolumn.sty}{%
  \RequirePackage{dcolumn}
  }{%
  \typeout{Package dcolumn.sty not found - you will not be able to
    align columns on decimal points in tables}}
\IfFileExists{booktabs.sty}{%
  \RequirePackage{booktabs}
  }{%
  \typeout{Package booktabs.sty not found - your tables will not look
    as pretty as they ought too.}
  }
\IfFileExists{makeidx.sty}{%
  \RequirePackage{makeidx}
}{%
  \typeout{Package makeidx.sty not found - Author Index may not work.}
  }
%% commented out for moment since requires more modern version of LaTeX
%% \RequirePackage[latin1]{inputenc}
%    \end{macrocode}
% \changes{4.0}{2019/11/13}{Make \texttt{hyperref} and
% \texttt{etoolbox} into hard requirements, since we now use them
% extensively.}  We use hyperref to get internal and external links.
% But I don't like the colored boxes so I turn those off.  We use
% etoolbox for patching and robustifying at the moment.  In the
% future, I want to port a lot of other stuff over to it.  For
% instance, replace the ifthen package and the TeX ifs too.
%
% We target the latest version (1.7) of PDF, rather than the default
% of 1.5.  This reduces the chance of getting warnings from included
% graphics.  This is not strictly related to hyperref, but we put it
% here anyway.
%    \begin{macrocode}
\RequirePackage{etoolbox}
\RequirePackage{hyperref}
\hypersetup{hidelinks}
\pdfminorversion=7
%    \end{macrocode}
% We patch the \cs{@combinedblfloats} following a suggestion by Heiko
% Oberdiek at \url{https://tex.stackexchange.com/a/249743/3843}, which
% avoids problems when hyperlinks are split across
% pages. \changes{4.0}{2019/11/12}{Fix for issue where breaking
% hyperlink across pages gave an error.}
%    \begin{macrocode}
\urlstyle{same}
\newcommand\email{\begingroup \urlstyle{rm}\Url}
\patchcmd\@combinedblfloats{\box\@outputbox}{\unvbox\@outputbox}{}{%
  \errmessage{\noexpand\@combinedblfloats could not be patched}%
}%
%    \end{macrocode}
% \subsection{Other setup stuff}\label{sec:setup}
% Some declarations: 
%    \begin{macrocode}
\newif\if@needaddresses
\@needaddressestrue
\newif\if@frontmatter
\@frontmatterfalse
\newif\if@inpreamble
\@inpreambletrue
\newif\if@firstauthor
\@firstauthortrue
\newif\if@firstaffil
\@firstaffiltrue
\newif\if@firstaltaffil
\@firstaltaffiltrue
\newif\if@firstkwd
\@firstkwdtrue
\newif\if@noexpandADS
\@noexpandADSfalse
\newif\if@sessiontitlepages
\@sessiontitlepagesfalse
%    \end{macrocode}
% \begin{macro}{\@rmaaEmpty}
% \begin{macro}{\@rmaaArg}
% These are used in various commands to check if arg was left blank
%    \begin{macrocode}
\newcommand{\@rmaaEmpty}{}
\newcommand{\@rmaaArg}{}
%    \end{macrocode}  
% \end{macro}
% \end{macro}
%
% \subsection{Compatibility with old \LaTeX{} versions}
% \label{sec:compat}
% This package is based on the 1998/12/01 release of \LaTeX{}.
% Older releases may not define all the commands that we need.
% Luckily, we can use |\providecommand| to define these if necessary. 
% The following commands are not defined in the 1994/06/01 release. 
%
% \begin{macro}{\MakeUpperCase}
% \begin{macro}{\textsuperscript}
% \begin{macro}{\textbackslash}
% User commands:
%    \begin{macrocode}
\providecommand{\MakeUpperCase}[1]{\uppercase{#1}}
\providecommand{\textsuperscript}[1]{\ensuremath{^\textrm{#1}}}
\providecommand{\textbackslash}{\ensuremath{\backslash}}
%    \end{macrocode}  
% \end{macro}
% \end{macro}
% \end{macro}
% \begin{macro}{\protected@xdef}
% \begin{macro}{\restore@protect}
% \begin{macro}{\@unexpandable@protect}
% \begin{macro}{\hb@xt@}
% \LaTeX{} internal commands:
%    \begin{macrocode}
\providecommand{\protected@xdef}{%
        \let\@@protect\protect
        \let\protect\@unexpandable@protect
        \afterassignment\restore@protect
        \xdef}
\providecommand{\restore@protect}{\let\protect\@@protect}
\providecommand{\@unexpandable@protect}{\noexpand\protect\noexpand}
\providecommand{\hb@xt@}{\hbox to}
%    \end{macrocode}
% \end{macro}
% \end{macro}
% \end{macro}
% \end{macro}
%
% \subsection{Page layout}\label{sec:page-layout}
%
% The dimensions of the text were defined by Elizabeth. Text width is
% 169\,mm with equal margins on the two sides: 
%    \begin{macrocode}
\setlength{\textwidth}{169mm}
\setlength{\@tempdima}{\paperwidth}
\addtolength{\@tempdima}{-\textwidth}
\setlength{\oddsidemargin}{0.5\@tempdima}
\addtolength{\oddsidemargin}{-1in}
\setlength{\evensidemargin}{\oddsidemargin}
%    \end{macrocode}
% Text height is 225\,mm, top margin is 14\,mm. 
%    \begin{macrocode}
\setlength{\textheight}{225mm} 
\setlength{\topmargin}{-1in} 
\addtolength{\topmargin}{14mm}
%    \end{macrocode}
% Space between columns in two-column styles, header, footer,
% paragraph indentation:
%    \begin{macrocode}
\setlength{\columnsep}{9.5mm}
\setlength{\headsep}{0.7cm}
\setlength{\footskip}{0.75cm}
\setlength{\parindent}{1.5em}
%    \end{macrocode}
% 
% \begin{macro}{\@set@scalable@length}
% Try out a new scheme for dealing with various spacing dimensions and
% allowing the user to easily modify them. The command
% \begin{quote}
% |\set@scalable@length|\marg{name}\marg{defaultlength}\marg{stretch}\marg{shrink}
% \end{quote}
% does a 
% |\setlength|\marg{\cs{name}}\marg{defaultlength plus stretch minus shrink}
% and also defines the user command
% \textbackslash\meta{name}|Scale|\marg{factor} that can 
% be used to rescale length by \meta{factor}. Additionally, now
% defines the package command
% \textbackslash\meta{name}|@set@scale|\marg{FACTOR} 
% that also rescales the length by \meta{FACTOR}, so that the length
% becomes \meta{factor} $\times$ \meta{FACTOR} $\times$
% \meta{defaultlength}.   
%    \begin{macrocode}
\newcommand\@set@scalable@length[4]{%
  %% arg #1: NAME of length without leading `\'
  %%
  %% arg #2: default VALUE  
  \expandafter\newlength\csname @#1@default\endcsname
  \expandafter\setlength\csname @#1@default\endcsname{#2}%
  \expandafter\newlength\csname @@#1@@default\endcsname
  \expandafter\setlength\csname @@#1@@default\endcsname{#2}%
  %% arg #3: default STRETCH  
  \expandafter\newlength\csname @#1@stretch\endcsname
  \expandafter\setlength\csname @#1@stretch\endcsname{#3}%
  \expandafter\newlength\csname @@#1@@stretch\endcsname
  \expandafter\setlength\csname @@#1@@stretch\endcsname{#3}%
  %% arg #4: default SHRINK  
  \expandafter\newlength\csname @#1@shrink\endcsname
  \expandafter\setlength\csname @#1@shrink\endcsname{#4}%
  \expandafter\newlength\csname @@#1@@shrink\endcsname
  \expandafter\setlength\csname @@#1@@shrink\endcsname{#4}%
  %%
  %% User-defined scale \@NAME@scale is unity by default
  \expandafter\gdef\csname @#1@scale\endcsname{1}%
  %% Package-defined scale \@NAME@SCALE is unity by default
  \expandafter\gdef\csname @#1@SCALE\endcsname{1}%
  %%
  %% internal command \@reset@NAME should expand to 
  %% \setlength\@@NAME@@default{\@NAME@SCALE\@NAME@default}
  %% \setlength\NAME{\@NAME@scale\@@NAME@@default}
  %% in other words: scale x SCALE x default 
  \expandafter\gdef\csname @reset@#1\endcsname{%
    \expandafter\setlength\csname @@#1@@default\endcsname{%
      \csname @#1@SCALE\endcsname \csname @#1@default\endcsname
      }
    \expandafter\setlength\csname @@#1@@stretch\endcsname{%
      \csname @#1@SCALE\endcsname \csname @#1@stretch\endcsname
      }
    \expandafter\setlength\csname @@#1@@shrink\endcsname{%
      \csname @#1@SCALE\endcsname \csname @#1@shrink\endcsname
      }
    \expandafter\setlength\csname #1\endcsname{%
      \csname @#1@scale\endcsname \csname @@#1@@default\endcsname
      \@plus
      \csname @#1@scale\endcsname \csname @@#1@@stretch\endcsname
      \@minus
      \csname @#1@scale\endcsname \csname @@#1@@shrink\endcsname
      }}
  %%
  %% User command \NAMEScale should expand to 
  %% {\renewcommand\@NAME@scale{#1}\@reset@NAME}
  \expandafter\gdef\csname #1Scale\endcsname##1{%
    \expandafter\gdef\csname @#1@scale\endcsname{##1}%
    \csname @reset@#1\endcsname}
  %%
  %% Package command \NAME@set@SCALE should expand to 
  %% {\renewcommand\@NAME@SCALE{#1}\@reset@NAME}
  \expandafter\gdef\csname #1@set@SCALE\endcsname##1{%
    \expandafter\gdef\csname @#1@SCALE\endcsname{##1}%
    \csname @reset@#1\endcsname}
  %%
  %% initialize at default values with \@reset@NAME
  \csname @reset@#1\endcsname
  }
%    \end{macrocode}
% \end{macro}
% \begin{macro}{\parskip}
% First use of this new infrastructure: set interparagraph space to be
% zero but with a small amount of stretchiness.
%    \begin{macrocode}
%   \ifthenelse{\boolean{@manuscript}}{%
%     \@set@scalable@length{parskip}{0.5\baselineskip}{0.3\p@}{0.1\p@}
%   }{
    \@set@scalable@length{parskip}{0\p@}{0.3\p@}{0\p@}
%   }
%    \end{macrocode}
% \end{macro}
% 
% \begin{macro}{\abovesectionskip}
% \changes{4.0}{2019/11/12}{Use negative values for skip so that first paragraph after section headings is not indented.  This is something that I have been trying to fix for years!  It turns out that the anser was in the documentation for the standard command \cs{@startsection.}}
% \begin{macro}{\belowsectionskip}
% \begin{macro}{\aboveparaskip}
% \begin{macro}{\belowparaskip}
% Define some utility rubber lengths for use in
% section~\ref{sec:section-commands}. 
%    \begin{macrocode}
\newlength{\abovesectionskip}
\newlength{\belowsectionskip}
\newlength{\aboveparaskip}
\newlength{\belowparaskip}
%    \end{macrocode}
% The default values for these lengths are those traditionally used
% for the onecolumn style. 
%    \begin{macrocode}
\@set@scalable@length{abovesectionskip}{-3.5ex}{-4ex}{-2ex}
\@set@scalable@length{belowsectionskip}{2.3ex}{0ex}{2ex}
\@set@scalable@length{aboveparaskip}{3.25ex}{2ex}{0.5ex}
\@set@scalable@length{belowparaskip}{-1em}{\z@}{\z@}
%    \end{macrocode}
% \end{macro}\end{macro}\end{macro}\end{macro}
% We also make scalable some of the standard float parameters. These
% are the same defaults as in the standard article classes, but
% slightly more stretchy. For the twocolumn style, they are currently
% halved, see below. Note that I have also lumped |\abovecaptionskip|
% in with these, since it plays a similar role. There is also a
% |\belowcaptionskip| but it is set to zero in the standard classes. 
%    \begin{macrocode}
\@set@scalable@length{floatsep}{12\p@}{4\p@}{4\p@}
\@set@scalable@length{textfloatsep}{20\p@}{4\p@}{8\p@}
\@set@scalable@length{intextsep}{12\p@}{4\p@}{4\p@}
\@set@scalable@length{dblfloatsep}{12\p@}{4\p@}{4\p@}
\@set@scalable@length{dbltextfloatsep}{20\p@}{4\p@}{8\p@}
\@set@scalable@length{abovecaptionskip}{10\p@}{\z@}{\z@}
%    \end{macrocode}
% \begin{macro}{\@before@title@skip}         
% \begin{macro}{\@before@author@skip}        
% \begin{macro}{\@between@author@skip}       
% \begin{macro}{\@before@abstract@skip}      
% \begin{macro}{\@after@abstracthead@skip}   
% \begin{macro}{\@between@abstract@skip}     
% \begin{macro}{\@before@kwd@skip}           
% \begin{macro}{\@before@dates@skip}           
% Spacing parameters used in |\@maketitle|. These are not changed by
% default except by the user, so they don't need to go into
% |\@package@default@lengths|.  
%    \begin{macrocode}
\newlength{\@before@title@skip}         
\newlength{\@before@author@skip}        
\newlength{\@between@author@skip}       
\newlength{\@before@abstract@skip}      
\newlength{\@after@abstracthead@skip}   
\newlength{\@between@abstract@skip}     
\newlength{\@before@kwd@skip}           
\newlength{\@before@dates@skip}           
\@set@scalable@length{@before@title@skip}{2em}{\z@}{\z@}
\@set@scalable@length{@before@author@skip}{1.5em}{\z@}{\z@}
\@set@scalable@length{@between@author@skip}{0.5em}{\z@}{\z@}
\@set@scalable@length{@before@abstract@skip}{3ex}{1ex}{0.5ex}
\@set@scalable@length{@after@abstracthead@skip}{1ex}{\z@}{\z@}
\@set@scalable@length{@between@abstract@skip}{2ex}{1ex}{0.5ex}
\@set@scalable@length{@before@kwd@skip}{1ex}{0.5ex}{0.2ex}
\@set@scalable@length{@before@dates@skip}{2ex}{0.5ex}{0.2ex}
%    \end{macrocode}
% \end{macro}\end{macro}\end{macro}\end{macro}
% \end{macro}\end{macro}\end{macro}\end{macro}
% 
% Spacing above and below displayed equations also wants to be done
% like this, but it is more complicated since the standard article
% classes make these change with the fontsize. Something to do later,
% maybe. 
% 
%    \begin{macrocode}
%    \end{macrocode}
% \begin{macro}{\@package@default@lengths}
% The value actually used for these lengths is controlled by the
% \meta{SCALE} parameter, which depends on the style option that is in
% effect. This is set using the |\@package@default@lengths| command,
% which must be called every time the style changes (e.g., changing
% from onecolumn to twocolumn). 
%    \begin{macrocode}
\newcommand{\@package@default@lengths}{%
  \ifthenelse{\boolean{@rmtwocol}}{%
    %% twocolumn mode
    \abovesectionskip@set@SCALE{0.5}
    \belowsectionskip@set@SCALE{0.5}
    \aboveparaskip@set@SCALE{0.5}
    \belowparaskip@set@SCALE{1}
    \floatsep@set@SCALE{0.5}
    \textfloatsep@set@SCALE{0.5}
    \intextsep@set@SCALE{0.5}
    \dblfloatsep@set@SCALE{0.5}
    \dbltextfloatsep@set@SCALE{0.5}
    }{%
    %% onecolumn mode
    \abovesectionskip@set@SCALE{1}
    \belowsectionskip@set@SCALE{1}
    \aboveparaskip@set@SCALE{1}
    \belowparaskip@set@SCALE{1}
    \floatsep@set@SCALE{1}
    \textfloatsep@set@SCALE{1}
    \intextsep@set@SCALE{1}
    \dblfloatsep@set@SCALE{1}
    \dbltextfloatsep@set@SCALE{1}
    }%
  }
\@package@default@lengths
%    \end{macrocode}
% \end{macro}
% \begin{macro}{\RescaleLengths}
% \begin{macro}{\RescaleSecLengths}
% \begin{macro}{\RescaleFloatLengths}
% \begin{macro}{\RescaleTitleLengths}
% User commands that will fiddle with all the scalable lengths. 
%    \begin{macrocode}
\newcommand{\RescaleLengths}[1]{%
  \RescaleSecLengths{#1}\RescaleFloatLengths{#1}\RescaleTitleLengths{#1}}
\newcommand{\RescaleSecLengths}[1]{%
%% sections
  \abovesectionskipScale{#1}
  \belowsectionskipScale{#1}
  \aboveparaskipScale{#1}
  }
\newcommand{\RescaleFloatLengths}[1]{%
%% floats
  \floatsepScale{#1}
  \textfloatsepScale{#1}
  \intextsepScale{#1}
  \dblfloatsepScale{#1}
  \dbltextfloatsepScale{#1}
  \abovecaptionskipScale{#1}
  }
\newcommand{\RescaleTitleLengths}[1]{%
%% maketitle
  \@before@title@skipScale{#1}           
  \@before@author@skipScale{#1}        
  \@between@author@skipScale{#1}       
  \@before@abstract@skipScale{#1}      
  \@after@abstracthead@skipScale{#1}   
  \@between@abstract@skipScale{#1}     
  \@before@kwd@skipScale{#1}           
  \@before@dates@skipScale{#1}           
  }
%    \end{macrocode}
% \end{macro}\end{macro}\end{macro}\end{macro}
% \changes{2.2beta14}{2002/03/17}{Finished implementation of new
% infrastructure for scalable lengths. Converted the above/below
% section skips to use it. Also, all the floatsep parameters. Also,
% all the vertical lengths that are used by |\@maketitle|. }

% \begin{macro}{\MSextramargin}
% For the \texttt{manuscript} option we have an extra 0.3\,inch
% vertical margin and 1\,inch horizontal margin for marginal notes (by
% hand or with |\marginpar|. 
%
% \changes{3.23}{2007/09/09}{Reduce baselinestretch in
%   \texttt{manuscript} style to 1.2. No longer use raggedright for
%   manuscript.}
%
% \changes{3.24}{2007/09/14}{Further reduction of baselinestretch in
%   \texttt{manuscript} style to 1.1 since we now have wider margins.}
%
% \changes{3.24}{2007/09/14}{Split \texttt{manuscript}'s extra margin
%   into separate horizontal and vertical parts.}
% 
%    \begin{macrocode}
\newlength{\MSextraHmargin}\newlength{\MSextraVmargin}
\setlength{\MSextraHmargin}{1in} 
\setlength{\MSextraVmargin}{0.3in} 
\if@manuscript
   \renewcommand{\baselinestretch}{1.1}
   \addtolength{\textwidth}{-2\MSextraHmargin}
   \addtolength{\oddsidemargin}{\MSextraHmargin}
   \setlength{\evensidemargin}{\oddsidemargin}
   \addtolength{\textheight}{-2\MSextraVmargin}
   \addtolength{\topmargin}{\MSextraVmargin}
%    \raggedright
\fi
%    \end{macrocode}
% \end{macro}
% Float parameters: don't allow bottom floats by default. 
%    \begin{macrocode}
\renewcommand{\fps@figure}{tp}
\renewcommand{\fps@table}{tp}
%    \end{macrocode}
%
% \begin{macro}{\rmaa@AuthorWidthFrac}
% \begin{macro}{\rmaa@AbstractWidthFrac}
% These parameters control what fraction of the total text width can
% be used by the list of authors and the abstract/resumen,
% respectively.  
%    \begin{macrocode}
\newcommand{\rmaa@AuthorWidthFrac}{1.0}
\newcommand{\rmaa@AbstractWidthFrac}{0.75}
%    \end{macrocode}
% Use a wide abstract to save space in the new style.
%    \begin{macrocode}
\if@newstyle\renewcommand{\rmaa@AbstractWidthFrac}{1.0}\fi
%    \end{macrocode}
% \end{macro}
% \end{macro}
%
% Some things to consider in future versions. Trying to avoid too many
% widows, etc. Apparently, the following will work best with
% \cs{flushbottom}. 
%    \begin{macrocode}
%%\widowpenalty=10000
%%\clubpenalty=10000
%    \end{macrocode}
% Also, should perhaps try getting rid of underfull/overfull box warnings. 
% \changes{3.10}{2003/01/24}{Set large values for \cs{hbadness},
% \cs{vbadness}, \cs{hfuzz}, \cs{vfuzz}. To reduce amount of stuff
% that gets written to the terminal every time you compile something.}
%    \begin{macrocode}
\hbadness=10000
\vbadness=10000
\setlength\hfuzz{\maxdimen}
\setlength\vfuzz{\maxdimen}
%    \end{macrocode}
% \subsection{Various text styles for different contexts}
% \label{sec:text-styles} 
% \begin{macro}{\rmaa@AffilStyle}
% \begin{macro}{\rmaa@PagesStyle}
% \begin{macro}{\rmaa@YearStyle}
% \begin{macro}{\rmaa@VolumeStyle}
% \begin{macro}{\rmaa@VolumenStyle}
% \begin{macro}{\rmaa@JournalNameStyle}
% \begin{macro}{\rmaa@ReceivedStyle}
% \begin{macro}{\rmaa@AbstractHeaderStyle}
% \begin{macro}{\rmaa@JournalNameStyle}
% Style commands for the different fields used by |\maketitle| in
% constructing the title page of the article. These have the form
% |\rmaa@XYZStyle{TEXT}|. I'm not sure these were really a good idea,
% but never mind \dots 
%    \begin{macrocode}
\newcommand{\rmaa@AffilStyle}[1]{{\normalsize\rmfamily#1}}
\newcommand{\rmaa@PagesStyle}[1]{{\normalsize\rmfamily#1}}
\newcommand{\rmaa@YearStyle}[1]{{\normalsize\rmfamily#1}}
\newcommand{\rmaa@VolumeStyle}[1]{{\normalsize\rmfamily\bfseries#1}}
\newcommand{\rmaa@VolumenStyle}[1]{{\Large\sffamily\bfseries#1}}
\newcommand{\rmaa@JournalNameStyle}[1]{{\normalsize\itshape#1}}
\newcommand{\rmaa@ReceivedStyle}[1]{{\normalsize\rmfamily\itshape#1}}
\newcommand{\rmaa@AbstractHeaderStyle}[1]{%
  {\normalsize\rmfamily\MakeUpperCase{#1}}}
\newcommand{\rmaa@TitleStyle}[1]{%
  {\large\bfseries\rmfamily\boldmath\MakeUpperCase{#1}\par}}
%    \end{macrocode}
% \end{macro}
% \end{macro}
% \end{macro}
% \end{macro}
% \end{macro}
% \end{macro}
% \end{macro}
% \end{macro}
% \end{macro}
%
% \begin{macro}{\rmaa@FrontSecStyle}
% Headings of ``front sections'' (see section~\ref{sec:front-sections})
%    \begin{macrocode}
\newcommand{\rmaa@FrontSecStyle}[1]{%
  {\normalsize\bfseries\rmfamily\MakeUpperCase{#1}}}
%    \end{macrocode}
% \end{macro}
%
% \begin{macro}{\rmaa@TOCTitleStyle}
% \begin{macro}{\rmaa@TOCPartStyle}
% \begin{macro}{\rmaa@IndexTitleStyle}
% \begin{macro}{\rmaaTOCAuthorStyle}
% \begin{macro}{\rmaaIndexAuthorStyle}
% Commands to set style of TOC and index entries.
% Rewritten to avoid unnecessary font commands in the .toc and .ind files
%    \begin{macrocode}
\newcommand{\rmaa@TOCTitleStyle}[1]{\MakeUpperCase{#1}}
\newcommand{\rmaa@TOCPartStyle}[1]{\rmaa@TOCTitleStyle{#1}}
\newcommand{\rmaa@IndexTitleStyle}[1]{#1}
\newcommand{\rmaaTOCAuthorStyle}[1]{%
  {\itshape#1}}
\newcommand{\rmaaIndexAuthorStyle}[1]{%
  \protect\mbox{\bfseries#1}}
%    \end{macrocode}
% \end{macro}
% \end{macro}
% \end{macro}
% \end{macro}
% \end{macro}
%
% \begin{macro}{\rmaa@KeywordStyle}
% \begin{macro}{\rmaa@KeywordHeaderStyle}
% The style of the keywords depends on whether we are using the
% \texttt{newstyle} option or not. If we are, then they are small. 
%    \begin{macrocode}
\ifthenelse{\boolean{@newstyle}}{
  \newcommand{\rmaa@KeywordStyle}[1]{%
    {\small%
      \if@boldkeywords 
      \bfseries\MakeUpperCase{\raggedright #1}
      \else
      \raggedright #1
      \fi
    }}
  \newcommand{\rmaa@KeywordHeaderStyle}[1]{%
    {\small\itshape#1}}}
{
  \newcommand{\rmaa@KeywordStyle}[1]{%
    {\normalsize%
      \if@boldkeywords\bfseries\fi\raggedright #1}}
  \newcommand{\rmaa@KeywordHeaderStyle}[1]{%
    {\normalsize\itshape#1}}}
%    \end{macrocode}
% \end{macro}
% \end{macro}
%
% \subsection{Page styles}\label{sec:page-styles}
% \begin{macro}{\ps@rmaa}
% The running headers for normal pages:
%    \begin{macrocode}
\newcommand{\ps@rmaa}{%
%    \end{macrocode}
% nothing in either foot; 
%    \begin{macrocode}
  \renewcommand\@oddfoot{\relax}%
  \renewcommand\@evenfoot{\relax}%
%    \end{macrocode}
% page number and short author list in the even heads; 
%    \begin{macrocode}
  \renewcommand\@evenhead{\rlap{\thepage}\hfil\@shortauthor\hfil}%
%    \end{macrocode}
% short title and page number in the odd heads.
%    \begin{macrocode}
  \renewcommand\@oddhead{\hfil\@shorttitle\hfil\llap{\thepage}}
  }
%    \end{macrocode}
% \begin{macro}{\ps@rmaafirstpage}
% The header of the first page of each
% article:
%    \begin{macrocode}
\newcommand{\ps@rmaafirstpage}{%
  \renewcommand\@oddfoot{\hfil\thepage}%
%    \end{macrocode}
% footer corrected (06 Jun 2000) to have page on outside edge;
%    \begin{macrocode}
  \renewcommand\@evenfoot{\thepage\hfil}%
%    \end{macrocode}
% Header has journal name, volume, and page numbers for a regular
% article, or similar for special cases. 
% \changes{4.0}{2019/11/03}{Change to three-line header with copyright info and DOI.}
%    \begin{macrocode}
\if@preprint
\renewcommand\@oddhead{%
  \if@proceedings 
  \textsc{To appear in ``\@conftitle{} (\rmaa@YearStyle{\@year})''}\hfil
  \textit{RMxAC}
  \else
  \textsc{To appear in} \@journalname\hfil
  \fi}\fi%
\if@guide
\renewcommand\@oddhead{\textsc{Documentation} (\rmaa@YearStyle{\@year})\hfil}
\fi%
\if@manuscript
\renewcommand\@oddhead{Manuscript for \@journalname{} (\rmaa@YearStyle{\@year})\hfil}
\fi%
\if@journal
\renewcommand\@oddhead{%
  \parbox{\textwidth}{
    \@journalname, \rmaa@VolumeStyle{\@volume},
    \rmaa@PagesStyle{\@artpages}\
    (\rmaa@YearStyle{\@year})\hfil\\[1pt]
    {\small\textcopyright{} \@year:
      Instituto de Astronom\'\i{}a,
      Universidad Nacional Aut\'onoma de M\'exico}\hfil\\[1pt]
    {\small\@DOI@url}\hfil
  }%
}\fi%
  \renewcommand\@evenhead{\@oddhead}}
%    \end{macrocode}
% By default, use the |\ps@rmaa| page style.
%    \begin{macrocode}
\pagestyle{rmaa}
%    \end{macrocode}
% \end{macro}
% \end{macro}
%
% \subsection{Preamble commands}\label{sec:preamble} 
% Everything you need for setting the title page, TOC and index
% entries, etc.
% \subsubsection{Commands for the title page}
% \label{sec:title-commands}
% \begin{macro}{\resumen}
% \begin{macro}{\@resumen}
% \begin{macro}{\abstract}
% \begin{macro}{\@abstract}
% User commands for setting the resumen and abstract. 
% These should be used in the document header. The internal commands
% |\@resumen| and |\@abstract| are used to hold the actual text
% pending output via the |\maketitle| command. 
%    \begin{macrocode}
\newif\if@setresumen
\newcommand{\resumen}[1]{%
%%  \renewcommand{\@resumen}{{\selectlanguage{spanish}#1}}%
  \renewcommand{\@resumen}{{\ignorespaces#1}}%
  \@setresumentrue}
\renewcommand{\abstract}[1]{%
  \renewcommand{\@abstract}{{\ignorespaces#1}}
  \renewcommand\ADS@abstract{#1}
%% Not a good idea, since then author macros are not expanded and
%% cause errors later. However, by not doing it we are vulnerable to
%% fragile commands.....
  \if@noexpandADS\@onelevel@sanitize{\ADS@abstract}\fi
  \@abstractpresenttrue
  }
\newcommand{\ADSnoexpand}{\@noexpandADStrue}
%    \end{macrocode}
% \begin{macro}{\boldabstract}
% This is the version for doing 1-page conference proceedings. Still
% save the text for ADS but then write it out immediately as a bold
% paragraph. 
%    \begin{macrocode}
\newcommand{\boldabstract}[1]{%
  \gdef\ADS@abstract{#1}%
  \renewcommand{\@abstract}{#1}%
  \if@noexpandADS\@onelevel@sanitize{\ADS@abstract}\fi%
  \noindent\textbf{\boldmath #1}
  }
%    \end{macrocode}
% \end{macro}
% Default value in each case is to print out a warning if not set
%    \begin{macrocode}
  \newcommand{\@resumen}{%
    \@setresumenfalse \ClassWarningNoLine{rmaa}{No Spanish abstract
      (resumen) specified - if you do not speak Spanish, this will be
      translated by the editors}%
    \textbf{Favor de proporcionar un resumen en espa\~nol. If you are
      unable to translate your abstract into Spanish, the editors will
      do it for you.} }
  \newcommand{\@abstract}{\ClassWarningNoLine{rmaa}{No abstract
      specified}}
%    \end{macrocode}
% \end{macro}
% \end{macro}
% \end{macro}
% \end{macro}
%
% \begin{macro}{\title}
% \begin{macro}{\@rawtitle}
% \begin{macro}{\@TOCtitle}
% \begin{macro}{\rmaaIndexTitle}
% Similarly, |\title| is the preamble command to set the title of the
% article. This is used in various places, typeset in different
% styles, hence the use of the internal |\@rawtitle|, |\@title|,
% |\rmaaIndexTitle|, and |\@TOCtitle| commands, each of which uses the
% relevant style from section~\ref{sec:text-styles}
%    \begin{macrocode}
\newcommand{\@rawtitle}{Raw Title}
\newcommand{\@TOCtitle}{\rmaa@TOCTitleStyle{Title for TOC}}
\newcommand{\rmaaIndexTitle}{\rmaa@IndexTitleStyle{Title for index}.}
\newcommand{\@TOCaftertitle}{}
\renewcommand{\title}[1]{%
%%  {\ADS@sanitize\xdef\@rawtitle{#1}}
  \renewcommand\@rawtitle{\noexpand#1}
  \let\ADS@title=\@rawtitle
  \renewcommand{\@title}{\rmaa@TitleStyle{#1}}%
  \renewcommand{\@TOCtitle}{\rmaa@TOCTitleStyle{#1}\@TOCaftertitle}%
  \renewcommand{\rmaaIndexTitle}{\rmaa@IndexTitleStyle{#1\@TOCaftertitle}.}%
  }
%    \end{macrocode}
% \end{macro}
% \end{macro}
% \end{macro}
% \end{macro}
% 
% All the weird and whacky ways of indicating the author
% affiliations. 
% \begin{macro}{\affil}
% First, the way that sets it directly underneath the authors name.
% Note that |\affil| trumps all other affiliations for the purposes of
% getting into the ADS stuff (we use |\gdef|).   
% \dots 
%    \begin{macrocode}
\newcommand{\affil}[1]{%
  \\[1.5ex]\rmaa@AffilStyle{#1}\\[1.5ex]%
  \if@firstaffil\gdef\ADS@affil{#1}\@firstaffilfalse\fi%
  }
\renewcommand{\and}{and\\[1.5ex]}
%    \end{macrocode}
% \begin{macro}{\altaffil}
% Alternative affiliations are footnoted. 
%    \begin{macrocode}
\newcommand{\altaffil}[1]{%
  \if@firstaltaffil\gdef\ADS@affil{#1}\@firstaltaffilfalse\fi%
  \thanks{#1}}
%    \end{macrocode}
% \begin{macro}{\altaffilmark}
% \begin{macro}{\altaffiltext}
% Alternative-alternative mechanism for affiliations---should not be
% mixed with |\altaffil| (unless you know what you are doing) or there
% will be problems with numbering.
%    \begin{macrocode}
\newcommand{\altaffilmark}[1]{\textsuperscript{#1}}
\newcommand{\altaffiltext}[2]{%
  \if@firstaltaffil\gdef\ADS@affil{#2}\@firstaltaffilfalse\fi
  \stepcounter{footnote}%
  \protected@xdef\@thanks{\@thanks
    \protect\footnotetext[#1]{#2}}} 
%    \end{macrocode}
% \end{macro}
% \end{macro}
% \end{macro}
% \end{macro}
%
% \begin{macro}{\keywords}
% \begin{macro}{\@keywords}
% \begin{macro}{\@keywordheader}
% \label{keywords}Sundry minor preamble commands. All
% self-explanatory, I hope. |\keywords| is the old user interface
% (deprecated). New articles should use |\addkeyword| once for each
% keyword. That way the ADS record comes out right too (comma rather
% than em-dash separator).
% \changes{2.2beta12}{2001/12/07}{\cs{addkeyword} New user interface
%   for keywords.} 
% \changes{2.2beta13}{2002/03/11}{\cs{addkeyword} Improved behaviour
%   with .ads file. Use
%   \cs{ADS@sanitize} to change definitions of `~' and \cs{ion}.}
% \changes{3.23}{2007/09/09}{\cs{addkeyword} now signals a compilation
% error if the author tries to add more than 6 keywords.}
% \changes{3.24}{2007/09/14}{kwdcounter is now reset in
% \cs{@begofarticlehook} to prevent error in book style} 
%    \begin{macrocode}
\newcommand{\@kwdsep}{ --- }
\newcommand{\ADS@kwdsep}{, }
\newcounter{kwdcounter}\setcounter{kwdcounter}{0}
\newcommand{\addkeyword}[1]{%
  \addtocounter{kwdcounter}{1}
  \ifthenelse{\thekwdcounter > 6}{
    \ClassError{rmaa}{Too many keywords. A maximum of 6 keywords is allowed.}{}
  }{}
  \if@firstkwd%
  \renewcommand\@keywords{\if@boldkeywords\uppercase{#1}\else #1\fi}
  {\ADS@sanitize\def\ion{\plainion}\gdef\ADS@keywords{#1}}
  \@firstkwdfalse
  \else%
  \edef\@keywords{\@keywords\@kwdsep \if@boldkeywords\uppercase{#1}\else #1\fi}%
  {\ADS@sanitize\def\ion{\plainion}\xdef\ADS@keywords{\ADS@keywords\ADS@kwdsep #1}}%
  \fi}
\newcommand{\ADS@sanitize}{%
  \def\mathrm{}
  \def\@{\ }
  \def~{\ }
  \def\langle{<}
  \def\rangle{>}
  \def\log{log }
  \def\sin{sin }
  \def\cos{cos }
  \def\tan{tan }
  \def\fd{^d.}
  \def\fh{^h.}
  \def\fm{^m.}
  \def\fs{^s.}
  \def\fdg{�.}
  \def\farcm{'.}
  \def\farcs{''.}
  \def\arcdeg{�}
  \def\,{ }
  }
\newcommand{\keywords}[1]{%
  \renewcommand{\ADS@keywords}{#1}
  \renewcommand{\@keywords}{\uppercase{#1}}
  }
\newcommand{\@keywordheader}{\rmaa@KeywordHeaderStyle{Key Words:}}
\newcommand{\@keywords}{%
  \ClassWarningNoLine{rmaa}{No keywords specified}%
  Keywords go here
  }
%    \end{macrocode}
% \begin{macro}{\shorttitle}
% \begin{macro}{\shortauthor}
% \begin{macro}{\@shorttitle}
% \begin{macro}{\@shortauthor}
% Commands for setting the running headers.
%    \begin{macrocode}
\newcommand{\shorttitle}[1]{%
  \renewcommand\@shorttitle{\MakeUpperCase{#1}}}
\newcommand{\shortauthor}[1]{%
  \renewcommand\@shortauthor{\MakeUpperCase{#1}}}
\newcommand{\@shorttitle}{SHORT TITLE GOES HERE}
\newcommand{\@shortauthor}{SHORT AUTHOR GOES HERE}
%    \end{macrocode}
% \end{macro}
% \end{macro}
% \end{macro}
% \end{macro}
% \end{macro}
% \end{macro}
% \end{macro}


% \subsubsection{The full addresses}
% \label{sec:full-addresses}
% \begin{macro}{\fulladdresses}
% \begin{macro}{\@fulladdresses}
% \begin{macro}{\@fulladdheight}
% \begin{macro}{\@fulladddepth}
% \begin{macro}{\@fulladdbox}
% Aaghh, the full addresses! A perennial source of problems. 
%    \begin{macrocode}
\newlength{\@fulladdheight}
\newlength{\@fulladddepth}
\newsavebox{\@fulladdbox}
%% \fulladdresses command modified (11 May 2000) so that we can put
%% \fulladdresses{} in the article to suppress the output of the
%% addresses
\newcommand{\@FinalFigure}{\relax}
\newcommand{\FinalFigure}[1]{\renewcommand{\@FinalFigure}{#1}}
\newcommand{\fulladdresses}[1]{%
  \renewcommand{\@rmaaArg}{#1}
  \ifx\@rmaaArg\@rmaaEmpty
     \renewcommand{\@fulladdresses}{}
  \else
     \renewcommand{\@fulladdresses}{
       \normalsize%
       \ifx\@FinalFigure\relax\else\vfill\@FinalFigure\vfill\fi
%%       \vspace*{1ex}
       %% do addresses as a list (copied from thebibliography environment) 
       \list{\null}{\leftmargin 1.5em\topsep\z@\labelwidth\z@\labelsep\z@\itemindent -1.5em
         \itemsep\z@\parsep\z@\usecounter{enumi}}
       \sloppy\clubpenalty4000\widowpenalty4000
       \sfcode`\.=1000\relax
       #1
       \endlist
       }
  \fi
  }
\newcommand{\@fulladdresses}{Full addresses go here}
%    \end{macrocode}
% \begin{macro}{\suppressfulladdresses}
% New user-level command for suppressing full addresses should you so
% desire. Currently used in the single-page-poster style. 
%    \begin{macrocode}
\newcommand{\suppressfulladdresses}{\@suppressfulladdressestrue\@needaddressesfalse}
%    \end{macrocode}
% \end{macro}
% \end{macro}
% \end{macro}
% \end{macro}
% \end{macro}
% \end{macro}

% \begin{macro}{\printdates}
% \begin{macro}{\noprintdates}
% Command to activate the printing out of the received/accepted
% dates. This can also be done using the \texttt{printdates} option,
% which is the default for \texttt{journal}, but not for \texttt{proceedings}.
%    \begin{macrocode}
\newcommand{\printdates}{\@printdatestrue}
\newcommand{\noprintdates}{\@printdatesfalse}
%    \end{macrocode}
% \end{macro}
% \end{macro}
% \begin{macro}{\@outputaddressesifnecessary}
% Command to output the addresses if we haven't already done so
% This is called by: 
% \begin{enumerate}
% \item The end of |thebibliography| environment
% \item The |\outputfulladdresses| command
% \end{enumerate}
%    \begin{macrocode}
\newcommand{\@outputaddressesifnecessary}{%
  \ClassInfo{rmaa}{Calling @outputaddressesifnecessary}
  \ifthenelse{\boolean{@needaddresses}}{%
    %% If in one column style, or in the first column of a two column
    %% style then we can try to put it on this page 
    \ifthenelse{\boolean{@firstcolumn} \or \not\boolean{@rmtwocol}}{%
      \@outputaddresses}{%
      %% otherwise give warning it will be on a separate page
      \ClassWarning{rmaa}{Full addresses will go on a separate
        page (please see User Guide for how to remedy this)}}%
    %% If we haven't managed yet, put it on its own page
    \ifthenelse{\boolean{@needaddresses}}{%
      \if@debug{\ClassWarning{rmaa}{In this silly bit!}}\fi
      \clearpage\vfill\@outputaddresses}{}%
    }{
    %% If we don't need them, we still might need to do the eoa hook 
    \ifthenelse{\boolean{@suppressfulladdresses}}{%
      \if@debug{\ClassWarning{rmaa}{XXX}}\fi
%%      \if@book\else\@endofarticlehook\null\fi
      }{}
    }%
  }  
%    \end{macrocode}
% \end{macro}

% \begin{macro}{\outputfulladdresses}
% A user interface to the above command, in case there is no biblio
% (11 May 2000) 
%    \begin{macrocode}
\newcommand{\outputfulladdresses}{\@outputaddressesifnecessary}
%    \end{macrocode}
% \end{macro}

% \begin{macro}{\adjustfinalcols}
% This command starts a new column and should be 
% used to ensure that the columns balance on the last page. It also
% outputs the full 
% author addresses to ensure that this occurs in the left hand
% column. 
% The |\par| command is so that the current line will be taken account
% of when working out |\pagetotal|
%    \begin{macrocode}
\newcommand{\adjustfinalcols}{\if@twocolumn\goodbreak\@outputaddresses\newpage\noindent\fi}
%    \end{macrocode}
% \end{macro}

% \begin{macro}{\@outputaddresses}
% Command to do the work of outputting the addresses. 
% This can be called by 
% \begin{enumerate}
% \item an explicit |\adjustfinalcols|
% \item via |\@outputaddressesifnecessary|, q.v.
% \end{enumerate}
%
%    \begin{macrocode}
\newcommand{\@outputaddresses}{%
  % Allow option of suppressing the full addresses, in which case
  % we just run the hook if we are not in a book style
  \if@suppressfulladdresses
%%     \if@book\else\@endofarticlehook\null\fi
  \else
  %% save the full addresses in a box for output later
  \sbox{\@fulladdbox}{\parbox{\textwidth}{\@fulladdresses}}
  % calculate height of address box  
  \settoheight{\@fulladdheight}{\usebox{\@fulladdbox}}
  \settodepth{\@fulladddepth}{\usebox{\@fulladdbox}}
  \addtolength{\@fulladdheight}{\@fulladddepth}
  \if@debug
  { \typeout{Address total height: \the\@fulladdheight}
    \typeout{Text height: \the\pagetotal}
    \typeout{Goal height: \the\pagegoal}
    }
  \fi
  % reuse \@fulladdheight to mean the total height of text+addresses
  \addtolength{\@fulladdheight}{\pagetotal}
%%  \addtolength{\@fulladdheight}{-\baselineskip}
  \if@debug\typeout{Adjusted text + addresses height: \the\@fulladdheight}\fi
  % check it will fit on current page
  \ifdim\@fulladdheight<\pagegoal
     %% must shift it to the left
     %% to compensate for the hanging indentation in the reference list. 
     %% NOTE: this is not foolproof!
     %% NB. This will give overfull hbox warnings, which is a bit ugly.
     \vfill\noindent\if@biblist\hspace*{-\leftmargin}\fi
     \usebox{\@fulladdbox}
     % Set this to show we've done it
     \@needaddressesfalse
     % and assume we are at the end of the article
     % so that we do the EOA hook if not in the book style
%%     \if@book\else\@endofarticlehook\fi
  \else
     % if not, just do nothing. With luck, this function will get 
     % called again on the next page
     \if@debug{\typeout{Holding over full addresses to next page}}\fi
  \fi\fi
  }
%    \end{macrocode}
% \end{macro}

% \subsubsection{Commands for editors}
% \label{sec:editor-commands}
% Some user/editor commands for doing the TOC and index in the book
% option. 
% They are defined here rather in the book-specific part since they
% need to be always available in order to avoid breaking documents. 
% \begin{macro}{\listofauthors}
% \begin{macro}{\@rawauthors}
% \begin{macro}{\rmaaTOCauthors}
% The command |\listofauthors| formats an author list for use in the
% TOC. It should contain the same names as in the |\author| command,
% but without the affiliation garbage and using |\&| instead of
% |and|. 
%    \begin{macrocode}
\newcommand{\@rawauthors}{Raw Authors}
\newcommand{\rmaaTOCauthors}{\noexpand\rmaaTOCAuthorStyle{Author list for TOC}}
\newcommand{\listofauthors}[1]{%
  \renewcommand{\@rawauthors}{#1}
  \renewcommand\rmaaTOCauthors{\noexpand\rmaaTOCAuthorStyle{#1}}}
%    \end{macrocode}
% \begin{macro}{\indexauthor}\label{indexauthor}
% The command |\indexauthor| produces an index entry for each
% author. 
% This command must come \emph{after} the
% |\title| and the |\listofauthors| commands. Also, there should be one 
% |\indexauthor| command for \emph{each} author. We also add the
% author onto the end of the |\ADS@authors| macro and construct the
% |\ADS@authorinitial| by a most cunning method (see below), but only
% for the first author. 
% 
% No, No, NO! This was all wrong. Backed out completely: 
% \begin{quote}
% Completely rewritten (19 Dec 2002) since we now use the index class, so 
% as to be able to alphabetize accented characters. A side effect of this is
% that |\index| no longer expands out the argument after the |@|, so we 
% instead add all the |\index| commands to a list (|\indexauthor@deferlist|), 
% which is executed in the |\@begofarticlehook|. 
% 
% Note the stuff with |\toks0| is copied from Exercise 20.15 of the \TeX{}book---it is
% to make sure the latin1-accented name gets passed into the idx
% file. 
% \end{quote}
% We now go back to the old way of doing things, using an optional
% argument to the \cs{indexauthor} command that has the name without
% accents. 
% \changes{3.08}{2003/01/16}{Remerged changes from 3.06 that got
% dropped from 3.07 after a synchronization disaster. Basically, the
% accentless optional argument to \cs{indexauthor}. Must add the macro
% directories to my unison set-up to stop this sort of thing from happening.}
%    \begin{macrocode}
\newcommand\@isortkey{}
\newcommand\@emptykey{empty}
\newcommand\indexauthor[2][empty]{
  \renewcommand\@isortkey{#1}
  \ifx\@isortkey\@emptykey \renewcommand\@isortkey{#2}\fi
  \index{\@isortkey @{{\rmaaIndexAuthorStyle{#2}}\hskip1em plus1em minus0.4em
      \rmaaIndexTitle\hskip1em plus2em minus0.4em\penalty-100
      \rmaaTOCauthors}}
  \indexauthorADS{#2}
}
%%\def\indexauthor@deferlist{\relax}
%%\long\def\indexauthor#1{
%%  \def\@@name{#1} \toks0=\expandafter{\@@name}
%%  \toks1=\expandafter{\indexauthor@deferlist}
%%  \protected@xdef\indexauthor@deferlist{%
%%    \the\toks1
%%    \noexpand\index{\the\toks0@{\noexpand\rmaaIndexAuthorStyle{\the\toks0}\noexpand\hskip1em plus1em minus0.4em 
%%        \rmaaIndexTitle\noexpand\hskip1em plus2em minus0.4em\noexpand\penalty-100
%%        \rmaaTOCauthors}}
%%  }
%%  \indexauthorADS{#1}
%%  }
%%\xdef\indexauthor@deferlist{}
\newcommand\indexauthorADS[1]{%
  \if@firstauthor\protected@xdef\ADS@authorinitial{\firstletter{#1}}\fi
  \protected@xdef\ADS@authors{\ADS@authors\if@firstauthor\else ; \fi
    #1}
  \@firstauthorfalse
}
%    \end{macrocode}
% Note that the first |#1| sets the key for the index, making sure it
% actually starts with  
% first letter of the surname. Meanwhile, the |@{...}| expression,
% which starts with some  
% control sequences is what is written to the |.ind| file. 
% See the makeindex man page for details. 
% \end{macro}
% \end{macro}
% \end{macro}
% \end{macro}

% \begin{macro}{\KillAllAccents}
% \begin{macro}{\firstletter}
% These are utility functions used in the above. |lett| is a counter
% indicating the current letter in the string. |\firstletter| just
% prints the first letter of its argument.
%    \begin{macrocode}
\newcommand{\KillAllAccents}{\def\'{}\def\`{}\def\~{}\def\.{}\def\i{i}}
\newcommand\firstletter[1]{\@car #1\@nil}
%    \end{macrocode}
% \end{macro}
% \end{macro}
% \begin{macro}{\SetVolume}
% \begin{macro}{\SetMSnumber}
% \begin{macro}{\SetFirstPage}
% \begin{macro}{\SetYear}
% \begin{macro}{\ReceivedDate}
% \begin{macro}{\AcceptedDate}
% These are for the editors to set the volume number, etc. 
%    \begin{macrocode}
\newcommand{\SetVolume}[1]{%
  \renewcommand{\@volume}{#1}
  \edef\ADS@volume{\PadToFourWithDots{#1}}
  \renewcommand{\@volumen}{\rmaa@VolumenStyle{volumen #1}}%
  }
\newcommand{\SetMSnumber}[1]{\renewcommand{\@MSnum}{#1}}
%    \end{macrocode}
% Note that there is no real reason to use |\SetFirstPage| since all
% the page numbering is done automatically in the book style. 
%    \begin{macrocode}
\newcommand{\SetFirstPage}[1]{\setcounter{page}{#1}}
\newcommand{\SetYear}[1]{%
  \renewcommand{\@year}{#1}
  \renewcommand{\ADS@year}{#1}
  \renewcommand{\@anyo}{\rmaa@VolumenStyle{#1}}
  }
\newcommand{\ReceivedDate}[1]{\renewcommand{\@received}{#1}}
\newcommand{\AcceptedDate}[1]{\renewcommand{\@accepted}{#1}}
%    \end{macrocode}
%
% \begin{macro}{\DOI}
% User/editor command to enter the DOI, which will be typeset as a link.
% \changes{4.0}{2019/11/02}{New user command \cs{DOI} to typeset
% DOI in the header of first page and add it to ADS info.}
%    \begin{macrocode}
\newcommand\DOI[1]{%
  \def\@DOI@url{\url{https://doi.org/10.22201/ia.01851101p.#1}}%
  \def\ADS@doi{https://doi.org/10.22201/ia.01851101p.#1}%
}%
%    \end{macrocode}
% \end{macro}
%
% \begin{macro}{\ORCIDiD}
% \changes{4.0}{2019/11/12}{Add mechanism for giving ORCID iD links}
% \begin{macro}{\ORCIDiDfull}
% \begin{macro}{\@orcid@icon}
% \begin{macro}{\@orcid@iconbig}
% Commands to add icon links to ORCID iDs of authors. For example
% \begin{quote}
% |\ORCIDiD|\marg{0000-0000-0000-0000}
% \end{quote}
% to be used after each author name in the |\author|
% command and in the |\fulladdresses|.  Alternatively,
% |\OORCIDiDfull| will typeset the icon plus the full text of the
% link, but we don't use that at present.  The icon image is stored in |ORCIDiD_iconvector.pdf|, which needs to be on the input path.  In principle, we could make EPS versions if necessary.
%    \begin{macrocode}
\newcommand\@orcid@icon{\includegraphics[height=1.8ex]{ORCIDiD_iconvector}}
\newcommand\@orcid@iconbig{\includegraphics[height=2ex]{ORCIDiD_iconvector}}
\newcommand\ORCIDiD[1]{\href{https://orcid.org/#1}{\,\@orcid@icon}}
\newcommand\ORCIDiDfull[1]{%
  \href{https://orcid.org/#1}{\@orcid@iconbig}\,\url{https://orcid.org/#1}}
%    \end{macrocode}
% \end{macro}
% \end{macro}
% \end{macro}
% \end{macro}
%
% \begin{macro}{\@volume}
% \begin{macro}{\@volumen}
% \begin{macro}{\@artpages}
% \begin{macro}{\@year}
% \begin{macro}{\@anyo}
% \begin{macro}{\@received}
% \begin{macro}{\@accepted}
% \begin{macro}{\@MSnum}
% Default values for the containers used by above commands, plus a
% few more useful ones. 
%    \begin{macrocode}
\newcommand{\@volume}{00}
\newcommand{\@volumen}{\rmaa@VolumenStyle{00}}
\newcommand{\@artpages}{%
  \oldpageref{\articlename:RMAAFirstPage}--\oldpageref{\articlename:RMAALastPage}}
\newcommand{\@year}{\the\year}
\newcommand{\@anyo}{\rmaa@VolumenStyle{2007}}
\newcommand{\@journalname}{\rmaa@JournalNameStyle{%
    \if@proceedings\RMAASClong\else\RMAAlong\fi}}
\newcommand{\@received}{}
\newcommand{\@accepted}{}
\newcommand{\@MSnum}{0000}
\newcommand{\@DOI@prefix}{https://doi.org/10.22201/ia.01851101p}
\def\@DOI@url{\url{https://doi.org/10.22201/ia.01851101p.2099.00.00.00}}
%    \end{macrocode}
% \end{macro}
% \end{macro}
% \end{macro}
% \end{macro}
% \end{macro}
% \end{macro}
% \end{macro}
% \end{macro}
% \end{macro}
% \end{macro}
% \end{macro}
% \end{macro}
% \end{macro}
% \end{macro}
% 
% The following commands are most useful in the \texttt{book}
% style. However, the |\SetConfTitle| command can now be used in
% individual articles to set the header for the \texttt{preprint}
% style. 
%    \begin{macrocode}
\newcommand{\@reunionplace}{Lugar de la reuni\'on}
\newcommand{\@reuniondate}{Fecha de la reuni\'on}
\newcommand{\@conftitle}{Conference Title}
\newcommand{\SetReunion}[2]{\renewcommand{\@reunionplace}{#1}\renewcommand{\@reuniondate}{#2}}
\newcommand{\SetConfTitle}[1]{\renewcommand{\@conftitle}{#1}}
\newcommand{\@editors}{Juan Ram\'on de la Fuente and ``El Mosh''}
\newcommand{\SetEditors}[1]{\renewcommand{\@editors}{#1}}
\newcommand{\@editores}{Juan Ram\'on de la Fuente y ``El Mosh''}
\newcommand{\SetEditores}[1]{\renewcommand{\@editores}{#1}}
\newcommand{\@mes}{mes de la publicaci\'on}
\newcommand{\SetMes}[1]{\renewcommand{\@mes}{\rmaa@VolumenStyle{#1}}}
%    \end{macrocode}

% \subsection{Notes for Author Proofs and Communication between Editors}
% \label{sec:notes}
% We introduce a command \cs{Note} that should be used to describe
% any change made to the manuscript. These are a bit tricky to get
% right. In general we use \cs{marginpar} but this only works in outer
% paragraph mode, so we just bung the box where we are if we are in a
% figure caption or something. If we are in the preamble, then we just
% use a \marginpar as part of the \maketitle command.
%    \begin{macrocode}
\newif\if@arenotes
\@arenotesfalse
\newcounter{notecounter}
\setcounter{notecounter}{0}
\newcommand\NotesList{}
\newcommand\Note[1]{%
  \if@proof{%
    \@arenotestrue%
    \stepcounter{notecounter}%
    \if@inpreamble\else{%
      \NoteMark{\arabic{notecounter}}%
    }\fi
    \protected@xdef\NotesList{\NotesList 
      \protect\item #1}%
  }\fi%
}
\newcommand\Changed[2]{Changed ``#1'' to ``#2''.}
\newcommand\NoteMark[1]{%
  \ifinner
  \raisebox{1ex}[0pt][0pt]{\makebox[0pt][l]{\framebox{\footnotesize\bfseries
        N#1}}}%
  \else
  \marginpar{\framebox{\footnotesize\bfseries N#1}}%
  \fi
}
\newcommand\RefNote[1]{\if@proof\NoteMark{\ref{#1}}\relax\fi}
\newcommand\NoteLabel{\protect\label}
\newcommand\OutputNotes{%
  \bgroup
  \onecolumn
  \section*{EDITORIAL NOTES}
  \begin{enumerate}
    \NotesList
  \end{enumerate}
  \egroup
}
\DeclareRobustCommand\cs[1]{\texttt{\char`\\#1}}

%    \end{macrocode}
% \subsection{Bibliographic Records for ADS}
% \label{sec:ADS}
% \changes{2.2beta11}{2001/11/27}{Just about finished all the ADS
% stuff now. Except for the references, which I have given up on for
% the time being.}
% \begin{macro}{\ADS@BibCode}
% \begin{macro}{\ADS@year}
% \begin{macro}{\ADS@journal}
% \begin{macro}{\ADS@volume}
% \begin{macro}{\ADS@firstpage}
% \begin{macro}{\ADS@authorinitial}
% We use the |\ADS@BibCode| command to hold the bibliographic code for
% ADS. It is 
% used in the |%R| 
% field of the bibliographic record for each article. 
% The format is 
% \begin{quote}
% \texttt{YYYYJJJJJVVVVMPPPPA}  
% \end{quote}
% where \texttt{YYYY} is the year, \texttt{JJJJJ} is the journal code,
% which is \texttt{RMxAA} or \texttt{RMxAC}, \texttt{VVVV} is the
% volume number (right-justified), \texttt{M} is not used by us,
% \texttt{PPPP} is the first page number (right-justified), and
% \texttt{A} is the initial of the first author's surname. The volume
% and page fields are padded with periods (`\texttt{.}'), which causes
% us some fun (this is done with the macro |\PadToFourWithDots| described
% below) . 
% 
% All these individual bits are held in the commands |\ADS@year|,
% |\ADS@journal|, |\ADS@volume|, |\ADS@firstpage|,
% |\ADS@authorinitial|. 
%    \begin{macrocode}
\newcommand{\SetADSBibCode}[1]{\renewcommand{\ADS@BibCode}{#1}}
\newcommand{\ADS@BibCode}
{\ADS@year\ADS@journal\ADS@volume.\PadToFourWithDots{\ADS@firstpage}\ADS@authorinitial}
\newcommand{\ADS@year}{2001} 
\newcommand{\ADS@journal}{\if@proceedings RMxAC\else RMxAA\fi}
\newcommand{\ADS@volume}{....} 
%    \end{macrocode}
% These two are just set to default values here, |\ADS@firstpage| is set
% in the |document| environment for included articles, in the
% |frontsection| environment, |\tableofcontents| command,  |index| and
% |abstracts| environments. Also by |\AtBeginDocument| for stand-alone
% articles. 
% 
% God knows how we are going to do |\ADS@authorinitial|!
%    \begin{macrocode}
\newcommand{\ADS@firstpage}{....}
\newcommand{\ADS@authorinitial}{A}
%    \end{macrocode}
% \end{macro}
% \end{macro}
% \end{macro}
% \end{macro}
% \end{macro}
% \end{macro}
% We also have a whole slew of other fields that can be set. These are
% all listed in Table~\ref{tab:ADS}. Most of them still have to be
% defined. 
% \begin{table}[!ht]\raggedleft
% \caption{ADS Bibliographic Records}\label{tab:ADS}
% \begin{tabular}{lp{8em}lp{8em}}\hline
% Code & Description & Macro Name & Comment \\ \hline
% |%R| & Bibliographic Code & |\ADS@BibCode| & required \\
% |%A| & Author list & |\ADS@authors| & required, generate from |\indexauthor| \\
% |%F| & Author affiliation & |\ADS@affil| & format: AA(\dots)
% AB(\dots) \dots \\
% |%J| & Journal Name & |\ADS@journalname| &  \\
% |%D| & Publication Date & |\ADS@fulldate| &  \\
% |%L| & Last Page & |\ADS@lastpage| & generate from label \\
% |%T| & Title & |\ADS@title| & required, generate from |\title| \\
% |%C| & Abstract Copyright & |\ADS@copyright| & set to IA-UNAM \\
% |%U| & URL for electronic document & |\ADS@url| & set to PDF file \\
% |%d| & DOI for article & |\ADS@doi| & set by |\DOI| \\
% |%K| & Keywords & |\ADS@keywords| & set equal to |\@keywords| \\
% |%A| & Abstract Text & |\ADS@abstract| & set equal to |\@abstract| \\
% |%R| & References & |\ADS@references| & generate from the
% bibliography  \\ \hline
% \end{tabular}
% \end{table}
% A lot of these macros are set by other commands, as mentioned in the
% Table, but some need to be set  explicitly here. 
% 
% \begin{macro}{\ADS@doi}
% \changes{4.0}{2019/11/08}{Save DOI for ADS}
% The user/editor command \cs{DOI} is used to set this. 
%    \begin{macrocode}
\newcommand{\ADS@doi}{https://doi.org/10.22201/ia.01851101p.0000.00.00.00}
%    \end{macrocode}
% \end{macro}
% 
% \begin{macro}{\ADS@journalname}
% This
% does not contain any accents at the moment so as not to cause
% problems---this maybe should be changed though. 
%    \begin{macrocode}
\newcommand{\ADS@journalname}{%
  \if@proceedings\@conftitle\space (Eds. \@editors) \fi%
  Revista Mexicana de Astronom\'\i{}a y Astrof\'\i{}sica%
  \if@proceedings\space (Serie de Conferencias)\fi%
  \space Vol. \@volume, pp. \ADS@firstpage--\ADS@lastpage\space (\@year)%
  \space (http://www.astroscu.unam.mx/\~{}rmaa/)
  }
%    \end{macrocode}
% \end{macro}
% \begin{macro}{\ADS@keywords}
% \begin{macro}{\ADS@abstract}
% \begin{macro}{\ADS@title}
% |\ADS@abstract| is straightforward, since we already have the text
% saved as |\@abstract|. We could do the same with |\ADS@keywords| but
% to avoid unwanted formatting we explicitly set it in |\keywords|
% (p.~\pageref{keywords}). |\ADS@title| is just set to |\@rawtitle|,
% which we were already using for another purpose. 
%    \begin{macrocode}
\let\ADS@title=\@rawtitle
\newcommand{\ADS@keywords}{No keywords specified}
\newcommand{\ADS@abstract}{No abstract specified}
%    \end{macrocode}
% \end{macro}\end{macro}\end{macro}
% \begin{macro}{\ADS@authors}
% |\ADS@authors| wants to be in ``Other, A. N.'' format, but with
% semi-colons between authors. The most natural way to do this seems to
% be with |\indexauthor| (p.~\pageref{indexauthor}). We append each
% author to this macro, which means we must explicitly clear it for
% each new article (done in |\article|). .   
%    \begin{macrocode}
\newcommand{\ADS@authors}{}
%    \end{macrocode}
% \end{macro}
% \begin{macro}{\ADS@fulldate}
% \begin{macro}{\ADS@lastpage}
% \begin{macro}{\ADS@copyright}
% \begin{macro}{\ADS@website}
% \begin{macro}{\ADS@url}
% \begin{macro}{\SetWebsite}
% \begin{macro}{\SetMonthNumeric}
% Other sundry data, together with user commands for bits that can be
% set by the Editors. The |\ADS@url| is constructed from the website
% (currently |http://www.astrosmo.unam.mx/rmaa/|,
% the journal descriptor (e.g., |RMxAC..10/|), followed by the
% path to the PDF file (e.g., |PDF/RMxAC..10_HEW1.pdf|). 
% 
%    \begin{macrocode}
\newcommand{\ADS@website}{http://www.astrosmo.unam.mx/rmaa/}
\newcommand{\ADS@url}
{\ADS@website\ADS@journal\ADS@volume\ADS@number/PDF/\ADS@journal\ADS@volume\ADS@number\_\articlename.pdf}
\newcommand{\ADS@number}{}
\newcommand{\ADS@fulldate}{00/\ADS@year} 
\newcommand{\ADS@lastpage}{00} 
\newcommand{\ADS@copyright}{(C) \ADS@year; Instituto de Astronom\'\i{}a,
  Universidad Nacional Aut\'onoma de M\'exico} 
\newcommand{\SetNumber}[1]{\renewcommand{\ADS@number}{-#1}}
\newcommand{\SetWebsite}[1]{\renewcommand{\ADS@website}{#1}}
\newcommand{\SetMonthNumeric}[1]{\renewcommand{\ADS@fulldate}{#1/\ADS@year}}
%    \end{macrocode}
% \end{macro}\end{macro}\end{macro}\end{macro}\end{macro}\end{macro}\end{macro}
% \begin{macro}{\ADS@affil}
% \begin{macro}{\ADS@references}
% These two are somewhat more problematical. Affiliation currently
% only works for the first author (see \dots) pending redesign of the
% user interface. References are now done by wrapping the
% |thebibliography| environment to open another input stream on the
% same file and search through for the |\bibitems|. 
%    \begin{macrocode}
\newcommand{\ADS@references}{}
\newcommand{\ADS@affil}{}
%    \end{macrocode}
% \end{macro}\end{macro}
% 
% Define two macros that expand to the strings
% |\begin{thebibliography}| and |\end{thebibliography}|. Note that the
% idea for this and much else is inspired by the code for the docstrip
% package. Also define |\StartBibX| in case authors use the more
% standard |\begin{thebibliography}{}| syntax. 
%    \begin{macrocode}
\edef\EndBib{\expandafter\@gobble\string\\end\string{thebibliography\string}}
\edef\StartBib{\expandafter\@gobble\string\\begin\string{thebibliography\string}}
\edef\StartBibX{\expandafter\@gobble\string\\begin\string{thebibliography\string}\string{\string}}
%    \end{macrocode}
% Define two new logical variables. |\ifAnotherline| is true if we
% want to read another line from the input stream. It is set to false
% when we encounter the |\end{thebibliography}|. |\ifInBib| is set to
% true when we frind the |\begin{thebibliography}| and tells us to
% start saving the following lines. 
%    \begin{macrocode}
\newif\ifAnotherLine
\newif\ifInBib
%    \end{macrocode}
% \begin{macro}{makeOther}
% We need a macro to change the catcodes of lots of things to other. 
%    \begin{macrocode}
\def\makeOther#1{\catcode`#1=12\relax}
%    \end{macrocode}
% \end{macro}
% \begin{macro}{\RereadBiblio}
% Define a macro to do all the fiddly stuff. This is used at the
% beginning of |thebibliography| environment (Sec~\ref{sec:biblio}).
%    \begin{macrocode}
\newread\tf@biblio
\newcommand{\RereadBiblio}{%
%    \end{macrocode}
% Open the current file again for reading
%    \begin{macrocode}
  \begingroup
  \immediate\openin\tf@biblio=\articlename
%    \end{macrocode}
% Everything in a group to be safe. Initialize logicals. Black magic
% to change
% catcodes and avoid blank lines being interpreted as |\par|. 
%    \begin{macrocode}
  \AnotherLinetrue\InBibfalse
  \makeOther\ \makeOther\\\makeOther\$%
  \makeOther\#\makeOther\^\makeOther\^^K%
  \makeOther\_\makeOther\^^A\makeOther\%%
  \makeOther\~\makeOther\{\makeOther\}\makeOther\&%
  \endlinechar-1\relax
%    \end{macrocode}
% Start looping through the file. Read each line to the macro
% |\Line|. If we have got to EOF or end of biblio then set
% |\AnotherLinefalse|.
%    \begin{macrocode}
  \loop
  \read\tf@biblio to \Line
  \ifx\Line\EndBib\AnotherLinefalse\InBibfalse\fi
  \ifeof\tf@biblio \AnotherLinefalse\fi
%    \end{macrocode}
% If we are in the biblio and the line is not blank then append it to
% the |\ADS@references| macro together with a newline. 
%    \begin{macrocode}
  \ifInBib
  \ifx\Line\empty\else\xdef\ADS@references{\ADS@references\Line^^J}\fi
  \fi
%    \end{macrocode}
% Check if the line we just saw was the start of biblio, then repeat
% the loop if necessary. 
%    \begin{macrocode}
  \ifx\Line\StartBib\InBibtrue\fi
  \ifx\Line\StartBibX\InBibtrue\fi
  \ifAnotherLine\repeat
  \immediate\closein\tf@biblio
  \endgroup
  }
%    \end{macrocode}
% \end{macro}
% 
% \begin{macro}{\PadToFourWithDots}
% Utility function for left-padding a number to 4 characters with
% periods. Used in 
% constructing |\ADS@BibCode|. This is a very simple algorithm and
% \emph{requires} that the argument actually be a number---if it is
% not then dire consequences will result.  
%    \begin{macrocode}
\newcommand{\PadToFourWithDots}[1]{%
  \ifnum #1 > 999 #1%
  \else\ifnum #1 > 99 .#1%
  \else\ifnum #1 > 9 ..#1%
  \else ...#1%
  \fi\fi\fi}
%%  \ifnum \count0 > 999 #1%
%%  \else\ifnum \count0 > 99 .#1%
%%  \else\ifnum \count0 > 9 ..#1%
%%  \else ...#1%
%%  \fi\fi\fi}
%    \end{macrocode}
% \end{macro}
%
% \begin{macro}{\AuthorADS}
% \begin{macro}{\ADS@affiliations}
% New mechanism for setting the affiliations for ADS. The author labels (AA, AB, \dots) are generated using two nested counters.  We use the |\robustify| command from \texttt{etoolbox} so that accented characters in the affiliation addresses do not blow up.  The list of accents may need to be expanded. 
%    \begin{macrocode}
\newcounter{AUTHOR}
\newcounter{author}[AUTHOR]
\stepcounter{AUTHOR}
\newcommand\ADS@affiliations{}
\robustify{\'}\robustify{\i}\robustify{\~}\robustify{\`}
\def\AuthorADS#1#2{%
  \begingroup
  \def\ORCIDiD##1{<ID system="ORCID">##1</ID>}
  \def\Email##1{<EMAIL>##1</EMAIL>}
  \def\Affil##1{##1}
  \stepcounter{author}
  \ifthenelse{\value{author} > 26}{\stepcounter{AUTHOR}\stepcounter{author}}{}
  \ifthenelse{\equal{\ADS@affiliations}{\@rmaaEmpty}}{%
    \xdef\ADS@affiliations{\Alph{AUTHOR}\Alph{author}(#2)}
  }{%
    \xdef\ADS@affiliations{\ADS@affiliations, \Alph{AUTHOR}\Alph{author}(#2)}
  }%
  \endgroup
}
%    \end{macrocode}
% \end{macro}
% \end{macro}
% %
% \begin{macro}{\ADS@dumpdata}
% \changes{4.0}{2019/11/13}{Write all author affiliations for ADS}
% \begin{macro}{\ADS@write}
% \begin{macro}{\ADS@openfile}
% |\ADS@dumpdata| is the macro that actually does the job of writing out all the
% data to an auxiliary file, using |\ADS@write| to write each
% field. The file is set up with |\ADS@openfile|, which currently is
% done in the TOC of a book.  The new multi-author affiliations are used if |\ADS@affiliations| has been set, otherwise it falls back to the single-author version |\ADS@affil|. 
%    \begin{macrocode}
\newcommand{\ADSfield}[2]{\relax}
\newcommand{\ADSvfield}[2]{\relax}
%%\renewcommand{\ADSfield}[12]{\%#1 #2}
\newcommand{\ADS@dumpdata}
{ 
  \ADS@write{\string\begin{ADSarticle}}
  \ADS@WRITE{\string\ADSfield{R}{\ADS@BibCode}}
  \ADS@write{\string\ADSfield{A}{\ADS@authors}}
  \ifthenelse{\equal{\ADS@affiliations}{\@rmaaEmpty}}{%
    {\ADS@sanitize\ADS@write{\string\ADSfield{F}{AA(\ADS@affil)}}}
  }{%
    {\ADS@sanitize\ADS@write{\string\ADSfield{F}{\ADS@affiliations}}}
  }
  \ADS@write{\string\ADSfield{J}{\ADS@journalname}}
  \ADS@write{\string\ADSfield{D}{\ADS@fulldate}}
  \ADS@write{\string\ADSfield{L}{\ADS@lastpage}}
  {\ADS@sanitize\ADS@write{\string\ADSfield{T}{\ADS@title}}}
  \ADS@write{\string\ADSfield{C}{\ADS@copyright}}
  \ADS@write{\string\ADSfield{d}{\ADS@doi}}
  \ADS@write{\string\ADSfield{U}{\ADS@url}}
  \ADS@write{\string\ADSfield{K}{\ADS@keywords}}
  {\ADS@sanitize\ADS@write{\string\ADSfield{B}{\ADS@abstract}}}
  \ADS@write{\string\ADSfield{Z}{\ADS@references}}
  \ADS@write{\string\end{ADSarticle}}
  }
%%\newcommand{\ADS@write}[1]{\if@book\addtocontents{ads}{#1}\fi}
%%\newcommand{\ADS@openfile}{\@starttoc{ads}}
\newcommand{\ADS@write}[1]{\set@display@protect\immediate\write\tf@ads{#1}}
\newcommand{\ADS@WRITE}[1]{\immediate\write\tf@ads{#1}}
\newcommand{\ADS@openfile}{%
  \newwrite\tf@ads%
  \immediate\openout\tf@ads\jobname.ads\relax
  \ADS@WRITE{\ADS@preamble}}
\newcommand{\ADS@closefile}{\ADS@WRITE{\ADS@postamble}}
\def\ADS@preamble{%
  %%  \string\documentclass{article}
  %%  \string\begin{document}
  }
\def\ADS@postamble{%
  %%  \string\end{document}
  \string%\string% Local Variables:^J 
  \string%\string% mode: LaTeX^J
  \string%\string% mode: Font-lock^J
  \string%\string% End:^J
  }
%    \end{macrocode}
% \end{macro}
% \end{macro}
% \end{macro}
% 
% \subsection{Automatic calculation of the page range}
% \label{sec:page-range}
% Another issue that has cost me a lot of grief. Still, it finally
% seems to be sorted now. All hail Jeff Goldberg. 
%
% Cunning trick to get the page numbers from |\pageref|.  Put a label
% on the first and last page of the document using |\AtBeginDocument|
% and |\AtEndDocument|. .  This way won't work for the book option.
% In that case, we do it by defining |\@endofarticlehook| and
% |\@begofarticlehook| as parallels to |\AtEndDocument| and
% |\AtBeginDocument|.  Note that we have to defer adding
% |\@begofarticlehook| to |\AtBeginDocument| to make sure it comes
% after the redefinitions of |\ref| and |\label|.
%    \begin{macrocode}
\if@book\AtEndDocument{\ADS@closefile}\else
 %\AtBeginDocument{\@begofarticlehook}
\AtEndDocument{\@endofarticlehook}
\fi
%    \end{macrocode}
% \begin{macro}{\@begofarticlehook}
% Since we are getting more and more things that need to be done at
% the beginning of every article, frontsection, etc we define a sister
% hook to 
% |\@endofarticlehook|. This hook does the following: 
% \begin{enumerate}
% \item Labels the first page
% \item Sets |\ADS@firstpage|
% \item runs and then clears |\indexauthor@deferlist|
% \item resets the keyword counter
% \end{enumerate}
% It is executed in the following circumstances: 
% \begin{enumerate}
% \item At |\begin{document}| for standalone files
% \item At |\begin{document}| for each included file
% \item At beginning of front sections 
% \item At beginning of TOC
% \item At beginning of Author Index
% \item At beginning of abstracts
% \end{enumerate}
% For some reason, this is not always within the right group, so we
% use |\xdef| just to be sure. 
%    \begin{macrocode}
\newcommand{\@begofarticlehook}{%
  \if@frontmatter\frontlabel{RMAAFirstPage}\else\label{RMAAFirstPage}\fi%
  \xdef\ADS@firstpage{\@arabic{\c@page}}%
  \if@book\else\makeindex\fi % needed so default index is defined in stand-alone articles
%%  \indexauthor@deferlist
%%  \xdef\indexauthor@deferlist{}
  \@inpreamblefalse
  \setcounter{kwdcounter}{0}
  }
%    \end{macrocode}
% \end{macro}
% \begin{macro}{\@endofarticlehook}
% This is a hook to be run at the end of each article
% currently, it just labels the last page and writes out all the ADS
% data to an auxiliary file.  
% It should not produce any printed output. 
% It is modelled on the |\AtEndDocument| command, to which it reduces
% in the case of a single article. 
%    \begin{macrocode}
\newcommand{\@endofarticlehook}{%
  \if@debug{\ClassWarning{rmaa}{***Doing Hook***}}\fi
  \if@debug{\ClassWarning{rmaa}{ADSabstract is: \ADS@abstract}}\fi
  \clearpage\rmaa@lastpagelabel\if@book\ADS@dumpdata\fi
  \ifx\NotesList\@rmaaEmpty\else\OutputNotes\fi
  \if@debug{\ClassWarning{rmaa}{Notes list is: \NotesList}}\fi
}
%    \end{macrocode}
% \begin{macro}{\rmaa@lastpagelabel}
% New version of last page label command copied from Jeff Goldberg's
% lastpage.sty package - we can't just use |\label| here in case the
% output routine for the last page has already been called. Instead,
% we first explicitly clear the page, step back the page counter,
% explicitly write to the .aux file, then restep the page counter
% again.
%
% But this will conflict with the hyperref package, which does labels differently.  Compensate for that by using the five-argument form.
% 
% Also set the |\ADS@LastPage| macro to the appropriate number. 
%    \begin{macrocode}
\newcommand{\rmaa@lastpagelabel}{%
  \addtocounter{page}{-1}%
  \immediate\write\@auxout{\string
    \newlabel{\articlename:RMAALastPage}{{}{\thepage}{}{}{}}}%
  \edef\ADS@lastpage{\thepage}%
  \addtocounter{page}{1}}
%    \end{macrocode}
% \end{macro}
% \end{macro}

% \subsection{Deal with explicit change in number of columns}
% \label{sec:columns}
% Sometimes, when formatting a book, we need to switch between
% one-column articles and two-column articles. For this to work, we
% need to wrap the standard \LaTeX{} commands to set the |\@rmtwocol|
% flag appropriately and also reset the package default values of all
% scalable lengths. .  
%    \begin{macrocode}
\let\latexonecolumn=\onecolumn
\let\latextwocolumn=\twocolumn
\renewcommand{\onecolumn}{\@rmtwocolfalse\@package@default@lengths\latexonecolumn}
\renewcommand{\twocolumn}{\@rmtwocoltrue\@package@default@lengths\latextwocolumn}
%    \end{macrocode}
% 
% \subsection{The actual formatting of the title}
% \label{sec:title-format}

% Note, we have to redefine the |\maketitle| command to use the
% numbered footnotes. This means we needn't use the atbeginend stuff. 
%    \begin{macrocode}
\renewcommand{\maketitle}{\par
  \begingroup
%    \end{macrocode}
% Don't switch to marked footnotes. 
% This is consistent with existing style but I'm not sure I like it
%    \begin{macrocode}
  %%  \renewcommand\thefootnote{\fnsymbol{footnote}}%
%%  \def\@makefnmark{\hbox to\z@{$\m@th^{\@thefnmark}$\hss}}%
  \def\@makefnmark{\textsuperscript{\@thefnmark}\hspace{-0.2em}}%
  \long\def\@makefntext##1{\parindent 1em\noindent
    \hbox to1.8em{\hss$\m@th^{\@thefnmark}$}##1}%
  \if@rmtwocol
  { \if@debug{\ClassWarning{rmaa}{TWO COLUMN}}\fi
    \twocolumn[\@maketitle]%
    }
  \else
  { \if@debug{\ClassWarning{rmaa}{ONE COLUMN}}\fi
    \newpage
    \global\@topnum\z@   % Prevents figures from going at top of page.
    \if@manuscript \@MSmaketitle \else \@maketitle \fi
    }
  \fi
  \if@manuscript\else\thispagestyle{rmaafirstpage}\@thanks\fi
  \ifx\NotesList\@rmaaEmpty\else{%
    \NoteMark{1--\arabic{notecounter}}\vspace*{-2\baselineskip}%
  }\fi%
  \endgroup
%    \end{macrocode}
% Don't reset footnote counter (not sure whether this is good thing
% or not).
%    \begin{macrocode}
  %%  \setcounter{footnote}{0}%
  \let\thanks\relax\let\affil\relax\let\altaffil\relax
  \let\maketitle\relax\let\@maketitle\relax
  \gdef\@thanks{}\gdef\@author{}\gdef\@title{}}
%    \end{macrocode}

% \begin{macro}{\@kwhwidth}
% \begin{macro}{\@kwtwidth}
% Redefine \cs{@maketitle} to include the abstract, resumen and
% keywords. Note that in the manuscripty style we use
% \cs{@MSmaketitle} instead.  
% \changes{3.16}{2003/04/07}{
%  Added extra vertical space between authors and accepted date.}
% \changes{3.26}{2011/08/15}{Add extra label to review articles}
%    \begin{macrocode}
\newlength{\@kwhwidth}\newlength{\@kwtwidth}
\newcommand\reviewname{Review}
\renewcommand{\@maketitle}{%
  \newpage
  \null
  \if@reviewarticle
  \smallskip
  {\large\bfseries\itshape\reviewname}%
  \bigskip
  \else
  \vskip\@before@title@skip%
  \fi
  \begin{center}%
    {\@title \par}%
    \vskip\@before@author@skip%
    \parbox{\rmaa@AuthorWidthFrac\textwidth}{\centering
      \lineskip\@between@author@skip%
      \@author
      \par}
    %% this should be replaced with `Submitted to...' for preprints
    \if@printdates
    \vskip\@before@dates@skip%
    \if@preprint
    \rmaa@ReceivedStyle{Submitted to RMxAA \@received}
    \else
    \rmaa@ReceivedStyle{Received \@received; accepted \@accepted}
    \fi
    \fi
    \if@manuscript
    \rmaa@ReceivedStyle{Submitted to RMxAA \@received}
    \fi
    \if@guide
    \rmaa@ReceivedStyle{Version \rmaa@version, \rmaa@date}
    \fi
    \vskip\@before@abstract@skip%
    \if@abstractpresent
    \parbox{\rmaa@AbstractWidthFrac\textwidth}{%
      \setlength{\parindent}{0.9cm}%
      \centerline{\rmaa@AbstractHeaderStyle{Abstract}}\par%
      \vskip\@after@abstracthead@skip%
      \if@newstyle\noindent\fi\@abstract
      \vskip\@between@abstract@skip%
      \centerline{\rmaa@AbstractHeaderStyle{Resumen}}\par%
      \vskip\@after@abstracthead@skip%
      \if@newstyle\noindent\fi\@resumen
      \if@setresumen\else\@abstract\fi % append english abstract if
                                       % spanish one not set
      \if@guide\else            %do not use keywords for guide
      \vskip\@before@kwd@skip%
      % Format the Key Words with hanging indentation:
      % They are put in a \parbox with width equal to the line width
      % minus the width of the header. 
      { % Widths of keyword
        \settowidth{\@kwhwidth}{\@keywordheader~}   % header and
        \setlength{\@kwtwidth}{\linewidth}          % keyword text
        \addtolength{\@kwtwidth}{-\@kwhwidth}       %
        \noindent\@keywordheader~\parbox[t]{\@kwtwidth}{%
          \rmaa@KeywordStyle{\boldmath\uppercase{\@keywords}}}}%
      \fi
      }
      \fi
  \end{center}%
  \vspace*{\afterkwdskip}
  \par
  }
\newlength{\afterkwdskip}\setlength{\afterkwdskip}{\z@}
%    \end{macrocode}
% \end{macro}
% \end{macro}
%
%
% \begin{macro}{\reviewarticle}
% \begin{macro}{\reviewtoc}
% \begin{macro}{\reviewlof}
%   User command to specify a review article. Must be used in
%   preamble. Main effect is to place the text ``Review'' before the
%   article title.  Also add command for setting a table of contents
%   and a list of figures at the beginning of the article.
%    \begin{macrocode}
\newcommand{\reviewarticle}{\@reviewarticletrue}
\newcommand\reviewtoc{%
  \section*{CONTENTS}%
  \@starttoc{rtc}%
}
\newcommand\reviewlof{%
  \section*{LIST OF FIGURES}%
  \@starttoc{lof}%
}
%    \end{macrocode}
% \end{macro}
% \end{macro}
% \end{macro}
% \begin{macro}{\l@section}
%    \begin{macrocode}
\renewcommand*\l@section[2]{%
  \ifnum \c@tocdepth >\z@
    \addpenalty\@secpenalty
    \addvspace{0.8em \@plus\p@}%
    \setlength\@tempdima{1.5em}%
    \begingroup
      \parindent \z@ \rightskip \@pnumwidth
      \parfillskip -\@pnumwidth
      \leavevmode %\bfseries
      \advance\leftskip\@tempdima
      \hskip -\leftskip
      #1\nobreak\hfil \nobreak\hb@xt@\@pnumwidth{\hss #2}\par
    \endgroup
  \fi}
%    \end{macrocode}
% \end{macro}
%
% \begin{macro}{\@MSmaketitle}
% A special version of |\@maketitle| for the \texttt{manuscript}
% option. This can deal with a multi-page abstract/resumen. Not very
% well, mind. 
%    \begin{macrocode}
\newcommand{\@MSmaketitle}{%
  \if@debug\ClassInfo{rmaa}{Using special MS maketitle command}\fi
  \newpage
  {
    \thispagestyle{rmaafirstpage}
%     \setlength{\parindent}{0pt}
    \null\vskip 2ex
    \begin{center}
      \@title\par
      \@author\par
      \rmaa@ReceivedStyle{%
        \ifx\@received\@rmaaEmpty
          Draft version: \@date
        \else
          Received: \@received
          \ifx\@accepted\@rmaaEmpty
          \else
            ; Accepted: \@accepted
          \fi
        \fi}\par
    \end{center}
    \@thanks
    \centerline{\rmaa@AbstractHeaderStyle{Resumen}}
    \noindent\@resumen\par\vskip 2ex
    \centerline{\rmaa@AbstractHeaderStyle{Abstract}}
    \noindent\@abstract\par\vskip 2ex
    \noindent\@keywordheader~\rmaa@KeywordStyle{\@keywords}\par
%     \@keywordheader~\rmaa@KeywordStyle{\uppercase{\@keywords}}\par
  }%
  }
%    \end{macrocode}
% \end{macro}

% \subsection{The sectioning commands}
% \label{sec:section-commands}

% \begin{macro}{\section}
% \begin{macro}{\subsection}
% \begin{macro}{\subsubsection}
% \begin{macro}{\paragraph}
% \begin{macro}{\subparagraph}
% Redefine all the standard commands to use these lengths. 
%    \begin{macrocode}
\renewcommand\section{\@startsection {section}{1}{\z@}%
                                   {\abovesectionskip}%
                                   {\belowsectionskip}%
                                   %% how do we get upper case here
                                   {\centering\reset@font\normalsize}}
\renewcommand\subsection{\@startsection{subsection}{2}{\z@}%
                                     {\abovesectionskip}%
                                     {\belowsectionskip}%
                                     {\centering\reset@font\normalsize\itshape}}
\renewcommand\subsubsection{\@startsection{subsubsection}{3}{\z@}%
                                     {\abovesectionskip}%
                                     {\belowsectionskip}%
                                     {\centering\reset@font\normalsize\itshape}}
\renewcommand\paragraph{\@startsection{paragraph}{4}{\z@}%
                                    {\aboveparaskip}%
                                    {\belowparaskip}%
                                    {\reset@font\normalsize\bfseries}}
\renewcommand\subparagraph{\@startsection{subparagraph}{5}{\parindent}%
                                      {\aboveparaskip}%
                                      {\belowparaskip}%
                                      {\reset@font\normalsize\bfseries}}
%    \end{macrocode}
% \end{macro}
% \end{macro}
% \end{macro}
% \end{macro}
% \end{macro}

% \begin{macro}{secnumdepth}
% Make sure that we number the subsubsections---this is done by
% default in the |article| class but not the |book| class, which is
% why we must set it explicitly here.
%    \begin{macrocode}
\setcounter{secnumdepth}{3}
%    \end{macrocode}
% \end{macro}
% \changes{2.2beta10}{2001/11/07}{Fixed \cs{secnumdepth} for \cs{book}
% style.}

% \begin{macro}{\thesubsection}
% \begin{macro}{\thesubsubsection}
% For subsection and subsubsection, we want the numbers to be roman
% but the title to be italic. This seems the easiest way to go about it.
%    \begin{macrocode}
\renewcommand{\thesubsection}{{\upshape\thesection.\@arabic\c@subsection}}
\renewcommand{\thesubsubsection}{\thesubsection{\upshape.\@arabic\c@subsubsection}}
%    \end{macrocode}
% \end{macro}
% \end{macro}

% \begin{macro}{\@seccntformat}
% Add ``.'' to the number here so we can remove it from |\thesection|,
% etc. Reduce gap between number and title WJH 27 Jan 1995
%    \begin{macrocode}
\def\@seccntformat#1{\csname the#1\endcsname .\hskip 0.4em}
%    \end{macrocode}
% \end{macro}

% \begin{macro}{\@sect}
% I didn't really want to mess with the |\@sect| command, but it was
% the only way I could manage to force upper case. Other than the
% commented parts, it is identical to the standard version. 
%    \begin{macrocode}
\def\@sect#1#2#3#4#5#6[#7]#8{
  \ifnum #2>\c@secnumdepth
     \let\@svsec\@empty
  \else
     \refstepcounter{#1}%
     \let\@@protect\protect
     \def\protect{\noexpand\protect\noexpand}%
     \edef\@svsec{\@seccntformat{#1}}%
     \let\protect\@@protect
  \fi
%    \end{macrocode}
% This is the bit I changed (27 Jan 1995). Use upper case for
% |\section|, use mixed case for the rest. 
%    \begin{macrocode}
  \ifnum #2=1                   
     \def\sec@upcase##1{\uppercase{##1}}%
     \def\sec@shape{\upshape}%
  \else                         
     \def\sec@upcase##1{\relax{##1}}%
     \def\sec@shape{\itshape}%
  \fi
  \@tempskipa #5\relax
  \ifdim \@tempskipa>\z@
     \begingroup #6\relax
          \@hangfrom{\hskip #3\relax\@svsec}%
                    {\interlinepenalty \@M \sec@upcase{#8}\par}%
     \endgroup
     \csname #1mark\endcsname{#7}%
%    \end{macrocode}
% Only write TOC entries if doing the Guide style. 
%    \begin{macrocode}
     \if@guide 
        \addcontentsline{toc}{#1}{%
          \ifnum #2>\c@secnumdepth \else
             \protect\numberline{\csname the#1\endcsname}%
          \fi
          #7}%
     \fi
%    \end{macrocode}
% Also allow per-article TOC for review papers
%    \begin{macrocode}
     \if@reviewarticle
        \ifnum #2>2 \else
           \addcontentsline{rtc}{#1}{%
             \protect\numberline{\csname the#1\endcsname}%
             \sec@upcase{\sec@shape #7}}%
        \fi
     \fi
  \else
     \def\@svsechd{#6\hskip #3\relax  %% \relax added 2 May 90
       \@svsec #8\csname #1mark\endcsname{#7}%
       \if@guide
          \addcontentsline{toc}{#1}{%
            \ifnum #2>\c@secnumdepth \else
               \protect\numberline{\csname the#1\endcsname}%
            \fi
            #7}%
       \fi
       }%
  \fi
  \@xsect{#5}}
%    \end{macrocode}
% \end{macro}


% \subsection{Appendices}
% \label{sec:appendices}
% Make sure that appendices in individual articles work properly. We
% copy the definition from the article class (the book one doesn't do
% what we want).  
% \begin{environment}{appendices}
%   This uses an environment |appendices| inside of which each
%   |\section| will become a separate appendix.  Note that the
%   appendices must come before the references (otherwise the full
%   addresses mechanism will not work properly). We make the equation
%   numbering be subordinate to the appendix (equations A12, B6,
%   etc). Finally, we have a medskip followed by APPENDICES to
%   introduce them. \changes{3.17}{2002/04/08}{ Added appendices environment
%   that does the right thing.}
%    \begin{macrocode}
\newenvironment{appendices}{\par
  \begingroup
  \setcounter{section}{0}%
  \setcounter{subsection}{0}%
  \renewcommand\thesection{\@Alph\c@section}
  \renewcommand\theequation{\thesection\arabic{equation}}
  \medskip
  \section*{APPENDICES}}
{\endgroup}
%    \end{macrocode}
% \end{environment}
% \begin{environment}{appendix}
%   If there is only one appendix, then use |appendix| environment
%   instead.  \changes{4.0}{2019/11/12}{ Added
%   appendix environment for case where there is only one section.}
%    \begin{macrocode}
\renewenvironment{appendix}{\par
  \begingroup
  \setcounter{section}{0}%
  \setcounter{subsection}{0}%
  \renewcommand\thesection{\@Alph\c@section}
  \renewcommand\theequation{\thesection\arabic{equation}}
  \medskip
  \section*{APPENDIX}}
{\endgroup}
%    \end{macrocode}
% \end{environment}

% \subsection{Floating thingies}
% \label{sec:floats}   

%% 
%    \begin{macrocode}
\renewcommand{\thetable}{\@arabic\c@table}
\renewcommand{\thefigure}{\@arabic\c@figure}
\renewcommand{\figurename}{Fig.}
\renewcommand{\tablename}{TABLE}
%\renewcommand{\fnum@table}{\MakeUpperCase{\tablename}~\thetable}
%    \end{macrocode}
% For some reason, I have to write `TABLE' explicitly here
% If I use |\tablename|, it gets converted to mixed case---God knows
% why. 
% |\figurename| is also being ignored. 
%    \begin{macrocode}
\renewcommand{\fnum@table}{TABLE~\thetable}
\renewcommand{\fnum@figure}{Fig.~\thefigure}
%    \end{macrocode}

% Do captions in 9 pt face in the main journal style and in new-style
% anything. 
%    \begin{macrocode}
\long\def\@makefigcap#1#2{%
  \if@proceedings\else\small\fi
  \if@newstyle\small\fi 
  \vskip\abovecaptionskip
  \sbox\@tempboxa{#1. #2}%
  \ifdim \wd\@tempboxa >\hsize
    #1. #2 \par
  \else
    \global \@minipagefalse
    \hb@xt@\hsize{\hfil\box\@tempboxa\hfil}%
  \fi
  \vskip\belowcaptionskip}
%    \end{macrocode}

% Make sure table caption is upper case. 
%    \begin{macrocode}
\long\def\@maketabcap#1#2{%
  \hbox to\hsize{\hfil #1 \hfil}
  \vskip 6pt 
  \centering\MakeUpperCase{#2}\par
  \vskip 6pt}
%    \end{macrocode}

%    \begin{macrocode}
\newlength{\@rottabcapwidth}
\long\def\@makerottabcap#1#2{%
  \setbox\@tempboxa\hbox{\MakeUpperCase{#2}}%
  \ifdim\wd\@tempboxa >.8\vsize
     \setlength\@rottabcapwidth{.8\vsize}
  \else
     \setlength\@rottabcapwidth{\wd\@tempboxa}
  \fi
  \rotatebox{90}{%
    \begin{minipage}{\vsize}\centering
      \begin{minipage}{\@rottabcapwidth}\centering
        #1\\ \vskip 6pt \MakeUpperCase{#2}
      \end{minipage}%
    \end{minipage}
    }%
  \hskip 6pt
  }

\newenvironment{Sideways}
    {\begin{sideways}\begin{minipage}{\vsize}\centering}
    {\end{minipage}\end{sideways}}

\newenvironment{Sidewaystable}
{%
  \let\oldcaption=\caption%
  \def\caption{
    \def\@captype{table}\renewcommand{\@makecaption}{\@maketabcap}
    \oldcaption}
  \begin{sidewaystable}\centering}
  {\end{sidewaystable}\let\caption=\oldcaption}
\newenvironment{Sidewaystable*}
{%
  \let\oldcaption=\caption%
  \def\caption{
    \def\@captype{table}\renewcommand{\@makecaption}{\@maketabcap}
    \oldcaption}
  \begin{sidewaystable*}\centering}
  {\end{sidewaystable*}\let\caption=\oldcaption}
%    \end{macrocode}

% Not sure why we need to stick these lines in. 
%    \begin{macrocode}
\renewenvironment{figure}
                 {\renewcommand{\@makecaption}{\@makefigcap}\@float{figure}}
                 {\end@float}
\renewenvironment{figure*}
                 {\renewcommand{\@makecaption}{\@makefigcap}\@dblfloat{figure}}
                 {\end@dblfloat}
\renewenvironment{table}
                 {\renewcommand{\@makecaption}{\@maketabcap}\@float{table}}
                 {\end@float}
\renewenvironment{table*}
                 {\renewcommand{\@makecaption}{\@maketabcap}\@dblfloat{table}}
                 {\end@dblfloat}
%    \end{macrocode}

% Redefine |\@caption| so that we don't bother with the list of
% figures. This saves us having to worry about fragile commands in the
% captions. 
%    \begin{macrocode}
\long\def\@caption#1[#2]#3{%
  \par
%%  \addcontentsline{\csname ext@#1\endcsname}{#1}%
%%    {\protect\numberline{\csname the#1\endcsname}{\ignorespaces #2}}%
  \begingroup
    \@parboxrestore
    \if@minipage
      \@setminipage
    \fi
    \normalsize
    \@makecaption{\csname fnum@#1\endcsname}{\ignorespaces #3}\par
  \endgroup}
%    \end{macrocode}

% Make tables not quite so cramped. Allow |\hline\hline| to give a
% poor man's thick rule. 
%    \begin{macrocode}
\renewcommand{\arraystretch}{1.15}
\setlength{\doublerulesep}{0.4pt}
%    \end{macrocode}

% \begin{macro}{\tablecols}
% \begin{macro}{\@tablecols}
% \begin{macro}{\tabcaptioncontinued}
% |\tablecols| is a user command for setting the number of columns. It
% is used by commands that try and span the entire table (I think
% there must be a better way of acheiving
% this). |\tabcaptioncontinued| is 
% a user command for setting the caption of a continued table on the
% next page. 
%    \begin{macrocode}
\newcommand{\tablecols}[1]{\renewcommand{\@tablecols}{#1}}
\newcommand{\@tablecols}{1}
\newcommand{\tabcaptioncontinued}{%
  \multicolumn{\@tablecols}{c}{\addtocounter{table}{-1}TABLE~\thetable\ (CONTINUED)} \\ \\
  }
%    \end{macrocode}
% \end{macro}
% \end{macro}
% \end{macro}

% \subsubsection{Footnotes in tables}
% \label{sec:table-footnotes}
% \begin{macro}{\tablenotemark}
% \begin{macro}{\tablenote}
% \begin{macro}{\tablenoteskip}
% Old user commands |\tablenotemark{MARK}|,
% |\tablenote{MARK}{TEXT}|---now deprecated. 
%    \begin{macrocode}
\newcommand\tablenotemark[1]{$^\mathrm{#1}$} 
\newcommand\tablenote[2]{\if@proceedings\else\small\fi ${}^\mathrm{#1}$ #2} 
\newcommand\tablenoteskip{\\[0.5\baselineskip]} 
%    \end{macrocode}
% \end{macro}
% \end{macro}
% \end{macro}

% \begin{macro}{\tabnotemark}
% \begin{macro}{\tabnotetext}
% \begin{macro}{\tabnotetextTWO}
% Alternative way of doing footnotes in tables. 
%    \begin{macrocode}
\newcommand{\tabnotemark}[1]{\textsuperscript{\lowercase{#1}}}
%\newcommand{\tabnotetextTWO}[4]{\tabnotetext{#1}{#2}\tabnotetext{#3}{#4}}
\newcommand{\tabnotetextTWO}[4]{%
  \multicolumn{\@tablecols}{l}{%
    \parbox{0.5\tabnotewidth}{\small\textsuperscript{#1}#2}%
    \parbox{0.5\tabnotewidth}{\small\textsuperscript{#3}#4}%
    }%
  \\}
\newcommand{\tabnotetext}[2]{%
  \multicolumn{\@tablecols}{@{}p{\tabnotewidth}@{}}{\small\textsuperscript{#1}#2}\\}
\newlength{\tabnotewidth}
%    \end{macrocode}
% \end{macro}
% \end{macro}
% \end{macro}


% \begin{macro}{changemargin}
%    \begin{macrocode}
% Command to allow margins to be changed on the fly. Useful to allow
% tables and Figures to protrude into the margins. 
%
% \changes{3.24}{2007/09/14}{Added \texttt{changemargin} environment
% (copied from TeX FAQ).}
\newenvironment{changemargin}[2]{% 
  % From TeX FAQ http://www.tex.ac.uk/cgi-bin/texfaq2html?label=chngmargonfly
  \begin{list}{}{%
    \setlength{\topsep}{0pt}%
    \setlength{\leftmargin}{#1}%
    \setlength{\rightmargin}{#2}%
    \setlength{\listparindent}{\parindent}%
    \setlength{\itemindent}{\parindent}%
    \setlength{\parsep}{\parskip}%
  }%
  \item[]}{\end{list}}
%    \end{macrocode}
% \end{macro}

% \subsection{The bibliography and references}
% \label{sec:biblio}
% \begin{macro}{\beforebiblioskip}
% \begin{macro}{\afterbiblioskip}
% I'm not entirely sure why I introduced these two commands, since
% directly putting |\vspace| before and after
% |\begin{thebibliography}| seems to work well enough. 
%    \begin{macrocode}
\newlength{\beforebiblioskip}
\setlength{\beforebiblioskip}{0pt}
\newlength{\afterbiblioskip}
\setlength{\afterbiblioskip}{0pt}
%    \end{macrocode}
% \end{macro}
% \end{macro}
% \begin{macro}{thebibliography}
% The environment itself is done as starred section with extra
% vertical space after heading. This space is currently zero. 
% In the book style we do the cunning trick with rereading the file in
% order to get the references into the |\ADS@references| macro.
% 
% \changes{3.05}{2003/01/02}{Removed the \cs{beforebiblioskip} and
% \cs{afterbiblioskip} commands, since they cause unwanted vertical
% space in some instances.}  
%    \begin{macrocode}
\renewenvironment{thebibliography}{%
  \if@book\RereadBiblio\fi
%%  \vspace*{\beforebiblioskip}
  \section*{REFERENCES} \frenchspacing
%%  \vspace*{\afterbiblioskip}
  \small
%    \end{macrocode} 
% The references themselves are set as a list in order 
% to get the hanging indentation. They are done in a smaller font and
% with french spacing so we needn't worry about the distinction
% between sentence-ending periods and abbreviations. 
%    \begin{macrocode} 
  \list{\null}{\leftmargin 1.5em\labelwidth\z@\labelsep\z@\itemindent -1.5em
    \itemsep\z@\parsep\z@\usecounter{enumi}}
  \sloppy\clubpenalty4000\widowpenalty4000
  \sfcode`\.=1000\relax\@biblisttrue}
{\endlist\@biblistfalse\normalsize\@outputaddressesifnecessary}
%    \end{macrocode}
% \end{macro}
% \subsection{Uniquification of cross-references}
% \label{sec:uniquification}
% 
% \begin{macro}{\articlename}
% In order that the |\label{}| command in individual articles be
% guaranteed to give a unique anchor, we append the article name (held
% in the |\articlename| command) to each key. This requires
% redefinition of the following commands: |\cite|, |\nocite|, |\ref|,
% |\pageref|, |\label|. 
% \begin{macro}{\oldcite}
% \begin{macro}{\oldnocite}
% \begin{macro}{\oldref}
% \begin{macro}{\oldpageref}
% \begin{macro}{\oldlabel}
% The original versions of these commands are
% saved as |\oldcite|, |\oldnocite|, etc. 
%
% The uniquifying of cross-references is now always done, so 
% we don't get multiply-defined labels in the single article case
%    \begin{macrocode}
\newcommand{\articlename}{anon}
\AtBeginDocument{
  \let\oldcite=\cite 
  \let\oldnocite=\nocite 
  \let\oldref=\ref 
  \let\oldpageref=\pageref
  \let\oldlabel=\label
}
%    \end{macrocode}
% \end{macro}
% \end{macro}
% \end{macro}
% \end{macro}
% \end{macro}
% \end{macro}
%
% Uniquify all cross-references by prepending the article name.  This
% now has to go into |\AtBeginDocument| so it is working on the new
% versions that hyperref.sty has set up.  Then we need to put the
% |\@begofarticlehook| afterwards, so that the first page label is using
% the new definitions.  If amsmath is loaded, then we also need to patch
% the special label command that they use. \changes{4.0}{2019/11/18}{Fix labels for amsmath display equations, in case that package has been loaded.}
%    \begin{macrocode}
\AtBeginDocument{
  \renewcommand{\ref}[1]{\oldref{\articlename:#1}%
    \if@debug{\ClassInfo{rmaa}{Reference to \articlename:#1}}\fi}
  \renewcommand{\pageref}[1]{\oldpageref{\articlename:#1}%
    \if@debug{\ClassInfo{rmaa}{Page reference to \articlename:#1}}\fi}
  \renewcommand{\label}[1]{\oldlabel{\articlename:#1}%
    \if@debug{\ClassInfo{rmaa}{Label for \articlename:#1}}\fi}%
  \ifdef{\label@in@display}{%
    \let\oldlabel@in@display=\label@in@display
    \renewcommand{\label@in@display}[1]{\oldlabel@in@display{\articlename:#1}%
      \if@debug{\ClassInfo{rmaa}{AMSTeX Label for \articlename:#1}}\fi}%
  }{}
  \@begofarticlehook
}%
%    \end{macrocode}
% Citations are more complicated since they can contain a list of
% keys. These are really only here for use with the rmaacite package,
% which is utterly deprecated (02 Jan 2003) now that I have the natbib
% code working. The command \cs{cite} should \emph{not} be used with
% natbib: use \cs{citep}, \cs{citet}, etc.\@ instead. 
%    \begin{macrocode}
\renewcommand{\nocite}[1]{\@bsphack
  \@for\@citeb:=#1\do{%
    \edef\@citeb{\expandafter\@firstofone\@citeb:\articlename}%
    \if@filesw\immediate\write\@auxout{\string\citation{\@citeb}}\fi
    \@ifundefined{b@\@citeb:\articlename}{\G@refundefinedtrue
      \@latex@warning{Citation `\@citeb' undefined}}{}}%
  \@esphack}
\renewcommand{\cite}[1]{\oldcite{#1:\articlename}}
%   \end{macrocode}


% Horrible fiddling with the bibliography citation mechanism. 
% This has to be globally available in case rmaacite puts it in the
% \texttt{.aux} file. 
%    \begin{macrocode}
\def\lbibcite#1#2#3{\global\@namedef{b@#1}{#2}
                    \global\@namedef{by@#1}{#3}}
%    \end{macrocode}

% \begin{macro}{\oldbibitem}
% \begin{macro}{\@rmaaBibArg}
% \begin{macro}{\@rmaaSpace}
% \begin{macro}{rmaabib}
% Of course, |\bibitem| needs to get the uniquification treatment as
% well. While we are at it, we redefine it to use consecutive integers
% for the key if the user leaves it blank (using the counter
% |rmaabib|). This allows one to put 
% |\bibitem{}| for all your references without getting annoying
% messages about multiply-defined labels.
% 
% We also take the opportunity to insert the optional argument to
% \cs{bibitem} if it is missing in the form Author(DATE) so that natbib
% doesn't try to do numerical citations. Natbib already provides a
% hook \cs{@extra@b@citeb} to help uniquify the citation key, so we
% just plug into that. However, it puts it \emph{after} the original
% key, not before as I had had it. Easiest solution is to change the
% way that I do it. 02 Jan 2003. Note that the \cs{@extra@b@citeb}
% does not get written to the .aux  file in the \cs{citation} command
% but it does work. 
% \changes{3.14}{06 Feb 2003}{Reinstalled old definition of \cs{bibitem}
%  for the case where the nonatbib option is requested, so that the
%  rmaacite package will work. This means that you can't use rmaacite
%  and natbib in the same volume.}
%    \begin{macrocode}
\if@natbib\def\@extra@b@citeb{:\articlename}\fi
%%\else
\let\oldbibitem\bibitem
\newcommand{\@rmaaBibArg}{}
\newcommand{\@rmaaSpace}{ }
\newcounter{rmaabib}
\if@natbib
\renewcommand{\bibitem}[2][ABC(2003)]{%
  \renewcommand{\@rmaaBibArg}{#2}
  \ifx\@rmaaBibArg\@rmaaEmpty
     \addtocounter{rmaabib}{1}
     \oldbibitem[#1]{\thermaabib:\articlename}
  \else\ifx\@rmaaBibArg\@rmaaSpace
     \addtocounter{rmaabib}{1}
     \oldbibitem[#1]{\thermaabib:\articlename}
  \else 
     \oldbibitem[#1]{#2:\articlename}
  \fi\fi
  \renewcommand{\@rmaaBibArg}{}
  }
\else
\renewcommand{\bibitem}[1]{%
  \renewcommand{\@rmaaBibArg}{#1}
  \ifx\@rmaaBibArg\@rmaaEmpty
     \addtocounter{rmaabib}{1}
     \oldbibitem{\articlename:\thermaabib}
  \else\ifx\@rmaaBibArg\@rmaaSpace
     \addtocounter{rmaabib}{1}
     \oldbibitem{\articlename:\thermaabib}
  \else
     \oldbibitem{\articlename:#1}
  \fi\fi
  \renewcommand{\@rmaaBibArg}{}
  }
\fi
%%\fi
%    \end{macrocode}
% \end{macro}
% \end{macro}
% \end{macro}
% \end{macro}

% \begin{macro}{oldpcite}
% \begin{macro}{oldscite}
% Then, we need to repeat the above scheme, but for the |rmaacite|
% definition. 
%    \begin{macrocode}
\newcommand{\@redefinebibitem}{%
  \let\oldbibitem\bibitem
  \let\oldcite\cite
  \let\oldpcite\pcite
  \let\oldscite\scite
  \def\bibitem[##1<##2>]##3{%
    \renewcommand{\@rmaaBibArg}{##3}
    \ifx\@rmaaBibArg\@rmaaEmpty
       \addtocounter{rmaabib}{1}
       \oldbibitem[##1<##2>]{\thermaabib:\articlename}
    \else\ifx\@rmaaBibArg\@rmaaSpace
       \addtocounter{rmaabib}{1}
       \oldbibitem[##1<##2>]{\thermaabib:\articlename}
    \else 
       \oldbibitem[##1<##2>]{##3:\articlename}
    % trust the author here
%%       \oldbibitem[##1<##2>]{##3}
    \fi\fi
    }
  }
%    \end{macrocode}
% \end{macro}
% \end{macro}


% \subsection{HTML Table of Contents}
% \label{sec:hyper-TOC}
% \begin{macro}{\HyperTOCEntry}
% \begin{macro}{\HyperTOCpart}
% \begin{macro}{\HyperTOCmeeting}
% \begin{macro}{\HyperTOCeditors}
% \begin{macro}{\HyperTOCpreamble}
% These commands are used to write information to the \texttt{.htoc}
% file, which is used in generating a Web Page of the Table of
% Contents. Since they only ever appear preceded by |\protect| in the
% argument to an |\addtocontents| command,  they should never need to
% be expanded by the macros. However, we define them here just in case
% \dots
%    \begin{macrocode}
\newcommand{\HyperTOCEntry}[3]{\relax}
\newcommand{\HyperTOCpart}[1]{\relax}
\newcommand{\HyperTOCmeeting}[3]{\relax}
\newcommand{\HyperTOCeditors}[1]{\relax}
\newcommand{\HyperTOCpreamble}{\relax}
%    \end{macrocode}
% \end{macro}
% \end{macro}
% \end{macro}
% \end{macro}
% \end{macro}
%

% \subsection{Formatting the poster abstracts} 
% These can now be used without using the book option.
%    \begin{macrocode}

\newcommand{\@makeminititle}{%
  \newpage
  \null 
  \vskip 3em%
  \begin{center}%
    {\@title \par}%
  \end{center}
  \vskip 3em%
  }

\newif\if@abstractsads
\@abstractsadsfalse
\newcommand\abstractsADS{\@abstractsadstrue}

\if@book\RequirePackage{multicol}\fi
\newenvironment{abstracts}[1][Abstracts of Contributed Papers]{%
  %% First, make sure we start on an odd page (06 Jun 2000)
  \if@book\clearpage\thispagestyle{empty}\cleardoublepage\fi
  %%
  \renewcommand{\articlename}{ABSTRACTS}
  \title{#1}\shortauthor{Abstracts}\shorttitle{Abstracts}
  \listofauthors{}
  \if@book
  \addcontentsline{toc}{chapter}{\protect\@fiddleauthors{\@TOCtitle}{\rmaaTOCauthors}}%
  \addtocontents{htc}{\protect\HyperTOCEntry{\@rawtitle}{\@rawauthors}{\articlename}}
  \fi
  \pagestyle{rmaa} % this now seems necessary (06 Jun 2000)
  \thispagestyle{rmaafirstpage}
%%  \twocolumn[\@makeminititle]
  \renewcommand{\abstract}[1]{%
    \renewcommand\ADS@abstract{##1}%
    \renewcommand{\@abstract}{##1}\@writeabstract%
  }
  \renewcommand{\title}[2][]{%
    \gdef\ADS@authors{}%
    \gdef\ADS@affil{}%
    \gdef\ADS@references{}%
    \gdef\ADS@keywords{}%
    \renewcommand{\@rmaaArg}{##1}
    \ifx\@rmaaArg\@rmaaEmpty
       \renewcommand{\@TOCaftertitle}{}
    \else
       \renewcommand{\@TOCaftertitle}{ (\@rmaaArg)}
    \fi
    \renewcommand{\@title}{\rmaa@TOCTitleStyle{##2}}%
    \renewcommand{\@TOCtitle}{\rmaa@TOCTitleStyle{##2}\@TOCaftertitle}%
    \renewcommand{\rmaaIndexTitle}{\rmaa@IndexTitleStyle{##2\@TOCaftertitle.}}%
    \renewcommand\ADS@title{##2}%
    \xdef\ADS@firstpage{\@arabic{\c@page}}%
    }
  \renewcommand{\author}[1]{\let\@thanks\@empty\renewcommand\@author{##1} }
  \renewcommand{\footnotetext}[2][XX]{\par\noindent\textsuperscript{##1}~##2}
  \renewcommand{\footnoterule}{\kern-3\p@ \hrule \@width 3em \kern 2.6\p@}
  \if@book\@begofarticlehook\fi
  \if@book\begin{multicols}{2}[\@makeminititle]\else\@makeminititle\fi
  \begingroup
  \setcounter{footnote}{0}
  }% and at the end of the environment....
  {\endgroup\if@book\end{multicols}\listofabstracts\label{RMAALastPage}\clearpage\fi}
%% auxiliary commands used:
%% \@writeabstract does all the actual formatting 
\newcommand{\@writeabstract}{%
  \if@book%
  \addcontentsline{loa}{abstract}{\protect\@fiddleauthors{\@TOCtitle}{\rmaaTOCauthors}}%
  \fi%
  \begin{center}\@title\par\@author\end{center}
  \@abstract
  \par\noindent\footnoterule\nobreak\footnotesize\@thanks
  \par\noindent\hrulefill\goodbreak
  \vskip 1cm \@plus0.5cm
  %% write ADS entry
  \if@book\if@abstractsads\ADS@dumpdata\fi\fi%
  \endgroup % this pair is to make each abstract a self-contained group
  \begingroup
  \setcounter{footnote}{0}
  }

\if@book
%% List of abstracts (stuff written to .loa file)
\newcommand{\listofabstracts}{%
  \begin{multicols}{2}[%
    \vskip 1.5cm \vfill % leave at least 1.5 cm and go to bottom of page if possible
    \centerline{\MakeUpperCase{List of Abstracts}}
    \vskip 1.5cm \@plus0.5cm
    ][6cm] % this makes it start a new page if there is not at least 6cm of space
    \renewcommand{\@fiddleauthors}[2]{##1\hskip1em plus2em minus0.4em\penalty-100  ##2}
    \parindent\z@
%%    \parskip.5\p@ \@plus .7\p@\relax 
    \parskip3\p@ \@plus 1\p@ \@minus 1\p@\relax % allow quite a lot of
                                % leeway here to avoid breaks inside
                                % entries 
    \renewcommand{\\}{\relax} % temporarily turn off explicit line breaks just in case
    \renewcommand{\altaffilmark}[1]{\relax} % ditto with this
    \@starttoc{loa}
  \end{multicols}
  }
\newcommand*{\l@abstract}[2]{%
  \par\penalty-1000
  #1\hskip 1em\penalty0\null\nobreak\hfill#2}
\fi

%    \end{macrocode}
 
% \subsection{Code for typesetting an entire volume of conference
% proceedings} \label{sec:whole-volume}
%    \begin{macrocode}

\if@book
%% environment for including the individual articles
\newcommand{\@fiddleauthors}[2]{\relax}
\renewcommand{\theequation}{\@arabic\c@equation}
\newenvironment{articles}%
{ \begingroup % make it a group so that author definitions do not
              % spill over
  \gdef\ADS@authors{}
  \pagestyle{rmaa}
  \renewenvironment{document}{%
    \@begofarticlehook
    \addtocounter{chapter}{1}%
    \addcontentsline{toc}{chapter}{%
      \protect\@fiddleauthors{\@TOCtitle}{\rmaaTOCauthors}}%
    \addtocontents{htc}{\protect\HyperTOCEntry{\@rawtitle}{\@rawauthors}{\articlename}}
    }{\relax}
  \renewcommand{\documentclass}[2][foo]{\relax} 
  \renewcommand{\usepackage}[2][foo]{%
    \def\@packagename{##2}\def\@rmaacitename{rmaacite}%
    \ifx\@packagename\@rmaacitename
    \@rmcitetrue\ClassInfo{rmaa}{This paper uses rmaacite}
    \makeatletter\input{rmaacite.sty}\makeatother % read it in here so that all changes
    \else                                         % are local to paper
    \if@debug\ClassWarning{rmaa}{Package \@packagename\space was requested by \articlename\space}\fi
    \fi} 
  %% Don't allow individual articles to mess with the pages, volume or year
  \renewcommand{\SetFirstPage}[1]{\relax}
  \renewcommand{\SetYear}[1]{\relax}
  \renewcommand{\SetVolume}[1]{\relax}
  \renewcommand{\SetConfTitle}[1]{\relax}
  %% Don't allow individual articles to change these
  \renewcommand{\nonstopmode}{\relax}
  \renewcommand{\scrollmode}{\relax}
  \renewcommand{\outputplates}{\relax}
  %% XXXX need to check if using rmaacite and act accordingly
  %% avoid multiple declarations of the same length
  \let\oldnewlength=\newlength
  \renewcommand{\newlength}[1]{%
    \edef\reserved@a{\expandafter\@gobble\string ##1}%
    \@ifundefined{\reserved@a}{%
      \oldnewlength{##1}}{%
      \if@debug{\ClassInfo{rmaa}{\reserved@a\space already defined: doing nothing}}\fi%
      \relax}}%
  \renewcommand\newcommand{\providecommand}
  }%
{\endgroup\onecolumn}
     

%    \end{macrocode}
% Command to be used in `articles' environment for including an
% article
%    \begin{macrocode}
\newcommand{\article}[2][]{%
  % optional first argument is extra text to put in brackets after title in TOC
  % e.g. ``Invited Talk'' or ``Poster''
  \renewcommand{\articlename}{#2}
  % reset loads of counters
  \setcounter{footnote}{0}
  \setcounter{figure}{0}
  \setcounter{table}{0}
  \setcounter{section}{0}
  \setcounter{equation}{0}
  \setcounter{rmaabib}{0}
  \setcounter{author}{0}
  \setcounter{AUTHOR}{1}
  \gdef\ADS@authors{}%
  \gdef\ADS@affil{}%
  \gdef\ADS@affiliations{}%
  \gdef\ADS@references{}%
  \@firstauthortrue
  \@firstaffiltrue
  \@firstaltaffiltrue
  \@firstkwdtrue
  \@rmcitefalse
  \if@suppressfulladdresses\else\@needaddressestrue\fi
  \renewcommand{\@rmaaArg}{#1}
  \ifx\@rmaaArg\@rmaaEmpty
     \renewcommand{\@TOCaftertitle}{}
  \else
     \renewcommand{\@TOCaftertitle}{ (#1)}
  \fi
%    \end{macrocode}
% In the main journal articles must start on an odd page. 
%    \begin{macrocode}
  \if@journal\clearpage\thispagestyle{empty}\cleardoublepage\fi
%% Now input the file, then run the end-of-article hook (reinstated 28
  %% Sep 2001)
  {\input{\@articledir/\articlename.tex}\@endofarticlehook}}
%    \end{macrocode}
% Command to separate articles into sessions. Write an entry in the
% TOC, makes sure we start on an odd page (added 06 Jun 2000), then do
% a title page if \cs{SessionTitlePages} was called. 
% \changes{3.12}{04 Feb 2003}{Added option to have title page for each
% session. Turn it on with \cs{SessionTitlePages}.}
%    \begin{macrocode}
\newcommand{\session}[1]{
  %% this does nothing except write an entry in the TOC
  \addcontentsline{toc}{part}{\rmaa@TOCPartStyle{#1}}
  \addtocontents{htc}{\protect\HyperTOCpart{#1}}
  %% and also makes sure we start on an odd page (added 06 Jun 2000)
  \clearpage\thispagestyle{empty}\cleardoublepage
  \stepcounter{part}
  \if@sessiontitlepages{%
    \onecolumn
    \thispagestyle{empty}
    \renewcommand{\articlename}{Section\thepart TitlePage}
    \label{RMAAFirstPage}
    \make@session@title{#1}
    \label{RMAALastPage}
    \clearpage\thispagestyle{empty}\cleardoublepage
  }\fi
  }
\newcommand{\make@session@title}[1]{%
  {%
    \setlength\parindent{0pt}
    \vspace*{\stretch{1}}
    \rule{\linewidth}{1mm}
    \begin{flushright}\Huge\bfseries\boldmath
      #1
    \end{flushright}
    \rule{\linewidth}{1mm}
    \vspace*{\stretch{2}}
  }
}
%    \end{macrocode}
% Command to turn on separate title pages for each session. 
%    \begin{macrocode}
\newcommand{\SessionTitlePages}{\@sessiontitlepagestrue}
\newcommand{\BoldSessionTitles}{%
  \renewcommand{\rmaa@TOCPartStyle}[1]{\textbf{\rmaa@TOCTitleStyle{##1}}}
  }
\newcommand{\ArticleDir}[1]{\renewcommand{\@articledir}{#1/}\graphicspath{{1/}}}
\newcommand{\@articledir}{.}
%\newcommand{\ps@rmaafront}{\ps@rmaa}% page style for front matter (same as ever, actually)
\newcommand{\ps@rmaafront}{% page style for front matter:  page # in each foot
  \renewcommand\@oddfoot{\hfil{\thepage}}%
  \renewcommand\@evenfoot{{\thepage}\hfil}%
  % shorttitle in both heads
  \renewcommand\@evenhead{\hfil\@shorttitle\hfil}%
  \renewcommand\@oddhead{\@evenhead}
  }
\newcommand{\ps@rmaafrontfirstpage}{% page style for first page of front matter
  % page # in each foot, on outside edge
  \renewcommand\@oddfoot{\hfil\thepage}%
  \renewcommand\@evenfoot{\thepage\hfil}%
  % nothing in the even heads
  \renewcommand\@evenhead{\relax}%
  \renewcommand\@oddhead{\relax}%
  }
%% slight modification to the standard \frontmatter and \mainmatter
%% commands, just to set the \pagestyle
\renewcommand\frontmatter{%
  \latexonecolumn
  \cleardoublepage
  \@frontmattertrue
  \@mainmatterfalse
  \pagenumbering{roman}
  \pagestyle{rmaafront}
  \addtocounter{page}{2} % leave space for the title page
  }
\renewcommand\mainmatter{%
  \pagestyle{empty}
  \if@rmtwocol\twocolumn\fi
  \cleardoublepage
  \@frontmatterfalse
  \@mainmattertrue
  \pagenumbering{arabic}
  \pagestyle{rmaa}
  }
\renewcommand\backmatter{%
  \clearpage\pagestyle{empty}\cleardoublepage
  \@mainmatterfalse
  %% put extra space in the TOC
  \addtocontents{toc}{\bigskip}
  }
%    \end{macrocode}
%
% \subsection{Front sections}\label{sec:front-sections}
% \begin{macro}{\notitle}
% \begin{macro}{\shortname}
% \begin{macro}{frontsection}
% The environment \texttt{frontsection} is for doing a section in the
% front matter (preface, list of participants, or whatever). It has 
% one optional argument:
% \begin{enumerate}
% \item[\#1] Short title (no embedded spaces) for
% use in making PS files. E.g., ``list\_of\_participants''.   
% \end{enumerate}
% and two mandatory arguments: 
% \begin{enumerate}
% \item[\#2] Title of section (e.g., `Preface')
% \item[\#3] Commands to be evaluated before processing the contents
% of the environment. This is mainly useful for putting in the things
% that would go in the preamble of a normal article  (e.g.,
% |\listofauthors|, |\indexauthor|, |\notitle|) 
% \end{enumerate}
%    \begin{macrocode}
\newcommand{\notitle}{\@titlevisiblefalse} \newcommand{\shortname}{}
\newenvironment{frontsection}[3][]{%
  \begingroup
  \renewcommand{\shortname}{#1}
  \ifx\shortname\@rmaaEmpty
     \renewcommand{\articlename}{#2}
  \else
     \renewcommand{\articlename}{\shortname}
  \fi
  \@begofarticlehook
  \@titlevisibletrue % default
  \thispagestyle{rmaafrontfirstpage}
  \listofauthors{} % blank out this in case it is not specified
  \title{{#2}}\shorttitle{#2}\shortauthor{#2}
  #3 % evaluate the command given as second mandatory arg
  \addcontentsline{toc}{chapter}{\protect\@fiddleauthors{\@TOCtitle}{\rmaaTOCauthors}}%
  \addtocontents{htc}{\protect\HyperTOCEntry{\@rawtitle}{\@rawauthors}{\articlename}}
  \if@titlevisible
     \vspace*{1cm}\centerline{\rmaa@FrontSecStyle{\@title}}\vspace*{1cm}
  \fi
  \setlength{\parskip}{3ex}
  }{\frontlabel{RMAALastPage}\clearpage\endgroup}
%    \end{macrocode}
% \end{macro}
% \end{macro}
% \end{macro}
%    \begin{macrocode}
%% Definition of the command \frontlabel, which rights out something like `F4' to
%% the .aux file, instead of `\textlatin  {i}'. This is so it is easy to post-process
%% the file for making PS files of individual sections
\newcommand{\frontlabel}[1]{%
  \@bsphack
  \protected@write\@auxout{}%
     {\string\newlabel{\articlename:#1}{{}{F\@arabic{\c@page}}}}
  \@esphack
  }
%% User command to say X, Y, Z,\\ Editors of the Proceedings
\newcommand{\EditorSignature}{\vspace*{\parskip}\rightline{\@editors}%
  \rightline{Editors of the proceedings}}
%% Fiddle with \thesection, etc. so that they don't include the chapter number
%% Basically go back to the `article' definitions of these
\renewcommand{\thesection}{\@arabic\c@section}
\renewcommand{\thefigure}{\@arabic\c@figure}
\renewcommand{\thetable}{\@arabic\c@table}

\newcommand{\@oldargs}{}
\newcommand{\@newargs}{}
%% add TOC entry for each paper
%\renewcommand{\title}[1]%
%{\renewcommand{\@title}{\rmaa@TitleStyle{#1}}%
%\addtocounter{chapter}{1}\addcontentsline{toc}{chapter}{#1}%
%}

%% Formatting of the table of contents
%% dimensions of various parts
\newlength{\TOCtextwidth}\setlength{\TOCtextwidth}{12.65cm}
\newlength{\TOCpagenumwidth}\setlength{\TOCpagenumwidth}{1.3cm}
\newlength{\TOCentrywidth}\setlength{\TOCentrywidth}{\TOCtextwidth}
\addtolength{\TOCentrywidth}{-\TOCpagenumwidth} 
\newlength{\TOCextramargin}\setlength{\TOCextramargin}{1.25cm}
\newcommand{\TOCbox}[1]{\hspace{\TOCextramargin}\parbox{\TOCtextwidth}{#1}\par}
\newlength{\TOCextramarginx}\setlength{\TOCextramarginx}{\TOCextramargin}\addtolength{\TOCextramarginx}{4.3cm}
\newlength{\TOCtextwidthx}\setlength{\TOCtextwidthx}{\TOCtextwidth}\addtolength{\TOCtextwidthx}{-4.3cm}
\newcommand{\TOCboxx}[1]{\hspace{\TOCextramarginx}\parbox{\TOCtextwidthx}{#1}\par}

\newcommand{\RMAAbannerimage}{rmsc}
\renewcommand\tableofcontents{%
  \renewcommand{\articlename}{CONTENTS}
  \begingroup
  \@begofarticlehook
  \if@twocolumn
     \@restonecoltrue\onecolumn
  \else
     \@restonecolfalse
  \fi
  %% 9pt for journal style
  \if@proceedings\else\small\fi
  \if@newstyle\small\fi %% same for new style
  \if@proceedings\thispagestyle{rmaafrontfirstpage}\else\thispagestyle{empty}\fi
  \shortauthor{Contents}\shorttitle{Contents}
  \vspace*{18mm}
  \if@printbanner
  %% put in the banner header (need to get better image of this)
     \TOCbox{\hspace*{-0.15cm}\includegraphics[width=0.9\TOCtextwidth]{\RMAAbannerimage}}
     \vspace*{8mm}
  %% volume and publication date (should really be in funny font)
     \TOCboxx{\@volumen\\[2mm] \@mes\ \@anyo}
  \else
  %% put in the banner header (need to get better image of this)
     \TOCbox{\hspace*{-0.15cm}\includegraphics[width=0.9\TOCtextwidth]{\RMAAbannerimage}}
     \vspace*{12.5mm}
  %% volume and publication date (should really be in funny font)
     \TOCboxx{~\\[2mm] ~}
  \fi
  \vspace*{16mm}
  %% This is the command that fiddles with the TOC entry so that the
  %% title is flushleft but the authors are flushright
  \renewcommand{\@fiddleauthors}[2]{%
    % first put the paper title
    \renewcommand{\\}{\relax} % temporarily turn off explicit line breaks just in case
    \renewcommand{\altaffilmark}[1]{\relax} % ditto with this
    \renewcommand{\textsuperscript}[1]{\relax} % ditto with this 
    ##1
    % now save the author list in a box
    \sbox\@tempboxa{##2}%
    % check that box is shorter than the total width available for the entry
    \ifdim \wd\@tempboxa <\TOCentrywidth
       % if so, then put it right here, flush against right margin
       \penalty10000\hskip1em plus10em\penalty0\hbox{}\nobreak\hfill\box\@tempboxa
       % Note that we force a space of at least 1em between title and authorlist, 
       % stretchable to 11em since we may have very short titles but still with an 
       % authorlist too long to go on the same line.  The only linebreak allowed is 
       % directly before the authorlist.
    \else
       % if it is too long, we set the original authorlist text in its own parbox
       \par\vskip1ex\parbox[b]{\TOCentrywidth}{\raggedleft ##2}
    \fi
    }
  \@starttoc{toc}%
  \if@restonecol\twocolumn\fi
  \frontlabel{RMAALastPage}
  \clearpage
  \thispagestyle{empty}
  \cleardoublepage % make sure we start on an odd page
  \endgroup

  %%% ALSO, do the HTML TOC
  %% The file that will be used to make the HTML TOC
  \@starttoc{htc}
  %% All the stuff in this file does nothing in the book style
  \addtocontents{htc}{\protect\HyperTOCmeeting{\@conftitle}{\@reunionplace}{\@reuniondate}}
  \addtocontents{htc}{\protect\HyperTOCeditors{\@editors}}
  \addtocontents{htc}{\protect\HyperTOCpreamble}
  \addtocontents{htc}{\protect\HyperTOCEntry{Table of Contents}{}{\articlename}}

  %% and the ADS file
  \ADS@openfile
  }
%% This is now much simpler
\renewcommand\l@chapter[2]{%TOC entry for articles
  \ifnum \c@tocdepth >\m@ne
    \vskip 1.0em \@plus\p@
    \TOCbox{%
      \parbox[b]{\TOCentrywidth}{#1}\parbox[b]{\TOCpagenumwidth}{\hfill#2}
    }
  \fi} 
\renewcommand\l@part[2]{%TOC entry for sessions: centered, no page number
  \ifnum \c@tocdepth >\m@ne
    \vskip 1.5em \@plus1.5em \goodbreak
    \TOCbox{%
      \parbox[b]{\TOCentrywidth}{\hfil#1\hfil}\hfill
    }
  \fi\nobreak} 

%% Formatting theindex
%% strongly encourage page breaks between index entries
\renewcommand{\@idxitem}{\par\penalty-1000\hangindent.45cm} 
\renewcommand{\indexspace}{\par\penalty-1000\vskip 10\p@ \@plus15\p@ \@minus3\p@\relax}
\renewenvironment{theindex}{%
  %% first, make sure we start on an odd page (06 Jun 2000)
  \clearpage\thispagestyle{empty}\cleardoublepage
  %% this looks a bit like an article but is two columns and without abstract, etc
  \renewcommand{\articlename}{INDEX}
  \title{Author Index}\shortauthor{Author Index}\shorttitle{Author Index}
  \listofauthors{}
  \addcontentsline{toc}{chapter}{\protect\@fiddleauthors{\@TOCtitle}{\rmaaTOCauthors}}%
  \addtocontents{htc}{\protect\HyperTOCEntry{\@rawtitle}{\@rawauthors}{\articlename}}
  \pagestyle{rmaa}
  \thispagestyle{rmaafirstpage}
  \@begofarticlehook
  %% 9pt for journal style
  \if@proceedings\else\small\fi
  \if@newstyle\small\fi
  \twocolumn[\@makeminititle]
  %% vertical spacing (see also definition of \@idxitem and \indexspace
  \parskip\z@ \@plus .6\p@\relax % allow quite a lot of leaway here to avoid breaks inside entries
  %% horizontal spacing, breaking and glue parameters (see also the stuff in the definition 
  %% of \indexauthor
  \parindent\z@ 
  \exhyphenpenalty10000 % don't break on explicit hyphens, it always looks bad
  \rightskip=0pt plus.2em % allow a tiny bit of leeway at the ends of lines
  \spaceskip.3333em plus.5em minus.15em % a lot of leeway between words
  \xspaceskip.5em plus.5em minus.4em % this seems to have no effect
  \let\item\@idxitem
  \renewcommand{\\}{\relax} % temporarily turn off explicit line breaks just in case
  \renewcommand{\altaffilmark}[1]{\relax} % ditto with this
  \renewcommand{\textsuperscript}[1]{\relax} % ditto with this
  }{\label{RMAALastPage}\clearpage} % clearpage necessary so that we know about \shortauthor



%% end definitions for book styles
\fi

%    \end{macrocode}
% In order to do away with plates, we sometimes want a figure to
% appear in color only in the online version. The following command
% can bve used in a figure caption to typeset a note to this effect. 
%    \begin{macrocode}
\newcommand\ColorOnlineOnly{%
  \textsc{Note: this figure is available in color in the electronic
    version of this article, obtainable from}
  \texttt{http://\discretionary{}{}{}www.\discretionary{}{}{}astroscu.\discretionary{}{}{}unam.\discretionary{}{}{}mx/\discretionary{}{}{}\~{}rmaa/}.}
%% Plates
%% Define 2 counters: plate is the number in the whole document, plateloop is used 
%% internally by \outputplates
\newcounter{plate}
\newcounter{plateloop}
%% Currently limited to 15 plates in document
%% when we do a label, we need to write to \newlabel entries to the .aux file, one 
%% that says RMAAPlateX so that it can be referred to in \ouputplates, and one that
%% says articlename:RMAAPlateX so that the dvips program can find it and asociate it
%% with the right article
\newcommand{\@platelabeli}{\oldlabel{RMAAPlate1}\label{RMAAPlate1}}
\newcommand{\@platelabelii}{\oldlabel{RMAAPlate2}\label{RMAAPlate2}}
\newcommand{\@platelabeliii}{\oldlabel{RMAAPlate3}\label{RMAAPlate3}}
\newcommand{\@platelabeliv}{\oldlabel{RMAAPlate4}\label{RMAAPlate4}}
\newcommand{\@platelabelv}{\oldlabel{RMAAPlate5}\label{RMAAPlate5}}
\newcommand{\@platelabelvi}{\oldlabel{RMAAPlate6}\label{RMAAPlate6}}
\newcommand{\@platelabelvii}{\oldlabel{RMAAPlate7}\label{RMAAPlate7}}
\newcommand{\@platelabelviii}{\oldlabel{RMAAPlate8}\label{RMAAPlate8}}
\newcommand{\@platelabelix}{\oldlabel{RMAAPlate9}\label{RMAAPlate9}}
\newcommand{\@platelabelx}{\oldlabel{RMAAPlate10}\label{RMAAPlate10}}
\newcommand{\@platelabelxi}{\oldlabel{RMAAPlate11}\label{RMAAPlate11}}
\newcommand{\@platelabelxii}{\oldlabel{RMAAPlate12}\label{RMAAPlate12}}
\newcommand{\@platelabelxiii}{\oldlabel{RMAAPlate13}\label{RMAAPlate13}}
\newcommand{\@platelabelxiv}{\oldlabel{RMAAPlate14}\label{RMAAPlate14}}
\newcommand{\@platelabelxv}{\oldlabel{RMAAPlate15}\label{RMAAPlate15}}
\newcommand{\@outputplatelabeli}{\oldlabel{RMAAOutputPlate1}}
\newcommand{\@outputplatelabelii}{\oldlabel{RMAAOutputPlate2}}
\newcommand{\@outputplatelabeliii}{\oldlabel{RMAAOutputPlate3}}
\newcommand{\@outputplatelabeliv}{\oldlabel{RMAAOutputPlate4}}
\newcommand{\@outputplatelabelv}{\oldlabel{RMAAOutputPlate5}}
\newcommand{\@outputplatelabelvi}{\oldlabel{RMAAOutputPlate6}}
\newcommand{\@outputplatelabelvii}{\oldlabel{RMAAOutputPlate7}}
\newcommand{\@outputplatelabelviii}{\oldlabel{RMAAOutputPlate8}}
\newcommand{\@outputplatelabelix}{\oldlabel{RMAAOutputPlate9}}
\newcommand{\@outputplatelabelx}{\oldlabel{RMAAOutputPlate10}}
\newcommand{\@outputplatelabelxi}{\oldlabel{RMAAOutputPlate11}}
\newcommand{\@outputplatelabelxii}{\oldlabel{RMAAOutputPlate12}}
\newcommand{\@outputplatelabelxiii}{\oldlabel{RMAAOutputPlate13}}
\newcommand{\@outputplatelabelxiv}{\oldlabel{RMAAOutputPlate14}}
\newcommand{\@outputplatelabelxv}{\oldlabel{RMAAOutputPlate15}}
\newcommand{\@platerefi}{\oldpageref{RMAAPlate1}}
\newcommand{\@platerefii}{\oldpageref{RMAAPlate2}}
\newcommand{\@platerefiii}{\oldpageref{RMAAPlate3}}
\newcommand{\@platerefiv}{\oldpageref{RMAAPlate4}}
\newcommand{\@platerefv}{\oldpageref{RMAAPlate5}}
\newcommand{\@platerefvi}{\oldpageref{RMAAPlate6}}
\newcommand{\@platerefvii}{\oldpageref{RMAAPlate7}}
\newcommand{\@platerefviii}{\oldpageref{RMAAPlate8}}
\newcommand{\@platerefix}{\oldpageref{RMAAPlate9}}
\newcommand{\@platerefx}{\oldpageref{RMAAPlate10}}
\newcommand{\@platerefxi}{\oldpageref{RMAAPlate11}}
\newcommand{\@platerefxii}{\oldpageref{RMAAPlate12}}
\newcommand{\@platerefxiii}{\oldpageref{RMAAPlate13}}
\newcommand{\@platerefxiv}{\oldpageref{RMAAPlate14}}
\newcommand{\@platerefxv}{\oldpageref{RMAAPlate15}}
\newcounter{figplatei}
\newcounter{figplateii}
\newcounter{figplateiii}
\newcounter{figplateiv}
\newcounter{figplatev}
\newcounter{figplatevi}
\newcounter{figplatevii}
\newcounter{figplateviii}
\newcounter{figplateix}
\newcounter{figplatex}
\newcounter{figplatexi}
\newcounter{figplatexii}
\newcounter{figplatexiii}
\newcounter{figplatexiv}
\newcounter{figplatexv}
%% Printed output should use the counter `plate'
\renewcommand{\theplate}{\@arabic\c@plate}
\renewcommand{\theplateloop}{\@arabic\c@plateloop}
%% Command to output all the saved plates in this paper
\newcommand{\outputplates}{%
  \setcounter{plateloop}{0}
  \if@book\else\@endofarticlehook\fi % not entirely sure why, but this
                                     % seems to be needed
  \clearpage % first output the last page of text
  % now set pagestyle for blank even pages and make sure we start on an odd page
  \pagestyle{empty}
  \cleardoublepage
  \onecolumn
  \parindent\z@
  % first put in an extra sheet
  % the main purpose of this is so there is always a blank even page before 
  % each plate for filling in the PS file if the article has an odd number
  % of pages (see splitbook.sh)
  \null\vfil\centerline{PLATES}\vfil\clearpage\pagestyle{rmaaplate}\cleardoublepage
  % output each plate in turn
  \@whilenum\value{plateloop}<\value{plate}\do{%
    \addtocounter{plateloop}{1}
    % restore the figure counter from \@figcounti, \@figcountii, etc
    \setcounter{figure}{\value{figplate\roman{plateloop}}}
    % spew out the plate (in its own group for safety)
    {\csname @plate\roman{plateloop}\endcsname\par}
    }
  \pagestyle{rmaa}
  % reset the page number as if nothing had happened
  \addtocounter{page}{-\value{plate}}
  \addtocounter{page}{-\value{plate}}
  }
\newcommand\figcaption{%
  \def\@captype{figure}
  \renewcommand{\@makecaption}{\@makefigcap}
  \caption}
\newcommand\tabcaption{%
  \def\@captype{table}
  \renewcommand{\@makecaption}{\@maketabcap}
  \caption}

\newcommand{\plate}[2][1]{%
  %% User command to make a plate page: 
  %% Mandatory argument is what goes on the plate - this should not 
  %% include any floating environments (e.g. figure). Instead, \figcaption 
  %% should be used to do the captions. 
  %% Optional argument is the number of figures (i.e. captions) on the plate
  %% page (default is 1). This is necessary for incrementing the figure counter. 
  %%
  %% Note that we have to use a \gdef rather than an \xdef when we 
  %% save the plate output in a command. This means that in the book styles when it 
  %% comes to spewing out the plate, we will have forgotten about any user macros 
  %% that may have been defined. This can be got round by repeating the macros inside the plate.
  %% (Presumably, one would use a \newcommand rather than \renewcommand if they were new.)
  % increment global and local counters
  \addtocounter{plate}{1}
  % write appropriate label to .aux file
  \csname @platelabel\roman{plate}\endcsname
  % define a global macro called \@platei, \@plateii, \@plateiii, \@plateiv, ...
  \expandafter\gdef\csname @plate\roman{plate}\endcsname{%
    \csname @outputplatelabel\roman{plateloop}\endcsname
    \null\vfil #2 \vfil
    \cleardoublepage
    }
  % save the current shortauthor and shorttitle
  % need to be \xdef to take current values of \@shortauthor, etc 
  \expandafter\xdef\csname @shorta\roman{plate}\endcsname{\@shortauthor}
  \expandafter\xdef\csname @shortt\roman{plate}\endcsname{\@shorttitle}
  % save the current value of the figure counter for later use in \outputplates
  \setcounter{figplate\roman{plate}}{\value{figure}}
  % increment the figure counter according to the optional argument
  \addtocounter{figure}{#1}
  }
%% page style for plate pages 
\newcommand{\ps@rmaaplate}{%
  % Plate X (see page Y) in the foot
  \renewcommand\@oddfoot{\hfil\csname @shorta\roman{plateloop}\endcsname\ 
    (see page~\csname @plateref\roman{plateloop}\endcsname)\hfil}%
  \renewcommand\@evenfoot{\relax}%
  % nothing in even heads (blank pages)
  \renewcommand\@evenhead{\relax}
  % short title and page num in the odd heads
  \renewcommand\@oddhead{\hfil\csname @shortt\roman{plateloop}\endcsname%
    \hfil\llap{PLATE~\theplateloop}}
  }





%% Things for the User Guide
\if@guide
%% stolen from doc.sty
\def\MakeShortVerb#1{%
  \expandafter\ifx\csname cc\string#1\endcsname\relax
  \@shortvrbinfo{Made }{#1}%
  \add@special{#1}%
  \expandafter
  \xdef\csname cc\string#1\endcsname{\the\catcode`#1}%
  \begingroup
  \catcode`\~\active  \lccode`\~`#1%
  \lowercase{%
    \global\expandafter\let
    \csname ac\string#1\endcsname~%
    \gdef~{\verb~}}%
  \endgroup
  \global\catcode`#1\active
  \else
  \@shortvrbinfo\@empty{#1 already}%
  \fi}
\def\DeleteShortVerb#1{%
  \expandafter\ifx\csname cc\string#1\endcsname\relax
  \else
  \@shortvrbinfo{Deleted }{#1 as}%
  \rem@special{#1}%
  \global\catcode`#1\csname cc\string#1\endcsname
  \global \expandafter\let \csname cc\string#1\endcsname \relax
  \ifnum\catcode`#1=\active
  \begingroup
  \catcode`\~\active   \lccode`\~`#1%
  \lowercase{%
    \global\expandafter\let\expandafter~%
    \csname ac\string#1\endcsname}%
  \endgroup \fi \fi}
\def\@shortvrbinfo#1#2{%
  \PackageInfo{doc}{%
    #1\expandafter\@gobble\string#2 a short reference
    for \string\verb}}
\def\add@special#1{%
  \rem@special{#1}%
  \expandafter\gdef\expandafter\dospecials\expandafter
  {\dospecials \do #1}%
  \expandafter\gdef\expandafter\@sanitize\expandafter
  {\@sanitize \@makeother #1}}
\def\rem@special#1{%
  \def\do##1{%
    \ifnum`#1=`##1 \else \noexpand\do\noexpand##1\fi}%
  \xdef\dospecials{\dospecials}%
  \begingroup
  \def\@makeother##1{%
    \ifnum`#1=`##1 \else \noexpand\@makeother\noexpand##1\fi}%
  \xdef\@sanitize{\@sanitize}%
  \endgroup}
%% define | as shortcut for \verbatim
\MakeShortVerb{\|}
%% comand to specify arguments to LaTeX commands
\newcommand{\rmArg}[1]{%
  \ensuremath{\langle}{\scriptsize \textsf{\MakeUpperCase{#1}}}\ensuremath{\rangle}}
%\newcommand{\rmArg}[1]{\ensuremath{\langle}\textsc{#1}\ensuremath{\rangle}}
\providecommand{\BibTeX}{{\rmfamily B\kern-.05em%
    \textsc{i\kern-.025em b}\kern-.08em%
    T\kern-.1667em\lower.7ex\hbox{E}\kern-.125emX}}
%% Environment for code examples
\newenvironment{Code}{%
  \begin{list}{}{\setlength{\leftmargin}{0.05\columnwidth}} 
  \item[]}{\end{list}}
\fi
%</package>
%    \end{macrocode}

% \section{User commands for astronomical articles}
% \label{sec:user-commands}
% Some random definitions. These go both into the main class file and
% also into |adshtml.sty|.
% 
% Definitions of the journal name in different incarnations.  
%    \begin{macrocode}
%<*package|adshtml-style>
\newcommand{\RMAA}{RMxAA}
\newcommand{\RMAAlong}{Revista Mexicana de Astronom\'{\i}a y Astrof\'{\i}sica}
\newcommand{\RMAASC}{RMxAC}
\newcommand{\RMAASClong}{\RMAA\ (Serie de Conferencias)}
%    \end{macrocode}
% \begin{macro}{\ibidrule}
% This puts a horizontal rule with a period after it. It is used when
% the author list of a reference is identical to the preceding one. 
%    \begin{macrocode}
\newcommand{\ibidrule}{\rule[-0.03cm]{4.7em}{.01cm}.\ }
%    \end{macrocode}
% \end{macro}
% \begin{macro}{\ThisVolume}
% References to other articles in the same volume can be coded as,
% e.g., |\bibitem{} Bloggs, J. \ThisVolume{bloggs}|.
%    \begin{macrocode}
\newcommand{\ThisVolume}[1]{%
  \@year, \if@proceedings\RMAASC\else\RMAA\fi, \@volume,
  \if@book\oldpageref{#1:RMAAFirstPage}\else 000\fi\ (this volume)}
%    \end{macrocode}
% \end{macro}
% Other stuff that may be useful in ref lists.
%    \begin{macrocode}
\newcommand{\etal}{et al.\@}
%    \end{macrocode}
% Various new commands from AAS macros
%    \begin{macrocode}
\def\acknowledgments{\vskip 3ex plus .8ex minus .4ex}
\let\acknowledgements=\acknowledgments                  % second spelling
\let\jnl@style=\rm
\def\ref@jnl#1{{\jnl@style#1}}
\DeclareRobustCommand\aj{\ref@jnl{AJ}}                   
\DeclareRobustCommand\araa{\ref@jnl{ARA\&A}}             
\DeclareRobustCommand\apj{\ref@jnl{ApJ}}                 
\DeclareRobustCommand\apjl{\ref@jnl{ApJ}}                
\DeclareRobustCommand\apjs{\ref@jnl{ApJS}}               
\DeclareRobustCommand\ao{\ref@jnl{Appl.~Opt.}}           
\DeclareRobustCommand\apss{\ref@jnl{Ap\&SS}}             
\DeclareRobustCommand\aap{\ref@jnl{A\&A}}                
\DeclareRobustCommand\aapr{\ref@jnl{A\&A~Rev.}}          
\DeclareRobustCommand\aaps{\ref@jnl{A\&AS}}              
\DeclareRobustCommand\azh{\ref@jnl{AZh}}                 
\DeclareRobustCommand\baas{\ref@jnl{BAAS}}               
\DeclareRobustCommand\jrasc{\ref@jnl{JRASC}}             
\DeclareRobustCommand\memras{\ref@jnl{MmRAS}}            
\DeclareRobustCommand\mnras{\ref@jnl{MNRAS}}             
\DeclareRobustCommand\pra{\ref@jnl{Phys.~Rev.~A}}        
\DeclareRobustCommand\prb{\ref@jnl{Phys.~Rev.~B}}        
\DeclareRobustCommand\prc{\ref@jnl{Phys.~Rev.~C}}        
\DeclareRobustCommand\prd{\ref@jnl{Phys.~Rev.~D}}        
\DeclareRobustCommand\pre{\ref@jnl{Phys.~Rev.~E}}        
\DeclareRobustCommand\prl{\ref@jnl{Phys.~Rev.~Lett.}}    
\DeclareRobustCommand\pasp{\ref@jnl{PASP}}               
\DeclareRobustCommand\pasj{\ref@jnl{PASJ}}               
\DeclareRobustCommand\qjras{\ref@jnl{QJRAS}}             
\DeclareRobustCommand\skytel{\ref@jnl{S\&T}}             
\DeclareRobustCommand\solphys{\ref@jnl{Sol.~Phys.}}      
\DeclareRobustCommand\sovast{\ref@jnl{Soviet~Ast.}}      
\DeclareRobustCommand\ssr{\ref@jnl{Space~Sci.~Rev.}}     
\DeclareRobustCommand\zap{\ref@jnl{ZAp}}                 
\DeclareRobustCommand\nat{\ref@jnl{Nature}}              
\DeclareRobustCommand\iaucirc{\ref@jnl{IAU~Circ.}}
\DeclareRobustCommand\aplett{\ref@jnl{Astrophys.~Lett.}}
\DeclareRobustCommand\apspr{\ref@jnl{Astrophys.~Space~Phys.~Res.}}
\DeclareRobustCommand\bain{\ref@jnl{Bull.~Astron.~Inst.~Netherlands}}
\DeclareRobustCommand\fcp{\ref@jnl{Fund.~Cosmic~Phys.}}
\DeclareRobustCommand\gca{\ref@jnl{Geochim.~Cosmochim.~Acta}}
\DeclareRobustCommand\grl{\ref@jnl{Geophys.~Res.~Lett.}}
\DeclareRobustCommand\jcp{\ref@jnl{J.~Chem.~Phys.}}      
\DeclareRobustCommand\jgr{\ref@jnl{J.~Geophys.~Res.}}    
\DeclareRobustCommand\jqsrt{\ref@jnl{J.~Quant.~Spec.~Radiat.~Transf.}}
\DeclareRobustCommand\memsai{\ref@jnl{Mem.~Soc.~Astron.~Italiana}}
\DeclareRobustCommand\nphysa{\ref@jnl{Nucl.~Phys.~A}}
\DeclareRobustCommand\physrep{\ref@jnl{Phys.~Rep.}}
\DeclareRobustCommand\physscr{\ref@jnl{Phys.~Scr}}
\DeclareRobustCommand\planss{\ref@jnl{Planet.~Space~Sci.}}
\DeclareRobustCommand\procspie{\ref@jnl{Proc.~SPIE}}
\let\astap=\aap
\let\apjlett=\apjl
\let\apjsupp=\apjs
\let\applopt=\ao
\def\phn{\phantom{0}}
\def\phd{\phantom{.}}
\def\phs{\phantom{\ensuremath{-}}}
\def\phm#1{\phantom{#1}}
%    \end{macrocode}
% All these are redone as robust and using more modern commands where
% possible. The robustness is so they will not get expanded if used in
% abstracts and titles when these are passed to the ADS file. 
% \changes{3.16}{2003/04/07}{
%  Added \cs{Color} command for doing photometric colors.}
%    \begin{macrocode}
%%\def\sun{\hbox{$\odot$}}
\DeclareRobustCommand\sun{\ensuremath{\odot}}
%%\def\earth{\hbox{$\oplus$}}
\DeclareRobustCommand\earth{\ensuremath{\oplus}}
%%\def\lesssim{\mathrel{\hbox{\rlap{\hbox{\lower4pt\hbox{$\sim$}}}\hbox{$<$}}}}
%%\def\gtrsim{\mathrel{\hbox{\rlap{\hbox{\lower4pt\hbox{$\sim$}}}\hbox{$>$}}}}
\DeclareRobustCommand\lesssim{\mathrel{\hbox{\smaller\smaller\rlap{\hbox{\lower3pt\hbox{$\sim$}}}\raise1.5pt\hbox{$<$}}}}
\DeclareRobustCommand\gtrsim{\mathrel{\hbox{\smaller\smaller\rlap{\hbox{\lower3pt\hbox{$\sim$}}}\raise1.5pt\hbox{$>$}}}}
%%\def\lta{\mathrel{\spose{\lower 3pt\hbox{$\mathchar"218$}}
%%     \raise 2.0pt\hbox{$\mathchar"13C$}}}
%%\def\gta{\mathrel{\spose{\lower 3pt\hbox{$\mathchar"218$}}
%%     \raise 2.0pt\hbox{$\mathchar"13E$}}}
\DeclareRobustCommand\sq{\hbox{\rlap{$\sqcap$}$\sqcup$}}
%%\def\arcdeg{\hbox{$^\circ$}}
%%\def\arcmin{\hbox{$^\prime$}}
%%\def\arcsec{\hbox{$^{\prime\prime}$}}
\DeclareRobustCommand\arcdeg{\ensuremath{^\circ}}
\DeclareRobustCommand\arcmin{\ensuremath{^\prime}}
\DeclareRobustCommand\arcsec{\ensuremath{^{\prime\prime}}}
%%\def\fd{\hbox{$.\!\!^{\rm d}$}}
%%\def\fh{\hbox{$.\!\!^{\rm h}$}}
%%\def\fm{\hbox{$.\!\!^{\rm m}$}}
%%\def\fs{\hbox{$.\!\!^{\rm s}$}}
%%\def\fdg{\hbox{$.\!\!^\circ$}}
%%\def\farcm{\hbox{$.\mkern-4mu^\prime$}}
%%\def\farcs{\hbox{$.\!\!^{\prime\prime}$}}
%%\def\fp{\hbox{$.\!\!^{\scriptscriptstyle\rm p}$}}
\DeclareRobustCommand\fd{\ensuremath{.\!\!^\mathrm{d}}}
\DeclareRobustCommand\fh{\ensuremath{.\!\!^\mathrm{h}}}
\DeclareRobustCommand\fm{\ensuremath{.\!\!^\mathrm{m}}}
\DeclareRobustCommand\fs{\ensuremath{\mkern0.5mu.\mkern-5mu^\mathrm{s}\mkern-1mu}}
\DeclareRobustCommand\fdg{\ensuremath{.\!\!^\circ}}
\DeclareRobustCommand\farcm{\ensuremath{.\mkern-4mu^\prime}}
\DeclareRobustCommand\farcs{\ensuremath{\mkern1mu.\mkern-6mu^{\prime\prime}\mkern-1.5mu}}
%%\DeclareRobustCommand\farcs{\ensuremath{.\!\!^{\prime\prime}}}
\DeclareRobustCommand\fp{\ensuremath{.\!\!^{\scriptscriptstyle\rm p}}}
%%\def\micron{\hbox{$\mu$m}}
\DeclareRobustCommand\micron{\ensuremath{\mu\mathrm{m}}}
\let\la=\lesssim                        
\let\ga=\gtrsim
%%\def\case#1#2{\hbox{$\frac{#1}{#2}$}}
%%\def\slantfrac#1#2{\hbox{$\,^#1\!/_#2$}}
%%\def\onehalf{\slantfrac{1}{2}}
%%\def\onethird{\slantfrac{1}{3}}
%%\def\twothirds{\slantfrac{2}{3}}
%%\def\onequarter{\slantfrac{1}{4}}
%%\def\threequarters{\slantfrac{3}{4}}
\DeclareRobustCommand\case[2]{\ensuremath{\frac{#1}{#2}}}
\DeclareRobustCommand\slantfrac[2]{\ensuremath{\,^#1\!/_#2}}
\DeclareRobustCommand\onehalf{\slantfrac{1}{2}}
\DeclareRobustCommand\onethird{\slantfrac{1}{3}}
\DeclareRobustCommand\twothirds{\slantfrac{2}{3}}
\DeclareRobustCommand\onequarter{\slantfrac{1}{4}}
\DeclareRobustCommand\threequarters{\slantfrac{3}{4}}
%%\def\ubvr{\hbox{$U\!BV\!R$}}            
%%\def\ub{\hbox{$U\!-\!B$}}               
%%\def\bv{\hbox{$B\!-\!V$}}               
%%\def\vr{\hbox{$V\!-\!R$}}               
%%\def\ur{\hbox{$U\!-\!R$}}
%    \end{macrocode}
% New commands for ``photometric'' colors 
%    \begin{macrocode}
\DeclareRobustCommand\ubvr{\ensuremath{U\!BV\!R}}            
\DeclareRobustCommand\Color[2]{%
  \ensuremath{%
    #1\!\ifx#1V\!\fi-\!#2
  }}
\DeclareRobustCommand\ub{\Color{U}{B}}
\DeclareRobustCommand\bv{\Color{B}{V}}
\DeclareRobustCommand\vr{\Color{V}{R}}
\DeclareRobustCommand\ur{\Color{U}{R}}
%% New version of the \ion command, that actually works! 
%    \end{macrocode}
% \begin{macro}{\ion}
% This version of the |\ion| command allows you to write either
% |\ion{N}{4}| or |\ion{N}{iv}| in order to get N~IV. 
%    \begin{macrocode}
\DeclareRobustCommand\ion[2]{%
  \mbox{#1\kern0.2em%
    \smaller\rmfamily%
    \edef\@tempa{\@car#2\@nil}%
    \ifcat1\@tempa%
    \@Roman{#2}%
    \else%
    \uppercase{#2}%
    \fi}}
\newcounter{IonStage}
\DeclareRobustCommand\plainion[2]{\setcounter{IonStage}{#2}#1 \Roman{IonStage}}
%%\def\ion#1#2{\ensuremath{\mathrm{#1}\;}{\protect\small\rm\@Roman{#2}}}
%    \end{macrocode}
% \end{macro}
%    \begin{macrocode}
\DeclareRobustCommand\nodata{$\cdots$}
%    \end{macrocode}
% These are needed by the |\ion| command. Isn't there an easier way of
% doing this?
%    \begin{macrocode}
\let\@small=\small
\let\@normalsize=\normalsize
\let\@footnotesize=\footnotesize
\let\@scriptsize=\scriptsize
\let\@tiny=\tiny
\let\@large=\large
\let\@Large=\Large
\let\@LARGE=\LARGE
\let\@huge=\huge
\let\@Huge=\Huge
\newcommand{\smaller}{\small}
\renewcommand{\normalsize}{\renewcommand{\smaller}{\small}\@normalsize}
\renewcommand{\small}{\renewcommand{\smaller}{\footnotesize}\@small}
\renewcommand{\footnotesize}{\renewcommand{\smaller}{\scriptsize}\@footnotesize}
\renewcommand{\scriptsize}{\renewcommand{\smaller}{\tiny}\@scriptsize}
\renewcommand{\tiny}{\renewcommand{\smaller}{\tiny}\@tiny}
\renewcommand{\large}{\renewcommand{\smaller}{\normalsize}\@large}  
\renewcommand{\Large}{\renewcommand{\smaller}{\large}\@Large}  
\renewcommand{\LARGE}{\renewcommand{\smaller}{\Large}\@LARGE}  
\renewcommand{\huge}{\renewcommand{\smaller}{\LARGE}\@huge}  
\renewcommand{\Huge}{\renewcommand{\smaller}{\huge}\@Huge}  
%</package|adshtml-style>
%    \end{macrocode}


% \section{Other files}

% \subsection{The driver and make files}
% \label{sec:driver_make}
% These are put here for completeness, although they still need to be
% distributed separately since otherwise the user won't be able to
% unpack this file into its constituent parts.
% 
% The |Makefile| has two principal rules: |docs|, and |macros|. The
% |docs| rule produces a PDF and gzipped PS file of the documentation.
% I still haven't perfected the combination of rules required to get
% it to run \LaTeX{} enough times in all circumstances. The |macros|
% rule runs \LaTeX{} on |rmaa.ins| in order to unpack |rmaa.cls| and
% all the other little files (including the |Makefile| and the
% |rmaa.ins| files themselves\dots 
% 
% First set up lists of various types of file. 
%    \begin{macrocode}
%<*makefile>
DOCFILES=rmaa.pdf
MACRO_LTX_FILES=rmaa.cls rmaacite.sty adshtml.sty
SHELLSCRIPTS=makeads.sh
MACRO_OTHER_FILES=adshtml.perl sample.latex2html-init
MAKEFILE=Makefile
MACROFILES=$(MACRO_LTX_FILES) $(SHELLSCRIPTS) $(MACRO_OTHER_FILES)         
SKELETONFILES=ADS.tex 
%    \end{macrocode}
% Install the files read only into the package directory, specified via a path relative to the source directory. \changes{4.0}{2019/11/14}{Reorganize folder layout so that hardly anything is in the top level.  That way, things are clearer on github.}
%    \begin{macrocode}
INSTALL=install 
INSTALL_DIR=../packages
%    \end{macrocode}
% The default rule is to make everything. Installation must be
% explicitly asked for.
%    \begin{macrocode}
all: docs macros

install: install-docs install-macros

%    \end{macrocode}
% Most files are installed read-only to prevent editing. Skeleton
% files must be writeable. Must be careful not to overwrite the
% Makefile in the parent directory. 
%    \begin{macrocode}
install-macros: macros
        $(INSTALL) -m 444 $(MACRO_LTX_FILES) $(MACRO_OTHER_FILES) $(INSTALL_DIR)
        $(INSTALL) -m 555 $(SHELLSCRIPTS) $(INSTALL_DIR)
        $(INSTALL) -m 644 $(SKELETONFILES) $(INSTALL_DIR)

install-docs: docs
        $(INSTALL) rmaa.pdf $(INSTALL_DIR)/rm-fulldocs.pdf

%    \end{macrocode}
% Make the PDF and gzipped PS files. 
%    \begin{macrocode}
docs: $(DOCFILES)

%    \end{macrocode}
% Make everything else if either the |.dtx| or |.ins| files have
% changed. 
%    \begin{macrocode}
macros: rmaa.dtx rmaa.ins
%    \end{macrocode}
% Ensure that all the macro files can be written to then recreate
% them. The |touch| is there in case any of the files don't exist
% yet. 
%    \begin{macrocode}
        touch $(MACROFILES) $(MAKEFILE)
        chmod u+w $(MACROFILES) $(MAKEFILE) rmaa.ins
        latex "\input{rmaa.ins}"
%    \end{macrocode}
% \TeX{} automatically adds the |.tex| extension to extensionless
% files. We fix this and take the opportunity to convert
% back to TABs all that it has lovingly converted to spaces. 
%    \begin{macrocode}
        cp -f $(MAKEFILE) $(MAKEFILE).BCK && unexpand $(MAKEFILE).tex >| $(MAKEFILE) && rm $(MAKEFILE).tex
%    \end{macrocode}
% Make the shell scripts executable and everything unwriteable (so
% they don't get hand-edited by mistake). However, undo this for the
% files that the user is \emph{supposed} to modify. 
%    \begin{macrocode}
        chmod a+x $(SHELLSCRIPTS)
        chmod a-w $(MACROFILES) $(MAKEFILE)
        chmod u+w $(SKELETONFILES)

%    \end{macrocode}
% We now use PDF\LaTeX{} for this.
%    \begin{macrocode}
rmaa.idx: rmaa.dtx
        @echo "REGENERATING .idx FILE"
        pdflatex "\nonstopmode\input{rmaa.dtx}"

rmaa.aux: rmaa.dtx rmaa.ind rmaa.gls
        @echo "REGENERATING .aux FILE"
        pdflatex "\nonstopmode\input{rmaa.dtx}"

rmaa.ind: rmaa.idx
        @echo "REGENERATING INDEX"
        makeindex -s gind.ist rmaa

rmaa.gls: rmaa.glo
        @echo "REGENERATING CHANGELOG"
        makeindex -s gglo.ist -o rmaa.gls rmaa.glo

rmaa.pdf: rmaa.aux rmaa.ind rmaa.gls
        @echo "REGENERATING .pdf FILE"
        pdflatex "\nonstopmode\input{rmaa.dtx}"

clean: 
        rm -f rmaa.aux rmaa.dvi rmaa.ind rmaa.gls
%</makefile>
%    \end{macrocode}

% The |rmaa.ins| file uses the \textsf{docstrip} package to split out
% all the files and strip off this documentation.
%    \begin{macrocode}
%<*ins>
%%
%% Installation batch file for the RMAA LaTeX package
%%
\def\batchfile{rmaa.ins}
\input docstrip.tex
\preamble

PLEASE DO NOT EDIT THIS FILE! It is generated automatically from
the file rmaa.dtx using the accompanying Makefile and rmaa.ins 
Any hand-edited changes will be lost

A LaTeX document class for the RMxAA

Author: William Henney (w.henney@irya.unam.mx)

Latest version available from: http://www.irya.unam.mx/rmaa

Copyright (C) 1999-2019 Instituto de Astronomia, UNAM, Mexico

Permission is hereby granted, free of charge, to any person obtaining
a copy of this software and associated documentation files (the
"Software"), to deal in the Software without restriction, including
without limitation the rights to use, copy, modify, merge, publish,
distribute, sublicense, and/or sell copies of the Software, and to
permit persons to whom the Software is furnished to do so, subject to
the following conditions:

The above copyright notice and this permission notice shall be
included in all copies or substantial portions of the Software.

THE SOFTWARE IS PROVIDED "AS IS", WITHOUT WARRANTY OF ANY KIND,
EXPRESS OR IMPLIED, INCLUDING BUT NOT LIMITED TO THE WARRANTIES OF
MERCHANTABILITY, FITNESS FOR A PARTICULAR PURPOSE AND
NONINFRINGEMENT. IN NO EVENT SHALL THE AUTHORS OR COPYRIGHT HOLDERS BE
LIABLE FOR ANY CLAIM, DAMAGES OR OTHER LIABILITY, WHETHER IN AN ACTION
OF CONTRACT, TORT OR OTHERWISE, ARISING FROM, OUT OF OR IN CONNECTION
WITH THE SOFTWARE OR THE USE OR OTHER DEALINGS IN THE SOFTWARE.
\endpreamble
%    \end{macrocode}
% Clobber files without asking
%    \begin{macrocode}
\askforoverwritefalse
%    \end{macrocode}
% Make the \LaTeX{} macro files (class and styles)
%    \begin{macrocode}
\generate{% 
  \file{rmaa.cls}{\from{rmaa.dtx}{package}}
  \file{rmaa.ins}{\from{rmaa.dtx}{ins}}
  \file{rmaacite.sty}{\from{rmaa.dtx}{rmaacite}}
  \file{adshtml.sty}{\from{rmaa.dtx}{adshtml-style}}
  \file{ADS.tex}{\from{rmaa.dtx}{ads-container}}
  }
%    \end{macrocode}
% Make the other non-\LaTeX{} files. These are done separately since
% they can't use `\%\%' as a comment character for the
% preamble/postamble (and Perl files can't have the postamble I don't
% think). 
%    \begin{macrocode}
\generate{\nopreamble\nopostamble% 
  \file{Makefile}{\from{rmaa.dtx}{makefile}}
  \file{adshtml.perl}{\from{rmaa.dtx}{adshtml-perl}}
  \file{makeads.sh}{\from{rmaa.dtx}{ads-shell-script}}
  \file{sample.latex2html-init}{\from{rmaa.dtx}{latex2html-init}}
  }

%% Local Variables: 
%% mode: LaTeX
%% mode: Font-lock
%% End:
%</ins>
%    \end{macrocode}

% \subsection{Extended citation syntax}
% Modified from a file called mnrascite.sty
% I no longer remember where I stole the original version of this from
% 
% Alterations to the cite commands to make them do more sensible things.
% N.B. the bibliography style should produce entries which have labels where
% the date is delimited by angle brackets '<>' for these macros to work
% properly. They will not give very sensible error messages if you do not do 
% this properly.
% 
% A new command |\scite{key}| is defined which produces a citation which is 
% suitable for insertion in a sentence. i.e., it does not put brackets around 
% the whole citation, only the date. Also |\pcite{key}| produce output with no
% parentheses whatsoever. 
% 
% Should redefine |\@bibitem| really to have any hope of compatibility
% the command |\lbibcite| is defined to deal with the citations.
%    \begin{macrocode}
%<*rmaacite>
\def\bibitem[#1<#2>]#3{\par\hangindent=1.5em\hangafter=1
  \if@filesw
  {\def\protect##1{\string ##1\space}
    \immediate\write\@auxout{\string\lbibcite{#3}{#1}{#2}}}
  \fi}
\@redefinebibitem % what it says
%    \end{macrocode}
% Note that the |\lbibcite| defines two types of citation one contains the names,
% the other the year
%    \begin{macrocode}
\def\lbibcite#1#2#3{\global\@namedef{b@#1}{#2}
                    \global\@namedef{by@#1}{#3}}

\def\@bcite#1#2{({#1\if@tempswa , #2\fi})}
\def\@pcite#1#2{#1\if@tempswa , #2\fi}
%    \end{macrocode}
% This declaration contains what is put between the name and the date
% \changes{3.11}{2003/01/31}{Moved the \cs{articlename} to after the key for consistency
% with the new natbib stuff.}
%    \begin{macrocode}
\def\@brcite{\ (}
\def\@ebrcite{)}
\def\@fmttop#1{\@fmtcite[\csname #1\endcsname]}
\def\@citex[#1]#2{%
%%  \if@filesw\immediate\write\@auxout{\string\citation{#2}}\fi
%% WJH added 16 Jan 2002
  \@for\@citeb:=#2\do{%
    \edef\@citeb{\expandafter\@firstofone\@citeb:\articlename}%
    \if@filesw\immediate\write\@auxout{\string\citation{\@citeb}}\fi}%
%% to here
  \def\@citea{}%
  \@cite{\@for\@citeb:=#2\do
    {%
      \edef\@citeb{\expandafter\@firstofone\@citeb:\articlename}%
      \@citea\def\@citea{;\ }\@ifundefined
      {b@\@citeb}{{\bf ?}\@warning {Citation `\@citeb' on page \thepage \space undefined}}%
      {\csname b@\@citeb\endcsname\@citebro\csname by@\@citeb\endcsname\@citebrc}}}%
  {#1}}

\def\cite{\@ifnextchar [{\let\@citebro=\space
                          \let\@citebrc=\relax
                          \let\@cite=\@bcite\@tempswatrue\@citex}
                        {\let\@citebro=\space
                          \let\@citebrc=\relax
                          \let\@cite=\@bcite\@tempswafalse\@citex[]}}
\def\pcite{\@ifnextchar [{\let\@citebro=\space
                          \let\@citebrc=\relax
                          \let\@cite=\@pcite\@tempswatrue\@citex}
                        {\let\@citebro=\space
                          \let\@citebrc=\relax
                          \let\@cite=\@pcite\@tempswafalse\@citex[]}}
\def\scite{\@ifnextchar [{\let\@citebro=\@brcite
                          \let\@citebrc=\@ebrcite
                          \let\@cite=\@pcite\@tempswatrue\@citex}
                        {\let\@citebro=\@brcite
                          \let\@citebrc=\@ebrcite
                          \let\@cite=\@pcite\@tempswafalse\@citex[]}}

\renewenvironment{thebibliography}
  {\section*{REFERENCES}\vspace*{\afterbiblioskip}\small
    \if@book\RereadBiblio\fi
    \sloppy\clubpenalty4000\widowpenalty4000
    \sfcode`\.=1000\relax
    \setlength\parindent{\z@}\setlength\parskip{\z@}}
  {\relax\normalsize\@outputaddressesifnecessary} 
%</rmaacite>
%    \end{macrocode}

% \subsection{Auxilliary files for generating ADS data}

% The ADS data are written by \LaTeX{} to a |.ads| file. This file
% then has to be included in the following file (|ADS.tex|),
% which is run through \LaTeX{}2HTML (see below). 

%    \begin{macrocode}
%<*ads-container>
%%
%% Instructions:
%% Edit the \input{} line below to reflect the name of your file. 
\documentclass{article}
\usepackage[latin1]{inputenc}
\usepackage{adshtml}
\newcommand{\ADSfield}[2]
{\%#1  #2\\}
%%\newcommand{\ADSvfield}[2]
%%{\%#1 \begin{verbatim}#2\end{verbatim}}
\newenvironment{ADSarticle}{}{\par\bigskip}
\renewcommand{\bibitem}[2][]{\par}
\setlength{\parindent}{0pt}
\renewcommand{\ThisVolume}[1]{%  
  2001, \RMAASC, 12, \oldpageref{#1:RMAAFirstPage}}  
 

\begin{document}
\makeatletter
\input{rmbooksamp.ads}
\end{document}
%</ads-container>
%    \end{macrocode}

% This is the |adshtml.sty| style file used by |ADS.tex|. It just
% defines a few 
% macros that might be used in abstracts, etc (really it should define
% them all). This now gets all its code from
% Sec~\ref{sec:user-commands}. 
%    \begin{macrocode}
%<*adshtml-style>
%% Nothing else needs to go here now. 
%</adshtml-style>
%    \end{macrocode}
% The init file for latex2html.
%    \begin{macrocode}
%<*latex2html-init>
$PDFTOCAIRO = "pdftocairo";
$PDFLATEX = "pdflatex";
$LATEX = "latex";
$MAX_SPLIT_DEPTH = 0;
$HTML_VERSION = "4.0,math";
$NO_SIMPLE_MATH = 1;
$ADDRESS = "<p><address><a href=\"mailto:rmaa\@astro.unam.mx\">rmaa\@astro.unam.mx</address></p>"; 
$INFO = "<p>This document was created using latex2html and the rmaa latex macros</p>"; 
$NO_NAVIGATION = 1;
1; 
%</latex2html-init>
%    \end{macrocode}
% This is the |adshtml.perl| file that is used by \LaTeX{}2HTML. Some
% bizarre \TeX{}y command sometimes show up so we tell it to ignore
% these. All the astronomy macros that come from AAS are processed in
% \TeX{}. That way they get made to images but with an |ALT| tag that
% has the original code, which is what will be picked up by |lynx|,
% with the result that they get passed intact into the final |.txt|
% file. 
%    \begin{macrocode}
%<*adshtml-perl>
package main;

$dum = &do_cmd_makeatletter();

sub do_cmd_RMAA {
    local($_) = @_;
    "RMxAA" . $_;
}

sub do_cmd_RMAASC {
    local($_) = @_;
    "RMxAASC" . $_;
}

sub do_cmd_etal {
    local($_) = @_;
    "et al\." . $_;
}

sub do_cmd_ibidrule {
    local($_) = @_;
    "<img src=\"foo\" alt=\"\ibidrule\">" . $_;
}


sub do_cmd_bibitem {
    local($_) = @_;
    local($key);
    local($opt, $dummy)=&get_next_optional_argument; 
    $key = &missing_braces unless
        ((s/$next_pair_pr_rx/$key = $2; ''/eo)
        ||(s/$next_pair_rx/$key = $2; ''/eo));
    "<p>" . $_; 
}

&ignore_commands (<<_RAW_ARG_CMDS_);
lowercase
adjustfinalcols
textsuperscript # {}
newline # []
thinmuskip
mskip
smallskip
medskip
bigskip
hspace # {}
hspacestar # {}
vspace # {}
vspacestar # {}
mathrm
_at_mathrm
spacefactor
discretionary
unhbox 
voidb_at_x 
penalty 
_at_M
_at_m
null
IeC
_RAW_ARG_CMDS_

&process_commands_in_tex (<<_RAW_ARG_CMDS_);
sun
earth
lesssim
gtrsim
sq
argdeg
arcmin
arcsec
fd
fh
fm
fs
fdg
farcm
farcs
fp
micron
la
ga
case # {} # {}
slantfrac # {} # {}
onehalf
onethird
twothirds
onequarter
threequarters
ubvr
ub
bv
vr
ur
aj
araa
apj
apjl
apjs
ao
apss
aap
aapr
aaps
azh
baas
jrasc
memras
mnras
pra
prb
prc
prd
pre
prl
pasp
pasj
qjras
skytel
solphys
sovast
ssr
zap
nat
iaucirc
aplett
apspr
bain
fcp
gca
grl
jcp
jgr
jqsrt
memsai
nphysa
physrep
physscr
planss
procspie
_RAW_ARG_CMDS_

1;                              # This must be the last line
%</adshtml-perl> 
%    \end{macrocode}

% This is a shell script |makeads.sh| that generates the |_ads.txt| file from
% |_ads.tex| by using \LaTeX{}2HTML and |lynx|.  \changes{4.0}{2019/11/13}{Make sure that lynx writes utf-8 output instead of latin-1.}
% \changes{2.2beta13}{2002/03/12}{Complete rewrite of |makeads.sh|. Now much 
% improved I think}
% \changes{4.0}{2019/11/13}{Fix the at-sign substitution so that it only eliminates \cs{@} and does not mess with email addresses.}

%    \begin{macrocode}
%<*ads-shell-script>
#! /bin/sh
if [ -z $1 ]; then
    echo "Usage: $0 basename"
    exit
fi
BASENAME=$1
ADSBASENAME=${BASENAME}_ads
AUXFILE=$BASENAME.aux
TMPFILE=tmp.$$


echo "Generating $ADSBASENAME.tex"

# For some reason, must double-up \\ to \\\\ but not \ to \\ - odd!
cat > $ADSBASENAME.tex <<EOF
%% PLEASE DO NOT EDIT THIS FILE!  It is generated automatically from
%% the script makeads.sh.  Any hand-edited changes will be lost
%% 
\documentclass{article}
\usepackage[latin1]{inputenc}
\usepackage{html}
\usepackage{adshtml}
\newcommand{\ADSfield}[2]{\%#1  #2\\\\}
\newenvironment{ADSarticle}{}{\par\bigskip}
\setlength{\parindent}{0pt}
\newcounter{istage}
\renewcommand{\ion}[2]{\setcounter{istage}{#2}#1 \Roman{istage}}
\begin{document}
\makeatletter
EOF
echo >> $ADSBASENAME.tex "\input{$BASENAME.ads2}" 
echo >> $ADSBASENAME.tex "\end{document}"



echo "Generating $ADSBASENAME.sed" 

# Find type of journal
grep '\\documentclass\[.*proceedings.*\]' $BASENAME.tex >/dev/null \
     && JOURNAL='RMxAC' \
     || JOURNAL='RMxAA'

# Find the year and volume number
YEAR=`grep '\\SetYear' $BASENAME.tex | sed 's/.*\\SetYear{\(.*\)}.*/\1/'`
VOLUME=`grep '\\SetVolume' $BASENAME.tex | sed 's/.*\\SetVolume{\(.*\)}.*/\1/'`

# Generate an auxilliary sed script 
cat > $TMPFILE <<EOF
s|\\\\newlabel{\(.*\):RMAAFirstPage}{{.*}{\(.*\)}}|s/\\\\\\\\ThisVolume{\1}/$YEAR, RMxAC, $VOLUME, \2/g|g
EOF

# Now go through the .aux file, looking for \newlabels and use the sed 
# script just generated in order to make another sed script, one that 
# will replace \ThisVolume{LABEL} with the correct reference
grep '^\\newlabel{.*:RMAAFirstPage}' $AUXFILE \
     | sed -f $TMPFILE > $ADSBASENAME.sed && rm $TMPFILE

# Apend a further sed command to get rid of \@, since redefining 
# it in _ads.tex didn't work
echo >> $ADSBASENAME.sed 's/\\\\@//g'

echo "Generating $BASENAME.ads2 from $ADSBASENAME.sed and $BASENAME.ads" 

# Apply this to the .ads file, generating the .ads2 file, which is included
# by the _ads.tex file.
sed -f $ADSBASENAME.sed $BASENAME.ads > $BASENAME.ads2

echo "Generating $ADSBASENAME.txt from $ADSBASENAME.tex and $BASENAME.ads2"

# Now run latex2html on the _ads.tex file
latex2html -no_navigation -info 0 -address 0 -html_version 4.0,math $ADSBASENAME

# Then, convert to utf8 with lynx and tidy it up a bit
lynx --justify=off --width=100 --dump --nolist --display_charset=utf-8 \
     $ADSBASENAME/$ADSBASENAME.html | cut -c4- \
     | sed 's/ibidrule/\\ibidrule/g' > $ADSBASENAME.txt
%</ads-shell-script>
%    \end{macrocode}

% \section{Epilogue}
% Err, that's it \dots

\endinput


% \endinput
% Local Variables:
% mode: doctex
% TeX-master: t
% End:
"

rmaa.ind: rmaa.idx
        @echo "REGENERATING INDEX"
        makeindex -s gind.ist rmaa

rmaa.gls: rmaa.glo
        @echo "REGENERATING CHANGELOG"
        makeindex -s gglo.ist -o rmaa.gls rmaa.glo

rmaa.pdf: rmaa.aux rmaa.ind rmaa.gls
        @echo "REGENERATING .pdf FILE"
        pdflatex "\nonstopmode% \def\fileversion{4.0} \def\filedate{2019/11/12}
% \MakeShortVerb{\|}
% \title{A \LaTeX{} document class for the Revista Mexicana de
% Astronom\'\i a y Astrof\'\i sica\thanks{This file has version
% \fileversion{} last revised \filedate{}.}}
% \author{William J. Henney\thanks{w.henney@irya.unam.mx}}
% \date{\filedate}
% \maketitle
% \paragraph{Abstract}
% This document describes the \texttt{rmaa} \LaTeX{}
% document class for the journal ``Revista Mexicana de
% Astronom\'\i a y Astrof\'\i sica''. The macros can be used for 
% \begin{enumerate}
% \item Preparing author manuscripts for the journal (including
% conference proceedings to be published in the ``Serie de
% Conferencias'')
% \item Typesetting an entire volume of the journal by inputting the
% individual article files from a container file. Facilities are
% included for automatic generation of TOC and Author Index. 
% \end{enumerate}
% 
% \tableofcontents
% \section{The documentation driver file}
% This is the part that generates the documentation that you may or
% may not be reading, which is done with the help of the Makefile
% discussed in section~\ref{sec:driver_make}. 
%    \begin{macrocode}
%<*driver>
\documentclass{ltxdoc}
\usepackage{hyperref}
\usepackage{palatino}
\EnableCrossrefs \CodelineIndex \RecordChanges
\setlength{\parindent}{0pt}
\def\MacroFont{\footnotesize\tt}
\begin{document} \DocInput{rmaa.dtx} \PrintIndex \PrintChanges \end{document}
%</driver>
%    \end{macrocode}

% \section{The implementation}
% 
% \changes{3.24}{2007/09/14}{Added user access to the macro version
% and date via \cs{rmVersion} and \cs{rmDate}}
% 
%    \begin{macrocode}
%<*package>
\newcommand{\rmaa@version}{3.26}
\newcommand{\rmaa@date}{2011/08/15}
\newcommand\rmVersion{\rmaa@version}
\newcommand\rmDate{\rmaa@date}
\NeedsTeXFormat{LaTeX2e}[1995/06/01]
%    \end{macrocode}
% \changes{4.0}{2019/11/02}{Change abstract to appear before
% resumen. Change RevMexAA to RMxAA for consistency with ADS.}
%
% \changes{3.22}{2007/03/22}{Allow proceedings articles to have
%   received/accepted dates} \changes{3.21}{2006/05/15}{Changes for
%   Silvia - abstracts are now shipped out to ADS.}
% \changes{3.19}{2006/01/19}{Changes for Silvia - remove adresses from
%   end, new abstract style} \changes{2.2beta}{2001/10/30}{ Fixed bug
%   whereby twocolumn option was not passed through to book
%   class. Finally sorted out the ``label on last page'' issue. }
% \changes{2.1}{2001/01/23}{ Fixed some inconsistencies in
%   \cs{nocite} and \cs{bibitem} uniqification Fixed the definition of
%   \cs{ion} so it will work in math mode and in contexts where the
%   typesize is different from normal (e.g., title, fig captions).
%   Modified the definition of \cs{@caption}, so we don't write any
%   list-of-figures data (we never use it anyway, and it just causes
%   problems with fragile commands in fig captions). }
% \changes{2.0}{2000/11/30}{ Added `newstyle' option, which should be
%   used in all new proceedings volumes (wider abstract, smaller
%   typeface in fig captions, etc.) Added commands for formatting
%   table captions. New demonstation document rmpostsamp.tex for
%   1-page posters. This will be the last release under a PRI
%   government (for a while...)}  \changes{1.9}{2000/06/06}{
%   Corrected position of page number in footer of articles that open
%   on even pages. Ensured that first article in each session opens on
%   an odd page, as do list of abstracts and author index. Added extra
%   space in TOC before the backmatter. Added command for editors,
%   \cs{BoldSessionTitles}, that makes produces bold titles for
%   session in the TOC.}  \changes{1.8}{2000/05/16}{ Fixed bug in
%   last page label with \cs{suppressfulladdresses} }
% \changes{1.7}{2000/05/11}{ Fixed more little things. Page labels
%   now work, even if you have no bibliography. You can output the
%   full addresses by hand with the \cs{outputfulladdresses}
%   command. Editors' Guide updated but still not finished. Revamped
%   splitbook.sh now produces PS file of each article for web page and
%   also an EPS file of each page for sending to the printers.}
% \changes{1.6}{2000/04/15}{ More cleaning up: full stops after
%   titles in index, running headers and footers in frontmatter,
%   fontsize in title of index, etc} \changes{1.5}{2000/04/15}{
%   Eliminated superscripts from titles in TOC and index}
% \changes{1.4}{2000/04/14}{ Added copyright and license info (GPL)}
% \changes{1.3}{2000/03/01}{ Corrected \cs{figcaption} to do the
%   caption formatting right Corrected font size of sec headings in
%   front matter} \changes{1.2}{2000/02/16}{ Added \cs{plate}
%   command} \changes{1.1}{2000/02/14}{ Added manuscript option,
%   corrected bugs} \changes{1.0}{2000/02/08}{ Now does a whole book
%   - lots of stuff tidied up} \changes{0.7}{1999/10/04}{ Increased
%   spacing between sections in single column styles. Suppressed
%   floats at bottom of pages (didn't work though) }
% \changes{0.6}{1999/09/24}{ Tweaked \cs{adjustfinalcols} (still not
%   perfect) Added proceedings option Added onecolumn/twocolumn
%   options (orthogonal to journal/proceedings) }
% \changes{0.5}{1999/09/23}{ Added \cs{fulladdresses} and
%   \cs{adjustfinalcols} (just about working)}
% \changes{0.4}{1999/09/22}{ Minor stylistic changes}
% \changes{0.3}{1999/06/28}{ Better mechanism for alternative
%   affiliations Added `guide' option for user guide}
% \changes{0.2}{1999/05/27}{ All basic functionality present, I
%   think} \changes{0.1}{1999/05/26}{ Just about working - still
%   missing Fig/Tab/Refs} \changes{0.0}{1999/05/21}{ So bored today,
%   decided to rewrite RMAA LaTeX macros from scratch}
%    \begin{macrocode}
\ProvidesClass{rmaa}[\rmaa@date\space v\rmaa@version\space RMxAA article
class]
%    \end{macrocode}
% \subsection{Option Handling}
% \label{sec:options}
% Declarations that are needed in options
% \changes{3.05}{2003/10/22}{Added natbib option}
% \changes{3.06}{2003/01/03}{Natbib option turned on permanently}
% \changes{3.07}{2003/01/11}{Added proof option}
% \changes{3.14}{2003/01/11}{Option to turn off natbib}
% \changes{3.18}{2003/10/22}{trying out doctex mode} 
% \changes{3.23}{2007/09/09}{Manuscript is now on by default.}
% \changes{3.26}{2011/08/15}{New option for review articles}
%    \begin{macrocode}
\newif\if@reviewarticle
\@reviewarticlefalse
\newif\if@proof
\@prooffalse
\newif\if@journal
\@journaltrue
\newif\if@printdates
\@printdatesfalse
\newif\if@proceedings
\@proceedingsfalse
\newif\if@manuscript
\@manuscripttrue
\newif\if@bigfont
\@bigfonttrue
\newif\if@guide
\@guidefalse
\newif\if@preprint
\@preprintfalse
\newif\if@rmtwocol
\@rmtwocoltrue
\newif\if@biblist
\@biblistfalse
\newif\if@book
\@bookfalse
\newif\if@natbib
\@natbibtrue
\newif\if@rmcite
\@rmcitefalse
\newif\if@titlevisible
\@titlevisibletrue
\newif\if@abstractpresent
\@abstractpresentfalse
\newif\if@printbanner
\@printbannerfalse
\newif\if@suppressfulladdresses
\@suppressfulladdressesfalse
\newif\if@newstyle
\@newstylefalse
\newif\if@debug
\@debugfalse
\newif\if@boldkeywords
\@boldkeywordsfalse
%    \end{macrocode}
% The option handling---mainly setting various booleans.
% \changes{2.2beta11}{2001/11/26}{Made twocol and
% newstyle the default for proceedings.}
% \changes{3.23}{2007/09/09}{Manuscript is now 10 pt instead of 12pt}
% \changes{3.23}{2007/09/10}{Full addresses are now automatically
% suppressed for proceedings style.}
% \changes{3.24}{2007/09/14}{New option \texttt{boldkeywords} that
% does the keywords in the old ugly way. Now off by default.}
%    \begin{macrocode}
\DeclareOption{manuscript}{\typeout{Manuscript for referee} 
  \@manuscripttrue\@journalfalse\@rmtwocolfalse
  \PassOptionsToClass{10pt}{article}}
\DeclareOption{bigfont}{\typeout{Big font manuscript requested by Gary
  Ferland} 
  \@manuscripttrue\@journalfalse\@rmtwocolfalse\@bigfonttrue
  \PassOptionsToClass{11pt}{article}}
\DeclareOption{journal}{\typeout{Main journal option}%
  \@journaltrue\@rmtwocoltrue\@printdatestrue\@manuscriptfalse
  \PassOptionsToClass{10pt}{article}}
\DeclareOption{guide}{\typeout{User guide option}%
  \@guidetrue\@journalfalse\@rmtwocoltrue\@manuscriptfalse
  \PassOptionsToClass{10pt}{article}}
\DeclareOption{preprint}{\typeout{Preprint option}%
  \@preprinttrue\@manuscriptfalse\PassOptionsToClass{10pt}{article}}
\DeclareOption{proof}{\typeout{Author proof option}%
  \@prooftrue\@manuscriptfalse}
\DeclareOption{proceedings}{\typeout{Proceedings option}%
  \@proceedingstrue\@journalfalse\@rmtwocoltrue\@newstyletrue
  \@suppressfulladdressestrue
  \@printdatesfalse\@manuscriptfalse\PassOptionsToClass{10pt}{article}} 
\DeclareOption{printdates}{\typeout{Force printing of received/accepted dates}%
  \@printdatestrue} 
\DeclareOption{book}{\typeout{Whole issue option}%
  \@booktrue} 
\DeclareOption{nonatbib}{\typeout{Turn off natbib support (required
    for compatibility with old rmaacite package).}%
  \@natbibfalse} 
\DeclareOption{newstyle}{\typeout{Use new layout style}%
  \@newstyletrue}
\DeclareOption{oldstyle}{\typeout{Use old layout style}%
  \@newstylefalse}
\DeclareOption{debug}{\typeout{Loads of annoying messages}%
  \@debugtrue}
%    \end{macrocode}
% These two need to be declared explicitly so that |\ExecuteOptions|
% actually does something.
%    \begin{macrocode}
\DeclareOption{twocolumn}{\@rmtwocoltrue}
\DeclareOption{onecolumn}{\@rmtwocolfalse}
\DeclareOption{twoside}{\PassOptionsToClass{\CurrentOption}{article}}
\DeclareOption{boldkeywords}{\@boldkeywordstrue}
%    \end{macrocode}
% Bold, upper-case keywords used to be the default. Now they must be
% turned on explicitly with the \texttt{boldkeywords} option. 
%    \begin{macrocode}
\DeclareOption{boldkeywords}{\@boldkeywordstrue}
%    \end{macrocode}
% Pass any other option along to the \texttt{article} class.
%    \begin{macrocode}
\DeclareOption*{\PassOptionsToClass{\CurrentOption}{article}}
%    \end{macrocode}
% Default options are to be two-sided and based on the standard
% \LaTeX{} \texttt{article} class,
% unless we are processing the entire book, in which case we base
% ourselves on the standard \texttt{book} class. Default is to use
% manuscript option unless otherwise specified. 
%    \begin{macrocode}
\ExecuteOptions{twoside,manuscript} \ProcessOptions\relax
\if@book
\if@rmtwocol\LoadClass[twocolumn]{book}\else\LoadClass{book}\fi
\newcommand{\abstract}{\relax}
\newcommand{\@maketitle}{\relax}
\else
\if@rmtwocol\LoadClass[twocolumn]{article}\else\LoadClass{article}\fi
\fi
%    \end{macrocode}
% Attempt to use natbib for citations 02 Jan 2003. 
% \changes{3.05}{2003/01/02}{If the natbib option is used, then set
% things up so that natbib's version of the bibliography is identical
% to the standard one.}
% \changes{3.09}{2003/01/23}{Explicitly require a recent version of
% natbib, so that we get a decent error message if the package is too
% old.} 
% \changes{3.23}{2007/09/09}{Now allow Bibtex to work properly} 
% \changes{3.25}{2009/10/01}{Change \cs{newcommand} tp \cs{def} when setting \cs{bibfont}. This allows it to work with all versions of natbib.} 
% \changes{3.25}{2010/02/06}{Updated copyright year.} 
% \changes{3.25}{2010/02/06}{Updated email and web addresses to \texttt{crya.unam.mx}.} 
%    \begin{macrocode}
\if@natbib
\RequirePackage{natbib}[1999/02/25]
\bibpunct{(}{)}{;}{a}{}{,}
\def\bibfont{\small}
\renewcommand\bibsection{\section*{REFERENCES}}
\setlength\bibsep{0pt}
\setlength\bibhang{1.5em}
\bibliographystyle{rmaa}
\fi
%    \end{macrocode}
% We want the \texttt{graphicx} package for the extended syntax of
% |\includegraphics|. We want the \texttt{ifthen} package for using in
% these macros (although we don't much yet). We explicitly load
% \texttt{babel} since some \LaTeX{} distributions messed up the
% hyphenation if we didn't.  
%    \begin{macrocode}
\RequirePackage{graphicx}
\RequirePackage{ifthen}
\IfFileExists{babel.sty}{%
%   \let\@restoreindent=\@afterindentfalse
  \RequirePackage[english]{babel}
%   \@restoreindent
  %% reinstate the definition of \@roman so as not to break .idx files
  \@ifundefined{org@@roman}{\relax}{\renewcommand{\@roman}[1]{\org@@roman{##1}}}
  }{%
  \typeout{Package babel.sty not found - don't worry about this}}
%    \end{macrocode}
% Try to load the \texttt{dcolumn} and \texttt{booktabs} packages if
% they can be found. These can produce nicer-looking tables. 
%    \begin{macrocode}
\IfFileExists{dcolumn.sty}{%
  \RequirePackage{dcolumn}
  }{%
  \typeout{Package dcolumn.sty not found - you will not be able to
    align columns on decimal points in tables}}
\IfFileExists{booktabs.sty}{%
  \RequirePackage{booktabs}
  }{%
  \typeout{Package booktabs.sty not found - your tables will not look
    as pretty as they ought too.}
  }
\IfFileExists{makeidx.sty}{%
  \RequirePackage{makeidx}
}{%
  \typeout{Package makeidx.sty not found - Author Index may not work.}
  }
%% commented out for moment since requires more modern version of LaTeX
%% \RequirePackage[latin1]{inputenc}
%    \end{macrocode}
% \changes{4.0}{2019/11/13}{Make \texttt{hyperref} and
% \texttt{etoolbox} into hard requirements, since we now use them
% extensively.}  We use hyperref to get internal and external links.
% But I don't like the colored boxes so I turn those off.  We use
% etoolbox for patching and robustifying at the moment.  In the
% future, I want to port a lot of other stuff over to it.  For
% instance, replace the ifthen package and the TeX ifs too.
%
% We target the latest version (1.7) of PDF, rather than the default
% of 1.5.  This reduces the chance of getting warnings from included
% graphics.  This is not strictly related to hyperref, but we put it
% here anyway.
%    \begin{macrocode}
\RequirePackage{etoolbox}
\RequirePackage{hyperref}
\hypersetup{hidelinks}
\pdfminorversion=7
%    \end{macrocode}
% We patch the \cs{@combinedblfloats} following a suggestion by Heiko
% Oberdiek at \url{https://tex.stackexchange.com/a/249743/3843}, which
% avoids problems when hyperlinks are split across
% pages. \changes{4.0}{2019/11/12}{Fix for issue where breaking
% hyperlink across pages gave an error.}
%    \begin{macrocode}
\urlstyle{same}
\newcommand\email{\begingroup \urlstyle{rm}\Url}
\patchcmd\@combinedblfloats{\box\@outputbox}{\unvbox\@outputbox}{}{%
  \errmessage{\noexpand\@combinedblfloats could not be patched}%
}%
%    \end{macrocode}
% \subsection{Other setup stuff}\label{sec:setup}
% Some declarations: 
%    \begin{macrocode}
\newif\if@needaddresses
\@needaddressestrue
\newif\if@frontmatter
\@frontmatterfalse
\newif\if@inpreamble
\@inpreambletrue
\newif\if@firstauthor
\@firstauthortrue
\newif\if@firstaffil
\@firstaffiltrue
\newif\if@firstaltaffil
\@firstaltaffiltrue
\newif\if@firstkwd
\@firstkwdtrue
\newif\if@noexpandADS
\@noexpandADSfalse
\newif\if@sessiontitlepages
\@sessiontitlepagesfalse
%    \end{macrocode}
% \begin{macro}{\@rmaaEmpty}
% \begin{macro}{\@rmaaArg}
% These are used in various commands to check if arg was left blank
%    \begin{macrocode}
\newcommand{\@rmaaEmpty}{}
\newcommand{\@rmaaArg}{}
%    \end{macrocode}  
% \end{macro}
% \end{macro}
%
% \subsection{Compatibility with old \LaTeX{} versions}
% \label{sec:compat}
% This package is based on the 1998/12/01 release of \LaTeX{}.
% Older releases may not define all the commands that we need.
% Luckily, we can use |\providecommand| to define these if necessary. 
% The following commands are not defined in the 1994/06/01 release. 
%
% \begin{macro}{\MakeUpperCase}
% \begin{macro}{\textsuperscript}
% \begin{macro}{\textbackslash}
% User commands:
%    \begin{macrocode}
\providecommand{\MakeUpperCase}[1]{\uppercase{#1}}
\providecommand{\textsuperscript}[1]{\ensuremath{^\textrm{#1}}}
\providecommand{\textbackslash}{\ensuremath{\backslash}}
%    \end{macrocode}  
% \end{macro}
% \end{macro}
% \end{macro}
% \begin{macro}{\protected@xdef}
% \begin{macro}{\restore@protect}
% \begin{macro}{\@unexpandable@protect}
% \begin{macro}{\hb@xt@}
% \LaTeX{} internal commands:
%    \begin{macrocode}
\providecommand{\protected@xdef}{%
        \let\@@protect\protect
        \let\protect\@unexpandable@protect
        \afterassignment\restore@protect
        \xdef}
\providecommand{\restore@protect}{\let\protect\@@protect}
\providecommand{\@unexpandable@protect}{\noexpand\protect\noexpand}
\providecommand{\hb@xt@}{\hbox to}
%    \end{macrocode}
% \end{macro}
% \end{macro}
% \end{macro}
% \end{macro}
%
% \subsection{Page layout}\label{sec:page-layout}
%
% The dimensions of the text were defined by Elizabeth. Text width is
% 169\,mm with equal margins on the two sides: 
%    \begin{macrocode}
\setlength{\textwidth}{169mm}
\setlength{\@tempdima}{\paperwidth}
\addtolength{\@tempdima}{-\textwidth}
\setlength{\oddsidemargin}{0.5\@tempdima}
\addtolength{\oddsidemargin}{-1in}
\setlength{\evensidemargin}{\oddsidemargin}
%    \end{macrocode}
% Text height is 225\,mm, top margin is 14\,mm. 
%    \begin{macrocode}
\setlength{\textheight}{225mm} 
\setlength{\topmargin}{-1in} 
\addtolength{\topmargin}{14mm}
%    \end{macrocode}
% Space between columns in two-column styles, header, footer,
% paragraph indentation:
%    \begin{macrocode}
\setlength{\columnsep}{9.5mm}
\setlength{\headsep}{0.7cm}
\setlength{\footskip}{0.75cm}
\setlength{\parindent}{1.5em}
%    \end{macrocode}
% 
% \begin{macro}{\@set@scalable@length}
% Try out a new scheme for dealing with various spacing dimensions and
% allowing the user to easily modify them. The command
% \begin{quote}
% |\set@scalable@length|\marg{name}\marg{defaultlength}\marg{stretch}\marg{shrink}
% \end{quote}
% does a 
% |\setlength|\marg{\cs{name}}\marg{defaultlength plus stretch minus shrink}
% and also defines the user command
% \textbackslash\meta{name}|Scale|\marg{factor} that can 
% be used to rescale length by \meta{factor}. Additionally, now
% defines the package command
% \textbackslash\meta{name}|@set@scale|\marg{FACTOR} 
% that also rescales the length by \meta{FACTOR}, so that the length
% becomes \meta{factor} $\times$ \meta{FACTOR} $\times$
% \meta{defaultlength}.   
%    \begin{macrocode}
\newcommand\@set@scalable@length[4]{%
  %% arg #1: NAME of length without leading `\'
  %%
  %% arg #2: default VALUE  
  \expandafter\newlength\csname @#1@default\endcsname
  \expandafter\setlength\csname @#1@default\endcsname{#2}%
  \expandafter\newlength\csname @@#1@@default\endcsname
  \expandafter\setlength\csname @@#1@@default\endcsname{#2}%
  %% arg #3: default STRETCH  
  \expandafter\newlength\csname @#1@stretch\endcsname
  \expandafter\setlength\csname @#1@stretch\endcsname{#3}%
  \expandafter\newlength\csname @@#1@@stretch\endcsname
  \expandafter\setlength\csname @@#1@@stretch\endcsname{#3}%
  %% arg #4: default SHRINK  
  \expandafter\newlength\csname @#1@shrink\endcsname
  \expandafter\setlength\csname @#1@shrink\endcsname{#4}%
  \expandafter\newlength\csname @@#1@@shrink\endcsname
  \expandafter\setlength\csname @@#1@@shrink\endcsname{#4}%
  %%
  %% User-defined scale \@NAME@scale is unity by default
  \expandafter\gdef\csname @#1@scale\endcsname{1}%
  %% Package-defined scale \@NAME@SCALE is unity by default
  \expandafter\gdef\csname @#1@SCALE\endcsname{1}%
  %%
  %% internal command \@reset@NAME should expand to 
  %% \setlength\@@NAME@@default{\@NAME@SCALE\@NAME@default}
  %% \setlength\NAME{\@NAME@scale\@@NAME@@default}
  %% in other words: scale x SCALE x default 
  \expandafter\gdef\csname @reset@#1\endcsname{%
    \expandafter\setlength\csname @@#1@@default\endcsname{%
      \csname @#1@SCALE\endcsname \csname @#1@default\endcsname
      }
    \expandafter\setlength\csname @@#1@@stretch\endcsname{%
      \csname @#1@SCALE\endcsname \csname @#1@stretch\endcsname
      }
    \expandafter\setlength\csname @@#1@@shrink\endcsname{%
      \csname @#1@SCALE\endcsname \csname @#1@shrink\endcsname
      }
    \expandafter\setlength\csname #1\endcsname{%
      \csname @#1@scale\endcsname \csname @@#1@@default\endcsname
      \@plus
      \csname @#1@scale\endcsname \csname @@#1@@stretch\endcsname
      \@minus
      \csname @#1@scale\endcsname \csname @@#1@@shrink\endcsname
      }}
  %%
  %% User command \NAMEScale should expand to 
  %% {\renewcommand\@NAME@scale{#1}\@reset@NAME}
  \expandafter\gdef\csname #1Scale\endcsname##1{%
    \expandafter\gdef\csname @#1@scale\endcsname{##1}%
    \csname @reset@#1\endcsname}
  %%
  %% Package command \NAME@set@SCALE should expand to 
  %% {\renewcommand\@NAME@SCALE{#1}\@reset@NAME}
  \expandafter\gdef\csname #1@set@SCALE\endcsname##1{%
    \expandafter\gdef\csname @#1@SCALE\endcsname{##1}%
    \csname @reset@#1\endcsname}
  %%
  %% initialize at default values with \@reset@NAME
  \csname @reset@#1\endcsname
  }
%    \end{macrocode}
% \end{macro}
% \begin{macro}{\parskip}
% First use of this new infrastructure: set interparagraph space to be
% zero but with a small amount of stretchiness.
%    \begin{macrocode}
%   \ifthenelse{\boolean{@manuscript}}{%
%     \@set@scalable@length{parskip}{0.5\baselineskip}{0.3\p@}{0.1\p@}
%   }{
    \@set@scalable@length{parskip}{0\p@}{0.3\p@}{0\p@}
%   }
%    \end{macrocode}
% \end{macro}
% 
% \begin{macro}{\abovesectionskip}
% \changes{4.0}{2019/11/12}{Use negative values for skip so that first paragraph after section headings is not indented.  This is something that I have been trying to fix for years!  It turns out that the anser was in the documentation for the standard command \cs{@startsection.}}
% \begin{macro}{\belowsectionskip}
% \begin{macro}{\aboveparaskip}
% \begin{macro}{\belowparaskip}
% Define some utility rubber lengths for use in
% section~\ref{sec:section-commands}. 
%    \begin{macrocode}
\newlength{\abovesectionskip}
\newlength{\belowsectionskip}
\newlength{\aboveparaskip}
\newlength{\belowparaskip}
%    \end{macrocode}
% The default values for these lengths are those traditionally used
% for the onecolumn style. 
%    \begin{macrocode}
\@set@scalable@length{abovesectionskip}{-3.5ex}{-4ex}{-2ex}
\@set@scalable@length{belowsectionskip}{2.3ex}{0ex}{2ex}
\@set@scalable@length{aboveparaskip}{3.25ex}{2ex}{0.5ex}
\@set@scalable@length{belowparaskip}{-1em}{\z@}{\z@}
%    \end{macrocode}
% \end{macro}\end{macro}\end{macro}\end{macro}
% We also make scalable some of the standard float parameters. These
% are the same defaults as in the standard article classes, but
% slightly more stretchy. For the twocolumn style, they are currently
% halved, see below. Note that I have also lumped |\abovecaptionskip|
% in with these, since it plays a similar role. There is also a
% |\belowcaptionskip| but it is set to zero in the standard classes. 
%    \begin{macrocode}
\@set@scalable@length{floatsep}{12\p@}{4\p@}{4\p@}
\@set@scalable@length{textfloatsep}{20\p@}{4\p@}{8\p@}
\@set@scalable@length{intextsep}{12\p@}{4\p@}{4\p@}
\@set@scalable@length{dblfloatsep}{12\p@}{4\p@}{4\p@}
\@set@scalable@length{dbltextfloatsep}{20\p@}{4\p@}{8\p@}
\@set@scalable@length{abovecaptionskip}{10\p@}{\z@}{\z@}
%    \end{macrocode}
% \begin{macro}{\@before@title@skip}         
% \begin{macro}{\@before@author@skip}        
% \begin{macro}{\@between@author@skip}       
% \begin{macro}{\@before@abstract@skip}      
% \begin{macro}{\@after@abstracthead@skip}   
% \begin{macro}{\@between@abstract@skip}     
% \begin{macro}{\@before@kwd@skip}           
% \begin{macro}{\@before@dates@skip}           
% Spacing parameters used in |\@maketitle|. These are not changed by
% default except by the user, so they don't need to go into
% |\@package@default@lengths|.  
%    \begin{macrocode}
\newlength{\@before@title@skip}         
\newlength{\@before@author@skip}        
\newlength{\@between@author@skip}       
\newlength{\@before@abstract@skip}      
\newlength{\@after@abstracthead@skip}   
\newlength{\@between@abstract@skip}     
\newlength{\@before@kwd@skip}           
\newlength{\@before@dates@skip}           
\@set@scalable@length{@before@title@skip}{2em}{\z@}{\z@}
\@set@scalable@length{@before@author@skip}{1.5em}{\z@}{\z@}
\@set@scalable@length{@between@author@skip}{0.5em}{\z@}{\z@}
\@set@scalable@length{@before@abstract@skip}{3ex}{1ex}{0.5ex}
\@set@scalable@length{@after@abstracthead@skip}{1ex}{\z@}{\z@}
\@set@scalable@length{@between@abstract@skip}{2ex}{1ex}{0.5ex}
\@set@scalable@length{@before@kwd@skip}{1ex}{0.5ex}{0.2ex}
\@set@scalable@length{@before@dates@skip}{2ex}{0.5ex}{0.2ex}
%    \end{macrocode}
% \end{macro}\end{macro}\end{macro}\end{macro}
% \end{macro}\end{macro}\end{macro}\end{macro}
% 
% Spacing above and below displayed equations also wants to be done
% like this, but it is more complicated since the standard article
% classes make these change with the fontsize. Something to do later,
% maybe. 
% 
%    \begin{macrocode}
%    \end{macrocode}
% \begin{macro}{\@package@default@lengths}
% The value actually used for these lengths is controlled by the
% \meta{SCALE} parameter, which depends on the style option that is in
% effect. This is set using the |\@package@default@lengths| command,
% which must be called every time the style changes (e.g., changing
% from onecolumn to twocolumn). 
%    \begin{macrocode}
\newcommand{\@package@default@lengths}{%
  \ifthenelse{\boolean{@rmtwocol}}{%
    %% twocolumn mode
    \abovesectionskip@set@SCALE{0.5}
    \belowsectionskip@set@SCALE{0.5}
    \aboveparaskip@set@SCALE{0.5}
    \belowparaskip@set@SCALE{1}
    \floatsep@set@SCALE{0.5}
    \textfloatsep@set@SCALE{0.5}
    \intextsep@set@SCALE{0.5}
    \dblfloatsep@set@SCALE{0.5}
    \dbltextfloatsep@set@SCALE{0.5}
    }{%
    %% onecolumn mode
    \abovesectionskip@set@SCALE{1}
    \belowsectionskip@set@SCALE{1}
    \aboveparaskip@set@SCALE{1}
    \belowparaskip@set@SCALE{1}
    \floatsep@set@SCALE{1}
    \textfloatsep@set@SCALE{1}
    \intextsep@set@SCALE{1}
    \dblfloatsep@set@SCALE{1}
    \dbltextfloatsep@set@SCALE{1}
    }%
  }
\@package@default@lengths
%    \end{macrocode}
% \end{macro}
% \begin{macro}{\RescaleLengths}
% \begin{macro}{\RescaleSecLengths}
% \begin{macro}{\RescaleFloatLengths}
% \begin{macro}{\RescaleTitleLengths}
% User commands that will fiddle with all the scalable lengths. 
%    \begin{macrocode}
\newcommand{\RescaleLengths}[1]{%
  \RescaleSecLengths{#1}\RescaleFloatLengths{#1}\RescaleTitleLengths{#1}}
\newcommand{\RescaleSecLengths}[1]{%
%% sections
  \abovesectionskipScale{#1}
  \belowsectionskipScale{#1}
  \aboveparaskipScale{#1}
  }
\newcommand{\RescaleFloatLengths}[1]{%
%% floats
  \floatsepScale{#1}
  \textfloatsepScale{#1}
  \intextsepScale{#1}
  \dblfloatsepScale{#1}
  \dbltextfloatsepScale{#1}
  \abovecaptionskipScale{#1}
  }
\newcommand{\RescaleTitleLengths}[1]{%
%% maketitle
  \@before@title@skipScale{#1}           
  \@before@author@skipScale{#1}        
  \@between@author@skipScale{#1}       
  \@before@abstract@skipScale{#1}      
  \@after@abstracthead@skipScale{#1}   
  \@between@abstract@skipScale{#1}     
  \@before@kwd@skipScale{#1}           
  \@before@dates@skipScale{#1}           
  }
%    \end{macrocode}
% \end{macro}\end{macro}\end{macro}\end{macro}
% \changes{2.2beta14}{2002/03/17}{Finished implementation of new
% infrastructure for scalable lengths. Converted the above/below
% section skips to use it. Also, all the floatsep parameters. Also,
% all the vertical lengths that are used by |\@maketitle|. }

% \begin{macro}{\MSextramargin}
% For the \texttt{manuscript} option we have an extra 0.3\,inch
% vertical margin and 1\,inch horizontal margin for marginal notes (by
% hand or with |\marginpar|. 
%
% \changes{3.23}{2007/09/09}{Reduce baselinestretch in
%   \texttt{manuscript} style to 1.2. No longer use raggedright for
%   manuscript.}
%
% \changes{3.24}{2007/09/14}{Further reduction of baselinestretch in
%   \texttt{manuscript} style to 1.1 since we now have wider margins.}
%
% \changes{3.24}{2007/09/14}{Split \texttt{manuscript}'s extra margin
%   into separate horizontal and vertical parts.}
% 
%    \begin{macrocode}
\newlength{\MSextraHmargin}\newlength{\MSextraVmargin}
\setlength{\MSextraHmargin}{1in} 
\setlength{\MSextraVmargin}{0.3in} 
\if@manuscript
   \renewcommand{\baselinestretch}{1.1}
   \addtolength{\textwidth}{-2\MSextraHmargin}
   \addtolength{\oddsidemargin}{\MSextraHmargin}
   \setlength{\evensidemargin}{\oddsidemargin}
   \addtolength{\textheight}{-2\MSextraVmargin}
   \addtolength{\topmargin}{\MSextraVmargin}
%    \raggedright
\fi
%    \end{macrocode}
% \end{macro}
% Float parameters: don't allow bottom floats by default. 
%    \begin{macrocode}
\renewcommand{\fps@figure}{tp}
\renewcommand{\fps@table}{tp}
%    \end{macrocode}
%
% \begin{macro}{\rmaa@AuthorWidthFrac}
% \begin{macro}{\rmaa@AbstractWidthFrac}
% These parameters control what fraction of the total text width can
% be used by the list of authors and the abstract/resumen,
% respectively.  
%    \begin{macrocode}
\newcommand{\rmaa@AuthorWidthFrac}{1.0}
\newcommand{\rmaa@AbstractWidthFrac}{0.75}
%    \end{macrocode}
% Use a wide abstract to save space in the new style.
%    \begin{macrocode}
\if@newstyle\renewcommand{\rmaa@AbstractWidthFrac}{1.0}\fi
%    \end{macrocode}
% \end{macro}
% \end{macro}
%
% Some things to consider in future versions. Trying to avoid too many
% widows, etc. Apparently, the following will work best with
% \cs{flushbottom}. 
%    \begin{macrocode}
%%\widowpenalty=10000
%%\clubpenalty=10000
%    \end{macrocode}
% Also, should perhaps try getting rid of underfull/overfull box warnings. 
% \changes{3.10}{2003/01/24}{Set large values for \cs{hbadness},
% \cs{vbadness}, \cs{hfuzz}, \cs{vfuzz}. To reduce amount of stuff
% that gets written to the terminal every time you compile something.}
%    \begin{macrocode}
\hbadness=10000
\vbadness=10000
\setlength\hfuzz{\maxdimen}
\setlength\vfuzz{\maxdimen}
%    \end{macrocode}
% \subsection{Various text styles for different contexts}
% \label{sec:text-styles} 
% \begin{macro}{\rmaa@AffilStyle}
% \begin{macro}{\rmaa@PagesStyle}
% \begin{macro}{\rmaa@YearStyle}
% \begin{macro}{\rmaa@VolumeStyle}
% \begin{macro}{\rmaa@VolumenStyle}
% \begin{macro}{\rmaa@JournalNameStyle}
% \begin{macro}{\rmaa@ReceivedStyle}
% \begin{macro}{\rmaa@AbstractHeaderStyle}
% \begin{macro}{\rmaa@JournalNameStyle}
% Style commands for the different fields used by |\maketitle| in
% constructing the title page of the article. These have the form
% |\rmaa@XYZStyle{TEXT}|. I'm not sure these were really a good idea,
% but never mind \dots 
%    \begin{macrocode}
\newcommand{\rmaa@AffilStyle}[1]{{\normalsize\rmfamily#1}}
\newcommand{\rmaa@PagesStyle}[1]{{\normalsize\rmfamily#1}}
\newcommand{\rmaa@YearStyle}[1]{{\normalsize\rmfamily#1}}
\newcommand{\rmaa@VolumeStyle}[1]{{\normalsize\rmfamily\bfseries#1}}
\newcommand{\rmaa@VolumenStyle}[1]{{\Large\sffamily\bfseries#1}}
\newcommand{\rmaa@JournalNameStyle}[1]{{\normalsize\itshape#1}}
\newcommand{\rmaa@ReceivedStyle}[1]{{\normalsize\rmfamily\itshape#1}}
\newcommand{\rmaa@AbstractHeaderStyle}[1]{%
  {\normalsize\rmfamily\MakeUpperCase{#1}}}
\newcommand{\rmaa@TitleStyle}[1]{%
  {\large\bfseries\rmfamily\boldmath\MakeUpperCase{#1}\par}}
%    \end{macrocode}
% \end{macro}
% \end{macro}
% \end{macro}
% \end{macro}
% \end{macro}
% \end{macro}
% \end{macro}
% \end{macro}
% \end{macro}
%
% \begin{macro}{\rmaa@FrontSecStyle}
% Headings of ``front sections'' (see section~\ref{sec:front-sections})
%    \begin{macrocode}
\newcommand{\rmaa@FrontSecStyle}[1]{%
  {\normalsize\bfseries\rmfamily\MakeUpperCase{#1}}}
%    \end{macrocode}
% \end{macro}
%
% \begin{macro}{\rmaa@TOCTitleStyle}
% \begin{macro}{\rmaa@TOCPartStyle}
% \begin{macro}{\rmaa@IndexTitleStyle}
% \begin{macro}{\rmaaTOCAuthorStyle}
% \begin{macro}{\rmaaIndexAuthorStyle}
% Commands to set style of TOC and index entries.
% Rewritten to avoid unnecessary font commands in the .toc and .ind files
%    \begin{macrocode}
\newcommand{\rmaa@TOCTitleStyle}[1]{\MakeUpperCase{#1}}
\newcommand{\rmaa@TOCPartStyle}[1]{\rmaa@TOCTitleStyle{#1}}
\newcommand{\rmaa@IndexTitleStyle}[1]{#1}
\newcommand{\rmaaTOCAuthorStyle}[1]{%
  {\itshape#1}}
\newcommand{\rmaaIndexAuthorStyle}[1]{%
  \protect\mbox{\bfseries#1}}
%    \end{macrocode}
% \end{macro}
% \end{macro}
% \end{macro}
% \end{macro}
% \end{macro}
%
% \begin{macro}{\rmaa@KeywordStyle}
% \begin{macro}{\rmaa@KeywordHeaderStyle}
% The style of the keywords depends on whether we are using the
% \texttt{newstyle} option or not. If we are, then they are small. 
%    \begin{macrocode}
\ifthenelse{\boolean{@newstyle}}{
  \newcommand{\rmaa@KeywordStyle}[1]{%
    {\small%
      \if@boldkeywords 
      \bfseries\MakeUpperCase{\raggedright #1}
      \else
      \raggedright #1
      \fi
    }}
  \newcommand{\rmaa@KeywordHeaderStyle}[1]{%
    {\small\itshape#1}}}
{
  \newcommand{\rmaa@KeywordStyle}[1]{%
    {\normalsize%
      \if@boldkeywords\bfseries\fi\raggedright #1}}
  \newcommand{\rmaa@KeywordHeaderStyle}[1]{%
    {\normalsize\itshape#1}}}
%    \end{macrocode}
% \end{macro}
% \end{macro}
%
% \subsection{Page styles}\label{sec:page-styles}
% \begin{macro}{\ps@rmaa}
% The running headers for normal pages:
%    \begin{macrocode}
\newcommand{\ps@rmaa}{%
%    \end{macrocode}
% nothing in either foot; 
%    \begin{macrocode}
  \renewcommand\@oddfoot{\relax}%
  \renewcommand\@evenfoot{\relax}%
%    \end{macrocode}
% page number and short author list in the even heads; 
%    \begin{macrocode}
  \renewcommand\@evenhead{\rlap{\thepage}\hfil\@shortauthor\hfil}%
%    \end{macrocode}
% short title and page number in the odd heads.
%    \begin{macrocode}
  \renewcommand\@oddhead{\hfil\@shorttitle\hfil\llap{\thepage}}
  }
%    \end{macrocode}
% \begin{macro}{\ps@rmaafirstpage}
% The header of the first page of each
% article:
%    \begin{macrocode}
\newcommand{\ps@rmaafirstpage}{%
  \renewcommand\@oddfoot{\hfil\thepage}%
%    \end{macrocode}
% footer corrected (06 Jun 2000) to have page on outside edge;
%    \begin{macrocode}
  \renewcommand\@evenfoot{\thepage\hfil}%
%    \end{macrocode}
% Header has journal name, volume, and page numbers for a regular
% article, or similar for special cases. 
% \changes{4.0}{2019/11/03}{Change to three-line header with copyright info and DOI.}
%    \begin{macrocode}
\if@preprint
\renewcommand\@oddhead{%
  \if@proceedings 
  \textsc{To appear in ``\@conftitle{} (\rmaa@YearStyle{\@year})''}\hfil
  \textit{RMxAC}
  \else
  \textsc{To appear in} \@journalname\hfil
  \fi}\fi%
\if@guide
\renewcommand\@oddhead{\textsc{Documentation} (\rmaa@YearStyle{\@year})\hfil}
\fi%
\if@manuscript
\renewcommand\@oddhead{Manuscript for \@journalname{} (\rmaa@YearStyle{\@year})\hfil}
\fi%
\if@journal
\renewcommand\@oddhead{%
  \parbox{\textwidth}{
    \@journalname, \rmaa@VolumeStyle{\@volume},
    \rmaa@PagesStyle{\@artpages}\
    (\rmaa@YearStyle{\@year})\hfil\\[1pt]
    {\small\textcopyright{} \@year:
      Instituto de Astronom\'\i{}a,
      Universidad Nacional Aut\'onoma de M\'exico}\hfil\\[1pt]
    {\small\@DOI@url}\hfil
  }%
}\fi%
  \renewcommand\@evenhead{\@oddhead}}
%    \end{macrocode}
% By default, use the |\ps@rmaa| page style.
%    \begin{macrocode}
\pagestyle{rmaa}
%    \end{macrocode}
% \end{macro}
% \end{macro}
%
% \subsection{Preamble commands}\label{sec:preamble} 
% Everything you need for setting the title page, TOC and index
% entries, etc.
% \subsubsection{Commands for the title page}
% \label{sec:title-commands}
% \begin{macro}{\resumen}
% \begin{macro}{\@resumen}
% \begin{macro}{\abstract}
% \begin{macro}{\@abstract}
% User commands for setting the resumen and abstract. 
% These should be used in the document header. The internal commands
% |\@resumen| and |\@abstract| are used to hold the actual text
% pending output via the |\maketitle| command. 
%    \begin{macrocode}
\newif\if@setresumen
\newcommand{\resumen}[1]{%
%%  \renewcommand{\@resumen}{{\selectlanguage{spanish}#1}}%
  \renewcommand{\@resumen}{{\ignorespaces#1}}%
  \@setresumentrue}
\renewcommand{\abstract}[1]{%
  \renewcommand{\@abstract}{{\ignorespaces#1}}
  \renewcommand\ADS@abstract{#1}
%% Not a good idea, since then author macros are not expanded and
%% cause errors later. However, by not doing it we are vulnerable to
%% fragile commands.....
  \if@noexpandADS\@onelevel@sanitize{\ADS@abstract}\fi
  \@abstractpresenttrue
  }
\newcommand{\ADSnoexpand}{\@noexpandADStrue}
%    \end{macrocode}
% \begin{macro}{\boldabstract}
% This is the version for doing 1-page conference proceedings. Still
% save the text for ADS but then write it out immediately as a bold
% paragraph. 
%    \begin{macrocode}
\newcommand{\boldabstract}[1]{%
  \gdef\ADS@abstract{#1}%
  \renewcommand{\@abstract}{#1}%
  \if@noexpandADS\@onelevel@sanitize{\ADS@abstract}\fi%
  \noindent\textbf{\boldmath #1}
  }
%    \end{macrocode}
% \end{macro}
% Default value in each case is to print out a warning if not set
%    \begin{macrocode}
  \newcommand{\@resumen}{%
    \@setresumenfalse \ClassWarningNoLine{rmaa}{No Spanish abstract
      (resumen) specified - if you do not speak Spanish, this will be
      translated by the editors}%
    \textbf{Favor de proporcionar un resumen en espa\~nol. If you are
      unable to translate your abstract into Spanish, the editors will
      do it for you.} }
  \newcommand{\@abstract}{\ClassWarningNoLine{rmaa}{No abstract
      specified}}
%    \end{macrocode}
% \end{macro}
% \end{macro}
% \end{macro}
% \end{macro}
%
% \begin{macro}{\title}
% \begin{macro}{\@rawtitle}
% \begin{macro}{\@TOCtitle}
% \begin{macro}{\rmaaIndexTitle}
% Similarly, |\title| is the preamble command to set the title of the
% article. This is used in various places, typeset in different
% styles, hence the use of the internal |\@rawtitle|, |\@title|,
% |\rmaaIndexTitle|, and |\@TOCtitle| commands, each of which uses the
% relevant style from section~\ref{sec:text-styles}
%    \begin{macrocode}
\newcommand{\@rawtitle}{Raw Title}
\newcommand{\@TOCtitle}{\rmaa@TOCTitleStyle{Title for TOC}}
\newcommand{\rmaaIndexTitle}{\rmaa@IndexTitleStyle{Title for index}.}
\newcommand{\@TOCaftertitle}{}
\renewcommand{\title}[1]{%
%%  {\ADS@sanitize\xdef\@rawtitle{#1}}
  \renewcommand\@rawtitle{\noexpand#1}
  \let\ADS@title=\@rawtitle
  \renewcommand{\@title}{\rmaa@TitleStyle{#1}}%
  \renewcommand{\@TOCtitle}{\rmaa@TOCTitleStyle{#1}\@TOCaftertitle}%
  \renewcommand{\rmaaIndexTitle}{\rmaa@IndexTitleStyle{#1\@TOCaftertitle}.}%
  }
%    \end{macrocode}
% \end{macro}
% \end{macro}
% \end{macro}
% \end{macro}
% 
% All the weird and whacky ways of indicating the author
% affiliations. 
% \begin{macro}{\affil}
% First, the way that sets it directly underneath the authors name.
% Note that |\affil| trumps all other affiliations for the purposes of
% getting into the ADS stuff (we use |\gdef|).   
% \dots 
%    \begin{macrocode}
\newcommand{\affil}[1]{%
  \\[1.5ex]\rmaa@AffilStyle{#1}\\[1.5ex]%
  \if@firstaffil\gdef\ADS@affil{#1}\@firstaffilfalse\fi%
  }
\renewcommand{\and}{and\\[1.5ex]}
%    \end{macrocode}
% \begin{macro}{\altaffil}
% Alternative affiliations are footnoted. 
%    \begin{macrocode}
\newcommand{\altaffil}[1]{%
  \if@firstaltaffil\gdef\ADS@affil{#1}\@firstaltaffilfalse\fi%
  \thanks{#1}}
%    \end{macrocode}
% \begin{macro}{\altaffilmark}
% \begin{macro}{\altaffiltext}
% Alternative-alternative mechanism for affiliations---should not be
% mixed with |\altaffil| (unless you know what you are doing) or there
% will be problems with numbering.
%    \begin{macrocode}
\newcommand{\altaffilmark}[1]{\textsuperscript{#1}}
\newcommand{\altaffiltext}[2]{%
  \if@firstaltaffil\gdef\ADS@affil{#2}\@firstaltaffilfalse\fi
  \stepcounter{footnote}%
  \protected@xdef\@thanks{\@thanks
    \protect\footnotetext[#1]{#2}}} 
%    \end{macrocode}
% \end{macro}
% \end{macro}
% \end{macro}
% \end{macro}
%
% \begin{macro}{\keywords}
% \begin{macro}{\@keywords}
% \begin{macro}{\@keywordheader}
% \label{keywords}Sundry minor preamble commands. All
% self-explanatory, I hope. |\keywords| is the old user interface
% (deprecated). New articles should use |\addkeyword| once for each
% keyword. That way the ADS record comes out right too (comma rather
% than em-dash separator).
% \changes{2.2beta12}{2001/12/07}{\cs{addkeyword} New user interface
%   for keywords.} 
% \changes{2.2beta13}{2002/03/11}{\cs{addkeyword} Improved behaviour
%   with .ads file. Use
%   \cs{ADS@sanitize} to change definitions of `~' and \cs{ion}.}
% \changes{3.23}{2007/09/09}{\cs{addkeyword} now signals a compilation
% error if the author tries to add more than 6 keywords.}
% \changes{3.24}{2007/09/14}{kwdcounter is now reset in
% \cs{@begofarticlehook} to prevent error in book style} 
%    \begin{macrocode}
\newcommand{\@kwdsep}{ --- }
\newcommand{\ADS@kwdsep}{, }
\newcounter{kwdcounter}\setcounter{kwdcounter}{0}
\newcommand{\addkeyword}[1]{%
  \addtocounter{kwdcounter}{1}
  \ifthenelse{\thekwdcounter > 6}{
    \ClassError{rmaa}{Too many keywords. A maximum of 6 keywords is allowed.}{}
  }{}
  \if@firstkwd%
  \renewcommand\@keywords{\if@boldkeywords\uppercase{#1}\else #1\fi}
  {\ADS@sanitize\def\ion{\plainion}\gdef\ADS@keywords{#1}}
  \@firstkwdfalse
  \else%
  \edef\@keywords{\@keywords\@kwdsep \if@boldkeywords\uppercase{#1}\else #1\fi}%
  {\ADS@sanitize\def\ion{\plainion}\xdef\ADS@keywords{\ADS@keywords\ADS@kwdsep #1}}%
  \fi}
\newcommand{\ADS@sanitize}{%
  \def\mathrm{}
  \def\@{\ }
  \def~{\ }
  \def\langle{<}
  \def\rangle{>}
  \def\log{log }
  \def\sin{sin }
  \def\cos{cos }
  \def\tan{tan }
  \def\fd{^d.}
  \def\fh{^h.}
  \def\fm{^m.}
  \def\fs{^s.}
  \def\fdg{�.}
  \def\farcm{'.}
  \def\farcs{''.}
  \def\arcdeg{�}
  \def\,{ }
  }
\newcommand{\keywords}[1]{%
  \renewcommand{\ADS@keywords}{#1}
  \renewcommand{\@keywords}{\uppercase{#1}}
  }
\newcommand{\@keywordheader}{\rmaa@KeywordHeaderStyle{Key Words:}}
\newcommand{\@keywords}{%
  \ClassWarningNoLine{rmaa}{No keywords specified}%
  Keywords go here
  }
%    \end{macrocode}
% \begin{macro}{\shorttitle}
% \begin{macro}{\shortauthor}
% \begin{macro}{\@shorttitle}
% \begin{macro}{\@shortauthor}
% Commands for setting the running headers.
%    \begin{macrocode}
\newcommand{\shorttitle}[1]{%
  \renewcommand\@shorttitle{\MakeUpperCase{#1}}}
\newcommand{\shortauthor}[1]{%
  \renewcommand\@shortauthor{\MakeUpperCase{#1}}}
\newcommand{\@shorttitle}{SHORT TITLE GOES HERE}
\newcommand{\@shortauthor}{SHORT AUTHOR GOES HERE}
%    \end{macrocode}
% \end{macro}
% \end{macro}
% \end{macro}
% \end{macro}
% \end{macro}
% \end{macro}
% \end{macro}


% \subsubsection{The full addresses}
% \label{sec:full-addresses}
% \begin{macro}{\fulladdresses}
% \begin{macro}{\@fulladdresses}
% \begin{macro}{\@fulladdheight}
% \begin{macro}{\@fulladddepth}
% \begin{macro}{\@fulladdbox}
% Aaghh, the full addresses! A perennial source of problems. 
%    \begin{macrocode}
\newlength{\@fulladdheight}
\newlength{\@fulladddepth}
\newsavebox{\@fulladdbox}
%% \fulladdresses command modified (11 May 2000) so that we can put
%% \fulladdresses{} in the article to suppress the output of the
%% addresses
\newcommand{\@FinalFigure}{\relax}
\newcommand{\FinalFigure}[1]{\renewcommand{\@FinalFigure}{#1}}
\newcommand{\fulladdresses}[1]{%
  \renewcommand{\@rmaaArg}{#1}
  \ifx\@rmaaArg\@rmaaEmpty
     \renewcommand{\@fulladdresses}{}
  \else
     \renewcommand{\@fulladdresses}{
       \normalsize%
       \ifx\@FinalFigure\relax\else\vfill\@FinalFigure\vfill\fi
%%       \vspace*{1ex}
       %% do addresses as a list (copied from thebibliography environment) 
       \list{\null}{\leftmargin 1.5em\topsep\z@\labelwidth\z@\labelsep\z@\itemindent -1.5em
         \itemsep\z@\parsep\z@\usecounter{enumi}}
       \sloppy\clubpenalty4000\widowpenalty4000
       \sfcode`\.=1000\relax
       #1
       \endlist
       }
  \fi
  }
\newcommand{\@fulladdresses}{Full addresses go here}
%    \end{macrocode}
% \begin{macro}{\suppressfulladdresses}
% New user-level command for suppressing full addresses should you so
% desire. Currently used in the single-page-poster style. 
%    \begin{macrocode}
\newcommand{\suppressfulladdresses}{\@suppressfulladdressestrue\@needaddressesfalse}
%    \end{macrocode}
% \end{macro}
% \end{macro}
% \end{macro}
% \end{macro}
% \end{macro}
% \end{macro}

% \begin{macro}{\printdates}
% \begin{macro}{\noprintdates}
% Command to activate the printing out of the received/accepted
% dates. This can also be done using the \texttt{printdates} option,
% which is the default for \texttt{journal}, but not for \texttt{proceedings}.
%    \begin{macrocode}
\newcommand{\printdates}{\@printdatestrue}
\newcommand{\noprintdates}{\@printdatesfalse}
%    \end{macrocode}
% \end{macro}
% \end{macro}
% \begin{macro}{\@outputaddressesifnecessary}
% Command to output the addresses if we haven't already done so
% This is called by: 
% \begin{enumerate}
% \item The end of |thebibliography| environment
% \item The |\outputfulladdresses| command
% \end{enumerate}
%    \begin{macrocode}
\newcommand{\@outputaddressesifnecessary}{%
  \ClassInfo{rmaa}{Calling @outputaddressesifnecessary}
  \ifthenelse{\boolean{@needaddresses}}{%
    %% If in one column style, or in the first column of a two column
    %% style then we can try to put it on this page 
    \ifthenelse{\boolean{@firstcolumn} \or \not\boolean{@rmtwocol}}{%
      \@outputaddresses}{%
      %% otherwise give warning it will be on a separate page
      \ClassWarning{rmaa}{Full addresses will go on a separate
        page (please see User Guide for how to remedy this)}}%
    %% If we haven't managed yet, put it on its own page
    \ifthenelse{\boolean{@needaddresses}}{%
      \if@debug{\ClassWarning{rmaa}{In this silly bit!}}\fi
      \clearpage\vfill\@outputaddresses}{}%
    }{
    %% If we don't need them, we still might need to do the eoa hook 
    \ifthenelse{\boolean{@suppressfulladdresses}}{%
      \if@debug{\ClassWarning{rmaa}{XXX}}\fi
%%      \if@book\else\@endofarticlehook\null\fi
      }{}
    }%
  }  
%    \end{macrocode}
% \end{macro}

% \begin{macro}{\outputfulladdresses}
% A user interface to the above command, in case there is no biblio
% (11 May 2000) 
%    \begin{macrocode}
\newcommand{\outputfulladdresses}{\@outputaddressesifnecessary}
%    \end{macrocode}
% \end{macro}

% \begin{macro}{\adjustfinalcols}
% This command starts a new column and should be 
% used to ensure that the columns balance on the last page. It also
% outputs the full 
% author addresses to ensure that this occurs in the left hand
% column. 
% The |\par| command is so that the current line will be taken account
% of when working out |\pagetotal|
%    \begin{macrocode}
\newcommand{\adjustfinalcols}{\if@twocolumn\goodbreak\@outputaddresses\newpage\noindent\fi}
%    \end{macrocode}
% \end{macro}

% \begin{macro}{\@outputaddresses}
% Command to do the work of outputting the addresses. 
% This can be called by 
% \begin{enumerate}
% \item an explicit |\adjustfinalcols|
% \item via |\@outputaddressesifnecessary|, q.v.
% \end{enumerate}
%
%    \begin{macrocode}
\newcommand{\@outputaddresses}{%
  % Allow option of suppressing the full addresses, in which case
  % we just run the hook if we are not in a book style
  \if@suppressfulladdresses
%%     \if@book\else\@endofarticlehook\null\fi
  \else
  %% save the full addresses in a box for output later
  \sbox{\@fulladdbox}{\parbox{\textwidth}{\@fulladdresses}}
  % calculate height of address box  
  \settoheight{\@fulladdheight}{\usebox{\@fulladdbox}}
  \settodepth{\@fulladddepth}{\usebox{\@fulladdbox}}
  \addtolength{\@fulladdheight}{\@fulladddepth}
  \if@debug
  { \typeout{Address total height: \the\@fulladdheight}
    \typeout{Text height: \the\pagetotal}
    \typeout{Goal height: \the\pagegoal}
    }
  \fi
  % reuse \@fulladdheight to mean the total height of text+addresses
  \addtolength{\@fulladdheight}{\pagetotal}
%%  \addtolength{\@fulladdheight}{-\baselineskip}
  \if@debug\typeout{Adjusted text + addresses height: \the\@fulladdheight}\fi
  % check it will fit on current page
  \ifdim\@fulladdheight<\pagegoal
     %% must shift it to the left
     %% to compensate for the hanging indentation in the reference list. 
     %% NOTE: this is not foolproof!
     %% NB. This will give overfull hbox warnings, which is a bit ugly.
     \vfill\noindent\if@biblist\hspace*{-\leftmargin}\fi
     \usebox{\@fulladdbox}
     % Set this to show we've done it
     \@needaddressesfalse
     % and assume we are at the end of the article
     % so that we do the EOA hook if not in the book style
%%     \if@book\else\@endofarticlehook\fi
  \else
     % if not, just do nothing. With luck, this function will get 
     % called again on the next page
     \if@debug{\typeout{Holding over full addresses to next page}}\fi
  \fi\fi
  }
%    \end{macrocode}
% \end{macro}

% \subsubsection{Commands for editors}
% \label{sec:editor-commands}
% Some user/editor commands for doing the TOC and index in the book
% option. 
% They are defined here rather in the book-specific part since they
% need to be always available in order to avoid breaking documents. 
% \begin{macro}{\listofauthors}
% \begin{macro}{\@rawauthors}
% \begin{macro}{\rmaaTOCauthors}
% The command |\listofauthors| formats an author list for use in the
% TOC. It should contain the same names as in the |\author| command,
% but without the affiliation garbage and using |\&| instead of
% |and|. 
%    \begin{macrocode}
\newcommand{\@rawauthors}{Raw Authors}
\newcommand{\rmaaTOCauthors}{\noexpand\rmaaTOCAuthorStyle{Author list for TOC}}
\newcommand{\listofauthors}[1]{%
  \renewcommand{\@rawauthors}{#1}
  \renewcommand\rmaaTOCauthors{\noexpand\rmaaTOCAuthorStyle{#1}}}
%    \end{macrocode}
% \begin{macro}{\indexauthor}\label{indexauthor}
% The command |\indexauthor| produces an index entry for each
% author. 
% This command must come \emph{after} the
% |\title| and the |\listofauthors| commands. Also, there should be one 
% |\indexauthor| command for \emph{each} author. We also add the
% author onto the end of the |\ADS@authors| macro and construct the
% |\ADS@authorinitial| by a most cunning method (see below), but only
% for the first author. 
% 
% No, No, NO! This was all wrong. Backed out completely: 
% \begin{quote}
% Completely rewritten (19 Dec 2002) since we now use the index class, so 
% as to be able to alphabetize accented characters. A side effect of this is
% that |\index| no longer expands out the argument after the |@|, so we 
% instead add all the |\index| commands to a list (|\indexauthor@deferlist|), 
% which is executed in the |\@begofarticlehook|. 
% 
% Note the stuff with |\toks0| is copied from Exercise 20.15 of the \TeX{}book---it is
% to make sure the latin1-accented name gets passed into the idx
% file. 
% \end{quote}
% We now go back to the old way of doing things, using an optional
% argument to the \cs{indexauthor} command that has the name without
% accents. 
% \changes{3.08}{2003/01/16}{Remerged changes from 3.06 that got
% dropped from 3.07 after a synchronization disaster. Basically, the
% accentless optional argument to \cs{indexauthor}. Must add the macro
% directories to my unison set-up to stop this sort of thing from happening.}
%    \begin{macrocode}
\newcommand\@isortkey{}
\newcommand\@emptykey{empty}
\newcommand\indexauthor[2][empty]{
  \renewcommand\@isortkey{#1}
  \ifx\@isortkey\@emptykey \renewcommand\@isortkey{#2}\fi
  \index{\@isortkey @{{\rmaaIndexAuthorStyle{#2}}\hskip1em plus1em minus0.4em
      \rmaaIndexTitle\hskip1em plus2em minus0.4em\penalty-100
      \rmaaTOCauthors}}
  \indexauthorADS{#2}
}
%%\def\indexauthor@deferlist{\relax}
%%\long\def\indexauthor#1{
%%  \def\@@name{#1} \toks0=\expandafter{\@@name}
%%  \toks1=\expandafter{\indexauthor@deferlist}
%%  \protected@xdef\indexauthor@deferlist{%
%%    \the\toks1
%%    \noexpand\index{\the\toks0@{\noexpand\rmaaIndexAuthorStyle{\the\toks0}\noexpand\hskip1em plus1em minus0.4em 
%%        \rmaaIndexTitle\noexpand\hskip1em plus2em minus0.4em\noexpand\penalty-100
%%        \rmaaTOCauthors}}
%%  }
%%  \indexauthorADS{#1}
%%  }
%%\xdef\indexauthor@deferlist{}
\newcommand\indexauthorADS[1]{%
  \if@firstauthor\protected@xdef\ADS@authorinitial{\firstletter{#1}}\fi
  \protected@xdef\ADS@authors{\ADS@authors\if@firstauthor\else ; \fi
    #1}
  \@firstauthorfalse
}
%    \end{macrocode}
% Note that the first |#1| sets the key for the index, making sure it
% actually starts with  
% first letter of the surname. Meanwhile, the |@{...}| expression,
% which starts with some  
% control sequences is what is written to the |.ind| file. 
% See the makeindex man page for details. 
% \end{macro}
% \end{macro}
% \end{macro}
% \end{macro}

% \begin{macro}{\KillAllAccents}
% \begin{macro}{\firstletter}
% These are utility functions used in the above. |lett| is a counter
% indicating the current letter in the string. |\firstletter| just
% prints the first letter of its argument.
%    \begin{macrocode}
\newcommand{\KillAllAccents}{\def\'{}\def\`{}\def\~{}\def\.{}\def\i{i}}
\newcommand\firstletter[1]{\@car #1\@nil}
%    \end{macrocode}
% \end{macro}
% \end{macro}
% \begin{macro}{\SetVolume}
% \begin{macro}{\SetMSnumber}
% \begin{macro}{\SetFirstPage}
% \begin{macro}{\SetYear}
% \begin{macro}{\ReceivedDate}
% \begin{macro}{\AcceptedDate}
% These are for the editors to set the volume number, etc. 
%    \begin{macrocode}
\newcommand{\SetVolume}[1]{%
  \renewcommand{\@volume}{#1}
  \edef\ADS@volume{\PadToFourWithDots{#1}}
  \renewcommand{\@volumen}{\rmaa@VolumenStyle{volumen #1}}%
  }
\newcommand{\SetMSnumber}[1]{\renewcommand{\@MSnum}{#1}}
%    \end{macrocode}
% Note that there is no real reason to use |\SetFirstPage| since all
% the page numbering is done automatically in the book style. 
%    \begin{macrocode}
\newcommand{\SetFirstPage}[1]{\setcounter{page}{#1}}
\newcommand{\SetYear}[1]{%
  \renewcommand{\@year}{#1}
  \renewcommand{\ADS@year}{#1}
  \renewcommand{\@anyo}{\rmaa@VolumenStyle{#1}}
  }
\newcommand{\ReceivedDate}[1]{\renewcommand{\@received}{#1}}
\newcommand{\AcceptedDate}[1]{\renewcommand{\@accepted}{#1}}
%    \end{macrocode}
%
% \begin{macro}{\DOI}
% User/editor command to enter the DOI, which will be typeset as a link.
% \changes{4.0}{2019/11/02}{New user command \cs{DOI} to typeset
% DOI in the header of first page and add it to ADS info.}
%    \begin{macrocode}
\newcommand\DOI[1]{%
  \def\@DOI@url{\url{https://doi.org/10.22201/ia.01851101p.#1}}%
  \def\ADS@doi{https://doi.org/10.22201/ia.01851101p.#1}%
}%
%    \end{macrocode}
% \end{macro}
%
% \begin{macro}{\ORCIDiD}
% \changes{4.0}{2019/11/12}{Add mechanism for giving ORCID iD links}
% \begin{macro}{\ORCIDiDfull}
% \begin{macro}{\@orcid@icon}
% \begin{macro}{\@orcid@iconbig}
% Commands to add icon links to ORCID iDs of authors. For example
% \begin{quote}
% |\ORCIDiD|\marg{0000-0000-0000-0000}
% \end{quote}
% to be used after each author name in the |\author|
% command and in the |\fulladdresses|.  Alternatively,
% |\OORCIDiDfull| will typeset the icon plus the full text of the
% link, but we don't use that at present.  The icon image is stored in |ORCIDiD_iconvector.pdf|, which needs to be on the input path.  In principle, we could make EPS versions if necessary.
%    \begin{macrocode}
\newcommand\@orcid@icon{\includegraphics[height=1.8ex]{ORCIDiD_iconvector}}
\newcommand\@orcid@iconbig{\includegraphics[height=2ex]{ORCIDiD_iconvector}}
\newcommand\ORCIDiD[1]{\href{https://orcid.org/#1}{\,\@orcid@icon}}
\newcommand\ORCIDiDfull[1]{%
  \href{https://orcid.org/#1}{\@orcid@iconbig}\,\url{https://orcid.org/#1}}
%    \end{macrocode}
% \end{macro}
% \end{macro}
% \end{macro}
% \end{macro}
%
% \begin{macro}{\@volume}
% \begin{macro}{\@volumen}
% \begin{macro}{\@artpages}
% \begin{macro}{\@year}
% \begin{macro}{\@anyo}
% \begin{macro}{\@received}
% \begin{macro}{\@accepted}
% \begin{macro}{\@MSnum}
% Default values for the containers used by above commands, plus a
% few more useful ones. 
%    \begin{macrocode}
\newcommand{\@volume}{00}
\newcommand{\@volumen}{\rmaa@VolumenStyle{00}}
\newcommand{\@artpages}{%
  \oldpageref{\articlename:RMAAFirstPage}--\oldpageref{\articlename:RMAALastPage}}
\newcommand{\@year}{\the\year}
\newcommand{\@anyo}{\rmaa@VolumenStyle{2007}}
\newcommand{\@journalname}{\rmaa@JournalNameStyle{%
    \if@proceedings\RMAASClong\else\RMAAlong\fi}}
\newcommand{\@received}{}
\newcommand{\@accepted}{}
\newcommand{\@MSnum}{0000}
\newcommand{\@DOI@prefix}{https://doi.org/10.22201/ia.01851101p}
\def\@DOI@url{\url{https://doi.org/10.22201/ia.01851101p.2099.00.00.00}}
%    \end{macrocode}
% \end{macro}
% \end{macro}
% \end{macro}
% \end{macro}
% \end{macro}
% \end{macro}
% \end{macro}
% \end{macro}
% \end{macro}
% \end{macro}
% \end{macro}
% \end{macro}
% \end{macro}
% \end{macro}
% 
% The following commands are most useful in the \texttt{book}
% style. However, the |\SetConfTitle| command can now be used in
% individual articles to set the header for the \texttt{preprint}
% style. 
%    \begin{macrocode}
\newcommand{\@reunionplace}{Lugar de la reuni\'on}
\newcommand{\@reuniondate}{Fecha de la reuni\'on}
\newcommand{\@conftitle}{Conference Title}
\newcommand{\SetReunion}[2]{\renewcommand{\@reunionplace}{#1}\renewcommand{\@reuniondate}{#2}}
\newcommand{\SetConfTitle}[1]{\renewcommand{\@conftitle}{#1}}
\newcommand{\@editors}{Juan Ram\'on de la Fuente and ``El Mosh''}
\newcommand{\SetEditors}[1]{\renewcommand{\@editors}{#1}}
\newcommand{\@editores}{Juan Ram\'on de la Fuente y ``El Mosh''}
\newcommand{\SetEditores}[1]{\renewcommand{\@editores}{#1}}
\newcommand{\@mes}{mes de la publicaci\'on}
\newcommand{\SetMes}[1]{\renewcommand{\@mes}{\rmaa@VolumenStyle{#1}}}
%    \end{macrocode}

% \subsection{Notes for Author Proofs and Communication between Editors}
% \label{sec:notes}
% We introduce a command \cs{Note} that should be used to describe
% any change made to the manuscript. These are a bit tricky to get
% right. In general we use \cs{marginpar} but this only works in outer
% paragraph mode, so we just bung the box where we are if we are in a
% figure caption or something. If we are in the preamble, then we just
% use a \marginpar as part of the \maketitle command.
%    \begin{macrocode}
\newif\if@arenotes
\@arenotesfalse
\newcounter{notecounter}
\setcounter{notecounter}{0}
\newcommand\NotesList{}
\newcommand\Note[1]{%
  \if@proof{%
    \@arenotestrue%
    \stepcounter{notecounter}%
    \if@inpreamble\else{%
      \NoteMark{\arabic{notecounter}}%
    }\fi
    \protected@xdef\NotesList{\NotesList 
      \protect\item #1}%
  }\fi%
}
\newcommand\Changed[2]{Changed ``#1'' to ``#2''.}
\newcommand\NoteMark[1]{%
  \ifinner
  \raisebox{1ex}[0pt][0pt]{\makebox[0pt][l]{\framebox{\footnotesize\bfseries
        N#1}}}%
  \else
  \marginpar{\framebox{\footnotesize\bfseries N#1}}%
  \fi
}
\newcommand\RefNote[1]{\if@proof\NoteMark{\ref{#1}}\relax\fi}
\newcommand\NoteLabel{\protect\label}
\newcommand\OutputNotes{%
  \bgroup
  \onecolumn
  \section*{EDITORIAL NOTES}
  \begin{enumerate}
    \NotesList
  \end{enumerate}
  \egroup
}
\DeclareRobustCommand\cs[1]{\texttt{\char`\\#1}}

%    \end{macrocode}
% \subsection{Bibliographic Records for ADS}
% \label{sec:ADS}
% \changes{2.2beta11}{2001/11/27}{Just about finished all the ADS
% stuff now. Except for the references, which I have given up on for
% the time being.}
% \begin{macro}{\ADS@BibCode}
% \begin{macro}{\ADS@year}
% \begin{macro}{\ADS@journal}
% \begin{macro}{\ADS@volume}
% \begin{macro}{\ADS@firstpage}
% \begin{macro}{\ADS@authorinitial}
% We use the |\ADS@BibCode| command to hold the bibliographic code for
% ADS. It is 
% used in the |%R| 
% field of the bibliographic record for each article. 
% The format is 
% \begin{quote}
% \texttt{YYYYJJJJJVVVVMPPPPA}  
% \end{quote}
% where \texttt{YYYY} is the year, \texttt{JJJJJ} is the journal code,
% which is \texttt{RMxAA} or \texttt{RMxAC}, \texttt{VVVV} is the
% volume number (right-justified), \texttt{M} is not used by us,
% \texttt{PPPP} is the first page number (right-justified), and
% \texttt{A} is the initial of the first author's surname. The volume
% and page fields are padded with periods (`\texttt{.}'), which causes
% us some fun (this is done with the macro |\PadToFourWithDots| described
% below) . 
% 
% All these individual bits are held in the commands |\ADS@year|,
% |\ADS@journal|, |\ADS@volume|, |\ADS@firstpage|,
% |\ADS@authorinitial|. 
%    \begin{macrocode}
\newcommand{\SetADSBibCode}[1]{\renewcommand{\ADS@BibCode}{#1}}
\newcommand{\ADS@BibCode}
{\ADS@year\ADS@journal\ADS@volume.\PadToFourWithDots{\ADS@firstpage}\ADS@authorinitial}
\newcommand{\ADS@year}{2001} 
\newcommand{\ADS@journal}{\if@proceedings RMxAC\else RMxAA\fi}
\newcommand{\ADS@volume}{....} 
%    \end{macrocode}
% These two are just set to default values here, |\ADS@firstpage| is set
% in the |document| environment for included articles, in the
% |frontsection| environment, |\tableofcontents| command,  |index| and
% |abstracts| environments. Also by |\AtBeginDocument| for stand-alone
% articles. 
% 
% God knows how we are going to do |\ADS@authorinitial|!
%    \begin{macrocode}
\newcommand{\ADS@firstpage}{....}
\newcommand{\ADS@authorinitial}{A}
%    \end{macrocode}
% \end{macro}
% \end{macro}
% \end{macro}
% \end{macro}
% \end{macro}
% \end{macro}
% We also have a whole slew of other fields that can be set. These are
% all listed in Table~\ref{tab:ADS}. Most of them still have to be
% defined. 
% \begin{table}[!ht]\raggedleft
% \caption{ADS Bibliographic Records}\label{tab:ADS}
% \begin{tabular}{lp{8em}lp{8em}}\hline
% Code & Description & Macro Name & Comment \\ \hline
% |%R| & Bibliographic Code & |\ADS@BibCode| & required \\
% |%A| & Author list & |\ADS@authors| & required, generate from |\indexauthor| \\
% |%F| & Author affiliation & |\ADS@affil| & format: AA(\dots)
% AB(\dots) \dots \\
% |%J| & Journal Name & |\ADS@journalname| &  \\
% |%D| & Publication Date & |\ADS@fulldate| &  \\
% |%L| & Last Page & |\ADS@lastpage| & generate from label \\
% |%T| & Title & |\ADS@title| & required, generate from |\title| \\
% |%C| & Abstract Copyright & |\ADS@copyright| & set to IA-UNAM \\
% |%U| & URL for electronic document & |\ADS@url| & set to PDF file \\
% |%d| & DOI for article & |\ADS@doi| & set by |\DOI| \\
% |%K| & Keywords & |\ADS@keywords| & set equal to |\@keywords| \\
% |%A| & Abstract Text & |\ADS@abstract| & set equal to |\@abstract| \\
% |%R| & References & |\ADS@references| & generate from the
% bibliography  \\ \hline
% \end{tabular}
% \end{table}
% A lot of these macros are set by other commands, as mentioned in the
% Table, but some need to be set  explicitly here. 
% 
% \begin{macro}{\ADS@doi}
% \changes{4.0}{2019/11/08}{Save DOI for ADS}
% The user/editor command \cs{DOI} is used to set this. 
%    \begin{macrocode}
\newcommand{\ADS@doi}{https://doi.org/10.22201/ia.01851101p.0000.00.00.00}
%    \end{macrocode}
% \end{macro}
% 
% \begin{macro}{\ADS@journalname}
% This
% does not contain any accents at the moment so as not to cause
% problems---this maybe should be changed though. 
%    \begin{macrocode}
\newcommand{\ADS@journalname}{%
  \if@proceedings\@conftitle\space (Eds. \@editors) \fi%
  Revista Mexicana de Astronom\'\i{}a y Astrof\'\i{}sica%
  \if@proceedings\space (Serie de Conferencias)\fi%
  \space Vol. \@volume, pp. \ADS@firstpage--\ADS@lastpage\space (\@year)%
  \space (http://www.astroscu.unam.mx/\~{}rmaa/)
  }
%    \end{macrocode}
% \end{macro}
% \begin{macro}{\ADS@keywords}
% \begin{macro}{\ADS@abstract}
% \begin{macro}{\ADS@title}
% |\ADS@abstract| is straightforward, since we already have the text
% saved as |\@abstract|. We could do the same with |\ADS@keywords| but
% to avoid unwanted formatting we explicitly set it in |\keywords|
% (p.~\pageref{keywords}). |\ADS@title| is just set to |\@rawtitle|,
% which we were already using for another purpose. 
%    \begin{macrocode}
\let\ADS@title=\@rawtitle
\newcommand{\ADS@keywords}{No keywords specified}
\newcommand{\ADS@abstract}{No abstract specified}
%    \end{macrocode}
% \end{macro}\end{macro}\end{macro}
% \begin{macro}{\ADS@authors}
% |\ADS@authors| wants to be in ``Other, A. N.'' format, but with
% semi-colons between authors. The most natural way to do this seems to
% be with |\indexauthor| (p.~\pageref{indexauthor}). We append each
% author to this macro, which means we must explicitly clear it for
% each new article (done in |\article|). .   
%    \begin{macrocode}
\newcommand{\ADS@authors}{}
%    \end{macrocode}
% \end{macro}
% \begin{macro}{\ADS@fulldate}
% \begin{macro}{\ADS@lastpage}
% \begin{macro}{\ADS@copyright}
% \begin{macro}{\ADS@website}
% \begin{macro}{\ADS@url}
% \begin{macro}{\SetWebsite}
% \begin{macro}{\SetMonthNumeric}
% Other sundry data, together with user commands for bits that can be
% set by the Editors. The |\ADS@url| is constructed from the website
% (currently |http://www.astrosmo.unam.mx/rmaa/|,
% the journal descriptor (e.g., |RMxAC..10/|), followed by the
% path to the PDF file (e.g., |PDF/RMxAC..10_HEW1.pdf|). 
% 
%    \begin{macrocode}
\newcommand{\ADS@website}{http://www.astrosmo.unam.mx/rmaa/}
\newcommand{\ADS@url}
{\ADS@website\ADS@journal\ADS@volume\ADS@number/PDF/\ADS@journal\ADS@volume\ADS@number\_\articlename.pdf}
\newcommand{\ADS@number}{}
\newcommand{\ADS@fulldate}{00/\ADS@year} 
\newcommand{\ADS@lastpage}{00} 
\newcommand{\ADS@copyright}{(C) \ADS@year; Instituto de Astronom\'\i{}a,
  Universidad Nacional Aut\'onoma de M\'exico} 
\newcommand{\SetNumber}[1]{\renewcommand{\ADS@number}{-#1}}
\newcommand{\SetWebsite}[1]{\renewcommand{\ADS@website}{#1}}
\newcommand{\SetMonthNumeric}[1]{\renewcommand{\ADS@fulldate}{#1/\ADS@year}}
%    \end{macrocode}
% \end{macro}\end{macro}\end{macro}\end{macro}\end{macro}\end{macro}\end{macro}
% \begin{macro}{\ADS@affil}
% \begin{macro}{\ADS@references}
% These two are somewhat more problematical. Affiliation currently
% only works for the first author (see \dots) pending redesign of the
% user interface. References are now done by wrapping the
% |thebibliography| environment to open another input stream on the
% same file and search through for the |\bibitems|. 
%    \begin{macrocode}
\newcommand{\ADS@references}{}
\newcommand{\ADS@affil}{}
%    \end{macrocode}
% \end{macro}\end{macro}
% 
% Define two macros that expand to the strings
% |\begin{thebibliography}| and |\end{thebibliography}|. Note that the
% idea for this and much else is inspired by the code for the docstrip
% package. Also define |\StartBibX| in case authors use the more
% standard |\begin{thebibliography}{}| syntax. 
%    \begin{macrocode}
\edef\EndBib{\expandafter\@gobble\string\\end\string{thebibliography\string}}
\edef\StartBib{\expandafter\@gobble\string\\begin\string{thebibliography\string}}
\edef\StartBibX{\expandafter\@gobble\string\\begin\string{thebibliography\string}\string{\string}}
%    \end{macrocode}
% Define two new logical variables. |\ifAnotherline| is true if we
% want to read another line from the input stream. It is set to false
% when we encounter the |\end{thebibliography}|. |\ifInBib| is set to
% true when we frind the |\begin{thebibliography}| and tells us to
% start saving the following lines. 
%    \begin{macrocode}
\newif\ifAnotherLine
\newif\ifInBib
%    \end{macrocode}
% \begin{macro}{makeOther}
% We need a macro to change the catcodes of lots of things to other. 
%    \begin{macrocode}
\def\makeOther#1{\catcode`#1=12\relax}
%    \end{macrocode}
% \end{macro}
% \begin{macro}{\RereadBiblio}
% Define a macro to do all the fiddly stuff. This is used at the
% beginning of |thebibliography| environment (Sec~\ref{sec:biblio}).
%    \begin{macrocode}
\newread\tf@biblio
\newcommand{\RereadBiblio}{%
%    \end{macrocode}
% Open the current file again for reading
%    \begin{macrocode}
  \begingroup
  \immediate\openin\tf@biblio=\articlename
%    \end{macrocode}
% Everything in a group to be safe. Initialize logicals. Black magic
% to change
% catcodes and avoid blank lines being interpreted as |\par|. 
%    \begin{macrocode}
  \AnotherLinetrue\InBibfalse
  \makeOther\ \makeOther\\\makeOther\$%
  \makeOther\#\makeOther\^\makeOther\^^K%
  \makeOther\_\makeOther\^^A\makeOther\%%
  \makeOther\~\makeOther\{\makeOther\}\makeOther\&%
  \endlinechar-1\relax
%    \end{macrocode}
% Start looping through the file. Read each line to the macro
% |\Line|. If we have got to EOF or end of biblio then set
% |\AnotherLinefalse|.
%    \begin{macrocode}
  \loop
  \read\tf@biblio to \Line
  \ifx\Line\EndBib\AnotherLinefalse\InBibfalse\fi
  \ifeof\tf@biblio \AnotherLinefalse\fi
%    \end{macrocode}
% If we are in the biblio and the line is not blank then append it to
% the |\ADS@references| macro together with a newline. 
%    \begin{macrocode}
  \ifInBib
  \ifx\Line\empty\else\xdef\ADS@references{\ADS@references\Line^^J}\fi
  \fi
%    \end{macrocode}
% Check if the line we just saw was the start of biblio, then repeat
% the loop if necessary. 
%    \begin{macrocode}
  \ifx\Line\StartBib\InBibtrue\fi
  \ifx\Line\StartBibX\InBibtrue\fi
  \ifAnotherLine\repeat
  \immediate\closein\tf@biblio
  \endgroup
  }
%    \end{macrocode}
% \end{macro}
% 
% \begin{macro}{\PadToFourWithDots}
% Utility function for left-padding a number to 4 characters with
% periods. Used in 
% constructing |\ADS@BibCode|. This is a very simple algorithm and
% \emph{requires} that the argument actually be a number---if it is
% not then dire consequences will result.  
%    \begin{macrocode}
\newcommand{\PadToFourWithDots}[1]{%
  \ifnum #1 > 999 #1%
  \else\ifnum #1 > 99 .#1%
  \else\ifnum #1 > 9 ..#1%
  \else ...#1%
  \fi\fi\fi}
%%  \ifnum \count0 > 999 #1%
%%  \else\ifnum \count0 > 99 .#1%
%%  \else\ifnum \count0 > 9 ..#1%
%%  \else ...#1%
%%  \fi\fi\fi}
%    \end{macrocode}
% \end{macro}
%
% \begin{macro}{\AuthorADS}
% \begin{macro}{\ADS@affiliations}
% New mechanism for setting the affiliations for ADS. The author labels (AA, AB, \dots) are generated using two nested counters.  We use the |\robustify| command from \texttt{etoolbox} so that accented characters in the affiliation addresses do not blow up.  The list of accents may need to be expanded. 
%    \begin{macrocode}
\newcounter{AUTHOR}
\newcounter{author}[AUTHOR]
\stepcounter{AUTHOR}
\newcommand\ADS@affiliations{}
\robustify{\'}\robustify{\i}\robustify{\~}\robustify{\`}
\def\AuthorADS#1#2{%
  \begingroup
  \def\ORCIDiD##1{<ID system="ORCID">##1</ID>}
  \def\Email##1{<EMAIL>##1</EMAIL>}
  \def\Affil##1{##1}
  \stepcounter{author}
  \ifthenelse{\value{author} > 26}{\stepcounter{AUTHOR}\stepcounter{author}}{}
  \ifthenelse{\equal{\ADS@affiliations}{\@rmaaEmpty}}{%
    \xdef\ADS@affiliations{\Alph{AUTHOR}\Alph{author}(#2)}
  }{%
    \xdef\ADS@affiliations{\ADS@affiliations, \Alph{AUTHOR}\Alph{author}(#2)}
  }%
  \endgroup
}
%    \end{macrocode}
% \end{macro}
% \end{macro}
% %
% \begin{macro}{\ADS@dumpdata}
% \changes{4.0}{2019/11/13}{Write all author affiliations for ADS}
% \begin{macro}{\ADS@write}
% \begin{macro}{\ADS@openfile}
% |\ADS@dumpdata| is the macro that actually does the job of writing out all the
% data to an auxiliary file, using |\ADS@write| to write each
% field. The file is set up with |\ADS@openfile|, which currently is
% done in the TOC of a book.  The new multi-author affiliations are used if |\ADS@affiliations| has been set, otherwise it falls back to the single-author version |\ADS@affil|. 
%    \begin{macrocode}
\newcommand{\ADSfield}[2]{\relax}
\newcommand{\ADSvfield}[2]{\relax}
%%\renewcommand{\ADSfield}[12]{\%#1 #2}
\newcommand{\ADS@dumpdata}
{ 
  \ADS@write{\string\begin{ADSarticle}}
  \ADS@WRITE{\string\ADSfield{R}{\ADS@BibCode}}
  \ADS@write{\string\ADSfield{A}{\ADS@authors}}
  \ifthenelse{\equal{\ADS@affiliations}{\@rmaaEmpty}}{%
    {\ADS@sanitize\ADS@write{\string\ADSfield{F}{AA(\ADS@affil)}}}
  }{%
    {\ADS@sanitize\ADS@write{\string\ADSfield{F}{\ADS@affiliations}}}
  }
  \ADS@write{\string\ADSfield{J}{\ADS@journalname}}
  \ADS@write{\string\ADSfield{D}{\ADS@fulldate}}
  \ADS@write{\string\ADSfield{L}{\ADS@lastpage}}
  {\ADS@sanitize\ADS@write{\string\ADSfield{T}{\ADS@title}}}
  \ADS@write{\string\ADSfield{C}{\ADS@copyright}}
  \ADS@write{\string\ADSfield{d}{\ADS@doi}}
  \ADS@write{\string\ADSfield{U}{\ADS@url}}
  \ADS@write{\string\ADSfield{K}{\ADS@keywords}}
  {\ADS@sanitize\ADS@write{\string\ADSfield{B}{\ADS@abstract}}}
  \ADS@write{\string\ADSfield{Z}{\ADS@references}}
  \ADS@write{\string\end{ADSarticle}}
  }
%%\newcommand{\ADS@write}[1]{\if@book\addtocontents{ads}{#1}\fi}
%%\newcommand{\ADS@openfile}{\@starttoc{ads}}
\newcommand{\ADS@write}[1]{\set@display@protect\immediate\write\tf@ads{#1}}
\newcommand{\ADS@WRITE}[1]{\immediate\write\tf@ads{#1}}
\newcommand{\ADS@openfile}{%
  \newwrite\tf@ads%
  \immediate\openout\tf@ads\jobname.ads\relax
  \ADS@WRITE{\ADS@preamble}}
\newcommand{\ADS@closefile}{\ADS@WRITE{\ADS@postamble}}
\def\ADS@preamble{%
  %%  \string\documentclass{article}
  %%  \string\begin{document}
  }
\def\ADS@postamble{%
  %%  \string\end{document}
  \string%\string% Local Variables:^J 
  \string%\string% mode: LaTeX^J
  \string%\string% mode: Font-lock^J
  \string%\string% End:^J
  }
%    \end{macrocode}
% \end{macro}
% \end{macro}
% \end{macro}
% 
% \subsection{Automatic calculation of the page range}
% \label{sec:page-range}
% Another issue that has cost me a lot of grief. Still, it finally
% seems to be sorted now. All hail Jeff Goldberg. 
%
% Cunning trick to get the page numbers from |\pageref|.  Put a label
% on the first and last page of the document using |\AtBeginDocument|
% and |\AtEndDocument|. .  This way won't work for the book option.
% In that case, we do it by defining |\@endofarticlehook| and
% |\@begofarticlehook| as parallels to |\AtEndDocument| and
% |\AtBeginDocument|.  Note that we have to defer adding
% |\@begofarticlehook| to |\AtBeginDocument| to make sure it comes
% after the redefinitions of |\ref| and |\label|.
%    \begin{macrocode}
\if@book\AtEndDocument{\ADS@closefile}\else
 %\AtBeginDocument{\@begofarticlehook}
\AtEndDocument{\@endofarticlehook}
\fi
%    \end{macrocode}
% \begin{macro}{\@begofarticlehook}
% Since we are getting more and more things that need to be done at
% the beginning of every article, frontsection, etc we define a sister
% hook to 
% |\@endofarticlehook|. This hook does the following: 
% \begin{enumerate}
% \item Labels the first page
% \item Sets |\ADS@firstpage|
% \item runs and then clears |\indexauthor@deferlist|
% \item resets the keyword counter
% \end{enumerate}
% It is executed in the following circumstances: 
% \begin{enumerate}
% \item At |\begin{document}| for standalone files
% \item At |\begin{document}| for each included file
% \item At beginning of front sections 
% \item At beginning of TOC
% \item At beginning of Author Index
% \item At beginning of abstracts
% \end{enumerate}
% For some reason, this is not always within the right group, so we
% use |\xdef| just to be sure. 
%    \begin{macrocode}
\newcommand{\@begofarticlehook}{%
  \if@frontmatter\frontlabel{RMAAFirstPage}\else\label{RMAAFirstPage}\fi%
  \xdef\ADS@firstpage{\@arabic{\c@page}}%
  \if@book\else\makeindex\fi % needed so default index is defined in stand-alone articles
%%  \indexauthor@deferlist
%%  \xdef\indexauthor@deferlist{}
  \@inpreamblefalse
  \setcounter{kwdcounter}{0}
  }
%    \end{macrocode}
% \end{macro}
% \begin{macro}{\@endofarticlehook}
% This is a hook to be run at the end of each article
% currently, it just labels the last page and writes out all the ADS
% data to an auxiliary file.  
% It should not produce any printed output. 
% It is modelled on the |\AtEndDocument| command, to which it reduces
% in the case of a single article. 
%    \begin{macrocode}
\newcommand{\@endofarticlehook}{%
  \if@debug{\ClassWarning{rmaa}{***Doing Hook***}}\fi
  \if@debug{\ClassWarning{rmaa}{ADSabstract is: \ADS@abstract}}\fi
  \clearpage\rmaa@lastpagelabel\if@book\ADS@dumpdata\fi
  \ifx\NotesList\@rmaaEmpty\else\OutputNotes\fi
  \if@debug{\ClassWarning{rmaa}{Notes list is: \NotesList}}\fi
}
%    \end{macrocode}
% \begin{macro}{\rmaa@lastpagelabel}
% New version of last page label command copied from Jeff Goldberg's
% lastpage.sty package - we can't just use |\label| here in case the
% output routine for the last page has already been called. Instead,
% we first explicitly clear the page, step back the page counter,
% explicitly write to the .aux file, then restep the page counter
% again.
%
% But this will conflict with the hyperref package, which does labels differently.  Compensate for that by using the five-argument form.
% 
% Also set the |\ADS@LastPage| macro to the appropriate number. 
%    \begin{macrocode}
\newcommand{\rmaa@lastpagelabel}{%
  \addtocounter{page}{-1}%
  \immediate\write\@auxout{\string
    \newlabel{\articlename:RMAALastPage}{{}{\thepage}{}{}{}}}%
  \edef\ADS@lastpage{\thepage}%
  \addtocounter{page}{1}}
%    \end{macrocode}
% \end{macro}
% \end{macro}

% \subsection{Deal with explicit change in number of columns}
% \label{sec:columns}
% Sometimes, when formatting a book, we need to switch between
% one-column articles and two-column articles. For this to work, we
% need to wrap the standard \LaTeX{} commands to set the |\@rmtwocol|
% flag appropriately and also reset the package default values of all
% scalable lengths. .  
%    \begin{macrocode}
\let\latexonecolumn=\onecolumn
\let\latextwocolumn=\twocolumn
\renewcommand{\onecolumn}{\@rmtwocolfalse\@package@default@lengths\latexonecolumn}
\renewcommand{\twocolumn}{\@rmtwocoltrue\@package@default@lengths\latextwocolumn}
%    \end{macrocode}
% 
% \subsection{The actual formatting of the title}
% \label{sec:title-format}

% Note, we have to redefine the |\maketitle| command to use the
% numbered footnotes. This means we needn't use the atbeginend stuff. 
%    \begin{macrocode}
\renewcommand{\maketitle}{\par
  \begingroup
%    \end{macrocode}
% Don't switch to marked footnotes. 
% This is consistent with existing style but I'm not sure I like it
%    \begin{macrocode}
  %%  \renewcommand\thefootnote{\fnsymbol{footnote}}%
%%  \def\@makefnmark{\hbox to\z@{$\m@th^{\@thefnmark}$\hss}}%
  \def\@makefnmark{\textsuperscript{\@thefnmark}\hspace{-0.2em}}%
  \long\def\@makefntext##1{\parindent 1em\noindent
    \hbox to1.8em{\hss$\m@th^{\@thefnmark}$}##1}%
  \if@rmtwocol
  { \if@debug{\ClassWarning{rmaa}{TWO COLUMN}}\fi
    \twocolumn[\@maketitle]%
    }
  \else
  { \if@debug{\ClassWarning{rmaa}{ONE COLUMN}}\fi
    \newpage
    \global\@topnum\z@   % Prevents figures from going at top of page.
    \if@manuscript \@MSmaketitle \else \@maketitle \fi
    }
  \fi
  \if@manuscript\else\thispagestyle{rmaafirstpage}\@thanks\fi
  \ifx\NotesList\@rmaaEmpty\else{%
    \NoteMark{1--\arabic{notecounter}}\vspace*{-2\baselineskip}%
  }\fi%
  \endgroup
%    \end{macrocode}
% Don't reset footnote counter (not sure whether this is good thing
% or not).
%    \begin{macrocode}
  %%  \setcounter{footnote}{0}%
  \let\thanks\relax\let\affil\relax\let\altaffil\relax
  \let\maketitle\relax\let\@maketitle\relax
  \gdef\@thanks{}\gdef\@author{}\gdef\@title{}}
%    \end{macrocode}

% \begin{macro}{\@kwhwidth}
% \begin{macro}{\@kwtwidth}
% Redefine \cs{@maketitle} to include the abstract, resumen and
% keywords. Note that in the manuscripty style we use
% \cs{@MSmaketitle} instead.  
% \changes{3.16}{2003/04/07}{
%  Added extra vertical space between authors and accepted date.}
% \changes{3.26}{2011/08/15}{Add extra label to review articles}
%    \begin{macrocode}
\newlength{\@kwhwidth}\newlength{\@kwtwidth}
\newcommand\reviewname{Review}
\renewcommand{\@maketitle}{%
  \newpage
  \null
  \if@reviewarticle
  \smallskip
  {\large\bfseries\itshape\reviewname}%
  \bigskip
  \else
  \vskip\@before@title@skip%
  \fi
  \begin{center}%
    {\@title \par}%
    \vskip\@before@author@skip%
    \parbox{\rmaa@AuthorWidthFrac\textwidth}{\centering
      \lineskip\@between@author@skip%
      \@author
      \par}
    %% this should be replaced with `Submitted to...' for preprints
    \if@printdates
    \vskip\@before@dates@skip%
    \if@preprint
    \rmaa@ReceivedStyle{Submitted to RMxAA \@received}
    \else
    \rmaa@ReceivedStyle{Received \@received; accepted \@accepted}
    \fi
    \fi
    \if@manuscript
    \rmaa@ReceivedStyle{Submitted to RMxAA \@received}
    \fi
    \if@guide
    \rmaa@ReceivedStyle{Version \rmaa@version, \rmaa@date}
    \fi
    \vskip\@before@abstract@skip%
    \if@abstractpresent
    \parbox{\rmaa@AbstractWidthFrac\textwidth}{%
      \setlength{\parindent}{0.9cm}%
      \centerline{\rmaa@AbstractHeaderStyle{Abstract}}\par%
      \vskip\@after@abstracthead@skip%
      \if@newstyle\noindent\fi\@abstract
      \vskip\@between@abstract@skip%
      \centerline{\rmaa@AbstractHeaderStyle{Resumen}}\par%
      \vskip\@after@abstracthead@skip%
      \if@newstyle\noindent\fi\@resumen
      \if@setresumen\else\@abstract\fi % append english abstract if
                                       % spanish one not set
      \if@guide\else            %do not use keywords for guide
      \vskip\@before@kwd@skip%
      % Format the Key Words with hanging indentation:
      % They are put in a \parbox with width equal to the line width
      % minus the width of the header. 
      { % Widths of keyword
        \settowidth{\@kwhwidth}{\@keywordheader~}   % header and
        \setlength{\@kwtwidth}{\linewidth}          % keyword text
        \addtolength{\@kwtwidth}{-\@kwhwidth}       %
        \noindent\@keywordheader~\parbox[t]{\@kwtwidth}{%
          \rmaa@KeywordStyle{\boldmath\uppercase{\@keywords}}}}%
      \fi
      }
      \fi
  \end{center}%
  \vspace*{\afterkwdskip}
  \par
  }
\newlength{\afterkwdskip}\setlength{\afterkwdskip}{\z@}
%    \end{macrocode}
% \end{macro}
% \end{macro}
%
%
% \begin{macro}{\reviewarticle}
% \begin{macro}{\reviewtoc}
% \begin{macro}{\reviewlof}
%   User command to specify a review article. Must be used in
%   preamble. Main effect is to place the text ``Review'' before the
%   article title.  Also add command for setting a table of contents
%   and a list of figures at the beginning of the article.
%    \begin{macrocode}
\newcommand{\reviewarticle}{\@reviewarticletrue}
\newcommand\reviewtoc{%
  \section*{CONTENTS}%
  \@starttoc{rtc}%
}
\newcommand\reviewlof{%
  \section*{LIST OF FIGURES}%
  \@starttoc{lof}%
}
%    \end{macrocode}
% \end{macro}
% \end{macro}
% \end{macro}
% \begin{macro}{\l@section}
%    \begin{macrocode}
\renewcommand*\l@section[2]{%
  \ifnum \c@tocdepth >\z@
    \addpenalty\@secpenalty
    \addvspace{0.8em \@plus\p@}%
    \setlength\@tempdima{1.5em}%
    \begingroup
      \parindent \z@ \rightskip \@pnumwidth
      \parfillskip -\@pnumwidth
      \leavevmode %\bfseries
      \advance\leftskip\@tempdima
      \hskip -\leftskip
      #1\nobreak\hfil \nobreak\hb@xt@\@pnumwidth{\hss #2}\par
    \endgroup
  \fi}
%    \end{macrocode}
% \end{macro}
%
% \begin{macro}{\@MSmaketitle}
% A special version of |\@maketitle| for the \texttt{manuscript}
% option. This can deal with a multi-page abstract/resumen. Not very
% well, mind. 
%    \begin{macrocode}
\newcommand{\@MSmaketitle}{%
  \if@debug\ClassInfo{rmaa}{Using special MS maketitle command}\fi
  \newpage
  {
    \thispagestyle{rmaafirstpage}
%     \setlength{\parindent}{0pt}
    \null\vskip 2ex
    \begin{center}
      \@title\par
      \@author\par
      \rmaa@ReceivedStyle{%
        \ifx\@received\@rmaaEmpty
          Draft version: \@date
        \else
          Received: \@received
          \ifx\@accepted\@rmaaEmpty
          \else
            ; Accepted: \@accepted
          \fi
        \fi}\par
    \end{center}
    \@thanks
    \centerline{\rmaa@AbstractHeaderStyle{Resumen}}
    \noindent\@resumen\par\vskip 2ex
    \centerline{\rmaa@AbstractHeaderStyle{Abstract}}
    \noindent\@abstract\par\vskip 2ex
    \noindent\@keywordheader~\rmaa@KeywordStyle{\@keywords}\par
%     \@keywordheader~\rmaa@KeywordStyle{\uppercase{\@keywords}}\par
  }%
  }
%    \end{macrocode}
% \end{macro}

% \subsection{The sectioning commands}
% \label{sec:section-commands}

% \begin{macro}{\section}
% \begin{macro}{\subsection}
% \begin{macro}{\subsubsection}
% \begin{macro}{\paragraph}
% \begin{macro}{\subparagraph}
% Redefine all the standard commands to use these lengths. 
%    \begin{macrocode}
\renewcommand\section{\@startsection {section}{1}{\z@}%
                                   {\abovesectionskip}%
                                   {\belowsectionskip}%
                                   %% how do we get upper case here
                                   {\centering\reset@font\normalsize}}
\renewcommand\subsection{\@startsection{subsection}{2}{\z@}%
                                     {\abovesectionskip}%
                                     {\belowsectionskip}%
                                     {\centering\reset@font\normalsize\itshape}}
\renewcommand\subsubsection{\@startsection{subsubsection}{3}{\z@}%
                                     {\abovesectionskip}%
                                     {\belowsectionskip}%
                                     {\centering\reset@font\normalsize\itshape}}
\renewcommand\paragraph{\@startsection{paragraph}{4}{\z@}%
                                    {\aboveparaskip}%
                                    {\belowparaskip}%
                                    {\reset@font\normalsize\bfseries}}
\renewcommand\subparagraph{\@startsection{subparagraph}{5}{\parindent}%
                                      {\aboveparaskip}%
                                      {\belowparaskip}%
                                      {\reset@font\normalsize\bfseries}}
%    \end{macrocode}
% \end{macro}
% \end{macro}
% \end{macro}
% \end{macro}
% \end{macro}

% \begin{macro}{secnumdepth}
% Make sure that we number the subsubsections---this is done by
% default in the |article| class but not the |book| class, which is
% why we must set it explicitly here.
%    \begin{macrocode}
\setcounter{secnumdepth}{3}
%    \end{macrocode}
% \end{macro}
% \changes{2.2beta10}{2001/11/07}{Fixed \cs{secnumdepth} for \cs{book}
% style.}

% \begin{macro}{\thesubsection}
% \begin{macro}{\thesubsubsection}
% For subsection and subsubsection, we want the numbers to be roman
% but the title to be italic. This seems the easiest way to go about it.
%    \begin{macrocode}
\renewcommand{\thesubsection}{{\upshape\thesection.\@arabic\c@subsection}}
\renewcommand{\thesubsubsection}{\thesubsection{\upshape.\@arabic\c@subsubsection}}
%    \end{macrocode}
% \end{macro}
% \end{macro}

% \begin{macro}{\@seccntformat}
% Add ``.'' to the number here so we can remove it from |\thesection|,
% etc. Reduce gap between number and title WJH 27 Jan 1995
%    \begin{macrocode}
\def\@seccntformat#1{\csname the#1\endcsname .\hskip 0.4em}
%    \end{macrocode}
% \end{macro}

% \begin{macro}{\@sect}
% I didn't really want to mess with the |\@sect| command, but it was
% the only way I could manage to force upper case. Other than the
% commented parts, it is identical to the standard version. 
%    \begin{macrocode}
\def\@sect#1#2#3#4#5#6[#7]#8{
  \ifnum #2>\c@secnumdepth
     \let\@svsec\@empty
  \else
     \refstepcounter{#1}%
     \let\@@protect\protect
     \def\protect{\noexpand\protect\noexpand}%
     \edef\@svsec{\@seccntformat{#1}}%
     \let\protect\@@protect
  \fi
%    \end{macrocode}
% This is the bit I changed (27 Jan 1995). Use upper case for
% |\section|, use mixed case for the rest. 
%    \begin{macrocode}
  \ifnum #2=1                   
     \def\sec@upcase##1{\uppercase{##1}}%
     \def\sec@shape{\upshape}%
  \else                         
     \def\sec@upcase##1{\relax{##1}}%
     \def\sec@shape{\itshape}%
  \fi
  \@tempskipa #5\relax
  \ifdim \@tempskipa>\z@
     \begingroup #6\relax
          \@hangfrom{\hskip #3\relax\@svsec}%
                    {\interlinepenalty \@M \sec@upcase{#8}\par}%
     \endgroup
     \csname #1mark\endcsname{#7}%
%    \end{macrocode}
% Only write TOC entries if doing the Guide style. 
%    \begin{macrocode}
     \if@guide 
        \addcontentsline{toc}{#1}{%
          \ifnum #2>\c@secnumdepth \else
             \protect\numberline{\csname the#1\endcsname}%
          \fi
          #7}%
     \fi
%    \end{macrocode}
% Also allow per-article TOC for review papers
%    \begin{macrocode}
     \if@reviewarticle
        \ifnum #2>2 \else
           \addcontentsline{rtc}{#1}{%
             \protect\numberline{\csname the#1\endcsname}%
             \sec@upcase{\sec@shape #7}}%
        \fi
     \fi
  \else
     \def\@svsechd{#6\hskip #3\relax  %% \relax added 2 May 90
       \@svsec #8\csname #1mark\endcsname{#7}%
       \if@guide
          \addcontentsline{toc}{#1}{%
            \ifnum #2>\c@secnumdepth \else
               \protect\numberline{\csname the#1\endcsname}%
            \fi
            #7}%
       \fi
       }%
  \fi
  \@xsect{#5}}
%    \end{macrocode}
% \end{macro}


% \subsection{Appendices}
% \label{sec:appendices}
% Make sure that appendices in individual articles work properly. We
% copy the definition from the article class (the book one doesn't do
% what we want).  
% \begin{environment}{appendices}
%   This uses an environment |appendices| inside of which each
%   |\section| will become a separate appendix.  Note that the
%   appendices must come before the references (otherwise the full
%   addresses mechanism will not work properly). We make the equation
%   numbering be subordinate to the appendix (equations A12, B6,
%   etc). Finally, we have a medskip followed by APPENDICES to
%   introduce them. \changes{3.17}{2002/04/08}{ Added appendices environment
%   that does the right thing.}
%    \begin{macrocode}
\newenvironment{appendices}{\par
  \begingroup
  \setcounter{section}{0}%
  \setcounter{subsection}{0}%
  \renewcommand\thesection{\@Alph\c@section}
  \renewcommand\theequation{\thesection\arabic{equation}}
  \medskip
  \section*{APPENDICES}}
{\endgroup}
%    \end{macrocode}
% \end{environment}
% \begin{environment}{appendix}
%   If there is only one appendix, then use |appendix| environment
%   instead.  \changes{4.0}{2019/11/12}{ Added
%   appendix environment for case where there is only one section.}
%    \begin{macrocode}
\renewenvironment{appendix}{\par
  \begingroup
  \setcounter{section}{0}%
  \setcounter{subsection}{0}%
  \renewcommand\thesection{\@Alph\c@section}
  \renewcommand\theequation{\thesection\arabic{equation}}
  \medskip
  \section*{APPENDIX}}
{\endgroup}
%    \end{macrocode}
% \end{environment}

% \subsection{Floating thingies}
% \label{sec:floats}   

%% 
%    \begin{macrocode}
\renewcommand{\thetable}{\@arabic\c@table}
\renewcommand{\thefigure}{\@arabic\c@figure}
\renewcommand{\figurename}{Fig.}
\renewcommand{\tablename}{TABLE}
%\renewcommand{\fnum@table}{\MakeUpperCase{\tablename}~\thetable}
%    \end{macrocode}
% For some reason, I have to write `TABLE' explicitly here
% If I use |\tablename|, it gets converted to mixed case---God knows
% why. 
% |\figurename| is also being ignored. 
%    \begin{macrocode}
\renewcommand{\fnum@table}{TABLE~\thetable}
\renewcommand{\fnum@figure}{Fig.~\thefigure}
%    \end{macrocode}

% Do captions in 9 pt face in the main journal style and in new-style
% anything. 
%    \begin{macrocode}
\long\def\@makefigcap#1#2{%
  \if@proceedings\else\small\fi
  \if@newstyle\small\fi 
  \vskip\abovecaptionskip
  \sbox\@tempboxa{#1. #2}%
  \ifdim \wd\@tempboxa >\hsize
    #1. #2 \par
  \else
    \global \@minipagefalse
    \hb@xt@\hsize{\hfil\box\@tempboxa\hfil}%
  \fi
  \vskip\belowcaptionskip}
%    \end{macrocode}

% Make sure table caption is upper case. 
%    \begin{macrocode}
\long\def\@maketabcap#1#2{%
  \hbox to\hsize{\hfil #1 \hfil}
  \vskip 6pt 
  \centering\MakeUpperCase{#2}\par
  \vskip 6pt}
%    \end{macrocode}

%    \begin{macrocode}
\newlength{\@rottabcapwidth}
\long\def\@makerottabcap#1#2{%
  \setbox\@tempboxa\hbox{\MakeUpperCase{#2}}%
  \ifdim\wd\@tempboxa >.8\vsize
     \setlength\@rottabcapwidth{.8\vsize}
  \else
     \setlength\@rottabcapwidth{\wd\@tempboxa}
  \fi
  \rotatebox{90}{%
    \begin{minipage}{\vsize}\centering
      \begin{minipage}{\@rottabcapwidth}\centering
        #1\\ \vskip 6pt \MakeUpperCase{#2}
      \end{minipage}%
    \end{minipage}
    }%
  \hskip 6pt
  }

\newenvironment{Sideways}
    {\begin{sideways}\begin{minipage}{\vsize}\centering}
    {\end{minipage}\end{sideways}}

\newenvironment{Sidewaystable}
{%
  \let\oldcaption=\caption%
  \def\caption{
    \def\@captype{table}\renewcommand{\@makecaption}{\@maketabcap}
    \oldcaption}
  \begin{sidewaystable}\centering}
  {\end{sidewaystable}\let\caption=\oldcaption}
\newenvironment{Sidewaystable*}
{%
  \let\oldcaption=\caption%
  \def\caption{
    \def\@captype{table}\renewcommand{\@makecaption}{\@maketabcap}
    \oldcaption}
  \begin{sidewaystable*}\centering}
  {\end{sidewaystable*}\let\caption=\oldcaption}
%    \end{macrocode}

% Not sure why we need to stick these lines in. 
%    \begin{macrocode}
\renewenvironment{figure}
                 {\renewcommand{\@makecaption}{\@makefigcap}\@float{figure}}
                 {\end@float}
\renewenvironment{figure*}
                 {\renewcommand{\@makecaption}{\@makefigcap}\@dblfloat{figure}}
                 {\end@dblfloat}
\renewenvironment{table}
                 {\renewcommand{\@makecaption}{\@maketabcap}\@float{table}}
                 {\end@float}
\renewenvironment{table*}
                 {\renewcommand{\@makecaption}{\@maketabcap}\@dblfloat{table}}
                 {\end@dblfloat}
%    \end{macrocode}

% Redefine |\@caption| so that we don't bother with the list of
% figures. This saves us having to worry about fragile commands in the
% captions. 
%    \begin{macrocode}
\long\def\@caption#1[#2]#3{%
  \par
%%  \addcontentsline{\csname ext@#1\endcsname}{#1}%
%%    {\protect\numberline{\csname the#1\endcsname}{\ignorespaces #2}}%
  \begingroup
    \@parboxrestore
    \if@minipage
      \@setminipage
    \fi
    \normalsize
    \@makecaption{\csname fnum@#1\endcsname}{\ignorespaces #3}\par
  \endgroup}
%    \end{macrocode}

% Make tables not quite so cramped. Allow |\hline\hline| to give a
% poor man's thick rule. 
%    \begin{macrocode}
\renewcommand{\arraystretch}{1.15}
\setlength{\doublerulesep}{0.4pt}
%    \end{macrocode}

% \begin{macro}{\tablecols}
% \begin{macro}{\@tablecols}
% \begin{macro}{\tabcaptioncontinued}
% |\tablecols| is a user command for setting the number of columns. It
% is used by commands that try and span the entire table (I think
% there must be a better way of acheiving
% this). |\tabcaptioncontinued| is 
% a user command for setting the caption of a continued table on the
% next page. 
%    \begin{macrocode}
\newcommand{\tablecols}[1]{\renewcommand{\@tablecols}{#1}}
\newcommand{\@tablecols}{1}
\newcommand{\tabcaptioncontinued}{%
  \multicolumn{\@tablecols}{c}{\addtocounter{table}{-1}TABLE~\thetable\ (CONTINUED)} \\ \\
  }
%    \end{macrocode}
% \end{macro}
% \end{macro}
% \end{macro}

% \subsubsection{Footnotes in tables}
% \label{sec:table-footnotes}
% \begin{macro}{\tablenotemark}
% \begin{macro}{\tablenote}
% \begin{macro}{\tablenoteskip}
% Old user commands |\tablenotemark{MARK}|,
% |\tablenote{MARK}{TEXT}|---now deprecated. 
%    \begin{macrocode}
\newcommand\tablenotemark[1]{$^\mathrm{#1}$} 
\newcommand\tablenote[2]{\if@proceedings\else\small\fi ${}^\mathrm{#1}$ #2} 
\newcommand\tablenoteskip{\\[0.5\baselineskip]} 
%    \end{macrocode}
% \end{macro}
% \end{macro}
% \end{macro}

% \begin{macro}{\tabnotemark}
% \begin{macro}{\tabnotetext}
% \begin{macro}{\tabnotetextTWO}
% Alternative way of doing footnotes in tables. 
%    \begin{macrocode}
\newcommand{\tabnotemark}[1]{\textsuperscript{\lowercase{#1}}}
%\newcommand{\tabnotetextTWO}[4]{\tabnotetext{#1}{#2}\tabnotetext{#3}{#4}}
\newcommand{\tabnotetextTWO}[4]{%
  \multicolumn{\@tablecols}{l}{%
    \parbox{0.5\tabnotewidth}{\small\textsuperscript{#1}#2}%
    \parbox{0.5\tabnotewidth}{\small\textsuperscript{#3}#4}%
    }%
  \\}
\newcommand{\tabnotetext}[2]{%
  \multicolumn{\@tablecols}{@{}p{\tabnotewidth}@{}}{\small\textsuperscript{#1}#2}\\}
\newlength{\tabnotewidth}
%    \end{macrocode}
% \end{macro}
% \end{macro}
% \end{macro}


% \begin{macro}{changemargin}
%    \begin{macrocode}
% Command to allow margins to be changed on the fly. Useful to allow
% tables and Figures to protrude into the margins. 
%
% \changes{3.24}{2007/09/14}{Added \texttt{changemargin} environment
% (copied from TeX FAQ).}
\newenvironment{changemargin}[2]{% 
  % From TeX FAQ http://www.tex.ac.uk/cgi-bin/texfaq2html?label=chngmargonfly
  \begin{list}{}{%
    \setlength{\topsep}{0pt}%
    \setlength{\leftmargin}{#1}%
    \setlength{\rightmargin}{#2}%
    \setlength{\listparindent}{\parindent}%
    \setlength{\itemindent}{\parindent}%
    \setlength{\parsep}{\parskip}%
  }%
  \item[]}{\end{list}}
%    \end{macrocode}
% \end{macro}

% \subsection{The bibliography and references}
% \label{sec:biblio}
% \begin{macro}{\beforebiblioskip}
% \begin{macro}{\afterbiblioskip}
% I'm not entirely sure why I introduced these two commands, since
% directly putting |\vspace| before and after
% |\begin{thebibliography}| seems to work well enough. 
%    \begin{macrocode}
\newlength{\beforebiblioskip}
\setlength{\beforebiblioskip}{0pt}
\newlength{\afterbiblioskip}
\setlength{\afterbiblioskip}{0pt}
%    \end{macrocode}
% \end{macro}
% \end{macro}
% \begin{macro}{thebibliography}
% The environment itself is done as starred section with extra
% vertical space after heading. This space is currently zero. 
% In the book style we do the cunning trick with rereading the file in
% order to get the references into the |\ADS@references| macro.
% 
% \changes{3.05}{2003/01/02}{Removed the \cs{beforebiblioskip} and
% \cs{afterbiblioskip} commands, since they cause unwanted vertical
% space in some instances.}  
%    \begin{macrocode}
\renewenvironment{thebibliography}{%
  \if@book\RereadBiblio\fi
%%  \vspace*{\beforebiblioskip}
  \section*{REFERENCES} \frenchspacing
%%  \vspace*{\afterbiblioskip}
  \small
%    \end{macrocode} 
% The references themselves are set as a list in order 
% to get the hanging indentation. They are done in a smaller font and
% with french spacing so we needn't worry about the distinction
% between sentence-ending periods and abbreviations. 
%    \begin{macrocode} 
  \list{\null}{\leftmargin 1.5em\labelwidth\z@\labelsep\z@\itemindent -1.5em
    \itemsep\z@\parsep\z@\usecounter{enumi}}
  \sloppy\clubpenalty4000\widowpenalty4000
  \sfcode`\.=1000\relax\@biblisttrue}
{\endlist\@biblistfalse\normalsize\@outputaddressesifnecessary}
%    \end{macrocode}
% \end{macro}
% \subsection{Uniquification of cross-references}
% \label{sec:uniquification}
% 
% \begin{macro}{\articlename}
% In order that the |\label{}| command in individual articles be
% guaranteed to give a unique anchor, we append the article name (held
% in the |\articlename| command) to each key. This requires
% redefinition of the following commands: |\cite|, |\nocite|, |\ref|,
% |\pageref|, |\label|. 
% \begin{macro}{\oldcite}
% \begin{macro}{\oldnocite}
% \begin{macro}{\oldref}
% \begin{macro}{\oldpageref}
% \begin{macro}{\oldlabel}
% The original versions of these commands are
% saved as |\oldcite|, |\oldnocite|, etc. 
%
% The uniquifying of cross-references is now always done, so 
% we don't get multiply-defined labels in the single article case
%    \begin{macrocode}
\newcommand{\articlename}{anon}
\AtBeginDocument{
  \let\oldcite=\cite 
  \let\oldnocite=\nocite 
  \let\oldref=\ref 
  \let\oldpageref=\pageref
  \let\oldlabel=\label
}
%    \end{macrocode}
% \end{macro}
% \end{macro}
% \end{macro}
% \end{macro}
% \end{macro}
% \end{macro}
%
% Uniquify all cross-references by prepending the article name.  This
% now has to go into |\AtBeginDocument| so it is working on the new
% versions that hyperref.sty has set up.  Then we need to put the
% |\@begofarticlehook| afterwards, so that the first page label is using
% the new definitions.  If amsmath is loaded, then we also need to patch
% the special label command that they use. \changes{4.0}{2019/11/18}{Fix labels for amsmath display equations, in case that package has been loaded.}
%    \begin{macrocode}
\AtBeginDocument{
  \renewcommand{\ref}[1]{\oldref{\articlename:#1}%
    \if@debug{\ClassInfo{rmaa}{Reference to \articlename:#1}}\fi}
  \renewcommand{\pageref}[1]{\oldpageref{\articlename:#1}%
    \if@debug{\ClassInfo{rmaa}{Page reference to \articlename:#1}}\fi}
  \renewcommand{\label}[1]{\oldlabel{\articlename:#1}%
    \if@debug{\ClassInfo{rmaa}{Label for \articlename:#1}}\fi}%
  \ifdef{\label@in@display}{%
    \let\oldlabel@in@display=\label@in@display
    \renewcommand{\label@in@display}[1]{\oldlabel@in@display{\articlename:#1}%
      \if@debug{\ClassInfo{rmaa}{AMSTeX Label for \articlename:#1}}\fi}%
  }{}
  \@begofarticlehook
}%
%    \end{macrocode}
% Citations are more complicated since they can contain a list of
% keys. These are really only here for use with the rmaacite package,
% which is utterly deprecated (02 Jan 2003) now that I have the natbib
% code working. The command \cs{cite} should \emph{not} be used with
% natbib: use \cs{citep}, \cs{citet}, etc.\@ instead. 
%    \begin{macrocode}
\renewcommand{\nocite}[1]{\@bsphack
  \@for\@citeb:=#1\do{%
    \edef\@citeb{\expandafter\@firstofone\@citeb:\articlename}%
    \if@filesw\immediate\write\@auxout{\string\citation{\@citeb}}\fi
    \@ifundefined{b@\@citeb:\articlename}{\G@refundefinedtrue
      \@latex@warning{Citation `\@citeb' undefined}}{}}%
  \@esphack}
\renewcommand{\cite}[1]{\oldcite{#1:\articlename}}
%   \end{macrocode}


% Horrible fiddling with the bibliography citation mechanism. 
% This has to be globally available in case rmaacite puts it in the
% \texttt{.aux} file. 
%    \begin{macrocode}
\def\lbibcite#1#2#3{\global\@namedef{b@#1}{#2}
                    \global\@namedef{by@#1}{#3}}
%    \end{macrocode}

% \begin{macro}{\oldbibitem}
% \begin{macro}{\@rmaaBibArg}
% \begin{macro}{\@rmaaSpace}
% \begin{macro}{rmaabib}
% Of course, |\bibitem| needs to get the uniquification treatment as
% well. While we are at it, we redefine it to use consecutive integers
% for the key if the user leaves it blank (using the counter
% |rmaabib|). This allows one to put 
% |\bibitem{}| for all your references without getting annoying
% messages about multiply-defined labels.
% 
% We also take the opportunity to insert the optional argument to
% \cs{bibitem} if it is missing in the form Author(DATE) so that natbib
% doesn't try to do numerical citations. Natbib already provides a
% hook \cs{@extra@b@citeb} to help uniquify the citation key, so we
% just plug into that. However, it puts it \emph{after} the original
% key, not before as I had had it. Easiest solution is to change the
% way that I do it. 02 Jan 2003. Note that the \cs{@extra@b@citeb}
% does not get written to the .aux  file in the \cs{citation} command
% but it does work. 
% \changes{3.14}{06 Feb 2003}{Reinstalled old definition of \cs{bibitem}
%  for the case where the nonatbib option is requested, so that the
%  rmaacite package will work. This means that you can't use rmaacite
%  and natbib in the same volume.}
%    \begin{macrocode}
\if@natbib\def\@extra@b@citeb{:\articlename}\fi
%%\else
\let\oldbibitem\bibitem
\newcommand{\@rmaaBibArg}{}
\newcommand{\@rmaaSpace}{ }
\newcounter{rmaabib}
\if@natbib
\renewcommand{\bibitem}[2][ABC(2003)]{%
  \renewcommand{\@rmaaBibArg}{#2}
  \ifx\@rmaaBibArg\@rmaaEmpty
     \addtocounter{rmaabib}{1}
     \oldbibitem[#1]{\thermaabib:\articlename}
  \else\ifx\@rmaaBibArg\@rmaaSpace
     \addtocounter{rmaabib}{1}
     \oldbibitem[#1]{\thermaabib:\articlename}
  \else 
     \oldbibitem[#1]{#2:\articlename}
  \fi\fi
  \renewcommand{\@rmaaBibArg}{}
  }
\else
\renewcommand{\bibitem}[1]{%
  \renewcommand{\@rmaaBibArg}{#1}
  \ifx\@rmaaBibArg\@rmaaEmpty
     \addtocounter{rmaabib}{1}
     \oldbibitem{\articlename:\thermaabib}
  \else\ifx\@rmaaBibArg\@rmaaSpace
     \addtocounter{rmaabib}{1}
     \oldbibitem{\articlename:\thermaabib}
  \else
     \oldbibitem{\articlename:#1}
  \fi\fi
  \renewcommand{\@rmaaBibArg}{}
  }
\fi
%%\fi
%    \end{macrocode}
% \end{macro}
% \end{macro}
% \end{macro}
% \end{macro}

% \begin{macro}{oldpcite}
% \begin{macro}{oldscite}
% Then, we need to repeat the above scheme, but for the |rmaacite|
% definition. 
%    \begin{macrocode}
\newcommand{\@redefinebibitem}{%
  \let\oldbibitem\bibitem
  \let\oldcite\cite
  \let\oldpcite\pcite
  \let\oldscite\scite
  \def\bibitem[##1<##2>]##3{%
    \renewcommand{\@rmaaBibArg}{##3}
    \ifx\@rmaaBibArg\@rmaaEmpty
       \addtocounter{rmaabib}{1}
       \oldbibitem[##1<##2>]{\thermaabib:\articlename}
    \else\ifx\@rmaaBibArg\@rmaaSpace
       \addtocounter{rmaabib}{1}
       \oldbibitem[##1<##2>]{\thermaabib:\articlename}
    \else 
       \oldbibitem[##1<##2>]{##3:\articlename}
    % trust the author here
%%       \oldbibitem[##1<##2>]{##3}
    \fi\fi
    }
  }
%    \end{macrocode}
% \end{macro}
% \end{macro}


% \subsection{HTML Table of Contents}
% \label{sec:hyper-TOC}
% \begin{macro}{\HyperTOCEntry}
% \begin{macro}{\HyperTOCpart}
% \begin{macro}{\HyperTOCmeeting}
% \begin{macro}{\HyperTOCeditors}
% \begin{macro}{\HyperTOCpreamble}
% These commands are used to write information to the \texttt{.htoc}
% file, which is used in generating a Web Page of the Table of
% Contents. Since they only ever appear preceded by |\protect| in the
% argument to an |\addtocontents| command,  they should never need to
% be expanded by the macros. However, we define them here just in case
% \dots
%    \begin{macrocode}
\newcommand{\HyperTOCEntry}[3]{\relax}
\newcommand{\HyperTOCpart}[1]{\relax}
\newcommand{\HyperTOCmeeting}[3]{\relax}
\newcommand{\HyperTOCeditors}[1]{\relax}
\newcommand{\HyperTOCpreamble}{\relax}
%    \end{macrocode}
% \end{macro}
% \end{macro}
% \end{macro}
% \end{macro}
% \end{macro}
%

% \subsection{Formatting the poster abstracts} 
% These can now be used without using the book option.
%    \begin{macrocode}

\newcommand{\@makeminititle}{%
  \newpage
  \null 
  \vskip 3em%
  \begin{center}%
    {\@title \par}%
  \end{center}
  \vskip 3em%
  }

\newif\if@abstractsads
\@abstractsadsfalse
\newcommand\abstractsADS{\@abstractsadstrue}

\if@book\RequirePackage{multicol}\fi
\newenvironment{abstracts}[1][Abstracts of Contributed Papers]{%
  %% First, make sure we start on an odd page (06 Jun 2000)
  \if@book\clearpage\thispagestyle{empty}\cleardoublepage\fi
  %%
  \renewcommand{\articlename}{ABSTRACTS}
  \title{#1}\shortauthor{Abstracts}\shorttitle{Abstracts}
  \listofauthors{}
  \if@book
  \addcontentsline{toc}{chapter}{\protect\@fiddleauthors{\@TOCtitle}{\rmaaTOCauthors}}%
  \addtocontents{htc}{\protect\HyperTOCEntry{\@rawtitle}{\@rawauthors}{\articlename}}
  \fi
  \pagestyle{rmaa} % this now seems necessary (06 Jun 2000)
  \thispagestyle{rmaafirstpage}
%%  \twocolumn[\@makeminititle]
  \renewcommand{\abstract}[1]{%
    \renewcommand\ADS@abstract{##1}%
    \renewcommand{\@abstract}{##1}\@writeabstract%
  }
  \renewcommand{\title}[2][]{%
    \gdef\ADS@authors{}%
    \gdef\ADS@affil{}%
    \gdef\ADS@references{}%
    \gdef\ADS@keywords{}%
    \renewcommand{\@rmaaArg}{##1}
    \ifx\@rmaaArg\@rmaaEmpty
       \renewcommand{\@TOCaftertitle}{}
    \else
       \renewcommand{\@TOCaftertitle}{ (\@rmaaArg)}
    \fi
    \renewcommand{\@title}{\rmaa@TOCTitleStyle{##2}}%
    \renewcommand{\@TOCtitle}{\rmaa@TOCTitleStyle{##2}\@TOCaftertitle}%
    \renewcommand{\rmaaIndexTitle}{\rmaa@IndexTitleStyle{##2\@TOCaftertitle.}}%
    \renewcommand\ADS@title{##2}%
    \xdef\ADS@firstpage{\@arabic{\c@page}}%
    }
  \renewcommand{\author}[1]{\let\@thanks\@empty\renewcommand\@author{##1} }
  \renewcommand{\footnotetext}[2][XX]{\par\noindent\textsuperscript{##1}~##2}
  \renewcommand{\footnoterule}{\kern-3\p@ \hrule \@width 3em \kern 2.6\p@}
  \if@book\@begofarticlehook\fi
  \if@book\begin{multicols}{2}[\@makeminititle]\else\@makeminititle\fi
  \begingroup
  \setcounter{footnote}{0}
  }% and at the end of the environment....
  {\endgroup\if@book\end{multicols}\listofabstracts\label{RMAALastPage}\clearpage\fi}
%% auxiliary commands used:
%% \@writeabstract does all the actual formatting 
\newcommand{\@writeabstract}{%
  \if@book%
  \addcontentsline{loa}{abstract}{\protect\@fiddleauthors{\@TOCtitle}{\rmaaTOCauthors}}%
  \fi%
  \begin{center}\@title\par\@author\end{center}
  \@abstract
  \par\noindent\footnoterule\nobreak\footnotesize\@thanks
  \par\noindent\hrulefill\goodbreak
  \vskip 1cm \@plus0.5cm
  %% write ADS entry
  \if@book\if@abstractsads\ADS@dumpdata\fi\fi%
  \endgroup % this pair is to make each abstract a self-contained group
  \begingroup
  \setcounter{footnote}{0}
  }

\if@book
%% List of abstracts (stuff written to .loa file)
\newcommand{\listofabstracts}{%
  \begin{multicols}{2}[%
    \vskip 1.5cm \vfill % leave at least 1.5 cm and go to bottom of page if possible
    \centerline{\MakeUpperCase{List of Abstracts}}
    \vskip 1.5cm \@plus0.5cm
    ][6cm] % this makes it start a new page if there is not at least 6cm of space
    \renewcommand{\@fiddleauthors}[2]{##1\hskip1em plus2em minus0.4em\penalty-100  ##2}
    \parindent\z@
%%    \parskip.5\p@ \@plus .7\p@\relax 
    \parskip3\p@ \@plus 1\p@ \@minus 1\p@\relax % allow quite a lot of
                                % leeway here to avoid breaks inside
                                % entries 
    \renewcommand{\\}{\relax} % temporarily turn off explicit line breaks just in case
    \renewcommand{\altaffilmark}[1]{\relax} % ditto with this
    \@starttoc{loa}
  \end{multicols}
  }
\newcommand*{\l@abstract}[2]{%
  \par\penalty-1000
  #1\hskip 1em\penalty0\null\nobreak\hfill#2}
\fi

%    \end{macrocode}
 
% \subsection{Code for typesetting an entire volume of conference
% proceedings} \label{sec:whole-volume}
%    \begin{macrocode}

\if@book
%% environment for including the individual articles
\newcommand{\@fiddleauthors}[2]{\relax}
\renewcommand{\theequation}{\@arabic\c@equation}
\newenvironment{articles}%
{ \begingroup % make it a group so that author definitions do not
              % spill over
  \gdef\ADS@authors{}
  \pagestyle{rmaa}
  \renewenvironment{document}{%
    \@begofarticlehook
    \addtocounter{chapter}{1}%
    \addcontentsline{toc}{chapter}{%
      \protect\@fiddleauthors{\@TOCtitle}{\rmaaTOCauthors}}%
    \addtocontents{htc}{\protect\HyperTOCEntry{\@rawtitle}{\@rawauthors}{\articlename}}
    }{\relax}
  \renewcommand{\documentclass}[2][foo]{\relax} 
  \renewcommand{\usepackage}[2][foo]{%
    \def\@packagename{##2}\def\@rmaacitename{rmaacite}%
    \ifx\@packagename\@rmaacitename
    \@rmcitetrue\ClassInfo{rmaa}{This paper uses rmaacite}
    \makeatletter\input{rmaacite.sty}\makeatother % read it in here so that all changes
    \else                                         % are local to paper
    \if@debug\ClassWarning{rmaa}{Package \@packagename\space was requested by \articlename\space}\fi
    \fi} 
  %% Don't allow individual articles to mess with the pages, volume or year
  \renewcommand{\SetFirstPage}[1]{\relax}
  \renewcommand{\SetYear}[1]{\relax}
  \renewcommand{\SetVolume}[1]{\relax}
  \renewcommand{\SetConfTitle}[1]{\relax}
  %% Don't allow individual articles to change these
  \renewcommand{\nonstopmode}{\relax}
  \renewcommand{\scrollmode}{\relax}
  \renewcommand{\outputplates}{\relax}
  %% XXXX need to check if using rmaacite and act accordingly
  %% avoid multiple declarations of the same length
  \let\oldnewlength=\newlength
  \renewcommand{\newlength}[1]{%
    \edef\reserved@a{\expandafter\@gobble\string ##1}%
    \@ifundefined{\reserved@a}{%
      \oldnewlength{##1}}{%
      \if@debug{\ClassInfo{rmaa}{\reserved@a\space already defined: doing nothing}}\fi%
      \relax}}%
  \renewcommand\newcommand{\providecommand}
  }%
{\endgroup\onecolumn}
     

%    \end{macrocode}
% Command to be used in `articles' environment for including an
% article
%    \begin{macrocode}
\newcommand{\article}[2][]{%
  % optional first argument is extra text to put in brackets after title in TOC
  % e.g. ``Invited Talk'' or ``Poster''
  \renewcommand{\articlename}{#2}
  % reset loads of counters
  \setcounter{footnote}{0}
  \setcounter{figure}{0}
  \setcounter{table}{0}
  \setcounter{section}{0}
  \setcounter{equation}{0}
  \setcounter{rmaabib}{0}
  \setcounter{author}{0}
  \setcounter{AUTHOR}{1}
  \gdef\ADS@authors{}%
  \gdef\ADS@affil{}%
  \gdef\ADS@affiliations{}%
  \gdef\ADS@references{}%
  \@firstauthortrue
  \@firstaffiltrue
  \@firstaltaffiltrue
  \@firstkwdtrue
  \@rmcitefalse
  \if@suppressfulladdresses\else\@needaddressestrue\fi
  \renewcommand{\@rmaaArg}{#1}
  \ifx\@rmaaArg\@rmaaEmpty
     \renewcommand{\@TOCaftertitle}{}
  \else
     \renewcommand{\@TOCaftertitle}{ (#1)}
  \fi
%    \end{macrocode}
% In the main journal articles must start on an odd page. 
%    \begin{macrocode}
  \if@journal\clearpage\thispagestyle{empty}\cleardoublepage\fi
%% Now input the file, then run the end-of-article hook (reinstated 28
  %% Sep 2001)
  {\input{\@articledir/\articlename.tex}\@endofarticlehook}}
%    \end{macrocode}
% Command to separate articles into sessions. Write an entry in the
% TOC, makes sure we start on an odd page (added 06 Jun 2000), then do
% a title page if \cs{SessionTitlePages} was called. 
% \changes{3.12}{04 Feb 2003}{Added option to have title page for each
% session. Turn it on with \cs{SessionTitlePages}.}
%    \begin{macrocode}
\newcommand{\session}[1]{
  %% this does nothing except write an entry in the TOC
  \addcontentsline{toc}{part}{\rmaa@TOCPartStyle{#1}}
  \addtocontents{htc}{\protect\HyperTOCpart{#1}}
  %% and also makes sure we start on an odd page (added 06 Jun 2000)
  \clearpage\thispagestyle{empty}\cleardoublepage
  \stepcounter{part}
  \if@sessiontitlepages{%
    \onecolumn
    \thispagestyle{empty}
    \renewcommand{\articlename}{Section\thepart TitlePage}
    \label{RMAAFirstPage}
    \make@session@title{#1}
    \label{RMAALastPage}
    \clearpage\thispagestyle{empty}\cleardoublepage
  }\fi
  }
\newcommand{\make@session@title}[1]{%
  {%
    \setlength\parindent{0pt}
    \vspace*{\stretch{1}}
    \rule{\linewidth}{1mm}
    \begin{flushright}\Huge\bfseries\boldmath
      #1
    \end{flushright}
    \rule{\linewidth}{1mm}
    \vspace*{\stretch{2}}
  }
}
%    \end{macrocode}
% Command to turn on separate title pages for each session. 
%    \begin{macrocode}
\newcommand{\SessionTitlePages}{\@sessiontitlepagestrue}
\newcommand{\BoldSessionTitles}{%
  \renewcommand{\rmaa@TOCPartStyle}[1]{\textbf{\rmaa@TOCTitleStyle{##1}}}
  }
\newcommand{\ArticleDir}[1]{\renewcommand{\@articledir}{#1/}\graphicspath{{1/}}}
\newcommand{\@articledir}{.}
%\newcommand{\ps@rmaafront}{\ps@rmaa}% page style for front matter (same as ever, actually)
\newcommand{\ps@rmaafront}{% page style for front matter:  page # in each foot
  \renewcommand\@oddfoot{\hfil{\thepage}}%
  \renewcommand\@evenfoot{{\thepage}\hfil}%
  % shorttitle in both heads
  \renewcommand\@evenhead{\hfil\@shorttitle\hfil}%
  \renewcommand\@oddhead{\@evenhead}
  }
\newcommand{\ps@rmaafrontfirstpage}{% page style for first page of front matter
  % page # in each foot, on outside edge
  \renewcommand\@oddfoot{\hfil\thepage}%
  \renewcommand\@evenfoot{\thepage\hfil}%
  % nothing in the even heads
  \renewcommand\@evenhead{\relax}%
  \renewcommand\@oddhead{\relax}%
  }
%% slight modification to the standard \frontmatter and \mainmatter
%% commands, just to set the \pagestyle
\renewcommand\frontmatter{%
  \latexonecolumn
  \cleardoublepage
  \@frontmattertrue
  \@mainmatterfalse
  \pagenumbering{roman}
  \pagestyle{rmaafront}
  \addtocounter{page}{2} % leave space for the title page
  }
\renewcommand\mainmatter{%
  \pagestyle{empty}
  \if@rmtwocol\twocolumn\fi
  \cleardoublepage
  \@frontmatterfalse
  \@mainmattertrue
  \pagenumbering{arabic}
  \pagestyle{rmaa}
  }
\renewcommand\backmatter{%
  \clearpage\pagestyle{empty}\cleardoublepage
  \@mainmatterfalse
  %% put extra space in the TOC
  \addtocontents{toc}{\bigskip}
  }
%    \end{macrocode}
%
% \subsection{Front sections}\label{sec:front-sections}
% \begin{macro}{\notitle}
% \begin{macro}{\shortname}
% \begin{macro}{frontsection}
% The environment \texttt{frontsection} is for doing a section in the
% front matter (preface, list of participants, or whatever). It has 
% one optional argument:
% \begin{enumerate}
% \item[\#1] Short title (no embedded spaces) for
% use in making PS files. E.g., ``list\_of\_participants''.   
% \end{enumerate}
% and two mandatory arguments: 
% \begin{enumerate}
% \item[\#2] Title of section (e.g., `Preface')
% \item[\#3] Commands to be evaluated before processing the contents
% of the environment. This is mainly useful for putting in the things
% that would go in the preamble of a normal article  (e.g.,
% |\listofauthors|, |\indexauthor|, |\notitle|) 
% \end{enumerate}
%    \begin{macrocode}
\newcommand{\notitle}{\@titlevisiblefalse} \newcommand{\shortname}{}
\newenvironment{frontsection}[3][]{%
  \begingroup
  \renewcommand{\shortname}{#1}
  \ifx\shortname\@rmaaEmpty
     \renewcommand{\articlename}{#2}
  \else
     \renewcommand{\articlename}{\shortname}
  \fi
  \@begofarticlehook
  \@titlevisibletrue % default
  \thispagestyle{rmaafrontfirstpage}
  \listofauthors{} % blank out this in case it is not specified
  \title{{#2}}\shorttitle{#2}\shortauthor{#2}
  #3 % evaluate the command given as second mandatory arg
  \addcontentsline{toc}{chapter}{\protect\@fiddleauthors{\@TOCtitle}{\rmaaTOCauthors}}%
  \addtocontents{htc}{\protect\HyperTOCEntry{\@rawtitle}{\@rawauthors}{\articlename}}
  \if@titlevisible
     \vspace*{1cm}\centerline{\rmaa@FrontSecStyle{\@title}}\vspace*{1cm}
  \fi
  \setlength{\parskip}{3ex}
  }{\frontlabel{RMAALastPage}\clearpage\endgroup}
%    \end{macrocode}
% \end{macro}
% \end{macro}
% \end{macro}
%    \begin{macrocode}
%% Definition of the command \frontlabel, which rights out something like `F4' to
%% the .aux file, instead of `\textlatin  {i}'. This is so it is easy to post-process
%% the file for making PS files of individual sections
\newcommand{\frontlabel}[1]{%
  \@bsphack
  \protected@write\@auxout{}%
     {\string\newlabel{\articlename:#1}{{}{F\@arabic{\c@page}}}}
  \@esphack
  }
%% User command to say X, Y, Z,\\ Editors of the Proceedings
\newcommand{\EditorSignature}{\vspace*{\parskip}\rightline{\@editors}%
  \rightline{Editors of the proceedings}}
%% Fiddle with \thesection, etc. so that they don't include the chapter number
%% Basically go back to the `article' definitions of these
\renewcommand{\thesection}{\@arabic\c@section}
\renewcommand{\thefigure}{\@arabic\c@figure}
\renewcommand{\thetable}{\@arabic\c@table}

\newcommand{\@oldargs}{}
\newcommand{\@newargs}{}
%% add TOC entry for each paper
%\renewcommand{\title}[1]%
%{\renewcommand{\@title}{\rmaa@TitleStyle{#1}}%
%\addtocounter{chapter}{1}\addcontentsline{toc}{chapter}{#1}%
%}

%% Formatting of the table of contents
%% dimensions of various parts
\newlength{\TOCtextwidth}\setlength{\TOCtextwidth}{12.65cm}
\newlength{\TOCpagenumwidth}\setlength{\TOCpagenumwidth}{1.3cm}
\newlength{\TOCentrywidth}\setlength{\TOCentrywidth}{\TOCtextwidth}
\addtolength{\TOCentrywidth}{-\TOCpagenumwidth} 
\newlength{\TOCextramargin}\setlength{\TOCextramargin}{1.25cm}
\newcommand{\TOCbox}[1]{\hspace{\TOCextramargin}\parbox{\TOCtextwidth}{#1}\par}
\newlength{\TOCextramarginx}\setlength{\TOCextramarginx}{\TOCextramargin}\addtolength{\TOCextramarginx}{4.3cm}
\newlength{\TOCtextwidthx}\setlength{\TOCtextwidthx}{\TOCtextwidth}\addtolength{\TOCtextwidthx}{-4.3cm}
\newcommand{\TOCboxx}[1]{\hspace{\TOCextramarginx}\parbox{\TOCtextwidthx}{#1}\par}

\newcommand{\RMAAbannerimage}{rmsc}
\renewcommand\tableofcontents{%
  \renewcommand{\articlename}{CONTENTS}
  \begingroup
  \@begofarticlehook
  \if@twocolumn
     \@restonecoltrue\onecolumn
  \else
     \@restonecolfalse
  \fi
  %% 9pt for journal style
  \if@proceedings\else\small\fi
  \if@newstyle\small\fi %% same for new style
  \if@proceedings\thispagestyle{rmaafrontfirstpage}\else\thispagestyle{empty}\fi
  \shortauthor{Contents}\shorttitle{Contents}
  \vspace*{18mm}
  \if@printbanner
  %% put in the banner header (need to get better image of this)
     \TOCbox{\hspace*{-0.15cm}\includegraphics[width=0.9\TOCtextwidth]{\RMAAbannerimage}}
     \vspace*{8mm}
  %% volume and publication date (should really be in funny font)
     \TOCboxx{\@volumen\\[2mm] \@mes\ \@anyo}
  \else
  %% put in the banner header (need to get better image of this)
     \TOCbox{\hspace*{-0.15cm}\includegraphics[width=0.9\TOCtextwidth]{\RMAAbannerimage}}
     \vspace*{12.5mm}
  %% volume and publication date (should really be in funny font)
     \TOCboxx{~\\[2mm] ~}
  \fi
  \vspace*{16mm}
  %% This is the command that fiddles with the TOC entry so that the
  %% title is flushleft but the authors are flushright
  \renewcommand{\@fiddleauthors}[2]{%
    % first put the paper title
    \renewcommand{\\}{\relax} % temporarily turn off explicit line breaks just in case
    \renewcommand{\altaffilmark}[1]{\relax} % ditto with this
    \renewcommand{\textsuperscript}[1]{\relax} % ditto with this 
    ##1
    % now save the author list in a box
    \sbox\@tempboxa{##2}%
    % check that box is shorter than the total width available for the entry
    \ifdim \wd\@tempboxa <\TOCentrywidth
       % if so, then put it right here, flush against right margin
       \penalty10000\hskip1em plus10em\penalty0\hbox{}\nobreak\hfill\box\@tempboxa
       % Note that we force a space of at least 1em between title and authorlist, 
       % stretchable to 11em since we may have very short titles but still with an 
       % authorlist too long to go on the same line.  The only linebreak allowed is 
       % directly before the authorlist.
    \else
       % if it is too long, we set the original authorlist text in its own parbox
       \par\vskip1ex\parbox[b]{\TOCentrywidth}{\raggedleft ##2}
    \fi
    }
  \@starttoc{toc}%
  \if@restonecol\twocolumn\fi
  \frontlabel{RMAALastPage}
  \clearpage
  \thispagestyle{empty}
  \cleardoublepage % make sure we start on an odd page
  \endgroup

  %%% ALSO, do the HTML TOC
  %% The file that will be used to make the HTML TOC
  \@starttoc{htc}
  %% All the stuff in this file does nothing in the book style
  \addtocontents{htc}{\protect\HyperTOCmeeting{\@conftitle}{\@reunionplace}{\@reuniondate}}
  \addtocontents{htc}{\protect\HyperTOCeditors{\@editors}}
  \addtocontents{htc}{\protect\HyperTOCpreamble}
  \addtocontents{htc}{\protect\HyperTOCEntry{Table of Contents}{}{\articlename}}

  %% and the ADS file
  \ADS@openfile
  }
%% This is now much simpler
\renewcommand\l@chapter[2]{%TOC entry for articles
  \ifnum \c@tocdepth >\m@ne
    \vskip 1.0em \@plus\p@
    \TOCbox{%
      \parbox[b]{\TOCentrywidth}{#1}\parbox[b]{\TOCpagenumwidth}{\hfill#2}
    }
  \fi} 
\renewcommand\l@part[2]{%TOC entry for sessions: centered, no page number
  \ifnum \c@tocdepth >\m@ne
    \vskip 1.5em \@plus1.5em \goodbreak
    \TOCbox{%
      \parbox[b]{\TOCentrywidth}{\hfil#1\hfil}\hfill
    }
  \fi\nobreak} 

%% Formatting theindex
%% strongly encourage page breaks between index entries
\renewcommand{\@idxitem}{\par\penalty-1000\hangindent.45cm} 
\renewcommand{\indexspace}{\par\penalty-1000\vskip 10\p@ \@plus15\p@ \@minus3\p@\relax}
\renewenvironment{theindex}{%
  %% first, make sure we start on an odd page (06 Jun 2000)
  \clearpage\thispagestyle{empty}\cleardoublepage
  %% this looks a bit like an article but is two columns and without abstract, etc
  \renewcommand{\articlename}{INDEX}
  \title{Author Index}\shortauthor{Author Index}\shorttitle{Author Index}
  \listofauthors{}
  \addcontentsline{toc}{chapter}{\protect\@fiddleauthors{\@TOCtitle}{\rmaaTOCauthors}}%
  \addtocontents{htc}{\protect\HyperTOCEntry{\@rawtitle}{\@rawauthors}{\articlename}}
  \pagestyle{rmaa}
  \thispagestyle{rmaafirstpage}
  \@begofarticlehook
  %% 9pt for journal style
  \if@proceedings\else\small\fi
  \if@newstyle\small\fi
  \twocolumn[\@makeminititle]
  %% vertical spacing (see also definition of \@idxitem and \indexspace
  \parskip\z@ \@plus .6\p@\relax % allow quite a lot of leaway here to avoid breaks inside entries
  %% horizontal spacing, breaking and glue parameters (see also the stuff in the definition 
  %% of \indexauthor
  \parindent\z@ 
  \exhyphenpenalty10000 % don't break on explicit hyphens, it always looks bad
  \rightskip=0pt plus.2em % allow a tiny bit of leeway at the ends of lines
  \spaceskip.3333em plus.5em minus.15em % a lot of leeway between words
  \xspaceskip.5em plus.5em minus.4em % this seems to have no effect
  \let\item\@idxitem
  \renewcommand{\\}{\relax} % temporarily turn off explicit line breaks just in case
  \renewcommand{\altaffilmark}[1]{\relax} % ditto with this
  \renewcommand{\textsuperscript}[1]{\relax} % ditto with this
  }{\label{RMAALastPage}\clearpage} % clearpage necessary so that we know about \shortauthor



%% end definitions for book styles
\fi

%    \end{macrocode}
% In order to do away with plates, we sometimes want a figure to
% appear in color only in the online version. The following command
% can bve used in a figure caption to typeset a note to this effect. 
%    \begin{macrocode}
\newcommand\ColorOnlineOnly{%
  \textsc{Note: this figure is available in color in the electronic
    version of this article, obtainable from}
  \texttt{http://\discretionary{}{}{}www.\discretionary{}{}{}astroscu.\discretionary{}{}{}unam.\discretionary{}{}{}mx/\discretionary{}{}{}\~{}rmaa/}.}
%% Plates
%% Define 2 counters: plate is the number in the whole document, plateloop is used 
%% internally by \outputplates
\newcounter{plate}
\newcounter{plateloop}
%% Currently limited to 15 plates in document
%% when we do a label, we need to write to \newlabel entries to the .aux file, one 
%% that says RMAAPlateX so that it can be referred to in \ouputplates, and one that
%% says articlename:RMAAPlateX so that the dvips program can find it and asociate it
%% with the right article
\newcommand{\@platelabeli}{\oldlabel{RMAAPlate1}\label{RMAAPlate1}}
\newcommand{\@platelabelii}{\oldlabel{RMAAPlate2}\label{RMAAPlate2}}
\newcommand{\@platelabeliii}{\oldlabel{RMAAPlate3}\label{RMAAPlate3}}
\newcommand{\@platelabeliv}{\oldlabel{RMAAPlate4}\label{RMAAPlate4}}
\newcommand{\@platelabelv}{\oldlabel{RMAAPlate5}\label{RMAAPlate5}}
\newcommand{\@platelabelvi}{\oldlabel{RMAAPlate6}\label{RMAAPlate6}}
\newcommand{\@platelabelvii}{\oldlabel{RMAAPlate7}\label{RMAAPlate7}}
\newcommand{\@platelabelviii}{\oldlabel{RMAAPlate8}\label{RMAAPlate8}}
\newcommand{\@platelabelix}{\oldlabel{RMAAPlate9}\label{RMAAPlate9}}
\newcommand{\@platelabelx}{\oldlabel{RMAAPlate10}\label{RMAAPlate10}}
\newcommand{\@platelabelxi}{\oldlabel{RMAAPlate11}\label{RMAAPlate11}}
\newcommand{\@platelabelxii}{\oldlabel{RMAAPlate12}\label{RMAAPlate12}}
\newcommand{\@platelabelxiii}{\oldlabel{RMAAPlate13}\label{RMAAPlate13}}
\newcommand{\@platelabelxiv}{\oldlabel{RMAAPlate14}\label{RMAAPlate14}}
\newcommand{\@platelabelxv}{\oldlabel{RMAAPlate15}\label{RMAAPlate15}}
\newcommand{\@outputplatelabeli}{\oldlabel{RMAAOutputPlate1}}
\newcommand{\@outputplatelabelii}{\oldlabel{RMAAOutputPlate2}}
\newcommand{\@outputplatelabeliii}{\oldlabel{RMAAOutputPlate3}}
\newcommand{\@outputplatelabeliv}{\oldlabel{RMAAOutputPlate4}}
\newcommand{\@outputplatelabelv}{\oldlabel{RMAAOutputPlate5}}
\newcommand{\@outputplatelabelvi}{\oldlabel{RMAAOutputPlate6}}
\newcommand{\@outputplatelabelvii}{\oldlabel{RMAAOutputPlate7}}
\newcommand{\@outputplatelabelviii}{\oldlabel{RMAAOutputPlate8}}
\newcommand{\@outputplatelabelix}{\oldlabel{RMAAOutputPlate9}}
\newcommand{\@outputplatelabelx}{\oldlabel{RMAAOutputPlate10}}
\newcommand{\@outputplatelabelxi}{\oldlabel{RMAAOutputPlate11}}
\newcommand{\@outputplatelabelxii}{\oldlabel{RMAAOutputPlate12}}
\newcommand{\@outputplatelabelxiii}{\oldlabel{RMAAOutputPlate13}}
\newcommand{\@outputplatelabelxiv}{\oldlabel{RMAAOutputPlate14}}
\newcommand{\@outputplatelabelxv}{\oldlabel{RMAAOutputPlate15}}
\newcommand{\@platerefi}{\oldpageref{RMAAPlate1}}
\newcommand{\@platerefii}{\oldpageref{RMAAPlate2}}
\newcommand{\@platerefiii}{\oldpageref{RMAAPlate3}}
\newcommand{\@platerefiv}{\oldpageref{RMAAPlate4}}
\newcommand{\@platerefv}{\oldpageref{RMAAPlate5}}
\newcommand{\@platerefvi}{\oldpageref{RMAAPlate6}}
\newcommand{\@platerefvii}{\oldpageref{RMAAPlate7}}
\newcommand{\@platerefviii}{\oldpageref{RMAAPlate8}}
\newcommand{\@platerefix}{\oldpageref{RMAAPlate9}}
\newcommand{\@platerefx}{\oldpageref{RMAAPlate10}}
\newcommand{\@platerefxi}{\oldpageref{RMAAPlate11}}
\newcommand{\@platerefxii}{\oldpageref{RMAAPlate12}}
\newcommand{\@platerefxiii}{\oldpageref{RMAAPlate13}}
\newcommand{\@platerefxiv}{\oldpageref{RMAAPlate14}}
\newcommand{\@platerefxv}{\oldpageref{RMAAPlate15}}
\newcounter{figplatei}
\newcounter{figplateii}
\newcounter{figplateiii}
\newcounter{figplateiv}
\newcounter{figplatev}
\newcounter{figplatevi}
\newcounter{figplatevii}
\newcounter{figplateviii}
\newcounter{figplateix}
\newcounter{figplatex}
\newcounter{figplatexi}
\newcounter{figplatexii}
\newcounter{figplatexiii}
\newcounter{figplatexiv}
\newcounter{figplatexv}
%% Printed output should use the counter `plate'
\renewcommand{\theplate}{\@arabic\c@plate}
\renewcommand{\theplateloop}{\@arabic\c@plateloop}
%% Command to output all the saved plates in this paper
\newcommand{\outputplates}{%
  \setcounter{plateloop}{0}
  \if@book\else\@endofarticlehook\fi % not entirely sure why, but this
                                     % seems to be needed
  \clearpage % first output the last page of text
  % now set pagestyle for blank even pages and make sure we start on an odd page
  \pagestyle{empty}
  \cleardoublepage
  \onecolumn
  \parindent\z@
  % first put in an extra sheet
  % the main purpose of this is so there is always a blank even page before 
  % each plate for filling in the PS file if the article has an odd number
  % of pages (see splitbook.sh)
  \null\vfil\centerline{PLATES}\vfil\clearpage\pagestyle{rmaaplate}\cleardoublepage
  % output each plate in turn
  \@whilenum\value{plateloop}<\value{plate}\do{%
    \addtocounter{plateloop}{1}
    % restore the figure counter from \@figcounti, \@figcountii, etc
    \setcounter{figure}{\value{figplate\roman{plateloop}}}
    % spew out the plate (in its own group for safety)
    {\csname @plate\roman{plateloop}\endcsname\par}
    }
  \pagestyle{rmaa}
  % reset the page number as if nothing had happened
  \addtocounter{page}{-\value{plate}}
  \addtocounter{page}{-\value{plate}}
  }
\newcommand\figcaption{%
  \def\@captype{figure}
  \renewcommand{\@makecaption}{\@makefigcap}
  \caption}
\newcommand\tabcaption{%
  \def\@captype{table}
  \renewcommand{\@makecaption}{\@maketabcap}
  \caption}

\newcommand{\plate}[2][1]{%
  %% User command to make a plate page: 
  %% Mandatory argument is what goes on the plate - this should not 
  %% include any floating environments (e.g. figure). Instead, \figcaption 
  %% should be used to do the captions. 
  %% Optional argument is the number of figures (i.e. captions) on the plate
  %% page (default is 1). This is necessary for incrementing the figure counter. 
  %%
  %% Note that we have to use a \gdef rather than an \xdef when we 
  %% save the plate output in a command. This means that in the book styles when it 
  %% comes to spewing out the plate, we will have forgotten about any user macros 
  %% that may have been defined. This can be got round by repeating the macros inside the plate.
  %% (Presumably, one would use a \newcommand rather than \renewcommand if they were new.)
  % increment global and local counters
  \addtocounter{plate}{1}
  % write appropriate label to .aux file
  \csname @platelabel\roman{plate}\endcsname
  % define a global macro called \@platei, \@plateii, \@plateiii, \@plateiv, ...
  \expandafter\gdef\csname @plate\roman{plate}\endcsname{%
    \csname @outputplatelabel\roman{plateloop}\endcsname
    \null\vfil #2 \vfil
    \cleardoublepage
    }
  % save the current shortauthor and shorttitle
  % need to be \xdef to take current values of \@shortauthor, etc 
  \expandafter\xdef\csname @shorta\roman{plate}\endcsname{\@shortauthor}
  \expandafter\xdef\csname @shortt\roman{plate}\endcsname{\@shorttitle}
  % save the current value of the figure counter for later use in \outputplates
  \setcounter{figplate\roman{plate}}{\value{figure}}
  % increment the figure counter according to the optional argument
  \addtocounter{figure}{#1}
  }
%% page style for plate pages 
\newcommand{\ps@rmaaplate}{%
  % Plate X (see page Y) in the foot
  \renewcommand\@oddfoot{\hfil\csname @shorta\roman{plateloop}\endcsname\ 
    (see page~\csname @plateref\roman{plateloop}\endcsname)\hfil}%
  \renewcommand\@evenfoot{\relax}%
  % nothing in even heads (blank pages)
  \renewcommand\@evenhead{\relax}
  % short title and page num in the odd heads
  \renewcommand\@oddhead{\hfil\csname @shortt\roman{plateloop}\endcsname%
    \hfil\llap{PLATE~\theplateloop}}
  }





%% Things for the User Guide
\if@guide
%% stolen from doc.sty
\def\MakeShortVerb#1{%
  \expandafter\ifx\csname cc\string#1\endcsname\relax
  \@shortvrbinfo{Made }{#1}%
  \add@special{#1}%
  \expandafter
  \xdef\csname cc\string#1\endcsname{\the\catcode`#1}%
  \begingroup
  \catcode`\~\active  \lccode`\~`#1%
  \lowercase{%
    \global\expandafter\let
    \csname ac\string#1\endcsname~%
    \gdef~{\verb~}}%
  \endgroup
  \global\catcode`#1\active
  \else
  \@shortvrbinfo\@empty{#1 already}%
  \fi}
\def\DeleteShortVerb#1{%
  \expandafter\ifx\csname cc\string#1\endcsname\relax
  \else
  \@shortvrbinfo{Deleted }{#1 as}%
  \rem@special{#1}%
  \global\catcode`#1\csname cc\string#1\endcsname
  \global \expandafter\let \csname cc\string#1\endcsname \relax
  \ifnum\catcode`#1=\active
  \begingroup
  \catcode`\~\active   \lccode`\~`#1%
  \lowercase{%
    \global\expandafter\let\expandafter~%
    \csname ac\string#1\endcsname}%
  \endgroup \fi \fi}
\def\@shortvrbinfo#1#2{%
  \PackageInfo{doc}{%
    #1\expandafter\@gobble\string#2 a short reference
    for \string\verb}}
\def\add@special#1{%
  \rem@special{#1}%
  \expandafter\gdef\expandafter\dospecials\expandafter
  {\dospecials \do #1}%
  \expandafter\gdef\expandafter\@sanitize\expandafter
  {\@sanitize \@makeother #1}}
\def\rem@special#1{%
  \def\do##1{%
    \ifnum`#1=`##1 \else \noexpand\do\noexpand##1\fi}%
  \xdef\dospecials{\dospecials}%
  \begingroup
  \def\@makeother##1{%
    \ifnum`#1=`##1 \else \noexpand\@makeother\noexpand##1\fi}%
  \xdef\@sanitize{\@sanitize}%
  \endgroup}
%% define | as shortcut for \verbatim
\MakeShortVerb{\|}
%% comand to specify arguments to LaTeX commands
\newcommand{\rmArg}[1]{%
  \ensuremath{\langle}{\scriptsize \textsf{\MakeUpperCase{#1}}}\ensuremath{\rangle}}
%\newcommand{\rmArg}[1]{\ensuremath{\langle}\textsc{#1}\ensuremath{\rangle}}
\providecommand{\BibTeX}{{\rmfamily B\kern-.05em%
    \textsc{i\kern-.025em b}\kern-.08em%
    T\kern-.1667em\lower.7ex\hbox{E}\kern-.125emX}}
%% Environment for code examples
\newenvironment{Code}{%
  \begin{list}{}{\setlength{\leftmargin}{0.05\columnwidth}} 
  \item[]}{\end{list}}
\fi
%</package>
%    \end{macrocode}

% \section{User commands for astronomical articles}
% \label{sec:user-commands}
% Some random definitions. These go both into the main class file and
% also into |adshtml.sty|.
% 
% Definitions of the journal name in different incarnations.  
%    \begin{macrocode}
%<*package|adshtml-style>
\newcommand{\RMAA}{RMxAA}
\newcommand{\RMAAlong}{Revista Mexicana de Astronom\'{\i}a y Astrof\'{\i}sica}
\newcommand{\RMAASC}{RMxAC}
\newcommand{\RMAASClong}{\RMAA\ (Serie de Conferencias)}
%    \end{macrocode}
% \begin{macro}{\ibidrule}
% This puts a horizontal rule with a period after it. It is used when
% the author list of a reference is identical to the preceding one. 
%    \begin{macrocode}
\newcommand{\ibidrule}{\rule[-0.03cm]{4.7em}{.01cm}.\ }
%    \end{macrocode}
% \end{macro}
% \begin{macro}{\ThisVolume}
% References to other articles in the same volume can be coded as,
% e.g., |\bibitem{} Bloggs, J. \ThisVolume{bloggs}|.
%    \begin{macrocode}
\newcommand{\ThisVolume}[1]{%
  \@year, \if@proceedings\RMAASC\else\RMAA\fi, \@volume,
  \if@book\oldpageref{#1:RMAAFirstPage}\else 000\fi\ (this volume)}
%    \end{macrocode}
% \end{macro}
% Other stuff that may be useful in ref lists.
%    \begin{macrocode}
\newcommand{\etal}{et al.\@}
%    \end{macrocode}
% Various new commands from AAS macros
%    \begin{macrocode}
\def\acknowledgments{\vskip 3ex plus .8ex minus .4ex}
\let\acknowledgements=\acknowledgments                  % second spelling
\let\jnl@style=\rm
\def\ref@jnl#1{{\jnl@style#1}}
\DeclareRobustCommand\aj{\ref@jnl{AJ}}                   
\DeclareRobustCommand\araa{\ref@jnl{ARA\&A}}             
\DeclareRobustCommand\apj{\ref@jnl{ApJ}}                 
\DeclareRobustCommand\apjl{\ref@jnl{ApJ}}                
\DeclareRobustCommand\apjs{\ref@jnl{ApJS}}               
\DeclareRobustCommand\ao{\ref@jnl{Appl.~Opt.}}           
\DeclareRobustCommand\apss{\ref@jnl{Ap\&SS}}             
\DeclareRobustCommand\aap{\ref@jnl{A\&A}}                
\DeclareRobustCommand\aapr{\ref@jnl{A\&A~Rev.}}          
\DeclareRobustCommand\aaps{\ref@jnl{A\&AS}}              
\DeclareRobustCommand\azh{\ref@jnl{AZh}}                 
\DeclareRobustCommand\baas{\ref@jnl{BAAS}}               
\DeclareRobustCommand\jrasc{\ref@jnl{JRASC}}             
\DeclareRobustCommand\memras{\ref@jnl{MmRAS}}            
\DeclareRobustCommand\mnras{\ref@jnl{MNRAS}}             
\DeclareRobustCommand\pra{\ref@jnl{Phys.~Rev.~A}}        
\DeclareRobustCommand\prb{\ref@jnl{Phys.~Rev.~B}}        
\DeclareRobustCommand\prc{\ref@jnl{Phys.~Rev.~C}}        
\DeclareRobustCommand\prd{\ref@jnl{Phys.~Rev.~D}}        
\DeclareRobustCommand\pre{\ref@jnl{Phys.~Rev.~E}}        
\DeclareRobustCommand\prl{\ref@jnl{Phys.~Rev.~Lett.}}    
\DeclareRobustCommand\pasp{\ref@jnl{PASP}}               
\DeclareRobustCommand\pasj{\ref@jnl{PASJ}}               
\DeclareRobustCommand\qjras{\ref@jnl{QJRAS}}             
\DeclareRobustCommand\skytel{\ref@jnl{S\&T}}             
\DeclareRobustCommand\solphys{\ref@jnl{Sol.~Phys.}}      
\DeclareRobustCommand\sovast{\ref@jnl{Soviet~Ast.}}      
\DeclareRobustCommand\ssr{\ref@jnl{Space~Sci.~Rev.}}     
\DeclareRobustCommand\zap{\ref@jnl{ZAp}}                 
\DeclareRobustCommand\nat{\ref@jnl{Nature}}              
\DeclareRobustCommand\iaucirc{\ref@jnl{IAU~Circ.}}
\DeclareRobustCommand\aplett{\ref@jnl{Astrophys.~Lett.}}
\DeclareRobustCommand\apspr{\ref@jnl{Astrophys.~Space~Phys.~Res.}}
\DeclareRobustCommand\bain{\ref@jnl{Bull.~Astron.~Inst.~Netherlands}}
\DeclareRobustCommand\fcp{\ref@jnl{Fund.~Cosmic~Phys.}}
\DeclareRobustCommand\gca{\ref@jnl{Geochim.~Cosmochim.~Acta}}
\DeclareRobustCommand\grl{\ref@jnl{Geophys.~Res.~Lett.}}
\DeclareRobustCommand\jcp{\ref@jnl{J.~Chem.~Phys.}}      
\DeclareRobustCommand\jgr{\ref@jnl{J.~Geophys.~Res.}}    
\DeclareRobustCommand\jqsrt{\ref@jnl{J.~Quant.~Spec.~Radiat.~Transf.}}
\DeclareRobustCommand\memsai{\ref@jnl{Mem.~Soc.~Astron.~Italiana}}
\DeclareRobustCommand\nphysa{\ref@jnl{Nucl.~Phys.~A}}
\DeclareRobustCommand\physrep{\ref@jnl{Phys.~Rep.}}
\DeclareRobustCommand\physscr{\ref@jnl{Phys.~Scr}}
\DeclareRobustCommand\planss{\ref@jnl{Planet.~Space~Sci.}}
\DeclareRobustCommand\procspie{\ref@jnl{Proc.~SPIE}}
\let\astap=\aap
\let\apjlett=\apjl
\let\apjsupp=\apjs
\let\applopt=\ao
\def\phn{\phantom{0}}
\def\phd{\phantom{.}}
\def\phs{\phantom{\ensuremath{-}}}
\def\phm#1{\phantom{#1}}
%    \end{macrocode}
% All these are redone as robust and using more modern commands where
% possible. The robustness is so they will not get expanded if used in
% abstracts and titles when these are passed to the ADS file. 
% \changes{3.16}{2003/04/07}{
%  Added \cs{Color} command for doing photometric colors.}
%    \begin{macrocode}
%%\def\sun{\hbox{$\odot$}}
\DeclareRobustCommand\sun{\ensuremath{\odot}}
%%\def\earth{\hbox{$\oplus$}}
\DeclareRobustCommand\earth{\ensuremath{\oplus}}
%%\def\lesssim{\mathrel{\hbox{\rlap{\hbox{\lower4pt\hbox{$\sim$}}}\hbox{$<$}}}}
%%\def\gtrsim{\mathrel{\hbox{\rlap{\hbox{\lower4pt\hbox{$\sim$}}}\hbox{$>$}}}}
\DeclareRobustCommand\lesssim{\mathrel{\hbox{\smaller\smaller\rlap{\hbox{\lower3pt\hbox{$\sim$}}}\raise1.5pt\hbox{$<$}}}}
\DeclareRobustCommand\gtrsim{\mathrel{\hbox{\smaller\smaller\rlap{\hbox{\lower3pt\hbox{$\sim$}}}\raise1.5pt\hbox{$>$}}}}
%%\def\lta{\mathrel{\spose{\lower 3pt\hbox{$\mathchar"218$}}
%%     \raise 2.0pt\hbox{$\mathchar"13C$}}}
%%\def\gta{\mathrel{\spose{\lower 3pt\hbox{$\mathchar"218$}}
%%     \raise 2.0pt\hbox{$\mathchar"13E$}}}
\DeclareRobustCommand\sq{\hbox{\rlap{$\sqcap$}$\sqcup$}}
%%\def\arcdeg{\hbox{$^\circ$}}
%%\def\arcmin{\hbox{$^\prime$}}
%%\def\arcsec{\hbox{$^{\prime\prime}$}}
\DeclareRobustCommand\arcdeg{\ensuremath{^\circ}}
\DeclareRobustCommand\arcmin{\ensuremath{^\prime}}
\DeclareRobustCommand\arcsec{\ensuremath{^{\prime\prime}}}
%%\def\fd{\hbox{$.\!\!^{\rm d}$}}
%%\def\fh{\hbox{$.\!\!^{\rm h}$}}
%%\def\fm{\hbox{$.\!\!^{\rm m}$}}
%%\def\fs{\hbox{$.\!\!^{\rm s}$}}
%%\def\fdg{\hbox{$.\!\!^\circ$}}
%%\def\farcm{\hbox{$.\mkern-4mu^\prime$}}
%%\def\farcs{\hbox{$.\!\!^{\prime\prime}$}}
%%\def\fp{\hbox{$.\!\!^{\scriptscriptstyle\rm p}$}}
\DeclareRobustCommand\fd{\ensuremath{.\!\!^\mathrm{d}}}
\DeclareRobustCommand\fh{\ensuremath{.\!\!^\mathrm{h}}}
\DeclareRobustCommand\fm{\ensuremath{.\!\!^\mathrm{m}}}
\DeclareRobustCommand\fs{\ensuremath{\mkern0.5mu.\mkern-5mu^\mathrm{s}\mkern-1mu}}
\DeclareRobustCommand\fdg{\ensuremath{.\!\!^\circ}}
\DeclareRobustCommand\farcm{\ensuremath{.\mkern-4mu^\prime}}
\DeclareRobustCommand\farcs{\ensuremath{\mkern1mu.\mkern-6mu^{\prime\prime}\mkern-1.5mu}}
%%\DeclareRobustCommand\farcs{\ensuremath{.\!\!^{\prime\prime}}}
\DeclareRobustCommand\fp{\ensuremath{.\!\!^{\scriptscriptstyle\rm p}}}
%%\def\micron{\hbox{$\mu$m}}
\DeclareRobustCommand\micron{\ensuremath{\mu\mathrm{m}}}
\let\la=\lesssim                        
\let\ga=\gtrsim
%%\def\case#1#2{\hbox{$\frac{#1}{#2}$}}
%%\def\slantfrac#1#2{\hbox{$\,^#1\!/_#2$}}
%%\def\onehalf{\slantfrac{1}{2}}
%%\def\onethird{\slantfrac{1}{3}}
%%\def\twothirds{\slantfrac{2}{3}}
%%\def\onequarter{\slantfrac{1}{4}}
%%\def\threequarters{\slantfrac{3}{4}}
\DeclareRobustCommand\case[2]{\ensuremath{\frac{#1}{#2}}}
\DeclareRobustCommand\slantfrac[2]{\ensuremath{\,^#1\!/_#2}}
\DeclareRobustCommand\onehalf{\slantfrac{1}{2}}
\DeclareRobustCommand\onethird{\slantfrac{1}{3}}
\DeclareRobustCommand\twothirds{\slantfrac{2}{3}}
\DeclareRobustCommand\onequarter{\slantfrac{1}{4}}
\DeclareRobustCommand\threequarters{\slantfrac{3}{4}}
%%\def\ubvr{\hbox{$U\!BV\!R$}}            
%%\def\ub{\hbox{$U\!-\!B$}}               
%%\def\bv{\hbox{$B\!-\!V$}}               
%%\def\vr{\hbox{$V\!-\!R$}}               
%%\def\ur{\hbox{$U\!-\!R$}}
%    \end{macrocode}
% New commands for ``photometric'' colors 
%    \begin{macrocode}
\DeclareRobustCommand\ubvr{\ensuremath{U\!BV\!R}}            
\DeclareRobustCommand\Color[2]{%
  \ensuremath{%
    #1\!\ifx#1V\!\fi-\!#2
  }}
\DeclareRobustCommand\ub{\Color{U}{B}}
\DeclareRobustCommand\bv{\Color{B}{V}}
\DeclareRobustCommand\vr{\Color{V}{R}}
\DeclareRobustCommand\ur{\Color{U}{R}}
%% New version of the \ion command, that actually works! 
%    \end{macrocode}
% \begin{macro}{\ion}
% This version of the |\ion| command allows you to write either
% |\ion{N}{4}| or |\ion{N}{iv}| in order to get N~IV. 
%    \begin{macrocode}
\DeclareRobustCommand\ion[2]{%
  \mbox{#1\kern0.2em%
    \smaller\rmfamily%
    \edef\@tempa{\@car#2\@nil}%
    \ifcat1\@tempa%
    \@Roman{#2}%
    \else%
    \uppercase{#2}%
    \fi}}
\newcounter{IonStage}
\DeclareRobustCommand\plainion[2]{\setcounter{IonStage}{#2}#1 \Roman{IonStage}}
%%\def\ion#1#2{\ensuremath{\mathrm{#1}\;}{\protect\small\rm\@Roman{#2}}}
%    \end{macrocode}
% \end{macro}
%    \begin{macrocode}
\DeclareRobustCommand\nodata{$\cdots$}
%    \end{macrocode}
% These are needed by the |\ion| command. Isn't there an easier way of
% doing this?
%    \begin{macrocode}
\let\@small=\small
\let\@normalsize=\normalsize
\let\@footnotesize=\footnotesize
\let\@scriptsize=\scriptsize
\let\@tiny=\tiny
\let\@large=\large
\let\@Large=\Large
\let\@LARGE=\LARGE
\let\@huge=\huge
\let\@Huge=\Huge
\newcommand{\smaller}{\small}
\renewcommand{\normalsize}{\renewcommand{\smaller}{\small}\@normalsize}
\renewcommand{\small}{\renewcommand{\smaller}{\footnotesize}\@small}
\renewcommand{\footnotesize}{\renewcommand{\smaller}{\scriptsize}\@footnotesize}
\renewcommand{\scriptsize}{\renewcommand{\smaller}{\tiny}\@scriptsize}
\renewcommand{\tiny}{\renewcommand{\smaller}{\tiny}\@tiny}
\renewcommand{\large}{\renewcommand{\smaller}{\normalsize}\@large}  
\renewcommand{\Large}{\renewcommand{\smaller}{\large}\@Large}  
\renewcommand{\LARGE}{\renewcommand{\smaller}{\Large}\@LARGE}  
\renewcommand{\huge}{\renewcommand{\smaller}{\LARGE}\@huge}  
\renewcommand{\Huge}{\renewcommand{\smaller}{\huge}\@Huge}  
%</package|adshtml-style>
%    \end{macrocode}


% \section{Other files}

% \subsection{The driver and make files}
% \label{sec:driver_make}
% These are put here for completeness, although they still need to be
% distributed separately since otherwise the user won't be able to
% unpack this file into its constituent parts.
% 
% The |Makefile| has two principal rules: |docs|, and |macros|. The
% |docs| rule produces a PDF and gzipped PS file of the documentation.
% I still haven't perfected the combination of rules required to get
% it to run \LaTeX{} enough times in all circumstances. The |macros|
% rule runs \LaTeX{} on |rmaa.ins| in order to unpack |rmaa.cls| and
% all the other little files (including the |Makefile| and the
% |rmaa.ins| files themselves\dots 
% 
% First set up lists of various types of file. 
%    \begin{macrocode}
%<*makefile>
DOCFILES=rmaa.pdf
MACRO_LTX_FILES=rmaa.cls rmaacite.sty adshtml.sty
SHELLSCRIPTS=makeads.sh
MACRO_OTHER_FILES=adshtml.perl sample.latex2html-init
MAKEFILE=Makefile
MACROFILES=$(MACRO_LTX_FILES) $(SHELLSCRIPTS) $(MACRO_OTHER_FILES)         
SKELETONFILES=ADS.tex 
%    \end{macrocode}
% Install the files read only into the package directory, specified via a path relative to the source directory. \changes{4.0}{2019/11/14}{Reorganize folder layout so that hardly anything is in the top level.  That way, things are clearer on github.}
%    \begin{macrocode}
INSTALL=install 
INSTALL_DIR=../packages
%    \end{macrocode}
% The default rule is to make everything. Installation must be
% explicitly asked for.
%    \begin{macrocode}
all: docs macros

install: install-docs install-macros

%    \end{macrocode}
% Most files are installed read-only to prevent editing. Skeleton
% files must be writeable. Must be careful not to overwrite the
% Makefile in the parent directory. 
%    \begin{macrocode}
install-macros: macros
        $(INSTALL) -m 444 $(MACRO_LTX_FILES) $(MACRO_OTHER_FILES) $(INSTALL_DIR)
        $(INSTALL) -m 555 $(SHELLSCRIPTS) $(INSTALL_DIR)
        $(INSTALL) -m 644 $(SKELETONFILES) $(INSTALL_DIR)

install-docs: docs
        $(INSTALL) rmaa.pdf $(INSTALL_DIR)/rm-fulldocs.pdf

%    \end{macrocode}
% Make the PDF and gzipped PS files. 
%    \begin{macrocode}
docs: $(DOCFILES)

%    \end{macrocode}
% Make everything else if either the |.dtx| or |.ins| files have
% changed. 
%    \begin{macrocode}
macros: rmaa.dtx rmaa.ins
%    \end{macrocode}
% Ensure that all the macro files can be written to then recreate
% them. The |touch| is there in case any of the files don't exist
% yet. 
%    \begin{macrocode}
        touch $(MACROFILES) $(MAKEFILE)
        chmod u+w $(MACROFILES) $(MAKEFILE) rmaa.ins
        latex "\input{rmaa.ins}"
%    \end{macrocode}
% \TeX{} automatically adds the |.tex| extension to extensionless
% files. We fix this and take the opportunity to convert
% back to TABs all that it has lovingly converted to spaces. 
%    \begin{macrocode}
        cp -f $(MAKEFILE) $(MAKEFILE).BCK && unexpand $(MAKEFILE).tex >| $(MAKEFILE) && rm $(MAKEFILE).tex
%    \end{macrocode}
% Make the shell scripts executable and everything unwriteable (so
% they don't get hand-edited by mistake). However, undo this for the
% files that the user is \emph{supposed} to modify. 
%    \begin{macrocode}
        chmod a+x $(SHELLSCRIPTS)
        chmod a-w $(MACROFILES) $(MAKEFILE)
        chmod u+w $(SKELETONFILES)

%    \end{macrocode}
% We now use PDF\LaTeX{} for this.
%    \begin{macrocode}
rmaa.idx: rmaa.dtx
        @echo "REGENERATING .idx FILE"
        pdflatex "\nonstopmode\input{rmaa.dtx}"

rmaa.aux: rmaa.dtx rmaa.ind rmaa.gls
        @echo "REGENERATING .aux FILE"
        pdflatex "\nonstopmode\input{rmaa.dtx}"

rmaa.ind: rmaa.idx
        @echo "REGENERATING INDEX"
        makeindex -s gind.ist rmaa

rmaa.gls: rmaa.glo
        @echo "REGENERATING CHANGELOG"
        makeindex -s gglo.ist -o rmaa.gls rmaa.glo

rmaa.pdf: rmaa.aux rmaa.ind rmaa.gls
        @echo "REGENERATING .pdf FILE"
        pdflatex "\nonstopmode\input{rmaa.dtx}"

clean: 
        rm -f rmaa.aux rmaa.dvi rmaa.ind rmaa.gls
%</makefile>
%    \end{macrocode}

% The |rmaa.ins| file uses the \textsf{docstrip} package to split out
% all the files and strip off this documentation.
%    \begin{macrocode}
%<*ins>
%%
%% Installation batch file for the RMAA LaTeX package
%%
\def\batchfile{rmaa.ins}
\input docstrip.tex
\preamble

PLEASE DO NOT EDIT THIS FILE! It is generated automatically from
the file rmaa.dtx using the accompanying Makefile and rmaa.ins 
Any hand-edited changes will be lost

A LaTeX document class for the RMxAA

Author: William Henney (w.henney@irya.unam.mx)

Latest version available from: http://www.irya.unam.mx/rmaa

Copyright (C) 1999-2019 Instituto de Astronomia, UNAM, Mexico

Permission is hereby granted, free of charge, to any person obtaining
a copy of this software and associated documentation files (the
"Software"), to deal in the Software without restriction, including
without limitation the rights to use, copy, modify, merge, publish,
distribute, sublicense, and/or sell copies of the Software, and to
permit persons to whom the Software is furnished to do so, subject to
the following conditions:

The above copyright notice and this permission notice shall be
included in all copies or substantial portions of the Software.

THE SOFTWARE IS PROVIDED "AS IS", WITHOUT WARRANTY OF ANY KIND,
EXPRESS OR IMPLIED, INCLUDING BUT NOT LIMITED TO THE WARRANTIES OF
MERCHANTABILITY, FITNESS FOR A PARTICULAR PURPOSE AND
NONINFRINGEMENT. IN NO EVENT SHALL THE AUTHORS OR COPYRIGHT HOLDERS BE
LIABLE FOR ANY CLAIM, DAMAGES OR OTHER LIABILITY, WHETHER IN AN ACTION
OF CONTRACT, TORT OR OTHERWISE, ARISING FROM, OUT OF OR IN CONNECTION
WITH THE SOFTWARE OR THE USE OR OTHER DEALINGS IN THE SOFTWARE.
\endpreamble
%    \end{macrocode}
% Clobber files without asking
%    \begin{macrocode}
\askforoverwritefalse
%    \end{macrocode}
% Make the \LaTeX{} macro files (class and styles)
%    \begin{macrocode}
\generate{% 
  \file{rmaa.cls}{\from{rmaa.dtx}{package}}
  \file{rmaa.ins}{\from{rmaa.dtx}{ins}}
  \file{rmaacite.sty}{\from{rmaa.dtx}{rmaacite}}
  \file{adshtml.sty}{\from{rmaa.dtx}{adshtml-style}}
  \file{ADS.tex}{\from{rmaa.dtx}{ads-container}}
  }
%    \end{macrocode}
% Make the other non-\LaTeX{} files. These are done separately since
% they can't use `\%\%' as a comment character for the
% preamble/postamble (and Perl files can't have the postamble I don't
% think). 
%    \begin{macrocode}
\generate{\nopreamble\nopostamble% 
  \file{Makefile}{\from{rmaa.dtx}{makefile}}
  \file{adshtml.perl}{\from{rmaa.dtx}{adshtml-perl}}
  \file{makeads.sh}{\from{rmaa.dtx}{ads-shell-script}}
  \file{sample.latex2html-init}{\from{rmaa.dtx}{latex2html-init}}
  }

%% Local Variables: 
%% mode: LaTeX
%% mode: Font-lock
%% End:
%</ins>
%    \end{macrocode}

% \subsection{Extended citation syntax}
% Modified from a file called mnrascite.sty
% I no longer remember where I stole the original version of this from
% 
% Alterations to the cite commands to make them do more sensible things.
% N.B. the bibliography style should produce entries which have labels where
% the date is delimited by angle brackets '<>' for these macros to work
% properly. They will not give very sensible error messages if you do not do 
% this properly.
% 
% A new command |\scite{key}| is defined which produces a citation which is 
% suitable for insertion in a sentence. i.e., it does not put brackets around 
% the whole citation, only the date. Also |\pcite{key}| produce output with no
% parentheses whatsoever. 
% 
% Should redefine |\@bibitem| really to have any hope of compatibility
% the command |\lbibcite| is defined to deal with the citations.
%    \begin{macrocode}
%<*rmaacite>
\def\bibitem[#1<#2>]#3{\par\hangindent=1.5em\hangafter=1
  \if@filesw
  {\def\protect##1{\string ##1\space}
    \immediate\write\@auxout{\string\lbibcite{#3}{#1}{#2}}}
  \fi}
\@redefinebibitem % what it says
%    \end{macrocode}
% Note that the |\lbibcite| defines two types of citation one contains the names,
% the other the year
%    \begin{macrocode}
\def\lbibcite#1#2#3{\global\@namedef{b@#1}{#2}
                    \global\@namedef{by@#1}{#3}}

\def\@bcite#1#2{({#1\if@tempswa , #2\fi})}
\def\@pcite#1#2{#1\if@tempswa , #2\fi}
%    \end{macrocode}
% This declaration contains what is put between the name and the date
% \changes{3.11}{2003/01/31}{Moved the \cs{articlename} to after the key for consistency
% with the new natbib stuff.}
%    \begin{macrocode}
\def\@brcite{\ (}
\def\@ebrcite{)}
\def\@fmttop#1{\@fmtcite[\csname #1\endcsname]}
\def\@citex[#1]#2{%
%%  \if@filesw\immediate\write\@auxout{\string\citation{#2}}\fi
%% WJH added 16 Jan 2002
  \@for\@citeb:=#2\do{%
    \edef\@citeb{\expandafter\@firstofone\@citeb:\articlename}%
    \if@filesw\immediate\write\@auxout{\string\citation{\@citeb}}\fi}%
%% to here
  \def\@citea{}%
  \@cite{\@for\@citeb:=#2\do
    {%
      \edef\@citeb{\expandafter\@firstofone\@citeb:\articlename}%
      \@citea\def\@citea{;\ }\@ifundefined
      {b@\@citeb}{{\bf ?}\@warning {Citation `\@citeb' on page \thepage \space undefined}}%
      {\csname b@\@citeb\endcsname\@citebro\csname by@\@citeb\endcsname\@citebrc}}}%
  {#1}}

\def\cite{\@ifnextchar [{\let\@citebro=\space
                          \let\@citebrc=\relax
                          \let\@cite=\@bcite\@tempswatrue\@citex}
                        {\let\@citebro=\space
                          \let\@citebrc=\relax
                          \let\@cite=\@bcite\@tempswafalse\@citex[]}}
\def\pcite{\@ifnextchar [{\let\@citebro=\space
                          \let\@citebrc=\relax
                          \let\@cite=\@pcite\@tempswatrue\@citex}
                        {\let\@citebro=\space
                          \let\@citebrc=\relax
                          \let\@cite=\@pcite\@tempswafalse\@citex[]}}
\def\scite{\@ifnextchar [{\let\@citebro=\@brcite
                          \let\@citebrc=\@ebrcite
                          \let\@cite=\@pcite\@tempswatrue\@citex}
                        {\let\@citebro=\@brcite
                          \let\@citebrc=\@ebrcite
                          \let\@cite=\@pcite\@tempswafalse\@citex[]}}

\renewenvironment{thebibliography}
  {\section*{REFERENCES}\vspace*{\afterbiblioskip}\small
    \if@book\RereadBiblio\fi
    \sloppy\clubpenalty4000\widowpenalty4000
    \sfcode`\.=1000\relax
    \setlength\parindent{\z@}\setlength\parskip{\z@}}
  {\relax\normalsize\@outputaddressesifnecessary} 
%</rmaacite>
%    \end{macrocode}

% \subsection{Auxilliary files for generating ADS data}

% The ADS data are written by \LaTeX{} to a |.ads| file. This file
% then has to be included in the following file (|ADS.tex|),
% which is run through \LaTeX{}2HTML (see below). 

%    \begin{macrocode}
%<*ads-container>
%%
%% Instructions:
%% Edit the \input{} line below to reflect the name of your file. 
\documentclass{article}
\usepackage[latin1]{inputenc}
\usepackage{adshtml}
\newcommand{\ADSfield}[2]
{\%#1  #2\\}
%%\newcommand{\ADSvfield}[2]
%%{\%#1 \begin{verbatim}#2\end{verbatim}}
\newenvironment{ADSarticle}{}{\par\bigskip}
\renewcommand{\bibitem}[2][]{\par}
\setlength{\parindent}{0pt}
\renewcommand{\ThisVolume}[1]{%  
  2001, \RMAASC, 12, \oldpageref{#1:RMAAFirstPage}}  
 

\begin{document}
\makeatletter
\input{rmbooksamp.ads}
\end{document}
%</ads-container>
%    \end{macrocode}

% This is the |adshtml.sty| style file used by |ADS.tex|. It just
% defines a few 
% macros that might be used in abstracts, etc (really it should define
% them all). This now gets all its code from
% Sec~\ref{sec:user-commands}. 
%    \begin{macrocode}
%<*adshtml-style>
%% Nothing else needs to go here now. 
%</adshtml-style>
%    \end{macrocode}
% The init file for latex2html.
%    \begin{macrocode}
%<*latex2html-init>
$PDFTOCAIRO = "pdftocairo";
$PDFLATEX = "pdflatex";
$LATEX = "latex";
$MAX_SPLIT_DEPTH = 0;
$HTML_VERSION = "4.0,math";
$NO_SIMPLE_MATH = 1;
$ADDRESS = "<p><address><a href=\"mailto:rmaa\@astro.unam.mx\">rmaa\@astro.unam.mx</address></p>"; 
$INFO = "<p>This document was created using latex2html and the rmaa latex macros</p>"; 
$NO_NAVIGATION = 1;
1; 
%</latex2html-init>
%    \end{macrocode}
% This is the |adshtml.perl| file that is used by \LaTeX{}2HTML. Some
% bizarre \TeX{}y command sometimes show up so we tell it to ignore
% these. All the astronomy macros that come from AAS are processed in
% \TeX{}. That way they get made to images but with an |ALT| tag that
% has the original code, which is what will be picked up by |lynx|,
% with the result that they get passed intact into the final |.txt|
% file. 
%    \begin{macrocode}
%<*adshtml-perl>
package main;

$dum = &do_cmd_makeatletter();

sub do_cmd_RMAA {
    local($_) = @_;
    "RMxAA" . $_;
}

sub do_cmd_RMAASC {
    local($_) = @_;
    "RMxAASC" . $_;
}

sub do_cmd_etal {
    local($_) = @_;
    "et al\." . $_;
}

sub do_cmd_ibidrule {
    local($_) = @_;
    "<img src=\"foo\" alt=\"\ibidrule\">" . $_;
}


sub do_cmd_bibitem {
    local($_) = @_;
    local($key);
    local($opt, $dummy)=&get_next_optional_argument; 
    $key = &missing_braces unless
        ((s/$next_pair_pr_rx/$key = $2; ''/eo)
        ||(s/$next_pair_rx/$key = $2; ''/eo));
    "<p>" . $_; 
}

&ignore_commands (<<_RAW_ARG_CMDS_);
lowercase
adjustfinalcols
textsuperscript # {}
newline # []
thinmuskip
mskip
smallskip
medskip
bigskip
hspace # {}
hspacestar # {}
vspace # {}
vspacestar # {}
mathrm
_at_mathrm
spacefactor
discretionary
unhbox 
voidb_at_x 
penalty 
_at_M
_at_m
null
IeC
_RAW_ARG_CMDS_

&process_commands_in_tex (<<_RAW_ARG_CMDS_);
sun
earth
lesssim
gtrsim
sq
argdeg
arcmin
arcsec
fd
fh
fm
fs
fdg
farcm
farcs
fp
micron
la
ga
case # {} # {}
slantfrac # {} # {}
onehalf
onethird
twothirds
onequarter
threequarters
ubvr
ub
bv
vr
ur
aj
araa
apj
apjl
apjs
ao
apss
aap
aapr
aaps
azh
baas
jrasc
memras
mnras
pra
prb
prc
prd
pre
prl
pasp
pasj
qjras
skytel
solphys
sovast
ssr
zap
nat
iaucirc
aplett
apspr
bain
fcp
gca
grl
jcp
jgr
jqsrt
memsai
nphysa
physrep
physscr
planss
procspie
_RAW_ARG_CMDS_

1;                              # This must be the last line
%</adshtml-perl> 
%    \end{macrocode}

% This is a shell script |makeads.sh| that generates the |_ads.txt| file from
% |_ads.tex| by using \LaTeX{}2HTML and |lynx|.  \changes{4.0}{2019/11/13}{Make sure that lynx writes utf-8 output instead of latin-1.}
% \changes{2.2beta13}{2002/03/12}{Complete rewrite of |makeads.sh|. Now much 
% improved I think}
% \changes{4.0}{2019/11/13}{Fix the at-sign substitution so that it only eliminates \cs{@} and does not mess with email addresses.}

%    \begin{macrocode}
%<*ads-shell-script>
#! /bin/sh
if [ -z $1 ]; then
    echo "Usage: $0 basename"
    exit
fi
BASENAME=$1
ADSBASENAME=${BASENAME}_ads
AUXFILE=$BASENAME.aux
TMPFILE=tmp.$$


echo "Generating $ADSBASENAME.tex"

# For some reason, must double-up \\ to \\\\ but not \ to \\ - odd!
cat > $ADSBASENAME.tex <<EOF
%% PLEASE DO NOT EDIT THIS FILE!  It is generated automatically from
%% the script makeads.sh.  Any hand-edited changes will be lost
%% 
\documentclass{article}
\usepackage[latin1]{inputenc}
\usepackage{html}
\usepackage{adshtml}
\newcommand{\ADSfield}[2]{\%#1  #2\\\\}
\newenvironment{ADSarticle}{}{\par\bigskip}
\setlength{\parindent}{0pt}
\newcounter{istage}
\renewcommand{\ion}[2]{\setcounter{istage}{#2}#1 \Roman{istage}}
\begin{document}
\makeatletter
EOF
echo >> $ADSBASENAME.tex "\input{$BASENAME.ads2}" 
echo >> $ADSBASENAME.tex "\end{document}"



echo "Generating $ADSBASENAME.sed" 

# Find type of journal
grep '\\documentclass\[.*proceedings.*\]' $BASENAME.tex >/dev/null \
     && JOURNAL='RMxAC' \
     || JOURNAL='RMxAA'

# Find the year and volume number
YEAR=`grep '\\SetYear' $BASENAME.tex | sed 's/.*\\SetYear{\(.*\)}.*/\1/'`
VOLUME=`grep '\\SetVolume' $BASENAME.tex | sed 's/.*\\SetVolume{\(.*\)}.*/\1/'`

# Generate an auxilliary sed script 
cat > $TMPFILE <<EOF
s|\\\\newlabel{\(.*\):RMAAFirstPage}{{.*}{\(.*\)}}|s/\\\\\\\\ThisVolume{\1}/$YEAR, RMxAC, $VOLUME, \2/g|g
EOF

# Now go through the .aux file, looking for \newlabels and use the sed 
# script just generated in order to make another sed script, one that 
# will replace \ThisVolume{LABEL} with the correct reference
grep '^\\newlabel{.*:RMAAFirstPage}' $AUXFILE \
     | sed -f $TMPFILE > $ADSBASENAME.sed && rm $TMPFILE

# Apend a further sed command to get rid of \@, since redefining 
# it in _ads.tex didn't work
echo >> $ADSBASENAME.sed 's/\\\\@//g'

echo "Generating $BASENAME.ads2 from $ADSBASENAME.sed and $BASENAME.ads" 

# Apply this to the .ads file, generating the .ads2 file, which is included
# by the _ads.tex file.
sed -f $ADSBASENAME.sed $BASENAME.ads > $BASENAME.ads2

echo "Generating $ADSBASENAME.txt from $ADSBASENAME.tex and $BASENAME.ads2"

# Now run latex2html on the _ads.tex file
latex2html -no_navigation -info 0 -address 0 -html_version 4.0,math $ADSBASENAME

# Then, convert to utf8 with lynx and tidy it up a bit
lynx --justify=off --width=100 --dump --nolist --display_charset=utf-8 \
     $ADSBASENAME/$ADSBASENAME.html | cut -c4- \
     | sed 's/ibidrule/\\ibidrule/g' > $ADSBASENAME.txt
%</ads-shell-script>
%    \end{macrocode}

% \section{Epilogue}
% Err, that's it \dots

\endinput


% \endinput
% Local Variables:
% mode: doctex
% TeX-master: t
% End:
"

clean: 
        rm -f rmaa.aux rmaa.dvi rmaa.ind rmaa.gls
%</makefile>
%    \end{macrocode}

% The |rmaa.ins| file uses the \textsf{docstrip} package to split out
% all the files and strip off this documentation.
%    \begin{macrocode}
%<*ins>
%%
%% Installation batch file for the RMAA LaTeX package
%%
\def\batchfile{rmaa.ins}
\input docstrip.tex
\preamble

PLEASE DO NOT EDIT THIS FILE! It is generated automatically from
the file rmaa.dtx using the accompanying Makefile and rmaa.ins 
Any hand-edited changes will be lost

A LaTeX document class for the RMxAA

Author: William Henney (w.henney@irya.unam.mx)

Latest version available from: http://www.irya.unam.mx/rmaa

Copyright (C) 1999-2019 Instituto de Astronomia, UNAM, Mexico

Permission is hereby granted, free of charge, to any person obtaining
a copy of this software and associated documentation files (the
"Software"), to deal in the Software without restriction, including
without limitation the rights to use, copy, modify, merge, publish,
distribute, sublicense, and/or sell copies of the Software, and to
permit persons to whom the Software is furnished to do so, subject to
the following conditions:

The above copyright notice and this permission notice shall be
included in all copies or substantial portions of the Software.

THE SOFTWARE IS PROVIDED "AS IS", WITHOUT WARRANTY OF ANY KIND,
EXPRESS OR IMPLIED, INCLUDING BUT NOT LIMITED TO THE WARRANTIES OF
MERCHANTABILITY, FITNESS FOR A PARTICULAR PURPOSE AND
NONINFRINGEMENT. IN NO EVENT SHALL THE AUTHORS OR COPYRIGHT HOLDERS BE
LIABLE FOR ANY CLAIM, DAMAGES OR OTHER LIABILITY, WHETHER IN AN ACTION
OF CONTRACT, TORT OR OTHERWISE, ARISING FROM, OUT OF OR IN CONNECTION
WITH THE SOFTWARE OR THE USE OR OTHER DEALINGS IN THE SOFTWARE.
\endpreamble
%    \end{macrocode}
% Clobber files without asking
%    \begin{macrocode}
\askforoverwritefalse
%    \end{macrocode}
% Make the \LaTeX{} macro files (class and styles)
%    \begin{macrocode}
\generate{% 
  \file{rmaa.cls}{\from{rmaa.dtx}{package}}
  \file{rmaa.ins}{\from{rmaa.dtx}{ins}}
  \file{rmaacite.sty}{\from{rmaa.dtx}{rmaacite}}
  \file{adshtml.sty}{\from{rmaa.dtx}{adshtml-style}}
  \file{ADS.tex}{\from{rmaa.dtx}{ads-container}}
  }
%    \end{macrocode}
% Make the other non-\LaTeX{} files. These are done separately since
% they can't use `\%\%' as a comment character for the
% preamble/postamble (and Perl files can't have the postamble I don't
% think). 
%    \begin{macrocode}
\generate{\nopreamble\nopostamble% 
  \file{Makefile}{\from{rmaa.dtx}{makefile}}
  \file{adshtml.perl}{\from{rmaa.dtx}{adshtml-perl}}
  \file{makeads.sh}{\from{rmaa.dtx}{ads-shell-script}}
  \file{sample.latex2html-init}{\from{rmaa.dtx}{latex2html-init}}
  }

%% Local Variables: 
%% mode: LaTeX
%% mode: Font-lock
%% End:
%</ins>
%    \end{macrocode}

% \subsection{Extended citation syntax}
% Modified from a file called mnrascite.sty
% I no longer remember where I stole the original version of this from
% 
% Alterations to the cite commands to make them do more sensible things.
% N.B. the bibliography style should produce entries which have labels where
% the date is delimited by angle brackets '<>' for these macros to work
% properly. They will not give very sensible error messages if you do not do 
% this properly.
% 
% A new command |\scite{key}| is defined which produces a citation which is 
% suitable for insertion in a sentence. i.e., it does not put brackets around 
% the whole citation, only the date. Also |\pcite{key}| produce output with no
% parentheses whatsoever. 
% 
% Should redefine |\@bibitem| really to have any hope of compatibility
% the command |\lbibcite| is defined to deal with the citations.
%    \begin{macrocode}
%<*rmaacite>
\def\bibitem[#1<#2>]#3{\par\hangindent=1.5em\hangafter=1
  \if@filesw
  {\def\protect##1{\string ##1\space}
    \immediate\write\@auxout{\string\lbibcite{#3}{#1}{#2}}}
  \fi}
\@redefinebibitem % what it says
%    \end{macrocode}
% Note that the |\lbibcite| defines two types of citation one contains the names,
% the other the year
%    \begin{macrocode}
\def\lbibcite#1#2#3{\global\@namedef{b@#1}{#2}
                    \global\@namedef{by@#1}{#3}}

\def\@bcite#1#2{({#1\if@tempswa , #2\fi})}
\def\@pcite#1#2{#1\if@tempswa , #2\fi}
%    \end{macrocode}
% This declaration contains what is put between the name and the date
% \changes{3.11}{2003/01/31}{Moved the \cs{articlename} to after the key for consistency
% with the new natbib stuff.}
%    \begin{macrocode}
\def\@brcite{\ (}
\def\@ebrcite{)}
\def\@fmttop#1{\@fmtcite[\csname #1\endcsname]}
\def\@citex[#1]#2{%
%%  \if@filesw\immediate\write\@auxout{\string\citation{#2}}\fi
%% WJH added 16 Jan 2002
  \@for\@citeb:=#2\do{%
    \edef\@citeb{\expandafter\@firstofone\@citeb:\articlename}%
    \if@filesw\immediate\write\@auxout{\string\citation{\@citeb}}\fi}%
%% to here
  \def\@citea{}%
  \@cite{\@for\@citeb:=#2\do
    {%
      \edef\@citeb{\expandafter\@firstofone\@citeb:\articlename}%
      \@citea\def\@citea{;\ }\@ifundefined
      {b@\@citeb}{{\bf ?}\@warning {Citation `\@citeb' on page \thepage \space undefined}}%
      {\csname b@\@citeb\endcsname\@citebro\csname by@\@citeb\endcsname\@citebrc}}}%
  {#1}}

\def\cite{\@ifnextchar [{\let\@citebro=\space
                          \let\@citebrc=\relax
                          \let\@cite=\@bcite\@tempswatrue\@citex}
                        {\let\@citebro=\space
                          \let\@citebrc=\relax
                          \let\@cite=\@bcite\@tempswafalse\@citex[]}}
\def\pcite{\@ifnextchar [{\let\@citebro=\space
                          \let\@citebrc=\relax
                          \let\@cite=\@pcite\@tempswatrue\@citex}
                        {\let\@citebro=\space
                          \let\@citebrc=\relax
                          \let\@cite=\@pcite\@tempswafalse\@citex[]}}
\def\scite{\@ifnextchar [{\let\@citebro=\@brcite
                          \let\@citebrc=\@ebrcite
                          \let\@cite=\@pcite\@tempswatrue\@citex}
                        {\let\@citebro=\@brcite
                          \let\@citebrc=\@ebrcite
                          \let\@cite=\@pcite\@tempswafalse\@citex[]}}

\renewenvironment{thebibliography}
  {\section*{REFERENCES}\vspace*{\afterbiblioskip}\small
    \if@book\RereadBiblio\fi
    \sloppy\clubpenalty4000\widowpenalty4000
    \sfcode`\.=1000\relax
    \setlength\parindent{\z@}\setlength\parskip{\z@}}
  {\relax\normalsize\@outputaddressesifnecessary} 
%</rmaacite>
%    \end{macrocode}

% \subsection{Auxilliary files for generating ADS data}

% The ADS data are written by \LaTeX{} to a |.ads| file. This file
% then has to be included in the following file (|ADS.tex|),
% which is run through \LaTeX{}2HTML (see below). 

%    \begin{macrocode}
%<*ads-container>
%%
%% Instructions:
%% Edit the \input{} line below to reflect the name of your file. 
\documentclass{article}
\usepackage[latin1]{inputenc}
\usepackage{adshtml}
\newcommand{\ADSfield}[2]
{\%#1  #2\\}
%%\newcommand{\ADSvfield}[2]
%%{\%#1 \begin{verbatim}#2\end{verbatim}}
\newenvironment{ADSarticle}{}{\par\bigskip}
\renewcommand{\bibitem}[2][]{\par}
\setlength{\parindent}{0pt}
\renewcommand{\ThisVolume}[1]{%  
  2001, \RMAASC, 12, \oldpageref{#1:RMAAFirstPage}}  
 

\begin{document}
\makeatletter
%% Demonstration file for formatting an entire volume
\documentclass[book,journal,twocolumn]{rmaa}
%%%%%%%%%%%%%%%%%%%%%%%%%%%%%%%%%%%%%%%%%%%%%%%%%%%%%%%%%%%%%%%%
%%%%%%%%%%%%%%%%            PREAMBLE            %%%%%%%%%%%%%%%%
%%%%%%%%%%%%%%%%%%%%%%%%%%%%%%%%%%%%%%%%%%%%%%%%%%%%%%%%%%%%%%%%

\usepackage{paralist}
\usepackage{psfrag,color}
% Needed for landscape tables
\usepackage{pdflscape}

% Used in rm-minimal-test
\usepackage{kantlipsum}

% publication details of volume
\SetVolume{9}
\SetYear{2000}
\SetMonthNumeric{06}
\SetMes{junio}
% details of meeting
\SetConfTitle{Astrophysical Plasmas: Codes, Models, \& Observations}
\SetReunion{M\'exico, D. F.}{octubre 25--29, 1999}
\SetEditors{S. Jane Arthur, Nancy S. Brickhouse, \& Jos\'e Franco}

% this writes the .idx file that is used to generate the index 
\makeindex

%% This command makes the session headings come out bold in the
%% Table-of-Contents. Comment it out if you don't like it.
\BoldSessionTitles

%% This command makes a title page for each session
%% Comment it out if you don't like it.
\SessionTitlePages

%% Use this command if you want the abstract-only papers to be
%% indexed by ADS
\abstractsADS

%% At the special request of Alberto Carramin~ana, the following is a
%% way to suppress the author adresses at the end of each article
%\suppressfulladdresses

%% Use this command to add Received/accepted dates for each paper
% \printdates


\begin{document}

%%%%%%%%%%%%%%%%%%%%%%%%%%%%%%%%%%%%%%%%%%%%%%%%%%%%%%%%%%%%%%%%b
%%%%%%%%%%%%%%%%          FRONT MATTER          %%%%%%%%%%%%%%%%
%%%%%%%%%%%%%%%%%%%%%%%%%%%%%%%%%%%%%%%%%%%%%%%%%%%%%%%%%%%%%%%%
\frontmatter % start roman numbering

% format the TOC
\renewcommand{\RMAAbannerimage}{rmsc_blank}
\tableofcontents

%% preface, group photo, list of participants go here
%% use the `frontsection' environment (see editors' guide)
\begin{frontsection}{Preface}{%
    \listofauthors{S.~J.~Arthur, N.~S.~Brickhouse, J.~Franco}
    \indexauthor{Arthur, S. J.}
    \indexauthor{Brickhouse, N.}
    \indexauthor{Franco, J.}
    }
  These are the proceedings of a very nice conference that Jane and
  Pepe organised. This is the preface, which explains what a wonderful
  time was had by all and how grateful they are for all the people who
  gave them money or helped out.
  
  This is another paragraph of the preface, in case they have anything
  else to say.  For example, the Scientific Organizing Committee
  consisted of
  \begin{quote}
    J. Franco (Mexico, Chair), D. Cox (USA) and various other notables
  \end{quote}
  while the Local Organizing Committee comprised
  \begin{quote}
    S. J. Arthur (Chair), S. Kurtz, etc.
  \end{quote}

  \EditorSignature

\end{frontsection}

\begin{frontsection}[Photo]{Group Photograph}{\notitle}
\null\vfill
\centerline{Here goes the photograph.}
\vfill
\end{frontsection}

\begin{frontsection}[Participants]{List of Participants}{}
%% Here you put a formatted list of participants
\vfil
\dotfill First page of list of participants \dotfill
\vfil\clearpage

\null\vfil
\dotfill Second page of list of participants \dotfill
\vfil

\end{frontsection}

\begin{frontsection}[XXX]{Another front section}{}

\end{frontsection}

%%%%%%%%%%%%%%%%%%%%%%%%%%%%%%%%%%%%%%%%%%%%%%%%%%%%%%%%%%%%%%%%
%%%%%%%%%%%%%%%%          MAIN MATTER           %%%%%%%%%%%%%%%%
%%%%%%%%%%%%%%%%%%%%%%%%%%%%%%%%%%%%%%%%%%%%%%%%%%%%%%%%%%%%%%%%
\mainmatter % start arabic numbering

\begin{articles}
  %% Main articles go here
  %% Use \article{ARTICLENAME} for each individual article
  %% If you want any text to appear in brackets after the article
  %% title in the TOC and index, then put it as an optional first
  %% argument, e.g.,  \article[Invited Talk]{ARTICLENAME} 

  \session{The Contributed Papers}
  \article{rm-journal-example} % An illustrative article
  \article{rm-onepage} % another article, which has a plate
  \article{rm-minimal-test}
  \session{Conference Summary}
  \article[Invited Review]{rmsummsamp} % An illustrative conference summary

\end{articles}

%%%%%%%%%%%%%%%%%%%%%%%%%%%%%%%%%%%%%%%%%%%%%%%%%%%%%%%%%%%%%%%%
%%%%%%%%%%%%%%%%          BACK MATTER           %%%%%%%%%%%%%%%%
%%%%%%%%%%%%%%%%%%%%%%%%%%%%%%%%%%%%%%%%%%%%%%%%%%%%%%%%%%%%%%%%
\backmatter

%% Abstracts of presentations for which the
%% authors have not submitted a paper
\begin{abstracts}

  \title{An abstract of a poster}
  \author{W. J. Henney\altaffilmark{1}}
  \listofauthors{W.~J.~Henney}
  \indexauthor{Henney, W. J}
  \altaffiltext{1}{Where I am from done with altaffilmark/text}
  \abstract{This is a test abstract.   This is a test abstract.   This
    is a test abstract.   This is a test abstract.}
  

  \title{An abstract of another  poster}
  \author{W.~J.~Henney\altaffil{This is done with altaffil} and
    A.~N.~Other\altaffil{So is this} }
  \listofauthors{W.~J.~Henney \& A.~N.~Other}
  \indexauthor{Henney, W. J}
  \indexauthor{Other, A. N.}
  \abstract{This is a test abstract.   This is a test abstract.   This
    is a test abstract.   This is a test abstract.}

  \title{An abstract of yet another  poster}
  \author{W.~J.~Henney\altaffil{This is done with altaffil} and
    A.~N.~Other\altaffil{So is this} }
  \listofauthors{W.~J.~Henney \& A.~N.~Other}
  \indexauthor{Henney, W. J}
  \indexauthor{Other, A. N.}
  \abstract{This is a test abstract.   This is a test abstract.   This
    is a test abstract.   This is a test abstract.}

  \title{The abstract of a poster with a very long title to see how
    well it gets split up into lines}
  \author{A.~N.~Other\altaffil{The affiliation of A.~N.~Other is also
      very long so that it will go over several lines. His email
      address is anon@nowhere.org}}
  \listofauthors{A.~N.~Other}
  \indexauthor{Other, A. N.}
  \abstract{We will try and make this abstract longer so there is more
    stuff to put on the page and it doesn't look too empty. However, I
    have nothing whatsoever to say so will have to stop here. Let's
    repeat that a few times. We will try and make this abstract longer
    so there is more stuff to put on the page and it doesn't look too
    empty. However, I have nothing whatsoever to say so will have to
    stop here. Let's repeat that a few times. We will try and make
    this abstract longer so there is more stuff to put on the page and
    it doesn't look too empty. However, I have nothing whatsoever to
    say so will have to stop here. Let's repeat that a few times. We
    will try and make this abstract longer so there is more stuff to
    put on the page and it doesn't look too empty. However, I have
    nothing whatsoever to say so will have to stop here. Let's repeat
    that a few times.}
    
  \title[Comment to go in TOC]{The abstract of a second poster with a
    very long title to see how well it gets split up into lines}
  \author{A.~N.~Other\altaffil{The affiliation of A.~N.~Other is also
      very long so that it will go over several lines. His email
      address is anon@nowhere.org}}
  \listofauthors{A.~N.~Other}
  \indexauthor{Other, A. N.}
  \abstract{We will try and make this abstract longer so there is more
    stuff to put on the page and it doesn't look too empty. However, I
    have nothing whatsoever to say so will have to stop here. Let's
    repeat that a few times. We will try and make this abstract longer
    so there is more stuff to put on the page and it doesn't look too
    empty. However, I have nothing whatsoever to say so will have to
    stop here. Let's repeat that a few times. We will try and make
    this abstract longer so there is more stuff to put on the page and
    it doesn't look too empty. However, I have nothing whatsoever to
    say so will have to stop here. Let's repeat that a few times. We
    will try and make this abstract longer so there is more stuff to
    put on the page and it doesn't look too empty. However, I have
    nothing whatsoever to say so will have to stop here. Let's repeat
    that a few times.}
    
  \title{The abstract of a third poster with a very long title to see how
    well it gets split up into lines}
  \author{A.~N.~Other\altaffil{The affiliation of A.~N.~Other is also
      very long so that it will go over several lines. His email
      address is anon@nowhere.org}}
  \listofauthors{A.~N.~Other}
  \indexauthor{Other, A. N.}
  \abstract{We will try and make this abstract longer so there is more
    stuff to put on the page and it doesn't look too empty. However, I
    have nothing whatsoever to say so will have to stop here. Let's
    repeat that a few times. We will try and make this abstract longer
    so there is more stuff to put on the page and it doesn't look too
    empty. However, I have nothing whatsoever to say so will have to
    stop here. Let's repeat that a few times. We will try and make
    this abstract longer so there is more stuff to put on the page and
    it doesn't look too empty. However, I have nothing whatsoever to
    say so will have to stop here. Let's repeat that a few times. We
    will try and make this abstract longer so there is more stuff to
    put on the page and it doesn't look too empty. However, I have
    nothing whatsoever to say so will have to stop here. Let's repeat
    that a few times.}
    
  \title{The abstract of a fourth poster with a very long title to see how
    well it gets split up into lines}
  \author{A.~N.~Other\altaffil{The affiliation of A.~N.~Other is also
      very long so that it will go over several lines. His email
      address is anon@nowhere.org}}
  \listofauthors{A.~N.~Other}
  \indexauthor{Other, A. N.}
  \abstract{We will try and make this abstract longer so there is more
    stuff to put on the page and it doesn't look too empty. However, I
    have nothing whatsoever to say so will have to stop here. Let's
    repeat that a few times. We will try and make this abstract longer
    so there is more stuff to put on the page and it doesn't look too
    empty. However, I have nothing whatsoever to say so will have to
    stop here. Let's repeat that a few times. We will try and make
    this abstract longer so there is more stuff to put on the page and
    it doesn't look too empty. However, I have nothing whatsoever to
    say so will have to stop here. Let's repeat that a few times. We
    will try and make this abstract longer so there is more stuff to
    put on the page and it doesn't look too empty. However, I have
    nothing whatsoever to say so will have to stop here. Let's repeat
    that a few times.}
    
  \title{The abstract of a fifth poster with a very long title to see how
    well it gets split up into lines}
  \author{A.~N.~Other\altaffil{The affiliation of A.~N.~Other is also
      very long so that it will go over several lines. His email
      address is anon@nowhere.org}}
  \listofauthors{A.~N.~Other}
  \indexauthor{Other, A. N.}
  \abstract{We will try and make this abstract longer so there is more
    stuff to put on the page and it doesn't look too empty. However, I
    have nothing whatsoever to say so will have to stop here. Let's
    repeat that a few times. We will try and make this abstract longer
    so there is more stuff to put on the page and it doesn't look too
    empty. However, I have nothing whatsoever to say so will have to
    stop here. Let's repeat that a few times. We will try and make
    this abstract longer so there is more stuff to put on the page and
    it doesn't look too empty. However, I have nothing whatsoever to
    say so will have to stop here. Let's repeat that a few times. We
    will try and make this abstract longer so there is more stuff to
    put on the page and it doesn't look too empty. However, I have
    nothing whatsoever to say so will have to stop here. Let's repeat
    that a few times.}
    
  \title{The abstract of a sixth poster with a very long title to see how
    well it gets split up into lines}
  \author{A.~N.~Other\altaffil{The affiliation of A.~N.~Other is also
      very long so that it will go over several lines. His email
      address is anon@nowhere.org}}
  \listofauthors{A.~N.~Other}
  \indexauthor{Other, A. N.}
  \abstract{We will try and make this abstract longer so there is more
    stuff to put on the page and it doesn't look too empty. However, I
    have nothing whatsoever to say so will have to stop here. Let's
    repeat that a few times. We will try and make this abstract longer
    so there is more stuff to put on the page and it doesn't look too
    empty. However, I have nothing whatsoever to say so will have to
    stop here. Let's repeat that a few times. We will try and make
    this abstract longer so there is more stuff to put on the page and
    it doesn't look too empty. However, I have nothing whatsoever to
    say so will have to stop here. Let's repeat that a few times. We
    will try and make this abstract longer so there is more stuff to
    put on the page and it doesn't look too empty. However, I have
    nothing whatsoever to say so will have to stop here. Let's repeat
    that a few times.}
    
  \title{The abstract of a seventh poster with a very long title to see how
    well it gets split up into lines}
  \author{A.~N.~Other\altaffil{The affiliation of A.~N.~Other is also
      very long so that it will go over several lines. His email
      address is anon@nowhere.org}}
  \listofauthors{A.~N.~Other}
  \indexauthor{Other, A. N.}
  \abstract{We will try and make this abstract longer so there is more
    stuff to put on the page and it doesn't look too empty. However, I
    have nothing whatsoever to say so will have to stop here. Let's
    repeat that a few times. We will try and make this abstract longer
    so there is more stuff to put on the page and it doesn't look too
    empty. However, I have nothing whatsoever to say so will have to
    stop here. Let's repeat that a few times. We will try and make
    this abstract longer so there is more stuff to put on the page and
    it doesn't look too empty. However, I have nothing whatsoever to
    say so will have to stop here. Let's repeat that a few times. We
    will try and make this abstract longer so there is more stuff to
    put on the page and it doesn't look too empty. However, I have
    nothing whatsoever to say so will have to stop here. Let's repeat
    that a few times.}
  
  \title{An abstract to go in the list of abstracts, but not in the
    index of authors}
  \author{I.~A.~M. Hidden\altaffil{Puerto Escondido}}
  \listofauthors{I.~A.~M. Hidden}
  \abstract{The idea is that by omitting the
    \textbackslash{}indexauthor command, we can avoid an entry in the
    author index.}

\end{abstracts}


%% The Index: after latexing twice, run `makeindex -s rmaa' on this
%% file, then latex twice more
\printindex

%% if you rename this file to something other than `rmbooksamp', then
%% you must change the above line accordingly

%%
%% Plates: If you have any, then the following line will print
%% them out:

\outputplates

\end{document}

%%% Local Variables:
%%% mode: latex
%%% TeX-master: t
%%% End:

\end{document}
%</ads-container>
%    \end{macrocode}

% This is the |adshtml.sty| style file used by |ADS.tex|. It just
% defines a few 
% macros that might be used in abstracts, etc (really it should define
% them all). This now gets all its code from
% Sec~\ref{sec:user-commands}. 
%    \begin{macrocode}
%<*adshtml-style>
%% Nothing else needs to go here now. 
%</adshtml-style>
%    \end{macrocode}
% The init file for latex2html.
%    \begin{macrocode}
%<*latex2html-init>
$PDFTOCAIRO = "pdftocairo";
$PDFLATEX = "pdflatex";
$LATEX = "latex";
$MAX_SPLIT_DEPTH = 0;
$HTML_VERSION = "4.0,math";
$NO_SIMPLE_MATH = 1;
$ADDRESS = "<p><address><a href=\"mailto:rmaa\@astro.unam.mx\">rmaa\@astro.unam.mx</address></p>"; 
$INFO = "<p>This document was created using latex2html and the rmaa latex macros</p>"; 
$NO_NAVIGATION = 1;
1; 
%</latex2html-init>
%    \end{macrocode}
% This is the |adshtml.perl| file that is used by \LaTeX{}2HTML. Some
% bizarre \TeX{}y command sometimes show up so we tell it to ignore
% these. All the astronomy macros that come from AAS are processed in
% \TeX{}. That way they get made to images but with an |ALT| tag that
% has the original code, which is what will be picked up by |lynx|,
% with the result that they get passed intact into the final |.txt|
% file. 
%    \begin{macrocode}
%<*adshtml-perl>
package main;

$dum = &do_cmd_makeatletter();

sub do_cmd_RMAA {
    local($_) = @_;
    "RMxAA" . $_;
}

sub do_cmd_RMAASC {
    local($_) = @_;
    "RMxAASC" . $_;
}

sub do_cmd_etal {
    local($_) = @_;
    "et al\." . $_;
}

sub do_cmd_ibidrule {
    local($_) = @_;
    "<img src=\"foo\" alt=\"\ibidrule\">" . $_;
}


sub do_cmd_bibitem {
    local($_) = @_;
    local($key);
    local($opt, $dummy)=&get_next_optional_argument; 
    $key = &missing_braces unless
        ((s/$next_pair_pr_rx/$key = $2; ''/eo)
        ||(s/$next_pair_rx/$key = $2; ''/eo));
    "<p>" . $_; 
}

&ignore_commands (<<_RAW_ARG_CMDS_);
lowercase
adjustfinalcols
textsuperscript # {}
newline # []
thinmuskip
mskip
smallskip
medskip
bigskip
hspace # {}
hspacestar # {}
vspace # {}
vspacestar # {}
mathrm
_at_mathrm
spacefactor
discretionary
unhbox 
voidb_at_x 
penalty 
_at_M
_at_m
null
IeC
_RAW_ARG_CMDS_

&process_commands_in_tex (<<_RAW_ARG_CMDS_);
sun
earth
lesssim
gtrsim
sq
argdeg
arcmin
arcsec
fd
fh
fm
fs
fdg
farcm
farcs
fp
micron
la
ga
case # {} # {}
slantfrac # {} # {}
onehalf
onethird
twothirds
onequarter
threequarters
ubvr
ub
bv
vr
ur
aj
araa
apj
apjl
apjs
ao
apss
aap
aapr
aaps
azh
baas
jrasc
memras
mnras
pra
prb
prc
prd
pre
prl
pasp
pasj
qjras
skytel
solphys
sovast
ssr
zap
nat
iaucirc
aplett
apspr
bain
fcp
gca
grl
jcp
jgr
jqsrt
memsai
nphysa
physrep
physscr
planss
procspie
_RAW_ARG_CMDS_

1;                              # This must be the last line
%</adshtml-perl> 
%    \end{macrocode}

% This is a shell script |makeads.sh| that generates the |_ads.txt| file from
% |_ads.tex| by using \LaTeX{}2HTML and |lynx|.  \changes{4.0}{2019/11/13}{Make sure that lynx writes utf-8 output instead of latin-1.}
% \changes{2.2beta13}{2002/03/12}{Complete rewrite of |makeads.sh|. Now much 
% improved I think}
% \changes{4.0}{2019/11/13}{Fix the at-sign substitution so that it only eliminates \cs{@} and does not mess with email addresses.}

%    \begin{macrocode}
%<*ads-shell-script>
#! /bin/sh
if [ -z $1 ]; then
    echo "Usage: $0 basename"
    exit
fi
BASENAME=$1
ADSBASENAME=${BASENAME}_ads
AUXFILE=$BASENAME.aux
TMPFILE=tmp.$$


echo "Generating $ADSBASENAME.tex"

# For some reason, must double-up \\ to \\\\ but not \ to \\ - odd!
cat > $ADSBASENAME.tex <<EOF
%% PLEASE DO NOT EDIT THIS FILE!  It is generated automatically from
%% the script makeads.sh.  Any hand-edited changes will be lost
%% 
\documentclass{article}
\usepackage[latin1]{inputenc}
\usepackage{html}
\usepackage{adshtml}
\newcommand{\ADSfield}[2]{\%#1  #2\\\\}
\newenvironment{ADSarticle}{}{\par\bigskip}
\setlength{\parindent}{0pt}
\newcounter{istage}
\renewcommand{\ion}[2]{\setcounter{istage}{#2}#1 \Roman{istage}}
\begin{document}
\makeatletter
EOF
echo >> $ADSBASENAME.tex "\input{$BASENAME.ads2}" 
echo >> $ADSBASENAME.tex "\end{document}"



echo "Generating $ADSBASENAME.sed" 

# Find type of journal
grep '\\documentclass\[.*proceedings.*\]' $BASENAME.tex >/dev/null \
     && JOURNAL='RMxAC' \
     || JOURNAL='RMxAA'

# Find the year and volume number
YEAR=`grep '\\SetYear' $BASENAME.tex | sed 's/.*\\SetYear{\(.*\)}.*/\1/'`
VOLUME=`grep '\\SetVolume' $BASENAME.tex | sed 's/.*\\SetVolume{\(.*\)}.*/\1/'`

# Generate an auxilliary sed script 
cat > $TMPFILE <<EOF
s|\\\\newlabel{\(.*\):RMAAFirstPage}{{.*}{\(.*\)}}|s/\\\\\\\\ThisVolume{\1}/$YEAR, RMxAC, $VOLUME, \2/g|g
EOF

# Now go through the .aux file, looking for \newlabels and use the sed 
# script just generated in order to make another sed script, one that 
# will replace \ThisVolume{LABEL} with the correct reference
grep '^\\newlabel{.*:RMAAFirstPage}' $AUXFILE \
     | sed -f $TMPFILE > $ADSBASENAME.sed && rm $TMPFILE

# Apend a further sed command to get rid of \@, since redefining 
# it in _ads.tex didn't work
echo >> $ADSBASENAME.sed 's/\\\\@//g'

echo "Generating $BASENAME.ads2 from $ADSBASENAME.sed and $BASENAME.ads" 

# Apply this to the .ads file, generating the .ads2 file, which is included
# by the _ads.tex file.
sed -f $ADSBASENAME.sed $BASENAME.ads > $BASENAME.ads2

echo "Generating $ADSBASENAME.txt from $ADSBASENAME.tex and $BASENAME.ads2"

# Now run latex2html on the _ads.tex file
latex2html -no_navigation -info 0 -address 0 -html_version 4.0,math $ADSBASENAME

# Then, convert to utf8 with lynx and tidy it up a bit
lynx --justify=off --width=100 --dump --nolist --display_charset=utf-8 \
     $ADSBASENAME/$ADSBASENAME.html | cut -c4- \
     | sed 's/ibidrule/\\ibidrule/g' > $ADSBASENAME.txt
%</ads-shell-script>
%    \end{macrocode}

% \section{Epilogue}
% Err, that's it \dots

\endinput


% \endinput
% Local Variables:
% mode: doctex
% TeX-master: t
% End:
"

rmaa.aux: rmaa.dtx rmaa.ind rmaa.gls
        @echo "REGENERATING .aux FILE"
        pdflatex "\nonstopmode% \def\fileversion{4.0} \def\filedate{2019/11/12}
% \MakeShortVerb{\|}
% \title{A \LaTeX{} document class for the Revista Mexicana de
% Astronom\'\i a y Astrof\'\i sica\thanks{This file has version
% \fileversion{} last revised \filedate{}.}}
% \author{William J. Henney\thanks{w.henney@irya.unam.mx}}
% \date{\filedate}
% \maketitle
% \paragraph{Abstract}
% This document describes the \texttt{rmaa} \LaTeX{}
% document class for the journal ``Revista Mexicana de
% Astronom\'\i a y Astrof\'\i sica''. The macros can be used for 
% \begin{enumerate}
% \item Preparing author manuscripts for the journal (including
% conference proceedings to be published in the ``Serie de
% Conferencias'')
% \item Typesetting an entire volume of the journal by inputting the
% individual article files from a container file. Facilities are
% included for automatic generation of TOC and Author Index. 
% \end{enumerate}
% 
% \tableofcontents
% \section{The documentation driver file}
% This is the part that generates the documentation that you may or
% may not be reading, which is done with the help of the Makefile
% discussed in section~\ref{sec:driver_make}. 
%    \begin{macrocode}
%<*driver>
\documentclass{ltxdoc}
\usepackage{hyperref}
\usepackage{palatino}
\EnableCrossrefs \CodelineIndex \RecordChanges
\setlength{\parindent}{0pt}
\def\MacroFont{\footnotesize\tt}
\begin{document} \DocInput{rmaa.dtx} \PrintIndex \PrintChanges \end{document}
%</driver>
%    \end{macrocode}

% \section{The implementation}
% 
% \changes{3.24}{2007/09/14}{Added user access to the macro version
% and date via \cs{rmVersion} and \cs{rmDate}}
% 
%    \begin{macrocode}
%<*package>
\newcommand{\rmaa@version}{3.26}
\newcommand{\rmaa@date}{2011/08/15}
\newcommand\rmVersion{\rmaa@version}
\newcommand\rmDate{\rmaa@date}
\NeedsTeXFormat{LaTeX2e}[1995/06/01]
%    \end{macrocode}
% \changes{4.0}{2019/11/02}{Change abstract to appear before
% resumen. Change RevMexAA to RMxAA for consistency with ADS.}
%
% \changes{3.22}{2007/03/22}{Allow proceedings articles to have
%   received/accepted dates} \changes{3.21}{2006/05/15}{Changes for
%   Silvia - abstracts are now shipped out to ADS.}
% \changes{3.19}{2006/01/19}{Changes for Silvia - remove adresses from
%   end, new abstract style} \changes{2.2beta}{2001/10/30}{ Fixed bug
%   whereby twocolumn option was not passed through to book
%   class. Finally sorted out the ``label on last page'' issue. }
% \changes{2.1}{2001/01/23}{ Fixed some inconsistencies in
%   \cs{nocite} and \cs{bibitem} uniqification Fixed the definition of
%   \cs{ion} so it will work in math mode and in contexts where the
%   typesize is different from normal (e.g., title, fig captions).
%   Modified the definition of \cs{@caption}, so we don't write any
%   list-of-figures data (we never use it anyway, and it just causes
%   problems with fragile commands in fig captions). }
% \changes{2.0}{2000/11/30}{ Added `newstyle' option, which should be
%   used in all new proceedings volumes (wider abstract, smaller
%   typeface in fig captions, etc.) Added commands for formatting
%   table captions. New demonstation document rmpostsamp.tex for
%   1-page posters. This will be the last release under a PRI
%   government (for a while...)}  \changes{1.9}{2000/06/06}{
%   Corrected position of page number in footer of articles that open
%   on even pages. Ensured that first article in each session opens on
%   an odd page, as do list of abstracts and author index. Added extra
%   space in TOC before the backmatter. Added command for editors,
%   \cs{BoldSessionTitles}, that makes produces bold titles for
%   session in the TOC.}  \changes{1.8}{2000/05/16}{ Fixed bug in
%   last page label with \cs{suppressfulladdresses} }
% \changes{1.7}{2000/05/11}{ Fixed more little things. Page labels
%   now work, even if you have no bibliography. You can output the
%   full addresses by hand with the \cs{outputfulladdresses}
%   command. Editors' Guide updated but still not finished. Revamped
%   splitbook.sh now produces PS file of each article for web page and
%   also an EPS file of each page for sending to the printers.}
% \changes{1.6}{2000/04/15}{ More cleaning up: full stops after
%   titles in index, running headers and footers in frontmatter,
%   fontsize in title of index, etc} \changes{1.5}{2000/04/15}{
%   Eliminated superscripts from titles in TOC and index}
% \changes{1.4}{2000/04/14}{ Added copyright and license info (GPL)}
% \changes{1.3}{2000/03/01}{ Corrected \cs{figcaption} to do the
%   caption formatting right Corrected font size of sec headings in
%   front matter} \changes{1.2}{2000/02/16}{ Added \cs{plate}
%   command} \changes{1.1}{2000/02/14}{ Added manuscript option,
%   corrected bugs} \changes{1.0}{2000/02/08}{ Now does a whole book
%   - lots of stuff tidied up} \changes{0.7}{1999/10/04}{ Increased
%   spacing between sections in single column styles. Suppressed
%   floats at bottom of pages (didn't work though) }
% \changes{0.6}{1999/09/24}{ Tweaked \cs{adjustfinalcols} (still not
%   perfect) Added proceedings option Added onecolumn/twocolumn
%   options (orthogonal to journal/proceedings) }
% \changes{0.5}{1999/09/23}{ Added \cs{fulladdresses} and
%   \cs{adjustfinalcols} (just about working)}
% \changes{0.4}{1999/09/22}{ Minor stylistic changes}
% \changes{0.3}{1999/06/28}{ Better mechanism for alternative
%   affiliations Added `guide' option for user guide}
% \changes{0.2}{1999/05/27}{ All basic functionality present, I
%   think} \changes{0.1}{1999/05/26}{ Just about working - still
%   missing Fig/Tab/Refs} \changes{0.0}{1999/05/21}{ So bored today,
%   decided to rewrite RMAA LaTeX macros from scratch}
%    \begin{macrocode}
\ProvidesClass{rmaa}[\rmaa@date\space v\rmaa@version\space RMxAA article
class]
%    \end{macrocode}
% \subsection{Option Handling}
% \label{sec:options}
% Declarations that are needed in options
% \changes{3.05}{2003/10/22}{Added natbib option}
% \changes{3.06}{2003/01/03}{Natbib option turned on permanently}
% \changes{3.07}{2003/01/11}{Added proof option}
% \changes{3.14}{2003/01/11}{Option to turn off natbib}
% \changes{3.18}{2003/10/22}{trying out doctex mode} 
% \changes{3.23}{2007/09/09}{Manuscript is now on by default.}
% \changes{3.26}{2011/08/15}{New option for review articles}
%    \begin{macrocode}
\newif\if@reviewarticle
\@reviewarticlefalse
\newif\if@proof
\@prooffalse
\newif\if@journal
\@journaltrue
\newif\if@printdates
\@printdatesfalse
\newif\if@proceedings
\@proceedingsfalse
\newif\if@manuscript
\@manuscripttrue
\newif\if@bigfont
\@bigfonttrue
\newif\if@guide
\@guidefalse
\newif\if@preprint
\@preprintfalse
\newif\if@rmtwocol
\@rmtwocoltrue
\newif\if@biblist
\@biblistfalse
\newif\if@book
\@bookfalse
\newif\if@natbib
\@natbibtrue
\newif\if@rmcite
\@rmcitefalse
\newif\if@titlevisible
\@titlevisibletrue
\newif\if@abstractpresent
\@abstractpresentfalse
\newif\if@printbanner
\@printbannerfalse
\newif\if@suppressfulladdresses
\@suppressfulladdressesfalse
\newif\if@newstyle
\@newstylefalse
\newif\if@debug
\@debugfalse
\newif\if@boldkeywords
\@boldkeywordsfalse
%    \end{macrocode}
% The option handling---mainly setting various booleans.
% \changes{2.2beta11}{2001/11/26}{Made twocol and
% newstyle the default for proceedings.}
% \changes{3.23}{2007/09/09}{Manuscript is now 10 pt instead of 12pt}
% \changes{3.23}{2007/09/10}{Full addresses are now automatically
% suppressed for proceedings style.}
% \changes{3.24}{2007/09/14}{New option \texttt{boldkeywords} that
% does the keywords in the old ugly way. Now off by default.}
%    \begin{macrocode}
\DeclareOption{manuscript}{\typeout{Manuscript for referee} 
  \@manuscripttrue\@journalfalse\@rmtwocolfalse
  \PassOptionsToClass{10pt}{article}}
\DeclareOption{bigfont}{\typeout{Big font manuscript requested by Gary
  Ferland} 
  \@manuscripttrue\@journalfalse\@rmtwocolfalse\@bigfonttrue
  \PassOptionsToClass{11pt}{article}}
\DeclareOption{journal}{\typeout{Main journal option}%
  \@journaltrue\@rmtwocoltrue\@printdatestrue\@manuscriptfalse
  \PassOptionsToClass{10pt}{article}}
\DeclareOption{guide}{\typeout{User guide option}%
  \@guidetrue\@journalfalse\@rmtwocoltrue\@manuscriptfalse
  \PassOptionsToClass{10pt}{article}}
\DeclareOption{preprint}{\typeout{Preprint option}%
  \@preprinttrue\@manuscriptfalse\PassOptionsToClass{10pt}{article}}
\DeclareOption{proof}{\typeout{Author proof option}%
  \@prooftrue\@manuscriptfalse}
\DeclareOption{proceedings}{\typeout{Proceedings option}%
  \@proceedingstrue\@journalfalse\@rmtwocoltrue\@newstyletrue
  \@suppressfulladdressestrue
  \@printdatesfalse\@manuscriptfalse\PassOptionsToClass{10pt}{article}} 
\DeclareOption{printdates}{\typeout{Force printing of received/accepted dates}%
  \@printdatestrue} 
\DeclareOption{book}{\typeout{Whole issue option}%
  \@booktrue} 
\DeclareOption{nonatbib}{\typeout{Turn off natbib support (required
    for compatibility with old rmaacite package).}%
  \@natbibfalse} 
\DeclareOption{newstyle}{\typeout{Use new layout style}%
  \@newstyletrue}
\DeclareOption{oldstyle}{\typeout{Use old layout style}%
  \@newstylefalse}
\DeclareOption{debug}{\typeout{Loads of annoying messages}%
  \@debugtrue}
%    \end{macrocode}
% These two need to be declared explicitly so that |\ExecuteOptions|
% actually does something.
%    \begin{macrocode}
\DeclareOption{twocolumn}{\@rmtwocoltrue}
\DeclareOption{onecolumn}{\@rmtwocolfalse}
\DeclareOption{twoside}{\PassOptionsToClass{\CurrentOption}{article}}
\DeclareOption{boldkeywords}{\@boldkeywordstrue}
%    \end{macrocode}
% Bold, upper-case keywords used to be the default. Now they must be
% turned on explicitly with the \texttt{boldkeywords} option. 
%    \begin{macrocode}
\DeclareOption{boldkeywords}{\@boldkeywordstrue}
%    \end{macrocode}
% Pass any other option along to the \texttt{article} class.
%    \begin{macrocode}
\DeclareOption*{\PassOptionsToClass{\CurrentOption}{article}}
%    \end{macrocode}
% Default options are to be two-sided and based on the standard
% \LaTeX{} \texttt{article} class,
% unless we are processing the entire book, in which case we base
% ourselves on the standard \texttt{book} class. Default is to use
% manuscript option unless otherwise specified. 
%    \begin{macrocode}
\ExecuteOptions{twoside,manuscript} \ProcessOptions\relax
\if@book
\if@rmtwocol\LoadClass[twocolumn]{book}\else\LoadClass{book}\fi
\newcommand{\abstract}{\relax}
\newcommand{\@maketitle}{\relax}
\else
\if@rmtwocol\LoadClass[twocolumn]{article}\else\LoadClass{article}\fi
\fi
%    \end{macrocode}
% Attempt to use natbib for citations 02 Jan 2003. 
% \changes{3.05}{2003/01/02}{If the natbib option is used, then set
% things up so that natbib's version of the bibliography is identical
% to the standard one.}
% \changes{3.09}{2003/01/23}{Explicitly require a recent version of
% natbib, so that we get a decent error message if the package is too
% old.} 
% \changes{3.23}{2007/09/09}{Now allow Bibtex to work properly} 
% \changes{3.25}{2009/10/01}{Change \cs{newcommand} tp \cs{def} when setting \cs{bibfont}. This allows it to work with all versions of natbib.} 
% \changes{3.25}{2010/02/06}{Updated copyright year.} 
% \changes{3.25}{2010/02/06}{Updated email and web addresses to \texttt{crya.unam.mx}.} 
%    \begin{macrocode}
\if@natbib
\RequirePackage{natbib}[1999/02/25]
\bibpunct{(}{)}{;}{a}{}{,}
\def\bibfont{\small}
\renewcommand\bibsection{\section*{REFERENCES}}
\setlength\bibsep{0pt}
\setlength\bibhang{1.5em}
\bibliographystyle{rmaa}
\fi
%    \end{macrocode}
% We want the \texttt{graphicx} package for the extended syntax of
% |\includegraphics|. We want the \texttt{ifthen} package for using in
% these macros (although we don't much yet). We explicitly load
% \texttt{babel} since some \LaTeX{} distributions messed up the
% hyphenation if we didn't.  
%    \begin{macrocode}
\RequirePackage{graphicx}
\RequirePackage{ifthen}
\IfFileExists{babel.sty}{%
%   \let\@restoreindent=\@afterindentfalse
  \RequirePackage[english]{babel}
%   \@restoreindent
  %% reinstate the definition of \@roman so as not to break .idx files
  \@ifundefined{org@@roman}{\relax}{\renewcommand{\@roman}[1]{\org@@roman{##1}}}
  }{%
  \typeout{Package babel.sty not found - don't worry about this}}
%    \end{macrocode}
% Try to load the \texttt{dcolumn} and \texttt{booktabs} packages if
% they can be found. These can produce nicer-looking tables. 
%    \begin{macrocode}
\IfFileExists{dcolumn.sty}{%
  \RequirePackage{dcolumn}
  }{%
  \typeout{Package dcolumn.sty not found - you will not be able to
    align columns on decimal points in tables}}
\IfFileExists{booktabs.sty}{%
  \RequirePackage{booktabs}
  }{%
  \typeout{Package booktabs.sty not found - your tables will not look
    as pretty as they ought too.}
  }
\IfFileExists{makeidx.sty}{%
  \RequirePackage{makeidx}
}{%
  \typeout{Package makeidx.sty not found - Author Index may not work.}
  }
%% commented out for moment since requires more modern version of LaTeX
%% \RequirePackage[latin1]{inputenc}
%    \end{macrocode}
% \changes{4.0}{2019/11/13}{Make \texttt{hyperref} and
% \texttt{etoolbox} into hard requirements, since we now use them
% extensively.}  We use hyperref to get internal and external links.
% But I don't like the colored boxes so I turn those off.  We use
% etoolbox for patching and robustifying at the moment.  In the
% future, I want to port a lot of other stuff over to it.  For
% instance, replace the ifthen package and the TeX ifs too.
%
% We target the latest version (1.7) of PDF, rather than the default
% of 1.5.  This reduces the chance of getting warnings from included
% graphics.  This is not strictly related to hyperref, but we put it
% here anyway.
%    \begin{macrocode}
\RequirePackage{etoolbox}
\RequirePackage{hyperref}
\hypersetup{hidelinks}
\pdfminorversion=7
%    \end{macrocode}
% We patch the \cs{@combinedblfloats} following a suggestion by Heiko
% Oberdiek at \url{https://tex.stackexchange.com/a/249743/3843}, which
% avoids problems when hyperlinks are split across
% pages. \changes{4.0}{2019/11/12}{Fix for issue where breaking
% hyperlink across pages gave an error.}
%    \begin{macrocode}
\urlstyle{same}
\newcommand\email{\begingroup \urlstyle{rm}\Url}
\patchcmd\@combinedblfloats{\box\@outputbox}{\unvbox\@outputbox}{}{%
  \errmessage{\noexpand\@combinedblfloats could not be patched}%
}%
%    \end{macrocode}
% \subsection{Other setup stuff}\label{sec:setup}
% Some declarations: 
%    \begin{macrocode}
\newif\if@needaddresses
\@needaddressestrue
\newif\if@frontmatter
\@frontmatterfalse
\newif\if@inpreamble
\@inpreambletrue
\newif\if@firstauthor
\@firstauthortrue
\newif\if@firstaffil
\@firstaffiltrue
\newif\if@firstaltaffil
\@firstaltaffiltrue
\newif\if@firstkwd
\@firstkwdtrue
\newif\if@noexpandADS
\@noexpandADSfalse
\newif\if@sessiontitlepages
\@sessiontitlepagesfalse
%    \end{macrocode}
% \begin{macro}{\@rmaaEmpty}
% \begin{macro}{\@rmaaArg}
% These are used in various commands to check if arg was left blank
%    \begin{macrocode}
\newcommand{\@rmaaEmpty}{}
\newcommand{\@rmaaArg}{}
%    \end{macrocode}  
% \end{macro}
% \end{macro}
%
% \subsection{Compatibility with old \LaTeX{} versions}
% \label{sec:compat}
% This package is based on the 1998/12/01 release of \LaTeX{}.
% Older releases may not define all the commands that we need.
% Luckily, we can use |\providecommand| to define these if necessary. 
% The following commands are not defined in the 1994/06/01 release. 
%
% \begin{macro}{\MakeUpperCase}
% \begin{macro}{\textsuperscript}
% \begin{macro}{\textbackslash}
% User commands:
%    \begin{macrocode}
\providecommand{\MakeUpperCase}[1]{\uppercase{#1}}
\providecommand{\textsuperscript}[1]{\ensuremath{^\textrm{#1}}}
\providecommand{\textbackslash}{\ensuremath{\backslash}}
%    \end{macrocode}  
% \end{macro}
% \end{macro}
% \end{macro}
% \begin{macro}{\protected@xdef}
% \begin{macro}{\restore@protect}
% \begin{macro}{\@unexpandable@protect}
% \begin{macro}{\hb@xt@}
% \LaTeX{} internal commands:
%    \begin{macrocode}
\providecommand{\protected@xdef}{%
        \let\@@protect\protect
        \let\protect\@unexpandable@protect
        \afterassignment\restore@protect
        \xdef}
\providecommand{\restore@protect}{\let\protect\@@protect}
\providecommand{\@unexpandable@protect}{\noexpand\protect\noexpand}
\providecommand{\hb@xt@}{\hbox to}
%    \end{macrocode}
% \end{macro}
% \end{macro}
% \end{macro}
% \end{macro}
%
% \subsection{Page layout}\label{sec:page-layout}
%
% The dimensions of the text were defined by Elizabeth. Text width is
% 169\,mm with equal margins on the two sides: 
%    \begin{macrocode}
\setlength{\textwidth}{169mm}
\setlength{\@tempdima}{\paperwidth}
\addtolength{\@tempdima}{-\textwidth}
\setlength{\oddsidemargin}{0.5\@tempdima}
\addtolength{\oddsidemargin}{-1in}
\setlength{\evensidemargin}{\oddsidemargin}
%    \end{macrocode}
% Text height is 225\,mm, top margin is 14\,mm. 
%    \begin{macrocode}
\setlength{\textheight}{225mm} 
\setlength{\topmargin}{-1in} 
\addtolength{\topmargin}{14mm}
%    \end{macrocode}
% Space between columns in two-column styles, header, footer,
% paragraph indentation:
%    \begin{macrocode}
\setlength{\columnsep}{9.5mm}
\setlength{\headsep}{0.7cm}
\setlength{\footskip}{0.75cm}
\setlength{\parindent}{1.5em}
%    \end{macrocode}
% 
% \begin{macro}{\@set@scalable@length}
% Try out a new scheme for dealing with various spacing dimensions and
% allowing the user to easily modify them. The command
% \begin{quote}
% |\set@scalable@length|\marg{name}\marg{defaultlength}\marg{stretch}\marg{shrink}
% \end{quote}
% does a 
% |\setlength|\marg{\cs{name}}\marg{defaultlength plus stretch minus shrink}
% and also defines the user command
% \textbackslash\meta{name}|Scale|\marg{factor} that can 
% be used to rescale length by \meta{factor}. Additionally, now
% defines the package command
% \textbackslash\meta{name}|@set@scale|\marg{FACTOR} 
% that also rescales the length by \meta{FACTOR}, so that the length
% becomes \meta{factor} $\times$ \meta{FACTOR} $\times$
% \meta{defaultlength}.   
%    \begin{macrocode}
\newcommand\@set@scalable@length[4]{%
  %% arg #1: NAME of length without leading `\'
  %%
  %% arg #2: default VALUE  
  \expandafter\newlength\csname @#1@default\endcsname
  \expandafter\setlength\csname @#1@default\endcsname{#2}%
  \expandafter\newlength\csname @@#1@@default\endcsname
  \expandafter\setlength\csname @@#1@@default\endcsname{#2}%
  %% arg #3: default STRETCH  
  \expandafter\newlength\csname @#1@stretch\endcsname
  \expandafter\setlength\csname @#1@stretch\endcsname{#3}%
  \expandafter\newlength\csname @@#1@@stretch\endcsname
  \expandafter\setlength\csname @@#1@@stretch\endcsname{#3}%
  %% arg #4: default SHRINK  
  \expandafter\newlength\csname @#1@shrink\endcsname
  \expandafter\setlength\csname @#1@shrink\endcsname{#4}%
  \expandafter\newlength\csname @@#1@@shrink\endcsname
  \expandafter\setlength\csname @@#1@@shrink\endcsname{#4}%
  %%
  %% User-defined scale \@NAME@scale is unity by default
  \expandafter\gdef\csname @#1@scale\endcsname{1}%
  %% Package-defined scale \@NAME@SCALE is unity by default
  \expandafter\gdef\csname @#1@SCALE\endcsname{1}%
  %%
  %% internal command \@reset@NAME should expand to 
  %% \setlength\@@NAME@@default{\@NAME@SCALE\@NAME@default}
  %% \setlength\NAME{\@NAME@scale\@@NAME@@default}
  %% in other words: scale x SCALE x default 
  \expandafter\gdef\csname @reset@#1\endcsname{%
    \expandafter\setlength\csname @@#1@@default\endcsname{%
      \csname @#1@SCALE\endcsname \csname @#1@default\endcsname
      }
    \expandafter\setlength\csname @@#1@@stretch\endcsname{%
      \csname @#1@SCALE\endcsname \csname @#1@stretch\endcsname
      }
    \expandafter\setlength\csname @@#1@@shrink\endcsname{%
      \csname @#1@SCALE\endcsname \csname @#1@shrink\endcsname
      }
    \expandafter\setlength\csname #1\endcsname{%
      \csname @#1@scale\endcsname \csname @@#1@@default\endcsname
      \@plus
      \csname @#1@scale\endcsname \csname @@#1@@stretch\endcsname
      \@minus
      \csname @#1@scale\endcsname \csname @@#1@@shrink\endcsname
      }}
  %%
  %% User command \NAMEScale should expand to 
  %% {\renewcommand\@NAME@scale{#1}\@reset@NAME}
  \expandafter\gdef\csname #1Scale\endcsname##1{%
    \expandafter\gdef\csname @#1@scale\endcsname{##1}%
    \csname @reset@#1\endcsname}
  %%
  %% Package command \NAME@set@SCALE should expand to 
  %% {\renewcommand\@NAME@SCALE{#1}\@reset@NAME}
  \expandafter\gdef\csname #1@set@SCALE\endcsname##1{%
    \expandafter\gdef\csname @#1@SCALE\endcsname{##1}%
    \csname @reset@#1\endcsname}
  %%
  %% initialize at default values with \@reset@NAME
  \csname @reset@#1\endcsname
  }
%    \end{macrocode}
% \end{macro}
% \begin{macro}{\parskip}
% First use of this new infrastructure: set interparagraph space to be
% zero but with a small amount of stretchiness.
%    \begin{macrocode}
%   \ifthenelse{\boolean{@manuscript}}{%
%     \@set@scalable@length{parskip}{0.5\baselineskip}{0.3\p@}{0.1\p@}
%   }{
    \@set@scalable@length{parskip}{0\p@}{0.3\p@}{0\p@}
%   }
%    \end{macrocode}
% \end{macro}
% 
% \begin{macro}{\abovesectionskip}
% \changes{4.0}{2019/11/12}{Use negative values for skip so that first paragraph after section headings is not indented.  This is something that I have been trying to fix for years!  It turns out that the anser was in the documentation for the standard command \cs{@startsection.}}
% \begin{macro}{\belowsectionskip}
% \begin{macro}{\aboveparaskip}
% \begin{macro}{\belowparaskip}
% Define some utility rubber lengths for use in
% section~\ref{sec:section-commands}. 
%    \begin{macrocode}
\newlength{\abovesectionskip}
\newlength{\belowsectionskip}
\newlength{\aboveparaskip}
\newlength{\belowparaskip}
%    \end{macrocode}
% The default values for these lengths are those traditionally used
% for the onecolumn style. 
%    \begin{macrocode}
\@set@scalable@length{abovesectionskip}{-3.5ex}{-4ex}{-2ex}
\@set@scalable@length{belowsectionskip}{2.3ex}{0ex}{2ex}
\@set@scalable@length{aboveparaskip}{3.25ex}{2ex}{0.5ex}
\@set@scalable@length{belowparaskip}{-1em}{\z@}{\z@}
%    \end{macrocode}
% \end{macro}\end{macro}\end{macro}\end{macro}
% We also make scalable some of the standard float parameters. These
% are the same defaults as in the standard article classes, but
% slightly more stretchy. For the twocolumn style, they are currently
% halved, see below. Note that I have also lumped |\abovecaptionskip|
% in with these, since it plays a similar role. There is also a
% |\belowcaptionskip| but it is set to zero in the standard classes. 
%    \begin{macrocode}
\@set@scalable@length{floatsep}{12\p@}{4\p@}{4\p@}
\@set@scalable@length{textfloatsep}{20\p@}{4\p@}{8\p@}
\@set@scalable@length{intextsep}{12\p@}{4\p@}{4\p@}
\@set@scalable@length{dblfloatsep}{12\p@}{4\p@}{4\p@}
\@set@scalable@length{dbltextfloatsep}{20\p@}{4\p@}{8\p@}
\@set@scalable@length{abovecaptionskip}{10\p@}{\z@}{\z@}
%    \end{macrocode}
% \begin{macro}{\@before@title@skip}         
% \begin{macro}{\@before@author@skip}        
% \begin{macro}{\@between@author@skip}       
% \begin{macro}{\@before@abstract@skip}      
% \begin{macro}{\@after@abstracthead@skip}   
% \begin{macro}{\@between@abstract@skip}     
% \begin{macro}{\@before@kwd@skip}           
% \begin{macro}{\@before@dates@skip}           
% Spacing parameters used in |\@maketitle|. These are not changed by
% default except by the user, so they don't need to go into
% |\@package@default@lengths|.  
%    \begin{macrocode}
\newlength{\@before@title@skip}         
\newlength{\@before@author@skip}        
\newlength{\@between@author@skip}       
\newlength{\@before@abstract@skip}      
\newlength{\@after@abstracthead@skip}   
\newlength{\@between@abstract@skip}     
\newlength{\@before@kwd@skip}           
\newlength{\@before@dates@skip}           
\@set@scalable@length{@before@title@skip}{2em}{\z@}{\z@}
\@set@scalable@length{@before@author@skip}{1.5em}{\z@}{\z@}
\@set@scalable@length{@between@author@skip}{0.5em}{\z@}{\z@}
\@set@scalable@length{@before@abstract@skip}{3ex}{1ex}{0.5ex}
\@set@scalable@length{@after@abstracthead@skip}{1ex}{\z@}{\z@}
\@set@scalable@length{@between@abstract@skip}{2ex}{1ex}{0.5ex}
\@set@scalable@length{@before@kwd@skip}{1ex}{0.5ex}{0.2ex}
\@set@scalable@length{@before@dates@skip}{2ex}{0.5ex}{0.2ex}
%    \end{macrocode}
% \end{macro}\end{macro}\end{macro}\end{macro}
% \end{macro}\end{macro}\end{macro}\end{macro}
% 
% Spacing above and below displayed equations also wants to be done
% like this, but it is more complicated since the standard article
% classes make these change with the fontsize. Something to do later,
% maybe. 
% 
%    \begin{macrocode}
%    \end{macrocode}
% \begin{macro}{\@package@default@lengths}
% The value actually used for these lengths is controlled by the
% \meta{SCALE} parameter, which depends on the style option that is in
% effect. This is set using the |\@package@default@lengths| command,
% which must be called every time the style changes (e.g., changing
% from onecolumn to twocolumn). 
%    \begin{macrocode}
\newcommand{\@package@default@lengths}{%
  \ifthenelse{\boolean{@rmtwocol}}{%
    %% twocolumn mode
    \abovesectionskip@set@SCALE{0.5}
    \belowsectionskip@set@SCALE{0.5}
    \aboveparaskip@set@SCALE{0.5}
    \belowparaskip@set@SCALE{1}
    \floatsep@set@SCALE{0.5}
    \textfloatsep@set@SCALE{0.5}
    \intextsep@set@SCALE{0.5}
    \dblfloatsep@set@SCALE{0.5}
    \dbltextfloatsep@set@SCALE{0.5}
    }{%
    %% onecolumn mode
    \abovesectionskip@set@SCALE{1}
    \belowsectionskip@set@SCALE{1}
    \aboveparaskip@set@SCALE{1}
    \belowparaskip@set@SCALE{1}
    \floatsep@set@SCALE{1}
    \textfloatsep@set@SCALE{1}
    \intextsep@set@SCALE{1}
    \dblfloatsep@set@SCALE{1}
    \dbltextfloatsep@set@SCALE{1}
    }%
  }
\@package@default@lengths
%    \end{macrocode}
% \end{macro}
% \begin{macro}{\RescaleLengths}
% \begin{macro}{\RescaleSecLengths}
% \begin{macro}{\RescaleFloatLengths}
% \begin{macro}{\RescaleTitleLengths}
% User commands that will fiddle with all the scalable lengths. 
%    \begin{macrocode}
\newcommand{\RescaleLengths}[1]{%
  \RescaleSecLengths{#1}\RescaleFloatLengths{#1}\RescaleTitleLengths{#1}}
\newcommand{\RescaleSecLengths}[1]{%
%% sections
  \abovesectionskipScale{#1}
  \belowsectionskipScale{#1}
  \aboveparaskipScale{#1}
  }
\newcommand{\RescaleFloatLengths}[1]{%
%% floats
  \floatsepScale{#1}
  \textfloatsepScale{#1}
  \intextsepScale{#1}
  \dblfloatsepScale{#1}
  \dbltextfloatsepScale{#1}
  \abovecaptionskipScale{#1}
  }
\newcommand{\RescaleTitleLengths}[1]{%
%% maketitle
  \@before@title@skipScale{#1}           
  \@before@author@skipScale{#1}        
  \@between@author@skipScale{#1}       
  \@before@abstract@skipScale{#1}      
  \@after@abstracthead@skipScale{#1}   
  \@between@abstract@skipScale{#1}     
  \@before@kwd@skipScale{#1}           
  \@before@dates@skipScale{#1}           
  }
%    \end{macrocode}
% \end{macro}\end{macro}\end{macro}\end{macro}
% \changes{2.2beta14}{2002/03/17}{Finished implementation of new
% infrastructure for scalable lengths. Converted the above/below
% section skips to use it. Also, all the floatsep parameters. Also,
% all the vertical lengths that are used by |\@maketitle|. }

% \begin{macro}{\MSextramargin}
% For the \texttt{manuscript} option we have an extra 0.3\,inch
% vertical margin and 1\,inch horizontal margin for marginal notes (by
% hand or with |\marginpar|. 
%
% \changes{3.23}{2007/09/09}{Reduce baselinestretch in
%   \texttt{manuscript} style to 1.2. No longer use raggedright for
%   manuscript.}
%
% \changes{3.24}{2007/09/14}{Further reduction of baselinestretch in
%   \texttt{manuscript} style to 1.1 since we now have wider margins.}
%
% \changes{3.24}{2007/09/14}{Split \texttt{manuscript}'s extra margin
%   into separate horizontal and vertical parts.}
% 
%    \begin{macrocode}
\newlength{\MSextraHmargin}\newlength{\MSextraVmargin}
\setlength{\MSextraHmargin}{1in} 
\setlength{\MSextraVmargin}{0.3in} 
\if@manuscript
   \renewcommand{\baselinestretch}{1.1}
   \addtolength{\textwidth}{-2\MSextraHmargin}
   \addtolength{\oddsidemargin}{\MSextraHmargin}
   \setlength{\evensidemargin}{\oddsidemargin}
   \addtolength{\textheight}{-2\MSextraVmargin}
   \addtolength{\topmargin}{\MSextraVmargin}
%    \raggedright
\fi
%    \end{macrocode}
% \end{macro}
% Float parameters: don't allow bottom floats by default. 
%    \begin{macrocode}
\renewcommand{\fps@figure}{tp}
\renewcommand{\fps@table}{tp}
%    \end{macrocode}
%
% \begin{macro}{\rmaa@AuthorWidthFrac}
% \begin{macro}{\rmaa@AbstractWidthFrac}
% These parameters control what fraction of the total text width can
% be used by the list of authors and the abstract/resumen,
% respectively.  
%    \begin{macrocode}
\newcommand{\rmaa@AuthorWidthFrac}{1.0}
\newcommand{\rmaa@AbstractWidthFrac}{0.75}
%    \end{macrocode}
% Use a wide abstract to save space in the new style.
%    \begin{macrocode}
\if@newstyle\renewcommand{\rmaa@AbstractWidthFrac}{1.0}\fi
%    \end{macrocode}
% \end{macro}
% \end{macro}
%
% Some things to consider in future versions. Trying to avoid too many
% widows, etc. Apparently, the following will work best with
% \cs{flushbottom}. 
%    \begin{macrocode}
%%\widowpenalty=10000
%%\clubpenalty=10000
%    \end{macrocode}
% Also, should perhaps try getting rid of underfull/overfull box warnings. 
% \changes{3.10}{2003/01/24}{Set large values for \cs{hbadness},
% \cs{vbadness}, \cs{hfuzz}, \cs{vfuzz}. To reduce amount of stuff
% that gets written to the terminal every time you compile something.}
%    \begin{macrocode}
\hbadness=10000
\vbadness=10000
\setlength\hfuzz{\maxdimen}
\setlength\vfuzz{\maxdimen}
%    \end{macrocode}
% \subsection{Various text styles for different contexts}
% \label{sec:text-styles} 
% \begin{macro}{\rmaa@AffilStyle}
% \begin{macro}{\rmaa@PagesStyle}
% \begin{macro}{\rmaa@YearStyle}
% \begin{macro}{\rmaa@VolumeStyle}
% \begin{macro}{\rmaa@VolumenStyle}
% \begin{macro}{\rmaa@JournalNameStyle}
% \begin{macro}{\rmaa@ReceivedStyle}
% \begin{macro}{\rmaa@AbstractHeaderStyle}
% \begin{macro}{\rmaa@JournalNameStyle}
% Style commands for the different fields used by |\maketitle| in
% constructing the title page of the article. These have the form
% |\rmaa@XYZStyle{TEXT}|. I'm not sure these were really a good idea,
% but never mind \dots 
%    \begin{macrocode}
\newcommand{\rmaa@AffilStyle}[1]{{\normalsize\rmfamily#1}}
\newcommand{\rmaa@PagesStyle}[1]{{\normalsize\rmfamily#1}}
\newcommand{\rmaa@YearStyle}[1]{{\normalsize\rmfamily#1}}
\newcommand{\rmaa@VolumeStyle}[1]{{\normalsize\rmfamily\bfseries#1}}
\newcommand{\rmaa@VolumenStyle}[1]{{\Large\sffamily\bfseries#1}}
\newcommand{\rmaa@JournalNameStyle}[1]{{\normalsize\itshape#1}}
\newcommand{\rmaa@ReceivedStyle}[1]{{\normalsize\rmfamily\itshape#1}}
\newcommand{\rmaa@AbstractHeaderStyle}[1]{%
  {\normalsize\rmfamily\MakeUpperCase{#1}}}
\newcommand{\rmaa@TitleStyle}[1]{%
  {\large\bfseries\rmfamily\boldmath\MakeUpperCase{#1}\par}}
%    \end{macrocode}
% \end{macro}
% \end{macro}
% \end{macro}
% \end{macro}
% \end{macro}
% \end{macro}
% \end{macro}
% \end{macro}
% \end{macro}
%
% \begin{macro}{\rmaa@FrontSecStyle}
% Headings of ``front sections'' (see section~\ref{sec:front-sections})
%    \begin{macrocode}
\newcommand{\rmaa@FrontSecStyle}[1]{%
  {\normalsize\bfseries\rmfamily\MakeUpperCase{#1}}}
%    \end{macrocode}
% \end{macro}
%
% \begin{macro}{\rmaa@TOCTitleStyle}
% \begin{macro}{\rmaa@TOCPartStyle}
% \begin{macro}{\rmaa@IndexTitleStyle}
% \begin{macro}{\rmaaTOCAuthorStyle}
% \begin{macro}{\rmaaIndexAuthorStyle}
% Commands to set style of TOC and index entries.
% Rewritten to avoid unnecessary font commands in the .toc and .ind files
%    \begin{macrocode}
\newcommand{\rmaa@TOCTitleStyle}[1]{\MakeUpperCase{#1}}
\newcommand{\rmaa@TOCPartStyle}[1]{\rmaa@TOCTitleStyle{#1}}
\newcommand{\rmaa@IndexTitleStyle}[1]{#1}
\newcommand{\rmaaTOCAuthorStyle}[1]{%
  {\itshape#1}}
\newcommand{\rmaaIndexAuthorStyle}[1]{%
  \protect\mbox{\bfseries#1}}
%    \end{macrocode}
% \end{macro}
% \end{macro}
% \end{macro}
% \end{macro}
% \end{macro}
%
% \begin{macro}{\rmaa@KeywordStyle}
% \begin{macro}{\rmaa@KeywordHeaderStyle}
% The style of the keywords depends on whether we are using the
% \texttt{newstyle} option or not. If we are, then they are small. 
%    \begin{macrocode}
\ifthenelse{\boolean{@newstyle}}{
  \newcommand{\rmaa@KeywordStyle}[1]{%
    {\small%
      \if@boldkeywords 
      \bfseries\MakeUpperCase{\raggedright #1}
      \else
      \raggedright #1
      \fi
    }}
  \newcommand{\rmaa@KeywordHeaderStyle}[1]{%
    {\small\itshape#1}}}
{
  \newcommand{\rmaa@KeywordStyle}[1]{%
    {\normalsize%
      \if@boldkeywords\bfseries\fi\raggedright #1}}
  \newcommand{\rmaa@KeywordHeaderStyle}[1]{%
    {\normalsize\itshape#1}}}
%    \end{macrocode}
% \end{macro}
% \end{macro}
%
% \subsection{Page styles}\label{sec:page-styles}
% \begin{macro}{\ps@rmaa}
% The running headers for normal pages:
%    \begin{macrocode}
\newcommand{\ps@rmaa}{%
%    \end{macrocode}
% nothing in either foot; 
%    \begin{macrocode}
  \renewcommand\@oddfoot{\relax}%
  \renewcommand\@evenfoot{\relax}%
%    \end{macrocode}
% page number and short author list in the even heads; 
%    \begin{macrocode}
  \renewcommand\@evenhead{\rlap{\thepage}\hfil\@shortauthor\hfil}%
%    \end{macrocode}
% short title and page number in the odd heads.
%    \begin{macrocode}
  \renewcommand\@oddhead{\hfil\@shorttitle\hfil\llap{\thepage}}
  }
%    \end{macrocode}
% \begin{macro}{\ps@rmaafirstpage}
% The header of the first page of each
% article:
%    \begin{macrocode}
\newcommand{\ps@rmaafirstpage}{%
  \renewcommand\@oddfoot{\hfil\thepage}%
%    \end{macrocode}
% footer corrected (06 Jun 2000) to have page on outside edge;
%    \begin{macrocode}
  \renewcommand\@evenfoot{\thepage\hfil}%
%    \end{macrocode}
% Header has journal name, volume, and page numbers for a regular
% article, or similar for special cases. 
% \changes{4.0}{2019/11/03}{Change to three-line header with copyright info and DOI.}
%    \begin{macrocode}
\if@preprint
\renewcommand\@oddhead{%
  \if@proceedings 
  \textsc{To appear in ``\@conftitle{} (\rmaa@YearStyle{\@year})''}\hfil
  \textit{RMxAC}
  \else
  \textsc{To appear in} \@journalname\hfil
  \fi}\fi%
\if@guide
\renewcommand\@oddhead{\textsc{Documentation} (\rmaa@YearStyle{\@year})\hfil}
\fi%
\if@manuscript
\renewcommand\@oddhead{Manuscript for \@journalname{} (\rmaa@YearStyle{\@year})\hfil}
\fi%
\if@journal
\renewcommand\@oddhead{%
  \parbox{\textwidth}{
    \@journalname, \rmaa@VolumeStyle{\@volume},
    \rmaa@PagesStyle{\@artpages}\
    (\rmaa@YearStyle{\@year})\hfil\\[1pt]
    {\small\textcopyright{} \@year:
      Instituto de Astronom\'\i{}a,
      Universidad Nacional Aut\'onoma de M\'exico}\hfil\\[1pt]
    {\small\@DOI@url}\hfil
  }%
}\fi%
  \renewcommand\@evenhead{\@oddhead}}
%    \end{macrocode}
% By default, use the |\ps@rmaa| page style.
%    \begin{macrocode}
\pagestyle{rmaa}
%    \end{macrocode}
% \end{macro}
% \end{macro}
%
% \subsection{Preamble commands}\label{sec:preamble} 
% Everything you need for setting the title page, TOC and index
% entries, etc.
% \subsubsection{Commands for the title page}
% \label{sec:title-commands}
% \begin{macro}{\resumen}
% \begin{macro}{\@resumen}
% \begin{macro}{\abstract}
% \begin{macro}{\@abstract}
% User commands for setting the resumen and abstract. 
% These should be used in the document header. The internal commands
% |\@resumen| and |\@abstract| are used to hold the actual text
% pending output via the |\maketitle| command. 
%    \begin{macrocode}
\newif\if@setresumen
\newcommand{\resumen}[1]{%
%%  \renewcommand{\@resumen}{{\selectlanguage{spanish}#1}}%
  \renewcommand{\@resumen}{{\ignorespaces#1}}%
  \@setresumentrue}
\renewcommand{\abstract}[1]{%
  \renewcommand{\@abstract}{{\ignorespaces#1}}
  \renewcommand\ADS@abstract{#1}
%% Not a good idea, since then author macros are not expanded and
%% cause errors later. However, by not doing it we are vulnerable to
%% fragile commands.....
  \if@noexpandADS\@onelevel@sanitize{\ADS@abstract}\fi
  \@abstractpresenttrue
  }
\newcommand{\ADSnoexpand}{\@noexpandADStrue}
%    \end{macrocode}
% \begin{macro}{\boldabstract}
% This is the version for doing 1-page conference proceedings. Still
% save the text for ADS but then write it out immediately as a bold
% paragraph. 
%    \begin{macrocode}
\newcommand{\boldabstract}[1]{%
  \gdef\ADS@abstract{#1}%
  \renewcommand{\@abstract}{#1}%
  \if@noexpandADS\@onelevel@sanitize{\ADS@abstract}\fi%
  \noindent\textbf{\boldmath #1}
  }
%    \end{macrocode}
% \end{macro}
% Default value in each case is to print out a warning if not set
%    \begin{macrocode}
  \newcommand{\@resumen}{%
    \@setresumenfalse \ClassWarningNoLine{rmaa}{No Spanish abstract
      (resumen) specified - if you do not speak Spanish, this will be
      translated by the editors}%
    \textbf{Favor de proporcionar un resumen en espa\~nol. If you are
      unable to translate your abstract into Spanish, the editors will
      do it for you.} }
  \newcommand{\@abstract}{\ClassWarningNoLine{rmaa}{No abstract
      specified}}
%    \end{macrocode}
% \end{macro}
% \end{macro}
% \end{macro}
% \end{macro}
%
% \begin{macro}{\title}
% \begin{macro}{\@rawtitle}
% \begin{macro}{\@TOCtitle}
% \begin{macro}{\rmaaIndexTitle}
% Similarly, |\title| is the preamble command to set the title of the
% article. This is used in various places, typeset in different
% styles, hence the use of the internal |\@rawtitle|, |\@title|,
% |\rmaaIndexTitle|, and |\@TOCtitle| commands, each of which uses the
% relevant style from section~\ref{sec:text-styles}
%    \begin{macrocode}
\newcommand{\@rawtitle}{Raw Title}
\newcommand{\@TOCtitle}{\rmaa@TOCTitleStyle{Title for TOC}}
\newcommand{\rmaaIndexTitle}{\rmaa@IndexTitleStyle{Title for index}.}
\newcommand{\@TOCaftertitle}{}
\renewcommand{\title}[1]{%
%%  {\ADS@sanitize\xdef\@rawtitle{#1}}
  \renewcommand\@rawtitle{\noexpand#1}
  \let\ADS@title=\@rawtitle
  \renewcommand{\@title}{\rmaa@TitleStyle{#1}}%
  \renewcommand{\@TOCtitle}{\rmaa@TOCTitleStyle{#1}\@TOCaftertitle}%
  \renewcommand{\rmaaIndexTitle}{\rmaa@IndexTitleStyle{#1\@TOCaftertitle}.}%
  }
%    \end{macrocode}
% \end{macro}
% \end{macro}
% \end{macro}
% \end{macro}
% 
% All the weird and whacky ways of indicating the author
% affiliations. 
% \begin{macro}{\affil}
% First, the way that sets it directly underneath the authors name.
% Note that |\affil| trumps all other affiliations for the purposes of
% getting into the ADS stuff (we use |\gdef|).   
% \dots 
%    \begin{macrocode}
\newcommand{\affil}[1]{%
  \\[1.5ex]\rmaa@AffilStyle{#1}\\[1.5ex]%
  \if@firstaffil\gdef\ADS@affil{#1}\@firstaffilfalse\fi%
  }
\renewcommand{\and}{and\\[1.5ex]}
%    \end{macrocode}
% \begin{macro}{\altaffil}
% Alternative affiliations are footnoted. 
%    \begin{macrocode}
\newcommand{\altaffil}[1]{%
  \if@firstaltaffil\gdef\ADS@affil{#1}\@firstaltaffilfalse\fi%
  \thanks{#1}}
%    \end{macrocode}
% \begin{macro}{\altaffilmark}
% \begin{macro}{\altaffiltext}
% Alternative-alternative mechanism for affiliations---should not be
% mixed with |\altaffil| (unless you know what you are doing) or there
% will be problems with numbering.
%    \begin{macrocode}
\newcommand{\altaffilmark}[1]{\textsuperscript{#1}}
\newcommand{\altaffiltext}[2]{%
  \if@firstaltaffil\gdef\ADS@affil{#2}\@firstaltaffilfalse\fi
  \stepcounter{footnote}%
  \protected@xdef\@thanks{\@thanks
    \protect\footnotetext[#1]{#2}}} 
%    \end{macrocode}
% \end{macro}
% \end{macro}
% \end{macro}
% \end{macro}
%
% \begin{macro}{\keywords}
% \begin{macro}{\@keywords}
% \begin{macro}{\@keywordheader}
% \label{keywords}Sundry minor preamble commands. All
% self-explanatory, I hope. |\keywords| is the old user interface
% (deprecated). New articles should use |\addkeyword| once for each
% keyword. That way the ADS record comes out right too (comma rather
% than em-dash separator).
% \changes{2.2beta12}{2001/12/07}{\cs{addkeyword} New user interface
%   for keywords.} 
% \changes{2.2beta13}{2002/03/11}{\cs{addkeyword} Improved behaviour
%   with .ads file. Use
%   \cs{ADS@sanitize} to change definitions of `~' and \cs{ion}.}
% \changes{3.23}{2007/09/09}{\cs{addkeyword} now signals a compilation
% error if the author tries to add more than 6 keywords.}
% \changes{3.24}{2007/09/14}{kwdcounter is now reset in
% \cs{@begofarticlehook} to prevent error in book style} 
%    \begin{macrocode}
\newcommand{\@kwdsep}{ --- }
\newcommand{\ADS@kwdsep}{, }
\newcounter{kwdcounter}\setcounter{kwdcounter}{0}
\newcommand{\addkeyword}[1]{%
  \addtocounter{kwdcounter}{1}
  \ifthenelse{\thekwdcounter > 6}{
    \ClassError{rmaa}{Too many keywords. A maximum of 6 keywords is allowed.}{}
  }{}
  \if@firstkwd%
  \renewcommand\@keywords{\if@boldkeywords\uppercase{#1}\else #1\fi}
  {\ADS@sanitize\def\ion{\plainion}\gdef\ADS@keywords{#1}}
  \@firstkwdfalse
  \else%
  \edef\@keywords{\@keywords\@kwdsep \if@boldkeywords\uppercase{#1}\else #1\fi}%
  {\ADS@sanitize\def\ion{\plainion}\xdef\ADS@keywords{\ADS@keywords\ADS@kwdsep #1}}%
  \fi}
\newcommand{\ADS@sanitize}{%
  \def\mathrm{}
  \def\@{\ }
  \def~{\ }
  \def\langle{<}
  \def\rangle{>}
  \def\log{log }
  \def\sin{sin }
  \def\cos{cos }
  \def\tan{tan }
  \def\fd{^d.}
  \def\fh{^h.}
  \def\fm{^m.}
  \def\fs{^s.}
  \def\fdg{�.}
  \def\farcm{'.}
  \def\farcs{''.}
  \def\arcdeg{�}
  \def\,{ }
  }
\newcommand{\keywords}[1]{%
  \renewcommand{\ADS@keywords}{#1}
  \renewcommand{\@keywords}{\uppercase{#1}}
  }
\newcommand{\@keywordheader}{\rmaa@KeywordHeaderStyle{Key Words:}}
\newcommand{\@keywords}{%
  \ClassWarningNoLine{rmaa}{No keywords specified}%
  Keywords go here
  }
%    \end{macrocode}
% \begin{macro}{\shorttitle}
% \begin{macro}{\shortauthor}
% \begin{macro}{\@shorttitle}
% \begin{macro}{\@shortauthor}
% Commands for setting the running headers.
%    \begin{macrocode}
\newcommand{\shorttitle}[1]{%
  \renewcommand\@shorttitle{\MakeUpperCase{#1}}}
\newcommand{\shortauthor}[1]{%
  \renewcommand\@shortauthor{\MakeUpperCase{#1}}}
\newcommand{\@shorttitle}{SHORT TITLE GOES HERE}
\newcommand{\@shortauthor}{SHORT AUTHOR GOES HERE}
%    \end{macrocode}
% \end{macro}
% \end{macro}
% \end{macro}
% \end{macro}
% \end{macro}
% \end{macro}
% \end{macro}


% \subsubsection{The full addresses}
% \label{sec:full-addresses}
% \begin{macro}{\fulladdresses}
% \begin{macro}{\@fulladdresses}
% \begin{macro}{\@fulladdheight}
% \begin{macro}{\@fulladddepth}
% \begin{macro}{\@fulladdbox}
% Aaghh, the full addresses! A perennial source of problems. 
%    \begin{macrocode}
\newlength{\@fulladdheight}
\newlength{\@fulladddepth}
\newsavebox{\@fulladdbox}
%% \fulladdresses command modified (11 May 2000) so that we can put
%% \fulladdresses{} in the article to suppress the output of the
%% addresses
\newcommand{\@FinalFigure}{\relax}
\newcommand{\FinalFigure}[1]{\renewcommand{\@FinalFigure}{#1}}
\newcommand{\fulladdresses}[1]{%
  \renewcommand{\@rmaaArg}{#1}
  \ifx\@rmaaArg\@rmaaEmpty
     \renewcommand{\@fulladdresses}{}
  \else
     \renewcommand{\@fulladdresses}{
       \normalsize%
       \ifx\@FinalFigure\relax\else\vfill\@FinalFigure\vfill\fi
%%       \vspace*{1ex}
       %% do addresses as a list (copied from thebibliography environment) 
       \list{\null}{\leftmargin 1.5em\topsep\z@\labelwidth\z@\labelsep\z@\itemindent -1.5em
         \itemsep\z@\parsep\z@\usecounter{enumi}}
       \sloppy\clubpenalty4000\widowpenalty4000
       \sfcode`\.=1000\relax
       #1
       \endlist
       }
  \fi
  }
\newcommand{\@fulladdresses}{Full addresses go here}
%    \end{macrocode}
% \begin{macro}{\suppressfulladdresses}
% New user-level command for suppressing full addresses should you so
% desire. Currently used in the single-page-poster style. 
%    \begin{macrocode}
\newcommand{\suppressfulladdresses}{\@suppressfulladdressestrue\@needaddressesfalse}
%    \end{macrocode}
% \end{macro}
% \end{macro}
% \end{macro}
% \end{macro}
% \end{macro}
% \end{macro}

% \begin{macro}{\printdates}
% \begin{macro}{\noprintdates}
% Command to activate the printing out of the received/accepted
% dates. This can also be done using the \texttt{printdates} option,
% which is the default for \texttt{journal}, but not for \texttt{proceedings}.
%    \begin{macrocode}
\newcommand{\printdates}{\@printdatestrue}
\newcommand{\noprintdates}{\@printdatesfalse}
%    \end{macrocode}
% \end{macro}
% \end{macro}
% \begin{macro}{\@outputaddressesifnecessary}
% Command to output the addresses if we haven't already done so
% This is called by: 
% \begin{enumerate}
% \item The end of |thebibliography| environment
% \item The |\outputfulladdresses| command
% \end{enumerate}
%    \begin{macrocode}
\newcommand{\@outputaddressesifnecessary}{%
  \ClassInfo{rmaa}{Calling @outputaddressesifnecessary}
  \ifthenelse{\boolean{@needaddresses}}{%
    %% If in one column style, or in the first column of a two column
    %% style then we can try to put it on this page 
    \ifthenelse{\boolean{@firstcolumn} \or \not\boolean{@rmtwocol}}{%
      \@outputaddresses}{%
      %% otherwise give warning it will be on a separate page
      \ClassWarning{rmaa}{Full addresses will go on a separate
        page (please see User Guide for how to remedy this)}}%
    %% If we haven't managed yet, put it on its own page
    \ifthenelse{\boolean{@needaddresses}}{%
      \if@debug{\ClassWarning{rmaa}{In this silly bit!}}\fi
      \clearpage\vfill\@outputaddresses}{}%
    }{
    %% If we don't need them, we still might need to do the eoa hook 
    \ifthenelse{\boolean{@suppressfulladdresses}}{%
      \if@debug{\ClassWarning{rmaa}{XXX}}\fi
%%      \if@book\else\@endofarticlehook\null\fi
      }{}
    }%
  }  
%    \end{macrocode}
% \end{macro}

% \begin{macro}{\outputfulladdresses}
% A user interface to the above command, in case there is no biblio
% (11 May 2000) 
%    \begin{macrocode}
\newcommand{\outputfulladdresses}{\@outputaddressesifnecessary}
%    \end{macrocode}
% \end{macro}

% \begin{macro}{\adjustfinalcols}
% This command starts a new column and should be 
% used to ensure that the columns balance on the last page. It also
% outputs the full 
% author addresses to ensure that this occurs in the left hand
% column. 
% The |\par| command is so that the current line will be taken account
% of when working out |\pagetotal|
%    \begin{macrocode}
\newcommand{\adjustfinalcols}{\if@twocolumn\goodbreak\@outputaddresses\newpage\noindent\fi}
%    \end{macrocode}
% \end{macro}

% \begin{macro}{\@outputaddresses}
% Command to do the work of outputting the addresses. 
% This can be called by 
% \begin{enumerate}
% \item an explicit |\adjustfinalcols|
% \item via |\@outputaddressesifnecessary|, q.v.
% \end{enumerate}
%
%    \begin{macrocode}
\newcommand{\@outputaddresses}{%
  % Allow option of suppressing the full addresses, in which case
  % we just run the hook if we are not in a book style
  \if@suppressfulladdresses
%%     \if@book\else\@endofarticlehook\null\fi
  \else
  %% save the full addresses in a box for output later
  \sbox{\@fulladdbox}{\parbox{\textwidth}{\@fulladdresses}}
  % calculate height of address box  
  \settoheight{\@fulladdheight}{\usebox{\@fulladdbox}}
  \settodepth{\@fulladddepth}{\usebox{\@fulladdbox}}
  \addtolength{\@fulladdheight}{\@fulladddepth}
  \if@debug
  { \typeout{Address total height: \the\@fulladdheight}
    \typeout{Text height: \the\pagetotal}
    \typeout{Goal height: \the\pagegoal}
    }
  \fi
  % reuse \@fulladdheight to mean the total height of text+addresses
  \addtolength{\@fulladdheight}{\pagetotal}
%%  \addtolength{\@fulladdheight}{-\baselineskip}
  \if@debug\typeout{Adjusted text + addresses height: \the\@fulladdheight}\fi
  % check it will fit on current page
  \ifdim\@fulladdheight<\pagegoal
     %% must shift it to the left
     %% to compensate for the hanging indentation in the reference list. 
     %% NOTE: this is not foolproof!
     %% NB. This will give overfull hbox warnings, which is a bit ugly.
     \vfill\noindent\if@biblist\hspace*{-\leftmargin}\fi
     \usebox{\@fulladdbox}
     % Set this to show we've done it
     \@needaddressesfalse
     % and assume we are at the end of the article
     % so that we do the EOA hook if not in the book style
%%     \if@book\else\@endofarticlehook\fi
  \else
     % if not, just do nothing. With luck, this function will get 
     % called again on the next page
     \if@debug{\typeout{Holding over full addresses to next page}}\fi
  \fi\fi
  }
%    \end{macrocode}
% \end{macro}

% \subsubsection{Commands for editors}
% \label{sec:editor-commands}
% Some user/editor commands for doing the TOC and index in the book
% option. 
% They are defined here rather in the book-specific part since they
% need to be always available in order to avoid breaking documents. 
% \begin{macro}{\listofauthors}
% \begin{macro}{\@rawauthors}
% \begin{macro}{\rmaaTOCauthors}
% The command |\listofauthors| formats an author list for use in the
% TOC. It should contain the same names as in the |\author| command,
% but without the affiliation garbage and using |\&| instead of
% |and|. 
%    \begin{macrocode}
\newcommand{\@rawauthors}{Raw Authors}
\newcommand{\rmaaTOCauthors}{\noexpand\rmaaTOCAuthorStyle{Author list for TOC}}
\newcommand{\listofauthors}[1]{%
  \renewcommand{\@rawauthors}{#1}
  \renewcommand\rmaaTOCauthors{\noexpand\rmaaTOCAuthorStyle{#1}}}
%    \end{macrocode}
% \begin{macro}{\indexauthor}\label{indexauthor}
% The command |\indexauthor| produces an index entry for each
% author. 
% This command must come \emph{after} the
% |\title| and the |\listofauthors| commands. Also, there should be one 
% |\indexauthor| command for \emph{each} author. We also add the
% author onto the end of the |\ADS@authors| macro and construct the
% |\ADS@authorinitial| by a most cunning method (see below), but only
% for the first author. 
% 
% No, No, NO! This was all wrong. Backed out completely: 
% \begin{quote}
% Completely rewritten (19 Dec 2002) since we now use the index class, so 
% as to be able to alphabetize accented characters. A side effect of this is
% that |\index| no longer expands out the argument after the |@|, so we 
% instead add all the |\index| commands to a list (|\indexauthor@deferlist|), 
% which is executed in the |\@begofarticlehook|. 
% 
% Note the stuff with |\toks0| is copied from Exercise 20.15 of the \TeX{}book---it is
% to make sure the latin1-accented name gets passed into the idx
% file. 
% \end{quote}
% We now go back to the old way of doing things, using an optional
% argument to the \cs{indexauthor} command that has the name without
% accents. 
% \changes{3.08}{2003/01/16}{Remerged changes from 3.06 that got
% dropped from 3.07 after a synchronization disaster. Basically, the
% accentless optional argument to \cs{indexauthor}. Must add the macro
% directories to my unison set-up to stop this sort of thing from happening.}
%    \begin{macrocode}
\newcommand\@isortkey{}
\newcommand\@emptykey{empty}
\newcommand\indexauthor[2][empty]{
  \renewcommand\@isortkey{#1}
  \ifx\@isortkey\@emptykey \renewcommand\@isortkey{#2}\fi
  \index{\@isortkey @{{\rmaaIndexAuthorStyle{#2}}\hskip1em plus1em minus0.4em
      \rmaaIndexTitle\hskip1em plus2em minus0.4em\penalty-100
      \rmaaTOCauthors}}
  \indexauthorADS{#2}
}
%%\def\indexauthor@deferlist{\relax}
%%\long\def\indexauthor#1{
%%  \def\@@name{#1} \toks0=\expandafter{\@@name}
%%  \toks1=\expandafter{\indexauthor@deferlist}
%%  \protected@xdef\indexauthor@deferlist{%
%%    \the\toks1
%%    \noexpand\index{\the\toks0@{\noexpand\rmaaIndexAuthorStyle{\the\toks0}\noexpand\hskip1em plus1em minus0.4em 
%%        \rmaaIndexTitle\noexpand\hskip1em plus2em minus0.4em\noexpand\penalty-100
%%        \rmaaTOCauthors}}
%%  }
%%  \indexauthorADS{#1}
%%  }
%%\xdef\indexauthor@deferlist{}
\newcommand\indexauthorADS[1]{%
  \if@firstauthor\protected@xdef\ADS@authorinitial{\firstletter{#1}}\fi
  \protected@xdef\ADS@authors{\ADS@authors\if@firstauthor\else ; \fi
    #1}
  \@firstauthorfalse
}
%    \end{macrocode}
% Note that the first |#1| sets the key for the index, making sure it
% actually starts with  
% first letter of the surname. Meanwhile, the |@{...}| expression,
% which starts with some  
% control sequences is what is written to the |.ind| file. 
% See the makeindex man page for details. 
% \end{macro}
% \end{macro}
% \end{macro}
% \end{macro}

% \begin{macro}{\KillAllAccents}
% \begin{macro}{\firstletter}
% These are utility functions used in the above. |lett| is a counter
% indicating the current letter in the string. |\firstletter| just
% prints the first letter of its argument.
%    \begin{macrocode}
\newcommand{\KillAllAccents}{\def\'{}\def\`{}\def\~{}\def\.{}\def\i{i}}
\newcommand\firstletter[1]{\@car #1\@nil}
%    \end{macrocode}
% \end{macro}
% \end{macro}
% \begin{macro}{\SetVolume}
% \begin{macro}{\SetMSnumber}
% \begin{macro}{\SetFirstPage}
% \begin{macro}{\SetYear}
% \begin{macro}{\ReceivedDate}
% \begin{macro}{\AcceptedDate}
% These are for the editors to set the volume number, etc. 
%    \begin{macrocode}
\newcommand{\SetVolume}[1]{%
  \renewcommand{\@volume}{#1}
  \edef\ADS@volume{\PadToFourWithDots{#1}}
  \renewcommand{\@volumen}{\rmaa@VolumenStyle{volumen #1}}%
  }
\newcommand{\SetMSnumber}[1]{\renewcommand{\@MSnum}{#1}}
%    \end{macrocode}
% Note that there is no real reason to use |\SetFirstPage| since all
% the page numbering is done automatically in the book style. 
%    \begin{macrocode}
\newcommand{\SetFirstPage}[1]{\setcounter{page}{#1}}
\newcommand{\SetYear}[1]{%
  \renewcommand{\@year}{#1}
  \renewcommand{\ADS@year}{#1}
  \renewcommand{\@anyo}{\rmaa@VolumenStyle{#1}}
  }
\newcommand{\ReceivedDate}[1]{\renewcommand{\@received}{#1}}
\newcommand{\AcceptedDate}[1]{\renewcommand{\@accepted}{#1}}
%    \end{macrocode}
%
% \begin{macro}{\DOI}
% User/editor command to enter the DOI, which will be typeset as a link.
% \changes{4.0}{2019/11/02}{New user command \cs{DOI} to typeset
% DOI in the header of first page and add it to ADS info.}
%    \begin{macrocode}
\newcommand\DOI[1]{%
  \def\@DOI@url{\url{https://doi.org/10.22201/ia.01851101p.#1}}%
  \def\ADS@doi{https://doi.org/10.22201/ia.01851101p.#1}%
}%
%    \end{macrocode}
% \end{macro}
%
% \begin{macro}{\ORCIDiD}
% \changes{4.0}{2019/11/12}{Add mechanism for giving ORCID iD links}
% \begin{macro}{\ORCIDiDfull}
% \begin{macro}{\@orcid@icon}
% \begin{macro}{\@orcid@iconbig}
% Commands to add icon links to ORCID iDs of authors. For example
% \begin{quote}
% |\ORCIDiD|\marg{0000-0000-0000-0000}
% \end{quote}
% to be used after each author name in the |\author|
% command and in the |\fulladdresses|.  Alternatively,
% |\OORCIDiDfull| will typeset the icon plus the full text of the
% link, but we don't use that at present.  The icon image is stored in |ORCIDiD_iconvector.pdf|, which needs to be on the input path.  In principle, we could make EPS versions if necessary.
%    \begin{macrocode}
\newcommand\@orcid@icon{\includegraphics[height=1.8ex]{ORCIDiD_iconvector}}
\newcommand\@orcid@iconbig{\includegraphics[height=2ex]{ORCIDiD_iconvector}}
\newcommand\ORCIDiD[1]{\href{https://orcid.org/#1}{\,\@orcid@icon}}
\newcommand\ORCIDiDfull[1]{%
  \href{https://orcid.org/#1}{\@orcid@iconbig}\,\url{https://orcid.org/#1}}
%    \end{macrocode}
% \end{macro}
% \end{macro}
% \end{macro}
% \end{macro}
%
% \begin{macro}{\@volume}
% \begin{macro}{\@volumen}
% \begin{macro}{\@artpages}
% \begin{macro}{\@year}
% \begin{macro}{\@anyo}
% \begin{macro}{\@received}
% \begin{macro}{\@accepted}
% \begin{macro}{\@MSnum}
% Default values for the containers used by above commands, plus a
% few more useful ones. 
%    \begin{macrocode}
\newcommand{\@volume}{00}
\newcommand{\@volumen}{\rmaa@VolumenStyle{00}}
\newcommand{\@artpages}{%
  \oldpageref{\articlename:RMAAFirstPage}--\oldpageref{\articlename:RMAALastPage}}
\newcommand{\@year}{\the\year}
\newcommand{\@anyo}{\rmaa@VolumenStyle{2007}}
\newcommand{\@journalname}{\rmaa@JournalNameStyle{%
    \if@proceedings\RMAASClong\else\RMAAlong\fi}}
\newcommand{\@received}{}
\newcommand{\@accepted}{}
\newcommand{\@MSnum}{0000}
\newcommand{\@DOI@prefix}{https://doi.org/10.22201/ia.01851101p}
\def\@DOI@url{\url{https://doi.org/10.22201/ia.01851101p.2099.00.00.00}}
%    \end{macrocode}
% \end{macro}
% \end{macro}
% \end{macro}
% \end{macro}
% \end{macro}
% \end{macro}
% \end{macro}
% \end{macro}
% \end{macro}
% \end{macro}
% \end{macro}
% \end{macro}
% \end{macro}
% \end{macro}
% 
% The following commands are most useful in the \texttt{book}
% style. However, the |\SetConfTitle| command can now be used in
% individual articles to set the header for the \texttt{preprint}
% style. 
%    \begin{macrocode}
\newcommand{\@reunionplace}{Lugar de la reuni\'on}
\newcommand{\@reuniondate}{Fecha de la reuni\'on}
\newcommand{\@conftitle}{Conference Title}
\newcommand{\SetReunion}[2]{\renewcommand{\@reunionplace}{#1}\renewcommand{\@reuniondate}{#2}}
\newcommand{\SetConfTitle}[1]{\renewcommand{\@conftitle}{#1}}
\newcommand{\@editors}{Juan Ram\'on de la Fuente and ``El Mosh''}
\newcommand{\SetEditors}[1]{\renewcommand{\@editors}{#1}}
\newcommand{\@editores}{Juan Ram\'on de la Fuente y ``El Mosh''}
\newcommand{\SetEditores}[1]{\renewcommand{\@editores}{#1}}
\newcommand{\@mes}{mes de la publicaci\'on}
\newcommand{\SetMes}[1]{\renewcommand{\@mes}{\rmaa@VolumenStyle{#1}}}
%    \end{macrocode}

% \subsection{Notes for Author Proofs and Communication between Editors}
% \label{sec:notes}
% We introduce a command \cs{Note} that should be used to describe
% any change made to the manuscript. These are a bit tricky to get
% right. In general we use \cs{marginpar} but this only works in outer
% paragraph mode, so we just bung the box where we are if we are in a
% figure caption or something. If we are in the preamble, then we just
% use a \marginpar as part of the \maketitle command.
%    \begin{macrocode}
\newif\if@arenotes
\@arenotesfalse
\newcounter{notecounter}
\setcounter{notecounter}{0}
\newcommand\NotesList{}
\newcommand\Note[1]{%
  \if@proof{%
    \@arenotestrue%
    \stepcounter{notecounter}%
    \if@inpreamble\else{%
      \NoteMark{\arabic{notecounter}}%
    }\fi
    \protected@xdef\NotesList{\NotesList 
      \protect\item #1}%
  }\fi%
}
\newcommand\Changed[2]{Changed ``#1'' to ``#2''.}
\newcommand\NoteMark[1]{%
  \ifinner
  \raisebox{1ex}[0pt][0pt]{\makebox[0pt][l]{\framebox{\footnotesize\bfseries
        N#1}}}%
  \else
  \marginpar{\framebox{\footnotesize\bfseries N#1}}%
  \fi
}
\newcommand\RefNote[1]{\if@proof\NoteMark{\ref{#1}}\relax\fi}
\newcommand\NoteLabel{\protect\label}
\newcommand\OutputNotes{%
  \bgroup
  \onecolumn
  \section*{EDITORIAL NOTES}
  \begin{enumerate}
    \NotesList
  \end{enumerate}
  \egroup
}
\DeclareRobustCommand\cs[1]{\texttt{\char`\\#1}}

%    \end{macrocode}
% \subsection{Bibliographic Records for ADS}
% \label{sec:ADS}
% \changes{2.2beta11}{2001/11/27}{Just about finished all the ADS
% stuff now. Except for the references, which I have given up on for
% the time being.}
% \begin{macro}{\ADS@BibCode}
% \begin{macro}{\ADS@year}
% \begin{macro}{\ADS@journal}
% \begin{macro}{\ADS@volume}
% \begin{macro}{\ADS@firstpage}
% \begin{macro}{\ADS@authorinitial}
% We use the |\ADS@BibCode| command to hold the bibliographic code for
% ADS. It is 
% used in the |%R| 
% field of the bibliographic record for each article. 
% The format is 
% \begin{quote}
% \texttt{YYYYJJJJJVVVVMPPPPA}  
% \end{quote}
% where \texttt{YYYY} is the year, \texttt{JJJJJ} is the journal code,
% which is \texttt{RMxAA} or \texttt{RMxAC}, \texttt{VVVV} is the
% volume number (right-justified), \texttt{M} is not used by us,
% \texttt{PPPP} is the first page number (right-justified), and
% \texttt{A} is the initial of the first author's surname. The volume
% and page fields are padded with periods (`\texttt{.}'), which causes
% us some fun (this is done with the macro |\PadToFourWithDots| described
% below) . 
% 
% All these individual bits are held in the commands |\ADS@year|,
% |\ADS@journal|, |\ADS@volume|, |\ADS@firstpage|,
% |\ADS@authorinitial|. 
%    \begin{macrocode}
\newcommand{\SetADSBibCode}[1]{\renewcommand{\ADS@BibCode}{#1}}
\newcommand{\ADS@BibCode}
{\ADS@year\ADS@journal\ADS@volume.\PadToFourWithDots{\ADS@firstpage}\ADS@authorinitial}
\newcommand{\ADS@year}{2001} 
\newcommand{\ADS@journal}{\if@proceedings RMxAC\else RMxAA\fi}
\newcommand{\ADS@volume}{....} 
%    \end{macrocode}
% These two are just set to default values here, |\ADS@firstpage| is set
% in the |document| environment for included articles, in the
% |frontsection| environment, |\tableofcontents| command,  |index| and
% |abstracts| environments. Also by |\AtBeginDocument| for stand-alone
% articles. 
% 
% God knows how we are going to do |\ADS@authorinitial|!
%    \begin{macrocode}
\newcommand{\ADS@firstpage}{....}
\newcommand{\ADS@authorinitial}{A}
%    \end{macrocode}
% \end{macro}
% \end{macro}
% \end{macro}
% \end{macro}
% \end{macro}
% \end{macro}
% We also have a whole slew of other fields that can be set. These are
% all listed in Table~\ref{tab:ADS}. Most of them still have to be
% defined. 
% \begin{table}[!ht]\raggedleft
% \caption{ADS Bibliographic Records}\label{tab:ADS}
% \begin{tabular}{lp{8em}lp{8em}}\hline
% Code & Description & Macro Name & Comment \\ \hline
% |%R| & Bibliographic Code & |\ADS@BibCode| & required \\
% |%A| & Author list & |\ADS@authors| & required, generate from |\indexauthor| \\
% |%F| & Author affiliation & |\ADS@affil| & format: AA(\dots)
% AB(\dots) \dots \\
% |%J| & Journal Name & |\ADS@journalname| &  \\
% |%D| & Publication Date & |\ADS@fulldate| &  \\
% |%L| & Last Page & |\ADS@lastpage| & generate from label \\
% |%T| & Title & |\ADS@title| & required, generate from |\title| \\
% |%C| & Abstract Copyright & |\ADS@copyright| & set to IA-UNAM \\
% |%U| & URL for electronic document & |\ADS@url| & set to PDF file \\
% |%d| & DOI for article & |\ADS@doi| & set by |\DOI| \\
% |%K| & Keywords & |\ADS@keywords| & set equal to |\@keywords| \\
% |%A| & Abstract Text & |\ADS@abstract| & set equal to |\@abstract| \\
% |%R| & References & |\ADS@references| & generate from the
% bibliography  \\ \hline
% \end{tabular}
% \end{table}
% A lot of these macros are set by other commands, as mentioned in the
% Table, but some need to be set  explicitly here. 
% 
% \begin{macro}{\ADS@doi}
% \changes{4.0}{2019/11/08}{Save DOI for ADS}
% The user/editor command \cs{DOI} is used to set this. 
%    \begin{macrocode}
\newcommand{\ADS@doi}{https://doi.org/10.22201/ia.01851101p.0000.00.00.00}
%    \end{macrocode}
% \end{macro}
% 
% \begin{macro}{\ADS@journalname}
% This
% does not contain any accents at the moment so as not to cause
% problems---this maybe should be changed though. 
%    \begin{macrocode}
\newcommand{\ADS@journalname}{%
  \if@proceedings\@conftitle\space (Eds. \@editors) \fi%
  Revista Mexicana de Astronom\'\i{}a y Astrof\'\i{}sica%
  \if@proceedings\space (Serie de Conferencias)\fi%
  \space Vol. \@volume, pp. \ADS@firstpage--\ADS@lastpage\space (\@year)%
  \space (http://www.astroscu.unam.mx/\~{}rmaa/)
  }
%    \end{macrocode}
% \end{macro}
% \begin{macro}{\ADS@keywords}
% \begin{macro}{\ADS@abstract}
% \begin{macro}{\ADS@title}
% |\ADS@abstract| is straightforward, since we already have the text
% saved as |\@abstract|. We could do the same with |\ADS@keywords| but
% to avoid unwanted formatting we explicitly set it in |\keywords|
% (p.~\pageref{keywords}). |\ADS@title| is just set to |\@rawtitle|,
% which we were already using for another purpose. 
%    \begin{macrocode}
\let\ADS@title=\@rawtitle
\newcommand{\ADS@keywords}{No keywords specified}
\newcommand{\ADS@abstract}{No abstract specified}
%    \end{macrocode}
% \end{macro}\end{macro}\end{macro}
% \begin{macro}{\ADS@authors}
% |\ADS@authors| wants to be in ``Other, A. N.'' format, but with
% semi-colons between authors. The most natural way to do this seems to
% be with |\indexauthor| (p.~\pageref{indexauthor}). We append each
% author to this macro, which means we must explicitly clear it for
% each new article (done in |\article|). .   
%    \begin{macrocode}
\newcommand{\ADS@authors}{}
%    \end{macrocode}
% \end{macro}
% \begin{macro}{\ADS@fulldate}
% \begin{macro}{\ADS@lastpage}
% \begin{macro}{\ADS@copyright}
% \begin{macro}{\ADS@website}
% \begin{macro}{\ADS@url}
% \begin{macro}{\SetWebsite}
% \begin{macro}{\SetMonthNumeric}
% Other sundry data, together with user commands for bits that can be
% set by the Editors. The |\ADS@url| is constructed from the website
% (currently |http://www.astrosmo.unam.mx/rmaa/|,
% the journal descriptor (e.g., |RMxAC..10/|), followed by the
% path to the PDF file (e.g., |PDF/RMxAC..10_HEW1.pdf|). 
% 
%    \begin{macrocode}
\newcommand{\ADS@website}{http://www.astrosmo.unam.mx/rmaa/}
\newcommand{\ADS@url}
{\ADS@website\ADS@journal\ADS@volume\ADS@number/PDF/\ADS@journal\ADS@volume\ADS@number\_\articlename.pdf}
\newcommand{\ADS@number}{}
\newcommand{\ADS@fulldate}{00/\ADS@year} 
\newcommand{\ADS@lastpage}{00} 
\newcommand{\ADS@copyright}{(C) \ADS@year; Instituto de Astronom\'\i{}a,
  Universidad Nacional Aut\'onoma de M\'exico} 
\newcommand{\SetNumber}[1]{\renewcommand{\ADS@number}{-#1}}
\newcommand{\SetWebsite}[1]{\renewcommand{\ADS@website}{#1}}
\newcommand{\SetMonthNumeric}[1]{\renewcommand{\ADS@fulldate}{#1/\ADS@year}}
%    \end{macrocode}
% \end{macro}\end{macro}\end{macro}\end{macro}\end{macro}\end{macro}\end{macro}
% \begin{macro}{\ADS@affil}
% \begin{macro}{\ADS@references}
% These two are somewhat more problematical. Affiliation currently
% only works for the first author (see \dots) pending redesign of the
% user interface. References are now done by wrapping the
% |thebibliography| environment to open another input stream on the
% same file and search through for the |\bibitems|. 
%    \begin{macrocode}
\newcommand{\ADS@references}{}
\newcommand{\ADS@affil}{}
%    \end{macrocode}
% \end{macro}\end{macro}
% 
% Define two macros that expand to the strings
% |\begin{thebibliography}| and |\end{thebibliography}|. Note that the
% idea for this and much else is inspired by the code for the docstrip
% package. Also define |\StartBibX| in case authors use the more
% standard |\begin{thebibliography}{}| syntax. 
%    \begin{macrocode}
\edef\EndBib{\expandafter\@gobble\string\\end\string{thebibliography\string}}
\edef\StartBib{\expandafter\@gobble\string\\begin\string{thebibliography\string}}
\edef\StartBibX{\expandafter\@gobble\string\\begin\string{thebibliography\string}\string{\string}}
%    \end{macrocode}
% Define two new logical variables. |\ifAnotherline| is true if we
% want to read another line from the input stream. It is set to false
% when we encounter the |\end{thebibliography}|. |\ifInBib| is set to
% true when we frind the |\begin{thebibliography}| and tells us to
% start saving the following lines. 
%    \begin{macrocode}
\newif\ifAnotherLine
\newif\ifInBib
%    \end{macrocode}
% \begin{macro}{makeOther}
% We need a macro to change the catcodes of lots of things to other. 
%    \begin{macrocode}
\def\makeOther#1{\catcode`#1=12\relax}
%    \end{macrocode}
% \end{macro}
% \begin{macro}{\RereadBiblio}
% Define a macro to do all the fiddly stuff. This is used at the
% beginning of |thebibliography| environment (Sec~\ref{sec:biblio}).
%    \begin{macrocode}
\newread\tf@biblio
\newcommand{\RereadBiblio}{%
%    \end{macrocode}
% Open the current file again for reading
%    \begin{macrocode}
  \begingroup
  \immediate\openin\tf@biblio=\articlename
%    \end{macrocode}
% Everything in a group to be safe. Initialize logicals. Black magic
% to change
% catcodes and avoid blank lines being interpreted as |\par|. 
%    \begin{macrocode}
  \AnotherLinetrue\InBibfalse
  \makeOther\ \makeOther\\\makeOther\$%
  \makeOther\#\makeOther\^\makeOther\^^K%
  \makeOther\_\makeOther\^^A\makeOther\%%
  \makeOther\~\makeOther\{\makeOther\}\makeOther\&%
  \endlinechar-1\relax
%    \end{macrocode}
% Start looping through the file. Read each line to the macro
% |\Line|. If we have got to EOF or end of biblio then set
% |\AnotherLinefalse|.
%    \begin{macrocode}
  \loop
  \read\tf@biblio to \Line
  \ifx\Line\EndBib\AnotherLinefalse\InBibfalse\fi
  \ifeof\tf@biblio \AnotherLinefalse\fi
%    \end{macrocode}
% If we are in the biblio and the line is not blank then append it to
% the |\ADS@references| macro together with a newline. 
%    \begin{macrocode}
  \ifInBib
  \ifx\Line\empty\else\xdef\ADS@references{\ADS@references\Line^^J}\fi
  \fi
%    \end{macrocode}
% Check if the line we just saw was the start of biblio, then repeat
% the loop if necessary. 
%    \begin{macrocode}
  \ifx\Line\StartBib\InBibtrue\fi
  \ifx\Line\StartBibX\InBibtrue\fi
  \ifAnotherLine\repeat
  \immediate\closein\tf@biblio
  \endgroup
  }
%    \end{macrocode}
% \end{macro}
% 
% \begin{macro}{\PadToFourWithDots}
% Utility function for left-padding a number to 4 characters with
% periods. Used in 
% constructing |\ADS@BibCode|. This is a very simple algorithm and
% \emph{requires} that the argument actually be a number---if it is
% not then dire consequences will result.  
%    \begin{macrocode}
\newcommand{\PadToFourWithDots}[1]{%
  \ifnum #1 > 999 #1%
  \else\ifnum #1 > 99 .#1%
  \else\ifnum #1 > 9 ..#1%
  \else ...#1%
  \fi\fi\fi}
%%  \ifnum \count0 > 999 #1%
%%  \else\ifnum \count0 > 99 .#1%
%%  \else\ifnum \count0 > 9 ..#1%
%%  \else ...#1%
%%  \fi\fi\fi}
%    \end{macrocode}
% \end{macro}
%
% \begin{macro}{\AuthorADS}
% \begin{macro}{\ADS@affiliations}
% New mechanism for setting the affiliations for ADS. The author labels (AA, AB, \dots) are generated using two nested counters.  We use the |\robustify| command from \texttt{etoolbox} so that accented characters in the affiliation addresses do not blow up.  The list of accents may need to be expanded. 
%    \begin{macrocode}
\newcounter{AUTHOR}
\newcounter{author}[AUTHOR]
\stepcounter{AUTHOR}
\newcommand\ADS@affiliations{}
\robustify{\'}\robustify{\i}\robustify{\~}\robustify{\`}
\def\AuthorADS#1#2{%
  \begingroup
  \def\ORCIDiD##1{<ID system="ORCID">##1</ID>}
  \def\Email##1{<EMAIL>##1</EMAIL>}
  \def\Affil##1{##1}
  \stepcounter{author}
  \ifthenelse{\value{author} > 26}{\stepcounter{AUTHOR}\stepcounter{author}}{}
  \ifthenelse{\equal{\ADS@affiliations}{\@rmaaEmpty}}{%
    \xdef\ADS@affiliations{\Alph{AUTHOR}\Alph{author}(#2)}
  }{%
    \xdef\ADS@affiliations{\ADS@affiliations, \Alph{AUTHOR}\Alph{author}(#2)}
  }%
  \endgroup
}
%    \end{macrocode}
% \end{macro}
% \end{macro}
% %
% \begin{macro}{\ADS@dumpdata}
% \changes{4.0}{2019/11/13}{Write all author affiliations for ADS}
% \begin{macro}{\ADS@write}
% \begin{macro}{\ADS@openfile}
% |\ADS@dumpdata| is the macro that actually does the job of writing out all the
% data to an auxiliary file, using |\ADS@write| to write each
% field. The file is set up with |\ADS@openfile|, which currently is
% done in the TOC of a book.  The new multi-author affiliations are used if |\ADS@affiliations| has been set, otherwise it falls back to the single-author version |\ADS@affil|. 
%    \begin{macrocode}
\newcommand{\ADSfield}[2]{\relax}
\newcommand{\ADSvfield}[2]{\relax}
%%\renewcommand{\ADSfield}[12]{\%#1 #2}
\newcommand{\ADS@dumpdata}
{ 
  \ADS@write{\string\begin{ADSarticle}}
  \ADS@WRITE{\string\ADSfield{R}{\ADS@BibCode}}
  \ADS@write{\string\ADSfield{A}{\ADS@authors}}
  \ifthenelse{\equal{\ADS@affiliations}{\@rmaaEmpty}}{%
    {\ADS@sanitize\ADS@write{\string\ADSfield{F}{AA(\ADS@affil)}}}
  }{%
    {\ADS@sanitize\ADS@write{\string\ADSfield{F}{\ADS@affiliations}}}
  }
  \ADS@write{\string\ADSfield{J}{\ADS@journalname}}
  \ADS@write{\string\ADSfield{D}{\ADS@fulldate}}
  \ADS@write{\string\ADSfield{L}{\ADS@lastpage}}
  {\ADS@sanitize\ADS@write{\string\ADSfield{T}{\ADS@title}}}
  \ADS@write{\string\ADSfield{C}{\ADS@copyright}}
  \ADS@write{\string\ADSfield{d}{\ADS@doi}}
  \ADS@write{\string\ADSfield{U}{\ADS@url}}
  \ADS@write{\string\ADSfield{K}{\ADS@keywords}}
  {\ADS@sanitize\ADS@write{\string\ADSfield{B}{\ADS@abstract}}}
  \ADS@write{\string\ADSfield{Z}{\ADS@references}}
  \ADS@write{\string\end{ADSarticle}}
  }
%%\newcommand{\ADS@write}[1]{\if@book\addtocontents{ads}{#1}\fi}
%%\newcommand{\ADS@openfile}{\@starttoc{ads}}
\newcommand{\ADS@write}[1]{\set@display@protect\immediate\write\tf@ads{#1}}
\newcommand{\ADS@WRITE}[1]{\immediate\write\tf@ads{#1}}
\newcommand{\ADS@openfile}{%
  \newwrite\tf@ads%
  \immediate\openout\tf@ads\jobname.ads\relax
  \ADS@WRITE{\ADS@preamble}}
\newcommand{\ADS@closefile}{\ADS@WRITE{\ADS@postamble}}
\def\ADS@preamble{%
  %%  \string\documentclass{article}
  %%  \string\begin{document}
  }
\def\ADS@postamble{%
  %%  \string\end{document}
  \string%\string% Local Variables:^J 
  \string%\string% mode: LaTeX^J
  \string%\string% mode: Font-lock^J
  \string%\string% End:^J
  }
%    \end{macrocode}
% \end{macro}
% \end{macro}
% \end{macro}
% 
% \subsection{Automatic calculation of the page range}
% \label{sec:page-range}
% Another issue that has cost me a lot of grief. Still, it finally
% seems to be sorted now. All hail Jeff Goldberg. 
%
% Cunning trick to get the page numbers from |\pageref|.  Put a label
% on the first and last page of the document using |\AtBeginDocument|
% and |\AtEndDocument|. .  This way won't work for the book option.
% In that case, we do it by defining |\@endofarticlehook| and
% |\@begofarticlehook| as parallels to |\AtEndDocument| and
% |\AtBeginDocument|.  Note that we have to defer adding
% |\@begofarticlehook| to |\AtBeginDocument| to make sure it comes
% after the redefinitions of |\ref| and |\label|.
%    \begin{macrocode}
\if@book\AtEndDocument{\ADS@closefile}\else
 %\AtBeginDocument{\@begofarticlehook}
\AtEndDocument{\@endofarticlehook}
\fi
%    \end{macrocode}
% \begin{macro}{\@begofarticlehook}
% Since we are getting more and more things that need to be done at
% the beginning of every article, frontsection, etc we define a sister
% hook to 
% |\@endofarticlehook|. This hook does the following: 
% \begin{enumerate}
% \item Labels the first page
% \item Sets |\ADS@firstpage|
% \item runs and then clears |\indexauthor@deferlist|
% \item resets the keyword counter
% \end{enumerate}
% It is executed in the following circumstances: 
% \begin{enumerate}
% \item At |\begin{document}| for standalone files
% \item At |\begin{document}| for each included file
% \item At beginning of front sections 
% \item At beginning of TOC
% \item At beginning of Author Index
% \item At beginning of abstracts
% \end{enumerate}
% For some reason, this is not always within the right group, so we
% use |\xdef| just to be sure. 
%    \begin{macrocode}
\newcommand{\@begofarticlehook}{%
  \if@frontmatter\frontlabel{RMAAFirstPage}\else\label{RMAAFirstPage}\fi%
  \xdef\ADS@firstpage{\@arabic{\c@page}}%
  \if@book\else\makeindex\fi % needed so default index is defined in stand-alone articles
%%  \indexauthor@deferlist
%%  \xdef\indexauthor@deferlist{}
  \@inpreamblefalse
  \setcounter{kwdcounter}{0}
  }
%    \end{macrocode}
% \end{macro}
% \begin{macro}{\@endofarticlehook}
% This is a hook to be run at the end of each article
% currently, it just labels the last page and writes out all the ADS
% data to an auxiliary file.  
% It should not produce any printed output. 
% It is modelled on the |\AtEndDocument| command, to which it reduces
% in the case of a single article. 
%    \begin{macrocode}
\newcommand{\@endofarticlehook}{%
  \if@debug{\ClassWarning{rmaa}{***Doing Hook***}}\fi
  \if@debug{\ClassWarning{rmaa}{ADSabstract is: \ADS@abstract}}\fi
  \clearpage\rmaa@lastpagelabel\if@book\ADS@dumpdata\fi
  \ifx\NotesList\@rmaaEmpty\else\OutputNotes\fi
  \if@debug{\ClassWarning{rmaa}{Notes list is: \NotesList}}\fi
}
%    \end{macrocode}
% \begin{macro}{\rmaa@lastpagelabel}
% New version of last page label command copied from Jeff Goldberg's
% lastpage.sty package - we can't just use |\label| here in case the
% output routine for the last page has already been called. Instead,
% we first explicitly clear the page, step back the page counter,
% explicitly write to the .aux file, then restep the page counter
% again.
%
% But this will conflict with the hyperref package, which does labels differently.  Compensate for that by using the five-argument form.
% 
% Also set the |\ADS@LastPage| macro to the appropriate number. 
%    \begin{macrocode}
\newcommand{\rmaa@lastpagelabel}{%
  \addtocounter{page}{-1}%
  \immediate\write\@auxout{\string
    \newlabel{\articlename:RMAALastPage}{{}{\thepage}{}{}{}}}%
  \edef\ADS@lastpage{\thepage}%
  \addtocounter{page}{1}}
%    \end{macrocode}
% \end{macro}
% \end{macro}

% \subsection{Deal with explicit change in number of columns}
% \label{sec:columns}
% Sometimes, when formatting a book, we need to switch between
% one-column articles and two-column articles. For this to work, we
% need to wrap the standard \LaTeX{} commands to set the |\@rmtwocol|
% flag appropriately and also reset the package default values of all
% scalable lengths. .  
%    \begin{macrocode}
\let\latexonecolumn=\onecolumn
\let\latextwocolumn=\twocolumn
\renewcommand{\onecolumn}{\@rmtwocolfalse\@package@default@lengths\latexonecolumn}
\renewcommand{\twocolumn}{\@rmtwocoltrue\@package@default@lengths\latextwocolumn}
%    \end{macrocode}
% 
% \subsection{The actual formatting of the title}
% \label{sec:title-format}

% Note, we have to redefine the |\maketitle| command to use the
% numbered footnotes. This means we needn't use the atbeginend stuff. 
%    \begin{macrocode}
\renewcommand{\maketitle}{\par
  \begingroup
%    \end{macrocode}
% Don't switch to marked footnotes. 
% This is consistent with existing style but I'm not sure I like it
%    \begin{macrocode}
  %%  \renewcommand\thefootnote{\fnsymbol{footnote}}%
%%  \def\@makefnmark{\hbox to\z@{$\m@th^{\@thefnmark}$\hss}}%
  \def\@makefnmark{\textsuperscript{\@thefnmark}\hspace{-0.2em}}%
  \long\def\@makefntext##1{\parindent 1em\noindent
    \hbox to1.8em{\hss$\m@th^{\@thefnmark}$}##1}%
  \if@rmtwocol
  { \if@debug{\ClassWarning{rmaa}{TWO COLUMN}}\fi
    \twocolumn[\@maketitle]%
    }
  \else
  { \if@debug{\ClassWarning{rmaa}{ONE COLUMN}}\fi
    \newpage
    \global\@topnum\z@   % Prevents figures from going at top of page.
    \if@manuscript \@MSmaketitle \else \@maketitle \fi
    }
  \fi
  \if@manuscript\else\thispagestyle{rmaafirstpage}\@thanks\fi
  \ifx\NotesList\@rmaaEmpty\else{%
    \NoteMark{1--\arabic{notecounter}}\vspace*{-2\baselineskip}%
  }\fi%
  \endgroup
%    \end{macrocode}
% Don't reset footnote counter (not sure whether this is good thing
% or not).
%    \begin{macrocode}
  %%  \setcounter{footnote}{0}%
  \let\thanks\relax\let\affil\relax\let\altaffil\relax
  \let\maketitle\relax\let\@maketitle\relax
  \gdef\@thanks{}\gdef\@author{}\gdef\@title{}}
%    \end{macrocode}

% \begin{macro}{\@kwhwidth}
% \begin{macro}{\@kwtwidth}
% Redefine \cs{@maketitle} to include the abstract, resumen and
% keywords. Note that in the manuscripty style we use
% \cs{@MSmaketitle} instead.  
% \changes{3.16}{2003/04/07}{
%  Added extra vertical space between authors and accepted date.}
% \changes{3.26}{2011/08/15}{Add extra label to review articles}
%    \begin{macrocode}
\newlength{\@kwhwidth}\newlength{\@kwtwidth}
\newcommand\reviewname{Review}
\renewcommand{\@maketitle}{%
  \newpage
  \null
  \if@reviewarticle
  \smallskip
  {\large\bfseries\itshape\reviewname}%
  \bigskip
  \else
  \vskip\@before@title@skip%
  \fi
  \begin{center}%
    {\@title \par}%
    \vskip\@before@author@skip%
    \parbox{\rmaa@AuthorWidthFrac\textwidth}{\centering
      \lineskip\@between@author@skip%
      \@author
      \par}
    %% this should be replaced with `Submitted to...' for preprints
    \if@printdates
    \vskip\@before@dates@skip%
    \if@preprint
    \rmaa@ReceivedStyle{Submitted to RMxAA \@received}
    \else
    \rmaa@ReceivedStyle{Received \@received; accepted \@accepted}
    \fi
    \fi
    \if@manuscript
    \rmaa@ReceivedStyle{Submitted to RMxAA \@received}
    \fi
    \if@guide
    \rmaa@ReceivedStyle{Version \rmaa@version, \rmaa@date}
    \fi
    \vskip\@before@abstract@skip%
    \if@abstractpresent
    \parbox{\rmaa@AbstractWidthFrac\textwidth}{%
      \setlength{\parindent}{0.9cm}%
      \centerline{\rmaa@AbstractHeaderStyle{Abstract}}\par%
      \vskip\@after@abstracthead@skip%
      \if@newstyle\noindent\fi\@abstract
      \vskip\@between@abstract@skip%
      \centerline{\rmaa@AbstractHeaderStyle{Resumen}}\par%
      \vskip\@after@abstracthead@skip%
      \if@newstyle\noindent\fi\@resumen
      \if@setresumen\else\@abstract\fi % append english abstract if
                                       % spanish one not set
      \if@guide\else            %do not use keywords for guide
      \vskip\@before@kwd@skip%
      % Format the Key Words with hanging indentation:
      % They are put in a \parbox with width equal to the line width
      % minus the width of the header. 
      { % Widths of keyword
        \settowidth{\@kwhwidth}{\@keywordheader~}   % header and
        \setlength{\@kwtwidth}{\linewidth}          % keyword text
        \addtolength{\@kwtwidth}{-\@kwhwidth}       %
        \noindent\@keywordheader~\parbox[t]{\@kwtwidth}{%
          \rmaa@KeywordStyle{\boldmath\uppercase{\@keywords}}}}%
      \fi
      }
      \fi
  \end{center}%
  \vspace*{\afterkwdskip}
  \par
  }
\newlength{\afterkwdskip}\setlength{\afterkwdskip}{\z@}
%    \end{macrocode}
% \end{macro}
% \end{macro}
%
%
% \begin{macro}{\reviewarticle}
% \begin{macro}{\reviewtoc}
% \begin{macro}{\reviewlof}
%   User command to specify a review article. Must be used in
%   preamble. Main effect is to place the text ``Review'' before the
%   article title.  Also add command for setting a table of contents
%   and a list of figures at the beginning of the article.
%    \begin{macrocode}
\newcommand{\reviewarticle}{\@reviewarticletrue}
\newcommand\reviewtoc{%
  \section*{CONTENTS}%
  \@starttoc{rtc}%
}
\newcommand\reviewlof{%
  \section*{LIST OF FIGURES}%
  \@starttoc{lof}%
}
%    \end{macrocode}
% \end{macro}
% \end{macro}
% \end{macro}
% \begin{macro}{\l@section}
%    \begin{macrocode}
\renewcommand*\l@section[2]{%
  \ifnum \c@tocdepth >\z@
    \addpenalty\@secpenalty
    \addvspace{0.8em \@plus\p@}%
    \setlength\@tempdima{1.5em}%
    \begingroup
      \parindent \z@ \rightskip \@pnumwidth
      \parfillskip -\@pnumwidth
      \leavevmode %\bfseries
      \advance\leftskip\@tempdima
      \hskip -\leftskip
      #1\nobreak\hfil \nobreak\hb@xt@\@pnumwidth{\hss #2}\par
    \endgroup
  \fi}
%    \end{macrocode}
% \end{macro}
%
% \begin{macro}{\@MSmaketitle}
% A special version of |\@maketitle| for the \texttt{manuscript}
% option. This can deal with a multi-page abstract/resumen. Not very
% well, mind. 
%    \begin{macrocode}
\newcommand{\@MSmaketitle}{%
  \if@debug\ClassInfo{rmaa}{Using special MS maketitle command}\fi
  \newpage
  {
    \thispagestyle{rmaafirstpage}
%     \setlength{\parindent}{0pt}
    \null\vskip 2ex
    \begin{center}
      \@title\par
      \@author\par
      \rmaa@ReceivedStyle{%
        \ifx\@received\@rmaaEmpty
          Draft version: \@date
        \else
          Received: \@received
          \ifx\@accepted\@rmaaEmpty
          \else
            ; Accepted: \@accepted
          \fi
        \fi}\par
    \end{center}
    \@thanks
    \centerline{\rmaa@AbstractHeaderStyle{Resumen}}
    \noindent\@resumen\par\vskip 2ex
    \centerline{\rmaa@AbstractHeaderStyle{Abstract}}
    \noindent\@abstract\par\vskip 2ex
    \noindent\@keywordheader~\rmaa@KeywordStyle{\@keywords}\par
%     \@keywordheader~\rmaa@KeywordStyle{\uppercase{\@keywords}}\par
  }%
  }
%    \end{macrocode}
% \end{macro}

% \subsection{The sectioning commands}
% \label{sec:section-commands}

% \begin{macro}{\section}
% \begin{macro}{\subsection}
% \begin{macro}{\subsubsection}
% \begin{macro}{\paragraph}
% \begin{macro}{\subparagraph}
% Redefine all the standard commands to use these lengths. 
%    \begin{macrocode}
\renewcommand\section{\@startsection {section}{1}{\z@}%
                                   {\abovesectionskip}%
                                   {\belowsectionskip}%
                                   %% how do we get upper case here
                                   {\centering\reset@font\normalsize}}
\renewcommand\subsection{\@startsection{subsection}{2}{\z@}%
                                     {\abovesectionskip}%
                                     {\belowsectionskip}%
                                     {\centering\reset@font\normalsize\itshape}}
\renewcommand\subsubsection{\@startsection{subsubsection}{3}{\z@}%
                                     {\abovesectionskip}%
                                     {\belowsectionskip}%
                                     {\centering\reset@font\normalsize\itshape}}
\renewcommand\paragraph{\@startsection{paragraph}{4}{\z@}%
                                    {\aboveparaskip}%
                                    {\belowparaskip}%
                                    {\reset@font\normalsize\bfseries}}
\renewcommand\subparagraph{\@startsection{subparagraph}{5}{\parindent}%
                                      {\aboveparaskip}%
                                      {\belowparaskip}%
                                      {\reset@font\normalsize\bfseries}}
%    \end{macrocode}
% \end{macro}
% \end{macro}
% \end{macro}
% \end{macro}
% \end{macro}

% \begin{macro}{secnumdepth}
% Make sure that we number the subsubsections---this is done by
% default in the |article| class but not the |book| class, which is
% why we must set it explicitly here.
%    \begin{macrocode}
\setcounter{secnumdepth}{3}
%    \end{macrocode}
% \end{macro}
% \changes{2.2beta10}{2001/11/07}{Fixed \cs{secnumdepth} for \cs{book}
% style.}

% \begin{macro}{\thesubsection}
% \begin{macro}{\thesubsubsection}
% For subsection and subsubsection, we want the numbers to be roman
% but the title to be italic. This seems the easiest way to go about it.
%    \begin{macrocode}
\renewcommand{\thesubsection}{{\upshape\thesection.\@arabic\c@subsection}}
\renewcommand{\thesubsubsection}{\thesubsection{\upshape.\@arabic\c@subsubsection}}
%    \end{macrocode}
% \end{macro}
% \end{macro}

% \begin{macro}{\@seccntformat}
% Add ``.'' to the number here so we can remove it from |\thesection|,
% etc. Reduce gap between number and title WJH 27 Jan 1995
%    \begin{macrocode}
\def\@seccntformat#1{\csname the#1\endcsname .\hskip 0.4em}
%    \end{macrocode}
% \end{macro}

% \begin{macro}{\@sect}
% I didn't really want to mess with the |\@sect| command, but it was
% the only way I could manage to force upper case. Other than the
% commented parts, it is identical to the standard version. 
%    \begin{macrocode}
\def\@sect#1#2#3#4#5#6[#7]#8{
  \ifnum #2>\c@secnumdepth
     \let\@svsec\@empty
  \else
     \refstepcounter{#1}%
     \let\@@protect\protect
     \def\protect{\noexpand\protect\noexpand}%
     \edef\@svsec{\@seccntformat{#1}}%
     \let\protect\@@protect
  \fi
%    \end{macrocode}
% This is the bit I changed (27 Jan 1995). Use upper case for
% |\section|, use mixed case for the rest. 
%    \begin{macrocode}
  \ifnum #2=1                   
     \def\sec@upcase##1{\uppercase{##1}}%
     \def\sec@shape{\upshape}%
  \else                         
     \def\sec@upcase##1{\relax{##1}}%
     \def\sec@shape{\itshape}%
  \fi
  \@tempskipa #5\relax
  \ifdim \@tempskipa>\z@
     \begingroup #6\relax
          \@hangfrom{\hskip #3\relax\@svsec}%
                    {\interlinepenalty \@M \sec@upcase{#8}\par}%
     \endgroup
     \csname #1mark\endcsname{#7}%
%    \end{macrocode}
% Only write TOC entries if doing the Guide style. 
%    \begin{macrocode}
     \if@guide 
        \addcontentsline{toc}{#1}{%
          \ifnum #2>\c@secnumdepth \else
             \protect\numberline{\csname the#1\endcsname}%
          \fi
          #7}%
     \fi
%    \end{macrocode}
% Also allow per-article TOC for review papers
%    \begin{macrocode}
     \if@reviewarticle
        \ifnum #2>2 \else
           \addcontentsline{rtc}{#1}{%
             \protect\numberline{\csname the#1\endcsname}%
             \sec@upcase{\sec@shape #7}}%
        \fi
     \fi
  \else
     \def\@svsechd{#6\hskip #3\relax  %% \relax added 2 May 90
       \@svsec #8\csname #1mark\endcsname{#7}%
       \if@guide
          \addcontentsline{toc}{#1}{%
            \ifnum #2>\c@secnumdepth \else
               \protect\numberline{\csname the#1\endcsname}%
            \fi
            #7}%
       \fi
       }%
  \fi
  \@xsect{#5}}
%    \end{macrocode}
% \end{macro}


% \subsection{Appendices}
% \label{sec:appendices}
% Make sure that appendices in individual articles work properly. We
% copy the definition from the article class (the book one doesn't do
% what we want).  
% \begin{environment}{appendices}
%   This uses an environment |appendices| inside of which each
%   |\section| will become a separate appendix.  Note that the
%   appendices must come before the references (otherwise the full
%   addresses mechanism will not work properly). We make the equation
%   numbering be subordinate to the appendix (equations A12, B6,
%   etc). Finally, we have a medskip followed by APPENDICES to
%   introduce them. \changes{3.17}{2002/04/08}{ Added appendices environment
%   that does the right thing.}
%    \begin{macrocode}
\newenvironment{appendices}{\par
  \begingroup
  \setcounter{section}{0}%
  \setcounter{subsection}{0}%
  \renewcommand\thesection{\@Alph\c@section}
  \renewcommand\theequation{\thesection\arabic{equation}}
  \medskip
  \section*{APPENDICES}}
{\endgroup}
%    \end{macrocode}
% \end{environment}
% \begin{environment}{appendix}
%   If there is only one appendix, then use |appendix| environment
%   instead.  \changes{4.0}{2019/11/12}{ Added
%   appendix environment for case where there is only one section.}
%    \begin{macrocode}
\renewenvironment{appendix}{\par
  \begingroup
  \setcounter{section}{0}%
  \setcounter{subsection}{0}%
  \renewcommand\thesection{\@Alph\c@section}
  \renewcommand\theequation{\thesection\arabic{equation}}
  \medskip
  \section*{APPENDIX}}
{\endgroup}
%    \end{macrocode}
% \end{environment}

% \subsection{Floating thingies}
% \label{sec:floats}   

%% 
%    \begin{macrocode}
\renewcommand{\thetable}{\@arabic\c@table}
\renewcommand{\thefigure}{\@arabic\c@figure}
\renewcommand{\figurename}{Fig.}
\renewcommand{\tablename}{TABLE}
%\renewcommand{\fnum@table}{\MakeUpperCase{\tablename}~\thetable}
%    \end{macrocode}
% For some reason, I have to write `TABLE' explicitly here
% If I use |\tablename|, it gets converted to mixed case---God knows
% why. 
% |\figurename| is also being ignored. 
%    \begin{macrocode}
\renewcommand{\fnum@table}{TABLE~\thetable}
\renewcommand{\fnum@figure}{Fig.~\thefigure}
%    \end{macrocode}

% Do captions in 9 pt face in the main journal style and in new-style
% anything. 
%    \begin{macrocode}
\long\def\@makefigcap#1#2{%
  \if@proceedings\else\small\fi
  \if@newstyle\small\fi 
  \vskip\abovecaptionskip
  \sbox\@tempboxa{#1. #2}%
  \ifdim \wd\@tempboxa >\hsize
    #1. #2 \par
  \else
    \global \@minipagefalse
    \hb@xt@\hsize{\hfil\box\@tempboxa\hfil}%
  \fi
  \vskip\belowcaptionskip}
%    \end{macrocode}

% Make sure table caption is upper case. 
%    \begin{macrocode}
\long\def\@maketabcap#1#2{%
  \hbox to\hsize{\hfil #1 \hfil}
  \vskip 6pt 
  \centering\MakeUpperCase{#2}\par
  \vskip 6pt}
%    \end{macrocode}

%    \begin{macrocode}
\newlength{\@rottabcapwidth}
\long\def\@makerottabcap#1#2{%
  \setbox\@tempboxa\hbox{\MakeUpperCase{#2}}%
  \ifdim\wd\@tempboxa >.8\vsize
     \setlength\@rottabcapwidth{.8\vsize}
  \else
     \setlength\@rottabcapwidth{\wd\@tempboxa}
  \fi
  \rotatebox{90}{%
    \begin{minipage}{\vsize}\centering
      \begin{minipage}{\@rottabcapwidth}\centering
        #1\\ \vskip 6pt \MakeUpperCase{#2}
      \end{minipage}%
    \end{minipage}
    }%
  \hskip 6pt
  }

\newenvironment{Sideways}
    {\begin{sideways}\begin{minipage}{\vsize}\centering}
    {\end{minipage}\end{sideways}}

\newenvironment{Sidewaystable}
{%
  \let\oldcaption=\caption%
  \def\caption{
    \def\@captype{table}\renewcommand{\@makecaption}{\@maketabcap}
    \oldcaption}
  \begin{sidewaystable}\centering}
  {\end{sidewaystable}\let\caption=\oldcaption}
\newenvironment{Sidewaystable*}
{%
  \let\oldcaption=\caption%
  \def\caption{
    \def\@captype{table}\renewcommand{\@makecaption}{\@maketabcap}
    \oldcaption}
  \begin{sidewaystable*}\centering}
  {\end{sidewaystable*}\let\caption=\oldcaption}
%    \end{macrocode}

% Not sure why we need to stick these lines in. 
%    \begin{macrocode}
\renewenvironment{figure}
                 {\renewcommand{\@makecaption}{\@makefigcap}\@float{figure}}
                 {\end@float}
\renewenvironment{figure*}
                 {\renewcommand{\@makecaption}{\@makefigcap}\@dblfloat{figure}}
                 {\end@dblfloat}
\renewenvironment{table}
                 {\renewcommand{\@makecaption}{\@maketabcap}\@float{table}}
                 {\end@float}
\renewenvironment{table*}
                 {\renewcommand{\@makecaption}{\@maketabcap}\@dblfloat{table}}
                 {\end@dblfloat}
%    \end{macrocode}

% Redefine |\@caption| so that we don't bother with the list of
% figures. This saves us having to worry about fragile commands in the
% captions. 
%    \begin{macrocode}
\long\def\@caption#1[#2]#3{%
  \par
%%  \addcontentsline{\csname ext@#1\endcsname}{#1}%
%%    {\protect\numberline{\csname the#1\endcsname}{\ignorespaces #2}}%
  \begingroup
    \@parboxrestore
    \if@minipage
      \@setminipage
    \fi
    \normalsize
    \@makecaption{\csname fnum@#1\endcsname}{\ignorespaces #3}\par
  \endgroup}
%    \end{macrocode}

% Make tables not quite so cramped. Allow |\hline\hline| to give a
% poor man's thick rule. 
%    \begin{macrocode}
\renewcommand{\arraystretch}{1.15}
\setlength{\doublerulesep}{0.4pt}
%    \end{macrocode}

% \begin{macro}{\tablecols}
% \begin{macro}{\@tablecols}
% \begin{macro}{\tabcaptioncontinued}
% |\tablecols| is a user command for setting the number of columns. It
% is used by commands that try and span the entire table (I think
% there must be a better way of acheiving
% this). |\tabcaptioncontinued| is 
% a user command for setting the caption of a continued table on the
% next page. 
%    \begin{macrocode}
\newcommand{\tablecols}[1]{\renewcommand{\@tablecols}{#1}}
\newcommand{\@tablecols}{1}
\newcommand{\tabcaptioncontinued}{%
  \multicolumn{\@tablecols}{c}{\addtocounter{table}{-1}TABLE~\thetable\ (CONTINUED)} \\ \\
  }
%    \end{macrocode}
% \end{macro}
% \end{macro}
% \end{macro}

% \subsubsection{Footnotes in tables}
% \label{sec:table-footnotes}
% \begin{macro}{\tablenotemark}
% \begin{macro}{\tablenote}
% \begin{macro}{\tablenoteskip}
% Old user commands |\tablenotemark{MARK}|,
% |\tablenote{MARK}{TEXT}|---now deprecated. 
%    \begin{macrocode}
\newcommand\tablenotemark[1]{$^\mathrm{#1}$} 
\newcommand\tablenote[2]{\if@proceedings\else\small\fi ${}^\mathrm{#1}$ #2} 
\newcommand\tablenoteskip{\\[0.5\baselineskip]} 
%    \end{macrocode}
% \end{macro}
% \end{macro}
% \end{macro}

% \begin{macro}{\tabnotemark}
% \begin{macro}{\tabnotetext}
% \begin{macro}{\tabnotetextTWO}
% Alternative way of doing footnotes in tables. 
%    \begin{macrocode}
\newcommand{\tabnotemark}[1]{\textsuperscript{\lowercase{#1}}}
%\newcommand{\tabnotetextTWO}[4]{\tabnotetext{#1}{#2}\tabnotetext{#3}{#4}}
\newcommand{\tabnotetextTWO}[4]{%
  \multicolumn{\@tablecols}{l}{%
    \parbox{0.5\tabnotewidth}{\small\textsuperscript{#1}#2}%
    \parbox{0.5\tabnotewidth}{\small\textsuperscript{#3}#4}%
    }%
  \\}
\newcommand{\tabnotetext}[2]{%
  \multicolumn{\@tablecols}{@{}p{\tabnotewidth}@{}}{\small\textsuperscript{#1}#2}\\}
\newlength{\tabnotewidth}
%    \end{macrocode}
% \end{macro}
% \end{macro}
% \end{macro}


% \begin{macro}{changemargin}
%    \begin{macrocode}
% Command to allow margins to be changed on the fly. Useful to allow
% tables and Figures to protrude into the margins. 
%
% \changes{3.24}{2007/09/14}{Added \texttt{changemargin} environment
% (copied from TeX FAQ).}
\newenvironment{changemargin}[2]{% 
  % From TeX FAQ http://www.tex.ac.uk/cgi-bin/texfaq2html?label=chngmargonfly
  \begin{list}{}{%
    \setlength{\topsep}{0pt}%
    \setlength{\leftmargin}{#1}%
    \setlength{\rightmargin}{#2}%
    \setlength{\listparindent}{\parindent}%
    \setlength{\itemindent}{\parindent}%
    \setlength{\parsep}{\parskip}%
  }%
  \item[]}{\end{list}}
%    \end{macrocode}
% \end{macro}

% \subsection{The bibliography and references}
% \label{sec:biblio}
% \begin{macro}{\beforebiblioskip}
% \begin{macro}{\afterbiblioskip}
% I'm not entirely sure why I introduced these two commands, since
% directly putting |\vspace| before and after
% |\begin{thebibliography}| seems to work well enough. 
%    \begin{macrocode}
\newlength{\beforebiblioskip}
\setlength{\beforebiblioskip}{0pt}
\newlength{\afterbiblioskip}
\setlength{\afterbiblioskip}{0pt}
%    \end{macrocode}
% \end{macro}
% \end{macro}
% \begin{macro}{thebibliography}
% The environment itself is done as starred section with extra
% vertical space after heading. This space is currently zero. 
% In the book style we do the cunning trick with rereading the file in
% order to get the references into the |\ADS@references| macro.
% 
% \changes{3.05}{2003/01/02}{Removed the \cs{beforebiblioskip} and
% \cs{afterbiblioskip} commands, since they cause unwanted vertical
% space in some instances.}  
%    \begin{macrocode}
\renewenvironment{thebibliography}{%
  \if@book\RereadBiblio\fi
%%  \vspace*{\beforebiblioskip}
  \section*{REFERENCES} \frenchspacing
%%  \vspace*{\afterbiblioskip}
  \small
%    \end{macrocode} 
% The references themselves are set as a list in order 
% to get the hanging indentation. They are done in a smaller font and
% with french spacing so we needn't worry about the distinction
% between sentence-ending periods and abbreviations. 
%    \begin{macrocode} 
  \list{\null}{\leftmargin 1.5em\labelwidth\z@\labelsep\z@\itemindent -1.5em
    \itemsep\z@\parsep\z@\usecounter{enumi}}
  \sloppy\clubpenalty4000\widowpenalty4000
  \sfcode`\.=1000\relax\@biblisttrue}
{\endlist\@biblistfalse\normalsize\@outputaddressesifnecessary}
%    \end{macrocode}
% \end{macro}
% \subsection{Uniquification of cross-references}
% \label{sec:uniquification}
% 
% \begin{macro}{\articlename}
% In order that the |\label{}| command in individual articles be
% guaranteed to give a unique anchor, we append the article name (held
% in the |\articlename| command) to each key. This requires
% redefinition of the following commands: |\cite|, |\nocite|, |\ref|,
% |\pageref|, |\label|. 
% \begin{macro}{\oldcite}
% \begin{macro}{\oldnocite}
% \begin{macro}{\oldref}
% \begin{macro}{\oldpageref}
% \begin{macro}{\oldlabel}
% The original versions of these commands are
% saved as |\oldcite|, |\oldnocite|, etc. 
%
% The uniquifying of cross-references is now always done, so 
% we don't get multiply-defined labels in the single article case
%    \begin{macrocode}
\newcommand{\articlename}{anon}
\AtBeginDocument{
  \let\oldcite=\cite 
  \let\oldnocite=\nocite 
  \let\oldref=\ref 
  \let\oldpageref=\pageref
  \let\oldlabel=\label
}
%    \end{macrocode}
% \end{macro}
% \end{macro}
% \end{macro}
% \end{macro}
% \end{macro}
% \end{macro}
%
% Uniquify all cross-references by prepending the article name.  This
% now has to go into |\AtBeginDocument| so it is working on the new
% versions that hyperref.sty has set up.  Then we need to put the
% |\@begofarticlehook| afterwards, so that the first page label is using
% the new definitions.  If amsmath is loaded, then we also need to patch
% the special label command that they use. \changes{4.0}{2019/11/18}{Fix labels for amsmath display equations, in case that package has been loaded.}
%    \begin{macrocode}
\AtBeginDocument{
  \renewcommand{\ref}[1]{\oldref{\articlename:#1}%
    \if@debug{\ClassInfo{rmaa}{Reference to \articlename:#1}}\fi}
  \renewcommand{\pageref}[1]{\oldpageref{\articlename:#1}%
    \if@debug{\ClassInfo{rmaa}{Page reference to \articlename:#1}}\fi}
  \renewcommand{\label}[1]{\oldlabel{\articlename:#1}%
    \if@debug{\ClassInfo{rmaa}{Label for \articlename:#1}}\fi}%
  \ifdef{\label@in@display}{%
    \let\oldlabel@in@display=\label@in@display
    \renewcommand{\label@in@display}[1]{\oldlabel@in@display{\articlename:#1}%
      \if@debug{\ClassInfo{rmaa}{AMSTeX Label for \articlename:#1}}\fi}%
  }{}
  \@begofarticlehook
}%
%    \end{macrocode}
% Citations are more complicated since they can contain a list of
% keys. These are really only here for use with the rmaacite package,
% which is utterly deprecated (02 Jan 2003) now that I have the natbib
% code working. The command \cs{cite} should \emph{not} be used with
% natbib: use \cs{citep}, \cs{citet}, etc.\@ instead. 
%    \begin{macrocode}
\renewcommand{\nocite}[1]{\@bsphack
  \@for\@citeb:=#1\do{%
    \edef\@citeb{\expandafter\@firstofone\@citeb:\articlename}%
    \if@filesw\immediate\write\@auxout{\string\citation{\@citeb}}\fi
    \@ifundefined{b@\@citeb:\articlename}{\G@refundefinedtrue
      \@latex@warning{Citation `\@citeb' undefined}}{}}%
  \@esphack}
\renewcommand{\cite}[1]{\oldcite{#1:\articlename}}
%   \end{macrocode}


% Horrible fiddling with the bibliography citation mechanism. 
% This has to be globally available in case rmaacite puts it in the
% \texttt{.aux} file. 
%    \begin{macrocode}
\def\lbibcite#1#2#3{\global\@namedef{b@#1}{#2}
                    \global\@namedef{by@#1}{#3}}
%    \end{macrocode}

% \begin{macro}{\oldbibitem}
% \begin{macro}{\@rmaaBibArg}
% \begin{macro}{\@rmaaSpace}
% \begin{macro}{rmaabib}
% Of course, |\bibitem| needs to get the uniquification treatment as
% well. While we are at it, we redefine it to use consecutive integers
% for the key if the user leaves it blank (using the counter
% |rmaabib|). This allows one to put 
% |\bibitem{}| for all your references without getting annoying
% messages about multiply-defined labels.
% 
% We also take the opportunity to insert the optional argument to
% \cs{bibitem} if it is missing in the form Author(DATE) so that natbib
% doesn't try to do numerical citations. Natbib already provides a
% hook \cs{@extra@b@citeb} to help uniquify the citation key, so we
% just plug into that. However, it puts it \emph{after} the original
% key, not before as I had had it. Easiest solution is to change the
% way that I do it. 02 Jan 2003. Note that the \cs{@extra@b@citeb}
% does not get written to the .aux  file in the \cs{citation} command
% but it does work. 
% \changes{3.14}{06 Feb 2003}{Reinstalled old definition of \cs{bibitem}
%  for the case where the nonatbib option is requested, so that the
%  rmaacite package will work. This means that you can't use rmaacite
%  and natbib in the same volume.}
%    \begin{macrocode}
\if@natbib\def\@extra@b@citeb{:\articlename}\fi
%%\else
\let\oldbibitem\bibitem
\newcommand{\@rmaaBibArg}{}
\newcommand{\@rmaaSpace}{ }
\newcounter{rmaabib}
\if@natbib
\renewcommand{\bibitem}[2][ABC(2003)]{%
  \renewcommand{\@rmaaBibArg}{#2}
  \ifx\@rmaaBibArg\@rmaaEmpty
     \addtocounter{rmaabib}{1}
     \oldbibitem[#1]{\thermaabib:\articlename}
  \else\ifx\@rmaaBibArg\@rmaaSpace
     \addtocounter{rmaabib}{1}
     \oldbibitem[#1]{\thermaabib:\articlename}
  \else 
     \oldbibitem[#1]{#2:\articlename}
  \fi\fi
  \renewcommand{\@rmaaBibArg}{}
  }
\else
\renewcommand{\bibitem}[1]{%
  \renewcommand{\@rmaaBibArg}{#1}
  \ifx\@rmaaBibArg\@rmaaEmpty
     \addtocounter{rmaabib}{1}
     \oldbibitem{\articlename:\thermaabib}
  \else\ifx\@rmaaBibArg\@rmaaSpace
     \addtocounter{rmaabib}{1}
     \oldbibitem{\articlename:\thermaabib}
  \else
     \oldbibitem{\articlename:#1}
  \fi\fi
  \renewcommand{\@rmaaBibArg}{}
  }
\fi
%%\fi
%    \end{macrocode}
% \end{macro}
% \end{macro}
% \end{macro}
% \end{macro}

% \begin{macro}{oldpcite}
% \begin{macro}{oldscite}
% Then, we need to repeat the above scheme, but for the |rmaacite|
% definition. 
%    \begin{macrocode}
\newcommand{\@redefinebibitem}{%
  \let\oldbibitem\bibitem
  \let\oldcite\cite
  \let\oldpcite\pcite
  \let\oldscite\scite
  \def\bibitem[##1<##2>]##3{%
    \renewcommand{\@rmaaBibArg}{##3}
    \ifx\@rmaaBibArg\@rmaaEmpty
       \addtocounter{rmaabib}{1}
       \oldbibitem[##1<##2>]{\thermaabib:\articlename}
    \else\ifx\@rmaaBibArg\@rmaaSpace
       \addtocounter{rmaabib}{1}
       \oldbibitem[##1<##2>]{\thermaabib:\articlename}
    \else 
       \oldbibitem[##1<##2>]{##3:\articlename}
    % trust the author here
%%       \oldbibitem[##1<##2>]{##3}
    \fi\fi
    }
  }
%    \end{macrocode}
% \end{macro}
% \end{macro}


% \subsection{HTML Table of Contents}
% \label{sec:hyper-TOC}
% \begin{macro}{\HyperTOCEntry}
% \begin{macro}{\HyperTOCpart}
% \begin{macro}{\HyperTOCmeeting}
% \begin{macro}{\HyperTOCeditors}
% \begin{macro}{\HyperTOCpreamble}
% These commands are used to write information to the \texttt{.htoc}
% file, which is used in generating a Web Page of the Table of
% Contents. Since they only ever appear preceded by |\protect| in the
% argument to an |\addtocontents| command,  they should never need to
% be expanded by the macros. However, we define them here just in case
% \dots
%    \begin{macrocode}
\newcommand{\HyperTOCEntry}[3]{\relax}
\newcommand{\HyperTOCpart}[1]{\relax}
\newcommand{\HyperTOCmeeting}[3]{\relax}
\newcommand{\HyperTOCeditors}[1]{\relax}
\newcommand{\HyperTOCpreamble}{\relax}
%    \end{macrocode}
% \end{macro}
% \end{macro}
% \end{macro}
% \end{macro}
% \end{macro}
%

% \subsection{Formatting the poster abstracts} 
% These can now be used without using the book option.
%    \begin{macrocode}

\newcommand{\@makeminititle}{%
  \newpage
  \null 
  \vskip 3em%
  \begin{center}%
    {\@title \par}%
  \end{center}
  \vskip 3em%
  }

\newif\if@abstractsads
\@abstractsadsfalse
\newcommand\abstractsADS{\@abstractsadstrue}

\if@book\RequirePackage{multicol}\fi
\newenvironment{abstracts}[1][Abstracts of Contributed Papers]{%
  %% First, make sure we start on an odd page (06 Jun 2000)
  \if@book\clearpage\thispagestyle{empty}\cleardoublepage\fi
  %%
  \renewcommand{\articlename}{ABSTRACTS}
  \title{#1}\shortauthor{Abstracts}\shorttitle{Abstracts}
  \listofauthors{}
  \if@book
  \addcontentsline{toc}{chapter}{\protect\@fiddleauthors{\@TOCtitle}{\rmaaTOCauthors}}%
  \addtocontents{htc}{\protect\HyperTOCEntry{\@rawtitle}{\@rawauthors}{\articlename}}
  \fi
  \pagestyle{rmaa} % this now seems necessary (06 Jun 2000)
  \thispagestyle{rmaafirstpage}
%%  \twocolumn[\@makeminititle]
  \renewcommand{\abstract}[1]{%
    \renewcommand\ADS@abstract{##1}%
    \renewcommand{\@abstract}{##1}\@writeabstract%
  }
  \renewcommand{\title}[2][]{%
    \gdef\ADS@authors{}%
    \gdef\ADS@affil{}%
    \gdef\ADS@references{}%
    \gdef\ADS@keywords{}%
    \renewcommand{\@rmaaArg}{##1}
    \ifx\@rmaaArg\@rmaaEmpty
       \renewcommand{\@TOCaftertitle}{}
    \else
       \renewcommand{\@TOCaftertitle}{ (\@rmaaArg)}
    \fi
    \renewcommand{\@title}{\rmaa@TOCTitleStyle{##2}}%
    \renewcommand{\@TOCtitle}{\rmaa@TOCTitleStyle{##2}\@TOCaftertitle}%
    \renewcommand{\rmaaIndexTitle}{\rmaa@IndexTitleStyle{##2\@TOCaftertitle.}}%
    \renewcommand\ADS@title{##2}%
    \xdef\ADS@firstpage{\@arabic{\c@page}}%
    }
  \renewcommand{\author}[1]{\let\@thanks\@empty\renewcommand\@author{##1} }
  \renewcommand{\footnotetext}[2][XX]{\par\noindent\textsuperscript{##1}~##2}
  \renewcommand{\footnoterule}{\kern-3\p@ \hrule \@width 3em \kern 2.6\p@}
  \if@book\@begofarticlehook\fi
  \if@book\begin{multicols}{2}[\@makeminititle]\else\@makeminititle\fi
  \begingroup
  \setcounter{footnote}{0}
  }% and at the end of the environment....
  {\endgroup\if@book\end{multicols}\listofabstracts\label{RMAALastPage}\clearpage\fi}
%% auxiliary commands used:
%% \@writeabstract does all the actual formatting 
\newcommand{\@writeabstract}{%
  \if@book%
  \addcontentsline{loa}{abstract}{\protect\@fiddleauthors{\@TOCtitle}{\rmaaTOCauthors}}%
  \fi%
  \begin{center}\@title\par\@author\end{center}
  \@abstract
  \par\noindent\footnoterule\nobreak\footnotesize\@thanks
  \par\noindent\hrulefill\goodbreak
  \vskip 1cm \@plus0.5cm
  %% write ADS entry
  \if@book\if@abstractsads\ADS@dumpdata\fi\fi%
  \endgroup % this pair is to make each abstract a self-contained group
  \begingroup
  \setcounter{footnote}{0}
  }

\if@book
%% List of abstracts (stuff written to .loa file)
\newcommand{\listofabstracts}{%
  \begin{multicols}{2}[%
    \vskip 1.5cm \vfill % leave at least 1.5 cm and go to bottom of page if possible
    \centerline{\MakeUpperCase{List of Abstracts}}
    \vskip 1.5cm \@plus0.5cm
    ][6cm] % this makes it start a new page if there is not at least 6cm of space
    \renewcommand{\@fiddleauthors}[2]{##1\hskip1em plus2em minus0.4em\penalty-100  ##2}
    \parindent\z@
%%    \parskip.5\p@ \@plus .7\p@\relax 
    \parskip3\p@ \@plus 1\p@ \@minus 1\p@\relax % allow quite a lot of
                                % leeway here to avoid breaks inside
                                % entries 
    \renewcommand{\\}{\relax} % temporarily turn off explicit line breaks just in case
    \renewcommand{\altaffilmark}[1]{\relax} % ditto with this
    \@starttoc{loa}
  \end{multicols}
  }
\newcommand*{\l@abstract}[2]{%
  \par\penalty-1000
  #1\hskip 1em\penalty0\null\nobreak\hfill#2}
\fi

%    \end{macrocode}
 
% \subsection{Code for typesetting an entire volume of conference
% proceedings} \label{sec:whole-volume}
%    \begin{macrocode}

\if@book
%% environment for including the individual articles
\newcommand{\@fiddleauthors}[2]{\relax}
\renewcommand{\theequation}{\@arabic\c@equation}
\newenvironment{articles}%
{ \begingroup % make it a group so that author definitions do not
              % spill over
  \gdef\ADS@authors{}
  \pagestyle{rmaa}
  \renewenvironment{document}{%
    \@begofarticlehook
    \addtocounter{chapter}{1}%
    \addcontentsline{toc}{chapter}{%
      \protect\@fiddleauthors{\@TOCtitle}{\rmaaTOCauthors}}%
    \addtocontents{htc}{\protect\HyperTOCEntry{\@rawtitle}{\@rawauthors}{\articlename}}
    }{\relax}
  \renewcommand{\documentclass}[2][foo]{\relax} 
  \renewcommand{\usepackage}[2][foo]{%
    \def\@packagename{##2}\def\@rmaacitename{rmaacite}%
    \ifx\@packagename\@rmaacitename
    \@rmcitetrue\ClassInfo{rmaa}{This paper uses rmaacite}
    \makeatletter%%
%% This is file `rmaacite.sty',
%% generated with the docstrip utility.
%%
%% The original source files were:
%%
%% rmaa.dtx  (with options: `rmaacite')
%% 
%% PLEASE DO NOT EDIT THIS FILE! It is generated automatically from
%% the file rmaa.dtx using the accompanying Makefile and rmaa.ins
%% Any hand-edited changes will be lost
%% 
%% A LaTeX document class for the RMxAA
%% 
%% Author: William Henney (w.henney@crya.unam.mx)
%% 
%% Latest version available from: http://www.crya.unam.mx/rmaa
%% 
%% Copyright (C) 1999-2010 Instituto de Astronomia, UNAM, Mexico
%% 
%% This program is free software; you can redistribute it and/or
%% modify it under the terms of the GNU General Public License
%% as published by the Free Software Foundation; either version 2
%% of the License, or (at your option) any later version.
%% 
%% This program is distributed in the hope that it will be useful,
%% but WITHOUT ANY WARRANTY; without even the implied warranty of
%% MERCHANTABILITY or FITNESS FOR A PARTICULAR PURPOSE.  See the
%% GNU General Public License for more details.
%% 
%% You should have received a copy of the GNU General Public License
%% along with this program; if not, write to the Free Software
%% Foundation, Inc., 59 Temple Place - Suite 330, Boston, MA  02111-1307, USA.






\def\bibitem[#1<#2>]#3{\par\hangindent=1.5em\hangafter=1
  \if@filesw
  {\def\protect##1{\string ##1\space}
    \immediate\write\@auxout{\string\lbibcite{#3}{#1}{#2}}}
  \fi}
\@redefinebibitem % what it says
\def\lbibcite#1#2#3{\global\@namedef{b@#1}{#2}
                    \global\@namedef{by@#1}{#3}}

\def\@bcite#1#2{({#1\if@tempswa , #2\fi})}
\def\@pcite#1#2{#1\if@tempswa , #2\fi}
\def\@brcite{\ (}
\def\@ebrcite{)}
\def\@fmttop#1{\@fmtcite[\csname #1\endcsname]}
\def\@citex[#1]#2{%
%%  \if@filesw\immediate\write\@auxout{\string\citation{#2}}\fi
%% WJH added 16 Jan 2002
  \@for\@citeb:=#2\do{%
    \edef\@citeb{\expandafter\@firstofone\@citeb:\articlename}%
    \if@filesw\immediate\write\@auxout{\string\citation{\@citeb}}\fi}%
%% to here
  \def\@citea{}%
  \@cite{\@for\@citeb:=#2\do
    {%
      \edef\@citeb{\expandafter\@firstofone\@citeb:\articlename}%
      \@citea\def\@citea{;\ }\@ifundefined
      {b@\@citeb}{{\bf ?}\@warning {Citation `\@citeb' on page \thepage \space undefined}}%
      {\csname b@\@citeb\endcsname\@citebro\csname by@\@citeb\endcsname\@citebrc}}}%
  {#1}}

\def\cite{\@ifnextchar [{\let\@citebro=\space
                          \let\@citebrc=\relax
                          \let\@cite=\@bcite\@tempswatrue\@citex}
                        {\let\@citebro=\space
                          \let\@citebrc=\relax
                          \let\@cite=\@bcite\@tempswafalse\@citex[]}}
\def\pcite{\@ifnextchar [{\let\@citebro=\space
                          \let\@citebrc=\relax
                          \let\@cite=\@pcite\@tempswatrue\@citex}
                        {\let\@citebro=\space
                          \let\@citebrc=\relax
                          \let\@cite=\@pcite\@tempswafalse\@citex[]}}
\def\scite{\@ifnextchar [{\let\@citebro=\@brcite
                          \let\@citebrc=\@ebrcite
                          \let\@cite=\@pcite\@tempswatrue\@citex}
                        {\let\@citebro=\@brcite
                          \let\@citebrc=\@ebrcite
                          \let\@cite=\@pcite\@tempswafalse\@citex[]}}

\renewenvironment{thebibliography}
  {\section*{REFERENCES}\vspace*{\afterbiblioskip}\small
    \if@book\RereadBiblio\fi
    \sloppy\clubpenalty4000\widowpenalty4000
    \sfcode`\.=1000\relax
    \setlength\parindent{\z@}\setlength\parskip{\z@}}
  {\relax\normalsize\@outputaddressesifnecessary}








\endinput
%%
%% End of file `rmaacite.sty'.
\makeatother % read it in here so that all changes
    \else                                         % are local to paper
    \if@debug\ClassWarning{rmaa}{Package \@packagename\space was requested by \articlename\space}\fi
    \fi} 
  %% Don't allow individual articles to mess with the pages, volume or year
  \renewcommand{\SetFirstPage}[1]{\relax}
  \renewcommand{\SetYear}[1]{\relax}
  \renewcommand{\SetVolume}[1]{\relax}
  \renewcommand{\SetConfTitle}[1]{\relax}
  %% Don't allow individual articles to change these
  \renewcommand{\nonstopmode}{\relax}
  \renewcommand{\scrollmode}{\relax}
  \renewcommand{\outputplates}{\relax}
  %% XXXX need to check if using rmaacite and act accordingly
  %% avoid multiple declarations of the same length
  \let\oldnewlength=\newlength
  \renewcommand{\newlength}[1]{%
    \edef\reserved@a{\expandafter\@gobble\string ##1}%
    \@ifundefined{\reserved@a}{%
      \oldnewlength{##1}}{%
      \if@debug{\ClassInfo{rmaa}{\reserved@a\space already defined: doing nothing}}\fi%
      \relax}}%
  \renewcommand\newcommand{\providecommand}
  }%
{\endgroup\onecolumn}
     

%    \end{macrocode}
% Command to be used in `articles' environment for including an
% article
%    \begin{macrocode}
\newcommand{\article}[2][]{%
  % optional first argument is extra text to put in brackets after title in TOC
  % e.g. ``Invited Talk'' or ``Poster''
  \renewcommand{\articlename}{#2}
  % reset loads of counters
  \setcounter{footnote}{0}
  \setcounter{figure}{0}
  \setcounter{table}{0}
  \setcounter{section}{0}
  \setcounter{equation}{0}
  \setcounter{rmaabib}{0}
  \setcounter{author}{0}
  \setcounter{AUTHOR}{1}
  \gdef\ADS@authors{}%
  \gdef\ADS@affil{}%
  \gdef\ADS@affiliations{}%
  \gdef\ADS@references{}%
  \@firstauthortrue
  \@firstaffiltrue
  \@firstaltaffiltrue
  \@firstkwdtrue
  \@rmcitefalse
  \if@suppressfulladdresses\else\@needaddressestrue\fi
  \renewcommand{\@rmaaArg}{#1}
  \ifx\@rmaaArg\@rmaaEmpty
     \renewcommand{\@TOCaftertitle}{}
  \else
     \renewcommand{\@TOCaftertitle}{ (#1)}
  \fi
%    \end{macrocode}
% In the main journal articles must start on an odd page. 
%    \begin{macrocode}
  \if@journal\clearpage\thispagestyle{empty}\cleardoublepage\fi
%% Now input the file, then run the end-of-article hook (reinstated 28
  %% Sep 2001)
  {\input{\@articledir/\articlename.tex}\@endofarticlehook}}
%    \end{macrocode}
% Command to separate articles into sessions. Write an entry in the
% TOC, makes sure we start on an odd page (added 06 Jun 2000), then do
% a title page if \cs{SessionTitlePages} was called. 
% \changes{3.12}{04 Feb 2003}{Added option to have title page for each
% session. Turn it on with \cs{SessionTitlePages}.}
%    \begin{macrocode}
\newcommand{\session}[1]{
  %% this does nothing except write an entry in the TOC
  \addcontentsline{toc}{part}{\rmaa@TOCPartStyle{#1}}
  \addtocontents{htc}{\protect\HyperTOCpart{#1}}
  %% and also makes sure we start on an odd page (added 06 Jun 2000)
  \clearpage\thispagestyle{empty}\cleardoublepage
  \stepcounter{part}
  \if@sessiontitlepages{%
    \onecolumn
    \thispagestyle{empty}
    \renewcommand{\articlename}{Section\thepart TitlePage}
    \label{RMAAFirstPage}
    \make@session@title{#1}
    \label{RMAALastPage}
    \clearpage\thispagestyle{empty}\cleardoublepage
  }\fi
  }
\newcommand{\make@session@title}[1]{%
  {%
    \setlength\parindent{0pt}
    \vspace*{\stretch{1}}
    \rule{\linewidth}{1mm}
    \begin{flushright}\Huge\bfseries\boldmath
      #1
    \end{flushright}
    \rule{\linewidth}{1mm}
    \vspace*{\stretch{2}}
  }
}
%    \end{macrocode}
% Command to turn on separate title pages for each session. 
%    \begin{macrocode}
\newcommand{\SessionTitlePages}{\@sessiontitlepagestrue}
\newcommand{\BoldSessionTitles}{%
  \renewcommand{\rmaa@TOCPartStyle}[1]{\textbf{\rmaa@TOCTitleStyle{##1}}}
  }
\newcommand{\ArticleDir}[1]{\renewcommand{\@articledir}{#1/}\graphicspath{{1/}}}
\newcommand{\@articledir}{.}
%\newcommand{\ps@rmaafront}{\ps@rmaa}% page style for front matter (same as ever, actually)
\newcommand{\ps@rmaafront}{% page style for front matter:  page # in each foot
  \renewcommand\@oddfoot{\hfil{\thepage}}%
  \renewcommand\@evenfoot{{\thepage}\hfil}%
  % shorttitle in both heads
  \renewcommand\@evenhead{\hfil\@shorttitle\hfil}%
  \renewcommand\@oddhead{\@evenhead}
  }
\newcommand{\ps@rmaafrontfirstpage}{% page style for first page of front matter
  % page # in each foot, on outside edge
  \renewcommand\@oddfoot{\hfil\thepage}%
  \renewcommand\@evenfoot{\thepage\hfil}%
  % nothing in the even heads
  \renewcommand\@evenhead{\relax}%
  \renewcommand\@oddhead{\relax}%
  }
%% slight modification to the standard \frontmatter and \mainmatter
%% commands, just to set the \pagestyle
\renewcommand\frontmatter{%
  \latexonecolumn
  \cleardoublepage
  \@frontmattertrue
  \@mainmatterfalse
  \pagenumbering{roman}
  \pagestyle{rmaafront}
  \addtocounter{page}{2} % leave space for the title page
  }
\renewcommand\mainmatter{%
  \pagestyle{empty}
  \if@rmtwocol\twocolumn\fi
  \cleardoublepage
  \@frontmatterfalse
  \@mainmattertrue
  \pagenumbering{arabic}
  \pagestyle{rmaa}
  }
\renewcommand\backmatter{%
  \clearpage\pagestyle{empty}\cleardoublepage
  \@mainmatterfalse
  %% put extra space in the TOC
  \addtocontents{toc}{\bigskip}
  }
%    \end{macrocode}
%
% \subsection{Front sections}\label{sec:front-sections}
% \begin{macro}{\notitle}
% \begin{macro}{\shortname}
% \begin{macro}{frontsection}
% The environment \texttt{frontsection} is for doing a section in the
% front matter (preface, list of participants, or whatever). It has 
% one optional argument:
% \begin{enumerate}
% \item[\#1] Short title (no embedded spaces) for
% use in making PS files. E.g., ``list\_of\_participants''.   
% \end{enumerate}
% and two mandatory arguments: 
% \begin{enumerate}
% \item[\#2] Title of section (e.g., `Preface')
% \item[\#3] Commands to be evaluated before processing the contents
% of the environment. This is mainly useful for putting in the things
% that would go in the preamble of a normal article  (e.g.,
% |\listofauthors|, |\indexauthor|, |\notitle|) 
% \end{enumerate}
%    \begin{macrocode}
\newcommand{\notitle}{\@titlevisiblefalse} \newcommand{\shortname}{}
\newenvironment{frontsection}[3][]{%
  \begingroup
  \renewcommand{\shortname}{#1}
  \ifx\shortname\@rmaaEmpty
     \renewcommand{\articlename}{#2}
  \else
     \renewcommand{\articlename}{\shortname}
  \fi
  \@begofarticlehook
  \@titlevisibletrue % default
  \thispagestyle{rmaafrontfirstpage}
  \listofauthors{} % blank out this in case it is not specified
  \title{{#2}}\shorttitle{#2}\shortauthor{#2}
  #3 % evaluate the command given as second mandatory arg
  \addcontentsline{toc}{chapter}{\protect\@fiddleauthors{\@TOCtitle}{\rmaaTOCauthors}}%
  \addtocontents{htc}{\protect\HyperTOCEntry{\@rawtitle}{\@rawauthors}{\articlename}}
  \if@titlevisible
     \vspace*{1cm}\centerline{\rmaa@FrontSecStyle{\@title}}\vspace*{1cm}
  \fi
  \setlength{\parskip}{3ex}
  }{\frontlabel{RMAALastPage}\clearpage\endgroup}
%    \end{macrocode}
% \end{macro}
% \end{macro}
% \end{macro}
%    \begin{macrocode}
%% Definition of the command \frontlabel, which rights out something like `F4' to
%% the .aux file, instead of `\textlatin  {i}'. This is so it is easy to post-process
%% the file for making PS files of individual sections
\newcommand{\frontlabel}[1]{%
  \@bsphack
  \protected@write\@auxout{}%
     {\string\newlabel{\articlename:#1}{{}{F\@arabic{\c@page}}}}
  \@esphack
  }
%% User command to say X, Y, Z,\\ Editors of the Proceedings
\newcommand{\EditorSignature}{\vspace*{\parskip}\rightline{\@editors}%
  \rightline{Editors of the proceedings}}
%% Fiddle with \thesection, etc. so that they don't include the chapter number
%% Basically go back to the `article' definitions of these
\renewcommand{\thesection}{\@arabic\c@section}
\renewcommand{\thefigure}{\@arabic\c@figure}
\renewcommand{\thetable}{\@arabic\c@table}

\newcommand{\@oldargs}{}
\newcommand{\@newargs}{}
%% add TOC entry for each paper
%\renewcommand{\title}[1]%
%{\renewcommand{\@title}{\rmaa@TitleStyle{#1}}%
%\addtocounter{chapter}{1}\addcontentsline{toc}{chapter}{#1}%
%}

%% Formatting of the table of contents
%% dimensions of various parts
\newlength{\TOCtextwidth}\setlength{\TOCtextwidth}{12.65cm}
\newlength{\TOCpagenumwidth}\setlength{\TOCpagenumwidth}{1.3cm}
\newlength{\TOCentrywidth}\setlength{\TOCentrywidth}{\TOCtextwidth}
\addtolength{\TOCentrywidth}{-\TOCpagenumwidth} 
\newlength{\TOCextramargin}\setlength{\TOCextramargin}{1.25cm}
\newcommand{\TOCbox}[1]{\hspace{\TOCextramargin}\parbox{\TOCtextwidth}{#1}\par}
\newlength{\TOCextramarginx}\setlength{\TOCextramarginx}{\TOCextramargin}\addtolength{\TOCextramarginx}{4.3cm}
\newlength{\TOCtextwidthx}\setlength{\TOCtextwidthx}{\TOCtextwidth}\addtolength{\TOCtextwidthx}{-4.3cm}
\newcommand{\TOCboxx}[1]{\hspace{\TOCextramarginx}\parbox{\TOCtextwidthx}{#1}\par}

\newcommand{\RMAAbannerimage}{rmsc}
\renewcommand\tableofcontents{%
  \renewcommand{\articlename}{CONTENTS}
  \begingroup
  \@begofarticlehook
  \if@twocolumn
     \@restonecoltrue\onecolumn
  \else
     \@restonecolfalse
  \fi
  %% 9pt for journal style
  \if@proceedings\else\small\fi
  \if@newstyle\small\fi %% same for new style
  \if@proceedings\thispagestyle{rmaafrontfirstpage}\else\thispagestyle{empty}\fi
  \shortauthor{Contents}\shorttitle{Contents}
  \vspace*{18mm}
  \if@printbanner
  %% put in the banner header (need to get better image of this)
     \TOCbox{\hspace*{-0.15cm}\includegraphics[width=0.9\TOCtextwidth]{\RMAAbannerimage}}
     \vspace*{8mm}
  %% volume and publication date (should really be in funny font)
     \TOCboxx{\@volumen\\[2mm] \@mes\ \@anyo}
  \else
  %% put in the banner header (need to get better image of this)
     \TOCbox{\hspace*{-0.15cm}\includegraphics[width=0.9\TOCtextwidth]{\RMAAbannerimage}}
     \vspace*{12.5mm}
  %% volume and publication date (should really be in funny font)
     \TOCboxx{~\\[2mm] ~}
  \fi
  \vspace*{16mm}
  %% This is the command that fiddles with the TOC entry so that the
  %% title is flushleft but the authors are flushright
  \renewcommand{\@fiddleauthors}[2]{%
    % first put the paper title
    \renewcommand{\\}{\relax} % temporarily turn off explicit line breaks just in case
    \renewcommand{\altaffilmark}[1]{\relax} % ditto with this
    \renewcommand{\textsuperscript}[1]{\relax} % ditto with this 
    ##1
    % now save the author list in a box
    \sbox\@tempboxa{##2}%
    % check that box is shorter than the total width available for the entry
    \ifdim \wd\@tempboxa <\TOCentrywidth
       % if so, then put it right here, flush against right margin
       \penalty10000\hskip1em plus10em\penalty0\hbox{}\nobreak\hfill\box\@tempboxa
       % Note that we force a space of at least 1em between title and authorlist, 
       % stretchable to 11em since we may have very short titles but still with an 
       % authorlist too long to go on the same line.  The only linebreak allowed is 
       % directly before the authorlist.
    \else
       % if it is too long, we set the original authorlist text in its own parbox
       \par\vskip1ex\parbox[b]{\TOCentrywidth}{\raggedleft ##2}
    \fi
    }
  \@starttoc{toc}%
  \if@restonecol\twocolumn\fi
  \frontlabel{RMAALastPage}
  \clearpage
  \thispagestyle{empty}
  \cleardoublepage % make sure we start on an odd page
  \endgroup

  %%% ALSO, do the HTML TOC
  %% The file that will be used to make the HTML TOC
  \@starttoc{htc}
  %% All the stuff in this file does nothing in the book style
  \addtocontents{htc}{\protect\HyperTOCmeeting{\@conftitle}{\@reunionplace}{\@reuniondate}}
  \addtocontents{htc}{\protect\HyperTOCeditors{\@editors}}
  \addtocontents{htc}{\protect\HyperTOCpreamble}
  \addtocontents{htc}{\protect\HyperTOCEntry{Table of Contents}{}{\articlename}}

  %% and the ADS file
  \ADS@openfile
  }
%% This is now much simpler
\renewcommand\l@chapter[2]{%TOC entry for articles
  \ifnum \c@tocdepth >\m@ne
    \vskip 1.0em \@plus\p@
    \TOCbox{%
      \parbox[b]{\TOCentrywidth}{#1}\parbox[b]{\TOCpagenumwidth}{\hfill#2}
    }
  \fi} 
\renewcommand\l@part[2]{%TOC entry for sessions: centered, no page number
  \ifnum \c@tocdepth >\m@ne
    \vskip 1.5em \@plus1.5em \goodbreak
    \TOCbox{%
      \parbox[b]{\TOCentrywidth}{\hfil#1\hfil}\hfill
    }
  \fi\nobreak} 

%% Formatting theindex
%% strongly encourage page breaks between index entries
\renewcommand{\@idxitem}{\par\penalty-1000\hangindent.45cm} 
\renewcommand{\indexspace}{\par\penalty-1000\vskip 10\p@ \@plus15\p@ \@minus3\p@\relax}
\renewenvironment{theindex}{%
  %% first, make sure we start on an odd page (06 Jun 2000)
  \clearpage\thispagestyle{empty}\cleardoublepage
  %% this looks a bit like an article but is two columns and without abstract, etc
  \renewcommand{\articlename}{INDEX}
  \title{Author Index}\shortauthor{Author Index}\shorttitle{Author Index}
  \listofauthors{}
  \addcontentsline{toc}{chapter}{\protect\@fiddleauthors{\@TOCtitle}{\rmaaTOCauthors}}%
  \addtocontents{htc}{\protect\HyperTOCEntry{\@rawtitle}{\@rawauthors}{\articlename}}
  \pagestyle{rmaa}
  \thispagestyle{rmaafirstpage}
  \@begofarticlehook
  %% 9pt for journal style
  \if@proceedings\else\small\fi
  \if@newstyle\small\fi
  \twocolumn[\@makeminititle]
  %% vertical spacing (see also definition of \@idxitem and \indexspace
  \parskip\z@ \@plus .6\p@\relax % allow quite a lot of leaway here to avoid breaks inside entries
  %% horizontal spacing, breaking and glue parameters (see also the stuff in the definition 
  %% of \indexauthor
  \parindent\z@ 
  \exhyphenpenalty10000 % don't break on explicit hyphens, it always looks bad
  \rightskip=0pt plus.2em % allow a tiny bit of leeway at the ends of lines
  \spaceskip.3333em plus.5em minus.15em % a lot of leeway between words
  \xspaceskip.5em plus.5em minus.4em % this seems to have no effect
  \let\item\@idxitem
  \renewcommand{\\}{\relax} % temporarily turn off explicit line breaks just in case
  \renewcommand{\altaffilmark}[1]{\relax} % ditto with this
  \renewcommand{\textsuperscript}[1]{\relax} % ditto with this
  }{\label{RMAALastPage}\clearpage} % clearpage necessary so that we know about \shortauthor



%% end definitions for book styles
\fi

%    \end{macrocode}
% In order to do away with plates, we sometimes want a figure to
% appear in color only in the online version. The following command
% can bve used in a figure caption to typeset a note to this effect. 
%    \begin{macrocode}
\newcommand\ColorOnlineOnly{%
  \textsc{Note: this figure is available in color in the electronic
    version of this article, obtainable from}
  \texttt{http://\discretionary{}{}{}www.\discretionary{}{}{}astroscu.\discretionary{}{}{}unam.\discretionary{}{}{}mx/\discretionary{}{}{}\~{}rmaa/}.}
%% Plates
%% Define 2 counters: plate is the number in the whole document, plateloop is used 
%% internally by \outputplates
\newcounter{plate}
\newcounter{plateloop}
%% Currently limited to 15 plates in document
%% when we do a label, we need to write to \newlabel entries to the .aux file, one 
%% that says RMAAPlateX so that it can be referred to in \ouputplates, and one that
%% says articlename:RMAAPlateX so that the dvips program can find it and asociate it
%% with the right article
\newcommand{\@platelabeli}{\oldlabel{RMAAPlate1}\label{RMAAPlate1}}
\newcommand{\@platelabelii}{\oldlabel{RMAAPlate2}\label{RMAAPlate2}}
\newcommand{\@platelabeliii}{\oldlabel{RMAAPlate3}\label{RMAAPlate3}}
\newcommand{\@platelabeliv}{\oldlabel{RMAAPlate4}\label{RMAAPlate4}}
\newcommand{\@platelabelv}{\oldlabel{RMAAPlate5}\label{RMAAPlate5}}
\newcommand{\@platelabelvi}{\oldlabel{RMAAPlate6}\label{RMAAPlate6}}
\newcommand{\@platelabelvii}{\oldlabel{RMAAPlate7}\label{RMAAPlate7}}
\newcommand{\@platelabelviii}{\oldlabel{RMAAPlate8}\label{RMAAPlate8}}
\newcommand{\@platelabelix}{\oldlabel{RMAAPlate9}\label{RMAAPlate9}}
\newcommand{\@platelabelx}{\oldlabel{RMAAPlate10}\label{RMAAPlate10}}
\newcommand{\@platelabelxi}{\oldlabel{RMAAPlate11}\label{RMAAPlate11}}
\newcommand{\@platelabelxii}{\oldlabel{RMAAPlate12}\label{RMAAPlate12}}
\newcommand{\@platelabelxiii}{\oldlabel{RMAAPlate13}\label{RMAAPlate13}}
\newcommand{\@platelabelxiv}{\oldlabel{RMAAPlate14}\label{RMAAPlate14}}
\newcommand{\@platelabelxv}{\oldlabel{RMAAPlate15}\label{RMAAPlate15}}
\newcommand{\@outputplatelabeli}{\oldlabel{RMAAOutputPlate1}}
\newcommand{\@outputplatelabelii}{\oldlabel{RMAAOutputPlate2}}
\newcommand{\@outputplatelabeliii}{\oldlabel{RMAAOutputPlate3}}
\newcommand{\@outputplatelabeliv}{\oldlabel{RMAAOutputPlate4}}
\newcommand{\@outputplatelabelv}{\oldlabel{RMAAOutputPlate5}}
\newcommand{\@outputplatelabelvi}{\oldlabel{RMAAOutputPlate6}}
\newcommand{\@outputplatelabelvii}{\oldlabel{RMAAOutputPlate7}}
\newcommand{\@outputplatelabelviii}{\oldlabel{RMAAOutputPlate8}}
\newcommand{\@outputplatelabelix}{\oldlabel{RMAAOutputPlate9}}
\newcommand{\@outputplatelabelx}{\oldlabel{RMAAOutputPlate10}}
\newcommand{\@outputplatelabelxi}{\oldlabel{RMAAOutputPlate11}}
\newcommand{\@outputplatelabelxii}{\oldlabel{RMAAOutputPlate12}}
\newcommand{\@outputplatelabelxiii}{\oldlabel{RMAAOutputPlate13}}
\newcommand{\@outputplatelabelxiv}{\oldlabel{RMAAOutputPlate14}}
\newcommand{\@outputplatelabelxv}{\oldlabel{RMAAOutputPlate15}}
\newcommand{\@platerefi}{\oldpageref{RMAAPlate1}}
\newcommand{\@platerefii}{\oldpageref{RMAAPlate2}}
\newcommand{\@platerefiii}{\oldpageref{RMAAPlate3}}
\newcommand{\@platerefiv}{\oldpageref{RMAAPlate4}}
\newcommand{\@platerefv}{\oldpageref{RMAAPlate5}}
\newcommand{\@platerefvi}{\oldpageref{RMAAPlate6}}
\newcommand{\@platerefvii}{\oldpageref{RMAAPlate7}}
\newcommand{\@platerefviii}{\oldpageref{RMAAPlate8}}
\newcommand{\@platerefix}{\oldpageref{RMAAPlate9}}
\newcommand{\@platerefx}{\oldpageref{RMAAPlate10}}
\newcommand{\@platerefxi}{\oldpageref{RMAAPlate11}}
\newcommand{\@platerefxii}{\oldpageref{RMAAPlate12}}
\newcommand{\@platerefxiii}{\oldpageref{RMAAPlate13}}
\newcommand{\@platerefxiv}{\oldpageref{RMAAPlate14}}
\newcommand{\@platerefxv}{\oldpageref{RMAAPlate15}}
\newcounter{figplatei}
\newcounter{figplateii}
\newcounter{figplateiii}
\newcounter{figplateiv}
\newcounter{figplatev}
\newcounter{figplatevi}
\newcounter{figplatevii}
\newcounter{figplateviii}
\newcounter{figplateix}
\newcounter{figplatex}
\newcounter{figplatexi}
\newcounter{figplatexii}
\newcounter{figplatexiii}
\newcounter{figplatexiv}
\newcounter{figplatexv}
%% Printed output should use the counter `plate'
\renewcommand{\theplate}{\@arabic\c@plate}
\renewcommand{\theplateloop}{\@arabic\c@plateloop}
%% Command to output all the saved plates in this paper
\newcommand{\outputplates}{%
  \setcounter{plateloop}{0}
  \if@book\else\@endofarticlehook\fi % not entirely sure why, but this
                                     % seems to be needed
  \clearpage % first output the last page of text
  % now set pagestyle for blank even pages and make sure we start on an odd page
  \pagestyle{empty}
  \cleardoublepage
  \onecolumn
  \parindent\z@
  % first put in an extra sheet
  % the main purpose of this is so there is always a blank even page before 
  % each plate for filling in the PS file if the article has an odd number
  % of pages (see splitbook.sh)
  \null\vfil\centerline{PLATES}\vfil\clearpage\pagestyle{rmaaplate}\cleardoublepage
  % output each plate in turn
  \@whilenum\value{plateloop}<\value{plate}\do{%
    \addtocounter{plateloop}{1}
    % restore the figure counter from \@figcounti, \@figcountii, etc
    \setcounter{figure}{\value{figplate\roman{plateloop}}}
    % spew out the plate (in its own group for safety)
    {\csname @plate\roman{plateloop}\endcsname\par}
    }
  \pagestyle{rmaa}
  % reset the page number as if nothing had happened
  \addtocounter{page}{-\value{plate}}
  \addtocounter{page}{-\value{plate}}
  }
\newcommand\figcaption{%
  \def\@captype{figure}
  \renewcommand{\@makecaption}{\@makefigcap}
  \caption}
\newcommand\tabcaption{%
  \def\@captype{table}
  \renewcommand{\@makecaption}{\@maketabcap}
  \caption}

\newcommand{\plate}[2][1]{%
  %% User command to make a plate page: 
  %% Mandatory argument is what goes on the plate - this should not 
  %% include any floating environments (e.g. figure). Instead, \figcaption 
  %% should be used to do the captions. 
  %% Optional argument is the number of figures (i.e. captions) on the plate
  %% page (default is 1). This is necessary for incrementing the figure counter. 
  %%
  %% Note that we have to use a \gdef rather than an \xdef when we 
  %% save the plate output in a command. This means that in the book styles when it 
  %% comes to spewing out the plate, we will have forgotten about any user macros 
  %% that may have been defined. This can be got round by repeating the macros inside the plate.
  %% (Presumably, one would use a \newcommand rather than \renewcommand if they were new.)
  % increment global and local counters
  \addtocounter{plate}{1}
  % write appropriate label to .aux file
  \csname @platelabel\roman{plate}\endcsname
  % define a global macro called \@platei, \@plateii, \@plateiii, \@plateiv, ...
  \expandafter\gdef\csname @plate\roman{plate}\endcsname{%
    \csname @outputplatelabel\roman{plateloop}\endcsname
    \null\vfil #2 \vfil
    \cleardoublepage
    }
  % save the current shortauthor and shorttitle
  % need to be \xdef to take current values of \@shortauthor, etc 
  \expandafter\xdef\csname @shorta\roman{plate}\endcsname{\@shortauthor}
  \expandafter\xdef\csname @shortt\roman{plate}\endcsname{\@shorttitle}
  % save the current value of the figure counter for later use in \outputplates
  \setcounter{figplate\roman{plate}}{\value{figure}}
  % increment the figure counter according to the optional argument
  \addtocounter{figure}{#1}
  }
%% page style for plate pages 
\newcommand{\ps@rmaaplate}{%
  % Plate X (see page Y) in the foot
  \renewcommand\@oddfoot{\hfil\csname @shorta\roman{plateloop}\endcsname\ 
    (see page~\csname @plateref\roman{plateloop}\endcsname)\hfil}%
  \renewcommand\@evenfoot{\relax}%
  % nothing in even heads (blank pages)
  \renewcommand\@evenhead{\relax}
  % short title and page num in the odd heads
  \renewcommand\@oddhead{\hfil\csname @shortt\roman{plateloop}\endcsname%
    \hfil\llap{PLATE~\theplateloop}}
  }





%% Things for the User Guide
\if@guide
%% stolen from doc.sty
\def\MakeShortVerb#1{%
  \expandafter\ifx\csname cc\string#1\endcsname\relax
  \@shortvrbinfo{Made }{#1}%
  \add@special{#1}%
  \expandafter
  \xdef\csname cc\string#1\endcsname{\the\catcode`#1}%
  \begingroup
  \catcode`\~\active  \lccode`\~`#1%
  \lowercase{%
    \global\expandafter\let
    \csname ac\string#1\endcsname~%
    \gdef~{\verb~}}%
  \endgroup
  \global\catcode`#1\active
  \else
  \@shortvrbinfo\@empty{#1 already}%
  \fi}
\def\DeleteShortVerb#1{%
  \expandafter\ifx\csname cc\string#1\endcsname\relax
  \else
  \@shortvrbinfo{Deleted }{#1 as}%
  \rem@special{#1}%
  \global\catcode`#1\csname cc\string#1\endcsname
  \global \expandafter\let \csname cc\string#1\endcsname \relax
  \ifnum\catcode`#1=\active
  \begingroup
  \catcode`\~\active   \lccode`\~`#1%
  \lowercase{%
    \global\expandafter\let\expandafter~%
    \csname ac\string#1\endcsname}%
  \endgroup \fi \fi}
\def\@shortvrbinfo#1#2{%
  \PackageInfo{doc}{%
    #1\expandafter\@gobble\string#2 a short reference
    for \string\verb}}
\def\add@special#1{%
  \rem@special{#1}%
  \expandafter\gdef\expandafter\dospecials\expandafter
  {\dospecials \do #1}%
  \expandafter\gdef\expandafter\@sanitize\expandafter
  {\@sanitize \@makeother #1}}
\def\rem@special#1{%
  \def\do##1{%
    \ifnum`#1=`##1 \else \noexpand\do\noexpand##1\fi}%
  \xdef\dospecials{\dospecials}%
  \begingroup
  \def\@makeother##1{%
    \ifnum`#1=`##1 \else \noexpand\@makeother\noexpand##1\fi}%
  \xdef\@sanitize{\@sanitize}%
  \endgroup}
%% define | as shortcut for \verbatim
\MakeShortVerb{\|}
%% comand to specify arguments to LaTeX commands
\newcommand{\rmArg}[1]{%
  \ensuremath{\langle}{\scriptsize \textsf{\MakeUpperCase{#1}}}\ensuremath{\rangle}}
%\newcommand{\rmArg}[1]{\ensuremath{\langle}\textsc{#1}\ensuremath{\rangle}}
\providecommand{\BibTeX}{{\rmfamily B\kern-.05em%
    \textsc{i\kern-.025em b}\kern-.08em%
    T\kern-.1667em\lower.7ex\hbox{E}\kern-.125emX}}
%% Environment for code examples
\newenvironment{Code}{%
  \begin{list}{}{\setlength{\leftmargin}{0.05\columnwidth}} 
  \item[]}{\end{list}}
\fi
%</package>
%    \end{macrocode}

% \section{User commands for astronomical articles}
% \label{sec:user-commands}
% Some random definitions. These go both into the main class file and
% also into |adshtml.sty|.
% 
% Definitions of the journal name in different incarnations.  
%    \begin{macrocode}
%<*package|adshtml-style>
\newcommand{\RMAA}{RMxAA}
\newcommand{\RMAAlong}{Revista Mexicana de Astronom\'{\i}a y Astrof\'{\i}sica}
\newcommand{\RMAASC}{RMxAC}
\newcommand{\RMAASClong}{\RMAA\ (Serie de Conferencias)}
%    \end{macrocode}
% \begin{macro}{\ibidrule}
% This puts a horizontal rule with a period after it. It is used when
% the author list of a reference is identical to the preceding one. 
%    \begin{macrocode}
\newcommand{\ibidrule}{\rule[-0.03cm]{4.7em}{.01cm}.\ }
%    \end{macrocode}
% \end{macro}
% \begin{macro}{\ThisVolume}
% References to other articles in the same volume can be coded as,
% e.g., |\bibitem{} Bloggs, J. \ThisVolume{bloggs}|.
%    \begin{macrocode}
\newcommand{\ThisVolume}[1]{%
  \@year, \if@proceedings\RMAASC\else\RMAA\fi, \@volume,
  \if@book\oldpageref{#1:RMAAFirstPage}\else 000\fi\ (this volume)}
%    \end{macrocode}
% \end{macro}
% Other stuff that may be useful in ref lists.
%    \begin{macrocode}
\newcommand{\etal}{et al.\@}
%    \end{macrocode}
% Various new commands from AAS macros
%    \begin{macrocode}
\def\acknowledgments{\vskip 3ex plus .8ex minus .4ex}
\let\acknowledgements=\acknowledgments                  % second spelling
\let\jnl@style=\rm
\def\ref@jnl#1{{\jnl@style#1}}
\DeclareRobustCommand\aj{\ref@jnl{AJ}}                   
\DeclareRobustCommand\araa{\ref@jnl{ARA\&A}}             
\DeclareRobustCommand\apj{\ref@jnl{ApJ}}                 
\DeclareRobustCommand\apjl{\ref@jnl{ApJ}}                
\DeclareRobustCommand\apjs{\ref@jnl{ApJS}}               
\DeclareRobustCommand\ao{\ref@jnl{Appl.~Opt.}}           
\DeclareRobustCommand\apss{\ref@jnl{Ap\&SS}}             
\DeclareRobustCommand\aap{\ref@jnl{A\&A}}                
\DeclareRobustCommand\aapr{\ref@jnl{A\&A~Rev.}}          
\DeclareRobustCommand\aaps{\ref@jnl{A\&AS}}              
\DeclareRobustCommand\azh{\ref@jnl{AZh}}                 
\DeclareRobustCommand\baas{\ref@jnl{BAAS}}               
\DeclareRobustCommand\jrasc{\ref@jnl{JRASC}}             
\DeclareRobustCommand\memras{\ref@jnl{MmRAS}}            
\DeclareRobustCommand\mnras{\ref@jnl{MNRAS}}             
\DeclareRobustCommand\pra{\ref@jnl{Phys.~Rev.~A}}        
\DeclareRobustCommand\prb{\ref@jnl{Phys.~Rev.~B}}        
\DeclareRobustCommand\prc{\ref@jnl{Phys.~Rev.~C}}        
\DeclareRobustCommand\prd{\ref@jnl{Phys.~Rev.~D}}        
\DeclareRobustCommand\pre{\ref@jnl{Phys.~Rev.~E}}        
\DeclareRobustCommand\prl{\ref@jnl{Phys.~Rev.~Lett.}}    
\DeclareRobustCommand\pasp{\ref@jnl{PASP}}               
\DeclareRobustCommand\pasj{\ref@jnl{PASJ}}               
\DeclareRobustCommand\qjras{\ref@jnl{QJRAS}}             
\DeclareRobustCommand\skytel{\ref@jnl{S\&T}}             
\DeclareRobustCommand\solphys{\ref@jnl{Sol.~Phys.}}      
\DeclareRobustCommand\sovast{\ref@jnl{Soviet~Ast.}}      
\DeclareRobustCommand\ssr{\ref@jnl{Space~Sci.~Rev.}}     
\DeclareRobustCommand\zap{\ref@jnl{ZAp}}                 
\DeclareRobustCommand\nat{\ref@jnl{Nature}}              
\DeclareRobustCommand\iaucirc{\ref@jnl{IAU~Circ.}}
\DeclareRobustCommand\aplett{\ref@jnl{Astrophys.~Lett.}}
\DeclareRobustCommand\apspr{\ref@jnl{Astrophys.~Space~Phys.~Res.}}
\DeclareRobustCommand\bain{\ref@jnl{Bull.~Astron.~Inst.~Netherlands}}
\DeclareRobustCommand\fcp{\ref@jnl{Fund.~Cosmic~Phys.}}
\DeclareRobustCommand\gca{\ref@jnl{Geochim.~Cosmochim.~Acta}}
\DeclareRobustCommand\grl{\ref@jnl{Geophys.~Res.~Lett.}}
\DeclareRobustCommand\jcp{\ref@jnl{J.~Chem.~Phys.}}      
\DeclareRobustCommand\jgr{\ref@jnl{J.~Geophys.~Res.}}    
\DeclareRobustCommand\jqsrt{\ref@jnl{J.~Quant.~Spec.~Radiat.~Transf.}}
\DeclareRobustCommand\memsai{\ref@jnl{Mem.~Soc.~Astron.~Italiana}}
\DeclareRobustCommand\nphysa{\ref@jnl{Nucl.~Phys.~A}}
\DeclareRobustCommand\physrep{\ref@jnl{Phys.~Rep.}}
\DeclareRobustCommand\physscr{\ref@jnl{Phys.~Scr}}
\DeclareRobustCommand\planss{\ref@jnl{Planet.~Space~Sci.}}
\DeclareRobustCommand\procspie{\ref@jnl{Proc.~SPIE}}
\let\astap=\aap
\let\apjlett=\apjl
\let\apjsupp=\apjs
\let\applopt=\ao
\def\phn{\phantom{0}}
\def\phd{\phantom{.}}
\def\phs{\phantom{\ensuremath{-}}}
\def\phm#1{\phantom{#1}}
%    \end{macrocode}
% All these are redone as robust and using more modern commands where
% possible. The robustness is so they will not get expanded if used in
% abstracts and titles when these are passed to the ADS file. 
% \changes{3.16}{2003/04/07}{
%  Added \cs{Color} command for doing photometric colors.}
%    \begin{macrocode}
%%\def\sun{\hbox{$\odot$}}
\DeclareRobustCommand\sun{\ensuremath{\odot}}
%%\def\earth{\hbox{$\oplus$}}
\DeclareRobustCommand\earth{\ensuremath{\oplus}}
%%\def\lesssim{\mathrel{\hbox{\rlap{\hbox{\lower4pt\hbox{$\sim$}}}\hbox{$<$}}}}
%%\def\gtrsim{\mathrel{\hbox{\rlap{\hbox{\lower4pt\hbox{$\sim$}}}\hbox{$>$}}}}
\DeclareRobustCommand\lesssim{\mathrel{\hbox{\smaller\smaller\rlap{\hbox{\lower3pt\hbox{$\sim$}}}\raise1.5pt\hbox{$<$}}}}
\DeclareRobustCommand\gtrsim{\mathrel{\hbox{\smaller\smaller\rlap{\hbox{\lower3pt\hbox{$\sim$}}}\raise1.5pt\hbox{$>$}}}}
%%\def\lta{\mathrel{\spose{\lower 3pt\hbox{$\mathchar"218$}}
%%     \raise 2.0pt\hbox{$\mathchar"13C$}}}
%%\def\gta{\mathrel{\spose{\lower 3pt\hbox{$\mathchar"218$}}
%%     \raise 2.0pt\hbox{$\mathchar"13E$}}}
\DeclareRobustCommand\sq{\hbox{\rlap{$\sqcap$}$\sqcup$}}
%%\def\arcdeg{\hbox{$^\circ$}}
%%\def\arcmin{\hbox{$^\prime$}}
%%\def\arcsec{\hbox{$^{\prime\prime}$}}
\DeclareRobustCommand\arcdeg{\ensuremath{^\circ}}
\DeclareRobustCommand\arcmin{\ensuremath{^\prime}}
\DeclareRobustCommand\arcsec{\ensuremath{^{\prime\prime}}}
%%\def\fd{\hbox{$.\!\!^{\rm d}$}}
%%\def\fh{\hbox{$.\!\!^{\rm h}$}}
%%\def\fm{\hbox{$.\!\!^{\rm m}$}}
%%\def\fs{\hbox{$.\!\!^{\rm s}$}}
%%\def\fdg{\hbox{$.\!\!^\circ$}}
%%\def\farcm{\hbox{$.\mkern-4mu^\prime$}}
%%\def\farcs{\hbox{$.\!\!^{\prime\prime}$}}
%%\def\fp{\hbox{$.\!\!^{\scriptscriptstyle\rm p}$}}
\DeclareRobustCommand\fd{\ensuremath{.\!\!^\mathrm{d}}}
\DeclareRobustCommand\fh{\ensuremath{.\!\!^\mathrm{h}}}
\DeclareRobustCommand\fm{\ensuremath{.\!\!^\mathrm{m}}}
\DeclareRobustCommand\fs{\ensuremath{\mkern0.5mu.\mkern-5mu^\mathrm{s}\mkern-1mu}}
\DeclareRobustCommand\fdg{\ensuremath{.\!\!^\circ}}
\DeclareRobustCommand\farcm{\ensuremath{.\mkern-4mu^\prime}}
\DeclareRobustCommand\farcs{\ensuremath{\mkern1mu.\mkern-6mu^{\prime\prime}\mkern-1.5mu}}
%%\DeclareRobustCommand\farcs{\ensuremath{.\!\!^{\prime\prime}}}
\DeclareRobustCommand\fp{\ensuremath{.\!\!^{\scriptscriptstyle\rm p}}}
%%\def\micron{\hbox{$\mu$m}}
\DeclareRobustCommand\micron{\ensuremath{\mu\mathrm{m}}}
\let\la=\lesssim                        
\let\ga=\gtrsim
%%\def\case#1#2{\hbox{$\frac{#1}{#2}$}}
%%\def\slantfrac#1#2{\hbox{$\,^#1\!/_#2$}}
%%\def\onehalf{\slantfrac{1}{2}}
%%\def\onethird{\slantfrac{1}{3}}
%%\def\twothirds{\slantfrac{2}{3}}
%%\def\onequarter{\slantfrac{1}{4}}
%%\def\threequarters{\slantfrac{3}{4}}
\DeclareRobustCommand\case[2]{\ensuremath{\frac{#1}{#2}}}
\DeclareRobustCommand\slantfrac[2]{\ensuremath{\,^#1\!/_#2}}
\DeclareRobustCommand\onehalf{\slantfrac{1}{2}}
\DeclareRobustCommand\onethird{\slantfrac{1}{3}}
\DeclareRobustCommand\twothirds{\slantfrac{2}{3}}
\DeclareRobustCommand\onequarter{\slantfrac{1}{4}}
\DeclareRobustCommand\threequarters{\slantfrac{3}{4}}
%%\def\ubvr{\hbox{$U\!BV\!R$}}            
%%\def\ub{\hbox{$U\!-\!B$}}               
%%\def\bv{\hbox{$B\!-\!V$}}               
%%\def\vr{\hbox{$V\!-\!R$}}               
%%\def\ur{\hbox{$U\!-\!R$}}
%    \end{macrocode}
% New commands for ``photometric'' colors 
%    \begin{macrocode}
\DeclareRobustCommand\ubvr{\ensuremath{U\!BV\!R}}            
\DeclareRobustCommand\Color[2]{%
  \ensuremath{%
    #1\!\ifx#1V\!\fi-\!#2
  }}
\DeclareRobustCommand\ub{\Color{U}{B}}
\DeclareRobustCommand\bv{\Color{B}{V}}
\DeclareRobustCommand\vr{\Color{V}{R}}
\DeclareRobustCommand\ur{\Color{U}{R}}
%% New version of the \ion command, that actually works! 
%    \end{macrocode}
% \begin{macro}{\ion}
% This version of the |\ion| command allows you to write either
% |\ion{N}{4}| or |\ion{N}{iv}| in order to get N~IV. 
%    \begin{macrocode}
\DeclareRobustCommand\ion[2]{%
  \mbox{#1\kern0.2em%
    \smaller\rmfamily%
    \edef\@tempa{\@car#2\@nil}%
    \ifcat1\@tempa%
    \@Roman{#2}%
    \else%
    \uppercase{#2}%
    \fi}}
\newcounter{IonStage}
\DeclareRobustCommand\plainion[2]{\setcounter{IonStage}{#2}#1 \Roman{IonStage}}
%%\def\ion#1#2{\ensuremath{\mathrm{#1}\;}{\protect\small\rm\@Roman{#2}}}
%    \end{macrocode}
% \end{macro}
%    \begin{macrocode}
\DeclareRobustCommand\nodata{$\cdots$}
%    \end{macrocode}
% These are needed by the |\ion| command. Isn't there an easier way of
% doing this?
%    \begin{macrocode}
\let\@small=\small
\let\@normalsize=\normalsize
\let\@footnotesize=\footnotesize
\let\@scriptsize=\scriptsize
\let\@tiny=\tiny
\let\@large=\large
\let\@Large=\Large
\let\@LARGE=\LARGE
\let\@huge=\huge
\let\@Huge=\Huge
\newcommand{\smaller}{\small}
\renewcommand{\normalsize}{\renewcommand{\smaller}{\small}\@normalsize}
\renewcommand{\small}{\renewcommand{\smaller}{\footnotesize}\@small}
\renewcommand{\footnotesize}{\renewcommand{\smaller}{\scriptsize}\@footnotesize}
\renewcommand{\scriptsize}{\renewcommand{\smaller}{\tiny}\@scriptsize}
\renewcommand{\tiny}{\renewcommand{\smaller}{\tiny}\@tiny}
\renewcommand{\large}{\renewcommand{\smaller}{\normalsize}\@large}  
\renewcommand{\Large}{\renewcommand{\smaller}{\large}\@Large}  
\renewcommand{\LARGE}{\renewcommand{\smaller}{\Large}\@LARGE}  
\renewcommand{\huge}{\renewcommand{\smaller}{\LARGE}\@huge}  
\renewcommand{\Huge}{\renewcommand{\smaller}{\huge}\@Huge}  
%</package|adshtml-style>
%    \end{macrocode}


% \section{Other files}

% \subsection{The driver and make files}
% \label{sec:driver_make}
% These are put here for completeness, although they still need to be
% distributed separately since otherwise the user won't be able to
% unpack this file into its constituent parts.
% 
% The |Makefile| has two principal rules: |docs|, and |macros|. The
% |docs| rule produces a PDF and gzipped PS file of the documentation.
% I still haven't perfected the combination of rules required to get
% it to run \LaTeX{} enough times in all circumstances. The |macros|
% rule runs \LaTeX{} on |rmaa.ins| in order to unpack |rmaa.cls| and
% all the other little files (including the |Makefile| and the
% |rmaa.ins| files themselves\dots 
% 
% First set up lists of various types of file. 
%    \begin{macrocode}
%<*makefile>
DOCFILES=rmaa.pdf
MACRO_LTX_FILES=rmaa.cls rmaacite.sty adshtml.sty
SHELLSCRIPTS=makeads.sh
MACRO_OTHER_FILES=adshtml.perl sample.latex2html-init
MAKEFILE=Makefile
MACROFILES=$(MACRO_LTX_FILES) $(SHELLSCRIPTS) $(MACRO_OTHER_FILES)         
SKELETONFILES=ADS.tex 
%    \end{macrocode}
% Install the files read only into the package directory, specified via a path relative to the source directory. \changes{4.0}{2019/11/14}{Reorganize folder layout so that hardly anything is in the top level.  That way, things are clearer on github.}
%    \begin{macrocode}
INSTALL=install 
INSTALL_DIR=../packages
%    \end{macrocode}
% The default rule is to make everything. Installation must be
% explicitly asked for.
%    \begin{macrocode}
all: docs macros

install: install-docs install-macros

%    \end{macrocode}
% Most files are installed read-only to prevent editing. Skeleton
% files must be writeable. Must be careful not to overwrite the
% Makefile in the parent directory. 
%    \begin{macrocode}
install-macros: macros
        $(INSTALL) -m 444 $(MACRO_LTX_FILES) $(MACRO_OTHER_FILES) $(INSTALL_DIR)
        $(INSTALL) -m 555 $(SHELLSCRIPTS) $(INSTALL_DIR)
        $(INSTALL) -m 644 $(SKELETONFILES) $(INSTALL_DIR)

install-docs: docs
        $(INSTALL) rmaa.pdf $(INSTALL_DIR)/rm-fulldocs.pdf

%    \end{macrocode}
% Make the PDF and gzipped PS files. 
%    \begin{macrocode}
docs: $(DOCFILES)

%    \end{macrocode}
% Make everything else if either the |.dtx| or |.ins| files have
% changed. 
%    \begin{macrocode}
macros: rmaa.dtx rmaa.ins
%    \end{macrocode}
% Ensure that all the macro files can be written to then recreate
% them. The |touch| is there in case any of the files don't exist
% yet. 
%    \begin{macrocode}
        touch $(MACROFILES) $(MAKEFILE)
        chmod u+w $(MACROFILES) $(MAKEFILE) rmaa.ins
        latex "%%
%% This is file `rmaa.ins',
%% generated with the docstrip utility.
%%
%% The original source files were:
%%
%% rmaa.dtx  (with options: `ins')
%% 
%% PLEASE DO NOT EDIT THIS FILE! It is generated automatically from
%% the file rmaa.dtx using the accompanying Makefile and rmaa.ins
%% Any hand-edited changes will be lost
%% 
%% A LaTeX document class for the RMxAA
%% 
%% Author: William Henney (w.henney@crya.unam.mx)
%% 
%% Latest version available from: http://www.crya.unam.mx/rmaa
%% 
%% Copyright (C) 1999-2010 Instituto de Astronomia, UNAM, Mexico
%% 
%% This program is free software; you can redistribute it and/or
%% modify it under the terms of the GNU General Public License
%% as published by the Free Software Foundation; either version 2
%% of the License, or (at your option) any later version.
%% 
%% This program is distributed in the hope that it will be useful,
%% but WITHOUT ANY WARRANTY; without even the implied warranty of
%% MERCHANTABILITY or FITNESS FOR A PARTICULAR PURPOSE.  See the
%% GNU General Public License for more details.
%% 
%% You should have received a copy of the GNU General Public License
%% along with this program; if not, write to the Free Software
%% Foundation, Inc., 59 Temple Place - Suite 330, Boston, MA  02111-1307, USA.





%%
%% Installation batch file for the RMAA LaTeX package
%%
\def\batchfile{rmaa.ins}
\input docstrip.tex
\preamble

PLEASE DO NOT EDIT THIS FILE! It is generated automatically from
the file rmaa.dtx using the accompanying Makefile and rmaa.ins
Any hand-edited changes will be lost

A LaTeX document class for the RMxAA

Author: William Henney (w.henney@crya.unam.mx)

Latest version available from: http://www.crya.unam.mx/rmaa

Copyright (C) 1999-2010 Instituto de Astronomia, UNAM, Mexico

This program is free software; you can redistribute it and/or
modify it under the terms of the GNU General Public License
as published by the Free Software Foundation; either version 2
of the License, or (at your option) any later version.

This program is distributed in the hope that it will be useful,
but WITHOUT ANY WARRANTY; without even the implied warranty of
MERCHANTABILITY or FITNESS FOR A PARTICULAR PURPOSE.  See the
GNU General Public License for more details.

You should have received a copy of the GNU General Public License
along with this program; if not, write to the Free Software
Foundation, Inc., 59 Temple Place - Suite 330, Boston, MA  02111-1307, USA.
\endpreamble
\askforoverwritefalse
\generate{%
  \file{rmaa.cls}{\from{rmaa.dtx}{package}}
  \file{rmaa.ins}{\from{rmaa.dtx}{ins}}
  \file{rmaacite.sty}{\from{rmaa.dtx}{rmaacite}}
  \file{adshtml.sty}{\from{rmaa.dtx}{adshtml-style}}
  \file{ADS.tex}{\from{rmaa.dtx}{ads-container}}
  }
\generate{\nopreamble\nopostamble%
  \file{Makefile}{\from{rmaa.dtx}{makefile}}
  \file{adshtml.perl}{\from{rmaa.dtx}{adshtml-perl}}
  \file{makeads.sh}{\from{rmaa.dtx}{ads-shell-script}}
  \file{sample.latex2html-init}{\from{rmaa.dtx}{latex2html-init}}
  }

%% Local Variables:
%% mode: LaTeX
%% mode: Font-lock
%% End:









\endinput
%%
%% End of file `rmaa.ins'.
"
%    \end{macrocode}
% \TeX{} automatically adds the |.tex| extension to extensionless
% files. We fix this and take the opportunity to convert
% back to TABs all that it has lovingly converted to spaces. 
%    \begin{macrocode}
        cp -f $(MAKEFILE) $(MAKEFILE).BCK && unexpand $(MAKEFILE).tex >| $(MAKEFILE) && rm $(MAKEFILE).tex
%    \end{macrocode}
% Make the shell scripts executable and everything unwriteable (so
% they don't get hand-edited by mistake). However, undo this for the
% files that the user is \emph{supposed} to modify. 
%    \begin{macrocode}
        chmod a+x $(SHELLSCRIPTS)
        chmod a-w $(MACROFILES) $(MAKEFILE)
        chmod u+w $(SKELETONFILES)

%    \end{macrocode}
% We now use PDF\LaTeX{} for this.
%    \begin{macrocode}
rmaa.idx: rmaa.dtx
        @echo "REGENERATING .idx FILE"
        pdflatex "\nonstopmode% \def\fileversion{4.0} \def\filedate{2019/11/12}
% \MakeShortVerb{\|}
% \title{A \LaTeX{} document class for the Revista Mexicana de
% Astronom\'\i a y Astrof\'\i sica\thanks{This file has version
% \fileversion{} last revised \filedate{}.}}
% \author{William J. Henney\thanks{w.henney@irya.unam.mx}}
% \date{\filedate}
% \maketitle
% \paragraph{Abstract}
% This document describes the \texttt{rmaa} \LaTeX{}
% document class for the journal ``Revista Mexicana de
% Astronom\'\i a y Astrof\'\i sica''. The macros can be used for 
% \begin{enumerate}
% \item Preparing author manuscripts for the journal (including
% conference proceedings to be published in the ``Serie de
% Conferencias'')
% \item Typesetting an entire volume of the journal by inputting the
% individual article files from a container file. Facilities are
% included for automatic generation of TOC and Author Index. 
% \end{enumerate}
% 
% \tableofcontents
% \section{The documentation driver file}
% This is the part that generates the documentation that you may or
% may not be reading, which is done with the help of the Makefile
% discussed in section~\ref{sec:driver_make}. 
%    \begin{macrocode}
%<*driver>
\documentclass{ltxdoc}
\usepackage{hyperref}
\usepackage{palatino}
\EnableCrossrefs \CodelineIndex \RecordChanges
\setlength{\parindent}{0pt}
\def\MacroFont{\footnotesize\tt}
\begin{document} \DocInput{rmaa.dtx} \PrintIndex \PrintChanges \end{document}
%</driver>
%    \end{macrocode}

% \section{The implementation}
% 
% \changes{3.24}{2007/09/14}{Added user access to the macro version
% and date via \cs{rmVersion} and \cs{rmDate}}
% 
%    \begin{macrocode}
%<*package>
\newcommand{\rmaa@version}{3.26}
\newcommand{\rmaa@date}{2011/08/15}
\newcommand\rmVersion{\rmaa@version}
\newcommand\rmDate{\rmaa@date}
\NeedsTeXFormat{LaTeX2e}[1995/06/01]
%    \end{macrocode}
% \changes{4.0}{2019/11/02}{Change abstract to appear before
% resumen. Change RevMexAA to RMxAA for consistency with ADS.}
%
% \changes{3.22}{2007/03/22}{Allow proceedings articles to have
%   received/accepted dates} \changes{3.21}{2006/05/15}{Changes for
%   Silvia - abstracts are now shipped out to ADS.}
% \changes{3.19}{2006/01/19}{Changes for Silvia - remove adresses from
%   end, new abstract style} \changes{2.2beta}{2001/10/30}{ Fixed bug
%   whereby twocolumn option was not passed through to book
%   class. Finally sorted out the ``label on last page'' issue. }
% \changes{2.1}{2001/01/23}{ Fixed some inconsistencies in
%   \cs{nocite} and \cs{bibitem} uniqification Fixed the definition of
%   \cs{ion} so it will work in math mode and in contexts where the
%   typesize is different from normal (e.g., title, fig captions).
%   Modified the definition of \cs{@caption}, so we don't write any
%   list-of-figures data (we never use it anyway, and it just causes
%   problems with fragile commands in fig captions). }
% \changes{2.0}{2000/11/30}{ Added `newstyle' option, which should be
%   used in all new proceedings volumes (wider abstract, smaller
%   typeface in fig captions, etc.) Added commands for formatting
%   table captions. New demonstation document rmpostsamp.tex for
%   1-page posters. This will be the last release under a PRI
%   government (for a while...)}  \changes{1.9}{2000/06/06}{
%   Corrected position of page number in footer of articles that open
%   on even pages. Ensured that first article in each session opens on
%   an odd page, as do list of abstracts and author index. Added extra
%   space in TOC before the backmatter. Added command for editors,
%   \cs{BoldSessionTitles}, that makes produces bold titles for
%   session in the TOC.}  \changes{1.8}{2000/05/16}{ Fixed bug in
%   last page label with \cs{suppressfulladdresses} }
% \changes{1.7}{2000/05/11}{ Fixed more little things. Page labels
%   now work, even if you have no bibliography. You can output the
%   full addresses by hand with the \cs{outputfulladdresses}
%   command. Editors' Guide updated but still not finished. Revamped
%   splitbook.sh now produces PS file of each article for web page and
%   also an EPS file of each page for sending to the printers.}
% \changes{1.6}{2000/04/15}{ More cleaning up: full stops after
%   titles in index, running headers and footers in frontmatter,
%   fontsize in title of index, etc} \changes{1.5}{2000/04/15}{
%   Eliminated superscripts from titles in TOC and index}
% \changes{1.4}{2000/04/14}{ Added copyright and license info (GPL)}
% \changes{1.3}{2000/03/01}{ Corrected \cs{figcaption} to do the
%   caption formatting right Corrected font size of sec headings in
%   front matter} \changes{1.2}{2000/02/16}{ Added \cs{plate}
%   command} \changes{1.1}{2000/02/14}{ Added manuscript option,
%   corrected bugs} \changes{1.0}{2000/02/08}{ Now does a whole book
%   - lots of stuff tidied up} \changes{0.7}{1999/10/04}{ Increased
%   spacing between sections in single column styles. Suppressed
%   floats at bottom of pages (didn't work though) }
% \changes{0.6}{1999/09/24}{ Tweaked \cs{adjustfinalcols} (still not
%   perfect) Added proceedings option Added onecolumn/twocolumn
%   options (orthogonal to journal/proceedings) }
% \changes{0.5}{1999/09/23}{ Added \cs{fulladdresses} and
%   \cs{adjustfinalcols} (just about working)}
% \changes{0.4}{1999/09/22}{ Minor stylistic changes}
% \changes{0.3}{1999/06/28}{ Better mechanism for alternative
%   affiliations Added `guide' option for user guide}
% \changes{0.2}{1999/05/27}{ All basic functionality present, I
%   think} \changes{0.1}{1999/05/26}{ Just about working - still
%   missing Fig/Tab/Refs} \changes{0.0}{1999/05/21}{ So bored today,
%   decided to rewrite RMAA LaTeX macros from scratch}
%    \begin{macrocode}
\ProvidesClass{rmaa}[\rmaa@date\space v\rmaa@version\space RMxAA article
class]
%    \end{macrocode}
% \subsection{Option Handling}
% \label{sec:options}
% Declarations that are needed in options
% \changes{3.05}{2003/10/22}{Added natbib option}
% \changes{3.06}{2003/01/03}{Natbib option turned on permanently}
% \changes{3.07}{2003/01/11}{Added proof option}
% \changes{3.14}{2003/01/11}{Option to turn off natbib}
% \changes{3.18}{2003/10/22}{trying out doctex mode} 
% \changes{3.23}{2007/09/09}{Manuscript is now on by default.}
% \changes{3.26}{2011/08/15}{New option for review articles}
%    \begin{macrocode}
\newif\if@reviewarticle
\@reviewarticlefalse
\newif\if@proof
\@prooffalse
\newif\if@journal
\@journaltrue
\newif\if@printdates
\@printdatesfalse
\newif\if@proceedings
\@proceedingsfalse
\newif\if@manuscript
\@manuscripttrue
\newif\if@bigfont
\@bigfonttrue
\newif\if@guide
\@guidefalse
\newif\if@preprint
\@preprintfalse
\newif\if@rmtwocol
\@rmtwocoltrue
\newif\if@biblist
\@biblistfalse
\newif\if@book
\@bookfalse
\newif\if@natbib
\@natbibtrue
\newif\if@rmcite
\@rmcitefalse
\newif\if@titlevisible
\@titlevisibletrue
\newif\if@abstractpresent
\@abstractpresentfalse
\newif\if@printbanner
\@printbannerfalse
\newif\if@suppressfulladdresses
\@suppressfulladdressesfalse
\newif\if@newstyle
\@newstylefalse
\newif\if@debug
\@debugfalse
\newif\if@boldkeywords
\@boldkeywordsfalse
%    \end{macrocode}
% The option handling---mainly setting various booleans.
% \changes{2.2beta11}{2001/11/26}{Made twocol and
% newstyle the default for proceedings.}
% \changes{3.23}{2007/09/09}{Manuscript is now 10 pt instead of 12pt}
% \changes{3.23}{2007/09/10}{Full addresses are now automatically
% suppressed for proceedings style.}
% \changes{3.24}{2007/09/14}{New option \texttt{boldkeywords} that
% does the keywords in the old ugly way. Now off by default.}
%    \begin{macrocode}
\DeclareOption{manuscript}{\typeout{Manuscript for referee} 
  \@manuscripttrue\@journalfalse\@rmtwocolfalse
  \PassOptionsToClass{10pt}{article}}
\DeclareOption{bigfont}{\typeout{Big font manuscript requested by Gary
  Ferland} 
  \@manuscripttrue\@journalfalse\@rmtwocolfalse\@bigfonttrue
  \PassOptionsToClass{11pt}{article}}
\DeclareOption{journal}{\typeout{Main journal option}%
  \@journaltrue\@rmtwocoltrue\@printdatestrue\@manuscriptfalse
  \PassOptionsToClass{10pt}{article}}
\DeclareOption{guide}{\typeout{User guide option}%
  \@guidetrue\@journalfalse\@rmtwocoltrue\@manuscriptfalse
  \PassOptionsToClass{10pt}{article}}
\DeclareOption{preprint}{\typeout{Preprint option}%
  \@preprinttrue\@manuscriptfalse\PassOptionsToClass{10pt}{article}}
\DeclareOption{proof}{\typeout{Author proof option}%
  \@prooftrue\@manuscriptfalse}
\DeclareOption{proceedings}{\typeout{Proceedings option}%
  \@proceedingstrue\@journalfalse\@rmtwocoltrue\@newstyletrue
  \@suppressfulladdressestrue
  \@printdatesfalse\@manuscriptfalse\PassOptionsToClass{10pt}{article}} 
\DeclareOption{printdates}{\typeout{Force printing of received/accepted dates}%
  \@printdatestrue} 
\DeclareOption{book}{\typeout{Whole issue option}%
  \@booktrue} 
\DeclareOption{nonatbib}{\typeout{Turn off natbib support (required
    for compatibility with old rmaacite package).}%
  \@natbibfalse} 
\DeclareOption{newstyle}{\typeout{Use new layout style}%
  \@newstyletrue}
\DeclareOption{oldstyle}{\typeout{Use old layout style}%
  \@newstylefalse}
\DeclareOption{debug}{\typeout{Loads of annoying messages}%
  \@debugtrue}
%    \end{macrocode}
% These two need to be declared explicitly so that |\ExecuteOptions|
% actually does something.
%    \begin{macrocode}
\DeclareOption{twocolumn}{\@rmtwocoltrue}
\DeclareOption{onecolumn}{\@rmtwocolfalse}
\DeclareOption{twoside}{\PassOptionsToClass{\CurrentOption}{article}}
\DeclareOption{boldkeywords}{\@boldkeywordstrue}
%    \end{macrocode}
% Bold, upper-case keywords used to be the default. Now they must be
% turned on explicitly with the \texttt{boldkeywords} option. 
%    \begin{macrocode}
\DeclareOption{boldkeywords}{\@boldkeywordstrue}
%    \end{macrocode}
% Pass any other option along to the \texttt{article} class.
%    \begin{macrocode}
\DeclareOption*{\PassOptionsToClass{\CurrentOption}{article}}
%    \end{macrocode}
% Default options are to be two-sided and based on the standard
% \LaTeX{} \texttt{article} class,
% unless we are processing the entire book, in which case we base
% ourselves on the standard \texttt{book} class. Default is to use
% manuscript option unless otherwise specified. 
%    \begin{macrocode}
\ExecuteOptions{twoside,manuscript} \ProcessOptions\relax
\if@book
\if@rmtwocol\LoadClass[twocolumn]{book}\else\LoadClass{book}\fi
\newcommand{\abstract}{\relax}
\newcommand{\@maketitle}{\relax}
\else
\if@rmtwocol\LoadClass[twocolumn]{article}\else\LoadClass{article}\fi
\fi
%    \end{macrocode}
% Attempt to use natbib for citations 02 Jan 2003. 
% \changes{3.05}{2003/01/02}{If the natbib option is used, then set
% things up so that natbib's version of the bibliography is identical
% to the standard one.}
% \changes{3.09}{2003/01/23}{Explicitly require a recent version of
% natbib, so that we get a decent error message if the package is too
% old.} 
% \changes{3.23}{2007/09/09}{Now allow Bibtex to work properly} 
% \changes{3.25}{2009/10/01}{Change \cs{newcommand} tp \cs{def} when setting \cs{bibfont}. This allows it to work with all versions of natbib.} 
% \changes{3.25}{2010/02/06}{Updated copyright year.} 
% \changes{3.25}{2010/02/06}{Updated email and web addresses to \texttt{crya.unam.mx}.} 
%    \begin{macrocode}
\if@natbib
\RequirePackage{natbib}[1999/02/25]
\bibpunct{(}{)}{;}{a}{}{,}
\def\bibfont{\small}
\renewcommand\bibsection{\section*{REFERENCES}}
\setlength\bibsep{0pt}
\setlength\bibhang{1.5em}
\bibliographystyle{rmaa}
\fi
%    \end{macrocode}
% We want the \texttt{graphicx} package for the extended syntax of
% |\includegraphics|. We want the \texttt{ifthen} package for using in
% these macros (although we don't much yet). We explicitly load
% \texttt{babel} since some \LaTeX{} distributions messed up the
% hyphenation if we didn't.  
%    \begin{macrocode}
\RequirePackage{graphicx}
\RequirePackage{ifthen}
\IfFileExists{babel.sty}{%
%   \let\@restoreindent=\@afterindentfalse
  \RequirePackage[english]{babel}
%   \@restoreindent
  %% reinstate the definition of \@roman so as not to break .idx files
  \@ifundefined{org@@roman}{\relax}{\renewcommand{\@roman}[1]{\org@@roman{##1}}}
  }{%
  \typeout{Package babel.sty not found - don't worry about this}}
%    \end{macrocode}
% Try to load the \texttt{dcolumn} and \texttt{booktabs} packages if
% they can be found. These can produce nicer-looking tables. 
%    \begin{macrocode}
\IfFileExists{dcolumn.sty}{%
  \RequirePackage{dcolumn}
  }{%
  \typeout{Package dcolumn.sty not found - you will not be able to
    align columns on decimal points in tables}}
\IfFileExists{booktabs.sty}{%
  \RequirePackage{booktabs}
  }{%
  \typeout{Package booktabs.sty not found - your tables will not look
    as pretty as they ought too.}
  }
\IfFileExists{makeidx.sty}{%
  \RequirePackage{makeidx}
}{%
  \typeout{Package makeidx.sty not found - Author Index may not work.}
  }
%% commented out for moment since requires more modern version of LaTeX
%% \RequirePackage[latin1]{inputenc}
%    \end{macrocode}
% \changes{4.0}{2019/11/13}{Make \texttt{hyperref} and
% \texttt{etoolbox} into hard requirements, since we now use them
% extensively.}  We use hyperref to get internal and external links.
% But I don't like the colored boxes so I turn those off.  We use
% etoolbox for patching and robustifying at the moment.  In the
% future, I want to port a lot of other stuff over to it.  For
% instance, replace the ifthen package and the TeX ifs too.
%
% We target the latest version (1.7) of PDF, rather than the default
% of 1.5.  This reduces the chance of getting warnings from included
% graphics.  This is not strictly related to hyperref, but we put it
% here anyway.
%    \begin{macrocode}
\RequirePackage{etoolbox}
\RequirePackage{hyperref}
\hypersetup{hidelinks}
\pdfminorversion=7
%    \end{macrocode}
% We patch the \cs{@combinedblfloats} following a suggestion by Heiko
% Oberdiek at \url{https://tex.stackexchange.com/a/249743/3843}, which
% avoids problems when hyperlinks are split across
% pages. \changes{4.0}{2019/11/12}{Fix for issue where breaking
% hyperlink across pages gave an error.}
%    \begin{macrocode}
\urlstyle{same}
\newcommand\email{\begingroup \urlstyle{rm}\Url}
\patchcmd\@combinedblfloats{\box\@outputbox}{\unvbox\@outputbox}{}{%
  \errmessage{\noexpand\@combinedblfloats could not be patched}%
}%
%    \end{macrocode}
% \subsection{Other setup stuff}\label{sec:setup}
% Some declarations: 
%    \begin{macrocode}
\newif\if@needaddresses
\@needaddressestrue
\newif\if@frontmatter
\@frontmatterfalse
\newif\if@inpreamble
\@inpreambletrue
\newif\if@firstauthor
\@firstauthortrue
\newif\if@firstaffil
\@firstaffiltrue
\newif\if@firstaltaffil
\@firstaltaffiltrue
\newif\if@firstkwd
\@firstkwdtrue
\newif\if@noexpandADS
\@noexpandADSfalse
\newif\if@sessiontitlepages
\@sessiontitlepagesfalse
%    \end{macrocode}
% \begin{macro}{\@rmaaEmpty}
% \begin{macro}{\@rmaaArg}
% These are used in various commands to check if arg was left blank
%    \begin{macrocode}
\newcommand{\@rmaaEmpty}{}
\newcommand{\@rmaaArg}{}
%    \end{macrocode}  
% \end{macro}
% \end{macro}
%
% \subsection{Compatibility with old \LaTeX{} versions}
% \label{sec:compat}
% This package is based on the 1998/12/01 release of \LaTeX{}.
% Older releases may not define all the commands that we need.
% Luckily, we can use |\providecommand| to define these if necessary. 
% The following commands are not defined in the 1994/06/01 release. 
%
% \begin{macro}{\MakeUpperCase}
% \begin{macro}{\textsuperscript}
% \begin{macro}{\textbackslash}
% User commands:
%    \begin{macrocode}
\providecommand{\MakeUpperCase}[1]{\uppercase{#1}}
\providecommand{\textsuperscript}[1]{\ensuremath{^\textrm{#1}}}
\providecommand{\textbackslash}{\ensuremath{\backslash}}
%    \end{macrocode}  
% \end{macro}
% \end{macro}
% \end{macro}
% \begin{macro}{\protected@xdef}
% \begin{macro}{\restore@protect}
% \begin{macro}{\@unexpandable@protect}
% \begin{macro}{\hb@xt@}
% \LaTeX{} internal commands:
%    \begin{macrocode}
\providecommand{\protected@xdef}{%
        \let\@@protect\protect
        \let\protect\@unexpandable@protect
        \afterassignment\restore@protect
        \xdef}
\providecommand{\restore@protect}{\let\protect\@@protect}
\providecommand{\@unexpandable@protect}{\noexpand\protect\noexpand}
\providecommand{\hb@xt@}{\hbox to}
%    \end{macrocode}
% \end{macro}
% \end{macro}
% \end{macro}
% \end{macro}
%
% \subsection{Page layout}\label{sec:page-layout}
%
% The dimensions of the text were defined by Elizabeth. Text width is
% 169\,mm with equal margins on the two sides: 
%    \begin{macrocode}
\setlength{\textwidth}{169mm}
\setlength{\@tempdima}{\paperwidth}
\addtolength{\@tempdima}{-\textwidth}
\setlength{\oddsidemargin}{0.5\@tempdima}
\addtolength{\oddsidemargin}{-1in}
\setlength{\evensidemargin}{\oddsidemargin}
%    \end{macrocode}
% Text height is 225\,mm, top margin is 14\,mm. 
%    \begin{macrocode}
\setlength{\textheight}{225mm} 
\setlength{\topmargin}{-1in} 
\addtolength{\topmargin}{14mm}
%    \end{macrocode}
% Space between columns in two-column styles, header, footer,
% paragraph indentation:
%    \begin{macrocode}
\setlength{\columnsep}{9.5mm}
\setlength{\headsep}{0.7cm}
\setlength{\footskip}{0.75cm}
\setlength{\parindent}{1.5em}
%    \end{macrocode}
% 
% \begin{macro}{\@set@scalable@length}
% Try out a new scheme for dealing with various spacing dimensions and
% allowing the user to easily modify them. The command
% \begin{quote}
% |\set@scalable@length|\marg{name}\marg{defaultlength}\marg{stretch}\marg{shrink}
% \end{quote}
% does a 
% |\setlength|\marg{\cs{name}}\marg{defaultlength plus stretch minus shrink}
% and also defines the user command
% \textbackslash\meta{name}|Scale|\marg{factor} that can 
% be used to rescale length by \meta{factor}. Additionally, now
% defines the package command
% \textbackslash\meta{name}|@set@scale|\marg{FACTOR} 
% that also rescales the length by \meta{FACTOR}, so that the length
% becomes \meta{factor} $\times$ \meta{FACTOR} $\times$
% \meta{defaultlength}.   
%    \begin{macrocode}
\newcommand\@set@scalable@length[4]{%
  %% arg #1: NAME of length without leading `\'
  %%
  %% arg #2: default VALUE  
  \expandafter\newlength\csname @#1@default\endcsname
  \expandafter\setlength\csname @#1@default\endcsname{#2}%
  \expandafter\newlength\csname @@#1@@default\endcsname
  \expandafter\setlength\csname @@#1@@default\endcsname{#2}%
  %% arg #3: default STRETCH  
  \expandafter\newlength\csname @#1@stretch\endcsname
  \expandafter\setlength\csname @#1@stretch\endcsname{#3}%
  \expandafter\newlength\csname @@#1@@stretch\endcsname
  \expandafter\setlength\csname @@#1@@stretch\endcsname{#3}%
  %% arg #4: default SHRINK  
  \expandafter\newlength\csname @#1@shrink\endcsname
  \expandafter\setlength\csname @#1@shrink\endcsname{#4}%
  \expandafter\newlength\csname @@#1@@shrink\endcsname
  \expandafter\setlength\csname @@#1@@shrink\endcsname{#4}%
  %%
  %% User-defined scale \@NAME@scale is unity by default
  \expandafter\gdef\csname @#1@scale\endcsname{1}%
  %% Package-defined scale \@NAME@SCALE is unity by default
  \expandafter\gdef\csname @#1@SCALE\endcsname{1}%
  %%
  %% internal command \@reset@NAME should expand to 
  %% \setlength\@@NAME@@default{\@NAME@SCALE\@NAME@default}
  %% \setlength\NAME{\@NAME@scale\@@NAME@@default}
  %% in other words: scale x SCALE x default 
  \expandafter\gdef\csname @reset@#1\endcsname{%
    \expandafter\setlength\csname @@#1@@default\endcsname{%
      \csname @#1@SCALE\endcsname \csname @#1@default\endcsname
      }
    \expandafter\setlength\csname @@#1@@stretch\endcsname{%
      \csname @#1@SCALE\endcsname \csname @#1@stretch\endcsname
      }
    \expandafter\setlength\csname @@#1@@shrink\endcsname{%
      \csname @#1@SCALE\endcsname \csname @#1@shrink\endcsname
      }
    \expandafter\setlength\csname #1\endcsname{%
      \csname @#1@scale\endcsname \csname @@#1@@default\endcsname
      \@plus
      \csname @#1@scale\endcsname \csname @@#1@@stretch\endcsname
      \@minus
      \csname @#1@scale\endcsname \csname @@#1@@shrink\endcsname
      }}
  %%
  %% User command \NAMEScale should expand to 
  %% {\renewcommand\@NAME@scale{#1}\@reset@NAME}
  \expandafter\gdef\csname #1Scale\endcsname##1{%
    \expandafter\gdef\csname @#1@scale\endcsname{##1}%
    \csname @reset@#1\endcsname}
  %%
  %% Package command \NAME@set@SCALE should expand to 
  %% {\renewcommand\@NAME@SCALE{#1}\@reset@NAME}
  \expandafter\gdef\csname #1@set@SCALE\endcsname##1{%
    \expandafter\gdef\csname @#1@SCALE\endcsname{##1}%
    \csname @reset@#1\endcsname}
  %%
  %% initialize at default values with \@reset@NAME
  \csname @reset@#1\endcsname
  }
%    \end{macrocode}
% \end{macro}
% \begin{macro}{\parskip}
% First use of this new infrastructure: set interparagraph space to be
% zero but with a small amount of stretchiness.
%    \begin{macrocode}
%   \ifthenelse{\boolean{@manuscript}}{%
%     \@set@scalable@length{parskip}{0.5\baselineskip}{0.3\p@}{0.1\p@}
%   }{
    \@set@scalable@length{parskip}{0\p@}{0.3\p@}{0\p@}
%   }
%    \end{macrocode}
% \end{macro}
% 
% \begin{macro}{\abovesectionskip}
% \changes{4.0}{2019/11/12}{Use negative values for skip so that first paragraph after section headings is not indented.  This is something that I have been trying to fix for years!  It turns out that the anser was in the documentation for the standard command \cs{@startsection.}}
% \begin{macro}{\belowsectionskip}
% \begin{macro}{\aboveparaskip}
% \begin{macro}{\belowparaskip}
% Define some utility rubber lengths for use in
% section~\ref{sec:section-commands}. 
%    \begin{macrocode}
\newlength{\abovesectionskip}
\newlength{\belowsectionskip}
\newlength{\aboveparaskip}
\newlength{\belowparaskip}
%    \end{macrocode}
% The default values for these lengths are those traditionally used
% for the onecolumn style. 
%    \begin{macrocode}
\@set@scalable@length{abovesectionskip}{-3.5ex}{-4ex}{-2ex}
\@set@scalable@length{belowsectionskip}{2.3ex}{0ex}{2ex}
\@set@scalable@length{aboveparaskip}{3.25ex}{2ex}{0.5ex}
\@set@scalable@length{belowparaskip}{-1em}{\z@}{\z@}
%    \end{macrocode}
% \end{macro}\end{macro}\end{macro}\end{macro}
% We also make scalable some of the standard float parameters. These
% are the same defaults as in the standard article classes, but
% slightly more stretchy. For the twocolumn style, they are currently
% halved, see below. Note that I have also lumped |\abovecaptionskip|
% in with these, since it plays a similar role. There is also a
% |\belowcaptionskip| but it is set to zero in the standard classes. 
%    \begin{macrocode}
\@set@scalable@length{floatsep}{12\p@}{4\p@}{4\p@}
\@set@scalable@length{textfloatsep}{20\p@}{4\p@}{8\p@}
\@set@scalable@length{intextsep}{12\p@}{4\p@}{4\p@}
\@set@scalable@length{dblfloatsep}{12\p@}{4\p@}{4\p@}
\@set@scalable@length{dbltextfloatsep}{20\p@}{4\p@}{8\p@}
\@set@scalable@length{abovecaptionskip}{10\p@}{\z@}{\z@}
%    \end{macrocode}
% \begin{macro}{\@before@title@skip}         
% \begin{macro}{\@before@author@skip}        
% \begin{macro}{\@between@author@skip}       
% \begin{macro}{\@before@abstract@skip}      
% \begin{macro}{\@after@abstracthead@skip}   
% \begin{macro}{\@between@abstract@skip}     
% \begin{macro}{\@before@kwd@skip}           
% \begin{macro}{\@before@dates@skip}           
% Spacing parameters used in |\@maketitle|. These are not changed by
% default except by the user, so they don't need to go into
% |\@package@default@lengths|.  
%    \begin{macrocode}
\newlength{\@before@title@skip}         
\newlength{\@before@author@skip}        
\newlength{\@between@author@skip}       
\newlength{\@before@abstract@skip}      
\newlength{\@after@abstracthead@skip}   
\newlength{\@between@abstract@skip}     
\newlength{\@before@kwd@skip}           
\newlength{\@before@dates@skip}           
\@set@scalable@length{@before@title@skip}{2em}{\z@}{\z@}
\@set@scalable@length{@before@author@skip}{1.5em}{\z@}{\z@}
\@set@scalable@length{@between@author@skip}{0.5em}{\z@}{\z@}
\@set@scalable@length{@before@abstract@skip}{3ex}{1ex}{0.5ex}
\@set@scalable@length{@after@abstracthead@skip}{1ex}{\z@}{\z@}
\@set@scalable@length{@between@abstract@skip}{2ex}{1ex}{0.5ex}
\@set@scalable@length{@before@kwd@skip}{1ex}{0.5ex}{0.2ex}
\@set@scalable@length{@before@dates@skip}{2ex}{0.5ex}{0.2ex}
%    \end{macrocode}
% \end{macro}\end{macro}\end{macro}\end{macro}
% \end{macro}\end{macro}\end{macro}\end{macro}
% 
% Spacing above and below displayed equations also wants to be done
% like this, but it is more complicated since the standard article
% classes make these change with the fontsize. Something to do later,
% maybe. 
% 
%    \begin{macrocode}
%    \end{macrocode}
% \begin{macro}{\@package@default@lengths}
% The value actually used for these lengths is controlled by the
% \meta{SCALE} parameter, which depends on the style option that is in
% effect. This is set using the |\@package@default@lengths| command,
% which must be called every time the style changes (e.g., changing
% from onecolumn to twocolumn). 
%    \begin{macrocode}
\newcommand{\@package@default@lengths}{%
  \ifthenelse{\boolean{@rmtwocol}}{%
    %% twocolumn mode
    \abovesectionskip@set@SCALE{0.5}
    \belowsectionskip@set@SCALE{0.5}
    \aboveparaskip@set@SCALE{0.5}
    \belowparaskip@set@SCALE{1}
    \floatsep@set@SCALE{0.5}
    \textfloatsep@set@SCALE{0.5}
    \intextsep@set@SCALE{0.5}
    \dblfloatsep@set@SCALE{0.5}
    \dbltextfloatsep@set@SCALE{0.5}
    }{%
    %% onecolumn mode
    \abovesectionskip@set@SCALE{1}
    \belowsectionskip@set@SCALE{1}
    \aboveparaskip@set@SCALE{1}
    \belowparaskip@set@SCALE{1}
    \floatsep@set@SCALE{1}
    \textfloatsep@set@SCALE{1}
    \intextsep@set@SCALE{1}
    \dblfloatsep@set@SCALE{1}
    \dbltextfloatsep@set@SCALE{1}
    }%
  }
\@package@default@lengths
%    \end{macrocode}
% \end{macro}
% \begin{macro}{\RescaleLengths}
% \begin{macro}{\RescaleSecLengths}
% \begin{macro}{\RescaleFloatLengths}
% \begin{macro}{\RescaleTitleLengths}
% User commands that will fiddle with all the scalable lengths. 
%    \begin{macrocode}
\newcommand{\RescaleLengths}[1]{%
  \RescaleSecLengths{#1}\RescaleFloatLengths{#1}\RescaleTitleLengths{#1}}
\newcommand{\RescaleSecLengths}[1]{%
%% sections
  \abovesectionskipScale{#1}
  \belowsectionskipScale{#1}
  \aboveparaskipScale{#1}
  }
\newcommand{\RescaleFloatLengths}[1]{%
%% floats
  \floatsepScale{#1}
  \textfloatsepScale{#1}
  \intextsepScale{#1}
  \dblfloatsepScale{#1}
  \dbltextfloatsepScale{#1}
  \abovecaptionskipScale{#1}
  }
\newcommand{\RescaleTitleLengths}[1]{%
%% maketitle
  \@before@title@skipScale{#1}           
  \@before@author@skipScale{#1}        
  \@between@author@skipScale{#1}       
  \@before@abstract@skipScale{#1}      
  \@after@abstracthead@skipScale{#1}   
  \@between@abstract@skipScale{#1}     
  \@before@kwd@skipScale{#1}           
  \@before@dates@skipScale{#1}           
  }
%    \end{macrocode}
% \end{macro}\end{macro}\end{macro}\end{macro}
% \changes{2.2beta14}{2002/03/17}{Finished implementation of new
% infrastructure for scalable lengths. Converted the above/below
% section skips to use it. Also, all the floatsep parameters. Also,
% all the vertical lengths that are used by |\@maketitle|. }

% \begin{macro}{\MSextramargin}
% For the \texttt{manuscript} option we have an extra 0.3\,inch
% vertical margin and 1\,inch horizontal margin for marginal notes (by
% hand or with |\marginpar|. 
%
% \changes{3.23}{2007/09/09}{Reduce baselinestretch in
%   \texttt{manuscript} style to 1.2. No longer use raggedright for
%   manuscript.}
%
% \changes{3.24}{2007/09/14}{Further reduction of baselinestretch in
%   \texttt{manuscript} style to 1.1 since we now have wider margins.}
%
% \changes{3.24}{2007/09/14}{Split \texttt{manuscript}'s extra margin
%   into separate horizontal and vertical parts.}
% 
%    \begin{macrocode}
\newlength{\MSextraHmargin}\newlength{\MSextraVmargin}
\setlength{\MSextraHmargin}{1in} 
\setlength{\MSextraVmargin}{0.3in} 
\if@manuscript
   \renewcommand{\baselinestretch}{1.1}
   \addtolength{\textwidth}{-2\MSextraHmargin}
   \addtolength{\oddsidemargin}{\MSextraHmargin}
   \setlength{\evensidemargin}{\oddsidemargin}
   \addtolength{\textheight}{-2\MSextraVmargin}
   \addtolength{\topmargin}{\MSextraVmargin}
%    \raggedright
\fi
%    \end{macrocode}
% \end{macro}
% Float parameters: don't allow bottom floats by default. 
%    \begin{macrocode}
\renewcommand{\fps@figure}{tp}
\renewcommand{\fps@table}{tp}
%    \end{macrocode}
%
% \begin{macro}{\rmaa@AuthorWidthFrac}
% \begin{macro}{\rmaa@AbstractWidthFrac}
% These parameters control what fraction of the total text width can
% be used by the list of authors and the abstract/resumen,
% respectively.  
%    \begin{macrocode}
\newcommand{\rmaa@AuthorWidthFrac}{1.0}
\newcommand{\rmaa@AbstractWidthFrac}{0.75}
%    \end{macrocode}
% Use a wide abstract to save space in the new style.
%    \begin{macrocode}
\if@newstyle\renewcommand{\rmaa@AbstractWidthFrac}{1.0}\fi
%    \end{macrocode}
% \end{macro}
% \end{macro}
%
% Some things to consider in future versions. Trying to avoid too many
% widows, etc. Apparently, the following will work best with
% \cs{flushbottom}. 
%    \begin{macrocode}
%%\widowpenalty=10000
%%\clubpenalty=10000
%    \end{macrocode}
% Also, should perhaps try getting rid of underfull/overfull box warnings. 
% \changes{3.10}{2003/01/24}{Set large values for \cs{hbadness},
% \cs{vbadness}, \cs{hfuzz}, \cs{vfuzz}. To reduce amount of stuff
% that gets written to the terminal every time you compile something.}
%    \begin{macrocode}
\hbadness=10000
\vbadness=10000
\setlength\hfuzz{\maxdimen}
\setlength\vfuzz{\maxdimen}
%    \end{macrocode}
% \subsection{Various text styles for different contexts}
% \label{sec:text-styles} 
% \begin{macro}{\rmaa@AffilStyle}
% \begin{macro}{\rmaa@PagesStyle}
% \begin{macro}{\rmaa@YearStyle}
% \begin{macro}{\rmaa@VolumeStyle}
% \begin{macro}{\rmaa@VolumenStyle}
% \begin{macro}{\rmaa@JournalNameStyle}
% \begin{macro}{\rmaa@ReceivedStyle}
% \begin{macro}{\rmaa@AbstractHeaderStyle}
% \begin{macro}{\rmaa@JournalNameStyle}
% Style commands for the different fields used by |\maketitle| in
% constructing the title page of the article. These have the form
% |\rmaa@XYZStyle{TEXT}|. I'm not sure these were really a good idea,
% but never mind \dots 
%    \begin{macrocode}
\newcommand{\rmaa@AffilStyle}[1]{{\normalsize\rmfamily#1}}
\newcommand{\rmaa@PagesStyle}[1]{{\normalsize\rmfamily#1}}
\newcommand{\rmaa@YearStyle}[1]{{\normalsize\rmfamily#1}}
\newcommand{\rmaa@VolumeStyle}[1]{{\normalsize\rmfamily\bfseries#1}}
\newcommand{\rmaa@VolumenStyle}[1]{{\Large\sffamily\bfseries#1}}
\newcommand{\rmaa@JournalNameStyle}[1]{{\normalsize\itshape#1}}
\newcommand{\rmaa@ReceivedStyle}[1]{{\normalsize\rmfamily\itshape#1}}
\newcommand{\rmaa@AbstractHeaderStyle}[1]{%
  {\normalsize\rmfamily\MakeUpperCase{#1}}}
\newcommand{\rmaa@TitleStyle}[1]{%
  {\large\bfseries\rmfamily\boldmath\MakeUpperCase{#1}\par}}
%    \end{macrocode}
% \end{macro}
% \end{macro}
% \end{macro}
% \end{macro}
% \end{macro}
% \end{macro}
% \end{macro}
% \end{macro}
% \end{macro}
%
% \begin{macro}{\rmaa@FrontSecStyle}
% Headings of ``front sections'' (see section~\ref{sec:front-sections})
%    \begin{macrocode}
\newcommand{\rmaa@FrontSecStyle}[1]{%
  {\normalsize\bfseries\rmfamily\MakeUpperCase{#1}}}
%    \end{macrocode}
% \end{macro}
%
% \begin{macro}{\rmaa@TOCTitleStyle}
% \begin{macro}{\rmaa@TOCPartStyle}
% \begin{macro}{\rmaa@IndexTitleStyle}
% \begin{macro}{\rmaaTOCAuthorStyle}
% \begin{macro}{\rmaaIndexAuthorStyle}
% Commands to set style of TOC and index entries.
% Rewritten to avoid unnecessary font commands in the .toc and .ind files
%    \begin{macrocode}
\newcommand{\rmaa@TOCTitleStyle}[1]{\MakeUpperCase{#1}}
\newcommand{\rmaa@TOCPartStyle}[1]{\rmaa@TOCTitleStyle{#1}}
\newcommand{\rmaa@IndexTitleStyle}[1]{#1}
\newcommand{\rmaaTOCAuthorStyle}[1]{%
  {\itshape#1}}
\newcommand{\rmaaIndexAuthorStyle}[1]{%
  \protect\mbox{\bfseries#1}}
%    \end{macrocode}
% \end{macro}
% \end{macro}
% \end{macro}
% \end{macro}
% \end{macro}
%
% \begin{macro}{\rmaa@KeywordStyle}
% \begin{macro}{\rmaa@KeywordHeaderStyle}
% The style of the keywords depends on whether we are using the
% \texttt{newstyle} option or not. If we are, then they are small. 
%    \begin{macrocode}
\ifthenelse{\boolean{@newstyle}}{
  \newcommand{\rmaa@KeywordStyle}[1]{%
    {\small%
      \if@boldkeywords 
      \bfseries\MakeUpperCase{\raggedright #1}
      \else
      \raggedright #1
      \fi
    }}
  \newcommand{\rmaa@KeywordHeaderStyle}[1]{%
    {\small\itshape#1}}}
{
  \newcommand{\rmaa@KeywordStyle}[1]{%
    {\normalsize%
      \if@boldkeywords\bfseries\fi\raggedright #1}}
  \newcommand{\rmaa@KeywordHeaderStyle}[1]{%
    {\normalsize\itshape#1}}}
%    \end{macrocode}
% \end{macro}
% \end{macro}
%
% \subsection{Page styles}\label{sec:page-styles}
% \begin{macro}{\ps@rmaa}
% The running headers for normal pages:
%    \begin{macrocode}
\newcommand{\ps@rmaa}{%
%    \end{macrocode}
% nothing in either foot; 
%    \begin{macrocode}
  \renewcommand\@oddfoot{\relax}%
  \renewcommand\@evenfoot{\relax}%
%    \end{macrocode}
% page number and short author list in the even heads; 
%    \begin{macrocode}
  \renewcommand\@evenhead{\rlap{\thepage}\hfil\@shortauthor\hfil}%
%    \end{macrocode}
% short title and page number in the odd heads.
%    \begin{macrocode}
  \renewcommand\@oddhead{\hfil\@shorttitle\hfil\llap{\thepage}}
  }
%    \end{macrocode}
% \begin{macro}{\ps@rmaafirstpage}
% The header of the first page of each
% article:
%    \begin{macrocode}
\newcommand{\ps@rmaafirstpage}{%
  \renewcommand\@oddfoot{\hfil\thepage}%
%    \end{macrocode}
% footer corrected (06 Jun 2000) to have page on outside edge;
%    \begin{macrocode}
  \renewcommand\@evenfoot{\thepage\hfil}%
%    \end{macrocode}
% Header has journal name, volume, and page numbers for a regular
% article, or similar for special cases. 
% \changes{4.0}{2019/11/03}{Change to three-line header with copyright info and DOI.}
%    \begin{macrocode}
\if@preprint
\renewcommand\@oddhead{%
  \if@proceedings 
  \textsc{To appear in ``\@conftitle{} (\rmaa@YearStyle{\@year})''}\hfil
  \textit{RMxAC}
  \else
  \textsc{To appear in} \@journalname\hfil
  \fi}\fi%
\if@guide
\renewcommand\@oddhead{\textsc{Documentation} (\rmaa@YearStyle{\@year})\hfil}
\fi%
\if@manuscript
\renewcommand\@oddhead{Manuscript for \@journalname{} (\rmaa@YearStyle{\@year})\hfil}
\fi%
\if@journal
\renewcommand\@oddhead{%
  \parbox{\textwidth}{
    \@journalname, \rmaa@VolumeStyle{\@volume},
    \rmaa@PagesStyle{\@artpages}\
    (\rmaa@YearStyle{\@year})\hfil\\[1pt]
    {\small\textcopyright{} \@year:
      Instituto de Astronom\'\i{}a,
      Universidad Nacional Aut\'onoma de M\'exico}\hfil\\[1pt]
    {\small\@DOI@url}\hfil
  }%
}\fi%
  \renewcommand\@evenhead{\@oddhead}}
%    \end{macrocode}
% By default, use the |\ps@rmaa| page style.
%    \begin{macrocode}
\pagestyle{rmaa}
%    \end{macrocode}
% \end{macro}
% \end{macro}
%
% \subsection{Preamble commands}\label{sec:preamble} 
% Everything you need for setting the title page, TOC and index
% entries, etc.
% \subsubsection{Commands for the title page}
% \label{sec:title-commands}
% \begin{macro}{\resumen}
% \begin{macro}{\@resumen}
% \begin{macro}{\abstract}
% \begin{macro}{\@abstract}
% User commands for setting the resumen and abstract. 
% These should be used in the document header. The internal commands
% |\@resumen| and |\@abstract| are used to hold the actual text
% pending output via the |\maketitle| command. 
%    \begin{macrocode}
\newif\if@setresumen
\newcommand{\resumen}[1]{%
%%  \renewcommand{\@resumen}{{\selectlanguage{spanish}#1}}%
  \renewcommand{\@resumen}{{\ignorespaces#1}}%
  \@setresumentrue}
\renewcommand{\abstract}[1]{%
  \renewcommand{\@abstract}{{\ignorespaces#1}}
  \renewcommand\ADS@abstract{#1}
%% Not a good idea, since then author macros are not expanded and
%% cause errors later. However, by not doing it we are vulnerable to
%% fragile commands.....
  \if@noexpandADS\@onelevel@sanitize{\ADS@abstract}\fi
  \@abstractpresenttrue
  }
\newcommand{\ADSnoexpand}{\@noexpandADStrue}
%    \end{macrocode}
% \begin{macro}{\boldabstract}
% This is the version for doing 1-page conference proceedings. Still
% save the text for ADS but then write it out immediately as a bold
% paragraph. 
%    \begin{macrocode}
\newcommand{\boldabstract}[1]{%
  \gdef\ADS@abstract{#1}%
  \renewcommand{\@abstract}{#1}%
  \if@noexpandADS\@onelevel@sanitize{\ADS@abstract}\fi%
  \noindent\textbf{\boldmath #1}
  }
%    \end{macrocode}
% \end{macro}
% Default value in each case is to print out a warning if not set
%    \begin{macrocode}
  \newcommand{\@resumen}{%
    \@setresumenfalse \ClassWarningNoLine{rmaa}{No Spanish abstract
      (resumen) specified - if you do not speak Spanish, this will be
      translated by the editors}%
    \textbf{Favor de proporcionar un resumen en espa\~nol. If you are
      unable to translate your abstract into Spanish, the editors will
      do it for you.} }
  \newcommand{\@abstract}{\ClassWarningNoLine{rmaa}{No abstract
      specified}}
%    \end{macrocode}
% \end{macro}
% \end{macro}
% \end{macro}
% \end{macro}
%
% \begin{macro}{\title}
% \begin{macro}{\@rawtitle}
% \begin{macro}{\@TOCtitle}
% \begin{macro}{\rmaaIndexTitle}
% Similarly, |\title| is the preamble command to set the title of the
% article. This is used in various places, typeset in different
% styles, hence the use of the internal |\@rawtitle|, |\@title|,
% |\rmaaIndexTitle|, and |\@TOCtitle| commands, each of which uses the
% relevant style from section~\ref{sec:text-styles}
%    \begin{macrocode}
\newcommand{\@rawtitle}{Raw Title}
\newcommand{\@TOCtitle}{\rmaa@TOCTitleStyle{Title for TOC}}
\newcommand{\rmaaIndexTitle}{\rmaa@IndexTitleStyle{Title for index}.}
\newcommand{\@TOCaftertitle}{}
\renewcommand{\title}[1]{%
%%  {\ADS@sanitize\xdef\@rawtitle{#1}}
  \renewcommand\@rawtitle{\noexpand#1}
  \let\ADS@title=\@rawtitle
  \renewcommand{\@title}{\rmaa@TitleStyle{#1}}%
  \renewcommand{\@TOCtitle}{\rmaa@TOCTitleStyle{#1}\@TOCaftertitle}%
  \renewcommand{\rmaaIndexTitle}{\rmaa@IndexTitleStyle{#1\@TOCaftertitle}.}%
  }
%    \end{macrocode}
% \end{macro}
% \end{macro}
% \end{macro}
% \end{macro}
% 
% All the weird and whacky ways of indicating the author
% affiliations. 
% \begin{macro}{\affil}
% First, the way that sets it directly underneath the authors name.
% Note that |\affil| trumps all other affiliations for the purposes of
% getting into the ADS stuff (we use |\gdef|).   
% \dots 
%    \begin{macrocode}
\newcommand{\affil}[1]{%
  \\[1.5ex]\rmaa@AffilStyle{#1}\\[1.5ex]%
  \if@firstaffil\gdef\ADS@affil{#1}\@firstaffilfalse\fi%
  }
\renewcommand{\and}{and\\[1.5ex]}
%    \end{macrocode}
% \begin{macro}{\altaffil}
% Alternative affiliations are footnoted. 
%    \begin{macrocode}
\newcommand{\altaffil}[1]{%
  \if@firstaltaffil\gdef\ADS@affil{#1}\@firstaltaffilfalse\fi%
  \thanks{#1}}
%    \end{macrocode}
% \begin{macro}{\altaffilmark}
% \begin{macro}{\altaffiltext}
% Alternative-alternative mechanism for affiliations---should not be
% mixed with |\altaffil| (unless you know what you are doing) or there
% will be problems with numbering.
%    \begin{macrocode}
\newcommand{\altaffilmark}[1]{\textsuperscript{#1}}
\newcommand{\altaffiltext}[2]{%
  \if@firstaltaffil\gdef\ADS@affil{#2}\@firstaltaffilfalse\fi
  \stepcounter{footnote}%
  \protected@xdef\@thanks{\@thanks
    \protect\footnotetext[#1]{#2}}} 
%    \end{macrocode}
% \end{macro}
% \end{macro}
% \end{macro}
% \end{macro}
%
% \begin{macro}{\keywords}
% \begin{macro}{\@keywords}
% \begin{macro}{\@keywordheader}
% \label{keywords}Sundry minor preamble commands. All
% self-explanatory, I hope. |\keywords| is the old user interface
% (deprecated). New articles should use |\addkeyword| once for each
% keyword. That way the ADS record comes out right too (comma rather
% than em-dash separator).
% \changes{2.2beta12}{2001/12/07}{\cs{addkeyword} New user interface
%   for keywords.} 
% \changes{2.2beta13}{2002/03/11}{\cs{addkeyword} Improved behaviour
%   with .ads file. Use
%   \cs{ADS@sanitize} to change definitions of `~' and \cs{ion}.}
% \changes{3.23}{2007/09/09}{\cs{addkeyword} now signals a compilation
% error if the author tries to add more than 6 keywords.}
% \changes{3.24}{2007/09/14}{kwdcounter is now reset in
% \cs{@begofarticlehook} to prevent error in book style} 
%    \begin{macrocode}
\newcommand{\@kwdsep}{ --- }
\newcommand{\ADS@kwdsep}{, }
\newcounter{kwdcounter}\setcounter{kwdcounter}{0}
\newcommand{\addkeyword}[1]{%
  \addtocounter{kwdcounter}{1}
  \ifthenelse{\thekwdcounter > 6}{
    \ClassError{rmaa}{Too many keywords. A maximum of 6 keywords is allowed.}{}
  }{}
  \if@firstkwd%
  \renewcommand\@keywords{\if@boldkeywords\uppercase{#1}\else #1\fi}
  {\ADS@sanitize\def\ion{\plainion}\gdef\ADS@keywords{#1}}
  \@firstkwdfalse
  \else%
  \edef\@keywords{\@keywords\@kwdsep \if@boldkeywords\uppercase{#1}\else #1\fi}%
  {\ADS@sanitize\def\ion{\plainion}\xdef\ADS@keywords{\ADS@keywords\ADS@kwdsep #1}}%
  \fi}
\newcommand{\ADS@sanitize}{%
  \def\mathrm{}
  \def\@{\ }
  \def~{\ }
  \def\langle{<}
  \def\rangle{>}
  \def\log{log }
  \def\sin{sin }
  \def\cos{cos }
  \def\tan{tan }
  \def\fd{^d.}
  \def\fh{^h.}
  \def\fm{^m.}
  \def\fs{^s.}
  \def\fdg{�.}
  \def\farcm{'.}
  \def\farcs{''.}
  \def\arcdeg{�}
  \def\,{ }
  }
\newcommand{\keywords}[1]{%
  \renewcommand{\ADS@keywords}{#1}
  \renewcommand{\@keywords}{\uppercase{#1}}
  }
\newcommand{\@keywordheader}{\rmaa@KeywordHeaderStyle{Key Words:}}
\newcommand{\@keywords}{%
  \ClassWarningNoLine{rmaa}{No keywords specified}%
  Keywords go here
  }
%    \end{macrocode}
% \begin{macro}{\shorttitle}
% \begin{macro}{\shortauthor}
% \begin{macro}{\@shorttitle}
% \begin{macro}{\@shortauthor}
% Commands for setting the running headers.
%    \begin{macrocode}
\newcommand{\shorttitle}[1]{%
  \renewcommand\@shorttitle{\MakeUpperCase{#1}}}
\newcommand{\shortauthor}[1]{%
  \renewcommand\@shortauthor{\MakeUpperCase{#1}}}
\newcommand{\@shorttitle}{SHORT TITLE GOES HERE}
\newcommand{\@shortauthor}{SHORT AUTHOR GOES HERE}
%    \end{macrocode}
% \end{macro}
% \end{macro}
% \end{macro}
% \end{macro}
% \end{macro}
% \end{macro}
% \end{macro}


% \subsubsection{The full addresses}
% \label{sec:full-addresses}
% \begin{macro}{\fulladdresses}
% \begin{macro}{\@fulladdresses}
% \begin{macro}{\@fulladdheight}
% \begin{macro}{\@fulladddepth}
% \begin{macro}{\@fulladdbox}
% Aaghh, the full addresses! A perennial source of problems. 
%    \begin{macrocode}
\newlength{\@fulladdheight}
\newlength{\@fulladddepth}
\newsavebox{\@fulladdbox}
%% \fulladdresses command modified (11 May 2000) so that we can put
%% \fulladdresses{} in the article to suppress the output of the
%% addresses
\newcommand{\@FinalFigure}{\relax}
\newcommand{\FinalFigure}[1]{\renewcommand{\@FinalFigure}{#1}}
\newcommand{\fulladdresses}[1]{%
  \renewcommand{\@rmaaArg}{#1}
  \ifx\@rmaaArg\@rmaaEmpty
     \renewcommand{\@fulladdresses}{}
  \else
     \renewcommand{\@fulladdresses}{
       \normalsize%
       \ifx\@FinalFigure\relax\else\vfill\@FinalFigure\vfill\fi
%%       \vspace*{1ex}
       %% do addresses as a list (copied from thebibliography environment) 
       \list{\null}{\leftmargin 1.5em\topsep\z@\labelwidth\z@\labelsep\z@\itemindent -1.5em
         \itemsep\z@\parsep\z@\usecounter{enumi}}
       \sloppy\clubpenalty4000\widowpenalty4000
       \sfcode`\.=1000\relax
       #1
       \endlist
       }
  \fi
  }
\newcommand{\@fulladdresses}{Full addresses go here}
%    \end{macrocode}
% \begin{macro}{\suppressfulladdresses}
% New user-level command for suppressing full addresses should you so
% desire. Currently used in the single-page-poster style. 
%    \begin{macrocode}
\newcommand{\suppressfulladdresses}{\@suppressfulladdressestrue\@needaddressesfalse}
%    \end{macrocode}
% \end{macro}
% \end{macro}
% \end{macro}
% \end{macro}
% \end{macro}
% \end{macro}

% \begin{macro}{\printdates}
% \begin{macro}{\noprintdates}
% Command to activate the printing out of the received/accepted
% dates. This can also be done using the \texttt{printdates} option,
% which is the default for \texttt{journal}, but not for \texttt{proceedings}.
%    \begin{macrocode}
\newcommand{\printdates}{\@printdatestrue}
\newcommand{\noprintdates}{\@printdatesfalse}
%    \end{macrocode}
% \end{macro}
% \end{macro}
% \begin{macro}{\@outputaddressesifnecessary}
% Command to output the addresses if we haven't already done so
% This is called by: 
% \begin{enumerate}
% \item The end of |thebibliography| environment
% \item The |\outputfulladdresses| command
% \end{enumerate}
%    \begin{macrocode}
\newcommand{\@outputaddressesifnecessary}{%
  \ClassInfo{rmaa}{Calling @outputaddressesifnecessary}
  \ifthenelse{\boolean{@needaddresses}}{%
    %% If in one column style, or in the first column of a two column
    %% style then we can try to put it on this page 
    \ifthenelse{\boolean{@firstcolumn} \or \not\boolean{@rmtwocol}}{%
      \@outputaddresses}{%
      %% otherwise give warning it will be on a separate page
      \ClassWarning{rmaa}{Full addresses will go on a separate
        page (please see User Guide for how to remedy this)}}%
    %% If we haven't managed yet, put it on its own page
    \ifthenelse{\boolean{@needaddresses}}{%
      \if@debug{\ClassWarning{rmaa}{In this silly bit!}}\fi
      \clearpage\vfill\@outputaddresses}{}%
    }{
    %% If we don't need them, we still might need to do the eoa hook 
    \ifthenelse{\boolean{@suppressfulladdresses}}{%
      \if@debug{\ClassWarning{rmaa}{XXX}}\fi
%%      \if@book\else\@endofarticlehook\null\fi
      }{}
    }%
  }  
%    \end{macrocode}
% \end{macro}

% \begin{macro}{\outputfulladdresses}
% A user interface to the above command, in case there is no biblio
% (11 May 2000) 
%    \begin{macrocode}
\newcommand{\outputfulladdresses}{\@outputaddressesifnecessary}
%    \end{macrocode}
% \end{macro}

% \begin{macro}{\adjustfinalcols}
% This command starts a new column and should be 
% used to ensure that the columns balance on the last page. It also
% outputs the full 
% author addresses to ensure that this occurs in the left hand
% column. 
% The |\par| command is so that the current line will be taken account
% of when working out |\pagetotal|
%    \begin{macrocode}
\newcommand{\adjustfinalcols}{\if@twocolumn\goodbreak\@outputaddresses\newpage\noindent\fi}
%    \end{macrocode}
% \end{macro}

% \begin{macro}{\@outputaddresses}
% Command to do the work of outputting the addresses. 
% This can be called by 
% \begin{enumerate}
% \item an explicit |\adjustfinalcols|
% \item via |\@outputaddressesifnecessary|, q.v.
% \end{enumerate}
%
%    \begin{macrocode}
\newcommand{\@outputaddresses}{%
  % Allow option of suppressing the full addresses, in which case
  % we just run the hook if we are not in a book style
  \if@suppressfulladdresses
%%     \if@book\else\@endofarticlehook\null\fi
  \else
  %% save the full addresses in a box for output later
  \sbox{\@fulladdbox}{\parbox{\textwidth}{\@fulladdresses}}
  % calculate height of address box  
  \settoheight{\@fulladdheight}{\usebox{\@fulladdbox}}
  \settodepth{\@fulladddepth}{\usebox{\@fulladdbox}}
  \addtolength{\@fulladdheight}{\@fulladddepth}
  \if@debug
  { \typeout{Address total height: \the\@fulladdheight}
    \typeout{Text height: \the\pagetotal}
    \typeout{Goal height: \the\pagegoal}
    }
  \fi
  % reuse \@fulladdheight to mean the total height of text+addresses
  \addtolength{\@fulladdheight}{\pagetotal}
%%  \addtolength{\@fulladdheight}{-\baselineskip}
  \if@debug\typeout{Adjusted text + addresses height: \the\@fulladdheight}\fi
  % check it will fit on current page
  \ifdim\@fulladdheight<\pagegoal
     %% must shift it to the left
     %% to compensate for the hanging indentation in the reference list. 
     %% NOTE: this is not foolproof!
     %% NB. This will give overfull hbox warnings, which is a bit ugly.
     \vfill\noindent\if@biblist\hspace*{-\leftmargin}\fi
     \usebox{\@fulladdbox}
     % Set this to show we've done it
     \@needaddressesfalse
     % and assume we are at the end of the article
     % so that we do the EOA hook if not in the book style
%%     \if@book\else\@endofarticlehook\fi
  \else
     % if not, just do nothing. With luck, this function will get 
     % called again on the next page
     \if@debug{\typeout{Holding over full addresses to next page}}\fi
  \fi\fi
  }
%    \end{macrocode}
% \end{macro}

% \subsubsection{Commands for editors}
% \label{sec:editor-commands}
% Some user/editor commands for doing the TOC and index in the book
% option. 
% They are defined here rather in the book-specific part since they
% need to be always available in order to avoid breaking documents. 
% \begin{macro}{\listofauthors}
% \begin{macro}{\@rawauthors}
% \begin{macro}{\rmaaTOCauthors}
% The command |\listofauthors| formats an author list for use in the
% TOC. It should contain the same names as in the |\author| command,
% but without the affiliation garbage and using |\&| instead of
% |and|. 
%    \begin{macrocode}
\newcommand{\@rawauthors}{Raw Authors}
\newcommand{\rmaaTOCauthors}{\noexpand\rmaaTOCAuthorStyle{Author list for TOC}}
\newcommand{\listofauthors}[1]{%
  \renewcommand{\@rawauthors}{#1}
  \renewcommand\rmaaTOCauthors{\noexpand\rmaaTOCAuthorStyle{#1}}}
%    \end{macrocode}
% \begin{macro}{\indexauthor}\label{indexauthor}
% The command |\indexauthor| produces an index entry for each
% author. 
% This command must come \emph{after} the
% |\title| and the |\listofauthors| commands. Also, there should be one 
% |\indexauthor| command for \emph{each} author. We also add the
% author onto the end of the |\ADS@authors| macro and construct the
% |\ADS@authorinitial| by a most cunning method (see below), but only
% for the first author. 
% 
% No, No, NO! This was all wrong. Backed out completely: 
% \begin{quote}
% Completely rewritten (19 Dec 2002) since we now use the index class, so 
% as to be able to alphabetize accented characters. A side effect of this is
% that |\index| no longer expands out the argument after the |@|, so we 
% instead add all the |\index| commands to a list (|\indexauthor@deferlist|), 
% which is executed in the |\@begofarticlehook|. 
% 
% Note the stuff with |\toks0| is copied from Exercise 20.15 of the \TeX{}book---it is
% to make sure the latin1-accented name gets passed into the idx
% file. 
% \end{quote}
% We now go back to the old way of doing things, using an optional
% argument to the \cs{indexauthor} command that has the name without
% accents. 
% \changes{3.08}{2003/01/16}{Remerged changes from 3.06 that got
% dropped from 3.07 after a synchronization disaster. Basically, the
% accentless optional argument to \cs{indexauthor}. Must add the macro
% directories to my unison set-up to stop this sort of thing from happening.}
%    \begin{macrocode}
\newcommand\@isortkey{}
\newcommand\@emptykey{empty}
\newcommand\indexauthor[2][empty]{
  \renewcommand\@isortkey{#1}
  \ifx\@isortkey\@emptykey \renewcommand\@isortkey{#2}\fi
  \index{\@isortkey @{{\rmaaIndexAuthorStyle{#2}}\hskip1em plus1em minus0.4em
      \rmaaIndexTitle\hskip1em plus2em minus0.4em\penalty-100
      \rmaaTOCauthors}}
  \indexauthorADS{#2}
}
%%\def\indexauthor@deferlist{\relax}
%%\long\def\indexauthor#1{
%%  \def\@@name{#1} \toks0=\expandafter{\@@name}
%%  \toks1=\expandafter{\indexauthor@deferlist}
%%  \protected@xdef\indexauthor@deferlist{%
%%    \the\toks1
%%    \noexpand\index{\the\toks0@{\noexpand\rmaaIndexAuthorStyle{\the\toks0}\noexpand\hskip1em plus1em minus0.4em 
%%        \rmaaIndexTitle\noexpand\hskip1em plus2em minus0.4em\noexpand\penalty-100
%%        \rmaaTOCauthors}}
%%  }
%%  \indexauthorADS{#1}
%%  }
%%\xdef\indexauthor@deferlist{}
\newcommand\indexauthorADS[1]{%
  \if@firstauthor\protected@xdef\ADS@authorinitial{\firstletter{#1}}\fi
  \protected@xdef\ADS@authors{\ADS@authors\if@firstauthor\else ; \fi
    #1}
  \@firstauthorfalse
}
%    \end{macrocode}
% Note that the first |#1| sets the key for the index, making sure it
% actually starts with  
% first letter of the surname. Meanwhile, the |@{...}| expression,
% which starts with some  
% control sequences is what is written to the |.ind| file. 
% See the makeindex man page for details. 
% \end{macro}
% \end{macro}
% \end{macro}
% \end{macro}

% \begin{macro}{\KillAllAccents}
% \begin{macro}{\firstletter}
% These are utility functions used in the above. |lett| is a counter
% indicating the current letter in the string. |\firstletter| just
% prints the first letter of its argument.
%    \begin{macrocode}
\newcommand{\KillAllAccents}{\def\'{}\def\`{}\def\~{}\def\.{}\def\i{i}}
\newcommand\firstletter[1]{\@car #1\@nil}
%    \end{macrocode}
% \end{macro}
% \end{macro}
% \begin{macro}{\SetVolume}
% \begin{macro}{\SetMSnumber}
% \begin{macro}{\SetFirstPage}
% \begin{macro}{\SetYear}
% \begin{macro}{\ReceivedDate}
% \begin{macro}{\AcceptedDate}
% These are for the editors to set the volume number, etc. 
%    \begin{macrocode}
\newcommand{\SetVolume}[1]{%
  \renewcommand{\@volume}{#1}
  \edef\ADS@volume{\PadToFourWithDots{#1}}
  \renewcommand{\@volumen}{\rmaa@VolumenStyle{volumen #1}}%
  }
\newcommand{\SetMSnumber}[1]{\renewcommand{\@MSnum}{#1}}
%    \end{macrocode}
% Note that there is no real reason to use |\SetFirstPage| since all
% the page numbering is done automatically in the book style. 
%    \begin{macrocode}
\newcommand{\SetFirstPage}[1]{\setcounter{page}{#1}}
\newcommand{\SetYear}[1]{%
  \renewcommand{\@year}{#1}
  \renewcommand{\ADS@year}{#1}
  \renewcommand{\@anyo}{\rmaa@VolumenStyle{#1}}
  }
\newcommand{\ReceivedDate}[1]{\renewcommand{\@received}{#1}}
\newcommand{\AcceptedDate}[1]{\renewcommand{\@accepted}{#1}}
%    \end{macrocode}
%
% \begin{macro}{\DOI}
% User/editor command to enter the DOI, which will be typeset as a link.
% \changes{4.0}{2019/11/02}{New user command \cs{DOI} to typeset
% DOI in the header of first page and add it to ADS info.}
%    \begin{macrocode}
\newcommand\DOI[1]{%
  \def\@DOI@url{\url{https://doi.org/10.22201/ia.01851101p.#1}}%
  \def\ADS@doi{https://doi.org/10.22201/ia.01851101p.#1}%
}%
%    \end{macrocode}
% \end{macro}
%
% \begin{macro}{\ORCIDiD}
% \changes{4.0}{2019/11/12}{Add mechanism for giving ORCID iD links}
% \begin{macro}{\ORCIDiDfull}
% \begin{macro}{\@orcid@icon}
% \begin{macro}{\@orcid@iconbig}
% Commands to add icon links to ORCID iDs of authors. For example
% \begin{quote}
% |\ORCIDiD|\marg{0000-0000-0000-0000}
% \end{quote}
% to be used after each author name in the |\author|
% command and in the |\fulladdresses|.  Alternatively,
% |\OORCIDiDfull| will typeset the icon plus the full text of the
% link, but we don't use that at present.  The icon image is stored in |ORCIDiD_iconvector.pdf|, which needs to be on the input path.  In principle, we could make EPS versions if necessary.
%    \begin{macrocode}
\newcommand\@orcid@icon{\includegraphics[height=1.8ex]{ORCIDiD_iconvector}}
\newcommand\@orcid@iconbig{\includegraphics[height=2ex]{ORCIDiD_iconvector}}
\newcommand\ORCIDiD[1]{\href{https://orcid.org/#1}{\,\@orcid@icon}}
\newcommand\ORCIDiDfull[1]{%
  \href{https://orcid.org/#1}{\@orcid@iconbig}\,\url{https://orcid.org/#1}}
%    \end{macrocode}
% \end{macro}
% \end{macro}
% \end{macro}
% \end{macro}
%
% \begin{macro}{\@volume}
% \begin{macro}{\@volumen}
% \begin{macro}{\@artpages}
% \begin{macro}{\@year}
% \begin{macro}{\@anyo}
% \begin{macro}{\@received}
% \begin{macro}{\@accepted}
% \begin{macro}{\@MSnum}
% Default values for the containers used by above commands, plus a
% few more useful ones. 
%    \begin{macrocode}
\newcommand{\@volume}{00}
\newcommand{\@volumen}{\rmaa@VolumenStyle{00}}
\newcommand{\@artpages}{%
  \oldpageref{\articlename:RMAAFirstPage}--\oldpageref{\articlename:RMAALastPage}}
\newcommand{\@year}{\the\year}
\newcommand{\@anyo}{\rmaa@VolumenStyle{2007}}
\newcommand{\@journalname}{\rmaa@JournalNameStyle{%
    \if@proceedings\RMAASClong\else\RMAAlong\fi}}
\newcommand{\@received}{}
\newcommand{\@accepted}{}
\newcommand{\@MSnum}{0000}
\newcommand{\@DOI@prefix}{https://doi.org/10.22201/ia.01851101p}
\def\@DOI@url{\url{https://doi.org/10.22201/ia.01851101p.2099.00.00.00}}
%    \end{macrocode}
% \end{macro}
% \end{macro}
% \end{macro}
% \end{macro}
% \end{macro}
% \end{macro}
% \end{macro}
% \end{macro}
% \end{macro}
% \end{macro}
% \end{macro}
% \end{macro}
% \end{macro}
% \end{macro}
% 
% The following commands are most useful in the \texttt{book}
% style. However, the |\SetConfTitle| command can now be used in
% individual articles to set the header for the \texttt{preprint}
% style. 
%    \begin{macrocode}
\newcommand{\@reunionplace}{Lugar de la reuni\'on}
\newcommand{\@reuniondate}{Fecha de la reuni\'on}
\newcommand{\@conftitle}{Conference Title}
\newcommand{\SetReunion}[2]{\renewcommand{\@reunionplace}{#1}\renewcommand{\@reuniondate}{#2}}
\newcommand{\SetConfTitle}[1]{\renewcommand{\@conftitle}{#1}}
\newcommand{\@editors}{Juan Ram\'on de la Fuente and ``El Mosh''}
\newcommand{\SetEditors}[1]{\renewcommand{\@editors}{#1}}
\newcommand{\@editores}{Juan Ram\'on de la Fuente y ``El Mosh''}
\newcommand{\SetEditores}[1]{\renewcommand{\@editores}{#1}}
\newcommand{\@mes}{mes de la publicaci\'on}
\newcommand{\SetMes}[1]{\renewcommand{\@mes}{\rmaa@VolumenStyle{#1}}}
%    \end{macrocode}

% \subsection{Notes for Author Proofs and Communication between Editors}
% \label{sec:notes}
% We introduce a command \cs{Note} that should be used to describe
% any change made to the manuscript. These are a bit tricky to get
% right. In general we use \cs{marginpar} but this only works in outer
% paragraph mode, so we just bung the box where we are if we are in a
% figure caption or something. If we are in the preamble, then we just
% use a \marginpar as part of the \maketitle command.
%    \begin{macrocode}
\newif\if@arenotes
\@arenotesfalse
\newcounter{notecounter}
\setcounter{notecounter}{0}
\newcommand\NotesList{}
\newcommand\Note[1]{%
  \if@proof{%
    \@arenotestrue%
    \stepcounter{notecounter}%
    \if@inpreamble\else{%
      \NoteMark{\arabic{notecounter}}%
    }\fi
    \protected@xdef\NotesList{\NotesList 
      \protect\item #1}%
  }\fi%
}
\newcommand\Changed[2]{Changed ``#1'' to ``#2''.}
\newcommand\NoteMark[1]{%
  \ifinner
  \raisebox{1ex}[0pt][0pt]{\makebox[0pt][l]{\framebox{\footnotesize\bfseries
        N#1}}}%
  \else
  \marginpar{\framebox{\footnotesize\bfseries N#1}}%
  \fi
}
\newcommand\RefNote[1]{\if@proof\NoteMark{\ref{#1}}\relax\fi}
\newcommand\NoteLabel{\protect\label}
\newcommand\OutputNotes{%
  \bgroup
  \onecolumn
  \section*{EDITORIAL NOTES}
  \begin{enumerate}
    \NotesList
  \end{enumerate}
  \egroup
}
\DeclareRobustCommand\cs[1]{\texttt{\char`\\#1}}

%    \end{macrocode}
% \subsection{Bibliographic Records for ADS}
% \label{sec:ADS}
% \changes{2.2beta11}{2001/11/27}{Just about finished all the ADS
% stuff now. Except for the references, which I have given up on for
% the time being.}
% \begin{macro}{\ADS@BibCode}
% \begin{macro}{\ADS@year}
% \begin{macro}{\ADS@journal}
% \begin{macro}{\ADS@volume}
% \begin{macro}{\ADS@firstpage}
% \begin{macro}{\ADS@authorinitial}
% We use the |\ADS@BibCode| command to hold the bibliographic code for
% ADS. It is 
% used in the |%R| 
% field of the bibliographic record for each article. 
% The format is 
% \begin{quote}
% \texttt{YYYYJJJJJVVVVMPPPPA}  
% \end{quote}
% where \texttt{YYYY} is the year, \texttt{JJJJJ} is the journal code,
% which is \texttt{RMxAA} or \texttt{RMxAC}, \texttt{VVVV} is the
% volume number (right-justified), \texttt{M} is not used by us,
% \texttt{PPPP} is the first page number (right-justified), and
% \texttt{A} is the initial of the first author's surname. The volume
% and page fields are padded with periods (`\texttt{.}'), which causes
% us some fun (this is done with the macro |\PadToFourWithDots| described
% below) . 
% 
% All these individual bits are held in the commands |\ADS@year|,
% |\ADS@journal|, |\ADS@volume|, |\ADS@firstpage|,
% |\ADS@authorinitial|. 
%    \begin{macrocode}
\newcommand{\SetADSBibCode}[1]{\renewcommand{\ADS@BibCode}{#1}}
\newcommand{\ADS@BibCode}
{\ADS@year\ADS@journal\ADS@volume.\PadToFourWithDots{\ADS@firstpage}\ADS@authorinitial}
\newcommand{\ADS@year}{2001} 
\newcommand{\ADS@journal}{\if@proceedings RMxAC\else RMxAA\fi}
\newcommand{\ADS@volume}{....} 
%    \end{macrocode}
% These two are just set to default values here, |\ADS@firstpage| is set
% in the |document| environment for included articles, in the
% |frontsection| environment, |\tableofcontents| command,  |index| and
% |abstracts| environments. Also by |\AtBeginDocument| for stand-alone
% articles. 
% 
% God knows how we are going to do |\ADS@authorinitial|!
%    \begin{macrocode}
\newcommand{\ADS@firstpage}{....}
\newcommand{\ADS@authorinitial}{A}
%    \end{macrocode}
% \end{macro}
% \end{macro}
% \end{macro}
% \end{macro}
% \end{macro}
% \end{macro}
% We also have a whole slew of other fields that can be set. These are
% all listed in Table~\ref{tab:ADS}. Most of them still have to be
% defined. 
% \begin{table}[!ht]\raggedleft
% \caption{ADS Bibliographic Records}\label{tab:ADS}
% \begin{tabular}{lp{8em}lp{8em}}\hline
% Code & Description & Macro Name & Comment \\ \hline
% |%R| & Bibliographic Code & |\ADS@BibCode| & required \\
% |%A| & Author list & |\ADS@authors| & required, generate from |\indexauthor| \\
% |%F| & Author affiliation & |\ADS@affil| & format: AA(\dots)
% AB(\dots) \dots \\
% |%J| & Journal Name & |\ADS@journalname| &  \\
% |%D| & Publication Date & |\ADS@fulldate| &  \\
% |%L| & Last Page & |\ADS@lastpage| & generate from label \\
% |%T| & Title & |\ADS@title| & required, generate from |\title| \\
% |%C| & Abstract Copyright & |\ADS@copyright| & set to IA-UNAM \\
% |%U| & URL for electronic document & |\ADS@url| & set to PDF file \\
% |%d| & DOI for article & |\ADS@doi| & set by |\DOI| \\
% |%K| & Keywords & |\ADS@keywords| & set equal to |\@keywords| \\
% |%A| & Abstract Text & |\ADS@abstract| & set equal to |\@abstract| \\
% |%R| & References & |\ADS@references| & generate from the
% bibliography  \\ \hline
% \end{tabular}
% \end{table}
% A lot of these macros are set by other commands, as mentioned in the
% Table, but some need to be set  explicitly here. 
% 
% \begin{macro}{\ADS@doi}
% \changes{4.0}{2019/11/08}{Save DOI for ADS}
% The user/editor command \cs{DOI} is used to set this. 
%    \begin{macrocode}
\newcommand{\ADS@doi}{https://doi.org/10.22201/ia.01851101p.0000.00.00.00}
%    \end{macrocode}
% \end{macro}
% 
% \begin{macro}{\ADS@journalname}
% This
% does not contain any accents at the moment so as not to cause
% problems---this maybe should be changed though. 
%    \begin{macrocode}
\newcommand{\ADS@journalname}{%
  \if@proceedings\@conftitle\space (Eds. \@editors) \fi%
  Revista Mexicana de Astronom\'\i{}a y Astrof\'\i{}sica%
  \if@proceedings\space (Serie de Conferencias)\fi%
  \space Vol. \@volume, pp. \ADS@firstpage--\ADS@lastpage\space (\@year)%
  \space (http://www.astroscu.unam.mx/\~{}rmaa/)
  }
%    \end{macrocode}
% \end{macro}
% \begin{macro}{\ADS@keywords}
% \begin{macro}{\ADS@abstract}
% \begin{macro}{\ADS@title}
% |\ADS@abstract| is straightforward, since we already have the text
% saved as |\@abstract|. We could do the same with |\ADS@keywords| but
% to avoid unwanted formatting we explicitly set it in |\keywords|
% (p.~\pageref{keywords}). |\ADS@title| is just set to |\@rawtitle|,
% which we were already using for another purpose. 
%    \begin{macrocode}
\let\ADS@title=\@rawtitle
\newcommand{\ADS@keywords}{No keywords specified}
\newcommand{\ADS@abstract}{No abstract specified}
%    \end{macrocode}
% \end{macro}\end{macro}\end{macro}
% \begin{macro}{\ADS@authors}
% |\ADS@authors| wants to be in ``Other, A. N.'' format, but with
% semi-colons between authors. The most natural way to do this seems to
% be with |\indexauthor| (p.~\pageref{indexauthor}). We append each
% author to this macro, which means we must explicitly clear it for
% each new article (done in |\article|). .   
%    \begin{macrocode}
\newcommand{\ADS@authors}{}
%    \end{macrocode}
% \end{macro}
% \begin{macro}{\ADS@fulldate}
% \begin{macro}{\ADS@lastpage}
% \begin{macro}{\ADS@copyright}
% \begin{macro}{\ADS@website}
% \begin{macro}{\ADS@url}
% \begin{macro}{\SetWebsite}
% \begin{macro}{\SetMonthNumeric}
% Other sundry data, together with user commands for bits that can be
% set by the Editors. The |\ADS@url| is constructed from the website
% (currently |http://www.astrosmo.unam.mx/rmaa/|,
% the journal descriptor (e.g., |RMxAC..10/|), followed by the
% path to the PDF file (e.g., |PDF/RMxAC..10_HEW1.pdf|). 
% 
%    \begin{macrocode}
\newcommand{\ADS@website}{http://www.astrosmo.unam.mx/rmaa/}
\newcommand{\ADS@url}
{\ADS@website\ADS@journal\ADS@volume\ADS@number/PDF/\ADS@journal\ADS@volume\ADS@number\_\articlename.pdf}
\newcommand{\ADS@number}{}
\newcommand{\ADS@fulldate}{00/\ADS@year} 
\newcommand{\ADS@lastpage}{00} 
\newcommand{\ADS@copyright}{(C) \ADS@year; Instituto de Astronom\'\i{}a,
  Universidad Nacional Aut\'onoma de M\'exico} 
\newcommand{\SetNumber}[1]{\renewcommand{\ADS@number}{-#1}}
\newcommand{\SetWebsite}[1]{\renewcommand{\ADS@website}{#1}}
\newcommand{\SetMonthNumeric}[1]{\renewcommand{\ADS@fulldate}{#1/\ADS@year}}
%    \end{macrocode}
% \end{macro}\end{macro}\end{macro}\end{macro}\end{macro}\end{macro}\end{macro}
% \begin{macro}{\ADS@affil}
% \begin{macro}{\ADS@references}
% These two are somewhat more problematical. Affiliation currently
% only works for the first author (see \dots) pending redesign of the
% user interface. References are now done by wrapping the
% |thebibliography| environment to open another input stream on the
% same file and search through for the |\bibitems|. 
%    \begin{macrocode}
\newcommand{\ADS@references}{}
\newcommand{\ADS@affil}{}
%    \end{macrocode}
% \end{macro}\end{macro}
% 
% Define two macros that expand to the strings
% |\begin{thebibliography}| and |\end{thebibliography}|. Note that the
% idea for this and much else is inspired by the code for the docstrip
% package. Also define |\StartBibX| in case authors use the more
% standard |\begin{thebibliography}{}| syntax. 
%    \begin{macrocode}
\edef\EndBib{\expandafter\@gobble\string\\end\string{thebibliography\string}}
\edef\StartBib{\expandafter\@gobble\string\\begin\string{thebibliography\string}}
\edef\StartBibX{\expandafter\@gobble\string\\begin\string{thebibliography\string}\string{\string}}
%    \end{macrocode}
% Define two new logical variables. |\ifAnotherline| is true if we
% want to read another line from the input stream. It is set to false
% when we encounter the |\end{thebibliography}|. |\ifInBib| is set to
% true when we frind the |\begin{thebibliography}| and tells us to
% start saving the following lines. 
%    \begin{macrocode}
\newif\ifAnotherLine
\newif\ifInBib
%    \end{macrocode}
% \begin{macro}{makeOther}
% We need a macro to change the catcodes of lots of things to other. 
%    \begin{macrocode}
\def\makeOther#1{\catcode`#1=12\relax}
%    \end{macrocode}
% \end{macro}
% \begin{macro}{\RereadBiblio}
% Define a macro to do all the fiddly stuff. This is used at the
% beginning of |thebibliography| environment (Sec~\ref{sec:biblio}).
%    \begin{macrocode}
\newread\tf@biblio
\newcommand{\RereadBiblio}{%
%    \end{macrocode}
% Open the current file again for reading
%    \begin{macrocode}
  \begingroup
  \immediate\openin\tf@biblio=\articlename
%    \end{macrocode}
% Everything in a group to be safe. Initialize logicals. Black magic
% to change
% catcodes and avoid blank lines being interpreted as |\par|. 
%    \begin{macrocode}
  \AnotherLinetrue\InBibfalse
  \makeOther\ \makeOther\\\makeOther\$%
  \makeOther\#\makeOther\^\makeOther\^^K%
  \makeOther\_\makeOther\^^A\makeOther\%%
  \makeOther\~\makeOther\{\makeOther\}\makeOther\&%
  \endlinechar-1\relax
%    \end{macrocode}
% Start looping through the file. Read each line to the macro
% |\Line|. If we have got to EOF or end of biblio then set
% |\AnotherLinefalse|.
%    \begin{macrocode}
  \loop
  \read\tf@biblio to \Line
  \ifx\Line\EndBib\AnotherLinefalse\InBibfalse\fi
  \ifeof\tf@biblio \AnotherLinefalse\fi
%    \end{macrocode}
% If we are in the biblio and the line is not blank then append it to
% the |\ADS@references| macro together with a newline. 
%    \begin{macrocode}
  \ifInBib
  \ifx\Line\empty\else\xdef\ADS@references{\ADS@references\Line^^J}\fi
  \fi
%    \end{macrocode}
% Check if the line we just saw was the start of biblio, then repeat
% the loop if necessary. 
%    \begin{macrocode}
  \ifx\Line\StartBib\InBibtrue\fi
  \ifx\Line\StartBibX\InBibtrue\fi
  \ifAnotherLine\repeat
  \immediate\closein\tf@biblio
  \endgroup
  }
%    \end{macrocode}
% \end{macro}
% 
% \begin{macro}{\PadToFourWithDots}
% Utility function for left-padding a number to 4 characters with
% periods. Used in 
% constructing |\ADS@BibCode|. This is a very simple algorithm and
% \emph{requires} that the argument actually be a number---if it is
% not then dire consequences will result.  
%    \begin{macrocode}
\newcommand{\PadToFourWithDots}[1]{%
  \ifnum #1 > 999 #1%
  \else\ifnum #1 > 99 .#1%
  \else\ifnum #1 > 9 ..#1%
  \else ...#1%
  \fi\fi\fi}
%%  \ifnum \count0 > 999 #1%
%%  \else\ifnum \count0 > 99 .#1%
%%  \else\ifnum \count0 > 9 ..#1%
%%  \else ...#1%
%%  \fi\fi\fi}
%    \end{macrocode}
% \end{macro}
%
% \begin{macro}{\AuthorADS}
% \begin{macro}{\ADS@affiliations}
% New mechanism for setting the affiliations for ADS. The author labels (AA, AB, \dots) are generated using two nested counters.  We use the |\robustify| command from \texttt{etoolbox} so that accented characters in the affiliation addresses do not blow up.  The list of accents may need to be expanded. 
%    \begin{macrocode}
\newcounter{AUTHOR}
\newcounter{author}[AUTHOR]
\stepcounter{AUTHOR}
\newcommand\ADS@affiliations{}
\robustify{\'}\robustify{\i}\robustify{\~}\robustify{\`}
\def\AuthorADS#1#2{%
  \begingroup
  \def\ORCIDiD##1{<ID system="ORCID">##1</ID>}
  \def\Email##1{<EMAIL>##1</EMAIL>}
  \def\Affil##1{##1}
  \stepcounter{author}
  \ifthenelse{\value{author} > 26}{\stepcounter{AUTHOR}\stepcounter{author}}{}
  \ifthenelse{\equal{\ADS@affiliations}{\@rmaaEmpty}}{%
    \xdef\ADS@affiliations{\Alph{AUTHOR}\Alph{author}(#2)}
  }{%
    \xdef\ADS@affiliations{\ADS@affiliations, \Alph{AUTHOR}\Alph{author}(#2)}
  }%
  \endgroup
}
%    \end{macrocode}
% \end{macro}
% \end{macro}
% %
% \begin{macro}{\ADS@dumpdata}
% \changes{4.0}{2019/11/13}{Write all author affiliations for ADS}
% \begin{macro}{\ADS@write}
% \begin{macro}{\ADS@openfile}
% |\ADS@dumpdata| is the macro that actually does the job of writing out all the
% data to an auxiliary file, using |\ADS@write| to write each
% field. The file is set up with |\ADS@openfile|, which currently is
% done in the TOC of a book.  The new multi-author affiliations are used if |\ADS@affiliations| has been set, otherwise it falls back to the single-author version |\ADS@affil|. 
%    \begin{macrocode}
\newcommand{\ADSfield}[2]{\relax}
\newcommand{\ADSvfield}[2]{\relax}
%%\renewcommand{\ADSfield}[12]{\%#1 #2}
\newcommand{\ADS@dumpdata}
{ 
  \ADS@write{\string\begin{ADSarticle}}
  \ADS@WRITE{\string\ADSfield{R}{\ADS@BibCode}}
  \ADS@write{\string\ADSfield{A}{\ADS@authors}}
  \ifthenelse{\equal{\ADS@affiliations}{\@rmaaEmpty}}{%
    {\ADS@sanitize\ADS@write{\string\ADSfield{F}{AA(\ADS@affil)}}}
  }{%
    {\ADS@sanitize\ADS@write{\string\ADSfield{F}{\ADS@affiliations}}}
  }
  \ADS@write{\string\ADSfield{J}{\ADS@journalname}}
  \ADS@write{\string\ADSfield{D}{\ADS@fulldate}}
  \ADS@write{\string\ADSfield{L}{\ADS@lastpage}}
  {\ADS@sanitize\ADS@write{\string\ADSfield{T}{\ADS@title}}}
  \ADS@write{\string\ADSfield{C}{\ADS@copyright}}
  \ADS@write{\string\ADSfield{d}{\ADS@doi}}
  \ADS@write{\string\ADSfield{U}{\ADS@url}}
  \ADS@write{\string\ADSfield{K}{\ADS@keywords}}
  {\ADS@sanitize\ADS@write{\string\ADSfield{B}{\ADS@abstract}}}
  \ADS@write{\string\ADSfield{Z}{\ADS@references}}
  \ADS@write{\string\end{ADSarticle}}
  }
%%\newcommand{\ADS@write}[1]{\if@book\addtocontents{ads}{#1}\fi}
%%\newcommand{\ADS@openfile}{\@starttoc{ads}}
\newcommand{\ADS@write}[1]{\set@display@protect\immediate\write\tf@ads{#1}}
\newcommand{\ADS@WRITE}[1]{\immediate\write\tf@ads{#1}}
\newcommand{\ADS@openfile}{%
  \newwrite\tf@ads%
  \immediate\openout\tf@ads\jobname.ads\relax
  \ADS@WRITE{\ADS@preamble}}
\newcommand{\ADS@closefile}{\ADS@WRITE{\ADS@postamble}}
\def\ADS@preamble{%
  %%  \string\documentclass{article}
  %%  \string\begin{document}
  }
\def\ADS@postamble{%
  %%  \string\end{document}
  \string%\string% Local Variables:^J 
  \string%\string% mode: LaTeX^J
  \string%\string% mode: Font-lock^J
  \string%\string% End:^J
  }
%    \end{macrocode}
% \end{macro}
% \end{macro}
% \end{macro}
% 
% \subsection{Automatic calculation of the page range}
% \label{sec:page-range}
% Another issue that has cost me a lot of grief. Still, it finally
% seems to be sorted now. All hail Jeff Goldberg. 
%
% Cunning trick to get the page numbers from |\pageref|.  Put a label
% on the first and last page of the document using |\AtBeginDocument|
% and |\AtEndDocument|. .  This way won't work for the book option.
% In that case, we do it by defining |\@endofarticlehook| and
% |\@begofarticlehook| as parallels to |\AtEndDocument| and
% |\AtBeginDocument|.  Note that we have to defer adding
% |\@begofarticlehook| to |\AtBeginDocument| to make sure it comes
% after the redefinitions of |\ref| and |\label|.
%    \begin{macrocode}
\if@book\AtEndDocument{\ADS@closefile}\else
 %\AtBeginDocument{\@begofarticlehook}
\AtEndDocument{\@endofarticlehook}
\fi
%    \end{macrocode}
% \begin{macro}{\@begofarticlehook}
% Since we are getting more and more things that need to be done at
% the beginning of every article, frontsection, etc we define a sister
% hook to 
% |\@endofarticlehook|. This hook does the following: 
% \begin{enumerate}
% \item Labels the first page
% \item Sets |\ADS@firstpage|
% \item runs and then clears |\indexauthor@deferlist|
% \item resets the keyword counter
% \end{enumerate}
% It is executed in the following circumstances: 
% \begin{enumerate}
% \item At |\begin{document}| for standalone files
% \item At |\begin{document}| for each included file
% \item At beginning of front sections 
% \item At beginning of TOC
% \item At beginning of Author Index
% \item At beginning of abstracts
% \end{enumerate}
% For some reason, this is not always within the right group, so we
% use |\xdef| just to be sure. 
%    \begin{macrocode}
\newcommand{\@begofarticlehook}{%
  \if@frontmatter\frontlabel{RMAAFirstPage}\else\label{RMAAFirstPage}\fi%
  \xdef\ADS@firstpage{\@arabic{\c@page}}%
  \if@book\else\makeindex\fi % needed so default index is defined in stand-alone articles
%%  \indexauthor@deferlist
%%  \xdef\indexauthor@deferlist{}
  \@inpreamblefalse
  \setcounter{kwdcounter}{0}
  }
%    \end{macrocode}
% \end{macro}
% \begin{macro}{\@endofarticlehook}
% This is a hook to be run at the end of each article
% currently, it just labels the last page and writes out all the ADS
% data to an auxiliary file.  
% It should not produce any printed output. 
% It is modelled on the |\AtEndDocument| command, to which it reduces
% in the case of a single article. 
%    \begin{macrocode}
\newcommand{\@endofarticlehook}{%
  \if@debug{\ClassWarning{rmaa}{***Doing Hook***}}\fi
  \if@debug{\ClassWarning{rmaa}{ADSabstract is: \ADS@abstract}}\fi
  \clearpage\rmaa@lastpagelabel\if@book\ADS@dumpdata\fi
  \ifx\NotesList\@rmaaEmpty\else\OutputNotes\fi
  \if@debug{\ClassWarning{rmaa}{Notes list is: \NotesList}}\fi
}
%    \end{macrocode}
% \begin{macro}{\rmaa@lastpagelabel}
% New version of last page label command copied from Jeff Goldberg's
% lastpage.sty package - we can't just use |\label| here in case the
% output routine for the last page has already been called. Instead,
% we first explicitly clear the page, step back the page counter,
% explicitly write to the .aux file, then restep the page counter
% again.
%
% But this will conflict with the hyperref package, which does labels differently.  Compensate for that by using the five-argument form.
% 
% Also set the |\ADS@LastPage| macro to the appropriate number. 
%    \begin{macrocode}
\newcommand{\rmaa@lastpagelabel}{%
  \addtocounter{page}{-1}%
  \immediate\write\@auxout{\string
    \newlabel{\articlename:RMAALastPage}{{}{\thepage}{}{}{}}}%
  \edef\ADS@lastpage{\thepage}%
  \addtocounter{page}{1}}
%    \end{macrocode}
% \end{macro}
% \end{macro}

% \subsection{Deal with explicit change in number of columns}
% \label{sec:columns}
% Sometimes, when formatting a book, we need to switch between
% one-column articles and two-column articles. For this to work, we
% need to wrap the standard \LaTeX{} commands to set the |\@rmtwocol|
% flag appropriately and also reset the package default values of all
% scalable lengths. .  
%    \begin{macrocode}
\let\latexonecolumn=\onecolumn
\let\latextwocolumn=\twocolumn
\renewcommand{\onecolumn}{\@rmtwocolfalse\@package@default@lengths\latexonecolumn}
\renewcommand{\twocolumn}{\@rmtwocoltrue\@package@default@lengths\latextwocolumn}
%    \end{macrocode}
% 
% \subsection{The actual formatting of the title}
% \label{sec:title-format}

% Note, we have to redefine the |\maketitle| command to use the
% numbered footnotes. This means we needn't use the atbeginend stuff. 
%    \begin{macrocode}
\renewcommand{\maketitle}{\par
  \begingroup
%    \end{macrocode}
% Don't switch to marked footnotes. 
% This is consistent with existing style but I'm not sure I like it
%    \begin{macrocode}
  %%  \renewcommand\thefootnote{\fnsymbol{footnote}}%
%%  \def\@makefnmark{\hbox to\z@{$\m@th^{\@thefnmark}$\hss}}%
  \def\@makefnmark{\textsuperscript{\@thefnmark}\hspace{-0.2em}}%
  \long\def\@makefntext##1{\parindent 1em\noindent
    \hbox to1.8em{\hss$\m@th^{\@thefnmark}$}##1}%
  \if@rmtwocol
  { \if@debug{\ClassWarning{rmaa}{TWO COLUMN}}\fi
    \twocolumn[\@maketitle]%
    }
  \else
  { \if@debug{\ClassWarning{rmaa}{ONE COLUMN}}\fi
    \newpage
    \global\@topnum\z@   % Prevents figures from going at top of page.
    \if@manuscript \@MSmaketitle \else \@maketitle \fi
    }
  \fi
  \if@manuscript\else\thispagestyle{rmaafirstpage}\@thanks\fi
  \ifx\NotesList\@rmaaEmpty\else{%
    \NoteMark{1--\arabic{notecounter}}\vspace*{-2\baselineskip}%
  }\fi%
  \endgroup
%    \end{macrocode}
% Don't reset footnote counter (not sure whether this is good thing
% or not).
%    \begin{macrocode}
  %%  \setcounter{footnote}{0}%
  \let\thanks\relax\let\affil\relax\let\altaffil\relax
  \let\maketitle\relax\let\@maketitle\relax
  \gdef\@thanks{}\gdef\@author{}\gdef\@title{}}
%    \end{macrocode}

% \begin{macro}{\@kwhwidth}
% \begin{macro}{\@kwtwidth}
% Redefine \cs{@maketitle} to include the abstract, resumen and
% keywords. Note that in the manuscripty style we use
% \cs{@MSmaketitle} instead.  
% \changes{3.16}{2003/04/07}{
%  Added extra vertical space between authors and accepted date.}
% \changes{3.26}{2011/08/15}{Add extra label to review articles}
%    \begin{macrocode}
\newlength{\@kwhwidth}\newlength{\@kwtwidth}
\newcommand\reviewname{Review}
\renewcommand{\@maketitle}{%
  \newpage
  \null
  \if@reviewarticle
  \smallskip
  {\large\bfseries\itshape\reviewname}%
  \bigskip
  \else
  \vskip\@before@title@skip%
  \fi
  \begin{center}%
    {\@title \par}%
    \vskip\@before@author@skip%
    \parbox{\rmaa@AuthorWidthFrac\textwidth}{\centering
      \lineskip\@between@author@skip%
      \@author
      \par}
    %% this should be replaced with `Submitted to...' for preprints
    \if@printdates
    \vskip\@before@dates@skip%
    \if@preprint
    \rmaa@ReceivedStyle{Submitted to RMxAA \@received}
    \else
    \rmaa@ReceivedStyle{Received \@received; accepted \@accepted}
    \fi
    \fi
    \if@manuscript
    \rmaa@ReceivedStyle{Submitted to RMxAA \@received}
    \fi
    \if@guide
    \rmaa@ReceivedStyle{Version \rmaa@version, \rmaa@date}
    \fi
    \vskip\@before@abstract@skip%
    \if@abstractpresent
    \parbox{\rmaa@AbstractWidthFrac\textwidth}{%
      \setlength{\parindent}{0.9cm}%
      \centerline{\rmaa@AbstractHeaderStyle{Abstract}}\par%
      \vskip\@after@abstracthead@skip%
      \if@newstyle\noindent\fi\@abstract
      \vskip\@between@abstract@skip%
      \centerline{\rmaa@AbstractHeaderStyle{Resumen}}\par%
      \vskip\@after@abstracthead@skip%
      \if@newstyle\noindent\fi\@resumen
      \if@setresumen\else\@abstract\fi % append english abstract if
                                       % spanish one not set
      \if@guide\else            %do not use keywords for guide
      \vskip\@before@kwd@skip%
      % Format the Key Words with hanging indentation:
      % They are put in a \parbox with width equal to the line width
      % minus the width of the header. 
      { % Widths of keyword
        \settowidth{\@kwhwidth}{\@keywordheader~}   % header and
        \setlength{\@kwtwidth}{\linewidth}          % keyword text
        \addtolength{\@kwtwidth}{-\@kwhwidth}       %
        \noindent\@keywordheader~\parbox[t]{\@kwtwidth}{%
          \rmaa@KeywordStyle{\boldmath\uppercase{\@keywords}}}}%
      \fi
      }
      \fi
  \end{center}%
  \vspace*{\afterkwdskip}
  \par
  }
\newlength{\afterkwdskip}\setlength{\afterkwdskip}{\z@}
%    \end{macrocode}
% \end{macro}
% \end{macro}
%
%
% \begin{macro}{\reviewarticle}
% \begin{macro}{\reviewtoc}
% \begin{macro}{\reviewlof}
%   User command to specify a review article. Must be used in
%   preamble. Main effect is to place the text ``Review'' before the
%   article title.  Also add command for setting a table of contents
%   and a list of figures at the beginning of the article.
%    \begin{macrocode}
\newcommand{\reviewarticle}{\@reviewarticletrue}
\newcommand\reviewtoc{%
  \section*{CONTENTS}%
  \@starttoc{rtc}%
}
\newcommand\reviewlof{%
  \section*{LIST OF FIGURES}%
  \@starttoc{lof}%
}
%    \end{macrocode}
% \end{macro}
% \end{macro}
% \end{macro}
% \begin{macro}{\l@section}
%    \begin{macrocode}
\renewcommand*\l@section[2]{%
  \ifnum \c@tocdepth >\z@
    \addpenalty\@secpenalty
    \addvspace{0.8em \@plus\p@}%
    \setlength\@tempdima{1.5em}%
    \begingroup
      \parindent \z@ \rightskip \@pnumwidth
      \parfillskip -\@pnumwidth
      \leavevmode %\bfseries
      \advance\leftskip\@tempdima
      \hskip -\leftskip
      #1\nobreak\hfil \nobreak\hb@xt@\@pnumwidth{\hss #2}\par
    \endgroup
  \fi}
%    \end{macrocode}
% \end{macro}
%
% \begin{macro}{\@MSmaketitle}
% A special version of |\@maketitle| for the \texttt{manuscript}
% option. This can deal with a multi-page abstract/resumen. Not very
% well, mind. 
%    \begin{macrocode}
\newcommand{\@MSmaketitle}{%
  \if@debug\ClassInfo{rmaa}{Using special MS maketitle command}\fi
  \newpage
  {
    \thispagestyle{rmaafirstpage}
%     \setlength{\parindent}{0pt}
    \null\vskip 2ex
    \begin{center}
      \@title\par
      \@author\par
      \rmaa@ReceivedStyle{%
        \ifx\@received\@rmaaEmpty
          Draft version: \@date
        \else
          Received: \@received
          \ifx\@accepted\@rmaaEmpty
          \else
            ; Accepted: \@accepted
          \fi
        \fi}\par
    \end{center}
    \@thanks
    \centerline{\rmaa@AbstractHeaderStyle{Resumen}}
    \noindent\@resumen\par\vskip 2ex
    \centerline{\rmaa@AbstractHeaderStyle{Abstract}}
    \noindent\@abstract\par\vskip 2ex
    \noindent\@keywordheader~\rmaa@KeywordStyle{\@keywords}\par
%     \@keywordheader~\rmaa@KeywordStyle{\uppercase{\@keywords}}\par
  }%
  }
%    \end{macrocode}
% \end{macro}

% \subsection{The sectioning commands}
% \label{sec:section-commands}

% \begin{macro}{\section}
% \begin{macro}{\subsection}
% \begin{macro}{\subsubsection}
% \begin{macro}{\paragraph}
% \begin{macro}{\subparagraph}
% Redefine all the standard commands to use these lengths. 
%    \begin{macrocode}
\renewcommand\section{\@startsection {section}{1}{\z@}%
                                   {\abovesectionskip}%
                                   {\belowsectionskip}%
                                   %% how do we get upper case here
                                   {\centering\reset@font\normalsize}}
\renewcommand\subsection{\@startsection{subsection}{2}{\z@}%
                                     {\abovesectionskip}%
                                     {\belowsectionskip}%
                                     {\centering\reset@font\normalsize\itshape}}
\renewcommand\subsubsection{\@startsection{subsubsection}{3}{\z@}%
                                     {\abovesectionskip}%
                                     {\belowsectionskip}%
                                     {\centering\reset@font\normalsize\itshape}}
\renewcommand\paragraph{\@startsection{paragraph}{4}{\z@}%
                                    {\aboveparaskip}%
                                    {\belowparaskip}%
                                    {\reset@font\normalsize\bfseries}}
\renewcommand\subparagraph{\@startsection{subparagraph}{5}{\parindent}%
                                      {\aboveparaskip}%
                                      {\belowparaskip}%
                                      {\reset@font\normalsize\bfseries}}
%    \end{macrocode}
% \end{macro}
% \end{macro}
% \end{macro}
% \end{macro}
% \end{macro}

% \begin{macro}{secnumdepth}
% Make sure that we number the subsubsections---this is done by
% default in the |article| class but not the |book| class, which is
% why we must set it explicitly here.
%    \begin{macrocode}
\setcounter{secnumdepth}{3}
%    \end{macrocode}
% \end{macro}
% \changes{2.2beta10}{2001/11/07}{Fixed \cs{secnumdepth} for \cs{book}
% style.}

% \begin{macro}{\thesubsection}
% \begin{macro}{\thesubsubsection}
% For subsection and subsubsection, we want the numbers to be roman
% but the title to be italic. This seems the easiest way to go about it.
%    \begin{macrocode}
\renewcommand{\thesubsection}{{\upshape\thesection.\@arabic\c@subsection}}
\renewcommand{\thesubsubsection}{\thesubsection{\upshape.\@arabic\c@subsubsection}}
%    \end{macrocode}
% \end{macro}
% \end{macro}

% \begin{macro}{\@seccntformat}
% Add ``.'' to the number here so we can remove it from |\thesection|,
% etc. Reduce gap between number and title WJH 27 Jan 1995
%    \begin{macrocode}
\def\@seccntformat#1{\csname the#1\endcsname .\hskip 0.4em}
%    \end{macrocode}
% \end{macro}

% \begin{macro}{\@sect}
% I didn't really want to mess with the |\@sect| command, but it was
% the only way I could manage to force upper case. Other than the
% commented parts, it is identical to the standard version. 
%    \begin{macrocode}
\def\@sect#1#2#3#4#5#6[#7]#8{
  \ifnum #2>\c@secnumdepth
     \let\@svsec\@empty
  \else
     \refstepcounter{#1}%
     \let\@@protect\protect
     \def\protect{\noexpand\protect\noexpand}%
     \edef\@svsec{\@seccntformat{#1}}%
     \let\protect\@@protect
  \fi
%    \end{macrocode}
% This is the bit I changed (27 Jan 1995). Use upper case for
% |\section|, use mixed case for the rest. 
%    \begin{macrocode}
  \ifnum #2=1                   
     \def\sec@upcase##1{\uppercase{##1}}%
     \def\sec@shape{\upshape}%
  \else                         
     \def\sec@upcase##1{\relax{##1}}%
     \def\sec@shape{\itshape}%
  \fi
  \@tempskipa #5\relax
  \ifdim \@tempskipa>\z@
     \begingroup #6\relax
          \@hangfrom{\hskip #3\relax\@svsec}%
                    {\interlinepenalty \@M \sec@upcase{#8}\par}%
     \endgroup
     \csname #1mark\endcsname{#7}%
%    \end{macrocode}
% Only write TOC entries if doing the Guide style. 
%    \begin{macrocode}
     \if@guide 
        \addcontentsline{toc}{#1}{%
          \ifnum #2>\c@secnumdepth \else
             \protect\numberline{\csname the#1\endcsname}%
          \fi
          #7}%
     \fi
%    \end{macrocode}
% Also allow per-article TOC for review papers
%    \begin{macrocode}
     \if@reviewarticle
        \ifnum #2>2 \else
           \addcontentsline{rtc}{#1}{%
             \protect\numberline{\csname the#1\endcsname}%
             \sec@upcase{\sec@shape #7}}%
        \fi
     \fi
  \else
     \def\@svsechd{#6\hskip #3\relax  %% \relax added 2 May 90
       \@svsec #8\csname #1mark\endcsname{#7}%
       \if@guide
          \addcontentsline{toc}{#1}{%
            \ifnum #2>\c@secnumdepth \else
               \protect\numberline{\csname the#1\endcsname}%
            \fi
            #7}%
       \fi
       }%
  \fi
  \@xsect{#5}}
%    \end{macrocode}
% \end{macro}


% \subsection{Appendices}
% \label{sec:appendices}
% Make sure that appendices in individual articles work properly. We
% copy the definition from the article class (the book one doesn't do
% what we want).  
% \begin{environment}{appendices}
%   This uses an environment |appendices| inside of which each
%   |\section| will become a separate appendix.  Note that the
%   appendices must come before the references (otherwise the full
%   addresses mechanism will not work properly). We make the equation
%   numbering be subordinate to the appendix (equations A12, B6,
%   etc). Finally, we have a medskip followed by APPENDICES to
%   introduce them. \changes{3.17}{2002/04/08}{ Added appendices environment
%   that does the right thing.}
%    \begin{macrocode}
\newenvironment{appendices}{\par
  \begingroup
  \setcounter{section}{0}%
  \setcounter{subsection}{0}%
  \renewcommand\thesection{\@Alph\c@section}
  \renewcommand\theequation{\thesection\arabic{equation}}
  \medskip
  \section*{APPENDICES}}
{\endgroup}
%    \end{macrocode}
% \end{environment}
% \begin{environment}{appendix}
%   If there is only one appendix, then use |appendix| environment
%   instead.  \changes{4.0}{2019/11/12}{ Added
%   appendix environment for case where there is only one section.}
%    \begin{macrocode}
\renewenvironment{appendix}{\par
  \begingroup
  \setcounter{section}{0}%
  \setcounter{subsection}{0}%
  \renewcommand\thesection{\@Alph\c@section}
  \renewcommand\theequation{\thesection\arabic{equation}}
  \medskip
  \section*{APPENDIX}}
{\endgroup}
%    \end{macrocode}
% \end{environment}

% \subsection{Floating thingies}
% \label{sec:floats}   

%% 
%    \begin{macrocode}
\renewcommand{\thetable}{\@arabic\c@table}
\renewcommand{\thefigure}{\@arabic\c@figure}
\renewcommand{\figurename}{Fig.}
\renewcommand{\tablename}{TABLE}
%\renewcommand{\fnum@table}{\MakeUpperCase{\tablename}~\thetable}
%    \end{macrocode}
% For some reason, I have to write `TABLE' explicitly here
% If I use |\tablename|, it gets converted to mixed case---God knows
% why. 
% |\figurename| is also being ignored. 
%    \begin{macrocode}
\renewcommand{\fnum@table}{TABLE~\thetable}
\renewcommand{\fnum@figure}{Fig.~\thefigure}
%    \end{macrocode}

% Do captions in 9 pt face in the main journal style and in new-style
% anything. 
%    \begin{macrocode}
\long\def\@makefigcap#1#2{%
  \if@proceedings\else\small\fi
  \if@newstyle\small\fi 
  \vskip\abovecaptionskip
  \sbox\@tempboxa{#1. #2}%
  \ifdim \wd\@tempboxa >\hsize
    #1. #2 \par
  \else
    \global \@minipagefalse
    \hb@xt@\hsize{\hfil\box\@tempboxa\hfil}%
  \fi
  \vskip\belowcaptionskip}
%    \end{macrocode}

% Make sure table caption is upper case. 
%    \begin{macrocode}
\long\def\@maketabcap#1#2{%
  \hbox to\hsize{\hfil #1 \hfil}
  \vskip 6pt 
  \centering\MakeUpperCase{#2}\par
  \vskip 6pt}
%    \end{macrocode}

%    \begin{macrocode}
\newlength{\@rottabcapwidth}
\long\def\@makerottabcap#1#2{%
  \setbox\@tempboxa\hbox{\MakeUpperCase{#2}}%
  \ifdim\wd\@tempboxa >.8\vsize
     \setlength\@rottabcapwidth{.8\vsize}
  \else
     \setlength\@rottabcapwidth{\wd\@tempboxa}
  \fi
  \rotatebox{90}{%
    \begin{minipage}{\vsize}\centering
      \begin{minipage}{\@rottabcapwidth}\centering
        #1\\ \vskip 6pt \MakeUpperCase{#2}
      \end{minipage}%
    \end{minipage}
    }%
  \hskip 6pt
  }

\newenvironment{Sideways}
    {\begin{sideways}\begin{minipage}{\vsize}\centering}
    {\end{minipage}\end{sideways}}

\newenvironment{Sidewaystable}
{%
  \let\oldcaption=\caption%
  \def\caption{
    \def\@captype{table}\renewcommand{\@makecaption}{\@maketabcap}
    \oldcaption}
  \begin{sidewaystable}\centering}
  {\end{sidewaystable}\let\caption=\oldcaption}
\newenvironment{Sidewaystable*}
{%
  \let\oldcaption=\caption%
  \def\caption{
    \def\@captype{table}\renewcommand{\@makecaption}{\@maketabcap}
    \oldcaption}
  \begin{sidewaystable*}\centering}
  {\end{sidewaystable*}\let\caption=\oldcaption}
%    \end{macrocode}

% Not sure why we need to stick these lines in. 
%    \begin{macrocode}
\renewenvironment{figure}
                 {\renewcommand{\@makecaption}{\@makefigcap}\@float{figure}}
                 {\end@float}
\renewenvironment{figure*}
                 {\renewcommand{\@makecaption}{\@makefigcap}\@dblfloat{figure}}
                 {\end@dblfloat}
\renewenvironment{table}
                 {\renewcommand{\@makecaption}{\@maketabcap}\@float{table}}
                 {\end@float}
\renewenvironment{table*}
                 {\renewcommand{\@makecaption}{\@maketabcap}\@dblfloat{table}}
                 {\end@dblfloat}
%    \end{macrocode}

% Redefine |\@caption| so that we don't bother with the list of
% figures. This saves us having to worry about fragile commands in the
% captions. 
%    \begin{macrocode}
\long\def\@caption#1[#2]#3{%
  \par
%%  \addcontentsline{\csname ext@#1\endcsname}{#1}%
%%    {\protect\numberline{\csname the#1\endcsname}{\ignorespaces #2}}%
  \begingroup
    \@parboxrestore
    \if@minipage
      \@setminipage
    \fi
    \normalsize
    \@makecaption{\csname fnum@#1\endcsname}{\ignorespaces #3}\par
  \endgroup}
%    \end{macrocode}

% Make tables not quite so cramped. Allow |\hline\hline| to give a
% poor man's thick rule. 
%    \begin{macrocode}
\renewcommand{\arraystretch}{1.15}
\setlength{\doublerulesep}{0.4pt}
%    \end{macrocode}

% \begin{macro}{\tablecols}
% \begin{macro}{\@tablecols}
% \begin{macro}{\tabcaptioncontinued}
% |\tablecols| is a user command for setting the number of columns. It
% is used by commands that try and span the entire table (I think
% there must be a better way of acheiving
% this). |\tabcaptioncontinued| is 
% a user command for setting the caption of a continued table on the
% next page. 
%    \begin{macrocode}
\newcommand{\tablecols}[1]{\renewcommand{\@tablecols}{#1}}
\newcommand{\@tablecols}{1}
\newcommand{\tabcaptioncontinued}{%
  \multicolumn{\@tablecols}{c}{\addtocounter{table}{-1}TABLE~\thetable\ (CONTINUED)} \\ \\
  }
%    \end{macrocode}
% \end{macro}
% \end{macro}
% \end{macro}

% \subsubsection{Footnotes in tables}
% \label{sec:table-footnotes}
% \begin{macro}{\tablenotemark}
% \begin{macro}{\tablenote}
% \begin{macro}{\tablenoteskip}
% Old user commands |\tablenotemark{MARK}|,
% |\tablenote{MARK}{TEXT}|---now deprecated. 
%    \begin{macrocode}
\newcommand\tablenotemark[1]{$^\mathrm{#1}$} 
\newcommand\tablenote[2]{\if@proceedings\else\small\fi ${}^\mathrm{#1}$ #2} 
\newcommand\tablenoteskip{\\[0.5\baselineskip]} 
%    \end{macrocode}
% \end{macro}
% \end{macro}
% \end{macro}

% \begin{macro}{\tabnotemark}
% \begin{macro}{\tabnotetext}
% \begin{macro}{\tabnotetextTWO}
% Alternative way of doing footnotes in tables. 
%    \begin{macrocode}
\newcommand{\tabnotemark}[1]{\textsuperscript{\lowercase{#1}}}
%\newcommand{\tabnotetextTWO}[4]{\tabnotetext{#1}{#2}\tabnotetext{#3}{#4}}
\newcommand{\tabnotetextTWO}[4]{%
  \multicolumn{\@tablecols}{l}{%
    \parbox{0.5\tabnotewidth}{\small\textsuperscript{#1}#2}%
    \parbox{0.5\tabnotewidth}{\small\textsuperscript{#3}#4}%
    }%
  \\}
\newcommand{\tabnotetext}[2]{%
  \multicolumn{\@tablecols}{@{}p{\tabnotewidth}@{}}{\small\textsuperscript{#1}#2}\\}
\newlength{\tabnotewidth}
%    \end{macrocode}
% \end{macro}
% \end{macro}
% \end{macro}


% \begin{macro}{changemargin}
%    \begin{macrocode}
% Command to allow margins to be changed on the fly. Useful to allow
% tables and Figures to protrude into the margins. 
%
% \changes{3.24}{2007/09/14}{Added \texttt{changemargin} environment
% (copied from TeX FAQ).}
\newenvironment{changemargin}[2]{% 
  % From TeX FAQ http://www.tex.ac.uk/cgi-bin/texfaq2html?label=chngmargonfly
  \begin{list}{}{%
    \setlength{\topsep}{0pt}%
    \setlength{\leftmargin}{#1}%
    \setlength{\rightmargin}{#2}%
    \setlength{\listparindent}{\parindent}%
    \setlength{\itemindent}{\parindent}%
    \setlength{\parsep}{\parskip}%
  }%
  \item[]}{\end{list}}
%    \end{macrocode}
% \end{macro}

% \subsection{The bibliography and references}
% \label{sec:biblio}
% \begin{macro}{\beforebiblioskip}
% \begin{macro}{\afterbiblioskip}
% I'm not entirely sure why I introduced these two commands, since
% directly putting |\vspace| before and after
% |\begin{thebibliography}| seems to work well enough. 
%    \begin{macrocode}
\newlength{\beforebiblioskip}
\setlength{\beforebiblioskip}{0pt}
\newlength{\afterbiblioskip}
\setlength{\afterbiblioskip}{0pt}
%    \end{macrocode}
% \end{macro}
% \end{macro}
% \begin{macro}{thebibliography}
% The environment itself is done as starred section with extra
% vertical space after heading. This space is currently zero. 
% In the book style we do the cunning trick with rereading the file in
% order to get the references into the |\ADS@references| macro.
% 
% \changes{3.05}{2003/01/02}{Removed the \cs{beforebiblioskip} and
% \cs{afterbiblioskip} commands, since they cause unwanted vertical
% space in some instances.}  
%    \begin{macrocode}
\renewenvironment{thebibliography}{%
  \if@book\RereadBiblio\fi
%%  \vspace*{\beforebiblioskip}
  \section*{REFERENCES} \frenchspacing
%%  \vspace*{\afterbiblioskip}
  \small
%    \end{macrocode} 
% The references themselves are set as a list in order 
% to get the hanging indentation. They are done in a smaller font and
% with french spacing so we needn't worry about the distinction
% between sentence-ending periods and abbreviations. 
%    \begin{macrocode} 
  \list{\null}{\leftmargin 1.5em\labelwidth\z@\labelsep\z@\itemindent -1.5em
    \itemsep\z@\parsep\z@\usecounter{enumi}}
  \sloppy\clubpenalty4000\widowpenalty4000
  \sfcode`\.=1000\relax\@biblisttrue}
{\endlist\@biblistfalse\normalsize\@outputaddressesifnecessary}
%    \end{macrocode}
% \end{macro}
% \subsection{Uniquification of cross-references}
% \label{sec:uniquification}
% 
% \begin{macro}{\articlename}
% In order that the |\label{}| command in individual articles be
% guaranteed to give a unique anchor, we append the article name (held
% in the |\articlename| command) to each key. This requires
% redefinition of the following commands: |\cite|, |\nocite|, |\ref|,
% |\pageref|, |\label|. 
% \begin{macro}{\oldcite}
% \begin{macro}{\oldnocite}
% \begin{macro}{\oldref}
% \begin{macro}{\oldpageref}
% \begin{macro}{\oldlabel}
% The original versions of these commands are
% saved as |\oldcite|, |\oldnocite|, etc. 
%
% The uniquifying of cross-references is now always done, so 
% we don't get multiply-defined labels in the single article case
%    \begin{macrocode}
\newcommand{\articlename}{anon}
\AtBeginDocument{
  \let\oldcite=\cite 
  \let\oldnocite=\nocite 
  \let\oldref=\ref 
  \let\oldpageref=\pageref
  \let\oldlabel=\label
}
%    \end{macrocode}
% \end{macro}
% \end{macro}
% \end{macro}
% \end{macro}
% \end{macro}
% \end{macro}
%
% Uniquify all cross-references by prepending the article name.  This
% now has to go into |\AtBeginDocument| so it is working on the new
% versions that hyperref.sty has set up.  Then we need to put the
% |\@begofarticlehook| afterwards, so that the first page label is using
% the new definitions.  If amsmath is loaded, then we also need to patch
% the special label command that they use. \changes{4.0}{2019/11/18}{Fix labels for amsmath display equations, in case that package has been loaded.}
%    \begin{macrocode}
\AtBeginDocument{
  \renewcommand{\ref}[1]{\oldref{\articlename:#1}%
    \if@debug{\ClassInfo{rmaa}{Reference to \articlename:#1}}\fi}
  \renewcommand{\pageref}[1]{\oldpageref{\articlename:#1}%
    \if@debug{\ClassInfo{rmaa}{Page reference to \articlename:#1}}\fi}
  \renewcommand{\label}[1]{\oldlabel{\articlename:#1}%
    \if@debug{\ClassInfo{rmaa}{Label for \articlename:#1}}\fi}%
  \ifdef{\label@in@display}{%
    \let\oldlabel@in@display=\label@in@display
    \renewcommand{\label@in@display}[1]{\oldlabel@in@display{\articlename:#1}%
      \if@debug{\ClassInfo{rmaa}{AMSTeX Label for \articlename:#1}}\fi}%
  }{}
  \@begofarticlehook
}%
%    \end{macrocode}
% Citations are more complicated since they can contain a list of
% keys. These are really only here for use with the rmaacite package,
% which is utterly deprecated (02 Jan 2003) now that I have the natbib
% code working. The command \cs{cite} should \emph{not} be used with
% natbib: use \cs{citep}, \cs{citet}, etc.\@ instead. 
%    \begin{macrocode}
\renewcommand{\nocite}[1]{\@bsphack
  \@for\@citeb:=#1\do{%
    \edef\@citeb{\expandafter\@firstofone\@citeb:\articlename}%
    \if@filesw\immediate\write\@auxout{\string\citation{\@citeb}}\fi
    \@ifundefined{b@\@citeb:\articlename}{\G@refundefinedtrue
      \@latex@warning{Citation `\@citeb' undefined}}{}}%
  \@esphack}
\renewcommand{\cite}[1]{\oldcite{#1:\articlename}}
%   \end{macrocode}


% Horrible fiddling with the bibliography citation mechanism. 
% This has to be globally available in case rmaacite puts it in the
% \texttt{.aux} file. 
%    \begin{macrocode}
\def\lbibcite#1#2#3{\global\@namedef{b@#1}{#2}
                    \global\@namedef{by@#1}{#3}}
%    \end{macrocode}

% \begin{macro}{\oldbibitem}
% \begin{macro}{\@rmaaBibArg}
% \begin{macro}{\@rmaaSpace}
% \begin{macro}{rmaabib}
% Of course, |\bibitem| needs to get the uniquification treatment as
% well. While we are at it, we redefine it to use consecutive integers
% for the key if the user leaves it blank (using the counter
% |rmaabib|). This allows one to put 
% |\bibitem{}| for all your references without getting annoying
% messages about multiply-defined labels.
% 
% We also take the opportunity to insert the optional argument to
% \cs{bibitem} if it is missing in the form Author(DATE) so that natbib
% doesn't try to do numerical citations. Natbib already provides a
% hook \cs{@extra@b@citeb} to help uniquify the citation key, so we
% just plug into that. However, it puts it \emph{after} the original
% key, not before as I had had it. Easiest solution is to change the
% way that I do it. 02 Jan 2003. Note that the \cs{@extra@b@citeb}
% does not get written to the .aux  file in the \cs{citation} command
% but it does work. 
% \changes{3.14}{06 Feb 2003}{Reinstalled old definition of \cs{bibitem}
%  for the case where the nonatbib option is requested, so that the
%  rmaacite package will work. This means that you can't use rmaacite
%  and natbib in the same volume.}
%    \begin{macrocode}
\if@natbib\def\@extra@b@citeb{:\articlename}\fi
%%\else
\let\oldbibitem\bibitem
\newcommand{\@rmaaBibArg}{}
\newcommand{\@rmaaSpace}{ }
\newcounter{rmaabib}
\if@natbib
\renewcommand{\bibitem}[2][ABC(2003)]{%
  \renewcommand{\@rmaaBibArg}{#2}
  \ifx\@rmaaBibArg\@rmaaEmpty
     \addtocounter{rmaabib}{1}
     \oldbibitem[#1]{\thermaabib:\articlename}
  \else\ifx\@rmaaBibArg\@rmaaSpace
     \addtocounter{rmaabib}{1}
     \oldbibitem[#1]{\thermaabib:\articlename}
  \else 
     \oldbibitem[#1]{#2:\articlename}
  \fi\fi
  \renewcommand{\@rmaaBibArg}{}
  }
\else
\renewcommand{\bibitem}[1]{%
  \renewcommand{\@rmaaBibArg}{#1}
  \ifx\@rmaaBibArg\@rmaaEmpty
     \addtocounter{rmaabib}{1}
     \oldbibitem{\articlename:\thermaabib}
  \else\ifx\@rmaaBibArg\@rmaaSpace
     \addtocounter{rmaabib}{1}
     \oldbibitem{\articlename:\thermaabib}
  \else
     \oldbibitem{\articlename:#1}
  \fi\fi
  \renewcommand{\@rmaaBibArg}{}
  }
\fi
%%\fi
%    \end{macrocode}
% \end{macro}
% \end{macro}
% \end{macro}
% \end{macro}

% \begin{macro}{oldpcite}
% \begin{macro}{oldscite}
% Then, we need to repeat the above scheme, but for the |rmaacite|
% definition. 
%    \begin{macrocode}
\newcommand{\@redefinebibitem}{%
  \let\oldbibitem\bibitem
  \let\oldcite\cite
  \let\oldpcite\pcite
  \let\oldscite\scite
  \def\bibitem[##1<##2>]##3{%
    \renewcommand{\@rmaaBibArg}{##3}
    \ifx\@rmaaBibArg\@rmaaEmpty
       \addtocounter{rmaabib}{1}
       \oldbibitem[##1<##2>]{\thermaabib:\articlename}
    \else\ifx\@rmaaBibArg\@rmaaSpace
       \addtocounter{rmaabib}{1}
       \oldbibitem[##1<##2>]{\thermaabib:\articlename}
    \else 
       \oldbibitem[##1<##2>]{##3:\articlename}
    % trust the author here
%%       \oldbibitem[##1<##2>]{##3}
    \fi\fi
    }
  }
%    \end{macrocode}
% \end{macro}
% \end{macro}


% \subsection{HTML Table of Contents}
% \label{sec:hyper-TOC}
% \begin{macro}{\HyperTOCEntry}
% \begin{macro}{\HyperTOCpart}
% \begin{macro}{\HyperTOCmeeting}
% \begin{macro}{\HyperTOCeditors}
% \begin{macro}{\HyperTOCpreamble}
% These commands are used to write information to the \texttt{.htoc}
% file, which is used in generating a Web Page of the Table of
% Contents. Since they only ever appear preceded by |\protect| in the
% argument to an |\addtocontents| command,  they should never need to
% be expanded by the macros. However, we define them here just in case
% \dots
%    \begin{macrocode}
\newcommand{\HyperTOCEntry}[3]{\relax}
\newcommand{\HyperTOCpart}[1]{\relax}
\newcommand{\HyperTOCmeeting}[3]{\relax}
\newcommand{\HyperTOCeditors}[1]{\relax}
\newcommand{\HyperTOCpreamble}{\relax}
%    \end{macrocode}
% \end{macro}
% \end{macro}
% \end{macro}
% \end{macro}
% \end{macro}
%

% \subsection{Formatting the poster abstracts} 
% These can now be used without using the book option.
%    \begin{macrocode}

\newcommand{\@makeminititle}{%
  \newpage
  \null 
  \vskip 3em%
  \begin{center}%
    {\@title \par}%
  \end{center}
  \vskip 3em%
  }

\newif\if@abstractsads
\@abstractsadsfalse
\newcommand\abstractsADS{\@abstractsadstrue}

\if@book\RequirePackage{multicol}\fi
\newenvironment{abstracts}[1][Abstracts of Contributed Papers]{%
  %% First, make sure we start on an odd page (06 Jun 2000)
  \if@book\clearpage\thispagestyle{empty}\cleardoublepage\fi
  %%
  \renewcommand{\articlename}{ABSTRACTS}
  \title{#1}\shortauthor{Abstracts}\shorttitle{Abstracts}
  \listofauthors{}
  \if@book
  \addcontentsline{toc}{chapter}{\protect\@fiddleauthors{\@TOCtitle}{\rmaaTOCauthors}}%
  \addtocontents{htc}{\protect\HyperTOCEntry{\@rawtitle}{\@rawauthors}{\articlename}}
  \fi
  \pagestyle{rmaa} % this now seems necessary (06 Jun 2000)
  \thispagestyle{rmaafirstpage}
%%  \twocolumn[\@makeminititle]
  \renewcommand{\abstract}[1]{%
    \renewcommand\ADS@abstract{##1}%
    \renewcommand{\@abstract}{##1}\@writeabstract%
  }
  \renewcommand{\title}[2][]{%
    \gdef\ADS@authors{}%
    \gdef\ADS@affil{}%
    \gdef\ADS@references{}%
    \gdef\ADS@keywords{}%
    \renewcommand{\@rmaaArg}{##1}
    \ifx\@rmaaArg\@rmaaEmpty
       \renewcommand{\@TOCaftertitle}{}
    \else
       \renewcommand{\@TOCaftertitle}{ (\@rmaaArg)}
    \fi
    \renewcommand{\@title}{\rmaa@TOCTitleStyle{##2}}%
    \renewcommand{\@TOCtitle}{\rmaa@TOCTitleStyle{##2}\@TOCaftertitle}%
    \renewcommand{\rmaaIndexTitle}{\rmaa@IndexTitleStyle{##2\@TOCaftertitle.}}%
    \renewcommand\ADS@title{##2}%
    \xdef\ADS@firstpage{\@arabic{\c@page}}%
    }
  \renewcommand{\author}[1]{\let\@thanks\@empty\renewcommand\@author{##1} }
  \renewcommand{\footnotetext}[2][XX]{\par\noindent\textsuperscript{##1}~##2}
  \renewcommand{\footnoterule}{\kern-3\p@ \hrule \@width 3em \kern 2.6\p@}
  \if@book\@begofarticlehook\fi
  \if@book\begin{multicols}{2}[\@makeminititle]\else\@makeminititle\fi
  \begingroup
  \setcounter{footnote}{0}
  }% and at the end of the environment....
  {\endgroup\if@book\end{multicols}\listofabstracts\label{RMAALastPage}\clearpage\fi}
%% auxiliary commands used:
%% \@writeabstract does all the actual formatting 
\newcommand{\@writeabstract}{%
  \if@book%
  \addcontentsline{loa}{abstract}{\protect\@fiddleauthors{\@TOCtitle}{\rmaaTOCauthors}}%
  \fi%
  \begin{center}\@title\par\@author\end{center}
  \@abstract
  \par\noindent\footnoterule\nobreak\footnotesize\@thanks
  \par\noindent\hrulefill\goodbreak
  \vskip 1cm \@plus0.5cm
  %% write ADS entry
  \if@book\if@abstractsads\ADS@dumpdata\fi\fi%
  \endgroup % this pair is to make each abstract a self-contained group
  \begingroup
  \setcounter{footnote}{0}
  }

\if@book
%% List of abstracts (stuff written to .loa file)
\newcommand{\listofabstracts}{%
  \begin{multicols}{2}[%
    \vskip 1.5cm \vfill % leave at least 1.5 cm and go to bottom of page if possible
    \centerline{\MakeUpperCase{List of Abstracts}}
    \vskip 1.5cm \@plus0.5cm
    ][6cm] % this makes it start a new page if there is not at least 6cm of space
    \renewcommand{\@fiddleauthors}[2]{##1\hskip1em plus2em minus0.4em\penalty-100  ##2}
    \parindent\z@
%%    \parskip.5\p@ \@plus .7\p@\relax 
    \parskip3\p@ \@plus 1\p@ \@minus 1\p@\relax % allow quite a lot of
                                % leeway here to avoid breaks inside
                                % entries 
    \renewcommand{\\}{\relax} % temporarily turn off explicit line breaks just in case
    \renewcommand{\altaffilmark}[1]{\relax} % ditto with this
    \@starttoc{loa}
  \end{multicols}
  }
\newcommand*{\l@abstract}[2]{%
  \par\penalty-1000
  #1\hskip 1em\penalty0\null\nobreak\hfill#2}
\fi

%    \end{macrocode}
 
% \subsection{Code for typesetting an entire volume of conference
% proceedings} \label{sec:whole-volume}
%    \begin{macrocode}

\if@book
%% environment for including the individual articles
\newcommand{\@fiddleauthors}[2]{\relax}
\renewcommand{\theequation}{\@arabic\c@equation}
\newenvironment{articles}%
{ \begingroup % make it a group so that author definitions do not
              % spill over
  \gdef\ADS@authors{}
  \pagestyle{rmaa}
  \renewenvironment{document}{%
    \@begofarticlehook
    \addtocounter{chapter}{1}%
    \addcontentsline{toc}{chapter}{%
      \protect\@fiddleauthors{\@TOCtitle}{\rmaaTOCauthors}}%
    \addtocontents{htc}{\protect\HyperTOCEntry{\@rawtitle}{\@rawauthors}{\articlename}}
    }{\relax}
  \renewcommand{\documentclass}[2][foo]{\relax} 
  \renewcommand{\usepackage}[2][foo]{%
    \def\@packagename{##2}\def\@rmaacitename{rmaacite}%
    \ifx\@packagename\@rmaacitename
    \@rmcitetrue\ClassInfo{rmaa}{This paper uses rmaacite}
    \makeatletter\input{rmaacite.sty}\makeatother % read it in here so that all changes
    \else                                         % are local to paper
    \if@debug\ClassWarning{rmaa}{Package \@packagename\space was requested by \articlename\space}\fi
    \fi} 
  %% Don't allow individual articles to mess with the pages, volume or year
  \renewcommand{\SetFirstPage}[1]{\relax}
  \renewcommand{\SetYear}[1]{\relax}
  \renewcommand{\SetVolume}[1]{\relax}
  \renewcommand{\SetConfTitle}[1]{\relax}
  %% Don't allow individual articles to change these
  \renewcommand{\nonstopmode}{\relax}
  \renewcommand{\scrollmode}{\relax}
  \renewcommand{\outputplates}{\relax}
  %% XXXX need to check if using rmaacite and act accordingly
  %% avoid multiple declarations of the same length
  \let\oldnewlength=\newlength
  \renewcommand{\newlength}[1]{%
    \edef\reserved@a{\expandafter\@gobble\string ##1}%
    \@ifundefined{\reserved@a}{%
      \oldnewlength{##1}}{%
      \if@debug{\ClassInfo{rmaa}{\reserved@a\space already defined: doing nothing}}\fi%
      \relax}}%
  \renewcommand\newcommand{\providecommand}
  }%
{\endgroup\onecolumn}
     

%    \end{macrocode}
% Command to be used in `articles' environment for including an
% article
%    \begin{macrocode}
\newcommand{\article}[2][]{%
  % optional first argument is extra text to put in brackets after title in TOC
  % e.g. ``Invited Talk'' or ``Poster''
  \renewcommand{\articlename}{#2}
  % reset loads of counters
  \setcounter{footnote}{0}
  \setcounter{figure}{0}
  \setcounter{table}{0}
  \setcounter{section}{0}
  \setcounter{equation}{0}
  \setcounter{rmaabib}{0}
  \setcounter{author}{0}
  \setcounter{AUTHOR}{1}
  \gdef\ADS@authors{}%
  \gdef\ADS@affil{}%
  \gdef\ADS@affiliations{}%
  \gdef\ADS@references{}%
  \@firstauthortrue
  \@firstaffiltrue
  \@firstaltaffiltrue
  \@firstkwdtrue
  \@rmcitefalse
  \if@suppressfulladdresses\else\@needaddressestrue\fi
  \renewcommand{\@rmaaArg}{#1}
  \ifx\@rmaaArg\@rmaaEmpty
     \renewcommand{\@TOCaftertitle}{}
  \else
     \renewcommand{\@TOCaftertitle}{ (#1)}
  \fi
%    \end{macrocode}
% In the main journal articles must start on an odd page. 
%    \begin{macrocode}
  \if@journal\clearpage\thispagestyle{empty}\cleardoublepage\fi
%% Now input the file, then run the end-of-article hook (reinstated 28
  %% Sep 2001)
  {\input{\@articledir/\articlename.tex}\@endofarticlehook}}
%    \end{macrocode}
% Command to separate articles into sessions. Write an entry in the
% TOC, makes sure we start on an odd page (added 06 Jun 2000), then do
% a title page if \cs{SessionTitlePages} was called. 
% \changes{3.12}{04 Feb 2003}{Added option to have title page for each
% session. Turn it on with \cs{SessionTitlePages}.}
%    \begin{macrocode}
\newcommand{\session}[1]{
  %% this does nothing except write an entry in the TOC
  \addcontentsline{toc}{part}{\rmaa@TOCPartStyle{#1}}
  \addtocontents{htc}{\protect\HyperTOCpart{#1}}
  %% and also makes sure we start on an odd page (added 06 Jun 2000)
  \clearpage\thispagestyle{empty}\cleardoublepage
  \stepcounter{part}
  \if@sessiontitlepages{%
    \onecolumn
    \thispagestyle{empty}
    \renewcommand{\articlename}{Section\thepart TitlePage}
    \label{RMAAFirstPage}
    \make@session@title{#1}
    \label{RMAALastPage}
    \clearpage\thispagestyle{empty}\cleardoublepage
  }\fi
  }
\newcommand{\make@session@title}[1]{%
  {%
    \setlength\parindent{0pt}
    \vspace*{\stretch{1}}
    \rule{\linewidth}{1mm}
    \begin{flushright}\Huge\bfseries\boldmath
      #1
    \end{flushright}
    \rule{\linewidth}{1mm}
    \vspace*{\stretch{2}}
  }
}
%    \end{macrocode}
% Command to turn on separate title pages for each session. 
%    \begin{macrocode}
\newcommand{\SessionTitlePages}{\@sessiontitlepagestrue}
\newcommand{\BoldSessionTitles}{%
  \renewcommand{\rmaa@TOCPartStyle}[1]{\textbf{\rmaa@TOCTitleStyle{##1}}}
  }
\newcommand{\ArticleDir}[1]{\renewcommand{\@articledir}{#1/}\graphicspath{{1/}}}
\newcommand{\@articledir}{.}
%\newcommand{\ps@rmaafront}{\ps@rmaa}% page style for front matter (same as ever, actually)
\newcommand{\ps@rmaafront}{% page style for front matter:  page # in each foot
  \renewcommand\@oddfoot{\hfil{\thepage}}%
  \renewcommand\@evenfoot{{\thepage}\hfil}%
  % shorttitle in both heads
  \renewcommand\@evenhead{\hfil\@shorttitle\hfil}%
  \renewcommand\@oddhead{\@evenhead}
  }
\newcommand{\ps@rmaafrontfirstpage}{% page style for first page of front matter
  % page # in each foot, on outside edge
  \renewcommand\@oddfoot{\hfil\thepage}%
  \renewcommand\@evenfoot{\thepage\hfil}%
  % nothing in the even heads
  \renewcommand\@evenhead{\relax}%
  \renewcommand\@oddhead{\relax}%
  }
%% slight modification to the standard \frontmatter and \mainmatter
%% commands, just to set the \pagestyle
\renewcommand\frontmatter{%
  \latexonecolumn
  \cleardoublepage
  \@frontmattertrue
  \@mainmatterfalse
  \pagenumbering{roman}
  \pagestyle{rmaafront}
  \addtocounter{page}{2} % leave space for the title page
  }
\renewcommand\mainmatter{%
  \pagestyle{empty}
  \if@rmtwocol\twocolumn\fi
  \cleardoublepage
  \@frontmatterfalse
  \@mainmattertrue
  \pagenumbering{arabic}
  \pagestyle{rmaa}
  }
\renewcommand\backmatter{%
  \clearpage\pagestyle{empty}\cleardoublepage
  \@mainmatterfalse
  %% put extra space in the TOC
  \addtocontents{toc}{\bigskip}
  }
%    \end{macrocode}
%
% \subsection{Front sections}\label{sec:front-sections}
% \begin{macro}{\notitle}
% \begin{macro}{\shortname}
% \begin{macro}{frontsection}
% The environment \texttt{frontsection} is for doing a section in the
% front matter (preface, list of participants, or whatever). It has 
% one optional argument:
% \begin{enumerate}
% \item[\#1] Short title (no embedded spaces) for
% use in making PS files. E.g., ``list\_of\_participants''.   
% \end{enumerate}
% and two mandatory arguments: 
% \begin{enumerate}
% \item[\#2] Title of section (e.g., `Preface')
% \item[\#3] Commands to be evaluated before processing the contents
% of the environment. This is mainly useful for putting in the things
% that would go in the preamble of a normal article  (e.g.,
% |\listofauthors|, |\indexauthor|, |\notitle|) 
% \end{enumerate}
%    \begin{macrocode}
\newcommand{\notitle}{\@titlevisiblefalse} \newcommand{\shortname}{}
\newenvironment{frontsection}[3][]{%
  \begingroup
  \renewcommand{\shortname}{#1}
  \ifx\shortname\@rmaaEmpty
     \renewcommand{\articlename}{#2}
  \else
     \renewcommand{\articlename}{\shortname}
  \fi
  \@begofarticlehook
  \@titlevisibletrue % default
  \thispagestyle{rmaafrontfirstpage}
  \listofauthors{} % blank out this in case it is not specified
  \title{{#2}}\shorttitle{#2}\shortauthor{#2}
  #3 % evaluate the command given as second mandatory arg
  \addcontentsline{toc}{chapter}{\protect\@fiddleauthors{\@TOCtitle}{\rmaaTOCauthors}}%
  \addtocontents{htc}{\protect\HyperTOCEntry{\@rawtitle}{\@rawauthors}{\articlename}}
  \if@titlevisible
     \vspace*{1cm}\centerline{\rmaa@FrontSecStyle{\@title}}\vspace*{1cm}
  \fi
  \setlength{\parskip}{3ex}
  }{\frontlabel{RMAALastPage}\clearpage\endgroup}
%    \end{macrocode}
% \end{macro}
% \end{macro}
% \end{macro}
%    \begin{macrocode}
%% Definition of the command \frontlabel, which rights out something like `F4' to
%% the .aux file, instead of `\textlatin  {i}'. This is so it is easy to post-process
%% the file for making PS files of individual sections
\newcommand{\frontlabel}[1]{%
  \@bsphack
  \protected@write\@auxout{}%
     {\string\newlabel{\articlename:#1}{{}{F\@arabic{\c@page}}}}
  \@esphack
  }
%% User command to say X, Y, Z,\\ Editors of the Proceedings
\newcommand{\EditorSignature}{\vspace*{\parskip}\rightline{\@editors}%
  \rightline{Editors of the proceedings}}
%% Fiddle with \thesection, etc. so that they don't include the chapter number
%% Basically go back to the `article' definitions of these
\renewcommand{\thesection}{\@arabic\c@section}
\renewcommand{\thefigure}{\@arabic\c@figure}
\renewcommand{\thetable}{\@arabic\c@table}

\newcommand{\@oldargs}{}
\newcommand{\@newargs}{}
%% add TOC entry for each paper
%\renewcommand{\title}[1]%
%{\renewcommand{\@title}{\rmaa@TitleStyle{#1}}%
%\addtocounter{chapter}{1}\addcontentsline{toc}{chapter}{#1}%
%}

%% Formatting of the table of contents
%% dimensions of various parts
\newlength{\TOCtextwidth}\setlength{\TOCtextwidth}{12.65cm}
\newlength{\TOCpagenumwidth}\setlength{\TOCpagenumwidth}{1.3cm}
\newlength{\TOCentrywidth}\setlength{\TOCentrywidth}{\TOCtextwidth}
\addtolength{\TOCentrywidth}{-\TOCpagenumwidth} 
\newlength{\TOCextramargin}\setlength{\TOCextramargin}{1.25cm}
\newcommand{\TOCbox}[1]{\hspace{\TOCextramargin}\parbox{\TOCtextwidth}{#1}\par}
\newlength{\TOCextramarginx}\setlength{\TOCextramarginx}{\TOCextramargin}\addtolength{\TOCextramarginx}{4.3cm}
\newlength{\TOCtextwidthx}\setlength{\TOCtextwidthx}{\TOCtextwidth}\addtolength{\TOCtextwidthx}{-4.3cm}
\newcommand{\TOCboxx}[1]{\hspace{\TOCextramarginx}\parbox{\TOCtextwidthx}{#1}\par}

\newcommand{\RMAAbannerimage}{rmsc}
\renewcommand\tableofcontents{%
  \renewcommand{\articlename}{CONTENTS}
  \begingroup
  \@begofarticlehook
  \if@twocolumn
     \@restonecoltrue\onecolumn
  \else
     \@restonecolfalse
  \fi
  %% 9pt for journal style
  \if@proceedings\else\small\fi
  \if@newstyle\small\fi %% same for new style
  \if@proceedings\thispagestyle{rmaafrontfirstpage}\else\thispagestyle{empty}\fi
  \shortauthor{Contents}\shorttitle{Contents}
  \vspace*{18mm}
  \if@printbanner
  %% put in the banner header (need to get better image of this)
     \TOCbox{\hspace*{-0.15cm}\includegraphics[width=0.9\TOCtextwidth]{\RMAAbannerimage}}
     \vspace*{8mm}
  %% volume and publication date (should really be in funny font)
     \TOCboxx{\@volumen\\[2mm] \@mes\ \@anyo}
  \else
  %% put in the banner header (need to get better image of this)
     \TOCbox{\hspace*{-0.15cm}\includegraphics[width=0.9\TOCtextwidth]{\RMAAbannerimage}}
     \vspace*{12.5mm}
  %% volume and publication date (should really be in funny font)
     \TOCboxx{~\\[2mm] ~}
  \fi
  \vspace*{16mm}
  %% This is the command that fiddles with the TOC entry so that the
  %% title is flushleft but the authors are flushright
  \renewcommand{\@fiddleauthors}[2]{%
    % first put the paper title
    \renewcommand{\\}{\relax} % temporarily turn off explicit line breaks just in case
    \renewcommand{\altaffilmark}[1]{\relax} % ditto with this
    \renewcommand{\textsuperscript}[1]{\relax} % ditto with this 
    ##1
    % now save the author list in a box
    \sbox\@tempboxa{##2}%
    % check that box is shorter than the total width available for the entry
    \ifdim \wd\@tempboxa <\TOCentrywidth
       % if so, then put it right here, flush against right margin
       \penalty10000\hskip1em plus10em\penalty0\hbox{}\nobreak\hfill\box\@tempboxa
       % Note that we force a space of at least 1em between title and authorlist, 
       % stretchable to 11em since we may have very short titles but still with an 
       % authorlist too long to go on the same line.  The only linebreak allowed is 
       % directly before the authorlist.
    \else
       % if it is too long, we set the original authorlist text in its own parbox
       \par\vskip1ex\parbox[b]{\TOCentrywidth}{\raggedleft ##2}
    \fi
    }
  \@starttoc{toc}%
  \if@restonecol\twocolumn\fi
  \frontlabel{RMAALastPage}
  \clearpage
  \thispagestyle{empty}
  \cleardoublepage % make sure we start on an odd page
  \endgroup

  %%% ALSO, do the HTML TOC
  %% The file that will be used to make the HTML TOC
  \@starttoc{htc}
  %% All the stuff in this file does nothing in the book style
  \addtocontents{htc}{\protect\HyperTOCmeeting{\@conftitle}{\@reunionplace}{\@reuniondate}}
  \addtocontents{htc}{\protect\HyperTOCeditors{\@editors}}
  \addtocontents{htc}{\protect\HyperTOCpreamble}
  \addtocontents{htc}{\protect\HyperTOCEntry{Table of Contents}{}{\articlename}}

  %% and the ADS file
  \ADS@openfile
  }
%% This is now much simpler
\renewcommand\l@chapter[2]{%TOC entry for articles
  \ifnum \c@tocdepth >\m@ne
    \vskip 1.0em \@plus\p@
    \TOCbox{%
      \parbox[b]{\TOCentrywidth}{#1}\parbox[b]{\TOCpagenumwidth}{\hfill#2}
    }
  \fi} 
\renewcommand\l@part[2]{%TOC entry for sessions: centered, no page number
  \ifnum \c@tocdepth >\m@ne
    \vskip 1.5em \@plus1.5em \goodbreak
    \TOCbox{%
      \parbox[b]{\TOCentrywidth}{\hfil#1\hfil}\hfill
    }
  \fi\nobreak} 

%% Formatting theindex
%% strongly encourage page breaks between index entries
\renewcommand{\@idxitem}{\par\penalty-1000\hangindent.45cm} 
\renewcommand{\indexspace}{\par\penalty-1000\vskip 10\p@ \@plus15\p@ \@minus3\p@\relax}
\renewenvironment{theindex}{%
  %% first, make sure we start on an odd page (06 Jun 2000)
  \clearpage\thispagestyle{empty}\cleardoublepage
  %% this looks a bit like an article but is two columns and without abstract, etc
  \renewcommand{\articlename}{INDEX}
  \title{Author Index}\shortauthor{Author Index}\shorttitle{Author Index}
  \listofauthors{}
  \addcontentsline{toc}{chapter}{\protect\@fiddleauthors{\@TOCtitle}{\rmaaTOCauthors}}%
  \addtocontents{htc}{\protect\HyperTOCEntry{\@rawtitle}{\@rawauthors}{\articlename}}
  \pagestyle{rmaa}
  \thispagestyle{rmaafirstpage}
  \@begofarticlehook
  %% 9pt for journal style
  \if@proceedings\else\small\fi
  \if@newstyle\small\fi
  \twocolumn[\@makeminititle]
  %% vertical spacing (see also definition of \@idxitem and \indexspace
  \parskip\z@ \@plus .6\p@\relax % allow quite a lot of leaway here to avoid breaks inside entries
  %% horizontal spacing, breaking and glue parameters (see also the stuff in the definition 
  %% of \indexauthor
  \parindent\z@ 
  \exhyphenpenalty10000 % don't break on explicit hyphens, it always looks bad
  \rightskip=0pt plus.2em % allow a tiny bit of leeway at the ends of lines
  \spaceskip.3333em plus.5em minus.15em % a lot of leeway between words
  \xspaceskip.5em plus.5em minus.4em % this seems to have no effect
  \let\item\@idxitem
  \renewcommand{\\}{\relax} % temporarily turn off explicit line breaks just in case
  \renewcommand{\altaffilmark}[1]{\relax} % ditto with this
  \renewcommand{\textsuperscript}[1]{\relax} % ditto with this
  }{\label{RMAALastPage}\clearpage} % clearpage necessary so that we know about \shortauthor



%% end definitions for book styles
\fi

%    \end{macrocode}
% In order to do away with plates, we sometimes want a figure to
% appear in color only in the online version. The following command
% can bve used in a figure caption to typeset a note to this effect. 
%    \begin{macrocode}
\newcommand\ColorOnlineOnly{%
  \textsc{Note: this figure is available in color in the electronic
    version of this article, obtainable from}
  \texttt{http://\discretionary{}{}{}www.\discretionary{}{}{}astroscu.\discretionary{}{}{}unam.\discretionary{}{}{}mx/\discretionary{}{}{}\~{}rmaa/}.}
%% Plates
%% Define 2 counters: plate is the number in the whole document, plateloop is used 
%% internally by \outputplates
\newcounter{plate}
\newcounter{plateloop}
%% Currently limited to 15 plates in document
%% when we do a label, we need to write to \newlabel entries to the .aux file, one 
%% that says RMAAPlateX so that it can be referred to in \ouputplates, and one that
%% says articlename:RMAAPlateX so that the dvips program can find it and asociate it
%% with the right article
\newcommand{\@platelabeli}{\oldlabel{RMAAPlate1}\label{RMAAPlate1}}
\newcommand{\@platelabelii}{\oldlabel{RMAAPlate2}\label{RMAAPlate2}}
\newcommand{\@platelabeliii}{\oldlabel{RMAAPlate3}\label{RMAAPlate3}}
\newcommand{\@platelabeliv}{\oldlabel{RMAAPlate4}\label{RMAAPlate4}}
\newcommand{\@platelabelv}{\oldlabel{RMAAPlate5}\label{RMAAPlate5}}
\newcommand{\@platelabelvi}{\oldlabel{RMAAPlate6}\label{RMAAPlate6}}
\newcommand{\@platelabelvii}{\oldlabel{RMAAPlate7}\label{RMAAPlate7}}
\newcommand{\@platelabelviii}{\oldlabel{RMAAPlate8}\label{RMAAPlate8}}
\newcommand{\@platelabelix}{\oldlabel{RMAAPlate9}\label{RMAAPlate9}}
\newcommand{\@platelabelx}{\oldlabel{RMAAPlate10}\label{RMAAPlate10}}
\newcommand{\@platelabelxi}{\oldlabel{RMAAPlate11}\label{RMAAPlate11}}
\newcommand{\@platelabelxii}{\oldlabel{RMAAPlate12}\label{RMAAPlate12}}
\newcommand{\@platelabelxiii}{\oldlabel{RMAAPlate13}\label{RMAAPlate13}}
\newcommand{\@platelabelxiv}{\oldlabel{RMAAPlate14}\label{RMAAPlate14}}
\newcommand{\@platelabelxv}{\oldlabel{RMAAPlate15}\label{RMAAPlate15}}
\newcommand{\@outputplatelabeli}{\oldlabel{RMAAOutputPlate1}}
\newcommand{\@outputplatelabelii}{\oldlabel{RMAAOutputPlate2}}
\newcommand{\@outputplatelabeliii}{\oldlabel{RMAAOutputPlate3}}
\newcommand{\@outputplatelabeliv}{\oldlabel{RMAAOutputPlate4}}
\newcommand{\@outputplatelabelv}{\oldlabel{RMAAOutputPlate5}}
\newcommand{\@outputplatelabelvi}{\oldlabel{RMAAOutputPlate6}}
\newcommand{\@outputplatelabelvii}{\oldlabel{RMAAOutputPlate7}}
\newcommand{\@outputplatelabelviii}{\oldlabel{RMAAOutputPlate8}}
\newcommand{\@outputplatelabelix}{\oldlabel{RMAAOutputPlate9}}
\newcommand{\@outputplatelabelx}{\oldlabel{RMAAOutputPlate10}}
\newcommand{\@outputplatelabelxi}{\oldlabel{RMAAOutputPlate11}}
\newcommand{\@outputplatelabelxii}{\oldlabel{RMAAOutputPlate12}}
\newcommand{\@outputplatelabelxiii}{\oldlabel{RMAAOutputPlate13}}
\newcommand{\@outputplatelabelxiv}{\oldlabel{RMAAOutputPlate14}}
\newcommand{\@outputplatelabelxv}{\oldlabel{RMAAOutputPlate15}}
\newcommand{\@platerefi}{\oldpageref{RMAAPlate1}}
\newcommand{\@platerefii}{\oldpageref{RMAAPlate2}}
\newcommand{\@platerefiii}{\oldpageref{RMAAPlate3}}
\newcommand{\@platerefiv}{\oldpageref{RMAAPlate4}}
\newcommand{\@platerefv}{\oldpageref{RMAAPlate5}}
\newcommand{\@platerefvi}{\oldpageref{RMAAPlate6}}
\newcommand{\@platerefvii}{\oldpageref{RMAAPlate7}}
\newcommand{\@platerefviii}{\oldpageref{RMAAPlate8}}
\newcommand{\@platerefix}{\oldpageref{RMAAPlate9}}
\newcommand{\@platerefx}{\oldpageref{RMAAPlate10}}
\newcommand{\@platerefxi}{\oldpageref{RMAAPlate11}}
\newcommand{\@platerefxii}{\oldpageref{RMAAPlate12}}
\newcommand{\@platerefxiii}{\oldpageref{RMAAPlate13}}
\newcommand{\@platerefxiv}{\oldpageref{RMAAPlate14}}
\newcommand{\@platerefxv}{\oldpageref{RMAAPlate15}}
\newcounter{figplatei}
\newcounter{figplateii}
\newcounter{figplateiii}
\newcounter{figplateiv}
\newcounter{figplatev}
\newcounter{figplatevi}
\newcounter{figplatevii}
\newcounter{figplateviii}
\newcounter{figplateix}
\newcounter{figplatex}
\newcounter{figplatexi}
\newcounter{figplatexii}
\newcounter{figplatexiii}
\newcounter{figplatexiv}
\newcounter{figplatexv}
%% Printed output should use the counter `plate'
\renewcommand{\theplate}{\@arabic\c@plate}
\renewcommand{\theplateloop}{\@arabic\c@plateloop}
%% Command to output all the saved plates in this paper
\newcommand{\outputplates}{%
  \setcounter{plateloop}{0}
  \if@book\else\@endofarticlehook\fi % not entirely sure why, but this
                                     % seems to be needed
  \clearpage % first output the last page of text
  % now set pagestyle for blank even pages and make sure we start on an odd page
  \pagestyle{empty}
  \cleardoublepage
  \onecolumn
  \parindent\z@
  % first put in an extra sheet
  % the main purpose of this is so there is always a blank even page before 
  % each plate for filling in the PS file if the article has an odd number
  % of pages (see splitbook.sh)
  \null\vfil\centerline{PLATES}\vfil\clearpage\pagestyle{rmaaplate}\cleardoublepage
  % output each plate in turn
  \@whilenum\value{plateloop}<\value{plate}\do{%
    \addtocounter{plateloop}{1}
    % restore the figure counter from \@figcounti, \@figcountii, etc
    \setcounter{figure}{\value{figplate\roman{plateloop}}}
    % spew out the plate (in its own group for safety)
    {\csname @plate\roman{plateloop}\endcsname\par}
    }
  \pagestyle{rmaa}
  % reset the page number as if nothing had happened
  \addtocounter{page}{-\value{plate}}
  \addtocounter{page}{-\value{plate}}
  }
\newcommand\figcaption{%
  \def\@captype{figure}
  \renewcommand{\@makecaption}{\@makefigcap}
  \caption}
\newcommand\tabcaption{%
  \def\@captype{table}
  \renewcommand{\@makecaption}{\@maketabcap}
  \caption}

\newcommand{\plate}[2][1]{%
  %% User command to make a plate page: 
  %% Mandatory argument is what goes on the plate - this should not 
  %% include any floating environments (e.g. figure). Instead, \figcaption 
  %% should be used to do the captions. 
  %% Optional argument is the number of figures (i.e. captions) on the plate
  %% page (default is 1). This is necessary for incrementing the figure counter. 
  %%
  %% Note that we have to use a \gdef rather than an \xdef when we 
  %% save the plate output in a command. This means that in the book styles when it 
  %% comes to spewing out the plate, we will have forgotten about any user macros 
  %% that may have been defined. This can be got round by repeating the macros inside the plate.
  %% (Presumably, one would use a \newcommand rather than \renewcommand if they were new.)
  % increment global and local counters
  \addtocounter{plate}{1}
  % write appropriate label to .aux file
  \csname @platelabel\roman{plate}\endcsname
  % define a global macro called \@platei, \@plateii, \@plateiii, \@plateiv, ...
  \expandafter\gdef\csname @plate\roman{plate}\endcsname{%
    \csname @outputplatelabel\roman{plateloop}\endcsname
    \null\vfil #2 \vfil
    \cleardoublepage
    }
  % save the current shortauthor and shorttitle
  % need to be \xdef to take current values of \@shortauthor, etc 
  \expandafter\xdef\csname @shorta\roman{plate}\endcsname{\@shortauthor}
  \expandafter\xdef\csname @shortt\roman{plate}\endcsname{\@shorttitle}
  % save the current value of the figure counter for later use in \outputplates
  \setcounter{figplate\roman{plate}}{\value{figure}}
  % increment the figure counter according to the optional argument
  \addtocounter{figure}{#1}
  }
%% page style for plate pages 
\newcommand{\ps@rmaaplate}{%
  % Plate X (see page Y) in the foot
  \renewcommand\@oddfoot{\hfil\csname @shorta\roman{plateloop}\endcsname\ 
    (see page~\csname @plateref\roman{plateloop}\endcsname)\hfil}%
  \renewcommand\@evenfoot{\relax}%
  % nothing in even heads (blank pages)
  \renewcommand\@evenhead{\relax}
  % short title and page num in the odd heads
  \renewcommand\@oddhead{\hfil\csname @shortt\roman{plateloop}\endcsname%
    \hfil\llap{PLATE~\theplateloop}}
  }





%% Things for the User Guide
\if@guide
%% stolen from doc.sty
\def\MakeShortVerb#1{%
  \expandafter\ifx\csname cc\string#1\endcsname\relax
  \@shortvrbinfo{Made }{#1}%
  \add@special{#1}%
  \expandafter
  \xdef\csname cc\string#1\endcsname{\the\catcode`#1}%
  \begingroup
  \catcode`\~\active  \lccode`\~`#1%
  \lowercase{%
    \global\expandafter\let
    \csname ac\string#1\endcsname~%
    \gdef~{\verb~}}%
  \endgroup
  \global\catcode`#1\active
  \else
  \@shortvrbinfo\@empty{#1 already}%
  \fi}
\def\DeleteShortVerb#1{%
  \expandafter\ifx\csname cc\string#1\endcsname\relax
  \else
  \@shortvrbinfo{Deleted }{#1 as}%
  \rem@special{#1}%
  \global\catcode`#1\csname cc\string#1\endcsname
  \global \expandafter\let \csname cc\string#1\endcsname \relax
  \ifnum\catcode`#1=\active
  \begingroup
  \catcode`\~\active   \lccode`\~`#1%
  \lowercase{%
    \global\expandafter\let\expandafter~%
    \csname ac\string#1\endcsname}%
  \endgroup \fi \fi}
\def\@shortvrbinfo#1#2{%
  \PackageInfo{doc}{%
    #1\expandafter\@gobble\string#2 a short reference
    for \string\verb}}
\def\add@special#1{%
  \rem@special{#1}%
  \expandafter\gdef\expandafter\dospecials\expandafter
  {\dospecials \do #1}%
  \expandafter\gdef\expandafter\@sanitize\expandafter
  {\@sanitize \@makeother #1}}
\def\rem@special#1{%
  \def\do##1{%
    \ifnum`#1=`##1 \else \noexpand\do\noexpand##1\fi}%
  \xdef\dospecials{\dospecials}%
  \begingroup
  \def\@makeother##1{%
    \ifnum`#1=`##1 \else \noexpand\@makeother\noexpand##1\fi}%
  \xdef\@sanitize{\@sanitize}%
  \endgroup}
%% define | as shortcut for \verbatim
\MakeShortVerb{\|}
%% comand to specify arguments to LaTeX commands
\newcommand{\rmArg}[1]{%
  \ensuremath{\langle}{\scriptsize \textsf{\MakeUpperCase{#1}}}\ensuremath{\rangle}}
%\newcommand{\rmArg}[1]{\ensuremath{\langle}\textsc{#1}\ensuremath{\rangle}}
\providecommand{\BibTeX}{{\rmfamily B\kern-.05em%
    \textsc{i\kern-.025em b}\kern-.08em%
    T\kern-.1667em\lower.7ex\hbox{E}\kern-.125emX}}
%% Environment for code examples
\newenvironment{Code}{%
  \begin{list}{}{\setlength{\leftmargin}{0.05\columnwidth}} 
  \item[]}{\end{list}}
\fi
%</package>
%    \end{macrocode}

% \section{User commands for astronomical articles}
% \label{sec:user-commands}
% Some random definitions. These go both into the main class file and
% also into |adshtml.sty|.
% 
% Definitions of the journal name in different incarnations.  
%    \begin{macrocode}
%<*package|adshtml-style>
\newcommand{\RMAA}{RMxAA}
\newcommand{\RMAAlong}{Revista Mexicana de Astronom\'{\i}a y Astrof\'{\i}sica}
\newcommand{\RMAASC}{RMxAC}
\newcommand{\RMAASClong}{\RMAA\ (Serie de Conferencias)}
%    \end{macrocode}
% \begin{macro}{\ibidrule}
% This puts a horizontal rule with a period after it. It is used when
% the author list of a reference is identical to the preceding one. 
%    \begin{macrocode}
\newcommand{\ibidrule}{\rule[-0.03cm]{4.7em}{.01cm}.\ }
%    \end{macrocode}
% \end{macro}
% \begin{macro}{\ThisVolume}
% References to other articles in the same volume can be coded as,
% e.g., |\bibitem{} Bloggs, J. \ThisVolume{bloggs}|.
%    \begin{macrocode}
\newcommand{\ThisVolume}[1]{%
  \@year, \if@proceedings\RMAASC\else\RMAA\fi, \@volume,
  \if@book\oldpageref{#1:RMAAFirstPage}\else 000\fi\ (this volume)}
%    \end{macrocode}
% \end{macro}
% Other stuff that may be useful in ref lists.
%    \begin{macrocode}
\newcommand{\etal}{et al.\@}
%    \end{macrocode}
% Various new commands from AAS macros
%    \begin{macrocode}
\def\acknowledgments{\vskip 3ex plus .8ex minus .4ex}
\let\acknowledgements=\acknowledgments                  % second spelling
\let\jnl@style=\rm
\def\ref@jnl#1{{\jnl@style#1}}
\DeclareRobustCommand\aj{\ref@jnl{AJ}}                   
\DeclareRobustCommand\araa{\ref@jnl{ARA\&A}}             
\DeclareRobustCommand\apj{\ref@jnl{ApJ}}                 
\DeclareRobustCommand\apjl{\ref@jnl{ApJ}}                
\DeclareRobustCommand\apjs{\ref@jnl{ApJS}}               
\DeclareRobustCommand\ao{\ref@jnl{Appl.~Opt.}}           
\DeclareRobustCommand\apss{\ref@jnl{Ap\&SS}}             
\DeclareRobustCommand\aap{\ref@jnl{A\&A}}                
\DeclareRobustCommand\aapr{\ref@jnl{A\&A~Rev.}}          
\DeclareRobustCommand\aaps{\ref@jnl{A\&AS}}              
\DeclareRobustCommand\azh{\ref@jnl{AZh}}                 
\DeclareRobustCommand\baas{\ref@jnl{BAAS}}               
\DeclareRobustCommand\jrasc{\ref@jnl{JRASC}}             
\DeclareRobustCommand\memras{\ref@jnl{MmRAS}}            
\DeclareRobustCommand\mnras{\ref@jnl{MNRAS}}             
\DeclareRobustCommand\pra{\ref@jnl{Phys.~Rev.~A}}        
\DeclareRobustCommand\prb{\ref@jnl{Phys.~Rev.~B}}        
\DeclareRobustCommand\prc{\ref@jnl{Phys.~Rev.~C}}        
\DeclareRobustCommand\prd{\ref@jnl{Phys.~Rev.~D}}        
\DeclareRobustCommand\pre{\ref@jnl{Phys.~Rev.~E}}        
\DeclareRobustCommand\prl{\ref@jnl{Phys.~Rev.~Lett.}}    
\DeclareRobustCommand\pasp{\ref@jnl{PASP}}               
\DeclareRobustCommand\pasj{\ref@jnl{PASJ}}               
\DeclareRobustCommand\qjras{\ref@jnl{QJRAS}}             
\DeclareRobustCommand\skytel{\ref@jnl{S\&T}}             
\DeclareRobustCommand\solphys{\ref@jnl{Sol.~Phys.}}      
\DeclareRobustCommand\sovast{\ref@jnl{Soviet~Ast.}}      
\DeclareRobustCommand\ssr{\ref@jnl{Space~Sci.~Rev.}}     
\DeclareRobustCommand\zap{\ref@jnl{ZAp}}                 
\DeclareRobustCommand\nat{\ref@jnl{Nature}}              
\DeclareRobustCommand\iaucirc{\ref@jnl{IAU~Circ.}}
\DeclareRobustCommand\aplett{\ref@jnl{Astrophys.~Lett.}}
\DeclareRobustCommand\apspr{\ref@jnl{Astrophys.~Space~Phys.~Res.}}
\DeclareRobustCommand\bain{\ref@jnl{Bull.~Astron.~Inst.~Netherlands}}
\DeclareRobustCommand\fcp{\ref@jnl{Fund.~Cosmic~Phys.}}
\DeclareRobustCommand\gca{\ref@jnl{Geochim.~Cosmochim.~Acta}}
\DeclareRobustCommand\grl{\ref@jnl{Geophys.~Res.~Lett.}}
\DeclareRobustCommand\jcp{\ref@jnl{J.~Chem.~Phys.}}      
\DeclareRobustCommand\jgr{\ref@jnl{J.~Geophys.~Res.}}    
\DeclareRobustCommand\jqsrt{\ref@jnl{J.~Quant.~Spec.~Radiat.~Transf.}}
\DeclareRobustCommand\memsai{\ref@jnl{Mem.~Soc.~Astron.~Italiana}}
\DeclareRobustCommand\nphysa{\ref@jnl{Nucl.~Phys.~A}}
\DeclareRobustCommand\physrep{\ref@jnl{Phys.~Rep.}}
\DeclareRobustCommand\physscr{\ref@jnl{Phys.~Scr}}
\DeclareRobustCommand\planss{\ref@jnl{Planet.~Space~Sci.}}
\DeclareRobustCommand\procspie{\ref@jnl{Proc.~SPIE}}
\let\astap=\aap
\let\apjlett=\apjl
\let\apjsupp=\apjs
\let\applopt=\ao
\def\phn{\phantom{0}}
\def\phd{\phantom{.}}
\def\phs{\phantom{\ensuremath{-}}}
\def\phm#1{\phantom{#1}}
%    \end{macrocode}
% All these are redone as robust and using more modern commands where
% possible. The robustness is so they will not get expanded if used in
% abstracts and titles when these are passed to the ADS file. 
% \changes{3.16}{2003/04/07}{
%  Added \cs{Color} command for doing photometric colors.}
%    \begin{macrocode}
%%\def\sun{\hbox{$\odot$}}
\DeclareRobustCommand\sun{\ensuremath{\odot}}
%%\def\earth{\hbox{$\oplus$}}
\DeclareRobustCommand\earth{\ensuremath{\oplus}}
%%\def\lesssim{\mathrel{\hbox{\rlap{\hbox{\lower4pt\hbox{$\sim$}}}\hbox{$<$}}}}
%%\def\gtrsim{\mathrel{\hbox{\rlap{\hbox{\lower4pt\hbox{$\sim$}}}\hbox{$>$}}}}
\DeclareRobustCommand\lesssim{\mathrel{\hbox{\smaller\smaller\rlap{\hbox{\lower3pt\hbox{$\sim$}}}\raise1.5pt\hbox{$<$}}}}
\DeclareRobustCommand\gtrsim{\mathrel{\hbox{\smaller\smaller\rlap{\hbox{\lower3pt\hbox{$\sim$}}}\raise1.5pt\hbox{$>$}}}}
%%\def\lta{\mathrel{\spose{\lower 3pt\hbox{$\mathchar"218$}}
%%     \raise 2.0pt\hbox{$\mathchar"13C$}}}
%%\def\gta{\mathrel{\spose{\lower 3pt\hbox{$\mathchar"218$}}
%%     \raise 2.0pt\hbox{$\mathchar"13E$}}}
\DeclareRobustCommand\sq{\hbox{\rlap{$\sqcap$}$\sqcup$}}
%%\def\arcdeg{\hbox{$^\circ$}}
%%\def\arcmin{\hbox{$^\prime$}}
%%\def\arcsec{\hbox{$^{\prime\prime}$}}
\DeclareRobustCommand\arcdeg{\ensuremath{^\circ}}
\DeclareRobustCommand\arcmin{\ensuremath{^\prime}}
\DeclareRobustCommand\arcsec{\ensuremath{^{\prime\prime}}}
%%\def\fd{\hbox{$.\!\!^{\rm d}$}}
%%\def\fh{\hbox{$.\!\!^{\rm h}$}}
%%\def\fm{\hbox{$.\!\!^{\rm m}$}}
%%\def\fs{\hbox{$.\!\!^{\rm s}$}}
%%\def\fdg{\hbox{$.\!\!^\circ$}}
%%\def\farcm{\hbox{$.\mkern-4mu^\prime$}}
%%\def\farcs{\hbox{$.\!\!^{\prime\prime}$}}
%%\def\fp{\hbox{$.\!\!^{\scriptscriptstyle\rm p}$}}
\DeclareRobustCommand\fd{\ensuremath{.\!\!^\mathrm{d}}}
\DeclareRobustCommand\fh{\ensuremath{.\!\!^\mathrm{h}}}
\DeclareRobustCommand\fm{\ensuremath{.\!\!^\mathrm{m}}}
\DeclareRobustCommand\fs{\ensuremath{\mkern0.5mu.\mkern-5mu^\mathrm{s}\mkern-1mu}}
\DeclareRobustCommand\fdg{\ensuremath{.\!\!^\circ}}
\DeclareRobustCommand\farcm{\ensuremath{.\mkern-4mu^\prime}}
\DeclareRobustCommand\farcs{\ensuremath{\mkern1mu.\mkern-6mu^{\prime\prime}\mkern-1.5mu}}
%%\DeclareRobustCommand\farcs{\ensuremath{.\!\!^{\prime\prime}}}
\DeclareRobustCommand\fp{\ensuremath{.\!\!^{\scriptscriptstyle\rm p}}}
%%\def\micron{\hbox{$\mu$m}}
\DeclareRobustCommand\micron{\ensuremath{\mu\mathrm{m}}}
\let\la=\lesssim                        
\let\ga=\gtrsim
%%\def\case#1#2{\hbox{$\frac{#1}{#2}$}}
%%\def\slantfrac#1#2{\hbox{$\,^#1\!/_#2$}}
%%\def\onehalf{\slantfrac{1}{2}}
%%\def\onethird{\slantfrac{1}{3}}
%%\def\twothirds{\slantfrac{2}{3}}
%%\def\onequarter{\slantfrac{1}{4}}
%%\def\threequarters{\slantfrac{3}{4}}
\DeclareRobustCommand\case[2]{\ensuremath{\frac{#1}{#2}}}
\DeclareRobustCommand\slantfrac[2]{\ensuremath{\,^#1\!/_#2}}
\DeclareRobustCommand\onehalf{\slantfrac{1}{2}}
\DeclareRobustCommand\onethird{\slantfrac{1}{3}}
\DeclareRobustCommand\twothirds{\slantfrac{2}{3}}
\DeclareRobustCommand\onequarter{\slantfrac{1}{4}}
\DeclareRobustCommand\threequarters{\slantfrac{3}{4}}
%%\def\ubvr{\hbox{$U\!BV\!R$}}            
%%\def\ub{\hbox{$U\!-\!B$}}               
%%\def\bv{\hbox{$B\!-\!V$}}               
%%\def\vr{\hbox{$V\!-\!R$}}               
%%\def\ur{\hbox{$U\!-\!R$}}
%    \end{macrocode}
% New commands for ``photometric'' colors 
%    \begin{macrocode}
\DeclareRobustCommand\ubvr{\ensuremath{U\!BV\!R}}            
\DeclareRobustCommand\Color[2]{%
  \ensuremath{%
    #1\!\ifx#1V\!\fi-\!#2
  }}
\DeclareRobustCommand\ub{\Color{U}{B}}
\DeclareRobustCommand\bv{\Color{B}{V}}
\DeclareRobustCommand\vr{\Color{V}{R}}
\DeclareRobustCommand\ur{\Color{U}{R}}
%% New version of the \ion command, that actually works! 
%    \end{macrocode}
% \begin{macro}{\ion}
% This version of the |\ion| command allows you to write either
% |\ion{N}{4}| or |\ion{N}{iv}| in order to get N~IV. 
%    \begin{macrocode}
\DeclareRobustCommand\ion[2]{%
  \mbox{#1\kern0.2em%
    \smaller\rmfamily%
    \edef\@tempa{\@car#2\@nil}%
    \ifcat1\@tempa%
    \@Roman{#2}%
    \else%
    \uppercase{#2}%
    \fi}}
\newcounter{IonStage}
\DeclareRobustCommand\plainion[2]{\setcounter{IonStage}{#2}#1 \Roman{IonStage}}
%%\def\ion#1#2{\ensuremath{\mathrm{#1}\;}{\protect\small\rm\@Roman{#2}}}
%    \end{macrocode}
% \end{macro}
%    \begin{macrocode}
\DeclareRobustCommand\nodata{$\cdots$}
%    \end{macrocode}
% These are needed by the |\ion| command. Isn't there an easier way of
% doing this?
%    \begin{macrocode}
\let\@small=\small
\let\@normalsize=\normalsize
\let\@footnotesize=\footnotesize
\let\@scriptsize=\scriptsize
\let\@tiny=\tiny
\let\@large=\large
\let\@Large=\Large
\let\@LARGE=\LARGE
\let\@huge=\huge
\let\@Huge=\Huge
\newcommand{\smaller}{\small}
\renewcommand{\normalsize}{\renewcommand{\smaller}{\small}\@normalsize}
\renewcommand{\small}{\renewcommand{\smaller}{\footnotesize}\@small}
\renewcommand{\footnotesize}{\renewcommand{\smaller}{\scriptsize}\@footnotesize}
\renewcommand{\scriptsize}{\renewcommand{\smaller}{\tiny}\@scriptsize}
\renewcommand{\tiny}{\renewcommand{\smaller}{\tiny}\@tiny}
\renewcommand{\large}{\renewcommand{\smaller}{\normalsize}\@large}  
\renewcommand{\Large}{\renewcommand{\smaller}{\large}\@Large}  
\renewcommand{\LARGE}{\renewcommand{\smaller}{\Large}\@LARGE}  
\renewcommand{\huge}{\renewcommand{\smaller}{\LARGE}\@huge}  
\renewcommand{\Huge}{\renewcommand{\smaller}{\huge}\@Huge}  
%</package|adshtml-style>
%    \end{macrocode}


% \section{Other files}

% \subsection{The driver and make files}
% \label{sec:driver_make}
% These are put here for completeness, although they still need to be
% distributed separately since otherwise the user won't be able to
% unpack this file into its constituent parts.
% 
% The |Makefile| has two principal rules: |docs|, and |macros|. The
% |docs| rule produces a PDF and gzipped PS file of the documentation.
% I still haven't perfected the combination of rules required to get
% it to run \LaTeX{} enough times in all circumstances. The |macros|
% rule runs \LaTeX{} on |rmaa.ins| in order to unpack |rmaa.cls| and
% all the other little files (including the |Makefile| and the
% |rmaa.ins| files themselves\dots 
% 
% First set up lists of various types of file. 
%    \begin{macrocode}
%<*makefile>
DOCFILES=rmaa.pdf
MACRO_LTX_FILES=rmaa.cls rmaacite.sty adshtml.sty
SHELLSCRIPTS=makeads.sh
MACRO_OTHER_FILES=adshtml.perl sample.latex2html-init
MAKEFILE=Makefile
MACROFILES=$(MACRO_LTX_FILES) $(SHELLSCRIPTS) $(MACRO_OTHER_FILES)         
SKELETONFILES=ADS.tex 
%    \end{macrocode}
% Install the files read only into the package directory, specified via a path relative to the source directory. \changes{4.0}{2019/11/14}{Reorganize folder layout so that hardly anything is in the top level.  That way, things are clearer on github.}
%    \begin{macrocode}
INSTALL=install 
INSTALL_DIR=../packages
%    \end{macrocode}
% The default rule is to make everything. Installation must be
% explicitly asked for.
%    \begin{macrocode}
all: docs macros

install: install-docs install-macros

%    \end{macrocode}
% Most files are installed read-only to prevent editing. Skeleton
% files must be writeable. Must be careful not to overwrite the
% Makefile in the parent directory. 
%    \begin{macrocode}
install-macros: macros
        $(INSTALL) -m 444 $(MACRO_LTX_FILES) $(MACRO_OTHER_FILES) $(INSTALL_DIR)
        $(INSTALL) -m 555 $(SHELLSCRIPTS) $(INSTALL_DIR)
        $(INSTALL) -m 644 $(SKELETONFILES) $(INSTALL_DIR)

install-docs: docs
        $(INSTALL) rmaa.pdf $(INSTALL_DIR)/rm-fulldocs.pdf

%    \end{macrocode}
% Make the PDF and gzipped PS files. 
%    \begin{macrocode}
docs: $(DOCFILES)

%    \end{macrocode}
% Make everything else if either the |.dtx| or |.ins| files have
% changed. 
%    \begin{macrocode}
macros: rmaa.dtx rmaa.ins
%    \end{macrocode}
% Ensure that all the macro files can be written to then recreate
% them. The |touch| is there in case any of the files don't exist
% yet. 
%    \begin{macrocode}
        touch $(MACROFILES) $(MAKEFILE)
        chmod u+w $(MACROFILES) $(MAKEFILE) rmaa.ins
        latex "\input{rmaa.ins}"
%    \end{macrocode}
% \TeX{} automatically adds the |.tex| extension to extensionless
% files. We fix this and take the opportunity to convert
% back to TABs all that it has lovingly converted to spaces. 
%    \begin{macrocode}
        cp -f $(MAKEFILE) $(MAKEFILE).BCK && unexpand $(MAKEFILE).tex >| $(MAKEFILE) && rm $(MAKEFILE).tex
%    \end{macrocode}
% Make the shell scripts executable and everything unwriteable (so
% they don't get hand-edited by mistake). However, undo this for the
% files that the user is \emph{supposed} to modify. 
%    \begin{macrocode}
        chmod a+x $(SHELLSCRIPTS)
        chmod a-w $(MACROFILES) $(MAKEFILE)
        chmod u+w $(SKELETONFILES)

%    \end{macrocode}
% We now use PDF\LaTeX{} for this.
%    \begin{macrocode}
rmaa.idx: rmaa.dtx
        @echo "REGENERATING .idx FILE"
        pdflatex "\nonstopmode\input{rmaa.dtx}"

rmaa.aux: rmaa.dtx rmaa.ind rmaa.gls
        @echo "REGENERATING .aux FILE"
        pdflatex "\nonstopmode\input{rmaa.dtx}"

rmaa.ind: rmaa.idx
        @echo "REGENERATING INDEX"
        makeindex -s gind.ist rmaa

rmaa.gls: rmaa.glo
        @echo "REGENERATING CHANGELOG"
        makeindex -s gglo.ist -o rmaa.gls rmaa.glo

rmaa.pdf: rmaa.aux rmaa.ind rmaa.gls
        @echo "REGENERATING .pdf FILE"
        pdflatex "\nonstopmode\input{rmaa.dtx}"

clean: 
        rm -f rmaa.aux rmaa.dvi rmaa.ind rmaa.gls
%</makefile>
%    \end{macrocode}

% The |rmaa.ins| file uses the \textsf{docstrip} package to split out
% all the files and strip off this documentation.
%    \begin{macrocode}
%<*ins>
%%
%% Installation batch file for the RMAA LaTeX package
%%
\def\batchfile{rmaa.ins}
\input docstrip.tex
\preamble

PLEASE DO NOT EDIT THIS FILE! It is generated automatically from
the file rmaa.dtx using the accompanying Makefile and rmaa.ins 
Any hand-edited changes will be lost

A LaTeX document class for the RMxAA

Author: William Henney (w.henney@irya.unam.mx)

Latest version available from: http://www.irya.unam.mx/rmaa

Copyright (C) 1999-2019 Instituto de Astronomia, UNAM, Mexico

Permission is hereby granted, free of charge, to any person obtaining
a copy of this software and associated documentation files (the
"Software"), to deal in the Software without restriction, including
without limitation the rights to use, copy, modify, merge, publish,
distribute, sublicense, and/or sell copies of the Software, and to
permit persons to whom the Software is furnished to do so, subject to
the following conditions:

The above copyright notice and this permission notice shall be
included in all copies or substantial portions of the Software.

THE SOFTWARE IS PROVIDED "AS IS", WITHOUT WARRANTY OF ANY KIND,
EXPRESS OR IMPLIED, INCLUDING BUT NOT LIMITED TO THE WARRANTIES OF
MERCHANTABILITY, FITNESS FOR A PARTICULAR PURPOSE AND
NONINFRINGEMENT. IN NO EVENT SHALL THE AUTHORS OR COPYRIGHT HOLDERS BE
LIABLE FOR ANY CLAIM, DAMAGES OR OTHER LIABILITY, WHETHER IN AN ACTION
OF CONTRACT, TORT OR OTHERWISE, ARISING FROM, OUT OF OR IN CONNECTION
WITH THE SOFTWARE OR THE USE OR OTHER DEALINGS IN THE SOFTWARE.
\endpreamble
%    \end{macrocode}
% Clobber files without asking
%    \begin{macrocode}
\askforoverwritefalse
%    \end{macrocode}
% Make the \LaTeX{} macro files (class and styles)
%    \begin{macrocode}
\generate{% 
  \file{rmaa.cls}{\from{rmaa.dtx}{package}}
  \file{rmaa.ins}{\from{rmaa.dtx}{ins}}
  \file{rmaacite.sty}{\from{rmaa.dtx}{rmaacite}}
  \file{adshtml.sty}{\from{rmaa.dtx}{adshtml-style}}
  \file{ADS.tex}{\from{rmaa.dtx}{ads-container}}
  }
%    \end{macrocode}
% Make the other non-\LaTeX{} files. These are done separately since
% they can't use `\%\%' as a comment character for the
% preamble/postamble (and Perl files can't have the postamble I don't
% think). 
%    \begin{macrocode}
\generate{\nopreamble\nopostamble% 
  \file{Makefile}{\from{rmaa.dtx}{makefile}}
  \file{adshtml.perl}{\from{rmaa.dtx}{adshtml-perl}}
  \file{makeads.sh}{\from{rmaa.dtx}{ads-shell-script}}
  \file{sample.latex2html-init}{\from{rmaa.dtx}{latex2html-init}}
  }

%% Local Variables: 
%% mode: LaTeX
%% mode: Font-lock
%% End:
%</ins>
%    \end{macrocode}

% \subsection{Extended citation syntax}
% Modified from a file called mnrascite.sty
% I no longer remember where I stole the original version of this from
% 
% Alterations to the cite commands to make them do more sensible things.
% N.B. the bibliography style should produce entries which have labels where
% the date is delimited by angle brackets '<>' for these macros to work
% properly. They will not give very sensible error messages if you do not do 
% this properly.
% 
% A new command |\scite{key}| is defined which produces a citation which is 
% suitable for insertion in a sentence. i.e., it does not put brackets around 
% the whole citation, only the date. Also |\pcite{key}| produce output with no
% parentheses whatsoever. 
% 
% Should redefine |\@bibitem| really to have any hope of compatibility
% the command |\lbibcite| is defined to deal with the citations.
%    \begin{macrocode}
%<*rmaacite>
\def\bibitem[#1<#2>]#3{\par\hangindent=1.5em\hangafter=1
  \if@filesw
  {\def\protect##1{\string ##1\space}
    \immediate\write\@auxout{\string\lbibcite{#3}{#1}{#2}}}
  \fi}
\@redefinebibitem % what it says
%    \end{macrocode}
% Note that the |\lbibcite| defines two types of citation one contains the names,
% the other the year
%    \begin{macrocode}
\def\lbibcite#1#2#3{\global\@namedef{b@#1}{#2}
                    \global\@namedef{by@#1}{#3}}

\def\@bcite#1#2{({#1\if@tempswa , #2\fi})}
\def\@pcite#1#2{#1\if@tempswa , #2\fi}
%    \end{macrocode}
% This declaration contains what is put between the name and the date
% \changes{3.11}{2003/01/31}{Moved the \cs{articlename} to after the key for consistency
% with the new natbib stuff.}
%    \begin{macrocode}
\def\@brcite{\ (}
\def\@ebrcite{)}
\def\@fmttop#1{\@fmtcite[\csname #1\endcsname]}
\def\@citex[#1]#2{%
%%  \if@filesw\immediate\write\@auxout{\string\citation{#2}}\fi
%% WJH added 16 Jan 2002
  \@for\@citeb:=#2\do{%
    \edef\@citeb{\expandafter\@firstofone\@citeb:\articlename}%
    \if@filesw\immediate\write\@auxout{\string\citation{\@citeb}}\fi}%
%% to here
  \def\@citea{}%
  \@cite{\@for\@citeb:=#2\do
    {%
      \edef\@citeb{\expandafter\@firstofone\@citeb:\articlename}%
      \@citea\def\@citea{;\ }\@ifundefined
      {b@\@citeb}{{\bf ?}\@warning {Citation `\@citeb' on page \thepage \space undefined}}%
      {\csname b@\@citeb\endcsname\@citebro\csname by@\@citeb\endcsname\@citebrc}}}%
  {#1}}

\def\cite{\@ifnextchar [{\let\@citebro=\space
                          \let\@citebrc=\relax
                          \let\@cite=\@bcite\@tempswatrue\@citex}
                        {\let\@citebro=\space
                          \let\@citebrc=\relax
                          \let\@cite=\@bcite\@tempswafalse\@citex[]}}
\def\pcite{\@ifnextchar [{\let\@citebro=\space
                          \let\@citebrc=\relax
                          \let\@cite=\@pcite\@tempswatrue\@citex}
                        {\let\@citebro=\space
                          \let\@citebrc=\relax
                          \let\@cite=\@pcite\@tempswafalse\@citex[]}}
\def\scite{\@ifnextchar [{\let\@citebro=\@brcite
                          \let\@citebrc=\@ebrcite
                          \let\@cite=\@pcite\@tempswatrue\@citex}
                        {\let\@citebro=\@brcite
                          \let\@citebrc=\@ebrcite
                          \let\@cite=\@pcite\@tempswafalse\@citex[]}}

\renewenvironment{thebibliography}
  {\section*{REFERENCES}\vspace*{\afterbiblioskip}\small
    \if@book\RereadBiblio\fi
    \sloppy\clubpenalty4000\widowpenalty4000
    \sfcode`\.=1000\relax
    \setlength\parindent{\z@}\setlength\parskip{\z@}}
  {\relax\normalsize\@outputaddressesifnecessary} 
%</rmaacite>
%    \end{macrocode}

% \subsection{Auxilliary files for generating ADS data}

% The ADS data are written by \LaTeX{} to a |.ads| file. This file
% then has to be included in the following file (|ADS.tex|),
% which is run through \LaTeX{}2HTML (see below). 

%    \begin{macrocode}
%<*ads-container>
%%
%% Instructions:
%% Edit the \input{} line below to reflect the name of your file. 
\documentclass{article}
\usepackage[latin1]{inputenc}
\usepackage{adshtml}
\newcommand{\ADSfield}[2]
{\%#1  #2\\}
%%\newcommand{\ADSvfield}[2]
%%{\%#1 \begin{verbatim}#2\end{verbatim}}
\newenvironment{ADSarticle}{}{\par\bigskip}
\renewcommand{\bibitem}[2][]{\par}
\setlength{\parindent}{0pt}
\renewcommand{\ThisVolume}[1]{%  
  2001, \RMAASC, 12, \oldpageref{#1:RMAAFirstPage}}  
 

\begin{document}
\makeatletter
\input{rmbooksamp.ads}
\end{document}
%</ads-container>
%    \end{macrocode}

% This is the |adshtml.sty| style file used by |ADS.tex|. It just
% defines a few 
% macros that might be used in abstracts, etc (really it should define
% them all). This now gets all its code from
% Sec~\ref{sec:user-commands}. 
%    \begin{macrocode}
%<*adshtml-style>
%% Nothing else needs to go here now. 
%</adshtml-style>
%    \end{macrocode}
% The init file for latex2html.
%    \begin{macrocode}
%<*latex2html-init>
$PDFTOCAIRO = "pdftocairo";
$PDFLATEX = "pdflatex";
$LATEX = "latex";
$MAX_SPLIT_DEPTH = 0;
$HTML_VERSION = "4.0,math";
$NO_SIMPLE_MATH = 1;
$ADDRESS = "<p><address><a href=\"mailto:rmaa\@astro.unam.mx\">rmaa\@astro.unam.mx</address></p>"; 
$INFO = "<p>This document was created using latex2html and the rmaa latex macros</p>"; 
$NO_NAVIGATION = 1;
1; 
%</latex2html-init>
%    \end{macrocode}
% This is the |adshtml.perl| file that is used by \LaTeX{}2HTML. Some
% bizarre \TeX{}y command sometimes show up so we tell it to ignore
% these. All the astronomy macros that come from AAS are processed in
% \TeX{}. That way they get made to images but with an |ALT| tag that
% has the original code, which is what will be picked up by |lynx|,
% with the result that they get passed intact into the final |.txt|
% file. 
%    \begin{macrocode}
%<*adshtml-perl>
package main;

$dum = &do_cmd_makeatletter();

sub do_cmd_RMAA {
    local($_) = @_;
    "RMxAA" . $_;
}

sub do_cmd_RMAASC {
    local($_) = @_;
    "RMxAASC" . $_;
}

sub do_cmd_etal {
    local($_) = @_;
    "et al\." . $_;
}

sub do_cmd_ibidrule {
    local($_) = @_;
    "<img src=\"foo\" alt=\"\ibidrule\">" . $_;
}


sub do_cmd_bibitem {
    local($_) = @_;
    local($key);
    local($opt, $dummy)=&get_next_optional_argument; 
    $key = &missing_braces unless
        ((s/$next_pair_pr_rx/$key = $2; ''/eo)
        ||(s/$next_pair_rx/$key = $2; ''/eo));
    "<p>" . $_; 
}

&ignore_commands (<<_RAW_ARG_CMDS_);
lowercase
adjustfinalcols
textsuperscript # {}
newline # []
thinmuskip
mskip
smallskip
medskip
bigskip
hspace # {}
hspacestar # {}
vspace # {}
vspacestar # {}
mathrm
_at_mathrm
spacefactor
discretionary
unhbox 
voidb_at_x 
penalty 
_at_M
_at_m
null
IeC
_RAW_ARG_CMDS_

&process_commands_in_tex (<<_RAW_ARG_CMDS_);
sun
earth
lesssim
gtrsim
sq
argdeg
arcmin
arcsec
fd
fh
fm
fs
fdg
farcm
farcs
fp
micron
la
ga
case # {} # {}
slantfrac # {} # {}
onehalf
onethird
twothirds
onequarter
threequarters
ubvr
ub
bv
vr
ur
aj
araa
apj
apjl
apjs
ao
apss
aap
aapr
aaps
azh
baas
jrasc
memras
mnras
pra
prb
prc
prd
pre
prl
pasp
pasj
qjras
skytel
solphys
sovast
ssr
zap
nat
iaucirc
aplett
apspr
bain
fcp
gca
grl
jcp
jgr
jqsrt
memsai
nphysa
physrep
physscr
planss
procspie
_RAW_ARG_CMDS_

1;                              # This must be the last line
%</adshtml-perl> 
%    \end{macrocode}

% This is a shell script |makeads.sh| that generates the |_ads.txt| file from
% |_ads.tex| by using \LaTeX{}2HTML and |lynx|.  \changes{4.0}{2019/11/13}{Make sure that lynx writes utf-8 output instead of latin-1.}
% \changes{2.2beta13}{2002/03/12}{Complete rewrite of |makeads.sh|. Now much 
% improved I think}
% \changes{4.0}{2019/11/13}{Fix the at-sign substitution so that it only eliminates \cs{@} and does not mess with email addresses.}

%    \begin{macrocode}
%<*ads-shell-script>
#! /bin/sh
if [ -z $1 ]; then
    echo "Usage: $0 basename"
    exit
fi
BASENAME=$1
ADSBASENAME=${BASENAME}_ads
AUXFILE=$BASENAME.aux
TMPFILE=tmp.$$


echo "Generating $ADSBASENAME.tex"

# For some reason, must double-up \\ to \\\\ but not \ to \\ - odd!
cat > $ADSBASENAME.tex <<EOF
%% PLEASE DO NOT EDIT THIS FILE!  It is generated automatically from
%% the script makeads.sh.  Any hand-edited changes will be lost
%% 
\documentclass{article}
\usepackage[latin1]{inputenc}
\usepackage{html}
\usepackage{adshtml}
\newcommand{\ADSfield}[2]{\%#1  #2\\\\}
\newenvironment{ADSarticle}{}{\par\bigskip}
\setlength{\parindent}{0pt}
\newcounter{istage}
\renewcommand{\ion}[2]{\setcounter{istage}{#2}#1 \Roman{istage}}
\begin{document}
\makeatletter
EOF
echo >> $ADSBASENAME.tex "\input{$BASENAME.ads2}" 
echo >> $ADSBASENAME.tex "\end{document}"



echo "Generating $ADSBASENAME.sed" 

# Find type of journal
grep '\\documentclass\[.*proceedings.*\]' $BASENAME.tex >/dev/null \
     && JOURNAL='RMxAC' \
     || JOURNAL='RMxAA'

# Find the year and volume number
YEAR=`grep '\\SetYear' $BASENAME.tex | sed 's/.*\\SetYear{\(.*\)}.*/\1/'`
VOLUME=`grep '\\SetVolume' $BASENAME.tex | sed 's/.*\\SetVolume{\(.*\)}.*/\1/'`

# Generate an auxilliary sed script 
cat > $TMPFILE <<EOF
s|\\\\newlabel{\(.*\):RMAAFirstPage}{{.*}{\(.*\)}}|s/\\\\\\\\ThisVolume{\1}/$YEAR, RMxAC, $VOLUME, \2/g|g
EOF

# Now go through the .aux file, looking for \newlabels and use the sed 
# script just generated in order to make another sed script, one that 
# will replace \ThisVolume{LABEL} with the correct reference
grep '^\\newlabel{.*:RMAAFirstPage}' $AUXFILE \
     | sed -f $TMPFILE > $ADSBASENAME.sed && rm $TMPFILE

# Apend a further sed command to get rid of \@, since redefining 
# it in _ads.tex didn't work
echo >> $ADSBASENAME.sed 's/\\\\@//g'

echo "Generating $BASENAME.ads2 from $ADSBASENAME.sed and $BASENAME.ads" 

# Apply this to the .ads file, generating the .ads2 file, which is included
# by the _ads.tex file.
sed -f $ADSBASENAME.sed $BASENAME.ads > $BASENAME.ads2

echo "Generating $ADSBASENAME.txt from $ADSBASENAME.tex and $BASENAME.ads2"

# Now run latex2html on the _ads.tex file
latex2html -no_navigation -info 0 -address 0 -html_version 4.0,math $ADSBASENAME

# Then, convert to utf8 with lynx and tidy it up a bit
lynx --justify=off --width=100 --dump --nolist --display_charset=utf-8 \
     $ADSBASENAME/$ADSBASENAME.html | cut -c4- \
     | sed 's/ibidrule/\\ibidrule/g' > $ADSBASENAME.txt
%</ads-shell-script>
%    \end{macrocode}

% \section{Epilogue}
% Err, that's it \dots

\endinput


% \endinput
% Local Variables:
% mode: doctex
% TeX-master: t
% End:
"

rmaa.aux: rmaa.dtx rmaa.ind rmaa.gls
        @echo "REGENERATING .aux FILE"
        pdflatex "\nonstopmode% \def\fileversion{4.0} \def\filedate{2019/11/12}
% \MakeShortVerb{\|}
% \title{A \LaTeX{} document class for the Revista Mexicana de
% Astronom\'\i a y Astrof\'\i sica\thanks{This file has version
% \fileversion{} last revised \filedate{}.}}
% \author{William J. Henney\thanks{w.henney@irya.unam.mx}}
% \date{\filedate}
% \maketitle
% \paragraph{Abstract}
% This document describes the \texttt{rmaa} \LaTeX{}
% document class for the journal ``Revista Mexicana de
% Astronom\'\i a y Astrof\'\i sica''. The macros can be used for 
% \begin{enumerate}
% \item Preparing author manuscripts for the journal (including
% conference proceedings to be published in the ``Serie de
% Conferencias'')
% \item Typesetting an entire volume of the journal by inputting the
% individual article files from a container file. Facilities are
% included for automatic generation of TOC and Author Index. 
% \end{enumerate}
% 
% \tableofcontents
% \section{The documentation driver file}
% This is the part that generates the documentation that you may or
% may not be reading, which is done with the help of the Makefile
% discussed in section~\ref{sec:driver_make}. 
%    \begin{macrocode}
%<*driver>
\documentclass{ltxdoc}
\usepackage{hyperref}
\usepackage{palatino}
\EnableCrossrefs \CodelineIndex \RecordChanges
\setlength{\parindent}{0pt}
\def\MacroFont{\footnotesize\tt}
\begin{document} \DocInput{rmaa.dtx} \PrintIndex \PrintChanges \end{document}
%</driver>
%    \end{macrocode}

% \section{The implementation}
% 
% \changes{3.24}{2007/09/14}{Added user access to the macro version
% and date via \cs{rmVersion} and \cs{rmDate}}
% 
%    \begin{macrocode}
%<*package>
\newcommand{\rmaa@version}{3.26}
\newcommand{\rmaa@date}{2011/08/15}
\newcommand\rmVersion{\rmaa@version}
\newcommand\rmDate{\rmaa@date}
\NeedsTeXFormat{LaTeX2e}[1995/06/01]
%    \end{macrocode}
% \changes{4.0}{2019/11/02}{Change abstract to appear before
% resumen. Change RevMexAA to RMxAA for consistency with ADS.}
%
% \changes{3.22}{2007/03/22}{Allow proceedings articles to have
%   received/accepted dates} \changes{3.21}{2006/05/15}{Changes for
%   Silvia - abstracts are now shipped out to ADS.}
% \changes{3.19}{2006/01/19}{Changes for Silvia - remove adresses from
%   end, new abstract style} \changes{2.2beta}{2001/10/30}{ Fixed bug
%   whereby twocolumn option was not passed through to book
%   class. Finally sorted out the ``label on last page'' issue. }
% \changes{2.1}{2001/01/23}{ Fixed some inconsistencies in
%   \cs{nocite} and \cs{bibitem} uniqification Fixed the definition of
%   \cs{ion} so it will work in math mode and in contexts where the
%   typesize is different from normal (e.g., title, fig captions).
%   Modified the definition of \cs{@caption}, so we don't write any
%   list-of-figures data (we never use it anyway, and it just causes
%   problems with fragile commands in fig captions). }
% \changes{2.0}{2000/11/30}{ Added `newstyle' option, which should be
%   used in all new proceedings volumes (wider abstract, smaller
%   typeface in fig captions, etc.) Added commands for formatting
%   table captions. New demonstation document rmpostsamp.tex for
%   1-page posters. This will be the last release under a PRI
%   government (for a while...)}  \changes{1.9}{2000/06/06}{
%   Corrected position of page number in footer of articles that open
%   on even pages. Ensured that first article in each session opens on
%   an odd page, as do list of abstracts and author index. Added extra
%   space in TOC before the backmatter. Added command for editors,
%   \cs{BoldSessionTitles}, that makes produces bold titles for
%   session in the TOC.}  \changes{1.8}{2000/05/16}{ Fixed bug in
%   last page label with \cs{suppressfulladdresses} }
% \changes{1.7}{2000/05/11}{ Fixed more little things. Page labels
%   now work, even if you have no bibliography. You can output the
%   full addresses by hand with the \cs{outputfulladdresses}
%   command. Editors' Guide updated but still not finished. Revamped
%   splitbook.sh now produces PS file of each article for web page and
%   also an EPS file of each page for sending to the printers.}
% \changes{1.6}{2000/04/15}{ More cleaning up: full stops after
%   titles in index, running headers and footers in frontmatter,
%   fontsize in title of index, etc} \changes{1.5}{2000/04/15}{
%   Eliminated superscripts from titles in TOC and index}
% \changes{1.4}{2000/04/14}{ Added copyright and license info (GPL)}
% \changes{1.3}{2000/03/01}{ Corrected \cs{figcaption} to do the
%   caption formatting right Corrected font size of sec headings in
%   front matter} \changes{1.2}{2000/02/16}{ Added \cs{plate}
%   command} \changes{1.1}{2000/02/14}{ Added manuscript option,
%   corrected bugs} \changes{1.0}{2000/02/08}{ Now does a whole book
%   - lots of stuff tidied up} \changes{0.7}{1999/10/04}{ Increased
%   spacing between sections in single column styles. Suppressed
%   floats at bottom of pages (didn't work though) }
% \changes{0.6}{1999/09/24}{ Tweaked \cs{adjustfinalcols} (still not
%   perfect) Added proceedings option Added onecolumn/twocolumn
%   options (orthogonal to journal/proceedings) }
% \changes{0.5}{1999/09/23}{ Added \cs{fulladdresses} and
%   \cs{adjustfinalcols} (just about working)}
% \changes{0.4}{1999/09/22}{ Minor stylistic changes}
% \changes{0.3}{1999/06/28}{ Better mechanism for alternative
%   affiliations Added `guide' option for user guide}
% \changes{0.2}{1999/05/27}{ All basic functionality present, I
%   think} \changes{0.1}{1999/05/26}{ Just about working - still
%   missing Fig/Tab/Refs} \changes{0.0}{1999/05/21}{ So bored today,
%   decided to rewrite RMAA LaTeX macros from scratch}
%    \begin{macrocode}
\ProvidesClass{rmaa}[\rmaa@date\space v\rmaa@version\space RMxAA article
class]
%    \end{macrocode}
% \subsection{Option Handling}
% \label{sec:options}
% Declarations that are needed in options
% \changes{3.05}{2003/10/22}{Added natbib option}
% \changes{3.06}{2003/01/03}{Natbib option turned on permanently}
% \changes{3.07}{2003/01/11}{Added proof option}
% \changes{3.14}{2003/01/11}{Option to turn off natbib}
% \changes{3.18}{2003/10/22}{trying out doctex mode} 
% \changes{3.23}{2007/09/09}{Manuscript is now on by default.}
% \changes{3.26}{2011/08/15}{New option for review articles}
%    \begin{macrocode}
\newif\if@reviewarticle
\@reviewarticlefalse
\newif\if@proof
\@prooffalse
\newif\if@journal
\@journaltrue
\newif\if@printdates
\@printdatesfalse
\newif\if@proceedings
\@proceedingsfalse
\newif\if@manuscript
\@manuscripttrue
\newif\if@bigfont
\@bigfonttrue
\newif\if@guide
\@guidefalse
\newif\if@preprint
\@preprintfalse
\newif\if@rmtwocol
\@rmtwocoltrue
\newif\if@biblist
\@biblistfalse
\newif\if@book
\@bookfalse
\newif\if@natbib
\@natbibtrue
\newif\if@rmcite
\@rmcitefalse
\newif\if@titlevisible
\@titlevisibletrue
\newif\if@abstractpresent
\@abstractpresentfalse
\newif\if@printbanner
\@printbannerfalse
\newif\if@suppressfulladdresses
\@suppressfulladdressesfalse
\newif\if@newstyle
\@newstylefalse
\newif\if@debug
\@debugfalse
\newif\if@boldkeywords
\@boldkeywordsfalse
%    \end{macrocode}
% The option handling---mainly setting various booleans.
% \changes{2.2beta11}{2001/11/26}{Made twocol and
% newstyle the default for proceedings.}
% \changes{3.23}{2007/09/09}{Manuscript is now 10 pt instead of 12pt}
% \changes{3.23}{2007/09/10}{Full addresses are now automatically
% suppressed for proceedings style.}
% \changes{3.24}{2007/09/14}{New option \texttt{boldkeywords} that
% does the keywords in the old ugly way. Now off by default.}
%    \begin{macrocode}
\DeclareOption{manuscript}{\typeout{Manuscript for referee} 
  \@manuscripttrue\@journalfalse\@rmtwocolfalse
  \PassOptionsToClass{10pt}{article}}
\DeclareOption{bigfont}{\typeout{Big font manuscript requested by Gary
  Ferland} 
  \@manuscripttrue\@journalfalse\@rmtwocolfalse\@bigfonttrue
  \PassOptionsToClass{11pt}{article}}
\DeclareOption{journal}{\typeout{Main journal option}%
  \@journaltrue\@rmtwocoltrue\@printdatestrue\@manuscriptfalse
  \PassOptionsToClass{10pt}{article}}
\DeclareOption{guide}{\typeout{User guide option}%
  \@guidetrue\@journalfalse\@rmtwocoltrue\@manuscriptfalse
  \PassOptionsToClass{10pt}{article}}
\DeclareOption{preprint}{\typeout{Preprint option}%
  \@preprinttrue\@manuscriptfalse\PassOptionsToClass{10pt}{article}}
\DeclareOption{proof}{\typeout{Author proof option}%
  \@prooftrue\@manuscriptfalse}
\DeclareOption{proceedings}{\typeout{Proceedings option}%
  \@proceedingstrue\@journalfalse\@rmtwocoltrue\@newstyletrue
  \@suppressfulladdressestrue
  \@printdatesfalse\@manuscriptfalse\PassOptionsToClass{10pt}{article}} 
\DeclareOption{printdates}{\typeout{Force printing of received/accepted dates}%
  \@printdatestrue} 
\DeclareOption{book}{\typeout{Whole issue option}%
  \@booktrue} 
\DeclareOption{nonatbib}{\typeout{Turn off natbib support (required
    for compatibility with old rmaacite package).}%
  \@natbibfalse} 
\DeclareOption{newstyle}{\typeout{Use new layout style}%
  \@newstyletrue}
\DeclareOption{oldstyle}{\typeout{Use old layout style}%
  \@newstylefalse}
\DeclareOption{debug}{\typeout{Loads of annoying messages}%
  \@debugtrue}
%    \end{macrocode}
% These two need to be declared explicitly so that |\ExecuteOptions|
% actually does something.
%    \begin{macrocode}
\DeclareOption{twocolumn}{\@rmtwocoltrue}
\DeclareOption{onecolumn}{\@rmtwocolfalse}
\DeclareOption{twoside}{\PassOptionsToClass{\CurrentOption}{article}}
\DeclareOption{boldkeywords}{\@boldkeywordstrue}
%    \end{macrocode}
% Bold, upper-case keywords used to be the default. Now they must be
% turned on explicitly with the \texttt{boldkeywords} option. 
%    \begin{macrocode}
\DeclareOption{boldkeywords}{\@boldkeywordstrue}
%    \end{macrocode}
% Pass any other option along to the \texttt{article} class.
%    \begin{macrocode}
\DeclareOption*{\PassOptionsToClass{\CurrentOption}{article}}
%    \end{macrocode}
% Default options are to be two-sided and based on the standard
% \LaTeX{} \texttt{article} class,
% unless we are processing the entire book, in which case we base
% ourselves on the standard \texttt{book} class. Default is to use
% manuscript option unless otherwise specified. 
%    \begin{macrocode}
\ExecuteOptions{twoside,manuscript} \ProcessOptions\relax
\if@book
\if@rmtwocol\LoadClass[twocolumn]{book}\else\LoadClass{book}\fi
\newcommand{\abstract}{\relax}
\newcommand{\@maketitle}{\relax}
\else
\if@rmtwocol\LoadClass[twocolumn]{article}\else\LoadClass{article}\fi
\fi
%    \end{macrocode}
% Attempt to use natbib for citations 02 Jan 2003. 
% \changes{3.05}{2003/01/02}{If the natbib option is used, then set
% things up so that natbib's version of the bibliography is identical
% to the standard one.}
% \changes{3.09}{2003/01/23}{Explicitly require a recent version of
% natbib, so that we get a decent error message if the package is too
% old.} 
% \changes{3.23}{2007/09/09}{Now allow Bibtex to work properly} 
% \changes{3.25}{2009/10/01}{Change \cs{newcommand} tp \cs{def} when setting \cs{bibfont}. This allows it to work with all versions of natbib.} 
% \changes{3.25}{2010/02/06}{Updated copyright year.} 
% \changes{3.25}{2010/02/06}{Updated email and web addresses to \texttt{crya.unam.mx}.} 
%    \begin{macrocode}
\if@natbib
\RequirePackage{natbib}[1999/02/25]
\bibpunct{(}{)}{;}{a}{}{,}
\def\bibfont{\small}
\renewcommand\bibsection{\section*{REFERENCES}}
\setlength\bibsep{0pt}
\setlength\bibhang{1.5em}
\bibliographystyle{rmaa}
\fi
%    \end{macrocode}
% We want the \texttt{graphicx} package for the extended syntax of
% |\includegraphics|. We want the \texttt{ifthen} package for using in
% these macros (although we don't much yet). We explicitly load
% \texttt{babel} since some \LaTeX{} distributions messed up the
% hyphenation if we didn't.  
%    \begin{macrocode}
\RequirePackage{graphicx}
\RequirePackage{ifthen}
\IfFileExists{babel.sty}{%
%   \let\@restoreindent=\@afterindentfalse
  \RequirePackage[english]{babel}
%   \@restoreindent
  %% reinstate the definition of \@roman so as not to break .idx files
  \@ifundefined{org@@roman}{\relax}{\renewcommand{\@roman}[1]{\org@@roman{##1}}}
  }{%
  \typeout{Package babel.sty not found - don't worry about this}}
%    \end{macrocode}
% Try to load the \texttt{dcolumn} and \texttt{booktabs} packages if
% they can be found. These can produce nicer-looking tables. 
%    \begin{macrocode}
\IfFileExists{dcolumn.sty}{%
  \RequirePackage{dcolumn}
  }{%
  \typeout{Package dcolumn.sty not found - you will not be able to
    align columns on decimal points in tables}}
\IfFileExists{booktabs.sty}{%
  \RequirePackage{booktabs}
  }{%
  \typeout{Package booktabs.sty not found - your tables will not look
    as pretty as they ought too.}
  }
\IfFileExists{makeidx.sty}{%
  \RequirePackage{makeidx}
}{%
  \typeout{Package makeidx.sty not found - Author Index may not work.}
  }
%% commented out for moment since requires more modern version of LaTeX
%% \RequirePackage[latin1]{inputenc}
%    \end{macrocode}
% \changes{4.0}{2019/11/13}{Make \texttt{hyperref} and
% \texttt{etoolbox} into hard requirements, since we now use them
% extensively.}  We use hyperref to get internal and external links.
% But I don't like the colored boxes so I turn those off.  We use
% etoolbox for patching and robustifying at the moment.  In the
% future, I want to port a lot of other stuff over to it.  For
% instance, replace the ifthen package and the TeX ifs too.
%
% We target the latest version (1.7) of PDF, rather than the default
% of 1.5.  This reduces the chance of getting warnings from included
% graphics.  This is not strictly related to hyperref, but we put it
% here anyway.
%    \begin{macrocode}
\RequirePackage{etoolbox}
\RequirePackage{hyperref}
\hypersetup{hidelinks}
\pdfminorversion=7
%    \end{macrocode}
% We patch the \cs{@combinedblfloats} following a suggestion by Heiko
% Oberdiek at \url{https://tex.stackexchange.com/a/249743/3843}, which
% avoids problems when hyperlinks are split across
% pages. \changes{4.0}{2019/11/12}{Fix for issue where breaking
% hyperlink across pages gave an error.}
%    \begin{macrocode}
\urlstyle{same}
\newcommand\email{\begingroup \urlstyle{rm}\Url}
\patchcmd\@combinedblfloats{\box\@outputbox}{\unvbox\@outputbox}{}{%
  \errmessage{\noexpand\@combinedblfloats could not be patched}%
}%
%    \end{macrocode}
% \subsection{Other setup stuff}\label{sec:setup}
% Some declarations: 
%    \begin{macrocode}
\newif\if@needaddresses
\@needaddressestrue
\newif\if@frontmatter
\@frontmatterfalse
\newif\if@inpreamble
\@inpreambletrue
\newif\if@firstauthor
\@firstauthortrue
\newif\if@firstaffil
\@firstaffiltrue
\newif\if@firstaltaffil
\@firstaltaffiltrue
\newif\if@firstkwd
\@firstkwdtrue
\newif\if@noexpandADS
\@noexpandADSfalse
\newif\if@sessiontitlepages
\@sessiontitlepagesfalse
%    \end{macrocode}
% \begin{macro}{\@rmaaEmpty}
% \begin{macro}{\@rmaaArg}
% These are used in various commands to check if arg was left blank
%    \begin{macrocode}
\newcommand{\@rmaaEmpty}{}
\newcommand{\@rmaaArg}{}
%    \end{macrocode}  
% \end{macro}
% \end{macro}
%
% \subsection{Compatibility with old \LaTeX{} versions}
% \label{sec:compat}
% This package is based on the 1998/12/01 release of \LaTeX{}.
% Older releases may not define all the commands that we need.
% Luckily, we can use |\providecommand| to define these if necessary. 
% The following commands are not defined in the 1994/06/01 release. 
%
% \begin{macro}{\MakeUpperCase}
% \begin{macro}{\textsuperscript}
% \begin{macro}{\textbackslash}
% User commands:
%    \begin{macrocode}
\providecommand{\MakeUpperCase}[1]{\uppercase{#1}}
\providecommand{\textsuperscript}[1]{\ensuremath{^\textrm{#1}}}
\providecommand{\textbackslash}{\ensuremath{\backslash}}
%    \end{macrocode}  
% \end{macro}
% \end{macro}
% \end{macro}
% \begin{macro}{\protected@xdef}
% \begin{macro}{\restore@protect}
% \begin{macro}{\@unexpandable@protect}
% \begin{macro}{\hb@xt@}
% \LaTeX{} internal commands:
%    \begin{macrocode}
\providecommand{\protected@xdef}{%
        \let\@@protect\protect
        \let\protect\@unexpandable@protect
        \afterassignment\restore@protect
        \xdef}
\providecommand{\restore@protect}{\let\protect\@@protect}
\providecommand{\@unexpandable@protect}{\noexpand\protect\noexpand}
\providecommand{\hb@xt@}{\hbox to}
%    \end{macrocode}
% \end{macro}
% \end{macro}
% \end{macro}
% \end{macro}
%
% \subsection{Page layout}\label{sec:page-layout}
%
% The dimensions of the text were defined by Elizabeth. Text width is
% 169\,mm with equal margins on the two sides: 
%    \begin{macrocode}
\setlength{\textwidth}{169mm}
\setlength{\@tempdima}{\paperwidth}
\addtolength{\@tempdima}{-\textwidth}
\setlength{\oddsidemargin}{0.5\@tempdima}
\addtolength{\oddsidemargin}{-1in}
\setlength{\evensidemargin}{\oddsidemargin}
%    \end{macrocode}
% Text height is 225\,mm, top margin is 14\,mm. 
%    \begin{macrocode}
\setlength{\textheight}{225mm} 
\setlength{\topmargin}{-1in} 
\addtolength{\topmargin}{14mm}
%    \end{macrocode}
% Space between columns in two-column styles, header, footer,
% paragraph indentation:
%    \begin{macrocode}
\setlength{\columnsep}{9.5mm}
\setlength{\headsep}{0.7cm}
\setlength{\footskip}{0.75cm}
\setlength{\parindent}{1.5em}
%    \end{macrocode}
% 
% \begin{macro}{\@set@scalable@length}
% Try out a new scheme for dealing with various spacing dimensions and
% allowing the user to easily modify them. The command
% \begin{quote}
% |\set@scalable@length|\marg{name}\marg{defaultlength}\marg{stretch}\marg{shrink}
% \end{quote}
% does a 
% |\setlength|\marg{\cs{name}}\marg{defaultlength plus stretch minus shrink}
% and also defines the user command
% \textbackslash\meta{name}|Scale|\marg{factor} that can 
% be used to rescale length by \meta{factor}. Additionally, now
% defines the package command
% \textbackslash\meta{name}|@set@scale|\marg{FACTOR} 
% that also rescales the length by \meta{FACTOR}, so that the length
% becomes \meta{factor} $\times$ \meta{FACTOR} $\times$
% \meta{defaultlength}.   
%    \begin{macrocode}
\newcommand\@set@scalable@length[4]{%
  %% arg #1: NAME of length without leading `\'
  %%
  %% arg #2: default VALUE  
  \expandafter\newlength\csname @#1@default\endcsname
  \expandafter\setlength\csname @#1@default\endcsname{#2}%
  \expandafter\newlength\csname @@#1@@default\endcsname
  \expandafter\setlength\csname @@#1@@default\endcsname{#2}%
  %% arg #3: default STRETCH  
  \expandafter\newlength\csname @#1@stretch\endcsname
  \expandafter\setlength\csname @#1@stretch\endcsname{#3}%
  \expandafter\newlength\csname @@#1@@stretch\endcsname
  \expandafter\setlength\csname @@#1@@stretch\endcsname{#3}%
  %% arg #4: default SHRINK  
  \expandafter\newlength\csname @#1@shrink\endcsname
  \expandafter\setlength\csname @#1@shrink\endcsname{#4}%
  \expandafter\newlength\csname @@#1@@shrink\endcsname
  \expandafter\setlength\csname @@#1@@shrink\endcsname{#4}%
  %%
  %% User-defined scale \@NAME@scale is unity by default
  \expandafter\gdef\csname @#1@scale\endcsname{1}%
  %% Package-defined scale \@NAME@SCALE is unity by default
  \expandafter\gdef\csname @#1@SCALE\endcsname{1}%
  %%
  %% internal command \@reset@NAME should expand to 
  %% \setlength\@@NAME@@default{\@NAME@SCALE\@NAME@default}
  %% \setlength\NAME{\@NAME@scale\@@NAME@@default}
  %% in other words: scale x SCALE x default 
  \expandafter\gdef\csname @reset@#1\endcsname{%
    \expandafter\setlength\csname @@#1@@default\endcsname{%
      \csname @#1@SCALE\endcsname \csname @#1@default\endcsname
      }
    \expandafter\setlength\csname @@#1@@stretch\endcsname{%
      \csname @#1@SCALE\endcsname \csname @#1@stretch\endcsname
      }
    \expandafter\setlength\csname @@#1@@shrink\endcsname{%
      \csname @#1@SCALE\endcsname \csname @#1@shrink\endcsname
      }
    \expandafter\setlength\csname #1\endcsname{%
      \csname @#1@scale\endcsname \csname @@#1@@default\endcsname
      \@plus
      \csname @#1@scale\endcsname \csname @@#1@@stretch\endcsname
      \@minus
      \csname @#1@scale\endcsname \csname @@#1@@shrink\endcsname
      }}
  %%
  %% User command \NAMEScale should expand to 
  %% {\renewcommand\@NAME@scale{#1}\@reset@NAME}
  \expandafter\gdef\csname #1Scale\endcsname##1{%
    \expandafter\gdef\csname @#1@scale\endcsname{##1}%
    \csname @reset@#1\endcsname}
  %%
  %% Package command \NAME@set@SCALE should expand to 
  %% {\renewcommand\@NAME@SCALE{#1}\@reset@NAME}
  \expandafter\gdef\csname #1@set@SCALE\endcsname##1{%
    \expandafter\gdef\csname @#1@SCALE\endcsname{##1}%
    \csname @reset@#1\endcsname}
  %%
  %% initialize at default values with \@reset@NAME
  \csname @reset@#1\endcsname
  }
%    \end{macrocode}
% \end{macro}
% \begin{macro}{\parskip}
% First use of this new infrastructure: set interparagraph space to be
% zero but with a small amount of stretchiness.
%    \begin{macrocode}
%   \ifthenelse{\boolean{@manuscript}}{%
%     \@set@scalable@length{parskip}{0.5\baselineskip}{0.3\p@}{0.1\p@}
%   }{
    \@set@scalable@length{parskip}{0\p@}{0.3\p@}{0\p@}
%   }
%    \end{macrocode}
% \end{macro}
% 
% \begin{macro}{\abovesectionskip}
% \changes{4.0}{2019/11/12}{Use negative values for skip so that first paragraph after section headings is not indented.  This is something that I have been trying to fix for years!  It turns out that the anser was in the documentation for the standard command \cs{@startsection.}}
% \begin{macro}{\belowsectionskip}
% \begin{macro}{\aboveparaskip}
% \begin{macro}{\belowparaskip}
% Define some utility rubber lengths for use in
% section~\ref{sec:section-commands}. 
%    \begin{macrocode}
\newlength{\abovesectionskip}
\newlength{\belowsectionskip}
\newlength{\aboveparaskip}
\newlength{\belowparaskip}
%    \end{macrocode}
% The default values for these lengths are those traditionally used
% for the onecolumn style. 
%    \begin{macrocode}
\@set@scalable@length{abovesectionskip}{-3.5ex}{-4ex}{-2ex}
\@set@scalable@length{belowsectionskip}{2.3ex}{0ex}{2ex}
\@set@scalable@length{aboveparaskip}{3.25ex}{2ex}{0.5ex}
\@set@scalable@length{belowparaskip}{-1em}{\z@}{\z@}
%    \end{macrocode}
% \end{macro}\end{macro}\end{macro}\end{macro}
% We also make scalable some of the standard float parameters. These
% are the same defaults as in the standard article classes, but
% slightly more stretchy. For the twocolumn style, they are currently
% halved, see below. Note that I have also lumped |\abovecaptionskip|
% in with these, since it plays a similar role. There is also a
% |\belowcaptionskip| but it is set to zero in the standard classes. 
%    \begin{macrocode}
\@set@scalable@length{floatsep}{12\p@}{4\p@}{4\p@}
\@set@scalable@length{textfloatsep}{20\p@}{4\p@}{8\p@}
\@set@scalable@length{intextsep}{12\p@}{4\p@}{4\p@}
\@set@scalable@length{dblfloatsep}{12\p@}{4\p@}{4\p@}
\@set@scalable@length{dbltextfloatsep}{20\p@}{4\p@}{8\p@}
\@set@scalable@length{abovecaptionskip}{10\p@}{\z@}{\z@}
%    \end{macrocode}
% \begin{macro}{\@before@title@skip}         
% \begin{macro}{\@before@author@skip}        
% \begin{macro}{\@between@author@skip}       
% \begin{macro}{\@before@abstract@skip}      
% \begin{macro}{\@after@abstracthead@skip}   
% \begin{macro}{\@between@abstract@skip}     
% \begin{macro}{\@before@kwd@skip}           
% \begin{macro}{\@before@dates@skip}           
% Spacing parameters used in |\@maketitle|. These are not changed by
% default except by the user, so they don't need to go into
% |\@package@default@lengths|.  
%    \begin{macrocode}
\newlength{\@before@title@skip}         
\newlength{\@before@author@skip}        
\newlength{\@between@author@skip}       
\newlength{\@before@abstract@skip}      
\newlength{\@after@abstracthead@skip}   
\newlength{\@between@abstract@skip}     
\newlength{\@before@kwd@skip}           
\newlength{\@before@dates@skip}           
\@set@scalable@length{@before@title@skip}{2em}{\z@}{\z@}
\@set@scalable@length{@before@author@skip}{1.5em}{\z@}{\z@}
\@set@scalable@length{@between@author@skip}{0.5em}{\z@}{\z@}
\@set@scalable@length{@before@abstract@skip}{3ex}{1ex}{0.5ex}
\@set@scalable@length{@after@abstracthead@skip}{1ex}{\z@}{\z@}
\@set@scalable@length{@between@abstract@skip}{2ex}{1ex}{0.5ex}
\@set@scalable@length{@before@kwd@skip}{1ex}{0.5ex}{0.2ex}
\@set@scalable@length{@before@dates@skip}{2ex}{0.5ex}{0.2ex}
%    \end{macrocode}
% \end{macro}\end{macro}\end{macro}\end{macro}
% \end{macro}\end{macro}\end{macro}\end{macro}
% 
% Spacing above and below displayed equations also wants to be done
% like this, but it is more complicated since the standard article
% classes make these change with the fontsize. Something to do later,
% maybe. 
% 
%    \begin{macrocode}
%    \end{macrocode}
% \begin{macro}{\@package@default@lengths}
% The value actually used for these lengths is controlled by the
% \meta{SCALE} parameter, which depends on the style option that is in
% effect. This is set using the |\@package@default@lengths| command,
% which must be called every time the style changes (e.g., changing
% from onecolumn to twocolumn). 
%    \begin{macrocode}
\newcommand{\@package@default@lengths}{%
  \ifthenelse{\boolean{@rmtwocol}}{%
    %% twocolumn mode
    \abovesectionskip@set@SCALE{0.5}
    \belowsectionskip@set@SCALE{0.5}
    \aboveparaskip@set@SCALE{0.5}
    \belowparaskip@set@SCALE{1}
    \floatsep@set@SCALE{0.5}
    \textfloatsep@set@SCALE{0.5}
    \intextsep@set@SCALE{0.5}
    \dblfloatsep@set@SCALE{0.5}
    \dbltextfloatsep@set@SCALE{0.5}
    }{%
    %% onecolumn mode
    \abovesectionskip@set@SCALE{1}
    \belowsectionskip@set@SCALE{1}
    \aboveparaskip@set@SCALE{1}
    \belowparaskip@set@SCALE{1}
    \floatsep@set@SCALE{1}
    \textfloatsep@set@SCALE{1}
    \intextsep@set@SCALE{1}
    \dblfloatsep@set@SCALE{1}
    \dbltextfloatsep@set@SCALE{1}
    }%
  }
\@package@default@lengths
%    \end{macrocode}
% \end{macro}
% \begin{macro}{\RescaleLengths}
% \begin{macro}{\RescaleSecLengths}
% \begin{macro}{\RescaleFloatLengths}
% \begin{macro}{\RescaleTitleLengths}
% User commands that will fiddle with all the scalable lengths. 
%    \begin{macrocode}
\newcommand{\RescaleLengths}[1]{%
  \RescaleSecLengths{#1}\RescaleFloatLengths{#1}\RescaleTitleLengths{#1}}
\newcommand{\RescaleSecLengths}[1]{%
%% sections
  \abovesectionskipScale{#1}
  \belowsectionskipScale{#1}
  \aboveparaskipScale{#1}
  }
\newcommand{\RescaleFloatLengths}[1]{%
%% floats
  \floatsepScale{#1}
  \textfloatsepScale{#1}
  \intextsepScale{#1}
  \dblfloatsepScale{#1}
  \dbltextfloatsepScale{#1}
  \abovecaptionskipScale{#1}
  }
\newcommand{\RescaleTitleLengths}[1]{%
%% maketitle
  \@before@title@skipScale{#1}           
  \@before@author@skipScale{#1}        
  \@between@author@skipScale{#1}       
  \@before@abstract@skipScale{#1}      
  \@after@abstracthead@skipScale{#1}   
  \@between@abstract@skipScale{#1}     
  \@before@kwd@skipScale{#1}           
  \@before@dates@skipScale{#1}           
  }
%    \end{macrocode}
% \end{macro}\end{macro}\end{macro}\end{macro}
% \changes{2.2beta14}{2002/03/17}{Finished implementation of new
% infrastructure for scalable lengths. Converted the above/below
% section skips to use it. Also, all the floatsep parameters. Also,
% all the vertical lengths that are used by |\@maketitle|. }

% \begin{macro}{\MSextramargin}
% For the \texttt{manuscript} option we have an extra 0.3\,inch
% vertical margin and 1\,inch horizontal margin for marginal notes (by
% hand or with |\marginpar|. 
%
% \changes{3.23}{2007/09/09}{Reduce baselinestretch in
%   \texttt{manuscript} style to 1.2. No longer use raggedright for
%   manuscript.}
%
% \changes{3.24}{2007/09/14}{Further reduction of baselinestretch in
%   \texttt{manuscript} style to 1.1 since we now have wider margins.}
%
% \changes{3.24}{2007/09/14}{Split \texttt{manuscript}'s extra margin
%   into separate horizontal and vertical parts.}
% 
%    \begin{macrocode}
\newlength{\MSextraHmargin}\newlength{\MSextraVmargin}
\setlength{\MSextraHmargin}{1in} 
\setlength{\MSextraVmargin}{0.3in} 
\if@manuscript
   \renewcommand{\baselinestretch}{1.1}
   \addtolength{\textwidth}{-2\MSextraHmargin}
   \addtolength{\oddsidemargin}{\MSextraHmargin}
   \setlength{\evensidemargin}{\oddsidemargin}
   \addtolength{\textheight}{-2\MSextraVmargin}
   \addtolength{\topmargin}{\MSextraVmargin}
%    \raggedright
\fi
%    \end{macrocode}
% \end{macro}
% Float parameters: don't allow bottom floats by default. 
%    \begin{macrocode}
\renewcommand{\fps@figure}{tp}
\renewcommand{\fps@table}{tp}
%    \end{macrocode}
%
% \begin{macro}{\rmaa@AuthorWidthFrac}
% \begin{macro}{\rmaa@AbstractWidthFrac}
% These parameters control what fraction of the total text width can
% be used by the list of authors and the abstract/resumen,
% respectively.  
%    \begin{macrocode}
\newcommand{\rmaa@AuthorWidthFrac}{1.0}
\newcommand{\rmaa@AbstractWidthFrac}{0.75}
%    \end{macrocode}
% Use a wide abstract to save space in the new style.
%    \begin{macrocode}
\if@newstyle\renewcommand{\rmaa@AbstractWidthFrac}{1.0}\fi
%    \end{macrocode}
% \end{macro}
% \end{macro}
%
% Some things to consider in future versions. Trying to avoid too many
% widows, etc. Apparently, the following will work best with
% \cs{flushbottom}. 
%    \begin{macrocode}
%%\widowpenalty=10000
%%\clubpenalty=10000
%    \end{macrocode}
% Also, should perhaps try getting rid of underfull/overfull box warnings. 
% \changes{3.10}{2003/01/24}{Set large values for \cs{hbadness},
% \cs{vbadness}, \cs{hfuzz}, \cs{vfuzz}. To reduce amount of stuff
% that gets written to the terminal every time you compile something.}
%    \begin{macrocode}
\hbadness=10000
\vbadness=10000
\setlength\hfuzz{\maxdimen}
\setlength\vfuzz{\maxdimen}
%    \end{macrocode}
% \subsection{Various text styles for different contexts}
% \label{sec:text-styles} 
% \begin{macro}{\rmaa@AffilStyle}
% \begin{macro}{\rmaa@PagesStyle}
% \begin{macro}{\rmaa@YearStyle}
% \begin{macro}{\rmaa@VolumeStyle}
% \begin{macro}{\rmaa@VolumenStyle}
% \begin{macro}{\rmaa@JournalNameStyle}
% \begin{macro}{\rmaa@ReceivedStyle}
% \begin{macro}{\rmaa@AbstractHeaderStyle}
% \begin{macro}{\rmaa@JournalNameStyle}
% Style commands for the different fields used by |\maketitle| in
% constructing the title page of the article. These have the form
% |\rmaa@XYZStyle{TEXT}|. I'm not sure these were really a good idea,
% but never mind \dots 
%    \begin{macrocode}
\newcommand{\rmaa@AffilStyle}[1]{{\normalsize\rmfamily#1}}
\newcommand{\rmaa@PagesStyle}[1]{{\normalsize\rmfamily#1}}
\newcommand{\rmaa@YearStyle}[1]{{\normalsize\rmfamily#1}}
\newcommand{\rmaa@VolumeStyle}[1]{{\normalsize\rmfamily\bfseries#1}}
\newcommand{\rmaa@VolumenStyle}[1]{{\Large\sffamily\bfseries#1}}
\newcommand{\rmaa@JournalNameStyle}[1]{{\normalsize\itshape#1}}
\newcommand{\rmaa@ReceivedStyle}[1]{{\normalsize\rmfamily\itshape#1}}
\newcommand{\rmaa@AbstractHeaderStyle}[1]{%
  {\normalsize\rmfamily\MakeUpperCase{#1}}}
\newcommand{\rmaa@TitleStyle}[1]{%
  {\large\bfseries\rmfamily\boldmath\MakeUpperCase{#1}\par}}
%    \end{macrocode}
% \end{macro}
% \end{macro}
% \end{macro}
% \end{macro}
% \end{macro}
% \end{macro}
% \end{macro}
% \end{macro}
% \end{macro}
%
% \begin{macro}{\rmaa@FrontSecStyle}
% Headings of ``front sections'' (see section~\ref{sec:front-sections})
%    \begin{macrocode}
\newcommand{\rmaa@FrontSecStyle}[1]{%
  {\normalsize\bfseries\rmfamily\MakeUpperCase{#1}}}
%    \end{macrocode}
% \end{macro}
%
% \begin{macro}{\rmaa@TOCTitleStyle}
% \begin{macro}{\rmaa@TOCPartStyle}
% \begin{macro}{\rmaa@IndexTitleStyle}
% \begin{macro}{\rmaaTOCAuthorStyle}
% \begin{macro}{\rmaaIndexAuthorStyle}
% Commands to set style of TOC and index entries.
% Rewritten to avoid unnecessary font commands in the .toc and .ind files
%    \begin{macrocode}
\newcommand{\rmaa@TOCTitleStyle}[1]{\MakeUpperCase{#1}}
\newcommand{\rmaa@TOCPartStyle}[1]{\rmaa@TOCTitleStyle{#1}}
\newcommand{\rmaa@IndexTitleStyle}[1]{#1}
\newcommand{\rmaaTOCAuthorStyle}[1]{%
  {\itshape#1}}
\newcommand{\rmaaIndexAuthorStyle}[1]{%
  \protect\mbox{\bfseries#1}}
%    \end{macrocode}
% \end{macro}
% \end{macro}
% \end{macro}
% \end{macro}
% \end{macro}
%
% \begin{macro}{\rmaa@KeywordStyle}
% \begin{macro}{\rmaa@KeywordHeaderStyle}
% The style of the keywords depends on whether we are using the
% \texttt{newstyle} option or not. If we are, then they are small. 
%    \begin{macrocode}
\ifthenelse{\boolean{@newstyle}}{
  \newcommand{\rmaa@KeywordStyle}[1]{%
    {\small%
      \if@boldkeywords 
      \bfseries\MakeUpperCase{\raggedright #1}
      \else
      \raggedright #1
      \fi
    }}
  \newcommand{\rmaa@KeywordHeaderStyle}[1]{%
    {\small\itshape#1}}}
{
  \newcommand{\rmaa@KeywordStyle}[1]{%
    {\normalsize%
      \if@boldkeywords\bfseries\fi\raggedright #1}}
  \newcommand{\rmaa@KeywordHeaderStyle}[1]{%
    {\normalsize\itshape#1}}}
%    \end{macrocode}
% \end{macro}
% \end{macro}
%
% \subsection{Page styles}\label{sec:page-styles}
% \begin{macro}{\ps@rmaa}
% The running headers for normal pages:
%    \begin{macrocode}
\newcommand{\ps@rmaa}{%
%    \end{macrocode}
% nothing in either foot; 
%    \begin{macrocode}
  \renewcommand\@oddfoot{\relax}%
  \renewcommand\@evenfoot{\relax}%
%    \end{macrocode}
% page number and short author list in the even heads; 
%    \begin{macrocode}
  \renewcommand\@evenhead{\rlap{\thepage}\hfil\@shortauthor\hfil}%
%    \end{macrocode}
% short title and page number in the odd heads.
%    \begin{macrocode}
  \renewcommand\@oddhead{\hfil\@shorttitle\hfil\llap{\thepage}}
  }
%    \end{macrocode}
% \begin{macro}{\ps@rmaafirstpage}
% The header of the first page of each
% article:
%    \begin{macrocode}
\newcommand{\ps@rmaafirstpage}{%
  \renewcommand\@oddfoot{\hfil\thepage}%
%    \end{macrocode}
% footer corrected (06 Jun 2000) to have page on outside edge;
%    \begin{macrocode}
  \renewcommand\@evenfoot{\thepage\hfil}%
%    \end{macrocode}
% Header has journal name, volume, and page numbers for a regular
% article, or similar for special cases. 
% \changes{4.0}{2019/11/03}{Change to three-line header with copyright info and DOI.}
%    \begin{macrocode}
\if@preprint
\renewcommand\@oddhead{%
  \if@proceedings 
  \textsc{To appear in ``\@conftitle{} (\rmaa@YearStyle{\@year})''}\hfil
  \textit{RMxAC}
  \else
  \textsc{To appear in} \@journalname\hfil
  \fi}\fi%
\if@guide
\renewcommand\@oddhead{\textsc{Documentation} (\rmaa@YearStyle{\@year})\hfil}
\fi%
\if@manuscript
\renewcommand\@oddhead{Manuscript for \@journalname{} (\rmaa@YearStyle{\@year})\hfil}
\fi%
\if@journal
\renewcommand\@oddhead{%
  \parbox{\textwidth}{
    \@journalname, \rmaa@VolumeStyle{\@volume},
    \rmaa@PagesStyle{\@artpages}\
    (\rmaa@YearStyle{\@year})\hfil\\[1pt]
    {\small\textcopyright{} \@year:
      Instituto de Astronom\'\i{}a,
      Universidad Nacional Aut\'onoma de M\'exico}\hfil\\[1pt]
    {\small\@DOI@url}\hfil
  }%
}\fi%
  \renewcommand\@evenhead{\@oddhead}}
%    \end{macrocode}
% By default, use the |\ps@rmaa| page style.
%    \begin{macrocode}
\pagestyle{rmaa}
%    \end{macrocode}
% \end{macro}
% \end{macro}
%
% \subsection{Preamble commands}\label{sec:preamble} 
% Everything you need for setting the title page, TOC and index
% entries, etc.
% \subsubsection{Commands for the title page}
% \label{sec:title-commands}
% \begin{macro}{\resumen}
% \begin{macro}{\@resumen}
% \begin{macro}{\abstract}
% \begin{macro}{\@abstract}
% User commands for setting the resumen and abstract. 
% These should be used in the document header. The internal commands
% |\@resumen| and |\@abstract| are used to hold the actual text
% pending output via the |\maketitle| command. 
%    \begin{macrocode}
\newif\if@setresumen
\newcommand{\resumen}[1]{%
%%  \renewcommand{\@resumen}{{\selectlanguage{spanish}#1}}%
  \renewcommand{\@resumen}{{\ignorespaces#1}}%
  \@setresumentrue}
\renewcommand{\abstract}[1]{%
  \renewcommand{\@abstract}{{\ignorespaces#1}}
  \renewcommand\ADS@abstract{#1}
%% Not a good idea, since then author macros are not expanded and
%% cause errors later. However, by not doing it we are vulnerable to
%% fragile commands.....
  \if@noexpandADS\@onelevel@sanitize{\ADS@abstract}\fi
  \@abstractpresenttrue
  }
\newcommand{\ADSnoexpand}{\@noexpandADStrue}
%    \end{macrocode}
% \begin{macro}{\boldabstract}
% This is the version for doing 1-page conference proceedings. Still
% save the text for ADS but then write it out immediately as a bold
% paragraph. 
%    \begin{macrocode}
\newcommand{\boldabstract}[1]{%
  \gdef\ADS@abstract{#1}%
  \renewcommand{\@abstract}{#1}%
  \if@noexpandADS\@onelevel@sanitize{\ADS@abstract}\fi%
  \noindent\textbf{\boldmath #1}
  }
%    \end{macrocode}
% \end{macro}
% Default value in each case is to print out a warning if not set
%    \begin{macrocode}
  \newcommand{\@resumen}{%
    \@setresumenfalse \ClassWarningNoLine{rmaa}{No Spanish abstract
      (resumen) specified - if you do not speak Spanish, this will be
      translated by the editors}%
    \textbf{Favor de proporcionar un resumen en espa\~nol. If you are
      unable to translate your abstract into Spanish, the editors will
      do it for you.} }
  \newcommand{\@abstract}{\ClassWarningNoLine{rmaa}{No abstract
      specified}}
%    \end{macrocode}
% \end{macro}
% \end{macro}
% \end{macro}
% \end{macro}
%
% \begin{macro}{\title}
% \begin{macro}{\@rawtitle}
% \begin{macro}{\@TOCtitle}
% \begin{macro}{\rmaaIndexTitle}
% Similarly, |\title| is the preamble command to set the title of the
% article. This is used in various places, typeset in different
% styles, hence the use of the internal |\@rawtitle|, |\@title|,
% |\rmaaIndexTitle|, and |\@TOCtitle| commands, each of which uses the
% relevant style from section~\ref{sec:text-styles}
%    \begin{macrocode}
\newcommand{\@rawtitle}{Raw Title}
\newcommand{\@TOCtitle}{\rmaa@TOCTitleStyle{Title for TOC}}
\newcommand{\rmaaIndexTitle}{\rmaa@IndexTitleStyle{Title for index}.}
\newcommand{\@TOCaftertitle}{}
\renewcommand{\title}[1]{%
%%  {\ADS@sanitize\xdef\@rawtitle{#1}}
  \renewcommand\@rawtitle{\noexpand#1}
  \let\ADS@title=\@rawtitle
  \renewcommand{\@title}{\rmaa@TitleStyle{#1}}%
  \renewcommand{\@TOCtitle}{\rmaa@TOCTitleStyle{#1}\@TOCaftertitle}%
  \renewcommand{\rmaaIndexTitle}{\rmaa@IndexTitleStyle{#1\@TOCaftertitle}.}%
  }
%    \end{macrocode}
% \end{macro}
% \end{macro}
% \end{macro}
% \end{macro}
% 
% All the weird and whacky ways of indicating the author
% affiliations. 
% \begin{macro}{\affil}
% First, the way that sets it directly underneath the authors name.
% Note that |\affil| trumps all other affiliations for the purposes of
% getting into the ADS stuff (we use |\gdef|).   
% \dots 
%    \begin{macrocode}
\newcommand{\affil}[1]{%
  \\[1.5ex]\rmaa@AffilStyle{#1}\\[1.5ex]%
  \if@firstaffil\gdef\ADS@affil{#1}\@firstaffilfalse\fi%
  }
\renewcommand{\and}{and\\[1.5ex]}
%    \end{macrocode}
% \begin{macro}{\altaffil}
% Alternative affiliations are footnoted. 
%    \begin{macrocode}
\newcommand{\altaffil}[1]{%
  \if@firstaltaffil\gdef\ADS@affil{#1}\@firstaltaffilfalse\fi%
  \thanks{#1}}
%    \end{macrocode}
% \begin{macro}{\altaffilmark}
% \begin{macro}{\altaffiltext}
% Alternative-alternative mechanism for affiliations---should not be
% mixed with |\altaffil| (unless you know what you are doing) or there
% will be problems with numbering.
%    \begin{macrocode}
\newcommand{\altaffilmark}[1]{\textsuperscript{#1}}
\newcommand{\altaffiltext}[2]{%
  \if@firstaltaffil\gdef\ADS@affil{#2}\@firstaltaffilfalse\fi
  \stepcounter{footnote}%
  \protected@xdef\@thanks{\@thanks
    \protect\footnotetext[#1]{#2}}} 
%    \end{macrocode}
% \end{macro}
% \end{macro}
% \end{macro}
% \end{macro}
%
% \begin{macro}{\keywords}
% \begin{macro}{\@keywords}
% \begin{macro}{\@keywordheader}
% \label{keywords}Sundry minor preamble commands. All
% self-explanatory, I hope. |\keywords| is the old user interface
% (deprecated). New articles should use |\addkeyword| once for each
% keyword. That way the ADS record comes out right too (comma rather
% than em-dash separator).
% \changes{2.2beta12}{2001/12/07}{\cs{addkeyword} New user interface
%   for keywords.} 
% \changes{2.2beta13}{2002/03/11}{\cs{addkeyword} Improved behaviour
%   with .ads file. Use
%   \cs{ADS@sanitize} to change definitions of `~' and \cs{ion}.}
% \changes{3.23}{2007/09/09}{\cs{addkeyword} now signals a compilation
% error if the author tries to add more than 6 keywords.}
% \changes{3.24}{2007/09/14}{kwdcounter is now reset in
% \cs{@begofarticlehook} to prevent error in book style} 
%    \begin{macrocode}
\newcommand{\@kwdsep}{ --- }
\newcommand{\ADS@kwdsep}{, }
\newcounter{kwdcounter}\setcounter{kwdcounter}{0}
\newcommand{\addkeyword}[1]{%
  \addtocounter{kwdcounter}{1}
  \ifthenelse{\thekwdcounter > 6}{
    \ClassError{rmaa}{Too many keywords. A maximum of 6 keywords is allowed.}{}
  }{}
  \if@firstkwd%
  \renewcommand\@keywords{\if@boldkeywords\uppercase{#1}\else #1\fi}
  {\ADS@sanitize\def\ion{\plainion}\gdef\ADS@keywords{#1}}
  \@firstkwdfalse
  \else%
  \edef\@keywords{\@keywords\@kwdsep \if@boldkeywords\uppercase{#1}\else #1\fi}%
  {\ADS@sanitize\def\ion{\plainion}\xdef\ADS@keywords{\ADS@keywords\ADS@kwdsep #1}}%
  \fi}
\newcommand{\ADS@sanitize}{%
  \def\mathrm{}
  \def\@{\ }
  \def~{\ }
  \def\langle{<}
  \def\rangle{>}
  \def\log{log }
  \def\sin{sin }
  \def\cos{cos }
  \def\tan{tan }
  \def\fd{^d.}
  \def\fh{^h.}
  \def\fm{^m.}
  \def\fs{^s.}
  \def\fdg{�.}
  \def\farcm{'.}
  \def\farcs{''.}
  \def\arcdeg{�}
  \def\,{ }
  }
\newcommand{\keywords}[1]{%
  \renewcommand{\ADS@keywords}{#1}
  \renewcommand{\@keywords}{\uppercase{#1}}
  }
\newcommand{\@keywordheader}{\rmaa@KeywordHeaderStyle{Key Words:}}
\newcommand{\@keywords}{%
  \ClassWarningNoLine{rmaa}{No keywords specified}%
  Keywords go here
  }
%    \end{macrocode}
% \begin{macro}{\shorttitle}
% \begin{macro}{\shortauthor}
% \begin{macro}{\@shorttitle}
% \begin{macro}{\@shortauthor}
% Commands for setting the running headers.
%    \begin{macrocode}
\newcommand{\shorttitle}[1]{%
  \renewcommand\@shorttitle{\MakeUpperCase{#1}}}
\newcommand{\shortauthor}[1]{%
  \renewcommand\@shortauthor{\MakeUpperCase{#1}}}
\newcommand{\@shorttitle}{SHORT TITLE GOES HERE}
\newcommand{\@shortauthor}{SHORT AUTHOR GOES HERE}
%    \end{macrocode}
% \end{macro}
% \end{macro}
% \end{macro}
% \end{macro}
% \end{macro}
% \end{macro}
% \end{macro}


% \subsubsection{The full addresses}
% \label{sec:full-addresses}
% \begin{macro}{\fulladdresses}
% \begin{macro}{\@fulladdresses}
% \begin{macro}{\@fulladdheight}
% \begin{macro}{\@fulladddepth}
% \begin{macro}{\@fulladdbox}
% Aaghh, the full addresses! A perennial source of problems. 
%    \begin{macrocode}
\newlength{\@fulladdheight}
\newlength{\@fulladddepth}
\newsavebox{\@fulladdbox}
%% \fulladdresses command modified (11 May 2000) so that we can put
%% \fulladdresses{} in the article to suppress the output of the
%% addresses
\newcommand{\@FinalFigure}{\relax}
\newcommand{\FinalFigure}[1]{\renewcommand{\@FinalFigure}{#1}}
\newcommand{\fulladdresses}[1]{%
  \renewcommand{\@rmaaArg}{#1}
  \ifx\@rmaaArg\@rmaaEmpty
     \renewcommand{\@fulladdresses}{}
  \else
     \renewcommand{\@fulladdresses}{
       \normalsize%
       \ifx\@FinalFigure\relax\else\vfill\@FinalFigure\vfill\fi
%%       \vspace*{1ex}
       %% do addresses as a list (copied from thebibliography environment) 
       \list{\null}{\leftmargin 1.5em\topsep\z@\labelwidth\z@\labelsep\z@\itemindent -1.5em
         \itemsep\z@\parsep\z@\usecounter{enumi}}
       \sloppy\clubpenalty4000\widowpenalty4000
       \sfcode`\.=1000\relax
       #1
       \endlist
       }
  \fi
  }
\newcommand{\@fulladdresses}{Full addresses go here}
%    \end{macrocode}
% \begin{macro}{\suppressfulladdresses}
% New user-level command for suppressing full addresses should you so
% desire. Currently used in the single-page-poster style. 
%    \begin{macrocode}
\newcommand{\suppressfulladdresses}{\@suppressfulladdressestrue\@needaddressesfalse}
%    \end{macrocode}
% \end{macro}
% \end{macro}
% \end{macro}
% \end{macro}
% \end{macro}
% \end{macro}

% \begin{macro}{\printdates}
% \begin{macro}{\noprintdates}
% Command to activate the printing out of the received/accepted
% dates. This can also be done using the \texttt{printdates} option,
% which is the default for \texttt{journal}, but not for \texttt{proceedings}.
%    \begin{macrocode}
\newcommand{\printdates}{\@printdatestrue}
\newcommand{\noprintdates}{\@printdatesfalse}
%    \end{macrocode}
% \end{macro}
% \end{macro}
% \begin{macro}{\@outputaddressesifnecessary}
% Command to output the addresses if we haven't already done so
% This is called by: 
% \begin{enumerate}
% \item The end of |thebibliography| environment
% \item The |\outputfulladdresses| command
% \end{enumerate}
%    \begin{macrocode}
\newcommand{\@outputaddressesifnecessary}{%
  \ClassInfo{rmaa}{Calling @outputaddressesifnecessary}
  \ifthenelse{\boolean{@needaddresses}}{%
    %% If in one column style, or in the first column of a two column
    %% style then we can try to put it on this page 
    \ifthenelse{\boolean{@firstcolumn} \or \not\boolean{@rmtwocol}}{%
      \@outputaddresses}{%
      %% otherwise give warning it will be on a separate page
      \ClassWarning{rmaa}{Full addresses will go on a separate
        page (please see User Guide for how to remedy this)}}%
    %% If we haven't managed yet, put it on its own page
    \ifthenelse{\boolean{@needaddresses}}{%
      \if@debug{\ClassWarning{rmaa}{In this silly bit!}}\fi
      \clearpage\vfill\@outputaddresses}{}%
    }{
    %% If we don't need them, we still might need to do the eoa hook 
    \ifthenelse{\boolean{@suppressfulladdresses}}{%
      \if@debug{\ClassWarning{rmaa}{XXX}}\fi
%%      \if@book\else\@endofarticlehook\null\fi
      }{}
    }%
  }  
%    \end{macrocode}
% \end{macro}

% \begin{macro}{\outputfulladdresses}
% A user interface to the above command, in case there is no biblio
% (11 May 2000) 
%    \begin{macrocode}
\newcommand{\outputfulladdresses}{\@outputaddressesifnecessary}
%    \end{macrocode}
% \end{macro}

% \begin{macro}{\adjustfinalcols}
% This command starts a new column and should be 
% used to ensure that the columns balance on the last page. It also
% outputs the full 
% author addresses to ensure that this occurs in the left hand
% column. 
% The |\par| command is so that the current line will be taken account
% of when working out |\pagetotal|
%    \begin{macrocode}
\newcommand{\adjustfinalcols}{\if@twocolumn\goodbreak\@outputaddresses\newpage\noindent\fi}
%    \end{macrocode}
% \end{macro}

% \begin{macro}{\@outputaddresses}
% Command to do the work of outputting the addresses. 
% This can be called by 
% \begin{enumerate}
% \item an explicit |\adjustfinalcols|
% \item via |\@outputaddressesifnecessary|, q.v.
% \end{enumerate}
%
%    \begin{macrocode}
\newcommand{\@outputaddresses}{%
  % Allow option of suppressing the full addresses, in which case
  % we just run the hook if we are not in a book style
  \if@suppressfulladdresses
%%     \if@book\else\@endofarticlehook\null\fi
  \else
  %% save the full addresses in a box for output later
  \sbox{\@fulladdbox}{\parbox{\textwidth}{\@fulladdresses}}
  % calculate height of address box  
  \settoheight{\@fulladdheight}{\usebox{\@fulladdbox}}
  \settodepth{\@fulladddepth}{\usebox{\@fulladdbox}}
  \addtolength{\@fulladdheight}{\@fulladddepth}
  \if@debug
  { \typeout{Address total height: \the\@fulladdheight}
    \typeout{Text height: \the\pagetotal}
    \typeout{Goal height: \the\pagegoal}
    }
  \fi
  % reuse \@fulladdheight to mean the total height of text+addresses
  \addtolength{\@fulladdheight}{\pagetotal}
%%  \addtolength{\@fulladdheight}{-\baselineskip}
  \if@debug\typeout{Adjusted text + addresses height: \the\@fulladdheight}\fi
  % check it will fit on current page
  \ifdim\@fulladdheight<\pagegoal
     %% must shift it to the left
     %% to compensate for the hanging indentation in the reference list. 
     %% NOTE: this is not foolproof!
     %% NB. This will give overfull hbox warnings, which is a bit ugly.
     \vfill\noindent\if@biblist\hspace*{-\leftmargin}\fi
     \usebox{\@fulladdbox}
     % Set this to show we've done it
     \@needaddressesfalse
     % and assume we are at the end of the article
     % so that we do the EOA hook if not in the book style
%%     \if@book\else\@endofarticlehook\fi
  \else
     % if not, just do nothing. With luck, this function will get 
     % called again on the next page
     \if@debug{\typeout{Holding over full addresses to next page}}\fi
  \fi\fi
  }
%    \end{macrocode}
% \end{macro}

% \subsubsection{Commands for editors}
% \label{sec:editor-commands}
% Some user/editor commands for doing the TOC and index in the book
% option. 
% They are defined here rather in the book-specific part since they
% need to be always available in order to avoid breaking documents. 
% \begin{macro}{\listofauthors}
% \begin{macro}{\@rawauthors}
% \begin{macro}{\rmaaTOCauthors}
% The command |\listofauthors| formats an author list for use in the
% TOC. It should contain the same names as in the |\author| command,
% but without the affiliation garbage and using |\&| instead of
% |and|. 
%    \begin{macrocode}
\newcommand{\@rawauthors}{Raw Authors}
\newcommand{\rmaaTOCauthors}{\noexpand\rmaaTOCAuthorStyle{Author list for TOC}}
\newcommand{\listofauthors}[1]{%
  \renewcommand{\@rawauthors}{#1}
  \renewcommand\rmaaTOCauthors{\noexpand\rmaaTOCAuthorStyle{#1}}}
%    \end{macrocode}
% \begin{macro}{\indexauthor}\label{indexauthor}
% The command |\indexauthor| produces an index entry for each
% author. 
% This command must come \emph{after} the
% |\title| and the |\listofauthors| commands. Also, there should be one 
% |\indexauthor| command for \emph{each} author. We also add the
% author onto the end of the |\ADS@authors| macro and construct the
% |\ADS@authorinitial| by a most cunning method (see below), but only
% for the first author. 
% 
% No, No, NO! This was all wrong. Backed out completely: 
% \begin{quote}
% Completely rewritten (19 Dec 2002) since we now use the index class, so 
% as to be able to alphabetize accented characters. A side effect of this is
% that |\index| no longer expands out the argument after the |@|, so we 
% instead add all the |\index| commands to a list (|\indexauthor@deferlist|), 
% which is executed in the |\@begofarticlehook|. 
% 
% Note the stuff with |\toks0| is copied from Exercise 20.15 of the \TeX{}book---it is
% to make sure the latin1-accented name gets passed into the idx
% file. 
% \end{quote}
% We now go back to the old way of doing things, using an optional
% argument to the \cs{indexauthor} command that has the name without
% accents. 
% \changes{3.08}{2003/01/16}{Remerged changes from 3.06 that got
% dropped from 3.07 after a synchronization disaster. Basically, the
% accentless optional argument to \cs{indexauthor}. Must add the macro
% directories to my unison set-up to stop this sort of thing from happening.}
%    \begin{macrocode}
\newcommand\@isortkey{}
\newcommand\@emptykey{empty}
\newcommand\indexauthor[2][empty]{
  \renewcommand\@isortkey{#1}
  \ifx\@isortkey\@emptykey \renewcommand\@isortkey{#2}\fi
  \index{\@isortkey @{{\rmaaIndexAuthorStyle{#2}}\hskip1em plus1em minus0.4em
      \rmaaIndexTitle\hskip1em plus2em minus0.4em\penalty-100
      \rmaaTOCauthors}}
  \indexauthorADS{#2}
}
%%\def\indexauthor@deferlist{\relax}
%%\long\def\indexauthor#1{
%%  \def\@@name{#1} \toks0=\expandafter{\@@name}
%%  \toks1=\expandafter{\indexauthor@deferlist}
%%  \protected@xdef\indexauthor@deferlist{%
%%    \the\toks1
%%    \noexpand\index{\the\toks0@{\noexpand\rmaaIndexAuthorStyle{\the\toks0}\noexpand\hskip1em plus1em minus0.4em 
%%        \rmaaIndexTitle\noexpand\hskip1em plus2em minus0.4em\noexpand\penalty-100
%%        \rmaaTOCauthors}}
%%  }
%%  \indexauthorADS{#1}
%%  }
%%\xdef\indexauthor@deferlist{}
\newcommand\indexauthorADS[1]{%
  \if@firstauthor\protected@xdef\ADS@authorinitial{\firstletter{#1}}\fi
  \protected@xdef\ADS@authors{\ADS@authors\if@firstauthor\else ; \fi
    #1}
  \@firstauthorfalse
}
%    \end{macrocode}
% Note that the first |#1| sets the key for the index, making sure it
% actually starts with  
% first letter of the surname. Meanwhile, the |@{...}| expression,
% which starts with some  
% control sequences is what is written to the |.ind| file. 
% See the makeindex man page for details. 
% \end{macro}
% \end{macro}
% \end{macro}
% \end{macro}

% \begin{macro}{\KillAllAccents}
% \begin{macro}{\firstletter}
% These are utility functions used in the above. |lett| is a counter
% indicating the current letter in the string. |\firstletter| just
% prints the first letter of its argument.
%    \begin{macrocode}
\newcommand{\KillAllAccents}{\def\'{}\def\`{}\def\~{}\def\.{}\def\i{i}}
\newcommand\firstletter[1]{\@car #1\@nil}
%    \end{macrocode}
% \end{macro}
% \end{macro}
% \begin{macro}{\SetVolume}
% \begin{macro}{\SetMSnumber}
% \begin{macro}{\SetFirstPage}
% \begin{macro}{\SetYear}
% \begin{macro}{\ReceivedDate}
% \begin{macro}{\AcceptedDate}
% These are for the editors to set the volume number, etc. 
%    \begin{macrocode}
\newcommand{\SetVolume}[1]{%
  \renewcommand{\@volume}{#1}
  \edef\ADS@volume{\PadToFourWithDots{#1}}
  \renewcommand{\@volumen}{\rmaa@VolumenStyle{volumen #1}}%
  }
\newcommand{\SetMSnumber}[1]{\renewcommand{\@MSnum}{#1}}
%    \end{macrocode}
% Note that there is no real reason to use |\SetFirstPage| since all
% the page numbering is done automatically in the book style. 
%    \begin{macrocode}
\newcommand{\SetFirstPage}[1]{\setcounter{page}{#1}}
\newcommand{\SetYear}[1]{%
  \renewcommand{\@year}{#1}
  \renewcommand{\ADS@year}{#1}
  \renewcommand{\@anyo}{\rmaa@VolumenStyle{#1}}
  }
\newcommand{\ReceivedDate}[1]{\renewcommand{\@received}{#1}}
\newcommand{\AcceptedDate}[1]{\renewcommand{\@accepted}{#1}}
%    \end{macrocode}
%
% \begin{macro}{\DOI}
% User/editor command to enter the DOI, which will be typeset as a link.
% \changes{4.0}{2019/11/02}{New user command \cs{DOI} to typeset
% DOI in the header of first page and add it to ADS info.}
%    \begin{macrocode}
\newcommand\DOI[1]{%
  \def\@DOI@url{\url{https://doi.org/10.22201/ia.01851101p.#1}}%
  \def\ADS@doi{https://doi.org/10.22201/ia.01851101p.#1}%
}%
%    \end{macrocode}
% \end{macro}
%
% \begin{macro}{\ORCIDiD}
% \changes{4.0}{2019/11/12}{Add mechanism for giving ORCID iD links}
% \begin{macro}{\ORCIDiDfull}
% \begin{macro}{\@orcid@icon}
% \begin{macro}{\@orcid@iconbig}
% Commands to add icon links to ORCID iDs of authors. For example
% \begin{quote}
% |\ORCIDiD|\marg{0000-0000-0000-0000}
% \end{quote}
% to be used after each author name in the |\author|
% command and in the |\fulladdresses|.  Alternatively,
% |\OORCIDiDfull| will typeset the icon plus the full text of the
% link, but we don't use that at present.  The icon image is stored in |ORCIDiD_iconvector.pdf|, which needs to be on the input path.  In principle, we could make EPS versions if necessary.
%    \begin{macrocode}
\newcommand\@orcid@icon{\includegraphics[height=1.8ex]{ORCIDiD_iconvector}}
\newcommand\@orcid@iconbig{\includegraphics[height=2ex]{ORCIDiD_iconvector}}
\newcommand\ORCIDiD[1]{\href{https://orcid.org/#1}{\,\@orcid@icon}}
\newcommand\ORCIDiDfull[1]{%
  \href{https://orcid.org/#1}{\@orcid@iconbig}\,\url{https://orcid.org/#1}}
%    \end{macrocode}
% \end{macro}
% \end{macro}
% \end{macro}
% \end{macro}
%
% \begin{macro}{\@volume}
% \begin{macro}{\@volumen}
% \begin{macro}{\@artpages}
% \begin{macro}{\@year}
% \begin{macro}{\@anyo}
% \begin{macro}{\@received}
% \begin{macro}{\@accepted}
% \begin{macro}{\@MSnum}
% Default values for the containers used by above commands, plus a
% few more useful ones. 
%    \begin{macrocode}
\newcommand{\@volume}{00}
\newcommand{\@volumen}{\rmaa@VolumenStyle{00}}
\newcommand{\@artpages}{%
  \oldpageref{\articlename:RMAAFirstPage}--\oldpageref{\articlename:RMAALastPage}}
\newcommand{\@year}{\the\year}
\newcommand{\@anyo}{\rmaa@VolumenStyle{2007}}
\newcommand{\@journalname}{\rmaa@JournalNameStyle{%
    \if@proceedings\RMAASClong\else\RMAAlong\fi}}
\newcommand{\@received}{}
\newcommand{\@accepted}{}
\newcommand{\@MSnum}{0000}
\newcommand{\@DOI@prefix}{https://doi.org/10.22201/ia.01851101p}
\def\@DOI@url{\url{https://doi.org/10.22201/ia.01851101p.2099.00.00.00}}
%    \end{macrocode}
% \end{macro}
% \end{macro}
% \end{macro}
% \end{macro}
% \end{macro}
% \end{macro}
% \end{macro}
% \end{macro}
% \end{macro}
% \end{macro}
% \end{macro}
% \end{macro}
% \end{macro}
% \end{macro}
% 
% The following commands are most useful in the \texttt{book}
% style. However, the |\SetConfTitle| command can now be used in
% individual articles to set the header for the \texttt{preprint}
% style. 
%    \begin{macrocode}
\newcommand{\@reunionplace}{Lugar de la reuni\'on}
\newcommand{\@reuniondate}{Fecha de la reuni\'on}
\newcommand{\@conftitle}{Conference Title}
\newcommand{\SetReunion}[2]{\renewcommand{\@reunionplace}{#1}\renewcommand{\@reuniondate}{#2}}
\newcommand{\SetConfTitle}[1]{\renewcommand{\@conftitle}{#1}}
\newcommand{\@editors}{Juan Ram\'on de la Fuente and ``El Mosh''}
\newcommand{\SetEditors}[1]{\renewcommand{\@editors}{#1}}
\newcommand{\@editores}{Juan Ram\'on de la Fuente y ``El Mosh''}
\newcommand{\SetEditores}[1]{\renewcommand{\@editores}{#1}}
\newcommand{\@mes}{mes de la publicaci\'on}
\newcommand{\SetMes}[1]{\renewcommand{\@mes}{\rmaa@VolumenStyle{#1}}}
%    \end{macrocode}

% \subsection{Notes for Author Proofs and Communication between Editors}
% \label{sec:notes}
% We introduce a command \cs{Note} that should be used to describe
% any change made to the manuscript. These are a bit tricky to get
% right. In general we use \cs{marginpar} but this only works in outer
% paragraph mode, so we just bung the box where we are if we are in a
% figure caption or something. If we are in the preamble, then we just
% use a \marginpar as part of the \maketitle command.
%    \begin{macrocode}
\newif\if@arenotes
\@arenotesfalse
\newcounter{notecounter}
\setcounter{notecounter}{0}
\newcommand\NotesList{}
\newcommand\Note[1]{%
  \if@proof{%
    \@arenotestrue%
    \stepcounter{notecounter}%
    \if@inpreamble\else{%
      \NoteMark{\arabic{notecounter}}%
    }\fi
    \protected@xdef\NotesList{\NotesList 
      \protect\item #1}%
  }\fi%
}
\newcommand\Changed[2]{Changed ``#1'' to ``#2''.}
\newcommand\NoteMark[1]{%
  \ifinner
  \raisebox{1ex}[0pt][0pt]{\makebox[0pt][l]{\framebox{\footnotesize\bfseries
        N#1}}}%
  \else
  \marginpar{\framebox{\footnotesize\bfseries N#1}}%
  \fi
}
\newcommand\RefNote[1]{\if@proof\NoteMark{\ref{#1}}\relax\fi}
\newcommand\NoteLabel{\protect\label}
\newcommand\OutputNotes{%
  \bgroup
  \onecolumn
  \section*{EDITORIAL NOTES}
  \begin{enumerate}
    \NotesList
  \end{enumerate}
  \egroup
}
\DeclareRobustCommand\cs[1]{\texttt{\char`\\#1}}

%    \end{macrocode}
% \subsection{Bibliographic Records for ADS}
% \label{sec:ADS}
% \changes{2.2beta11}{2001/11/27}{Just about finished all the ADS
% stuff now. Except for the references, which I have given up on for
% the time being.}
% \begin{macro}{\ADS@BibCode}
% \begin{macro}{\ADS@year}
% \begin{macro}{\ADS@journal}
% \begin{macro}{\ADS@volume}
% \begin{macro}{\ADS@firstpage}
% \begin{macro}{\ADS@authorinitial}
% We use the |\ADS@BibCode| command to hold the bibliographic code for
% ADS. It is 
% used in the |%R| 
% field of the bibliographic record for each article. 
% The format is 
% \begin{quote}
% \texttt{YYYYJJJJJVVVVMPPPPA}  
% \end{quote}
% where \texttt{YYYY} is the year, \texttt{JJJJJ} is the journal code,
% which is \texttt{RMxAA} or \texttt{RMxAC}, \texttt{VVVV} is the
% volume number (right-justified), \texttt{M} is not used by us,
% \texttt{PPPP} is the first page number (right-justified), and
% \texttt{A} is the initial of the first author's surname. The volume
% and page fields are padded with periods (`\texttt{.}'), which causes
% us some fun (this is done with the macro |\PadToFourWithDots| described
% below) . 
% 
% All these individual bits are held in the commands |\ADS@year|,
% |\ADS@journal|, |\ADS@volume|, |\ADS@firstpage|,
% |\ADS@authorinitial|. 
%    \begin{macrocode}
\newcommand{\SetADSBibCode}[1]{\renewcommand{\ADS@BibCode}{#1}}
\newcommand{\ADS@BibCode}
{\ADS@year\ADS@journal\ADS@volume.\PadToFourWithDots{\ADS@firstpage}\ADS@authorinitial}
\newcommand{\ADS@year}{2001} 
\newcommand{\ADS@journal}{\if@proceedings RMxAC\else RMxAA\fi}
\newcommand{\ADS@volume}{....} 
%    \end{macrocode}
% These two are just set to default values here, |\ADS@firstpage| is set
% in the |document| environment for included articles, in the
% |frontsection| environment, |\tableofcontents| command,  |index| and
% |abstracts| environments. Also by |\AtBeginDocument| for stand-alone
% articles. 
% 
% God knows how we are going to do |\ADS@authorinitial|!
%    \begin{macrocode}
\newcommand{\ADS@firstpage}{....}
\newcommand{\ADS@authorinitial}{A}
%    \end{macrocode}
% \end{macro}
% \end{macro}
% \end{macro}
% \end{macro}
% \end{macro}
% \end{macro}
% We also have a whole slew of other fields that can be set. These are
% all listed in Table~\ref{tab:ADS}. Most of them still have to be
% defined. 
% \begin{table}[!ht]\raggedleft
% \caption{ADS Bibliographic Records}\label{tab:ADS}
% \begin{tabular}{lp{8em}lp{8em}}\hline
% Code & Description & Macro Name & Comment \\ \hline
% |%R| & Bibliographic Code & |\ADS@BibCode| & required \\
% |%A| & Author list & |\ADS@authors| & required, generate from |\indexauthor| \\
% |%F| & Author affiliation & |\ADS@affil| & format: AA(\dots)
% AB(\dots) \dots \\
% |%J| & Journal Name & |\ADS@journalname| &  \\
% |%D| & Publication Date & |\ADS@fulldate| &  \\
% |%L| & Last Page & |\ADS@lastpage| & generate from label \\
% |%T| & Title & |\ADS@title| & required, generate from |\title| \\
% |%C| & Abstract Copyright & |\ADS@copyright| & set to IA-UNAM \\
% |%U| & URL for electronic document & |\ADS@url| & set to PDF file \\
% |%d| & DOI for article & |\ADS@doi| & set by |\DOI| \\
% |%K| & Keywords & |\ADS@keywords| & set equal to |\@keywords| \\
% |%A| & Abstract Text & |\ADS@abstract| & set equal to |\@abstract| \\
% |%R| & References & |\ADS@references| & generate from the
% bibliography  \\ \hline
% \end{tabular}
% \end{table}
% A lot of these macros are set by other commands, as mentioned in the
% Table, but some need to be set  explicitly here. 
% 
% \begin{macro}{\ADS@doi}
% \changes{4.0}{2019/11/08}{Save DOI for ADS}
% The user/editor command \cs{DOI} is used to set this. 
%    \begin{macrocode}
\newcommand{\ADS@doi}{https://doi.org/10.22201/ia.01851101p.0000.00.00.00}
%    \end{macrocode}
% \end{macro}
% 
% \begin{macro}{\ADS@journalname}
% This
% does not contain any accents at the moment so as not to cause
% problems---this maybe should be changed though. 
%    \begin{macrocode}
\newcommand{\ADS@journalname}{%
  \if@proceedings\@conftitle\space (Eds. \@editors) \fi%
  Revista Mexicana de Astronom\'\i{}a y Astrof\'\i{}sica%
  \if@proceedings\space (Serie de Conferencias)\fi%
  \space Vol. \@volume, pp. \ADS@firstpage--\ADS@lastpage\space (\@year)%
  \space (http://www.astroscu.unam.mx/\~{}rmaa/)
  }
%    \end{macrocode}
% \end{macro}
% \begin{macro}{\ADS@keywords}
% \begin{macro}{\ADS@abstract}
% \begin{macro}{\ADS@title}
% |\ADS@abstract| is straightforward, since we already have the text
% saved as |\@abstract|. We could do the same with |\ADS@keywords| but
% to avoid unwanted formatting we explicitly set it in |\keywords|
% (p.~\pageref{keywords}). |\ADS@title| is just set to |\@rawtitle|,
% which we were already using for another purpose. 
%    \begin{macrocode}
\let\ADS@title=\@rawtitle
\newcommand{\ADS@keywords}{No keywords specified}
\newcommand{\ADS@abstract}{No abstract specified}
%    \end{macrocode}
% \end{macro}\end{macro}\end{macro}
% \begin{macro}{\ADS@authors}
% |\ADS@authors| wants to be in ``Other, A. N.'' format, but with
% semi-colons between authors. The most natural way to do this seems to
% be with |\indexauthor| (p.~\pageref{indexauthor}). We append each
% author to this macro, which means we must explicitly clear it for
% each new article (done in |\article|). .   
%    \begin{macrocode}
\newcommand{\ADS@authors}{}
%    \end{macrocode}
% \end{macro}
% \begin{macro}{\ADS@fulldate}
% \begin{macro}{\ADS@lastpage}
% \begin{macro}{\ADS@copyright}
% \begin{macro}{\ADS@website}
% \begin{macro}{\ADS@url}
% \begin{macro}{\SetWebsite}
% \begin{macro}{\SetMonthNumeric}
% Other sundry data, together with user commands for bits that can be
% set by the Editors. The |\ADS@url| is constructed from the website
% (currently |http://www.astrosmo.unam.mx/rmaa/|,
% the journal descriptor (e.g., |RMxAC..10/|), followed by the
% path to the PDF file (e.g., |PDF/RMxAC..10_HEW1.pdf|). 
% 
%    \begin{macrocode}
\newcommand{\ADS@website}{http://www.astrosmo.unam.mx/rmaa/}
\newcommand{\ADS@url}
{\ADS@website\ADS@journal\ADS@volume\ADS@number/PDF/\ADS@journal\ADS@volume\ADS@number\_\articlename.pdf}
\newcommand{\ADS@number}{}
\newcommand{\ADS@fulldate}{00/\ADS@year} 
\newcommand{\ADS@lastpage}{00} 
\newcommand{\ADS@copyright}{(C) \ADS@year; Instituto de Astronom\'\i{}a,
  Universidad Nacional Aut\'onoma de M\'exico} 
\newcommand{\SetNumber}[1]{\renewcommand{\ADS@number}{-#1}}
\newcommand{\SetWebsite}[1]{\renewcommand{\ADS@website}{#1}}
\newcommand{\SetMonthNumeric}[1]{\renewcommand{\ADS@fulldate}{#1/\ADS@year}}
%    \end{macrocode}
% \end{macro}\end{macro}\end{macro}\end{macro}\end{macro}\end{macro}\end{macro}
% \begin{macro}{\ADS@affil}
% \begin{macro}{\ADS@references}
% These two are somewhat more problematical. Affiliation currently
% only works for the first author (see \dots) pending redesign of the
% user interface. References are now done by wrapping the
% |thebibliography| environment to open another input stream on the
% same file and search through for the |\bibitems|. 
%    \begin{macrocode}
\newcommand{\ADS@references}{}
\newcommand{\ADS@affil}{}
%    \end{macrocode}
% \end{macro}\end{macro}
% 
% Define two macros that expand to the strings
% |\begin{thebibliography}| and |\end{thebibliography}|. Note that the
% idea for this and much else is inspired by the code for the docstrip
% package. Also define |\StartBibX| in case authors use the more
% standard |\begin{thebibliography}{}| syntax. 
%    \begin{macrocode}
\edef\EndBib{\expandafter\@gobble\string\\end\string{thebibliography\string}}
\edef\StartBib{\expandafter\@gobble\string\\begin\string{thebibliography\string}}
\edef\StartBibX{\expandafter\@gobble\string\\begin\string{thebibliography\string}\string{\string}}
%    \end{macrocode}
% Define two new logical variables. |\ifAnotherline| is true if we
% want to read another line from the input stream. It is set to false
% when we encounter the |\end{thebibliography}|. |\ifInBib| is set to
% true when we frind the |\begin{thebibliography}| and tells us to
% start saving the following lines. 
%    \begin{macrocode}
\newif\ifAnotherLine
\newif\ifInBib
%    \end{macrocode}
% \begin{macro}{makeOther}
% We need a macro to change the catcodes of lots of things to other. 
%    \begin{macrocode}
\def\makeOther#1{\catcode`#1=12\relax}
%    \end{macrocode}
% \end{macro}
% \begin{macro}{\RereadBiblio}
% Define a macro to do all the fiddly stuff. This is used at the
% beginning of |thebibliography| environment (Sec~\ref{sec:biblio}).
%    \begin{macrocode}
\newread\tf@biblio
\newcommand{\RereadBiblio}{%
%    \end{macrocode}
% Open the current file again for reading
%    \begin{macrocode}
  \begingroup
  \immediate\openin\tf@biblio=\articlename
%    \end{macrocode}
% Everything in a group to be safe. Initialize logicals. Black magic
% to change
% catcodes and avoid blank lines being interpreted as |\par|. 
%    \begin{macrocode}
  \AnotherLinetrue\InBibfalse
  \makeOther\ \makeOther\\\makeOther\$%
  \makeOther\#\makeOther\^\makeOther\^^K%
  \makeOther\_\makeOther\^^A\makeOther\%%
  \makeOther\~\makeOther\{\makeOther\}\makeOther\&%
  \endlinechar-1\relax
%    \end{macrocode}
% Start looping through the file. Read each line to the macro
% |\Line|. If we have got to EOF or end of biblio then set
% |\AnotherLinefalse|.
%    \begin{macrocode}
  \loop
  \read\tf@biblio to \Line
  \ifx\Line\EndBib\AnotherLinefalse\InBibfalse\fi
  \ifeof\tf@biblio \AnotherLinefalse\fi
%    \end{macrocode}
% If we are in the biblio and the line is not blank then append it to
% the |\ADS@references| macro together with a newline. 
%    \begin{macrocode}
  \ifInBib
  \ifx\Line\empty\else\xdef\ADS@references{\ADS@references\Line^^J}\fi
  \fi
%    \end{macrocode}
% Check if the line we just saw was the start of biblio, then repeat
% the loop if necessary. 
%    \begin{macrocode}
  \ifx\Line\StartBib\InBibtrue\fi
  \ifx\Line\StartBibX\InBibtrue\fi
  \ifAnotherLine\repeat
  \immediate\closein\tf@biblio
  \endgroup
  }
%    \end{macrocode}
% \end{macro}
% 
% \begin{macro}{\PadToFourWithDots}
% Utility function for left-padding a number to 4 characters with
% periods. Used in 
% constructing |\ADS@BibCode|. This is a very simple algorithm and
% \emph{requires} that the argument actually be a number---if it is
% not then dire consequences will result.  
%    \begin{macrocode}
\newcommand{\PadToFourWithDots}[1]{%
  \ifnum #1 > 999 #1%
  \else\ifnum #1 > 99 .#1%
  \else\ifnum #1 > 9 ..#1%
  \else ...#1%
  \fi\fi\fi}
%%  \ifnum \count0 > 999 #1%
%%  \else\ifnum \count0 > 99 .#1%
%%  \else\ifnum \count0 > 9 ..#1%
%%  \else ...#1%
%%  \fi\fi\fi}
%    \end{macrocode}
% \end{macro}
%
% \begin{macro}{\AuthorADS}
% \begin{macro}{\ADS@affiliations}
% New mechanism for setting the affiliations for ADS. The author labels (AA, AB, \dots) are generated using two nested counters.  We use the |\robustify| command from \texttt{etoolbox} so that accented characters in the affiliation addresses do not blow up.  The list of accents may need to be expanded. 
%    \begin{macrocode}
\newcounter{AUTHOR}
\newcounter{author}[AUTHOR]
\stepcounter{AUTHOR}
\newcommand\ADS@affiliations{}
\robustify{\'}\robustify{\i}\robustify{\~}\robustify{\`}
\def\AuthorADS#1#2{%
  \begingroup
  \def\ORCIDiD##1{<ID system="ORCID">##1</ID>}
  \def\Email##1{<EMAIL>##1</EMAIL>}
  \def\Affil##1{##1}
  \stepcounter{author}
  \ifthenelse{\value{author} > 26}{\stepcounter{AUTHOR}\stepcounter{author}}{}
  \ifthenelse{\equal{\ADS@affiliations}{\@rmaaEmpty}}{%
    \xdef\ADS@affiliations{\Alph{AUTHOR}\Alph{author}(#2)}
  }{%
    \xdef\ADS@affiliations{\ADS@affiliations, \Alph{AUTHOR}\Alph{author}(#2)}
  }%
  \endgroup
}
%    \end{macrocode}
% \end{macro}
% \end{macro}
% %
% \begin{macro}{\ADS@dumpdata}
% \changes{4.0}{2019/11/13}{Write all author affiliations for ADS}
% \begin{macro}{\ADS@write}
% \begin{macro}{\ADS@openfile}
% |\ADS@dumpdata| is the macro that actually does the job of writing out all the
% data to an auxiliary file, using |\ADS@write| to write each
% field. The file is set up with |\ADS@openfile|, which currently is
% done in the TOC of a book.  The new multi-author affiliations are used if |\ADS@affiliations| has been set, otherwise it falls back to the single-author version |\ADS@affil|. 
%    \begin{macrocode}
\newcommand{\ADSfield}[2]{\relax}
\newcommand{\ADSvfield}[2]{\relax}
%%\renewcommand{\ADSfield}[12]{\%#1 #2}
\newcommand{\ADS@dumpdata}
{ 
  \ADS@write{\string\begin{ADSarticle}}
  \ADS@WRITE{\string\ADSfield{R}{\ADS@BibCode}}
  \ADS@write{\string\ADSfield{A}{\ADS@authors}}
  \ifthenelse{\equal{\ADS@affiliations}{\@rmaaEmpty}}{%
    {\ADS@sanitize\ADS@write{\string\ADSfield{F}{AA(\ADS@affil)}}}
  }{%
    {\ADS@sanitize\ADS@write{\string\ADSfield{F}{\ADS@affiliations}}}
  }
  \ADS@write{\string\ADSfield{J}{\ADS@journalname}}
  \ADS@write{\string\ADSfield{D}{\ADS@fulldate}}
  \ADS@write{\string\ADSfield{L}{\ADS@lastpage}}
  {\ADS@sanitize\ADS@write{\string\ADSfield{T}{\ADS@title}}}
  \ADS@write{\string\ADSfield{C}{\ADS@copyright}}
  \ADS@write{\string\ADSfield{d}{\ADS@doi}}
  \ADS@write{\string\ADSfield{U}{\ADS@url}}
  \ADS@write{\string\ADSfield{K}{\ADS@keywords}}
  {\ADS@sanitize\ADS@write{\string\ADSfield{B}{\ADS@abstract}}}
  \ADS@write{\string\ADSfield{Z}{\ADS@references}}
  \ADS@write{\string\end{ADSarticle}}
  }
%%\newcommand{\ADS@write}[1]{\if@book\addtocontents{ads}{#1}\fi}
%%\newcommand{\ADS@openfile}{\@starttoc{ads}}
\newcommand{\ADS@write}[1]{\set@display@protect\immediate\write\tf@ads{#1}}
\newcommand{\ADS@WRITE}[1]{\immediate\write\tf@ads{#1}}
\newcommand{\ADS@openfile}{%
  \newwrite\tf@ads%
  \immediate\openout\tf@ads\jobname.ads\relax
  \ADS@WRITE{\ADS@preamble}}
\newcommand{\ADS@closefile}{\ADS@WRITE{\ADS@postamble}}
\def\ADS@preamble{%
  %%  \string\documentclass{article}
  %%  \string\begin{document}
  }
\def\ADS@postamble{%
  %%  \string\end{document}
  \string%\string% Local Variables:^J 
  \string%\string% mode: LaTeX^J
  \string%\string% mode: Font-lock^J
  \string%\string% End:^J
  }
%    \end{macrocode}
% \end{macro}
% \end{macro}
% \end{macro}
% 
% \subsection{Automatic calculation of the page range}
% \label{sec:page-range}
% Another issue that has cost me a lot of grief. Still, it finally
% seems to be sorted now. All hail Jeff Goldberg. 
%
% Cunning trick to get the page numbers from |\pageref|.  Put a label
% on the first and last page of the document using |\AtBeginDocument|
% and |\AtEndDocument|. .  This way won't work for the book option.
% In that case, we do it by defining |\@endofarticlehook| and
% |\@begofarticlehook| as parallels to |\AtEndDocument| and
% |\AtBeginDocument|.  Note that we have to defer adding
% |\@begofarticlehook| to |\AtBeginDocument| to make sure it comes
% after the redefinitions of |\ref| and |\label|.
%    \begin{macrocode}
\if@book\AtEndDocument{\ADS@closefile}\else
 %\AtBeginDocument{\@begofarticlehook}
\AtEndDocument{\@endofarticlehook}
\fi
%    \end{macrocode}
% \begin{macro}{\@begofarticlehook}
% Since we are getting more and more things that need to be done at
% the beginning of every article, frontsection, etc we define a sister
% hook to 
% |\@endofarticlehook|. This hook does the following: 
% \begin{enumerate}
% \item Labels the first page
% \item Sets |\ADS@firstpage|
% \item runs and then clears |\indexauthor@deferlist|
% \item resets the keyword counter
% \end{enumerate}
% It is executed in the following circumstances: 
% \begin{enumerate}
% \item At |\begin{document}| for standalone files
% \item At |\begin{document}| for each included file
% \item At beginning of front sections 
% \item At beginning of TOC
% \item At beginning of Author Index
% \item At beginning of abstracts
% \end{enumerate}
% For some reason, this is not always within the right group, so we
% use |\xdef| just to be sure. 
%    \begin{macrocode}
\newcommand{\@begofarticlehook}{%
  \if@frontmatter\frontlabel{RMAAFirstPage}\else\label{RMAAFirstPage}\fi%
  \xdef\ADS@firstpage{\@arabic{\c@page}}%
  \if@book\else\makeindex\fi % needed so default index is defined in stand-alone articles
%%  \indexauthor@deferlist
%%  \xdef\indexauthor@deferlist{}
  \@inpreamblefalse
  \setcounter{kwdcounter}{0}
  }
%    \end{macrocode}
% \end{macro}
% \begin{macro}{\@endofarticlehook}
% This is a hook to be run at the end of each article
% currently, it just labels the last page and writes out all the ADS
% data to an auxiliary file.  
% It should not produce any printed output. 
% It is modelled on the |\AtEndDocument| command, to which it reduces
% in the case of a single article. 
%    \begin{macrocode}
\newcommand{\@endofarticlehook}{%
  \if@debug{\ClassWarning{rmaa}{***Doing Hook***}}\fi
  \if@debug{\ClassWarning{rmaa}{ADSabstract is: \ADS@abstract}}\fi
  \clearpage\rmaa@lastpagelabel\if@book\ADS@dumpdata\fi
  \ifx\NotesList\@rmaaEmpty\else\OutputNotes\fi
  \if@debug{\ClassWarning{rmaa}{Notes list is: \NotesList}}\fi
}
%    \end{macrocode}
% \begin{macro}{\rmaa@lastpagelabel}
% New version of last page label command copied from Jeff Goldberg's
% lastpage.sty package - we can't just use |\label| here in case the
% output routine for the last page has already been called. Instead,
% we first explicitly clear the page, step back the page counter,
% explicitly write to the .aux file, then restep the page counter
% again.
%
% But this will conflict with the hyperref package, which does labels differently.  Compensate for that by using the five-argument form.
% 
% Also set the |\ADS@LastPage| macro to the appropriate number. 
%    \begin{macrocode}
\newcommand{\rmaa@lastpagelabel}{%
  \addtocounter{page}{-1}%
  \immediate\write\@auxout{\string
    \newlabel{\articlename:RMAALastPage}{{}{\thepage}{}{}{}}}%
  \edef\ADS@lastpage{\thepage}%
  \addtocounter{page}{1}}
%    \end{macrocode}
% \end{macro}
% \end{macro}

% \subsection{Deal with explicit change in number of columns}
% \label{sec:columns}
% Sometimes, when formatting a book, we need to switch between
% one-column articles and two-column articles. For this to work, we
% need to wrap the standard \LaTeX{} commands to set the |\@rmtwocol|
% flag appropriately and also reset the package default values of all
% scalable lengths. .  
%    \begin{macrocode}
\let\latexonecolumn=\onecolumn
\let\latextwocolumn=\twocolumn
\renewcommand{\onecolumn}{\@rmtwocolfalse\@package@default@lengths\latexonecolumn}
\renewcommand{\twocolumn}{\@rmtwocoltrue\@package@default@lengths\latextwocolumn}
%    \end{macrocode}
% 
% \subsection{The actual formatting of the title}
% \label{sec:title-format}

% Note, we have to redefine the |\maketitle| command to use the
% numbered footnotes. This means we needn't use the atbeginend stuff. 
%    \begin{macrocode}
\renewcommand{\maketitle}{\par
  \begingroup
%    \end{macrocode}
% Don't switch to marked footnotes. 
% This is consistent with existing style but I'm not sure I like it
%    \begin{macrocode}
  %%  \renewcommand\thefootnote{\fnsymbol{footnote}}%
%%  \def\@makefnmark{\hbox to\z@{$\m@th^{\@thefnmark}$\hss}}%
  \def\@makefnmark{\textsuperscript{\@thefnmark}\hspace{-0.2em}}%
  \long\def\@makefntext##1{\parindent 1em\noindent
    \hbox to1.8em{\hss$\m@th^{\@thefnmark}$}##1}%
  \if@rmtwocol
  { \if@debug{\ClassWarning{rmaa}{TWO COLUMN}}\fi
    \twocolumn[\@maketitle]%
    }
  \else
  { \if@debug{\ClassWarning{rmaa}{ONE COLUMN}}\fi
    \newpage
    \global\@topnum\z@   % Prevents figures from going at top of page.
    \if@manuscript \@MSmaketitle \else \@maketitle \fi
    }
  \fi
  \if@manuscript\else\thispagestyle{rmaafirstpage}\@thanks\fi
  \ifx\NotesList\@rmaaEmpty\else{%
    \NoteMark{1--\arabic{notecounter}}\vspace*{-2\baselineskip}%
  }\fi%
  \endgroup
%    \end{macrocode}
% Don't reset footnote counter (not sure whether this is good thing
% or not).
%    \begin{macrocode}
  %%  \setcounter{footnote}{0}%
  \let\thanks\relax\let\affil\relax\let\altaffil\relax
  \let\maketitle\relax\let\@maketitle\relax
  \gdef\@thanks{}\gdef\@author{}\gdef\@title{}}
%    \end{macrocode}

% \begin{macro}{\@kwhwidth}
% \begin{macro}{\@kwtwidth}
% Redefine \cs{@maketitle} to include the abstract, resumen and
% keywords. Note that in the manuscripty style we use
% \cs{@MSmaketitle} instead.  
% \changes{3.16}{2003/04/07}{
%  Added extra vertical space between authors and accepted date.}
% \changes{3.26}{2011/08/15}{Add extra label to review articles}
%    \begin{macrocode}
\newlength{\@kwhwidth}\newlength{\@kwtwidth}
\newcommand\reviewname{Review}
\renewcommand{\@maketitle}{%
  \newpage
  \null
  \if@reviewarticle
  \smallskip
  {\large\bfseries\itshape\reviewname}%
  \bigskip
  \else
  \vskip\@before@title@skip%
  \fi
  \begin{center}%
    {\@title \par}%
    \vskip\@before@author@skip%
    \parbox{\rmaa@AuthorWidthFrac\textwidth}{\centering
      \lineskip\@between@author@skip%
      \@author
      \par}
    %% this should be replaced with `Submitted to...' for preprints
    \if@printdates
    \vskip\@before@dates@skip%
    \if@preprint
    \rmaa@ReceivedStyle{Submitted to RMxAA \@received}
    \else
    \rmaa@ReceivedStyle{Received \@received; accepted \@accepted}
    \fi
    \fi
    \if@manuscript
    \rmaa@ReceivedStyle{Submitted to RMxAA \@received}
    \fi
    \if@guide
    \rmaa@ReceivedStyle{Version \rmaa@version, \rmaa@date}
    \fi
    \vskip\@before@abstract@skip%
    \if@abstractpresent
    \parbox{\rmaa@AbstractWidthFrac\textwidth}{%
      \setlength{\parindent}{0.9cm}%
      \centerline{\rmaa@AbstractHeaderStyle{Abstract}}\par%
      \vskip\@after@abstracthead@skip%
      \if@newstyle\noindent\fi\@abstract
      \vskip\@between@abstract@skip%
      \centerline{\rmaa@AbstractHeaderStyle{Resumen}}\par%
      \vskip\@after@abstracthead@skip%
      \if@newstyle\noindent\fi\@resumen
      \if@setresumen\else\@abstract\fi % append english abstract if
                                       % spanish one not set
      \if@guide\else            %do not use keywords for guide
      \vskip\@before@kwd@skip%
      % Format the Key Words with hanging indentation:
      % They are put in a \parbox with width equal to the line width
      % minus the width of the header. 
      { % Widths of keyword
        \settowidth{\@kwhwidth}{\@keywordheader~}   % header and
        \setlength{\@kwtwidth}{\linewidth}          % keyword text
        \addtolength{\@kwtwidth}{-\@kwhwidth}       %
        \noindent\@keywordheader~\parbox[t]{\@kwtwidth}{%
          \rmaa@KeywordStyle{\boldmath\uppercase{\@keywords}}}}%
      \fi
      }
      \fi
  \end{center}%
  \vspace*{\afterkwdskip}
  \par
  }
\newlength{\afterkwdskip}\setlength{\afterkwdskip}{\z@}
%    \end{macrocode}
% \end{macro}
% \end{macro}
%
%
% \begin{macro}{\reviewarticle}
% \begin{macro}{\reviewtoc}
% \begin{macro}{\reviewlof}
%   User command to specify a review article. Must be used in
%   preamble. Main effect is to place the text ``Review'' before the
%   article title.  Also add command for setting a table of contents
%   and a list of figures at the beginning of the article.
%    \begin{macrocode}
\newcommand{\reviewarticle}{\@reviewarticletrue}
\newcommand\reviewtoc{%
  \section*{CONTENTS}%
  \@starttoc{rtc}%
}
\newcommand\reviewlof{%
  \section*{LIST OF FIGURES}%
  \@starttoc{lof}%
}
%    \end{macrocode}
% \end{macro}
% \end{macro}
% \end{macro}
% \begin{macro}{\l@section}
%    \begin{macrocode}
\renewcommand*\l@section[2]{%
  \ifnum \c@tocdepth >\z@
    \addpenalty\@secpenalty
    \addvspace{0.8em \@plus\p@}%
    \setlength\@tempdima{1.5em}%
    \begingroup
      \parindent \z@ \rightskip \@pnumwidth
      \parfillskip -\@pnumwidth
      \leavevmode %\bfseries
      \advance\leftskip\@tempdima
      \hskip -\leftskip
      #1\nobreak\hfil \nobreak\hb@xt@\@pnumwidth{\hss #2}\par
    \endgroup
  \fi}
%    \end{macrocode}
% \end{macro}
%
% \begin{macro}{\@MSmaketitle}
% A special version of |\@maketitle| for the \texttt{manuscript}
% option. This can deal with a multi-page abstract/resumen. Not very
% well, mind. 
%    \begin{macrocode}
\newcommand{\@MSmaketitle}{%
  \if@debug\ClassInfo{rmaa}{Using special MS maketitle command}\fi
  \newpage
  {
    \thispagestyle{rmaafirstpage}
%     \setlength{\parindent}{0pt}
    \null\vskip 2ex
    \begin{center}
      \@title\par
      \@author\par
      \rmaa@ReceivedStyle{%
        \ifx\@received\@rmaaEmpty
          Draft version: \@date
        \else
          Received: \@received
          \ifx\@accepted\@rmaaEmpty
          \else
            ; Accepted: \@accepted
          \fi
        \fi}\par
    \end{center}
    \@thanks
    \centerline{\rmaa@AbstractHeaderStyle{Resumen}}
    \noindent\@resumen\par\vskip 2ex
    \centerline{\rmaa@AbstractHeaderStyle{Abstract}}
    \noindent\@abstract\par\vskip 2ex
    \noindent\@keywordheader~\rmaa@KeywordStyle{\@keywords}\par
%     \@keywordheader~\rmaa@KeywordStyle{\uppercase{\@keywords}}\par
  }%
  }
%    \end{macrocode}
% \end{macro}

% \subsection{The sectioning commands}
% \label{sec:section-commands}

% \begin{macro}{\section}
% \begin{macro}{\subsection}
% \begin{macro}{\subsubsection}
% \begin{macro}{\paragraph}
% \begin{macro}{\subparagraph}
% Redefine all the standard commands to use these lengths. 
%    \begin{macrocode}
\renewcommand\section{\@startsection {section}{1}{\z@}%
                                   {\abovesectionskip}%
                                   {\belowsectionskip}%
                                   %% how do we get upper case here
                                   {\centering\reset@font\normalsize}}
\renewcommand\subsection{\@startsection{subsection}{2}{\z@}%
                                     {\abovesectionskip}%
                                     {\belowsectionskip}%
                                     {\centering\reset@font\normalsize\itshape}}
\renewcommand\subsubsection{\@startsection{subsubsection}{3}{\z@}%
                                     {\abovesectionskip}%
                                     {\belowsectionskip}%
                                     {\centering\reset@font\normalsize\itshape}}
\renewcommand\paragraph{\@startsection{paragraph}{4}{\z@}%
                                    {\aboveparaskip}%
                                    {\belowparaskip}%
                                    {\reset@font\normalsize\bfseries}}
\renewcommand\subparagraph{\@startsection{subparagraph}{5}{\parindent}%
                                      {\aboveparaskip}%
                                      {\belowparaskip}%
                                      {\reset@font\normalsize\bfseries}}
%    \end{macrocode}
% \end{macro}
% \end{macro}
% \end{macro}
% \end{macro}
% \end{macro}

% \begin{macro}{secnumdepth}
% Make sure that we number the subsubsections---this is done by
% default in the |article| class but not the |book| class, which is
% why we must set it explicitly here.
%    \begin{macrocode}
\setcounter{secnumdepth}{3}
%    \end{macrocode}
% \end{macro}
% \changes{2.2beta10}{2001/11/07}{Fixed \cs{secnumdepth} for \cs{book}
% style.}

% \begin{macro}{\thesubsection}
% \begin{macro}{\thesubsubsection}
% For subsection and subsubsection, we want the numbers to be roman
% but the title to be italic. This seems the easiest way to go about it.
%    \begin{macrocode}
\renewcommand{\thesubsection}{{\upshape\thesection.\@arabic\c@subsection}}
\renewcommand{\thesubsubsection}{\thesubsection{\upshape.\@arabic\c@subsubsection}}
%    \end{macrocode}
% \end{macro}
% \end{macro}

% \begin{macro}{\@seccntformat}
% Add ``.'' to the number here so we can remove it from |\thesection|,
% etc. Reduce gap between number and title WJH 27 Jan 1995
%    \begin{macrocode}
\def\@seccntformat#1{\csname the#1\endcsname .\hskip 0.4em}
%    \end{macrocode}
% \end{macro}

% \begin{macro}{\@sect}
% I didn't really want to mess with the |\@sect| command, but it was
% the only way I could manage to force upper case. Other than the
% commented parts, it is identical to the standard version. 
%    \begin{macrocode}
\def\@sect#1#2#3#4#5#6[#7]#8{
  \ifnum #2>\c@secnumdepth
     \let\@svsec\@empty
  \else
     \refstepcounter{#1}%
     \let\@@protect\protect
     \def\protect{\noexpand\protect\noexpand}%
     \edef\@svsec{\@seccntformat{#1}}%
     \let\protect\@@protect
  \fi
%    \end{macrocode}
% This is the bit I changed (27 Jan 1995). Use upper case for
% |\section|, use mixed case for the rest. 
%    \begin{macrocode}
  \ifnum #2=1                   
     \def\sec@upcase##1{\uppercase{##1}}%
     \def\sec@shape{\upshape}%
  \else                         
     \def\sec@upcase##1{\relax{##1}}%
     \def\sec@shape{\itshape}%
  \fi
  \@tempskipa #5\relax
  \ifdim \@tempskipa>\z@
     \begingroup #6\relax
          \@hangfrom{\hskip #3\relax\@svsec}%
                    {\interlinepenalty \@M \sec@upcase{#8}\par}%
     \endgroup
     \csname #1mark\endcsname{#7}%
%    \end{macrocode}
% Only write TOC entries if doing the Guide style. 
%    \begin{macrocode}
     \if@guide 
        \addcontentsline{toc}{#1}{%
          \ifnum #2>\c@secnumdepth \else
             \protect\numberline{\csname the#1\endcsname}%
          \fi
          #7}%
     \fi
%    \end{macrocode}
% Also allow per-article TOC for review papers
%    \begin{macrocode}
     \if@reviewarticle
        \ifnum #2>2 \else
           \addcontentsline{rtc}{#1}{%
             \protect\numberline{\csname the#1\endcsname}%
             \sec@upcase{\sec@shape #7}}%
        \fi
     \fi
  \else
     \def\@svsechd{#6\hskip #3\relax  %% \relax added 2 May 90
       \@svsec #8\csname #1mark\endcsname{#7}%
       \if@guide
          \addcontentsline{toc}{#1}{%
            \ifnum #2>\c@secnumdepth \else
               \protect\numberline{\csname the#1\endcsname}%
            \fi
            #7}%
       \fi
       }%
  \fi
  \@xsect{#5}}
%    \end{macrocode}
% \end{macro}


% \subsection{Appendices}
% \label{sec:appendices}
% Make sure that appendices in individual articles work properly. We
% copy the definition from the article class (the book one doesn't do
% what we want).  
% \begin{environment}{appendices}
%   This uses an environment |appendices| inside of which each
%   |\section| will become a separate appendix.  Note that the
%   appendices must come before the references (otherwise the full
%   addresses mechanism will not work properly). We make the equation
%   numbering be subordinate to the appendix (equations A12, B6,
%   etc). Finally, we have a medskip followed by APPENDICES to
%   introduce them. \changes{3.17}{2002/04/08}{ Added appendices environment
%   that does the right thing.}
%    \begin{macrocode}
\newenvironment{appendices}{\par
  \begingroup
  \setcounter{section}{0}%
  \setcounter{subsection}{0}%
  \renewcommand\thesection{\@Alph\c@section}
  \renewcommand\theequation{\thesection\arabic{equation}}
  \medskip
  \section*{APPENDICES}}
{\endgroup}
%    \end{macrocode}
% \end{environment}
% \begin{environment}{appendix}
%   If there is only one appendix, then use |appendix| environment
%   instead.  \changes{4.0}{2019/11/12}{ Added
%   appendix environment for case where there is only one section.}
%    \begin{macrocode}
\renewenvironment{appendix}{\par
  \begingroup
  \setcounter{section}{0}%
  \setcounter{subsection}{0}%
  \renewcommand\thesection{\@Alph\c@section}
  \renewcommand\theequation{\thesection\arabic{equation}}
  \medskip
  \section*{APPENDIX}}
{\endgroup}
%    \end{macrocode}
% \end{environment}

% \subsection{Floating thingies}
% \label{sec:floats}   

%% 
%    \begin{macrocode}
\renewcommand{\thetable}{\@arabic\c@table}
\renewcommand{\thefigure}{\@arabic\c@figure}
\renewcommand{\figurename}{Fig.}
\renewcommand{\tablename}{TABLE}
%\renewcommand{\fnum@table}{\MakeUpperCase{\tablename}~\thetable}
%    \end{macrocode}
% For some reason, I have to write `TABLE' explicitly here
% If I use |\tablename|, it gets converted to mixed case---God knows
% why. 
% |\figurename| is also being ignored. 
%    \begin{macrocode}
\renewcommand{\fnum@table}{TABLE~\thetable}
\renewcommand{\fnum@figure}{Fig.~\thefigure}
%    \end{macrocode}

% Do captions in 9 pt face in the main journal style and in new-style
% anything. 
%    \begin{macrocode}
\long\def\@makefigcap#1#2{%
  \if@proceedings\else\small\fi
  \if@newstyle\small\fi 
  \vskip\abovecaptionskip
  \sbox\@tempboxa{#1. #2}%
  \ifdim \wd\@tempboxa >\hsize
    #1. #2 \par
  \else
    \global \@minipagefalse
    \hb@xt@\hsize{\hfil\box\@tempboxa\hfil}%
  \fi
  \vskip\belowcaptionskip}
%    \end{macrocode}

% Make sure table caption is upper case. 
%    \begin{macrocode}
\long\def\@maketabcap#1#2{%
  \hbox to\hsize{\hfil #1 \hfil}
  \vskip 6pt 
  \centering\MakeUpperCase{#2}\par
  \vskip 6pt}
%    \end{macrocode}

%    \begin{macrocode}
\newlength{\@rottabcapwidth}
\long\def\@makerottabcap#1#2{%
  \setbox\@tempboxa\hbox{\MakeUpperCase{#2}}%
  \ifdim\wd\@tempboxa >.8\vsize
     \setlength\@rottabcapwidth{.8\vsize}
  \else
     \setlength\@rottabcapwidth{\wd\@tempboxa}
  \fi
  \rotatebox{90}{%
    \begin{minipage}{\vsize}\centering
      \begin{minipage}{\@rottabcapwidth}\centering
        #1\\ \vskip 6pt \MakeUpperCase{#2}
      \end{minipage}%
    \end{minipage}
    }%
  \hskip 6pt
  }

\newenvironment{Sideways}
    {\begin{sideways}\begin{minipage}{\vsize}\centering}
    {\end{minipage}\end{sideways}}

\newenvironment{Sidewaystable}
{%
  \let\oldcaption=\caption%
  \def\caption{
    \def\@captype{table}\renewcommand{\@makecaption}{\@maketabcap}
    \oldcaption}
  \begin{sidewaystable}\centering}
  {\end{sidewaystable}\let\caption=\oldcaption}
\newenvironment{Sidewaystable*}
{%
  \let\oldcaption=\caption%
  \def\caption{
    \def\@captype{table}\renewcommand{\@makecaption}{\@maketabcap}
    \oldcaption}
  \begin{sidewaystable*}\centering}
  {\end{sidewaystable*}\let\caption=\oldcaption}
%    \end{macrocode}

% Not sure why we need to stick these lines in. 
%    \begin{macrocode}
\renewenvironment{figure}
                 {\renewcommand{\@makecaption}{\@makefigcap}\@float{figure}}
                 {\end@float}
\renewenvironment{figure*}
                 {\renewcommand{\@makecaption}{\@makefigcap}\@dblfloat{figure}}
                 {\end@dblfloat}
\renewenvironment{table}
                 {\renewcommand{\@makecaption}{\@maketabcap}\@float{table}}
                 {\end@float}
\renewenvironment{table*}
                 {\renewcommand{\@makecaption}{\@maketabcap}\@dblfloat{table}}
                 {\end@dblfloat}
%    \end{macrocode}

% Redefine |\@caption| so that we don't bother with the list of
% figures. This saves us having to worry about fragile commands in the
% captions. 
%    \begin{macrocode}
\long\def\@caption#1[#2]#3{%
  \par
%%  \addcontentsline{\csname ext@#1\endcsname}{#1}%
%%    {\protect\numberline{\csname the#1\endcsname}{\ignorespaces #2}}%
  \begingroup
    \@parboxrestore
    \if@minipage
      \@setminipage
    \fi
    \normalsize
    \@makecaption{\csname fnum@#1\endcsname}{\ignorespaces #3}\par
  \endgroup}
%    \end{macrocode}

% Make tables not quite so cramped. Allow |\hline\hline| to give a
% poor man's thick rule. 
%    \begin{macrocode}
\renewcommand{\arraystretch}{1.15}
\setlength{\doublerulesep}{0.4pt}
%    \end{macrocode}

% \begin{macro}{\tablecols}
% \begin{macro}{\@tablecols}
% \begin{macro}{\tabcaptioncontinued}
% |\tablecols| is a user command for setting the number of columns. It
% is used by commands that try and span the entire table (I think
% there must be a better way of acheiving
% this). |\tabcaptioncontinued| is 
% a user command for setting the caption of a continued table on the
% next page. 
%    \begin{macrocode}
\newcommand{\tablecols}[1]{\renewcommand{\@tablecols}{#1}}
\newcommand{\@tablecols}{1}
\newcommand{\tabcaptioncontinued}{%
  \multicolumn{\@tablecols}{c}{\addtocounter{table}{-1}TABLE~\thetable\ (CONTINUED)} \\ \\
  }
%    \end{macrocode}
% \end{macro}
% \end{macro}
% \end{macro}

% \subsubsection{Footnotes in tables}
% \label{sec:table-footnotes}
% \begin{macro}{\tablenotemark}
% \begin{macro}{\tablenote}
% \begin{macro}{\tablenoteskip}
% Old user commands |\tablenotemark{MARK}|,
% |\tablenote{MARK}{TEXT}|---now deprecated. 
%    \begin{macrocode}
\newcommand\tablenotemark[1]{$^\mathrm{#1}$} 
\newcommand\tablenote[2]{\if@proceedings\else\small\fi ${}^\mathrm{#1}$ #2} 
\newcommand\tablenoteskip{\\[0.5\baselineskip]} 
%    \end{macrocode}
% \end{macro}
% \end{macro}
% \end{macro}

% \begin{macro}{\tabnotemark}
% \begin{macro}{\tabnotetext}
% \begin{macro}{\tabnotetextTWO}
% Alternative way of doing footnotes in tables. 
%    \begin{macrocode}
\newcommand{\tabnotemark}[1]{\textsuperscript{\lowercase{#1}}}
%\newcommand{\tabnotetextTWO}[4]{\tabnotetext{#1}{#2}\tabnotetext{#3}{#4}}
\newcommand{\tabnotetextTWO}[4]{%
  \multicolumn{\@tablecols}{l}{%
    \parbox{0.5\tabnotewidth}{\small\textsuperscript{#1}#2}%
    \parbox{0.5\tabnotewidth}{\small\textsuperscript{#3}#4}%
    }%
  \\}
\newcommand{\tabnotetext}[2]{%
  \multicolumn{\@tablecols}{@{}p{\tabnotewidth}@{}}{\small\textsuperscript{#1}#2}\\}
\newlength{\tabnotewidth}
%    \end{macrocode}
% \end{macro}
% \end{macro}
% \end{macro}


% \begin{macro}{changemargin}
%    \begin{macrocode}
% Command to allow margins to be changed on the fly. Useful to allow
% tables and Figures to protrude into the margins. 
%
% \changes{3.24}{2007/09/14}{Added \texttt{changemargin} environment
% (copied from TeX FAQ).}
\newenvironment{changemargin}[2]{% 
  % From TeX FAQ http://www.tex.ac.uk/cgi-bin/texfaq2html?label=chngmargonfly
  \begin{list}{}{%
    \setlength{\topsep}{0pt}%
    \setlength{\leftmargin}{#1}%
    \setlength{\rightmargin}{#2}%
    \setlength{\listparindent}{\parindent}%
    \setlength{\itemindent}{\parindent}%
    \setlength{\parsep}{\parskip}%
  }%
  \item[]}{\end{list}}
%    \end{macrocode}
% \end{macro}

% \subsection{The bibliography and references}
% \label{sec:biblio}
% \begin{macro}{\beforebiblioskip}
% \begin{macro}{\afterbiblioskip}
% I'm not entirely sure why I introduced these two commands, since
% directly putting |\vspace| before and after
% |\begin{thebibliography}| seems to work well enough. 
%    \begin{macrocode}
\newlength{\beforebiblioskip}
\setlength{\beforebiblioskip}{0pt}
\newlength{\afterbiblioskip}
\setlength{\afterbiblioskip}{0pt}
%    \end{macrocode}
% \end{macro}
% \end{macro}
% \begin{macro}{thebibliography}
% The environment itself is done as starred section with extra
% vertical space after heading. This space is currently zero. 
% In the book style we do the cunning trick with rereading the file in
% order to get the references into the |\ADS@references| macro.
% 
% \changes{3.05}{2003/01/02}{Removed the \cs{beforebiblioskip} and
% \cs{afterbiblioskip} commands, since they cause unwanted vertical
% space in some instances.}  
%    \begin{macrocode}
\renewenvironment{thebibliography}{%
  \if@book\RereadBiblio\fi
%%  \vspace*{\beforebiblioskip}
  \section*{REFERENCES} \frenchspacing
%%  \vspace*{\afterbiblioskip}
  \small
%    \end{macrocode} 
% The references themselves are set as a list in order 
% to get the hanging indentation. They are done in a smaller font and
% with french spacing so we needn't worry about the distinction
% between sentence-ending periods and abbreviations. 
%    \begin{macrocode} 
  \list{\null}{\leftmargin 1.5em\labelwidth\z@\labelsep\z@\itemindent -1.5em
    \itemsep\z@\parsep\z@\usecounter{enumi}}
  \sloppy\clubpenalty4000\widowpenalty4000
  \sfcode`\.=1000\relax\@biblisttrue}
{\endlist\@biblistfalse\normalsize\@outputaddressesifnecessary}
%    \end{macrocode}
% \end{macro}
% \subsection{Uniquification of cross-references}
% \label{sec:uniquification}
% 
% \begin{macro}{\articlename}
% In order that the |\label{}| command in individual articles be
% guaranteed to give a unique anchor, we append the article name (held
% in the |\articlename| command) to each key. This requires
% redefinition of the following commands: |\cite|, |\nocite|, |\ref|,
% |\pageref|, |\label|. 
% \begin{macro}{\oldcite}
% \begin{macro}{\oldnocite}
% \begin{macro}{\oldref}
% \begin{macro}{\oldpageref}
% \begin{macro}{\oldlabel}
% The original versions of these commands are
% saved as |\oldcite|, |\oldnocite|, etc. 
%
% The uniquifying of cross-references is now always done, so 
% we don't get multiply-defined labels in the single article case
%    \begin{macrocode}
\newcommand{\articlename}{anon}
\AtBeginDocument{
  \let\oldcite=\cite 
  \let\oldnocite=\nocite 
  \let\oldref=\ref 
  \let\oldpageref=\pageref
  \let\oldlabel=\label
}
%    \end{macrocode}
% \end{macro}
% \end{macro}
% \end{macro}
% \end{macro}
% \end{macro}
% \end{macro}
%
% Uniquify all cross-references by prepending the article name.  This
% now has to go into |\AtBeginDocument| so it is working on the new
% versions that hyperref.sty has set up.  Then we need to put the
% |\@begofarticlehook| afterwards, so that the first page label is using
% the new definitions.  If amsmath is loaded, then we also need to patch
% the special label command that they use. \changes{4.0}{2019/11/18}{Fix labels for amsmath display equations, in case that package has been loaded.}
%    \begin{macrocode}
\AtBeginDocument{
  \renewcommand{\ref}[1]{\oldref{\articlename:#1}%
    \if@debug{\ClassInfo{rmaa}{Reference to \articlename:#1}}\fi}
  \renewcommand{\pageref}[1]{\oldpageref{\articlename:#1}%
    \if@debug{\ClassInfo{rmaa}{Page reference to \articlename:#1}}\fi}
  \renewcommand{\label}[1]{\oldlabel{\articlename:#1}%
    \if@debug{\ClassInfo{rmaa}{Label for \articlename:#1}}\fi}%
  \ifdef{\label@in@display}{%
    \let\oldlabel@in@display=\label@in@display
    \renewcommand{\label@in@display}[1]{\oldlabel@in@display{\articlename:#1}%
      \if@debug{\ClassInfo{rmaa}{AMSTeX Label for \articlename:#1}}\fi}%
  }{}
  \@begofarticlehook
}%
%    \end{macrocode}
% Citations are more complicated since they can contain a list of
% keys. These are really only here for use with the rmaacite package,
% which is utterly deprecated (02 Jan 2003) now that I have the natbib
% code working. The command \cs{cite} should \emph{not} be used with
% natbib: use \cs{citep}, \cs{citet}, etc.\@ instead. 
%    \begin{macrocode}
\renewcommand{\nocite}[1]{\@bsphack
  \@for\@citeb:=#1\do{%
    \edef\@citeb{\expandafter\@firstofone\@citeb:\articlename}%
    \if@filesw\immediate\write\@auxout{\string\citation{\@citeb}}\fi
    \@ifundefined{b@\@citeb:\articlename}{\G@refundefinedtrue
      \@latex@warning{Citation `\@citeb' undefined}}{}}%
  \@esphack}
\renewcommand{\cite}[1]{\oldcite{#1:\articlename}}
%   \end{macrocode}


% Horrible fiddling with the bibliography citation mechanism. 
% This has to be globally available in case rmaacite puts it in the
% \texttt{.aux} file. 
%    \begin{macrocode}
\def\lbibcite#1#2#3{\global\@namedef{b@#1}{#2}
                    \global\@namedef{by@#1}{#3}}
%    \end{macrocode}

% \begin{macro}{\oldbibitem}
% \begin{macro}{\@rmaaBibArg}
% \begin{macro}{\@rmaaSpace}
% \begin{macro}{rmaabib}
% Of course, |\bibitem| needs to get the uniquification treatment as
% well. While we are at it, we redefine it to use consecutive integers
% for the key if the user leaves it blank (using the counter
% |rmaabib|). This allows one to put 
% |\bibitem{}| for all your references without getting annoying
% messages about multiply-defined labels.
% 
% We also take the opportunity to insert the optional argument to
% \cs{bibitem} if it is missing in the form Author(DATE) so that natbib
% doesn't try to do numerical citations. Natbib already provides a
% hook \cs{@extra@b@citeb} to help uniquify the citation key, so we
% just plug into that. However, it puts it \emph{after} the original
% key, not before as I had had it. Easiest solution is to change the
% way that I do it. 02 Jan 2003. Note that the \cs{@extra@b@citeb}
% does not get written to the .aux  file in the \cs{citation} command
% but it does work. 
% \changes{3.14}{06 Feb 2003}{Reinstalled old definition of \cs{bibitem}
%  for the case where the nonatbib option is requested, so that the
%  rmaacite package will work. This means that you can't use rmaacite
%  and natbib in the same volume.}
%    \begin{macrocode}
\if@natbib\def\@extra@b@citeb{:\articlename}\fi
%%\else
\let\oldbibitem\bibitem
\newcommand{\@rmaaBibArg}{}
\newcommand{\@rmaaSpace}{ }
\newcounter{rmaabib}
\if@natbib
\renewcommand{\bibitem}[2][ABC(2003)]{%
  \renewcommand{\@rmaaBibArg}{#2}
  \ifx\@rmaaBibArg\@rmaaEmpty
     \addtocounter{rmaabib}{1}
     \oldbibitem[#1]{\thermaabib:\articlename}
  \else\ifx\@rmaaBibArg\@rmaaSpace
     \addtocounter{rmaabib}{1}
     \oldbibitem[#1]{\thermaabib:\articlename}
  \else 
     \oldbibitem[#1]{#2:\articlename}
  \fi\fi
  \renewcommand{\@rmaaBibArg}{}
  }
\else
\renewcommand{\bibitem}[1]{%
  \renewcommand{\@rmaaBibArg}{#1}
  \ifx\@rmaaBibArg\@rmaaEmpty
     \addtocounter{rmaabib}{1}
     \oldbibitem{\articlename:\thermaabib}
  \else\ifx\@rmaaBibArg\@rmaaSpace
     \addtocounter{rmaabib}{1}
     \oldbibitem{\articlename:\thermaabib}
  \else
     \oldbibitem{\articlename:#1}
  \fi\fi
  \renewcommand{\@rmaaBibArg}{}
  }
\fi
%%\fi
%    \end{macrocode}
% \end{macro}
% \end{macro}
% \end{macro}
% \end{macro}

% \begin{macro}{oldpcite}
% \begin{macro}{oldscite}
% Then, we need to repeat the above scheme, but for the |rmaacite|
% definition. 
%    \begin{macrocode}
\newcommand{\@redefinebibitem}{%
  \let\oldbibitem\bibitem
  \let\oldcite\cite
  \let\oldpcite\pcite
  \let\oldscite\scite
  \def\bibitem[##1<##2>]##3{%
    \renewcommand{\@rmaaBibArg}{##3}
    \ifx\@rmaaBibArg\@rmaaEmpty
       \addtocounter{rmaabib}{1}
       \oldbibitem[##1<##2>]{\thermaabib:\articlename}
    \else\ifx\@rmaaBibArg\@rmaaSpace
       \addtocounter{rmaabib}{1}
       \oldbibitem[##1<##2>]{\thermaabib:\articlename}
    \else 
       \oldbibitem[##1<##2>]{##3:\articlename}
    % trust the author here
%%       \oldbibitem[##1<##2>]{##3}
    \fi\fi
    }
  }
%    \end{macrocode}
% \end{macro}
% \end{macro}


% \subsection{HTML Table of Contents}
% \label{sec:hyper-TOC}
% \begin{macro}{\HyperTOCEntry}
% \begin{macro}{\HyperTOCpart}
% \begin{macro}{\HyperTOCmeeting}
% \begin{macro}{\HyperTOCeditors}
% \begin{macro}{\HyperTOCpreamble}
% These commands are used to write information to the \texttt{.htoc}
% file, which is used in generating a Web Page of the Table of
% Contents. Since they only ever appear preceded by |\protect| in the
% argument to an |\addtocontents| command,  they should never need to
% be expanded by the macros. However, we define them here just in case
% \dots
%    \begin{macrocode}
\newcommand{\HyperTOCEntry}[3]{\relax}
\newcommand{\HyperTOCpart}[1]{\relax}
\newcommand{\HyperTOCmeeting}[3]{\relax}
\newcommand{\HyperTOCeditors}[1]{\relax}
\newcommand{\HyperTOCpreamble}{\relax}
%    \end{macrocode}
% \end{macro}
% \end{macro}
% \end{macro}
% \end{macro}
% \end{macro}
%

% \subsection{Formatting the poster abstracts} 
% These can now be used without using the book option.
%    \begin{macrocode}

\newcommand{\@makeminititle}{%
  \newpage
  \null 
  \vskip 3em%
  \begin{center}%
    {\@title \par}%
  \end{center}
  \vskip 3em%
  }

\newif\if@abstractsads
\@abstractsadsfalse
\newcommand\abstractsADS{\@abstractsadstrue}

\if@book\RequirePackage{multicol}\fi
\newenvironment{abstracts}[1][Abstracts of Contributed Papers]{%
  %% First, make sure we start on an odd page (06 Jun 2000)
  \if@book\clearpage\thispagestyle{empty}\cleardoublepage\fi
  %%
  \renewcommand{\articlename}{ABSTRACTS}
  \title{#1}\shortauthor{Abstracts}\shorttitle{Abstracts}
  \listofauthors{}
  \if@book
  \addcontentsline{toc}{chapter}{\protect\@fiddleauthors{\@TOCtitle}{\rmaaTOCauthors}}%
  \addtocontents{htc}{\protect\HyperTOCEntry{\@rawtitle}{\@rawauthors}{\articlename}}
  \fi
  \pagestyle{rmaa} % this now seems necessary (06 Jun 2000)
  \thispagestyle{rmaafirstpage}
%%  \twocolumn[\@makeminititle]
  \renewcommand{\abstract}[1]{%
    \renewcommand\ADS@abstract{##1}%
    \renewcommand{\@abstract}{##1}\@writeabstract%
  }
  \renewcommand{\title}[2][]{%
    \gdef\ADS@authors{}%
    \gdef\ADS@affil{}%
    \gdef\ADS@references{}%
    \gdef\ADS@keywords{}%
    \renewcommand{\@rmaaArg}{##1}
    \ifx\@rmaaArg\@rmaaEmpty
       \renewcommand{\@TOCaftertitle}{}
    \else
       \renewcommand{\@TOCaftertitle}{ (\@rmaaArg)}
    \fi
    \renewcommand{\@title}{\rmaa@TOCTitleStyle{##2}}%
    \renewcommand{\@TOCtitle}{\rmaa@TOCTitleStyle{##2}\@TOCaftertitle}%
    \renewcommand{\rmaaIndexTitle}{\rmaa@IndexTitleStyle{##2\@TOCaftertitle.}}%
    \renewcommand\ADS@title{##2}%
    \xdef\ADS@firstpage{\@arabic{\c@page}}%
    }
  \renewcommand{\author}[1]{\let\@thanks\@empty\renewcommand\@author{##1} }
  \renewcommand{\footnotetext}[2][XX]{\par\noindent\textsuperscript{##1}~##2}
  \renewcommand{\footnoterule}{\kern-3\p@ \hrule \@width 3em \kern 2.6\p@}
  \if@book\@begofarticlehook\fi
  \if@book\begin{multicols}{2}[\@makeminititle]\else\@makeminititle\fi
  \begingroup
  \setcounter{footnote}{0}
  }% and at the end of the environment....
  {\endgroup\if@book\end{multicols}\listofabstracts\label{RMAALastPage}\clearpage\fi}
%% auxiliary commands used:
%% \@writeabstract does all the actual formatting 
\newcommand{\@writeabstract}{%
  \if@book%
  \addcontentsline{loa}{abstract}{\protect\@fiddleauthors{\@TOCtitle}{\rmaaTOCauthors}}%
  \fi%
  \begin{center}\@title\par\@author\end{center}
  \@abstract
  \par\noindent\footnoterule\nobreak\footnotesize\@thanks
  \par\noindent\hrulefill\goodbreak
  \vskip 1cm \@plus0.5cm
  %% write ADS entry
  \if@book\if@abstractsads\ADS@dumpdata\fi\fi%
  \endgroup % this pair is to make each abstract a self-contained group
  \begingroup
  \setcounter{footnote}{0}
  }

\if@book
%% List of abstracts (stuff written to .loa file)
\newcommand{\listofabstracts}{%
  \begin{multicols}{2}[%
    \vskip 1.5cm \vfill % leave at least 1.5 cm and go to bottom of page if possible
    \centerline{\MakeUpperCase{List of Abstracts}}
    \vskip 1.5cm \@plus0.5cm
    ][6cm] % this makes it start a new page if there is not at least 6cm of space
    \renewcommand{\@fiddleauthors}[2]{##1\hskip1em plus2em minus0.4em\penalty-100  ##2}
    \parindent\z@
%%    \parskip.5\p@ \@plus .7\p@\relax 
    \parskip3\p@ \@plus 1\p@ \@minus 1\p@\relax % allow quite a lot of
                                % leeway here to avoid breaks inside
                                % entries 
    \renewcommand{\\}{\relax} % temporarily turn off explicit line breaks just in case
    \renewcommand{\altaffilmark}[1]{\relax} % ditto with this
    \@starttoc{loa}
  \end{multicols}
  }
\newcommand*{\l@abstract}[2]{%
  \par\penalty-1000
  #1\hskip 1em\penalty0\null\nobreak\hfill#2}
\fi

%    \end{macrocode}
 
% \subsection{Code for typesetting an entire volume of conference
% proceedings} \label{sec:whole-volume}
%    \begin{macrocode}

\if@book
%% environment for including the individual articles
\newcommand{\@fiddleauthors}[2]{\relax}
\renewcommand{\theequation}{\@arabic\c@equation}
\newenvironment{articles}%
{ \begingroup % make it a group so that author definitions do not
              % spill over
  \gdef\ADS@authors{}
  \pagestyle{rmaa}
  \renewenvironment{document}{%
    \@begofarticlehook
    \addtocounter{chapter}{1}%
    \addcontentsline{toc}{chapter}{%
      \protect\@fiddleauthors{\@TOCtitle}{\rmaaTOCauthors}}%
    \addtocontents{htc}{\protect\HyperTOCEntry{\@rawtitle}{\@rawauthors}{\articlename}}
    }{\relax}
  \renewcommand{\documentclass}[2][foo]{\relax} 
  \renewcommand{\usepackage}[2][foo]{%
    \def\@packagename{##2}\def\@rmaacitename{rmaacite}%
    \ifx\@packagename\@rmaacitename
    \@rmcitetrue\ClassInfo{rmaa}{This paper uses rmaacite}
    \makeatletter\input{rmaacite.sty}\makeatother % read it in here so that all changes
    \else                                         % are local to paper
    \if@debug\ClassWarning{rmaa}{Package \@packagename\space was requested by \articlename\space}\fi
    \fi} 
  %% Don't allow individual articles to mess with the pages, volume or year
  \renewcommand{\SetFirstPage}[1]{\relax}
  \renewcommand{\SetYear}[1]{\relax}
  \renewcommand{\SetVolume}[1]{\relax}
  \renewcommand{\SetConfTitle}[1]{\relax}
  %% Don't allow individual articles to change these
  \renewcommand{\nonstopmode}{\relax}
  \renewcommand{\scrollmode}{\relax}
  \renewcommand{\outputplates}{\relax}
  %% XXXX need to check if using rmaacite and act accordingly
  %% avoid multiple declarations of the same length
  \let\oldnewlength=\newlength
  \renewcommand{\newlength}[1]{%
    \edef\reserved@a{\expandafter\@gobble\string ##1}%
    \@ifundefined{\reserved@a}{%
      \oldnewlength{##1}}{%
      \if@debug{\ClassInfo{rmaa}{\reserved@a\space already defined: doing nothing}}\fi%
      \relax}}%
  \renewcommand\newcommand{\providecommand}
  }%
{\endgroup\onecolumn}
     

%    \end{macrocode}
% Command to be used in `articles' environment for including an
% article
%    \begin{macrocode}
\newcommand{\article}[2][]{%
  % optional first argument is extra text to put in brackets after title in TOC
  % e.g. ``Invited Talk'' or ``Poster''
  \renewcommand{\articlename}{#2}
  % reset loads of counters
  \setcounter{footnote}{0}
  \setcounter{figure}{0}
  \setcounter{table}{0}
  \setcounter{section}{0}
  \setcounter{equation}{0}
  \setcounter{rmaabib}{0}
  \setcounter{author}{0}
  \setcounter{AUTHOR}{1}
  \gdef\ADS@authors{}%
  \gdef\ADS@affil{}%
  \gdef\ADS@affiliations{}%
  \gdef\ADS@references{}%
  \@firstauthortrue
  \@firstaffiltrue
  \@firstaltaffiltrue
  \@firstkwdtrue
  \@rmcitefalse
  \if@suppressfulladdresses\else\@needaddressestrue\fi
  \renewcommand{\@rmaaArg}{#1}
  \ifx\@rmaaArg\@rmaaEmpty
     \renewcommand{\@TOCaftertitle}{}
  \else
     \renewcommand{\@TOCaftertitle}{ (#1)}
  \fi
%    \end{macrocode}
% In the main journal articles must start on an odd page. 
%    \begin{macrocode}
  \if@journal\clearpage\thispagestyle{empty}\cleardoublepage\fi
%% Now input the file, then run the end-of-article hook (reinstated 28
  %% Sep 2001)
  {\input{\@articledir/\articlename.tex}\@endofarticlehook}}
%    \end{macrocode}
% Command to separate articles into sessions. Write an entry in the
% TOC, makes sure we start on an odd page (added 06 Jun 2000), then do
% a title page if \cs{SessionTitlePages} was called. 
% \changes{3.12}{04 Feb 2003}{Added option to have title page for each
% session. Turn it on with \cs{SessionTitlePages}.}
%    \begin{macrocode}
\newcommand{\session}[1]{
  %% this does nothing except write an entry in the TOC
  \addcontentsline{toc}{part}{\rmaa@TOCPartStyle{#1}}
  \addtocontents{htc}{\protect\HyperTOCpart{#1}}
  %% and also makes sure we start on an odd page (added 06 Jun 2000)
  \clearpage\thispagestyle{empty}\cleardoublepage
  \stepcounter{part}
  \if@sessiontitlepages{%
    \onecolumn
    \thispagestyle{empty}
    \renewcommand{\articlename}{Section\thepart TitlePage}
    \label{RMAAFirstPage}
    \make@session@title{#1}
    \label{RMAALastPage}
    \clearpage\thispagestyle{empty}\cleardoublepage
  }\fi
  }
\newcommand{\make@session@title}[1]{%
  {%
    \setlength\parindent{0pt}
    \vspace*{\stretch{1}}
    \rule{\linewidth}{1mm}
    \begin{flushright}\Huge\bfseries\boldmath
      #1
    \end{flushright}
    \rule{\linewidth}{1mm}
    \vspace*{\stretch{2}}
  }
}
%    \end{macrocode}
% Command to turn on separate title pages for each session. 
%    \begin{macrocode}
\newcommand{\SessionTitlePages}{\@sessiontitlepagestrue}
\newcommand{\BoldSessionTitles}{%
  \renewcommand{\rmaa@TOCPartStyle}[1]{\textbf{\rmaa@TOCTitleStyle{##1}}}
  }
\newcommand{\ArticleDir}[1]{\renewcommand{\@articledir}{#1/}\graphicspath{{1/}}}
\newcommand{\@articledir}{.}
%\newcommand{\ps@rmaafront}{\ps@rmaa}% page style for front matter (same as ever, actually)
\newcommand{\ps@rmaafront}{% page style for front matter:  page # in each foot
  \renewcommand\@oddfoot{\hfil{\thepage}}%
  \renewcommand\@evenfoot{{\thepage}\hfil}%
  % shorttitle in both heads
  \renewcommand\@evenhead{\hfil\@shorttitle\hfil}%
  \renewcommand\@oddhead{\@evenhead}
  }
\newcommand{\ps@rmaafrontfirstpage}{% page style for first page of front matter
  % page # in each foot, on outside edge
  \renewcommand\@oddfoot{\hfil\thepage}%
  \renewcommand\@evenfoot{\thepage\hfil}%
  % nothing in the even heads
  \renewcommand\@evenhead{\relax}%
  \renewcommand\@oddhead{\relax}%
  }
%% slight modification to the standard \frontmatter and \mainmatter
%% commands, just to set the \pagestyle
\renewcommand\frontmatter{%
  \latexonecolumn
  \cleardoublepage
  \@frontmattertrue
  \@mainmatterfalse
  \pagenumbering{roman}
  \pagestyle{rmaafront}
  \addtocounter{page}{2} % leave space for the title page
  }
\renewcommand\mainmatter{%
  \pagestyle{empty}
  \if@rmtwocol\twocolumn\fi
  \cleardoublepage
  \@frontmatterfalse
  \@mainmattertrue
  \pagenumbering{arabic}
  \pagestyle{rmaa}
  }
\renewcommand\backmatter{%
  \clearpage\pagestyle{empty}\cleardoublepage
  \@mainmatterfalse
  %% put extra space in the TOC
  \addtocontents{toc}{\bigskip}
  }
%    \end{macrocode}
%
% \subsection{Front sections}\label{sec:front-sections}
% \begin{macro}{\notitle}
% \begin{macro}{\shortname}
% \begin{macro}{frontsection}
% The environment \texttt{frontsection} is for doing a section in the
% front matter (preface, list of participants, or whatever). It has 
% one optional argument:
% \begin{enumerate}
% \item[\#1] Short title (no embedded spaces) for
% use in making PS files. E.g., ``list\_of\_participants''.   
% \end{enumerate}
% and two mandatory arguments: 
% \begin{enumerate}
% \item[\#2] Title of section (e.g., `Preface')
% \item[\#3] Commands to be evaluated before processing the contents
% of the environment. This is mainly useful for putting in the things
% that would go in the preamble of a normal article  (e.g.,
% |\listofauthors|, |\indexauthor|, |\notitle|) 
% \end{enumerate}
%    \begin{macrocode}
\newcommand{\notitle}{\@titlevisiblefalse} \newcommand{\shortname}{}
\newenvironment{frontsection}[3][]{%
  \begingroup
  \renewcommand{\shortname}{#1}
  \ifx\shortname\@rmaaEmpty
     \renewcommand{\articlename}{#2}
  \else
     \renewcommand{\articlename}{\shortname}
  \fi
  \@begofarticlehook
  \@titlevisibletrue % default
  \thispagestyle{rmaafrontfirstpage}
  \listofauthors{} % blank out this in case it is not specified
  \title{{#2}}\shorttitle{#2}\shortauthor{#2}
  #3 % evaluate the command given as second mandatory arg
  \addcontentsline{toc}{chapter}{\protect\@fiddleauthors{\@TOCtitle}{\rmaaTOCauthors}}%
  \addtocontents{htc}{\protect\HyperTOCEntry{\@rawtitle}{\@rawauthors}{\articlename}}
  \if@titlevisible
     \vspace*{1cm}\centerline{\rmaa@FrontSecStyle{\@title}}\vspace*{1cm}
  \fi
  \setlength{\parskip}{3ex}
  }{\frontlabel{RMAALastPage}\clearpage\endgroup}
%    \end{macrocode}
% \end{macro}
% \end{macro}
% \end{macro}
%    \begin{macrocode}
%% Definition of the command \frontlabel, which rights out something like `F4' to
%% the .aux file, instead of `\textlatin  {i}'. This is so it is easy to post-process
%% the file for making PS files of individual sections
\newcommand{\frontlabel}[1]{%
  \@bsphack
  \protected@write\@auxout{}%
     {\string\newlabel{\articlename:#1}{{}{F\@arabic{\c@page}}}}
  \@esphack
  }
%% User command to say X, Y, Z,\\ Editors of the Proceedings
\newcommand{\EditorSignature}{\vspace*{\parskip}\rightline{\@editors}%
  \rightline{Editors of the proceedings}}
%% Fiddle with \thesection, etc. so that they don't include the chapter number
%% Basically go back to the `article' definitions of these
\renewcommand{\thesection}{\@arabic\c@section}
\renewcommand{\thefigure}{\@arabic\c@figure}
\renewcommand{\thetable}{\@arabic\c@table}

\newcommand{\@oldargs}{}
\newcommand{\@newargs}{}
%% add TOC entry for each paper
%\renewcommand{\title}[1]%
%{\renewcommand{\@title}{\rmaa@TitleStyle{#1}}%
%\addtocounter{chapter}{1}\addcontentsline{toc}{chapter}{#1}%
%}

%% Formatting of the table of contents
%% dimensions of various parts
\newlength{\TOCtextwidth}\setlength{\TOCtextwidth}{12.65cm}
\newlength{\TOCpagenumwidth}\setlength{\TOCpagenumwidth}{1.3cm}
\newlength{\TOCentrywidth}\setlength{\TOCentrywidth}{\TOCtextwidth}
\addtolength{\TOCentrywidth}{-\TOCpagenumwidth} 
\newlength{\TOCextramargin}\setlength{\TOCextramargin}{1.25cm}
\newcommand{\TOCbox}[1]{\hspace{\TOCextramargin}\parbox{\TOCtextwidth}{#1}\par}
\newlength{\TOCextramarginx}\setlength{\TOCextramarginx}{\TOCextramargin}\addtolength{\TOCextramarginx}{4.3cm}
\newlength{\TOCtextwidthx}\setlength{\TOCtextwidthx}{\TOCtextwidth}\addtolength{\TOCtextwidthx}{-4.3cm}
\newcommand{\TOCboxx}[1]{\hspace{\TOCextramarginx}\parbox{\TOCtextwidthx}{#1}\par}

\newcommand{\RMAAbannerimage}{rmsc}
\renewcommand\tableofcontents{%
  \renewcommand{\articlename}{CONTENTS}
  \begingroup
  \@begofarticlehook
  \if@twocolumn
     \@restonecoltrue\onecolumn
  \else
     \@restonecolfalse
  \fi
  %% 9pt for journal style
  \if@proceedings\else\small\fi
  \if@newstyle\small\fi %% same for new style
  \if@proceedings\thispagestyle{rmaafrontfirstpage}\else\thispagestyle{empty}\fi
  \shortauthor{Contents}\shorttitle{Contents}
  \vspace*{18mm}
  \if@printbanner
  %% put in the banner header (need to get better image of this)
     \TOCbox{\hspace*{-0.15cm}\includegraphics[width=0.9\TOCtextwidth]{\RMAAbannerimage}}
     \vspace*{8mm}
  %% volume and publication date (should really be in funny font)
     \TOCboxx{\@volumen\\[2mm] \@mes\ \@anyo}
  \else
  %% put in the banner header (need to get better image of this)
     \TOCbox{\hspace*{-0.15cm}\includegraphics[width=0.9\TOCtextwidth]{\RMAAbannerimage}}
     \vspace*{12.5mm}
  %% volume and publication date (should really be in funny font)
     \TOCboxx{~\\[2mm] ~}
  \fi
  \vspace*{16mm}
  %% This is the command that fiddles with the TOC entry so that the
  %% title is flushleft but the authors are flushright
  \renewcommand{\@fiddleauthors}[2]{%
    % first put the paper title
    \renewcommand{\\}{\relax} % temporarily turn off explicit line breaks just in case
    \renewcommand{\altaffilmark}[1]{\relax} % ditto with this
    \renewcommand{\textsuperscript}[1]{\relax} % ditto with this 
    ##1
    % now save the author list in a box
    \sbox\@tempboxa{##2}%
    % check that box is shorter than the total width available for the entry
    \ifdim \wd\@tempboxa <\TOCentrywidth
       % if so, then put it right here, flush against right margin
       \penalty10000\hskip1em plus10em\penalty0\hbox{}\nobreak\hfill\box\@tempboxa
       % Note that we force a space of at least 1em between title and authorlist, 
       % stretchable to 11em since we may have very short titles but still with an 
       % authorlist too long to go on the same line.  The only linebreak allowed is 
       % directly before the authorlist.
    \else
       % if it is too long, we set the original authorlist text in its own parbox
       \par\vskip1ex\parbox[b]{\TOCentrywidth}{\raggedleft ##2}
    \fi
    }
  \@starttoc{toc}%
  \if@restonecol\twocolumn\fi
  \frontlabel{RMAALastPage}
  \clearpage
  \thispagestyle{empty}
  \cleardoublepage % make sure we start on an odd page
  \endgroup

  %%% ALSO, do the HTML TOC
  %% The file that will be used to make the HTML TOC
  \@starttoc{htc}
  %% All the stuff in this file does nothing in the book style
  \addtocontents{htc}{\protect\HyperTOCmeeting{\@conftitle}{\@reunionplace}{\@reuniondate}}
  \addtocontents{htc}{\protect\HyperTOCeditors{\@editors}}
  \addtocontents{htc}{\protect\HyperTOCpreamble}
  \addtocontents{htc}{\protect\HyperTOCEntry{Table of Contents}{}{\articlename}}

  %% and the ADS file
  \ADS@openfile
  }
%% This is now much simpler
\renewcommand\l@chapter[2]{%TOC entry for articles
  \ifnum \c@tocdepth >\m@ne
    \vskip 1.0em \@plus\p@
    \TOCbox{%
      \parbox[b]{\TOCentrywidth}{#1}\parbox[b]{\TOCpagenumwidth}{\hfill#2}
    }
  \fi} 
\renewcommand\l@part[2]{%TOC entry for sessions: centered, no page number
  \ifnum \c@tocdepth >\m@ne
    \vskip 1.5em \@plus1.5em \goodbreak
    \TOCbox{%
      \parbox[b]{\TOCentrywidth}{\hfil#1\hfil}\hfill
    }
  \fi\nobreak} 

%% Formatting theindex
%% strongly encourage page breaks between index entries
\renewcommand{\@idxitem}{\par\penalty-1000\hangindent.45cm} 
\renewcommand{\indexspace}{\par\penalty-1000\vskip 10\p@ \@plus15\p@ \@minus3\p@\relax}
\renewenvironment{theindex}{%
  %% first, make sure we start on an odd page (06 Jun 2000)
  \clearpage\thispagestyle{empty}\cleardoublepage
  %% this looks a bit like an article but is two columns and without abstract, etc
  \renewcommand{\articlename}{INDEX}
  \title{Author Index}\shortauthor{Author Index}\shorttitle{Author Index}
  \listofauthors{}
  \addcontentsline{toc}{chapter}{\protect\@fiddleauthors{\@TOCtitle}{\rmaaTOCauthors}}%
  \addtocontents{htc}{\protect\HyperTOCEntry{\@rawtitle}{\@rawauthors}{\articlename}}
  \pagestyle{rmaa}
  \thispagestyle{rmaafirstpage}
  \@begofarticlehook
  %% 9pt for journal style
  \if@proceedings\else\small\fi
  \if@newstyle\small\fi
  \twocolumn[\@makeminititle]
  %% vertical spacing (see also definition of \@idxitem and \indexspace
  \parskip\z@ \@plus .6\p@\relax % allow quite a lot of leaway here to avoid breaks inside entries
  %% horizontal spacing, breaking and glue parameters (see also the stuff in the definition 
  %% of \indexauthor
  \parindent\z@ 
  \exhyphenpenalty10000 % don't break on explicit hyphens, it always looks bad
  \rightskip=0pt plus.2em % allow a tiny bit of leeway at the ends of lines
  \spaceskip.3333em plus.5em minus.15em % a lot of leeway between words
  \xspaceskip.5em plus.5em minus.4em % this seems to have no effect
  \let\item\@idxitem
  \renewcommand{\\}{\relax} % temporarily turn off explicit line breaks just in case
  \renewcommand{\altaffilmark}[1]{\relax} % ditto with this
  \renewcommand{\textsuperscript}[1]{\relax} % ditto with this
  }{\label{RMAALastPage}\clearpage} % clearpage necessary so that we know about \shortauthor



%% end definitions for book styles
\fi

%    \end{macrocode}
% In order to do away with plates, we sometimes want a figure to
% appear in color only in the online version. The following command
% can bve used in a figure caption to typeset a note to this effect. 
%    \begin{macrocode}
\newcommand\ColorOnlineOnly{%
  \textsc{Note: this figure is available in color in the electronic
    version of this article, obtainable from}
  \texttt{http://\discretionary{}{}{}www.\discretionary{}{}{}astroscu.\discretionary{}{}{}unam.\discretionary{}{}{}mx/\discretionary{}{}{}\~{}rmaa/}.}
%% Plates
%% Define 2 counters: plate is the number in the whole document, plateloop is used 
%% internally by \outputplates
\newcounter{plate}
\newcounter{plateloop}
%% Currently limited to 15 plates in document
%% when we do a label, we need to write to \newlabel entries to the .aux file, one 
%% that says RMAAPlateX so that it can be referred to in \ouputplates, and one that
%% says articlename:RMAAPlateX so that the dvips program can find it and asociate it
%% with the right article
\newcommand{\@platelabeli}{\oldlabel{RMAAPlate1}\label{RMAAPlate1}}
\newcommand{\@platelabelii}{\oldlabel{RMAAPlate2}\label{RMAAPlate2}}
\newcommand{\@platelabeliii}{\oldlabel{RMAAPlate3}\label{RMAAPlate3}}
\newcommand{\@platelabeliv}{\oldlabel{RMAAPlate4}\label{RMAAPlate4}}
\newcommand{\@platelabelv}{\oldlabel{RMAAPlate5}\label{RMAAPlate5}}
\newcommand{\@platelabelvi}{\oldlabel{RMAAPlate6}\label{RMAAPlate6}}
\newcommand{\@platelabelvii}{\oldlabel{RMAAPlate7}\label{RMAAPlate7}}
\newcommand{\@platelabelviii}{\oldlabel{RMAAPlate8}\label{RMAAPlate8}}
\newcommand{\@platelabelix}{\oldlabel{RMAAPlate9}\label{RMAAPlate9}}
\newcommand{\@platelabelx}{\oldlabel{RMAAPlate10}\label{RMAAPlate10}}
\newcommand{\@platelabelxi}{\oldlabel{RMAAPlate11}\label{RMAAPlate11}}
\newcommand{\@platelabelxii}{\oldlabel{RMAAPlate12}\label{RMAAPlate12}}
\newcommand{\@platelabelxiii}{\oldlabel{RMAAPlate13}\label{RMAAPlate13}}
\newcommand{\@platelabelxiv}{\oldlabel{RMAAPlate14}\label{RMAAPlate14}}
\newcommand{\@platelabelxv}{\oldlabel{RMAAPlate15}\label{RMAAPlate15}}
\newcommand{\@outputplatelabeli}{\oldlabel{RMAAOutputPlate1}}
\newcommand{\@outputplatelabelii}{\oldlabel{RMAAOutputPlate2}}
\newcommand{\@outputplatelabeliii}{\oldlabel{RMAAOutputPlate3}}
\newcommand{\@outputplatelabeliv}{\oldlabel{RMAAOutputPlate4}}
\newcommand{\@outputplatelabelv}{\oldlabel{RMAAOutputPlate5}}
\newcommand{\@outputplatelabelvi}{\oldlabel{RMAAOutputPlate6}}
\newcommand{\@outputplatelabelvii}{\oldlabel{RMAAOutputPlate7}}
\newcommand{\@outputplatelabelviii}{\oldlabel{RMAAOutputPlate8}}
\newcommand{\@outputplatelabelix}{\oldlabel{RMAAOutputPlate9}}
\newcommand{\@outputplatelabelx}{\oldlabel{RMAAOutputPlate10}}
\newcommand{\@outputplatelabelxi}{\oldlabel{RMAAOutputPlate11}}
\newcommand{\@outputplatelabelxii}{\oldlabel{RMAAOutputPlate12}}
\newcommand{\@outputplatelabelxiii}{\oldlabel{RMAAOutputPlate13}}
\newcommand{\@outputplatelabelxiv}{\oldlabel{RMAAOutputPlate14}}
\newcommand{\@outputplatelabelxv}{\oldlabel{RMAAOutputPlate15}}
\newcommand{\@platerefi}{\oldpageref{RMAAPlate1}}
\newcommand{\@platerefii}{\oldpageref{RMAAPlate2}}
\newcommand{\@platerefiii}{\oldpageref{RMAAPlate3}}
\newcommand{\@platerefiv}{\oldpageref{RMAAPlate4}}
\newcommand{\@platerefv}{\oldpageref{RMAAPlate5}}
\newcommand{\@platerefvi}{\oldpageref{RMAAPlate6}}
\newcommand{\@platerefvii}{\oldpageref{RMAAPlate7}}
\newcommand{\@platerefviii}{\oldpageref{RMAAPlate8}}
\newcommand{\@platerefix}{\oldpageref{RMAAPlate9}}
\newcommand{\@platerefx}{\oldpageref{RMAAPlate10}}
\newcommand{\@platerefxi}{\oldpageref{RMAAPlate11}}
\newcommand{\@platerefxii}{\oldpageref{RMAAPlate12}}
\newcommand{\@platerefxiii}{\oldpageref{RMAAPlate13}}
\newcommand{\@platerefxiv}{\oldpageref{RMAAPlate14}}
\newcommand{\@platerefxv}{\oldpageref{RMAAPlate15}}
\newcounter{figplatei}
\newcounter{figplateii}
\newcounter{figplateiii}
\newcounter{figplateiv}
\newcounter{figplatev}
\newcounter{figplatevi}
\newcounter{figplatevii}
\newcounter{figplateviii}
\newcounter{figplateix}
\newcounter{figplatex}
\newcounter{figplatexi}
\newcounter{figplatexii}
\newcounter{figplatexiii}
\newcounter{figplatexiv}
\newcounter{figplatexv}
%% Printed output should use the counter `plate'
\renewcommand{\theplate}{\@arabic\c@plate}
\renewcommand{\theplateloop}{\@arabic\c@plateloop}
%% Command to output all the saved plates in this paper
\newcommand{\outputplates}{%
  \setcounter{plateloop}{0}
  \if@book\else\@endofarticlehook\fi % not entirely sure why, but this
                                     % seems to be needed
  \clearpage % first output the last page of text
  % now set pagestyle for blank even pages and make sure we start on an odd page
  \pagestyle{empty}
  \cleardoublepage
  \onecolumn
  \parindent\z@
  % first put in an extra sheet
  % the main purpose of this is so there is always a blank even page before 
  % each plate for filling in the PS file if the article has an odd number
  % of pages (see splitbook.sh)
  \null\vfil\centerline{PLATES}\vfil\clearpage\pagestyle{rmaaplate}\cleardoublepage
  % output each plate in turn
  \@whilenum\value{plateloop}<\value{plate}\do{%
    \addtocounter{plateloop}{1}
    % restore the figure counter from \@figcounti, \@figcountii, etc
    \setcounter{figure}{\value{figplate\roman{plateloop}}}
    % spew out the plate (in its own group for safety)
    {\csname @plate\roman{plateloop}\endcsname\par}
    }
  \pagestyle{rmaa}
  % reset the page number as if nothing had happened
  \addtocounter{page}{-\value{plate}}
  \addtocounter{page}{-\value{plate}}
  }
\newcommand\figcaption{%
  \def\@captype{figure}
  \renewcommand{\@makecaption}{\@makefigcap}
  \caption}
\newcommand\tabcaption{%
  \def\@captype{table}
  \renewcommand{\@makecaption}{\@maketabcap}
  \caption}

\newcommand{\plate}[2][1]{%
  %% User command to make a plate page: 
  %% Mandatory argument is what goes on the plate - this should not 
  %% include any floating environments (e.g. figure). Instead, \figcaption 
  %% should be used to do the captions. 
  %% Optional argument is the number of figures (i.e. captions) on the plate
  %% page (default is 1). This is necessary for incrementing the figure counter. 
  %%
  %% Note that we have to use a \gdef rather than an \xdef when we 
  %% save the plate output in a command. This means that in the book styles when it 
  %% comes to spewing out the plate, we will have forgotten about any user macros 
  %% that may have been defined. This can be got round by repeating the macros inside the plate.
  %% (Presumably, one would use a \newcommand rather than \renewcommand if they were new.)
  % increment global and local counters
  \addtocounter{plate}{1}
  % write appropriate label to .aux file
  \csname @platelabel\roman{plate}\endcsname
  % define a global macro called \@platei, \@plateii, \@plateiii, \@plateiv, ...
  \expandafter\gdef\csname @plate\roman{plate}\endcsname{%
    \csname @outputplatelabel\roman{plateloop}\endcsname
    \null\vfil #2 \vfil
    \cleardoublepage
    }
  % save the current shortauthor and shorttitle
  % need to be \xdef to take current values of \@shortauthor, etc 
  \expandafter\xdef\csname @shorta\roman{plate}\endcsname{\@shortauthor}
  \expandafter\xdef\csname @shortt\roman{plate}\endcsname{\@shorttitle}
  % save the current value of the figure counter for later use in \outputplates
  \setcounter{figplate\roman{plate}}{\value{figure}}
  % increment the figure counter according to the optional argument
  \addtocounter{figure}{#1}
  }
%% page style for plate pages 
\newcommand{\ps@rmaaplate}{%
  % Plate X (see page Y) in the foot
  \renewcommand\@oddfoot{\hfil\csname @shorta\roman{plateloop}\endcsname\ 
    (see page~\csname @plateref\roman{plateloop}\endcsname)\hfil}%
  \renewcommand\@evenfoot{\relax}%
  % nothing in even heads (blank pages)
  \renewcommand\@evenhead{\relax}
  % short title and page num in the odd heads
  \renewcommand\@oddhead{\hfil\csname @shortt\roman{plateloop}\endcsname%
    \hfil\llap{PLATE~\theplateloop}}
  }





%% Things for the User Guide
\if@guide
%% stolen from doc.sty
\def\MakeShortVerb#1{%
  \expandafter\ifx\csname cc\string#1\endcsname\relax
  \@shortvrbinfo{Made }{#1}%
  \add@special{#1}%
  \expandafter
  \xdef\csname cc\string#1\endcsname{\the\catcode`#1}%
  \begingroup
  \catcode`\~\active  \lccode`\~`#1%
  \lowercase{%
    \global\expandafter\let
    \csname ac\string#1\endcsname~%
    \gdef~{\verb~}}%
  \endgroup
  \global\catcode`#1\active
  \else
  \@shortvrbinfo\@empty{#1 already}%
  \fi}
\def\DeleteShortVerb#1{%
  \expandafter\ifx\csname cc\string#1\endcsname\relax
  \else
  \@shortvrbinfo{Deleted }{#1 as}%
  \rem@special{#1}%
  \global\catcode`#1\csname cc\string#1\endcsname
  \global \expandafter\let \csname cc\string#1\endcsname \relax
  \ifnum\catcode`#1=\active
  \begingroup
  \catcode`\~\active   \lccode`\~`#1%
  \lowercase{%
    \global\expandafter\let\expandafter~%
    \csname ac\string#1\endcsname}%
  \endgroup \fi \fi}
\def\@shortvrbinfo#1#2{%
  \PackageInfo{doc}{%
    #1\expandafter\@gobble\string#2 a short reference
    for \string\verb}}
\def\add@special#1{%
  \rem@special{#1}%
  \expandafter\gdef\expandafter\dospecials\expandafter
  {\dospecials \do #1}%
  \expandafter\gdef\expandafter\@sanitize\expandafter
  {\@sanitize \@makeother #1}}
\def\rem@special#1{%
  \def\do##1{%
    \ifnum`#1=`##1 \else \noexpand\do\noexpand##1\fi}%
  \xdef\dospecials{\dospecials}%
  \begingroup
  \def\@makeother##1{%
    \ifnum`#1=`##1 \else \noexpand\@makeother\noexpand##1\fi}%
  \xdef\@sanitize{\@sanitize}%
  \endgroup}
%% define | as shortcut for \verbatim
\MakeShortVerb{\|}
%% comand to specify arguments to LaTeX commands
\newcommand{\rmArg}[1]{%
  \ensuremath{\langle}{\scriptsize \textsf{\MakeUpperCase{#1}}}\ensuremath{\rangle}}
%\newcommand{\rmArg}[1]{\ensuremath{\langle}\textsc{#1}\ensuremath{\rangle}}
\providecommand{\BibTeX}{{\rmfamily B\kern-.05em%
    \textsc{i\kern-.025em b}\kern-.08em%
    T\kern-.1667em\lower.7ex\hbox{E}\kern-.125emX}}
%% Environment for code examples
\newenvironment{Code}{%
  \begin{list}{}{\setlength{\leftmargin}{0.05\columnwidth}} 
  \item[]}{\end{list}}
\fi
%</package>
%    \end{macrocode}

% \section{User commands for astronomical articles}
% \label{sec:user-commands}
% Some random definitions. These go both into the main class file and
% also into |adshtml.sty|.
% 
% Definitions of the journal name in different incarnations.  
%    \begin{macrocode}
%<*package|adshtml-style>
\newcommand{\RMAA}{RMxAA}
\newcommand{\RMAAlong}{Revista Mexicana de Astronom\'{\i}a y Astrof\'{\i}sica}
\newcommand{\RMAASC}{RMxAC}
\newcommand{\RMAASClong}{\RMAA\ (Serie de Conferencias)}
%    \end{macrocode}
% \begin{macro}{\ibidrule}
% This puts a horizontal rule with a period after it. It is used when
% the author list of a reference is identical to the preceding one. 
%    \begin{macrocode}
\newcommand{\ibidrule}{\rule[-0.03cm]{4.7em}{.01cm}.\ }
%    \end{macrocode}
% \end{macro}
% \begin{macro}{\ThisVolume}
% References to other articles in the same volume can be coded as,
% e.g., |\bibitem{} Bloggs, J. \ThisVolume{bloggs}|.
%    \begin{macrocode}
\newcommand{\ThisVolume}[1]{%
  \@year, \if@proceedings\RMAASC\else\RMAA\fi, \@volume,
  \if@book\oldpageref{#1:RMAAFirstPage}\else 000\fi\ (this volume)}
%    \end{macrocode}
% \end{macro}
% Other stuff that may be useful in ref lists.
%    \begin{macrocode}
\newcommand{\etal}{et al.\@}
%    \end{macrocode}
% Various new commands from AAS macros
%    \begin{macrocode}
\def\acknowledgments{\vskip 3ex plus .8ex minus .4ex}
\let\acknowledgements=\acknowledgments                  % second spelling
\let\jnl@style=\rm
\def\ref@jnl#1{{\jnl@style#1}}
\DeclareRobustCommand\aj{\ref@jnl{AJ}}                   
\DeclareRobustCommand\araa{\ref@jnl{ARA\&A}}             
\DeclareRobustCommand\apj{\ref@jnl{ApJ}}                 
\DeclareRobustCommand\apjl{\ref@jnl{ApJ}}                
\DeclareRobustCommand\apjs{\ref@jnl{ApJS}}               
\DeclareRobustCommand\ao{\ref@jnl{Appl.~Opt.}}           
\DeclareRobustCommand\apss{\ref@jnl{Ap\&SS}}             
\DeclareRobustCommand\aap{\ref@jnl{A\&A}}                
\DeclareRobustCommand\aapr{\ref@jnl{A\&A~Rev.}}          
\DeclareRobustCommand\aaps{\ref@jnl{A\&AS}}              
\DeclareRobustCommand\azh{\ref@jnl{AZh}}                 
\DeclareRobustCommand\baas{\ref@jnl{BAAS}}               
\DeclareRobustCommand\jrasc{\ref@jnl{JRASC}}             
\DeclareRobustCommand\memras{\ref@jnl{MmRAS}}            
\DeclareRobustCommand\mnras{\ref@jnl{MNRAS}}             
\DeclareRobustCommand\pra{\ref@jnl{Phys.~Rev.~A}}        
\DeclareRobustCommand\prb{\ref@jnl{Phys.~Rev.~B}}        
\DeclareRobustCommand\prc{\ref@jnl{Phys.~Rev.~C}}        
\DeclareRobustCommand\prd{\ref@jnl{Phys.~Rev.~D}}        
\DeclareRobustCommand\pre{\ref@jnl{Phys.~Rev.~E}}        
\DeclareRobustCommand\prl{\ref@jnl{Phys.~Rev.~Lett.}}    
\DeclareRobustCommand\pasp{\ref@jnl{PASP}}               
\DeclareRobustCommand\pasj{\ref@jnl{PASJ}}               
\DeclareRobustCommand\qjras{\ref@jnl{QJRAS}}             
\DeclareRobustCommand\skytel{\ref@jnl{S\&T}}             
\DeclareRobustCommand\solphys{\ref@jnl{Sol.~Phys.}}      
\DeclareRobustCommand\sovast{\ref@jnl{Soviet~Ast.}}      
\DeclareRobustCommand\ssr{\ref@jnl{Space~Sci.~Rev.}}     
\DeclareRobustCommand\zap{\ref@jnl{ZAp}}                 
\DeclareRobustCommand\nat{\ref@jnl{Nature}}              
\DeclareRobustCommand\iaucirc{\ref@jnl{IAU~Circ.}}
\DeclareRobustCommand\aplett{\ref@jnl{Astrophys.~Lett.}}
\DeclareRobustCommand\apspr{\ref@jnl{Astrophys.~Space~Phys.~Res.}}
\DeclareRobustCommand\bain{\ref@jnl{Bull.~Astron.~Inst.~Netherlands}}
\DeclareRobustCommand\fcp{\ref@jnl{Fund.~Cosmic~Phys.}}
\DeclareRobustCommand\gca{\ref@jnl{Geochim.~Cosmochim.~Acta}}
\DeclareRobustCommand\grl{\ref@jnl{Geophys.~Res.~Lett.}}
\DeclareRobustCommand\jcp{\ref@jnl{J.~Chem.~Phys.}}      
\DeclareRobustCommand\jgr{\ref@jnl{J.~Geophys.~Res.}}    
\DeclareRobustCommand\jqsrt{\ref@jnl{J.~Quant.~Spec.~Radiat.~Transf.}}
\DeclareRobustCommand\memsai{\ref@jnl{Mem.~Soc.~Astron.~Italiana}}
\DeclareRobustCommand\nphysa{\ref@jnl{Nucl.~Phys.~A}}
\DeclareRobustCommand\physrep{\ref@jnl{Phys.~Rep.}}
\DeclareRobustCommand\physscr{\ref@jnl{Phys.~Scr}}
\DeclareRobustCommand\planss{\ref@jnl{Planet.~Space~Sci.}}
\DeclareRobustCommand\procspie{\ref@jnl{Proc.~SPIE}}
\let\astap=\aap
\let\apjlett=\apjl
\let\apjsupp=\apjs
\let\applopt=\ao
\def\phn{\phantom{0}}
\def\phd{\phantom{.}}
\def\phs{\phantom{\ensuremath{-}}}
\def\phm#1{\phantom{#1}}
%    \end{macrocode}
% All these are redone as robust and using more modern commands where
% possible. The robustness is so they will not get expanded if used in
% abstracts and titles when these are passed to the ADS file. 
% \changes{3.16}{2003/04/07}{
%  Added \cs{Color} command for doing photometric colors.}
%    \begin{macrocode}
%%\def\sun{\hbox{$\odot$}}
\DeclareRobustCommand\sun{\ensuremath{\odot}}
%%\def\earth{\hbox{$\oplus$}}
\DeclareRobustCommand\earth{\ensuremath{\oplus}}
%%\def\lesssim{\mathrel{\hbox{\rlap{\hbox{\lower4pt\hbox{$\sim$}}}\hbox{$<$}}}}
%%\def\gtrsim{\mathrel{\hbox{\rlap{\hbox{\lower4pt\hbox{$\sim$}}}\hbox{$>$}}}}
\DeclareRobustCommand\lesssim{\mathrel{\hbox{\smaller\smaller\rlap{\hbox{\lower3pt\hbox{$\sim$}}}\raise1.5pt\hbox{$<$}}}}
\DeclareRobustCommand\gtrsim{\mathrel{\hbox{\smaller\smaller\rlap{\hbox{\lower3pt\hbox{$\sim$}}}\raise1.5pt\hbox{$>$}}}}
%%\def\lta{\mathrel{\spose{\lower 3pt\hbox{$\mathchar"218$}}
%%     \raise 2.0pt\hbox{$\mathchar"13C$}}}
%%\def\gta{\mathrel{\spose{\lower 3pt\hbox{$\mathchar"218$}}
%%     \raise 2.0pt\hbox{$\mathchar"13E$}}}
\DeclareRobustCommand\sq{\hbox{\rlap{$\sqcap$}$\sqcup$}}
%%\def\arcdeg{\hbox{$^\circ$}}
%%\def\arcmin{\hbox{$^\prime$}}
%%\def\arcsec{\hbox{$^{\prime\prime}$}}
\DeclareRobustCommand\arcdeg{\ensuremath{^\circ}}
\DeclareRobustCommand\arcmin{\ensuremath{^\prime}}
\DeclareRobustCommand\arcsec{\ensuremath{^{\prime\prime}}}
%%\def\fd{\hbox{$.\!\!^{\rm d}$}}
%%\def\fh{\hbox{$.\!\!^{\rm h}$}}
%%\def\fm{\hbox{$.\!\!^{\rm m}$}}
%%\def\fs{\hbox{$.\!\!^{\rm s}$}}
%%\def\fdg{\hbox{$.\!\!^\circ$}}
%%\def\farcm{\hbox{$.\mkern-4mu^\prime$}}
%%\def\farcs{\hbox{$.\!\!^{\prime\prime}$}}
%%\def\fp{\hbox{$.\!\!^{\scriptscriptstyle\rm p}$}}
\DeclareRobustCommand\fd{\ensuremath{.\!\!^\mathrm{d}}}
\DeclareRobustCommand\fh{\ensuremath{.\!\!^\mathrm{h}}}
\DeclareRobustCommand\fm{\ensuremath{.\!\!^\mathrm{m}}}
\DeclareRobustCommand\fs{\ensuremath{\mkern0.5mu.\mkern-5mu^\mathrm{s}\mkern-1mu}}
\DeclareRobustCommand\fdg{\ensuremath{.\!\!^\circ}}
\DeclareRobustCommand\farcm{\ensuremath{.\mkern-4mu^\prime}}
\DeclareRobustCommand\farcs{\ensuremath{\mkern1mu.\mkern-6mu^{\prime\prime}\mkern-1.5mu}}
%%\DeclareRobustCommand\farcs{\ensuremath{.\!\!^{\prime\prime}}}
\DeclareRobustCommand\fp{\ensuremath{.\!\!^{\scriptscriptstyle\rm p}}}
%%\def\micron{\hbox{$\mu$m}}
\DeclareRobustCommand\micron{\ensuremath{\mu\mathrm{m}}}
\let\la=\lesssim                        
\let\ga=\gtrsim
%%\def\case#1#2{\hbox{$\frac{#1}{#2}$}}
%%\def\slantfrac#1#2{\hbox{$\,^#1\!/_#2$}}
%%\def\onehalf{\slantfrac{1}{2}}
%%\def\onethird{\slantfrac{1}{3}}
%%\def\twothirds{\slantfrac{2}{3}}
%%\def\onequarter{\slantfrac{1}{4}}
%%\def\threequarters{\slantfrac{3}{4}}
\DeclareRobustCommand\case[2]{\ensuremath{\frac{#1}{#2}}}
\DeclareRobustCommand\slantfrac[2]{\ensuremath{\,^#1\!/_#2}}
\DeclareRobustCommand\onehalf{\slantfrac{1}{2}}
\DeclareRobustCommand\onethird{\slantfrac{1}{3}}
\DeclareRobustCommand\twothirds{\slantfrac{2}{3}}
\DeclareRobustCommand\onequarter{\slantfrac{1}{4}}
\DeclareRobustCommand\threequarters{\slantfrac{3}{4}}
%%\def\ubvr{\hbox{$U\!BV\!R$}}            
%%\def\ub{\hbox{$U\!-\!B$}}               
%%\def\bv{\hbox{$B\!-\!V$}}               
%%\def\vr{\hbox{$V\!-\!R$}}               
%%\def\ur{\hbox{$U\!-\!R$}}
%    \end{macrocode}
% New commands for ``photometric'' colors 
%    \begin{macrocode}
\DeclareRobustCommand\ubvr{\ensuremath{U\!BV\!R}}            
\DeclareRobustCommand\Color[2]{%
  \ensuremath{%
    #1\!\ifx#1V\!\fi-\!#2
  }}
\DeclareRobustCommand\ub{\Color{U}{B}}
\DeclareRobustCommand\bv{\Color{B}{V}}
\DeclareRobustCommand\vr{\Color{V}{R}}
\DeclareRobustCommand\ur{\Color{U}{R}}
%% New version of the \ion command, that actually works! 
%    \end{macrocode}
% \begin{macro}{\ion}
% This version of the |\ion| command allows you to write either
% |\ion{N}{4}| or |\ion{N}{iv}| in order to get N~IV. 
%    \begin{macrocode}
\DeclareRobustCommand\ion[2]{%
  \mbox{#1\kern0.2em%
    \smaller\rmfamily%
    \edef\@tempa{\@car#2\@nil}%
    \ifcat1\@tempa%
    \@Roman{#2}%
    \else%
    \uppercase{#2}%
    \fi}}
\newcounter{IonStage}
\DeclareRobustCommand\plainion[2]{\setcounter{IonStage}{#2}#1 \Roman{IonStage}}
%%\def\ion#1#2{\ensuremath{\mathrm{#1}\;}{\protect\small\rm\@Roman{#2}}}
%    \end{macrocode}
% \end{macro}
%    \begin{macrocode}
\DeclareRobustCommand\nodata{$\cdots$}
%    \end{macrocode}
% These are needed by the |\ion| command. Isn't there an easier way of
% doing this?
%    \begin{macrocode}
\let\@small=\small
\let\@normalsize=\normalsize
\let\@footnotesize=\footnotesize
\let\@scriptsize=\scriptsize
\let\@tiny=\tiny
\let\@large=\large
\let\@Large=\Large
\let\@LARGE=\LARGE
\let\@huge=\huge
\let\@Huge=\Huge
\newcommand{\smaller}{\small}
\renewcommand{\normalsize}{\renewcommand{\smaller}{\small}\@normalsize}
\renewcommand{\small}{\renewcommand{\smaller}{\footnotesize}\@small}
\renewcommand{\footnotesize}{\renewcommand{\smaller}{\scriptsize}\@footnotesize}
\renewcommand{\scriptsize}{\renewcommand{\smaller}{\tiny}\@scriptsize}
\renewcommand{\tiny}{\renewcommand{\smaller}{\tiny}\@tiny}
\renewcommand{\large}{\renewcommand{\smaller}{\normalsize}\@large}  
\renewcommand{\Large}{\renewcommand{\smaller}{\large}\@Large}  
\renewcommand{\LARGE}{\renewcommand{\smaller}{\Large}\@LARGE}  
\renewcommand{\huge}{\renewcommand{\smaller}{\LARGE}\@huge}  
\renewcommand{\Huge}{\renewcommand{\smaller}{\huge}\@Huge}  
%</package|adshtml-style>
%    \end{macrocode}


% \section{Other files}

% \subsection{The driver and make files}
% \label{sec:driver_make}
% These are put here for completeness, although they still need to be
% distributed separately since otherwise the user won't be able to
% unpack this file into its constituent parts.
% 
% The |Makefile| has two principal rules: |docs|, and |macros|. The
% |docs| rule produces a PDF and gzipped PS file of the documentation.
% I still haven't perfected the combination of rules required to get
% it to run \LaTeX{} enough times in all circumstances. The |macros|
% rule runs \LaTeX{} on |rmaa.ins| in order to unpack |rmaa.cls| and
% all the other little files (including the |Makefile| and the
% |rmaa.ins| files themselves\dots 
% 
% First set up lists of various types of file. 
%    \begin{macrocode}
%<*makefile>
DOCFILES=rmaa.pdf
MACRO_LTX_FILES=rmaa.cls rmaacite.sty adshtml.sty
SHELLSCRIPTS=makeads.sh
MACRO_OTHER_FILES=adshtml.perl sample.latex2html-init
MAKEFILE=Makefile
MACROFILES=$(MACRO_LTX_FILES) $(SHELLSCRIPTS) $(MACRO_OTHER_FILES)         
SKELETONFILES=ADS.tex 
%    \end{macrocode}
% Install the files read only into the package directory, specified via a path relative to the source directory. \changes{4.0}{2019/11/14}{Reorganize folder layout so that hardly anything is in the top level.  That way, things are clearer on github.}
%    \begin{macrocode}
INSTALL=install 
INSTALL_DIR=../packages
%    \end{macrocode}
% The default rule is to make everything. Installation must be
% explicitly asked for.
%    \begin{macrocode}
all: docs macros

install: install-docs install-macros

%    \end{macrocode}
% Most files are installed read-only to prevent editing. Skeleton
% files must be writeable. Must be careful not to overwrite the
% Makefile in the parent directory. 
%    \begin{macrocode}
install-macros: macros
        $(INSTALL) -m 444 $(MACRO_LTX_FILES) $(MACRO_OTHER_FILES) $(INSTALL_DIR)
        $(INSTALL) -m 555 $(SHELLSCRIPTS) $(INSTALL_DIR)
        $(INSTALL) -m 644 $(SKELETONFILES) $(INSTALL_DIR)

install-docs: docs
        $(INSTALL) rmaa.pdf $(INSTALL_DIR)/rm-fulldocs.pdf

%    \end{macrocode}
% Make the PDF and gzipped PS files. 
%    \begin{macrocode}
docs: $(DOCFILES)

%    \end{macrocode}
% Make everything else if either the |.dtx| or |.ins| files have
% changed. 
%    \begin{macrocode}
macros: rmaa.dtx rmaa.ins
%    \end{macrocode}
% Ensure that all the macro files can be written to then recreate
% them. The |touch| is there in case any of the files don't exist
% yet. 
%    \begin{macrocode}
        touch $(MACROFILES) $(MAKEFILE)
        chmod u+w $(MACROFILES) $(MAKEFILE) rmaa.ins
        latex "\input{rmaa.ins}"
%    \end{macrocode}
% \TeX{} automatically adds the |.tex| extension to extensionless
% files. We fix this and take the opportunity to convert
% back to TABs all that it has lovingly converted to spaces. 
%    \begin{macrocode}
        cp -f $(MAKEFILE) $(MAKEFILE).BCK && unexpand $(MAKEFILE).tex >| $(MAKEFILE) && rm $(MAKEFILE).tex
%    \end{macrocode}
% Make the shell scripts executable and everything unwriteable (so
% they don't get hand-edited by mistake). However, undo this for the
% files that the user is \emph{supposed} to modify. 
%    \begin{macrocode}
        chmod a+x $(SHELLSCRIPTS)
        chmod a-w $(MACROFILES) $(MAKEFILE)
        chmod u+w $(SKELETONFILES)

%    \end{macrocode}
% We now use PDF\LaTeX{} for this.
%    \begin{macrocode}
rmaa.idx: rmaa.dtx
        @echo "REGENERATING .idx FILE"
        pdflatex "\nonstopmode\input{rmaa.dtx}"

rmaa.aux: rmaa.dtx rmaa.ind rmaa.gls
        @echo "REGENERATING .aux FILE"
        pdflatex "\nonstopmode\input{rmaa.dtx}"

rmaa.ind: rmaa.idx
        @echo "REGENERATING INDEX"
        makeindex -s gind.ist rmaa

rmaa.gls: rmaa.glo
        @echo "REGENERATING CHANGELOG"
        makeindex -s gglo.ist -o rmaa.gls rmaa.glo

rmaa.pdf: rmaa.aux rmaa.ind rmaa.gls
        @echo "REGENERATING .pdf FILE"
        pdflatex "\nonstopmode\input{rmaa.dtx}"

clean: 
        rm -f rmaa.aux rmaa.dvi rmaa.ind rmaa.gls
%</makefile>
%    \end{macrocode}

% The |rmaa.ins| file uses the \textsf{docstrip} package to split out
% all the files and strip off this documentation.
%    \begin{macrocode}
%<*ins>
%%
%% Installation batch file for the RMAA LaTeX package
%%
\def\batchfile{rmaa.ins}
\input docstrip.tex
\preamble

PLEASE DO NOT EDIT THIS FILE! It is generated automatically from
the file rmaa.dtx using the accompanying Makefile and rmaa.ins 
Any hand-edited changes will be lost

A LaTeX document class for the RMxAA

Author: William Henney (w.henney@irya.unam.mx)

Latest version available from: http://www.irya.unam.mx/rmaa

Copyright (C) 1999-2019 Instituto de Astronomia, UNAM, Mexico

Permission is hereby granted, free of charge, to any person obtaining
a copy of this software and associated documentation files (the
"Software"), to deal in the Software without restriction, including
without limitation the rights to use, copy, modify, merge, publish,
distribute, sublicense, and/or sell copies of the Software, and to
permit persons to whom the Software is furnished to do so, subject to
the following conditions:

The above copyright notice and this permission notice shall be
included in all copies or substantial portions of the Software.

THE SOFTWARE IS PROVIDED "AS IS", WITHOUT WARRANTY OF ANY KIND,
EXPRESS OR IMPLIED, INCLUDING BUT NOT LIMITED TO THE WARRANTIES OF
MERCHANTABILITY, FITNESS FOR A PARTICULAR PURPOSE AND
NONINFRINGEMENT. IN NO EVENT SHALL THE AUTHORS OR COPYRIGHT HOLDERS BE
LIABLE FOR ANY CLAIM, DAMAGES OR OTHER LIABILITY, WHETHER IN AN ACTION
OF CONTRACT, TORT OR OTHERWISE, ARISING FROM, OUT OF OR IN CONNECTION
WITH THE SOFTWARE OR THE USE OR OTHER DEALINGS IN THE SOFTWARE.
\endpreamble
%    \end{macrocode}
% Clobber files without asking
%    \begin{macrocode}
\askforoverwritefalse
%    \end{macrocode}
% Make the \LaTeX{} macro files (class and styles)
%    \begin{macrocode}
\generate{% 
  \file{rmaa.cls}{\from{rmaa.dtx}{package}}
  \file{rmaa.ins}{\from{rmaa.dtx}{ins}}
  \file{rmaacite.sty}{\from{rmaa.dtx}{rmaacite}}
  \file{adshtml.sty}{\from{rmaa.dtx}{adshtml-style}}
  \file{ADS.tex}{\from{rmaa.dtx}{ads-container}}
  }
%    \end{macrocode}
% Make the other non-\LaTeX{} files. These are done separately since
% they can't use `\%\%' as a comment character for the
% preamble/postamble (and Perl files can't have the postamble I don't
% think). 
%    \begin{macrocode}
\generate{\nopreamble\nopostamble% 
  \file{Makefile}{\from{rmaa.dtx}{makefile}}
  \file{adshtml.perl}{\from{rmaa.dtx}{adshtml-perl}}
  \file{makeads.sh}{\from{rmaa.dtx}{ads-shell-script}}
  \file{sample.latex2html-init}{\from{rmaa.dtx}{latex2html-init}}
  }

%% Local Variables: 
%% mode: LaTeX
%% mode: Font-lock
%% End:
%</ins>
%    \end{macrocode}

% \subsection{Extended citation syntax}
% Modified from a file called mnrascite.sty
% I no longer remember where I stole the original version of this from
% 
% Alterations to the cite commands to make them do more sensible things.
% N.B. the bibliography style should produce entries which have labels where
% the date is delimited by angle brackets '<>' for these macros to work
% properly. They will not give very sensible error messages if you do not do 
% this properly.
% 
% A new command |\scite{key}| is defined which produces a citation which is 
% suitable for insertion in a sentence. i.e., it does not put brackets around 
% the whole citation, only the date. Also |\pcite{key}| produce output with no
% parentheses whatsoever. 
% 
% Should redefine |\@bibitem| really to have any hope of compatibility
% the command |\lbibcite| is defined to deal with the citations.
%    \begin{macrocode}
%<*rmaacite>
\def\bibitem[#1<#2>]#3{\par\hangindent=1.5em\hangafter=1
  \if@filesw
  {\def\protect##1{\string ##1\space}
    \immediate\write\@auxout{\string\lbibcite{#3}{#1}{#2}}}
  \fi}
\@redefinebibitem % what it says
%    \end{macrocode}
% Note that the |\lbibcite| defines two types of citation one contains the names,
% the other the year
%    \begin{macrocode}
\def\lbibcite#1#2#3{\global\@namedef{b@#1}{#2}
                    \global\@namedef{by@#1}{#3}}

\def\@bcite#1#2{({#1\if@tempswa , #2\fi})}
\def\@pcite#1#2{#1\if@tempswa , #2\fi}
%    \end{macrocode}
% This declaration contains what is put between the name and the date
% \changes{3.11}{2003/01/31}{Moved the \cs{articlename} to after the key for consistency
% with the new natbib stuff.}
%    \begin{macrocode}
\def\@brcite{\ (}
\def\@ebrcite{)}
\def\@fmttop#1{\@fmtcite[\csname #1\endcsname]}
\def\@citex[#1]#2{%
%%  \if@filesw\immediate\write\@auxout{\string\citation{#2}}\fi
%% WJH added 16 Jan 2002
  \@for\@citeb:=#2\do{%
    \edef\@citeb{\expandafter\@firstofone\@citeb:\articlename}%
    \if@filesw\immediate\write\@auxout{\string\citation{\@citeb}}\fi}%
%% to here
  \def\@citea{}%
  \@cite{\@for\@citeb:=#2\do
    {%
      \edef\@citeb{\expandafter\@firstofone\@citeb:\articlename}%
      \@citea\def\@citea{;\ }\@ifundefined
      {b@\@citeb}{{\bf ?}\@warning {Citation `\@citeb' on page \thepage \space undefined}}%
      {\csname b@\@citeb\endcsname\@citebro\csname by@\@citeb\endcsname\@citebrc}}}%
  {#1}}

\def\cite{\@ifnextchar [{\let\@citebro=\space
                          \let\@citebrc=\relax
                          \let\@cite=\@bcite\@tempswatrue\@citex}
                        {\let\@citebro=\space
                          \let\@citebrc=\relax
                          \let\@cite=\@bcite\@tempswafalse\@citex[]}}
\def\pcite{\@ifnextchar [{\let\@citebro=\space
                          \let\@citebrc=\relax
                          \let\@cite=\@pcite\@tempswatrue\@citex}
                        {\let\@citebro=\space
                          \let\@citebrc=\relax
                          \let\@cite=\@pcite\@tempswafalse\@citex[]}}
\def\scite{\@ifnextchar [{\let\@citebro=\@brcite
                          \let\@citebrc=\@ebrcite
                          \let\@cite=\@pcite\@tempswatrue\@citex}
                        {\let\@citebro=\@brcite
                          \let\@citebrc=\@ebrcite
                          \let\@cite=\@pcite\@tempswafalse\@citex[]}}

\renewenvironment{thebibliography}
  {\section*{REFERENCES}\vspace*{\afterbiblioskip}\small
    \if@book\RereadBiblio\fi
    \sloppy\clubpenalty4000\widowpenalty4000
    \sfcode`\.=1000\relax
    \setlength\parindent{\z@}\setlength\parskip{\z@}}
  {\relax\normalsize\@outputaddressesifnecessary} 
%</rmaacite>
%    \end{macrocode}

% \subsection{Auxilliary files for generating ADS data}

% The ADS data are written by \LaTeX{} to a |.ads| file. This file
% then has to be included in the following file (|ADS.tex|),
% which is run through \LaTeX{}2HTML (see below). 

%    \begin{macrocode}
%<*ads-container>
%%
%% Instructions:
%% Edit the \input{} line below to reflect the name of your file. 
\documentclass{article}
\usepackage[latin1]{inputenc}
\usepackage{adshtml}
\newcommand{\ADSfield}[2]
{\%#1  #2\\}
%%\newcommand{\ADSvfield}[2]
%%{\%#1 \begin{verbatim}#2\end{verbatim}}
\newenvironment{ADSarticle}{}{\par\bigskip}
\renewcommand{\bibitem}[2][]{\par}
\setlength{\parindent}{0pt}
\renewcommand{\ThisVolume}[1]{%  
  2001, \RMAASC, 12, \oldpageref{#1:RMAAFirstPage}}  
 

\begin{document}
\makeatletter
\input{rmbooksamp.ads}
\end{document}
%</ads-container>
%    \end{macrocode}

% This is the |adshtml.sty| style file used by |ADS.tex|. It just
% defines a few 
% macros that might be used in abstracts, etc (really it should define
% them all). This now gets all its code from
% Sec~\ref{sec:user-commands}. 
%    \begin{macrocode}
%<*adshtml-style>
%% Nothing else needs to go here now. 
%</adshtml-style>
%    \end{macrocode}
% The init file for latex2html.
%    \begin{macrocode}
%<*latex2html-init>
$PDFTOCAIRO = "pdftocairo";
$PDFLATEX = "pdflatex";
$LATEX = "latex";
$MAX_SPLIT_DEPTH = 0;
$HTML_VERSION = "4.0,math";
$NO_SIMPLE_MATH = 1;
$ADDRESS = "<p><address><a href=\"mailto:rmaa\@astro.unam.mx\">rmaa\@astro.unam.mx</address></p>"; 
$INFO = "<p>This document was created using latex2html and the rmaa latex macros</p>"; 
$NO_NAVIGATION = 1;
1; 
%</latex2html-init>
%    \end{macrocode}
% This is the |adshtml.perl| file that is used by \LaTeX{}2HTML. Some
% bizarre \TeX{}y command sometimes show up so we tell it to ignore
% these. All the astronomy macros that come from AAS are processed in
% \TeX{}. That way they get made to images but with an |ALT| tag that
% has the original code, which is what will be picked up by |lynx|,
% with the result that they get passed intact into the final |.txt|
% file. 
%    \begin{macrocode}
%<*adshtml-perl>
package main;

$dum = &do_cmd_makeatletter();

sub do_cmd_RMAA {
    local($_) = @_;
    "RMxAA" . $_;
}

sub do_cmd_RMAASC {
    local($_) = @_;
    "RMxAASC" . $_;
}

sub do_cmd_etal {
    local($_) = @_;
    "et al\." . $_;
}

sub do_cmd_ibidrule {
    local($_) = @_;
    "<img src=\"foo\" alt=\"\ibidrule\">" . $_;
}


sub do_cmd_bibitem {
    local($_) = @_;
    local($key);
    local($opt, $dummy)=&get_next_optional_argument; 
    $key = &missing_braces unless
        ((s/$next_pair_pr_rx/$key = $2; ''/eo)
        ||(s/$next_pair_rx/$key = $2; ''/eo));
    "<p>" . $_; 
}

&ignore_commands (<<_RAW_ARG_CMDS_);
lowercase
adjustfinalcols
textsuperscript # {}
newline # []
thinmuskip
mskip
smallskip
medskip
bigskip
hspace # {}
hspacestar # {}
vspace # {}
vspacestar # {}
mathrm
_at_mathrm
spacefactor
discretionary
unhbox 
voidb_at_x 
penalty 
_at_M
_at_m
null
IeC
_RAW_ARG_CMDS_

&process_commands_in_tex (<<_RAW_ARG_CMDS_);
sun
earth
lesssim
gtrsim
sq
argdeg
arcmin
arcsec
fd
fh
fm
fs
fdg
farcm
farcs
fp
micron
la
ga
case # {} # {}
slantfrac # {} # {}
onehalf
onethird
twothirds
onequarter
threequarters
ubvr
ub
bv
vr
ur
aj
araa
apj
apjl
apjs
ao
apss
aap
aapr
aaps
azh
baas
jrasc
memras
mnras
pra
prb
prc
prd
pre
prl
pasp
pasj
qjras
skytel
solphys
sovast
ssr
zap
nat
iaucirc
aplett
apspr
bain
fcp
gca
grl
jcp
jgr
jqsrt
memsai
nphysa
physrep
physscr
planss
procspie
_RAW_ARG_CMDS_

1;                              # This must be the last line
%</adshtml-perl> 
%    \end{macrocode}

% This is a shell script |makeads.sh| that generates the |_ads.txt| file from
% |_ads.tex| by using \LaTeX{}2HTML and |lynx|.  \changes{4.0}{2019/11/13}{Make sure that lynx writes utf-8 output instead of latin-1.}
% \changes{2.2beta13}{2002/03/12}{Complete rewrite of |makeads.sh|. Now much 
% improved I think}
% \changes{4.0}{2019/11/13}{Fix the at-sign substitution so that it only eliminates \cs{@} and does not mess with email addresses.}

%    \begin{macrocode}
%<*ads-shell-script>
#! /bin/sh
if [ -z $1 ]; then
    echo "Usage: $0 basename"
    exit
fi
BASENAME=$1
ADSBASENAME=${BASENAME}_ads
AUXFILE=$BASENAME.aux
TMPFILE=tmp.$$


echo "Generating $ADSBASENAME.tex"

# For some reason, must double-up \\ to \\\\ but not \ to \\ - odd!
cat > $ADSBASENAME.tex <<EOF
%% PLEASE DO NOT EDIT THIS FILE!  It is generated automatically from
%% the script makeads.sh.  Any hand-edited changes will be lost
%% 
\documentclass{article}
\usepackage[latin1]{inputenc}
\usepackage{html}
\usepackage{adshtml}
\newcommand{\ADSfield}[2]{\%#1  #2\\\\}
\newenvironment{ADSarticle}{}{\par\bigskip}
\setlength{\parindent}{0pt}
\newcounter{istage}
\renewcommand{\ion}[2]{\setcounter{istage}{#2}#1 \Roman{istage}}
\begin{document}
\makeatletter
EOF
echo >> $ADSBASENAME.tex "\input{$BASENAME.ads2}" 
echo >> $ADSBASENAME.tex "\end{document}"



echo "Generating $ADSBASENAME.sed" 

# Find type of journal
grep '\\documentclass\[.*proceedings.*\]' $BASENAME.tex >/dev/null \
     && JOURNAL='RMxAC' \
     || JOURNAL='RMxAA'

# Find the year and volume number
YEAR=`grep '\\SetYear' $BASENAME.tex | sed 's/.*\\SetYear{\(.*\)}.*/\1/'`
VOLUME=`grep '\\SetVolume' $BASENAME.tex | sed 's/.*\\SetVolume{\(.*\)}.*/\1/'`

# Generate an auxilliary sed script 
cat > $TMPFILE <<EOF
s|\\\\newlabel{\(.*\):RMAAFirstPage}{{.*}{\(.*\)}}|s/\\\\\\\\ThisVolume{\1}/$YEAR, RMxAC, $VOLUME, \2/g|g
EOF

# Now go through the .aux file, looking for \newlabels and use the sed 
# script just generated in order to make another sed script, one that 
# will replace \ThisVolume{LABEL} with the correct reference
grep '^\\newlabel{.*:RMAAFirstPage}' $AUXFILE \
     | sed -f $TMPFILE > $ADSBASENAME.sed && rm $TMPFILE

# Apend a further sed command to get rid of \@, since redefining 
# it in _ads.tex didn't work
echo >> $ADSBASENAME.sed 's/\\\\@//g'

echo "Generating $BASENAME.ads2 from $ADSBASENAME.sed and $BASENAME.ads" 

# Apply this to the .ads file, generating the .ads2 file, which is included
# by the _ads.tex file.
sed -f $ADSBASENAME.sed $BASENAME.ads > $BASENAME.ads2

echo "Generating $ADSBASENAME.txt from $ADSBASENAME.tex and $BASENAME.ads2"

# Now run latex2html on the _ads.tex file
latex2html -no_navigation -info 0 -address 0 -html_version 4.0,math $ADSBASENAME

# Then, convert to utf8 with lynx and tidy it up a bit
lynx --justify=off --width=100 --dump --nolist --display_charset=utf-8 \
     $ADSBASENAME/$ADSBASENAME.html | cut -c4- \
     | sed 's/ibidrule/\\ibidrule/g' > $ADSBASENAME.txt
%</ads-shell-script>
%    \end{macrocode}

% \section{Epilogue}
% Err, that's it \dots

\endinput


% \endinput
% Local Variables:
% mode: doctex
% TeX-master: t
% End:
"

rmaa.ind: rmaa.idx
        @echo "REGENERATING INDEX"
        makeindex -s gind.ist rmaa

rmaa.gls: rmaa.glo
        @echo "REGENERATING CHANGELOG"
        makeindex -s gglo.ist -o rmaa.gls rmaa.glo

rmaa.pdf: rmaa.aux rmaa.ind rmaa.gls
        @echo "REGENERATING .pdf FILE"
        pdflatex "\nonstopmode% \def\fileversion{4.0} \def\filedate{2019/11/12}
% \MakeShortVerb{\|}
% \title{A \LaTeX{} document class for the Revista Mexicana de
% Astronom\'\i a y Astrof\'\i sica\thanks{This file has version
% \fileversion{} last revised \filedate{}.}}
% \author{William J. Henney\thanks{w.henney@irya.unam.mx}}
% \date{\filedate}
% \maketitle
% \paragraph{Abstract}
% This document describes the \texttt{rmaa} \LaTeX{}
% document class for the journal ``Revista Mexicana de
% Astronom\'\i a y Astrof\'\i sica''. The macros can be used for 
% \begin{enumerate}
% \item Preparing author manuscripts for the journal (including
% conference proceedings to be published in the ``Serie de
% Conferencias'')
% \item Typesetting an entire volume of the journal by inputting the
% individual article files from a container file. Facilities are
% included for automatic generation of TOC and Author Index. 
% \end{enumerate}
% 
% \tableofcontents
% \section{The documentation driver file}
% This is the part that generates the documentation that you may or
% may not be reading, which is done with the help of the Makefile
% discussed in section~\ref{sec:driver_make}. 
%    \begin{macrocode}
%<*driver>
\documentclass{ltxdoc}
\usepackage{hyperref}
\usepackage{palatino}
\EnableCrossrefs \CodelineIndex \RecordChanges
\setlength{\parindent}{0pt}
\def\MacroFont{\footnotesize\tt}
\begin{document} \DocInput{rmaa.dtx} \PrintIndex \PrintChanges \end{document}
%</driver>
%    \end{macrocode}

% \section{The implementation}
% 
% \changes{3.24}{2007/09/14}{Added user access to the macro version
% and date via \cs{rmVersion} and \cs{rmDate}}
% 
%    \begin{macrocode}
%<*package>
\newcommand{\rmaa@version}{3.26}
\newcommand{\rmaa@date}{2011/08/15}
\newcommand\rmVersion{\rmaa@version}
\newcommand\rmDate{\rmaa@date}
\NeedsTeXFormat{LaTeX2e}[1995/06/01]
%    \end{macrocode}
% \changes{4.0}{2019/11/02}{Change abstract to appear before
% resumen. Change RevMexAA to RMxAA for consistency with ADS.}
%
% \changes{3.22}{2007/03/22}{Allow proceedings articles to have
%   received/accepted dates} \changes{3.21}{2006/05/15}{Changes for
%   Silvia - abstracts are now shipped out to ADS.}
% \changes{3.19}{2006/01/19}{Changes for Silvia - remove adresses from
%   end, new abstract style} \changes{2.2beta}{2001/10/30}{ Fixed bug
%   whereby twocolumn option was not passed through to book
%   class. Finally sorted out the ``label on last page'' issue. }
% \changes{2.1}{2001/01/23}{ Fixed some inconsistencies in
%   \cs{nocite} and \cs{bibitem} uniqification Fixed the definition of
%   \cs{ion} so it will work in math mode and in contexts where the
%   typesize is different from normal (e.g., title, fig captions).
%   Modified the definition of \cs{@caption}, so we don't write any
%   list-of-figures data (we never use it anyway, and it just causes
%   problems with fragile commands in fig captions). }
% \changes{2.0}{2000/11/30}{ Added `newstyle' option, which should be
%   used in all new proceedings volumes (wider abstract, smaller
%   typeface in fig captions, etc.) Added commands for formatting
%   table captions. New demonstation document rmpostsamp.tex for
%   1-page posters. This will be the last release under a PRI
%   government (for a while...)}  \changes{1.9}{2000/06/06}{
%   Corrected position of page number in footer of articles that open
%   on even pages. Ensured that first article in each session opens on
%   an odd page, as do list of abstracts and author index. Added extra
%   space in TOC before the backmatter. Added command for editors,
%   \cs{BoldSessionTitles}, that makes produces bold titles for
%   session in the TOC.}  \changes{1.8}{2000/05/16}{ Fixed bug in
%   last page label with \cs{suppressfulladdresses} }
% \changes{1.7}{2000/05/11}{ Fixed more little things. Page labels
%   now work, even if you have no bibliography. You can output the
%   full addresses by hand with the \cs{outputfulladdresses}
%   command. Editors' Guide updated but still not finished. Revamped
%   splitbook.sh now produces PS file of each article for web page and
%   also an EPS file of each page for sending to the printers.}
% \changes{1.6}{2000/04/15}{ More cleaning up: full stops after
%   titles in index, running headers and footers in frontmatter,
%   fontsize in title of index, etc} \changes{1.5}{2000/04/15}{
%   Eliminated superscripts from titles in TOC and index}
% \changes{1.4}{2000/04/14}{ Added copyright and license info (GPL)}
% \changes{1.3}{2000/03/01}{ Corrected \cs{figcaption} to do the
%   caption formatting right Corrected font size of sec headings in
%   front matter} \changes{1.2}{2000/02/16}{ Added \cs{plate}
%   command} \changes{1.1}{2000/02/14}{ Added manuscript option,
%   corrected bugs} \changes{1.0}{2000/02/08}{ Now does a whole book
%   - lots of stuff tidied up} \changes{0.7}{1999/10/04}{ Increased
%   spacing between sections in single column styles. Suppressed
%   floats at bottom of pages (didn't work though) }
% \changes{0.6}{1999/09/24}{ Tweaked \cs{adjustfinalcols} (still not
%   perfect) Added proceedings option Added onecolumn/twocolumn
%   options (orthogonal to journal/proceedings) }
% \changes{0.5}{1999/09/23}{ Added \cs{fulladdresses} and
%   \cs{adjustfinalcols} (just about working)}
% \changes{0.4}{1999/09/22}{ Minor stylistic changes}
% \changes{0.3}{1999/06/28}{ Better mechanism for alternative
%   affiliations Added `guide' option for user guide}
% \changes{0.2}{1999/05/27}{ All basic functionality present, I
%   think} \changes{0.1}{1999/05/26}{ Just about working - still
%   missing Fig/Tab/Refs} \changes{0.0}{1999/05/21}{ So bored today,
%   decided to rewrite RMAA LaTeX macros from scratch}
%    \begin{macrocode}
\ProvidesClass{rmaa}[\rmaa@date\space v\rmaa@version\space RMxAA article
class]
%    \end{macrocode}
% \subsection{Option Handling}
% \label{sec:options}
% Declarations that are needed in options
% \changes{3.05}{2003/10/22}{Added natbib option}
% \changes{3.06}{2003/01/03}{Natbib option turned on permanently}
% \changes{3.07}{2003/01/11}{Added proof option}
% \changes{3.14}{2003/01/11}{Option to turn off natbib}
% \changes{3.18}{2003/10/22}{trying out doctex mode} 
% \changes{3.23}{2007/09/09}{Manuscript is now on by default.}
% \changes{3.26}{2011/08/15}{New option for review articles}
%    \begin{macrocode}
\newif\if@reviewarticle
\@reviewarticlefalse
\newif\if@proof
\@prooffalse
\newif\if@journal
\@journaltrue
\newif\if@printdates
\@printdatesfalse
\newif\if@proceedings
\@proceedingsfalse
\newif\if@manuscript
\@manuscripttrue
\newif\if@bigfont
\@bigfonttrue
\newif\if@guide
\@guidefalse
\newif\if@preprint
\@preprintfalse
\newif\if@rmtwocol
\@rmtwocoltrue
\newif\if@biblist
\@biblistfalse
\newif\if@book
\@bookfalse
\newif\if@natbib
\@natbibtrue
\newif\if@rmcite
\@rmcitefalse
\newif\if@titlevisible
\@titlevisibletrue
\newif\if@abstractpresent
\@abstractpresentfalse
\newif\if@printbanner
\@printbannerfalse
\newif\if@suppressfulladdresses
\@suppressfulladdressesfalse
\newif\if@newstyle
\@newstylefalse
\newif\if@debug
\@debugfalse
\newif\if@boldkeywords
\@boldkeywordsfalse
%    \end{macrocode}
% The option handling---mainly setting various booleans.
% \changes{2.2beta11}{2001/11/26}{Made twocol and
% newstyle the default for proceedings.}
% \changes{3.23}{2007/09/09}{Manuscript is now 10 pt instead of 12pt}
% \changes{3.23}{2007/09/10}{Full addresses are now automatically
% suppressed for proceedings style.}
% \changes{3.24}{2007/09/14}{New option \texttt{boldkeywords} that
% does the keywords in the old ugly way. Now off by default.}
%    \begin{macrocode}
\DeclareOption{manuscript}{\typeout{Manuscript for referee} 
  \@manuscripttrue\@journalfalse\@rmtwocolfalse
  \PassOptionsToClass{10pt}{article}}
\DeclareOption{bigfont}{\typeout{Big font manuscript requested by Gary
  Ferland} 
  \@manuscripttrue\@journalfalse\@rmtwocolfalse\@bigfonttrue
  \PassOptionsToClass{11pt}{article}}
\DeclareOption{journal}{\typeout{Main journal option}%
  \@journaltrue\@rmtwocoltrue\@printdatestrue\@manuscriptfalse
  \PassOptionsToClass{10pt}{article}}
\DeclareOption{guide}{\typeout{User guide option}%
  \@guidetrue\@journalfalse\@rmtwocoltrue\@manuscriptfalse
  \PassOptionsToClass{10pt}{article}}
\DeclareOption{preprint}{\typeout{Preprint option}%
  \@preprinttrue\@manuscriptfalse\PassOptionsToClass{10pt}{article}}
\DeclareOption{proof}{\typeout{Author proof option}%
  \@prooftrue\@manuscriptfalse}
\DeclareOption{proceedings}{\typeout{Proceedings option}%
  \@proceedingstrue\@journalfalse\@rmtwocoltrue\@newstyletrue
  \@suppressfulladdressestrue
  \@printdatesfalse\@manuscriptfalse\PassOptionsToClass{10pt}{article}} 
\DeclareOption{printdates}{\typeout{Force printing of received/accepted dates}%
  \@printdatestrue} 
\DeclareOption{book}{\typeout{Whole issue option}%
  \@booktrue} 
\DeclareOption{nonatbib}{\typeout{Turn off natbib support (required
    for compatibility with old rmaacite package).}%
  \@natbibfalse} 
\DeclareOption{newstyle}{\typeout{Use new layout style}%
  \@newstyletrue}
\DeclareOption{oldstyle}{\typeout{Use old layout style}%
  \@newstylefalse}
\DeclareOption{debug}{\typeout{Loads of annoying messages}%
  \@debugtrue}
%    \end{macrocode}
% These two need to be declared explicitly so that |\ExecuteOptions|
% actually does something.
%    \begin{macrocode}
\DeclareOption{twocolumn}{\@rmtwocoltrue}
\DeclareOption{onecolumn}{\@rmtwocolfalse}
\DeclareOption{twoside}{\PassOptionsToClass{\CurrentOption}{article}}
\DeclareOption{boldkeywords}{\@boldkeywordstrue}
%    \end{macrocode}
% Bold, upper-case keywords used to be the default. Now they must be
% turned on explicitly with the \texttt{boldkeywords} option. 
%    \begin{macrocode}
\DeclareOption{boldkeywords}{\@boldkeywordstrue}
%    \end{macrocode}
% Pass any other option along to the \texttt{article} class.
%    \begin{macrocode}
\DeclareOption*{\PassOptionsToClass{\CurrentOption}{article}}
%    \end{macrocode}
% Default options are to be two-sided and based on the standard
% \LaTeX{} \texttt{article} class,
% unless we are processing the entire book, in which case we base
% ourselves on the standard \texttt{book} class. Default is to use
% manuscript option unless otherwise specified. 
%    \begin{macrocode}
\ExecuteOptions{twoside,manuscript} \ProcessOptions\relax
\if@book
\if@rmtwocol\LoadClass[twocolumn]{book}\else\LoadClass{book}\fi
\newcommand{\abstract}{\relax}
\newcommand{\@maketitle}{\relax}
\else
\if@rmtwocol\LoadClass[twocolumn]{article}\else\LoadClass{article}\fi
\fi
%    \end{macrocode}
% Attempt to use natbib for citations 02 Jan 2003. 
% \changes{3.05}{2003/01/02}{If the natbib option is used, then set
% things up so that natbib's version of the bibliography is identical
% to the standard one.}
% \changes{3.09}{2003/01/23}{Explicitly require a recent version of
% natbib, so that we get a decent error message if the package is too
% old.} 
% \changes{3.23}{2007/09/09}{Now allow Bibtex to work properly} 
% \changes{3.25}{2009/10/01}{Change \cs{newcommand} tp \cs{def} when setting \cs{bibfont}. This allows it to work with all versions of natbib.} 
% \changes{3.25}{2010/02/06}{Updated copyright year.} 
% \changes{3.25}{2010/02/06}{Updated email and web addresses to \texttt{crya.unam.mx}.} 
%    \begin{macrocode}
\if@natbib
\RequirePackage{natbib}[1999/02/25]
\bibpunct{(}{)}{;}{a}{}{,}
\def\bibfont{\small}
\renewcommand\bibsection{\section*{REFERENCES}}
\setlength\bibsep{0pt}
\setlength\bibhang{1.5em}
\bibliographystyle{rmaa}
\fi
%    \end{macrocode}
% We want the \texttt{graphicx} package for the extended syntax of
% |\includegraphics|. We want the \texttt{ifthen} package for using in
% these macros (although we don't much yet). We explicitly load
% \texttt{babel} since some \LaTeX{} distributions messed up the
% hyphenation if we didn't.  
%    \begin{macrocode}
\RequirePackage{graphicx}
\RequirePackage{ifthen}
\IfFileExists{babel.sty}{%
%   \let\@restoreindent=\@afterindentfalse
  \RequirePackage[english]{babel}
%   \@restoreindent
  %% reinstate the definition of \@roman so as not to break .idx files
  \@ifundefined{org@@roman}{\relax}{\renewcommand{\@roman}[1]{\org@@roman{##1}}}
  }{%
  \typeout{Package babel.sty not found - don't worry about this}}
%    \end{macrocode}
% Try to load the \texttt{dcolumn} and \texttt{booktabs} packages if
% they can be found. These can produce nicer-looking tables. 
%    \begin{macrocode}
\IfFileExists{dcolumn.sty}{%
  \RequirePackage{dcolumn}
  }{%
  \typeout{Package dcolumn.sty not found - you will not be able to
    align columns on decimal points in tables}}
\IfFileExists{booktabs.sty}{%
  \RequirePackage{booktabs}
  }{%
  \typeout{Package booktabs.sty not found - your tables will not look
    as pretty as they ought too.}
  }
\IfFileExists{makeidx.sty}{%
  \RequirePackage{makeidx}
}{%
  \typeout{Package makeidx.sty not found - Author Index may not work.}
  }
%% commented out for moment since requires more modern version of LaTeX
%% \RequirePackage[latin1]{inputenc}
%    \end{macrocode}
% \changes{4.0}{2019/11/13}{Make \texttt{hyperref} and
% \texttt{etoolbox} into hard requirements, since we now use them
% extensively.}  We use hyperref to get internal and external links.
% But I don't like the colored boxes so I turn those off.  We use
% etoolbox for patching and robustifying at the moment.  In the
% future, I want to port a lot of other stuff over to it.  For
% instance, replace the ifthen package and the TeX ifs too.
%
% We target the latest version (1.7) of PDF, rather than the default
% of 1.5.  This reduces the chance of getting warnings from included
% graphics.  This is not strictly related to hyperref, but we put it
% here anyway.
%    \begin{macrocode}
\RequirePackage{etoolbox}
\RequirePackage{hyperref}
\hypersetup{hidelinks}
\pdfminorversion=7
%    \end{macrocode}
% We patch the \cs{@combinedblfloats} following a suggestion by Heiko
% Oberdiek at \url{https://tex.stackexchange.com/a/249743/3843}, which
% avoids problems when hyperlinks are split across
% pages. \changes{4.0}{2019/11/12}{Fix for issue where breaking
% hyperlink across pages gave an error.}
%    \begin{macrocode}
\urlstyle{same}
\newcommand\email{\begingroup \urlstyle{rm}\Url}
\patchcmd\@combinedblfloats{\box\@outputbox}{\unvbox\@outputbox}{}{%
  \errmessage{\noexpand\@combinedblfloats could not be patched}%
}%
%    \end{macrocode}
% \subsection{Other setup stuff}\label{sec:setup}
% Some declarations: 
%    \begin{macrocode}
\newif\if@needaddresses
\@needaddressestrue
\newif\if@frontmatter
\@frontmatterfalse
\newif\if@inpreamble
\@inpreambletrue
\newif\if@firstauthor
\@firstauthortrue
\newif\if@firstaffil
\@firstaffiltrue
\newif\if@firstaltaffil
\@firstaltaffiltrue
\newif\if@firstkwd
\@firstkwdtrue
\newif\if@noexpandADS
\@noexpandADSfalse
\newif\if@sessiontitlepages
\@sessiontitlepagesfalse
%    \end{macrocode}
% \begin{macro}{\@rmaaEmpty}
% \begin{macro}{\@rmaaArg}
% These are used in various commands to check if arg was left blank
%    \begin{macrocode}
\newcommand{\@rmaaEmpty}{}
\newcommand{\@rmaaArg}{}
%    \end{macrocode}  
% \end{macro}
% \end{macro}
%
% \subsection{Compatibility with old \LaTeX{} versions}
% \label{sec:compat}
% This package is based on the 1998/12/01 release of \LaTeX{}.
% Older releases may not define all the commands that we need.
% Luckily, we can use |\providecommand| to define these if necessary. 
% The following commands are not defined in the 1994/06/01 release. 
%
% \begin{macro}{\MakeUpperCase}
% \begin{macro}{\textsuperscript}
% \begin{macro}{\textbackslash}
% User commands:
%    \begin{macrocode}
\providecommand{\MakeUpperCase}[1]{\uppercase{#1}}
\providecommand{\textsuperscript}[1]{\ensuremath{^\textrm{#1}}}
\providecommand{\textbackslash}{\ensuremath{\backslash}}
%    \end{macrocode}  
% \end{macro}
% \end{macro}
% \end{macro}
% \begin{macro}{\protected@xdef}
% \begin{macro}{\restore@protect}
% \begin{macro}{\@unexpandable@protect}
% \begin{macro}{\hb@xt@}
% \LaTeX{} internal commands:
%    \begin{macrocode}
\providecommand{\protected@xdef}{%
        \let\@@protect\protect
        \let\protect\@unexpandable@protect
        \afterassignment\restore@protect
        \xdef}
\providecommand{\restore@protect}{\let\protect\@@protect}
\providecommand{\@unexpandable@protect}{\noexpand\protect\noexpand}
\providecommand{\hb@xt@}{\hbox to}
%    \end{macrocode}
% \end{macro}
% \end{macro}
% \end{macro}
% \end{macro}
%
% \subsection{Page layout}\label{sec:page-layout}
%
% The dimensions of the text were defined by Elizabeth. Text width is
% 169\,mm with equal margins on the two sides: 
%    \begin{macrocode}
\setlength{\textwidth}{169mm}
\setlength{\@tempdima}{\paperwidth}
\addtolength{\@tempdima}{-\textwidth}
\setlength{\oddsidemargin}{0.5\@tempdima}
\addtolength{\oddsidemargin}{-1in}
\setlength{\evensidemargin}{\oddsidemargin}
%    \end{macrocode}
% Text height is 225\,mm, top margin is 14\,mm. 
%    \begin{macrocode}
\setlength{\textheight}{225mm} 
\setlength{\topmargin}{-1in} 
\addtolength{\topmargin}{14mm}
%    \end{macrocode}
% Space between columns in two-column styles, header, footer,
% paragraph indentation:
%    \begin{macrocode}
\setlength{\columnsep}{9.5mm}
\setlength{\headsep}{0.7cm}
\setlength{\footskip}{0.75cm}
\setlength{\parindent}{1.5em}
%    \end{macrocode}
% 
% \begin{macro}{\@set@scalable@length}
% Try out a new scheme for dealing with various spacing dimensions and
% allowing the user to easily modify them. The command
% \begin{quote}
% |\set@scalable@length|\marg{name}\marg{defaultlength}\marg{stretch}\marg{shrink}
% \end{quote}
% does a 
% |\setlength|\marg{\cs{name}}\marg{defaultlength plus stretch minus shrink}
% and also defines the user command
% \textbackslash\meta{name}|Scale|\marg{factor} that can 
% be used to rescale length by \meta{factor}. Additionally, now
% defines the package command
% \textbackslash\meta{name}|@set@scale|\marg{FACTOR} 
% that also rescales the length by \meta{FACTOR}, so that the length
% becomes \meta{factor} $\times$ \meta{FACTOR} $\times$
% \meta{defaultlength}.   
%    \begin{macrocode}
\newcommand\@set@scalable@length[4]{%
  %% arg #1: NAME of length without leading `\'
  %%
  %% arg #2: default VALUE  
  \expandafter\newlength\csname @#1@default\endcsname
  \expandafter\setlength\csname @#1@default\endcsname{#2}%
  \expandafter\newlength\csname @@#1@@default\endcsname
  \expandafter\setlength\csname @@#1@@default\endcsname{#2}%
  %% arg #3: default STRETCH  
  \expandafter\newlength\csname @#1@stretch\endcsname
  \expandafter\setlength\csname @#1@stretch\endcsname{#3}%
  \expandafter\newlength\csname @@#1@@stretch\endcsname
  \expandafter\setlength\csname @@#1@@stretch\endcsname{#3}%
  %% arg #4: default SHRINK  
  \expandafter\newlength\csname @#1@shrink\endcsname
  \expandafter\setlength\csname @#1@shrink\endcsname{#4}%
  \expandafter\newlength\csname @@#1@@shrink\endcsname
  \expandafter\setlength\csname @@#1@@shrink\endcsname{#4}%
  %%
  %% User-defined scale \@NAME@scale is unity by default
  \expandafter\gdef\csname @#1@scale\endcsname{1}%
  %% Package-defined scale \@NAME@SCALE is unity by default
  \expandafter\gdef\csname @#1@SCALE\endcsname{1}%
  %%
  %% internal command \@reset@NAME should expand to 
  %% \setlength\@@NAME@@default{\@NAME@SCALE\@NAME@default}
  %% \setlength\NAME{\@NAME@scale\@@NAME@@default}
  %% in other words: scale x SCALE x default 
  \expandafter\gdef\csname @reset@#1\endcsname{%
    \expandafter\setlength\csname @@#1@@default\endcsname{%
      \csname @#1@SCALE\endcsname \csname @#1@default\endcsname
      }
    \expandafter\setlength\csname @@#1@@stretch\endcsname{%
      \csname @#1@SCALE\endcsname \csname @#1@stretch\endcsname
      }
    \expandafter\setlength\csname @@#1@@shrink\endcsname{%
      \csname @#1@SCALE\endcsname \csname @#1@shrink\endcsname
      }
    \expandafter\setlength\csname #1\endcsname{%
      \csname @#1@scale\endcsname \csname @@#1@@default\endcsname
      \@plus
      \csname @#1@scale\endcsname \csname @@#1@@stretch\endcsname
      \@minus
      \csname @#1@scale\endcsname \csname @@#1@@shrink\endcsname
      }}
  %%
  %% User command \NAMEScale should expand to 
  %% {\renewcommand\@NAME@scale{#1}\@reset@NAME}
  \expandafter\gdef\csname #1Scale\endcsname##1{%
    \expandafter\gdef\csname @#1@scale\endcsname{##1}%
    \csname @reset@#1\endcsname}
  %%
  %% Package command \NAME@set@SCALE should expand to 
  %% {\renewcommand\@NAME@SCALE{#1}\@reset@NAME}
  \expandafter\gdef\csname #1@set@SCALE\endcsname##1{%
    \expandafter\gdef\csname @#1@SCALE\endcsname{##1}%
    \csname @reset@#1\endcsname}
  %%
  %% initialize at default values with \@reset@NAME
  \csname @reset@#1\endcsname
  }
%    \end{macrocode}
% \end{macro}
% \begin{macro}{\parskip}
% First use of this new infrastructure: set interparagraph space to be
% zero but with a small amount of stretchiness.
%    \begin{macrocode}
%   \ifthenelse{\boolean{@manuscript}}{%
%     \@set@scalable@length{parskip}{0.5\baselineskip}{0.3\p@}{0.1\p@}
%   }{
    \@set@scalable@length{parskip}{0\p@}{0.3\p@}{0\p@}
%   }
%    \end{macrocode}
% \end{macro}
% 
% \begin{macro}{\abovesectionskip}
% \changes{4.0}{2019/11/12}{Use negative values for skip so that first paragraph after section headings is not indented.  This is something that I have been trying to fix for years!  It turns out that the anser was in the documentation for the standard command \cs{@startsection.}}
% \begin{macro}{\belowsectionskip}
% \begin{macro}{\aboveparaskip}
% \begin{macro}{\belowparaskip}
% Define some utility rubber lengths for use in
% section~\ref{sec:section-commands}. 
%    \begin{macrocode}
\newlength{\abovesectionskip}
\newlength{\belowsectionskip}
\newlength{\aboveparaskip}
\newlength{\belowparaskip}
%    \end{macrocode}
% The default values for these lengths are those traditionally used
% for the onecolumn style. 
%    \begin{macrocode}
\@set@scalable@length{abovesectionskip}{-3.5ex}{-4ex}{-2ex}
\@set@scalable@length{belowsectionskip}{2.3ex}{0ex}{2ex}
\@set@scalable@length{aboveparaskip}{3.25ex}{2ex}{0.5ex}
\@set@scalable@length{belowparaskip}{-1em}{\z@}{\z@}
%    \end{macrocode}
% \end{macro}\end{macro}\end{macro}\end{macro}
% We also make scalable some of the standard float parameters. These
% are the same defaults as in the standard article classes, but
% slightly more stretchy. For the twocolumn style, they are currently
% halved, see below. Note that I have also lumped |\abovecaptionskip|
% in with these, since it plays a similar role. There is also a
% |\belowcaptionskip| but it is set to zero in the standard classes. 
%    \begin{macrocode}
\@set@scalable@length{floatsep}{12\p@}{4\p@}{4\p@}
\@set@scalable@length{textfloatsep}{20\p@}{4\p@}{8\p@}
\@set@scalable@length{intextsep}{12\p@}{4\p@}{4\p@}
\@set@scalable@length{dblfloatsep}{12\p@}{4\p@}{4\p@}
\@set@scalable@length{dbltextfloatsep}{20\p@}{4\p@}{8\p@}
\@set@scalable@length{abovecaptionskip}{10\p@}{\z@}{\z@}
%    \end{macrocode}
% \begin{macro}{\@before@title@skip}         
% \begin{macro}{\@before@author@skip}        
% \begin{macro}{\@between@author@skip}       
% \begin{macro}{\@before@abstract@skip}      
% \begin{macro}{\@after@abstracthead@skip}   
% \begin{macro}{\@between@abstract@skip}     
% \begin{macro}{\@before@kwd@skip}           
% \begin{macro}{\@before@dates@skip}           
% Spacing parameters used in |\@maketitle|. These are not changed by
% default except by the user, so they don't need to go into
% |\@package@default@lengths|.  
%    \begin{macrocode}
\newlength{\@before@title@skip}         
\newlength{\@before@author@skip}        
\newlength{\@between@author@skip}       
\newlength{\@before@abstract@skip}      
\newlength{\@after@abstracthead@skip}   
\newlength{\@between@abstract@skip}     
\newlength{\@before@kwd@skip}           
\newlength{\@before@dates@skip}           
\@set@scalable@length{@before@title@skip}{2em}{\z@}{\z@}
\@set@scalable@length{@before@author@skip}{1.5em}{\z@}{\z@}
\@set@scalable@length{@between@author@skip}{0.5em}{\z@}{\z@}
\@set@scalable@length{@before@abstract@skip}{3ex}{1ex}{0.5ex}
\@set@scalable@length{@after@abstracthead@skip}{1ex}{\z@}{\z@}
\@set@scalable@length{@between@abstract@skip}{2ex}{1ex}{0.5ex}
\@set@scalable@length{@before@kwd@skip}{1ex}{0.5ex}{0.2ex}
\@set@scalable@length{@before@dates@skip}{2ex}{0.5ex}{0.2ex}
%    \end{macrocode}
% \end{macro}\end{macro}\end{macro}\end{macro}
% \end{macro}\end{macro}\end{macro}\end{macro}
% 
% Spacing above and below displayed equations also wants to be done
% like this, but it is more complicated since the standard article
% classes make these change with the fontsize. Something to do later,
% maybe. 
% 
%    \begin{macrocode}
%    \end{macrocode}
% \begin{macro}{\@package@default@lengths}
% The value actually used for these lengths is controlled by the
% \meta{SCALE} parameter, which depends on the style option that is in
% effect. This is set using the |\@package@default@lengths| command,
% which must be called every time the style changes (e.g., changing
% from onecolumn to twocolumn). 
%    \begin{macrocode}
\newcommand{\@package@default@lengths}{%
  \ifthenelse{\boolean{@rmtwocol}}{%
    %% twocolumn mode
    \abovesectionskip@set@SCALE{0.5}
    \belowsectionskip@set@SCALE{0.5}
    \aboveparaskip@set@SCALE{0.5}
    \belowparaskip@set@SCALE{1}
    \floatsep@set@SCALE{0.5}
    \textfloatsep@set@SCALE{0.5}
    \intextsep@set@SCALE{0.5}
    \dblfloatsep@set@SCALE{0.5}
    \dbltextfloatsep@set@SCALE{0.5}
    }{%
    %% onecolumn mode
    \abovesectionskip@set@SCALE{1}
    \belowsectionskip@set@SCALE{1}
    \aboveparaskip@set@SCALE{1}
    \belowparaskip@set@SCALE{1}
    \floatsep@set@SCALE{1}
    \textfloatsep@set@SCALE{1}
    \intextsep@set@SCALE{1}
    \dblfloatsep@set@SCALE{1}
    \dbltextfloatsep@set@SCALE{1}
    }%
  }
\@package@default@lengths
%    \end{macrocode}
% \end{macro}
% \begin{macro}{\RescaleLengths}
% \begin{macro}{\RescaleSecLengths}
% \begin{macro}{\RescaleFloatLengths}
% \begin{macro}{\RescaleTitleLengths}
% User commands that will fiddle with all the scalable lengths. 
%    \begin{macrocode}
\newcommand{\RescaleLengths}[1]{%
  \RescaleSecLengths{#1}\RescaleFloatLengths{#1}\RescaleTitleLengths{#1}}
\newcommand{\RescaleSecLengths}[1]{%
%% sections
  \abovesectionskipScale{#1}
  \belowsectionskipScale{#1}
  \aboveparaskipScale{#1}
  }
\newcommand{\RescaleFloatLengths}[1]{%
%% floats
  \floatsepScale{#1}
  \textfloatsepScale{#1}
  \intextsepScale{#1}
  \dblfloatsepScale{#1}
  \dbltextfloatsepScale{#1}
  \abovecaptionskipScale{#1}
  }
\newcommand{\RescaleTitleLengths}[1]{%
%% maketitle
  \@before@title@skipScale{#1}           
  \@before@author@skipScale{#1}        
  \@between@author@skipScale{#1}       
  \@before@abstract@skipScale{#1}      
  \@after@abstracthead@skipScale{#1}   
  \@between@abstract@skipScale{#1}     
  \@before@kwd@skipScale{#1}           
  \@before@dates@skipScale{#1}           
  }
%    \end{macrocode}
% \end{macro}\end{macro}\end{macro}\end{macro}
% \changes{2.2beta14}{2002/03/17}{Finished implementation of new
% infrastructure for scalable lengths. Converted the above/below
% section skips to use it. Also, all the floatsep parameters. Also,
% all the vertical lengths that are used by |\@maketitle|. }

% \begin{macro}{\MSextramargin}
% For the \texttt{manuscript} option we have an extra 0.3\,inch
% vertical margin and 1\,inch horizontal margin for marginal notes (by
% hand or with |\marginpar|. 
%
% \changes{3.23}{2007/09/09}{Reduce baselinestretch in
%   \texttt{manuscript} style to 1.2. No longer use raggedright for
%   manuscript.}
%
% \changes{3.24}{2007/09/14}{Further reduction of baselinestretch in
%   \texttt{manuscript} style to 1.1 since we now have wider margins.}
%
% \changes{3.24}{2007/09/14}{Split \texttt{manuscript}'s extra margin
%   into separate horizontal and vertical parts.}
% 
%    \begin{macrocode}
\newlength{\MSextraHmargin}\newlength{\MSextraVmargin}
\setlength{\MSextraHmargin}{1in} 
\setlength{\MSextraVmargin}{0.3in} 
\if@manuscript
   \renewcommand{\baselinestretch}{1.1}
   \addtolength{\textwidth}{-2\MSextraHmargin}
   \addtolength{\oddsidemargin}{\MSextraHmargin}
   \setlength{\evensidemargin}{\oddsidemargin}
   \addtolength{\textheight}{-2\MSextraVmargin}
   \addtolength{\topmargin}{\MSextraVmargin}
%    \raggedright
\fi
%    \end{macrocode}
% \end{macro}
% Float parameters: don't allow bottom floats by default. 
%    \begin{macrocode}
\renewcommand{\fps@figure}{tp}
\renewcommand{\fps@table}{tp}
%    \end{macrocode}
%
% \begin{macro}{\rmaa@AuthorWidthFrac}
% \begin{macro}{\rmaa@AbstractWidthFrac}
% These parameters control what fraction of the total text width can
% be used by the list of authors and the abstract/resumen,
% respectively.  
%    \begin{macrocode}
\newcommand{\rmaa@AuthorWidthFrac}{1.0}
\newcommand{\rmaa@AbstractWidthFrac}{0.75}
%    \end{macrocode}
% Use a wide abstract to save space in the new style.
%    \begin{macrocode}
\if@newstyle\renewcommand{\rmaa@AbstractWidthFrac}{1.0}\fi
%    \end{macrocode}
% \end{macro}
% \end{macro}
%
% Some things to consider in future versions. Trying to avoid too many
% widows, etc. Apparently, the following will work best with
% \cs{flushbottom}. 
%    \begin{macrocode}
%%\widowpenalty=10000
%%\clubpenalty=10000
%    \end{macrocode}
% Also, should perhaps try getting rid of underfull/overfull box warnings. 
% \changes{3.10}{2003/01/24}{Set large values for \cs{hbadness},
% \cs{vbadness}, \cs{hfuzz}, \cs{vfuzz}. To reduce amount of stuff
% that gets written to the terminal every time you compile something.}
%    \begin{macrocode}
\hbadness=10000
\vbadness=10000
\setlength\hfuzz{\maxdimen}
\setlength\vfuzz{\maxdimen}
%    \end{macrocode}
% \subsection{Various text styles for different contexts}
% \label{sec:text-styles} 
% \begin{macro}{\rmaa@AffilStyle}
% \begin{macro}{\rmaa@PagesStyle}
% \begin{macro}{\rmaa@YearStyle}
% \begin{macro}{\rmaa@VolumeStyle}
% \begin{macro}{\rmaa@VolumenStyle}
% \begin{macro}{\rmaa@JournalNameStyle}
% \begin{macro}{\rmaa@ReceivedStyle}
% \begin{macro}{\rmaa@AbstractHeaderStyle}
% \begin{macro}{\rmaa@JournalNameStyle}
% Style commands for the different fields used by |\maketitle| in
% constructing the title page of the article. These have the form
% |\rmaa@XYZStyle{TEXT}|. I'm not sure these were really a good idea,
% but never mind \dots 
%    \begin{macrocode}
\newcommand{\rmaa@AffilStyle}[1]{{\normalsize\rmfamily#1}}
\newcommand{\rmaa@PagesStyle}[1]{{\normalsize\rmfamily#1}}
\newcommand{\rmaa@YearStyle}[1]{{\normalsize\rmfamily#1}}
\newcommand{\rmaa@VolumeStyle}[1]{{\normalsize\rmfamily\bfseries#1}}
\newcommand{\rmaa@VolumenStyle}[1]{{\Large\sffamily\bfseries#1}}
\newcommand{\rmaa@JournalNameStyle}[1]{{\normalsize\itshape#1}}
\newcommand{\rmaa@ReceivedStyle}[1]{{\normalsize\rmfamily\itshape#1}}
\newcommand{\rmaa@AbstractHeaderStyle}[1]{%
  {\normalsize\rmfamily\MakeUpperCase{#1}}}
\newcommand{\rmaa@TitleStyle}[1]{%
  {\large\bfseries\rmfamily\boldmath\MakeUpperCase{#1}\par}}
%    \end{macrocode}
% \end{macro}
% \end{macro}
% \end{macro}
% \end{macro}
% \end{macro}
% \end{macro}
% \end{macro}
% \end{macro}
% \end{macro}
%
% \begin{macro}{\rmaa@FrontSecStyle}
% Headings of ``front sections'' (see section~\ref{sec:front-sections})
%    \begin{macrocode}
\newcommand{\rmaa@FrontSecStyle}[1]{%
  {\normalsize\bfseries\rmfamily\MakeUpperCase{#1}}}
%    \end{macrocode}
% \end{macro}
%
% \begin{macro}{\rmaa@TOCTitleStyle}
% \begin{macro}{\rmaa@TOCPartStyle}
% \begin{macro}{\rmaa@IndexTitleStyle}
% \begin{macro}{\rmaaTOCAuthorStyle}
% \begin{macro}{\rmaaIndexAuthorStyle}
% Commands to set style of TOC and index entries.
% Rewritten to avoid unnecessary font commands in the .toc and .ind files
%    \begin{macrocode}
\newcommand{\rmaa@TOCTitleStyle}[1]{\MakeUpperCase{#1}}
\newcommand{\rmaa@TOCPartStyle}[1]{\rmaa@TOCTitleStyle{#1}}
\newcommand{\rmaa@IndexTitleStyle}[1]{#1}
\newcommand{\rmaaTOCAuthorStyle}[1]{%
  {\itshape#1}}
\newcommand{\rmaaIndexAuthorStyle}[1]{%
  \protect\mbox{\bfseries#1}}
%    \end{macrocode}
% \end{macro}
% \end{macro}
% \end{macro}
% \end{macro}
% \end{macro}
%
% \begin{macro}{\rmaa@KeywordStyle}
% \begin{macro}{\rmaa@KeywordHeaderStyle}
% The style of the keywords depends on whether we are using the
% \texttt{newstyle} option or not. If we are, then they are small. 
%    \begin{macrocode}
\ifthenelse{\boolean{@newstyle}}{
  \newcommand{\rmaa@KeywordStyle}[1]{%
    {\small%
      \if@boldkeywords 
      \bfseries\MakeUpperCase{\raggedright #1}
      \else
      \raggedright #1
      \fi
    }}
  \newcommand{\rmaa@KeywordHeaderStyle}[1]{%
    {\small\itshape#1}}}
{
  \newcommand{\rmaa@KeywordStyle}[1]{%
    {\normalsize%
      \if@boldkeywords\bfseries\fi\raggedright #1}}
  \newcommand{\rmaa@KeywordHeaderStyle}[1]{%
    {\normalsize\itshape#1}}}
%    \end{macrocode}
% \end{macro}
% \end{macro}
%
% \subsection{Page styles}\label{sec:page-styles}
% \begin{macro}{\ps@rmaa}
% The running headers for normal pages:
%    \begin{macrocode}
\newcommand{\ps@rmaa}{%
%    \end{macrocode}
% nothing in either foot; 
%    \begin{macrocode}
  \renewcommand\@oddfoot{\relax}%
  \renewcommand\@evenfoot{\relax}%
%    \end{macrocode}
% page number and short author list in the even heads; 
%    \begin{macrocode}
  \renewcommand\@evenhead{\rlap{\thepage}\hfil\@shortauthor\hfil}%
%    \end{macrocode}
% short title and page number in the odd heads.
%    \begin{macrocode}
  \renewcommand\@oddhead{\hfil\@shorttitle\hfil\llap{\thepage}}
  }
%    \end{macrocode}
% \begin{macro}{\ps@rmaafirstpage}
% The header of the first page of each
% article:
%    \begin{macrocode}
\newcommand{\ps@rmaafirstpage}{%
  \renewcommand\@oddfoot{\hfil\thepage}%
%    \end{macrocode}
% footer corrected (06 Jun 2000) to have page on outside edge;
%    \begin{macrocode}
  \renewcommand\@evenfoot{\thepage\hfil}%
%    \end{macrocode}
% Header has journal name, volume, and page numbers for a regular
% article, or similar for special cases. 
% \changes{4.0}{2019/11/03}{Change to three-line header with copyright info and DOI.}
%    \begin{macrocode}
\if@preprint
\renewcommand\@oddhead{%
  \if@proceedings 
  \textsc{To appear in ``\@conftitle{} (\rmaa@YearStyle{\@year})''}\hfil
  \textit{RMxAC}
  \else
  \textsc{To appear in} \@journalname\hfil
  \fi}\fi%
\if@guide
\renewcommand\@oddhead{\textsc{Documentation} (\rmaa@YearStyle{\@year})\hfil}
\fi%
\if@manuscript
\renewcommand\@oddhead{Manuscript for \@journalname{} (\rmaa@YearStyle{\@year})\hfil}
\fi%
\if@journal
\renewcommand\@oddhead{%
  \parbox{\textwidth}{
    \@journalname, \rmaa@VolumeStyle{\@volume},
    \rmaa@PagesStyle{\@artpages}\
    (\rmaa@YearStyle{\@year})\hfil\\[1pt]
    {\small\textcopyright{} \@year:
      Instituto de Astronom\'\i{}a,
      Universidad Nacional Aut\'onoma de M\'exico}\hfil\\[1pt]
    {\small\@DOI@url}\hfil
  }%
}\fi%
  \renewcommand\@evenhead{\@oddhead}}
%    \end{macrocode}
% By default, use the |\ps@rmaa| page style.
%    \begin{macrocode}
\pagestyle{rmaa}
%    \end{macrocode}
% \end{macro}
% \end{macro}
%
% \subsection{Preamble commands}\label{sec:preamble} 
% Everything you need for setting the title page, TOC and index
% entries, etc.
% \subsubsection{Commands for the title page}
% \label{sec:title-commands}
% \begin{macro}{\resumen}
% \begin{macro}{\@resumen}
% \begin{macro}{\abstract}
% \begin{macro}{\@abstract}
% User commands for setting the resumen and abstract. 
% These should be used in the document header. The internal commands
% |\@resumen| and |\@abstract| are used to hold the actual text
% pending output via the |\maketitle| command. 
%    \begin{macrocode}
\newif\if@setresumen
\newcommand{\resumen}[1]{%
%%  \renewcommand{\@resumen}{{\selectlanguage{spanish}#1}}%
  \renewcommand{\@resumen}{{\ignorespaces#1}}%
  \@setresumentrue}
\renewcommand{\abstract}[1]{%
  \renewcommand{\@abstract}{{\ignorespaces#1}}
  \renewcommand\ADS@abstract{#1}
%% Not a good idea, since then author macros are not expanded and
%% cause errors later. However, by not doing it we are vulnerable to
%% fragile commands.....
  \if@noexpandADS\@onelevel@sanitize{\ADS@abstract}\fi
  \@abstractpresenttrue
  }
\newcommand{\ADSnoexpand}{\@noexpandADStrue}
%    \end{macrocode}
% \begin{macro}{\boldabstract}
% This is the version for doing 1-page conference proceedings. Still
% save the text for ADS but then write it out immediately as a bold
% paragraph. 
%    \begin{macrocode}
\newcommand{\boldabstract}[1]{%
  \gdef\ADS@abstract{#1}%
  \renewcommand{\@abstract}{#1}%
  \if@noexpandADS\@onelevel@sanitize{\ADS@abstract}\fi%
  \noindent\textbf{\boldmath #1}
  }
%    \end{macrocode}
% \end{macro}
% Default value in each case is to print out a warning if not set
%    \begin{macrocode}
  \newcommand{\@resumen}{%
    \@setresumenfalse \ClassWarningNoLine{rmaa}{No Spanish abstract
      (resumen) specified - if you do not speak Spanish, this will be
      translated by the editors}%
    \textbf{Favor de proporcionar un resumen en espa\~nol. If you are
      unable to translate your abstract into Spanish, the editors will
      do it for you.} }
  \newcommand{\@abstract}{\ClassWarningNoLine{rmaa}{No abstract
      specified}}
%    \end{macrocode}
% \end{macro}
% \end{macro}
% \end{macro}
% \end{macro}
%
% \begin{macro}{\title}
% \begin{macro}{\@rawtitle}
% \begin{macro}{\@TOCtitle}
% \begin{macro}{\rmaaIndexTitle}
% Similarly, |\title| is the preamble command to set the title of the
% article. This is used in various places, typeset in different
% styles, hence the use of the internal |\@rawtitle|, |\@title|,
% |\rmaaIndexTitle|, and |\@TOCtitle| commands, each of which uses the
% relevant style from section~\ref{sec:text-styles}
%    \begin{macrocode}
\newcommand{\@rawtitle}{Raw Title}
\newcommand{\@TOCtitle}{\rmaa@TOCTitleStyle{Title for TOC}}
\newcommand{\rmaaIndexTitle}{\rmaa@IndexTitleStyle{Title for index}.}
\newcommand{\@TOCaftertitle}{}
\renewcommand{\title}[1]{%
%%  {\ADS@sanitize\xdef\@rawtitle{#1}}
  \renewcommand\@rawtitle{\noexpand#1}
  \let\ADS@title=\@rawtitle
  \renewcommand{\@title}{\rmaa@TitleStyle{#1}}%
  \renewcommand{\@TOCtitle}{\rmaa@TOCTitleStyle{#1}\@TOCaftertitle}%
  \renewcommand{\rmaaIndexTitle}{\rmaa@IndexTitleStyle{#1\@TOCaftertitle}.}%
  }
%    \end{macrocode}
% \end{macro}
% \end{macro}
% \end{macro}
% \end{macro}
% 
% All the weird and whacky ways of indicating the author
% affiliations. 
% \begin{macro}{\affil}
% First, the way that sets it directly underneath the authors name.
% Note that |\affil| trumps all other affiliations for the purposes of
% getting into the ADS stuff (we use |\gdef|).   
% \dots 
%    \begin{macrocode}
\newcommand{\affil}[1]{%
  \\[1.5ex]\rmaa@AffilStyle{#1}\\[1.5ex]%
  \if@firstaffil\gdef\ADS@affil{#1}\@firstaffilfalse\fi%
  }
\renewcommand{\and}{and\\[1.5ex]}
%    \end{macrocode}
% \begin{macro}{\altaffil}
% Alternative affiliations are footnoted. 
%    \begin{macrocode}
\newcommand{\altaffil}[1]{%
  \if@firstaltaffil\gdef\ADS@affil{#1}\@firstaltaffilfalse\fi%
  \thanks{#1}}
%    \end{macrocode}
% \begin{macro}{\altaffilmark}
% \begin{macro}{\altaffiltext}
% Alternative-alternative mechanism for affiliations---should not be
% mixed with |\altaffil| (unless you know what you are doing) or there
% will be problems with numbering.
%    \begin{macrocode}
\newcommand{\altaffilmark}[1]{\textsuperscript{#1}}
\newcommand{\altaffiltext}[2]{%
  \if@firstaltaffil\gdef\ADS@affil{#2}\@firstaltaffilfalse\fi
  \stepcounter{footnote}%
  \protected@xdef\@thanks{\@thanks
    \protect\footnotetext[#1]{#2}}} 
%    \end{macrocode}
% \end{macro}
% \end{macro}
% \end{macro}
% \end{macro}
%
% \begin{macro}{\keywords}
% \begin{macro}{\@keywords}
% \begin{macro}{\@keywordheader}
% \label{keywords}Sundry minor preamble commands. All
% self-explanatory, I hope. |\keywords| is the old user interface
% (deprecated). New articles should use |\addkeyword| once for each
% keyword. That way the ADS record comes out right too (comma rather
% than em-dash separator).
% \changes{2.2beta12}{2001/12/07}{\cs{addkeyword} New user interface
%   for keywords.} 
% \changes{2.2beta13}{2002/03/11}{\cs{addkeyword} Improved behaviour
%   with .ads file. Use
%   \cs{ADS@sanitize} to change definitions of `~' and \cs{ion}.}
% \changes{3.23}{2007/09/09}{\cs{addkeyword} now signals a compilation
% error if the author tries to add more than 6 keywords.}
% \changes{3.24}{2007/09/14}{kwdcounter is now reset in
% \cs{@begofarticlehook} to prevent error in book style} 
%    \begin{macrocode}
\newcommand{\@kwdsep}{ --- }
\newcommand{\ADS@kwdsep}{, }
\newcounter{kwdcounter}\setcounter{kwdcounter}{0}
\newcommand{\addkeyword}[1]{%
  \addtocounter{kwdcounter}{1}
  \ifthenelse{\thekwdcounter > 6}{
    \ClassError{rmaa}{Too many keywords. A maximum of 6 keywords is allowed.}{}
  }{}
  \if@firstkwd%
  \renewcommand\@keywords{\if@boldkeywords\uppercase{#1}\else #1\fi}
  {\ADS@sanitize\def\ion{\plainion}\gdef\ADS@keywords{#1}}
  \@firstkwdfalse
  \else%
  \edef\@keywords{\@keywords\@kwdsep \if@boldkeywords\uppercase{#1}\else #1\fi}%
  {\ADS@sanitize\def\ion{\plainion}\xdef\ADS@keywords{\ADS@keywords\ADS@kwdsep #1}}%
  \fi}
\newcommand{\ADS@sanitize}{%
  \def\mathrm{}
  \def\@{\ }
  \def~{\ }
  \def\langle{<}
  \def\rangle{>}
  \def\log{log }
  \def\sin{sin }
  \def\cos{cos }
  \def\tan{tan }
  \def\fd{^d.}
  \def\fh{^h.}
  \def\fm{^m.}
  \def\fs{^s.}
  \def\fdg{�.}
  \def\farcm{'.}
  \def\farcs{''.}
  \def\arcdeg{�}
  \def\,{ }
  }
\newcommand{\keywords}[1]{%
  \renewcommand{\ADS@keywords}{#1}
  \renewcommand{\@keywords}{\uppercase{#1}}
  }
\newcommand{\@keywordheader}{\rmaa@KeywordHeaderStyle{Key Words:}}
\newcommand{\@keywords}{%
  \ClassWarningNoLine{rmaa}{No keywords specified}%
  Keywords go here
  }
%    \end{macrocode}
% \begin{macro}{\shorttitle}
% \begin{macro}{\shortauthor}
% \begin{macro}{\@shorttitle}
% \begin{macro}{\@shortauthor}
% Commands for setting the running headers.
%    \begin{macrocode}
\newcommand{\shorttitle}[1]{%
  \renewcommand\@shorttitle{\MakeUpperCase{#1}}}
\newcommand{\shortauthor}[1]{%
  \renewcommand\@shortauthor{\MakeUpperCase{#1}}}
\newcommand{\@shorttitle}{SHORT TITLE GOES HERE}
\newcommand{\@shortauthor}{SHORT AUTHOR GOES HERE}
%    \end{macrocode}
% \end{macro}
% \end{macro}
% \end{macro}
% \end{macro}
% \end{macro}
% \end{macro}
% \end{macro}


% \subsubsection{The full addresses}
% \label{sec:full-addresses}
% \begin{macro}{\fulladdresses}
% \begin{macro}{\@fulladdresses}
% \begin{macro}{\@fulladdheight}
% \begin{macro}{\@fulladddepth}
% \begin{macro}{\@fulladdbox}
% Aaghh, the full addresses! A perennial source of problems. 
%    \begin{macrocode}
\newlength{\@fulladdheight}
\newlength{\@fulladddepth}
\newsavebox{\@fulladdbox}
%% \fulladdresses command modified (11 May 2000) so that we can put
%% \fulladdresses{} in the article to suppress the output of the
%% addresses
\newcommand{\@FinalFigure}{\relax}
\newcommand{\FinalFigure}[1]{\renewcommand{\@FinalFigure}{#1}}
\newcommand{\fulladdresses}[1]{%
  \renewcommand{\@rmaaArg}{#1}
  \ifx\@rmaaArg\@rmaaEmpty
     \renewcommand{\@fulladdresses}{}
  \else
     \renewcommand{\@fulladdresses}{
       \normalsize%
       \ifx\@FinalFigure\relax\else\vfill\@FinalFigure\vfill\fi
%%       \vspace*{1ex}
       %% do addresses as a list (copied from thebibliography environment) 
       \list{\null}{\leftmargin 1.5em\topsep\z@\labelwidth\z@\labelsep\z@\itemindent -1.5em
         \itemsep\z@\parsep\z@\usecounter{enumi}}
       \sloppy\clubpenalty4000\widowpenalty4000
       \sfcode`\.=1000\relax
       #1
       \endlist
       }
  \fi
  }
\newcommand{\@fulladdresses}{Full addresses go here}
%    \end{macrocode}
% \begin{macro}{\suppressfulladdresses}
% New user-level command for suppressing full addresses should you so
% desire. Currently used in the single-page-poster style. 
%    \begin{macrocode}
\newcommand{\suppressfulladdresses}{\@suppressfulladdressestrue\@needaddressesfalse}
%    \end{macrocode}
% \end{macro}
% \end{macro}
% \end{macro}
% \end{macro}
% \end{macro}
% \end{macro}

% \begin{macro}{\printdates}
% \begin{macro}{\noprintdates}
% Command to activate the printing out of the received/accepted
% dates. This can also be done using the \texttt{printdates} option,
% which is the default for \texttt{journal}, but not for \texttt{proceedings}.
%    \begin{macrocode}
\newcommand{\printdates}{\@printdatestrue}
\newcommand{\noprintdates}{\@printdatesfalse}
%    \end{macrocode}
% \end{macro}
% \end{macro}
% \begin{macro}{\@outputaddressesifnecessary}
% Command to output the addresses if we haven't already done so
% This is called by: 
% \begin{enumerate}
% \item The end of |thebibliography| environment
% \item The |\outputfulladdresses| command
% \end{enumerate}
%    \begin{macrocode}
\newcommand{\@outputaddressesifnecessary}{%
  \ClassInfo{rmaa}{Calling @outputaddressesifnecessary}
  \ifthenelse{\boolean{@needaddresses}}{%
    %% If in one column style, or in the first column of a two column
    %% style then we can try to put it on this page 
    \ifthenelse{\boolean{@firstcolumn} \or \not\boolean{@rmtwocol}}{%
      \@outputaddresses}{%
      %% otherwise give warning it will be on a separate page
      \ClassWarning{rmaa}{Full addresses will go on a separate
        page (please see User Guide for how to remedy this)}}%
    %% If we haven't managed yet, put it on its own page
    \ifthenelse{\boolean{@needaddresses}}{%
      \if@debug{\ClassWarning{rmaa}{In this silly bit!}}\fi
      \clearpage\vfill\@outputaddresses}{}%
    }{
    %% If we don't need them, we still might need to do the eoa hook 
    \ifthenelse{\boolean{@suppressfulladdresses}}{%
      \if@debug{\ClassWarning{rmaa}{XXX}}\fi
%%      \if@book\else\@endofarticlehook\null\fi
      }{}
    }%
  }  
%    \end{macrocode}
% \end{macro}

% \begin{macro}{\outputfulladdresses}
% A user interface to the above command, in case there is no biblio
% (11 May 2000) 
%    \begin{macrocode}
\newcommand{\outputfulladdresses}{\@outputaddressesifnecessary}
%    \end{macrocode}
% \end{macro}

% \begin{macro}{\adjustfinalcols}
% This command starts a new column and should be 
% used to ensure that the columns balance on the last page. It also
% outputs the full 
% author addresses to ensure that this occurs in the left hand
% column. 
% The |\par| command is so that the current line will be taken account
% of when working out |\pagetotal|
%    \begin{macrocode}
\newcommand{\adjustfinalcols}{\if@twocolumn\goodbreak\@outputaddresses\newpage\noindent\fi}
%    \end{macrocode}
% \end{macro}

% \begin{macro}{\@outputaddresses}
% Command to do the work of outputting the addresses. 
% This can be called by 
% \begin{enumerate}
% \item an explicit |\adjustfinalcols|
% \item via |\@outputaddressesifnecessary|, q.v.
% \end{enumerate}
%
%    \begin{macrocode}
\newcommand{\@outputaddresses}{%
  % Allow option of suppressing the full addresses, in which case
  % we just run the hook if we are not in a book style
  \if@suppressfulladdresses
%%     \if@book\else\@endofarticlehook\null\fi
  \else
  %% save the full addresses in a box for output later
  \sbox{\@fulladdbox}{\parbox{\textwidth}{\@fulladdresses}}
  % calculate height of address box  
  \settoheight{\@fulladdheight}{\usebox{\@fulladdbox}}
  \settodepth{\@fulladddepth}{\usebox{\@fulladdbox}}
  \addtolength{\@fulladdheight}{\@fulladddepth}
  \if@debug
  { \typeout{Address total height: \the\@fulladdheight}
    \typeout{Text height: \the\pagetotal}
    \typeout{Goal height: \the\pagegoal}
    }
  \fi
  % reuse \@fulladdheight to mean the total height of text+addresses
  \addtolength{\@fulladdheight}{\pagetotal}
%%  \addtolength{\@fulladdheight}{-\baselineskip}
  \if@debug\typeout{Adjusted text + addresses height: \the\@fulladdheight}\fi
  % check it will fit on current page
  \ifdim\@fulladdheight<\pagegoal
     %% must shift it to the left
     %% to compensate for the hanging indentation in the reference list. 
     %% NOTE: this is not foolproof!
     %% NB. This will give overfull hbox warnings, which is a bit ugly.
     \vfill\noindent\if@biblist\hspace*{-\leftmargin}\fi
     \usebox{\@fulladdbox}
     % Set this to show we've done it
     \@needaddressesfalse
     % and assume we are at the end of the article
     % so that we do the EOA hook if not in the book style
%%     \if@book\else\@endofarticlehook\fi
  \else
     % if not, just do nothing. With luck, this function will get 
     % called again on the next page
     \if@debug{\typeout{Holding over full addresses to next page}}\fi
  \fi\fi
  }
%    \end{macrocode}
% \end{macro}

% \subsubsection{Commands for editors}
% \label{sec:editor-commands}
% Some user/editor commands for doing the TOC and index in the book
% option. 
% They are defined here rather in the book-specific part since they
% need to be always available in order to avoid breaking documents. 
% \begin{macro}{\listofauthors}
% \begin{macro}{\@rawauthors}
% \begin{macro}{\rmaaTOCauthors}
% The command |\listofauthors| formats an author list for use in the
% TOC. It should contain the same names as in the |\author| command,
% but without the affiliation garbage and using |\&| instead of
% |and|. 
%    \begin{macrocode}
\newcommand{\@rawauthors}{Raw Authors}
\newcommand{\rmaaTOCauthors}{\noexpand\rmaaTOCAuthorStyle{Author list for TOC}}
\newcommand{\listofauthors}[1]{%
  \renewcommand{\@rawauthors}{#1}
  \renewcommand\rmaaTOCauthors{\noexpand\rmaaTOCAuthorStyle{#1}}}
%    \end{macrocode}
% \begin{macro}{\indexauthor}\label{indexauthor}
% The command |\indexauthor| produces an index entry for each
% author. 
% This command must come \emph{after} the
% |\title| and the |\listofauthors| commands. Also, there should be one 
% |\indexauthor| command for \emph{each} author. We also add the
% author onto the end of the |\ADS@authors| macro and construct the
% |\ADS@authorinitial| by a most cunning method (see below), but only
% for the first author. 
% 
% No, No, NO! This was all wrong. Backed out completely: 
% \begin{quote}
% Completely rewritten (19 Dec 2002) since we now use the index class, so 
% as to be able to alphabetize accented characters. A side effect of this is
% that |\index| no longer expands out the argument after the |@|, so we 
% instead add all the |\index| commands to a list (|\indexauthor@deferlist|), 
% which is executed in the |\@begofarticlehook|. 
% 
% Note the stuff with |\toks0| is copied from Exercise 20.15 of the \TeX{}book---it is
% to make sure the latin1-accented name gets passed into the idx
% file. 
% \end{quote}
% We now go back to the old way of doing things, using an optional
% argument to the \cs{indexauthor} command that has the name without
% accents. 
% \changes{3.08}{2003/01/16}{Remerged changes from 3.06 that got
% dropped from 3.07 after a synchronization disaster. Basically, the
% accentless optional argument to \cs{indexauthor}. Must add the macro
% directories to my unison set-up to stop this sort of thing from happening.}
%    \begin{macrocode}
\newcommand\@isortkey{}
\newcommand\@emptykey{empty}
\newcommand\indexauthor[2][empty]{
  \renewcommand\@isortkey{#1}
  \ifx\@isortkey\@emptykey \renewcommand\@isortkey{#2}\fi
  \index{\@isortkey @{{\rmaaIndexAuthorStyle{#2}}\hskip1em plus1em minus0.4em
      \rmaaIndexTitle\hskip1em plus2em minus0.4em\penalty-100
      \rmaaTOCauthors}}
  \indexauthorADS{#2}
}
%%\def\indexauthor@deferlist{\relax}
%%\long\def\indexauthor#1{
%%  \def\@@name{#1} \toks0=\expandafter{\@@name}
%%  \toks1=\expandafter{\indexauthor@deferlist}
%%  \protected@xdef\indexauthor@deferlist{%
%%    \the\toks1
%%    \noexpand\index{\the\toks0@{\noexpand\rmaaIndexAuthorStyle{\the\toks0}\noexpand\hskip1em plus1em minus0.4em 
%%        \rmaaIndexTitle\noexpand\hskip1em plus2em minus0.4em\noexpand\penalty-100
%%        \rmaaTOCauthors}}
%%  }
%%  \indexauthorADS{#1}
%%  }
%%\xdef\indexauthor@deferlist{}
\newcommand\indexauthorADS[1]{%
  \if@firstauthor\protected@xdef\ADS@authorinitial{\firstletter{#1}}\fi
  \protected@xdef\ADS@authors{\ADS@authors\if@firstauthor\else ; \fi
    #1}
  \@firstauthorfalse
}
%    \end{macrocode}
% Note that the first |#1| sets the key for the index, making sure it
% actually starts with  
% first letter of the surname. Meanwhile, the |@{...}| expression,
% which starts with some  
% control sequences is what is written to the |.ind| file. 
% See the makeindex man page for details. 
% \end{macro}
% \end{macro}
% \end{macro}
% \end{macro}

% \begin{macro}{\KillAllAccents}
% \begin{macro}{\firstletter}
% These are utility functions used in the above. |lett| is a counter
% indicating the current letter in the string. |\firstletter| just
% prints the first letter of its argument.
%    \begin{macrocode}
\newcommand{\KillAllAccents}{\def\'{}\def\`{}\def\~{}\def\.{}\def\i{i}}
\newcommand\firstletter[1]{\@car #1\@nil}
%    \end{macrocode}
% \end{macro}
% \end{macro}
% \begin{macro}{\SetVolume}
% \begin{macro}{\SetMSnumber}
% \begin{macro}{\SetFirstPage}
% \begin{macro}{\SetYear}
% \begin{macro}{\ReceivedDate}
% \begin{macro}{\AcceptedDate}
% These are for the editors to set the volume number, etc. 
%    \begin{macrocode}
\newcommand{\SetVolume}[1]{%
  \renewcommand{\@volume}{#1}
  \edef\ADS@volume{\PadToFourWithDots{#1}}
  \renewcommand{\@volumen}{\rmaa@VolumenStyle{volumen #1}}%
  }
\newcommand{\SetMSnumber}[1]{\renewcommand{\@MSnum}{#1}}
%    \end{macrocode}
% Note that there is no real reason to use |\SetFirstPage| since all
% the page numbering is done automatically in the book style. 
%    \begin{macrocode}
\newcommand{\SetFirstPage}[1]{\setcounter{page}{#1}}
\newcommand{\SetYear}[1]{%
  \renewcommand{\@year}{#1}
  \renewcommand{\ADS@year}{#1}
  \renewcommand{\@anyo}{\rmaa@VolumenStyle{#1}}
  }
\newcommand{\ReceivedDate}[1]{\renewcommand{\@received}{#1}}
\newcommand{\AcceptedDate}[1]{\renewcommand{\@accepted}{#1}}
%    \end{macrocode}
%
% \begin{macro}{\DOI}
% User/editor command to enter the DOI, which will be typeset as a link.
% \changes{4.0}{2019/11/02}{New user command \cs{DOI} to typeset
% DOI in the header of first page and add it to ADS info.}
%    \begin{macrocode}
\newcommand\DOI[1]{%
  \def\@DOI@url{\url{https://doi.org/10.22201/ia.01851101p.#1}}%
  \def\ADS@doi{https://doi.org/10.22201/ia.01851101p.#1}%
}%
%    \end{macrocode}
% \end{macro}
%
% \begin{macro}{\ORCIDiD}
% \changes{4.0}{2019/11/12}{Add mechanism for giving ORCID iD links}
% \begin{macro}{\ORCIDiDfull}
% \begin{macro}{\@orcid@icon}
% \begin{macro}{\@orcid@iconbig}
% Commands to add icon links to ORCID iDs of authors. For example
% \begin{quote}
% |\ORCIDiD|\marg{0000-0000-0000-0000}
% \end{quote}
% to be used after each author name in the |\author|
% command and in the |\fulladdresses|.  Alternatively,
% |\OORCIDiDfull| will typeset the icon plus the full text of the
% link, but we don't use that at present.  The icon image is stored in |ORCIDiD_iconvector.pdf|, which needs to be on the input path.  In principle, we could make EPS versions if necessary.
%    \begin{macrocode}
\newcommand\@orcid@icon{\includegraphics[height=1.8ex]{ORCIDiD_iconvector}}
\newcommand\@orcid@iconbig{\includegraphics[height=2ex]{ORCIDiD_iconvector}}
\newcommand\ORCIDiD[1]{\href{https://orcid.org/#1}{\,\@orcid@icon}}
\newcommand\ORCIDiDfull[1]{%
  \href{https://orcid.org/#1}{\@orcid@iconbig}\,\url{https://orcid.org/#1}}
%    \end{macrocode}
% \end{macro}
% \end{macro}
% \end{macro}
% \end{macro}
%
% \begin{macro}{\@volume}
% \begin{macro}{\@volumen}
% \begin{macro}{\@artpages}
% \begin{macro}{\@year}
% \begin{macro}{\@anyo}
% \begin{macro}{\@received}
% \begin{macro}{\@accepted}
% \begin{macro}{\@MSnum}
% Default values for the containers used by above commands, plus a
% few more useful ones. 
%    \begin{macrocode}
\newcommand{\@volume}{00}
\newcommand{\@volumen}{\rmaa@VolumenStyle{00}}
\newcommand{\@artpages}{%
  \oldpageref{\articlename:RMAAFirstPage}--\oldpageref{\articlename:RMAALastPage}}
\newcommand{\@year}{\the\year}
\newcommand{\@anyo}{\rmaa@VolumenStyle{2007}}
\newcommand{\@journalname}{\rmaa@JournalNameStyle{%
    \if@proceedings\RMAASClong\else\RMAAlong\fi}}
\newcommand{\@received}{}
\newcommand{\@accepted}{}
\newcommand{\@MSnum}{0000}
\newcommand{\@DOI@prefix}{https://doi.org/10.22201/ia.01851101p}
\def\@DOI@url{\url{https://doi.org/10.22201/ia.01851101p.2099.00.00.00}}
%    \end{macrocode}
% \end{macro}
% \end{macro}
% \end{macro}
% \end{macro}
% \end{macro}
% \end{macro}
% \end{macro}
% \end{macro}
% \end{macro}
% \end{macro}
% \end{macro}
% \end{macro}
% \end{macro}
% \end{macro}
% 
% The following commands are most useful in the \texttt{book}
% style. However, the |\SetConfTitle| command can now be used in
% individual articles to set the header for the \texttt{preprint}
% style. 
%    \begin{macrocode}
\newcommand{\@reunionplace}{Lugar de la reuni\'on}
\newcommand{\@reuniondate}{Fecha de la reuni\'on}
\newcommand{\@conftitle}{Conference Title}
\newcommand{\SetReunion}[2]{\renewcommand{\@reunionplace}{#1}\renewcommand{\@reuniondate}{#2}}
\newcommand{\SetConfTitle}[1]{\renewcommand{\@conftitle}{#1}}
\newcommand{\@editors}{Juan Ram\'on de la Fuente and ``El Mosh''}
\newcommand{\SetEditors}[1]{\renewcommand{\@editors}{#1}}
\newcommand{\@editores}{Juan Ram\'on de la Fuente y ``El Mosh''}
\newcommand{\SetEditores}[1]{\renewcommand{\@editores}{#1}}
\newcommand{\@mes}{mes de la publicaci\'on}
\newcommand{\SetMes}[1]{\renewcommand{\@mes}{\rmaa@VolumenStyle{#1}}}
%    \end{macrocode}

% \subsection{Notes for Author Proofs and Communication between Editors}
% \label{sec:notes}
% We introduce a command \cs{Note} that should be used to describe
% any change made to the manuscript. These are a bit tricky to get
% right. In general we use \cs{marginpar} but this only works in outer
% paragraph mode, so we just bung the box where we are if we are in a
% figure caption or something. If we are in the preamble, then we just
% use a \marginpar as part of the \maketitle command.
%    \begin{macrocode}
\newif\if@arenotes
\@arenotesfalse
\newcounter{notecounter}
\setcounter{notecounter}{0}
\newcommand\NotesList{}
\newcommand\Note[1]{%
  \if@proof{%
    \@arenotestrue%
    \stepcounter{notecounter}%
    \if@inpreamble\else{%
      \NoteMark{\arabic{notecounter}}%
    }\fi
    \protected@xdef\NotesList{\NotesList 
      \protect\item #1}%
  }\fi%
}
\newcommand\Changed[2]{Changed ``#1'' to ``#2''.}
\newcommand\NoteMark[1]{%
  \ifinner
  \raisebox{1ex}[0pt][0pt]{\makebox[0pt][l]{\framebox{\footnotesize\bfseries
        N#1}}}%
  \else
  \marginpar{\framebox{\footnotesize\bfseries N#1}}%
  \fi
}
\newcommand\RefNote[1]{\if@proof\NoteMark{\ref{#1}}\relax\fi}
\newcommand\NoteLabel{\protect\label}
\newcommand\OutputNotes{%
  \bgroup
  \onecolumn
  \section*{EDITORIAL NOTES}
  \begin{enumerate}
    \NotesList
  \end{enumerate}
  \egroup
}
\DeclareRobustCommand\cs[1]{\texttt{\char`\\#1}}

%    \end{macrocode}
% \subsection{Bibliographic Records for ADS}
% \label{sec:ADS}
% \changes{2.2beta11}{2001/11/27}{Just about finished all the ADS
% stuff now. Except for the references, which I have given up on for
% the time being.}
% \begin{macro}{\ADS@BibCode}
% \begin{macro}{\ADS@year}
% \begin{macro}{\ADS@journal}
% \begin{macro}{\ADS@volume}
% \begin{macro}{\ADS@firstpage}
% \begin{macro}{\ADS@authorinitial}
% We use the |\ADS@BibCode| command to hold the bibliographic code for
% ADS. It is 
% used in the |%R| 
% field of the bibliographic record for each article. 
% The format is 
% \begin{quote}
% \texttt{YYYYJJJJJVVVVMPPPPA}  
% \end{quote}
% where \texttt{YYYY} is the year, \texttt{JJJJJ} is the journal code,
% which is \texttt{RMxAA} or \texttt{RMxAC}, \texttt{VVVV} is the
% volume number (right-justified), \texttt{M} is not used by us,
% \texttt{PPPP} is the first page number (right-justified), and
% \texttt{A} is the initial of the first author's surname. The volume
% and page fields are padded with periods (`\texttt{.}'), which causes
% us some fun (this is done with the macro |\PadToFourWithDots| described
% below) . 
% 
% All these individual bits are held in the commands |\ADS@year|,
% |\ADS@journal|, |\ADS@volume|, |\ADS@firstpage|,
% |\ADS@authorinitial|. 
%    \begin{macrocode}
\newcommand{\SetADSBibCode}[1]{\renewcommand{\ADS@BibCode}{#1}}
\newcommand{\ADS@BibCode}
{\ADS@year\ADS@journal\ADS@volume.\PadToFourWithDots{\ADS@firstpage}\ADS@authorinitial}
\newcommand{\ADS@year}{2001} 
\newcommand{\ADS@journal}{\if@proceedings RMxAC\else RMxAA\fi}
\newcommand{\ADS@volume}{....} 
%    \end{macrocode}
% These two are just set to default values here, |\ADS@firstpage| is set
% in the |document| environment for included articles, in the
% |frontsection| environment, |\tableofcontents| command,  |index| and
% |abstracts| environments. Also by |\AtBeginDocument| for stand-alone
% articles. 
% 
% God knows how we are going to do |\ADS@authorinitial|!
%    \begin{macrocode}
\newcommand{\ADS@firstpage}{....}
\newcommand{\ADS@authorinitial}{A}
%    \end{macrocode}
% \end{macro}
% \end{macro}
% \end{macro}
% \end{macro}
% \end{macro}
% \end{macro}
% We also have a whole slew of other fields that can be set. These are
% all listed in Table~\ref{tab:ADS}. Most of them still have to be
% defined. 
% \begin{table}[!ht]\raggedleft
% \caption{ADS Bibliographic Records}\label{tab:ADS}
% \begin{tabular}{lp{8em}lp{8em}}\hline
% Code & Description & Macro Name & Comment \\ \hline
% |%R| & Bibliographic Code & |\ADS@BibCode| & required \\
% |%A| & Author list & |\ADS@authors| & required, generate from |\indexauthor| \\
% |%F| & Author affiliation & |\ADS@affil| & format: AA(\dots)
% AB(\dots) \dots \\
% |%J| & Journal Name & |\ADS@journalname| &  \\
% |%D| & Publication Date & |\ADS@fulldate| &  \\
% |%L| & Last Page & |\ADS@lastpage| & generate from label \\
% |%T| & Title & |\ADS@title| & required, generate from |\title| \\
% |%C| & Abstract Copyright & |\ADS@copyright| & set to IA-UNAM \\
% |%U| & URL for electronic document & |\ADS@url| & set to PDF file \\
% |%d| & DOI for article & |\ADS@doi| & set by |\DOI| \\
% |%K| & Keywords & |\ADS@keywords| & set equal to |\@keywords| \\
% |%A| & Abstract Text & |\ADS@abstract| & set equal to |\@abstract| \\
% |%R| & References & |\ADS@references| & generate from the
% bibliography  \\ \hline
% \end{tabular}
% \end{table}
% A lot of these macros are set by other commands, as mentioned in the
% Table, but some need to be set  explicitly here. 
% 
% \begin{macro}{\ADS@doi}
% \changes{4.0}{2019/11/08}{Save DOI for ADS}
% The user/editor command \cs{DOI} is used to set this. 
%    \begin{macrocode}
\newcommand{\ADS@doi}{https://doi.org/10.22201/ia.01851101p.0000.00.00.00}
%    \end{macrocode}
% \end{macro}
% 
% \begin{macro}{\ADS@journalname}
% This
% does not contain any accents at the moment so as not to cause
% problems---this maybe should be changed though. 
%    \begin{macrocode}
\newcommand{\ADS@journalname}{%
  \if@proceedings\@conftitle\space (Eds. \@editors) \fi%
  Revista Mexicana de Astronom\'\i{}a y Astrof\'\i{}sica%
  \if@proceedings\space (Serie de Conferencias)\fi%
  \space Vol. \@volume, pp. \ADS@firstpage--\ADS@lastpage\space (\@year)%
  \space (http://www.astroscu.unam.mx/\~{}rmaa/)
  }
%    \end{macrocode}
% \end{macro}
% \begin{macro}{\ADS@keywords}
% \begin{macro}{\ADS@abstract}
% \begin{macro}{\ADS@title}
% |\ADS@abstract| is straightforward, since we already have the text
% saved as |\@abstract|. We could do the same with |\ADS@keywords| but
% to avoid unwanted formatting we explicitly set it in |\keywords|
% (p.~\pageref{keywords}). |\ADS@title| is just set to |\@rawtitle|,
% which we were already using for another purpose. 
%    \begin{macrocode}
\let\ADS@title=\@rawtitle
\newcommand{\ADS@keywords}{No keywords specified}
\newcommand{\ADS@abstract}{No abstract specified}
%    \end{macrocode}
% \end{macro}\end{macro}\end{macro}
% \begin{macro}{\ADS@authors}
% |\ADS@authors| wants to be in ``Other, A. N.'' format, but with
% semi-colons between authors. The most natural way to do this seems to
% be with |\indexauthor| (p.~\pageref{indexauthor}). We append each
% author to this macro, which means we must explicitly clear it for
% each new article (done in |\article|). .   
%    \begin{macrocode}
\newcommand{\ADS@authors}{}
%    \end{macrocode}
% \end{macro}
% \begin{macro}{\ADS@fulldate}
% \begin{macro}{\ADS@lastpage}
% \begin{macro}{\ADS@copyright}
% \begin{macro}{\ADS@website}
% \begin{macro}{\ADS@url}
% \begin{macro}{\SetWebsite}
% \begin{macro}{\SetMonthNumeric}
% Other sundry data, together with user commands for bits that can be
% set by the Editors. The |\ADS@url| is constructed from the website
% (currently |http://www.astrosmo.unam.mx/rmaa/|,
% the journal descriptor (e.g., |RMxAC..10/|), followed by the
% path to the PDF file (e.g., |PDF/RMxAC..10_HEW1.pdf|). 
% 
%    \begin{macrocode}
\newcommand{\ADS@website}{http://www.astrosmo.unam.mx/rmaa/}
\newcommand{\ADS@url}
{\ADS@website\ADS@journal\ADS@volume\ADS@number/PDF/\ADS@journal\ADS@volume\ADS@number\_\articlename.pdf}
\newcommand{\ADS@number}{}
\newcommand{\ADS@fulldate}{00/\ADS@year} 
\newcommand{\ADS@lastpage}{00} 
\newcommand{\ADS@copyright}{(C) \ADS@year; Instituto de Astronom\'\i{}a,
  Universidad Nacional Aut\'onoma de M\'exico} 
\newcommand{\SetNumber}[1]{\renewcommand{\ADS@number}{-#1}}
\newcommand{\SetWebsite}[1]{\renewcommand{\ADS@website}{#1}}
\newcommand{\SetMonthNumeric}[1]{\renewcommand{\ADS@fulldate}{#1/\ADS@year}}
%    \end{macrocode}
% \end{macro}\end{macro}\end{macro}\end{macro}\end{macro}\end{macro}\end{macro}
% \begin{macro}{\ADS@affil}
% \begin{macro}{\ADS@references}
% These two are somewhat more problematical. Affiliation currently
% only works for the first author (see \dots) pending redesign of the
% user interface. References are now done by wrapping the
% |thebibliography| environment to open another input stream on the
% same file and search through for the |\bibitems|. 
%    \begin{macrocode}
\newcommand{\ADS@references}{}
\newcommand{\ADS@affil}{}
%    \end{macrocode}
% \end{macro}\end{macro}
% 
% Define two macros that expand to the strings
% |\begin{thebibliography}| and |\end{thebibliography}|. Note that the
% idea for this and much else is inspired by the code for the docstrip
% package. Also define |\StartBibX| in case authors use the more
% standard |\begin{thebibliography}{}| syntax. 
%    \begin{macrocode}
\edef\EndBib{\expandafter\@gobble\string\\end\string{thebibliography\string}}
\edef\StartBib{\expandafter\@gobble\string\\begin\string{thebibliography\string}}
\edef\StartBibX{\expandafter\@gobble\string\\begin\string{thebibliography\string}\string{\string}}
%    \end{macrocode}
% Define two new logical variables. |\ifAnotherline| is true if we
% want to read another line from the input stream. It is set to false
% when we encounter the |\end{thebibliography}|. |\ifInBib| is set to
% true when we frind the |\begin{thebibliography}| and tells us to
% start saving the following lines. 
%    \begin{macrocode}
\newif\ifAnotherLine
\newif\ifInBib
%    \end{macrocode}
% \begin{macro}{makeOther}
% We need a macro to change the catcodes of lots of things to other. 
%    \begin{macrocode}
\def\makeOther#1{\catcode`#1=12\relax}
%    \end{macrocode}
% \end{macro}
% \begin{macro}{\RereadBiblio}
% Define a macro to do all the fiddly stuff. This is used at the
% beginning of |thebibliography| environment (Sec~\ref{sec:biblio}).
%    \begin{macrocode}
\newread\tf@biblio
\newcommand{\RereadBiblio}{%
%    \end{macrocode}
% Open the current file again for reading
%    \begin{macrocode}
  \begingroup
  \immediate\openin\tf@biblio=\articlename
%    \end{macrocode}
% Everything in a group to be safe. Initialize logicals. Black magic
% to change
% catcodes and avoid blank lines being interpreted as |\par|. 
%    \begin{macrocode}
  \AnotherLinetrue\InBibfalse
  \makeOther\ \makeOther\\\makeOther\$%
  \makeOther\#\makeOther\^\makeOther\^^K%
  \makeOther\_\makeOther\^^A\makeOther\%%
  \makeOther\~\makeOther\{\makeOther\}\makeOther\&%
  \endlinechar-1\relax
%    \end{macrocode}
% Start looping through the file. Read each line to the macro
% |\Line|. If we have got to EOF or end of biblio then set
% |\AnotherLinefalse|.
%    \begin{macrocode}
  \loop
  \read\tf@biblio to \Line
  \ifx\Line\EndBib\AnotherLinefalse\InBibfalse\fi
  \ifeof\tf@biblio \AnotherLinefalse\fi
%    \end{macrocode}
% If we are in the biblio and the line is not blank then append it to
% the |\ADS@references| macro together with a newline. 
%    \begin{macrocode}
  \ifInBib
  \ifx\Line\empty\else\xdef\ADS@references{\ADS@references\Line^^J}\fi
  \fi
%    \end{macrocode}
% Check if the line we just saw was the start of biblio, then repeat
% the loop if necessary. 
%    \begin{macrocode}
  \ifx\Line\StartBib\InBibtrue\fi
  \ifx\Line\StartBibX\InBibtrue\fi
  \ifAnotherLine\repeat
  \immediate\closein\tf@biblio
  \endgroup
  }
%    \end{macrocode}
% \end{macro}
% 
% \begin{macro}{\PadToFourWithDots}
% Utility function for left-padding a number to 4 characters with
% periods. Used in 
% constructing |\ADS@BibCode|. This is a very simple algorithm and
% \emph{requires} that the argument actually be a number---if it is
% not then dire consequences will result.  
%    \begin{macrocode}
\newcommand{\PadToFourWithDots}[1]{%
  \ifnum #1 > 999 #1%
  \else\ifnum #1 > 99 .#1%
  \else\ifnum #1 > 9 ..#1%
  \else ...#1%
  \fi\fi\fi}
%%  \ifnum \count0 > 999 #1%
%%  \else\ifnum \count0 > 99 .#1%
%%  \else\ifnum \count0 > 9 ..#1%
%%  \else ...#1%
%%  \fi\fi\fi}
%    \end{macrocode}
% \end{macro}
%
% \begin{macro}{\AuthorADS}
% \begin{macro}{\ADS@affiliations}
% New mechanism for setting the affiliations for ADS. The author labels (AA, AB, \dots) are generated using two nested counters.  We use the |\robustify| command from \texttt{etoolbox} so that accented characters in the affiliation addresses do not blow up.  The list of accents may need to be expanded. 
%    \begin{macrocode}
\newcounter{AUTHOR}
\newcounter{author}[AUTHOR]
\stepcounter{AUTHOR}
\newcommand\ADS@affiliations{}
\robustify{\'}\robustify{\i}\robustify{\~}\robustify{\`}
\def\AuthorADS#1#2{%
  \begingroup
  \def\ORCIDiD##1{<ID system="ORCID">##1</ID>}
  \def\Email##1{<EMAIL>##1</EMAIL>}
  \def\Affil##1{##1}
  \stepcounter{author}
  \ifthenelse{\value{author} > 26}{\stepcounter{AUTHOR}\stepcounter{author}}{}
  \ifthenelse{\equal{\ADS@affiliations}{\@rmaaEmpty}}{%
    \xdef\ADS@affiliations{\Alph{AUTHOR}\Alph{author}(#2)}
  }{%
    \xdef\ADS@affiliations{\ADS@affiliations, \Alph{AUTHOR}\Alph{author}(#2)}
  }%
  \endgroup
}
%    \end{macrocode}
% \end{macro}
% \end{macro}
% %
% \begin{macro}{\ADS@dumpdata}
% \changes{4.0}{2019/11/13}{Write all author affiliations for ADS}
% \begin{macro}{\ADS@write}
% \begin{macro}{\ADS@openfile}
% |\ADS@dumpdata| is the macro that actually does the job of writing out all the
% data to an auxiliary file, using |\ADS@write| to write each
% field. The file is set up with |\ADS@openfile|, which currently is
% done in the TOC of a book.  The new multi-author affiliations are used if |\ADS@affiliations| has been set, otherwise it falls back to the single-author version |\ADS@affil|. 
%    \begin{macrocode}
\newcommand{\ADSfield}[2]{\relax}
\newcommand{\ADSvfield}[2]{\relax}
%%\renewcommand{\ADSfield}[12]{\%#1 #2}
\newcommand{\ADS@dumpdata}
{ 
  \ADS@write{\string\begin{ADSarticle}}
  \ADS@WRITE{\string\ADSfield{R}{\ADS@BibCode}}
  \ADS@write{\string\ADSfield{A}{\ADS@authors}}
  \ifthenelse{\equal{\ADS@affiliations}{\@rmaaEmpty}}{%
    {\ADS@sanitize\ADS@write{\string\ADSfield{F}{AA(\ADS@affil)}}}
  }{%
    {\ADS@sanitize\ADS@write{\string\ADSfield{F}{\ADS@affiliations}}}
  }
  \ADS@write{\string\ADSfield{J}{\ADS@journalname}}
  \ADS@write{\string\ADSfield{D}{\ADS@fulldate}}
  \ADS@write{\string\ADSfield{L}{\ADS@lastpage}}
  {\ADS@sanitize\ADS@write{\string\ADSfield{T}{\ADS@title}}}
  \ADS@write{\string\ADSfield{C}{\ADS@copyright}}
  \ADS@write{\string\ADSfield{d}{\ADS@doi}}
  \ADS@write{\string\ADSfield{U}{\ADS@url}}
  \ADS@write{\string\ADSfield{K}{\ADS@keywords}}
  {\ADS@sanitize\ADS@write{\string\ADSfield{B}{\ADS@abstract}}}
  \ADS@write{\string\ADSfield{Z}{\ADS@references}}
  \ADS@write{\string\end{ADSarticle}}
  }
%%\newcommand{\ADS@write}[1]{\if@book\addtocontents{ads}{#1}\fi}
%%\newcommand{\ADS@openfile}{\@starttoc{ads}}
\newcommand{\ADS@write}[1]{\set@display@protect\immediate\write\tf@ads{#1}}
\newcommand{\ADS@WRITE}[1]{\immediate\write\tf@ads{#1}}
\newcommand{\ADS@openfile}{%
  \newwrite\tf@ads%
  \immediate\openout\tf@ads\jobname.ads\relax
  \ADS@WRITE{\ADS@preamble}}
\newcommand{\ADS@closefile}{\ADS@WRITE{\ADS@postamble}}
\def\ADS@preamble{%
  %%  \string\documentclass{article}
  %%  \string\begin{document}
  }
\def\ADS@postamble{%
  %%  \string\end{document}
  \string%\string% Local Variables:^J 
  \string%\string% mode: LaTeX^J
  \string%\string% mode: Font-lock^J
  \string%\string% End:^J
  }
%    \end{macrocode}
% \end{macro}
% \end{macro}
% \end{macro}
% 
% \subsection{Automatic calculation of the page range}
% \label{sec:page-range}
% Another issue that has cost me a lot of grief. Still, it finally
% seems to be sorted now. All hail Jeff Goldberg. 
%
% Cunning trick to get the page numbers from |\pageref|.  Put a label
% on the first and last page of the document using |\AtBeginDocument|
% and |\AtEndDocument|. .  This way won't work for the book option.
% In that case, we do it by defining |\@endofarticlehook| and
% |\@begofarticlehook| as parallels to |\AtEndDocument| and
% |\AtBeginDocument|.  Note that we have to defer adding
% |\@begofarticlehook| to |\AtBeginDocument| to make sure it comes
% after the redefinitions of |\ref| and |\label|.
%    \begin{macrocode}
\if@book\AtEndDocument{\ADS@closefile}\else
 %\AtBeginDocument{\@begofarticlehook}
\AtEndDocument{\@endofarticlehook}
\fi
%    \end{macrocode}
% \begin{macro}{\@begofarticlehook}
% Since we are getting more and more things that need to be done at
% the beginning of every article, frontsection, etc we define a sister
% hook to 
% |\@endofarticlehook|. This hook does the following: 
% \begin{enumerate}
% \item Labels the first page
% \item Sets |\ADS@firstpage|
% \item runs and then clears |\indexauthor@deferlist|
% \item resets the keyword counter
% \end{enumerate}
% It is executed in the following circumstances: 
% \begin{enumerate}
% \item At |\begin{document}| for standalone files
% \item At |\begin{document}| for each included file
% \item At beginning of front sections 
% \item At beginning of TOC
% \item At beginning of Author Index
% \item At beginning of abstracts
% \end{enumerate}
% For some reason, this is not always within the right group, so we
% use |\xdef| just to be sure. 
%    \begin{macrocode}
\newcommand{\@begofarticlehook}{%
  \if@frontmatter\frontlabel{RMAAFirstPage}\else\label{RMAAFirstPage}\fi%
  \xdef\ADS@firstpage{\@arabic{\c@page}}%
  \if@book\else\makeindex\fi % needed so default index is defined in stand-alone articles
%%  \indexauthor@deferlist
%%  \xdef\indexauthor@deferlist{}
  \@inpreamblefalse
  \setcounter{kwdcounter}{0}
  }
%    \end{macrocode}
% \end{macro}
% \begin{macro}{\@endofarticlehook}
% This is a hook to be run at the end of each article
% currently, it just labels the last page and writes out all the ADS
% data to an auxiliary file.  
% It should not produce any printed output. 
% It is modelled on the |\AtEndDocument| command, to which it reduces
% in the case of a single article. 
%    \begin{macrocode}
\newcommand{\@endofarticlehook}{%
  \if@debug{\ClassWarning{rmaa}{***Doing Hook***}}\fi
  \if@debug{\ClassWarning{rmaa}{ADSabstract is: \ADS@abstract}}\fi
  \clearpage\rmaa@lastpagelabel\if@book\ADS@dumpdata\fi
  \ifx\NotesList\@rmaaEmpty\else\OutputNotes\fi
  \if@debug{\ClassWarning{rmaa}{Notes list is: \NotesList}}\fi
}
%    \end{macrocode}
% \begin{macro}{\rmaa@lastpagelabel}
% New version of last page label command copied from Jeff Goldberg's
% lastpage.sty package - we can't just use |\label| here in case the
% output routine for the last page has already been called. Instead,
% we first explicitly clear the page, step back the page counter,
% explicitly write to the .aux file, then restep the page counter
% again.
%
% But this will conflict with the hyperref package, which does labels differently.  Compensate for that by using the five-argument form.
% 
% Also set the |\ADS@LastPage| macro to the appropriate number. 
%    \begin{macrocode}
\newcommand{\rmaa@lastpagelabel}{%
  \addtocounter{page}{-1}%
  \immediate\write\@auxout{\string
    \newlabel{\articlename:RMAALastPage}{{}{\thepage}{}{}{}}}%
  \edef\ADS@lastpage{\thepage}%
  \addtocounter{page}{1}}
%    \end{macrocode}
% \end{macro}
% \end{macro}

% \subsection{Deal with explicit change in number of columns}
% \label{sec:columns}
% Sometimes, when formatting a book, we need to switch between
% one-column articles and two-column articles. For this to work, we
% need to wrap the standard \LaTeX{} commands to set the |\@rmtwocol|
% flag appropriately and also reset the package default values of all
% scalable lengths. .  
%    \begin{macrocode}
\let\latexonecolumn=\onecolumn
\let\latextwocolumn=\twocolumn
\renewcommand{\onecolumn}{\@rmtwocolfalse\@package@default@lengths\latexonecolumn}
\renewcommand{\twocolumn}{\@rmtwocoltrue\@package@default@lengths\latextwocolumn}
%    \end{macrocode}
% 
% \subsection{The actual formatting of the title}
% \label{sec:title-format}

% Note, we have to redefine the |\maketitle| command to use the
% numbered footnotes. This means we needn't use the atbeginend stuff. 
%    \begin{macrocode}
\renewcommand{\maketitle}{\par
  \begingroup
%    \end{macrocode}
% Don't switch to marked footnotes. 
% This is consistent with existing style but I'm not sure I like it
%    \begin{macrocode}
  %%  \renewcommand\thefootnote{\fnsymbol{footnote}}%
%%  \def\@makefnmark{\hbox to\z@{$\m@th^{\@thefnmark}$\hss}}%
  \def\@makefnmark{\textsuperscript{\@thefnmark}\hspace{-0.2em}}%
  \long\def\@makefntext##1{\parindent 1em\noindent
    \hbox to1.8em{\hss$\m@th^{\@thefnmark}$}##1}%
  \if@rmtwocol
  { \if@debug{\ClassWarning{rmaa}{TWO COLUMN}}\fi
    \twocolumn[\@maketitle]%
    }
  \else
  { \if@debug{\ClassWarning{rmaa}{ONE COLUMN}}\fi
    \newpage
    \global\@topnum\z@   % Prevents figures from going at top of page.
    \if@manuscript \@MSmaketitle \else \@maketitle \fi
    }
  \fi
  \if@manuscript\else\thispagestyle{rmaafirstpage}\@thanks\fi
  \ifx\NotesList\@rmaaEmpty\else{%
    \NoteMark{1--\arabic{notecounter}}\vspace*{-2\baselineskip}%
  }\fi%
  \endgroup
%    \end{macrocode}
% Don't reset footnote counter (not sure whether this is good thing
% or not).
%    \begin{macrocode}
  %%  \setcounter{footnote}{0}%
  \let\thanks\relax\let\affil\relax\let\altaffil\relax
  \let\maketitle\relax\let\@maketitle\relax
  \gdef\@thanks{}\gdef\@author{}\gdef\@title{}}
%    \end{macrocode}

% \begin{macro}{\@kwhwidth}
% \begin{macro}{\@kwtwidth}
% Redefine \cs{@maketitle} to include the abstract, resumen and
% keywords. Note that in the manuscripty style we use
% \cs{@MSmaketitle} instead.  
% \changes{3.16}{2003/04/07}{
%  Added extra vertical space between authors and accepted date.}
% \changes{3.26}{2011/08/15}{Add extra label to review articles}
%    \begin{macrocode}
\newlength{\@kwhwidth}\newlength{\@kwtwidth}
\newcommand\reviewname{Review}
\renewcommand{\@maketitle}{%
  \newpage
  \null
  \if@reviewarticle
  \smallskip
  {\large\bfseries\itshape\reviewname}%
  \bigskip
  \else
  \vskip\@before@title@skip%
  \fi
  \begin{center}%
    {\@title \par}%
    \vskip\@before@author@skip%
    \parbox{\rmaa@AuthorWidthFrac\textwidth}{\centering
      \lineskip\@between@author@skip%
      \@author
      \par}
    %% this should be replaced with `Submitted to...' for preprints
    \if@printdates
    \vskip\@before@dates@skip%
    \if@preprint
    \rmaa@ReceivedStyle{Submitted to RMxAA \@received}
    \else
    \rmaa@ReceivedStyle{Received \@received; accepted \@accepted}
    \fi
    \fi
    \if@manuscript
    \rmaa@ReceivedStyle{Submitted to RMxAA \@received}
    \fi
    \if@guide
    \rmaa@ReceivedStyle{Version \rmaa@version, \rmaa@date}
    \fi
    \vskip\@before@abstract@skip%
    \if@abstractpresent
    \parbox{\rmaa@AbstractWidthFrac\textwidth}{%
      \setlength{\parindent}{0.9cm}%
      \centerline{\rmaa@AbstractHeaderStyle{Abstract}}\par%
      \vskip\@after@abstracthead@skip%
      \if@newstyle\noindent\fi\@abstract
      \vskip\@between@abstract@skip%
      \centerline{\rmaa@AbstractHeaderStyle{Resumen}}\par%
      \vskip\@after@abstracthead@skip%
      \if@newstyle\noindent\fi\@resumen
      \if@setresumen\else\@abstract\fi % append english abstract if
                                       % spanish one not set
      \if@guide\else            %do not use keywords for guide
      \vskip\@before@kwd@skip%
      % Format the Key Words with hanging indentation:
      % They are put in a \parbox with width equal to the line width
      % minus the width of the header. 
      { % Widths of keyword
        \settowidth{\@kwhwidth}{\@keywordheader~}   % header and
        \setlength{\@kwtwidth}{\linewidth}          % keyword text
        \addtolength{\@kwtwidth}{-\@kwhwidth}       %
        \noindent\@keywordheader~\parbox[t]{\@kwtwidth}{%
          \rmaa@KeywordStyle{\boldmath\uppercase{\@keywords}}}}%
      \fi
      }
      \fi
  \end{center}%
  \vspace*{\afterkwdskip}
  \par
  }
\newlength{\afterkwdskip}\setlength{\afterkwdskip}{\z@}
%    \end{macrocode}
% \end{macro}
% \end{macro}
%
%
% \begin{macro}{\reviewarticle}
% \begin{macro}{\reviewtoc}
% \begin{macro}{\reviewlof}
%   User command to specify a review article. Must be used in
%   preamble. Main effect is to place the text ``Review'' before the
%   article title.  Also add command for setting a table of contents
%   and a list of figures at the beginning of the article.
%    \begin{macrocode}
\newcommand{\reviewarticle}{\@reviewarticletrue}
\newcommand\reviewtoc{%
  \section*{CONTENTS}%
  \@starttoc{rtc}%
}
\newcommand\reviewlof{%
  \section*{LIST OF FIGURES}%
  \@starttoc{lof}%
}
%    \end{macrocode}
% \end{macro}
% \end{macro}
% \end{macro}
% \begin{macro}{\l@section}
%    \begin{macrocode}
\renewcommand*\l@section[2]{%
  \ifnum \c@tocdepth >\z@
    \addpenalty\@secpenalty
    \addvspace{0.8em \@plus\p@}%
    \setlength\@tempdima{1.5em}%
    \begingroup
      \parindent \z@ \rightskip \@pnumwidth
      \parfillskip -\@pnumwidth
      \leavevmode %\bfseries
      \advance\leftskip\@tempdima
      \hskip -\leftskip
      #1\nobreak\hfil \nobreak\hb@xt@\@pnumwidth{\hss #2}\par
    \endgroup
  \fi}
%    \end{macrocode}
% \end{macro}
%
% \begin{macro}{\@MSmaketitle}
% A special version of |\@maketitle| for the \texttt{manuscript}
% option. This can deal with a multi-page abstract/resumen. Not very
% well, mind. 
%    \begin{macrocode}
\newcommand{\@MSmaketitle}{%
  \if@debug\ClassInfo{rmaa}{Using special MS maketitle command}\fi
  \newpage
  {
    \thispagestyle{rmaafirstpage}
%     \setlength{\parindent}{0pt}
    \null\vskip 2ex
    \begin{center}
      \@title\par
      \@author\par
      \rmaa@ReceivedStyle{%
        \ifx\@received\@rmaaEmpty
          Draft version: \@date
        \else
          Received: \@received
          \ifx\@accepted\@rmaaEmpty
          \else
            ; Accepted: \@accepted
          \fi
        \fi}\par
    \end{center}
    \@thanks
    \centerline{\rmaa@AbstractHeaderStyle{Resumen}}
    \noindent\@resumen\par\vskip 2ex
    \centerline{\rmaa@AbstractHeaderStyle{Abstract}}
    \noindent\@abstract\par\vskip 2ex
    \noindent\@keywordheader~\rmaa@KeywordStyle{\@keywords}\par
%     \@keywordheader~\rmaa@KeywordStyle{\uppercase{\@keywords}}\par
  }%
  }
%    \end{macrocode}
% \end{macro}

% \subsection{The sectioning commands}
% \label{sec:section-commands}

% \begin{macro}{\section}
% \begin{macro}{\subsection}
% \begin{macro}{\subsubsection}
% \begin{macro}{\paragraph}
% \begin{macro}{\subparagraph}
% Redefine all the standard commands to use these lengths. 
%    \begin{macrocode}
\renewcommand\section{\@startsection {section}{1}{\z@}%
                                   {\abovesectionskip}%
                                   {\belowsectionskip}%
                                   %% how do we get upper case here
                                   {\centering\reset@font\normalsize}}
\renewcommand\subsection{\@startsection{subsection}{2}{\z@}%
                                     {\abovesectionskip}%
                                     {\belowsectionskip}%
                                     {\centering\reset@font\normalsize\itshape}}
\renewcommand\subsubsection{\@startsection{subsubsection}{3}{\z@}%
                                     {\abovesectionskip}%
                                     {\belowsectionskip}%
                                     {\centering\reset@font\normalsize\itshape}}
\renewcommand\paragraph{\@startsection{paragraph}{4}{\z@}%
                                    {\aboveparaskip}%
                                    {\belowparaskip}%
                                    {\reset@font\normalsize\bfseries}}
\renewcommand\subparagraph{\@startsection{subparagraph}{5}{\parindent}%
                                      {\aboveparaskip}%
                                      {\belowparaskip}%
                                      {\reset@font\normalsize\bfseries}}
%    \end{macrocode}
% \end{macro}
% \end{macro}
% \end{macro}
% \end{macro}
% \end{macro}

% \begin{macro}{secnumdepth}
% Make sure that we number the subsubsections---this is done by
% default in the |article| class but not the |book| class, which is
% why we must set it explicitly here.
%    \begin{macrocode}
\setcounter{secnumdepth}{3}
%    \end{macrocode}
% \end{macro}
% \changes{2.2beta10}{2001/11/07}{Fixed \cs{secnumdepth} for \cs{book}
% style.}

% \begin{macro}{\thesubsection}
% \begin{macro}{\thesubsubsection}
% For subsection and subsubsection, we want the numbers to be roman
% but the title to be italic. This seems the easiest way to go about it.
%    \begin{macrocode}
\renewcommand{\thesubsection}{{\upshape\thesection.\@arabic\c@subsection}}
\renewcommand{\thesubsubsection}{\thesubsection{\upshape.\@arabic\c@subsubsection}}
%    \end{macrocode}
% \end{macro}
% \end{macro}

% \begin{macro}{\@seccntformat}
% Add ``.'' to the number here so we can remove it from |\thesection|,
% etc. Reduce gap between number and title WJH 27 Jan 1995
%    \begin{macrocode}
\def\@seccntformat#1{\csname the#1\endcsname .\hskip 0.4em}
%    \end{macrocode}
% \end{macro}

% \begin{macro}{\@sect}
% I didn't really want to mess with the |\@sect| command, but it was
% the only way I could manage to force upper case. Other than the
% commented parts, it is identical to the standard version. 
%    \begin{macrocode}
\def\@sect#1#2#3#4#5#6[#7]#8{
  \ifnum #2>\c@secnumdepth
     \let\@svsec\@empty
  \else
     \refstepcounter{#1}%
     \let\@@protect\protect
     \def\protect{\noexpand\protect\noexpand}%
     \edef\@svsec{\@seccntformat{#1}}%
     \let\protect\@@protect
  \fi
%    \end{macrocode}
% This is the bit I changed (27 Jan 1995). Use upper case for
% |\section|, use mixed case for the rest. 
%    \begin{macrocode}
  \ifnum #2=1                   
     \def\sec@upcase##1{\uppercase{##1}}%
     \def\sec@shape{\upshape}%
  \else                         
     \def\sec@upcase##1{\relax{##1}}%
     \def\sec@shape{\itshape}%
  \fi
  \@tempskipa #5\relax
  \ifdim \@tempskipa>\z@
     \begingroup #6\relax
          \@hangfrom{\hskip #3\relax\@svsec}%
                    {\interlinepenalty \@M \sec@upcase{#8}\par}%
     \endgroup
     \csname #1mark\endcsname{#7}%
%    \end{macrocode}
% Only write TOC entries if doing the Guide style. 
%    \begin{macrocode}
     \if@guide 
        \addcontentsline{toc}{#1}{%
          \ifnum #2>\c@secnumdepth \else
             \protect\numberline{\csname the#1\endcsname}%
          \fi
          #7}%
     \fi
%    \end{macrocode}
% Also allow per-article TOC for review papers
%    \begin{macrocode}
     \if@reviewarticle
        \ifnum #2>2 \else
           \addcontentsline{rtc}{#1}{%
             \protect\numberline{\csname the#1\endcsname}%
             \sec@upcase{\sec@shape #7}}%
        \fi
     \fi
  \else
     \def\@svsechd{#6\hskip #3\relax  %% \relax added 2 May 90
       \@svsec #8\csname #1mark\endcsname{#7}%
       \if@guide
          \addcontentsline{toc}{#1}{%
            \ifnum #2>\c@secnumdepth \else
               \protect\numberline{\csname the#1\endcsname}%
            \fi
            #7}%
       \fi
       }%
  \fi
  \@xsect{#5}}
%    \end{macrocode}
% \end{macro}


% \subsection{Appendices}
% \label{sec:appendices}
% Make sure that appendices in individual articles work properly. We
% copy the definition from the article class (the book one doesn't do
% what we want).  
% \begin{environment}{appendices}
%   This uses an environment |appendices| inside of which each
%   |\section| will become a separate appendix.  Note that the
%   appendices must come before the references (otherwise the full
%   addresses mechanism will not work properly). We make the equation
%   numbering be subordinate to the appendix (equations A12, B6,
%   etc). Finally, we have a medskip followed by APPENDICES to
%   introduce them. \changes{3.17}{2002/04/08}{ Added appendices environment
%   that does the right thing.}
%    \begin{macrocode}
\newenvironment{appendices}{\par
  \begingroup
  \setcounter{section}{0}%
  \setcounter{subsection}{0}%
  \renewcommand\thesection{\@Alph\c@section}
  \renewcommand\theequation{\thesection\arabic{equation}}
  \medskip
  \section*{APPENDICES}}
{\endgroup}
%    \end{macrocode}
% \end{environment}
% \begin{environment}{appendix}
%   If there is only one appendix, then use |appendix| environment
%   instead.  \changes{4.0}{2019/11/12}{ Added
%   appendix environment for case where there is only one section.}
%    \begin{macrocode}
\renewenvironment{appendix}{\par
  \begingroup
  \setcounter{section}{0}%
  \setcounter{subsection}{0}%
  \renewcommand\thesection{\@Alph\c@section}
  \renewcommand\theequation{\thesection\arabic{equation}}
  \medskip
  \section*{APPENDIX}}
{\endgroup}
%    \end{macrocode}
% \end{environment}

% \subsection{Floating thingies}
% \label{sec:floats}   

%% 
%    \begin{macrocode}
\renewcommand{\thetable}{\@arabic\c@table}
\renewcommand{\thefigure}{\@arabic\c@figure}
\renewcommand{\figurename}{Fig.}
\renewcommand{\tablename}{TABLE}
%\renewcommand{\fnum@table}{\MakeUpperCase{\tablename}~\thetable}
%    \end{macrocode}
% For some reason, I have to write `TABLE' explicitly here
% If I use |\tablename|, it gets converted to mixed case---God knows
% why. 
% |\figurename| is also being ignored. 
%    \begin{macrocode}
\renewcommand{\fnum@table}{TABLE~\thetable}
\renewcommand{\fnum@figure}{Fig.~\thefigure}
%    \end{macrocode}

% Do captions in 9 pt face in the main journal style and in new-style
% anything. 
%    \begin{macrocode}
\long\def\@makefigcap#1#2{%
  \if@proceedings\else\small\fi
  \if@newstyle\small\fi 
  \vskip\abovecaptionskip
  \sbox\@tempboxa{#1. #2}%
  \ifdim \wd\@tempboxa >\hsize
    #1. #2 \par
  \else
    \global \@minipagefalse
    \hb@xt@\hsize{\hfil\box\@tempboxa\hfil}%
  \fi
  \vskip\belowcaptionskip}
%    \end{macrocode}

% Make sure table caption is upper case. 
%    \begin{macrocode}
\long\def\@maketabcap#1#2{%
  \hbox to\hsize{\hfil #1 \hfil}
  \vskip 6pt 
  \centering\MakeUpperCase{#2}\par
  \vskip 6pt}
%    \end{macrocode}

%    \begin{macrocode}
\newlength{\@rottabcapwidth}
\long\def\@makerottabcap#1#2{%
  \setbox\@tempboxa\hbox{\MakeUpperCase{#2}}%
  \ifdim\wd\@tempboxa >.8\vsize
     \setlength\@rottabcapwidth{.8\vsize}
  \else
     \setlength\@rottabcapwidth{\wd\@tempboxa}
  \fi
  \rotatebox{90}{%
    \begin{minipage}{\vsize}\centering
      \begin{minipage}{\@rottabcapwidth}\centering
        #1\\ \vskip 6pt \MakeUpperCase{#2}
      \end{minipage}%
    \end{minipage}
    }%
  \hskip 6pt
  }

\newenvironment{Sideways}
    {\begin{sideways}\begin{minipage}{\vsize}\centering}
    {\end{minipage}\end{sideways}}

\newenvironment{Sidewaystable}
{%
  \let\oldcaption=\caption%
  \def\caption{
    \def\@captype{table}\renewcommand{\@makecaption}{\@maketabcap}
    \oldcaption}
  \begin{sidewaystable}\centering}
  {\end{sidewaystable}\let\caption=\oldcaption}
\newenvironment{Sidewaystable*}
{%
  \let\oldcaption=\caption%
  \def\caption{
    \def\@captype{table}\renewcommand{\@makecaption}{\@maketabcap}
    \oldcaption}
  \begin{sidewaystable*}\centering}
  {\end{sidewaystable*}\let\caption=\oldcaption}
%    \end{macrocode}

% Not sure why we need to stick these lines in. 
%    \begin{macrocode}
\renewenvironment{figure}
                 {\renewcommand{\@makecaption}{\@makefigcap}\@float{figure}}
                 {\end@float}
\renewenvironment{figure*}
                 {\renewcommand{\@makecaption}{\@makefigcap}\@dblfloat{figure}}
                 {\end@dblfloat}
\renewenvironment{table}
                 {\renewcommand{\@makecaption}{\@maketabcap}\@float{table}}
                 {\end@float}
\renewenvironment{table*}
                 {\renewcommand{\@makecaption}{\@maketabcap}\@dblfloat{table}}
                 {\end@dblfloat}
%    \end{macrocode}

% Redefine |\@caption| so that we don't bother with the list of
% figures. This saves us having to worry about fragile commands in the
% captions. 
%    \begin{macrocode}
\long\def\@caption#1[#2]#3{%
  \par
%%  \addcontentsline{\csname ext@#1\endcsname}{#1}%
%%    {\protect\numberline{\csname the#1\endcsname}{\ignorespaces #2}}%
  \begingroup
    \@parboxrestore
    \if@minipage
      \@setminipage
    \fi
    \normalsize
    \@makecaption{\csname fnum@#1\endcsname}{\ignorespaces #3}\par
  \endgroup}
%    \end{macrocode}

% Make tables not quite so cramped. Allow |\hline\hline| to give a
% poor man's thick rule. 
%    \begin{macrocode}
\renewcommand{\arraystretch}{1.15}
\setlength{\doublerulesep}{0.4pt}
%    \end{macrocode}

% \begin{macro}{\tablecols}
% \begin{macro}{\@tablecols}
% \begin{macro}{\tabcaptioncontinued}
% |\tablecols| is a user command for setting the number of columns. It
% is used by commands that try and span the entire table (I think
% there must be a better way of acheiving
% this). |\tabcaptioncontinued| is 
% a user command for setting the caption of a continued table on the
% next page. 
%    \begin{macrocode}
\newcommand{\tablecols}[1]{\renewcommand{\@tablecols}{#1}}
\newcommand{\@tablecols}{1}
\newcommand{\tabcaptioncontinued}{%
  \multicolumn{\@tablecols}{c}{\addtocounter{table}{-1}TABLE~\thetable\ (CONTINUED)} \\ \\
  }
%    \end{macrocode}
% \end{macro}
% \end{macro}
% \end{macro}

% \subsubsection{Footnotes in tables}
% \label{sec:table-footnotes}
% \begin{macro}{\tablenotemark}
% \begin{macro}{\tablenote}
% \begin{macro}{\tablenoteskip}
% Old user commands |\tablenotemark{MARK}|,
% |\tablenote{MARK}{TEXT}|---now deprecated. 
%    \begin{macrocode}
\newcommand\tablenotemark[1]{$^\mathrm{#1}$} 
\newcommand\tablenote[2]{\if@proceedings\else\small\fi ${}^\mathrm{#1}$ #2} 
\newcommand\tablenoteskip{\\[0.5\baselineskip]} 
%    \end{macrocode}
% \end{macro}
% \end{macro}
% \end{macro}

% \begin{macro}{\tabnotemark}
% \begin{macro}{\tabnotetext}
% \begin{macro}{\tabnotetextTWO}
% Alternative way of doing footnotes in tables. 
%    \begin{macrocode}
\newcommand{\tabnotemark}[1]{\textsuperscript{\lowercase{#1}}}
%\newcommand{\tabnotetextTWO}[4]{\tabnotetext{#1}{#2}\tabnotetext{#3}{#4}}
\newcommand{\tabnotetextTWO}[4]{%
  \multicolumn{\@tablecols}{l}{%
    \parbox{0.5\tabnotewidth}{\small\textsuperscript{#1}#2}%
    \parbox{0.5\tabnotewidth}{\small\textsuperscript{#3}#4}%
    }%
  \\}
\newcommand{\tabnotetext}[2]{%
  \multicolumn{\@tablecols}{@{}p{\tabnotewidth}@{}}{\small\textsuperscript{#1}#2}\\}
\newlength{\tabnotewidth}
%    \end{macrocode}
% \end{macro}
% \end{macro}
% \end{macro}


% \begin{macro}{changemargin}
%    \begin{macrocode}
% Command to allow margins to be changed on the fly. Useful to allow
% tables and Figures to protrude into the margins. 
%
% \changes{3.24}{2007/09/14}{Added \texttt{changemargin} environment
% (copied from TeX FAQ).}
\newenvironment{changemargin}[2]{% 
  % From TeX FAQ http://www.tex.ac.uk/cgi-bin/texfaq2html?label=chngmargonfly
  \begin{list}{}{%
    \setlength{\topsep}{0pt}%
    \setlength{\leftmargin}{#1}%
    \setlength{\rightmargin}{#2}%
    \setlength{\listparindent}{\parindent}%
    \setlength{\itemindent}{\parindent}%
    \setlength{\parsep}{\parskip}%
  }%
  \item[]}{\end{list}}
%    \end{macrocode}
% \end{macro}

% \subsection{The bibliography and references}
% \label{sec:biblio}
% \begin{macro}{\beforebiblioskip}
% \begin{macro}{\afterbiblioskip}
% I'm not entirely sure why I introduced these two commands, since
% directly putting |\vspace| before and after
% |\begin{thebibliography}| seems to work well enough. 
%    \begin{macrocode}
\newlength{\beforebiblioskip}
\setlength{\beforebiblioskip}{0pt}
\newlength{\afterbiblioskip}
\setlength{\afterbiblioskip}{0pt}
%    \end{macrocode}
% \end{macro}
% \end{macro}
% \begin{macro}{thebibliography}
% The environment itself is done as starred section with extra
% vertical space after heading. This space is currently zero. 
% In the book style we do the cunning trick with rereading the file in
% order to get the references into the |\ADS@references| macro.
% 
% \changes{3.05}{2003/01/02}{Removed the \cs{beforebiblioskip} and
% \cs{afterbiblioskip} commands, since they cause unwanted vertical
% space in some instances.}  
%    \begin{macrocode}
\renewenvironment{thebibliography}{%
  \if@book\RereadBiblio\fi
%%  \vspace*{\beforebiblioskip}
  \section*{REFERENCES} \frenchspacing
%%  \vspace*{\afterbiblioskip}
  \small
%    \end{macrocode} 
% The references themselves are set as a list in order 
% to get the hanging indentation. They are done in a smaller font and
% with french spacing so we needn't worry about the distinction
% between sentence-ending periods and abbreviations. 
%    \begin{macrocode} 
  \list{\null}{\leftmargin 1.5em\labelwidth\z@\labelsep\z@\itemindent -1.5em
    \itemsep\z@\parsep\z@\usecounter{enumi}}
  \sloppy\clubpenalty4000\widowpenalty4000
  \sfcode`\.=1000\relax\@biblisttrue}
{\endlist\@biblistfalse\normalsize\@outputaddressesifnecessary}
%    \end{macrocode}
% \end{macro}
% \subsection{Uniquification of cross-references}
% \label{sec:uniquification}
% 
% \begin{macro}{\articlename}
% In order that the |\label{}| command in individual articles be
% guaranteed to give a unique anchor, we append the article name (held
% in the |\articlename| command) to each key. This requires
% redefinition of the following commands: |\cite|, |\nocite|, |\ref|,
% |\pageref|, |\label|. 
% \begin{macro}{\oldcite}
% \begin{macro}{\oldnocite}
% \begin{macro}{\oldref}
% \begin{macro}{\oldpageref}
% \begin{macro}{\oldlabel}
% The original versions of these commands are
% saved as |\oldcite|, |\oldnocite|, etc. 
%
% The uniquifying of cross-references is now always done, so 
% we don't get multiply-defined labels in the single article case
%    \begin{macrocode}
\newcommand{\articlename}{anon}
\AtBeginDocument{
  \let\oldcite=\cite 
  \let\oldnocite=\nocite 
  \let\oldref=\ref 
  \let\oldpageref=\pageref
  \let\oldlabel=\label
}
%    \end{macrocode}
% \end{macro}
% \end{macro}
% \end{macro}
% \end{macro}
% \end{macro}
% \end{macro}
%
% Uniquify all cross-references by prepending the article name.  This
% now has to go into |\AtBeginDocument| so it is working on the new
% versions that hyperref.sty has set up.  Then we need to put the
% |\@begofarticlehook| afterwards, so that the first page label is using
% the new definitions.  If amsmath is loaded, then we also need to patch
% the special label command that they use. \changes{4.0}{2019/11/18}{Fix labels for amsmath display equations, in case that package has been loaded.}
%    \begin{macrocode}
\AtBeginDocument{
  \renewcommand{\ref}[1]{\oldref{\articlename:#1}%
    \if@debug{\ClassInfo{rmaa}{Reference to \articlename:#1}}\fi}
  \renewcommand{\pageref}[1]{\oldpageref{\articlename:#1}%
    \if@debug{\ClassInfo{rmaa}{Page reference to \articlename:#1}}\fi}
  \renewcommand{\label}[1]{\oldlabel{\articlename:#1}%
    \if@debug{\ClassInfo{rmaa}{Label for \articlename:#1}}\fi}%
  \ifdef{\label@in@display}{%
    \let\oldlabel@in@display=\label@in@display
    \renewcommand{\label@in@display}[1]{\oldlabel@in@display{\articlename:#1}%
      \if@debug{\ClassInfo{rmaa}{AMSTeX Label for \articlename:#1}}\fi}%
  }{}
  \@begofarticlehook
}%
%    \end{macrocode}
% Citations are more complicated since they can contain a list of
% keys. These are really only here for use with the rmaacite package,
% which is utterly deprecated (02 Jan 2003) now that I have the natbib
% code working. The command \cs{cite} should \emph{not} be used with
% natbib: use \cs{citep}, \cs{citet}, etc.\@ instead. 
%    \begin{macrocode}
\renewcommand{\nocite}[1]{\@bsphack
  \@for\@citeb:=#1\do{%
    \edef\@citeb{\expandafter\@firstofone\@citeb:\articlename}%
    \if@filesw\immediate\write\@auxout{\string\citation{\@citeb}}\fi
    \@ifundefined{b@\@citeb:\articlename}{\G@refundefinedtrue
      \@latex@warning{Citation `\@citeb' undefined}}{}}%
  \@esphack}
\renewcommand{\cite}[1]{\oldcite{#1:\articlename}}
%   \end{macrocode}


% Horrible fiddling with the bibliography citation mechanism. 
% This has to be globally available in case rmaacite puts it in the
% \texttt{.aux} file. 
%    \begin{macrocode}
\def\lbibcite#1#2#3{\global\@namedef{b@#1}{#2}
                    \global\@namedef{by@#1}{#3}}
%    \end{macrocode}

% \begin{macro}{\oldbibitem}
% \begin{macro}{\@rmaaBibArg}
% \begin{macro}{\@rmaaSpace}
% \begin{macro}{rmaabib}
% Of course, |\bibitem| needs to get the uniquification treatment as
% well. While we are at it, we redefine it to use consecutive integers
% for the key if the user leaves it blank (using the counter
% |rmaabib|). This allows one to put 
% |\bibitem{}| for all your references without getting annoying
% messages about multiply-defined labels.
% 
% We also take the opportunity to insert the optional argument to
% \cs{bibitem} if it is missing in the form Author(DATE) so that natbib
% doesn't try to do numerical citations. Natbib already provides a
% hook \cs{@extra@b@citeb} to help uniquify the citation key, so we
% just plug into that. However, it puts it \emph{after} the original
% key, not before as I had had it. Easiest solution is to change the
% way that I do it. 02 Jan 2003. Note that the \cs{@extra@b@citeb}
% does not get written to the .aux  file in the \cs{citation} command
% but it does work. 
% \changes{3.14}{06 Feb 2003}{Reinstalled old definition of \cs{bibitem}
%  for the case where the nonatbib option is requested, so that the
%  rmaacite package will work. This means that you can't use rmaacite
%  and natbib in the same volume.}
%    \begin{macrocode}
\if@natbib\def\@extra@b@citeb{:\articlename}\fi
%%\else
\let\oldbibitem\bibitem
\newcommand{\@rmaaBibArg}{}
\newcommand{\@rmaaSpace}{ }
\newcounter{rmaabib}
\if@natbib
\renewcommand{\bibitem}[2][ABC(2003)]{%
  \renewcommand{\@rmaaBibArg}{#2}
  \ifx\@rmaaBibArg\@rmaaEmpty
     \addtocounter{rmaabib}{1}
     \oldbibitem[#1]{\thermaabib:\articlename}
  \else\ifx\@rmaaBibArg\@rmaaSpace
     \addtocounter{rmaabib}{1}
     \oldbibitem[#1]{\thermaabib:\articlename}
  \else 
     \oldbibitem[#1]{#2:\articlename}
  \fi\fi
  \renewcommand{\@rmaaBibArg}{}
  }
\else
\renewcommand{\bibitem}[1]{%
  \renewcommand{\@rmaaBibArg}{#1}
  \ifx\@rmaaBibArg\@rmaaEmpty
     \addtocounter{rmaabib}{1}
     \oldbibitem{\articlename:\thermaabib}
  \else\ifx\@rmaaBibArg\@rmaaSpace
     \addtocounter{rmaabib}{1}
     \oldbibitem{\articlename:\thermaabib}
  \else
     \oldbibitem{\articlename:#1}
  \fi\fi
  \renewcommand{\@rmaaBibArg}{}
  }
\fi
%%\fi
%    \end{macrocode}
% \end{macro}
% \end{macro}
% \end{macro}
% \end{macro}

% \begin{macro}{oldpcite}
% \begin{macro}{oldscite}
% Then, we need to repeat the above scheme, but for the |rmaacite|
% definition. 
%    \begin{macrocode}
\newcommand{\@redefinebibitem}{%
  \let\oldbibitem\bibitem
  \let\oldcite\cite
  \let\oldpcite\pcite
  \let\oldscite\scite
  \def\bibitem[##1<##2>]##3{%
    \renewcommand{\@rmaaBibArg}{##3}
    \ifx\@rmaaBibArg\@rmaaEmpty
       \addtocounter{rmaabib}{1}
       \oldbibitem[##1<##2>]{\thermaabib:\articlename}
    \else\ifx\@rmaaBibArg\@rmaaSpace
       \addtocounter{rmaabib}{1}
       \oldbibitem[##1<##2>]{\thermaabib:\articlename}
    \else 
       \oldbibitem[##1<##2>]{##3:\articlename}
    % trust the author here
%%       \oldbibitem[##1<##2>]{##3}
    \fi\fi
    }
  }
%    \end{macrocode}
% \end{macro}
% \end{macro}


% \subsection{HTML Table of Contents}
% \label{sec:hyper-TOC}
% \begin{macro}{\HyperTOCEntry}
% \begin{macro}{\HyperTOCpart}
% \begin{macro}{\HyperTOCmeeting}
% \begin{macro}{\HyperTOCeditors}
% \begin{macro}{\HyperTOCpreamble}
% These commands are used to write information to the \texttt{.htoc}
% file, which is used in generating a Web Page of the Table of
% Contents. Since they only ever appear preceded by |\protect| in the
% argument to an |\addtocontents| command,  they should never need to
% be expanded by the macros. However, we define them here just in case
% \dots
%    \begin{macrocode}
\newcommand{\HyperTOCEntry}[3]{\relax}
\newcommand{\HyperTOCpart}[1]{\relax}
\newcommand{\HyperTOCmeeting}[3]{\relax}
\newcommand{\HyperTOCeditors}[1]{\relax}
\newcommand{\HyperTOCpreamble}{\relax}
%    \end{macrocode}
% \end{macro}
% \end{macro}
% \end{macro}
% \end{macro}
% \end{macro}
%

% \subsection{Formatting the poster abstracts} 
% These can now be used without using the book option.
%    \begin{macrocode}

\newcommand{\@makeminititle}{%
  \newpage
  \null 
  \vskip 3em%
  \begin{center}%
    {\@title \par}%
  \end{center}
  \vskip 3em%
  }

\newif\if@abstractsads
\@abstractsadsfalse
\newcommand\abstractsADS{\@abstractsadstrue}

\if@book\RequirePackage{multicol}\fi
\newenvironment{abstracts}[1][Abstracts of Contributed Papers]{%
  %% First, make sure we start on an odd page (06 Jun 2000)
  \if@book\clearpage\thispagestyle{empty}\cleardoublepage\fi
  %%
  \renewcommand{\articlename}{ABSTRACTS}
  \title{#1}\shortauthor{Abstracts}\shorttitle{Abstracts}
  \listofauthors{}
  \if@book
  \addcontentsline{toc}{chapter}{\protect\@fiddleauthors{\@TOCtitle}{\rmaaTOCauthors}}%
  \addtocontents{htc}{\protect\HyperTOCEntry{\@rawtitle}{\@rawauthors}{\articlename}}
  \fi
  \pagestyle{rmaa} % this now seems necessary (06 Jun 2000)
  \thispagestyle{rmaafirstpage}
%%  \twocolumn[\@makeminititle]
  \renewcommand{\abstract}[1]{%
    \renewcommand\ADS@abstract{##1}%
    \renewcommand{\@abstract}{##1}\@writeabstract%
  }
  \renewcommand{\title}[2][]{%
    \gdef\ADS@authors{}%
    \gdef\ADS@affil{}%
    \gdef\ADS@references{}%
    \gdef\ADS@keywords{}%
    \renewcommand{\@rmaaArg}{##1}
    \ifx\@rmaaArg\@rmaaEmpty
       \renewcommand{\@TOCaftertitle}{}
    \else
       \renewcommand{\@TOCaftertitle}{ (\@rmaaArg)}
    \fi
    \renewcommand{\@title}{\rmaa@TOCTitleStyle{##2}}%
    \renewcommand{\@TOCtitle}{\rmaa@TOCTitleStyle{##2}\@TOCaftertitle}%
    \renewcommand{\rmaaIndexTitle}{\rmaa@IndexTitleStyle{##2\@TOCaftertitle.}}%
    \renewcommand\ADS@title{##2}%
    \xdef\ADS@firstpage{\@arabic{\c@page}}%
    }
  \renewcommand{\author}[1]{\let\@thanks\@empty\renewcommand\@author{##1} }
  \renewcommand{\footnotetext}[2][XX]{\par\noindent\textsuperscript{##1}~##2}
  \renewcommand{\footnoterule}{\kern-3\p@ \hrule \@width 3em \kern 2.6\p@}
  \if@book\@begofarticlehook\fi
  \if@book\begin{multicols}{2}[\@makeminititle]\else\@makeminititle\fi
  \begingroup
  \setcounter{footnote}{0}
  }% and at the end of the environment....
  {\endgroup\if@book\end{multicols}\listofabstracts\label{RMAALastPage}\clearpage\fi}
%% auxiliary commands used:
%% \@writeabstract does all the actual formatting 
\newcommand{\@writeabstract}{%
  \if@book%
  \addcontentsline{loa}{abstract}{\protect\@fiddleauthors{\@TOCtitle}{\rmaaTOCauthors}}%
  \fi%
  \begin{center}\@title\par\@author\end{center}
  \@abstract
  \par\noindent\footnoterule\nobreak\footnotesize\@thanks
  \par\noindent\hrulefill\goodbreak
  \vskip 1cm \@plus0.5cm
  %% write ADS entry
  \if@book\if@abstractsads\ADS@dumpdata\fi\fi%
  \endgroup % this pair is to make each abstract a self-contained group
  \begingroup
  \setcounter{footnote}{0}
  }

\if@book
%% List of abstracts (stuff written to .loa file)
\newcommand{\listofabstracts}{%
  \begin{multicols}{2}[%
    \vskip 1.5cm \vfill % leave at least 1.5 cm and go to bottom of page if possible
    \centerline{\MakeUpperCase{List of Abstracts}}
    \vskip 1.5cm \@plus0.5cm
    ][6cm] % this makes it start a new page if there is not at least 6cm of space
    \renewcommand{\@fiddleauthors}[2]{##1\hskip1em plus2em minus0.4em\penalty-100  ##2}
    \parindent\z@
%%    \parskip.5\p@ \@plus .7\p@\relax 
    \parskip3\p@ \@plus 1\p@ \@minus 1\p@\relax % allow quite a lot of
                                % leeway here to avoid breaks inside
                                % entries 
    \renewcommand{\\}{\relax} % temporarily turn off explicit line breaks just in case
    \renewcommand{\altaffilmark}[1]{\relax} % ditto with this
    \@starttoc{loa}
  \end{multicols}
  }
\newcommand*{\l@abstract}[2]{%
  \par\penalty-1000
  #1\hskip 1em\penalty0\null\nobreak\hfill#2}
\fi

%    \end{macrocode}
 
% \subsection{Code for typesetting an entire volume of conference
% proceedings} \label{sec:whole-volume}
%    \begin{macrocode}

\if@book
%% environment for including the individual articles
\newcommand{\@fiddleauthors}[2]{\relax}
\renewcommand{\theequation}{\@arabic\c@equation}
\newenvironment{articles}%
{ \begingroup % make it a group so that author definitions do not
              % spill over
  \gdef\ADS@authors{}
  \pagestyle{rmaa}
  \renewenvironment{document}{%
    \@begofarticlehook
    \addtocounter{chapter}{1}%
    \addcontentsline{toc}{chapter}{%
      \protect\@fiddleauthors{\@TOCtitle}{\rmaaTOCauthors}}%
    \addtocontents{htc}{\protect\HyperTOCEntry{\@rawtitle}{\@rawauthors}{\articlename}}
    }{\relax}
  \renewcommand{\documentclass}[2][foo]{\relax} 
  \renewcommand{\usepackage}[2][foo]{%
    \def\@packagename{##2}\def\@rmaacitename{rmaacite}%
    \ifx\@packagename\@rmaacitename
    \@rmcitetrue\ClassInfo{rmaa}{This paper uses rmaacite}
    \makeatletter\input{rmaacite.sty}\makeatother % read it in here so that all changes
    \else                                         % are local to paper
    \if@debug\ClassWarning{rmaa}{Package \@packagename\space was requested by \articlename\space}\fi
    \fi} 
  %% Don't allow individual articles to mess with the pages, volume or year
  \renewcommand{\SetFirstPage}[1]{\relax}
  \renewcommand{\SetYear}[1]{\relax}
  \renewcommand{\SetVolume}[1]{\relax}
  \renewcommand{\SetConfTitle}[1]{\relax}
  %% Don't allow individual articles to change these
  \renewcommand{\nonstopmode}{\relax}
  \renewcommand{\scrollmode}{\relax}
  \renewcommand{\outputplates}{\relax}
  %% XXXX need to check if using rmaacite and act accordingly
  %% avoid multiple declarations of the same length
  \let\oldnewlength=\newlength
  \renewcommand{\newlength}[1]{%
    \edef\reserved@a{\expandafter\@gobble\string ##1}%
    \@ifundefined{\reserved@a}{%
      \oldnewlength{##1}}{%
      \if@debug{\ClassInfo{rmaa}{\reserved@a\space already defined: doing nothing}}\fi%
      \relax}}%
  \renewcommand\newcommand{\providecommand}
  }%
{\endgroup\onecolumn}
     

%    \end{macrocode}
% Command to be used in `articles' environment for including an
% article
%    \begin{macrocode}
\newcommand{\article}[2][]{%
  % optional first argument is extra text to put in brackets after title in TOC
  % e.g. ``Invited Talk'' or ``Poster''
  \renewcommand{\articlename}{#2}
  % reset loads of counters
  \setcounter{footnote}{0}
  \setcounter{figure}{0}
  \setcounter{table}{0}
  \setcounter{section}{0}
  \setcounter{equation}{0}
  \setcounter{rmaabib}{0}
  \setcounter{author}{0}
  \setcounter{AUTHOR}{1}
  \gdef\ADS@authors{}%
  \gdef\ADS@affil{}%
  \gdef\ADS@affiliations{}%
  \gdef\ADS@references{}%
  \@firstauthortrue
  \@firstaffiltrue
  \@firstaltaffiltrue
  \@firstkwdtrue
  \@rmcitefalse
  \if@suppressfulladdresses\else\@needaddressestrue\fi
  \renewcommand{\@rmaaArg}{#1}
  \ifx\@rmaaArg\@rmaaEmpty
     \renewcommand{\@TOCaftertitle}{}
  \else
     \renewcommand{\@TOCaftertitle}{ (#1)}
  \fi
%    \end{macrocode}
% In the main journal articles must start on an odd page. 
%    \begin{macrocode}
  \if@journal\clearpage\thispagestyle{empty}\cleardoublepage\fi
%% Now input the file, then run the end-of-article hook (reinstated 28
  %% Sep 2001)
  {\input{\@articledir/\articlename.tex}\@endofarticlehook}}
%    \end{macrocode}
% Command to separate articles into sessions. Write an entry in the
% TOC, makes sure we start on an odd page (added 06 Jun 2000), then do
% a title page if \cs{SessionTitlePages} was called. 
% \changes{3.12}{04 Feb 2003}{Added option to have title page for each
% session. Turn it on with \cs{SessionTitlePages}.}
%    \begin{macrocode}
\newcommand{\session}[1]{
  %% this does nothing except write an entry in the TOC
  \addcontentsline{toc}{part}{\rmaa@TOCPartStyle{#1}}
  \addtocontents{htc}{\protect\HyperTOCpart{#1}}
  %% and also makes sure we start on an odd page (added 06 Jun 2000)
  \clearpage\thispagestyle{empty}\cleardoublepage
  \stepcounter{part}
  \if@sessiontitlepages{%
    \onecolumn
    \thispagestyle{empty}
    \renewcommand{\articlename}{Section\thepart TitlePage}
    \label{RMAAFirstPage}
    \make@session@title{#1}
    \label{RMAALastPage}
    \clearpage\thispagestyle{empty}\cleardoublepage
  }\fi
  }
\newcommand{\make@session@title}[1]{%
  {%
    \setlength\parindent{0pt}
    \vspace*{\stretch{1}}
    \rule{\linewidth}{1mm}
    \begin{flushright}\Huge\bfseries\boldmath
      #1
    \end{flushright}
    \rule{\linewidth}{1mm}
    \vspace*{\stretch{2}}
  }
}
%    \end{macrocode}
% Command to turn on separate title pages for each session. 
%    \begin{macrocode}
\newcommand{\SessionTitlePages}{\@sessiontitlepagestrue}
\newcommand{\BoldSessionTitles}{%
  \renewcommand{\rmaa@TOCPartStyle}[1]{\textbf{\rmaa@TOCTitleStyle{##1}}}
  }
\newcommand{\ArticleDir}[1]{\renewcommand{\@articledir}{#1/}\graphicspath{{1/}}}
\newcommand{\@articledir}{.}
%\newcommand{\ps@rmaafront}{\ps@rmaa}% page style for front matter (same as ever, actually)
\newcommand{\ps@rmaafront}{% page style for front matter:  page # in each foot
  \renewcommand\@oddfoot{\hfil{\thepage}}%
  \renewcommand\@evenfoot{{\thepage}\hfil}%
  % shorttitle in both heads
  \renewcommand\@evenhead{\hfil\@shorttitle\hfil}%
  \renewcommand\@oddhead{\@evenhead}
  }
\newcommand{\ps@rmaafrontfirstpage}{% page style for first page of front matter
  % page # in each foot, on outside edge
  \renewcommand\@oddfoot{\hfil\thepage}%
  \renewcommand\@evenfoot{\thepage\hfil}%
  % nothing in the even heads
  \renewcommand\@evenhead{\relax}%
  \renewcommand\@oddhead{\relax}%
  }
%% slight modification to the standard \frontmatter and \mainmatter
%% commands, just to set the \pagestyle
\renewcommand\frontmatter{%
  \latexonecolumn
  \cleardoublepage
  \@frontmattertrue
  \@mainmatterfalse
  \pagenumbering{roman}
  \pagestyle{rmaafront}
  \addtocounter{page}{2} % leave space for the title page
  }
\renewcommand\mainmatter{%
  \pagestyle{empty}
  \if@rmtwocol\twocolumn\fi
  \cleardoublepage
  \@frontmatterfalse
  \@mainmattertrue
  \pagenumbering{arabic}
  \pagestyle{rmaa}
  }
\renewcommand\backmatter{%
  \clearpage\pagestyle{empty}\cleardoublepage
  \@mainmatterfalse
  %% put extra space in the TOC
  \addtocontents{toc}{\bigskip}
  }
%    \end{macrocode}
%
% \subsection{Front sections}\label{sec:front-sections}
% \begin{macro}{\notitle}
% \begin{macro}{\shortname}
% \begin{macro}{frontsection}
% The environment \texttt{frontsection} is for doing a section in the
% front matter (preface, list of participants, or whatever). It has 
% one optional argument:
% \begin{enumerate}
% \item[\#1] Short title (no embedded spaces) for
% use in making PS files. E.g., ``list\_of\_participants''.   
% \end{enumerate}
% and two mandatory arguments: 
% \begin{enumerate}
% \item[\#2] Title of section (e.g., `Preface')
% \item[\#3] Commands to be evaluated before processing the contents
% of the environment. This is mainly useful for putting in the things
% that would go in the preamble of a normal article  (e.g.,
% |\listofauthors|, |\indexauthor|, |\notitle|) 
% \end{enumerate}
%    \begin{macrocode}
\newcommand{\notitle}{\@titlevisiblefalse} \newcommand{\shortname}{}
\newenvironment{frontsection}[3][]{%
  \begingroup
  \renewcommand{\shortname}{#1}
  \ifx\shortname\@rmaaEmpty
     \renewcommand{\articlename}{#2}
  \else
     \renewcommand{\articlename}{\shortname}
  \fi
  \@begofarticlehook
  \@titlevisibletrue % default
  \thispagestyle{rmaafrontfirstpage}
  \listofauthors{} % blank out this in case it is not specified
  \title{{#2}}\shorttitle{#2}\shortauthor{#2}
  #3 % evaluate the command given as second mandatory arg
  \addcontentsline{toc}{chapter}{\protect\@fiddleauthors{\@TOCtitle}{\rmaaTOCauthors}}%
  \addtocontents{htc}{\protect\HyperTOCEntry{\@rawtitle}{\@rawauthors}{\articlename}}
  \if@titlevisible
     \vspace*{1cm}\centerline{\rmaa@FrontSecStyle{\@title}}\vspace*{1cm}
  \fi
  \setlength{\parskip}{3ex}
  }{\frontlabel{RMAALastPage}\clearpage\endgroup}
%    \end{macrocode}
% \end{macro}
% \end{macro}
% \end{macro}
%    \begin{macrocode}
%% Definition of the command \frontlabel, which rights out something like `F4' to
%% the .aux file, instead of `\textlatin  {i}'. This is so it is easy to post-process
%% the file for making PS files of individual sections
\newcommand{\frontlabel}[1]{%
  \@bsphack
  \protected@write\@auxout{}%
     {\string\newlabel{\articlename:#1}{{}{F\@arabic{\c@page}}}}
  \@esphack
  }
%% User command to say X, Y, Z,\\ Editors of the Proceedings
\newcommand{\EditorSignature}{\vspace*{\parskip}\rightline{\@editors}%
  \rightline{Editors of the proceedings}}
%% Fiddle with \thesection, etc. so that they don't include the chapter number
%% Basically go back to the `article' definitions of these
\renewcommand{\thesection}{\@arabic\c@section}
\renewcommand{\thefigure}{\@arabic\c@figure}
\renewcommand{\thetable}{\@arabic\c@table}

\newcommand{\@oldargs}{}
\newcommand{\@newargs}{}
%% add TOC entry for each paper
%\renewcommand{\title}[1]%
%{\renewcommand{\@title}{\rmaa@TitleStyle{#1}}%
%\addtocounter{chapter}{1}\addcontentsline{toc}{chapter}{#1}%
%}

%% Formatting of the table of contents
%% dimensions of various parts
\newlength{\TOCtextwidth}\setlength{\TOCtextwidth}{12.65cm}
\newlength{\TOCpagenumwidth}\setlength{\TOCpagenumwidth}{1.3cm}
\newlength{\TOCentrywidth}\setlength{\TOCentrywidth}{\TOCtextwidth}
\addtolength{\TOCentrywidth}{-\TOCpagenumwidth} 
\newlength{\TOCextramargin}\setlength{\TOCextramargin}{1.25cm}
\newcommand{\TOCbox}[1]{\hspace{\TOCextramargin}\parbox{\TOCtextwidth}{#1}\par}
\newlength{\TOCextramarginx}\setlength{\TOCextramarginx}{\TOCextramargin}\addtolength{\TOCextramarginx}{4.3cm}
\newlength{\TOCtextwidthx}\setlength{\TOCtextwidthx}{\TOCtextwidth}\addtolength{\TOCtextwidthx}{-4.3cm}
\newcommand{\TOCboxx}[1]{\hspace{\TOCextramarginx}\parbox{\TOCtextwidthx}{#1}\par}

\newcommand{\RMAAbannerimage}{rmsc}
\renewcommand\tableofcontents{%
  \renewcommand{\articlename}{CONTENTS}
  \begingroup
  \@begofarticlehook
  \if@twocolumn
     \@restonecoltrue\onecolumn
  \else
     \@restonecolfalse
  \fi
  %% 9pt for journal style
  \if@proceedings\else\small\fi
  \if@newstyle\small\fi %% same for new style
  \if@proceedings\thispagestyle{rmaafrontfirstpage}\else\thispagestyle{empty}\fi
  \shortauthor{Contents}\shorttitle{Contents}
  \vspace*{18mm}
  \if@printbanner
  %% put in the banner header (need to get better image of this)
     \TOCbox{\hspace*{-0.15cm}\includegraphics[width=0.9\TOCtextwidth]{\RMAAbannerimage}}
     \vspace*{8mm}
  %% volume and publication date (should really be in funny font)
     \TOCboxx{\@volumen\\[2mm] \@mes\ \@anyo}
  \else
  %% put in the banner header (need to get better image of this)
     \TOCbox{\hspace*{-0.15cm}\includegraphics[width=0.9\TOCtextwidth]{\RMAAbannerimage}}
     \vspace*{12.5mm}
  %% volume and publication date (should really be in funny font)
     \TOCboxx{~\\[2mm] ~}
  \fi
  \vspace*{16mm}
  %% This is the command that fiddles with the TOC entry so that the
  %% title is flushleft but the authors are flushright
  \renewcommand{\@fiddleauthors}[2]{%
    % first put the paper title
    \renewcommand{\\}{\relax} % temporarily turn off explicit line breaks just in case
    \renewcommand{\altaffilmark}[1]{\relax} % ditto with this
    \renewcommand{\textsuperscript}[1]{\relax} % ditto with this 
    ##1
    % now save the author list in a box
    \sbox\@tempboxa{##2}%
    % check that box is shorter than the total width available for the entry
    \ifdim \wd\@tempboxa <\TOCentrywidth
       % if so, then put it right here, flush against right margin
       \penalty10000\hskip1em plus10em\penalty0\hbox{}\nobreak\hfill\box\@tempboxa
       % Note that we force a space of at least 1em between title and authorlist, 
       % stretchable to 11em since we may have very short titles but still with an 
       % authorlist too long to go on the same line.  The only linebreak allowed is 
       % directly before the authorlist.
    \else
       % if it is too long, we set the original authorlist text in its own parbox
       \par\vskip1ex\parbox[b]{\TOCentrywidth}{\raggedleft ##2}
    \fi
    }
  \@starttoc{toc}%
  \if@restonecol\twocolumn\fi
  \frontlabel{RMAALastPage}
  \clearpage
  \thispagestyle{empty}
  \cleardoublepage % make sure we start on an odd page
  \endgroup

  %%% ALSO, do the HTML TOC
  %% The file that will be used to make the HTML TOC
  \@starttoc{htc}
  %% All the stuff in this file does nothing in the book style
  \addtocontents{htc}{\protect\HyperTOCmeeting{\@conftitle}{\@reunionplace}{\@reuniondate}}
  \addtocontents{htc}{\protect\HyperTOCeditors{\@editors}}
  \addtocontents{htc}{\protect\HyperTOCpreamble}
  \addtocontents{htc}{\protect\HyperTOCEntry{Table of Contents}{}{\articlename}}

  %% and the ADS file
  \ADS@openfile
  }
%% This is now much simpler
\renewcommand\l@chapter[2]{%TOC entry for articles
  \ifnum \c@tocdepth >\m@ne
    \vskip 1.0em \@plus\p@
    \TOCbox{%
      \parbox[b]{\TOCentrywidth}{#1}\parbox[b]{\TOCpagenumwidth}{\hfill#2}
    }
  \fi} 
\renewcommand\l@part[2]{%TOC entry for sessions: centered, no page number
  \ifnum \c@tocdepth >\m@ne
    \vskip 1.5em \@plus1.5em \goodbreak
    \TOCbox{%
      \parbox[b]{\TOCentrywidth}{\hfil#1\hfil}\hfill
    }
  \fi\nobreak} 

%% Formatting theindex
%% strongly encourage page breaks between index entries
\renewcommand{\@idxitem}{\par\penalty-1000\hangindent.45cm} 
\renewcommand{\indexspace}{\par\penalty-1000\vskip 10\p@ \@plus15\p@ \@minus3\p@\relax}
\renewenvironment{theindex}{%
  %% first, make sure we start on an odd page (06 Jun 2000)
  \clearpage\thispagestyle{empty}\cleardoublepage
  %% this looks a bit like an article but is two columns and without abstract, etc
  \renewcommand{\articlename}{INDEX}
  \title{Author Index}\shortauthor{Author Index}\shorttitle{Author Index}
  \listofauthors{}
  \addcontentsline{toc}{chapter}{\protect\@fiddleauthors{\@TOCtitle}{\rmaaTOCauthors}}%
  \addtocontents{htc}{\protect\HyperTOCEntry{\@rawtitle}{\@rawauthors}{\articlename}}
  \pagestyle{rmaa}
  \thispagestyle{rmaafirstpage}
  \@begofarticlehook
  %% 9pt for journal style
  \if@proceedings\else\small\fi
  \if@newstyle\small\fi
  \twocolumn[\@makeminititle]
  %% vertical spacing (see also definition of \@idxitem and \indexspace
  \parskip\z@ \@plus .6\p@\relax % allow quite a lot of leaway here to avoid breaks inside entries
  %% horizontal spacing, breaking and glue parameters (see also the stuff in the definition 
  %% of \indexauthor
  \parindent\z@ 
  \exhyphenpenalty10000 % don't break on explicit hyphens, it always looks bad
  \rightskip=0pt plus.2em % allow a tiny bit of leeway at the ends of lines
  \spaceskip.3333em plus.5em minus.15em % a lot of leeway between words
  \xspaceskip.5em plus.5em minus.4em % this seems to have no effect
  \let\item\@idxitem
  \renewcommand{\\}{\relax} % temporarily turn off explicit line breaks just in case
  \renewcommand{\altaffilmark}[1]{\relax} % ditto with this
  \renewcommand{\textsuperscript}[1]{\relax} % ditto with this
  }{\label{RMAALastPage}\clearpage} % clearpage necessary so that we know about \shortauthor



%% end definitions for book styles
\fi

%    \end{macrocode}
% In order to do away with plates, we sometimes want a figure to
% appear in color only in the online version. The following command
% can bve used in a figure caption to typeset a note to this effect. 
%    \begin{macrocode}
\newcommand\ColorOnlineOnly{%
  \textsc{Note: this figure is available in color in the electronic
    version of this article, obtainable from}
  \texttt{http://\discretionary{}{}{}www.\discretionary{}{}{}astroscu.\discretionary{}{}{}unam.\discretionary{}{}{}mx/\discretionary{}{}{}\~{}rmaa/}.}
%% Plates
%% Define 2 counters: plate is the number in the whole document, plateloop is used 
%% internally by \outputplates
\newcounter{plate}
\newcounter{plateloop}
%% Currently limited to 15 plates in document
%% when we do a label, we need to write to \newlabel entries to the .aux file, one 
%% that says RMAAPlateX so that it can be referred to in \ouputplates, and one that
%% says articlename:RMAAPlateX so that the dvips program can find it and asociate it
%% with the right article
\newcommand{\@platelabeli}{\oldlabel{RMAAPlate1}\label{RMAAPlate1}}
\newcommand{\@platelabelii}{\oldlabel{RMAAPlate2}\label{RMAAPlate2}}
\newcommand{\@platelabeliii}{\oldlabel{RMAAPlate3}\label{RMAAPlate3}}
\newcommand{\@platelabeliv}{\oldlabel{RMAAPlate4}\label{RMAAPlate4}}
\newcommand{\@platelabelv}{\oldlabel{RMAAPlate5}\label{RMAAPlate5}}
\newcommand{\@platelabelvi}{\oldlabel{RMAAPlate6}\label{RMAAPlate6}}
\newcommand{\@platelabelvii}{\oldlabel{RMAAPlate7}\label{RMAAPlate7}}
\newcommand{\@platelabelviii}{\oldlabel{RMAAPlate8}\label{RMAAPlate8}}
\newcommand{\@platelabelix}{\oldlabel{RMAAPlate9}\label{RMAAPlate9}}
\newcommand{\@platelabelx}{\oldlabel{RMAAPlate10}\label{RMAAPlate10}}
\newcommand{\@platelabelxi}{\oldlabel{RMAAPlate11}\label{RMAAPlate11}}
\newcommand{\@platelabelxii}{\oldlabel{RMAAPlate12}\label{RMAAPlate12}}
\newcommand{\@platelabelxiii}{\oldlabel{RMAAPlate13}\label{RMAAPlate13}}
\newcommand{\@platelabelxiv}{\oldlabel{RMAAPlate14}\label{RMAAPlate14}}
\newcommand{\@platelabelxv}{\oldlabel{RMAAPlate15}\label{RMAAPlate15}}
\newcommand{\@outputplatelabeli}{\oldlabel{RMAAOutputPlate1}}
\newcommand{\@outputplatelabelii}{\oldlabel{RMAAOutputPlate2}}
\newcommand{\@outputplatelabeliii}{\oldlabel{RMAAOutputPlate3}}
\newcommand{\@outputplatelabeliv}{\oldlabel{RMAAOutputPlate4}}
\newcommand{\@outputplatelabelv}{\oldlabel{RMAAOutputPlate5}}
\newcommand{\@outputplatelabelvi}{\oldlabel{RMAAOutputPlate6}}
\newcommand{\@outputplatelabelvii}{\oldlabel{RMAAOutputPlate7}}
\newcommand{\@outputplatelabelviii}{\oldlabel{RMAAOutputPlate8}}
\newcommand{\@outputplatelabelix}{\oldlabel{RMAAOutputPlate9}}
\newcommand{\@outputplatelabelx}{\oldlabel{RMAAOutputPlate10}}
\newcommand{\@outputplatelabelxi}{\oldlabel{RMAAOutputPlate11}}
\newcommand{\@outputplatelabelxii}{\oldlabel{RMAAOutputPlate12}}
\newcommand{\@outputplatelabelxiii}{\oldlabel{RMAAOutputPlate13}}
\newcommand{\@outputplatelabelxiv}{\oldlabel{RMAAOutputPlate14}}
\newcommand{\@outputplatelabelxv}{\oldlabel{RMAAOutputPlate15}}
\newcommand{\@platerefi}{\oldpageref{RMAAPlate1}}
\newcommand{\@platerefii}{\oldpageref{RMAAPlate2}}
\newcommand{\@platerefiii}{\oldpageref{RMAAPlate3}}
\newcommand{\@platerefiv}{\oldpageref{RMAAPlate4}}
\newcommand{\@platerefv}{\oldpageref{RMAAPlate5}}
\newcommand{\@platerefvi}{\oldpageref{RMAAPlate6}}
\newcommand{\@platerefvii}{\oldpageref{RMAAPlate7}}
\newcommand{\@platerefviii}{\oldpageref{RMAAPlate8}}
\newcommand{\@platerefix}{\oldpageref{RMAAPlate9}}
\newcommand{\@platerefx}{\oldpageref{RMAAPlate10}}
\newcommand{\@platerefxi}{\oldpageref{RMAAPlate11}}
\newcommand{\@platerefxii}{\oldpageref{RMAAPlate12}}
\newcommand{\@platerefxiii}{\oldpageref{RMAAPlate13}}
\newcommand{\@platerefxiv}{\oldpageref{RMAAPlate14}}
\newcommand{\@platerefxv}{\oldpageref{RMAAPlate15}}
\newcounter{figplatei}
\newcounter{figplateii}
\newcounter{figplateiii}
\newcounter{figplateiv}
\newcounter{figplatev}
\newcounter{figplatevi}
\newcounter{figplatevii}
\newcounter{figplateviii}
\newcounter{figplateix}
\newcounter{figplatex}
\newcounter{figplatexi}
\newcounter{figplatexii}
\newcounter{figplatexiii}
\newcounter{figplatexiv}
\newcounter{figplatexv}
%% Printed output should use the counter `plate'
\renewcommand{\theplate}{\@arabic\c@plate}
\renewcommand{\theplateloop}{\@arabic\c@plateloop}
%% Command to output all the saved plates in this paper
\newcommand{\outputplates}{%
  \setcounter{plateloop}{0}
  \if@book\else\@endofarticlehook\fi % not entirely sure why, but this
                                     % seems to be needed
  \clearpage % first output the last page of text
  % now set pagestyle for blank even pages and make sure we start on an odd page
  \pagestyle{empty}
  \cleardoublepage
  \onecolumn
  \parindent\z@
  % first put in an extra sheet
  % the main purpose of this is so there is always a blank even page before 
  % each plate for filling in the PS file if the article has an odd number
  % of pages (see splitbook.sh)
  \null\vfil\centerline{PLATES}\vfil\clearpage\pagestyle{rmaaplate}\cleardoublepage
  % output each plate in turn
  \@whilenum\value{plateloop}<\value{plate}\do{%
    \addtocounter{plateloop}{1}
    % restore the figure counter from \@figcounti, \@figcountii, etc
    \setcounter{figure}{\value{figplate\roman{plateloop}}}
    % spew out the plate (in its own group for safety)
    {\csname @plate\roman{plateloop}\endcsname\par}
    }
  \pagestyle{rmaa}
  % reset the page number as if nothing had happened
  \addtocounter{page}{-\value{plate}}
  \addtocounter{page}{-\value{plate}}
  }
\newcommand\figcaption{%
  \def\@captype{figure}
  \renewcommand{\@makecaption}{\@makefigcap}
  \caption}
\newcommand\tabcaption{%
  \def\@captype{table}
  \renewcommand{\@makecaption}{\@maketabcap}
  \caption}

\newcommand{\plate}[2][1]{%
  %% User command to make a plate page: 
  %% Mandatory argument is what goes on the plate - this should not 
  %% include any floating environments (e.g. figure). Instead, \figcaption 
  %% should be used to do the captions. 
  %% Optional argument is the number of figures (i.e. captions) on the plate
  %% page (default is 1). This is necessary for incrementing the figure counter. 
  %%
  %% Note that we have to use a \gdef rather than an \xdef when we 
  %% save the plate output in a command. This means that in the book styles when it 
  %% comes to spewing out the plate, we will have forgotten about any user macros 
  %% that may have been defined. This can be got round by repeating the macros inside the plate.
  %% (Presumably, one would use a \newcommand rather than \renewcommand if they were new.)
  % increment global and local counters
  \addtocounter{plate}{1}
  % write appropriate label to .aux file
  \csname @platelabel\roman{plate}\endcsname
  % define a global macro called \@platei, \@plateii, \@plateiii, \@plateiv, ...
  \expandafter\gdef\csname @plate\roman{plate}\endcsname{%
    \csname @outputplatelabel\roman{plateloop}\endcsname
    \null\vfil #2 \vfil
    \cleardoublepage
    }
  % save the current shortauthor and shorttitle
  % need to be \xdef to take current values of \@shortauthor, etc 
  \expandafter\xdef\csname @shorta\roman{plate}\endcsname{\@shortauthor}
  \expandafter\xdef\csname @shortt\roman{plate}\endcsname{\@shorttitle}
  % save the current value of the figure counter for later use in \outputplates
  \setcounter{figplate\roman{plate}}{\value{figure}}
  % increment the figure counter according to the optional argument
  \addtocounter{figure}{#1}
  }
%% page style for plate pages 
\newcommand{\ps@rmaaplate}{%
  % Plate X (see page Y) in the foot
  \renewcommand\@oddfoot{\hfil\csname @shorta\roman{plateloop}\endcsname\ 
    (see page~\csname @plateref\roman{plateloop}\endcsname)\hfil}%
  \renewcommand\@evenfoot{\relax}%
  % nothing in even heads (blank pages)
  \renewcommand\@evenhead{\relax}
  % short title and page num in the odd heads
  \renewcommand\@oddhead{\hfil\csname @shortt\roman{plateloop}\endcsname%
    \hfil\llap{PLATE~\theplateloop}}
  }





%% Things for the User Guide
\if@guide
%% stolen from doc.sty
\def\MakeShortVerb#1{%
  \expandafter\ifx\csname cc\string#1\endcsname\relax
  \@shortvrbinfo{Made }{#1}%
  \add@special{#1}%
  \expandafter
  \xdef\csname cc\string#1\endcsname{\the\catcode`#1}%
  \begingroup
  \catcode`\~\active  \lccode`\~`#1%
  \lowercase{%
    \global\expandafter\let
    \csname ac\string#1\endcsname~%
    \gdef~{\verb~}}%
  \endgroup
  \global\catcode`#1\active
  \else
  \@shortvrbinfo\@empty{#1 already}%
  \fi}
\def\DeleteShortVerb#1{%
  \expandafter\ifx\csname cc\string#1\endcsname\relax
  \else
  \@shortvrbinfo{Deleted }{#1 as}%
  \rem@special{#1}%
  \global\catcode`#1\csname cc\string#1\endcsname
  \global \expandafter\let \csname cc\string#1\endcsname \relax
  \ifnum\catcode`#1=\active
  \begingroup
  \catcode`\~\active   \lccode`\~`#1%
  \lowercase{%
    \global\expandafter\let\expandafter~%
    \csname ac\string#1\endcsname}%
  \endgroup \fi \fi}
\def\@shortvrbinfo#1#2{%
  \PackageInfo{doc}{%
    #1\expandafter\@gobble\string#2 a short reference
    for \string\verb}}
\def\add@special#1{%
  \rem@special{#1}%
  \expandafter\gdef\expandafter\dospecials\expandafter
  {\dospecials \do #1}%
  \expandafter\gdef\expandafter\@sanitize\expandafter
  {\@sanitize \@makeother #1}}
\def\rem@special#1{%
  \def\do##1{%
    \ifnum`#1=`##1 \else \noexpand\do\noexpand##1\fi}%
  \xdef\dospecials{\dospecials}%
  \begingroup
  \def\@makeother##1{%
    \ifnum`#1=`##1 \else \noexpand\@makeother\noexpand##1\fi}%
  \xdef\@sanitize{\@sanitize}%
  \endgroup}
%% define | as shortcut for \verbatim
\MakeShortVerb{\|}
%% comand to specify arguments to LaTeX commands
\newcommand{\rmArg}[1]{%
  \ensuremath{\langle}{\scriptsize \textsf{\MakeUpperCase{#1}}}\ensuremath{\rangle}}
%\newcommand{\rmArg}[1]{\ensuremath{\langle}\textsc{#1}\ensuremath{\rangle}}
\providecommand{\BibTeX}{{\rmfamily B\kern-.05em%
    \textsc{i\kern-.025em b}\kern-.08em%
    T\kern-.1667em\lower.7ex\hbox{E}\kern-.125emX}}
%% Environment for code examples
\newenvironment{Code}{%
  \begin{list}{}{\setlength{\leftmargin}{0.05\columnwidth}} 
  \item[]}{\end{list}}
\fi
%</package>
%    \end{macrocode}

% \section{User commands for astronomical articles}
% \label{sec:user-commands}
% Some random definitions. These go both into the main class file and
% also into |adshtml.sty|.
% 
% Definitions of the journal name in different incarnations.  
%    \begin{macrocode}
%<*package|adshtml-style>
\newcommand{\RMAA}{RMxAA}
\newcommand{\RMAAlong}{Revista Mexicana de Astronom\'{\i}a y Astrof\'{\i}sica}
\newcommand{\RMAASC}{RMxAC}
\newcommand{\RMAASClong}{\RMAA\ (Serie de Conferencias)}
%    \end{macrocode}
% \begin{macro}{\ibidrule}
% This puts a horizontal rule with a period after it. It is used when
% the author list of a reference is identical to the preceding one. 
%    \begin{macrocode}
\newcommand{\ibidrule}{\rule[-0.03cm]{4.7em}{.01cm}.\ }
%    \end{macrocode}
% \end{macro}
% \begin{macro}{\ThisVolume}
% References to other articles in the same volume can be coded as,
% e.g., |\bibitem{} Bloggs, J. \ThisVolume{bloggs}|.
%    \begin{macrocode}
\newcommand{\ThisVolume}[1]{%
  \@year, \if@proceedings\RMAASC\else\RMAA\fi, \@volume,
  \if@book\oldpageref{#1:RMAAFirstPage}\else 000\fi\ (this volume)}
%    \end{macrocode}
% \end{macro}
% Other stuff that may be useful in ref lists.
%    \begin{macrocode}
\newcommand{\etal}{et al.\@}
%    \end{macrocode}
% Various new commands from AAS macros
%    \begin{macrocode}
\def\acknowledgments{\vskip 3ex plus .8ex minus .4ex}
\let\acknowledgements=\acknowledgments                  % second spelling
\let\jnl@style=\rm
\def\ref@jnl#1{{\jnl@style#1}}
\DeclareRobustCommand\aj{\ref@jnl{AJ}}                   
\DeclareRobustCommand\araa{\ref@jnl{ARA\&A}}             
\DeclareRobustCommand\apj{\ref@jnl{ApJ}}                 
\DeclareRobustCommand\apjl{\ref@jnl{ApJ}}                
\DeclareRobustCommand\apjs{\ref@jnl{ApJS}}               
\DeclareRobustCommand\ao{\ref@jnl{Appl.~Opt.}}           
\DeclareRobustCommand\apss{\ref@jnl{Ap\&SS}}             
\DeclareRobustCommand\aap{\ref@jnl{A\&A}}                
\DeclareRobustCommand\aapr{\ref@jnl{A\&A~Rev.}}          
\DeclareRobustCommand\aaps{\ref@jnl{A\&AS}}              
\DeclareRobustCommand\azh{\ref@jnl{AZh}}                 
\DeclareRobustCommand\baas{\ref@jnl{BAAS}}               
\DeclareRobustCommand\jrasc{\ref@jnl{JRASC}}             
\DeclareRobustCommand\memras{\ref@jnl{MmRAS}}            
\DeclareRobustCommand\mnras{\ref@jnl{MNRAS}}             
\DeclareRobustCommand\pra{\ref@jnl{Phys.~Rev.~A}}        
\DeclareRobustCommand\prb{\ref@jnl{Phys.~Rev.~B}}        
\DeclareRobustCommand\prc{\ref@jnl{Phys.~Rev.~C}}        
\DeclareRobustCommand\prd{\ref@jnl{Phys.~Rev.~D}}        
\DeclareRobustCommand\pre{\ref@jnl{Phys.~Rev.~E}}        
\DeclareRobustCommand\prl{\ref@jnl{Phys.~Rev.~Lett.}}    
\DeclareRobustCommand\pasp{\ref@jnl{PASP}}               
\DeclareRobustCommand\pasj{\ref@jnl{PASJ}}               
\DeclareRobustCommand\qjras{\ref@jnl{QJRAS}}             
\DeclareRobustCommand\skytel{\ref@jnl{S\&T}}             
\DeclareRobustCommand\solphys{\ref@jnl{Sol.~Phys.}}      
\DeclareRobustCommand\sovast{\ref@jnl{Soviet~Ast.}}      
\DeclareRobustCommand\ssr{\ref@jnl{Space~Sci.~Rev.}}     
\DeclareRobustCommand\zap{\ref@jnl{ZAp}}                 
\DeclareRobustCommand\nat{\ref@jnl{Nature}}              
\DeclareRobustCommand\iaucirc{\ref@jnl{IAU~Circ.}}
\DeclareRobustCommand\aplett{\ref@jnl{Astrophys.~Lett.}}
\DeclareRobustCommand\apspr{\ref@jnl{Astrophys.~Space~Phys.~Res.}}
\DeclareRobustCommand\bain{\ref@jnl{Bull.~Astron.~Inst.~Netherlands}}
\DeclareRobustCommand\fcp{\ref@jnl{Fund.~Cosmic~Phys.}}
\DeclareRobustCommand\gca{\ref@jnl{Geochim.~Cosmochim.~Acta}}
\DeclareRobustCommand\grl{\ref@jnl{Geophys.~Res.~Lett.}}
\DeclareRobustCommand\jcp{\ref@jnl{J.~Chem.~Phys.}}      
\DeclareRobustCommand\jgr{\ref@jnl{J.~Geophys.~Res.}}    
\DeclareRobustCommand\jqsrt{\ref@jnl{J.~Quant.~Spec.~Radiat.~Transf.}}
\DeclareRobustCommand\memsai{\ref@jnl{Mem.~Soc.~Astron.~Italiana}}
\DeclareRobustCommand\nphysa{\ref@jnl{Nucl.~Phys.~A}}
\DeclareRobustCommand\physrep{\ref@jnl{Phys.~Rep.}}
\DeclareRobustCommand\physscr{\ref@jnl{Phys.~Scr}}
\DeclareRobustCommand\planss{\ref@jnl{Planet.~Space~Sci.}}
\DeclareRobustCommand\procspie{\ref@jnl{Proc.~SPIE}}
\let\astap=\aap
\let\apjlett=\apjl
\let\apjsupp=\apjs
\let\applopt=\ao
\def\phn{\phantom{0}}
\def\phd{\phantom{.}}
\def\phs{\phantom{\ensuremath{-}}}
\def\phm#1{\phantom{#1}}
%    \end{macrocode}
% All these are redone as robust and using more modern commands where
% possible. The robustness is so they will not get expanded if used in
% abstracts and titles when these are passed to the ADS file. 
% \changes{3.16}{2003/04/07}{
%  Added \cs{Color} command for doing photometric colors.}
%    \begin{macrocode}
%%\def\sun{\hbox{$\odot$}}
\DeclareRobustCommand\sun{\ensuremath{\odot}}
%%\def\earth{\hbox{$\oplus$}}
\DeclareRobustCommand\earth{\ensuremath{\oplus}}
%%\def\lesssim{\mathrel{\hbox{\rlap{\hbox{\lower4pt\hbox{$\sim$}}}\hbox{$<$}}}}
%%\def\gtrsim{\mathrel{\hbox{\rlap{\hbox{\lower4pt\hbox{$\sim$}}}\hbox{$>$}}}}
\DeclareRobustCommand\lesssim{\mathrel{\hbox{\smaller\smaller\rlap{\hbox{\lower3pt\hbox{$\sim$}}}\raise1.5pt\hbox{$<$}}}}
\DeclareRobustCommand\gtrsim{\mathrel{\hbox{\smaller\smaller\rlap{\hbox{\lower3pt\hbox{$\sim$}}}\raise1.5pt\hbox{$>$}}}}
%%\def\lta{\mathrel{\spose{\lower 3pt\hbox{$\mathchar"218$}}
%%     \raise 2.0pt\hbox{$\mathchar"13C$}}}
%%\def\gta{\mathrel{\spose{\lower 3pt\hbox{$\mathchar"218$}}
%%     \raise 2.0pt\hbox{$\mathchar"13E$}}}
\DeclareRobustCommand\sq{\hbox{\rlap{$\sqcap$}$\sqcup$}}
%%\def\arcdeg{\hbox{$^\circ$}}
%%\def\arcmin{\hbox{$^\prime$}}
%%\def\arcsec{\hbox{$^{\prime\prime}$}}
\DeclareRobustCommand\arcdeg{\ensuremath{^\circ}}
\DeclareRobustCommand\arcmin{\ensuremath{^\prime}}
\DeclareRobustCommand\arcsec{\ensuremath{^{\prime\prime}}}
%%\def\fd{\hbox{$.\!\!^{\rm d}$}}
%%\def\fh{\hbox{$.\!\!^{\rm h}$}}
%%\def\fm{\hbox{$.\!\!^{\rm m}$}}
%%\def\fs{\hbox{$.\!\!^{\rm s}$}}
%%\def\fdg{\hbox{$.\!\!^\circ$}}
%%\def\farcm{\hbox{$.\mkern-4mu^\prime$}}
%%\def\farcs{\hbox{$.\!\!^{\prime\prime}$}}
%%\def\fp{\hbox{$.\!\!^{\scriptscriptstyle\rm p}$}}
\DeclareRobustCommand\fd{\ensuremath{.\!\!^\mathrm{d}}}
\DeclareRobustCommand\fh{\ensuremath{.\!\!^\mathrm{h}}}
\DeclareRobustCommand\fm{\ensuremath{.\!\!^\mathrm{m}}}
\DeclareRobustCommand\fs{\ensuremath{\mkern0.5mu.\mkern-5mu^\mathrm{s}\mkern-1mu}}
\DeclareRobustCommand\fdg{\ensuremath{.\!\!^\circ}}
\DeclareRobustCommand\farcm{\ensuremath{.\mkern-4mu^\prime}}
\DeclareRobustCommand\farcs{\ensuremath{\mkern1mu.\mkern-6mu^{\prime\prime}\mkern-1.5mu}}
%%\DeclareRobustCommand\farcs{\ensuremath{.\!\!^{\prime\prime}}}
\DeclareRobustCommand\fp{\ensuremath{.\!\!^{\scriptscriptstyle\rm p}}}
%%\def\micron{\hbox{$\mu$m}}
\DeclareRobustCommand\micron{\ensuremath{\mu\mathrm{m}}}
\let\la=\lesssim                        
\let\ga=\gtrsim
%%\def\case#1#2{\hbox{$\frac{#1}{#2}$}}
%%\def\slantfrac#1#2{\hbox{$\,^#1\!/_#2$}}
%%\def\onehalf{\slantfrac{1}{2}}
%%\def\onethird{\slantfrac{1}{3}}
%%\def\twothirds{\slantfrac{2}{3}}
%%\def\onequarter{\slantfrac{1}{4}}
%%\def\threequarters{\slantfrac{3}{4}}
\DeclareRobustCommand\case[2]{\ensuremath{\frac{#1}{#2}}}
\DeclareRobustCommand\slantfrac[2]{\ensuremath{\,^#1\!/_#2}}
\DeclareRobustCommand\onehalf{\slantfrac{1}{2}}
\DeclareRobustCommand\onethird{\slantfrac{1}{3}}
\DeclareRobustCommand\twothirds{\slantfrac{2}{3}}
\DeclareRobustCommand\onequarter{\slantfrac{1}{4}}
\DeclareRobustCommand\threequarters{\slantfrac{3}{4}}
%%\def\ubvr{\hbox{$U\!BV\!R$}}            
%%\def\ub{\hbox{$U\!-\!B$}}               
%%\def\bv{\hbox{$B\!-\!V$}}               
%%\def\vr{\hbox{$V\!-\!R$}}               
%%\def\ur{\hbox{$U\!-\!R$}}
%    \end{macrocode}
% New commands for ``photometric'' colors 
%    \begin{macrocode}
\DeclareRobustCommand\ubvr{\ensuremath{U\!BV\!R}}            
\DeclareRobustCommand\Color[2]{%
  \ensuremath{%
    #1\!\ifx#1V\!\fi-\!#2
  }}
\DeclareRobustCommand\ub{\Color{U}{B}}
\DeclareRobustCommand\bv{\Color{B}{V}}
\DeclareRobustCommand\vr{\Color{V}{R}}
\DeclareRobustCommand\ur{\Color{U}{R}}
%% New version of the \ion command, that actually works! 
%    \end{macrocode}
% \begin{macro}{\ion}
% This version of the |\ion| command allows you to write either
% |\ion{N}{4}| or |\ion{N}{iv}| in order to get N~IV. 
%    \begin{macrocode}
\DeclareRobustCommand\ion[2]{%
  \mbox{#1\kern0.2em%
    \smaller\rmfamily%
    \edef\@tempa{\@car#2\@nil}%
    \ifcat1\@tempa%
    \@Roman{#2}%
    \else%
    \uppercase{#2}%
    \fi}}
\newcounter{IonStage}
\DeclareRobustCommand\plainion[2]{\setcounter{IonStage}{#2}#1 \Roman{IonStage}}
%%\def\ion#1#2{\ensuremath{\mathrm{#1}\;}{\protect\small\rm\@Roman{#2}}}
%    \end{macrocode}
% \end{macro}
%    \begin{macrocode}
\DeclareRobustCommand\nodata{$\cdots$}
%    \end{macrocode}
% These are needed by the |\ion| command. Isn't there an easier way of
% doing this?
%    \begin{macrocode}
\let\@small=\small
\let\@normalsize=\normalsize
\let\@footnotesize=\footnotesize
\let\@scriptsize=\scriptsize
\let\@tiny=\tiny
\let\@large=\large
\let\@Large=\Large
\let\@LARGE=\LARGE
\let\@huge=\huge
\let\@Huge=\Huge
\newcommand{\smaller}{\small}
\renewcommand{\normalsize}{\renewcommand{\smaller}{\small}\@normalsize}
\renewcommand{\small}{\renewcommand{\smaller}{\footnotesize}\@small}
\renewcommand{\footnotesize}{\renewcommand{\smaller}{\scriptsize}\@footnotesize}
\renewcommand{\scriptsize}{\renewcommand{\smaller}{\tiny}\@scriptsize}
\renewcommand{\tiny}{\renewcommand{\smaller}{\tiny}\@tiny}
\renewcommand{\large}{\renewcommand{\smaller}{\normalsize}\@large}  
\renewcommand{\Large}{\renewcommand{\smaller}{\large}\@Large}  
\renewcommand{\LARGE}{\renewcommand{\smaller}{\Large}\@LARGE}  
\renewcommand{\huge}{\renewcommand{\smaller}{\LARGE}\@huge}  
\renewcommand{\Huge}{\renewcommand{\smaller}{\huge}\@Huge}  
%</package|adshtml-style>
%    \end{macrocode}


% \section{Other files}

% \subsection{The driver and make files}
% \label{sec:driver_make}
% These are put here for completeness, although they still need to be
% distributed separately since otherwise the user won't be able to
% unpack this file into its constituent parts.
% 
% The |Makefile| has two principal rules: |docs|, and |macros|. The
% |docs| rule produces a PDF and gzipped PS file of the documentation.
% I still haven't perfected the combination of rules required to get
% it to run \LaTeX{} enough times in all circumstances. The |macros|
% rule runs \LaTeX{} on |rmaa.ins| in order to unpack |rmaa.cls| and
% all the other little files (including the |Makefile| and the
% |rmaa.ins| files themselves\dots 
% 
% First set up lists of various types of file. 
%    \begin{macrocode}
%<*makefile>
DOCFILES=rmaa.pdf
MACRO_LTX_FILES=rmaa.cls rmaacite.sty adshtml.sty
SHELLSCRIPTS=makeads.sh
MACRO_OTHER_FILES=adshtml.perl sample.latex2html-init
MAKEFILE=Makefile
MACROFILES=$(MACRO_LTX_FILES) $(SHELLSCRIPTS) $(MACRO_OTHER_FILES)         
SKELETONFILES=ADS.tex 
%    \end{macrocode}
% Install the files read only into the package directory, specified via a path relative to the source directory. \changes{4.0}{2019/11/14}{Reorganize folder layout so that hardly anything is in the top level.  That way, things are clearer on github.}
%    \begin{macrocode}
INSTALL=install 
INSTALL_DIR=../packages
%    \end{macrocode}
% The default rule is to make everything. Installation must be
% explicitly asked for.
%    \begin{macrocode}
all: docs macros

install: install-docs install-macros

%    \end{macrocode}
% Most files are installed read-only to prevent editing. Skeleton
% files must be writeable. Must be careful not to overwrite the
% Makefile in the parent directory. 
%    \begin{macrocode}
install-macros: macros
        $(INSTALL) -m 444 $(MACRO_LTX_FILES) $(MACRO_OTHER_FILES) $(INSTALL_DIR)
        $(INSTALL) -m 555 $(SHELLSCRIPTS) $(INSTALL_DIR)
        $(INSTALL) -m 644 $(SKELETONFILES) $(INSTALL_DIR)

install-docs: docs
        $(INSTALL) rmaa.pdf $(INSTALL_DIR)/rm-fulldocs.pdf

%    \end{macrocode}
% Make the PDF and gzipped PS files. 
%    \begin{macrocode}
docs: $(DOCFILES)

%    \end{macrocode}
% Make everything else if either the |.dtx| or |.ins| files have
% changed. 
%    \begin{macrocode}
macros: rmaa.dtx rmaa.ins
%    \end{macrocode}
% Ensure that all the macro files can be written to then recreate
% them. The |touch| is there in case any of the files don't exist
% yet. 
%    \begin{macrocode}
        touch $(MACROFILES) $(MAKEFILE)
        chmod u+w $(MACROFILES) $(MAKEFILE) rmaa.ins
        latex "\input{rmaa.ins}"
%    \end{macrocode}
% \TeX{} automatically adds the |.tex| extension to extensionless
% files. We fix this and take the opportunity to convert
% back to TABs all that it has lovingly converted to spaces. 
%    \begin{macrocode}
        cp -f $(MAKEFILE) $(MAKEFILE).BCK && unexpand $(MAKEFILE).tex >| $(MAKEFILE) && rm $(MAKEFILE).tex
%    \end{macrocode}
% Make the shell scripts executable and everything unwriteable (so
% they don't get hand-edited by mistake). However, undo this for the
% files that the user is \emph{supposed} to modify. 
%    \begin{macrocode}
        chmod a+x $(SHELLSCRIPTS)
        chmod a-w $(MACROFILES) $(MAKEFILE)
        chmod u+w $(SKELETONFILES)

%    \end{macrocode}
% We now use PDF\LaTeX{} for this.
%    \begin{macrocode}
rmaa.idx: rmaa.dtx
        @echo "REGENERATING .idx FILE"
        pdflatex "\nonstopmode\input{rmaa.dtx}"

rmaa.aux: rmaa.dtx rmaa.ind rmaa.gls
        @echo "REGENERATING .aux FILE"
        pdflatex "\nonstopmode\input{rmaa.dtx}"

rmaa.ind: rmaa.idx
        @echo "REGENERATING INDEX"
        makeindex -s gind.ist rmaa

rmaa.gls: rmaa.glo
        @echo "REGENERATING CHANGELOG"
        makeindex -s gglo.ist -o rmaa.gls rmaa.glo

rmaa.pdf: rmaa.aux rmaa.ind rmaa.gls
        @echo "REGENERATING .pdf FILE"
        pdflatex "\nonstopmode\input{rmaa.dtx}"

clean: 
        rm -f rmaa.aux rmaa.dvi rmaa.ind rmaa.gls
%</makefile>
%    \end{macrocode}

% The |rmaa.ins| file uses the \textsf{docstrip} package to split out
% all the files and strip off this documentation.
%    \begin{macrocode}
%<*ins>
%%
%% Installation batch file for the RMAA LaTeX package
%%
\def\batchfile{rmaa.ins}
\input docstrip.tex
\preamble

PLEASE DO NOT EDIT THIS FILE! It is generated automatically from
the file rmaa.dtx using the accompanying Makefile and rmaa.ins 
Any hand-edited changes will be lost

A LaTeX document class for the RMxAA

Author: William Henney (w.henney@irya.unam.mx)

Latest version available from: http://www.irya.unam.mx/rmaa

Copyright (C) 1999-2019 Instituto de Astronomia, UNAM, Mexico

Permission is hereby granted, free of charge, to any person obtaining
a copy of this software and associated documentation files (the
"Software"), to deal in the Software without restriction, including
without limitation the rights to use, copy, modify, merge, publish,
distribute, sublicense, and/or sell copies of the Software, and to
permit persons to whom the Software is furnished to do so, subject to
the following conditions:

The above copyright notice and this permission notice shall be
included in all copies or substantial portions of the Software.

THE SOFTWARE IS PROVIDED "AS IS", WITHOUT WARRANTY OF ANY KIND,
EXPRESS OR IMPLIED, INCLUDING BUT NOT LIMITED TO THE WARRANTIES OF
MERCHANTABILITY, FITNESS FOR A PARTICULAR PURPOSE AND
NONINFRINGEMENT. IN NO EVENT SHALL THE AUTHORS OR COPYRIGHT HOLDERS BE
LIABLE FOR ANY CLAIM, DAMAGES OR OTHER LIABILITY, WHETHER IN AN ACTION
OF CONTRACT, TORT OR OTHERWISE, ARISING FROM, OUT OF OR IN CONNECTION
WITH THE SOFTWARE OR THE USE OR OTHER DEALINGS IN THE SOFTWARE.
\endpreamble
%    \end{macrocode}
% Clobber files without asking
%    \begin{macrocode}
\askforoverwritefalse
%    \end{macrocode}
% Make the \LaTeX{} macro files (class and styles)
%    \begin{macrocode}
\generate{% 
  \file{rmaa.cls}{\from{rmaa.dtx}{package}}
  \file{rmaa.ins}{\from{rmaa.dtx}{ins}}
  \file{rmaacite.sty}{\from{rmaa.dtx}{rmaacite}}
  \file{adshtml.sty}{\from{rmaa.dtx}{adshtml-style}}
  \file{ADS.tex}{\from{rmaa.dtx}{ads-container}}
  }
%    \end{macrocode}
% Make the other non-\LaTeX{} files. These are done separately since
% they can't use `\%\%' as a comment character for the
% preamble/postamble (and Perl files can't have the postamble I don't
% think). 
%    \begin{macrocode}
\generate{\nopreamble\nopostamble% 
  \file{Makefile}{\from{rmaa.dtx}{makefile}}
  \file{adshtml.perl}{\from{rmaa.dtx}{adshtml-perl}}
  \file{makeads.sh}{\from{rmaa.dtx}{ads-shell-script}}
  \file{sample.latex2html-init}{\from{rmaa.dtx}{latex2html-init}}
  }

%% Local Variables: 
%% mode: LaTeX
%% mode: Font-lock
%% End:
%</ins>
%    \end{macrocode}

% \subsection{Extended citation syntax}
% Modified from a file called mnrascite.sty
% I no longer remember where I stole the original version of this from
% 
% Alterations to the cite commands to make them do more sensible things.
% N.B. the bibliography style should produce entries which have labels where
% the date is delimited by angle brackets '<>' for these macros to work
% properly. They will not give very sensible error messages if you do not do 
% this properly.
% 
% A new command |\scite{key}| is defined which produces a citation which is 
% suitable for insertion in a sentence. i.e., it does not put brackets around 
% the whole citation, only the date. Also |\pcite{key}| produce output with no
% parentheses whatsoever. 
% 
% Should redefine |\@bibitem| really to have any hope of compatibility
% the command |\lbibcite| is defined to deal with the citations.
%    \begin{macrocode}
%<*rmaacite>
\def\bibitem[#1<#2>]#3{\par\hangindent=1.5em\hangafter=1
  \if@filesw
  {\def\protect##1{\string ##1\space}
    \immediate\write\@auxout{\string\lbibcite{#3}{#1}{#2}}}
  \fi}
\@redefinebibitem % what it says
%    \end{macrocode}
% Note that the |\lbibcite| defines two types of citation one contains the names,
% the other the year
%    \begin{macrocode}
\def\lbibcite#1#2#3{\global\@namedef{b@#1}{#2}
                    \global\@namedef{by@#1}{#3}}

\def\@bcite#1#2{({#1\if@tempswa , #2\fi})}
\def\@pcite#1#2{#1\if@tempswa , #2\fi}
%    \end{macrocode}
% This declaration contains what is put between the name and the date
% \changes{3.11}{2003/01/31}{Moved the \cs{articlename} to after the key for consistency
% with the new natbib stuff.}
%    \begin{macrocode}
\def\@brcite{\ (}
\def\@ebrcite{)}
\def\@fmttop#1{\@fmtcite[\csname #1\endcsname]}
\def\@citex[#1]#2{%
%%  \if@filesw\immediate\write\@auxout{\string\citation{#2}}\fi
%% WJH added 16 Jan 2002
  \@for\@citeb:=#2\do{%
    \edef\@citeb{\expandafter\@firstofone\@citeb:\articlename}%
    \if@filesw\immediate\write\@auxout{\string\citation{\@citeb}}\fi}%
%% to here
  \def\@citea{}%
  \@cite{\@for\@citeb:=#2\do
    {%
      \edef\@citeb{\expandafter\@firstofone\@citeb:\articlename}%
      \@citea\def\@citea{;\ }\@ifundefined
      {b@\@citeb}{{\bf ?}\@warning {Citation `\@citeb' on page \thepage \space undefined}}%
      {\csname b@\@citeb\endcsname\@citebro\csname by@\@citeb\endcsname\@citebrc}}}%
  {#1}}

\def\cite{\@ifnextchar [{\let\@citebro=\space
                          \let\@citebrc=\relax
                          \let\@cite=\@bcite\@tempswatrue\@citex}
                        {\let\@citebro=\space
                          \let\@citebrc=\relax
                          \let\@cite=\@bcite\@tempswafalse\@citex[]}}
\def\pcite{\@ifnextchar [{\let\@citebro=\space
                          \let\@citebrc=\relax
                          \let\@cite=\@pcite\@tempswatrue\@citex}
                        {\let\@citebro=\space
                          \let\@citebrc=\relax
                          \let\@cite=\@pcite\@tempswafalse\@citex[]}}
\def\scite{\@ifnextchar [{\let\@citebro=\@brcite
                          \let\@citebrc=\@ebrcite
                          \let\@cite=\@pcite\@tempswatrue\@citex}
                        {\let\@citebro=\@brcite
                          \let\@citebrc=\@ebrcite
                          \let\@cite=\@pcite\@tempswafalse\@citex[]}}

\renewenvironment{thebibliography}
  {\section*{REFERENCES}\vspace*{\afterbiblioskip}\small
    \if@book\RereadBiblio\fi
    \sloppy\clubpenalty4000\widowpenalty4000
    \sfcode`\.=1000\relax
    \setlength\parindent{\z@}\setlength\parskip{\z@}}
  {\relax\normalsize\@outputaddressesifnecessary} 
%</rmaacite>
%    \end{macrocode}

% \subsection{Auxilliary files for generating ADS data}

% The ADS data are written by \LaTeX{} to a |.ads| file. This file
% then has to be included in the following file (|ADS.tex|),
% which is run through \LaTeX{}2HTML (see below). 

%    \begin{macrocode}
%<*ads-container>
%%
%% Instructions:
%% Edit the \input{} line below to reflect the name of your file. 
\documentclass{article}
\usepackage[latin1]{inputenc}
\usepackage{adshtml}
\newcommand{\ADSfield}[2]
{\%#1  #2\\}
%%\newcommand{\ADSvfield}[2]
%%{\%#1 \begin{verbatim}#2\end{verbatim}}
\newenvironment{ADSarticle}{}{\par\bigskip}
\renewcommand{\bibitem}[2][]{\par}
\setlength{\parindent}{0pt}
\renewcommand{\ThisVolume}[1]{%  
  2001, \RMAASC, 12, \oldpageref{#1:RMAAFirstPage}}  
 

\begin{document}
\makeatletter
\input{rmbooksamp.ads}
\end{document}
%</ads-container>
%    \end{macrocode}

% This is the |adshtml.sty| style file used by |ADS.tex|. It just
% defines a few 
% macros that might be used in abstracts, etc (really it should define
% them all). This now gets all its code from
% Sec~\ref{sec:user-commands}. 
%    \begin{macrocode}
%<*adshtml-style>
%% Nothing else needs to go here now. 
%</adshtml-style>
%    \end{macrocode}
% The init file for latex2html.
%    \begin{macrocode}
%<*latex2html-init>
$PDFTOCAIRO = "pdftocairo";
$PDFLATEX = "pdflatex";
$LATEX = "latex";
$MAX_SPLIT_DEPTH = 0;
$HTML_VERSION = "4.0,math";
$NO_SIMPLE_MATH = 1;
$ADDRESS = "<p><address><a href=\"mailto:rmaa\@astro.unam.mx\">rmaa\@astro.unam.mx</address></p>"; 
$INFO = "<p>This document was created using latex2html and the rmaa latex macros</p>"; 
$NO_NAVIGATION = 1;
1; 
%</latex2html-init>
%    \end{macrocode}
% This is the |adshtml.perl| file that is used by \LaTeX{}2HTML. Some
% bizarre \TeX{}y command sometimes show up so we tell it to ignore
% these. All the astronomy macros that come from AAS are processed in
% \TeX{}. That way they get made to images but with an |ALT| tag that
% has the original code, which is what will be picked up by |lynx|,
% with the result that they get passed intact into the final |.txt|
% file. 
%    \begin{macrocode}
%<*adshtml-perl>
package main;

$dum = &do_cmd_makeatletter();

sub do_cmd_RMAA {
    local($_) = @_;
    "RMxAA" . $_;
}

sub do_cmd_RMAASC {
    local($_) = @_;
    "RMxAASC" . $_;
}

sub do_cmd_etal {
    local($_) = @_;
    "et al\." . $_;
}

sub do_cmd_ibidrule {
    local($_) = @_;
    "<img src=\"foo\" alt=\"\ibidrule\">" . $_;
}


sub do_cmd_bibitem {
    local($_) = @_;
    local($key);
    local($opt, $dummy)=&get_next_optional_argument; 
    $key = &missing_braces unless
        ((s/$next_pair_pr_rx/$key = $2; ''/eo)
        ||(s/$next_pair_rx/$key = $2; ''/eo));
    "<p>" . $_; 
}

&ignore_commands (<<_RAW_ARG_CMDS_);
lowercase
adjustfinalcols
textsuperscript # {}
newline # []
thinmuskip
mskip
smallskip
medskip
bigskip
hspace # {}
hspacestar # {}
vspace # {}
vspacestar # {}
mathrm
_at_mathrm
spacefactor
discretionary
unhbox 
voidb_at_x 
penalty 
_at_M
_at_m
null
IeC
_RAW_ARG_CMDS_

&process_commands_in_tex (<<_RAW_ARG_CMDS_);
sun
earth
lesssim
gtrsim
sq
argdeg
arcmin
arcsec
fd
fh
fm
fs
fdg
farcm
farcs
fp
micron
la
ga
case # {} # {}
slantfrac # {} # {}
onehalf
onethird
twothirds
onequarter
threequarters
ubvr
ub
bv
vr
ur
aj
araa
apj
apjl
apjs
ao
apss
aap
aapr
aaps
azh
baas
jrasc
memras
mnras
pra
prb
prc
prd
pre
prl
pasp
pasj
qjras
skytel
solphys
sovast
ssr
zap
nat
iaucirc
aplett
apspr
bain
fcp
gca
grl
jcp
jgr
jqsrt
memsai
nphysa
physrep
physscr
planss
procspie
_RAW_ARG_CMDS_

1;                              # This must be the last line
%</adshtml-perl> 
%    \end{macrocode}

% This is a shell script |makeads.sh| that generates the |_ads.txt| file from
% |_ads.tex| by using \LaTeX{}2HTML and |lynx|.  \changes{4.0}{2019/11/13}{Make sure that lynx writes utf-8 output instead of latin-1.}
% \changes{2.2beta13}{2002/03/12}{Complete rewrite of |makeads.sh|. Now much 
% improved I think}
% \changes{4.0}{2019/11/13}{Fix the at-sign substitution so that it only eliminates \cs{@} and does not mess with email addresses.}

%    \begin{macrocode}
%<*ads-shell-script>
#! /bin/sh
if [ -z $1 ]; then
    echo "Usage: $0 basename"
    exit
fi
BASENAME=$1
ADSBASENAME=${BASENAME}_ads
AUXFILE=$BASENAME.aux
TMPFILE=tmp.$$


echo "Generating $ADSBASENAME.tex"

# For some reason, must double-up \\ to \\\\ but not \ to \\ - odd!
cat > $ADSBASENAME.tex <<EOF
%% PLEASE DO NOT EDIT THIS FILE!  It is generated automatically from
%% the script makeads.sh.  Any hand-edited changes will be lost
%% 
\documentclass{article}
\usepackage[latin1]{inputenc}
\usepackage{html}
\usepackage{adshtml}
\newcommand{\ADSfield}[2]{\%#1  #2\\\\}
\newenvironment{ADSarticle}{}{\par\bigskip}
\setlength{\parindent}{0pt}
\newcounter{istage}
\renewcommand{\ion}[2]{\setcounter{istage}{#2}#1 \Roman{istage}}
\begin{document}
\makeatletter
EOF
echo >> $ADSBASENAME.tex "\input{$BASENAME.ads2}" 
echo >> $ADSBASENAME.tex "\end{document}"



echo "Generating $ADSBASENAME.sed" 

# Find type of journal
grep '\\documentclass\[.*proceedings.*\]' $BASENAME.tex >/dev/null \
     && JOURNAL='RMxAC' \
     || JOURNAL='RMxAA'

# Find the year and volume number
YEAR=`grep '\\SetYear' $BASENAME.tex | sed 's/.*\\SetYear{\(.*\)}.*/\1/'`
VOLUME=`grep '\\SetVolume' $BASENAME.tex | sed 's/.*\\SetVolume{\(.*\)}.*/\1/'`

# Generate an auxilliary sed script 
cat > $TMPFILE <<EOF
s|\\\\newlabel{\(.*\):RMAAFirstPage}{{.*}{\(.*\)}}|s/\\\\\\\\ThisVolume{\1}/$YEAR, RMxAC, $VOLUME, \2/g|g
EOF

# Now go through the .aux file, looking for \newlabels and use the sed 
# script just generated in order to make another sed script, one that 
# will replace \ThisVolume{LABEL} with the correct reference
grep '^\\newlabel{.*:RMAAFirstPage}' $AUXFILE \
     | sed -f $TMPFILE > $ADSBASENAME.sed && rm $TMPFILE

# Apend a further sed command to get rid of \@, since redefining 
# it in _ads.tex didn't work
echo >> $ADSBASENAME.sed 's/\\\\@//g'

echo "Generating $BASENAME.ads2 from $ADSBASENAME.sed and $BASENAME.ads" 

# Apply this to the .ads file, generating the .ads2 file, which is included
# by the _ads.tex file.
sed -f $ADSBASENAME.sed $BASENAME.ads > $BASENAME.ads2

echo "Generating $ADSBASENAME.txt from $ADSBASENAME.tex and $BASENAME.ads2"

# Now run latex2html on the _ads.tex file
latex2html -no_navigation -info 0 -address 0 -html_version 4.0,math $ADSBASENAME

# Then, convert to utf8 with lynx and tidy it up a bit
lynx --justify=off --width=100 --dump --nolist --display_charset=utf-8 \
     $ADSBASENAME/$ADSBASENAME.html | cut -c4- \
     | sed 's/ibidrule/\\ibidrule/g' > $ADSBASENAME.txt
%</ads-shell-script>
%    \end{macrocode}

% \section{Epilogue}
% Err, that's it \dots

\endinput


% \endinput
% Local Variables:
% mode: doctex
% TeX-master: t
% End:
"

clean: 
        rm -f rmaa.aux rmaa.dvi rmaa.ind rmaa.gls
%</makefile>
%    \end{macrocode}

% The |rmaa.ins| file uses the \textsf{docstrip} package to split out
% all the files and strip off this documentation.
%    \begin{macrocode}
%<*ins>
%%
%% Installation batch file for the RMAA LaTeX package
%%
\def\batchfile{rmaa.ins}
\input docstrip.tex
\preamble

PLEASE DO NOT EDIT THIS FILE! It is generated automatically from
the file rmaa.dtx using the accompanying Makefile and rmaa.ins 
Any hand-edited changes will be lost

A LaTeX document class for the RMxAA

Author: William Henney (w.henney@irya.unam.mx)

Latest version available from: http://www.irya.unam.mx/rmaa

Copyright (C) 1999-2019 Instituto de Astronomia, UNAM, Mexico

Permission is hereby granted, free of charge, to any person obtaining
a copy of this software and associated documentation files (the
"Software"), to deal in the Software without restriction, including
without limitation the rights to use, copy, modify, merge, publish,
distribute, sublicense, and/or sell copies of the Software, and to
permit persons to whom the Software is furnished to do so, subject to
the following conditions:

The above copyright notice and this permission notice shall be
included in all copies or substantial portions of the Software.

THE SOFTWARE IS PROVIDED "AS IS", WITHOUT WARRANTY OF ANY KIND,
EXPRESS OR IMPLIED, INCLUDING BUT NOT LIMITED TO THE WARRANTIES OF
MERCHANTABILITY, FITNESS FOR A PARTICULAR PURPOSE AND
NONINFRINGEMENT. IN NO EVENT SHALL THE AUTHORS OR COPYRIGHT HOLDERS BE
LIABLE FOR ANY CLAIM, DAMAGES OR OTHER LIABILITY, WHETHER IN AN ACTION
OF CONTRACT, TORT OR OTHERWISE, ARISING FROM, OUT OF OR IN CONNECTION
WITH THE SOFTWARE OR THE USE OR OTHER DEALINGS IN THE SOFTWARE.
\endpreamble
%    \end{macrocode}
% Clobber files without asking
%    \begin{macrocode}
\askforoverwritefalse
%    \end{macrocode}
% Make the \LaTeX{} macro files (class and styles)
%    \begin{macrocode}
\generate{% 
  \file{rmaa.cls}{\from{rmaa.dtx}{package}}
  \file{rmaa.ins}{\from{rmaa.dtx}{ins}}
  \file{rmaacite.sty}{\from{rmaa.dtx}{rmaacite}}
  \file{adshtml.sty}{\from{rmaa.dtx}{adshtml-style}}
  \file{ADS.tex}{\from{rmaa.dtx}{ads-container}}
  }
%    \end{macrocode}
% Make the other non-\LaTeX{} files. These are done separately since
% they can't use `\%\%' as a comment character for the
% preamble/postamble (and Perl files can't have the postamble I don't
% think). 
%    \begin{macrocode}
\generate{\nopreamble\nopostamble% 
  \file{Makefile}{\from{rmaa.dtx}{makefile}}
  \file{adshtml.perl}{\from{rmaa.dtx}{adshtml-perl}}
  \file{makeads.sh}{\from{rmaa.dtx}{ads-shell-script}}
  \file{sample.latex2html-init}{\from{rmaa.dtx}{latex2html-init}}
  }

%% Local Variables: 
%% mode: LaTeX
%% mode: Font-lock
%% End:
%</ins>
%    \end{macrocode}

% \subsection{Extended citation syntax}
% Modified from a file called mnrascite.sty
% I no longer remember where I stole the original version of this from
% 
% Alterations to the cite commands to make them do more sensible things.
% N.B. the bibliography style should produce entries which have labels where
% the date is delimited by angle brackets '<>' for these macros to work
% properly. They will not give very sensible error messages if you do not do 
% this properly.
% 
% A new command |\scite{key}| is defined which produces a citation which is 
% suitable for insertion in a sentence. i.e., it does not put brackets around 
% the whole citation, only the date. Also |\pcite{key}| produce output with no
% parentheses whatsoever. 
% 
% Should redefine |\@bibitem| really to have any hope of compatibility
% the command |\lbibcite| is defined to deal with the citations.
%    \begin{macrocode}
%<*rmaacite>
\def\bibitem[#1<#2>]#3{\par\hangindent=1.5em\hangafter=1
  \if@filesw
  {\def\protect##1{\string ##1\space}
    \immediate\write\@auxout{\string\lbibcite{#3}{#1}{#2}}}
  \fi}
\@redefinebibitem % what it says
%    \end{macrocode}
% Note that the |\lbibcite| defines two types of citation one contains the names,
% the other the year
%    \begin{macrocode}
\def\lbibcite#1#2#3{\global\@namedef{b@#1}{#2}
                    \global\@namedef{by@#1}{#3}}

\def\@bcite#1#2{({#1\if@tempswa , #2\fi})}
\def\@pcite#1#2{#1\if@tempswa , #2\fi}
%    \end{macrocode}
% This declaration contains what is put between the name and the date
% \changes{3.11}{2003/01/31}{Moved the \cs{articlename} to after the key for consistency
% with the new natbib stuff.}
%    \begin{macrocode}
\def\@brcite{\ (}
\def\@ebrcite{)}
\def\@fmttop#1{\@fmtcite[\csname #1\endcsname]}
\def\@citex[#1]#2{%
%%  \if@filesw\immediate\write\@auxout{\string\citation{#2}}\fi
%% WJH added 16 Jan 2002
  \@for\@citeb:=#2\do{%
    \edef\@citeb{\expandafter\@firstofone\@citeb:\articlename}%
    \if@filesw\immediate\write\@auxout{\string\citation{\@citeb}}\fi}%
%% to here
  \def\@citea{}%
  \@cite{\@for\@citeb:=#2\do
    {%
      \edef\@citeb{\expandafter\@firstofone\@citeb:\articlename}%
      \@citea\def\@citea{;\ }\@ifundefined
      {b@\@citeb}{{\bf ?}\@warning {Citation `\@citeb' on page \thepage \space undefined}}%
      {\csname b@\@citeb\endcsname\@citebro\csname by@\@citeb\endcsname\@citebrc}}}%
  {#1}}

\def\cite{\@ifnextchar [{\let\@citebro=\space
                          \let\@citebrc=\relax
                          \let\@cite=\@bcite\@tempswatrue\@citex}
                        {\let\@citebro=\space
                          \let\@citebrc=\relax
                          \let\@cite=\@bcite\@tempswafalse\@citex[]}}
\def\pcite{\@ifnextchar [{\let\@citebro=\space
                          \let\@citebrc=\relax
                          \let\@cite=\@pcite\@tempswatrue\@citex}
                        {\let\@citebro=\space
                          \let\@citebrc=\relax
                          \let\@cite=\@pcite\@tempswafalse\@citex[]}}
\def\scite{\@ifnextchar [{\let\@citebro=\@brcite
                          \let\@citebrc=\@ebrcite
                          \let\@cite=\@pcite\@tempswatrue\@citex}
                        {\let\@citebro=\@brcite
                          \let\@citebrc=\@ebrcite
                          \let\@cite=\@pcite\@tempswafalse\@citex[]}}

\renewenvironment{thebibliography}
  {\section*{REFERENCES}\vspace*{\afterbiblioskip}\small
    \if@book\RereadBiblio\fi
    \sloppy\clubpenalty4000\widowpenalty4000
    \sfcode`\.=1000\relax
    \setlength\parindent{\z@}\setlength\parskip{\z@}}
  {\relax\normalsize\@outputaddressesifnecessary} 
%</rmaacite>
%    \end{macrocode}

% \subsection{Auxilliary files for generating ADS data}

% The ADS data are written by \LaTeX{} to a |.ads| file. This file
% then has to be included in the following file (|ADS.tex|),
% which is run through \LaTeX{}2HTML (see below). 

%    \begin{macrocode}
%<*ads-container>
%%
%% Instructions:
%% Edit the \input{} line below to reflect the name of your file. 
\documentclass{article}
\usepackage[latin1]{inputenc}
\usepackage{adshtml}
\newcommand{\ADSfield}[2]
{\%#1  #2\\}
%%\newcommand{\ADSvfield}[2]
%%{\%#1 \begin{verbatim}#2\end{verbatim}}
\newenvironment{ADSarticle}{}{\par\bigskip}
\renewcommand{\bibitem}[2][]{\par}
\setlength{\parindent}{0pt}
\renewcommand{\ThisVolume}[1]{%  
  2001, \RMAASC, 12, \oldpageref{#1:RMAAFirstPage}}  
 

\begin{document}
\makeatletter
%% Demonstration file for formatting an entire volume
\documentclass[book,journal,twocolumn]{rmaa}
%%%%%%%%%%%%%%%%%%%%%%%%%%%%%%%%%%%%%%%%%%%%%%%%%%%%%%%%%%%%%%%%
%%%%%%%%%%%%%%%%            PREAMBLE            %%%%%%%%%%%%%%%%
%%%%%%%%%%%%%%%%%%%%%%%%%%%%%%%%%%%%%%%%%%%%%%%%%%%%%%%%%%%%%%%%

\usepackage{paralist}
\usepackage{psfrag,color}
% Needed for landscape tables
\usepackage{pdflscape}

% Used in rm-minimal-test
\usepackage{kantlipsum}

% publication details of volume
\SetVolume{9}
\SetYear{2000}
\SetMonthNumeric{06}
\SetMes{junio}
% details of meeting
\SetConfTitle{Astrophysical Plasmas: Codes, Models, \& Observations}
\SetReunion{M\'exico, D. F.}{octubre 25--29, 1999}
\SetEditors{S. Jane Arthur, Nancy S. Brickhouse, \& Jos\'e Franco}

% this writes the .idx file that is used to generate the index 
\makeindex

%% This command makes the session headings come out bold in the
%% Table-of-Contents. Comment it out if you don't like it.
\BoldSessionTitles

%% This command makes a title page for each session
%% Comment it out if you don't like it.
\SessionTitlePages

%% Use this command if you want the abstract-only papers to be
%% indexed by ADS
\abstractsADS

%% At the special request of Alberto Carramin~ana, the following is a
%% way to suppress the author adresses at the end of each article
%\suppressfulladdresses

%% Use this command to add Received/accepted dates for each paper
% \printdates


\begin{document}

%%%%%%%%%%%%%%%%%%%%%%%%%%%%%%%%%%%%%%%%%%%%%%%%%%%%%%%%%%%%%%%%b
%%%%%%%%%%%%%%%%          FRONT MATTER          %%%%%%%%%%%%%%%%
%%%%%%%%%%%%%%%%%%%%%%%%%%%%%%%%%%%%%%%%%%%%%%%%%%%%%%%%%%%%%%%%
\frontmatter % start roman numbering

% format the TOC
\renewcommand{\RMAAbannerimage}{rmsc_blank}
\tableofcontents

%% preface, group photo, list of participants go here
%% use the `frontsection' environment (see editors' guide)
\begin{frontsection}{Preface}{%
    \listofauthors{S.~J.~Arthur, N.~S.~Brickhouse, J.~Franco}
    \indexauthor{Arthur, S. J.}
    \indexauthor{Brickhouse, N.}
    \indexauthor{Franco, J.}
    }
  These are the proceedings of a very nice conference that Jane and
  Pepe organised. This is the preface, which explains what a wonderful
  time was had by all and how grateful they are for all the people who
  gave them money or helped out.
  
  This is another paragraph of the preface, in case they have anything
  else to say.  For example, the Scientific Organizing Committee
  consisted of
  \begin{quote}
    J. Franco (Mexico, Chair), D. Cox (USA) and various other notables
  \end{quote}
  while the Local Organizing Committee comprised
  \begin{quote}
    S. J. Arthur (Chair), S. Kurtz, etc.
  \end{quote}

  \EditorSignature

\end{frontsection}

\begin{frontsection}[Photo]{Group Photograph}{\notitle}
\null\vfill
\centerline{Here goes the photograph.}
\vfill
\end{frontsection}

\begin{frontsection}[Participants]{List of Participants}{}
%% Here you put a formatted list of participants
\vfil
\dotfill First page of list of participants \dotfill
\vfil\clearpage

\null\vfil
\dotfill Second page of list of participants \dotfill
\vfil

\end{frontsection}

\begin{frontsection}[XXX]{Another front section}{}

\end{frontsection}

%%%%%%%%%%%%%%%%%%%%%%%%%%%%%%%%%%%%%%%%%%%%%%%%%%%%%%%%%%%%%%%%
%%%%%%%%%%%%%%%%          MAIN MATTER           %%%%%%%%%%%%%%%%
%%%%%%%%%%%%%%%%%%%%%%%%%%%%%%%%%%%%%%%%%%%%%%%%%%%%%%%%%%%%%%%%
\mainmatter % start arabic numbering

\begin{articles}
  %% Main articles go here
  %% Use \article{ARTICLENAME} for each individual article
  %% If you want any text to appear in brackets after the article
  %% title in the TOC and index, then put it as an optional first
  %% argument, e.g.,  \article[Invited Talk]{ARTICLENAME} 

  \session{The Contributed Papers}
  \article{rm-journal-example} % An illustrative article
  \article{rm-onepage} % another article, which has a plate
  \article{rm-minimal-test}
  \session{Conference Summary}
  \article[Invited Review]{rmsummsamp} % An illustrative conference summary

\end{articles}

%%%%%%%%%%%%%%%%%%%%%%%%%%%%%%%%%%%%%%%%%%%%%%%%%%%%%%%%%%%%%%%%
%%%%%%%%%%%%%%%%          BACK MATTER           %%%%%%%%%%%%%%%%
%%%%%%%%%%%%%%%%%%%%%%%%%%%%%%%%%%%%%%%%%%%%%%%%%%%%%%%%%%%%%%%%
\backmatter

%% Abstracts of presentations for which the
%% authors have not submitted a paper
\begin{abstracts}

  \title{An abstract of a poster}
  \author{W. J. Henney\altaffilmark{1}}
  \listofauthors{W.~J.~Henney}
  \indexauthor{Henney, W. J}
  \altaffiltext{1}{Where I am from done with altaffilmark/text}
  \abstract{This is a test abstract.   This is a test abstract.   This
    is a test abstract.   This is a test abstract.}
  

  \title{An abstract of another  poster}
  \author{W.~J.~Henney\altaffil{This is done with altaffil} and
    A.~N.~Other\altaffil{So is this} }
  \listofauthors{W.~J.~Henney \& A.~N.~Other}
  \indexauthor{Henney, W. J}
  \indexauthor{Other, A. N.}
  \abstract{This is a test abstract.   This is a test abstract.   This
    is a test abstract.   This is a test abstract.}

  \title{An abstract of yet another  poster}
  \author{W.~J.~Henney\altaffil{This is done with altaffil} and
    A.~N.~Other\altaffil{So is this} }
  \listofauthors{W.~J.~Henney \& A.~N.~Other}
  \indexauthor{Henney, W. J}
  \indexauthor{Other, A. N.}
  \abstract{This is a test abstract.   This is a test abstract.   This
    is a test abstract.   This is a test abstract.}

  \title{The abstract of a poster with a very long title to see how
    well it gets split up into lines}
  \author{A.~N.~Other\altaffil{The affiliation of A.~N.~Other is also
      very long so that it will go over several lines. His email
      address is anon@nowhere.org}}
  \listofauthors{A.~N.~Other}
  \indexauthor{Other, A. N.}
  \abstract{We will try and make this abstract longer so there is more
    stuff to put on the page and it doesn't look too empty. However, I
    have nothing whatsoever to say so will have to stop here. Let's
    repeat that a few times. We will try and make this abstract longer
    so there is more stuff to put on the page and it doesn't look too
    empty. However, I have nothing whatsoever to say so will have to
    stop here. Let's repeat that a few times. We will try and make
    this abstract longer so there is more stuff to put on the page and
    it doesn't look too empty. However, I have nothing whatsoever to
    say so will have to stop here. Let's repeat that a few times. We
    will try and make this abstract longer so there is more stuff to
    put on the page and it doesn't look too empty. However, I have
    nothing whatsoever to say so will have to stop here. Let's repeat
    that a few times.}
    
  \title[Comment to go in TOC]{The abstract of a second poster with a
    very long title to see how well it gets split up into lines}
  \author{A.~N.~Other\altaffil{The affiliation of A.~N.~Other is also
      very long so that it will go over several lines. His email
      address is anon@nowhere.org}}
  \listofauthors{A.~N.~Other}
  \indexauthor{Other, A. N.}
  \abstract{We will try and make this abstract longer so there is more
    stuff to put on the page and it doesn't look too empty. However, I
    have nothing whatsoever to say so will have to stop here. Let's
    repeat that a few times. We will try and make this abstract longer
    so there is more stuff to put on the page and it doesn't look too
    empty. However, I have nothing whatsoever to say so will have to
    stop here. Let's repeat that a few times. We will try and make
    this abstract longer so there is more stuff to put on the page and
    it doesn't look too empty. However, I have nothing whatsoever to
    say so will have to stop here. Let's repeat that a few times. We
    will try and make this abstract longer so there is more stuff to
    put on the page and it doesn't look too empty. However, I have
    nothing whatsoever to say so will have to stop here. Let's repeat
    that a few times.}
    
  \title{The abstract of a third poster with a very long title to see how
    well it gets split up into lines}
  \author{A.~N.~Other\altaffil{The affiliation of A.~N.~Other is also
      very long so that it will go over several lines. His email
      address is anon@nowhere.org}}
  \listofauthors{A.~N.~Other}
  \indexauthor{Other, A. N.}
  \abstract{We will try and make this abstract longer so there is more
    stuff to put on the page and it doesn't look too empty. However, I
    have nothing whatsoever to say so will have to stop here. Let's
    repeat that a few times. We will try and make this abstract longer
    so there is more stuff to put on the page and it doesn't look too
    empty. However, I have nothing whatsoever to say so will have to
    stop here. Let's repeat that a few times. We will try and make
    this abstract longer so there is more stuff to put on the page and
    it doesn't look too empty. However, I have nothing whatsoever to
    say so will have to stop here. Let's repeat that a few times. We
    will try and make this abstract longer so there is more stuff to
    put on the page and it doesn't look too empty. However, I have
    nothing whatsoever to say so will have to stop here. Let's repeat
    that a few times.}
    
  \title{The abstract of a fourth poster with a very long title to see how
    well it gets split up into lines}
  \author{A.~N.~Other\altaffil{The affiliation of A.~N.~Other is also
      very long so that it will go over several lines. His email
      address is anon@nowhere.org}}
  \listofauthors{A.~N.~Other}
  \indexauthor{Other, A. N.}
  \abstract{We will try and make this abstract longer so there is more
    stuff to put on the page and it doesn't look too empty. However, I
    have nothing whatsoever to say so will have to stop here. Let's
    repeat that a few times. We will try and make this abstract longer
    so there is more stuff to put on the page and it doesn't look too
    empty. However, I have nothing whatsoever to say so will have to
    stop here. Let's repeat that a few times. We will try and make
    this abstract longer so there is more stuff to put on the page and
    it doesn't look too empty. However, I have nothing whatsoever to
    say so will have to stop here. Let's repeat that a few times. We
    will try and make this abstract longer so there is more stuff to
    put on the page and it doesn't look too empty. However, I have
    nothing whatsoever to say so will have to stop here. Let's repeat
    that a few times.}
    
  \title{The abstract of a fifth poster with a very long title to see how
    well it gets split up into lines}
  \author{A.~N.~Other\altaffil{The affiliation of A.~N.~Other is also
      very long so that it will go over several lines. His email
      address is anon@nowhere.org}}
  \listofauthors{A.~N.~Other}
  \indexauthor{Other, A. N.}
  \abstract{We will try and make this abstract longer so there is more
    stuff to put on the page and it doesn't look too empty. However, I
    have nothing whatsoever to say so will have to stop here. Let's
    repeat that a few times. We will try and make this abstract longer
    so there is more stuff to put on the page and it doesn't look too
    empty. However, I have nothing whatsoever to say so will have to
    stop here. Let's repeat that a few times. We will try and make
    this abstract longer so there is more stuff to put on the page and
    it doesn't look too empty. However, I have nothing whatsoever to
    say so will have to stop here. Let's repeat that a few times. We
    will try and make this abstract longer so there is more stuff to
    put on the page and it doesn't look too empty. However, I have
    nothing whatsoever to say so will have to stop here. Let's repeat
    that a few times.}
    
  \title{The abstract of a sixth poster with a very long title to see how
    well it gets split up into lines}
  \author{A.~N.~Other\altaffil{The affiliation of A.~N.~Other is also
      very long so that it will go over several lines. His email
      address is anon@nowhere.org}}
  \listofauthors{A.~N.~Other}
  \indexauthor{Other, A. N.}
  \abstract{We will try and make this abstract longer so there is more
    stuff to put on the page and it doesn't look too empty. However, I
    have nothing whatsoever to say so will have to stop here. Let's
    repeat that a few times. We will try and make this abstract longer
    so there is more stuff to put on the page and it doesn't look too
    empty. However, I have nothing whatsoever to say so will have to
    stop here. Let's repeat that a few times. We will try and make
    this abstract longer so there is more stuff to put on the page and
    it doesn't look too empty. However, I have nothing whatsoever to
    say so will have to stop here. Let's repeat that a few times. We
    will try and make this abstract longer so there is more stuff to
    put on the page and it doesn't look too empty. However, I have
    nothing whatsoever to say so will have to stop here. Let's repeat
    that a few times.}
    
  \title{The abstract of a seventh poster with a very long title to see how
    well it gets split up into lines}
  \author{A.~N.~Other\altaffil{The affiliation of A.~N.~Other is also
      very long so that it will go over several lines. His email
      address is anon@nowhere.org}}
  \listofauthors{A.~N.~Other}
  \indexauthor{Other, A. N.}
  \abstract{We will try and make this abstract longer so there is more
    stuff to put on the page and it doesn't look too empty. However, I
    have nothing whatsoever to say so will have to stop here. Let's
    repeat that a few times. We will try and make this abstract longer
    so there is more stuff to put on the page and it doesn't look too
    empty. However, I have nothing whatsoever to say so will have to
    stop here. Let's repeat that a few times. We will try and make
    this abstract longer so there is more stuff to put on the page and
    it doesn't look too empty. However, I have nothing whatsoever to
    say so will have to stop here. Let's repeat that a few times. We
    will try and make this abstract longer so there is more stuff to
    put on the page and it doesn't look too empty. However, I have
    nothing whatsoever to say so will have to stop here. Let's repeat
    that a few times.}
  
  \title{An abstract to go in the list of abstracts, but not in the
    index of authors}
  \author{I.~A.~M. Hidden\altaffil{Puerto Escondido}}
  \listofauthors{I.~A.~M. Hidden}
  \abstract{The idea is that by omitting the
    \textbackslash{}indexauthor command, we can avoid an entry in the
    author index.}

\end{abstracts}


%% The Index: after latexing twice, run `makeindex -s rmaa' on this
%% file, then latex twice more
\printindex

%% if you rename this file to something other than `rmbooksamp', then
%% you must change the above line accordingly

%%
%% Plates: If you have any, then the following line will print
%% them out:

\outputplates

\end{document}

%%% Local Variables:
%%% mode: latex
%%% TeX-master: t
%%% End:

\end{document}
%</ads-container>
%    \end{macrocode}

% This is the |adshtml.sty| style file used by |ADS.tex|. It just
% defines a few 
% macros that might be used in abstracts, etc (really it should define
% them all). This now gets all its code from
% Sec~\ref{sec:user-commands}. 
%    \begin{macrocode}
%<*adshtml-style>
%% Nothing else needs to go here now. 
%</adshtml-style>
%    \end{macrocode}
% The init file for latex2html.
%    \begin{macrocode}
%<*latex2html-init>
$PDFTOCAIRO = "pdftocairo";
$PDFLATEX = "pdflatex";
$LATEX = "latex";
$MAX_SPLIT_DEPTH = 0;
$HTML_VERSION = "4.0,math";
$NO_SIMPLE_MATH = 1;
$ADDRESS = "<p><address><a href=\"mailto:rmaa\@astro.unam.mx\">rmaa\@astro.unam.mx</address></p>"; 
$INFO = "<p>This document was created using latex2html and the rmaa latex macros</p>"; 
$NO_NAVIGATION = 1;
1; 
%</latex2html-init>
%    \end{macrocode}
% This is the |adshtml.perl| file that is used by \LaTeX{}2HTML. Some
% bizarre \TeX{}y command sometimes show up so we tell it to ignore
% these. All the astronomy macros that come from AAS are processed in
% \TeX{}. That way they get made to images but with an |ALT| tag that
% has the original code, which is what will be picked up by |lynx|,
% with the result that they get passed intact into the final |.txt|
% file. 
%    \begin{macrocode}
%<*adshtml-perl>
package main;

$dum = &do_cmd_makeatletter();

sub do_cmd_RMAA {
    local($_) = @_;
    "RMxAA" . $_;
}

sub do_cmd_RMAASC {
    local($_) = @_;
    "RMxAASC" . $_;
}

sub do_cmd_etal {
    local($_) = @_;
    "et al\." . $_;
}

sub do_cmd_ibidrule {
    local($_) = @_;
    "<img src=\"foo\" alt=\"\ibidrule\">" . $_;
}


sub do_cmd_bibitem {
    local($_) = @_;
    local($key);
    local($opt, $dummy)=&get_next_optional_argument; 
    $key = &missing_braces unless
        ((s/$next_pair_pr_rx/$key = $2; ''/eo)
        ||(s/$next_pair_rx/$key = $2; ''/eo));
    "<p>" . $_; 
}

&ignore_commands (<<_RAW_ARG_CMDS_);
lowercase
adjustfinalcols
textsuperscript # {}
newline # []
thinmuskip
mskip
smallskip
medskip
bigskip
hspace # {}
hspacestar # {}
vspace # {}
vspacestar # {}
mathrm
_at_mathrm
spacefactor
discretionary
unhbox 
voidb_at_x 
penalty 
_at_M
_at_m
null
IeC
_RAW_ARG_CMDS_

&process_commands_in_tex (<<_RAW_ARG_CMDS_);
sun
earth
lesssim
gtrsim
sq
argdeg
arcmin
arcsec
fd
fh
fm
fs
fdg
farcm
farcs
fp
micron
la
ga
case # {} # {}
slantfrac # {} # {}
onehalf
onethird
twothirds
onequarter
threequarters
ubvr
ub
bv
vr
ur
aj
araa
apj
apjl
apjs
ao
apss
aap
aapr
aaps
azh
baas
jrasc
memras
mnras
pra
prb
prc
prd
pre
prl
pasp
pasj
qjras
skytel
solphys
sovast
ssr
zap
nat
iaucirc
aplett
apspr
bain
fcp
gca
grl
jcp
jgr
jqsrt
memsai
nphysa
physrep
physscr
planss
procspie
_RAW_ARG_CMDS_

1;                              # This must be the last line
%</adshtml-perl> 
%    \end{macrocode}

% This is a shell script |makeads.sh| that generates the |_ads.txt| file from
% |_ads.tex| by using \LaTeX{}2HTML and |lynx|.  \changes{4.0}{2019/11/13}{Make sure that lynx writes utf-8 output instead of latin-1.}
% \changes{2.2beta13}{2002/03/12}{Complete rewrite of |makeads.sh|. Now much 
% improved I think}
% \changes{4.0}{2019/11/13}{Fix the at-sign substitution so that it only eliminates \cs{@} and does not mess with email addresses.}

%    \begin{macrocode}
%<*ads-shell-script>
#! /bin/sh
if [ -z $1 ]; then
    echo "Usage: $0 basename"
    exit
fi
BASENAME=$1
ADSBASENAME=${BASENAME}_ads
AUXFILE=$BASENAME.aux
TMPFILE=tmp.$$


echo "Generating $ADSBASENAME.tex"

# For some reason, must double-up \\ to \\\\ but not \ to \\ - odd!
cat > $ADSBASENAME.tex <<EOF
%% PLEASE DO NOT EDIT THIS FILE!  It is generated automatically from
%% the script makeads.sh.  Any hand-edited changes will be lost
%% 
\documentclass{article}
\usepackage[latin1]{inputenc}
\usepackage{html}
\usepackage{adshtml}
\newcommand{\ADSfield}[2]{\%#1  #2\\\\}
\newenvironment{ADSarticle}{}{\par\bigskip}
\setlength{\parindent}{0pt}
\newcounter{istage}
\renewcommand{\ion}[2]{\setcounter{istage}{#2}#1 \Roman{istage}}
\begin{document}
\makeatletter
EOF
echo >> $ADSBASENAME.tex "\input{$BASENAME.ads2}" 
echo >> $ADSBASENAME.tex "\end{document}"



echo "Generating $ADSBASENAME.sed" 

# Find type of journal
grep '\\documentclass\[.*proceedings.*\]' $BASENAME.tex >/dev/null \
     && JOURNAL='RMxAC' \
     || JOURNAL='RMxAA'

# Find the year and volume number
YEAR=`grep '\\SetYear' $BASENAME.tex | sed 's/.*\\SetYear{\(.*\)}.*/\1/'`
VOLUME=`grep '\\SetVolume' $BASENAME.tex | sed 's/.*\\SetVolume{\(.*\)}.*/\1/'`

# Generate an auxilliary sed script 
cat > $TMPFILE <<EOF
s|\\\\newlabel{\(.*\):RMAAFirstPage}{{.*}{\(.*\)}}|s/\\\\\\\\ThisVolume{\1}/$YEAR, RMxAC, $VOLUME, \2/g|g
EOF

# Now go through the .aux file, looking for \newlabels and use the sed 
# script just generated in order to make another sed script, one that 
# will replace \ThisVolume{LABEL} with the correct reference
grep '^\\newlabel{.*:RMAAFirstPage}' $AUXFILE \
     | sed -f $TMPFILE > $ADSBASENAME.sed && rm $TMPFILE

# Apend a further sed command to get rid of \@, since redefining 
# it in _ads.tex didn't work
echo >> $ADSBASENAME.sed 's/\\\\@//g'

echo "Generating $BASENAME.ads2 from $ADSBASENAME.sed and $BASENAME.ads" 

# Apply this to the .ads file, generating the .ads2 file, which is included
# by the _ads.tex file.
sed -f $ADSBASENAME.sed $BASENAME.ads > $BASENAME.ads2

echo "Generating $ADSBASENAME.txt from $ADSBASENAME.tex and $BASENAME.ads2"

# Now run latex2html on the _ads.tex file
latex2html -no_navigation -info 0 -address 0 -html_version 4.0,math $ADSBASENAME

# Then, convert to utf8 with lynx and tidy it up a bit
lynx --justify=off --width=100 --dump --nolist --display_charset=utf-8 \
     $ADSBASENAME/$ADSBASENAME.html | cut -c4- \
     | sed 's/ibidrule/\\ibidrule/g' > $ADSBASENAME.txt
%</ads-shell-script>
%    \end{macrocode}

% \section{Epilogue}
% Err, that's it \dots

\endinput


% \endinput
% Local Variables:
% mode: doctex
% TeX-master: t
% End:
"

rmaa.ind: rmaa.idx
        @echo "REGENERATING INDEX"
        makeindex -s gind.ist rmaa

rmaa.gls: rmaa.glo
        @echo "REGENERATING CHANGELOG"
        makeindex -s gglo.ist -o rmaa.gls rmaa.glo

rmaa.pdf: rmaa.aux rmaa.ind rmaa.gls
        @echo "REGENERATING .pdf FILE"
        pdflatex "\nonstopmode% \def\fileversion{4.0} \def\filedate{2019/11/12}
% \MakeShortVerb{\|}
% \title{A \LaTeX{} document class for the Revista Mexicana de
% Astronom\'\i a y Astrof\'\i sica\thanks{This file has version
% \fileversion{} last revised \filedate{}.}}
% \author{William J. Henney\thanks{w.henney@irya.unam.mx}}
% \date{\filedate}
% \maketitle
% \paragraph{Abstract}
% This document describes the \texttt{rmaa} \LaTeX{}
% document class for the journal ``Revista Mexicana de
% Astronom\'\i a y Astrof\'\i sica''. The macros can be used for 
% \begin{enumerate}
% \item Preparing author manuscripts for the journal (including
% conference proceedings to be published in the ``Serie de
% Conferencias'')
% \item Typesetting an entire volume of the journal by inputting the
% individual article files from a container file. Facilities are
% included for automatic generation of TOC and Author Index. 
% \end{enumerate}
% 
% \tableofcontents
% \section{The documentation driver file}
% This is the part that generates the documentation that you may or
% may not be reading, which is done with the help of the Makefile
% discussed in section~\ref{sec:driver_make}. 
%    \begin{macrocode}
%<*driver>
\documentclass{ltxdoc}
\usepackage{hyperref}
\usepackage{palatino}
\EnableCrossrefs \CodelineIndex \RecordChanges
\setlength{\parindent}{0pt}
\def\MacroFont{\footnotesize\tt}
\begin{document} \DocInput{rmaa.dtx} \PrintIndex \PrintChanges \end{document}
%</driver>
%    \end{macrocode}

% \section{The implementation}
% 
% \changes{3.24}{2007/09/14}{Added user access to the macro version
% and date via \cs{rmVersion} and \cs{rmDate}}
% 
%    \begin{macrocode}
%<*package>
\newcommand{\rmaa@version}{3.26}
\newcommand{\rmaa@date}{2011/08/15}
\newcommand\rmVersion{\rmaa@version}
\newcommand\rmDate{\rmaa@date}
\NeedsTeXFormat{LaTeX2e}[1995/06/01]
%    \end{macrocode}
% \changes{4.0}{2019/11/02}{Change abstract to appear before
% resumen. Change RevMexAA to RMxAA for consistency with ADS.}
%
% \changes{3.22}{2007/03/22}{Allow proceedings articles to have
%   received/accepted dates} \changes{3.21}{2006/05/15}{Changes for
%   Silvia - abstracts are now shipped out to ADS.}
% \changes{3.19}{2006/01/19}{Changes for Silvia - remove adresses from
%   end, new abstract style} \changes{2.2beta}{2001/10/30}{ Fixed bug
%   whereby twocolumn option was not passed through to book
%   class. Finally sorted out the ``label on last page'' issue. }
% \changes{2.1}{2001/01/23}{ Fixed some inconsistencies in
%   \cs{nocite} and \cs{bibitem} uniqification Fixed the definition of
%   \cs{ion} so it will work in math mode and in contexts where the
%   typesize is different from normal (e.g., title, fig captions).
%   Modified the definition of \cs{@caption}, so we don't write any
%   list-of-figures data (we never use it anyway, and it just causes
%   problems with fragile commands in fig captions). }
% \changes{2.0}{2000/11/30}{ Added `newstyle' option, which should be
%   used in all new proceedings volumes (wider abstract, smaller
%   typeface in fig captions, etc.) Added commands for formatting
%   table captions. New demonstation document rmpostsamp.tex for
%   1-page posters. This will be the last release under a PRI
%   government (for a while...)}  \changes{1.9}{2000/06/06}{
%   Corrected position of page number in footer of articles that open
%   on even pages. Ensured that first article in each session opens on
%   an odd page, as do list of abstracts and author index. Added extra
%   space in TOC before the backmatter. Added command for editors,
%   \cs{BoldSessionTitles}, that makes produces bold titles for
%   session in the TOC.}  \changes{1.8}{2000/05/16}{ Fixed bug in
%   last page label with \cs{suppressfulladdresses} }
% \changes{1.7}{2000/05/11}{ Fixed more little things. Page labels
%   now work, even if you have no bibliography. You can output the
%   full addresses by hand with the \cs{outputfulladdresses}
%   command. Editors' Guide updated but still not finished. Revamped
%   splitbook.sh now produces PS file of each article for web page and
%   also an EPS file of each page for sending to the printers.}
% \changes{1.6}{2000/04/15}{ More cleaning up: full stops after
%   titles in index, running headers and footers in frontmatter,
%   fontsize in title of index, etc} \changes{1.5}{2000/04/15}{
%   Eliminated superscripts from titles in TOC and index}
% \changes{1.4}{2000/04/14}{ Added copyright and license info (GPL)}
% \changes{1.3}{2000/03/01}{ Corrected \cs{figcaption} to do the
%   caption formatting right Corrected font size of sec headings in
%   front matter} \changes{1.2}{2000/02/16}{ Added \cs{plate}
%   command} \changes{1.1}{2000/02/14}{ Added manuscript option,
%   corrected bugs} \changes{1.0}{2000/02/08}{ Now does a whole book
%   - lots of stuff tidied up} \changes{0.7}{1999/10/04}{ Increased
%   spacing between sections in single column styles. Suppressed
%   floats at bottom of pages (didn't work though) }
% \changes{0.6}{1999/09/24}{ Tweaked \cs{adjustfinalcols} (still not
%   perfect) Added proceedings option Added onecolumn/twocolumn
%   options (orthogonal to journal/proceedings) }
% \changes{0.5}{1999/09/23}{ Added \cs{fulladdresses} and
%   \cs{adjustfinalcols} (just about working)}
% \changes{0.4}{1999/09/22}{ Minor stylistic changes}
% \changes{0.3}{1999/06/28}{ Better mechanism for alternative
%   affiliations Added `guide' option for user guide}
% \changes{0.2}{1999/05/27}{ All basic functionality present, I
%   think} \changes{0.1}{1999/05/26}{ Just about working - still
%   missing Fig/Tab/Refs} \changes{0.0}{1999/05/21}{ So bored today,
%   decided to rewrite RMAA LaTeX macros from scratch}
%    \begin{macrocode}
\ProvidesClass{rmaa}[\rmaa@date\space v\rmaa@version\space RMxAA article
class]
%    \end{macrocode}
% \subsection{Option Handling}
% \label{sec:options}
% Declarations that are needed in options
% \changes{3.05}{2003/10/22}{Added natbib option}
% \changes{3.06}{2003/01/03}{Natbib option turned on permanently}
% \changes{3.07}{2003/01/11}{Added proof option}
% \changes{3.14}{2003/01/11}{Option to turn off natbib}
% \changes{3.18}{2003/10/22}{trying out doctex mode} 
% \changes{3.23}{2007/09/09}{Manuscript is now on by default.}
% \changes{3.26}{2011/08/15}{New option for review articles}
%    \begin{macrocode}
\newif\if@reviewarticle
\@reviewarticlefalse
\newif\if@proof
\@prooffalse
\newif\if@journal
\@journaltrue
\newif\if@printdates
\@printdatesfalse
\newif\if@proceedings
\@proceedingsfalse
\newif\if@manuscript
\@manuscripttrue
\newif\if@bigfont
\@bigfonttrue
\newif\if@guide
\@guidefalse
\newif\if@preprint
\@preprintfalse
\newif\if@rmtwocol
\@rmtwocoltrue
\newif\if@biblist
\@biblistfalse
\newif\if@book
\@bookfalse
\newif\if@natbib
\@natbibtrue
\newif\if@rmcite
\@rmcitefalse
\newif\if@titlevisible
\@titlevisibletrue
\newif\if@abstractpresent
\@abstractpresentfalse
\newif\if@printbanner
\@printbannerfalse
\newif\if@suppressfulladdresses
\@suppressfulladdressesfalse
\newif\if@newstyle
\@newstylefalse
\newif\if@debug
\@debugfalse
\newif\if@boldkeywords
\@boldkeywordsfalse
%    \end{macrocode}
% The option handling---mainly setting various booleans.
% \changes{2.2beta11}{2001/11/26}{Made twocol and
% newstyle the default for proceedings.}
% \changes{3.23}{2007/09/09}{Manuscript is now 10 pt instead of 12pt}
% \changes{3.23}{2007/09/10}{Full addresses are now automatically
% suppressed for proceedings style.}
% \changes{3.24}{2007/09/14}{New option \texttt{boldkeywords} that
% does the keywords in the old ugly way. Now off by default.}
%    \begin{macrocode}
\DeclareOption{manuscript}{\typeout{Manuscript for referee} 
  \@manuscripttrue\@journalfalse\@rmtwocolfalse
  \PassOptionsToClass{10pt}{article}}
\DeclareOption{bigfont}{\typeout{Big font manuscript requested by Gary
  Ferland} 
  \@manuscripttrue\@journalfalse\@rmtwocolfalse\@bigfonttrue
  \PassOptionsToClass{11pt}{article}}
\DeclareOption{journal}{\typeout{Main journal option}%
  \@journaltrue\@rmtwocoltrue\@printdatestrue\@manuscriptfalse
  \PassOptionsToClass{10pt}{article}}
\DeclareOption{guide}{\typeout{User guide option}%
  \@guidetrue\@journalfalse\@rmtwocoltrue\@manuscriptfalse
  \PassOptionsToClass{10pt}{article}}
\DeclareOption{preprint}{\typeout{Preprint option}%
  \@preprinttrue\@manuscriptfalse\PassOptionsToClass{10pt}{article}}
\DeclareOption{proof}{\typeout{Author proof option}%
  \@prooftrue\@manuscriptfalse}
\DeclareOption{proceedings}{\typeout{Proceedings option}%
  \@proceedingstrue\@journalfalse\@rmtwocoltrue\@newstyletrue
  \@suppressfulladdressestrue
  \@printdatesfalse\@manuscriptfalse\PassOptionsToClass{10pt}{article}} 
\DeclareOption{printdates}{\typeout{Force printing of received/accepted dates}%
  \@printdatestrue} 
\DeclareOption{book}{\typeout{Whole issue option}%
  \@booktrue} 
\DeclareOption{nonatbib}{\typeout{Turn off natbib support (required
    for compatibility with old rmaacite package).}%
  \@natbibfalse} 
\DeclareOption{newstyle}{\typeout{Use new layout style}%
  \@newstyletrue}
\DeclareOption{oldstyle}{\typeout{Use old layout style}%
  \@newstylefalse}
\DeclareOption{debug}{\typeout{Loads of annoying messages}%
  \@debugtrue}
%    \end{macrocode}
% These two need to be declared explicitly so that |\ExecuteOptions|
% actually does something.
%    \begin{macrocode}
\DeclareOption{twocolumn}{\@rmtwocoltrue}
\DeclareOption{onecolumn}{\@rmtwocolfalse}
\DeclareOption{twoside}{\PassOptionsToClass{\CurrentOption}{article}}
\DeclareOption{boldkeywords}{\@boldkeywordstrue}
%    \end{macrocode}
% Bold, upper-case keywords used to be the default. Now they must be
% turned on explicitly with the \texttt{boldkeywords} option. 
%    \begin{macrocode}
\DeclareOption{boldkeywords}{\@boldkeywordstrue}
%    \end{macrocode}
% Pass any other option along to the \texttt{article} class.
%    \begin{macrocode}
\DeclareOption*{\PassOptionsToClass{\CurrentOption}{article}}
%    \end{macrocode}
% Default options are to be two-sided and based on the standard
% \LaTeX{} \texttt{article} class,
% unless we are processing the entire book, in which case we base
% ourselves on the standard \texttt{book} class. Default is to use
% manuscript option unless otherwise specified. 
%    \begin{macrocode}
\ExecuteOptions{twoside,manuscript} \ProcessOptions\relax
\if@book
\if@rmtwocol\LoadClass[twocolumn]{book}\else\LoadClass{book}\fi
\newcommand{\abstract}{\relax}
\newcommand{\@maketitle}{\relax}
\else
\if@rmtwocol\LoadClass[twocolumn]{article}\else\LoadClass{article}\fi
\fi
%    \end{macrocode}
% Attempt to use natbib for citations 02 Jan 2003. 
% \changes{3.05}{2003/01/02}{If the natbib option is used, then set
% things up so that natbib's version of the bibliography is identical
% to the standard one.}
% \changes{3.09}{2003/01/23}{Explicitly require a recent version of
% natbib, so that we get a decent error message if the package is too
% old.} 
% \changes{3.23}{2007/09/09}{Now allow Bibtex to work properly} 
% \changes{3.25}{2009/10/01}{Change \cs{newcommand} tp \cs{def} when setting \cs{bibfont}. This allows it to work with all versions of natbib.} 
% \changes{3.25}{2010/02/06}{Updated copyright year.} 
% \changes{3.25}{2010/02/06}{Updated email and web addresses to \texttt{crya.unam.mx}.} 
%    \begin{macrocode}
\if@natbib
\RequirePackage{natbib}[1999/02/25]
\bibpunct{(}{)}{;}{a}{}{,}
\def\bibfont{\small}
\renewcommand\bibsection{\section*{REFERENCES}}
\setlength\bibsep{0pt}
\setlength\bibhang{1.5em}
\bibliographystyle{rmaa}
\fi
%    \end{macrocode}
% We want the \texttt{graphicx} package for the extended syntax of
% |\includegraphics|. We want the \texttt{ifthen} package for using in
% these macros (although we don't much yet). We explicitly load
% \texttt{babel} since some \LaTeX{} distributions messed up the
% hyphenation if we didn't.  
%    \begin{macrocode}
\RequirePackage{graphicx}
\RequirePackage{ifthen}
\IfFileExists{babel.sty}{%
%   \let\@restoreindent=\@afterindentfalse
  \RequirePackage[english]{babel}
%   \@restoreindent
  %% reinstate the definition of \@roman so as not to break .idx files
  \@ifundefined{org@@roman}{\relax}{\renewcommand{\@roman}[1]{\org@@roman{##1}}}
  }{%
  \typeout{Package babel.sty not found - don't worry about this}}
%    \end{macrocode}
% Try to load the \texttt{dcolumn} and \texttt{booktabs} packages if
% they can be found. These can produce nicer-looking tables. 
%    \begin{macrocode}
\IfFileExists{dcolumn.sty}{%
  \RequirePackage{dcolumn}
  }{%
  \typeout{Package dcolumn.sty not found - you will not be able to
    align columns on decimal points in tables}}
\IfFileExists{booktabs.sty}{%
  \RequirePackage{booktabs}
  }{%
  \typeout{Package booktabs.sty not found - your tables will not look
    as pretty as they ought too.}
  }
\IfFileExists{makeidx.sty}{%
  \RequirePackage{makeidx}
}{%
  \typeout{Package makeidx.sty not found - Author Index may not work.}
  }
%% commented out for moment since requires more modern version of LaTeX
%% \RequirePackage[latin1]{inputenc}
%    \end{macrocode}
% \changes{4.0}{2019/11/13}{Make \texttt{hyperref} and
% \texttt{etoolbox} into hard requirements, since we now use them
% extensively.}  We use hyperref to get internal and external links.
% But I don't like the colored boxes so I turn those off.  We use
% etoolbox for patching and robustifying at the moment.  In the
% future, I want to port a lot of other stuff over to it.  For
% instance, replace the ifthen package and the TeX ifs too.
%
% We target the latest version (1.7) of PDF, rather than the default
% of 1.5.  This reduces the chance of getting warnings from included
% graphics.  This is not strictly related to hyperref, but we put it
% here anyway.
%    \begin{macrocode}
\RequirePackage{etoolbox}
\RequirePackage{hyperref}
\hypersetup{hidelinks}
\pdfminorversion=7
%    \end{macrocode}
% We patch the \cs{@combinedblfloats} following a suggestion by Heiko
% Oberdiek at \url{https://tex.stackexchange.com/a/249743/3843}, which
% avoids problems when hyperlinks are split across
% pages. \changes{4.0}{2019/11/12}{Fix for issue where breaking
% hyperlink across pages gave an error.}
%    \begin{macrocode}
\urlstyle{same}
\newcommand\email{\begingroup \urlstyle{rm}\Url}
\patchcmd\@combinedblfloats{\box\@outputbox}{\unvbox\@outputbox}{}{%
  \errmessage{\noexpand\@combinedblfloats could not be patched}%
}%
%    \end{macrocode}
% \subsection{Other setup stuff}\label{sec:setup}
% Some declarations: 
%    \begin{macrocode}
\newif\if@needaddresses
\@needaddressestrue
\newif\if@frontmatter
\@frontmatterfalse
\newif\if@inpreamble
\@inpreambletrue
\newif\if@firstauthor
\@firstauthortrue
\newif\if@firstaffil
\@firstaffiltrue
\newif\if@firstaltaffil
\@firstaltaffiltrue
\newif\if@firstkwd
\@firstkwdtrue
\newif\if@noexpandADS
\@noexpandADSfalse
\newif\if@sessiontitlepages
\@sessiontitlepagesfalse
%    \end{macrocode}
% \begin{macro}{\@rmaaEmpty}
% \begin{macro}{\@rmaaArg}
% These are used in various commands to check if arg was left blank
%    \begin{macrocode}
\newcommand{\@rmaaEmpty}{}
\newcommand{\@rmaaArg}{}
%    \end{macrocode}  
% \end{macro}
% \end{macro}
%
% \subsection{Compatibility with old \LaTeX{} versions}
% \label{sec:compat}
% This package is based on the 1998/12/01 release of \LaTeX{}.
% Older releases may not define all the commands that we need.
% Luckily, we can use |\providecommand| to define these if necessary. 
% The following commands are not defined in the 1994/06/01 release. 
%
% \begin{macro}{\MakeUpperCase}
% \begin{macro}{\textsuperscript}
% \begin{macro}{\textbackslash}
% User commands:
%    \begin{macrocode}
\providecommand{\MakeUpperCase}[1]{\uppercase{#1}}
\providecommand{\textsuperscript}[1]{\ensuremath{^\textrm{#1}}}
\providecommand{\textbackslash}{\ensuremath{\backslash}}
%    \end{macrocode}  
% \end{macro}
% \end{macro}
% \end{macro}
% \begin{macro}{\protected@xdef}
% \begin{macro}{\restore@protect}
% \begin{macro}{\@unexpandable@protect}
% \begin{macro}{\hb@xt@}
% \LaTeX{} internal commands:
%    \begin{macrocode}
\providecommand{\protected@xdef}{%
        \let\@@protect\protect
        \let\protect\@unexpandable@protect
        \afterassignment\restore@protect
        \xdef}
\providecommand{\restore@protect}{\let\protect\@@protect}
\providecommand{\@unexpandable@protect}{\noexpand\protect\noexpand}
\providecommand{\hb@xt@}{\hbox to}
%    \end{macrocode}
% \end{macro}
% \end{macro}
% \end{macro}
% \end{macro}
%
% \subsection{Page layout}\label{sec:page-layout}
%
% The dimensions of the text were defined by Elizabeth. Text width is
% 169\,mm with equal margins on the two sides: 
%    \begin{macrocode}
\setlength{\textwidth}{169mm}
\setlength{\@tempdima}{\paperwidth}
\addtolength{\@tempdima}{-\textwidth}
\setlength{\oddsidemargin}{0.5\@tempdima}
\addtolength{\oddsidemargin}{-1in}
\setlength{\evensidemargin}{\oddsidemargin}
%    \end{macrocode}
% Text height is 225\,mm, top margin is 14\,mm. 
%    \begin{macrocode}
\setlength{\textheight}{225mm} 
\setlength{\topmargin}{-1in} 
\addtolength{\topmargin}{14mm}
%    \end{macrocode}
% Space between columns in two-column styles, header, footer,
% paragraph indentation:
%    \begin{macrocode}
\setlength{\columnsep}{9.5mm}
\setlength{\headsep}{0.7cm}
\setlength{\footskip}{0.75cm}
\setlength{\parindent}{1.5em}
%    \end{macrocode}
% 
% \begin{macro}{\@set@scalable@length}
% Try out a new scheme for dealing with various spacing dimensions and
% allowing the user to easily modify them. The command
% \begin{quote}
% |\set@scalable@length|\marg{name}\marg{defaultlength}\marg{stretch}\marg{shrink}
% \end{quote}
% does a 
% |\setlength|\marg{\cs{name}}\marg{defaultlength plus stretch minus shrink}
% and also defines the user command
% \textbackslash\meta{name}|Scale|\marg{factor} that can 
% be used to rescale length by \meta{factor}. Additionally, now
% defines the package command
% \textbackslash\meta{name}|@set@scale|\marg{FACTOR} 
% that also rescales the length by \meta{FACTOR}, so that the length
% becomes \meta{factor} $\times$ \meta{FACTOR} $\times$
% \meta{defaultlength}.   
%    \begin{macrocode}
\newcommand\@set@scalable@length[4]{%
  %% arg #1: NAME of length without leading `\'
  %%
  %% arg #2: default VALUE  
  \expandafter\newlength\csname @#1@default\endcsname
  \expandafter\setlength\csname @#1@default\endcsname{#2}%
  \expandafter\newlength\csname @@#1@@default\endcsname
  \expandafter\setlength\csname @@#1@@default\endcsname{#2}%
  %% arg #3: default STRETCH  
  \expandafter\newlength\csname @#1@stretch\endcsname
  \expandafter\setlength\csname @#1@stretch\endcsname{#3}%
  \expandafter\newlength\csname @@#1@@stretch\endcsname
  \expandafter\setlength\csname @@#1@@stretch\endcsname{#3}%
  %% arg #4: default SHRINK  
  \expandafter\newlength\csname @#1@shrink\endcsname
  \expandafter\setlength\csname @#1@shrink\endcsname{#4}%
  \expandafter\newlength\csname @@#1@@shrink\endcsname
  \expandafter\setlength\csname @@#1@@shrink\endcsname{#4}%
  %%
  %% User-defined scale \@NAME@scale is unity by default
  \expandafter\gdef\csname @#1@scale\endcsname{1}%
  %% Package-defined scale \@NAME@SCALE is unity by default
  \expandafter\gdef\csname @#1@SCALE\endcsname{1}%
  %%
  %% internal command \@reset@NAME should expand to 
  %% \setlength\@@NAME@@default{\@NAME@SCALE\@NAME@default}
  %% \setlength\NAME{\@NAME@scale\@@NAME@@default}
  %% in other words: scale x SCALE x default 
  \expandafter\gdef\csname @reset@#1\endcsname{%
    \expandafter\setlength\csname @@#1@@default\endcsname{%
      \csname @#1@SCALE\endcsname \csname @#1@default\endcsname
      }
    \expandafter\setlength\csname @@#1@@stretch\endcsname{%
      \csname @#1@SCALE\endcsname \csname @#1@stretch\endcsname
      }
    \expandafter\setlength\csname @@#1@@shrink\endcsname{%
      \csname @#1@SCALE\endcsname \csname @#1@shrink\endcsname
      }
    \expandafter\setlength\csname #1\endcsname{%
      \csname @#1@scale\endcsname \csname @@#1@@default\endcsname
      \@plus
      \csname @#1@scale\endcsname \csname @@#1@@stretch\endcsname
      \@minus
      \csname @#1@scale\endcsname \csname @@#1@@shrink\endcsname
      }}
  %%
  %% User command \NAMEScale should expand to 
  %% {\renewcommand\@NAME@scale{#1}\@reset@NAME}
  \expandafter\gdef\csname #1Scale\endcsname##1{%
    \expandafter\gdef\csname @#1@scale\endcsname{##1}%
    \csname @reset@#1\endcsname}
  %%
  %% Package command \NAME@set@SCALE should expand to 
  %% {\renewcommand\@NAME@SCALE{#1}\@reset@NAME}
  \expandafter\gdef\csname #1@set@SCALE\endcsname##1{%
    \expandafter\gdef\csname @#1@SCALE\endcsname{##1}%
    \csname @reset@#1\endcsname}
  %%
  %% initialize at default values with \@reset@NAME
  \csname @reset@#1\endcsname
  }
%    \end{macrocode}
% \end{macro}
% \begin{macro}{\parskip}
% First use of this new infrastructure: set interparagraph space to be
% zero but with a small amount of stretchiness.
%    \begin{macrocode}
%   \ifthenelse{\boolean{@manuscript}}{%
%     \@set@scalable@length{parskip}{0.5\baselineskip}{0.3\p@}{0.1\p@}
%   }{
    \@set@scalable@length{parskip}{0\p@}{0.3\p@}{0\p@}
%   }
%    \end{macrocode}
% \end{macro}
% 
% \begin{macro}{\abovesectionskip}
% \changes{4.0}{2019/11/12}{Use negative values for skip so that first paragraph after section headings is not indented.  This is something that I have been trying to fix for years!  It turns out that the anser was in the documentation for the standard command \cs{@startsection.}}
% \begin{macro}{\belowsectionskip}
% \begin{macro}{\aboveparaskip}
% \begin{macro}{\belowparaskip}
% Define some utility rubber lengths for use in
% section~\ref{sec:section-commands}. 
%    \begin{macrocode}
\newlength{\abovesectionskip}
\newlength{\belowsectionskip}
\newlength{\aboveparaskip}
\newlength{\belowparaskip}
%    \end{macrocode}
% The default values for these lengths are those traditionally used
% for the onecolumn style. 
%    \begin{macrocode}
\@set@scalable@length{abovesectionskip}{-3.5ex}{-4ex}{-2ex}
\@set@scalable@length{belowsectionskip}{2.3ex}{0ex}{2ex}
\@set@scalable@length{aboveparaskip}{3.25ex}{2ex}{0.5ex}
\@set@scalable@length{belowparaskip}{-1em}{\z@}{\z@}
%    \end{macrocode}
% \end{macro}\end{macro}\end{macro}\end{macro}
% We also make scalable some of the standard float parameters. These
% are the same defaults as in the standard article classes, but
% slightly more stretchy. For the twocolumn style, they are currently
% halved, see below. Note that I have also lumped |\abovecaptionskip|
% in with these, since it plays a similar role. There is also a
% |\belowcaptionskip| but it is set to zero in the standard classes. 
%    \begin{macrocode}
\@set@scalable@length{floatsep}{12\p@}{4\p@}{4\p@}
\@set@scalable@length{textfloatsep}{20\p@}{4\p@}{8\p@}
\@set@scalable@length{intextsep}{12\p@}{4\p@}{4\p@}
\@set@scalable@length{dblfloatsep}{12\p@}{4\p@}{4\p@}
\@set@scalable@length{dbltextfloatsep}{20\p@}{4\p@}{8\p@}
\@set@scalable@length{abovecaptionskip}{10\p@}{\z@}{\z@}
%    \end{macrocode}
% \begin{macro}{\@before@title@skip}         
% \begin{macro}{\@before@author@skip}        
% \begin{macro}{\@between@author@skip}       
% \begin{macro}{\@before@abstract@skip}      
% \begin{macro}{\@after@abstracthead@skip}   
% \begin{macro}{\@between@abstract@skip}     
% \begin{macro}{\@before@kwd@skip}           
% \begin{macro}{\@before@dates@skip}           
% Spacing parameters used in |\@maketitle|. These are not changed by
% default except by the user, so they don't need to go into
% |\@package@default@lengths|.  
%    \begin{macrocode}
\newlength{\@before@title@skip}         
\newlength{\@before@author@skip}        
\newlength{\@between@author@skip}       
\newlength{\@before@abstract@skip}      
\newlength{\@after@abstracthead@skip}   
\newlength{\@between@abstract@skip}     
\newlength{\@before@kwd@skip}           
\newlength{\@before@dates@skip}           
\@set@scalable@length{@before@title@skip}{2em}{\z@}{\z@}
\@set@scalable@length{@before@author@skip}{1.5em}{\z@}{\z@}
\@set@scalable@length{@between@author@skip}{0.5em}{\z@}{\z@}
\@set@scalable@length{@before@abstract@skip}{3ex}{1ex}{0.5ex}
\@set@scalable@length{@after@abstracthead@skip}{1ex}{\z@}{\z@}
\@set@scalable@length{@between@abstract@skip}{2ex}{1ex}{0.5ex}
\@set@scalable@length{@before@kwd@skip}{1ex}{0.5ex}{0.2ex}
\@set@scalable@length{@before@dates@skip}{2ex}{0.5ex}{0.2ex}
%    \end{macrocode}
% \end{macro}\end{macro}\end{macro}\end{macro}
% \end{macro}\end{macro}\end{macro}\end{macro}
% 
% Spacing above and below displayed equations also wants to be done
% like this, but it is more complicated since the standard article
% classes make these change with the fontsize. Something to do later,
% maybe. 
% 
%    \begin{macrocode}
%    \end{macrocode}
% \begin{macro}{\@package@default@lengths}
% The value actually used for these lengths is controlled by the
% \meta{SCALE} parameter, which depends on the style option that is in
% effect. This is set using the |\@package@default@lengths| command,
% which must be called every time the style changes (e.g., changing
% from onecolumn to twocolumn). 
%    \begin{macrocode}
\newcommand{\@package@default@lengths}{%
  \ifthenelse{\boolean{@rmtwocol}}{%
    %% twocolumn mode
    \abovesectionskip@set@SCALE{0.5}
    \belowsectionskip@set@SCALE{0.5}
    \aboveparaskip@set@SCALE{0.5}
    \belowparaskip@set@SCALE{1}
    \floatsep@set@SCALE{0.5}
    \textfloatsep@set@SCALE{0.5}
    \intextsep@set@SCALE{0.5}
    \dblfloatsep@set@SCALE{0.5}
    \dbltextfloatsep@set@SCALE{0.5}
    }{%
    %% onecolumn mode
    \abovesectionskip@set@SCALE{1}
    \belowsectionskip@set@SCALE{1}
    \aboveparaskip@set@SCALE{1}
    \belowparaskip@set@SCALE{1}
    \floatsep@set@SCALE{1}
    \textfloatsep@set@SCALE{1}
    \intextsep@set@SCALE{1}
    \dblfloatsep@set@SCALE{1}
    \dbltextfloatsep@set@SCALE{1}
    }%
  }
\@package@default@lengths
%    \end{macrocode}
% \end{macro}
% \begin{macro}{\RescaleLengths}
% \begin{macro}{\RescaleSecLengths}
% \begin{macro}{\RescaleFloatLengths}
% \begin{macro}{\RescaleTitleLengths}
% User commands that will fiddle with all the scalable lengths. 
%    \begin{macrocode}
\newcommand{\RescaleLengths}[1]{%
  \RescaleSecLengths{#1}\RescaleFloatLengths{#1}\RescaleTitleLengths{#1}}
\newcommand{\RescaleSecLengths}[1]{%
%% sections
  \abovesectionskipScale{#1}
  \belowsectionskipScale{#1}
  \aboveparaskipScale{#1}
  }
\newcommand{\RescaleFloatLengths}[1]{%
%% floats
  \floatsepScale{#1}
  \textfloatsepScale{#1}
  \intextsepScale{#1}
  \dblfloatsepScale{#1}
  \dbltextfloatsepScale{#1}
  \abovecaptionskipScale{#1}
  }
\newcommand{\RescaleTitleLengths}[1]{%
%% maketitle
  \@before@title@skipScale{#1}           
  \@before@author@skipScale{#1}        
  \@between@author@skipScale{#1}       
  \@before@abstract@skipScale{#1}      
  \@after@abstracthead@skipScale{#1}   
  \@between@abstract@skipScale{#1}     
  \@before@kwd@skipScale{#1}           
  \@before@dates@skipScale{#1}           
  }
%    \end{macrocode}
% \end{macro}\end{macro}\end{macro}\end{macro}
% \changes{2.2beta14}{2002/03/17}{Finished implementation of new
% infrastructure for scalable lengths. Converted the above/below
% section skips to use it. Also, all the floatsep parameters. Also,
% all the vertical lengths that are used by |\@maketitle|. }

% \begin{macro}{\MSextramargin}
% For the \texttt{manuscript} option we have an extra 0.3\,inch
% vertical margin and 1\,inch horizontal margin for marginal notes (by
% hand or with |\marginpar|. 
%
% \changes{3.23}{2007/09/09}{Reduce baselinestretch in
%   \texttt{manuscript} style to 1.2. No longer use raggedright for
%   manuscript.}
%
% \changes{3.24}{2007/09/14}{Further reduction of baselinestretch in
%   \texttt{manuscript} style to 1.1 since we now have wider margins.}
%
% \changes{3.24}{2007/09/14}{Split \texttt{manuscript}'s extra margin
%   into separate horizontal and vertical parts.}
% 
%    \begin{macrocode}
\newlength{\MSextraHmargin}\newlength{\MSextraVmargin}
\setlength{\MSextraHmargin}{1in} 
\setlength{\MSextraVmargin}{0.3in} 
\if@manuscript
   \renewcommand{\baselinestretch}{1.1}
   \addtolength{\textwidth}{-2\MSextraHmargin}
   \addtolength{\oddsidemargin}{\MSextraHmargin}
   \setlength{\evensidemargin}{\oddsidemargin}
   \addtolength{\textheight}{-2\MSextraVmargin}
   \addtolength{\topmargin}{\MSextraVmargin}
%    \raggedright
\fi
%    \end{macrocode}
% \end{macro}
% Float parameters: don't allow bottom floats by default. 
%    \begin{macrocode}
\renewcommand{\fps@figure}{tp}
\renewcommand{\fps@table}{tp}
%    \end{macrocode}
%
% \begin{macro}{\rmaa@AuthorWidthFrac}
% \begin{macro}{\rmaa@AbstractWidthFrac}
% These parameters control what fraction of the total text width can
% be used by the list of authors and the abstract/resumen,
% respectively.  
%    \begin{macrocode}
\newcommand{\rmaa@AuthorWidthFrac}{1.0}
\newcommand{\rmaa@AbstractWidthFrac}{0.75}
%    \end{macrocode}
% Use a wide abstract to save space in the new style.
%    \begin{macrocode}
\if@newstyle\renewcommand{\rmaa@AbstractWidthFrac}{1.0}\fi
%    \end{macrocode}
% \end{macro}
% \end{macro}
%
% Some things to consider in future versions. Trying to avoid too many
% widows, etc. Apparently, the following will work best with
% \cs{flushbottom}. 
%    \begin{macrocode}
%%\widowpenalty=10000
%%\clubpenalty=10000
%    \end{macrocode}
% Also, should perhaps try getting rid of underfull/overfull box warnings. 
% \changes{3.10}{2003/01/24}{Set large values for \cs{hbadness},
% \cs{vbadness}, \cs{hfuzz}, \cs{vfuzz}. To reduce amount of stuff
% that gets written to the terminal every time you compile something.}
%    \begin{macrocode}
\hbadness=10000
\vbadness=10000
\setlength\hfuzz{\maxdimen}
\setlength\vfuzz{\maxdimen}
%    \end{macrocode}
% \subsection{Various text styles for different contexts}
% \label{sec:text-styles} 
% \begin{macro}{\rmaa@AffilStyle}
% \begin{macro}{\rmaa@PagesStyle}
% \begin{macro}{\rmaa@YearStyle}
% \begin{macro}{\rmaa@VolumeStyle}
% \begin{macro}{\rmaa@VolumenStyle}
% \begin{macro}{\rmaa@JournalNameStyle}
% \begin{macro}{\rmaa@ReceivedStyle}
% \begin{macro}{\rmaa@AbstractHeaderStyle}
% \begin{macro}{\rmaa@JournalNameStyle}
% Style commands for the different fields used by |\maketitle| in
% constructing the title page of the article. These have the form
% |\rmaa@XYZStyle{TEXT}|. I'm not sure these were really a good idea,
% but never mind \dots 
%    \begin{macrocode}
\newcommand{\rmaa@AffilStyle}[1]{{\normalsize\rmfamily#1}}
\newcommand{\rmaa@PagesStyle}[1]{{\normalsize\rmfamily#1}}
\newcommand{\rmaa@YearStyle}[1]{{\normalsize\rmfamily#1}}
\newcommand{\rmaa@VolumeStyle}[1]{{\normalsize\rmfamily\bfseries#1}}
\newcommand{\rmaa@VolumenStyle}[1]{{\Large\sffamily\bfseries#1}}
\newcommand{\rmaa@JournalNameStyle}[1]{{\normalsize\itshape#1}}
\newcommand{\rmaa@ReceivedStyle}[1]{{\normalsize\rmfamily\itshape#1}}
\newcommand{\rmaa@AbstractHeaderStyle}[1]{%
  {\normalsize\rmfamily\MakeUpperCase{#1}}}
\newcommand{\rmaa@TitleStyle}[1]{%
  {\large\bfseries\rmfamily\boldmath\MakeUpperCase{#1}\par}}
%    \end{macrocode}
% \end{macro}
% \end{macro}
% \end{macro}
% \end{macro}
% \end{macro}
% \end{macro}
% \end{macro}
% \end{macro}
% \end{macro}
%
% \begin{macro}{\rmaa@FrontSecStyle}
% Headings of ``front sections'' (see section~\ref{sec:front-sections})
%    \begin{macrocode}
\newcommand{\rmaa@FrontSecStyle}[1]{%
  {\normalsize\bfseries\rmfamily\MakeUpperCase{#1}}}
%    \end{macrocode}
% \end{macro}
%
% \begin{macro}{\rmaa@TOCTitleStyle}
% \begin{macro}{\rmaa@TOCPartStyle}
% \begin{macro}{\rmaa@IndexTitleStyle}
% \begin{macro}{\rmaaTOCAuthorStyle}
% \begin{macro}{\rmaaIndexAuthorStyle}
% Commands to set style of TOC and index entries.
% Rewritten to avoid unnecessary font commands in the .toc and .ind files
%    \begin{macrocode}
\newcommand{\rmaa@TOCTitleStyle}[1]{\MakeUpperCase{#1}}
\newcommand{\rmaa@TOCPartStyle}[1]{\rmaa@TOCTitleStyle{#1}}
\newcommand{\rmaa@IndexTitleStyle}[1]{#1}
\newcommand{\rmaaTOCAuthorStyle}[1]{%
  {\itshape#1}}
\newcommand{\rmaaIndexAuthorStyle}[1]{%
  \protect\mbox{\bfseries#1}}
%    \end{macrocode}
% \end{macro}
% \end{macro}
% \end{macro}
% \end{macro}
% \end{macro}
%
% \begin{macro}{\rmaa@KeywordStyle}
% \begin{macro}{\rmaa@KeywordHeaderStyle}
% The style of the keywords depends on whether we are using the
% \texttt{newstyle} option or not. If we are, then they are small. 
%    \begin{macrocode}
\ifthenelse{\boolean{@newstyle}}{
  \newcommand{\rmaa@KeywordStyle}[1]{%
    {\small%
      \if@boldkeywords 
      \bfseries\MakeUpperCase{\raggedright #1}
      \else
      \raggedright #1
      \fi
    }}
  \newcommand{\rmaa@KeywordHeaderStyle}[1]{%
    {\small\itshape#1}}}
{
  \newcommand{\rmaa@KeywordStyle}[1]{%
    {\normalsize%
      \if@boldkeywords\bfseries\fi\raggedright #1}}
  \newcommand{\rmaa@KeywordHeaderStyle}[1]{%
    {\normalsize\itshape#1}}}
%    \end{macrocode}
% \end{macro}
% \end{macro}
%
% \subsection{Page styles}\label{sec:page-styles}
% \begin{macro}{\ps@rmaa}
% The running headers for normal pages:
%    \begin{macrocode}
\newcommand{\ps@rmaa}{%
%    \end{macrocode}
% nothing in either foot; 
%    \begin{macrocode}
  \renewcommand\@oddfoot{\relax}%
  \renewcommand\@evenfoot{\relax}%
%    \end{macrocode}
% page number and short author list in the even heads; 
%    \begin{macrocode}
  \renewcommand\@evenhead{\rlap{\thepage}\hfil\@shortauthor\hfil}%
%    \end{macrocode}
% short title and page number in the odd heads.
%    \begin{macrocode}
  \renewcommand\@oddhead{\hfil\@shorttitle\hfil\llap{\thepage}}
  }
%    \end{macrocode}
% \begin{macro}{\ps@rmaafirstpage}
% The header of the first page of each
% article:
%    \begin{macrocode}
\newcommand{\ps@rmaafirstpage}{%
  \renewcommand\@oddfoot{\hfil\thepage}%
%    \end{macrocode}
% footer corrected (06 Jun 2000) to have page on outside edge;
%    \begin{macrocode}
  \renewcommand\@evenfoot{\thepage\hfil}%
%    \end{macrocode}
% Header has journal name, volume, and page numbers for a regular
% article, or similar for special cases. 
% \changes{4.0}{2019/11/03}{Change to three-line header with copyright info and DOI.}
%    \begin{macrocode}
\if@preprint
\renewcommand\@oddhead{%
  \if@proceedings 
  \textsc{To appear in ``\@conftitle{} (\rmaa@YearStyle{\@year})''}\hfil
  \textit{RMxAC}
  \else
  \textsc{To appear in} \@journalname\hfil
  \fi}\fi%
\if@guide
\renewcommand\@oddhead{\textsc{Documentation} (\rmaa@YearStyle{\@year})\hfil}
\fi%
\if@manuscript
\renewcommand\@oddhead{Manuscript for \@journalname{} (\rmaa@YearStyle{\@year})\hfil}
\fi%
\if@journal
\renewcommand\@oddhead{%
  \parbox{\textwidth}{
    \@journalname, \rmaa@VolumeStyle{\@volume},
    \rmaa@PagesStyle{\@artpages}\
    (\rmaa@YearStyle{\@year})\hfil\\[1pt]
    {\small\textcopyright{} \@year:
      Instituto de Astronom\'\i{}a,
      Universidad Nacional Aut\'onoma de M\'exico}\hfil\\[1pt]
    {\small\@DOI@url}\hfil
  }%
}\fi%
  \renewcommand\@evenhead{\@oddhead}}
%    \end{macrocode}
% By default, use the |\ps@rmaa| page style.
%    \begin{macrocode}
\pagestyle{rmaa}
%    \end{macrocode}
% \end{macro}
% \end{macro}
%
% \subsection{Preamble commands}\label{sec:preamble} 
% Everything you need for setting the title page, TOC and index
% entries, etc.
% \subsubsection{Commands for the title page}
% \label{sec:title-commands}
% \begin{macro}{\resumen}
% \begin{macro}{\@resumen}
% \begin{macro}{\abstract}
% \begin{macro}{\@abstract}
% User commands for setting the resumen and abstract. 
% These should be used in the document header. The internal commands
% |\@resumen| and |\@abstract| are used to hold the actual text
% pending output via the |\maketitle| command. 
%    \begin{macrocode}
\newif\if@setresumen
\newcommand{\resumen}[1]{%
%%  \renewcommand{\@resumen}{{\selectlanguage{spanish}#1}}%
  \renewcommand{\@resumen}{{\ignorespaces#1}}%
  \@setresumentrue}
\renewcommand{\abstract}[1]{%
  \renewcommand{\@abstract}{{\ignorespaces#1}}
  \renewcommand\ADS@abstract{#1}
%% Not a good idea, since then author macros are not expanded and
%% cause errors later. However, by not doing it we are vulnerable to
%% fragile commands.....
  \if@noexpandADS\@onelevel@sanitize{\ADS@abstract}\fi
  \@abstractpresenttrue
  }
\newcommand{\ADSnoexpand}{\@noexpandADStrue}
%    \end{macrocode}
% \begin{macro}{\boldabstract}
% This is the version for doing 1-page conference proceedings. Still
% save the text for ADS but then write it out immediately as a bold
% paragraph. 
%    \begin{macrocode}
\newcommand{\boldabstract}[1]{%
  \gdef\ADS@abstract{#1}%
  \renewcommand{\@abstract}{#1}%
  \if@noexpandADS\@onelevel@sanitize{\ADS@abstract}\fi%
  \noindent\textbf{\boldmath #1}
  }
%    \end{macrocode}
% \end{macro}
% Default value in each case is to print out a warning if not set
%    \begin{macrocode}
  \newcommand{\@resumen}{%
    \@setresumenfalse \ClassWarningNoLine{rmaa}{No Spanish abstract
      (resumen) specified - if you do not speak Spanish, this will be
      translated by the editors}%
    \textbf{Favor de proporcionar un resumen en espa\~nol. If you are
      unable to translate your abstract into Spanish, the editors will
      do it for you.} }
  \newcommand{\@abstract}{\ClassWarningNoLine{rmaa}{No abstract
      specified}}
%    \end{macrocode}
% \end{macro}
% \end{macro}
% \end{macro}
% \end{macro}
%
% \begin{macro}{\title}
% \begin{macro}{\@rawtitle}
% \begin{macro}{\@TOCtitle}
% \begin{macro}{\rmaaIndexTitle}
% Similarly, |\title| is the preamble command to set the title of the
% article. This is used in various places, typeset in different
% styles, hence the use of the internal |\@rawtitle|, |\@title|,
% |\rmaaIndexTitle|, and |\@TOCtitle| commands, each of which uses the
% relevant style from section~\ref{sec:text-styles}
%    \begin{macrocode}
\newcommand{\@rawtitle}{Raw Title}
\newcommand{\@TOCtitle}{\rmaa@TOCTitleStyle{Title for TOC}}
\newcommand{\rmaaIndexTitle}{\rmaa@IndexTitleStyle{Title for index}.}
\newcommand{\@TOCaftertitle}{}
\renewcommand{\title}[1]{%
%%  {\ADS@sanitize\xdef\@rawtitle{#1}}
  \renewcommand\@rawtitle{\noexpand#1}
  \let\ADS@title=\@rawtitle
  \renewcommand{\@title}{\rmaa@TitleStyle{#1}}%
  \renewcommand{\@TOCtitle}{\rmaa@TOCTitleStyle{#1}\@TOCaftertitle}%
  \renewcommand{\rmaaIndexTitle}{\rmaa@IndexTitleStyle{#1\@TOCaftertitle}.}%
  }
%    \end{macrocode}
% \end{macro}
% \end{macro}
% \end{macro}
% \end{macro}
% 
% All the weird and whacky ways of indicating the author
% affiliations. 
% \begin{macro}{\affil}
% First, the way that sets it directly underneath the authors name.
% Note that |\affil| trumps all other affiliations for the purposes of
% getting into the ADS stuff (we use |\gdef|).   
% \dots 
%    \begin{macrocode}
\newcommand{\affil}[1]{%
  \\[1.5ex]\rmaa@AffilStyle{#1}\\[1.5ex]%
  \if@firstaffil\gdef\ADS@affil{#1}\@firstaffilfalse\fi%
  }
\renewcommand{\and}{and\\[1.5ex]}
%    \end{macrocode}
% \begin{macro}{\altaffil}
% Alternative affiliations are footnoted. 
%    \begin{macrocode}
\newcommand{\altaffil}[1]{%
  \if@firstaltaffil\gdef\ADS@affil{#1}\@firstaltaffilfalse\fi%
  \thanks{#1}}
%    \end{macrocode}
% \begin{macro}{\altaffilmark}
% \begin{macro}{\altaffiltext}
% Alternative-alternative mechanism for affiliations---should not be
% mixed with |\altaffil| (unless you know what you are doing) or there
% will be problems with numbering.
%    \begin{macrocode}
\newcommand{\altaffilmark}[1]{\textsuperscript{#1}}
\newcommand{\altaffiltext}[2]{%
  \if@firstaltaffil\gdef\ADS@affil{#2}\@firstaltaffilfalse\fi
  \stepcounter{footnote}%
  \protected@xdef\@thanks{\@thanks
    \protect\footnotetext[#1]{#2}}} 
%    \end{macrocode}
% \end{macro}
% \end{macro}
% \end{macro}
% \end{macro}
%
% \begin{macro}{\keywords}
% \begin{macro}{\@keywords}
% \begin{macro}{\@keywordheader}
% \label{keywords}Sundry minor preamble commands. All
% self-explanatory, I hope. |\keywords| is the old user interface
% (deprecated). New articles should use |\addkeyword| once for each
% keyword. That way the ADS record comes out right too (comma rather
% than em-dash separator).
% \changes{2.2beta12}{2001/12/07}{\cs{addkeyword} New user interface
%   for keywords.} 
% \changes{2.2beta13}{2002/03/11}{\cs{addkeyword} Improved behaviour
%   with .ads file. Use
%   \cs{ADS@sanitize} to change definitions of `~' and \cs{ion}.}
% \changes{3.23}{2007/09/09}{\cs{addkeyword} now signals a compilation
% error if the author tries to add more than 6 keywords.}
% \changes{3.24}{2007/09/14}{kwdcounter is now reset in
% \cs{@begofarticlehook} to prevent error in book style} 
%    \begin{macrocode}
\newcommand{\@kwdsep}{ --- }
\newcommand{\ADS@kwdsep}{, }
\newcounter{kwdcounter}\setcounter{kwdcounter}{0}
\newcommand{\addkeyword}[1]{%
  \addtocounter{kwdcounter}{1}
  \ifthenelse{\thekwdcounter > 6}{
    \ClassError{rmaa}{Too many keywords. A maximum of 6 keywords is allowed.}{}
  }{}
  \if@firstkwd%
  \renewcommand\@keywords{\if@boldkeywords\uppercase{#1}\else #1\fi}
  {\ADS@sanitize\def\ion{\plainion}\gdef\ADS@keywords{#1}}
  \@firstkwdfalse
  \else%
  \edef\@keywords{\@keywords\@kwdsep \if@boldkeywords\uppercase{#1}\else #1\fi}%
  {\ADS@sanitize\def\ion{\plainion}\xdef\ADS@keywords{\ADS@keywords\ADS@kwdsep #1}}%
  \fi}
\newcommand{\ADS@sanitize}{%
  \def\mathrm{}
  \def\@{\ }
  \def~{\ }
  \def\langle{<}
  \def\rangle{>}
  \def\log{log }
  \def\sin{sin }
  \def\cos{cos }
  \def\tan{tan }
  \def\fd{^d.}
  \def\fh{^h.}
  \def\fm{^m.}
  \def\fs{^s.}
  \def\fdg{�.}
  \def\farcm{'.}
  \def\farcs{''.}
  \def\arcdeg{�}
  \def\,{ }
  }
\newcommand{\keywords}[1]{%
  \renewcommand{\ADS@keywords}{#1}
  \renewcommand{\@keywords}{\uppercase{#1}}
  }
\newcommand{\@keywordheader}{\rmaa@KeywordHeaderStyle{Key Words:}}
\newcommand{\@keywords}{%
  \ClassWarningNoLine{rmaa}{No keywords specified}%
  Keywords go here
  }
%    \end{macrocode}
% \begin{macro}{\shorttitle}
% \begin{macro}{\shortauthor}
% \begin{macro}{\@shorttitle}
% \begin{macro}{\@shortauthor}
% Commands for setting the running headers.
%    \begin{macrocode}
\newcommand{\shorttitle}[1]{%
  \renewcommand\@shorttitle{\MakeUpperCase{#1}}}
\newcommand{\shortauthor}[1]{%
  \renewcommand\@shortauthor{\MakeUpperCase{#1}}}
\newcommand{\@shorttitle}{SHORT TITLE GOES HERE}
\newcommand{\@shortauthor}{SHORT AUTHOR GOES HERE}
%    \end{macrocode}
% \end{macro}
% \end{macro}
% \end{macro}
% \end{macro}
% \end{macro}
% \end{macro}
% \end{macro}


% \subsubsection{The full addresses}
% \label{sec:full-addresses}
% \begin{macro}{\fulladdresses}
% \begin{macro}{\@fulladdresses}
% \begin{macro}{\@fulladdheight}
% \begin{macro}{\@fulladddepth}
% \begin{macro}{\@fulladdbox}
% Aaghh, the full addresses! A perennial source of problems. 
%    \begin{macrocode}
\newlength{\@fulladdheight}
\newlength{\@fulladddepth}
\newsavebox{\@fulladdbox}
%% \fulladdresses command modified (11 May 2000) so that we can put
%% \fulladdresses{} in the article to suppress the output of the
%% addresses
\newcommand{\@FinalFigure}{\relax}
\newcommand{\FinalFigure}[1]{\renewcommand{\@FinalFigure}{#1}}
\newcommand{\fulladdresses}[1]{%
  \renewcommand{\@rmaaArg}{#1}
  \ifx\@rmaaArg\@rmaaEmpty
     \renewcommand{\@fulladdresses}{}
  \else
     \renewcommand{\@fulladdresses}{
       \normalsize%
       \ifx\@FinalFigure\relax\else\vfill\@FinalFigure\vfill\fi
%%       \vspace*{1ex}
       %% do addresses as a list (copied from thebibliography environment) 
       \list{\null}{\leftmargin 1.5em\topsep\z@\labelwidth\z@\labelsep\z@\itemindent -1.5em
         \itemsep\z@\parsep\z@\usecounter{enumi}}
       \sloppy\clubpenalty4000\widowpenalty4000
       \sfcode`\.=1000\relax
       #1
       \endlist
       }
  \fi
  }
\newcommand{\@fulladdresses}{Full addresses go here}
%    \end{macrocode}
% \begin{macro}{\suppressfulladdresses}
% New user-level command for suppressing full addresses should you so
% desire. Currently used in the single-page-poster style. 
%    \begin{macrocode}
\newcommand{\suppressfulladdresses}{\@suppressfulladdressestrue\@needaddressesfalse}
%    \end{macrocode}
% \end{macro}
% \end{macro}
% \end{macro}
% \end{macro}
% \end{macro}
% \end{macro}

% \begin{macro}{\printdates}
% \begin{macro}{\noprintdates}
% Command to activate the printing out of the received/accepted
% dates. This can also be done using the \texttt{printdates} option,
% which is the default for \texttt{journal}, but not for \texttt{proceedings}.
%    \begin{macrocode}
\newcommand{\printdates}{\@printdatestrue}
\newcommand{\noprintdates}{\@printdatesfalse}
%    \end{macrocode}
% \end{macro}
% \end{macro}
% \begin{macro}{\@outputaddressesifnecessary}
% Command to output the addresses if we haven't already done so
% This is called by: 
% \begin{enumerate}
% \item The end of |thebibliography| environment
% \item The |\outputfulladdresses| command
% \end{enumerate}
%    \begin{macrocode}
\newcommand{\@outputaddressesifnecessary}{%
  \ClassInfo{rmaa}{Calling @outputaddressesifnecessary}
  \ifthenelse{\boolean{@needaddresses}}{%
    %% If in one column style, or in the first column of a two column
    %% style then we can try to put it on this page 
    \ifthenelse{\boolean{@firstcolumn} \or \not\boolean{@rmtwocol}}{%
      \@outputaddresses}{%
      %% otherwise give warning it will be on a separate page
      \ClassWarning{rmaa}{Full addresses will go on a separate
        page (please see User Guide for how to remedy this)}}%
    %% If we haven't managed yet, put it on its own page
    \ifthenelse{\boolean{@needaddresses}}{%
      \if@debug{\ClassWarning{rmaa}{In this silly bit!}}\fi
      \clearpage\vfill\@outputaddresses}{}%
    }{
    %% If we don't need them, we still might need to do the eoa hook 
    \ifthenelse{\boolean{@suppressfulladdresses}}{%
      \if@debug{\ClassWarning{rmaa}{XXX}}\fi
%%      \if@book\else\@endofarticlehook\null\fi
      }{}
    }%
  }  
%    \end{macrocode}
% \end{macro}

% \begin{macro}{\outputfulladdresses}
% A user interface to the above command, in case there is no biblio
% (11 May 2000) 
%    \begin{macrocode}
\newcommand{\outputfulladdresses}{\@outputaddressesifnecessary}
%    \end{macrocode}
% \end{macro}

% \begin{macro}{\adjustfinalcols}
% This command starts a new column and should be 
% used to ensure that the columns balance on the last page. It also
% outputs the full 
% author addresses to ensure that this occurs in the left hand
% column. 
% The |\par| command is so that the current line will be taken account
% of when working out |\pagetotal|
%    \begin{macrocode}
\newcommand{\adjustfinalcols}{\if@twocolumn\goodbreak\@outputaddresses\newpage\noindent\fi}
%    \end{macrocode}
% \end{macro}

% \begin{macro}{\@outputaddresses}
% Command to do the work of outputting the addresses. 
% This can be called by 
% \begin{enumerate}
% \item an explicit |\adjustfinalcols|
% \item via |\@outputaddressesifnecessary|, q.v.
% \end{enumerate}
%
%    \begin{macrocode}
\newcommand{\@outputaddresses}{%
  % Allow option of suppressing the full addresses, in which case
  % we just run the hook if we are not in a book style
  \if@suppressfulladdresses
%%     \if@book\else\@endofarticlehook\null\fi
  \else
  %% save the full addresses in a box for output later
  \sbox{\@fulladdbox}{\parbox{\textwidth}{\@fulladdresses}}
  % calculate height of address box  
  \settoheight{\@fulladdheight}{\usebox{\@fulladdbox}}
  \settodepth{\@fulladddepth}{\usebox{\@fulladdbox}}
  \addtolength{\@fulladdheight}{\@fulladddepth}
  \if@debug
  { \typeout{Address total height: \the\@fulladdheight}
    \typeout{Text height: \the\pagetotal}
    \typeout{Goal height: \the\pagegoal}
    }
  \fi
  % reuse \@fulladdheight to mean the total height of text+addresses
  \addtolength{\@fulladdheight}{\pagetotal}
%%  \addtolength{\@fulladdheight}{-\baselineskip}
  \if@debug\typeout{Adjusted text + addresses height: \the\@fulladdheight}\fi
  % check it will fit on current page
  \ifdim\@fulladdheight<\pagegoal
     %% must shift it to the left
     %% to compensate for the hanging indentation in the reference list. 
     %% NOTE: this is not foolproof!
     %% NB. This will give overfull hbox warnings, which is a bit ugly.
     \vfill\noindent\if@biblist\hspace*{-\leftmargin}\fi
     \usebox{\@fulladdbox}
     % Set this to show we've done it
     \@needaddressesfalse
     % and assume we are at the end of the article
     % so that we do the EOA hook if not in the book style
%%     \if@book\else\@endofarticlehook\fi
  \else
     % if not, just do nothing. With luck, this function will get 
     % called again on the next page
     \if@debug{\typeout{Holding over full addresses to next page}}\fi
  \fi\fi
  }
%    \end{macrocode}
% \end{macro}

% \subsubsection{Commands for editors}
% \label{sec:editor-commands}
% Some user/editor commands for doing the TOC and index in the book
% option. 
% They are defined here rather in the book-specific part since they
% need to be always available in order to avoid breaking documents. 
% \begin{macro}{\listofauthors}
% \begin{macro}{\@rawauthors}
% \begin{macro}{\rmaaTOCauthors}
% The command |\listofauthors| formats an author list for use in the
% TOC. It should contain the same names as in the |\author| command,
% but without the affiliation garbage and using |\&| instead of
% |and|. 
%    \begin{macrocode}
\newcommand{\@rawauthors}{Raw Authors}
\newcommand{\rmaaTOCauthors}{\noexpand\rmaaTOCAuthorStyle{Author list for TOC}}
\newcommand{\listofauthors}[1]{%
  \renewcommand{\@rawauthors}{#1}
  \renewcommand\rmaaTOCauthors{\noexpand\rmaaTOCAuthorStyle{#1}}}
%    \end{macrocode}
% \begin{macro}{\indexauthor}\label{indexauthor}
% The command |\indexauthor| produces an index entry for each
% author. 
% This command must come \emph{after} the
% |\title| and the |\listofauthors| commands. Also, there should be one 
% |\indexauthor| command for \emph{each} author. We also add the
% author onto the end of the |\ADS@authors| macro and construct the
% |\ADS@authorinitial| by a most cunning method (see below), but only
% for the first author. 
% 
% No, No, NO! This was all wrong. Backed out completely: 
% \begin{quote}
% Completely rewritten (19 Dec 2002) since we now use the index class, so 
% as to be able to alphabetize accented characters. A side effect of this is
% that |\index| no longer expands out the argument after the |@|, so we 
% instead add all the |\index| commands to a list (|\indexauthor@deferlist|), 
% which is executed in the |\@begofarticlehook|. 
% 
% Note the stuff with |\toks0| is copied from Exercise 20.15 of the \TeX{}book---it is
% to make sure the latin1-accented name gets passed into the idx
% file. 
% \end{quote}
% We now go back to the old way of doing things, using an optional
% argument to the \cs{indexauthor} command that has the name without
% accents. 
% \changes{3.08}{2003/01/16}{Remerged changes from 3.06 that got
% dropped from 3.07 after a synchronization disaster. Basically, the
% accentless optional argument to \cs{indexauthor}. Must add the macro
% directories to my unison set-up to stop this sort of thing from happening.}
%    \begin{macrocode}
\newcommand\@isortkey{}
\newcommand\@emptykey{empty}
\newcommand\indexauthor[2][empty]{
  \renewcommand\@isortkey{#1}
  \ifx\@isortkey\@emptykey \renewcommand\@isortkey{#2}\fi
  \index{\@isortkey @{{\rmaaIndexAuthorStyle{#2}}\hskip1em plus1em minus0.4em
      \rmaaIndexTitle\hskip1em plus2em minus0.4em\penalty-100
      \rmaaTOCauthors}}
  \indexauthorADS{#2}
}
%%\def\indexauthor@deferlist{\relax}
%%\long\def\indexauthor#1{
%%  \def\@@name{#1} \toks0=\expandafter{\@@name}
%%  \toks1=\expandafter{\indexauthor@deferlist}
%%  \protected@xdef\indexauthor@deferlist{%
%%    \the\toks1
%%    \noexpand\index{\the\toks0@{\noexpand\rmaaIndexAuthorStyle{\the\toks0}\noexpand\hskip1em plus1em minus0.4em 
%%        \rmaaIndexTitle\noexpand\hskip1em plus2em minus0.4em\noexpand\penalty-100
%%        \rmaaTOCauthors}}
%%  }
%%  \indexauthorADS{#1}
%%  }
%%\xdef\indexauthor@deferlist{}
\newcommand\indexauthorADS[1]{%
  \if@firstauthor\protected@xdef\ADS@authorinitial{\firstletter{#1}}\fi
  \protected@xdef\ADS@authors{\ADS@authors\if@firstauthor\else ; \fi
    #1}
  \@firstauthorfalse
}
%    \end{macrocode}
% Note that the first |#1| sets the key for the index, making sure it
% actually starts with  
% first letter of the surname. Meanwhile, the |@{...}| expression,
% which starts with some  
% control sequences is what is written to the |.ind| file. 
% See the makeindex man page for details. 
% \end{macro}
% \end{macro}
% \end{macro}
% \end{macro}

% \begin{macro}{\KillAllAccents}
% \begin{macro}{\firstletter}
% These are utility functions used in the above. |lett| is a counter
% indicating the current letter in the string. |\firstletter| just
% prints the first letter of its argument.
%    \begin{macrocode}
\newcommand{\KillAllAccents}{\def\'{}\def\`{}\def\~{}\def\.{}\def\i{i}}
\newcommand\firstletter[1]{\@car #1\@nil}
%    \end{macrocode}
% \end{macro}
% \end{macro}
% \begin{macro}{\SetVolume}
% \begin{macro}{\SetMSnumber}
% \begin{macro}{\SetFirstPage}
% \begin{macro}{\SetYear}
% \begin{macro}{\ReceivedDate}
% \begin{macro}{\AcceptedDate}
% These are for the editors to set the volume number, etc. 
%    \begin{macrocode}
\newcommand{\SetVolume}[1]{%
  \renewcommand{\@volume}{#1}
  \edef\ADS@volume{\PadToFourWithDots{#1}}
  \renewcommand{\@volumen}{\rmaa@VolumenStyle{volumen #1}}%
  }
\newcommand{\SetMSnumber}[1]{\renewcommand{\@MSnum}{#1}}
%    \end{macrocode}
% Note that there is no real reason to use |\SetFirstPage| since all
% the page numbering is done automatically in the book style. 
%    \begin{macrocode}
\newcommand{\SetFirstPage}[1]{\setcounter{page}{#1}}
\newcommand{\SetYear}[1]{%
  \renewcommand{\@year}{#1}
  \renewcommand{\ADS@year}{#1}
  \renewcommand{\@anyo}{\rmaa@VolumenStyle{#1}}
  }
\newcommand{\ReceivedDate}[1]{\renewcommand{\@received}{#1}}
\newcommand{\AcceptedDate}[1]{\renewcommand{\@accepted}{#1}}
%    \end{macrocode}
%
% \begin{macro}{\DOI}
% User/editor command to enter the DOI, which will be typeset as a link.
% \changes{4.0}{2019/11/02}{New user command \cs{DOI} to typeset
% DOI in the header of first page and add it to ADS info.}
%    \begin{macrocode}
\newcommand\DOI[1]{%
  \def\@DOI@url{\url{https://doi.org/10.22201/ia.01851101p.#1}}%
  \def\ADS@doi{https://doi.org/10.22201/ia.01851101p.#1}%
}%
%    \end{macrocode}
% \end{macro}
%
% \begin{macro}{\ORCIDiD}
% \changes{4.0}{2019/11/12}{Add mechanism for giving ORCID iD links}
% \begin{macro}{\ORCIDiDfull}
% \begin{macro}{\@orcid@icon}
% \begin{macro}{\@orcid@iconbig}
% Commands to add icon links to ORCID iDs of authors. For example
% \begin{quote}
% |\ORCIDiD|\marg{0000-0000-0000-0000}
% \end{quote}
% to be used after each author name in the |\author|
% command and in the |\fulladdresses|.  Alternatively,
% |\OORCIDiDfull| will typeset the icon plus the full text of the
% link, but we don't use that at present.  The icon image is stored in |ORCIDiD_iconvector.pdf|, which needs to be on the input path.  In principle, we could make EPS versions if necessary.
%    \begin{macrocode}
\newcommand\@orcid@icon{\includegraphics[height=1.8ex]{ORCIDiD_iconvector}}
\newcommand\@orcid@iconbig{\includegraphics[height=2ex]{ORCIDiD_iconvector}}
\newcommand\ORCIDiD[1]{\href{https://orcid.org/#1}{\,\@orcid@icon}}
\newcommand\ORCIDiDfull[1]{%
  \href{https://orcid.org/#1}{\@orcid@iconbig}\,\url{https://orcid.org/#1}}
%    \end{macrocode}
% \end{macro}
% \end{macro}
% \end{macro}
% \end{macro}
%
% \begin{macro}{\@volume}
% \begin{macro}{\@volumen}
% \begin{macro}{\@artpages}
% \begin{macro}{\@year}
% \begin{macro}{\@anyo}
% \begin{macro}{\@received}
% \begin{macro}{\@accepted}
% \begin{macro}{\@MSnum}
% Default values for the containers used by above commands, plus a
% few more useful ones. 
%    \begin{macrocode}
\newcommand{\@volume}{00}
\newcommand{\@volumen}{\rmaa@VolumenStyle{00}}
\newcommand{\@artpages}{%
  \oldpageref{\articlename:RMAAFirstPage}--\oldpageref{\articlename:RMAALastPage}}
\newcommand{\@year}{\the\year}
\newcommand{\@anyo}{\rmaa@VolumenStyle{2007}}
\newcommand{\@journalname}{\rmaa@JournalNameStyle{%
    \if@proceedings\RMAASClong\else\RMAAlong\fi}}
\newcommand{\@received}{}
\newcommand{\@accepted}{}
\newcommand{\@MSnum}{0000}
\newcommand{\@DOI@prefix}{https://doi.org/10.22201/ia.01851101p}
\def\@DOI@url{\url{https://doi.org/10.22201/ia.01851101p.2099.00.00.00}}
%    \end{macrocode}
% \end{macro}
% \end{macro}
% \end{macro}
% \end{macro}
% \end{macro}
% \end{macro}
% \end{macro}
% \end{macro}
% \end{macro}
% \end{macro}
% \end{macro}
% \end{macro}
% \end{macro}
% \end{macro}
% 
% The following commands are most useful in the \texttt{book}
% style. However, the |\SetConfTitle| command can now be used in
% individual articles to set the header for the \texttt{preprint}
% style. 
%    \begin{macrocode}
\newcommand{\@reunionplace}{Lugar de la reuni\'on}
\newcommand{\@reuniondate}{Fecha de la reuni\'on}
\newcommand{\@conftitle}{Conference Title}
\newcommand{\SetReunion}[2]{\renewcommand{\@reunionplace}{#1}\renewcommand{\@reuniondate}{#2}}
\newcommand{\SetConfTitle}[1]{\renewcommand{\@conftitle}{#1}}
\newcommand{\@editors}{Juan Ram\'on de la Fuente and ``El Mosh''}
\newcommand{\SetEditors}[1]{\renewcommand{\@editors}{#1}}
\newcommand{\@editores}{Juan Ram\'on de la Fuente y ``El Mosh''}
\newcommand{\SetEditores}[1]{\renewcommand{\@editores}{#1}}
\newcommand{\@mes}{mes de la publicaci\'on}
\newcommand{\SetMes}[1]{\renewcommand{\@mes}{\rmaa@VolumenStyle{#1}}}
%    \end{macrocode}

% \subsection{Notes for Author Proofs and Communication between Editors}
% \label{sec:notes}
% We introduce a command \cs{Note} that should be used to describe
% any change made to the manuscript. These are a bit tricky to get
% right. In general we use \cs{marginpar} but this only works in outer
% paragraph mode, so we just bung the box where we are if we are in a
% figure caption or something. If we are in the preamble, then we just
% use a \marginpar as part of the \maketitle command.
%    \begin{macrocode}
\newif\if@arenotes
\@arenotesfalse
\newcounter{notecounter}
\setcounter{notecounter}{0}
\newcommand\NotesList{}
\newcommand\Note[1]{%
  \if@proof{%
    \@arenotestrue%
    \stepcounter{notecounter}%
    \if@inpreamble\else{%
      \NoteMark{\arabic{notecounter}}%
    }\fi
    \protected@xdef\NotesList{\NotesList 
      \protect\item #1}%
  }\fi%
}
\newcommand\Changed[2]{Changed ``#1'' to ``#2''.}
\newcommand\NoteMark[1]{%
  \ifinner
  \raisebox{1ex}[0pt][0pt]{\makebox[0pt][l]{\framebox{\footnotesize\bfseries
        N#1}}}%
  \else
  \marginpar{\framebox{\footnotesize\bfseries N#1}}%
  \fi
}
\newcommand\RefNote[1]{\if@proof\NoteMark{\ref{#1}}\relax\fi}
\newcommand\NoteLabel{\protect\label}
\newcommand\OutputNotes{%
  \bgroup
  \onecolumn
  \section*{EDITORIAL NOTES}
  \begin{enumerate}
    \NotesList
  \end{enumerate}
  \egroup
}
\DeclareRobustCommand\cs[1]{\texttt{\char`\\#1}}

%    \end{macrocode}
% \subsection{Bibliographic Records for ADS}
% \label{sec:ADS}
% \changes{2.2beta11}{2001/11/27}{Just about finished all the ADS
% stuff now. Except for the references, which I have given up on for
% the time being.}
% \begin{macro}{\ADS@BibCode}
% \begin{macro}{\ADS@year}
% \begin{macro}{\ADS@journal}
% \begin{macro}{\ADS@volume}
% \begin{macro}{\ADS@firstpage}
% \begin{macro}{\ADS@authorinitial}
% We use the |\ADS@BibCode| command to hold the bibliographic code for
% ADS. It is 
% used in the |%R| 
% field of the bibliographic record for each article. 
% The format is 
% \begin{quote}
% \texttt{YYYYJJJJJVVVVMPPPPA}  
% \end{quote}
% where \texttt{YYYY} is the year, \texttt{JJJJJ} is the journal code,
% which is \texttt{RMxAA} or \texttt{RMxAC}, \texttt{VVVV} is the
% volume number (right-justified), \texttt{M} is not used by us,
% \texttt{PPPP} is the first page number (right-justified), and
% \texttt{A} is the initial of the first author's surname. The volume
% and page fields are padded with periods (`\texttt{.}'), which causes
% us some fun (this is done with the macro |\PadToFourWithDots| described
% below) . 
% 
% All these individual bits are held in the commands |\ADS@year|,
% |\ADS@journal|, |\ADS@volume|, |\ADS@firstpage|,
% |\ADS@authorinitial|. 
%    \begin{macrocode}
\newcommand{\SetADSBibCode}[1]{\renewcommand{\ADS@BibCode}{#1}}
\newcommand{\ADS@BibCode}
{\ADS@year\ADS@journal\ADS@volume.\PadToFourWithDots{\ADS@firstpage}\ADS@authorinitial}
\newcommand{\ADS@year}{2001} 
\newcommand{\ADS@journal}{\if@proceedings RMxAC\else RMxAA\fi}
\newcommand{\ADS@volume}{....} 
%    \end{macrocode}
% These two are just set to default values here, |\ADS@firstpage| is set
% in the |document| environment for included articles, in the
% |frontsection| environment, |\tableofcontents| command,  |index| and
% |abstracts| environments. Also by |\AtBeginDocument| for stand-alone
% articles. 
% 
% God knows how we are going to do |\ADS@authorinitial|!
%    \begin{macrocode}
\newcommand{\ADS@firstpage}{....}
\newcommand{\ADS@authorinitial}{A}
%    \end{macrocode}
% \end{macro}
% \end{macro}
% \end{macro}
% \end{macro}
% \end{macro}
% \end{macro}
% We also have a whole slew of other fields that can be set. These are
% all listed in Table~\ref{tab:ADS}. Most of them still have to be
% defined. 
% \begin{table}[!ht]\raggedleft
% \caption{ADS Bibliographic Records}\label{tab:ADS}
% \begin{tabular}{lp{8em}lp{8em}}\hline
% Code & Description & Macro Name & Comment \\ \hline
% |%R| & Bibliographic Code & |\ADS@BibCode| & required \\
% |%A| & Author list & |\ADS@authors| & required, generate from |\indexauthor| \\
% |%F| & Author affiliation & |\ADS@affil| & format: AA(\dots)
% AB(\dots) \dots \\
% |%J| & Journal Name & |\ADS@journalname| &  \\
% |%D| & Publication Date & |\ADS@fulldate| &  \\
% |%L| & Last Page & |\ADS@lastpage| & generate from label \\
% |%T| & Title & |\ADS@title| & required, generate from |\title| \\
% |%C| & Abstract Copyright & |\ADS@copyright| & set to IA-UNAM \\
% |%U| & URL for electronic document & |\ADS@url| & set to PDF file \\
% |%d| & DOI for article & |\ADS@doi| & set by |\DOI| \\
% |%K| & Keywords & |\ADS@keywords| & set equal to |\@keywords| \\
% |%A| & Abstract Text & |\ADS@abstract| & set equal to |\@abstract| \\
% |%R| & References & |\ADS@references| & generate from the
% bibliography  \\ \hline
% \end{tabular}
% \end{table}
% A lot of these macros are set by other commands, as mentioned in the
% Table, but some need to be set  explicitly here. 
% 
% \begin{macro}{\ADS@doi}
% \changes{4.0}{2019/11/08}{Save DOI for ADS}
% The user/editor command \cs{DOI} is used to set this. 
%    \begin{macrocode}
\newcommand{\ADS@doi}{https://doi.org/10.22201/ia.01851101p.0000.00.00.00}
%    \end{macrocode}
% \end{macro}
% 
% \begin{macro}{\ADS@journalname}
% This
% does not contain any accents at the moment so as not to cause
% problems---this maybe should be changed though. 
%    \begin{macrocode}
\newcommand{\ADS@journalname}{%
  \if@proceedings\@conftitle\space (Eds. \@editors) \fi%
  Revista Mexicana de Astronom\'\i{}a y Astrof\'\i{}sica%
  \if@proceedings\space (Serie de Conferencias)\fi%
  \space Vol. \@volume, pp. \ADS@firstpage--\ADS@lastpage\space (\@year)%
  \space (http://www.astroscu.unam.mx/\~{}rmaa/)
  }
%    \end{macrocode}
% \end{macro}
% \begin{macro}{\ADS@keywords}
% \begin{macro}{\ADS@abstract}
% \begin{macro}{\ADS@title}
% |\ADS@abstract| is straightforward, since we already have the text
% saved as |\@abstract|. We could do the same with |\ADS@keywords| but
% to avoid unwanted formatting we explicitly set it in |\keywords|
% (p.~\pageref{keywords}). |\ADS@title| is just set to |\@rawtitle|,
% which we were already using for another purpose. 
%    \begin{macrocode}
\let\ADS@title=\@rawtitle
\newcommand{\ADS@keywords}{No keywords specified}
\newcommand{\ADS@abstract}{No abstract specified}
%    \end{macrocode}
% \end{macro}\end{macro}\end{macro}
% \begin{macro}{\ADS@authors}
% |\ADS@authors| wants to be in ``Other, A. N.'' format, but with
% semi-colons between authors. The most natural way to do this seems to
% be with |\indexauthor| (p.~\pageref{indexauthor}). We append each
% author to this macro, which means we must explicitly clear it for
% each new article (done in |\article|). .   
%    \begin{macrocode}
\newcommand{\ADS@authors}{}
%    \end{macrocode}
% \end{macro}
% \begin{macro}{\ADS@fulldate}
% \begin{macro}{\ADS@lastpage}
% \begin{macro}{\ADS@copyright}
% \begin{macro}{\ADS@website}
% \begin{macro}{\ADS@url}
% \begin{macro}{\SetWebsite}
% \begin{macro}{\SetMonthNumeric}
% Other sundry data, together with user commands for bits that can be
% set by the Editors. The |\ADS@url| is constructed from the website
% (currently |http://www.astrosmo.unam.mx/rmaa/|,
% the journal descriptor (e.g., |RMxAC..10/|), followed by the
% path to the PDF file (e.g., |PDF/RMxAC..10_HEW1.pdf|). 
% 
%    \begin{macrocode}
\newcommand{\ADS@website}{http://www.astrosmo.unam.mx/rmaa/}
\newcommand{\ADS@url}
{\ADS@website\ADS@journal\ADS@volume\ADS@number/PDF/\ADS@journal\ADS@volume\ADS@number\_\articlename.pdf}
\newcommand{\ADS@number}{}
\newcommand{\ADS@fulldate}{00/\ADS@year} 
\newcommand{\ADS@lastpage}{00} 
\newcommand{\ADS@copyright}{(C) \ADS@year; Instituto de Astronom\'\i{}a,
  Universidad Nacional Aut\'onoma de M\'exico} 
\newcommand{\SetNumber}[1]{\renewcommand{\ADS@number}{-#1}}
\newcommand{\SetWebsite}[1]{\renewcommand{\ADS@website}{#1}}
\newcommand{\SetMonthNumeric}[1]{\renewcommand{\ADS@fulldate}{#1/\ADS@year}}
%    \end{macrocode}
% \end{macro}\end{macro}\end{macro}\end{macro}\end{macro}\end{macro}\end{macro}
% \begin{macro}{\ADS@affil}
% \begin{macro}{\ADS@references}
% These two are somewhat more problematical. Affiliation currently
% only works for the first author (see \dots) pending redesign of the
% user interface. References are now done by wrapping the
% |thebibliography| environment to open another input stream on the
% same file and search through for the |\bibitems|. 
%    \begin{macrocode}
\newcommand{\ADS@references}{}
\newcommand{\ADS@affil}{}
%    \end{macrocode}
% \end{macro}\end{macro}
% 
% Define two macros that expand to the strings
% |\begin{thebibliography}| and |\end{thebibliography}|. Note that the
% idea for this and much else is inspired by the code for the docstrip
% package. Also define |\StartBibX| in case authors use the more
% standard |\begin{thebibliography}{}| syntax. 
%    \begin{macrocode}
\edef\EndBib{\expandafter\@gobble\string\\end\string{thebibliography\string}}
\edef\StartBib{\expandafter\@gobble\string\\begin\string{thebibliography\string}}
\edef\StartBibX{\expandafter\@gobble\string\\begin\string{thebibliography\string}\string{\string}}
%    \end{macrocode}
% Define two new logical variables. |\ifAnotherline| is true if we
% want to read another line from the input stream. It is set to false
% when we encounter the |\end{thebibliography}|. |\ifInBib| is set to
% true when we frind the |\begin{thebibliography}| and tells us to
% start saving the following lines. 
%    \begin{macrocode}
\newif\ifAnotherLine
\newif\ifInBib
%    \end{macrocode}
% \begin{macro}{makeOther}
% We need a macro to change the catcodes of lots of things to other. 
%    \begin{macrocode}
\def\makeOther#1{\catcode`#1=12\relax}
%    \end{macrocode}
% \end{macro}
% \begin{macro}{\RereadBiblio}
% Define a macro to do all the fiddly stuff. This is used at the
% beginning of |thebibliography| environment (Sec~\ref{sec:biblio}).
%    \begin{macrocode}
\newread\tf@biblio
\newcommand{\RereadBiblio}{%
%    \end{macrocode}
% Open the current file again for reading
%    \begin{macrocode}
  \begingroup
  \immediate\openin\tf@biblio=\articlename
%    \end{macrocode}
% Everything in a group to be safe. Initialize logicals. Black magic
% to change
% catcodes and avoid blank lines being interpreted as |\par|. 
%    \begin{macrocode}
  \AnotherLinetrue\InBibfalse
  \makeOther\ \makeOther\\\makeOther\$%
  \makeOther\#\makeOther\^\makeOther\^^K%
  \makeOther\_\makeOther\^^A\makeOther\%%
  \makeOther\~\makeOther\{\makeOther\}\makeOther\&%
  \endlinechar-1\relax
%    \end{macrocode}
% Start looping through the file. Read each line to the macro
% |\Line|. If we have got to EOF or end of biblio then set
% |\AnotherLinefalse|.
%    \begin{macrocode}
  \loop
  \read\tf@biblio to \Line
  \ifx\Line\EndBib\AnotherLinefalse\InBibfalse\fi
  \ifeof\tf@biblio \AnotherLinefalse\fi
%    \end{macrocode}
% If we are in the biblio and the line is not blank then append it to
% the |\ADS@references| macro together with a newline. 
%    \begin{macrocode}
  \ifInBib
  \ifx\Line\empty\else\xdef\ADS@references{\ADS@references\Line^^J}\fi
  \fi
%    \end{macrocode}
% Check if the line we just saw was the start of biblio, then repeat
% the loop if necessary. 
%    \begin{macrocode}
  \ifx\Line\StartBib\InBibtrue\fi
  \ifx\Line\StartBibX\InBibtrue\fi
  \ifAnotherLine\repeat
  \immediate\closein\tf@biblio
  \endgroup
  }
%    \end{macrocode}
% \end{macro}
% 
% \begin{macro}{\PadToFourWithDots}
% Utility function for left-padding a number to 4 characters with
% periods. Used in 
% constructing |\ADS@BibCode|. This is a very simple algorithm and
% \emph{requires} that the argument actually be a number---if it is
% not then dire consequences will result.  
%    \begin{macrocode}
\newcommand{\PadToFourWithDots}[1]{%
  \ifnum #1 > 999 #1%
  \else\ifnum #1 > 99 .#1%
  \else\ifnum #1 > 9 ..#1%
  \else ...#1%
  \fi\fi\fi}
%%  \ifnum \count0 > 999 #1%
%%  \else\ifnum \count0 > 99 .#1%
%%  \else\ifnum \count0 > 9 ..#1%
%%  \else ...#1%
%%  \fi\fi\fi}
%    \end{macrocode}
% \end{macro}
%
% \begin{macro}{\AuthorADS}
% \begin{macro}{\ADS@affiliations}
% New mechanism for setting the affiliations for ADS. The author labels (AA, AB, \dots) are generated using two nested counters.  We use the |\robustify| command from \texttt{etoolbox} so that accented characters in the affiliation addresses do not blow up.  The list of accents may need to be expanded. 
%    \begin{macrocode}
\newcounter{AUTHOR}
\newcounter{author}[AUTHOR]
\stepcounter{AUTHOR}
\newcommand\ADS@affiliations{}
\robustify{\'}\robustify{\i}\robustify{\~}\robustify{\`}
\def\AuthorADS#1#2{%
  \begingroup
  \def\ORCIDiD##1{<ID system="ORCID">##1</ID>}
  \def\Email##1{<EMAIL>##1</EMAIL>}
  \def\Affil##1{##1}
  \stepcounter{author}
  \ifthenelse{\value{author} > 26}{\stepcounter{AUTHOR}\stepcounter{author}}{}
  \ifthenelse{\equal{\ADS@affiliations}{\@rmaaEmpty}}{%
    \xdef\ADS@affiliations{\Alph{AUTHOR}\Alph{author}(#2)}
  }{%
    \xdef\ADS@affiliations{\ADS@affiliations, \Alph{AUTHOR}\Alph{author}(#2)}
  }%
  \endgroup
}
%    \end{macrocode}
% \end{macro}
% \end{macro}
% %
% \begin{macro}{\ADS@dumpdata}
% \changes{4.0}{2019/11/13}{Write all author affiliations for ADS}
% \begin{macro}{\ADS@write}
% \begin{macro}{\ADS@openfile}
% |\ADS@dumpdata| is the macro that actually does the job of writing out all the
% data to an auxiliary file, using |\ADS@write| to write each
% field. The file is set up with |\ADS@openfile|, which currently is
% done in the TOC of a book.  The new multi-author affiliations are used if |\ADS@affiliations| has been set, otherwise it falls back to the single-author version |\ADS@affil|. 
%    \begin{macrocode}
\newcommand{\ADSfield}[2]{\relax}
\newcommand{\ADSvfield}[2]{\relax}
%%\renewcommand{\ADSfield}[12]{\%#1 #2}
\newcommand{\ADS@dumpdata}
{ 
  \ADS@write{\string\begin{ADSarticle}}
  \ADS@WRITE{\string\ADSfield{R}{\ADS@BibCode}}
  \ADS@write{\string\ADSfield{A}{\ADS@authors}}
  \ifthenelse{\equal{\ADS@affiliations}{\@rmaaEmpty}}{%
    {\ADS@sanitize\ADS@write{\string\ADSfield{F}{AA(\ADS@affil)}}}
  }{%
    {\ADS@sanitize\ADS@write{\string\ADSfield{F}{\ADS@affiliations}}}
  }
  \ADS@write{\string\ADSfield{J}{\ADS@journalname}}
  \ADS@write{\string\ADSfield{D}{\ADS@fulldate}}
  \ADS@write{\string\ADSfield{L}{\ADS@lastpage}}
  {\ADS@sanitize\ADS@write{\string\ADSfield{T}{\ADS@title}}}
  \ADS@write{\string\ADSfield{C}{\ADS@copyright}}
  \ADS@write{\string\ADSfield{d}{\ADS@doi}}
  \ADS@write{\string\ADSfield{U}{\ADS@url}}
  \ADS@write{\string\ADSfield{K}{\ADS@keywords}}
  {\ADS@sanitize\ADS@write{\string\ADSfield{B}{\ADS@abstract}}}
  \ADS@write{\string\ADSfield{Z}{\ADS@references}}
  \ADS@write{\string\end{ADSarticle}}
  }
%%\newcommand{\ADS@write}[1]{\if@book\addtocontents{ads}{#1}\fi}
%%\newcommand{\ADS@openfile}{\@starttoc{ads}}
\newcommand{\ADS@write}[1]{\set@display@protect\immediate\write\tf@ads{#1}}
\newcommand{\ADS@WRITE}[1]{\immediate\write\tf@ads{#1}}
\newcommand{\ADS@openfile}{%
  \newwrite\tf@ads%
  \immediate\openout\tf@ads\jobname.ads\relax
  \ADS@WRITE{\ADS@preamble}}
\newcommand{\ADS@closefile}{\ADS@WRITE{\ADS@postamble}}
\def\ADS@preamble{%
  %%  \string\documentclass{article}
  %%  \string\begin{document}
  }
\def\ADS@postamble{%
  %%  \string\end{document}
  \string%\string% Local Variables:^J 
  \string%\string% mode: LaTeX^J
  \string%\string% mode: Font-lock^J
  \string%\string% End:^J
  }
%    \end{macrocode}
% \end{macro}
% \end{macro}
% \end{macro}
% 
% \subsection{Automatic calculation of the page range}
% \label{sec:page-range}
% Another issue that has cost me a lot of grief. Still, it finally
% seems to be sorted now. All hail Jeff Goldberg. 
%
% Cunning trick to get the page numbers from |\pageref|.  Put a label
% on the first and last page of the document using |\AtBeginDocument|
% and |\AtEndDocument|. .  This way won't work for the book option.
% In that case, we do it by defining |\@endofarticlehook| and
% |\@begofarticlehook| as parallels to |\AtEndDocument| and
% |\AtBeginDocument|.  Note that we have to defer adding
% |\@begofarticlehook| to |\AtBeginDocument| to make sure it comes
% after the redefinitions of |\ref| and |\label|.
%    \begin{macrocode}
\if@book\AtEndDocument{\ADS@closefile}\else
 %\AtBeginDocument{\@begofarticlehook}
\AtEndDocument{\@endofarticlehook}
\fi
%    \end{macrocode}
% \begin{macro}{\@begofarticlehook}
% Since we are getting more and more things that need to be done at
% the beginning of every article, frontsection, etc we define a sister
% hook to 
% |\@endofarticlehook|. This hook does the following: 
% \begin{enumerate}
% \item Labels the first page
% \item Sets |\ADS@firstpage|
% \item runs and then clears |\indexauthor@deferlist|
% \item resets the keyword counter
% \end{enumerate}
% It is executed in the following circumstances: 
% \begin{enumerate}
% \item At |\begin{document}| for standalone files
% \item At |\begin{document}| for each included file
% \item At beginning of front sections 
% \item At beginning of TOC
% \item At beginning of Author Index
% \item At beginning of abstracts
% \end{enumerate}
% For some reason, this is not always within the right group, so we
% use |\xdef| just to be sure. 
%    \begin{macrocode}
\newcommand{\@begofarticlehook}{%
  \if@frontmatter\frontlabel{RMAAFirstPage}\else\label{RMAAFirstPage}\fi%
  \xdef\ADS@firstpage{\@arabic{\c@page}}%
  \if@book\else\makeindex\fi % needed so default index is defined in stand-alone articles
%%  \indexauthor@deferlist
%%  \xdef\indexauthor@deferlist{}
  \@inpreamblefalse
  \setcounter{kwdcounter}{0}
  }
%    \end{macrocode}
% \end{macro}
% \begin{macro}{\@endofarticlehook}
% This is a hook to be run at the end of each article
% currently, it just labels the last page and writes out all the ADS
% data to an auxiliary file.  
% It should not produce any printed output. 
% It is modelled on the |\AtEndDocument| command, to which it reduces
% in the case of a single article. 
%    \begin{macrocode}
\newcommand{\@endofarticlehook}{%
  \if@debug{\ClassWarning{rmaa}{***Doing Hook***}}\fi
  \if@debug{\ClassWarning{rmaa}{ADSabstract is: \ADS@abstract}}\fi
  \clearpage\rmaa@lastpagelabel\if@book\ADS@dumpdata\fi
  \ifx\NotesList\@rmaaEmpty\else\OutputNotes\fi
  \if@debug{\ClassWarning{rmaa}{Notes list is: \NotesList}}\fi
}
%    \end{macrocode}
% \begin{macro}{\rmaa@lastpagelabel}
% New version of last page label command copied from Jeff Goldberg's
% lastpage.sty package - we can't just use |\label| here in case the
% output routine for the last page has already been called. Instead,
% we first explicitly clear the page, step back the page counter,
% explicitly write to the .aux file, then restep the page counter
% again.
%
% But this will conflict with the hyperref package, which does labels differently.  Compensate for that by using the five-argument form.
% 
% Also set the |\ADS@LastPage| macro to the appropriate number. 
%    \begin{macrocode}
\newcommand{\rmaa@lastpagelabel}{%
  \addtocounter{page}{-1}%
  \immediate\write\@auxout{\string
    \newlabel{\articlename:RMAALastPage}{{}{\thepage}{}{}{}}}%
  \edef\ADS@lastpage{\thepage}%
  \addtocounter{page}{1}}
%    \end{macrocode}
% \end{macro}
% \end{macro}

% \subsection{Deal with explicit change in number of columns}
% \label{sec:columns}
% Sometimes, when formatting a book, we need to switch between
% one-column articles and two-column articles. For this to work, we
% need to wrap the standard \LaTeX{} commands to set the |\@rmtwocol|
% flag appropriately and also reset the package default values of all
% scalable lengths. .  
%    \begin{macrocode}
\let\latexonecolumn=\onecolumn
\let\latextwocolumn=\twocolumn
\renewcommand{\onecolumn}{\@rmtwocolfalse\@package@default@lengths\latexonecolumn}
\renewcommand{\twocolumn}{\@rmtwocoltrue\@package@default@lengths\latextwocolumn}
%    \end{macrocode}
% 
% \subsection{The actual formatting of the title}
% \label{sec:title-format}

% Note, we have to redefine the |\maketitle| command to use the
% numbered footnotes. This means we needn't use the atbeginend stuff. 
%    \begin{macrocode}
\renewcommand{\maketitle}{\par
  \begingroup
%    \end{macrocode}
% Don't switch to marked footnotes. 
% This is consistent with existing style but I'm not sure I like it
%    \begin{macrocode}
  %%  \renewcommand\thefootnote{\fnsymbol{footnote}}%
%%  \def\@makefnmark{\hbox to\z@{$\m@th^{\@thefnmark}$\hss}}%
  \def\@makefnmark{\textsuperscript{\@thefnmark}\hspace{-0.2em}}%
  \long\def\@makefntext##1{\parindent 1em\noindent
    \hbox to1.8em{\hss$\m@th^{\@thefnmark}$}##1}%
  \if@rmtwocol
  { \if@debug{\ClassWarning{rmaa}{TWO COLUMN}}\fi
    \twocolumn[\@maketitle]%
    }
  \else
  { \if@debug{\ClassWarning{rmaa}{ONE COLUMN}}\fi
    \newpage
    \global\@topnum\z@   % Prevents figures from going at top of page.
    \if@manuscript \@MSmaketitle \else \@maketitle \fi
    }
  \fi
  \if@manuscript\else\thispagestyle{rmaafirstpage}\@thanks\fi
  \ifx\NotesList\@rmaaEmpty\else{%
    \NoteMark{1--\arabic{notecounter}}\vspace*{-2\baselineskip}%
  }\fi%
  \endgroup
%    \end{macrocode}
% Don't reset footnote counter (not sure whether this is good thing
% or not).
%    \begin{macrocode}
  %%  \setcounter{footnote}{0}%
  \let\thanks\relax\let\affil\relax\let\altaffil\relax
  \let\maketitle\relax\let\@maketitle\relax
  \gdef\@thanks{}\gdef\@author{}\gdef\@title{}}
%    \end{macrocode}

% \begin{macro}{\@kwhwidth}
% \begin{macro}{\@kwtwidth}
% Redefine \cs{@maketitle} to include the abstract, resumen and
% keywords. Note that in the manuscripty style we use
% \cs{@MSmaketitle} instead.  
% \changes{3.16}{2003/04/07}{
%  Added extra vertical space between authors and accepted date.}
% \changes{3.26}{2011/08/15}{Add extra label to review articles}
%    \begin{macrocode}
\newlength{\@kwhwidth}\newlength{\@kwtwidth}
\newcommand\reviewname{Review}
\renewcommand{\@maketitle}{%
  \newpage
  \null
  \if@reviewarticle
  \smallskip
  {\large\bfseries\itshape\reviewname}%
  \bigskip
  \else
  \vskip\@before@title@skip%
  \fi
  \begin{center}%
    {\@title \par}%
    \vskip\@before@author@skip%
    \parbox{\rmaa@AuthorWidthFrac\textwidth}{\centering
      \lineskip\@between@author@skip%
      \@author
      \par}
    %% this should be replaced with `Submitted to...' for preprints
    \if@printdates
    \vskip\@before@dates@skip%
    \if@preprint
    \rmaa@ReceivedStyle{Submitted to RMxAA \@received}
    \else
    \rmaa@ReceivedStyle{Received \@received; accepted \@accepted}
    \fi
    \fi
    \if@manuscript
    \rmaa@ReceivedStyle{Submitted to RMxAA \@received}
    \fi
    \if@guide
    \rmaa@ReceivedStyle{Version \rmaa@version, \rmaa@date}
    \fi
    \vskip\@before@abstract@skip%
    \if@abstractpresent
    \parbox{\rmaa@AbstractWidthFrac\textwidth}{%
      \setlength{\parindent}{0.9cm}%
      \centerline{\rmaa@AbstractHeaderStyle{Abstract}}\par%
      \vskip\@after@abstracthead@skip%
      \if@newstyle\noindent\fi\@abstract
      \vskip\@between@abstract@skip%
      \centerline{\rmaa@AbstractHeaderStyle{Resumen}}\par%
      \vskip\@after@abstracthead@skip%
      \if@newstyle\noindent\fi\@resumen
      \if@setresumen\else\@abstract\fi % append english abstract if
                                       % spanish one not set
      \if@guide\else            %do not use keywords for guide
      \vskip\@before@kwd@skip%
      % Format the Key Words with hanging indentation:
      % They are put in a \parbox with width equal to the line width
      % minus the width of the header. 
      { % Widths of keyword
        \settowidth{\@kwhwidth}{\@keywordheader~}   % header and
        \setlength{\@kwtwidth}{\linewidth}          % keyword text
        \addtolength{\@kwtwidth}{-\@kwhwidth}       %
        \noindent\@keywordheader~\parbox[t]{\@kwtwidth}{%
          \rmaa@KeywordStyle{\boldmath\uppercase{\@keywords}}}}%
      \fi
      }
      \fi
  \end{center}%
  \vspace*{\afterkwdskip}
  \par
  }
\newlength{\afterkwdskip}\setlength{\afterkwdskip}{\z@}
%    \end{macrocode}
% \end{macro}
% \end{macro}
%
%
% \begin{macro}{\reviewarticle}
% \begin{macro}{\reviewtoc}
% \begin{macro}{\reviewlof}
%   User command to specify a review article. Must be used in
%   preamble. Main effect is to place the text ``Review'' before the
%   article title.  Also add command for setting a table of contents
%   and a list of figures at the beginning of the article.
%    \begin{macrocode}
\newcommand{\reviewarticle}{\@reviewarticletrue}
\newcommand\reviewtoc{%
  \section*{CONTENTS}%
  \@starttoc{rtc}%
}
\newcommand\reviewlof{%
  \section*{LIST OF FIGURES}%
  \@starttoc{lof}%
}
%    \end{macrocode}
% \end{macro}
% \end{macro}
% \end{macro}
% \begin{macro}{\l@section}
%    \begin{macrocode}
\renewcommand*\l@section[2]{%
  \ifnum \c@tocdepth >\z@
    \addpenalty\@secpenalty
    \addvspace{0.8em \@plus\p@}%
    \setlength\@tempdima{1.5em}%
    \begingroup
      \parindent \z@ \rightskip \@pnumwidth
      \parfillskip -\@pnumwidth
      \leavevmode %\bfseries
      \advance\leftskip\@tempdima
      \hskip -\leftskip
      #1\nobreak\hfil \nobreak\hb@xt@\@pnumwidth{\hss #2}\par
    \endgroup
  \fi}
%    \end{macrocode}
% \end{macro}
%
% \begin{macro}{\@MSmaketitle}
% A special version of |\@maketitle| for the \texttt{manuscript}
% option. This can deal with a multi-page abstract/resumen. Not very
% well, mind. 
%    \begin{macrocode}
\newcommand{\@MSmaketitle}{%
  \if@debug\ClassInfo{rmaa}{Using special MS maketitle command}\fi
  \newpage
  {
    \thispagestyle{rmaafirstpage}
%     \setlength{\parindent}{0pt}
    \null\vskip 2ex
    \begin{center}
      \@title\par
      \@author\par
      \rmaa@ReceivedStyle{%
        \ifx\@received\@rmaaEmpty
          Draft version: \@date
        \else
          Received: \@received
          \ifx\@accepted\@rmaaEmpty
          \else
            ; Accepted: \@accepted
          \fi
        \fi}\par
    \end{center}
    \@thanks
    \centerline{\rmaa@AbstractHeaderStyle{Resumen}}
    \noindent\@resumen\par\vskip 2ex
    \centerline{\rmaa@AbstractHeaderStyle{Abstract}}
    \noindent\@abstract\par\vskip 2ex
    \noindent\@keywordheader~\rmaa@KeywordStyle{\@keywords}\par
%     \@keywordheader~\rmaa@KeywordStyle{\uppercase{\@keywords}}\par
  }%
  }
%    \end{macrocode}
% \end{macro}

% \subsection{The sectioning commands}
% \label{sec:section-commands}

% \begin{macro}{\section}
% \begin{macro}{\subsection}
% \begin{macro}{\subsubsection}
% \begin{macro}{\paragraph}
% \begin{macro}{\subparagraph}
% Redefine all the standard commands to use these lengths. 
%    \begin{macrocode}
\renewcommand\section{\@startsection {section}{1}{\z@}%
                                   {\abovesectionskip}%
                                   {\belowsectionskip}%
                                   %% how do we get upper case here
                                   {\centering\reset@font\normalsize}}
\renewcommand\subsection{\@startsection{subsection}{2}{\z@}%
                                     {\abovesectionskip}%
                                     {\belowsectionskip}%
                                     {\centering\reset@font\normalsize\itshape}}
\renewcommand\subsubsection{\@startsection{subsubsection}{3}{\z@}%
                                     {\abovesectionskip}%
                                     {\belowsectionskip}%
                                     {\centering\reset@font\normalsize\itshape}}
\renewcommand\paragraph{\@startsection{paragraph}{4}{\z@}%
                                    {\aboveparaskip}%
                                    {\belowparaskip}%
                                    {\reset@font\normalsize\bfseries}}
\renewcommand\subparagraph{\@startsection{subparagraph}{5}{\parindent}%
                                      {\aboveparaskip}%
                                      {\belowparaskip}%
                                      {\reset@font\normalsize\bfseries}}
%    \end{macrocode}
% \end{macro}
% \end{macro}
% \end{macro}
% \end{macro}
% \end{macro}

% \begin{macro}{secnumdepth}
% Make sure that we number the subsubsections---this is done by
% default in the |article| class but not the |book| class, which is
% why we must set it explicitly here.
%    \begin{macrocode}
\setcounter{secnumdepth}{3}
%    \end{macrocode}
% \end{macro}
% \changes{2.2beta10}{2001/11/07}{Fixed \cs{secnumdepth} for \cs{book}
% style.}

% \begin{macro}{\thesubsection}
% \begin{macro}{\thesubsubsection}
% For subsection and subsubsection, we want the numbers to be roman
% but the title to be italic. This seems the easiest way to go about it.
%    \begin{macrocode}
\renewcommand{\thesubsection}{{\upshape\thesection.\@arabic\c@subsection}}
\renewcommand{\thesubsubsection}{\thesubsection{\upshape.\@arabic\c@subsubsection}}
%    \end{macrocode}
% \end{macro}
% \end{macro}

% \begin{macro}{\@seccntformat}
% Add ``.'' to the number here so we can remove it from |\thesection|,
% etc. Reduce gap between number and title WJH 27 Jan 1995
%    \begin{macrocode}
\def\@seccntformat#1{\csname the#1\endcsname .\hskip 0.4em}
%    \end{macrocode}
% \end{macro}

% \begin{macro}{\@sect}
% I didn't really want to mess with the |\@sect| command, but it was
% the only way I could manage to force upper case. Other than the
% commented parts, it is identical to the standard version. 
%    \begin{macrocode}
\def\@sect#1#2#3#4#5#6[#7]#8{
  \ifnum #2>\c@secnumdepth
     \let\@svsec\@empty
  \else
     \refstepcounter{#1}%
     \let\@@protect\protect
     \def\protect{\noexpand\protect\noexpand}%
     \edef\@svsec{\@seccntformat{#1}}%
     \let\protect\@@protect
  \fi
%    \end{macrocode}
% This is the bit I changed (27 Jan 1995). Use upper case for
% |\section|, use mixed case for the rest. 
%    \begin{macrocode}
  \ifnum #2=1                   
     \def\sec@upcase##1{\uppercase{##1}}%
     \def\sec@shape{\upshape}%
  \else                         
     \def\sec@upcase##1{\relax{##1}}%
     \def\sec@shape{\itshape}%
  \fi
  \@tempskipa #5\relax
  \ifdim \@tempskipa>\z@
     \begingroup #6\relax
          \@hangfrom{\hskip #3\relax\@svsec}%
                    {\interlinepenalty \@M \sec@upcase{#8}\par}%
     \endgroup
     \csname #1mark\endcsname{#7}%
%    \end{macrocode}
% Only write TOC entries if doing the Guide style. 
%    \begin{macrocode}
     \if@guide 
        \addcontentsline{toc}{#1}{%
          \ifnum #2>\c@secnumdepth \else
             \protect\numberline{\csname the#1\endcsname}%
          \fi
          #7}%
     \fi
%    \end{macrocode}
% Also allow per-article TOC for review papers
%    \begin{macrocode}
     \if@reviewarticle
        \ifnum #2>2 \else
           \addcontentsline{rtc}{#1}{%
             \protect\numberline{\csname the#1\endcsname}%
             \sec@upcase{\sec@shape #7}}%
        \fi
     \fi
  \else
     \def\@svsechd{#6\hskip #3\relax  %% \relax added 2 May 90
       \@svsec #8\csname #1mark\endcsname{#7}%
       \if@guide
          \addcontentsline{toc}{#1}{%
            \ifnum #2>\c@secnumdepth \else
               \protect\numberline{\csname the#1\endcsname}%
            \fi
            #7}%
       \fi
       }%
  \fi
  \@xsect{#5}}
%    \end{macrocode}
% \end{macro}


% \subsection{Appendices}
% \label{sec:appendices}
% Make sure that appendices in individual articles work properly. We
% copy the definition from the article class (the book one doesn't do
% what we want).  
% \begin{environment}{appendices}
%   This uses an environment |appendices| inside of which each
%   |\section| will become a separate appendix.  Note that the
%   appendices must come before the references (otherwise the full
%   addresses mechanism will not work properly). We make the equation
%   numbering be subordinate to the appendix (equations A12, B6,
%   etc). Finally, we have a medskip followed by APPENDICES to
%   introduce them. \changes{3.17}{2002/04/08}{ Added appendices environment
%   that does the right thing.}
%    \begin{macrocode}
\newenvironment{appendices}{\par
  \begingroup
  \setcounter{section}{0}%
  \setcounter{subsection}{0}%
  \renewcommand\thesection{\@Alph\c@section}
  \renewcommand\theequation{\thesection\arabic{equation}}
  \medskip
  \section*{APPENDICES}}
{\endgroup}
%    \end{macrocode}
% \end{environment}
% \begin{environment}{appendix}
%   If there is only one appendix, then use |appendix| environment
%   instead.  \changes{4.0}{2019/11/12}{ Added
%   appendix environment for case where there is only one section.}
%    \begin{macrocode}
\renewenvironment{appendix}{\par
  \begingroup
  \setcounter{section}{0}%
  \setcounter{subsection}{0}%
  \renewcommand\thesection{\@Alph\c@section}
  \renewcommand\theequation{\thesection\arabic{equation}}
  \medskip
  \section*{APPENDIX}}
{\endgroup}
%    \end{macrocode}
% \end{environment}

% \subsection{Floating thingies}
% \label{sec:floats}   

%% 
%    \begin{macrocode}
\renewcommand{\thetable}{\@arabic\c@table}
\renewcommand{\thefigure}{\@arabic\c@figure}
\renewcommand{\figurename}{Fig.}
\renewcommand{\tablename}{TABLE}
%\renewcommand{\fnum@table}{\MakeUpperCase{\tablename}~\thetable}
%    \end{macrocode}
% For some reason, I have to write `TABLE' explicitly here
% If I use |\tablename|, it gets converted to mixed case---God knows
% why. 
% |\figurename| is also being ignored. 
%    \begin{macrocode}
\renewcommand{\fnum@table}{TABLE~\thetable}
\renewcommand{\fnum@figure}{Fig.~\thefigure}
%    \end{macrocode}

% Do captions in 9 pt face in the main journal style and in new-style
% anything. 
%    \begin{macrocode}
\long\def\@makefigcap#1#2{%
  \if@proceedings\else\small\fi
  \if@newstyle\small\fi 
  \vskip\abovecaptionskip
  \sbox\@tempboxa{#1. #2}%
  \ifdim \wd\@tempboxa >\hsize
    #1. #2 \par
  \else
    \global \@minipagefalse
    \hb@xt@\hsize{\hfil\box\@tempboxa\hfil}%
  \fi
  \vskip\belowcaptionskip}
%    \end{macrocode}

% Make sure table caption is upper case. 
%    \begin{macrocode}
\long\def\@maketabcap#1#2{%
  \hbox to\hsize{\hfil #1 \hfil}
  \vskip 6pt 
  \centering\MakeUpperCase{#2}\par
  \vskip 6pt}
%    \end{macrocode}

%    \begin{macrocode}
\newlength{\@rottabcapwidth}
\long\def\@makerottabcap#1#2{%
  \setbox\@tempboxa\hbox{\MakeUpperCase{#2}}%
  \ifdim\wd\@tempboxa >.8\vsize
     \setlength\@rottabcapwidth{.8\vsize}
  \else
     \setlength\@rottabcapwidth{\wd\@tempboxa}
  \fi
  \rotatebox{90}{%
    \begin{minipage}{\vsize}\centering
      \begin{minipage}{\@rottabcapwidth}\centering
        #1\\ \vskip 6pt \MakeUpperCase{#2}
      \end{minipage}%
    \end{minipage}
    }%
  \hskip 6pt
  }

\newenvironment{Sideways}
    {\begin{sideways}\begin{minipage}{\vsize}\centering}
    {\end{minipage}\end{sideways}}

\newenvironment{Sidewaystable}
{%
  \let\oldcaption=\caption%
  \def\caption{
    \def\@captype{table}\renewcommand{\@makecaption}{\@maketabcap}
    \oldcaption}
  \begin{sidewaystable}\centering}
  {\end{sidewaystable}\let\caption=\oldcaption}
\newenvironment{Sidewaystable*}
{%
  \let\oldcaption=\caption%
  \def\caption{
    \def\@captype{table}\renewcommand{\@makecaption}{\@maketabcap}
    \oldcaption}
  \begin{sidewaystable*}\centering}
  {\end{sidewaystable*}\let\caption=\oldcaption}
%    \end{macrocode}

% Not sure why we need to stick these lines in. 
%    \begin{macrocode}
\renewenvironment{figure}
                 {\renewcommand{\@makecaption}{\@makefigcap}\@float{figure}}
                 {\end@float}
\renewenvironment{figure*}
                 {\renewcommand{\@makecaption}{\@makefigcap}\@dblfloat{figure}}
                 {\end@dblfloat}
\renewenvironment{table}
                 {\renewcommand{\@makecaption}{\@maketabcap}\@float{table}}
                 {\end@float}
\renewenvironment{table*}
                 {\renewcommand{\@makecaption}{\@maketabcap}\@dblfloat{table}}
                 {\end@dblfloat}
%    \end{macrocode}

% Redefine |\@caption| so that we don't bother with the list of
% figures. This saves us having to worry about fragile commands in the
% captions. 
%    \begin{macrocode}
\long\def\@caption#1[#2]#3{%
  \par
%%  \addcontentsline{\csname ext@#1\endcsname}{#1}%
%%    {\protect\numberline{\csname the#1\endcsname}{\ignorespaces #2}}%
  \begingroup
    \@parboxrestore
    \if@minipage
      \@setminipage
    \fi
    \normalsize
    \@makecaption{\csname fnum@#1\endcsname}{\ignorespaces #3}\par
  \endgroup}
%    \end{macrocode}

% Make tables not quite so cramped. Allow |\hline\hline| to give a
% poor man's thick rule. 
%    \begin{macrocode}
\renewcommand{\arraystretch}{1.15}
\setlength{\doublerulesep}{0.4pt}
%    \end{macrocode}

% \begin{macro}{\tablecols}
% \begin{macro}{\@tablecols}
% \begin{macro}{\tabcaptioncontinued}
% |\tablecols| is a user command for setting the number of columns. It
% is used by commands that try and span the entire table (I think
% there must be a better way of acheiving
% this). |\tabcaptioncontinued| is 
% a user command for setting the caption of a continued table on the
% next page. 
%    \begin{macrocode}
\newcommand{\tablecols}[1]{\renewcommand{\@tablecols}{#1}}
\newcommand{\@tablecols}{1}
\newcommand{\tabcaptioncontinued}{%
  \multicolumn{\@tablecols}{c}{\addtocounter{table}{-1}TABLE~\thetable\ (CONTINUED)} \\ \\
  }
%    \end{macrocode}
% \end{macro}
% \end{macro}
% \end{macro}

% \subsubsection{Footnotes in tables}
% \label{sec:table-footnotes}
% \begin{macro}{\tablenotemark}
% \begin{macro}{\tablenote}
% \begin{macro}{\tablenoteskip}
% Old user commands |\tablenotemark{MARK}|,
% |\tablenote{MARK}{TEXT}|---now deprecated. 
%    \begin{macrocode}
\newcommand\tablenotemark[1]{$^\mathrm{#1}$} 
\newcommand\tablenote[2]{\if@proceedings\else\small\fi ${}^\mathrm{#1}$ #2} 
\newcommand\tablenoteskip{\\[0.5\baselineskip]} 
%    \end{macrocode}
% \end{macro}
% \end{macro}
% \end{macro}

% \begin{macro}{\tabnotemark}
% \begin{macro}{\tabnotetext}
% \begin{macro}{\tabnotetextTWO}
% Alternative way of doing footnotes in tables. 
%    \begin{macrocode}
\newcommand{\tabnotemark}[1]{\textsuperscript{\lowercase{#1}}}
%\newcommand{\tabnotetextTWO}[4]{\tabnotetext{#1}{#2}\tabnotetext{#3}{#4}}
\newcommand{\tabnotetextTWO}[4]{%
  \multicolumn{\@tablecols}{l}{%
    \parbox{0.5\tabnotewidth}{\small\textsuperscript{#1}#2}%
    \parbox{0.5\tabnotewidth}{\small\textsuperscript{#3}#4}%
    }%
  \\}
\newcommand{\tabnotetext}[2]{%
  \multicolumn{\@tablecols}{@{}p{\tabnotewidth}@{}}{\small\textsuperscript{#1}#2}\\}
\newlength{\tabnotewidth}
%    \end{macrocode}
% \end{macro}
% \end{macro}
% \end{macro}


% \begin{macro}{changemargin}
%    \begin{macrocode}
% Command to allow margins to be changed on the fly. Useful to allow
% tables and Figures to protrude into the margins. 
%
% \changes{3.24}{2007/09/14}{Added \texttt{changemargin} environment
% (copied from TeX FAQ).}
\newenvironment{changemargin}[2]{% 
  % From TeX FAQ http://www.tex.ac.uk/cgi-bin/texfaq2html?label=chngmargonfly
  \begin{list}{}{%
    \setlength{\topsep}{0pt}%
    \setlength{\leftmargin}{#1}%
    \setlength{\rightmargin}{#2}%
    \setlength{\listparindent}{\parindent}%
    \setlength{\itemindent}{\parindent}%
    \setlength{\parsep}{\parskip}%
  }%
  \item[]}{\end{list}}
%    \end{macrocode}
% \end{macro}

% \subsection{The bibliography and references}
% \label{sec:biblio}
% \begin{macro}{\beforebiblioskip}
% \begin{macro}{\afterbiblioskip}
% I'm not entirely sure why I introduced these two commands, since
% directly putting |\vspace| before and after
% |\begin{thebibliography}| seems to work well enough. 
%    \begin{macrocode}
\newlength{\beforebiblioskip}
\setlength{\beforebiblioskip}{0pt}
\newlength{\afterbiblioskip}
\setlength{\afterbiblioskip}{0pt}
%    \end{macrocode}
% \end{macro}
% \end{macro}
% \begin{macro}{thebibliography}
% The environment itself is done as starred section with extra
% vertical space after heading. This space is currently zero. 
% In the book style we do the cunning trick with rereading the file in
% order to get the references into the |\ADS@references| macro.
% 
% \changes{3.05}{2003/01/02}{Removed the \cs{beforebiblioskip} and
% \cs{afterbiblioskip} commands, since they cause unwanted vertical
% space in some instances.}  
%    \begin{macrocode}
\renewenvironment{thebibliography}{%
  \if@book\RereadBiblio\fi
%%  \vspace*{\beforebiblioskip}
  \section*{REFERENCES} \frenchspacing
%%  \vspace*{\afterbiblioskip}
  \small
%    \end{macrocode} 
% The references themselves are set as a list in order 
% to get the hanging indentation. They are done in a smaller font and
% with french spacing so we needn't worry about the distinction
% between sentence-ending periods and abbreviations. 
%    \begin{macrocode} 
  \list{\null}{\leftmargin 1.5em\labelwidth\z@\labelsep\z@\itemindent -1.5em
    \itemsep\z@\parsep\z@\usecounter{enumi}}
  \sloppy\clubpenalty4000\widowpenalty4000
  \sfcode`\.=1000\relax\@biblisttrue}
{\endlist\@biblistfalse\normalsize\@outputaddressesifnecessary}
%    \end{macrocode}
% \end{macro}
% \subsection{Uniquification of cross-references}
% \label{sec:uniquification}
% 
% \begin{macro}{\articlename}
% In order that the |\label{}| command in individual articles be
% guaranteed to give a unique anchor, we append the article name (held
% in the |\articlename| command) to each key. This requires
% redefinition of the following commands: |\cite|, |\nocite|, |\ref|,
% |\pageref|, |\label|. 
% \begin{macro}{\oldcite}
% \begin{macro}{\oldnocite}
% \begin{macro}{\oldref}
% \begin{macro}{\oldpageref}
% \begin{macro}{\oldlabel}
% The original versions of these commands are
% saved as |\oldcite|, |\oldnocite|, etc. 
%
% The uniquifying of cross-references is now always done, so 
% we don't get multiply-defined labels in the single article case
%    \begin{macrocode}
\newcommand{\articlename}{anon}
\AtBeginDocument{
  \let\oldcite=\cite 
  \let\oldnocite=\nocite 
  \let\oldref=\ref 
  \let\oldpageref=\pageref
  \let\oldlabel=\label
}
%    \end{macrocode}
% \end{macro}
% \end{macro}
% \end{macro}
% \end{macro}
% \end{macro}
% \end{macro}
%
% Uniquify all cross-references by prepending the article name.  This
% now has to go into |\AtBeginDocument| so it is working on the new
% versions that hyperref.sty has set up.  Then we need to put the
% |\@begofarticlehook| afterwards, so that the first page label is using
% the new definitions.  If amsmath is loaded, then we also need to patch
% the special label command that they use. \changes{4.0}{2019/11/18}{Fix labels for amsmath display equations, in case that package has been loaded.}
%    \begin{macrocode}
\AtBeginDocument{
  \renewcommand{\ref}[1]{\oldref{\articlename:#1}%
    \if@debug{\ClassInfo{rmaa}{Reference to \articlename:#1}}\fi}
  \renewcommand{\pageref}[1]{\oldpageref{\articlename:#1}%
    \if@debug{\ClassInfo{rmaa}{Page reference to \articlename:#1}}\fi}
  \renewcommand{\label}[1]{\oldlabel{\articlename:#1}%
    \if@debug{\ClassInfo{rmaa}{Label for \articlename:#1}}\fi}%
  \ifdef{\label@in@display}{%
    \let\oldlabel@in@display=\label@in@display
    \renewcommand{\label@in@display}[1]{\oldlabel@in@display{\articlename:#1}%
      \if@debug{\ClassInfo{rmaa}{AMSTeX Label for \articlename:#1}}\fi}%
  }{}
  \@begofarticlehook
}%
%    \end{macrocode}
% Citations are more complicated since they can contain a list of
% keys. These are really only here for use with the rmaacite package,
% which is utterly deprecated (02 Jan 2003) now that I have the natbib
% code working. The command \cs{cite} should \emph{not} be used with
% natbib: use \cs{citep}, \cs{citet}, etc.\@ instead. 
%    \begin{macrocode}
\renewcommand{\nocite}[1]{\@bsphack
  \@for\@citeb:=#1\do{%
    \edef\@citeb{\expandafter\@firstofone\@citeb:\articlename}%
    \if@filesw\immediate\write\@auxout{\string\citation{\@citeb}}\fi
    \@ifundefined{b@\@citeb:\articlename}{\G@refundefinedtrue
      \@latex@warning{Citation `\@citeb' undefined}}{}}%
  \@esphack}
\renewcommand{\cite}[1]{\oldcite{#1:\articlename}}
%   \end{macrocode}


% Horrible fiddling with the bibliography citation mechanism. 
% This has to be globally available in case rmaacite puts it in the
% \texttt{.aux} file. 
%    \begin{macrocode}
\def\lbibcite#1#2#3{\global\@namedef{b@#1}{#2}
                    \global\@namedef{by@#1}{#3}}
%    \end{macrocode}

% \begin{macro}{\oldbibitem}
% \begin{macro}{\@rmaaBibArg}
% \begin{macro}{\@rmaaSpace}
% \begin{macro}{rmaabib}
% Of course, |\bibitem| needs to get the uniquification treatment as
% well. While we are at it, we redefine it to use consecutive integers
% for the key if the user leaves it blank (using the counter
% |rmaabib|). This allows one to put 
% |\bibitem{}| for all your references without getting annoying
% messages about multiply-defined labels.
% 
% We also take the opportunity to insert the optional argument to
% \cs{bibitem} if it is missing in the form Author(DATE) so that natbib
% doesn't try to do numerical citations. Natbib already provides a
% hook \cs{@extra@b@citeb} to help uniquify the citation key, so we
% just plug into that. However, it puts it \emph{after} the original
% key, not before as I had had it. Easiest solution is to change the
% way that I do it. 02 Jan 2003. Note that the \cs{@extra@b@citeb}
% does not get written to the .aux  file in the \cs{citation} command
% but it does work. 
% \changes{3.14}{06 Feb 2003}{Reinstalled old definition of \cs{bibitem}
%  for the case where the nonatbib option is requested, so that the
%  rmaacite package will work. This means that you can't use rmaacite
%  and natbib in the same volume.}
%    \begin{macrocode}
\if@natbib\def\@extra@b@citeb{:\articlename}\fi
%%\else
\let\oldbibitem\bibitem
\newcommand{\@rmaaBibArg}{}
\newcommand{\@rmaaSpace}{ }
\newcounter{rmaabib}
\if@natbib
\renewcommand{\bibitem}[2][ABC(2003)]{%
  \renewcommand{\@rmaaBibArg}{#2}
  \ifx\@rmaaBibArg\@rmaaEmpty
     \addtocounter{rmaabib}{1}
     \oldbibitem[#1]{\thermaabib:\articlename}
  \else\ifx\@rmaaBibArg\@rmaaSpace
     \addtocounter{rmaabib}{1}
     \oldbibitem[#1]{\thermaabib:\articlename}
  \else 
     \oldbibitem[#1]{#2:\articlename}
  \fi\fi
  \renewcommand{\@rmaaBibArg}{}
  }
\else
\renewcommand{\bibitem}[1]{%
  \renewcommand{\@rmaaBibArg}{#1}
  \ifx\@rmaaBibArg\@rmaaEmpty
     \addtocounter{rmaabib}{1}
     \oldbibitem{\articlename:\thermaabib}
  \else\ifx\@rmaaBibArg\@rmaaSpace
     \addtocounter{rmaabib}{1}
     \oldbibitem{\articlename:\thermaabib}
  \else
     \oldbibitem{\articlename:#1}
  \fi\fi
  \renewcommand{\@rmaaBibArg}{}
  }
\fi
%%\fi
%    \end{macrocode}
% \end{macro}
% \end{macro}
% \end{macro}
% \end{macro}

% \begin{macro}{oldpcite}
% \begin{macro}{oldscite}
% Then, we need to repeat the above scheme, but for the |rmaacite|
% definition. 
%    \begin{macrocode}
\newcommand{\@redefinebibitem}{%
  \let\oldbibitem\bibitem
  \let\oldcite\cite
  \let\oldpcite\pcite
  \let\oldscite\scite
  \def\bibitem[##1<##2>]##3{%
    \renewcommand{\@rmaaBibArg}{##3}
    \ifx\@rmaaBibArg\@rmaaEmpty
       \addtocounter{rmaabib}{1}
       \oldbibitem[##1<##2>]{\thermaabib:\articlename}
    \else\ifx\@rmaaBibArg\@rmaaSpace
       \addtocounter{rmaabib}{1}
       \oldbibitem[##1<##2>]{\thermaabib:\articlename}
    \else 
       \oldbibitem[##1<##2>]{##3:\articlename}
    % trust the author here
%%       \oldbibitem[##1<##2>]{##3}
    \fi\fi
    }
  }
%    \end{macrocode}
% \end{macro}
% \end{macro}


% \subsection{HTML Table of Contents}
% \label{sec:hyper-TOC}
% \begin{macro}{\HyperTOCEntry}
% \begin{macro}{\HyperTOCpart}
% \begin{macro}{\HyperTOCmeeting}
% \begin{macro}{\HyperTOCeditors}
% \begin{macro}{\HyperTOCpreamble}
% These commands are used to write information to the \texttt{.htoc}
% file, which is used in generating a Web Page of the Table of
% Contents. Since they only ever appear preceded by |\protect| in the
% argument to an |\addtocontents| command,  they should never need to
% be expanded by the macros. However, we define them here just in case
% \dots
%    \begin{macrocode}
\newcommand{\HyperTOCEntry}[3]{\relax}
\newcommand{\HyperTOCpart}[1]{\relax}
\newcommand{\HyperTOCmeeting}[3]{\relax}
\newcommand{\HyperTOCeditors}[1]{\relax}
\newcommand{\HyperTOCpreamble}{\relax}
%    \end{macrocode}
% \end{macro}
% \end{macro}
% \end{macro}
% \end{macro}
% \end{macro}
%

% \subsection{Formatting the poster abstracts} 
% These can now be used without using the book option.
%    \begin{macrocode}

\newcommand{\@makeminititle}{%
  \newpage
  \null 
  \vskip 3em%
  \begin{center}%
    {\@title \par}%
  \end{center}
  \vskip 3em%
  }

\newif\if@abstractsads
\@abstractsadsfalse
\newcommand\abstractsADS{\@abstractsadstrue}

\if@book\RequirePackage{multicol}\fi
\newenvironment{abstracts}[1][Abstracts of Contributed Papers]{%
  %% First, make sure we start on an odd page (06 Jun 2000)
  \if@book\clearpage\thispagestyle{empty}\cleardoublepage\fi
  %%
  \renewcommand{\articlename}{ABSTRACTS}
  \title{#1}\shortauthor{Abstracts}\shorttitle{Abstracts}
  \listofauthors{}
  \if@book
  \addcontentsline{toc}{chapter}{\protect\@fiddleauthors{\@TOCtitle}{\rmaaTOCauthors}}%
  \addtocontents{htc}{\protect\HyperTOCEntry{\@rawtitle}{\@rawauthors}{\articlename}}
  \fi
  \pagestyle{rmaa} % this now seems necessary (06 Jun 2000)
  \thispagestyle{rmaafirstpage}
%%  \twocolumn[\@makeminititle]
  \renewcommand{\abstract}[1]{%
    \renewcommand\ADS@abstract{##1}%
    \renewcommand{\@abstract}{##1}\@writeabstract%
  }
  \renewcommand{\title}[2][]{%
    \gdef\ADS@authors{}%
    \gdef\ADS@affil{}%
    \gdef\ADS@references{}%
    \gdef\ADS@keywords{}%
    \renewcommand{\@rmaaArg}{##1}
    \ifx\@rmaaArg\@rmaaEmpty
       \renewcommand{\@TOCaftertitle}{}
    \else
       \renewcommand{\@TOCaftertitle}{ (\@rmaaArg)}
    \fi
    \renewcommand{\@title}{\rmaa@TOCTitleStyle{##2}}%
    \renewcommand{\@TOCtitle}{\rmaa@TOCTitleStyle{##2}\@TOCaftertitle}%
    \renewcommand{\rmaaIndexTitle}{\rmaa@IndexTitleStyle{##2\@TOCaftertitle.}}%
    \renewcommand\ADS@title{##2}%
    \xdef\ADS@firstpage{\@arabic{\c@page}}%
    }
  \renewcommand{\author}[1]{\let\@thanks\@empty\renewcommand\@author{##1} }
  \renewcommand{\footnotetext}[2][XX]{\par\noindent\textsuperscript{##1}~##2}
  \renewcommand{\footnoterule}{\kern-3\p@ \hrule \@width 3em \kern 2.6\p@}
  \if@book\@begofarticlehook\fi
  \if@book\begin{multicols}{2}[\@makeminititle]\else\@makeminititle\fi
  \begingroup
  \setcounter{footnote}{0}
  }% and at the end of the environment....
  {\endgroup\if@book\end{multicols}\listofabstracts\label{RMAALastPage}\clearpage\fi}
%% auxiliary commands used:
%% \@writeabstract does all the actual formatting 
\newcommand{\@writeabstract}{%
  \if@book%
  \addcontentsline{loa}{abstract}{\protect\@fiddleauthors{\@TOCtitle}{\rmaaTOCauthors}}%
  \fi%
  \begin{center}\@title\par\@author\end{center}
  \@abstract
  \par\noindent\footnoterule\nobreak\footnotesize\@thanks
  \par\noindent\hrulefill\goodbreak
  \vskip 1cm \@plus0.5cm
  %% write ADS entry
  \if@book\if@abstractsads\ADS@dumpdata\fi\fi%
  \endgroup % this pair is to make each abstract a self-contained group
  \begingroup
  \setcounter{footnote}{0}
  }

\if@book
%% List of abstracts (stuff written to .loa file)
\newcommand{\listofabstracts}{%
  \begin{multicols}{2}[%
    \vskip 1.5cm \vfill % leave at least 1.5 cm and go to bottom of page if possible
    \centerline{\MakeUpperCase{List of Abstracts}}
    \vskip 1.5cm \@plus0.5cm
    ][6cm] % this makes it start a new page if there is not at least 6cm of space
    \renewcommand{\@fiddleauthors}[2]{##1\hskip1em plus2em minus0.4em\penalty-100  ##2}
    \parindent\z@
%%    \parskip.5\p@ \@plus .7\p@\relax 
    \parskip3\p@ \@plus 1\p@ \@minus 1\p@\relax % allow quite a lot of
                                % leeway here to avoid breaks inside
                                % entries 
    \renewcommand{\\}{\relax} % temporarily turn off explicit line breaks just in case
    \renewcommand{\altaffilmark}[1]{\relax} % ditto with this
    \@starttoc{loa}
  \end{multicols}
  }
\newcommand*{\l@abstract}[2]{%
  \par\penalty-1000
  #1\hskip 1em\penalty0\null\nobreak\hfill#2}
\fi

%    \end{macrocode}
 
% \subsection{Code for typesetting an entire volume of conference
% proceedings} \label{sec:whole-volume}
%    \begin{macrocode}

\if@book
%% environment for including the individual articles
\newcommand{\@fiddleauthors}[2]{\relax}
\renewcommand{\theequation}{\@arabic\c@equation}
\newenvironment{articles}%
{ \begingroup % make it a group so that author definitions do not
              % spill over
  \gdef\ADS@authors{}
  \pagestyle{rmaa}
  \renewenvironment{document}{%
    \@begofarticlehook
    \addtocounter{chapter}{1}%
    \addcontentsline{toc}{chapter}{%
      \protect\@fiddleauthors{\@TOCtitle}{\rmaaTOCauthors}}%
    \addtocontents{htc}{\protect\HyperTOCEntry{\@rawtitle}{\@rawauthors}{\articlename}}
    }{\relax}
  \renewcommand{\documentclass}[2][foo]{\relax} 
  \renewcommand{\usepackage}[2][foo]{%
    \def\@packagename{##2}\def\@rmaacitename{rmaacite}%
    \ifx\@packagename\@rmaacitename
    \@rmcitetrue\ClassInfo{rmaa}{This paper uses rmaacite}
    \makeatletter%%
%% This is file `rmaacite.sty',
%% generated with the docstrip utility.
%%
%% The original source files were:
%%
%% rmaa.dtx  (with options: `rmaacite')
%% 
%% PLEASE DO NOT EDIT THIS FILE! It is generated automatically from
%% the file rmaa.dtx using the accompanying Makefile and rmaa.ins
%% Any hand-edited changes will be lost
%% 
%% A LaTeX document class for the RMxAA
%% 
%% Author: William Henney (w.henney@crya.unam.mx)
%% 
%% Latest version available from: http://www.crya.unam.mx/rmaa
%% 
%% Copyright (C) 1999-2010 Instituto de Astronomia, UNAM, Mexico
%% 
%% This program is free software; you can redistribute it and/or
%% modify it under the terms of the GNU General Public License
%% as published by the Free Software Foundation; either version 2
%% of the License, or (at your option) any later version.
%% 
%% This program is distributed in the hope that it will be useful,
%% but WITHOUT ANY WARRANTY; without even the implied warranty of
%% MERCHANTABILITY or FITNESS FOR A PARTICULAR PURPOSE.  See the
%% GNU General Public License for more details.
%% 
%% You should have received a copy of the GNU General Public License
%% along with this program; if not, write to the Free Software
%% Foundation, Inc., 59 Temple Place - Suite 330, Boston, MA  02111-1307, USA.






\def\bibitem[#1<#2>]#3{\par\hangindent=1.5em\hangafter=1
  \if@filesw
  {\def\protect##1{\string ##1\space}
    \immediate\write\@auxout{\string\lbibcite{#3}{#1}{#2}}}
  \fi}
\@redefinebibitem % what it says
\def\lbibcite#1#2#3{\global\@namedef{b@#1}{#2}
                    \global\@namedef{by@#1}{#3}}

\def\@bcite#1#2{({#1\if@tempswa , #2\fi})}
\def\@pcite#1#2{#1\if@tempswa , #2\fi}
\def\@brcite{\ (}
\def\@ebrcite{)}
\def\@fmttop#1{\@fmtcite[\csname #1\endcsname]}
\def\@citex[#1]#2{%
%%  \if@filesw\immediate\write\@auxout{\string\citation{#2}}\fi
%% WJH added 16 Jan 2002
  \@for\@citeb:=#2\do{%
    \edef\@citeb{\expandafter\@firstofone\@citeb:\articlename}%
    \if@filesw\immediate\write\@auxout{\string\citation{\@citeb}}\fi}%
%% to here
  \def\@citea{}%
  \@cite{\@for\@citeb:=#2\do
    {%
      \edef\@citeb{\expandafter\@firstofone\@citeb:\articlename}%
      \@citea\def\@citea{;\ }\@ifundefined
      {b@\@citeb}{{\bf ?}\@warning {Citation `\@citeb' on page \thepage \space undefined}}%
      {\csname b@\@citeb\endcsname\@citebro\csname by@\@citeb\endcsname\@citebrc}}}%
  {#1}}

\def\cite{\@ifnextchar [{\let\@citebro=\space
                          \let\@citebrc=\relax
                          \let\@cite=\@bcite\@tempswatrue\@citex}
                        {\let\@citebro=\space
                          \let\@citebrc=\relax
                          \let\@cite=\@bcite\@tempswafalse\@citex[]}}
\def\pcite{\@ifnextchar [{\let\@citebro=\space
                          \let\@citebrc=\relax
                          \let\@cite=\@pcite\@tempswatrue\@citex}
                        {\let\@citebro=\space
                          \let\@citebrc=\relax
                          \let\@cite=\@pcite\@tempswafalse\@citex[]}}
\def\scite{\@ifnextchar [{\let\@citebro=\@brcite
                          \let\@citebrc=\@ebrcite
                          \let\@cite=\@pcite\@tempswatrue\@citex}
                        {\let\@citebro=\@brcite
                          \let\@citebrc=\@ebrcite
                          \let\@cite=\@pcite\@tempswafalse\@citex[]}}

\renewenvironment{thebibliography}
  {\section*{REFERENCES}\vspace*{\afterbiblioskip}\small
    \if@book\RereadBiblio\fi
    \sloppy\clubpenalty4000\widowpenalty4000
    \sfcode`\.=1000\relax
    \setlength\parindent{\z@}\setlength\parskip{\z@}}
  {\relax\normalsize\@outputaddressesifnecessary}








\endinput
%%
%% End of file `rmaacite.sty'.
\makeatother % read it in here so that all changes
    \else                                         % are local to paper
    \if@debug\ClassWarning{rmaa}{Package \@packagename\space was requested by \articlename\space}\fi
    \fi} 
  %% Don't allow individual articles to mess with the pages, volume or year
  \renewcommand{\SetFirstPage}[1]{\relax}
  \renewcommand{\SetYear}[1]{\relax}
  \renewcommand{\SetVolume}[1]{\relax}
  \renewcommand{\SetConfTitle}[1]{\relax}
  %% Don't allow individual articles to change these
  \renewcommand{\nonstopmode}{\relax}
  \renewcommand{\scrollmode}{\relax}
  \renewcommand{\outputplates}{\relax}
  %% XXXX need to check if using rmaacite and act accordingly
  %% avoid multiple declarations of the same length
  \let\oldnewlength=\newlength
  \renewcommand{\newlength}[1]{%
    \edef\reserved@a{\expandafter\@gobble\string ##1}%
    \@ifundefined{\reserved@a}{%
      \oldnewlength{##1}}{%
      \if@debug{\ClassInfo{rmaa}{\reserved@a\space already defined: doing nothing}}\fi%
      \relax}}%
  \renewcommand\newcommand{\providecommand}
  }%
{\endgroup\onecolumn}
     

%    \end{macrocode}
% Command to be used in `articles' environment for including an
% article
%    \begin{macrocode}
\newcommand{\article}[2][]{%
  % optional first argument is extra text to put in brackets after title in TOC
  % e.g. ``Invited Talk'' or ``Poster''
  \renewcommand{\articlename}{#2}
  % reset loads of counters
  \setcounter{footnote}{0}
  \setcounter{figure}{0}
  \setcounter{table}{0}
  \setcounter{section}{0}
  \setcounter{equation}{0}
  \setcounter{rmaabib}{0}
  \setcounter{author}{0}
  \setcounter{AUTHOR}{1}
  \gdef\ADS@authors{}%
  \gdef\ADS@affil{}%
  \gdef\ADS@affiliations{}%
  \gdef\ADS@references{}%
  \@firstauthortrue
  \@firstaffiltrue
  \@firstaltaffiltrue
  \@firstkwdtrue
  \@rmcitefalse
  \if@suppressfulladdresses\else\@needaddressestrue\fi
  \renewcommand{\@rmaaArg}{#1}
  \ifx\@rmaaArg\@rmaaEmpty
     \renewcommand{\@TOCaftertitle}{}
  \else
     \renewcommand{\@TOCaftertitle}{ (#1)}
  \fi
%    \end{macrocode}
% In the main journal articles must start on an odd page. 
%    \begin{macrocode}
  \if@journal\clearpage\thispagestyle{empty}\cleardoublepage\fi
%% Now input the file, then run the end-of-article hook (reinstated 28
  %% Sep 2001)
  {\input{\@articledir/\articlename.tex}\@endofarticlehook}}
%    \end{macrocode}
% Command to separate articles into sessions. Write an entry in the
% TOC, makes sure we start on an odd page (added 06 Jun 2000), then do
% a title page if \cs{SessionTitlePages} was called. 
% \changes{3.12}{04 Feb 2003}{Added option to have title page for each
% session. Turn it on with \cs{SessionTitlePages}.}
%    \begin{macrocode}
\newcommand{\session}[1]{
  %% this does nothing except write an entry in the TOC
  \addcontentsline{toc}{part}{\rmaa@TOCPartStyle{#1}}
  \addtocontents{htc}{\protect\HyperTOCpart{#1}}
  %% and also makes sure we start on an odd page (added 06 Jun 2000)
  \clearpage\thispagestyle{empty}\cleardoublepage
  \stepcounter{part}
  \if@sessiontitlepages{%
    \onecolumn
    \thispagestyle{empty}
    \renewcommand{\articlename}{Section\thepart TitlePage}
    \label{RMAAFirstPage}
    \make@session@title{#1}
    \label{RMAALastPage}
    \clearpage\thispagestyle{empty}\cleardoublepage
  }\fi
  }
\newcommand{\make@session@title}[1]{%
  {%
    \setlength\parindent{0pt}
    \vspace*{\stretch{1}}
    \rule{\linewidth}{1mm}
    \begin{flushright}\Huge\bfseries\boldmath
      #1
    \end{flushright}
    \rule{\linewidth}{1mm}
    \vspace*{\stretch{2}}
  }
}
%    \end{macrocode}
% Command to turn on separate title pages for each session. 
%    \begin{macrocode}
\newcommand{\SessionTitlePages}{\@sessiontitlepagestrue}
\newcommand{\BoldSessionTitles}{%
  \renewcommand{\rmaa@TOCPartStyle}[1]{\textbf{\rmaa@TOCTitleStyle{##1}}}
  }
\newcommand{\ArticleDir}[1]{\renewcommand{\@articledir}{#1/}\graphicspath{{1/}}}
\newcommand{\@articledir}{.}
%\newcommand{\ps@rmaafront}{\ps@rmaa}% page style for front matter (same as ever, actually)
\newcommand{\ps@rmaafront}{% page style for front matter:  page # in each foot
  \renewcommand\@oddfoot{\hfil{\thepage}}%
  \renewcommand\@evenfoot{{\thepage}\hfil}%
  % shorttitle in both heads
  \renewcommand\@evenhead{\hfil\@shorttitle\hfil}%
  \renewcommand\@oddhead{\@evenhead}
  }
\newcommand{\ps@rmaafrontfirstpage}{% page style for first page of front matter
  % page # in each foot, on outside edge
  \renewcommand\@oddfoot{\hfil\thepage}%
  \renewcommand\@evenfoot{\thepage\hfil}%
  % nothing in the even heads
  \renewcommand\@evenhead{\relax}%
  \renewcommand\@oddhead{\relax}%
  }
%% slight modification to the standard \frontmatter and \mainmatter
%% commands, just to set the \pagestyle
\renewcommand\frontmatter{%
  \latexonecolumn
  \cleardoublepage
  \@frontmattertrue
  \@mainmatterfalse
  \pagenumbering{roman}
  \pagestyle{rmaafront}
  \addtocounter{page}{2} % leave space for the title page
  }
\renewcommand\mainmatter{%
  \pagestyle{empty}
  \if@rmtwocol\twocolumn\fi
  \cleardoublepage
  \@frontmatterfalse
  \@mainmattertrue
  \pagenumbering{arabic}
  \pagestyle{rmaa}
  }
\renewcommand\backmatter{%
  \clearpage\pagestyle{empty}\cleardoublepage
  \@mainmatterfalse
  %% put extra space in the TOC
  \addtocontents{toc}{\bigskip}
  }
%    \end{macrocode}
%
% \subsection{Front sections}\label{sec:front-sections}
% \begin{macro}{\notitle}
% \begin{macro}{\shortname}
% \begin{macro}{frontsection}
% The environment \texttt{frontsection} is for doing a section in the
% front matter (preface, list of participants, or whatever). It has 
% one optional argument:
% \begin{enumerate}
% \item[\#1] Short title (no embedded spaces) for
% use in making PS files. E.g., ``list\_of\_participants''.   
% \end{enumerate}
% and two mandatory arguments: 
% \begin{enumerate}
% \item[\#2] Title of section (e.g., `Preface')
% \item[\#3] Commands to be evaluated before processing the contents
% of the environment. This is mainly useful for putting in the things
% that would go in the preamble of a normal article  (e.g.,
% |\listofauthors|, |\indexauthor|, |\notitle|) 
% \end{enumerate}
%    \begin{macrocode}
\newcommand{\notitle}{\@titlevisiblefalse} \newcommand{\shortname}{}
\newenvironment{frontsection}[3][]{%
  \begingroup
  \renewcommand{\shortname}{#1}
  \ifx\shortname\@rmaaEmpty
     \renewcommand{\articlename}{#2}
  \else
     \renewcommand{\articlename}{\shortname}
  \fi
  \@begofarticlehook
  \@titlevisibletrue % default
  \thispagestyle{rmaafrontfirstpage}
  \listofauthors{} % blank out this in case it is not specified
  \title{{#2}}\shorttitle{#2}\shortauthor{#2}
  #3 % evaluate the command given as second mandatory arg
  \addcontentsline{toc}{chapter}{\protect\@fiddleauthors{\@TOCtitle}{\rmaaTOCauthors}}%
  \addtocontents{htc}{\protect\HyperTOCEntry{\@rawtitle}{\@rawauthors}{\articlename}}
  \if@titlevisible
     \vspace*{1cm}\centerline{\rmaa@FrontSecStyle{\@title}}\vspace*{1cm}
  \fi
  \setlength{\parskip}{3ex}
  }{\frontlabel{RMAALastPage}\clearpage\endgroup}
%    \end{macrocode}
% \end{macro}
% \end{macro}
% \end{macro}
%    \begin{macrocode}
%% Definition of the command \frontlabel, which rights out something like `F4' to
%% the .aux file, instead of `\textlatin  {i}'. This is so it is easy to post-process
%% the file for making PS files of individual sections
\newcommand{\frontlabel}[1]{%
  \@bsphack
  \protected@write\@auxout{}%
     {\string\newlabel{\articlename:#1}{{}{F\@arabic{\c@page}}}}
  \@esphack
  }
%% User command to say X, Y, Z,\\ Editors of the Proceedings
\newcommand{\EditorSignature}{\vspace*{\parskip}\rightline{\@editors}%
  \rightline{Editors of the proceedings}}
%% Fiddle with \thesection, etc. so that they don't include the chapter number
%% Basically go back to the `article' definitions of these
\renewcommand{\thesection}{\@arabic\c@section}
\renewcommand{\thefigure}{\@arabic\c@figure}
\renewcommand{\thetable}{\@arabic\c@table}

\newcommand{\@oldargs}{}
\newcommand{\@newargs}{}
%% add TOC entry for each paper
%\renewcommand{\title}[1]%
%{\renewcommand{\@title}{\rmaa@TitleStyle{#1}}%
%\addtocounter{chapter}{1}\addcontentsline{toc}{chapter}{#1}%
%}

%% Formatting of the table of contents
%% dimensions of various parts
\newlength{\TOCtextwidth}\setlength{\TOCtextwidth}{12.65cm}
\newlength{\TOCpagenumwidth}\setlength{\TOCpagenumwidth}{1.3cm}
\newlength{\TOCentrywidth}\setlength{\TOCentrywidth}{\TOCtextwidth}
\addtolength{\TOCentrywidth}{-\TOCpagenumwidth} 
\newlength{\TOCextramargin}\setlength{\TOCextramargin}{1.25cm}
\newcommand{\TOCbox}[1]{\hspace{\TOCextramargin}\parbox{\TOCtextwidth}{#1}\par}
\newlength{\TOCextramarginx}\setlength{\TOCextramarginx}{\TOCextramargin}\addtolength{\TOCextramarginx}{4.3cm}
\newlength{\TOCtextwidthx}\setlength{\TOCtextwidthx}{\TOCtextwidth}\addtolength{\TOCtextwidthx}{-4.3cm}
\newcommand{\TOCboxx}[1]{\hspace{\TOCextramarginx}\parbox{\TOCtextwidthx}{#1}\par}

\newcommand{\RMAAbannerimage}{rmsc}
\renewcommand\tableofcontents{%
  \renewcommand{\articlename}{CONTENTS}
  \begingroup
  \@begofarticlehook
  \if@twocolumn
     \@restonecoltrue\onecolumn
  \else
     \@restonecolfalse
  \fi
  %% 9pt for journal style
  \if@proceedings\else\small\fi
  \if@newstyle\small\fi %% same for new style
  \if@proceedings\thispagestyle{rmaafrontfirstpage}\else\thispagestyle{empty}\fi
  \shortauthor{Contents}\shorttitle{Contents}
  \vspace*{18mm}
  \if@printbanner
  %% put in the banner header (need to get better image of this)
     \TOCbox{\hspace*{-0.15cm}\includegraphics[width=0.9\TOCtextwidth]{\RMAAbannerimage}}
     \vspace*{8mm}
  %% volume and publication date (should really be in funny font)
     \TOCboxx{\@volumen\\[2mm] \@mes\ \@anyo}
  \else
  %% put in the banner header (need to get better image of this)
     \TOCbox{\hspace*{-0.15cm}\includegraphics[width=0.9\TOCtextwidth]{\RMAAbannerimage}}
     \vspace*{12.5mm}
  %% volume and publication date (should really be in funny font)
     \TOCboxx{~\\[2mm] ~}
  \fi
  \vspace*{16mm}
  %% This is the command that fiddles with the TOC entry so that the
  %% title is flushleft but the authors are flushright
  \renewcommand{\@fiddleauthors}[2]{%
    % first put the paper title
    \renewcommand{\\}{\relax} % temporarily turn off explicit line breaks just in case
    \renewcommand{\altaffilmark}[1]{\relax} % ditto with this
    \renewcommand{\textsuperscript}[1]{\relax} % ditto with this 
    ##1
    % now save the author list in a box
    \sbox\@tempboxa{##2}%
    % check that box is shorter than the total width available for the entry
    \ifdim \wd\@tempboxa <\TOCentrywidth
       % if so, then put it right here, flush against right margin
       \penalty10000\hskip1em plus10em\penalty0\hbox{}\nobreak\hfill\box\@tempboxa
       % Note that we force a space of at least 1em between title and authorlist, 
       % stretchable to 11em since we may have very short titles but still with an 
       % authorlist too long to go on the same line.  The only linebreak allowed is 
       % directly before the authorlist.
    \else
       % if it is too long, we set the original authorlist text in its own parbox
       \par\vskip1ex\parbox[b]{\TOCentrywidth}{\raggedleft ##2}
    \fi
    }
  \@starttoc{toc}%
  \if@restonecol\twocolumn\fi
  \frontlabel{RMAALastPage}
  \clearpage
  \thispagestyle{empty}
  \cleardoublepage % make sure we start on an odd page
  \endgroup

  %%% ALSO, do the HTML TOC
  %% The file that will be used to make the HTML TOC
  \@starttoc{htc}
  %% All the stuff in this file does nothing in the book style
  \addtocontents{htc}{\protect\HyperTOCmeeting{\@conftitle}{\@reunionplace}{\@reuniondate}}
  \addtocontents{htc}{\protect\HyperTOCeditors{\@editors}}
  \addtocontents{htc}{\protect\HyperTOCpreamble}
  \addtocontents{htc}{\protect\HyperTOCEntry{Table of Contents}{}{\articlename}}

  %% and the ADS file
  \ADS@openfile
  }
%% This is now much simpler
\renewcommand\l@chapter[2]{%TOC entry for articles
  \ifnum \c@tocdepth >\m@ne
    \vskip 1.0em \@plus\p@
    \TOCbox{%
      \parbox[b]{\TOCentrywidth}{#1}\parbox[b]{\TOCpagenumwidth}{\hfill#2}
    }
  \fi} 
\renewcommand\l@part[2]{%TOC entry for sessions: centered, no page number
  \ifnum \c@tocdepth >\m@ne
    \vskip 1.5em \@plus1.5em \goodbreak
    \TOCbox{%
      \parbox[b]{\TOCentrywidth}{\hfil#1\hfil}\hfill
    }
  \fi\nobreak} 

%% Formatting theindex
%% strongly encourage page breaks between index entries
\renewcommand{\@idxitem}{\par\penalty-1000\hangindent.45cm} 
\renewcommand{\indexspace}{\par\penalty-1000\vskip 10\p@ \@plus15\p@ \@minus3\p@\relax}
\renewenvironment{theindex}{%
  %% first, make sure we start on an odd page (06 Jun 2000)
  \clearpage\thispagestyle{empty}\cleardoublepage
  %% this looks a bit like an article but is two columns and without abstract, etc
  \renewcommand{\articlename}{INDEX}
  \title{Author Index}\shortauthor{Author Index}\shorttitle{Author Index}
  \listofauthors{}
  \addcontentsline{toc}{chapter}{\protect\@fiddleauthors{\@TOCtitle}{\rmaaTOCauthors}}%
  \addtocontents{htc}{\protect\HyperTOCEntry{\@rawtitle}{\@rawauthors}{\articlename}}
  \pagestyle{rmaa}
  \thispagestyle{rmaafirstpage}
  \@begofarticlehook
  %% 9pt for journal style
  \if@proceedings\else\small\fi
  \if@newstyle\small\fi
  \twocolumn[\@makeminititle]
  %% vertical spacing (see also definition of \@idxitem and \indexspace
  \parskip\z@ \@plus .6\p@\relax % allow quite a lot of leaway here to avoid breaks inside entries
  %% horizontal spacing, breaking and glue parameters (see also the stuff in the definition 
  %% of \indexauthor
  \parindent\z@ 
  \exhyphenpenalty10000 % don't break on explicit hyphens, it always looks bad
  \rightskip=0pt plus.2em % allow a tiny bit of leeway at the ends of lines
  \spaceskip.3333em plus.5em minus.15em % a lot of leeway between words
  \xspaceskip.5em plus.5em minus.4em % this seems to have no effect
  \let\item\@idxitem
  \renewcommand{\\}{\relax} % temporarily turn off explicit line breaks just in case
  \renewcommand{\altaffilmark}[1]{\relax} % ditto with this
  \renewcommand{\textsuperscript}[1]{\relax} % ditto with this
  }{\label{RMAALastPage}\clearpage} % clearpage necessary so that we know about \shortauthor



%% end definitions for book styles
\fi

%    \end{macrocode}
% In order to do away with plates, we sometimes want a figure to
% appear in color only in the online version. The following command
% can bve used in a figure caption to typeset a note to this effect. 
%    \begin{macrocode}
\newcommand\ColorOnlineOnly{%
  \textsc{Note: this figure is available in color in the electronic
    version of this article, obtainable from}
  \texttt{http://\discretionary{}{}{}www.\discretionary{}{}{}astroscu.\discretionary{}{}{}unam.\discretionary{}{}{}mx/\discretionary{}{}{}\~{}rmaa/}.}
%% Plates
%% Define 2 counters: plate is the number in the whole document, plateloop is used 
%% internally by \outputplates
\newcounter{plate}
\newcounter{plateloop}
%% Currently limited to 15 plates in document
%% when we do a label, we need to write to \newlabel entries to the .aux file, one 
%% that says RMAAPlateX so that it can be referred to in \ouputplates, and one that
%% says articlename:RMAAPlateX so that the dvips program can find it and asociate it
%% with the right article
\newcommand{\@platelabeli}{\oldlabel{RMAAPlate1}\label{RMAAPlate1}}
\newcommand{\@platelabelii}{\oldlabel{RMAAPlate2}\label{RMAAPlate2}}
\newcommand{\@platelabeliii}{\oldlabel{RMAAPlate3}\label{RMAAPlate3}}
\newcommand{\@platelabeliv}{\oldlabel{RMAAPlate4}\label{RMAAPlate4}}
\newcommand{\@platelabelv}{\oldlabel{RMAAPlate5}\label{RMAAPlate5}}
\newcommand{\@platelabelvi}{\oldlabel{RMAAPlate6}\label{RMAAPlate6}}
\newcommand{\@platelabelvii}{\oldlabel{RMAAPlate7}\label{RMAAPlate7}}
\newcommand{\@platelabelviii}{\oldlabel{RMAAPlate8}\label{RMAAPlate8}}
\newcommand{\@platelabelix}{\oldlabel{RMAAPlate9}\label{RMAAPlate9}}
\newcommand{\@platelabelx}{\oldlabel{RMAAPlate10}\label{RMAAPlate10}}
\newcommand{\@platelabelxi}{\oldlabel{RMAAPlate11}\label{RMAAPlate11}}
\newcommand{\@platelabelxii}{\oldlabel{RMAAPlate12}\label{RMAAPlate12}}
\newcommand{\@platelabelxiii}{\oldlabel{RMAAPlate13}\label{RMAAPlate13}}
\newcommand{\@platelabelxiv}{\oldlabel{RMAAPlate14}\label{RMAAPlate14}}
\newcommand{\@platelabelxv}{\oldlabel{RMAAPlate15}\label{RMAAPlate15}}
\newcommand{\@outputplatelabeli}{\oldlabel{RMAAOutputPlate1}}
\newcommand{\@outputplatelabelii}{\oldlabel{RMAAOutputPlate2}}
\newcommand{\@outputplatelabeliii}{\oldlabel{RMAAOutputPlate3}}
\newcommand{\@outputplatelabeliv}{\oldlabel{RMAAOutputPlate4}}
\newcommand{\@outputplatelabelv}{\oldlabel{RMAAOutputPlate5}}
\newcommand{\@outputplatelabelvi}{\oldlabel{RMAAOutputPlate6}}
\newcommand{\@outputplatelabelvii}{\oldlabel{RMAAOutputPlate7}}
\newcommand{\@outputplatelabelviii}{\oldlabel{RMAAOutputPlate8}}
\newcommand{\@outputplatelabelix}{\oldlabel{RMAAOutputPlate9}}
\newcommand{\@outputplatelabelx}{\oldlabel{RMAAOutputPlate10}}
\newcommand{\@outputplatelabelxi}{\oldlabel{RMAAOutputPlate11}}
\newcommand{\@outputplatelabelxii}{\oldlabel{RMAAOutputPlate12}}
\newcommand{\@outputplatelabelxiii}{\oldlabel{RMAAOutputPlate13}}
\newcommand{\@outputplatelabelxiv}{\oldlabel{RMAAOutputPlate14}}
\newcommand{\@outputplatelabelxv}{\oldlabel{RMAAOutputPlate15}}
\newcommand{\@platerefi}{\oldpageref{RMAAPlate1}}
\newcommand{\@platerefii}{\oldpageref{RMAAPlate2}}
\newcommand{\@platerefiii}{\oldpageref{RMAAPlate3}}
\newcommand{\@platerefiv}{\oldpageref{RMAAPlate4}}
\newcommand{\@platerefv}{\oldpageref{RMAAPlate5}}
\newcommand{\@platerefvi}{\oldpageref{RMAAPlate6}}
\newcommand{\@platerefvii}{\oldpageref{RMAAPlate7}}
\newcommand{\@platerefviii}{\oldpageref{RMAAPlate8}}
\newcommand{\@platerefix}{\oldpageref{RMAAPlate9}}
\newcommand{\@platerefx}{\oldpageref{RMAAPlate10}}
\newcommand{\@platerefxi}{\oldpageref{RMAAPlate11}}
\newcommand{\@platerefxii}{\oldpageref{RMAAPlate12}}
\newcommand{\@platerefxiii}{\oldpageref{RMAAPlate13}}
\newcommand{\@platerefxiv}{\oldpageref{RMAAPlate14}}
\newcommand{\@platerefxv}{\oldpageref{RMAAPlate15}}
\newcounter{figplatei}
\newcounter{figplateii}
\newcounter{figplateiii}
\newcounter{figplateiv}
\newcounter{figplatev}
\newcounter{figplatevi}
\newcounter{figplatevii}
\newcounter{figplateviii}
\newcounter{figplateix}
\newcounter{figplatex}
\newcounter{figplatexi}
\newcounter{figplatexii}
\newcounter{figplatexiii}
\newcounter{figplatexiv}
\newcounter{figplatexv}
%% Printed output should use the counter `plate'
\renewcommand{\theplate}{\@arabic\c@plate}
\renewcommand{\theplateloop}{\@arabic\c@plateloop}
%% Command to output all the saved plates in this paper
\newcommand{\outputplates}{%
  \setcounter{plateloop}{0}
  \if@book\else\@endofarticlehook\fi % not entirely sure why, but this
                                     % seems to be needed
  \clearpage % first output the last page of text
  % now set pagestyle for blank even pages and make sure we start on an odd page
  \pagestyle{empty}
  \cleardoublepage
  \onecolumn
  \parindent\z@
  % first put in an extra sheet
  % the main purpose of this is so there is always a blank even page before 
  % each plate for filling in the PS file if the article has an odd number
  % of pages (see splitbook.sh)
  \null\vfil\centerline{PLATES}\vfil\clearpage\pagestyle{rmaaplate}\cleardoublepage
  % output each plate in turn
  \@whilenum\value{plateloop}<\value{plate}\do{%
    \addtocounter{plateloop}{1}
    % restore the figure counter from \@figcounti, \@figcountii, etc
    \setcounter{figure}{\value{figplate\roman{plateloop}}}
    % spew out the plate (in its own group for safety)
    {\csname @plate\roman{plateloop}\endcsname\par}
    }
  \pagestyle{rmaa}
  % reset the page number as if nothing had happened
  \addtocounter{page}{-\value{plate}}
  \addtocounter{page}{-\value{plate}}
  }
\newcommand\figcaption{%
  \def\@captype{figure}
  \renewcommand{\@makecaption}{\@makefigcap}
  \caption}
\newcommand\tabcaption{%
  \def\@captype{table}
  \renewcommand{\@makecaption}{\@maketabcap}
  \caption}

\newcommand{\plate}[2][1]{%
  %% User command to make a plate page: 
  %% Mandatory argument is what goes on the plate - this should not 
  %% include any floating environments (e.g. figure). Instead, \figcaption 
  %% should be used to do the captions. 
  %% Optional argument is the number of figures (i.e. captions) on the plate
  %% page (default is 1). This is necessary for incrementing the figure counter. 
  %%
  %% Note that we have to use a \gdef rather than an \xdef when we 
  %% save the plate output in a command. This means that in the book styles when it 
  %% comes to spewing out the plate, we will have forgotten about any user macros 
  %% that may have been defined. This can be got round by repeating the macros inside the plate.
  %% (Presumably, one would use a \newcommand rather than \renewcommand if they were new.)
  % increment global and local counters
  \addtocounter{plate}{1}
  % write appropriate label to .aux file
  \csname @platelabel\roman{plate}\endcsname
  % define a global macro called \@platei, \@plateii, \@plateiii, \@plateiv, ...
  \expandafter\gdef\csname @plate\roman{plate}\endcsname{%
    \csname @outputplatelabel\roman{plateloop}\endcsname
    \null\vfil #2 \vfil
    \cleardoublepage
    }
  % save the current shortauthor and shorttitle
  % need to be \xdef to take current values of \@shortauthor, etc 
  \expandafter\xdef\csname @shorta\roman{plate}\endcsname{\@shortauthor}
  \expandafter\xdef\csname @shortt\roman{plate}\endcsname{\@shorttitle}
  % save the current value of the figure counter for later use in \outputplates
  \setcounter{figplate\roman{plate}}{\value{figure}}
  % increment the figure counter according to the optional argument
  \addtocounter{figure}{#1}
  }
%% page style for plate pages 
\newcommand{\ps@rmaaplate}{%
  % Plate X (see page Y) in the foot
  \renewcommand\@oddfoot{\hfil\csname @shorta\roman{plateloop}\endcsname\ 
    (see page~\csname @plateref\roman{plateloop}\endcsname)\hfil}%
  \renewcommand\@evenfoot{\relax}%
  % nothing in even heads (blank pages)
  \renewcommand\@evenhead{\relax}
  % short title and page num in the odd heads
  \renewcommand\@oddhead{\hfil\csname @shortt\roman{plateloop}\endcsname%
    \hfil\llap{PLATE~\theplateloop}}
  }





%% Things for the User Guide
\if@guide
%% stolen from doc.sty
\def\MakeShortVerb#1{%
  \expandafter\ifx\csname cc\string#1\endcsname\relax
  \@shortvrbinfo{Made }{#1}%
  \add@special{#1}%
  \expandafter
  \xdef\csname cc\string#1\endcsname{\the\catcode`#1}%
  \begingroup
  \catcode`\~\active  \lccode`\~`#1%
  \lowercase{%
    \global\expandafter\let
    \csname ac\string#1\endcsname~%
    \gdef~{\verb~}}%
  \endgroup
  \global\catcode`#1\active
  \else
  \@shortvrbinfo\@empty{#1 already}%
  \fi}
\def\DeleteShortVerb#1{%
  \expandafter\ifx\csname cc\string#1\endcsname\relax
  \else
  \@shortvrbinfo{Deleted }{#1 as}%
  \rem@special{#1}%
  \global\catcode`#1\csname cc\string#1\endcsname
  \global \expandafter\let \csname cc\string#1\endcsname \relax
  \ifnum\catcode`#1=\active
  \begingroup
  \catcode`\~\active   \lccode`\~`#1%
  \lowercase{%
    \global\expandafter\let\expandafter~%
    \csname ac\string#1\endcsname}%
  \endgroup \fi \fi}
\def\@shortvrbinfo#1#2{%
  \PackageInfo{doc}{%
    #1\expandafter\@gobble\string#2 a short reference
    for \string\verb}}
\def\add@special#1{%
  \rem@special{#1}%
  \expandafter\gdef\expandafter\dospecials\expandafter
  {\dospecials \do #1}%
  \expandafter\gdef\expandafter\@sanitize\expandafter
  {\@sanitize \@makeother #1}}
\def\rem@special#1{%
  \def\do##1{%
    \ifnum`#1=`##1 \else \noexpand\do\noexpand##1\fi}%
  \xdef\dospecials{\dospecials}%
  \begingroup
  \def\@makeother##1{%
    \ifnum`#1=`##1 \else \noexpand\@makeother\noexpand##1\fi}%
  \xdef\@sanitize{\@sanitize}%
  \endgroup}
%% define | as shortcut for \verbatim
\MakeShortVerb{\|}
%% comand to specify arguments to LaTeX commands
\newcommand{\rmArg}[1]{%
  \ensuremath{\langle}{\scriptsize \textsf{\MakeUpperCase{#1}}}\ensuremath{\rangle}}
%\newcommand{\rmArg}[1]{\ensuremath{\langle}\textsc{#1}\ensuremath{\rangle}}
\providecommand{\BibTeX}{{\rmfamily B\kern-.05em%
    \textsc{i\kern-.025em b}\kern-.08em%
    T\kern-.1667em\lower.7ex\hbox{E}\kern-.125emX}}
%% Environment for code examples
\newenvironment{Code}{%
  \begin{list}{}{\setlength{\leftmargin}{0.05\columnwidth}} 
  \item[]}{\end{list}}
\fi
%</package>
%    \end{macrocode}

% \section{User commands for astronomical articles}
% \label{sec:user-commands}
% Some random definitions. These go both into the main class file and
% also into |adshtml.sty|.
% 
% Definitions of the journal name in different incarnations.  
%    \begin{macrocode}
%<*package|adshtml-style>
\newcommand{\RMAA}{RMxAA}
\newcommand{\RMAAlong}{Revista Mexicana de Astronom\'{\i}a y Astrof\'{\i}sica}
\newcommand{\RMAASC}{RMxAC}
\newcommand{\RMAASClong}{\RMAA\ (Serie de Conferencias)}
%    \end{macrocode}
% \begin{macro}{\ibidrule}
% This puts a horizontal rule with a period after it. It is used when
% the author list of a reference is identical to the preceding one. 
%    \begin{macrocode}
\newcommand{\ibidrule}{\rule[-0.03cm]{4.7em}{.01cm}.\ }
%    \end{macrocode}
% \end{macro}
% \begin{macro}{\ThisVolume}
% References to other articles in the same volume can be coded as,
% e.g., |\bibitem{} Bloggs, J. \ThisVolume{bloggs}|.
%    \begin{macrocode}
\newcommand{\ThisVolume}[1]{%
  \@year, \if@proceedings\RMAASC\else\RMAA\fi, \@volume,
  \if@book\oldpageref{#1:RMAAFirstPage}\else 000\fi\ (this volume)}
%    \end{macrocode}
% \end{macro}
% Other stuff that may be useful in ref lists.
%    \begin{macrocode}
\newcommand{\etal}{et al.\@}
%    \end{macrocode}
% Various new commands from AAS macros
%    \begin{macrocode}
\def\acknowledgments{\vskip 3ex plus .8ex minus .4ex}
\let\acknowledgements=\acknowledgments                  % second spelling
\let\jnl@style=\rm
\def\ref@jnl#1{{\jnl@style#1}}
\DeclareRobustCommand\aj{\ref@jnl{AJ}}                   
\DeclareRobustCommand\araa{\ref@jnl{ARA\&A}}             
\DeclareRobustCommand\apj{\ref@jnl{ApJ}}                 
\DeclareRobustCommand\apjl{\ref@jnl{ApJ}}                
\DeclareRobustCommand\apjs{\ref@jnl{ApJS}}               
\DeclareRobustCommand\ao{\ref@jnl{Appl.~Opt.}}           
\DeclareRobustCommand\apss{\ref@jnl{Ap\&SS}}             
\DeclareRobustCommand\aap{\ref@jnl{A\&A}}                
\DeclareRobustCommand\aapr{\ref@jnl{A\&A~Rev.}}          
\DeclareRobustCommand\aaps{\ref@jnl{A\&AS}}              
\DeclareRobustCommand\azh{\ref@jnl{AZh}}                 
\DeclareRobustCommand\baas{\ref@jnl{BAAS}}               
\DeclareRobustCommand\jrasc{\ref@jnl{JRASC}}             
\DeclareRobustCommand\memras{\ref@jnl{MmRAS}}            
\DeclareRobustCommand\mnras{\ref@jnl{MNRAS}}             
\DeclareRobustCommand\pra{\ref@jnl{Phys.~Rev.~A}}        
\DeclareRobustCommand\prb{\ref@jnl{Phys.~Rev.~B}}        
\DeclareRobustCommand\prc{\ref@jnl{Phys.~Rev.~C}}        
\DeclareRobustCommand\prd{\ref@jnl{Phys.~Rev.~D}}        
\DeclareRobustCommand\pre{\ref@jnl{Phys.~Rev.~E}}        
\DeclareRobustCommand\prl{\ref@jnl{Phys.~Rev.~Lett.}}    
\DeclareRobustCommand\pasp{\ref@jnl{PASP}}               
\DeclareRobustCommand\pasj{\ref@jnl{PASJ}}               
\DeclareRobustCommand\qjras{\ref@jnl{QJRAS}}             
\DeclareRobustCommand\skytel{\ref@jnl{S\&T}}             
\DeclareRobustCommand\solphys{\ref@jnl{Sol.~Phys.}}      
\DeclareRobustCommand\sovast{\ref@jnl{Soviet~Ast.}}      
\DeclareRobustCommand\ssr{\ref@jnl{Space~Sci.~Rev.}}     
\DeclareRobustCommand\zap{\ref@jnl{ZAp}}                 
\DeclareRobustCommand\nat{\ref@jnl{Nature}}              
\DeclareRobustCommand\iaucirc{\ref@jnl{IAU~Circ.}}
\DeclareRobustCommand\aplett{\ref@jnl{Astrophys.~Lett.}}
\DeclareRobustCommand\apspr{\ref@jnl{Astrophys.~Space~Phys.~Res.}}
\DeclareRobustCommand\bain{\ref@jnl{Bull.~Astron.~Inst.~Netherlands}}
\DeclareRobustCommand\fcp{\ref@jnl{Fund.~Cosmic~Phys.}}
\DeclareRobustCommand\gca{\ref@jnl{Geochim.~Cosmochim.~Acta}}
\DeclareRobustCommand\grl{\ref@jnl{Geophys.~Res.~Lett.}}
\DeclareRobustCommand\jcp{\ref@jnl{J.~Chem.~Phys.}}      
\DeclareRobustCommand\jgr{\ref@jnl{J.~Geophys.~Res.}}    
\DeclareRobustCommand\jqsrt{\ref@jnl{J.~Quant.~Spec.~Radiat.~Transf.}}
\DeclareRobustCommand\memsai{\ref@jnl{Mem.~Soc.~Astron.~Italiana}}
\DeclareRobustCommand\nphysa{\ref@jnl{Nucl.~Phys.~A}}
\DeclareRobustCommand\physrep{\ref@jnl{Phys.~Rep.}}
\DeclareRobustCommand\physscr{\ref@jnl{Phys.~Scr}}
\DeclareRobustCommand\planss{\ref@jnl{Planet.~Space~Sci.}}
\DeclareRobustCommand\procspie{\ref@jnl{Proc.~SPIE}}
\let\astap=\aap
\let\apjlett=\apjl
\let\apjsupp=\apjs
\let\applopt=\ao
\def\phn{\phantom{0}}
\def\phd{\phantom{.}}
\def\phs{\phantom{\ensuremath{-}}}
\def\phm#1{\phantom{#1}}
%    \end{macrocode}
% All these are redone as robust and using more modern commands where
% possible. The robustness is so they will not get expanded if used in
% abstracts and titles when these are passed to the ADS file. 
% \changes{3.16}{2003/04/07}{
%  Added \cs{Color} command for doing photometric colors.}
%    \begin{macrocode}
%%\def\sun{\hbox{$\odot$}}
\DeclareRobustCommand\sun{\ensuremath{\odot}}
%%\def\earth{\hbox{$\oplus$}}
\DeclareRobustCommand\earth{\ensuremath{\oplus}}
%%\def\lesssim{\mathrel{\hbox{\rlap{\hbox{\lower4pt\hbox{$\sim$}}}\hbox{$<$}}}}
%%\def\gtrsim{\mathrel{\hbox{\rlap{\hbox{\lower4pt\hbox{$\sim$}}}\hbox{$>$}}}}
\DeclareRobustCommand\lesssim{\mathrel{\hbox{\smaller\smaller\rlap{\hbox{\lower3pt\hbox{$\sim$}}}\raise1.5pt\hbox{$<$}}}}
\DeclareRobustCommand\gtrsim{\mathrel{\hbox{\smaller\smaller\rlap{\hbox{\lower3pt\hbox{$\sim$}}}\raise1.5pt\hbox{$>$}}}}
%%\def\lta{\mathrel{\spose{\lower 3pt\hbox{$\mathchar"218$}}
%%     \raise 2.0pt\hbox{$\mathchar"13C$}}}
%%\def\gta{\mathrel{\spose{\lower 3pt\hbox{$\mathchar"218$}}
%%     \raise 2.0pt\hbox{$\mathchar"13E$}}}
\DeclareRobustCommand\sq{\hbox{\rlap{$\sqcap$}$\sqcup$}}
%%\def\arcdeg{\hbox{$^\circ$}}
%%\def\arcmin{\hbox{$^\prime$}}
%%\def\arcsec{\hbox{$^{\prime\prime}$}}
\DeclareRobustCommand\arcdeg{\ensuremath{^\circ}}
\DeclareRobustCommand\arcmin{\ensuremath{^\prime}}
\DeclareRobustCommand\arcsec{\ensuremath{^{\prime\prime}}}
%%\def\fd{\hbox{$.\!\!^{\rm d}$}}
%%\def\fh{\hbox{$.\!\!^{\rm h}$}}
%%\def\fm{\hbox{$.\!\!^{\rm m}$}}
%%\def\fs{\hbox{$.\!\!^{\rm s}$}}
%%\def\fdg{\hbox{$.\!\!^\circ$}}
%%\def\farcm{\hbox{$.\mkern-4mu^\prime$}}
%%\def\farcs{\hbox{$.\!\!^{\prime\prime}$}}
%%\def\fp{\hbox{$.\!\!^{\scriptscriptstyle\rm p}$}}
\DeclareRobustCommand\fd{\ensuremath{.\!\!^\mathrm{d}}}
\DeclareRobustCommand\fh{\ensuremath{.\!\!^\mathrm{h}}}
\DeclareRobustCommand\fm{\ensuremath{.\!\!^\mathrm{m}}}
\DeclareRobustCommand\fs{\ensuremath{\mkern0.5mu.\mkern-5mu^\mathrm{s}\mkern-1mu}}
\DeclareRobustCommand\fdg{\ensuremath{.\!\!^\circ}}
\DeclareRobustCommand\farcm{\ensuremath{.\mkern-4mu^\prime}}
\DeclareRobustCommand\farcs{\ensuremath{\mkern1mu.\mkern-6mu^{\prime\prime}\mkern-1.5mu}}
%%\DeclareRobustCommand\farcs{\ensuremath{.\!\!^{\prime\prime}}}
\DeclareRobustCommand\fp{\ensuremath{.\!\!^{\scriptscriptstyle\rm p}}}
%%\def\micron{\hbox{$\mu$m}}
\DeclareRobustCommand\micron{\ensuremath{\mu\mathrm{m}}}
\let\la=\lesssim                        
\let\ga=\gtrsim
%%\def\case#1#2{\hbox{$\frac{#1}{#2}$}}
%%\def\slantfrac#1#2{\hbox{$\,^#1\!/_#2$}}
%%\def\onehalf{\slantfrac{1}{2}}
%%\def\onethird{\slantfrac{1}{3}}
%%\def\twothirds{\slantfrac{2}{3}}
%%\def\onequarter{\slantfrac{1}{4}}
%%\def\threequarters{\slantfrac{3}{4}}
\DeclareRobustCommand\case[2]{\ensuremath{\frac{#1}{#2}}}
\DeclareRobustCommand\slantfrac[2]{\ensuremath{\,^#1\!/_#2}}
\DeclareRobustCommand\onehalf{\slantfrac{1}{2}}
\DeclareRobustCommand\onethird{\slantfrac{1}{3}}
\DeclareRobustCommand\twothirds{\slantfrac{2}{3}}
\DeclareRobustCommand\onequarter{\slantfrac{1}{4}}
\DeclareRobustCommand\threequarters{\slantfrac{3}{4}}
%%\def\ubvr{\hbox{$U\!BV\!R$}}            
%%\def\ub{\hbox{$U\!-\!B$}}               
%%\def\bv{\hbox{$B\!-\!V$}}               
%%\def\vr{\hbox{$V\!-\!R$}}               
%%\def\ur{\hbox{$U\!-\!R$}}
%    \end{macrocode}
% New commands for ``photometric'' colors 
%    \begin{macrocode}
\DeclareRobustCommand\ubvr{\ensuremath{U\!BV\!R}}            
\DeclareRobustCommand\Color[2]{%
  \ensuremath{%
    #1\!\ifx#1V\!\fi-\!#2
  }}
\DeclareRobustCommand\ub{\Color{U}{B}}
\DeclareRobustCommand\bv{\Color{B}{V}}
\DeclareRobustCommand\vr{\Color{V}{R}}
\DeclareRobustCommand\ur{\Color{U}{R}}
%% New version of the \ion command, that actually works! 
%    \end{macrocode}
% \begin{macro}{\ion}
% This version of the |\ion| command allows you to write either
% |\ion{N}{4}| or |\ion{N}{iv}| in order to get N~IV. 
%    \begin{macrocode}
\DeclareRobustCommand\ion[2]{%
  \mbox{#1\kern0.2em%
    \smaller\rmfamily%
    \edef\@tempa{\@car#2\@nil}%
    \ifcat1\@tempa%
    \@Roman{#2}%
    \else%
    \uppercase{#2}%
    \fi}}
\newcounter{IonStage}
\DeclareRobustCommand\plainion[2]{\setcounter{IonStage}{#2}#1 \Roman{IonStage}}
%%\def\ion#1#2{\ensuremath{\mathrm{#1}\;}{\protect\small\rm\@Roman{#2}}}
%    \end{macrocode}
% \end{macro}
%    \begin{macrocode}
\DeclareRobustCommand\nodata{$\cdots$}
%    \end{macrocode}
% These are needed by the |\ion| command. Isn't there an easier way of
% doing this?
%    \begin{macrocode}
\let\@small=\small
\let\@normalsize=\normalsize
\let\@footnotesize=\footnotesize
\let\@scriptsize=\scriptsize
\let\@tiny=\tiny
\let\@large=\large
\let\@Large=\Large
\let\@LARGE=\LARGE
\let\@huge=\huge
\let\@Huge=\Huge
\newcommand{\smaller}{\small}
\renewcommand{\normalsize}{\renewcommand{\smaller}{\small}\@normalsize}
\renewcommand{\small}{\renewcommand{\smaller}{\footnotesize}\@small}
\renewcommand{\footnotesize}{\renewcommand{\smaller}{\scriptsize}\@footnotesize}
\renewcommand{\scriptsize}{\renewcommand{\smaller}{\tiny}\@scriptsize}
\renewcommand{\tiny}{\renewcommand{\smaller}{\tiny}\@tiny}
\renewcommand{\large}{\renewcommand{\smaller}{\normalsize}\@large}  
\renewcommand{\Large}{\renewcommand{\smaller}{\large}\@Large}  
\renewcommand{\LARGE}{\renewcommand{\smaller}{\Large}\@LARGE}  
\renewcommand{\huge}{\renewcommand{\smaller}{\LARGE}\@huge}  
\renewcommand{\Huge}{\renewcommand{\smaller}{\huge}\@Huge}  
%</package|adshtml-style>
%    \end{macrocode}


% \section{Other files}

% \subsection{The driver and make files}
% \label{sec:driver_make}
% These are put here for completeness, although they still need to be
% distributed separately since otherwise the user won't be able to
% unpack this file into its constituent parts.
% 
% The |Makefile| has two principal rules: |docs|, and |macros|. The
% |docs| rule produces a PDF and gzipped PS file of the documentation.
% I still haven't perfected the combination of rules required to get
% it to run \LaTeX{} enough times in all circumstances. The |macros|
% rule runs \LaTeX{} on |rmaa.ins| in order to unpack |rmaa.cls| and
% all the other little files (including the |Makefile| and the
% |rmaa.ins| files themselves\dots 
% 
% First set up lists of various types of file. 
%    \begin{macrocode}
%<*makefile>
DOCFILES=rmaa.pdf
MACRO_LTX_FILES=rmaa.cls rmaacite.sty adshtml.sty
SHELLSCRIPTS=makeads.sh
MACRO_OTHER_FILES=adshtml.perl sample.latex2html-init
MAKEFILE=Makefile
MACROFILES=$(MACRO_LTX_FILES) $(SHELLSCRIPTS) $(MACRO_OTHER_FILES)         
SKELETONFILES=ADS.tex 
%    \end{macrocode}
% Install the files read only into the package directory, specified via a path relative to the source directory. \changes{4.0}{2019/11/14}{Reorganize folder layout so that hardly anything is in the top level.  That way, things are clearer on github.}
%    \begin{macrocode}
INSTALL=install 
INSTALL_DIR=../packages
%    \end{macrocode}
% The default rule is to make everything. Installation must be
% explicitly asked for.
%    \begin{macrocode}
all: docs macros

install: install-docs install-macros

%    \end{macrocode}
% Most files are installed read-only to prevent editing. Skeleton
% files must be writeable. Must be careful not to overwrite the
% Makefile in the parent directory. 
%    \begin{macrocode}
install-macros: macros
        $(INSTALL) -m 444 $(MACRO_LTX_FILES) $(MACRO_OTHER_FILES) $(INSTALL_DIR)
        $(INSTALL) -m 555 $(SHELLSCRIPTS) $(INSTALL_DIR)
        $(INSTALL) -m 644 $(SKELETONFILES) $(INSTALL_DIR)

install-docs: docs
        $(INSTALL) rmaa.pdf $(INSTALL_DIR)/rm-fulldocs.pdf

%    \end{macrocode}
% Make the PDF and gzipped PS files. 
%    \begin{macrocode}
docs: $(DOCFILES)

%    \end{macrocode}
% Make everything else if either the |.dtx| or |.ins| files have
% changed. 
%    \begin{macrocode}
macros: rmaa.dtx rmaa.ins
%    \end{macrocode}
% Ensure that all the macro files can be written to then recreate
% them. The |touch| is there in case any of the files don't exist
% yet. 
%    \begin{macrocode}
        touch $(MACROFILES) $(MAKEFILE)
        chmod u+w $(MACROFILES) $(MAKEFILE) rmaa.ins
        latex "%%
%% This is file `rmaa.ins',
%% generated with the docstrip utility.
%%
%% The original source files were:
%%
%% rmaa.dtx  (with options: `ins')
%% 
%% PLEASE DO NOT EDIT THIS FILE! It is generated automatically from
%% the file rmaa.dtx using the accompanying Makefile and rmaa.ins
%% Any hand-edited changes will be lost
%% 
%% A LaTeX document class for the RMxAA
%% 
%% Author: William Henney (w.henney@crya.unam.mx)
%% 
%% Latest version available from: http://www.crya.unam.mx/rmaa
%% 
%% Copyright (C) 1999-2010 Instituto de Astronomia, UNAM, Mexico
%% 
%% This program is free software; you can redistribute it and/or
%% modify it under the terms of the GNU General Public License
%% as published by the Free Software Foundation; either version 2
%% of the License, or (at your option) any later version.
%% 
%% This program is distributed in the hope that it will be useful,
%% but WITHOUT ANY WARRANTY; without even the implied warranty of
%% MERCHANTABILITY or FITNESS FOR A PARTICULAR PURPOSE.  See the
%% GNU General Public License for more details.
%% 
%% You should have received a copy of the GNU General Public License
%% along with this program; if not, write to the Free Software
%% Foundation, Inc., 59 Temple Place - Suite 330, Boston, MA  02111-1307, USA.





%%
%% Installation batch file for the RMAA LaTeX package
%%
\def\batchfile{rmaa.ins}
\input docstrip.tex
\preamble

PLEASE DO NOT EDIT THIS FILE! It is generated automatically from
the file rmaa.dtx using the accompanying Makefile and rmaa.ins
Any hand-edited changes will be lost

A LaTeX document class for the RMxAA

Author: William Henney (w.henney@crya.unam.mx)

Latest version available from: http://www.crya.unam.mx/rmaa

Copyright (C) 1999-2010 Instituto de Astronomia, UNAM, Mexico

This program is free software; you can redistribute it and/or
modify it under the terms of the GNU General Public License
as published by the Free Software Foundation; either version 2
of the License, or (at your option) any later version.

This program is distributed in the hope that it will be useful,
but WITHOUT ANY WARRANTY; without even the implied warranty of
MERCHANTABILITY or FITNESS FOR A PARTICULAR PURPOSE.  See the
GNU General Public License for more details.

You should have received a copy of the GNU General Public License
along with this program; if not, write to the Free Software
Foundation, Inc., 59 Temple Place - Suite 330, Boston, MA  02111-1307, USA.
\endpreamble
\askforoverwritefalse
\generate{%
  \file{rmaa.cls}{\from{rmaa.dtx}{package}}
  \file{rmaa.ins}{\from{rmaa.dtx}{ins}}
  \file{rmaacite.sty}{\from{rmaa.dtx}{rmaacite}}
  \file{adshtml.sty}{\from{rmaa.dtx}{adshtml-style}}
  \file{ADS.tex}{\from{rmaa.dtx}{ads-container}}
  }
\generate{\nopreamble\nopostamble%
  \file{Makefile}{\from{rmaa.dtx}{makefile}}
  \file{adshtml.perl}{\from{rmaa.dtx}{adshtml-perl}}
  \file{makeads.sh}{\from{rmaa.dtx}{ads-shell-script}}
  \file{sample.latex2html-init}{\from{rmaa.dtx}{latex2html-init}}
  }

%% Local Variables:
%% mode: LaTeX
%% mode: Font-lock
%% End:









\endinput
%%
%% End of file `rmaa.ins'.
"
%    \end{macrocode}
% \TeX{} automatically adds the |.tex| extension to extensionless
% files. We fix this and take the opportunity to convert
% back to TABs all that it has lovingly converted to spaces. 
%    \begin{macrocode}
        cp -f $(MAKEFILE) $(MAKEFILE).BCK && unexpand $(MAKEFILE).tex >| $(MAKEFILE) && rm $(MAKEFILE).tex
%    \end{macrocode}
% Make the shell scripts executable and everything unwriteable (so
% they don't get hand-edited by mistake). However, undo this for the
% files that the user is \emph{supposed} to modify. 
%    \begin{macrocode}
        chmod a+x $(SHELLSCRIPTS)
        chmod a-w $(MACROFILES) $(MAKEFILE)
        chmod u+w $(SKELETONFILES)

%    \end{macrocode}
% We now use PDF\LaTeX{} for this.
%    \begin{macrocode}
rmaa.idx: rmaa.dtx
        @echo "REGENERATING .idx FILE"
        pdflatex "\nonstopmode% \def\fileversion{4.0} \def\filedate{2019/11/12}
% \MakeShortVerb{\|}
% \title{A \LaTeX{} document class for the Revista Mexicana de
% Astronom\'\i a y Astrof\'\i sica\thanks{This file has version
% \fileversion{} last revised \filedate{}.}}
% \author{William J. Henney\thanks{w.henney@irya.unam.mx}}
% \date{\filedate}
% \maketitle
% \paragraph{Abstract}
% This document describes the \texttt{rmaa} \LaTeX{}
% document class for the journal ``Revista Mexicana de
% Astronom\'\i a y Astrof\'\i sica''. The macros can be used for 
% \begin{enumerate}
% \item Preparing author manuscripts for the journal (including
% conference proceedings to be published in the ``Serie de
% Conferencias'')
% \item Typesetting an entire volume of the journal by inputting the
% individual article files from a container file. Facilities are
% included for automatic generation of TOC and Author Index. 
% \end{enumerate}
% 
% \tableofcontents
% \section{The documentation driver file}
% This is the part that generates the documentation that you may or
% may not be reading, which is done with the help of the Makefile
% discussed in section~\ref{sec:driver_make}. 
%    \begin{macrocode}
%<*driver>
\documentclass{ltxdoc}
\usepackage{hyperref}
\usepackage{palatino}
\EnableCrossrefs \CodelineIndex \RecordChanges
\setlength{\parindent}{0pt}
\def\MacroFont{\footnotesize\tt}
\begin{document} \DocInput{rmaa.dtx} \PrintIndex \PrintChanges \end{document}
%</driver>
%    \end{macrocode}

% \section{The implementation}
% 
% \changes{3.24}{2007/09/14}{Added user access to the macro version
% and date via \cs{rmVersion} and \cs{rmDate}}
% 
%    \begin{macrocode}
%<*package>
\newcommand{\rmaa@version}{3.26}
\newcommand{\rmaa@date}{2011/08/15}
\newcommand\rmVersion{\rmaa@version}
\newcommand\rmDate{\rmaa@date}
\NeedsTeXFormat{LaTeX2e}[1995/06/01]
%    \end{macrocode}
% \changes{4.0}{2019/11/02}{Change abstract to appear before
% resumen. Change RevMexAA to RMxAA for consistency with ADS.}
%
% \changes{3.22}{2007/03/22}{Allow proceedings articles to have
%   received/accepted dates} \changes{3.21}{2006/05/15}{Changes for
%   Silvia - abstracts are now shipped out to ADS.}
% \changes{3.19}{2006/01/19}{Changes for Silvia - remove adresses from
%   end, new abstract style} \changes{2.2beta}{2001/10/30}{ Fixed bug
%   whereby twocolumn option was not passed through to book
%   class. Finally sorted out the ``label on last page'' issue. }
% \changes{2.1}{2001/01/23}{ Fixed some inconsistencies in
%   \cs{nocite} and \cs{bibitem} uniqification Fixed the definition of
%   \cs{ion} so it will work in math mode and in contexts where the
%   typesize is different from normal (e.g., title, fig captions).
%   Modified the definition of \cs{@caption}, so we don't write any
%   list-of-figures data (we never use it anyway, and it just causes
%   problems with fragile commands in fig captions). }
% \changes{2.0}{2000/11/30}{ Added `newstyle' option, which should be
%   used in all new proceedings volumes (wider abstract, smaller
%   typeface in fig captions, etc.) Added commands for formatting
%   table captions. New demonstation document rmpostsamp.tex for
%   1-page posters. This will be the last release under a PRI
%   government (for a while...)}  \changes{1.9}{2000/06/06}{
%   Corrected position of page number in footer of articles that open
%   on even pages. Ensured that first article in each session opens on
%   an odd page, as do list of abstracts and author index. Added extra
%   space in TOC before the backmatter. Added command for editors,
%   \cs{BoldSessionTitles}, that makes produces bold titles for
%   session in the TOC.}  \changes{1.8}{2000/05/16}{ Fixed bug in
%   last page label with \cs{suppressfulladdresses} }
% \changes{1.7}{2000/05/11}{ Fixed more little things. Page labels
%   now work, even if you have no bibliography. You can output the
%   full addresses by hand with the \cs{outputfulladdresses}
%   command. Editors' Guide updated but still not finished. Revamped
%   splitbook.sh now produces PS file of each article for web page and
%   also an EPS file of each page for sending to the printers.}
% \changes{1.6}{2000/04/15}{ More cleaning up: full stops after
%   titles in index, running headers and footers in frontmatter,
%   fontsize in title of index, etc} \changes{1.5}{2000/04/15}{
%   Eliminated superscripts from titles in TOC and index}
% \changes{1.4}{2000/04/14}{ Added copyright and license info (GPL)}
% \changes{1.3}{2000/03/01}{ Corrected \cs{figcaption} to do the
%   caption formatting right Corrected font size of sec headings in
%   front matter} \changes{1.2}{2000/02/16}{ Added \cs{plate}
%   command} \changes{1.1}{2000/02/14}{ Added manuscript option,
%   corrected bugs} \changes{1.0}{2000/02/08}{ Now does a whole book
%   - lots of stuff tidied up} \changes{0.7}{1999/10/04}{ Increased
%   spacing between sections in single column styles. Suppressed
%   floats at bottom of pages (didn't work though) }
% \changes{0.6}{1999/09/24}{ Tweaked \cs{adjustfinalcols} (still not
%   perfect) Added proceedings option Added onecolumn/twocolumn
%   options (orthogonal to journal/proceedings) }
% \changes{0.5}{1999/09/23}{ Added \cs{fulladdresses} and
%   \cs{adjustfinalcols} (just about working)}
% \changes{0.4}{1999/09/22}{ Minor stylistic changes}
% \changes{0.3}{1999/06/28}{ Better mechanism for alternative
%   affiliations Added `guide' option for user guide}
% \changes{0.2}{1999/05/27}{ All basic functionality present, I
%   think} \changes{0.1}{1999/05/26}{ Just about working - still
%   missing Fig/Tab/Refs} \changes{0.0}{1999/05/21}{ So bored today,
%   decided to rewrite RMAA LaTeX macros from scratch}
%    \begin{macrocode}
\ProvidesClass{rmaa}[\rmaa@date\space v\rmaa@version\space RMxAA article
class]
%    \end{macrocode}
% \subsection{Option Handling}
% \label{sec:options}
% Declarations that are needed in options
% \changes{3.05}{2003/10/22}{Added natbib option}
% \changes{3.06}{2003/01/03}{Natbib option turned on permanently}
% \changes{3.07}{2003/01/11}{Added proof option}
% \changes{3.14}{2003/01/11}{Option to turn off natbib}
% \changes{3.18}{2003/10/22}{trying out doctex mode} 
% \changes{3.23}{2007/09/09}{Manuscript is now on by default.}
% \changes{3.26}{2011/08/15}{New option for review articles}
%    \begin{macrocode}
\newif\if@reviewarticle
\@reviewarticlefalse
\newif\if@proof
\@prooffalse
\newif\if@journal
\@journaltrue
\newif\if@printdates
\@printdatesfalse
\newif\if@proceedings
\@proceedingsfalse
\newif\if@manuscript
\@manuscripttrue
\newif\if@bigfont
\@bigfonttrue
\newif\if@guide
\@guidefalse
\newif\if@preprint
\@preprintfalse
\newif\if@rmtwocol
\@rmtwocoltrue
\newif\if@biblist
\@biblistfalse
\newif\if@book
\@bookfalse
\newif\if@natbib
\@natbibtrue
\newif\if@rmcite
\@rmcitefalse
\newif\if@titlevisible
\@titlevisibletrue
\newif\if@abstractpresent
\@abstractpresentfalse
\newif\if@printbanner
\@printbannerfalse
\newif\if@suppressfulladdresses
\@suppressfulladdressesfalse
\newif\if@newstyle
\@newstylefalse
\newif\if@debug
\@debugfalse
\newif\if@boldkeywords
\@boldkeywordsfalse
%    \end{macrocode}
% The option handling---mainly setting various booleans.
% \changes{2.2beta11}{2001/11/26}{Made twocol and
% newstyle the default for proceedings.}
% \changes{3.23}{2007/09/09}{Manuscript is now 10 pt instead of 12pt}
% \changes{3.23}{2007/09/10}{Full addresses are now automatically
% suppressed for proceedings style.}
% \changes{3.24}{2007/09/14}{New option \texttt{boldkeywords} that
% does the keywords in the old ugly way. Now off by default.}
%    \begin{macrocode}
\DeclareOption{manuscript}{\typeout{Manuscript for referee} 
  \@manuscripttrue\@journalfalse\@rmtwocolfalse
  \PassOptionsToClass{10pt}{article}}
\DeclareOption{bigfont}{\typeout{Big font manuscript requested by Gary
  Ferland} 
  \@manuscripttrue\@journalfalse\@rmtwocolfalse\@bigfonttrue
  \PassOptionsToClass{11pt}{article}}
\DeclareOption{journal}{\typeout{Main journal option}%
  \@journaltrue\@rmtwocoltrue\@printdatestrue\@manuscriptfalse
  \PassOptionsToClass{10pt}{article}}
\DeclareOption{guide}{\typeout{User guide option}%
  \@guidetrue\@journalfalse\@rmtwocoltrue\@manuscriptfalse
  \PassOptionsToClass{10pt}{article}}
\DeclareOption{preprint}{\typeout{Preprint option}%
  \@preprinttrue\@manuscriptfalse\PassOptionsToClass{10pt}{article}}
\DeclareOption{proof}{\typeout{Author proof option}%
  \@prooftrue\@manuscriptfalse}
\DeclareOption{proceedings}{\typeout{Proceedings option}%
  \@proceedingstrue\@journalfalse\@rmtwocoltrue\@newstyletrue
  \@suppressfulladdressestrue
  \@printdatesfalse\@manuscriptfalse\PassOptionsToClass{10pt}{article}} 
\DeclareOption{printdates}{\typeout{Force printing of received/accepted dates}%
  \@printdatestrue} 
\DeclareOption{book}{\typeout{Whole issue option}%
  \@booktrue} 
\DeclareOption{nonatbib}{\typeout{Turn off natbib support (required
    for compatibility with old rmaacite package).}%
  \@natbibfalse} 
\DeclareOption{newstyle}{\typeout{Use new layout style}%
  \@newstyletrue}
\DeclareOption{oldstyle}{\typeout{Use old layout style}%
  \@newstylefalse}
\DeclareOption{debug}{\typeout{Loads of annoying messages}%
  \@debugtrue}
%    \end{macrocode}
% These two need to be declared explicitly so that |\ExecuteOptions|
% actually does something.
%    \begin{macrocode}
\DeclareOption{twocolumn}{\@rmtwocoltrue}
\DeclareOption{onecolumn}{\@rmtwocolfalse}
\DeclareOption{twoside}{\PassOptionsToClass{\CurrentOption}{article}}
\DeclareOption{boldkeywords}{\@boldkeywordstrue}
%    \end{macrocode}
% Bold, upper-case keywords used to be the default. Now they must be
% turned on explicitly with the \texttt{boldkeywords} option. 
%    \begin{macrocode}
\DeclareOption{boldkeywords}{\@boldkeywordstrue}
%    \end{macrocode}
% Pass any other option along to the \texttt{article} class.
%    \begin{macrocode}
\DeclareOption*{\PassOptionsToClass{\CurrentOption}{article}}
%    \end{macrocode}
% Default options are to be two-sided and based on the standard
% \LaTeX{} \texttt{article} class,
% unless we are processing the entire book, in which case we base
% ourselves on the standard \texttt{book} class. Default is to use
% manuscript option unless otherwise specified. 
%    \begin{macrocode}
\ExecuteOptions{twoside,manuscript} \ProcessOptions\relax
\if@book
\if@rmtwocol\LoadClass[twocolumn]{book}\else\LoadClass{book}\fi
\newcommand{\abstract}{\relax}
\newcommand{\@maketitle}{\relax}
\else
\if@rmtwocol\LoadClass[twocolumn]{article}\else\LoadClass{article}\fi
\fi
%    \end{macrocode}
% Attempt to use natbib for citations 02 Jan 2003. 
% \changes{3.05}{2003/01/02}{If the natbib option is used, then set
% things up so that natbib's version of the bibliography is identical
% to the standard one.}
% \changes{3.09}{2003/01/23}{Explicitly require a recent version of
% natbib, so that we get a decent error message if the package is too
% old.} 
% \changes{3.23}{2007/09/09}{Now allow Bibtex to work properly} 
% \changes{3.25}{2009/10/01}{Change \cs{newcommand} tp \cs{def} when setting \cs{bibfont}. This allows it to work with all versions of natbib.} 
% \changes{3.25}{2010/02/06}{Updated copyright year.} 
% \changes{3.25}{2010/02/06}{Updated email and web addresses to \texttt{crya.unam.mx}.} 
%    \begin{macrocode}
\if@natbib
\RequirePackage{natbib}[1999/02/25]
\bibpunct{(}{)}{;}{a}{}{,}
\def\bibfont{\small}
\renewcommand\bibsection{\section*{REFERENCES}}
\setlength\bibsep{0pt}
\setlength\bibhang{1.5em}
\bibliographystyle{rmaa}
\fi
%    \end{macrocode}
% We want the \texttt{graphicx} package for the extended syntax of
% |\includegraphics|. We want the \texttt{ifthen} package for using in
% these macros (although we don't much yet). We explicitly load
% \texttt{babel} since some \LaTeX{} distributions messed up the
% hyphenation if we didn't.  
%    \begin{macrocode}
\RequirePackage{graphicx}
\RequirePackage{ifthen}
\IfFileExists{babel.sty}{%
%   \let\@restoreindent=\@afterindentfalse
  \RequirePackage[english]{babel}
%   \@restoreindent
  %% reinstate the definition of \@roman so as not to break .idx files
  \@ifundefined{org@@roman}{\relax}{\renewcommand{\@roman}[1]{\org@@roman{##1}}}
  }{%
  \typeout{Package babel.sty not found - don't worry about this}}
%    \end{macrocode}
% Try to load the \texttt{dcolumn} and \texttt{booktabs} packages if
% they can be found. These can produce nicer-looking tables. 
%    \begin{macrocode}
\IfFileExists{dcolumn.sty}{%
  \RequirePackage{dcolumn}
  }{%
  \typeout{Package dcolumn.sty not found - you will not be able to
    align columns on decimal points in tables}}
\IfFileExists{booktabs.sty}{%
  \RequirePackage{booktabs}
  }{%
  \typeout{Package booktabs.sty not found - your tables will not look
    as pretty as they ought too.}
  }
\IfFileExists{makeidx.sty}{%
  \RequirePackage{makeidx}
}{%
  \typeout{Package makeidx.sty not found - Author Index may not work.}
  }
%% commented out for moment since requires more modern version of LaTeX
%% \RequirePackage[latin1]{inputenc}
%    \end{macrocode}
% \changes{4.0}{2019/11/13}{Make \texttt{hyperref} and
% \texttt{etoolbox} into hard requirements, since we now use them
% extensively.}  We use hyperref to get internal and external links.
% But I don't like the colored boxes so I turn those off.  We use
% etoolbox for patching and robustifying at the moment.  In the
% future, I want to port a lot of other stuff over to it.  For
% instance, replace the ifthen package and the TeX ifs too.
%
% We target the latest version (1.7) of PDF, rather than the default
% of 1.5.  This reduces the chance of getting warnings from included
% graphics.  This is not strictly related to hyperref, but we put it
% here anyway.
%    \begin{macrocode}
\RequirePackage{etoolbox}
\RequirePackage{hyperref}
\hypersetup{hidelinks}
\pdfminorversion=7
%    \end{macrocode}
% We patch the \cs{@combinedblfloats} following a suggestion by Heiko
% Oberdiek at \url{https://tex.stackexchange.com/a/249743/3843}, which
% avoids problems when hyperlinks are split across
% pages. \changes{4.0}{2019/11/12}{Fix for issue where breaking
% hyperlink across pages gave an error.}
%    \begin{macrocode}
\urlstyle{same}
\newcommand\email{\begingroup \urlstyle{rm}\Url}
\patchcmd\@combinedblfloats{\box\@outputbox}{\unvbox\@outputbox}{}{%
  \errmessage{\noexpand\@combinedblfloats could not be patched}%
}%
%    \end{macrocode}
% \subsection{Other setup stuff}\label{sec:setup}
% Some declarations: 
%    \begin{macrocode}
\newif\if@needaddresses
\@needaddressestrue
\newif\if@frontmatter
\@frontmatterfalse
\newif\if@inpreamble
\@inpreambletrue
\newif\if@firstauthor
\@firstauthortrue
\newif\if@firstaffil
\@firstaffiltrue
\newif\if@firstaltaffil
\@firstaltaffiltrue
\newif\if@firstkwd
\@firstkwdtrue
\newif\if@noexpandADS
\@noexpandADSfalse
\newif\if@sessiontitlepages
\@sessiontitlepagesfalse
%    \end{macrocode}
% \begin{macro}{\@rmaaEmpty}
% \begin{macro}{\@rmaaArg}
% These are used in various commands to check if arg was left blank
%    \begin{macrocode}
\newcommand{\@rmaaEmpty}{}
\newcommand{\@rmaaArg}{}
%    \end{macrocode}  
% \end{macro}
% \end{macro}
%
% \subsection{Compatibility with old \LaTeX{} versions}
% \label{sec:compat}
% This package is based on the 1998/12/01 release of \LaTeX{}.
% Older releases may not define all the commands that we need.
% Luckily, we can use |\providecommand| to define these if necessary. 
% The following commands are not defined in the 1994/06/01 release. 
%
% \begin{macro}{\MakeUpperCase}
% \begin{macro}{\textsuperscript}
% \begin{macro}{\textbackslash}
% User commands:
%    \begin{macrocode}
\providecommand{\MakeUpperCase}[1]{\uppercase{#1}}
\providecommand{\textsuperscript}[1]{\ensuremath{^\textrm{#1}}}
\providecommand{\textbackslash}{\ensuremath{\backslash}}
%    \end{macrocode}  
% \end{macro}
% \end{macro}
% \end{macro}
% \begin{macro}{\protected@xdef}
% \begin{macro}{\restore@protect}
% \begin{macro}{\@unexpandable@protect}
% \begin{macro}{\hb@xt@}
% \LaTeX{} internal commands:
%    \begin{macrocode}
\providecommand{\protected@xdef}{%
        \let\@@protect\protect
        \let\protect\@unexpandable@protect
        \afterassignment\restore@protect
        \xdef}
\providecommand{\restore@protect}{\let\protect\@@protect}
\providecommand{\@unexpandable@protect}{\noexpand\protect\noexpand}
\providecommand{\hb@xt@}{\hbox to}
%    \end{macrocode}
% \end{macro}
% \end{macro}
% \end{macro}
% \end{macro}
%
% \subsection{Page layout}\label{sec:page-layout}
%
% The dimensions of the text were defined by Elizabeth. Text width is
% 169\,mm with equal margins on the two sides: 
%    \begin{macrocode}
\setlength{\textwidth}{169mm}
\setlength{\@tempdima}{\paperwidth}
\addtolength{\@tempdima}{-\textwidth}
\setlength{\oddsidemargin}{0.5\@tempdima}
\addtolength{\oddsidemargin}{-1in}
\setlength{\evensidemargin}{\oddsidemargin}
%    \end{macrocode}
% Text height is 225\,mm, top margin is 14\,mm. 
%    \begin{macrocode}
\setlength{\textheight}{225mm} 
\setlength{\topmargin}{-1in} 
\addtolength{\topmargin}{14mm}
%    \end{macrocode}
% Space between columns in two-column styles, header, footer,
% paragraph indentation:
%    \begin{macrocode}
\setlength{\columnsep}{9.5mm}
\setlength{\headsep}{0.7cm}
\setlength{\footskip}{0.75cm}
\setlength{\parindent}{1.5em}
%    \end{macrocode}
% 
% \begin{macro}{\@set@scalable@length}
% Try out a new scheme for dealing with various spacing dimensions and
% allowing the user to easily modify them. The command
% \begin{quote}
% |\set@scalable@length|\marg{name}\marg{defaultlength}\marg{stretch}\marg{shrink}
% \end{quote}
% does a 
% |\setlength|\marg{\cs{name}}\marg{defaultlength plus stretch minus shrink}
% and also defines the user command
% \textbackslash\meta{name}|Scale|\marg{factor} that can 
% be used to rescale length by \meta{factor}. Additionally, now
% defines the package command
% \textbackslash\meta{name}|@set@scale|\marg{FACTOR} 
% that also rescales the length by \meta{FACTOR}, so that the length
% becomes \meta{factor} $\times$ \meta{FACTOR} $\times$
% \meta{defaultlength}.   
%    \begin{macrocode}
\newcommand\@set@scalable@length[4]{%
  %% arg #1: NAME of length without leading `\'
  %%
  %% arg #2: default VALUE  
  \expandafter\newlength\csname @#1@default\endcsname
  \expandafter\setlength\csname @#1@default\endcsname{#2}%
  \expandafter\newlength\csname @@#1@@default\endcsname
  \expandafter\setlength\csname @@#1@@default\endcsname{#2}%
  %% arg #3: default STRETCH  
  \expandafter\newlength\csname @#1@stretch\endcsname
  \expandafter\setlength\csname @#1@stretch\endcsname{#3}%
  \expandafter\newlength\csname @@#1@@stretch\endcsname
  \expandafter\setlength\csname @@#1@@stretch\endcsname{#3}%
  %% arg #4: default SHRINK  
  \expandafter\newlength\csname @#1@shrink\endcsname
  \expandafter\setlength\csname @#1@shrink\endcsname{#4}%
  \expandafter\newlength\csname @@#1@@shrink\endcsname
  \expandafter\setlength\csname @@#1@@shrink\endcsname{#4}%
  %%
  %% User-defined scale \@NAME@scale is unity by default
  \expandafter\gdef\csname @#1@scale\endcsname{1}%
  %% Package-defined scale \@NAME@SCALE is unity by default
  \expandafter\gdef\csname @#1@SCALE\endcsname{1}%
  %%
  %% internal command \@reset@NAME should expand to 
  %% \setlength\@@NAME@@default{\@NAME@SCALE\@NAME@default}
  %% \setlength\NAME{\@NAME@scale\@@NAME@@default}
  %% in other words: scale x SCALE x default 
  \expandafter\gdef\csname @reset@#1\endcsname{%
    \expandafter\setlength\csname @@#1@@default\endcsname{%
      \csname @#1@SCALE\endcsname \csname @#1@default\endcsname
      }
    \expandafter\setlength\csname @@#1@@stretch\endcsname{%
      \csname @#1@SCALE\endcsname \csname @#1@stretch\endcsname
      }
    \expandafter\setlength\csname @@#1@@shrink\endcsname{%
      \csname @#1@SCALE\endcsname \csname @#1@shrink\endcsname
      }
    \expandafter\setlength\csname #1\endcsname{%
      \csname @#1@scale\endcsname \csname @@#1@@default\endcsname
      \@plus
      \csname @#1@scale\endcsname \csname @@#1@@stretch\endcsname
      \@minus
      \csname @#1@scale\endcsname \csname @@#1@@shrink\endcsname
      }}
  %%
  %% User command \NAMEScale should expand to 
  %% {\renewcommand\@NAME@scale{#1}\@reset@NAME}
  \expandafter\gdef\csname #1Scale\endcsname##1{%
    \expandafter\gdef\csname @#1@scale\endcsname{##1}%
    \csname @reset@#1\endcsname}
  %%
  %% Package command \NAME@set@SCALE should expand to 
  %% {\renewcommand\@NAME@SCALE{#1}\@reset@NAME}
  \expandafter\gdef\csname #1@set@SCALE\endcsname##1{%
    \expandafter\gdef\csname @#1@SCALE\endcsname{##1}%
    \csname @reset@#1\endcsname}
  %%
  %% initialize at default values with \@reset@NAME
  \csname @reset@#1\endcsname
  }
%    \end{macrocode}
% \end{macro}
% \begin{macro}{\parskip}
% First use of this new infrastructure: set interparagraph space to be
% zero but with a small amount of stretchiness.
%    \begin{macrocode}
%   \ifthenelse{\boolean{@manuscript}}{%
%     \@set@scalable@length{parskip}{0.5\baselineskip}{0.3\p@}{0.1\p@}
%   }{
    \@set@scalable@length{parskip}{0\p@}{0.3\p@}{0\p@}
%   }
%    \end{macrocode}
% \end{macro}
% 
% \begin{macro}{\abovesectionskip}
% \changes{4.0}{2019/11/12}{Use negative values for skip so that first paragraph after section headings is not indented.  This is something that I have been trying to fix for years!  It turns out that the anser was in the documentation for the standard command \cs{@startsection.}}
% \begin{macro}{\belowsectionskip}
% \begin{macro}{\aboveparaskip}
% \begin{macro}{\belowparaskip}
% Define some utility rubber lengths for use in
% section~\ref{sec:section-commands}. 
%    \begin{macrocode}
\newlength{\abovesectionskip}
\newlength{\belowsectionskip}
\newlength{\aboveparaskip}
\newlength{\belowparaskip}
%    \end{macrocode}
% The default values for these lengths are those traditionally used
% for the onecolumn style. 
%    \begin{macrocode}
\@set@scalable@length{abovesectionskip}{-3.5ex}{-4ex}{-2ex}
\@set@scalable@length{belowsectionskip}{2.3ex}{0ex}{2ex}
\@set@scalable@length{aboveparaskip}{3.25ex}{2ex}{0.5ex}
\@set@scalable@length{belowparaskip}{-1em}{\z@}{\z@}
%    \end{macrocode}
% \end{macro}\end{macro}\end{macro}\end{macro}
% We also make scalable some of the standard float parameters. These
% are the same defaults as in the standard article classes, but
% slightly more stretchy. For the twocolumn style, they are currently
% halved, see below. Note that I have also lumped |\abovecaptionskip|
% in with these, since it plays a similar role. There is also a
% |\belowcaptionskip| but it is set to zero in the standard classes. 
%    \begin{macrocode}
\@set@scalable@length{floatsep}{12\p@}{4\p@}{4\p@}
\@set@scalable@length{textfloatsep}{20\p@}{4\p@}{8\p@}
\@set@scalable@length{intextsep}{12\p@}{4\p@}{4\p@}
\@set@scalable@length{dblfloatsep}{12\p@}{4\p@}{4\p@}
\@set@scalable@length{dbltextfloatsep}{20\p@}{4\p@}{8\p@}
\@set@scalable@length{abovecaptionskip}{10\p@}{\z@}{\z@}
%    \end{macrocode}
% \begin{macro}{\@before@title@skip}         
% \begin{macro}{\@before@author@skip}        
% \begin{macro}{\@between@author@skip}       
% \begin{macro}{\@before@abstract@skip}      
% \begin{macro}{\@after@abstracthead@skip}   
% \begin{macro}{\@between@abstract@skip}     
% \begin{macro}{\@before@kwd@skip}           
% \begin{macro}{\@before@dates@skip}           
% Spacing parameters used in |\@maketitle|. These are not changed by
% default except by the user, so they don't need to go into
% |\@package@default@lengths|.  
%    \begin{macrocode}
\newlength{\@before@title@skip}         
\newlength{\@before@author@skip}        
\newlength{\@between@author@skip}       
\newlength{\@before@abstract@skip}      
\newlength{\@after@abstracthead@skip}   
\newlength{\@between@abstract@skip}     
\newlength{\@before@kwd@skip}           
\newlength{\@before@dates@skip}           
\@set@scalable@length{@before@title@skip}{2em}{\z@}{\z@}
\@set@scalable@length{@before@author@skip}{1.5em}{\z@}{\z@}
\@set@scalable@length{@between@author@skip}{0.5em}{\z@}{\z@}
\@set@scalable@length{@before@abstract@skip}{3ex}{1ex}{0.5ex}
\@set@scalable@length{@after@abstracthead@skip}{1ex}{\z@}{\z@}
\@set@scalable@length{@between@abstract@skip}{2ex}{1ex}{0.5ex}
\@set@scalable@length{@before@kwd@skip}{1ex}{0.5ex}{0.2ex}
\@set@scalable@length{@before@dates@skip}{2ex}{0.5ex}{0.2ex}
%    \end{macrocode}
% \end{macro}\end{macro}\end{macro}\end{macro}
% \end{macro}\end{macro}\end{macro}\end{macro}
% 
% Spacing above and below displayed equations also wants to be done
% like this, but it is more complicated since the standard article
% classes make these change with the fontsize. Something to do later,
% maybe. 
% 
%    \begin{macrocode}
%    \end{macrocode}
% \begin{macro}{\@package@default@lengths}
% The value actually used for these lengths is controlled by the
% \meta{SCALE} parameter, which depends on the style option that is in
% effect. This is set using the |\@package@default@lengths| command,
% which must be called every time the style changes (e.g., changing
% from onecolumn to twocolumn). 
%    \begin{macrocode}
\newcommand{\@package@default@lengths}{%
  \ifthenelse{\boolean{@rmtwocol}}{%
    %% twocolumn mode
    \abovesectionskip@set@SCALE{0.5}
    \belowsectionskip@set@SCALE{0.5}
    \aboveparaskip@set@SCALE{0.5}
    \belowparaskip@set@SCALE{1}
    \floatsep@set@SCALE{0.5}
    \textfloatsep@set@SCALE{0.5}
    \intextsep@set@SCALE{0.5}
    \dblfloatsep@set@SCALE{0.5}
    \dbltextfloatsep@set@SCALE{0.5}
    }{%
    %% onecolumn mode
    \abovesectionskip@set@SCALE{1}
    \belowsectionskip@set@SCALE{1}
    \aboveparaskip@set@SCALE{1}
    \belowparaskip@set@SCALE{1}
    \floatsep@set@SCALE{1}
    \textfloatsep@set@SCALE{1}
    \intextsep@set@SCALE{1}
    \dblfloatsep@set@SCALE{1}
    \dbltextfloatsep@set@SCALE{1}
    }%
  }
\@package@default@lengths
%    \end{macrocode}
% \end{macro}
% \begin{macro}{\RescaleLengths}
% \begin{macro}{\RescaleSecLengths}
% \begin{macro}{\RescaleFloatLengths}
% \begin{macro}{\RescaleTitleLengths}
% User commands that will fiddle with all the scalable lengths. 
%    \begin{macrocode}
\newcommand{\RescaleLengths}[1]{%
  \RescaleSecLengths{#1}\RescaleFloatLengths{#1}\RescaleTitleLengths{#1}}
\newcommand{\RescaleSecLengths}[1]{%
%% sections
  \abovesectionskipScale{#1}
  \belowsectionskipScale{#1}
  \aboveparaskipScale{#1}
  }
\newcommand{\RescaleFloatLengths}[1]{%
%% floats
  \floatsepScale{#1}
  \textfloatsepScale{#1}
  \intextsepScale{#1}
  \dblfloatsepScale{#1}
  \dbltextfloatsepScale{#1}
  \abovecaptionskipScale{#1}
  }
\newcommand{\RescaleTitleLengths}[1]{%
%% maketitle
  \@before@title@skipScale{#1}           
  \@before@author@skipScale{#1}        
  \@between@author@skipScale{#1}       
  \@before@abstract@skipScale{#1}      
  \@after@abstracthead@skipScale{#1}   
  \@between@abstract@skipScale{#1}     
  \@before@kwd@skipScale{#1}           
  \@before@dates@skipScale{#1}           
  }
%    \end{macrocode}
% \end{macro}\end{macro}\end{macro}\end{macro}
% \changes{2.2beta14}{2002/03/17}{Finished implementation of new
% infrastructure for scalable lengths. Converted the above/below
% section skips to use it. Also, all the floatsep parameters. Also,
% all the vertical lengths that are used by |\@maketitle|. }

% \begin{macro}{\MSextramargin}
% For the \texttt{manuscript} option we have an extra 0.3\,inch
% vertical margin and 1\,inch horizontal margin for marginal notes (by
% hand or with |\marginpar|. 
%
% \changes{3.23}{2007/09/09}{Reduce baselinestretch in
%   \texttt{manuscript} style to 1.2. No longer use raggedright for
%   manuscript.}
%
% \changes{3.24}{2007/09/14}{Further reduction of baselinestretch in
%   \texttt{manuscript} style to 1.1 since we now have wider margins.}
%
% \changes{3.24}{2007/09/14}{Split \texttt{manuscript}'s extra margin
%   into separate horizontal and vertical parts.}
% 
%    \begin{macrocode}
\newlength{\MSextraHmargin}\newlength{\MSextraVmargin}
\setlength{\MSextraHmargin}{1in} 
\setlength{\MSextraVmargin}{0.3in} 
\if@manuscript
   \renewcommand{\baselinestretch}{1.1}
   \addtolength{\textwidth}{-2\MSextraHmargin}
   \addtolength{\oddsidemargin}{\MSextraHmargin}
   \setlength{\evensidemargin}{\oddsidemargin}
   \addtolength{\textheight}{-2\MSextraVmargin}
   \addtolength{\topmargin}{\MSextraVmargin}
%    \raggedright
\fi
%    \end{macrocode}
% \end{macro}
% Float parameters: don't allow bottom floats by default. 
%    \begin{macrocode}
\renewcommand{\fps@figure}{tp}
\renewcommand{\fps@table}{tp}
%    \end{macrocode}
%
% \begin{macro}{\rmaa@AuthorWidthFrac}
% \begin{macro}{\rmaa@AbstractWidthFrac}
% These parameters control what fraction of the total text width can
% be used by the list of authors and the abstract/resumen,
% respectively.  
%    \begin{macrocode}
\newcommand{\rmaa@AuthorWidthFrac}{1.0}
\newcommand{\rmaa@AbstractWidthFrac}{0.75}
%    \end{macrocode}
% Use a wide abstract to save space in the new style.
%    \begin{macrocode}
\if@newstyle\renewcommand{\rmaa@AbstractWidthFrac}{1.0}\fi
%    \end{macrocode}
% \end{macro}
% \end{macro}
%
% Some things to consider in future versions. Trying to avoid too many
% widows, etc. Apparently, the following will work best with
% \cs{flushbottom}. 
%    \begin{macrocode}
%%\widowpenalty=10000
%%\clubpenalty=10000
%    \end{macrocode}
% Also, should perhaps try getting rid of underfull/overfull box warnings. 
% \changes{3.10}{2003/01/24}{Set large values for \cs{hbadness},
% \cs{vbadness}, \cs{hfuzz}, \cs{vfuzz}. To reduce amount of stuff
% that gets written to the terminal every time you compile something.}
%    \begin{macrocode}
\hbadness=10000
\vbadness=10000
\setlength\hfuzz{\maxdimen}
\setlength\vfuzz{\maxdimen}
%    \end{macrocode}
% \subsection{Various text styles for different contexts}
% \label{sec:text-styles} 
% \begin{macro}{\rmaa@AffilStyle}
% \begin{macro}{\rmaa@PagesStyle}
% \begin{macro}{\rmaa@YearStyle}
% \begin{macro}{\rmaa@VolumeStyle}
% \begin{macro}{\rmaa@VolumenStyle}
% \begin{macro}{\rmaa@JournalNameStyle}
% \begin{macro}{\rmaa@ReceivedStyle}
% \begin{macro}{\rmaa@AbstractHeaderStyle}
% \begin{macro}{\rmaa@JournalNameStyle}
% Style commands for the different fields used by |\maketitle| in
% constructing the title page of the article. These have the form
% |\rmaa@XYZStyle{TEXT}|. I'm not sure these were really a good idea,
% but never mind \dots 
%    \begin{macrocode}
\newcommand{\rmaa@AffilStyle}[1]{{\normalsize\rmfamily#1}}
\newcommand{\rmaa@PagesStyle}[1]{{\normalsize\rmfamily#1}}
\newcommand{\rmaa@YearStyle}[1]{{\normalsize\rmfamily#1}}
\newcommand{\rmaa@VolumeStyle}[1]{{\normalsize\rmfamily\bfseries#1}}
\newcommand{\rmaa@VolumenStyle}[1]{{\Large\sffamily\bfseries#1}}
\newcommand{\rmaa@JournalNameStyle}[1]{{\normalsize\itshape#1}}
\newcommand{\rmaa@ReceivedStyle}[1]{{\normalsize\rmfamily\itshape#1}}
\newcommand{\rmaa@AbstractHeaderStyle}[1]{%
  {\normalsize\rmfamily\MakeUpperCase{#1}}}
\newcommand{\rmaa@TitleStyle}[1]{%
  {\large\bfseries\rmfamily\boldmath\MakeUpperCase{#1}\par}}
%    \end{macrocode}
% \end{macro}
% \end{macro}
% \end{macro}
% \end{macro}
% \end{macro}
% \end{macro}
% \end{macro}
% \end{macro}
% \end{macro}
%
% \begin{macro}{\rmaa@FrontSecStyle}
% Headings of ``front sections'' (see section~\ref{sec:front-sections})
%    \begin{macrocode}
\newcommand{\rmaa@FrontSecStyle}[1]{%
  {\normalsize\bfseries\rmfamily\MakeUpperCase{#1}}}
%    \end{macrocode}
% \end{macro}
%
% \begin{macro}{\rmaa@TOCTitleStyle}
% \begin{macro}{\rmaa@TOCPartStyle}
% \begin{macro}{\rmaa@IndexTitleStyle}
% \begin{macro}{\rmaaTOCAuthorStyle}
% \begin{macro}{\rmaaIndexAuthorStyle}
% Commands to set style of TOC and index entries.
% Rewritten to avoid unnecessary font commands in the .toc and .ind files
%    \begin{macrocode}
\newcommand{\rmaa@TOCTitleStyle}[1]{\MakeUpperCase{#1}}
\newcommand{\rmaa@TOCPartStyle}[1]{\rmaa@TOCTitleStyle{#1}}
\newcommand{\rmaa@IndexTitleStyle}[1]{#1}
\newcommand{\rmaaTOCAuthorStyle}[1]{%
  {\itshape#1}}
\newcommand{\rmaaIndexAuthorStyle}[1]{%
  \protect\mbox{\bfseries#1}}
%    \end{macrocode}
% \end{macro}
% \end{macro}
% \end{macro}
% \end{macro}
% \end{macro}
%
% \begin{macro}{\rmaa@KeywordStyle}
% \begin{macro}{\rmaa@KeywordHeaderStyle}
% The style of the keywords depends on whether we are using the
% \texttt{newstyle} option or not. If we are, then they are small. 
%    \begin{macrocode}
\ifthenelse{\boolean{@newstyle}}{
  \newcommand{\rmaa@KeywordStyle}[1]{%
    {\small%
      \if@boldkeywords 
      \bfseries\MakeUpperCase{\raggedright #1}
      \else
      \raggedright #1
      \fi
    }}
  \newcommand{\rmaa@KeywordHeaderStyle}[1]{%
    {\small\itshape#1}}}
{
  \newcommand{\rmaa@KeywordStyle}[1]{%
    {\normalsize%
      \if@boldkeywords\bfseries\fi\raggedright #1}}
  \newcommand{\rmaa@KeywordHeaderStyle}[1]{%
    {\normalsize\itshape#1}}}
%    \end{macrocode}
% \end{macro}
% \end{macro}
%
% \subsection{Page styles}\label{sec:page-styles}
% \begin{macro}{\ps@rmaa}
% The running headers for normal pages:
%    \begin{macrocode}
\newcommand{\ps@rmaa}{%
%    \end{macrocode}
% nothing in either foot; 
%    \begin{macrocode}
  \renewcommand\@oddfoot{\relax}%
  \renewcommand\@evenfoot{\relax}%
%    \end{macrocode}
% page number and short author list in the even heads; 
%    \begin{macrocode}
  \renewcommand\@evenhead{\rlap{\thepage}\hfil\@shortauthor\hfil}%
%    \end{macrocode}
% short title and page number in the odd heads.
%    \begin{macrocode}
  \renewcommand\@oddhead{\hfil\@shorttitle\hfil\llap{\thepage}}
  }
%    \end{macrocode}
% \begin{macro}{\ps@rmaafirstpage}
% The header of the first page of each
% article:
%    \begin{macrocode}
\newcommand{\ps@rmaafirstpage}{%
  \renewcommand\@oddfoot{\hfil\thepage}%
%    \end{macrocode}
% footer corrected (06 Jun 2000) to have page on outside edge;
%    \begin{macrocode}
  \renewcommand\@evenfoot{\thepage\hfil}%
%    \end{macrocode}
% Header has journal name, volume, and page numbers for a regular
% article, or similar for special cases. 
% \changes{4.0}{2019/11/03}{Change to three-line header with copyright info and DOI.}
%    \begin{macrocode}
\if@preprint
\renewcommand\@oddhead{%
  \if@proceedings 
  \textsc{To appear in ``\@conftitle{} (\rmaa@YearStyle{\@year})''}\hfil
  \textit{RMxAC}
  \else
  \textsc{To appear in} \@journalname\hfil
  \fi}\fi%
\if@guide
\renewcommand\@oddhead{\textsc{Documentation} (\rmaa@YearStyle{\@year})\hfil}
\fi%
\if@manuscript
\renewcommand\@oddhead{Manuscript for \@journalname{} (\rmaa@YearStyle{\@year})\hfil}
\fi%
\if@journal
\renewcommand\@oddhead{%
  \parbox{\textwidth}{
    \@journalname, \rmaa@VolumeStyle{\@volume},
    \rmaa@PagesStyle{\@artpages}\
    (\rmaa@YearStyle{\@year})\hfil\\[1pt]
    {\small\textcopyright{} \@year:
      Instituto de Astronom\'\i{}a,
      Universidad Nacional Aut\'onoma de M\'exico}\hfil\\[1pt]
    {\small\@DOI@url}\hfil
  }%
}\fi%
  \renewcommand\@evenhead{\@oddhead}}
%    \end{macrocode}
% By default, use the |\ps@rmaa| page style.
%    \begin{macrocode}
\pagestyle{rmaa}
%    \end{macrocode}
% \end{macro}
% \end{macro}
%
% \subsection{Preamble commands}\label{sec:preamble} 
% Everything you need for setting the title page, TOC and index
% entries, etc.
% \subsubsection{Commands for the title page}
% \label{sec:title-commands}
% \begin{macro}{\resumen}
% \begin{macro}{\@resumen}
% \begin{macro}{\abstract}
% \begin{macro}{\@abstract}
% User commands for setting the resumen and abstract. 
% These should be used in the document header. The internal commands
% |\@resumen| and |\@abstract| are used to hold the actual text
% pending output via the |\maketitle| command. 
%    \begin{macrocode}
\newif\if@setresumen
\newcommand{\resumen}[1]{%
%%  \renewcommand{\@resumen}{{\selectlanguage{spanish}#1}}%
  \renewcommand{\@resumen}{{\ignorespaces#1}}%
  \@setresumentrue}
\renewcommand{\abstract}[1]{%
  \renewcommand{\@abstract}{{\ignorespaces#1}}
  \renewcommand\ADS@abstract{#1}
%% Not a good idea, since then author macros are not expanded and
%% cause errors later. However, by not doing it we are vulnerable to
%% fragile commands.....
  \if@noexpandADS\@onelevel@sanitize{\ADS@abstract}\fi
  \@abstractpresenttrue
  }
\newcommand{\ADSnoexpand}{\@noexpandADStrue}
%    \end{macrocode}
% \begin{macro}{\boldabstract}
% This is the version for doing 1-page conference proceedings. Still
% save the text for ADS but then write it out immediately as a bold
% paragraph. 
%    \begin{macrocode}
\newcommand{\boldabstract}[1]{%
  \gdef\ADS@abstract{#1}%
  \renewcommand{\@abstract}{#1}%
  \if@noexpandADS\@onelevel@sanitize{\ADS@abstract}\fi%
  \noindent\textbf{\boldmath #1}
  }
%    \end{macrocode}
% \end{macro}
% Default value in each case is to print out a warning if not set
%    \begin{macrocode}
  \newcommand{\@resumen}{%
    \@setresumenfalse \ClassWarningNoLine{rmaa}{No Spanish abstract
      (resumen) specified - if you do not speak Spanish, this will be
      translated by the editors}%
    \textbf{Favor de proporcionar un resumen en espa\~nol. If you are
      unable to translate your abstract into Spanish, the editors will
      do it for you.} }
  \newcommand{\@abstract}{\ClassWarningNoLine{rmaa}{No abstract
      specified}}
%    \end{macrocode}
% \end{macro}
% \end{macro}
% \end{macro}
% \end{macro}
%
% \begin{macro}{\title}
% \begin{macro}{\@rawtitle}
% \begin{macro}{\@TOCtitle}
% \begin{macro}{\rmaaIndexTitle}
% Similarly, |\title| is the preamble command to set the title of the
% article. This is used in various places, typeset in different
% styles, hence the use of the internal |\@rawtitle|, |\@title|,
% |\rmaaIndexTitle|, and |\@TOCtitle| commands, each of which uses the
% relevant style from section~\ref{sec:text-styles}
%    \begin{macrocode}
\newcommand{\@rawtitle}{Raw Title}
\newcommand{\@TOCtitle}{\rmaa@TOCTitleStyle{Title for TOC}}
\newcommand{\rmaaIndexTitle}{\rmaa@IndexTitleStyle{Title for index}.}
\newcommand{\@TOCaftertitle}{}
\renewcommand{\title}[1]{%
%%  {\ADS@sanitize\xdef\@rawtitle{#1}}
  \renewcommand\@rawtitle{\noexpand#1}
  \let\ADS@title=\@rawtitle
  \renewcommand{\@title}{\rmaa@TitleStyle{#1}}%
  \renewcommand{\@TOCtitle}{\rmaa@TOCTitleStyle{#1}\@TOCaftertitle}%
  \renewcommand{\rmaaIndexTitle}{\rmaa@IndexTitleStyle{#1\@TOCaftertitle}.}%
  }
%    \end{macrocode}
% \end{macro}
% \end{macro}
% \end{macro}
% \end{macro}
% 
% All the weird and whacky ways of indicating the author
% affiliations. 
% \begin{macro}{\affil}
% First, the way that sets it directly underneath the authors name.
% Note that |\affil| trumps all other affiliations for the purposes of
% getting into the ADS stuff (we use |\gdef|).   
% \dots 
%    \begin{macrocode}
\newcommand{\affil}[1]{%
  \\[1.5ex]\rmaa@AffilStyle{#1}\\[1.5ex]%
  \if@firstaffil\gdef\ADS@affil{#1}\@firstaffilfalse\fi%
  }
\renewcommand{\and}{and\\[1.5ex]}
%    \end{macrocode}
% \begin{macro}{\altaffil}
% Alternative affiliations are footnoted. 
%    \begin{macrocode}
\newcommand{\altaffil}[1]{%
  \if@firstaltaffil\gdef\ADS@affil{#1}\@firstaltaffilfalse\fi%
  \thanks{#1}}
%    \end{macrocode}
% \begin{macro}{\altaffilmark}
% \begin{macro}{\altaffiltext}
% Alternative-alternative mechanism for affiliations---should not be
% mixed with |\altaffil| (unless you know what you are doing) or there
% will be problems with numbering.
%    \begin{macrocode}
\newcommand{\altaffilmark}[1]{\textsuperscript{#1}}
\newcommand{\altaffiltext}[2]{%
  \if@firstaltaffil\gdef\ADS@affil{#2}\@firstaltaffilfalse\fi
  \stepcounter{footnote}%
  \protected@xdef\@thanks{\@thanks
    \protect\footnotetext[#1]{#2}}} 
%    \end{macrocode}
% \end{macro}
% \end{macro}
% \end{macro}
% \end{macro}
%
% \begin{macro}{\keywords}
% \begin{macro}{\@keywords}
% \begin{macro}{\@keywordheader}
% \label{keywords}Sundry minor preamble commands. All
% self-explanatory, I hope. |\keywords| is the old user interface
% (deprecated). New articles should use |\addkeyword| once for each
% keyword. That way the ADS record comes out right too (comma rather
% than em-dash separator).
% \changes{2.2beta12}{2001/12/07}{\cs{addkeyword} New user interface
%   for keywords.} 
% \changes{2.2beta13}{2002/03/11}{\cs{addkeyword} Improved behaviour
%   with .ads file. Use
%   \cs{ADS@sanitize} to change definitions of `~' and \cs{ion}.}
% \changes{3.23}{2007/09/09}{\cs{addkeyword} now signals a compilation
% error if the author tries to add more than 6 keywords.}
% \changes{3.24}{2007/09/14}{kwdcounter is now reset in
% \cs{@begofarticlehook} to prevent error in book style} 
%    \begin{macrocode}
\newcommand{\@kwdsep}{ --- }
\newcommand{\ADS@kwdsep}{, }
\newcounter{kwdcounter}\setcounter{kwdcounter}{0}
\newcommand{\addkeyword}[1]{%
  \addtocounter{kwdcounter}{1}
  \ifthenelse{\thekwdcounter > 6}{
    \ClassError{rmaa}{Too many keywords. A maximum of 6 keywords is allowed.}{}
  }{}
  \if@firstkwd%
  \renewcommand\@keywords{\if@boldkeywords\uppercase{#1}\else #1\fi}
  {\ADS@sanitize\def\ion{\plainion}\gdef\ADS@keywords{#1}}
  \@firstkwdfalse
  \else%
  \edef\@keywords{\@keywords\@kwdsep \if@boldkeywords\uppercase{#1}\else #1\fi}%
  {\ADS@sanitize\def\ion{\plainion}\xdef\ADS@keywords{\ADS@keywords\ADS@kwdsep #1}}%
  \fi}
\newcommand{\ADS@sanitize}{%
  \def\mathrm{}
  \def\@{\ }
  \def~{\ }
  \def\langle{<}
  \def\rangle{>}
  \def\log{log }
  \def\sin{sin }
  \def\cos{cos }
  \def\tan{tan }
  \def\fd{^d.}
  \def\fh{^h.}
  \def\fm{^m.}
  \def\fs{^s.}
  \def\fdg{�.}
  \def\farcm{'.}
  \def\farcs{''.}
  \def\arcdeg{�}
  \def\,{ }
  }
\newcommand{\keywords}[1]{%
  \renewcommand{\ADS@keywords}{#1}
  \renewcommand{\@keywords}{\uppercase{#1}}
  }
\newcommand{\@keywordheader}{\rmaa@KeywordHeaderStyle{Key Words:}}
\newcommand{\@keywords}{%
  \ClassWarningNoLine{rmaa}{No keywords specified}%
  Keywords go here
  }
%    \end{macrocode}
% \begin{macro}{\shorttitle}
% \begin{macro}{\shortauthor}
% \begin{macro}{\@shorttitle}
% \begin{macro}{\@shortauthor}
% Commands for setting the running headers.
%    \begin{macrocode}
\newcommand{\shorttitle}[1]{%
  \renewcommand\@shorttitle{\MakeUpperCase{#1}}}
\newcommand{\shortauthor}[1]{%
  \renewcommand\@shortauthor{\MakeUpperCase{#1}}}
\newcommand{\@shorttitle}{SHORT TITLE GOES HERE}
\newcommand{\@shortauthor}{SHORT AUTHOR GOES HERE}
%    \end{macrocode}
% \end{macro}
% \end{macro}
% \end{macro}
% \end{macro}
% \end{macro}
% \end{macro}
% \end{macro}


% \subsubsection{The full addresses}
% \label{sec:full-addresses}
% \begin{macro}{\fulladdresses}
% \begin{macro}{\@fulladdresses}
% \begin{macro}{\@fulladdheight}
% \begin{macro}{\@fulladddepth}
% \begin{macro}{\@fulladdbox}
% Aaghh, the full addresses! A perennial source of problems. 
%    \begin{macrocode}
\newlength{\@fulladdheight}
\newlength{\@fulladddepth}
\newsavebox{\@fulladdbox}
%% \fulladdresses command modified (11 May 2000) so that we can put
%% \fulladdresses{} in the article to suppress the output of the
%% addresses
\newcommand{\@FinalFigure}{\relax}
\newcommand{\FinalFigure}[1]{\renewcommand{\@FinalFigure}{#1}}
\newcommand{\fulladdresses}[1]{%
  \renewcommand{\@rmaaArg}{#1}
  \ifx\@rmaaArg\@rmaaEmpty
     \renewcommand{\@fulladdresses}{}
  \else
     \renewcommand{\@fulladdresses}{
       \normalsize%
       \ifx\@FinalFigure\relax\else\vfill\@FinalFigure\vfill\fi
%%       \vspace*{1ex}
       %% do addresses as a list (copied from thebibliography environment) 
       \list{\null}{\leftmargin 1.5em\topsep\z@\labelwidth\z@\labelsep\z@\itemindent -1.5em
         \itemsep\z@\parsep\z@\usecounter{enumi}}
       \sloppy\clubpenalty4000\widowpenalty4000
       \sfcode`\.=1000\relax
       #1
       \endlist
       }
  \fi
  }
\newcommand{\@fulladdresses}{Full addresses go here}
%    \end{macrocode}
% \begin{macro}{\suppressfulladdresses}
% New user-level command for suppressing full addresses should you so
% desire. Currently used in the single-page-poster style. 
%    \begin{macrocode}
\newcommand{\suppressfulladdresses}{\@suppressfulladdressestrue\@needaddressesfalse}
%    \end{macrocode}
% \end{macro}
% \end{macro}
% \end{macro}
% \end{macro}
% \end{macro}
% \end{macro}

% \begin{macro}{\printdates}
% \begin{macro}{\noprintdates}
% Command to activate the printing out of the received/accepted
% dates. This can also be done using the \texttt{printdates} option,
% which is the default for \texttt{journal}, but not for \texttt{proceedings}.
%    \begin{macrocode}
\newcommand{\printdates}{\@printdatestrue}
\newcommand{\noprintdates}{\@printdatesfalse}
%    \end{macrocode}
% \end{macro}
% \end{macro}
% \begin{macro}{\@outputaddressesifnecessary}
% Command to output the addresses if we haven't already done so
% This is called by: 
% \begin{enumerate}
% \item The end of |thebibliography| environment
% \item The |\outputfulladdresses| command
% \end{enumerate}
%    \begin{macrocode}
\newcommand{\@outputaddressesifnecessary}{%
  \ClassInfo{rmaa}{Calling @outputaddressesifnecessary}
  \ifthenelse{\boolean{@needaddresses}}{%
    %% If in one column style, or in the first column of a two column
    %% style then we can try to put it on this page 
    \ifthenelse{\boolean{@firstcolumn} \or \not\boolean{@rmtwocol}}{%
      \@outputaddresses}{%
      %% otherwise give warning it will be on a separate page
      \ClassWarning{rmaa}{Full addresses will go on a separate
        page (please see User Guide for how to remedy this)}}%
    %% If we haven't managed yet, put it on its own page
    \ifthenelse{\boolean{@needaddresses}}{%
      \if@debug{\ClassWarning{rmaa}{In this silly bit!}}\fi
      \clearpage\vfill\@outputaddresses}{}%
    }{
    %% If we don't need them, we still might need to do the eoa hook 
    \ifthenelse{\boolean{@suppressfulladdresses}}{%
      \if@debug{\ClassWarning{rmaa}{XXX}}\fi
%%      \if@book\else\@endofarticlehook\null\fi
      }{}
    }%
  }  
%    \end{macrocode}
% \end{macro}

% \begin{macro}{\outputfulladdresses}
% A user interface to the above command, in case there is no biblio
% (11 May 2000) 
%    \begin{macrocode}
\newcommand{\outputfulladdresses}{\@outputaddressesifnecessary}
%    \end{macrocode}
% \end{macro}

% \begin{macro}{\adjustfinalcols}
% This command starts a new column and should be 
% used to ensure that the columns balance on the last page. It also
% outputs the full 
% author addresses to ensure that this occurs in the left hand
% column. 
% The |\par| command is so that the current line will be taken account
% of when working out |\pagetotal|
%    \begin{macrocode}
\newcommand{\adjustfinalcols}{\if@twocolumn\goodbreak\@outputaddresses\newpage\noindent\fi}
%    \end{macrocode}
% \end{macro}

% \begin{macro}{\@outputaddresses}
% Command to do the work of outputting the addresses. 
% This can be called by 
% \begin{enumerate}
% \item an explicit |\adjustfinalcols|
% \item via |\@outputaddressesifnecessary|, q.v.
% \end{enumerate}
%
%    \begin{macrocode}
\newcommand{\@outputaddresses}{%
  % Allow option of suppressing the full addresses, in which case
  % we just run the hook if we are not in a book style
  \if@suppressfulladdresses
%%     \if@book\else\@endofarticlehook\null\fi
  \else
  %% save the full addresses in a box for output later
  \sbox{\@fulladdbox}{\parbox{\textwidth}{\@fulladdresses}}
  % calculate height of address box  
  \settoheight{\@fulladdheight}{\usebox{\@fulladdbox}}
  \settodepth{\@fulladddepth}{\usebox{\@fulladdbox}}
  \addtolength{\@fulladdheight}{\@fulladddepth}
  \if@debug
  { \typeout{Address total height: \the\@fulladdheight}
    \typeout{Text height: \the\pagetotal}
    \typeout{Goal height: \the\pagegoal}
    }
  \fi
  % reuse \@fulladdheight to mean the total height of text+addresses
  \addtolength{\@fulladdheight}{\pagetotal}
%%  \addtolength{\@fulladdheight}{-\baselineskip}
  \if@debug\typeout{Adjusted text + addresses height: \the\@fulladdheight}\fi
  % check it will fit on current page
  \ifdim\@fulladdheight<\pagegoal
     %% must shift it to the left
     %% to compensate for the hanging indentation in the reference list. 
     %% NOTE: this is not foolproof!
     %% NB. This will give overfull hbox warnings, which is a bit ugly.
     \vfill\noindent\if@biblist\hspace*{-\leftmargin}\fi
     \usebox{\@fulladdbox}
     % Set this to show we've done it
     \@needaddressesfalse
     % and assume we are at the end of the article
     % so that we do the EOA hook if not in the book style
%%     \if@book\else\@endofarticlehook\fi
  \else
     % if not, just do nothing. With luck, this function will get 
     % called again on the next page
     \if@debug{\typeout{Holding over full addresses to next page}}\fi
  \fi\fi
  }
%    \end{macrocode}
% \end{macro}

% \subsubsection{Commands for editors}
% \label{sec:editor-commands}
% Some user/editor commands for doing the TOC and index in the book
% option. 
% They are defined here rather in the book-specific part since they
% need to be always available in order to avoid breaking documents. 
% \begin{macro}{\listofauthors}
% \begin{macro}{\@rawauthors}
% \begin{macro}{\rmaaTOCauthors}
% The command |\listofauthors| formats an author list for use in the
% TOC. It should contain the same names as in the |\author| command,
% but without the affiliation garbage and using |\&| instead of
% |and|. 
%    \begin{macrocode}
\newcommand{\@rawauthors}{Raw Authors}
\newcommand{\rmaaTOCauthors}{\noexpand\rmaaTOCAuthorStyle{Author list for TOC}}
\newcommand{\listofauthors}[1]{%
  \renewcommand{\@rawauthors}{#1}
  \renewcommand\rmaaTOCauthors{\noexpand\rmaaTOCAuthorStyle{#1}}}
%    \end{macrocode}
% \begin{macro}{\indexauthor}\label{indexauthor}
% The command |\indexauthor| produces an index entry for each
% author. 
% This command must come \emph{after} the
% |\title| and the |\listofauthors| commands. Also, there should be one 
% |\indexauthor| command for \emph{each} author. We also add the
% author onto the end of the |\ADS@authors| macro and construct the
% |\ADS@authorinitial| by a most cunning method (see below), but only
% for the first author. 
% 
% No, No, NO! This was all wrong. Backed out completely: 
% \begin{quote}
% Completely rewritten (19 Dec 2002) since we now use the index class, so 
% as to be able to alphabetize accented characters. A side effect of this is
% that |\index| no longer expands out the argument after the |@|, so we 
% instead add all the |\index| commands to a list (|\indexauthor@deferlist|), 
% which is executed in the |\@begofarticlehook|. 
% 
% Note the stuff with |\toks0| is copied from Exercise 20.15 of the \TeX{}book---it is
% to make sure the latin1-accented name gets passed into the idx
% file. 
% \end{quote}
% We now go back to the old way of doing things, using an optional
% argument to the \cs{indexauthor} command that has the name without
% accents. 
% \changes{3.08}{2003/01/16}{Remerged changes from 3.06 that got
% dropped from 3.07 after a synchronization disaster. Basically, the
% accentless optional argument to \cs{indexauthor}. Must add the macro
% directories to my unison set-up to stop this sort of thing from happening.}
%    \begin{macrocode}
\newcommand\@isortkey{}
\newcommand\@emptykey{empty}
\newcommand\indexauthor[2][empty]{
  \renewcommand\@isortkey{#1}
  \ifx\@isortkey\@emptykey \renewcommand\@isortkey{#2}\fi
  \index{\@isortkey @{{\rmaaIndexAuthorStyle{#2}}\hskip1em plus1em minus0.4em
      \rmaaIndexTitle\hskip1em plus2em minus0.4em\penalty-100
      \rmaaTOCauthors}}
  \indexauthorADS{#2}
}
%%\def\indexauthor@deferlist{\relax}
%%\long\def\indexauthor#1{
%%  \def\@@name{#1} \toks0=\expandafter{\@@name}
%%  \toks1=\expandafter{\indexauthor@deferlist}
%%  \protected@xdef\indexauthor@deferlist{%
%%    \the\toks1
%%    \noexpand\index{\the\toks0@{\noexpand\rmaaIndexAuthorStyle{\the\toks0}\noexpand\hskip1em plus1em minus0.4em 
%%        \rmaaIndexTitle\noexpand\hskip1em plus2em minus0.4em\noexpand\penalty-100
%%        \rmaaTOCauthors}}
%%  }
%%  \indexauthorADS{#1}
%%  }
%%\xdef\indexauthor@deferlist{}
\newcommand\indexauthorADS[1]{%
  \if@firstauthor\protected@xdef\ADS@authorinitial{\firstletter{#1}}\fi
  \protected@xdef\ADS@authors{\ADS@authors\if@firstauthor\else ; \fi
    #1}
  \@firstauthorfalse
}
%    \end{macrocode}
% Note that the first |#1| sets the key for the index, making sure it
% actually starts with  
% first letter of the surname. Meanwhile, the |@{...}| expression,
% which starts with some  
% control sequences is what is written to the |.ind| file. 
% See the makeindex man page for details. 
% \end{macro}
% \end{macro}
% \end{macro}
% \end{macro}

% \begin{macro}{\KillAllAccents}
% \begin{macro}{\firstletter}
% These are utility functions used in the above. |lett| is a counter
% indicating the current letter in the string. |\firstletter| just
% prints the first letter of its argument.
%    \begin{macrocode}
\newcommand{\KillAllAccents}{\def\'{}\def\`{}\def\~{}\def\.{}\def\i{i}}
\newcommand\firstletter[1]{\@car #1\@nil}
%    \end{macrocode}
% \end{macro}
% \end{macro}
% \begin{macro}{\SetVolume}
% \begin{macro}{\SetMSnumber}
% \begin{macro}{\SetFirstPage}
% \begin{macro}{\SetYear}
% \begin{macro}{\ReceivedDate}
% \begin{macro}{\AcceptedDate}
% These are for the editors to set the volume number, etc. 
%    \begin{macrocode}
\newcommand{\SetVolume}[1]{%
  \renewcommand{\@volume}{#1}
  \edef\ADS@volume{\PadToFourWithDots{#1}}
  \renewcommand{\@volumen}{\rmaa@VolumenStyle{volumen #1}}%
  }
\newcommand{\SetMSnumber}[1]{\renewcommand{\@MSnum}{#1}}
%    \end{macrocode}
% Note that there is no real reason to use |\SetFirstPage| since all
% the page numbering is done automatically in the book style. 
%    \begin{macrocode}
\newcommand{\SetFirstPage}[1]{\setcounter{page}{#1}}
\newcommand{\SetYear}[1]{%
  \renewcommand{\@year}{#1}
  \renewcommand{\ADS@year}{#1}
  \renewcommand{\@anyo}{\rmaa@VolumenStyle{#1}}
  }
\newcommand{\ReceivedDate}[1]{\renewcommand{\@received}{#1}}
\newcommand{\AcceptedDate}[1]{\renewcommand{\@accepted}{#1}}
%    \end{macrocode}
%
% \begin{macro}{\DOI}
% User/editor command to enter the DOI, which will be typeset as a link.
% \changes{4.0}{2019/11/02}{New user command \cs{DOI} to typeset
% DOI in the header of first page and add it to ADS info.}
%    \begin{macrocode}
\newcommand\DOI[1]{%
  \def\@DOI@url{\url{https://doi.org/10.22201/ia.01851101p.#1}}%
  \def\ADS@doi{https://doi.org/10.22201/ia.01851101p.#1}%
}%
%    \end{macrocode}
% \end{macro}
%
% \begin{macro}{\ORCIDiD}
% \changes{4.0}{2019/11/12}{Add mechanism for giving ORCID iD links}
% \begin{macro}{\ORCIDiDfull}
% \begin{macro}{\@orcid@icon}
% \begin{macro}{\@orcid@iconbig}
% Commands to add icon links to ORCID iDs of authors. For example
% \begin{quote}
% |\ORCIDiD|\marg{0000-0000-0000-0000}
% \end{quote}
% to be used after each author name in the |\author|
% command and in the |\fulladdresses|.  Alternatively,
% |\OORCIDiDfull| will typeset the icon plus the full text of the
% link, but we don't use that at present.  The icon image is stored in |ORCIDiD_iconvector.pdf|, which needs to be on the input path.  In principle, we could make EPS versions if necessary.
%    \begin{macrocode}
\newcommand\@orcid@icon{\includegraphics[height=1.8ex]{ORCIDiD_iconvector}}
\newcommand\@orcid@iconbig{\includegraphics[height=2ex]{ORCIDiD_iconvector}}
\newcommand\ORCIDiD[1]{\href{https://orcid.org/#1}{\,\@orcid@icon}}
\newcommand\ORCIDiDfull[1]{%
  \href{https://orcid.org/#1}{\@orcid@iconbig}\,\url{https://orcid.org/#1}}
%    \end{macrocode}
% \end{macro}
% \end{macro}
% \end{macro}
% \end{macro}
%
% \begin{macro}{\@volume}
% \begin{macro}{\@volumen}
% \begin{macro}{\@artpages}
% \begin{macro}{\@year}
% \begin{macro}{\@anyo}
% \begin{macro}{\@received}
% \begin{macro}{\@accepted}
% \begin{macro}{\@MSnum}
% Default values for the containers used by above commands, plus a
% few more useful ones. 
%    \begin{macrocode}
\newcommand{\@volume}{00}
\newcommand{\@volumen}{\rmaa@VolumenStyle{00}}
\newcommand{\@artpages}{%
  \oldpageref{\articlename:RMAAFirstPage}--\oldpageref{\articlename:RMAALastPage}}
\newcommand{\@year}{\the\year}
\newcommand{\@anyo}{\rmaa@VolumenStyle{2007}}
\newcommand{\@journalname}{\rmaa@JournalNameStyle{%
    \if@proceedings\RMAASClong\else\RMAAlong\fi}}
\newcommand{\@received}{}
\newcommand{\@accepted}{}
\newcommand{\@MSnum}{0000}
\newcommand{\@DOI@prefix}{https://doi.org/10.22201/ia.01851101p}
\def\@DOI@url{\url{https://doi.org/10.22201/ia.01851101p.2099.00.00.00}}
%    \end{macrocode}
% \end{macro}
% \end{macro}
% \end{macro}
% \end{macro}
% \end{macro}
% \end{macro}
% \end{macro}
% \end{macro}
% \end{macro}
% \end{macro}
% \end{macro}
% \end{macro}
% \end{macro}
% \end{macro}
% 
% The following commands are most useful in the \texttt{book}
% style. However, the |\SetConfTitle| command can now be used in
% individual articles to set the header for the \texttt{preprint}
% style. 
%    \begin{macrocode}
\newcommand{\@reunionplace}{Lugar de la reuni\'on}
\newcommand{\@reuniondate}{Fecha de la reuni\'on}
\newcommand{\@conftitle}{Conference Title}
\newcommand{\SetReunion}[2]{\renewcommand{\@reunionplace}{#1}\renewcommand{\@reuniondate}{#2}}
\newcommand{\SetConfTitle}[1]{\renewcommand{\@conftitle}{#1}}
\newcommand{\@editors}{Juan Ram\'on de la Fuente and ``El Mosh''}
\newcommand{\SetEditors}[1]{\renewcommand{\@editors}{#1}}
\newcommand{\@editores}{Juan Ram\'on de la Fuente y ``El Mosh''}
\newcommand{\SetEditores}[1]{\renewcommand{\@editores}{#1}}
\newcommand{\@mes}{mes de la publicaci\'on}
\newcommand{\SetMes}[1]{\renewcommand{\@mes}{\rmaa@VolumenStyle{#1}}}
%    \end{macrocode}

% \subsection{Notes for Author Proofs and Communication between Editors}
% \label{sec:notes}
% We introduce a command \cs{Note} that should be used to describe
% any change made to the manuscript. These are a bit tricky to get
% right. In general we use \cs{marginpar} but this only works in outer
% paragraph mode, so we just bung the box where we are if we are in a
% figure caption or something. If we are in the preamble, then we just
% use a \marginpar as part of the \maketitle command.
%    \begin{macrocode}
\newif\if@arenotes
\@arenotesfalse
\newcounter{notecounter}
\setcounter{notecounter}{0}
\newcommand\NotesList{}
\newcommand\Note[1]{%
  \if@proof{%
    \@arenotestrue%
    \stepcounter{notecounter}%
    \if@inpreamble\else{%
      \NoteMark{\arabic{notecounter}}%
    }\fi
    \protected@xdef\NotesList{\NotesList 
      \protect\item #1}%
  }\fi%
}
\newcommand\Changed[2]{Changed ``#1'' to ``#2''.}
\newcommand\NoteMark[1]{%
  \ifinner
  \raisebox{1ex}[0pt][0pt]{\makebox[0pt][l]{\framebox{\footnotesize\bfseries
        N#1}}}%
  \else
  \marginpar{\framebox{\footnotesize\bfseries N#1}}%
  \fi
}
\newcommand\RefNote[1]{\if@proof\NoteMark{\ref{#1}}\relax\fi}
\newcommand\NoteLabel{\protect\label}
\newcommand\OutputNotes{%
  \bgroup
  \onecolumn
  \section*{EDITORIAL NOTES}
  \begin{enumerate}
    \NotesList
  \end{enumerate}
  \egroup
}
\DeclareRobustCommand\cs[1]{\texttt{\char`\\#1}}

%    \end{macrocode}
% \subsection{Bibliographic Records for ADS}
% \label{sec:ADS}
% \changes{2.2beta11}{2001/11/27}{Just about finished all the ADS
% stuff now. Except for the references, which I have given up on for
% the time being.}
% \begin{macro}{\ADS@BibCode}
% \begin{macro}{\ADS@year}
% \begin{macro}{\ADS@journal}
% \begin{macro}{\ADS@volume}
% \begin{macro}{\ADS@firstpage}
% \begin{macro}{\ADS@authorinitial}
% We use the |\ADS@BibCode| command to hold the bibliographic code for
% ADS. It is 
% used in the |%R| 
% field of the bibliographic record for each article. 
% The format is 
% \begin{quote}
% \texttt{YYYYJJJJJVVVVMPPPPA}  
% \end{quote}
% where \texttt{YYYY} is the year, \texttt{JJJJJ} is the journal code,
% which is \texttt{RMxAA} or \texttt{RMxAC}, \texttt{VVVV} is the
% volume number (right-justified), \texttt{M} is not used by us,
% \texttt{PPPP} is the first page number (right-justified), and
% \texttt{A} is the initial of the first author's surname. The volume
% and page fields are padded with periods (`\texttt{.}'), which causes
% us some fun (this is done with the macro |\PadToFourWithDots| described
% below) . 
% 
% All these individual bits are held in the commands |\ADS@year|,
% |\ADS@journal|, |\ADS@volume|, |\ADS@firstpage|,
% |\ADS@authorinitial|. 
%    \begin{macrocode}
\newcommand{\SetADSBibCode}[1]{\renewcommand{\ADS@BibCode}{#1}}
\newcommand{\ADS@BibCode}
{\ADS@year\ADS@journal\ADS@volume.\PadToFourWithDots{\ADS@firstpage}\ADS@authorinitial}
\newcommand{\ADS@year}{2001} 
\newcommand{\ADS@journal}{\if@proceedings RMxAC\else RMxAA\fi}
\newcommand{\ADS@volume}{....} 
%    \end{macrocode}
% These two are just set to default values here, |\ADS@firstpage| is set
% in the |document| environment for included articles, in the
% |frontsection| environment, |\tableofcontents| command,  |index| and
% |abstracts| environments. Also by |\AtBeginDocument| for stand-alone
% articles. 
% 
% God knows how we are going to do |\ADS@authorinitial|!
%    \begin{macrocode}
\newcommand{\ADS@firstpage}{....}
\newcommand{\ADS@authorinitial}{A}
%    \end{macrocode}
% \end{macro}
% \end{macro}
% \end{macro}
% \end{macro}
% \end{macro}
% \end{macro}
% We also have a whole slew of other fields that can be set. These are
% all listed in Table~\ref{tab:ADS}. Most of them still have to be
% defined. 
% \begin{table}[!ht]\raggedleft
% \caption{ADS Bibliographic Records}\label{tab:ADS}
% \begin{tabular}{lp{8em}lp{8em}}\hline
% Code & Description & Macro Name & Comment \\ \hline
% |%R| & Bibliographic Code & |\ADS@BibCode| & required \\
% |%A| & Author list & |\ADS@authors| & required, generate from |\indexauthor| \\
% |%F| & Author affiliation & |\ADS@affil| & format: AA(\dots)
% AB(\dots) \dots \\
% |%J| & Journal Name & |\ADS@journalname| &  \\
% |%D| & Publication Date & |\ADS@fulldate| &  \\
% |%L| & Last Page & |\ADS@lastpage| & generate from label \\
% |%T| & Title & |\ADS@title| & required, generate from |\title| \\
% |%C| & Abstract Copyright & |\ADS@copyright| & set to IA-UNAM \\
% |%U| & URL for electronic document & |\ADS@url| & set to PDF file \\
% |%d| & DOI for article & |\ADS@doi| & set by |\DOI| \\
% |%K| & Keywords & |\ADS@keywords| & set equal to |\@keywords| \\
% |%A| & Abstract Text & |\ADS@abstract| & set equal to |\@abstract| \\
% |%R| & References & |\ADS@references| & generate from the
% bibliography  \\ \hline
% \end{tabular}
% \end{table}
% A lot of these macros are set by other commands, as mentioned in the
% Table, but some need to be set  explicitly here. 
% 
% \begin{macro}{\ADS@doi}
% \changes{4.0}{2019/11/08}{Save DOI for ADS}
% The user/editor command \cs{DOI} is used to set this. 
%    \begin{macrocode}
\newcommand{\ADS@doi}{https://doi.org/10.22201/ia.01851101p.0000.00.00.00}
%    \end{macrocode}
% \end{macro}
% 
% \begin{macro}{\ADS@journalname}
% This
% does not contain any accents at the moment so as not to cause
% problems---this maybe should be changed though. 
%    \begin{macrocode}
\newcommand{\ADS@journalname}{%
  \if@proceedings\@conftitle\space (Eds. \@editors) \fi%
  Revista Mexicana de Astronom\'\i{}a y Astrof\'\i{}sica%
  \if@proceedings\space (Serie de Conferencias)\fi%
  \space Vol. \@volume, pp. \ADS@firstpage--\ADS@lastpage\space (\@year)%
  \space (http://www.astroscu.unam.mx/\~{}rmaa/)
  }
%    \end{macrocode}
% \end{macro}
% \begin{macro}{\ADS@keywords}
% \begin{macro}{\ADS@abstract}
% \begin{macro}{\ADS@title}
% |\ADS@abstract| is straightforward, since we already have the text
% saved as |\@abstract|. We could do the same with |\ADS@keywords| but
% to avoid unwanted formatting we explicitly set it in |\keywords|
% (p.~\pageref{keywords}). |\ADS@title| is just set to |\@rawtitle|,
% which we were already using for another purpose. 
%    \begin{macrocode}
\let\ADS@title=\@rawtitle
\newcommand{\ADS@keywords}{No keywords specified}
\newcommand{\ADS@abstract}{No abstract specified}
%    \end{macrocode}
% \end{macro}\end{macro}\end{macro}
% \begin{macro}{\ADS@authors}
% |\ADS@authors| wants to be in ``Other, A. N.'' format, but with
% semi-colons between authors. The most natural way to do this seems to
% be with |\indexauthor| (p.~\pageref{indexauthor}). We append each
% author to this macro, which means we must explicitly clear it for
% each new article (done in |\article|). .   
%    \begin{macrocode}
\newcommand{\ADS@authors}{}
%    \end{macrocode}
% \end{macro}
% \begin{macro}{\ADS@fulldate}
% \begin{macro}{\ADS@lastpage}
% \begin{macro}{\ADS@copyright}
% \begin{macro}{\ADS@website}
% \begin{macro}{\ADS@url}
% \begin{macro}{\SetWebsite}
% \begin{macro}{\SetMonthNumeric}
% Other sundry data, together with user commands for bits that can be
% set by the Editors. The |\ADS@url| is constructed from the website
% (currently |http://www.astrosmo.unam.mx/rmaa/|,
% the journal descriptor (e.g., |RMxAC..10/|), followed by the
% path to the PDF file (e.g., |PDF/RMxAC..10_HEW1.pdf|). 
% 
%    \begin{macrocode}
\newcommand{\ADS@website}{http://www.astrosmo.unam.mx/rmaa/}
\newcommand{\ADS@url}
{\ADS@website\ADS@journal\ADS@volume\ADS@number/PDF/\ADS@journal\ADS@volume\ADS@number\_\articlename.pdf}
\newcommand{\ADS@number}{}
\newcommand{\ADS@fulldate}{00/\ADS@year} 
\newcommand{\ADS@lastpage}{00} 
\newcommand{\ADS@copyright}{(C) \ADS@year; Instituto de Astronom\'\i{}a,
  Universidad Nacional Aut\'onoma de M\'exico} 
\newcommand{\SetNumber}[1]{\renewcommand{\ADS@number}{-#1}}
\newcommand{\SetWebsite}[1]{\renewcommand{\ADS@website}{#1}}
\newcommand{\SetMonthNumeric}[1]{\renewcommand{\ADS@fulldate}{#1/\ADS@year}}
%    \end{macrocode}
% \end{macro}\end{macro}\end{macro}\end{macro}\end{macro}\end{macro}\end{macro}
% \begin{macro}{\ADS@affil}
% \begin{macro}{\ADS@references}
% These two are somewhat more problematical. Affiliation currently
% only works for the first author (see \dots) pending redesign of the
% user interface. References are now done by wrapping the
% |thebibliography| environment to open another input stream on the
% same file and search through for the |\bibitems|. 
%    \begin{macrocode}
\newcommand{\ADS@references}{}
\newcommand{\ADS@affil}{}
%    \end{macrocode}
% \end{macro}\end{macro}
% 
% Define two macros that expand to the strings
% |\begin{thebibliography}| and |\end{thebibliography}|. Note that the
% idea for this and much else is inspired by the code for the docstrip
% package. Also define |\StartBibX| in case authors use the more
% standard |\begin{thebibliography}{}| syntax. 
%    \begin{macrocode}
\edef\EndBib{\expandafter\@gobble\string\\end\string{thebibliography\string}}
\edef\StartBib{\expandafter\@gobble\string\\begin\string{thebibliography\string}}
\edef\StartBibX{\expandafter\@gobble\string\\begin\string{thebibliography\string}\string{\string}}
%    \end{macrocode}
% Define two new logical variables. |\ifAnotherline| is true if we
% want to read another line from the input stream. It is set to false
% when we encounter the |\end{thebibliography}|. |\ifInBib| is set to
% true when we frind the |\begin{thebibliography}| and tells us to
% start saving the following lines. 
%    \begin{macrocode}
\newif\ifAnotherLine
\newif\ifInBib
%    \end{macrocode}
% \begin{macro}{makeOther}
% We need a macro to change the catcodes of lots of things to other. 
%    \begin{macrocode}
\def\makeOther#1{\catcode`#1=12\relax}
%    \end{macrocode}
% \end{macro}
% \begin{macro}{\RereadBiblio}
% Define a macro to do all the fiddly stuff. This is used at the
% beginning of |thebibliography| environment (Sec~\ref{sec:biblio}).
%    \begin{macrocode}
\newread\tf@biblio
\newcommand{\RereadBiblio}{%
%    \end{macrocode}
% Open the current file again for reading
%    \begin{macrocode}
  \begingroup
  \immediate\openin\tf@biblio=\articlename
%    \end{macrocode}
% Everything in a group to be safe. Initialize logicals. Black magic
% to change
% catcodes and avoid blank lines being interpreted as |\par|. 
%    \begin{macrocode}
  \AnotherLinetrue\InBibfalse
  \makeOther\ \makeOther\\\makeOther\$%
  \makeOther\#\makeOther\^\makeOther\^^K%
  \makeOther\_\makeOther\^^A\makeOther\%%
  \makeOther\~\makeOther\{\makeOther\}\makeOther\&%
  \endlinechar-1\relax
%    \end{macrocode}
% Start looping through the file. Read each line to the macro
% |\Line|. If we have got to EOF or end of biblio then set
% |\AnotherLinefalse|.
%    \begin{macrocode}
  \loop
  \read\tf@biblio to \Line
  \ifx\Line\EndBib\AnotherLinefalse\InBibfalse\fi
  \ifeof\tf@biblio \AnotherLinefalse\fi
%    \end{macrocode}
% If we are in the biblio and the line is not blank then append it to
% the |\ADS@references| macro together with a newline. 
%    \begin{macrocode}
  \ifInBib
  \ifx\Line\empty\else\xdef\ADS@references{\ADS@references\Line^^J}\fi
  \fi
%    \end{macrocode}
% Check if the line we just saw was the start of biblio, then repeat
% the loop if necessary. 
%    \begin{macrocode}
  \ifx\Line\StartBib\InBibtrue\fi
  \ifx\Line\StartBibX\InBibtrue\fi
  \ifAnotherLine\repeat
  \immediate\closein\tf@biblio
  \endgroup
  }
%    \end{macrocode}
% \end{macro}
% 
% \begin{macro}{\PadToFourWithDots}
% Utility function for left-padding a number to 4 characters with
% periods. Used in 
% constructing |\ADS@BibCode|. This is a very simple algorithm and
% \emph{requires} that the argument actually be a number---if it is
% not then dire consequences will result.  
%    \begin{macrocode}
\newcommand{\PadToFourWithDots}[1]{%
  \ifnum #1 > 999 #1%
  \else\ifnum #1 > 99 .#1%
  \else\ifnum #1 > 9 ..#1%
  \else ...#1%
  \fi\fi\fi}
%%  \ifnum \count0 > 999 #1%
%%  \else\ifnum \count0 > 99 .#1%
%%  \else\ifnum \count0 > 9 ..#1%
%%  \else ...#1%
%%  \fi\fi\fi}
%    \end{macrocode}
% \end{macro}
%
% \begin{macro}{\AuthorADS}
% \begin{macro}{\ADS@affiliations}
% New mechanism for setting the affiliations for ADS. The author labels (AA, AB, \dots) are generated using two nested counters.  We use the |\robustify| command from \texttt{etoolbox} so that accented characters in the affiliation addresses do not blow up.  The list of accents may need to be expanded. 
%    \begin{macrocode}
\newcounter{AUTHOR}
\newcounter{author}[AUTHOR]
\stepcounter{AUTHOR}
\newcommand\ADS@affiliations{}
\robustify{\'}\robustify{\i}\robustify{\~}\robustify{\`}
\def\AuthorADS#1#2{%
  \begingroup
  \def\ORCIDiD##1{<ID system="ORCID">##1</ID>}
  \def\Email##1{<EMAIL>##1</EMAIL>}
  \def\Affil##1{##1}
  \stepcounter{author}
  \ifthenelse{\value{author} > 26}{\stepcounter{AUTHOR}\stepcounter{author}}{}
  \ifthenelse{\equal{\ADS@affiliations}{\@rmaaEmpty}}{%
    \xdef\ADS@affiliations{\Alph{AUTHOR}\Alph{author}(#2)}
  }{%
    \xdef\ADS@affiliations{\ADS@affiliations, \Alph{AUTHOR}\Alph{author}(#2)}
  }%
  \endgroup
}
%    \end{macrocode}
% \end{macro}
% \end{macro}
% %
% \begin{macro}{\ADS@dumpdata}
% \changes{4.0}{2019/11/13}{Write all author affiliations for ADS}
% \begin{macro}{\ADS@write}
% \begin{macro}{\ADS@openfile}
% |\ADS@dumpdata| is the macro that actually does the job of writing out all the
% data to an auxiliary file, using |\ADS@write| to write each
% field. The file is set up with |\ADS@openfile|, which currently is
% done in the TOC of a book.  The new multi-author affiliations are used if |\ADS@affiliations| has been set, otherwise it falls back to the single-author version |\ADS@affil|. 
%    \begin{macrocode}
\newcommand{\ADSfield}[2]{\relax}
\newcommand{\ADSvfield}[2]{\relax}
%%\renewcommand{\ADSfield}[12]{\%#1 #2}
\newcommand{\ADS@dumpdata}
{ 
  \ADS@write{\string\begin{ADSarticle}}
  \ADS@WRITE{\string\ADSfield{R}{\ADS@BibCode}}
  \ADS@write{\string\ADSfield{A}{\ADS@authors}}
  \ifthenelse{\equal{\ADS@affiliations}{\@rmaaEmpty}}{%
    {\ADS@sanitize\ADS@write{\string\ADSfield{F}{AA(\ADS@affil)}}}
  }{%
    {\ADS@sanitize\ADS@write{\string\ADSfield{F}{\ADS@affiliations}}}
  }
  \ADS@write{\string\ADSfield{J}{\ADS@journalname}}
  \ADS@write{\string\ADSfield{D}{\ADS@fulldate}}
  \ADS@write{\string\ADSfield{L}{\ADS@lastpage}}
  {\ADS@sanitize\ADS@write{\string\ADSfield{T}{\ADS@title}}}
  \ADS@write{\string\ADSfield{C}{\ADS@copyright}}
  \ADS@write{\string\ADSfield{d}{\ADS@doi}}
  \ADS@write{\string\ADSfield{U}{\ADS@url}}
  \ADS@write{\string\ADSfield{K}{\ADS@keywords}}
  {\ADS@sanitize\ADS@write{\string\ADSfield{B}{\ADS@abstract}}}
  \ADS@write{\string\ADSfield{Z}{\ADS@references}}
  \ADS@write{\string\end{ADSarticle}}
  }
%%\newcommand{\ADS@write}[1]{\if@book\addtocontents{ads}{#1}\fi}
%%\newcommand{\ADS@openfile}{\@starttoc{ads}}
\newcommand{\ADS@write}[1]{\set@display@protect\immediate\write\tf@ads{#1}}
\newcommand{\ADS@WRITE}[1]{\immediate\write\tf@ads{#1}}
\newcommand{\ADS@openfile}{%
  \newwrite\tf@ads%
  \immediate\openout\tf@ads\jobname.ads\relax
  \ADS@WRITE{\ADS@preamble}}
\newcommand{\ADS@closefile}{\ADS@WRITE{\ADS@postamble}}
\def\ADS@preamble{%
  %%  \string\documentclass{article}
  %%  \string\begin{document}
  }
\def\ADS@postamble{%
  %%  \string\end{document}
  \string%\string% Local Variables:^J 
  \string%\string% mode: LaTeX^J
  \string%\string% mode: Font-lock^J
  \string%\string% End:^J
  }
%    \end{macrocode}
% \end{macro}
% \end{macro}
% \end{macro}
% 
% \subsection{Automatic calculation of the page range}
% \label{sec:page-range}
% Another issue that has cost me a lot of grief. Still, it finally
% seems to be sorted now. All hail Jeff Goldberg. 
%
% Cunning trick to get the page numbers from |\pageref|.  Put a label
% on the first and last page of the document using |\AtBeginDocument|
% and |\AtEndDocument|. .  This way won't work for the book option.
% In that case, we do it by defining |\@endofarticlehook| and
% |\@begofarticlehook| as parallels to |\AtEndDocument| and
% |\AtBeginDocument|.  Note that we have to defer adding
% |\@begofarticlehook| to |\AtBeginDocument| to make sure it comes
% after the redefinitions of |\ref| and |\label|.
%    \begin{macrocode}
\if@book\AtEndDocument{\ADS@closefile}\else
 %\AtBeginDocument{\@begofarticlehook}
\AtEndDocument{\@endofarticlehook}
\fi
%    \end{macrocode}
% \begin{macro}{\@begofarticlehook}
% Since we are getting more and more things that need to be done at
% the beginning of every article, frontsection, etc we define a sister
% hook to 
% |\@endofarticlehook|. This hook does the following: 
% \begin{enumerate}
% \item Labels the first page
% \item Sets |\ADS@firstpage|
% \item runs and then clears |\indexauthor@deferlist|
% \item resets the keyword counter
% \end{enumerate}
% It is executed in the following circumstances: 
% \begin{enumerate}
% \item At |\begin{document}| for standalone files
% \item At |\begin{document}| for each included file
% \item At beginning of front sections 
% \item At beginning of TOC
% \item At beginning of Author Index
% \item At beginning of abstracts
% \end{enumerate}
% For some reason, this is not always within the right group, so we
% use |\xdef| just to be sure. 
%    \begin{macrocode}
\newcommand{\@begofarticlehook}{%
  \if@frontmatter\frontlabel{RMAAFirstPage}\else\label{RMAAFirstPage}\fi%
  \xdef\ADS@firstpage{\@arabic{\c@page}}%
  \if@book\else\makeindex\fi % needed so default index is defined in stand-alone articles
%%  \indexauthor@deferlist
%%  \xdef\indexauthor@deferlist{}
  \@inpreamblefalse
  \setcounter{kwdcounter}{0}
  }
%    \end{macrocode}
% \end{macro}
% \begin{macro}{\@endofarticlehook}
% This is a hook to be run at the end of each article
% currently, it just labels the last page and writes out all the ADS
% data to an auxiliary file.  
% It should not produce any printed output. 
% It is modelled on the |\AtEndDocument| command, to which it reduces
% in the case of a single article. 
%    \begin{macrocode}
\newcommand{\@endofarticlehook}{%
  \if@debug{\ClassWarning{rmaa}{***Doing Hook***}}\fi
  \if@debug{\ClassWarning{rmaa}{ADSabstract is: \ADS@abstract}}\fi
  \clearpage\rmaa@lastpagelabel\if@book\ADS@dumpdata\fi
  \ifx\NotesList\@rmaaEmpty\else\OutputNotes\fi
  \if@debug{\ClassWarning{rmaa}{Notes list is: \NotesList}}\fi
}
%    \end{macrocode}
% \begin{macro}{\rmaa@lastpagelabel}
% New version of last page label command copied from Jeff Goldberg's
% lastpage.sty package - we can't just use |\label| here in case the
% output routine for the last page has already been called. Instead,
% we first explicitly clear the page, step back the page counter,
% explicitly write to the .aux file, then restep the page counter
% again.
%
% But this will conflict with the hyperref package, which does labels differently.  Compensate for that by using the five-argument form.
% 
% Also set the |\ADS@LastPage| macro to the appropriate number. 
%    \begin{macrocode}
\newcommand{\rmaa@lastpagelabel}{%
  \addtocounter{page}{-1}%
  \immediate\write\@auxout{\string
    \newlabel{\articlename:RMAALastPage}{{}{\thepage}{}{}{}}}%
  \edef\ADS@lastpage{\thepage}%
  \addtocounter{page}{1}}
%    \end{macrocode}
% \end{macro}
% \end{macro}

% \subsection{Deal with explicit change in number of columns}
% \label{sec:columns}
% Sometimes, when formatting a book, we need to switch between
% one-column articles and two-column articles. For this to work, we
% need to wrap the standard \LaTeX{} commands to set the |\@rmtwocol|
% flag appropriately and also reset the package default values of all
% scalable lengths. .  
%    \begin{macrocode}
\let\latexonecolumn=\onecolumn
\let\latextwocolumn=\twocolumn
\renewcommand{\onecolumn}{\@rmtwocolfalse\@package@default@lengths\latexonecolumn}
\renewcommand{\twocolumn}{\@rmtwocoltrue\@package@default@lengths\latextwocolumn}
%    \end{macrocode}
% 
% \subsection{The actual formatting of the title}
% \label{sec:title-format}

% Note, we have to redefine the |\maketitle| command to use the
% numbered footnotes. This means we needn't use the atbeginend stuff. 
%    \begin{macrocode}
\renewcommand{\maketitle}{\par
  \begingroup
%    \end{macrocode}
% Don't switch to marked footnotes. 
% This is consistent with existing style but I'm not sure I like it
%    \begin{macrocode}
  %%  \renewcommand\thefootnote{\fnsymbol{footnote}}%
%%  \def\@makefnmark{\hbox to\z@{$\m@th^{\@thefnmark}$\hss}}%
  \def\@makefnmark{\textsuperscript{\@thefnmark}\hspace{-0.2em}}%
  \long\def\@makefntext##1{\parindent 1em\noindent
    \hbox to1.8em{\hss$\m@th^{\@thefnmark}$}##1}%
  \if@rmtwocol
  { \if@debug{\ClassWarning{rmaa}{TWO COLUMN}}\fi
    \twocolumn[\@maketitle]%
    }
  \else
  { \if@debug{\ClassWarning{rmaa}{ONE COLUMN}}\fi
    \newpage
    \global\@topnum\z@   % Prevents figures from going at top of page.
    \if@manuscript \@MSmaketitle \else \@maketitle \fi
    }
  \fi
  \if@manuscript\else\thispagestyle{rmaafirstpage}\@thanks\fi
  \ifx\NotesList\@rmaaEmpty\else{%
    \NoteMark{1--\arabic{notecounter}}\vspace*{-2\baselineskip}%
  }\fi%
  \endgroup
%    \end{macrocode}
% Don't reset footnote counter (not sure whether this is good thing
% or not).
%    \begin{macrocode}
  %%  \setcounter{footnote}{0}%
  \let\thanks\relax\let\affil\relax\let\altaffil\relax
  \let\maketitle\relax\let\@maketitle\relax
  \gdef\@thanks{}\gdef\@author{}\gdef\@title{}}
%    \end{macrocode}

% \begin{macro}{\@kwhwidth}
% \begin{macro}{\@kwtwidth}
% Redefine \cs{@maketitle} to include the abstract, resumen and
% keywords. Note that in the manuscripty style we use
% \cs{@MSmaketitle} instead.  
% \changes{3.16}{2003/04/07}{
%  Added extra vertical space between authors and accepted date.}
% \changes{3.26}{2011/08/15}{Add extra label to review articles}
%    \begin{macrocode}
\newlength{\@kwhwidth}\newlength{\@kwtwidth}
\newcommand\reviewname{Review}
\renewcommand{\@maketitle}{%
  \newpage
  \null
  \if@reviewarticle
  \smallskip
  {\large\bfseries\itshape\reviewname}%
  \bigskip
  \else
  \vskip\@before@title@skip%
  \fi
  \begin{center}%
    {\@title \par}%
    \vskip\@before@author@skip%
    \parbox{\rmaa@AuthorWidthFrac\textwidth}{\centering
      \lineskip\@between@author@skip%
      \@author
      \par}
    %% this should be replaced with `Submitted to...' for preprints
    \if@printdates
    \vskip\@before@dates@skip%
    \if@preprint
    \rmaa@ReceivedStyle{Submitted to RMxAA \@received}
    \else
    \rmaa@ReceivedStyle{Received \@received; accepted \@accepted}
    \fi
    \fi
    \if@manuscript
    \rmaa@ReceivedStyle{Submitted to RMxAA \@received}
    \fi
    \if@guide
    \rmaa@ReceivedStyle{Version \rmaa@version, \rmaa@date}
    \fi
    \vskip\@before@abstract@skip%
    \if@abstractpresent
    \parbox{\rmaa@AbstractWidthFrac\textwidth}{%
      \setlength{\parindent}{0.9cm}%
      \centerline{\rmaa@AbstractHeaderStyle{Abstract}}\par%
      \vskip\@after@abstracthead@skip%
      \if@newstyle\noindent\fi\@abstract
      \vskip\@between@abstract@skip%
      \centerline{\rmaa@AbstractHeaderStyle{Resumen}}\par%
      \vskip\@after@abstracthead@skip%
      \if@newstyle\noindent\fi\@resumen
      \if@setresumen\else\@abstract\fi % append english abstract if
                                       % spanish one not set
      \if@guide\else            %do not use keywords for guide
      \vskip\@before@kwd@skip%
      % Format the Key Words with hanging indentation:
      % They are put in a \parbox with width equal to the line width
      % minus the width of the header. 
      { % Widths of keyword
        \settowidth{\@kwhwidth}{\@keywordheader~}   % header and
        \setlength{\@kwtwidth}{\linewidth}          % keyword text
        \addtolength{\@kwtwidth}{-\@kwhwidth}       %
        \noindent\@keywordheader~\parbox[t]{\@kwtwidth}{%
          \rmaa@KeywordStyle{\boldmath\uppercase{\@keywords}}}}%
      \fi
      }
      \fi
  \end{center}%
  \vspace*{\afterkwdskip}
  \par
  }
\newlength{\afterkwdskip}\setlength{\afterkwdskip}{\z@}
%    \end{macrocode}
% \end{macro}
% \end{macro}
%
%
% \begin{macro}{\reviewarticle}
% \begin{macro}{\reviewtoc}
% \begin{macro}{\reviewlof}
%   User command to specify a review article. Must be used in
%   preamble. Main effect is to place the text ``Review'' before the
%   article title.  Also add command for setting a table of contents
%   and a list of figures at the beginning of the article.
%    \begin{macrocode}
\newcommand{\reviewarticle}{\@reviewarticletrue}
\newcommand\reviewtoc{%
  \section*{CONTENTS}%
  \@starttoc{rtc}%
}
\newcommand\reviewlof{%
  \section*{LIST OF FIGURES}%
  \@starttoc{lof}%
}
%    \end{macrocode}
% \end{macro}
% \end{macro}
% \end{macro}
% \begin{macro}{\l@section}
%    \begin{macrocode}
\renewcommand*\l@section[2]{%
  \ifnum \c@tocdepth >\z@
    \addpenalty\@secpenalty
    \addvspace{0.8em \@plus\p@}%
    \setlength\@tempdima{1.5em}%
    \begingroup
      \parindent \z@ \rightskip \@pnumwidth
      \parfillskip -\@pnumwidth
      \leavevmode %\bfseries
      \advance\leftskip\@tempdima
      \hskip -\leftskip
      #1\nobreak\hfil \nobreak\hb@xt@\@pnumwidth{\hss #2}\par
    \endgroup
  \fi}
%    \end{macrocode}
% \end{macro}
%
% \begin{macro}{\@MSmaketitle}
% A special version of |\@maketitle| for the \texttt{manuscript}
% option. This can deal with a multi-page abstract/resumen. Not very
% well, mind. 
%    \begin{macrocode}
\newcommand{\@MSmaketitle}{%
  \if@debug\ClassInfo{rmaa}{Using special MS maketitle command}\fi
  \newpage
  {
    \thispagestyle{rmaafirstpage}
%     \setlength{\parindent}{0pt}
    \null\vskip 2ex
    \begin{center}
      \@title\par
      \@author\par
      \rmaa@ReceivedStyle{%
        \ifx\@received\@rmaaEmpty
          Draft version: \@date
        \else
          Received: \@received
          \ifx\@accepted\@rmaaEmpty
          \else
            ; Accepted: \@accepted
          \fi
        \fi}\par
    \end{center}
    \@thanks
    \centerline{\rmaa@AbstractHeaderStyle{Resumen}}
    \noindent\@resumen\par\vskip 2ex
    \centerline{\rmaa@AbstractHeaderStyle{Abstract}}
    \noindent\@abstract\par\vskip 2ex
    \noindent\@keywordheader~\rmaa@KeywordStyle{\@keywords}\par
%     \@keywordheader~\rmaa@KeywordStyle{\uppercase{\@keywords}}\par
  }%
  }
%    \end{macrocode}
% \end{macro}

% \subsection{The sectioning commands}
% \label{sec:section-commands}

% \begin{macro}{\section}
% \begin{macro}{\subsection}
% \begin{macro}{\subsubsection}
% \begin{macro}{\paragraph}
% \begin{macro}{\subparagraph}
% Redefine all the standard commands to use these lengths. 
%    \begin{macrocode}
\renewcommand\section{\@startsection {section}{1}{\z@}%
                                   {\abovesectionskip}%
                                   {\belowsectionskip}%
                                   %% how do we get upper case here
                                   {\centering\reset@font\normalsize}}
\renewcommand\subsection{\@startsection{subsection}{2}{\z@}%
                                     {\abovesectionskip}%
                                     {\belowsectionskip}%
                                     {\centering\reset@font\normalsize\itshape}}
\renewcommand\subsubsection{\@startsection{subsubsection}{3}{\z@}%
                                     {\abovesectionskip}%
                                     {\belowsectionskip}%
                                     {\centering\reset@font\normalsize\itshape}}
\renewcommand\paragraph{\@startsection{paragraph}{4}{\z@}%
                                    {\aboveparaskip}%
                                    {\belowparaskip}%
                                    {\reset@font\normalsize\bfseries}}
\renewcommand\subparagraph{\@startsection{subparagraph}{5}{\parindent}%
                                      {\aboveparaskip}%
                                      {\belowparaskip}%
                                      {\reset@font\normalsize\bfseries}}
%    \end{macrocode}
% \end{macro}
% \end{macro}
% \end{macro}
% \end{macro}
% \end{macro}

% \begin{macro}{secnumdepth}
% Make sure that we number the subsubsections---this is done by
% default in the |article| class but not the |book| class, which is
% why we must set it explicitly here.
%    \begin{macrocode}
\setcounter{secnumdepth}{3}
%    \end{macrocode}
% \end{macro}
% \changes{2.2beta10}{2001/11/07}{Fixed \cs{secnumdepth} for \cs{book}
% style.}

% \begin{macro}{\thesubsection}
% \begin{macro}{\thesubsubsection}
% For subsection and subsubsection, we want the numbers to be roman
% but the title to be italic. This seems the easiest way to go about it.
%    \begin{macrocode}
\renewcommand{\thesubsection}{{\upshape\thesection.\@arabic\c@subsection}}
\renewcommand{\thesubsubsection}{\thesubsection{\upshape.\@arabic\c@subsubsection}}
%    \end{macrocode}
% \end{macro}
% \end{macro}

% \begin{macro}{\@seccntformat}
% Add ``.'' to the number here so we can remove it from |\thesection|,
% etc. Reduce gap between number and title WJH 27 Jan 1995
%    \begin{macrocode}
\def\@seccntformat#1{\csname the#1\endcsname .\hskip 0.4em}
%    \end{macrocode}
% \end{macro}

% \begin{macro}{\@sect}
% I didn't really want to mess with the |\@sect| command, but it was
% the only way I could manage to force upper case. Other than the
% commented parts, it is identical to the standard version. 
%    \begin{macrocode}
\def\@sect#1#2#3#4#5#6[#7]#8{
  \ifnum #2>\c@secnumdepth
     \let\@svsec\@empty
  \else
     \refstepcounter{#1}%
     \let\@@protect\protect
     \def\protect{\noexpand\protect\noexpand}%
     \edef\@svsec{\@seccntformat{#1}}%
     \let\protect\@@protect
  \fi
%    \end{macrocode}
% This is the bit I changed (27 Jan 1995). Use upper case for
% |\section|, use mixed case for the rest. 
%    \begin{macrocode}
  \ifnum #2=1                   
     \def\sec@upcase##1{\uppercase{##1}}%
     \def\sec@shape{\upshape}%
  \else                         
     \def\sec@upcase##1{\relax{##1}}%
     \def\sec@shape{\itshape}%
  \fi
  \@tempskipa #5\relax
  \ifdim \@tempskipa>\z@
     \begingroup #6\relax
          \@hangfrom{\hskip #3\relax\@svsec}%
                    {\interlinepenalty \@M \sec@upcase{#8}\par}%
     \endgroup
     \csname #1mark\endcsname{#7}%
%    \end{macrocode}
% Only write TOC entries if doing the Guide style. 
%    \begin{macrocode}
     \if@guide 
        \addcontentsline{toc}{#1}{%
          \ifnum #2>\c@secnumdepth \else
             \protect\numberline{\csname the#1\endcsname}%
          \fi
          #7}%
     \fi
%    \end{macrocode}
% Also allow per-article TOC for review papers
%    \begin{macrocode}
     \if@reviewarticle
        \ifnum #2>2 \else
           \addcontentsline{rtc}{#1}{%
             \protect\numberline{\csname the#1\endcsname}%
             \sec@upcase{\sec@shape #7}}%
        \fi
     \fi
  \else
     \def\@svsechd{#6\hskip #3\relax  %% \relax added 2 May 90
       \@svsec #8\csname #1mark\endcsname{#7}%
       \if@guide
          \addcontentsline{toc}{#1}{%
            \ifnum #2>\c@secnumdepth \else
               \protect\numberline{\csname the#1\endcsname}%
            \fi
            #7}%
       \fi
       }%
  \fi
  \@xsect{#5}}
%    \end{macrocode}
% \end{macro}


% \subsection{Appendices}
% \label{sec:appendices}
% Make sure that appendices in individual articles work properly. We
% copy the definition from the article class (the book one doesn't do
% what we want).  
% \begin{environment}{appendices}
%   This uses an environment |appendices| inside of which each
%   |\section| will become a separate appendix.  Note that the
%   appendices must come before the references (otherwise the full
%   addresses mechanism will not work properly). We make the equation
%   numbering be subordinate to the appendix (equations A12, B6,
%   etc). Finally, we have a medskip followed by APPENDICES to
%   introduce them. \changes{3.17}{2002/04/08}{ Added appendices environment
%   that does the right thing.}
%    \begin{macrocode}
\newenvironment{appendices}{\par
  \begingroup
  \setcounter{section}{0}%
  \setcounter{subsection}{0}%
  \renewcommand\thesection{\@Alph\c@section}
  \renewcommand\theequation{\thesection\arabic{equation}}
  \medskip
  \section*{APPENDICES}}
{\endgroup}
%    \end{macrocode}
% \end{environment}
% \begin{environment}{appendix}
%   If there is only one appendix, then use |appendix| environment
%   instead.  \changes{4.0}{2019/11/12}{ Added
%   appendix environment for case where there is only one section.}
%    \begin{macrocode}
\renewenvironment{appendix}{\par
  \begingroup
  \setcounter{section}{0}%
  \setcounter{subsection}{0}%
  \renewcommand\thesection{\@Alph\c@section}
  \renewcommand\theequation{\thesection\arabic{equation}}
  \medskip
  \section*{APPENDIX}}
{\endgroup}
%    \end{macrocode}
% \end{environment}

% \subsection{Floating thingies}
% \label{sec:floats}   

%% 
%    \begin{macrocode}
\renewcommand{\thetable}{\@arabic\c@table}
\renewcommand{\thefigure}{\@arabic\c@figure}
\renewcommand{\figurename}{Fig.}
\renewcommand{\tablename}{TABLE}
%\renewcommand{\fnum@table}{\MakeUpperCase{\tablename}~\thetable}
%    \end{macrocode}
% For some reason, I have to write `TABLE' explicitly here
% If I use |\tablename|, it gets converted to mixed case---God knows
% why. 
% |\figurename| is also being ignored. 
%    \begin{macrocode}
\renewcommand{\fnum@table}{TABLE~\thetable}
\renewcommand{\fnum@figure}{Fig.~\thefigure}
%    \end{macrocode}

% Do captions in 9 pt face in the main journal style and in new-style
% anything. 
%    \begin{macrocode}
\long\def\@makefigcap#1#2{%
  \if@proceedings\else\small\fi
  \if@newstyle\small\fi 
  \vskip\abovecaptionskip
  \sbox\@tempboxa{#1. #2}%
  \ifdim \wd\@tempboxa >\hsize
    #1. #2 \par
  \else
    \global \@minipagefalse
    \hb@xt@\hsize{\hfil\box\@tempboxa\hfil}%
  \fi
  \vskip\belowcaptionskip}
%    \end{macrocode}

% Make sure table caption is upper case. 
%    \begin{macrocode}
\long\def\@maketabcap#1#2{%
  \hbox to\hsize{\hfil #1 \hfil}
  \vskip 6pt 
  \centering\MakeUpperCase{#2}\par
  \vskip 6pt}
%    \end{macrocode}

%    \begin{macrocode}
\newlength{\@rottabcapwidth}
\long\def\@makerottabcap#1#2{%
  \setbox\@tempboxa\hbox{\MakeUpperCase{#2}}%
  \ifdim\wd\@tempboxa >.8\vsize
     \setlength\@rottabcapwidth{.8\vsize}
  \else
     \setlength\@rottabcapwidth{\wd\@tempboxa}
  \fi
  \rotatebox{90}{%
    \begin{minipage}{\vsize}\centering
      \begin{minipage}{\@rottabcapwidth}\centering
        #1\\ \vskip 6pt \MakeUpperCase{#2}
      \end{minipage}%
    \end{minipage}
    }%
  \hskip 6pt
  }

\newenvironment{Sideways}
    {\begin{sideways}\begin{minipage}{\vsize}\centering}
    {\end{minipage}\end{sideways}}

\newenvironment{Sidewaystable}
{%
  \let\oldcaption=\caption%
  \def\caption{
    \def\@captype{table}\renewcommand{\@makecaption}{\@maketabcap}
    \oldcaption}
  \begin{sidewaystable}\centering}
  {\end{sidewaystable}\let\caption=\oldcaption}
\newenvironment{Sidewaystable*}
{%
  \let\oldcaption=\caption%
  \def\caption{
    \def\@captype{table}\renewcommand{\@makecaption}{\@maketabcap}
    \oldcaption}
  \begin{sidewaystable*}\centering}
  {\end{sidewaystable*}\let\caption=\oldcaption}
%    \end{macrocode}

% Not sure why we need to stick these lines in. 
%    \begin{macrocode}
\renewenvironment{figure}
                 {\renewcommand{\@makecaption}{\@makefigcap}\@float{figure}}
                 {\end@float}
\renewenvironment{figure*}
                 {\renewcommand{\@makecaption}{\@makefigcap}\@dblfloat{figure}}
                 {\end@dblfloat}
\renewenvironment{table}
                 {\renewcommand{\@makecaption}{\@maketabcap}\@float{table}}
                 {\end@float}
\renewenvironment{table*}
                 {\renewcommand{\@makecaption}{\@maketabcap}\@dblfloat{table}}
                 {\end@dblfloat}
%    \end{macrocode}

% Redefine |\@caption| so that we don't bother with the list of
% figures. This saves us having to worry about fragile commands in the
% captions. 
%    \begin{macrocode}
\long\def\@caption#1[#2]#3{%
  \par
%%  \addcontentsline{\csname ext@#1\endcsname}{#1}%
%%    {\protect\numberline{\csname the#1\endcsname}{\ignorespaces #2}}%
  \begingroup
    \@parboxrestore
    \if@minipage
      \@setminipage
    \fi
    \normalsize
    \@makecaption{\csname fnum@#1\endcsname}{\ignorespaces #3}\par
  \endgroup}
%    \end{macrocode}

% Make tables not quite so cramped. Allow |\hline\hline| to give a
% poor man's thick rule. 
%    \begin{macrocode}
\renewcommand{\arraystretch}{1.15}
\setlength{\doublerulesep}{0.4pt}
%    \end{macrocode}

% \begin{macro}{\tablecols}
% \begin{macro}{\@tablecols}
% \begin{macro}{\tabcaptioncontinued}
% |\tablecols| is a user command for setting the number of columns. It
% is used by commands that try and span the entire table (I think
% there must be a better way of acheiving
% this). |\tabcaptioncontinued| is 
% a user command for setting the caption of a continued table on the
% next page. 
%    \begin{macrocode}
\newcommand{\tablecols}[1]{\renewcommand{\@tablecols}{#1}}
\newcommand{\@tablecols}{1}
\newcommand{\tabcaptioncontinued}{%
  \multicolumn{\@tablecols}{c}{\addtocounter{table}{-1}TABLE~\thetable\ (CONTINUED)} \\ \\
  }
%    \end{macrocode}
% \end{macro}
% \end{macro}
% \end{macro}

% \subsubsection{Footnotes in tables}
% \label{sec:table-footnotes}
% \begin{macro}{\tablenotemark}
% \begin{macro}{\tablenote}
% \begin{macro}{\tablenoteskip}
% Old user commands |\tablenotemark{MARK}|,
% |\tablenote{MARK}{TEXT}|---now deprecated. 
%    \begin{macrocode}
\newcommand\tablenotemark[1]{$^\mathrm{#1}$} 
\newcommand\tablenote[2]{\if@proceedings\else\small\fi ${}^\mathrm{#1}$ #2} 
\newcommand\tablenoteskip{\\[0.5\baselineskip]} 
%    \end{macrocode}
% \end{macro}
% \end{macro}
% \end{macro}

% \begin{macro}{\tabnotemark}
% \begin{macro}{\tabnotetext}
% \begin{macro}{\tabnotetextTWO}
% Alternative way of doing footnotes in tables. 
%    \begin{macrocode}
\newcommand{\tabnotemark}[1]{\textsuperscript{\lowercase{#1}}}
%\newcommand{\tabnotetextTWO}[4]{\tabnotetext{#1}{#2}\tabnotetext{#3}{#4}}
\newcommand{\tabnotetextTWO}[4]{%
  \multicolumn{\@tablecols}{l}{%
    \parbox{0.5\tabnotewidth}{\small\textsuperscript{#1}#2}%
    \parbox{0.5\tabnotewidth}{\small\textsuperscript{#3}#4}%
    }%
  \\}
\newcommand{\tabnotetext}[2]{%
  \multicolumn{\@tablecols}{@{}p{\tabnotewidth}@{}}{\small\textsuperscript{#1}#2}\\}
\newlength{\tabnotewidth}
%    \end{macrocode}
% \end{macro}
% \end{macro}
% \end{macro}


% \begin{macro}{changemargin}
%    \begin{macrocode}
% Command to allow margins to be changed on the fly. Useful to allow
% tables and Figures to protrude into the margins. 
%
% \changes{3.24}{2007/09/14}{Added \texttt{changemargin} environment
% (copied from TeX FAQ).}
\newenvironment{changemargin}[2]{% 
  % From TeX FAQ http://www.tex.ac.uk/cgi-bin/texfaq2html?label=chngmargonfly
  \begin{list}{}{%
    \setlength{\topsep}{0pt}%
    \setlength{\leftmargin}{#1}%
    \setlength{\rightmargin}{#2}%
    \setlength{\listparindent}{\parindent}%
    \setlength{\itemindent}{\parindent}%
    \setlength{\parsep}{\parskip}%
  }%
  \item[]}{\end{list}}
%    \end{macrocode}
% \end{macro}

% \subsection{The bibliography and references}
% \label{sec:biblio}
% \begin{macro}{\beforebiblioskip}
% \begin{macro}{\afterbiblioskip}
% I'm not entirely sure why I introduced these two commands, since
% directly putting |\vspace| before and after
% |\begin{thebibliography}| seems to work well enough. 
%    \begin{macrocode}
\newlength{\beforebiblioskip}
\setlength{\beforebiblioskip}{0pt}
\newlength{\afterbiblioskip}
\setlength{\afterbiblioskip}{0pt}
%    \end{macrocode}
% \end{macro}
% \end{macro}
% \begin{macro}{thebibliography}
% The environment itself is done as starred section with extra
% vertical space after heading. This space is currently zero. 
% In the book style we do the cunning trick with rereading the file in
% order to get the references into the |\ADS@references| macro.
% 
% \changes{3.05}{2003/01/02}{Removed the \cs{beforebiblioskip} and
% \cs{afterbiblioskip} commands, since they cause unwanted vertical
% space in some instances.}  
%    \begin{macrocode}
\renewenvironment{thebibliography}{%
  \if@book\RereadBiblio\fi
%%  \vspace*{\beforebiblioskip}
  \section*{REFERENCES} \frenchspacing
%%  \vspace*{\afterbiblioskip}
  \small
%    \end{macrocode} 
% The references themselves are set as a list in order 
% to get the hanging indentation. They are done in a smaller font and
% with french spacing so we needn't worry about the distinction
% between sentence-ending periods and abbreviations. 
%    \begin{macrocode} 
  \list{\null}{\leftmargin 1.5em\labelwidth\z@\labelsep\z@\itemindent -1.5em
    \itemsep\z@\parsep\z@\usecounter{enumi}}
  \sloppy\clubpenalty4000\widowpenalty4000
  \sfcode`\.=1000\relax\@biblisttrue}
{\endlist\@biblistfalse\normalsize\@outputaddressesifnecessary}
%    \end{macrocode}
% \end{macro}
% \subsection{Uniquification of cross-references}
% \label{sec:uniquification}
% 
% \begin{macro}{\articlename}
% In order that the |\label{}| command in individual articles be
% guaranteed to give a unique anchor, we append the article name (held
% in the |\articlename| command) to each key. This requires
% redefinition of the following commands: |\cite|, |\nocite|, |\ref|,
% |\pageref|, |\label|. 
% \begin{macro}{\oldcite}
% \begin{macro}{\oldnocite}
% \begin{macro}{\oldref}
% \begin{macro}{\oldpageref}
% \begin{macro}{\oldlabel}
% The original versions of these commands are
% saved as |\oldcite|, |\oldnocite|, etc. 
%
% The uniquifying of cross-references is now always done, so 
% we don't get multiply-defined labels in the single article case
%    \begin{macrocode}
\newcommand{\articlename}{anon}
\AtBeginDocument{
  \let\oldcite=\cite 
  \let\oldnocite=\nocite 
  \let\oldref=\ref 
  \let\oldpageref=\pageref
  \let\oldlabel=\label
}
%    \end{macrocode}
% \end{macro}
% \end{macro}
% \end{macro}
% \end{macro}
% \end{macro}
% \end{macro}
%
% Uniquify all cross-references by prepending the article name.  This
% now has to go into |\AtBeginDocument| so it is working on the new
% versions that hyperref.sty has set up.  Then we need to put the
% |\@begofarticlehook| afterwards, so that the first page label is using
% the new definitions.  If amsmath is loaded, then we also need to patch
% the special label command that they use. \changes{4.0}{2019/11/18}{Fix labels for amsmath display equations, in case that package has been loaded.}
%    \begin{macrocode}
\AtBeginDocument{
  \renewcommand{\ref}[1]{\oldref{\articlename:#1}%
    \if@debug{\ClassInfo{rmaa}{Reference to \articlename:#1}}\fi}
  \renewcommand{\pageref}[1]{\oldpageref{\articlename:#1}%
    \if@debug{\ClassInfo{rmaa}{Page reference to \articlename:#1}}\fi}
  \renewcommand{\label}[1]{\oldlabel{\articlename:#1}%
    \if@debug{\ClassInfo{rmaa}{Label for \articlename:#1}}\fi}%
  \ifdef{\label@in@display}{%
    \let\oldlabel@in@display=\label@in@display
    \renewcommand{\label@in@display}[1]{\oldlabel@in@display{\articlename:#1}%
      \if@debug{\ClassInfo{rmaa}{AMSTeX Label for \articlename:#1}}\fi}%
  }{}
  \@begofarticlehook
}%
%    \end{macrocode}
% Citations are more complicated since they can contain a list of
% keys. These are really only here for use with the rmaacite package,
% which is utterly deprecated (02 Jan 2003) now that I have the natbib
% code working. The command \cs{cite} should \emph{not} be used with
% natbib: use \cs{citep}, \cs{citet}, etc.\@ instead. 
%    \begin{macrocode}
\renewcommand{\nocite}[1]{\@bsphack
  \@for\@citeb:=#1\do{%
    \edef\@citeb{\expandafter\@firstofone\@citeb:\articlename}%
    \if@filesw\immediate\write\@auxout{\string\citation{\@citeb}}\fi
    \@ifundefined{b@\@citeb:\articlename}{\G@refundefinedtrue
      \@latex@warning{Citation `\@citeb' undefined}}{}}%
  \@esphack}
\renewcommand{\cite}[1]{\oldcite{#1:\articlename}}
%   \end{macrocode}


% Horrible fiddling with the bibliography citation mechanism. 
% This has to be globally available in case rmaacite puts it in the
% \texttt{.aux} file. 
%    \begin{macrocode}
\def\lbibcite#1#2#3{\global\@namedef{b@#1}{#2}
                    \global\@namedef{by@#1}{#3}}
%    \end{macrocode}

% \begin{macro}{\oldbibitem}
% \begin{macro}{\@rmaaBibArg}
% \begin{macro}{\@rmaaSpace}
% \begin{macro}{rmaabib}
% Of course, |\bibitem| needs to get the uniquification treatment as
% well. While we are at it, we redefine it to use consecutive integers
% for the key if the user leaves it blank (using the counter
% |rmaabib|). This allows one to put 
% |\bibitem{}| for all your references without getting annoying
% messages about multiply-defined labels.
% 
% We also take the opportunity to insert the optional argument to
% \cs{bibitem} if it is missing in the form Author(DATE) so that natbib
% doesn't try to do numerical citations. Natbib already provides a
% hook \cs{@extra@b@citeb} to help uniquify the citation key, so we
% just plug into that. However, it puts it \emph{after} the original
% key, not before as I had had it. Easiest solution is to change the
% way that I do it. 02 Jan 2003. Note that the \cs{@extra@b@citeb}
% does not get written to the .aux  file in the \cs{citation} command
% but it does work. 
% \changes{3.14}{06 Feb 2003}{Reinstalled old definition of \cs{bibitem}
%  for the case where the nonatbib option is requested, so that the
%  rmaacite package will work. This means that you can't use rmaacite
%  and natbib in the same volume.}
%    \begin{macrocode}
\if@natbib\def\@extra@b@citeb{:\articlename}\fi
%%\else
\let\oldbibitem\bibitem
\newcommand{\@rmaaBibArg}{}
\newcommand{\@rmaaSpace}{ }
\newcounter{rmaabib}
\if@natbib
\renewcommand{\bibitem}[2][ABC(2003)]{%
  \renewcommand{\@rmaaBibArg}{#2}
  \ifx\@rmaaBibArg\@rmaaEmpty
     \addtocounter{rmaabib}{1}
     \oldbibitem[#1]{\thermaabib:\articlename}
  \else\ifx\@rmaaBibArg\@rmaaSpace
     \addtocounter{rmaabib}{1}
     \oldbibitem[#1]{\thermaabib:\articlename}
  \else 
     \oldbibitem[#1]{#2:\articlename}
  \fi\fi
  \renewcommand{\@rmaaBibArg}{}
  }
\else
\renewcommand{\bibitem}[1]{%
  \renewcommand{\@rmaaBibArg}{#1}
  \ifx\@rmaaBibArg\@rmaaEmpty
     \addtocounter{rmaabib}{1}
     \oldbibitem{\articlename:\thermaabib}
  \else\ifx\@rmaaBibArg\@rmaaSpace
     \addtocounter{rmaabib}{1}
     \oldbibitem{\articlename:\thermaabib}
  \else
     \oldbibitem{\articlename:#1}
  \fi\fi
  \renewcommand{\@rmaaBibArg}{}
  }
\fi
%%\fi
%    \end{macrocode}
% \end{macro}
% \end{macro}
% \end{macro}
% \end{macro}

% \begin{macro}{oldpcite}
% \begin{macro}{oldscite}
% Then, we need to repeat the above scheme, but for the |rmaacite|
% definition. 
%    \begin{macrocode}
\newcommand{\@redefinebibitem}{%
  \let\oldbibitem\bibitem
  \let\oldcite\cite
  \let\oldpcite\pcite
  \let\oldscite\scite
  \def\bibitem[##1<##2>]##3{%
    \renewcommand{\@rmaaBibArg}{##3}
    \ifx\@rmaaBibArg\@rmaaEmpty
       \addtocounter{rmaabib}{1}
       \oldbibitem[##1<##2>]{\thermaabib:\articlename}
    \else\ifx\@rmaaBibArg\@rmaaSpace
       \addtocounter{rmaabib}{1}
       \oldbibitem[##1<##2>]{\thermaabib:\articlename}
    \else 
       \oldbibitem[##1<##2>]{##3:\articlename}
    % trust the author here
%%       \oldbibitem[##1<##2>]{##3}
    \fi\fi
    }
  }
%    \end{macrocode}
% \end{macro}
% \end{macro}


% \subsection{HTML Table of Contents}
% \label{sec:hyper-TOC}
% \begin{macro}{\HyperTOCEntry}
% \begin{macro}{\HyperTOCpart}
% \begin{macro}{\HyperTOCmeeting}
% \begin{macro}{\HyperTOCeditors}
% \begin{macro}{\HyperTOCpreamble}
% These commands are used to write information to the \texttt{.htoc}
% file, which is used in generating a Web Page of the Table of
% Contents. Since they only ever appear preceded by |\protect| in the
% argument to an |\addtocontents| command,  they should never need to
% be expanded by the macros. However, we define them here just in case
% \dots
%    \begin{macrocode}
\newcommand{\HyperTOCEntry}[3]{\relax}
\newcommand{\HyperTOCpart}[1]{\relax}
\newcommand{\HyperTOCmeeting}[3]{\relax}
\newcommand{\HyperTOCeditors}[1]{\relax}
\newcommand{\HyperTOCpreamble}{\relax}
%    \end{macrocode}
% \end{macro}
% \end{macro}
% \end{macro}
% \end{macro}
% \end{macro}
%

% \subsection{Formatting the poster abstracts} 
% These can now be used without using the book option.
%    \begin{macrocode}

\newcommand{\@makeminititle}{%
  \newpage
  \null 
  \vskip 3em%
  \begin{center}%
    {\@title \par}%
  \end{center}
  \vskip 3em%
  }

\newif\if@abstractsads
\@abstractsadsfalse
\newcommand\abstractsADS{\@abstractsadstrue}

\if@book\RequirePackage{multicol}\fi
\newenvironment{abstracts}[1][Abstracts of Contributed Papers]{%
  %% First, make sure we start on an odd page (06 Jun 2000)
  \if@book\clearpage\thispagestyle{empty}\cleardoublepage\fi
  %%
  \renewcommand{\articlename}{ABSTRACTS}
  \title{#1}\shortauthor{Abstracts}\shorttitle{Abstracts}
  \listofauthors{}
  \if@book
  \addcontentsline{toc}{chapter}{\protect\@fiddleauthors{\@TOCtitle}{\rmaaTOCauthors}}%
  \addtocontents{htc}{\protect\HyperTOCEntry{\@rawtitle}{\@rawauthors}{\articlename}}
  \fi
  \pagestyle{rmaa} % this now seems necessary (06 Jun 2000)
  \thispagestyle{rmaafirstpage}
%%  \twocolumn[\@makeminititle]
  \renewcommand{\abstract}[1]{%
    \renewcommand\ADS@abstract{##1}%
    \renewcommand{\@abstract}{##1}\@writeabstract%
  }
  \renewcommand{\title}[2][]{%
    \gdef\ADS@authors{}%
    \gdef\ADS@affil{}%
    \gdef\ADS@references{}%
    \gdef\ADS@keywords{}%
    \renewcommand{\@rmaaArg}{##1}
    \ifx\@rmaaArg\@rmaaEmpty
       \renewcommand{\@TOCaftertitle}{}
    \else
       \renewcommand{\@TOCaftertitle}{ (\@rmaaArg)}
    \fi
    \renewcommand{\@title}{\rmaa@TOCTitleStyle{##2}}%
    \renewcommand{\@TOCtitle}{\rmaa@TOCTitleStyle{##2}\@TOCaftertitle}%
    \renewcommand{\rmaaIndexTitle}{\rmaa@IndexTitleStyle{##2\@TOCaftertitle.}}%
    \renewcommand\ADS@title{##2}%
    \xdef\ADS@firstpage{\@arabic{\c@page}}%
    }
  \renewcommand{\author}[1]{\let\@thanks\@empty\renewcommand\@author{##1} }
  \renewcommand{\footnotetext}[2][XX]{\par\noindent\textsuperscript{##1}~##2}
  \renewcommand{\footnoterule}{\kern-3\p@ \hrule \@width 3em \kern 2.6\p@}
  \if@book\@begofarticlehook\fi
  \if@book\begin{multicols}{2}[\@makeminititle]\else\@makeminititle\fi
  \begingroup
  \setcounter{footnote}{0}
  }% and at the end of the environment....
  {\endgroup\if@book\end{multicols}\listofabstracts\label{RMAALastPage}\clearpage\fi}
%% auxiliary commands used:
%% \@writeabstract does all the actual formatting 
\newcommand{\@writeabstract}{%
  \if@book%
  \addcontentsline{loa}{abstract}{\protect\@fiddleauthors{\@TOCtitle}{\rmaaTOCauthors}}%
  \fi%
  \begin{center}\@title\par\@author\end{center}
  \@abstract
  \par\noindent\footnoterule\nobreak\footnotesize\@thanks
  \par\noindent\hrulefill\goodbreak
  \vskip 1cm \@plus0.5cm
  %% write ADS entry
  \if@book\if@abstractsads\ADS@dumpdata\fi\fi%
  \endgroup % this pair is to make each abstract a self-contained group
  \begingroup
  \setcounter{footnote}{0}
  }

\if@book
%% List of abstracts (stuff written to .loa file)
\newcommand{\listofabstracts}{%
  \begin{multicols}{2}[%
    \vskip 1.5cm \vfill % leave at least 1.5 cm and go to bottom of page if possible
    \centerline{\MakeUpperCase{List of Abstracts}}
    \vskip 1.5cm \@plus0.5cm
    ][6cm] % this makes it start a new page if there is not at least 6cm of space
    \renewcommand{\@fiddleauthors}[2]{##1\hskip1em plus2em minus0.4em\penalty-100  ##2}
    \parindent\z@
%%    \parskip.5\p@ \@plus .7\p@\relax 
    \parskip3\p@ \@plus 1\p@ \@minus 1\p@\relax % allow quite a lot of
                                % leeway here to avoid breaks inside
                                % entries 
    \renewcommand{\\}{\relax} % temporarily turn off explicit line breaks just in case
    \renewcommand{\altaffilmark}[1]{\relax} % ditto with this
    \@starttoc{loa}
  \end{multicols}
  }
\newcommand*{\l@abstract}[2]{%
  \par\penalty-1000
  #1\hskip 1em\penalty0\null\nobreak\hfill#2}
\fi

%    \end{macrocode}
 
% \subsection{Code for typesetting an entire volume of conference
% proceedings} \label{sec:whole-volume}
%    \begin{macrocode}

\if@book
%% environment for including the individual articles
\newcommand{\@fiddleauthors}[2]{\relax}
\renewcommand{\theequation}{\@arabic\c@equation}
\newenvironment{articles}%
{ \begingroup % make it a group so that author definitions do not
              % spill over
  \gdef\ADS@authors{}
  \pagestyle{rmaa}
  \renewenvironment{document}{%
    \@begofarticlehook
    \addtocounter{chapter}{1}%
    \addcontentsline{toc}{chapter}{%
      \protect\@fiddleauthors{\@TOCtitle}{\rmaaTOCauthors}}%
    \addtocontents{htc}{\protect\HyperTOCEntry{\@rawtitle}{\@rawauthors}{\articlename}}
    }{\relax}
  \renewcommand{\documentclass}[2][foo]{\relax} 
  \renewcommand{\usepackage}[2][foo]{%
    \def\@packagename{##2}\def\@rmaacitename{rmaacite}%
    \ifx\@packagename\@rmaacitename
    \@rmcitetrue\ClassInfo{rmaa}{This paper uses rmaacite}
    \makeatletter\input{rmaacite.sty}\makeatother % read it in here so that all changes
    \else                                         % are local to paper
    \if@debug\ClassWarning{rmaa}{Package \@packagename\space was requested by \articlename\space}\fi
    \fi} 
  %% Don't allow individual articles to mess with the pages, volume or year
  \renewcommand{\SetFirstPage}[1]{\relax}
  \renewcommand{\SetYear}[1]{\relax}
  \renewcommand{\SetVolume}[1]{\relax}
  \renewcommand{\SetConfTitle}[1]{\relax}
  %% Don't allow individual articles to change these
  \renewcommand{\nonstopmode}{\relax}
  \renewcommand{\scrollmode}{\relax}
  \renewcommand{\outputplates}{\relax}
  %% XXXX need to check if using rmaacite and act accordingly
  %% avoid multiple declarations of the same length
  \let\oldnewlength=\newlength
  \renewcommand{\newlength}[1]{%
    \edef\reserved@a{\expandafter\@gobble\string ##1}%
    \@ifundefined{\reserved@a}{%
      \oldnewlength{##1}}{%
      \if@debug{\ClassInfo{rmaa}{\reserved@a\space already defined: doing nothing}}\fi%
      \relax}}%
  \renewcommand\newcommand{\providecommand}
  }%
{\endgroup\onecolumn}
     

%    \end{macrocode}
% Command to be used in `articles' environment for including an
% article
%    \begin{macrocode}
\newcommand{\article}[2][]{%
  % optional first argument is extra text to put in brackets after title in TOC
  % e.g. ``Invited Talk'' or ``Poster''
  \renewcommand{\articlename}{#2}
  % reset loads of counters
  \setcounter{footnote}{0}
  \setcounter{figure}{0}
  \setcounter{table}{0}
  \setcounter{section}{0}
  \setcounter{equation}{0}
  \setcounter{rmaabib}{0}
  \setcounter{author}{0}
  \setcounter{AUTHOR}{1}
  \gdef\ADS@authors{}%
  \gdef\ADS@affil{}%
  \gdef\ADS@affiliations{}%
  \gdef\ADS@references{}%
  \@firstauthortrue
  \@firstaffiltrue
  \@firstaltaffiltrue
  \@firstkwdtrue
  \@rmcitefalse
  \if@suppressfulladdresses\else\@needaddressestrue\fi
  \renewcommand{\@rmaaArg}{#1}
  \ifx\@rmaaArg\@rmaaEmpty
     \renewcommand{\@TOCaftertitle}{}
  \else
     \renewcommand{\@TOCaftertitle}{ (#1)}
  \fi
%    \end{macrocode}
% In the main journal articles must start on an odd page. 
%    \begin{macrocode}
  \if@journal\clearpage\thispagestyle{empty}\cleardoublepage\fi
%% Now input the file, then run the end-of-article hook (reinstated 28
  %% Sep 2001)
  {\input{\@articledir/\articlename.tex}\@endofarticlehook}}
%    \end{macrocode}
% Command to separate articles into sessions. Write an entry in the
% TOC, makes sure we start on an odd page (added 06 Jun 2000), then do
% a title page if \cs{SessionTitlePages} was called. 
% \changes{3.12}{04 Feb 2003}{Added option to have title page for each
% session. Turn it on with \cs{SessionTitlePages}.}
%    \begin{macrocode}
\newcommand{\session}[1]{
  %% this does nothing except write an entry in the TOC
  \addcontentsline{toc}{part}{\rmaa@TOCPartStyle{#1}}
  \addtocontents{htc}{\protect\HyperTOCpart{#1}}
  %% and also makes sure we start on an odd page (added 06 Jun 2000)
  \clearpage\thispagestyle{empty}\cleardoublepage
  \stepcounter{part}
  \if@sessiontitlepages{%
    \onecolumn
    \thispagestyle{empty}
    \renewcommand{\articlename}{Section\thepart TitlePage}
    \label{RMAAFirstPage}
    \make@session@title{#1}
    \label{RMAALastPage}
    \clearpage\thispagestyle{empty}\cleardoublepage
  }\fi
  }
\newcommand{\make@session@title}[1]{%
  {%
    \setlength\parindent{0pt}
    \vspace*{\stretch{1}}
    \rule{\linewidth}{1mm}
    \begin{flushright}\Huge\bfseries\boldmath
      #1
    \end{flushright}
    \rule{\linewidth}{1mm}
    \vspace*{\stretch{2}}
  }
}
%    \end{macrocode}
% Command to turn on separate title pages for each session. 
%    \begin{macrocode}
\newcommand{\SessionTitlePages}{\@sessiontitlepagestrue}
\newcommand{\BoldSessionTitles}{%
  \renewcommand{\rmaa@TOCPartStyle}[1]{\textbf{\rmaa@TOCTitleStyle{##1}}}
  }
\newcommand{\ArticleDir}[1]{\renewcommand{\@articledir}{#1/}\graphicspath{{1/}}}
\newcommand{\@articledir}{.}
%\newcommand{\ps@rmaafront}{\ps@rmaa}% page style for front matter (same as ever, actually)
\newcommand{\ps@rmaafront}{% page style for front matter:  page # in each foot
  \renewcommand\@oddfoot{\hfil{\thepage}}%
  \renewcommand\@evenfoot{{\thepage}\hfil}%
  % shorttitle in both heads
  \renewcommand\@evenhead{\hfil\@shorttitle\hfil}%
  \renewcommand\@oddhead{\@evenhead}
  }
\newcommand{\ps@rmaafrontfirstpage}{% page style for first page of front matter
  % page # in each foot, on outside edge
  \renewcommand\@oddfoot{\hfil\thepage}%
  \renewcommand\@evenfoot{\thepage\hfil}%
  % nothing in the even heads
  \renewcommand\@evenhead{\relax}%
  \renewcommand\@oddhead{\relax}%
  }
%% slight modification to the standard \frontmatter and \mainmatter
%% commands, just to set the \pagestyle
\renewcommand\frontmatter{%
  \latexonecolumn
  \cleardoublepage
  \@frontmattertrue
  \@mainmatterfalse
  \pagenumbering{roman}
  \pagestyle{rmaafront}
  \addtocounter{page}{2} % leave space for the title page
  }
\renewcommand\mainmatter{%
  \pagestyle{empty}
  \if@rmtwocol\twocolumn\fi
  \cleardoublepage
  \@frontmatterfalse
  \@mainmattertrue
  \pagenumbering{arabic}
  \pagestyle{rmaa}
  }
\renewcommand\backmatter{%
  \clearpage\pagestyle{empty}\cleardoublepage
  \@mainmatterfalse
  %% put extra space in the TOC
  \addtocontents{toc}{\bigskip}
  }
%    \end{macrocode}
%
% \subsection{Front sections}\label{sec:front-sections}
% \begin{macro}{\notitle}
% \begin{macro}{\shortname}
% \begin{macro}{frontsection}
% The environment \texttt{frontsection} is for doing a section in the
% front matter (preface, list of participants, or whatever). It has 
% one optional argument:
% \begin{enumerate}
% \item[\#1] Short title (no embedded spaces) for
% use in making PS files. E.g., ``list\_of\_participants''.   
% \end{enumerate}
% and two mandatory arguments: 
% \begin{enumerate}
% \item[\#2] Title of section (e.g., `Preface')
% \item[\#3] Commands to be evaluated before processing the contents
% of the environment. This is mainly useful for putting in the things
% that would go in the preamble of a normal article  (e.g.,
% |\listofauthors|, |\indexauthor|, |\notitle|) 
% \end{enumerate}
%    \begin{macrocode}
\newcommand{\notitle}{\@titlevisiblefalse} \newcommand{\shortname}{}
\newenvironment{frontsection}[3][]{%
  \begingroup
  \renewcommand{\shortname}{#1}
  \ifx\shortname\@rmaaEmpty
     \renewcommand{\articlename}{#2}
  \else
     \renewcommand{\articlename}{\shortname}
  \fi
  \@begofarticlehook
  \@titlevisibletrue % default
  \thispagestyle{rmaafrontfirstpage}
  \listofauthors{} % blank out this in case it is not specified
  \title{{#2}}\shorttitle{#2}\shortauthor{#2}
  #3 % evaluate the command given as second mandatory arg
  \addcontentsline{toc}{chapter}{\protect\@fiddleauthors{\@TOCtitle}{\rmaaTOCauthors}}%
  \addtocontents{htc}{\protect\HyperTOCEntry{\@rawtitle}{\@rawauthors}{\articlename}}
  \if@titlevisible
     \vspace*{1cm}\centerline{\rmaa@FrontSecStyle{\@title}}\vspace*{1cm}
  \fi
  \setlength{\parskip}{3ex}
  }{\frontlabel{RMAALastPage}\clearpage\endgroup}
%    \end{macrocode}
% \end{macro}
% \end{macro}
% \end{macro}
%    \begin{macrocode}
%% Definition of the command \frontlabel, which rights out something like `F4' to
%% the .aux file, instead of `\textlatin  {i}'. This is so it is easy to post-process
%% the file for making PS files of individual sections
\newcommand{\frontlabel}[1]{%
  \@bsphack
  \protected@write\@auxout{}%
     {\string\newlabel{\articlename:#1}{{}{F\@arabic{\c@page}}}}
  \@esphack
  }
%% User command to say X, Y, Z,\\ Editors of the Proceedings
\newcommand{\EditorSignature}{\vspace*{\parskip}\rightline{\@editors}%
  \rightline{Editors of the proceedings}}
%% Fiddle with \thesection, etc. so that they don't include the chapter number
%% Basically go back to the `article' definitions of these
\renewcommand{\thesection}{\@arabic\c@section}
\renewcommand{\thefigure}{\@arabic\c@figure}
\renewcommand{\thetable}{\@arabic\c@table}

\newcommand{\@oldargs}{}
\newcommand{\@newargs}{}
%% add TOC entry for each paper
%\renewcommand{\title}[1]%
%{\renewcommand{\@title}{\rmaa@TitleStyle{#1}}%
%\addtocounter{chapter}{1}\addcontentsline{toc}{chapter}{#1}%
%}

%% Formatting of the table of contents
%% dimensions of various parts
\newlength{\TOCtextwidth}\setlength{\TOCtextwidth}{12.65cm}
\newlength{\TOCpagenumwidth}\setlength{\TOCpagenumwidth}{1.3cm}
\newlength{\TOCentrywidth}\setlength{\TOCentrywidth}{\TOCtextwidth}
\addtolength{\TOCentrywidth}{-\TOCpagenumwidth} 
\newlength{\TOCextramargin}\setlength{\TOCextramargin}{1.25cm}
\newcommand{\TOCbox}[1]{\hspace{\TOCextramargin}\parbox{\TOCtextwidth}{#1}\par}
\newlength{\TOCextramarginx}\setlength{\TOCextramarginx}{\TOCextramargin}\addtolength{\TOCextramarginx}{4.3cm}
\newlength{\TOCtextwidthx}\setlength{\TOCtextwidthx}{\TOCtextwidth}\addtolength{\TOCtextwidthx}{-4.3cm}
\newcommand{\TOCboxx}[1]{\hspace{\TOCextramarginx}\parbox{\TOCtextwidthx}{#1}\par}

\newcommand{\RMAAbannerimage}{rmsc}
\renewcommand\tableofcontents{%
  \renewcommand{\articlename}{CONTENTS}
  \begingroup
  \@begofarticlehook
  \if@twocolumn
     \@restonecoltrue\onecolumn
  \else
     \@restonecolfalse
  \fi
  %% 9pt for journal style
  \if@proceedings\else\small\fi
  \if@newstyle\small\fi %% same for new style
  \if@proceedings\thispagestyle{rmaafrontfirstpage}\else\thispagestyle{empty}\fi
  \shortauthor{Contents}\shorttitle{Contents}
  \vspace*{18mm}
  \if@printbanner
  %% put in the banner header (need to get better image of this)
     \TOCbox{\hspace*{-0.15cm}\includegraphics[width=0.9\TOCtextwidth]{\RMAAbannerimage}}
     \vspace*{8mm}
  %% volume and publication date (should really be in funny font)
     \TOCboxx{\@volumen\\[2mm] \@mes\ \@anyo}
  \else
  %% put in the banner header (need to get better image of this)
     \TOCbox{\hspace*{-0.15cm}\includegraphics[width=0.9\TOCtextwidth]{\RMAAbannerimage}}
     \vspace*{12.5mm}
  %% volume and publication date (should really be in funny font)
     \TOCboxx{~\\[2mm] ~}
  \fi
  \vspace*{16mm}
  %% This is the command that fiddles with the TOC entry so that the
  %% title is flushleft but the authors are flushright
  \renewcommand{\@fiddleauthors}[2]{%
    % first put the paper title
    \renewcommand{\\}{\relax} % temporarily turn off explicit line breaks just in case
    \renewcommand{\altaffilmark}[1]{\relax} % ditto with this
    \renewcommand{\textsuperscript}[1]{\relax} % ditto with this 
    ##1
    % now save the author list in a box
    \sbox\@tempboxa{##2}%
    % check that box is shorter than the total width available for the entry
    \ifdim \wd\@tempboxa <\TOCentrywidth
       % if so, then put it right here, flush against right margin
       \penalty10000\hskip1em plus10em\penalty0\hbox{}\nobreak\hfill\box\@tempboxa
       % Note that we force a space of at least 1em between title and authorlist, 
       % stretchable to 11em since we may have very short titles but still with an 
       % authorlist too long to go on the same line.  The only linebreak allowed is 
       % directly before the authorlist.
    \else
       % if it is too long, we set the original authorlist text in its own parbox
       \par\vskip1ex\parbox[b]{\TOCentrywidth}{\raggedleft ##2}
    \fi
    }
  \@starttoc{toc}%
  \if@restonecol\twocolumn\fi
  \frontlabel{RMAALastPage}
  \clearpage
  \thispagestyle{empty}
  \cleardoublepage % make sure we start on an odd page
  \endgroup

  %%% ALSO, do the HTML TOC
  %% The file that will be used to make the HTML TOC
  \@starttoc{htc}
  %% All the stuff in this file does nothing in the book style
  \addtocontents{htc}{\protect\HyperTOCmeeting{\@conftitle}{\@reunionplace}{\@reuniondate}}
  \addtocontents{htc}{\protect\HyperTOCeditors{\@editors}}
  \addtocontents{htc}{\protect\HyperTOCpreamble}
  \addtocontents{htc}{\protect\HyperTOCEntry{Table of Contents}{}{\articlename}}

  %% and the ADS file
  \ADS@openfile
  }
%% This is now much simpler
\renewcommand\l@chapter[2]{%TOC entry for articles
  \ifnum \c@tocdepth >\m@ne
    \vskip 1.0em \@plus\p@
    \TOCbox{%
      \parbox[b]{\TOCentrywidth}{#1}\parbox[b]{\TOCpagenumwidth}{\hfill#2}
    }
  \fi} 
\renewcommand\l@part[2]{%TOC entry for sessions: centered, no page number
  \ifnum \c@tocdepth >\m@ne
    \vskip 1.5em \@plus1.5em \goodbreak
    \TOCbox{%
      \parbox[b]{\TOCentrywidth}{\hfil#1\hfil}\hfill
    }
  \fi\nobreak} 

%% Formatting theindex
%% strongly encourage page breaks between index entries
\renewcommand{\@idxitem}{\par\penalty-1000\hangindent.45cm} 
\renewcommand{\indexspace}{\par\penalty-1000\vskip 10\p@ \@plus15\p@ \@minus3\p@\relax}
\renewenvironment{theindex}{%
  %% first, make sure we start on an odd page (06 Jun 2000)
  \clearpage\thispagestyle{empty}\cleardoublepage
  %% this looks a bit like an article but is two columns and without abstract, etc
  \renewcommand{\articlename}{INDEX}
  \title{Author Index}\shortauthor{Author Index}\shorttitle{Author Index}
  \listofauthors{}
  \addcontentsline{toc}{chapter}{\protect\@fiddleauthors{\@TOCtitle}{\rmaaTOCauthors}}%
  \addtocontents{htc}{\protect\HyperTOCEntry{\@rawtitle}{\@rawauthors}{\articlename}}
  \pagestyle{rmaa}
  \thispagestyle{rmaafirstpage}
  \@begofarticlehook
  %% 9pt for journal style
  \if@proceedings\else\small\fi
  \if@newstyle\small\fi
  \twocolumn[\@makeminititle]
  %% vertical spacing (see also definition of \@idxitem and \indexspace
  \parskip\z@ \@plus .6\p@\relax % allow quite a lot of leaway here to avoid breaks inside entries
  %% horizontal spacing, breaking and glue parameters (see also the stuff in the definition 
  %% of \indexauthor
  \parindent\z@ 
  \exhyphenpenalty10000 % don't break on explicit hyphens, it always looks bad
  \rightskip=0pt plus.2em % allow a tiny bit of leeway at the ends of lines
  \spaceskip.3333em plus.5em minus.15em % a lot of leeway between words
  \xspaceskip.5em plus.5em minus.4em % this seems to have no effect
  \let\item\@idxitem
  \renewcommand{\\}{\relax} % temporarily turn off explicit line breaks just in case
  \renewcommand{\altaffilmark}[1]{\relax} % ditto with this
  \renewcommand{\textsuperscript}[1]{\relax} % ditto with this
  }{\label{RMAALastPage}\clearpage} % clearpage necessary so that we know about \shortauthor



%% end definitions for book styles
\fi

%    \end{macrocode}
% In order to do away with plates, we sometimes want a figure to
% appear in color only in the online version. The following command
% can bve used in a figure caption to typeset a note to this effect. 
%    \begin{macrocode}
\newcommand\ColorOnlineOnly{%
  \textsc{Note: this figure is available in color in the electronic
    version of this article, obtainable from}
  \texttt{http://\discretionary{}{}{}www.\discretionary{}{}{}astroscu.\discretionary{}{}{}unam.\discretionary{}{}{}mx/\discretionary{}{}{}\~{}rmaa/}.}
%% Plates
%% Define 2 counters: plate is the number in the whole document, plateloop is used 
%% internally by \outputplates
\newcounter{plate}
\newcounter{plateloop}
%% Currently limited to 15 plates in document
%% when we do a label, we need to write to \newlabel entries to the .aux file, one 
%% that says RMAAPlateX so that it can be referred to in \ouputplates, and one that
%% says articlename:RMAAPlateX so that the dvips program can find it and asociate it
%% with the right article
\newcommand{\@platelabeli}{\oldlabel{RMAAPlate1}\label{RMAAPlate1}}
\newcommand{\@platelabelii}{\oldlabel{RMAAPlate2}\label{RMAAPlate2}}
\newcommand{\@platelabeliii}{\oldlabel{RMAAPlate3}\label{RMAAPlate3}}
\newcommand{\@platelabeliv}{\oldlabel{RMAAPlate4}\label{RMAAPlate4}}
\newcommand{\@platelabelv}{\oldlabel{RMAAPlate5}\label{RMAAPlate5}}
\newcommand{\@platelabelvi}{\oldlabel{RMAAPlate6}\label{RMAAPlate6}}
\newcommand{\@platelabelvii}{\oldlabel{RMAAPlate7}\label{RMAAPlate7}}
\newcommand{\@platelabelviii}{\oldlabel{RMAAPlate8}\label{RMAAPlate8}}
\newcommand{\@platelabelix}{\oldlabel{RMAAPlate9}\label{RMAAPlate9}}
\newcommand{\@platelabelx}{\oldlabel{RMAAPlate10}\label{RMAAPlate10}}
\newcommand{\@platelabelxi}{\oldlabel{RMAAPlate11}\label{RMAAPlate11}}
\newcommand{\@platelabelxii}{\oldlabel{RMAAPlate12}\label{RMAAPlate12}}
\newcommand{\@platelabelxiii}{\oldlabel{RMAAPlate13}\label{RMAAPlate13}}
\newcommand{\@platelabelxiv}{\oldlabel{RMAAPlate14}\label{RMAAPlate14}}
\newcommand{\@platelabelxv}{\oldlabel{RMAAPlate15}\label{RMAAPlate15}}
\newcommand{\@outputplatelabeli}{\oldlabel{RMAAOutputPlate1}}
\newcommand{\@outputplatelabelii}{\oldlabel{RMAAOutputPlate2}}
\newcommand{\@outputplatelabeliii}{\oldlabel{RMAAOutputPlate3}}
\newcommand{\@outputplatelabeliv}{\oldlabel{RMAAOutputPlate4}}
\newcommand{\@outputplatelabelv}{\oldlabel{RMAAOutputPlate5}}
\newcommand{\@outputplatelabelvi}{\oldlabel{RMAAOutputPlate6}}
\newcommand{\@outputplatelabelvii}{\oldlabel{RMAAOutputPlate7}}
\newcommand{\@outputplatelabelviii}{\oldlabel{RMAAOutputPlate8}}
\newcommand{\@outputplatelabelix}{\oldlabel{RMAAOutputPlate9}}
\newcommand{\@outputplatelabelx}{\oldlabel{RMAAOutputPlate10}}
\newcommand{\@outputplatelabelxi}{\oldlabel{RMAAOutputPlate11}}
\newcommand{\@outputplatelabelxii}{\oldlabel{RMAAOutputPlate12}}
\newcommand{\@outputplatelabelxiii}{\oldlabel{RMAAOutputPlate13}}
\newcommand{\@outputplatelabelxiv}{\oldlabel{RMAAOutputPlate14}}
\newcommand{\@outputplatelabelxv}{\oldlabel{RMAAOutputPlate15}}
\newcommand{\@platerefi}{\oldpageref{RMAAPlate1}}
\newcommand{\@platerefii}{\oldpageref{RMAAPlate2}}
\newcommand{\@platerefiii}{\oldpageref{RMAAPlate3}}
\newcommand{\@platerefiv}{\oldpageref{RMAAPlate4}}
\newcommand{\@platerefv}{\oldpageref{RMAAPlate5}}
\newcommand{\@platerefvi}{\oldpageref{RMAAPlate6}}
\newcommand{\@platerefvii}{\oldpageref{RMAAPlate7}}
\newcommand{\@platerefviii}{\oldpageref{RMAAPlate8}}
\newcommand{\@platerefix}{\oldpageref{RMAAPlate9}}
\newcommand{\@platerefx}{\oldpageref{RMAAPlate10}}
\newcommand{\@platerefxi}{\oldpageref{RMAAPlate11}}
\newcommand{\@platerefxii}{\oldpageref{RMAAPlate12}}
\newcommand{\@platerefxiii}{\oldpageref{RMAAPlate13}}
\newcommand{\@platerefxiv}{\oldpageref{RMAAPlate14}}
\newcommand{\@platerefxv}{\oldpageref{RMAAPlate15}}
\newcounter{figplatei}
\newcounter{figplateii}
\newcounter{figplateiii}
\newcounter{figplateiv}
\newcounter{figplatev}
\newcounter{figplatevi}
\newcounter{figplatevii}
\newcounter{figplateviii}
\newcounter{figplateix}
\newcounter{figplatex}
\newcounter{figplatexi}
\newcounter{figplatexii}
\newcounter{figplatexiii}
\newcounter{figplatexiv}
\newcounter{figplatexv}
%% Printed output should use the counter `plate'
\renewcommand{\theplate}{\@arabic\c@plate}
\renewcommand{\theplateloop}{\@arabic\c@plateloop}
%% Command to output all the saved plates in this paper
\newcommand{\outputplates}{%
  \setcounter{plateloop}{0}
  \if@book\else\@endofarticlehook\fi % not entirely sure why, but this
                                     % seems to be needed
  \clearpage % first output the last page of text
  % now set pagestyle for blank even pages and make sure we start on an odd page
  \pagestyle{empty}
  \cleardoublepage
  \onecolumn
  \parindent\z@
  % first put in an extra sheet
  % the main purpose of this is so there is always a blank even page before 
  % each plate for filling in the PS file if the article has an odd number
  % of pages (see splitbook.sh)
  \null\vfil\centerline{PLATES}\vfil\clearpage\pagestyle{rmaaplate}\cleardoublepage
  % output each plate in turn
  \@whilenum\value{plateloop}<\value{plate}\do{%
    \addtocounter{plateloop}{1}
    % restore the figure counter from \@figcounti, \@figcountii, etc
    \setcounter{figure}{\value{figplate\roman{plateloop}}}
    % spew out the plate (in its own group for safety)
    {\csname @plate\roman{plateloop}\endcsname\par}
    }
  \pagestyle{rmaa}
  % reset the page number as if nothing had happened
  \addtocounter{page}{-\value{plate}}
  \addtocounter{page}{-\value{plate}}
  }
\newcommand\figcaption{%
  \def\@captype{figure}
  \renewcommand{\@makecaption}{\@makefigcap}
  \caption}
\newcommand\tabcaption{%
  \def\@captype{table}
  \renewcommand{\@makecaption}{\@maketabcap}
  \caption}

\newcommand{\plate}[2][1]{%
  %% User command to make a plate page: 
  %% Mandatory argument is what goes on the plate - this should not 
  %% include any floating environments (e.g. figure). Instead, \figcaption 
  %% should be used to do the captions. 
  %% Optional argument is the number of figures (i.e. captions) on the plate
  %% page (default is 1). This is necessary for incrementing the figure counter. 
  %%
  %% Note that we have to use a \gdef rather than an \xdef when we 
  %% save the plate output in a command. This means that in the book styles when it 
  %% comes to spewing out the plate, we will have forgotten about any user macros 
  %% that may have been defined. This can be got round by repeating the macros inside the plate.
  %% (Presumably, one would use a \newcommand rather than \renewcommand if they were new.)
  % increment global and local counters
  \addtocounter{plate}{1}
  % write appropriate label to .aux file
  \csname @platelabel\roman{plate}\endcsname
  % define a global macro called \@platei, \@plateii, \@plateiii, \@plateiv, ...
  \expandafter\gdef\csname @plate\roman{plate}\endcsname{%
    \csname @outputplatelabel\roman{plateloop}\endcsname
    \null\vfil #2 \vfil
    \cleardoublepage
    }
  % save the current shortauthor and shorttitle
  % need to be \xdef to take current values of \@shortauthor, etc 
  \expandafter\xdef\csname @shorta\roman{plate}\endcsname{\@shortauthor}
  \expandafter\xdef\csname @shortt\roman{plate}\endcsname{\@shorttitle}
  % save the current value of the figure counter for later use in \outputplates
  \setcounter{figplate\roman{plate}}{\value{figure}}
  % increment the figure counter according to the optional argument
  \addtocounter{figure}{#1}
  }
%% page style for plate pages 
\newcommand{\ps@rmaaplate}{%
  % Plate X (see page Y) in the foot
  \renewcommand\@oddfoot{\hfil\csname @shorta\roman{plateloop}\endcsname\ 
    (see page~\csname @plateref\roman{plateloop}\endcsname)\hfil}%
  \renewcommand\@evenfoot{\relax}%
  % nothing in even heads (blank pages)
  \renewcommand\@evenhead{\relax}
  % short title and page num in the odd heads
  \renewcommand\@oddhead{\hfil\csname @shortt\roman{plateloop}\endcsname%
    \hfil\llap{PLATE~\theplateloop}}
  }





%% Things for the User Guide
\if@guide
%% stolen from doc.sty
\def\MakeShortVerb#1{%
  \expandafter\ifx\csname cc\string#1\endcsname\relax
  \@shortvrbinfo{Made }{#1}%
  \add@special{#1}%
  \expandafter
  \xdef\csname cc\string#1\endcsname{\the\catcode`#1}%
  \begingroup
  \catcode`\~\active  \lccode`\~`#1%
  \lowercase{%
    \global\expandafter\let
    \csname ac\string#1\endcsname~%
    \gdef~{\verb~}}%
  \endgroup
  \global\catcode`#1\active
  \else
  \@shortvrbinfo\@empty{#1 already}%
  \fi}
\def\DeleteShortVerb#1{%
  \expandafter\ifx\csname cc\string#1\endcsname\relax
  \else
  \@shortvrbinfo{Deleted }{#1 as}%
  \rem@special{#1}%
  \global\catcode`#1\csname cc\string#1\endcsname
  \global \expandafter\let \csname cc\string#1\endcsname \relax
  \ifnum\catcode`#1=\active
  \begingroup
  \catcode`\~\active   \lccode`\~`#1%
  \lowercase{%
    \global\expandafter\let\expandafter~%
    \csname ac\string#1\endcsname}%
  \endgroup \fi \fi}
\def\@shortvrbinfo#1#2{%
  \PackageInfo{doc}{%
    #1\expandafter\@gobble\string#2 a short reference
    for \string\verb}}
\def\add@special#1{%
  \rem@special{#1}%
  \expandafter\gdef\expandafter\dospecials\expandafter
  {\dospecials \do #1}%
  \expandafter\gdef\expandafter\@sanitize\expandafter
  {\@sanitize \@makeother #1}}
\def\rem@special#1{%
  \def\do##1{%
    \ifnum`#1=`##1 \else \noexpand\do\noexpand##1\fi}%
  \xdef\dospecials{\dospecials}%
  \begingroup
  \def\@makeother##1{%
    \ifnum`#1=`##1 \else \noexpand\@makeother\noexpand##1\fi}%
  \xdef\@sanitize{\@sanitize}%
  \endgroup}
%% define | as shortcut for \verbatim
\MakeShortVerb{\|}
%% comand to specify arguments to LaTeX commands
\newcommand{\rmArg}[1]{%
  \ensuremath{\langle}{\scriptsize \textsf{\MakeUpperCase{#1}}}\ensuremath{\rangle}}
%\newcommand{\rmArg}[1]{\ensuremath{\langle}\textsc{#1}\ensuremath{\rangle}}
\providecommand{\BibTeX}{{\rmfamily B\kern-.05em%
    \textsc{i\kern-.025em b}\kern-.08em%
    T\kern-.1667em\lower.7ex\hbox{E}\kern-.125emX}}
%% Environment for code examples
\newenvironment{Code}{%
  \begin{list}{}{\setlength{\leftmargin}{0.05\columnwidth}} 
  \item[]}{\end{list}}
\fi
%</package>
%    \end{macrocode}

% \section{User commands for astronomical articles}
% \label{sec:user-commands}
% Some random definitions. These go both into the main class file and
% also into |adshtml.sty|.
% 
% Definitions of the journal name in different incarnations.  
%    \begin{macrocode}
%<*package|adshtml-style>
\newcommand{\RMAA}{RMxAA}
\newcommand{\RMAAlong}{Revista Mexicana de Astronom\'{\i}a y Astrof\'{\i}sica}
\newcommand{\RMAASC}{RMxAC}
\newcommand{\RMAASClong}{\RMAA\ (Serie de Conferencias)}
%    \end{macrocode}
% \begin{macro}{\ibidrule}
% This puts a horizontal rule with a period after it. It is used when
% the author list of a reference is identical to the preceding one. 
%    \begin{macrocode}
\newcommand{\ibidrule}{\rule[-0.03cm]{4.7em}{.01cm}.\ }
%    \end{macrocode}
% \end{macro}
% \begin{macro}{\ThisVolume}
% References to other articles in the same volume can be coded as,
% e.g., |\bibitem{} Bloggs, J. \ThisVolume{bloggs}|.
%    \begin{macrocode}
\newcommand{\ThisVolume}[1]{%
  \@year, \if@proceedings\RMAASC\else\RMAA\fi, \@volume,
  \if@book\oldpageref{#1:RMAAFirstPage}\else 000\fi\ (this volume)}
%    \end{macrocode}
% \end{macro}
% Other stuff that may be useful in ref lists.
%    \begin{macrocode}
\newcommand{\etal}{et al.\@}
%    \end{macrocode}
% Various new commands from AAS macros
%    \begin{macrocode}
\def\acknowledgments{\vskip 3ex plus .8ex minus .4ex}
\let\acknowledgements=\acknowledgments                  % second spelling
\let\jnl@style=\rm
\def\ref@jnl#1{{\jnl@style#1}}
\DeclareRobustCommand\aj{\ref@jnl{AJ}}                   
\DeclareRobustCommand\araa{\ref@jnl{ARA\&A}}             
\DeclareRobustCommand\apj{\ref@jnl{ApJ}}                 
\DeclareRobustCommand\apjl{\ref@jnl{ApJ}}                
\DeclareRobustCommand\apjs{\ref@jnl{ApJS}}               
\DeclareRobustCommand\ao{\ref@jnl{Appl.~Opt.}}           
\DeclareRobustCommand\apss{\ref@jnl{Ap\&SS}}             
\DeclareRobustCommand\aap{\ref@jnl{A\&A}}                
\DeclareRobustCommand\aapr{\ref@jnl{A\&A~Rev.}}          
\DeclareRobustCommand\aaps{\ref@jnl{A\&AS}}              
\DeclareRobustCommand\azh{\ref@jnl{AZh}}                 
\DeclareRobustCommand\baas{\ref@jnl{BAAS}}               
\DeclareRobustCommand\jrasc{\ref@jnl{JRASC}}             
\DeclareRobustCommand\memras{\ref@jnl{MmRAS}}            
\DeclareRobustCommand\mnras{\ref@jnl{MNRAS}}             
\DeclareRobustCommand\pra{\ref@jnl{Phys.~Rev.~A}}        
\DeclareRobustCommand\prb{\ref@jnl{Phys.~Rev.~B}}        
\DeclareRobustCommand\prc{\ref@jnl{Phys.~Rev.~C}}        
\DeclareRobustCommand\prd{\ref@jnl{Phys.~Rev.~D}}        
\DeclareRobustCommand\pre{\ref@jnl{Phys.~Rev.~E}}        
\DeclareRobustCommand\prl{\ref@jnl{Phys.~Rev.~Lett.}}    
\DeclareRobustCommand\pasp{\ref@jnl{PASP}}               
\DeclareRobustCommand\pasj{\ref@jnl{PASJ}}               
\DeclareRobustCommand\qjras{\ref@jnl{QJRAS}}             
\DeclareRobustCommand\skytel{\ref@jnl{S\&T}}             
\DeclareRobustCommand\solphys{\ref@jnl{Sol.~Phys.}}      
\DeclareRobustCommand\sovast{\ref@jnl{Soviet~Ast.}}      
\DeclareRobustCommand\ssr{\ref@jnl{Space~Sci.~Rev.}}     
\DeclareRobustCommand\zap{\ref@jnl{ZAp}}                 
\DeclareRobustCommand\nat{\ref@jnl{Nature}}              
\DeclareRobustCommand\iaucirc{\ref@jnl{IAU~Circ.}}
\DeclareRobustCommand\aplett{\ref@jnl{Astrophys.~Lett.}}
\DeclareRobustCommand\apspr{\ref@jnl{Astrophys.~Space~Phys.~Res.}}
\DeclareRobustCommand\bain{\ref@jnl{Bull.~Astron.~Inst.~Netherlands}}
\DeclareRobustCommand\fcp{\ref@jnl{Fund.~Cosmic~Phys.}}
\DeclareRobustCommand\gca{\ref@jnl{Geochim.~Cosmochim.~Acta}}
\DeclareRobustCommand\grl{\ref@jnl{Geophys.~Res.~Lett.}}
\DeclareRobustCommand\jcp{\ref@jnl{J.~Chem.~Phys.}}      
\DeclareRobustCommand\jgr{\ref@jnl{J.~Geophys.~Res.}}    
\DeclareRobustCommand\jqsrt{\ref@jnl{J.~Quant.~Spec.~Radiat.~Transf.}}
\DeclareRobustCommand\memsai{\ref@jnl{Mem.~Soc.~Astron.~Italiana}}
\DeclareRobustCommand\nphysa{\ref@jnl{Nucl.~Phys.~A}}
\DeclareRobustCommand\physrep{\ref@jnl{Phys.~Rep.}}
\DeclareRobustCommand\physscr{\ref@jnl{Phys.~Scr}}
\DeclareRobustCommand\planss{\ref@jnl{Planet.~Space~Sci.}}
\DeclareRobustCommand\procspie{\ref@jnl{Proc.~SPIE}}
\let\astap=\aap
\let\apjlett=\apjl
\let\apjsupp=\apjs
\let\applopt=\ao
\def\phn{\phantom{0}}
\def\phd{\phantom{.}}
\def\phs{\phantom{\ensuremath{-}}}
\def\phm#1{\phantom{#1}}
%    \end{macrocode}
% All these are redone as robust and using more modern commands where
% possible. The robustness is so they will not get expanded if used in
% abstracts and titles when these are passed to the ADS file. 
% \changes{3.16}{2003/04/07}{
%  Added \cs{Color} command for doing photometric colors.}
%    \begin{macrocode}
%%\def\sun{\hbox{$\odot$}}
\DeclareRobustCommand\sun{\ensuremath{\odot}}
%%\def\earth{\hbox{$\oplus$}}
\DeclareRobustCommand\earth{\ensuremath{\oplus}}
%%\def\lesssim{\mathrel{\hbox{\rlap{\hbox{\lower4pt\hbox{$\sim$}}}\hbox{$<$}}}}
%%\def\gtrsim{\mathrel{\hbox{\rlap{\hbox{\lower4pt\hbox{$\sim$}}}\hbox{$>$}}}}
\DeclareRobustCommand\lesssim{\mathrel{\hbox{\smaller\smaller\rlap{\hbox{\lower3pt\hbox{$\sim$}}}\raise1.5pt\hbox{$<$}}}}
\DeclareRobustCommand\gtrsim{\mathrel{\hbox{\smaller\smaller\rlap{\hbox{\lower3pt\hbox{$\sim$}}}\raise1.5pt\hbox{$>$}}}}
%%\def\lta{\mathrel{\spose{\lower 3pt\hbox{$\mathchar"218$}}
%%     \raise 2.0pt\hbox{$\mathchar"13C$}}}
%%\def\gta{\mathrel{\spose{\lower 3pt\hbox{$\mathchar"218$}}
%%     \raise 2.0pt\hbox{$\mathchar"13E$}}}
\DeclareRobustCommand\sq{\hbox{\rlap{$\sqcap$}$\sqcup$}}
%%\def\arcdeg{\hbox{$^\circ$}}
%%\def\arcmin{\hbox{$^\prime$}}
%%\def\arcsec{\hbox{$^{\prime\prime}$}}
\DeclareRobustCommand\arcdeg{\ensuremath{^\circ}}
\DeclareRobustCommand\arcmin{\ensuremath{^\prime}}
\DeclareRobustCommand\arcsec{\ensuremath{^{\prime\prime}}}
%%\def\fd{\hbox{$.\!\!^{\rm d}$}}
%%\def\fh{\hbox{$.\!\!^{\rm h}$}}
%%\def\fm{\hbox{$.\!\!^{\rm m}$}}
%%\def\fs{\hbox{$.\!\!^{\rm s}$}}
%%\def\fdg{\hbox{$.\!\!^\circ$}}
%%\def\farcm{\hbox{$.\mkern-4mu^\prime$}}
%%\def\farcs{\hbox{$.\!\!^{\prime\prime}$}}
%%\def\fp{\hbox{$.\!\!^{\scriptscriptstyle\rm p}$}}
\DeclareRobustCommand\fd{\ensuremath{.\!\!^\mathrm{d}}}
\DeclareRobustCommand\fh{\ensuremath{.\!\!^\mathrm{h}}}
\DeclareRobustCommand\fm{\ensuremath{.\!\!^\mathrm{m}}}
\DeclareRobustCommand\fs{\ensuremath{\mkern0.5mu.\mkern-5mu^\mathrm{s}\mkern-1mu}}
\DeclareRobustCommand\fdg{\ensuremath{.\!\!^\circ}}
\DeclareRobustCommand\farcm{\ensuremath{.\mkern-4mu^\prime}}
\DeclareRobustCommand\farcs{\ensuremath{\mkern1mu.\mkern-6mu^{\prime\prime}\mkern-1.5mu}}
%%\DeclareRobustCommand\farcs{\ensuremath{.\!\!^{\prime\prime}}}
\DeclareRobustCommand\fp{\ensuremath{.\!\!^{\scriptscriptstyle\rm p}}}
%%\def\micron{\hbox{$\mu$m}}
\DeclareRobustCommand\micron{\ensuremath{\mu\mathrm{m}}}
\let\la=\lesssim                        
\let\ga=\gtrsim
%%\def\case#1#2{\hbox{$\frac{#1}{#2}$}}
%%\def\slantfrac#1#2{\hbox{$\,^#1\!/_#2$}}
%%\def\onehalf{\slantfrac{1}{2}}
%%\def\onethird{\slantfrac{1}{3}}
%%\def\twothirds{\slantfrac{2}{3}}
%%\def\onequarter{\slantfrac{1}{4}}
%%\def\threequarters{\slantfrac{3}{4}}
\DeclareRobustCommand\case[2]{\ensuremath{\frac{#1}{#2}}}
\DeclareRobustCommand\slantfrac[2]{\ensuremath{\,^#1\!/_#2}}
\DeclareRobustCommand\onehalf{\slantfrac{1}{2}}
\DeclareRobustCommand\onethird{\slantfrac{1}{3}}
\DeclareRobustCommand\twothirds{\slantfrac{2}{3}}
\DeclareRobustCommand\onequarter{\slantfrac{1}{4}}
\DeclareRobustCommand\threequarters{\slantfrac{3}{4}}
%%\def\ubvr{\hbox{$U\!BV\!R$}}            
%%\def\ub{\hbox{$U\!-\!B$}}               
%%\def\bv{\hbox{$B\!-\!V$}}               
%%\def\vr{\hbox{$V\!-\!R$}}               
%%\def\ur{\hbox{$U\!-\!R$}}
%    \end{macrocode}
% New commands for ``photometric'' colors 
%    \begin{macrocode}
\DeclareRobustCommand\ubvr{\ensuremath{U\!BV\!R}}            
\DeclareRobustCommand\Color[2]{%
  \ensuremath{%
    #1\!\ifx#1V\!\fi-\!#2
  }}
\DeclareRobustCommand\ub{\Color{U}{B}}
\DeclareRobustCommand\bv{\Color{B}{V}}
\DeclareRobustCommand\vr{\Color{V}{R}}
\DeclareRobustCommand\ur{\Color{U}{R}}
%% New version of the \ion command, that actually works! 
%    \end{macrocode}
% \begin{macro}{\ion}
% This version of the |\ion| command allows you to write either
% |\ion{N}{4}| or |\ion{N}{iv}| in order to get N~IV. 
%    \begin{macrocode}
\DeclareRobustCommand\ion[2]{%
  \mbox{#1\kern0.2em%
    \smaller\rmfamily%
    \edef\@tempa{\@car#2\@nil}%
    \ifcat1\@tempa%
    \@Roman{#2}%
    \else%
    \uppercase{#2}%
    \fi}}
\newcounter{IonStage}
\DeclareRobustCommand\plainion[2]{\setcounter{IonStage}{#2}#1 \Roman{IonStage}}
%%\def\ion#1#2{\ensuremath{\mathrm{#1}\;}{\protect\small\rm\@Roman{#2}}}
%    \end{macrocode}
% \end{macro}
%    \begin{macrocode}
\DeclareRobustCommand\nodata{$\cdots$}
%    \end{macrocode}
% These are needed by the |\ion| command. Isn't there an easier way of
% doing this?
%    \begin{macrocode}
\let\@small=\small
\let\@normalsize=\normalsize
\let\@footnotesize=\footnotesize
\let\@scriptsize=\scriptsize
\let\@tiny=\tiny
\let\@large=\large
\let\@Large=\Large
\let\@LARGE=\LARGE
\let\@huge=\huge
\let\@Huge=\Huge
\newcommand{\smaller}{\small}
\renewcommand{\normalsize}{\renewcommand{\smaller}{\small}\@normalsize}
\renewcommand{\small}{\renewcommand{\smaller}{\footnotesize}\@small}
\renewcommand{\footnotesize}{\renewcommand{\smaller}{\scriptsize}\@footnotesize}
\renewcommand{\scriptsize}{\renewcommand{\smaller}{\tiny}\@scriptsize}
\renewcommand{\tiny}{\renewcommand{\smaller}{\tiny}\@tiny}
\renewcommand{\large}{\renewcommand{\smaller}{\normalsize}\@large}  
\renewcommand{\Large}{\renewcommand{\smaller}{\large}\@Large}  
\renewcommand{\LARGE}{\renewcommand{\smaller}{\Large}\@LARGE}  
\renewcommand{\huge}{\renewcommand{\smaller}{\LARGE}\@huge}  
\renewcommand{\Huge}{\renewcommand{\smaller}{\huge}\@Huge}  
%</package|adshtml-style>
%    \end{macrocode}


% \section{Other files}

% \subsection{The driver and make files}
% \label{sec:driver_make}
% These are put here for completeness, although they still need to be
% distributed separately since otherwise the user won't be able to
% unpack this file into its constituent parts.
% 
% The |Makefile| has two principal rules: |docs|, and |macros|. The
% |docs| rule produces a PDF and gzipped PS file of the documentation.
% I still haven't perfected the combination of rules required to get
% it to run \LaTeX{} enough times in all circumstances. The |macros|
% rule runs \LaTeX{} on |rmaa.ins| in order to unpack |rmaa.cls| and
% all the other little files (including the |Makefile| and the
% |rmaa.ins| files themselves\dots 
% 
% First set up lists of various types of file. 
%    \begin{macrocode}
%<*makefile>
DOCFILES=rmaa.pdf
MACRO_LTX_FILES=rmaa.cls rmaacite.sty adshtml.sty
SHELLSCRIPTS=makeads.sh
MACRO_OTHER_FILES=adshtml.perl sample.latex2html-init
MAKEFILE=Makefile
MACROFILES=$(MACRO_LTX_FILES) $(SHELLSCRIPTS) $(MACRO_OTHER_FILES)         
SKELETONFILES=ADS.tex 
%    \end{macrocode}
% Install the files read only into the package directory, specified via a path relative to the source directory. \changes{4.0}{2019/11/14}{Reorganize folder layout so that hardly anything is in the top level.  That way, things are clearer on github.}
%    \begin{macrocode}
INSTALL=install 
INSTALL_DIR=../packages
%    \end{macrocode}
% The default rule is to make everything. Installation must be
% explicitly asked for.
%    \begin{macrocode}
all: docs macros

install: install-docs install-macros

%    \end{macrocode}
% Most files are installed read-only to prevent editing. Skeleton
% files must be writeable. Must be careful not to overwrite the
% Makefile in the parent directory. 
%    \begin{macrocode}
install-macros: macros
        $(INSTALL) -m 444 $(MACRO_LTX_FILES) $(MACRO_OTHER_FILES) $(INSTALL_DIR)
        $(INSTALL) -m 555 $(SHELLSCRIPTS) $(INSTALL_DIR)
        $(INSTALL) -m 644 $(SKELETONFILES) $(INSTALL_DIR)

install-docs: docs
        $(INSTALL) rmaa.pdf $(INSTALL_DIR)/rm-fulldocs.pdf

%    \end{macrocode}
% Make the PDF and gzipped PS files. 
%    \begin{macrocode}
docs: $(DOCFILES)

%    \end{macrocode}
% Make everything else if either the |.dtx| or |.ins| files have
% changed. 
%    \begin{macrocode}
macros: rmaa.dtx rmaa.ins
%    \end{macrocode}
% Ensure that all the macro files can be written to then recreate
% them. The |touch| is there in case any of the files don't exist
% yet. 
%    \begin{macrocode}
        touch $(MACROFILES) $(MAKEFILE)
        chmod u+w $(MACROFILES) $(MAKEFILE) rmaa.ins
        latex "\input{rmaa.ins}"
%    \end{macrocode}
% \TeX{} automatically adds the |.tex| extension to extensionless
% files. We fix this and take the opportunity to convert
% back to TABs all that it has lovingly converted to spaces. 
%    \begin{macrocode}
        cp -f $(MAKEFILE) $(MAKEFILE).BCK && unexpand $(MAKEFILE).tex >| $(MAKEFILE) && rm $(MAKEFILE).tex
%    \end{macrocode}
% Make the shell scripts executable and everything unwriteable (so
% they don't get hand-edited by mistake). However, undo this for the
% files that the user is \emph{supposed} to modify. 
%    \begin{macrocode}
        chmod a+x $(SHELLSCRIPTS)
        chmod a-w $(MACROFILES) $(MAKEFILE)
        chmod u+w $(SKELETONFILES)

%    \end{macrocode}
% We now use PDF\LaTeX{} for this.
%    \begin{macrocode}
rmaa.idx: rmaa.dtx
        @echo "REGENERATING .idx FILE"
        pdflatex "\nonstopmode\input{rmaa.dtx}"

rmaa.aux: rmaa.dtx rmaa.ind rmaa.gls
        @echo "REGENERATING .aux FILE"
        pdflatex "\nonstopmode\input{rmaa.dtx}"

rmaa.ind: rmaa.idx
        @echo "REGENERATING INDEX"
        makeindex -s gind.ist rmaa

rmaa.gls: rmaa.glo
        @echo "REGENERATING CHANGELOG"
        makeindex -s gglo.ist -o rmaa.gls rmaa.glo

rmaa.pdf: rmaa.aux rmaa.ind rmaa.gls
        @echo "REGENERATING .pdf FILE"
        pdflatex "\nonstopmode\input{rmaa.dtx}"

clean: 
        rm -f rmaa.aux rmaa.dvi rmaa.ind rmaa.gls
%</makefile>
%    \end{macrocode}

% The |rmaa.ins| file uses the \textsf{docstrip} package to split out
% all the files and strip off this documentation.
%    \begin{macrocode}
%<*ins>
%%
%% Installation batch file for the RMAA LaTeX package
%%
\def\batchfile{rmaa.ins}
\input docstrip.tex
\preamble

PLEASE DO NOT EDIT THIS FILE! It is generated automatically from
the file rmaa.dtx using the accompanying Makefile and rmaa.ins 
Any hand-edited changes will be lost

A LaTeX document class for the RMxAA

Author: William Henney (w.henney@irya.unam.mx)

Latest version available from: http://www.irya.unam.mx/rmaa

Copyright (C) 1999-2019 Instituto de Astronomia, UNAM, Mexico

Permission is hereby granted, free of charge, to any person obtaining
a copy of this software and associated documentation files (the
"Software"), to deal in the Software without restriction, including
without limitation the rights to use, copy, modify, merge, publish,
distribute, sublicense, and/or sell copies of the Software, and to
permit persons to whom the Software is furnished to do so, subject to
the following conditions:

The above copyright notice and this permission notice shall be
included in all copies or substantial portions of the Software.

THE SOFTWARE IS PROVIDED "AS IS", WITHOUT WARRANTY OF ANY KIND,
EXPRESS OR IMPLIED, INCLUDING BUT NOT LIMITED TO THE WARRANTIES OF
MERCHANTABILITY, FITNESS FOR A PARTICULAR PURPOSE AND
NONINFRINGEMENT. IN NO EVENT SHALL THE AUTHORS OR COPYRIGHT HOLDERS BE
LIABLE FOR ANY CLAIM, DAMAGES OR OTHER LIABILITY, WHETHER IN AN ACTION
OF CONTRACT, TORT OR OTHERWISE, ARISING FROM, OUT OF OR IN CONNECTION
WITH THE SOFTWARE OR THE USE OR OTHER DEALINGS IN THE SOFTWARE.
\endpreamble
%    \end{macrocode}
% Clobber files without asking
%    \begin{macrocode}
\askforoverwritefalse
%    \end{macrocode}
% Make the \LaTeX{} macro files (class and styles)
%    \begin{macrocode}
\generate{% 
  \file{rmaa.cls}{\from{rmaa.dtx}{package}}
  \file{rmaa.ins}{\from{rmaa.dtx}{ins}}
  \file{rmaacite.sty}{\from{rmaa.dtx}{rmaacite}}
  \file{adshtml.sty}{\from{rmaa.dtx}{adshtml-style}}
  \file{ADS.tex}{\from{rmaa.dtx}{ads-container}}
  }
%    \end{macrocode}
% Make the other non-\LaTeX{} files. These are done separately since
% they can't use `\%\%' as a comment character for the
% preamble/postamble (and Perl files can't have the postamble I don't
% think). 
%    \begin{macrocode}
\generate{\nopreamble\nopostamble% 
  \file{Makefile}{\from{rmaa.dtx}{makefile}}
  \file{adshtml.perl}{\from{rmaa.dtx}{adshtml-perl}}
  \file{makeads.sh}{\from{rmaa.dtx}{ads-shell-script}}
  \file{sample.latex2html-init}{\from{rmaa.dtx}{latex2html-init}}
  }

%% Local Variables: 
%% mode: LaTeX
%% mode: Font-lock
%% End:
%</ins>
%    \end{macrocode}

% \subsection{Extended citation syntax}
% Modified from a file called mnrascite.sty
% I no longer remember where I stole the original version of this from
% 
% Alterations to the cite commands to make them do more sensible things.
% N.B. the bibliography style should produce entries which have labels where
% the date is delimited by angle brackets '<>' for these macros to work
% properly. They will not give very sensible error messages if you do not do 
% this properly.
% 
% A new command |\scite{key}| is defined which produces a citation which is 
% suitable for insertion in a sentence. i.e., it does not put brackets around 
% the whole citation, only the date. Also |\pcite{key}| produce output with no
% parentheses whatsoever. 
% 
% Should redefine |\@bibitem| really to have any hope of compatibility
% the command |\lbibcite| is defined to deal with the citations.
%    \begin{macrocode}
%<*rmaacite>
\def\bibitem[#1<#2>]#3{\par\hangindent=1.5em\hangafter=1
  \if@filesw
  {\def\protect##1{\string ##1\space}
    \immediate\write\@auxout{\string\lbibcite{#3}{#1}{#2}}}
  \fi}
\@redefinebibitem % what it says
%    \end{macrocode}
% Note that the |\lbibcite| defines two types of citation one contains the names,
% the other the year
%    \begin{macrocode}
\def\lbibcite#1#2#3{\global\@namedef{b@#1}{#2}
                    \global\@namedef{by@#1}{#3}}

\def\@bcite#1#2{({#1\if@tempswa , #2\fi})}
\def\@pcite#1#2{#1\if@tempswa , #2\fi}
%    \end{macrocode}
% This declaration contains what is put between the name and the date
% \changes{3.11}{2003/01/31}{Moved the \cs{articlename} to after the key for consistency
% with the new natbib stuff.}
%    \begin{macrocode}
\def\@brcite{\ (}
\def\@ebrcite{)}
\def\@fmttop#1{\@fmtcite[\csname #1\endcsname]}
\def\@citex[#1]#2{%
%%  \if@filesw\immediate\write\@auxout{\string\citation{#2}}\fi
%% WJH added 16 Jan 2002
  \@for\@citeb:=#2\do{%
    \edef\@citeb{\expandafter\@firstofone\@citeb:\articlename}%
    \if@filesw\immediate\write\@auxout{\string\citation{\@citeb}}\fi}%
%% to here
  \def\@citea{}%
  \@cite{\@for\@citeb:=#2\do
    {%
      \edef\@citeb{\expandafter\@firstofone\@citeb:\articlename}%
      \@citea\def\@citea{;\ }\@ifundefined
      {b@\@citeb}{{\bf ?}\@warning {Citation `\@citeb' on page \thepage \space undefined}}%
      {\csname b@\@citeb\endcsname\@citebro\csname by@\@citeb\endcsname\@citebrc}}}%
  {#1}}

\def\cite{\@ifnextchar [{\let\@citebro=\space
                          \let\@citebrc=\relax
                          \let\@cite=\@bcite\@tempswatrue\@citex}
                        {\let\@citebro=\space
                          \let\@citebrc=\relax
                          \let\@cite=\@bcite\@tempswafalse\@citex[]}}
\def\pcite{\@ifnextchar [{\let\@citebro=\space
                          \let\@citebrc=\relax
                          \let\@cite=\@pcite\@tempswatrue\@citex}
                        {\let\@citebro=\space
                          \let\@citebrc=\relax
                          \let\@cite=\@pcite\@tempswafalse\@citex[]}}
\def\scite{\@ifnextchar [{\let\@citebro=\@brcite
                          \let\@citebrc=\@ebrcite
                          \let\@cite=\@pcite\@tempswatrue\@citex}
                        {\let\@citebro=\@brcite
                          \let\@citebrc=\@ebrcite
                          \let\@cite=\@pcite\@tempswafalse\@citex[]}}

\renewenvironment{thebibliography}
  {\section*{REFERENCES}\vspace*{\afterbiblioskip}\small
    \if@book\RereadBiblio\fi
    \sloppy\clubpenalty4000\widowpenalty4000
    \sfcode`\.=1000\relax
    \setlength\parindent{\z@}\setlength\parskip{\z@}}
  {\relax\normalsize\@outputaddressesifnecessary} 
%</rmaacite>
%    \end{macrocode}

% \subsection{Auxilliary files for generating ADS data}

% The ADS data are written by \LaTeX{} to a |.ads| file. This file
% then has to be included in the following file (|ADS.tex|),
% which is run through \LaTeX{}2HTML (see below). 

%    \begin{macrocode}
%<*ads-container>
%%
%% Instructions:
%% Edit the \input{} line below to reflect the name of your file. 
\documentclass{article}
\usepackage[latin1]{inputenc}
\usepackage{adshtml}
\newcommand{\ADSfield}[2]
{\%#1  #2\\}
%%\newcommand{\ADSvfield}[2]
%%{\%#1 \begin{verbatim}#2\end{verbatim}}
\newenvironment{ADSarticle}{}{\par\bigskip}
\renewcommand{\bibitem}[2][]{\par}
\setlength{\parindent}{0pt}
\renewcommand{\ThisVolume}[1]{%  
  2001, \RMAASC, 12, \oldpageref{#1:RMAAFirstPage}}  
 

\begin{document}
\makeatletter
\input{rmbooksamp.ads}
\end{document}
%</ads-container>
%    \end{macrocode}

% This is the |adshtml.sty| style file used by |ADS.tex|. It just
% defines a few 
% macros that might be used in abstracts, etc (really it should define
% them all). This now gets all its code from
% Sec~\ref{sec:user-commands}. 
%    \begin{macrocode}
%<*adshtml-style>
%% Nothing else needs to go here now. 
%</adshtml-style>
%    \end{macrocode}
% The init file for latex2html.
%    \begin{macrocode}
%<*latex2html-init>
$PDFTOCAIRO = "pdftocairo";
$PDFLATEX = "pdflatex";
$LATEX = "latex";
$MAX_SPLIT_DEPTH = 0;
$HTML_VERSION = "4.0,math";
$NO_SIMPLE_MATH = 1;
$ADDRESS = "<p><address><a href=\"mailto:rmaa\@astro.unam.mx\">rmaa\@astro.unam.mx</address></p>"; 
$INFO = "<p>This document was created using latex2html and the rmaa latex macros</p>"; 
$NO_NAVIGATION = 1;
1; 
%</latex2html-init>
%    \end{macrocode}
% This is the |adshtml.perl| file that is used by \LaTeX{}2HTML. Some
% bizarre \TeX{}y command sometimes show up so we tell it to ignore
% these. All the astronomy macros that come from AAS are processed in
% \TeX{}. That way they get made to images but with an |ALT| tag that
% has the original code, which is what will be picked up by |lynx|,
% with the result that they get passed intact into the final |.txt|
% file. 
%    \begin{macrocode}
%<*adshtml-perl>
package main;

$dum = &do_cmd_makeatletter();

sub do_cmd_RMAA {
    local($_) = @_;
    "RMxAA" . $_;
}

sub do_cmd_RMAASC {
    local($_) = @_;
    "RMxAASC" . $_;
}

sub do_cmd_etal {
    local($_) = @_;
    "et al\." . $_;
}

sub do_cmd_ibidrule {
    local($_) = @_;
    "<img src=\"foo\" alt=\"\ibidrule\">" . $_;
}


sub do_cmd_bibitem {
    local($_) = @_;
    local($key);
    local($opt, $dummy)=&get_next_optional_argument; 
    $key = &missing_braces unless
        ((s/$next_pair_pr_rx/$key = $2; ''/eo)
        ||(s/$next_pair_rx/$key = $2; ''/eo));
    "<p>" . $_; 
}

&ignore_commands (<<_RAW_ARG_CMDS_);
lowercase
adjustfinalcols
textsuperscript # {}
newline # []
thinmuskip
mskip
smallskip
medskip
bigskip
hspace # {}
hspacestar # {}
vspace # {}
vspacestar # {}
mathrm
_at_mathrm
spacefactor
discretionary
unhbox 
voidb_at_x 
penalty 
_at_M
_at_m
null
IeC
_RAW_ARG_CMDS_

&process_commands_in_tex (<<_RAW_ARG_CMDS_);
sun
earth
lesssim
gtrsim
sq
argdeg
arcmin
arcsec
fd
fh
fm
fs
fdg
farcm
farcs
fp
micron
la
ga
case # {} # {}
slantfrac # {} # {}
onehalf
onethird
twothirds
onequarter
threequarters
ubvr
ub
bv
vr
ur
aj
araa
apj
apjl
apjs
ao
apss
aap
aapr
aaps
azh
baas
jrasc
memras
mnras
pra
prb
prc
prd
pre
prl
pasp
pasj
qjras
skytel
solphys
sovast
ssr
zap
nat
iaucirc
aplett
apspr
bain
fcp
gca
grl
jcp
jgr
jqsrt
memsai
nphysa
physrep
physscr
planss
procspie
_RAW_ARG_CMDS_

1;                              # This must be the last line
%</adshtml-perl> 
%    \end{macrocode}

% This is a shell script |makeads.sh| that generates the |_ads.txt| file from
% |_ads.tex| by using \LaTeX{}2HTML and |lynx|.  \changes{4.0}{2019/11/13}{Make sure that lynx writes utf-8 output instead of latin-1.}
% \changes{2.2beta13}{2002/03/12}{Complete rewrite of |makeads.sh|. Now much 
% improved I think}
% \changes{4.0}{2019/11/13}{Fix the at-sign substitution so that it only eliminates \cs{@} and does not mess with email addresses.}

%    \begin{macrocode}
%<*ads-shell-script>
#! /bin/sh
if [ -z $1 ]; then
    echo "Usage: $0 basename"
    exit
fi
BASENAME=$1
ADSBASENAME=${BASENAME}_ads
AUXFILE=$BASENAME.aux
TMPFILE=tmp.$$


echo "Generating $ADSBASENAME.tex"

# For some reason, must double-up \\ to \\\\ but not \ to \\ - odd!
cat > $ADSBASENAME.tex <<EOF
%% PLEASE DO NOT EDIT THIS FILE!  It is generated automatically from
%% the script makeads.sh.  Any hand-edited changes will be lost
%% 
\documentclass{article}
\usepackage[latin1]{inputenc}
\usepackage{html}
\usepackage{adshtml}
\newcommand{\ADSfield}[2]{\%#1  #2\\\\}
\newenvironment{ADSarticle}{}{\par\bigskip}
\setlength{\parindent}{0pt}
\newcounter{istage}
\renewcommand{\ion}[2]{\setcounter{istage}{#2}#1 \Roman{istage}}
\begin{document}
\makeatletter
EOF
echo >> $ADSBASENAME.tex "\input{$BASENAME.ads2}" 
echo >> $ADSBASENAME.tex "\end{document}"



echo "Generating $ADSBASENAME.sed" 

# Find type of journal
grep '\\documentclass\[.*proceedings.*\]' $BASENAME.tex >/dev/null \
     && JOURNAL='RMxAC' \
     || JOURNAL='RMxAA'

# Find the year and volume number
YEAR=`grep '\\SetYear' $BASENAME.tex | sed 's/.*\\SetYear{\(.*\)}.*/\1/'`
VOLUME=`grep '\\SetVolume' $BASENAME.tex | sed 's/.*\\SetVolume{\(.*\)}.*/\1/'`

# Generate an auxilliary sed script 
cat > $TMPFILE <<EOF
s|\\\\newlabel{\(.*\):RMAAFirstPage}{{.*}{\(.*\)}}|s/\\\\\\\\ThisVolume{\1}/$YEAR, RMxAC, $VOLUME, \2/g|g
EOF

# Now go through the .aux file, looking for \newlabels and use the sed 
# script just generated in order to make another sed script, one that 
# will replace \ThisVolume{LABEL} with the correct reference
grep '^\\newlabel{.*:RMAAFirstPage}' $AUXFILE \
     | sed -f $TMPFILE > $ADSBASENAME.sed && rm $TMPFILE

# Apend a further sed command to get rid of \@, since redefining 
# it in _ads.tex didn't work
echo >> $ADSBASENAME.sed 's/\\\\@//g'

echo "Generating $BASENAME.ads2 from $ADSBASENAME.sed and $BASENAME.ads" 

# Apply this to the .ads file, generating the .ads2 file, which is included
# by the _ads.tex file.
sed -f $ADSBASENAME.sed $BASENAME.ads > $BASENAME.ads2

echo "Generating $ADSBASENAME.txt from $ADSBASENAME.tex and $BASENAME.ads2"

# Now run latex2html on the _ads.tex file
latex2html -no_navigation -info 0 -address 0 -html_version 4.0,math $ADSBASENAME

# Then, convert to utf8 with lynx and tidy it up a bit
lynx --justify=off --width=100 --dump --nolist --display_charset=utf-8 \
     $ADSBASENAME/$ADSBASENAME.html | cut -c4- \
     | sed 's/ibidrule/\\ibidrule/g' > $ADSBASENAME.txt
%</ads-shell-script>
%    \end{macrocode}

% \section{Epilogue}
% Err, that's it \dots

\endinput


% \endinput
% Local Variables:
% mode: doctex
% TeX-master: t
% End:
"

rmaa.aux: rmaa.dtx rmaa.ind rmaa.gls
        @echo "REGENERATING .aux FILE"
        pdflatex "\nonstopmode% \def\fileversion{4.0} \def\filedate{2019/11/12}
% \MakeShortVerb{\|}
% \title{A \LaTeX{} document class for the Revista Mexicana de
% Astronom\'\i a y Astrof\'\i sica\thanks{This file has version
% \fileversion{} last revised \filedate{}.}}
% \author{William J. Henney\thanks{w.henney@irya.unam.mx}}
% \date{\filedate}
% \maketitle
% \paragraph{Abstract}
% This document describes the \texttt{rmaa} \LaTeX{}
% document class for the journal ``Revista Mexicana de
% Astronom\'\i a y Astrof\'\i sica''. The macros can be used for 
% \begin{enumerate}
% \item Preparing author manuscripts for the journal (including
% conference proceedings to be published in the ``Serie de
% Conferencias'')
% \item Typesetting an entire volume of the journal by inputting the
% individual article files from a container file. Facilities are
% included for automatic generation of TOC and Author Index. 
% \end{enumerate}
% 
% \tableofcontents
% \section{The documentation driver file}
% This is the part that generates the documentation that you may or
% may not be reading, which is done with the help of the Makefile
% discussed in section~\ref{sec:driver_make}. 
%    \begin{macrocode}
%<*driver>
\documentclass{ltxdoc}
\usepackage{hyperref}
\usepackage{palatino}
\EnableCrossrefs \CodelineIndex \RecordChanges
\setlength{\parindent}{0pt}
\def\MacroFont{\footnotesize\tt}
\begin{document} \DocInput{rmaa.dtx} \PrintIndex \PrintChanges \end{document}
%</driver>
%    \end{macrocode}

% \section{The implementation}
% 
% \changes{3.24}{2007/09/14}{Added user access to the macro version
% and date via \cs{rmVersion} and \cs{rmDate}}
% 
%    \begin{macrocode}
%<*package>
\newcommand{\rmaa@version}{3.26}
\newcommand{\rmaa@date}{2011/08/15}
\newcommand\rmVersion{\rmaa@version}
\newcommand\rmDate{\rmaa@date}
\NeedsTeXFormat{LaTeX2e}[1995/06/01]
%    \end{macrocode}
% \changes{4.0}{2019/11/02}{Change abstract to appear before
% resumen. Change RevMexAA to RMxAA for consistency with ADS.}
%
% \changes{3.22}{2007/03/22}{Allow proceedings articles to have
%   received/accepted dates} \changes{3.21}{2006/05/15}{Changes for
%   Silvia - abstracts are now shipped out to ADS.}
% \changes{3.19}{2006/01/19}{Changes for Silvia - remove adresses from
%   end, new abstract style} \changes{2.2beta}{2001/10/30}{ Fixed bug
%   whereby twocolumn option was not passed through to book
%   class. Finally sorted out the ``label on last page'' issue. }
% \changes{2.1}{2001/01/23}{ Fixed some inconsistencies in
%   \cs{nocite} and \cs{bibitem} uniqification Fixed the definition of
%   \cs{ion} so it will work in math mode and in contexts where the
%   typesize is different from normal (e.g., title, fig captions).
%   Modified the definition of \cs{@caption}, so we don't write any
%   list-of-figures data (we never use it anyway, and it just causes
%   problems with fragile commands in fig captions). }
% \changes{2.0}{2000/11/30}{ Added `newstyle' option, which should be
%   used in all new proceedings volumes (wider abstract, smaller
%   typeface in fig captions, etc.) Added commands for formatting
%   table captions. New demonstation document rmpostsamp.tex for
%   1-page posters. This will be the last release under a PRI
%   government (for a while...)}  \changes{1.9}{2000/06/06}{
%   Corrected position of page number in footer of articles that open
%   on even pages. Ensured that first article in each session opens on
%   an odd page, as do list of abstracts and author index. Added extra
%   space in TOC before the backmatter. Added command for editors,
%   \cs{BoldSessionTitles}, that makes produces bold titles for
%   session in the TOC.}  \changes{1.8}{2000/05/16}{ Fixed bug in
%   last page label with \cs{suppressfulladdresses} }
% \changes{1.7}{2000/05/11}{ Fixed more little things. Page labels
%   now work, even if you have no bibliography. You can output the
%   full addresses by hand with the \cs{outputfulladdresses}
%   command. Editors' Guide updated but still not finished. Revamped
%   splitbook.sh now produces PS file of each article for web page and
%   also an EPS file of each page for sending to the printers.}
% \changes{1.6}{2000/04/15}{ More cleaning up: full stops after
%   titles in index, running headers and footers in frontmatter,
%   fontsize in title of index, etc} \changes{1.5}{2000/04/15}{
%   Eliminated superscripts from titles in TOC and index}
% \changes{1.4}{2000/04/14}{ Added copyright and license info (GPL)}
% \changes{1.3}{2000/03/01}{ Corrected \cs{figcaption} to do the
%   caption formatting right Corrected font size of sec headings in
%   front matter} \changes{1.2}{2000/02/16}{ Added \cs{plate}
%   command} \changes{1.1}{2000/02/14}{ Added manuscript option,
%   corrected bugs} \changes{1.0}{2000/02/08}{ Now does a whole book
%   - lots of stuff tidied up} \changes{0.7}{1999/10/04}{ Increased
%   spacing between sections in single column styles. Suppressed
%   floats at bottom of pages (didn't work though) }
% \changes{0.6}{1999/09/24}{ Tweaked \cs{adjustfinalcols} (still not
%   perfect) Added proceedings option Added onecolumn/twocolumn
%   options (orthogonal to journal/proceedings) }
% \changes{0.5}{1999/09/23}{ Added \cs{fulladdresses} and
%   \cs{adjustfinalcols} (just about working)}
% \changes{0.4}{1999/09/22}{ Minor stylistic changes}
% \changes{0.3}{1999/06/28}{ Better mechanism for alternative
%   affiliations Added `guide' option for user guide}
% \changes{0.2}{1999/05/27}{ All basic functionality present, I
%   think} \changes{0.1}{1999/05/26}{ Just about working - still
%   missing Fig/Tab/Refs} \changes{0.0}{1999/05/21}{ So bored today,
%   decided to rewrite RMAA LaTeX macros from scratch}
%    \begin{macrocode}
\ProvidesClass{rmaa}[\rmaa@date\space v\rmaa@version\space RMxAA article
class]
%    \end{macrocode}
% \subsection{Option Handling}
% \label{sec:options}
% Declarations that are needed in options
% \changes{3.05}{2003/10/22}{Added natbib option}
% \changes{3.06}{2003/01/03}{Natbib option turned on permanently}
% \changes{3.07}{2003/01/11}{Added proof option}
% \changes{3.14}{2003/01/11}{Option to turn off natbib}
% \changes{3.18}{2003/10/22}{trying out doctex mode} 
% \changes{3.23}{2007/09/09}{Manuscript is now on by default.}
% \changes{3.26}{2011/08/15}{New option for review articles}
%    \begin{macrocode}
\newif\if@reviewarticle
\@reviewarticlefalse
\newif\if@proof
\@prooffalse
\newif\if@journal
\@journaltrue
\newif\if@printdates
\@printdatesfalse
\newif\if@proceedings
\@proceedingsfalse
\newif\if@manuscript
\@manuscripttrue
\newif\if@bigfont
\@bigfonttrue
\newif\if@guide
\@guidefalse
\newif\if@preprint
\@preprintfalse
\newif\if@rmtwocol
\@rmtwocoltrue
\newif\if@biblist
\@biblistfalse
\newif\if@book
\@bookfalse
\newif\if@natbib
\@natbibtrue
\newif\if@rmcite
\@rmcitefalse
\newif\if@titlevisible
\@titlevisibletrue
\newif\if@abstractpresent
\@abstractpresentfalse
\newif\if@printbanner
\@printbannerfalse
\newif\if@suppressfulladdresses
\@suppressfulladdressesfalse
\newif\if@newstyle
\@newstylefalse
\newif\if@debug
\@debugfalse
\newif\if@boldkeywords
\@boldkeywordsfalse
%    \end{macrocode}
% The option handling---mainly setting various booleans.
% \changes{2.2beta11}{2001/11/26}{Made twocol and
% newstyle the default for proceedings.}
% \changes{3.23}{2007/09/09}{Manuscript is now 10 pt instead of 12pt}
% \changes{3.23}{2007/09/10}{Full addresses are now automatically
% suppressed for proceedings style.}
% \changes{3.24}{2007/09/14}{New option \texttt{boldkeywords} that
% does the keywords in the old ugly way. Now off by default.}
%    \begin{macrocode}
\DeclareOption{manuscript}{\typeout{Manuscript for referee} 
  \@manuscripttrue\@journalfalse\@rmtwocolfalse
  \PassOptionsToClass{10pt}{article}}
\DeclareOption{bigfont}{\typeout{Big font manuscript requested by Gary
  Ferland} 
  \@manuscripttrue\@journalfalse\@rmtwocolfalse\@bigfonttrue
  \PassOptionsToClass{11pt}{article}}
\DeclareOption{journal}{\typeout{Main journal option}%
  \@journaltrue\@rmtwocoltrue\@printdatestrue\@manuscriptfalse
  \PassOptionsToClass{10pt}{article}}
\DeclareOption{guide}{\typeout{User guide option}%
  \@guidetrue\@journalfalse\@rmtwocoltrue\@manuscriptfalse
  \PassOptionsToClass{10pt}{article}}
\DeclareOption{preprint}{\typeout{Preprint option}%
  \@preprinttrue\@manuscriptfalse\PassOptionsToClass{10pt}{article}}
\DeclareOption{proof}{\typeout{Author proof option}%
  \@prooftrue\@manuscriptfalse}
\DeclareOption{proceedings}{\typeout{Proceedings option}%
  \@proceedingstrue\@journalfalse\@rmtwocoltrue\@newstyletrue
  \@suppressfulladdressestrue
  \@printdatesfalse\@manuscriptfalse\PassOptionsToClass{10pt}{article}} 
\DeclareOption{printdates}{\typeout{Force printing of received/accepted dates}%
  \@printdatestrue} 
\DeclareOption{book}{\typeout{Whole issue option}%
  \@booktrue} 
\DeclareOption{nonatbib}{\typeout{Turn off natbib support (required
    for compatibility with old rmaacite package).}%
  \@natbibfalse} 
\DeclareOption{newstyle}{\typeout{Use new layout style}%
  \@newstyletrue}
\DeclareOption{oldstyle}{\typeout{Use old layout style}%
  \@newstylefalse}
\DeclareOption{debug}{\typeout{Loads of annoying messages}%
  \@debugtrue}
%    \end{macrocode}
% These two need to be declared explicitly so that |\ExecuteOptions|
% actually does something.
%    \begin{macrocode}
\DeclareOption{twocolumn}{\@rmtwocoltrue}
\DeclareOption{onecolumn}{\@rmtwocolfalse}
\DeclareOption{twoside}{\PassOptionsToClass{\CurrentOption}{article}}
\DeclareOption{boldkeywords}{\@boldkeywordstrue}
%    \end{macrocode}
% Bold, upper-case keywords used to be the default. Now they must be
% turned on explicitly with the \texttt{boldkeywords} option. 
%    \begin{macrocode}
\DeclareOption{boldkeywords}{\@boldkeywordstrue}
%    \end{macrocode}
% Pass any other option along to the \texttt{article} class.
%    \begin{macrocode}
\DeclareOption*{\PassOptionsToClass{\CurrentOption}{article}}
%    \end{macrocode}
% Default options are to be two-sided and based on the standard
% \LaTeX{} \texttt{article} class,
% unless we are processing the entire book, in which case we base
% ourselves on the standard \texttt{book} class. Default is to use
% manuscript option unless otherwise specified. 
%    \begin{macrocode}
\ExecuteOptions{twoside,manuscript} \ProcessOptions\relax
\if@book
\if@rmtwocol\LoadClass[twocolumn]{book}\else\LoadClass{book}\fi
\newcommand{\abstract}{\relax}
\newcommand{\@maketitle}{\relax}
\else
\if@rmtwocol\LoadClass[twocolumn]{article}\else\LoadClass{article}\fi
\fi
%    \end{macrocode}
% Attempt to use natbib for citations 02 Jan 2003. 
% \changes{3.05}{2003/01/02}{If the natbib option is used, then set
% things up so that natbib's version of the bibliography is identical
% to the standard one.}
% \changes{3.09}{2003/01/23}{Explicitly require a recent version of
% natbib, so that we get a decent error message if the package is too
% old.} 
% \changes{3.23}{2007/09/09}{Now allow Bibtex to work properly} 
% \changes{3.25}{2009/10/01}{Change \cs{newcommand} tp \cs{def} when setting \cs{bibfont}. This allows it to work with all versions of natbib.} 
% \changes{3.25}{2010/02/06}{Updated copyright year.} 
% \changes{3.25}{2010/02/06}{Updated email and web addresses to \texttt{crya.unam.mx}.} 
%    \begin{macrocode}
\if@natbib
\RequirePackage{natbib}[1999/02/25]
\bibpunct{(}{)}{;}{a}{}{,}
\def\bibfont{\small}
\renewcommand\bibsection{\section*{REFERENCES}}
\setlength\bibsep{0pt}
\setlength\bibhang{1.5em}
\bibliographystyle{rmaa}
\fi
%    \end{macrocode}
% We want the \texttt{graphicx} package for the extended syntax of
% |\includegraphics|. We want the \texttt{ifthen} package for using in
% these macros (although we don't much yet). We explicitly load
% \texttt{babel} since some \LaTeX{} distributions messed up the
% hyphenation if we didn't.  
%    \begin{macrocode}
\RequirePackage{graphicx}
\RequirePackage{ifthen}
\IfFileExists{babel.sty}{%
%   \let\@restoreindent=\@afterindentfalse
  \RequirePackage[english]{babel}
%   \@restoreindent
  %% reinstate the definition of \@roman so as not to break .idx files
  \@ifundefined{org@@roman}{\relax}{\renewcommand{\@roman}[1]{\org@@roman{##1}}}
  }{%
  \typeout{Package babel.sty not found - don't worry about this}}
%    \end{macrocode}
% Try to load the \texttt{dcolumn} and \texttt{booktabs} packages if
% they can be found. These can produce nicer-looking tables. 
%    \begin{macrocode}
\IfFileExists{dcolumn.sty}{%
  \RequirePackage{dcolumn}
  }{%
  \typeout{Package dcolumn.sty not found - you will not be able to
    align columns on decimal points in tables}}
\IfFileExists{booktabs.sty}{%
  \RequirePackage{booktabs}
  }{%
  \typeout{Package booktabs.sty not found - your tables will not look
    as pretty as they ought too.}
  }
\IfFileExists{makeidx.sty}{%
  \RequirePackage{makeidx}
}{%
  \typeout{Package makeidx.sty not found - Author Index may not work.}
  }
%% commented out for moment since requires more modern version of LaTeX
%% \RequirePackage[latin1]{inputenc}
%    \end{macrocode}
% \changes{4.0}{2019/11/13}{Make \texttt{hyperref} and
% \texttt{etoolbox} into hard requirements, since we now use them
% extensively.}  We use hyperref to get internal and external links.
% But I don't like the colored boxes so I turn those off.  We use
% etoolbox for patching and robustifying at the moment.  In the
% future, I want to port a lot of other stuff over to it.  For
% instance, replace the ifthen package and the TeX ifs too.
%
% We target the latest version (1.7) of PDF, rather than the default
% of 1.5.  This reduces the chance of getting warnings from included
% graphics.  This is not strictly related to hyperref, but we put it
% here anyway.
%    \begin{macrocode}
\RequirePackage{etoolbox}
\RequirePackage{hyperref}
\hypersetup{hidelinks}
\pdfminorversion=7
%    \end{macrocode}
% We patch the \cs{@combinedblfloats} following a suggestion by Heiko
% Oberdiek at \url{https://tex.stackexchange.com/a/249743/3843}, which
% avoids problems when hyperlinks are split across
% pages. \changes{4.0}{2019/11/12}{Fix for issue where breaking
% hyperlink across pages gave an error.}
%    \begin{macrocode}
\urlstyle{same}
\newcommand\email{\begingroup \urlstyle{rm}\Url}
\patchcmd\@combinedblfloats{\box\@outputbox}{\unvbox\@outputbox}{}{%
  \errmessage{\noexpand\@combinedblfloats could not be patched}%
}%
%    \end{macrocode}
% \subsection{Other setup stuff}\label{sec:setup}
% Some declarations: 
%    \begin{macrocode}
\newif\if@needaddresses
\@needaddressestrue
\newif\if@frontmatter
\@frontmatterfalse
\newif\if@inpreamble
\@inpreambletrue
\newif\if@firstauthor
\@firstauthortrue
\newif\if@firstaffil
\@firstaffiltrue
\newif\if@firstaltaffil
\@firstaltaffiltrue
\newif\if@firstkwd
\@firstkwdtrue
\newif\if@noexpandADS
\@noexpandADSfalse
\newif\if@sessiontitlepages
\@sessiontitlepagesfalse
%    \end{macrocode}
% \begin{macro}{\@rmaaEmpty}
% \begin{macro}{\@rmaaArg}
% These are used in various commands to check if arg was left blank
%    \begin{macrocode}
\newcommand{\@rmaaEmpty}{}
\newcommand{\@rmaaArg}{}
%    \end{macrocode}  
% \end{macro}
% \end{macro}
%
% \subsection{Compatibility with old \LaTeX{} versions}
% \label{sec:compat}
% This package is based on the 1998/12/01 release of \LaTeX{}.
% Older releases may not define all the commands that we need.
% Luckily, we can use |\providecommand| to define these if necessary. 
% The following commands are not defined in the 1994/06/01 release. 
%
% \begin{macro}{\MakeUpperCase}
% \begin{macro}{\textsuperscript}
% \begin{macro}{\textbackslash}
% User commands:
%    \begin{macrocode}
\providecommand{\MakeUpperCase}[1]{\uppercase{#1}}
\providecommand{\textsuperscript}[1]{\ensuremath{^\textrm{#1}}}
\providecommand{\textbackslash}{\ensuremath{\backslash}}
%    \end{macrocode}  
% \end{macro}
% \end{macro}
% \end{macro}
% \begin{macro}{\protected@xdef}
% \begin{macro}{\restore@protect}
% \begin{macro}{\@unexpandable@protect}
% \begin{macro}{\hb@xt@}
% \LaTeX{} internal commands:
%    \begin{macrocode}
\providecommand{\protected@xdef}{%
        \let\@@protect\protect
        \let\protect\@unexpandable@protect
        \afterassignment\restore@protect
        \xdef}
\providecommand{\restore@protect}{\let\protect\@@protect}
\providecommand{\@unexpandable@protect}{\noexpand\protect\noexpand}
\providecommand{\hb@xt@}{\hbox to}
%    \end{macrocode}
% \end{macro}
% \end{macro}
% \end{macro}
% \end{macro}
%
% \subsection{Page layout}\label{sec:page-layout}
%
% The dimensions of the text were defined by Elizabeth. Text width is
% 169\,mm with equal margins on the two sides: 
%    \begin{macrocode}
\setlength{\textwidth}{169mm}
\setlength{\@tempdima}{\paperwidth}
\addtolength{\@tempdima}{-\textwidth}
\setlength{\oddsidemargin}{0.5\@tempdima}
\addtolength{\oddsidemargin}{-1in}
\setlength{\evensidemargin}{\oddsidemargin}
%    \end{macrocode}
% Text height is 225\,mm, top margin is 14\,mm. 
%    \begin{macrocode}
\setlength{\textheight}{225mm} 
\setlength{\topmargin}{-1in} 
\addtolength{\topmargin}{14mm}
%    \end{macrocode}
% Space between columns in two-column styles, header, footer,
% paragraph indentation:
%    \begin{macrocode}
\setlength{\columnsep}{9.5mm}
\setlength{\headsep}{0.7cm}
\setlength{\footskip}{0.75cm}
\setlength{\parindent}{1.5em}
%    \end{macrocode}
% 
% \begin{macro}{\@set@scalable@length}
% Try out a new scheme for dealing with various spacing dimensions and
% allowing the user to easily modify them. The command
% \begin{quote}
% |\set@scalable@length|\marg{name}\marg{defaultlength}\marg{stretch}\marg{shrink}
% \end{quote}
% does a 
% |\setlength|\marg{\cs{name}}\marg{defaultlength plus stretch minus shrink}
% and also defines the user command
% \textbackslash\meta{name}|Scale|\marg{factor} that can 
% be used to rescale length by \meta{factor}. Additionally, now
% defines the package command
% \textbackslash\meta{name}|@set@scale|\marg{FACTOR} 
% that also rescales the length by \meta{FACTOR}, so that the length
% becomes \meta{factor} $\times$ \meta{FACTOR} $\times$
% \meta{defaultlength}.   
%    \begin{macrocode}
\newcommand\@set@scalable@length[4]{%
  %% arg #1: NAME of length without leading `\'
  %%
  %% arg #2: default VALUE  
  \expandafter\newlength\csname @#1@default\endcsname
  \expandafter\setlength\csname @#1@default\endcsname{#2}%
  \expandafter\newlength\csname @@#1@@default\endcsname
  \expandafter\setlength\csname @@#1@@default\endcsname{#2}%
  %% arg #3: default STRETCH  
  \expandafter\newlength\csname @#1@stretch\endcsname
  \expandafter\setlength\csname @#1@stretch\endcsname{#3}%
  \expandafter\newlength\csname @@#1@@stretch\endcsname
  \expandafter\setlength\csname @@#1@@stretch\endcsname{#3}%
  %% arg #4: default SHRINK  
  \expandafter\newlength\csname @#1@shrink\endcsname
  \expandafter\setlength\csname @#1@shrink\endcsname{#4}%
  \expandafter\newlength\csname @@#1@@shrink\endcsname
  \expandafter\setlength\csname @@#1@@shrink\endcsname{#4}%
  %%
  %% User-defined scale \@NAME@scale is unity by default
  \expandafter\gdef\csname @#1@scale\endcsname{1}%
  %% Package-defined scale \@NAME@SCALE is unity by default
  \expandafter\gdef\csname @#1@SCALE\endcsname{1}%
  %%
  %% internal command \@reset@NAME should expand to 
  %% \setlength\@@NAME@@default{\@NAME@SCALE\@NAME@default}
  %% \setlength\NAME{\@NAME@scale\@@NAME@@default}
  %% in other words: scale x SCALE x default 
  \expandafter\gdef\csname @reset@#1\endcsname{%
    \expandafter\setlength\csname @@#1@@default\endcsname{%
      \csname @#1@SCALE\endcsname \csname @#1@default\endcsname
      }
    \expandafter\setlength\csname @@#1@@stretch\endcsname{%
      \csname @#1@SCALE\endcsname \csname @#1@stretch\endcsname
      }
    \expandafter\setlength\csname @@#1@@shrink\endcsname{%
      \csname @#1@SCALE\endcsname \csname @#1@shrink\endcsname
      }
    \expandafter\setlength\csname #1\endcsname{%
      \csname @#1@scale\endcsname \csname @@#1@@default\endcsname
      \@plus
      \csname @#1@scale\endcsname \csname @@#1@@stretch\endcsname
      \@minus
      \csname @#1@scale\endcsname \csname @@#1@@shrink\endcsname
      }}
  %%
  %% User command \NAMEScale should expand to 
  %% {\renewcommand\@NAME@scale{#1}\@reset@NAME}
  \expandafter\gdef\csname #1Scale\endcsname##1{%
    \expandafter\gdef\csname @#1@scale\endcsname{##1}%
    \csname @reset@#1\endcsname}
  %%
  %% Package command \NAME@set@SCALE should expand to 
  %% {\renewcommand\@NAME@SCALE{#1}\@reset@NAME}
  \expandafter\gdef\csname #1@set@SCALE\endcsname##1{%
    \expandafter\gdef\csname @#1@SCALE\endcsname{##1}%
    \csname @reset@#1\endcsname}
  %%
  %% initialize at default values with \@reset@NAME
  \csname @reset@#1\endcsname
  }
%    \end{macrocode}
% \end{macro}
% \begin{macro}{\parskip}
% First use of this new infrastructure: set interparagraph space to be
% zero but with a small amount of stretchiness.
%    \begin{macrocode}
%   \ifthenelse{\boolean{@manuscript}}{%
%     \@set@scalable@length{parskip}{0.5\baselineskip}{0.3\p@}{0.1\p@}
%   }{
    \@set@scalable@length{parskip}{0\p@}{0.3\p@}{0\p@}
%   }
%    \end{macrocode}
% \end{macro}
% 
% \begin{macro}{\abovesectionskip}
% \changes{4.0}{2019/11/12}{Use negative values for skip so that first paragraph after section headings is not indented.  This is something that I have been trying to fix for years!  It turns out that the anser was in the documentation for the standard command \cs{@startsection.}}
% \begin{macro}{\belowsectionskip}
% \begin{macro}{\aboveparaskip}
% \begin{macro}{\belowparaskip}
% Define some utility rubber lengths for use in
% section~\ref{sec:section-commands}. 
%    \begin{macrocode}
\newlength{\abovesectionskip}
\newlength{\belowsectionskip}
\newlength{\aboveparaskip}
\newlength{\belowparaskip}
%    \end{macrocode}
% The default values for these lengths are those traditionally used
% for the onecolumn style. 
%    \begin{macrocode}
\@set@scalable@length{abovesectionskip}{-3.5ex}{-4ex}{-2ex}
\@set@scalable@length{belowsectionskip}{2.3ex}{0ex}{2ex}
\@set@scalable@length{aboveparaskip}{3.25ex}{2ex}{0.5ex}
\@set@scalable@length{belowparaskip}{-1em}{\z@}{\z@}
%    \end{macrocode}
% \end{macro}\end{macro}\end{macro}\end{macro}
% We also make scalable some of the standard float parameters. These
% are the same defaults as in the standard article classes, but
% slightly more stretchy. For the twocolumn style, they are currently
% halved, see below. Note that I have also lumped |\abovecaptionskip|
% in with these, since it plays a similar role. There is also a
% |\belowcaptionskip| but it is set to zero in the standard classes. 
%    \begin{macrocode}
\@set@scalable@length{floatsep}{12\p@}{4\p@}{4\p@}
\@set@scalable@length{textfloatsep}{20\p@}{4\p@}{8\p@}
\@set@scalable@length{intextsep}{12\p@}{4\p@}{4\p@}
\@set@scalable@length{dblfloatsep}{12\p@}{4\p@}{4\p@}
\@set@scalable@length{dbltextfloatsep}{20\p@}{4\p@}{8\p@}
\@set@scalable@length{abovecaptionskip}{10\p@}{\z@}{\z@}
%    \end{macrocode}
% \begin{macro}{\@before@title@skip}         
% \begin{macro}{\@before@author@skip}        
% \begin{macro}{\@between@author@skip}       
% \begin{macro}{\@before@abstract@skip}      
% \begin{macro}{\@after@abstracthead@skip}   
% \begin{macro}{\@between@abstract@skip}     
% \begin{macro}{\@before@kwd@skip}           
% \begin{macro}{\@before@dates@skip}           
% Spacing parameters used in |\@maketitle|. These are not changed by
% default except by the user, so they don't need to go into
% |\@package@default@lengths|.  
%    \begin{macrocode}
\newlength{\@before@title@skip}         
\newlength{\@before@author@skip}        
\newlength{\@between@author@skip}       
\newlength{\@before@abstract@skip}      
\newlength{\@after@abstracthead@skip}   
\newlength{\@between@abstract@skip}     
\newlength{\@before@kwd@skip}           
\newlength{\@before@dates@skip}           
\@set@scalable@length{@before@title@skip}{2em}{\z@}{\z@}
\@set@scalable@length{@before@author@skip}{1.5em}{\z@}{\z@}
\@set@scalable@length{@between@author@skip}{0.5em}{\z@}{\z@}
\@set@scalable@length{@before@abstract@skip}{3ex}{1ex}{0.5ex}
\@set@scalable@length{@after@abstracthead@skip}{1ex}{\z@}{\z@}
\@set@scalable@length{@between@abstract@skip}{2ex}{1ex}{0.5ex}
\@set@scalable@length{@before@kwd@skip}{1ex}{0.5ex}{0.2ex}
\@set@scalable@length{@before@dates@skip}{2ex}{0.5ex}{0.2ex}
%    \end{macrocode}
% \end{macro}\end{macro}\end{macro}\end{macro}
% \end{macro}\end{macro}\end{macro}\end{macro}
% 
% Spacing above and below displayed equations also wants to be done
% like this, but it is more complicated since the standard article
% classes make these change with the fontsize. Something to do later,
% maybe. 
% 
%    \begin{macrocode}
%    \end{macrocode}
% \begin{macro}{\@package@default@lengths}
% The value actually used for these lengths is controlled by the
% \meta{SCALE} parameter, which depends on the style option that is in
% effect. This is set using the |\@package@default@lengths| command,
% which must be called every time the style changes (e.g., changing
% from onecolumn to twocolumn). 
%    \begin{macrocode}
\newcommand{\@package@default@lengths}{%
  \ifthenelse{\boolean{@rmtwocol}}{%
    %% twocolumn mode
    \abovesectionskip@set@SCALE{0.5}
    \belowsectionskip@set@SCALE{0.5}
    \aboveparaskip@set@SCALE{0.5}
    \belowparaskip@set@SCALE{1}
    \floatsep@set@SCALE{0.5}
    \textfloatsep@set@SCALE{0.5}
    \intextsep@set@SCALE{0.5}
    \dblfloatsep@set@SCALE{0.5}
    \dbltextfloatsep@set@SCALE{0.5}
    }{%
    %% onecolumn mode
    \abovesectionskip@set@SCALE{1}
    \belowsectionskip@set@SCALE{1}
    \aboveparaskip@set@SCALE{1}
    \belowparaskip@set@SCALE{1}
    \floatsep@set@SCALE{1}
    \textfloatsep@set@SCALE{1}
    \intextsep@set@SCALE{1}
    \dblfloatsep@set@SCALE{1}
    \dbltextfloatsep@set@SCALE{1}
    }%
  }
\@package@default@lengths
%    \end{macrocode}
% \end{macro}
% \begin{macro}{\RescaleLengths}
% \begin{macro}{\RescaleSecLengths}
% \begin{macro}{\RescaleFloatLengths}
% \begin{macro}{\RescaleTitleLengths}
% User commands that will fiddle with all the scalable lengths. 
%    \begin{macrocode}
\newcommand{\RescaleLengths}[1]{%
  \RescaleSecLengths{#1}\RescaleFloatLengths{#1}\RescaleTitleLengths{#1}}
\newcommand{\RescaleSecLengths}[1]{%
%% sections
  \abovesectionskipScale{#1}
  \belowsectionskipScale{#1}
  \aboveparaskipScale{#1}
  }
\newcommand{\RescaleFloatLengths}[1]{%
%% floats
  \floatsepScale{#1}
  \textfloatsepScale{#1}
  \intextsepScale{#1}
  \dblfloatsepScale{#1}
  \dbltextfloatsepScale{#1}
  \abovecaptionskipScale{#1}
  }
\newcommand{\RescaleTitleLengths}[1]{%
%% maketitle
  \@before@title@skipScale{#1}           
  \@before@author@skipScale{#1}        
  \@between@author@skipScale{#1}       
  \@before@abstract@skipScale{#1}      
  \@after@abstracthead@skipScale{#1}   
  \@between@abstract@skipScale{#1}     
  \@before@kwd@skipScale{#1}           
  \@before@dates@skipScale{#1}           
  }
%    \end{macrocode}
% \end{macro}\end{macro}\end{macro}\end{macro}
% \changes{2.2beta14}{2002/03/17}{Finished implementation of new
% infrastructure for scalable lengths. Converted the above/below
% section skips to use it. Also, all the floatsep parameters. Also,
% all the vertical lengths that are used by |\@maketitle|. }

% \begin{macro}{\MSextramargin}
% For the \texttt{manuscript} option we have an extra 0.3\,inch
% vertical margin and 1\,inch horizontal margin for marginal notes (by
% hand or with |\marginpar|. 
%
% \changes{3.23}{2007/09/09}{Reduce baselinestretch in
%   \texttt{manuscript} style to 1.2. No longer use raggedright for
%   manuscript.}
%
% \changes{3.24}{2007/09/14}{Further reduction of baselinestretch in
%   \texttt{manuscript} style to 1.1 since we now have wider margins.}
%
% \changes{3.24}{2007/09/14}{Split \texttt{manuscript}'s extra margin
%   into separate horizontal and vertical parts.}
% 
%    \begin{macrocode}
\newlength{\MSextraHmargin}\newlength{\MSextraVmargin}
\setlength{\MSextraHmargin}{1in} 
\setlength{\MSextraVmargin}{0.3in} 
\if@manuscript
   \renewcommand{\baselinestretch}{1.1}
   \addtolength{\textwidth}{-2\MSextraHmargin}
   \addtolength{\oddsidemargin}{\MSextraHmargin}
   \setlength{\evensidemargin}{\oddsidemargin}
   \addtolength{\textheight}{-2\MSextraVmargin}
   \addtolength{\topmargin}{\MSextraVmargin}
%    \raggedright
\fi
%    \end{macrocode}
% \end{macro}
% Float parameters: don't allow bottom floats by default. 
%    \begin{macrocode}
\renewcommand{\fps@figure}{tp}
\renewcommand{\fps@table}{tp}
%    \end{macrocode}
%
% \begin{macro}{\rmaa@AuthorWidthFrac}
% \begin{macro}{\rmaa@AbstractWidthFrac}
% These parameters control what fraction of the total text width can
% be used by the list of authors and the abstract/resumen,
% respectively.  
%    \begin{macrocode}
\newcommand{\rmaa@AuthorWidthFrac}{1.0}
\newcommand{\rmaa@AbstractWidthFrac}{0.75}
%    \end{macrocode}
% Use a wide abstract to save space in the new style.
%    \begin{macrocode}
\if@newstyle\renewcommand{\rmaa@AbstractWidthFrac}{1.0}\fi
%    \end{macrocode}
% \end{macro}
% \end{macro}
%
% Some things to consider in future versions. Trying to avoid too many
% widows, etc. Apparently, the following will work best with
% \cs{flushbottom}. 
%    \begin{macrocode}
%%\widowpenalty=10000
%%\clubpenalty=10000
%    \end{macrocode}
% Also, should perhaps try getting rid of underfull/overfull box warnings. 
% \changes{3.10}{2003/01/24}{Set large values for \cs{hbadness},
% \cs{vbadness}, \cs{hfuzz}, \cs{vfuzz}. To reduce amount of stuff
% that gets written to the terminal every time you compile something.}
%    \begin{macrocode}
\hbadness=10000
\vbadness=10000
\setlength\hfuzz{\maxdimen}
\setlength\vfuzz{\maxdimen}
%    \end{macrocode}
% \subsection{Various text styles for different contexts}
% \label{sec:text-styles} 
% \begin{macro}{\rmaa@AffilStyle}
% \begin{macro}{\rmaa@PagesStyle}
% \begin{macro}{\rmaa@YearStyle}
% \begin{macro}{\rmaa@VolumeStyle}
% \begin{macro}{\rmaa@VolumenStyle}
% \begin{macro}{\rmaa@JournalNameStyle}
% \begin{macro}{\rmaa@ReceivedStyle}
% \begin{macro}{\rmaa@AbstractHeaderStyle}
% \begin{macro}{\rmaa@JournalNameStyle}
% Style commands for the different fields used by |\maketitle| in
% constructing the title page of the article. These have the form
% |\rmaa@XYZStyle{TEXT}|. I'm not sure these were really a good idea,
% but never mind \dots 
%    \begin{macrocode}
\newcommand{\rmaa@AffilStyle}[1]{{\normalsize\rmfamily#1}}
\newcommand{\rmaa@PagesStyle}[1]{{\normalsize\rmfamily#1}}
\newcommand{\rmaa@YearStyle}[1]{{\normalsize\rmfamily#1}}
\newcommand{\rmaa@VolumeStyle}[1]{{\normalsize\rmfamily\bfseries#1}}
\newcommand{\rmaa@VolumenStyle}[1]{{\Large\sffamily\bfseries#1}}
\newcommand{\rmaa@JournalNameStyle}[1]{{\normalsize\itshape#1}}
\newcommand{\rmaa@ReceivedStyle}[1]{{\normalsize\rmfamily\itshape#1}}
\newcommand{\rmaa@AbstractHeaderStyle}[1]{%
  {\normalsize\rmfamily\MakeUpperCase{#1}}}
\newcommand{\rmaa@TitleStyle}[1]{%
  {\large\bfseries\rmfamily\boldmath\MakeUpperCase{#1}\par}}
%    \end{macrocode}
% \end{macro}
% \end{macro}
% \end{macro}
% \end{macro}
% \end{macro}
% \end{macro}
% \end{macro}
% \end{macro}
% \end{macro}
%
% \begin{macro}{\rmaa@FrontSecStyle}
% Headings of ``front sections'' (see section~\ref{sec:front-sections})
%    \begin{macrocode}
\newcommand{\rmaa@FrontSecStyle}[1]{%
  {\normalsize\bfseries\rmfamily\MakeUpperCase{#1}}}
%    \end{macrocode}
% \end{macro}
%
% \begin{macro}{\rmaa@TOCTitleStyle}
% \begin{macro}{\rmaa@TOCPartStyle}
% \begin{macro}{\rmaa@IndexTitleStyle}
% \begin{macro}{\rmaaTOCAuthorStyle}
% \begin{macro}{\rmaaIndexAuthorStyle}
% Commands to set style of TOC and index entries.
% Rewritten to avoid unnecessary font commands in the .toc and .ind files
%    \begin{macrocode}
\newcommand{\rmaa@TOCTitleStyle}[1]{\MakeUpperCase{#1}}
\newcommand{\rmaa@TOCPartStyle}[1]{\rmaa@TOCTitleStyle{#1}}
\newcommand{\rmaa@IndexTitleStyle}[1]{#1}
\newcommand{\rmaaTOCAuthorStyle}[1]{%
  {\itshape#1}}
\newcommand{\rmaaIndexAuthorStyle}[1]{%
  \protect\mbox{\bfseries#1}}
%    \end{macrocode}
% \end{macro}
% \end{macro}
% \end{macro}
% \end{macro}
% \end{macro}
%
% \begin{macro}{\rmaa@KeywordStyle}
% \begin{macro}{\rmaa@KeywordHeaderStyle}
% The style of the keywords depends on whether we are using the
% \texttt{newstyle} option or not. If we are, then they are small. 
%    \begin{macrocode}
\ifthenelse{\boolean{@newstyle}}{
  \newcommand{\rmaa@KeywordStyle}[1]{%
    {\small%
      \if@boldkeywords 
      \bfseries\MakeUpperCase{\raggedright #1}
      \else
      \raggedright #1
      \fi
    }}
  \newcommand{\rmaa@KeywordHeaderStyle}[1]{%
    {\small\itshape#1}}}
{
  \newcommand{\rmaa@KeywordStyle}[1]{%
    {\normalsize%
      \if@boldkeywords\bfseries\fi\raggedright #1}}
  \newcommand{\rmaa@KeywordHeaderStyle}[1]{%
    {\normalsize\itshape#1}}}
%    \end{macrocode}
% \end{macro}
% \end{macro}
%
% \subsection{Page styles}\label{sec:page-styles}
% \begin{macro}{\ps@rmaa}
% The running headers for normal pages:
%    \begin{macrocode}
\newcommand{\ps@rmaa}{%
%    \end{macrocode}
% nothing in either foot; 
%    \begin{macrocode}
  \renewcommand\@oddfoot{\relax}%
  \renewcommand\@evenfoot{\relax}%
%    \end{macrocode}
% page number and short author list in the even heads; 
%    \begin{macrocode}
  \renewcommand\@evenhead{\rlap{\thepage}\hfil\@shortauthor\hfil}%
%    \end{macrocode}
% short title and page number in the odd heads.
%    \begin{macrocode}
  \renewcommand\@oddhead{\hfil\@shorttitle\hfil\llap{\thepage}}
  }
%    \end{macrocode}
% \begin{macro}{\ps@rmaafirstpage}
% The header of the first page of each
% article:
%    \begin{macrocode}
\newcommand{\ps@rmaafirstpage}{%
  \renewcommand\@oddfoot{\hfil\thepage}%
%    \end{macrocode}
% footer corrected (06 Jun 2000) to have page on outside edge;
%    \begin{macrocode}
  \renewcommand\@evenfoot{\thepage\hfil}%
%    \end{macrocode}
% Header has journal name, volume, and page numbers for a regular
% article, or similar for special cases. 
% \changes{4.0}{2019/11/03}{Change to three-line header with copyright info and DOI.}
%    \begin{macrocode}
\if@preprint
\renewcommand\@oddhead{%
  \if@proceedings 
  \textsc{To appear in ``\@conftitle{} (\rmaa@YearStyle{\@year})''}\hfil
  \textit{RMxAC}
  \else
  \textsc{To appear in} \@journalname\hfil
  \fi}\fi%
\if@guide
\renewcommand\@oddhead{\textsc{Documentation} (\rmaa@YearStyle{\@year})\hfil}
\fi%
\if@manuscript
\renewcommand\@oddhead{Manuscript for \@journalname{} (\rmaa@YearStyle{\@year})\hfil}
\fi%
\if@journal
\renewcommand\@oddhead{%
  \parbox{\textwidth}{
    \@journalname, \rmaa@VolumeStyle{\@volume},
    \rmaa@PagesStyle{\@artpages}\
    (\rmaa@YearStyle{\@year})\hfil\\[1pt]
    {\small\textcopyright{} \@year:
      Instituto de Astronom\'\i{}a,
      Universidad Nacional Aut\'onoma de M\'exico}\hfil\\[1pt]
    {\small\@DOI@url}\hfil
  }%
}\fi%
  \renewcommand\@evenhead{\@oddhead}}
%    \end{macrocode}
% By default, use the |\ps@rmaa| page style.
%    \begin{macrocode}
\pagestyle{rmaa}
%    \end{macrocode}
% \end{macro}
% \end{macro}
%
% \subsection{Preamble commands}\label{sec:preamble} 
% Everything you need for setting the title page, TOC and index
% entries, etc.
% \subsubsection{Commands for the title page}
% \label{sec:title-commands}
% \begin{macro}{\resumen}
% \begin{macro}{\@resumen}
% \begin{macro}{\abstract}
% \begin{macro}{\@abstract}
% User commands for setting the resumen and abstract. 
% These should be used in the document header. The internal commands
% |\@resumen| and |\@abstract| are used to hold the actual text
% pending output via the |\maketitle| command. 
%    \begin{macrocode}
\newif\if@setresumen
\newcommand{\resumen}[1]{%
%%  \renewcommand{\@resumen}{{\selectlanguage{spanish}#1}}%
  \renewcommand{\@resumen}{{\ignorespaces#1}}%
  \@setresumentrue}
\renewcommand{\abstract}[1]{%
  \renewcommand{\@abstract}{{\ignorespaces#1}}
  \renewcommand\ADS@abstract{#1}
%% Not a good idea, since then author macros are not expanded and
%% cause errors later. However, by not doing it we are vulnerable to
%% fragile commands.....
  \if@noexpandADS\@onelevel@sanitize{\ADS@abstract}\fi
  \@abstractpresenttrue
  }
\newcommand{\ADSnoexpand}{\@noexpandADStrue}
%    \end{macrocode}
% \begin{macro}{\boldabstract}
% This is the version for doing 1-page conference proceedings. Still
% save the text for ADS but then write it out immediately as a bold
% paragraph. 
%    \begin{macrocode}
\newcommand{\boldabstract}[1]{%
  \gdef\ADS@abstract{#1}%
  \renewcommand{\@abstract}{#1}%
  \if@noexpandADS\@onelevel@sanitize{\ADS@abstract}\fi%
  \noindent\textbf{\boldmath #1}
  }
%    \end{macrocode}
% \end{macro}
% Default value in each case is to print out a warning if not set
%    \begin{macrocode}
  \newcommand{\@resumen}{%
    \@setresumenfalse \ClassWarningNoLine{rmaa}{No Spanish abstract
      (resumen) specified - if you do not speak Spanish, this will be
      translated by the editors}%
    \textbf{Favor de proporcionar un resumen en espa\~nol. If you are
      unable to translate your abstract into Spanish, the editors will
      do it for you.} }
  \newcommand{\@abstract}{\ClassWarningNoLine{rmaa}{No abstract
      specified}}
%    \end{macrocode}
% \end{macro}
% \end{macro}
% \end{macro}
% \end{macro}
%
% \begin{macro}{\title}
% \begin{macro}{\@rawtitle}
% \begin{macro}{\@TOCtitle}
% \begin{macro}{\rmaaIndexTitle}
% Similarly, |\title| is the preamble command to set the title of the
% article. This is used in various places, typeset in different
% styles, hence the use of the internal |\@rawtitle|, |\@title|,
% |\rmaaIndexTitle|, and |\@TOCtitle| commands, each of which uses the
% relevant style from section~\ref{sec:text-styles}
%    \begin{macrocode}
\newcommand{\@rawtitle}{Raw Title}
\newcommand{\@TOCtitle}{\rmaa@TOCTitleStyle{Title for TOC}}
\newcommand{\rmaaIndexTitle}{\rmaa@IndexTitleStyle{Title for index}.}
\newcommand{\@TOCaftertitle}{}
\renewcommand{\title}[1]{%
%%  {\ADS@sanitize\xdef\@rawtitle{#1}}
  \renewcommand\@rawtitle{\noexpand#1}
  \let\ADS@title=\@rawtitle
  \renewcommand{\@title}{\rmaa@TitleStyle{#1}}%
  \renewcommand{\@TOCtitle}{\rmaa@TOCTitleStyle{#1}\@TOCaftertitle}%
  \renewcommand{\rmaaIndexTitle}{\rmaa@IndexTitleStyle{#1\@TOCaftertitle}.}%
  }
%    \end{macrocode}
% \end{macro}
% \end{macro}
% \end{macro}
% \end{macro}
% 
% All the weird and whacky ways of indicating the author
% affiliations. 
% \begin{macro}{\affil}
% First, the way that sets it directly underneath the authors name.
% Note that |\affil| trumps all other affiliations for the purposes of
% getting into the ADS stuff (we use |\gdef|).   
% \dots 
%    \begin{macrocode}
\newcommand{\affil}[1]{%
  \\[1.5ex]\rmaa@AffilStyle{#1}\\[1.5ex]%
  \if@firstaffil\gdef\ADS@affil{#1}\@firstaffilfalse\fi%
  }
\renewcommand{\and}{and\\[1.5ex]}
%    \end{macrocode}
% \begin{macro}{\altaffil}
% Alternative affiliations are footnoted. 
%    \begin{macrocode}
\newcommand{\altaffil}[1]{%
  \if@firstaltaffil\gdef\ADS@affil{#1}\@firstaltaffilfalse\fi%
  \thanks{#1}}
%    \end{macrocode}
% \begin{macro}{\altaffilmark}
% \begin{macro}{\altaffiltext}
% Alternative-alternative mechanism for affiliations---should not be
% mixed with |\altaffil| (unless you know what you are doing) or there
% will be problems with numbering.
%    \begin{macrocode}
\newcommand{\altaffilmark}[1]{\textsuperscript{#1}}
\newcommand{\altaffiltext}[2]{%
  \if@firstaltaffil\gdef\ADS@affil{#2}\@firstaltaffilfalse\fi
  \stepcounter{footnote}%
  \protected@xdef\@thanks{\@thanks
    \protect\footnotetext[#1]{#2}}} 
%    \end{macrocode}
% \end{macro}
% \end{macro}
% \end{macro}
% \end{macro}
%
% \begin{macro}{\keywords}
% \begin{macro}{\@keywords}
% \begin{macro}{\@keywordheader}
% \label{keywords}Sundry minor preamble commands. All
% self-explanatory, I hope. |\keywords| is the old user interface
% (deprecated). New articles should use |\addkeyword| once for each
% keyword. That way the ADS record comes out right too (comma rather
% than em-dash separator).
% \changes{2.2beta12}{2001/12/07}{\cs{addkeyword} New user interface
%   for keywords.} 
% \changes{2.2beta13}{2002/03/11}{\cs{addkeyword} Improved behaviour
%   with .ads file. Use
%   \cs{ADS@sanitize} to change definitions of `~' and \cs{ion}.}
% \changes{3.23}{2007/09/09}{\cs{addkeyword} now signals a compilation
% error if the author tries to add more than 6 keywords.}
% \changes{3.24}{2007/09/14}{kwdcounter is now reset in
% \cs{@begofarticlehook} to prevent error in book style} 
%    \begin{macrocode}
\newcommand{\@kwdsep}{ --- }
\newcommand{\ADS@kwdsep}{, }
\newcounter{kwdcounter}\setcounter{kwdcounter}{0}
\newcommand{\addkeyword}[1]{%
  \addtocounter{kwdcounter}{1}
  \ifthenelse{\thekwdcounter > 6}{
    \ClassError{rmaa}{Too many keywords. A maximum of 6 keywords is allowed.}{}
  }{}
  \if@firstkwd%
  \renewcommand\@keywords{\if@boldkeywords\uppercase{#1}\else #1\fi}
  {\ADS@sanitize\def\ion{\plainion}\gdef\ADS@keywords{#1}}
  \@firstkwdfalse
  \else%
  \edef\@keywords{\@keywords\@kwdsep \if@boldkeywords\uppercase{#1}\else #1\fi}%
  {\ADS@sanitize\def\ion{\plainion}\xdef\ADS@keywords{\ADS@keywords\ADS@kwdsep #1}}%
  \fi}
\newcommand{\ADS@sanitize}{%
  \def\mathrm{}
  \def\@{\ }
  \def~{\ }
  \def\langle{<}
  \def\rangle{>}
  \def\log{log }
  \def\sin{sin }
  \def\cos{cos }
  \def\tan{tan }
  \def\fd{^d.}
  \def\fh{^h.}
  \def\fm{^m.}
  \def\fs{^s.}
  \def\fdg{�.}
  \def\farcm{'.}
  \def\farcs{''.}
  \def\arcdeg{�}
  \def\,{ }
  }
\newcommand{\keywords}[1]{%
  \renewcommand{\ADS@keywords}{#1}
  \renewcommand{\@keywords}{\uppercase{#1}}
  }
\newcommand{\@keywordheader}{\rmaa@KeywordHeaderStyle{Key Words:}}
\newcommand{\@keywords}{%
  \ClassWarningNoLine{rmaa}{No keywords specified}%
  Keywords go here
  }
%    \end{macrocode}
% \begin{macro}{\shorttitle}
% \begin{macro}{\shortauthor}
% \begin{macro}{\@shorttitle}
% \begin{macro}{\@shortauthor}
% Commands for setting the running headers.
%    \begin{macrocode}
\newcommand{\shorttitle}[1]{%
  \renewcommand\@shorttitle{\MakeUpperCase{#1}}}
\newcommand{\shortauthor}[1]{%
  \renewcommand\@shortauthor{\MakeUpperCase{#1}}}
\newcommand{\@shorttitle}{SHORT TITLE GOES HERE}
\newcommand{\@shortauthor}{SHORT AUTHOR GOES HERE}
%    \end{macrocode}
% \end{macro}
% \end{macro}
% \end{macro}
% \end{macro}
% \end{macro}
% \end{macro}
% \end{macro}


% \subsubsection{The full addresses}
% \label{sec:full-addresses}
% \begin{macro}{\fulladdresses}
% \begin{macro}{\@fulladdresses}
% \begin{macro}{\@fulladdheight}
% \begin{macro}{\@fulladddepth}
% \begin{macro}{\@fulladdbox}
% Aaghh, the full addresses! A perennial source of problems. 
%    \begin{macrocode}
\newlength{\@fulladdheight}
\newlength{\@fulladddepth}
\newsavebox{\@fulladdbox}
%% \fulladdresses command modified (11 May 2000) so that we can put
%% \fulladdresses{} in the article to suppress the output of the
%% addresses
\newcommand{\@FinalFigure}{\relax}
\newcommand{\FinalFigure}[1]{\renewcommand{\@FinalFigure}{#1}}
\newcommand{\fulladdresses}[1]{%
  \renewcommand{\@rmaaArg}{#1}
  \ifx\@rmaaArg\@rmaaEmpty
     \renewcommand{\@fulladdresses}{}
  \else
     \renewcommand{\@fulladdresses}{
       \normalsize%
       \ifx\@FinalFigure\relax\else\vfill\@FinalFigure\vfill\fi
%%       \vspace*{1ex}
       %% do addresses as a list (copied from thebibliography environment) 
       \list{\null}{\leftmargin 1.5em\topsep\z@\labelwidth\z@\labelsep\z@\itemindent -1.5em
         \itemsep\z@\parsep\z@\usecounter{enumi}}
       \sloppy\clubpenalty4000\widowpenalty4000
       \sfcode`\.=1000\relax
       #1
       \endlist
       }
  \fi
  }
\newcommand{\@fulladdresses}{Full addresses go here}
%    \end{macrocode}
% \begin{macro}{\suppressfulladdresses}
% New user-level command for suppressing full addresses should you so
% desire. Currently used in the single-page-poster style. 
%    \begin{macrocode}
\newcommand{\suppressfulladdresses}{\@suppressfulladdressestrue\@needaddressesfalse}
%    \end{macrocode}
% \end{macro}
% \end{macro}
% \end{macro}
% \end{macro}
% \end{macro}
% \end{macro}

% \begin{macro}{\printdates}
% \begin{macro}{\noprintdates}
% Command to activate the printing out of the received/accepted
% dates. This can also be done using the \texttt{printdates} option,
% which is the default for \texttt{journal}, but not for \texttt{proceedings}.
%    \begin{macrocode}
\newcommand{\printdates}{\@printdatestrue}
\newcommand{\noprintdates}{\@printdatesfalse}
%    \end{macrocode}
% \end{macro}
% \end{macro}
% \begin{macro}{\@outputaddressesifnecessary}
% Command to output the addresses if we haven't already done so
% This is called by: 
% \begin{enumerate}
% \item The end of |thebibliography| environment
% \item The |\outputfulladdresses| command
% \end{enumerate}
%    \begin{macrocode}
\newcommand{\@outputaddressesifnecessary}{%
  \ClassInfo{rmaa}{Calling @outputaddressesifnecessary}
  \ifthenelse{\boolean{@needaddresses}}{%
    %% If in one column style, or in the first column of a two column
    %% style then we can try to put it on this page 
    \ifthenelse{\boolean{@firstcolumn} \or \not\boolean{@rmtwocol}}{%
      \@outputaddresses}{%
      %% otherwise give warning it will be on a separate page
      \ClassWarning{rmaa}{Full addresses will go on a separate
        page (please see User Guide for how to remedy this)}}%
    %% If we haven't managed yet, put it on its own page
    \ifthenelse{\boolean{@needaddresses}}{%
      \if@debug{\ClassWarning{rmaa}{In this silly bit!}}\fi
      \clearpage\vfill\@outputaddresses}{}%
    }{
    %% If we don't need them, we still might need to do the eoa hook 
    \ifthenelse{\boolean{@suppressfulladdresses}}{%
      \if@debug{\ClassWarning{rmaa}{XXX}}\fi
%%      \if@book\else\@endofarticlehook\null\fi
      }{}
    }%
  }  
%    \end{macrocode}
% \end{macro}

% \begin{macro}{\outputfulladdresses}
% A user interface to the above command, in case there is no biblio
% (11 May 2000) 
%    \begin{macrocode}
\newcommand{\outputfulladdresses}{\@outputaddressesifnecessary}
%    \end{macrocode}
% \end{macro}

% \begin{macro}{\adjustfinalcols}
% This command starts a new column and should be 
% used to ensure that the columns balance on the last page. It also
% outputs the full 
% author addresses to ensure that this occurs in the left hand
% column. 
% The |\par| command is so that the current line will be taken account
% of when working out |\pagetotal|
%    \begin{macrocode}
\newcommand{\adjustfinalcols}{\if@twocolumn\goodbreak\@outputaddresses\newpage\noindent\fi}
%    \end{macrocode}
% \end{macro}

% \begin{macro}{\@outputaddresses}
% Command to do the work of outputting the addresses. 
% This can be called by 
% \begin{enumerate}
% \item an explicit |\adjustfinalcols|
% \item via |\@outputaddressesifnecessary|, q.v.
% \end{enumerate}
%
%    \begin{macrocode}
\newcommand{\@outputaddresses}{%
  % Allow option of suppressing the full addresses, in which case
  % we just run the hook if we are not in a book style
  \if@suppressfulladdresses
%%     \if@book\else\@endofarticlehook\null\fi
  \else
  %% save the full addresses in a box for output later
  \sbox{\@fulladdbox}{\parbox{\textwidth}{\@fulladdresses}}
  % calculate height of address box  
  \settoheight{\@fulladdheight}{\usebox{\@fulladdbox}}
  \settodepth{\@fulladddepth}{\usebox{\@fulladdbox}}
  \addtolength{\@fulladdheight}{\@fulladddepth}
  \if@debug
  { \typeout{Address total height: \the\@fulladdheight}
    \typeout{Text height: \the\pagetotal}
    \typeout{Goal height: \the\pagegoal}
    }
  \fi
  % reuse \@fulladdheight to mean the total height of text+addresses
  \addtolength{\@fulladdheight}{\pagetotal}
%%  \addtolength{\@fulladdheight}{-\baselineskip}
  \if@debug\typeout{Adjusted text + addresses height: \the\@fulladdheight}\fi
  % check it will fit on current page
  \ifdim\@fulladdheight<\pagegoal
     %% must shift it to the left
     %% to compensate for the hanging indentation in the reference list. 
     %% NOTE: this is not foolproof!
     %% NB. This will give overfull hbox warnings, which is a bit ugly.
     \vfill\noindent\if@biblist\hspace*{-\leftmargin}\fi
     \usebox{\@fulladdbox}
     % Set this to show we've done it
     \@needaddressesfalse
     % and assume we are at the end of the article
     % so that we do the EOA hook if not in the book style
%%     \if@book\else\@endofarticlehook\fi
  \else
     % if not, just do nothing. With luck, this function will get 
     % called again on the next page
     \if@debug{\typeout{Holding over full addresses to next page}}\fi
  \fi\fi
  }
%    \end{macrocode}
% \end{macro}

% \subsubsection{Commands for editors}
% \label{sec:editor-commands}
% Some user/editor commands for doing the TOC and index in the book
% option. 
% They are defined here rather in the book-specific part since they
% need to be always available in order to avoid breaking documents. 
% \begin{macro}{\listofauthors}
% \begin{macro}{\@rawauthors}
% \begin{macro}{\rmaaTOCauthors}
% The command |\listofauthors| formats an author list for use in the
% TOC. It should contain the same names as in the |\author| command,
% but without the affiliation garbage and using |\&| instead of
% |and|. 
%    \begin{macrocode}
\newcommand{\@rawauthors}{Raw Authors}
\newcommand{\rmaaTOCauthors}{\noexpand\rmaaTOCAuthorStyle{Author list for TOC}}
\newcommand{\listofauthors}[1]{%
  \renewcommand{\@rawauthors}{#1}
  \renewcommand\rmaaTOCauthors{\noexpand\rmaaTOCAuthorStyle{#1}}}
%    \end{macrocode}
% \begin{macro}{\indexauthor}\label{indexauthor}
% The command |\indexauthor| produces an index entry for each
% author. 
% This command must come \emph{after} the
% |\title| and the |\listofauthors| commands. Also, there should be one 
% |\indexauthor| command for \emph{each} author. We also add the
% author onto the end of the |\ADS@authors| macro and construct the
% |\ADS@authorinitial| by a most cunning method (see below), but only
% for the first author. 
% 
% No, No, NO! This was all wrong. Backed out completely: 
% \begin{quote}
% Completely rewritten (19 Dec 2002) since we now use the index class, so 
% as to be able to alphabetize accented characters. A side effect of this is
% that |\index| no longer expands out the argument after the |@|, so we 
% instead add all the |\index| commands to a list (|\indexauthor@deferlist|), 
% which is executed in the |\@begofarticlehook|. 
% 
% Note the stuff with |\toks0| is copied from Exercise 20.15 of the \TeX{}book---it is
% to make sure the latin1-accented name gets passed into the idx
% file. 
% \end{quote}
% We now go back to the old way of doing things, using an optional
% argument to the \cs{indexauthor} command that has the name without
% accents. 
% \changes{3.08}{2003/01/16}{Remerged changes from 3.06 that got
% dropped from 3.07 after a synchronization disaster. Basically, the
% accentless optional argument to \cs{indexauthor}. Must add the macro
% directories to my unison set-up to stop this sort of thing from happening.}
%    \begin{macrocode}
\newcommand\@isortkey{}
\newcommand\@emptykey{empty}
\newcommand\indexauthor[2][empty]{
  \renewcommand\@isortkey{#1}
  \ifx\@isortkey\@emptykey \renewcommand\@isortkey{#2}\fi
  \index{\@isortkey @{{\rmaaIndexAuthorStyle{#2}}\hskip1em plus1em minus0.4em
      \rmaaIndexTitle\hskip1em plus2em minus0.4em\penalty-100
      \rmaaTOCauthors}}
  \indexauthorADS{#2}
}
%%\def\indexauthor@deferlist{\relax}
%%\long\def\indexauthor#1{
%%  \def\@@name{#1} \toks0=\expandafter{\@@name}
%%  \toks1=\expandafter{\indexauthor@deferlist}
%%  \protected@xdef\indexauthor@deferlist{%
%%    \the\toks1
%%    \noexpand\index{\the\toks0@{\noexpand\rmaaIndexAuthorStyle{\the\toks0}\noexpand\hskip1em plus1em minus0.4em 
%%        \rmaaIndexTitle\noexpand\hskip1em plus2em minus0.4em\noexpand\penalty-100
%%        \rmaaTOCauthors}}
%%  }
%%  \indexauthorADS{#1}
%%  }
%%\xdef\indexauthor@deferlist{}
\newcommand\indexauthorADS[1]{%
  \if@firstauthor\protected@xdef\ADS@authorinitial{\firstletter{#1}}\fi
  \protected@xdef\ADS@authors{\ADS@authors\if@firstauthor\else ; \fi
    #1}
  \@firstauthorfalse
}
%    \end{macrocode}
% Note that the first |#1| sets the key for the index, making sure it
% actually starts with  
% first letter of the surname. Meanwhile, the |@{...}| expression,
% which starts with some  
% control sequences is what is written to the |.ind| file. 
% See the makeindex man page for details. 
% \end{macro}
% \end{macro}
% \end{macro}
% \end{macro}

% \begin{macro}{\KillAllAccents}
% \begin{macro}{\firstletter}
% These are utility functions used in the above. |lett| is a counter
% indicating the current letter in the string. |\firstletter| just
% prints the first letter of its argument.
%    \begin{macrocode}
\newcommand{\KillAllAccents}{\def\'{}\def\`{}\def\~{}\def\.{}\def\i{i}}
\newcommand\firstletter[1]{\@car #1\@nil}
%    \end{macrocode}
% \end{macro}
% \end{macro}
% \begin{macro}{\SetVolume}
% \begin{macro}{\SetMSnumber}
% \begin{macro}{\SetFirstPage}
% \begin{macro}{\SetYear}
% \begin{macro}{\ReceivedDate}
% \begin{macro}{\AcceptedDate}
% These are for the editors to set the volume number, etc. 
%    \begin{macrocode}
\newcommand{\SetVolume}[1]{%
  \renewcommand{\@volume}{#1}
  \edef\ADS@volume{\PadToFourWithDots{#1}}
  \renewcommand{\@volumen}{\rmaa@VolumenStyle{volumen #1}}%
  }
\newcommand{\SetMSnumber}[1]{\renewcommand{\@MSnum}{#1}}
%    \end{macrocode}
% Note that there is no real reason to use |\SetFirstPage| since all
% the page numbering is done automatically in the book style. 
%    \begin{macrocode}
\newcommand{\SetFirstPage}[1]{\setcounter{page}{#1}}
\newcommand{\SetYear}[1]{%
  \renewcommand{\@year}{#1}
  \renewcommand{\ADS@year}{#1}
  \renewcommand{\@anyo}{\rmaa@VolumenStyle{#1}}
  }
\newcommand{\ReceivedDate}[1]{\renewcommand{\@received}{#1}}
\newcommand{\AcceptedDate}[1]{\renewcommand{\@accepted}{#1}}
%    \end{macrocode}
%
% \begin{macro}{\DOI}
% User/editor command to enter the DOI, which will be typeset as a link.
% \changes{4.0}{2019/11/02}{New user command \cs{DOI} to typeset
% DOI in the header of first page and add it to ADS info.}
%    \begin{macrocode}
\newcommand\DOI[1]{%
  \def\@DOI@url{\url{https://doi.org/10.22201/ia.01851101p.#1}}%
  \def\ADS@doi{https://doi.org/10.22201/ia.01851101p.#1}%
}%
%    \end{macrocode}
% \end{macro}
%
% \begin{macro}{\ORCIDiD}
% \changes{4.0}{2019/11/12}{Add mechanism for giving ORCID iD links}
% \begin{macro}{\ORCIDiDfull}
% \begin{macro}{\@orcid@icon}
% \begin{macro}{\@orcid@iconbig}
% Commands to add icon links to ORCID iDs of authors. For example
% \begin{quote}
% |\ORCIDiD|\marg{0000-0000-0000-0000}
% \end{quote}
% to be used after each author name in the |\author|
% command and in the |\fulladdresses|.  Alternatively,
% |\OORCIDiDfull| will typeset the icon plus the full text of the
% link, but we don't use that at present.  The icon image is stored in |ORCIDiD_iconvector.pdf|, which needs to be on the input path.  In principle, we could make EPS versions if necessary.
%    \begin{macrocode}
\newcommand\@orcid@icon{\includegraphics[height=1.8ex]{ORCIDiD_iconvector}}
\newcommand\@orcid@iconbig{\includegraphics[height=2ex]{ORCIDiD_iconvector}}
\newcommand\ORCIDiD[1]{\href{https://orcid.org/#1}{\,\@orcid@icon}}
\newcommand\ORCIDiDfull[1]{%
  \href{https://orcid.org/#1}{\@orcid@iconbig}\,\url{https://orcid.org/#1}}
%    \end{macrocode}
% \end{macro}
% \end{macro}
% \end{macro}
% \end{macro}
%
% \begin{macro}{\@volume}
% \begin{macro}{\@volumen}
% \begin{macro}{\@artpages}
% \begin{macro}{\@year}
% \begin{macro}{\@anyo}
% \begin{macro}{\@received}
% \begin{macro}{\@accepted}
% \begin{macro}{\@MSnum}
% Default values for the containers used by above commands, plus a
% few more useful ones. 
%    \begin{macrocode}
\newcommand{\@volume}{00}
\newcommand{\@volumen}{\rmaa@VolumenStyle{00}}
\newcommand{\@artpages}{%
  \oldpageref{\articlename:RMAAFirstPage}--\oldpageref{\articlename:RMAALastPage}}
\newcommand{\@year}{\the\year}
\newcommand{\@anyo}{\rmaa@VolumenStyle{2007}}
\newcommand{\@journalname}{\rmaa@JournalNameStyle{%
    \if@proceedings\RMAASClong\else\RMAAlong\fi}}
\newcommand{\@received}{}
\newcommand{\@accepted}{}
\newcommand{\@MSnum}{0000}
\newcommand{\@DOI@prefix}{https://doi.org/10.22201/ia.01851101p}
\def\@DOI@url{\url{https://doi.org/10.22201/ia.01851101p.2099.00.00.00}}
%    \end{macrocode}
% \end{macro}
% \end{macro}
% \end{macro}
% \end{macro}
% \end{macro}
% \end{macro}
% \end{macro}
% \end{macro}
% \end{macro}
% \end{macro}
% \end{macro}
% \end{macro}
% \end{macro}
% \end{macro}
% 
% The following commands are most useful in the \texttt{book}
% style. However, the |\SetConfTitle| command can now be used in
% individual articles to set the header for the \texttt{preprint}
% style. 
%    \begin{macrocode}
\newcommand{\@reunionplace}{Lugar de la reuni\'on}
\newcommand{\@reuniondate}{Fecha de la reuni\'on}
\newcommand{\@conftitle}{Conference Title}
\newcommand{\SetReunion}[2]{\renewcommand{\@reunionplace}{#1}\renewcommand{\@reuniondate}{#2}}
\newcommand{\SetConfTitle}[1]{\renewcommand{\@conftitle}{#1}}
\newcommand{\@editors}{Juan Ram\'on de la Fuente and ``El Mosh''}
\newcommand{\SetEditors}[1]{\renewcommand{\@editors}{#1}}
\newcommand{\@editores}{Juan Ram\'on de la Fuente y ``El Mosh''}
\newcommand{\SetEditores}[1]{\renewcommand{\@editores}{#1}}
\newcommand{\@mes}{mes de la publicaci\'on}
\newcommand{\SetMes}[1]{\renewcommand{\@mes}{\rmaa@VolumenStyle{#1}}}
%    \end{macrocode}

% \subsection{Notes for Author Proofs and Communication between Editors}
% \label{sec:notes}
% We introduce a command \cs{Note} that should be used to describe
% any change made to the manuscript. These are a bit tricky to get
% right. In general we use \cs{marginpar} but this only works in outer
% paragraph mode, so we just bung the box where we are if we are in a
% figure caption or something. If we are in the preamble, then we just
% use a \marginpar as part of the \maketitle command.
%    \begin{macrocode}
\newif\if@arenotes
\@arenotesfalse
\newcounter{notecounter}
\setcounter{notecounter}{0}
\newcommand\NotesList{}
\newcommand\Note[1]{%
  \if@proof{%
    \@arenotestrue%
    \stepcounter{notecounter}%
    \if@inpreamble\else{%
      \NoteMark{\arabic{notecounter}}%
    }\fi
    \protected@xdef\NotesList{\NotesList 
      \protect\item #1}%
  }\fi%
}
\newcommand\Changed[2]{Changed ``#1'' to ``#2''.}
\newcommand\NoteMark[1]{%
  \ifinner
  \raisebox{1ex}[0pt][0pt]{\makebox[0pt][l]{\framebox{\footnotesize\bfseries
        N#1}}}%
  \else
  \marginpar{\framebox{\footnotesize\bfseries N#1}}%
  \fi
}
\newcommand\RefNote[1]{\if@proof\NoteMark{\ref{#1}}\relax\fi}
\newcommand\NoteLabel{\protect\label}
\newcommand\OutputNotes{%
  \bgroup
  \onecolumn
  \section*{EDITORIAL NOTES}
  \begin{enumerate}
    \NotesList
  \end{enumerate}
  \egroup
}
\DeclareRobustCommand\cs[1]{\texttt{\char`\\#1}}

%    \end{macrocode}
% \subsection{Bibliographic Records for ADS}
% \label{sec:ADS}
% \changes{2.2beta11}{2001/11/27}{Just about finished all the ADS
% stuff now. Except for the references, which I have given up on for
% the time being.}
% \begin{macro}{\ADS@BibCode}
% \begin{macro}{\ADS@year}
% \begin{macro}{\ADS@journal}
% \begin{macro}{\ADS@volume}
% \begin{macro}{\ADS@firstpage}
% \begin{macro}{\ADS@authorinitial}
% We use the |\ADS@BibCode| command to hold the bibliographic code for
% ADS. It is 
% used in the |%R| 
% field of the bibliographic record for each article. 
% The format is 
% \begin{quote}
% \texttt{YYYYJJJJJVVVVMPPPPA}  
% \end{quote}
% where \texttt{YYYY} is the year, \texttt{JJJJJ} is the journal code,
% which is \texttt{RMxAA} or \texttt{RMxAC}, \texttt{VVVV} is the
% volume number (right-justified), \texttt{M} is not used by us,
% \texttt{PPPP} is the first page number (right-justified), and
% \texttt{A} is the initial of the first author's surname. The volume
% and page fields are padded with periods (`\texttt{.}'), which causes
% us some fun (this is done with the macro |\PadToFourWithDots| described
% below) . 
% 
% All these individual bits are held in the commands |\ADS@year|,
% |\ADS@journal|, |\ADS@volume|, |\ADS@firstpage|,
% |\ADS@authorinitial|. 
%    \begin{macrocode}
\newcommand{\SetADSBibCode}[1]{\renewcommand{\ADS@BibCode}{#1}}
\newcommand{\ADS@BibCode}
{\ADS@year\ADS@journal\ADS@volume.\PadToFourWithDots{\ADS@firstpage}\ADS@authorinitial}
\newcommand{\ADS@year}{2001} 
\newcommand{\ADS@journal}{\if@proceedings RMxAC\else RMxAA\fi}
\newcommand{\ADS@volume}{....} 
%    \end{macrocode}
% These two are just set to default values here, |\ADS@firstpage| is set
% in the |document| environment for included articles, in the
% |frontsection| environment, |\tableofcontents| command,  |index| and
% |abstracts| environments. Also by |\AtBeginDocument| for stand-alone
% articles. 
% 
% God knows how we are going to do |\ADS@authorinitial|!
%    \begin{macrocode}
\newcommand{\ADS@firstpage}{....}
\newcommand{\ADS@authorinitial}{A}
%    \end{macrocode}
% \end{macro}
% \end{macro}
% \end{macro}
% \end{macro}
% \end{macro}
% \end{macro}
% We also have a whole slew of other fields that can be set. These are
% all listed in Table~\ref{tab:ADS}. Most of them still have to be
% defined. 
% \begin{table}[!ht]\raggedleft
% \caption{ADS Bibliographic Records}\label{tab:ADS}
% \begin{tabular}{lp{8em}lp{8em}}\hline
% Code & Description & Macro Name & Comment \\ \hline
% |%R| & Bibliographic Code & |\ADS@BibCode| & required \\
% |%A| & Author list & |\ADS@authors| & required, generate from |\indexauthor| \\
% |%F| & Author affiliation & |\ADS@affil| & format: AA(\dots)
% AB(\dots) \dots \\
% |%J| & Journal Name & |\ADS@journalname| &  \\
% |%D| & Publication Date & |\ADS@fulldate| &  \\
% |%L| & Last Page & |\ADS@lastpage| & generate from label \\
% |%T| & Title & |\ADS@title| & required, generate from |\title| \\
% |%C| & Abstract Copyright & |\ADS@copyright| & set to IA-UNAM \\
% |%U| & URL for electronic document & |\ADS@url| & set to PDF file \\
% |%d| & DOI for article & |\ADS@doi| & set by |\DOI| \\
% |%K| & Keywords & |\ADS@keywords| & set equal to |\@keywords| \\
% |%A| & Abstract Text & |\ADS@abstract| & set equal to |\@abstract| \\
% |%R| & References & |\ADS@references| & generate from the
% bibliography  \\ \hline
% \end{tabular}
% \end{table}
% A lot of these macros are set by other commands, as mentioned in the
% Table, but some need to be set  explicitly here. 
% 
% \begin{macro}{\ADS@doi}
% \changes{4.0}{2019/11/08}{Save DOI for ADS}
% The user/editor command \cs{DOI} is used to set this. 
%    \begin{macrocode}
\newcommand{\ADS@doi}{https://doi.org/10.22201/ia.01851101p.0000.00.00.00}
%    \end{macrocode}
% \end{macro}
% 
% \begin{macro}{\ADS@journalname}
% This
% does not contain any accents at the moment so as not to cause
% problems---this maybe should be changed though. 
%    \begin{macrocode}
\newcommand{\ADS@journalname}{%
  \if@proceedings\@conftitle\space (Eds. \@editors) \fi%
  Revista Mexicana de Astronom\'\i{}a y Astrof\'\i{}sica%
  \if@proceedings\space (Serie de Conferencias)\fi%
  \space Vol. \@volume, pp. \ADS@firstpage--\ADS@lastpage\space (\@year)%
  \space (http://www.astroscu.unam.mx/\~{}rmaa/)
  }
%    \end{macrocode}
% \end{macro}
% \begin{macro}{\ADS@keywords}
% \begin{macro}{\ADS@abstract}
% \begin{macro}{\ADS@title}
% |\ADS@abstract| is straightforward, since we already have the text
% saved as |\@abstract|. We could do the same with |\ADS@keywords| but
% to avoid unwanted formatting we explicitly set it in |\keywords|
% (p.~\pageref{keywords}). |\ADS@title| is just set to |\@rawtitle|,
% which we were already using for another purpose. 
%    \begin{macrocode}
\let\ADS@title=\@rawtitle
\newcommand{\ADS@keywords}{No keywords specified}
\newcommand{\ADS@abstract}{No abstract specified}
%    \end{macrocode}
% \end{macro}\end{macro}\end{macro}
% \begin{macro}{\ADS@authors}
% |\ADS@authors| wants to be in ``Other, A. N.'' format, but with
% semi-colons between authors. The most natural way to do this seems to
% be with |\indexauthor| (p.~\pageref{indexauthor}). We append each
% author to this macro, which means we must explicitly clear it for
% each new article (done in |\article|). .   
%    \begin{macrocode}
\newcommand{\ADS@authors}{}
%    \end{macrocode}
% \end{macro}
% \begin{macro}{\ADS@fulldate}
% \begin{macro}{\ADS@lastpage}
% \begin{macro}{\ADS@copyright}
% \begin{macro}{\ADS@website}
% \begin{macro}{\ADS@url}
% \begin{macro}{\SetWebsite}
% \begin{macro}{\SetMonthNumeric}
% Other sundry data, together with user commands for bits that can be
% set by the Editors. The |\ADS@url| is constructed from the website
% (currently |http://www.astrosmo.unam.mx/rmaa/|,
% the journal descriptor (e.g., |RMxAC..10/|), followed by the
% path to the PDF file (e.g., |PDF/RMxAC..10_HEW1.pdf|). 
% 
%    \begin{macrocode}
\newcommand{\ADS@website}{http://www.astrosmo.unam.mx/rmaa/}
\newcommand{\ADS@url}
{\ADS@website\ADS@journal\ADS@volume\ADS@number/PDF/\ADS@journal\ADS@volume\ADS@number\_\articlename.pdf}
\newcommand{\ADS@number}{}
\newcommand{\ADS@fulldate}{00/\ADS@year} 
\newcommand{\ADS@lastpage}{00} 
\newcommand{\ADS@copyright}{(C) \ADS@year; Instituto de Astronom\'\i{}a,
  Universidad Nacional Aut\'onoma de M\'exico} 
\newcommand{\SetNumber}[1]{\renewcommand{\ADS@number}{-#1}}
\newcommand{\SetWebsite}[1]{\renewcommand{\ADS@website}{#1}}
\newcommand{\SetMonthNumeric}[1]{\renewcommand{\ADS@fulldate}{#1/\ADS@year}}
%    \end{macrocode}
% \end{macro}\end{macro}\end{macro}\end{macro}\end{macro}\end{macro}\end{macro}
% \begin{macro}{\ADS@affil}
% \begin{macro}{\ADS@references}
% These two are somewhat more problematical. Affiliation currently
% only works for the first author (see \dots) pending redesign of the
% user interface. References are now done by wrapping the
% |thebibliography| environment to open another input stream on the
% same file and search through for the |\bibitems|. 
%    \begin{macrocode}
\newcommand{\ADS@references}{}
\newcommand{\ADS@affil}{}
%    \end{macrocode}
% \end{macro}\end{macro}
% 
% Define two macros that expand to the strings
% |\begin{thebibliography}| and |\end{thebibliography}|. Note that the
% idea for this and much else is inspired by the code for the docstrip
% package. Also define |\StartBibX| in case authors use the more
% standard |\begin{thebibliography}{}| syntax. 
%    \begin{macrocode}
\edef\EndBib{\expandafter\@gobble\string\\end\string{thebibliography\string}}
\edef\StartBib{\expandafter\@gobble\string\\begin\string{thebibliography\string}}
\edef\StartBibX{\expandafter\@gobble\string\\begin\string{thebibliography\string}\string{\string}}
%    \end{macrocode}
% Define two new logical variables. |\ifAnotherline| is true if we
% want to read another line from the input stream. It is set to false
% when we encounter the |\end{thebibliography}|. |\ifInBib| is set to
% true when we frind the |\begin{thebibliography}| and tells us to
% start saving the following lines. 
%    \begin{macrocode}
\newif\ifAnotherLine
\newif\ifInBib
%    \end{macrocode}
% \begin{macro}{makeOther}
% We need a macro to change the catcodes of lots of things to other. 
%    \begin{macrocode}
\def\makeOther#1{\catcode`#1=12\relax}
%    \end{macrocode}
% \end{macro}
% \begin{macro}{\RereadBiblio}
% Define a macro to do all the fiddly stuff. This is used at the
% beginning of |thebibliography| environment (Sec~\ref{sec:biblio}).
%    \begin{macrocode}
\newread\tf@biblio
\newcommand{\RereadBiblio}{%
%    \end{macrocode}
% Open the current file again for reading
%    \begin{macrocode}
  \begingroup
  \immediate\openin\tf@biblio=\articlename
%    \end{macrocode}
% Everything in a group to be safe. Initialize logicals. Black magic
% to change
% catcodes and avoid blank lines being interpreted as |\par|. 
%    \begin{macrocode}
  \AnotherLinetrue\InBibfalse
  \makeOther\ \makeOther\\\makeOther\$%
  \makeOther\#\makeOther\^\makeOther\^^K%
  \makeOther\_\makeOther\^^A\makeOther\%%
  \makeOther\~\makeOther\{\makeOther\}\makeOther\&%
  \endlinechar-1\relax
%    \end{macrocode}
% Start looping through the file. Read each line to the macro
% |\Line|. If we have got to EOF or end of biblio then set
% |\AnotherLinefalse|.
%    \begin{macrocode}
  \loop
  \read\tf@biblio to \Line
  \ifx\Line\EndBib\AnotherLinefalse\InBibfalse\fi
  \ifeof\tf@biblio \AnotherLinefalse\fi
%    \end{macrocode}
% If we are in the biblio and the line is not blank then append it to
% the |\ADS@references| macro together with a newline. 
%    \begin{macrocode}
  \ifInBib
  \ifx\Line\empty\else\xdef\ADS@references{\ADS@references\Line^^J}\fi
  \fi
%    \end{macrocode}
% Check if the line we just saw was the start of biblio, then repeat
% the loop if necessary. 
%    \begin{macrocode}
  \ifx\Line\StartBib\InBibtrue\fi
  \ifx\Line\StartBibX\InBibtrue\fi
  \ifAnotherLine\repeat
  \immediate\closein\tf@biblio
  \endgroup
  }
%    \end{macrocode}
% \end{macro}
% 
% \begin{macro}{\PadToFourWithDots}
% Utility function for left-padding a number to 4 characters with
% periods. Used in 
% constructing |\ADS@BibCode|. This is a very simple algorithm and
% \emph{requires} that the argument actually be a number---if it is
% not then dire consequences will result.  
%    \begin{macrocode}
\newcommand{\PadToFourWithDots}[1]{%
  \ifnum #1 > 999 #1%
  \else\ifnum #1 > 99 .#1%
  \else\ifnum #1 > 9 ..#1%
  \else ...#1%
  \fi\fi\fi}
%%  \ifnum \count0 > 999 #1%
%%  \else\ifnum \count0 > 99 .#1%
%%  \else\ifnum \count0 > 9 ..#1%
%%  \else ...#1%
%%  \fi\fi\fi}
%    \end{macrocode}
% \end{macro}
%
% \begin{macro}{\AuthorADS}
% \begin{macro}{\ADS@affiliations}
% New mechanism for setting the affiliations for ADS. The author labels (AA, AB, \dots) are generated using two nested counters.  We use the |\robustify| command from \texttt{etoolbox} so that accented characters in the affiliation addresses do not blow up.  The list of accents may need to be expanded. 
%    \begin{macrocode}
\newcounter{AUTHOR}
\newcounter{author}[AUTHOR]
\stepcounter{AUTHOR}
\newcommand\ADS@affiliations{}
\robustify{\'}\robustify{\i}\robustify{\~}\robustify{\`}
\def\AuthorADS#1#2{%
  \begingroup
  \def\ORCIDiD##1{<ID system="ORCID">##1</ID>}
  \def\Email##1{<EMAIL>##1</EMAIL>}
  \def\Affil##1{##1}
  \stepcounter{author}
  \ifthenelse{\value{author} > 26}{\stepcounter{AUTHOR}\stepcounter{author}}{}
  \ifthenelse{\equal{\ADS@affiliations}{\@rmaaEmpty}}{%
    \xdef\ADS@affiliations{\Alph{AUTHOR}\Alph{author}(#2)}
  }{%
    \xdef\ADS@affiliations{\ADS@affiliations, \Alph{AUTHOR}\Alph{author}(#2)}
  }%
  \endgroup
}
%    \end{macrocode}
% \end{macro}
% \end{macro}
% %
% \begin{macro}{\ADS@dumpdata}
% \changes{4.0}{2019/11/13}{Write all author affiliations for ADS}
% \begin{macro}{\ADS@write}
% \begin{macro}{\ADS@openfile}
% |\ADS@dumpdata| is the macro that actually does the job of writing out all the
% data to an auxiliary file, using |\ADS@write| to write each
% field. The file is set up with |\ADS@openfile|, which currently is
% done in the TOC of a book.  The new multi-author affiliations are used if |\ADS@affiliations| has been set, otherwise it falls back to the single-author version |\ADS@affil|. 
%    \begin{macrocode}
\newcommand{\ADSfield}[2]{\relax}
\newcommand{\ADSvfield}[2]{\relax}
%%\renewcommand{\ADSfield}[12]{\%#1 #2}
\newcommand{\ADS@dumpdata}
{ 
  \ADS@write{\string\begin{ADSarticle}}
  \ADS@WRITE{\string\ADSfield{R}{\ADS@BibCode}}
  \ADS@write{\string\ADSfield{A}{\ADS@authors}}
  \ifthenelse{\equal{\ADS@affiliations}{\@rmaaEmpty}}{%
    {\ADS@sanitize\ADS@write{\string\ADSfield{F}{AA(\ADS@affil)}}}
  }{%
    {\ADS@sanitize\ADS@write{\string\ADSfield{F}{\ADS@affiliations}}}
  }
  \ADS@write{\string\ADSfield{J}{\ADS@journalname}}
  \ADS@write{\string\ADSfield{D}{\ADS@fulldate}}
  \ADS@write{\string\ADSfield{L}{\ADS@lastpage}}
  {\ADS@sanitize\ADS@write{\string\ADSfield{T}{\ADS@title}}}
  \ADS@write{\string\ADSfield{C}{\ADS@copyright}}
  \ADS@write{\string\ADSfield{d}{\ADS@doi}}
  \ADS@write{\string\ADSfield{U}{\ADS@url}}
  \ADS@write{\string\ADSfield{K}{\ADS@keywords}}
  {\ADS@sanitize\ADS@write{\string\ADSfield{B}{\ADS@abstract}}}
  \ADS@write{\string\ADSfield{Z}{\ADS@references}}
  \ADS@write{\string\end{ADSarticle}}
  }
%%\newcommand{\ADS@write}[1]{\if@book\addtocontents{ads}{#1}\fi}
%%\newcommand{\ADS@openfile}{\@starttoc{ads}}
\newcommand{\ADS@write}[1]{\set@display@protect\immediate\write\tf@ads{#1}}
\newcommand{\ADS@WRITE}[1]{\immediate\write\tf@ads{#1}}
\newcommand{\ADS@openfile}{%
  \newwrite\tf@ads%
  \immediate\openout\tf@ads\jobname.ads\relax
  \ADS@WRITE{\ADS@preamble}}
\newcommand{\ADS@closefile}{\ADS@WRITE{\ADS@postamble}}
\def\ADS@preamble{%
  %%  \string\documentclass{article}
  %%  \string\begin{document}
  }
\def\ADS@postamble{%
  %%  \string\end{document}
  \string%\string% Local Variables:^J 
  \string%\string% mode: LaTeX^J
  \string%\string% mode: Font-lock^J
  \string%\string% End:^J
  }
%    \end{macrocode}
% \end{macro}
% \end{macro}
% \end{macro}
% 
% \subsection{Automatic calculation of the page range}
% \label{sec:page-range}
% Another issue that has cost me a lot of grief. Still, it finally
% seems to be sorted now. All hail Jeff Goldberg. 
%
% Cunning trick to get the page numbers from |\pageref|.  Put a label
% on the first and last page of the document using |\AtBeginDocument|
% and |\AtEndDocument|. .  This way won't work for the book option.
% In that case, we do it by defining |\@endofarticlehook| and
% |\@begofarticlehook| as parallels to |\AtEndDocument| and
% |\AtBeginDocument|.  Note that we have to defer adding
% |\@begofarticlehook| to |\AtBeginDocument| to make sure it comes
% after the redefinitions of |\ref| and |\label|.
%    \begin{macrocode}
\if@book\AtEndDocument{\ADS@closefile}\else
 %\AtBeginDocument{\@begofarticlehook}
\AtEndDocument{\@endofarticlehook}
\fi
%    \end{macrocode}
% \begin{macro}{\@begofarticlehook}
% Since we are getting more and more things that need to be done at
% the beginning of every article, frontsection, etc we define a sister
% hook to 
% |\@endofarticlehook|. This hook does the following: 
% \begin{enumerate}
% \item Labels the first page
% \item Sets |\ADS@firstpage|
% \item runs and then clears |\indexauthor@deferlist|
% \item resets the keyword counter
% \end{enumerate}
% It is executed in the following circumstances: 
% \begin{enumerate}
% \item At |\begin{document}| for standalone files
% \item At |\begin{document}| for each included file
% \item At beginning of front sections 
% \item At beginning of TOC
% \item At beginning of Author Index
% \item At beginning of abstracts
% \end{enumerate}
% For some reason, this is not always within the right group, so we
% use |\xdef| just to be sure. 
%    \begin{macrocode}
\newcommand{\@begofarticlehook}{%
  \if@frontmatter\frontlabel{RMAAFirstPage}\else\label{RMAAFirstPage}\fi%
  \xdef\ADS@firstpage{\@arabic{\c@page}}%
  \if@book\else\makeindex\fi % needed so default index is defined in stand-alone articles
%%  \indexauthor@deferlist
%%  \xdef\indexauthor@deferlist{}
  \@inpreamblefalse
  \setcounter{kwdcounter}{0}
  }
%    \end{macrocode}
% \end{macro}
% \begin{macro}{\@endofarticlehook}
% This is a hook to be run at the end of each article
% currently, it just labels the last page and writes out all the ADS
% data to an auxiliary file.  
% It should not produce any printed output. 
% It is modelled on the |\AtEndDocument| command, to which it reduces
% in the case of a single article. 
%    \begin{macrocode}
\newcommand{\@endofarticlehook}{%
  \if@debug{\ClassWarning{rmaa}{***Doing Hook***}}\fi
  \if@debug{\ClassWarning{rmaa}{ADSabstract is: \ADS@abstract}}\fi
  \clearpage\rmaa@lastpagelabel\if@book\ADS@dumpdata\fi
  \ifx\NotesList\@rmaaEmpty\else\OutputNotes\fi
  \if@debug{\ClassWarning{rmaa}{Notes list is: \NotesList}}\fi
}
%    \end{macrocode}
% \begin{macro}{\rmaa@lastpagelabel}
% New version of last page label command copied from Jeff Goldberg's
% lastpage.sty package - we can't just use |\label| here in case the
% output routine for the last page has already been called. Instead,
% we first explicitly clear the page, step back the page counter,
% explicitly write to the .aux file, then restep the page counter
% again.
%
% But this will conflict with the hyperref package, which does labels differently.  Compensate for that by using the five-argument form.
% 
% Also set the |\ADS@LastPage| macro to the appropriate number. 
%    \begin{macrocode}
\newcommand{\rmaa@lastpagelabel}{%
  \addtocounter{page}{-1}%
  \immediate\write\@auxout{\string
    \newlabel{\articlename:RMAALastPage}{{}{\thepage}{}{}{}}}%
  \edef\ADS@lastpage{\thepage}%
  \addtocounter{page}{1}}
%    \end{macrocode}
% \end{macro}
% \end{macro}

% \subsection{Deal with explicit change in number of columns}
% \label{sec:columns}
% Sometimes, when formatting a book, we need to switch between
% one-column articles and two-column articles. For this to work, we
% need to wrap the standard \LaTeX{} commands to set the |\@rmtwocol|
% flag appropriately and also reset the package default values of all
% scalable lengths. .  
%    \begin{macrocode}
\let\latexonecolumn=\onecolumn
\let\latextwocolumn=\twocolumn
\renewcommand{\onecolumn}{\@rmtwocolfalse\@package@default@lengths\latexonecolumn}
\renewcommand{\twocolumn}{\@rmtwocoltrue\@package@default@lengths\latextwocolumn}
%    \end{macrocode}
% 
% \subsection{The actual formatting of the title}
% \label{sec:title-format}

% Note, we have to redefine the |\maketitle| command to use the
% numbered footnotes. This means we needn't use the atbeginend stuff. 
%    \begin{macrocode}
\renewcommand{\maketitle}{\par
  \begingroup
%    \end{macrocode}
% Don't switch to marked footnotes. 
% This is consistent with existing style but I'm not sure I like it
%    \begin{macrocode}
  %%  \renewcommand\thefootnote{\fnsymbol{footnote}}%
%%  \def\@makefnmark{\hbox to\z@{$\m@th^{\@thefnmark}$\hss}}%
  \def\@makefnmark{\textsuperscript{\@thefnmark}\hspace{-0.2em}}%
  \long\def\@makefntext##1{\parindent 1em\noindent
    \hbox to1.8em{\hss$\m@th^{\@thefnmark}$}##1}%
  \if@rmtwocol
  { \if@debug{\ClassWarning{rmaa}{TWO COLUMN}}\fi
    \twocolumn[\@maketitle]%
    }
  \else
  { \if@debug{\ClassWarning{rmaa}{ONE COLUMN}}\fi
    \newpage
    \global\@topnum\z@   % Prevents figures from going at top of page.
    \if@manuscript \@MSmaketitle \else \@maketitle \fi
    }
  \fi
  \if@manuscript\else\thispagestyle{rmaafirstpage}\@thanks\fi
  \ifx\NotesList\@rmaaEmpty\else{%
    \NoteMark{1--\arabic{notecounter}}\vspace*{-2\baselineskip}%
  }\fi%
  \endgroup
%    \end{macrocode}
% Don't reset footnote counter (not sure whether this is good thing
% or not).
%    \begin{macrocode}
  %%  \setcounter{footnote}{0}%
  \let\thanks\relax\let\affil\relax\let\altaffil\relax
  \let\maketitle\relax\let\@maketitle\relax
  \gdef\@thanks{}\gdef\@author{}\gdef\@title{}}
%    \end{macrocode}

% \begin{macro}{\@kwhwidth}
% \begin{macro}{\@kwtwidth}
% Redefine \cs{@maketitle} to include the abstract, resumen and
% keywords. Note that in the manuscripty style we use
% \cs{@MSmaketitle} instead.  
% \changes{3.16}{2003/04/07}{
%  Added extra vertical space between authors and accepted date.}
% \changes{3.26}{2011/08/15}{Add extra label to review articles}
%    \begin{macrocode}
\newlength{\@kwhwidth}\newlength{\@kwtwidth}
\newcommand\reviewname{Review}
\renewcommand{\@maketitle}{%
  \newpage
  \null
  \if@reviewarticle
  \smallskip
  {\large\bfseries\itshape\reviewname}%
  \bigskip
  \else
  \vskip\@before@title@skip%
  \fi
  \begin{center}%
    {\@title \par}%
    \vskip\@before@author@skip%
    \parbox{\rmaa@AuthorWidthFrac\textwidth}{\centering
      \lineskip\@between@author@skip%
      \@author
      \par}
    %% this should be replaced with `Submitted to...' for preprints
    \if@printdates
    \vskip\@before@dates@skip%
    \if@preprint
    \rmaa@ReceivedStyle{Submitted to RMxAA \@received}
    \else
    \rmaa@ReceivedStyle{Received \@received; accepted \@accepted}
    \fi
    \fi
    \if@manuscript
    \rmaa@ReceivedStyle{Submitted to RMxAA \@received}
    \fi
    \if@guide
    \rmaa@ReceivedStyle{Version \rmaa@version, \rmaa@date}
    \fi
    \vskip\@before@abstract@skip%
    \if@abstractpresent
    \parbox{\rmaa@AbstractWidthFrac\textwidth}{%
      \setlength{\parindent}{0.9cm}%
      \centerline{\rmaa@AbstractHeaderStyle{Abstract}}\par%
      \vskip\@after@abstracthead@skip%
      \if@newstyle\noindent\fi\@abstract
      \vskip\@between@abstract@skip%
      \centerline{\rmaa@AbstractHeaderStyle{Resumen}}\par%
      \vskip\@after@abstracthead@skip%
      \if@newstyle\noindent\fi\@resumen
      \if@setresumen\else\@abstract\fi % append english abstract if
                                       % spanish one not set
      \if@guide\else            %do not use keywords for guide
      \vskip\@before@kwd@skip%
      % Format the Key Words with hanging indentation:
      % They are put in a \parbox with width equal to the line width
      % minus the width of the header. 
      { % Widths of keyword
        \settowidth{\@kwhwidth}{\@keywordheader~}   % header and
        \setlength{\@kwtwidth}{\linewidth}          % keyword text
        \addtolength{\@kwtwidth}{-\@kwhwidth}       %
        \noindent\@keywordheader~\parbox[t]{\@kwtwidth}{%
          \rmaa@KeywordStyle{\boldmath\uppercase{\@keywords}}}}%
      \fi
      }
      \fi
  \end{center}%
  \vspace*{\afterkwdskip}
  \par
  }
\newlength{\afterkwdskip}\setlength{\afterkwdskip}{\z@}
%    \end{macrocode}
% \end{macro}
% \end{macro}
%
%
% \begin{macro}{\reviewarticle}
% \begin{macro}{\reviewtoc}
% \begin{macro}{\reviewlof}
%   User command to specify a review article. Must be used in
%   preamble. Main effect is to place the text ``Review'' before the
%   article title.  Also add command for setting a table of contents
%   and a list of figures at the beginning of the article.
%    \begin{macrocode}
\newcommand{\reviewarticle}{\@reviewarticletrue}
\newcommand\reviewtoc{%
  \section*{CONTENTS}%
  \@starttoc{rtc}%
}
\newcommand\reviewlof{%
  \section*{LIST OF FIGURES}%
  \@starttoc{lof}%
}
%    \end{macrocode}
% \end{macro}
% \end{macro}
% \end{macro}
% \begin{macro}{\l@section}
%    \begin{macrocode}
\renewcommand*\l@section[2]{%
  \ifnum \c@tocdepth >\z@
    \addpenalty\@secpenalty
    \addvspace{0.8em \@plus\p@}%
    \setlength\@tempdima{1.5em}%
    \begingroup
      \parindent \z@ \rightskip \@pnumwidth
      \parfillskip -\@pnumwidth
      \leavevmode %\bfseries
      \advance\leftskip\@tempdima
      \hskip -\leftskip
      #1\nobreak\hfil \nobreak\hb@xt@\@pnumwidth{\hss #2}\par
    \endgroup
  \fi}
%    \end{macrocode}
% \end{macro}
%
% \begin{macro}{\@MSmaketitle}
% A special version of |\@maketitle| for the \texttt{manuscript}
% option. This can deal with a multi-page abstract/resumen. Not very
% well, mind. 
%    \begin{macrocode}
\newcommand{\@MSmaketitle}{%
  \if@debug\ClassInfo{rmaa}{Using special MS maketitle command}\fi
  \newpage
  {
    \thispagestyle{rmaafirstpage}
%     \setlength{\parindent}{0pt}
    \null\vskip 2ex
    \begin{center}
      \@title\par
      \@author\par
      \rmaa@ReceivedStyle{%
        \ifx\@received\@rmaaEmpty
          Draft version: \@date
        \else
          Received: \@received
          \ifx\@accepted\@rmaaEmpty
          \else
            ; Accepted: \@accepted
          \fi
        \fi}\par
    \end{center}
    \@thanks
    \centerline{\rmaa@AbstractHeaderStyle{Resumen}}
    \noindent\@resumen\par\vskip 2ex
    \centerline{\rmaa@AbstractHeaderStyle{Abstract}}
    \noindent\@abstract\par\vskip 2ex
    \noindent\@keywordheader~\rmaa@KeywordStyle{\@keywords}\par
%     \@keywordheader~\rmaa@KeywordStyle{\uppercase{\@keywords}}\par
  }%
  }
%    \end{macrocode}
% \end{macro}

% \subsection{The sectioning commands}
% \label{sec:section-commands}

% \begin{macro}{\section}
% \begin{macro}{\subsection}
% \begin{macro}{\subsubsection}
% \begin{macro}{\paragraph}
% \begin{macro}{\subparagraph}
% Redefine all the standard commands to use these lengths. 
%    \begin{macrocode}
\renewcommand\section{\@startsection {section}{1}{\z@}%
                                   {\abovesectionskip}%
                                   {\belowsectionskip}%
                                   %% how do we get upper case here
                                   {\centering\reset@font\normalsize}}
\renewcommand\subsection{\@startsection{subsection}{2}{\z@}%
                                     {\abovesectionskip}%
                                     {\belowsectionskip}%
                                     {\centering\reset@font\normalsize\itshape}}
\renewcommand\subsubsection{\@startsection{subsubsection}{3}{\z@}%
                                     {\abovesectionskip}%
                                     {\belowsectionskip}%
                                     {\centering\reset@font\normalsize\itshape}}
\renewcommand\paragraph{\@startsection{paragraph}{4}{\z@}%
                                    {\aboveparaskip}%
                                    {\belowparaskip}%
                                    {\reset@font\normalsize\bfseries}}
\renewcommand\subparagraph{\@startsection{subparagraph}{5}{\parindent}%
                                      {\aboveparaskip}%
                                      {\belowparaskip}%
                                      {\reset@font\normalsize\bfseries}}
%    \end{macrocode}
% \end{macro}
% \end{macro}
% \end{macro}
% \end{macro}
% \end{macro}

% \begin{macro}{secnumdepth}
% Make sure that we number the subsubsections---this is done by
% default in the |article| class but not the |book| class, which is
% why we must set it explicitly here.
%    \begin{macrocode}
\setcounter{secnumdepth}{3}
%    \end{macrocode}
% \end{macro}
% \changes{2.2beta10}{2001/11/07}{Fixed \cs{secnumdepth} for \cs{book}
% style.}

% \begin{macro}{\thesubsection}
% \begin{macro}{\thesubsubsection}
% For subsection and subsubsection, we want the numbers to be roman
% but the title to be italic. This seems the easiest way to go about it.
%    \begin{macrocode}
\renewcommand{\thesubsection}{{\upshape\thesection.\@arabic\c@subsection}}
\renewcommand{\thesubsubsection}{\thesubsection{\upshape.\@arabic\c@subsubsection}}
%    \end{macrocode}
% \end{macro}
% \end{macro}

% \begin{macro}{\@seccntformat}
% Add ``.'' to the number here so we can remove it from |\thesection|,
% etc. Reduce gap between number and title WJH 27 Jan 1995
%    \begin{macrocode}
\def\@seccntformat#1{\csname the#1\endcsname .\hskip 0.4em}
%    \end{macrocode}
% \end{macro}

% \begin{macro}{\@sect}
% I didn't really want to mess with the |\@sect| command, but it was
% the only way I could manage to force upper case. Other than the
% commented parts, it is identical to the standard version. 
%    \begin{macrocode}
\def\@sect#1#2#3#4#5#6[#7]#8{
  \ifnum #2>\c@secnumdepth
     \let\@svsec\@empty
  \else
     \refstepcounter{#1}%
     \let\@@protect\protect
     \def\protect{\noexpand\protect\noexpand}%
     \edef\@svsec{\@seccntformat{#1}}%
     \let\protect\@@protect
  \fi
%    \end{macrocode}
% This is the bit I changed (27 Jan 1995). Use upper case for
% |\section|, use mixed case for the rest. 
%    \begin{macrocode}
  \ifnum #2=1                   
     \def\sec@upcase##1{\uppercase{##1}}%
     \def\sec@shape{\upshape}%
  \else                         
     \def\sec@upcase##1{\relax{##1}}%
     \def\sec@shape{\itshape}%
  \fi
  \@tempskipa #5\relax
  \ifdim \@tempskipa>\z@
     \begingroup #6\relax
          \@hangfrom{\hskip #3\relax\@svsec}%
                    {\interlinepenalty \@M \sec@upcase{#8}\par}%
     \endgroup
     \csname #1mark\endcsname{#7}%
%    \end{macrocode}
% Only write TOC entries if doing the Guide style. 
%    \begin{macrocode}
     \if@guide 
        \addcontentsline{toc}{#1}{%
          \ifnum #2>\c@secnumdepth \else
             \protect\numberline{\csname the#1\endcsname}%
          \fi
          #7}%
     \fi
%    \end{macrocode}
% Also allow per-article TOC for review papers
%    \begin{macrocode}
     \if@reviewarticle
        \ifnum #2>2 \else
           \addcontentsline{rtc}{#1}{%
             \protect\numberline{\csname the#1\endcsname}%
             \sec@upcase{\sec@shape #7}}%
        \fi
     \fi
  \else
     \def\@svsechd{#6\hskip #3\relax  %% \relax added 2 May 90
       \@svsec #8\csname #1mark\endcsname{#7}%
       \if@guide
          \addcontentsline{toc}{#1}{%
            \ifnum #2>\c@secnumdepth \else
               \protect\numberline{\csname the#1\endcsname}%
            \fi
            #7}%
       \fi
       }%
  \fi
  \@xsect{#5}}
%    \end{macrocode}
% \end{macro}


% \subsection{Appendices}
% \label{sec:appendices}
% Make sure that appendices in individual articles work properly. We
% copy the definition from the article class (the book one doesn't do
% what we want).  
% \begin{environment}{appendices}
%   This uses an environment |appendices| inside of which each
%   |\section| will become a separate appendix.  Note that the
%   appendices must come before the references (otherwise the full
%   addresses mechanism will not work properly). We make the equation
%   numbering be subordinate to the appendix (equations A12, B6,
%   etc). Finally, we have a medskip followed by APPENDICES to
%   introduce them. \changes{3.17}{2002/04/08}{ Added appendices environment
%   that does the right thing.}
%    \begin{macrocode}
\newenvironment{appendices}{\par
  \begingroup
  \setcounter{section}{0}%
  \setcounter{subsection}{0}%
  \renewcommand\thesection{\@Alph\c@section}
  \renewcommand\theequation{\thesection\arabic{equation}}
  \medskip
  \section*{APPENDICES}}
{\endgroup}
%    \end{macrocode}
% \end{environment}
% \begin{environment}{appendix}
%   If there is only one appendix, then use |appendix| environment
%   instead.  \changes{4.0}{2019/11/12}{ Added
%   appendix environment for case where there is only one section.}
%    \begin{macrocode}
\renewenvironment{appendix}{\par
  \begingroup
  \setcounter{section}{0}%
  \setcounter{subsection}{0}%
  \renewcommand\thesection{\@Alph\c@section}
  \renewcommand\theequation{\thesection\arabic{equation}}
  \medskip
  \section*{APPENDIX}}
{\endgroup}
%    \end{macrocode}
% \end{environment}

% \subsection{Floating thingies}
% \label{sec:floats}   

%% 
%    \begin{macrocode}
\renewcommand{\thetable}{\@arabic\c@table}
\renewcommand{\thefigure}{\@arabic\c@figure}
\renewcommand{\figurename}{Fig.}
\renewcommand{\tablename}{TABLE}
%\renewcommand{\fnum@table}{\MakeUpperCase{\tablename}~\thetable}
%    \end{macrocode}
% For some reason, I have to write `TABLE' explicitly here
% If I use |\tablename|, it gets converted to mixed case---God knows
% why. 
% |\figurename| is also being ignored. 
%    \begin{macrocode}
\renewcommand{\fnum@table}{TABLE~\thetable}
\renewcommand{\fnum@figure}{Fig.~\thefigure}
%    \end{macrocode}

% Do captions in 9 pt face in the main journal style and in new-style
% anything. 
%    \begin{macrocode}
\long\def\@makefigcap#1#2{%
  \if@proceedings\else\small\fi
  \if@newstyle\small\fi 
  \vskip\abovecaptionskip
  \sbox\@tempboxa{#1. #2}%
  \ifdim \wd\@tempboxa >\hsize
    #1. #2 \par
  \else
    \global \@minipagefalse
    \hb@xt@\hsize{\hfil\box\@tempboxa\hfil}%
  \fi
  \vskip\belowcaptionskip}
%    \end{macrocode}

% Make sure table caption is upper case. 
%    \begin{macrocode}
\long\def\@maketabcap#1#2{%
  \hbox to\hsize{\hfil #1 \hfil}
  \vskip 6pt 
  \centering\MakeUpperCase{#2}\par
  \vskip 6pt}
%    \end{macrocode}

%    \begin{macrocode}
\newlength{\@rottabcapwidth}
\long\def\@makerottabcap#1#2{%
  \setbox\@tempboxa\hbox{\MakeUpperCase{#2}}%
  \ifdim\wd\@tempboxa >.8\vsize
     \setlength\@rottabcapwidth{.8\vsize}
  \else
     \setlength\@rottabcapwidth{\wd\@tempboxa}
  \fi
  \rotatebox{90}{%
    \begin{minipage}{\vsize}\centering
      \begin{minipage}{\@rottabcapwidth}\centering
        #1\\ \vskip 6pt \MakeUpperCase{#2}
      \end{minipage}%
    \end{minipage}
    }%
  \hskip 6pt
  }

\newenvironment{Sideways}
    {\begin{sideways}\begin{minipage}{\vsize}\centering}
    {\end{minipage}\end{sideways}}

\newenvironment{Sidewaystable}
{%
  \let\oldcaption=\caption%
  \def\caption{
    \def\@captype{table}\renewcommand{\@makecaption}{\@maketabcap}
    \oldcaption}
  \begin{sidewaystable}\centering}
  {\end{sidewaystable}\let\caption=\oldcaption}
\newenvironment{Sidewaystable*}
{%
  \let\oldcaption=\caption%
  \def\caption{
    \def\@captype{table}\renewcommand{\@makecaption}{\@maketabcap}
    \oldcaption}
  \begin{sidewaystable*}\centering}
  {\end{sidewaystable*}\let\caption=\oldcaption}
%    \end{macrocode}

% Not sure why we need to stick these lines in. 
%    \begin{macrocode}
\renewenvironment{figure}
                 {\renewcommand{\@makecaption}{\@makefigcap}\@float{figure}}
                 {\end@float}
\renewenvironment{figure*}
                 {\renewcommand{\@makecaption}{\@makefigcap}\@dblfloat{figure}}
                 {\end@dblfloat}
\renewenvironment{table}
                 {\renewcommand{\@makecaption}{\@maketabcap}\@float{table}}
                 {\end@float}
\renewenvironment{table*}
                 {\renewcommand{\@makecaption}{\@maketabcap}\@dblfloat{table}}
                 {\end@dblfloat}
%    \end{macrocode}

% Redefine |\@caption| so that we don't bother with the list of
% figures. This saves us having to worry about fragile commands in the
% captions. 
%    \begin{macrocode}
\long\def\@caption#1[#2]#3{%
  \par
%%  \addcontentsline{\csname ext@#1\endcsname}{#1}%
%%    {\protect\numberline{\csname the#1\endcsname}{\ignorespaces #2}}%
  \begingroup
    \@parboxrestore
    \if@minipage
      \@setminipage
    \fi
    \normalsize
    \@makecaption{\csname fnum@#1\endcsname}{\ignorespaces #3}\par
  \endgroup}
%    \end{macrocode}

% Make tables not quite so cramped. Allow |\hline\hline| to give a
% poor man's thick rule. 
%    \begin{macrocode}
\renewcommand{\arraystretch}{1.15}
\setlength{\doublerulesep}{0.4pt}
%    \end{macrocode}

% \begin{macro}{\tablecols}
% \begin{macro}{\@tablecols}
% \begin{macro}{\tabcaptioncontinued}
% |\tablecols| is a user command for setting the number of columns. It
% is used by commands that try and span the entire table (I think
% there must be a better way of acheiving
% this). |\tabcaptioncontinued| is 
% a user command for setting the caption of a continued table on the
% next page. 
%    \begin{macrocode}
\newcommand{\tablecols}[1]{\renewcommand{\@tablecols}{#1}}
\newcommand{\@tablecols}{1}
\newcommand{\tabcaptioncontinued}{%
  \multicolumn{\@tablecols}{c}{\addtocounter{table}{-1}TABLE~\thetable\ (CONTINUED)} \\ \\
  }
%    \end{macrocode}
% \end{macro}
% \end{macro}
% \end{macro}

% \subsubsection{Footnotes in tables}
% \label{sec:table-footnotes}
% \begin{macro}{\tablenotemark}
% \begin{macro}{\tablenote}
% \begin{macro}{\tablenoteskip}
% Old user commands |\tablenotemark{MARK}|,
% |\tablenote{MARK}{TEXT}|---now deprecated. 
%    \begin{macrocode}
\newcommand\tablenotemark[1]{$^\mathrm{#1}$} 
\newcommand\tablenote[2]{\if@proceedings\else\small\fi ${}^\mathrm{#1}$ #2} 
\newcommand\tablenoteskip{\\[0.5\baselineskip]} 
%    \end{macrocode}
% \end{macro}
% \end{macro}
% \end{macro}

% \begin{macro}{\tabnotemark}
% \begin{macro}{\tabnotetext}
% \begin{macro}{\tabnotetextTWO}
% Alternative way of doing footnotes in tables. 
%    \begin{macrocode}
\newcommand{\tabnotemark}[1]{\textsuperscript{\lowercase{#1}}}
%\newcommand{\tabnotetextTWO}[4]{\tabnotetext{#1}{#2}\tabnotetext{#3}{#4}}
\newcommand{\tabnotetextTWO}[4]{%
  \multicolumn{\@tablecols}{l}{%
    \parbox{0.5\tabnotewidth}{\small\textsuperscript{#1}#2}%
    \parbox{0.5\tabnotewidth}{\small\textsuperscript{#3}#4}%
    }%
  \\}
\newcommand{\tabnotetext}[2]{%
  \multicolumn{\@tablecols}{@{}p{\tabnotewidth}@{}}{\small\textsuperscript{#1}#2}\\}
\newlength{\tabnotewidth}
%    \end{macrocode}
% \end{macro}
% \end{macro}
% \end{macro}


% \begin{macro}{changemargin}
%    \begin{macrocode}
% Command to allow margins to be changed on the fly. Useful to allow
% tables and Figures to protrude into the margins. 
%
% \changes{3.24}{2007/09/14}{Added \texttt{changemargin} environment
% (copied from TeX FAQ).}
\newenvironment{changemargin}[2]{% 
  % From TeX FAQ http://www.tex.ac.uk/cgi-bin/texfaq2html?label=chngmargonfly
  \begin{list}{}{%
    \setlength{\topsep}{0pt}%
    \setlength{\leftmargin}{#1}%
    \setlength{\rightmargin}{#2}%
    \setlength{\listparindent}{\parindent}%
    \setlength{\itemindent}{\parindent}%
    \setlength{\parsep}{\parskip}%
  }%
  \item[]}{\end{list}}
%    \end{macrocode}
% \end{macro}

% \subsection{The bibliography and references}
% \label{sec:biblio}
% \begin{macro}{\beforebiblioskip}
% \begin{macro}{\afterbiblioskip}
% I'm not entirely sure why I introduced these two commands, since
% directly putting |\vspace| before and after
% |\begin{thebibliography}| seems to work well enough. 
%    \begin{macrocode}
\newlength{\beforebiblioskip}
\setlength{\beforebiblioskip}{0pt}
\newlength{\afterbiblioskip}
\setlength{\afterbiblioskip}{0pt}
%    \end{macrocode}
% \end{macro}
% \end{macro}
% \begin{macro}{thebibliography}
% The environment itself is done as starred section with extra
% vertical space after heading. This space is currently zero. 
% In the book style we do the cunning trick with rereading the file in
% order to get the references into the |\ADS@references| macro.
% 
% \changes{3.05}{2003/01/02}{Removed the \cs{beforebiblioskip} and
% \cs{afterbiblioskip} commands, since they cause unwanted vertical
% space in some instances.}  
%    \begin{macrocode}
\renewenvironment{thebibliography}{%
  \if@book\RereadBiblio\fi
%%  \vspace*{\beforebiblioskip}
  \section*{REFERENCES} \frenchspacing
%%  \vspace*{\afterbiblioskip}
  \small
%    \end{macrocode} 
% The references themselves are set as a list in order 
% to get the hanging indentation. They are done in a smaller font and
% with french spacing so we needn't worry about the distinction
% between sentence-ending periods and abbreviations. 
%    \begin{macrocode} 
  \list{\null}{\leftmargin 1.5em\labelwidth\z@\labelsep\z@\itemindent -1.5em
    \itemsep\z@\parsep\z@\usecounter{enumi}}
  \sloppy\clubpenalty4000\widowpenalty4000
  \sfcode`\.=1000\relax\@biblisttrue}
{\endlist\@biblistfalse\normalsize\@outputaddressesifnecessary}
%    \end{macrocode}
% \end{macro}
% \subsection{Uniquification of cross-references}
% \label{sec:uniquification}
% 
% \begin{macro}{\articlename}
% In order that the |\label{}| command in individual articles be
% guaranteed to give a unique anchor, we append the article name (held
% in the |\articlename| command) to each key. This requires
% redefinition of the following commands: |\cite|, |\nocite|, |\ref|,
% |\pageref|, |\label|. 
% \begin{macro}{\oldcite}
% \begin{macro}{\oldnocite}
% \begin{macro}{\oldref}
% \begin{macro}{\oldpageref}
% \begin{macro}{\oldlabel}
% The original versions of these commands are
% saved as |\oldcite|, |\oldnocite|, etc. 
%
% The uniquifying of cross-references is now always done, so 
% we don't get multiply-defined labels in the single article case
%    \begin{macrocode}
\newcommand{\articlename}{anon}
\AtBeginDocument{
  \let\oldcite=\cite 
  \let\oldnocite=\nocite 
  \let\oldref=\ref 
  \let\oldpageref=\pageref
  \let\oldlabel=\label
}
%    \end{macrocode}
% \end{macro}
% \end{macro}
% \end{macro}
% \end{macro}
% \end{macro}
% \end{macro}
%
% Uniquify all cross-references by prepending the article name.  This
% now has to go into |\AtBeginDocument| so it is working on the new
% versions that hyperref.sty has set up.  Then we need to put the
% |\@begofarticlehook| afterwards, so that the first page label is using
% the new definitions.  If amsmath is loaded, then we also need to patch
% the special label command that they use. \changes{4.0}{2019/11/18}{Fix labels for amsmath display equations, in case that package has been loaded.}
%    \begin{macrocode}
\AtBeginDocument{
  \renewcommand{\ref}[1]{\oldref{\articlename:#1}%
    \if@debug{\ClassInfo{rmaa}{Reference to \articlename:#1}}\fi}
  \renewcommand{\pageref}[1]{\oldpageref{\articlename:#1}%
    \if@debug{\ClassInfo{rmaa}{Page reference to \articlename:#1}}\fi}
  \renewcommand{\label}[1]{\oldlabel{\articlename:#1}%
    \if@debug{\ClassInfo{rmaa}{Label for \articlename:#1}}\fi}%
  \ifdef{\label@in@display}{%
    \let\oldlabel@in@display=\label@in@display
    \renewcommand{\label@in@display}[1]{\oldlabel@in@display{\articlename:#1}%
      \if@debug{\ClassInfo{rmaa}{AMSTeX Label for \articlename:#1}}\fi}%
  }{}
  \@begofarticlehook
}%
%    \end{macrocode}
% Citations are more complicated since they can contain a list of
% keys. These are really only here for use with the rmaacite package,
% which is utterly deprecated (02 Jan 2003) now that I have the natbib
% code working. The command \cs{cite} should \emph{not} be used with
% natbib: use \cs{citep}, \cs{citet}, etc.\@ instead. 
%    \begin{macrocode}
\renewcommand{\nocite}[1]{\@bsphack
  \@for\@citeb:=#1\do{%
    \edef\@citeb{\expandafter\@firstofone\@citeb:\articlename}%
    \if@filesw\immediate\write\@auxout{\string\citation{\@citeb}}\fi
    \@ifundefined{b@\@citeb:\articlename}{\G@refundefinedtrue
      \@latex@warning{Citation `\@citeb' undefined}}{}}%
  \@esphack}
\renewcommand{\cite}[1]{\oldcite{#1:\articlename}}
%   \end{macrocode}


% Horrible fiddling with the bibliography citation mechanism. 
% This has to be globally available in case rmaacite puts it in the
% \texttt{.aux} file. 
%    \begin{macrocode}
\def\lbibcite#1#2#3{\global\@namedef{b@#1}{#2}
                    \global\@namedef{by@#1}{#3}}
%    \end{macrocode}

% \begin{macro}{\oldbibitem}
% \begin{macro}{\@rmaaBibArg}
% \begin{macro}{\@rmaaSpace}
% \begin{macro}{rmaabib}
% Of course, |\bibitem| needs to get the uniquification treatment as
% well. While we are at it, we redefine it to use consecutive integers
% for the key if the user leaves it blank (using the counter
% |rmaabib|). This allows one to put 
% |\bibitem{}| for all your references without getting annoying
% messages about multiply-defined labels.
% 
% We also take the opportunity to insert the optional argument to
% \cs{bibitem} if it is missing in the form Author(DATE) so that natbib
% doesn't try to do numerical citations. Natbib already provides a
% hook \cs{@extra@b@citeb} to help uniquify the citation key, so we
% just plug into that. However, it puts it \emph{after} the original
% key, not before as I had had it. Easiest solution is to change the
% way that I do it. 02 Jan 2003. Note that the \cs{@extra@b@citeb}
% does not get written to the .aux  file in the \cs{citation} command
% but it does work. 
% \changes{3.14}{06 Feb 2003}{Reinstalled old definition of \cs{bibitem}
%  for the case where the nonatbib option is requested, so that the
%  rmaacite package will work. This means that you can't use rmaacite
%  and natbib in the same volume.}
%    \begin{macrocode}
\if@natbib\def\@extra@b@citeb{:\articlename}\fi
%%\else
\let\oldbibitem\bibitem
\newcommand{\@rmaaBibArg}{}
\newcommand{\@rmaaSpace}{ }
\newcounter{rmaabib}
\if@natbib
\renewcommand{\bibitem}[2][ABC(2003)]{%
  \renewcommand{\@rmaaBibArg}{#2}
  \ifx\@rmaaBibArg\@rmaaEmpty
     \addtocounter{rmaabib}{1}
     \oldbibitem[#1]{\thermaabib:\articlename}
  \else\ifx\@rmaaBibArg\@rmaaSpace
     \addtocounter{rmaabib}{1}
     \oldbibitem[#1]{\thermaabib:\articlename}
  \else 
     \oldbibitem[#1]{#2:\articlename}
  \fi\fi
  \renewcommand{\@rmaaBibArg}{}
  }
\else
\renewcommand{\bibitem}[1]{%
  \renewcommand{\@rmaaBibArg}{#1}
  \ifx\@rmaaBibArg\@rmaaEmpty
     \addtocounter{rmaabib}{1}
     \oldbibitem{\articlename:\thermaabib}
  \else\ifx\@rmaaBibArg\@rmaaSpace
     \addtocounter{rmaabib}{1}
     \oldbibitem{\articlename:\thermaabib}
  \else
     \oldbibitem{\articlename:#1}
  \fi\fi
  \renewcommand{\@rmaaBibArg}{}
  }
\fi
%%\fi
%    \end{macrocode}
% \end{macro}
% \end{macro}
% \end{macro}
% \end{macro}

% \begin{macro}{oldpcite}
% \begin{macro}{oldscite}
% Then, we need to repeat the above scheme, but for the |rmaacite|
% definition. 
%    \begin{macrocode}
\newcommand{\@redefinebibitem}{%
  \let\oldbibitem\bibitem
  \let\oldcite\cite
  \let\oldpcite\pcite
  \let\oldscite\scite
  \def\bibitem[##1<##2>]##3{%
    \renewcommand{\@rmaaBibArg}{##3}
    \ifx\@rmaaBibArg\@rmaaEmpty
       \addtocounter{rmaabib}{1}
       \oldbibitem[##1<##2>]{\thermaabib:\articlename}
    \else\ifx\@rmaaBibArg\@rmaaSpace
       \addtocounter{rmaabib}{1}
       \oldbibitem[##1<##2>]{\thermaabib:\articlename}
    \else 
       \oldbibitem[##1<##2>]{##3:\articlename}
    % trust the author here
%%       \oldbibitem[##1<##2>]{##3}
    \fi\fi
    }
  }
%    \end{macrocode}
% \end{macro}
% \end{macro}


% \subsection{HTML Table of Contents}
% \label{sec:hyper-TOC}
% \begin{macro}{\HyperTOCEntry}
% \begin{macro}{\HyperTOCpart}
% \begin{macro}{\HyperTOCmeeting}
% \begin{macro}{\HyperTOCeditors}
% \begin{macro}{\HyperTOCpreamble}
% These commands are used to write information to the \texttt{.htoc}
% file, which is used in generating a Web Page of the Table of
% Contents. Since they only ever appear preceded by |\protect| in the
% argument to an |\addtocontents| command,  they should never need to
% be expanded by the macros. However, we define them here just in case
% \dots
%    \begin{macrocode}
\newcommand{\HyperTOCEntry}[3]{\relax}
\newcommand{\HyperTOCpart}[1]{\relax}
\newcommand{\HyperTOCmeeting}[3]{\relax}
\newcommand{\HyperTOCeditors}[1]{\relax}
\newcommand{\HyperTOCpreamble}{\relax}
%    \end{macrocode}
% \end{macro}
% \end{macro}
% \end{macro}
% \end{macro}
% \end{macro}
%

% \subsection{Formatting the poster abstracts} 
% These can now be used without using the book option.
%    \begin{macrocode}

\newcommand{\@makeminititle}{%
  \newpage
  \null 
  \vskip 3em%
  \begin{center}%
    {\@title \par}%
  \end{center}
  \vskip 3em%
  }

\newif\if@abstractsads
\@abstractsadsfalse
\newcommand\abstractsADS{\@abstractsadstrue}

\if@book\RequirePackage{multicol}\fi
\newenvironment{abstracts}[1][Abstracts of Contributed Papers]{%
  %% First, make sure we start on an odd page (06 Jun 2000)
  \if@book\clearpage\thispagestyle{empty}\cleardoublepage\fi
  %%
  \renewcommand{\articlename}{ABSTRACTS}
  \title{#1}\shortauthor{Abstracts}\shorttitle{Abstracts}
  \listofauthors{}
  \if@book
  \addcontentsline{toc}{chapter}{\protect\@fiddleauthors{\@TOCtitle}{\rmaaTOCauthors}}%
  \addtocontents{htc}{\protect\HyperTOCEntry{\@rawtitle}{\@rawauthors}{\articlename}}
  \fi
  \pagestyle{rmaa} % this now seems necessary (06 Jun 2000)
  \thispagestyle{rmaafirstpage}
%%  \twocolumn[\@makeminititle]
  \renewcommand{\abstract}[1]{%
    \renewcommand\ADS@abstract{##1}%
    \renewcommand{\@abstract}{##1}\@writeabstract%
  }
  \renewcommand{\title}[2][]{%
    \gdef\ADS@authors{}%
    \gdef\ADS@affil{}%
    \gdef\ADS@references{}%
    \gdef\ADS@keywords{}%
    \renewcommand{\@rmaaArg}{##1}
    \ifx\@rmaaArg\@rmaaEmpty
       \renewcommand{\@TOCaftertitle}{}
    \else
       \renewcommand{\@TOCaftertitle}{ (\@rmaaArg)}
    \fi
    \renewcommand{\@title}{\rmaa@TOCTitleStyle{##2}}%
    \renewcommand{\@TOCtitle}{\rmaa@TOCTitleStyle{##2}\@TOCaftertitle}%
    \renewcommand{\rmaaIndexTitle}{\rmaa@IndexTitleStyle{##2\@TOCaftertitle.}}%
    \renewcommand\ADS@title{##2}%
    \xdef\ADS@firstpage{\@arabic{\c@page}}%
    }
  \renewcommand{\author}[1]{\let\@thanks\@empty\renewcommand\@author{##1} }
  \renewcommand{\footnotetext}[2][XX]{\par\noindent\textsuperscript{##1}~##2}
  \renewcommand{\footnoterule}{\kern-3\p@ \hrule \@width 3em \kern 2.6\p@}
  \if@book\@begofarticlehook\fi
  \if@book\begin{multicols}{2}[\@makeminititle]\else\@makeminititle\fi
  \begingroup
  \setcounter{footnote}{0}
  }% and at the end of the environment....
  {\endgroup\if@book\end{multicols}\listofabstracts\label{RMAALastPage}\clearpage\fi}
%% auxiliary commands used:
%% \@writeabstract does all the actual formatting 
\newcommand{\@writeabstract}{%
  \if@book%
  \addcontentsline{loa}{abstract}{\protect\@fiddleauthors{\@TOCtitle}{\rmaaTOCauthors}}%
  \fi%
  \begin{center}\@title\par\@author\end{center}
  \@abstract
  \par\noindent\footnoterule\nobreak\footnotesize\@thanks
  \par\noindent\hrulefill\goodbreak
  \vskip 1cm \@plus0.5cm
  %% write ADS entry
  \if@book\if@abstractsads\ADS@dumpdata\fi\fi%
  \endgroup % this pair is to make each abstract a self-contained group
  \begingroup
  \setcounter{footnote}{0}
  }

\if@book
%% List of abstracts (stuff written to .loa file)
\newcommand{\listofabstracts}{%
  \begin{multicols}{2}[%
    \vskip 1.5cm \vfill % leave at least 1.5 cm and go to bottom of page if possible
    \centerline{\MakeUpperCase{List of Abstracts}}
    \vskip 1.5cm \@plus0.5cm
    ][6cm] % this makes it start a new page if there is not at least 6cm of space
    \renewcommand{\@fiddleauthors}[2]{##1\hskip1em plus2em minus0.4em\penalty-100  ##2}
    \parindent\z@
%%    \parskip.5\p@ \@plus .7\p@\relax 
    \parskip3\p@ \@plus 1\p@ \@minus 1\p@\relax % allow quite a lot of
                                % leeway here to avoid breaks inside
                                % entries 
    \renewcommand{\\}{\relax} % temporarily turn off explicit line breaks just in case
    \renewcommand{\altaffilmark}[1]{\relax} % ditto with this
    \@starttoc{loa}
  \end{multicols}
  }
\newcommand*{\l@abstract}[2]{%
  \par\penalty-1000
  #1\hskip 1em\penalty0\null\nobreak\hfill#2}
\fi

%    \end{macrocode}
 
% \subsection{Code for typesetting an entire volume of conference
% proceedings} \label{sec:whole-volume}
%    \begin{macrocode}

\if@book
%% environment for including the individual articles
\newcommand{\@fiddleauthors}[2]{\relax}
\renewcommand{\theequation}{\@arabic\c@equation}
\newenvironment{articles}%
{ \begingroup % make it a group so that author definitions do not
              % spill over
  \gdef\ADS@authors{}
  \pagestyle{rmaa}
  \renewenvironment{document}{%
    \@begofarticlehook
    \addtocounter{chapter}{1}%
    \addcontentsline{toc}{chapter}{%
      \protect\@fiddleauthors{\@TOCtitle}{\rmaaTOCauthors}}%
    \addtocontents{htc}{\protect\HyperTOCEntry{\@rawtitle}{\@rawauthors}{\articlename}}
    }{\relax}
  \renewcommand{\documentclass}[2][foo]{\relax} 
  \renewcommand{\usepackage}[2][foo]{%
    \def\@packagename{##2}\def\@rmaacitename{rmaacite}%
    \ifx\@packagename\@rmaacitename
    \@rmcitetrue\ClassInfo{rmaa}{This paper uses rmaacite}
    \makeatletter\input{rmaacite.sty}\makeatother % read it in here so that all changes
    \else                                         % are local to paper
    \if@debug\ClassWarning{rmaa}{Package \@packagename\space was requested by \articlename\space}\fi
    \fi} 
  %% Don't allow individual articles to mess with the pages, volume or year
  \renewcommand{\SetFirstPage}[1]{\relax}
  \renewcommand{\SetYear}[1]{\relax}
  \renewcommand{\SetVolume}[1]{\relax}
  \renewcommand{\SetConfTitle}[1]{\relax}
  %% Don't allow individual articles to change these
  \renewcommand{\nonstopmode}{\relax}
  \renewcommand{\scrollmode}{\relax}
  \renewcommand{\outputplates}{\relax}
  %% XXXX need to check if using rmaacite and act accordingly
  %% avoid multiple declarations of the same length
  \let\oldnewlength=\newlength
  \renewcommand{\newlength}[1]{%
    \edef\reserved@a{\expandafter\@gobble\string ##1}%
    \@ifundefined{\reserved@a}{%
      \oldnewlength{##1}}{%
      \if@debug{\ClassInfo{rmaa}{\reserved@a\space already defined: doing nothing}}\fi%
      \relax}}%
  \renewcommand\newcommand{\providecommand}
  }%
{\endgroup\onecolumn}
     

%    \end{macrocode}
% Command to be used in `articles' environment for including an
% article
%    \begin{macrocode}
\newcommand{\article}[2][]{%
  % optional first argument is extra text to put in brackets after title in TOC
  % e.g. ``Invited Talk'' or ``Poster''
  \renewcommand{\articlename}{#2}
  % reset loads of counters
  \setcounter{footnote}{0}
  \setcounter{figure}{0}
  \setcounter{table}{0}
  \setcounter{section}{0}
  \setcounter{equation}{0}
  \setcounter{rmaabib}{0}
  \setcounter{author}{0}
  \setcounter{AUTHOR}{1}
  \gdef\ADS@authors{}%
  \gdef\ADS@affil{}%
  \gdef\ADS@affiliations{}%
  \gdef\ADS@references{}%
  \@firstauthortrue
  \@firstaffiltrue
  \@firstaltaffiltrue
  \@firstkwdtrue
  \@rmcitefalse
  \if@suppressfulladdresses\else\@needaddressestrue\fi
  \renewcommand{\@rmaaArg}{#1}
  \ifx\@rmaaArg\@rmaaEmpty
     \renewcommand{\@TOCaftertitle}{}
  \else
     \renewcommand{\@TOCaftertitle}{ (#1)}
  \fi
%    \end{macrocode}
% In the main journal articles must start on an odd page. 
%    \begin{macrocode}
  \if@journal\clearpage\thispagestyle{empty}\cleardoublepage\fi
%% Now input the file, then run the end-of-article hook (reinstated 28
  %% Sep 2001)
  {\input{\@articledir/\articlename.tex}\@endofarticlehook}}
%    \end{macrocode}
% Command to separate articles into sessions. Write an entry in the
% TOC, makes sure we start on an odd page (added 06 Jun 2000), then do
% a title page if \cs{SessionTitlePages} was called. 
% \changes{3.12}{04 Feb 2003}{Added option to have title page for each
% session. Turn it on with \cs{SessionTitlePages}.}
%    \begin{macrocode}
\newcommand{\session}[1]{
  %% this does nothing except write an entry in the TOC
  \addcontentsline{toc}{part}{\rmaa@TOCPartStyle{#1}}
  \addtocontents{htc}{\protect\HyperTOCpart{#1}}
  %% and also makes sure we start on an odd page (added 06 Jun 2000)
  \clearpage\thispagestyle{empty}\cleardoublepage
  \stepcounter{part}
  \if@sessiontitlepages{%
    \onecolumn
    \thispagestyle{empty}
    \renewcommand{\articlename}{Section\thepart TitlePage}
    \label{RMAAFirstPage}
    \make@session@title{#1}
    \label{RMAALastPage}
    \clearpage\thispagestyle{empty}\cleardoublepage
  }\fi
  }
\newcommand{\make@session@title}[1]{%
  {%
    \setlength\parindent{0pt}
    \vspace*{\stretch{1}}
    \rule{\linewidth}{1mm}
    \begin{flushright}\Huge\bfseries\boldmath
      #1
    \end{flushright}
    \rule{\linewidth}{1mm}
    \vspace*{\stretch{2}}
  }
}
%    \end{macrocode}
% Command to turn on separate title pages for each session. 
%    \begin{macrocode}
\newcommand{\SessionTitlePages}{\@sessiontitlepagestrue}
\newcommand{\BoldSessionTitles}{%
  \renewcommand{\rmaa@TOCPartStyle}[1]{\textbf{\rmaa@TOCTitleStyle{##1}}}
  }
\newcommand{\ArticleDir}[1]{\renewcommand{\@articledir}{#1/}\graphicspath{{1/}}}
\newcommand{\@articledir}{.}
%\newcommand{\ps@rmaafront}{\ps@rmaa}% page style for front matter (same as ever, actually)
\newcommand{\ps@rmaafront}{% page style for front matter:  page # in each foot
  \renewcommand\@oddfoot{\hfil{\thepage}}%
  \renewcommand\@evenfoot{{\thepage}\hfil}%
  % shorttitle in both heads
  \renewcommand\@evenhead{\hfil\@shorttitle\hfil}%
  \renewcommand\@oddhead{\@evenhead}
  }
\newcommand{\ps@rmaafrontfirstpage}{% page style for first page of front matter
  % page # in each foot, on outside edge
  \renewcommand\@oddfoot{\hfil\thepage}%
  \renewcommand\@evenfoot{\thepage\hfil}%
  % nothing in the even heads
  \renewcommand\@evenhead{\relax}%
  \renewcommand\@oddhead{\relax}%
  }
%% slight modification to the standard \frontmatter and \mainmatter
%% commands, just to set the \pagestyle
\renewcommand\frontmatter{%
  \latexonecolumn
  \cleardoublepage
  \@frontmattertrue
  \@mainmatterfalse
  \pagenumbering{roman}
  \pagestyle{rmaafront}
  \addtocounter{page}{2} % leave space for the title page
  }
\renewcommand\mainmatter{%
  \pagestyle{empty}
  \if@rmtwocol\twocolumn\fi
  \cleardoublepage
  \@frontmatterfalse
  \@mainmattertrue
  \pagenumbering{arabic}
  \pagestyle{rmaa}
  }
\renewcommand\backmatter{%
  \clearpage\pagestyle{empty}\cleardoublepage
  \@mainmatterfalse
  %% put extra space in the TOC
  \addtocontents{toc}{\bigskip}
  }
%    \end{macrocode}
%
% \subsection{Front sections}\label{sec:front-sections}
% \begin{macro}{\notitle}
% \begin{macro}{\shortname}
% \begin{macro}{frontsection}
% The environment \texttt{frontsection} is for doing a section in the
% front matter (preface, list of participants, or whatever). It has 
% one optional argument:
% \begin{enumerate}
% \item[\#1] Short title (no embedded spaces) for
% use in making PS files. E.g., ``list\_of\_participants''.   
% \end{enumerate}
% and two mandatory arguments: 
% \begin{enumerate}
% \item[\#2] Title of section (e.g., `Preface')
% \item[\#3] Commands to be evaluated before processing the contents
% of the environment. This is mainly useful for putting in the things
% that would go in the preamble of a normal article  (e.g.,
% |\listofauthors|, |\indexauthor|, |\notitle|) 
% \end{enumerate}
%    \begin{macrocode}
\newcommand{\notitle}{\@titlevisiblefalse} \newcommand{\shortname}{}
\newenvironment{frontsection}[3][]{%
  \begingroup
  \renewcommand{\shortname}{#1}
  \ifx\shortname\@rmaaEmpty
     \renewcommand{\articlename}{#2}
  \else
     \renewcommand{\articlename}{\shortname}
  \fi
  \@begofarticlehook
  \@titlevisibletrue % default
  \thispagestyle{rmaafrontfirstpage}
  \listofauthors{} % blank out this in case it is not specified
  \title{{#2}}\shorttitle{#2}\shortauthor{#2}
  #3 % evaluate the command given as second mandatory arg
  \addcontentsline{toc}{chapter}{\protect\@fiddleauthors{\@TOCtitle}{\rmaaTOCauthors}}%
  \addtocontents{htc}{\protect\HyperTOCEntry{\@rawtitle}{\@rawauthors}{\articlename}}
  \if@titlevisible
     \vspace*{1cm}\centerline{\rmaa@FrontSecStyle{\@title}}\vspace*{1cm}
  \fi
  \setlength{\parskip}{3ex}
  }{\frontlabel{RMAALastPage}\clearpage\endgroup}
%    \end{macrocode}
% \end{macro}
% \end{macro}
% \end{macro}
%    \begin{macrocode}
%% Definition of the command \frontlabel, which rights out something like `F4' to
%% the .aux file, instead of `\textlatin  {i}'. This is so it is easy to post-process
%% the file for making PS files of individual sections
\newcommand{\frontlabel}[1]{%
  \@bsphack
  \protected@write\@auxout{}%
     {\string\newlabel{\articlename:#1}{{}{F\@arabic{\c@page}}}}
  \@esphack
  }
%% User command to say X, Y, Z,\\ Editors of the Proceedings
\newcommand{\EditorSignature}{\vspace*{\parskip}\rightline{\@editors}%
  \rightline{Editors of the proceedings}}
%% Fiddle with \thesection, etc. so that they don't include the chapter number
%% Basically go back to the `article' definitions of these
\renewcommand{\thesection}{\@arabic\c@section}
\renewcommand{\thefigure}{\@arabic\c@figure}
\renewcommand{\thetable}{\@arabic\c@table}

\newcommand{\@oldargs}{}
\newcommand{\@newargs}{}
%% add TOC entry for each paper
%\renewcommand{\title}[1]%
%{\renewcommand{\@title}{\rmaa@TitleStyle{#1}}%
%\addtocounter{chapter}{1}\addcontentsline{toc}{chapter}{#1}%
%}

%% Formatting of the table of contents
%% dimensions of various parts
\newlength{\TOCtextwidth}\setlength{\TOCtextwidth}{12.65cm}
\newlength{\TOCpagenumwidth}\setlength{\TOCpagenumwidth}{1.3cm}
\newlength{\TOCentrywidth}\setlength{\TOCentrywidth}{\TOCtextwidth}
\addtolength{\TOCentrywidth}{-\TOCpagenumwidth} 
\newlength{\TOCextramargin}\setlength{\TOCextramargin}{1.25cm}
\newcommand{\TOCbox}[1]{\hspace{\TOCextramargin}\parbox{\TOCtextwidth}{#1}\par}
\newlength{\TOCextramarginx}\setlength{\TOCextramarginx}{\TOCextramargin}\addtolength{\TOCextramarginx}{4.3cm}
\newlength{\TOCtextwidthx}\setlength{\TOCtextwidthx}{\TOCtextwidth}\addtolength{\TOCtextwidthx}{-4.3cm}
\newcommand{\TOCboxx}[1]{\hspace{\TOCextramarginx}\parbox{\TOCtextwidthx}{#1}\par}

\newcommand{\RMAAbannerimage}{rmsc}
\renewcommand\tableofcontents{%
  \renewcommand{\articlename}{CONTENTS}
  \begingroup
  \@begofarticlehook
  \if@twocolumn
     \@restonecoltrue\onecolumn
  \else
     \@restonecolfalse
  \fi
  %% 9pt for journal style
  \if@proceedings\else\small\fi
  \if@newstyle\small\fi %% same for new style
  \if@proceedings\thispagestyle{rmaafrontfirstpage}\else\thispagestyle{empty}\fi
  \shortauthor{Contents}\shorttitle{Contents}
  \vspace*{18mm}
  \if@printbanner
  %% put in the banner header (need to get better image of this)
     \TOCbox{\hspace*{-0.15cm}\includegraphics[width=0.9\TOCtextwidth]{\RMAAbannerimage}}
     \vspace*{8mm}
  %% volume and publication date (should really be in funny font)
     \TOCboxx{\@volumen\\[2mm] \@mes\ \@anyo}
  \else
  %% put in the banner header (need to get better image of this)
     \TOCbox{\hspace*{-0.15cm}\includegraphics[width=0.9\TOCtextwidth]{\RMAAbannerimage}}
     \vspace*{12.5mm}
  %% volume and publication date (should really be in funny font)
     \TOCboxx{~\\[2mm] ~}
  \fi
  \vspace*{16mm}
  %% This is the command that fiddles with the TOC entry so that the
  %% title is flushleft but the authors are flushright
  \renewcommand{\@fiddleauthors}[2]{%
    % first put the paper title
    \renewcommand{\\}{\relax} % temporarily turn off explicit line breaks just in case
    \renewcommand{\altaffilmark}[1]{\relax} % ditto with this
    \renewcommand{\textsuperscript}[1]{\relax} % ditto with this 
    ##1
    % now save the author list in a box
    \sbox\@tempboxa{##2}%
    % check that box is shorter than the total width available for the entry
    \ifdim \wd\@tempboxa <\TOCentrywidth
       % if so, then put it right here, flush against right margin
       \penalty10000\hskip1em plus10em\penalty0\hbox{}\nobreak\hfill\box\@tempboxa
       % Note that we force a space of at least 1em between title and authorlist, 
       % stretchable to 11em since we may have very short titles but still with an 
       % authorlist too long to go on the same line.  The only linebreak allowed is 
       % directly before the authorlist.
    \else
       % if it is too long, we set the original authorlist text in its own parbox
       \par\vskip1ex\parbox[b]{\TOCentrywidth}{\raggedleft ##2}
    \fi
    }
  \@starttoc{toc}%
  \if@restonecol\twocolumn\fi
  \frontlabel{RMAALastPage}
  \clearpage
  \thispagestyle{empty}
  \cleardoublepage % make sure we start on an odd page
  \endgroup

  %%% ALSO, do the HTML TOC
  %% The file that will be used to make the HTML TOC
  \@starttoc{htc}
  %% All the stuff in this file does nothing in the book style
  \addtocontents{htc}{\protect\HyperTOCmeeting{\@conftitle}{\@reunionplace}{\@reuniondate}}
  \addtocontents{htc}{\protect\HyperTOCeditors{\@editors}}
  \addtocontents{htc}{\protect\HyperTOCpreamble}
  \addtocontents{htc}{\protect\HyperTOCEntry{Table of Contents}{}{\articlename}}

  %% and the ADS file
  \ADS@openfile
  }
%% This is now much simpler
\renewcommand\l@chapter[2]{%TOC entry for articles
  \ifnum \c@tocdepth >\m@ne
    \vskip 1.0em \@plus\p@
    \TOCbox{%
      \parbox[b]{\TOCentrywidth}{#1}\parbox[b]{\TOCpagenumwidth}{\hfill#2}
    }
  \fi} 
\renewcommand\l@part[2]{%TOC entry for sessions: centered, no page number
  \ifnum \c@tocdepth >\m@ne
    \vskip 1.5em \@plus1.5em \goodbreak
    \TOCbox{%
      \parbox[b]{\TOCentrywidth}{\hfil#1\hfil}\hfill
    }
  \fi\nobreak} 

%% Formatting theindex
%% strongly encourage page breaks between index entries
\renewcommand{\@idxitem}{\par\penalty-1000\hangindent.45cm} 
\renewcommand{\indexspace}{\par\penalty-1000\vskip 10\p@ \@plus15\p@ \@minus3\p@\relax}
\renewenvironment{theindex}{%
  %% first, make sure we start on an odd page (06 Jun 2000)
  \clearpage\thispagestyle{empty}\cleardoublepage
  %% this looks a bit like an article but is two columns and without abstract, etc
  \renewcommand{\articlename}{INDEX}
  \title{Author Index}\shortauthor{Author Index}\shorttitle{Author Index}
  \listofauthors{}
  \addcontentsline{toc}{chapter}{\protect\@fiddleauthors{\@TOCtitle}{\rmaaTOCauthors}}%
  \addtocontents{htc}{\protect\HyperTOCEntry{\@rawtitle}{\@rawauthors}{\articlename}}
  \pagestyle{rmaa}
  \thispagestyle{rmaafirstpage}
  \@begofarticlehook
  %% 9pt for journal style
  \if@proceedings\else\small\fi
  \if@newstyle\small\fi
  \twocolumn[\@makeminititle]
  %% vertical spacing (see also definition of \@idxitem and \indexspace
  \parskip\z@ \@plus .6\p@\relax % allow quite a lot of leaway here to avoid breaks inside entries
  %% horizontal spacing, breaking and glue parameters (see also the stuff in the definition 
  %% of \indexauthor
  \parindent\z@ 
  \exhyphenpenalty10000 % don't break on explicit hyphens, it always looks bad
  \rightskip=0pt plus.2em % allow a tiny bit of leeway at the ends of lines
  \spaceskip.3333em plus.5em minus.15em % a lot of leeway between words
  \xspaceskip.5em plus.5em minus.4em % this seems to have no effect
  \let\item\@idxitem
  \renewcommand{\\}{\relax} % temporarily turn off explicit line breaks just in case
  \renewcommand{\altaffilmark}[1]{\relax} % ditto with this
  \renewcommand{\textsuperscript}[1]{\relax} % ditto with this
  }{\label{RMAALastPage}\clearpage} % clearpage necessary so that we know about \shortauthor



%% end definitions for book styles
\fi

%    \end{macrocode}
% In order to do away with plates, we sometimes want a figure to
% appear in color only in the online version. The following command
% can bve used in a figure caption to typeset a note to this effect. 
%    \begin{macrocode}
\newcommand\ColorOnlineOnly{%
  \textsc{Note: this figure is available in color in the electronic
    version of this article, obtainable from}
  \texttt{http://\discretionary{}{}{}www.\discretionary{}{}{}astroscu.\discretionary{}{}{}unam.\discretionary{}{}{}mx/\discretionary{}{}{}\~{}rmaa/}.}
%% Plates
%% Define 2 counters: plate is the number in the whole document, plateloop is used 
%% internally by \outputplates
\newcounter{plate}
\newcounter{plateloop}
%% Currently limited to 15 plates in document
%% when we do a label, we need to write to \newlabel entries to the .aux file, one 
%% that says RMAAPlateX so that it can be referred to in \ouputplates, and one that
%% says articlename:RMAAPlateX so that the dvips program can find it and asociate it
%% with the right article
\newcommand{\@platelabeli}{\oldlabel{RMAAPlate1}\label{RMAAPlate1}}
\newcommand{\@platelabelii}{\oldlabel{RMAAPlate2}\label{RMAAPlate2}}
\newcommand{\@platelabeliii}{\oldlabel{RMAAPlate3}\label{RMAAPlate3}}
\newcommand{\@platelabeliv}{\oldlabel{RMAAPlate4}\label{RMAAPlate4}}
\newcommand{\@platelabelv}{\oldlabel{RMAAPlate5}\label{RMAAPlate5}}
\newcommand{\@platelabelvi}{\oldlabel{RMAAPlate6}\label{RMAAPlate6}}
\newcommand{\@platelabelvii}{\oldlabel{RMAAPlate7}\label{RMAAPlate7}}
\newcommand{\@platelabelviii}{\oldlabel{RMAAPlate8}\label{RMAAPlate8}}
\newcommand{\@platelabelix}{\oldlabel{RMAAPlate9}\label{RMAAPlate9}}
\newcommand{\@platelabelx}{\oldlabel{RMAAPlate10}\label{RMAAPlate10}}
\newcommand{\@platelabelxi}{\oldlabel{RMAAPlate11}\label{RMAAPlate11}}
\newcommand{\@platelabelxii}{\oldlabel{RMAAPlate12}\label{RMAAPlate12}}
\newcommand{\@platelabelxiii}{\oldlabel{RMAAPlate13}\label{RMAAPlate13}}
\newcommand{\@platelabelxiv}{\oldlabel{RMAAPlate14}\label{RMAAPlate14}}
\newcommand{\@platelabelxv}{\oldlabel{RMAAPlate15}\label{RMAAPlate15}}
\newcommand{\@outputplatelabeli}{\oldlabel{RMAAOutputPlate1}}
\newcommand{\@outputplatelabelii}{\oldlabel{RMAAOutputPlate2}}
\newcommand{\@outputplatelabeliii}{\oldlabel{RMAAOutputPlate3}}
\newcommand{\@outputplatelabeliv}{\oldlabel{RMAAOutputPlate4}}
\newcommand{\@outputplatelabelv}{\oldlabel{RMAAOutputPlate5}}
\newcommand{\@outputplatelabelvi}{\oldlabel{RMAAOutputPlate6}}
\newcommand{\@outputplatelabelvii}{\oldlabel{RMAAOutputPlate7}}
\newcommand{\@outputplatelabelviii}{\oldlabel{RMAAOutputPlate8}}
\newcommand{\@outputplatelabelix}{\oldlabel{RMAAOutputPlate9}}
\newcommand{\@outputplatelabelx}{\oldlabel{RMAAOutputPlate10}}
\newcommand{\@outputplatelabelxi}{\oldlabel{RMAAOutputPlate11}}
\newcommand{\@outputplatelabelxii}{\oldlabel{RMAAOutputPlate12}}
\newcommand{\@outputplatelabelxiii}{\oldlabel{RMAAOutputPlate13}}
\newcommand{\@outputplatelabelxiv}{\oldlabel{RMAAOutputPlate14}}
\newcommand{\@outputplatelabelxv}{\oldlabel{RMAAOutputPlate15}}
\newcommand{\@platerefi}{\oldpageref{RMAAPlate1}}
\newcommand{\@platerefii}{\oldpageref{RMAAPlate2}}
\newcommand{\@platerefiii}{\oldpageref{RMAAPlate3}}
\newcommand{\@platerefiv}{\oldpageref{RMAAPlate4}}
\newcommand{\@platerefv}{\oldpageref{RMAAPlate5}}
\newcommand{\@platerefvi}{\oldpageref{RMAAPlate6}}
\newcommand{\@platerefvii}{\oldpageref{RMAAPlate7}}
\newcommand{\@platerefviii}{\oldpageref{RMAAPlate8}}
\newcommand{\@platerefix}{\oldpageref{RMAAPlate9}}
\newcommand{\@platerefx}{\oldpageref{RMAAPlate10}}
\newcommand{\@platerefxi}{\oldpageref{RMAAPlate11}}
\newcommand{\@platerefxii}{\oldpageref{RMAAPlate12}}
\newcommand{\@platerefxiii}{\oldpageref{RMAAPlate13}}
\newcommand{\@platerefxiv}{\oldpageref{RMAAPlate14}}
\newcommand{\@platerefxv}{\oldpageref{RMAAPlate15}}
\newcounter{figplatei}
\newcounter{figplateii}
\newcounter{figplateiii}
\newcounter{figplateiv}
\newcounter{figplatev}
\newcounter{figplatevi}
\newcounter{figplatevii}
\newcounter{figplateviii}
\newcounter{figplateix}
\newcounter{figplatex}
\newcounter{figplatexi}
\newcounter{figplatexii}
\newcounter{figplatexiii}
\newcounter{figplatexiv}
\newcounter{figplatexv}
%% Printed output should use the counter `plate'
\renewcommand{\theplate}{\@arabic\c@plate}
\renewcommand{\theplateloop}{\@arabic\c@plateloop}
%% Command to output all the saved plates in this paper
\newcommand{\outputplates}{%
  \setcounter{plateloop}{0}
  \if@book\else\@endofarticlehook\fi % not entirely sure why, but this
                                     % seems to be needed
  \clearpage % first output the last page of text
  % now set pagestyle for blank even pages and make sure we start on an odd page
  \pagestyle{empty}
  \cleardoublepage
  \onecolumn
  \parindent\z@
  % first put in an extra sheet
  % the main purpose of this is so there is always a blank even page before 
  % each plate for filling in the PS file if the article has an odd number
  % of pages (see splitbook.sh)
  \null\vfil\centerline{PLATES}\vfil\clearpage\pagestyle{rmaaplate}\cleardoublepage
  % output each plate in turn
  \@whilenum\value{plateloop}<\value{plate}\do{%
    \addtocounter{plateloop}{1}
    % restore the figure counter from \@figcounti, \@figcountii, etc
    \setcounter{figure}{\value{figplate\roman{plateloop}}}
    % spew out the plate (in its own group for safety)
    {\csname @plate\roman{plateloop}\endcsname\par}
    }
  \pagestyle{rmaa}
  % reset the page number as if nothing had happened
  \addtocounter{page}{-\value{plate}}
  \addtocounter{page}{-\value{plate}}
  }
\newcommand\figcaption{%
  \def\@captype{figure}
  \renewcommand{\@makecaption}{\@makefigcap}
  \caption}
\newcommand\tabcaption{%
  \def\@captype{table}
  \renewcommand{\@makecaption}{\@maketabcap}
  \caption}

\newcommand{\plate}[2][1]{%
  %% User command to make a plate page: 
  %% Mandatory argument is what goes on the plate - this should not 
  %% include any floating environments (e.g. figure). Instead, \figcaption 
  %% should be used to do the captions. 
  %% Optional argument is the number of figures (i.e. captions) on the plate
  %% page (default is 1). This is necessary for incrementing the figure counter. 
  %%
  %% Note that we have to use a \gdef rather than an \xdef when we 
  %% save the plate output in a command. This means that in the book styles when it 
  %% comes to spewing out the plate, we will have forgotten about any user macros 
  %% that may have been defined. This can be got round by repeating the macros inside the plate.
  %% (Presumably, one would use a \newcommand rather than \renewcommand if they were new.)
  % increment global and local counters
  \addtocounter{plate}{1}
  % write appropriate label to .aux file
  \csname @platelabel\roman{plate}\endcsname
  % define a global macro called \@platei, \@plateii, \@plateiii, \@plateiv, ...
  \expandafter\gdef\csname @plate\roman{plate}\endcsname{%
    \csname @outputplatelabel\roman{plateloop}\endcsname
    \null\vfil #2 \vfil
    \cleardoublepage
    }
  % save the current shortauthor and shorttitle
  % need to be \xdef to take current values of \@shortauthor, etc 
  \expandafter\xdef\csname @shorta\roman{plate}\endcsname{\@shortauthor}
  \expandafter\xdef\csname @shortt\roman{plate}\endcsname{\@shorttitle}
  % save the current value of the figure counter for later use in \outputplates
  \setcounter{figplate\roman{plate}}{\value{figure}}
  % increment the figure counter according to the optional argument
  \addtocounter{figure}{#1}
  }
%% page style for plate pages 
\newcommand{\ps@rmaaplate}{%
  % Plate X (see page Y) in the foot
  \renewcommand\@oddfoot{\hfil\csname @shorta\roman{plateloop}\endcsname\ 
    (see page~\csname @plateref\roman{plateloop}\endcsname)\hfil}%
  \renewcommand\@evenfoot{\relax}%
  % nothing in even heads (blank pages)
  \renewcommand\@evenhead{\relax}
  % short title and page num in the odd heads
  \renewcommand\@oddhead{\hfil\csname @shortt\roman{plateloop}\endcsname%
    \hfil\llap{PLATE~\theplateloop}}
  }





%% Things for the User Guide
\if@guide
%% stolen from doc.sty
\def\MakeShortVerb#1{%
  \expandafter\ifx\csname cc\string#1\endcsname\relax
  \@shortvrbinfo{Made }{#1}%
  \add@special{#1}%
  \expandafter
  \xdef\csname cc\string#1\endcsname{\the\catcode`#1}%
  \begingroup
  \catcode`\~\active  \lccode`\~`#1%
  \lowercase{%
    \global\expandafter\let
    \csname ac\string#1\endcsname~%
    \gdef~{\verb~}}%
  \endgroup
  \global\catcode`#1\active
  \else
  \@shortvrbinfo\@empty{#1 already}%
  \fi}
\def\DeleteShortVerb#1{%
  \expandafter\ifx\csname cc\string#1\endcsname\relax
  \else
  \@shortvrbinfo{Deleted }{#1 as}%
  \rem@special{#1}%
  \global\catcode`#1\csname cc\string#1\endcsname
  \global \expandafter\let \csname cc\string#1\endcsname \relax
  \ifnum\catcode`#1=\active
  \begingroup
  \catcode`\~\active   \lccode`\~`#1%
  \lowercase{%
    \global\expandafter\let\expandafter~%
    \csname ac\string#1\endcsname}%
  \endgroup \fi \fi}
\def\@shortvrbinfo#1#2{%
  \PackageInfo{doc}{%
    #1\expandafter\@gobble\string#2 a short reference
    for \string\verb}}
\def\add@special#1{%
  \rem@special{#1}%
  \expandafter\gdef\expandafter\dospecials\expandafter
  {\dospecials \do #1}%
  \expandafter\gdef\expandafter\@sanitize\expandafter
  {\@sanitize \@makeother #1}}
\def\rem@special#1{%
  \def\do##1{%
    \ifnum`#1=`##1 \else \noexpand\do\noexpand##1\fi}%
  \xdef\dospecials{\dospecials}%
  \begingroup
  \def\@makeother##1{%
    \ifnum`#1=`##1 \else \noexpand\@makeother\noexpand##1\fi}%
  \xdef\@sanitize{\@sanitize}%
  \endgroup}
%% define | as shortcut for \verbatim
\MakeShortVerb{\|}
%% comand to specify arguments to LaTeX commands
\newcommand{\rmArg}[1]{%
  \ensuremath{\langle}{\scriptsize \textsf{\MakeUpperCase{#1}}}\ensuremath{\rangle}}
%\newcommand{\rmArg}[1]{\ensuremath{\langle}\textsc{#1}\ensuremath{\rangle}}
\providecommand{\BibTeX}{{\rmfamily B\kern-.05em%
    \textsc{i\kern-.025em b}\kern-.08em%
    T\kern-.1667em\lower.7ex\hbox{E}\kern-.125emX}}
%% Environment for code examples
\newenvironment{Code}{%
  \begin{list}{}{\setlength{\leftmargin}{0.05\columnwidth}} 
  \item[]}{\end{list}}
\fi
%</package>
%    \end{macrocode}

% \section{User commands for astronomical articles}
% \label{sec:user-commands}
% Some random definitions. These go both into the main class file and
% also into |adshtml.sty|.
% 
% Definitions of the journal name in different incarnations.  
%    \begin{macrocode}
%<*package|adshtml-style>
\newcommand{\RMAA}{RMxAA}
\newcommand{\RMAAlong}{Revista Mexicana de Astronom\'{\i}a y Astrof\'{\i}sica}
\newcommand{\RMAASC}{RMxAC}
\newcommand{\RMAASClong}{\RMAA\ (Serie de Conferencias)}
%    \end{macrocode}
% \begin{macro}{\ibidrule}
% This puts a horizontal rule with a period after it. It is used when
% the author list of a reference is identical to the preceding one. 
%    \begin{macrocode}
\newcommand{\ibidrule}{\rule[-0.03cm]{4.7em}{.01cm}.\ }
%    \end{macrocode}
% \end{macro}
% \begin{macro}{\ThisVolume}
% References to other articles in the same volume can be coded as,
% e.g., |\bibitem{} Bloggs, J. \ThisVolume{bloggs}|.
%    \begin{macrocode}
\newcommand{\ThisVolume}[1]{%
  \@year, \if@proceedings\RMAASC\else\RMAA\fi, \@volume,
  \if@book\oldpageref{#1:RMAAFirstPage}\else 000\fi\ (this volume)}
%    \end{macrocode}
% \end{macro}
% Other stuff that may be useful in ref lists.
%    \begin{macrocode}
\newcommand{\etal}{et al.\@}
%    \end{macrocode}
% Various new commands from AAS macros
%    \begin{macrocode}
\def\acknowledgments{\vskip 3ex plus .8ex minus .4ex}
\let\acknowledgements=\acknowledgments                  % second spelling
\let\jnl@style=\rm
\def\ref@jnl#1{{\jnl@style#1}}
\DeclareRobustCommand\aj{\ref@jnl{AJ}}                   
\DeclareRobustCommand\araa{\ref@jnl{ARA\&A}}             
\DeclareRobustCommand\apj{\ref@jnl{ApJ}}                 
\DeclareRobustCommand\apjl{\ref@jnl{ApJ}}                
\DeclareRobustCommand\apjs{\ref@jnl{ApJS}}               
\DeclareRobustCommand\ao{\ref@jnl{Appl.~Opt.}}           
\DeclareRobustCommand\apss{\ref@jnl{Ap\&SS}}             
\DeclareRobustCommand\aap{\ref@jnl{A\&A}}                
\DeclareRobustCommand\aapr{\ref@jnl{A\&A~Rev.}}          
\DeclareRobustCommand\aaps{\ref@jnl{A\&AS}}              
\DeclareRobustCommand\azh{\ref@jnl{AZh}}                 
\DeclareRobustCommand\baas{\ref@jnl{BAAS}}               
\DeclareRobustCommand\jrasc{\ref@jnl{JRASC}}             
\DeclareRobustCommand\memras{\ref@jnl{MmRAS}}            
\DeclareRobustCommand\mnras{\ref@jnl{MNRAS}}             
\DeclareRobustCommand\pra{\ref@jnl{Phys.~Rev.~A}}        
\DeclareRobustCommand\prb{\ref@jnl{Phys.~Rev.~B}}        
\DeclareRobustCommand\prc{\ref@jnl{Phys.~Rev.~C}}        
\DeclareRobustCommand\prd{\ref@jnl{Phys.~Rev.~D}}        
\DeclareRobustCommand\pre{\ref@jnl{Phys.~Rev.~E}}        
\DeclareRobustCommand\prl{\ref@jnl{Phys.~Rev.~Lett.}}    
\DeclareRobustCommand\pasp{\ref@jnl{PASP}}               
\DeclareRobustCommand\pasj{\ref@jnl{PASJ}}               
\DeclareRobustCommand\qjras{\ref@jnl{QJRAS}}             
\DeclareRobustCommand\skytel{\ref@jnl{S\&T}}             
\DeclareRobustCommand\solphys{\ref@jnl{Sol.~Phys.}}      
\DeclareRobustCommand\sovast{\ref@jnl{Soviet~Ast.}}      
\DeclareRobustCommand\ssr{\ref@jnl{Space~Sci.~Rev.}}     
\DeclareRobustCommand\zap{\ref@jnl{ZAp}}                 
\DeclareRobustCommand\nat{\ref@jnl{Nature}}              
\DeclareRobustCommand\iaucirc{\ref@jnl{IAU~Circ.}}
\DeclareRobustCommand\aplett{\ref@jnl{Astrophys.~Lett.}}
\DeclareRobustCommand\apspr{\ref@jnl{Astrophys.~Space~Phys.~Res.}}
\DeclareRobustCommand\bain{\ref@jnl{Bull.~Astron.~Inst.~Netherlands}}
\DeclareRobustCommand\fcp{\ref@jnl{Fund.~Cosmic~Phys.}}
\DeclareRobustCommand\gca{\ref@jnl{Geochim.~Cosmochim.~Acta}}
\DeclareRobustCommand\grl{\ref@jnl{Geophys.~Res.~Lett.}}
\DeclareRobustCommand\jcp{\ref@jnl{J.~Chem.~Phys.}}      
\DeclareRobustCommand\jgr{\ref@jnl{J.~Geophys.~Res.}}    
\DeclareRobustCommand\jqsrt{\ref@jnl{J.~Quant.~Spec.~Radiat.~Transf.}}
\DeclareRobustCommand\memsai{\ref@jnl{Mem.~Soc.~Astron.~Italiana}}
\DeclareRobustCommand\nphysa{\ref@jnl{Nucl.~Phys.~A}}
\DeclareRobustCommand\physrep{\ref@jnl{Phys.~Rep.}}
\DeclareRobustCommand\physscr{\ref@jnl{Phys.~Scr}}
\DeclareRobustCommand\planss{\ref@jnl{Planet.~Space~Sci.}}
\DeclareRobustCommand\procspie{\ref@jnl{Proc.~SPIE}}
\let\astap=\aap
\let\apjlett=\apjl
\let\apjsupp=\apjs
\let\applopt=\ao
\def\phn{\phantom{0}}
\def\phd{\phantom{.}}
\def\phs{\phantom{\ensuremath{-}}}
\def\phm#1{\phantom{#1}}
%    \end{macrocode}
% All these are redone as robust and using more modern commands where
% possible. The robustness is so they will not get expanded if used in
% abstracts and titles when these are passed to the ADS file. 
% \changes{3.16}{2003/04/07}{
%  Added \cs{Color} command for doing photometric colors.}
%    \begin{macrocode}
%%\def\sun{\hbox{$\odot$}}
\DeclareRobustCommand\sun{\ensuremath{\odot}}
%%\def\earth{\hbox{$\oplus$}}
\DeclareRobustCommand\earth{\ensuremath{\oplus}}
%%\def\lesssim{\mathrel{\hbox{\rlap{\hbox{\lower4pt\hbox{$\sim$}}}\hbox{$<$}}}}
%%\def\gtrsim{\mathrel{\hbox{\rlap{\hbox{\lower4pt\hbox{$\sim$}}}\hbox{$>$}}}}
\DeclareRobustCommand\lesssim{\mathrel{\hbox{\smaller\smaller\rlap{\hbox{\lower3pt\hbox{$\sim$}}}\raise1.5pt\hbox{$<$}}}}
\DeclareRobustCommand\gtrsim{\mathrel{\hbox{\smaller\smaller\rlap{\hbox{\lower3pt\hbox{$\sim$}}}\raise1.5pt\hbox{$>$}}}}
%%\def\lta{\mathrel{\spose{\lower 3pt\hbox{$\mathchar"218$}}
%%     \raise 2.0pt\hbox{$\mathchar"13C$}}}
%%\def\gta{\mathrel{\spose{\lower 3pt\hbox{$\mathchar"218$}}
%%     \raise 2.0pt\hbox{$\mathchar"13E$}}}
\DeclareRobustCommand\sq{\hbox{\rlap{$\sqcap$}$\sqcup$}}
%%\def\arcdeg{\hbox{$^\circ$}}
%%\def\arcmin{\hbox{$^\prime$}}
%%\def\arcsec{\hbox{$^{\prime\prime}$}}
\DeclareRobustCommand\arcdeg{\ensuremath{^\circ}}
\DeclareRobustCommand\arcmin{\ensuremath{^\prime}}
\DeclareRobustCommand\arcsec{\ensuremath{^{\prime\prime}}}
%%\def\fd{\hbox{$.\!\!^{\rm d}$}}
%%\def\fh{\hbox{$.\!\!^{\rm h}$}}
%%\def\fm{\hbox{$.\!\!^{\rm m}$}}
%%\def\fs{\hbox{$.\!\!^{\rm s}$}}
%%\def\fdg{\hbox{$.\!\!^\circ$}}
%%\def\farcm{\hbox{$.\mkern-4mu^\prime$}}
%%\def\farcs{\hbox{$.\!\!^{\prime\prime}$}}
%%\def\fp{\hbox{$.\!\!^{\scriptscriptstyle\rm p}$}}
\DeclareRobustCommand\fd{\ensuremath{.\!\!^\mathrm{d}}}
\DeclareRobustCommand\fh{\ensuremath{.\!\!^\mathrm{h}}}
\DeclareRobustCommand\fm{\ensuremath{.\!\!^\mathrm{m}}}
\DeclareRobustCommand\fs{\ensuremath{\mkern0.5mu.\mkern-5mu^\mathrm{s}\mkern-1mu}}
\DeclareRobustCommand\fdg{\ensuremath{.\!\!^\circ}}
\DeclareRobustCommand\farcm{\ensuremath{.\mkern-4mu^\prime}}
\DeclareRobustCommand\farcs{\ensuremath{\mkern1mu.\mkern-6mu^{\prime\prime}\mkern-1.5mu}}
%%\DeclareRobustCommand\farcs{\ensuremath{.\!\!^{\prime\prime}}}
\DeclareRobustCommand\fp{\ensuremath{.\!\!^{\scriptscriptstyle\rm p}}}
%%\def\micron{\hbox{$\mu$m}}
\DeclareRobustCommand\micron{\ensuremath{\mu\mathrm{m}}}
\let\la=\lesssim                        
\let\ga=\gtrsim
%%\def\case#1#2{\hbox{$\frac{#1}{#2}$}}
%%\def\slantfrac#1#2{\hbox{$\,^#1\!/_#2$}}
%%\def\onehalf{\slantfrac{1}{2}}
%%\def\onethird{\slantfrac{1}{3}}
%%\def\twothirds{\slantfrac{2}{3}}
%%\def\onequarter{\slantfrac{1}{4}}
%%\def\threequarters{\slantfrac{3}{4}}
\DeclareRobustCommand\case[2]{\ensuremath{\frac{#1}{#2}}}
\DeclareRobustCommand\slantfrac[2]{\ensuremath{\,^#1\!/_#2}}
\DeclareRobustCommand\onehalf{\slantfrac{1}{2}}
\DeclareRobustCommand\onethird{\slantfrac{1}{3}}
\DeclareRobustCommand\twothirds{\slantfrac{2}{3}}
\DeclareRobustCommand\onequarter{\slantfrac{1}{4}}
\DeclareRobustCommand\threequarters{\slantfrac{3}{4}}
%%\def\ubvr{\hbox{$U\!BV\!R$}}            
%%\def\ub{\hbox{$U\!-\!B$}}               
%%\def\bv{\hbox{$B\!-\!V$}}               
%%\def\vr{\hbox{$V\!-\!R$}}               
%%\def\ur{\hbox{$U\!-\!R$}}
%    \end{macrocode}
% New commands for ``photometric'' colors 
%    \begin{macrocode}
\DeclareRobustCommand\ubvr{\ensuremath{U\!BV\!R}}            
\DeclareRobustCommand\Color[2]{%
  \ensuremath{%
    #1\!\ifx#1V\!\fi-\!#2
  }}
\DeclareRobustCommand\ub{\Color{U}{B}}
\DeclareRobustCommand\bv{\Color{B}{V}}
\DeclareRobustCommand\vr{\Color{V}{R}}
\DeclareRobustCommand\ur{\Color{U}{R}}
%% New version of the \ion command, that actually works! 
%    \end{macrocode}
% \begin{macro}{\ion}
% This version of the |\ion| command allows you to write either
% |\ion{N}{4}| or |\ion{N}{iv}| in order to get N~IV. 
%    \begin{macrocode}
\DeclareRobustCommand\ion[2]{%
  \mbox{#1\kern0.2em%
    \smaller\rmfamily%
    \edef\@tempa{\@car#2\@nil}%
    \ifcat1\@tempa%
    \@Roman{#2}%
    \else%
    \uppercase{#2}%
    \fi}}
\newcounter{IonStage}
\DeclareRobustCommand\plainion[2]{\setcounter{IonStage}{#2}#1 \Roman{IonStage}}
%%\def\ion#1#2{\ensuremath{\mathrm{#1}\;}{\protect\small\rm\@Roman{#2}}}
%    \end{macrocode}
% \end{macro}
%    \begin{macrocode}
\DeclareRobustCommand\nodata{$\cdots$}
%    \end{macrocode}
% These are needed by the |\ion| command. Isn't there an easier way of
% doing this?
%    \begin{macrocode}
\let\@small=\small
\let\@normalsize=\normalsize
\let\@footnotesize=\footnotesize
\let\@scriptsize=\scriptsize
\let\@tiny=\tiny
\let\@large=\large
\let\@Large=\Large
\let\@LARGE=\LARGE
\let\@huge=\huge
\let\@Huge=\Huge
\newcommand{\smaller}{\small}
\renewcommand{\normalsize}{\renewcommand{\smaller}{\small}\@normalsize}
\renewcommand{\small}{\renewcommand{\smaller}{\footnotesize}\@small}
\renewcommand{\footnotesize}{\renewcommand{\smaller}{\scriptsize}\@footnotesize}
\renewcommand{\scriptsize}{\renewcommand{\smaller}{\tiny}\@scriptsize}
\renewcommand{\tiny}{\renewcommand{\smaller}{\tiny}\@tiny}
\renewcommand{\large}{\renewcommand{\smaller}{\normalsize}\@large}  
\renewcommand{\Large}{\renewcommand{\smaller}{\large}\@Large}  
\renewcommand{\LARGE}{\renewcommand{\smaller}{\Large}\@LARGE}  
\renewcommand{\huge}{\renewcommand{\smaller}{\LARGE}\@huge}  
\renewcommand{\Huge}{\renewcommand{\smaller}{\huge}\@Huge}  
%</package|adshtml-style>
%    \end{macrocode}


% \section{Other files}

% \subsection{The driver and make files}
% \label{sec:driver_make}
% These are put here for completeness, although they still need to be
% distributed separately since otherwise the user won't be able to
% unpack this file into its constituent parts.
% 
% The |Makefile| has two principal rules: |docs|, and |macros|. The
% |docs| rule produces a PDF and gzipped PS file of the documentation.
% I still haven't perfected the combination of rules required to get
% it to run \LaTeX{} enough times in all circumstances. The |macros|
% rule runs \LaTeX{} on |rmaa.ins| in order to unpack |rmaa.cls| and
% all the other little files (including the |Makefile| and the
% |rmaa.ins| files themselves\dots 
% 
% First set up lists of various types of file. 
%    \begin{macrocode}
%<*makefile>
DOCFILES=rmaa.pdf
MACRO_LTX_FILES=rmaa.cls rmaacite.sty adshtml.sty
SHELLSCRIPTS=makeads.sh
MACRO_OTHER_FILES=adshtml.perl sample.latex2html-init
MAKEFILE=Makefile
MACROFILES=$(MACRO_LTX_FILES) $(SHELLSCRIPTS) $(MACRO_OTHER_FILES)         
SKELETONFILES=ADS.tex 
%    \end{macrocode}
% Install the files read only into the package directory, specified via a path relative to the source directory. \changes{4.0}{2019/11/14}{Reorganize folder layout so that hardly anything is in the top level.  That way, things are clearer on github.}
%    \begin{macrocode}
INSTALL=install 
INSTALL_DIR=../packages
%    \end{macrocode}
% The default rule is to make everything. Installation must be
% explicitly asked for.
%    \begin{macrocode}
all: docs macros

install: install-docs install-macros

%    \end{macrocode}
% Most files are installed read-only to prevent editing. Skeleton
% files must be writeable. Must be careful not to overwrite the
% Makefile in the parent directory. 
%    \begin{macrocode}
install-macros: macros
        $(INSTALL) -m 444 $(MACRO_LTX_FILES) $(MACRO_OTHER_FILES) $(INSTALL_DIR)
        $(INSTALL) -m 555 $(SHELLSCRIPTS) $(INSTALL_DIR)
        $(INSTALL) -m 644 $(SKELETONFILES) $(INSTALL_DIR)

install-docs: docs
        $(INSTALL) rmaa.pdf $(INSTALL_DIR)/rm-fulldocs.pdf

%    \end{macrocode}
% Make the PDF and gzipped PS files. 
%    \begin{macrocode}
docs: $(DOCFILES)

%    \end{macrocode}
% Make everything else if either the |.dtx| or |.ins| files have
% changed. 
%    \begin{macrocode}
macros: rmaa.dtx rmaa.ins
%    \end{macrocode}
% Ensure that all the macro files can be written to then recreate
% them. The |touch| is there in case any of the files don't exist
% yet. 
%    \begin{macrocode}
        touch $(MACROFILES) $(MAKEFILE)
        chmod u+w $(MACROFILES) $(MAKEFILE) rmaa.ins
        latex "\input{rmaa.ins}"
%    \end{macrocode}
% \TeX{} automatically adds the |.tex| extension to extensionless
% files. We fix this and take the opportunity to convert
% back to TABs all that it has lovingly converted to spaces. 
%    \begin{macrocode}
        cp -f $(MAKEFILE) $(MAKEFILE).BCK && unexpand $(MAKEFILE).tex >| $(MAKEFILE) && rm $(MAKEFILE).tex
%    \end{macrocode}
% Make the shell scripts executable and everything unwriteable (so
% they don't get hand-edited by mistake). However, undo this for the
% files that the user is \emph{supposed} to modify. 
%    \begin{macrocode}
        chmod a+x $(SHELLSCRIPTS)
        chmod a-w $(MACROFILES) $(MAKEFILE)
        chmod u+w $(SKELETONFILES)

%    \end{macrocode}
% We now use PDF\LaTeX{} for this.
%    \begin{macrocode}
rmaa.idx: rmaa.dtx
        @echo "REGENERATING .idx FILE"
        pdflatex "\nonstopmode\input{rmaa.dtx}"

rmaa.aux: rmaa.dtx rmaa.ind rmaa.gls
        @echo "REGENERATING .aux FILE"
        pdflatex "\nonstopmode\input{rmaa.dtx}"

rmaa.ind: rmaa.idx
        @echo "REGENERATING INDEX"
        makeindex -s gind.ist rmaa

rmaa.gls: rmaa.glo
        @echo "REGENERATING CHANGELOG"
        makeindex -s gglo.ist -o rmaa.gls rmaa.glo

rmaa.pdf: rmaa.aux rmaa.ind rmaa.gls
        @echo "REGENERATING .pdf FILE"
        pdflatex "\nonstopmode\input{rmaa.dtx}"

clean: 
        rm -f rmaa.aux rmaa.dvi rmaa.ind rmaa.gls
%</makefile>
%    \end{macrocode}

% The |rmaa.ins| file uses the \textsf{docstrip} package to split out
% all the files and strip off this documentation.
%    \begin{macrocode}
%<*ins>
%%
%% Installation batch file for the RMAA LaTeX package
%%
\def\batchfile{rmaa.ins}
\input docstrip.tex
\preamble

PLEASE DO NOT EDIT THIS FILE! It is generated automatically from
the file rmaa.dtx using the accompanying Makefile and rmaa.ins 
Any hand-edited changes will be lost

A LaTeX document class for the RMxAA

Author: William Henney (w.henney@irya.unam.mx)

Latest version available from: http://www.irya.unam.mx/rmaa

Copyright (C) 1999-2019 Instituto de Astronomia, UNAM, Mexico

Permission is hereby granted, free of charge, to any person obtaining
a copy of this software and associated documentation files (the
"Software"), to deal in the Software without restriction, including
without limitation the rights to use, copy, modify, merge, publish,
distribute, sublicense, and/or sell copies of the Software, and to
permit persons to whom the Software is furnished to do so, subject to
the following conditions:

The above copyright notice and this permission notice shall be
included in all copies or substantial portions of the Software.

THE SOFTWARE IS PROVIDED "AS IS", WITHOUT WARRANTY OF ANY KIND,
EXPRESS OR IMPLIED, INCLUDING BUT NOT LIMITED TO THE WARRANTIES OF
MERCHANTABILITY, FITNESS FOR A PARTICULAR PURPOSE AND
NONINFRINGEMENT. IN NO EVENT SHALL THE AUTHORS OR COPYRIGHT HOLDERS BE
LIABLE FOR ANY CLAIM, DAMAGES OR OTHER LIABILITY, WHETHER IN AN ACTION
OF CONTRACT, TORT OR OTHERWISE, ARISING FROM, OUT OF OR IN CONNECTION
WITH THE SOFTWARE OR THE USE OR OTHER DEALINGS IN THE SOFTWARE.
\endpreamble
%    \end{macrocode}
% Clobber files without asking
%    \begin{macrocode}
\askforoverwritefalse
%    \end{macrocode}
% Make the \LaTeX{} macro files (class and styles)
%    \begin{macrocode}
\generate{% 
  \file{rmaa.cls}{\from{rmaa.dtx}{package}}
  \file{rmaa.ins}{\from{rmaa.dtx}{ins}}
  \file{rmaacite.sty}{\from{rmaa.dtx}{rmaacite}}
  \file{adshtml.sty}{\from{rmaa.dtx}{adshtml-style}}
  \file{ADS.tex}{\from{rmaa.dtx}{ads-container}}
  }
%    \end{macrocode}
% Make the other non-\LaTeX{} files. These are done separately since
% they can't use `\%\%' as a comment character for the
% preamble/postamble (and Perl files can't have the postamble I don't
% think). 
%    \begin{macrocode}
\generate{\nopreamble\nopostamble% 
  \file{Makefile}{\from{rmaa.dtx}{makefile}}
  \file{adshtml.perl}{\from{rmaa.dtx}{adshtml-perl}}
  \file{makeads.sh}{\from{rmaa.dtx}{ads-shell-script}}
  \file{sample.latex2html-init}{\from{rmaa.dtx}{latex2html-init}}
  }

%% Local Variables: 
%% mode: LaTeX
%% mode: Font-lock
%% End:
%</ins>
%    \end{macrocode}

% \subsection{Extended citation syntax}
% Modified from a file called mnrascite.sty
% I no longer remember where I stole the original version of this from
% 
% Alterations to the cite commands to make them do more sensible things.
% N.B. the bibliography style should produce entries which have labels where
% the date is delimited by angle brackets '<>' for these macros to work
% properly. They will not give very sensible error messages if you do not do 
% this properly.
% 
% A new command |\scite{key}| is defined which produces a citation which is 
% suitable for insertion in a sentence. i.e., it does not put brackets around 
% the whole citation, only the date. Also |\pcite{key}| produce output with no
% parentheses whatsoever. 
% 
% Should redefine |\@bibitem| really to have any hope of compatibility
% the command |\lbibcite| is defined to deal with the citations.
%    \begin{macrocode}
%<*rmaacite>
\def\bibitem[#1<#2>]#3{\par\hangindent=1.5em\hangafter=1
  \if@filesw
  {\def\protect##1{\string ##1\space}
    \immediate\write\@auxout{\string\lbibcite{#3}{#1}{#2}}}
  \fi}
\@redefinebibitem % what it says
%    \end{macrocode}
% Note that the |\lbibcite| defines two types of citation one contains the names,
% the other the year
%    \begin{macrocode}
\def\lbibcite#1#2#3{\global\@namedef{b@#1}{#2}
                    \global\@namedef{by@#1}{#3}}

\def\@bcite#1#2{({#1\if@tempswa , #2\fi})}
\def\@pcite#1#2{#1\if@tempswa , #2\fi}
%    \end{macrocode}
% This declaration contains what is put between the name and the date
% \changes{3.11}{2003/01/31}{Moved the \cs{articlename} to after the key for consistency
% with the new natbib stuff.}
%    \begin{macrocode}
\def\@brcite{\ (}
\def\@ebrcite{)}
\def\@fmttop#1{\@fmtcite[\csname #1\endcsname]}
\def\@citex[#1]#2{%
%%  \if@filesw\immediate\write\@auxout{\string\citation{#2}}\fi
%% WJH added 16 Jan 2002
  \@for\@citeb:=#2\do{%
    \edef\@citeb{\expandafter\@firstofone\@citeb:\articlename}%
    \if@filesw\immediate\write\@auxout{\string\citation{\@citeb}}\fi}%
%% to here
  \def\@citea{}%
  \@cite{\@for\@citeb:=#2\do
    {%
      \edef\@citeb{\expandafter\@firstofone\@citeb:\articlename}%
      \@citea\def\@citea{;\ }\@ifundefined
      {b@\@citeb}{{\bf ?}\@warning {Citation `\@citeb' on page \thepage \space undefined}}%
      {\csname b@\@citeb\endcsname\@citebro\csname by@\@citeb\endcsname\@citebrc}}}%
  {#1}}

\def\cite{\@ifnextchar [{\let\@citebro=\space
                          \let\@citebrc=\relax
                          \let\@cite=\@bcite\@tempswatrue\@citex}
                        {\let\@citebro=\space
                          \let\@citebrc=\relax
                          \let\@cite=\@bcite\@tempswafalse\@citex[]}}
\def\pcite{\@ifnextchar [{\let\@citebro=\space
                          \let\@citebrc=\relax
                          \let\@cite=\@pcite\@tempswatrue\@citex}
                        {\let\@citebro=\space
                          \let\@citebrc=\relax
                          \let\@cite=\@pcite\@tempswafalse\@citex[]}}
\def\scite{\@ifnextchar [{\let\@citebro=\@brcite
                          \let\@citebrc=\@ebrcite
                          \let\@cite=\@pcite\@tempswatrue\@citex}
                        {\let\@citebro=\@brcite
                          \let\@citebrc=\@ebrcite
                          \let\@cite=\@pcite\@tempswafalse\@citex[]}}

\renewenvironment{thebibliography}
  {\section*{REFERENCES}\vspace*{\afterbiblioskip}\small
    \if@book\RereadBiblio\fi
    \sloppy\clubpenalty4000\widowpenalty4000
    \sfcode`\.=1000\relax
    \setlength\parindent{\z@}\setlength\parskip{\z@}}
  {\relax\normalsize\@outputaddressesifnecessary} 
%</rmaacite>
%    \end{macrocode}

% \subsection{Auxilliary files for generating ADS data}

% The ADS data are written by \LaTeX{} to a |.ads| file. This file
% then has to be included in the following file (|ADS.tex|),
% which is run through \LaTeX{}2HTML (see below). 

%    \begin{macrocode}
%<*ads-container>
%%
%% Instructions:
%% Edit the \input{} line below to reflect the name of your file. 
\documentclass{article}
\usepackage[latin1]{inputenc}
\usepackage{adshtml}
\newcommand{\ADSfield}[2]
{\%#1  #2\\}
%%\newcommand{\ADSvfield}[2]
%%{\%#1 \begin{verbatim}#2\end{verbatim}}
\newenvironment{ADSarticle}{}{\par\bigskip}
\renewcommand{\bibitem}[2][]{\par}
\setlength{\parindent}{0pt}
\renewcommand{\ThisVolume}[1]{%  
  2001, \RMAASC, 12, \oldpageref{#1:RMAAFirstPage}}  
 

\begin{document}
\makeatletter
\input{rmbooksamp.ads}
\end{document}
%</ads-container>
%    \end{macrocode}

% This is the |adshtml.sty| style file used by |ADS.tex|. It just
% defines a few 
% macros that might be used in abstracts, etc (really it should define
% them all). This now gets all its code from
% Sec~\ref{sec:user-commands}. 
%    \begin{macrocode}
%<*adshtml-style>
%% Nothing else needs to go here now. 
%</adshtml-style>
%    \end{macrocode}
% The init file for latex2html.
%    \begin{macrocode}
%<*latex2html-init>
$PDFTOCAIRO = "pdftocairo";
$PDFLATEX = "pdflatex";
$LATEX = "latex";
$MAX_SPLIT_DEPTH = 0;
$HTML_VERSION = "4.0,math";
$NO_SIMPLE_MATH = 1;
$ADDRESS = "<p><address><a href=\"mailto:rmaa\@astro.unam.mx\">rmaa\@astro.unam.mx</address></p>"; 
$INFO = "<p>This document was created using latex2html and the rmaa latex macros</p>"; 
$NO_NAVIGATION = 1;
1; 
%</latex2html-init>
%    \end{macrocode}
% This is the |adshtml.perl| file that is used by \LaTeX{}2HTML. Some
% bizarre \TeX{}y command sometimes show up so we tell it to ignore
% these. All the astronomy macros that come from AAS are processed in
% \TeX{}. That way they get made to images but with an |ALT| tag that
% has the original code, which is what will be picked up by |lynx|,
% with the result that they get passed intact into the final |.txt|
% file. 
%    \begin{macrocode}
%<*adshtml-perl>
package main;

$dum = &do_cmd_makeatletter();

sub do_cmd_RMAA {
    local($_) = @_;
    "RMxAA" . $_;
}

sub do_cmd_RMAASC {
    local($_) = @_;
    "RMxAASC" . $_;
}

sub do_cmd_etal {
    local($_) = @_;
    "et al\." . $_;
}

sub do_cmd_ibidrule {
    local($_) = @_;
    "<img src=\"foo\" alt=\"\ibidrule\">" . $_;
}


sub do_cmd_bibitem {
    local($_) = @_;
    local($key);
    local($opt, $dummy)=&get_next_optional_argument; 
    $key = &missing_braces unless
        ((s/$next_pair_pr_rx/$key = $2; ''/eo)
        ||(s/$next_pair_rx/$key = $2; ''/eo));
    "<p>" . $_; 
}

&ignore_commands (<<_RAW_ARG_CMDS_);
lowercase
adjustfinalcols
textsuperscript # {}
newline # []
thinmuskip
mskip
smallskip
medskip
bigskip
hspace # {}
hspacestar # {}
vspace # {}
vspacestar # {}
mathrm
_at_mathrm
spacefactor
discretionary
unhbox 
voidb_at_x 
penalty 
_at_M
_at_m
null
IeC
_RAW_ARG_CMDS_

&process_commands_in_tex (<<_RAW_ARG_CMDS_);
sun
earth
lesssim
gtrsim
sq
argdeg
arcmin
arcsec
fd
fh
fm
fs
fdg
farcm
farcs
fp
micron
la
ga
case # {} # {}
slantfrac # {} # {}
onehalf
onethird
twothirds
onequarter
threequarters
ubvr
ub
bv
vr
ur
aj
araa
apj
apjl
apjs
ao
apss
aap
aapr
aaps
azh
baas
jrasc
memras
mnras
pra
prb
prc
prd
pre
prl
pasp
pasj
qjras
skytel
solphys
sovast
ssr
zap
nat
iaucirc
aplett
apspr
bain
fcp
gca
grl
jcp
jgr
jqsrt
memsai
nphysa
physrep
physscr
planss
procspie
_RAW_ARG_CMDS_

1;                              # This must be the last line
%</adshtml-perl> 
%    \end{macrocode}

% This is a shell script |makeads.sh| that generates the |_ads.txt| file from
% |_ads.tex| by using \LaTeX{}2HTML and |lynx|.  \changes{4.0}{2019/11/13}{Make sure that lynx writes utf-8 output instead of latin-1.}
% \changes{2.2beta13}{2002/03/12}{Complete rewrite of |makeads.sh|. Now much 
% improved I think}
% \changes{4.0}{2019/11/13}{Fix the at-sign substitution so that it only eliminates \cs{@} and does not mess with email addresses.}

%    \begin{macrocode}
%<*ads-shell-script>
#! /bin/sh
if [ -z $1 ]; then
    echo "Usage: $0 basename"
    exit
fi
BASENAME=$1
ADSBASENAME=${BASENAME}_ads
AUXFILE=$BASENAME.aux
TMPFILE=tmp.$$


echo "Generating $ADSBASENAME.tex"

# For some reason, must double-up \\ to \\\\ but not \ to \\ - odd!
cat > $ADSBASENAME.tex <<EOF
%% PLEASE DO NOT EDIT THIS FILE!  It is generated automatically from
%% the script makeads.sh.  Any hand-edited changes will be lost
%% 
\documentclass{article}
\usepackage[latin1]{inputenc}
\usepackage{html}
\usepackage{adshtml}
\newcommand{\ADSfield}[2]{\%#1  #2\\\\}
\newenvironment{ADSarticle}{}{\par\bigskip}
\setlength{\parindent}{0pt}
\newcounter{istage}
\renewcommand{\ion}[2]{\setcounter{istage}{#2}#1 \Roman{istage}}
\begin{document}
\makeatletter
EOF
echo >> $ADSBASENAME.tex "\input{$BASENAME.ads2}" 
echo >> $ADSBASENAME.tex "\end{document}"



echo "Generating $ADSBASENAME.sed" 

# Find type of journal
grep '\\documentclass\[.*proceedings.*\]' $BASENAME.tex >/dev/null \
     && JOURNAL='RMxAC' \
     || JOURNAL='RMxAA'

# Find the year and volume number
YEAR=`grep '\\SetYear' $BASENAME.tex | sed 's/.*\\SetYear{\(.*\)}.*/\1/'`
VOLUME=`grep '\\SetVolume' $BASENAME.tex | sed 's/.*\\SetVolume{\(.*\)}.*/\1/'`

# Generate an auxilliary sed script 
cat > $TMPFILE <<EOF
s|\\\\newlabel{\(.*\):RMAAFirstPage}{{.*}{\(.*\)}}|s/\\\\\\\\ThisVolume{\1}/$YEAR, RMxAC, $VOLUME, \2/g|g
EOF

# Now go through the .aux file, looking for \newlabels and use the sed 
# script just generated in order to make another sed script, one that 
# will replace \ThisVolume{LABEL} with the correct reference
grep '^\\newlabel{.*:RMAAFirstPage}' $AUXFILE \
     | sed -f $TMPFILE > $ADSBASENAME.sed && rm $TMPFILE

# Apend a further sed command to get rid of \@, since redefining 
# it in _ads.tex didn't work
echo >> $ADSBASENAME.sed 's/\\\\@//g'

echo "Generating $BASENAME.ads2 from $ADSBASENAME.sed and $BASENAME.ads" 

# Apply this to the .ads file, generating the .ads2 file, which is included
# by the _ads.tex file.
sed -f $ADSBASENAME.sed $BASENAME.ads > $BASENAME.ads2

echo "Generating $ADSBASENAME.txt from $ADSBASENAME.tex and $BASENAME.ads2"

# Now run latex2html on the _ads.tex file
latex2html -no_navigation -info 0 -address 0 -html_version 4.0,math $ADSBASENAME

# Then, convert to utf8 with lynx and tidy it up a bit
lynx --justify=off --width=100 --dump --nolist --display_charset=utf-8 \
     $ADSBASENAME/$ADSBASENAME.html | cut -c4- \
     | sed 's/ibidrule/\\ibidrule/g' > $ADSBASENAME.txt
%</ads-shell-script>
%    \end{macrocode}

% \section{Epilogue}
% Err, that's it \dots

\endinput


% \endinput
% Local Variables:
% mode: doctex
% TeX-master: t
% End:
"

rmaa.ind: rmaa.idx
        @echo "REGENERATING INDEX"
        makeindex -s gind.ist rmaa

rmaa.gls: rmaa.glo
        @echo "REGENERATING CHANGELOG"
        makeindex -s gglo.ist -o rmaa.gls rmaa.glo

rmaa.pdf: rmaa.aux rmaa.ind rmaa.gls
        @echo "REGENERATING .pdf FILE"
        pdflatex "\nonstopmode% \def\fileversion{4.0} \def\filedate{2019/11/12}
% \MakeShortVerb{\|}
% \title{A \LaTeX{} document class for the Revista Mexicana de
% Astronom\'\i a y Astrof\'\i sica\thanks{This file has version
% \fileversion{} last revised \filedate{}.}}
% \author{William J. Henney\thanks{w.henney@irya.unam.mx}}
% \date{\filedate}
% \maketitle
% \paragraph{Abstract}
% This document describes the \texttt{rmaa} \LaTeX{}
% document class for the journal ``Revista Mexicana de
% Astronom\'\i a y Astrof\'\i sica''. The macros can be used for 
% \begin{enumerate}
% \item Preparing author manuscripts for the journal (including
% conference proceedings to be published in the ``Serie de
% Conferencias'')
% \item Typesetting an entire volume of the journal by inputting the
% individual article files from a container file. Facilities are
% included for automatic generation of TOC and Author Index. 
% \end{enumerate}
% 
% \tableofcontents
% \section{The documentation driver file}
% This is the part that generates the documentation that you may or
% may not be reading, which is done with the help of the Makefile
% discussed in section~\ref{sec:driver_make}. 
%    \begin{macrocode}
%<*driver>
\documentclass{ltxdoc}
\usepackage{hyperref}
\usepackage{palatino}
\EnableCrossrefs \CodelineIndex \RecordChanges
\setlength{\parindent}{0pt}
\def\MacroFont{\footnotesize\tt}
\begin{document} \DocInput{rmaa.dtx} \PrintIndex \PrintChanges \end{document}
%</driver>
%    \end{macrocode}

% \section{The implementation}
% 
% \changes{3.24}{2007/09/14}{Added user access to the macro version
% and date via \cs{rmVersion} and \cs{rmDate}}
% 
%    \begin{macrocode}
%<*package>
\newcommand{\rmaa@version}{3.26}
\newcommand{\rmaa@date}{2011/08/15}
\newcommand\rmVersion{\rmaa@version}
\newcommand\rmDate{\rmaa@date}
\NeedsTeXFormat{LaTeX2e}[1995/06/01]
%    \end{macrocode}
% \changes{4.0}{2019/11/02}{Change abstract to appear before
% resumen. Change RevMexAA to RMxAA for consistency with ADS.}
%
% \changes{3.22}{2007/03/22}{Allow proceedings articles to have
%   received/accepted dates} \changes{3.21}{2006/05/15}{Changes for
%   Silvia - abstracts are now shipped out to ADS.}
% \changes{3.19}{2006/01/19}{Changes for Silvia - remove adresses from
%   end, new abstract style} \changes{2.2beta}{2001/10/30}{ Fixed bug
%   whereby twocolumn option was not passed through to book
%   class. Finally sorted out the ``label on last page'' issue. }
% \changes{2.1}{2001/01/23}{ Fixed some inconsistencies in
%   \cs{nocite} and \cs{bibitem} uniqification Fixed the definition of
%   \cs{ion} so it will work in math mode and in contexts where the
%   typesize is different from normal (e.g., title, fig captions).
%   Modified the definition of \cs{@caption}, so we don't write any
%   list-of-figures data (we never use it anyway, and it just causes
%   problems with fragile commands in fig captions). }
% \changes{2.0}{2000/11/30}{ Added `newstyle' option, which should be
%   used in all new proceedings volumes (wider abstract, smaller
%   typeface in fig captions, etc.) Added commands for formatting
%   table captions. New demonstation document rmpostsamp.tex for
%   1-page posters. This will be the last release under a PRI
%   government (for a while...)}  \changes{1.9}{2000/06/06}{
%   Corrected position of page number in footer of articles that open
%   on even pages. Ensured that first article in each session opens on
%   an odd page, as do list of abstracts and author index. Added extra
%   space in TOC before the backmatter. Added command for editors,
%   \cs{BoldSessionTitles}, that makes produces bold titles for
%   session in the TOC.}  \changes{1.8}{2000/05/16}{ Fixed bug in
%   last page label with \cs{suppressfulladdresses} }
% \changes{1.7}{2000/05/11}{ Fixed more little things. Page labels
%   now work, even if you have no bibliography. You can output the
%   full addresses by hand with the \cs{outputfulladdresses}
%   command. Editors' Guide updated but still not finished. Revamped
%   splitbook.sh now produces PS file of each article for web page and
%   also an EPS file of each page for sending to the printers.}
% \changes{1.6}{2000/04/15}{ More cleaning up: full stops after
%   titles in index, running headers and footers in frontmatter,
%   fontsize in title of index, etc} \changes{1.5}{2000/04/15}{
%   Eliminated superscripts from titles in TOC and index}
% \changes{1.4}{2000/04/14}{ Added copyright and license info (GPL)}
% \changes{1.3}{2000/03/01}{ Corrected \cs{figcaption} to do the
%   caption formatting right Corrected font size of sec headings in
%   front matter} \changes{1.2}{2000/02/16}{ Added \cs{plate}
%   command} \changes{1.1}{2000/02/14}{ Added manuscript option,
%   corrected bugs} \changes{1.0}{2000/02/08}{ Now does a whole book
%   - lots of stuff tidied up} \changes{0.7}{1999/10/04}{ Increased
%   spacing between sections in single column styles. Suppressed
%   floats at bottom of pages (didn't work though) }
% \changes{0.6}{1999/09/24}{ Tweaked \cs{adjustfinalcols} (still not
%   perfect) Added proceedings option Added onecolumn/twocolumn
%   options (orthogonal to journal/proceedings) }
% \changes{0.5}{1999/09/23}{ Added \cs{fulladdresses} and
%   \cs{adjustfinalcols} (just about working)}
% \changes{0.4}{1999/09/22}{ Minor stylistic changes}
% \changes{0.3}{1999/06/28}{ Better mechanism for alternative
%   affiliations Added `guide' option for user guide}
% \changes{0.2}{1999/05/27}{ All basic functionality present, I
%   think} \changes{0.1}{1999/05/26}{ Just about working - still
%   missing Fig/Tab/Refs} \changes{0.0}{1999/05/21}{ So bored today,
%   decided to rewrite RMAA LaTeX macros from scratch}
%    \begin{macrocode}
\ProvidesClass{rmaa}[\rmaa@date\space v\rmaa@version\space RMxAA article
class]
%    \end{macrocode}
% \subsection{Option Handling}
% \label{sec:options}
% Declarations that are needed in options
% \changes{3.05}{2003/10/22}{Added natbib option}
% \changes{3.06}{2003/01/03}{Natbib option turned on permanently}
% \changes{3.07}{2003/01/11}{Added proof option}
% \changes{3.14}{2003/01/11}{Option to turn off natbib}
% \changes{3.18}{2003/10/22}{trying out doctex mode} 
% \changes{3.23}{2007/09/09}{Manuscript is now on by default.}
% \changes{3.26}{2011/08/15}{New option for review articles}
%    \begin{macrocode}
\newif\if@reviewarticle
\@reviewarticlefalse
\newif\if@proof
\@prooffalse
\newif\if@journal
\@journaltrue
\newif\if@printdates
\@printdatesfalse
\newif\if@proceedings
\@proceedingsfalse
\newif\if@manuscript
\@manuscripttrue
\newif\if@bigfont
\@bigfonttrue
\newif\if@guide
\@guidefalse
\newif\if@preprint
\@preprintfalse
\newif\if@rmtwocol
\@rmtwocoltrue
\newif\if@biblist
\@biblistfalse
\newif\if@book
\@bookfalse
\newif\if@natbib
\@natbibtrue
\newif\if@rmcite
\@rmcitefalse
\newif\if@titlevisible
\@titlevisibletrue
\newif\if@abstractpresent
\@abstractpresentfalse
\newif\if@printbanner
\@printbannerfalse
\newif\if@suppressfulladdresses
\@suppressfulladdressesfalse
\newif\if@newstyle
\@newstylefalse
\newif\if@debug
\@debugfalse
\newif\if@boldkeywords
\@boldkeywordsfalse
%    \end{macrocode}
% The option handling---mainly setting various booleans.
% \changes{2.2beta11}{2001/11/26}{Made twocol and
% newstyle the default for proceedings.}
% \changes{3.23}{2007/09/09}{Manuscript is now 10 pt instead of 12pt}
% \changes{3.23}{2007/09/10}{Full addresses are now automatically
% suppressed for proceedings style.}
% \changes{3.24}{2007/09/14}{New option \texttt{boldkeywords} that
% does the keywords in the old ugly way. Now off by default.}
%    \begin{macrocode}
\DeclareOption{manuscript}{\typeout{Manuscript for referee} 
  \@manuscripttrue\@journalfalse\@rmtwocolfalse
  \PassOptionsToClass{10pt}{article}}
\DeclareOption{bigfont}{\typeout{Big font manuscript requested by Gary
  Ferland} 
  \@manuscripttrue\@journalfalse\@rmtwocolfalse\@bigfonttrue
  \PassOptionsToClass{11pt}{article}}
\DeclareOption{journal}{\typeout{Main journal option}%
  \@journaltrue\@rmtwocoltrue\@printdatestrue\@manuscriptfalse
  \PassOptionsToClass{10pt}{article}}
\DeclareOption{guide}{\typeout{User guide option}%
  \@guidetrue\@journalfalse\@rmtwocoltrue\@manuscriptfalse
  \PassOptionsToClass{10pt}{article}}
\DeclareOption{preprint}{\typeout{Preprint option}%
  \@preprinttrue\@manuscriptfalse\PassOptionsToClass{10pt}{article}}
\DeclareOption{proof}{\typeout{Author proof option}%
  \@prooftrue\@manuscriptfalse}
\DeclareOption{proceedings}{\typeout{Proceedings option}%
  \@proceedingstrue\@journalfalse\@rmtwocoltrue\@newstyletrue
  \@suppressfulladdressestrue
  \@printdatesfalse\@manuscriptfalse\PassOptionsToClass{10pt}{article}} 
\DeclareOption{printdates}{\typeout{Force printing of received/accepted dates}%
  \@printdatestrue} 
\DeclareOption{book}{\typeout{Whole issue option}%
  \@booktrue} 
\DeclareOption{nonatbib}{\typeout{Turn off natbib support (required
    for compatibility with old rmaacite package).}%
  \@natbibfalse} 
\DeclareOption{newstyle}{\typeout{Use new layout style}%
  \@newstyletrue}
\DeclareOption{oldstyle}{\typeout{Use old layout style}%
  \@newstylefalse}
\DeclareOption{debug}{\typeout{Loads of annoying messages}%
  \@debugtrue}
%    \end{macrocode}
% These two need to be declared explicitly so that |\ExecuteOptions|
% actually does something.
%    \begin{macrocode}
\DeclareOption{twocolumn}{\@rmtwocoltrue}
\DeclareOption{onecolumn}{\@rmtwocolfalse}
\DeclareOption{twoside}{\PassOptionsToClass{\CurrentOption}{article}}
\DeclareOption{boldkeywords}{\@boldkeywordstrue}
%    \end{macrocode}
% Bold, upper-case keywords used to be the default. Now they must be
% turned on explicitly with the \texttt{boldkeywords} option. 
%    \begin{macrocode}
\DeclareOption{boldkeywords}{\@boldkeywordstrue}
%    \end{macrocode}
% Pass any other option along to the \texttt{article} class.
%    \begin{macrocode}
\DeclareOption*{\PassOptionsToClass{\CurrentOption}{article}}
%    \end{macrocode}
% Default options are to be two-sided and based on the standard
% \LaTeX{} \texttt{article} class,
% unless we are processing the entire book, in which case we base
% ourselves on the standard \texttt{book} class. Default is to use
% manuscript option unless otherwise specified. 
%    \begin{macrocode}
\ExecuteOptions{twoside,manuscript} \ProcessOptions\relax
\if@book
\if@rmtwocol\LoadClass[twocolumn]{book}\else\LoadClass{book}\fi
\newcommand{\abstract}{\relax}
\newcommand{\@maketitle}{\relax}
\else
\if@rmtwocol\LoadClass[twocolumn]{article}\else\LoadClass{article}\fi
\fi
%    \end{macrocode}
% Attempt to use natbib for citations 02 Jan 2003. 
% \changes{3.05}{2003/01/02}{If the natbib option is used, then set
% things up so that natbib's version of the bibliography is identical
% to the standard one.}
% \changes{3.09}{2003/01/23}{Explicitly require a recent version of
% natbib, so that we get a decent error message if the package is too
% old.} 
% \changes{3.23}{2007/09/09}{Now allow Bibtex to work properly} 
% \changes{3.25}{2009/10/01}{Change \cs{newcommand} tp \cs{def} when setting \cs{bibfont}. This allows it to work with all versions of natbib.} 
% \changes{3.25}{2010/02/06}{Updated copyright year.} 
% \changes{3.25}{2010/02/06}{Updated email and web addresses to \texttt{crya.unam.mx}.} 
%    \begin{macrocode}
\if@natbib
\RequirePackage{natbib}[1999/02/25]
\bibpunct{(}{)}{;}{a}{}{,}
\def\bibfont{\small}
\renewcommand\bibsection{\section*{REFERENCES}}
\setlength\bibsep{0pt}
\setlength\bibhang{1.5em}
\bibliographystyle{rmaa}
\fi
%    \end{macrocode}
% We want the \texttt{graphicx} package for the extended syntax of
% |\includegraphics|. We want the \texttt{ifthen} package for using in
% these macros (although we don't much yet). We explicitly load
% \texttt{babel} since some \LaTeX{} distributions messed up the
% hyphenation if we didn't.  
%    \begin{macrocode}
\RequirePackage{graphicx}
\RequirePackage{ifthen}
\IfFileExists{babel.sty}{%
%   \let\@restoreindent=\@afterindentfalse
  \RequirePackage[english]{babel}
%   \@restoreindent
  %% reinstate the definition of \@roman so as not to break .idx files
  \@ifundefined{org@@roman}{\relax}{\renewcommand{\@roman}[1]{\org@@roman{##1}}}
  }{%
  \typeout{Package babel.sty not found - don't worry about this}}
%    \end{macrocode}
% Try to load the \texttt{dcolumn} and \texttt{booktabs} packages if
% they can be found. These can produce nicer-looking tables. 
%    \begin{macrocode}
\IfFileExists{dcolumn.sty}{%
  \RequirePackage{dcolumn}
  }{%
  \typeout{Package dcolumn.sty not found - you will not be able to
    align columns on decimal points in tables}}
\IfFileExists{booktabs.sty}{%
  \RequirePackage{booktabs}
  }{%
  \typeout{Package booktabs.sty not found - your tables will not look
    as pretty as they ought too.}
  }
\IfFileExists{makeidx.sty}{%
  \RequirePackage{makeidx}
}{%
  \typeout{Package makeidx.sty not found - Author Index may not work.}
  }
%% commented out for moment since requires more modern version of LaTeX
%% \RequirePackage[latin1]{inputenc}
%    \end{macrocode}
% \changes{4.0}{2019/11/13}{Make \texttt{hyperref} and
% \texttt{etoolbox} into hard requirements, since we now use them
% extensively.}  We use hyperref to get internal and external links.
% But I don't like the colored boxes so I turn those off.  We use
% etoolbox for patching and robustifying at the moment.  In the
% future, I want to port a lot of other stuff over to it.  For
% instance, replace the ifthen package and the TeX ifs too.
%
% We target the latest version (1.7) of PDF, rather than the default
% of 1.5.  This reduces the chance of getting warnings from included
% graphics.  This is not strictly related to hyperref, but we put it
% here anyway.
%    \begin{macrocode}
\RequirePackage{etoolbox}
\RequirePackage{hyperref}
\hypersetup{hidelinks}
\pdfminorversion=7
%    \end{macrocode}
% We patch the \cs{@combinedblfloats} following a suggestion by Heiko
% Oberdiek at \url{https://tex.stackexchange.com/a/249743/3843}, which
% avoids problems when hyperlinks are split across
% pages. \changes{4.0}{2019/11/12}{Fix for issue where breaking
% hyperlink across pages gave an error.}
%    \begin{macrocode}
\urlstyle{same}
\newcommand\email{\begingroup \urlstyle{rm}\Url}
\patchcmd\@combinedblfloats{\box\@outputbox}{\unvbox\@outputbox}{}{%
  \errmessage{\noexpand\@combinedblfloats could not be patched}%
}%
%    \end{macrocode}
% \subsection{Other setup stuff}\label{sec:setup}
% Some declarations: 
%    \begin{macrocode}
\newif\if@needaddresses
\@needaddressestrue
\newif\if@frontmatter
\@frontmatterfalse
\newif\if@inpreamble
\@inpreambletrue
\newif\if@firstauthor
\@firstauthortrue
\newif\if@firstaffil
\@firstaffiltrue
\newif\if@firstaltaffil
\@firstaltaffiltrue
\newif\if@firstkwd
\@firstkwdtrue
\newif\if@noexpandADS
\@noexpandADSfalse
\newif\if@sessiontitlepages
\@sessiontitlepagesfalse
%    \end{macrocode}
% \begin{macro}{\@rmaaEmpty}
% \begin{macro}{\@rmaaArg}
% These are used in various commands to check if arg was left blank
%    \begin{macrocode}
\newcommand{\@rmaaEmpty}{}
\newcommand{\@rmaaArg}{}
%    \end{macrocode}  
% \end{macro}
% \end{macro}
%
% \subsection{Compatibility with old \LaTeX{} versions}
% \label{sec:compat}
% This package is based on the 1998/12/01 release of \LaTeX{}.
% Older releases may not define all the commands that we need.
% Luckily, we can use |\providecommand| to define these if necessary. 
% The following commands are not defined in the 1994/06/01 release. 
%
% \begin{macro}{\MakeUpperCase}
% \begin{macro}{\textsuperscript}
% \begin{macro}{\textbackslash}
% User commands:
%    \begin{macrocode}
\providecommand{\MakeUpperCase}[1]{\uppercase{#1}}
\providecommand{\textsuperscript}[1]{\ensuremath{^\textrm{#1}}}
\providecommand{\textbackslash}{\ensuremath{\backslash}}
%    \end{macrocode}  
% \end{macro}
% \end{macro}
% \end{macro}
% \begin{macro}{\protected@xdef}
% \begin{macro}{\restore@protect}
% \begin{macro}{\@unexpandable@protect}
% \begin{macro}{\hb@xt@}
% \LaTeX{} internal commands:
%    \begin{macrocode}
\providecommand{\protected@xdef}{%
        \let\@@protect\protect
        \let\protect\@unexpandable@protect
        \afterassignment\restore@protect
        \xdef}
\providecommand{\restore@protect}{\let\protect\@@protect}
\providecommand{\@unexpandable@protect}{\noexpand\protect\noexpand}
\providecommand{\hb@xt@}{\hbox to}
%    \end{macrocode}
% \end{macro}
% \end{macro}
% \end{macro}
% \end{macro}
%
% \subsection{Page layout}\label{sec:page-layout}
%
% The dimensions of the text were defined by Elizabeth. Text width is
% 169\,mm with equal margins on the two sides: 
%    \begin{macrocode}
\setlength{\textwidth}{169mm}
\setlength{\@tempdima}{\paperwidth}
\addtolength{\@tempdima}{-\textwidth}
\setlength{\oddsidemargin}{0.5\@tempdima}
\addtolength{\oddsidemargin}{-1in}
\setlength{\evensidemargin}{\oddsidemargin}
%    \end{macrocode}
% Text height is 225\,mm, top margin is 14\,mm. 
%    \begin{macrocode}
\setlength{\textheight}{225mm} 
\setlength{\topmargin}{-1in} 
\addtolength{\topmargin}{14mm}
%    \end{macrocode}
% Space between columns in two-column styles, header, footer,
% paragraph indentation:
%    \begin{macrocode}
\setlength{\columnsep}{9.5mm}
\setlength{\headsep}{0.7cm}
\setlength{\footskip}{0.75cm}
\setlength{\parindent}{1.5em}
%    \end{macrocode}
% 
% \begin{macro}{\@set@scalable@length}
% Try out a new scheme for dealing with various spacing dimensions and
% allowing the user to easily modify them. The command
% \begin{quote}
% |\set@scalable@length|\marg{name}\marg{defaultlength}\marg{stretch}\marg{shrink}
% \end{quote}
% does a 
% |\setlength|\marg{\cs{name}}\marg{defaultlength plus stretch minus shrink}
% and also defines the user command
% \textbackslash\meta{name}|Scale|\marg{factor} that can 
% be used to rescale length by \meta{factor}. Additionally, now
% defines the package command
% \textbackslash\meta{name}|@set@scale|\marg{FACTOR} 
% that also rescales the length by \meta{FACTOR}, so that the length
% becomes \meta{factor} $\times$ \meta{FACTOR} $\times$
% \meta{defaultlength}.   
%    \begin{macrocode}
\newcommand\@set@scalable@length[4]{%
  %% arg #1: NAME of length without leading `\'
  %%
  %% arg #2: default VALUE  
  \expandafter\newlength\csname @#1@default\endcsname
  \expandafter\setlength\csname @#1@default\endcsname{#2}%
  \expandafter\newlength\csname @@#1@@default\endcsname
  \expandafter\setlength\csname @@#1@@default\endcsname{#2}%
  %% arg #3: default STRETCH  
  \expandafter\newlength\csname @#1@stretch\endcsname
  \expandafter\setlength\csname @#1@stretch\endcsname{#3}%
  \expandafter\newlength\csname @@#1@@stretch\endcsname
  \expandafter\setlength\csname @@#1@@stretch\endcsname{#3}%
  %% arg #4: default SHRINK  
  \expandafter\newlength\csname @#1@shrink\endcsname
  \expandafter\setlength\csname @#1@shrink\endcsname{#4}%
  \expandafter\newlength\csname @@#1@@shrink\endcsname
  \expandafter\setlength\csname @@#1@@shrink\endcsname{#4}%
  %%
  %% User-defined scale \@NAME@scale is unity by default
  \expandafter\gdef\csname @#1@scale\endcsname{1}%
  %% Package-defined scale \@NAME@SCALE is unity by default
  \expandafter\gdef\csname @#1@SCALE\endcsname{1}%
  %%
  %% internal command \@reset@NAME should expand to 
  %% \setlength\@@NAME@@default{\@NAME@SCALE\@NAME@default}
  %% \setlength\NAME{\@NAME@scale\@@NAME@@default}
  %% in other words: scale x SCALE x default 
  \expandafter\gdef\csname @reset@#1\endcsname{%
    \expandafter\setlength\csname @@#1@@default\endcsname{%
      \csname @#1@SCALE\endcsname \csname @#1@default\endcsname
      }
    \expandafter\setlength\csname @@#1@@stretch\endcsname{%
      \csname @#1@SCALE\endcsname \csname @#1@stretch\endcsname
      }
    \expandafter\setlength\csname @@#1@@shrink\endcsname{%
      \csname @#1@SCALE\endcsname \csname @#1@shrink\endcsname
      }
    \expandafter\setlength\csname #1\endcsname{%
      \csname @#1@scale\endcsname \csname @@#1@@default\endcsname
      \@plus
      \csname @#1@scale\endcsname \csname @@#1@@stretch\endcsname
      \@minus
      \csname @#1@scale\endcsname \csname @@#1@@shrink\endcsname
      }}
  %%
  %% User command \NAMEScale should expand to 
  %% {\renewcommand\@NAME@scale{#1}\@reset@NAME}
  \expandafter\gdef\csname #1Scale\endcsname##1{%
    \expandafter\gdef\csname @#1@scale\endcsname{##1}%
    \csname @reset@#1\endcsname}
  %%
  %% Package command \NAME@set@SCALE should expand to 
  %% {\renewcommand\@NAME@SCALE{#1}\@reset@NAME}
  \expandafter\gdef\csname #1@set@SCALE\endcsname##1{%
    \expandafter\gdef\csname @#1@SCALE\endcsname{##1}%
    \csname @reset@#1\endcsname}
  %%
  %% initialize at default values with \@reset@NAME
  \csname @reset@#1\endcsname
  }
%    \end{macrocode}
% \end{macro}
% \begin{macro}{\parskip}
% First use of this new infrastructure: set interparagraph space to be
% zero but with a small amount of stretchiness.
%    \begin{macrocode}
%   \ifthenelse{\boolean{@manuscript}}{%
%     \@set@scalable@length{parskip}{0.5\baselineskip}{0.3\p@}{0.1\p@}
%   }{
    \@set@scalable@length{parskip}{0\p@}{0.3\p@}{0\p@}
%   }
%    \end{macrocode}
% \end{macro}
% 
% \begin{macro}{\abovesectionskip}
% \changes{4.0}{2019/11/12}{Use negative values for skip so that first paragraph after section headings is not indented.  This is something that I have been trying to fix for years!  It turns out that the anser was in the documentation for the standard command \cs{@startsection.}}
% \begin{macro}{\belowsectionskip}
% \begin{macro}{\aboveparaskip}
% \begin{macro}{\belowparaskip}
% Define some utility rubber lengths for use in
% section~\ref{sec:section-commands}. 
%    \begin{macrocode}
\newlength{\abovesectionskip}
\newlength{\belowsectionskip}
\newlength{\aboveparaskip}
\newlength{\belowparaskip}
%    \end{macrocode}
% The default values for these lengths are those traditionally used
% for the onecolumn style. 
%    \begin{macrocode}
\@set@scalable@length{abovesectionskip}{-3.5ex}{-4ex}{-2ex}
\@set@scalable@length{belowsectionskip}{2.3ex}{0ex}{2ex}
\@set@scalable@length{aboveparaskip}{3.25ex}{2ex}{0.5ex}
\@set@scalable@length{belowparaskip}{-1em}{\z@}{\z@}
%    \end{macrocode}
% \end{macro}\end{macro}\end{macro}\end{macro}
% We also make scalable some of the standard float parameters. These
% are the same defaults as in the standard article classes, but
% slightly more stretchy. For the twocolumn style, they are currently
% halved, see below. Note that I have also lumped |\abovecaptionskip|
% in with these, since it plays a similar role. There is also a
% |\belowcaptionskip| but it is set to zero in the standard classes. 
%    \begin{macrocode}
\@set@scalable@length{floatsep}{12\p@}{4\p@}{4\p@}
\@set@scalable@length{textfloatsep}{20\p@}{4\p@}{8\p@}
\@set@scalable@length{intextsep}{12\p@}{4\p@}{4\p@}
\@set@scalable@length{dblfloatsep}{12\p@}{4\p@}{4\p@}
\@set@scalable@length{dbltextfloatsep}{20\p@}{4\p@}{8\p@}
\@set@scalable@length{abovecaptionskip}{10\p@}{\z@}{\z@}
%    \end{macrocode}
% \begin{macro}{\@before@title@skip}         
% \begin{macro}{\@before@author@skip}        
% \begin{macro}{\@between@author@skip}       
% \begin{macro}{\@before@abstract@skip}      
% \begin{macro}{\@after@abstracthead@skip}   
% \begin{macro}{\@between@abstract@skip}     
% \begin{macro}{\@before@kwd@skip}           
% \begin{macro}{\@before@dates@skip}           
% Spacing parameters used in |\@maketitle|. These are not changed by
% default except by the user, so they don't need to go into
% |\@package@default@lengths|.  
%    \begin{macrocode}
\newlength{\@before@title@skip}         
\newlength{\@before@author@skip}        
\newlength{\@between@author@skip}       
\newlength{\@before@abstract@skip}      
\newlength{\@after@abstracthead@skip}   
\newlength{\@between@abstract@skip}     
\newlength{\@before@kwd@skip}           
\newlength{\@before@dates@skip}           
\@set@scalable@length{@before@title@skip}{2em}{\z@}{\z@}
\@set@scalable@length{@before@author@skip}{1.5em}{\z@}{\z@}
\@set@scalable@length{@between@author@skip}{0.5em}{\z@}{\z@}
\@set@scalable@length{@before@abstract@skip}{3ex}{1ex}{0.5ex}
\@set@scalable@length{@after@abstracthead@skip}{1ex}{\z@}{\z@}
\@set@scalable@length{@between@abstract@skip}{2ex}{1ex}{0.5ex}
\@set@scalable@length{@before@kwd@skip}{1ex}{0.5ex}{0.2ex}
\@set@scalable@length{@before@dates@skip}{2ex}{0.5ex}{0.2ex}
%    \end{macrocode}
% \end{macro}\end{macro}\end{macro}\end{macro}
% \end{macro}\end{macro}\end{macro}\end{macro}
% 
% Spacing above and below displayed equations also wants to be done
% like this, but it is more complicated since the standard article
% classes make these change with the fontsize. Something to do later,
% maybe. 
% 
%    \begin{macrocode}
%    \end{macrocode}
% \begin{macro}{\@package@default@lengths}
% The value actually used for these lengths is controlled by the
% \meta{SCALE} parameter, which depends on the style option that is in
% effect. This is set using the |\@package@default@lengths| command,
% which must be called every time the style changes (e.g., changing
% from onecolumn to twocolumn). 
%    \begin{macrocode}
\newcommand{\@package@default@lengths}{%
  \ifthenelse{\boolean{@rmtwocol}}{%
    %% twocolumn mode
    \abovesectionskip@set@SCALE{0.5}
    \belowsectionskip@set@SCALE{0.5}
    \aboveparaskip@set@SCALE{0.5}
    \belowparaskip@set@SCALE{1}
    \floatsep@set@SCALE{0.5}
    \textfloatsep@set@SCALE{0.5}
    \intextsep@set@SCALE{0.5}
    \dblfloatsep@set@SCALE{0.5}
    \dbltextfloatsep@set@SCALE{0.5}
    }{%
    %% onecolumn mode
    \abovesectionskip@set@SCALE{1}
    \belowsectionskip@set@SCALE{1}
    \aboveparaskip@set@SCALE{1}
    \belowparaskip@set@SCALE{1}
    \floatsep@set@SCALE{1}
    \textfloatsep@set@SCALE{1}
    \intextsep@set@SCALE{1}
    \dblfloatsep@set@SCALE{1}
    \dbltextfloatsep@set@SCALE{1}
    }%
  }
\@package@default@lengths
%    \end{macrocode}
% \end{macro}
% \begin{macro}{\RescaleLengths}
% \begin{macro}{\RescaleSecLengths}
% \begin{macro}{\RescaleFloatLengths}
% \begin{macro}{\RescaleTitleLengths}
% User commands that will fiddle with all the scalable lengths. 
%    \begin{macrocode}
\newcommand{\RescaleLengths}[1]{%
  \RescaleSecLengths{#1}\RescaleFloatLengths{#1}\RescaleTitleLengths{#1}}
\newcommand{\RescaleSecLengths}[1]{%
%% sections
  \abovesectionskipScale{#1}
  \belowsectionskipScale{#1}
  \aboveparaskipScale{#1}
  }
\newcommand{\RescaleFloatLengths}[1]{%
%% floats
  \floatsepScale{#1}
  \textfloatsepScale{#1}
  \intextsepScale{#1}
  \dblfloatsepScale{#1}
  \dbltextfloatsepScale{#1}
  \abovecaptionskipScale{#1}
  }
\newcommand{\RescaleTitleLengths}[1]{%
%% maketitle
  \@before@title@skipScale{#1}           
  \@before@author@skipScale{#1}        
  \@between@author@skipScale{#1}       
  \@before@abstract@skipScale{#1}      
  \@after@abstracthead@skipScale{#1}   
  \@between@abstract@skipScale{#1}     
  \@before@kwd@skipScale{#1}           
  \@before@dates@skipScale{#1}           
  }
%    \end{macrocode}
% \end{macro}\end{macro}\end{macro}\end{macro}
% \changes{2.2beta14}{2002/03/17}{Finished implementation of new
% infrastructure for scalable lengths. Converted the above/below
% section skips to use it. Also, all the floatsep parameters. Also,
% all the vertical lengths that are used by |\@maketitle|. }

% \begin{macro}{\MSextramargin}
% For the \texttt{manuscript} option we have an extra 0.3\,inch
% vertical margin and 1\,inch horizontal margin for marginal notes (by
% hand or with |\marginpar|. 
%
% \changes{3.23}{2007/09/09}{Reduce baselinestretch in
%   \texttt{manuscript} style to 1.2. No longer use raggedright for
%   manuscript.}
%
% \changes{3.24}{2007/09/14}{Further reduction of baselinestretch in
%   \texttt{manuscript} style to 1.1 since we now have wider margins.}
%
% \changes{3.24}{2007/09/14}{Split \texttt{manuscript}'s extra margin
%   into separate horizontal and vertical parts.}
% 
%    \begin{macrocode}
\newlength{\MSextraHmargin}\newlength{\MSextraVmargin}
\setlength{\MSextraHmargin}{1in} 
\setlength{\MSextraVmargin}{0.3in} 
\if@manuscript
   \renewcommand{\baselinestretch}{1.1}
   \addtolength{\textwidth}{-2\MSextraHmargin}
   \addtolength{\oddsidemargin}{\MSextraHmargin}
   \setlength{\evensidemargin}{\oddsidemargin}
   \addtolength{\textheight}{-2\MSextraVmargin}
   \addtolength{\topmargin}{\MSextraVmargin}
%    \raggedright
\fi
%    \end{macrocode}
% \end{macro}
% Float parameters: don't allow bottom floats by default. 
%    \begin{macrocode}
\renewcommand{\fps@figure}{tp}
\renewcommand{\fps@table}{tp}
%    \end{macrocode}
%
% \begin{macro}{\rmaa@AuthorWidthFrac}
% \begin{macro}{\rmaa@AbstractWidthFrac}
% These parameters control what fraction of the total text width can
% be used by the list of authors and the abstract/resumen,
% respectively.  
%    \begin{macrocode}
\newcommand{\rmaa@AuthorWidthFrac}{1.0}
\newcommand{\rmaa@AbstractWidthFrac}{0.75}
%    \end{macrocode}
% Use a wide abstract to save space in the new style.
%    \begin{macrocode}
\if@newstyle\renewcommand{\rmaa@AbstractWidthFrac}{1.0}\fi
%    \end{macrocode}
% \end{macro}
% \end{macro}
%
% Some things to consider in future versions. Trying to avoid too many
% widows, etc. Apparently, the following will work best with
% \cs{flushbottom}. 
%    \begin{macrocode}
%%\widowpenalty=10000
%%\clubpenalty=10000
%    \end{macrocode}
% Also, should perhaps try getting rid of underfull/overfull box warnings. 
% \changes{3.10}{2003/01/24}{Set large values for \cs{hbadness},
% \cs{vbadness}, \cs{hfuzz}, \cs{vfuzz}. To reduce amount of stuff
% that gets written to the terminal every time you compile something.}
%    \begin{macrocode}
\hbadness=10000
\vbadness=10000
\setlength\hfuzz{\maxdimen}
\setlength\vfuzz{\maxdimen}
%    \end{macrocode}
% \subsection{Various text styles for different contexts}
% \label{sec:text-styles} 
% \begin{macro}{\rmaa@AffilStyle}
% \begin{macro}{\rmaa@PagesStyle}
% \begin{macro}{\rmaa@YearStyle}
% \begin{macro}{\rmaa@VolumeStyle}
% \begin{macro}{\rmaa@VolumenStyle}
% \begin{macro}{\rmaa@JournalNameStyle}
% \begin{macro}{\rmaa@ReceivedStyle}
% \begin{macro}{\rmaa@AbstractHeaderStyle}
% \begin{macro}{\rmaa@JournalNameStyle}
% Style commands for the different fields used by |\maketitle| in
% constructing the title page of the article. These have the form
% |\rmaa@XYZStyle{TEXT}|. I'm not sure these were really a good idea,
% but never mind \dots 
%    \begin{macrocode}
\newcommand{\rmaa@AffilStyle}[1]{{\normalsize\rmfamily#1}}
\newcommand{\rmaa@PagesStyle}[1]{{\normalsize\rmfamily#1}}
\newcommand{\rmaa@YearStyle}[1]{{\normalsize\rmfamily#1}}
\newcommand{\rmaa@VolumeStyle}[1]{{\normalsize\rmfamily\bfseries#1}}
\newcommand{\rmaa@VolumenStyle}[1]{{\Large\sffamily\bfseries#1}}
\newcommand{\rmaa@JournalNameStyle}[1]{{\normalsize\itshape#1}}
\newcommand{\rmaa@ReceivedStyle}[1]{{\normalsize\rmfamily\itshape#1}}
\newcommand{\rmaa@AbstractHeaderStyle}[1]{%
  {\normalsize\rmfamily\MakeUpperCase{#1}}}
\newcommand{\rmaa@TitleStyle}[1]{%
  {\large\bfseries\rmfamily\boldmath\MakeUpperCase{#1}\par}}
%    \end{macrocode}
% \end{macro}
% \end{macro}
% \end{macro}
% \end{macro}
% \end{macro}
% \end{macro}
% \end{macro}
% \end{macro}
% \end{macro}
%
% \begin{macro}{\rmaa@FrontSecStyle}
% Headings of ``front sections'' (see section~\ref{sec:front-sections})
%    \begin{macrocode}
\newcommand{\rmaa@FrontSecStyle}[1]{%
  {\normalsize\bfseries\rmfamily\MakeUpperCase{#1}}}
%    \end{macrocode}
% \end{macro}
%
% \begin{macro}{\rmaa@TOCTitleStyle}
% \begin{macro}{\rmaa@TOCPartStyle}
% \begin{macro}{\rmaa@IndexTitleStyle}
% \begin{macro}{\rmaaTOCAuthorStyle}
% \begin{macro}{\rmaaIndexAuthorStyle}
% Commands to set style of TOC and index entries.
% Rewritten to avoid unnecessary font commands in the .toc and .ind files
%    \begin{macrocode}
\newcommand{\rmaa@TOCTitleStyle}[1]{\MakeUpperCase{#1}}
\newcommand{\rmaa@TOCPartStyle}[1]{\rmaa@TOCTitleStyle{#1}}
\newcommand{\rmaa@IndexTitleStyle}[1]{#1}
\newcommand{\rmaaTOCAuthorStyle}[1]{%
  {\itshape#1}}
\newcommand{\rmaaIndexAuthorStyle}[1]{%
  \protect\mbox{\bfseries#1}}
%    \end{macrocode}
% \end{macro}
% \end{macro}
% \end{macro}
% \end{macro}
% \end{macro}
%
% \begin{macro}{\rmaa@KeywordStyle}
% \begin{macro}{\rmaa@KeywordHeaderStyle}
% The style of the keywords depends on whether we are using the
% \texttt{newstyle} option or not. If we are, then they are small. 
%    \begin{macrocode}
\ifthenelse{\boolean{@newstyle}}{
  \newcommand{\rmaa@KeywordStyle}[1]{%
    {\small%
      \if@boldkeywords 
      \bfseries\MakeUpperCase{\raggedright #1}
      \else
      \raggedright #1
      \fi
    }}
  \newcommand{\rmaa@KeywordHeaderStyle}[1]{%
    {\small\itshape#1}}}
{
  \newcommand{\rmaa@KeywordStyle}[1]{%
    {\normalsize%
      \if@boldkeywords\bfseries\fi\raggedright #1}}
  \newcommand{\rmaa@KeywordHeaderStyle}[1]{%
    {\normalsize\itshape#1}}}
%    \end{macrocode}
% \end{macro}
% \end{macro}
%
% \subsection{Page styles}\label{sec:page-styles}
% \begin{macro}{\ps@rmaa}
% The running headers for normal pages:
%    \begin{macrocode}
\newcommand{\ps@rmaa}{%
%    \end{macrocode}
% nothing in either foot; 
%    \begin{macrocode}
  \renewcommand\@oddfoot{\relax}%
  \renewcommand\@evenfoot{\relax}%
%    \end{macrocode}
% page number and short author list in the even heads; 
%    \begin{macrocode}
  \renewcommand\@evenhead{\rlap{\thepage}\hfil\@shortauthor\hfil}%
%    \end{macrocode}
% short title and page number in the odd heads.
%    \begin{macrocode}
  \renewcommand\@oddhead{\hfil\@shorttitle\hfil\llap{\thepage}}
  }
%    \end{macrocode}
% \begin{macro}{\ps@rmaafirstpage}
% The header of the first page of each
% article:
%    \begin{macrocode}
\newcommand{\ps@rmaafirstpage}{%
  \renewcommand\@oddfoot{\hfil\thepage}%
%    \end{macrocode}
% footer corrected (06 Jun 2000) to have page on outside edge;
%    \begin{macrocode}
  \renewcommand\@evenfoot{\thepage\hfil}%
%    \end{macrocode}
% Header has journal name, volume, and page numbers for a regular
% article, or similar for special cases. 
% \changes{4.0}{2019/11/03}{Change to three-line header with copyright info and DOI.}
%    \begin{macrocode}
\if@preprint
\renewcommand\@oddhead{%
  \if@proceedings 
  \textsc{To appear in ``\@conftitle{} (\rmaa@YearStyle{\@year})''}\hfil
  \textit{RMxAC}
  \else
  \textsc{To appear in} \@journalname\hfil
  \fi}\fi%
\if@guide
\renewcommand\@oddhead{\textsc{Documentation} (\rmaa@YearStyle{\@year})\hfil}
\fi%
\if@manuscript
\renewcommand\@oddhead{Manuscript for \@journalname{} (\rmaa@YearStyle{\@year})\hfil}
\fi%
\if@journal
\renewcommand\@oddhead{%
  \parbox{\textwidth}{
    \@journalname, \rmaa@VolumeStyle{\@volume},
    \rmaa@PagesStyle{\@artpages}\
    (\rmaa@YearStyle{\@year})\hfil\\[1pt]
    {\small\textcopyright{} \@year:
      Instituto de Astronom\'\i{}a,
      Universidad Nacional Aut\'onoma de M\'exico}\hfil\\[1pt]
    {\small\@DOI@url}\hfil
  }%
}\fi%
  \renewcommand\@evenhead{\@oddhead}}
%    \end{macrocode}
% By default, use the |\ps@rmaa| page style.
%    \begin{macrocode}
\pagestyle{rmaa}
%    \end{macrocode}
% \end{macro}
% \end{macro}
%
% \subsection{Preamble commands}\label{sec:preamble} 
% Everything you need for setting the title page, TOC and index
% entries, etc.
% \subsubsection{Commands for the title page}
% \label{sec:title-commands}
% \begin{macro}{\resumen}
% \begin{macro}{\@resumen}
% \begin{macro}{\abstract}
% \begin{macro}{\@abstract}
% User commands for setting the resumen and abstract. 
% These should be used in the document header. The internal commands
% |\@resumen| and |\@abstract| are used to hold the actual text
% pending output via the |\maketitle| command. 
%    \begin{macrocode}
\newif\if@setresumen
\newcommand{\resumen}[1]{%
%%  \renewcommand{\@resumen}{{\selectlanguage{spanish}#1}}%
  \renewcommand{\@resumen}{{\ignorespaces#1}}%
  \@setresumentrue}
\renewcommand{\abstract}[1]{%
  \renewcommand{\@abstract}{{\ignorespaces#1}}
  \renewcommand\ADS@abstract{#1}
%% Not a good idea, since then author macros are not expanded and
%% cause errors later. However, by not doing it we are vulnerable to
%% fragile commands.....
  \if@noexpandADS\@onelevel@sanitize{\ADS@abstract}\fi
  \@abstractpresenttrue
  }
\newcommand{\ADSnoexpand}{\@noexpandADStrue}
%    \end{macrocode}
% \begin{macro}{\boldabstract}
% This is the version for doing 1-page conference proceedings. Still
% save the text for ADS but then write it out immediately as a bold
% paragraph. 
%    \begin{macrocode}
\newcommand{\boldabstract}[1]{%
  \gdef\ADS@abstract{#1}%
  \renewcommand{\@abstract}{#1}%
  \if@noexpandADS\@onelevel@sanitize{\ADS@abstract}\fi%
  \noindent\textbf{\boldmath #1}
  }
%    \end{macrocode}
% \end{macro}
% Default value in each case is to print out a warning if not set
%    \begin{macrocode}
  \newcommand{\@resumen}{%
    \@setresumenfalse \ClassWarningNoLine{rmaa}{No Spanish abstract
      (resumen) specified - if you do not speak Spanish, this will be
      translated by the editors}%
    \textbf{Favor de proporcionar un resumen en espa\~nol. If you are
      unable to translate your abstract into Spanish, the editors will
      do it for you.} }
  \newcommand{\@abstract}{\ClassWarningNoLine{rmaa}{No abstract
      specified}}
%    \end{macrocode}
% \end{macro}
% \end{macro}
% \end{macro}
% \end{macro}
%
% \begin{macro}{\title}
% \begin{macro}{\@rawtitle}
% \begin{macro}{\@TOCtitle}
% \begin{macro}{\rmaaIndexTitle}
% Similarly, |\title| is the preamble command to set the title of the
% article. This is used in various places, typeset in different
% styles, hence the use of the internal |\@rawtitle|, |\@title|,
% |\rmaaIndexTitle|, and |\@TOCtitle| commands, each of which uses the
% relevant style from section~\ref{sec:text-styles}
%    \begin{macrocode}
\newcommand{\@rawtitle}{Raw Title}
\newcommand{\@TOCtitle}{\rmaa@TOCTitleStyle{Title for TOC}}
\newcommand{\rmaaIndexTitle}{\rmaa@IndexTitleStyle{Title for index}.}
\newcommand{\@TOCaftertitle}{}
\renewcommand{\title}[1]{%
%%  {\ADS@sanitize\xdef\@rawtitle{#1}}
  \renewcommand\@rawtitle{\noexpand#1}
  \let\ADS@title=\@rawtitle
  \renewcommand{\@title}{\rmaa@TitleStyle{#1}}%
  \renewcommand{\@TOCtitle}{\rmaa@TOCTitleStyle{#1}\@TOCaftertitle}%
  \renewcommand{\rmaaIndexTitle}{\rmaa@IndexTitleStyle{#1\@TOCaftertitle}.}%
  }
%    \end{macrocode}
% \end{macro}
% \end{macro}
% \end{macro}
% \end{macro}
% 
% All the weird and whacky ways of indicating the author
% affiliations. 
% \begin{macro}{\affil}
% First, the way that sets it directly underneath the authors name.
% Note that |\affil| trumps all other affiliations for the purposes of
% getting into the ADS stuff (we use |\gdef|).   
% \dots 
%    \begin{macrocode}
\newcommand{\affil}[1]{%
  \\[1.5ex]\rmaa@AffilStyle{#1}\\[1.5ex]%
  \if@firstaffil\gdef\ADS@affil{#1}\@firstaffilfalse\fi%
  }
\renewcommand{\and}{and\\[1.5ex]}
%    \end{macrocode}
% \begin{macro}{\altaffil}
% Alternative affiliations are footnoted. 
%    \begin{macrocode}
\newcommand{\altaffil}[1]{%
  \if@firstaltaffil\gdef\ADS@affil{#1}\@firstaltaffilfalse\fi%
  \thanks{#1}}
%    \end{macrocode}
% \begin{macro}{\altaffilmark}
% \begin{macro}{\altaffiltext}
% Alternative-alternative mechanism for affiliations---should not be
% mixed with |\altaffil| (unless you know what you are doing) or there
% will be problems with numbering.
%    \begin{macrocode}
\newcommand{\altaffilmark}[1]{\textsuperscript{#1}}
\newcommand{\altaffiltext}[2]{%
  \if@firstaltaffil\gdef\ADS@affil{#2}\@firstaltaffilfalse\fi
  \stepcounter{footnote}%
  \protected@xdef\@thanks{\@thanks
    \protect\footnotetext[#1]{#2}}} 
%    \end{macrocode}
% \end{macro}
% \end{macro}
% \end{macro}
% \end{macro}
%
% \begin{macro}{\keywords}
% \begin{macro}{\@keywords}
% \begin{macro}{\@keywordheader}
% \label{keywords}Sundry minor preamble commands. All
% self-explanatory, I hope. |\keywords| is the old user interface
% (deprecated). New articles should use |\addkeyword| once for each
% keyword. That way the ADS record comes out right too (comma rather
% than em-dash separator).
% \changes{2.2beta12}{2001/12/07}{\cs{addkeyword} New user interface
%   for keywords.} 
% \changes{2.2beta13}{2002/03/11}{\cs{addkeyword} Improved behaviour
%   with .ads file. Use
%   \cs{ADS@sanitize} to change definitions of `~' and \cs{ion}.}
% \changes{3.23}{2007/09/09}{\cs{addkeyword} now signals a compilation
% error if the author tries to add more than 6 keywords.}
% \changes{3.24}{2007/09/14}{kwdcounter is now reset in
% \cs{@begofarticlehook} to prevent error in book style} 
%    \begin{macrocode}
\newcommand{\@kwdsep}{ --- }
\newcommand{\ADS@kwdsep}{, }
\newcounter{kwdcounter}\setcounter{kwdcounter}{0}
\newcommand{\addkeyword}[1]{%
  \addtocounter{kwdcounter}{1}
  \ifthenelse{\thekwdcounter > 6}{
    \ClassError{rmaa}{Too many keywords. A maximum of 6 keywords is allowed.}{}
  }{}
  \if@firstkwd%
  \renewcommand\@keywords{\if@boldkeywords\uppercase{#1}\else #1\fi}
  {\ADS@sanitize\def\ion{\plainion}\gdef\ADS@keywords{#1}}
  \@firstkwdfalse
  \else%
  \edef\@keywords{\@keywords\@kwdsep \if@boldkeywords\uppercase{#1}\else #1\fi}%
  {\ADS@sanitize\def\ion{\plainion}\xdef\ADS@keywords{\ADS@keywords\ADS@kwdsep #1}}%
  \fi}
\newcommand{\ADS@sanitize}{%
  \def\mathrm{}
  \def\@{\ }
  \def~{\ }
  \def\langle{<}
  \def\rangle{>}
  \def\log{log }
  \def\sin{sin }
  \def\cos{cos }
  \def\tan{tan }
  \def\fd{^d.}
  \def\fh{^h.}
  \def\fm{^m.}
  \def\fs{^s.}
  \def\fdg{�.}
  \def\farcm{'.}
  \def\farcs{''.}
  \def\arcdeg{�}
  \def\,{ }
  }
\newcommand{\keywords}[1]{%
  \renewcommand{\ADS@keywords}{#1}
  \renewcommand{\@keywords}{\uppercase{#1}}
  }
\newcommand{\@keywordheader}{\rmaa@KeywordHeaderStyle{Key Words:}}
\newcommand{\@keywords}{%
  \ClassWarningNoLine{rmaa}{No keywords specified}%
  Keywords go here
  }
%    \end{macrocode}
% \begin{macro}{\shorttitle}
% \begin{macro}{\shortauthor}
% \begin{macro}{\@shorttitle}
% \begin{macro}{\@shortauthor}
% Commands for setting the running headers.
%    \begin{macrocode}
\newcommand{\shorttitle}[1]{%
  \renewcommand\@shorttitle{\MakeUpperCase{#1}}}
\newcommand{\shortauthor}[1]{%
  \renewcommand\@shortauthor{\MakeUpperCase{#1}}}
\newcommand{\@shorttitle}{SHORT TITLE GOES HERE}
\newcommand{\@shortauthor}{SHORT AUTHOR GOES HERE}
%    \end{macrocode}
% \end{macro}
% \end{macro}
% \end{macro}
% \end{macro}
% \end{macro}
% \end{macro}
% \end{macro}


% \subsubsection{The full addresses}
% \label{sec:full-addresses}
% \begin{macro}{\fulladdresses}
% \begin{macro}{\@fulladdresses}
% \begin{macro}{\@fulladdheight}
% \begin{macro}{\@fulladddepth}
% \begin{macro}{\@fulladdbox}
% Aaghh, the full addresses! A perennial source of problems. 
%    \begin{macrocode}
\newlength{\@fulladdheight}
\newlength{\@fulladddepth}
\newsavebox{\@fulladdbox}
%% \fulladdresses command modified (11 May 2000) so that we can put
%% \fulladdresses{} in the article to suppress the output of the
%% addresses
\newcommand{\@FinalFigure}{\relax}
\newcommand{\FinalFigure}[1]{\renewcommand{\@FinalFigure}{#1}}
\newcommand{\fulladdresses}[1]{%
  \renewcommand{\@rmaaArg}{#1}
  \ifx\@rmaaArg\@rmaaEmpty
     \renewcommand{\@fulladdresses}{}
  \else
     \renewcommand{\@fulladdresses}{
       \normalsize%
       \ifx\@FinalFigure\relax\else\vfill\@FinalFigure\vfill\fi
%%       \vspace*{1ex}
       %% do addresses as a list (copied from thebibliography environment) 
       \list{\null}{\leftmargin 1.5em\topsep\z@\labelwidth\z@\labelsep\z@\itemindent -1.5em
         \itemsep\z@\parsep\z@\usecounter{enumi}}
       \sloppy\clubpenalty4000\widowpenalty4000
       \sfcode`\.=1000\relax
       #1
       \endlist
       }
  \fi
  }
\newcommand{\@fulladdresses}{Full addresses go here}
%    \end{macrocode}
% \begin{macro}{\suppressfulladdresses}
% New user-level command for suppressing full addresses should you so
% desire. Currently used in the single-page-poster style. 
%    \begin{macrocode}
\newcommand{\suppressfulladdresses}{\@suppressfulladdressestrue\@needaddressesfalse}
%    \end{macrocode}
% \end{macro}
% \end{macro}
% \end{macro}
% \end{macro}
% \end{macro}
% \end{macro}

% \begin{macro}{\printdates}
% \begin{macro}{\noprintdates}
% Command to activate the printing out of the received/accepted
% dates. This can also be done using the \texttt{printdates} option,
% which is the default for \texttt{journal}, but not for \texttt{proceedings}.
%    \begin{macrocode}
\newcommand{\printdates}{\@printdatestrue}
\newcommand{\noprintdates}{\@printdatesfalse}
%    \end{macrocode}
% \end{macro}
% \end{macro}
% \begin{macro}{\@outputaddressesifnecessary}
% Command to output the addresses if we haven't already done so
% This is called by: 
% \begin{enumerate}
% \item The end of |thebibliography| environment
% \item The |\outputfulladdresses| command
% \end{enumerate}
%    \begin{macrocode}
\newcommand{\@outputaddressesifnecessary}{%
  \ClassInfo{rmaa}{Calling @outputaddressesifnecessary}
  \ifthenelse{\boolean{@needaddresses}}{%
    %% If in one column style, or in the first column of a two column
    %% style then we can try to put it on this page 
    \ifthenelse{\boolean{@firstcolumn} \or \not\boolean{@rmtwocol}}{%
      \@outputaddresses}{%
      %% otherwise give warning it will be on a separate page
      \ClassWarning{rmaa}{Full addresses will go on a separate
        page (please see User Guide for how to remedy this)}}%
    %% If we haven't managed yet, put it on its own page
    \ifthenelse{\boolean{@needaddresses}}{%
      \if@debug{\ClassWarning{rmaa}{In this silly bit!}}\fi
      \clearpage\vfill\@outputaddresses}{}%
    }{
    %% If we don't need them, we still might need to do the eoa hook 
    \ifthenelse{\boolean{@suppressfulladdresses}}{%
      \if@debug{\ClassWarning{rmaa}{XXX}}\fi
%%      \if@book\else\@endofarticlehook\null\fi
      }{}
    }%
  }  
%    \end{macrocode}
% \end{macro}

% \begin{macro}{\outputfulladdresses}
% A user interface to the above command, in case there is no biblio
% (11 May 2000) 
%    \begin{macrocode}
\newcommand{\outputfulladdresses}{\@outputaddressesifnecessary}
%    \end{macrocode}
% \end{macro}

% \begin{macro}{\adjustfinalcols}
% This command starts a new column and should be 
% used to ensure that the columns balance on the last page. It also
% outputs the full 
% author addresses to ensure that this occurs in the left hand
% column. 
% The |\par| command is so that the current line will be taken account
% of when working out |\pagetotal|
%    \begin{macrocode}
\newcommand{\adjustfinalcols}{\if@twocolumn\goodbreak\@outputaddresses\newpage\noindent\fi}
%    \end{macrocode}
% \end{macro}

% \begin{macro}{\@outputaddresses}
% Command to do the work of outputting the addresses. 
% This can be called by 
% \begin{enumerate}
% \item an explicit |\adjustfinalcols|
% \item via |\@outputaddressesifnecessary|, q.v.
% \end{enumerate}
%
%    \begin{macrocode}
\newcommand{\@outputaddresses}{%
  % Allow option of suppressing the full addresses, in which case
  % we just run the hook if we are not in a book style
  \if@suppressfulladdresses
%%     \if@book\else\@endofarticlehook\null\fi
  \else
  %% save the full addresses in a box for output later
  \sbox{\@fulladdbox}{\parbox{\textwidth}{\@fulladdresses}}
  % calculate height of address box  
  \settoheight{\@fulladdheight}{\usebox{\@fulladdbox}}
  \settodepth{\@fulladddepth}{\usebox{\@fulladdbox}}
  \addtolength{\@fulladdheight}{\@fulladddepth}
  \if@debug
  { \typeout{Address total height: \the\@fulladdheight}
    \typeout{Text height: \the\pagetotal}
    \typeout{Goal height: \the\pagegoal}
    }
  \fi
  % reuse \@fulladdheight to mean the total height of text+addresses
  \addtolength{\@fulladdheight}{\pagetotal}
%%  \addtolength{\@fulladdheight}{-\baselineskip}
  \if@debug\typeout{Adjusted text + addresses height: \the\@fulladdheight}\fi
  % check it will fit on current page
  \ifdim\@fulladdheight<\pagegoal
     %% must shift it to the left
     %% to compensate for the hanging indentation in the reference list. 
     %% NOTE: this is not foolproof!
     %% NB. This will give overfull hbox warnings, which is a bit ugly.
     \vfill\noindent\if@biblist\hspace*{-\leftmargin}\fi
     \usebox{\@fulladdbox}
     % Set this to show we've done it
     \@needaddressesfalse
     % and assume we are at the end of the article
     % so that we do the EOA hook if not in the book style
%%     \if@book\else\@endofarticlehook\fi
  \else
     % if not, just do nothing. With luck, this function will get 
     % called again on the next page
     \if@debug{\typeout{Holding over full addresses to next page}}\fi
  \fi\fi
  }
%    \end{macrocode}
% \end{macro}

% \subsubsection{Commands for editors}
% \label{sec:editor-commands}
% Some user/editor commands for doing the TOC and index in the book
% option. 
% They are defined here rather in the book-specific part since they
% need to be always available in order to avoid breaking documents. 
% \begin{macro}{\listofauthors}
% \begin{macro}{\@rawauthors}
% \begin{macro}{\rmaaTOCauthors}
% The command |\listofauthors| formats an author list for use in the
% TOC. It should contain the same names as in the |\author| command,
% but without the affiliation garbage and using |\&| instead of
% |and|. 
%    \begin{macrocode}
\newcommand{\@rawauthors}{Raw Authors}
\newcommand{\rmaaTOCauthors}{\noexpand\rmaaTOCAuthorStyle{Author list for TOC}}
\newcommand{\listofauthors}[1]{%
  \renewcommand{\@rawauthors}{#1}
  \renewcommand\rmaaTOCauthors{\noexpand\rmaaTOCAuthorStyle{#1}}}
%    \end{macrocode}
% \begin{macro}{\indexauthor}\label{indexauthor}
% The command |\indexauthor| produces an index entry for each
% author. 
% This command must come \emph{after} the
% |\title| and the |\listofauthors| commands. Also, there should be one 
% |\indexauthor| command for \emph{each} author. We also add the
% author onto the end of the |\ADS@authors| macro and construct the
% |\ADS@authorinitial| by a most cunning method (see below), but only
% for the first author. 
% 
% No, No, NO! This was all wrong. Backed out completely: 
% \begin{quote}
% Completely rewritten (19 Dec 2002) since we now use the index class, so 
% as to be able to alphabetize accented characters. A side effect of this is
% that |\index| no longer expands out the argument after the |@|, so we 
% instead add all the |\index| commands to a list (|\indexauthor@deferlist|), 
% which is executed in the |\@begofarticlehook|. 
% 
% Note the stuff with |\toks0| is copied from Exercise 20.15 of the \TeX{}book---it is
% to make sure the latin1-accented name gets passed into the idx
% file. 
% \end{quote}
% We now go back to the old way of doing things, using an optional
% argument to the \cs{indexauthor} command that has the name without
% accents. 
% \changes{3.08}{2003/01/16}{Remerged changes from 3.06 that got
% dropped from 3.07 after a synchronization disaster. Basically, the
% accentless optional argument to \cs{indexauthor}. Must add the macro
% directories to my unison set-up to stop this sort of thing from happening.}
%    \begin{macrocode}
\newcommand\@isortkey{}
\newcommand\@emptykey{empty}
\newcommand\indexauthor[2][empty]{
  \renewcommand\@isortkey{#1}
  \ifx\@isortkey\@emptykey \renewcommand\@isortkey{#2}\fi
  \index{\@isortkey @{{\rmaaIndexAuthorStyle{#2}}\hskip1em plus1em minus0.4em
      \rmaaIndexTitle\hskip1em plus2em minus0.4em\penalty-100
      \rmaaTOCauthors}}
  \indexauthorADS{#2}
}
%%\def\indexauthor@deferlist{\relax}
%%\long\def\indexauthor#1{
%%  \def\@@name{#1} \toks0=\expandafter{\@@name}
%%  \toks1=\expandafter{\indexauthor@deferlist}
%%  \protected@xdef\indexauthor@deferlist{%
%%    \the\toks1
%%    \noexpand\index{\the\toks0@{\noexpand\rmaaIndexAuthorStyle{\the\toks0}\noexpand\hskip1em plus1em minus0.4em 
%%        \rmaaIndexTitle\noexpand\hskip1em plus2em minus0.4em\noexpand\penalty-100
%%        \rmaaTOCauthors}}
%%  }
%%  \indexauthorADS{#1}
%%  }
%%\xdef\indexauthor@deferlist{}
\newcommand\indexauthorADS[1]{%
  \if@firstauthor\protected@xdef\ADS@authorinitial{\firstletter{#1}}\fi
  \protected@xdef\ADS@authors{\ADS@authors\if@firstauthor\else ; \fi
    #1}
  \@firstauthorfalse
}
%    \end{macrocode}
% Note that the first |#1| sets the key for the index, making sure it
% actually starts with  
% first letter of the surname. Meanwhile, the |@{...}| expression,
% which starts with some  
% control sequences is what is written to the |.ind| file. 
% See the makeindex man page for details. 
% \end{macro}
% \end{macro}
% \end{macro}
% \end{macro}

% \begin{macro}{\KillAllAccents}
% \begin{macro}{\firstletter}
% These are utility functions used in the above. |lett| is a counter
% indicating the current letter in the string. |\firstletter| just
% prints the first letter of its argument.
%    \begin{macrocode}
\newcommand{\KillAllAccents}{\def\'{}\def\`{}\def\~{}\def\.{}\def\i{i}}
\newcommand\firstletter[1]{\@car #1\@nil}
%    \end{macrocode}
% \end{macro}
% \end{macro}
% \begin{macro}{\SetVolume}
% \begin{macro}{\SetMSnumber}
% \begin{macro}{\SetFirstPage}
% \begin{macro}{\SetYear}
% \begin{macro}{\ReceivedDate}
% \begin{macro}{\AcceptedDate}
% These are for the editors to set the volume number, etc. 
%    \begin{macrocode}
\newcommand{\SetVolume}[1]{%
  \renewcommand{\@volume}{#1}
  \edef\ADS@volume{\PadToFourWithDots{#1}}
  \renewcommand{\@volumen}{\rmaa@VolumenStyle{volumen #1}}%
  }
\newcommand{\SetMSnumber}[1]{\renewcommand{\@MSnum}{#1}}
%    \end{macrocode}
% Note that there is no real reason to use |\SetFirstPage| since all
% the page numbering is done automatically in the book style. 
%    \begin{macrocode}
\newcommand{\SetFirstPage}[1]{\setcounter{page}{#1}}
\newcommand{\SetYear}[1]{%
  \renewcommand{\@year}{#1}
  \renewcommand{\ADS@year}{#1}
  \renewcommand{\@anyo}{\rmaa@VolumenStyle{#1}}
  }
\newcommand{\ReceivedDate}[1]{\renewcommand{\@received}{#1}}
\newcommand{\AcceptedDate}[1]{\renewcommand{\@accepted}{#1}}
%    \end{macrocode}
%
% \begin{macro}{\DOI}
% User/editor command to enter the DOI, which will be typeset as a link.
% \changes{4.0}{2019/11/02}{New user command \cs{DOI} to typeset
% DOI in the header of first page and add it to ADS info.}
%    \begin{macrocode}
\newcommand\DOI[1]{%
  \def\@DOI@url{\url{https://doi.org/10.22201/ia.01851101p.#1}}%
  \def\ADS@doi{https://doi.org/10.22201/ia.01851101p.#1}%
}%
%    \end{macrocode}
% \end{macro}
%
% \begin{macro}{\ORCIDiD}
% \changes{4.0}{2019/11/12}{Add mechanism for giving ORCID iD links}
% \begin{macro}{\ORCIDiDfull}
% \begin{macro}{\@orcid@icon}
% \begin{macro}{\@orcid@iconbig}
% Commands to add icon links to ORCID iDs of authors. For example
% \begin{quote}
% |\ORCIDiD|\marg{0000-0000-0000-0000}
% \end{quote}
% to be used after each author name in the |\author|
% command and in the |\fulladdresses|.  Alternatively,
% |\OORCIDiDfull| will typeset the icon plus the full text of the
% link, but we don't use that at present.  The icon image is stored in |ORCIDiD_iconvector.pdf|, which needs to be on the input path.  In principle, we could make EPS versions if necessary.
%    \begin{macrocode}
\newcommand\@orcid@icon{\includegraphics[height=1.8ex]{ORCIDiD_iconvector}}
\newcommand\@orcid@iconbig{\includegraphics[height=2ex]{ORCIDiD_iconvector}}
\newcommand\ORCIDiD[1]{\href{https://orcid.org/#1}{\,\@orcid@icon}}
\newcommand\ORCIDiDfull[1]{%
  \href{https://orcid.org/#1}{\@orcid@iconbig}\,\url{https://orcid.org/#1}}
%    \end{macrocode}
% \end{macro}
% \end{macro}
% \end{macro}
% \end{macro}
%
% \begin{macro}{\@volume}
% \begin{macro}{\@volumen}
% \begin{macro}{\@artpages}
% \begin{macro}{\@year}
% \begin{macro}{\@anyo}
% \begin{macro}{\@received}
% \begin{macro}{\@accepted}
% \begin{macro}{\@MSnum}
% Default values for the containers used by above commands, plus a
% few more useful ones. 
%    \begin{macrocode}
\newcommand{\@volume}{00}
\newcommand{\@volumen}{\rmaa@VolumenStyle{00}}
\newcommand{\@artpages}{%
  \oldpageref{\articlename:RMAAFirstPage}--\oldpageref{\articlename:RMAALastPage}}
\newcommand{\@year}{\the\year}
\newcommand{\@anyo}{\rmaa@VolumenStyle{2007}}
\newcommand{\@journalname}{\rmaa@JournalNameStyle{%
    \if@proceedings\RMAASClong\else\RMAAlong\fi}}
\newcommand{\@received}{}
\newcommand{\@accepted}{}
\newcommand{\@MSnum}{0000}
\newcommand{\@DOI@prefix}{https://doi.org/10.22201/ia.01851101p}
\def\@DOI@url{\url{https://doi.org/10.22201/ia.01851101p.2099.00.00.00}}
%    \end{macrocode}
% \end{macro}
% \end{macro}
% \end{macro}
% \end{macro}
% \end{macro}
% \end{macro}
% \end{macro}
% \end{macro}
% \end{macro}
% \end{macro}
% \end{macro}
% \end{macro}
% \end{macro}
% \end{macro}
% 
% The following commands are most useful in the \texttt{book}
% style. However, the |\SetConfTitle| command can now be used in
% individual articles to set the header for the \texttt{preprint}
% style. 
%    \begin{macrocode}
\newcommand{\@reunionplace}{Lugar de la reuni\'on}
\newcommand{\@reuniondate}{Fecha de la reuni\'on}
\newcommand{\@conftitle}{Conference Title}
\newcommand{\SetReunion}[2]{\renewcommand{\@reunionplace}{#1}\renewcommand{\@reuniondate}{#2}}
\newcommand{\SetConfTitle}[1]{\renewcommand{\@conftitle}{#1}}
\newcommand{\@editors}{Juan Ram\'on de la Fuente and ``El Mosh''}
\newcommand{\SetEditors}[1]{\renewcommand{\@editors}{#1}}
\newcommand{\@editores}{Juan Ram\'on de la Fuente y ``El Mosh''}
\newcommand{\SetEditores}[1]{\renewcommand{\@editores}{#1}}
\newcommand{\@mes}{mes de la publicaci\'on}
\newcommand{\SetMes}[1]{\renewcommand{\@mes}{\rmaa@VolumenStyle{#1}}}
%    \end{macrocode}

% \subsection{Notes for Author Proofs and Communication between Editors}
% \label{sec:notes}
% We introduce a command \cs{Note} that should be used to describe
% any change made to the manuscript. These are a bit tricky to get
% right. In general we use \cs{marginpar} but this only works in outer
% paragraph mode, so we just bung the box where we are if we are in a
% figure caption or something. If we are in the preamble, then we just
% use a \marginpar as part of the \maketitle command.
%    \begin{macrocode}
\newif\if@arenotes
\@arenotesfalse
\newcounter{notecounter}
\setcounter{notecounter}{0}
\newcommand\NotesList{}
\newcommand\Note[1]{%
  \if@proof{%
    \@arenotestrue%
    \stepcounter{notecounter}%
    \if@inpreamble\else{%
      \NoteMark{\arabic{notecounter}}%
    }\fi
    \protected@xdef\NotesList{\NotesList 
      \protect\item #1}%
  }\fi%
}
\newcommand\Changed[2]{Changed ``#1'' to ``#2''.}
\newcommand\NoteMark[1]{%
  \ifinner
  \raisebox{1ex}[0pt][0pt]{\makebox[0pt][l]{\framebox{\footnotesize\bfseries
        N#1}}}%
  \else
  \marginpar{\framebox{\footnotesize\bfseries N#1}}%
  \fi
}
\newcommand\RefNote[1]{\if@proof\NoteMark{\ref{#1}}\relax\fi}
\newcommand\NoteLabel{\protect\label}
\newcommand\OutputNotes{%
  \bgroup
  \onecolumn
  \section*{EDITORIAL NOTES}
  \begin{enumerate}
    \NotesList
  \end{enumerate}
  \egroup
}
\DeclareRobustCommand\cs[1]{\texttt{\char`\\#1}}

%    \end{macrocode}
% \subsection{Bibliographic Records for ADS}
% \label{sec:ADS}
% \changes{2.2beta11}{2001/11/27}{Just about finished all the ADS
% stuff now. Except for the references, which I have given up on for
% the time being.}
% \begin{macro}{\ADS@BibCode}
% \begin{macro}{\ADS@year}
% \begin{macro}{\ADS@journal}
% \begin{macro}{\ADS@volume}
% \begin{macro}{\ADS@firstpage}
% \begin{macro}{\ADS@authorinitial}
% We use the |\ADS@BibCode| command to hold the bibliographic code for
% ADS. It is 
% used in the |%R| 
% field of the bibliographic record for each article. 
% The format is 
% \begin{quote}
% \texttt{YYYYJJJJJVVVVMPPPPA}  
% \end{quote}
% where \texttt{YYYY} is the year, \texttt{JJJJJ} is the journal code,
% which is \texttt{RMxAA} or \texttt{RMxAC}, \texttt{VVVV} is the
% volume number (right-justified), \texttt{M} is not used by us,
% \texttt{PPPP} is the first page number (right-justified), and
% \texttt{A} is the initial of the first author's surname. The volume
% and page fields are padded with periods (`\texttt{.}'), which causes
% us some fun (this is done with the macro |\PadToFourWithDots| described
% below) . 
% 
% All these individual bits are held in the commands |\ADS@year|,
% |\ADS@journal|, |\ADS@volume|, |\ADS@firstpage|,
% |\ADS@authorinitial|. 
%    \begin{macrocode}
\newcommand{\SetADSBibCode}[1]{\renewcommand{\ADS@BibCode}{#1}}
\newcommand{\ADS@BibCode}
{\ADS@year\ADS@journal\ADS@volume.\PadToFourWithDots{\ADS@firstpage}\ADS@authorinitial}
\newcommand{\ADS@year}{2001} 
\newcommand{\ADS@journal}{\if@proceedings RMxAC\else RMxAA\fi}
\newcommand{\ADS@volume}{....} 
%    \end{macrocode}
% These two are just set to default values here, |\ADS@firstpage| is set
% in the |document| environment for included articles, in the
% |frontsection| environment, |\tableofcontents| command,  |index| and
% |abstracts| environments. Also by |\AtBeginDocument| for stand-alone
% articles. 
% 
% God knows how we are going to do |\ADS@authorinitial|!
%    \begin{macrocode}
\newcommand{\ADS@firstpage}{....}
\newcommand{\ADS@authorinitial}{A}
%    \end{macrocode}
% \end{macro}
% \end{macro}
% \end{macro}
% \end{macro}
% \end{macro}
% \end{macro}
% We also have a whole slew of other fields that can be set. These are
% all listed in Table~\ref{tab:ADS}. Most of them still have to be
% defined. 
% \begin{table}[!ht]\raggedleft
% \caption{ADS Bibliographic Records}\label{tab:ADS}
% \begin{tabular}{lp{8em}lp{8em}}\hline
% Code & Description & Macro Name & Comment \\ \hline
% |%R| & Bibliographic Code & |\ADS@BibCode| & required \\
% |%A| & Author list & |\ADS@authors| & required, generate from |\indexauthor| \\
% |%F| & Author affiliation & |\ADS@affil| & format: AA(\dots)
% AB(\dots) \dots \\
% |%J| & Journal Name & |\ADS@journalname| &  \\
% |%D| & Publication Date & |\ADS@fulldate| &  \\
% |%L| & Last Page & |\ADS@lastpage| & generate from label \\
% |%T| & Title & |\ADS@title| & required, generate from |\title| \\
% |%C| & Abstract Copyright & |\ADS@copyright| & set to IA-UNAM \\
% |%U| & URL for electronic document & |\ADS@url| & set to PDF file \\
% |%d| & DOI for article & |\ADS@doi| & set by |\DOI| \\
% |%K| & Keywords & |\ADS@keywords| & set equal to |\@keywords| \\
% |%A| & Abstract Text & |\ADS@abstract| & set equal to |\@abstract| \\
% |%R| & References & |\ADS@references| & generate from the
% bibliography  \\ \hline
% \end{tabular}
% \end{table}
% A lot of these macros are set by other commands, as mentioned in the
% Table, but some need to be set  explicitly here. 
% 
% \begin{macro}{\ADS@doi}
% \changes{4.0}{2019/11/08}{Save DOI for ADS}
% The user/editor command \cs{DOI} is used to set this. 
%    \begin{macrocode}
\newcommand{\ADS@doi}{https://doi.org/10.22201/ia.01851101p.0000.00.00.00}
%    \end{macrocode}
% \end{macro}
% 
% \begin{macro}{\ADS@journalname}
% This
% does not contain any accents at the moment so as not to cause
% problems---this maybe should be changed though. 
%    \begin{macrocode}
\newcommand{\ADS@journalname}{%
  \if@proceedings\@conftitle\space (Eds. \@editors) \fi%
  Revista Mexicana de Astronom\'\i{}a y Astrof\'\i{}sica%
  \if@proceedings\space (Serie de Conferencias)\fi%
  \space Vol. \@volume, pp. \ADS@firstpage--\ADS@lastpage\space (\@year)%
  \space (http://www.astroscu.unam.mx/\~{}rmaa/)
  }
%    \end{macrocode}
% \end{macro}
% \begin{macro}{\ADS@keywords}
% \begin{macro}{\ADS@abstract}
% \begin{macro}{\ADS@title}
% |\ADS@abstract| is straightforward, since we already have the text
% saved as |\@abstract|. We could do the same with |\ADS@keywords| but
% to avoid unwanted formatting we explicitly set it in |\keywords|
% (p.~\pageref{keywords}). |\ADS@title| is just set to |\@rawtitle|,
% which we were already using for another purpose. 
%    \begin{macrocode}
\let\ADS@title=\@rawtitle
\newcommand{\ADS@keywords}{No keywords specified}
\newcommand{\ADS@abstract}{No abstract specified}
%    \end{macrocode}
% \end{macro}\end{macro}\end{macro}
% \begin{macro}{\ADS@authors}
% |\ADS@authors| wants to be in ``Other, A. N.'' format, but with
% semi-colons between authors. The most natural way to do this seems to
% be with |\indexauthor| (p.~\pageref{indexauthor}). We append each
% author to this macro, which means we must explicitly clear it for
% each new article (done in |\article|). .   
%    \begin{macrocode}
\newcommand{\ADS@authors}{}
%    \end{macrocode}
% \end{macro}
% \begin{macro}{\ADS@fulldate}
% \begin{macro}{\ADS@lastpage}
% \begin{macro}{\ADS@copyright}
% \begin{macro}{\ADS@website}
% \begin{macro}{\ADS@url}
% \begin{macro}{\SetWebsite}
% \begin{macro}{\SetMonthNumeric}
% Other sundry data, together with user commands for bits that can be
% set by the Editors. The |\ADS@url| is constructed from the website
% (currently |http://www.astrosmo.unam.mx/rmaa/|,
% the journal descriptor (e.g., |RMxAC..10/|), followed by the
% path to the PDF file (e.g., |PDF/RMxAC..10_HEW1.pdf|). 
% 
%    \begin{macrocode}
\newcommand{\ADS@website}{http://www.astrosmo.unam.mx/rmaa/}
\newcommand{\ADS@url}
{\ADS@website\ADS@journal\ADS@volume\ADS@number/PDF/\ADS@journal\ADS@volume\ADS@number\_\articlename.pdf}
\newcommand{\ADS@number}{}
\newcommand{\ADS@fulldate}{00/\ADS@year} 
\newcommand{\ADS@lastpage}{00} 
\newcommand{\ADS@copyright}{(C) \ADS@year; Instituto de Astronom\'\i{}a,
  Universidad Nacional Aut\'onoma de M\'exico} 
\newcommand{\SetNumber}[1]{\renewcommand{\ADS@number}{-#1}}
\newcommand{\SetWebsite}[1]{\renewcommand{\ADS@website}{#1}}
\newcommand{\SetMonthNumeric}[1]{\renewcommand{\ADS@fulldate}{#1/\ADS@year}}
%    \end{macrocode}
% \end{macro}\end{macro}\end{macro}\end{macro}\end{macro}\end{macro}\end{macro}
% \begin{macro}{\ADS@affil}
% \begin{macro}{\ADS@references}
% These two are somewhat more problematical. Affiliation currently
% only works for the first author (see \dots) pending redesign of the
% user interface. References are now done by wrapping the
% |thebibliography| environment to open another input stream on the
% same file and search through for the |\bibitems|. 
%    \begin{macrocode}
\newcommand{\ADS@references}{}
\newcommand{\ADS@affil}{}
%    \end{macrocode}
% \end{macro}\end{macro}
% 
% Define two macros that expand to the strings
% |\begin{thebibliography}| and |\end{thebibliography}|. Note that the
% idea for this and much else is inspired by the code for the docstrip
% package. Also define |\StartBibX| in case authors use the more
% standard |\begin{thebibliography}{}| syntax. 
%    \begin{macrocode}
\edef\EndBib{\expandafter\@gobble\string\\end\string{thebibliography\string}}
\edef\StartBib{\expandafter\@gobble\string\\begin\string{thebibliography\string}}
\edef\StartBibX{\expandafter\@gobble\string\\begin\string{thebibliography\string}\string{\string}}
%    \end{macrocode}
% Define two new logical variables. |\ifAnotherline| is true if we
% want to read another line from the input stream. It is set to false
% when we encounter the |\end{thebibliography}|. |\ifInBib| is set to
% true when we frind the |\begin{thebibliography}| and tells us to
% start saving the following lines. 
%    \begin{macrocode}
\newif\ifAnotherLine
\newif\ifInBib
%    \end{macrocode}
% \begin{macro}{makeOther}
% We need a macro to change the catcodes of lots of things to other. 
%    \begin{macrocode}
\def\makeOther#1{\catcode`#1=12\relax}
%    \end{macrocode}
% \end{macro}
% \begin{macro}{\RereadBiblio}
% Define a macro to do all the fiddly stuff. This is used at the
% beginning of |thebibliography| environment (Sec~\ref{sec:biblio}).
%    \begin{macrocode}
\newread\tf@biblio
\newcommand{\RereadBiblio}{%
%    \end{macrocode}
% Open the current file again for reading
%    \begin{macrocode}
  \begingroup
  \immediate\openin\tf@biblio=\articlename
%    \end{macrocode}
% Everything in a group to be safe. Initialize logicals. Black magic
% to change
% catcodes and avoid blank lines being interpreted as |\par|. 
%    \begin{macrocode}
  \AnotherLinetrue\InBibfalse
  \makeOther\ \makeOther\\\makeOther\$%
  \makeOther\#\makeOther\^\makeOther\^^K%
  \makeOther\_\makeOther\^^A\makeOther\%%
  \makeOther\~\makeOther\{\makeOther\}\makeOther\&%
  \endlinechar-1\relax
%    \end{macrocode}
% Start looping through the file. Read each line to the macro
% |\Line|. If we have got to EOF or end of biblio then set
% |\AnotherLinefalse|.
%    \begin{macrocode}
  \loop
  \read\tf@biblio to \Line
  \ifx\Line\EndBib\AnotherLinefalse\InBibfalse\fi
  \ifeof\tf@biblio \AnotherLinefalse\fi
%    \end{macrocode}
% If we are in the biblio and the line is not blank then append it to
% the |\ADS@references| macro together with a newline. 
%    \begin{macrocode}
  \ifInBib
  \ifx\Line\empty\else\xdef\ADS@references{\ADS@references\Line^^J}\fi
  \fi
%    \end{macrocode}
% Check if the line we just saw was the start of biblio, then repeat
% the loop if necessary. 
%    \begin{macrocode}
  \ifx\Line\StartBib\InBibtrue\fi
  \ifx\Line\StartBibX\InBibtrue\fi
  \ifAnotherLine\repeat
  \immediate\closein\tf@biblio
  \endgroup
  }
%    \end{macrocode}
% \end{macro}
% 
% \begin{macro}{\PadToFourWithDots}
% Utility function for left-padding a number to 4 characters with
% periods. Used in 
% constructing |\ADS@BibCode|. This is a very simple algorithm and
% \emph{requires} that the argument actually be a number---if it is
% not then dire consequences will result.  
%    \begin{macrocode}
\newcommand{\PadToFourWithDots}[1]{%
  \ifnum #1 > 999 #1%
  \else\ifnum #1 > 99 .#1%
  \else\ifnum #1 > 9 ..#1%
  \else ...#1%
  \fi\fi\fi}
%%  \ifnum \count0 > 999 #1%
%%  \else\ifnum \count0 > 99 .#1%
%%  \else\ifnum \count0 > 9 ..#1%
%%  \else ...#1%
%%  \fi\fi\fi}
%    \end{macrocode}
% \end{macro}
%
% \begin{macro}{\AuthorADS}
% \begin{macro}{\ADS@affiliations}
% New mechanism for setting the affiliations for ADS. The author labels (AA, AB, \dots) are generated using two nested counters.  We use the |\robustify| command from \texttt{etoolbox} so that accented characters in the affiliation addresses do not blow up.  The list of accents may need to be expanded. 
%    \begin{macrocode}
\newcounter{AUTHOR}
\newcounter{author}[AUTHOR]
\stepcounter{AUTHOR}
\newcommand\ADS@affiliations{}
\robustify{\'}\robustify{\i}\robustify{\~}\robustify{\`}
\def\AuthorADS#1#2{%
  \begingroup
  \def\ORCIDiD##1{<ID system="ORCID">##1</ID>}
  \def\Email##1{<EMAIL>##1</EMAIL>}
  \def\Affil##1{##1}
  \stepcounter{author}
  \ifthenelse{\value{author} > 26}{\stepcounter{AUTHOR}\stepcounter{author}}{}
  \ifthenelse{\equal{\ADS@affiliations}{\@rmaaEmpty}}{%
    \xdef\ADS@affiliations{\Alph{AUTHOR}\Alph{author}(#2)}
  }{%
    \xdef\ADS@affiliations{\ADS@affiliations, \Alph{AUTHOR}\Alph{author}(#2)}
  }%
  \endgroup
}
%    \end{macrocode}
% \end{macro}
% \end{macro}
% %
% \begin{macro}{\ADS@dumpdata}
% \changes{4.0}{2019/11/13}{Write all author affiliations for ADS}
% \begin{macro}{\ADS@write}
% \begin{macro}{\ADS@openfile}
% |\ADS@dumpdata| is the macro that actually does the job of writing out all the
% data to an auxiliary file, using |\ADS@write| to write each
% field. The file is set up with |\ADS@openfile|, which currently is
% done in the TOC of a book.  The new multi-author affiliations are used if |\ADS@affiliations| has been set, otherwise it falls back to the single-author version |\ADS@affil|. 
%    \begin{macrocode}
\newcommand{\ADSfield}[2]{\relax}
\newcommand{\ADSvfield}[2]{\relax}
%%\renewcommand{\ADSfield}[12]{\%#1 #2}
\newcommand{\ADS@dumpdata}
{ 
  \ADS@write{\string\begin{ADSarticle}}
  \ADS@WRITE{\string\ADSfield{R}{\ADS@BibCode}}
  \ADS@write{\string\ADSfield{A}{\ADS@authors}}
  \ifthenelse{\equal{\ADS@affiliations}{\@rmaaEmpty}}{%
    {\ADS@sanitize\ADS@write{\string\ADSfield{F}{AA(\ADS@affil)}}}
  }{%
    {\ADS@sanitize\ADS@write{\string\ADSfield{F}{\ADS@affiliations}}}
  }
  \ADS@write{\string\ADSfield{J}{\ADS@journalname}}
  \ADS@write{\string\ADSfield{D}{\ADS@fulldate}}
  \ADS@write{\string\ADSfield{L}{\ADS@lastpage}}
  {\ADS@sanitize\ADS@write{\string\ADSfield{T}{\ADS@title}}}
  \ADS@write{\string\ADSfield{C}{\ADS@copyright}}
  \ADS@write{\string\ADSfield{d}{\ADS@doi}}
  \ADS@write{\string\ADSfield{U}{\ADS@url}}
  \ADS@write{\string\ADSfield{K}{\ADS@keywords}}
  {\ADS@sanitize\ADS@write{\string\ADSfield{B}{\ADS@abstract}}}
  \ADS@write{\string\ADSfield{Z}{\ADS@references}}
  \ADS@write{\string\end{ADSarticle}}
  }
%%\newcommand{\ADS@write}[1]{\if@book\addtocontents{ads}{#1}\fi}
%%\newcommand{\ADS@openfile}{\@starttoc{ads}}
\newcommand{\ADS@write}[1]{\set@display@protect\immediate\write\tf@ads{#1}}
\newcommand{\ADS@WRITE}[1]{\immediate\write\tf@ads{#1}}
\newcommand{\ADS@openfile}{%
  \newwrite\tf@ads%
  \immediate\openout\tf@ads\jobname.ads\relax
  \ADS@WRITE{\ADS@preamble}}
\newcommand{\ADS@closefile}{\ADS@WRITE{\ADS@postamble}}
\def\ADS@preamble{%
  %%  \string\documentclass{article}
  %%  \string\begin{document}
  }
\def\ADS@postamble{%
  %%  \string\end{document}
  \string%\string% Local Variables:^J 
  \string%\string% mode: LaTeX^J
  \string%\string% mode: Font-lock^J
  \string%\string% End:^J
  }
%    \end{macrocode}
% \end{macro}
% \end{macro}
% \end{macro}
% 
% \subsection{Automatic calculation of the page range}
% \label{sec:page-range}
% Another issue that has cost me a lot of grief. Still, it finally
% seems to be sorted now. All hail Jeff Goldberg. 
%
% Cunning trick to get the page numbers from |\pageref|.  Put a label
% on the first and last page of the document using |\AtBeginDocument|
% and |\AtEndDocument|. .  This way won't work for the book option.
% In that case, we do it by defining |\@endofarticlehook| and
% |\@begofarticlehook| as parallels to |\AtEndDocument| and
% |\AtBeginDocument|.  Note that we have to defer adding
% |\@begofarticlehook| to |\AtBeginDocument| to make sure it comes
% after the redefinitions of |\ref| and |\label|.
%    \begin{macrocode}
\if@book\AtEndDocument{\ADS@closefile}\else
 %\AtBeginDocument{\@begofarticlehook}
\AtEndDocument{\@endofarticlehook}
\fi
%    \end{macrocode}
% \begin{macro}{\@begofarticlehook}
% Since we are getting more and more things that need to be done at
% the beginning of every article, frontsection, etc we define a sister
% hook to 
% |\@endofarticlehook|. This hook does the following: 
% \begin{enumerate}
% \item Labels the first page
% \item Sets |\ADS@firstpage|
% \item runs and then clears |\indexauthor@deferlist|
% \item resets the keyword counter
% \end{enumerate}
% It is executed in the following circumstances: 
% \begin{enumerate}
% \item At |\begin{document}| for standalone files
% \item At |\begin{document}| for each included file
% \item At beginning of front sections 
% \item At beginning of TOC
% \item At beginning of Author Index
% \item At beginning of abstracts
% \end{enumerate}
% For some reason, this is not always within the right group, so we
% use |\xdef| just to be sure. 
%    \begin{macrocode}
\newcommand{\@begofarticlehook}{%
  \if@frontmatter\frontlabel{RMAAFirstPage}\else\label{RMAAFirstPage}\fi%
  \xdef\ADS@firstpage{\@arabic{\c@page}}%
  \if@book\else\makeindex\fi % needed so default index is defined in stand-alone articles
%%  \indexauthor@deferlist
%%  \xdef\indexauthor@deferlist{}
  \@inpreamblefalse
  \setcounter{kwdcounter}{0}
  }
%    \end{macrocode}
% \end{macro}
% \begin{macro}{\@endofarticlehook}
% This is a hook to be run at the end of each article
% currently, it just labels the last page and writes out all the ADS
% data to an auxiliary file.  
% It should not produce any printed output. 
% It is modelled on the |\AtEndDocument| command, to which it reduces
% in the case of a single article. 
%    \begin{macrocode}
\newcommand{\@endofarticlehook}{%
  \if@debug{\ClassWarning{rmaa}{***Doing Hook***}}\fi
  \if@debug{\ClassWarning{rmaa}{ADSabstract is: \ADS@abstract}}\fi
  \clearpage\rmaa@lastpagelabel\if@book\ADS@dumpdata\fi
  \ifx\NotesList\@rmaaEmpty\else\OutputNotes\fi
  \if@debug{\ClassWarning{rmaa}{Notes list is: \NotesList}}\fi
}
%    \end{macrocode}
% \begin{macro}{\rmaa@lastpagelabel}
% New version of last page label command copied from Jeff Goldberg's
% lastpage.sty package - we can't just use |\label| here in case the
% output routine for the last page has already been called. Instead,
% we first explicitly clear the page, step back the page counter,
% explicitly write to the .aux file, then restep the page counter
% again.
%
% But this will conflict with the hyperref package, which does labels differently.  Compensate for that by using the five-argument form.
% 
% Also set the |\ADS@LastPage| macro to the appropriate number. 
%    \begin{macrocode}
\newcommand{\rmaa@lastpagelabel}{%
  \addtocounter{page}{-1}%
  \immediate\write\@auxout{\string
    \newlabel{\articlename:RMAALastPage}{{}{\thepage}{}{}{}}}%
  \edef\ADS@lastpage{\thepage}%
  \addtocounter{page}{1}}
%    \end{macrocode}
% \end{macro}
% \end{macro}

% \subsection{Deal with explicit change in number of columns}
% \label{sec:columns}
% Sometimes, when formatting a book, we need to switch between
% one-column articles and two-column articles. For this to work, we
% need to wrap the standard \LaTeX{} commands to set the |\@rmtwocol|
% flag appropriately and also reset the package default values of all
% scalable lengths. .  
%    \begin{macrocode}
\let\latexonecolumn=\onecolumn
\let\latextwocolumn=\twocolumn
\renewcommand{\onecolumn}{\@rmtwocolfalse\@package@default@lengths\latexonecolumn}
\renewcommand{\twocolumn}{\@rmtwocoltrue\@package@default@lengths\latextwocolumn}
%    \end{macrocode}
% 
% \subsection{The actual formatting of the title}
% \label{sec:title-format}

% Note, we have to redefine the |\maketitle| command to use the
% numbered footnotes. This means we needn't use the atbeginend stuff. 
%    \begin{macrocode}
\renewcommand{\maketitle}{\par
  \begingroup
%    \end{macrocode}
% Don't switch to marked footnotes. 
% This is consistent with existing style but I'm not sure I like it
%    \begin{macrocode}
  %%  \renewcommand\thefootnote{\fnsymbol{footnote}}%
%%  \def\@makefnmark{\hbox to\z@{$\m@th^{\@thefnmark}$\hss}}%
  \def\@makefnmark{\textsuperscript{\@thefnmark}\hspace{-0.2em}}%
  \long\def\@makefntext##1{\parindent 1em\noindent
    \hbox to1.8em{\hss$\m@th^{\@thefnmark}$}##1}%
  \if@rmtwocol
  { \if@debug{\ClassWarning{rmaa}{TWO COLUMN}}\fi
    \twocolumn[\@maketitle]%
    }
  \else
  { \if@debug{\ClassWarning{rmaa}{ONE COLUMN}}\fi
    \newpage
    \global\@topnum\z@   % Prevents figures from going at top of page.
    \if@manuscript \@MSmaketitle \else \@maketitle \fi
    }
  \fi
  \if@manuscript\else\thispagestyle{rmaafirstpage}\@thanks\fi
  \ifx\NotesList\@rmaaEmpty\else{%
    \NoteMark{1--\arabic{notecounter}}\vspace*{-2\baselineskip}%
  }\fi%
  \endgroup
%    \end{macrocode}
% Don't reset footnote counter (not sure whether this is good thing
% or not).
%    \begin{macrocode}
  %%  \setcounter{footnote}{0}%
  \let\thanks\relax\let\affil\relax\let\altaffil\relax
  \let\maketitle\relax\let\@maketitle\relax
  \gdef\@thanks{}\gdef\@author{}\gdef\@title{}}
%    \end{macrocode}

% \begin{macro}{\@kwhwidth}
% \begin{macro}{\@kwtwidth}
% Redefine \cs{@maketitle} to include the abstract, resumen and
% keywords. Note that in the manuscripty style we use
% \cs{@MSmaketitle} instead.  
% \changes{3.16}{2003/04/07}{
%  Added extra vertical space between authors and accepted date.}
% \changes{3.26}{2011/08/15}{Add extra label to review articles}
%    \begin{macrocode}
\newlength{\@kwhwidth}\newlength{\@kwtwidth}
\newcommand\reviewname{Review}
\renewcommand{\@maketitle}{%
  \newpage
  \null
  \if@reviewarticle
  \smallskip
  {\large\bfseries\itshape\reviewname}%
  \bigskip
  \else
  \vskip\@before@title@skip%
  \fi
  \begin{center}%
    {\@title \par}%
    \vskip\@before@author@skip%
    \parbox{\rmaa@AuthorWidthFrac\textwidth}{\centering
      \lineskip\@between@author@skip%
      \@author
      \par}
    %% this should be replaced with `Submitted to...' for preprints
    \if@printdates
    \vskip\@before@dates@skip%
    \if@preprint
    \rmaa@ReceivedStyle{Submitted to RMxAA \@received}
    \else
    \rmaa@ReceivedStyle{Received \@received; accepted \@accepted}
    \fi
    \fi
    \if@manuscript
    \rmaa@ReceivedStyle{Submitted to RMxAA \@received}
    \fi
    \if@guide
    \rmaa@ReceivedStyle{Version \rmaa@version, \rmaa@date}
    \fi
    \vskip\@before@abstract@skip%
    \if@abstractpresent
    \parbox{\rmaa@AbstractWidthFrac\textwidth}{%
      \setlength{\parindent}{0.9cm}%
      \centerline{\rmaa@AbstractHeaderStyle{Abstract}}\par%
      \vskip\@after@abstracthead@skip%
      \if@newstyle\noindent\fi\@abstract
      \vskip\@between@abstract@skip%
      \centerline{\rmaa@AbstractHeaderStyle{Resumen}}\par%
      \vskip\@after@abstracthead@skip%
      \if@newstyle\noindent\fi\@resumen
      \if@setresumen\else\@abstract\fi % append english abstract if
                                       % spanish one not set
      \if@guide\else            %do not use keywords for guide
      \vskip\@before@kwd@skip%
      % Format the Key Words with hanging indentation:
      % They are put in a \parbox with width equal to the line width
      % minus the width of the header. 
      { % Widths of keyword
        \settowidth{\@kwhwidth}{\@keywordheader~}   % header and
        \setlength{\@kwtwidth}{\linewidth}          % keyword text
        \addtolength{\@kwtwidth}{-\@kwhwidth}       %
        \noindent\@keywordheader~\parbox[t]{\@kwtwidth}{%
          \rmaa@KeywordStyle{\boldmath\uppercase{\@keywords}}}}%
      \fi
      }
      \fi
  \end{center}%
  \vspace*{\afterkwdskip}
  \par
  }
\newlength{\afterkwdskip}\setlength{\afterkwdskip}{\z@}
%    \end{macrocode}
% \end{macro}
% \end{macro}
%
%
% \begin{macro}{\reviewarticle}
% \begin{macro}{\reviewtoc}
% \begin{macro}{\reviewlof}
%   User command to specify a review article. Must be used in
%   preamble. Main effect is to place the text ``Review'' before the
%   article title.  Also add command for setting a table of contents
%   and a list of figures at the beginning of the article.
%    \begin{macrocode}
\newcommand{\reviewarticle}{\@reviewarticletrue}
\newcommand\reviewtoc{%
  \section*{CONTENTS}%
  \@starttoc{rtc}%
}
\newcommand\reviewlof{%
  \section*{LIST OF FIGURES}%
  \@starttoc{lof}%
}
%    \end{macrocode}
% \end{macro}
% \end{macro}
% \end{macro}
% \begin{macro}{\l@section}
%    \begin{macrocode}
\renewcommand*\l@section[2]{%
  \ifnum \c@tocdepth >\z@
    \addpenalty\@secpenalty
    \addvspace{0.8em \@plus\p@}%
    \setlength\@tempdima{1.5em}%
    \begingroup
      \parindent \z@ \rightskip \@pnumwidth
      \parfillskip -\@pnumwidth
      \leavevmode %\bfseries
      \advance\leftskip\@tempdima
      \hskip -\leftskip
      #1\nobreak\hfil \nobreak\hb@xt@\@pnumwidth{\hss #2}\par
    \endgroup
  \fi}
%    \end{macrocode}
% \end{macro}
%
% \begin{macro}{\@MSmaketitle}
% A special version of |\@maketitle| for the \texttt{manuscript}
% option. This can deal with a multi-page abstract/resumen. Not very
% well, mind. 
%    \begin{macrocode}
\newcommand{\@MSmaketitle}{%
  \if@debug\ClassInfo{rmaa}{Using special MS maketitle command}\fi
  \newpage
  {
    \thispagestyle{rmaafirstpage}
%     \setlength{\parindent}{0pt}
    \null\vskip 2ex
    \begin{center}
      \@title\par
      \@author\par
      \rmaa@ReceivedStyle{%
        \ifx\@received\@rmaaEmpty
          Draft version: \@date
        \else
          Received: \@received
          \ifx\@accepted\@rmaaEmpty
          \else
            ; Accepted: \@accepted
          \fi
        \fi}\par
    \end{center}
    \@thanks
    \centerline{\rmaa@AbstractHeaderStyle{Resumen}}
    \noindent\@resumen\par\vskip 2ex
    \centerline{\rmaa@AbstractHeaderStyle{Abstract}}
    \noindent\@abstract\par\vskip 2ex
    \noindent\@keywordheader~\rmaa@KeywordStyle{\@keywords}\par
%     \@keywordheader~\rmaa@KeywordStyle{\uppercase{\@keywords}}\par
  }%
  }
%    \end{macrocode}
% \end{macro}

% \subsection{The sectioning commands}
% \label{sec:section-commands}

% \begin{macro}{\section}
% \begin{macro}{\subsection}
% \begin{macro}{\subsubsection}
% \begin{macro}{\paragraph}
% \begin{macro}{\subparagraph}
% Redefine all the standard commands to use these lengths. 
%    \begin{macrocode}
\renewcommand\section{\@startsection {section}{1}{\z@}%
                                   {\abovesectionskip}%
                                   {\belowsectionskip}%
                                   %% how do we get upper case here
                                   {\centering\reset@font\normalsize}}
\renewcommand\subsection{\@startsection{subsection}{2}{\z@}%
                                     {\abovesectionskip}%
                                     {\belowsectionskip}%
                                     {\centering\reset@font\normalsize\itshape}}
\renewcommand\subsubsection{\@startsection{subsubsection}{3}{\z@}%
                                     {\abovesectionskip}%
                                     {\belowsectionskip}%
                                     {\centering\reset@font\normalsize\itshape}}
\renewcommand\paragraph{\@startsection{paragraph}{4}{\z@}%
                                    {\aboveparaskip}%
                                    {\belowparaskip}%
                                    {\reset@font\normalsize\bfseries}}
\renewcommand\subparagraph{\@startsection{subparagraph}{5}{\parindent}%
                                      {\aboveparaskip}%
                                      {\belowparaskip}%
                                      {\reset@font\normalsize\bfseries}}
%    \end{macrocode}
% \end{macro}
% \end{macro}
% \end{macro}
% \end{macro}
% \end{macro}

% \begin{macro}{secnumdepth}
% Make sure that we number the subsubsections---this is done by
% default in the |article| class but not the |book| class, which is
% why we must set it explicitly here.
%    \begin{macrocode}
\setcounter{secnumdepth}{3}
%    \end{macrocode}
% \end{macro}
% \changes{2.2beta10}{2001/11/07}{Fixed \cs{secnumdepth} for \cs{book}
% style.}

% \begin{macro}{\thesubsection}
% \begin{macro}{\thesubsubsection}
% For subsection and subsubsection, we want the numbers to be roman
% but the title to be italic. This seems the easiest way to go about it.
%    \begin{macrocode}
\renewcommand{\thesubsection}{{\upshape\thesection.\@arabic\c@subsection}}
\renewcommand{\thesubsubsection}{\thesubsection{\upshape.\@arabic\c@subsubsection}}
%    \end{macrocode}
% \end{macro}
% \end{macro}

% \begin{macro}{\@seccntformat}
% Add ``.'' to the number here so we can remove it from |\thesection|,
% etc. Reduce gap between number and title WJH 27 Jan 1995
%    \begin{macrocode}
\def\@seccntformat#1{\csname the#1\endcsname .\hskip 0.4em}
%    \end{macrocode}
% \end{macro}

% \begin{macro}{\@sect}
% I didn't really want to mess with the |\@sect| command, but it was
% the only way I could manage to force upper case. Other than the
% commented parts, it is identical to the standard version. 
%    \begin{macrocode}
\def\@sect#1#2#3#4#5#6[#7]#8{
  \ifnum #2>\c@secnumdepth
     \let\@svsec\@empty
  \else
     \refstepcounter{#1}%
     \let\@@protect\protect
     \def\protect{\noexpand\protect\noexpand}%
     \edef\@svsec{\@seccntformat{#1}}%
     \let\protect\@@protect
  \fi
%    \end{macrocode}
% This is the bit I changed (27 Jan 1995). Use upper case for
% |\section|, use mixed case for the rest. 
%    \begin{macrocode}
  \ifnum #2=1                   
     \def\sec@upcase##1{\uppercase{##1}}%
     \def\sec@shape{\upshape}%
  \else                         
     \def\sec@upcase##1{\relax{##1}}%
     \def\sec@shape{\itshape}%
  \fi
  \@tempskipa #5\relax
  \ifdim \@tempskipa>\z@
     \begingroup #6\relax
          \@hangfrom{\hskip #3\relax\@svsec}%
                    {\interlinepenalty \@M \sec@upcase{#8}\par}%
     \endgroup
     \csname #1mark\endcsname{#7}%
%    \end{macrocode}
% Only write TOC entries if doing the Guide style. 
%    \begin{macrocode}
     \if@guide 
        \addcontentsline{toc}{#1}{%
          \ifnum #2>\c@secnumdepth \else
             \protect\numberline{\csname the#1\endcsname}%
          \fi
          #7}%
     \fi
%    \end{macrocode}
% Also allow per-article TOC for review papers
%    \begin{macrocode}
     \if@reviewarticle
        \ifnum #2>2 \else
           \addcontentsline{rtc}{#1}{%
             \protect\numberline{\csname the#1\endcsname}%
             \sec@upcase{\sec@shape #7}}%
        \fi
     \fi
  \else
     \def\@svsechd{#6\hskip #3\relax  %% \relax added 2 May 90
       \@svsec #8\csname #1mark\endcsname{#7}%
       \if@guide
          \addcontentsline{toc}{#1}{%
            \ifnum #2>\c@secnumdepth \else
               \protect\numberline{\csname the#1\endcsname}%
            \fi
            #7}%
       \fi
       }%
  \fi
  \@xsect{#5}}
%    \end{macrocode}
% \end{macro}


% \subsection{Appendices}
% \label{sec:appendices}
% Make sure that appendices in individual articles work properly. We
% copy the definition from the article class (the book one doesn't do
% what we want).  
% \begin{environment}{appendices}
%   This uses an environment |appendices| inside of which each
%   |\section| will become a separate appendix.  Note that the
%   appendices must come before the references (otherwise the full
%   addresses mechanism will not work properly). We make the equation
%   numbering be subordinate to the appendix (equations A12, B6,
%   etc). Finally, we have a medskip followed by APPENDICES to
%   introduce them. \changes{3.17}{2002/04/08}{ Added appendices environment
%   that does the right thing.}
%    \begin{macrocode}
\newenvironment{appendices}{\par
  \begingroup
  \setcounter{section}{0}%
  \setcounter{subsection}{0}%
  \renewcommand\thesection{\@Alph\c@section}
  \renewcommand\theequation{\thesection\arabic{equation}}
  \medskip
  \section*{APPENDICES}}
{\endgroup}
%    \end{macrocode}
% \end{environment}
% \begin{environment}{appendix}
%   If there is only one appendix, then use |appendix| environment
%   instead.  \changes{4.0}{2019/11/12}{ Added
%   appendix environment for case where there is only one section.}
%    \begin{macrocode}
\renewenvironment{appendix}{\par
  \begingroup
  \setcounter{section}{0}%
  \setcounter{subsection}{0}%
  \renewcommand\thesection{\@Alph\c@section}
  \renewcommand\theequation{\thesection\arabic{equation}}
  \medskip
  \section*{APPENDIX}}
{\endgroup}
%    \end{macrocode}
% \end{environment}

% \subsection{Floating thingies}
% \label{sec:floats}   

%% 
%    \begin{macrocode}
\renewcommand{\thetable}{\@arabic\c@table}
\renewcommand{\thefigure}{\@arabic\c@figure}
\renewcommand{\figurename}{Fig.}
\renewcommand{\tablename}{TABLE}
%\renewcommand{\fnum@table}{\MakeUpperCase{\tablename}~\thetable}
%    \end{macrocode}
% For some reason, I have to write `TABLE' explicitly here
% If I use |\tablename|, it gets converted to mixed case---God knows
% why. 
% |\figurename| is also being ignored. 
%    \begin{macrocode}
\renewcommand{\fnum@table}{TABLE~\thetable}
\renewcommand{\fnum@figure}{Fig.~\thefigure}
%    \end{macrocode}

% Do captions in 9 pt face in the main journal style and in new-style
% anything. 
%    \begin{macrocode}
\long\def\@makefigcap#1#2{%
  \if@proceedings\else\small\fi
  \if@newstyle\small\fi 
  \vskip\abovecaptionskip
  \sbox\@tempboxa{#1. #2}%
  \ifdim \wd\@tempboxa >\hsize
    #1. #2 \par
  \else
    \global \@minipagefalse
    \hb@xt@\hsize{\hfil\box\@tempboxa\hfil}%
  \fi
  \vskip\belowcaptionskip}
%    \end{macrocode}

% Make sure table caption is upper case. 
%    \begin{macrocode}
\long\def\@maketabcap#1#2{%
  \hbox to\hsize{\hfil #1 \hfil}
  \vskip 6pt 
  \centering\MakeUpperCase{#2}\par
  \vskip 6pt}
%    \end{macrocode}

%    \begin{macrocode}
\newlength{\@rottabcapwidth}
\long\def\@makerottabcap#1#2{%
  \setbox\@tempboxa\hbox{\MakeUpperCase{#2}}%
  \ifdim\wd\@tempboxa >.8\vsize
     \setlength\@rottabcapwidth{.8\vsize}
  \else
     \setlength\@rottabcapwidth{\wd\@tempboxa}
  \fi
  \rotatebox{90}{%
    \begin{minipage}{\vsize}\centering
      \begin{minipage}{\@rottabcapwidth}\centering
        #1\\ \vskip 6pt \MakeUpperCase{#2}
      \end{minipage}%
    \end{minipage}
    }%
  \hskip 6pt
  }

\newenvironment{Sideways}
    {\begin{sideways}\begin{minipage}{\vsize}\centering}
    {\end{minipage}\end{sideways}}

\newenvironment{Sidewaystable}
{%
  \let\oldcaption=\caption%
  \def\caption{
    \def\@captype{table}\renewcommand{\@makecaption}{\@maketabcap}
    \oldcaption}
  \begin{sidewaystable}\centering}
  {\end{sidewaystable}\let\caption=\oldcaption}
\newenvironment{Sidewaystable*}
{%
  \let\oldcaption=\caption%
  \def\caption{
    \def\@captype{table}\renewcommand{\@makecaption}{\@maketabcap}
    \oldcaption}
  \begin{sidewaystable*}\centering}
  {\end{sidewaystable*}\let\caption=\oldcaption}
%    \end{macrocode}

% Not sure why we need to stick these lines in. 
%    \begin{macrocode}
\renewenvironment{figure}
                 {\renewcommand{\@makecaption}{\@makefigcap}\@float{figure}}
                 {\end@float}
\renewenvironment{figure*}
                 {\renewcommand{\@makecaption}{\@makefigcap}\@dblfloat{figure}}
                 {\end@dblfloat}
\renewenvironment{table}
                 {\renewcommand{\@makecaption}{\@maketabcap}\@float{table}}
                 {\end@float}
\renewenvironment{table*}
                 {\renewcommand{\@makecaption}{\@maketabcap}\@dblfloat{table}}
                 {\end@dblfloat}
%    \end{macrocode}

% Redefine |\@caption| so that we don't bother with the list of
% figures. This saves us having to worry about fragile commands in the
% captions. 
%    \begin{macrocode}
\long\def\@caption#1[#2]#3{%
  \par
%%  \addcontentsline{\csname ext@#1\endcsname}{#1}%
%%    {\protect\numberline{\csname the#1\endcsname}{\ignorespaces #2}}%
  \begingroup
    \@parboxrestore
    \if@minipage
      \@setminipage
    \fi
    \normalsize
    \@makecaption{\csname fnum@#1\endcsname}{\ignorespaces #3}\par
  \endgroup}
%    \end{macrocode}

% Make tables not quite so cramped. Allow |\hline\hline| to give a
% poor man's thick rule. 
%    \begin{macrocode}
\renewcommand{\arraystretch}{1.15}
\setlength{\doublerulesep}{0.4pt}
%    \end{macrocode}

% \begin{macro}{\tablecols}
% \begin{macro}{\@tablecols}
% \begin{macro}{\tabcaptioncontinued}
% |\tablecols| is a user command for setting the number of columns. It
% is used by commands that try and span the entire table (I think
% there must be a better way of acheiving
% this). |\tabcaptioncontinued| is 
% a user command for setting the caption of a continued table on the
% next page. 
%    \begin{macrocode}
\newcommand{\tablecols}[1]{\renewcommand{\@tablecols}{#1}}
\newcommand{\@tablecols}{1}
\newcommand{\tabcaptioncontinued}{%
  \multicolumn{\@tablecols}{c}{\addtocounter{table}{-1}TABLE~\thetable\ (CONTINUED)} \\ \\
  }
%    \end{macrocode}
% \end{macro}
% \end{macro}
% \end{macro}

% \subsubsection{Footnotes in tables}
% \label{sec:table-footnotes}
% \begin{macro}{\tablenotemark}
% \begin{macro}{\tablenote}
% \begin{macro}{\tablenoteskip}
% Old user commands |\tablenotemark{MARK}|,
% |\tablenote{MARK}{TEXT}|---now deprecated. 
%    \begin{macrocode}
\newcommand\tablenotemark[1]{$^\mathrm{#1}$} 
\newcommand\tablenote[2]{\if@proceedings\else\small\fi ${}^\mathrm{#1}$ #2} 
\newcommand\tablenoteskip{\\[0.5\baselineskip]} 
%    \end{macrocode}
% \end{macro}
% \end{macro}
% \end{macro}

% \begin{macro}{\tabnotemark}
% \begin{macro}{\tabnotetext}
% \begin{macro}{\tabnotetextTWO}
% Alternative way of doing footnotes in tables. 
%    \begin{macrocode}
\newcommand{\tabnotemark}[1]{\textsuperscript{\lowercase{#1}}}
%\newcommand{\tabnotetextTWO}[4]{\tabnotetext{#1}{#2}\tabnotetext{#3}{#4}}
\newcommand{\tabnotetextTWO}[4]{%
  \multicolumn{\@tablecols}{l}{%
    \parbox{0.5\tabnotewidth}{\small\textsuperscript{#1}#2}%
    \parbox{0.5\tabnotewidth}{\small\textsuperscript{#3}#4}%
    }%
  \\}
\newcommand{\tabnotetext}[2]{%
  \multicolumn{\@tablecols}{@{}p{\tabnotewidth}@{}}{\small\textsuperscript{#1}#2}\\}
\newlength{\tabnotewidth}
%    \end{macrocode}
% \end{macro}
% \end{macro}
% \end{macro}


% \begin{macro}{changemargin}
%    \begin{macrocode}
% Command to allow margins to be changed on the fly. Useful to allow
% tables and Figures to protrude into the margins. 
%
% \changes{3.24}{2007/09/14}{Added \texttt{changemargin} environment
% (copied from TeX FAQ).}
\newenvironment{changemargin}[2]{% 
  % From TeX FAQ http://www.tex.ac.uk/cgi-bin/texfaq2html?label=chngmargonfly
  \begin{list}{}{%
    \setlength{\topsep}{0pt}%
    \setlength{\leftmargin}{#1}%
    \setlength{\rightmargin}{#2}%
    \setlength{\listparindent}{\parindent}%
    \setlength{\itemindent}{\parindent}%
    \setlength{\parsep}{\parskip}%
  }%
  \item[]}{\end{list}}
%    \end{macrocode}
% \end{macro}

% \subsection{The bibliography and references}
% \label{sec:biblio}
% \begin{macro}{\beforebiblioskip}
% \begin{macro}{\afterbiblioskip}
% I'm not entirely sure why I introduced these two commands, since
% directly putting |\vspace| before and after
% |\begin{thebibliography}| seems to work well enough. 
%    \begin{macrocode}
\newlength{\beforebiblioskip}
\setlength{\beforebiblioskip}{0pt}
\newlength{\afterbiblioskip}
\setlength{\afterbiblioskip}{0pt}
%    \end{macrocode}
% \end{macro}
% \end{macro}
% \begin{macro}{thebibliography}
% The environment itself is done as starred section with extra
% vertical space after heading. This space is currently zero. 
% In the book style we do the cunning trick with rereading the file in
% order to get the references into the |\ADS@references| macro.
% 
% \changes{3.05}{2003/01/02}{Removed the \cs{beforebiblioskip} and
% \cs{afterbiblioskip} commands, since they cause unwanted vertical
% space in some instances.}  
%    \begin{macrocode}
\renewenvironment{thebibliography}{%
  \if@book\RereadBiblio\fi
%%  \vspace*{\beforebiblioskip}
  \section*{REFERENCES} \frenchspacing
%%  \vspace*{\afterbiblioskip}
  \small
%    \end{macrocode} 
% The references themselves are set as a list in order 
% to get the hanging indentation. They are done in a smaller font and
% with french spacing so we needn't worry about the distinction
% between sentence-ending periods and abbreviations. 
%    \begin{macrocode} 
  \list{\null}{\leftmargin 1.5em\labelwidth\z@\labelsep\z@\itemindent -1.5em
    \itemsep\z@\parsep\z@\usecounter{enumi}}
  \sloppy\clubpenalty4000\widowpenalty4000
  \sfcode`\.=1000\relax\@biblisttrue}
{\endlist\@biblistfalse\normalsize\@outputaddressesifnecessary}
%    \end{macrocode}
% \end{macro}
% \subsection{Uniquification of cross-references}
% \label{sec:uniquification}
% 
% \begin{macro}{\articlename}
% In order that the |\label{}| command in individual articles be
% guaranteed to give a unique anchor, we append the article name (held
% in the |\articlename| command) to each key. This requires
% redefinition of the following commands: |\cite|, |\nocite|, |\ref|,
% |\pageref|, |\label|. 
% \begin{macro}{\oldcite}
% \begin{macro}{\oldnocite}
% \begin{macro}{\oldref}
% \begin{macro}{\oldpageref}
% \begin{macro}{\oldlabel}
% The original versions of these commands are
% saved as |\oldcite|, |\oldnocite|, etc. 
%
% The uniquifying of cross-references is now always done, so 
% we don't get multiply-defined labels in the single article case
%    \begin{macrocode}
\newcommand{\articlename}{anon}
\AtBeginDocument{
  \let\oldcite=\cite 
  \let\oldnocite=\nocite 
  \let\oldref=\ref 
  \let\oldpageref=\pageref
  \let\oldlabel=\label
}
%    \end{macrocode}
% \end{macro}
% \end{macro}
% \end{macro}
% \end{macro}
% \end{macro}
% \end{macro}
%
% Uniquify all cross-references by prepending the article name.  This
% now has to go into |\AtBeginDocument| so it is working on the new
% versions that hyperref.sty has set up.  Then we need to put the
% |\@begofarticlehook| afterwards, so that the first page label is using
% the new definitions.  If amsmath is loaded, then we also need to patch
% the special label command that they use. \changes{4.0}{2019/11/18}{Fix labels for amsmath display equations, in case that package has been loaded.}
%    \begin{macrocode}
\AtBeginDocument{
  \renewcommand{\ref}[1]{\oldref{\articlename:#1}%
    \if@debug{\ClassInfo{rmaa}{Reference to \articlename:#1}}\fi}
  \renewcommand{\pageref}[1]{\oldpageref{\articlename:#1}%
    \if@debug{\ClassInfo{rmaa}{Page reference to \articlename:#1}}\fi}
  \renewcommand{\label}[1]{\oldlabel{\articlename:#1}%
    \if@debug{\ClassInfo{rmaa}{Label for \articlename:#1}}\fi}%
  \ifdef{\label@in@display}{%
    \let\oldlabel@in@display=\label@in@display
    \renewcommand{\label@in@display}[1]{\oldlabel@in@display{\articlename:#1}%
      \if@debug{\ClassInfo{rmaa}{AMSTeX Label for \articlename:#1}}\fi}%
  }{}
  \@begofarticlehook
}%
%    \end{macrocode}
% Citations are more complicated since they can contain a list of
% keys. These are really only here for use with the rmaacite package,
% which is utterly deprecated (02 Jan 2003) now that I have the natbib
% code working. The command \cs{cite} should \emph{not} be used with
% natbib: use \cs{citep}, \cs{citet}, etc.\@ instead. 
%    \begin{macrocode}
\renewcommand{\nocite}[1]{\@bsphack
  \@for\@citeb:=#1\do{%
    \edef\@citeb{\expandafter\@firstofone\@citeb:\articlename}%
    \if@filesw\immediate\write\@auxout{\string\citation{\@citeb}}\fi
    \@ifundefined{b@\@citeb:\articlename}{\G@refundefinedtrue
      \@latex@warning{Citation `\@citeb' undefined}}{}}%
  \@esphack}
\renewcommand{\cite}[1]{\oldcite{#1:\articlename}}
%   \end{macrocode}


% Horrible fiddling with the bibliography citation mechanism. 
% This has to be globally available in case rmaacite puts it in the
% \texttt{.aux} file. 
%    \begin{macrocode}
\def\lbibcite#1#2#3{\global\@namedef{b@#1}{#2}
                    \global\@namedef{by@#1}{#3}}
%    \end{macrocode}

% \begin{macro}{\oldbibitem}
% \begin{macro}{\@rmaaBibArg}
% \begin{macro}{\@rmaaSpace}
% \begin{macro}{rmaabib}
% Of course, |\bibitem| needs to get the uniquification treatment as
% well. While we are at it, we redefine it to use consecutive integers
% for the key if the user leaves it blank (using the counter
% |rmaabib|). This allows one to put 
% |\bibitem{}| for all your references without getting annoying
% messages about multiply-defined labels.
% 
% We also take the opportunity to insert the optional argument to
% \cs{bibitem} if it is missing in the form Author(DATE) so that natbib
% doesn't try to do numerical citations. Natbib already provides a
% hook \cs{@extra@b@citeb} to help uniquify the citation key, so we
% just plug into that. However, it puts it \emph{after} the original
% key, not before as I had had it. Easiest solution is to change the
% way that I do it. 02 Jan 2003. Note that the \cs{@extra@b@citeb}
% does not get written to the .aux  file in the \cs{citation} command
% but it does work. 
% \changes{3.14}{06 Feb 2003}{Reinstalled old definition of \cs{bibitem}
%  for the case where the nonatbib option is requested, so that the
%  rmaacite package will work. This means that you can't use rmaacite
%  and natbib in the same volume.}
%    \begin{macrocode}
\if@natbib\def\@extra@b@citeb{:\articlename}\fi
%%\else
\let\oldbibitem\bibitem
\newcommand{\@rmaaBibArg}{}
\newcommand{\@rmaaSpace}{ }
\newcounter{rmaabib}
\if@natbib
\renewcommand{\bibitem}[2][ABC(2003)]{%
  \renewcommand{\@rmaaBibArg}{#2}
  \ifx\@rmaaBibArg\@rmaaEmpty
     \addtocounter{rmaabib}{1}
     \oldbibitem[#1]{\thermaabib:\articlename}
  \else\ifx\@rmaaBibArg\@rmaaSpace
     \addtocounter{rmaabib}{1}
     \oldbibitem[#1]{\thermaabib:\articlename}
  \else 
     \oldbibitem[#1]{#2:\articlename}
  \fi\fi
  \renewcommand{\@rmaaBibArg}{}
  }
\else
\renewcommand{\bibitem}[1]{%
  \renewcommand{\@rmaaBibArg}{#1}
  \ifx\@rmaaBibArg\@rmaaEmpty
     \addtocounter{rmaabib}{1}
     \oldbibitem{\articlename:\thermaabib}
  \else\ifx\@rmaaBibArg\@rmaaSpace
     \addtocounter{rmaabib}{1}
     \oldbibitem{\articlename:\thermaabib}
  \else
     \oldbibitem{\articlename:#1}
  \fi\fi
  \renewcommand{\@rmaaBibArg}{}
  }
\fi
%%\fi
%    \end{macrocode}
% \end{macro}
% \end{macro}
% \end{macro}
% \end{macro}

% \begin{macro}{oldpcite}
% \begin{macro}{oldscite}
% Then, we need to repeat the above scheme, but for the |rmaacite|
% definition. 
%    \begin{macrocode}
\newcommand{\@redefinebibitem}{%
  \let\oldbibitem\bibitem
  \let\oldcite\cite
  \let\oldpcite\pcite
  \let\oldscite\scite
  \def\bibitem[##1<##2>]##3{%
    \renewcommand{\@rmaaBibArg}{##3}
    \ifx\@rmaaBibArg\@rmaaEmpty
       \addtocounter{rmaabib}{1}
       \oldbibitem[##1<##2>]{\thermaabib:\articlename}
    \else\ifx\@rmaaBibArg\@rmaaSpace
       \addtocounter{rmaabib}{1}
       \oldbibitem[##1<##2>]{\thermaabib:\articlename}
    \else 
       \oldbibitem[##1<##2>]{##3:\articlename}
    % trust the author here
%%       \oldbibitem[##1<##2>]{##3}
    \fi\fi
    }
  }
%    \end{macrocode}
% \end{macro}
% \end{macro}


% \subsection{HTML Table of Contents}
% \label{sec:hyper-TOC}
% \begin{macro}{\HyperTOCEntry}
% \begin{macro}{\HyperTOCpart}
% \begin{macro}{\HyperTOCmeeting}
% \begin{macro}{\HyperTOCeditors}
% \begin{macro}{\HyperTOCpreamble}
% These commands are used to write information to the \texttt{.htoc}
% file, which is used in generating a Web Page of the Table of
% Contents. Since they only ever appear preceded by |\protect| in the
% argument to an |\addtocontents| command,  they should never need to
% be expanded by the macros. However, we define them here just in case
% \dots
%    \begin{macrocode}
\newcommand{\HyperTOCEntry}[3]{\relax}
\newcommand{\HyperTOCpart}[1]{\relax}
\newcommand{\HyperTOCmeeting}[3]{\relax}
\newcommand{\HyperTOCeditors}[1]{\relax}
\newcommand{\HyperTOCpreamble}{\relax}
%    \end{macrocode}
% \end{macro}
% \end{macro}
% \end{macro}
% \end{macro}
% \end{macro}
%

% \subsection{Formatting the poster abstracts} 
% These can now be used without using the book option.
%    \begin{macrocode}

\newcommand{\@makeminititle}{%
  \newpage
  \null 
  \vskip 3em%
  \begin{center}%
    {\@title \par}%
  \end{center}
  \vskip 3em%
  }

\newif\if@abstractsads
\@abstractsadsfalse
\newcommand\abstractsADS{\@abstractsadstrue}

\if@book\RequirePackage{multicol}\fi
\newenvironment{abstracts}[1][Abstracts of Contributed Papers]{%
  %% First, make sure we start on an odd page (06 Jun 2000)
  \if@book\clearpage\thispagestyle{empty}\cleardoublepage\fi
  %%
  \renewcommand{\articlename}{ABSTRACTS}
  \title{#1}\shortauthor{Abstracts}\shorttitle{Abstracts}
  \listofauthors{}
  \if@book
  \addcontentsline{toc}{chapter}{\protect\@fiddleauthors{\@TOCtitle}{\rmaaTOCauthors}}%
  \addtocontents{htc}{\protect\HyperTOCEntry{\@rawtitle}{\@rawauthors}{\articlename}}
  \fi
  \pagestyle{rmaa} % this now seems necessary (06 Jun 2000)
  \thispagestyle{rmaafirstpage}
%%  \twocolumn[\@makeminititle]
  \renewcommand{\abstract}[1]{%
    \renewcommand\ADS@abstract{##1}%
    \renewcommand{\@abstract}{##1}\@writeabstract%
  }
  \renewcommand{\title}[2][]{%
    \gdef\ADS@authors{}%
    \gdef\ADS@affil{}%
    \gdef\ADS@references{}%
    \gdef\ADS@keywords{}%
    \renewcommand{\@rmaaArg}{##1}
    \ifx\@rmaaArg\@rmaaEmpty
       \renewcommand{\@TOCaftertitle}{}
    \else
       \renewcommand{\@TOCaftertitle}{ (\@rmaaArg)}
    \fi
    \renewcommand{\@title}{\rmaa@TOCTitleStyle{##2}}%
    \renewcommand{\@TOCtitle}{\rmaa@TOCTitleStyle{##2}\@TOCaftertitle}%
    \renewcommand{\rmaaIndexTitle}{\rmaa@IndexTitleStyle{##2\@TOCaftertitle.}}%
    \renewcommand\ADS@title{##2}%
    \xdef\ADS@firstpage{\@arabic{\c@page}}%
    }
  \renewcommand{\author}[1]{\let\@thanks\@empty\renewcommand\@author{##1} }
  \renewcommand{\footnotetext}[2][XX]{\par\noindent\textsuperscript{##1}~##2}
  \renewcommand{\footnoterule}{\kern-3\p@ \hrule \@width 3em \kern 2.6\p@}
  \if@book\@begofarticlehook\fi
  \if@book\begin{multicols}{2}[\@makeminititle]\else\@makeminititle\fi
  \begingroup
  \setcounter{footnote}{0}
  }% and at the end of the environment....
  {\endgroup\if@book\end{multicols}\listofabstracts\label{RMAALastPage}\clearpage\fi}
%% auxiliary commands used:
%% \@writeabstract does all the actual formatting 
\newcommand{\@writeabstract}{%
  \if@book%
  \addcontentsline{loa}{abstract}{\protect\@fiddleauthors{\@TOCtitle}{\rmaaTOCauthors}}%
  \fi%
  \begin{center}\@title\par\@author\end{center}
  \@abstract
  \par\noindent\footnoterule\nobreak\footnotesize\@thanks
  \par\noindent\hrulefill\goodbreak
  \vskip 1cm \@plus0.5cm
  %% write ADS entry
  \if@book\if@abstractsads\ADS@dumpdata\fi\fi%
  \endgroup % this pair is to make each abstract a self-contained group
  \begingroup
  \setcounter{footnote}{0}
  }

\if@book
%% List of abstracts (stuff written to .loa file)
\newcommand{\listofabstracts}{%
  \begin{multicols}{2}[%
    \vskip 1.5cm \vfill % leave at least 1.5 cm and go to bottom of page if possible
    \centerline{\MakeUpperCase{List of Abstracts}}
    \vskip 1.5cm \@plus0.5cm
    ][6cm] % this makes it start a new page if there is not at least 6cm of space
    \renewcommand{\@fiddleauthors}[2]{##1\hskip1em plus2em minus0.4em\penalty-100  ##2}
    \parindent\z@
%%    \parskip.5\p@ \@plus .7\p@\relax 
    \parskip3\p@ \@plus 1\p@ \@minus 1\p@\relax % allow quite a lot of
                                % leeway here to avoid breaks inside
                                % entries 
    \renewcommand{\\}{\relax} % temporarily turn off explicit line breaks just in case
    \renewcommand{\altaffilmark}[1]{\relax} % ditto with this
    \@starttoc{loa}
  \end{multicols}
  }
\newcommand*{\l@abstract}[2]{%
  \par\penalty-1000
  #1\hskip 1em\penalty0\null\nobreak\hfill#2}
\fi

%    \end{macrocode}
 
% \subsection{Code for typesetting an entire volume of conference
% proceedings} \label{sec:whole-volume}
%    \begin{macrocode}

\if@book
%% environment for including the individual articles
\newcommand{\@fiddleauthors}[2]{\relax}
\renewcommand{\theequation}{\@arabic\c@equation}
\newenvironment{articles}%
{ \begingroup % make it a group so that author definitions do not
              % spill over
  \gdef\ADS@authors{}
  \pagestyle{rmaa}
  \renewenvironment{document}{%
    \@begofarticlehook
    \addtocounter{chapter}{1}%
    \addcontentsline{toc}{chapter}{%
      \protect\@fiddleauthors{\@TOCtitle}{\rmaaTOCauthors}}%
    \addtocontents{htc}{\protect\HyperTOCEntry{\@rawtitle}{\@rawauthors}{\articlename}}
    }{\relax}
  \renewcommand{\documentclass}[2][foo]{\relax} 
  \renewcommand{\usepackage}[2][foo]{%
    \def\@packagename{##2}\def\@rmaacitename{rmaacite}%
    \ifx\@packagename\@rmaacitename
    \@rmcitetrue\ClassInfo{rmaa}{This paper uses rmaacite}
    \makeatletter\input{rmaacite.sty}\makeatother % read it in here so that all changes
    \else                                         % are local to paper
    \if@debug\ClassWarning{rmaa}{Package \@packagename\space was requested by \articlename\space}\fi
    \fi} 
  %% Don't allow individual articles to mess with the pages, volume or year
  \renewcommand{\SetFirstPage}[1]{\relax}
  \renewcommand{\SetYear}[1]{\relax}
  \renewcommand{\SetVolume}[1]{\relax}
  \renewcommand{\SetConfTitle}[1]{\relax}
  %% Don't allow individual articles to change these
  \renewcommand{\nonstopmode}{\relax}
  \renewcommand{\scrollmode}{\relax}
  \renewcommand{\outputplates}{\relax}
  %% XXXX need to check if using rmaacite and act accordingly
  %% avoid multiple declarations of the same length
  \let\oldnewlength=\newlength
  \renewcommand{\newlength}[1]{%
    \edef\reserved@a{\expandafter\@gobble\string ##1}%
    \@ifundefined{\reserved@a}{%
      \oldnewlength{##1}}{%
      \if@debug{\ClassInfo{rmaa}{\reserved@a\space already defined: doing nothing}}\fi%
      \relax}}%
  \renewcommand\newcommand{\providecommand}
  }%
{\endgroup\onecolumn}
     

%    \end{macrocode}
% Command to be used in `articles' environment for including an
% article
%    \begin{macrocode}
\newcommand{\article}[2][]{%
  % optional first argument is extra text to put in brackets after title in TOC
  % e.g. ``Invited Talk'' or ``Poster''
  \renewcommand{\articlename}{#2}
  % reset loads of counters
  \setcounter{footnote}{0}
  \setcounter{figure}{0}
  \setcounter{table}{0}
  \setcounter{section}{0}
  \setcounter{equation}{0}
  \setcounter{rmaabib}{0}
  \setcounter{author}{0}
  \setcounter{AUTHOR}{1}
  \gdef\ADS@authors{}%
  \gdef\ADS@affil{}%
  \gdef\ADS@affiliations{}%
  \gdef\ADS@references{}%
  \@firstauthortrue
  \@firstaffiltrue
  \@firstaltaffiltrue
  \@firstkwdtrue
  \@rmcitefalse
  \if@suppressfulladdresses\else\@needaddressestrue\fi
  \renewcommand{\@rmaaArg}{#1}
  \ifx\@rmaaArg\@rmaaEmpty
     \renewcommand{\@TOCaftertitle}{}
  \else
     \renewcommand{\@TOCaftertitle}{ (#1)}
  \fi
%    \end{macrocode}
% In the main journal articles must start on an odd page. 
%    \begin{macrocode}
  \if@journal\clearpage\thispagestyle{empty}\cleardoublepage\fi
%% Now input the file, then run the end-of-article hook (reinstated 28
  %% Sep 2001)
  {\input{\@articledir/\articlename.tex}\@endofarticlehook}}
%    \end{macrocode}
% Command to separate articles into sessions. Write an entry in the
% TOC, makes sure we start on an odd page (added 06 Jun 2000), then do
% a title page if \cs{SessionTitlePages} was called. 
% \changes{3.12}{04 Feb 2003}{Added option to have title page for each
% session. Turn it on with \cs{SessionTitlePages}.}
%    \begin{macrocode}
\newcommand{\session}[1]{
  %% this does nothing except write an entry in the TOC
  \addcontentsline{toc}{part}{\rmaa@TOCPartStyle{#1}}
  \addtocontents{htc}{\protect\HyperTOCpart{#1}}
  %% and also makes sure we start on an odd page (added 06 Jun 2000)
  \clearpage\thispagestyle{empty}\cleardoublepage
  \stepcounter{part}
  \if@sessiontitlepages{%
    \onecolumn
    \thispagestyle{empty}
    \renewcommand{\articlename}{Section\thepart TitlePage}
    \label{RMAAFirstPage}
    \make@session@title{#1}
    \label{RMAALastPage}
    \clearpage\thispagestyle{empty}\cleardoublepage
  }\fi
  }
\newcommand{\make@session@title}[1]{%
  {%
    \setlength\parindent{0pt}
    \vspace*{\stretch{1}}
    \rule{\linewidth}{1mm}
    \begin{flushright}\Huge\bfseries\boldmath
      #1
    \end{flushright}
    \rule{\linewidth}{1mm}
    \vspace*{\stretch{2}}
  }
}
%    \end{macrocode}
% Command to turn on separate title pages for each session. 
%    \begin{macrocode}
\newcommand{\SessionTitlePages}{\@sessiontitlepagestrue}
\newcommand{\BoldSessionTitles}{%
  \renewcommand{\rmaa@TOCPartStyle}[1]{\textbf{\rmaa@TOCTitleStyle{##1}}}
  }
\newcommand{\ArticleDir}[1]{\renewcommand{\@articledir}{#1/}\graphicspath{{1/}}}
\newcommand{\@articledir}{.}
%\newcommand{\ps@rmaafront}{\ps@rmaa}% page style for front matter (same as ever, actually)
\newcommand{\ps@rmaafront}{% page style for front matter:  page # in each foot
  \renewcommand\@oddfoot{\hfil{\thepage}}%
  \renewcommand\@evenfoot{{\thepage}\hfil}%
  % shorttitle in both heads
  \renewcommand\@evenhead{\hfil\@shorttitle\hfil}%
  \renewcommand\@oddhead{\@evenhead}
  }
\newcommand{\ps@rmaafrontfirstpage}{% page style for first page of front matter
  % page # in each foot, on outside edge
  \renewcommand\@oddfoot{\hfil\thepage}%
  \renewcommand\@evenfoot{\thepage\hfil}%
  % nothing in the even heads
  \renewcommand\@evenhead{\relax}%
  \renewcommand\@oddhead{\relax}%
  }
%% slight modification to the standard \frontmatter and \mainmatter
%% commands, just to set the \pagestyle
\renewcommand\frontmatter{%
  \latexonecolumn
  \cleardoublepage
  \@frontmattertrue
  \@mainmatterfalse
  \pagenumbering{roman}
  \pagestyle{rmaafront}
  \addtocounter{page}{2} % leave space for the title page
  }
\renewcommand\mainmatter{%
  \pagestyle{empty}
  \if@rmtwocol\twocolumn\fi
  \cleardoublepage
  \@frontmatterfalse
  \@mainmattertrue
  \pagenumbering{arabic}
  \pagestyle{rmaa}
  }
\renewcommand\backmatter{%
  \clearpage\pagestyle{empty}\cleardoublepage
  \@mainmatterfalse
  %% put extra space in the TOC
  \addtocontents{toc}{\bigskip}
  }
%    \end{macrocode}
%
% \subsection{Front sections}\label{sec:front-sections}
% \begin{macro}{\notitle}
% \begin{macro}{\shortname}
% \begin{macro}{frontsection}
% The environment \texttt{frontsection} is for doing a section in the
% front matter (preface, list of participants, or whatever). It has 
% one optional argument:
% \begin{enumerate}
% \item[\#1] Short title (no embedded spaces) for
% use in making PS files. E.g., ``list\_of\_participants''.   
% \end{enumerate}
% and two mandatory arguments: 
% \begin{enumerate}
% \item[\#2] Title of section (e.g., `Preface')
% \item[\#3] Commands to be evaluated before processing the contents
% of the environment. This is mainly useful for putting in the things
% that would go in the preamble of a normal article  (e.g.,
% |\listofauthors|, |\indexauthor|, |\notitle|) 
% \end{enumerate}
%    \begin{macrocode}
\newcommand{\notitle}{\@titlevisiblefalse} \newcommand{\shortname}{}
\newenvironment{frontsection}[3][]{%
  \begingroup
  \renewcommand{\shortname}{#1}
  \ifx\shortname\@rmaaEmpty
     \renewcommand{\articlename}{#2}
  \else
     \renewcommand{\articlename}{\shortname}
  \fi
  \@begofarticlehook
  \@titlevisibletrue % default
  \thispagestyle{rmaafrontfirstpage}
  \listofauthors{} % blank out this in case it is not specified
  \title{{#2}}\shorttitle{#2}\shortauthor{#2}
  #3 % evaluate the command given as second mandatory arg
  \addcontentsline{toc}{chapter}{\protect\@fiddleauthors{\@TOCtitle}{\rmaaTOCauthors}}%
  \addtocontents{htc}{\protect\HyperTOCEntry{\@rawtitle}{\@rawauthors}{\articlename}}
  \if@titlevisible
     \vspace*{1cm}\centerline{\rmaa@FrontSecStyle{\@title}}\vspace*{1cm}
  \fi
  \setlength{\parskip}{3ex}
  }{\frontlabel{RMAALastPage}\clearpage\endgroup}
%    \end{macrocode}
% \end{macro}
% \end{macro}
% \end{macro}
%    \begin{macrocode}
%% Definition of the command \frontlabel, which rights out something like `F4' to
%% the .aux file, instead of `\textlatin  {i}'. This is so it is easy to post-process
%% the file for making PS files of individual sections
\newcommand{\frontlabel}[1]{%
  \@bsphack
  \protected@write\@auxout{}%
     {\string\newlabel{\articlename:#1}{{}{F\@arabic{\c@page}}}}
  \@esphack
  }
%% User command to say X, Y, Z,\\ Editors of the Proceedings
\newcommand{\EditorSignature}{\vspace*{\parskip}\rightline{\@editors}%
  \rightline{Editors of the proceedings}}
%% Fiddle with \thesection, etc. so that they don't include the chapter number
%% Basically go back to the `article' definitions of these
\renewcommand{\thesection}{\@arabic\c@section}
\renewcommand{\thefigure}{\@arabic\c@figure}
\renewcommand{\thetable}{\@arabic\c@table}

\newcommand{\@oldargs}{}
\newcommand{\@newargs}{}
%% add TOC entry for each paper
%\renewcommand{\title}[1]%
%{\renewcommand{\@title}{\rmaa@TitleStyle{#1}}%
%\addtocounter{chapter}{1}\addcontentsline{toc}{chapter}{#1}%
%}

%% Formatting of the table of contents
%% dimensions of various parts
\newlength{\TOCtextwidth}\setlength{\TOCtextwidth}{12.65cm}
\newlength{\TOCpagenumwidth}\setlength{\TOCpagenumwidth}{1.3cm}
\newlength{\TOCentrywidth}\setlength{\TOCentrywidth}{\TOCtextwidth}
\addtolength{\TOCentrywidth}{-\TOCpagenumwidth} 
\newlength{\TOCextramargin}\setlength{\TOCextramargin}{1.25cm}
\newcommand{\TOCbox}[1]{\hspace{\TOCextramargin}\parbox{\TOCtextwidth}{#1}\par}
\newlength{\TOCextramarginx}\setlength{\TOCextramarginx}{\TOCextramargin}\addtolength{\TOCextramarginx}{4.3cm}
\newlength{\TOCtextwidthx}\setlength{\TOCtextwidthx}{\TOCtextwidth}\addtolength{\TOCtextwidthx}{-4.3cm}
\newcommand{\TOCboxx}[1]{\hspace{\TOCextramarginx}\parbox{\TOCtextwidthx}{#1}\par}

\newcommand{\RMAAbannerimage}{rmsc}
\renewcommand\tableofcontents{%
  \renewcommand{\articlename}{CONTENTS}
  \begingroup
  \@begofarticlehook
  \if@twocolumn
     \@restonecoltrue\onecolumn
  \else
     \@restonecolfalse
  \fi
  %% 9pt for journal style
  \if@proceedings\else\small\fi
  \if@newstyle\small\fi %% same for new style
  \if@proceedings\thispagestyle{rmaafrontfirstpage}\else\thispagestyle{empty}\fi
  \shortauthor{Contents}\shorttitle{Contents}
  \vspace*{18mm}
  \if@printbanner
  %% put in the banner header (need to get better image of this)
     \TOCbox{\hspace*{-0.15cm}\includegraphics[width=0.9\TOCtextwidth]{\RMAAbannerimage}}
     \vspace*{8mm}
  %% volume and publication date (should really be in funny font)
     \TOCboxx{\@volumen\\[2mm] \@mes\ \@anyo}
  \else
  %% put in the banner header (need to get better image of this)
     \TOCbox{\hspace*{-0.15cm}\includegraphics[width=0.9\TOCtextwidth]{\RMAAbannerimage}}
     \vspace*{12.5mm}
  %% volume and publication date (should really be in funny font)
     \TOCboxx{~\\[2mm] ~}
  \fi
  \vspace*{16mm}
  %% This is the command that fiddles with the TOC entry so that the
  %% title is flushleft but the authors are flushright
  \renewcommand{\@fiddleauthors}[2]{%
    % first put the paper title
    \renewcommand{\\}{\relax} % temporarily turn off explicit line breaks just in case
    \renewcommand{\altaffilmark}[1]{\relax} % ditto with this
    \renewcommand{\textsuperscript}[1]{\relax} % ditto with this 
    ##1
    % now save the author list in a box
    \sbox\@tempboxa{##2}%
    % check that box is shorter than the total width available for the entry
    \ifdim \wd\@tempboxa <\TOCentrywidth
       % if so, then put it right here, flush against right margin
       \penalty10000\hskip1em plus10em\penalty0\hbox{}\nobreak\hfill\box\@tempboxa
       % Note that we force a space of at least 1em between title and authorlist, 
       % stretchable to 11em since we may have very short titles but still with an 
       % authorlist too long to go on the same line.  The only linebreak allowed is 
       % directly before the authorlist.
    \else
       % if it is too long, we set the original authorlist text in its own parbox
       \par\vskip1ex\parbox[b]{\TOCentrywidth}{\raggedleft ##2}
    \fi
    }
  \@starttoc{toc}%
  \if@restonecol\twocolumn\fi
  \frontlabel{RMAALastPage}
  \clearpage
  \thispagestyle{empty}
  \cleardoublepage % make sure we start on an odd page
  \endgroup

  %%% ALSO, do the HTML TOC
  %% The file that will be used to make the HTML TOC
  \@starttoc{htc}
  %% All the stuff in this file does nothing in the book style
  \addtocontents{htc}{\protect\HyperTOCmeeting{\@conftitle}{\@reunionplace}{\@reuniondate}}
  \addtocontents{htc}{\protect\HyperTOCeditors{\@editors}}
  \addtocontents{htc}{\protect\HyperTOCpreamble}
  \addtocontents{htc}{\protect\HyperTOCEntry{Table of Contents}{}{\articlename}}

  %% and the ADS file
  \ADS@openfile
  }
%% This is now much simpler
\renewcommand\l@chapter[2]{%TOC entry for articles
  \ifnum \c@tocdepth >\m@ne
    \vskip 1.0em \@plus\p@
    \TOCbox{%
      \parbox[b]{\TOCentrywidth}{#1}\parbox[b]{\TOCpagenumwidth}{\hfill#2}
    }
  \fi} 
\renewcommand\l@part[2]{%TOC entry for sessions: centered, no page number
  \ifnum \c@tocdepth >\m@ne
    \vskip 1.5em \@plus1.5em \goodbreak
    \TOCbox{%
      \parbox[b]{\TOCentrywidth}{\hfil#1\hfil}\hfill
    }
  \fi\nobreak} 

%% Formatting theindex
%% strongly encourage page breaks between index entries
\renewcommand{\@idxitem}{\par\penalty-1000\hangindent.45cm} 
\renewcommand{\indexspace}{\par\penalty-1000\vskip 10\p@ \@plus15\p@ \@minus3\p@\relax}
\renewenvironment{theindex}{%
  %% first, make sure we start on an odd page (06 Jun 2000)
  \clearpage\thispagestyle{empty}\cleardoublepage
  %% this looks a bit like an article but is two columns and without abstract, etc
  \renewcommand{\articlename}{INDEX}
  \title{Author Index}\shortauthor{Author Index}\shorttitle{Author Index}
  \listofauthors{}
  \addcontentsline{toc}{chapter}{\protect\@fiddleauthors{\@TOCtitle}{\rmaaTOCauthors}}%
  \addtocontents{htc}{\protect\HyperTOCEntry{\@rawtitle}{\@rawauthors}{\articlename}}
  \pagestyle{rmaa}
  \thispagestyle{rmaafirstpage}
  \@begofarticlehook
  %% 9pt for journal style
  \if@proceedings\else\small\fi
  \if@newstyle\small\fi
  \twocolumn[\@makeminititle]
  %% vertical spacing (see also definition of \@idxitem and \indexspace
  \parskip\z@ \@plus .6\p@\relax % allow quite a lot of leaway here to avoid breaks inside entries
  %% horizontal spacing, breaking and glue parameters (see also the stuff in the definition 
  %% of \indexauthor
  \parindent\z@ 
  \exhyphenpenalty10000 % don't break on explicit hyphens, it always looks bad
  \rightskip=0pt plus.2em % allow a tiny bit of leeway at the ends of lines
  \spaceskip.3333em plus.5em minus.15em % a lot of leeway between words
  \xspaceskip.5em plus.5em minus.4em % this seems to have no effect
  \let\item\@idxitem
  \renewcommand{\\}{\relax} % temporarily turn off explicit line breaks just in case
  \renewcommand{\altaffilmark}[1]{\relax} % ditto with this
  \renewcommand{\textsuperscript}[1]{\relax} % ditto with this
  }{\label{RMAALastPage}\clearpage} % clearpage necessary so that we know about \shortauthor



%% end definitions for book styles
\fi

%    \end{macrocode}
% In order to do away with plates, we sometimes want a figure to
% appear in color only in the online version. The following command
% can bve used in a figure caption to typeset a note to this effect. 
%    \begin{macrocode}
\newcommand\ColorOnlineOnly{%
  \textsc{Note: this figure is available in color in the electronic
    version of this article, obtainable from}
  \texttt{http://\discretionary{}{}{}www.\discretionary{}{}{}astroscu.\discretionary{}{}{}unam.\discretionary{}{}{}mx/\discretionary{}{}{}\~{}rmaa/}.}
%% Plates
%% Define 2 counters: plate is the number in the whole document, plateloop is used 
%% internally by \outputplates
\newcounter{plate}
\newcounter{plateloop}
%% Currently limited to 15 plates in document
%% when we do a label, we need to write to \newlabel entries to the .aux file, one 
%% that says RMAAPlateX so that it can be referred to in \ouputplates, and one that
%% says articlename:RMAAPlateX so that the dvips program can find it and asociate it
%% with the right article
\newcommand{\@platelabeli}{\oldlabel{RMAAPlate1}\label{RMAAPlate1}}
\newcommand{\@platelabelii}{\oldlabel{RMAAPlate2}\label{RMAAPlate2}}
\newcommand{\@platelabeliii}{\oldlabel{RMAAPlate3}\label{RMAAPlate3}}
\newcommand{\@platelabeliv}{\oldlabel{RMAAPlate4}\label{RMAAPlate4}}
\newcommand{\@platelabelv}{\oldlabel{RMAAPlate5}\label{RMAAPlate5}}
\newcommand{\@platelabelvi}{\oldlabel{RMAAPlate6}\label{RMAAPlate6}}
\newcommand{\@platelabelvii}{\oldlabel{RMAAPlate7}\label{RMAAPlate7}}
\newcommand{\@platelabelviii}{\oldlabel{RMAAPlate8}\label{RMAAPlate8}}
\newcommand{\@platelabelix}{\oldlabel{RMAAPlate9}\label{RMAAPlate9}}
\newcommand{\@platelabelx}{\oldlabel{RMAAPlate10}\label{RMAAPlate10}}
\newcommand{\@platelabelxi}{\oldlabel{RMAAPlate11}\label{RMAAPlate11}}
\newcommand{\@platelabelxii}{\oldlabel{RMAAPlate12}\label{RMAAPlate12}}
\newcommand{\@platelabelxiii}{\oldlabel{RMAAPlate13}\label{RMAAPlate13}}
\newcommand{\@platelabelxiv}{\oldlabel{RMAAPlate14}\label{RMAAPlate14}}
\newcommand{\@platelabelxv}{\oldlabel{RMAAPlate15}\label{RMAAPlate15}}
\newcommand{\@outputplatelabeli}{\oldlabel{RMAAOutputPlate1}}
\newcommand{\@outputplatelabelii}{\oldlabel{RMAAOutputPlate2}}
\newcommand{\@outputplatelabeliii}{\oldlabel{RMAAOutputPlate3}}
\newcommand{\@outputplatelabeliv}{\oldlabel{RMAAOutputPlate4}}
\newcommand{\@outputplatelabelv}{\oldlabel{RMAAOutputPlate5}}
\newcommand{\@outputplatelabelvi}{\oldlabel{RMAAOutputPlate6}}
\newcommand{\@outputplatelabelvii}{\oldlabel{RMAAOutputPlate7}}
\newcommand{\@outputplatelabelviii}{\oldlabel{RMAAOutputPlate8}}
\newcommand{\@outputplatelabelix}{\oldlabel{RMAAOutputPlate9}}
\newcommand{\@outputplatelabelx}{\oldlabel{RMAAOutputPlate10}}
\newcommand{\@outputplatelabelxi}{\oldlabel{RMAAOutputPlate11}}
\newcommand{\@outputplatelabelxii}{\oldlabel{RMAAOutputPlate12}}
\newcommand{\@outputplatelabelxiii}{\oldlabel{RMAAOutputPlate13}}
\newcommand{\@outputplatelabelxiv}{\oldlabel{RMAAOutputPlate14}}
\newcommand{\@outputplatelabelxv}{\oldlabel{RMAAOutputPlate15}}
\newcommand{\@platerefi}{\oldpageref{RMAAPlate1}}
\newcommand{\@platerefii}{\oldpageref{RMAAPlate2}}
\newcommand{\@platerefiii}{\oldpageref{RMAAPlate3}}
\newcommand{\@platerefiv}{\oldpageref{RMAAPlate4}}
\newcommand{\@platerefv}{\oldpageref{RMAAPlate5}}
\newcommand{\@platerefvi}{\oldpageref{RMAAPlate6}}
\newcommand{\@platerefvii}{\oldpageref{RMAAPlate7}}
\newcommand{\@platerefviii}{\oldpageref{RMAAPlate8}}
\newcommand{\@platerefix}{\oldpageref{RMAAPlate9}}
\newcommand{\@platerefx}{\oldpageref{RMAAPlate10}}
\newcommand{\@platerefxi}{\oldpageref{RMAAPlate11}}
\newcommand{\@platerefxii}{\oldpageref{RMAAPlate12}}
\newcommand{\@platerefxiii}{\oldpageref{RMAAPlate13}}
\newcommand{\@platerefxiv}{\oldpageref{RMAAPlate14}}
\newcommand{\@platerefxv}{\oldpageref{RMAAPlate15}}
\newcounter{figplatei}
\newcounter{figplateii}
\newcounter{figplateiii}
\newcounter{figplateiv}
\newcounter{figplatev}
\newcounter{figplatevi}
\newcounter{figplatevii}
\newcounter{figplateviii}
\newcounter{figplateix}
\newcounter{figplatex}
\newcounter{figplatexi}
\newcounter{figplatexii}
\newcounter{figplatexiii}
\newcounter{figplatexiv}
\newcounter{figplatexv}
%% Printed output should use the counter `plate'
\renewcommand{\theplate}{\@arabic\c@plate}
\renewcommand{\theplateloop}{\@arabic\c@plateloop}
%% Command to output all the saved plates in this paper
\newcommand{\outputplates}{%
  \setcounter{plateloop}{0}
  \if@book\else\@endofarticlehook\fi % not entirely sure why, but this
                                     % seems to be needed
  \clearpage % first output the last page of text
  % now set pagestyle for blank even pages and make sure we start on an odd page
  \pagestyle{empty}
  \cleardoublepage
  \onecolumn
  \parindent\z@
  % first put in an extra sheet
  % the main purpose of this is so there is always a blank even page before 
  % each plate for filling in the PS file if the article has an odd number
  % of pages (see splitbook.sh)
  \null\vfil\centerline{PLATES}\vfil\clearpage\pagestyle{rmaaplate}\cleardoublepage
  % output each plate in turn
  \@whilenum\value{plateloop}<\value{plate}\do{%
    \addtocounter{plateloop}{1}
    % restore the figure counter from \@figcounti, \@figcountii, etc
    \setcounter{figure}{\value{figplate\roman{plateloop}}}
    % spew out the plate (in its own group for safety)
    {\csname @plate\roman{plateloop}\endcsname\par}
    }
  \pagestyle{rmaa}
  % reset the page number as if nothing had happened
  \addtocounter{page}{-\value{plate}}
  \addtocounter{page}{-\value{plate}}
  }
\newcommand\figcaption{%
  \def\@captype{figure}
  \renewcommand{\@makecaption}{\@makefigcap}
  \caption}
\newcommand\tabcaption{%
  \def\@captype{table}
  \renewcommand{\@makecaption}{\@maketabcap}
  \caption}

\newcommand{\plate}[2][1]{%
  %% User command to make a plate page: 
  %% Mandatory argument is what goes on the plate - this should not 
  %% include any floating environments (e.g. figure). Instead, \figcaption 
  %% should be used to do the captions. 
  %% Optional argument is the number of figures (i.e. captions) on the plate
  %% page (default is 1). This is necessary for incrementing the figure counter. 
  %%
  %% Note that we have to use a \gdef rather than an \xdef when we 
  %% save the plate output in a command. This means that in the book styles when it 
  %% comes to spewing out the plate, we will have forgotten about any user macros 
  %% that may have been defined. This can be got round by repeating the macros inside the plate.
  %% (Presumably, one would use a \newcommand rather than \renewcommand if they were new.)
  % increment global and local counters
  \addtocounter{plate}{1}
  % write appropriate label to .aux file
  \csname @platelabel\roman{plate}\endcsname
  % define a global macro called \@platei, \@plateii, \@plateiii, \@plateiv, ...
  \expandafter\gdef\csname @plate\roman{plate}\endcsname{%
    \csname @outputplatelabel\roman{plateloop}\endcsname
    \null\vfil #2 \vfil
    \cleardoublepage
    }
  % save the current shortauthor and shorttitle
  % need to be \xdef to take current values of \@shortauthor, etc 
  \expandafter\xdef\csname @shorta\roman{plate}\endcsname{\@shortauthor}
  \expandafter\xdef\csname @shortt\roman{plate}\endcsname{\@shorttitle}
  % save the current value of the figure counter for later use in \outputplates
  \setcounter{figplate\roman{plate}}{\value{figure}}
  % increment the figure counter according to the optional argument
  \addtocounter{figure}{#1}
  }
%% page style for plate pages 
\newcommand{\ps@rmaaplate}{%
  % Plate X (see page Y) in the foot
  \renewcommand\@oddfoot{\hfil\csname @shorta\roman{plateloop}\endcsname\ 
    (see page~\csname @plateref\roman{plateloop}\endcsname)\hfil}%
  \renewcommand\@evenfoot{\relax}%
  % nothing in even heads (blank pages)
  \renewcommand\@evenhead{\relax}
  % short title and page num in the odd heads
  \renewcommand\@oddhead{\hfil\csname @shortt\roman{plateloop}\endcsname%
    \hfil\llap{PLATE~\theplateloop}}
  }





%% Things for the User Guide
\if@guide
%% stolen from doc.sty
\def\MakeShortVerb#1{%
  \expandafter\ifx\csname cc\string#1\endcsname\relax
  \@shortvrbinfo{Made }{#1}%
  \add@special{#1}%
  \expandafter
  \xdef\csname cc\string#1\endcsname{\the\catcode`#1}%
  \begingroup
  \catcode`\~\active  \lccode`\~`#1%
  \lowercase{%
    \global\expandafter\let
    \csname ac\string#1\endcsname~%
    \gdef~{\verb~}}%
  \endgroup
  \global\catcode`#1\active
  \else
  \@shortvrbinfo\@empty{#1 already}%
  \fi}
\def\DeleteShortVerb#1{%
  \expandafter\ifx\csname cc\string#1\endcsname\relax
  \else
  \@shortvrbinfo{Deleted }{#1 as}%
  \rem@special{#1}%
  \global\catcode`#1\csname cc\string#1\endcsname
  \global \expandafter\let \csname cc\string#1\endcsname \relax
  \ifnum\catcode`#1=\active
  \begingroup
  \catcode`\~\active   \lccode`\~`#1%
  \lowercase{%
    \global\expandafter\let\expandafter~%
    \csname ac\string#1\endcsname}%
  \endgroup \fi \fi}
\def\@shortvrbinfo#1#2{%
  \PackageInfo{doc}{%
    #1\expandafter\@gobble\string#2 a short reference
    for \string\verb}}
\def\add@special#1{%
  \rem@special{#1}%
  \expandafter\gdef\expandafter\dospecials\expandafter
  {\dospecials \do #1}%
  \expandafter\gdef\expandafter\@sanitize\expandafter
  {\@sanitize \@makeother #1}}
\def\rem@special#1{%
  \def\do##1{%
    \ifnum`#1=`##1 \else \noexpand\do\noexpand##1\fi}%
  \xdef\dospecials{\dospecials}%
  \begingroup
  \def\@makeother##1{%
    \ifnum`#1=`##1 \else \noexpand\@makeother\noexpand##1\fi}%
  \xdef\@sanitize{\@sanitize}%
  \endgroup}
%% define | as shortcut for \verbatim
\MakeShortVerb{\|}
%% comand to specify arguments to LaTeX commands
\newcommand{\rmArg}[1]{%
  \ensuremath{\langle}{\scriptsize \textsf{\MakeUpperCase{#1}}}\ensuremath{\rangle}}
%\newcommand{\rmArg}[1]{\ensuremath{\langle}\textsc{#1}\ensuremath{\rangle}}
\providecommand{\BibTeX}{{\rmfamily B\kern-.05em%
    \textsc{i\kern-.025em b}\kern-.08em%
    T\kern-.1667em\lower.7ex\hbox{E}\kern-.125emX}}
%% Environment for code examples
\newenvironment{Code}{%
  \begin{list}{}{\setlength{\leftmargin}{0.05\columnwidth}} 
  \item[]}{\end{list}}
\fi
%</package>
%    \end{macrocode}

% \section{User commands for astronomical articles}
% \label{sec:user-commands}
% Some random definitions. These go both into the main class file and
% also into |adshtml.sty|.
% 
% Definitions of the journal name in different incarnations.  
%    \begin{macrocode}
%<*package|adshtml-style>
\newcommand{\RMAA}{RMxAA}
\newcommand{\RMAAlong}{Revista Mexicana de Astronom\'{\i}a y Astrof\'{\i}sica}
\newcommand{\RMAASC}{RMxAC}
\newcommand{\RMAASClong}{\RMAA\ (Serie de Conferencias)}
%    \end{macrocode}
% \begin{macro}{\ibidrule}
% This puts a horizontal rule with a period after it. It is used when
% the author list of a reference is identical to the preceding one. 
%    \begin{macrocode}
\newcommand{\ibidrule}{\rule[-0.03cm]{4.7em}{.01cm}.\ }
%    \end{macrocode}
% \end{macro}
% \begin{macro}{\ThisVolume}
% References to other articles in the same volume can be coded as,
% e.g., |\bibitem{} Bloggs, J. \ThisVolume{bloggs}|.
%    \begin{macrocode}
\newcommand{\ThisVolume}[1]{%
  \@year, \if@proceedings\RMAASC\else\RMAA\fi, \@volume,
  \if@book\oldpageref{#1:RMAAFirstPage}\else 000\fi\ (this volume)}
%    \end{macrocode}
% \end{macro}
% Other stuff that may be useful in ref lists.
%    \begin{macrocode}
\newcommand{\etal}{et al.\@}
%    \end{macrocode}
% Various new commands from AAS macros
%    \begin{macrocode}
\def\acknowledgments{\vskip 3ex plus .8ex minus .4ex}
\let\acknowledgements=\acknowledgments                  % second spelling
\let\jnl@style=\rm
\def\ref@jnl#1{{\jnl@style#1}}
\DeclareRobustCommand\aj{\ref@jnl{AJ}}                   
\DeclareRobustCommand\araa{\ref@jnl{ARA\&A}}             
\DeclareRobustCommand\apj{\ref@jnl{ApJ}}                 
\DeclareRobustCommand\apjl{\ref@jnl{ApJ}}                
\DeclareRobustCommand\apjs{\ref@jnl{ApJS}}               
\DeclareRobustCommand\ao{\ref@jnl{Appl.~Opt.}}           
\DeclareRobustCommand\apss{\ref@jnl{Ap\&SS}}             
\DeclareRobustCommand\aap{\ref@jnl{A\&A}}                
\DeclareRobustCommand\aapr{\ref@jnl{A\&A~Rev.}}          
\DeclareRobustCommand\aaps{\ref@jnl{A\&AS}}              
\DeclareRobustCommand\azh{\ref@jnl{AZh}}                 
\DeclareRobustCommand\baas{\ref@jnl{BAAS}}               
\DeclareRobustCommand\jrasc{\ref@jnl{JRASC}}             
\DeclareRobustCommand\memras{\ref@jnl{MmRAS}}            
\DeclareRobustCommand\mnras{\ref@jnl{MNRAS}}             
\DeclareRobustCommand\pra{\ref@jnl{Phys.~Rev.~A}}        
\DeclareRobustCommand\prb{\ref@jnl{Phys.~Rev.~B}}        
\DeclareRobustCommand\prc{\ref@jnl{Phys.~Rev.~C}}        
\DeclareRobustCommand\prd{\ref@jnl{Phys.~Rev.~D}}        
\DeclareRobustCommand\pre{\ref@jnl{Phys.~Rev.~E}}        
\DeclareRobustCommand\prl{\ref@jnl{Phys.~Rev.~Lett.}}    
\DeclareRobustCommand\pasp{\ref@jnl{PASP}}               
\DeclareRobustCommand\pasj{\ref@jnl{PASJ}}               
\DeclareRobustCommand\qjras{\ref@jnl{QJRAS}}             
\DeclareRobustCommand\skytel{\ref@jnl{S\&T}}             
\DeclareRobustCommand\solphys{\ref@jnl{Sol.~Phys.}}      
\DeclareRobustCommand\sovast{\ref@jnl{Soviet~Ast.}}      
\DeclareRobustCommand\ssr{\ref@jnl{Space~Sci.~Rev.}}     
\DeclareRobustCommand\zap{\ref@jnl{ZAp}}                 
\DeclareRobustCommand\nat{\ref@jnl{Nature}}              
\DeclareRobustCommand\iaucirc{\ref@jnl{IAU~Circ.}}
\DeclareRobustCommand\aplett{\ref@jnl{Astrophys.~Lett.}}
\DeclareRobustCommand\apspr{\ref@jnl{Astrophys.~Space~Phys.~Res.}}
\DeclareRobustCommand\bain{\ref@jnl{Bull.~Astron.~Inst.~Netherlands}}
\DeclareRobustCommand\fcp{\ref@jnl{Fund.~Cosmic~Phys.}}
\DeclareRobustCommand\gca{\ref@jnl{Geochim.~Cosmochim.~Acta}}
\DeclareRobustCommand\grl{\ref@jnl{Geophys.~Res.~Lett.}}
\DeclareRobustCommand\jcp{\ref@jnl{J.~Chem.~Phys.}}      
\DeclareRobustCommand\jgr{\ref@jnl{J.~Geophys.~Res.}}    
\DeclareRobustCommand\jqsrt{\ref@jnl{J.~Quant.~Spec.~Radiat.~Transf.}}
\DeclareRobustCommand\memsai{\ref@jnl{Mem.~Soc.~Astron.~Italiana}}
\DeclareRobustCommand\nphysa{\ref@jnl{Nucl.~Phys.~A}}
\DeclareRobustCommand\physrep{\ref@jnl{Phys.~Rep.}}
\DeclareRobustCommand\physscr{\ref@jnl{Phys.~Scr}}
\DeclareRobustCommand\planss{\ref@jnl{Planet.~Space~Sci.}}
\DeclareRobustCommand\procspie{\ref@jnl{Proc.~SPIE}}
\let\astap=\aap
\let\apjlett=\apjl
\let\apjsupp=\apjs
\let\applopt=\ao
\def\phn{\phantom{0}}
\def\phd{\phantom{.}}
\def\phs{\phantom{\ensuremath{-}}}
\def\phm#1{\phantom{#1}}
%    \end{macrocode}
% All these are redone as robust and using more modern commands where
% possible. The robustness is so they will not get expanded if used in
% abstracts and titles when these are passed to the ADS file. 
% \changes{3.16}{2003/04/07}{
%  Added \cs{Color} command for doing photometric colors.}
%    \begin{macrocode}
%%\def\sun{\hbox{$\odot$}}
\DeclareRobustCommand\sun{\ensuremath{\odot}}
%%\def\earth{\hbox{$\oplus$}}
\DeclareRobustCommand\earth{\ensuremath{\oplus}}
%%\def\lesssim{\mathrel{\hbox{\rlap{\hbox{\lower4pt\hbox{$\sim$}}}\hbox{$<$}}}}
%%\def\gtrsim{\mathrel{\hbox{\rlap{\hbox{\lower4pt\hbox{$\sim$}}}\hbox{$>$}}}}
\DeclareRobustCommand\lesssim{\mathrel{\hbox{\smaller\smaller\rlap{\hbox{\lower3pt\hbox{$\sim$}}}\raise1.5pt\hbox{$<$}}}}
\DeclareRobustCommand\gtrsim{\mathrel{\hbox{\smaller\smaller\rlap{\hbox{\lower3pt\hbox{$\sim$}}}\raise1.5pt\hbox{$>$}}}}
%%\def\lta{\mathrel{\spose{\lower 3pt\hbox{$\mathchar"218$}}
%%     \raise 2.0pt\hbox{$\mathchar"13C$}}}
%%\def\gta{\mathrel{\spose{\lower 3pt\hbox{$\mathchar"218$}}
%%     \raise 2.0pt\hbox{$\mathchar"13E$}}}
\DeclareRobustCommand\sq{\hbox{\rlap{$\sqcap$}$\sqcup$}}
%%\def\arcdeg{\hbox{$^\circ$}}
%%\def\arcmin{\hbox{$^\prime$}}
%%\def\arcsec{\hbox{$^{\prime\prime}$}}
\DeclareRobustCommand\arcdeg{\ensuremath{^\circ}}
\DeclareRobustCommand\arcmin{\ensuremath{^\prime}}
\DeclareRobustCommand\arcsec{\ensuremath{^{\prime\prime}}}
%%\def\fd{\hbox{$.\!\!^{\rm d}$}}
%%\def\fh{\hbox{$.\!\!^{\rm h}$}}
%%\def\fm{\hbox{$.\!\!^{\rm m}$}}
%%\def\fs{\hbox{$.\!\!^{\rm s}$}}
%%\def\fdg{\hbox{$.\!\!^\circ$}}
%%\def\farcm{\hbox{$.\mkern-4mu^\prime$}}
%%\def\farcs{\hbox{$.\!\!^{\prime\prime}$}}
%%\def\fp{\hbox{$.\!\!^{\scriptscriptstyle\rm p}$}}
\DeclareRobustCommand\fd{\ensuremath{.\!\!^\mathrm{d}}}
\DeclareRobustCommand\fh{\ensuremath{.\!\!^\mathrm{h}}}
\DeclareRobustCommand\fm{\ensuremath{.\!\!^\mathrm{m}}}
\DeclareRobustCommand\fs{\ensuremath{\mkern0.5mu.\mkern-5mu^\mathrm{s}\mkern-1mu}}
\DeclareRobustCommand\fdg{\ensuremath{.\!\!^\circ}}
\DeclareRobustCommand\farcm{\ensuremath{.\mkern-4mu^\prime}}
\DeclareRobustCommand\farcs{\ensuremath{\mkern1mu.\mkern-6mu^{\prime\prime}\mkern-1.5mu}}
%%\DeclareRobustCommand\farcs{\ensuremath{.\!\!^{\prime\prime}}}
\DeclareRobustCommand\fp{\ensuremath{.\!\!^{\scriptscriptstyle\rm p}}}
%%\def\micron{\hbox{$\mu$m}}
\DeclareRobustCommand\micron{\ensuremath{\mu\mathrm{m}}}
\let\la=\lesssim                        
\let\ga=\gtrsim
%%\def\case#1#2{\hbox{$\frac{#1}{#2}$}}
%%\def\slantfrac#1#2{\hbox{$\,^#1\!/_#2$}}
%%\def\onehalf{\slantfrac{1}{2}}
%%\def\onethird{\slantfrac{1}{3}}
%%\def\twothirds{\slantfrac{2}{3}}
%%\def\onequarter{\slantfrac{1}{4}}
%%\def\threequarters{\slantfrac{3}{4}}
\DeclareRobustCommand\case[2]{\ensuremath{\frac{#1}{#2}}}
\DeclareRobustCommand\slantfrac[2]{\ensuremath{\,^#1\!/_#2}}
\DeclareRobustCommand\onehalf{\slantfrac{1}{2}}
\DeclareRobustCommand\onethird{\slantfrac{1}{3}}
\DeclareRobustCommand\twothirds{\slantfrac{2}{3}}
\DeclareRobustCommand\onequarter{\slantfrac{1}{4}}
\DeclareRobustCommand\threequarters{\slantfrac{3}{4}}
%%\def\ubvr{\hbox{$U\!BV\!R$}}            
%%\def\ub{\hbox{$U\!-\!B$}}               
%%\def\bv{\hbox{$B\!-\!V$}}               
%%\def\vr{\hbox{$V\!-\!R$}}               
%%\def\ur{\hbox{$U\!-\!R$}}
%    \end{macrocode}
% New commands for ``photometric'' colors 
%    \begin{macrocode}
\DeclareRobustCommand\ubvr{\ensuremath{U\!BV\!R}}            
\DeclareRobustCommand\Color[2]{%
  \ensuremath{%
    #1\!\ifx#1V\!\fi-\!#2
  }}
\DeclareRobustCommand\ub{\Color{U}{B}}
\DeclareRobustCommand\bv{\Color{B}{V}}
\DeclareRobustCommand\vr{\Color{V}{R}}
\DeclareRobustCommand\ur{\Color{U}{R}}
%% New version of the \ion command, that actually works! 
%    \end{macrocode}
% \begin{macro}{\ion}
% This version of the |\ion| command allows you to write either
% |\ion{N}{4}| or |\ion{N}{iv}| in order to get N~IV. 
%    \begin{macrocode}
\DeclareRobustCommand\ion[2]{%
  \mbox{#1\kern0.2em%
    \smaller\rmfamily%
    \edef\@tempa{\@car#2\@nil}%
    \ifcat1\@tempa%
    \@Roman{#2}%
    \else%
    \uppercase{#2}%
    \fi}}
\newcounter{IonStage}
\DeclareRobustCommand\plainion[2]{\setcounter{IonStage}{#2}#1 \Roman{IonStage}}
%%\def\ion#1#2{\ensuremath{\mathrm{#1}\;}{\protect\small\rm\@Roman{#2}}}
%    \end{macrocode}
% \end{macro}
%    \begin{macrocode}
\DeclareRobustCommand\nodata{$\cdots$}
%    \end{macrocode}
% These are needed by the |\ion| command. Isn't there an easier way of
% doing this?
%    \begin{macrocode}
\let\@small=\small
\let\@normalsize=\normalsize
\let\@footnotesize=\footnotesize
\let\@scriptsize=\scriptsize
\let\@tiny=\tiny
\let\@large=\large
\let\@Large=\Large
\let\@LARGE=\LARGE
\let\@huge=\huge
\let\@Huge=\Huge
\newcommand{\smaller}{\small}
\renewcommand{\normalsize}{\renewcommand{\smaller}{\small}\@normalsize}
\renewcommand{\small}{\renewcommand{\smaller}{\footnotesize}\@small}
\renewcommand{\footnotesize}{\renewcommand{\smaller}{\scriptsize}\@footnotesize}
\renewcommand{\scriptsize}{\renewcommand{\smaller}{\tiny}\@scriptsize}
\renewcommand{\tiny}{\renewcommand{\smaller}{\tiny}\@tiny}
\renewcommand{\large}{\renewcommand{\smaller}{\normalsize}\@large}  
\renewcommand{\Large}{\renewcommand{\smaller}{\large}\@Large}  
\renewcommand{\LARGE}{\renewcommand{\smaller}{\Large}\@LARGE}  
\renewcommand{\huge}{\renewcommand{\smaller}{\LARGE}\@huge}  
\renewcommand{\Huge}{\renewcommand{\smaller}{\huge}\@Huge}  
%</package|adshtml-style>
%    \end{macrocode}


% \section{Other files}

% \subsection{The driver and make files}
% \label{sec:driver_make}
% These are put here for completeness, although they still need to be
% distributed separately since otherwise the user won't be able to
% unpack this file into its constituent parts.
% 
% The |Makefile| has two principal rules: |docs|, and |macros|. The
% |docs| rule produces a PDF and gzipped PS file of the documentation.
% I still haven't perfected the combination of rules required to get
% it to run \LaTeX{} enough times in all circumstances. The |macros|
% rule runs \LaTeX{} on |rmaa.ins| in order to unpack |rmaa.cls| and
% all the other little files (including the |Makefile| and the
% |rmaa.ins| files themselves\dots 
% 
% First set up lists of various types of file. 
%    \begin{macrocode}
%<*makefile>
DOCFILES=rmaa.pdf
MACRO_LTX_FILES=rmaa.cls rmaacite.sty adshtml.sty
SHELLSCRIPTS=makeads.sh
MACRO_OTHER_FILES=adshtml.perl sample.latex2html-init
MAKEFILE=Makefile
MACROFILES=$(MACRO_LTX_FILES) $(SHELLSCRIPTS) $(MACRO_OTHER_FILES)         
SKELETONFILES=ADS.tex 
%    \end{macrocode}
% Install the files read only into the package directory, specified via a path relative to the source directory. \changes{4.0}{2019/11/14}{Reorganize folder layout so that hardly anything is in the top level.  That way, things are clearer on github.}
%    \begin{macrocode}
INSTALL=install 
INSTALL_DIR=../packages
%    \end{macrocode}
% The default rule is to make everything. Installation must be
% explicitly asked for.
%    \begin{macrocode}
all: docs macros

install: install-docs install-macros

%    \end{macrocode}
% Most files are installed read-only to prevent editing. Skeleton
% files must be writeable. Must be careful not to overwrite the
% Makefile in the parent directory. 
%    \begin{macrocode}
install-macros: macros
        $(INSTALL) -m 444 $(MACRO_LTX_FILES) $(MACRO_OTHER_FILES) $(INSTALL_DIR)
        $(INSTALL) -m 555 $(SHELLSCRIPTS) $(INSTALL_DIR)
        $(INSTALL) -m 644 $(SKELETONFILES) $(INSTALL_DIR)

install-docs: docs
        $(INSTALL) rmaa.pdf $(INSTALL_DIR)/rm-fulldocs.pdf

%    \end{macrocode}
% Make the PDF and gzipped PS files. 
%    \begin{macrocode}
docs: $(DOCFILES)

%    \end{macrocode}
% Make everything else if either the |.dtx| or |.ins| files have
% changed. 
%    \begin{macrocode}
macros: rmaa.dtx rmaa.ins
%    \end{macrocode}
% Ensure that all the macro files can be written to then recreate
% them. The |touch| is there in case any of the files don't exist
% yet. 
%    \begin{macrocode}
        touch $(MACROFILES) $(MAKEFILE)
        chmod u+w $(MACROFILES) $(MAKEFILE) rmaa.ins
        latex "\input{rmaa.ins}"
%    \end{macrocode}
% \TeX{} automatically adds the |.tex| extension to extensionless
% files. We fix this and take the opportunity to convert
% back to TABs all that it has lovingly converted to spaces. 
%    \begin{macrocode}
        cp -f $(MAKEFILE) $(MAKEFILE).BCK && unexpand $(MAKEFILE).tex >| $(MAKEFILE) && rm $(MAKEFILE).tex
%    \end{macrocode}
% Make the shell scripts executable and everything unwriteable (so
% they don't get hand-edited by mistake). However, undo this for the
% files that the user is \emph{supposed} to modify. 
%    \begin{macrocode}
        chmod a+x $(SHELLSCRIPTS)
        chmod a-w $(MACROFILES) $(MAKEFILE)
        chmod u+w $(SKELETONFILES)

%    \end{macrocode}
% We now use PDF\LaTeX{} for this.
%    \begin{macrocode}
rmaa.idx: rmaa.dtx
        @echo "REGENERATING .idx FILE"
        pdflatex "\nonstopmode\input{rmaa.dtx}"

rmaa.aux: rmaa.dtx rmaa.ind rmaa.gls
        @echo "REGENERATING .aux FILE"
        pdflatex "\nonstopmode\input{rmaa.dtx}"

rmaa.ind: rmaa.idx
        @echo "REGENERATING INDEX"
        makeindex -s gind.ist rmaa

rmaa.gls: rmaa.glo
        @echo "REGENERATING CHANGELOG"
        makeindex -s gglo.ist -o rmaa.gls rmaa.glo

rmaa.pdf: rmaa.aux rmaa.ind rmaa.gls
        @echo "REGENERATING .pdf FILE"
        pdflatex "\nonstopmode\input{rmaa.dtx}"

clean: 
        rm -f rmaa.aux rmaa.dvi rmaa.ind rmaa.gls
%</makefile>
%    \end{macrocode}

% The |rmaa.ins| file uses the \textsf{docstrip} package to split out
% all the files and strip off this documentation.
%    \begin{macrocode}
%<*ins>
%%
%% Installation batch file for the RMAA LaTeX package
%%
\def\batchfile{rmaa.ins}
\input docstrip.tex
\preamble

PLEASE DO NOT EDIT THIS FILE! It is generated automatically from
the file rmaa.dtx using the accompanying Makefile and rmaa.ins 
Any hand-edited changes will be lost

A LaTeX document class for the RMxAA

Author: William Henney (w.henney@irya.unam.mx)

Latest version available from: http://www.irya.unam.mx/rmaa

Copyright (C) 1999-2019 Instituto de Astronomia, UNAM, Mexico

Permission is hereby granted, free of charge, to any person obtaining
a copy of this software and associated documentation files (the
"Software"), to deal in the Software without restriction, including
without limitation the rights to use, copy, modify, merge, publish,
distribute, sublicense, and/or sell copies of the Software, and to
permit persons to whom the Software is furnished to do so, subject to
the following conditions:

The above copyright notice and this permission notice shall be
included in all copies or substantial portions of the Software.

THE SOFTWARE IS PROVIDED "AS IS", WITHOUT WARRANTY OF ANY KIND,
EXPRESS OR IMPLIED, INCLUDING BUT NOT LIMITED TO THE WARRANTIES OF
MERCHANTABILITY, FITNESS FOR A PARTICULAR PURPOSE AND
NONINFRINGEMENT. IN NO EVENT SHALL THE AUTHORS OR COPYRIGHT HOLDERS BE
LIABLE FOR ANY CLAIM, DAMAGES OR OTHER LIABILITY, WHETHER IN AN ACTION
OF CONTRACT, TORT OR OTHERWISE, ARISING FROM, OUT OF OR IN CONNECTION
WITH THE SOFTWARE OR THE USE OR OTHER DEALINGS IN THE SOFTWARE.
\endpreamble
%    \end{macrocode}
% Clobber files without asking
%    \begin{macrocode}
\askforoverwritefalse
%    \end{macrocode}
% Make the \LaTeX{} macro files (class and styles)
%    \begin{macrocode}
\generate{% 
  \file{rmaa.cls}{\from{rmaa.dtx}{package}}
  \file{rmaa.ins}{\from{rmaa.dtx}{ins}}
  \file{rmaacite.sty}{\from{rmaa.dtx}{rmaacite}}
  \file{adshtml.sty}{\from{rmaa.dtx}{adshtml-style}}
  \file{ADS.tex}{\from{rmaa.dtx}{ads-container}}
  }
%    \end{macrocode}
% Make the other non-\LaTeX{} files. These are done separately since
% they can't use `\%\%' as a comment character for the
% preamble/postamble (and Perl files can't have the postamble I don't
% think). 
%    \begin{macrocode}
\generate{\nopreamble\nopostamble% 
  \file{Makefile}{\from{rmaa.dtx}{makefile}}
  \file{adshtml.perl}{\from{rmaa.dtx}{adshtml-perl}}
  \file{makeads.sh}{\from{rmaa.dtx}{ads-shell-script}}
  \file{sample.latex2html-init}{\from{rmaa.dtx}{latex2html-init}}
  }

%% Local Variables: 
%% mode: LaTeX
%% mode: Font-lock
%% End:
%</ins>
%    \end{macrocode}

% \subsection{Extended citation syntax}
% Modified from a file called mnrascite.sty
% I no longer remember where I stole the original version of this from
% 
% Alterations to the cite commands to make them do more sensible things.
% N.B. the bibliography style should produce entries which have labels where
% the date is delimited by angle brackets '<>' for these macros to work
% properly. They will not give very sensible error messages if you do not do 
% this properly.
% 
% A new command |\scite{key}| is defined which produces a citation which is 
% suitable for insertion in a sentence. i.e., it does not put brackets around 
% the whole citation, only the date. Also |\pcite{key}| produce output with no
% parentheses whatsoever. 
% 
% Should redefine |\@bibitem| really to have any hope of compatibility
% the command |\lbibcite| is defined to deal with the citations.
%    \begin{macrocode}
%<*rmaacite>
\def\bibitem[#1<#2>]#3{\par\hangindent=1.5em\hangafter=1
  \if@filesw
  {\def\protect##1{\string ##1\space}
    \immediate\write\@auxout{\string\lbibcite{#3}{#1}{#2}}}
  \fi}
\@redefinebibitem % what it says
%    \end{macrocode}
% Note that the |\lbibcite| defines two types of citation one contains the names,
% the other the year
%    \begin{macrocode}
\def\lbibcite#1#2#3{\global\@namedef{b@#1}{#2}
                    \global\@namedef{by@#1}{#3}}

\def\@bcite#1#2{({#1\if@tempswa , #2\fi})}
\def\@pcite#1#2{#1\if@tempswa , #2\fi}
%    \end{macrocode}
% This declaration contains what is put between the name and the date
% \changes{3.11}{2003/01/31}{Moved the \cs{articlename} to after the key for consistency
% with the new natbib stuff.}
%    \begin{macrocode}
\def\@brcite{\ (}
\def\@ebrcite{)}
\def\@fmttop#1{\@fmtcite[\csname #1\endcsname]}
\def\@citex[#1]#2{%
%%  \if@filesw\immediate\write\@auxout{\string\citation{#2}}\fi
%% WJH added 16 Jan 2002
  \@for\@citeb:=#2\do{%
    \edef\@citeb{\expandafter\@firstofone\@citeb:\articlename}%
    \if@filesw\immediate\write\@auxout{\string\citation{\@citeb}}\fi}%
%% to here
  \def\@citea{}%
  \@cite{\@for\@citeb:=#2\do
    {%
      \edef\@citeb{\expandafter\@firstofone\@citeb:\articlename}%
      \@citea\def\@citea{;\ }\@ifundefined
      {b@\@citeb}{{\bf ?}\@warning {Citation `\@citeb' on page \thepage \space undefined}}%
      {\csname b@\@citeb\endcsname\@citebro\csname by@\@citeb\endcsname\@citebrc}}}%
  {#1}}

\def\cite{\@ifnextchar [{\let\@citebro=\space
                          \let\@citebrc=\relax
                          \let\@cite=\@bcite\@tempswatrue\@citex}
                        {\let\@citebro=\space
                          \let\@citebrc=\relax
                          \let\@cite=\@bcite\@tempswafalse\@citex[]}}
\def\pcite{\@ifnextchar [{\let\@citebro=\space
                          \let\@citebrc=\relax
                          \let\@cite=\@pcite\@tempswatrue\@citex}
                        {\let\@citebro=\space
                          \let\@citebrc=\relax
                          \let\@cite=\@pcite\@tempswafalse\@citex[]}}
\def\scite{\@ifnextchar [{\let\@citebro=\@brcite
                          \let\@citebrc=\@ebrcite
                          \let\@cite=\@pcite\@tempswatrue\@citex}
                        {\let\@citebro=\@brcite
                          \let\@citebrc=\@ebrcite
                          \let\@cite=\@pcite\@tempswafalse\@citex[]}}

\renewenvironment{thebibliography}
  {\section*{REFERENCES}\vspace*{\afterbiblioskip}\small
    \if@book\RereadBiblio\fi
    \sloppy\clubpenalty4000\widowpenalty4000
    \sfcode`\.=1000\relax
    \setlength\parindent{\z@}\setlength\parskip{\z@}}
  {\relax\normalsize\@outputaddressesifnecessary} 
%</rmaacite>
%    \end{macrocode}

% \subsection{Auxilliary files for generating ADS data}

% The ADS data are written by \LaTeX{} to a |.ads| file. This file
% then has to be included in the following file (|ADS.tex|),
% which is run through \LaTeX{}2HTML (see below). 

%    \begin{macrocode}
%<*ads-container>
%%
%% Instructions:
%% Edit the \input{} line below to reflect the name of your file. 
\documentclass{article}
\usepackage[latin1]{inputenc}
\usepackage{adshtml}
\newcommand{\ADSfield}[2]
{\%#1  #2\\}
%%\newcommand{\ADSvfield}[2]
%%{\%#1 \begin{verbatim}#2\end{verbatim}}
\newenvironment{ADSarticle}{}{\par\bigskip}
\renewcommand{\bibitem}[2][]{\par}
\setlength{\parindent}{0pt}
\renewcommand{\ThisVolume}[1]{%  
  2001, \RMAASC, 12, \oldpageref{#1:RMAAFirstPage}}  
 

\begin{document}
\makeatletter
\input{rmbooksamp.ads}
\end{document}
%</ads-container>
%    \end{macrocode}

% This is the |adshtml.sty| style file used by |ADS.tex|. It just
% defines a few 
% macros that might be used in abstracts, etc (really it should define
% them all). This now gets all its code from
% Sec~\ref{sec:user-commands}. 
%    \begin{macrocode}
%<*adshtml-style>
%% Nothing else needs to go here now. 
%</adshtml-style>
%    \end{macrocode}
% The init file for latex2html.
%    \begin{macrocode}
%<*latex2html-init>
$PDFTOCAIRO = "pdftocairo";
$PDFLATEX = "pdflatex";
$LATEX = "latex";
$MAX_SPLIT_DEPTH = 0;
$HTML_VERSION = "4.0,math";
$NO_SIMPLE_MATH = 1;
$ADDRESS = "<p><address><a href=\"mailto:rmaa\@astro.unam.mx\">rmaa\@astro.unam.mx</address></p>"; 
$INFO = "<p>This document was created using latex2html and the rmaa latex macros</p>"; 
$NO_NAVIGATION = 1;
1; 
%</latex2html-init>
%    \end{macrocode}
% This is the |adshtml.perl| file that is used by \LaTeX{}2HTML. Some
% bizarre \TeX{}y command sometimes show up so we tell it to ignore
% these. All the astronomy macros that come from AAS are processed in
% \TeX{}. That way they get made to images but with an |ALT| tag that
% has the original code, which is what will be picked up by |lynx|,
% with the result that they get passed intact into the final |.txt|
% file. 
%    \begin{macrocode}
%<*adshtml-perl>
package main;

$dum = &do_cmd_makeatletter();

sub do_cmd_RMAA {
    local($_) = @_;
    "RMxAA" . $_;
}

sub do_cmd_RMAASC {
    local($_) = @_;
    "RMxAASC" . $_;
}

sub do_cmd_etal {
    local($_) = @_;
    "et al\." . $_;
}

sub do_cmd_ibidrule {
    local($_) = @_;
    "<img src=\"foo\" alt=\"\ibidrule\">" . $_;
}


sub do_cmd_bibitem {
    local($_) = @_;
    local($key);
    local($opt, $dummy)=&get_next_optional_argument; 
    $key = &missing_braces unless
        ((s/$next_pair_pr_rx/$key = $2; ''/eo)
        ||(s/$next_pair_rx/$key = $2; ''/eo));
    "<p>" . $_; 
}

&ignore_commands (<<_RAW_ARG_CMDS_);
lowercase
adjustfinalcols
textsuperscript # {}
newline # []
thinmuskip
mskip
smallskip
medskip
bigskip
hspace # {}
hspacestar # {}
vspace # {}
vspacestar # {}
mathrm
_at_mathrm
spacefactor
discretionary
unhbox 
voidb_at_x 
penalty 
_at_M
_at_m
null
IeC
_RAW_ARG_CMDS_

&process_commands_in_tex (<<_RAW_ARG_CMDS_);
sun
earth
lesssim
gtrsim
sq
argdeg
arcmin
arcsec
fd
fh
fm
fs
fdg
farcm
farcs
fp
micron
la
ga
case # {} # {}
slantfrac # {} # {}
onehalf
onethird
twothirds
onequarter
threequarters
ubvr
ub
bv
vr
ur
aj
araa
apj
apjl
apjs
ao
apss
aap
aapr
aaps
azh
baas
jrasc
memras
mnras
pra
prb
prc
prd
pre
prl
pasp
pasj
qjras
skytel
solphys
sovast
ssr
zap
nat
iaucirc
aplett
apspr
bain
fcp
gca
grl
jcp
jgr
jqsrt
memsai
nphysa
physrep
physscr
planss
procspie
_RAW_ARG_CMDS_

1;                              # This must be the last line
%</adshtml-perl> 
%    \end{macrocode}

% This is a shell script |makeads.sh| that generates the |_ads.txt| file from
% |_ads.tex| by using \LaTeX{}2HTML and |lynx|.  \changes{4.0}{2019/11/13}{Make sure that lynx writes utf-8 output instead of latin-1.}
% \changes{2.2beta13}{2002/03/12}{Complete rewrite of |makeads.sh|. Now much 
% improved I think}
% \changes{4.0}{2019/11/13}{Fix the at-sign substitution so that it only eliminates \cs{@} and does not mess with email addresses.}

%    \begin{macrocode}
%<*ads-shell-script>
#! /bin/sh
if [ -z $1 ]; then
    echo "Usage: $0 basename"
    exit
fi
BASENAME=$1
ADSBASENAME=${BASENAME}_ads
AUXFILE=$BASENAME.aux
TMPFILE=tmp.$$


echo "Generating $ADSBASENAME.tex"

# For some reason, must double-up \\ to \\\\ but not \ to \\ - odd!
cat > $ADSBASENAME.tex <<EOF
%% PLEASE DO NOT EDIT THIS FILE!  It is generated automatically from
%% the script makeads.sh.  Any hand-edited changes will be lost
%% 
\documentclass{article}
\usepackage[latin1]{inputenc}
\usepackage{html}
\usepackage{adshtml}
\newcommand{\ADSfield}[2]{\%#1  #2\\\\}
\newenvironment{ADSarticle}{}{\par\bigskip}
\setlength{\parindent}{0pt}
\newcounter{istage}
\renewcommand{\ion}[2]{\setcounter{istage}{#2}#1 \Roman{istage}}
\begin{document}
\makeatletter
EOF
echo >> $ADSBASENAME.tex "\input{$BASENAME.ads2}" 
echo >> $ADSBASENAME.tex "\end{document}"



echo "Generating $ADSBASENAME.sed" 

# Find type of journal
grep '\\documentclass\[.*proceedings.*\]' $BASENAME.tex >/dev/null \
     && JOURNAL='RMxAC' \
     || JOURNAL='RMxAA'

# Find the year and volume number
YEAR=`grep '\\SetYear' $BASENAME.tex | sed 's/.*\\SetYear{\(.*\)}.*/\1/'`
VOLUME=`grep '\\SetVolume' $BASENAME.tex | sed 's/.*\\SetVolume{\(.*\)}.*/\1/'`

# Generate an auxilliary sed script 
cat > $TMPFILE <<EOF
s|\\\\newlabel{\(.*\):RMAAFirstPage}{{.*}{\(.*\)}}|s/\\\\\\\\ThisVolume{\1}/$YEAR, RMxAC, $VOLUME, \2/g|g
EOF

# Now go through the .aux file, looking for \newlabels and use the sed 
# script just generated in order to make another sed script, one that 
# will replace \ThisVolume{LABEL} with the correct reference
grep '^\\newlabel{.*:RMAAFirstPage}' $AUXFILE \
     | sed -f $TMPFILE > $ADSBASENAME.sed && rm $TMPFILE

# Apend a further sed command to get rid of \@, since redefining 
# it in _ads.tex didn't work
echo >> $ADSBASENAME.sed 's/\\\\@//g'

echo "Generating $BASENAME.ads2 from $ADSBASENAME.sed and $BASENAME.ads" 

# Apply this to the .ads file, generating the .ads2 file, which is included
# by the _ads.tex file.
sed -f $ADSBASENAME.sed $BASENAME.ads > $BASENAME.ads2

echo "Generating $ADSBASENAME.txt from $ADSBASENAME.tex and $BASENAME.ads2"

# Now run latex2html on the _ads.tex file
latex2html -no_navigation -info 0 -address 0 -html_version 4.0,math $ADSBASENAME

# Then, convert to utf8 with lynx and tidy it up a bit
lynx --justify=off --width=100 --dump --nolist --display_charset=utf-8 \
     $ADSBASENAME/$ADSBASENAME.html | cut -c4- \
     | sed 's/ibidrule/\\ibidrule/g' > $ADSBASENAME.txt
%</ads-shell-script>
%    \end{macrocode}

% \section{Epilogue}
% Err, that's it \dots

\endinput


% \endinput
% Local Variables:
% mode: doctex
% TeX-master: t
% End:
"

clean: 
        rm -f rmaa.aux rmaa.dvi rmaa.ind rmaa.gls
%</makefile>
%    \end{macrocode}

% The |rmaa.ins| file uses the \textsf{docstrip} package to split out
% all the files and strip off this documentation.
%    \begin{macrocode}
%<*ins>
%%
%% Installation batch file for the RMAA LaTeX package
%%
\def\batchfile{rmaa.ins}
\input docstrip.tex
\preamble

PLEASE DO NOT EDIT THIS FILE! It is generated automatically from
the file rmaa.dtx using the accompanying Makefile and rmaa.ins 
Any hand-edited changes will be lost

A LaTeX document class for the RMxAA

Author: William Henney (w.henney@irya.unam.mx)

Latest version available from: http://www.irya.unam.mx/rmaa

Copyright (C) 1999-2019 Instituto de Astronomia, UNAM, Mexico

Permission is hereby granted, free of charge, to any person obtaining
a copy of this software and associated documentation files (the
"Software"), to deal in the Software without restriction, including
without limitation the rights to use, copy, modify, merge, publish,
distribute, sublicense, and/or sell copies of the Software, and to
permit persons to whom the Software is furnished to do so, subject to
the following conditions:

The above copyright notice and this permission notice shall be
included in all copies or substantial portions of the Software.

THE SOFTWARE IS PROVIDED "AS IS", WITHOUT WARRANTY OF ANY KIND,
EXPRESS OR IMPLIED, INCLUDING BUT NOT LIMITED TO THE WARRANTIES OF
MERCHANTABILITY, FITNESS FOR A PARTICULAR PURPOSE AND
NONINFRINGEMENT. IN NO EVENT SHALL THE AUTHORS OR COPYRIGHT HOLDERS BE
LIABLE FOR ANY CLAIM, DAMAGES OR OTHER LIABILITY, WHETHER IN AN ACTION
OF CONTRACT, TORT OR OTHERWISE, ARISING FROM, OUT OF OR IN CONNECTION
WITH THE SOFTWARE OR THE USE OR OTHER DEALINGS IN THE SOFTWARE.
\endpreamble
%    \end{macrocode}
% Clobber files without asking
%    \begin{macrocode}
\askforoverwritefalse
%    \end{macrocode}
% Make the \LaTeX{} macro files (class and styles)
%    \begin{macrocode}
\generate{% 
  \file{rmaa.cls}{\from{rmaa.dtx}{package}}
  \file{rmaa.ins}{\from{rmaa.dtx}{ins}}
  \file{rmaacite.sty}{\from{rmaa.dtx}{rmaacite}}
  \file{adshtml.sty}{\from{rmaa.dtx}{adshtml-style}}
  \file{ADS.tex}{\from{rmaa.dtx}{ads-container}}
  }
%    \end{macrocode}
% Make the other non-\LaTeX{} files. These are done separately since
% they can't use `\%\%' as a comment character for the
% preamble/postamble (and Perl files can't have the postamble I don't
% think). 
%    \begin{macrocode}
\generate{\nopreamble\nopostamble% 
  \file{Makefile}{\from{rmaa.dtx}{makefile}}
  \file{adshtml.perl}{\from{rmaa.dtx}{adshtml-perl}}
  \file{makeads.sh}{\from{rmaa.dtx}{ads-shell-script}}
  \file{sample.latex2html-init}{\from{rmaa.dtx}{latex2html-init}}
  }

%% Local Variables: 
%% mode: LaTeX
%% mode: Font-lock
%% End:
%</ins>
%    \end{macrocode}

% \subsection{Extended citation syntax}
% Modified from a file called mnrascite.sty
% I no longer remember where I stole the original version of this from
% 
% Alterations to the cite commands to make them do more sensible things.
% N.B. the bibliography style should produce entries which have labels where
% the date is delimited by angle brackets '<>' for these macros to work
% properly. They will not give very sensible error messages if you do not do 
% this properly.
% 
% A new command |\scite{key}| is defined which produces a citation which is 
% suitable for insertion in a sentence. i.e., it does not put brackets around 
% the whole citation, only the date. Also |\pcite{key}| produce output with no
% parentheses whatsoever. 
% 
% Should redefine |\@bibitem| really to have any hope of compatibility
% the command |\lbibcite| is defined to deal with the citations.
%    \begin{macrocode}
%<*rmaacite>
\def\bibitem[#1<#2>]#3{\par\hangindent=1.5em\hangafter=1
  \if@filesw
  {\def\protect##1{\string ##1\space}
    \immediate\write\@auxout{\string\lbibcite{#3}{#1}{#2}}}
  \fi}
\@redefinebibitem % what it says
%    \end{macrocode}
% Note that the |\lbibcite| defines two types of citation one contains the names,
% the other the year
%    \begin{macrocode}
\def\lbibcite#1#2#3{\global\@namedef{b@#1}{#2}
                    \global\@namedef{by@#1}{#3}}

\def\@bcite#1#2{({#1\if@tempswa , #2\fi})}
\def\@pcite#1#2{#1\if@tempswa , #2\fi}
%    \end{macrocode}
% This declaration contains what is put between the name and the date
% \changes{3.11}{2003/01/31}{Moved the \cs{articlename} to after the key for consistency
% with the new natbib stuff.}
%    \begin{macrocode}
\def\@brcite{\ (}
\def\@ebrcite{)}
\def\@fmttop#1{\@fmtcite[\csname #1\endcsname]}
\def\@citex[#1]#2{%
%%  \if@filesw\immediate\write\@auxout{\string\citation{#2}}\fi
%% WJH added 16 Jan 2002
  \@for\@citeb:=#2\do{%
    \edef\@citeb{\expandafter\@firstofone\@citeb:\articlename}%
    \if@filesw\immediate\write\@auxout{\string\citation{\@citeb}}\fi}%
%% to here
  \def\@citea{}%
  \@cite{\@for\@citeb:=#2\do
    {%
      \edef\@citeb{\expandafter\@firstofone\@citeb:\articlename}%
      \@citea\def\@citea{;\ }\@ifundefined
      {b@\@citeb}{{\bf ?}\@warning {Citation `\@citeb' on page \thepage \space undefined}}%
      {\csname b@\@citeb\endcsname\@citebro\csname by@\@citeb\endcsname\@citebrc}}}%
  {#1}}

\def\cite{\@ifnextchar [{\let\@citebro=\space
                          \let\@citebrc=\relax
                          \let\@cite=\@bcite\@tempswatrue\@citex}
                        {\let\@citebro=\space
                          \let\@citebrc=\relax
                          \let\@cite=\@bcite\@tempswafalse\@citex[]}}
\def\pcite{\@ifnextchar [{\let\@citebro=\space
                          \let\@citebrc=\relax
                          \let\@cite=\@pcite\@tempswatrue\@citex}
                        {\let\@citebro=\space
                          \let\@citebrc=\relax
                          \let\@cite=\@pcite\@tempswafalse\@citex[]}}
\def\scite{\@ifnextchar [{\let\@citebro=\@brcite
                          \let\@citebrc=\@ebrcite
                          \let\@cite=\@pcite\@tempswatrue\@citex}
                        {\let\@citebro=\@brcite
                          \let\@citebrc=\@ebrcite
                          \let\@cite=\@pcite\@tempswafalse\@citex[]}}

\renewenvironment{thebibliography}
  {\section*{REFERENCES}\vspace*{\afterbiblioskip}\small
    \if@book\RereadBiblio\fi
    \sloppy\clubpenalty4000\widowpenalty4000
    \sfcode`\.=1000\relax
    \setlength\parindent{\z@}\setlength\parskip{\z@}}
  {\relax\normalsize\@outputaddressesifnecessary} 
%</rmaacite>
%    \end{macrocode}

% \subsection{Auxilliary files for generating ADS data}

% The ADS data are written by \LaTeX{} to a |.ads| file. This file
% then has to be included in the following file (|ADS.tex|),
% which is run through \LaTeX{}2HTML (see below). 

%    \begin{macrocode}
%<*ads-container>
%%
%% Instructions:
%% Edit the \input{} line below to reflect the name of your file. 
\documentclass{article}
\usepackage[latin1]{inputenc}
\usepackage{adshtml}
\newcommand{\ADSfield}[2]
{\%#1  #2\\}
%%\newcommand{\ADSvfield}[2]
%%{\%#1 \begin{verbatim}#2\end{verbatim}}
\newenvironment{ADSarticle}{}{\par\bigskip}
\renewcommand{\bibitem}[2][]{\par}
\setlength{\parindent}{0pt}
\renewcommand{\ThisVolume}[1]{%  
  2001, \RMAASC, 12, \oldpageref{#1:RMAAFirstPage}}  
 

\begin{document}
\makeatletter
%% Demonstration file for formatting an entire volume
\documentclass[book,journal,twocolumn]{rmaa}
%%%%%%%%%%%%%%%%%%%%%%%%%%%%%%%%%%%%%%%%%%%%%%%%%%%%%%%%%%%%%%%%
%%%%%%%%%%%%%%%%            PREAMBLE            %%%%%%%%%%%%%%%%
%%%%%%%%%%%%%%%%%%%%%%%%%%%%%%%%%%%%%%%%%%%%%%%%%%%%%%%%%%%%%%%%

\usepackage{paralist}
\usepackage{psfrag,color}
% Needed for landscape tables
\usepackage{pdflscape}

% Used in rm-minimal-test
\usepackage{kantlipsum}

% publication details of volume
\SetVolume{9}
\SetYear{2000}
\SetMonthNumeric{06}
\SetMes{junio}
% details of meeting
\SetConfTitle{Astrophysical Plasmas: Codes, Models, \& Observations}
\SetReunion{M\'exico, D. F.}{octubre 25--29, 1999}
\SetEditors{S. Jane Arthur, Nancy S. Brickhouse, \& Jos\'e Franco}

% this writes the .idx file that is used to generate the index 
\makeindex

%% This command makes the session headings come out bold in the
%% Table-of-Contents. Comment it out if you don't like it.
\BoldSessionTitles

%% This command makes a title page for each session
%% Comment it out if you don't like it.
\SessionTitlePages

%% Use this command if you want the abstract-only papers to be
%% indexed by ADS
\abstractsADS

%% At the special request of Alberto Carramin~ana, the following is a
%% way to suppress the author adresses at the end of each article
%\suppressfulladdresses

%% Use this command to add Received/accepted dates for each paper
% \printdates


\begin{document}

%%%%%%%%%%%%%%%%%%%%%%%%%%%%%%%%%%%%%%%%%%%%%%%%%%%%%%%%%%%%%%%%b
%%%%%%%%%%%%%%%%          FRONT MATTER          %%%%%%%%%%%%%%%%
%%%%%%%%%%%%%%%%%%%%%%%%%%%%%%%%%%%%%%%%%%%%%%%%%%%%%%%%%%%%%%%%
\frontmatter % start roman numbering

% format the TOC
\renewcommand{\RMAAbannerimage}{rmsc_blank}
\tableofcontents

%% preface, group photo, list of participants go here
%% use the `frontsection' environment (see editors' guide)
\begin{frontsection}{Preface}{%
    \listofauthors{S.~J.~Arthur, N.~S.~Brickhouse, J.~Franco}
    \indexauthor{Arthur, S. J.}
    \indexauthor{Brickhouse, N.}
    \indexauthor{Franco, J.}
    }
  These are the proceedings of a very nice conference that Jane and
  Pepe organised. This is the preface, which explains what a wonderful
  time was had by all and how grateful they are for all the people who
  gave them money or helped out.
  
  This is another paragraph of the preface, in case they have anything
  else to say.  For example, the Scientific Organizing Committee
  consisted of
  \begin{quote}
    J. Franco (Mexico, Chair), D. Cox (USA) and various other notables
  \end{quote}
  while the Local Organizing Committee comprised
  \begin{quote}
    S. J. Arthur (Chair), S. Kurtz, etc.
  \end{quote}

  \EditorSignature

\end{frontsection}

\begin{frontsection}[Photo]{Group Photograph}{\notitle}
\null\vfill
\centerline{Here goes the photograph.}
\vfill
\end{frontsection}

\begin{frontsection}[Participants]{List of Participants}{}
%% Here you put a formatted list of participants
\vfil
\dotfill First page of list of participants \dotfill
\vfil\clearpage

\null\vfil
\dotfill Second page of list of participants \dotfill
\vfil

\end{frontsection}

\begin{frontsection}[XXX]{Another front section}{}

\end{frontsection}

%%%%%%%%%%%%%%%%%%%%%%%%%%%%%%%%%%%%%%%%%%%%%%%%%%%%%%%%%%%%%%%%
%%%%%%%%%%%%%%%%          MAIN MATTER           %%%%%%%%%%%%%%%%
%%%%%%%%%%%%%%%%%%%%%%%%%%%%%%%%%%%%%%%%%%%%%%%%%%%%%%%%%%%%%%%%
\mainmatter % start arabic numbering

\begin{articles}
  %% Main articles go here
  %% Use \article{ARTICLENAME} for each individual article
  %% If you want any text to appear in brackets after the article
  %% title in the TOC and index, then put it as an optional first
  %% argument, e.g.,  \article[Invited Talk]{ARTICLENAME} 

  \session{The Contributed Papers}
  \article{rm-journal-example} % An illustrative article
  \article{rm-onepage} % another article, which has a plate
  \article{rm-minimal-test}
  \session{Conference Summary}
  \article[Invited Review]{rmsummsamp} % An illustrative conference summary

\end{articles}

%%%%%%%%%%%%%%%%%%%%%%%%%%%%%%%%%%%%%%%%%%%%%%%%%%%%%%%%%%%%%%%%
%%%%%%%%%%%%%%%%          BACK MATTER           %%%%%%%%%%%%%%%%
%%%%%%%%%%%%%%%%%%%%%%%%%%%%%%%%%%%%%%%%%%%%%%%%%%%%%%%%%%%%%%%%
\backmatter

%% Abstracts of presentations for which the
%% authors have not submitted a paper
\begin{abstracts}

  \title{An abstract of a poster}
  \author{W. J. Henney\altaffilmark{1}}
  \listofauthors{W.~J.~Henney}
  \indexauthor{Henney, W. J}
  \altaffiltext{1}{Where I am from done with altaffilmark/text}
  \abstract{This is a test abstract.   This is a test abstract.   This
    is a test abstract.   This is a test abstract.}
  

  \title{An abstract of another  poster}
  \author{W.~J.~Henney\altaffil{This is done with altaffil} and
    A.~N.~Other\altaffil{So is this} }
  \listofauthors{W.~J.~Henney \& A.~N.~Other}
  \indexauthor{Henney, W. J}
  \indexauthor{Other, A. N.}
  \abstract{This is a test abstract.   This is a test abstract.   This
    is a test abstract.   This is a test abstract.}

  \title{An abstract of yet another  poster}
  \author{W.~J.~Henney\altaffil{This is done with altaffil} and
    A.~N.~Other\altaffil{So is this} }
  \listofauthors{W.~J.~Henney \& A.~N.~Other}
  \indexauthor{Henney, W. J}
  \indexauthor{Other, A. N.}
  \abstract{This is a test abstract.   This is a test abstract.   This
    is a test abstract.   This is a test abstract.}

  \title{The abstract of a poster with a very long title to see how
    well it gets split up into lines}
  \author{A.~N.~Other\altaffil{The affiliation of A.~N.~Other is also
      very long so that it will go over several lines. His email
      address is anon@nowhere.org}}
  \listofauthors{A.~N.~Other}
  \indexauthor{Other, A. N.}
  \abstract{We will try and make this abstract longer so there is more
    stuff to put on the page and it doesn't look too empty. However, I
    have nothing whatsoever to say so will have to stop here. Let's
    repeat that a few times. We will try and make this abstract longer
    so there is more stuff to put on the page and it doesn't look too
    empty. However, I have nothing whatsoever to say so will have to
    stop here. Let's repeat that a few times. We will try and make
    this abstract longer so there is more stuff to put on the page and
    it doesn't look too empty. However, I have nothing whatsoever to
    say so will have to stop here. Let's repeat that a few times. We
    will try and make this abstract longer so there is more stuff to
    put on the page and it doesn't look too empty. However, I have
    nothing whatsoever to say so will have to stop here. Let's repeat
    that a few times.}
    
  \title[Comment to go in TOC]{The abstract of a second poster with a
    very long title to see how well it gets split up into lines}
  \author{A.~N.~Other\altaffil{The affiliation of A.~N.~Other is also
      very long so that it will go over several lines. His email
      address is anon@nowhere.org}}
  \listofauthors{A.~N.~Other}
  \indexauthor{Other, A. N.}
  \abstract{We will try and make this abstract longer so there is more
    stuff to put on the page and it doesn't look too empty. However, I
    have nothing whatsoever to say so will have to stop here. Let's
    repeat that a few times. We will try and make this abstract longer
    so there is more stuff to put on the page and it doesn't look too
    empty. However, I have nothing whatsoever to say so will have to
    stop here. Let's repeat that a few times. We will try and make
    this abstract longer so there is more stuff to put on the page and
    it doesn't look too empty. However, I have nothing whatsoever to
    say so will have to stop here. Let's repeat that a few times. We
    will try and make this abstract longer so there is more stuff to
    put on the page and it doesn't look too empty. However, I have
    nothing whatsoever to say so will have to stop here. Let's repeat
    that a few times.}
    
  \title{The abstract of a third poster with a very long title to see how
    well it gets split up into lines}
  \author{A.~N.~Other\altaffil{The affiliation of A.~N.~Other is also
      very long so that it will go over several lines. His email
      address is anon@nowhere.org}}
  \listofauthors{A.~N.~Other}
  \indexauthor{Other, A. N.}
  \abstract{We will try and make this abstract longer so there is more
    stuff to put on the page and it doesn't look too empty. However, I
    have nothing whatsoever to say so will have to stop here. Let's
    repeat that a few times. We will try and make this abstract longer
    so there is more stuff to put on the page and it doesn't look too
    empty. However, I have nothing whatsoever to say so will have to
    stop here. Let's repeat that a few times. We will try and make
    this abstract longer so there is more stuff to put on the page and
    it doesn't look too empty. However, I have nothing whatsoever to
    say so will have to stop here. Let's repeat that a few times. We
    will try and make this abstract longer so there is more stuff to
    put on the page and it doesn't look too empty. However, I have
    nothing whatsoever to say so will have to stop here. Let's repeat
    that a few times.}
    
  \title{The abstract of a fourth poster with a very long title to see how
    well it gets split up into lines}
  \author{A.~N.~Other\altaffil{The affiliation of A.~N.~Other is also
      very long so that it will go over several lines. His email
      address is anon@nowhere.org}}
  \listofauthors{A.~N.~Other}
  \indexauthor{Other, A. N.}
  \abstract{We will try and make this abstract longer so there is more
    stuff to put on the page and it doesn't look too empty. However, I
    have nothing whatsoever to say so will have to stop here. Let's
    repeat that a few times. We will try and make this abstract longer
    so there is more stuff to put on the page and it doesn't look too
    empty. However, I have nothing whatsoever to say so will have to
    stop here. Let's repeat that a few times. We will try and make
    this abstract longer so there is more stuff to put on the page and
    it doesn't look too empty. However, I have nothing whatsoever to
    say so will have to stop here. Let's repeat that a few times. We
    will try and make this abstract longer so there is more stuff to
    put on the page and it doesn't look too empty. However, I have
    nothing whatsoever to say so will have to stop here. Let's repeat
    that a few times.}
    
  \title{The abstract of a fifth poster with a very long title to see how
    well it gets split up into lines}
  \author{A.~N.~Other\altaffil{The affiliation of A.~N.~Other is also
      very long so that it will go over several lines. His email
      address is anon@nowhere.org}}
  \listofauthors{A.~N.~Other}
  \indexauthor{Other, A. N.}
  \abstract{We will try and make this abstract longer so there is more
    stuff to put on the page and it doesn't look too empty. However, I
    have nothing whatsoever to say so will have to stop here. Let's
    repeat that a few times. We will try and make this abstract longer
    so there is more stuff to put on the page and it doesn't look too
    empty. However, I have nothing whatsoever to say so will have to
    stop here. Let's repeat that a few times. We will try and make
    this abstract longer so there is more stuff to put on the page and
    it doesn't look too empty. However, I have nothing whatsoever to
    say so will have to stop here. Let's repeat that a few times. We
    will try and make this abstract longer so there is more stuff to
    put on the page and it doesn't look too empty. However, I have
    nothing whatsoever to say so will have to stop here. Let's repeat
    that a few times.}
    
  \title{The abstract of a sixth poster with a very long title to see how
    well it gets split up into lines}
  \author{A.~N.~Other\altaffil{The affiliation of A.~N.~Other is also
      very long so that it will go over several lines. His email
      address is anon@nowhere.org}}
  \listofauthors{A.~N.~Other}
  \indexauthor{Other, A. N.}
  \abstract{We will try and make this abstract longer so there is more
    stuff to put on the page and it doesn't look too empty. However, I
    have nothing whatsoever to say so will have to stop here. Let's
    repeat that a few times. We will try and make this abstract longer
    so there is more stuff to put on the page and it doesn't look too
    empty. However, I have nothing whatsoever to say so will have to
    stop here. Let's repeat that a few times. We will try and make
    this abstract longer so there is more stuff to put on the page and
    it doesn't look too empty. However, I have nothing whatsoever to
    say so will have to stop here. Let's repeat that a few times. We
    will try and make this abstract longer so there is more stuff to
    put on the page and it doesn't look too empty. However, I have
    nothing whatsoever to say so will have to stop here. Let's repeat
    that a few times.}
    
  \title{The abstract of a seventh poster with a very long title to see how
    well it gets split up into lines}
  \author{A.~N.~Other\altaffil{The affiliation of A.~N.~Other is also
      very long so that it will go over several lines. His email
      address is anon@nowhere.org}}
  \listofauthors{A.~N.~Other}
  \indexauthor{Other, A. N.}
  \abstract{We will try and make this abstract longer so there is more
    stuff to put on the page and it doesn't look too empty. However, I
    have nothing whatsoever to say so will have to stop here. Let's
    repeat that a few times. We will try and make this abstract longer
    so there is more stuff to put on the page and it doesn't look too
    empty. However, I have nothing whatsoever to say so will have to
    stop here. Let's repeat that a few times. We will try and make
    this abstract longer so there is more stuff to put on the page and
    it doesn't look too empty. However, I have nothing whatsoever to
    say so will have to stop here. Let's repeat that a few times. We
    will try and make this abstract longer so there is more stuff to
    put on the page and it doesn't look too empty. However, I have
    nothing whatsoever to say so will have to stop here. Let's repeat
    that a few times.}
  
  \title{An abstract to go in the list of abstracts, but not in the
    index of authors}
  \author{I.~A.~M. Hidden\altaffil{Puerto Escondido}}
  \listofauthors{I.~A.~M. Hidden}
  \abstract{The idea is that by omitting the
    \textbackslash{}indexauthor command, we can avoid an entry in the
    author index.}

\end{abstracts}


%% The Index: after latexing twice, run `makeindex -s rmaa' on this
%% file, then latex twice more
\printindex

%% if you rename this file to something other than `rmbooksamp', then
%% you must change the above line accordingly

%%
%% Plates: If you have any, then the following line will print
%% them out:

\outputplates

\end{document}

%%% Local Variables:
%%% mode: latex
%%% TeX-master: t
%%% End:

\end{document}
%</ads-container>
%    \end{macrocode}

% This is the |adshtml.sty| style file used by |ADS.tex|. It just
% defines a few 
% macros that might be used in abstracts, etc (really it should define
% them all). This now gets all its code from
% Sec~\ref{sec:user-commands}. 
%    \begin{macrocode}
%<*adshtml-style>
%% Nothing else needs to go here now. 
%</adshtml-style>
%    \end{macrocode}
% The init file for latex2html.
%    \begin{macrocode}
%<*latex2html-init>
$PDFTOCAIRO = "pdftocairo";
$PDFLATEX = "pdflatex";
$LATEX = "latex";
$MAX_SPLIT_DEPTH = 0;
$HTML_VERSION = "4.0,math";
$NO_SIMPLE_MATH = 1;
$ADDRESS = "<p><address><a href=\"mailto:rmaa\@astro.unam.mx\">rmaa\@astro.unam.mx</address></p>"; 
$INFO = "<p>This document was created using latex2html and the rmaa latex macros</p>"; 
$NO_NAVIGATION = 1;
1; 
%</latex2html-init>
%    \end{macrocode}
% This is the |adshtml.perl| file that is used by \LaTeX{}2HTML. Some
% bizarre \TeX{}y command sometimes show up so we tell it to ignore
% these. All the astronomy macros that come from AAS are processed in
% \TeX{}. That way they get made to images but with an |ALT| tag that
% has the original code, which is what will be picked up by |lynx|,
% with the result that they get passed intact into the final |.txt|
% file. 
%    \begin{macrocode}
%<*adshtml-perl>
package main;

$dum = &do_cmd_makeatletter();

sub do_cmd_RMAA {
    local($_) = @_;
    "RMxAA" . $_;
}

sub do_cmd_RMAASC {
    local($_) = @_;
    "RMxAASC" . $_;
}

sub do_cmd_etal {
    local($_) = @_;
    "et al\." . $_;
}

sub do_cmd_ibidrule {
    local($_) = @_;
    "<img src=\"foo\" alt=\"\ibidrule\">" . $_;
}


sub do_cmd_bibitem {
    local($_) = @_;
    local($key);
    local($opt, $dummy)=&get_next_optional_argument; 
    $key = &missing_braces unless
        ((s/$next_pair_pr_rx/$key = $2; ''/eo)
        ||(s/$next_pair_rx/$key = $2; ''/eo));
    "<p>" . $_; 
}

&ignore_commands (<<_RAW_ARG_CMDS_);
lowercase
adjustfinalcols
textsuperscript # {}
newline # []
thinmuskip
mskip
smallskip
medskip
bigskip
hspace # {}
hspacestar # {}
vspace # {}
vspacestar # {}
mathrm
_at_mathrm
spacefactor
discretionary
unhbox 
voidb_at_x 
penalty 
_at_M
_at_m
null
IeC
_RAW_ARG_CMDS_

&process_commands_in_tex (<<_RAW_ARG_CMDS_);
sun
earth
lesssim
gtrsim
sq
argdeg
arcmin
arcsec
fd
fh
fm
fs
fdg
farcm
farcs
fp
micron
la
ga
case # {} # {}
slantfrac # {} # {}
onehalf
onethird
twothirds
onequarter
threequarters
ubvr
ub
bv
vr
ur
aj
araa
apj
apjl
apjs
ao
apss
aap
aapr
aaps
azh
baas
jrasc
memras
mnras
pra
prb
prc
prd
pre
prl
pasp
pasj
qjras
skytel
solphys
sovast
ssr
zap
nat
iaucirc
aplett
apspr
bain
fcp
gca
grl
jcp
jgr
jqsrt
memsai
nphysa
physrep
physscr
planss
procspie
_RAW_ARG_CMDS_

1;                              # This must be the last line
%</adshtml-perl> 
%    \end{macrocode}

% This is a shell script |makeads.sh| that generates the |_ads.txt| file from
% |_ads.tex| by using \LaTeX{}2HTML and |lynx|.  \changes{4.0}{2019/11/13}{Make sure that lynx writes utf-8 output instead of latin-1.}
% \changes{2.2beta13}{2002/03/12}{Complete rewrite of |makeads.sh|. Now much 
% improved I think}
% \changes{4.0}{2019/11/13}{Fix the at-sign substitution so that it only eliminates \cs{@} and does not mess with email addresses.}

%    \begin{macrocode}
%<*ads-shell-script>
#! /bin/sh
if [ -z $1 ]; then
    echo "Usage: $0 basename"
    exit
fi
BASENAME=$1
ADSBASENAME=${BASENAME}_ads
AUXFILE=$BASENAME.aux
TMPFILE=tmp.$$


echo "Generating $ADSBASENAME.tex"

# For some reason, must double-up \\ to \\\\ but not \ to \\ - odd!
cat > $ADSBASENAME.tex <<EOF
%% PLEASE DO NOT EDIT THIS FILE!  It is generated automatically from
%% the script makeads.sh.  Any hand-edited changes will be lost
%% 
\documentclass{article}
\usepackage[latin1]{inputenc}
\usepackage{html}
\usepackage{adshtml}
\newcommand{\ADSfield}[2]{\%#1  #2\\\\}
\newenvironment{ADSarticle}{}{\par\bigskip}
\setlength{\parindent}{0pt}
\newcounter{istage}
\renewcommand{\ion}[2]{\setcounter{istage}{#2}#1 \Roman{istage}}
\begin{document}
\makeatletter
EOF
echo >> $ADSBASENAME.tex "\input{$BASENAME.ads2}" 
echo >> $ADSBASENAME.tex "\end{document}"



echo "Generating $ADSBASENAME.sed" 

# Find type of journal
grep '\\documentclass\[.*proceedings.*\]' $BASENAME.tex >/dev/null \
     && JOURNAL='RMxAC' \
     || JOURNAL='RMxAA'

# Find the year and volume number
YEAR=`grep '\\SetYear' $BASENAME.tex | sed 's/.*\\SetYear{\(.*\)}.*/\1/'`
VOLUME=`grep '\\SetVolume' $BASENAME.tex | sed 's/.*\\SetVolume{\(.*\)}.*/\1/'`

# Generate an auxilliary sed script 
cat > $TMPFILE <<EOF
s|\\\\newlabel{\(.*\):RMAAFirstPage}{{.*}{\(.*\)}}|s/\\\\\\\\ThisVolume{\1}/$YEAR, RMxAC, $VOLUME, \2/g|g
EOF

# Now go through the .aux file, looking for \newlabels and use the sed 
# script just generated in order to make another sed script, one that 
# will replace \ThisVolume{LABEL} with the correct reference
grep '^\\newlabel{.*:RMAAFirstPage}' $AUXFILE \
     | sed -f $TMPFILE > $ADSBASENAME.sed && rm $TMPFILE

# Apend a further sed command to get rid of \@, since redefining 
# it in _ads.tex didn't work
echo >> $ADSBASENAME.sed 's/\\\\@//g'

echo "Generating $BASENAME.ads2 from $ADSBASENAME.sed and $BASENAME.ads" 

# Apply this to the .ads file, generating the .ads2 file, which is included
# by the _ads.tex file.
sed -f $ADSBASENAME.sed $BASENAME.ads > $BASENAME.ads2

echo "Generating $ADSBASENAME.txt from $ADSBASENAME.tex and $BASENAME.ads2"

# Now run latex2html on the _ads.tex file
latex2html -no_navigation -info 0 -address 0 -html_version 4.0,math $ADSBASENAME

# Then, convert to utf8 with lynx and tidy it up a bit
lynx --justify=off --width=100 --dump --nolist --display_charset=utf-8 \
     $ADSBASENAME/$ADSBASENAME.html | cut -c4- \
     | sed 's/ibidrule/\\ibidrule/g' > $ADSBASENAME.txt
%</ads-shell-script>
%    \end{macrocode}

% \section{Epilogue}
% Err, that's it \dots

\endinput


% \endinput
% Local Variables:
% mode: doctex
% TeX-master: t
% End:
"

clean: 
        rm -f rmaa.aux rmaa.dvi rmaa.ind rmaa.gls
%</makefile>
%    \end{macrocode}

% The |rmaa.ins| file uses the \textsf{docstrip} package to split out
% all the files and strip off this documentation.
%    \begin{macrocode}
%<*ins>
%%
%% Installation batch file for the RMAA LaTeX package
%%
\def\batchfile{rmaa.ins}
\input docstrip.tex
\preamble

PLEASE DO NOT EDIT THIS FILE! It is generated automatically from
the file rmaa.dtx using the accompanying Makefile and rmaa.ins 
Any hand-edited changes will be lost

A LaTeX document class for the RMxAA

Author: William Henney (w.henney@irya.unam.mx)

Latest version available from: http://www.irya.unam.mx/rmaa

Copyright (C) 1999-2019 Instituto de Astronomia, UNAM, Mexico

Permission is hereby granted, free of charge, to any person obtaining
a copy of this software and associated documentation files (the
"Software"), to deal in the Software without restriction, including
without limitation the rights to use, copy, modify, merge, publish,
distribute, sublicense, and/or sell copies of the Software, and to
permit persons to whom the Software is furnished to do so, subject to
the following conditions:

The above copyright notice and this permission notice shall be
included in all copies or substantial portions of the Software.

THE SOFTWARE IS PROVIDED "AS IS", WITHOUT WARRANTY OF ANY KIND,
EXPRESS OR IMPLIED, INCLUDING BUT NOT LIMITED TO THE WARRANTIES OF
MERCHANTABILITY, FITNESS FOR A PARTICULAR PURPOSE AND
NONINFRINGEMENT. IN NO EVENT SHALL THE AUTHORS OR COPYRIGHT HOLDERS BE
LIABLE FOR ANY CLAIM, DAMAGES OR OTHER LIABILITY, WHETHER IN AN ACTION
OF CONTRACT, TORT OR OTHERWISE, ARISING FROM, OUT OF OR IN CONNECTION
WITH THE SOFTWARE OR THE USE OR OTHER DEALINGS IN THE SOFTWARE.
\endpreamble
%    \end{macrocode}
% Clobber files without asking
%    \begin{macrocode}
\askforoverwritefalse
%    \end{macrocode}
% Make the \LaTeX{} macro files (class and styles)
%    \begin{macrocode}
\generate{% 
  \file{rmaa.cls}{\from{rmaa.dtx}{package}}
  \file{rmaa.ins}{\from{rmaa.dtx}{ins}}
  \file{rmaacite.sty}{\from{rmaa.dtx}{rmaacite}}
  \file{adshtml.sty}{\from{rmaa.dtx}{adshtml-style}}
  \file{ADS.tex}{\from{rmaa.dtx}{ads-container}}
  }
%    \end{macrocode}
% Make the other non-\LaTeX{} files. These are done separately since
% they can't use `\%\%' as a comment character for the
% preamble/postamble (and Perl files can't have the postamble I don't
% think). 
%    \begin{macrocode}
\generate{\nopreamble\nopostamble% 
  \file{Makefile}{\from{rmaa.dtx}{makefile}}
  \file{adshtml.perl}{\from{rmaa.dtx}{adshtml-perl}}
  \file{makeads.sh}{\from{rmaa.dtx}{ads-shell-script}}
  \file{sample.latex2html-init}{\from{rmaa.dtx}{latex2html-init}}
  }

%% Local Variables: 
%% mode: LaTeX
%% mode: Font-lock
%% End:
%</ins>
%    \end{macrocode}

% \subsection{Extended citation syntax}
% Modified from a file called mnrascite.sty
% I no longer remember where I stole the original version of this from
% 
% Alterations to the cite commands to make them do more sensible things.
% N.B. the bibliography style should produce entries which have labels where
% the date is delimited by angle brackets '<>' for these macros to work
% properly. They will not give very sensible error messages if you do not do 
% this properly.
% 
% A new command |\scite{key}| is defined which produces a citation which is 
% suitable for insertion in a sentence. i.e., it does not put brackets around 
% the whole citation, only the date. Also |\pcite{key}| produce output with no
% parentheses whatsoever. 
% 
% Should redefine |\@bibitem| really to have any hope of compatibility
% the command |\lbibcite| is defined to deal with the citations.
%    \begin{macrocode}
%<*rmaacite>
\def\bibitem[#1<#2>]#3{\par\hangindent=1.5em\hangafter=1
  \if@filesw
  {\def\protect##1{\string ##1\space}
    \immediate\write\@auxout{\string\lbibcite{#3}{#1}{#2}}}
  \fi}
\@redefinebibitem % what it says
%    \end{macrocode}
% Note that the |\lbibcite| defines two types of citation one contains the names,
% the other the year
%    \begin{macrocode}
\def\lbibcite#1#2#3{\global\@namedef{b@#1}{#2}
                    \global\@namedef{by@#1}{#3}}

\def\@bcite#1#2{({#1\if@tempswa , #2\fi})}
\def\@pcite#1#2{#1\if@tempswa , #2\fi}
%    \end{macrocode}
% This declaration contains what is put between the name and the date
% \changes{3.11}{2003/01/31}{Moved the \cs{articlename} to after the key for consistency
% with the new natbib stuff.}
%    \begin{macrocode}
\def\@brcite{\ (}
\def\@ebrcite{)}
\def\@fmttop#1{\@fmtcite[\csname #1\endcsname]}
\def\@citex[#1]#2{%
%%  \if@filesw\immediate\write\@auxout{\string\citation{#2}}\fi
%% WJH added 16 Jan 2002
  \@for\@citeb:=#2\do{%
    \edef\@citeb{\expandafter\@firstofone\@citeb:\articlename}%
    \if@filesw\immediate\write\@auxout{\string\citation{\@citeb}}\fi}%
%% to here
  \def\@citea{}%
  \@cite{\@for\@citeb:=#2\do
    {%
      \edef\@citeb{\expandafter\@firstofone\@citeb:\articlename}%
      \@citea\def\@citea{;\ }\@ifundefined
      {b@\@citeb}{{\bf ?}\@warning {Citation `\@citeb' on page \thepage \space undefined}}%
      {\csname b@\@citeb\endcsname\@citebro\csname by@\@citeb\endcsname\@citebrc}}}%
  {#1}}

\def\cite{\@ifnextchar [{\let\@citebro=\space
                          \let\@citebrc=\relax
                          \let\@cite=\@bcite\@tempswatrue\@citex}
                        {\let\@citebro=\space
                          \let\@citebrc=\relax
                          \let\@cite=\@bcite\@tempswafalse\@citex[]}}
\def\pcite{\@ifnextchar [{\let\@citebro=\space
                          \let\@citebrc=\relax
                          \let\@cite=\@pcite\@tempswatrue\@citex}
                        {\let\@citebro=\space
                          \let\@citebrc=\relax
                          \let\@cite=\@pcite\@tempswafalse\@citex[]}}
\def\scite{\@ifnextchar [{\let\@citebro=\@brcite
                          \let\@citebrc=\@ebrcite
                          \let\@cite=\@pcite\@tempswatrue\@citex}
                        {\let\@citebro=\@brcite
                          \let\@citebrc=\@ebrcite
                          \let\@cite=\@pcite\@tempswafalse\@citex[]}}

\renewenvironment{thebibliography}
  {\section*{REFERENCES}\vspace*{\afterbiblioskip}\small
    \if@book\RereadBiblio\fi
    \sloppy\clubpenalty4000\widowpenalty4000
    \sfcode`\.=1000\relax
    \setlength\parindent{\z@}\setlength\parskip{\z@}}
  {\relax\normalsize\@outputaddressesifnecessary} 
%</rmaacite>
%    \end{macrocode}

% \subsection{Auxilliary files for generating ADS data}

% The ADS data are written by \LaTeX{} to a |.ads| file. This file
% then has to be included in the following file (|ADS.tex|),
% which is run through \LaTeX{}2HTML (see below). 

%    \begin{macrocode}
%<*ads-container>
%%
%% Instructions:
%% Edit the \input{} line below to reflect the name of your file. 
\documentclass{article}
\usepackage[latin1]{inputenc}
\usepackage{adshtml}
\newcommand{\ADSfield}[2]
{\%#1  #2\\}
%%\newcommand{\ADSvfield}[2]
%%{\%#1 \begin{verbatim}#2\end{verbatim}}
\newenvironment{ADSarticle}{}{\par\bigskip}
\renewcommand{\bibitem}[2][]{\par}
\setlength{\parindent}{0pt}
\renewcommand{\ThisVolume}[1]{%  
  2001, \RMAASC, 12, \oldpageref{#1:RMAAFirstPage}}  
 

\begin{document}
\makeatletter
%% Demonstration file for formatting an entire volume
\documentclass[book,journal,twocolumn]{rmaa}
%%%%%%%%%%%%%%%%%%%%%%%%%%%%%%%%%%%%%%%%%%%%%%%%%%%%%%%%%%%%%%%%
%%%%%%%%%%%%%%%%            PREAMBLE            %%%%%%%%%%%%%%%%
%%%%%%%%%%%%%%%%%%%%%%%%%%%%%%%%%%%%%%%%%%%%%%%%%%%%%%%%%%%%%%%%

\usepackage{paralist}
\usepackage{psfrag,color}
% Needed for landscape tables
\usepackage{pdflscape}

% Used in rm-minimal-test
\usepackage{kantlipsum}

% publication details of volume
\SetVolume{9}
\SetYear{2000}
\SetMonthNumeric{06}
\SetMes{junio}
% details of meeting
\SetConfTitle{Astrophysical Plasmas: Codes, Models, \& Observations}
\SetReunion{M\'exico, D. F.}{octubre 25--29, 1999}
\SetEditors{S. Jane Arthur, Nancy S. Brickhouse, \& Jos\'e Franco}

% this writes the .idx file that is used to generate the index 
\makeindex

%% This command makes the session headings come out bold in the
%% Table-of-Contents. Comment it out if you don't like it.
\BoldSessionTitles

%% This command makes a title page for each session
%% Comment it out if you don't like it.
\SessionTitlePages

%% Use this command if you want the abstract-only papers to be
%% indexed by ADS
\abstractsADS

%% At the special request of Alberto Carramin~ana, the following is a
%% way to suppress the author adresses at the end of each article
%\suppressfulladdresses

%% Use this command to add Received/accepted dates for each paper
% \printdates


\begin{document}

%%%%%%%%%%%%%%%%%%%%%%%%%%%%%%%%%%%%%%%%%%%%%%%%%%%%%%%%%%%%%%%%b
%%%%%%%%%%%%%%%%          FRONT MATTER          %%%%%%%%%%%%%%%%
%%%%%%%%%%%%%%%%%%%%%%%%%%%%%%%%%%%%%%%%%%%%%%%%%%%%%%%%%%%%%%%%
\frontmatter % start roman numbering

% format the TOC
\renewcommand{\RMAAbannerimage}{rmsc_blank}
\tableofcontents

%% preface, group photo, list of participants go here
%% use the `frontsection' environment (see editors' guide)
\begin{frontsection}{Preface}{%
    \listofauthors{S.~J.~Arthur, N.~S.~Brickhouse, J.~Franco}
    \indexauthor{Arthur, S. J.}
    \indexauthor{Brickhouse, N.}
    \indexauthor{Franco, J.}
    }
  These are the proceedings of a very nice conference that Jane and
  Pepe organised. This is the preface, which explains what a wonderful
  time was had by all and how grateful they are for all the people who
  gave them money or helped out.
  
  This is another paragraph of the preface, in case they have anything
  else to say.  For example, the Scientific Organizing Committee
  consisted of
  \begin{quote}
    J. Franco (Mexico, Chair), D. Cox (USA) and various other notables
  \end{quote}
  while the Local Organizing Committee comprised
  \begin{quote}
    S. J. Arthur (Chair), S. Kurtz, etc.
  \end{quote}

  \EditorSignature

\end{frontsection}

\begin{frontsection}[Photo]{Group Photograph}{\notitle}
\null\vfill
\centerline{Here goes the photograph.}
\vfill
\end{frontsection}

\begin{frontsection}[Participants]{List of Participants}{}
%% Here you put a formatted list of participants
\vfil
\dotfill First page of list of participants \dotfill
\vfil\clearpage

\null\vfil
\dotfill Second page of list of participants \dotfill
\vfil

\end{frontsection}

\begin{frontsection}[XXX]{Another front section}{}

\end{frontsection}

%%%%%%%%%%%%%%%%%%%%%%%%%%%%%%%%%%%%%%%%%%%%%%%%%%%%%%%%%%%%%%%%
%%%%%%%%%%%%%%%%          MAIN MATTER           %%%%%%%%%%%%%%%%
%%%%%%%%%%%%%%%%%%%%%%%%%%%%%%%%%%%%%%%%%%%%%%%%%%%%%%%%%%%%%%%%
\mainmatter % start arabic numbering

\begin{articles}
  %% Main articles go here
  %% Use \article{ARTICLENAME} for each individual article
  %% If you want any text to appear in brackets after the article
  %% title in the TOC and index, then put it as an optional first
  %% argument, e.g.,  \article[Invited Talk]{ARTICLENAME} 

  \session{The Contributed Papers}
  \article{rm-journal-example} % An illustrative article
  \article{rm-onepage} % another article, which has a plate
  \article{rm-minimal-test}
  \session{Conference Summary}
  \article[Invited Review]{rmsummsamp} % An illustrative conference summary

\end{articles}

%%%%%%%%%%%%%%%%%%%%%%%%%%%%%%%%%%%%%%%%%%%%%%%%%%%%%%%%%%%%%%%%
%%%%%%%%%%%%%%%%          BACK MATTER           %%%%%%%%%%%%%%%%
%%%%%%%%%%%%%%%%%%%%%%%%%%%%%%%%%%%%%%%%%%%%%%%%%%%%%%%%%%%%%%%%
\backmatter

%% Abstracts of presentations for which the
%% authors have not submitted a paper
\begin{abstracts}

  \title{An abstract of a poster}
  \author{W. J. Henney\altaffilmark{1}}
  \listofauthors{W.~J.~Henney}
  \indexauthor{Henney, W. J}
  \altaffiltext{1}{Where I am from done with altaffilmark/text}
  \abstract{This is a test abstract.   This is a test abstract.   This
    is a test abstract.   This is a test abstract.}
  

  \title{An abstract of another  poster}
  \author{W.~J.~Henney\altaffil{This is done with altaffil} and
    A.~N.~Other\altaffil{So is this} }
  \listofauthors{W.~J.~Henney \& A.~N.~Other}
  \indexauthor{Henney, W. J}
  \indexauthor{Other, A. N.}
  \abstract{This is a test abstract.   This is a test abstract.   This
    is a test abstract.   This is a test abstract.}

  \title{An abstract of yet another  poster}
  \author{W.~J.~Henney\altaffil{This is done with altaffil} and
    A.~N.~Other\altaffil{So is this} }
  \listofauthors{W.~J.~Henney \& A.~N.~Other}
  \indexauthor{Henney, W. J}
  \indexauthor{Other, A. N.}
  \abstract{This is a test abstract.   This is a test abstract.   This
    is a test abstract.   This is a test abstract.}

  \title{The abstract of a poster with a very long title to see how
    well it gets split up into lines}
  \author{A.~N.~Other\altaffil{The affiliation of A.~N.~Other is also
      very long so that it will go over several lines. His email
      address is anon@nowhere.org}}
  \listofauthors{A.~N.~Other}
  \indexauthor{Other, A. N.}
  \abstract{We will try and make this abstract longer so there is more
    stuff to put on the page and it doesn't look too empty. However, I
    have nothing whatsoever to say so will have to stop here. Let's
    repeat that a few times. We will try and make this abstract longer
    so there is more stuff to put on the page and it doesn't look too
    empty. However, I have nothing whatsoever to say so will have to
    stop here. Let's repeat that a few times. We will try and make
    this abstract longer so there is more stuff to put on the page and
    it doesn't look too empty. However, I have nothing whatsoever to
    say so will have to stop here. Let's repeat that a few times. We
    will try and make this abstract longer so there is more stuff to
    put on the page and it doesn't look too empty. However, I have
    nothing whatsoever to say so will have to stop here. Let's repeat
    that a few times.}
    
  \title[Comment to go in TOC]{The abstract of a second poster with a
    very long title to see how well it gets split up into lines}
  \author{A.~N.~Other\altaffil{The affiliation of A.~N.~Other is also
      very long so that it will go over several lines. His email
      address is anon@nowhere.org}}
  \listofauthors{A.~N.~Other}
  \indexauthor{Other, A. N.}
  \abstract{We will try and make this abstract longer so there is more
    stuff to put on the page and it doesn't look too empty. However, I
    have nothing whatsoever to say so will have to stop here. Let's
    repeat that a few times. We will try and make this abstract longer
    so there is more stuff to put on the page and it doesn't look too
    empty. However, I have nothing whatsoever to say so will have to
    stop here. Let's repeat that a few times. We will try and make
    this abstract longer so there is more stuff to put on the page and
    it doesn't look too empty. However, I have nothing whatsoever to
    say so will have to stop here. Let's repeat that a few times. We
    will try and make this abstract longer so there is more stuff to
    put on the page and it doesn't look too empty. However, I have
    nothing whatsoever to say so will have to stop here. Let's repeat
    that a few times.}
    
  \title{The abstract of a third poster with a very long title to see how
    well it gets split up into lines}
  \author{A.~N.~Other\altaffil{The affiliation of A.~N.~Other is also
      very long so that it will go over several lines. His email
      address is anon@nowhere.org}}
  \listofauthors{A.~N.~Other}
  \indexauthor{Other, A. N.}
  \abstract{We will try and make this abstract longer so there is more
    stuff to put on the page and it doesn't look too empty. However, I
    have nothing whatsoever to say so will have to stop here. Let's
    repeat that a few times. We will try and make this abstract longer
    so there is more stuff to put on the page and it doesn't look too
    empty. However, I have nothing whatsoever to say so will have to
    stop here. Let's repeat that a few times. We will try and make
    this abstract longer so there is more stuff to put on the page and
    it doesn't look too empty. However, I have nothing whatsoever to
    say so will have to stop here. Let's repeat that a few times. We
    will try and make this abstract longer so there is more stuff to
    put on the page and it doesn't look too empty. However, I have
    nothing whatsoever to say so will have to stop here. Let's repeat
    that a few times.}
    
  \title{The abstract of a fourth poster with a very long title to see how
    well it gets split up into lines}
  \author{A.~N.~Other\altaffil{The affiliation of A.~N.~Other is also
      very long so that it will go over several lines. His email
      address is anon@nowhere.org}}
  \listofauthors{A.~N.~Other}
  \indexauthor{Other, A. N.}
  \abstract{We will try and make this abstract longer so there is more
    stuff to put on the page and it doesn't look too empty. However, I
    have nothing whatsoever to say so will have to stop here. Let's
    repeat that a few times. We will try and make this abstract longer
    so there is more stuff to put on the page and it doesn't look too
    empty. However, I have nothing whatsoever to say so will have to
    stop here. Let's repeat that a few times. We will try and make
    this abstract longer so there is more stuff to put on the page and
    it doesn't look too empty. However, I have nothing whatsoever to
    say so will have to stop here. Let's repeat that a few times. We
    will try and make this abstract longer so there is more stuff to
    put on the page and it doesn't look too empty. However, I have
    nothing whatsoever to say so will have to stop here. Let's repeat
    that a few times.}
    
  \title{The abstract of a fifth poster with a very long title to see how
    well it gets split up into lines}
  \author{A.~N.~Other\altaffil{The affiliation of A.~N.~Other is also
      very long so that it will go over several lines. His email
      address is anon@nowhere.org}}
  \listofauthors{A.~N.~Other}
  \indexauthor{Other, A. N.}
  \abstract{We will try and make this abstract longer so there is more
    stuff to put on the page and it doesn't look too empty. However, I
    have nothing whatsoever to say so will have to stop here. Let's
    repeat that a few times. We will try and make this abstract longer
    so there is more stuff to put on the page and it doesn't look too
    empty. However, I have nothing whatsoever to say so will have to
    stop here. Let's repeat that a few times. We will try and make
    this abstract longer so there is more stuff to put on the page and
    it doesn't look too empty. However, I have nothing whatsoever to
    say so will have to stop here. Let's repeat that a few times. We
    will try and make this abstract longer so there is more stuff to
    put on the page and it doesn't look too empty. However, I have
    nothing whatsoever to say so will have to stop here. Let's repeat
    that a few times.}
    
  \title{The abstract of a sixth poster with a very long title to see how
    well it gets split up into lines}
  \author{A.~N.~Other\altaffil{The affiliation of A.~N.~Other is also
      very long so that it will go over several lines. His email
      address is anon@nowhere.org}}
  \listofauthors{A.~N.~Other}
  \indexauthor{Other, A. N.}
  \abstract{We will try and make this abstract longer so there is more
    stuff to put on the page and it doesn't look too empty. However, I
    have nothing whatsoever to say so will have to stop here. Let's
    repeat that a few times. We will try and make this abstract longer
    so there is more stuff to put on the page and it doesn't look too
    empty. However, I have nothing whatsoever to say so will have to
    stop here. Let's repeat that a few times. We will try and make
    this abstract longer so there is more stuff to put on the page and
    it doesn't look too empty. However, I have nothing whatsoever to
    say so will have to stop here. Let's repeat that a few times. We
    will try and make this abstract longer so there is more stuff to
    put on the page and it doesn't look too empty. However, I have
    nothing whatsoever to say so will have to stop here. Let's repeat
    that a few times.}
    
  \title{The abstract of a seventh poster with a very long title to see how
    well it gets split up into lines}
  \author{A.~N.~Other\altaffil{The affiliation of A.~N.~Other is also
      very long so that it will go over several lines. His email
      address is anon@nowhere.org}}
  \listofauthors{A.~N.~Other}
  \indexauthor{Other, A. N.}
  \abstract{We will try and make this abstract longer so there is more
    stuff to put on the page and it doesn't look too empty. However, I
    have nothing whatsoever to say so will have to stop here. Let's
    repeat that a few times. We will try and make this abstract longer
    so there is more stuff to put on the page and it doesn't look too
    empty. However, I have nothing whatsoever to say so will have to
    stop here. Let's repeat that a few times. We will try and make
    this abstract longer so there is more stuff to put on the page and
    it doesn't look too empty. However, I have nothing whatsoever to
    say so will have to stop here. Let's repeat that a few times. We
    will try and make this abstract longer so there is more stuff to
    put on the page and it doesn't look too empty. However, I have
    nothing whatsoever to say so will have to stop here. Let's repeat
    that a few times.}
  
  \title{An abstract to go in the list of abstracts, but not in the
    index of authors}
  \author{I.~A.~M. Hidden\altaffil{Puerto Escondido}}
  \listofauthors{I.~A.~M. Hidden}
  \abstract{The idea is that by omitting the
    \textbackslash{}indexauthor command, we can avoid an entry in the
    author index.}

\end{abstracts}


%% The Index: after latexing twice, run `makeindex -s rmaa' on this
%% file, then latex twice more
\printindex

%% if you rename this file to something other than `rmbooksamp', then
%% you must change the above line accordingly

%%
%% Plates: If you have any, then the following line will print
%% them out:

\outputplates

\end{document}

%%% Local Variables:
%%% mode: latex
%%% TeX-master: t
%%% End:

\end{document}
%</ads-container>
%    \end{macrocode}

% This is the |adshtml.sty| style file used by |ADS.tex|. It just
% defines a few 
% macros that might be used in abstracts, etc (really it should define
% them all). This now gets all its code from
% Sec~\ref{sec:user-commands}. 
%    \begin{macrocode}
%<*adshtml-style>
%% Nothing else needs to go here now. 
%</adshtml-style>
%    \end{macrocode}
% The init file for latex2html.
%    \begin{macrocode}
%<*latex2html-init>
$PDFTOCAIRO = "pdftocairo";
$PDFLATEX = "pdflatex";
$LATEX = "latex";
$MAX_SPLIT_DEPTH = 0;
$HTML_VERSION = "4.0,math";
$NO_SIMPLE_MATH = 1;
$ADDRESS = "<p><address><a href=\"mailto:rmaa\@astro.unam.mx\">rmaa\@astro.unam.mx</address></p>"; 
$INFO = "<p>This document was created using latex2html and the rmaa latex macros</p>"; 
$NO_NAVIGATION = 1;
1; 
%</latex2html-init>
%    \end{macrocode}
% This is the |adshtml.perl| file that is used by \LaTeX{}2HTML. Some
% bizarre \TeX{}y command sometimes show up so we tell it to ignore
% these. All the astronomy macros that come from AAS are processed in
% \TeX{}. That way they get made to images but with an |ALT| tag that
% has the original code, which is what will be picked up by |lynx|,
% with the result that they get passed intact into the final |.txt|
% file. 
%    \begin{macrocode}
%<*adshtml-perl>
package main;

$dum = &do_cmd_makeatletter();

sub do_cmd_RMAA {
    local($_) = @_;
    "RMxAA" . $_;
}

sub do_cmd_RMAASC {
    local($_) = @_;
    "RMxAASC" . $_;
}

sub do_cmd_etal {
    local($_) = @_;
    "et al\." . $_;
}

sub do_cmd_ibidrule {
    local($_) = @_;
    "<img src=\"foo\" alt=\"\ibidrule\">" . $_;
}


sub do_cmd_bibitem {
    local($_) = @_;
    local($key);
    local($opt, $dummy)=&get_next_optional_argument; 
    $key = &missing_braces unless
        ((s/$next_pair_pr_rx/$key = $2; ''/eo)
        ||(s/$next_pair_rx/$key = $2; ''/eo));
    "<p>" . $_; 
}

&ignore_commands (<<_RAW_ARG_CMDS_);
lowercase
adjustfinalcols
textsuperscript # {}
newline # []
thinmuskip
mskip
smallskip
medskip
bigskip
hspace # {}
hspacestar # {}
vspace # {}
vspacestar # {}
mathrm
_at_mathrm
spacefactor
discretionary
unhbox 
voidb_at_x 
penalty 
_at_M
_at_m
null
IeC
_RAW_ARG_CMDS_

&process_commands_in_tex (<<_RAW_ARG_CMDS_);
sun
earth
lesssim
gtrsim
sq
argdeg
arcmin
arcsec
fd
fh
fm
fs
fdg
farcm
farcs
fp
micron
la
ga
case # {} # {}
slantfrac # {} # {}
onehalf
onethird
twothirds
onequarter
threequarters
ubvr
ub
bv
vr
ur
aj
araa
apj
apjl
apjs
ao
apss
aap
aapr
aaps
azh
baas
jrasc
memras
mnras
pra
prb
prc
prd
pre
prl
pasp
pasj
qjras
skytel
solphys
sovast
ssr
zap
nat
iaucirc
aplett
apspr
bain
fcp
gca
grl
jcp
jgr
jqsrt
memsai
nphysa
physrep
physscr
planss
procspie
_RAW_ARG_CMDS_

1;                              # This must be the last line
%</adshtml-perl> 
%    \end{macrocode}

% This is a shell script |makeads.sh| that generates the |_ads.txt| file from
% |_ads.tex| by using \LaTeX{}2HTML and |lynx|.  \changes{4.0}{2019/11/13}{Make sure that lynx writes utf-8 output instead of latin-1.}
% \changes{2.2beta13}{2002/03/12}{Complete rewrite of |makeads.sh|. Now much 
% improved I think}
% \changes{4.0}{2019/11/13}{Fix the at-sign substitution so that it only eliminates \cs{@} and does not mess with email addresses.}

%    \begin{macrocode}
%<*ads-shell-script>
#! /bin/sh
if [ -z $1 ]; then
    echo "Usage: $0 basename"
    exit
fi
BASENAME=$1
ADSBASENAME=${BASENAME}_ads
AUXFILE=$BASENAME.aux
TMPFILE=tmp.$$


echo "Generating $ADSBASENAME.tex"

# For some reason, must double-up \\ to \\\\ but not \ to \\ - odd!
cat > $ADSBASENAME.tex <<EOF
%% PLEASE DO NOT EDIT THIS FILE!  It is generated automatically from
%% the script makeads.sh.  Any hand-edited changes will be lost
%% 
\documentclass{article}
\usepackage[latin1]{inputenc}
\usepackage{html}
\usepackage{adshtml}
\newcommand{\ADSfield}[2]{\%#1  #2\\\\}
\newenvironment{ADSarticle}{}{\par\bigskip}
\setlength{\parindent}{0pt}
\newcounter{istage}
\renewcommand{\ion}[2]{\setcounter{istage}{#2}#1 \Roman{istage}}
\begin{document}
\makeatletter
EOF
echo >> $ADSBASENAME.tex "\input{$BASENAME.ads2}" 
echo >> $ADSBASENAME.tex "\end{document}"



echo "Generating $ADSBASENAME.sed" 

# Find type of journal
grep '\\documentclass\[.*proceedings.*\]' $BASENAME.tex >/dev/null \
     && JOURNAL='RMxAC' \
     || JOURNAL='RMxAA'

# Find the year and volume number
YEAR=`grep '\\SetYear' $BASENAME.tex | sed 's/.*\\SetYear{\(.*\)}.*/\1/'`
VOLUME=`grep '\\SetVolume' $BASENAME.tex | sed 's/.*\\SetVolume{\(.*\)}.*/\1/'`

# Generate an auxilliary sed script 
cat > $TMPFILE <<EOF
s|\\\\newlabel{\(.*\):RMAAFirstPage}{{.*}{\(.*\)}}|s/\\\\\\\\ThisVolume{\1}/$YEAR, RMxAC, $VOLUME, \2/g|g
EOF

# Now go through the .aux file, looking for \newlabels and use the sed 
# script just generated in order to make another sed script, one that 
# will replace \ThisVolume{LABEL} with the correct reference
grep '^\\newlabel{.*:RMAAFirstPage}' $AUXFILE \
     | sed -f $TMPFILE > $ADSBASENAME.sed && rm $TMPFILE

# Apend a further sed command to get rid of \@, since redefining 
# it in _ads.tex didn't work
echo >> $ADSBASENAME.sed 's/\\\\@//g'

echo "Generating $BASENAME.ads2 from $ADSBASENAME.sed and $BASENAME.ads" 

# Apply this to the .ads file, generating the .ads2 file, which is included
# by the _ads.tex file.
sed -f $ADSBASENAME.sed $BASENAME.ads > $BASENAME.ads2

echo "Generating $ADSBASENAME.txt from $ADSBASENAME.tex and $BASENAME.ads2"

# Now run latex2html on the _ads.tex file
latex2html -no_navigation -info 0 -address 0 -html_version 4.0,math $ADSBASENAME

# Then, convert to utf8 with lynx and tidy it up a bit
lynx --justify=off --width=100 --dump --nolist --display_charset=utf-8 \
     $ADSBASENAME/$ADSBASENAME.html | cut -c4- \
     | sed 's/ibidrule/\\ibidrule/g' > $ADSBASENAME.txt
%</ads-shell-script>
%    \end{macrocode}

% \section{Epilogue}
% Err, that's it \dots

\endinput


% \endinput
% Local Variables:
% mode: doctex
% TeX-master: t
% End:
"

rmaa.ind: rmaa.idx
        @echo "REGENERATING INDEX"
        makeindex -s gind.ist rmaa

rmaa.gls: rmaa.glo
        @echo "REGENERATING CHANGELOG"
        makeindex -s gglo.ist -o rmaa.gls rmaa.glo

clean: 
        rm -f rmaa.aux rmaa.dvi rmaa.ind rmaa.gls
%</makefile>
%    \end{macrocode}

% The |rmaa.ins| file uses the \textsf{docstrip} package to split out
% all the files and strip off this documentation.
%    \begin{macrocode}
%<*ins>
%%
%% Installation batch file for the RMAA LaTeX package
%%
\def\batchfile{rmaa.ins}
\input docstrip.tex
\preamble

PLEASE DO NOT EDIT THIS FILE! It is generated automatically from
the file rmaa.dtx using the accompanying Makefile and rmaa.ins 
Any hand-edited changes will be lost

A LaTeX document class for the RMxAA

Author: William Henney (w.henney@crya.unam.mx)

Latest version available from: http://www.crya.unam.mx/rmaa

Copyright (C) 1999-2010 Instituto de Astronomia, UNAM, Mexico

This program is free software; you can redistribute it and/or
modify it under the terms of the GNU General Public License
as published by the Free Software Foundation; either version 2
of the License, or (at your option) any later version.

This program is distributed in the hope that it will be useful,
but WITHOUT ANY WARRANTY; without even the implied warranty of
MERCHANTABILITY or FITNESS FOR A PARTICULAR PURPOSE.  See the
GNU General Public License for more details.

You should have received a copy of the GNU General Public License
along with this program; if not, write to the Free Software
Foundation, Inc., 59 Temple Place - Suite 330, Boston, MA  02111-1307, USA.
\endpreamble
%    \end{macrocode}
% Clobber files without asking
%    \begin{macrocode}
\askforoverwritefalse
%    \end{macrocode}
% Make the \LaTeX{} macro files (class and styles)
%    \begin{macrocode}
\generate{% 
  \file{rmaa.cls}{\from{rmaa.dtx}{package}}
  \file{rmaa.ins}{\from{rmaa.dtx}{ins}}
  \file{rmaacite.sty}{\from{rmaa.dtx}{rmaacite}}
  \file{adshtml.sty}{\from{rmaa.dtx}{adshtml-style}}
  \file{ADS.tex}{\from{rmaa.dtx}{ads-container}}
  }
%    \end{macrocode}
% Make the other non-\LaTeX{} files. These are done separately since
% they can't use `\%\%' as a comment character for the
% preamble/postamble (and Perl files can't have the postamble I don't
% think). 
%    \begin{macrocode}
\generate{\nopreamble\nopostamble% 
  \file{Makefile}{\from{rmaa.dtx}{makefile}}
  \file{adshtml.perl}{\from{rmaa.dtx}{adshtml-perl}}
  \file{makeads.sh}{\from{rmaa.dtx}{ads-shell-script}}
  \file{sample.latex2html-init}{\from{rmaa.dtx}{latex2html-init}}
  }

%% Local Variables: 
%% mode: LaTeX
%% mode: Font-lock
%% End:
%</ins>
%    \end{macrocode}

% \subsection{Extended citation syntax}
% Modified from a file called mnrascite.sty
% I no longer remember where I stole the original version of this from
% 
% Alterations to the cite commands to make them do more sensible things.
% N.B. the bibliography style should produce entries which have labels where
% the date is delimited by angle brackets '<>' for these macros to work
% properly. They will not give very sensible error messages if you do not do 
% this properly.
% 
% A new command |\scite{key}| is defined which produces a citation which is 
% suitable for insertion in a sentence. i.e., it does not put brackets around 
% the whole citation, only the date. Also |\pcite{key}| produce output with no
% parentheses whatsoever. 
% 
% Should redefine |\@bibitem| really to have any hope of compatibility
% the command |\lbibcite| is defined to deal with the citations.
%    \begin{macrocode}
%<*rmaacite>
\def\bibitem[#1<#2>]#3{\par\hangindent=1.5em\hangafter=1
  \if@filesw
  {\def\protect##1{\string ##1\space}
    \immediate\write\@auxout{\string\lbibcite{#3}{#1}{#2}}}
  \fi}
\@redefinebibitem % what it says
%    \end{macrocode}
% Note that the |\lbibcite| defines two types of citation one contains the names,
% the other the year
%    \begin{macrocode}
\def\lbibcite#1#2#3{\global\@namedef{b@#1}{#2}
                    \global\@namedef{by@#1}{#3}}

\def\@bcite#1#2{({#1\if@tempswa , #2\fi})}
\def\@pcite#1#2{#1\if@tempswa , #2\fi}
%    \end{macrocode}
% This declaration contains what is put between the name and the date
% \changes{3.11}{2003/01/31}{Moved the \cs{articlename} to after the key for consistency
% with the new natbib stuff.}
%    \begin{macrocode}
\def\@brcite{\ (}
\def\@ebrcite{)}
\def\@fmttop#1{\@fmtcite[\csname #1\endcsname]}
\def\@citex[#1]#2{%
%%  \if@filesw\immediate\write\@auxout{\string\citation{#2}}\fi
%% WJH added 16 Jan 2002
  \@for\@citeb:=#2\do{%
    \edef\@citeb{\expandafter\@firstofone\@citeb:\articlename}%
    \if@filesw\immediate\write\@auxout{\string\citation{\@citeb}}\fi}%
%% to here
  \def\@citea{}%
  \@cite{\@for\@citeb:=#2\do
    {%
      \edef\@citeb{\expandafter\@firstofone\@citeb:\articlename}%
      \@citea\def\@citea{;\ }\@ifundefined
      {b@\@citeb}{{\bf ?}\@warning {Citation `\@citeb' on page \thepage \space undefined}}%
      {\csname b@\@citeb\endcsname\@citebro\csname by@\@citeb\endcsname\@citebrc}}}%
  {#1}}

\def\cite{\@ifnextchar [{\let\@citebro=\space
                          \let\@citebrc=\relax
                          \let\@cite=\@bcite\@tempswatrue\@citex}
                        {\let\@citebro=\space
                          \let\@citebrc=\relax
                          \let\@cite=\@bcite\@tempswafalse\@citex[]}}
\def\pcite{\@ifnextchar [{\let\@citebro=\space
                          \let\@citebrc=\relax
                          \let\@cite=\@pcite\@tempswatrue\@citex}
                        {\let\@citebro=\space
                          \let\@citebrc=\relax
                          \let\@cite=\@pcite\@tempswafalse\@citex[]}}
\def\scite{\@ifnextchar [{\let\@citebro=\@brcite
                          \let\@citebrc=\@ebrcite
                          \let\@cite=\@pcite\@tempswatrue\@citex}
                        {\let\@citebro=\@brcite
                          \let\@citebrc=\@ebrcite
                          \let\@cite=\@pcite\@tempswafalse\@citex[]}}

\renewenvironment{thebibliography}
  {\section*{REFERENCES}\vspace*{\afterbiblioskip}\small
    \if@book\RereadBiblio\fi
    \sloppy\clubpenalty4000\widowpenalty4000
    \sfcode`\.=1000\relax
    \setlength\parindent{\z@}\setlength\parskip{\z@}}
  {\relax\normalsize\@outputaddressesifnecessary} 
%</rmaacite>
%    \end{macrocode}

% \subsection{Auxilliary files for generating ADS data}

% The ADS data are written by \LaTeX{} to a |.ads| file. This file
% then has to be included in the following file (|ADS.tex|),
% which is run through \LaTeX{}2HTML (see below). 

%    \begin{macrocode}
%<*ads-container>
%%
%% Instructions:
%% Edit the \input{} line below to reflect the name of your file. 
\documentclass{article}
\usepackage[latin1]{inputenc}
\usepackage{adshtml}
\newcommand{\ADSfield}[2]
{\%#1  #2\\}
%%\newcommand{\ADSvfield}[2]
%%{\%#1 \begin{verbatim}#2\end{verbatim}}
\newenvironment{ADSarticle}{}{\par\bigskip}
\renewcommand{\bibitem}[2][]{\par}
\setlength{\parindent}{0pt}
\renewcommand{\ThisVolume}[1]{%  
  2001, \RMAASC, 12, \oldpageref{#1:RMAAFirstPage}}  
 

\begin{document}
\makeatletter
%% Demonstration file for formatting an entire volume
\documentclass[book,journal,twocolumn]{rmaa}
%%%%%%%%%%%%%%%%%%%%%%%%%%%%%%%%%%%%%%%%%%%%%%%%%%%%%%%%%%%%%%%%
%%%%%%%%%%%%%%%%            PREAMBLE            %%%%%%%%%%%%%%%%
%%%%%%%%%%%%%%%%%%%%%%%%%%%%%%%%%%%%%%%%%%%%%%%%%%%%%%%%%%%%%%%%

\usepackage{paralist}
\usepackage{psfrag,color}
% Needed for landscape tables
\usepackage{pdflscape}

% Used in rm-minimal-test
\usepackage{kantlipsum}

% publication details of volume
\SetVolume{9}
\SetYear{2000}
\SetMonthNumeric{06}
\SetMes{junio}
% details of meeting
\SetConfTitle{Astrophysical Plasmas: Codes, Models, \& Observations}
\SetReunion{M\'exico, D. F.}{octubre 25--29, 1999}
\SetEditors{S. Jane Arthur, Nancy S. Brickhouse, \& Jos\'e Franco}

% this writes the .idx file that is used to generate the index 
\makeindex

%% This command makes the session headings come out bold in the
%% Table-of-Contents. Comment it out if you don't like it.
\BoldSessionTitles

%% This command makes a title page for each session
%% Comment it out if you don't like it.
\SessionTitlePages

%% Use this command if you want the abstract-only papers to be
%% indexed by ADS
\abstractsADS

%% At the special request of Alberto Carramin~ana, the following is a
%% way to suppress the author adresses at the end of each article
%\suppressfulladdresses

%% Use this command to add Received/accepted dates for each paper
% \printdates


\begin{document}

%%%%%%%%%%%%%%%%%%%%%%%%%%%%%%%%%%%%%%%%%%%%%%%%%%%%%%%%%%%%%%%%b
%%%%%%%%%%%%%%%%          FRONT MATTER          %%%%%%%%%%%%%%%%
%%%%%%%%%%%%%%%%%%%%%%%%%%%%%%%%%%%%%%%%%%%%%%%%%%%%%%%%%%%%%%%%
\frontmatter % start roman numbering

% format the TOC
\renewcommand{\RMAAbannerimage}{rmsc_blank}
\tableofcontents

%% preface, group photo, list of participants go here
%% use the `frontsection' environment (see editors' guide)
\begin{frontsection}{Preface}{%
    \listofauthors{S.~J.~Arthur, N.~S.~Brickhouse, J.~Franco}
    \indexauthor{Arthur, S. J.}
    \indexauthor{Brickhouse, N.}
    \indexauthor{Franco, J.}
    }
  These are the proceedings of a very nice conference that Jane and
  Pepe organised. This is the preface, which explains what a wonderful
  time was had by all and how grateful they are for all the people who
  gave them money or helped out.
  
  This is another paragraph of the preface, in case they have anything
  else to say.  For example, the Scientific Organizing Committee
  consisted of
  \begin{quote}
    J. Franco (Mexico, Chair), D. Cox (USA) and various other notables
  \end{quote}
  while the Local Organizing Committee comprised
  \begin{quote}
    S. J. Arthur (Chair), S. Kurtz, etc.
  \end{quote}

  \EditorSignature

\end{frontsection}

\begin{frontsection}[Photo]{Group Photograph}{\notitle}
\null\vfill
\centerline{Here goes the photograph.}
\vfill
\end{frontsection}

\begin{frontsection}[Participants]{List of Participants}{}
%% Here you put a formatted list of participants
\vfil
\dotfill First page of list of participants \dotfill
\vfil\clearpage

\null\vfil
\dotfill Second page of list of participants \dotfill
\vfil

\end{frontsection}

\begin{frontsection}[XXX]{Another front section}{}

\end{frontsection}

%%%%%%%%%%%%%%%%%%%%%%%%%%%%%%%%%%%%%%%%%%%%%%%%%%%%%%%%%%%%%%%%
%%%%%%%%%%%%%%%%          MAIN MATTER           %%%%%%%%%%%%%%%%
%%%%%%%%%%%%%%%%%%%%%%%%%%%%%%%%%%%%%%%%%%%%%%%%%%%%%%%%%%%%%%%%
\mainmatter % start arabic numbering

\begin{articles}
  %% Main articles go here
  %% Use \article{ARTICLENAME} for each individual article
  %% If you want any text to appear in brackets after the article
  %% title in the TOC and index, then put it as an optional first
  %% argument, e.g.,  \article[Invited Talk]{ARTICLENAME} 

  \session{The Contributed Papers}
  \article{rm-journal-example} % An illustrative article
  \article{rm-onepage} % another article, which has a plate
  \article{rm-minimal-test}
  \session{Conference Summary}
  \article[Invited Review]{rmsummsamp} % An illustrative conference summary

\end{articles}

%%%%%%%%%%%%%%%%%%%%%%%%%%%%%%%%%%%%%%%%%%%%%%%%%%%%%%%%%%%%%%%%
%%%%%%%%%%%%%%%%          BACK MATTER           %%%%%%%%%%%%%%%%
%%%%%%%%%%%%%%%%%%%%%%%%%%%%%%%%%%%%%%%%%%%%%%%%%%%%%%%%%%%%%%%%
\backmatter

%% Abstracts of presentations for which the
%% authors have not submitted a paper
\begin{abstracts}

  \title{An abstract of a poster}
  \author{W. J. Henney\altaffilmark{1}}
  \listofauthors{W.~J.~Henney}
  \indexauthor{Henney, W. J}
  \altaffiltext{1}{Where I am from done with altaffilmark/text}
  \abstract{This is a test abstract.   This is a test abstract.   This
    is a test abstract.   This is a test abstract.}
  

  \title{An abstract of another  poster}
  \author{W.~J.~Henney\altaffil{This is done with altaffil} and
    A.~N.~Other\altaffil{So is this} }
  \listofauthors{W.~J.~Henney \& A.~N.~Other}
  \indexauthor{Henney, W. J}
  \indexauthor{Other, A. N.}
  \abstract{This is a test abstract.   This is a test abstract.   This
    is a test abstract.   This is a test abstract.}

  \title{An abstract of yet another  poster}
  \author{W.~J.~Henney\altaffil{This is done with altaffil} and
    A.~N.~Other\altaffil{So is this} }
  \listofauthors{W.~J.~Henney \& A.~N.~Other}
  \indexauthor{Henney, W. J}
  \indexauthor{Other, A. N.}
  \abstract{This is a test abstract.   This is a test abstract.   This
    is a test abstract.   This is a test abstract.}

  \title{The abstract of a poster with a very long title to see how
    well it gets split up into lines}
  \author{A.~N.~Other\altaffil{The affiliation of A.~N.~Other is also
      very long so that it will go over several lines. His email
      address is anon@nowhere.org}}
  \listofauthors{A.~N.~Other}
  \indexauthor{Other, A. N.}
  \abstract{We will try and make this abstract longer so there is more
    stuff to put on the page and it doesn't look too empty. However, I
    have nothing whatsoever to say so will have to stop here. Let's
    repeat that a few times. We will try and make this abstract longer
    so there is more stuff to put on the page and it doesn't look too
    empty. However, I have nothing whatsoever to say so will have to
    stop here. Let's repeat that a few times. We will try and make
    this abstract longer so there is more stuff to put on the page and
    it doesn't look too empty. However, I have nothing whatsoever to
    say so will have to stop here. Let's repeat that a few times. We
    will try and make this abstract longer so there is more stuff to
    put on the page and it doesn't look too empty. However, I have
    nothing whatsoever to say so will have to stop here. Let's repeat
    that a few times.}
    
  \title[Comment to go in TOC]{The abstract of a second poster with a
    very long title to see how well it gets split up into lines}
  \author{A.~N.~Other\altaffil{The affiliation of A.~N.~Other is also
      very long so that it will go over several lines. His email
      address is anon@nowhere.org}}
  \listofauthors{A.~N.~Other}
  \indexauthor{Other, A. N.}
  \abstract{We will try and make this abstract longer so there is more
    stuff to put on the page and it doesn't look too empty. However, I
    have nothing whatsoever to say so will have to stop here. Let's
    repeat that a few times. We will try and make this abstract longer
    so there is more stuff to put on the page and it doesn't look too
    empty. However, I have nothing whatsoever to say so will have to
    stop here. Let's repeat that a few times. We will try and make
    this abstract longer so there is more stuff to put on the page and
    it doesn't look too empty. However, I have nothing whatsoever to
    say so will have to stop here. Let's repeat that a few times. We
    will try and make this abstract longer so there is more stuff to
    put on the page and it doesn't look too empty. However, I have
    nothing whatsoever to say so will have to stop here. Let's repeat
    that a few times.}
    
  \title{The abstract of a third poster with a very long title to see how
    well it gets split up into lines}
  \author{A.~N.~Other\altaffil{The affiliation of A.~N.~Other is also
      very long so that it will go over several lines. His email
      address is anon@nowhere.org}}
  \listofauthors{A.~N.~Other}
  \indexauthor{Other, A. N.}
  \abstract{We will try and make this abstract longer so there is more
    stuff to put on the page and it doesn't look too empty. However, I
    have nothing whatsoever to say so will have to stop here. Let's
    repeat that a few times. We will try and make this abstract longer
    so there is more stuff to put on the page and it doesn't look too
    empty. However, I have nothing whatsoever to say so will have to
    stop here. Let's repeat that a few times. We will try and make
    this abstract longer so there is more stuff to put on the page and
    it doesn't look too empty. However, I have nothing whatsoever to
    say so will have to stop here. Let's repeat that a few times. We
    will try and make this abstract longer so there is more stuff to
    put on the page and it doesn't look too empty. However, I have
    nothing whatsoever to say so will have to stop here. Let's repeat
    that a few times.}
    
  \title{The abstract of a fourth poster with a very long title to see how
    well it gets split up into lines}
  \author{A.~N.~Other\altaffil{The affiliation of A.~N.~Other is also
      very long so that it will go over several lines. His email
      address is anon@nowhere.org}}
  \listofauthors{A.~N.~Other}
  \indexauthor{Other, A. N.}
  \abstract{We will try and make this abstract longer so there is more
    stuff to put on the page and it doesn't look too empty. However, I
    have nothing whatsoever to say so will have to stop here. Let's
    repeat that a few times. We will try and make this abstract longer
    so there is more stuff to put on the page and it doesn't look too
    empty. However, I have nothing whatsoever to say so will have to
    stop here. Let's repeat that a few times. We will try and make
    this abstract longer so there is more stuff to put on the page and
    it doesn't look too empty. However, I have nothing whatsoever to
    say so will have to stop here. Let's repeat that a few times. We
    will try and make this abstract longer so there is more stuff to
    put on the page and it doesn't look too empty. However, I have
    nothing whatsoever to say so will have to stop here. Let's repeat
    that a few times.}
    
  \title{The abstract of a fifth poster with a very long title to see how
    well it gets split up into lines}
  \author{A.~N.~Other\altaffil{The affiliation of A.~N.~Other is also
      very long so that it will go over several lines. His email
      address is anon@nowhere.org}}
  \listofauthors{A.~N.~Other}
  \indexauthor{Other, A. N.}
  \abstract{We will try and make this abstract longer so there is more
    stuff to put on the page and it doesn't look too empty. However, I
    have nothing whatsoever to say so will have to stop here. Let's
    repeat that a few times. We will try and make this abstract longer
    so there is more stuff to put on the page and it doesn't look too
    empty. However, I have nothing whatsoever to say so will have to
    stop here. Let's repeat that a few times. We will try and make
    this abstract longer so there is more stuff to put on the page and
    it doesn't look too empty. However, I have nothing whatsoever to
    say so will have to stop here. Let's repeat that a few times. We
    will try and make this abstract longer so there is more stuff to
    put on the page and it doesn't look too empty. However, I have
    nothing whatsoever to say so will have to stop here. Let's repeat
    that a few times.}
    
  \title{The abstract of a sixth poster with a very long title to see how
    well it gets split up into lines}
  \author{A.~N.~Other\altaffil{The affiliation of A.~N.~Other is also
      very long so that it will go over several lines. His email
      address is anon@nowhere.org}}
  \listofauthors{A.~N.~Other}
  \indexauthor{Other, A. N.}
  \abstract{We will try and make this abstract longer so there is more
    stuff to put on the page and it doesn't look too empty. However, I
    have nothing whatsoever to say so will have to stop here. Let's
    repeat that a few times. We will try and make this abstract longer
    so there is more stuff to put on the page and it doesn't look too
    empty. However, I have nothing whatsoever to say so will have to
    stop here. Let's repeat that a few times. We will try and make
    this abstract longer so there is more stuff to put on the page and
    it doesn't look too empty. However, I have nothing whatsoever to
    say so will have to stop here. Let's repeat that a few times. We
    will try and make this abstract longer so there is more stuff to
    put on the page and it doesn't look too empty. However, I have
    nothing whatsoever to say so will have to stop here. Let's repeat
    that a few times.}
    
  \title{The abstract of a seventh poster with a very long title to see how
    well it gets split up into lines}
  \author{A.~N.~Other\altaffil{The affiliation of A.~N.~Other is also
      very long so that it will go over several lines. His email
      address is anon@nowhere.org}}
  \listofauthors{A.~N.~Other}
  \indexauthor{Other, A. N.}
  \abstract{We will try and make this abstract longer so there is more
    stuff to put on the page and it doesn't look too empty. However, I
    have nothing whatsoever to say so will have to stop here. Let's
    repeat that a few times. We will try and make this abstract longer
    so there is more stuff to put on the page and it doesn't look too
    empty. However, I have nothing whatsoever to say so will have to
    stop here. Let's repeat that a few times. We will try and make
    this abstract longer so there is more stuff to put on the page and
    it doesn't look too empty. However, I have nothing whatsoever to
    say so will have to stop here. Let's repeat that a few times. We
    will try and make this abstract longer so there is more stuff to
    put on the page and it doesn't look too empty. However, I have
    nothing whatsoever to say so will have to stop here. Let's repeat
    that a few times.}
  
  \title{An abstract to go in the list of abstracts, but not in the
    index of authors}
  \author{I.~A.~M. Hidden\altaffil{Puerto Escondido}}
  \listofauthors{I.~A.~M. Hidden}
  \abstract{The idea is that by omitting the
    \textbackslash{}indexauthor command, we can avoid an entry in the
    author index.}

\end{abstracts}


%% The Index: after latexing twice, run `makeindex -s rmaa' on this
%% file, then latex twice more
\printindex

%% if you rename this file to something other than `rmbooksamp', then
%% you must change the above line accordingly

%%
%% Plates: If you have any, then the following line will print
%% them out:

\outputplates

\end{document}

%%% Local Variables:
%%% mode: latex
%%% TeX-master: t
%%% End:

\end{document}
%</ads-container>
%    \end{macrocode}

% This is the |adshtml.sty| style file used by |ADS.tex|. It just
% defines a few 
% macros that might be used in abstracts, etc (really it should define
% them all). This now gets all its code from
% Sec~\ref{sec:user-commands}. 
%    \begin{macrocode}
%<*adshtml-style>
%% Nothing else needs to go here now. 
%</adshtml-style>
%    \end{macrocode}
% The init file for latex2html.
%    \begin{macrocode}
%<*latex2html-init>
$MAX_SPLIT_DEPTH = 0;
$HTML_VERSION = "4.0,math";
$NO_SIMPLE_MATH = 1;
$ADDRESS = "<p><address><a href=\"mailto:rmaa\@astroscu.unam.mx\">rmaa\@astroscu.unam.mx</address></p>"; 
$INFO = "<p>This document was created using latex2html and the rmaa latex macros</p>"; 
$NO_NAVIGATION = 1;
1; 
%</latex2html-init>
%    \end{macrocode}
% This is the |adshtml.perl| file that is used by \LaTeX{}2HTML. Some
% bizarre \TeX{}y command sometimes show up so we tell it to ignore
% these. All the astronomy macros that come from AAS are processed in
% \TeX{}. That way they get made to images but with an |ALT| tag that
% has the original code, which is what will be picked up by |lynx|,
% with the result that they get passed intact into the final |.txt|
% file. 
%    \begin{macrocode}
%<*adshtml-perl>
package main;

$dum = &do_cmd_makeatletter();

sub do_cmd_RMAA {
    local($_) = @_;
    "RMxAA" . $_;
}

sub do_cmd_RMAASC {
    local($_) = @_;
    "RMxAASC" . $_;
}

sub do_cmd_etal {
    local($_) = @_;
    "et al\." . $_;
}

sub do_cmd_ibidrule {
    local($_) = @_;
    "<img src=\"foo\" alt=\"\ibidrule\">" . $_;
}


sub do_cmd_bibitem {
    local($_) = @_;
    local($key);
    local($opt, $dummy)=&get_next_optional_argument; 
    $key = &missing_braces unless
        ((s/$next_pair_pr_rx/$key = $2; ''/eo)
        ||(s/$next_pair_rx/$key = $2; ''/eo));
    "<p>" . $_; 
}

&ignore_commands (<<_RAW_ARG_CMDS_);
lowercase
adjustfinalcols
textsuperscript # {}
newline # []
thinmuskip
mskip
smallskip
medskip
bigskip
hspace # {}
hspacestar # {}
vspace # {}
vspacestar # {}
mathrm
_at_mathrm
spacefactor
discretionary
unhbox 
voidb_at_x 
penalty 
_at_M
_at_m
null
IeC
_RAW_ARG_CMDS_

&process_commands_in_tex (<<_RAW_ARG_CMDS_);
sun
earth
lesssim
gtrsim
sq
argdeg
arcmin
arcsec
fd
fh
fm
fs
fdg
farcm
farcs
fp
micron
la
ga
case # {} # {}
slantfrac # {} # {}
onehalf
onethird
twothirds
onequarter
threequarters
ubvr
ub
bv
vr
ur
aj
araa
apj
apjl
apjs
ao
apss
aap
aapr
aaps
azh
baas
jrasc
memras
mnras
pra
prb
prc
prd
pre
prl
pasp
pasj
qjras
skytel
solphys
sovast
ssr
zap
nat
iaucirc
aplett
apspr
bain
fcp
gca
grl
jcp
jgr
jqsrt
memsai
nphysa
physrep
physscr
planss
procspie
_RAW_ARG_CMDS_

1;                              # This must be the last line
%</adshtml-perl> 
%    \end{macrocode}

% This is a shell script |makeads.sh| that generates the |_ads.txt| file from
% |_ads.tex| by using \LaTeX{}2HTML and |lynx|. 
% \changes{2.2beta13}{2002/03/12}{Complete rewrite of |makeads.sh|. Now much 
% improved I think}

%    \begin{macrocode}
%<*ads-shell-script>
#! /bin/sh
if [ -z $1 ]; then
    echo "Usage: $0 basename"
    exit
fi
BASENAME=$1
ADSBASENAME=${BASENAME}_ads
AUXFILE=$BASENAME.aux
TMPFILE=tmp.$$


echo "Generating $ADSBASENAME.tex"

# For some reason, must double-up \\ to \\\\ but not \ to \\ - odd!
cat > $ADSBASENAME.tex <<EOF
%% PLEASE DO NOT EDIT THIS FILE!  It is generated automatically from
%% the script makeads.sh.  Any hand-edited changes will be lost
%% 
\documentclass{article}
\usepackage[latin1]{inputenc}
\usepackage{html}
\usepackage{adshtml}
\newcommand{\ADSfield}[2]{\%#1  #2\\\\}
\newenvironment{ADSarticle}{}{\par\bigskip}
\setlength{\parindent}{0pt}
\newcounter{istage}
\renewcommand{\ion}[2]{\setcounter{istage}{#2}#1 \Roman{istage}}
\begin{document}
\makeatletter
EOF
echo >> $ADSBASENAME.tex "\input{$BASENAME.ads2}" 
echo >> $ADSBASENAME.tex "\end{document}"



echo "Generating $ADSBASENAME.sed" 

# Find type of journal
grep '\\documentclass\[.*proceedings.*\]' $BASENAME.tex >/dev/null \
     && JOURNAL='RMxAC' \
     || JOURNAL='RMxAA'

# Find the year and volume number
YEAR=`grep '\\SetYear' $BASENAME.tex | sed 's/.*\\SetYear{\(.*\)}.*/\1/'`
VOLUME=`grep '\\SetVolume' $BASENAME.tex | sed 's/.*\\SetVolume{\(.*\)}.*/\1/'`

# Generate an auxilliary sed script 
cat > $TMPFILE <<EOF
s|\\\\newlabel{\(.*\):RMAAFirstPage}{{.*}{\(.*\)}}|s/\\\\\\\\ThisVolume{\1}/$YEAR, RMxAC, $VOLUME, \2/g|g
EOF

# Now go through the .aux file, looking for \newlabels and use the sed 
# script just generated in order to make another sed script, one that 
# will replace \ThisVolume{LABEL} with the correct reference
grep '^\\newlabel{.*:RMAAFirstPage}' $AUXFILE \
     | sed -f $TMPFILE > $ADSBASENAME.sed && rm $TMPFILE

# Apend a further sed command to get rid of \@, since redefining 
# it in _ads.tex didn't work
echo >> $ADSBASENAME.sed 's/\\@//g'

echo "Generating $BASENAME.ads2 from $ADSBASENAME.sed and $BASENAME.ads" 

# Apply this to the .ads file, generating the .ads2 file, which is included
# by the _ads.tex file.
sed -f $ADSBASENAME.sed $BASENAME.ads > $BASENAME.ads2

echo "Generating $ADSBASENAME.txt from $ADSBASENAME.tex and $BASENAME.ads2"

# Now run latex2html on the _ads.tex file
latex2html -no_navigation -info 0 -address 0 -html_version 4.0,math $ADSBASENAME

# Then, convert to plaintext with lynx and tidy it up a bit
lynx --justify=off --width=100 --dump --nolist \
     $ADSBASENAME/$ADSBASENAME.html | cut -c4- \
     | sed 's/ibidrule/\\ibidrule/g' > $ADSBASENAME.txt
%</ads-shell-script>
%    \end{macrocode}

% \section{Epilogue}
% Err, that's it \dots

\endinput


% \endinput
% Local Variables:
% mode: doctex
% TeX-master: t
% End:
